\documentclass[11pt,a4paper]{article}
\usepackage[margin=25mm,headheight=15pt]{geometry}
\usepackage[T1]{fontenc}
\usepackage[utf8]{inputenc}
\usepackage{lmodern}
\usepackage{amsmath,amssymb,amsthm,amscd,mathtools,mathrsfs}
\usepackage{microtype,booktabs,longtable,array,xcolor,graphicx,enumitem,needspace,fancyhdr}
\usepackage[round,authoryear]{natbib}
\usepackage[unicode,colorlinks=true,linkcolor=blue!40!black,citecolor=blue!40!black,urlcolor=blue!40!black]{hyperref}
\usepackage{xurl}
\usepackage{listings}
\usepackage{tabularx}
\lstset{basicstyle=\ttfamily\small,breaklines=true,columns=fullflexible,keepspaces=true}
\usepackage{tikz}
\usepackage{tikz-cd}
\usetikzlibrary{arrows.meta,calc,patterns,positioning}
\usepackage{aliascnt}
\usepackage[nameinlink,noabbrev,capitalise]{cleveref}
\makeatletter
\newcommand{\SsixDisplayAliases}[2]{%
  \begingroup
  \def\label##1{%
    \@ifundefined{r@#1}{}{%
      \protected@write\@auxout{}{%
        \string\newlabel{##1}{\csname r@#1\endcsname}}%
    }%
    \@ifundefined{r@#1@cref}{}{%
      \protected@write\@auxout{}{%
        \string\newlabel{##1@cref}{\csname r@#1@cref\endcsname}}%
    }%
  }%
  #2%
  \endgroup
}
\makeatother

\allowdisplaybreaks[2]
\numberwithin{equation}{section}
\setcounter{MaxMatrixCols}{32}
\setlength{\emergencystretch}{3em}
\setlength{\parskip}{2pt}
\newtheorem{theorem}{Theorem}[section]
\newaliascnt{proposition}{theorem}
\newtheorem{proposition}[proposition]{Proposition}
\aliascntresetthe{proposition}
\crefname{proposition}{Proposition}{Propositions}
\newaliascnt{lemma}{theorem}
\newtheorem{lemma}[lemma]{Lemma}
\aliascntresetthe{lemma}
\crefname{lemma}{Lemma}{Lemmas}
\newaliascnt{corollary}{theorem}
\newtheorem{corollary}[corollary]{Corollary}
\aliascntresetthe{corollary}
\crefname{corollary}{Corollary}{Corollaries}
\newaliascnt{claim}{theorem}
\newtheorem{claim}[claim]{Claim}
\aliascntresetthe{claim}
\crefname{claim}{Claim}{Claims}
\newaliascnt{statement}{theorem}
\newtheorem{statement}[statement]{Statement}
\aliascntresetthe{statement}
\crefname{statement}{Statement}{Statements}
\theoremstyle{definition}
\newaliascnt{definition}{theorem}
\newtheorem{definition}[definition]{Definition}
\aliascntresetthe{definition}
\crefname{definition}{Definition}{Definitions}
\newaliascnt{remark}{theorem}
\newtheorem{remark}[remark]{Remark}
\aliascntresetthe{remark}
\crefname{remark}{Remark}{Remarks}
\newaliascnt{example}{theorem}
\newtheorem{example}[example]{Example}
\aliascntresetthe{example}
\crefname{example}{Example}{Examples}
\newaliascnt{problem}{theorem}
\newtheorem{problem}[problem]{Open problem}
\aliascntresetthe{problem}
\crefname{problem}{Open problem}{Open problems}
\newaliascnt{question}{theorem}
\newtheorem{question}[question]{Question}
\aliascntresetthe{question}
\crefname{question}{Question}{Questions}
\newaliascnt{auditrecord}{theorem}
\newtheorem{auditrecord}[auditrecord]{Verification record}
\aliascntresetthe{auditrecord}
\crefname{auditrecord}{Verification record}{Verification records}
\newaliascnt{test}{theorem}
\newtheorem{test}[test]{Verification}
\aliascntresetthe{test}
\crefname{test}{Verification}{Verifications}
\newaliascnt{conjecture}{theorem}
\newtheorem{conjecture}[conjecture]{Conjecture}
\aliascntresetthe{conjecture}
\crefname{conjecture}{Conjecture}{Conjectures}
\newcommand{\C}{\mathbb C}\newcommand{\R}{\mathbb R}\newcommand{\Z}{\mathbb Z}\newcommand{\Q}{\mathbb Q}\newcommand{\N}{\mathbb N}\newcommand{\F}{\mathbb F}
\newcommand{\CC}{\mathbb C}\newcommand{\RR}{\mathbb R}\newcommand{\ZZ}{\mathbb Z}\newcommand{\QQ}{\mathbb Q}\newcommand{\NN}{\mathbb N}
\newcommand{\PP}{\mathbb P}\newcommand{\Pp}{\mathbb P}\newcommand{\Pone}{\mathbb P^1}\newcommand{\HH}{\mathbb H}
\newcommand{\RP}{\mathbb{RP}}\newcommand{\CP}{\mathbb{CP}}
\newcommand{\OO}{\mathcal O}\newcommand{\cO}{\mathcal O}\newcommand{\cM}{\mathcal M}\newcommand{\cF}{\mathcal F}
\newcommand{\h}{\mathfrak h}\newcommand{\hz}{\mathfrak h_z}\newcommand{\Acon}{\mathfrak A}\newcommand{\wt}[1]{\widehat{#1}}
\newcommand{\red}{\mathrm{red}}\newcommand{\tor}{\mathrm{tor}}\newcommand{\dP}{\mathrm{dP}}
\newcommand{\gh}{\widehat\gamma}\newcommand{\uh}{\widehat u}\newcommand{\wh}{\widehat w}\newcommand{\dhat}{\widehat\delta}\renewcommand{\dh}{\widehat\delta}
\newcommand{\dd}{\mathrm d}\newcommand{\ii}{\mathrm i}\newcommand{\eps}{\varepsilon}\newcommand{\partt}{\mathrm{part}}\newcommand{\Bcirc}{B^\circ}\newcommand{\ddbar}{\partial\bar\partial}
\DeclareMathOperator{\Pic}{Pic}\DeclareMathOperator{\Alb}{Alb}\DeclareMathOperator{\Aut}{Aut}\DeclareMathOperator{\NS}{NS}
\DeclareMathOperator{\Hom}{Hom}\DeclareMathOperator{\End}{End}\DeclareMathOperator{\Ext}{Ext}\DeclareMathOperator{\Tor}{Tor}
\DeclareMathOperator{\Spec}{Spec}\DeclareMathOperator{\Proj}{Proj}\DeclareMathOperator{\Sym}{Sym}\DeclareMathOperator{\Diff}{Diff}
\DeclareMathOperator{\GL}{GL}\DeclareMathOperator{\SL}{SL}\DeclareMathOperator{\PSL}{PSL}\DeclareMathOperator{\Sp}{Sp}
\DeclareMathOperator{\im}{im}\DeclareMathOperator{\coker}{coker}\DeclareMathOperator{\rank}{rank}\DeclareMathOperator{\diag}{diag}
\DeclareMathOperator{\ord}{ord}\DeclareMathOperator{\vol}{vol}\DeclareMathOperator{\pr}{pr}\DeclareMathOperator{\id}{id}
\DeclareMathOperator{\divisor}{div}\DeclareMathOperator{\Tr}{Tr}\DeclareMathOperator{\Nm}{Nm}\DeclareMathOperator{\Sha}{Sha}
\DeclareMathOperator{\Imn}{Im}\DeclareMathOperator{\Ren}{Re}\DeclareMathOperator{\trdeg}{tr.deg}
\newaliascnt{cited}{theorem}
\newtheorem{cited}[cited]{Cited result}
\aliascntresetthe{cited}
\crefname{cited}{Cited result}{Cited results}
\newaliascnt{refuted}{theorem}
\newtheorem{refuted}[refuted]{Correction}
\aliascntresetthe{refuted}
\crefname{refuted}{Correction}{Corrections}
\newaliascnt{computed}{theorem}
\newtheorem{computed}[computed]{Computation}
\aliascntresetthe{computed}
\crefname{computed}{Computation}{Computations}
\newaliascnt{openq}{theorem}
\newtheorem{openq}[openq]{Open question}
\aliascntresetthe{openq}
\crefname{openq}{Open question}{Open questions}

\providecommand{\rk}{\operatorname{rk}}
\providecommand{\Ann}{\operatorname{Ann}}
\providecommand{\Log}{\operatorname{Log}}
\providecommand{\vir}{\mathrm{vir}}

\usepackage{etoolbox}
\hypersetup{bookmarksnumbered=true,bookmarksdepth=2}
\setcounter{tocdepth}{2}
\makeatletter\renewcommand{\@pnumwidth}{2.5em}\renewcommand{\@tocrmarg}{3.5em}\renewcommand*{\l@subsection}{\@dottedtocline{2}{1.5em}{3.5em}}\makeatother
\newtheorem{maintheorem}{Theorem}
\renewcommand{\themaintheorem}{\Alph{maintheorem}}
\crefname{maintheorem}{Theorem}{Theorems}
\newcommand{\SsixPublishRecords}{https://github.com/KokunoYumeto/s6-complex-structure-workbench/blob/main/RECORDS.md\#paper-iii}
\newcommand{\SsixRepositoryCommit}{a02f2eacf5a8c3db7cf8cdfcdba50c0002ee91c7}
\newcommand{\SeriesRefText}[1]{\ifcsname SRef@#1\endcsname\csname SRef@#1\endcsname\else\textbf{UNRESOLVED}\fi}
\newcommand{\SeriesRef}[1]{\ifmmode\text{\SeriesRefText{#1}}\else\SeriesRefText{#1}\fi}

\newcommand{\CurrentPaper}{III}
\expandafter\def\csname SRef@def:II:marked-conventions\endcsname{\href{S6_SERIES_PUBLISH_II.pdf\#nameddest=tag-P202}{II.2.1}}
\expandafter\def\csname SRef@def:II:analytic-picard-data\endcsname{\href{S6_SERIES_PUBLISH_II.pdf\#nameddest=tag-P309}{II.35.1}}
\expandafter\def\csname SRef@push:src:A.4.2\endcsname{\href{S6_SERIES_PUBLISH_I.pdf\#nameddest=tag-02ZS}{I.21.12}}
\expandafter\def\csname SRef@mapsI:src:A.1.1\endcsname{\href{S6_SERIES_PUBLISH_I.pdf\#nameddest=tag-02PA}{I.16.2}}
\expandafter\def\csname SRef@mapsI:src:A.1.2\endcsname{\href{S6_SERIES_PUBLISH_I.pdf\#nameddest=tag-02PB}{I.16.4}}
\expandafter\def\csname SRef@mapsI:src:A.1.3\endcsname{\href{S6_SERIES_PUBLISH_I.pdf\#nameddest=tag-02PC}{I.16.6}}
\expandafter\def\csname SRef@fam:src:B.1.1\endcsname{\href{S6_SERIES_PUBLISH_I.pdf\#nameddest=tag-01Q7}{I.2.11}}
\expandafter\def\csname SRef@fam:src:B.1.4\endcsname{\href{S6_SERIES_PUBLISH_I.pdf\#nameddest=tag-01QA}{I.2.13}}
\expandafter\def\csname SRef@fam:src:A.3.1\endcsname{\href{S6_SERIES_PUBLISH_I.pdf\#nameddest=tag-01Q2}{I.7.5}}
\expandafter\def\csname SRef@fam:src:B.2.1\endcsname{\href{S6_SERIES_PUBLISH_I.pdf\#nameddest=tag-01QB}{I.8.1}}
\expandafter\def\csname SRef@fam:src:B.2.2\endcsname{\href{S6_SERIES_PUBLISH_I.pdf\#nameddest=tag-01QC}{I.8.2}}
\expandafter\def\csname SRef@push:src:E.1.1\endcsname{\href{S6_SERIES_PUBLISH_I.pdf\#nameddest=tag-032R}{I.9.2}}
\expandafter\def\csname SRef@push:src:A.1.2\endcsname{\href{S6_SERIES_PUBLISH_I.pdf\#nameddest=tag-02ZI}{I.20.6}}
\expandafter\def\csname SRef@push:src:A.2.2\endcsname{\href{S6_SERIES_PUBLISH_I.pdf\#nameddest=tag-02ZK}{I.21.2}}
\expandafter\def\csname SRef@mapsI:src:A.2.1\endcsname{\href{S6_SERIES_PUBLISH_I.pdf\#nameddest=tag-02PD}{I.16.7}}
\expandafter\def\csname SRef@fam:src:C.2.3\endcsname{\href{S6_SERIES_PUBLISH_I.pdf\#nameddest=tag-01QL}{I.6.6}}
\expandafter\def\csname SRef@push:src:E.5.2\endcsname{\href{S6_SERIES_PUBLISH_II.pdf\#nameddest=tag-033O}{II.30.1}}
\expandafter\def\csname SRef@push:src:E.5.3\endcsname{\href{S6_SERIES_PUBLISH_II.pdf\#nameddest=tag-033P}{II.30.4}}
\expandafter\def\csname SRef@maps:src:D.2.12\endcsname{\hyperlink{tag-02KM}{III.\ref*{maps:src:D.2.12}}}
\expandafter\def\csname SRef@maps:src:D.2.13\endcsname{\hyperlink{tag-02KN}{III.\ref*{maps:src:D.2.13}}}
\expandafter\def\csname SRef@maps:src:D.2.14\endcsname{\hyperlink{tag-02KO}{III.\ref*{maps:src:D.2.14}}}
\expandafter\def\csname SRef@maps:src:D.2.16\endcsname{\hyperlink{tag-02KQ}{III.\ref*{maps:src:D.2.16}}}
\expandafter\def\csname SRef@maps:src:D.2.22\endcsname{\hyperlink{tag-02KX}{III.\ref*{maps:src:D.2.22}}}
\expandafter\def\csname SRef@maps:src:D.2.23\endcsname{\hyperlink{tag-02KY}{III.\ref*{maps:src:D.2.23}}}
\expandafter\def\csname SRef@night:P3F.26\endcsname{\hyperlink{tag-02S2}{III.\ref*{night:P3F.26}}}
\expandafter\def\csname SRef@night:P3F.27\endcsname{\hyperlink{tag-02S3}{III.\ref*{night:P3F.27}}}
\expandafter\def\csname SRef@night:P3F.28\endcsname{\hyperlink{tag-02S4}{III.\ref*{night:P3F.28}}}
\expandafter\def\csname SRef@night:P3F.29\endcsname{\hyperlink{tag-02S5}{III.\ref*{night:P3F.29}}}
\expandafter\def\csname SRef@night:P3F.31\endcsname{\hyperlink{tag-02S8}{III.\ref*{night:P3F.31}}}
\expandafter\def\csname SRef@iv:correction:P3K.6\endcsname{\href{S6_SERIES_PUBLISH_IV.pdf\#nameddest=tag-01US}{IV.8.2}}
\expandafter\def\csname SRef@night:P3K.7\endcsname{\href{S6_SERIES_PUBLISH_IV.pdf\#nameddest=tag-02VH}{IV.4.8}}
\expandafter\def\csname SRef@night:P3K.65\endcsname{\href{S6_SERIES_PUBLISH_IV.pdf\#nameddest=tag-02VC}{IV.5.13}}
\expandafter\def\csname SRef@night:P3K.68\endcsname{\href{S6_SERIES_PUBLISH_IV.pdf\#nameddest=tag-02VF}{IV.6.1}}

\begin{document}
\title{The $(3,4,\infty)$ family of complex six-spheres\\[.5em]III. Maps and the resulting manifolds}
\author{AI-produced mathematical text}
\date{25 November 2026}
\maketitle

\begin{abstract}
We compute the homotopy classes of three maps associated with the
$(3,4,\infty)$ family of complex six-spheres. The torus quotient
represents the element of order two in $\pi_6(S^3)$, the holomorphic
map to $\mathbb P^1$ is null-homotopic, and the cusp map that keeps
the angular coordinate represents the nonzero element of $\pi_6(S^4)$.
Normal framings and integral attaching maps give these classes and a
finite cell model. For every cyclic universal cover we calculate the
integral cohomology ring and the complete smooth classification
invariant. We classify the natural section-line sphere fillings over
twenty-one elliptic profiles: there are $59$ marked patterns on $50$
initial data, with twelve profiles admitting unbounded multiplicity.
We also determine weighted smooth types and the homology of every free
affine cap over the elliptic surface with fibres $IV^*,II,I_2$,
including its additional three-primary torsion. A weighted member over
the surface with fibres $I_8,4I_1$ has zero refined Pontryagin class, as follows
from an embedded Hopf sphere and its critical-cycle linking number.
The local Kummer cover has no smooth compact global extension at
degree greater than one; an explicit finite flat normal extension
has additional components above the cusp. Finally, moving the zero
of the line-bundle linearization gives a conifold smoothing
diffeomorphic to $S^3\times S^3$, of algebraic dimension one.
\end{abstract}

\clearpage
\tableofcontents
\clearpage
\section{Results and their proofs in outline}\label{sec:results}
\subsection{The series and this paper}
The five papers study the periods, complex structures, maps and
degenerations of the $(3,4,\infty)$ torus-fibred family and the
manifolds obtained from its elliptic-surface construction. This paper
determines sphere-map classes, integral attachment invariants and
smooth types. The constructions include an ordered period lattice,
finite affine twists, positive normal circles and a cusp fan. The
table lists the principal results and their cross-paper inputs;
for Paper~III, the indicated uses are the quotient construction,
the section-line family and the return-map identification.
\begin{center}\small
\begin{tabularx}{\textwidth}{@{}lXX@{}}\toprule
Paper & Main result & Cross-paper inputs\\\midrule
I & Periods, moduli, fibre strata and the two ends of the family & ---\\
II & Cotangent obstruction and the sign theorem for degenerations & I, IV\\
III & Sphere-map classes, integral covers and elliptic smooth types &
I: \SeriesRef{mapsI:src:A.2.1}; II: \SeriesRef{push:src:A.1.2},
\SeriesRef{push:src:E.5.2}; IV: \SeriesRef{iv:correction:P3K.6},
\SeriesRef{night:P3K.7}\\
IV & Smooth return maps in $\Theta_7$ and cusp-fan calculations & I, II, III, V\\
V & Resolved fibre products and the marked rank-six fourfold family & I, III, IV\\
\bottomrule\end{tabularx}\end{center}

\subsection{The main theorems}
\hypertarget{tag-03KA}{}\label{series:III:main-theorems}
The marked family and its admissible twists are those of
Definition~\ref{mapsIII:def:marked-family}; its parameter range is
Definition~\ref{mapsIII:def:period-parameters}. We use the complex
orientation, positive normal circles and the outward stable framing.
The maps in Theorems~\ref{series:III:A}--\ref{series:III:C} have the
domains and based sphere identifications stated there. Each elliptic
construction in Theorem~\ref{series:III:F} includes its section line,
linearization zero, fan and affine cap data.

\hypertarget{main-III-A}{}
\begin{maintheorem}[Relative realization of the sphere quotient]
\hypertarget{tag-03KB}{}\label{series:III:A}
{\normalfont\scriptsize tag 03KB}\par
Let $X_p$ be an original sphere member, with its fixed oriented
identification with $S^6$, let $N$ be an open neighbourhood of
its cusp, and fix the quotient target $Q$
and a homeomorphism $\varphi:Q\to S^3$ as in
\SeriesRef{mapsI:src:A.1.3}. The cusp replacement of
\SeriesRef{mapsI:src:A.2.1}, with its prescribed outer collar and
image in the designated solid torus, defines the quotient homotopy
class used in this paper.

Every element of $[S^6,S^3]=\mathbb Z/12$ is represented by
$\varphi q$, where $q:X_p\to Q$ is smooth and agrees with the
torus quotient on $X_p\setminus N$. For the designated cusp
solid torus $T\subset Q$ and every $y\in T$,
\[
 \operatorname{im}\bigl(\pi_6(T,y)\longrightarrow
          \pi_6(S^3,\varphi(y))\bigr)=\{0\}.
\]
These relative assertions are proved in \cref{maps:src:A.2.2}.
\end{maintheorem}
\textit{Mechanism.} A map supported in an interior six-ball changes
the homotopy class by any element of $\pi_6(S^3)$, while the chosen
solid-torus target has zero sixth homotopy group.
\textit{Literature status.} New relative realization for the quotient
of \cite[Remark~8.3, pp.~22--23]{E26}.

\hypertarget{main-III-B}{}
\begin{maintheorem}[The quotient and algebraic-reduction classes]
\hypertarget{tag-03KC}{}\label{series:III:B}
{\normalfont\scriptsize tag 03KC}\par
For a sphere member of the marked $(3,4,\infty)$ family, let
$e:X_p\to S^3$ be the quotient replacement with the target and
collar data of \SeriesRef{mapsI:src:A.2.1}. Then
\[
 [e]=6\nu'=\eta_3\eta_4\eta_5\ne0
       \quad\text{in }\pi_6(S^3)=\mathbb Z/12.
\]
Its stabilization is $12\nu\in\pi_3^{\mathrm S}=\mathbb Z/24$,
and its real Adams invariant is $1/2\in\mathbb Q/\mathbb Z$.
A regular inverse image is a three-torus with its Lie stable
framing, by \cref{maps:src:B.1.1,maps:src:B.1.2}.

The holomorphic map $f:X_p\to\mathbb P^1$ is null-homotopic.
Compatible frames of the two divisor line bundles give functions
$a,b$ with no common zero and
\[
 \begin{gathered}
 M f=[a^3:b^4]\simeq\eta_2 P_{3,4}G,\\
 G=\frac{(a,b)}{\sqrt{|a|^2+|b|^2}},
 \qquad M[x:y]=[x:y-x],
 \end{gathered}
\]
where $P_{3,4}(z_1,z_2)=(z_1^3,z_2^4)/\|(z_1^3,z_2^4)\|$
has degree twelve. The frame changes, homotopy and inverse
$M^{-1}[a:b]=[a:a+b]$ are given in \cref{maps:src:B.2.1}.
\end{maintheorem}
\textit{Mechanism.} The quotient framing gives the Lie-framed
three-torus, and the divisor-section homotopy factors the base map
through a degree-twelve map that kills $\pi_6(S^3)$.
\textit{Literature status.} New computations for these maps, using
\cite[Theorem~4.1 and Corollaries~4.2--4.3, pp.~76--77]{mapsToda52}
and \cite[Theorem~4.14, p.~423]{mapsAPS2}.

\hypertarget{main-III-C}{}
\begin{maintheorem}[The cusp map and an integral cell model]
\hypertarget{tag-03KD}{}\label{series:III:C}
{\normalfont\scriptsize tag 03KD}\par
For every sphere member in
Definition~\ref{mapsIII:def:retained-angle-cusp-data}, let
$H:X_p\to S^4$ be the based cusp map of
Definition~\ref{mapsIII:def:cusp-suspension-maps}. It keeps the
base angle and collapses the coordinate two-skeleton of the
three-torus before reduced suspension. Then
\[
 [H]=\eta_4\eta_5\ne0\in\pi_6(S^4)=\mathbb Z/2.
\]
Its regular inverse image is the vanishing two-torus with Lie stable
framing and Arf invariant one, as proved in \cref{maps:src:B.3.1}.

For the $p=-1$ three-filling datum of
Definition~\ref{mapsIII:def:three-filling-datum}, with zero cusp
translation drift, the mapping cylinders of the regular torus bundle,
the finite affine caps and the triangular cusp have a finite CW
homotopy model with cell counts
\[
 (c_0,c_1,c_2,c_3,c_4,c_5,c_6)=(5,20,45,56,42,17,3).
\]
Its integral attaching maps give $\pi_1=1$ and the integral homology
of $S^6$, hence the homotopy type of $S^6$
(\cref{maps:src:B.5.1,maps:src:B.5.2}).
\end{maintheorem}
\textit{Mechanism.} A periodic contraction of the normal-frame
comparison identifies the Lie torus, and the marked finite and cusp
attachments determine the integral cell complex.
\textit{Literature status.} New framing and cell calculations for
the sphere construction of \cite[Theorem~4.1, p.~12]{E26}.

\hypertarget{main-III-D}{}
\begin{maintheorem}[Local Kummer maps and global extensions]
\hypertarget{tag-03KE}{}\label{series:III:D}
{\normalfont\scriptsize tag 03KE}\par
For the marked cusp of Definition~\ref{mapsIII:def:kummer-data}
and every integer $N\ge1$, the coordinate power map gives
\[
 \kappa_N:N_0^{\langle N\rangle}\longrightarrow N_0,
 \qquad (\xi_1,\xi_2,v)\longmapsto(\xi_1^N,\xi_2^N,v^N),
\]
a finite flat map of degree $N^3$ between smooth threefolds.
The source uses the correction $C_N(v)=C(v^N)/N$ and the same
triangular fan. Its factor over the normal base change $q=v^N$ is
\[
 Q_N:N_0^{\langle N\rangle}\longrightarrow
 N_0\mathbin{\times}_{\Delta_q}\Delta_v,
 \qquad \deg Q_N=N^2.
\]
For $N>1$, this factor is quasi-\'etale and its nonflat and
non-\'etale loci are the inverse images of the double and triple
strata; for $N=1$ both maps are identities.

For $N>1$, $\kappa_N$ is never the restriction, even to one cusp
component, of a finite map to $X$ from a connected smooth compact
complex source. It also has no normal compact extension whose
entire cusp preimage is this single prescribed map. There is,
however, a connected finite flat cover $Y_N\to X$ with normal
compact source, branched only over $S_1\cup S_2\cup W$, of degree
\[
 N^2|G_N|,\qquad
 G_N=\langle(A_1,-1),(A_2,0)\rangle
       \subset\operatorname{GL}_4(\mathbb Z/N)\times\mathbb Z/N.
\]
One component over the cusp is exactly $\kappa_N$; for $N>1$
another component contains a singular point. These statements,
including the identity extension for $N=1$, are
\cref{maps:src:C.1.1,maps:src:C.1.2,maps:src:C.1.3,maps:src:C.1.4,maps:src:C.1.5}.
\end{maintheorem}
\textit{Mechanism.} Smoothness on adjacent fan cones forces a scalar
toric lattice, while monodromy changes the Kummer lattice and a
finite semidirect quotient supplies the normal global cover.
\textit{Literature status.} New obstruction and completion for this
cusp, using the covering classification in
\cite[Theorem~1.38, pp.~67--68]{Hatcher2002}.

\hypertarget{main-III-E}{}
\begin{maintheorem}[The cyclic universal covers]
\hypertarget{tag-03KF}{}\label{series:III:E}
{\normalfont\scriptsize tag 03KF}\par
For every admissible member $X_p$ of the marked family, put $d=|p|$;
then $d\ge1$ and $\gcd(d,6)=1$. Its fundamental group is
$\mathbb Z/d$, and its simply connected universal cover has
\[
 H_2(\widetilde X_p)=\mathbb Z^{d-1}\oplus\mathbb Z/d,
 \qquad H_3(\widetilde X_p)=\mathbb Z/d,
 \qquad \chi(\widetilde X_p)=2d.
\]
It is spin, $c_2=p_1=0$ integrally, and its refined Pontryagin
class is zero. For $d>1$, the $d$ cusp-divisor classes $x_i$ have
sole relation $\sum_i x_i=0$ and satisfy
\[
 x_i x_j=C_{ij}\kappa,\qquad
 C_{ij}=\mathbf1_{j=i+1}+\mathbf1_{j=i-1}-2\mathbf1_{j=i},
\]
with indices modulo $d$ and $\kappa$ a generator of the order-$d$
subgroup of $H^4$. There is a marked cyclic coordinate
$u\in H^2(\widetilde X_p;\mathbb Z/d)$ for which the integral
cubic is zero and
\[
 \langle ux_ix_j,[\widetilde X_p]\rangle=C_{ij},\qquad
 \langle u^2x_i,[\widetilde X_p]\rangle=-12,\qquad
 \langle u^3,[\widetilde X_p]\rangle=0\pmod d.
\]
These give the complete fundamental-class invariant, including
all zero coordinates, in
Definition~\ref{mapsIII:def:cyclic-classifying-data}.

The canonical manifold $M_d$ of
Definition~\ref{mapsIII:def:canonical-cyclic-model} realizes these
data and is the unique oriented smooth type with them among closed
simply connected spin six-manifolds. Every admissible cover satisfies
$\widetilde X_p\cong M_d$, and $M_1=S^6$
(\cref{maps:src:D.1.4,maps:src:D.1.6,maps:src:D.1.8,maps:src:D.1.9,maps:src:D.1.10,maps:src:D.1.11}).
\end{maintheorem}
\textit{Mechanism.} Integral cusp and cap maps determine the groups,
and the divisor products and finite-coefficient cubic determine the
complete marked spin-bordism invariant.
\textit{Literature status.} New cover invariants for the family of
\cite[Theorems~7.17 and~7.22, pp.~57, 60]{AC26}, recognized by
\cite[Theorems~3 and~5, pp.~795--796; \S4.1, p.~809]{manZhubr75}.

\hypertarget{main-III-F}{}
\begin{maintheorem}[Elliptic sphere fillings and weighted smooth types]
\hypertarget{tag-03KG}{}\label{series:III:F}
{\normalfont\scriptsize tag 03KG}\par
\begin{enumerate}[label=(\alph*),leftmargin=*]
\item For the natural line $M=\mathcal O(P-O)$, allow all section
pairs, component-supported linearization zeros, original and Mumford
fillings, additive and regular logarithms, multiplicities and affine
twists in the sector of \cref{maps:src:D.2.22}. With the marked
patterns of Definition~\ref{maps:def:marked-filling-pattern}, there
are exactly $59$ sphere patterns on $50$ initial rows over the
twenty-one profiles in \cref{maps:src:D.2.28}. Twelve profiles admit
unbounded sphere multiplicities, four require every logarithmic
order to be one, and five admit none. Every threefold satisfying these sphere conditions is
orientation-preservingly diffeomorphic to $S^6$; the complete affine
and multiplicity conditions are \eqref{maps:eq:D.HY3}--\eqref{maps:eq:D.HY5}.

\item For every positive weight $d$, the Engel data of
Definition~\ref{maps:def:weighted-engel-data}, with
$m_0=3k$, $m_1=4\ell$, $k,\ell>0$ and
$m_1n_0+m_0n_1=\pm1$, give one oriented smooth type $M_d^{\rm w}$ for
all the allowed affine twists. It has
\[
 H_2=\mathbb Z^{d-1}\oplus\mathbb Z/d,\quad
 H_3=\mathbb Z/d,\quad c_1=c_2=\lambda=P_{\rm ref}=0.
\]
Its full fundamental-class invariant is
\eqref{maps:eq:D.WI15}, including the Pontryagin square modulo
$d\gcd(d,2)$ and third power modulo $d\gcd(d,3)$, with values
$-12$ and $0$. For $d$ prime to six,
$M_d^{\rm w}\cong M_d$ with the complex orientations and
the marking proved in \cref{maps:src:D.2.48}, where $M_d$ is
the cyclic-cover model of Theorem~\ref{series:III:E}.

\item For the OS45 data of Definition~\ref{maps:def:weighted-os45-data},
where OS45 denotes entry~45 of
\cite[Main Theorem table, p.~85]{manOguisoShioda91}
(root lattice $A_7$, Mordell--Weil lattice $\langle1/8\rangle$),
namely $P=2G$, $Q=8G$, weights $(2,0,0,0,0)$ at $I_8,4I_1$,
the triangular Mumford caps and ordinary clutch $e=1$, the
manifold is simply connected and satisfies
\[
 H_2=\mathbb Z^{15}\oplus\mathbb Z/8,
 \qquad H_3=\mathbb Z/8,
 \qquad c_1=c_2=\lambda=P_{\rm ref}=0.
\]
Its primitive integral cubic and complete fundamental-class
invariant are given in \cref{maps:src:D.2.53}.

\item For the affine cap of entry~49 of the same table
(OS49, with root lattice $E_6\oplus A_1$ and Mordell--Weil
lattice $\langle1/6\rangle$), let $m=3k$, $k>0$ and
$\ell=n(e_\gamma-e_u)+d\delta$, $n,d\in\mathbb Z$. Freeness
is equivalent to
\[
 \gcd(k,n,d)=1\quad\text{and}\quad(3\nmid n\ \text{or}\ 3\mid k).
\]
Its fundamental group is $\mathbb Z/|n|$, with
$\mathbb Z/0=\mathbb Z$. For $n\ne0$, putting
$g=\gcd(3k,n,3d)\in\{1,3\}$ gives
\[
 H_1=H_4=\mathbb Z/|n|,\qquad
 H_2=H_3=\mathbb Z/|n|\oplus\mathbb Z/g,\qquad H_5=0.
\]
For $n=0$, freeness is $3\mid k$ and $\gcd(k,d)=1$, and
\[
 H_1=H_5=\mathbb Z,\quad
 H_2=H_3=\mathbb Z^2\oplus\mathbb Z/3,\quad H_4=\mathbb Z^2.
\]
In both cases $H_0=H_6=\mathbb Z$. The sphere sector is exactly
$n=\pm1$. In the coprime sector $\gcd(n,3k)=1$ one has $g=1$;
the free triple $(m,n,d)=(9,3,1)$ has $g=3$ and
$H_2=H_3=(\mathbb Z/3)^2$
(\cref{night:P3F.33,night:P3F.34}).
\end{enumerate}
\end{maintheorem}
\textit{Mechanism.} Primitive specialization and oriented normal-torus
maps give the integral attachment groups, while cusp intersections
and framed critical-cycle linkings determine their smooth invariants.
\textit{Literature status.} New classifications for these explicit
data within the section-line construction of
\cite[Remark~8.5, p.~24]{E26}.

\hypertarget{main-III-G}{}
\begin{maintheorem}[A conifold smoothing on $S^3\times S^3$]
\hypertarget{tag-03KH}{}\label{series:III:G}
{\normalfont\scriptsize tag 03KH}\par
Let $S\to\mathbb P^1$ have fibres $IV^*,III,I_1$ at $0,1,\infty$,
let $L=\mathcal O_S(P-O)$, and use $T=t_{6P}$ and
$L\otimes T^*L^{-1}\simeq\pi^*\mathcal O(1)$ with the direction
and sphere twists of \cref{maps:src:D.3.1}. Choose sections
$\phi_0,\phi_1$ vanishing at $\infty,0$, with
$\phi_0=q\phi_1$ near infinity, and put
$\phi_s=\pi^*(\phi_0-s\phi_1)$.

For every sufficiently small fixed $u\ne0$ there is a proper flat
family $\mathcal Y_u\to\Delta_s$, also proper and flat over
$\mathbb P^1\times\Delta_s$. Its central fibre contracts one
$(-1,-1)$-curve of the original sphere to an ordinary double point;
its local equation is
\[
 xy=u(zv-s),\qquad q=zv.
\]
For $0<|s|\ll1$ its complex-oriented smooth fibre is
orientation-preservingly diffeomorphic to $S^3\times S^3$ and has
algebraic dimension one. The surgery replaces $S^2\times D^4$ by
$D^3\times S^3$, with primitive boundary maps. Its middle form is
the standard integral skew form and $w_2=p_1=0$
(\cref{maps:src:D.3.1,maps:src:D.3.2,maps:src:D.3.3}).
\end{maintheorem}
\textit{Mechanism.} The moving section gives the node equation,
and the oriented boundary maps of its smoothing determine the
middle homology and smooth surgery type.
\textit{Literature status.} New moving-zero family from
\cite[Proposition~2.4 and Propositions~2.7--2.8, pp.~5--7]{E26},
with smooth recognition by \cite[Theorem~2.6, p.~8]{BHS19}.

\paragraph{Further results.}
\begin{itemize}[leftmargin=*]
\item \Cref{maps:src:B.3.2,maps:src:B.3.3}: angle-forgetting maps are
null-homotopic, and the cusp-core map has class $(6\nu',0,0)$ in its
specified $S^3\vee S^2$ cofiber coordinates.
\item \Cref{mapsaudit:finite-degree-three}: the two finite pullback
lattices in degree three have indices one and two, with all transfer maps.
\item \Cref{maps:src:D.1.7,maps:src:D.1.12}: the cyclic models have
the parameter disc $0<|W|<216$ and explicit cyclic quotient actions.
\item \Cref{maps:src:D.2.22}: height and component tests give
$172$ natural section-line initial data and $226$ with vertical corrections.
\item \Cref{maps:src:D.2.18}: the three-component cusp gives the
flag-manifold summand in the oriented type of its sphere-group sector.
\item \Cref{maps:src:D.2.20,maps:src:D.2.21}: integral groups and
oriented smooth types for six primitive original families and mixed fillings.
\item \Cref{maps:src:D.2.29,maps:src:D.2.30}: finite additive caps
with arbitrary contracting-character index and affine twists.
\item \Cref{maps:src:D.2.31,maps:src:D.2.32,maps:src:D.2.33}:
integral cap maps for the full natural catalogue with original fillings.
\item \Cref{maps:src:D.2.35,maps:src:D.2.36,maps:src:D.2.38,maps:src:D.2.42,maps:src:D.2.44,maps:src:D.2.45}: smooth types of
the enumerated minimal and residual families from their full integral invariants.
\item \Cref{maps:src:D.2.46,maps:src:D.2.47}: whole-fibre
attachment groups and the trilinear divisor formula for finite-index cusp lattices.
\item \Cref{maps:src:D.2.51,maps:src:D.2.52}: invariant-Chern
original caps, their marked kernels and height-integral section pairs.
\item \Cref{night:P3F.31,night:P3F.32}: zero smooth return for
the OS49 sphere caps, with product mapping tori.
\item \Cref{wb:cusp-cover-extension-theorem,wb:cusp-cover-naive-action-failure}:
normalization, conductor and deck action for arbitrary holomorphic cusp corrections.
\item \Cref{wb:cusp-spec-globalization-obstruction}: global torus
quotients, their intrinsic distances and their extension obstructions.
\item \Cref{wb:cusp-cone-target-homotopy-types}: the compact cusp
carriers have explicit wedge and smash decompositions with marked homology generators.
\item \Cref{wb:cusp-null:thm,wb:cusp-s12-one-copy-invariant}:
a framed cusp nullbordism and the coefficient calculation for the sphere carriers.
\item \Cref{wb:sphere-metric-scaling-trichotomy,wb:cusp-cone-sharp-global-R4-scope}: radial scaling limits and the
correspondence between Euclidean cells, compact targets and boundary fibres.
\item \Cref{maps:src:E.2.2}: transverse pullbacks realize
$S^6\times N$ with the stated complex and topological data.
\end{itemize}

\subsection{The proofs in outline}

\hypertarget{outline-III-A}{}
\subsubsection*{Theorem A: relative sphere quotients}
\begin{enumerate}[label=\arabic*.,leftmargin=*]
\item\hypertarget{step-III-A-1}{} The primitive invariant character
gives the fibrewise circle quotient, and the integral period maps fix
its signed meridian (\SeriesRef{mapsI:src:A.1.1}).
\item\hypertarget{step-III-A-2}{} The cusp replacement has its image
in the designated solid torus and is fixed on the outer collar;
its relative homotopy assertion is \SeriesRef{mapsI:src:A.2.1}.
\item\hypertarget{step-III-A-3}{} A based collapse $c:S^6\to S^6$,
homotopic to the identity, is constant off an interior six-ball.
Quaternionic multiplication gives
$e_g(x)=e(x)e(c(x))^{-1}g(c(x))$, which equals $e$ on the exterior
and is homotopic to any chosen $g$ (\cref{maps:src:A.2.2}).
\item\hypertarget{step-III-A-4}{} Relative smooth approximation
keeps the exterior fixed. A based map into $S^1\times D^2$
contracts by lifting its circle coordinate and contracting its disc
coordinate, proving the zero image in degree six
(\cref{maps:src:A.2.2}).
\end{enumerate}

\hypertarget{outline-III-B}{}
\subsubsection*{Theorem B: the two map classes}
\begin{enumerate}[label=\arabic*.,leftmargin=*]
\item\hypertarget{step-III-B-1}{} The common invariant mod-two
covectors are generated by $\gamma$. The framing difference therefore
lifts to $\operatorname{Spin}(7)$ on every coset of $\ker\gamma$
(\cref{maps:src:B.1.1}).
\item\hypertarget{step-III-B-2}{} Integral duality identifies the
common degree-three invariant with $\gamma u w$, whose restriction
to that kernel is zero. Transfer and torsion-freeness of $H^3(T^3)$
kill the remaining lifted framing obstruction
(\cref{maps:src:B.1.1}).
\item\hypertarget{step-III-B-3}{} The framed regular fibre is the
Lie torus, with stable class $\eta^3=12\nu$; its periodic Dirac
operator has paired nonzero eigenvalues and a two-dimensional kernel,
giving real Adams invariant $1/2$ (\cref{maps:src:B.1.2}).
\item\hypertarget{step-III-B-4}{} Stabilization sends
$\nu'\mapsto2\nu$ and is injective on $\mathbb Z/12$.
It therefore identifies the quotient class with $6\nu'$ and
$\eta_3\eta_4\eta_5$ (\cref{maps:src:B.1.2}).
\item\hypertarget{step-III-B-5}{} The divisor multiplicities
$3S_1,4S_2$ and $H^1(X)=H^2(X)=0$ give compatible line frames:
$e'_j=e^{-b_j/m_j}e_j$, with inverse $e_j=e^{b_j/m_j}e'_j$.
Their sections satisfy $M f=[a^3:b^4]$ (\cref{maps:src:B.2.1}).
\item\hypertarget{step-III-B-6}{} With
$r=(|a|^2+|b|^2)^{1/2}$, the homotopy
$[a^3e^{-3s\log r}:b^4e^{-4s\log r}]$ joins $M f$ to
$\eta_2P_{3,4}G$. The radial coordinate of $P_{3,4}$ is strictly
increasing and its phase degrees are three and four, so its degree
is twelve (\cref{maps:src:B.2.1}).
\item\hypertarget{step-III-B-7}{} Degree classification identifies
$P_{3,4}$ with the quaternionic power map $q\mapsto q^{12}$,
whose induced action is multiplication by twelve on $\pi_6(S^3)$.
Composing the resulting null-homotopy with $\eta_2$ and then
$M^{-1}$ proves $[f]=0$ (\cref{maps:src:B.2.1}).
\end{enumerate}

\hypertarget{outline-III-C}{}
\subsubsection*{Theorem C: the cusp frame and cell attachments}
\begin{enumerate}[label=\arabic*.,leftmargin=*]
\item\hypertarget{step-III-C-1}{} The period-dependent unit change
and its inverse place the full regular inverse-image torus in one
triple-point chart; differentiation includes the bounded period term
(\cref{wb:h24:ret:chart-lemma,mapsIII:lem:retained-angle-differential}).
\item\hypertarget{step-III-C-2}{} The ordered normal and tangent
frame is $D(a,b)G_0$, with
$D=\operatorname{diag}(e^{-i(a+b)},e^{ia},e^{ib})$.
Two periodic $\operatorname{SU}(2)$ contractions identify its stable
framing with the Lie framing (\cref{wb:h24:ret:framing-proposition}).
\item\hypertarget{step-III-C-3}{} The Lie torus has Arf invariant
one. Based smoothing protects the full regular fibre, and
Freudenthal injectivity gives $[H]=\eta_4\eta_5$
(\cref{mapsIII:lem:lie-torus-spin-invariant,maps:src:B.3.1}).
\item\hypertarget{step-III-C-4}{} The torus-extension homotopy
contracts both angle-forgetting maps. The cusp-core collapse is
homotopic to the quotient followed by its core cofiber map, giving
$(6\nu',0,0)$ in the specified wedge coordinates
(\cref{maps:src:B.3.2,maps:src:B.3.3}).
\item\hypertarget{step-III-C-5}{} For the $p=-1$ datum, the finite
map $[x,\theta]\mapsto[x+\theta v_j/m_j]$ includes the logarithmic
drift, and the cusp seam is
$[a,b,\theta]\mapsto[a,b-\theta(B_0a-v(a))]$.
The coordinate identifications give the three mapping cylinders
(\cref{maps:src:B.5.1}).
\item\hypertarget{step-III-C-6}{} The marked boundary tori meet
the product-collar criterion of
\cite[Theorem~6.1, p.~77]{mapsWaldhausen68}. Cellular approximation
of the torus return and attaching maps gives the displayed cell
counts, whose alternating sum is two (\cref{maps:src:B.5.1}).
\item\hypertarget{step-III-C-7}{} The integral finite and cusp maps
kill the fundamental group and the degree-two attachment cokernel.
Duality and the Euler number give sphere homology; Hurewicz and
Whitehead give the homotopy sphere (\cref{maps:src:B.5.2}).
\end{enumerate}

\hypertarget{outline-III-D}{}
\subsubsection*{Theorem D: Kummer lattices and normal completion}
\begin{enumerate}[label=\arabic*.,leftmargin=*]
\item\hypertarget{step-III-D-1}{} On $q=v^N$ the source correction
$C(v^N)/N$ makes the power map equivariant. The chart basis gives
absolute degree $N^3$; the invariant-ring factor over the base change
has degree $N^2$ and image lattice
$\langle\widehat\gamma,\widehat u,N\widehat w,N\widehat\delta\rangle$
(\cref{maps:src:C.1.1}).
\item\hypertarget{step-III-D-2}{} Adjacent unimodular triangles
give reciprocal divisibility conditions on the primitive ray indices.
Their equality forces the cocharacter lattice of a smooth cover of
the entire cusp fan to be $n\mathbb Z^3$
(\cref{maps:src:C.1.2}).
\item\hypertarget{step-III-D-3}{} The order-four meridian carries
$L_N$ to a lattice with toric intersection
$\mathbb Z\widehat w\oplus N\mathbb Z\widehat\delta$.
The primitive first generator forces a scalar index one, while the
second has index $N$; this excludes a smooth compact extension for
$N>1$ (\cref{maps:src:C.1.3}).
\item\hypertarget{step-III-D-4}{} The lattices
$D^aM_0^bL_N$, $0\le a,b<N$, for
$D=A_2^2A_1A_2^2A_1^{-1}$, are distinct. A normal extension with
only the prescribed cusp component would have Stein base degree
$N$, contradicting its orbit of at least $N^2$ sheets
(\cref{maps:src:C.1.4}).
\item\hypertarget{step-III-D-5}{} The subgroup of
$(\mathbb Z/N)^4\rtimes G_N$ generated by the two long periods
has index $N^2|G_N|$ and the required cusp intersection.
Fractional monomials in each half-open fundamental box form a free
basis of its normal completion; uniqueness of normalization glues
the charts and their inverse overlaps (\cref{maps:src:C.1.5}).
\end{enumerate}

\hypertarget{outline-III-E}{}
\subsubsection*{Theorem E: integral invariants of the cyclic covers}
\begin{enumerate}[label=\arabic*.,leftmargin=*]
\item\hypertarget{step-III-E-1}{} The cyclic character identifies
the cusp cover with the period sublattice
$\mathbb Z\widehat u\oplus d\mathbb Z\widehat\gamma$.
Its normalization has $d$ components and its Laurent relation
matrix reduces to $\operatorname{diag}(1,t-1,t-1,0)$
(\cref{maps:src:D.1.2}).
\item\hypertarget{step-III-E-2}{} The nontrivial deck characters
give acyclic regular-fibre complexes. The integral finite pullbacks
and cusp sweeps determine the free groups and the order-$d$
torsion in degrees two and three
(\cref{maps:src:D.1.3,maps:src:D.1.4}).
\item\hypertarget{step-III-E-3}{} Positive normal Thom classes
give the divisor relation $\sum x_i=0$; restrictions to each
normalized component give $x_ix_j=C_{ij}\kappa$
(\cref{maps:src:D.1.6}).
\item\hypertarget{step-III-E-4}{} A class evaluating to one on
the fixed $\delta$-curve pulls back to the cyclic coordinate $u$.
The component restrictions give the mixed cubic values
$C_{ij},-12,0$, and the full $H_6(K(H_2,2))$ decomposition
identifies every coordinate (\cref{maps:src:D.1.9,maps:src:D.1.10}).
\item\hypertarget{step-III-E-5}{} Degree-four pullback kills
$c_2$ and $p_1$; multiplication by four is injective on
$\mathbb Z^{d-1}\oplus\mathbb Z/d$, so the refined class is zero.
The marked spin-bordism classification with $b_3=0$ identifies
the cover with the constructed canonical $M_d$
(\cref{maps:src:D.1.8,maps:src:D.1.11}).
\item\hypertarget{step-III-E-6}{} Euler multiplicativity gives
$\chi=2d$. A connected sum of $g$ copies of $S^3\times S^3$
has Euler number $2-2g$; equality for $d>1$ would require the
negative value $g=1-d$. Thus these covers are excluded from
those connected-sum types (\cref{maps:src:D.1.1}).
\end{enumerate}

\hypertarget{outline-III-F}{}
\subsubsection*{Theorem F: specialization and characteristic invariants}
\begin{enumerate}[label=\arabic*.,leftmargin=*]
\item\hypertarget{step-III-F-1}{} The height equality, component
equations and narrowness at every other reducible site give the
finite natural section-line initial set
(\cref{maps:src:D.2.22}).
\item\hypertarget{step-III-F-2}{} Primitive specialization and
the positive normal-torus maps identify the full integral cap
cokernels. For simply connected natural threefolds their order is
$cR^2$, so the sphere condition is $c=R=1$ with the stated
fundamental-group criterion (\cref{maps:src:D.2.26}).
\item\hypertarget{step-III-F-3}{} The $I_3^*$ and $I_1^*$
obstructions eliminate five profiles. The remaining integral
presentations give $59$ patterns on $50$ initial rows and the
twelve/four/five multiplicity alternatives
(\cref{maps:src:D.2.24,maps:src:D.2.27,maps:src:D.2.28}).
\item\hypertarget{step-III-F-4}{} In the minimal and residual
families, the primitive cubic, refined class and full
fundamental-class image distinguish the stated surgeries and
smooth representatives (\cref{maps:src:D.2.35,maps:src:D.2.36,maps:src:D.2.38,maps:src:D.2.42,maps:src:D.2.44,maps:src:D.2.45}).
\item\hypertarget{step-III-F-5}{} Component restrictions compute
the weighted cusp cubic, including loops and repeated incidences.
For the OS45 datum, the source-line Hopf sphere has differential
reference parity $\varepsilon_\rho=1$ and complete critical-cycle
linking number $-1$; the framed formula gives $P_{\rm ref}=0$
(\cref{maps:src:D.2.47,maps:src:D.2.53}).
\item\hypertarget{step-III-F-6}{} For the simply connected
weight-$d$ Engel data, the determinant $P=\pm1$ gives the primitive
groups. The refined class is zero and the complete coefficient
invariant contains the cyclic Cartan matrix, Pontryagin-square
value $-12$ and third-power value zero, proving the single type
$M_d^{\rm w}$ for every $d>0$ (\cref{maps:src:D.2.48}).
\item\hypertarget{step-III-F-7}{} In the full free OS49 sector,
the invariant quotient coordinates $(\psi,-w+3\delta)$ give the
exact freeness criterion. The oriented normal circle and primitive
Gysin lattice have Smith factors $g,|n|$ for $n\ne0$; the
$n=0$ pair sequence gives the displayed free and order-three
summands (\cref{night:P3F.33,night:P3F.34}).
\end{enumerate}

\hypertarget{outline-III-G}{}
\subsubsection*{Theorem G: the conifold smoothing}
\begin{enumerate}[label=\arabic*.,leftmargin=*]
\item\hypertarget{step-III-G-1}{} The section
$\phi_s=\phi_0-s\phi_1$ gives annular core and neck charts whose
overlap maps have reciprocal inverses and satisfy the line-bundle
cocycle (\cref{maps:src:D.3.1}).
\item\hypertarget{step-III-G-2}{} Their local equation is
$xy=u(zv-s)$. Uniform collars include the order-three and
order-four root caps in a proper flat family; its central node
has the stated small resolution (\cref{maps:src:D.3.1}).
\item\hypertarget{step-III-G-3}{} The quadric smoothing has an
explicit orientation-preserving identification with $TS^3$.
The complement of the exceptional sphere is two-connected,
and the primitive boundary column $(1,\pm1)^t$ gives
$H_3=\mathbb Z^2$, $H_2=H_4=0$ (\cref{maps:src:D.3.2}).
\item\hypertarget{step-III-G-4}{} The middle pairing is the
standard unimodular skew form and the normal invariant vanishes.
Wall's classification and $\Theta_6=0$ identify the complex-oriented
smoothing with $S^3\times S^3$ (\cref{maps:src:D.3.2}).
\item\hypertarget{step-III-G-5}{} The root canonical formula and
varying torus periods give algebraic dimension one; formal
cohomology across the exceptional curve gives
$(h^{0,0},h^{0,1},h^{0,2},h^{0,3})=(1,1,0,0)$.
The algebraic-reduction fibre maps also give the comparison with
Calabi--Eckmann deformations
(\cref{maps:src:D.3.3,maps:src:D.3.4}).
\end{enumerate}

\subsection{Known and new}
\begin{center}\small
\begin{tabularx}{\textwidth}{@{}lX@{}}\toprule
Theorem & Contribution and principal source relationship\\\midrule
A & New relative realization for Engel's quotient
\cite[Remark~8.3, pp.~22--23]{E26}.\\
B & New map classes; unstable suspension and the real Adams
formula are \cite[Theorem~4.1, pp.~76--77]{mapsToda52} and
\cite[Theorem~4.14, p.~423]{mapsAPS2}.\\
C & New cusp framing and marked cell presentation;
the sphere construction is \cite[Theorem~4.1, p.~12]{E26}.\\
D & New full-fan obstruction and normal extension;
covering classification is \cite[Theorem~1.38, pp.~67--68]{Hatcher2002}.\\
E & New complete cover invariant; marked smooth recognition uses
\cite[Theorems~3 and~5, pp.~795--796; \S4.1, p.~809]{manZhubr75}.\\
F & New natural-line and weighted classifications within
\cite[Remark~8.5, p.~24]{E26}.\\
G & New moving-zero family and its smooth type; the surgery
classification is \cite[Theorem~2.6, p.~8]{BHS19}.\\
\bottomrule\end{tabularx}\end{center}

\subsection{Methodology}
This text is AI-authored. GPT-5.6 Sol (OpenAI), in Codex at
reasoning effort ultra, produced the working edition during
24 August--9 September 2026. GPT-6 Astra (OpenAI), in ChatGPT Pro
mode, carried out the research of 12--16 November, independent
re-derivations during 16--18 November, writing during 16--20 November,
and publication work during 20--26 November 2026.

Claude Fable 5.2 (Anthropic), at effort high and at effort max,
and Claude Opus 5.5 (Anthropic), at effort max, performed readings,
intake checks and recomputations during 12--20 November 2026.
Claude Opus 5.5 at effort max rewrote Paper~V's abstract and
introduction on 20 November, read all of Paper~I on 21--22 November
and adjusted its deposited text and links, and read all of
Paper~II on 24 November, and began its reading of Paper~III on
25 November. It also checked the publication material received
during 23--25 November, including arithmetic and source-scope
questions; its Paper~III reading and computations continued on
26 November.

The models worked from explicit statements, marked geometric data
and primary sources. Proposed proofs and computations were passed
between them for comparison; disagreements were resolved at the
affected statement, map or parameter range. The checks include
derivations written before consulting earlier proofs, independent
proofs, exact arithmetic, certified rational bounds and comparison with
the cited literature. Each computation is described mathematically
where it is used, and the programs are listed at the end of the paper.
During the publication work, existing proofs were reread in full
against their replacements, with coordinate maps, inverse maps and
exceptional parameters checked explicitly. Those readings record
their use and rechecking of existing arguments. The records state
the scope of each check and preserve every corrected assertion.
The records associate 1,253 entries with the current statements of this
paper: 422 examinations of proof steps, 2 independent re-derivations, 58
computations or formal certificates, 651 readings of proofs and
applications, 90 uses of background results without a fault found, and 30
refutation records associated with their corrected statements. These entries
belong to 964 recorded evidence groups. Each event is counted once; entries
sharing an evidence group retain that common group. One argument or
calculation can be associated with several statements and kinds of
examination. Each outcome specifies its scope and prior proof exposure.

No human has checked the mathematics.

We dedicate this text to the public domain under
\href{https://creativecommons.org/publicdomain/zero/1.0/}{CC0 1.0 Universal}.
The standing address for statement checks, provenance and programs
is the \href{\SsixPublishRecords}{records page for Paper~III}.
The repository commit is \texttt{\SsixRepositoryCommit}.

\subsection{How the paper is organized}
The marked periods and relative sphere quotient come first. Integral
homology of the cyclic covers supplies the sphere framing; the base
map, angle-forgetting maps and cusp-frame calculation then determine
the homotopy classes. The three-filling cell model and finite-band
maps precede the Kummer construction and the complete cyclic
invariants. The elliptic-surface sections develop the initial data,
cap maps, sphere patterns and weighted smooth types, followed by the
OS49 family and the conifold smoothing. Global torus quotients lead
to the compact cusp carriers, their framed nullbordism and sphere-map
coefficients, and the metric scaling limits. The final part contains
the general section-line and multiplicity questions, auxiliary-cover
compactifications, the two raw carrier classes, and the remaining
cap-map and smooth-type classification in \cref{maps:src:D.2.4}.

\begin{figure}[htbp]\centering
\begin{tikzpicture}[>=stealth,every node/.style={font=\small,align=center},
 box/.style={draw,rounded corners,minimum width=2.25cm,minimum height=.7cm}]
\node[box] (P) at (1.5,0) {Marked periods\\and fillings};
\node[box] (L) at (6.5,0) {Elliptic surface\\and section line};
\node[box] (A) at (0,-1.35) {\hyperlink{main-III-A}{Theorem A}};
\node[box] (E) at (3.25,-1.35) {\hyperlink{main-III-E}{Theorem E}};
\node[box] (D) at (6.5,-1.35) {\hyperlink{main-III-D}{Theorem D}};
\node[box] (B) at (0,-2.75) {\hyperlink{main-III-B}{Theorem B}};
\node[box] (F) at (3.25,-2.75) {\hyperlink{main-III-F}{Theorem F}};
\node[box] (G) at (6.5,-2.75) {\hyperlink{main-III-G}{Theorem G}};
\node[box] (C) at (0,-4.1) {\hyperlink{main-III-C}{Theorem C}};
\draw[->,semithick] (P) -- (A);
\draw[->,semithick] (P) -- (E);
\draw[->,semithick] (P) to[bend left=13] (D);
\draw[->,semithick] (A) -- (B);
\draw[->,semithick] (E) -- node[above,sloped,font=\scriptsize]{homology} (B);
\draw[->,semithick] (B) -- (C);
\draw[->,semithick] (E) -- node[right,font=\scriptsize]{coprime\\weights} (F);
\draw[->,semithick] (L.south west) -- (4.9,-.7) -- (4.9,-2.75) -- (F.east);
\draw[->,semithick] (L.east) to[bend left=30] (G.east);
\end{tikzpicture}
\caption{Dependencies of the indicated arguments. The arrow from E
to B is the integral homology calculation in
\cref{maps:src:D.1.4} used by \cref{maps:src:B.1.1}; the arrow
from E to F is the marked coprime-weight identification in
\cref{maps:src:D.2.48}. The arrow from B to C uses the same
Hopf-circle class in \cref{maps:src:B.3.1}.}
\end{figure}
\clearpage

\section{Conventions and the marked family}\hypertarget{tag-034J}{}\label{sec:III:conventions}
\subsection{Markings and period parameters}
\hypertarget{tag-Q004}{}\label{mapsIII:sec:marked-conventions}

\begin{definition}[The marked family]
\hypertarget{tag-Q000}{}\label{mapsIII:def:marked-family}
{\normalfont\scriptsize tag Q000}\par
We use the marked $(3,4,\infty)$ family over
$B=\mathbb P^1_t$, with $(p_1,p_2,p_0)=(0,1,\infty)$.
Its regular-fibre lattice is
\[
 \Lambda=\mathbb Z\langle\widehat\gamma,\widehat u,
                              \widehat w,\widehat\delta\rangle,
\]
and $\gamma$ is the covector dual to $\widehat\gamma$.
The fibre-linear matrices of the analytic deck generators are
$A_1,A_2,M_0$, with
$A_1A_2M_0=I$. Base meridians are counterclockwise; the analytic
root generators satisfy
\[
 s_j(g_jz)=e^{-2\pi i/m_j}s_j(z),\qquad
 (m_1,m_2)=(3,4),\qquad g_1g_2g_0=1.
\]
The multiple fibres have $f^*p_j=m_jS_j$, and
$W=f^{-1}(p_0)$ is the reduced cusp. Write
\[
 \varepsilon_1=(1,2,-4,0)^t,\qquad
 \varepsilon_2=(1,3,-3,0)^t,\qquad
 v_j=\ell_j\varepsilon_j+k_j\widehat\delta\quad(j=1,2),
\]
where $\ell_j,k_j\in\mathbb Z$. In the real period marking,
the finite root generator acts by
$(s,x)\mapsto(e^{-2\pi i/m_j}s,A_jx+v_j/m_j)$.
The integer $\ell_0$ is the cusp clutching winding: after the
vanishing fibre cycles are killed, the remaining fibre generator
$z$ and clockwise peripherals $x,y$ satisfy
$x^3=z^{\ell_1}$, $y^4=z^{\ell_2}$ and $xy=z^{\ell_0}$.
Put $p=12\ell_0-4\ell_1-3\ell_2$. The standard twists have
$(\ell_1,\ell_2,k_1,k_2)=(1,-1,0,0)$, hence
$v_1=\varepsilon_1$ and $v_2=-\varepsilon_2$.
Admissibility requires $3\nmid\ell_1$ and $2\nmid\ell_2$;
hence $p$ is prime to six.
All homology and cohomology groups have integral coefficients
unless another coefficient group is displayed. Orientations are
the complex orientation, outward stabilization, and the stated
order of normal coordinates.
\end{definition}

\begin{remark}[Transport and dual coordinates]
\hypertarget{tag-Q001}{}\label{mapsIII:rem:dual-transport}
{\normalfont\scriptsize tag Q001}\par
For a homology transport $x\mapsto Ax$, parallel transport of
a covector column is $\alpha\mapsto T\alpha$, where
$T=A^{-t}$: the identity
\[
 (T\alpha)^t(Ax)=\alpha^tx
\]
preserves evaluation, and the inverse transport is $\alpha\mapsto
A^t\alpha$. The pullback by the map inducing $A$ on homology
is also written explicitly as $\alpha\mapsto A^t\alpha$.
In fibre calculations, a coordinate tuple $(g,u,w,d)$ denotes
$g\widehat\gamma+u\widehat u+w\widehat w+d\widehat\delta$;
its inverse sends a cycle to its four coefficients in this basis.
\end{remark}

\begin{definition}[Period parameters]
\hypertarget{tag-Q002}{}\label{mapsIII:def:period-parameters}
{\normalfont\scriptsize tag Q002}\par
The period matrix and its scalar parameter are
\[
 \Pi_c=\begin{pmatrix}6\mu&\tau&1&0\\
             \beta^{\rm part}+c&\mu&0&1\end{pmatrix},
 \qquad w=e^{-\pi i c}.
\]
We use the normalization $(\beta^{\rm part}+\tau)(p_0)=0$ and set
\[
 D_{\rm part}=\operatorname{Im}\beta^{\rm part}
       -\frac{6(\operatorname{Im}\mu)^2}{\operatorname{Im}\tau},
 \qquad M=-\frac{3\log6}{\pi},\qquad
 D=D_{\rm part}+\operatorname{Im}c.
\]
The period calculation in
\SeriesRef{fam:src:C.2.3} gives
$\max D_{\rm part}=M$; the admissible range is
$\operatorname{Im}c<-M$. A positive loop in the $w$-plane lifts
to $c(r)=c_0-2r$, since
$e^{-\pi i(c_0-2r)}=e^{-\pi i c_0}e^{2\pi i r}$.
\end{definition}

\begin{definition}[The change of marking]
\hypertarget{tag-Q003}{}\label{mapsIII:def:engel-marking}
{\normalfont\scriptsize tag Q003}\par
Engel's indices $(0,1,\infty)$ correspond to $(1,2,0)$ here
\cite[Remark~8.3, pp.~22--23]{E26}. We denote his line
bundle by $M_S=\mathcal O_S(P-O)$. The coordinate change is
$x\mapsto P_Ex$, where
\[
 P_E=\begin{pmatrix}0&-2&-1&-2\\2&-1&0&-1\\
                      0&2&1&3\\-1&0&0&0\end{pmatrix},
 \qquad
 P_E^{-1}=\begin{pmatrix}0&0&0&-1\\-1&-1&-1&-2\\
                         -1&2&0&4\\1&0&1&0\end{pmatrix}.
\]
Matrix multiplication gives $P_EP_E^{-1}=P_E^{-1}P_E=I$.
For $\psi=(0,0,0,1)$ the character and twist coordinates are
\[
 \psi P_E=-\gamma,\qquad
 a_{j-1}=-P_Ev_j/m_j,\qquad
 v_j=-m_jP_E^{-1}a_{j-1}.
\]
\end{definition}

\subsection{Radial period coordinates and the compact quotient}
\hypertarget{tag-02QA}{}\label{mapsIII:carrier-radial-prerequisites}

\begin{definition}[Radial cusp data]
\hypertarget{tag-Q005}{}\label{mapsIII:def:radial-cusp-data}
{\normalfont\scriptsize tag Q005}\par
Use the cusp coordinate $t_c=1/t$ on the marked base and the
logarithmic coordinate $t_c=e^{2\pi i s}$. Put
$\beta=\beta^{\rm part}+c$, $h=\tau-s$, and $b=\beta+\tau$.
The inverse substitutions are $\tau=s+h$ and $\beta=b-s-h$.
The cusp expansion of \cite[Proposition~3.21, p.~22]{AC26} gives
\[
 \Pi=[Z\mid I_2],\qquad
 Z=s\begin{pmatrix}0&1\\-1&0\end{pmatrix}
       +\begin{pmatrix}6\mu(t_c)&h(t_c)\\
                       b(t_c)-h(t_c)&\mu(t_c)\end{pmatrix}.
\]
The functions $\mu,b,h$ are holomorphic on a disc containing zero.
On the positive radial ray put $t_c=e^{-2\pi T}$ and $s=iT$, and set
\begin{equation}
 \begin{aligned}
 \mu_T&=\alpha_T+im_T,& h_T&=\eta_T+iH_T,& b_T&=\xi_T+iB_T,\\
 c_T&=\xi_T-\eta_T,& L_T&=T+H_T,& q_T&=L_T-B_T,\\
 D_T&=L_Tq_T+6m_T^2.
 \end{aligned}
 \hypertarget{tag-03A3}{}\label{wb:cusp-spec-radial-notation}
\end{equation}
The components $\alpha_T,m_T,\eta_T,H_T,\xi_T,B_T$ and
the combinations $c_T,L_T,q_T,D_T$ are real. Write
$m=\operatorname{Im}\mu(0)$, $H=\operatorname{Im}h(0)$ and
$B=\operatorname{Im}b(0)$. Choose $T_0$ so that the radial ray is in
the holomorphic coordinate disc and $D_T>0$ for $T\geq T_0$.
We identify $\mathbb C^2$ with $\mathbb R^4$ in the order
$(x_1,y_1,x_2,y_2)$ and use the Euclidean metric and its flat quotient
metrics.
\end{definition}

\begin{proof}
For a holomorphic function $g$ near zero, the function
$(g(t)-g(0))/t$ extends holomorphically to zero and is bounded on
a smaller closed disc. Substitution of $t_c=e^{-2\pi T}$ gives
\begin{equation}
 m_T=m+O(e^{-2\pi T}),\qquad
 H_T=H+O(e^{-2\pi T}),\qquad
 B_T=B+O(e^{-2\pi T}),
 \hypertarget{tag-039V}{}\label{wb:cusp-spec-holomorphic-errors}
\end{equation}
and the same estimates for the real parts. Hence
$D_T=T^2+O(T)$, which is positive for all sufficiently large $T$.
This proves the existence of the stated choice of $T_0$.
\end{proof}

\begin{proposition}[Real periods and their dual coordinates]
\hypertarget{tag-02QB}{}\label{mapsIII:carrier-real-periods}
{\normalfont\scriptsize tag 02QB}\par
The radial fibre and its real period matrix are
\[
 F_T=\mathbb R^4/P_T\mathbb Z^4,\qquad
 P_T=\begin{pmatrix}
 6\alpha_T&\eta_T&1&0\\
 6m_T&L_T&0&0\\
 c_T&\alpha_T&0&1\\
 -q_T&m_T&0&0
 \end{pmatrix}.
\]
Its determinant is $-D_T$. Its complete inverse transpose is
\begin{equation}
 P_T^{-\mathsf T}=
 \begin{pmatrix}
 0&0&1&0\\[2pt]
 \dfrac{m_T}{D_T}&\dfrac{q_T}{D_T}&
 \dfrac{-6m_T\alpha_T-q_T\eta_T}{D_T}&
 \dfrac{-m_Tc_T-q_T\alpha_T}{D_T}\\[7pt]
 0&0&0&1\\[2pt]
 -\dfrac{L_T}{D_T}&\dfrac{6m_T}{D_T}&
 \dfrac{6(L_T\alpha_T-m_T\eta_T)}{D_T}&
 \dfrac{L_Tc_T-6m_T\alpha_T}{D_T}
 \end{pmatrix}.
 \hypertarget{tag-039O}{}\label{wb:cusp-spec-full-inverse-transpose}
\end{equation}
Thus $(P_T\mathbb Z^4)^*=P_T^{-\mathsf T}\mathbb Z^4$.
\end{proposition}

\begin{proof}
At $s=iT$ the complex period matrix is
\[
 Z_T=\begin{pmatrix}6\mu_T&h_T+iT\\
                          b_T-h_T-iT&\mu_T\end{pmatrix}.
\]
The real-linear isomorphism
\[
 R:\mathbb C^2\longrightarrow\mathbb R^4,
 \qquad R(z_1,z_2)=(\operatorname{Re}z_1,
 \operatorname{Im}z_1,\operatorname{Re}z_2,
 \operatorname{Im}z_2)
\]
sends the four columns of $[Z_T\mid I_2]$ to the four
columns of $P_T$ in the stated order. It therefore induces
the quotient isomorphism $[z]\mapsto[R(z)]$, whose inverse
is induced by $(x_1,y_1,x_2,y_2)\mapsto(x_1+iy_1,x_2+iy_2)$.
Expansion along the last two columns gives
$\det P_T=-L_Tq_T-6m_T^2=-D_T$. The radial choice
$T\geq T_0$ gives $D_T>0$.

To solve $P_T^{\mathsf T}x=(r,s,n_1,n_2)^{\mathsf T}$,
the last two equations give $x_1=n_1$ and $x_3=n_2$.
Set
\[
 R_0=r-6\alpha_Tn_1-c_Tn_2,\qquad
 S_0=s-\eta_Tn_1-\alpha_Tn_2.
\]
The other two equations and their solution are
\[
 \begin{pmatrix}6m_T&-q_T\\L_T&m_T\end{pmatrix}
 \binom{x_2}{x_4}=\binom{R_0}{S_0},
 \qquad
 \binom{x_2}{x_4}=
 \frac1{D_T}\begin{pmatrix}m_T&q_T\\-L_T&6m_T\end{pmatrix}
 \binom{R_0}{S_0}.
\]
Substituting $R_0,S_0$ proves every entry of
\eqref{wb:cusp-spec-full-inverse-transpose}.

Under the Euclidean pairing, a vector $\ell$ is integral
on $P_T\mathbb Z^4$ precisely when
$P_T^{\mathsf T}\ell\in\mathbb Z^4$. The maps
\[
 n\longmapsto P_T^{-\mathsf T}n,
 \qquad \ell\longmapsto P_T^{\mathsf T}\ell
\]
are inverse isomorphisms between $\mathbb Z^4$ and the
dual lattice. This proves its asserted formula.
\end{proof}

\begin{proposition}[Exact compact quotient of the native cusp fibre]
\hypertarget{tag-039Y}{}\label{wb:cusp-spec-native-quotient}
{\normalfont\scriptsize tag 039Y}\par
Let $e_1,\ldots,e_4$ be the coordinate unit vectors and
$v_{\gamma,T},v_{u,T}$ the first two columns of $P_T$.
Define the orthogonal projection and its projected period matrix by
\[
 \rho(x_1,x_2,x_3,x_4)=(x_2,x_4),\qquad
 Y_T=\begin{pmatrix}6m_T&L_T\\-q_T&m_T\end{pmatrix}.
\]
Then $\det Y_T=D_T>0$, so $Q_T=\mathbb R^2/Y_T\mathbb Z^2$
is a compact real torus, and projection induces the exact sequence
of compact real Lie groups
\[
 0\longrightarrow
 \mathbb T^2_{\mathrm{re}}=
 \frac{\mathbb Re_1\oplus\mathbb Re_3}
      {\mathbb Ze_1\oplus\mathbb Ze_3}
 \longrightarrow F_T\xrightarrow{\ \pi_T\ }Q_T
 \longrightarrow0.
\]
This is a Riemannian quotient for the displayed flat metrics.
There is a constant $C$, independent of $T\geq T_0$, such that
for every $(a,b)\in\mathbb R^2$,
\[
 \|Y_T(a,b)^{\mathsf T}\|\geq(T-C)\sqrt{a^2+b^2},
 \qquad \operatorname{sys}(Q_T)\geq T-C.
\]
Give $Q_T^2$ the Euclidean $\ell^2$ product metric.
For every $R>0$, the radius-$R$ balls centred at the identities
of $Q_T$ and $Q_T^2$ are exactly Euclidean once $T-C>4R$.
Consequently $(Q_T,0)$ and $(Q_T^2,0)$ converge to
$(\mathbb R^2,0)$ and $(\mathbb R^4,0)$, respectively,
locally isometrically, smoothly, and in pointed
Gromov--Hausdorff topology.
\end{proposition}

\begin{proof}
Put $\Gamma_T=P_T\mathbb Z^4$ and
$V=\mathbb Re_1\oplus\mathbb Re_3$. The columns of $P_T$ give
\[
 \Gamma_T=\mathbb Zv_{\gamma,T}\oplus\mathbb Zv_{u,T}
                 \oplus\mathbb Ze_1\oplus\mathbb Ze_3.
\]
 Multiplication gives
\[
 \rho P_T=[Y_T\mid0],\qquad
 \det Y_T=6m_T^2+L_Tq_T=D_T>0.
\]
Thus $\rho(\Gamma_T)=Y_T\mathbb Z^2$, and
$\pi_T([x])=[\rho x]$ is a surjective smooth homomorphism.
For a period $P_Tn$ lying in $V$, its first two coefficients
satisfy $Y_T(n_1,n_2)^{\mathsf T}=0$. Hence $n_1=n_2=0$,
and
\[
 \Gamma_T\cap V=\mathbb Ze_1\oplus\mathbb Ze_3.
\]
If $\pi_T([x])=0$, write
$\rho x=Y_T(a,b)^{\mathsf T}$ with $a,b\in\mathbb Z$.
Subtracting $P_T(a,b,0,0)^{\mathsf T}$ gives a representative
in $V$. Therefore the kernel is exactly the image of
$V/(\mathbb Ze_1\oplus\mathbb Ze_3)$ under $[v]\mapsto[v]$,
and this inclusion is injective.

The differential of $\rho$ takes $V^\perp$ isometrically
onto $\mathbb R^2$. The distance between two kernel-torus
orbits, represented by $x,y\in\mathbb R^4$, is
\[
 \begin{aligned}
 \inf_{v\in V,\,\lambda\in\Gamma_T}
       \|x-y-v-\lambda\|
 &=\inf_{\lambda\in\Gamma_T}
       \|\rho(x-y-\lambda)\|\\
 &=\inf_{n\in\mathbb Z^2}
       \|\rho x-\rho y-Y_Tn\|.
 \end{aligned}
\]
This is the displayed flat quotient distance on $Q_T$.

Write
\[
 Y_T=TJ+E_T,\qquad
 J=\begin{pmatrix}0&1\\-1&0\end{pmatrix},\qquad
 E_T=\begin{pmatrix}6m_T&H_T\\B_T-H_T&m_T\end{pmatrix}.
\]
The residual functions are bounded on the radial ray by
holomorphy. Thus one may take a finite constant
\[
 C=\sup_{T\geq T_0}
     \sqrt{37m_T^2+H_T^2+(B_T-H_T)^2},
\]
which bounds the operator norm of $E_T$ by its Frobenius
norm. Since $J$ is orthogonal, the reverse triangle
inequality gives $\|Y_Tv\|\geq(T-C)\|v\|$.
For a nonzero integer vector it gives
\[
 \|Y_Tn\|\geq\max\{T-C,0\}\|n\|\geq T-C.
\]
The systole of a Euclidean torus is the length of its
shortest nonzero lattice vector. The product lattice
$Y_T\mathbb Z^2\oplus Y_T\mathbb Z^2$ has the same
shortest length, proving both systolic bounds.

Let $x,y$ lie in a Euclidean radius-$R$ ball in dimension
two or four. For every nonzero vector $\lambda$ of the
corresponding lattice, the assumption $T-C>4R$ gives
\[
 \|x-y+\lambda\|\geq\|\lambda\|-\|x-y\|
 \geq(T-C)-2R>2R\geq\|x-y\|.
\]
The quotient distance between these points is therefore
$\|x-y\|$. Every point of the quotient ball has a
representative in the Euclidean ball, obtained by taking
a shortest representative of its lattice coset. The
covering projection is consequently an isometry onto
the entire radius-$R$ ball. These maps are compatible
on smaller balls and pull back the metrics exactly to
the Euclidean metric. They give the stated pointed
local isometries and smooth convergence; the closed
balls are eventually exactly isometric, which also
proves pointed Gromov--Hausdorff convergence.
\end{proof}

\section{Sphere maps, coverings and framings}\hypertarget{tag-034K}{}\label{sec:III:maps}
\subsection{Relative realization of sphere maps}
\begin{proposition}[Relative realization of the twelve homotopy classes]\hypertarget{tag-02J3}{}\label{maps:src:A.2.2}
{\normalfont\scriptsize tag 02J3}\par


Let $N$ be an open neighbourhood of the cusp $W$ in an original
sphere member, and fix the quotient homeomorphism $\varphi:Q\to S^3$.
Every class in $[S^6,S^3]\cong\mathbb Z/12$ has a representative
$\varphi q$, where $q:S^6\to Q$ is smooth and agrees with the literal
quotient on $S^6\setminus N$.
For the designated cusp solid torus $T\subset Q$ and any basepoint
$y\in T$, the inclusion followed by $\varphi$ induces
\[
 (\varphi|_T)_*:\pi_6(T,y)\longrightarrow\pi_6(S^3,\varphi(y)),
 \qquad \operatorname{im}(\varphi|_T)_*=\{0\}.
\]

\end{proposition}


\begin{proof}
By
\SeriesRef{mapsI:src:A.2.1}, choose the smooth quotient replacement
$e_{\rm q}:S^6\to Q$ agreeing with the literal quotient outside a
cusp neighbourhood contained in $N$.  Fix a homeomorphism
$\varphi:Q\to S^3$ supplied by the original $|p|=1$ case of
\SeriesRef{mapsI:src:A.1.3}, and put $e=\varphi e_{\rm q}$.
Use the unit-quaternion model of $S^3$, with its continuous
multiplication and inverse.

Choose $x_0\in S^6$, put $y_0=e(x_0)$, and take stereographic
coordinates $S^6\setminus\{x_0\}\cong\mathbb R^6$, with
$x_0$ corresponding to $\infty$.  Choose an open coordinate ball
$D_0=D_R(b)$ and a larger open ball $D_1$ such that
$\overline{D_0}\subset D_1$ and $\overline{D_1}\subset N$.
For $r>0$, $v\in S^5$ and $0\le t<1$, define in these coordinates
\[
 c_t(b+rv)=b+\frac{r}{1-t\min(r/R,1)}v,
 \qquad c_t(b)=b,\qquad c_t(\infty)=\infty.
\]
At $t=1$, use the same formula for $r<R$ and send every point
with $r\ge R$ to $\infty$.  Transport these maps back to $S^6$.
They form a continuous based homotopy: at $(r,t)=(R,1)$ the
radius tends to infinity from either side, and at the point at
infinity it tends to infinity uniformly in $t$, since the displayed
radius is at least $r$.  At $r=0$ continuity is immediate from
the formula.  Thus, with $c=c_1$,
\[
 c_0=1_{S^6},\qquad c_t(x_0)=x_0,
 \qquad c|_{S^6\setminus D_0}=x_0.
\]

For any based representative $g:(S^6,x_0)\to(S^3,y_0)$, set
\[
 d_g(x)=e(x)^{-1}g(x),\qquad
 e_g(x)=e(x)d_g(c(x)).
\]
Since $d_g(x_0)=1$, the map $e_g$ equals $e$ at every point of
$S^6\setminus D_0$, hence on the whole exterior $S^6\setminus N$.
The continuous homotopy
\[
 H_g(x,t)=e(x)d_g(c_t(x))
\]
is based at $y_0$ and has endpoints $g$ and $e_g$.  Therefore
$e_g$ represents exactly the prescribed based class $[g]$.

The based and unbased conventions give the same class set here.
For an arbitrary map $f:S^6\to S^3$, the formula
\[
 f\longmapsto \bigl(x\longmapsto y_0 f(x_0)^{-1}f(x)\bigr)
\]
makes it based.  Left multiplication along a path from $1$ to
$y_0f(x_0)^{-1}$ gives an unbased homotopy to that based map.
For an unbased homotopy $H$, the formula
$y_0H(x_0,t)^{-1}H(x,t)$ is a based homotopy of its normalized
endpoints.  Normalization and forgetting the basepoint are
therefore inverse bijections.  Left translation by $y_0^{-1}$
identifies the based target with $(S^3,1)$.  Since
$\pi_6(S^3)=\mathbb Z/12$
\cite[Theorem~5.41(b), p.~573]{HatcherSpectralSequences},
the maps $e_g$ realize all twelve distinct unbased classes while
having the same prescribed exterior restriction.

The smooth realization assertion takes place in the original
quotient target $Q$.  Put $A=S^6\setminus D_1$.  The map
$\varphi^{-1}e_g:S^6\to Q$ equals the smooth map $e_{\rm q}$ on
an open neighbourhood of the closed set $A$, because
$\overline{D_0}\subset D_1$.  Relative Whitney approximation
gives a smooth map into $Q$ and a homotopy stationary on $A$
\cite[Theorem~2.1, arXiv version~4, p.~2]{IwaseWhitney2022}.
This preserves the entire exterior $S^6\setminus N\subset A$.
Composing this continuous homotopy with $\varphi$ preserves the
class in $[S^6,S^3]$.

Finally let $T\subset Q$ be the designated cusp solid torus,
choose product coordinates $T=S^1\times D^2$, and let
$j:T\hookrightarrow Q\xrightarrow{\varphi}S^3$.
For a basepoint $y=(e^{2\pi i\theta_0},z_0)$, the maps
$(u,z)\mapsto(u,(1-s)z+sz_0)$ deform $T$ onto
$S^1\times\{z_0\}$.  The circle component of a based map
$S^6\to T$ lifts to a function $\theta:S^6\to\mathbb R$ with
$\theta(x_0)=\theta_0$, by the covering-space lifting property
\cite[Proposition~1.33, p.~61]{Hatcher2002}.
Writing the other component as $z(x)$, the based homotopy
\[
 (x,s)\longmapsto
 \left(e^{2\pi i((1-s)\theta(x)+s\theta_0)},
       (1-s)z(x)+sz_0\right)
\]
contracts the map to $y$.  Consequently $\pi_6(T,y)=0$, and
\[
 j_*:\pi_6(T,y)\longrightarrow\pi_6(S^3,j(y))
\]
has image $\{0\}$, represented by the constant map in $T$.
\end{proof}

\Needspace{6\baselineskip}
\begin{remark}[The cusp replacement]
\hypertarget{tag-Q007}{}\label{mapsIII:rem:cusp-replacement}
{\normalfont\scriptsize tag Q007}\par
The quotient construction of \cite[Remark~8.3, pp.~22--23]{E26}
is literal over the regular part and the two finite fillings.
At the cusp, \SeriesRef{mapsI:src:A.2.1} constructs its replacement
with prescribed outer collar and image in the designated solid torus,
and proves independence of that replacement in its stated relative
homotopy class. That relative result supplies the homotopy class
of the quotient map used below.
\end{remark}


\phantomsection\hypertarget{tag-033S}{}\label{qa:III:layout:12}
\subsection{Cyclic covers and integral homology}
\begin{definition}[The cyclic-cover notation]
\hypertarget{tag-Q008}{}\label{mapsIII:def:cyclic-cover}
{\normalfont\scriptsize tag Q008}\par
For an admissible marked member $X_p$, put
$p=12\ell_0-4\ell_1-3\ell_2$ and $d=|p|$.
Write $\widetilde X_p\to X_p$ for its universal covering and
$\widetilde W_p$ for the inverse image of the cusp $W$.
We write $dP_6$ for the blowup of $\mathbb P^2$ at its three coordinate points.
\end{definition}

\begin{theorem}[Euler characteristic of cyclic covers]\hypertarget{tag-02K7}{}\label{maps:src:D.1.1}
{\normalfont\scriptsize tag 02K7}\par


For every admissible member $X_p$,
the universal cover $\widetilde X_p$ is a compact complex threefold of
Euler characteristic $2d$. Its Betti numbers necessarily satisfy
\[
 b_3(\widetilde X_p)=2+2b_2(\widetilde X_p)-2d,
 \qquad b_2(\widetilde X_p)\geq d-1.
\]
Consequently, for $d>1$ its second integral homology cannot vanish,
and it cannot be diffeomorphic to $\#^{d-1}(S^3\times S^3)$.
\end{theorem}


\begin{proof}
Admissibility gives
$3\nmid\ell_1$ and $2\nmid\ell_2$. Since
$p\equiv-\ell_1\pmod3$ and $p\equiv\ell_2\pmod2$,
we have $p\ne0$ and $d=|p|\geq1$.

Use the marked gluings of
\SeriesRef{def:II:marked-conventions} and
\SeriesRef{def:II:analytic-picard-data}.
The cusp kills the two toric fibre cycles. The first finite
conjugation then kills the $\widehat u$-cycle, since it sends that
cycle to $-\widehat w+\widehat\delta$.
The remaining fibre generator $z$ is central. The clockwise finite
peripherals $x,y$ satisfy $x^3=z^{\ell_1}$ and $y^4=z^{\ell_2}$.
The cusp map is exterior-to-local and is composed with translation
by the section. Its positive-coordinate winding generator is the
inverse of the clockwise peripheral, so that peripheral maps to
$z^{-\ell_0}$. Their ordered product therefore gives $xy=z^{\ell_0}$.
Thus van Kampen gives
\[
\pi_1(X_p)\simeq
\left\langle z,x,y:\ [z,x]=[z,y]=1,\quad
x^3=z^{\ell_1},\quad y^4=z^{\ell_2},\quad xy=z^{\ell_0}\right\rangle.
\]
Indeed
\[
x=z^{4\ell_0-\ell_1-\ell_2},\qquad
y=z^{\ell_1+\ell_2-3\ell_0},
\]
and the remaining relation is $z^p=1$. Conversely, sending $z$ to
$1$, $x$ to $4\ell_0-\ell_1-\ell_2$, and $y$ to
$\ell_1+\ell_2-3\ell_0$ modulo $p$ respects all the displayed
relations and is inverse to $1\mapsto z$.
Consequently the universal cover has exactly $d$ sheets.

The prescribed cusp fan has one component orbit, three open
$\mathbb C^*$ double strata and two triple points. Its component
interior is $(\mathbb C^*)^2$, so its reduced central fibre has
Euler number $2$. The finite fibres have Euler number zero by their
free finite torus covers. The regular torus bundle has compactly
supported Euler number zero by a cell stratification of the
punctured base, including arbitrary monodromy. Hence $\chi(X_p)=2$.

The finite universal covering is unramified and pulls back the
complex atlas. It is therefore a compact complex threefold with
its complex orientation. Lifting a finite CW structure gives
exactly $d$ lifts of every cell, so
\[
\chi(\widetilde X_p)=d\,\chi(X_p)=2d.
\]
Simple connectivity gives $b_1=0$. Oriented Poincar\'e duality gives
$b_5=b_1=0$, $b_4=b_2$, and $b_0=b_6=1$. Thus
\[
2d=2+2b_2(\widetilde X_p)-b_3(\widetilde X_p).
\]
The inequality follows from $b_3\geq0$. For $d>1$, the resulting
positive value of $b_2$ is the free rank of
$H_2(\widetilde X_p;\mathbb Z)$, so that integral group cannot vanish.

A connected sum of $r$ copies of $S^3\times S^3$ has Euler
characteristic $2-2r$: each summand has Euler characteristic zero,
and a connected sum of closed six-manifolds subtracts two from the
sum of their Euler characteristics. For $r=d-1$, its Euler number
is $4-2d$, which equals $2d$ only for $d=1$.

For a concrete admissible example take
$(\ell_0,\ell_1,\ell_2)=(0,2,-1)$, $k_1=k_2=0$, $c=-i$, and the zero
cusp section. The finite coefficients satisfy $3\nmid2$ and
$2\nmid-1$, and $\operatorname{Im}c=-1<(3/\pi)\log6=-M$.
The marked finite deck and overlap constructions therefore give
this admissible member. Here $p=-5$ and $d=5$, so the universal
cover has Euler characteristic $10$, whereas
$\#^4(S^3\times S^3)$ has Euler characteristic $-6$.
\end{proof}

\begin{theorem}[The cusp of a cyclic cover]\hypertarget{tag-02KB}{}\label{maps:src:D.1.2}
{\normalfont\scriptsize tag 02KB}\par


The inverse image $\widetilde W_p$ of the cusp is connected. Its
normalization has $d$ components, each isomorphic to $dP_6$. Its
integral homology is torsion free, with ranks
\[
 (b_0,b_1,b_2,b_3,b_4)(\widetilde W_p)=(1,2,d+3,2,d).
\]
Here the covering of the dual torus is the sublattice
$\mathbb Z\widehat u\oplus d\mathbb Z\widehat\gamma$ of
$\overline\Lambda=\mathbb Z\widehat\gamma\oplus\mathbb Z\widehat u$.
\end{theorem}


\begin{proof}
\begin{enumerate}[label=\arabic*.,leftmargin=*]
\item The presentation and inverse cyclic maps in
\cref{maps:src:D.1.1} send the regular-fibre lattice onto
$\mathbb Z/d$ by its primitive $\gamma$ coordinate.
Admissibility makes $d$ coprime to six.
The period cover of the cusp is $Y_\epsilon\to N_0$ with period group
$\overline\Lambda$. Its central divisor $Y_0$ is connected: its
$dP_6$ components meet along the edges of the connected triangular
lattice. Moreover its components and double curves are simply
connected, and its triple points give exactly the triangle relations
in van Kampen's presentation. That presentation is the edge-path
presentation of the triangulated plane, hence is trivial. Thus the
period endpoint map identifies $\pi_1(W)$ with $\overline\Lambda$.

The restriction of the cyclic covering is therefore the
quotient map
\[
 Y_\epsilon/\Gamma'\longrightarrow Y_\epsilon/\overline\Lambda,
 \qquad \Gamma'=\mathbb Z\widehat u\oplus
 d\mathbb Z\widehat\gamma.
\]
Its identification with the pullback cover sends a point to the
endpoint of the lift of a path from a fixed basepoint; changing the
path by $\Gamma'$ preserves that endpoint. Restricting to $q=0$
gives $\widetilde W_p=Y_0/\Gamma'$, which is connected. Since
$B_0(\widehat u)=e_1$ and $B_0(\widehat\gamma)=-e_2$, the vertex
periods are $\mathbb Ze_1+d\mathbb Ze_2$.

The star fan of each divisor has rays
$e_1,e_2,e_2-e_1,-e_1,-e_2,e_1-e_2$. It is obtained by inserting
the sums of adjacent rays of the $\mathbb P^2$ fan
$e_1,e_2-e_1,-e_2$. The chart maps
$(u,v)\mapsto(u,uv)$ and $(u,v)\mapsto(uv,v)$ identify each
insertion with a blowup at the corresponding coordinate point.
Thus each divisor is $dP_6$. No nonzero period fixes its vertex,
so the normalization has exactly $d$ such components. Its map is
the quotient of the disjoint divisor union; locally it separates
the coordinate-plane branches of $z_0z_1z_2=0$, giving the finite
integral-closure map also at the double curves and triple points.

\item We compute the normalization resolution, retaining integral coefficients.
Let $R=\mathbb Z[t]/(t^d-1)$, where $t$ shifts a
component-supported cochain from index $j$ to $j+1$. This is pullback
by the inverse of the positive $e_2$ deck translation; use the same
shift on double-curve cochains. The degree-zero row of the resolution is the cellular
cochain complex of the dual torus with this $d$-sheeted marking; its
cohomology is $(\mathbb Z,\mathbb Z^2,\mathbb Z)$ in degrees $0,1,2$.
The degree-four row is $R$ in horizontal degree zero. In degree two,
write the six boundary curves on $dP_6$ in cyclic order as
\[
 E_1,\quad H-E_1-E_2,\quad E_2,\quad H-E_2-E_3,
 \quad E_3,\quad H-E_1-E_3.
\]
Here $E_1,E_2,E_3$ are the exceptional rays
$-e_1,e_1-e_2,e_2$, respectively. At the vertex $je_2$, orient the
three edge types by $e_1,e_2,e_1-e_2$, and denote them by
$e_j^{\mathrm h},e_j^{\mathrm v},e_j^{\mathrm a}$. Their endpoints have indices $j,j+1,j-1$.
Choose the double-curve edge basis $-e_j^{\mathrm h},-e_{j-1}^{\mathrm v},e_j^{\mathrm a}$.
For comparison, in the exceptional basis $(E_A,E_B,E_C)$ belonging
to the rays $(e_2,-e_1,e_1-e_2)$, the restriction matrix for
$(e_j^{\mathrm h},e_j^{\mathrm v},e_j^{\mathrm a})$ is
\[
 B_f(t)=\begin{pmatrix}
 1&1&1&1\\
 -t^{-1}&-1&-t^{-1}&-t^{-1}\\
 -t&-t&-t&-1
 \end{pmatrix}.
\]
Indeed restriction to a boundary curve is intersection with its
complex-oriented class, with component pairing
$\operatorname{diag}(1,-1,-1,-1)$. The two branches of $e_j^{\mathrm v}$ are
$E_A$ at $j$ and $H-E_B-E_C$ at $j+1$, giving the second row;
the other two rows follow from their two indicated rays in the
same way. The changes of component and edge bases are
\[
 P=\begin{pmatrix}1&0&0&0\\0&0&0&1\\0&1&0&0\\0&0&1&0\end{pmatrix},
 \qquad U=\operatorname{diag}(-1,-t,1),
\]
with inverses $P^{\mathsf t}$ and
$\operatorname{diag}(-1,-t^{-1},1)$. Thus $B=UB_fP$.
The restriction differential from the $d$ normalizations to the $3d$
double curves, in the basis $(H,E_1,E_2,E_3)$, is
\[
 B(t)=
 \begin{pmatrix}
 -1&-1&-1&-1\\
 1&1&1&t\\
 -t&-t&-1&-t
 \end{pmatrix}:R^4\longrightarrow R^3.
\]
Replacing the second row by its sum with the first and the third by
its difference with $t$ times the first gives rows with sole entries
$t-1$ in columns four and three. The first-row unit eliminates the
remaining first-row entries. Thus this matrix is equivalent over $R$
to a matrix with diagonal entries $1,t-1,t-1$ and one zero column.
Multiplication by $t-1$ on $R$ has kernel $\mathbb Z(1+t+\cdots+t^{d-1})$
and cokernel $\mathbb Z$. Therefore $\ker B$ is free abelian of rank
$d+2$, and $\operatorname{coker}B\simeq\mathbb Z^2$.
There can be no higher differential: all rows with odd vertical degree
are zero and the resolution has horizontal degrees only $0,1,2$.
The only extension in degree two is an extension of free abelian
groups, hence splits as an abelian group. This gives cohomology ranks
$(1,2,d+3,2,d)$ with no torsion. The universal coefficient theorem,
applied from the top degree downwards, gives the same assertion for
homology. The normalization resolution used here is exact on stalks:
at a point with $r$ local branches it is the augmented cochain complex
of the $(r-1)$-simplex, and hence computes ordinary cohomology even at
the self-identifications and the triple points.
\end{enumerate}
\end{proof}

\begin{theorem}[Rational homology and deck characters]\hypertarget{tag-02KC}{}\label{maps:src:D.1.3}
{\normalfont\scriptsize tag 02KC}\par


For every admissible $p$, the rational homology of $\widetilde X_p$ has
ranks
\[
 (1,0,d-1,0,d-1,0,1).
\]
The deck group on either $H^2(\widetilde X_p;\mathbb C)$ or
$H^4(\widetilde X_p;\mathbb C)$ is the direct sum of its nontrivial
characters, each with multiplicity one.
\end{theorem}


\begin{proof}
\emph{Character cochains.}
Let $\pi:\widetilde X_p\to X_p$ be the covering, and let $g$
be the deck generator given by the image of $\widehat\gamma$ in
\cref{maps:src:D.1.1}. For $\lambda^d=1$, let
$\mathcal L_\lambda$ be the local system of parallel sections of
the flat line bundle
\[
 (\widetilde X_p\times\mathbb C)/
 \bigl((x,z)\sim(gx,\lambda z)\bigr).
\]
In this convention the parallel transport around the loop whose lift
ends at $gx$ multiplies the coefficient by $\lambda^{-1}$.
A cochain with these coefficients is equivalently a complex-valued
cochain $\varphi$ on $\widetilde X_p$ satisfying
$\varphi(g\sigma)=\lambda\varphi(\sigma)$: its value on a simplex
downstairs is the orbit of $(\sigma,\varphi(\sigma))$ for a chosen
lift, and this orbit is independent of that lift. Conversely the
coefficient lift determines these values on all lifts. The two
constructions commute with restrictions to faces and are inverse
cochain maps
\cite[\S3.H, pp.~328 and 333--334]{Hatcher2002}.
Consequently
\[
 H^*(X_p;\mathcal L_\lambda)
   \simeq H^*(\widetilde X_p;\mathbb C)_{g^*=\lambda}.
\]
The projectors
$d^{-1}\sum_{r=0}^{d-1}\lambda^{-r}(g^*)^r$ give the direct sum
of these subcomplexes for all $d$th roots of unity.

\emph{The regular region and the finite fillings.}
Take $\lambda\ne1$. On the regular fibre the equivariance
character is $\lambda$ on $\widehat\gamma$ and one on
$\widehat u,\widehat w,\widehat\delta$.
The opposite faces of a lifted cubical cell of the torus differ
by the corresponding lattice generator. Its cellular cochain
complex is therefore the exterior algebra on the four coordinate
one-cochains, with differential
$D=(\lambda-1)e_\gamma\wedge(-)$.
The operator
\[
 h=(\lambda-1)^{-1}\iota_\gamma
 \qquad\text{satisfies}\qquad Dh+hD=1.
\]
Thus the fibre cohomology vanishes.

Write $E$ for the bundle over a closed pair of pants, $N_j$ for
the three fillings and $B_j$ for their boundary bundles, taking
small collars when applying Mayer--Vietoris.
The flat real marking lifts a deformation of the pair of pants
onto a graph spine: lift each base path and keep the marked real
vector fixed along that path.
Cut a graph spine of the pair of pants into contractible stars
and intervals. The inverse images of these sets and of every
overlap are products with the regular torus, and have zero
cohomology with the restricted coefficients by the preceding
contraction. Mayer--Vietoris gives the same vanishing for $E$.
Cutting a boundary circle proves it for every $B_j$.

For $j=1,2$, the real marking of the finite filling is
\[
 (D\times(\Lambda_{\mathbb R}/\Lambda))/
 \left\langle(s,x)\longmapsto
 (e^{-2\pi i/m_j}s,A_jx+v_j/m_j)\right\rangle.
\]
The homotopy $(s,x)\mapsto((1-r)s,x)$ commutes with this action
and retracts $N_j$ onto its reduced fibre $S_j$.
The degree-$m_j$ torus cover of $S_j$ induces the fibre-lattice
inclusion on fundamental groups. Its composite with the cyclic
quotient sends $\widehat\gamma$ to $g$ and the other three marked
generators to the identity, by \cref{maps:src:D.1.1}.
It therefore pulls back $\mathcal L_\lambda$ to the system just
computed.
The cochain transfer with these pulled-back coefficients sums
the values on all lifts of a simplex; composed with pullback it
is multiplication by $m_j$. Hence pullback is injective over
$\mathbb C$, and $H^*(S_j;\mathcal L_\lambda|_{S_j})=0$.
This also proves the vanishing for $N_1,N_2$ and for
$U=E\cup N_1\cup N_2$.

\emph{The cusp neighbourhood.}
We construct the comparison between a closed cusp tube
$N_0$ and $W$. On a coordinate chart with
$q=z_1\cdots z_k$, $1\le k\le3$, define the real vector field
on $q\ne0$ by
\[
 a_i=\frac{|z_i|^{-2}}{\sum_{\nu=1}^k|z_\nu|^{-2}},
 \qquad
 \mathcal V_i=-a_i z_i\ (i\le k),\qquad
 \mathcal V_i=0\ (i>k),
\]
and set it equal to zero on $q=0$.
Writing $r_i=|z_i|$ and $r_{\min}=\min r_i$, one has
$|\mathcal V_i|\le r_{\min}$. Differentiation gives
\[
 \frac{\partial a_i}{\partial r_j}
 =-2\delta_{ij}\frac{a_i}{r_i}+2\frac{a_i a_j}{r_j}.
\]
The real derivatives of $\mathcal V$ are bounded
off the divisor: the diagonal terms are bounded, and for $i\ne j$
the cross terms are bounded by a constant times
\[
 \frac{r_i}{r_j}a_i a_j
 \le\frac{(r_i/r_j)^3}{(1+(r_i/r_j)^2)^2}.
\]
Angular derivatives in unit tangent directions are bounded by
$a_i\le1$. A line segment in a coordinate ball either lies in
a coordinate hyperplane or meets the coordinate hyperplanes
in finitely many points. Integrating the derivative bounds on
its remaining subsegments and using continuity proves local
Lipschitz continuity across the divisor.
Moreover
\[
 dq(\mathcal V)=-q,\qquad
 |\mathcal V|\le C r_{\min}\le C|q|^{1/3}
 \quad (|q|\le1).
\]
A smooth partition of unity on finitely many charts covering $W$
preserves the equation $dq(\mathcal V)=-q$ and gives a uniform
bound of this form in a smooth metric.
Choose the compact closed tube inside that neighbourhood.
Local existence, uniqueness and continuous dependence give the
local flow $\Phi_t$
\cite[author's preliminary version, Theorems~2.2 and~2.9,
Corollary~2.15, pp.~38, 44--45, 52--53]{TeschlODE2012}.
Along the flow $q(\Phi_t(x))=e^{-t}q(x)$, so it stays in the
compact tube and the continuation criterion gives the flow for
all $t\ge0$.
The remaining path length after time $t$ is at most
$3C|q(x)|^{1/3}e^{-t/3}$.
Thus $\Phi_t$ converges uniformly to a continuous retraction
$r:N_0\to W$. The homotopy
\[
 H(x,s)=\Phi_{-\log(1-s)}(x)\quad(s<1),\qquad H(x,1)=r(x)
\]
is continuous and fixes $W$ throughout.
Lifting the vector field through $\pi$ gives a deck-invariant
field, and uniqueness makes the lifted homotopy equivariant.
The same length bound in the pullback metric gives its uniform
limit. Thus it induces cochain homotopy equivalences on the
corresponding eigensubcomplexes. In particular restriction
along $W\hookrightarrow N_0$ induces an isomorphism with
$\mathcal L_\lambda$ coefficients.

The cohomology of $U$ and of its overlap $B_0$ with the cusp tube
vanishes. Mayer--Vietoris and the preceding retraction give the
restriction isomorphisms
\[
 H^*(X_p;\mathcal L_\lambda)
 \xrightarrow{\ \simeq\ }H^*(N_0;\mathcal L_\lambda|_{N_0})
 \xrightarrow{\ \simeq\ }H^*(W;\mathcal L_\lambda|_W).
\]
The inverse of the second map is induced by the retraction with
coefficient fibres identified by transport along $H$; the first
is the isomorphism in the Mayer--Vietoris sequence.

\emph{The character calculation on the cusp.}
Use the branch normalization complex and its augmentation
in \cref{maps:src:D.1.2}. Index its covering components by
$j\in\mathbb Z/d$ and the vertices $je_2$.
The period lattice map $\mathsf B_0$ is given by
$\mathsf B_0\widehat\gamma=-e_2$ and
$\mathsf B_0\widehat u=e_1$. Hence $g$ sends a
component of index $j$ to the one of index $j-1$.
Thus $g^*$ sends a component-supported cochain of index $j$ to
index $j+1$. This is the operator $t$ in the preceding restriction
matrix $B(t)$, so its $\lambda$ eigenspace is obtained by
substituting $t=\lambda$.

The degree-zero row is the cellular cochain complex of
$\mathbb R^2/(\mathbb Ze_1+d\mathbb Ze_2)$.
The deck translation $x\mapsto x-e_2$ is homotopic to the identity
through $x\mapsto x-re_2$. Its action on the cohomology of that
row is therefore the identity, and its $\lambda\ne1$
eigensubcomplex is acyclic.
In the degree-two row, the operations
$R_2\mapsto R_2+R_1$ and $R_3\mapsto R_3-tR_1$
on $B(t)$ give
\[
 \begin{pmatrix}
 -1&-1&-1&-1\\
 0&0&0&t-1\\
 0&0&t-1&0
 \end{pmatrix}.
\]
At $t=\lambda\ne1$ this map is onto, with kernel
$\mathbb C(-1,1,0,0)$.
The degree-four row is one copy of the regular representation
in horizontal degree zero, so its $\lambda$ eigenspace has
dimension one. The normalization complex has horizontal degrees
$0,1,2$ and even vertical rows. The differential $d_r$ has
bidegree $(r,1-r)$: for $r=2$ its target is an odd row, and
for $r\ge3$ its target lies outside these columns. Hence every
higher differential vanishes.
It follows that
\[
 H^i(X_p;\mathcal L_\lambda)=
 \begin{cases}\mathbb C,&i=2,4,\\0,&i\ne2,4.\end{cases}
\]

\emph{The trivial character and component classes.}
For $\lambda=1$ the coefficient system is constant.
Admissibility gives $p\ne0$, and
\cite[Theorem~7.22, p.~60]{AC26} gives
$H^i(X_p;\mathbb C)=\mathbb C$ for $i=0,6$ and zero otherwise.
The character decomposition above gives the asserted
homology ranks and exactly one copy of every nontrivial character
on each of $H^2$ and $H^4$. When $d=1$ the character sum is empty.

The normalization maps also describe the degree-four generators.
Let $\eta_j:D_j\to\widetilde W_p$ be its component maps.
Their fundamental classes form the regular basis in
$H_4(\widetilde W_p;\mathbb Q)$.
After extension of scalars to $\mathbb C$, set
$\zeta=e^{2\pi i/d}$. For $1\le k<d$, the classes
\[
 \sum_{j=0}^{d-1}\zeta^{kj}(\eta_j)_*[D_j]
\]
have $g_*$ eigenvalue $\zeta^k$.
The duals of the restriction isomorphisms above show that inclusion
$\widetilde W_p\hookrightarrow\widetilde X_p$ maps each of these
lines isomorphically to the corresponding degree-four line of the
cover. The invariant part of $H_4(\widetilde X_p;\mathbb C)$ is zero
by the character decomposition and the constant-coefficient
calculation. Hence the component
fundamental classes generate $H_4(\widetilde X_p;\mathbb Q)$;
the rational cokernel is zero because it becomes zero over
$\mathbb C$.
\end{proof}

\begin{definition}[The lifted fibre marking]
\hypertarget{tag-Q009}{}\label{mapsIII:def:lifted-fibre}
{\normalfont\scriptsize tag Q009}\par
Put
\[
 \Lambda_d=\mathbb Z\langle d\widehat\gamma,\widehat u,
                   \widehat w,\widehat\delta\rangle,
 \qquad F_d=\Lambda_{\mathbb R}/\Lambda_d.
\]
The coefficient map $L=\operatorname{diag}(d,1,1,1)$ identifies
$\mathbb Z^4$ with $\Lambda_d$; its inverse on that sublattice
divides the first coefficient by $d$. In this basis,
\[
 A_1^{(d)}=\begin{pmatrix}1&0&0&0\\6d&0&1&0\\-6d&-1&-1&0\\-2d&1&0&1\end{pmatrix},
 \qquad
 A_2^{(d)}=\begin{pmatrix}1&0&0&0\\0&0&-1&0\\-6d&1&0&0\\3d&0&1&1\end{pmatrix},
 \qquad A_j^{(d)}=L^{-1}A_jL.
\]
\end{definition}

\begin{theorem}[Integral homology of the universal cover]\hypertarget{tag-02KD}{}\label{maps:src:D.1.4}
{\normalfont\scriptsize tag 02KD}\par


For every admissible triple, with $d=|p|$, one has
\[
 H_i(\widetilde X_p;\mathbb Z)=
 \begin{cases}
 \mathbb Z,&i=0,6,\\
 \mathbb Z^{d-1}\oplus\mathbb Z/d,&i=2,\\
 \mathbb Z/d,&i=3,\\
 \mathbb Z^{d-1},&i=4,\\
 0,&i=1,5.
 \end{cases}
\]
For $d>1$ the cover is not a connected sum of copies of
$S^2\times S^4$.
\end{theorem}


\begin{proof}
Put $e=p/d\in\{1,-1\}$. Within this proof the symbols
$\gamma,u,w,\delta$ denote respectively the homology vectors
$\widehat\gamma,\widehat u,\widehat w,\widehat\delta$; dual
covectors carry a superscript $*$. Admissibility implies that $d$
is odd and prime to three. We use the covering character of
\cref{maps:src:D.1.1}, and the component-chain construction of
\cref{maps:src:D.1.2}.

\begin{enumerate}[label=\arabic*.,leftmargin=*]
\item
Set
\[
 b_1=4\ell_0-\ell_1-\ell_2,\qquad b_2=\ell_1+\ell_2-3\ell_0.
\]
The covering character sends $\gamma$ to $1$, the other fibre vectors
to zero, and the clockwise zero-section peripherals $\rho_j$ to $b_j$
modulo $d$. Thus the covered fibre is
\[
 F_d=\mathbb R^4/\Lambda_d,\qquad
 \Lambda_d=\mathbb Z\langle g,u,w,\delta\rangle,\quad g=d\gamma.
\]
Write $\varepsilon_1=(1,2,-4,0)^t$,
$\varepsilon_2=(1,3,-3,0)^t$ and
$v_j=\ell_j\varepsilon_j+k_j\delta$.
The real marked finite root deck is
$(s,x)\mapsto(e^{-2\pi i/m_j}s,A_jx+v_j/m_j)$.
The local-to-exterior overlap adds
$(\log s)/(2\pi i)\,\Pi_c(s)v_j$; its change under the clockwise
deck is $-\Pi_c(g_js)v_j/m_j$.
Consequently its boundary map sends $\rho_j$ to the affine
generator $r_j$ with $r_j^{m_j}=a^{v_j}$.

In the kernel of the covering character take
$t_j=a^{-b_j\gamma}r_j$. Its affine action on $F_d$ is
$x\mapsto A_jx+c_j$, where $c_j=v_j/m_j-b_j\gamma$.
The identities
\[
 (I+A_1+A_1^2)\gamma=3\varepsilon_1,\qquad
 (I+A_2+A_2^2+A_2^3)\gamma=4\varepsilon_2
\]
and $3b_1-\ell_1=p$, $4b_2-\ell_2=-p$ give
\[
 t_1^3=a^{\nu_1},\quad \nu_1=-p\varepsilon_1+k_1\delta,
 \qquad
 t_2^4=a^{\nu_2},\quad \nu_2=p\varepsilon_2+k_2\delta.
\]
Their first-coordinate shifts in $\Lambda_d$ are $-e/3$ and $e/4$.
Hence their orders are exactly three and four and every nonidentity
power acts freely. The covered reduced finite fibres are
$S'_j=F_d/\langle t_j\rangle$, with quotient maps $Q_j$.
The subgroup generated by $\Lambda_d,t_j$ is the exact covering
kernel: the normal form $a^\lambda r_j^k$ and its zero character
condition express it in these generators. The inverse is inclusion
of these generators in the original affine deck group.

In the real marking, $[(s,x)]\mapsto[((1-r)s,x)]$ is an equivariant
retraction of the finite piece to $S'_j$. Let $U$ be the complement
of a smaller cusp tube in the universal cover. The joining band
between the two finite discs therefore gives the homotopy pushout
$S'_1\leftarrow F_d\to S'_2$ with maps $Q_1,Q_2$.
More precisely, at a fixed overlap point its map is
$Q_j(x-\eta_j)$, where
\[
 \eta_j=\bigl(\Pi_c(s_j)|_{\mathbb R^4}\bigr)^{-1}
 \left(\frac{\log s_j}{2\pi i}\Pi_c(s_j)v_j\right)\in\mathbb R^4.
\]
The homotopy $Q_j(x-r\eta_j)$ gives the comparison of these attaching
maps. A map between their double mapping cylinders follows the
first homotopy backwards, the middle cylinder, and the second
homotopy forwards, on three consecutive subintervals. Reversing
the homotopies gives its inverse up to cancellation of the
backtracking homotopies. This specifies the homotopy equivalence
and the Mayer--Vietoris map $(Q_{1*},-Q_{2*})$.

\item
The images of $A_1-I$ and $A_2-I$ in $\Lambda_d$ are respectively
\[
 \langle-u-w+\delta,\ u-2w\rangle,\qquad
 \langle-u+w,\ -u-w+\delta\rangle.
\]
Their first columns are respectively
$-2d(-u-w+\delta)+4d(u-2w)$ and
$3d(-u-w+\delta)-3d(-u+w)$, so these are the entire integral
images. The quotient coordinates are $(g,w)$, with
$(u,\delta)=(2w,3w)$ and $(w,2w)$ respectively.
The affine power relations eliminate $g$ with a unit coefficient.
Thus $H_1(S'_j)=\mathbb Z\langle t_j,w\rangle$, and
\[
 \begin{array}{c|rrrr}
 &g&u&w&\delta\\\hline
 Q_{1*}&-3e t_1+3e k_1w&2w&w&3w\\
 Q_{2*}&4e t_2-2e k_2w&w&w&2w.
 \end{array}
\]
The map $(Q_{1*},-Q_{2*})$ has kernel
$\mathbb Z b$, $b=u+w-\delta$. Its cokernel is
\[
 H_1(U)=\mathbb Z z_0,\qquad
 t_1=4z_0,\quad t_2=-3z_0,\quad z_0=t_1+t_2,\quad
 g=-12e z_0,\quad u=w=\delta=0.
\]
Indeed the columns $u-w$ and $u-2w$ kill the two $w$ generators;
the remaining relation is $3t_1+4t_2=0$.
The quotient functional is $(4,0,-3,0)$ and its section is
$1\mapsto(1,0,1,0)$.

Let $g^*,u^*,w^*,\delta^*$ denote the dual fibre basis and put
\[
 \alpha_1=g^*\wedge(2u^*+w^*+3\delta^*),\quad
 \alpha_2=g^*\wedge(u^*+w^*+2\delta^*),\quad
 Q=u^*\wedge w^*+6d g^*\wedge\delta^*.
\]
For coefficients $(a,b,c,f,h,k)$ in the order
$(gu,gw,g\delta,uw,u\delta,w\delta)$, the invariant equations are
\[
 \begin{array}{ll}
 j=1:&h=k=0,\ a=2b,\ c=3b+6df,\\
 j=2:&h=k=0,\ a=b,\ c=2b+6df.
 \end{array}
\]
Substitution proves both directions, so the full integral invariant
lattices are $\mathbb Z\alpha_j\oplus\mathbb ZQ$.
The free $H_1$ just computed, Euler number zero, duality and universal
coefficients give $H_2(S'_j)=\mathbb Z^2$ with unimodular intersection
pairing. In the ordered real fibre orientation the two invariant
Gram matrices are
\[
 \begin{pmatrix}0&3\\3&12d\end{pmatrix},\qquad
 \begin{pmatrix}0&2\\2&12d\end{pmatrix}.
\]
Their determinants are $-9$ and $-4$. Give $S'_j$ the quotient
orientation induced by this ordered real fibre orientation. Pullback
is injective by transfer, and its rational image
is the invariant subspace by averaging deck lifts. Degree multiplies
pairings by three or four. The determinant equation therefore gives
integral image indices one and two, respectively.

To identify the second subgroup, use the torus
$(r,z)\mapsto[rc_2+z\delta]\in S'_2$. Its edges $r=0,1$ are identified
by $t_2$, and $z=0,1$ by $\delta$. With the meridian parameter first,
an invariant two-form has integral $\omega(\nu_2,\delta)/4$:
sum the four transforms of $c_2$, using $A_2\delta=\delta$.
For $a\alpha_2+bQ$ this is $(ea+3pb)/2$.
Thus $a+b$ is even. That subgroup has the required index two, proving
\[
 \begin{gathered}
 Q_1^*H^2(S'_1)=\langle\alpha_1,Q\rangle,\\
 Q_2^*H^2(S'_2)=\langle\beta_2,C_2^{\rm coh}\rangle,\\
 \beta_2=2\alpha_2,\ C_2^{\rm coh}=\alpha_2+Q .
 \end{gathered}
\]
Its integral inverse coordinates are $((a-b)/2,b)$.

Let $A,B,C_{\rm fin},D$ be homology bases dual to these cohomology
preimages. Evaluation gives the full fibre map
\[
 M_2=
 \begin{pmatrix}
 2&1&3&0&0&0\\
 0&0&6d&1&0&0\\
 -2&-2&-4&0&0&0\\
 -1&-1&-2-6d&-1&0&0
 \end{pmatrix}.
\]
Its image is exactly $\langle A,B-D,2C_{\rm fin}+D\rangle$:
preimages are the first column minus the second, the fourth
column, and the first minus twice the second; the third column
is $3A+6d(B-D)-2(2C_{\rm fin}+D)$.
Define
\[
 \lambda(aA+bB+cC_{\rm fin}+fD)=2b-c+2f,\qquad h=-C_{\rm fin}.
\]
Then $\lambda(h)=1$ and
$v-\lambda(v)h=aA+b(B-D)+(b+f)(2C_{\rm fin}+D)$.
These are the integral quotient and inverse maps, giving
\begin{equation}\label{maps:eq:finite_h2_extension}
 0\longrightarrow\mathbb Z h\longrightarrow H_2(U)
 \xrightarrow{\partial_{\rm MV}}\mathbb Z b\longrightarrow0.
\end{equation}

\item
For the cusp pair, write $Y=Y_\epsilon$, $Y_0=\{q=0\}$ and
$T=Y\setminus Y_0=(\mathbb C^*)^2\times D^*$.
The lifted pair is
\[
 (N,E)=(Y,T)/\Gamma',\qquad
 \Gamma'=\mathbb Z u\oplus d\mathbb Z\gamma,\quad
 B_0\Gamma'=\mathbb Z e_1\oplus d\mathbb Z e_2 .
\]
This is the covering identification of \cref{maps:src:D.1.2};
its map is given by lifting a path from matching basepoints.

The inclusion $Y_0\to Y$ is an integral homology equivalence.
Cover $Y$ by maximal triangle charts intersected with $|q|<\epsilon$. An intersection with
nonzero common cone of dimension $k$ has coordinates
$\mathbb C^k\times(\mathbb C^*)^{3-k}$, with $q$ the product of the
affine coordinates times an invertible monomial. Scaling the affine
coordinates to zero retracts both the intersection and its central
divisor to $\{0\}\times(\mathbb C^*)^{3-k}$.
An intersection with common cone zero is exactly $T$ and has empty
central divisor. In the relative Čech chain complex the only
nonzero intersection homology terms are therefore copies of $H_*(T)$,
with identity maps. Their horizontal complex is that of the full
simplex on all triangles relative to the nerve $K$ of tuples with
nonempty intersection in $\mathbb R^2$.
The cellular Čech complex of the closed triangles computes
$H_*(K)=H_*(\mathbb R^2)$: over each open cell its augmented row is
the full simplex on its containing triangles, contracted by adjoining
one such triangle, while every nonempty intersection is contractible.
Thus the relative full-simplex complex is acyclic. The same row
contraction after singular subdivision justifies the open-cover
augmentation, and its first-quadrant spectral sequence gives the
claimed homology equivalence using finite-support chains.

The normalization complex from the proof of
\cref{maps:src:D.1.2}, applied on the whole lattice with finite support,
has $H_3(Y_0)$ equal to the kernel of the edge-curve map.
The exact local curve relations are the injective character maps
$R_v(x_v,y_v)=\sum_r\langle(x_v,y_v),r\rangle c_r$.
Matching half-edges forces $x_{i,j}=x_{i+1,j}$ and
$y_{i,j}=y_{i,j+1}$. Finite support forces both to vanish.
The normalization complex has columns zero, one and two, with even
component and edge homology. Its total-degree-three terms in bidegrees
$(0,3)$ and $(2,1)$ vanish, and the remaining term $(1,2)$ has
zero source and target for every higher differential.
Hence $H_3(Y)=H_3(Y_0)=0$.
The three ray meridians $(0,0,1),(1,0,1),(0,1,1)$ bound their
normal discs and generate $\pi_1(T)$; general position makes
$\pi_1(T)\to\pi_1(Y)$ surjective. Thus $H_0(Y,T)=H_1(Y,T)=0$.

Removing double intersections, of real codimension at least four,
changes neither relative two-cycles nor their three-dimensional
bounding chains. Excision and the complex normal-disc Thom map give
$H_2(Y,T)=\bigoplus_{v\in\mathbb Z^2}\mathbb Z n_v$.
Its inverse is signed intersection with the smooth divisor strata.
The period action sends $n_v$ to $n_{v+B_0\lambda}$, so this
is a sum of $d$ regular $\mathbb Z\Gamma'$-modules.
The dense torus lies in the contractible chart of the triangle
$0,e_1,e_2$; simultaneous radial scaling contracts this chart within
$|q|<\epsilon$. Thus $H_2(T)\to H_2(Y)$ is zero and the relative
sequence gives an equivariant isomorphism
\[
 \partial:H_3(Y,T)\xrightarrow{\sim}
 \bigwedge^2\mathbb Z\langle w,\delta,q_{\rm circ}\rangle .
\]
Inverse representatives are the cones on the corresponding tori in
that chart. Here $q_{\rm circ}$ denotes the positive coordinate circle.

On $H_1(T)$ the prescribed period action is
$u:q_{\rm circ}\mapsto q_{\rm circ}+w$, $d\gamma:q_{\rm circ}\mapsto q_{\rm circ}-d\delta$,
fixing $w,\delta$. The factor $\exp(2\pi iC(q)\lambda)$ is homotopic
to one by multiplying its exponent by $r\in[0,1]$.
On the exterior square the coinvariants therefore kill precisely
$w\wedge\delta$, retaining freely $q_{\rm circ}\wedge w,q_{\rm circ}\wedge\delta$.
To pass to the quotient pair, tensor its relative singular chains
with the length-two free resolution for $\mathbb Z^2$.
Write $\mathsf a,\mathsf b$ for the group-ring generators corresponding
to $u,d\gamma$; its differentials are
$(\mathsf a-1,\mathsf b-1)$ and $(1-\mathsf b,\mathsf a-1)^t$.
The chain modules are free on deck orbits. In total degrees two and
three the only surviving terms are the coinvariants just described:
the higher group homology of the regular $H_2(Y,T)$ module vanishes,
and relative degrees zero and one vanish. Projection of the 
cycles consequently gives
\[
 H_2(N,E)=\mathbb Z\langle n_0,\ldots,n_{d-1}\rangle,\qquad
 H_3(N,E)=\mathbb Z\langle c_w,c_\delta\rangle ,
\]
where $n_j$ comes from the ray $(je_2,1)$, and
$\partial c_w=(q_{\rm circ},w)$, $\partial c_\delta=(q_{\rm circ},\delta)$,
with the base-circle parameter first.

\item
The cusp arrow is $G_0T_\sigma$ from exterior to local.
Use local vector basepoint zero and the exterior basepoint
$-\widetilde\sigma(s_*)$, connected to its zero section by the straight
vector path. The lift of the positive $q_{\rm circ}$ loop is
$-\widetilde\sigma(s)$ in exterior coordinates. The defining winding
relation gives its exact based word
\[
 \mu=a^{-\lambda_\sigma}\rho_0^{-1}
     =a^{-\lambda_\sigma}\rho_1\rho_2
     =a^\kappa t_1t_2,\qquad
 \kappa=-\lambda_\sigma+b_1\gamma+b_2A_1\gamma.
\]
The endpoint comparison uses the image of the stated basepoint path
under the deck; its terminal translation is $-\lambda_\sigma$,
as in the displayed based word.
Since $b_1+b_2=\ell_0$, the first coordinate of $\kappa$ is zero.
Hence $[\mu]=z_0$ in $H_1(U)$. In particular the functional
$\Theta=-e z_0^*$ restricts to $12g^*$, and
$\Theta(\mu)=-e$, while $\Theta(\mu^{-1})=e$.

The ray disc $n_j$ is the toric extension of
$z\mapsto(1,z^j,rz)$, $|z|\leq1$, in the coordinates $(x_1,x_2,q)$,
with $0<r<\epsilon$. Its positive boundary is $q_{\rm circ}+j\delta$.
Thus the relative boundary to $H_1(U)$ is augmentation:
$n_j\mapsto z_0$.

The matrix $A_1A_2$ fixes $w$. Applying the rightmost deck first,
the two finite traces are
$w\xrightarrow{A_2}-u+\delta\xrightarrow{A_1}w$.
With path parameter first, the boundary of the $S'_1$ trace is
$w-A_2w=b$. The fibre translations in the overlaps contribute
cycles with zero Mayer--Vietoris boundary. Therefore the 
image $z$ of $(q_{\rm circ},w)$ maps to $b$ in \eqref{maps:eq:finite_h2_extension}.
The map $(r,s)\mapsto rh+sz$ is an integral isomorphism onto
$H_2(U)$: its inverse reads the $b$ coefficient, subtracts that
multiple of $z$, and applies $\lambda$ to a finite-piece preimage.
The displayed identity for $\lambda$ proves independence of the
preimage.

Translation by $r\delta$ gives a circle action
$U\times S^1_\delta\to U$: $A_j\delta=\delta$,
$\Pi_c\delta=(0,1)$, and every overlap translation commutes with it.
The action is free even at the finite central quotients, since every
nonidentity finite deck has nonintegral first-coordinate shift.
Its orbit operation on $H_1(U)$ identifies the class of $(q_{\rm circ},\delta)$
with the sum of the finite orbit tori of $t_1,t_2$, as
$[\mu]=[t_1]+[t_2]$.
For either such torus the integral of an invariant form is
$\omega(\nu_j,\delta)/m_j$, by the square used above.
Their coordinates are consequently
\[
 -e A-2pB,\qquad e C_{\rm fin}+\frac{e+3p}{2}D.
\]
The second coefficient is integral because $d$ is odd. Applying
$\lambda$ to their sum gives $-4p-e+(e+3p)=-p$.
The complete signed attachment is therefore
\[
 H_3(N,E)\longrightarrow H_2(U),\qquad
 (c_w,c_\delta)\longmapsto(z,-ph),\qquad
 [\,\partial\,]_{(h,z)}=
 \begin{pmatrix}0&-p\\1&0\end{pmatrix}.
\]
Its cokernel map is $rh+sz\mapsto[r]\in\mathbb Z/p\mathbb Z$,
with inverse $[1]\mapsto[h]$.
Changes of the section lift alter $\kappa$ by a vector in
$\langle w,\delta\rangle$, whose $H_1(U)$ image is zero; the
circle-orbit operation factors through that same $H_1$ map.

\item
The cusp collar inclusion into $E$ is a homotopy equivalence:
move the base radius to the collar in the real marking, keeping
the marked fibre coordinate fixed. Constant monodromy makes this
homotopy equivariant under the base deck. Excision and the 
relative sequence thus give
\[
 0\longrightarrow\mathbb Z/d\langle[h]\rangle
 \longrightarrow H_2(\widetilde X_p)
 \longrightarrow
 I:=\ker\{\mathbb Z^d\xrightarrow{\sum}\mathbb Z\}
 \longrightarrow0.
\]
The following identity supplies an explicit splitting:
\[
 \delta=(-u-w+\delta)-2(u-2w)-3(-u+w).
\]
The trace cylinders along $t_j$, with linear triangles in the
fibre correcting sums of lattice loops, give chains
$C_j^{\rm tr}(\lambda)$ with boundary $(A_j-I)\lambda$.
Thus $D_\delta=C_1^{\rm tr}(u)-2C_1^{\rm tr}(w)-3C_2^{\rm tr}(u)$ has boundary $\delta$.
The triangle in the $(q_{\rm circ},\delta)$ covering plane with vertices
$0,q_{\rm circ},q_{\rm circ}+j\delta$, subdivided at the integral points of its vertical
edge and oppositely oriented, gives a torus chain $B_j$ with
boundary $\partial n_j-\partial n_0-j\delta$.
The cycles
\[
 s_j=n_j-n_0-B_j-jD_\delta,\qquad 1\leq j<d,
\]
map to the basis $n_j-n_0$ of $I$.
Therefore $([r],m_j)\mapsto r[h]+\sum m_j[s_j]$ is an isomorphism
$\mathbb Z/d\oplus\mathbb Z^{d-1}\to H_2(\widetilde X_p)$.
Its inverse reads the relative coordinates, subtracts those $s_j$,
and reads the remaining $h$ coefficient modulo $p$.

\item
The cover is compact, simply connected, complex oriented, and has
Euler number $2d$ by \cref{maps:src:D.1.1}.
Thus $H_0=H_6=\mathbb Z$ and $H_1=H_5=0$.
Universal coefficients and cap product with its fundamental class
give
$H_4=H^2=\operatorname{Hom}(H_2,\mathbb Z)=\mathbb Z^{d-1}$.
Euler characteristic gives $b_3=2+2(d-1)-2d=0$.
The degree-three universal coefficient sequence therefore identifies
$H^3$ with $\operatorname{Ext}(H_2,\mathbb Z)=\mathbb Z/d$:
apply Hom to $0\to\mathbb Z\xrightarrow{d}\mathbb Z
\to\mathbb Z/d\to0$.
Cap product identifies this group with $H_3$.
For $d=1$ the attachment matrix is unimodular and the displayed
free splitting is empty, so the same computation includes that case.

Finally, K\"unneth and Mayer--Vietoris along the connected-sum sphere
$S^5$ give free $H_2$ and zero $H_3$ for every connected sum of
$S^2\times S^4$ factors. Removing each gluing six-ball leaves
these degrees unchanged by the sequence of $(D^6,S^5)$.
For $d>1$ the class $[h]$ has exact order $d$.
Any homeomorphism to that connected sum would induce a homology
isomorphism sending it to zero in a free group, contradicting the
inverse induced map.
\end{enumerate}
\end{proof}

\subsection{Framings and the quotient homotopy class}
\begin{definition}[The character kernel and stable generators]
\hypertarget{tag-Q010}{}\label{mapsIII:def:framing-conventions}
{\normalfont\scriptsize tag Q010}\par
Let $(\gamma,u,w,\delta)$ be the basis dual to
$(\widehat\gamma,\widehat u,\widehat w,\widehat\delta)$, and put
\[
 L=\ker\gamma=\mathbb Z\langle\widehat u,\widehat w,
                                    \widehat\delta\rangle,
 \qquad T_L=L_{\mathbb R}/L\subset\Lambda_{\mathbb R}/\Lambda.
\]
We write $\pi_3^{\rm S}=\mathbb Z/24\langle\nu\rangle$, with
$e_{\mathbb R}(\nu)=1/24$, and
$\pi_6(S^3)=\mathbb Z/12\langle\nu'\rangle$ in the convention
of \cite[Theorem~4.1 and Corollaries~4.2--4.3, pp.~76--77]{mapsToda52}.
The map $\eta_n:S^{n+1}\to S^n$ is the $(n-2)$-fold suspension
of the complex Hopf map, for $n\ge2$.
We use outward stabilization, with the normal frame before the
tangential frame and a constant orientation-preserving adjustment
of any permutation of those fields.
\end{definition}

\begin{lemma}[Framing on the character kernel]
\hypertarget{tag-02J4}{}\label{maps:src:B.1.1}
{\normalfont\scriptsize tag 02J4}\par
For every admissible member with $|p|=1$, the manifold $X_p$ is
an oriented smooth six-sphere. Let $F$ be a regular fibre.
Compare its bounding stable framing with the frame on $TX_p|_F$
formed by four translation-invariant vertical fields and an
oriented lifted base frame. The restriction of their stable
difference to every coset of $T_L\subset F$ is null-homotopic.
\end{lemma}

\begin{proof}
Put $d=|p|=1$.  The fundamental-group formula
$\pi_1(X_p)=\mathbb Z/d$ and
Theorem~\ref{maps:src:D.1.4} show that $X_p$ is a simply
connected integral homology six-sphere.  Successive applications
of the Hurewicz theorem give $\pi_i(X_p)=0$ for $1\le i\le5$
and an isomorphism $\pi_6(X_p)\to H_6(X_p;\mathbb Z)$.
A map $S^6\to X_p$ representing the positive generator is a
homology equivalence, hence a homotopy equivalence by
\cite[Theorem~4.32 and Corollary~4.33, pp.~366--367]{Hatcher2002}.
The equality $\Theta_6=0$ and the smooth interpretation of this
group in \cite[pp.~504--505]{manKM63} give an oriented
diffeomorphism $\chi:X_p\to S^6$.

The global stable framing is
$\tau_x(v,t)=d\chi_x(v)+t\chi(x)$.
Two oriented stable trivializations differ by a map $S^6\to\mathrm{SO}$.
The equality $\pi_6(\mathrm{SO})=0$
\cite[Corollary to Theorem~II and the table on p.~315]{mapsBott59}
makes its stable homotopy class independent of $\chi$.

The affine return maps have the form $h_j(x)=A_jx+a_j$.
Choose a constant change $C$ of the four vertical fields
so that, together with an oriented lifted base frame and the
positive stabilizing vector, they form an oriented frame.
Let $D:F\to GL^+(7,\mathbb R)$ compare that frame with $\tau$.
The transported-frame homotopy has endpoint
\[
 D_1(x)=D_0(h_jx)\operatorname{diag}(C^{-1}A_jC,I_3).
\]
The constant last factor is connected to the identity.
Polar decomposition gives $g\simeq g\circ h_j$ for
$g:F\to SO(7)$. The isotopy $x\mapsto A_jx+t a_j$
identifies the cohomology action with $A_j^t=T_j^{-1}$.

Use the matrices of Definition~\ref{mapsIII:def:lifted-fibre} at $d=1$.
For a mod-two covector $a\gamma+bu+cw+f\delta$,
$A_1$-invariance gives $b=0$, $f=c$; $A_2$-invariance gives
$f=0$, $c=b$. Conversely those equations make both matrix
differences vanish. Thus the common invariant group is
$(\mathbb Z/2)\gamma$.
The integral fixed-vector equations give, in both directions,
\[
 \begin{aligned}
 \ker(A_1-I)&=\mathbb Z\langle\widehat\gamma+2\widehat u-4\widehat w,
                                      \widehat\delta\rangle,\\
 \ker(A_2-I)&=\mathbb Z\langle\widehat\gamma+3\widehat u-3\widehat w,
                                      \widehat\delta\rangle.
 \end{aligned}
\]
Their intersection is $\mathbb Z\widehat\delta$.
For $\Omega=\gamma\wedge u\wedge w\wedge\delta$, the integral map
$\operatorname{PD}(v)=\iota_v\Omega$ is
\[
 a\widehat\gamma+b\widehat u+c\widehat w+f\widehat\delta
 \longmapsto
 a\,uw\delta-b\,\gamma w\delta+c\,\gamma u\delta-f\,\gamma uw.
\]
Its inverse sends the coefficients $(A,B,C,D)$ in the ordered
basis $(uw\delta,\gamma w\delta,\gamma u\delta,\gamma uw)$
to $A\widehat\gamma-B\widehat u+C\widehat w-D\widehat\delta$;
both composites fix every coefficient. Since $\det A_j=1$,
evaluation on three vectors gives
\[
 (\textstyle\bigwedge^3 A_j^t)\operatorname{PD}(v)
       =\operatorname{PD}(A_j^{-1}v),
 \qquad \operatorname{PD}(\widehat\delta)=-\gamma\wedge u\wedge w.
\]
Indeed both evaluations equal
$\Omega(v,A_jv_1,A_jv_2,A_jv_3)
 =\Omega(A_j^{-1}v,v_1,v_2,v_3)$.
It follows that
\[
 H^1(F;\mathbb Z/2)^{h_1^*,h_2^*}=(\mathbb Z/2)\gamma,
 \qquad
 H^3(F;\mathbb Q)^{h_1^*,h_2^*}=\mathbb Q\,\gamma u w.
\]
For the coset inclusion $j_a(y)=a+y$, one has $j_a^*\gamma=0$.
Thus $g j_a$ lifts to $\ell_a:T_L\to\operatorname{Spin}(7)$.
Let $c:\operatorname{Spin}(7)\to SO(7)$ be the double cover
and let $u_3$ generate its integral third cohomology.
Transfer gives $\beta=\frac12\operatorname{tr}(u_3)$ with
$c^*\beta=u_3\otimes1$, since the deck translation is homotopic
to the identity. Hence
\[
 (\ell_a^*u_3)\otimes1=j_a^*g^*\beta=0.
\]
The group $H^3(T_L;\mathbb Z)=\mathbb Z$ is torsion free, so
this class vanishes integrally. Spin$(7)$ is two-connected;
the lifted map can be made constant on the two-skeleton and
then factors through $T_L/T_L^{(2)}=S^3$.
Its remaining class is zero by the degree-three Hurewicz isomorphism.
The two lifts differ by multiplication by the central element $-1$,
which is homotopic to the identity along a path in Spin$(7)$.
This proves the framing assertion for every coset.
\end{proof}

\begin{theorem}[The quotient homotopy class]\hypertarget{tag-02J5}{}\label{maps:src:B.1.2}
{\normalfont\scriptsize tag 02J5}\par


For the quotient replacement of \SeriesRef{mapsI:src:A.2.1} with $|p|=1$,
\begin{equation}
 [e]=6\nu'=\eta_3\eta_4\eta_5\ne0
       \quad\text{in }\pi_6(S^3)=\mathbb Z/12.
 \hypertarget{tag-01UV}{}\label{maps:eq:B.1.1}
\end{equation}
Its stable class is $12\nu$ and its real Adams invariant
is $1/2\in\mathbb Q/\mathbb Z$.

\end{theorem}


\begin{proof}
\begin{enumerate}[label=\arabic*.,leftmargin=*]
\item By Lemma~\ref{maps:src:B.1.1}, choose an oriented diffeomorphism
$\chi:X_p\to S^6$ and its bounding stable framing $\tau$.
The target of \SeriesRef{mapsI:src:A.1.3} is an oriented three-sphere.
Choose an oriented diffeomorphism $\psi:Q\to S^3$ and write
$e=\psi\circ e_{\rm q}\circ\chi^{-1}$.


Choose a point of $Q$ over a regular base value outside
the cusp disc and the two finite root discs.
Because $e_{\rm q}$ is over $B$, its entire inverse image
is the single coset $a+T_L$.
Order a target frame as the positive real base frame
$(b_1,b_2)$ followed by $\partial_\theta$ in the
positive $\gamma$-direction.  If
$de_{\rm q}(K_i)=b_i+\alpha_i\partial_\theta$
for a lifted base vector, put
$H_i=K_i-\alpha_iV_{\widehat\gamma}$.
The derivative normal frame is
$(H_1,H_2,V_{\widehat\gamma})$.
Complete it by the translation fields
$(V_{\widehat u},V_{\widehat w},V_{\widehat\delta})$,
choosing the sign of the last field to give the prescribed
orientation.  Changing the lifts is a shear homotopy;
reordering these fields is a constant change of basis.
Lemma~\ref{maps:src:B.1.1} identifies the induced stable
tangent framing with the Lie framing of $T^3$.

For its circle factor, use the Hopf map
$(z_1,z_2)\mapsto[z_1:z_2]$.
At the fibre $(e^{i\theta},0)$, the derivative normal frame
in the coordinate $z_2/z_1$ is
$(0,e^{i\theta}),(0,i e^{i\theta})$.
Together with the tangent and radial vectors, its frame
differs from the constant quaternionic frame by
left multiplication by $e^{i\theta}$ and a constant permutation.
The loop of left multiplications extends over the disk by
\[
 (r,\theta)\longmapsto
 L_{\,r\cos\theta+i r\sin\theta+j\sqrt{1-r^2}} .
\]
It therefore gives the translation-framed circle, whose
stable class is $\eta$.
The product tubular collapse is the smash product of the
three circle collapses.  Pontryagin--Thom and the integral
stable relation of \cite[p.~385]{Hatcher2002} give
\begin{equation}
 \operatorname{st}[e]=[T^3,\mathrm{Lie}]
 = [S^1,\mathrm{Lie}]^3=\eta^3=12\nu.
 \hypertarget{tag-01UW}{}\label{maps:eq:B.1.2}
\end{equation}
The constant orientation changes act by additive inversion;
the element $12\nu$ has order two.

\item To compute the real Adams invariant, put the flat
product metric on $\mathbb R^3/\mathbb Z^3$ with the spin
structure determined by its Lie framing.
Its spinors are periodic $\mathbb C^2$-valued functions.
On the Fourier mode $e^{2\pi i k\cdot x}$, $k\in\mathbb Z^3$,
the Dirac operator is the Hermitian matrix
$2\pi\sum k_j\sigma_j$.
The Pauli relations
$\sigma_i\sigma_j+\sigma_j\sigma_i=2\delta_{ij}I$
give eigenvalues $\pm2\pi|k|$ for $k\ne0$ and a
two-dimensional complex kernel for $k=0$.
Thus the eta function is zero in its convergent half-plane
and hence at zero, while $h=2$.
The Levi-Civita connection and the connection induced
by the translation framing coincide, so their transgression
term is zero.  The real formula of
\cite[Theorem~(4.14), p.~423]{mapsAPS2}, with
$\epsilon(1)=1/2$, gives
\begin{equation}
 e_{\mathbb R}(T^3,\mathrm{Lie})
       =\frac12\frac{0+2}{2}=\frac12\pmod{\mathbb Z}.
 \hypertarget{tag-01UX}{}\label{maps:eq:B.1.3}
\end{equation}
In this dimension the real and complex normalizations are
related by $e_{\mathbb C}=2e_{\mathbb R}$,
as in \cite[p.~424]{mapsAPS2}.
Thus the same class has complex invariant zero.

\item For the unstable-to-stable map, take $\nu_4$ to be the
quaternionic Hopf class.  The element $\alpha_3$ in
\cite[Theorem~4.1 and Corollaries~4.2--4.3, pp.~76--77]{mapsToda52}
gives the generator $\nu'$ with
\[
 [\iota_4,\iota_4]=2\nu_4-E\nu'.
\]
Suspension kills the Whitehead square, so
$E^2\nu'=2E\nu_4$.
The group $\pi_8(S^5)=\mathbb Z/24$ is already stable,
and $E\nu_4$ is its generator $\nu$.
Together with $\pi_6(S^3)=\mathbb Z/12$
\cite[Theorem~5.41(b), p.~573]{HatcherSpectralSequences},
this shows that stabilization sends a generator to the
order-twelve element $2\nu$ and is injective.
Consequently the class with stable image $12\nu$ is $6\nu'$.
For $\eta_n=\Sigma^{n-2}\eta_2$, the composite
$\eta_3\eta_4\eta_5$ has stable image $\eta^3=12\nu$;
the same injection identifies it with $6\nu'$.
Changing either cyclic generator multiplies it by a unit,
which fixes the unique order-two element in that cyclic group.
\end{enumerate}
\end{proof}

\subsection{Based smoothing at a regular fibre}

\begin{lemma}[Based smoothing at a regular fibre]\hypertarget{tag-Q255}{}\label{mapsIII:based-regular-fibre-smoothing}
{\normalfont\scriptsize tag Q255}\par
For $n\geq1$, let $X$ be a compact smooth manifold without boundary, let
$f:(X,x_0)\to(S^n,*)$ be continuous, and let $y\ne *$ have nonempty
inverse image $N=f^{-1}(y)$. Suppose $f$ is smooth on an open
neighbourhood of $N$ and its differential has rank $n$ along $N$.
There is a based homotopy from $f$ to a smooth map $f_{\mathrm{sm}}$
which is constant near $x_0$. The homotopy is fixed on a closed
tubular neighbourhood of $N$, has full inverse image $N$ at $y$
at every time, and preserves the differential there. In particular,
the ordered normal framing induced by any fixed target chart at
$y$ is unchanged. The framed inverse image of $f_{\mathrm{sm}}$
therefore represents the based homotopy class of $f$ under the
Pontryagin--Thom construction.
\end{lemma}

\begin{proof}
Choose a target chart $\kappa:O\to B\subset\mathbb R^n$ about $*$,
where $B$ is a convex ball centred at $0$, $\kappa(*)=0$ and
$y\notin O$. Continuity gives an open neighbourhood $U$ of $x_0$
whose closure maps into $O$ and is disjoint from $N$.
Choose a smooth function $\beta:X\to[0,1]$, supported in $U$ and
equal to $1$ near $x_0$. The first homotopy is
\[
 f_s(x)=
 \begin{cases}
 \kappa^{-1}\bigl((1-s\beta(x))\kappa(f(x))\bigr),&x\in U,\\
 f(x),&x\notin U,
 \end{cases}
 \qquad 0\leq s\leq1.
\]
The scalar multiplication preserves $B$. The two formulas agree
near the boundary of the support of $\beta$, so they define a
continuous homotopy. At $x_0$ the chart vector is $0$ at every time.
Every changed value remains in $O$, which excludes $y$. The endpoint
$f_b=f_1$ is constant near $x_0$ and agrees with $f$ on a neighbourhood
of the full inverse image $N$.

By the regular-value theorem, $N$ is a compact smooth submanifold.
Choose a closed tubular neighbourhood $T$ with
$N\subset\operatorname{int}T$, contained in the unchanged smooth
region and disjoint from $\overline U$. Choose also a closed
neighbourhood $K_0$ of $x_0$ inside the region where $f_b=*$.
In the round target metric, compactness gives
\[
 d_0=\min_{x\in X\setminus\operatorname{int}T}d(f_b(x),y)>0.
\]
This bound uses the complement of this same protected neighbourhood
$T$. Choose a smooth cutoff $\chi:X\to[0,1]$ which equals $1$
on a neighbourhood of $T\cup K_0$ and has support in the open set
where $f_b$ is smooth. Regard $f_b$ as a map into
$S^n\subset\mathbb R^{n+1}$. Choose
\[
 0<\eta<\min\{d_0/2,\pi/2\},\qquad
 0<\varepsilon<\min\{1/4,\sin(\eta/2)\}.
\]
Here $\pi$ is the injectivity radius of the round sphere.

There is a smooth $a:X\to\mathbb R^{n+1}$ with
$\|a-f_b\|_\infty<\varepsilon$. To construct it, choose finitely
many coordinate charts $\varphi_j$ and a smooth subordinate partition
of unity $\psi_j$, with compact supports inside the charts. In chart
$j$, choose a smooth cutoff $\zeta_j$ equal to $1$ near
$\varphi_j(\operatorname{supp}\psi_j)$ and supported inside the
coordinate domain. Extend
$b_j=\zeta_j(f_b\circ\varphi_j^{-1})$ by zero to the coordinate
Euclidean space. For a nonnegative smooth kernel $K$ supported in
the unit ball with integral one, put
$K_\delta(u)=\delta^{-\dim X}K(u/\delta)$ and
\[
 a(x)=\sum_j\psi_j(x)(K_{\delta_j}*b_j)(\varphi_j(x)),
\]
where each summand is zero off its chart. Uniform continuity on
the finitely many compact supports permits $\delta_j>0$ so small
that each convolution differs from $f_b$ by less than
$\varepsilon$ on $\operatorname{supp}\psi_j$. The weighted sum
has that same strict bound and is smooth.
Define
\[
 v=\chi f_b+(1-\chi)a,\qquad
 f_{\mathrm{sm}}=\frac{v}{\|v\|}.
\]
The product $\chi f_b$ is smooth because its support lies in the
smooth region of $f_b$; outside that support it is zero.
Thus $v$ is smooth. Moreover
$\|v-f_b\|<\varepsilon$ and $\|v\|>1-\varepsilon>0$, so radial
projection is defined and smooth. On a neighbourhood of
$T\cup K_0$ it equals $f_b$ exactly. The chordal estimate
\[
 \|f_{\mathrm{sm}}-f_b\|
 \leq |1-\|v\||+\|v-f_b\|<2\varepsilon
\]
gives round distance less than $2\arcsin\varepsilon<\eta$.

For $x\in X$, let $w(x)=\exp_{f_b(x)}^{-1}(f_{\mathrm{sm}}(x))$
be the unique tangent vector of length less than $\eta$.
The inverse exists and varies continuously because the distance
is below the injectivity radius. The second homotopy is
\[
 G_s(x)=\exp_{f_b(x)}(s\,w(x)),\qquad 0\leq s\leq1.
\]
It is constant on $T\cup K_0$, since $w=0$ there.
For $x\notin\operatorname{int}T$,
\[
 d(G_s(x),y)\geq d(f_b(x),y)-d(G_s(x),f_b(x))
 >d_0-\eta>d_0/2.
\]
Thus this homotopy acquires no additional inverse image of $y$.
On $T$ the map and its differential remain those of $f$.
Concatenate the two homotopies, using $f_{2s}$ for $0\leq s\leq1/2$
and $G_{2s-1}$ for $1/2\leq s\leq1$. They agree at the seam,
fix $x_0$ throughout, start at $f$ and end at $f_{\mathrm{sm}}$.
The whole inverse image and every ordered inverse-differential
normal frame are unchanged.

Finally the framed tubular collapse for $f_{\mathrm{sm}}$ may be
supported in $T$, disjoint from $x_0$. The inverse-differential
framing and its Pontryagin--Thom interpretation are precisely
\cite[\S7, Theorem~B and Lemma~4, pp.~43, 48--50]{CuspNullMilnor1965}.
In the comparison outside the tube, stereographic interpolation
in $S^n\setminus\{y\}$ keeps $x_0$ fixed: for a chart
$h:S^n\setminus\{y\}\to\mathbb R^n$, the interpolation of two
maps taking value $*$ at $x_0$ is
$h^{-1}((1-s)h(f_0(x))+s h(f_1(x)))$, whose value there is always
$*$. Hence the cited construction has the required based meaning.
\end{proof}
\subsection{The holomorphic map to the base}

\begin{theorem}[Null-homotopy of the base map]
\hypertarget{tag-02J6}{}\label{maps:src:B.2.1}
{\normalfont\scriptsize tag 02J6}\par
For every admissible member of Definition~\ref{mapsIII:def:marked-family}
with $|p|=1$, the holomorphic base map
$f:X_p\to\mathbb P^1$ is null-homotopic.
\end{theorem}

\begin{proof}
Fix such a member and write $X=X_p$. Its two reduced multiple fibres satisfy
$f^*(0)=3S_1$ and $f^*(1)=4S_2$.
Use the base isomorphism
\[
 M[x:y]=[x:y-x],\qquad M^{-1}[a:b]=[a:a+b],
 \qquad \widetilde f=M\circ f.
\]
The two displayed composites are the identity.
For $L=\widetilde f^*\mathcal O_{\mathbb P^1}(1)$,
the coordinate sections have divisors $3S_1$ and $4S_2$.
The divisor identifications are bundle isomorphisms
\[
 \phi_j:\mathcal O_X(S_j)^{\otimes m_j}\longrightarrow L,
 \qquad (m_1,m_2)=(3,4),
\]
taking the powers of the canonical divisor sections to those
coordinate sections.

The smooth sphere $X$ has $H^1(X;\mathbb Z)=H^2(X;\mathbb Z)=0$.
The first Chern class classification
\cite[Proposition~3.10, p.~86]{CuspInvariantHatcherVB}
therefore gives trivializations of $L$ and both divisor line bundles.
Choose smooth nonvanishing frames $\ell,e_1,e_2$.
Such frames follow from continuous ones by smooth approximation:
in a Hermitian norm keep the error below half the positive minimum
norm of the continuous frame on compact $X$.
Define $u_j:X\to\mathbb C^*$ by
$\phi_j(e_j^{\otimes m_j})=u_j\ell$.
The induced map $\pi_1(X)\to\pi_1(\mathbb C^*)=\mathbb Z$
is zero, so the exponential covering has a lift
$b_j:X\to\mathbb C$ with $e^{b_j}=u_j$
\cite[Proposition~1.33, pp.~61--62]{Hatcher2002}.
Local inverse branches of the exponential make the lift smooth.
The compatible frames and their inverse changes are
\[
 e'_j=e^{-b_j/m_j}e_j,\qquad e_j=e^{b_j/m_j}e'_j,
 \qquad \phi_j((e'_j)^{\otimes m_j})=\ell.
\]
Write the canonical divisor sections as $a e'_1$ and $b e'_2$.
They have no common zero, and the two coordinate sections of $L$
give $\widetilde f=[a^3:b^4]$.

Put $r=(|a|^2+|b|^2)^{1/2}>0$ and $G=(a,b)/r:X\to S^3$.
Let $\eta_2:S^3\to\mathbb P^1$ be the Hopf map and set
\begin{equation}
 P_{3,4}(z_1,z_2)=
 \frac{(z_1^3,z_2^4)}{\sqrt{|z_1|^6+|z_2|^8}}.
 \hypertarget{tag-01UY}{}\label{maps:eq:B.2.1}
\end{equation}
The denominator is positive on $S^3$.
The everywhere-defined smooth homotopy
\[
 F_s(x)=[a(x)^3e^{-3s\log r(x)}:
          b(x)^4e^{-4s\log r(x)}],\qquad 0\le s\le1,
\]
has $F_0=\widetilde f$ and $F_1=\eta_2P_{3,4}G$.

For $t=|z_1|^2$, the first squared radius under $P_{3,4}$ is
\[
 R(t)=\frac{t^3}{t^3+(1-t)^4},\qquad
 R'(t)=\frac{t^2(1-t)^3(3+t)}{(t^3+(1-t)^4)^2}>0
 \quad(0<t<1).
\]
Since $R(0)=0$ and $R(1)=1$, a target with both coordinates
nonzero has one source radius and $3\cdot4=12$ phase preimages.
Every local degree is positive, with sign
$\operatorname{sgn}(3\cdot4\,R'(t))$ in the same oriented
radius-and-phase coordinates on source and target. Thus
$\deg P_{3,4}=12$.

For quaternionic multiplication $m:S^3\times S^3\to S^3$
and a generator $\omega\in H^3(S^3;\mathbb Z)$,
the K\"unneth decomposition and the two unit inclusions give
$m^*\omega=\operatorname{pr}_1^*\omega+
\operatorname{pr}_2^*\omega$.
Iteration and pullback along the diagonal show that
$p_{12}(q)=q^{12}$ has degree twelve.
Degree classification
\cite[Corollary~4.25, p.~361]{Hatcher2002}
gives $P_{3,4}\simeq p_{12}$.
On based cube maps to $S^3$, concatenation and pointwise
multiplication have the same unit and satisfy the interchange identity;
hence the latter gives addition in $\pi_6(S^3)$.
Consequently $(p_{12})_*$ is multiplication by twelve and vanishes
on $\pi_6(S^3)=\mathbb Z/12$
\cite[Theorem~5.41(b), p.~573]{HatcherSpectralSequences}.
To pass from $G$ to a based representative, choose $x_0\in X$
and use $G_0(x)=G(x_0)^{-1}G(x)$, with inverse change
$G(x)=G(x_0)G_0(x)$.
A path from $1$ to $G(x_0)^{-1}$ gives a homotopy from $G$ to $G_0$.
It follows that $\eta_2P_{3,4}G$ is null-homotopic.
The homotopy $F_s$ and the displayed inverse $M^{-1}$ prove
that $f$ is null-homotopic.
\end{proof}

\subsection{Suspension homotopies from torus extensions}

\begin{definition}[Torus extension data]
\hypertarget{tag-Q013}{}\label{mapsIII:def:torus-extension-data}
{\normalfont\scriptsize tag Q013}\par
A torus extension datum consists of a compact Hausdorff space $X$,
a compact manifold with boundary $N\subset X$ which is a closed
neighbourhood with $\operatorname{int}_X(N)=N\setminus\partial N$,
and a continuous $r:N\to[0,\epsilon]$ such that
$W=r^{-1}(0)$, $\partial N=r^{-1}(\epsilon)$ and
$U=N\setminus W$ is connected and locally path connected.
It also contains continuous maps
$\alpha:U\to T^d=\mathbb R^d/\mathbb Z^d$ and $h:N\to T^d$
such that $\alpha$ and $h|_U$ induce the same homomorphism on
$\pi_1(U)$, a point $x_0\in W$
with $h(x_0)=0$, and a based continuous map
$c:(T^d,0)\to(Z,z_0)$ to a compact Hausdorff based space.
\end{definition}

\begin{lemma}[Suspension of a torus extension]
\hypertarget{tag-02Q8}{}\label{mapsIII:carrier-extension-lemma}
{\normalfont\scriptsize tag 02Q8}\par
For the data of Definition~\ref{mapsIII:def:torus-extension-data}, set
\[
 G(x)=
 \begin{cases}
 [r(x),c(\alpha(x))],&x\in U,\\
 *,&x\in W\cup(X\setminus\operatorname{int}N),
 \end{cases}
 \qquad
 \Sigma_r Z=
 \frac{[0,\epsilon]\times Z}
 {(\{0,\epsilon\}\times Z)\cup([0,\epsilon]\times\{z_0\})}.
\]
The map $G:X\to\Sigma_rZ$ is based null-homotopic, by a homotopy
stationary on $(X\setminus\operatorname{int}N)\cup\{x_0\}$.
\end{lemma}




\begin{proof}
\begin{enumerate}[label=\arabic*.,leftmargin=*]
\item The difference $\alpha-h|_U:U\to T^d$ induces the zero homomorphism
on fundamental groups. The lifting criterion for the covering
$\mathbb R^d\to T^d$
\cite[Proposition~1.33, pp.~61--62]{Hatcher2002} supplies a continuous function
$L:U\to\mathbb R^d$ with
\[
 \alpha(x)=h(x)+[L(x)].
\]
For $s\in[0,1]$, put
\[
 A_s(x)=h(x)+[(1-s)L(x)],\qquad
 G_s(x)=[r(x),c(A_s(x))]\quad(x\in U),
\]
and set $G_s=*$ on $W$ and outside $\operatorname{int}N$.
These formulas agree on the outer boundary. They are jointly continuous
at $W$ and at $\partial N$, even if $L$ is unbounded toward $W$:
every neighbourhood of the suspension basepoint contains the images of
$[0,\delta)\times Z$ and $(\epsilon-\delta,\epsilon]\times Z$
for some $\delta>0$, by compactness of $Z$. This estimate is uniform
in $s$ and in $A_s(x)$. Thus $G_s$ is a based homotopy from $G$
to the map $x\mapsto[r(x),c(h(x))]$ on $N$.

\item Now use that $h$ is defined on all of $N$. Set
\[
 H_s(x)=[(1-s)r(x)+s\epsilon,c(h(x))]
 \quad(x\in N),
 \qquad H_s(x)=*\quad(x\notin\operatorname{int}N).
\]
This is continuous by the closed pasting lemma. The height on
$\partial N$ is always $\epsilon$, so the formulas agree there.
At $x_0$ the second coordinate is $c(0)=z_0$, so $H_s(x_0)=*$
for every $s$. Finally $H_0=G_1$ and $H_1=*$.
Concatenating the two homotopies proves the assertion.
\end{enumerate}
\end{proof}

\section{Cusp quotient maps and cell attachments}
\subsection{Angle-forgetting and cusp-core quotient maps}

\begin{definition}[Suspension quotients of the cusp]
\hypertarget{tag-Q011}{}\label{mapsIII:def:cusp-suspension-maps}
{\normalfont\scriptsize tag Q011}\par
Use the intrinsic closed cusp neighbourhood $N_0$ in the original
sphere member, or in one of the defined admissible $|p|=1$ members.
Write $q=t_c=re^{i\phi}$, $0\le r\le\epsilon$, and
$q=e^{2\pi i s}$. In Definition~\ref{mapsIII:def:radial-cusp-data} put
\[
 B_0=\begin{pmatrix}0&1\\-1&0\end{pmatrix},\qquad
 C(q)=\begin{pmatrix}6\mu(q)&h(q)\\b(q)-h(q)&\mu(q)\end{pmatrix},
 \qquad Z=sB_0+C(q).
\]
Let $\overline\Lambda=\Lambda/
\mathbb Z\langle\widehat w,\widehat\delta\rangle$ and
$\overline T=\overline\Lambda_{\mathbb R}/\overline\Lambda$.
The ordered basis $([\widehat\gamma],[\widehat u])$ identifies
$\overline T$ with $\mathbb R^2/\mathbb Z^2$ by its two coefficients;
the inverse sends $(a_1,a_2)$ to
$a_1[\widehat\gamma]+a_2[\widehat u]$ modulo the lattice.
For a period lift $\zeta=Z(s)u_{\rm long}+v$, with real
$u_{\rm long},v$, define $\overline\pi=[u_{\rm long}]$.
The real period basis gives the unique coefficients, and period
translations change $u_{\rm long}$ by its integral deck vector.

Give every circle its zero-cell at the identity and use the product
cell structures on $\overline T=T^2$ and $S^1\times\overline T=T^3$.
The map $q_2$ sends $x\in N_0\setminus W$ to
$[r,e^{i\phi},\overline\pi(x)]$ in the unreduced suspension
$\Sigma^u T^3$. It sends $W$ to the lower vertex and the
exterior to the upper vertex. Collapsing $T^{3,(2)}$ and the
basepoint arc gives
\[
 H:X\longrightarrow\Sigma(T^3/T^{3,(2)})=S^4.
\]
Forgetting $e^{i\phi}$ and then collapsing $T^{2,(1)}$ and the
basepoint arc gives $\widehat q_3:X\to\Sigma(T^2/T^{2,(1)})=S^3$.
Forgetting $e^{i\phi}$ and applying the first-coordinate character
$[u_{\rm long}]\mapsto e^{2\pi i(u_{\rm long})_1}$ gives
$\widehat q_\gamma:X\to\Sigma S^1=S^2$.
The reduced maps are the quotients of the corresponding unreduced
suspension maps by their specified basepoint arcs.
\end{definition}

\begin{theorem}[Null-homotopies after forgetting the angle]\hypertarget{tag-02J8}{}\label{maps:src:B.3.2}
{\normalfont\scriptsize tag 02J8}\par


The two angle-forgetting maps of Definition~\ref{mapsIII:def:cusp-suspension-maps} satisfy
$[\widehat q_3]=0$ and $[\widehat q_\gamma]=0$.

\end{theorem}


\begin{proof}
\begin{enumerate}[label=\arabic*.,leftmargin=*]
\item Let $N$ be the closed cusp neighbourhood and let
$p:Y_N\to N$ be its free period cover, with group $\Z^2$.
On its punctured part the real long-period coordinate is
\[
 a(x,q)=\bigl((\log|q|)B_0-2\pi\operatorname{Im}C(q)\bigr)^{-1}
             (\log|x_1|,\log|x_2|)^t.
\]
The matrix is $-2\pi\operatorname{Im}Z$ and is invertible because
the columns of $[Z\mid I_2]$ form a real lattice basis.
The defining period action gives
$a(\Psi_\lambda y)=a(y)+\lambda$, and $[a]=\overline\pi$.

Choose a finite partition of unity $\rho_i$ subordinate to evenly
covered sets $U_i\subset N$, and sections $s_i:U_i\to Y_N$.
Write $y=\Psi_{\lambda_i(y)}s_i(p(y))$ and put
$F(y)=\sum_i\rho_i(p(y))\lambda_i(y)$, extending each summand
by zero off its support. Then $F(\Psi_\lambda y)=F(y)+\lambda$.
For a lift $y_0$ of the chosen $x_0\in W$, the map
\[
 h(p(y))=[F(y)-F(y_0)]
\]
is continuous on $N$ and has $h(x_0)=0$.
The difference $a-F+F(y_0)$ is invariant under every deck
translation, so it descends to a continuous
$\ell:N\setminus W\to\R^2$ with
$\overline\pi=h+[\ell]$.

Choose a radial function $b$ that is zero near $W$ and one on an
outer collar. The formulas
\[
 \pi'=h+[b\ell],\qquad \ell'=(1-b)\ell
\]
define a continuous $\pi'$ on all $N$, by setting $\pi'=h$ near
$W$. They give $\pi'=\overline\pi$ on the collar and
$\overline\pi=\pi'+[\ell']$ on the punctured cusp.


\item Apply Lemma~\ref{mapsIII:carrier-extension-lemma} with
$N=N_0$, $r=|q|$, $\alpha=\overline\pi$, $h=\pi'$ and
$L=\ell'$. The neighbourhood
has its product outer collar, and its punctured part is a torus
bundle with connected fibre over the punctured disc; hence the
frontier, connectedness and local path-connectedness conditions hold.
Also $\pi'(x_0)=0$, and
$\overline\pi=\pi'+[\ell']$ verifies the required equality on
fundamental groups. For the three-sphere target take the based
collapse $c_3:T^2\to T^2/T^{2,(1)}=S^2$. For the two-sphere
target take the first-coordinate character $c_\gamma:T^2\to S^1$.
Their suspended formulas are exactly those of
Definition~\ref{mapsIII:def:cusp-suspension-maps}.
The lemma therefore supplies both based null-homotopies, stationary
at $x_0$ and on the exterior.

Before collapsing the basepoint arc, the same two homotopy formulas
in that lemma give the corresponding unreduced homotopies. Their
continuity at both vertices uses the compact torus and is uniform
in the homotopy parameter; $\ell'=0$ fixes the outer collar in the
first stage. In the sphere coordinates
\[
 [r,v]\longmapsto
 (\sin(\pi r/\epsilon)\beta(v),\cos(\pi r/\epsilon)),
 \qquad\beta(0)=e_1,
\]
the second stage moves $x_0$ along
$(\sin(\pi s)e_1,\cos(\pi s))$.
Postcompose by the rotation
\[
 \begin{pmatrix}\cos\pi s&-\sin\pi s\\
                 \sin\pi s&\cos\pi s\end{pmatrix}
\]
on the first and last coordinate plane, fixing the other coordinates.
Matrix multiplication takes that path and the terminal constant to
the lower vertex. This gives the based unreduced null-homotopy.
In the reduced suspension the same arc is already collapsed.
\end{enumerate}
\end{proof}

\begin{definition}[The cusp-core collapse]
\hypertarget{tag-Q012}{}\label{mapsIII:def:core-collapse}
{\normalfont\scriptsize tag Q012}\par
For the original sphere quotient use $R=\epsilon$ and write
$Q=M\cup_\beta T$, where $T=D_R^2\times S^1$ is the designated
cusp solid torus and $K=\{0\}\times S^1$ is its core.
The boundary identification in retained quotient coordinates is
$\beta(Re^{i\phi},\theta)=(e^{i\phi},\theta)$.
Set $C=M\cup\operatorname{Cone}(\partial M)$, with cone vertex $c_0$.
Put $\theta_\gamma(x)=e^{2\pi i(u_{\rm long}(x))_1}$ on the punctured cusp.
The map $q_1:X\to C$ is the prescribed Seifert quotient on the
exterior. On the punctured cusp it is
$x\mapsto[r(x),\beta(Re^{i\phi(x)},\theta_\gamma(x))]$,
and it sends $W$ to $c_0$. The two formulas agree on the outer boundary.
\end{definition}

\begin{theorem}[The cusp-core collapse class]\hypertarget{tag-02J9}{}\label{maps:src:B.3.3}
{\normalfont\scriptsize tag 02J9}\par


For the map $q_1$ of Definition~\ref{mapsIII:def:core-collapse},
the core cofiber splitting $C\simeq S^3\vee S^2$ constructed below gives
\begin{equation}
 [q_1]=(j_3)_*(6\nu'),
 \hypertarget{tag-01V3}{}\label{maps:eq:B.3.5}
\end{equation}
with no $S^2$ or mixed Whitehead component.

\end{theorem}


\begin{proof}
\begin{enumerate}[label=\arabic*.,leftmargin=*]
\item In the notation of Definition~\ref{mapsIII:def:core-collapse},
define $\kappa:Q\to C$
to be the identity on $M$ and, on $T$, by
\[
 \kappa(re^{i\phi},\theta)
   =[r,\beta(Re^{i\phi},\theta)]\quad(r>0),\qquad
 \kappa(0,\theta)=c_0.
\]
The formula is continuous at the core because all cone heights tend to
zero there. Its only nontrivial fibre is $K$. It induces a continuous
bijection $Q/K\to C$ between compact Hausdorff spaces, hence a homeomorphism.

\item Choose the quotient replacement $e$ supported
inside a smaller cusp disc. For some $0<a<b<R$, its restriction to
$a\le r\le R$ is
$(re^{i\phi},\theta_\gamma(x))\in T$. Thus $q_1=\kappa e$ on this
annulus and on the whole exterior. Both maps send the entire cusp
neighbourhood into the same cone. Since $e$ is over $B$, it sends $W$ into
$K$, so both maps send $W$ to $c_0$.
Let $D_u([r,y])=[(1-u)r,y]$ be the radial cone contraction, and let
$\chi(r)$ be $1$ for $r\le a$, $(b-r)/(b-a)$ for $a\le r\le b$, and
$0$ for $r\ge b$. On the cusp define
\[
 G_s(x)=
 \begin{cases}
 D_{2s\chi(r(x))}(q_1(x)),&0\le s\le\tfrac12,\\
 D_{2(1-s)\chi(r(x))}(\kappa e(x)),&\tfrac12\le s\le1.
 \end{cases}
\]
At $s=1/2$ the two formulas agree: where $\chi=1$ both give the vertex,
and where $\chi<1$ the two original maps agree. Extend $G_s$ by the common
map on the exterior. This is continuous and fixed on the outer collar,
the whole exterior and $W$, with $G_0=q_1$ and $G_1=\kappa e$.

\item Fix the identification $\sigma:Q\to S^3$ used for the class of $e$,
and choose a basepoint $m_*\in\operatorname{int}M$ in the image of an
exterior source basepoint. A stereographic chart
$\tau:Q\setminus\{m_\infty\}\to\mathbb R^3$, with
$m_\infty\in\operatorname{int}M\setminus\{m_*\}$, gives the nullhomotopy
\[
 H(k,s)=\tau^{-1}\big((1-s)\tau(k)+s\tau(m_*)\big)
 \quad(k\in K).
\]
The inclusion $K=S^1\to S^3$
is null-homotopic, and its cofiber is
$S^3\vee\Sigma S^1$
\cite[Examples~0.13--0.15 and Proposition~0.18, pp.~14--16]{Hatcher2002}.
Here the core has the mapping-cylinder neighbourhood $T$, obtained from
the boundary projection $\partial T\to K$, so the quotient is its
homotopy cofiber. The relative form of the cited attachment theorem
specifies the required choice as follows.

Put $E=Q\cup_K CK$ and let $h:E\to C$ equal $\kappa$ on $Q$ and collapse
$CK$ to $c_0$. Write $[\theta,u]\in CK$ with boundary at $u=0$ and vertex
at $u=1$. A homotopy inverse $g:C\to E$ is the identity on $M$ and is
\[
 g([r,\beta(Re^{i\phi},\theta)])=
 \begin{cases}
 [\theta,1-2r/R]\in CK,&0\le r\le R/2,\\
 ((2r-R)e^{i\phi},\theta)\in T,&R/2\le r\le R.
 \end{cases}
\]
Indeed, put $a_s(r)=(1+s)r-sR$. The homotopy from $1_E$ to $gh$ is the
identity on $M$, sends $[\theta,u]\in CK$ to
$[\theta,s+(1-s)u]$, and sends a point of $T$ to
$(a_s(r)e^{i\phi},\theta)$ when $a_s(r)\ge0$ and to
$[\theta,-a_s(r)/R]$ when $a_s(r)\le0$. The formulas agree on the core
and fix the outer boundary. They descend to a homotopy from $1_C$ to
$hg$, with cone height $\max\{a_s(r),0\}$.

The attaching homotopy $H$ gives inverse equivalences
$A:E\to Q\vee S^2$ and $B:Q\vee S^2\to E$, both restricting to the
identity on $Q$, with their inverse homotopies fixed on $Q$.
Define
\[
 \Phi=(\sigma\vee1)Ag:C\longrightarrow S^3\vee S^2,
 \qquad J=hB(\sigma^{-1}\vee1):S^3\vee S^2\longrightarrow C.
\]
These are inverse equivalences. For the sphere inclusions $\iota_3,\iota_2$
in the wedge, they satisfy
\[
 j_3:=J\iota_3=\kappa\sigma^{-1},\qquad
 \Phi\kappa\simeq\iota_3\sigma.
\]
All homotopies fix $m_*$. For another nullhomotopy and its equivalence
$\Phi'$, the transition obeys
$(\Phi'J)\iota_3=\Phi'\kappa\sigma^{-1}\simeq\iota_3$.
Thus the actual collapse specifies the three-sphere inclusion in every
one of these core cofiber splittings.

\item The homotopies above give
$\Phi q_1\simeq\Phi\kappa e\simeq\iota_3\sigma e$.
Theorem~\ref{maps:src:B.1.2} now gives
$[\Phi q_1]=(\iota_3)_*(6\nu')$, and hence
\eqref{maps:eq:B.3.5}. This is an equality in the full homotopy group.
To specify the remaining components, let $p_3,p_2$ be the wedge retractions
and put $K_6=\ker(p_3)_*\cap\ker(p_2)_*$. The direct-sum coordinate map is
\[
 \alpha\longmapsto\big((p_3)_*\alpha,(p_2)_*\alpha,
 \alpha-(\iota_3)_*(p_3)_*\alpha-(\iota_2)_*(p_2)_*\alpha\big),
\]
with inverse $(u,v,k)\mapsto(\iota_3)_*u+(\iota_2)_*v+k$.
For $[\Phi q_1]$ it gives $(6\nu',0,0)$. Naturality of Whitehead products
places every mixed word in $K_6$: each wedge retraction kills one of its
entries, and a Whitehead product with a constant entry extends across
the product cell and is nullhomotopic. The entire $K_6$ coordinate is
zero, as is the $S^2$ coordinate. Finally $C\simeq S^3\vee S^2$ is simply
connected, so change of basepoint has no path ambiguity and the same
conclusion holds for the unbased class.
\end{enumerate}
\end{proof}

\subsection{The regular cusp torus}
\hypertarget{tag-03BF}{}\label{wb:h24:ret:section}

\begin{definition}[Retained-angle cusp data]
\hypertarget{tag-Q021}{}\label{mapsIII:def:retained-angle-cusp-data}
{\normalfont\scriptsize tag Q021}\par
Use the original sphere member or any defined admissible $|p|=1$
member in Definition~\ref{mapsIII:def:cusp-suspension-maps}, with the
fixed based sphere identification of Lemma~\ref{maps:src:B.1.1}
and its outward stable framing. The intrinsic cusp neighbourhood
$N_0$ and its marked quotient are the same in these members.
The overlap arrow is $G_0\circ T_\sigma$, where the prescribed
fibre translation has inverse $T_{-\sigma}$. These arrows preserve
the base parameter and the constant upper suspension value. Every
local map below therefore agrees on their common intrinsic cusp.

Write $t=q$ for the cusp parameter and $h_{\rm per}=h$ for
the period function in that definition. The radial convention
has $B_0=J$ in the same ordered bases. In the real period
decomposition below, write $u=u_{\rm long}$ and $k=v$.
Use the original toric chart of \SeriesRef{mapsI:src:A.1.2}.

The ordered quotient lattice is
\[
 \overline\Lambda=\mathbb Z\langle\overline{\widehat\gamma},
              \overline{\widehat u}\rangle,
 \qquad
 K_\Lambda=\mathbb Z\langle\widehat w,\widehat\delta\rangle,
 \qquad
 \overline T^2=\mathbb R^2/\mathbb Z^2.
\]
In these ordered bases the original period data are
\begin{equation}
 B_0=\begin{pmatrix}0&1\\-1&0\end{pmatrix},\qquad
 Z(s)=sB_0+C(t),\qquad t=e^{2\pi i s},\qquad
 C(t)=\begin{pmatrix}6\mu(t)&h_{\rm per}(t)\\
 b(t)-h_{\rm per}(t)&\mu(t)\end{pmatrix}.
 \hypertarget{tag-03B4}{}\label{wb:h24:ret:eq:original-periods}
\end{equation}
Here \(b=\beta+\tau\), \(s=\tau-h_{\rm per}(t)\), and every entry of
\(C\) is holomorphic at \(t=0\).
Write
\begin{equation}
 R(t)=-2\pi\operatorname{Im}C(t),\qquad
 \ell=\log|t|,\qquad A(t)=\ell B_0+R(t).
 \hypertarget{tag-03AF}{}\label{wb:h24:ret:eq:A}
\end{equation}
Choose a closed cusp disc \(|t|\leq\epsilon\) inside the original
coordinate disc and small enough for the quotient construction.  On a
slightly larger closed disc put \(M_R=\sup\|R(t)\|<\infty\).  We will
choose a regular value with \(0<r_0<\epsilon\) and
\(|\log r_0|>M_R\), with the map $H$ and its cutoff $\epsilon$ fixed.

The original covering model on its dense torus is
\begin{equation}
 x_j=e^{2\pi i\zeta_j},\qquad
 \Psi_n(x,t)=
 \bigl(e^{2\pi i C(t)n}t^{B_0n}x,t\bigr),\qquad n\in\mathbb Z^2.
 \hypertarget{tag-03AS}{}\label{wb:h24:ret:eq:deck}
\end{equation}
For a branch of \(s\), writing
\(\zeta=Z(s)u+k\) with \(u,k\in\mathbb R^2\) gives
\begin{equation}
 \log|x|=A(t)u.
 \hypertarget{tag-03B0}{}\label{wb:h24:ret:eq:logx}
\end{equation}
Indeed, \(\operatorname{Im}s=-\ell/(2\pi)\), so
\(-2\pi\operatorname{Im}(Z(s)u+k)=\ell B_0u+R(t)u\).
Since \(B_0\) is orthogonal,
\begin{equation}
 \bigl|\bigl(\ell B_0+vR(t)\bigr)\xi\bigr|
 \geq (|\ell|-vM_R)|\xi|>0
 \quad(0\leq v\leq1,\ \xi\ne0)
 \hypertarget{tag-03AZ}{}\label{wb:h24:ret:eq:invertibility}
\end{equation}
whenever \(|\ell|>M_R\).  Thus \(A(t)\) is invertible there.  Equation
\eqref{wb:h24:ret:eq:deck} changes \(A(t)^{-1}\log|x|\) by \(n\), so the native
quotient is exactly
\begin{equation}
 \overline\pi_f([x,t])=[A(t)^{-1}\log|x|]\in\mathbb R^2/\mathbb Z^2.
 \hypertarget{tag-03B1}{}\label{wb:h24:ret:eq:native}
\end{equation}
This also follows from the flat marking in \cref{wb:cusp-cone-native-descent}.

Let
\[
 T^3=S^1_\phi\times\overline T^2,
 \qquad S^1_\phi=\mathbb R/(2\pi\mathbb Z),
\]
with coordinate basepoints \(\phi=0\), \(u_1=0\), \(u_2=0\).
Its coordinate two-skeleton is the following specified closed subset:
\begin{equation}
 K^{(2)}=
 \bigl(\{0\}\times\overline T^2\bigr)
 \cup\bigl(S^1_\phi\times\{u_1=0\}\bigr)
 \cup\bigl(S^1_\phi\times\{u_2=0\}\bigr).
 \hypertarget{tag-03B9}{}\label{wb:h24:ret:eq:two-skeleton}
\end{equation}
The product CW structure has one open three-cell, with coordinates
\((\phi,u_1,u_2)\in(0,2\pi)\times(0,1)^2\).  Consequently
\[
 c_3:T^3\longrightarrow T^3/K^{(2)}\cong S^3
\]
is the degree-one quotient, where the orientation of \(S^3\) is fixed
by the displayed coordinate order.  One explicit identification of its
open cell with \(\mathbb R^3\) is
\[
 (\phi,u_1,u_2)\longmapsto
 \bigl(-\cot(\phi/2),-\cot(\pi u_1),-\cot(\pi u_2)\bigr).
\]
The inverse on $\mathbb R^3$ is
\[
 (v_\phi,v_1,v_2)\longmapsto
 \left(\pi+2\arctan v_\phi,
       \tfrac12+\frac1\pi\arctan v_1,
       \tfrac12+\frac1\pi\arctan v_2\right).
\]
The identities $-\cot(\pi/2+\arctan v)=v$ and the stated open
coordinate ranges verify both composites. Each derivative in its
own variable is strictly positive, and approach
to any boundary face leaves every compact subset of \(\mathbb R^3\).
Extending by its point at infinity gives the stated based sphere
identification and verifies its continuity on the collapsed skeleton.

The exact conical target and map are
\begin{align}
 \mathcal A&=
 \frac{[0,\epsilon]\times T^3}
 {\{0\}\times T^3\sim v_0,\quad
  \{\epsilon\}\times T^3\sim v_1},
 \hypertarget{tag-03AH}{}\label{wb:h24:ret:eq:A-target}\\
 q_2(x)&=[r,e^{i\phi},\overline\pi_f(x)]
       &&(t(f(x))=re^{i\phi},\ 0<r<\epsilon),\notag\\
 q_2(W)&=v_0,\qquad
 q_2\bigl(X\setminus\operatorname{int}N_0\bigr)=v_1.
 \hypertarget{tag-03B5}{}\label{wb:h24:ret:eq:q2}
\end{align}
Fix the same source basepoint \(x_{\rm base}\in W\).  Define
\begin{equation}
 p_3:\mathcal A\longrightarrow\Sigma_r(S^3,*),\qquad
 p_3([r,z])=[r,c_3(z)],\qquad
 H=p_3q_2:(X,x_{\rm base})\longrightarrow(S^4,*).
 \hypertarget{tag-03AN}{}\label{wb:h24:ret:eq:H}
\end{equation}
The reduced suspension here collapses both endpoint copies of \(S^3\)
and the entire interval over \(*\) to one point.  Thus \(p_3\) sends
both \(v_0\) and \(v_1\) to \(*\) and is well-defined and based.
It is continuous by the defining quotient topology.  An explicit
identification of its open four-cell with \(\mathbb R^4\) is
\begin{equation}
 (r,\phi,u_1,u_2)\longmapsto
 \bigl(-\cot(\pi r/\epsilon),-\cot(\phi/2),
       -\cot(\pi u_1),-\cot(\pi u_2)\bigr).
 \hypertarget{tag-03AW}{}\label{wb:h24:ret:eq:four-cell}
\end{equation}
Its inverse sends $(v_0,v_\phi,v_1,v_2)$ to
\[
 r=\epsilon\left(\tfrac12+\frac1\pi\arctan v_0\right),\quad
 \phi=\pi+2\arctan v_\phi,\quad
 u_i=\tfrac12+\frac1\pi\arctan v_i\quad(i=1,2).
\]
The same cotangent identities verify both inverse compositions.
This fixes the target orientation. Replacing \(r\) locally by
\(\ell=\log r\) preserves this orientation, since \(dr/d\ell=r>0\).
The ordered coordinates in \eqref{wb:h24:ret:eq:four-cell} specify the sphere
identification used throughout this calculation.

Choose the actual lift
\begin{equation}
 u_0=(-1/4,1/4)^{\mathsf T},\qquad
 p_0=B_0u_0=(1/4,1/4)^{\mathsf T},\qquad \phi_0=\pi.
 \hypertarget{tag-03AR}{}\label{wb:h24:ret:eq:chosen-value}
\end{equation}
Modulo \(\mathbb Z^2\), \(u_0\) is \((3/4,1/4)\) in the open cell.  Put
\(t_0=r_0e^{i\phi_0}\) and
\(y_0=[r_0,c_3(e^{i\phi_0},[u_0])]\).  The numbers \(r_0\) and
\(t_0\) will remain nonzero throughout the framing calculation.

The original triangle \(\{0,e_1,e_2\}\) has toric chart
\begin{equation}
 \mathcal U_\sigma\cong\mathbb C^3,\qquad
 (z_0,z_1,z_2)=(t/(x_1x_2),x_1,x_2),\qquad
 t=z_0z_1z_2.
 \hypertarget{tag-03BB}{}\label{wb:h24:ret:eq:z-chart}
\end{equation}
These characters are dual to \((0,0,1),(1,0,1),(0,1,1)\), respectively,
as direct pairing verifies.  This is the chart of
\SeriesRef{mapsI:src:A.1.2}, with its original coordinates.

For \(j=1,2\), introduce the nowhere vanishing holomorphic functions
\begin{equation}
 h_j(t)=\exp\bigl(-2\pi i(C(t)u_0)_j\bigr).
 \hypertarget{tag-03BA}{}\label{wb:h24:ret:eq:units}
\end{equation}
The following is an exact coordinate map, with the exact inverse
displayed alongside it:
\begin{align}
 F(z)&=(w_0,w_1,w_2)
 =\bigl(z_0/(h_1(t)h_2(t)),\ h_1(t)z_1,\ h_2(t)z_2\bigr),
 \hypertarget{tag-03AL}{}\label{wb:h24:ret:eq:F}\\
 F^{-1}(w)&=
 \bigl(w_0h_1(t)h_2(t),\ w_1/h_1(t),\ w_2/h_2(t)\bigr),
 \quad t=w_0w_1w_2.
 \hypertarget{tag-03AM}{}\label{wb:h24:ret:eq:Finv}
\end{align}
Both maps are holomorphic wherever \(|t|\) is in the period disc,
both preserve \(t\), and substitution proves that both composites are
the identity.  In particular the maps extend across the origin and
across every coordinate divisor. Both coordinate descriptions use
the same base parameter $t$ and the same holomorphic period germ $C$.

\end{definition}

\begin{lemma}[A cusp torus in one coordinate chart]
\hypertarget{tag-03AE}{}\label{wb:h24:ret:chart-lemma}
{\normalfont\scriptsize tag 03AE}\par
For sufficiently small \(r_0>0\), the whole inverse image
\(M=H^{-1}(y_0)\) is contained in one contractible coordinate neighborhood
\(U\subset X\) that lifts isomorphically to an open neighborhood of the
toric origin.
In the coordinates \(w=F(z)\) it is the embedded two-torus
\begin{equation}
 w(a,b)=
 \bigl(\rho_0e^{i(\phi_0-a-b)},\rho_1e^{ia},\rho_2e^{ib}\bigr),
 \qquad
 \rho_0=r_0^{1/2},\quad \rho_1=\rho_2=r_0^{1/4},
 \quad (a,b)\in\bigl(\mathbb R/(2\pi\mathbb Z)\bigr)^2.
 \hypertarget{tag-03AP}{}\label{wb:h24:ret:eq:M}
 \hypertarget{tag-01V0}{}
\end{equation}
\SsixDisplayAliases{wb:h24:ret:eq:M}{\label{maps:eq:B.3.2}}
\end{lemma}

\begin{proof}
Let \(r_C>0\) be a radius on which the period correction \(C(t)\) is
holomorphic, and let \(\chi:Y_\epsilon\to N_0\) be the toric quotient.
The functions \(h_1,h_2\) in \eqref{wb:h24:ret:eq:units} are holomorphic
and nowhere zero on that disc.  The map \(F\) preserves the coordinate
product \(t\).  Denote the map in \eqref{wb:h24:ret:eq:Finv} by \(G\);
substitution gives
\(G F=\operatorname{id}\) and \(F G=\operatorname{id}\).
Thus \(F\) and the displayed map \(G\) are holomorphic inverses on
\(\{|z_0z_1z_2|<a_0\}\) for every
\(0<a_0<\min(r_C,\epsilon)\), including the three central coordinate planes.

We construct a common injective neighborhood.  Use the maximum row-sum
matrix norm and set
\[
 a_0=\tfrac12\min(r_C,\epsilon,1),\qquad
 K=\max_{|t|\leq a_0}\|R(t)\|_\infty,\qquad
 L=1+2K-\tfrac13\log a_0,\qquad \delta=e^{-L}.
\]
Then \(L>2K\), \(\delta<1\), and
\(\delta^3=a_0e^{-3-6K}<a_0\).  Put
\[
 W_\delta=\{w\in\mathbb C^3:|w_i|<\delta\ (0\leq i\leq2)\},
 \qquad V=G(W_\delta).
\]
For a point of \(W_\delta\) with \(t\ne0\), write
\[
 S=-\log|t|>3L,\qquad
 p_i=\frac{\log|w_i|}{\log|t|}\quad(i=1,2),\qquad \eta=\frac{L}{S}.
\]
The three polydisc inequalities give
\[
 p_1>\eta,\qquad p_2>\eta,\qquad p_1+p_2<1-\eta.
\]
Since \(B_0\) and \(B_0^{-1}\) preserve the maximum norm and
\(\|R(t)/\log|t|\|_\infty<1/2\), the geometric series for the inverse
of \(I+B_0^{-1}R(t)/\log|t|\) gives
\(\|(B_0+R(t)/\log|t|)^{-1}\|_\infty<2\).
Consequently
\[
 s=A(t)^{-1}\binom{\log|w_1|}{\log|w_2|}
   =\left(B_0+\frac{R(t)}{\log|t|}\right)^{-1}p,
 \qquad
 \left\|s-\binom{-p_2}{p_1}\right\|_\infty
 \leq\frac{2K}{S}<\eta.
\]
It follows that \(-1<s_1<0\) and \(0<s_2<1\).

In the \(w\)-coordinates the deck map \(\Psi_n\) multiplies \(w_j\),
for \(j=1,2\), by the same factor
\(e^{2\pi i(C(t)n)_j}t^{(B_0n)_j}\) as it multiplies \(x_j\).
Its logarithmic absolute-value increment is exactly \(A(t)n\), so it
sends \(s\) to \(s+n\).  Two points of the open rectangle
\((-1,0)\times(0,1)\) have an integral difference only when they agree.
Thus \(V\cap\Psi_n(V)\) is empty over \(t\ne0\) for every \(n\ne0\).
It is empty also at \(t=0\): a nonempty intersection would be open in
the toric chart, hence would contain a point outside the union of its
three central coordinate planes, contradicting the preceding conclusion.

Therefore \(\chi|_V\) is a holomorphic coordinate isomorphism onto the
open set \(U=\chi(V)\subset X\).  Its inverse selects the unique
representative in \(V\), and \(F(\chi|_V)^{-1}\) identifies \(U\) with
\(W_\delta\).  The contraction is \(w\mapsto(1-\tau)w\),
\(0\leq\tau\leq1\); in the toric chart it is
\(z\mapsto G((1-\tau)F(z))\), with \(G\) evaluated at the new product.

We now identify the whole inverse image.  The point
\((e^{i\pi},[u_0])\) has cube coordinates \((1/2,3/4,1/4)\) in the
complement of \(K^{(2)}\).  The collapse \(c_3\) is injective there,
and the suspension collapse is injective at its interior nonbasepoints.
Equations \eqref{wb:h24:ret:eq:H} and \eqref{wb:h24:ret:eq:q2} therefore give
\[
 M=\{x:t(f(x))=-r_0,\ \overline\pi_f(x)=[u_0]\}.
\]
The exterior of the punctured cusp maps to the suspension basepoint and
contributes no point of this inverse image.

Take \(0<r_0<\delta^4\), put \(t_0=-r_0\), and choose
\(s_0=1/2+iT\), where \(T=-\log(r_0)/(2\pi)\).
Since \(\operatorname{Im}Z(s_0)=-A(t_0)/(2\pi)\) is invertible, the
marking \(\zeta=v+Z(s_0)u\) has inverse
\[
 u=(\operatorname{Im}Z(s_0))^{-1}\operatorname{Im}\zeta,
 \qquad v=\operatorname{Re}\zeta-(\operatorname{Re}Z(s_0))u.
\]
For each point of \(M\), applying the deck map associated with the
integer difference \(u_0-u\) gives the representative with \(u=u_0\).
Its remaining coordinate is \(v\bmod\mathbb Z^2\).  Since
\(B_0u_0=(1/4,1/4)^{\mathsf T}\), its exponential coordinates satisfy
\[
 x_j=e^{2\pi i v_j}e^{2\pi i s_0/4}
                  e^{2\pi i(C(t_0)u_0)_j},\qquad
 w_j=h_j(t_0)x_j
     =r_0^{1/4}e^{i(2\pi v_j+\pi/4)}\quad(j=1,2).
\]
Set \(a=2\pi v_1+\pi/4\) and \(b=2\pi v_2+\pi/4\).
Their inverse formulas are
\(v_1=(a-\pi/4)/(2\pi)\) and \(v_2=(b-\pi/4)/(2\pi)\), modulo one.
The product \(w_0w_1w_2=t_0\) gives exactly
\eqref{wb:h24:ret:eq:M}.  Conversely every pair of phases gives the
required values of \(t\) and \([u]\).  All three radii are less than
\(\delta\), so the entire torus lies in \(W_\delta\) and descends
injectively to \(U\).

The derivative columns of this parametrization are
\((-iw_0,iw_1,0)\) and \((-iw_0,0,iw_2)\); their second and third
coordinates prove real linear independence.  The inverse on the image
is \((w_1/r_0^{1/4},w_2/r_0^{1/4})\in S^1\times S^1\).
Thus the parametrization is the asserted embedding.
\end{proof}

\subsection{Differential and orientation of the cusp torus}

\begin{lemma}[The differential normal framing]
\hypertarget{tag-Q022}{}\label{mapsIII:lem:retained-angle-differential}
{\normalfont\scriptsize tag Q022}\par
For the torus of Lemma~\ref{wb:h24:ret:chart-lemma}, the map $H$
is a submersion along its full inverse image. Its target-coordinate
normal framing is given by the vectors in
\eqref{wb:h24:ret:eq:normal-first} and
\eqref{wb:h24:ret:eq:normal-u}. The literal four-cell framing and
this framing are joined by the positive diagonal path below, and
\eqref{wb:h24:ret:eq:A-homotopy} gives a normal-frame homotopy
on the same embedded torus for every parameter $0\leq v\leq1$.
\end{lemma}

\begin{proof}

We calculate in a neighborhood of the compact torus \eqref{wb:h24:ret:eq:M}.
The quotient coordinates admit a unique local lift near \(u_0\).
Set
\[
 v=(\log|w_1|,\log|w_2|)^{\mathsf T},\quad
 \ell=\log|w_0w_1w_2|,\quad
 \phi=\arg(w_0w_1w_2)
\]
using the branch of \(\phi\) near \(\phi_0=\pi\).
Equations \eqref{wb:h24:ret:eq:logx} and \eqref{wb:h24:ret:eq:F} give the exact identity
\begin{equation}
 v=\ell B_0u+R(t)(u-u_0),\qquad
 u=u_0+A(t)^{-1}(v-\ell p_0).
 \hypertarget{tag-03AV}{}\label{wb:h24:ret:eq:exact-u}
\end{equation}
Differentiate this identity before restricting to the inverse image:
\begin{equation}
 du=-A(t)^{-1}(dA(t))A(t)^{-1}(v-\ell p_0)
        +A(t)^{-1}(dv-p_0\,d\ell).
 \hypertarget{tag-03AX}{}\label{wb:h24:ret:eq:full-derivative}
\end{equation}
On \(M\) the vector \(v-\ell p_0\) is exactly zero.  Therefore
\begin{equation}
 du|_M=A_0^{-1}(dv-p_0\,d\ell),\qquad A_0=A(t_0).
 \hypertarget{tag-03AT}{}\label{wb:h24:ret:eq:derivative-on-M}
 \hypertarget{tag-01V1}{}
\end{equation}
\SsixDisplayAliases{wb:h24:ret:eq:derivative-on-M}{\label{maps:eq:B.3.3}}
The equality follows from the vanishing factor $v-\ell p_0$
in the full derivative, including the entire $dA(t)$ term.

Identify \(T_w\mathbb C^3\) with \(\mathbb C^3\) in the given coordinate chart,
and define real tangent vectors
\begin{equation}
 L_j=(0,\ldots,w_j,\ldots,0),\qquad
 Q_j=(0,\ldots,iw_j,\ldots,0)\qquad(0\leq j\leq2).
 \hypertarget{tag-03AO}{}\label{wb:h24:ret:eq:LQ}
\end{equation}
Thus \(d\log|w_i|(L_j)=\delta_{ij}\),
\(d\arg w_i(Q_j)=\delta_{ij}\), and the crossed evaluations vanish.
The ordered tangent vectors of the parametrization \eqref{wb:h24:ret:eq:M} are
\begin{equation}
 T_a=-Q_0+Q_1,\qquad T_b=-Q_0+Q_2.
 \hypertarget{tag-03B8}{}\label{wb:h24:ret:eq:tangents}
\end{equation}
They are independent because their \(w_1\) and \(w_2\) components,
respectively, are nonzero.  For the target order
\((\ell,\phi,u_1,u_2)\), define the following four lifts of its
coordinate vectors:
\begin{align}
 N_\ell&=\tfrac12L_0+\tfrac14L_1+\tfrac14L_2,
       &N_\phi&=Q_0,\hypertarget{tag-03B2}{}\label{wb:h24:ret:eq:normal-first}\\
 N_i&=-\bigl((A_0)_{1i}+(A_0)_{2i}\bigr)L_0
           +(A_0)_{1i}L_1+(A_0)_{2i}L_2
       &&(i=1,2).\hypertarget{tag-03B3}{}\label{wb:h24:ret:eq:normal-u}
 \hypertarget{tag-01V2}{}
\end{align}
\SsixDisplayAliases{wb:h24:ret:eq:normal-u}{\label{maps:eq:B.3.4}}
Let \(P=(\ell,\phi,u_1,u_2)\).  Evaluation using
\eqref{wb:h24:ret:eq:derivative-on-M} gives the entire differential table
\begin{equation}
\begin{array}{c|rrrrrr}
 &N_\ell&N_\phi&N_1&N_2&T_a&T_b\\\hline
 d\ell&1&0&0&0&0&0\\
 d\phi&0&1&0&0&0&0\\
 du_1&0&0&1&0&0&0\\
 du_2&0&0&0&1&0&0
\end{array}.
 \hypertarget{tag-03AU}{}\label{wb:h24:ret:eq:differential-table}
\end{equation}
In particular \(P\), and therefore \(H\) in its open target cell,
is a submersion along \(M\).  The classes of
\((N_\ell,N_\phi,N_1,N_2)\) in \(TX|_M/TM\) give the specified
normal four-framing.  Projecting the displayed lifts along
\(TM\) to any chosen orthogonal complement preserves the quotient
frame.

If the normal frame is specified using the literal four real coordinates
in \eqref{wb:h24:ret:eq:four-cell}, their differential at the chosen value is
\[
 dH_{\rm cell}=K\,dP,\qquad
 K=\operatorname{diag}\!\left(
 \frac{\pi r_0}{\epsilon}\csc^2\frac{\pi r_0}{\epsilon},
 \frac12,2\pi,2\pi\right).
\]
The first factor includes \(dr=r_0\,d\ell\), the second uses
\(\phi_0=\pi\), and the final two use the actual representatives
\((3/4,1/4)\) of \(u_0\). Thus, writing
\(\mathcal N=(N_\ell,N_\phi,N_1,N_2)\), the literal dual frame is
\(\mathcal N K^{-1}\). All four entries of \(K\) are positive
constants on \(M\). The explicitly positive diagonal matrices
\((1-v)K^{-1}+vI_4\), \(0\leq v\leq1\), give a homotopy
from that literal frame to \(\mathcal N\). This retains the original
target coordinates and establishes that either coordinate description
gives the same framed bordism class.

The period matrix gives an explicit homotopy of this normal framing.
Put $\ell_0=\log r_0$ and
\begin{equation}
 A_{0,v}=\ell_0B_0+vR(t_0),\qquad 0\leq v\leq1,
 \hypertarget{tag-03AG}{}\label{wb:h24:ret:eq:A-homotopy}
\end{equation}
and replace \(A_0\) by \(A_{0,v}\) only in
\eqref{wb:h24:ret:eq:normal-u}.  Equation \eqref{wb:h24:ret:eq:invertibility} proves
invertibility for every \(v\).  Their images under the actual
\((du_1,du_2)\) are the columns of \(A_0^{-1}A_{0,v}\), so these
are normal frames throughout the homotopy.  At every time the embedding, the period germ and the target map
remain the specified ones; only the displayed normal-frame matrix varies.

\end{proof}

\begin{lemma}[Orientation of the cusp torus]
\hypertarget{tag-03BD}{}\label{wb:h24:ret:orientation-lemma}
{\normalfont\scriptsize tag 03BD}\par
Take \(r_0>0\) sufficiently small in the chart of
\cref{wb:h24:ret:chart-lemma}, and use the ordered normal frame
\((N_\ell,N_\phi,N_1,N_2)\) defined in
\eqref{wb:h24:ret:eq:normal-first} and
\eqref{wb:h24:ret:eq:normal-u}.
In the complex real-coordinate orientation
\[
 (\operatorname{Re}w_0,\operatorname{Im}w_0,
  \operatorname{Re}w_1,\operatorname{Im}w_1,
  \operatorname{Re}w_2,\operatorname{Im}w_2)
\]
one has
\begin{equation}
 \det(N_\ell,N_\phi,N_1,N_2,T_a,T_b)
       =-\det(A_0)|w_0w_1w_2|^2<0.
 \hypertarget{tag-03B7}{}\label{wb:h24:ret:eq:sign}
\end{equation}
With the normal frame placed before the tangent frame, the induced
tangent orientation is the orientation of \((-T_a,T_b)\).
\end{lemma}

\begin{proof}
On a closed disc inside the common domain of the period germs, put
\(K=\max\|R(t)\|_\infty\).  Shrink the chart's radial range so that
\(\lvert\ell_0\rvert>2K+1\), where \(\ell_0=\log r_0\).
For \(0\leq v\leq1\), write
\[
 A_{0,v}=\ell_0B_0(I+E_v),\qquad
 E_v=\frac v{\ell_0}B_0^{-1}R(t_0),\qquad B_0^{-1}=-B_0.
\]
The bound \(\|E_v\|_\infty<1/2\) gives the continuous inverse
\[
 A_{0,v}^{-1}
   =\left(\sum_{j=0}^{\infty}(-E_v)^j\right)
      \frac{B_0^{-1}}{\ell_0}.
\]
The geometric series identity proves both inverse compositions.
Thus \(\det A_{0,v}\) never vanishes, and its value at \(v=0\)
is \(\ell_0^2>0\).  In particular \(\det A_0>0\).

Write \(w_j=x_j+iy_j\).  The maps from the coordinate pair to
the frame \((L_j,Q_j)\), and their inverses, are
\[
 D_j=\begin{pmatrix}x_j&-y_j\\y_j&x_j\end{pmatrix},
 \qquad
 D_j^{-1}=\frac1{|w_j|^2}
            \begin{pmatrix}x_j&y_j\\-y_j&x_j\end{pmatrix}.
\]
Consequently \(D=\operatorname{diag}(D_0,D_1,D_2)\) has determinant
\(\prod_j|w_j|^2>0\).  Write
\(A_0=\left(\begin{smallmatrix}a&b\\c&d\end{smallmatrix}\right)\)
in the following calculation.  In the ordered frame
\((L_0,Q_0,L_1,Q_1,L_2,Q_2)\), the six columns have coefficient matrix
\[
 H=\begin{pmatrix}
 1/2&0&-a-c&-b-d&0&0\\
 0&1&0&0&-1&-1\\
 1/4&0&a&b&0&0\\
 0&0&0&0&1&0\\
 1/4&0&c&d&0&0\\
 0&0&0&0&0&1
 \end{pmatrix}.
\]
Permute its rows into order \((1,3,5,2,4,6)\) and its columns
into order \((1,3,4,2,5,6)\).  The inverse maps are the transpose
permutation matrices; their signs are respectively \(-1\) and \(+1\).
The resulting matrix is \(\operatorname{diag}(U,V)\), where
\[
 U=\begin{pmatrix}1/2&-a-c&-b-d\\1/4&a&b\\1/4&c&d\end{pmatrix},
 \qquad
 V=\begin{pmatrix}1&-1&-1\\0&1&0\\0&0&1\end{pmatrix}.
\]
The determinant-one row map \(V^{-1}\), which adds the last two
rows to the first, gives \(V^{-1}U\) first row \((1,0,0)\) and
lower-right block \(A_0\); its inverse subtracts those rows.
Hence \(\det U=\det A_0\), \(\det V=1\), and
\(\det H=-\det A_0\).  The actual ambient matrix is \(DH\), proving
the displayed equality and inequality.

The derivatives of the torus parametrization are precisely
\(T_a=-Q_0+Q_1\) and \(T_b=-Q_0+Q_2\).
The nonzero six-frame determinant proves that the first four
columns project to a frame of the normal quotient.  The inverse
of that projection takes the first four coordinates under
\((DH)^{-1}\), which are unchanged on adding either tangent column.
Finally the tangent map \(\operatorname{diag}(-1,1)\), inverse
itself, sends \((T_a,T_b)\) to \((-T_a,T_b)\).  It changes the
concatenated determinant to \(\det A_0|w_0w_1w_2|^2>0\), giving
the asserted induced orientation.
\end{proof}

\subsection{Stable framing and the spin invariant}

\begin{proposition}[Stable framing of the cusp torus]
\hypertarget{tag-03BC}{}\label{wb:h24:ret:framing-proposition}
{\normalfont\scriptsize tag 03BC}\par
Use the complex real-coordinate orientation and place the ordered normal
four-frame before the tangent frame of \(M=H^{-1}(y_0)\).
The induced stable tangential framing is homotopic to the Lie framing
\((-\partial_a,\partial_b)\) on the parametrizing two-torus.
The reflection \((a,b)\mapsto(-a,b)\) identifies this convention with
the Lie framing \((\partial_a,\partial_b)\), with the orientation sign
of \eqref{wb:h24:ret:eq:sign}.
\end{proposition}

\begin{proof}
Define a complex-linear matrix, also regarded as a real-linear matrix,
\begin{equation}
 D(a,b)=\operatorname{diag}(e^{-i(a+b)},e^{ia},e^{ib})\in\operatorname{SU}(3).
 \hypertarget{tag-03AJ}{}\label{wb:h24:ret:eq:D}
\end{equation}
By \eqref{wb:h24:ret:eq:M}, every \(L_j\) and \(Q_j\) at \((a,b)\) is
\(D(a,b)\) times its value at \((0,0)\).  The coefficients in
\eqref{wb:h24:ret:eq:normal-first}, \eqref{wb:h24:ret:eq:normal-u}, and
\eqref{wb:h24:ret:eq:tangents} are constant on \(M\), since \(t=t_0\).
It follows that the positively oriented full frame is exactly
\begin{equation}
 (N_\ell,N_\phi,N_1,N_2,-T_a,T_b)_{(a,b)}
       =D(a,b)G_0,
 \hypertarget{tag-03AY}{}\label{wb:h24:ret:eq:full-frame}
\end{equation}
where \(G_0\in\operatorname{GL}^+(6,\mathbb R)\) is its single value at \((0,0)\).
This gives the complete phase dependence through the displayed
$D(a,b)$ and the fixed positive matrix $G_0$.

Here is a periodic contraction of \(D\), with all parameters
specified.  For \(v\in[0,1]\) put
\[
 s_v=\sin(\pi v/2),\qquad c_v=\cos(\pi v/2),\qquad
 Q_v(\theta)=
 \begin{pmatrix}s_ve^{-i\theta}&-c_v\\c_v&s_ve^{i\theta}\end{pmatrix},
 \qquad K_v(\theta)=Q_v(\theta)Q_v(0)^*.
\]
The two columns of \(Q_v\) have norm one and zero Hermitian inner
product, and \(\det Q_v=s_v^2+c_v^2=1\).  Thus
\(Q_v,K_v\in\operatorname{SU}(2)\).  Direct substitution gives
\[
 K_0(\theta)=I,\quad
 K_1(\theta)=\operatorname{diag}(e^{-i\theta},e^{i\theta}),\quad
 K_v(0)=I.
\]
Let \(\iota_{01},\iota_{02}:\operatorname{SU}(2)\to\operatorname{SU}(3)\) act in the
ordered coordinate planes \((0,1)\) and \((0,2)\), respectively,
and as the identity in the unused coordinate.  Then
\begin{equation}
 D_v(a,b)=\iota_{01}(K_v(a))\,\iota_{02}(K_v(b))
 \hypertarget{tag-03AK}{}\label{wb:h24:ret:eq:D-homotopy}
\end{equation}
is a well-defined continuous map on
\([0,1]\times(\mathbb R/(2\pi\mathbb Z))^2\), with
\(D_0=I\) and \(D_1=D\).  This proves the needed null-homotopy
directly, including both phase seams.  In particular
\(D_v(a,b)G_0\) contracts the comparison matrix in
\eqref{wb:h24:ret:eq:full-frame} to one constant positive matrix.

In the chart of \cref{wb:h24:ret:chart-lemma}, define
\[
 J:\mathbb R^4\oplus TM\longrightarrow T\mathbb C^3|_M,
 \qquad
 J(x,V)=x_1N_\ell+x_2N_\phi+x_3N_1+x_4N_2+V.
\]
In the tangent coordinates \((-\partial_a,\partial_b)\), its matrix
is the positive matrix
\[
 G(a,b)=D(a,b)G_0
\]
of \eqref{wb:h24:ret:eq:full-frame}, with complex matrices understood
as their realifications. Its inverse is
\(G_0^{-1}D(a,b)^{-1}\). The first four coordinates of this inverse
give the inverse of the normal-frame map to \(TX|_M/TM\): adding a
tangent vector changes only the last two coordinates. Thus \(J\)
is the stable tangent trivialization determined by the given normal
frame and the ambient coordinate frame.

The homotopy \eqref{wb:h24:ret:eq:D-homotopy} joins \(D\) to the
identity through periodic invertible matrices. Multiplication by
\(G_0\) therefore joins \(G\) to a constant positive matrix, with
the inverse at every time given by the reverse product of the inverses.
Here is an explicit path from that constant to the identity. Put
\[
 P=(G_0^{\mathsf T}G_0)^{1/2},\qquad O=G_0P^{-1}.
\]
The positive spectral calculus gives the path \(G_0P^{-t}\),
\(0\leq t\leq1\), with inverse \(P^tG_0^{-1}\).
Its endpoint \(O\) is orthogonal with determinant one. In a fixed real
orthogonal basis it is a sum of planar rotation blocks and identity
blocks; the minus-one eigenspaces pair into rotation blocks of angle
\(\pi\). Replacing each rotation angle \(\theta\) by
\((1-t)\theta\) joins \(O\) to the identity. Reversing the angles
gives the inverse matrices. This supplies a homotopy of stable
trivializations from \(J\) to the stated Lie trivialization.

The literal target four-cell coordinates have the positive diagonal
differential factor described after
\eqref{wb:h24:ret:eq:normal-u}. Denote it by \(K\).
The matrices \((1-t)K^{-1}+tI_4\) have positive diagonal entries
and inverse given by their reciprocals. Multiplying the normal frame
by this path carries the literal target-coordinate frame to the one
used in \(J\), preserving its induced stable-framing class.

For the comparison with the source sphere, pull back its usual stable
trivialization \(TS^6\oplus\mathbb R\cong\mathbb R^7\) by the fixed
based diffeomorphism. On the single contractible chart containing
\(M\), the transition from the coordinate frame plus one real line
is a continuous invertible matrix \(F(z)\).
Let \(\sigma=\operatorname{sign}\det F\), constant on this chart, and
put \(R_\sigma=\operatorname{diag}(\sigma,1,\ldots,1)\).
Then \(F(z)R_\sigma\) is positive and has inverse
\(R_\sigma F(z)^{-1}\). Contracting the chart to its centre contracts
this map to a constant positive matrix, which is joined to the
identity by the preceding spectral and rotation construction.
Thus the entire transition is homotopic to the constant
\(R_\sigma\). The orientation-preserving sphere identification
gives \(R_\sigma=I\); reversing that orientation gives exactly one
constant reflection. The latter negates the framed bordism class.
The torus-coordinate reflection in the statement has derivative
\(\operatorname{diag}(-1,1)\) and is its own inverse, proving the
specified orientation comparison.

The same periodic matrices give the following explicit circle-loop
identity. For the real-linear involution
$\mathcal J(z_1,z_2)=(z_1,\overline z_2)$, with inverse $\mathcal J$,
\[
 \mathcal J\operatorname{diag}(e^{ia},e^{ia})\mathcal J
 =\operatorname{diag}(e^{ia},e^{-ia})=K_1(-a).
\]
Thus $\mathcal J K_v(-a)\mathcal J$ joins the identity to
$\operatorname{diag}(e^{ia},e^{ia})$ and its reversed parameter
contracts that loop. All matrices are periodic in $a$ and all
inverse products are defined on the same real four-dimensional space.
\end{proof}

\begin{lemma}[The spin invariant of the Lie torus]
\hypertarget{tag-Q023}{}\label{mapsIII:lem:lie-torus-spin-invariant}
{\normalfont\scriptsize tag Q023}\par
For an oriented closed spin surface, the Arf sum of its spin quadratic
form is independent of the symplectic basis and invariant under spin
bordism. The Lie-framed two-torus has quadratic form
$q(xe+yf)=x+xy+y$ and Arf invariant one. Its stable framed class is
$\eta^2\ne0$ in $\pi_2^{\mathrm S}=\mathbb Z/2$; reversing one
Lie tangent vector changes its sign and hence preserves this class.
\end{lemma}

\begin{proof}

A stable framing gives a spin structure by lifting its oriented frame
to the trivial \(\operatorname{Spin}\)-bundle.  On the Lie-framed
circle this is the disconnected double cover, whereas the spin
structure induced from the boundary of a framed disc is the
connected double cover: as the boundary tangent turns once, its
\(2\pi\)-rotation lifts from \(1\) to \(-1\) in
\(\operatorname{Spin}(2)\).  We call the former nonbounding and the
latter bounding.

For a closed oriented spin surface \(S\), the spin quadratic
refinement
\[
 q_S:H_1(S;\mathbb F_2)\longrightarrow\mathbb F_2
\]
takes a class represented by an embedded circle to zero or one
according to whether the induced spin structure, using its normal
line in the surface, is bounding or nonbounding.  It satisfies
\begin{equation}
 q_S(x+y)=q_S(x)+q_S(y)+(x\mathbin{\cdot}y).
 \hypertarget{tag-03B6}{}\label{wb:h24:ret:eq:quadratic}
\end{equation}
This is the usual spin quadratic form: resolving transverse
intersections adds one tangent turn per intersection, and the
\(2\pi\)-spin-lift calculation above gives the last term.  The
well-defined correspondence between spin structures and such
quadratic refinements is given in
\cite[Theorem~3.2 and Remark~3.3, p.~6]{DebrayGunningham2018Root}; its original
surface proof is \cite[Theorem~3A, p.~371]{Johnson1980Root}.

For a symplectic basis \((e_i,f_i)_{i=1}^g\), put
\begin{equation}
 \operatorname{Arf}(S)=\sum_{i=1}^gq_S(e_i)q_S(f_i)\in\mathbb F_2.
 \hypertarget{tag-03AI}{}\label{wb:h24:ret:eq:Arf}
\end{equation}
The basis independence can be seen without selecting a preferred
basis: factor the finite sum over the symplectic planes to obtain
\[
 \sum_{x\in H_1(S;\mathbb F_2)}(-1)^{q_S(x)}
   =\prod_{i=1}^g
     \sum_{a,b\in\mathbb F_2}(-1)^{a q_S(e_i)+b q_S(f_i)+ab}
   =2^g(-1)^{\sum_i q_S(e_i)q_S(f_i)}.
\]
The left side has no choice of basis, proving the assertion.

Here is also the bordism-invariance argument needed for the present
obstruction.  A spin bordism has a handle decomposition.  An index
zero or three handle creates or removes a spin sphere, whose first
homology and Arf invariant are zero.  An index one handle either
connects two components, in which case the first-homology quadratic
forms add orthogonally, or increases the genus of one component.
In the latter case the new symplectic pair contains the meridian
of the handle.  That meridian bounds the cocore disc with its
induced spin structure, so its quadratic value is zero; the pair
contributes zero to \eqref{wb:h24:ret:eq:Arf}, regardless of the value on its
dual.  Existing cycles can be represented outside the attaching
discs, so their quadratic values are preserved.  An index two
handle is the reverse of an index one handle.  Therefore each
handle preserves the total Arf invariant of the boundary, and a
spin boundary has Arf invariant zero.  This also proves that a
stably framed boundary has Arf invariant zero, because forgetting
the framing gives a spin bordism.  The full spin bordism statement
is recorded in \cite[Theorem~3.6, p.~6]{DebrayGunningham2018Root}.

For the intrinsic Lie framing, \cite[Example~3.8, p.~6]{DebrayGunningham2018Root} gives the following quadratic values. On
\(M=(\mathbb R/(2\pi\mathbb Z))^2\), take
\[
 e=[a\mapsto(a,0)],\qquad f=[b\mapsto(0,b)].
\]
Both the tangent and normal vectors in the surface are constant in
the Lie frame along either circle.  Thus both circles have the
nonbounding spin structure, and
\begin{equation}
 e\mathbin{\cdot}f=1,\qquad q_M(e)=q_M(f)=1,
 \qquad q_M(xe+yf)=x+xy+y,
 \qquad \operatorname{Arf}(M)=1.
 \hypertarget{tag-03AQ}{}\label{wb:h24:ret:eq:actual-Arf}
\end{equation}
Reversing the first Lie tangent vector changes the orientation but
does not change these mod-two values.  In particular the
specifically framed inverse image computed above is not a framed
boundary.

Let $\eta\in\pi_1^{\mathrm S}$ be the stable complex Hopf class.
The circle-framing calculation in the proof of
Theorem~\ref{maps:src:B.1.2} applies in the same coordinates:
its map is $(z_1,z_2)\mapsto[z_1:z_2]$, its fibre is
$(e^{ia},0)$, and its ordered derivative normal vectors are
$(0,e^{ia})$ and $(0,ie^{ia})$ in the coordinate $z_2/z_1$.
These are the same complex coordinates, with circle parameter
$\theta=a$, so the translation-framed circle represents $\eta$. The additional explicit
$\mathcal J K_v(-a)\mathcal J$ contraction is proved in
Proposition~\ref{wb:h24:ret:framing-proposition}.

The product of two normally framed inverse images represents the
product in stable homotopy: its tubular collapse is the smash of
the two tubular collapse maps, as follows immediately from the
product tubular neighborhood \(\nu_1\times\nu_2\) and its ordered
normal coordinates.  Consequently the product of the two Lie
circles represents \(\eta^2\).  Equation \eqref{wb:h24:ret:eq:actual-Arf}
proves \(\eta^2\ne0\) by an explicit bordism invariant.  The
stable group \(\pi_2^{\rm S}=\mathbb Z/2\) is Pontryagin's computation
\cite[pp.~957--959]{Pontryagin1950Root} and \cite[Theorem~24, pp.~109--112]{Pontryagin1955Root}; the stable table in
\cite[opening table, p.~339, and Corollary~4.24, pp.~360--361]{Hatcher2002} gives the same group.  Thus the sign from
\eqref{wb:h24:ret:eq:sign} satisfies \(-\eta^2=\eta^2\).

\end{proof}

\subsection{The retained-angle homotopy class}

\begin{theorem}[The retained-angle homotopy class]\hypertarget{tag-02J7}{}\label{maps:src:B.3.1}
{\normalfont\scriptsize tag 02J7}\par


For the original sphere member and every defined admissible $|p|=1$
member in Definition~\ref{mapsIII:def:retained-angle-cusp-data},
the retained-angle map of Definition~\ref{mapsIII:def:cusp-suspension-maps}
satisfies
\begin{equation}
 [H]=\eta_4\eta_5\ne0\in\pi_6(S^4)=\mathbb Z/2.
 \hypertarget{tag-01UZ}{}\label{maps:eq:B.3.1}
\end{equation}
Its regular inverse image is the native vanishing
two-torus with its Lie stable framing and Arf invariant one.

\end{theorem}

\begin{proof}
\begin{enumerate}[label=\arabic*.,leftmargin=*]
\item The full inverse image and submersion assertions are
Lemma~\ref{wb:h24:ret:chart-lemma} and
Lemma~\ref{mapsIII:lem:retained-angle-differential}. Apply
Lemma~\ref{mapsIII:based-regular-fibre-smoothing} with
\[
 f=H,\qquad N=M,\qquad x_0=x_{\rm base},\qquad y=y_0,
 \qquad n=4.
\]
Write $H_b=f_b$ for the first endpoint in that lemma and
take the protected tube $T=U_1$;
its first chart $\kappa$ and bump $\beta$ make the map constant
near the basepoint before the positive gap
\[
 d_0=\min_{x\in X\setminus\operatorname{int}U_1}
                 d(H_b(x),y_0)>0
\]
is chosen. The lemma supplies the complete based smoothing and
short-geodesic homotopy, fixed on this same tube and near the
basepoint, with error less than $d_0/2$ on its complement.
Its endpoint has exactly the same full inverse image and
differential at $y_0$. The ordered target chart is
\eqref{wb:h24:ret:eq:four-cell}; transporting its frame by
the positive diagonal map in
Lemma~\ref{mapsIII:lem:retained-angle-differential}, with the
displayed reciprocal inverse, gives precisely the four normal
vectors used below.

\item Pontryagin--Thom therefore identifies the stabilization of
\([H]\) with the bordism class of \(M\) carrying the four normal
vectors \eqref{wb:h24:ret:eq:normal-first}--\eqref{wb:h24:ret:eq:normal-u}.  The comparison
\eqref{wb:h24:ret:eq:full-frame} and its explicit contraction identify that
class with the Lie torus, with the computed order-two orientation sign.  Lemma~\ref{mapsIII:lem:lie-torus-spin-invariant} therefore gives
\[
 E^\infty[H]=\eta^2\ne0.
\]
The same ordered stable product satisfies
$\eta\eta^2=\eta^3=12\nu\ne0$ by
Theorem~\ref{maps:src:B.1.2}; it uses the identical Hopf-circle
class just identified.
Freudenthal suspension is an isomorphism
\[
 E:\pi_{n+2}(S^n)\longrightarrow\pi_{n+3}(S^{n+1})
 \qquad(n\geq4),
\]
because \(n+2\leq2n-2\).  Therefore stabilization from
\(\pi_6(S^4)\) is an isomorphism onto \(\pi_2^{\rm S}\).
The composite \(\eta_4\circ\eta_5:S^6\to S^4\) stabilizes to
\(\eta^2\), so injectivity gives \eqref{maps:eq:B.3.1}.
\end{enumerate}
\end{proof}

\begin{corollary}[Order-two quotient detected by the cusp map]
\hypertarget{tag-03BE}{}\label{wb:h24:ret:q2-detected-quotient}
{\normalfont\scriptsize tag 03BE}\par
Let \(d:(S^6,*)\to(X,x_{\mathrm{base}})\) be the fixed based
diffeomorphism, and use the maps \(q_2\) and \(p_3\) of
\eqref{wb:h24:ret:eq:q2} and \eqref{wb:h24:ret:eq:H}. Put
\[
 G=\pi_6(\mathcal A,v_0),\qquad \alpha=[q_2\circ d]\in G.
\]
Postcomposition induces a surjective homomorphism
\[
 \varepsilon=(p_3)_*:G\longrightarrow\pi_6(S^4,*)\cong\mathbb Z/2,
 \qquad \varepsilon(\alpha)=\eta_4\circ\eta_5.
\]
Identify \(\eta_4\circ\eta_5\) with \(1\in\mathbb Z/2\).
For \(K=\ker\varepsilon\), the induced isomorphism and its inverse are
\[
 \begin{aligned}
 \overline\varepsilon:G/K&\longrightarrow\mathbb Z/2,
 &\beta+K&\longmapsto\varepsilon(\beta),\\
 (\overline\varepsilon)^{-1}:\mathbb Z/2&\longrightarrow G/K,
 &0&\longmapsto K,\quad 1\longmapsto\alpha+K.
 \end{aligned}
\]
Every odd multiple of \(\alpha\) is nonzero. Its order is infinite
or an even positive integer.
\end{corollary}

\begin{proof}
Postcomposition with the based map \(p_3\) carries based homotopies
to based homotopies and commutes with the source-sphere pinch map.
It therefore induces the displayed group homomorphism. The defining
identity \(H=p_3q_2\) and \cref{maps:src:B.3.1} give
\[
 \varepsilon(\alpha)=[p_3q_2d]=[Hd]=\eta_4\eta_5\ne0.
\]
This is the nonzero element of the target group of order two,
so \(\varepsilon\) is surjective. A based null-homotopy of
\(q_2\) would postcompose with \(p_3\) to a based null-homotopy
of \(H\), contradicting \eqref{maps:eq:B.3.1}.

The map \(\overline\varepsilon\) is well defined on cosets because
\(\varepsilon(\beta+k)=\varepsilon(\beta)\) for every \(k\in K\).
Moreover \(2\alpha\in K\), since \(\varepsilon(2\alpha)=0\),
so its displayed inverse respects addition in \(\mathbb Z/2\).
For \(\beta\in G\), let \(e\in\{0,1\}\) represent
\(\varepsilon(\beta)\). Then
\[
 \varepsilon(\beta-e\alpha)=0,\qquad
 \beta+K=e\alpha+K.
\]
This verifies the inverse on every coset; the other composite is
the identity on the two elements of \(\mathbb Z/2\).

For every integer \(n\),
\[
 \varepsilon(n\alpha)=(n\bmod2)\eta_4\eta_5.
\]
An odd multiple has nonzero image. Every integer annihilating
\(\alpha\) is even, which proves the order assertion.
\end{proof}

\subsection{Smash powers of the angle-forgetting maps}
\hypertarget{tag-02Q9}{}\label{mapsIII:carrier-extension-section}

\begin{proposition}[Smash powers of cusp quotient maps]
\hypertarget{tag-02Q7}{}\label{mapsIII:angle-forgotten-one-copy-null}
{\normalfont\scriptsize tag 02Q7}\par
For the original sphere member and $x_0\in W$, every positive
smash power of either map $\widehat q_3$ or $\widehat q_\gamma$
of Definition~\ref{mapsIII:def:cusp-suspension-maps} is based null-homotopic.
\end{proposition}

\begin{proof}
Use either based null-homotopy $H$ from Theorem~\ref{maps:src:B.3.2}.
For each integer
$a\geq1$, the formula
\[
 ([x_1,\ldots,x_a],s)
       \longmapsto[H(x_1,s),\ldots,H(x_a,s)]
\]
descends to the smash quotient because $H(x_0,s)=*$.
The quotient map from the compact Hausdorff product $X^a$
remains a quotient after taking its product with the
interval, so the descended homotopy is continuous.
At $s=0$ it is the specified smash power and at $s=1$
it is constant at the basepoint.
\end{proof}

\subsection{Quotient maps in three sphere-fibration families}

\begin{definition}[The three marked sphere families]
\hypertarget{tag-Q014}{}\label{mapsIII:def:three-sphere-families}
{\normalfont\scriptsize tag Q014}\par
Use the sphere members of Viaclovsky's two-parameter family,
the Liu--Rollenske family, and the sphere members $X_{3k,4l}$
with $k,l\ge1$ and $\gcd(3k,4l)=1$, with the sections,
integral logarithmic collar lifts, twists and cusp markings of
those constructions. Their marked lattices, primitive characters
and cusp vanishing lattices are the following.
\[
\begin{array}{c|c|c|c}
 &\Lambda&\chi&K\\\hline
 \text{two-parameter}&
 \langle\widehat\psi,\widehat u,\widehat w,\widehat\delta\rangle&
 \psi&\langle\widehat w,\widehat\delta\rangle\\
 \text{Liu--Rollenske}&\langle e_1,e_2,e_3,e_4\rangle&
 \xi=e_4^*&\langle e_1-e_2-e_3,\,2e_3-e_1\rangle\\
 \text{logarithmic}&\langle\delta,a,b,d\rangle&
 \psi=d^*&\langle\delta,a\rangle .
\end{array}
\]
The cusp quotient maps use the product-cell collapses and suspension
conventions of Definition~\ref{mapsIII:def:cusp-suspension-maps},
in each prescribed marked long-period basis.
Put $L=\ker\chi$. For the logarithmic family use its constructed
integral lifts and write
\[
 m_0=3k,\quad m_1=4l,\quad
 u_0=d+2a+2b,\quad u_1=d+3a+3b,\quad
 v_i=\rho_i\delta+N_i u_i.
\]
The sphere condition is
\[
 b_{\mathrm{sph}}=m_1N_0+m_0N_1=\pm1.
\]
\end{definition}

\begin{theorem}[Quotient and cusp maps in three families]\hypertarget{tag-02JA}{}\label{maps:src:B.4.1}
{\normalfont\scriptsize tag 02JA}\par
For the data of Definition~\ref{mapsIII:def:three-sphere-families},
the primitive invariant-circle quotient extends
over each degeneration neighbourhood into its target solid torus.
Its homotopy class in that solid torus, relative to the fixed
outer collar, is unique. The resulting target is diffeomorphic
to $S^3$, and the quotient map has class $6\nu'$.
In each case the map to $\mathbb P^1$ is null-homotopic.
The retained-angle map at the primitive rank-two cusp has class
$\eta_4\eta_5$, and the two corresponding angle-forgetting maps
are null-homotopic. 
\end{theorem}

\begin{proof}
\begin{enumerate}[leftmargin=*]
\item We use the sphere identifications in
\cite[Theorem~1.1, p.~2; Proposition~6.4, pp.~19--20]{manViaclovskyTwo26},
\cite[Theorem~5.1, p.~21]{manLR26}, and
\cite[Corollary~4.6, p.~14]{manViaclovskyLogs26}.
Reducing $m_1N_0+m_0N_1=\pm1$ modulo $m_i$ proves
$\gcd(N_i,m_i)=1$ for each logarithmic filling.
In the three rows of the input table the bases of $L$ are,
respectively, $(\widehat u,\widehat w,\widehat\delta)$,
$(e_1,e_2,e_3)$ and $(\delta,a,b)$.
Their coefficient maps $L\to\mathbb Z^3$ have inverse the
corresponding integral linear combinations. The complementary
vectors $\widehat\psi,e_4,d$ have character value one.
Thus the quotient $[x]\mapsto\chi(x)\bmod\mathbb Z$
has connected kernel $T_L=L_{\mathbb R}/L$.
The integral map from the common marking to
$(\delta,a,b,d)$ is
\[
 P_V=\begin{pmatrix}
 -2&1&0&1\\-6&0&-1&0\\0&-1&0&0\\-1&0&0&0
 \end{pmatrix},\qquad \det P_V=-1.
\]
Its columns give
\[
 \begin{aligned}
 P_V\widehat\gamma&=-2\delta-6a-d,&
 P_V\widehat u&=\delta-b,\\
 P_V\widehat w&=-a,&P_V\widehat\delta&=\delta.
 \end{aligned}
\]
The torus isomorphism is $[x]\mapsto[P_Vx]$, with inverse
$[y]\mapsto[P_V^{-1}y]$, where
\[
 P_V^{-1}=\begin{pmatrix}
 0&0&0&-1\\0&0&-1&0\\0&-1&0&6\\1&0&1&-2
 \end{pmatrix}.
\]
The two matrix products are the identity.
It intertwines the finite matrices of
\SeriesRef{mapsI:src:A.1.1} with those in
\cite[Lemma~2.1, pp.~3--4]{manViaclovskyLogs26}.
With the integral translation vectors of
\SeriesRef{push:src:E.1.1}, it sends $w_i$ to $-v_i$;
the deck-generator inversion gives
$p=-b_{\mathrm{sph}}$ and $d^*P_V=-\gamma$.


\item For the two-parameter family use
\SeriesRef{push:src:E.5.2}. Its dual homology matrices are
\[
\begin{aligned}
 H_*&=\begin{pmatrix}1&0&0&0\\-2&0&-1&0\\0&1&0&0\\1&1&0&1\end{pmatrix},&
 H_1&=\begin{pmatrix}1&0&0&0\\0&1&0&0\\-2&-1&1&0\\-2&-1&0&1\end{pmatrix},\\
 H_2&=\begin{pmatrix}1&0&0&0\\0&1&1&0\\0&0&1&0\\0&0&0&1\end{pmatrix},&
 H_3&=\begin{pmatrix}1&0&0&0\\0&1&0&0\\0&-1&1&0\\1&0&0&1\end{pmatrix}.
\end{aligned}
\]
The period identity $\Pi_iH_i=G_i\Pi$ identifies the geometric
return in period coefficients as $H_i^{-1}$; its cohomological
pullback is $C_i=H_i^{-t}$.
Solving the fixed-vector equations gives
\begin{equation}
 \Lambda^\Gamma=\mathbb Z\widehat\delta,\qquad
 (\Lambda^*)^\Gamma=\mathbb Z\psi,\qquad
 (\Lambda^*\otimes\mathbb F_2)^\Gamma=\mathbb F_2\psi.
 \hypertarget{tag-01V4}{}\label{maps:eq:B.4.1}
\end{equation}
The primitive vanishing lattices of the reduced fibres are
\begin{equation}
 K_1=\mathbb Z(\widehat w+\widehat\delta),\qquad
 K_2=\mathbb Z\widehat u,\qquad
 K_3=\mathbb Z\langle\widehat w,\widehat\delta\rangle.
 \hypertarget{tag-01V5}{}\label{maps:eq:B.4.2}
\end{equation}
For Liu--Rollenske the homological returns are
\[
 R_0=\begin{pmatrix}-1&-1&-1&0\\2&1&1&0\\0&0&1&0\\0&0&0&1\end{pmatrix},
 \qquad
 R_1=\begin{pmatrix}1&1&0&1\\-2&0&-1&-1\\0&-1&1&-1\\0&0&0&1\end{pmatrix},
\]
with $R_\infty=(R_1R_0)^{-1}$ and pullbacks $C_i=R_i^t$.
They give
\begin{equation}
 \Lambda^\Gamma=\mathbb Z(-e_1+2e_3),\qquad
 (\Lambda^*)^\Gamma=\mathbb Z e_4^*,\qquad
 (\Lambda^*\otimes\mathbb F_2)^\Gamma=\mathbb F_2e_4^*.
 \hypertarget{tag-01V6}{}\label{maps:eq:B.4.3}
\end{equation}
These are the markings in
\cite[Propositions~3.12--3.13, pp.~13--14]{manLR26}; in the
local basis at $1$ they satisfy
$e_1=2f_1+f_3$, $e_2=f_2$, $e_3=f_1+f_3$,
$e_4=f_2+f_4$. The inverse basis equations are
$f_1=e_1-e_3$, $f_2=e_2$, $f_3=2e_3-e_1$ and $f_4=e_4-e_2$.
Thus $\delta=f_3$ and $\xi=f_4^*$.
In the logarithmic marking the returns are
\[
\begin{aligned}
 A_0&=\begin{pmatrix}1&-1&0&2\\0&-1&-1&6\\0&1&0&0\\0&0&0&1\end{pmatrix},&
 A_1&=\begin{pmatrix}1&0&1&-3\\0&0&1&0\\0&-1&0&6\\0&0&0&1\end{pmatrix},\\
 A_\infty&=\begin{pmatrix}1&0&0&1\\0&1&1&0\\0&0&1&0\\0&0&0&1\end{pmatrix},&
 C_i&=A_i^t.
\end{aligned}
\]
Their common mod-two cohomology fixed space is $\mathbb F_2\psi$.
The degree-three fixed spaces, calculated by taking the
$3\times3$ minors of the displayed pullbacks, are
\[
\begin{array}{c|c}
 &\bigcap_i\ker(\wedge^3C_i-I)\subset H^3(F;\mathbb Q)\\\hline
 \text{two-parameter}&\mathbb Q(\psi\wedge u\wedge w)\\
 \text{Liu--Rollenske}&\mathbb Q(2e_{124}^*-e_{234}^*)\\
 \text{logarithmic}&\mathbb Q(a^*\wedge b^*\wedge\psi).
\end{array}
\]
For example, in the two-parameter exterior basis $123,124,134,234$,
$C_3-I$ on exterior cubes first sets the last two coefficients
to zero, and $C_2-I$ sets the second to zero. In the
Liu--Rollenske case $\wedge^3R_0^t-I$ sets the third
coefficient to zero and the second to minus twice the fourth;
$\wedge^3R_1^t-I$ then sets the first to zero. In the
logarithmic case the first matrix fixes vectors
$(r,-2r,2r,s)$ and the second forces $r=0$.
Every displayed class restricts to zero on $T_L$.

\item On the ordinary locus the chosen sections give the map
$(q,[x])\mapsto(q,\chi(x))$. For a finite logarithmic filling
its circle target is
\[
 Q_{m,N}=(D^2_s\times S^1_\theta)/
 ((s,\theta)\sim(e^{2\pi i/m}s,\theta+N/m)).
\]
Choose integers $a,c$ with $aN+cm=1$. The functions
\[
 z=s e^{-2\pi ia\theta},\qquad v=e^{2\pi im\theta}
\]
identify this quotient with $D^2\times S^1$.
For the inverse choose a solution of $e^{2\pi im\theta}=v$ and put
$s=z e^{2\pi ia\theta}$. The $m$ choices of $\theta$ are permuted
transitively by the deck generator, since $\gcd(N,m)=1$;
$aN+cm=1$ makes the corresponding $s$ values deck equivalent.
Substitution gives both inverse composites on the quotient.
Its base map is $q=z^mv^a$.
The comparison from the filling to the ordinary part subtracts
$N\arg s/(2\pi)$ from the circle coordinate. Thus, with $x$
the section lift of a positive base loop and $h$ the positive
circle, the meridian is
\[
 \mu=mx-Nh.
\]
This is the inverse logarithmic gluing of
\cite[equations~(2.15)--(2.17), pp.~5--6; Proposition~3.1,
pp.~6--7]{manViaclovskyLogs26}.
For Liu--Rollenske the data are $(m,N)=(4,1)$, because the
translation is $e_4/4$ and $\xi(e_4)=1$
\cite[Proposition~2.12 and Definition~2.13, pp.~7--8]{manLR26}.
For the two-parameter construction the translation is
$\ell/4$, $\ell=(1,-1,-1,0)^t$, with $\psi(\ell)=1$;
the collar lift is fixed by $\gamma_*^4=\ell$ in
\cite[equation~(3.14), p.~8]{manViaclovskyTwo26}. It gives the same
meridian $4x-h$.

At the two reduced rank-one fibres the boundary maps kill
$K_1$ and $K_2$ respectively, and the section disc kills
the base meridian. The nodal thimble kills the indicated
primitive fibre circle, while the other three fibre circles
give $\pi_1(N_i)=\Lambda/K_i$; these are the neighbourhood
maps of \cite[Proposition~3.1, pp.~6--7]{manViaclovskyTwo26}.
For Liu--Rollenske the neighbourhood at $1$ is
$S_{\Delta_1}\times S^1_\delta\times S^1_\xi$ over
$\Delta_1$, with $S_{\Delta_1}$ simply connected, by
\cite[Lemma~3.10, pp.~11--12]{manLR26}.
Its second-circle projection extends $\xi$.

At a rank-two cusp consider the space $\widetilde N$ before
the long-period quotient. Its punctured part is
$(\mathbb C^*)^2\times\Delta^*$, with coordinates
$(z_1,z_2,q)$, and a triple-point chart has $q=w_0w_1w_2$.
Moving a compact loop transversely off the central divisor
shows that the punctured inclusion surjects on fundamental
groups. In the chart $(z_1,z_2)=(w_1,w_2)$, the three
coordinate discs have boundary winding vectors
$(0,0,1),(1,0,1),(0,1,1)$. These form an integral basis
of $\pi_1((\mathbb C^*)^2\times\Delta^*)$ and bound the
displayed discs. For the other square diagonal the chart is
$(z_1,z_2)=(w_1,w_1w_2)$ and the vectors are
$(0,0,1),(1,1,1),(0,1,1)$, again an integral basis.
Multiplication by holomorphic units contributes zero winding.
Hence $\widetilde N$ is simply connected.

In logarithmic coordinates its periods are
$Z(q)=sI_2+C(q)$, $s=(2\pi i)^{-1}\log q$, with $C$
holomorphic on the disc. A long period $n\in\mathbb Z^2$
acts by $p_n(z,q)=(z\exp(2\pi iZ(q)n),q)$.
On the punctured part the real long coordinate
\[
 t(z,q)=(\operatorname{Im}Z(q))^{-1}
 \left(-\frac{\log|z|}{2\pi}\right)
\]
satisfies $t(p_n(z,q))=t(z,q)+n$.
In a compact triple-point polydisc the ratios
$\log|w_j|/\log|q|$ are bounded, so this identity permits
only finitely many translates meeting that polydisc.
On the central divisor a compact chart meets finitely many
indexed components, and translation of their lattice indices
gives the same finiteness. A stabilizer would translate a
finite incident cell onto itself and hence is zero.
The period quotient is therefore a covering, and it gives
\[
 \pi_1(N)=\Lambda/K,\qquad
 \pi_1(\partial N)=\Lambda\rtimes\mathbb Z
       \longrightarrow\Lambda/K,
 \qquad (\lambda,j)\longmapsto[\lambda].
\]
Here a fibre loop with long coefficient $n$ lifts from a point
to its $p_n$-translate, the toric circles bound the coordinate
discs, and the section-defined base loop bounds its section
disc. These are also the specified boundary maps in
\cite[Proposition~3.1, pp.~6--7]{manViaclovskyLogs26}.

In every reduced neighbourhood $\chi$ therefore factors
through $\pi_1(N)$. An equivariant real function on its
associated cover, obtained by a partition of unity, gives a
circle map with this homomorphism. On the outer collar its
ratio with the prescribed quotient has a real logarithm;
a cut-off of that logarithm makes the two maps equal there.
Together with the base coordinate it gives $N\to D^2\times S^1$.
The boundary surjects on $\pi_1(N)$ in all these models and
in a finite filling. Two extensions with the same outer-collar
values consequently have a circle-coordinate ratio with a
real logarithm zero on the collar. Multiplying by its scalar
homotopy and using the straight homotopy for the disc
coordinates proves uniqueness in the stated relative sense.

\item Fill the reduced target discs first. The two finite
data in the logarithmic family are $(m_i,N_i)$ above; for
each other family they are $(m_0,N_0;m_1,N_1)=(4,1;1,0)$.
On the remaining annular boundary use $x=x_0=-x_1$ and $h$.
Then
\[
 \mu_0=m_0x-N_0h,\qquad \mu_1=-m_1x-N_1h,
\]
and their determinant is $-b_{\mathrm{sph}}$, with
$b_{\mathrm{sph}}=1$ for the latter data. The inverse is
\[
 x=b_{\mathrm{sph}}(N_1\mu_0-N_0\mu_1),\qquad
 h=-b_{\mathrm{sph}}(m_1\mu_0+m_0\mu_1).
\]
Sending $\mu_0,\mu_1$ to the standard meridians of the
genus-one splitting of $S^3$ gives a torus diffeomorphism
extending over both solid tori. In its angular coordinates
$(\phi_1,\phi_2)$, the base phase is
$m_0\phi_1-m_1\phi_2$: it evaluates to one on $x$ and
to zero on $h$. Choosing the radial coordinates to agree
with the local base maps identifies the projection, up to
radial homotopy, with
\[
 \begin{aligned}
 \pi(z_1,z_2)=[z_1^{m_0}:z_2^{m_1}]
     &=\eta_2 F_{m_0,m_1}(z_1,z_2),\\
 F_{m_0,m_1}&=
 \frac{(z_1^{m_0},z_2^{m_1})}
 {\sqrt{|z_1|^{2m_0}+|z_2|^{2m_1}}}.
 \end{aligned}
\]
For $t=|z_1|^2$ the first squared radius is
\[
 R(t)=\frac{t^{m_0}}{t^{m_0}+(1-t)^{m_1}},\qquad
 R'(t)=\frac{t^{m_0-1}(1-t)^{m_1-1}
                (m_0(1-t)+m_1t)}
               {(t^{m_0}+(1-t)^{m_1})^2}>0\quad(0<t<1).
\]
It runs from zero to one. A regular value with both entries nonzero
has exactly $m_0m_1$ phase preimages, each with positive local degree.
Thus $\deg F_{m_0,m_1}=m_0m_1$.

\item Choose an oriented sphere identification $j:X\to S^6$
and the stable frame $\tau_x(v,a)=dj_x(v)+a j(x)$.
Compare it along a regular fibre with the four translation
fields, two oriented horizontal lifts and the stabilizing
field. Polar decomposition gives their difference
$g:F\to SO(7)$. Along a base loop with affine return
$h(x)=Ax+a$, the endpoint frames differ by
$\operatorname{diag}(A,I_3)$ and a horizontal shear.
The shear contracts by scalar multiplication, $A$ is joined
to the identity in $GL^+(4,\mathbb R)$, and translation by
$a$ is isotopic to the identity. Hence $g\simeq g h$.

Let $\alpha\in H^1(SO(7);\mathbb F_2)$ be the Spin-covering
class. The fixed-space calculations show that $g^*\alpha$
is a multiple of $\chi\bmod2$, so $g|_{T_L}$ lifts to
$\widetilde g:T_L\to\operatorname{Spin}(7)$.
For an integral generator $u_3$ of its third cohomology,
the covering transfer gives the rational class
$\beta=\tfrac12\operatorname{tr}(u_3)$ on
$SO(7)$ with $c^*\beta=u_3$: the deck transformation is
translation by a central element of the connected Spin
group and acts trivially on cohomology. The invariant
degree-three classes displayed above all restrict to zero
on $T_L$. Thus $\widetilde g^*u_3=0$ rationally and then
integrally, since $H^3(T^3;\mathbb Z)$ is torsion-free.
The group $\operatorname{Spin}(7)$ is two-connected and
its third homotopy group is detected by $u_3$. Contracting
the two-skeleton of $T^3$ leaves a map from $S^3$ with
zero integer class. This proves $g|_{T_L}\simeq *$.

A target value in the ordinary locus has exactly one
inverse-image coset of $T_L$. Its differential normal
frame consists of the lifted base frame and a vertical
vector $v_\chi$ with $\chi(v_\chi)=1$; changes of lifts
are the contracting shears just described. Its stable
tangent framing is therefore the Lie framing.
The product Pontryagin--Thom calculation and stabilization
map in the proof of \cref{maps:src:B.1.2} give
\[
 \begin{gathered}
 \sigma[e]=[T^3,\mathrm{Lie}]=\eta^3=12\nu,\\
 \sigma:\mathbb Z/12\langle\nu'\rangle
       \longrightarrow\mathbb Z/24\langle\nu\rangle,
 \quad \nu'\longmapsto2\nu.
 \end{gathered}
\]
This homomorphism is injective, so $[e]=6\nu'$.
Every orientation reversal occurring in the displayed
marking maps induces inversion of this order-two class
and leaves it unchanged.
On $S^3=SU(2)$ a degree-$d$ self-map is homotopic to the
pointwise $d$th-power map. Pointwise multiplication induces
addition on $\pi_6(S^3)$, so its action there is multiplication
by $d$. Here $d=4$ or $12kl$, both even. Consequently
$[F_{m_0,m_1}e]=0$ and $[f]=[\eta_2F_{m_0,m_1}e]=0$.

\item In the native toric order the period matrix is
$Z=sI_2+C(q)$ as above.
For the two-parameter family its native long periods are
$h=(\tau,z)$ and $k=(-2z,-\beta-c)$; in the ordered
periods of \SeriesRef{push:src:E.5.3} the full matrix is
\[
 P_2=\begin{pmatrix}0&-1&0&0\\1&0&0&0\\0&0&1&0\\0&0&0&1\end{pmatrix},
 \qquad
 P_2^{-1}=\begin{pmatrix}0&1&0&0\\-1&0&0&0\\0&0&1&0\\0&0&0&1\end{pmatrix}.
\]
For Liu--Rollenske the long periods are $(h_2,h_4)$ and
the positive toric circles are $(h_1,h_3)$; the marked long
order is $(h_4,h_2)$.
For the logarithmic family the native order $(b,d,a,\delta)$
maps to the common marking by
\[
 P_{\log}=\begin{pmatrix}0&-1&0&0\\-1&0&0&0\\0&6&-1&0\\1&-2&0&1\end{pmatrix},
 \qquad
 P_{\log}^{-1}=\begin{pmatrix}0&-1&0&0\\-1&0&0&0\\-6&0&-1&0\\-2&1&0&1\end{pmatrix}.
\]
These are maps of tori $[t]\mapsto[Pt]$, with inverses
$[v]\mapsto[P^{-1}v]$; multiplying each displayed pair gives the identity.
They include all toric shear entries. They follow from the multiplicative
periods in \cite[equation~(3.1), p.~6; Section~4.2,
p.~10]{manViaclovskyTwo26}, \cite[Proposition~3.12, pp.~13--14]{manLR26},
and \cite[equations~(2.8)--(2.10), p.~4]{manViaclovskyLogs26}.
Their induced maps on the ordered long torus are
\[
 J_2=\begin{pmatrix}0&-1\\1&0\end{pmatrix},\qquad
 J_{\mathrm{LR}}=\begin{pmatrix}0&1\\1&0\end{pmatrix},\qquad
 J_{\log}=\begin{pmatrix}0&-1\\-1&0\end{pmatrix}.
\]
Here $J_2^{-1}=\left(\begin{smallmatrix}0&1\\-1&0\end{smallmatrix}\right)$,
$J_{\mathrm{LR}}^{-1}=J_{\mathrm{LR}}$ and $J_{\log}^{-1}=J_{\log}$.
The retained three-torus map is $(\phi,t)\mapsto(\phi,Jt)$,
with inverse $(\phi,v)\mapsto(\phi,J^{-1}v)$.
After its two-skeleton is collapsed, its top-cell degree
is $\det J=1,-1,-1$ respectively. Indeed the maps from
$T^3$ to the two-connected target $S^3$ are classified by
their degree-three cohomology class; these torus maps
multiply that class by the displayed determinant.

Choose $t^0=(1/5,1/3)$, $q_0=r_0e^{i\pi}$ and a regular
open-cell value in the definition of
\cref{maps:src:B.3.1}. In the two diagonal presentations
of the square use respectively
\[
 M=I_2\quad\text{and}\quad
 M=\begin{pmatrix}1&0\\-1&1\end{pmatrix},
 \qquad (z_1,z_2)\longmapsto(z_1,z_2/z_1)
\]
for the latter change. Their triple-point coordinates
before the unit gauge are respectively
$(q/(z_1z_2),z_1,z_2)$ and $(q/z_2,z_1,z_2/z_1)$.
Their monomial exponent matrices have determinant one.
Set $p=Mt^0$ and
\[
 w_j=(z^M)_j\exp(-2\pi i(MC(q)t^0)_j),\quad j=1,2,
 \qquad w_0=q/(w_1w_2).
\]
The units extend holomorphically to the chart. The entire
regular inverse image is
\[
 (w_0,w_1,w_2)=
 \bigl(r_0^{1-p_1-p_2}e^{i(\pi-a-b)},
       r_0^{p_1}e^{ia},r_0^{p_2}e^{ib}\bigr),
 \qquad (a,b)\in(\mathbb R/2\pi\mathbb Z)^2.
\]
To see both completeness and injectivity, write a point of
the punctured cover as $x+Z(q)t$, with $x$ modulo
$\mathbb Z^2$. The chosen value gives $Jt=Jt^0$ modulo
$\mathbb Z^2$, hence $t=t^0+n$ for a unique integral
long-period translate. Applying $p_{-n}$ leaves exactly
the displayed coordinate-circle torus. Equality between
two of its representatives gives $x-x'=a+Z(q)n$;
imaginary parts force $n=0$.
The positive exponent triples are $(7/15,1/5,1/3)$ and
$(2/3,1/5,2/15)$. Thus the whole torus tends uniformly
to the triple point as $r_0\to0$.
The proper period action proved above has trivial
stabilizer there. Choose a small neighbourhood disjoint
from its nontrivial translates and then decrease $r_0$
until the whole torus lies in a coordinate ball inside it.
Put $\ell=\log|q|$ and $\rho_j=\log|w_j|$. Near this torus,
\[
 (\rho_1,\rho_2)^t
   =\ell p+R(q)(t-t^0),\qquad
 R(q)=\ell M-2\pi\operatorname{Im}(MC(q)).
\]
For small $q$, $R(q)$ is invertible, and on the torus
\[
 dt=R(q)^{-1}\bigl(d(\rho_1,\rho_2)^t-p\,d\ell\bigr).
\]
Together with $d\phi=d\arg(w_0w_1w_2)$ and the radial
derivative, these are the normal covectors of the retained
map. Their independence follows from the invertibility of $R(q)$
and the nonzero radial and angular coordinate derivatives.
Apply Lemma~\ref{mapsIII:based-regular-fibre-smoothing} to this
family's retained map $f=H:X\to S^4$, the entire displayed torus
$N=H^{-1}(y)$, and any fixed $x_0\in W$, where $H(x_0)=*$.
The regular value $y$ is the chosen open-cell value, and the
target coordinates are the
specified product-cell coordinates of
Definition~\ref{mapsIII:def:three-sphere-families}.
The full-fibre and smooth-neighbourhood assertions were proved
above, and the displayed independent covectors prove rank four.
That lemma supplies a based smoothing fixed on a closed tube
around this entire torus, preserving precisely this ordered
normal framing. The ordered tangent and normal frame, relative to
the coordinate-ball frame and a constant change of basis,
varies by
\[
 \operatorname{diag}(e^{-i(a+b)},e^{ia},e^{ib})\in SU(3).
\]
This is the matrix $D(a,b)$ of
Proposition~\ref{wb:h24:ret:framing-proposition}, under the identity
map on the two angular coordinates and the constant frame change
just described. Its two displayed $K_v$ factors give the periodic
contraction of this comparison matrix.
The ambient stable frame also extends over this ball.
The inverse-image torus therefore has its Lie framing;
both coordinate circles have the nonbounding spin structure,
and its Arf invariant is one. The product collapse gives
$\eta^2$, as in \cref{maps:src:B.3.1}; the Freudenthal
isomorphism $\pi_6(S^4)\simeq\pi_2^{\mathrm S}$ sends
$\eta_4\eta_5$ to this element. The degree signs of the
three marked transports preserve the order-two class.

\item Take the closed cusp neighbourhood $N$ with its product
outer collar. The boundary-to-neighbourhood map proved above is the
long-period projection to $\Lambda/K=\mathbb Z^2$.
The equivariant-cover construction in the proof of
Theorem~\ref{maps:src:B.3.2}, on this actual free period cover,
gives $h_0:N\to T^2$ inducing that projection, with $h_0(x_0)=0$.
The native long-period map $\alpha$ induces the same homomorphism;
hence $\alpha=h_0+[\ell]$ for a real lift $\ell$ on $N\setminus W$.
Using a radial cut-off $b_c$ equal to zero near $W$ and one on the
outer collar, put
\[
 h=h_0+[b_c\ell],\qquad L_{\rm ext}=(1-b_c)\ell.
\]
Then $h(x_0)=0$, $h=\alpha$ on the outer collar and
$\alpha=h+[L_{\rm ext}]$. These are the operations of the cited
construction in the same real native deck coordinates.

Apply Lemma~\ref{mapsIII:carrier-extension-lemma} with $r=|q|$,
this $\alpha,h$, the lift $L=L_{\rm ext}$ in that proof, and the based target maps
$c_3=\operatorname{collapse}\circ J:T^2\to S^2$ and
$c_\gamma=\operatorname{character}\circ J:T^2\to S^1$.
The matrices $J$ and their displayed inverses carry these formulas
exactly to the two prescribed marked angle-forgetting maps.
The punctured cusp is a torus bundle with connected fibre over a
punctured disc, so it is connected and locally path connected;
its product collar gives the required frontier equality.
All compactness and basepoint hypotheses of the lemma hold.
Its two explicit homotopies therefore give the claimed based
null-homotopies, including their fixed exterior and basepoint.


\end{enumerate}
\end{proof}

\begin{remark}[Two-parameter quotient classes]
\hypertarget{tag-Q015}{}\label{mapsIII:rem:three-family-quotient-class}
{\normalfont\scriptsize tag Q015}\par
\SeriesRef{push:src:E.5.3} supplies distinct two-parameter members
with different sufficiently small $\mu$.
Theorem~\ref{maps:src:B.4.1} gives quotient class $6\nu'$ for both.
\end{remark}

\subsection{A cell presentation from the three fillings}

\begin{definition}[The marked three-filling datum]
\hypertarget{tag-Q016}{}\label{mapsIII:def:three-filling-datum}
{\normalfont\scriptsize tag Q016}\par
Use Definition~\ref{mapsIII:def:marked-family} with $p=-1$,
\[
 (g,u,w,d)=(\widehat\gamma,\widehat u,\widehat w,\widehat\delta),
 \qquad (m_1,m_2)=(3,4),
\]
Thus $v_1=g+2u-4w$, $v_2=-g-3u+3w$, and the cusp
translation drift is zero.
The fibre is $F=\mathbb R^4/\Lambda$.
For a matrix $A$ on $F=\mathbb R^4/\Lambda$, our mapping-torus
convention is
\begin{equation}
 M_A=(F\times[0,1])/((x,1)\sim(A^{-1}x,0)).
 \hypertarget{tag-01V7}{}\label{maps:eq:B.5.1}
\end{equation}
The positive interval loop $\rho$ satisfies
$\rho t_\lambda\rho^{-1}=t_{A\lambda}$.
In the finite boundary chart below its base coordinate is
$q=e^{-2\pi i\theta}$, so this interval runs clockwise.
The counterclockwise loop is $\eta=\rho^{-1}$; it satisfies
$\eta t_\lambda\eta^{-1}=t_{A^{-1}\lambda}$.
The point-return map read from the endpoint relation is
$A^{-1}$ along $\rho$ and $A$ along $\eta$.
The reversal of roof coordinates and its inverse are
\[
 M_A\longrightarrow M_{A^{-1}},\quad
 [x,t]\longmapsto[A^{-1}x,1-t],
 \qquad [y,s]\longmapsto[Ay,1-s].
\]
At the endpoints the quotient relations make both composites
the identity and fix the marked fibre over the base point.
\end{definition}

\Needspace{6\baselineskip}
\begin{theorem}[The three-filling mapping-cylinder model]\hypertarget{tag-02JB}{}\label{maps:src:B.5.1}
{\normalfont\scriptsize tag 02JB}\par


For Definition~\ref{mapsIII:def:three-filling-datum}, the underlying
topological manifold of $X$ is obtained from
the flat $T^4$-bundle over a ribbon bouquet of two circles,
with monodromies $A_1,A_2$, and the three spaces $S_1,S_2,W$
by the three double mapping cylinders specified below.
Their maps admit finite polyhedral cell structures.  They
therefore give a CW structure adapted to the pair of pants
and the three special fibres, with all attaching maps
specified by integral affine formulae and primitive circle
quotients.  Replacing the torus return maps and the attaching maps by
cellular maps on the product torus cell structures gives a
finite CW homotopy model with
cell counts
\begin{equation}
 (c_0,c_1,c_2,c_3,c_4,c_5,c_6)
       =(5,20,45,56,42,17,3).
 \hypertarget{tag-01VH}{}\label{maps:eq:B.5.2}
\end{equation}

\end{theorem}


\begin{proof}
\begin{enumerate}[leftmargin=*]
\item Give a bouquet with oriented loops $\rho_1,\rho_2$ the ribbon
order in which the two ends of each loop are consecutive.
Its oriented thickening is a pair of pants.  Collapse the
widths of its two bands to the edges and the remaining
vertex disc to the vertex.  The induced boundary maps are
piecewise linear paths with words
$\rho_1,\rho_2,(\rho_1\rho_2)^{-1}$; intervals between
successive letters are mapped to the vertex.  The thickening
is the mapping cylinder of this boundary map.  This assertion
can be checked on each band, where the cylinder consists of
the two half-bands, and on the vertex disc, where the cones
on the intervening boundary intervals are its sectors.
Lift this construction to the flat $F$-bundle.  Write $U_*$
for the resulting bundle over the bouquet.  Its boundary
spaces are $E_1=M_{A_1}$, $E_2=M_{A_2}$ and
$E_0=M_{M_0}$. Their maps
$b_j:E_j\to U_*$ and $b_0:E_0\to U_*$ are the
lifted ribbon boundary maps.  They are the identity on a
marked fibre and send their interval loops respectively to
$\rho_1,\rho_2,\rho_2^{-1}\rho_1^{-1}$.
The equality $M_0=(A_1A_2)^{-1}$ verifies the last map
including its fibre monodromy.

\item At a finite fibre put
\begin{equation}
 h_j(x)=A_jx+v_j/m_j,\qquad S_j=F/\langle h_j\rangle,
 \qquad
 r_j([x,\theta])=[x+\theta v_j/m_j].
 \hypertarget{tag-01VI}{}\label{maps:eq:B.5.3}
\end{equation}
The last formula is a well-defined map $M_{A_j}\to S_j$:
at $\theta=1$ its value agrees with $[A_j^{-1}x]$, because
$h_j(A_j^{-1}x)=x+v_j/m_j$.  The finite neighbourhood is
the associated disc bundle
$(F\times D^2)/\langle(x,s)\mapsto(h_jx,e^{-2\pi i/m_j}s)\rangle$.
Its boundary point corresponding to $[x,\theta]$ is
\[
 [x+\theta v_j/m_j,\ e^{-2\pi i\theta/m_j}].
\]
Radial collapse is exactly \eqref{maps:eq:B.5.3}; its mapping cylinder is
the disc bundle.  The bundle projection is $[x,s]\mapsto[x]\in S_j$,
while the holomorphic base coordinate is $q=s^{m_j}$.
Thus this smooth normal-disc model has scheme-theoretic
fibre $m_jS_j$ for the original holomorphic map.  Write
$\sigma_j(s)=\log(s)\Pi v_j/(2\pi i)$.
The logarithmic chart map from the filling to the regular bundle is
$\zeta_N\mapsto\zeta_J=\zeta_N+\sigma_j(s)$;
its inverse is $\zeta_J\mapsto\zeta_N=\zeta_J-\sigma_j(s)$.
For $s=\rho e^{-2\pi i\theta/m_j}$ the latter translation is
$-\log(\rho)\Pi(s)v_j/(2\pi i)+\theta\Pi(s)v_j/m_j$.
It therefore gives exactly the positive drift in the displayed
boundary chart \cite[Lemma~6.5(ii), pp.~38--39]{AC26}.  Subtracting that displayed drift leaves a
periodic translation with a lift to $\mathbb R^4$; writing this periodic lift as $\ell$,
the collar isotopy is $x\mapsto x-t\ell$, with inverse
$x\mapsto x+t\ell$, for $0\le t\le1$.
It identifies the gluing with \eqref{maps:eq:B.5.3} after the
displayed nonzero drift has been kept.
The normal-circle meridian winds $m_j$ times around the base;
this is the kernel meridian in \cite[Lemma~7.4, proof, p.~42]{AC26}.

\item Here is the complete cusp boundary map.  Put
$a=(a_1,a_2)\in\mathbb R^2/\mathbb Z^2$ for the $(g,u)$
coordinates and $b\in T_f=\mathbb R^2/\mathbb Z^2$ for
the $(w,d)$ coordinates.  Let $\mathcal H$ be the hexagonal
honeycomb dual to the $A_2$ triangulation.  For $B_0a$ in
the interior of the hexagon centred at $v\in\mathbb Z^2$,
write $v(a)=v$.  The central fibre has the explicit quotient
description
\begin{equation}
 W_{\mathcal H}=(T_B\times T_f)/\sim,
 \qquad T_B=\mathbb R^2/\mathbb Z^2,
 \hypertarget{tag-01VJ}{}\label{maps:eq:B.5.4}
\end{equation}
where over a honeycomb edge separating centres $v,v'$ one
collapses the primitive circle on $v'-v$, and over each
honeycomb vertex one collapses all of $T_f$.  No circle is
collapsed over a hexagon interior.  The three oriented edge
directions can be chosen as
\begin{equation}
 n_1=w,\qquad n_2=d-w,\qquad n_3=-d.
 \hypertarget{tag-01VK}{}\label{maps:eq:B.5.5}
\end{equation}
The quotient is homeomorphic to the actual $W$: the positive
part of each toric component is its hexagon, the phase
stabilizers are exactly these primitive circles, and the
period action identifies hexagons by $B_0\mathbb Z^2$.
On the universal phase coordinates $\phi=\arg(x)/(2\pi)$,
the constant phase gauge is $b=\phi-\operatorname{Re}C(0)a$,
with inverse $\phi=b+\operatorname{Re}C(0)a$.
It conjugates the constant phase cocycle to the quotient just defined.

\item With the interval convention \eqref{maps:eq:B.5.1}, the cusp attachment is
\begin{equation}
 r_0([a,b,\theta])
   =[a,b-\theta(B_0a-v(a))]\ \in W_{\mathcal H}.
 \hypertarget{tag-01VL}{}\label{maps:eq:B.5.6}
\end{equation}
This formula includes its seam data.  On crossing an edge
the two expressions differ by $\theta(v-v')$, which lies
in the circle collapsed on that edge.  At a vertex every
phase is collapsed.  Replacing $a$ by $a+\lambda$ changes
$v(a)$ by $B_0\lambda$, so the expression is periodic.
At $\theta=1$, the integer vector $v(a)$ disappears modulo
$\mathbb Z^2$, leaving $[a,b-B_0a]$; this is the endpoint
identification by $M_0^{-1}$.  Hence \eqref{maps:eq:B.5.6} is continuous
on the entire mapping torus, not only on its open strata.

\item We justify that the cusp neighbourhood is its actual
mapping cylinder.  Straighten $C(q)$ to $C_0=C(0)$ by
the map on logarithmic coordinates
\begin{equation}
 \xi\longmapsto
 \xi-(C(q)-C_0)(\operatorname{Im}Z_q)^{-1}
                       \operatorname{Im}\xi.
 \hypertarget{tag-01VM}{}\label{maps:eq:B.5.7}
\end{equation}
Writing $\xi=Z_qa+b$ proves that this sends $Z_qa+b$ to
$Z_q^0a+b$, fixes the real lattice, and conjugates the
period actions. Its inverse is
\[
 \xi^0\longmapsto
 \xi^0+(C(q)-C_0)(\operatorname{Im}Z_q^0)^{-1}
                             \operatorname{Im}\xi^0.
\]
Indeed the two imaginary-part formulas recover the identical
real coefficient $a$, and both composites fix $a,b$.  In a bounded-position toric chart the
correction multiplying $x$ tends to one as $q\to0$;
the same holds for its inverse.  Thus it extends by the
identity on $W$ to a homeomorphism of the closed cusp tube.

\item Let $A$ be the nonnegative real part of the frozen toric
cover over $|q|\leq\eta$, and $A_0$ its central part.
The local orthant descriptions make $A$ a three-manifold
with boundary; the corners may be rounded.  Its interior
has coordinates $(y,\rho)\in\mathbb R^2\times(0,\eta)$,
$y=\log|x|/\log\rho$.  The positive period action is
$y\mapsto y+B_\rho^+\lambda$, where
$B_\rho^+=B_0-2\pi\operatorname{Im}C_0/\log\rho$.
For small $\eta$ these matrices are invertible.  Thus the
interior of the positive quotient $Q^+$ is explicitly
$T^2\times(0,\eta)$, using $a=(B_\rho^+)^{-1}y$;
the inverse is $y=B_\rho^+a$ on the universal positive coordinates.
The bounded fundamental parallelogram meets only finitely
many of the toric bounded-position charts, including at
$\rho=0$, so $Q^+$ is compact.  Its two boundary components
are tori.  The inner one is the honeycomb quotient, and
the outer one is the positive regular-fibre torus.

\item Both boundary inclusions induce the identity on the period
deck group $\mathbb Z^2$.  Indeed $A_0$ is the honeycomb
plane and hence contractible; $A$ is contractible because
its interior is contractible and pushing a boundary collar
inward gives a homotopy equivalence with that interior.
The quotient is irreducible.  An embedded locally flat
two-sphere lies in its interior $T^2\times(0,\eta)$.
Lift it to $\mathbb R^2\times\mathbb R$.  Its bounded
ball is disjoint from all its nonzero lattice translates:
disjoint boundary spheres would otherwise give nested
balls, impossible for translates of equal finite volume.
The ball therefore descends.  This also proves
irreducibility of the product target $T^2\times I$.
The boundary tori are incompressible because their
fundamental groups inject.

\item There is a map of pairs $Q^+\to T^2\times I$ inducing
the identity on $\pi_1$ and restricting to the prescribed
homeomorphisms on the two boundaries.  For its first
coordinate, extend the classifying map to $T^2=K(\mathbb
Z^2,1)$ and use the homotopy extension property on the
boundary; use a collar height for the second coordinate.
Apply \cite[Theorem 6.1, p.~77]{mapsWaldhausen68} to this
map.  The source is sufficiently large, since a boundary
torus pushed into its interior is incompressible.  The
product target has the same property, and all the other
hypotheses were just verified.  The homotopy is fixed on
the boundary.  Its alternative of mapping the connected
source into the boundary is impossible, since the two
source boundaries map to distinct target components.
It therefore ends at a covering map.  The induced
fundamental-group isomorphism makes that covering
one-sheeted.  We obtain a product homeomorphism relative
to both prescribed boundary markings.

\item Lift this product to $A$ and spread it by compact phases.
The full period action on a phase $(\phi,\vartheta)$ is
$(\phi,\vartheta)\mapsto
(u_\lambda\phi\vartheta^{B_0\lambda},\vartheta)$;
it does not depend on the positive radial point.  The
lifted product is therefore equivariant for the full
action.  Phase stabilizers occur only on $A_0$ and are
exactly those of \eqref{maps:eq:B.5.4}.  Consequently the resulting
quotient is the mapping cylinder of the boundary phase
map.  On the component with centre $v$, the $q$-phase
acts on $T_f$ by $-v$ times that phase, since $(v,1)$ is
the stabilizer ray.  Substituting the clockwise interval
phase and the marked period $-\theta B_0a$ gives precisely
\eqref{maps:eq:B.5.6}.

\item Gluing the three mapping cylinders of $r_j$ to the
mapping cylinder of the ribbon boundary maps gives
\begin{equation}
 U_*\ \cup\ \coprod_{j=0}^2(E_j\times[0,1])
       \ \cup\ (W\sqcup S_1\sqcup S_2),
 \hypertarget{tag-01VN}{}\label{maps:eq:B.5.8}
\end{equation}
with $(x,0)\sim b_j(x)$ and $(x,1)\sim r_j(x)$.
The preceding constructions identify \eqref{maps:eq:B.5.8}
homeomorphically with $X$.  The comparison need not
preserve the radial value of $f$ inside the cusp, but
it preserves the special strata and all three outer
collars; this is the sense of adaptation used here.

\item For a finite cell construction use rational hexagons,
unit cubes for tori, and the affine quotient \eqref{maps:eq:B.5.3}.
For \eqref{maps:eq:B.5.6} use toric phase coordinates
$b'=b-\theta B_0a$, with inverse $b=b'+\theta B_0a$.
In these coordinates its expression
on a hexagon is $[a,b'+\theta v]$, and its face
identifications are affine:
$a\mapsto a+\lambda$ changes $b'$ by
$-\theta B_0\lambda$.  There are finitely many faces
in the fundamental polyhedra.  Cut their images at the
integer phase hyperplanes and subdivide; this makes
every face map a finite piecewise affine map.  Primitive
edge-circle quotients are the integral linear maps
$b\mapsto\det(n_k,b)\bmod\mathbb Z$.
Triangulate each occurrence of a torus separately and
interpolate through a common subdivision along each
mapping cylinder.  Each occurrence uses its stated subdivision and the common
subdivision on the adjoining cylinder.  Products of simplices with intervals are
triangulated by their ordered vertex prisms.  These
rules supply finite characteristic maps for \eqref{maps:eq:B.5.8},
and pulling the cells back gives the asserted CW
structure of $X$.

\item For the smaller homotopy model put the usual product
cell structure on every torus. First replace each roof map by a
cellular approximation using \cite[Theorem~4.8, p.~349]{Hatcher2002}.
Here is the comparison on each of the actual mapping-torus pieces.
Write $f_0$ for its roof map ($A_j^{-1}$, $M_0^{-1}$, or
$B_j^{-1}$), $f_1$ for the cellular approximation, and
$H(x,s)$ for the homotopy from $f_0$ to $f_1$.
The map from the old mapping torus to the new one fixes its
marked fibre and sends its cylinder to
\[
 [x,t]\longmapsto
 \begin{cases}
 [x,2t],&0\le t\le1/2,\\
 H(x,2-2t),&1/2\le t\le1.
 \end{cases}
\]
At $t=1/2$ both values are $f_1(x)$, and at $t=1$ the
value is $f_0(x)$, so the quotient relations are respected.
The homotopy inverse uses the same formula with
$H(x,1-s)$. For $U_*$ perform this construction on both
branches; their maps agree on the common fibre.
For $W$ use its face-phase cells directly: the edge maps
$b\mapsto\det(n_k,b)$ take the common phase vertex to the
circle vertex and are cellular on the product torus cells.

Denote these object maps by $e_Y:Y\to Y^c$ and
$p_Y:Y^c\to Y$, for $Y=U_*,W,S_1,S_2,E_0,E_1,E_2$;
on $W$ they are the identity. Their composite homotopies
$H_Y:p_Ye_Y\simeq\operatorname{id}_Y$ and
$K_Y:e_Yp_Y\simeq\operatorname{id}_{Y^c}$ are obtained by
contracting each extra path followed by its reverse.
Explicitly, with $\lambda(v)=\min(2v,2-2v)$, the two
contractions are $H(x,(1-r)\lambda(v))$ and
$H(x,1-(1-r)\lambda(v))$, followed by linear
reparametrization of the remaining cylinder interval.
They fix the marked fibre throughout.

For each of the six arrows $f:E\to Y$ (the $b_j$ and $r_j$),
choose a cellular approximation $f^c:E^c\to Y^c$ of
$e_Yfp_E$, with homotopy $L_f$ from $e_Yfp_E$ to $f^c$.
The endpoint paths for the comparison and its inverse are
\[
 A_f(x,s)=
 \begin{cases}
 e_YfH_E(x,1-2s),&0\le s\le1/2,\\
 L_f(e_Ex,2s-1),&1/2\le s\le1,
 \end{cases}
\]
\[
 B_f(y,s)=
 \begin{cases}
 p_YL_f(y,1-2s),&0\le s\le1/2,\\
 H_Y(fp_Ey,2s-1),&1/2\le s\le1.
 \end{cases}
\]
Thus $A_f$ runs from $e_Yf$ to $f^ce_E$, while $B_f$
runs from $p_Yf^c$ to $fp_E$.
On target pieces use $e_Y$; on the $j$th double-cylinder
prism use
\[
 [x,t]\longmapsto
 \begin{cases}
 A_{b_j}(x,3t),&0\le t\le1/3,\\
 [e_{E_j}x,3t-1],&1/3\le t\le2/3,\\
 A_{r_j}(x,3-3t),&2/3\le t\le1.
 \end{cases}
\]
The endpoint identities give a continuous quotient map.
Its homotopy inverse uses $p_Y$ on target pieces,
$p_{E_j}$ on the middle prism, and $B_{b_j},B_{r_j}$
on the two end thirds with the same orientations.
In the composites, the $L_f$ paths retrace and contract
by the displayed $\lambda$ formula. The square
$(r,s)\mapsto H_Y(f(H_E(x,r)),s)$ interchanges the remaining
endpoint paths. The corresponding reversed prism square gives
the $K_Y,K_E$ contraction in the other composite.
The object homotopies and linear reparametrization of each
remaining prism therefore give the identity; they agree on
the marked fibres and target pieces. This proves the claimed
homotopy equivalence with the cellular double-cylinder model.

The two finite surfaces are mapping tori of $T_L$, with
primitive circle coordinate
$s=m_j\gamma/\gamma(v_j)$, where
$\gamma(v_1)=1$ and $\gamma(v_2)=-1$.
Their matrices are $B_j=A_j|_L$ in the roof convention of
Definition~\ref{mapsIII:def:three-filling-datum}.
The cellular mapping-torus models have counts $(1,4,6,4,1)$.
The counts for $U_*^c$ are $(1,6,14,16,9,2)$, and
each $E_j^c$ has counts $(1,5,10,10,5,1)$.
In \eqref{maps:eq:B.5.4} the quotient honeycomb has two vertices,
three edges, and one hexagon. The surviving phase spaces are
respectively a point, a circle, and a two-torus, giving counts
$(2,3,4,2,1)$ for $W$. Adding these cells and one prism cell
for each boundary cell gives \eqref{maps:eq:B.5.2}.
\end{enumerate}
\end{proof}

\begin{theorem}[Homology and framed fibre of the attaching model]
\hypertarget{tag-02JC}{}\label{maps:src:B.5.2}
{\normalfont\scriptsize tag 02JC}\par
The $p=-1$ attaching data of \cref{maps:src:B.5.1}, with
zero cusp translation drift, give $\pi_1(X)=1$ and
\[
 H_k(X;\mathbb Z)=
 \begin{cases}\mathbb Z,&k=0,6,\\0,&\text{otherwise}.
 \end{cases}
\]
They identify $X$ as a homotopy sphere and exhibit the
normally framed three-torus of \cref{maps:src:B.1.1}.
\end{theorem}

\begin{proof}
Write $\delta=d$ for the fourth marked lattice vector in
the attaching data. All homology and cohomology groups
in the proof have integral coefficients.
\begin{enumerate}[leftmargin=*]
\item The core has presentation
\[
 \pi_1(U_*)=
 \langle\Lambda,\rho_1,\rho_2\mid
 [\Lambda,\Lambda]=1,\quad
 \rho_jt_\lambda\rho_j^{-1}=t_{A_j\lambda}\rangle.
\]
The finite boundary map in \eqref{maps:eq:B.5.3} sends
$\rho_j$ to the affine deck generator and imposes
$\rho_j^{m_j}=t_{v_j}$. In the honeycomb quotient,
the phase circles contract towards the vertices,
and the face adds the base-torus relation. Thus
$\pi_1(W)=\mathbb Z\langle g,u\rangle$,
with fibre map the projection to these two
coordinates. The cusp interval maps to a point
at a hexagon centre. The map in
\eqref{maps:eq:B.5.6} therefore kills $w,\delta$ and
$a_0=\rho_2^{-1}\rho_1^{-1}$. Since $A_1w=u-w$,
its normal closure also kills $u$. If $c$ is the
image of $g$, the resulting presentation is
\begin{equation}
 \langle c,x,y\mid c\text{ central},\quad
 xy=1,\quad x^3=c,\quad y^4=c^{-1}\rangle.
 \hypertarget{tag-01VO}{}\label{maps:eq:B.5.9}
\end{equation}
Substitution of $y=x^{-1}$ gives $x^3=x^4$, and
hence $x=c=y=1$.

\item Let $X^\circ$ be the union of the core and the two
finite fillings. Retract the core onto its marked
torus bundle over the bouquet. Its two mapping-torus
branches meet in $F$; the finite attachment maps replace
these branches by $S_1,S_2$. Their restrictions to $F$
are the quotient maps $\pi_j:F\to S_j$. Thus the
resulting homotopy pushout is the double mapping cylinder
$S_1\leftarrow F\to S_2$.

Write $\gamma,u^*,w^*,\delta^*$ for the basis dual to
$g,u,w,\delta$, and put
\begin{equation}
 \eta_1=2u^*+w^*+3\delta^*,\qquad
 \eta_2=u^*+w^*+2\delta^*,\qquad
 q_2=u^*\wedge w^*+6\gamma\wedge\delta^*.
 \hypertarget{tag-01V8}{}\label{maps:eq:B.5.10}
\end{equation}
Abelianizing the finite deck presentations gives
\[
 \begin{array}{c|ccc}
 &g&u&\delta\\\hline
 H_1(S_1)&3\rho_1&2w&3w\\
 H_1(S_2)&-4\rho_2&w&2w.
 \end{array}
\]
In particular $(\rho_1,w)$ and $(\rho_2,w)$ are
free integral bases. In these bases the degree-one
map $(\pi_{1*},-\pi_{2*})$ is
\begin{equation}
 \alpha_1=
 \begin{pmatrix}
 3&0&0&0\\0&2&1&3\\4&0&0&0\\0&-1&-1&-2
 \end{pmatrix}.
 \hypertarget{tag-01V9}{}\label{maps:eq:B.5.11}
\end{equation}
Its rank is three, and two maximal minors have
absolute values three and four, so the gcd of
the maximal minors is one. Its nonzero Smith
factors are therefore $1,1,1$. Its kernel
is generated by $\nu_0=-u-w+\delta$, and its primitive
left kernel is $(4,0,-3,0)$. Therefore
$H_1(X^\circ)=\mathbb Z\zeta$, with
\begin{equation}
 g\longmapsto12\zeta,\qquad
 \rho_1\longmapsto4\zeta,\qquad
 \rho_2\longmapsto-3\zeta,\qquad
 a_0\longmapsto-\zeta.
 \hypertarget{tag-01VA}{}\label{maps:eq:B.5.12}
\end{equation}

\item The finite manifolds are mapping tori of
$L=\mathbb Z\langle u,w,\delta\rangle$. Indeed,
$(v_j,u,w,\delta)$ is an integral basis of $\Lambda$
and $A_jv_j=v_j$. In coordinates
$x=y+s v_j/m_j$, the affine generator acts by
$(y,s)\mapsto(B_jy,s+1)$, where
\[
 B_1=\begin{pmatrix}0&1&0\\-1&-1&0\\1&0&1\end{pmatrix},
 \qquad
 B_2=\begin{pmatrix}0&-1&0\\1&0&0\\0&1&1\end{pmatrix}.
\]
The inverse coordinates on the real covering space are
$s=m_j\gamma(x)/\gamma(v_j)$ and
$y=x-sv_j/m_j$; here $\gamma(v_1)=1$ and
$\gamma(v_2)=-1$. These formulas descend to the
stated mapping tori.

For $r=1,2$, each matrix $\bigwedge^rB_j-I$ has
rank two and nonzero Smith factors $1,1$; for
$r=0,3$ it is zero. The mapping-torus sequences
therefore give free homology groups for $S_j$, of
ranks $(1,2,2,2,1)$. Their degree-two intersection
pairings are unimodular.

Exterior multiplication gives the primitive invariant
lattices
\[
 I_1=\langle\gamma\wedge\eta_1,q_2\rangle,
 \qquad I_2=\langle\gamma\wedge\eta_2,q_2\rangle.
\]
Their intersection determinants have absolute values
$9$ and $4$. Since a covering of degree $m_j$
multiplies the intersection pairing by $m_j$, its
pullback image has determinant of absolute value
$m_j^2$. Its index in $I_1$ is consequently one,
and its index in $I_2$ is two.

The embedded torus generated by $v_2,\delta$ is
preserved by $h_2$ and covers its quotient with
degree four. Evaluation of
$a\gamma\wedge\eta_2+bq_2$ on this torus is
$-2a-6b$. Thus every pullback class has
$a+b\equiv0\pmod2$. This sublattice has index two
in $I_2$, so it is the pullback image. We obtain
\begin{equation}
 \begin{aligned}
 \pi_1^*H^2(S_1)&=\langle\gamma\wedge\eta_1,q_2\rangle,\\
 \pi_2^*H^2(S_2)&=\langle2\gamma\wedge\eta_2,
                              \gamma\wedge\eta_2+q_2\rangle.
 \end{aligned}
 \hypertarget{tag-01VB}{}\label{maps:eq:B.5.13}
\end{equation}

\item Choose cohomology bases with pullbacks
$(\gamma\wedge\eta_1,q_2-2\gamma\wedge\eta_1)$
and $(2\gamma\wedge\eta_2,q_2-3\gamma\wedge\eta_2)$,
and use their Kronecker-dual homology bases.
On $(gu,gw,g\delta,uw,u\delta,w\delta)$, with
exterior products abbreviated, the degree-two map is
\begin{equation}
 \alpha_2=
 \begin{pmatrix}
 2&1&3&0&0&0\\-4&-2&0&1&0&0\\
 -2&-2&-4&0&0&0\\3&3&0&-1&0&0
 \end{pmatrix}.
 \hypertarget{tag-01VC}{}\label{maps:eq:B.5.14}
\end{equation}
Its rank is three. The minors on columns
$(1,2,4)$ and rows $(1,2,3)$, respectively
$(1,2,4)$, are $2$ and $-3$. Thus its nonzero
Smith factors are $1,1,1$, and its
primitive left kernel is $\Phi=(4,2,3,2)$.
Its cokernel is therefore $\mathbb Ze$. The class
$\mathcal A\in H^2(X^\circ)$ represented on the
two finite pieces by $\Phi$ satisfies
$\langle\mathcal A,e\rangle=1$ and
$\mathcal A|_F=2q_2$. The double mapping cylinder
sequence gives
\begin{equation}
 0\longrightarrow\mathbb Ze\longrightarrow
 H_2(X^\circ)\longrightarrow\mathbb Z\nu_0
 \longrightarrow0.
 \hypertarget{tag-01VD}{}\label{maps:eq:B.5.15}
\end{equation}

\item The honeycomb quotient has two vertices, three
oriented edges, and one face. On its one-skeleton
$\partial_1(e_k)=Q-P$. Its face word is
$e_1e_2^{-1}e_3e_1^{-1}e_2e_3^{-1}$.
Over a vertex the phase fibre is a point, over an
edge it is a circle, and over the face it is $T_f$.
Opposite occurrences of each edge have identical
phase quotient maps and opposite orientations.
The face-edge boundary terms therefore cancel in
both phase degrees zero and one; the remaining
phase-torus cellular boundary is zero as well.
It follows that
\[
 H_*(W)=(\mathbb Z,\mathbb Z^2,\mathbb Z^4,
                              \mathbb Z^2,\mathbb Z).
\]

At $\theta=0$ the cusp map is $c(a,b)=[a,b]$.
It induces projection to $\langle g,u\rangle$ in
degree one. In degree two, with target basis
$(T_B,\overline C_1,\overline C_2,\overline C_3)$,
its matrix is
\begin{equation}
 c_2=\begin{pmatrix}
 1&0&0&0&0&0\\0&0&-1&0&0&0\\
 0&-1&-1&1&1&0\\0&0&0&1&0&0
 \end{pmatrix}.
 \hypertarget{tag-01VE}{}\label{maps:eq:B.5.16}
\end{equation}
Here the three phase directions are
$n_1=(1,0)$, $n_2=(-1,1)$, $n_3=(0,-1)$.
For a mixed torus $\beta v$, its $k$th coefficient
is
$\det(n_k,B_0\beta)\det(n_k,v)$.
The first determinants on $(g,u)$ are respectively
$(-1,0),(1,-1),(0,1)$, while the second determinants
are the phase functionals
$\delta^*,-w^*-\delta^*,w^*$.
These give all four mixed columns of
\eqref{maps:eq:B.5.16}. The base torus maps to
$T_B$, and the phase torus $w\delta$ maps to zero.

In degree three, $guw,gu\delta$ map to the two
face-phase generators, and $gw\delta,uw\delta$
map to zero. In degree four the map has degree one
on the face times $T_f$. Thus these maps are
integrally surjective in every degree. Their kernels
have the following primitive bases:
\[
 \begin{array}{c|c}
 r&\ker c_*\subset\bigwedge^r\Lambda\\\hline
 1&\langle w,\delta\rangle\\
 2&\langle w\delta,gw+u\delta\rangle\\
 3&\langle gw\delta,uw\delta\rangle\\
 4&0.
 \end{array}
\]
The identities
$M_0g=g-\delta$, $M_0u=u+w$, $M_0w=w$,
$M_0\delta=\delta$ give the same bases for
$\operatorname{im}(\bigwedge^rM_0-I)$.

\item Put $K_i=\ker(H_i(E_0)\to H_i(W))$.
The mapping-torus sequences and the preceding
surjections give
\begin{equation}
 K_1=\mathbb Za_0,\qquad
 K_2=\mathbb Z\xi_w\oplus\mathbb Z\xi_\delta,
 \qquad \xi_v=[S^1_v\times S^1_{a_0}].
 \hypertarget{tag-01VF}{}\label{maps:eq:B.5.17}
\end{equation}
The tori exist because $M_0$ fixes $w,\delta$.
The mapping-torus projection sends their classes
to the basis $w,\delta$ of $\Lambda^{M_0}$, so
they span a primitive subgroup of $H_2(E_0)$.
At a hexagon centre the cusp formula
\eqref{maps:eq:B.5.6} is independent of the
$a_0$ coordinate. Both tori therefore have zero
image in $H_2(W)$. Since $H_2(E_0)=\mathbb Z^6$
and $H_2(W)=\mathbb Z^4$, they span the full kernel.

\item Let $j:E_0\to X^\circ$ be the other cusp
boundary map. Orient the connecting map in
\eqref{maps:eq:B.5.15} so that
\[
 \partial j_*\xi_v=(I-A_2)v
                  =(I-A_1^{-1})v
 \qquad(v\in\langle w,\delta\rangle).
\]
Cutting the swept torus along the finite band
gives two rectangles with unmatched boundaries
$v$ and $-A_2v$, proving this formula.
Since $(I-A_2)w=-\nu_0$ and
$(I-A_2)\delta=0$, the classes
$e,z=j_*\xi_w$ form a basis of $H_2(X^\circ)$,
and $j_*\xi_\delta$ is a multiple of $e$.

Translation in the $\delta$ direction commutes
with both finite affine actions and extends over
$X^\circ$. Its cohomological sweep operator
restricts on $F$ to contraction by $\delta$.
Since $\mathcal A|_F=2q_2$, this gives
\[
 \iota_\delta(2q_2)=-12\gamma.
\]
By \eqref{maps:eq:B.5.12}, restriction of
$H^1(X^\circ)$ to $H^1(F)$ is injective with
image $\mathbb Z\,12\gamma$. If
$\Theta(\zeta)=1$, it follows that
$\iota_\delta\mathcal A=-\Theta$.
The torus orientation is the $\delta$ direction
followed by the $a_0$ direction. Thus
\[
 \langle\mathcal A,j_*\xi_\delta\rangle
 =\langle-\Theta,a_0\rangle=1,
 \qquad \Theta(a_0)=-1=p.
\]
Consequently
\begin{equation}
 j_*\xi_\delta=-pe=e,\qquad
 [j_*|_{K_2}]_{(\xi_\delta,\xi_w)}^{(e,z)}
       =\begin{pmatrix}1&0\\0&1\end{pmatrix}.
 \hypertarget{tag-01VG}{}\label{maps:eq:B.5.18}
\end{equation}

\item The degree-one cusp attachment map, with
domain basis $(g,u,a_0)$ and target basis
$(\bar g,\bar u,\zeta)$, is
\[
 D_1=\begin{pmatrix}1&0&0\\0&1&0\\-12&0&1\end{pmatrix},
 \qquad
 D_1^{-1}=\begin{pmatrix}1&0&0\\0&1&0\\12&0&1\end{pmatrix}.
\]
In degree two choose an integral section
$\sigma:H_2(W)\to H_2(E_0)$ of the surjection.
The decomposition map and inverse are
\[
 (v,k)\longmapsto\sigma(v)+k,\qquad
 x\longmapsto(r_{0*}x,x-\sigma(r_{0*}x)).
\]
Using $H_2(E_0)=\sigma H_2(W)\oplus K_2$ and the
bases in \eqref{maps:eq:B.5.18}, the map
$(r_{0*},-j_*)$ is
\[
 (v,k)\longmapsto(v,-j_*\sigma(v)-k).
\]
The integral target automorphism
$(v,h)\mapsto(v,-h-j_*\sigma(v))$ is its own inverse and carries
this map to $I_4\oplus I_2$. Thus both attachment
maps are unimodular, and Mayer--Vietoris gives
$H_1(X)=H_2(X)=0$.

The cell count in \eqref{maps:eq:B.5.2} yields
$\chi(X)=5-20+45-56+42-17+3=2$.
As $X$ is a closed oriented six-manifold,
Poincar\'e duality and the universal coefficient
theorem give $H_4(X)=H_5(X)=0$. Moreover
$\operatorname{Ext}(H_3(X),\mathbb Z)$ injects
into $H^4(X)=H_2(X)=0$, so $H_3(X)$ is free.
Its rank is zero by the Euler characteristic.
This proves the asserted integral homology groups.

\item Successive applications of the Hurewicz theorem
give $\pi_i(X)=0$ for $1\leq i\leq5$ and an
isomorphism $\pi_6(X)\to H_6(X)$.
Choose a representative $S^6\to X$ of the class
mapping to the oriented generator of $H_6(X)$.
It is an integral homology isomorphism between
simply connected CW spaces, hence a homotopy
equivalence by \cite[Theorem~4.32 and
Corollary~4.33, pp.~366--368]{Hatcher2002}.

\item Retain a small regular disc in the pair of pants
when taking its ribbon spine. The relative
construction contains the product block
$F\times D^2$. In its marked coordinates the
torus is
\[
 T_L=\{\gamma=\theta_0\}\subset F,
 \qquad L=\mathbb Z\langle u,w,\delta\rangle.
\]
If $(s_1,s_2)$ is a positive real base chart, its
tubular map is
\[
 (u,w,\delta;s_1,s_2,h)
 \longmapsto
 \bigl(\operatorname{base}(s_1,s_2),
              [\theta_0+h,u,w,\delta]\bigr).
\]
The quotient map sends this point to
$(\operatorname{base}(s_1,s_2),[\theta_0+h])$.
Its differential identifies the ordered normal
coordinates with the positive real base frame
followed by the positive $\gamma$ direction.
This is the normal framing of
\cref{maps:src:B.1.1}, with the identity map on
the torus and on its three normal coordinates.
\end{enumerate}
\end{proof}


\subsection{Integral degree-three maps of the finite band}

\begin{definition}[The marked finite-band coverings]
\hypertarget{tag-Q020}{}\label{mapsIII:def:finite-band-coverings}
{\normalfont\scriptsize tag Q020}\par
Use the three-filling datum of
Definition~\ref{mapsIII:def:three-filling-datum} and the actual
quotients $\pi_j:F\to S_j=F/\langle h_j\rangle$ of
\eqref{maps:eq:B.5.3}. On the ordered basis
$(g,u_h,w_h,d)=(\widehat\gamma,\widehat u,\widehat w,\widehat\delta)$
of $\Lambda$, their affine actions are
\[
 h_j(x)=A_jx+v_j/m_j,\qquad (m_1,m_2)=(3,4),
\]
\[
 A_1=\begin{pmatrix}1&0&0&0\\6&0&1&0\\-6&-1&-1&0\\-2&1&0&1\end{pmatrix},
 \quad
 A_2=\begin{pmatrix}1&0&0&0\\0&0&-1&0\\-6&1&0&0\\3&0&1&1\end{pmatrix},
 \quad
 v_1=\begin{pmatrix}1\\2\\-4\\0\end{pmatrix},\quad
 v_2=\begin{pmatrix}-1\\-3\\3\\0\end{pmatrix}.
\]
Write $(\gamma,u,w,\delta)$ for the dual basis of
$V=H^1(F;\mathbb Z)$; thus $u=u^*$, $w=w^*$ and
$\delta=\delta^*$ in the cell-presentation proof.
Orient $F$ by $\operatorname{vol}=\gamma uw\delta$ and $S_j$
by its quotient orientation. Define
\[
 b=\gamma uw,\qquad c_1=uw\delta-4\gamma u\delta-2\gamma w\delta,
 \qquad c_2=uw\delta-3\gamma u\delta-3\gamma w\delta.
\]
\end{definition}

\Needspace{6\baselineskip}
\begin{proposition}[The degree-three finite-band map]
\hypertarget{tag-02QE}{}\label{mapsaudit:finite-degree-three}
{\normalfont\scriptsize tag 02QE}\par
For Definition~\ref{mapsIII:def:finite-band-coverings},
\[
 \pi_1^*H^3(S_1;\mathbb Z)=\mathbb Zb\oplus\mathbb Zc_1,
 \qquad
 \pi_2^*H^3(S_2;\mathbb Z)=\mathbb Z(2b)\oplus\mathbb Zc_2.
\]
Choose the unique cohomology bases having these ordered
pullbacks and the Kronecker-dual bases of $H_3(S_j;\mathbb Z)$.
Relative to the homology basis dual to
$(\gamma uw,\gamma u\delta,\gamma w\delta,uw\delta)$,
the finite-band Mayer--Vietoris map is
\[
 \alpha_3=(\pi_{1*},-\pi_{2*})=
 \begin{pmatrix}1&0&0&0\\0&-4&-2&1\\-2&0&0&0\\0&3&3&-1\end{pmatrix}.
\]
Its nonzero Smith entries are $(1,1,1)$ and its integral
kernel is generated by $(0,1,1,6)^t$. The degree-three
pushforward images have indices one and two, respectively;
the degree-one images have indices three and four. The
corresponding pairs of covariant maps admit no commuting
identification by integral isomorphisms. Their pullback and
covering-transfer maps satisfy
$D_F\pi_j^*=\operatorname{tr}_jD_{S_j}$, where
$\langle\alpha,D_M\xi\rangle=\int_M\alpha\wedge\xi$.
\end{proposition}

\begin{proof}
\begin{enumerate}[label=\arabic*.,leftmargin=*]
\item The quotient actions, their free affine covering
presentations and the free integral homology of $S_j$ are
proved in Theorem~\ref{maps:src:B.5.2} for these identical
matrices and translations. In the same ordered marked lattices,
its bases $(\rho_1,w_h)$ and $(\rho_2,w_h)$ give the degree-one maps
\[
 Q_1=\begin{pmatrix}3&0&0&0\\0&2&1&3\end{pmatrix},\qquad
 Q_2=\begin{pmatrix}-4&0&0&0\\0&1&1&2\end{pmatrix}.
\]
The integral dual bases have pullbacks
$(3\gamma,\eta_1)$ and $(-4\gamma,\eta_2)$, where
$\eta_1=2u+w+3\delta$ and $\eta_2=u+w+2\delta$.
Replacing the second downstairs basis by $(-\rho_2,w_h)$
uses $\operatorname{diag}(-1,1)$, inverse itself, and gives
the positive pullback basis $(4\gamma,\eta_2)$.

For the degree-three basis in the statement, the duality map
$D_F:H^3(F)\to H_1(F)$ has matrix
\[
 J=\begin{pmatrix}0&0&0&1\\0&0&-1&0\\0&1&0&0\\-1&0&0&0\end{pmatrix},
 \qquad J^{-1}=-J.
\]
It sends $b$ to $-d$, $c_1$ to $v_1$ and $c_2$ to $-v_2$.
Change of variables for the orientation-preserving deck map
gives $D_Fh_j^*=A_j^{-1}D_F$. Solving $(A_j-I)x=0$ gives
\[
 x=(a,2a,-4a,e)\quad(j=1),\qquad
 x=(a,3a,-3a,e)\quad(j=2),\qquad a,e\in\mathbb Z.
\]
The inverse parameter maps read the first and fourth
coordinates. Applying $J^{-1}$ proves that the full integral
fixed lattices in degree three are
$I_1=\mathbb Z\langle b,c_1\rangle$ and
$I_2=\mathbb Z\langle b,c_2\rangle$.

Exterior multiplication pairs the positive degree-one bases
above with $(b,c_j)$ by
\[
 K_1=\begin{pmatrix}0&3\\-3&0\end{pmatrix},\qquad
 K_2=\begin{pmatrix}0&4\\-2&0\end{pmatrix}.
\]
The downstairs integral pairing is unimodular because these
oriented closed four-manifolds have free integral homology.
Pullback multiplies the pairing by $m_j$, as follows by
lifting an oriented fundamental triangulation. For the
coordinate matrix $C_j$ of a basis of the pullback image
in $I_j$, this says $K_jC_j=m_jU_j$ with $U_j$ unimodular.
The image has full rank: covering transfer gives
$\operatorname{tr}\pi_j^*=m_j$ on cohomology over
$\mathbb Q$, and all the groups used here are free.
Thus
\[
 [I_j:\pi_j^*H^3(S_j;\mathbb Z)]
 =|\det C_j|=m_j^2/|\det K_j|=1,2.
\]
For a descended class $rb+sc_2$, pairing with $\eta_2$
gives $-2r\in4\mathbb Z$. Hence its $b$ coefficient is
even. The lattice $\mathbb Z(2b)\oplus\mathbb Zc_2$
already has index two in $I_2$; equality of the indices
proves equality of the full sublattices. This gives both
pullback formulas and their sufficiency. Pullback is
injective over $\mathbb Z$, because its kernel is killed
by the degree and its domain is free. Consequently the
ordered downstairs bases in the statement exist uniquely.

\item Kronecker naturality gives
$\langle a,\pi_{j*}z\rangle=\langle\pi_j^*a,z\rangle$.
Thus each degree-three pushforward is the transpose of
its displayed pullback matrix, and the second summand of
the Mayer--Vietoris difference has the stated minus sign.
Row three of $\alpha_3$ is minus twice row one, while the
minor on rows $(1,2,4)$ and columns $(1,2,4)$ is one.
Its rank is three and all three positive Smith entries are
one. The kernel equations give $z_1=0$, $z_3=z_2$ and
$z_4=6z_2$, with $z_2$ an arbitrary integer.
The first pushforward sends the first and fourth basis
cycles to $(1,0)$ and $(0,1)$. The second sends them to
$(2,0)$ and $(0,1)$, and its first coordinate is always
even. These are the full images and prove the two indices.

The degree-one matrices $Q_j$ have images of indices
three and four. Their cokernels have orders three and
four; the corresponding degree-three cokernels have
orders one and two. A commuting square of these maps
with integral isomorphisms on source and target induces
an isomorphism of the cokernels, contradicting these
orders in each case. This is the stated obstruction to
that precise commuting identification.
The relation between degrees is instead given by the
actual duality and transfer maps: lifting every simplex
of an oriented fundamental chain gives
\[
 \begin{aligned}
 D_F\pi_j^*&=\operatorname{tr}_jD_{S_j}
       &&:H^3(S_j;\mathbb Z)\longrightarrow H_1(F;\mathbb Z),\\
 (\pi_j)_*\operatorname{tr}_j&=m_j\operatorname{id}
       &&\text{on }H_1(S_j;\mathbb Z),\\
 \operatorname{tr}_j(\pi_j)_*&=\sum_{k=0}^{m_j-1}A_j^k
       &&\text{on }H_1(F;\mathbb Z).
 \end{aligned}
\]
These identities, together with the Kronecker transpose
above, describe the maps in the fixed integral markings.
\end{enumerate}
\end{proof}

\section{Local Kummer maps and cyclic cover invariants}
\subsection{The local Kummer cover and its global obstruction}

\begin{definition}[The marked cusp and its root covers]
\hypertarget{tag-Q017}{}\label{mapsIII:def:kummer-data}
{\normalfont\scriptsize tag Q017}\par
Use the original family of Definition~\ref{mapsIII:def:marked-family}
and put $V_{\rm tor}=\mathbb Z\langle\widehat w,\widehat\delta\rangle$.
The cusp period matrix, in the ordered bases
$(\widehat\gamma,\widehat u)$ and $(\widehat w,\widehat\delta)$, is
\[
 Z(s)=sB_0+C(q),\qquad q=e^{2\pi i s},\qquad
 B_0=\begin{pmatrix}0&1\\-1&0\end{pmatrix},
\]
with the holomorphic correction of
Definition~\ref{mapsIII:def:cusp-suspension-maps}.
The fan in $N'=\mathbb Z^2\oplus\mathbb Z$ consists of the
cones over the height-one triangles
\[
 \{u,u+e_1,u+e_2\},\qquad
 \{u+e_1,u+e_2,u+e_1+e_2\},\qquad u\in\mathbb Z^2,
\]
and their faces; $e_1,e_2$ are the coordinates corresponding to
$\widehat w,\widehat\delta$.
For $N\ge1$, the base root is $q=v^N$ and the source correction is
$C_N(v)=C(v^N)/N$. Choose a sufficiently small root radius $\epsilon_N>0$
and represent the germ on the compatible discs $|v|<\epsilon_N$ and
$|q|<\epsilon_N^N$;
the source disc is the full inverse image of the target disc.
The source period cover uses that same fan.
In group presentations, $a_1,a_2$ denote the clockwise based loops
whose conjugation matrices are $A_1,A_2$, and
$a_0=(a_1a_2)^{-1}$. Their inverse loops and point-return maps are
those of Definition~\ref{mapsIII:def:three-filling-datum}.
\end{definition}

\begin{theorem}[The local Kummer construction]\hypertarget{tag-02K2}{}\label{maps:src:C.1.1}
{\normalfont\scriptsize tag 02K2}\par


For $N\geq1$ and the data of
Definition~\ref{mapsIII:def:kummer-data}, the map
\[
 P_N(\xi_1,\xi_2,v)=(\xi_1^N,\xi_2^N,v^N)
\]
descends to a finite flat holomorphic map
$\kappa_N:N_0^{\langle N\rangle}\to N_0$ of degree $N^3$ between
smooth threefolds. The source central fibre is reduced and irreducible,
with normalization $\mathrm{dP}_6$. Its ramification divisor is
$(N-1)W_N$, and $\kappa_N^*W=NW_N$.

The base change $N_0[N]=N_0\times_{\Delta_q}\Delta_v$ is already normal.
The induced finite map $Q_N:N_0^{\langle N\rangle}\to N_0[N]$ has
degree $N^2$. For $N>1$ it is quasi-\'etale, and its nonflat and non-\'etale
loci are precisely the inverse images of the double and triple strata.
At $N=1$, $Q_1$ is the identity. For $N>1$, $Q_N$ is not a normalization morphism. On a regular fibre its image lattice is
\begin{equation}
 L_N=\Z\langle\widehat\gamma,\widehat u,
                   N\widehat w,N\widehat\delta\rangle.
 \hypertarget{tag-01XE}{}\label{maps:eq:C.1.1}
\end{equation}

\end{theorem}


\begin{proof}
\begin{enumerate}[leftmargin=*]
\item On the target torus the period action is
$\Psi_\lambda(x,q)=(e^{2\pi iC(q)\lambda}q^{B_0\lambda}x,q)$,
where $B_0=\left(\begin{smallmatrix}0&1\\-1&0\end{smallmatrix}\right)$.
The source uses the same fan and the action
\[
 \Psi^N_\lambda(\xi,v)=
 \big(e^{2\pi iC(v^N)\lambda/N}v^{B_0\lambda}\xi,v\big).
\]
The integral shear preserves the $A_2$ triangulation and the
holomorphic unit factor acts in every toric chart. Additivity in $\lambda$
proves the action law. Raising the two torus coordinates to the $N$th
power proves $P_N\Psi^N_\lambda=\Psi_\lambda P_N$ exactly.
The specified holomorphic function $2\pi iC(v^N)\lambda/N$
defines the exponential factor on the whole root disc.

\item In a maximal smooth toric chart the map is
$(w_0,w_1,w_2)\mapsto(w_0^N,w_1^N,w_2^N)$. Its convergent power-series
ring is free over the target ring on the $N^3$ monomials whose three
exponents lie in $[0,N-1]$. Thus $P_N$ is finite flat. These charts cover
the fan, and their inverse images are the corresponding charts because
scalar multiplication preserves every cone. Freeness and proper
discontinuity of the source lattice action now follow directly from the
target action: a source fixed point would map to a target fixed point;
if two source compact sets meet under a translate, their target compact
images meet under the same translate. Only finitely many such translates
exist. Hence the source action is free and properly discontinuous.

\item Over an evenly covered target open set choose one lift to the toric
cover. Every source lattice orbit over it has a unique representative
whose image is in that lift. Hence the descended morphism is locally the
same finite flat map, of degree $N^3$. It is proper. Since $f_0$ is proper,
the source map to $\Delta_v$ is proper as well: the inverse image of a
compact $K\subset\Delta_v$ is closed in
$\kappa_N^{-1}(f_0^{-1}(\{v^N:v\in K\}))$, a compact set.
The free action on a smooth toric space gives a smooth source. The action sends the component indexed by $u$ to the one indexed by
$u+B_0\lambda$. Since $B_0$ is unimodular, these components form one orbit
with trivial component stabilizer. The component over $(0,0,1)$ has the
complete star fan with rays
$\pm e_1,\pm e_2,\pm(e_1-e_2)$; it is the blow-up of $\mathbb P^2$ at
its three torus fixed points, hence $\mathrm{dP}_6$
\cite[Propositions~4.28 and~4.45, pp.~49, 53]{TelenToric2022}.
Its map onto the quotient central fibre is finite and birational:
proper discontinuity gives finite fibres, and its open torus maps
isomorphically onto the open central stratum. It is therefore the
normalization. The central fibre is reduced by the local product equation. In this chart
$q=z_0z_1z_2$, $v=w_0w_1w_2$, and the Jacobian is
$N^3(w_0w_1w_2)^{N-1}$, proving the asserted divisors.

\item At a triple point the relative factor has the exact ring map
\begin{equation}
 \frac{\C\{z_0,z_1,z_2,v\}}{(z_0z_1z_2-v^N)}
 \longrightarrow\C\{w_0,w_1,w_2\},\qquad
 z_i\longmapsto w_i^N,\quad v\longmapsto w_0w_1w_2.
 \hypertarget{tag-01XF}{}\label{maps:eq:C.1.2}
\end{equation}
The source ring on the left is the invariant ring for
$G=\{(\epsilon_0,\epsilon_1,\epsilon_2)\in\mu_N^3:
\epsilon_0\epsilon_1\epsilon_2=1\}$ on the right. Indeed an invariant
monomial has three exponents congruent modulo $N$; factor out their
common residue as a power of $w_0w_1w_2$. The same expansion for convergent
series proves the analytic assertion. A finite-group invariant ring of a
normal domain is normal: an element of its fraction field integral over
the invariants is integral over the original ring, belongs to that ring
by normality, and is invariant. This proves normality here and on the
face charts, and the degree is $|G|=N^2$.

\item The stabilizer is trivial when at most one $w_i$ vanishes, has order $N$
when exactly two vanish, and is $G$ at the origin. The origin fibre has
length
\[
 N^3-(N-1)^3=3N^2-3N+1>N^2\quad(N>1).
\]
At a double point the two-variable factor has generic rank $N$ and
closed fibre length $2N-1>N$. These counts follow by excluding monomials
divisible by the product of the variables from the box of exponents
$0,\ldots,N-1$. They prove failure of flatness exactly where claimed;
elsewhere the free finite action gives \'etaleness and flatness.

\item Put $v=e^{2\pi ir}$, $q=v^N$, and $\Omega=Z(Nr)$.
The source period matrix is
$Z_N(r)=rB_0+C(v^N)/N=\Omega/N$.
Thus the regular-fibre map induced by $Q_N$ is
\[
 F_N:\frac{\C^2}{\Z^2+(\Omega/N)\Z^2}
 \longrightarrow\frac{\C^2}{\Z^2+\Omega\Z^2},\qquad
 [\zeta]\longmapsto[N\zeta].
\]
Write $L=N\Z^2+\Omega\Z^2$ and $\Lambda=\Z^2+\Omega\Z^2$.
The coordinate map
\[
 H_N:\frac{\C^2}{\Z^2+(\Omega/N)\Z^2}
 \longrightarrow\frac{\C^2}{L},\qquad
 [\zeta]\longmapsto[N\zeta]
\]
is an isomorphism with inverse $[\eta]\mapsto[\eta/N]$.
Under this isomorphism, $F_N$ is the identity-induced quotient map
$\C^2/L\to\C^2/\Lambda$, $[\eta]\mapsto[\eta]$.
The image lattice is therefore exactly \eqref{maps:eq:C.1.1}, with
index $[\Lambda:L]=N^2$. Multiplication by $N$ on lifted endpoints
identifies $\ker F_N=N^{-1}\Lambda/(\Z^2+(\Omega/N)\Z^2)$ with
$\Lambda/L\simeq(\Z/N)^2$, the group of fibre deck translations.
For $N=1$, the corrections and actions are identical and
\eqref{maps:eq:C.1.2} is the identity after eliminating $v$, so $Q_1$
is the identity. For $N>1$, the generic degree $N^2>1$ makes $Q_N$
nonbirational, so it is not a normalization morphism.

In the ordered marking, the period shift $s\mapsto s+1$ gives
\[
 M_0=\begin{pmatrix}I_2&0\\ B_0&I_2\end{pmatrix},\qquad
 S_N=\operatorname{diag}(I_2,NI_2).
\]
Here $R=M_0-I_4=\left(\begin{smallmatrix}0&0\\ B_0&0\end{smallmatrix}\right)$
satisfies $R^2=0$, so
$M_0^N=\left(\begin{smallmatrix}I_2&0\\ NB_0&I_2\end{smallmatrix}\right)$
and
\[
 S_N^{-1}M_0^NS_N
 =\begin{pmatrix}I_2&0\\ B_0&I_2\end{pmatrix}=M_0.
\]
This is the compatibility of the marked fibre lattice map with source
monodromy and the target monodromy pulled back by $q=v^N$.
\end{enumerate}
\end{proof}

\Needspace{6\baselineskip}
\begin{lemma}[Smooth covers of the entire cusp fan]\hypertarget{tag-02K3}{}\label{maps:src:C.1.2}
{\normalfont\scriptsize tag 02K3}\par


Let $\Sigma$ be the complete periodic $A_2$ triangulation at height one
used for the cusp, in $N'=\Z^2\oplus\Z$. Suppose a finite connected
holomorphic cover of its toric neighbourhood is \'etale over
$(\C^*)^2\times\Delta^*$ and has smooth total space. The cocharacter
sublattice defining its restriction to this torus is $nN'$ for some
$n\geq1$. In particular its intersection with the two-dimensional
vanishing lattice is $n\Z^2$.

\end{lemma}


\begin{proof}
\begin{enumerate}[leftmargin=*]
\item Write $f:Y\to X$ for the given cover, $U=(\C^*)^2\times\Delta^*$
and $D=X\setminus U$. The smooth source is normal, and $f^{-1}(D)$ is
closed analytic and nowhere dense. Hence $Y^\circ=f^{-1}(U)$ is connected
\cite[Expos\'e~XII, Lemma~2.5, p.~244]{SGA1Recomposed2004}.
The coordinate exponential universal cover of $U$ has deck lattice
$N'=\Z^3$, with $(a,b,k)$ represented by the loop
\[
 s\longmapsto
 (x_*e^{2\pi ias},y_*e^{2\pi ibs},t_*e^{2\pi iks}).
\]
Thus the connected finite cover $Y^\circ\to U$ corresponds to a
finite-index sublattice $K\subset N'$
\cite[Propositions~1.36--1.37 and Theorem~1.38, pp.~66--68]{Hatcher2002}.
The covering isomorphism is holomorphic on local inverse charts.

Put $M'=\operatorname{Hom}(N',\Z)$ and
$M_K=\operatorname{Hom}(K,\Z)\subset M'_{\mathbb Q}$.
For every maximal cone $\sigma$, the monomial inclusion
\[
 \C[\sigma^\vee\cap M']\longrightarrow
 \C[\sigma^\vee\cap M_K]
\]
extends that torus cover. Choose $r\geq1$ with $rN'\subset K$.
Fractional residues in the unit cube of a target dual ray basis give
finitely many module generators for the second ring over the first.
If $m_1,m_2,m_3$ is that dual basis, the second ring is also the invariant
ring of $\C[z_1,z_2,z_3]$, $z_i=\chi^{m_i/r}$, for the action
\[
 h\cdot z_i=
 \exp\!\left(2\pi i\langle m_i,h\rangle/r\right)z_i,
 \qquad h\in K/rN'.
\]
This invariant ring is normal: an element of its fraction field
integral over it is integral over the polynomial ring, hence belongs
to that ring, and invariance puts it in the invariant subring. Its
fraction field is $\C(M_K)$, so it is precisely the integral closure
of $\C[\sigma^\vee\cap M']$ in the finite torus field extension.
The inclusions commute with the monomial localizations defining face
charts, and therefore glue.

Analytification of each finite-type cone chart preserves normality and
finiteness \cite[Expos\'e~XII, Propositions~2.1(vi) and~3.2(vi),
pp.~242, 245--247]{SGA1Recomposed2004}.
The resulting finite normal cover $Z_K\to X$ has the same real cones
and lattice $K$, and its restriction over $U$ is isomorphic to
$Y^\circ\to U$. That isomorphism and its inverse extend over $X$ by
\cite[Expos\'e~XII, Proposition~5.3, pp.~254--255]{SGA1Recomposed2004}.
Uniqueness makes their compositions the identity, so $Y\cong Z_K$
over $X$. All constructions are local on the cone charts and apply
to the entire periodic fan.

\item For a target ray $r_v=(v,1)$ let $a_vr_v$ be its primitive generator in
$K$. In the prescribed fan choose
\[
 r_0=(0,0,1),\qquad r_1=(1,0,1),\qquad r_2=(0,1,1).
\]
The opposite rays across its three edges are
\[
 \begin{aligned}
 r_1+r_2-r_0&=(1,1,1),\\
 r_0+r_2-r_1&=(-1,1,1),\\
 r_0+r_1-r_2&=(1,-1,1).
 \end{aligned}
\]
They belong to the adjacent defining triangles translated by
$0,-e_1,-e_2$, respectively.
A smooth source chart over a triangle with vertices $v_0,v_1,v_2$
requires $a_0r_0,a_1r_1,a_2r_2$ to be a basis of $K$
\cite[Theorem~2.57, p.~19]{TelenToric2022}.
Indeed, at the fixed point the indecomposable elements of
$\sigma^\vee\cap M_K$ give a basis of the maximal ideal modulo its square.
Smoothness in dimension three makes their number three; they generate
$M_K$ and hence form its basis. The dual basis consists of the three
primitive ray generators in $K$. Algebraic and analytic regularity agree
by \cite[Expos\'e~XII, Proposition~2.1(iv), p.~242]{SGA1Recomposed2004}.
In the adjacent triangle across the edge $(r_1,r_2)$ the other vertex is
$r_3=r_1+r_2-r_0$. Both target triples are integral bases. Equality of
their indices gives
\[
 [N':K]=a_0a_1a_2=a_3a_1a_2,
 \qquad a_3=a_0.
\]
Expressing $a_3r_3$ in the first source basis gives coefficients
$-1,a_0/a_1,a_0/a_2$. These are integers. Applying the same argument
across the other two edges of the first triangle gives all the reverse
divisibilities. Thus $a_0=a_1=a_2=n$, and their three vectors generate
$K=nN'$. This already proves the result; connectivity of the triangles
also shows directly that every ray has ramification index $n$.

\item The local role of the neighbouring cones has an explicit example.
In the displayed initial basis take
$K=\mathbb Zr_0+\mathbb Zr_1+\mathbb Z(2r_2)
=\{(x,y,z):y\in2\mathbb Z\}$.
Its torus map is $(s_1,s_2,q)\mapsto(s_1,s_2^2,q)$.
The initial source chart has primitive basis $r_0,r_1,2r_2$
and map $(z_0,z_1,z_2)\mapsto(z_0,z_1,z_2^2)$.
Across the edge $(r_1,r_2)$, the target basis is
$r_3=(1,1,1),r_1,r_2$, with dual basis
\[
 e'_0=(1,1,-1),\quad e'_1=(0,-1,1),\quad e'_2=(-1,0,1).
\]
Since $(0,1/2,0)=(e'_0+e'_2)/2$, the corresponding normal
source chart has ring
\[
 \mathbb C\{x_0,x_1,x_2,B\}/(B^2-x_0x_2),
 \qquad B=\chi^{(e'_0+e'_2)/2},
\]
and its finite map forgets $B$. It is free of rank two over
the target ring, with basis $1,B$. At its origin the cotangent
space has dimension four while the chart has dimension three.
Equivalently, the primitive source rays $2r_3,r_1,2r_2$ span
an index-two sublattice of $K$. Thus the same entire-fan cover
has the displayed smooth chart and singular adjacent chart.
\end{enumerate}
\end{proof}

\begin{theorem}[The lattice obstruction to a smooth global extension]\hypertarget{tag-02K4}{}\label{maps:src:C.1.3}
{\normalfont\scriptsize tag 02K4}\par


For $N>1$, the local map $\kappa_N$ in Theorem \ref{maps:src:C.1.1} cannot be the
restriction, even to one component of the inverse image of a cusp
neighbourhood, of a finite holomorphic map $Y\to X$ from a connected
smooth compact complex manifold. For $N=1$ the
displayed local map is the identity and has the identity global extension.

\end{theorem}


\begin{proof}
\begin{enumerate}[leftmargin=*]
\item Suppose such a finite map $F:Y\to X$ exists. Form the Stein factorization
of the proper map $fF$, with connected fibres $g:Y\to C'$ and finite map
$h:C'\to B$ \cite[\S1, p.~458]{mapsLi1995Cycles}. Near the specified cusp point, the fibres of $g$ are tori
and their maps to the corresponding fibres of $f$ are \'etale isogenies,
with image lattice $L_N$ of \eqref{maps:eq:C.1.1}. No component of the ramification
divisor of $F$ can dominate $C'$: a proper dominating component would
meet every nearby regular fibre in this open set where $F$ is \'etale.
Thus all ramification is vertical. Delete its finitely many base values,
the branch values of $h$, and $0,1,\infty$. Over the remaining connected
base the map is a cover of torus bundles. On every sheet the image
of a fibre fundamental group is a finite-index sublattice of $\Lambda$,
and continuation along a loop with point-return matrix $A$ carries it
to $AL_N$. Indeed the actual bundle transports satisfy
$F_1\widetilde H=HF_0$, hence
$F_{1*}\widetilde H_*=H_*F_{0*}$; the inverse loop uses
$H^{-1},\widetilde H^{-1}$ and recovers the original subgroup.
Translations induce the identity on the abelian fibre group by a path
from the translation point to zero.

\item Follow the counterclockwise based meridian $a_2^{-1}$,
whose point-return map is $A_2$ by
Definition~\ref{mapsIII:def:three-filling-datum}, choosing it to avoid
the additional deleted points, and then return to the same small
punctured cusp disc. The resulting sheet has image lattice $A_2L_N$.
It lies in some component of $F^{-1}(N_0)$; it need not be the original
component. Pull that component back to the unramified toric covering
$Y_\Sigma\to N_0$. This is still a finite cover with smooth total space,
\'etale off the central fibre. Each connected component therefore satisfies
Lemma \ref{maps:src:C.1.2}. Consequently its fibre image lattice must meet
$V_{\rm tor}$ in a scalar sublattice $nV_{\rm tor}$.

\item The actual second monodromy gives
\[
 A_2\widehat u=\widehat w,\qquad
 A_2\widehat w=-\widehat u+\widehat\delta,\qquad
 A_2\widehat\delta=\widehat\delta.
\]
The first coordinate of $A_2\widehat\gamma$ is one. Taking the
intersection of $A_2L_N$ with $V_{\rm tor}$ therefore gives exactly
\begin{equation}
 A_2L_N\cap V_{\rm tor}
   =\Z\widehat w\oplus N\Z\widehat\delta.
 \hypertarget{tag-01XG}{}\label{maps:eq:C.1.3}
\end{equation}
It cannot equal $nV_{\rm tor}$ when $N>1$: the primitive first generator
forces $n=1$ but the second generator then excludes
$\widehat\delta$. This contradiction proves the assertion for $N>1$.
For the reverse based meridian one has
\[
 A_2^{-1}\widehat u=-\widehat w+\widehat\delta,\qquad
 A_2^{-1}\widehat w=\widehat u,\qquad
 A_2^{-1}\widehat\delta=\widehat\delta,
\]
and hence
\[
 (A_2^{-1}L_N)\cap V_{\rm tor}
 =\mathbb Z(-\widehat w+\widehat\delta)
       +N\mathbb Z\widehat\delta.
\]
The integral automorphism
$S(\widehat w)=-\widehat w+\widehat\delta$,
$S(\widehat\delta)=\widehat\delta$ satisfies $S^2=I$ and
carries \eqref{maps:eq:C.1.3} to this intersection while preserving
$nV_{\rm tor}$ for every $n$. Thus the scalar-lattice obstruction
also holds for the reverse meridian.
For $N=1$, the identity $\kappa_1=\operatorname{id}_{N_0}$ extends
as $\operatorname{id}_X$.
\end{enumerate}
\end{proof}

\Needspace{6\baselineskip}
\begin{proposition}[The monodromy-orbit degree obstruction]\hypertarget{tag-02K5}{}\label{maps:src:C.1.4}
{\normalfont\scriptsize tag 02K5}\par


For $N>1$, the local map $\kappa_N$ has no finite global extension
from a normal compact complex source whose whole inverse image of
$N_0$ is that single prescribed map. The monodromy orbit of $L_N$ has at least
$N^2$ elements, whereas such an extension would have a Stein base cover
of degree $N$.

\end{proposition}


\begin{proof}
Put $D=A_2^2A_1A_2^2A_1^{-1}$. Exact multiplication gives
\[
 D=\begin{pmatrix}1&0&0&0\\0&1&0&0\\6&0&1&0\\-6&1&0&1\end{pmatrix}.
\]
The $N^2$ lattices $D^aM_0^bL_N$, $0\leq a,b<N$, are distinct.
Indeed their intersections with $\ker\gamma$, reduced modulo $N\Lambda$,
are generated respectively by
$\widehat u+b\widehat w+a\widehat\delta$; the coefficient of
$\widehat u$ is one, so equality of these cyclic submodules forces
equality of both residues $a,b$.

Let $F:Z\to X$ be such an extension and take the connected
normal component containing the prescribed cusp source. Its finite
image contains the cusp neighbourhood and hence is all of $X$.
The proper map $fF$ has Stein factorization
\[
 Z\xrightarrow{g}C'\xrightarrow{h}B
\]
with connected fibres and finite $h$
\cite[\S1, p.~458]{mapsLi1995Cycles}.
The whole prescribed cusp preimage has base coordinate $q=v^N$;
over a small nonzero $q$ it has exactly $N$ connected torus fibres.
Thus $\deg h=N$.

The closed analytic set
$E=Z_{\rm sing}\cup\operatorname{Supp}\Omega_{Z/X}$ misses
the entire preimage of the punctured cusp disc. Its image under
the proper map $fF$ is a proper analytic subset of $B$, hence finite
\cite[Chapter~II, Theorem~8.8, p.~118]{mapsDemailly2012CADG}.
Delete that image, the three original exceptional values and the
finitely many branch values of $h$. Over the remaining base,
$F$ is a finite cover of torus bundles. A loop can be perturbed
away from the additional deleted points through a homotopy in the
original ordinary base, preserving its original monodromy.

For a sheet $y\in C'$ over the chosen base point put
$\mathcal L_y=(F|_{g^{-1}(y)})_*\pi_1(g^{-1}(y))\subset\Lambda$.
Fibre basepoint changes act trivially on this subgroup because
$\Lambda$ is abelian. For an actual bundle transport $H$ and its
lift $\widetilde H$, the equality
\[
 (F|_{g^{-1}(y')})_*\widetilde H_*
 =H_*(F|_{g^{-1}(y)})_*
\]
gives $\mathcal L_{y'}=H_*\mathcal L_y$.
The two inverse based meridians have point-return matrices
$A_1,A_2$. Replacing generators by their inverses leaves the
same generated subgroup, since each original generator is the
inverse of its replacement. Therefore the whole monodromy orbit
of $L_N$ is the image of the $N$ sheets under $y\mapsto\mathcal L_y$.
The already proved $N^2$ distinct lattices give $N^2\le N$,
contradicting $N>1$.

\end{proof}

\begin{theorem}[Normal global Kummer covers]\hypertarget{tag-02K6}{}\label{maps:src:C.1.5}
{\normalfont\scriptsize tag 02K6}\par


For every $N\geq1$, the prescribed local Kummer cover extends as one
component above the cusp of a connected finite flat cover with normal
compact complex source, branched only over
$S_1\cup S_2\cup W$, of degree
\begin{equation}
 N^2|G_N|,\qquad
 G_N=\langle(A_1,-1),(A_2,0)\rangle
 \subset\operatorname{GL}_4(\Z/N)\times\Z/N.
 \hypertarget{tag-01XH}{}\label{maps:eq:C.1.4}
\end{equation}
Here matrices and the second coordinates are reduced modulo $N$.
The component in question has degree $N^3$ over $N_0$ and is precisely
$\kappa_N$. For $N>1$ the singular locus of the source meets another
component above the cusp.

\end{theorem}


\begin{proof}
\begin{enumerate}[leftmargin=*]
\item Remove the three special fibres and write the result as $X^*\to B^*$.
The base retracts to a graph. We choose a continuous section that agrees
with the zero section of the marked torus family on a cusp collar.
Such a choice exists: the only obstruction is in
$H^2(B^*,\text{cusp collar};\Lambda)$, which is zero since the based
cusp loop is a free factor in the fundamental group of the punctured
sphere. Explicitly the free-basis map
$(a_1,a_2)\mapsto(a_0,a_1)$ has inverse
$a_2=a_1^{-1}a_0^{-1}$. A pair-of-pants spine containing the
cusp loop gives the relative graph retraction, along which the
specified section extends. The resulting group presentation is
\[
 \pi_1(X^*)=\Lambda\rtimes F(a_1,a_2),\qquad
 a_0=(a_1a_2)^{-1},\quad a_j\lambda a_j^{-1}=A_j\lambda.
\]
The restriction to the cusp collar is the specified splitting with
monodromy $M_0$. A topological finite covering of a complex manifold
has a unique lifted holomorphic atlas, so it is enough first to
construct its subgroup.

\item Map the free group to $G_N$ by the two generators in \eqref{maps:eq:C.1.4} and
map the fibre lattice to $(\Z/N)^4$. This gives an epimorphism
\[
 \pi_1(X^*)\longrightarrow
 K_N=(\Z/N)^4\rtimes G_N,
\]
where only the matrix coordinate acts on the first factor. Let
$U_N\subset(\Z/N)^4$ be the subgroup spanned by $g,u$, of order
$N^2$, and take the connected cover corresponding to the inverse
image $H_N$ of $U_N$. If $\varphi:\pi_1(X^*)\to K_N$ is
the displayed epimorphism, the coset bijection is
\[
 H_N\backslash\pi_1(X^*)\longrightarrow U_N\backslash K_N,
 \qquad H_N\gamma\longmapsto U_N\varphi(\gamma).
\]
Its inverse chooses any lift of a representative in $K_N$;
the relation $H_N=\varphi^{-1}(U_N)$ makes that inverse
independent of the lift. Thus the degree is
$|K_N|/|U_N|=N^2|G_N|$.
An element $\lambda a_0^k$ of the cusp group lies in this subgroup
exactly when $k\equiv0\pmod N$ and
$\lambda\in L_N$. Indeed $a_0$ has second coordinate $1$, and
$M_0^N\equiv I\pmod N$ because $(M_0-I)^2=0$.
Consequently its cusp subgroup is exactly
$L_N\rtimes\langle a_0^N\rangle$. On universal logarithmic coordinates the power map of
Theorem~\ref{maps:src:C.1.1} is
$(s_v,\zeta_v)\mapsto(Ns_v,N\zeta_v)$.
The identity $NZ_N(s_v)=Z(Ns_v)$ sends the first two period
generators to themselves, the last two to their $N$-multiples,
and the zero-section base lift to $a_0^N$.
Thus it has exactly the displayed subgroup.
The selected connected collar cover is therefore holomorphically
isomorphic to the punctured part of $\kappa_N$.

\item We construct the extension, including its normality and flatness.
At any point of the omitted divisor choose a polydisc with divisor
$z_1\cdots z_r=0$, $1\leq r\leq3$. Each connected component of
the restricted finite cover is classified by a finite-index subgroup
$K\subset\Z^r$. Write $M'=\operatorname{Hom}(K,\Z)$, containing
$\Z^r$ by restriction. Its completion is the normal toric analytic
space with ring generated by convergent monomials with exponents
\[
 M'\cap\R_{\geq0}^r
\]
over the remaining polydisc coordinates. This semigroup is saturated,
so the ring is normal. For each coset in $M'/\Z^r$ choose its unique
representative whose coordinates belong to $[0,1)$. The corresponding
monomials form a free basis over the original polydisc ring: subtracting
the representative leaves nonnegative integer coordinates, and distinct
cosets cannot cancel. Thus the extension is finite flat. It is also
the integral closure in the finite covering's local sheaf of meromorphic
functions. On each overlap the isomorphism of the unramified covers
and its inverse extend by
\cite[Expos\'e~XII, Proposition~5.3, pp.~254--255]{SGA1Recomposed2004}.
Uniqueness makes their composites the identity and gives the cocycle
equation. The local extensions therefore glue on all overlaps,
including the self-identifications of $W$.

\item The resulting space is normal and connected; finiteness over compact
$X$ makes it compact. The same uniqueness identifies the selected
cusp completion with the already normal smooth source in Theorem~\ref{maps:src:C.1.1}.
Above $S_1$ or $S_2$ the local construction has $r=1$, and its
invariant ring is $\C\{z^h,x_2,x_3\}$ for the order $h$ of the
corresponding diagonal subgroup. The map sending a new coordinate $t$
to $z^h$ and fixing $x_2,x_3$ identifies it with a convergent
three-variable power-series ring; the inverse sends the unique series
$\sum_{k\ge0}a_k(x_2,x_3)z^{hk}$ to
$\sum_{k\ge0}a_k(x_2,x_3)t^k$. Thus the source is smooth there,
as it is over $X^*$. The selected cusp component is smooth by
Theorem~\ref{maps:src:C.1.1}. For $N>1$, Theorem \ref{maps:src:C.1.3}
forces a singular point on the constructed compact source. That point
must therefore lie on another component over the cusp. This also
establishes the asserted location of the singularities.
\end{enumerate}
\end{proof}

\begin{remark}[Small congruence orbits and groups]
\hypertarget{tag-Q018}{}\label{mapsIII:rem:kummer-finite-values}
{\normalfont\scriptsize tag Q018}\par
For the marked matrices and $L_N$ above, the full monodromy orbit
has the following cardinalities:
\[
 \begin{array}{c|rrrrrrr}
 N&2&3&4&5&6&7&8\\\hline
 |\langle A_1,A_2\rangle L_N|&12&36&96&150&432&392&768.
 \end{array}
\]
For the exact matrix--phase groups in \eqref{maps:eq:C.1.4},
\[
 |G_2|=48,\qquad |G_3|=216,\qquad |G_4|=3072.
\]
These are finite congruence computations. The all-$N$ orbit bound
is the integral lattice argument of Proposition~\ref{maps:src:C.1.4},
and the all-$N$ finite cover is constructed in
Theorem~\ref{maps:src:C.1.5}.
\end{remark}


\subsection{Cohomology products and characteristic classes of cyclic covers}
\begin{definition}[Marked divisors and ruling classes]\hypertarget{tag-Q142}{}\label{mapsIII:def:cyclic-cohomology-data}
{\normalfont\scriptsize tag Q142}\par
Let $\pi:M=\widetilde X_p\to X_p$ be the cyclic cover of
Definition~\ref{mapsIII:def:cyclic-cover}, with $d=|p|$ and its
complex orientation. For $d>1$, index the cusp divisor cycles
$D_j$, $j\in\mathbb Z/d$, by positive $e_2$ translation, and
write $\nu_j:\overline D_j\to M$ for their normalization maps.
All component and curve classes have their complex orientations.
On $\overline D_j\cong dP_6$, let $H_j$ be the pullback of a line
and let $E_{A,j},E_{B,j},E_{C,j}$ be the exceptional classes of
rays $e_2,-e_1,e_1-e_2$, respectively. In the normalization basis
$(H,E_1,E_2,E_3)$ of Theorem~\ref{maps:src:D.1.2}, this is the
coordinate map
\[
 (a,b_1,b_2,b_3)\longmapsto(a,b_3,b_1,b_2)
 \quad\text{in the basis }(H,E_A,E_B,E_C).
\]
The inverse sends $(a,c_A,c_B,c_C)$ to $(a,c_B,c_C,c_A)$:
the ray identifications are $(E_1,E_2,E_3)=(E_B,E_C,E_A)$. Put
\[
 x_j=\operatorname{PD}\bigl((\nu_j)_*[\overline D_j]\bigr),
 \qquad r_j=(\nu_j)_*(H_j-E_{B,j}),
 \qquad \kappa=\operatorname{PD}(r_0).
\]
\end{definition}

\begin{proposition}[Characteristic classes and rational products]\hypertarget{tag-02KE}{}\label{maps:src:D.1.5}
{\normalfont\scriptsize tag 02KE}\par


For the complex orientation, $w_2(\widetilde X_p)=0$,
$c_1(\widetilde X_p)=0$, and $p_1(\widetilde X_p)$ is torsion of
order dividing $d$. The rational cup product
$H^2(\widetilde X_p;\mathbb Q)\otimes H^2(\widetilde X_p;\mathbb Q)
\to H^4(\widetilde X_p;\mathbb Q)$ is zero.
\end{proposition}


\begin{proof}
\emph{Characteristic classes.}
Use the cover of Definition~\ref{mapsIII:def:cyclic-cohomology-data}
and write $G=\mathbb Z/d$ for its deck group. The covering differential is the complex bundle
isomorphism $d\pi:TM\to\pi^*TX_p$. For every deck transformation $g$,
$dg:TM\to g^*TM$ likewise identifies the bundles. Naturality therefore
makes $c_1(TM)$ and $p_1(TM_{\mathbb R})$ invariant under $G$; see
\cite[Lemma~14.2, pp.~158--159, and Lemma~15.2, p.~174]{MilnorStasheff1974}.
By Theorem~\ref{maps:src:D.1.3}, their images in complex cohomology are zero.
The group $H^2(M;\mathbb Z)$ is free: the integral universal-coefficient
sequence identifies it with $\operatorname{Hom}(H_2(M;\mathbb Z),\mathbb Z)$,
since $M$ is simply connected. It follows that $c_1(TM)=0$ integrally.

For a complex vector bundle, $w_2$ of its underlying real bundle is
the mod-two reduction of $c_1$. One proof pulls back to its complete
flag bundle, where the bundle splits into line bundles and the
pullback in mod-two cohomology is injective by the iterated projective
bundle formula. For a complex line, its complex orientation identifies
its real Euler class with $c_1$; reducing that class modulo two gives
$w_2$. The Whitney formulas on the flag bundle give the asserted
identity for the original bundle, by injectivity. This proves the
identity recorded in \cite[Problem~14-B, p.~171]{MilnorStasheff1974}
and yields $w_2(M)=0$.

The zero complex image of $p_1(M)$ makes it an integral torsion class.
By Theorem~\ref{maps:src:D.1.4}, the integral universal-coefficient sequence
in degree four is
\[
0\longrightarrow\operatorname{Ext}^1_{\mathbb Z}(\mathbb Z/d,\mathbb Z)
\longrightarrow H^4(M;\mathbb Z)
\longrightarrow\operatorname{Hom}(\mathbb Z^{d-1},\mathbb Z)
\longrightarrow0.
\]
Applying $\operatorname{Hom}(-,\mathbb Z)$ to
$0\to\mathbb Z\xrightarrow{d}\mathbb Z\to\mathbb Z/d\to0$
identifies its left term with $\mathbb Z/d$. The right term is free,
so the left term is exactly the torsion subgroup of $H^4(M;\mathbb Z)$.
Consequently $d\,p_1(M)=0$. In particular $p_1(M)=0$ for $d=1$.

\emph{The rational product.}
Use the cusp tube $N_0$, its complement with finite fillings $U$,
their overlap $B_0$, and the central fibre from the proof of
Theorem~\ref{maps:src:D.1.3}. Set
$N=\pi^{-1}(N_0)$, $Z=\pi^{-1}(U)$, $E=\pi^{-1}(B_0)$,
and let $W$ denote the covered central fibre. These are the same
subspaces and covering restrictions as in that proof.
For a deck generator $g$ and $\lambda^d=1$, its cochain map
sends a local-system cochain to the lifted cochain
$\varphi$ with $\varphi(g\sigma)=\lambda\varphi(\sigma)$;
the inverse assigns its value to the coefficient fibre of the
downstairs simplex. Both maps commute with restriction.
Thus its contraction on the regular fibre and its finite-cover
transfer prove that the nontrivial-character cohomology of $Z$
and $E$ vanishes. Its lifted cusp homotopy gives a
$G$-equivariant retraction $r:N\to W$: restriction along
$W\hookrightarrow N$ and $r^*$ are inverse on cohomology.

For a nontrivial character $\chi$, applying the exact projector
$d^{-1}\sum_{g\in G}\chi(g)^{-1}g^*$ to Mayer--Vietoris therefore
makes restriction an isomorphism
\[
H^4(M;\mathbb C)_\chi\xrightarrow{\ \cong\ }
H^4(N;\mathbb C)_\chi\xrightarrow{\ \cong\ }
H^4(W;\mathbb C)_\chi.
\]
The normalization resolution in Theorem~\ref{maps:src:D.1.2} gives
$H^4(W;\mathbb C)\cong\bigoplus_jH^4(\overline D_j;\mathbb C)$ by
pullback along the normalization maps $\nu_j$ of
Definition~\ref{mapsIII:def:cyclic-cohomology-data}:
the double curves and triple points contribute no other term in
total degree four. Since Theorem~\ref{maps:src:D.1.3} gives no invariant
degree-four class on $M$, the resulting total restriction
\[
H^4(M;\mathbb C)\longrightarrow\bigoplus_jH^4(\overline D_j;\mathbb C)
\]
is injective.

Every degree-two class is a sum of nontrivial-character classes.
For a character with cyclic shift eigenvalue $t=\lambda\ne1$,
the restrictions to the two branches of every double curve agree.
Thus its component vector lies in the kernel of the
restriction matrix $B(\lambda)$ of Theorem~\ref{maps:src:D.1.2}.
Adding its first two row equations gives
$(\lambda-1)x_3=0$; subtracting its third row equation from
$\lambda$ times the first gives $(1-\lambda)x_2=0$.
The first row then gives $x_H+x_1=0$. Hence
\[
\ker B(\lambda)=\mathbb C(-H+E_1).
\]
Here $E_1$ is the exceptional divisor of ray $-e_1$ in the
normalization marking, and its intersection form gives
$(-H+E_1)^2=1-1=0$. On each component, the restrictions of any two
character classes are scalar multiples of this same vector,
including when their characters differ. Every product therefore
restricts to zero. Injectivity of the degree-four restriction proves
the complex cup product is zero. Extension of scalars
$H^4(M;\mathbb Q)\to H^4(M;\mathbb C)$ is injective, so the stated
rational cup product is zero as well. For $d=1$, the same conclusion
also follows directly from $H^2(M;\mathbb Q)=0$.
\end{proof}

\Needspace{6\baselineskip}
\begin{theorem}[Integral cohomology ring]\hypertarget{tag-02KF}{}\label{maps:src:D.1.6}
{\normalfont\scriptsize tag 02KF}\par


For $d>1$, the classes $x_j$ of
Definition~\ref{mapsIII:def:cyclic-cohomology-data} generate
$H^2(M;\mathbb Z)$ with the sole relation $\sum_jx_j=0$.
Here $d\ge5$, since admissibility gives $\gcd(d,6)=1$.
All ruling classes $r_j$ have the same image $r$, and
$\kappa=\operatorname{PD}(r)$ generates the torsion subgroup of
$H^4(M;\mathbb Z)$. The products are
\[
 x_i x_j=\begin{cases}
 -2\kappa,&i=j,\\
 \kappa,&i-j=1\text{ or }-1\pmod d,\\
 0,&\text{otherwise}.
 \end{cases}
\]
Together with the orientation pairing between the free parts of $H^2$
and $H^4$, these are every product in positive degree. For $d=1$ the
ring is the ring of $S^6$.
\end{theorem}


\begin{proof}
Write $M=\widetilde X_p$, let $N$ be its cusp tube, and let
$U$ be the exterior with its two finite fillings and a cusp collar.
Use the positive normal discs and the integral classes $h,z,z_0$
constructed in \cref{maps:src:D.1.4}. Their maps are
\[
 H_2(M,U)=\mathbb Z\langle n_0,\ldots,n_{d-1}\rangle
 \longrightarrow H_1(U)=\mathbb Zz_0,\qquad n_j\longmapsto z_0,
\]
and
\[
 H_3(M,U)\longrightarrow H_2(U)=\mathbb Zh\oplus\mathbb Zz,
 \qquad(c_w,c_\delta)\longmapsto(z,-ph).
\]
In particular $[h]$ generates the torsion subgroup of $H_2(M)$
and has exact order $d$.

\emph{Divisor generators.}
The relative Thom cocycle dual to $n_j$ maps to
$x_j=c_1(\mathcal O(D_j))$. At a self-intersection the Cartier
equation of $D_j$ is the product of its local branch coordinates;
the normal-disc evaluations add with their complex orientations.
This identifies the cocycle with the dual of the full divisor cycle.
The image of $H_2(U)$ in $H_2(M)$ is torsion. Every integral
degree-two cohomology class of $M$ therefore evaluates to zero on
that image. Since $H_1(U)$ is free, the universal coefficient map
$H^2(U)\to\operatorname{Hom}(H_2(U),\mathbb Z)$ is an
isomorphism. Thus $H^2(M)\to H^2(U)$ is zero.
The connecting map
$H^1(U)\to H^2(M,U)=\operatorname{Hom}(\mathbb Z^d,\mathbb Z)$
is the transpose of the displayed augmentation: its generator
evaluates to one on every $n_j$. Exactness gives
\[
 H^2(M)=\mathbb Z\langle x_0,\ldots,x_{d-1}\rangle
          /\mathbb Z(x_0+\cdots+x_{d-1}).
\]

\emph{The intersection curves.}
Use the marked exceptional classes of
Definition~\ref{mapsIII:def:cyclic-cohomology-data}.
The six boundary curves in cyclic ray order have classes
\[
 H-E_A-E_C,\quad E_A,\quad H-E_A-E_B,\quad
 E_B,\quad H-E_B-E_C,\quad E_C.
\]
Their branch identifications give
\[
 H_j-E_{A,j}-E_{C,j}=E_{B,j},\qquad
 E_{A,j}=E_{A,j+1},\qquad E_{C,j}=E_{C,j-1}
 \quad\text{in }H_2(N).
\]
Consequently the rulings $r_j=H_j-E_{B,j}$ all have the same
image $r=E_A+E_C$.
For $d>1$, admissibility implies $d\ge5$.
The intersection of $D_j$ with $D_{j+1}$ has the two curve
classes $E_A$ and $H-E_A-E_B$ on its normalization, whose
sum is $r_j$. For $D_{j-1}$ the corresponding sum is
$E_C+(H-E_B-E_C)=r_j$. All other distinct component pairs
have no common curve, by the edge directions of the fan.
The two Cartier equations have disjoint coordinate factors,
so each curve has intersection multiplicity one. Their complex
Thom classes therefore give
\[
 x_jx_{j+1}=x_jx_{j-1}=\operatorname{PD}(r),
 \qquad x_ix_j=0\quad(i\ne j,\ i-j\not\equiv\pm1\bmod d).
\]
Multiplication of $\sum_i x_i=0$ by $x_j$ gives
$x_j^2=-2\operatorname{PD}(r)$.

\emph{The integral ruling class.}
We identify its order by a chain calculation in the specified toric
cover. Put $Y=Y_\epsilon$, $T=Y\setminus\{q=0\}$, and
$\Gamma'=\mathbb Zu\oplus d\mathbb Z\gamma$, so that
$N=Y/\Gamma'$. The normal-disc calculation of
\cref{maps:src:D.1.4} gives
\[
 H_2(Y,T)=\bigoplus_{v\in\mathbb Z^2}\mathbb Zn_v,
 \qquad\partial n_v=q_{\rm circ}+v_1w+v_2\delta.
\]
Here $q_{\rm circ},w,\delta$ denote the positive coordinate circles.
The dense torus $T$ lies in a smooth triangle chart; radial
contraction of its three affine coordinates remains in
$|q|<\epsilon$. Hence $H_2(T)\to H_2(Y)$ is zero, and
$H_2(Y)\to H_2(Y,T)$ is injective.

Fix $0<\rho<\epsilon$ and $b=(1,1,\rho)\in T$.
Represent $n_0,n_{e_1}$ by the toric extensions of
$(1,1,\rho z)$ and $(z,1,\rho z)$ over $|z|\le1$.
Their boundary loops are $q_{\rm circ}$ and $q_{\rm circ}+w$ at $b$.
The linear triangle in the universal angle cover gives a chain
$L\subset T$ with
$\partial L=\partial n_{e_1}-\partial n_0-w$.
Thus
\[
 D=n_{e_1}-n_0-L,\qquad\partial D=w.
\]
The period action is
\[
 \Psi_u(x,q)=
 \bigl(\exp(2\pi iC(q)u)q^{e_1}x,q\bigr).
\]
At $q=\rho$, the path obtained by multiplying the two toric
coordinates by
\[
 \exp\bigl(t(\log\rho+2\pi i(C(\rho)u)_1)\bigr),
 \qquad
 \exp\bigl(2\pi it(C(\rho)u)_2\bigr)
 \quad(0\le t\le1)
\]
joins $b$ to $\Psi_u b$.
Rotate this path by the positive $w$-circle, orienting the
resulting cylinder $A$ by $t$ first. Then
$\partial A=\Psi_uw-w$, and its quotient is the marked
fibre torus $T(u,w)$. This follows directly from the
$u$-period column $Z(s)u$ with $\rho=e^{2\pi is}$.
The cycle $Z=\Psi_uD-D-A$ has relative class
\[
 n_{2e_1}-2n_{e_1}+n_0.
\]

A generic ruling $H-E_B$ in the component over $e_1$ is the
strict transform of a projective line through the point of $E_B$
and avoiding the other two blown-up points. Its intersections
with the six boundary curves above are $(1,0,0,1,0,0)$.
Thus it meets the components over $0$ and $2e_1$ once each and
no other distinct component. The total fibre divisor of $q$
is principal; its remaining intersection coefficient with its
own component is therefore $-2$. Its image in $H_2(Y,T)$ is
the same displayed vector. Injectivity of $H_2(Y)\to H_2(Y,T)$
identifies $[Z]$ with that complex ruling class.
Under the quotient $\pi_Y:Y\to N$ one has
$(\pi_Y)_*\Psi_uD=(\pi_Y)_*D$ as chains. Therefore
\[
 r=-[T(u,w)]\quad\text{in }H_2(N;\mathbb Z).
\]

In the finite-fibre bases of \cref{maps:src:D.1.4}, the same
torus maps to $B$, and the quotient functional has value
$\lambda(B)=2$. The overlap translations preserve its homology
class by the straight translation homotopy in the marked torus.
The gluing hence gives $r=-2[h]$ in $H_2(M)$.
Since $d$ is odd, the identity
\[
 -2\frac{d-1}{2}=1-d
\]
gives the inverse $[h]=\frac{d-1}{2}r$ in the cyclic subgroup
of order $d$. Thus $r$ generates this subgroup and has exact
order $d$. Poincar\'e duality with the complex orientation gives
the claimed generator $\kappa=\operatorname{PD}(r)$ of
$\operatorname{Tor}H^4(M)$.

\emph{The remaining products.}
Universal coefficients and \cref{maps:src:D.1.4} give
$H^3(M)=\mathbb Z/d$ and $H^1(M)=H^5(M)=0$.
Products of degree-two and degree-three classes land in $H^5=0$.
Products of two degree-three classes, and products involving
torsion and landing in $H^6=\mathbb Z$, are zero.
Poincar\'e duality gives the perfect orientation pairing between
the free quotients of $H^2$ and $H^4$. Together with the
degree-two products already computed, these exhaust the positive
degrees. For $d=1$, all intermediate cohomology groups vanish.
Sending the unit and complex orientation class to the corresponding
generators of $H^*(S^6;\mathbb Z)$ gives the asserted ring
isomorphism at that endpoint.
\end{proof}

\Needspace{6\baselineskip}
\begin{theorem}[Vanishing of the integral characteristic classes]\hypertarget{tag-02KH}{}\label{maps:src:D.1.8}
{\normalfont\scriptsize tag 02KH}\par


The pullback
\[
 \pi^*:H^4(X_p;\mathbb Z)\longrightarrow
 H^4(\widetilde X_p;\mathbb Z)
\]
is zero. In particular
$c_2(\widetilde X_p)=p_1(\widetilde X_p)=0$ integrally.
\end{theorem}


\begin{proof}
Use the complex orientations. In the marked regular fibre
$F=\Lambda_{\mathbb R}/\Lambda$, put
\[
 T:\mathbb R^2/\mathbb Z^2\longrightarrow F\longrightarrow X_p,
 \qquad [s,t]\longmapsto[s\widehat u+t\widehat w].
\]
This is an embedded torus, since $\widehat u,\widehat w$ span a
primitive direct summand of $\Lambda$. We first identify its integral
class downstairs. In the finite-filling Mayer--Vietoris calculation,
$\operatorname{coker}\alpha_2\simeq\mathbb Z$ has primitive dual
functional $\Phi$, represented by a class whose restriction to $F$ is
$2q$, where
\[
 q=u\wedge w+6\gamma\wedge\delta.
\]
The unhatted letters are dual to the marked fibre lattice. If $h$ is
the generator with $\Phi(h)=1$, the cusp attachment gives
$H_2(X_p;\mathbb Z)=\langle h\mid ph=0\rangle$;
these are the maps of \cite[Lemma~7.19(c), pp.~58--59;
Theorem~7.22, p.~60]{AC26}. The fibre torus has zero cut-boundary
component, and evaluation on $2q$ gives
\[
 [T]=\langle 2q,\widehat u\wedge\widehat w\rangle h=2h.
\]
Because $d=|p|$ is odd, this class generates $H_2(X_p;\mathbb Z)$.

The marked quotient of the fibre fundamental group sends
$\widehat\gamma$ to $1\in\mathbb Z/d$ and sends
$\widehat u,\widehat w,\widehat\delta$ to zero
\cite[Theorem~7.17, p.~57]{AC26}. Consequently the restricted covering is
\[
 F'=\Lambda_{\mathbb R}/\Lambda'\longrightarrow
 F=\Lambda_{\mathbb R}/\Lambda,
 \qquad
 \Lambda'=\mathbb Z\langle d\widehat\gamma,
 \widehat u,\widehat w,\widehat\delta\rangle.
\]
The inverse image of $T$ consists of the $d$ tori
\[
 T_j([s,t])=[j\widehat\gamma+s\widehat u+t\widehat w],
 \qquad j\in\mathbb Z/d,
\]
each mapping to $T$ with degree one. The homotopy
$([s,t],r)\mapsto[rj\widehat\gamma+s\widehat u+t\widehat w]$
in $F'$ identifies all their integral classes. Write $y$ for their
common image in $H_2(\widetilde X_p;\mathbb Z)$. It is fixed by
the deck group. By \cref{maps:src:D.1.3}, its rational image is zero;
by \cref{maps:src:D.1.4}, its torsion order divides $d$.
The chain transfer, which sums all lifts of every simplex
\cite[\S3.G, p.~321]{Hatcher2002}, therefore satisfies
\[
 \pi^![T]=\sum_{j=0}^{d-1}[T_j]=dy=0.
\]
Since $[T]=2h$ generates the downstairs group, $\pi^!$ vanishes on
all of $H_2(X_p;\mathbb Z)$.

For $a\in H^4(X_p;\mathbb Z)$, the cap-product chain identity gives
\[
 \operatorname{PD}_{\widetilde X_p}(\pi^*a)
 =\pi^!\bigl(\operatorname{PD}_{X_p}(a)\bigr).
\]
Indeed, lift each oriented simplex in a fundamental cycle to every
sheet. The pulled-back cochain has the same value on each lifted front
face, and the complementary faces are precisely the lifts of the
downstairs cap chain. Integral Poincar\'e duality
\cite[Theorem~3.30, p.~241]{Hatcher2002} now proves
$\pi^*=0$ in degree four. The covering derivative gives a complex
bundle isomorphism
$T\widetilde X_p\simeq\pi^*TX_p$, so naturality gives
\[
 c_2(\widetilde X_p)=\pi^*c_2(X_p)=0,
 \qquad
 p_1(\widetilde X_p)=\pi^*p_1(X_p)=0.
\]
For $d=1$, the same conclusion follows directly from
$H_2(X_p;\mathbb Z)=0$ and Poincar\'e duality.
\end{proof}

\Needspace{9\baselineskip}
\subsection{Finite-coefficient cubics and the fundamental class}
\begin{definition}[Coefficient cubic and the fixed edge]\hypertarget{tag-Q143}{}\label{mapsIII:def:cyclic-coefficient-data}
{\normalfont\scriptsize tag Q143}\par
For the cyclic cover of Definition~\ref{mapsIII:def:cyclic-cohomology-data}
with $d>1$, put $R=\mathbb Z/d$ and use $x_j$ also for the
reduction in $H^2(M;R)$ of its integral divisor class.
Let $C_\delta$ be the image in $X_p$ of the complex toric curve
on the edge $\{0,e_2\}$ and its translates in the cusp cover.
For indices in $\mathbb Z/d$, set
\[
 C_{ii}=-2,\qquad C_{i,i+1}=C_{i,i-1}=1,
 \qquad C_{ij}=0\quad(j\notin\{i,i-1,i+1\}).
\]
The coefficient cubic is the symmetric trilinear pairing
$(a,b,c)\mapsto\langle abc,[M]\rangle$ on $H^2(M;R)$,
using the complex fundamental class.
\end{definition}

\begin{theorem}[Finite-coefficient cubic]\hypertarget{tag-02KI}{}\label{maps:src:D.1.9}
{\normalfont\scriptsize tag 02KI}\par


For $d>1$, the curve $C_\delta$ of
Definition~\ref{mapsIII:def:cyclic-coefficient-data} is the
complex-oriented fixed curve of the $\delta$ action on $X_p$.
There is a class $v\in H^2(X_p;R)$ with $v(C_\delta)=1$.
For $u=\pi^*v$, the full coefficient cubic is
\[
 \begin{aligned}
 \langle x_i x_j x_k,[\widetilde X_p]\rangle&=0,
 &\langle u x_i x_j,[\widetilde X_p]\rangle&=C_{ij},\\
 \langle u^2x_i,[\widetilde X_p]\rangle&=-12,
 &\langle u^3,[\widetilde X_p]\rangle&=0
 \quad\text{in }\mathbb Z/d.
 \end{aligned}
\]
The $d-1$ classes $x_0,\ldots,x_{d-2}$ together with $u$ form a
$(\mathbb Z/d)$-basis of $H^2(\widetilde X_p;\mathbb Z/d)$.
In particular $u$ is deck invariant. 
\end{theorem}


\begin{proof}
Use the marked normalization and coefficient data of
Definitions~\ref{mapsIII:def:cyclic-cohomology-data} and
\ref{mapsIII:def:cyclic-coefficient-data}. The six boundary classes
are those in Theorem~\ref{maps:src:D.1.6}, through the identity
on the same component and its complex-oriented rays. Its branch
pushforwards give $H=E_A+E_B+E_C$ in the cusp and
$r=H-E_B=E_A+E_C$.

\emph{The fixed curve and its integral class.}
Write $n_v=(v,1)$ for the ray at $v$. In the chart of the triangle
$\{0,e_1,e_2\}$, the cocharacter of the delta action is
$(e_2,0)=n_{e_2}-n_0$. Its dual-coordinate action is
\[
 \lambda\cdot(z_0,z_1,z_2)=(\lambda^{-1}z_0,z_1,\lambda z_2).
\]
Thus the curve on the edge $\{0,e_2\}$ is fixed pointwise.
On any edge orbit the height-zero part of the stabilizer is the
primitive subgroup parallel to that edge; hence these are exactly
the fixed curves, together with their endpoints. Every triangle
contains one such edge. Their closures are disjoint copies of
$\mathbb P^1$, and the full period lattice acts transitively on them.
No nonzero vertex translation preserves an edge, so each maps
isomorphically to the single downstairs curve $C_\delta$.
On the normalized component it is the positive complex curve $E_A$.
The free period quotient introduces no additional fixed points,
since the correcting period element for a small connected arc of the
delta action is locally constant and initially the identity.
The delta circle is free on the regular fibres. On a finite filling,
applying the invariant covector $\gamma$ to a putative stabilizer
gives $a\ell_j/m_j\in\mathbb Z$, so admissibility forces the
affine power to be the identity and then forces the delta phase to
be integral. This identifies the full fixed curve.

We need two more integral cusp maps. Use the normal-disc classes
$N_v\in H_2(Y,T;\mathbb Z)$ of the toric tube, with
\[
 \partial N_v=\mu_q+v_1\widehat w+v_2\widehat\delta,
 \qquad H_2(Y,T;\mathbb Z)=\bigoplus_v\mathbb ZN_v.
\]
Here $\mu_q$ is the positive base meridian. Under the identity
of the pair $(Y,T)$ with the pair in Theorem~\ref{maps:src:D.1.6},
the classes $N_v$ are its $n_v$ and $\mu_q$ is its $q$-circle.
Thus the relative homology map is that theorem's injective map
$H_2(Y;\mathbb Z)\to H_2(Y,T;\mathbb Z)$.

At $b=(1,1,\rho)$, the discs $(1,1,\rho z)$ and $(1,z,\rho z)$
represent $N_0,N_{e_2}$. A triangle in the universal angle cover
subtracts the meridian difference and gives a chain $D$ with
$\partial D=\widehat\delta$ and relative class $N_{e_2}-N_0$.
For a period $\lambda$, join $b$ to $\Psi_\lambda b$ by the path
whose $i$th toric coordinate is multiplied by
\[
 \exp\!\left(t\bigl((B_0\lambda)_i\log\rho+
             2\pi i(C(\rho)\lambda)_i\bigr)\right),
 \qquad 0\le t\le1.
\]
Its positive delta-circle sweep $A_\lambda$, oriented with $t$
first, has boundary $\Psi_\lambda(\widehat\delta)-\widehat\delta$.
Thus $Z_\lambda=\Psi_\lambda D-D-A_\lambda$ is closed.

In the downstairs quotient take $\lambda=-\widehat\gamma$,
so $B_0\lambda=e_2$. The relative vector of $Z_\lambda$ is
$N_{2e_2}-2N_{e_2}+N_0$. A generic complex ruling $H-E_A$ on
the component at $e_2$ has boundary intersection numbers
$(0,1,0,0,1,0)$. It meets the components at zero and $2e_2$
once each; the principal total fibre gives self-component coefficient
$-2$. Its relative vector is therefore the same. Injectivity
identifies the absolute classes. In the downstairs quotient the
first two chains in $Z_\lambda$ cancel, while the sweep is the
marked torus $T(-\widehat\gamma,\widehat\delta)$. Consequently
\[
 H-E_A=E_B+E_C=T(\widehat\gamma,\widehat\delta)
 \quad\text{downstairs}.
\]
For $\lambda=\widehat u$, the relative vector is
$N_{e_1+e_2}-N_{e_1}-N_{e_2}+N_0$. The diagonal toric curve
has these coefficients: its self-intersection is $-1$ on each
adjacent component, and its two endpoint intersections are $1$.
The same injectivity and chain cancellation give
\[
 E_C=-T(\widehat u,\widehat\delta).
\]
This period belongs to both covering lattices, so the last identity
holds in both cusp quotients. The first ruling calculation in
\cref{maps:src:D.1.6} gives $r=-T(\widehat u,\widehat w)$.

Let $h$ be the downstairs generator defined by the primitive
finite-filling functional $\Phi(h)=1$. Its restriction to a regular
fibre is $2q_F=2u_F\wedge w_F+12\gamma_F\wedge\delta_F$,
where the subscript denotes the one-classes dual to the marked lattice
\cite[Lemma~7.19(c), pp.~58--59; Theorem~7.22, p.~60]{AC26}.
Evaluation on the fibre tori gives, respectively, $2h,0,12h$
for $(\widehat u,\widehat w)$, $(\widehat u,\widehat\delta)$
and $(\widehat\gamma,\widehat\delta)$. The curve maps above
therefore give
\[
 (H,E_A,E_B,E_C)\longmapsto(10h,-2h,12h,0),
 \qquad [C_\delta]=-2h.
\]
These fibre-cycle maps also apply to every cusp twist, since its
gluing translation is homotopic to the identity on a fixed fibre.

Put $b=(d-1)/2$, so $-2b=1-d$. The homomorphism
$H_2(X_p;\mathbb Z)\to R$ sending $h$ to $b$ takes
$[C_\delta]$ to one. The universal coefficient epimorphism
\cite[Theorem~3.2, pp.~195--196]{Hatcher2002} lifts it to
$v\in H^2(X_p;R)$. Set $u=\pi^*v$. The equality $\pi g=\pi$
makes $u$ deck invariant. Use the normalization map
$\nu_j:\overline D_j\to M$ of
Definition~\ref{mapsIII:def:cyclic-cohomology-data}; it is the
same component map used here. The period translation from $je_2$
to zero preserves its labelled rays and commutes with projection
downstairs. Hence
\[
 \bigl(u(H),u(E_A),u(E_B),u(E_C)\bigr)=(-5,1,-6,0),
 \qquad \nu_j^*u=-5H-E_A+6E_B
 \quad\text{in }H^2(\overline D_j;R).
\]
The latter equality uses the intersection form
$\operatorname{diag}(1,-1,-1,-1)$.

\emph{The basis and cubic.}
By \cref{maps:src:D.1.4}, $H_1(M;\mathbb Z)=0$ and
$H_2(M;\mathbb Z)=\mathbb Z^{d-1}\oplus\mathbb Z/d$.
Universal coefficients identify $H^2(M;R)$ with
$\operatorname{Hom}(H_2(M;\mathbb Z),R)$. The reductions of
integral degree-two classes are exactly the homomorphisms that kill
the torsion subgroup: lift their values on a basis of the free
quotient to integers. Thus the reductions of
$x_0,\ldots,x_{d-2}$ are a basis of that annihilator.
The ruling $r$ generates the torsion subgroup, and its downstairs
image is $[C_\delta]$, so $u(r)=1$. Every class $a$ decomposes
uniquely as $(a-a(r)u)+a(r)u$, proving the asserted basis over $R$.
In homology, subtracting $u(e_i)r$ from any lift $e_i$ of a free
quotient basis gives free lifts annihilated by $u$; an integer lift
of $u(e_i)$ gives the same element because $dr=0$.

The integral product $x_ix_j=C_{ij}\kappa$ is torsion, so its
product with $x_k$ is zero in $H^6(M;\mathbb Z)=\mathbb Z$.
Reduction gives the first cubic formula. Coefficient reduction
commutes with the cap-product definition of Poincar\'e duality;
the reduction of $\kappa$ is dual to $r$. Therefore
\[
 \langle ux_ix_j,[M]\rangle=C_{ij}u(r)=C_{ij}.
\]
The normalization pushes its positive fundamental class to $[D_j]$
with degree one. The projection formula gives
\[
 \langle u^2x_j,[M]\rangle
 =\langle(\nu_j^*u)^2,[\overline D_j]\rangle
 =25-1-36=-12\quad\text{in }R.
\]
Finally $\pi_*[M]=d[X_p]$, with the complex orientations, so
\[
 \langle u^3,[M]\rangle
 =\langle v^3,\pi_*[M]\rangle
 =d\langle v^3,[X_p]\rangle=0\quad\text{in }R.
\]
These values cover every basis triple. The positive $e_2$ deck
generator pulls $x_j$ back to $x_{j-1}$ and fixes $u$; thus the
basis also gives the explicit deck-module isomorphism
\[
 \left(R^d/R(1,\ldots,1)\right)\oplus R
 \longrightarrow H^2(M;R),\qquad
 ([a_j],b)\longmapsto\sum_ja_jx_j+bu.
\]
\end{proof}

\begin{definition}[Cubic coordinates and the fundamental-class invariant]\hypertarget{tag-Q144}{}\label{mapsIII:def:fundamental-class-invariant}
{\normalfont\scriptsize tag Q144}\par
Let $r\ge0$, let $d\in\mathbb Z$ satisfy $\gcd(d,6)=1$, and put
$A=\mathbb Z^r\oplus\mathbb Z/d$. Write $e_1,\ldots,e_r$
for its free basis and $t=[1]$ for its cyclic generator.
Under the degree-two universal coefficient isomorphism, the
integral classes $x_i$ on $K(A,2)$ are the homomorphisms dual
to $e_i$. The class $u_0\in H^2(K(A,2);\mathbb Z/d)$ is the
cyclic projection. For arbitrary $a_i\in\mathbb Z/d$, put
$u=u_0+\sum_i a_ix_i$, using the same symbols for the reductions
of the integral $x_i$. These are exactly the classes with $u(t)=1$.

For a closed simply connected oriented six-manifold $N$, let
$f_N:N\to K(H_2(N),2)$ be the homotopy class corresponding to
$\operatorname{id}_{H_2(N)}$ under the degree-two universal
coefficient isomorphism. Its fundamental-class invariant is
\[
 \mu(N)=(f_N)_*[N]
 \in H_6(K(H_2(N),2);\mathbb Z),
\]
as defined in \cite[Sections~1.1--1.2, p.~794]{manZhubr75}.
\end{definition}

\Needspace{6\baselineskip}
\begin{lemma}[Cubic detection of the fundamental-class invariant]\hypertarget{tag-02K8}{}\label{maps:src:D.1.10}
{\normalfont\scriptsize tag 02K8}\par


For the group and coordinates of
Definition~\ref{mapsIII:def:fundamental-class-invariant},
the integral homology group of $K(A,2)$ in degree six is
\[
 H_6(K(A,2);\mathbb Z)\simeq
 \mathbb Z^{\binom{r+2}{3}}\oplus
 (\mathbb Z/d)^{\binom{r+1}{2}+r+1}.
\]
Its free coordinates are detected by integral degree-two cubics;
its torsion coordinates are detected by the finite-coefficient
monomials $u x_i x_j$ ($i\leq j$), $u^2x_i$ and $u^3$.
For cyclic covers with $d>1$, Theorem~\ref{maps:src:D.1.9}
and the integral cubic determine the complete class
$\mu(\widetilde X_p)$. For the cover with $d=1$, this class is zero.
\end{lemma}


\begin{proof}
Write $D=|d|$; the identity on integers induces
$\mathbb Z/d\simeq\mathbb Z/D$ because $d\mathbb Z=D\mathbb Z$.
Put $L=\mathbb Z[1/6]$, $R=\mathbb Z/d$, and
$B=K(\mathbb Z,3)$. Pull back the path fibration $PB\to B$
along multiplication by $D$ on $B$, and denote its total space by $E$.
The homotopy exact sequence gives
\[
 0\longrightarrow\pi_3(E)\longrightarrow\mathbb Z
 \xrightarrow{D}\mathbb Z\longrightarrow\pi_2(E)
 \longrightarrow0,
\]
and every other positive homotopy group of $E$ vanishes.
Thus $E\simeq K(\mathbb Z/d,2)$. The fibre inclusion
$i:\mathbb CP^\infty\to E$ induces reduction modulo $d$ on $H_2$.

Use the Serre spectral sequences in
\cite[Theorems~5.3 and~5.15, pp.~527 and~542]{HatcherSpectralSequences}.
For the path fibration of $B$, let
$c\in H^2(\mathbb CP^\infty;L)$ be the positive generator and
choose $a\in H^3(B;L)$ by $d_3(c)=a$. Since $B$ is
$2$-connected and $PB$ is contractible, this transgression is
an isomorphism. The filtered cochain identity
\[
 d(\alpha\smile\beta)
 =d\alpha\smile\beta+(-1)^{|\alpha|}\alpha\smile d\beta
\]
induces the Leibniz rule on the first nonzero filtration
component, so $d_3(c^k)=kac^{k-1}$. Also $a^2=0$, since
$a$ has odd degree and $2$ is invertible in $L$.

Successively examining the bottom row shows
\[
 H^n(B;L)=
 \begin{cases}L,&n=0,3,\\0,&1\leq n\leq8,\ n\ne3.
 \end{cases}
\]
Indeed the possible incoming sources for degrees $4,5,6,7,8$,
after the previously determined zero columns are removed, are
respectively none, $c^2$ at $d_5$, $ac$ at $d_3$, $c^3$ at
$d_7$, and $ac^2$ at $d_5$. The classes $c^2,c^3$ have injective
$d_3$ maps, while $ac,ac^2$ are killed by $2c^2,3c^3$.
The map $d_3(ac)=-a^2$ is zero. Hence each possible additional
bottom-row group would survive to the cohomology of $PB$ and
must vanish.

Naturality for the pullback along $[D]$ gives $d_3(c)=Da$
in the sequence for $E$. Its cohomology in the required range is
therefore
\[
 \begin{gathered}
 H^2(E;L)=H^4(E;L)=H^6(E;L)=0,\\
 H^3(E;L)=L/(D),\qquad H^5(E;L)=L/(2D),\qquad
 H^7(E;L)=L/(3D).
 \end{gathered}
\]
Each nonzero diagonal has one surviving entry. The possible
later map from $(3,4)$ is $d_5$ into $(8,0)$, which vanishes
by the base calculation; the other later maps have zero source
or target. The three quotient ideals agree because $2,3$ are
units in $L$, and the residue map identifies each quotient with $R$.

To pass to integral homology, first observe that every positive
$H_n(E;\mathbb Z)$ is finite with prime divisors among those of $D$.
For $C_D=\langle g\rangle$, the group-ring resolution with
alternating maps $g-1$ and $N=1+g+\cdots+g^{D-1}$ is exact:
$\ker(g-1)=\mathbb ZN$ and
$\ker N=(g-1)\mathbb Z[C_D]$. Tensoring with the trivial
module gives
\[
 H_q(K(\mathbb Z/d,1);\mathbb Z)=
 \begin{cases}
 \mathbb Z,&q=0,\\R,&q>0\text{ odd},\\0,&q>0\text{ even}.
 \end{cases}
\]
Apply the homology Serre sequence to
$K(\mathbb Z/d,1)\to PE\to E$. Its bottom-row term
$H_n(E;\mathbb Z)$ has no incoming differential. By induction
on $n$, every outgoing target is a subquotient of
$H_{n-s}(E;R)$ for $2\leq s\leq n$, and these groups are
finite with prime divisors in $D$ by the homology universal
coefficient sequence. The final kernel is zero because $PE$
is contractible. The resulting finite filtration proves the
assertion, starting with $H_1(E;\mathbb Z)=0$.

Localization at $6$ preserves these finite groups. For such a
group $T$, $\operatorname{Hom}(T,L)=0$, while
$\operatorname{Ext}^1(\mathbb Z/m,L)=L/mL$ preserves each cyclic
invariant factor. The cohomology universal coefficient sequence
applied to the displayed cohomology groups consequently gives
\[
 H_2(E;\mathbb Z)=H_4(E;\mathbb Z)=H_6(E;\mathbb Z)=R,
 \qquad H_1(E;\mathbb Z)=H_3(E;\mathbb Z)=H_5(E;\mathbb Z)=0.
\]

The coefficient sequence $L\xrightarrow{D}L\to R$ gives
$H^p(B;R)=R$ for $p=0,3$ and zero for the other $p\leq7$.
In the sequence for $E$ with these coefficients, $d_3(c)=0$,
so $c,c^2,c^3$ survive. Their later differential targets vanish,
and each even diagonal through degree six has only its
column-zero term. Thus
\[
 i^*:H^{2k}(E;R)\xrightarrow{\sim}R c^k,
 \qquad 1\leq k\leq3.
\]
Let $u_0\in H^2(E;R)$ evaluate to one on the residue class
$[1]\in H_2(E;\mathbb Z)$. Then $i^*u_0=c$. The zero odd
homology groups make evaluation an isomorphism
$H^{2k}(E;R)\simeq\operatorname{Hom}(H_{2k}(E;\mathbb Z),R)$.
Consequently the complex-oriented classes
\[
 z_k=i_*[\mathbb CP^k],\qquad
 \langle u_0^k,z_k\rangle=1,
\]
are the exact cyclic generators in degrees two, four and six.
This argument applies to composite $D$ as well.

The projections and inclusions of $A$ induce a homotopy equivalence
\[
 K(A,2)\simeq P\times E,
 \qquad P=(\mathbb CP^\infty)^r.
\]
Their composites induce the identity on $\pi_2$, and all other
homotopy groups vanish. Under this equivalence, the classes in
Definition~\ref{mapsIII:def:fundamental-class-invariant} are the
pullbacks of the free-factor $x_i$ and cyclic-factor $u_0$ by
the two product projections. Restriction to the corresponding
factor is inverse on each of these summands. We use the same
symbols for these pullbacks. The integral K\"unneth cross-product map is
an isomorphism here because $H_*(P;\mathbb Z)$ is free in even
degrees. In total degree six it gives
\[
 \bigl(H_6(P)\otimes H_0(E)\bigr)\oplus
 \bigl(H_4(P)\otimes H_2(E)\bigr)\oplus
 \bigl(H_2(P)\otimes H_4(E)\bigr)\oplus
 \bigl(H_0(P)\otimes H_6(E)\bigr).
\]
The four multiplicities are respectively
$\binom{r+2}{3}$, $\binom{r+1}{2}$, $r$ and $1$.
For the first factor, write $x_i$ for its integral degree-two
generators, $\beta_i$ for $[\mathbb CP^1]$ in factor $i$, and
$\gamma_{ij}$ for $[\mathbb CP^2]$ when $i=j$ or
$[\mathbb CP^1\times\mathbb CP^1]$ in factors $i,j$ when $i<j$.
The finite-coordinate dual pairs are
\[
 \begin{array}{c|c}
 \gamma_{ij}\times z_1&u_0x_ix_j\\
 \beta_i\times z_2&u_0^2x_i\\
 {}[\mathrm{point}]\times z_3&u_0^3.
 \end{array}
\]
Each displayed pairing is one, and the pairings with the other
coordinates are zero by their degrees in the individual factors.
The free degree-six coordinates are likewise dual to the integral
monomials $x^\alpha$, $|\alpha|=3$, on the products
$\prod_i\mathbb CP^{\alpha_i}$. These pairings prove the claimed
group and the full evaluation isomorphism.

For the class $u$ in the statement, let $t$ be the cyclic generator
and $e_i$ the free basis of $A$. Choose integer representatives
$a_i$ of $u(e_i)$ and define
\[
 \alpha(e_i)=e_i-a_it,\qquad\alpha(t)=t.
\]
This is independent of the representatives because $dt=0$ and
has inverse $e_i\mapsto e_i+a_it$, $t\mapsto t$. It fixes each
integral $x_i$ and satisfies $K(\alpha,2)^*u=u_0$.
Naturality transports the preceding evaluation isomorphism to
the monomials formed with $u$.

Apply the classifying map of
Definition~\ref{mapsIII:def:fundamental-class-invariant} to $M$.
For every integral or finite-coefficient monomial $z$, naturality
gives $\langle z,\mu(M)\rangle=\langle f_M^*z,[M]\rangle$.
The evaluation isomorphism therefore recovers this full class
from the cubics in the statement. When $D=1$, the fibre $E$
is contractible and all finite coordinates are zero. For the
cover with $d=1$, Theorem~\ref{maps:src:D.1.4} gives $H_2=0$, so its
class $\mu$ is zero directly. For $r=0$, the same product
calculation has only the cyclic coordinate $z_3$.
\end{proof}

\subsection{Smooth recognition and cyclic complex families}
\begin{definition}[Canonical cyclic models and complex parameters]\hypertarget{tag-Q145}{}\label{mapsIII:def:canonical-cyclic-model}
{\normalfont\scriptsize tag Q145}\par
For each $d\ge1$ with $\gcd(d,6)=1$, take the marked datum
\[
 \mathcal A_d=(\ell_0,\ell_1,\ell_2,k_1,k_2,\sigma)
             =(0,-d,d,0,0,0).
\]
Write $\mathcal P=\{\operatorname{Im}c<-M\}$ for the admissible
half-plane of Definition~\ref{mapsIII:def:period-parameters},
and let $X_{d,c}=X(c;\mathcal A_d)$. Put $X_d=X_{d,-i}$ and
let $Z_d$ be its universal cover with the pulled-back complex
structure. Let $M_d$ be the underlying smooth manifold of $Z_d$,
with its complex orientation.

For an arbitrary marked datum and the parameter $w=e^{-\pi ic}$
of Definition~\ref{mapsIII:def:period-parameters}, set
\[
 W=(-1)^{k_2}w=e^{-\pi i(c+k_2)},\qquad
 u=w^2=e^{-2\pi ic}=W^2.
\]
For fixed integral $k_2$, the inverse coordinate change is
$w=(-1)^{k_2}W$. In the canonical datum $k_2=0$, so $W=w$.
\end{definition}

\Needspace{6\baselineskip}
\begin{theorem}[Smooth type and complex parameter]\hypertarget{tag-02KG}{}\label{maps:src:D.1.7}
{\normalfont\scriptsize tag 02KG}\par


Within the multiplicity-$(3,4)$ construction, the oriented
diffeomorphism type of $X_p$ depends exactly on $d=|p|$.
For each $d$ coprime to six it carries a punctured-disc family of
pairwise non-biholomorphic complex structures, parametrized by
$W=e^{-\pi i(c+k_2)}$, $0<|W|<216$. At each fixed nonzero admissible
linearization scale $u=e^{-2\pi ic}$ it carries exactly two of these
biholomorphism classes. Its universal covers likewise have one
oriented diffeomorphism type for each $d$.
\end{theorem}


\begin{proof}
Use the canonical data and parameter half-plane of
Definition~\ref{mapsIII:def:canonical-cyclic-model}. They satisfy
$p=4d-3d=d$, and $3\nmid d$, $2\nmid d$ give the two finite
admissibility conditions.
For an arbitrary datum $\mathcal A$ with $|p|=d$, the criterion in
\SeriesRef{fam:src:B.2.1}, with
\[
 c'=c+k_2,\qquad n=k_2,\qquad a=d/p\in\{1,-1\},
\]
gives a biholomorphism
\[
 B_{\mathcal A}:X(c;\mathcal A)\longrightarrow X_{d,c+k_2}
\]
over the marked base, with fibre-linear part $aI_2$. Indeed,
$p'=ap$ and $n\equiv k_2-k_2'\pmod2$, with $k_2'=0$.
The integer parameter shift preserves $\mathcal P$.

\emph{The canonical family.}
We construct a proper holomorphic submersion over $\mathcal P$.
Fix a parameter disc $Q$ with compact closure in $\mathcal P$.
On a regular base chart the quotient is defined by
\[
 (z,\zeta,c)\sim(z,\zeta+\Pi_c(z)\lambda,c),
 \qquad\lambda\in\Lambda.
\]
The real period map is nonsingular by \SeriesRef{fam:src:C.2.3}.
On compact base and parameter subsets, its least singular value is
bounded away from zero. The marking $\zeta=\Pi_c(z)x$ therefore
identifies the exterior family with a fixed real torus bundle times
the parameter; its integral transition matrices are constant in $c$.

The finite root pieces have real presentations
\[
 (D_j\times\Lambda_{\mathbb R}/\Lambda\times Q)/\langle r_j\rangle,
 \quad
 r_j(s,x,c)=(e^{-2\pi i/m_j}s,A_jx+v_j/m_j,c),
\]
with $v_1=-d\varepsilon_1$ and $v_2=d\varepsilon_2$.
The first-coordinate covector $\gamma$ is $A_j$-invariant.
A fixed point of $r_j^a$ on the central fibre, for $0<a<m_j$,
would give $a\ell_j/m_j\in\mathbb Z$, contrary to
$\gcd(\ell_j,m_j)=1$. Off the central fibre the root rotation
has no such fixed point. Thus these are free holomorphic finite
quotients. Their projections to $Q$ are submersions, and their closed
pieces over a compact parameter set are compact.

At the cusp use the defining toric space $Y$ and action
\[
 \Psi_{\lambda,c}(x,q)=
 \bigl(e^{2\pi iC_c(q)\lambda}q^{B_0\lambda}x,q\bigr),
 \qquad
 B_0=\begin{pmatrix}0&1\\-1&0\end{pmatrix}.
\]
On $\overline Q$ and a small closed $q$-disc, bound
$\|\operatorname{Im}C_c(q)\|$ by $M_0$. For $q\ne0$ set
\[
 y=-\frac{\log|q|}{2\pi},\qquad
 \ell=\frac{\log|x|}{\log|q|},\qquad
 L=B_0+\frac{\operatorname{Im}C_c(q)}y.
\]
The action on these real coordinates is
$\ell\mapsto\ell+L\lambda$. Shrink the disc so that
$y\ge2M_0+1$. Since $B_0$ is orthogonal, every singular value of
$L$ lies between $1/2$ and $3/2$. A fixed point at nonzero $q$
therefore has $\lambda=0$. At $q=0$ the finite nonempty set of
components through a fixed point would be invariant under translation
by $B_0\lambda$; summing its vertex labels again gives $\lambda=0$.

The action is properly discontinuous uniformly over $Q$. Indeed,
compact subsets lie in finitely many bounded toric charts. In the
chart of a triangle with vertices $v_0,v_1,v_2$, write
$b_i=\log|z_i|/\log|q|$. Then
\[
 \sum_i b_i=1,\qquad \ell=\sum_i b_i v_i,
 \qquad b_i\ge-\frac{\log R}{2\pi y}
\]
when $|z_i|\le R$ and $R\ge1$. Thus $\ell$ is uniformly bounded
on each compact set at nonzero $q$. A translate meeting another
compact set satisfies
$|\lambda|\le2(|\ell|+|\ell'|)$, allowing only finitely many
integral $\lambda$. On the central fibre the same conclusion follows
from its finitely many component labels in those charts.
The free quotient is therefore a complex manifold, locally a toric
chart times $Q$, and its parameter projection is a submersion.

For compactness, choose $\lambda$ with
$L^{-1}\ell+\lambda\in[0,1]^2$. The translated $\ell$ has norm
less than three. Only finitely many fan triangles meet this ball.
In a containing triangle its barycentric coefficients satisfy
$b_i\ge0$, and the dual integral chart coordinates obey
\[
 \log|z_i|=b_i\log|q|,
 \qquad |z_i|\le1,
 \qquad q=z_0z_1z_2.
\]
Hence finitely many closed unit polydiscs give representatives for
all noncentral orbits in the closed tube. Every central orbit has a
representative on the compact component $D_0\simeq\mathrm dP_6$,
because $B_0\mathbb Z^2=\mathbb Z^2$. These compact representatives
times a compact parameter set map onto the closed cusp quotient.

Glue these cusp, finite and exterior pieces by their defining
holomorphic overlap maps, with $\sigma=0$. They preserve the parameter
and satisfy their defining cocycles. Over the regular overlaps they
identify the same marked torus, and each special base point has its
own local piece; thus the glued space is Hausdorff. The closed pieces
with collars cover every fibre, so their compactness proves that the
parameter projection is proper. It is a holomorphic submersion by
the local charts. The constructions on parameter discs agree on
intersections through the identity of the defining models; shrinking
a cusp tube leaves the same regular overlap. They give a proper
holomorphic family on all of $\mathcal P$.

Along the segment from $c_0$ to $c_1$, pull back this submersion.
On the resulting smooth family choose a metric and the horizontal
field $V=\operatorname{grad}t/|\operatorname{grad}t|^2$.
It satisfies $dt(V)=1$. Local flow existence and compact continuation
\cite[author's preliminary version, Theorem~2.2, p.~38; Lemma~2.14, pp.~52--53]{TeschlODE2012}
over the closed segment make its
flow exist for the full interval, with inverse the negative-time
flow. Its fibre derivative preserves the continuous complex
orientation. Write $T_{c_0,c_1}:X_{d,c_0}\to X_{d,c_1}$ for this
oriented diffeomorphism. The comparison of two arbitrary data of
the same $d$ is therefore
\[
 B_{\mathcal A'}^{-1}\circ T_{c+k_2,c'+k_2'}\circ B_{\mathcal A}.
\]
Their fundamental groups have order $d$ by
\cref{maps:src:D.1.1}, which distinguishes different $d$.
The based lift of this comparison to universal covers is a
diffeomorphism: the based lift of its inverse is its inverse by
uniqueness of lifts. Evenly covered charts prove smoothness and
preservation of orientation. Different $d$ are distinguished upstairs
by $b_2=d-1$ from \cref{maps:src:D.1.3}.

\emph{Coherent descent to the punctured disc.}
We give the even shifts for the canonical data. Put
$S_n=I+n\widehat\delta\gamma$ for $n\in\mathbb Z$. The
periods and marked matrices satisfy
\[
 \Pi_{c+n}=\Pi_cS_n,\qquad \det S_n=1,
 \qquad S_nA_j=A_jS_n.
\]
Thus the regular lattices and their complex-linear transitions agree.
For even $n$, define
\[
\begin{aligned}
 r_1&=\frac{nd}{3}(\widehat u+\widehat w),&L_1&=nd\widehat w,\\
 r_2&=-\frac{nd}{4}(\widehat u+\widehat w),&L_2&=-\frac{nd}{2}\widehat u,
\end{aligned}
\]
and $b_j=\Pi_c r_j$. With $(a_1,a_2)=(-d,d)$, direct multiplication
gives $(I-A_j)r_j=na_j\widehat\delta/m_j+L_j$.
If $w_j=\log s_j/(2\pi i)$ in the marked clockwise root coordinate,
the finite overlap translation changes by $na_jw_je$, where
$e=\Pi_c\widehat\delta$. Hence its middle correction is
$a_j^{\rm loc}=b_j+na_jw_je$ and has integral affine cocycle $L_j$.

The cusp cocycle and fan translation are
\[
 L_0=-A_0(L_1+A_1L_2)
 =-\frac{3nd}{2}\widehat w+\frac{nd}{2}\widehat\delta,
 \qquad q_n=(3nd/2,-nd/2)^t.
\]
Since $s(g_0z)=s(z)-1$, the local section $s q_n$ has difference
$-q_n=\Pi_cL_0$. The map $(x,q)\mapsto(q^{q_n}x,q)$ extends
by the integral fan automorphism $(v,k)\mapsto(v+kq_n,k)$.
Also $C_{c+n}\lambda-C_c\lambda=(0,n\lambda_1)^t$, so the
entire cusp actions agree.

In the additive construction in the proof of
\SeriesRef{fam:src:A.3.1}, choose the standard global correction
$a_*^{(2)}$ and its local differences $\eta_{*,j}^{(2)}$ once.
Its vector sheaf fits into
$0\to\mathcal O_Be\to\mathscr V\to\mathcal O_B(-p_0)\to0$
and has $H^1(B,\mathscr V)=0$. The canonical cocycles and all
their prescribed local sections above are exactly $-d$ times the
standard ones, with the same vector sheaf. Consequently
\[
 a_d^{(2k)}=-dk a_*^{(2)},
 \qquad \eta_{d,j}^{(2k)}=-dk\eta_{*,j}^{(2)}
\]
give global corrections with those exact local extension conditions.
Every cocycle column has $\gamma=0$; the matrices preserve $\gamma$.
Thus the affine terms $\Pi_cL_g$ and these chosen corrections are
independent of $c$.

For $n=2k$ the middle and finite maps add, respectively,
$a_d^{(n)}$ and $b_j+\eta_{d,j}^{(n)}$. The cusp map is
\[
 (x,q,c)\longmapsto
 \bigl(e^{2\pi i\eta_{d,0}^{(n)}(q)}q^{q_n}x,q,c+n\bigr).
\]
The affine cocycles give descent of all three maps. Their finite
overlap equality is
$\sigma_{j,c}+a_d^{(n)}=\sigma_{j,c+n}+b_j+\eta_{d,j}^{(n)}$;
their cusp equality is
$e^{2\pi ia_d^{(n)}}=q^{q_n}e^{2\pi i\eta_{d,0}^{(n)}}$.
The inverses are the negative translations and inverse multipliers.
They glue to biholomorphisms $F_{n,c}:X_{d,c}\to X_{d,c+n}$.
Addition of the corrections and multiplication of the cusp units give
\[
 F_{2h,c+2k}F_{2k,c}=F_{2(h+k),c},\qquad F_{0,c}=\mathrm{id}.
\]
Their joint parameter dependence is holomorphic.

This action is free and properly discontinuous because it translates
the parameter by even integers. Local inverse logarithms identify
its quotient family with the proper local families above. The
parameter quotient is
\[
 W=e^{-\pi ic},\qquad 0<|W|<e^{-\pi M}=216.
\]
For canonical data the criterion of \SeriesRef{fam:src:B.2.1}
is exactly $c'-c\in2\mathbb Z$, so the quotient fibres are
pairwise non-biholomorphic. For arbitrary data their complete
invariant is $e^{-\pi i(c+k_2)}$, by the map $B_{\mathcal A}$.
Since $k_2$ is integral, $W^2=u$. At fixed admissible $u\ne0$
its two roots are distinct, have modulus $\sqrt{|u|}<216$, and
both occur in the family. The same classification criterion makes
these precisely two biholomorphism classes. All constructions include
$d=1$; the stated parameter discs have their strict endpoints.
\end{proof}

\begin{definition}[Complete invariant data of the cyclic model]\hypertarget{tag-Q146}{}\label{mapsIII:def:cyclic-classifying-data}
{\normalfont\scriptsize tag Q146}\par
For $d\ge1$ prime to six, put $A_d=\mathbb Z^{d-1}\oplus\mathbb Z/d$.
For $d>1$, use the cross-product generators of
Lemma~\ref{maps:src:D.1.10}, indexing the free basis by
$0,\ldots,d-2$. Index $i$ here corresponds to index $i+1$ in
Definition~\ref{mapsIII:def:fundamental-class-invariant}; this
reindexing preserves every group element and product.
Set
\[
 \mu_d=
 \sum_{0\le i\le j\le d-2}C_{ij}(\gamma_{ij}\times z_1)
 -12\sum_{i=0}^{d-2}(\beta_i\times z_2)
 \in H_6(K(A_d,2);\mathbb Z).
\]
Here $C$ is the principal submatrix on $0,\ldots,d-2$ of
Definition~\ref{mapsIII:def:cyclic-coefficient-data}: its diagonal
is $-2$, its two neighbouring diagonals are $1$, and its other
entries are zero. The free coordinates and the $z_3$ coefficient
of $\mu_d$ are zero. For $d=1$, put $A_1=0$ and $\mu_1=0$.
The cyclic classifying datum is
\[
 \mathcal I_d=(A_d,0,0,\mu_d).
\]
Its entries give second homology, the third Betti number, the
refined Pontryagin class $p(N)$ and the fundamental-class invariant,
respectively. The group-valued invariants are expressed by a
marking $H_2(N)\to A_d$.
\end{definition}

\begin{theorem}[Smooth recognition of the cyclic covers]\hypertarget{tag-02K9}{}\label{maps:src:D.1.11}
{\normalfont\scriptsize tag 02K9}\par
For each $d\ge1$ prime to six, the canonical manifold $M_d$ of
Definition~\ref{mapsIII:def:canonical-cyclic-model} is simply
connected and spin, and realizes the complete marked data
$\mathcal I_d$ of Definition~\ref{mapsIII:def:cyclic-classifying-data}.
These data specify a unique oriented diffeomorphism type among
closed simply connected spin six-manifolds. Every admissible
cover satisfies
\[
 \widetilde X_p\cong M_{|p|},\qquad M_1=S^6.
\]
\end{theorem}


\begin{proof}
Give $M=\widetilde X_p$ its complex orientation and put $d=|p|$.
It is a compact smooth simply connected six-manifold. Downstairs,
\cite[Theorems~7.17 and~7.22, pp.~57 and~60]{AC26} gives
$H_1(X_p;\mathbb Z)=H_2(X_p;\mathbb Z)=\mathbb Z/d$.
The degree-two universal coefficient sequence with coefficients
$\mathbb Z/2$ has outer terms
$\operatorname{Ext}^1(\mathbb Z/d,\mathbb Z/2)$ and
$\operatorname{Hom}(\mathbb Z/d,\mathbb Z/2)$.
They are the cokernel and kernel of multiplication by $d$ on
$\mathbb Z/2$, respectively, and vanish because $d$ is odd.
Thus $H^2(X_p;\mathbb Z/2)=0$ and $w_2(TX_p)=0$.
The covering derivative
$v_m\mapsto(m,D\pi_m(v_m))$, with inverse given in a covering
chart by the derivative of its local inverse, is the same complex
bundle map used in Proposition~\ref{maps:src:D.1.5}.
That proposition gives $w_2(TM)=0$, so $M$ is spin.
Its spin structure is unique since $H^1(M;\mathbb Z/2)=0$.

By Theorem~\ref{maps:src:D.1.4},
\[
 H_2(M;\mathbb Z)\simeq A_d,
 \qquad H_3(M;\mathbb Z)\simeq\mathbb Z/d.
\]
In particular $b_3(M)=0$. We use the marked classification
established in \cite[Theorem~3, p.~795; Section~4.1, p.~809]{manZhubr75}
for simply connected spin six-manifolds with zero free third homology.
That proof constructs the representative by killing
$H_3/\operatorname{Tor}H_3$ while preserving the prescribed
isomorphism on $H_2$. Its uniqueness argument changes the marked
spin bordism to an $h$-cobordism.
Theorem~5 of that paper, on p.~796, gives the injective map
\[
 (P,h):\Omega_6^{\mathrm{spin}}(K(G,2))
 \longrightarrow G\oplus H_6(K(G,2);\mathbb Z),
 \qquad h([N,f])=f_*[N].
\]
For closed simply connected spin six-manifolds $N_1,N_2$
with $b_3(N_i)=0$ and isomorphisms
$\phi_i:H_2(N_i;\mathbb Z)\to G$, equality of their marked
bordism classes therefore yields a degree $+1$ diffeomorphism
$F:N_1\to N_2$ satisfying $\phi_2F_*=\phi_1$.
Its inverse induces $\phi_1^{-1}\phi_2$ on second homology.
\footnote{The English rank conditions in
\cite[Sections~1.1--1.2, pp.~793--794]{manZhubr75} use $H_2$.
Section~4.1, p.~809, constructs the representative by killing
the free part of $H_3$ while preserving $H_2\simeq G$;
that construction and its marked uniqueness proof supply the
$b_3=0$ scope used here.}

For a group $G$ without $2$-torsion, the map $P$ is defined by
\[
 P([N,f])=\tfrac14 f_*\operatorname{PD}_N p_1(TN)
\]
in \cite[pp.~795--796]{manZhubr75}. Section~3.13, p.~807,
identifies $p(N)$ with $P(\gamma(N,\operatorname{id}_{H_2(N)}))$,
where $\gamma(N,\phi)$ denotes the bordism class of the map
inducing $\phi$ on $H_2$. Hence
\[
 4p(M)=\operatorname{PD}_M p_1(TM).
\]
Multiplication by four is injective on $A_d$: it is injective
on the free summand, and integers $a,b$ with $4a+db=1$ give
an inverse on $\mathbb Z/d$. The integral vanishing of $p_1$
in Theorem~\ref{maps:src:D.1.8} therefore gives $p(M)=0$.

Suppose first that $d>1$. Let $r\in H_2(M;\mathbb Z)$ be the
complex-oriented ruling generator of the torsion subgroup in
Theorem~\ref{maps:src:D.1.6}. The class $u$ in Theorem~\ref{maps:src:D.1.9}
satisfies $u(r)=1$. Choose integral lifts $e_i$ of the free
homology basis dual to $x_0,\ldots,x_{d-2}$, and choose integers
$a_i$ representing $u(e_i)$. Set
\[
 e_i'=e_i-a_ir,
 \qquad
 \psi:A_d\longrightarrow H_2(M;\mathbb Z),
 \quad (n_i,[a])\longmapsto\sum_i n_ie_i'+ar.
\]
The definition is independent of the integer representatives,
and $x_j(e_i')=\delta_{ij}$, $u(e_i')=0$. The map $\psi$ is
an isomorphism: the $x_i$ recover the free coordinates, and
$u$ recovers the remaining cyclic coordinate.

Use $f_M$ and $\mu(M)$ from
Definition~\ref{mapsIII:def:fundamental-class-invariant}.
The marked class is $K(\psi^{-1},2)_*\mu(M)$.
The pullbacks of the marked coordinate classes are exactly
$x_0,\ldots,x_{d-2},u$. Lemma~\ref{maps:src:D.1.10} identifies
their full integral and finite-coefficient evaluation maps with
the cross-product coordinates of
Definition~\ref{mapsIII:def:cyclic-classifying-data}.
The integral cubic vanishes by Theorem~\ref{maps:src:D.1.6}.
Theorem~\ref{maps:src:D.1.9} gives mixed coefficients $C_{ij}$
and $-12$, and zero coefficient for $z_3$. Each cross-product
generator pairs to one with its own monomial and to zero with
the others. Thus the complete marked class is
$K(\psi^{-1},2)_*\mu(M)=\mu_d$ as defined there, including all
zero free coordinates. It depends only on $d$.
For $d=1$, Theorem~\ref{maps:src:D.1.4} gives $H_2=H_3=0$;
use the unique marking $0\to0$ and $\mu_1=0$, since
$K(0,2)$ is a point.

For two manifolds with the stated data, choose the corresponding
markings $\psi_i:A_d\to H_2(N_i;\mathbb Z)$ and put
$\phi_i=\psi_i^{-1}$. Naturality of $P$ and the defining
pushforward of $\mu$ give
\[
 (P,h)\bigl(\gamma(N_i,\phi_i)\bigr)=(0,\mu_d).
\]
Theorem~5 of \cite{manZhubr75} makes the two bordism classes
equal, and Theorem~3 gives the orientation-preserving
diffeomorphism with
$F_*=\psi_2\psi_1^{-1}$ on second homology. This proves
uniqueness for the stated complete invariant.

For every positive $d$ prime to six, the canonical datum of
Definition~\ref{mapsIII:def:canonical-cyclic-model} is admissible
and has $p=4d-3d=d$. Its member at $c=-i$ has
$W=e^{-\pi}\in(0,216)$, in the parameter domain of
Theorem~\ref{maps:src:D.1.7}. Its cover $Z_d$, whose underlying
oriented manifold is $M_d$, realizes the data just computed. Every
other admissible cover with $|p|=d$ has the same data, so the
marked classification gives $\widetilde X_p\cong M_d$.

For $d=1$, the standard oriented sphere $S^6$ has the same
homology and fundamental-class data. Its stable tangent
bundle is trivial through
\[
 TS^6\oplus\mathbb R\longrightarrow S^6\times\mathbb R^7,
 \quad (x;v,a)\longmapsto(x;v+ax),
\]
with inverse $(x;y)\mapsto(x;y-\langle y,x\rangle x,
\langle y,x\rangle)$. Thus it is spin with $p_1=0$.
Applying the same marked classification with $G=0$ yields
the oriented diffeomorphism $M_1\cong S^6$.
\end{proof}

\begin{theorem}[A canonical cyclic action and its quotient]\hypertarget{tag-02KA}{}\label{maps:src:D.1.12}
{\normalfont\scriptsize tag 02KA}\par
For every $d\geq1$ prime to six, the manifold $M_d$ carries a free
cyclic action $\sigma_d$ defined by the affine and toric construction
in the proof. Every admissible multiplicity-$(3,4)$ member with
$|p|=d$ is orientation-preservingly diffeomorphic to
$M_d/\langle\sigma_d\rangle$. Its universal-cover deck action is
equivariantly diffeomorphic to $\sigma_d$, allowing an automorphism
of the abstract group $\mathbb Z/d$.
The complex structures on this quotient arising from the
multiplicity-$(3,4)$ construction form the punctured-disc family
of Theorem~\ref{maps:src:D.1.7}, with exactly two of its biholomorphism
classes at each fixed nonzero admissible scale $u$.
For $d=1$, the group is trivial and the quotient is $M_1=S^6$.
\end{theorem}


\begin{proof}
Use the canonical member $X_d$ and its cover $Z_d$ from
Definition~\ref{mapsIII:def:canonical-cyclic-model}. The
admissibility and strict parameter conditions for this same
datum are verified in Theorems~\ref{maps:src:D.1.7} and
\ref{maps:src:D.1.11}.

Use the original homology basis $(\gamma,u,w,\delta)$ and the
covering basis $(g,u,w,\delta)$, where $g=d\gamma$. Their inclusion
is $S_d=\operatorname{diag}(d,1,1,1)$. The covered root matrices are
\[
 A_1^{(d)}=
 \begin{pmatrix}
 1&0&0&0\\6d&0&1&0\\-6d&-1&-1&0\\-2d&1&0&1
 \end{pmatrix},\qquad
 A_2^{(d)}=
 \begin{pmatrix}
 1&0&0&0\\0&0&-1&0\\-6d&1&0&0\\3d&0&1&1
 \end{pmatrix}.
\]
They satisfy $S_dA_j^{(d)}=A_jS_d$. The finite affine vectors are
\[
 v'_1=(-1,-2d,4d,0)^t,\qquad v'_2=(1,3d,-3d,0)^t.
\]
They are fixed by their respective matrices, and the finite pieces
are the quotients by
\[
 R_j(s,x)=
 \left(e^{-2\pi i/m_j}s,A_j^{(d)}x+\frac{v'_j}{m_j}\right),
 \qquad(m_1,m_2)=(3,4).
\]
Their $m_j$-th powers are the integral translations $v'_j$.
The first-coordinate shifts $-1/3$ and $1/4$ prove freeness on the
central fibres; the root rotation proves it elsewhere. These are
the covering pieces of \cref{maps:src:D.1.4}. In complex coordinates
their period matrix is $\Pi_{d,c}=\Pi_cS_d$, so the displayed
intertwining relation also identifies their holomorphic root maps.

On a regular covered fibre define $D_\gamma$ by translation through
\[
 a=(1/d,0,0,0)^t.
\]
It has order $d$, and
\[
 A_1^{(d)}a-a=(0,6,-6,-2)^t,\qquad
 A_2^{(d)}a-a=(0,0,-6,3)^t.
\]
These are covering-lattice vectors. Hence $D_\gamma$ commutes with
$R_j$ on the torus and descends to each finite piece. Its formulas
agree on the regular overlaps. If $0<k<d$ and $D_\gamma^k$ fixed
a point of a finite quotient, comparison of the first coordinates
would give
\[
 \frac{k}{d}\equiv-\frac{j}{3}\pmod{\mathbb Z}
 \quad\text{or}\quad
 \frac{k}{d}\equiv\frac{j}{4}\pmod{\mathbb Z}.
\]
Coprimality forces $d\mid k$, which excludes such a stabilizer.

At the cusp let
$\Gamma=\mathbb Z\gamma\oplus\mathbb Z u$ and
$\Gamma_d=\mathbb Z d\gamma\oplus\mathbb Z u$.
In the defining toric coordinates the period action is
\[
 \Psi_\lambda(x,q)=
 \bigl(e^{2\pi iC_{-i}(q)\lambda}q^{B_0\lambda}x,q\bigr),
 \qquad B_0=\begin{pmatrix}0&1\\-1&0\end{pmatrix}.
\]
Thus $B_0\gamma=-e_2$ and $B_0u=e_1$, and the covering cusp is
$Y/\Gamma_d$. Define $D_\gamma$ there by $\Psi_\gamma$. It is the
fan translation by $-e_2$ with the holomorphic unit displayed above.
The identity $\Psi_\lambda\Psi_\mu=\Psi_{\lambda+\mu}$ proves that
it descends and that its $d$-th power is trivial on this quotient.
On the regular overlap it is the same marked period translation.

The action of $\Gamma$ on $Y$ is free. On the central fibre, a fixed
point would have a finite nonempty set of incident component labels
invariant under translation by $B_0\lambda$; their sum forces
$\lambda=0$. Off that fibre, the logarithmic-coordinate matrix
$B_0+\operatorname{Im}C_{-i}(q)/y$, with
$y=-\log|q|/(2\pi)$, is invertible on the cusp tube by the estimates
in \cref{maps:src:D.1.7}. A fixed point again forces $\lambda=0$.
Consequently a fixed point of $D_\gamma^k$ on $Y/\Gamma_d$ would
imply $k\gamma\in\Gamma_d$, or $d\mid k$. This proves freeness on
all of $Z_d$.

The lattice projection identifies the quotient by this action with
$X_d$. On regular fibres it is induced by
$\mathbb R^4/\Lambda_d\to\mathbb R^4/\Lambda$; its inverse on orbit
spaces sends $[x]_\Lambda$ to the orbit of any covering lift.
The maps $S_dA_j^{(d)}=A_jS_d$ and $S_dv'_j=v_j$ extend this
identification over the finite pieces. At the cusp it is
$(Y/\Gamma_d)/(\Gamma/\Gamma_d)\to Y/\Gamma$, with the same
orbit-lift inverse. The maps commute with the defining overlaps,
so they give a biholomorphism
\[
 Z_d/\langle D_\gamma\rangle\cong X_d.
\]
Use the identity from $Z_d$ to its underlying oriented manifold
$M_d$ in Definition~\ref{mapsIII:def:canonical-cyclic-model} to
set $\sigma_d=D_\gamma$. For any other oriented representative
$N_d$ of this type, Theorem~\ref{maps:src:D.1.11} gives an oriented
diffeomorphism $\iota_d:M_d\to N_d$. The transported action is
$\iota_d\sigma_d\iota_d^{-1}$, with inverse conjugacy
$\iota_d^{-1}$.

For an arbitrary member $X(c;\mathcal A)$ with $|p|=d$,
\cref{maps:src:D.1.7} gives the oriented diffeomorphism
\[
 F_{\mathcal A}=T_{c+k_2,-i}\circ B_{\mathcal A}:
 X(c;\mathcal A)\longrightarrow X_d.
\]
Choose matching basepoints. Its based lift between universal covers
has inverse the based lift of $F_{\mathcal A}^{-1}$, by uniqueness
of lifts, and is a smooth oriented diffeomorphism in evenly covered
charts. If $g$ is a deck element, the induced automorphism
$\alpha_{\mathcal A}=(F_{\mathcal A})_*$ satisfies
\[
 \widetilde F_{\mathcal A}D_g\widetilde F_{\mathcal A}^{-1}
 =D_{\alpha_{\mathcal A}(g)}.
\]
Indeed both sides have the same effect on the endpoint of every
based lifted loop. The map $\alpha_{\mathcal A}$ is an automorphism
of the cyclic group, with inverse induced by
$F_{\mathcal A}^{-1}$. This proves the equivariant classification
and the quotient diffeomorphism.

Transporting a member's complex charts through $F_{\mathcal A}$
identifies its complex structure with one on this quotient. Their
classification in \cref{maps:src:D.1.7} is by
$W=e^{-\pi i(c+k_2)}$, $0<|W|<216$. At fixed nonzero admissible
$u$, the equation $W^2=u$ has the two distinct admissible roots
$W,-W$, giving exactly the stated two biholomorphism classes.
For $d=1$ the two lattice quotients are trivial, $D_\gamma$ is the
identity and the quotient is $M_1=S^6$.
\end{proof}

\section{Marked elliptic models and triangular caps}
\hypertarget{tag-02OJ}{}\label{mapsD:repair:modular-collisions}

\subsection{Marked rational elliptic surfaces}

\begin{definition}[Marked monodromy and section cocycles]
\hypertarget{tag-Q100}{}\label{mapsIII:def:marked-elliptic-cocycles}
{\normalfont\scriptsize tag Q100}\par
We use positive base meridians and column vectors for integral elliptic
homology.  A positive nodal monodromy is
\[
 T=\begin{pmatrix}1&1\\0&1\end{pmatrix},\qquad
 S=\begin{pmatrix}0&-1\\1&0\end{pmatrix},\qquad
 R=ST=\begin{pmatrix}0&-1\\1&1\end{pmatrix}.
\]
Thus $s=[S]$ and $r=[R]$ generate
$G=\operatorname{PSL}_2(\mathbb Z)=\langle s,r\mid s^2=r^3=1\rangle$,
and $[T]=sr$.  Order the positive based meridians so that their
monodromy product is $A_1\cdots A_n=I$.  The section cocycle
$(q_i)$ satisfies
\begin{equation}\hypertarget{tag-02OF}{}\label{mapsD:repair:cocycle-relation}
 q_i\in\operatorname{sat}\operatorname{im}(A_i-I),\qquad
 \sum_{i=1}^n P_iq_i=0,\qquad P_i=A_1\cdots A_{i-1}.
\end{equation}
Changing one global logarithm by $v\in\mathbb Z^2$ changes $q_i$ by
$(A_i-I)v$.  Saturation here means intersection of the rational image
with $\mathbb Z^2$.
\end{definition}

\begin{remark}[The two fibre configurations]
\hypertarget{tag-Q101}{}\label{mapsIII:rem:collision-configurations}
{\normalfont\scriptsize tag Q101}\par
The fibre configurations and rank-one Mordell--Weil lattices below
occur in Miranda's existence list
\cite[\S4, pp.~206--210]{manMiranda90} and in Fukae's
\cite[Table~3, rows~50 and~56, p.~35]{manFukae02}.
Theorem~\ref{mapsD:repair:collision-realization} gives their positive
based monodromy and integral section cocycles explicitly.
\end{remark}

\begin{remark}[Positive projective periods]
\hypertarget{tag-Q102}{}\label{mapsIII:rem:positive-projective-periods}
{\normalfont\scriptsize tag Q102}\par
On a simply connected regular base chart, choose periods
$\omega_1,\omega_2$ with $\tau=\omega_1/\omega_2\in\mathbb H$.
The map from the period quotient to its normalized form and its inverse are
\[
 [z]\longmapsto[z/\omega_2],\qquad
 [w]\longmapsto[\omega_2w],\qquad
 \mathbb C/(\mathbb Z\omega_1+\mathbb Z\omega_2)
       \longleftrightarrow\mathbb C/(\mathbb Z\tau+\mathbb Z).
\]
An oriented integral period change
$g=\left(\begin{smallmatrix}a&b\\c&d\end{smallmatrix}\right)\in SL_2(\mathbb Z)$
acts on the period column by $g$, with inverse $g^{-1}$, and sends
$\tau$ to $\tau'=(a\tau+b)/(c\tau+d)$. The normalized fibre map is
$w\mapsto w/(c\tau+d)$, with inverse multiplication by $c\tau+d$.
It maps the period lattices bijectively because
\[
 \frac1{c\tau+d}=a-c\tau',\qquad
 \frac{\tau}{c\tau+d}=-b+d\tau',\qquad ad-bc=1.
\]
The period map and functional invariant described in
\cite[p.~2]{manHerfurtner90} therefore identify the projective period
system with the pullback of the modular system by $J$.
For positive loops around $J=0,1$ we use
\[
 A=\begin{pmatrix}0&1\\-1&1\end{pmatrix},\qquad
 B=\begin{pmatrix}0&1\\-1&0\end{pmatrix}.
\]
Their actions are $A(\tau)=1/(1-\tau)$ and $B(\tau)=-1/\tau$.
Their derivatives at $e^{\pi i/3}$ and $i$ are
$e^{2\pi i/3}$ and $-1$, fixing the positive local orientations.
The cusp conventions are related by the literal identities
\[
 A^{-1}B=T,\qquad
 -(AB)^{-1}=STS^{-1}=\begin{pmatrix}1&0\\-1&1\end{pmatrix}.
\]
Conjugation by $S$ has inverse conjugation by $S^{-1}$, so both
displayed cusp matrices are positive primitive Picard--Lefschetz
matrices in their indicated integral period bases.
\end{remark}

\begin{theorem}[Two marked rational elliptic surfaces]
\hypertarget{tag-02OG}{}\label{mapsD:repair:collision-realization}
{\normalfont\scriptsize tag 02OG}\par
There are rational elliptic surfaces with zero section over
$\mathbb P^1$ having, in the indicated order, the following singular
fibres and positive based monodromy matrices:
\[
\begin{array}{c|c|c}
\text{profile}&\text{fibre types}&(A_1,\ldots,A_n)\\ \hline
50/\mathrm{II}&(\mathrm{II},I_3,I_1^*)&
\left(
\begin{pmatrix}3&7\\-1&-2\end{pmatrix},
\begin{pmatrix}7&12\\-3&-5\end{pmatrix},
\begin{pmatrix}-2&-1\\1&0\end{pmatrix}\right)\\[5pt]
56/\mathrm{II}&(I_2,\mathrm{II},I_5,I_3)&
\left(
\begin{pmatrix}5&8\\-2&-3\end{pmatrix},
\begin{pmatrix}27&37\\-19&-26\end{pmatrix},
\begin{pmatrix}61&80\\-45&-59\end{pmatrix},
\begin{pmatrix}4&3\\-3&-2\end{pmatrix}\right).
\end{array}
\]
Their Mordell--Weil groups are infinite cyclic.  A generator can be
chosen with the respective integral logarithm cocycle
\[
 q^{50}=\bigl((-5,2),(-8,4),(0,1)\bigr),\qquad
 q^{56}=\bigl((-2,1),(0,0),(4,-3),(1,-1)\bigr).
\]
Its height is $1/12$ in profile $50/\mathrm{II}$ and $1/30$ in
profile $56/\mathrm{II}$.  The matrices act in an oriented integral
homology basis, and the cocycle uses a compatible integral logarithm branch.
\end{theorem}

\begin{proof}
We construct the two surfaces and then identify their full period and
section markings.  Write
\[
 y^2=4x^3-F(t)x-H(t),\qquad \Delta=F^3-27H^2,\qquad J=F^3/\Delta.
\]
For $50/\mathrm{II}$ take
\[
 F=3t(t+8),\qquad H=t(t^2-20t-8).
\]
For $56/\mathrm{II}$ put
\[
 \begin{split}
 P(t)&=81t^3-9t^2-53t-27,\\
 Q(t)&=729t^5-1215t^4-1350t^3-350t^2-675t-243,
 \end{split}
 \qquad F=3(t-3)P(t),\quad H=(t-3)Q(t).
\]
The latter pair is the last row of
\cite[Table III, p.~23]{manHerfurtner90}.  Direct multiplication gives
\[
 \Delta_{50}=1728t^2(t-1)^3,
 \qquad \Delta_{56}=-2^{14}3^4t^3(t-3)^2(9t+5)^2.
\]
The second base chart and its inverse are
\[
 s=t^{-1},\quad X=s^2x,\quad Y=s^3y;
 \qquad t=s^{-1},\quad x=s^{-2}X,\quad y=s^{-3}Y.
\]
The ordered valuation triples of $F,H,\Delta$ are
\[
 \begin{array}{c|c|c}
 \text{profile}&\text{places}&\text{valuation triples}\\\hline
 50/\mathrm{II}&0,1,\infty&(1,1,2),(0,0,3),(2,3,7)\\
 56/\mathrm{II}&-5/9,3,\infty,0&(0,0,2),(1,1,2),(0,0,5),(0,0,3).
 \end{array}
\]
These are precisely the stated minimal fibre types by
\cite[Table I, p.~6]{manHerfurtner90}; the discriminants have no other
zeros.  The map $(x,y)\mapsto(x,y/2)$, with inverse
$(x,y')\mapsto(x,2y')$, gives the monic convention of
\cite[\S8.4, p.~39]{manSS10}.  Its minimal degree criterion gives
rational elliptic surfaces, with the zero section $[0:1:0]$.
The coefficients glue as sections of $\mathcal O(4)$ and
$\mathcal O(6)$ in the displayed chart change.  The resulting
Weierstrass hypersurface in
$\mathbb P(\mathcal O\oplus\mathcal O(2)\oplus\mathcal O(3))$
is projective; resolving its Weierstrass singularities gives the
minimal elliptic surface with fundamental line bundle $\mathcal O(1)$.

\emph{The exact covering and peripheral marking.}
Their rational invariant maps are
\[
 \begin{aligned}
 J_{50}&=\frac{t(t+8)^3}{64(t-1)^3},&
 J_{50}-1&=\frac{(t^2-20t-8)^2}{64(t-1)^3},\\
 J_{56}&=-\frac{(t-3)P(t)^3}{49152t^3(9t+5)^2},&
 J_{56}-1&=-\frac{Q(t)^2}{49152t^3(9t+5)^2}.
 \end{aligned}
\]
The normalized modular coordinate is $J$; the conventional invariant is
$j=1728J$, as in \cite[Proposition 4.3 and Remark 4.4, p.~49]{mapsMilne2017}.
All the displayed distinct factors are relatively prime and squarefree.
For these numerators and denominators, the derivative numerators are
\[
 64(t-1)^2(t+8)^2(t^2-20t-8),\qquad
 -245760t^2(9t+5)P(t)^2Q(t).
\]
Thus the degrees are four and ten and the branch partitions over
$0,1,\infty$ are
\[
 (3,1),(2,2),(3,1),\qquad
 (3,3,3,1),(2,2,2,2,2),(5,3,2).
\]
There are no other branch values; the displayed ramification totals are
$6=2\cdot4-2$ and $18=2\cdot10-2$.

\phantomsection\hypertarget{tag-033V}{}\label{qa:III:layout:4}
Here is a finite identification of these coverings.  Fix the
$0$-permutation to $(012)(3)$ or $(012)(345)(678)(9)$, and enumerate
the perfect matchings for the $1$-permutation by pairing the least
unused label with each remaining label in turn.  There are three and
$945$ matchings.  Keep the connected pairs whose product has cusp
widths $(3,1)$ or $(5,3,2)$.  To compare two pairs, start at each
possible root, number vertices by breadth-first search in the order
$a,b$, and take the least resulting transition table.  Equality of
these tables gives the simultaneous label bijection and its inverse;
inequality excludes such a bijection.  The retained counts are three
and $162$, respectively, and each list has one table class.  Useful
representatives, in image-list notation starting at zero, are
\[
 \begin{array}{c|l|l}
 &a&b\\\hline
 50&(0,2,3,1)&(1,0,3,2)\\
 56&(1,2,0,6,8,5,7,3,9,4)&(3,4,5,0,1,2,9,8,7,6).
 \end{array}
\]
The loops $a,b$ are the following positive rectangles based at $J=1/2$:
\[
 \begin{split}
 a:&\ (1/2,1/2+i/2,-1/2+i/2,-1/2-i/2,1/2-i/2,1/2),\\
 b:&\ (1/2,1/2-i/2,3/2-i/2,3/2+i/2,1/2+i/2,1/2).
 \end{split}
\]

We specify the covering isomorphism, including its choice of sheets.
Write $J=N/D$ using the coprime numerators and denominators displayed
above.  Glue cut $J$-spheres by the table and complete each end by
$z\mapsto z^e$.  Order the roots of $N-D/2$ by real part and then
imaginary part.  For a root $r$, continue it along the tree words
$w_i$ below, obtaining $r_i$.  Take the first $r$ for which these
roots are distinct and continuation along $a,b$ sends $r_i$ to
$r_{a(i)},r_{b(i)}$.  The one-class enumeration proves existence of
this finite choice.  The map sends sheet $i$ to $r_i$ and a path
from $1/2$ to the endpoint of its lift under the rational map
$J$.  Equivariance proves independence of path representatives.
Its inverse lifts a chosen path backwards to the unique $r_i$ and
then forwards from sheet $i$ in the glued cover.  The two composites
are identity by path-lifting uniqueness.  In local coordinates with
$J-J_0=z^e$, or $1/J=z^e$ at a pole, the endpoint map is
$z\mapsto\zeta z$ and its inverse is $z\mapsto\zeta^{-1}z$,
where $\zeta^e=1$ is fixed by the sheet.  Hence the maps extend
holomorphically and preserve the positive peripheral orientation.

The degree-ten cover also determines the Weierstrass construction
directly.  Write $j=1728J_{56}$, $a_0=3$, let $b_1,b_2,b_3$ be the
roots of $P$, let $d_1,\ldots,d_5$ be the roots of $Q$, and put
$(p_5,p_3,p_2)=(\infty,0,-5/9)$.  The squarefreeness and coprimality
already checked make these points distinct.  As sections of
$\mathcal O(12)$,
\[
\begin{aligned}
 \operatorname{div}(\Delta)&=5p_5+3p_3+2p_2+2a_0,\\
 \operatorname{div}(j\Delta)&=3(a_0+b_1+b_2+b_3),\\
 \operatorname{div}((j-1728)\Delta)&=2(a_0+d_1+\cdots+d_5).
\end{aligned}
\]
The roots are explicitly $c_4=12F$ and $c_6=216H$:
$c_4^3=j\Delta$, $c_6^2=(j-1728)\Delta$.
More intrinsically, a section with either divided effective divisor
exists in $\mathcal O(4)$ or $\mathcal O(6)$ on $\mathbb P^1$;
its appropriate power differs from the required section by a nonzero
constant, which has the required root over $\mathbb C$.
All scalar choices are related explicitly.  For
$\Delta'=\gamma\Delta$, $c_4'=r_4c_4$, $c_6'=r_6c_6$ one has
$r_4^3=r_6^2=\gamma$.  Choose $\lambda^2=r_6/r_4$; then
$\lambda^4=r_4$ and $\lambda^6=r_6$.
The monic fibre map is $(x_m,y_m)\mapsto
(\lambda^2x_m,\lambda^3y_m)$, with inverse given by reciprocal
factors.  Pullback by the covering isomorphism above therefore
identifies this divisor construction with the explicit surface.
Its coefficient orders are never simultaneously at least $4,6$;
the singular fibres are exactly $II$ at $a_0$ and $I_m$ at $p_m$.

Use the projective period matrices $A,B$ and the positive
orientations of Remark~\ref{mapsIII:rem:positive-projective-periods}.
The elliptic projective local system is their pullback by $J$.
The tree words used for the covering map, in sheet order, are
\[
 (1,b,ba,ba^2),\qquad
 (1,a,a^2,b,ab,a^2b,ba,ba^2,aba,aba^2).
\]
They reach every displayed sheet, and reversing each path supplies
its inverse in the path-lifting construction.

Here is the modular cusp description of these same markings.
Set $S=-B$, $R=ST$ and
$G=\operatorname{PSL}_2(\mathbb Z)=\langle s,r\mid s^2=r^3=1\rangle$.
In the following right-action tables the lists give the images of
$0,1,\ldots$; the inverse permutations are $s$ and $r^2$:
\[
\begin{array}{c|l|l}
 &s&r\\\hline
 50&(1,0,3,2)&(2,1,3,0)\\
 56&(1,0,4,5,2,3,9,8,7,6)&(2,1,3,0,6,8,7,4,9,5).
\end{array}
\]
They satisfy the relations and are transitive, so the stabilizers of
$0$ have indices $4,10$ in $G$: the bijection is
$\Gamma g\mapsto0\cdot g$, with inverse obtained by any table
path from $0$ to the specified sheet.  Two such paths differ by
an element of the stabilizer, which proves independence and the
two inverse identities.  Their inverse images in
$SL_2(\mathbb Z)$ contain $-I$.  Compactification uses $z^e$ at
elliptic points and $q=\exp(2\pi i\tau/w)$ at a width-$w$ cusp;
these charts apply to the displayed finite-index subgroups
\cite[pp.~35--38]{mapsMilne2017}.  Their cycles are
\[
\begin{array}{c|l|l}
 &50&56\\\hline
 s&(01)(23)&(01)(24)(35)(69)(78)\\
 r&(023)(1)&(023)(1)(467)(589)\\
 sr&(012)(3)&(01265)(384)(79).
\end{array}
\]
Each has one elliptic point of order three and none of order two.
The cusp widths are $(3,1)$ and $(5,3,2)$, and the ramification
contributions over $0,1728,\infty$ are $(2,2,2)$ and $(6,5,7)$.
Riemann--Hurwitz gives genus zero in each case.
For the degree-four passport, the three matchings are also one orbit
under $b\mapsto rbr^{-1}$ with $r=(023)(1)$; the inverse is
conjugation by $r^{-1}$.  This gives the explicit relabellings
while fixing the order-three permutation.

The exact change of projective generators is
$a=sr^{-1}s^{-1}$, $b=s$, with inverse
$s=b$, $r=b^{-1}a^{-1}b$.
For $50$ the induced sheet bijection to the preceding $a,b$ table
is the identity.  For $56$ it and its inverse are
\[
 (5,2,4,3,1,0,6,7,8,9),\qquad
 (5,4,1,3,2,0,6,7,8,9).
\]
Checking each image against both transition lists proves equivariance.
In degree ten the forward bijection takes the cusp-table base sheet
to sheet $5$; the tree
path from the preceding base sheet is $a^2b$, whose literal period
matrix is $-T$.  Reverse paths and the inverse sheet list give the
inverse covering and basepoint changes.  At the branch endpoints
the maps extend by the same root-of-unity charts already described.

For the cusp graph put black vertices at the $r$-orbits, white
vertices at the $s$-orbits, and a subdivided edge with midpoint $q$
for each sheet $q$.  Use the clockwise $r$-order at black vertices
and remove the cusp discs and the fixed black vertex's disc.
This is the lifted arc between the two elliptic values.  Its
complementary faces are cusp discs: the local maps $z\mapsto z^e$
fill the smooth elliptic vertices, and the cusp chart fills each
boundary cycle.  After the indicated removals, radial collapse in
the faces retracts the surface onto the subdivided graph with a
small circle around the fixed black puncture.  Inclusion of this
graph and the radial collapse are inverse on based homotopy classes.
The positive cusp word is $sr=T$ projectively, and the positive
fixed-black meridian is $r^{-1}$: the derivative of $R$ at its
upper-half-plane fixed point is $e^{-2\pi i/3}$, while the cusp
coordinate identifies $\tau\mapsto\tau+1$ with one positive nodal turn.
Stems from $0$ to $1,3,7$ are $s,r^{-1},rsr^{-1}$,
with the last used in degree ten.
The corresponding integral matrices are
\begin{equation}\hypertarget{tag-02OE}{}\label{mapsD:repair:canonical-matrices}
\begin{aligned}
 H_c&=SR^{-1}S^{-1}=\begin{pmatrix}0&1\\-1&1\end{pmatrix},&
 V_c&=R^{-1}TR=\begin{pmatrix}2&1\\-1&0\end{pmatrix},\\
 B_c&=V_c^3=\begin{pmatrix}4&3\\-3&-2\end{pmatrix},&
 C_c&=(RSR^{-1})T^2(RSR^{-1})^{-1}
       =\begin{pmatrix}5&2\\-8&-3\end{pmatrix}.
\end{aligned}
\end{equation}
Orient half-edges from midpoints to black or white endpoints.
All half-edges except $3\to w_1$ form the degree-four tree;
in degree ten omit $8\to w_4,9\to w_3$.
Here white vertices are numbered in the order of the displayed
$s$-cycles.  The trees have respectively $8,19$ vertices and
$7,18$ edges and are connected.  Contracting them leaves chord
$\alpha$ or chords $\alpha,\beta$, in the order just stated,
and the clockwise loop $e$ at the fixed black puncture.
Insertion of the unique tree paths is the inverse on based loops.
Every $r$-step has trivial word except the fixed step, which reads
$e$; an $s$-step reads its outgoing omitted half-edge, if present,
followed by the inverse incoming omitted half-edge.  This gives
the following free words.  In this display each matrix symbol names
its specified based peripheral loop, before applying projective monodromy:
\begin{equation}\hypertarget{tag-02OH}{}\label{mapsD:repair:free-peripherals}
\begin{array}{c|ccc}
 50&T^3=e\alpha^{-1}&V_c=\alpha&H_c=e^{-1}\\
 56&T^5=e\beta^{-1}&B_c=\alpha&C_c=\alpha^{-1}\beta,\quad H_c=e^{-1}.
\end{array}
\end{equation}
The cycles $(012),(3)$, and respectively $(01265),(384),(79)$,
give these words directly.  Their products freely cancel to the
identity.  The inverse free generators are
$\alpha=V_c,e=H_c^{-1}$ in degree four and
$\alpha=B_c,\beta=B_cC_c,e=H_c^{-1}$ in degree ten.
Separating the tree stems inside its thin neighbourhood supplies
disjoint peripheral stems with this cyclic order.

Cyclically order the degree-four peripherals as $(H_c,T^3,V_c)$
and the degree-ten peripherals as $(C_c,H_c,T^5,B_c)$.
Uniformization on a simply connected set gives the period quotient
$\mathbb C/(\mathbb Z+\tau\mathbb Z)$ and its projective continuation
along the cover.  The Kodaira types fix the integral lifts:
$II$ has trace $1$, $I_m$ has the positive unipotent lift, and
$I_1^*$ has its negative.  Thus the two tuples are
$(H_c,T^3,-V_c)$ and $(C_c,H_c,T^5,B_c)$, both with product $I$.
Every peripheral sign is fixed, so the quotient of two such lifts
is a sign character trivial on a generating peripheral set and
therefore trivial everywhere.  The smooth filled faces have identity
monodromy.

The integral period-coordinate maps to the stated tuples are
$V_c$ and $B_c$, with inverse matrices
$(0,-1;1,2)$ and $(-2,-3;3,4)$, respectively.  Direct multiplication
therefore gives
\[
 (V_cH_cV_c^{-1},V_cT^3V_c^{-1},-V_c)=(A^{50}_1,A^{50}_2,A^{50}_3),
\]
\[
 (B_cC_cB_c^{-1},B_cH_cB_c^{-1},B_cT^5B_c^{-1},B_c)
 =(A^{56}_1,A^{56}_2,A^{56}_3,A^{56}_4).
\]
These are precisely the displayed integral matrices, with the stated
ordering of the singular fibres.  The sheet maps, path lifts and
period-coordinate changes give their marking on the explicit surfaces.

\emph{The complete integral section-cocycle quotient.}
Put $P_i=A_1\cdots A_{i-1}$, $d(h)=((A_i-I)h)_i$, and
$\Omega=\left(\begin{smallmatrix}0&1\\-1&0\end{smallmatrix}\right)$.
Use negative Abel--Jacobi increments.  The peripheral product gives
$\sum_iP_ic_i=0$, while adding an integral period $h$ to a logarithm
adds $(A_i-I)h$; subtraction is its inverse.  At $I_n$, the
local presentation is $\mathbb C^\times/(t^na(t))^{\mathbb Z}$ with
$a$ a unit.  A section representative $u=t^jb(t)$ has negative
increment $-j$ times the primitive invariant period.  The node maps
$U=u/t^j$, $V=t^{j+1}/u$, $UV=t$, and inverse $u=t^jU$ identify
the component charts.  Replacing $u$ by $u(t^na)^r$ gives the integral
coboundary $-nr$ in this period.  At $II$ and $I_1^*$ the matrix
$A-I$ has full rational rank, so its saturated image is all of
$\mathbb Z^2$.
We therefore use the integral lattice
\[
 \mathcal Z=\left\{(c_i):
 c_i\in\mathbb Z^2\cap\operatorname{im}_{\mathbb Q}(A_i-I),
 \quad\sum_iP_ic_i=0\right\},
 \qquad\mathcal Z/d(\mathbb Z^2),
\]
as in \eqref{mapsD:repair:cocycle-relation}.

The type-II differences have determinant one, with integral inverses
\[
 D_{50}=(A_1^{50}-I)^{-1}=\begin{pmatrix}-3&-7\\1&2\end{pmatrix},
 \qquad
 D_{56}=(A_2^{56}-I)^{-1}=\begin{pmatrix}-27&-37\\19&26\end{pmatrix}.
\]
For $50$, put $h=D_{50}c_1$ and $c^0=c-d(h)$.
Then $c_1^0=0$, and the remaining middle charge is
$k(-2,1)$, with inverse coefficient $k=(c_2^0)_2$.
The closing relation forces the last charge to be $k(2,-1)$,
since its prefix is unimodular and its local image has full rank.
Consequently
\[
 q^0=\bigl((0,0),(-2,1),(2,-1)\bigr),\qquad
 q^{50}=q^0+d(1,-1),\qquad
 c=d\bigl(h-k(1,-1)\bigr)+kq^{50}.
\]
This gives an integral assembly map and inverse for every closed
cocycle: $c$ returns $(h-k(1,-1),k)$, and $(j,k)$ returns
$d(j)+kq^{50}$.  Both composites are identity by the displayed gauge.

For $56$, put $h=D_{56}c_2$ and $c^0=c-d(h)$.
The remaining local images have primitive generators
$(-2,1),(4,-3),(1,-1)$; their inverse coefficient functions are
respectively the second coordinate, the sum of both coordinates,
and the first coordinate.  Their prefix transports are
\[
 (-2,1),\quad(1,0),\quad(1,-1).
\]
The first two have determinant $-1$.  The kernel of the resulting
$2$-by-$3$ matrix is exactly $\mathbb Z(1,1,1)$: its second row
equates the first and third coefficients, and its first then equates
the second.  Thus $c^0=kq^{56}$, where $k=(c_1^0)_2$.
The maps $c\mapsto(h,k)$ and $(h,k)\mapsto d(h)+kq^{56}$
are integral inverses.  In both profiles the quotient maps are
$[c]\mapsto k$ and $k\mapsto[kq]$, proving that the entire integral
cocycle quotient is infinite cyclic with the stated primitive class.

The quotient also has the following integral junction realization.
Let $L$ be the orthogonal complement of $\langle F,O\rangle$ in
integral $H_2$; its $(F,O)$ Gram matrix is
$(0,1;1,-1)$, of determinant $-1$.  Subtracting the integral
orthogonal projection gives the explicit inverse maps
\[
 [D]\longmapsto D-(D\cdot O+D\cdot F)F-(D\cdot F)O,
 \qquad D\in L\longmapsto[D],
\]
between $H_2/\langle F,O\rangle$ and $L$.
Let $V_{\mathrm{fib}}$ be generated by the components disjoint
from $O$.  Fukae's rational junction charges form $W_2$;
integral null sweeps form $U_1$, and their saturation is $U_2$.
The maps constructed in
\cite[Lemma 4.1, Theorem 4.2 and Corollaries 4.2--4.3,
pp.~21--22]{manFukae02} identify
\[
 W_2/U_1\simeq L/V_{\mathrm{fib}},\qquad
 W_2/U_2\simeq\operatorname{MW}_{\mathrm{free}}
       \text{ with opposite pairing},\qquad
 U_2/U_1\simeq\operatorname{MW}_{\mathrm{tors}}.
\]
Here is the exact charge and convention map in the present situation:
\[
 y_i=P_ic_i,\qquad K_i=P_iA_i^{-1}P_i^{-1},
 \qquad c_i=P_i^{-1}y_i,
\]
and list the pairs $(K_i,y_i)$ in reverse order.
Fukae's circuit is the inverse of ordinary monodromy
\cite[p.~3, footnote 1]{manFukae02}.
The original ordinary peripheral generators $g_1,\ldots,g_n$
are changed to
\[
 \bigl(p_ng_np_n^{-1},\ldots,p_1g_1p_1^{-1}\bigr),
 \qquad p_i=g_1\cdots g_{i-1},
\]
by the adjacent stem move $(g,h)\mapsto(ghg^{-1},g)$.
Its inverse is $(u,v)\mapsto(v,v^{-1}uv)$.
Use positions $(1,2,1)$ for three punctures and
$(1,2,3,1,2,1)$ for four; reversing these steps with the inverse
move recovers each original generator.  These moves take place in
a disc around two stems and preserve their positive peripheral
orientation.  Thus the reversed circuits describe peripheral stems.

Since $P_{i+1}=P_iA_i$ and $P_{n+1}=I$, multiplication gives
$K_1\cdots K_n=I$ and
$K_{i+1}\cdots K_n=P_{i+1}$.  The latter is exactly the product
of preceding circuits acting on a vector in Fukae's reversed list.
Also
\[
\begin{aligned}
 \operatorname{im}_{\mathbb Q}(I-K_i)
   &=P_i\operatorname{im}_{\mathbb Q}(A_i-I),&
 \sum_i y_i&=0,\\
 P_i(A_i-I)v&=(P_{i+1}-P_i)v
           =(I-K_i)P_{i+1}v.
\end{aligned}
\]
The unimodularity of $P_i$ makes both charge maps integral inverses
on the saturated images.  The last formula identifies integral
coboundaries with $U_1$ and rational $v$ giving integral charges
with $U_2$, exactly as in
\cite[Lemmas 3.1--3.2 and (11), pp.~11--13]{manFukae02}.

The geometric map agrees with the section's logarithm jumps.
Cut along the peripheral stems and contract the section to $O$
using its logarithm on the cut disc.  The swept three-chain has
the section difference as its horizontal boundary; its side
boundaries are the swept period cycles with charges $P_ic_i$.
Their total charge is zero, and local rational prongs fill the
rational images of $I-K_i$.  Keeping the integral charges gives
$W_2$, while a common logarithm shift gives the null sweep just
calculated.  Conversely Fukae's charge construction gives the
class in $L/V_{\mathrm{fib}}$ with these same jumps.
For a rational surface $H^1(\mathcal O)=H^2(\mathcal O)=0$;
the exponential sequence therefore identifies $H^2(\mathbb Z)$
with $\operatorname{Pic}$.  Restrict the corresponding divisor to
the generic elliptic fibre and subtract its degree times $O$.
The inverse of the Abel map gives its unique point; properness
extends this point over each base valuation ring to a section.
Its cut construction returns the original charge class.  In the
other direction, restricting the section divisor returns the
original point.  The vertical components and $F,O$ account
precisely for the quotient already taken.  This proves the full
integral section interpretation, including the torsion quotient;
the divisor-to-section map is also
\cite[Theorem~6.3, p.~23]{manSS10}.

\emph{The section generator and its height.}
The preceding integral isomorphism gives a section $G$ with class
$[q]$.  Its logarithm cocycle is $q+d(h_G)$ for an integral $h_G$;
changing the branch by $-h_G$ gives exactly $q$, with inverse branch
change $+h_G$.  The quotient maps prove that $G$ generates the
Mordell--Weil group and that this group has no torsion.

For closed cocycles $c,c'$ choose rational $z_i'$ with
$(A_i-I)z_i'=c_i'$, and put $\beta_i=-c_i^t\Omega A_i$.  In the
period-cocycle coordinates write the pairing as
\[
 \mathcal B(c,c')=-\sum_i\beta_i z_i'
 -\sum_{i<j}\beta_i A_{i+1}\cdots A_{j-1}c_j'.
\]
It is independent of the invariant-vector ambiguity in $z_i'$:
for $c_i=(A_i-I)v$ and $A_iw=w$, symplectic invariance gives
$c_i^t\Omega A_iw=0$.
Expanding the local terms with $(A_i-I)z_i=c_i$ also gives
\[
 \mathcal B(c,c')-\mathcal B(c',c)
 =\left(\sum_iP_ic_i\right)^t\Omega
  \left(\sum_iP_ic_i'\right)=0.
\]
Thus the diagonal formula determines this bilinear form by polarization.
For the diagonal pairing this is equivalently
\begin{equation}\hypertarget{tag-02OI}{}\label{mapsD:repair:height-formula}
 \langle[c],[c]\rangle_{\mathrm{MW}}
 =\mathcal B(c,c)
 =-\left\{\sum_i z_i\wedge c_i
       -\sum_{i<j}(P_ic_i)\wedge(P_jc_j)\right\},
 \qquad u\wedge v=u^t\Omega v.
\end{equation}
Indeed $-\beta_i z_i=c_i^t\Omega(z_i+c_i)=-z_i\wedge c_i$,
and symplectic invariance gives each cross term
$c_i^t\Omega A_i\cdots A_{j-1}c_j
=(P_ic_i)\wedge(P_jc_j)$.
In the junction coordinates a preimage of $y_i$ is
$P_i(z_i+c_i)$ under $I-K_i$, so its diagonal wedge is
$z_i\wedge c_i$.  Reversing the charge order reverses the sign of
the cross wedges.  This is Fukae's pairing (9), p.~10, with the
opposite sign for the Mordell--Weil height.
At a multiplicative place write
$A_i=I+n_i v(-v_2,v_1)$ and $c_i=t_i'v$, with primitive $v$.
Choose $w$ with $v\wedge w=1$.  The matrix $M=(v,w)$ has inverse
$\left(\begin{smallmatrix}w_2&-w_1\\-v_2&v_1\end{smallmatrix}\right)$
and satisfies $M^{-1}A_iM=T^{n_i}$ and
$M^{-1}c_i=(t_i',0)$.  A preimage in these coordinates is
$(0,t_i'/n_i)$, whose wedge
is $-t_i'^2/n_i$, the multiplicative term in (9).
The kernel and image of $A_i-I$ are the same isotropic line,
so changing this preimage leaves the wedge unchanged.

The fixed-point vectors and beta rows for the displayed cocycles are
\[
 \begin{array}{c|l|l}
 &z_i&\beta_i\\\hline
 50&(1,-1),(-4/3,0),(1/4,-3/4)&(1,4),(4,8),(-2,-1)\\
 56&(-1/2,0),(0,0),(1/15,0),(1/3,0)&(1,2),(0,0),(-3,-4),(-1,-1).
 \end{array}
\]
For $50$ the local terms sum to $3+16/3-1/4$ and the cross
terms are $-8,+8,-8$.  Equivalently, the prefix charges and
wedges are
\phantomsection\hypertarget{tag-033T}{}\label{qa:III:layout:2}
\[
\begin{gathered}
 (P_iq_i)=\bigl((-5,2),(4,0),(1,-2)\bigr),\\
 \sum_i z_i\wedge q_i=-3-\frac{16}{3}+\frac14=-\frac{97}{12},
 \qquad\sum_{i<j}(P_iq_i)\wedge(P_jq_j)=-8.
\end{gathered}
\]
The junction square is $-1/12$ and the height is $1/12$.
For $56$ the local terms are $1/2,0,1/5,1/3$, and the cross
terms are $0,-1,+1,0,0,-1$.  The alternative third preimage
$(-1/5,1/5)$ differs from $(1/15,0)$ by
$-(4,-3)/15$, which lies in $\ker(A_3-I)$; both give the
same local wedge.  Thus
\phantomsection\hypertarget{tag-033U}{}\label{qa:III:layout:3}
\[
\begin{gathered}
 (P_iq_i)=\bigl((-2,1),(0,0),(1,0),(1,-1)\bigr),\\
 \sum_i z_i\wedge q_i=-\frac12-\frac15-\frac13=-\frac{31}{30},
 \qquad\sum_{i<j}(P_iq_i)\wedge(P_jq_j)=-1.
\end{gathered}
\]
The junction square is $-1/30$ and the height is $1/30$.
The integral decomposition $c=kq+d(h)$ and the null-sweep formula
show that the pairing descends to the quotient.  In particular
$\langle nG,mG\rangle=nm/12$ or $nm/30$.
The mutually inverse group maps are $n\mapsto nG$ and the gauge
coefficient $P\mapsto k(c(P))$ given above.  This realizes the full
marked matrices, logarithm cocycles and heights in the statement.
\end{proof}

\begin{remark}[A second Weierstrass coordinate for the first surface]
\hypertarget{tag-Q103}{}\label{mapsIII:rem:collision-weierstrass-coordinate}
{\normalfont\scriptsize tag Q103}\par
An alternative coordinate $u$ for the first surface gives the model
\begin{equation}\hypertarget{tag-02OD}{}\label{mapsD:repair:50-weierstrass}
\begin{split}
 c_4&=u^2(u+27)(u+3),\qquad
 c_6=u^3(u+27)(u^2+18u-27),\\
 \Delta_o&=u^7(u+27)^2,\qquad
 y_o^2=x_o^3-c_4x_o/48-c_6/864.
\end{split}
\end{equation}
Here $c_4^3-c_6^2=1728\Delta_o$ by multiplication.  The exact
coordinate maps from the first model and their inverses are
\[
\begin{aligned}
 u&=\frac{27}{t-1},&t&=1+\frac{27}{u},\\
 x_o&=-\frac{81x}{2(t-1)^2},&x&=-\frac{2(t-1)^2x_o}{81},\\
 y_o&=\frac{729iy}{4\sqrt2(t-1)^3},&
 y&=\frac{4\sqrt2(t-1)^3y_o}{729i}.
\end{aligned}
\]
Substitution multiplies the first model's equation by
$-531441/[32(t-1)^6]$.  At $t=1$, use $s_o=1/u$ and
$(X_o,Y_o)=(s_o^2x_o,s_o^3y_o)$; the fibre map is
$(X_o,Y_o)=(-x/18,iy/(108\sqrt2))$.
At $t=\infty$, put $v=1/t$ and $(X,Y)=(v^2x,v^3y)$;
then $u=27v/(1-v)$ and the fibre factors are
$-81/[2(1-v)^2]$ and $729i/[4\sqrt2(1-v)^3]$.
They are units at $v=0$.  These formulas and their reciprocal
unit factors extend the isomorphism at both endpoints and preserve
the zero section.  They identify the local Weierstrass singularities
and their minimal resolutions.  Thus the orders at $u=0,-27,\infty$
are exactly $(2,3,7),(1,1,2),(0,0,3)$, and
\[
 j_o(u)=\frac{(u+27)(u+3)^3}{u},\qquad
 j_o(u)-1728=\frac{(u^2+18u-27)^2}{u},\qquad
 j_o\!\left(\frac{27}{t-1}\right)=1728J_{50}(t).
\]
\end{remark}


\subsection{Triangular caps and their integral homology}

\begin{definition}[Triangular cusp neighbourhood]
\hypertarget{tag-Q104}{}\label{mapsIII:def:triangular-cusp}
{\normalfont\scriptsize tag Q104}\par
Let \(L\subset\mathbb Z^2\) be a full sublattice of index \(D>0\).
Write \(B_0:\mathbb Z^2\longrightarrow\mathbb Z^2\) for an integer
basis matrix of \(L\), with positive determinant \(D\). The rays of
the three-dimensional toroidal fan are
\[
                         (v,1),\qquad v\in\mathbb Z^2.
\]
Its maximal cones are the cones over
\[
 [v,v+e_1,v+e_2],\qquad
 [v+e_1,v+e_1+e_2,v+e_2].
\]
Use the toroidal space of this fan, with the base character \(q\)
given by height in \(\mathbb Z^2\oplus\mathbb Z\).
The two period generators translate the vertex by the columns
of \(B_0\). On the dense torus their maps have the form
\[
             z\longmapsto q^{B_0 e_j}u_j(q)z,\qquad j=1,2,
\]
where \(u_j(q)\) is a pair of nonvanishing holomorphic units.
Multiplication and exponentiation here are componentwise.
For sufficiently small \(\epsilon>0\), let \(N_L\) denote the quotient
over \(|q|\leq\epsilon\) and put \(M_L=\partial N_L\).
The nearby torus has lattice \(\Lambda=\mathbb Z^4\), marked by
the two short angular circles and the two long period circles.
\end{definition}

\begin{theorem}[A triangular cusp and its boundary kernel]\hypertarget{tag-02LK}{}\label{maps:src:D.2.43}
{\normalfont\scriptsize tag 02LK}\par
The space \(N_L\) of Definition~\ref{mapsIII:def:triangular-cusp}
is a compact smooth neighbourhood of its central fibre.
Write \(s\mid t\) for the positive Smith factors of \(L\), so that
\(D=st\). Its nearby-torus return on \(\Lambda\) is \(A=I+N\), where
\(N^2=0\) and the long-to-short map has Smith factors \(s,t\).
Set \(F=\ker(A-I)\), the primitive short lattice.

Then \(\pi_1(N_L)=L\), with the long loops mapped to their period
displacements. Its integral homology is
\begin{equation}\hypertarget{tag-01YY}{}\label{maps:eq:D.GC1}
\begin{array}{c|ccccc}
 j&0&1&2&3&4\\ \hline
 H_j(N_L;\mathbb Z)&\mathbb Z&\mathbb Z^2&
 \mathbb Z^{D+3}\oplus\mathbb Z/s&\mathbb Z^2&\mathbb Z^D .
 \end{array}
\end{equation}
All higher groups vanish, and \(\chi(N_L)=2D\). The fibre kernel is
\begin{equation}\hypertarget{tag-01Z4}{}\label{maps:eq:D.GC2}
\ker\{H_2(T^4;\mathbb Z)\to H_2(N_L;\mathbb Z)\}
 =(\Lambda^2A-I)\Lambda^2\Lambda+\Lambda^2F.
\end{equation}
It is saturated exactly when \(s=1\). The fibre image is
\(\mathbb Z^4\oplus\mathbb Z/s\), and its cokernel in \(H_2(N_L)\)
is free of rank \(D-1\).

The two fixed angular normal tori bound and project to a basis of
\(F\) in the Wang sequence. Together with the fibre kernel they
give the full boundary kernel. In particular,
\begin{equation}\hypertarget{tag-01Z5}{}\label{maps:eq:D.GC3}
\begin{split}
 H_1(M_L;\mathbb Z)&=\mathbb Z^3\oplus\mathbb Z/s\oplus\mathbb Z/t,\\
 H_2(M_L;\mathbb Z)&=\mathbb Z^6\oplus(\mathbb Z/s)^2,\\
 \ker\{H_2(M_L;\mathbb Z)\to H_2(N_L;\mathbb Z)\}
                &\cong\mathbb Z^2\oplus\mathbb Z/s.
 \end{split}
\end{equation}
The killed torsion is the primitive angular area modulo the
monodromy image; the other boundary torsion coordinate maps
isomorphically to the cap torsion.

Let \(\theta\) be the positive \(q\)-meridian. Index the normalized
components by \(\bar v\in\mathbb Z^2/L\). The positive transverse
disc to that component has boundary
\begin{equation}\hypertarget{tag-01Z6}{}\label{maps:eq:D.GC4}
\theta+\bar v
 \quad\hbox{in }\mathbb Z\theta\oplus\mathbb Z^2/L
                  \subset H_1(M_L;\mathbb Z).
\end{equation}
These are all columns of
\(H_2(N_L,M_L)=H^4(N_L)=\mathbb Z^D\to H_1(M_L)\).
The component differences map to their differences in
\(\mathbb Z^2/L\).
\end{theorem}


\begin{proof}
\medskip
\noindent\emph{Fan charts, covers and units.}
The determinant of the three ray columns of either cone is one.
Thus its chart is smooth. In that chart the base character \(q\)
is the product of the three ray coordinates. Along a face the
same assertion has the corresponding invertible coordinates.
The inverse chart map is the character map from the inverse
integer ray matrix. These character maps also fix the boundary markings.

The ray at \(v\) gives a normalized central component \(S_v\).
The primitive outgoing directions around it are
\[
 e_1,e_2,e_2-e_1,-e_1,-e_2,e_1-e_2.
\]
Its fan is the smooth hexagon fan. The surface \(S_v\) is the
blowup of \(\mathbb P^2\) at its three coordinate points.
The two period maps of Definition~\ref{mapsIII:def:triangular-cusp}
commute and extend by the displayed character maps.
Their inverse uses \(-B_0 e_j\) and \(u_j(q)^{-1}\).

A nonzero period cannot fix a central point: it would preserve
the finite nonempty set of component vertices incident to that
point, whereas nonzero translation has no finite orbit in
\(\mathbb Z^2\). On \(q\ne0\), the real logarithms of the two
radial period vectors are independent for a sufficiently small
disc, uniformly for bounded unit factors. Consequently the
period action is free there too. A compact set in the toroidal
cover meets only finitely many central charts; translation of
their vertex labels proves properness near the central fibre.
Away from it properness is the ordinary discrete-period
calculation in logarithmic coordinates. The quotient is the
given smooth cusp, with inverse sheets specified by a chosen
period representative.

Here is a direct treatment of all units, including the outer
collar. Choose holomorphic logarithms \(b_j=\log u_j\) on the
base disc and replace \(u_j\) by
\[
                            u_j^\tau=\exp(\tau b_j),
                      \qquad 0\leq\tau\leq1.
\]
The same toroidal charts and the period maps above give a smooth
family over the closed interval. The vertex translation is
unchanged, so the preceding free-action and properness arguments
apply to the whole family. Each quotient has exactly \(D\)
component labels. Local charts are products with the interval;
their transition maps are the original monomials multiplied by
these units. Their inverses use the inverse integer matrices
and reciprocal units.

Compactness over a closed smaller disc can be checked in these
same coordinates. Divide the radial logarithms by \(\log|q|\)
and take a fundamental parallelogram for the period matrix
\(B_0+O(1/|\log|q||)\). These parallelograms lie in a fixed bounded
subset of \(\mathbb R^2\), uniformly in \(\tau\). Only finitely
many fan triangles meet that subset. Their character charts,
with compact phase coordinates, give convergent subsequences
as \(q\) tends to zero. For \(q\) bounded away from zero the
ordinary real-period parallelograms are compact. This proves
the required compactness of the whole quotient family.

For the explicit collar marking, choose a branch of \(\log q\)
and real period coordinates \(u,v\in\mathbb R^2\). Put
\[
 Z_\tau(q;u,v)=
 \exp\!\left(2\pi i u+(\log q)B_0v+
                    \tau\sum_jv_jb_j(q)\right).
\]
The real period matrix in this formula is invertible for small
\(|q|>0\): its radial leading block is \((\log|q|)B_0\), and the
unit correction is bounded. The inverse takes a logarithm of
\(z\), applies that real inverse matrix, and reads \(u,v\)
modulo their integer periods. Thus
\[
 Z_0(q;u,v)\longmapsto Z_\tau(q;u,v)
\]
is a collar diffeomorphism, with the indicated inverse.
Changing the branch of \(\log q\) changes \(u\) by \(B_0v\) on
both sides. The map therefore descends around the whole base
circle and fixes the marked short and long lattice coordinates.
It also fixes the origin section \(u=v=0\).

Extend the velocity of this collar map to a horizontal vector
field on the family. In each interior toroidal product chart
use a lift of \(\partial_\tau\) fixing its toric coordinates.
It preserves \(q\) and is tangent to every central stratum.
A partition of unity combines these lifts with the prescribed
collar lift; every summand still has \(\tau\)-component one,
has zero \(q\)-component and is tangent to each stratum on which
it is defined. Use a collar cutoff equal to one near the outer
boundary. Compactness of the quotient neighbourhood gives the
flow for the whole interval. Its time-one map is the required
oriented diffeomorphism, and its inverse is the same flow run
backwards. On the collar it is exactly the displayed period map.
It preserves the labelled component strata and positive normal
orientations. This proves the unit comparison with its inverse and boundary marking.

Changing a chosen logarithm by \(2\pi i k\) changes the collar
coordinates by the integral short shear \(u\mapsto u+kv\).
Its inverse subtracts that shear. The induced exterior-square
map is used on every fibre class. The unit-one calculations therefore transport through this
family map with their lattice and boundary markings.

These units also extend to the full short lattice. Put
\(K(q)=(b_1(q),b_2(q))\) and define, for \(v\in\mathbb Z^2\),
\[
 \Psi_v(z,q)=(q^v\exp(K(q)B_0^{-1}v)z,q).
\]
Then \(\Psi_{v+w}=\Psi_v\Psi_w\), and \(\Psi_{B_0e_j}\) is the
given period. The inverse is \(\Psi_{-v}\). The same character
charts extend this full action, whose leading logarithmic period
matrix is the identity; the preceding free-action and properness
proofs apply. The identity of covering charts induces the
degree-\(D\) map \(N_L\to N_{\mathbb Z^2}\).
Its local inverse sheets are labelled by \(\mathbb Z^2/L\) and
given by the maps \(\Psi_v\). Changing a logarithm changes this
extension by a finite unit character trivial on \(L\), and fixes
the given \(L\)-quotient and its boundary marking.

\medskip
\noindent\emph{The central fibre and branch resolution.}
Let \(Y_L\) be the central fibre and use the compact
tube \(|q|\leq\epsilon\), inside a slightly larger proper
open \(q\)-neighbourhood from the fan construction above. Its only boundary
is \(|q|=\epsilon\).
Choose a smooth Hermitian metric on the larger neighbourhood,
averaging it over the compact short torus. Put \({\rho_q}=|q|^2\).
In a normal-crossing chart \(q=z_1\cdots z_r\), \(r\leq3\),
the Euclidean gradient satisfies
\[
 \|\nabla {\rho_q}\|^2=4{\rho_q}^2\sum_{j=1}^r|z_j|^{-2}
       \ \geq\ 4r\,{\rho_q}^{\,2-1/r}.
\]
The inequality is AM--GM. On the finite compact chart cover,
comparison with the chosen Hermitian metric gives a uniform
constant \(c>0\) such that
\[
                        \|\nabla {\rho_q}\|\geq c {\rho_q}^{5/6}
                  \quad (0<{\rho_q}\leq\epsilon^2<1).
\]
An invertible local factor in \(q\) is absorbed into one
normal coordinate; the corresponding coordinate and metric
comparisons are uniformly bounded on this smaller tube.
There are no critical points off \(Y_L\).

On its complement use the vector field
\[
                         W_{\rm rad}=-\nabla {\rho_q}/\|\nabla {\rho_q}\|^2 .
\]
Its flow \(\Phi_t\) satisfies \({\rho_q}(\Phi_t(p))={\rho_q}(p)-t\).
At the outer boundary it points inward. For every
\(\eta>0\), the slab \(\eta\leq {\rho_q}\leq\epsilon^2\) is
compact, and this intrinsic vector field is smooth on it.
The finite atlas above gives the same uniform constant
throughout these slabs, so changing charts cannot end a
trajectory. It therefore exists until its remaining
\({\rho_q}\)-value tends to zero.
Its remaining length from value \(a\) to zero is at most
\[
                    \int_0^a c^{-1}r^{-5/6}\,dr
                                =6c^{-1}a^{1/6}.
\]
The path therefore has a unique limit in \(Y_L\). The same
bound is uniform and tends to zero with \(a\); continuous
dependence of finite-time flows, followed by this uniform
tail estimate, proves continuity of the limiting map.
At points already in \(Y_L\) define the map to be identity.

The explicit strong deformation retraction is
\[
 {\mathcal H}(p,\tau)=\Phi_{\tau {\rho_q}(p)}(p)\quad(0\leq\tau<1),
 \qquad {\mathcal H}(p,1)=\lim_{t\uparrow {\rho_q}(p)}\Phi_t(p).
\]
On \(Y_L\) it is constant in \(\tau\). The uniform length
bound proves joint continuity also at \(\tau=1\) and along
\(Y_L\). Its inverse homotopy-equivalence map is the central
inclusion: one composite is identity and the other is joined
to identity by \({\mathcal H}\).

The Hermitian identity \(dq(\nabla {\rho_q})=a(p)q\), \(a(p)>0\),
shows that this flow fixes \(\arg q\). Explicitly,
\[
 q(\Phi_t(p))=q(p)\sqrt{1-t/{\rho_q}(p)}.
\]
For \(0<\alpha<\beta\leq\epsilon^2\), the map
\(\Phi_{\beta-\alpha}:{\rho_q}^{-1}(\beta)\to {\rho_q}^{-1}(\alpha)\)
has inverse \(\Phi_{-(\beta-\alpha)}\), defined on that
positive-radius level and remaining in the compact slab.
This gives the exact radial collar and its inverse. Averaging
the metric made the flow short-torus equivariant. The same
construction applies to the given unit factors; alternatively
the unit family transports it with its prescribed collar marking.

For cellular calculations, triangulate the double curves with
their triple points as vertices, transport their cells by the
gluing maps, and extend across the transverse coordinate
disc neighbourhoods and compact component complements. The local branch construction in \ref{maps:src:D.1.2}
applies to these charts on the \(L\)-quotient.

For the normalization resolution, use the disjoint union of
the components in degree zero, of their double curves in degree
one, and of their triple points in degree two. At a point on
one, two or three local branches the augmented local complex
is the augmented simplex on those branches. Its contraction
adds a fixed branch as the first vertex, with the alternating
sign. It is therefore exact. This argument also applies when
different branches have the same label after quotienting by
\(L\): the branch positions remain distinct in the resolution.
Equivalently, insert the corresponding mapping cylinders;
their local branch simplices contract by this same formula.
This gives the normalization total complex for ordinary
integral homology of \(Y_L\), hence of \(N_L\).

There are \(D\) components, \(3D\) double curves and \(2D\)
triple points. Their nonzero homology groups are
\[
\begin{array}{c|ccc}
 &H_0&H_2&H_4\\ \hline
S_v&\mathbb Z&\mathbb Z^4&\mathbb Z\\
\text{double curve}&\mathbb Z&\mathbb Z&0\\
\text{triple point}&\mathbb Z&0&0 .
\end{array}
\]
The degree-zero row is precisely the triangulation of
\(\mathbb R^2/L\) defined above. The degree-four row
is \(\mathbb Z^D\). We compute the degree-two row with its integral relations below.

For the fundamental group, apply the same mapping-cylinder
branch resolution. Every component is simply connected and
every double curve is connected and simply connected.
The paths within a component change by nullhomotopies there;
the triple points impose exactly the triangular face
relations. Thus the van Kampen presentation is the edge-path
presentation of \(\mathbb R^2/L\). Its map records the
displacement of a lifted loop in the universal triangulation.
The inverse follows the marked paths for the two columns
of \(B_0\). The face relations and the component nullhomotopies
show that both composites fix the loop generators. Therefore
\[
                          \pi_1(N_L)=L\simeq\mathbb Z^2.
\]
The short angular loops vanish in their toric component.
A long period maps to its literal displacement \(B_0 e_j\).
The origin \(q\)-meridian bounds its transverse disc.
These specify the map from the entire boundary group.

\medskip
\noindent\emph{The normalization module.}
Let \(R=\mathbb Z[x^{\pm1},y^{\pm1}]\), \(I=(x-1,y-1)\),
and put \(a=x-1,b=y-1\). In the cyclic divisor order
\(({\mathcal E}_1,H-{\mathcal E}_1-{\mathcal E}_2,{\mathcal E}_2,H-{\mathcal E}_2-{\mathcal E}_3,{\mathcal E}_3,H-{\mathcal E}_1-{\mathcal E}_3)\),
take the second deck generator with reversed sign. The
opposite-pair weights are then \(x^{-1},y,xy\). The
coordinate map \(E=\operatorname{diag}(1,-1)\) sends
\(L\) to \(EL\), with component map \([v]\mapsto[Ev]\)
and inverse \([w]\mapsto[Ew]\). The resulting restriction
and cancelled matrices are
\begin{equation}\hypertarget{tag-01Z7}{}\label{maps:eq:D.GC5}
B=\begin{pmatrix}
 -x^{-1}&-1&-x^{-1}&-x^{-1}\\
 1&1&1&y\\
 -xy&-xy&-1&-xy
 \end{pmatrix}.
\end{equation}
\begin{equation}\hypertarget{tag-01Z8}{}\label{maps:eq:D.GC6}
\mathcal F=\begin{pmatrix}1-x&0\\0&xy-1\\y-1&xy(y-1)\end{pmatrix}.
\end{equation}
The three rows of \(B\) are the intersection restrictions of the
opposite boundary curves, with the displayed deck weights. The
intersection form in \((H,{\mathcal E}_1,{\mathcal E}_2,{\mathcal E}_3)\) is
\(\operatorname{diag}(1,-1,-1,-1)\). The deck involution sends \(x,y\) to \(x^{-1},y^{-1}\).
The homology row is
\(d=\overline B^{\,t}\) in the homology basis dual to that
cohomology basis. The diagonal matrix
\(J_4=\operatorname{diag}(1,-1,-1,-1)\), with inverse itself,
converts those coordinates to the divisor-class basis. The four curves \(D_0={\mathcal E}_1,D_1=H-{\mathcal E}_1-{\mathcal E}_2,D_2={\mathcal E}_2,D_3=H-{\mathcal E}_2-{\mathcal E}_3\)
also form an integral basis. Its inverse is
\({\mathcal E}_1=D_0,\ {\mathcal E}_2=D_2,\ H=D_0+D_1+D_2,\ {\mathcal E}_3=D_0+D_1-D_3\).

The cancellation is
\[
 U=\begin{pmatrix}
 1&0&0&0\\-1&1&0&0\\-1&0&1&0\\-y&0&0&1
 \end{pmatrix},
 \qquad
 V=\begin{pmatrix}0&x&0\\1&1&xy\\0&0&1\end{pmatrix},
 \qquad UB^tV=\operatorname{diag}(1,\mathcal F).
\]
Their determinants are \(1,-x\), so their adjugates give Laurent
integer inverses. The identity identifies the cancelled
coordinate, including its inverse representative.

Set
\[
 T=\begin{pmatrix}
 0&0&-x^{-1}\\
 -x^{-1}y^{-2}&x^{-1}y^{-1}&-x^{-2}y^{-1}\\
 -y^{-1}&0&0
 \end{pmatrix},\qquad
 S=\begin{pmatrix}y&xy^2\\x&0\end{pmatrix}.
\]
Let
\[
 C=\begin{pmatrix}-b&0\\a&-b\\0&a\end{pmatrix},
 \qquad g=(a^2,ab,b^2).
\]
Multiplication gives
\begin{equation}\hypertarget{tag-01Z9}{}\label{maps:eq:D.GC7}
T^{-1}\mathcal F=CS.
\end{equation}
The determinants are \(\det T=-x^{-2}y^{-2}\) and
\(\det S=-x^2y^2\).

Here are explicit maps for the homology differential, including
the involution. Put
\[
 {\mathsf D}_0=\operatorname{diag}(x^{-2},(xy)^{-1},y^{-2}),\quad
 {\mathsf E}_0=\operatorname{diag}(-x^{-2}y^{-1},-x^{-1}y^{-2}).
\]
Then \({\mathsf D}_0\overline C=C{\mathsf E}_0\).
Let \(p:R^4\to R^3\) omit the first coordinate and
\(\iota=p^t\). Define
\[
\begin{aligned}
 Q&={\mathsf D}_0\overline T^{-1}p\overline U,&
 \sigma&=\overline U^{-1}\iota\overline T {\mathsf D}_0^{-1},\\
 A_1&={\mathsf E}_0(0,\overline S)\overline V^{-1},&
 A_2&=\overline V
       \binom{0}{\overline S^{-1}{\mathsf E}_0^{-1}},\\
 {\mathsf H}&=\overline V e_0e_0^t\overline U .
\end{aligned}
\]
Here \(e_0\) denotes the first coordinate in its indicated rank,
so \({\mathsf H}\) has size \(3\times4\). Direct substitution proves
\[
\begin{gathered}
 Q\sigma=I_3,\quad \sigma Q=I_4-d{\mathsf H},\quad
 Qd=CA_1,\quad \sigma C=dA_2,\\
 A_1A_2=I_2,\qquad A_2A_1+{\mathsf H}d=I_3 .
\end{gathered}
\]
These are chain maps and both inverse chain homotopies. In
particular the quotient map to \(I^2\) is \(gQ\). Its inverse
takes a generator expression \(h=gz\) and uses \(\sigma z\)
modulo \(\operatorname{im}d\).

\(C\) gives every syzygy of \(g\).
If \(a^2p+abq+b^2r=0\), reduction modulo \(a\) gives \(r=ar_1\);
division by \(a\), followed by reduction modulo \(b\), gives
\(p=bp_1\). Then \(q=-ap_1-br_1\), so
\((p,q,r)^t=C(-p_1,r_1)^t\). The inverse relation coordinates are
\((-p/b,r/a)\). These are Laurent polynomials on the kernel.
The first and last coordinates also show \(\ker C=0\).
For \(h\in I^2\), an explicit generator expression is
\[
 z=\left(
 \frac{h(x,1)}{a^2},\
 \frac{\partial_yh(x,1)}a,\
 \frac{h(x,y)-h(x,1)-b\partial_yh(x,1)}{b^2}
 \right)^t.
\]
The augmentation and first-derivative vanishings make all three
entries Laurent integral. Changing this expression adds a
\(C\)-image, so the inverse quotient is independent of the expression.

The universal degree-zero branch row is the contractible plane,
and its degree-four row is \(R\). The degree-two row has injective
differential \(d\) and cokernel \(I^2\). The three-column
normalization resolution has no possible higher differential.
Thus the displayed morphism gives
\begin{equation}\hypertarget{tag-01ZA}{}\label{maps:eq:D.GC8}
H_2(\widetilde N_L;\mathbb Z)\cong I^2
\end{equation}
as an \(R\)-module.

Return to the original short coordinates by
\(\varphi(x)=x,\varphi(y)=y^{-1}\), carrying \(EL\) back to \(L\).
The component maps are the inverse reflections already displayed.
Set
\[
 D_r=\operatorname{diag}(1,-y^{-1},y^{-2}),\qquad
 E_r=\operatorname{diag}(-y^{-1},y^{-2}).
\]
Then \(D_r\varphi(C)=CE_r\) and \(gD_r=\varphi(g)\).
Thus \(D_r,E_r\) give the generator and relation maps from the
reflected presentation to the original one, with inverses
\(D_r^{-1},E_r^{-1}\); applying \(\varphi\) twice is identity.
The module isomorphism \eqref{maps:eq:D.GC8} therefore holds with
the original short-deck action as well. On an oriented lattice
basis \(u,v\), the reflection sends the oriented fundamental
class to its negative, and an oriented basis of \(EL\) is
\(Ev,Eu\). Accordingly its alternating tensor is
\(-(E\otimes E)(u\otimes v-v\otimes u)\), which specifies
the sign of this coordinate change on both Tor groups below.
All subsequent lattice coordinates use the original short basis.

\medskip
\noindent\emph{Integral finite-cover homology.}
Put \(G=\mathbb Z^2/L\), and let \(S_L=\mathbb Z[\mathbb Z^2/L]\). The Koszul complex
\[
 S_L\xrightarrow{(-b,a)}S_L^2\xrightarrow{(a,b)}S_L
\]
is the square-cell chain complex of \(\mathbb R^2/L\).
The map on \(H_1\) reads the signed displacement of a lifted
edge loop; its inverse follows an edge path for a vector of \(L\).
A zero-displacement loop lifts to a closed plane loop and is
filled by its square cells. The positive sum of all squares,
\(N_G=\sum_{v\in G}v\), gives the \(H_2=\mathbb Z\) generator.
These maps establish the Tor-to-\(H_*(L)\) identification.

The resolution \(R\xrightarrow{(-b,a)}R^2\to I\) shows that
\(I\otimes_R S_L\) is the edge-chain group modulo square boundaries.
Its boundary maps onto the free augmentation ideal of \(S_L\).
Its kernel is \(L\). Choose for each coset an edge path \(P_v\)
from zero to a fixed representative \(r_v\). Then
\[
 [z]\longmapsto
 \left(c=\partial z,\ 
       \operatorname{disp}\bigl(z-\sum_{v\ne0}c_vP_v\bigr)\right)
\]
is a map to the augmentation ideal plus \(L\). Its inverse uses
the indicated paths and a cycle for the lattice displacement.
The plane fillings show that the two composites differ only by
square boundaries. Thus \(I\otimes_R S_L\) is free of rank \(D+1\).

The map to \(I/I^2=\mathbb Z^2\) reads total horizontal and
vertical edge counts. In those coordinates it is
\[
                    (c,\ell)\longmapsto
                       \ell+\sum_{v\ne0}c_vr_v.
\]
The two edge generators at the origin provide a section. Its
kernel \(J_L\) is free of rank \(D-1\). More explicitly it is
the lattice of augmentation vectors satisfying
\(\sum c_v\bar v=0\) in \(G\); the inverse assigns
\(\ell=-\sum c_vr_v\). An integer Smith decomposition of this lattice constraint
gives its basis and inverse coordinates; its index in the
augmentation group is \(D\).

Use the exact sequences
\begin{equation}\hypertarget{tag-01ZB}{}\label{maps:eq:D.GC9}
\begin{gathered}
0\to I\to R\to\mathbb Z\to0,\\
0\to I^2\to I\to\mathbb Z^2\to0,
\end{gathered}
\end{equation}
where the last action is trivial. Their relevant Tor segment is
\begin{equation}\hypertarget{tag-01YZ}{}\label{maps:eq:D.GC10}
\mathbb Z\xrightarrow{\alpha_L}L\otimes\mathbb Z^2
 \longrightarrow I^2\otimes_R S_L
 \longrightarrow I\otimes_R S_L
 \longrightarrow\mathbb Z^2\longrightarrow0.
\end{equation}
The chain lift of \(I\to I/I^2\) in degree one is
\[
               1\longmapsto-e_y\otimes E_1+e_x\otimes E_2,
\]
whose Koszul boundary is \((-b,a)\).
Applying it to \(N_G\), the two norm edge cycles have respective
displacements \(D e_1,D e_2\). For an oriented basis \(u,v\) of
\(L\), the adjugate identities are
\[
 D e_1=v_2u-u_2v,\qquad D e_2=-v_1u+u_1v.
\]
Substitution gives exactly
\begin{equation}\hypertarget{tag-01Z0}{}\label{maps:eq:D.GC11}
\alpha_L(1)=u\otimes v-v\otimes u.
\end{equation}
The second factors here use the inclusion in the ambient lattice.
In the basis
\((u\otimes E_1,u\otimes E_2,v\otimes E_1,v\otimes E_2)\),
the column is \((v_1,v_2,-u_1,-u_2)^t\).
Its entries have gcd \(s\), so it is injective. This fixes its
sign from the positive square orientation.

Choose unimodular integer matrices with
\(U_\alpha\alpha_L V_\alpha=(s,0,0,0)^t\).
Then the quotient map reads the last three entries of
\(U_\alpha z\) and its first entry modulo \(s\); the inverse is
\(U_\alpha^{-1}(k,z_2,z_3,z_4)^t\). These compositions differ
exactly by \(\mathbb Z\alpha_L\), and
\(sU_\alpha^{-1}e_1=\alpha_L V_\alpha\).
For \(u=(s,0),v=(0,t)\), a direct inverse marking is
\[
 (z_1,z_2,z_3,z_4)\longmapsto
       (z_1,z_2+(t/s)z_3,z_4,\ [-z_3]_s).
\]
Its inverse representative is
\((r_1,r_2+(t/s)k,-k,r_3)\). Thus the exact quotient is
\(\mathbb Z^3\oplus\mathbb Z/s\).

For the middle arrow and the splitting, let \(\varepsilon\)
denote evaluation at \((1,1)\).
Write \(X^w=x^{w_1}y^{w_2}\),
\(A_u=X^u-1\), \(A_v=X^v-1\), and
\(\mathfrak a_L=(A_u,A_v)\). The unique coset decomposition
\[
 R=\bigoplus_{r\in G}X^{r}\mathbb Z[(X^u)^{\pm1},(X^v)^{\pm1}]
\]
has inverse given by the exponent vector
\((u,v)^{-1}(w-r)\). It proves
\(S_L=R/\mathfrak a_L\), with a regular sequence \(A_u,A_v\).
The middle connecting map has the positive product convention
\[
 \beta(\xi)=
 [A_u(\xi_1a+\xi_2b)+A_v(\xi_3a+\xi_4b)]
       \quad\hbox{in } I^2/\mathfrak a_LI^2 .
\]
This is obtained by lifting \(E_1,E_2\) to \(a,b\) in \(I\)
and taking their period boundary. If
\(R_w=A_w-w_1a-w_2b\in I^2\), then
\[
 \beta(\alpha_L)=[-A_uR_v+A_vR_u]=0.
\]

For a kernel element represented by \(h\in I^2\cap\mathfrak a_L\),
write \(h=A_uf+A_vg\) by the coset decomposition and Laurent
division. Its first derivatives vanish, so
\(\varepsilon(f)u+\varepsilon(g)v=0\); hence \(f,g\in I\).
The inverse to beta is
\[
 h\longmapsto
 [(\partial_xf,\partial_yf,\partial_xg,\partial_yg)(1,1)]
                    \quad\bmod\mathbb Z\alpha_L .
\]
Changing the decomposition by
\((f,g)\mapsto(f+A_vk,g-A_uk)\) adds
\(\varepsilon(k)\alpha_L\). Those are all changes by regularity.
Adding \(\mathfrak a_LI^2\) leaves the derivative class unchanged.
Subtracting the indicated linear parts from \(f,g\) leaves
elements of \(I^2\), proving both inverse composites.

The Laurent division itself has an explicit formula at every sign.
For \(t_u=X^u,t_v=X^v\), a monomial \(X^rt_u^it_v^j\) satisfies
\[
 X^rt_u^it_v^j-X^r
 =A_uX^r\frac{t_u^i-1}{t_u-1}t_v^j
       +A_vX^r\frac{t_v^j-1}{t_v-1}.
\]
Each quotient is the finite geometric sum for positive powers,
zero at power zero, and the negative sum of reciprocal powers
for negative powers. Extending linearly gives the required
decomposition and its integral inverse.

For \(c\in J_L\), let \(h_0=\sum c_rX^r\). Its first moment
\(w=\sum c_rr\) belongs to \(L\). Put
\((n_1,n_2)^t=(u,v)^{-1}w\), and define
\[
                    \sigma_J(c)=h_0-n_1A_u-n_2A_v.
\]
Its augmentation and both first derivatives vanish, so it belongs
to \(I^2\); its residue in \(S_L\) is exactly \(c\). This is an
integral section. The full forward map from \(I^2/\mathfrak a_LI^2\)
reads its residue \(c\), subtracts \(\sigma_J(c)\), and applies
the derivative inverse above. The inverse adds
\(\sigma_J(c)+\beta(\xi)\). Thus
\begin{equation}\hypertarget{tag-01Z1}{}\label{maps:eq:D.GC12}
I^2\otimes_R S_L
       \simeq J_L\oplus(\mathbb Z^4/\mathbb Z\alpha_L)
       \simeq \mathbb Z^{D+2}\oplus\mathbb Z/s
\end{equation}
with both inverse maps. Its finite coordinate is the image
under \(\beta\) of the primitive column \(\alpha_L/s\).


Finally \(\operatorname{pd}_R I=1\), so the preceding exact segment
is
\[
 0\to\operatorname{Tor}_2^R(\mathbb Z^2,S_L)
 \to\operatorname{Tor}_1^R(I^2,S_L)
 \to\operatorname{Tor}_1^R(I,S_L)
 \xrightarrow{\alpha_L}\operatorname{Tor}_1^R(\mathbb Z^2,S_L).
\]
Injectivity of \(\alpha_L\) makes the first arrow an isomorphism,
giving \(\operatorname{Tor}_1(I^2,S_L)=\mathbb Z^2\).
The connecting map has its usual explicit chain inverse: lift a
representative through the free resolutions, take its boundary,
and read the two coefficients of the positive norm 2-cycle.
Independence follows from the displayed exact segment.
Equivalently, in the \(C\)-resolution its kernel over \(S_L\) is
the two constant norm columns \(N_Ge_1,N_Ge_2\), with inverse
reading their common coefficients. Indeed the equations are
\(bf=0\), \(af=bg\), and \(ag=0\). Their common middle value
is invariant under both generators, hence is \(kN_G\); its
augmentation is zero, so \(Dk=0\) and \(k=0\). Thus both
\(f\) and \(g\) are constant norm multiples. The chain maps above carry these two generators to the
two degree-three normalization cycles.

\medskip
\noindent\emph{The cap homology.}
The normalization spectral sequence has only columns
zero, one and two and only even vertical rows. The second
differential has zero odd target; a higher differential
leaves these three columns. All its terms therefore persist.
The degree-zero row is the triangular torus \(\mathbb R^2/L\).
Subdividing each Koszul square into its two triangles, with
inverse cellular collapse, gives its groups \(\mathbb Z,L,\mathbb Z\).
The only remaining extension is a copy of \(\mathbb Z\) over
\eqref{maps:eq:D.GC12}, which splits integrally.

For an explicit splitting, first use units one. On each
toric component choose the closure of its positive real
torus and identify it with the hexagonal moment polygon.
The positive real interval on a double curve is identified
by the toric gluing. The hexagons therefore form the
dual cell decomposition of \(\mathbb R^2/L\). Their section
in \(Y_L\) lifts its oriented two-cycle. Equivalently,
choose the positive intervals between the two triple points
of each double curve and fill the resulting boundary
hexagon in the component's positive moment polygon.
This is the explicit total-complex lift. Subtract its
integer multiple and then use the normalization chain equivalence and
the integral quotient inverse above. That is the inverse
splitting in degree two. The unit-family flow transports
this entire chain and these maps.

Combining the normalization rows proves \eqref{maps:eq:D.GC1}.
Their total real dimension is four, so higher homology vanishes.
The displayed groups give \(\chi(N_L)=2D\).

\medskip
\noindent\emph{Boundary coordinates.}
The period exponential defined above identifies the boundary
with the mapping torus of the literal return
\[
 A=\begin{pmatrix}I&B_0\\0&I\end{pmatrix}
 \quad\hbox{on}\quad
 \Lambda=F\oplus\Lambda_{\rm long}.
\]
The inverse is the real period logarithm. The origin section
gives its positive meridian \(\theta\).
Choose integer Smith matrices \(U_L,V_L\) with
\[
               U_L B_0 V_L=\operatorname{diag}(s,t),\qquad s\mid t.
\]
Replace the short basis by \(U_L^{-1}\) and the long basis by
\(V_L\), retaining their inverse matrices. In the interleaved
basis \((f_1,\ell_1,f_2,\ell_2)\),
\[
                         A=I+sE_{12}+tE_{34}.
\]
This changes the marked periods through the real-period map
and its inverse, while retaining the same fan and component incidence.

In wedge order \(12,13,14,23,24,34\), put
\[
       a_0=e_{13},\qquad
       b_0=e_{14}+(t/s)e_{23}+t e_{13}.
\]
Expanding the exterior square of this return gives exactly
\begin{equation}\hypertarget{tag-01Z2}{}\label{maps:eq:D.GC13}
\operatorname{im}(\Lambda^2A-I)
                        =s\mathbb Z a_0+s\mathbb Z b_0.
\end{equation}
For example the images of \(e_{23}\), \(e_{14}\) and
\(e_{24}\) are \(s a_0\), \(t a_0\) and \(s b_0\).
Every other image is zero or their integer combination.
The six vectors
\[
                    a_0,b_0,e_{12},e_{23},e_{24},e_{34}
\]
are a unimodular basis: subtract \(t a_0+(t/s)e_{23}\)
from its second vector to recover \(e_{14}\).
This is its explicit inverse change of coordinates.

Thus the fibre coinvariant is
\[
 C_M=\mathbb Z^4\oplus(\mathbb Z/s)\langle a_0\rangle
                         \oplus(\mathbb Z/s)\langle b_0\rangle.
\]
The fixed tori \(f_i\times\theta\) give an integral
section of the degree-two Wang projection. The inverse
subtracts those tori using the projected two short
coordinates. Hence
\[
 H_2(M_L)=C_M\oplus F
       =\mathbb Z^6\oplus(\mathbb Z/s)^2.
\]
In degree one the same marked mapping-torus presentation
gives
\[
 H_1(M_L)=\Lambda_{\rm long}\oplus
                     \mathbb Z\theta\oplus(\mathbb Z^2/L)
       =\mathbb Z^3\oplus\mathbb Z/s\oplus\mathbb Z/t.
\]
The quotient map reads the long coordinates, the meridian,
and the short coordinates modulo \(B_0\mathbb Z^2\); its
inverse uses their marked loops. The short finite-coordinate
inverse uses the Smith representatives \(U_L^{-1}(i,j)\).

\medskip
\noindent\emph{The full boundary kernel.}
Use the central component at vertex zero and its toric normal
disc. Its positive boundary is the origin meridian \(\theta\).
The compact short torus acts on all toroidal charts and
commutes with their periods. Acting on that disc by a short
circle \(f_i\) gives a three-chain whose boundary is exactly
the fixed torus \(f_i\times\theta\). Its normal
orientation is the positive complex orientation.
Thus both fixed-normal tori bound, and their Wang
projections are the primitive basis of \(F\).

In the common universal covering, the preimage of a nearby fibre is
\(T^2\times\mathbb R^2\). A deck period translates its real long
coordinate and rotates its short torus, so its action on \(H_2\)
is identity, with the homotopy given by fractional translation.
The homology map into \(I^2\) is \(R\)-linear. An invariant element
of \(I^2\) obeys \((x-1)h=0\) in the Laurent domain, hence is zero.
Therefore the primitive angular area has zero image, and projection
gives the same nullhomology in every finite period cover.
We now determine the whole kernel, retaining the torsion.
Integral UCT, \eqref{maps:eq:D.GC1} and duality give
\begin{equation}\hypertarget{tag-01Z3}{}\label{maps:eq:D.GC14}
\begin{gathered}
H_3(N_L,M_L)=H^3(N_L)=\mathbb Z^2\oplus\mathbb Z/s,\\
H_2(N_L,M_L)=H^4(N_L)=\mathbb Z^D.
\end{gathered}
\end{equation}
Its boundary map followed by the Wang projection to \(F\)
is onto, because the two explicit three-chains above map
to the two basis vectors. A surjection from the first relative group to \(\mathbb Z^2\) has kernel exactly its torsion:
the free quotient has rank two and an onto integer
two-by-two map, whose inverse is integral.

Consequently, after subtracting the two marked normal
images, any element of the entire boundary kernel comes
from the torsion subgroup in the first relative group. Its order is at most
\(s\). On the other hand \(a_0\) has exact order \(s\)
in the fibre coinvariants, by the unimodular six-bivector coordinates, and it
already vanishes in the cap. It therefore accounts for
the entire remaining kernel. We have proved the exact
marked identity
\[
 \ker\{H_2(M_L)\to H_2(N_L)\}
 =\mathbb Z\langle f_1\times\theta,f_2\times\theta\rangle
                  \oplus(\mathbb Z/s)\langle a_0\rangle.
\]
The inverse reads the Wang coordinates and then the
angular torsion coordinate in the six-bivector basis. This argument also
includes \(s=1\), where the last summand is zero.

Since the nearby fibre maps into the \(C_M\) factor,
the return image and the boundary kernel give its entire kernel:
\[
 \ker\{H_2(T^4)\to H_2(N_L)\}
       =(\Lambda^2A-I)\Lambda^2\Lambda+\Lambda^2F.
\]
In the Smith coordinates it is
\(\mathbb Z a_0+s\mathbb Z b_0\).
Its inverse relation coordinates read the \(a_0\)
coordinate and divide the \(b_0\) coordinate by \(s\)
on this image. The quotient map kills \(a_0\), reduces
the \(b_0\) coordinate modulo \(s\), and keeps the four
free coordinates in the six-bivector basis. Its inverse uses those
six-bivector representatives. In particular the fibre
image is
\[
                   \mathbb Z^4\oplus\mathbb Z/s.
\]
Its torsion generator is the specialization of
\(b_0\). It has exact order \(s\) by the full boundary kernel, and so generates
the entire torsion of the cap. The other boundary torsion
generator \(a_0\) is killed. To compare this coordinate with \(\beta(\alpha_L/s)\),
use the normalization chain maps and branch augmentation,
then the boundary preimage rule below. The inverse inserts
the fibre representative in that same branch complex and
uses the period-ideal derivative inverse. Their composites
fix the respective cyclic classes, so these maps determine
the unit relating the two chosen cyclic coordinates.

For \(s>1\), the displayed fibre kernel is not saturated:
the class \(b_0\) is outside it but \(s b_0\) belongs
to it. For \(s=1\), the two displayed generators form
a primitive sublattice. This is the exact lattice
test, including the map exhibiting each assertion.

\medskip
\noindent\emph{Component discs and quotient inverses.}
Choose one lift \(v\in\mathbb Z^2\) of each component
label in \(G\). Its normal ray is \((v,1)\).
A character \((m,k)\) restricts to its positive
transverse disc as a nonzero unit times
\[
                             z^{\langle m,v\rangle+k}.
\]
Pairing the positive boundary circle with the height
and the two short characters therefore gives exactly
\[
                              \theta+\bar v
         \quad\hbox{in}\quad\mathbb Z\theta\oplus G.
\]
These pairings recover the ray and the positive disc
parameter, so they also give the inverse marking.
Replacing \(v\) by \(v+\ell\), \(\ell\in L\), changes
the representative by a short monodromy boundary.
A unit extends nonvanishingly across the disc and has
winding zero. Thus the component-meridian formula descends with all the original
period units retained.

The \(D\) fundamental component classes give the
integral basis of \(H_4(N_L)\). UCT and \eqref{maps:eq:D.GC1} give
\(H^4(N_L)=\mathbb Z^D\), with the dual basis of
positive transverse discs. Poincaré--Lefschetz
duality therefore identifies the whole relative map
as
\[
 \Gamma:\mathbb Z^G\longrightarrow
             \mathbb Z\theta\oplus G,\qquad
 (z_v)\longmapsto\left(\sum_vz_v,\sum_vz_v\bar v\right).
\]
Its long coordinates are zero. These columns give the integral relative map.

Here is an integral basis of its kernel with inverse
coordinates. Use Smith labels
\((i,j)\in\mathbb Z/s\oplus\mathbb Z/t\) and the
representatives \(U_L^{-1}(i,j)\). Let \(e_{ij}\) be
the component basis and \(w_{ij}=e_{ij}-e_{00}\).
Include \(s w_{10}\) when \(s>1\), and \(t w_{01}\)
when \(t>1\). For each other nonzero pair include
\[
                         w_{ij}-i w_{10}-j w_{01},
\]
omitting a term whose corresponding modulus is one.
There are \(D-1\) vectors.

For \(z\in\ker\Gamma\), the inverse reads \(z_{ij}\)
at the pairs in the residual-difference list, and reads
\[
               \frac{\sum_{ij}i z_{ij}}s,\qquad
               \frac{\sum_{ij}j z_{ij}}t
\]
at the two active special vectors. Divisibility is
exactly \(\Gamma(z)=0\). Substitution reconstructs every
coefficient, including \(z_{00}\) by the zero-sum
equation. This proves both inverse equations.
At \(D=1\) the kernel and this list are empty.

The corresponding absolute two-cycles can also be
specified. For \(s w_{10}\), take the indicated
difference of component discs and subtract the
boundary monodromy cylinder of the long vector
\(B_0^{-1}sU_L^{-1}e_1\), which is integral. For
\(t w_{01}\) use \(B_0^{-1}tU_L^{-1}e_2\). Its
short boundary is precisely the period in the component-meridian formula.
For a residual-difference vector, both the total meridian and the total
short representative are zero. The commuting short
and meridian loops bound their addition triangles
in the marked boundary torus. These triangles close
the indicated relative discs. Thus the construction
gives a splitting over every kernel basis
vector. Reading the displayed inverse coordinates, subtracting those closed cycles
and applying the fibre quotient coordinates gives
its inverse.

The full pair sequence and the computed boundary and disc maps now give
\[
 0\longrightarrow \mathbb Z^4\oplus\mathbb Z/s
 \longrightarrow H_2(N_L)
 \longrightarrow\ker\Gamma\longrightarrow0.
\]
The first arrow is the regular-fibre specialization,
with the quotient and inverse above. The
constructed closed cycles split the displayed pair sequence. Its
inverse can be implemented on the finite cellular chains
already used above. Read the relative coordinate of
an absolute two-cycle, use the inverse disc coordinates, and subtract the specified
closed cycles. The remaining relative class is zero.
Solve the integral equation
\[
            z-\sum_\alpha n_\alpha Z_\alpha
                       =\partial W+\iota(b)
\]
for a three-chain \(W\) and a boundary two-cycle \(b\).
Existence is exactly zero of that relative class; a reversible
Smith decomposition of this cellular boundary map chooses
an integer preimage, with inverse coordinates on its image.
Read the Wang coordinates of \(b\), discard the two normal
coordinates and \(a_0\), and read its \(b_0\) and four free
coordinates in the six-bivector basis. A different solution changes \(b\)
by precisely the full boundary kernel, so the resulting coordinate is independent
of the solution. This is the required inverse on the first
summand. It also gives the inverse comparison with the normalization
torsion coordinate specified above. The cokernel is therefore
free of rank \(D-1\).
\end{proof}




\subsection{Component intersections in triangular caps}

\begin{theorem}[Component intersections in a triangular cusp]\hypertarget{tag-02LO}{}\label{maps:src:D.2.47}
{\normalfont\scriptsize tag 02LO}\par
Let $N_L$ be the smooth triangular Mumford cap of Theorem~\ref{maps:src:D.2.43},
with period sublattice $L\subset\mathbb Z^2$ of finite index
$D>0$. Write $V=\mathbb Z^2/L$ for its vertices and
$E_i$, $i\in V$, for its reduced irreducible central divisors.
Their normalizations $\nu_i:P_i\to E_i\subset N_L$ are the
degree-six toric del Pezzo surface, even when a divisor has
self-gluings. Use the cyclic ray order
\begin{equation}
 (r_0,\ldots,r_5)=
       (e_1,e_2,e_2-e_1,-e_1,-e_2,e_1-e_2).                 \hypertarget{tag-0292}{}\label{maps:eq:D47.1}
\end{equation}
For a coefficient vector $u=(u_i)_{i\in V}$ put
$E(u)=\sum_i u_iE_i$ and
$a_{i,r}(u)=u_{i+r}-u_i$. The entire integral trilinear
form on these compact divisor classes is
\begin{equation}
\begin{split}
 I_L(u,v,w)=\sum_{i\in V}u_i\Bigl(
 &-\sum_{r=0}^5 a_{i,r}(v)a_{i,r}(w)\\
 &+\sum_{r=0}^5\bigl(a_{i,r}(v)a_{i,r+1}(w)
                    +a_{i,r+1}(v)a_{i,r}(w)\bigr)\Bigr),
\end{split}                                                   \hypertarget{tag-0297}{}\label{maps:eq:D47.2}
\end{equation}
where ray indices are cyclic. It is symmetric in all three
arguments, and the constant vector annihilates it.
Equivalently, the cubic is
\begin{equation}
\begin{split}
 I_L(u,u,u)={}&6\sum_{i\in V}u_i^3
       -3\sum_{\{i,j\}\in\mathcal E}
                      (u_i^2u_j+u_iu_j^2)\\
       &+6\sum_{(i,j,k)\in\mathcal T}u_iu_ju_k.
\end{split}                                                   \hypertarget{tag-0298}{}\label{maps:eq:D47.3}
\end{equation}
Here $\mathcal E$ is the set of the $3D$ quotient edge
\emph{cells}, indexed by $i$ and the first three directions
in \eqref{maps:eq:D47.1}, and $\mathcal T$ consists of the $2D$ triangle
cells. Loops, multiple edges and repeated vertex labels in
a triangle are retained with their actual cell multiplicities.
The Chern and Pontryagin evaluations on every component are
\begin{equation}
 \left\langle c_2(TN_L),\nu_{i*}[P_i]\right\rangle=0,
 \qquad
 \left\langle p_1(TN_L),\nu_{i*}[P_i]\right\rangle=0.         \hypertarget{tag-0299}{}\label{maps:eq:D47.4}
\end{equation}
All these formulas remain the same when the cap is included
in a closed smooth complex threefold.

For $L=\mathbb Ze_1\oplus d\mathbb Ze_2$, the entire
component trilinear form is zero for every $d>0$.
For $L=2\mathbb Z^2$ and its four vertices,
\begin{equation}
 E_i^3=6,\qquad E_i^2E_j=-2\ (i\ne j),\qquad
 E_iE_jE_k=2\ (i,j,k\text{ distinct}).                      \hypertarget{tag-029A}{}\label{maps:eq:D47.5}
\end{equation}
These formulas give the component products in the cohomology ring
and their indicated characteristic pairings.

\end{theorem}

\begin{proof}
We use the toric cover with lattice $\mathbb Z^2\oplus\mathbb Z$,
rays $(v,1)$ and cones on the two triangles in each unit square.
The quotient lattice at the ray $(v,1)$ is identified with
$\mathbb Z^2$ by
\[
 [(w,h)]\longmapsto w-hv,\qquad
 a\longmapsto[(a,0)].
\]
These inverse maps give the six rays in \eqref{maps:eq:D47.1}
on the normalization $P_i$. Denote its six boundary curves by
$C_{i,r}$. Consecutive primitive rays have determinant one, so
their dual monomials give a smooth chart in which the two curves
are coordinate axes. Nonadjacent boundary curves are disjoint.
Choose a character $m$ with $\langle m,r\rangle=1$.
Its principal divisor has coefficients $\langle m,r_s\rangle$;
its degree on $C_{i,r}$ is zero. Since
$r_{r-1}+r_{r+1}=r_r$, this gives
\begin{equation}
 C_{i,r}^2=-1,\qquad C_{i,r}C_{i,r+1}=1,\qquad
 C_{i,r}C_{i,s}=0\quad(s\not\equiv r,r\pm1).
 \hypertarget{tag-029B}{}\label{maps:eq:D47.6}
\end{equation}
The logarithmic two-form in these charts gives
$K_{P_i}=\mathcal O_{P_i}(-\sum_r C_{i,r})$.

On the dense torus of the covering threefold, put
\begin{equation}
 \omega=d\log X\wedge d\log Y\wedge dq.
 \hypertarget{tag-0293}{}\label{maps:eq:D47.10}
\end{equation}
A period $l\in L$ sends $(X,Y,q)$ to
$(c_1(q)q^{l_1}X,c_2(q)q^{l_2}Y,q)$, with nonvanishing
holomorphic units $c_1,c_2$. Its pullback adds multiples of $dq$
to the first two logarithmic differentials, so it fixes $\omega$.
For a height-one cone, let $B$ have its primitive rays as columns.
Its dual coordinates satisfy $q=z_0z_1z_2$, and direct exterior
multiplication gives
\begin{equation}
 \omega=\det(B)\,dz_0\wedge dz_1\wedge dz_2,
 \qquad \det(B)\in\{1,-1\}.
 \hypertarget{tag-0294}{}\label{maps:eq:D47.11}
\end{equation}
Thus $\omega$ descends to a nowhere-zero form on the cap.
Each $\nu_i$ is an immersion: a preimage of a crossing is one
coordinate branch in the toric cover. Its normal quotient line
maps to $K_{P_i}$ by
$[n]\mapsto((v,w)\mapsto\omega(n,v,w))$.
Adding a tangent vector to $n$ has zero effect. In the displayed
chart this map sends a transverse coordinate frame to a nonzero
two-form frame and has the inverse frame map. Hence
\begin{equation}
 0\longrightarrow TP_i\longrightarrow\nu_i^*TN_L
                  \longrightarrow K_{P_i}\longrightarrow0.
 \hypertarget{tag-0295}{}\label{maps:eq:D47.12}
\end{equation}

Lift a divisor chart to the cover and distinguish the branch over
$i$. Restricting its own divisor factor gives the normal transition
function; an adjacent branch gives the boundary curve with its
specified label. The local multiplication map, with division by
the same factors as inverse, gives
\begin{equation}
 \begin{split}
 \nu_i^*\mathcal O(E(u))
 &\simeq K_{P_i}^{\otimes u_i}\otimes
             \mathcal O_{P_i}\Bigl(\sum_r u_{i+r}C_{i,r}\Bigr)\\
 &\simeq\mathcal O_{P_i}\Bigl(
                      \sum_r(u_{i+r}-u_i)C_{i,r}\Bigr).
 \end{split}
 \hypertarget{tag-029C}{}\label{maps:eq:D47.7}
\end{equation}
The maps agree on chart overlaps because both products acquire
the same monomial unit. A period changes lifts and preserves all
coefficient labels. Every adjacent branch contributes its factor,
including a branch whose label is $i$. In particular
$\nu_i^*\mathcal O(E_i)\simeq K_{P_i}\otimes
\mathcal O_{P_i}(\sum_{r:i+r=i}C_{i,r})$.

Let $\rho:H^2(N_L,\partial N_L;\mathbb Z)\to H^2(N_L;\mathbb Z)$
be the natural map. A compact divisor defines a relative class
$e(u)$ with $\rho e(u)=c_1\mathcal O(E(u))$.
The product is evaluated as
$e(u)\smile\rho e(v)\smile\rho e(w)$ on the relative fundamental
class. The relative cup-product identities and their sign
$(-1)^{2\cdot2}=1$ give symmetry in the three arguments.
The normalization pushes forward with generic degree one, and
the projection formula gives
\begin{equation}
 I_L(u,v,w)=\sum_i u_i\int_{P_i}
       c_1\nu_i^*\mathcal O(E(v))\,
       c_1\nu_i^*\mathcal O(E(w)).
 \hypertarget{tag-029D}{}\label{maps:eq:D47.8}
\end{equation}
Equations \eqref{maps:eq:D47.6} and \eqref{maps:eq:D47.7} give
\eqref{maps:eq:D47.2}. The constant vector is the reduced divisor
of $q$; multiplication by $q$ trivializes its absolute divisor
line. Thus $\rho e(\mathbf1)=0$, and the displayed relative product
vanishes with this vector in every slot. The relative class itself
is the connecting image of the winding class of $q/|q|$ on the
boundary, through the long exact sequence of the pair.

To obtain the cubic, put $b_r=u_{i+r}$ and $a=u_i$. The local
summand expands to
\[
 6a^3-2\sum_r a^2b_r-\sum_r ab_r^2
                      +2\sum_r ab_rb_{r+1}.
\]
The directed-edge involution $(i,r)\mapsto(i+r,r+3)$ exchanges
the two mixed sums. It pairs the two directions of every edge
cell, including a loop, and gives its coefficient
$-3(u_i^2u_j+u_iu_j^2)$. The triangle cells are
\begin{equation}
 (i,i+e_1,i+e_2),\qquad
 (i+e_1+e_2,i+e_1,i+e_2)\quad(i\in V).
 \hypertarget{tag-029E}{}\label{maps:eq:D47.9}
\end{equation}
Each has three corner incidences; the last mixed sum counts each
with coefficient six. These incidence maps prove
\eqref{maps:eq:D47.3}, including every repeated label. Writing
$C(u)=I_L(u,u,u)$, the inverse relation is the integral identity
\[
 \begin{split}
 6I_L(u,v,w)={}&C(u+v+w)-C(u+v)-C(u+w)-C(v+w)\\
              &+C(u)+C(v)+C(w).
 \end{split}
\]
Expansion cancels the terms missing any argument and leaves the
six ordered mixed terms.

For $L=\mathbb Ze_1\oplus d\mathbb Ze_2$, the maps
$[(a,b)]\mapsto[b]\in\mathbb Z/d$ and $[j]\mapsto[(0,j)]$
identify the vertices. The ray images are $(0,1,1,0,-1,-1)$.
The difference vectors for $v,w$ have the forms
$(0,x,x,0,y,y)$ and $(0,x',x',0,y',y')$.
Their local pairing is $-2xx'-2yy'+2xx'+2yy'=0$ for every
positive $d$. For $L=2\mathbb Z^2$, reduction modulo two identifies
the vertices with $(\mathbb Z/2)^2$. Each pair of distinct vertices
has two edge cells, and each three-element vertex subset has two
triangle cells. The cubic coefficients are respectively $6,-6,12$;
division by the one, three and six ordered trilinear terms gives
\eqref{maps:eq:D47.5}.

The hexagonal toric surface has six point orbits; its
positive-dimensional torus orbits have Euler characteristic zero. Hence
$\int_{P_i}c_2(TP_i)=\chi(P_i)=6$, while
\eqref{maps:eq:D47.6} gives $\int_{P_i}c_1(TP_i)^2=-6+12=6$.
The tangent-normal sequence and the Whitney product formula give
\begin{equation}
 \nu_i^*c_2(TN_L)=c_2(TP_i)-c_1(TP_i)^2,\qquad
 \nu_i^*p_1(TN_L)=2c_1(TP_i)^2-2c_2(TP_i).
 \hypertarget{tag-0296}{}\label{maps:eq:D47.13}
\end{equation}
Both evaluations are zero. Finally, for the inclusion of the cap
in a closed smooth complex threefold $X$, its differential
identifies $TN_L$ with the restriction of $TX$, with inverse the
inverse coordinate differential. Divisor lines restrict through
the same inclusion, and the compact component cycles push forward.
Naturality therefore gives the same products and characteristic
evaluations in $X$.
\end{proof}


\subsection{Integral cohomology, Euler numbers and clutching}

\begin{definition}[Rank-one wheel]
\hypertarget{tag-Q105}{}\label{mapsIII:def:rank-one-wheel}
{\normalfont\scriptsize tag Q105}\par
A rank-one wheel is the cyclic gluing of finitely many ruled surfaces
over elliptic curves, with each infinity section identified with the
next zero section by an isomorphism of those elliptic curves. The zero
and infinity sections in each ruled surface are disjoint. This includes
one ruled surface with its two sections identified. The images of the
paired sections are called seams.
\end{definition}

\Needspace{6\baselineskip}
\begin{proposition}[Euler numbers and integral clutching]\hypertarget{tag-02L5}{}\label{maps:src:D.2.3}
{\normalfont\scriptsize tag 02L5}\par


For any of the semistable triangular rank-two Mumford fillings, its
reduced central fibre has Euler number twice the number of vertices
in the dual torus triangulation. A rank-one wheel over an elliptic
curve, a smooth torus, and a free finite quotient of a torus have
Euler number zero. Consequently the total Euler number is the sum
of these rank-two contributions whenever all other fillings are of
these types. There are two compact fibrations with the same local
fibre models and multiplicities whose fundamental groups are respectively
the trivial group and \(\mathbb Z/11\mathbb Z\). Their gluing maps differ
by an integral cusp translation.
\end{proposition}


\begin{proof}
Use compactly supported Euler characteristic on the open strata.
For a cone over a vertex, an edge or a triangle in the periodic
triangulation, the central-fibre orbit is respectively
$(\mathbb C^*)^2$, $\mathbb C^*$ or a point. A nonzero lattice
translation fixes no bounded cell, so the quotient map identifies
each chosen orbit isomorphically with its quotient stratum.
Consequently $\chi(W)=F$. Counting edge incidences with multiplicity
gives $3F=2E$, while the quotient torus gives $V-E+F=0$; hence
$\chi(W)=2V$.

The open component strata of a rank-one wheel are $\mathbb C^*$-bundles over elliptic curves,
and its seam strata are elliptic curves. Every boundary point
belongs to exactly one specified pair, including the one-component
case, so these are all its strata and all have Euler number zero.
A smooth torus has Euler number zero. For a free finite quotient,
the quotient projection is a finite covering and gives
$\chi(T)=|G|\chi(T/G)$, so its Euler number is also zero.

Over the complement of the finitely many filling values, the proper
smooth torus fibration is trivial over each open cell of a finite
base triangulation. Every such inverse-image cell has compactly
supported Euler number zero. Additivity therefore expresses
$\chi(X)$ as the sum of the reduced-fibre Euler numbers. Reduction
does not change topological support, so fibre multiplicities add
no factor. This proves the asserted total.

For the integral-clutching example, take the standard finite data
$(l_1,l_2)=(1,-1)$ and choose $c=0$. These data are admissible by \SeriesRef{fam:src:C.2.3}:
$3\nmid1$, $2\nmid-1$, and
$0<-M=(3/\pi)\log6$. In the marked cusp coordinates put
$q=e^{2\pi is}$, $Z(s)=sB_0+C(q)$ and $e_1=(1,0)^{\mathsf t}$.
For $n\in\mathbb Z$ the holomorphic vector
\[
 \widetilde\sigma_n(s)=
 n\left(\frac{s(s+1)}2B_0e_1+sC(q)e_1\right)
\]
satisfies
\[
 \widetilde\sigma_n(s+1)-\widetilde\sigma_n(s)
 =nZ(s+1)e_1.
\]
It therefore descends to a cusp collar section with winding
$n\overline\gamma$ and $l_0=n$, using the specified
$s\mapsto s+1$ winding convention. Replace the cusp arrow by
$G_0\circ T_{\sigma_n}$, retaining all local filling models.
Translation is a holomorphic base-preserving automorphism on the
overlap, so the local fibres and their multiplicities remain
identified by the identity maps of those models.

The marked peripheral presentation
\cite[Definition~6.6 and Proposition~6.7, pp.~39--40; Theorem~7.17, p.~57]{AC26} gives
\[
 p_n=12l_0-4l_1-3l_2=12n-1,
 \qquad \pi_1(X_n)\simeq\mathbb Z/|12n-1|\mathbb Z.
\]
Thus $X_0$ has trivial fundamental group and
$\pi_1(X_1)\simeq\mathbb Z/11\mathbb Z$. A homeomorphism induces
the isomorphism $[\alpha]\mapsto[h\circ\alpha]$ on based loop
classes, whereas every homomorphism from the one-element group
to $\mathbb Z/11\mathbb Z$ has image $\{0\}$. The two spaces
therefore have distinct homeomorphism types despite their
identified local fibre data. The integral clutching parameter is
therefore detected by the fundamental group.
\end{proof}

\begin{definition}[Finite marked-model input family]
\hypertarget{tag-Q106}{}\label{mapsIII:def:finite-filling-inputs}
{\normalfont\scriptsize tag Q106}\par
The finite input family \(\mathcal I\) uses the marked local cochain models of
\cite[documentation, ``How it works'']{manExplorer26}.
Their integral complexes, marked boundary maps and peripheral
relations form part of the input data.
Throughout the elliptic-profile tables, OS$j$ denotes entry $j$ of
\cite[Main Theorem table, p.~85]{manOguisoShioda91},
with its listed root lattice and Mordell--Weil lattice. For the
torsion-free rank-one entries used here, $G$ is a marked generator
and $n=1/\langle G,G\rangle$. The configurations and these
integers are:
\begin{center}\small
\begin{tabular}{c c p{9cm}}
OS&n&Configurations\\ \hline
43&2&\(III^*+3I_1\); \(III^*+II+I_1\)\\
45&8&\(I_8+4I_1\); \(I_8+II+2I_1\)\\
46&4&\(I_3^*+3I_1\); \(I_3^*+II+I_1\)\\
47&14&\(I_7+I_2+3I_1\); \(I_7+I_2+II+I_1\); \(I_7+III+2I_1\)\\
49&6&\(IV^*+I_2+2I_1\); \(IV^*+I_2+II\); \(IV^*+III+I_1\)\\
50&12&\(I_1^*+I_3+2I_1\); \(I_1^*+I_3+II\); \(I_1^*+IV+I_1\)\\
55&20&\(I_5+I_4+3I_1\); \(I_5+I_4+II+I_1\)\\
56&30&\(I_5+I_3+I_2+2I_1\); \(I_5+I_3+I_2+II\);
      \(I_5+I_3+III+I_1\); \(I_5+IV+I_2+I_1\)
\end{tabular}
\end{center}
For each listed configuration, take
\(P=G\), \(Q=nG\), a simple zero at the first \(I_1\) in its
marked fibre order, or at the multiplicative fibre of smallest index
when there is no \(I_1\), and the following two choices independently
at each additive fibre: its original filling, or the invariant
logarithmic vector \((0,1/m)\), where $m$ is its minimal
semistable-reduction order.
The set \(\mathcal I\) consists of all these marked choices.
\end{definition}

\Needspace{6\baselineskip}
\begin{computed}[Integral cohomology of a finite filling family]\hypertarget{tag-02KU}{}\label{maps:src:D.2.2}
{\normalfont\scriptsize tag 02KU}\par
The family \(\mathcal I\) of Definition~\ref{mapsIII:def:finite-filling-inputs}
has fifty inputs. For \(\iota\in\mathcal I\), write
\(\mathcal C_\iota^\bullet\) for its marked integral
Mayer--Vietoris complex and \(\Pi_\iota\) for the group defined
by its marked peripheral presentation. Both are specified in the
proof. The marked cochain and peripheral identifications give
\(H^*(\mathcal C_\iota)\) and \(\Pi_\iota\) as the integral
cohomology and fundamental group of the corresponding oriented
marked model. They have the following properties. The group \(\Pi_\iota\) is trivial for twenty-four
inputs and nontrivial for twenty-six.
Exactly the following three inputs have integral cohomology
\(H^0=H^6=\mathbb Z\), \(H^i=0\) for \(1\leq i\leq5\):
\[
\begin{array}{c|c|l}
\mathrm{OS}&\text{configuration}&\text{modified fibres}\\ \hline
43&III^*+3I_1&III^*\text{, of order four}\\
43&III^*+II+I_1&III^*\text{ only}\\
49&IV^*+III+I_1&IV^*\text{ and }III .
\end{array}
\]
Their groups \(\Pi_\iota\) are all trivial. The corresponding
homology groups are the degree-reversed cohomology lists, by the
cap-product map \(\alpha\mapsto\alpha\frown[M]\) for the oriented
six-dimensional models.
\end{computed}

\begin{proof}
For a configuration $C$, let \(A(C)\) be its set of additive fibre
occurrences. Repeated fibres are distinct occurrences. A vector
\(\epsilon\in\{0,1\}^{A(C)}\) chooses the original filling at a
zero entry and the logarithmic vector \((0,1/m)\) at a unit entry.
The orders are
\[
\begin{array}{c|rrrrrrr}
\text{type}&II&III&IV&III^*&IV^*&I_1^*&I_3^*\\ \hline
m&6&4&3&4&3&2&2 .
\end{array}
\]
They follow by taking powers of the Kodaira representatives
\[
 M_{II}=\begin{pmatrix}1&1\\-1&0\end{pmatrix},\qquad
 M_{III}=\begin{pmatrix}0&-1\\1&0\end{pmatrix},\qquad
 M_{IV}=\begin{pmatrix}0&-1\\1&-1\end{pmatrix}.
\]
The first two satisfy \(M_{II}^3=M_{III}^2=-I\), and the
third has order three. Their smaller positive powers have trace
different from two. The finite starred representatives are their
inverses. For \(I_j^*\), the representative is
\(-\left(\begin{smallmatrix}1&j\\0&1\end{smallmatrix}\right)\);
its odd powers have trace \(-2\) and its square is unipotent.
Thus the displayed integers are precisely the first unipotent powers.

The input set is the disjoint union
\[
 \mathcal I=\coprod_C\{0,1\}^{A(C)},\qquad
 |\mathcal I|=6+3+6+5+10+10+3+7=50,
\]
where the summands group the configurations by increasing OS entry.
The section coefficients, simple-zero rule and invariant coordinates
fix the remaining data of each input.

We specify the comparison between the marked fibre order and the
order printed in the configuration table. Match the first occurrence
of each type in the marked order with the first occurrence of that
type in the printed order, and continue in occurrence order.
This defines a permutation \(\pi\), from marked positions to printed
positions. The maps on the logarithmic vectors are
\[
 (\ell_i)_i\longmapsto(\ell_{\pi^{-1}(j)})_j,\qquad
 (\ell_j)_j\longmapsto(\ell_{\pi(i)})_i .
\]
Substitution shows that these maps are inverse. The simple zero is
chosen in the marked order before this relabelling.

In the table, \(\circ\) denotes the first configuration of an OS
entry above; the other profiles are named by their
additional additive fibre. All positions are numbered from one.
The integer \(\nu\) is the marked position of the simple zero.
The pair \((N_0,N_1)\) counts trivial and nontrivial
peripheral groups \(\Pi_\iota\), respectively. The last column gives all
sphere-cohomology choices, with bits in printed additive-fibre order.
\begin{center}\small
\begin{tabular}{c c c c c c}
OS&profile&\(\pi\)&\(\nu\)&\((N_0,N_1)\)&sphere choices\\ \hline
43&\(\circ\)&\((1,2,3,4)\)&2&\((1,1)\)&1\\
43&II&\((1,2,3)\)&3&\((2,2)\)&10\\
45&\(\circ\)&\((1,2,3,4,5)\)&2&\((0,1)\)&\(\varnothing\)\\
45&II&\((1,2,3,4)\)&3&\((1,1)\)&\(\varnothing\)\\
46&\(\circ\)&\((1,2,3,4)\)&2&\((1,1)\)&\(\varnothing\)\\
46&II&\((1,2,3)\)&3&\((2,2)\)&\(\varnothing\)\\
47&\(\circ\)&\((1,2,3,4,5)\)&3&\((0,1)\)&\(\varnothing\)\\
47&II&\((1,2,4,3)\)&3&\((1,1)\)&\(\varnothing\)\\
47&III&\((1,3,4,2)\)&2&\((1,1)\)&\(\varnothing\)\\
49&\(\circ\)&\((2,1,3,4)\)&3&\((1,1)\)&\(\varnothing\)\\
49&II&\((3,1,2)\)&3&\((2,2)\)&\(\varnothing\)\\
49&III&\((1,2,3)\)&3&\((3,1)\)&11\\
50&\(\circ\)&\((1,2,3,4)\)&3&\((1,1)\)&\(\varnothing\)\\
50&II&\((3,2,1)\)&2&\((2,2)\)&\(\varnothing\)\\
50&IV&\((1,2,3)\)&3&\((2,2)\)&\(\varnothing\)\\
55&\(\circ\)&\((1,2,3,4,5)\)&3&\((0,1)\)&\(\varnothing\)\\
55&II&\((1,2,3,4)\)&4&\((1,1)\)&\(\varnothing\)\\
56&\(\circ\)&\((1,2,3,4,5)\)&4&\((0,1)\)&\(\varnothing\)\\
56&II&\((3,4,1,2)\)&1&\((1,1)\)&\(\varnothing\)\\
56&III&\((1,2,3,4)\)&4&\((1,1)\)&\(\varnothing\)\\
56&IV&\((1,3,2,4)\)&4&\((1,1)\)&\(\varnothing\)
\end{tabular}
\end{center}

For each input, the section coefficients are \((1,n)\), the
linearization divisor has its sole unit entry at \(\nu\), and
every multiplicative logarithmic entry is null. At an additive
occurrence the entry is null or exactly \((0,1/m)\), as a rational
vector. The original fillings have multiplicity one. The modified
finite types have the indicated denominator and free quotient model;
the modified \(I_j^*\) types have semistable-quotient multiplicity
two. The weighted multiplicative fibre has its weight-one Mumford
model. Applying \(\pi\) carries these marked data to the
corresponding element of \(\mathcal I\); applying \(\pi^{-1}\)
recovers the ordered local data. The independent binary choices
at the additive occurrences give distinct data and exhaust
\(\mathcal I\). Thus these are exactly the fifty presentations
to be calculated.

Here is the integral presentation used in each case. Let
\(C_0^\bullet\) be the marked cochain complex of the regular
bundle, and let \(C_j^\bullet,B_j^\bullet\) be those of the
$j$-th filling and its boundary. Write
\[
 \mathsf R_j:C_0^\bullet\longrightarrow B_j^\bullet,
 \qquad \mathsf S_j:C_j^\bullet\longrightarrow B_j^\bullet
\]
for the integral restrictions. In \(\mathsf S_j\), the local restriction
is composed with its comparison to the common boundary marking,
including the prescribed clutching term. These are the marked
pairs of \cite{manExplorer26}. They include the original additive
fillings, finite good-reduction quotients, multiplicative Mumford
fillings, and the twofold semistable quotients at \(I_1^*,I_3^*\).
For the last two, the local cochain ranks in degrees zero to four
are respectively \((1,2,2,3,2)\) and \((1,2,4,7,4)\); the
boundary ranks are \((1,3,4,4,3,1)\). Their full boundary
comparisons are retained in \(\mathsf S_j\).

Put
\[
 \mathcal A^\bullet=C_0^\bullet\oplus\bigoplus_j C_j^\bullet,
 \qquad \mathcal B^\bullet=\bigoplus_j B_j^\bullet,
 \qquad \rho(a_0,(a_j)_j)=(\mathsf R_j a_0-\mathsf S_j a_j)_j.
\]
In the concatenated integral bases the Mayer--Vietoris complex is
\[
 \mathcal C_\iota^i=\mathcal A^i\oplus\mathcal B^{i-1},
 \qquad
 D_i=\begin{pmatrix}d_{\mathcal A}^i&0\\
                   \rho^i&-d_{\mathcal B}^{i-1}\end{pmatrix}.
\]
The cochain-map identities give
\(\rho^{i+1}d_{\mathcal A}^i=d_{\mathcal B}^i\rho^i\),
so multiplication gives \(D_{i+1}D_i=0\).

Choose an integral kernel basis \(Z_i\) for \(D_i\), with
coordinate map \(L_i\). In a Smith factorization
\(P_iD_iQ_i\), let \(t_i=\operatorname{rank}D_i\).
Take for \(Z_i\) the columns of \(Q_i\) after column \(t_i\),
and for \(L_i\) the corresponding rows of \(Q_i^{-1}\). Thus
\[
 L_iZ_i=I,\qquad Z_iL_i x=x\quad(x\in\ker D_i),\qquad
 A_i=L_iD_{i-1},\qquad D_{i-1}=Z_iA_i.
\]
The nonzero diagonal entries of the Smith form of \(D_i\)
force all other coordinates of a cycle to vanish, which proves
these identities over \(\mathbb Z\). A Smith form
\(U_iA_iV_i\) therefore gives
\[
 H^i(\mathcal C_\iota)=\operatorname{coker}A_i
 =\mathbb Z^{r_i}\oplus
       \bigoplus_{j=1}^{s_i}\mathbb Z/d_{ij},
 \qquad d_{i1}\mid\cdots\mid d_{is_i},\quad d_{ij}>1.
\]
The map from a cycle \(x\) is given by the free and nonunit
Smith coordinates of \(U_iL_i x\), with reduction modulo
\(d_{ij}\) in the latter coordinates. Its inverse sends the
corresponding coordinate vector \(e_j\) to the class of
\(Z_iU_i^{-1}e_j\). The two compositions differ from the
identity only by the displayed boundary relations.

For all fifty complexes, these reductions give
\(H^0=H^6=\mathbb Z\). The Euler characteristic is
\(\sum_i(-1)^ir_i\), the ranks in complementary degrees agree,
and the invariant factors in degrees \(i\) and \(7-i\) agree
for \(1\leq i\leq6\).
The abelianization of the peripheral presentation has free rank
\(r_1\) and the invariant factors of \(H^2\), in agreement
with the integral universal coefficient sequence.

For the group calculation, let \(e_1,\ldots,e_4\) be the
commuting fibre generators, put \(e^v=\prod_a e_a^{v_a}\),
and let \(x_j\) be the marked peripherals. The local data give
monodromies \(T_j\), vanishing lattices \(W_j\), multiplicities
\(\mu_j\) and peripheral periods \(v_j\). The presentation is
\[
 \begin{gathered}
 x_je^u x_j^{-1}=e^{T_ju},\qquad e^w=1\ (w\in W_j),\\
 x_j^{\mu_j}=e^{v_j},\qquad x_1\cdots x_r=1.
 \end{gathered}
\]
Close the vanishing lattices, together with
\((T_j-I)\mathbb Z^4\) at the multiplicity-one sites, under
all \(T_j^{\pm1}\), and call the resulting lattice \(L\).
Integral reduction gives \((T_j-I)\mathbb Z^4\subset L\)
for every site in all fifty cases. Hence the fibre subgroup is
central. There are at most two peripherals of multiplicity
greater than one. The other peripherals lie in the fibre
subgroup, and the product relation expresses either remaining
peripheral in terms of the other and a central element. Thus
\(\Pi_\iota\) is abelian, and the Smith form of the exponent-sum
relation matrix computes the full group. It gives twenty-four
trivial groups, twenty-one infinite cyclic groups, and one each of
\[
 \mathbb Z/10,\quad\mathbb Z/4,\quad\mathbb Z/3,\quad
 \mathbb Z/8,\quad\mathbb Z/5.
\]
The five finite cases are respectively profiles
\(43/II,46/II,49/II,50/II,50/IV\), with both additive
fibres modified. Applying the exact test
\[
 (r_0,s_0)=(r_6,s_6)=(1,0),\qquad
 (r_i,s_i)=(0,0)\quad(1\leq i\leq5)
\]
gives precisely the three choices in the last column of the table.
These are precisely the three inputs with the integral cohomology
of a six-sphere, and their groups \(\Pi_\iota\) are trivial.
In these three graded groups the
unit and the degree-six generator determine the whole unital ring,
since the square of the latter has degree twelve and is zero.

The Euler characteristic is two for forty inputs, four for six,
and six for four. The six inputs with Euler characteristic four
are the choices in profiles \(49/II\) and \(56/II\); their
simple zero lies at \(I_2\). The four inputs in profile
\(50/II\) have their simple zero at \(I_3\), giving Euler
characteristic six. These values also follow by the stated
alternating-rank calculation.

The integral clutch comparison has four pairs, with
\[
 (\mathrm{OS},n)=(45,8),(47,14),(55,20),(56,30).
\]
In each pair use the first configuration of the corresponding
OS entry, take \(P=nG\), \(Q=G\), place the simple zero at its
fifth fibre, and add a smooth sixth fibre. The first five
logarithmic entries are null. At the smooth fibre take the ambient
vector \((0,0,0,c)\), for \(c=0,1\). Addition of the fourth
ambient basis vector to this final entry, fixing all other data,
carries the first input to the second; subtraction is its inverse.
The complete local-model data agree within each pair. At every
\(I_j\) with \(j>1\) there is one component orbit of length $j$,
with the same local \(I_1\) model and zero $P$ component; the
smooth sixth fibre has the product model.

For each of the four values of $n$, the same cone and
peripheral presentation give
\[
\begin{array}{c|c|c}
c&\pi_1&(H^0,H^1,H^2,H^3,H^4,H^5,H^6)\\ \hline
0&\mathbb Z&
 (\mathbb Z,\mathbb Z,\mathbb Z^2,\mathbb Z^2,
  \mathbb Z^2,\mathbb Z,\mathbb Z)\\
1&1&(\mathbb Z,0,0,0,0,0,\mathbb Z).
\end{array}
\]
Thus the four unit-clutch inputs have the integral cohomology
of a six-sphere and trivial peripheral group. Their section coefficients \((n,1)\), with \(n>1\),
give a disjoint input set from the coefficient pairs \((1,n)\)
defining \(\mathcal I\).

The marked rational elliptic surfaces underlying profiles
\(50/II\) and \(56/II\), with fibre orders
\((II,I_3,I_1^*)\) and \((I_2,II,I_5,I_3)\), are constructed
in \cref{mapsD:repair:collision-realization}.
\end{proof}


\section{Logarithmic elliptic profiles}

\subsection{The \texorpdfstring{$IV^*+III+I_1$}{IV*+III+I1} profile}

\begin{definition}[Marking for two finite-monodromy fibres]
\hypertarget{tag-Q107}{}\label{mapsIII:def:two-finite-fibre-marking}
{\normalfont\scriptsize tag Q107}\par
The marked lattice for the $IV^*+III+I_1$ profile is
$\Lambda=\mathbb Z\langle\delta,a,b,d\rangle$. Its homological
monodromies and distinguished vectors are
\[
 A_0=\begin{pmatrix}1&-1&0&2\\0&-1&-1&6\\0&1&0&0\\0&0&0&1\end{pmatrix},
 \qquad
 A_1=\begin{pmatrix}1&0&1&-3\\0&0&1&0\\0&-1&0&6\\0&0&0&1\end{pmatrix},
\]
\[
 A_\infty=\begin{pmatrix}1&0&0&1\\0&1&1&0\\0&0&1&0\\0&0&0&1\end{pmatrix},
 \qquad
 u_0=d+2a+2b,\quad u_1=d+3a+3b.
\]
These are the marking and local models of
\cite[Lemma~2.1 and equations (2.11)--(2.12), pp.~3--4;
Propositions~2.3 and 3.1, pp.~5--7]{manViaclovskyLogs26}.
\end{definition}

\Needspace{6\baselineskip}
\begin{theorem}[Logarithmic multiplicities over $IV^*+III+I_1$]\hypertarget{tag-02LR}{}\label{maps:src:D.2.5}
{\normalfont\scriptsize tag 02LR}\par


Use the marking of Definition~\ref{mapsIII:def:two-finite-fibre-marking}
for the $IV^*,III,I_1$ configuration, its rank-two Mumford filling,
and the logarithmic data of
\cite[\S2.1, equations~(2.11)--(2.12), p.~4]{manViaclovskyLogs26}:
\[
 m_0=3k,\qquad m_1=4\ell,\qquad
 v_i=r_i\delta+n_i u_i,\qquad
 \gcd(n_i,m_i)=1,\qquad k,\ell\in\mathbb Z_{>0}.
\]
Here $r_i,n_i$ range over all integers. Put
$P=m_1n_0+m_0n_1$. The parameter ranges are unbounded.
If $P\ne0$, the resulting smooth compact threefold has
\[
 \pi_1(X)=\mathbb Z/|P|,\qquad
 H_j(X;\mathbb Z)=
 \begin{cases}
 \mathbb Z&j=0,6,\\
 \mathbb Z/|P|&j=1,2,3,4,\\
 0&j=5.
 \end{cases}
\]
If $P=0$, then $\pi_1(X)=\mathbb Z$ and all integral homology
is free, with Betti numbers $(1,1,2,2,2,1,1)$.
In particular, in this admissible family $X$ is simply connected
exactly when $|P|=1$, and then it is diffeomorphic to $S^6$.
Such sphere data exist exactly when $\gcd(3k,4\ell)=1$.
\end{theorem}


\begin{proof}
In the marking of Definition~\ref{mapsIII:def:two-finite-fibre-marking},
$A_0A_1A_\infty=I$, $A_0^3=A_1^4=I$, and
$\Lambda^{A_i}=\mathbb Z\delta\oplus\mathbb Zu_i$ for $i=0,1$.
On the degree-$m_i$ root disc the affine generator rotates the base
by $e^{2\pi i/m_i}$ and translates the torus by $v_i/m_i$.
Its $j$th power changes the invariant circle coordinate $d^\vee$
by $jn_i/m_i$. Coprimality proves freeness for
$0<j<m_i$, including on the central torus. The quotient is therefore
a smooth filling $N_i$ retracting onto its bielliptic surface $F_i$.
An invariant connection and radial transport give this retraction.
The rank-two Mumford filling supplies $N_\infty$.

Let $V$ be the inverse image of the complementary three-holed sphere
and put $M_i=\partial N_i$. The section-based loops satisfy
\[
 x_0x_1x_\infty=1,\qquad
 x_i\lambda x_i^{-1}=A_i\lambda.
\]
The logarithmic overlap subtracts
$(\log s)/(2\pi i)\,v_i$ in the root-disc coordinate.
Consequently the finite boundary maps send $x_i$ to a loop $g_i$
with $g_i^{m_i}=v_i$; their positive normal-circle kernels are
generated by $\mu_i=(-v_i,m_i)$. The finite torus cover of $F_i$
gives $\pi_2(F_i)=0$, so the normal-circle sequence identifies
this entire kernel. At infinity the toric discs kill
$K_\infty=\mathbb Z\delta\oplus\mathbb Za$, and the section disc
kills $x_\infty$. The period quotient identifies the remaining
map with
\[
 \pi_1(M_\infty)\longrightarrow\Lambda/K_\infty,
 \qquad(\lambda,j)\longmapsto[\lambda].
\]

Van Kampen now kills $\delta,a,x_\infty$. The conjugation
relation for $a$ kills $b$, and $x_1=x_0^{-1}$. The surviving
loops $d,x=x_0$ commute, since $(A_0-I)d=2\delta+6a$.
Conversely these two loops with the following two relations satisfy
every relation in the amalgamated presentation. Hence
\[
 \pi_1(X)=
 \mathbb Z\langle d,x\rangle/
 \langle-n_0d+m_0x,\,-n_1d-m_1x\rangle.
\]
Choose $\alpha m_0+\beta n_0=1$. For the column relation matrix,
\[
 \begin{pmatrix}-\beta&\alpha\\m_0&n_0\end{pmatrix}
 \begin{pmatrix}-n_0&-n_1\\m_0&-m_1\end{pmatrix}
 =
 \begin{pmatrix}1&\beta n_1-\alpha m_1\\0&-P\end{pmatrix}.
\]
The left matrix has determinant $-1$. Subtracting the indicated
multiple of the first column gives $\operatorname{diag}(1,-P)$.
The resulting isomorphism and its inverse are
\[
 d\longmapsto m_0,\quad x\longmapsto n_0
 \quad\text{in }\mathbb Z/P\mathbb Z,\qquad
 [1]\longmapsto\alpha d+\beta x.
\]
Here $\mathbb Z/0\mathbb Z=\mathbb Z$. These formulas prove both
fundamental-group assertions.

We next compute the degree-one kernel in Mayer--Vietoris.
The images of $A_0-I,A_1-I$ are respectively
\[
 \mathbb Z\langle\delta+3a,a+b\rangle,\qquad
 \mathbb Z\langle\delta+2a,a+b\rangle.
\]
Preimages of the displayed generators are $(-a-b,-b)$ and
$(b-a,-a)$, respectively. In the coordinate order
$(\delta,a,b,d)$ the quotient maps are
\[
 q_0=\begin{pmatrix}-3&1&-1&0\\0&0&0&1\end{pmatrix},
 \qquad
 q_1=\begin{pmatrix}-2&1&-1&0\\0&0&0&1\end{pmatrix};
\]
their sections use $a,d$. The common quotient is $\mathbb Zd$.
Thus $H_1(V)=\mathbb Z\langle d,x_0,x_1\rangle$, while each
$H_1(M_i)$ has rank three and is free.
Writing $\kappa_0=3,\kappa_1=2$, the finite surface groups give
\[
 H_1(F_i)=
 \mathbb Z\langle a,d,g_i\rangle/
 \mathbb Z(\kappa_i r_i a-n_i d+m_i g_i).
\]
Choose $m_i\alpha_i+n_i\beta_i=1$ and set
$z_i=\alpha_i d+\beta_i g_i$. The mutually inverse formulas
\[
 d=m_i z_i+\kappa_i r_i\beta_i a,\qquad
 g_i=n_i z_i-\kappa_i r_i\alpha_i a
\]
identify $H_1(F_i)$ with $\mathbb Z\langle a,z_i\rangle$.
The map from $H_1(M_i)$ has matrix
\[
 \begin{pmatrix}1&\kappa_i r_i\beta_i&-\kappa_i r_i\alpha_i\\
                 0&m_i&n_i\end{pmatrix}
\]
in the bases $(a,d,x_i)$ and $(a,z_i)$, with section
$a\mapsto a$, $z_i\mapsto\alpha_i d+\beta_i x_i$.
Its kernel is $\mathbb Z\mu_i$. At infinity,
$H_1(F_\infty)=\Lambda/K_\infty$ and the kernel is
$\mathbb Zx_\infty$.

Projection to $H_1(V)$ identifies the kernel of the full
Mayer--Vietoris map with the kernel of
\[
 G=\begin{pmatrix}-n_0&-n_1&0\\m_0&0&-1\\0&m_1&-1\end{pmatrix}
 :\mathbb Z\langle\mu_0,\mu_1,x_\infty\rangle
 \longrightarrow\mathbb Z\langle d,x_0,x_1\rangle.
\]
The local kernel descriptions give the inverse identification.
Since $\det G=-P$, this kernel is zero for $P\ne0$.
For $P=0$, equality of the reduced fractions
$n_0/m_0=-n_1/m_1$ gives
\[
 m_0=m_1=m=12q,\quad k=4q,\quad\ell=3q,\quad
 n_0=-n_1=n,\qquad \ker G=\mathbb Z(1,1,m).
\]

Each filling retracts to its central surface and satisfies
$H_2(N_i,M_i)=H^4(F_i)=\mathbb Z$ by
Poincar\'e--Lefschetz duality. At the cusp the top group is that
of its single normalized component. The map from $H_2(N_i)$
to this group is intersection with $F_i$. Since $m_iF_i$ is
principal, this homomorphism to a torsion-free group vanishes.
Therefore $H_2(M_i)\to H_2(N_i)$ is onto. Define
\[
 \tau_i:H^3(F_i)=H_3(N_i,M_i)
 \xrightarrow{\partial}H_2(M_i)\longrightarrow H_2(V).
\]
Moving a filling component of a Mayer--Vietoris representative
to its boundary identifies its degree-two cokernel with
$H_2(V)/\sum_i\operatorname{im}\tau_i$; the inverse includes
the $V$ component. Exactness gives
\[
 0\longrightarrow H_2(V)/\sum_i\operatorname{im}\tau_i
 \longrightarrow H_2(X)\longrightarrow\ker G\longrightarrow0.
\]

The cellular filtration over the two-loop graph gives
\[
 0\longrightarrow(\bigwedge\nolimits^2\Lambda)_{A_0,A_1}
 \longrightarrow H_2(V)\xrightarrow{p}
 \ker\bigl[(A_0-I)\ (A_1-I)\bigr]\longrightarrow0.
\]
The first map includes a fibre cycle; $p$ cuts a cycle along
the two chosen base loops. In the exterior basis
$(s_1,\ldots,s_6)=(\delta a,\delta b,\delta d,ab,ad,bd)$,
put ${\mathsf E}_i=\bigwedge^2 A_i-I$. The combined relation image is
$\mathbb Z\langle s_1,s_2,s_5,s_6,s_3+6s_4\rangle$.
The reverse inclusion is witnessed by
\[
 \begin{aligned}
 s_1&={\mathsf E}_0s_4,& s_2&={\mathsf E}_1s_4,\\
 s_5&=2{\mathsf E}_0s_1-2{\mathsf E}_0s_2+{\mathsf E}_0s_4-{\mathsf E}_0s_5-{\mathsf E}_1s_6,\\
 s_6&=-2{\mathsf E}_0s_1+2{\mathsf E}_0s_2+{\mathsf E}_0s_4+{\mathsf E}_0s_5-{\mathsf E}_0s_6+{\mathsf E}_1s_6,\\
 s_3+6s_4&=-5{\mathsf E}_0s_1+9{\mathsf E}_0s_2+2{\mathsf E}_0s_5-{\mathsf E}_0s_6+3{\mathsf E}_1s_6.
 \end{aligned}
\]
The functional
$\xi=a^\vee\wedge b^\vee+6d^\vee\wedge\delta^\vee$
annihilates every relation, has exactly this kernel, and evaluates
to one on $s_4$. Hence the fibre term is
$\mathbb Zs$, with $s=[a\wedge b]$ and quotient coordinate $\xi$.

Cutting the peripheral trace along $x_\infty=x_1^{-1}x_0^{-1}$
gives the edge map
\[
 \ell_\infty(\lambda)=-(A_1\lambda,\lambda),\qquad
 \ell_\infty(\delta)=(-\delta,-\delta),\quad
 \ell_\infty(a)=(b,-a).
\]
For $(v_0,v_1)$ in the edge kernel, applying $q_0$ to its
defining equation gives $b^\vee(v_1)=3d^\vee(v_1)$.
Set $h=d^\vee(v_1)$ and $\lambda=v_1-hu_1\in K_\infty$.
Adding $\ell_\infty(\lambda)$ leaves second component $hu_1$
and a first component fixed by $A_0$. Its unique coordinates
in $(\delta,u_0)$, together with $h$, define an integral
quotient isomorphism with inverse basis
\[
 (\delta,0),\qquad(u_0,0),\qquad(0,u_1).
\]
Comparison of second and then first components proves injectivity.

The cusp normalization is the hexagonal toric surface
$\operatorname{Bl}_3\mathbb P^2$. The differences of degrees
of $AH-\sum c_jE_j$ on its three opposite boundary pairs are
$(c_1+c_2+c_3-A)(1,1,1)$. The normalization sequence therefore gives
\[
 H^3(F_\infty)=\mathbb Z^3/\mathbb Z(1,1,1)
 \xrightarrow{\ \sim\ }\mathbb Z^2,\quad
 [y]\longmapsto(y_1-y_3,y_2-y_3),
\]
with inverse represented by $(w_1,w_2,0)$.
Sweeping the section disc by the two toric circles gives relative
chains whose boundaries are the $\delta$- and $a$-sweep tori.
Their edge coordinates are the two displayed values of
$\ell_\infty$. Their span $L$ is primitive: the preceding
quotient map leaves three free edge coordinates, and the fibre
term is $\mathbb Zs$. It follows that
\[
 H_2(V)/L=\mathbb Z\langle s,t,t_0,t_1\rangle,\qquad
 t=T_{\delta,0},\quad t_0=T_{u_0,0},\quad t_1=T_{u_1,1}.
\]
The rank-two group $H^3(F_\infty)$ maps onto a subgroup
containing $L$. Its quotient by $L$ has rank zero in this free
group, hence is zero. Thus $L=\operatorname{im}\tau_\infty$.
The flat vertical two-form $\Omega$ defined by $\xi$ evaluates
to one on $s$ and to zero on $t,t_0,t_1$; its integral supplies
the first coordinate in this basis.

Orient a finite normal torus by its loop first and its positive
normal circle second. For a covering-torus loop $\lambda$, the
root-disc trivialization describes it by
$(r,u)\mapsto(r\lambda-u v_i,m_i u)$, with the successive
segments identified by $A_i$. Cutting along a regular fibre
gives $\mathcal N_i\lambda$, where
$\mathcal N_i=\sum_{j=0}^{m_i-1}A_i^j$, and integrating
$\Omega$ gives $-\xi(\lambda,v_i)$.
For $g_i$ these coordinates are $v_i$ and zero: apply the
homomorphisms to $m_i[g_i]=[v_i]$ and cancel $m_i$ in their
free targets. These are the normal-circle boundary maps of
\cite[Propositions~4.3 and 4.5, pp.~11--14]{manViaclovskyLogs26},
with the finite bases established above for all the stated data.
The norms and form evaluations are
\[
 \mathcal N_0a=-k\delta,\quad\mathcal N_0d=m_0u_0,\qquad
 \mathcal N_1a=-2\ell\delta,\quad\mathcal N_1d=m_1u_1,
\]
\[
 \xi(a,u_0)=2,\quad\xi(a,u_1)=3,\quad
 \xi(d,\delta)=6,\quad\xi(d,u_i)=0.
\]
Modulo the cusp relations, $(0,\delta)=-(\delta,0)$.
Consequently the normal-circle map, after Poincar\'e duality
on the finite surfaces, has matrix
\[
 B=\begin{pmatrix}
 -2n_0&-6r_0\alpha_0&-3n_1&-6r_1\alpha_1\\
 -k&r_0\beta_0&2\ell&-r_1\beta_1\\
 0&1&0&0\\0&0&0&1
 \end{pmatrix}
\]
in the bases $(a,z_0,a,z_1)$ and $(s,t,t_0,t_1)$.

Let $C$ permute the columns to $(a_0,a_1,z_0,z_1)$, and put
\[
 K=\begin{pmatrix}-2n_0&-3n_1\\-k&2\ell\end{pmatrix},\qquad
 R_0=\begin{pmatrix}-6r_0\alpha_0&-6r_1\alpha_1\\
                    r_0\beta_0&-r_1\beta_1\end{pmatrix}.
\]
The matrices
\[
 BC=\begin{pmatrix}K&R_0\\0&I_2\end{pmatrix},\qquad
 U_4=\begin{pmatrix}I_2&-R_0\\0&I_2\end{pmatrix}
\]
satisfy $U_4BC=\operatorname{diag}(K,I_2)$.
The inverse of $U_4$ replaces $-R_0$ by $R_0$.
In particular the cokernel isomorphism is
\[
 [x]\longmapsto
 \left[
 \begin{pmatrix}
 x_s+6r_0\alpha_0x_0+6r_1\alpha_1x_1\\
 x_t-r_0\beta_0x_0+r_1\beta_1x_1
 \end{pmatrix}\right]\in\operatorname{coker}K,
\]
and its inverse appends two zero coordinates.
Now $\det K=-P$. A common divisor of its entries divides
$\gcd(k,2n_0)$, hence divides $2$; the odd integer $3n_1$
removes that factor. The gcd of all entries is therefore one.

For completeness the cyclic coordinate can be obtained by
integer Euclidean operations. Move a nonzero entry to the
upper-left corner and make it positive. Divide the lower-left
entry by it using a row addition, swapping rows for a nonzero
remainder, and do the same to the upper-right entry by columns.
Each nonzero remainder decreases the positive pivot. Once the
matrix is diagonal, add the second row to the first whenever
the pivot fails to divide the other diagonal entry; the next
column division again decreases it. This terminates with a
pivot generating the ideal of all entries, which is $\mathbb Z$.
Recording the row and column additions, swaps and sign changes
gives unimodular ${\mathsf U},{\mathsf V}$ with
\[
 {\mathsf U}K{\mathsf V}=\operatorname{diag}(1,|P|).
\]
Each operation has its inverse of the same kind. For $P\ne0$
the isomorphism and inverse are
\[
 [y]\longmapsto ({\mathsf U}y)_2\bmod |P|,\qquad
 [1]\longmapsto[{\mathsf U}^{-1}(0,1)^t].
\]
Indeed ${\mathsf U}K=\operatorname{diag}(1,|P|){\mathsf V}^{-1}$, and a vector
${\mathsf U}y=(a,|P|b)^t$ comes from $K{\mathsf V}(a,b)^t$.
Thus $H_2(X)=\mathbb Z/|P|$ for $P\ne0$.

At $P=0$ the matrix has the explicit factorization
\[
 K=\begin{pmatrix}n\\2q\end{pmatrix}\begin{pmatrix}-2&3\end{pmatrix}.
\]
Its two columns sum to the primitive vector $(n,2q)^t$.
Choose $n\eta+2q\zeta=1$. The cokernel coordinate and inverse are
\[
 (y_1,y_2)\longmapsto-2q y_1+n y_2,\qquad
 1\longmapsto(-\zeta,\eta).
\]
The matrix with rows $(\eta,\zeta)$ and $(-2q,n)$ has determinant
one and sends the image generator to $(1,0)^t$.
To split the remaining Mayer--Vietoris extension, take normal
discs $D_0,D_1$ with boundaries $\mu_0,\mu_1$, and the section
disc $D_\infty$ with boundary $x_\infty$. Their boundary sum is
\[
 \mu_0+\mu_1+m x_\infty
 =m(x_0+x_1+x_\infty)-(r_0+r_1)\delta+n(a+b).
\]
This equality of singular one-cycles uses the linear triangles
that add lattice loops. A section of the base pair of pants
fills the first parenthesis. Tracing a fibre loop $\lambda$
along $x_i$ gives a chain with boundary $(A_i-I)\lambda$.
The identities
\[
 a+b=-(A_0-I)b,\qquad
 \delta=2(A_0-I)(a+b)+3(A_1-I)(b-a)
\]
therefore give an integral chain $C_0$ in $V$ with the displayed
boundary. Specifically use $m$ section chains, $-n$ traces of
$b$ along $x_0$, and $-(r_0+r_1)$ times the two traces in the
second identity. The cycle
\[
 Z=D_0+D_1+mD_\infty-C_0
\]
maps to $(1,1,m)\in\ker G$. Let $h_{\rm cok}$ be the cokernel class
represented by $(-\zeta,\eta,0,0)$ in the four target coordinates.
Then $(a,b)\mapsto a h_{\rm cok}+b[Z]$ identifies $\mathbb Z^2$ with
$H_2(X)$. Its inverse reads $b$ under the connecting map,
subtracts $b[Z]$, and applies the preceding cokernel coordinate.
This proves the asserted extension and both inverse identities.

The smooth torus and finite bielliptic fibres have Euler number
zero. The cusp has two zero-dimensional toric strata; all its
other strata are positive-dimensional algebraic tori.
Additivity gives $\chi(X)=2$.
For $P\ne0$, set $D=|P|$. Cap product with the complex
fundamental class and the universal coefficient sequence give
\[
 H_5=H^1=0,\qquad
 H_4=H^2=\operatorname{Ext}(\mathbb Z/D,\mathbb Z)
 =\mathbb Z/D.
\]
Here $\operatorname{Ext}$ is the cokernel of multiplication
by $D$ obtained from the two-term free resolution of
$\mathbb Z/D$. Euler characteristic gives $b_3=0$, and the
degree-three universal coefficient sequence gives
$H_3=H^3=\operatorname{Ext}(H_2,\mathbb Z)=\mathbb Z/D$.
Together with $H_0=H_6=\mathbb Z$ these are all the stated groups.
For $P=0$, the same maps give $H_1=H_5=\mathbb Z$ and
$H_2=H_4=\mathbb Z^2$. The degree-three Ext term vanishes,
so cap product and evaluation identify $H_3$ with its integer
dual; it is free. Euler characteristic gives its rank as two.

The fundamental-group computation identifies simple connectivity
with $|P|=1$. In that case Hurewicz and Whitehead identify $X$
as a homotopy six-sphere: a representative of the generator
of $H_6$ is a homology equivalence $S^6\to X$.
The dimension-six entry and its diffeomorphism interpretation in
\cite[pp.~504--505]{manKM63} give $X\cong S^6$ smoothly.
Finally, every common divisor of $m_0,m_1$ divides $P$;
B\'ezout gives $m_1n_0+m_0n_1=1$ for coprime multiplicities,
and reduction modulo each $m_i$ proves local admissibility.
This proves the exact sphere-existence criterion.
For unbounded sphere multiplicities take $k=2j+1$ and
$\ell=3k+1$, $j\ge0$. For each such pair, replacing
$(n_0,n_1)$ by $(n_0+m_0t,n_1-m_1t)$ preserves $P=1$
for every $t\in\mathbb Z$; $r_0,r_1$ are arbitrary.
For example, $(m_0,m_1;n_0,n_1)=(9,4;-2,1)$ gives $S^6$,
$(9,4;1,1)$ gives the four groups $\mathbb Z/13$, and
$(12,12;1,-1)$ gives the infinite cyclic case.
\end{proof}

\begin{remark}[The initial logarithmic homology calculation]
\hypertarget{tag-Q109}{}\label{mapsIII:rem:initial-logarithmic-homology}
{\normalfont\scriptsize tag Q109}\par
The local logarithmic quotient is the construction of
\cite[Proposition~2.3, pp.~5--6]{manViaclovskyLogs26}. The second homology of the initial family is expressed by the
integral cap-map cokernel in
\cite[Proposition~4.1, p.~9]{manViaclovskyLogs26};
\cref{maps:src:D.2.5} gives the complete positive-integer
multiplicities and both determinant cases in the same marking.
\end{remark}

\subsection{The \texorpdfstring{$III^*+3I_1$}{III*+3I1} profile}

\begin{theorem}[Logarithmic multiplicities over $III^*+3I_1$]\hypertarget{tag-02LW}{}\label{maps:src:D.2.6}
{\normalfont\scriptsize tag 02LW}\par


Use the $III^*,I_1,I_1,I_1$ construction of \SeriesRef{push:src:E.5.2} and \SeriesRef{push:src:E.5.3}, with
two primitive rank-one Mumford fillings and one rank-two hexagonal
filling. In the ordinary homology basis $(e_0,e_1,e_2,e_3)$ dual
to $(\psi,u,w,\delta)$, put $\ell=e_0-e_1-e_2$.
At $III^*$ one may make a logarithmic transformation with
\[
 m=4k,\qquad v=n\ell+r e_3,\qquad
 k\in\mathbb Z_{>0},\quad n,r\in\mathbb Z,\quad \gcd(n,4k)=1.
\]
For every such datum the resulting compact complex threefold has
\[
 \pi_1(X)=\mathbb Z/|n|,\qquad
 H_0=H_6=\mathbb Z,\quad H_5=0,\quad
 H_1=H_2=H_3=H_4=\mathbb Z/|n|.
\]
Consequently $n=1$ or $n=-1$ gives a complex structure on $S^6$
for every $k>0$. The multiplicities $4k$ give an affirmative answer
to \cite[Question~4, p.~2]{manViaclovskyLogs26} for this configuration.
\end{theorem}


\begin{proof}
Use the marking and the analytic family of
\SeriesRef{push:src:E.5.2} and \SeriesRef{push:src:E.5.3}.
Their local models are described in
\cite[Proposition~3.1 and Lemma~3.2, pp.~6--8;
\S3.3, pp.~8--9]{manViaclovskyTwo26}.
Write $A_i=T_i^{-t}$ on
$\Lambda=\mathbb Z\langle e_0,e_1,e_2,e_3\rangle$.
In the ordered punctures $*,1,2,3$ these matrices are
\[
 A_*=\begin{pmatrix}1&0&0&0\\-2&0&-1&0\\0&1&0&0\\1&1&0&1\end{pmatrix},
 \qquad
 A_1=\begin{pmatrix}1&0&0&0\\0&1&0&0\\-2&-1&1&0\\-2&-1&0&1\end{pmatrix},
\]
\[
 A_2=\begin{pmatrix}1&0&0&0\\0&1&1&0\\0&0&1&0\\0&0&0&1\end{pmatrix},
 \qquad
 A_3=\begin{pmatrix}1&0&0&0\\0&1&0&0\\0&-1&1&0\\1&0&0&1\end{pmatrix}.
\]
They satisfy $A_*A_1A_2A_3=I$, $A_*^4=I$, and
$\Lambda^{A_*}=\mathbb Z\ell\oplus\mathbb Ze_3$.

Pull the good-reduction family back by the degree-$k$ map
of the covering disc. Write its root coordinate as $u$,
so the original base coordinate is $u^m$, $m=4k$.
The deck generator $\rho$ over
$u\mapsto e^{-2\pi i/m}u$ induces $A_*$.
Set $\widetilde\rho=t_{v/m}\rho$.
Invariance of $v$ makes its $m$th power the identity.
For $0<j<m$, its $j$th power changes the invariant circle
coordinate $\psi=e_0^\vee$ by $jn/m$; this is nonintegral.
The action is therefore free on the central torus, and the
base rotation gives freeness elsewhere. Its quotient is a
smooth filling with reduced bielliptic central surface $F_*$.

The overlap and its inverse are the translations
$h_u=t_{\sigma(u)}$ and $h_u^{-1}=t_{-\sigma(u)}$, where
\[
 \sigma(u)=\frac{\log u}{2\pi i}v,\qquad
 \sigma(e^{-2\pi i/m}u)=\rho(\sigma(u))-v/m.
\]
A change of logarithm adds a period. The displayed identity gives
\[
 h_{e^{-2\pi i/m}u}\,\widetilde\rho\,h_u^{-1}=\rho,
\]
so the quotient glues to the specified punctured family.
Together with the three Mumford fillings this gives a compact
complex threefold for every stated datum.

We fix geometric orientations in the boundary maps.
Let $g$ be the affine lift of $\widetilde\rho$ on the central
torus, so $g^m=v$ and its conjugation on $\Lambda$ is $A_*$.
The external section is represented on the twisted cover by
$-\sigma(u)$. A positive base loop ends at
$\widetilde\rho^{-1}(-\sigma(u))$; hence its section-based
meridian $x_*$ maps to $g^{-1}$.
In particular $x_*^m$ maps to $-v$.
The positive normal-circle kernel is generated by
$\mu=(v,m)$ in fibre and positive-base coordinates.
The finite torus cover of $F_*$ gives $\pi_2(F_*)=0$;
the normal-circle sequence therefore proves the kernel
description and boundary surjectivity.

The primitive vanishing lattices at the reduced fillings are
\[
 K_1=\mathbb Z(e_2+e_3),\qquad
 K_2=\mathbb Ze_1,\qquad
 K_3=\mathbb Ze_2\oplus\mathbb Ze_3.
\]
The boundary map at each of these fillings is
$(\lambda,j)\mapsto[\lambda]\in\Lambda/K_i$.
Its toric discs kill $K_i$, and the extended section disc kills
the positive base meridian $x_i$. The period quotient supplies
the surviving loops and proves surjectivity.
Van Kampen thus kills $e_1,e_2,e_3$ and all four base meridians.
The remaining relation is $ne_0=0$.
The homomorphism sending $e_0$ to $1\bmod n$ and all other
generators to zero has inverse $1\mapsto[e_0]$.
This proves $\pi_1(X)=\mathbb Z/|n|$.

For the finite coinvariants, $\operatorname{im}(A_*-I)$ has
basis $e_1+e_2,e_3-2e_1$, with preimages $-e_2,e_0$.
The quotient map is
\[
 q_*=\begin{pmatrix}1&0&0&0\\0&1&-1&2\end{pmatrix},
\]
with section $(a,b)\mapsto ae_0+be_1$.
Consequently
\[
 H_1(F_*)=
 \mathbb Z\langle e_0,e_1,g\rangle/
 \mathbb Z(mg-ne_0-2re_1).
\]
Choose $m\alpha+n\beta=1$ and put $z=\alpha e_0+\beta g$.
The inverse formulas
\[
 e_0=mz-2r\beta e_1,\qquad
 g=nz+2r\alpha e_1
\]
give the integral basis $(e_1,z)$.
In the boundary basis $(e_0,e_1,x_*)$, the map and a section are
\[
 \begin{pmatrix}-2r\beta&1&-2r\alpha\\m&0&-n\end{pmatrix},
 \qquad e_1\longmapsto e_1,\quad
 z\longmapsto\alpha e_0-\beta x_*.
\]
Its kernel is $\mathbb Z(n,2r,m)$.
At each reduced filling,
$H_1(M_i)=(\Lambda/K_i)\oplus\mathbb Zx_i$
maps by projection to $H_1(F_i)=\Lambda/K_i$.

The common coinvariants are $\mathbb Ze_0$:
$(A_2-I)e_2=e_1$, $(A_*-I)e_0=e_3-2e_1$, and
$-(A_1-I)e_1=e_2+e_3$ generate the other three coordinates.
Thus $H_1(V)=\mathbb Z\langle e_0,x_*,x_1,x_2\rangle$.
The kernel of the full degree-one Mayer--Vietoris map is
identified, using the local kernels just specified, with
the kernel of
\[
 G_+=\begin{pmatrix}
 n&0&0&0\\m&0&0&-1\\0&1&0&-1\\0&0&1&-1
 \end{pmatrix}
 :\mathbb Z\langle\mu,x_1,x_2,x_3\rangle
 \longrightarrow H_1(V).
\]
Its determinant is $-n$. Admissibility gives $n\ne0$, so
this kernel vanishes.

Every filling retracts to its irreducible central surface and has
$H_2(N_i,M_i)=H^4(F_i)=\mathbb Z$.
The map from $H_2(N_i)$ is intersection with that surface.
A multiple of the surface is principal, so this map vanishes.
Hence $H_2(M_i)\to H_2(N_i)$ is onto.
Write
\[
 \tau_i:H^3(F_i)=H_3(N_i,M_i)
 \xrightarrow{\partial}H_2(M_i)\longrightarrow H_2(V).
\]
In Mayer--Vietoris, move each filling representative to its
boundary and then to the $V$ component. The inverse includes
that component, and the resulting relations are precisely
the images of the $\tau_i$. The injectivity of $G_+$ gives
\[
 H_2(X)=H_2(V)/\sum_i\operatorname{im}\tau_i.
\]

To compute this quotient, use the graph with edges $*,1,2$.
The fibre inclusion and cutting map give
\[
 0\longrightarrow(\bigwedge\nolimits^2\Lambda)_{A_*,A_1,A_2}
 \longrightarrow H_2(V)\xrightarrow{p}\ker D\longrightarrow0,
 \qquad D=[A_*-I\quad A_1-I\quad A_2-I].
\]
Put $(s_1,\ldots,s_6)=(e_0e_1,e_0e_2,e_0e_3,e_1e_2,e_1e_3,e_2e_3)$
and $E_i=\bigwedge^2A_i-I$.
The relation image has basis $s_1,s_2,s_5,s_6,s_3-2s_4$.
The following preimages prove its integral inclusion:
\[
 \begin{aligned}
 s_1&=E_2s_2,&s_5&=E_2s_6,&s_6&=E_1s_4,\\
 s_2&=-E_*s_2-E_2s_2+E_2s_6,\\
 s_3-2s_4&=E_*s_1+E_*s_2+2E_2s_2+E_2s_6+E_1s_4.
 \end{aligned}
\]
Conversely
$\xi=e_1^\vee\wedge e_2^\vee+2e_0^\vee\wedge e_3^\vee$
annihilates every relation and has exactly this kernel.
Its value on $s_4$ is one, giving the fibre coordinate
with inverse generator $s=[e_1\wedge e_2]$.

The reduced invariant bases are
\[
 \mathsf B_1=(e_0-2e_1,e_2,e_3),\qquad
 \mathsf B_2=(e_0,e_1,e_3),\qquad \mathsf B_3=(e_2,e_3).
\]
For each invariant vector $\lambda$, the circle translation
$t\mapsto t\Pi(q)\lambda$ extends over its Mumford filling.
Indeed its logarithmic disc terms cancel by invariance.
The remaining unit factors act on the flat ruled normalization
and preserve its two sections and their translation gluing;
at the hexagonal fibre the two toric rotations extend in
each chart. The time-one translation is a period and hence
the identity. Sweeping the extended section disc therefore
gives a relative three-chain with boundary $T_{\lambda,i}$.
The integral specialization in these primitive models gives
$H^3(F_i)$ free abelian of ranks $3,3,2$, respectively
\cite[Proposition~3.1, pp.~6--7]{manViaclovskyTwo26}.

The eight sweep edge columns are
\[
 L=\begin{bmatrix}
 0&0&-A_1A_2\mathsf B_3\\
 \mathsf B_1&0&-A_2\mathsf B_3\\
 0&\mathsf B_2&-\mathsf B_3
 \end{bmatrix}.
\]
The last block follows by cutting the third peripheral trace:
$j_3(\lambda)=-(A_1A_2\lambda,A_2\lambda,\lambda)$.
For an edge-kernel vector with components
$a=(a_0,a_1,a_2,a_3)$, $b=(b_0,b_1,b_2,b_3)$,
$c=(c_0,c_1,c_2,c_3)$, its defining equations are
\[
 a_2=-a_0,\qquad b_1=a_0+a_1-2b_0,\qquad c_2=a_0+a_1.
\]
Set $h=a_0+a_1$. The exact decomposition is
\[
 (a,b,c)=a_0(\ell,0,0)+L\eta,
\]
\[
 \eta=(b_0,b_2-h,b_3-a_3-h,c_0,c_1,c_3-a_3-h,-h,-a_3-h)^t.
\]
These formulas and the displayed columns are inverse integral maps.
Thus $\ker D/\operatorname{im}L=\mathbb Z$, with quotient
coordinate $a_0$ and inverse $(\ell,0,0)$.

Lift the latter vector by $t=T_{\ell,*}$ and the other eight
vectors by the relative sweep tori. Together with $s$ they
form an integral basis of $H_2(V)$: cut, use the displayed
inverse coordinates, and subtract the sweeps to leave the
unique fibre coordinate.
The closed flat vertical form $\Omega$ associated with $\xi$
evaluates to one on $s$ and zero on every sweep.
Consequently, after these eight relations, the quotient and
its inverse are
\[
 Q\xrightarrow{\ \sim\ }\mathbb Z^2,\qquad
 [Z]\longmapsto\left(\int_Z\Omega,\ e_0^\vee(p(Z)_*)\right),
 \qquad(a,b)\longmapsto as+bt.
\]
The eight relative cycles belong to the three reduced
boundary images. Their span is primitive; those source groups
have total rank eight. The further image modulo their span
has rank zero in the free abelian group $Q$, hence vanishes.
This proves the exact reduced-filling quotient.

We now compute the finite boundary map using positive geometric
normal circles. In the twisted root-disc trivialization a
covering loop $\lambda$ times the normal circle maps through
$h_u$ to
\[
 (a,b)\longmapsto(a\lambda+bv,mb)
\]
in marked fibre and positive-base coordinates.
Its fibre cut is $\mathcal N_m\lambda$, and its integral
against $\Omega$ is $\xi(\lambda,v)$.
Here
\[
 \mathcal N_me_1=2ke_3,\qquad
 \mathcal N_me_0=m\ell,\qquad
 \xi(e_1,v)=-n,\qquad \xi(e_0,v)=2r.
\]
The $e_3$ sweep has zero coordinate in $Q$.
Thus the finite normal tori of $e_1,e_0$ have coordinates
$(-n,0),(2r,m)$.
Applying the homomorphism to $mg=ne_0+2re_1$ gives
$m\tau_*(g)=(0,mn)$; freeness of $Q$ gives
$\tau_*(g)=(0,n)$.
In the finite basis $(e_1,z)$ the matrix is therefore
\[
 B_+=\begin{pmatrix}-n&2r\alpha\\0&1\end{pmatrix},
 \qquad
 \begin{pmatrix}1&-2r\alpha\\0&1\end{pmatrix}B_+
 =\begin{pmatrix}-n&0\\0&1\end{pmatrix}.
\]
In the positive geometric coordinates the cyclic map and inverse are
\[
 [(a,b)]\longmapsto a-2r\alpha b\bmod n,\qquad
 [1]\longmapsto[s].
\]
Subtracting $b$ times the second column leaves
$(a-2r\alpha b,0)$, which belongs to the first column image
exactly when its first coordinate is divisible by $n$.
Hence $H_2(X)=\mathbb Z/|n|$ with the stated inverse generator.

The two rank-one fibres have Euler number zero, as do the
smooth tori and the bielliptic fibre. The hexagonal fibre has
Euler number $6-3\cdot2+2=2$. Thus $\chi(X)=2$.
Cap product with the complex fundamental class and the
universal coefficient sequences give
\[
 H_5=0,\qquad
 H_4=\operatorname{Ext}(H_1,\mathbb Z)=\mathbb Z/|n|.
\]
The Euler characteristic forces $b_3=0$, and the next
universal coefficient sequence gives
$H_3=\operatorname{Ext}(H_2,\mathbb Z)=\mathbb Z/|n|$.
The Ext coordinates are the cokernels of multiplication
by $|n|$ in the corresponding free resolutions.
Connectedness and orientation give $H_0=H_6=\mathbb Z$.
For $n=\pm1$, Hurewicz and Whitehead identify $X$ as a
homotopy six-sphere, and
\cite[table on p.~504 and first remark on p.~505]{manKM63} gives its diffeomorphism
with the standard $S^6$. Every positive $k$ is admissible
for these two values of $n$, giving the unbounded sequence
of multiplicities $4k$.
\end{proof}

\begin{remark}[Signed coordinates for the $III^*+3I_1$ boundary map]
\hypertarget{tag-Q111}{}\label{mapsIII:rem:iii-star-boundary-orientation}
{\normalfont\scriptsize tag Q111}\par
In the finite boundary calculation of \cref{maps:src:D.2.6}, choose
$m\alpha+n\beta=1$. The positive geometric matrix is
$B_+=\left(\begin{smallmatrix}-n&2r\alpha\\0&1\end{smallmatrix}\right)$.
Negating the relative source generators and replacing $t$
by $-t$ gives the exact alternative coordinate matrix
\[
 \begin{pmatrix}1&0\\0&-1\end{pmatrix}
 B_+(-I_2)
 =\begin{pmatrix}n&-2r\alpha\\0&1\end{pmatrix}.
\]
Both source and target changes are integral involutions.
The source map negates both relative generators; the target map fixes
the fibre class $s$ and negates its peripheral sweep $t$. Each map
is its own inverse.
\end{remark}

\subsection{Wang's primitive \texorpdfstring{$I_2$}{I2} filling}

\begin{theorem}[Logarithmic multiplicities with a primitive $I_2$ filling]\hypertarget{tag-02LX}{}\label{maps:src:D.2.7}
{\normalfont\scriptsize tag 02LX}\par


Use the rational elliptic surface
\[
 y^2=x^3-3(1+t)x+(2+3t+2t^2),
 \qquad P=(1,\sqrt2\,t),\qquad Q=6P.
\]
Its fibres are $IV^*,I_2,I_1,I_1$, and the degree-one
linearization vanishes at the second $I_1$ point. Use Wang's
primitive one-component rank-one Mumford model at $I_2$, the
rank-one model at the other $I_1$, and the hexagonal rank-two
model at the zero of the linearization. In the homology marking
dual to $(\psi,u,w,\delta)$ put
\[
 \begin{gathered}
 \ell=e_0-2e_1-4e_2,\qquad m=3k,\qquad v=n\ell+r e_3,\\
 k\in\mathbb Z_{>0},\qquad n,r\in\mathbb Z,\qquad \gcd(n,3k)=1.
 \end{gathered}
\]
The order-$m$ logarithmic quotient at $IV^*$ is smooth for every
integer $r$, and the resulting threefold satisfies
\[
 \begin{gathered}
 \pi_1(X)=\mathbb Z/|n|,\qquad H_0=H_6=\mathbb Z,
 \qquad H_5=0,\\
 H_1=H_4=\mathbb Z/|n|,\qquad
 H_2=H_3=\mathbb Z/(2|n|).
 \end{gathered}
\]
In particular every $k>0$, with $n=\pm1$, gives a simply
connected complex rational homology six-sphere with two-torsion.
At $(k,n,r)=(1,1,0)$ this is the construction of
\cite[Propositions~2.7--2.10, pp.~4--7]{manWang26}.
\end{theorem}


\begin{proof}
We use Wang's local models and period marking
\cite[Proposition~2.1, p.~2; Lemma~2.6,
Propositions~2.7--2.10, pp.~4--7;
Lemma~4.3 and Proposition~4.4, pp.~10--11]{manWang26}.
In the additive coordinates of that marking the periods are
\[
\begin{gathered}{}
 \alpha^{\rm per}_0=(0,1),\quad \alpha^{\rm per}_1=(1,0),\quad
 \alpha^{\rm per}_2=(\tau,-p),\quad \alpha^{\rm per}_3=(-6p,\kappa),
\\{}
 \qquad
 (e_0,e_1,e_2,e_3)=(\alpha^{\rm per}_3,-\alpha^{\rm per}_2,\alpha^{\rm per}_1,-\alpha^{\rm per}_0).
\end{gathered}
\]
The multiplicative line coordinate is
$v_{\rm line}=\exp(-2\pi i y)$, so its positive angular
cycle is $e_3$.
At $I_2$ the period change $g_{\rm per}=Fh^3$, with inverse
$F=g_{\rm per}h^{-3}$, has matrix
$\left(\begin{smallmatrix}1&3\\0&1\end{smallmatrix}\right)$.
The monomial map and its inverse are
\[
 x=v_{\rm line},\quad y=v_{\rm line}^2/z,
 \qquad z=x^2/y,\quad v_{\rm line}=x.
\]
Their exponent matrices are inverse integral matrices of
determinant one, and they send $(2,1)$ to $(1,0)$.
This gives the prescribed one-component rank-one filling.
The source construction extends the section across this
filling and the other two Mumford fillings.

Write $A_i=T_i^{-t}$ on
$\Lambda=\mathbb Z\langle e_0,e_1,e_2,e_3\rangle$.
For the ordered punctures $*,2,+,-$, the homology matrices are
\[
 A_*=\begin{pmatrix}1&0&0&0\\-6&0&-1&0\\-6&1&-1&0\\2&1&0&1\end{pmatrix},
 \qquad
 A_2=\begin{pmatrix}1&0&0&0\\0&1&0&0\\-6&-2&1&0\\-3&-1&0&1\end{pmatrix},
\]
\[
 A_+=\begin{pmatrix}1&0&0&0\\0&0&1&0\\0&-1&2&0\\0&0&0&1\end{pmatrix},
 \qquad
 A_-=\begin{pmatrix}1&0&0&0\\0&1&0&0\\0&-1&1&0\\1&0&0&1\end{pmatrix}.
\]
They satisfy $A_*A_2A_+A_-=I$, $A_*^3=I$, and
$\Lambda^{A_*}=\mathbb Z\ell\oplus\mathbb Ze_3$.

Pull the cubic good-reduction family back by the degree-$k$
map of its covering disc. In the resulting root coordinate $u$,
the original base coordinate is $u^m$, where $m=3k$.
The section-preserving deck map $\rho$ over
$u\mapsto e^{-2\pi i/m}u$ induces $A_*$.
The period $v=n\ell+re_3$ is invariant, so
$\widetilde\rho=t_{v/m}\rho$ has $m$th power the identity.
Its $j$th power changes the invariant circle coordinate
$\psi=e_0^\vee$ by $jn/m$.
For $0<j<m$ this is nonintegral, proving freeness on the
central torus. The base rotation proves freeness elsewhere.
The quotient is a smooth proper filling with central fibre
$mB$, where $B$ is bielliptic.

The punctured overlap is $h_u=t_{\sigma(u)}$, with inverse
$t_{-\sigma(u)}$, where
\[
 \sigma(u)=\frac{\log u}{2\pi i}v,\qquad
 \sigma(e^{-2\pi i/m}u)=\rho(\sigma(u))-v/m.
\]
Changing the logarithm adds a period, and
\[
 h_{e^{-2\pi i/m}u}\widetilde\rho h_u^{-1}=\rho.
\]
Thus this quotient glues holomorphically to the specified
punctured family. Together with the three Mumford fillings
it gives the compact complex threefold for every stated datum.

Let $g$ be the affine lift on the universal cover of the
central torus. It satisfies
\[
 gt_\eta g^{-1}=t_{A_*\eta},\qquad g^m=t_v.
\]
The positive section meridian $x_*$ is represented on the
twisted cover by continuation of $-\sigma(u)$.
Its endpoint is $\widetilde\rho^{-1}(-\sigma(u))$, so
\[
 x_*=g^{-1},\qquad x_*^m=t_{-v}.
\]
The boundary map is $(\eta,j)\mapsto t_\eta g^{-j}$.
Its kernel is generated by the positive normal circle
$\mu=(v,m)$: a kernel element has $j=qm$ and $\eta=qv$.
The normal forms $t_\eta g^a$, $0\leq a<m$, are unique
by the freeness test using $\psi$. This proves the full
presentation and boundary surjectivity.

The primitive vanishing lattices at the reduced fillings are
\[
 K_2=\mathbb Z(2e_2+e_3),\qquad
 K_+=\mathbb Z(e_1+e_2),\qquad
 K_-=\mathbb Ze_2\oplus\mathbb Ze_3.
\]
Their quotient maps and sections are
\[
 \begin{array}{c|c}
 q_2(x)=(x_0,x_1,x_2-2x_3)&(e_0,e_1,e_2)\\
 q_+(x)=(x_0,-x_1+x_2,x_3)&(e_0,e_2,e_3)\\
 q_-(x)=(x_0,x_1)&(e_0,e_1).
 \end{array}
\]
Subtracting each section from $x$ leaves, respectively,
$x_3(2e_2+e_3)$, $x_1(e_1+e_2)$ and $x_2e_2+x_3e_3$.
The local boundary map is $(\eta,j)\mapsto[\eta]\in\Lambda/K_i$.
The toric discs kill $K_i$, and the section disc kills the
positive base meridian $x_i$; the period quotient supplies
the remaining generators.

Let $V$ be the inverse image of the four-punctured base,
and let $X_D$ be obtained by adjoining the three reduced
fillings. Van Kampen kills $e_1,e_2,e_3$ and the three
reduced meridians; the sphere relation kills $x_*$.
Consequently $\pi_1(X_D)=\mathbb Z\langle e_0\rangle$.
Adjoining the finite filling imposes $ne_0=0$.
The maps $e_0\mapsto1\bmod n$ and $1\mapsto[e_0]$
are inverse homomorphisms on the full presentation, proving
$\pi_1(X)=\mathbb Z/|n|$.

The common fibre coinvariants in $V$ are $\mathbb Ze_0$.
Indeed, set
$p_{\rm rel}=-(A_+-I)e_1=e_1+e_2$,
$q_{\rm rel}=-(A_*-I)e_2=e_1+2e_2$ and
$w_{\rm rel}=-(A_2-I)e_1=2e_2+e_3$.
Then $e_2=q_{\rm rel}-p_{\rm rel}$, $e_1=2p_{\rm rel}-q_{\rm rel}$, and $e_3=w_{\rm rel}-2q_{\rm rel}+2p_{\rm rel}$.
The section therefore gives
\[
 H_1(V)=\mathbb Z\langle e_0,x_*,x_2,x_+\rangle,
 \qquad x_-=-x_*-x_2-x_+.
\]
At the reduced fillings the degree-one boundary maps are
the projections
$(\Lambda/K_i)\oplus\mathbb Zx_i\to\Lambda/K_i$.
Their combined kernel maps to $H_1(V)$ by
\[
 G_{\rm red}=\begin{pmatrix}
 0&0&0\\0&0&-1\\1&0&-1\\0&1&-1
 \end{pmatrix}.
\]
It is injective and its cokernel is the $e_0$ coordinate.
Explicitly, the base vector $(a,b,c)$ has inverse relation
coordinates $(b-a,c-a,-a)$.

For every filling $N_i$ with boundary $M_i$, duality gives
$H_2(N_i,M_i)=H^4(F_i)=\mathbb Z$.
The map from $H_2(N_i)$ is intersection with the irreducible
reduced central divisor. A multiple of that divisor is
principal, so its intersection homomorphism is zero.
Hence $H_2(M_i)\to H_2(N_i)$ is onto.
Let
\[
 \tau_i:H_3(N_i,M_i)=H^3(F_i)
 \xrightarrow{\partial}H_2(M_i)\longrightarrow H_2(V).
\]
In Mayer--Vietoris, choose a boundary preimage of each
filling class and move it to the $V$ component.
Changing this choice adds an image of $\tau_i$, and inclusion
of the $V$ component is the inverse map.
The injectivity of $G_{\rm red}$ proves
\[
 H_2(X_D)=H_2(V)/
 \bigl(\operatorname{im}\tau_2+\operatorname{im}\tau_+
                    +\operatorname{im}\tau_-\bigr).
\]

Cut the torus bundle along the three graph edges $*,2,+$.
Fibre inclusion and cutting give
\[
 0\longrightarrow(\bigwedge\nolimits^2\Lambda)_{A_*,A_2,A_+}
 \longrightarrow H_2(V)\xrightarrow{p_{\rm cut}}\ker D\longrightarrow0,
 \qquad D=[A_*-I\quad A_2-I\quad A_+-I].
\]
In the exterior basis
$(s_1,\ldots,s_6)=(e_0e_1,e_0e_2,e_0e_3,e_1e_2,e_1e_3,e_2e_3)$,
put $E_i=\bigwedge^2A_i-I$.
The relation image has basis
$s_1,s_2,s_5,s_6,s_3-6s_4$, with integer preimages
\[
 \begin{aligned}
 s_6&=E_2s_4,&s_5&=E_*s_4-E_2s_4,\\
 s_2&=-E_*s_2-E_+s_2+2E_*s_4,\\
 s_1&=E_+s_2-s_2,&s_3-6s_4&=-E_2s_1-2s_2+3s_5.
 \end{aligned}
\]
The last two formulas use the preceding relation preimages.
Conversely
\[
 \xi=e_1^\vee\wedge e_2^\vee+6e_0^\vee\wedge e_3^\vee
\]
annihilates every relation, and its kernel is spanned by
the five displayed relation vectors.
Its value on $s_4$ is one, so it gives the fibre quotient
with inverse generator $s=[e_1\wedge e_2]$.

The reduced fixed lattices have bases
\[
 \mathsf B_2=(e_0-3e_1,e_2,e_3),\qquad
 \mathsf B_+=(e_0,e_1+e_2,e_3),\qquad \mathsf B_-=(e_2,e_3).
\]
Every invariant period acts by a circle translation on its
Mumford filling. At $I_2$, the first vector is
$\alpha^{\rm per}_3+3\alpha^{\rm per}_2$, the zero-valuation period $Fh^3$;
the other two are angular periods.
At the other rank-one point, the elliptic basis
$(\alpha^{\rm per}_1-\alpha^{\rm per}_2,\alpha^{\rm per}_1)$ has determinant one and changes the
elliptic monodromy to
$\left(\begin{smallmatrix}1&1\\0&1\end{smallmatrix}\right)$.
Its first vector is $e_1+e_2$; $e_0$ is the zero-valuation
period and $e_3$ the other angular period.
The logarithmic disc terms cancel for these invariant
vectors. Their unit multipliers act on the ruled
normalization and commute with its boundary gluing.
At the hexagonal model the two toric rotations extend
in every chart. In each case the time-one map is a period.
Sweeping the extending section disc gives the relative
three-chains with these normal-torus boundaries.
The local integral specialization gives source ranks $3,3,2$.

The eight edge columns of these tori are
\[
 L=\begin{bmatrix}
 0&0&-A_2A_+\mathsf B_-\\
 \mathsf B_2&0&-A_+\mathsf B_-\\
 0&\mathsf B_+&-\mathsf B_-
 \end{bmatrix}.
\]
The last peripheral cut is
$j_-(\lambda)=-(A_2A_+\lambda,A_+\lambda,\lambda)$.
For a graph-kernel vector with components $a,b,c$,
the defining equations are
\[
 a_2=-4a_0,\qquad b_1=2a_0+a_1-3b_0,\qquad
 c_2=2a_0+a_1+c_1.
\]
Writing $h_{\rm edge}=2a_0+a_1$, its inverse coordinates are
\[
 (a,b,c)=a_0(\ell,0,0)+L\eta,
\]
\[
 \eta=(b_0,b_2-2h_{\rm edge},b_3-a_3-h_{\rm edge},c_0,c_1,
                 c_3-a_3-h_{\rm edge},-h_{\rm edge},-a_3-h_{\rm edge})^t.
\]
Substitution in either direction is the identity.
Thus $(\ell,0,0)$ and the eight $L$ columns form an
integral basis of $\ker D$. Also
$j_*(e_3)=-L_3-L_6-L_8$.

Lift $(\ell,0,0)$ by $t=T_{\ell,*}$ and the other columns
by the relative sweep tori. Together with $s$ these give
an integral basis of $H_2(V)$: cut, apply the displayed
inverse coordinates, and subtract the sweeps to leave
the unique fibre coordinate.
The closed flat vertical form $\Omega$ associated with
$\xi$ evaluates to one on $s$ and zero on every sweep.
The quotient by the eight reduced tori and its inverse are
\[
 \mathcal Q_{\rm red}\xrightarrow{\ \sim\ }\mathbb Z^2,\qquad
 [Z]\longmapsto\left(\int_Z\Omega,e_0^\vee(p_{\rm cut}(Z)_*)\right),
 \qquad (a,b)\longmapsto as+bt.
\]
Their span is primitive, and their relative source groups
have total rank eight. The further image in the free abelian group
$\mathcal Q_{\rm red}$ has rank zero and hence vanishes.
This proves $H_2(X_D)=\mathcal Q_{\rm red}=\mathbb Z\langle s,t\rangle$.

For the finite coinvariants,
$\operatorname{im}(A_*-I)$ has basis
$e_1+2e_2,-e_1+e_2+e_3$, with preimages $-e_2,e_1$.
Its quotient map is
\[
 q_*(x)=(x_0,-2x_1+x_2-3x_3),
\]
with section $(a,b)\mapsto ae_0+be_2$.
The difference from that section is
$(x_1+x_3)(e_1+2e_2)+x_3(-e_1+e_2+e_3)$.
Consequently
\[
 H_1(B)=
 \mathbb Z\langle e_0,e_2,g\rangle/
 \mathbb Z(-ne_0+3re_2+mg).
\]
Choose $m\alpha+n\beta=1$ and set
$z_0=\alpha e_0+\beta g$. The inverse formulas
\[
 e_0=mz_0+3r\beta e_2,\qquad
 g=nz_0-3r\alpha e_2
\]
give the integral basis $(e_2,z_0)$.
In the positive boundary basis $(e_0,e_2,x_*)$,
the map and its section are
\[
 J=\begin{pmatrix}3r\beta&1&3r\alpha\\m&0&-n\end{pmatrix},
 \qquad
 C_\partial=\begin{pmatrix}0&\alpha\\1&0\\0&-\beta\end{pmatrix}.
\]
Writing $\mu_1=(n,-3r,m)^t$, we have
$JC_\partial=I$, $J\mu_1=0$, and
$C_\partial J-I=-\mu_1(\beta,0,\alpha)$.
Thus the kernel is exactly $\mathbb Z\mu_1$.
The full degree-one Mayer--Vietoris kernel is now the
kernel of
\[
 G=\begin{pmatrix}
 n&0&0&0\\m&0&0&-1\\0&1&0&-1\\0&0&1&-1
 \end{pmatrix}.
\]
Its determinant is $-n\ne0$, so the pair quotient above gives
$H_2(X)=\mathcal Q_{\rm red}/\operatorname{im}\tau_*$.

Use positive normal-circle tori for the finite relative
generators. A loop $c$ is represented by the relative
class $-c\times D^2$, whose boundary is
$c\times\partial D^2$ in the order loop, normal circle.
Through the logarithmic overlap its covering torus has
parameterization
\[
 (a,b)\longmapsto(a\lambda+bv,mb).
\]
Its fibre cut is $\mathcal N_m\lambda$ and its integral
against $\Omega$ is $\xi(\lambda,v)$.
Here
\[
 \mathcal N_me_2=-ke_3,\qquad
 \mathcal N_me_0=m\ell,\qquad
 \xi(e_2,v)=2n,\qquad \xi(e_0,v)=6r.
\]
The $e_3$ sweep dies in $\mathcal Q_{\rm red}$, so
$\tau_*(e_2)=(2n,0)$ and $\tau_*(e_0)=(6r,m)$.
Applying the homomorphism to $mg=ne_0-3re_2$ gives
$m\tau_*(g)=(0,mn)$. Freeness of $\mathcal Q_{\rm red}$ gives
$\tau_*(g)=(0,n)$ and $\tau_*(z_0)=(6r\alpha,1)$.
The finite matrix is therefore
\[
 B_+=\begin{pmatrix}2n&6r\alpha\\0&1\end{pmatrix},
 \qquad
 \begin{pmatrix}1&-6r\alpha\\0&1\end{pmatrix}B_+
 =\begin{pmatrix}2n&0\\0&1\end{pmatrix}.
\]
In the positive coordinates, the cyclic map and its inverse are
\[
 [(a,b)]\longmapsto a-6r\alpha b\pmod{2n},\qquad
 [1]\longmapsto[s].
\]
Subtracting $b$ times the second column leaves
$(a-6r\alpha b,0)$; this is in the first-column image
exactly when its first coordinate is divisible by $2n$.
Hence $H_2(X)=\mathbb Z/(2|n|)$.

The two rank-one fibres, the bielliptic fibre and the
smooth tori have Euler number zero. The hexagonal fibre
has Euler number $6-3\cdot2+2=2$, giving $\chi(X)=2$.
Duality and universal coefficients give
$H_5=0$ and
$H_4=\operatorname{Ext}(H_1,\mathbb Z)=\mathbb Z/|n|$.
The Euler characteristic forces $b_3=0$, and then
$H_3=\operatorname{Ext}(H_2,\mathbb Z)=\mathbb Z/(2|n|)$.
Explicitly, the torsion characters taking $[e_0]$ to
$1/|n|$ and $[s]$ to $1/(2|n|)$ in $\mathbb Q/\mathbb Z$
map under the integral connecting homomorphism and cap
product with $[X]$ to the respective $H_4$ and $H_3$
generators. The inverse maps use inverse cap product
and these characters. Connectedness and orientation give
$H_0=H_6=\mathbb Z$.
Every $k>0$ is admissible for $n=\pm1$, giving the
stated simply connected rational homology spheres.

At $(k,n,r)=(1,1,0)$, the vector in Wang's construction
is exactly $\alpha^{\rm per}_3+2\alpha^{\rm per}_2-4\alpha^{\rm per}_1=\ell$.
The identity on the cubic root-disc family intertwines
the affine actions $t_{\ell/3}\rho$ and their logarithmic
overlaps. The identity maps on the three Mumford charts
intertwine their period quotients and agree on the overlaps.
They descend and glue to the isomorphism with Wang's model.
\end{proof}

\Needspace{6\baselineskip}
\begin{remark}[Signed coordinates for the primitive $I_2$ boundary map]
\hypertarget{tag-Q112}{}\label{mapsIII:rem:wang-boundary-orientation}
{\normalfont\scriptsize tag Q112}\par
In the finite boundary calculation of \cref{maps:src:D.2.7}, choose
$3k\alpha+n\beta=1$. The positive geometric matrix is
$B_+=\left(\begin{smallmatrix}2n&6r\alpha\\0&1\end{smallmatrix}\right)$.
Negating both relative generators and replacing $t$ by
$-t$ gives the exact alternative convention
\[
 \operatorname{diag}(1,-1)\,B_+(-I_2)
 =\begin{pmatrix}-2n&-6r\alpha\\0&1\end{pmatrix}.
\]
Both changes are integral involutions.
The source map negates both relative generators; the target map fixes
$s=[e_1\wedge e_2]$ and negates $t$. The reverse maps are the same
two involutions.
\end{remark}

\subsection{The oriented type of the Wang family}

\begin{definition}[Conic surgery]
\hypertarget{tag-Q108}{}\label{mapsIII:def:conic-surgery}
{\normalfont\scriptsize tag Q108}\par
The conic-surgery manifold is
\[
 M_{\rm conic}=
 \bigl(\mathbb{CP}^3\setminus\operatorname{int}(S^2\times D^4)\bigr)
 \cup_{S^2\times S^3}(D^3\times S^3),
\]
where the removed sphere is the conic
$[s:t]\mapsto[s^2:st:t^2:0]$, with its oriented normal framing.
\end{definition}

\begin{theorem}[The oriented type of the simply connected Wang members]\hypertarget{tag-02LY}{}\label{maps:src:D.2.8}
{\normalfont\scriptsize tag 02LY}\par
Every member of \cref{maps:src:D.2.7} with $n=\pm1$, $k>0$ and
$r\in\mathbb Z$ is orientation-preservingly diffeomorphic to
$M_{\rm conic}$ of Definition~\ref{mapsIII:def:conic-surgery}.
For the nonzero class $a\in H^2(X;\mathbb F_2)$,
\[
 \langle a^3,[X]_2\rangle=1.
\]
\end{theorem}


\begin{proof}
By \cref{maps:src:D.2.7}, the manifold $X$ is simply connected and
\[
 H_2(X;\mathbb Z)=H_3(X;\mathbb Z)=\mathbb Z/2,
 \qquad H_1(X;\mathbb Z)=H_4(X;\mathbb Z)=H_5(X;\mathbb Z)=0.
\]
The marked generator of $H_2$ is the regular-fibre torus
$s=[e_1\wedge e_2]$. Universal coefficients give
$H^2(X;\mathbb Z)=0$ and $H^4(X;\mathbb Z)=\mathbb Z/2$.
Consequently
\[
 w_2(TX)=\rho_2c_1(TX)=0,
 \qquad p_1(TX)=c_1(TX)^2-2c_2(TX)=0.
\]
The class $a$ is characterized by $a(s)=1$.

\emph{The involution and its fixed curves.}
In the multiplicative coordinates $(z,v)$ of the construction,
the positive phase circles are $e_2,e_3$. Multiplication in $v$
extends to a circle action on $X$: it commutes with the periods,
extends in the monomial charts, and commutes with the finite root
action, whose linear part fixes $e_3$. Let $\iota$ be multiplication
by $-1$ in $v$, or translation by $\tau=e_3/2$ in real markings.

At the primitive $I_2$ filling, the monomial change
\[
 x=v,\qquad y=v^2/z,\qquad z=x^2/y
\]
has determinant one. The vanishing phase and logarithmic period are
\[
 b=2e_2+e_3,\qquad h=-e_1,
\]
and the remaining elliptic periods are $-e_2$ and $e_0-3e_1$.
These are the rank-one coordinates of
\cite[Proposition~2.7, pp.~4--5]{manWang26}.
The involution is $(x,y)\mapsto(-x,y)$; it fixes exactly the elliptic
double curve $E$, and negates both coordinates in the transverse
node $x_+x_-=q$. At $\alpha_+$ it is a nontrivial half-period
translation on the surviving elliptic factor and has no fixed point.

At $\alpha_-$ the toric lattice has phase basis $(e_2,e_3)$ and the
periodic antidiagonal triangulation of
\cite[Proposition~2.8, p.~6]{manWang26}. In its two triangle types,
$e_3$ has weights $(-1,0,1)$ and $(-1,1,0)$. The vertical edges
form one orbit, giving one fixed curve $C\cong\mathbb{CP}^1$.
Its endpoints belong to the two distinct triangle orbits.
The horizontal and diagonal edge directions do not contain
$e_3/2$ modulo the phase lattice. Thus $C$ is the entire fixed
locus in this block, and both its normal characters are nontrivial.

There is no fixed point in the finite block. Write $m=3k$ and
$v_{\log}=n\ell+re_3$. A fixed central point would lift to an equation
\[
 x+\tau=A_*^jx+\frac{jv_{\log}}m\pmod\Lambda,
 \qquad 0\leq j<m.
\]
The invariant covector $e_0^*$ gives $jn/m\in\mathbb Z$.
Since $n=\pm1$, this forces $j=0$, after which the equation would
make $e_3/2$ an integral period. Off the central fibre, a nonzero
$j$ also rotates the root parameter. We have proved
\[
 X^\iota=E\sqcup C.
\]

\emph{A line class and its two weights.}
Let $Y$ be the complement of the interior of the finite block.
The exact reduced-boundary maps in \cref{maps:src:D.2.7} give
\[
 \pi_1(Y)=\mathbb Z\langle e_0\rangle,
 \qquad H_2(Y;\mathbb Z)=\mathbb Z\langle s,t\rangle.
\]
Here $t$ is the fixed sweep at the finite boundary. The exterior
maps onto this second homology, with kernel the eight reduced sweep
relations. The inclusion $Y\to X$ sends $s$ to its marked generator
and $t$ to $-6r\alpha s$, where $3k\alpha+n\beta=1$.
In particular $a|Y$ has values $(1,0)$ on $(s,t)$.

Put
\[
 \xi=e_1^*\wedge e_2^*+6e_0^*\wedge e_3^*,
 \qquad Q(x)=x_1x_2+x_1+x_2\pmod2.
\]
Substitution in the four monodromy matrices gives
$A_i^*\xi=\xi$ and $Q(A_ix)=Q(x)$, and
\[
 Q(\lambda+\mu)-Q(\lambda)-Q(\mu)
 =\xi(\lambda,\mu)\pmod2.
\]
Thus the factors
\[
 j_\lambda(x)=(-1)^{Q(\lambda)}
 e^{\pi i\xi(\lambda,x)}
\]
define a line bundle on the regular real torus family. The
monodromy lift is $(x,\zeta)\mapsto(A_ix,\zeta)$, where $\zeta$
is the line coordinate. Its first Chern class is
$\xi$; with $\nabla=d-\pi i\xi(x,dx)$ its curvature has first
Chern form $\xi$. Moreover,
\[
 T_\tau(x,\zeta)=
 \bigl(x+\tau,e^{\pi i\xi(\tau,x)}\zeta\bigr)
\]
commutes with the lattice factors, since
$\xi(\tau,\lambda)=-3\lambda_0$ is integral. Its square is the
lattice identification for $e_3$, because $Q(e_3)=0$. It gives an
involutive lift commuting with all monodromies.

We describe its extension and weight in a nodal chart. Let $b$ be
the primitive vanishing phase, $h$ its logarithmic period, and
$N=\xi(b,h)$ even. Write
\[
 x=ub+wh,\qquad z=e^{2\pi i(u+\omega w)},
 \qquad q=e^{2\pi i\omega}.
\]
The parameter monodromy identifies $(\omega+1,u,w)$ with
$(\omega,u+w,w)$. The change of frame
\[
 G(u,w,\omega)=e^{-\pi iNuw-\pi iN\omega w^2}
\]
is used by $\zeta'=G\zeta$, with inverse $\zeta=G^{-1}\zeta'$.
It respects this identification and changes the two period factors to
\[
 j'_b=1,\qquad
 j'_h=(-1)^{Q(h)}q^{-N/2}z^{-N}.
\]
The period $h$ sends $z$ to $qz$. The canonical lift of $b/2$
is the identity on the line scalar in this frame and sends
$z$ to $-z$. The integer Cartier values $f(j)=Nj^2/2$ give the
nodal extension. Indeed their period difference is $Nj+N/2$;
for $D=\sum f(j)D_j$, multiplication by the inverse character
$q^{-N/2}z^{-N}$ identifies $\mathcal O(D)$ with
$h^*\mathcal O(D)$. If $m_\sigma$ interpolates the Cartier values,
the regular frame and its fibre coordinate are
\[
 e_\sigma=\chi^{-m_\sigma}e,\qquad
 \zeta_\sigma=\chi^{m_\sigma}\zeta',
 \qquad \zeta'=\chi^{-m_\sigma}\zeta_\sigma.
\]
The half-phase lift is the identity on the dense line scalar $\zeta'$,
so its forward action on $\zeta_\sigma$ has character $+m_\sigma$;
pullback of the frame has the inverse character. On the edge
from $0$ to $1$, the fibre character is $+N/2$, with involution
weight $N/2$ modulo two.
For odd $N$ the phase sign is $Q(b)=1$. On the phase cover use
$\zeta_{\rm odd}=z^{1/2}\zeta_{\rm sym}$, with inverse
$\zeta_{\rm sym}=z^{-1/2}\zeta_{\rm odd}$. The phase factor becomes
one and the logarithmic factor becomes
$(-1)^{Q(h)}q^{-(N-1)/2}z^{-N}$, giving the integral Cartier
values $(Nj^2-j)/2$.

At the rank-two block these constructions give
\[
 \psi(r,s)=\frac{r(r-1)}2-3s^2.
\]
Its differences under the two period translations are $r$ and
$-6s-3$. Interpolation on each unimodular triangle gives integral
characters on every toric chart. The holomorphic unit multipliers
enter these same characters; their logarithms on the parameter disc
change none of the valuation exponents. The nodal frame calculation
identifies the involutions on the overlaps. At the other rank-one
block the vanishing phase is $e_1+e_2$, its degree is one, and the
odd-degree construction applies. These local line bundles glue to
an equivariant line bundle $L$ on $Y$.

For completeness the free reduced block uses the quotient lattice
$\Lambda'=\Lambda+\mathbb Z\tau$ with the semicharacter
\[
 Q'(\lambda+j\tau)=Q(\lambda)+j\xi(\lambda,\tau)\pmod2.
\]
Changing $(\lambda,j)$ to $(\lambda+ke_3,j-2k)$ changes this
expression by $-6k\lambda_0$, so it is well defined.
Its polar form is $\xi$ modulo two, and it is preserved by all
monodromies. Put
\[
 j'_v(x)=(-1)^{Q'(v)}e^{\pi i\xi(v,x)},\qquad v\in\Lambda'.
\]
The factor for $\lambda+j\tau$ is exactly the composite
$J_\lambda T_\tau^j$ of the original lattice and involution lifts.
Thus the quotient bundle map is
$\operatorname{orbit}([x,\zeta]_\Lambda)\mapsto[x,\zeta]_{\Lambda'}$,
with inverse given by the orbit of the same representative.
Its pullback under the torus covering is the original bundle by
$[x,\zeta]_\Lambda\mapsto([x]_\Lambda,[x,\zeta]_{\Lambda'})$.
The inverse sends
$([x]_\Lambda,[x+v,\eta]_{\Lambda'})$ to
$[x,j'_v(x)^{-1}\eta]_\Lambda$; the cocycle identity makes it
independent of the representative $v$.
This identifies the local quotient extension with the stated
equivariant extension on its overlap.

The Chern class evaluates to $1$ on $s$ and to zero on the section
sweeps, including $t$. The reduced boundaries have free $H_1$, so
their integral $H^2$ classes are determined by these periods; the
local extensions and the exterior class agree integrally. Hence
\[
 \rho_2c_1(L)=a|Y.
\]
At $I_2$ the nodal degree is
$N=\xi(2e_2+e_3,-e_1)=2$, and
\[
 \tau=\frac{2e_2+e_3}{2}-e_2.
\]
For an integral period $\lambda$, write its actual lattice
identification and its canonical translation lift as
\[
 \begin{aligned}
 J_\lambda(x,\zeta)
   &=(x+\lambda,j_\lambda(x)\zeta),\\
 T_\lambda(x,\zeta)
   &=(x+\lambda,e^{\pi i\xi(\lambda,x)}\zeta)
    =(-1)^{Q(\lambda)}J_\lambda(x,\zeta).
 \end{aligned}
\]
Here the scalar in the last expression multiplies the line coordinate.
Since $J_\lambda$ is the quotient identification, $T_\lambda$ acts on
the quotient line by $(-1)^{Q(\lambda)}$. The phase plane is
$\xi$-isotropic, so
$T_\tau=T_{(2e_2+e_3)/2}T_{-e_2}$ with no extra scalar.
The integral correction therefore adds $Q(-e_2)=Q(e_2)=1$ to the
half-phase weight.
The weight on $E$ is therefore $1+1=0$. At $C$ we have
$b=e_3$, $h=e_0$ and $N=-6$, giving weight $-3=1$ modulo two.

The ordinary degrees are odd on both fixed curves. On $E$ it is
\[
 \xi(-e_2,e_0-3e_1)=-3.
\]
On the normalized cusp, the six ray values of $\psi$ are
$0,-3,-2,1,-3,-3$. Summing the corresponding toric divisors gives
\[
 c_1(L|_{\mathrm{dP}_6})=-5H-E_A+6E_B.
\]
The vertical curve is $E_A$, so the degree on $C$ is $1$.

\emph{The cubic evaluation.}
The class $a$ extends to the Borel cohomology of the circle action:
the only possible outgoing differentials from bidegree $(0,2)$
have targets $H^2(BS^1;H^1(X;\mathbb F_2))$ and
$H^3(BS^1;\mathbb F_2)$, both zero. Restricting to
$G=\langle\iota\rangle\cong\mathbb Z/2$,
choose $A\in H^2_{\mathbb Z/2}(X;\mathbb F_2)$ with ordinary
restriction $a$. Put $t\in H^1(B(\mathbb Z/2);\mathbb F_2)$ for
the nonzero class.

The fixed point on $C$ splits the fundamental-group extension for
$Y_{\mathbb Z/2}$. The circle action induces the identity on
$\pi_1(Y)=\mathbb Z$, so
$\pi_1(Y_{\mathbb Z/2})=\mathbb Z\times\mathbb Z/2$.
Let $\lambda\in H^1(Y_{\mathbb Z/2};\mathbb Z)$ be projection to
the first factor. Its restriction to any fixed point is zero.
The degree-two Borel filtration consequently gives
\[
 A|Y-\rho_2c_1^{\mathbb Z/2}(L)
 =b t^2+c\rho_2(\lambda)t,
 \qquad b,c\in\mathbb F_2.
\]
These are all its kernel terms: $H^1(Y;\mathbb F_2)$ is
one-dimensional, and both displayed classes survive, by the fixed
point and the integral lift $\lambda$. Thus the two fixed-point
$t^2$ coefficients of $A$ have sum $1$. On the fixed curves write
\[
 A|E=a_E+b_Et+c_Et^2,
 \qquad A|C=a_C+c_Ct^2,
 \qquad c_E+c_C=1,
\]
where $a_E\in H^2(E;\mathbb F_2)$ and
$a_C\in H^2(C;\mathbb F_2)$ each evaluate to one on their curve,
and $b_E\in H^1(E;\mathbb F_2)$.
Every such $b_E$ lifts integrally, so $b_E^2=0$.

The involution acts by $-1$ on each complex normal fibre.
Over $F\times BG$ the equivariant normal bundle is therefore
$\operatorname{pr}_F^*\nu_F\otimes\operatorname{pr}_{BG}^*\varepsilon$,
where $\varepsilon$ is the complex sign line and
$\rho_2c_1(\varepsilon)=t^2$. Its equivariant Euler class is
$t^4+t^2\rho_2c_1(\nu_F)+\rho_2c_2(\nu_F)=t^4$:
$c_1(\nu_F)=-c_1(TF)$ is even, and two ordinary degree-two
classes have zero product on the curve.
After inverting $t$, the fixed-locus Thom map gives the localization
isomorphism. Indeed the involution is free off a tubular neighbourhood
of its fixed set; the Borel construction of that complement is its
finite-dimensional quotient, whose cohomology is killed after this
localization. The inverse of the Thom restriction divides each fixed
component by its Euler class. Therefore
\[
 \begin{aligned}
 \langle a^3,[X]_2\rangle
 &=\int_E\frac{(a_E+b_Et+c_Et^2)^3}{t^4}
   +\int_C\frac{(a_C+c_Ct^2)^3}{t^4}\\
 &=c_E+c_C=1.
 \end{aligned}
\]
In the second equality $b_E^2=0$, and the only term of curve degree
two in each numerator is $c_Fa_Ft^4$. Equivariant integration has
total degree zero here, so it equals the ordinary cubic evaluation.

\emph{The conic surgery and the oriented type.}
Let $\Gamma\subset\mathbb{CP}^3$ be the conic of
Definition~\ref{mapsIII:def:conic-surgery}.
It represents twice the positive line
class, since the pullback of $\mathcal O(1)$ is
$\mathcal O_{\mathbb{CP}^1}(2)$. The normal sequence for
$\Gamma\subset\mathbb{CP}^2\subset\mathbb{CP}^3$ is
$0\to\mathcal O(4)\to\nu_\Gamma\to\mathcal O(2)\to0$;
it splits because $H^1(\mathbb{CP}^1;\mathcal O(2))=0$.
Its underlying oriented real normal bundle has rank four and
$w_2=0$, hence is trivial. Its framing is unique
up to homotopy because $\pi_2SO(4)=0$. Let $V$ be the oriented
seven-dimensional index-three surgery trace from
$\mathbb{CP}^3$ to $M_{\rm conic}$.

Write $C_0=\mathbb{CP}^3\setminus\operatorname{int}N(\Gamma)$
for the conic complement. Thom excision gives the map
$H_4(\mathbb{CP}^3)\to H_4(\mathbb{CP}^3,C_0)$ as multiplication
by two. Thus $H_2(C_0)=\mathbb Z$ and $H_3(C_0)=\mathbb Z/2$.
The boundary $S^2$ maps to twice the $H_2$ generator, and the
boundary $S^3$ maps to the nonzero torsion class. The surgery
Mayer--Vietoris maps in these degrees are consequently
\[
 \mathbb Z\longrightarrow\mathbb Z,\quad j\longmapsto2j,
 \qquad
 \mathbb Z\longrightarrow\mathbb Z/2\oplus\mathbb Z,
 \quad j\longmapsto([j]_2,j).
\]
They give $H_2(M_{\rm conic})=H_3(M_{\rm conic})=\mathbb Z/2$
and the same remaining homology as $X$. The conic complement is
simply connected by transversality, and van Kampen gives the same
for the surgery. The spin structure on $\mathbb{CP}^3$ extends over
the framed handle; the sphere has its unique bounding spin structure.

The relative handle map in the trace is multiplication by two on
$H_2(\mathbb{CP}^3)$, so $H_2(V)=\mathbb Z/2$ and $H_3(V)=0$.
Viewed from the other boundary, the handle has index four, so
$H_2(M_{\rm conic})\to H_2(V)$ is an isomorphism. Hence a class
$A_V\in H^2(V;\mathbb F_2)$ restricts to the reduced hyperplane
class and to the nonzero class on $M_{\rm conic}$. Evaluation on
$\partial[V]_2$ gives
\[
 \langle a_{\rm conic}^3,[M_{\rm conic}]_2\rangle
 =\langle h^3,[\mathbb{CP}^3]_2\rangle=1.
\]
Also $H^4(V;\mathbb Z)\to H^4(\mathbb{CP}^3;\mathbb Z)$ is an
isomorphism. Let $B$ be the class restricting to $h^2$. Tangent-bundle naturality
and $p_1(\mathbb{CP}^3)=4h^2$ give $p_1(V)=4B$. Its restriction
to $H^4(M_{\rm conic};\mathbb Z)=\mathbb Z/2$ is zero.

Apply Andrus's oriented classification
\cite[Theorem~5.2, p.~42; definitions in \S\S3.1--3.2,
pp.~17--21]{manAndrus77} with cyclic homology order two. Its
augmented invariant consists here of the even mod-two cohomology
ring, the orientation reduction and $p_1$. The additional cubing
operation in that theorem concerns orders divisible by three.
For both manifolds the ring has basis $(1,a,d,z)$ with
$a^2=d$ and $ad=z$, and the distinguished $p_1$ is zero.
The map sending these four basis classes to their counterparts is
an orientation-preserving augmented-ring isomorphism, with inverse
given by the reverse assignments. The theorem yields the required
orientation-preserving diffeomorphism.

For Wang's sufficiently small nonzero parameters $\lambda$,
the construction identification in \cref{maps:src:D.2.7} at
$(k,n,r)=(1,1,0)$ identifies the member with $X_2(\lambda)$.
The resulting oriented type agrees with
\cite[Theorem~7.1, p.~15; Appendix~B, pp.~17--19]{manWang26}.
\end{proof}

\subsection{The \texorpdfstring{$III^*+II+I_1$}{III*+II+I1} profile}

\begin{theorem}[Logarithmic multiplicities over $III^*+II+I_1$]\hypertarget{tag-02LZ}{}\label{maps:src:D.2.9}
{\normalfont\scriptsize tag 02LZ}\par


Use the $III^*,II,I_1$ surface and the unchanged type-$II$
filling of Liu--Rollenske, with the degree-one linearization zero
at $I_1$ and the primitive hexagonal filling there.
In their homology basis $(e_1,e_2,e_3,e_4)$ put
\[
 \begin{gathered}
 \delta=-e_1+2e_3,\qquad \ell=e_4,\qquad
 m=4k,\qquad v=n\ell+r\delta,\\
 k\in\mathbb Z_{>0},\qquad n,r\in\mathbb Z,\qquad
 \gcd(n,4k)=1.
 \end{gathered}
\]
The logarithmic transformation at $III^*$ exists for every
such integer datum, and its total space satisfies
\[
 \pi_1(X)=\mathbb Z/|n|,\qquad
 H_0=H_6=\mathbb Z,\quad H_5=0,\quad
 H_1=H_2=H_3=H_4=\mathbb Z/|n|.
\]
Thus $n=\pm1$ gives $S^6$ for every $k>0$. In particular,
the Liu--Rollenske configuration also answers
\cite[Question~4, p.~2]{manViaclovskyLogs26} affirmatively.
\end{theorem}


\begin{proof}
Write $m=4k$ and $\eta=e_3-e_2$. In the ordinary homology marking of
\cite[v2, Propositions~3.6 and~3.13, pp.~9 and~14]{manLR26}, the
positive local monodromies are
\[
\begin{gathered}
A_0=\begin{pmatrix}-1&-1&-1&0\\2&1&1&0\\0&0&1&0\\0&0&0&1\end{pmatrix},
\qquad
A_1=\begin{pmatrix}1&1&0&1\\-2&0&-1&-1\\0&-1&1&-1\\0&0&0&1\end{pmatrix},\\
A_\infty=\begin{pmatrix}1&0&0&-1\\-2&0&-1&1\\2&1&2&1\\0&0&0&1\end{pmatrix},
\qquad A_\infty A_1A_0=I.
\end{gathered}
\]
Here $A_0^4=I$, while $N_\infty=A_\infty-I$ satisfies
$N_\infty^2=0$ and
\[
\operatorname{im}N_\infty=\ker N_\infty
 =\mathbb Z\delta+\mathbb Z\eta.
\]

\emph{The finite quotient and its attachment.}
Use the good reduction of
\cite[Lemma~2.7, Corollary~2.8 and Proposition~3.9,
pp.~6 and~10--11]{manLR26}, followed by the degree-$k$ pullback.
Write $b=u^m$ for the local base change. Its generator covers
$u\mapsto e^{2\pi i/m}u$ and induces the fourth-root deck map on $u^k$.
The positive-root deck action on the central product cover is
\[
\rho(c,w)=(-ic,w+1/4).
\]
Its periods are $e_1=(2,0)$, $e_2=(1+i,0)$,
$e_3=(1,1/2)$, $e_4=(0,\lambda)$ and
$\delta=(0,1)$. Thus its linear part is $R=A_0^{-1}$ and
its affine lift satisfies $\rho^m=t_{k\delta}$, where $t_a$
denotes translation by $a$. Since $v$ is invariant, the action
\[
g=t_{v/m}\rho
\]
has order $m$ on the compact family. For $0<j<m$, the translation
of $g^j$ detected by $e_4^*$ is $jn/m$, which is nonintegral.
Hence the action is free on the central torus. Its base action is
free away from the origin. The quotient is therefore smooth and
has central multiplicity $m$.

The period section
\[
\sigma_v(u)=\frac{\log u}{2\pi i}v(u)
\]
is single-valued as a torus section. Translation by $\sigma_v$
conjugates $\rho$ to $g$, so translation by $-\sigma_v$ identifies
the new punctured neighborhood with the original one.
Its inverse is translation by $\sigma_v$.

We compare this affine origin with the extending section of the
original construction. In the Weierstrass equation of
\cite[Proposition~2.1, p.~3]{manLR26}, good reduction uses
$x=u^{6k}X$, $y=u^{9k}Y$. Near the zero section the conormal
frames satisfy
\[
-x/y=u^{-3k}(-X/Y).
\]
The deck map sends $(X,Y)$ to $(-X,-iY)$, so its action on the
conormal line has coefficient $i$. Thus its affine translation is
$\delta/4$ throughout this marked collar.
The section from \cite[Proposition~2.4, p.~4]{manLR26} is
nonzero at $b=0$, since its sole zero is at infinity. Its coefficient
in the good-reduction frame is consequently $u^{-3k}p(u^m)$,
with $p$ holomorphic and nonvanishing on the disc.
One positive base loop changes its angular lift by $-3\delta/4$,
whereas $\rho$ changes that lift by $\delta/4$.
Thus the positive exterior meridian $x_0$ and the affine generator
are related by
\[
x_0=t_{-\delta}g,
\qquad x_0^m=t_{v'},
\qquad v'=n\ell+q\delta,
\qquad q=r-3k.
\]
The parameter change has inverse $r=q+3k$.

The radial and unit terms in this comparison also give an explicit
isotopy of the boundary map. In the flat real marking write
$u=\varrho e^{i\theta}$ and let $J_u$ be the fibre complex structure.
Put $\zeta=e^{2\pi i/m}$. Then $J_{\zeta u}R=RJ_u$.
Choose a holomorphic logarithm of $p$ on the disc, put
$a(b)=\log p(b)/(2\pi i)$, and set
\[
u_p(u)=\operatorname{Re}a(u^m)\delta
       +\operatorname{Im}a(u^m)J_u\delta,
\qquad
B_{\rm per}(u)=\frac{\log\varrho}{2\pi}J_uv'-u_p(u).
\]
The attachment in section-based coordinates is
\[
h_{\rm att}(y,u)=\left(y-\frac{\theta}{2\pi}v'+B_{\rm per}(u),u\right).
\]
The vector $B_{\rm per}$ is periodic in $\theta$ and satisfies
$B_{\rm per}(\zeta u)=RB_{\rm per}(u)$. Therefore
\[
h_{{\rm att},c}(y,u)=\left(y-\frac{\theta}{2\pi}v'+cB_{\rm per}(u),u\right),
\qquad 0\le c\le1,
\]
is an isotopy on the quotient collar, with inverse
$(z,u)\mapsto(z+\theta v'/(2\pi)-cB_{\rm per}(u),u)$.
Indeed its fibre-coordinate deck identity is
$h_{{\rm att},c}g=Rh_{{\rm att},c}+\delta$; the last term is an integral period.
Multiplying $B_{\rm per}$ by a radial cutoff extends the translation isotopy
by the identity to the exterior. Its inverse uses the negative
translation. The resulting integral prism homotopy preserves all
boundary homology classes. We may consequently calculate the
normal boundary using $h_{{\rm att},0}$, for which the positive normal disc has
meridian
\[
d_0=t_{-v'}x_0^m.
\]
The holomorphic attachment $h_{\rm att}=h_{{\rm att},1}$ gives the compact complex
threefold in the statement.

\emph{The two reduced fillings.}
Let $V$ be the torus bundle over the complementary pair of pants,
and write $B_1,B_\infty$ for the type-$II$ and hexagonal
neighborhoods, with boundaries $E_1,E_\infty$.
The global section extends over both of these neighborhoods by
\cite[Propositions~2.4--2.5, pp.~4--5]{manLR26}.
In the basis
\[
(f_1,f_2,f_3,f_4)=(e_1-e_3,e_2,\delta,\ell-e_2),
\]
the type-$II$ neighborhood is the product of its elliptic-surface
neighborhood and the two circles $f_3,f_4$;
see \cite[Lemma~3.10 and Proposition~3.11, pp.~11--12]{manLR26}.
The central curve is homeomorphic to $S^2$. The fundamental-group
map of the fibre is
\[
p_1(x_1,x_2,x_3,x_4)=(x_1+x_3,x_4),
\]
with section $(a,b)\mapsto a\delta+b(\ell-e_2)$ and kernel
$\mathbb Z(e_1-e_3)+\mathbb Ze_2$.
The fibre classes $f_1\wedge f_2$ and $f_3\wedge f_4$ map to the two generators
of $H_2(S^2\times T^2)$. This is an integral isomorphism from
the fibre coinvariants, since the elliptic block of $A_1-I$
has determinant one.
The two boundary sweeps of $\delta$ and $\ell-e_2$ bound the
section disc swept by these circle factors. The Wang sequence then
identifies them with the entire kernel of $H_2(E_1)\to H_2(B_1)$.

At infinity choose
\[
(h_1,h_2,h_3,h_4)=(\eta,e_2,\delta,e_2+\ell).
\]
This is an integral basis: the inverse coordinate map is
\[
(x_1,x_2,x_3,x_4)\longmapsto
(2x_1+x_3,\ 2x_1+x_2+x_3-x_4,\ -x_1,\ x_4).
\]
The monodromy sends $h_2$ to $h_2+h_1$ and $h_4$ to $h_4+h_3$.
The primitive Mumford model has a toric cover indexed by the
periodic triangulation of $\mathbb R^2$. Its toric components and
double curves are simply connected, and its triangle relations give
a simply connected cover. The deck group is $\mathbb Z^2$.
The fibre map is therefore
\[
p_\infty(x_1,x_2,x_3,x_4)
 =(2x_1+x_2+x_3-x_4,x_4),
\]
with section $(a,b)\mapsto ae_2+b(e_2+\ell)$ and kernel
$\mathbb Z\delta+\mathbb Z\eta$.

In the $h$-wedge basis $(12,13,14,23,24,34)$, the image of
$\bigwedge^2A_\infty-I$ is generated by
$h_{13}$ and $h_{14}+h_{23}+h_{13}$.
Its quotient map and section are
\[
(a,b,c,d,e,f)\mapsto(a,c-d,e,f),
\qquad (a,c,e,f)\mapsto(a,0,c,0,e,f).
\]
These sections together with the two relation vectors form an
integral basis. The local specialization isomorphism of
\cite[Proposition~3.12, p.~13]{manLR26} consequently identifies
this quotient with $H_2(B_\infty)=\mathbb Z^4$.
The angular torus action extends over the toric filling. Applying
its two loops $\delta,\eta$ to the extended section disc gives
the two boundary sweeps. Their projections form a basis of the
Wang invariant lattice, so they are the full kernel of
$H_2(E_\infty)\to H_2(B_\infty)$.

For each reduced filling, duality now gives
$H_3(B_i,E_i)=\mathbb Z^2$ and $H_2(B_i,E_i)=\mathbb Z$.
Also $H_1(B_i,E_i)=H^5(B_i)=0$, since the neighborhood retracts
onto its four-dimensional central fibre.
The boundary maps are respectively these two sweeps and the positive
section meridian $x_i$. We orient a sweep by its fibre loop followed
by its positive normal circle. The relative three-chain is the
negative of the loop times its complex-oriented normal disc,
whose boundary has that orientation. The degree-two relative
generators are the positively oriented discs themselves.

\emph{Integral homology of the exterior and the reduced attachment.}
The images of $A_0-I$ and $A_1-I$ generate
$\mathbb Ze_1+\mathbb Ze_2+\mathbb Ze_3$.
Indeed their columns give
\[
e_1=-(A_0-I)e_2,\quad
e_2=(A_0-I)(e_3-e_2),\quad
e_3=-(A_0-I)e_3-(A_1-I)e_2.
\]
The section thus gives
\[
H_1(V)=\mathbb Z\langle\ell,x_0,x_1\rangle,
\qquad x_\infty=-x_0-x_1.
\]
Write $e_{ij}=e_i\wedge e_j$ and
$e_{ij}^*=e_i^*\wedge e_j^*$. For the fibre part of $H_2(V)$ put
\[
\omega=2e_{12}^*-e_{23}^*-e_{34}^*,\qquad s=-e_{34},
\quad C_i=\bigwedge^2A_i-I.
\]
The form $\omega$ is monodromy invariant. The relations include
\[
\begin{aligned}
e_{13}&=-C_0e_{23},\\
e_{14}&=-C_0e_{24},\\
e_{24}&=C_0(e_{34}-e_{24}),\\
e_{12}+2e_{23}&=C_0(e_{13}-2e_{23}),\\
e_{23}-e_{34}&=C_1e_{24}+C_0(e_{13}-2e_{23}+e_{34}).
\end{aligned}
\]
They reduce every fibre class uniquely to $\omega(c)s$.
Hence its coinvariant quotient is $\mathbb Zs$, with quotient
$c\mapsto\omega(c)$ and inverse $1\mapsto s$.

The graph part is the kernel of
$D=[A_0-I\ \ A_1-I]$. With the fibre-then-base orientation its
cellular boundary is $-D$. Write a graph vector as $(p,q)$.
Its equations are
\[
p_3=-2p_1-p_2,\qquad q_3=-p_2-2q_1,\qquad q_4=-q_2.
\]
For a local invariant vector define
\[
j_0(x)=(x,0),\qquad j_1(x)=(0,x),\qquad
j_\infty(x)=(-x,-A_0x).
\]
Let
\[
E=j_0(\ell),\quad L_1=j_1(\delta),\quad
L_2=j_1(\ell-e_2),\quad
L_3=j_\infty(\delta),\quad L_4=j_\infty(\eta).
\]
These form an integral basis of $\ker D$: its coordinate inverse is
\[
(p,q)\longmapsto(p_4,p_1-q_1,-q_2,p_1,p_2).
\]
Choose their geometric sweep cycles as lifts, and denote the lift
of $E$ by $t$. The graph homology sequence then gives
\[
H_2(V)=\mathbb Z\langle s,t,L_1,L_2,L_3,L_4\rangle,
\qquad j_0(\delta)=-L_1-L_3.
\]
The invariant vertical form $\omega$ extends over the flat bundle
to a closed form $\Omega$, taking value one on $s$ and zero on
each chosen sweep. These cycles show that its periods are integral.
The collar translation isotopy preserves this integral class.

Set $X_D=V\cup B_1\cup B_\infty$. Its relative boundary in degree
three has columns $L_1,L_2,L_3,L_4$, and that in degree two is
\[
G_D=\begin{pmatrix}0&0\\0&-1\\1&-1\end{pmatrix}
\]
in the basis $(\ell,x_0,x_1)$. The inverse of its lower block sends
$(b,c)$ to $(c-b,-b)$. Exactness therefore gives
\[
H_2(X_D)=\mathbb Z\langle s,t\rangle,
\qquad H_1(X_D)=\mathbb Z\langle\ell\rangle.
\]
The first quotient drops the four $L_i$ coordinates and has section
given by the cycles $s,t$. On fundamental groups the two fillings
kill $e_1-e_3,e_2,\eta$ and $x_1,x_\infty$. The first three
vectors form an integral basis of $\mathbb Ze_1+\mathbb Ze_2+
\mathbb Ze_3$, and the base relation also kills $x_0$.
The boundary maps onto both filling groups are surjective, so
\[
\pi_1(X_D)=\mathbb Z\langle\ell\rangle.
\]

\emph{The finite relative maps.}
Let $B_0$ be the modified finite neighborhood and $G$ its reduced
central surface. Then $B_0$ is the normal disc bundle of the free torus
quotient $G$, so its relative groups are the Thom groups.
The group of $G$ has commuting lattice translations, conjugation
by $R$, and $x_0^m=t_{v'}$. The image of $R-I$ is
$\mathbb Ze_1+\mathbb Ze_2$, since
$R-I=-R(A_0-I)$. With $u_{\rm fib}=e_3$, its abelianization is
\[
mx_0=n\ell+2qu_{\rm fib}.
\]
Choose $\alpha,\beta\in\mathbb Z$ with $m\alpha+n\beta=1$,
and put $z=\alpha\ell+\beta x_0$. Then
\[
\ell=mz-2q\beta u_{\rm fib},\qquad x_0=nz+2q\alpha u_{\rm fib}.
\]
Thus $H_1(G)=\mathbb Z\langle u_{\rm fib},z\rangle$.
The boundary map on $H_1$, in the bases $(\ell,u_{\rm fib},x_0)$ and $(u_{\rm fib},z)$,
is
\[
J=\begin{pmatrix}-2q\beta&1&2q\alpha\\m&0&n\end{pmatrix}.
\]
Its section sends $u_{\rm fib}$ to $(0,1,0)$ and $z$ to $(\alpha,0,\beta)$,
and its kernel is generated by $(-n,-2q,m)$. The full inverse
coordinate map is
\[
(a,b,c)\longmapsto
(-2q\beta a+b+2q\alpha c,\ ma+nc,\ -\beta a+\alpha c),
\]
where the last coordinate is the coefficient of that kernel vector.
It follows that $H_3(B_0,E_0)=\mathbb Z^2$ and
$H_2(B_0,E_0)=\mathbb Z$, while $H_1(B_0,E_0)=0$,
with the signed Thom convention above.

For a fibre loop $\lambda$, the oriented normal transfer has graph
component $j_0(\mathcal N_m\lambda)$ and $\Omega$-value
$-\omega(\lambda,v')$. Indeed, the pure winding collar has fibre
coordinates $a\lambda-bv'$, with $0\le a,b\le1$; the translation isotopy
identifies its integral homology class with that for $h_{\rm att}$.
The norm identities are
\[
\mathcal N_m e_3=2k\delta,\qquad
\mathcal N_m\ell=m\ell.
\]
The transfer of $x_0$ has graph component $j_0(v')$ and zero
fibre value. Its $m$ lifts concatenate to $x_0^m=t_{v'}$.
In the commuting subgroup generated by $x_0$ and $v'$, the normal
torus is $x_0\wedge(mx_0-v')=v'\wedge x_0$; it is the invariant
sweep of $v'$, on which $\Omega$ evaluates to zero.
Consequently the transfers in $H_2(V)$ are
\[
\begin{aligned}
\tau(u_{\rm fib})&=ns-2k(L_1+L_3),\\
\tau(\ell)&=-2qs+mt,\\
\tau(x_0)&=nt-q(L_1+L_3),\\
\tau(z)&=-2q\alpha s+t-q\beta(L_1+L_3).
\end{aligned}
\]
They respect $m\tau(x_0)-n\tau(\ell)-2q\tau(u_{\rm fib})=0$.
After the reduced attachment the finite boundary matrix is
\[
B=\begin{pmatrix}n&-2q\alpha\\0&1\end{pmatrix},
\quad
\operatorname{coker}B\xrightarrow{\ \sim\ }\mathbb Z/|n|,
\quad [(a,b)]\longmapsto a+2q\alpha b,
\]
with inverse $[1]\mapsto[s]$. A vector in the kernel of this
quotient has preimage $((a+2q\alpha b)/n,b)$.

For the full pair $(X,V)$ the two relevant boundary matrices are
\[
\mathcal F=\begin{pmatrix}
n&-2q\alpha&0&0&0&0\\0&1&0&0&0&0\\
-2k&-q\beta&1&0&0&0\\0&0&0&1&0&0\\
-2k&-q\beta&0&0&1&0\\0&0&0&0&0&1
\end{pmatrix},
\qquad
\mathcal G=\begin{pmatrix}-n&0&0\\m&0&-1\\0&1&-1\end{pmatrix}.
\]
Adding $2k$ times columns three and five to the first column of
$\mathcal F$, and $q\beta$ times those columns to the second,
gives $\operatorname{diag}(B,I_4)$. Adding $2q\alpha$ times its
second row to its first gives $\operatorname{diag}(n,1,1,1,1,1)$.
All these changes have integral inverses obtained by subtraction.
Moreover $\det\mathcal G=-n\ne0$.
Exactness thus identifies $H_2(X)=\operatorname{coker}\mathcal F
=\mathbb Z/|n|$ and $H_1(X)=\operatorname{coker}\mathcal G
=\mathbb Z/|n|$.
Equivalently, for $(X,X_D)$ the final normal disc maps to
$-n\ell$, an injective map $\mathbb Z\to H_1(X_D)$.

The map $\pi_1(E_0)\to\pi_1(B_0)$ is surjective, with kernel
generated by the normal circle $t_{-v'}x_0^m$.
Van Kampen therefore adds precisely $\ell^n=1$ to
$\pi_1(X_D)$, proving $\pi_1(X)=\mathbb Z/|n|$.

The finite torus quotient and $S^2\times T^2$ have Euler
characteristic zero. The hexagonal fibre has two vertex classes
after opposite-edge gluing, while its open toric surface and open
double curves have Euler characteristic zero. Thus $\chi(X)=2$.
Rational duality gives $b_4=b_2=0$ and $b_5=b_1=0$; the Euler
equation gives $b_3=0$. Integral duality and the universal
coefficient theorem now give
\[
H_5=H^1=0,\qquad
H_4=H^2=\operatorname{Ext}(H_1,\mathbb Z)=\mathbb Z/|n|,
\qquad
H_3=H^3=\operatorname{Ext}(H_2,\mathbb Z)=\mathbb Z/|n|.
\]
Here the corresponding $\operatorname{Hom}$ terms vanish by the
computed ranks. The complex orientation gives $H_6=\mathbb Z$,
and connectedness gives $H_0=\mathbb Z$.
The torsion identifications can be specified by the characters
$\ell\mapsto1/|n|$ and $s\mapsto1/|n|$ in $\mathbb Q/\mathbb Z$:
their integral Bocksteins and Poincar\'e duals generate $H_4$ and
$H_3$, respectively.

For $n=\pm1$, Hurewicz and Whitehead identify $X$ with a homotopy
six-sphere. The dimension-six entry of
\cite[table on p.~504 and first remark on p.~505]{manKM63}
then gives a diffeomorphism with $S^6$.
Taking $n=1,r=0$ for each $k>0$ gives the unbounded multiplicities
$4k$ in this configuration and answers
\cite[Question~4, p.~2]{manViaclovskyLogs26} affirmatively.
\end{proof}

\subsection{Finite multiplicities in Yang's two families}

\begin{definition}[Marked Type B and Type C data]
\hypertarget{tag-Q113}{}\label{mapsIII:def:yang-finite-data}
{\normalfont\scriptsize tag Q113}\par
Fix positive integers $R,k,l$. Use Yang's Type B or Type C period
family with cusp index $R$ and its canonical cusp collar, and put
\[
 \begin{array}{c|ccc|cc}
 \text{type}&c&d&a=d/c&m_1&m_2\\\hline
 B&2&4&2&4k&6l\\
 C&3&3&1&3k&6l
 \end{array}.
\]
In the marking $\Lambda=\mathbb Z\langle\gamma,u,w,\delta\rangle$,
set $\ell_1=\gamma-R(u+w)$. 
In the stated marking the finite deck
matrices and forward transports are
\[
 A_1=\begin{pmatrix}
 1&0&0&0\\0&0&1&0\\-cR&-1&2-c&0\\-R&0&-1&1
 \end{pmatrix},\qquad
 A_2=\begin{pmatrix}
 1&0&0&0\\0&1&-1&0\\0&1&0&0\\0&0&0&1
 \end{pmatrix},\qquad P_j=A_j^{-1}.
\]
\end{definition}

\begin{theorem}[Finite multiplicities in the two Yang families]
\hypertarget{tag-02KK}{}\label{maps:src:D.2.10}
{\normalfont\scriptsize tag 02KK}\par
For the marked data of Definition~\ref{mapsIII:def:yang-finite-data},
choose integers $n_j,r_j$ with $\gcd(n_j,m_j)=1$ and take
the finite translation vectors
\[
 v_1=n_1\ell_1+r_1\delta,\qquad v_2=n_2\gamma+r_2\delta.
\]
These data define a compact smooth complex threefold $X$. Write
\[
 P=m_2n_1+m_1n_2,\qquad h=\gcd(R,ak,P)>0.
\]
For $P\ne0$, put $G=\mathbb Z/h\oplus\mathbb Z/(R|P|/h)$,
with trivial factors omitted. Then $\pi_1(X)=\mathbb Z/|P|$ and
\[
 \begin{array}{c|ccccccc}
 i&0&1&2&3&4&5&6\\\hline
 H_i(X;\mathbb Z)&\mathbb Z&\mathbb Z/|P|&
 \mathbb Z^{R-1}\oplus G&G&
 \mathbb Z^{R-1}\oplus\mathbb Z/|P|&0&\mathbb Z
 \end{array}.
\]
For $P=0$, one has $m_1=m_2$, $n_2=-n_1$,
$h=\gcd(R,ak)$, $\pi_1(X)=\mathbb Z$, and
\[
 \begin{array}{c|ccccccc}
 i&0&1&2&3&4&5&6\\\hline
 H_i(X;\mathbb Z)&\mathbb Z&\mathbb Z&
 \mathbb Z^{R+1}\oplus\mathbb Z/h&
 \mathbb Z^2\oplus\mathbb Z/h&\mathbb Z^{R+1}&
 \mathbb Z&\mathbb Z
 \end{array}.
\]
In both cases $\chi(X)=2R$, and these groups are independent of
$r_1,r_2$. Every free affine choice of the two fillings with these
multiplicities, including choices beyond the stated coprimality
condition, has a surjection
\[
 \pi_1(X)\longrightarrow\mathbb Z/\gcd(m_1,m_2).
\]
The quotient has order divisible by two in Type B and by three
in Type C, so none of these threefolds is simply connected.

For $k=l=1$, $(n_1,n_2)=(1,-1)$ and $r_1=r_2=0$, the construction
is Yang's original filling. Its groups are those of
\cite[Theorem~7.13, p.~50]{manYang26}: Type B has fundamental
group $\mathbb Z/2$ and middle torsion
$\mathbb Z/R\oplus\mathbb Z/2$; Type C has fundamental group
$\mathbb Z/3$ and middle torsion $\mathbb Z/(3R)$.
\end{theorem}

\begin{proof}
We use clockwise meridians. The matrices $A_1,A_2$ of
Definition~\ref{mapsIII:def:yang-finite-data} have orders $d,6$, and their fixed lattices are
$\mathbb Z\langle\ell_1,\delta\rangle$ and
$\mathbb Z\langle\gamma,\delta\rangle$.
Pull back the two good-reduction root discs by $\sigma=s^k$
and $\sigma=s^l$, respectively. The new coarse coordinates are
$t=s^{m_j}$. The period equivariance and logarithmic collars of
\cite[Lemmas~4.1 and 4.3, pp.~18--20]{manYang26} give the actions
\[
 (s,x)\longmapsto
 (e^{-2\pi i/m_j}s,A_jx+v_j/m_j).
\]
Their $m_j$th powers are integral translations. On the central
torus, applying the invariant character $\gamma^*$ to a fixed
point equation for the $q$th power gives $qn_j/m_j\in\mathbb Z$.
Coprimality excludes $0<q<m_j$. Off the central torus a
nontrivial power moves the root coordinate. Thus the actions
are free and their quotients are smooth, with central
multiplicities $m_j$.

The two inverse collar maps add and subtract
$(\log s)\Pi(s)v_j/(2\pi i)$. Changing the logarithm adds an
integral period; substituting the root rotation conjugates the
linear and affine actions. The three proper local fillings
therefore glue to the unchanged regular family over the compact
base sphere, giving the asserted compact complex threefold.

The finite boundary relation is $g_j^{m_j}=v_j$, where $g_j$
is the transferred clockwise meridian
\cite[Proposition~4.4, p.~20]{manYang26}.
The cusp kills $w,\delta$ and $g_1g_2$. The conjugation relation
$A_1w=u+(2-c)w-\delta$ then kills $u$.
The remaining fibre generator $\gamma$ is central, so van Kampen
gives the full presentation
\[
 \langle\gamma,g\mid [\gamma,g]=1,\quad
 m_1g=n_1\gamma,\quad -m_2g=n_2\gamma\rangle.
\]
Choose $m_1\alpha+n_1\beta=1$ and set $z=\alpha\gamma+\beta g$.
The first relation gives $\gamma=m_1z$, $g=n_1z$, with inverse
$z=\alpha\gamma+\beta g$; the second becomes $Pz=0$.
These formulas identify the group with $\mathbb Z/|P|$ for
$P\ne0$, and with $\mathbb Z$ for $P=0$.
For arbitrary free affine fillings, send every fibre loop to zero,
$g_1$ to $1$ and $g_2$ to $-1$ modulo $\gcd(m_1,m_2)$.
Every conjugation, cusp and finite power relation is preserved,
and $1$ is in the image. This proves the stated surjection.

Let $X^\circ$ be the union of the regular region and the two finite
fillings before the cusp is attached. Its two sides retract onto
the reduced surfaces $S_1,S_2$, and their intersection retracts
onto a regular torus $F$. The linear coinvariants give
$w=-u$, $\delta=cu$ on the first side, and $u=w=0$ on the second.
Thus
\[
 \begin{split}
 H_1(S_1)&=\langle\gamma,u,g_1\mid
       m_1g_1=n_1\gamma+cr_1u\rangle,\\
 H_1(S_2)&=\langle\gamma,\delta,g_2\mid
       m_2g_2=n_2\gamma+r_2\delta\rangle.
 \end{split}
\]
For $m_j\alpha_j+n_j\beta_j=1$, put
$z_j=\alpha_j\gamma+\beta_jg_j$.
On the first side the inverse formulas are
\[
 \gamma=m_1z_1-cr_1\beta_1u,\qquad
 g_1=n_1z_1+cr_1\alpha_1u;
\]
on the second they are
$\gamma=m_2z_2-r_2\beta_2\delta$ and
$g_2=n_2z_2+r_2\alpha_2\delta$.
Hence these groups are free of rank two. Each $S_j$ is an oriented
finite quotient of a torus and has Euler characteristic zero.
Duality and the universal coefficient theorem give
$H^2(S_j)=\mathbb Z^2$ with unimodular intersection pairing.

In the exterior cohomology basis
\[
 (\gamma^*u^*,\gamma^*w^*,\gamma^*\delta^*,
           u^*w^*,u^*\delta^*,w^*\delta^*),
\]
the first invariant lattice has basis
\[
 \omega_1=-R\gamma^*u^*+R\gamma^*w^*+u^*w^*,\qquad
 \omega_2=\gamma^*u^*-\gamma^*w^*+c\gamma^*\delta^*.
\]
The two squares vanish and $\omega_1\omega_2=c\,\mathrm{vol}_F$.
Indeed the fixed equations for coordinates $(x,y,z,t,r,s)$ are
$r=s=0$, $y=-x$ and $z=c(x+Rt)$; their inverse basis
coordinates are $(t,x+Rt)$.
The second invariant lattice is generated by
$u^*w^*,\gamma^*\delta^*$, with unimodular pairing: its fixed
equations are $x=y=r=s=0$.
For the degree-$m_j$ covering $p_j:F\to S_j$, transfer makes
$p_j^*$ injective, and its intersection pairing is multiplied
by $m_j$. Taking determinants gives
\[
 [H^2(F)^{A_1}:p_1^*H^2(S_1)]=m_1/c=ak,\qquad
 [H^2(F)^{A_2}:p_2^*H^2(S_2)]=m_2.
\]
The square $(q,t)\mapsto[q\delta+t v_j/m_j]$ defines a torus in
$S_j$. Evaluation of a pulled-back integral class on it is
$\alpha_{\rm inv}(\delta,v_j)/m_j\in\mathbb Z$.
For $\alpha_{\rm inv}=A\omega_1+B\omega_2$ this is $-cn_1B/m_1$,
so $B$ is a multiple of $ak$. At the second surface the coefficient
of $\gamma^*\delta^*$ is a multiple of $m_2$.
The resulting sublattices have the indices just computed, proving
that their basis matrices are exactly
\[
 \mathsf B_1=\begin{pmatrix}
 -R&ak\\R&-ak\\0&m_1\\1&0\\0&0\\0&0
 \end{pmatrix},\qquad
 \mathsf B_2=\begin{pmatrix}
 0&0\\0&0\\0&-m_2\\1&0\\0&0\\0&0
 \end{pmatrix}.
\]
In the dual homology bases the finite $\delta$ sweeps have
coordinates $(0,-n_1)$ and $(0,n_2)$. Their coordinates are used
for every admissible $n_j$.

With the Mayer--Vietoris sign $(i_{1*},-i_{2*})$, the map on
degree-two fibre classes is
\[
 \alpha_{\rm MV}=\begin{pmatrix}
 -R&R&0&1&0&0\\ak&-ak&m_1&0&0&0\\
 0&0&0&-1&0&0\\0&0&m_2&0&0&0
 \end{pmatrix}.
\]
The degree-one kernel is $\mathbb Z(u+w)$: vanishing on the
second side forces the $\gamma,\delta$ coefficients to vanish;
vanishing on the first then equates the $u,w$ coefficients.
The finite power relation creates no further kernel element from
$\Lambda$, since its $g_j$ coefficient is $m_j\ne0$.

Write $N$ for the cusp tube and $M=\partial N$.
Its integral cohomology, Euler number and full boundary kernel are
\cite[Theorem~7.8, pp.~44--45, and Lemma~7.10, pp.~48--49]{manYang26}
\[
 H^3(N)=\mathbb Z^2,\quad H^4(N)=\mathbb Z^R,\quad
 \chi(N)=2R,\qquad
 K_2=\ker(H_2(M)\to H_2(N))
       =\mathbb Z\xi_w\oplus\mathbb Z\xi_\delta.
\]
Here $\xi_\nu$ is the circle-first angular sweep of $\nu$ over
the clockwise cusp meridian; its bounding chain is the angular
circle action applied to the extended zero-section disc.
These local cusp data are unchanged by the finite translations.

Cutting the $w$ sweep along the two finite edges gives connecting
class $-(u+w)$: the forward transports send $w$ first to $-u$
and then back to $w$. It is a primitive generator of the
degree-one Mayer--Vietoris kernel. The $\delta$ sweep is the
negative sum of the two finite sweeps, so its negative image
has column $b=(0,-n_1,0,n_2)^{\mathsf T}$. Consequently
\[
 H_2(X^\circ)/j_*K_2=\operatorname{coker}[\alpha_{\rm MV}\mid b].
\]
For clarity, the elimination of its unit relation is integral.
On target coordinates use
$(a,b,c,d)\mapsto(a+c,b,d,a)$; its determinant is a unit.
On source coordinates $e_0,\ldots,e_6$, use the ordered basis
\[
 e_1-Re_3,\ e_2,\ e_6,\ e_3,\ e_0+e_1,\ e_4,\ e_5.
\]
The resulting matrix is the direct sum of the following $C$,
a unit, and three zero columns:
\[
 C=\begin{pmatrix}R&0&0\\-ak&m_1&-n_1\\0&m_2&n_2\end{pmatrix},
 \qquad \det C=RP.
\]
Both basis changes have integral inverses, and the unit coordinate
vanishes in the quotient. The entry gcd of $C$ is one. Its
nonzero two-by-two minors, up to signs, are
\[
 Rm_1,\ Rn_1,\ Rm_2,\ Rn_2,\ akm_2,\ akn_2,\ P.
\]
Their ideal is $(R,ak,P)$: B\'ezout combinations of $Rm_1,Rn_1$
give $R$, and those of $akm_2,akn_2$ give $ak$; every listed
minor belongs to that ideal. Thus their gcd is $h$.
Integral Euclidean row and column operations carry $C$ to
$\operatorname{diag}(1,h,R|P|/h)$ for $P\ne0$ and to
$\operatorname{diag}(1,h,0)$ for $P=0$.
These operations give quotient isomorphisms explicitly: keep
the left and right products $U,V$ of the elementary operations,
so $UCV$ is diagonal, and send $[x]$ to $[Ux]$, with inverse
$[y]\mapsto[U^{-1}y]$.
At each stage move a nonzero entry to the pivot position and use
integer division to clear its row and column; a nonzero remainder
strictly decreases the pivot. A pivot failing to divide a
remaining entry is combined with that entry and reduced again.
This terminates by positive integer descent. Entry ideals,
minor ideals and determinant give the displayed diagonal entries.

Excision and duality identify $H_2(N,M)$ with
$H^4(N)=\mathbb Z^R$. Clockwise normal discs at smooth points
of the $R$ components give a basis. In the branch resolution,
the component fundamental cycles form the full basis of $H_4(N)$;
$H_3(N)$ is a subgroup of the free double-curve group and is free.
Thus $H^4(N)$ is the integral dual of $H_4(N)$, and the discs
give its negative dual basis in the clockwise convention.
A component with ray
$(a,b,1)$ has meridian $a_0-aw-b\delta$, so every disc maps
to $a_0$ in $H_1(X^\circ)$. In that group $u,w,\delta$ have
already vanished, and
\[
 H_1(X^\circ)=\langle\gamma,g_1,g_2\mid
 m_1g_1=n_1\gamma,\ m_2g_2=n_2\gamma\rangle,
 \qquad a_0=-g_1-g_2.
\]
The relation matrix has entry gcd one and maximal-minor gcd
$d_0=\gcd(m_1,m_2)$, giving $H_1(X^\circ)=\mathbb Z\oplus\mathbb Z/d_0$
by the same integral operations. Its primitive free quotient is
\[
 \Phi=\frac1{d_0}(m_1m_2,\ m_2n_1,\ m_1n_2).
\]
The coefficient gcd is one by the two coprimality conditions,
and $\Phi(a_0)=-P/d_0$.
For $P\ne0$, $a_0$ has infinite order, and the kernel of the
relative boundary consists of the vectors with coordinate sum
zero. It has basis $q_j-q_0$, $1\leq j<R$.

For $P=0$, the reduced fractions $n_1/m_1=-n_2/m_2$ have the
same positive denominator: each denominator divides the other
by coprimality. Thus $m_1=m_2=m$, $n_2=-n_1$.
In the generators $\gamma,g_1,a_0$, the relations are
$mg_1=n_1\gamma$ and $ma_0=0$. The first is primitive,
so $a_0$ has exact order $m$. The relative kernel is
\[
 \{(y_0,\ldots,y_{R-1}):\textstyle\sum y_j\in m\mathbb Z\},
\]
with basis $mq_0,q_j-q_0$ for $1\leq j<R$.
Its rank is $R$ and its index is $m$.

The exact sequence of the attachment is
\[
 0\longrightarrow\operatorname{coker}C
 \longrightarrow H_2(X)
 \longrightarrow\ker\bigl(H_2(N,M)\to H_1(X^\circ)\bigr)
 \longrightarrow0.
\]
A splitting is obtained by representing each displayed kernel
basis vector by its normal discs, filling their null-homologous
boundary by an integral two-chain in $X^\circ$, and subtracting
that chain. Sending the basis to these closed cycles defines a
section; subtraction of that section gives the inverse direct-sum
coordinate map. This proves the degree-two assertions.

Every other gluing piece has Euler characteristic zero, so
$\chi(X)=\chi(N)=2R$. Duality gives $b_4=b_2$, $b_5=b_1$;
the Euler characteristic gives $b_3=0$ for $P\ne0$ and $b_3=2$
for $P=0$. The universal coefficient theorem and integral
duality give $\operatorname{Tor}H_3\simeq\operatorname{Tor}H_2$
and $\operatorname{Tor}H_4\simeq\operatorname{Tor}H_1$.
Concretely, a torsion character with values in $\mathbb Q/\mathbb Z$
maps by the Bockstein to integral cohomology and then by cap product
with $[X]$ to the paired torsion group. These maps, together with
the free ranks, give every entry of the tables.

At $k=l=1$, $n_1=1$, $n_2=-1$, $r_1=r_2=0$, the root
pullbacks and collar maps are Yang's original ones.
For Type B, $P=2$ and $h=\gcd(R,2)$. For even $R$,
interchanging the two cyclic coordinates identifies
$\mathbb Z/2\oplus\mathbb Z/R$ with the stated group, and
the same interchange is its inverse. For odd $R$, the map
$[x]_{2R}\mapsto([x]_R,[x]_2)$ has inverse
$([a]_R,[b]_2)\mapsto[2u_{\rm crt}a+Rv_{\rm crt}b]_{2R}$,
where $2u_{\rm crt}+Rv_{\rm crt}=1$.
Both composite maps are the identity modulo their indicated orders.
For Type C, $P=3$ and $h=1$, giving $\mathbb Z/(3R)$.
The identification of the local pieces and their collars proves
that these specializations have the stated original groups.
\end{proof}

\subsection{A component-supported zero on the \texorpdfstring{$IV^*+II+I_2$}{IV*+II+I2} profile}

\begin{definition}[A component-supported elliptic datum]
\hypertarget{tag-Q114}{}\label{mapsIII:def:os49-component-zero}
{\normalfont\scriptsize tag Q114}\par
Let $f:S\to\mathbb P^1$ be the minimal regular model of
\begin{equation}
 y^2=x^3-3t^3(t-1)x+t^4(t-1)(2t-1),
 \qquad P=(t^2,t^2).
 \hypertarget{tag-01Y6}{}\label{maps:eq:D.C1}
\end{equation}
Write $O$ for its zero section, put $M=\mathcal O(P-O)$ and
$Q=3P$, and let $C_0$ be the component met by $O$ in the fibre
over infinity.
\end{definition}

\Needspace{6\baselineskip}
\begin{theorem}[Logarithmic multiplicities with a component-supported zero]\hypertarget{tag-02KL}{}\label{maps:src:D.2.11}
{\normalfont\scriptsize tag 02KL}\par


The surface of Definition~\ref{mapsIII:def:os49-component-zero}
has fibres $IV^*,II,I_2$ at $0,1,\infty$ and Mordell--Weil group
$\mathbb Z$ generated by $P$. There is a linearization with divisor $C_0$, and its torus
family has Yang's index-one Type C marking. Keep the type-$II$
filling unchanged and use the index-one toroidal filling at
$I_2$. At $IV^*$ choose any logarithmic order and invariant vector
\[
 m=3k,\qquad v=n(\gamma-u-w)+r\delta,
 \qquad k\in\mathbb Z_{>0},\quad n,r\in\mathbb Z,
 \quad\gcd(n,3k)=1.
\]
The resulting compact smooth complex threefold has
\[
 \pi_1(X)=\mathbb Z/|n|,\qquad
 H_i(X;\mathbb Z)=
 \begin{cases}
 \mathbb Z&i=0,6,\\
 \mathbb Z/|n|&1\leq i\leq4,\\
 0&i=5.
 \end{cases}
\]
Consequently $n=\pm1$ gives the standard $S^6$ for every positive
$k$. This is an affirmative unbounded configuration for
\cite[Question~4, p.~2]{manViaclovskyLogs26}.
\end{theorem}


\begin{proof}
\emph{The elliptic surface and component-supported linearization.}
For the Weierstrass coefficients $a_4,a_6$ of
\eqref{maps:eq:D.C1}, direct calculation gives
\[
 4a_4^3+27a_6^2=27t^8(t-1)^2,\qquad
 j=6912t(1-t).
\]
The minimal valuations give $IV^*$ at zero and $II$ at one.
In the chart at infinity the discriminant has order two and
the cubic has a node, giving $I_2$. These fibres have Euler
sum twelve; the degree-four and degree-six Weierstrass
coefficients give the rational elliptic surface. Substitution
checks the section $P$, which is disjoint from $O$ and meets
the nonidentity components at $IV^*$ and $I_2$.
Its height is $2-4/3-1/2=1/6$. The Shioda--Tate formula gives
rank $10-2-6-1=1$. The N\'eron--Severi lattice of this rational
surface has absolute discriminant one, while the trivial lattice
has absolute discriminant $3\cdot2=6$. Let $G$ generate the free
Mordell--Weil quotient, let $t_0$ be the order of its torsion group,
and write $P=sG+T$ with $T$ torsion. The discriminant formula gives
\[
 h(G)=t_0^2/6,\qquad 1/6=h(P)=s^2t_0^2/6.
\]
Thus $|s|=t_0=1$, so $P$ generates the whole Mordell--Weil group.
The rank, unimodularity, height and discriminant identities used here
are \cite[Corollary~6.13, p.~29; \S8.8, p.~41; \S11.8 and Table~4, pp.~51--52; \S11.10, equation~(22), p.~52]{manSS10}.

The section $Q=3P$ meets the identity component of $IV^*$
and the nonidentity component of $I_2$. Its height is $3/2$,
so the height formula gives $Q\cdot O=0$.
The mixed height $\langle P,Q\rangle=1/2$, with the sole
local contribution $1/2$ at $I_2$, gives $P\cdot Q=0$.
Let $J=M\otimes t_Q^*M^{-1}$. Its multidegree vanishes at
$IV^*$ and $II$. At $I_2$, translation by $Q$ exchanges the
two components, so its multidegrees are $(-2,2)$ on
$(C_0,C_1)$, exactly those of $\mathcal O(C_0)$.
Their difference is thus a pullback from the base. Restriction
to $O$ gives
\[
 J\cdot O=Q\cdot O-P\cdot Q+1=1=C_0\cdot O,
\]
so the base degree of that difference is zero. Therefore
$J\simeq\mathcal O(C_0)$. The line isomorphism can be given explicitly.
Put $D=P-O-(P-Q)+(-Q)$. On the generic cubic, the line through
$P,-Q$ divided by the vertical line through $P-Q$ has divisor $D$.
The resulting rational function $F$ on $S$ has
$\operatorname{div}(F)=D+V$ with $V$ vertical. The multidegrees
above imply $V+C_0=\sum a_vF_v$, and the intersection with $O$
gives $\sum a_v=0$. Choose a rational base function $B$ with
divisor $\sum a_v[v]$. Then
\[
 \rho=F/(B\circ f),\qquad \operatorname{div}(\rho)=D-C_0,
 \qquad
 \mathcal O(D)\xrightarrow{\ \rho\ }\mathcal O(C_0),
 \qquad
 \mathcal O(C_0)\xrightarrow{\ \rho^{-1}\ }\mathcal O(D)
\]
are inverse line maps. Pulling back the canonical section of
$\mathcal O(C_0)$ supplies the effective linearization.
Its norm under the component exchange has divisor
$C_0+C_1=F_\infty$. This explains why a half full-fibre weight
appears in period coordinates while the geometric divisor is
integral.

\emph{The Type C period marking.}
The elliptic $j$-function is the one in
\cite[Proposition~3.1, p.~11]{manYang26}. Over $t=s^3$ the good
model has central curve $Y^2=X^3+1$, and $P$ specializes to
$(0,1)$, a nonzero fixed point of order three.
Over $t-1=s^6$, the same section specializes to the origin
of the good elliptic model. At $I_2$ it meets the nonidentity
component. Choose the sign of the elliptic marking so that its
Abel--Jacobi coordinate at the first point is $(\tau+1)/3$.
Its normalized lift therefore obeys the two Type C affine
laws and has $\mu-\tau/2$ regular at infinity.
The uniqueness proof of \cite[Proposition~3.4, p.~13]{manYang26},
using $H^0(\mathcal O(-1))=H^1(\mathcal O(-1))=0$, identifies
it with $\mu_C$.
The extension line $M$ has this same Picard coordinate,
whereas the translation has coordinate $3\mu_C$.
The norm divisor just calculated fixes the cusp central
period. The remaining scalar period is the affine constant
of \cite[Proposition~3.6, p.~14]{manYang26}. Consequently the open
two-torus family is the index-one Type C family, for a suitable
admissible constant, and its primitive toroidal cusp has one
normalized $dP_6$ component
\cite[Theorem~4.8, pp.~22--23; Corollary~4.13, p.~25]{manYang26}.

The period marking has the following integral comparison.
In the order $(II,IV^*,I_2)$, the homology matrices are
\[
 \begin{split}
 M_{II}&=\begin{pmatrix}1&1&0&0\\-1&0&0&0\\0&0&1&0\\0&0&0&1\end{pmatrix},\\
 M_{IV^*}&=\begin{pmatrix}0&-1&0&0\\1&-1&0&-3\\0&1&1&1\\0&0&0&1\end{pmatrix},\\
 M_{I_2}&=\begin{pmatrix}1&2&0&3\\0&1&0&0\\0&-1&1&-1\\0&0&0&1\end{pmatrix}.
 \end{split}
\]
The first two follow from the locally narrow finite-point
section cocycles $(0,0)$ and $(0,-1)$ for $P$, multiplied by
three for $Q$. Their product gives the displayed $I_2$ matrix
for the component divisor $C_0$.
Put ${\mathsf U}=M_{IV^*}$, ${\mathsf V}={\mathsf U}^{-1}M_{II}{\mathsf U}$ and
\[
 {\mathsf S}=\begin{pmatrix}
 0&0&-1&0\\0&-1&1&1\\0&0&0&1\\-1&0&3&0
 \end{pmatrix}.
\]
For the Type C matrices of Definition~\ref{mapsIII:def:yang-finite-data}
at $R=1$, put $T_j=A_j^{-t}$. Direct multiplication gives
$\det {\mathsf S}=-1$, ${\mathsf U}^{-t}{\mathsf S}={\mathsf S}T_1$ and ${\mathsf V}^{-t}{\mathsf S}={\mathsf S}T_2$. Its inverse is
\[
 {\mathsf S}^{-1}=\begin{pmatrix}
 -3&0&0&-1\\-1&-1&1&0\\-1&0&0&0\\0&0&1&0
 \end{pmatrix}.
\]
Thus ${\mathsf S}^{-t}$ sends the Type C homology coordinates to the displayed
period coordinates, with inverse ${\mathsf S}^t$, and
${\mathsf U}{\mathsf S}^{-t}={\mathsf S}^{-t}A_1$, ${\mathsf V}{\mathsf S}^{-t}={\mathsf S}^{-t}A_2$.
The preceding Hurwitz reorder sends the pair $(M_{II},M_{IV^*})$
to $({\mathsf U},{\mathsf U}^{-1}M_{II}{\mathsf U})$; its inverse sends $({\mathsf U},{\mathsf V})$ to
$({\mathsf U}{\mathsf V}{\mathsf U}^{-1},{\mathsf U})$. These maps give the ordered integral comparison.

\emph{The product type-$II$ filling and its integral marking.}
Here the product proof of \cite[Lemma~3.10, pp.~11--12]{manLR26} applies
with translation section $3P$, as the following explicit product
construction shows. A closed elliptic neighborhood $E$ of the
type-$II$ fibre retracts to its cuspidal cubic, homeomorphic to
$S^2$. The restriction of $M$ has degree zero and is therefore
differentiably trivial. In a smooth trivialization the lifted
contracting generator has form
\[
 (x,z)\longmapsto(t_{-Q}(x),\lambda(x)z),\qquad |\lambda|<1.
\]
As $E$ is simply connected, $\lambda$ has a smooth logarithm.
The relative exponential of the additive N\'eron group is an
isomorphism on the central smooth locus. After shrinking the
disc it provides a holomorphic logarithm of $-Q$ and hence a
fibre-preserving isotopy of $t_{-Q}$ to the identity.
The translations extend over the cuspidal cubic and fix its cusp.
Write $L=\log\lambda$, and let $q$ be the logarithm of $-Q$.
On $E\times S^1$, the return map of a fundamental annulus is
$h(x,\theta)=(t_{-Q}x,\theta+\operatorname{Im}L(x))$.
The isotopy
\[
 h_t(x,\theta)=
 (t_{\exp(tq)}x,\theta+t\operatorname{Im}L(x)),\qquad 0\le t\le1,
\]
has inverse
$(y,\phi)\mapsto(t_{-\exp(tq)}y,
\phi-t\operatorname{Im}L(t_{-\exp(tq)}y))$.
Here $\exp$ is the relative exponential and its group inverse is
written additively. The mapping-torus identification
$(z,1)\sim(h(z),0)$ gives the product map
\[
 [z,t]\longmapsto(h_t(z),e^{2\pi it}),\qquad
 (w,e^{2\pi it})\longmapsto[h_t^{-1}(w),t].
\]
The endpoint identities make these mutually inverse maps well defined.
Thus the filling is homeomorphic over its base disc to $E\times T^2$,
with angular circle $\delta$ and a primitive second circle of
$\gamma^*$-value one.

Its regular specialization sends the elliptic area $uw$ to
the fundamental class of the cuspidal cubic and $\gamma\delta$
to the area of this second torus; it kills the mixed classes.
Equivalently the exact degree-two pullback lattice is
$\langle u^*w^*,\gamma^*\delta^*\rangle$.
Both transverse normal tori bound products of a section disc
with a circle. In particular its $\delta$ meridian sweep is
zero. These assertions also follow by applying the marked
specialization proof of \cite[Proposition~3.11, p.~12]{manLR26}
to the product just constructed.

\emph{The logarithmic quotient and integral homology.}
At $IV^*$ the order-$3k$ free quotient and its marked collar
are those of \cref{maps:src:D.2.10} with $R=1$, $n_1=n$, $r_1=r$.
The two other compact fillings have just been specified,
so the same finite-core gluing proves compactness and
smoothness. The unchanged type-$II$ piece kills $u,w$ and
its meridian; the cusp kills $\delta$ and the remaining
meridian. Van Kampen leaves exactly $n\gamma=0$, proving
$\pi_1(X)=\mathbb Z/|n|$. The quotient map sends
$\gamma$ to $[1]_n$, and its inverse sends $[1]_n$ to
$[\gamma]$ in this presentation.

The first finite pullback matrix is unchanged. For the second
piece the primitive product lattice replaces its coefficient
$m_2$ by one, and its vanishing $\delta$ sweep replaces $n_2$
by zero. Thus the complete augmented integral map is
\[
 {\mathsf D}=\begin{pmatrix}
 -1&1&0&1&0&0&0\\
 k&-k&3k&0&0&0&-n\\
 0&0&0&-1&0&0&0\\
 0&0&1&0&0&0&0
 \end{pmatrix}.
\]
The primitive $w$ cusp sweep still kills the degree-one
Mayer--Vietoris kernel. Unit elimination reduces ${\mathsf D}$ to
one unit factor and
\[
 \begin{pmatrix}1&0&0\\-k&3k&-n\\0&1&0\end{pmatrix}.
\]
The preceding target/source maps for the augmented matrix specialize
at $R=1,m_1=3k,m_2=1,n_1=n,n_2=0$ to this unit elimination.
For the displayed three-by-three block $C$, the integral matrix
\[
 W=\begin{pmatrix}1&0&0\\0&0&1\\k&1&-3k\end{pmatrix},
 \qquad WC=\operatorname{diag}(1,1,-n),\qquad\det W=-1
\]
gives the cyclic quotient map
$[(a,b,c)]\mapsto[ka+b-3kc]_n$, with inverse
$[1]_n\mapsto[(0,1,0)]$. When $ka+b-3kc=nq$, the vector
$(a,c,-q)$ is its preimage under $C$. These formulas prove the
Smith factors $(1,1,|n|)$ for every allowed integer datum;
the displayed maps contain no $r$.
Before attaching the cusp, the remaining degree-one group is
$\langle\gamma,g\mid3kg=n\gamma\rangle=\mathbb Z$.
Its cusp meridian is $-g$, a nonzero multiple $-n$ of a
primitive generator. For the cusp tube $N$, the relative map
$H_2(N,\partial N)=\mathbb Z\to H_1(X^\circ)=\mathbb Z$ is injective,
so it contributes no extra free degree-two class.
We conclude $H_2(X)=\mathbb Z/|n|$.
The cusp has Euler number two and the other pieces have
Euler number zero. Duality and the universal coefficient
theorem give the remaining groups. When $|n|=1$, sphere
recognition applies exactly as in \cref{maps:src:D.2.6}.

\end{proof}

\section{Unbounded multiplicities and logarithmic modifications}
\subsection{Seven elliptic profiles with unbounded logarithmic multiplicity}

\begin{definition}[Seven marked elliptic profiles]
\hypertarget{tag-Q115}{}\label{mapsIII:def:seven-log-profiles}
{\normalfont\scriptsize tag Q115}\par
Let $G$ be the marked Mordell--Weil generator of the indicated
rational elliptic surface. In the table, the superscript $\bullet$
marks the fibre supporting the degree-one linearization, and $d$
is the linear order at the sole additive fibre.
\[
 \begin{array}{c|l|rr|r}
 \mathrm{OS}&\text{ordered fibre configuration}&P/G&Q/G&d\\\hline
 45&I_8,II,I_1^\bullet,I_1&8&1&6\\
 47&I_7,I_2,I_1^\bullet,II&14&1&6\\
 47&I_7,I_1^\bullet,I_1,III&7&2&4\\
 55&I_5,I_4,II,I_1^\bullet&20&1&6\\
 56&I_5,I_2,IV,I_1^\bullet&10&3&3\\
 56&I_5,I_3,III,I_1^\bullet&15&2&4\\
 49&I_2,IV^*,I_1^\bullet,I_1&2&3&3
 \end{array}
\]
Use $M=\mathcal O(P-O)$ and a degree-one linearization
$\phi:t_Q^*M\to M$, whose forward action covers $-Q$.
Keep the two weight-zero multiplicative fillings and the
primitive toroidal filling at $I_1^\bullet$.
In the period basis $(e_1,e_2,\delta,\gamma)$, with $\delta$ the
positive line circle and $\gamma$ the forward linearization period,
use the following additive fixed vectors, in table order:
\[
 \begin{array}{c|c}
 1&(-1,0,0,1)^t\\
 2&(0,-1,0,1)^t\\
 3&(-3,-2,0,1)^t\\
 4&(0,0,0,1)^t\\
 5&(-1,3,0,1)^t\\
 6&(-1,3,0,1)^t\\
 7&(-3,-1,0,1)^t
 \end{array}\qquad (\text{entry }i=\ell_i).
\]
\end{definition}

\begin{theorem}[Seven families with unbounded logarithmic multiplicity]\hypertarget{tag-02KM}{}\label{maps:src:D.2.12}
{\normalfont\scriptsize tag 02KM}\par


For each profile of Definition~\ref{mapsIII:def:seven-log-profiles}, put
\[
 m=dk,\qquad v=n\ell+r\delta,
 \qquad k\in\mathbb Z_{>0},\quad\gcd(n,dk)=1,
 \qquad n,r\in\mathbb Z.
\]
Every row admits these logarithmic fillings for every displayed
integer choice, and the resulting compact smooth threefold has
\[
 \begin{gathered}
 \pi_1(X)=\mathbb Z/|n|,\qquad
 H_i(X;\mathbb Z)=\mathbb Z/|n|\quad(1\leq i\leq4),\\
 H_5(X)=0,\qquad H_0(X)=H_6(X)=\mathbb Z.
 \end{gathered}
\]
In particular $n=\pm1$ gives $S^6$ for every $k>0$ in all seven
rows. For an arbitrary admissible affine generator in the same additive
reduction model, let $w$ be its integral translation period after
$m$ iterates and define $n=\gamma^*(w)$. Simple connectivity
requires $n=\pm1$; hence the displayed coprimality domain contains
all sphere cases of this one-log construction.
\end{theorem}


\begin{proof}
Write $P=pG$, $Q=qG$, and orient a peripheral circle by the
complex orientation of the base. The lattice of a smooth quotient
fibre is $\Lambda=\mathbb Z\langle e_1,e_2,\delta,\gamma\rangle$,
where $\delta$ is the positive line-bundle circle and $\gamma$
is the forward linearization period. Put $\chi=\gamma^*$.

\emph{Geometric realization of the marked elliptic surfaces.}
The marked elliptic surfaces have Weierstrass equations
\[
y^2=4x^3-g_2(t)x-g_3(t),\qquad
\Delta=g_2^3-27g_3^2,\qquad J=g_2^3/\Delta.
\]
The following coefficient pairs use the order in Definition~\ref{mapsIII:def:seven-log-profiles}:
\[
\begin{array}{c|l}
1&g_2=12t(t^3-6t^2+15t-12)\\
 &g_3=4t(2t^5-18t^4+72t^3-144t^2+135t-27)\\
2&g_2=12t(9t^3+36t^2+42t+14)\\
 &g_3=12t(18t^5+108t^4+234t^3+222t^2+87t+8)\\
3&g_2=12t(t^3+4t^2+10t+6)\\
 &g_3=4t^2(2t^4+12t^3+42t^2+70t+63)\\
4&g_2=3(8t-1)(8t^3+87t^2+96t-64)\\
 &g_3=(8t-1)(64t^5+1040t^4+3925t^3+100t^2+3200t-512)\\
5&g_2=12t^2(t^2+8t+10)\\
 &g_3=4t^2(2t^4+24t^3+78t^2+56t+27)\\
6&g_2=75(5t-1)(5t^3+45t^2+39t-25)\\
 &g_3=25(5t-1)^2(25t^4+340t^3+1086t^2+100t+625)\\
7&g_2=3,\qquad g_3=1+t^2.
\end{array}
\]
The first six are the equations of
\cite[Table III, pp.~23--24]{manHerfurtner90}; the corresponding
author data and fibre positions are identified in
\cite[Table~1, p.~10; \S9, p.~17]{manBesserLivne12}. On the second base chart
the exact map is $s=t^{-1}$, $X=s^2x$, $Y=s^3y$, with inverse
$x=s^{-2}X$, $y=s^{-3}Y$. The minimal discriminant orders give the
stated fibres. These degree-$(4,6)$ Weierstrass models have rational
minimal resolutions by \cite[\S8.4, p.~39]{manSS10}; the map to
its monic convention is $(x,y)\mapsto(x,y/2)$, with inverse
$(x,y')\mapsto(x,2y')$.

For row~7, the discriminant is $-27t^2(t^2+2)$ and its orders
at infinity are $(4,4,8)$, giving $IV^*$ by
\cite[Table I, p.~6]{manHerfurtner90}. Rationality also follows
from the inverse maps
\[
\begin{gathered}{}
U=y+it,\quad V=y-it,\quad UV=(x-1)(2x+1)^2,
\\{}
\qquad V=\frac{(x-1)(2x+1)^2}{U},\quad
y=\frac{U+V}{2},\quad t=\frac{U-V}{2i}.
\end{gathered}
\]
The infinity equation is $Y^2=4X^3-s^4(1+3X+s^2)$.
The local change $Z=s(1+3X+s^2)^{1/4}$, $X_1=4^{1/3}X$
has an analytic inverse and gives $Y^2=X_1^3-Z^4$.
Since $Z/s$ is a unit, it preserves the central fibre divisor.
All seven models have the section $[0:1:0]$.

For the integral markings, at $j=1/2$
label the real roots of $J(t)=j$ increasingly, followed by the nonreal
conjugate pairs in increasing real part, the negative imaginary root
first in each pair. Positive loops about $0$ and $1$ follow respectively
the rectangles
\[
\begin{split}
&(1/2,1/2+i/2,-1/2+i/2,-1/2-i/2,1/2-i/2,1/2),\\
&(1/2,1/2-i/2,3/2-i/2,3/2+i/2,1/2+i/2,1/2).
\end{split}
\]
Their sheet permutations are given below, as image lists beginning
with label~0. The displayed $f$ identifies the subsequent ribbon
labels with these analytic labels; its inverse is the inverse
permutation.
\[
\begin{aligned}
\sigma_0^{(1)}&=(0,2,3,1,6,9,8,5,4,7),&\sigma_1^{(1)}&=(1,0,4,5,2,3,8,9,6,7),\\
f_{1}&=(0,1,2,3,4,5,6,7,8,9).\\
\sigma_0^{(2)}&=(5,6,9,3,0,4,7,1,2,8),&\sigma_1^{(2)}&=(1,0,3,2,5,4,8,9,6,7),\\
f_{2}&=(0,1,2,3,4,5,6,7,8,9).\\
\sigma_0^{(3)}&=(8,3,6,5,2,1,4,0,7),&\sigma_1^{(3)}&=(0,3,4,1,2,7,8,5,6),\\
f_{3}&=(0,1,2,3,4,5,6,7,8).\\
\sigma_0^{(4)}&=(5,6,2,8,0,4,7,1,9,3),&\sigma_1^{(4)}&=(1,0,3,2,5,4,8,9,6,7),\\
f_{4}&=(5,4,0,1,6,7,8,9,3,2).\\
\sigma_0^{(5)}&=(3,4,0,2,5,1,7,6),&\sigma_1^{(5)}&=(1,0,3,2,6,7,4,5),\\
f_{5}&=(3,2,0,1,4,5,6,7).\\
\sigma_0^{(6)}&=(4,5,7,0,3,6,1,8,2),&\sigma_1^{(6)}&=(1,0,2,4,3,7,8,5,6),\\
f_{6}&=(2,7,8,5,6,1,0,4,3).\\
\sigma_0^{(7)}&=(1,3,0,2),&\sigma_1^{(7)}&=(2,3,0,1),\\
f_{7}&=(0,1,2,3).\\
\end{aligned}
\]

The following rational inequalities prove these finite continuations.
Write $J=N_J/D_J$ and $P_j(t)=N_J(t)-jD_J(t)$. Expand
$P_{j_0}(c+z)=\sum a_{\nu}z^\nu$ and $D_J(c+z)=\sum d_{\nu}z^\nu$
at a rational complex centre. A disk of radius $r$ contains one
root throughout a straight segment whenever
\[
|a_1|_-r-|a_0|_+-\sum_{\nu\geq2}|a_\nu|_+r^\nu
 > |j_1-j_0|_+\sum_\nu|d_\nu|_+r^\nu,
\]
where $|u+iv|_- =\max(|u|,|v|)$ and
$|u+iv|_+=|u|+|v|$. This is Rouch\'e's inequality with rational
bounds. Choose disjoint dyadic-radius disks, halve each segment
until this strict inequality holds, and enclose its endpoint roots
in disks of radius $10^{-45}$ contained in the preceding disks.
Applied to the displayed polynomials and rectangles, these rational
inequalities give the sheet lists above; the final original root
disk is contained in the last continuation disk. Thus each entry
specifies the actual analytic continuation, including its endpoint.
The covering isomorphism sends a sheet and a path to the unique
analytic lift beginning in the root disk indexed by $f$; its inverse
uses $f^{-1}$. At a branch point both maps extend by $z\mapsto z^e$.

Use the positively oriented projective period matrices $A_{\rm per},B_{\rm per}$ of
Remark~\ref{mapsIII:rem:positive-projective-periods}. Its period
isomorphism and inverse identify the projective elliptic local system
with the pullback by $J$, with the finite and cusp orientations fixed
there.

In ribbon labels use $f^{-1}\sigma_i f$, with inverse label map $f$.
Form the covering ribbon graph whose $a$- and $b$-edges are these
two permutations, with rotation
$(a_{\rm out},a_{\rm in},b_{\rm out},b_{\rm in})$.
Fill the smooth faces of degrees three over~0 and two over~1,
using $A_{\rm per}^3=B_{\rm per}^2=-I$ projectively. Delete a separating edge on each
such face and contract a spanning tree, conjugating the remaining
edge words by their tree paths. Restoring those edges and disks gives
the inverse maps on fundamental groups. The resulting three chords
$x,y,z$ have the following positive peripheral bases; the fourth
loop is $(uvw)^{-1}$.
\[
\begin{array}{c|ccc}
\text{row}&u&v&w\\\hline
1&z^{-1}yx^{-1}&x&y^{-1}\\
2&x^{-1}z^{-1}y&y^{-1}&z\\
3&z^{-1}x^{-1}y&y^{-1}xzx^{-1}y^{-1}xz^{-1}x^{-1}y&y^{-1}xzx^{-1}y\\
4&x^{-1}y&y^{-1}z^{-1}&z\\
5&x^{-1}y^{-1}&z^{-1}&zy\\
6&zx^{-1}&y^{-1}x&x^{-1}yx\\
7&zy^{-1}x^{-1}&x&yz^{-1}y^{-1}\\
\end{array}
\]
\[
\begin{array}{c|l}
\text{row}&(x,y,z)\\\hline
1&(v,w^{-1},w^{-1}v^{-1}u^{-1})\\
2&(w^{-1}v^{-1}u^{-1},v^{-1},w)\\
3&(w^{-1}v^{-1}u^{-1},w^{-1}v^{-1}w,uwu^{-1})\\
4&(w^{-1}v^{-1}u^{-1},w^{-1}v^{-1},w)\\
5&(w^{-1}v^{-1}u^{-1},vw,v^{-1})\\
6&(vw,vwv^{-1},uvw)\\
7&(v,w^{-1}v^{-1}u^{-1},uvw^{-1}v^{-1}u^{-1})\\
\end{array}
\]

Each indicated inverse substitutes back to $x,y,z$, and the reverse
substitution gives $u,v,w$. These are based free-group isomorphisms.
The word operations are the positive Hurwitz map
$(U,V)\mapsto(V,V^{-1}UV)$ and its inverse
$(U,V)\mapsto(UVU^{-1},U)$. The final lattice maps, with
$\mathsf M_i^{\rm ell}H=HA_i$ for all four peripheral matrices, are
\[
\begin{array}{c|c|c}
\text{row}&H&H^{-1}\\\hline
1&(0,-1;1,1)&(1,1;-1,0)\\
2&(-2,-1;1,0)&(0,1;-1,-2)\\
3&(-1,7;-1,6)&(6,-7;1,-1)\\
4&(1,-1;1,0)&(0,1;-1,1)\\
5&(1,-1;1,0)&(0,1;-1,1)\\
6&(2,1;1,1)&(1,-1;-1,2)\\
7&(1,-1;0,1)&(1,1;0,1).
\end{array}
\]
Here $(a,b;c,d)$ denotes a $2\times2$ matrix by rows.
All determinants are one. The actual elliptic local system supplies
an $SL_2(\mathbb Z)$ lift of the projective representation. Two
such lifts differ by a central sign character. The positive trace~2
at the three multiplicative places fixes their signs; the nonzero
additive trace, or for $III$ the product-one relation, fixes the last.
The smooth filled faces have actual identity monodromy. Hence these
are precisely the positive integral matrices used below.

\emph{The section cocycle and the normalized Poincar\'e line.}
Here is the local component map with its sign. At an $I_n$ fibre,
write its multiplicative period as $q(t)=t^n a(t)$ with the
holomorphic unit $a$ retained. The $n$ node annuli give
$E_t=\mathbb C^\times/q(t)^{\mathbb Z}$; their product gluing
factor is $t^n a(t)$. In the positive basis $(a_0,b_0)$, where
$a_0$ turns once in $\mathbb C^\times$ and
$b_0=\log q/(2\pi i)$, the matrix is
$A=\left(\begin{smallmatrix}1&n\\0&1\end{smallmatrix}\right)$.
A section of component $j$ has representative $u=t^j b(t)$ with
$b$ a unit. Its node chart and inverse are
\[
 U=u/t^j,\quad V=t^{j+1}/u,\quad UV=t,\qquad u=t^jU.
\]
The adjacent chart has $U_{j+1}=V_j^{-1}$ and
$V_{j+1}=tV_j$; the cyclic transition is $u\mapsto u/q$ with
inverse $u\mapsto uq$. Replacing $u$ by $uq^r$ sends
$(j,b)$ to $(j+nr,ba^r)$, with inverse obtained by $-r$.
Positive continuation adds $j$ to $\log u/(2\pi i)$, so the
negative-increment convention gives $c=-ja_0$. Consequently
\[
 S=\mathbb Za_0,\quad (A-I)\mathbb Z^2=n\mathbb Za_0,\qquad
 j=-\operatorname{coeff}_{a_0}(c)\pmod n.
\]
The representative change gives
$-nr a_0=(A-I)(-r b_0)$. For an integral marking $H=(v,w)$,
the forward maps are $a_0\mapsto v$, $b_0\mapsto w$,
$c\mapsto-jv$, and the inverse is $H^{-1}$. Thus narrowness is
exactly $c\in(A-I)\mathbb Z^2$. A raw $Q$ increment $\kappa_{\rm loc}v$ means
component $-\kappa_{\rm loc}$; the forward linearization covering $-Q$ advances
the component by $+\kappa_{\rm loc}$. At a finite additive place $A-I$ is
invertible over $\mathbb Q$, so its saturated image is $\mathbb Z^2$.

Negative Abel--Jacobi increments of an actual section give
$c_i\in\operatorname{sat}\operatorname{im}(T_i^{\rm ell}-I)$,
with $\sum_iT_1^{\rm ell}\cdots T_{i-1}^{\rm ell}c_i=0$.
Changing its logarithm by $h\in\mathbb Z^2$ adds
$(T_i^{\rm ell}-I)h$. In local saturated-image coordinates write
these two maps as $C$ and $R$. The integer kernel and coboundary
maps give $\ker R/\operatorname{im}C=\mathbb Z\langle g\rangle$.
Their full inverse decomposes any closed vector uniquely as
\[
c=a_G g+C h.
\]
Explicitly, for a kernel inclusion $K$, inverse $N$, and elementary
integer changes $U,V$ satisfying $UNCV=(I_2,0)^t$, if the displayed
$g$ has quotient coordinate $\sigma_g=\pm1$, then
\[
a_G=\sigma_g(UNc)_3,\qquad
h=V\bigl(UN(c-a_Gg)\bigr)_{1,2}.
\]
The following local bases $\mathcal U_i$ and complete integer matrices $\mathsf M_{\rm sec}$
give these maps in the original peripheral order. The first two columns
of $\mathsf M_{\rm sec}$ are $C$, its third is the local generator $g$, and
its last two form a complement $E$. Each has determinant $\pm1$,
and direct multiplication gives $R\mathsf M_{\rm sec}=(0_{2\times3},I_2)$. Thus the
first three rows of $\mathsf M_{\rm sec}^{-1}=\operatorname{adj}(\mathsf M_{\rm sec})/\det \mathsf M_{\rm sec}$ give
$(h_1,h_2,a_G)$ and the last two give $R$. In particular both
assembly and decomposition maps are integral and inverse.
\[
\begin{array}{c|cccc}
\text{row}&\mathcal U_1&\mathcal U_2&\mathcal U_3&\mathcal U_4\\\hline
1&\left(\begin{smallmatrix}1\\0\end{smallmatrix}\right)&\left(\begin{smallmatrix}1&0\\-1&1\end{smallmatrix}\right)&\left(\begin{smallmatrix}1\\-2\end{smallmatrix}\right)&\left(\begin{smallmatrix}-1\\-1\end{smallmatrix}\right)\\
2&\left(\begin{smallmatrix}1\\0\end{smallmatrix}\right)&\left(\begin{smallmatrix}1\\-1\end{smallmatrix}\right)&\left(\begin{smallmatrix}1\\-3\end{smallmatrix}\right)&\left(\begin{smallmatrix}-1&0\\-1&-1\end{smallmatrix}\right)\\
3&\left(\begin{smallmatrix}1\\0\end{smallmatrix}\right)&\left(\begin{smallmatrix}1\\-3\end{smallmatrix}\right)&\left(\begin{smallmatrix}-1\\-1\end{smallmatrix}\right)&\left(\begin{smallmatrix}1&-8\\0&-1\end{smallmatrix}\right)\\
4&\left(\begin{smallmatrix}1\\0\end{smallmatrix}\right)&\left(\begin{smallmatrix}1\\-1\end{smallmatrix}\right)&\left(\begin{smallmatrix}3&1\\-5&-2\end{smallmatrix}\right)&\left(\begin{smallmatrix}-1\\-1\end{smallmatrix}\right)\\
5&\left(\begin{smallmatrix}1\\0\end{smallmatrix}\right)&\left(\begin{smallmatrix}1\\-1\end{smallmatrix}\right)&\left(\begin{smallmatrix}1&0\\-3&-1\end{smallmatrix}\right)&\left(\begin{smallmatrix}-1\\-1\end{smallmatrix}\right)\\
6&\left(\begin{smallmatrix}1\\0\end{smallmatrix}\right)&\left(\begin{smallmatrix}1\\-2\end{smallmatrix}\right)&\left(\begin{smallmatrix}5&-2\\-13&5\end{smallmatrix}\right)&\left(\begin{smallmatrix}-1\\-1\end{smallmatrix}\right)\\
7&\left(\begin{smallmatrix}1\\0\end{smallmatrix}\right)&\left(\begin{smallmatrix}0&-1\\1&0\end{smallmatrix}\right)&\left(\begin{smallmatrix}-3\\-1\end{smallmatrix}\right)&\left(\begin{smallmatrix}1\\0\end{smallmatrix}\right)\\
\end{array}
\]
\[
\begin{array}{c|c}
\text{row}&\mathsf M_{\rm sec}\\\hline
1&\left(\begin{smallmatrix}0&8&3&1&-7\\1&3&1&0&-1\\0&1&0&0&0\\2&1&1&0&0\\1&-1&0&0&0\end{smallmatrix}\right)\\
2&\left(\begin{smallmatrix}0&7&5&1&-6\\2&2&1&0&-1\\3&1&0&0&0\\1&-1&-1&0&0\\0&1&1&0&0\end{smallmatrix}\right)\\
3&\left(\begin{smallmatrix}0&7&5&1&-7\\3&1&0&0&0\\1&-1&-1&0&1\\2&-11&-8&0&-5\\2&-12&-9&0&0\end{smallmatrix}\right)\\
4&\left(\begin{smallmatrix}0&5&-3&1&-4\\4&4&-1&0&-1\\1&1&0&0&0\\1&0&0&0&0\\1&-1&1&0&0\end{smallmatrix}\right)\\
5&\left(\begin{smallmatrix}0&5&-4&1&-4\\2&2&-1&0&-1\\0&1&-1&0&0\\3&0&1&0&0\\1&-1&1&0&0\end{smallmatrix}\right)\\
6&\left(\begin{smallmatrix}0&5&-4&1&-4\\6&3&-1&0&0\\3&1&0&0&0\\2&0&1&0&-1\\1&-1&1&0&0\end{smallmatrix}\right)\\
7&\left(\begin{smallmatrix}0&2&1&1&-2\\1&-3&0&0&1\\0&3&1&0&0\\1&-3&0&0&0\\0&1&0&0&0\end{smallmatrix}\right)\\
\end{array}
\]
The divisor-line isomorphism to the circle model determines the
Mordell--Weil generator's integer coordinate $\varepsilon_G$ in
$c_G=\varepsilon_Gg+Ch_G$. For $\operatorname{Im}\tau>0$ define
\[
 \vartheta(z,\tau)=\sum_{\nu\in\mathbb Z}
 \exp\!\left(\pi i(\nu+\tfrac12)^2\tau+
            2\pi i(\nu+\tfrac12)(z+\tfrac12)\right).
\]
The series and its $z$ derivatives converge normally on compact
sets. A nonzero Fourier coefficient shows it is not identically
zero. Index shifts give
\[
 \vartheta(z+1)=-\vartheta(z),\qquad
 \vartheta(z+\tau)=-e^{-\pi i\tau-2\pi iz}\vartheta(z).
\]
Pairing $\nu$ with $-\nu-1$ proves oddness. On a translated period
parallelogram with no boundary zero, the logarithmic derivative
is $1$-periodic and decreases by $2\pi i$ under $\tau$.
Its positive boundary integral is $2\pi i$, so there is one simple
zero per period cell, at the period lattice. Hence
\[
 F(z,w)=\frac{\vartheta_z(0)\vartheta(z-w)}
                  {\vartheta(z)\vartheta(-w)}
\]
has divisor $\Delta-O_1-O_2$, with
$\lim_{z\to0}zF=1$ and $\lim_{w\to0}(-w)F=1$.
Pairing these residues with the invariant differential fixes the
zero-axis rigidifications. The normal cotangent factor on either
axis is the Hodge line, so the normalized line over the family is
$\mathcal O(\Delta-O_1-O_2)\otimes\pi^*\mathcal L_{\rm Hodge}^{-1}$.

Write $z=x_1+\tau x_2$ and $w=y_1+\tau y_2$. The theta
identities give its holomorphic line coordinate $\xi$ the multiplier
\[
 (z,w,\xi)\longmapsto
 (z+m+n\tau,w+p+q\tau,
  e^{2\pi i(nw+qz+nq\tau)}\xi).
\]
The explicit smooth bundle map and its inverse are
\[
 \eta=e^{-2\pi iz y_2}\xi,\qquad
 \xi=e^{2\pi iz y_2}\eta.
\]
The resulting exponent is $ny_1-my_2-mq$; the last term is
integral. Thus $|\eta|=1$, with $\eta=e^{-2\pi i\theta}$,
gives precisely the circle quotient
\[
 (x,y,\theta)\longmapsto(x+r,y+s,\theta+r^tJy),
 \quad r,s\in\mathbb Z^2,\qquad
 \alpha=d\theta-x^tJdy.
\]
Both inverse maps preserve the two zero-axis rigidifications.
The connection is lattice invariant. Simultaneous elliptic
inversion is $(z,w,\xi)\mapsto(-z,-w,\xi)$ and is its own
inverse on this rigidified bundle.

For $A=\left(\begin{smallmatrix}a&b\\c_0&d_A\end{smallmatrix}\right)
\in SL_2(\mathbb Z)$ set $\Delta_A=d_A-c_0\tau$. Then
\[
 \tau'=\frac{a\tau-b}{\Delta_A},\quad x'=Ax,\quad y'=Ay,\quad
 z'=\frac z{\Delta_A},\quad w'=\frac w{\Delta_A},\quad
 \operatorname{Im}\tau'=\frac{\operatorname{Im}\tau}{|\Delta_A|^2}.
\]
Holding $\theta$ fixed gives the holomorphic line map
\[
 \xi'=\exp(2\pi i c_0zw/\Delta_A)\xi,
 \qquad z'y'_2-zy_2=c_0zw/\Delta_A.
\]
Its inverse uses $A^{-1}$; its denominator is
$a+c_0\tau'=\Delta_A^{-1}$ and its exponential factor cancels the first.
These maps respect the theta lattice factors and the zero-axis
rigidifications, so they identify the actual normalized divisor
line under a period change. Two such holomorphic rigidified
identifications differ by a constant on $E\times E$, fixed to one
on the axes.

More generally the affine circle map with integral increments
$r=(r_1,r_2)$ and $s=(s_1,s_2)$ is
\[
 (x,y,\theta)\longmapsto
 (Ax+r,Ay+s,\theta+r^tJAy+\kappa).
\]
Its inverse has parameters
$(A^{-1},-A^{-1}r,-A^{-1}s,r^tJs-\kappa)$.
With $z_A=z/\Delta_A,w_A=w/\Delta_A$, the holomorphic fibre multiplier is
\[
 \exp 2\pi i\bigl(c_0zw/\Delta_A+r_2w_A+s_2z_A+
                         r_2s_2\tau'-\kappa\bigr).
\]
The literal circle-gauge computation differs by the integer
$r_1s_2$. The forward and inverse exponents sum to
$-r_1s_2-(A^{-1}r)_1(A^{-1}s)_2$, again integral, so the
inverse is literal on the line as well as on the circle.

For the local affine increments $c,c'$ this gives the
connection-preserving lift
\[
\theta\longmapsto\theta+c^tJAy+\kappa.
\]
At fixed representatives $x=-z,y=-z'$ with
$(A-I)z=c$, $(A-I)z'=c'$, zero real displacement gives
$\kappa=c^tJAz'=-\beta z'$, where $\beta=-c^tJA$.
Composition over the positive peripheral loops gives the central
translation
\[
\mathcal B(c,c')=-\sum_i\beta_i z_i'
 -\sum_{i<j}\beta_iA_{i+1}\cdots A_{j-1}c_j'.
\]
Substitution of $c=a_Gg+Ch$, $c'=b_Gg+Cj$ gives
$\mathcal B(c,c')=a_Gb_G/(pq)$; this is an identity in the
arbitrary integer variables.

To identify its full integer normalization, multiply the actual
sections by $N_{\rm comp}=pq$, which kills every fibre component group.
The resulting sections are narrow. On the product of identity
N\'eron models the normalized line is
\[
\mathcal O(\Delta-O_1-O_2)\otimes\pi^*\mathcal L_{\rm Hodge}^{-1}.
\]
Its zero-axis rigidifications fix its extension: two generic
identifications differ by a multiple of the irreducible closed
fibre, and restriction to a zero axis forces that multiple to
vanish; see also \cite[Proposition 6.7.1, p.~306]{manEdixhovenLido23}.
Locally a narrow section has a holomorphic logarithm after an
integral period shift. The contraction
$(P,Q)\mapsto(\exp(t\log P),\exp(t\log Q))$, $0\leq t\leq1$,
identifies its disk frame with the zero rigidification. Parallel
transport along this contraction and along its reverse gives the
two inverse frame maps. For narrow increments choose integral
$z,z'$ with $c=(A-I)z$, $c'=(A-I)z'$. The literal quotient
circle map has $\kappa=0$ modulo $\mathbb Z$, whereas its
fixed-point prescription $\kappa=c^tJAz'$ is an integer.
They are therefore real lifts of the same circle map, and their
difference from the extending disk frame is integer-valued.
Shift the local branches by these integral periods, so their local
cocycles vanish, and perform the contraction there. The integer
discrepancy is continuous in the contraction parameter and zero
at the zero pair, hence zero throughout. Reversing the branch
shift and contraction gives the inverse frame identification.
For non-narrow rational fixed representatives we use the rational
pairing only; its identification with height follows after
multiplication to narrow sections below.

The positive clutching sign is fixed by $\eta=e^{-2\pi i\theta}$:
a section $t^r$ has theta winding $-r$ and divisor degree $r$.
Thus the degree of the pulled-back Poincar\'e line is $-\mathcal B$.
For a narrow pair its divisor expression has degree
$P\cdot Q-P\cdot O-Q\cdot O-1=-\langle P,Q\rangle$.
This includes a repeated section, since its diagonal normal line
has degree $P^2=-1$. The Shioda projection is a homomorphism
\cite[Proposition 11.4 and the height formula, pp.~50--51]{manSS10}.
Both its height and $\mathcal B$ scale by $N_{\rm comp}^2$ under
multiplication by $N_{\rm comp}$. Dividing the narrow equality
therefore gives $\mathcal B(c_G,c_G)=\langle G,G\rangle$ with
its entire integer part.

The seven root lattices have rank seven and determinants
$8,14,14,20,30,30,6$. Torsion is zero in the $II$ rows by
specialization to their additive identity group; in the other rows
its square divides the squarefree root determinant. The rational
surface's N\'eron--Severi lattice is unimodular, so the Schur
complement gives $\langle G,G\rangle=1/(pq)$, in agreement with
\cite[Main Theorem table, p.~85]{manOguisoShioda91}.
Consequently $\varepsilon_G^2=1$. The displayed coboundary inverse now gives
the exact integral change of logarithm. For $c_G=\varepsilon_Gg+(A-I)h_G$ put
\[
S_{h_G}=\begin{pmatrix}I_2&0&qh_G\\-ph_G^tJ&1&0\\0&0&1\end{pmatrix},\qquad
S_{h_G}^{-1}=\begin{pmatrix}I_2&0&-qh_G\\ph_G^tJ&1&0\\0&0&1\end{pmatrix}.
\]
With $C_{\varepsilon_G}=\operatorname{diag}(\varepsilon_G I_2,1,1)S_{h_G}$, the full lifted
matrices obey $T_{\rm actual}=C_{\varepsilon_G}^{-1}T_{\rm marked}C_{\varepsilon_G}$,
including every weight and central coefficient. Elliptic inversion
realizes $\varepsilon_G=-1$ on the bundles. These maps and their inverses carry
the chosen periods, affine vector and attachment maps to the stated
marking.


\emph{The line bundle, its linearization, and a section.}
The indicated elliptic surfaces have
$\langle G,G\rangle=1/(pq)$. At every fibre one of $P,Q$ is
narrow. At the two weight-zero multiplicative places $P$ is
narrow. The following table lists the negative component residues
$a(Q)$, namely the coefficients of the negative Abel--Jacobi
increments in the chosen invariant period. To relate them to the
geometric component label, use the explicit involution
\[
 j_{\rm raw}=-j_{\rm geom}\pmod n,\qquad
 j_{\rm geom}=-j_{\rm raw}\pmod n.
\]
Thus $[Q]_{\rm geom}=-a(Q)$; for example row~1 has
$a(Q)=3$ and geometric component $5$ modulo $8$.
At these nontrivial places choose $v=(1,-u)$, with $u=0$ at
upper-triangular matrices, $u=1$ at the other $I_2,I_4$ matrices, and $u=2$
at the $I_3$ of row~6. The literal period maps are
\[
 H_u=\begin{pmatrix}1&0\\-u&1\end{pmatrix},\quad
 H_u^{-1}=\begin{pmatrix}1&0\\u&1\end{pmatrix},\qquad
 A=H_u\begin{pmatrix}1&n\\0&1\end{pmatrix}H_u^{-1}.
\]
The raw $Q$ increment is $a(Q)v$ modulo $n\mathbb Zv$.
In these coordinates the residues are
\[
\begin{array}{c|l}
\text{row}&a(Q)\\\hline
1&3\pmod8\\
2&5\pmod7,\ 1\pmod2\\
3&3\pmod7\\
4&2\pmod5,\ 3\pmod4\\
5&3\pmod5,\ 1\pmod2\\
6&2\pmod5,\ 1\pmod3\\
7&1\pmod2.
\end{array}
\]
At the additive place $Q$ is narrow. The negative component
class of $P$ is zero at $II$ and nonzero modulo two at $III$.
For the displayed $IV$ and $IV^*$ matrices, the quotient maps
on an increment $c=(c_1,c_2)$ are respectively
$c\mapsto-c_2$ and $c\mapsto-c_1$ modulo three; their sections
are $j\mapsto(0,-j)$ and $j\mapsto(-j,0)$.
Their kernels are $(A-I)\mathbb Z^2$, since they contain that
image and both indices are three. They give the raw classes
$1$ and $2$ for $P$; the geometric classes are their negatives.

Set $\mathcal L=\mathcal Hom(t_Q^*M,M)$. Translation preserves
degree-zero line bundles on a smooth elliptic curve, so
$\mathcal L$ is vertical. Its component degrees vanish: $P$
is narrow at the multiplicative places, whereas $t_Q$ preserves
components at the additive place. The kernel of each fibre's
intersection matrix is generated by the full fibre. Since there
are no multiple fibres, $\mathcal L=\pi^*\mathcal O(b)$, where
\[
b=\mathcal L\cdot O
 =1+P\cdot O+Q\cdot O-P\cdot Q
 =\langle P,Q\rangle=1.
\]
Here the local mixed height terms vanish because one section is
narrow. Choose a section $\phi$ of $\mathcal L$ with its sole
simple zero at $I_1^\bullet$. Its action on $M$ covers $-Q$.
Multiplying it by a sufficiently small nonzero constant makes
its norm uniformly less than one in a Hermitian metric on the
compact surface. Off its zero fibre, iterates therefore act
freely and properly on $M^\times$; a closed annulus in the line
direction gives a compact fundamental region on each compact
base subset. This is the contraction construction of
\cite[Lemma~2.1 and Corollary~2.2, pp.~3--4; Proposition~2.4, p.~5]{E26}.

The restriction $M|_O$ has degree
\[
d_M=P\cdot O+1
 =\frac{p/q+\operatorname{contr}_{\rm add}(P)}2
 =(4,7,2,10,2,4,1)
\]
in the seven rows. The local terms are $0,1/2,2/3,4/3$
for $II,III,IV,IV^*$, respectively. Choose a section with its
only zero, of order $d_M$, at $I_1^\bullet$. Its quotient is a
section over the regular base and over the weight-zero fillings.
At the remaining reduced place, with parameter $t$, write
\[
\phi=tU,\qquad U:t_Q^*M\xrightarrow{\sim}M.
\]
Division by the full reduced Cartier fibre leaves a nowhere
zero section of the Hom line, including at the node. Thus
$U_j=U|_{jQ}:M|_{(j+1)Q}\to M|_{jQ}$ is invertible.
If the chosen section is $t^{d_M}v$ along $O$, its $d_M$ inverse
iterates have representative
\[
t^{-d_M}U_{d_M-1}^{-1}\cdots U_0^{-1}(t^{d_M}v)
 =U_{d_M-1}^{-1}\cdots U_0^{-1}v\quad\hbox{along }d_M Q.
\]
Each section $jQ$ meets the smooth locus of an $I_1$ fibre:
in $uv=t$ its two coordinates cannot both vanish at $t=0$.
The displayed representative is consequently nonzero in the
open chart of the toroidal filling. On the overlap it is in
the same $\phi$-orbit; $d_M$ forward iterates give the inverse
comparison. This proves extension of the section and fixes
all three reduced meridians. The local inverse period is also
the one in \cite[Propositions~2.7 and~2.9, pp.~6--7]{E26}.

\emph{The period marking.}
For an elliptic matrix $A$, let $g$ be the marked cocycle of $G$
and put $p_i=pg$, $q_i=qg$. With
$J=\left(\begin{smallmatrix}0&1\\-1&0\end{smallmatrix}\right)$,
choose $(A-I)h=p_i$, $(A-I)z=q_i$ over $\mathbb Q$ and put
$b=-h^tJ$, $\beta=-p_i^tJA$. The lift with local weight $w$ is
\[
T=\begin{pmatrix}A&0&q_i\\\beta&1&\beta z+w\\0&0&1\end{pmatrix}.
\]
This marking follows from the Poincar\'e line bundle.
By the identities for $\vartheta$ above, the quotient
$\vartheta(z-P)/\vartheta(z)$ has divisor $P-O$ and
$\tau$-multiplier $e^{2\pi iP}$. Continuation of $P$ by $1$
adds $\delta$ to the second elliptic period. Continuation by
$\tau$, followed by the inverse gauge $w\mapsto e^{-2\pi iz}w$,
subtracts $\delta$ from the first. Thus an Abel--Jacobi increment
$\Delta P$ contributes $\Delta P^tJ$. Realize the marked
cocycles as $p_i=-\Delta P$, $q_i=-\Delta Q$. Since the forward
period covers $-Q$, its column is $+q_i$ and its line winding
is $+w$. This gives precisely $T$ above.

The full change of coordinates and its inverse are
\[
S=\begin{pmatrix}I_2&0&z\\b&1&bz\\0&0&1\end{pmatrix},\qquad
S^{-1}=\begin{pmatrix}I_2&0&-z\\-b&1&0\\0&0&1\end{pmatrix},
\qquad T=S^{-1}(\operatorname{diag}(A,1,1)+wE_{\delta\gamma})S.
\]
These identify the torus with lattice $\Lambda$ with the torus
with lattice $S\Lambda$. At the additive place $z$ is integral
and $\ell=\gamma-z$, so $S\ell=\gamma$ and $S\delta=\delta$,
including their line coordinates. A different scalar $c_0$ in
the last column is carried to this choice by
$I+(bz-c_0)E_{\delta\gamma}$, with inverse
$I-(bz-c_0)E_{\delta\gamma}$, between the corresponding image
lattices. At a place where $Q$ is narrow, the integral change
$\gamma\mapsto\gamma-z$ leaves central weight $w$; where $P$
is narrow, the integral change $u\mapsto u+bx$ kills $\beta$
and again leaves weight $w$. Opposite additions give both
inverses. Simultaneously changing the generator's sign is
elliptic inversion: it pulls $\mathcal O(P-O)$ to
$\mathcal O(-P-O)$, carries the quotient action, and induces
$\operatorname{diag}(-I_2,1,1)$ on $\Lambda$.
In the geometric labels the forward action covering $-Q$ is
\[
 j_{\rm geom}\longmapsto j_{\rm geom}+a(Q),
 \qquad j_{\rm geom}\longmapsto j_{\rm geom}-a(Q)
 \quad\text{for its inverse}.
\]
Conjugation by the displayed involution
$j_{\rm raw}=-j_{\rm geom}$ gives
$j_{\rm raw}\mapsto j_{\rm raw}-a(Q)$, whose inverse adds
$a(Q)$. Elliptic inversion negates both $j_{\rm geom}$ and
$a(Q)$ and intertwines these displayed actions with the actions
for the opposite generator.

Use cyclic order with the additive place first and write the
lifts as $T_0,T_1,T_2,T_3$. The following table specifies each
elliptic matrix and generator cocycle in that order: an entry
$(a,b,c,d;u,v;w)$ means
$A=\left(\begin{smallmatrix}a&b\\c&d\end{smallmatrix}\right)$,
$g=(u,v)^t$, with weight $w$.
\[
\begin{array}{c|cc}
\text{row}&(A_0,g_0,\nu_0)&(A_1,g_1,\nu_1)\\\hline
1&(2,3,-1,-1;1,-1;0)&(3,1,-4,-1;1,-2;1)\\
2&(0,1,-1,1;1,0;0)&(1,7,0,1;5,0;0)\\
3&(-13,85,-2,13;64,9;0)&(1,7,0,1;5,0;0)\\
4&(5,3,-7,-4;0,0;0)&(0,1,-1,2;-1,-1;1)\\
5&(1,1,-3,-2;-1,2;0)&(0,1,-1,2;-1,-1;1)\\
6&(12,5,-29,-12;-2,5;0)&(0,1,-1,2;-1,-1;1)\\
7&(1,-3,1,-2;-1,0;0)&(-2,9,-1,4;0,0;1)\\
\end{array}
\]
\[
\begin{array}{c|cc}
\text{row}&(A_2,g_2,\nu_2)&(A_3,g_3,\nu_3)\\\hline
1&(0,1,-1,2;0,0;0)&(1,8,0,1;3,0;0)\\
2&(3,2,-2,-1;1,-1;0)&(4,1,-9,-2;0,0;1)\\
3&(4,1,-9,-2;0,0;1)&(0,1,-1,2;1,1;0)\\
4&(1,5,0,1;-3,0;0)&(5,4,-4,-3;-1,1;0)\\
5&(1,5,0,1;-4,0;0)&(3,2,-2,-1;-1,1;0)\\
6&(1,5,0,1;-4,0;0)&(7,3,-12,-5;-1,2;0)\\
7&(1,1,0,1;0,0;0)&(1,2,0,1;1,0;0)\\
\end{array}
\]
Substitution in the formula for $T$ gives
$T_0T_1T_2T_3=I$, $T_0^d=I$, and
$\chi T_i=\chi$, $T_i\delta=\delta$.
The form
\begin{equation}
\omega=p\,e_1^*\wedge e_2^*-q\,\delta^*\wedge\gamma^*
\hypertarget{tag-023I}{}\label{maps:eq:D.R2}
\end{equation}
is preserved by every $T_i$.

\emph{The three reduced fillings.}
For $i=1,2,3$ choose the columns $(w_i,b_i,\delta,c_i)$
listed below as a matrix $H_i$. Its determinant is $1$ or $-1$;
the inverse is $\operatorname{adj}(H_i)/\det H_i$.
At a weight-zero place
$H_i^{-1}T_iH_i=I+E_{12}$, and at the marked place it is
$I+E_{12}+E_{34}$. In particular the displayed $w_i$ are
primitive vanishing circles, and the fixed basis is
$(w_i,\delta,c_i)$ or $(w_i,\delta)$, respectively.
\[
\begin{array}{c|ccc}
\text{row}&w_1&b_1&c_1\\\hline
1&(1,-2,-8,0)^t&(0,1,0,0)^t&(0,-1,0,1)^t\\
2&(1,0,-10,0)^t&(0,-2,0,3)^t&(0,5,0,-7)^t\\
3&(1,0,-5,0)^t&(0,3,0,-2)^t&(0,-10,0,7)^t\\
4&(-1,-1,-20,0)^t&(1,0,0,0)^t&(-1,0,0,1)^t\\
5&(-1,-1,-10,0)^t&(1,0,0,0)^t&(-3,0,0,1)^t\\
6&(-1,-1,-15,0)^t&(1,0,0,0)^t&(-2,0,0,1)^t\\
7&(-3,-1,0,0)^t&(1,0,0,0)^t&(0,0,0,1)^t\\
\end{array}
\]
\[
\begin{array}{c|ccc}
\text{row}&w_2&b_2&c_2\\\hline
1&(-1,-1,0,0)^t&(1,0,0,0)^t&(0,0,0,1)^t\\
2&(1,-1,-7,0)^t&(0,0,0,1)^t&(1,0,0,-2)^t\\
3&(1,-3,0,0)^t&(0,1,0,0)^t&(0,0,0,1)^t\\
4&(-1,0,-12,0)^t&(0,1,0,2)^t&(0,-3,0,-5)^t\\
5&(1,0,8,0)^t&(0,5,0,2)^t&(0,-12,0,-5)^t\\
6&(1,0,12,0)^t&(0,-3,0,-2)^t&(0,8,0,5)^t\\
7&(1,0,0,0)^t&(0,1,0,0)^t&(0,0,0,1)^t\\
\end{array}
\]
\[
\begin{array}{c|ccc}
\text{row}&w_3&b_3&c_3\\\hline
1&(1,0,-3,0)^t&(0,-1,0,3)^t&(0,3,0,-8)^t\\
2&(1,-3,0,0)^t&(0,1,0,0)^t&(0,0,0,1)^t\\
3&(-1,-1,7,0)^t&(1,0,0,0)^t&(2,0,0,1)^t\\
4&(-1,1,-5,0)^t&(0,0,0,1)^t&(1,0,0,4)^t\\
5&(1,-1,5,0)^t&(2,0,0,1)^t&(-3,0,0,-2)^t\\
6&(-1,2,-5,0)^t&(0,1,0,2)^t&(0,-2,0,-3)^t\\
7&(1,0,-1,0)^t&(0,-1,0,1)^t&(0,3,0,-2)^t\\
\end{array}
\]
At a weight-zero $I_j$, $a(Q)$ is a unit modulo $j$.
The forward geometric permutation $r\mapsto r+a(Q)$ and
its inverse $r\mapsto r-a(Q)$ cycle through all $j$ components.
After $j$ steps the normalized quotient is
$\mathbb P_E(\mathcal O\oplus L)$, where $\deg L=0$.
For flat multipliers $\mu_1,\mu_2$, the map multiplying a
line coordinate by
$\exp(-t_1\log\mu_1-t_2\log\mu_2)$ trivializes the underlying
smooth line bundle; reversing the sign is its inverse.
The zero and infinity sections are glued by a translation
$x\mapsto x+a$ on $E$. Choose $f:\mathbb P^1\to[0,1]$
with $f(0)=0$, $f(\infty)=1$. Then
$(z,x)\mapsto(z,x-f(z)a)$, with inverse given by addition,
identifies the quotient with
$(\mathbb P^1/(0\sim\infty))\times T^2$.
An arc from $0$ to $\infty$ splits its normalization sequence,
giving $H_j=H_j(\mathbb P^1\times T^2)\oplus H_{j-1}(T^2)$.
The local equation $uv=t$ is unramified under the free
contracting quotient. Thus the nearby vanishing circle is
exactly $w_i$, with multiplicity one.

In coordinates $(w,b,\delta,c)$ its boundary maps are
\[
\begin{aligned}
H_1(T^4)&\longrightarrow\mathbb Z^3,&
(x_1,x_2,x_3,x_4)&\longmapsto(x_2,x_3,x_4),\\
\bigwedge^2\Lambda/(\bigwedge^2T-I)&\longrightarrow\mathbb Z^4,&
(x_{12},x_{13},x_{14},x_{23},x_{24},x_{34})
 &\longmapsto(x_{12},x_{23},x_{24},x_{34}).
\end{aligned}
\]
Their sections insert zeros into the omitted coordinates.
The two exterior relations are $w\wedge\delta,w\wedge c$.
At the marked $I_1$, the two periods in
\cite[Proposition~2.7, p.~6]{E26} have primitive windings
$w$ and $\delta$. The primitive hexagonal filling of
Theorem~\ref{maps:src:D.2.6}, equivalently
\cite[Proposition~3.1, pp.~6--7]{manViaclovskyTwo26}, has maps
\[
\begin{aligned}
H_1(T^4)&\longrightarrow\mathbb Z^2,&x&\longmapsto(x_2,x_4),\\
\bigwedge^2\Lambda/(\bigwedge^2T-I)&\longrightarrow\mathbb Z^4,&x
 &\longmapsto(x_{12},x_{14}-x_{23},x_{24},x_{34}).
\end{aligned}
\]
Sections are $(u,v)\mapsto(0,u,0,v)$ and
$(u,v,w,z)\mapsto(u,0,v,0,w,z)$.
The exterior relations are $w\wedge\delta$ and
$w\wedge c+b\wedge\delta+w\wedge\delta$.
Composing with $H_i^{-1}$ and $\bigwedge^2H_i^{-1}$ gives
the maps in the common marking; composing their sections with
$H_i$ and $\bigwedge^2H_i$ gives the inverse quotient maps.

The fixed circle actions extend to these fillings. Their
normal-circle tori form the full kernel of specialization in
degree two. The relative groups in degrees $3,2,1$ are
$(\mathbb Z^3,\mathbb Z,0)$ at a one-node filling and
$(\mathbb Z^2,\mathbb Z,0)$ at the hexagonal filling.
Indeed these neighbourhoods retract to their central surfaces.
The normalization splitting gives free $H_2,H_3$ of ranks
$4,3$ at the one-node surface; the primitive toric cell model
gives free $H_2,H_3$ of ranks $4,2$ at the hexagonal surface.
The universal coefficient theorem and Poincar\'e--Lefschetz
duality identify relative $H_3$ with the corresponding free
$H^3$, relative $H_2$ with $H^4=\mathbb Z$, and relative
$H_1$ with $H^5=0$. The displayed specialization sections
and the Wang sequence then identify relative $H_3$ with
the full sweep kernel. Orient a relative three-cycle as the negative
product of a fixed fibre loop with its positive complex normal
disc. Its boundary is the positive fibre-then-base sweep.
The extended section gives the positive relative two-disc.

\emph{The regular-family lattice and the reduced attachments.}
Let $V$ lie over the four-punctured sphere. The chosen section
and the real period map give a diffeomorphism from the flat
torus bundle to $V$: $[x]\mapsto\exp(\Pi_bx)$, whose inverse
is the group logarithm modulo $\Pi_b\Lambda$.
The graph complex has boundary $-\partial_{\rm gr}$, where
\[
\begin{gathered}{}
\partial_{\rm gr}=[T_0-I\ T_1-I\ T_2-I],\qquad
\\{}
j_0(x)=(x,0,0),\quad j_1(x)=(0,x,0),\quad
\\{}
j_2(x)=(0,0,x),\quad j_3(x)=-(T_1T_2x,T_2x,x).
\end{gathered}
\]
For $T_3x=x$, the fourth peripheral loop is the inverse of
the three-edge walk. Its successive transported vectors are
$x,T_2x,T_1T_2x$ and its terminal vector is
$T_0T_1T_2x=x$. This proves the formula for $j_3$ with its sign.

Write
$K=[j_0\ell\ j_1F_1\ j_2F_2\ j_3F_3]$, using the fixed
bases above; denote its columns by $t,L_1,\ldots,L_8$.
Here are three complement columns $C=(C_1,C_2,C_3)$ in the
standard basis $f_1,\ldots,f_{12}$ of the graph chain group.
\[
\begin{array}{c|ccc|r}
\text{row}&C_1&C_2&C_3&\det[C\ K]\\\hline
1&-f_{4}+f_{6}-f_{9}&f_{1}-f_{6}&-f_{6}+f_{8}&1\\
2&\substack{2f_{1}-f_{2}-f_{4}+f_{6}\\-f_{10}}&-f_{1}-f_{4}&f_{1}+f_{6}-f_{8}-f_{12}&1\\
3&\substack{-3f_{1}+4f_{2}-3f_{4}-f_{6}\\+f_{8}+f_{12}}&\substack{f_{1}-3f_{2}+2f_{4}+2f_{6}\\-f_{10}}&f_{12}&-1\\
4&\substack{-3f_{1}+3f_{2}+3f_{5}-3f_{6}\\+2f_{10}}&\substack{-2f_{1}+2f_{2}+f_{5}-2f_{6}\\+f_{10}}&-f_{5}+f_{8}&1\\
5&\substack{2f_{1}+f_{2}+f_{4}-2f_{5}\\+f_{6}}&f_{1}-f_{2}-f_{5}&f_{2}+f_{4}+f_{8}+f_{10}&1\\
6&\substack{f_{1}+2f_{4}-2f_{5}+f_{6}\\-3f_{8}+f_{10}+2f_{12}}&\substack{f_{2}+f_{4}-2f_{5}+2f_{6}\\-f_{10}}&-f_{5}+f_{6}+f_{8}&1\\
7&f_{10}&f_{1}-2f_{8}&f_{8}&-1\\
\end{array}
\]
Direct multiplication gives, with $\partial_{\rm gr}'$ the first three rows of $\partial_{\rm gr}$,
\begin{equation}
\partial_{\rm gr}'C=I_3,\qquad \partial_{\rm gr}K=0,\qquad
\det[C\ K]=\pm1,\qquad
[C\ K]^{-1}=\begin{pmatrix}\partial_{\rm gr}'\\L\end{pmatrix},\quad LK=I_9.
\hypertarget{tag-023H}{}\label{maps:eq:D.R1}
\end{equation}
Thus the inverse on the graph kernel is the last nine rows
of $\operatorname{adj}[C\ K]/\det[C\ K]$.
The coordinates of $j_0\delta$ are $(0,\mathbf c)^t$, where
\[
\begin{array}{c|c}
1,4,5,6,7&\mathbf c=(0,-1,0,-1,0,0,-1,0)^t\\
2&\mathbf c=(0,-1,0,0,-1,0,0,-1)^t\\
3&\mathbf c=(0,-1,0,0,-1,0,-1,0)^t.
\end{array}
\]

For the integral fibre quotient, order bivectors as
$(12,13,14,23,24,34)$ and choose
$s=a_s e_{12}+b_s e_{34}$ with
\[
(a_s,b_s)=(0,-1),(0,-1),(-1,-4),(0,-1),(1,3),(-1,-8),(2,1).
\]
Then $pa_s-qb_s=1$. Put
$\kappa=qe_{12}+pe_{34}$ and
$B_f=[s,\kappa,e_{13},e_{14},e_{23},e_{24}]$;
its inverse is integral and its first row is $\omega$.
Let
$E=[\bigwedge^2T_0-I\ \bigwedge^2T_1-I\ \bigwedge^2T_2-I]$
and let $W$ be the last five rows of $B_f^{-1}E$.
For a five-element ordered column set $I$, write $W_I$ for
that square submatrix and $\iota_I:\mathbb Z^5\to\mathbb Z^{18}$
for coordinate insertion. The following choices define an
integer right inverse
$Y=\sum_I a_I\iota_I\operatorname{adj}(W_I)$:
\[
\begin{array}{c|c|r|r}
\text{row}&I&a_I&\det W_I\\\hline
1&(2,3,4,5,11)&1&1\\
2&(2,3,4,5,15)&1&1\\
3&(1,2,3,5,11)&392&-56\\
3&(1,2,3,5,15)&-49&-462\\
3&(1,3,4,5,15)&1&-693\\
3&(2,3,4,5,11)&1&8\\
4&(2,3,4,5,9)&-1&-1\\
5&(1,3,4,5,9)&-1&-10\\
5&(2,3,4,5,9)&-3&3\\
6&(1,2,3,5,9)&3&90\\
6&(1,2,3,5,17)&1&-315\\
6&(1,3,4,5,9)&-1&-30\\
6&(2,3,4,5,9)&4&4\\
7&(1,2,3,5,17)&1&-2\\
7&(2,3,4,5,17)&1&3\\
\end{array}
\]
In each row $\sum_I a_I\det W_I=1$; consequently
$WY=I_5$ and $EY=[\kappa,e_{13},e_{14},e_{23},e_{24}]$.
Since $\omega E=0$, its image is exactly $\ker\omega$.
The fourth exterior relation also lies in this image, with
preimage $Y(B_f^{-1})_{2:6}(\bigwedge^2T_3-I)$.
The quotient is therefore $x\mapsto\omega(x)$, with inverse
$1\mapsto s$. 

The constant form $\omega$ on the flat family defines an
integral class $\Omega$: it has value one on $s$ and zero on
each actual fixed-loop sweep. The cellular complex and this
class give
\[
H_2(V)=\mathbb Z\langle s,t,L_1,\ldots,L_8\rangle.
\]
The fibre images of the three reduced vanishing lattices span
$\ker\chi=\mathbb Z\langle e_1,e_2,\delta\rangle$.
Explicitly, let $W_v$ have their four columns in cyclic order,
using $w_i$ and also $\delta$ at the hexagonal place. The
following $4\times3$ matrices give $W_vZ=(e_1,e_2,\delta)$.
\[
\begin{array}{c|c}
\text{row}&Z\\\hline
1&\begin{pmatrix}-1&-1&0\\4&-2&1\\2&1&0\\4&2&0\end{pmatrix}\\
2&\begin{pmatrix}21&7&2\\-30&-10&-3\\10&3&1\\0&0&0\end{pmatrix}\\
3&\begin{pmatrix}1&-1&0\\0&0&0\\5&2&1\\0&-1&0\end{pmatrix}\\
4&\begin{pmatrix}0&-1&0\\-12&-8&1\\-1&1&0\\0&0&0\end{pmatrix}\\
5&\begin{pmatrix}0&-1&0\\-8&-2&1\\1&-1&0\\0&0&0\end{pmatrix}\\
6&\begin{pmatrix}0&-1&0\\-12&-3&1\\1&-1&0\\0&0&0\end{pmatrix}\\
7&\begin{pmatrix}0&-1&0\\0&0&1\\1&-3&0\\0&0&0\end{pmatrix}\\
\end{array}
\]
Thus
$H_1(V)=\mathbb Z\langle\ell,x_0,x_1,x_2\rangle$,
with $x_3=-x_0-x_1-x_2$. The reduced relative two-discs map
to $x_1,x_2,x_3$, with matrix
\[
G_D=\begin{pmatrix}0&0&0\\0&0&-1\\1&0&-1\\0&1&-1\end{pmatrix}.
\]
On the base coordinates its inverse is
$(a,b,c)\mapsto(b-a,c-a,-a)$. Their relative three-cycles
map to $L_1,\ldots,L_8$. Hence attaching these fillings gives
\[
H_2(X_D)=\mathbb Z\langle s,t\rangle,\qquad
H_1(X_D)=\pi_1(X_D)=\mathbb Z\langle\ell\rangle.
\]
For the fundamental group, the four vanishing columns kill
exactly the abelian fibre subgroup $\ker\chi$; the three
section meridians and the base relation kill every base loop.
This gives the asserted group quotient as well as its abelianization.

\emph{The additive model and its attachment.}
At the additive place set $\lambda(x)=-\omega(x,\ell)$.
The quotient and section of $T_0-I$ are
\[
\pi_0(x)=(\lambda(x),\chi(x)),\qquad
(a,b)\longmapsto au+b\ell,
\]
where the vectors $u$ and matrices $Y_0$ below satisfy
$(T_0-I)Y_0=I-u\lambda-\ell\chi$.
\[
\begin{array}{c|c|c}
\text{row}&u&Y_0\\\hline
1&(0,0,1,0)^t&\begin{pmatrix}-1&-1&0&-1\\1&1&0&1\\0&0&0&0\\-1&-2&0&-1\end{pmatrix}\\
2&(0,0,1,0)^t&\begin{pmatrix}0&-1&0&-1\\1&0&0&0\\0&0&0&0\\0&-1&0&-1\end{pmatrix}\\
3&(0,1,11,0)^t&\begin{pmatrix}172&-269&25&-19\\-187&290&-26&21\\0&0&0&0\\143&-222&20&-16\end{pmatrix}\\
4&(0,0,1,0)^t&\begin{pmatrix}-5&-3&0&0\\7&4&0&0\\0&0&0&0\\0&0&0&0\end{pmatrix}\\
5&(0,-1,-3,0)^t&\begin{pmatrix}8&3&-1&0\\4&0&0&1\\0&0&0&0\\1&0&0&0\end{pmatrix}\\
6&(0,1,8,0)^t&\begin{pmatrix}-37&-12&2&0\\16&4&-2&1\\0&0&0&0\\-82&-28&3&1\end{pmatrix}\\
7&(-1,0,1,0)^t&\begin{pmatrix}0&1&0&1\\0&0&0&0\\0&0&0&0\\-1&2&-1&-1\end{pmatrix}\\
\end{array}
\]
This identity proves both directions of the quotient. Direct
summation gives
\begin{equation}
[\delta]=q[u],\quad
N_du=(d/q)\delta,\quad N_d\ell=d\ell,\quad
-\omega(u,\ell)=1,\quad -\omega(\ell,\delta)=-q,
\qquad N_d=\sum_{j=0}^{d-1}T_0^j.
\hypertarget{tag-023J}{}\label{maps:eq:D.R3}
\end{equation}
The bracketed identity is in the coinvariant quotient; its
relation preimage is $Y_0(\delta-qu)$. The norm identities
are equalities in $\Lambda$. The fixed lattice is
$\mathbb Z\langle\delta,\ell\rangle$.

For the positive root deck generator the linear part is
$R=T_0^{-1}$: the exterior mapping torus identifies
$(y,1)$ with $(T_0y,0)$. The elliptic coordinate changes
\[
H=(\begin{smallmatrix}1&-2\\0&1\end{smallmatrix}),\ I_2,\
(\begin{smallmatrix}6&-7\\1&-1\end{smallmatrix}),\
(\begin{smallmatrix}1&-2\\-1&3\end{smallmatrix}),\
(\begin{smallmatrix}0&-1\\1&2\end{smallmatrix}),\
(\begin{smallmatrix}1&-2\\-2&5\end{smallmatrix}),\
(\begin{smallmatrix}1&1\\0&1\end{smallmatrix})
\]
all have determinant one and identify the additive elliptic
blocks with the usual positive matrices
\[
II:\begin{pmatrix}0&1\\-1&1\end{pmatrix},\quad
III:\begin{pmatrix}0&1\\-1&0\end{pmatrix},\quad
IV:\begin{pmatrix}0&1\\-1&-1\end{pmatrix},\quad
IV^*:\begin{pmatrix}0&-1\\1&-1\end{pmatrix}.
\]
Their inverses are their adjugates. Together with $S$, these
maps carry both invariant periods and the section to the
good-reduction model of \cite[Proposition~3.1, pp.~8--9]{E26}.
The specialization of a section cocycle $g_i$ is
$-(A-I)^{-1}g_i\bmod\mathbb Z^2$, with inverse relation
$g_i=-(A-I)z_i$. This shows $Q'=0$ on the central curve;
$P'=0$ at $II$, while $P'$ is nonzero $2$-torsion at $III$
and nonzero $3$-torsion at $IV,IV^*$.
The birational comparison of $M$ with
$M'=\mathcal O(P'-O')$ identifies their linearizations off
the contracted curves. After contraction the missing locus is
a finite set in the smooth surface, so the line-bundle
isomorphism and its inverse extend by Hartogs. Since $Q'=0$
on the central curve, its restriction to $M'$ is a nonzero
scalar. Scaling $\phi$ makes its modulus less than one.
The flat multipliers of $M'$ and this scalar present the
central quotient as a complex two-torus with the stated
period lattice. This is the local extension argument in
\cite[Proposition~3.1 and Corollary~3.2, pp.~8--9]{E26},
with the displayed section specializations.

For $II,III,IV$ the good-reduction substitution is
$t=s^d$, $x=s^2X$, $y=s^3Y$; for $IV^*$ it is
$t=s^3$, $x=s^4X$, $y=s^6Y$.
The tangent multipliers at $O'$ are respectively $\zeta_d$
and $\zeta_3^2$. The forward action on the conormal is
their inverse. At $III,IV,IV^*$ the factor
$\mathcal O(P')$ is nonzero at $O'$, so this is also the
character of $\mathcal O(P'-O')|_{O'}$.
At $II$, both sections specialize to $O'$ and both Cartier
factors contribute. In a coordinate $z$ at $O'$,
\[
z_{\rm old}=s^{-k}z,\qquad
p_{\rm new}(s)=s^kp_{\rm old}(s^m),\qquad
\frac{z_{\rm old}}{z_{\rm old}-p_{\rm old}}
 =\frac{z}{z-p_{\rm new}}.
\]
This equality of Cartier frames holds also when
$p_{\rm old}(0)=0$. It cancels the two characters and poles.

Write $a$ for the pole order per root step, $\epsilon_{\rm line}$ for
the integral line winding used by the section, and
$c=d\epsilon_{\rm line}-a$. Their values are
\[
\begin{array}{c|rrrrrrr}
\text{row}&1&2&3&4&5&6&7\\\hline
a&0&0&1&0&1&1&2\\
\epsilon_{\rm line}&0&0&1&0&1&1&1\\
c&0&0&3&0&2&3&1.
\end{array}
\]
The natural root lift and the stated affine generator are
\[
\rho(y,s)=(Ry+(c/d)\delta,\zeta_m s),\qquad
g=t_{v/m}\rho,\qquad
g^m=t_{n\ell+(r+kc)\delta}.
\]
The value of $\chi$ changes by $n/m$, so no nonidentity
power fixes a point on the central torus when $\gcd(n,m)=1$.
The base rotation makes the action have order exactly $m$.
It is a finite free holomorphic action, and its quotient is
smooth and proper. The positive exterior meridian is
\[
x_0=t_{-\epsilon_{\rm line}\delta}g,\qquad
x_0^m=t_{v'},\qquad v'=n\ell+R_0\delta,\qquad
R_0=r-ak.
\]
The inverse parameter comparison is $r=R_0+ak$.

For the collar's radial and unit terms, write $s=\varrho e^{i\theta}$; the nonzero section factor is
$p(s^m)$, and let $b(t)=\log p(t)/(2\pi i)$ on a disc.
In flat real coordinates, with complex structure $J_s$, put
\[
u_p=\Re b\,\delta+\Im b\,J_s\delta,\quad
\mathcal B_{\rm col}(s)=\frac{\log\varrho}{2\pi}J_sv'-u_p,\quad
h_t(y,s)=y-\frac{\theta}{2\pi}v'+t\mathcal B_{\rm col}(s),\quad 0\le t\le1.
\]
The inverse adds $\theta v'/(2\pi)$ and subtracts $t\mathcal B_{\rm col}(s)$.
Here $\mathcal B_{\rm col}$ is periodic and $\mathcal B_{\rm col}(\zeta_ms)=R\mathcal B_{\rm col}(s)$, because
$J_{\zeta_ms}=RJ_sR^{-1}$ and $v',\delta$ are fixed.
Consequently $h_tg=Rh_t+\epsilon_{\rm line}\delta$ on the collar.
The integral last term gives the required deck compatibility.
At $t=1$ this is the holomorphic logarithmic attachment;
at $t=0$ it is the pure angular one. A radial cutoff on the
periodic residual $\mathcal B_{\rm col}$ gives the isotopy on $V$, with the
displayed inverse throughout. Its integral prism preserves
$\Omega$ and the graph norm. The positive normal circle is
therefore exactly $d_0=t_{-v'}x_0^m$ in the section marking.
Gluing these collars to the three reduced fillings gives
the claimed compact complex threefold for every stated datum.

\emph{The finite boundary and the integral extensions.}
Its central surface has fundamental group generated by
$\Lambda,x_0$, with $x_0\lambda x_0^{-1}=R\lambda$ and
$x_0^m=v'$. Choose $m\alpha+n\beta=1$ and set
$z=\alpha\ell+\beta x_0$. Abelianization gives
\[
H_1(G_0)=\mathbb Z\langle u,z\rangle,\quad
\ell=mz-qR_0\beta u,\qquad x_0=nz+qR_0\alpha u.
\]
The boundary map, its section, and its normal kernel are
\[
J_0=\begin{pmatrix}-qR_0\beta&1&qR_0\alpha\\m&0&n\end{pmatrix},
\qquad (u,z)\longmapsto(0,1,0)u+(\alpha,0,\beta)z,
\qquad\ker J_0=\mathbb Z(-n,-qR_0,m).
\]
The inverse coordinates after adjoining that kernel vector are
\[
(a,b,c)\longmapsto
(-qR_0\beta a+b+qR_0\alpha c,\ ma+nc,\ -\beta a+\alpha c).
\]
The normal character is zero on $\Lambda$ and $1/m$ on
$x_0$. It has the real lift $u\mapsto0$, $z\mapsto\beta/m$.
Explicitly, the normal line is smoothly trivialized by
\[
[(y,s)]\longmapsto([y],e^{-2\pi i\beta\chi(y)}s),\qquad
([y],w)\longmapsto[(y,e^{2\pi i\beta\chi(y)}w)].
\]
The factor under $x_0$ is
$e^{2\pi i(1-n\beta)/m}=e^{2\pi i\alpha}=1$;
lattice translations also give factor one. These maps are
inverse. Thom's isomorphism gives relative groups
$H_3=\mathbb Z^2$, $H_2=\mathbb Z$, $H_1=0$.

Here $N_m=kN_d=\sum_{j=0}^{m-1}R^j$.
With the relative orientation specified above, the angular
collar sends a fibre loop $\lambda$ to a normal torus of
area $-\omega(\lambda,v')$ and graph component
$j_0(N_m\lambda)$. For $x_0$, its $m$ lifted segments
concatenate to $v'$; the commuting torus
$x_0\wedge(mx_0-v')=v'\wedge x_0$ has $\Omega$-value zero.
Thus, using $j_0\delta=\mathbf c\cdot L$,
\[
\begin{aligned}
\tau(u)&=ns+(m/q)\mathbf c\cdot L,&
\tau(\ell)&=-qR_0s+mt,\\
\tau(x_0)&=nt+R_0\mathbf c\cdot L,&
\tau(z)&=-qR_0\alpha s+t+R_0\beta \mathbf c\cdot L.
\end{aligned}
\]
These respect
$m\tau(x_0)-n\tau(\ell)-qR_0\tau(u)=0$.
The full relative degree-three boundary, in bases
$(u,z,L_1,\ldots,L_8)$ and $(s,t,L_1,\ldots,L_8)$, is
\[
\mathcal F=\begin{pmatrix}
n&-qR_0\alpha&0\\0&1&0\\(m/q)\mathbf c&R_0\beta\mathbf c&I_8
\end{pmatrix}.
\]
Subtracting $(m/q)\mathbf c$ and $R_0\beta\mathbf c$ times the reduced
columns from its first and second columns, then adding
$qR_0\alpha$ times its second row to the first, gives
$\operatorname{diag}(n,1,I_8)$. Opposite additions give
the inverse changes. In the reduced quotient this is
\begin{equation}
B_{\rm rel}=\begin{pmatrix}n&-qR_0\alpha\\0&1\end{pmatrix},\qquad
[(a,b)]\longmapsto a+qR_0\alpha b\pmod{|n|},\qquad
1\longmapsto[s].
\hypertarget{tag-023K}{}\label{maps:eq:D.R4}
\end{equation}
A vector in the kernel of the displayed quotient has
preimage $((a+qR_0\alpha b)/n,b)$ under $B_{\rm rel}$.

The relative degree-two boundary in the basis
$(\ell,x_0,x_1,x_2)$ is
\[
\mathcal G=
\begin{pmatrix}-n&0&0&0\\m&0&0&-1\\0&1&0&-1\\0&0&1&-1\end{pmatrix},
\qquad\det\mathcal G=n.
\]
Its cokernel sends $\ell$ to $1\pmod{|n|}$ and the base
loops to zero. For $(a,b,c,d)$ with $n\mid a$, its preimage is
$t_0=-a/n$, $t_3=mt_0-b$, $t_1=c+t_3$, $t_2=d+t_3$.
Since $n\ne0$, this boundary is injective. The exact sequence
of $(X,V)$ therefore gives
$H_2(X)=\operatorname{coker}\mathcal F$ and
$H_1(X)=\operatorname{coker}\mathcal G$.

For the full Mayer--Vietoris map on $H_1$, the local quotient
sections just displayed give ten generators; adjoining the
four normal circles is an integral basis of its fourteen
boundary generators. If $S_1$ is the $4\times10$ matrix of
their images in $H_1(V)$, the map in these bases is
$\left(\begin{smallmatrix}S_1&\mathcal G\\-I_{10}&0\end{smallmatrix}\right)$.
Adding $S_1$ times the lower rows to the upper and then
permuting the rows gives
$\operatorname{diag}(-I_{10},\mathcal G)$; subtraction and
the inverse permutation undo it. This also proves directly
that the full map has zero kernel and cyclic cokernel.

\emph{The fundamental group and arbitrary affine vectors.}
After the three reduced fillings the group is
$\mathbb Z\langle\ell\rangle$, with every base loop killed.
The remaining normal circle $t_{-v'}x_0^m$ maps to
$\ell^{-n}$. The boundary surjects onto the finite filling
group and this circle generates its kernel, so van Kampen
gives $\pi_1(X)=\mathbb Z/|n|$.

For the endpoint concerning all affine choices, take the
given admissible generator $g_v=t_{v/m}\rho$, with the same
linear part and root degree, whose cyclic action is free
on the central torus. Its total period is
\[
w=\frac1m\sum_{j=0}^{m-1}R^jv+kc\delta\in\Lambda,
\qquad n=\chi(w)=\chi(v)\in\mathbb Z.
\]
This includes moving coordinates in $v$. Put
\[
\Pi=\frac1d\sum_{j=0}^{d-1}R^j,\qquad
B_0=\frac1d\sum_{j=0}^{d-1}jR^j
       -\frac{d-1}{2}\ell\chi.
\]
Telescoping gives $(R-I)B_0=I-\Pi$ and $\chi B_0=0$.
The translation $F_v=t_{B_0v/m}$, with inverse $t_{-B_0v/m}$,
intertwines the quotient actions:
\[
F_vg_vF_v^{-1}=t_{\Pi v/m}\rho.
\]
The real period vector has a holomorphic period image, so
this is a holomorphic map between the indicated root quotients.
Their boundary identification is $h_{\rm avg}\circ F_v$,
with inverse $F_v^{-1}\circ h_{\rm avg}^{-1}$, where
$h_{\rm avg}$ is the collar already constructed.
In this transported origin the exact section period is
\[
w_{\rm sec}=w-\epsilon_{\rm line}m\delta=\Pi v-ak\delta\in\Lambda,
\qquad d_0=t_{-w_{\rm sec}}x_0^m.
\]
It is fixed and integral because the original action has
the displayed integral total period. Its image after the
reduced fillings is $\ell^{-n}$. Van Kampen consequently
gives $\mathbb Z/n\mathbb Z$ for every admissible affine
choice in this construction, with $\mathbb Z/0\mathbb Z$
understood as $\mathbb Z$. Thus its simply connected members
have $n=\pm1$, which is automatically coprime to $m$.

\emph{The remaining homology and the smooth sphere.}
The finite torus quotient and both one-node fillings have
Euler characteristic zero. The hexagonal filling has Betti
numbers $(1,2,4,2,1)$ and Euler characteristic two. The
regular part and all boundaries have Euler characteristic
zero, so $\chi(X)=2$. The calculated finite $H_1,H_2$,
duality and this Euler characteristic give $b_3=0$.
The integral universal coefficient theorem now yields
\[
H_5=0,\quad H_4\cong\operatorname{Ext}(H_1,\mathbb Z)
 \cong\mathbb Z/|n|,\quad
H_3\cong\operatorname{Ext}(H_2,\mathbb Z)
 \cong\mathbb Z/|n|.
\]
More explicitly the characters sending $\ell$ and $s$ to
$1/|n|$ in $\mathbb Q/\mathbb Z$, followed by the integral
Bockstein and cap product with $[X]$, give the generators
of $H_4$ and $H_3$. The complex orientation gives
$H_0=H_6=\mathbb Z$. When $n=\pm1$, van Kampen and the
homology calculation make $X$ a simply connected integral
homology six-sphere. Hurewicz and Whitehead make it a
homotopy six-sphere, and
\cite[p.504, dimension-six entry, and p.505, first remark]{manKM63}
give its diffeomorphism with $S^6$.
\end{proof}

\subsection{Integral homology of smooth logarithmic modifications}

\begin{definition}[Four semistable logarithmic data]
\hypertarget{tag-Q117}{}\label{mapsIII:def:smooth-logarithmic-data}
{\normalfont\scriptsize tag Q117}\par
Use the four marked rational elliptic surfaces constructed from
rows $1,2,4,6$ of the Weierstrass models in the proof of
Theorem~\ref{maps:src:D.2.12}, with $G$ a Mordell--Weil generator:
\[
\begin{array}{c|l|r}
\mathrm{OS}&\text{ordered fibres}&p\\\hline
45&I_8,I_1,I_1,I_1,I_1^\bullet&8\\
47&I_7,I_2,I_1,I_1,I_1^\bullet&14\\
55&I_5,I_4,I_1,I_1,I_1^\bullet&20\\
56&I_5,I_3,I_2,I_1,I_1^\bullet&30.
\end{array}
\]
Use $P=pG$, $Q=G$, and the degree-one linearization with zero
at $I_1^\bullet$. Fill the four unmarked singular fibres with
primitive one-node neighbourhoods and $I_1^\bullet$ with a
one-component hexagonal cusp. In the period lattice
$\Lambda=\mathbb Z\langle e_1,e_2,\delta,\gamma\rangle$,
$\delta$ is the positive line circle and $\gamma$ is the forward
period of the linearization covering $-Q$.
At an additional smooth fibre, perform the logarithmic transformation of order
$m\in\mathbb Z_{>0}$ and translation vector
\begin{equation}
 v=a e_1+b e_2+c\delta+n\gamma,\quad a,b,c,n\in\mathbb Z,\qquad
 \gcd(m,a,b,c,n)=1.
 \hypertarget{tag-023L}{}\label{maps:eq:D.S1}
\end{equation}
Put
\begin{equation}
 t=\gcd(m,n,a,b)\,\gcd(m,n,p,c).
 \hypertarget{tag-023R}{}\label{maps:eq:D.S2}
\end{equation}
\end{definition}

\begin{theorem}[Homology of the four smooth logarithmic families]
\hypertarget{tag-02KN}{}\label{maps:src:D.2.13}
{\normalfont\scriptsize tag 02KN}\par
The logarithmic transformation of
Definition~\ref{mapsIII:def:smooth-logarithmic-data} is free.
The two factors in \eqref{maps:eq:D.S2} are coprime, and $t\mid m$.
For $n\ne0$, also $t\mid n$, and the result has
\begin{equation}
\begin{aligned}
 \pi_1(X)&=H_1(X)=\mathbb Z/|n|,\qquad H_5(X)=0,\\
 H_2(X)&=H_3(X)=G_{m,n,t},\qquad H_4(X)=\mathbb Z/|n|,\\
 G_{m,n,t}&=\mathbb Z/t\oplus\mathbb Z/m
               \oplus\mathbb Z/(m|n|/t).
\end{aligned}
\hypertarget{tag-023S}{}\label{maps:eq:D.S3}
\end{equation}
For $n=0$, the answer is instead
\begin{equation}
\begin{aligned}
 \pi_1(X)&=H_1(X)=H_5(X)=\mathbb Z,\\
 H_2(X)&=H_3(X)=\mathbb Z^2\oplus\mathbb Z/t\oplus\mathbb Z/m,
 \qquad H_4(X)=\mathbb Z^2.
\end{aligned}
\hypertarget{tag-023T}{}\label{maps:eq:D.S4}
\end{equation}
In both cases $H_0=H_6=\mathbb Z$ and $\chi(X)=2$.
The simply connected members are precisely $n=\pm1$; they
have $H_2=H_3=(\mathbb Z/m)^2$. Consequently the sphere
members in these four families are precisely $m=1$, $n=\pm1$.
The members with $m>1$ and $n=\pm1$ are simply connected complex
rational homology spheres whose two middle torsion groups are
$(\mathbb Z/m)^2$.
\end{theorem}


\begin{proof}
In the marking of Definition~\ref{mapsIII:def:smooth-logarithmic-data},
put $\chi=\gamma^*$ and
$J=\left(\begin{smallmatrix}0&1\\-1&0\end{smallmatrix}\right)$.

\emph{The marked elliptic surfaces.}
We construct the four five-fibre surfaces from rows $1,2,4,6$ of
the Weierstrass models in the proof of Theorem~\ref{maps:src:D.2.12}.
Write their coefficients as $F_0,H_0$, their equation as
$y^2=4x^3-F_0(t)x-H_0(t)$, and their discriminant as
$\Delta_0=F_0^3-27H_0^2$. The data needed for the deformation are
\[
\begin{array}{c|c|c|c|c|l}
\mathrm{OS}&N&b_{\rm site}&B&a_{\rm site}&\Delta_0\\\hline
45&8&--&0&0&-11664t^2(3t^2-14t+27)\\
47&7&-8/9&2&0&-432t^2(4t+9)(9t+8)^2\\
55&5&0&4&1/8&-114791256t^4(t+10)(8t-1)^2\\
56&5&0&3&1/5&-7464960000t^3(3t+25)(5t-1)^3.
\end{array}
\]
Here $N$ is the order at infinity, and $b_{\rm site},B$ specify a finite
multiplicative fibre to preserve. The fibre at $a_{\rm site}$ is $II$ in the
first three rows and $III$ in the fourth.

For $F=\sum_{i=0}^4 f_it^i$ and $H=\sum_{i=0}^6 h_it^i$ impose
\[
\begin{gathered}{}
 [t^j](F^3-27H^2)=0\quad(13-N\leq j\leq12),\qquad
\\{}
 (F^3-27H^2)^{(j)}(b_{\rm site})=0\quad(0\leq j<B),\qquad
\\{}
 f_4=f_4(0).
\end{gathered}
\]
The finite conditions are absent in row $45$, where we also fix
$f_3=f_3(0)$. In the first three rows impose
$F(a_{\rm site})=h^2$, $H(a_{\rm site})=0$. In row $56$ impose
\[
 F(a_{\rm site})=12h^2,\qquad H(a_{\rm site})=-8h^3,\qquad H'(a_{\rm site})+hF'(a_{\rm site})=0.
\]
These are twelve equations for twelve coefficients. Order their
variables as $(f_0,\ldots,f_4,h_0,\ldots,h_6)$ and their equations
as displayed, with the conditions at $a_{\rm site}$ before the fixed leading
coefficients. Since
$d\Delta=3F_0^2dF-54H_0dH$, the four Jacobian determinants are
\[
 -2^{30}3^{27},\qquad 2^{50}3^{31}7^5,\qquad
 -2^{66}3^{48},\qquad -2^{27}3^{28}5^{41}.
\]
Their inverses are their adjugates divided by these determinants.
The holomorphic inverse function theorem gives convergent coefficient
germs satisfying all the equations. The second base chart is
$s=t^{-1}$, $X=s^2x$, $Y=s^3y$, with inverse
$t=s^{-1}$, $x=s^{-2}X$, $y=s^{-3}Y$.
The nonzero leading discriminant coefficient and $F_0(b_{\rm site})$ preserve
exactly $I_N$ at infinity and $I_B$ at $b_{\rm site}$.

At a type-$II$ point put ${d_{\rm disc}}=H_0'(a_{\rm site})$; its values are
$-108,96,-6561/8$. In the coordinate $t=a_{\rm site}+h^3z$ the scaled
discriminant tends to $1-27{d_{\rm disc}}^2z^2$. Its two roots are simple, and
$F_h$ at either continued root is $h^2(1+O(h))$.
Thus the point splits into two $I_1$ fibres. For row $56$ put
$A=F_0'(a_{\rm site})=-5760$. Direct differentiation gives
\[
 \Delta''(a_{\rm site})=18h^2\bigl(F'(a_{\rm site})^2+24h^2F''(a_{\rm site})+24hH''(a_{\rm site})\bigr).
\]
Consequently $a_{\rm site}$ has exactly multiplicity two for small nonzero $h$.
After writing $\Delta_h(t)=(t-a_{\rm site})^2D_h(t)$, the other root is
$a_{\rm site}-9h^2/A+O(h^3)$, where $F_h=3h^2+O(h^3)$.
The two fibres are $I_2$ and $I_1$. Every simple root of $\Delta_0$ continues;
the displayed multiplicities exhaust the discriminant degrees.
The same minimal Weierstrass criterion as in
Theorem~\ref{maps:src:D.2.12} gives rational elliptic surfaces with
the zero section $[0:1:0]$ and the four stated configurations.

Here is their integral marking. For a primitive column $u$ write
$\mathsf K_{\rm PL}(u)=I+uu^tJ$. At the five ordered places set
$A_i=\mathsf K_{\rm PL}(u_i)^{\nu_i}$ and let $g_i=c_i u_i$ be the generator
section cocycle. All inputs are given by
\[
\begin{array}{c|c|c|c}
\mathrm{OS}&(u_1,u_2,u_3,u_4,u_5)&(\nu_1,\ldots,\nu_5)&(c_1,\ldots,c_5)\\\hline
45&((1,0),(1,0),(2,-1),(1,-2),(1,1))&(8,1,1,1,1)&(3,0,1,1,0)\\
47&((1,0),(1,-1),(1,-3),(1,1),(1,0))&(7,2,1,1,1)&(5,1,0,1,1)\\
55&((1,0),(1,0),(2,-1),(1,-2),(1,1))&(5,4,1,1,1)&(-3,-3,-1,0,-1)\\
56&((1,0),(1,0),(1,-1),(1,-3),(1,1))&(5,3,2,1,1)&(-4,-2,-1,0,-1).
\end{array}
\]
The pairs in the second column are column vectors. To identify this
marking with the actual deformation, transport the zero section and
the elliptic local system outside disjoint discs around the singular
fibres; reverse transport is the inverse. Inside a splitting disc the
two positive factors have the form $\mathsf K_{\rm PL}(v)^j\mathsf K_{\rm PL}(w)$, with $j=1$ for
$II$ and $j=2$ for $III$. Their trace is
$2-j\det(v,w)^2$, so their vanishing vectors form a unimodular basis.
In such a basis the product is
$E_j=\left(\begin{smallmatrix}1&j\\0&1\end{smallmatrix}\right)
\left(\begin{smallmatrix}1&0\\-1&1\end{smallmatrix}\right)$.
An integral matrix centralizing $E_j$ is
$\alpha_{\rm cen}I+\beta_{\rm cen}E_j$, and determinant one gives
$\alpha_{\rm cen}^2+(2-j)\alpha_{\rm cen}\beta_{\rm cen}+\beta_{\rm cen}^2=1$.
The integer solutions are powers
of $E_j$ up to sign. Supported full twists of the splitting disc,
whose inverses are the opposite twists, therefore match the local
ordered pair while fixing the exterior marking.

For clarity the positive based-loop move and its inverse are
\[
 (x,y)\longmapsto(y,y^{-1}xy),\qquad
 (x',y')\longmapsto(x'y'x'^{-1},x').
\]
On cocycles of $(U,V)$ the move is
$(p,q)\mapsto(q,V^{-1}(p+(U-I)q))$, with inverse
$(p',q')\mapsto(Vq'-(U-I)p',p')$.
A merger sends $(p,q)$ to $p+Uq$; its inverse on saturated
rank-one coordinates is $[v,Uw]^{-1}$. Numbering places from zero,
the four tuples reach the actual additive rows $1,2,4,6$ by
\[
\begin{array}{c|c|c|c}
\mathrm{OS}&\text{positive moves}&\text{merged pair}&[v,Uw]\\\hline
45&\text{none}&1,2&\left(\begin{smallmatrix}1&1\\0&-1\end{smallmatrix}\right)\\
47&\text{none}&3,4&\left(\begin{smallmatrix}-1&0\\-1&-1\end{smallmatrix}\right)\\
55&1,1&2,3&\left(\begin{smallmatrix}2&-1\\-3&1\end{smallmatrix}\right)\\
56&1,1&2,3&\left(\begin{smallmatrix}2&-3\\-5&7\end{smallmatrix}\right).
\end{array}
\]
Each merger determinant has absolute value one. Matrix multiplication
and the displayed affine cocycle map give exactly the additive
tuples of Theorem~\ref{maps:src:D.2.12}; reversing the moves
supplies the five-place marking.

We also identify the actual primitive section. The stacked matrices
$A_i-I$ have an integral left inverse. Thus they have no common
nonzero fixed element in $(\mathbb Q/\mathbb Z)^2$, and a torsion
section must vanish. The root lattice has rank seven and determinant
$p$, whereas the rational surface's N\'eron--Severi lattice is unimodular
of rank ten. The primitive rank-one Mordell--Weil lattice consequently
has height $1/p$. The section-cocycle complex is
\[
 \mathbb Z^2\xrightarrow{\,h_{\rm cob}\mapsto((A_i-I)h_{\rm cob})_i\,}
 \bigoplus_{i=1}^5\operatorname{sat}\operatorname{im}(A_i-I)
 \xrightarrow{\,(z_i)\mapsto z_1+A_1z_2+\cdots+A_1\cdots A_4z_5\,}
 \mathbb Z^2.
\]
Its middle cohomology is $\mathbb Z$, and the displayed $g$ has
primitive coordinate $1,-1,1,1$ in the four rows. Euclidean integer
row and column reduction, with the inverse elementary operations,
gives the decomposition $c=kg+dh_{\rm cob}$ of every integral cocycle.
The Poincar\'e-line and boundary-degree identification proved in
Theorem~\ref{maps:src:D.2.12} gives
$\langle kg+dh_{\rm cob},\ell g+dj\rangle=k\ell/p$ for these tuples.
The actual section of height $1/p$ therefore has $k=\pm1$.
After choosing its sign, its exact change of four-period coordinates is
\[
 S_{h_{\rm cob}}=\begin{pmatrix}I_2&0&h_{\rm cob}\\-ph_{\rm cob}^tJ&1&0\\0&0&1\end{pmatrix},
 \qquad
 S_{h_{\rm cob}}^{-1}=\begin{pmatrix}I_2&0&-h_{\rm cob}\\ph_{\rm cob}^tJ&1&0\\0&0&1\end{pmatrix}.
\]
The sign change uses $\operatorname{diag}(-I_2,1,1)$.
These maps identify the line circle and forward period, including
their central coefficients. In the geometric component labels a raw
residue $c_i$ is $-c_i$; negating labels and its inverse identify
translation by $-Q$ with the raw forward convention. All residues
at higher multiplicative fibres are units.

For use below, the actual four-period matrices are completely
determined by these inputs. Solve
$(A_i-I)h_{i,\rm loc}=pg_i$, $(A_i-I)z_i=g_i$ over $\mathbb Q$, and put
$b_i=-h_{i,\rm loc}^tJ$, $\beta_i=-pg_i^tJA_i$ and
$w=(0,0,0,0,1)$. Then
\[
 T_i=\begin{pmatrix}A_i&0&g_i\\\beta_i&1&\beta_i z_i+w_i\\0&0&1\end{pmatrix},
 \qquad
 S_i=\begin{pmatrix}I_2&0&z_i\\b_i&1&b_iz_i\\0&0&1\end{pmatrix},
 \qquad
 S_i^{-1}=\begin{pmatrix}I_2&0&-z_i\\-b_i&1&0\\0&0&1\end{pmatrix}.
\]
The formula for $T_i$ is independent of the kernel choices in the two
local equations. We have
$T_i=S_i^{-1}(\operatorname{diag}(A_i,1,1)+w_iE_{34})S_i$,
each $T_i$ is integral, and $T_1\cdots T_5=I$.
The mixed height is one, so the linearization has its simple zero
at the marked $I_1$. Moreover $\deg M|_O=p/2$. Dividing the
linearization there by its reduced fibre gives an invertible local
map $U_{\rm lin}$; $p/2$ inverse iterates carry $t^{p/2}v$ to
$U_{{\rm lin},p/2-1}^{-1}\cdots U_{{\rm lin},0}^{-1}v$. Forward iteration is the inverse.
This fixes the reduced meridian framing. The primitive one-node
and hexagonal local models are therefore those of
Theorem~\ref{maps:src:D.2.12}. The additional smooth place has zero
section cocycles and its unshifted section frame.

\emph{The smooth logarithmic filling.}
On the root disc $t=s^m$ its generator is
$g_{\rm log}(y,s)=(y+v/m,e^{2\pi i/m}s)$. The least positive integer $j$
for which $jv/m\in\Lambda$ is $m/\gcd(m,a,b,c,n)$.
This proves the exact freeness assertion on the central torus;
off it the base rotation is free. The period image of $v$ is
holomorphic, so this finite free quotient is a smooth proper
holomorphic filling.

Write $s=\varrho e^{i\theta}$, let $J_s$ be the complex structure
in real period coordinates, and let $q_{\rm sec}(s^m)$ denote the nonzero
section factor in the smooth frame. Choose its logarithm and set
\[
 b_{\rm sec}(s)=\frac{\log q_{\rm sec}(s^m)}{2\pi i},\quad
 u_{\rm sec}=\Re b_{\rm sec}\,\delta+\Im b_{\rm sec}\,J_s\delta,\quad
 B_{\rm col}(s)=\frac{\log\varrho}{2\pi}J_sv-u_{\rm sec}.
\]
The maps
$h_\tau(y,s)=y-\theta v/(2\pi)+\tau B_{\rm col}(s)$, $0\leq\tau\leq1$,
with inverse $y\mapsto y+\theta v/(2\pi)-\tau B_{\rm col}(s)$, identify the
holomorphic logarithmic collar with the angular collar. The residual
$B_{\rm col}$ is periodic and deck invariant. A radial cutoff on it extends
this isotopy over the exterior. Its normal circle is thus exactly
$d_0=mx_0-v$, where $x_0$ is the positive base loop.

The central torus has lattice
\begin{equation}
 \Lambda_v=\Lambda+\mathbb Z(v/m)
   \cong\mathbb Z^5/\mathbb Zq_0,
 \qquad q_0=(-a,-b,-c,-n,m)^t.
 \hypertarget{tag-023U}{}\label{maps:eq:D.S5}
\end{equation}
The quotient map sends $(z,k)$ to $z+kv/m$.
Choose $\alpha\in\mathbb Z$ and an integral covector $\ell$ with
$\alpha m+\ell(v)=1$. For $r_0=(-\ell,\alpha)$ the matrix
$P_0=I-q_0r_0$ has image $\ker r_0$, and the maps
$[z]\mapsto P_0z$ and $u\mapsto[u]$ are inverse.
The normal line has the explicit trivialization
\[
 [(y,s)]\longmapsto([y],e^{-2\pi i\ell(y)}s),\qquad
 ([y],w)\longmapsto[(y,e^{2\pi i\ell(y)}w)].
\]
Both maps descend: the deck factor is
$e^{2\pi i(1-\ell(v))/m}=e^{2\pi i\alpha}=1$.
Thom's isomorphism gives local relative groups
$H_3=\Lambda_v$, $H_2=\mathbb Z$, $H_1=0$.
We orient a relative three-generator as minus the fibre loop times
the positive normal disc.

\emph{The integral exterior and reduced fillings.}
Let $V$ be the regular torus bundle over the sphere with the five
singular-fibre discs and the smooth logarithmic disc removed.
Insert $T_0=I$ before the five matrices, and put
$D=[T_0-I\ \cdots\ T_4-I]$, with first three rows $D_0$.
For $i<5$ let $j_i$ be its four-coordinate block inclusion, and set
\[
 j_5(x)=-(T_1T_2T_3T_4x,T_2T_3T_4x,T_3T_4x,T_4x,x).
\]
Then $Dj_5(x)=(I-T_5^{-1})x$.
Choose integral unimodular ${\mathsf H}_i$, with ${\mathsf H}_ie_3=\delta$, such that
\[
 {\mathsf H}_i^{-1}T_i{\mathsf H}_i=I+E_{12}\quad(1\leq i\leq4),\qquad
 {\mathsf H}_5^{-1}T_5{\mathsf H}_5=I+E_{12}+E_{34}.
\]
Integer reduction of the specified $T_i-I$, keeping the displayed
primitive $\delta$ column, supplies these bases: the nonzero
elementary divisors are respectively $1$ and $(1,1)$, and at the
last place $\delta$ is a primitive element of the image.
The fixed-loop bases are
${\mathsf F}_i=({\mathsf H}_ie_1,{\mathsf H}_ie_3,{\mathsf H}_ie_4)$ and
${\mathsf F}_5=({\mathsf H}_5e_1,{\mathsf H}_5e_3)$.
Let ${\mathsf B}_{\rm gr}=[j_1{\mathsf F}_1\ \cdots\ j_5{\mathsf F}_5]$. Integer elimination on the
four displayed tuples gives
\begin{equation}
 K=[{\mathsf B}_{\rm gr}\ j_0(e_1)\ j_0(e_2)\ j_0(\gamma)],\qquad
 \operatorname{im}D=\ker\chi,\qquad
 \ker D=\operatorname{im}K.
 \hypertarget{tag-023V}{}\label{maps:eq:D.S6}
\end{equation}
Here $D_0$ has three unit elementary divisors, and $K$ has seventeen.
These finite reductions use the inputs $u_i,\nu_i,c_i,w_i$ and the
formula for $T_i$ above. To construct the integral inverse, choose
integer coefficients $a_I$ with
$\sum_Ia_I\det K_I=1$. The map
$L=\sum_Ia_I\operatorname{adj}(K_I)\operatorname{pr}_I$ satisfies
$LK=I$. Since the rational kernel of $D$ has rank seventeen,
$KL$ is identity on its integral kernel. The same minor construction
gives an integral right inverse $C$ of $D_0$.
The last three coordinates of $Lj_0(x)$ are $(x_1,x_2,x_4)$;
the first fourteen coordinates of $Lj_0(\delta)$ are
\[
 c_\delta=(0,-1,0,0,-1,0,0,-1,0,0,-1,0,0,-1)^t
\]
for the local bases specified by the reduced models. Equivalently,
the global $\delta$-circle over the punctured base gives
$j_0\delta+j_1\delta+\cdots+j_5\delta=0$ in the graph kernel.
Thus the smooth $\delta$ sweep is the negative sum of the reduced
ones, independent of the chosen fixed-loop bases.

In exterior-square order $(12,13,14,23,24,34)$ put
\[
 \omega=(p,0,0,0,0,-1),\quad s=-e_{34},\quad
 {\mathsf F}_{\rm ex}=(e_{12}+pe_{34},e_{13},e_{14},e_{23},e_{24}).
\]
For ${\mathsf B}_{\rm ex}=[\bigwedge^2T_0-I\ \cdots\ \bigwedge^2T_4-I]$
integer elimination gives an integral $Z$ with ${\mathsf B}_{\rm ex}Z={\mathsf F}_{\rm ex}$.
As $\omega {\mathsf B}_{\rm ex}=0$, $\omega s=1$ and $[s,{\mathsf F}_{\rm ex}]$ is unimodular,
the fibre quotient maps are $[x]\mapsto\omega x$ and
$k\mapsto[ks]$. The section and fibre-first sweep chains give
the splitting $H_2(V)=\mathbb Zs\oplus\mathbb Z^{17}$; torus
addition uses the triangle with vertices $0,x,x+y$ in real period
space, with inverse coordinate subtraction.
The six vanishing columns
$({\mathsf H}_1e_1,\ldots,{\mathsf H}_4e_1,{\mathsf H}_5e_1,{\mathsf H}_5e_3)$ have an integral right
inverse onto $\langle e_1,e_2,\delta\rangle$.
Their normal closure therefore kills precisely $\ker\chi$ in
the abelian fibre group. The section kills all reduced base meridians.
After these five fillings, the fundamental group is
$\mathbb Z\langle\gamma\rangle$ and the degree-two group has basis
$(s,t_{e_1},t_{e_2},t_\gamma)$.

\emph{The logarithmic boundary and its quotient.}
The angular collar sends a loop $\lambda\in\Lambda$ to
$-\omega(\lambda,v)s+mj_0(\lambda)$ and sends $x_0$ to $j_0(v)$.
The relation $mx_0=v$ is respected because $\omega(v,v)=0$.
Thus the reduced transfer, in source order
$(e_1,e_2,\delta,\gamma,x_0)$, is
\begin{equation}
 R=\begin{pmatrix}
 -pb&pa&n&-c&0\\m&0&0&0&a\\0&m&0&0&b\\0&0&0&m&n
 \end{pmatrix},\qquad Rq_0=0.
 \hypertarget{tag-023W}{}\label{maps:eq:D.S7}
\end{equation}
Before eliminating the reduced normal tori, the map is
$\left(\begin{smallmatrix}R&0\\C_{\rm low}&I_{14}\end{smallmatrix}\right)$.
Right multiplication by
$\left(\begin{smallmatrix}I_5&0\\-C_{\rm low}&I_{14}\end{smallmatrix}\right)$
gives $R\oplus I_{14}$; changing the minus sign to plus gives
the inverse change of source. This accounts for the entire relative
degree-three image.

Write $d_1=\gcd(m,n,a,b)$, $d_2=\gcd(m,n,p,c)$.
A common prime divisor would divide all five entries of the freeness
condition, so $\gcd(d_1,d_2)=1$ and $t=d_1d_2$ divides $m,n$.
We compute the determinantal ideals over each $\mathbb Z_{(q)}$.
When $m$ is a unit, pivoting on its three entries and replacing the
last column by itself minus $a/m,b/m,n/m$ times the corresponding
columns leaves the ideal $(n)$. When $n$ is a unit, its top-row
pivot and last-column pivot leave $m,m$.

When $q\mid m,n$, freeness makes $a$, $b$ or $c$ a unit.
In the first two cases respectively use
\[
 C_a=\begin{pmatrix}1/a&0\\-b&a\end{pmatrix},\qquad
 C_b=\begin{pmatrix}0&1/b\\-b&a\end{pmatrix}.
\]
These determinant-one maps send $(a,b)^t$ to $(1,0)^t$;
their inverses are their adjugates. Multiplication of $R$ by
$\operatorname{diag}(1,C_a,1)$ on the left and
$\operatorname{diag}(C_a^{-1},1,1,1)$ on the right, or the analogous
$C_b$ maps, preserves $p$ and gives the case $a=1,b=0$.
Pivoting on its last-column unit leaves
\[
 \begin{pmatrix}0&p&n&-c\\0&m&0&0\\-mn&0&0&m\end{pmatrix}.
\]
Its entry, two-minor and maximal-minor ideals are
$(m,n,p,c)$, $m(m,n,p,c)$ and $(m^2n)$.
When $c$ is a unit, pivot on $-c$, multiply the last remaining
row by $c$, add $pb$ times the first and subtract $pa$ times the
second. The result is
\[
 \begin{pmatrix}m&0&0&a\\0&m&0&b\\0&0&mn&cn\end{pmatrix}.
\]
Its corresponding ideals are $(m,n,a,b)$, $m(m,n,a,b)$ and
$(m^2n)$. Opposite row and column additions and inverse unit
multiplications undo every operation. These cases give
\begin{equation}
 (1,t,m,m|n|/t)\quad(n\ne0),\qquad
 (1,t,m,0)\quad(n=0)
 \hypertarget{tag-023X}{}\label{maps:eq:D.S8}
\end{equation}
as the Smith entries. In the latter case the matrix has rank three.

The isomorphism of the central lattice and each quotient is explicit.
Integer reduction of $q_0$ gives ${\mathsf U}_0q_0=e_1$. Write $r'_0$ for
its first row. If ${\mathsf H}_{\rm cen}$ is its last four rows and $E$ the last four
columns of ${\mathsf U}_0^{-1}$, then
${\mathsf H}_{\rm cen}E=I_4$, ${\mathsf H}_{\rm cen}q_0=0$ and $E{\mathsf H}_{\rm cen}+q_0r'_0=I_5$.
The inverse central maps are $[z]\mapsto {\mathsf H}_{\rm cen}z$ and $u\mapsto[Eu]$.
For ${\mathsf U}_R RE{\mathsf V}_R$ equal to the diagonal matrix in \eqref{maps:eq:D.S8},
the target quotient maps are $[x]\mapsto[{\mathsf U}_R x]$ and
$[y]\mapsto[{\mathsf U}_R^{-1}y]$. The Euclidean row and column operations
used above determine ${\mathsf U}_R,{\mathsf V}_R$ and their inverse operations integrally.

\emph{The fundamental group and the free relative term.}
The smooth boundary surjects onto $\Lambda_v$, with normal kernel
$\mathbb Z(mx_0-v)$. After the reduced fillings and the base relation,
this kernel maps to $-n\gamma$. Van Kampen gives
$\pi_1(X)=\mathbb Z/n\mathbb Z$, including $\mathbb Z$ at $n=0$.
In target order $(\gamma,x_0,x_1,x_2,x_3,x_4)$, the six normal-disc
boundaries are $(-n,m,0,0,0,0)^t$, the four unit $x_i$ columns,
and $(0,-1,-1,-1,-1,-1)^t$. Their determinant is $n$.
Writing this six-by-six matrix as ${\mathsf G}_{\rm rel}$, the local quotient sections
also identify the full Mayer--Vietoris map on $H_1$ with
$\left(\begin{smallmatrix}S&{\mathsf G}_{\rm rel}\\-I_{18}&0\end{smallmatrix}\right)$.
Adding $S$ times the lower rows to the upper rows and permuting
the rows gives $\operatorname{diag}(-I_{18},{\mathsf G}_{\rm rel})$; subtraction and
the inverse permutation undo these operations. Thus the reduced
relative calculation preserves both kernel and cokernel of the full map.
For $n\ne0$ there is no relative kernel, so $H_2(X)=\operatorname{coker}R$.
For $n=0$ its primitive kernel is $(1,m,m,m,m,m)$.
In the stated boundary order, let $\mathcal D_i$ denote the normal discs.
The chain $\mathcal D_0+m(\mathcal D_1+\cdots+\mathcal D_5)$ has boundary $-v$.
The monodromy cylinders specified by $C(a,b,c)^t$, with the
addition triangles above, bound $v$. Their sum is a closed
two-cycle mapping to that primitive relative generator and therefore
splits this free quotient.
The other free quotient of $\operatorname{coker}R$ is
\[
 (x_s,x_1,x_2,x_\gamma)\longmapsto
 \frac{mx_s+pbx_1-pax_2+cx_\gamma}{d_2}.
\]
Its coefficients have gcd one: freeness gives
$\gcd(m,pb,pa,c)=\gcd(m,p,c)=d_2$.
A B\'ezout preimage of one gives its section. These two sections
establish the full $n=0$ formula for $H_2$.

\emph{The remaining homology.}
The regular torus family, its boundaries, the multiple smooth torus
and the four primitive one-node fillings have Euler characteristic
zero. The one-component hexagonal fibre has Euler characteristic two,
by its normalization and conductor calculation. Hence $\chi(X)=2$.
The complex orientation gives $H_0=H_6=\mathbb Z$.
For $n\ne0$, duality and the integral universal coefficient sequence
give $H_5=0$, $H_4=\mathbb Z/|n|$, free middle rank zero, and
$H_3\cong\operatorname{Ext}(H_2,\mathbb Z)$.
For every Smith summand $\mathbb Z/d$, the character taking its
generator to $1/d\in\mathbb Q/\mathbb Z$, followed by the integral
Bockstein and cap product with $[X]$, gives its torsion generator
in $H_3$. Reading the linking coefficient is the inverse map.

For $n=0$ we have $H_5=\mathbb Z$, $H_4=\mathbb Z^2$ and
free middle rank two. Here are explicit choices for its splitting.
Put $d=\gcd(m,a,b)$, choose $cA+dB=1$, and in the central
five-generator presentation set
\[
 u=(-a/d,-b/d,0,0,m/d)^t,\qquad
 \lambda=Au+Be_\delta.
\]
The kernel of $R$ before its source relation has basis
$(u,e_\delta)$, and $q_0=du-ce_\delta$.
Thus $[\alpha u+\beta e_\delta]\mapsto c\alpha+d\beta$
is the kernel coordinate, with inverse $1\mapsto[\lambda]$.
In $\Lambda_v$, $\lambda=\delta/d$.
Its complete relative vector is
$(\lambda,-(m/d)c_\delta)$.
The fibre part of its boundary is
$\kappa=(a/d)e_{13}+(b/d)e_{23}$, and its exact bounding
monodromy cylinders are specified by
$Z(0,a/d,0,b/d,0)^t$.
If $\Sigma$ is the sphere with the six discs removed, the product
$\delta\times\Sigma$ has boundary
$j_0\delta+\cdots+j_5\delta$ with the fibre-first convention.
Subtracting the stated fibre bounding chain and
$(m/d)(\delta\times\Sigma)$ from the relative Thom chains,
and using the addition triangles, gives a closed three-cycle $Y$.
Its relative coefficient homomorphism $f_{\rm rel}:H_3(X)\to\mathbb Z$
satisfies $f_{\rm rel}(Y)=1$.

Choose a finite oriented cellular decomposition refining these chains.
For ${\mathsf U}_3\partial_3{\mathsf V}_3={\mathsf D}_3$ of rank $r$, put
$K_3=({\mathsf V}_3)_{[:,r+1:\,]}$ and $N_3=({\mathsf V}_3^{-1})_{[r+1:\,,\,]}$.
Then $K_3$ includes the full integral cycle lattice and $N_3K_3=I$.
The cochain $\eta_3=(f_{\rm rel}K_3)N_3$ agrees with $f_{\rm rel}$ on cycles and
vanishes on $\partial_4$, so it is an integral cocycle with
$\eta_3(Y)=1$. Its dual-cell cycle $Z_3=\eta_3\cap[X]$ has
$I(Y,Z_3)=1$, for the convention
$I(A,B)=\operatorname{PD}^{-1}(B)(A)$.
The degree-three intersection form is alternating. Since its free
rank is two, the inverse free-coordinate maps are
\[
 (r,s)\longmapsto rY+sZ_3,
 \qquad [x]\longmapsto\bigl(I(x,Z_3),I(Y,x)\bigr).
\]
Subtract these coordinates from a class and apply the torsion linking
maps already described. This yields
$H_3=\mathbb Z^2\oplus\mathbb Z/t\oplus\mathbb Z/m$
with the displayed free-coordinate and torsion-linking inverse maps.

Finally $\pi_1(X)$ vanishes exactly for $n=\pm1$. Then $t=1$
and both middle torsion groups are $(\mathbb Z/m)^2$.
For $m=1$ the manifold is a simply connected integral homology
six-sphere; Hurewicz and Whitehead give a homotopy sphere.
The dimension-six entry of \cite[p.~504]{manKM63}, together with
the first remark on p.~505, gives an oriented diffeomorphism to $S^6$.
A diffeomorphism $f:X\to S^6$ induces inverse homology maps
$f_*$ and $(f^{-1})_*$; for $m>1$ the two $\mathbb Z/m$ summands
in $H_2(X)$ preclude an isomorphism with $H_2(S^6)=0$.
\end{proof}

\subsection{Twisted suspensions and oriented smooth types}

\begin{definition}[Twisted Smale suspension]
\hypertarget{tag-Q118}{}\label{mapsIII:def:twisted-smale-suspension}
{\normalfont\scriptsize tag Q118}\par
For a positive integer $m$, let $N_m$ be the closed simply connected
spin five-manifold with
$H_2(N_m;\mathbb Z)=(\mathbb Z/m)^2$, determined up to oriented
diffeomorphism by \cite[Theorems 2.2--2.3, pp.~376--377]{manBarden65};
in particular $N_1=S^5$. Its twisted suspension is
\begin{equation}
 \widetilde\Sigma_0 N_m
 =\bigl((N_m\setminus\mathring D^5)\times S^1\bigr)
       \cup_{\widetilde\alpha}(S^4\times D^2),
 \qquad \widetilde\alpha(x,z)=(\alpha_zx,z),
 \hypertarget{tag-023Y}{}\label{maps:eq:D.S9}
\end{equation}
where $\alpha_{e^{2\pi it}}$ rotates the first coordinate plane of
$\mathbb R^5$ through $2\pi t$ and fixes its orthogonal complement.
The gluing map in \eqref{maps:eq:D.S9} goes from the boundary of the
attached $S^4\times D^2$ to that of the complement, as in
\cite[Definition 5.1, p.~17]{manGR25}.
\end{definition}

\begin{theorem}[Oriented smooth types of the four semistable families]
\hypertarget{tag-02KO}{}\label{maps:src:D.2.14}
{\normalfont\scriptsize tag 02KO}\par
Every simply connected member of the four families in
Theorem~\ref{maps:src:D.2.13}, with $m\geq1$ and $n=\pm1$, is oriented
diffeomorphic to
\begin{equation}
                         \widetilde\Sigma_0 N_m.
 \hypertarget{tag-023M}{}\label{maps:eq:D.S10}
\end{equation}
Its oriented smooth type depends only on $m$, and these types are
distinct for distinct positive $m$. It is spin, has $p_1=0$, and is
the standard oriented $S^6$ when $m=1$.
\end{theorem}


\begin{proof}
Let $X$ be such a member and put $G=(\mathbb Z/m)^2$.
By Theorem~\ref{maps:src:D.2.13}, $X$ is simply connected and
$H_2(X;\mathbb Z)=G$. We use the period marking of
Definition~\ref{mapsIII:def:smooth-logarithmic-data} and write
$J=\left(\begin{smallmatrix}0&1\\-1&0\end{smallmatrix}\right)$.

\emph{The angular circle.}
Multiplication of the line coordinate by $t\in S^1$ has period
direction $\delta$. It commutes with the line-bundle transitions,
the contracting quotient, and the logarithmic translations, and
therefore defines one smooth action on $X$.
At a rank-one filling write its marked data as
$A=I+\nu u u^tJ$ and $g=\kappa u$. The integral local model in
Theorem~\ref{maps:src:D.2.13} gives
\[
 T-I=\lambda(\nu u^tJ,0,\kappa),\qquad
 \lambda=(u_1,u_2,-p\kappa/\nu,0)^t,
 \qquad \nu\mid p,\quad \gcd(\nu,\kappa)=1.
\]
The row is primitive, and the elliptic projection $u$ of $\lambda$
is primitive. Complete $u$ to a basis of $\mathbb Z^2$.
Together with $\delta$, its lift gives a primitive rank-two
sublattice. Hence $\delta$ remains primitive in
$\ker(T-I)/\mathbb Z\lambda$, the period lattice of the elliptic
base of the ruled normalization. The action is free on that base,
and therefore on the ruled surface and its identified sections.
The local integral conjugacies preserve $\delta$ and carry this
action to the filling. On the added logarithmic torus the lattice is
\[
 \Lambda_v=\mathbb Z^5/
       \mathbb Z(-a,-b,-c,-n,m).
\]
The homomorphism $(x_1,x_2,x_3,x_4,x_5)\mapsto x_3-(c/n)x_4$
vanishes on the relation and takes $\delta$ to $1$.
Thus $\delta$ is primitive there for both $n=1$ and $n=-1$.
It is primitive on each regular fibre as well, so the action is
free away from the cusp.

The cusp marking can be read directly from its periods. Its last
elliptic monodromy has $\nu=1$. Choose $r$ with $\det(u,r)=1$,
and set
\[
 M=[u,r],\quad z=\kappa r,\quad b=-p z^tJ,\qquad
 S=\begin{pmatrix}I_2&0&z\\b&1&0\\0&0&1\end{pmatrix},
 \quad H=S^{-1}\operatorname{diag}(M,1,1).
\]
Here
$S^{-1}=\left(\begin{smallmatrix}I_2&0&-z\\-b&1&0\\0&0&1\end{smallmatrix}\right)$;
all these matrices are integral and unimodular, and direct
multiplication gives
\[
 H\delta=\delta,\qquad H^{-1}T_5H=I+E_{12}+E_{34}.
\]
The short periods are $\lambda=He_1$ and $\delta=He_3$, while
the long periods $He_2,He_4$ acquire these under monodromy.
The angular isomorphism sends $H(x,0,y,0)$ to $(x,y)$, with
inverse $(x,y)\mapsto H(x,0,y,0)$. In the original elliptic-line
angular coordinates it is the multiplicative map
\[
 (X_1,X_2)\longmapsto (X_1,X_1^{p\kappa}X_2),
 \qquad (X_1,X_2)\longmapsto (X_1,X_1^{-p\kappa}X_2)
\]
in the reverse direction. This follows from
$[\lambda,\delta]=[u,\delta]
 \left(\begin{smallmatrix}1&0\\-p\kappa&1\end{smallmatrix}\right)$.
For the four families the triples $(u,\kappa,p)$ are respectively
$((1,1),0,8)$, $((1,0),1,14)$, $((1,1),-1,20)$ and
$((1,1),-1,30)$, and $r=(0,1)$ works in each case.
Consequently $\delta=(0,1)$ in the angular marking of the
primitive triangular cusp in Theorem~\ref{maps:src:D.2.6},
with the fan and period quotient of
Definition~\ref{mapsIII:def:triangular-cusp}.

\emph{The fixed curve and its normal framing.}
In $N=\mathbb Z^2\oplus\mathbb Z$ the cusp fan has rays $(v,1)$
and maximal cones over
\[
 (v,v+e_1,v+e_2),\qquad
 (v+e_1+e_2,v+e_1,v+e_2).
\]
A period $\ell\in\mathbb Z^2$ acts by
\[
 (X_1,X_2,q)\longmapsto
 (c_1(q)q^{\ell_1}X_1,c_2(q)q^{\ell_2}X_2,q),
\]
where the nonzero holomorphic units $c_j$ commute with the circle.
It translates the fan cells by $\ell$. A nonzero translation
preserves no finite nonempty cell, so it creates no additional
isotropy on a central toric orbit.
For a cone with ray-column matrix $R$, the circle weights are
$R^{-1}(0,1,0)^t$. For the two displayed triangle types they are
$(-1,0,1)$ and $(1,-1,0)$. On an edge of direction $d$ the
residual weight is $\det(d,\delta)$, which is $1,0,-1$ for
$d=e_1,e_2,e_2-e_1$. On a ray orbit $\delta$ is primitive
modulo the ray. Thus every stabilizer is either trivial or the
whole circle, and the fixed curves correspond precisely to
$e_2$-edges. There is one such edge in each triangle.
Their complete curves are disjoint; translation is transitive on
them and preserves the two triangle types, so it identifies
neither the two endpoints of one curve nor two points on it.
The fixed set in $X$ is therefore one embedded $C\simeq\mathbb P^1$.

For the edge from $0$ to $e_2$, use the two adjacent charts
\[
 (x,y,w)=\bigl(q/(X_1X_2),X_2,X_1\bigr),\qquad
 (x',y',w')=\bigl(q/X_2,X_1X_2,1/X_1\bigr).
\]
Their inverses are
$(X_1,X_2,q)=(w,y,xyw)$ and
$(X_1,X_2,q)=(1/w',y'w',x'y'w')$.
The transition and its inverse are
\[
 (x',y',w')=(wx,wy,w^{-1}),\qquad
 (x,y,w)=(w'x',w'y',(w')^{-1}).
\]
The curve is $x=y=0$, and the weights of $(x,y)$ are $(-1,+1)$.
Both normal transitions multiply by $w$. In the order $(y,x)$
this gives
\begin{equation}
 \nu_C=\mathcal O(-1)_{+1}\oplus\mathcal O(-1)_{-1}.
 \hypertarget{tag-023N}{}\label{maps:eq:D.S11}
\end{equation}
Each nontrivial period translate of this complete curve is disjoint from it;
hence these charts also give its normal bundle in the quotient,
including the holomorphic unit factors in the period action.

Multiplying each normal coordinate by $\sqrt{1+|w|^2}$ changes
the overlap multiplier to $w/|w|$; the inverse divides by the
same positive function. In the coordinates $(\bar x,y)$ both
circle weights are $+1$, and the equatorial transition is
$g(\theta)=\operatorname{diag}(e^{-i\theta},e^{i\theta})$.
An explicit extension of this loop over a disc is
\[
 U(r,\theta)=
 \begin{pmatrix}
  \sin\beta(r)e^{-i\theta}&-\cos\beta(r)\\
  \cos\beta(r)&\sin\beta(r)e^{i\theta}
 \end{pmatrix}\in SU(2),
\]
where $\beta(r)=r$ near $0$ and $\beta(r)=\pi/2$ near $1$.
Its entries are smooth at the centre, its inverse is $U^*$,
and it equals $g$ on a boundary collar. Use the northern normal
coordinate unchanged and multiply the southern one by $U^{-1}$.
These formulas trivialize the bundle equivariantly.
Conjugation reverses the normal orientation; simultaneously
reverse the chosen parametrization $S^2\to C$. The resulting
trivialization preserves the orientation of $X$.
An invariant normal exponential and its inverse now identify a
tubular neighbourhood with $C\times D^4$, carrying scalar
circle multiplication on $D^4\subset\mathbb C^2$.

\emph{The smooth orbit space and its homology map.}
Put $X_0=X\setminus(C\times\mathring D^4)$ and let
$p:X_0\to Q_0$ be its principal circle bundle, with Euler class $e$.
Adjoin $C\times D^3$ along its Hopf boundary to form the smooth
closed orbit manifold $Q$. The normal quotient chart is
\[
 h(z_1,z_2)=
 (2\operatorname{Re}(z_1\bar z_2),2\operatorname{Im}(z_1\bar z_2),
                         |z_1|^2-|z_2|^2),\qquad
 \mathcal H(z)=\frac{h(z)}{|z|}\ (z\ne0),\quad \mathcal H(0)=0.
\]
It has $|\mathcal H(z)|=|z|$ and identifies exactly the circle
orbits. It specifies the smooth coordinates of $Q$ and a continuous
orbit map at the fixed curve. On $S^2$ the Hopf sections are
\[
 s_N(v)=\left(\sqrt{\frac{1+v_3}{2}},
                  \frac{v_1-iv_2}{\sqrt{2(1+v_3)}}\right),\qquad
 s_S(v)=\left(\frac{v_1+iv_2}{\sqrt{2(1-v_3)}},
                         \sqrt{\frac{1-v_3}{2}}\right).
\]
Their overlap multiplier is
$(v_1+iv_2)/\sqrt{v_1^2+v_2^2}$. The inverse local orbit chart
with circle coordinate $t$ is $|q|t s_N(q/|q|)$ or its southern
counterpart. Orient the meridian sphere $\sigma$ so that this
transition has winding $+1$; then $\langle e,\sigma\rangle=1$.
The orientation of $Q_0$ is fixed by placing its horizontal
five-plane before the positive circle direction in $TX_0$.
It extends over the displayed normal quotient charts.

Discs of dimensions at most three can be made transverse to $C$
and disjoint from it, since $3+2<6$. Thus the inclusion
$j:X_0\to X$ induces isomorphisms on $\pi_1,\pi_2$ and $H_2$.
The homotopy sequence and Hurewicz give $\pi_1(Q_0)=0$ and
\[
 0\longrightarrow G\xrightarrow{p_*}H_2(Q_0;\mathbb Z)
       \xrightarrow{\langle e,-\rangle}\mathbb Z\longrightarrow0,
 \qquad k\longmapsto k\sigma
\]
as an explicit splitting. Excision gives
$H_3(Q,Q_0)=\mathbb Z$, with boundary $\sigma$, and
$H_2(Q,Q_0)=0$. If $i:Q_0\to Q$ is inclusion, then
\[
 \Theta=i_*p_*j_*^{-1}:H_2(X;\mathbb Z)\xrightarrow{\cong}H_2(Q;\mathbb Z).
\]
To invert $\Theta$, lift a class to $a\in H_2(Q_0)$, replace it
by $a-\langle e,a\rangle\sigma$, lift the resulting sphere to
$X_0$, and apply $j_*$. The sphere lifts because its pulled-back
circle bundle has Euler number zero. Two choices of section differ
by a map $S^2\to S^1$ and are homotopic; a different first lift
differs by an integral multiple of $\sigma$ and gives the same
corrected class. Van Kampen also gives $\pi_1(Q)=0$.

The residue sequence contains
\[
 H^2(Q;\mathbb Z)\longrightarrow H^2(Q_0;\mathbb Z)
 \xrightarrow{\mathrm{res}}H^0(C;\mathbb Z)
 \longrightarrow H^3(Q;\mathbb Z).
\]
Since $\mathrm{res}(e)=1$, the image of $1$, which is the
Poincar\'e dual of $[C_Q]$, vanishes. Here $C_Q$ denotes the
image of $C$ in $Q$. Thus $[C_Q]=0$ in $H_2(Q;\mathbb Z)$.
Moreover $H^2(Q;\mathbb Z)=0$ because $H_2(Q)=G$ is finite;
the residue map is an isomorphism, with inverse $k\mapsto ke$.

The complex manifold $X$ is spin, since $H^2(X;\mathbb Z)=0$
implies $c_1(X)=0$. On its free part
$TX_0\cong p^*TQ_0\oplus\mathbb R$.
The mod-two Gysin sequence identifies $\ker p^*$ in degree two
with the span of $e\bmod2$. Its nonzero element evaluates to
$1$ on $\sigma$, whereas $w_2(Q_0)$ evaluates to $0$, because
$TQ_0|_\sigma\cong TS^2\oplus\mathbb R^3$.
Hence $w_2(Q_0)=0$. Excision makes
$H^2(Q;\mathbb Z/2)\to H^2(Q_0;\mathbb Z/2)$ injective, so $Q$
is spin.

\emph{The circle action on the twisted suspension.}
Write $C_m=N_m\setminus\mathring D^5$ and first use the model
\[
 M_-=(C_m\times S^1)\cup_g(S^4\times D^2),\qquad
 g(x,z)=(\alpha_zx,z),\quad g^{-1}(x,z)=(\alpha_z^{-1}x,z),
\]
where $g$ goes from the complement to the attached piece.
The actions are $t(x,z)=(x,tz)$ on the first piece and
$t(x,w)=(\alpha_tx,tw)$ on the second. The identity
$\alpha_t\alpha_z=\alpha_{tz}$ makes them agree along $g$.
The fixed set is $C_* = \{(0,0,u)\in S^4:u\in\mathbb R^3\}$
at $w=0$, with two normal scalar weights $+1$; every other
stabilizer is trivial.

When $N_m=S^5$, the map to $\partial(D^5\times D^2)$ is
$(x,z)\mapsto(\alpha_zx,z)$ on $D^5\times S^1$ and the
identity on $S^4\times D^2$. Its inverse uses $\alpha_z^{-1}$
on the first piece. The product boundary orientation is negative
on $D^5\times S^1$ and positive on $S^4\times D^2$; the
complementary disc coordinate has the former orientation.
The maps consequently preserve orientation. Invariant corner
rounding and radial projection identify this boundary with
$S^6\subset\mathbb R^3\oplus\mathbb C^2$, with the circle
acting scalarly on $\mathbb C^2$. Its orbit map is
$(u,z)\mapsto(u,\mathcal H(z))\in S^5$.
In particular the quotient of the first piece is a parametrized
$D^5$, and that of its complement is a second $D^5$.
The parametrized disc isotopy theorem
\cite[Theorem 2.12, p.~10]{manGR25}, followed by ambient isotopy
extension, identifies the latter disc with a standard complementary
disc with the same boundary parametrization: the isotopy carries
each point of the first embedded disc to the point with its
unchanged parameter in the standard disc.
On the attached boundary the orbit coordinate is
$\alpha_z^{-1}x$, and its composition with $g$ is $x$.
Thus for every $m$ the quotient of $M_-$ is $N_m$ with its
disc restored by the identity boundary map. Its fixed sphere lies
inside that disc and is nullhomotopic.

The loop $\alpha$ lifts to
$\cos(\pi t)+\sin(\pi t)e_1e_2$ in $\operatorname{Spin}(5)$,
where $(e_1e_2)^2=-1$. The endpoints are $1$ and $-1$,
so it represents the nonzero element of $\pi_1(SO(5))$.
Because the gluing directions are opposite, $M_-$ is the
construction of \eqref{maps:eq:D.S9} with representative
$\alpha^{-1}$. To identify the representatives explicitly, put
\[
 A=\operatorname{diag}(1,-1,-1,1,1),\qquad
 A\alpha_zA^{-1}=\alpha_z^{-1},\qquad
 A\alpha_z^{-1}=\alpha_z A.
\]
Let $A_s$ be rotation through $\pi s$ in the coordinate plane
$(2,3)$. On a boundary collar of $C_m$ set
$F_m(r,x)=(r,A_{\chi(r)}x)$, where $\chi=1$ near the boundary
and $\chi=0$ off the collar, and extend by the identity.
Its inverse replaces $A_{\chi(r)}$ by $A_{\chi(r)}^{-1}$.
The maps $F_m\times\mathrm{id}$ and $A\times\mathrm{id}$ on
the two pieces of $M_-$ agree in \eqref{maps:eq:D.S9} because
$A\alpha_z^{-1}u=\alpha_zAu$. They give an oriented
diffeomorphism
$\Psi:M_-\to\widetilde\Sigma_0N_m$ with these piecewise
inverses. It is equivariant for the attached target action
$(u,w)\mapsto(\alpha_t^{-1}u,tw)$.

\emph{Reconstructing the total space from its quotient.}
The torsion linking form of $Q$ is alternating by the spin case
of \cite[Lemma F, p.~372]{manBarden65}.
Choose $x\in G$ of order $m$ and $y$ with
$\ell(x,y)=1/m$. The map
$(r,s)\mapsto rx+sy$ has inverse
\[
 a\longmapsto
 \bigl(\operatorname{num}_m\ell(a,y),
              -\operatorname{num}_m\ell(a,x)\bigr),
 \qquad \operatorname{num}_m([j/m])=j\bmod m.
\]
It identifies the linking form with the standard alternating
form on $(\mathbb Z/m)^2$. Choose the same linking coordinates
on $N_m$. By \cite[Theorems 2.2--2.3, pp.~376--377]{manBarden65}
the resulting linking isomorphism is induced by an oriented
diffeomorphism $f_0:Q\to N_m$.

Hurewicz and $[C_Q]=0$ make $C_Q$ nullhomotopic. Choose its
parametrization by the orientation fixed in the normal
trivialization above, and choose the corresponding parametrization
of $C_*$ in $M_-$. Both maps into $N_m$ are nullhomotopic.
By \cite[Theorem B(b), pp.~367--368]{manBarden65}, they are
isotopic: the inequality is $2\cdot5>3\cdot2+3$, and the
required connectivity is connectedness and simple connectedness.
Ambient isotopy extension supplies $I_t$ with $I_1f_0$ carrying
the first parametrized sphere to the second. Set $f=I_1f_0$;
its inverse is $f_0^{-1}I_1^{-1}$. This map preserves the chosen
core orientation and hence the oriented normal three-plane
and the meridian of Euler residue $+1$.

Tubular-neighbourhood uniqueness adjusts $f$ near the core to
a fibrewise orthogonal normal map $R:C_Q\to SO(3)$.
The map $SU(2)\to SO(3)$ defined by
$h(Vz)=\rho(V)h(z)$ is the double cover associated to the Hopf
map. Choose a lift at one point; lifting paths on the simply
connected sphere gives a smooth lift $\widehat R:C_Q\to SU(2)$.
It and its conjugate transpose give mutually inverse equivariant
normal maps. Together with the core map they yield a tubular
diffeomorphism $\Phi_N:C\times D^4\to C_*\times D^4$ covering
$f$. The normal-framing adjustment is unobstructed because
$\pi_2(SO(3))=0$.

Identify the free bases by $f$. Their two circle bundles have
Euler residue $1$, and the isomorphism
$H^2(Q_0;\mathbb Z)\xrightarrow{\mathrm{res}}\mathbb Z$ proves
that their Euler classes agree. An isomorphism can be written
from their transitions. Take a common good cover and sections
$s_i,t_i$ with $s_j=s_i a_{ij}$ and $t_j=t_i b_{ij}$.
Choose logarithms $\ell_{ij}$ of $b_{ij}/a_{ij}$, subtracting an
integral cochain so that
$\ell_{ij}+\ell_{jk}=\ell_{ik}$; this is possible precisely
because their Chern classes agree. For a subordinate partition
of unity set $f_i=\sum_k\rho_k\ell_{ik}$. Then
$f_i-f_j=\ell_{ij}$, and
\[
 \Phi_0(s_i z)=t_i e^{2\pi if_i}z,\qquad
 \Phi_0^{-1}(t_i z)=s_i e^{-2\pi if_i}z
\]
are well-defined inverse bundle maps.
On a collar of $L=C\times S^2$, let
$\Phi_N=\Phi_0\circ R_h$, with $R_h$ the gauge multiplication
by $h$. Since $H^1(L\times[0,\epsilon);\mathbb Z)=0$, write
$h=e^{2\pi ik}$ with $k$ a smooth real function. Replace
$\Phi_0$ by $\Phi_0\circ R_{e^{2\pi i\chi k}}$, where $\chi$
is one near the boundary and zero off the collar.
The inverse is $R_{e^{-2\pi i\chi k}}\circ\Phi_0^{-1}$.
The modified map equals $\Phi_N$ on a full collar, so the two
maps glue smoothly, with their displayed inverses, to
$\Phi:X\to M_-$.
Their orientations agree with the horizontal-before-circle
convention and with the chosen oriented normal trivializations.
The oriented diffeomorphism $\Psi\Phi$ proves
\eqref{maps:eq:D.S10}.

\emph{The surgery descriptions.}
A constant nonzero normal vector pushes $C$ into $X_0$.
Under $\Theta$ its class maps to $[C_Q]=0$, because its quotient
is a normal push-off of $C_Q$. Thus $[C]=0$ in $H_2(X;\mathbb Z)$,
and Hurewicz makes its parametrization nullhomotopic in $X$.
Use the equivariant normal trivialization above to replace
$S^2\times D^4$ by $D^3\times S^3$, with scalar Hopf action
on $S^3$, and denote the resulting free circle manifold by $Y$.
By \cite[Theorem 2.9, p.~9]{manGR25}, homotopic embeddings
$S^2\to X^6$ are isotopic, since $6>2\cdot2+1$.
The framing has one homotopy class because $\pi_2(SO(4))=0$.
Transport the surgery by this parametrized ambient isotopy to
an unknotted sphere in a ball. In
$S^6=\partial(D^3\times D^4)$ the replacement gives the
double of $D^3\times S^3$, namely $S^3\times S^3$.
The identity on the exterior and these product maps in the
ball, transported by the isotopy and its inverse, give
\begin{equation}
                         Y\cong X\#(S^3\times S^3).
 \hypertarget{tag-023O}{}\label{maps:eq:D.S12}
\end{equation}

Let $B=Y/S^1$, and let $e_Y$ be its Euler class.
The quotient surgery replaces $S^2\times D^3$ in $Q$ by
$D^3\times S^2$. The parametrized isotopy of $C_Q$ used above
and $\pi_2(SO(3))=0$ identify it with the framed unknotted
surgery in a ball. Applying
$S^5=\partial(D^3\times D^3)$ identifies the result with
$Q\#(S^2\times S^3)$, and then with
$(S^2\times S^3)\#N_m$. The new $S^2$ has the Hopf Euler
number $1$, while the $N_m$ component has $H^2(N_m;\mathbb Z)=0$.
Thus, writing $e=e_Y$ in the pair below,
\begin{equation}
 (B,e)\cong\bigl((S^2\times S^3)\#N_m,(x,0)\bigr),
 \hypertarget{tag-023P}{}\label{maps:eq:D.S13}
\end{equation}
where $x$ evaluates to $1$ on the oriented Hopf base.
This identification also proves that $B$ is spin.
Finally use $\Psi\Phi$ on the punctured $X$ factor of
\eqref{maps:eq:D.S12}, transporting the removed disc by that
diffeomorphism. It extends by the identity over the other factor.
Interchanging the punctured factors and sending
$(t,u)\mapsto(-t,\rho u)$ on the connecting neck, with $\rho$
a reflection of $S^5$ and the corresponding boundary
parametrizations, is orientation preserving and is its own
inverse. This yields
\begin{equation}
          Y\cong(S^3\times S^3)\#\widetilde\Sigma_0 N_m.
 \hypertarget{tag-023Q}{}\label{maps:eq:D.S14}
\end{equation}

\emph{Characteristic classes and the parameter.}
For $P=C_m\times S^1\subset M_-$, excision gives
$H^k(M_-,P;R)=H^{k-2}(S^4;R)$ with any coefficient ring $R$.
The map $H^1(P;\mathbb Z/2)\to H^2(M_-,P;\mathbb Z/2)$
sends the circle generator to the normal-disc generator and
is an isomorphism. Thus restriction is injective on
$H^2(-;\mathbb Z/2)$, and $w_2(M_-)=0$ follows from
$TP\cong\operatorname{pr}_1^*TN_m\oplus\mathbb R$.
In degree four both adjacent relative groups vanish, so
$H^4(M_-;\mathbb Z)\to H^4(P;\mathbb Z)$ is an isomorphism.
Since $H^4(N_m;\mathbb Z)\cong H_1(N_m;\mathbb Z)=0$,
the restriction of $p_1(M_-)$ is zero, hence $p_1(M_-)=0$.
The derivative of $\Psi\Phi$ and its inverse carry these
characteristic classes to those of $X$.
All permitted $a,b,c$ and both signs of $n$ have the same
oriented model for fixed $m$. A diffeomorphism between two
values $m,m'$ would induce inverse isomorphisms of groups of
orders $m^2$ and $(m')^2$, forcing $m=m'$.
For $m=1$ the explicit boundary map and radial projection above,
composed with $\Psi^{-1}$, identify
$\widetilde\Sigma_0S^5$ with the standard oriented $S^6$.
\end{proof}

\section{Section-line families and elliptic catalogues}
\subsection{Marked elliptic profiles and whole-fibre linearizations}

\begin{definition}[Ordered elliptic profile data]
\hypertarget{tag-Q110}{}\label{mapsIII:def:elliptic-marked-data}
{\normalfont\scriptsize tag Q110}\par
For each of the twenty-one rational elliptic profiles in the tables,
let $f:S\to\mathbb P^1$ be its marked surface and let $G$ generate
its Mordell--Weil group. The ordered elliptic periods are $e_1,e_2$;
$\delta$ is the positive line circle and $\gamma$ is the forward
linearization period. Write
$\Lambda=\mathbb Z\langle e_1,e_2,\delta,\gamma\rangle$,
$\chi=\gamma^*$, $P=aG$, $Q=bG$ and
$\langle G,G\rangle=1/d$. The line is $\mathcal O(P-O)$ and
the forward linearization covers $-Q$.
The following tables specify the ordered elliptic matrices and
negative Abel--Jacobi cocycles of $G$, in the convention of
Theorem~\ref{maps:src:D.2.12}. A row $(\alpha,\beta,\zeta,\eta;p,q)$ means
$A_i=\left(\begin{smallmatrix}\alpha&\beta\\\zeta&\eta\end{smallmatrix}\right)$
and $g_i=(p,q)^t$. The symbol $\mathrm{base}$ denotes the profile
with the largest number of singular fibres in its OS entry; an
unadorned OS number denotes this same ordered profile. The index
$i$ is its peripheral order. The height denominators are
\[
\begin{array}{c|rrrrrrrr}
\mathrm{OS}&43&45&46&47&49&50&55&56\\\hline
d&2&8&4&14&6&12&20&30.
\end{array}
\]

\begingroup\small\setlength{\arraycolsep}{3pt}
\[
\begin{array}{clrl|rrrr|rr}
\mathrm{OS}&\mathrm{profile}&i&\mathrm{type}&\alpha&\beta&\zeta&\eta&p&q\\\hline
43&\mathrm{base}&1&III^{*}&2&-5&1&-2&-1&0\\
43&\mathrm{base}&2&I_{1}&-2&9&-1&4&0&0\\
43&\mathrm{base}&3&I_{1}&1&1&0&1&0&0\\
43&\mathrm{base}&4&I_{1}&1&1&0&1&1&0\\
43&\mathrm{II}&1&III^{*}&2&-5&1&-2&-1&0\\
43&\mathrm{II}&2&II&-2&7&-1&3&0&0\\
43&\mathrm{II}&3&I_{1}&1&1&0&1&1&0\\
45&\mathrm{base}&1&I_{8}&1&8&0&1&3&0\\
45&\mathrm{base}&2&I_{1}&1&1&0&1&0&0\\
45&\mathrm{base}&3&I_{1}&3&4&-1&-1&2&-1\\
45&\mathrm{base}&4&I_{1}&3&1&-4&-1&1&-2\\
45&\mathrm{base}&5&I_{1}&0&1&-1&2&0&0\\
45&\mathrm{II}&1&I_{8}&1&8&0&1&3&0\\
45&\mathrm{II}&2&II&2&3&-1&-1&1&-1\\
45&\mathrm{II}&3&I_{1}&3&1&-4&-1&1&-2\\
45&\mathrm{II}&4&I_{1}&0&1&-1&2&0&0\\
46&\mathrm{base}&1&I_{3}^{*}&-1&-3&0&-1&1&1\\
46&\mathrm{base}&2&I_{1}&1&1&0&1&0&0\\
46&\mathrm{base}&3&I_{1}&2&1&-1&0&0&0\\
46&\mathrm{base}&4&I_{1}&0&1&-1&2&-1&-1\\
46&\mathrm{II}&1&I_{3}^{*}&-1&-3&0&-1&1&1\\
46&\mathrm{II}&2&II&1&1&-1&0&0&0\\
46&\mathrm{II}&3&I_{1}&0&1&-1&2&-1&-1\\
47&\mathrm{base}&1&I_{7}&1&7&0&1&5&0\\
\end{array}
\]
\endgroup

\begingroup\small\setlength{\arraycolsep}{3pt}
\[
\begin{array}{clrl|rrrr|rr}
\mathrm{OS}&\mathrm{profile}&i&\mathrm{type}&\alpha&\beta&\zeta&\eta&p&q\\\hline
47&\mathrm{base}&2&I_{2}&3&2&-2&-1&1&-1\\
47&\mathrm{base}&3&I_{1}&4&1&-9&-2&0&0\\
47&\mathrm{base}&4&I_{1}&0&1&-1&2&1&1\\
47&\mathrm{base}&5&I_{1}&1&1&0&1&1&0\\
47&\mathrm{II}&1&I_{7}&1&7&0&1&5&0\\
47&\mathrm{II}&2&I_{2}&3&2&-2&-1&1&-1\\
47&\mathrm{II}&3&I_{1}&4&1&-9&-2&0&0\\
47&\mathrm{II}&4&II&0&1&-1&1&1&0\\
47&\mathrm{III}&1&I_{7}&1&7&0&1&5&0\\
47&\mathrm{III}&2&I_{1}&4&1&-9&-2&0&0\\
47&\mathrm{III}&3&I_{1}&0&1&-1&2&1&1\\
47&\mathrm{III}&4&III&-13&85&-2&13&64&9\\
49&\mathrm{base}&1&I_{2}&1&2&0&1&1&0\\
49&\mathrm{base}&2&IV^{*}&1&-3&1&-2&-1&0\\
49&\mathrm{base}&3&I_{1}&-2&9&-1&4&0&0\\
49&\mathrm{base}&4&I_{1}&1&1&0&1&0&0\\
49&\mathrm{II}&1&II&1&1&-1&0&0&0\\
49&\mathrm{II}&2&IV^{*}&0&-1&1&-1&0&-1\\
49&\mathrm{II}&3&I_{2}&1&2&0&1&1&0\\
49&\mathrm{III}&1&IV^{*}&1&-3&1&-2&-1&0\\
49&\mathrm{III}&2&III&-2&5&-1&2&-2&-1\\
49&\mathrm{III}&3&I_{1}&1&1&0&1&0&0\\
50&\mathrm{base}&1&I_{1}^{*}&-1&-1&0&-1&-1&1\\
50&\mathrm{base}&2&I_{3}&1&3&0&1&-2&0\\
\end{array}
\]
\endgroup

\begingroup\small\setlength{\arraycolsep}{3pt}
\[
\begin{array}{clrl|rrrr|rr}
\mathrm{OS}&\mathrm{profile}&i&\mathrm{type}&\alpha&\beta&\zeta&\eta&p&q\\\hline
50&\mathrm{base}&3&I_{1}&2&1&-1&0&0&0\\
50&\mathrm{base}&4&I_{1}&0&1&-1&2&-1&-1\\
50&\mathrm{II}&1&II&3&7&-1&-2&-5&2\\
50&\mathrm{II}&2&I_{3}&7&12&-3&-5&-8&4\\
50&\mathrm{II}&3&I_{1}^{*}&-2&-1&1&0&0&1\\
50&\mathrm{IV}&1&I_{1}^{*}&-1&-1&0&-1&-1&1\\
50&\mathrm{IV}&2&IV&-1&1&-1&0&-2&0\\
50&\mathrm{IV}&3&I_{1}&0&1&-1&2&-1&-1\\
55&\mathrm{base}&1&I_{5}&1&5&0&1&-3&0\\
55&\mathrm{base}&2&I_{4}&1&4&0&1&-3&0\\
55&\mathrm{base}&3&I_{1}&3&4&-1&-1&-2&1\\
55&\mathrm{base}&4&I_{1}&3&1&-4&-1&0&0\\
55&\mathrm{base}&5&I_{1}&0&1&-1&2&-1&-1\\
55&\mathrm{II}&1&I_{5}&1&5&0&1&-3&0\\
55&\mathrm{II}&2&I_{4}&5&4&-4&-3&-1&1\\
55&\mathrm{II}&3&II&5&3&-7&-4&0&0\\
55&\mathrm{II}&4&I_{1}&0&1&-1&2&-1&-1\\
56&\mathrm{base}&1&I_{5}&1&5&0&1&-4&0\\
56&\mathrm{base}&2&I_{3}&1&3&0&1&-2&0\\
56&\mathrm{base}&3&I_{2}&3&2&-2&-1&-1&1\\
56&\mathrm{base}&4&I_{1}&4&1&-9&-2&0&0\\
56&\mathrm{base}&5&I_{1}&0&1&-1&2&-1&-1\\
56&\mathrm{II}&1&I_{2}&5&8&-2&-3&-2&1\\
56&\mathrm{II}&2&II&27&37&-19&-26&0&0\\
\end{array}
\]
\endgroup

\begingroup\small\setlength{\arraycolsep}{3pt}
\[
\begin{array}{clrl|rrrr|rr}
\mathrm{OS}&\mathrm{profile}&i&\mathrm{type}&\alpha&\beta&\zeta&\eta&p&q\\\hline
56&\mathrm{II}&3&I_{5}&61&80&-45&-59&4&-3\\
56&\mathrm{II}&4&I_{3}&4&3&-3&-2&1&-1\\
56&\mathrm{III}&1&I_{5}&1&5&0&1&-4&0\\
56&\mathrm{III}&2&I_{3}&7&3&-12&-5&-1&2\\
56&\mathrm{III}&3&III&12&5&-29&-12&-2&5\\
56&\mathrm{III}&4&I_{1}&0&1&-1&2&-1&-1\\
56&\mathrm{IV}&1&I_{5}&1&5&0&1&-4&0\\
56&\mathrm{IV}&2&I_{2}&3&2&-2&-1&-1&1\\
56&\mathrm{IV}&3&IV&1&1&-3&-2&-1&2\\
56&\mathrm{IV}&4&I_{1}&0&1&-1&2&-1&-1\\
\end{array}
\]
\endgroup

\end{definition}

\begin{proposition}[Degree-one whole-fibre linearizations]\hypertarget{tag-02KT}{}\label{maps:src:D.2.19}
{\normalfont\scriptsize tag 02KT}\par


Use the section line
$M=\mathcal O(P-O)$ and require its linearization line
$J=M\otimes t_Q^*M^{-1}$ to be $f^*\mathcal O_{\mathbb P^1}(1)$.
This is the case of a simple zero on one whole fibre.
Write $P=aG$, $Q=bG$ and $h(G)=1/d$.
Up to simultaneous sign reversal, the possible ordered
pairs are exactly the positive pairs $ab=d$ for which,
at every reducible fibre, $P$ or $Q$ is narrow there.
For the marked profiles of Definition~\ref{mapsIII:def:elliptic-marked-data} they are
\begin{equation}
\begin{array}{c|r|l}
 \mathrm{OS}&d& (a,b)\\\hline
 43&2&(1,2),(2,1)\\
 45&8&(1,8),(8,1)\\
 46&4&(1,4),(4,1)\\
 47&14&(1,14),(2,7),(7,2),(14,1)\\
 49&6&(1,6),(2,3),(3,2),(6,1)\\
 50&12&(1,12),(3,4),(4,3),(12,1)\\
 55&20&(1,20),(4,5),(5,4),(20,1)\\
 56&30&(1,30),(2,15),(3,10),(5,6),\\
 &&(6,5),(10,3),(15,2),(30,1).
\end{array}
\hypertarget{tag-01YO}{}\label{maps:eq:D.F1}
\end{equation}
The same row applies to every profile of its OS entry.
Across all 21 marked profiles there are 104 positive
ordered factor pairs with $ab=d$, exactly 88 of which
satisfy the stated requirement. Including each marked
multiplicative fibre and one regular-fibre position for
the simple zero gives 348 finite initial data. All their
marked integral monodromy tuples are given by the formula in the proof.
Each datum specifies the marked profile, the ordered pair,
and the chosen whole-fibre zero for this section line.
The regular-fibre position is one fixed marked regular point.
\end{proposition}


\begin{proof}
Write \(J_0=\begin{psmallmatrix}0&1\\-1&0\end{psmallmatrix}\)
for the alternating matrix on the marked elliptic periods.
The forward line arrow covers translation by \(-Q\); with target
variable \(p\), it is a map \(M_{p+Q}\to M_p\).
Pulling back the section divisors by translation gives
\[
 D=P-O-(P-Q)+(-Q),\qquad J=\mathcal O(D).
\]
The inverse translation adds \(Q\). At a reducible fibre, denote
the component classes of \(P,Q\) by \(p,q\in\Phi\). The degree
of \(J\) on the multiplicity-one component indexed by \(j\) is
\[
 [j=p]-[j=0]-[j=p-q]+[j=-q].
\]
The vector vanishes for \(p=0\) and for \(q=0\). For nonzero
\(p,q\), its entry at the identity is \(-1-[p=q]\), so it
does not vanish. Sections meet multiplicity-one components, and
translation preserves component multiplicities. The remaining
component degrees are therefore zero. We have proved the exact
component condition
\begin{equation}
                  \bar P=0\quad\hbox{or}\quad\bar Q=0
 \hypertarget{tag-01YP}{}\label{maps:eq:D.F2}
\end{equation}
at every reducible fibre.

Here is the line isomorphism that realizes the condition.
On the generic Weierstrass fibre, divide the line through \(P,-Q\)
by the vertical line through \(P-Q\). The numerator has divisor
\(P+(-Q)+(Q-P)-3O\), and the denominator has divisor
\((P-Q)+(Q-P)-2O\). Their quotient has divisor \(D\).
Use the tangent for coincident points; for \(P=Q\), use
\(x-x(P)\), with divisor \(P+(-P)-2O\).
Thus a rational function \(F\) on the elliptic surface satisfies
\[
                         \operatorname{div}(F)=D+V
\]
for an integral vertical divisor \(V\).
Under \eqref{maps:eq:D.F2}, the intersection of \(V\) with
every fibre component is zero. The kernel of each fibre
intersection matrix is its multiplicity vector. That vector
contains an entry one, since the surface has a section, so
\(V=\sum_v b_vF_v\) with \(b_v\in\mathbb Z\).
Consequently
\[
 J\simeq f^*\mathcal O_{\mathbb P^1}\!\left(-\sum_v b_v[v]\right).
\]
Its degree follows by intersecting \(D\) with \(O\):
\[
 D\cdot O=1+P\cdot O+Q\cdot O-P\cdot Q
          =\langle P,Q\rangle=ab/d.
\]
Here \(O^2=-1\), translation gives
\((P-Q)\cdot O=P\cdot Q\), inversion gives
\((-Q)\cdot O=Q\cdot O\), and every local mixed height
correction vanishes by \eqref{maps:eq:D.F2}; see
\cite[Section~11.8, pp.~51--52]{manSS10}.
Pullback of a base line is injective on its degree because
restriction to \(O\) is the inverse. This proves both directions
of the stated criterion for \(J\simeq f^*\mathcal O(1)\).

Fix an allowed zero position \(z_0\), and put \(F_0=f^*[z_0]\).
For an accepted pair, the base divisor
\(\sum b_v[v]+[z_0]\) has degree zero. A product of powers of
linear factors on the base gives a rational function \(B\) with
that divisor. Then
\[
 \rho=F/(B\circ f),\qquad \operatorname{div}(\rho)=D-F_0,
 \qquad
 \mathcal O(D)\xrightarrow{\;\rho\;}\mathcal O(F_0),
 \qquad
 \mathcal O(F_0)\xrightarrow{\;\rho^{-1}\;}\mathcal O(D)
\]
are inverse line maps. Pulling back the canonical section of
\(\mathcal O(F_0)\) produces the required section with one
simple whole-fibre zero. Its norm as the displayed line arrow
has a finite maximum on the compact surface. Multiplication by
a sufficiently small nonzero scalar gives the contracting arrow.
This constructs every zero marker in the statement.

The equation \(ab=d>0\) gives equal signs for \(a,b\).
Elliptic inversion pulls \(\mathcal O(P-O)\) to
\(\mathcal O((-P)-O)\), conjugates translation by \(Q\) to
translation by \(-Q\), and transports \(\rho\) and its inverse.
Its inverse is itself, and its marked period map is
\(\operatorname{diag}(-I_2,1,1)\). Hence the simultaneous
negative pair is transported to the positive representative
without exchanging the two ordered factors.

For the finite list, use the ordered matrices and section vectors of
Definition~\ref{mapsIII:def:elliptic-marked-data}.
Let \(A_i\) be a local elliptic monodromy and \(g_i\) the
negative Abel--Jacobi cocycle of \(G\). The negative increment
for \(kG\) is \(kg_i\); a change of logarithm branch contributes
the coboundary \((A_i-I)v\), with inverse contribution
\(-(A_i-I)v\). In this marking the component class is
\[
 [kg_i]\in\mathbb Z^2/(A_i-I)\mathbb Z^2.
\]
The map \([x]\mapsto[-x]\), whose inverse is itself, gives the
positive-increment marking. It fixes the zero class and preserves
the order of every component class.
At an \(I_n\) place write
\(A_i=I+n u u^tJ_0\), with \(u\) primitive, and
\(g_i=\kappa u\). Choose \(r\) with \(\det(u,r)=1\).
The component coordinate is \(k\kappa\bmod n\), because
\((A_i-I)((k\kappa/n)r)=kg_i\). The order of \([g_i]\)
is \(n/\gcd(n,\kappa)\). At an additive place,
\(A_i-I\) is invertible and the solution is
\[
 (A_i-I)^{-1}g_i
       =\frac{\operatorname{adj}(A_i-I)g_i}{\det(A_i-I)}.
\]
The least common multiple of its coordinate denominators is
the order of the component class. Applying these exact formulas
to the defining matrices gives
\[
\begin{array}{c|r|l|r|r}
 \mathrm{OS}&d&\text{nontrivial orders of }[g_i]
    &\text{factor pairs}&\text{accepted pairs}\\\hline
 43&2&2&2&2\\
 45&8&8&4&2\\
 46&4&4&3&2\\
 47&14&7,2&4&4\\
 49&6&2,3&4&4\\
 50&12&4,3&6&4\\
 55&20&5,4&6&4\\
 56&30&5,3,2&8&8
\end{array}
\]
Each listed order must divide at least one of \(a,b\).
Listing \(a\mid d\) and taking \(b=d/a\) now gives precisely
\eqref{maps:eq:D.F1}. These formulas also specify every rejected
pair: at a failed component order the identity entry of the
degree vector above is \(-1-[p=q]\).

The profiles and zero markers give the following complete count;
in each profile one regular position has been chosen.
\[
\begin{array}{c|l|l|r|r}
 \mathrm{OS}&\text{profiles in order}&\text{marker counts}
    &\text{accepted per profile}&\text{initial data}\\\hline
 43&\mathrm{base},\mathrm{II}&4,2&2&12\\
 45&\mathrm{base},\mathrm{II}&6,4&2&20\\
 46&\mathrm{base},\mathrm{II}&4,2&2&12\\
 47&\mathrm{base},\mathrm{II},\mathrm{III}&6,4,4&4&56\\
 49&\mathrm{base},\mathrm{II},\mathrm{III}&4,2,2&4&32\\
 50&\mathrm{base},\mathrm{II},\mathrm{IV}&4,2,2&4&32\\
 55&\mathrm{base},\mathrm{II}&6,4&4&40\\
 56&\mathrm{base},\mathrm{II},\mathrm{III},\mathrm{IV}&6,4,4,4&8&144
\end{array}
\]
There are \(2,2,2,3,3,3,2,4\) profiles in these rows.
Multiplying by the factor-pair and accepted-pair columns gives
\(104\) and \(88\), respectively; the final column sums to
\(348\). Each object is indexed by its profile, its ordered pair,
and its chosen zero position, so the enumeration has no duplicate
index and exhausts those marked choices.

Here are all the period maps in a single exact formula. Use the
ordered lattice \(\mathbb Z\langle e_1,e_2,\delta,\gamma\rangle\)
and put
\[
 p_i=ag_i,\quad q_i=bg_i,\quad
 (A_i-I)x_i=q_i,\quad \beta_i=-p_i^tJ_0A_i,\quad
 c_i=\beta_ix_i+w_i.
\]
For a multiplicative marker, \(w_i\) equals one there and zero
elsewhere. At the regular marker use all \(w_i=0\) and append
\(I+E_{\delta\gamma}\) at its marked peripheral position.
The normalized Poincar\'e line construction in
Theorem~\ref{maps:src:D.2.12} gives
\[
 T_i=\begin{pmatrix}A_i&0&q_i\\ \beta_i&1&c_i\\0&0&1\end{pmatrix},
 \qquad
 T_i^{-1}=\begin{pmatrix}
 A_i^{-1}&0&-A_i^{-1}q_i\\
 -\beta_iA_i^{-1}&1&\beta_iA_i^{-1}q_i-c_i\\
 0&0&1
 \end{pmatrix}.
\]
Both products are the identity by block multiplication.
For narrow \(Q\), take \(x_i\) integral. For narrow \(P\),
write \(p_i=(A_i-I)y_i\) with \(y_i\) integral. Symplecticity
gives
\[
 \beta_i=y_i^tJ_0(A_i-I),\qquad
 \beta_ix_i=y_i^tJ_0q_i\in\mathbb Z.
\]
This identity also proves that an \(A_i\)-fixed change of \(x_i\)
leaves \(c_i\) unchanged. Thus the displayed maps and their
inverses are integral for every accepted pair. They have
determinant one, fix \(\delta\), and preserve \(\chi=\gamma^*\).
On complex universal covers the corresponding change of marked
periods uses the identity point map; the inverse coordinate
change is the displayed inverse matrix.

For clarity, the product identity can be read directly in these
blocks. The elliptic and section cocycles give
\(A_1\cdots A_s=I\),
\(\sum_j A_1\cdots A_{j-1}q_j=0\), and
\(\sum_i\beta_iA_{i+1}\cdots A_s=0\).
The remaining central entry is
\[
 \sum_i\beta_ix_i+
 \sum_{i<j}\beta_iA_{i+1}\cdots A_{j-1}q_j+
 \sum_iw_i.
\]
The first two sums are \(-D\cdot O=-ab/d\): in the normalized
Poincar\'e frames they are the total line-frame winding around the
ordered boundary of the punctured base. The inverse boundary of
the base discs has winding minus the degree of the base line
constructed above. The marked zero contributes \(\sum_iw_i=1\),
so the central entry is zero. The regular-marker matrix commutes
with each displayed block, giving the same calculation with the
additional regular peripheral loop. Hence every ordered product
is \(I\).

The additive places are retained in this formula. The orders of
their semisimple monodromy parts are
\[
\begin{array}{c|rrrrrrr}
 \text{type}&\mathrm{II}&\mathrm{III}&\mathrm{IV}&\mathrm{IV}^*
    &\mathrm{III}^*&I_1^*&I_3^*\\\hline
 \text{order}&6&4&3&3&4&2&2.
\end{array}
\]
They are the least positive \(r\) for which \(A_i^r-I\) is
nilpotent. The eigenvalues of \(T_i\) are those of \(A_i\)
together with two ones, so the same orders apply to the full
four-period maps. The formulas specify every entry of every
tuple from the marked data.

\end{proof}

\begin{remark}[Integral vertical changes of the section line]
\hypertarget{tag-Q116}{}\label{mapsIII:rem:vertical-section-lines}
{\normalfont\scriptsize tag Q116}\par
For an integral vertical divisor $V$, the line
$M'=M\otimes\mathcal O(V)$ has linearization line
\[
 J'=J\otimes\mathcal O(V-t_Q^*V).
\]
Tensoring by $\mathcal O(-V)$ reverses this change, and
intersection with each component gives its degree change.
For the OS49 datum $P=G,Q=3G$ of
Theorem~\ref{maps:src:D.2.11}, the mixed height is $1/2$ and
the linearization line is $\mathcal O(C_0)$. At its $I_2$ fibre
the degree vector is $(-2,2)$. The whole-fibre line
$\mathcal O(C_0+C_1)=f^*\mathcal O([\infty])$ has degree
vector $(0,0)$. Multiplication of local meromorphic sections gives
$\mathcal O(C_0)\otimes\mathcal O(C_1)\simeq
\mathcal O(C_0+C_1)$. The tensor functors
$L\mapsto L\otimes\mathcal O(C_1)$ and
$L\mapsto L\otimes\mathcal O(-C_1)$ are inverse under the
evaluation identification
$\mathcal O(C_1)\otimes\mathcal O(-C_1)\simeq\mathcal O$;
they carry these two line bundles to each other.
\end{remark}

\subsection{Finite section-line data for homotopy six-spheres}

\begin{theorem}[Finite section-line data for homotopy six-spheres]
\hypertarget{tag-02KX}{}\label{maps:src:D.2.22}
{\normalfont\scriptsize tag 02KX}\par
Let \(S\to\mathbb P^1\) be one of the twenty-one marked rational
elliptic profiles of Definition~\ref{mapsIII:def:elliptic-marked-data}, with Mordell--Weil group
\(\mathbb ZG\) and \(\langle G,G\rangle=1/d\).
Put \(P=aG\), \(Q=bG\), and \(M=\mathcal O(P-O)\).
Let a nonzero holomorphic linearization
\[
 s:M_{p+Q}\longrightarrow M_p,\qquad
 J=M\otimes t_Q^*M^{-1},\qquad \operatorname{div}(s)=D\ge0
\]
define the contracting quotient. Use a zero-free original quotient
or a free affine logarithmic filling at every additive fibre, an
original quotient or a smooth Mumford filling at every multiplicative
fibre, and finitely many free logarithmic fillings at regular fibres.
The linearization is nonvanishing at every additive fibre.
At a regular zero of \(s\), use the rank-one Mumford filling.
These local types include arbitrary free logarithmic orders.

For every homotopy six-sphere obtained in this way, \(D\) is one
reduced component of one multiplicative fibre \(I_n\), and \(Q\)
acts transitively on its components. Moreover,
\begin{equation}
 \langle P,Q\rangle=\frac1n,\qquad ab=\frac dn.
 \hypertarget{tag-01YT}{}\label{maps:eq:D.FI1}
\end{equation}
Elliptic inversion simultaneously changes the signs of \(P,Q\);
choose the representative with \(a,b>0\).
There are exactly \(172\) marked triples consisting of a profile,
an ordered section pair, and a zero position satisfying the
integral component equation. The tables and component rule in the
proof specify this complete finite set. Subsequent filling choices
retain their logarithmic and gluing parameters.

For \(M=\mathcal O(P-O+V)\), with \(V\) any integral vertical
divisor, the same necessary reduction holds. Up to the simultaneous
translation of the zero component and its line data described below,
there are \(226\) marked choices of profile, ordered section pair,
and zero fibre. Every one has an integral solution \(V_0\), and
all solutions for a fixed zero component are exactly
\[
                 V_0+\operatorname{Div}_{\mathrm{vert}}(S)^Q .
\]
Here the invariant lattice consists of integral vertical divisors
with finite support. Its passage to line-bundle classes is the
quotient by principal divisors pulled back from \(\mathbb P^1\).
\end{theorem}


\begin{proof}
\begin{enumerate}[label=\arabic*.,leftmargin=*]
\item \emph{Line frames and local Euler number.}
Near \(I_n\), choose a Tate parameter \(q=t^n\).
If initially \(q=t^nu(t)\), the coordinate
\(t'=t\exp(\log u(t)/n)\), with its inverse holomorphic germ,
gives this parameter and transports the line and its units.
Represent \(P=t^p\alpha(t)\), \(0\le p<n\), with \(\alpha\ne0\).
The theta product
\[
 \theta(z;q)=(1-z)\prod_{k\ge1}(1-q^kz)(1-q^kz^{-1})
\]
gives, by reindexing,
\[
 \theta(qz;q)=-z^{-1}\theta(z;q),\qquad
 \sigma(z)=\frac{\theta(z/P;q)}{\theta(z;q)},\qquad
 \sigma(qz)=P\sigma(z).
\]
Thus \(\sigma\) is the meromorphic section of the line with
factor of automorphy \(P\), with divisor \(P-O\).
At a generic point \(z=t^i\beta(t)\) of component \(i\),
\[
 f_i=\operatorname{ord}_t\sigma=\min(0,i-p),\qquad
 f_{i+n}=f_i+p.
\]
Only \(1-z/P\) can have negative valuation in the displayed range;
the generic unit \(\beta\) avoids the section points.
For \(M=\mathcal O(P-O+\sum v_iC_i)\), the valuation of the flat
frame in an extending frame is
\[
 F_i=-f_i+v_i,\qquad F_{i+n}=F_i+A,\qquad A=-p.
\]
Let \(+Q\) shift the Tate component index by \(B\), and let \(w\)
be the valuation of the flat forward multiplier. Comparing the
source and target frames gives
\[
 z_i=\operatorname{ord}_{C_i}D=w+F_i-F_{i+B}.
\]
The sum of the frame differences is \(A\) for \(B=1\).
Composing shifts, and then reversing them, proves for every
integer \(B\) that
\begin{equation}
 \sum_{i=0}^{n-1}(F_{i+B}-F_i)=AB,\qquad
 \sum_{i=0}^{n-1}z_i=nw-AB.
 \hypertarget{tag-01YU}{}\label{maps:eq:D.FI2}
\end{equation}

The elliptic deck map \((z,v)\mapsto(qz,Pv)\) has valuation
column \((n,-A)\). The inverse of the forward line arrow has
column \((B,-w)\): its multiplier is the reciprocal of the
forward multiplier at the translated point. The coordinate map
\((z,v,t)\mapsto(z,v^{-1},t)\), whose inverse is itself, changes
these columns to \((n,A),(B,w)\), reciprocates every vertical
unit, and carries the fan by \(\operatorname{diag}(1,-1)\).
Thus the period index is \(|nw-AB|\).
For an effective nonzero local divisor it equals \(\sum_i z_i\).
A unimodular triangular cellulation of this lattice quotient has
twice as many triangles as vertices. Each triangle gives one
zero-dimensional torus stratum; every positive-dimensional torus
stratum has Euler number zero. The same count applies after the
displayed fan transport. Consequently
\begin{equation}
                         \chi(U_{I_n})=2\sum_i z_i.
 \hypertarget{tag-01YV}{}\label{maps:eq:D.FI3}
\end{equation}
For example \(n=2,p=B=1,w=0\) gives
\(F=(1,0)\), \(A=-1\), and \(z=(1,0)\), with index one.
This is the component-supported zero of Theorem~\ref{maps:src:D.2.11}.

\item \emph{Rational homology of the local fillings.}
Write \(F=T^4\), let \(T_i\) be the monodromy of a deleted
fibre disc, and put \(L_i=\partial U_i\). Set
\[
 f_i=\dim H_1(F;\mathbb Q)_{T_i},\qquad
 g_i=\dim H_2(F;\mathbb Q)_{T_i},\qquad h_i=b_1(U_i).
\]
A loop in \(U_i\) can be moved off the central support. The
punctured tube is generated by fibre loops and a meridian.
A transverse disc at a smooth point of a central component gives
the relation \(m\mu-v=0\), with \(m>0\) its multiplicity and
\(v\) a fibre loop in the collar. Over \(\mathbb Q\), therefore,
the nearby fibre generates \(H_1(U_i)\), and its conjugation
relations give
\[
                              h_i\le f_i.
\]
For a free finite quotient the same disc is read in its defining
cyclic cover; the relation includes its full order \(m\).

For a rank-one Mumford cap, the normalized components are ruled
surfaces over the residual elliptic curve. The ruling line has
degree zero. Multiplication by a topological unit frame, with
division as inverse, identifies each ruled bundle with its
product. The gluing translations and line units are joined to
the identity by translation paths and logarithms on their collars.
Inserting those paths on the normalization cylinders gives a
polygon of spheres times the residual two-torus; the reverse
paths give the inverse maps. The specialized fibre area, the
two products of the polygon cycle with a residual circle, and
the residual torus area are four independent classes. They
are the images of a basis of \(H_2(F;\mathbb Q)_{T_i}\).
Thus restriction
\[
                  H^2(U_i;\mathbb Q)\longrightarrow
                  H^2(F;\mathbb Q)^{T_i}
\]
is onto.

For a rank-two triangular cap, write its period lattice as
\(L\subset\mathbb Z^2\). On the universal toric space its periods are
\[
 \Psi_\ell(X_1,X_2,t)=
 (u_{\ell,1}(t)t^{\ell_1}X_1,
  u_{\ell,2}(t)t^{\ell_2}X_2,t),\qquad \ell\in L.
\]
The unit map is a character of \(L\). Smith coordinates for
\(L\hookrightarrow\mathbb Z^2\), followed by holomorphic logarithms
and roots of its two unit values, extend this character to
\(\mathbb Z^2\). Both actions use the same triangular fan.
The identity of the universal space induces a finite period
cover from the given cap to the minimal hexagonal cap.
If a nonzero deck coset fixed a central point, a nonzero lattice
translation would preserve its bounded toric cell. This is
impossible: iterates would make that cell unbounded.
On the punctured part the dominant matrix \(\log|t|\,I\), plus
the bounded logarithms of the unit character, is invertible
on a sufficiently small disc, proving freeness there as well.
The local inverse of the covering is the unchanged chart on
the chosen sheet.

For the minimal quotient \(U_0\), the conductor maps of
Theorem~\ref{maps:src:D.1.2} at \(d=1\) give
\[
 (H_0,H_1,H_2,H_3,H_4)(U_0;\mathbb Q)
       =(\mathbb Q,\mathbb Q^2,\mathbb Q^4,\mathbb Q^2,\mathbb Q).
\]
The cap retracts onto its compact normal-crossing core. The
face maps are the branch inclusions and unit multiplications
of complex degree one, so these groups apply to the unit
character just constructed. The boundary Wang sequence has
\[
 b_1(L_0)=3,\qquad b_2(L_0)=6,\qquad f_0=2,\qquad g_0=4.
\]
Its map to the cap kills the two short loops and the normal
meridian, and preserves the two long loops; lifting the long
periods gives the inverse on their images.
Poincar\'e--Lefschetz duality gives
\(H_2(U_0,L_0;\mathbb Q)=H^4(U_0;\mathbb Q)=\mathbb Q\).
The positive normal disc maps onto the meridian kernel in
\(H_1(L_0)\), so \(H_2(L_0)\to H_2(U_0)\) is onto.
Sweep this disc by the two short angular circles. Their boundary
tori bound the parametrized \(D^2\times S^1\) sweeps in \(U_0\),
and project to a basis of \(\ker(T_0-I)\) in the Wang sequence.
They are independent, hence span the two-dimensional kernel of
the cap map. They complement the fibre summand, proving
\[
                  H_2(F;\mathbb Q)_{T_0}
                    \xrightarrow{\;\cong\;}H_2(U_0;\mathbb Q).
\]
The inverse lifts a cap class to \(L_0\) and subtracts its two
normal sweeps. Dualizing gives the four specialization classes.
Pullback through the finite period cover preserves their
independence: its regular-fibre map is a finite torus covering
and a rational cohomology isomorphism intertwining monodromy.
A subdivision has the proper toric map induced by the identity
lattice and the refined fan; it is the identity on \(F\), so
its pullback preserves these classes as well.

For a free finite affine quotient, average a lifted class over
the finite deck group. Rational transfer identifies the average
with a class on the quotient. For a base rotation, the deck
action on the cover's monodromy-invariant space is induced by
the original monodromy; its invariants are exactly those of
the quotient boundary. This proves the same restriction
surjectivity for the good-reduction caps, the semistable covers
of odd starred fibres, and every permitted logarithmic order.
The regular good-reduction cover is \(D^2\times T^4\), where
restriction is already an isomorphism.

\item \emph{The original caps and their component degrees.}
Use the original elliptic-surface neighbourhood \(N\), its line
\(M\), and its component translation. The unit-circle bundle of
\(M\) over the translation mapping torus models the cap.
Indeed, write the contraction as
\((y,w)\mapsto(\tau y,a(y)w)\).
Its positive roof is \(r(y)=-\log|a(\tau^{-1}y)|\).
For \(k\max r<1\), choose a nonnegative middle-supported function
\(b_y\) on \([0,r(y)]\), with integral one. The increasing map
\[
 t\longmapsto\int_0^t
       [k+(1-kr(y))b_y(s)]\,ds
\]
takes this interval to \([0,1]\) and has derivative \(k\) at
both end collars. Its increasing inverse straightens the
mapping-torus roof smoothly. The angular gluing retains
\(a/|a|\); the inverse uses the reciprocal multiplier at the
inverse translated point.

The central-curve inclusion in \(N\) is a homotopy equivalence.
Normal and crossing collapses on a local embedded resolution,
followed by contraction of its exceptional trees, give this
map for the nodal, tangent, and triple-point curves. Inclusion
commutes with translation, so it also gives a homotopy equivalence
of the mapping tori. Write \(B=T(t_{-Q}:N\to N)\).
The additive curve has \(H_1(N;\mathbb Q)=0\); a multiplicative
curve has the oriented polygon cycle \(H_1(N;\mathbb Q)=\mathbb Q\).
The reduced component classes form a basis of \(H_2(N;\mathbb Q)\),
and \(H_j(N)=0\) for \(j\ge3\). If there are \(\rho\) component
orbits, the Wang sequence gives \(b_3(B)=\rho\) and
\[
 b_1(B)=1\quad\hbox{at an additive fibre},\qquad
 b_1(B)=2\quad\hbox{at a multiplicative fibre}.
\]

The multidegree row \(\mathbf c\) of \(M\) is invariant under \(Q\),
since the linearization is a unit. If \(\mathbf c\ne0\), the Euler
class of the circle bundle over \(B\) is nonzero on a component,
and its Gysin sequence gives \(h=b_1(B)\).
If \(\mathbf c=0\), choose a holomorphic frame of the degree-zero
line on each normalized rational component. The node
identifications and unit translation lifts are constant
nonzero scalars in these frames. They give a flat line on
the mapping-torus core, including its polygon--orbit two-cycle.
For the tangent and cuspidal models, the normalization and
intersection-point quotient give the same topological line.
Thus the rational Euler class is zero. The permutation Wang
sequence makes \(H^2(B;\mathbb Z)\) free, so it is also zero
integrally. The Gysin sequence then gives \(h=b_1(B)+1\).

At a zero-free additive place the fibre coinvariant dimension
is \(f=2\); at \(I_n\) it is \(f=3\).
Here is the central-entry check also for a vertical correction.
Choose an iterate \(NQ\) preserving every component. Tensoring
the first \(N\) arrows cancels \(V\) exactly and leaves the
natural line \(\mathcal O(P-O)\) with the unit linearization
for \(NQ\). In the Poincar\'e period formula of
Proposition~\ref{maps:src:D.2.19}, its central entry is
\(\beta x\), where \((A_{\rm ell}-I)x=Nb g\);
the additional valuation is zero because this arrow is a unit.
For an additive fibre put \(\ell=(A_{\rm ell}-I)^{-1}g\).
The invertible rational coordinate change
\[
 x'=x+b\ell\gamma,\qquad
 \delta'=\delta-a\ell^tJx,\qquad\gamma'=\gamma
\]
has inverse \(x=x'-b\ell\gamma'\),
\(\delta=\delta'+a\ell^tJx'\).
It takes the one-step central entry to
\(s=c_{\mathrm{per}}-ab \ell^tJg\), where
\(c_{\mathrm{per}}\) is the line--orbit matrix entry,
and the iterated entry to \(Ns\).
The unit iterated formula gives \(Ns=0\), hence \(s=0\).
The elliptic block has no eigenvalue one, proving \(f=2\).
At \(I_n\), the same iterated unit formula and
\eqref{maps:eq:D.FI2} give a rank-one full monodromy difference;
its elliptic transvection is nonzero. This proves \(f=3\).
Consequently \(h=f-1\) when \(\mathbf c\ne0\), and
\(h=f\) when \(\mathbf c=0\).

In the latter case Gysin gives \(b_4(U)=b_3(B)=\rho\).
Since \(H_1(L)\to H_1(U)\) is onto, the pair sequence and duality
give the exact defect
\[
 \dim\operatorname{coker}[H_2(L;\mathbb Q)\to H_2(U;\mathbb Q)]
       =b_4(U)-b_1(L)+b_1(U)=\rho-1.
\]
Geometrically, pair component two-cycles with the orbit sums
of component four-cycles. Their pairing is the Kodaira
intersection form on the invariant component subspace.
Its radical is the whole-fibre multiplicity vector, and its
rank is \(\rho-1\).
When \(\rho=1\), the weighted fibre area and the two-circle
torus are independent specialized fibre classes. At \(I_n\)
the two products of its oriented polygon cycle with the
circle factors give two further independent classes.
To construct them, take the finite mapping-torus cover fixing
components and insert paths of its component translations
and logarithms of its constant line gluing on the cylinders.
The reverse paths give the inverse product maps. Transfer
descends the classes to \(U\). Their ranks are respectively
two and four, exactly \(g\); thus the restriction to
\(H^2(F;\mathbb Q)^T\) is onto in every one-orbit case.

\item \emph{The global rank bound and the support of \(D\).}
If \(a=0\) or \(b=0\), an iterate preserving all components makes
the effective divisor a sum of whole fibres of degree
\(N\langle P,Q\rangle=Nab/d=0\). Every coefficient is then zero.
All the remaining local types have Euler number zero, so these
cases give no homotopy sphere. Hence take \(ab\ne0\).

Delete \(N\) fibre discs containing all singular fibres, zeros,
and logarithmic sites, and denote the exterior by \(E\).
The group-law function identifies \(\mathcal O(aG-O)\) with
\(\mathcal O(G-O)^{\otimes a}\) times a vertical line; its
reciprocal is the inverse. The canonical section of that
vertical line is nonzero over the punctured base.
In this common frame the monodromy has blocks
\[
 T_i=\begin{pmatrix}
 B_i&0&b g_i\\-a g_i^tJB_i&1&c_i\\0&0&1
 \end{pmatrix},\qquad
 J=\begin{pmatrix}0&1\\-1&0\end{pmatrix},\qquad
 B_i\in SL_2(\mathbb Z).
\]
The entry \(c_i\) includes the central winding of the line.
For the twenty-one marked tuples the concatenated three-row
matrix with blocks \((B_i-I;-g_i^tJB_i)\) has rank three.
For precision, the following table gives a nonzero
three-column determinant for each tuple, with columns numbered
from zero in their peripheral order:
\[
\begin{array}{c|c|r@{\qquad}c|c|r}
43&(0,1,5)&1&43/\mathrm{II}&(0,1,3)&1\\
45&(1,3,4)&3&45/\mathrm{II}&(1,2,3)&-3\\
46&(0,1,3)&-2&46/\mathrm{II}&(0,1,2)&1\\
47&(1,2,4)&-3&47/\mathrm{II}&(1,2,4)&-3\\
47/\mathrm{III}&(1,2,4)&-3&49&(1,2,3)&-1\\
49/\mathrm{II}&(0,1,3)&1&49/\mathrm{III}&(0,1,3)&-1\\
50&(0,1,3)&2&50/\mathrm{II}&(0,1,2)&1\\
50/\mathrm{IV}&(0,1,2)&-1&55&(1,3,4)&3\\
55/\mathrm{II}&(1,2,4)&-1&56&(1,3,4)&-4\\
56/\mathrm{II}&(0,2,3)&1&56/\mathrm{III}&(1,2,4)&2\\
56/\mathrm{IV}&(1,2,4)&1&&&
\end{array}
\]
Each is computed from the marked \(B_i,g_i\) by the displayed
three-row formula; its inverse is its adjugate divided by that
nonzero determinant. Scaling the last row by \(a\ne0\) preserves
rank. All full differences have zero fourth coordinate, so
their common degree-one coinvariant has dimension one.

Put \(v_2=\dim H_2(F;\mathbb Q)_{\pi_1(E)}\).
Direct multiplication gives
\[
 T_i^t\Omega T_i=\Omega,\qquad
 \Omega=\operatorname{diag}
       \left(aJ,\begin{pmatrix}0&-b\\b&0\end{pmatrix}\right).
\]
Before imposing \(\det B_i=1\), the residual is
\(a(\det B_i-1)\operatorname{diag}(J,0_2)\);
in particular all central and cross terms cancel for every
\(c_i,g_i\). The nonzero invariant two-form proves \(v_2\ge1\).
The punctured base retracts to a graph with \(N-1\) loops.
Its local-system chain map on \(H_1(F)\) has rank three, and
its graph spectral sequence therefore gives
\[
 b_1(E)=N,\qquad b_2(E)=v_2+4N-7.
\]
The boundary Wang sequences give
\(b_1(L_i)=f_i+1\), \(b_2(L_i)=g_i+f_i\).

For a rational homology six-sphere the first Mayer--Vietoris
map is onto, with kernel dimension \(\sum_i(f_i-h_i)\).
This kernel is a quotient of \(H_2(X;\mathbb Q)=0\).
Since each \(h_i\le f_i\), equality holds at every place.
The original-cap calculation forces zero line multidegree.
The degree-two map is onto each cap as well, so every original
cap has \(\rho=1\). All caps now have the restriction
surjectivity proved above, and hence \(b_2(U_i)\ge g_i\).
After projecting onto their degree-two groups, the remaining
boundary kernel has dimension at most \(\sum_i f_i\).
It must map onto \(H_2(E)\), yielding
\begin{equation}
 v_2+4N-7\le\sum_i f_i,\qquad
 \sum_i\operatorname{rank}(T_i-I)\le7-v_2\le6.
 \hypertarget{tag-01YW}{}\label{maps:eq:D.FI4}
\end{equation}

Let \(F_a,F_m\) count additive and multiplicative fibres of \(S\).
Shioda--Tate gives root rank seven, since \(\rho(S)=10\) and
the Mordell--Weil rank is one. The Euler number is twelve.
Subtracting the component root ranks from the Euler numbers
of the Kodaira fibres gives
\[
                             2F_a+F_m=5;
\]
see \cite[Theorem 6.3, p.~23; Corollary 6.13,
pp.~29--30; Section 8.3, p.~39]{manSS10}.
The listed fibre configurations also give these identities
directly. Additive monodromy differences have rank at least
two; multiplicative ones have rank at least one.
A multiplicative zero raises the rank to two by its positive
determinant in \eqref{maps:eq:D.FI2}. A regular zero supplies
an additional rank one.
Original zero-free caps have a free angular line circle:
a nonzero iterate of a strict contraction changes the norm,
so cannot equal a circle rotation. Their Euler number is zero.
Rank-one caps have elliptic strata of Euler number zero,
and the good-reduction and free finite cases have the same
property by their finite covering maps. Additivity over
the base and \eqref{maps:eq:D.FI3} show that Euler number two
requires at least one multiplicative zero. It consumes the
single extra rank permitted by \eqref{maps:eq:D.FI4}.
There is therefore one such place and no regular zero.
Equation \eqref{maps:eq:D.FI3} makes its total coefficient one:
\(D=C_h\) is a reduced component.

\item \emph{Iteration, height, and the sign map.}
Choose an integer \(N>0\) for which \(NQ\) meets the identity
component of every reducible fibre. Tensoring the first \(N\)
arrows gives
\[
 J_N=M\otimes t_{NQ}^*M^{-1},\qquad
 D_N=\sum_{k=0}^{N-1}t_{kQ}^*D.
\]
In frames the multiplier is the product of the original
ones at \(p,p+Q,\ldots,p+(N-1)Q\); reciprocal evaluation
gives the inverse tensor map.
Because \(t_{NQ}\) preserves every component, \(J_N\) has
zero component degrees. The kernel of the full Kodaira
intersection matrix is the positive multiplicity vector
with a component equal to one. Thus an integral vertical
divisor of zero degrees is an integral whole-fibre multiple.
The effective \(D_N\), supported on the orbit of \(C_h\),
can be such a multiple only when that orbit consists of
every component of \(I_n\). Consequently \(Q\) is transitive
and \(D_N=(N/n)F\).

The factors of \(V\) cancel because \(t_{NQ}^*V=V\).
The group-law line map and the narrow height formula of
Proposition~\ref{maps:src:D.2.19} identify the base degree of \(J_N\)
with \(\langle P,NQ\rangle=Nab/d\); restriction to the zero
section is the inverse on base-line degrees. All local mixed
corrections vanish for the narrow second section. Comparing
degrees proves \eqref{maps:eq:D.FI1}, with its positive sign.
Elliptic inversion \(\iota(p)=-p\), whose inverse is itself,
pulls back the line and arrow to
\[
 (P,Q,V,D)\longmapsto(-P,-Q,\iota^*V,\iota^*D).
\]
Its fibre-period map is \(\operatorname{diag}(-I_2,1,1)\).
It transports the multiplier, fan and collars with their
inverse maps. Since \(ab>0\), every simultaneous sign orbit
contains exactly one positive ordered pair.

\item \emph{Integral component equations.}
Let \(C_v\) be the full component intersection matrix at a
fibre, including the zero component. Let \(\sigma_v\) be
the component action of \(+Q\), and define its coefficient
pullback by \((T_vx)_j=x_{\sigma_v(j)}\).
The forward line arrow covers \(-Q\), with inverse
permutation on points. Translation preserves intersection,
so \(C_vT_v=T_vC_v\).
For a simple component index \(p\), put \(m(p)=e_p-e_0\).
The component degrees of \(\mathcal O(P-O)\) are \(m(p)\).
At \(I_s\), with \(P,Q\) indices \(p,q\), their difference
line has degree vector
\[
 d_j=[j=p]-[j=0]-[j=p-q]+[j=-q].
\]
At every fibre \(d=(1-T)m(p)\). Thus the full equation
for the divisor vector \(z\) and integral vertical vector \(V\) is
\begin{equation}
                   (1-T)CV=Cz-(1-T)m(p).
 \hypertarget{tag-01YX}{}\label{maps:eq:D.FI5}
\end{equation}
At the chosen zero fibre \(z=e_h\); at every other place \(z=0\).

The component maps follow directly from the marked periods.
At \(I_s\), write
\[
 A=I+suu^tJ,\qquad g=\kappa u,\qquad
 \det(u,r)=1.
\]
The row \((r_2,-r_1)\) takes \(u\) to one and \(r\) to
zero; its inverse on the torsion component takes one to \(u\).
Hence \(p=a\kappa\), \(q=b\kappa\) modulo \(s\).
These cocycle indices are negative Tate valuation: the section
cocycle is the negative elliptic-logarithm increment. The
comparison \(j\mapsto-j\), its own inverse, carries the Tate
component matrix to this cyclic component matrix. In this
marking \(+Q\) sends \(j\) to \(j+q\), and the forward arrow
sends \(j\) to \(j-q\).
At an additive fibre the component quotient is
\(\mathbb Z^2/(A-I)\mathbb Z^2\).
An onto row \(h\) modulo \(s=|\det(A-I)|\), satisfying
\(h(A-I)=0\), identifies it with \(\mathbb Z/s\);
an \(r\) with \(hr=1\pmod s\) supplies its inverse.
The relation lattice has index \(s\), proving both inverse
quotient maps. The generator index is \(hg\), in the
component marking of the profile.

For clarity the matrices and component permutations used in
\eqref{maps:eq:D.FI5} have the following complete description.
At \(I_s\),
\((Cx)_j=-2x_j+x_{j-1}+x_{j+1}\), counting both neighbours
when \(s=2\); \(C=0\) for \(s=1\).
Types III and IV have the same two- and three-component
intersection matrices on their respective divisors.
For IV*, use vertices \(O_0,O_1,O_2,B_0,B_1,B_2,C\),
edges \(O_iB_i,B_iC\), and multiplicities
\((1,1,1,2,2,2,3)\); a component generator cycles the arms.
For III*, use \(O_0,O_1,A_0,A_1,B_0,B_1,C,S\),
chains \(O_i-A_i-B_i-C\), and edge \(CS\), with
multiplicities \((1,1,2,2,3,3,4,2)\); the generator
swaps the long arms and fixes \(C,S\).
For \(I_k^*\), \(k=1,3\), use
\(L_0,R_0,L_1,R_1,B_0,\ldots,B_k\), attach
\(L_0,L_1\) to \(B_0\), \(R_0,R_1\) to \(B_k\),
and join successive \(B_i\).
The generator cycles \((L_0,R_0,L_1,R_1)\) and reverses
the central chain. Outer multiplicities are one and the
central multiplicities two. Every tree matrix has diagonal
\(-2\) and one on each edge. These are the full affine
diagrams obtained by restoring the marked zero component;
see \cite[Section 6.5 and Figure 6, pp.~25--26;
Sections 7.2--7.6, pp.~32--34]{manSS10}.

\item \emph{The natural pairs and their line isomorphisms.}
At an unselected fibre \(d=0\) precisely when \(p=0\) or
\(q=0\). If both are nonzero, its identity entry is
\(-1-[p=q]\), which proves necessity.
At the selected \(I_n\), \(q\) is a unit. For \(n>1\),
the negative identity entry in \(d=Ce_h\) forces \(h=0\);
its value \(-2\) then forces \(p=q\ne0\).
The two positive entries must be the two neighbours of zero,
including their doubled multiplicity at \(I_2\). Thus
\[
                         p=q=1\ \hbox{or}\ -1\pmod n.
\]
For \(n=1\) the equation has its unique component.

The following table gives every positive ordered pair for
each selected fibre type in the indicated OS family. Each
occurs at every such marked position of that profile listed
in the final table.
\[
\begin{array}{c|c|l}
\mathrm{OS}&\text{zero type}&(a,b)\\\hline
43&I_1&(1,2),(2,1)\\
45&I_1&(1,8),(8,1)\\
46&I_1&(1,4),(4,1)\\
47&I_1&(1,14),(2,7),(7,2),(14,1)\\
47&I_2&(1,7),(7,1)\\
49&I_1&(1,6),(2,3),(3,2),(6,1)\\
49&I_2&(1,3),(3,1)\\
50&I_1&(1,12),(3,4),(4,3),(12,1)\\
50&I_3&(1,4),(4,1)\\
55&I_1&(1,20),(4,5),(5,4),(20,1)\\
55&I_4&(1,5),(5,1)\\
56&I_1&(1,30),(2,15),(3,10),(5,6),
            (6,5),(10,3),(15,2),(30,1)\\
56&I_2&(1,15),(3,5),(5,3),(15,1)\\
56&I_3&(1,10),(2,5),(5,2),(10,1)\\
56&I_5&(1,6),(6,1)
\end{array}
\]
To obtain this table, list \(a\mid d/n\), set \(b=d/(na)\),
and apply the displayed component criterion at every fibre.
Narrowness at an additive fibre is integrality of
\((A-I)^{-1}ag\) or \((A-I)^{-1}bg\); at \(I_s\) it is
\(s\mid a\kappa\) or \(s\mid b\kappa\).
The nontrivial component orders, by increasing OS family,
are \(2;8;4;(7,2);(3,2);(4,3);(5,4);(5,3,2)\).
These divisibilities and the selected adjacency condition
give exactly the table. In particular the other multiplicative
types give no additional natural pairs.

The component equations imply a global line isomorphism.
On the generic Weierstrass fibre, divide the line through
\(P,-Q\) by the vertical line through \(P-Q\).
Use the tangent in the coincident case; when \(P=Q\) use
\(x-x(P)\). This rational surface function \(F_0\) has
\[
 \operatorname{div}(F_0)=D_0+W_{\mathrm{div}},\qquad
 D_0=P-O-(P-Q)+(-Q),\qquad
 W_{\mathrm{div}}\ \text{integral vertical}.
\]
For \(V=0\), \eqref{maps:eq:D.FI5} gives
\(W_{\mathrm{div}}+C_0=\sum_v k_vF_v\).
The norm and \eqref{maps:eq:D.FI1} give \(\sum_v k_v=0\).
Choose a rational base function \(B\) with divisor
\(\sum_vk_v[v]\). Then
\[
 \rho=\frac{F_0}{B\circ f},\qquad
 \operatorname{div}(\rho)=D_0-C_0,\qquad
 \mathcal O(D_0)\xrightarrow{\ \rho\ }\mathcal O(C_0)
\]
is an isomorphism, with inverse multiplication by \(\rho^{-1}\).
Pull back the canonical section of \(\mathcal O(C_0)\).
It has exactly the required holomorphic zero. A sufficiently
small nonzero scalar makes its line map contracting on the
compact surface.

\item \emph{Integral vertical solutions.}
For a full component matrix \(C\), define its rational potential by
\[
 Cz(p)=e_p-e_0,\qquad z(p)_0=0.
\]
Delete row and column zero. The remaining negative Cartan
matrix is invertible, and its inverse gives these coordinates;
the omitted equation follows by pairing with the multiplicity
kernel, whose zero entry is one. At \(I_s\) this is explicitly
\[
 z(p)_j=-\frac{\min(p,j)(s-\max(p,j))}{s},\qquad
 z(p)_j\equiv\frac{pj}{s}\pmod{\mathbb Z}.
\]
For each selected \(I_n\) in the profile table,
\[
 \gcd(n,d/n)=1,\qquad (d/n)\kappa^2\equiv1\pmod n.
\]
For the nontrivial cyclic places the values
\((\mathrm{OS},n,\kappa)\) are
\[
\begin{gathered}
 (45,8,3),\ (47,7,5),\ (47,2,1),\ (49,2,1),\\
 (50,3,1\text{ or }2),\ (55,5,2),\
 (55,4,1\text{ or }3),\\
 (56,5,1),\ (56,3,1\text{ or }2),\ (56,2,1).
\end{gathered}
\]
The alternatives reflect their marked primitive elliptic bases.
Substitution proves both identities, so every factor pair
\(ab=d/n\) has \(q=b\kappa\) a unit and \(pq=1\pmod n\).
For \(n>1\), define
\[
 W_0=0,\qquad W_{kq}=\frac{k}{n}-1\ (1\le k<n),\qquad
                        V=W-z(p).
\]
Its recurrence is \((1-T)W=e_0-\frac1n\mathbf1\).
Reading successive values around the \(q\)-orbit is its
inverse with \(W_0=0\). The congruence \(pq=1\) and the
potential formula show that \(V\) is integral.
Applying \(C\) proves \eqref{maps:eq:D.FI5}.
For \(n=1\), take \(V=0\).

At an unselected fibre choose one anchor \(j_{\mathcal O}\)
in each \(Q\)-orbit and set
\[
 V_j=z(p)_{j_{\mathcal O}}-z(p)_j,\qquad j\in\mathcal O.
\]
Then \(z(p)+V\) is constant on each orbit, which proves
the zero right side of \eqref{maps:eq:D.FI5}.
The other component orders divide \(ab=d/n\).
At prime orders this forces \(p=0\) or \(q=0\), and \(V=0\).
The complete nonzero unselected possibilities and solutions are
\[
\begin{array}{c|c|l}
\text{type}&(p,q)&V\\\hline
I_k^*,\ k=1,3&(2,2)&(0,0,1,0,0,\ldots,0)\\
I_4&(2,2)&(0,0,1,0)\\
I_8&(6,4)&(0,0,0,0,1,1,1,0)\\
I_8&(4,6)&(0,0,1,1,2,1,1,0)
\end{array}
\]
For the star rows the potential at \(p=2\) is
\((0,-1/2,-1,-1/2,-1,\ldots,-1)\) in the component
order above. The orbit formula gives the printed integer
vector. The cycle formula gives the other rows.
The complete nonzero selected solutions are
\[
\begin{array}{c|c|l}
I_5&(2,3)&(0,0,1,0,0)\\
I_5&(3,2)&(0,0,0,1,0)\\
I_7&(5,3)&(0,0,0,0,1,1,0)\\
I_7&(3,5)&(0,0,1,1,1,0,0)\\
I_8&(3,3)&(0,0,1,1,1,1,0,0)
\end{array}
\]
In every remaining local case this construction gives \(V=0\).
Thus the displayed vectors prove integral solvability for every
ordered factor pair.

For their sum \(V\), the same group-law function satisfies
\[
 W_{\mathrm{div}}-(V-t_Q^*V)+C_0=\sum_v k_vF_v.
\]
The norm cancels \(V\) and gives
\(Nab/d-N/n=0\), so \(\sum_v k_v=0\).
Choose the corresponding rational base function \(B\).
Now multiplication by \(\rho=F_0/(B\circ f)\), with reciprocal
inverse, gives
\[
 \operatorname{div}(\rho)=D_0+(V-t_Q^*V)-C_0,\qquad
                     J\xrightarrow{\ \rho\ }\mathcal O(C_0).
\]
The pulled canonical section and a small nonzero scalar
give the desired contracting arrow.

\item \emph{Translation of the zero and complete solution freedom.}
For a zero \(C_h\), choose \(R=kQ\) with \(kq=h\pmod n\).
The translation \(t_R\), with inverse \(t_{-R}\), commutes
with \(t_Q\) and pulls \(C_h\) back to \(C_0\).
Its pulled line is
\(\mathcal O((P-R)-(-R)+t_R^*V)\).
The divisor
\[
                  A_R=(P-R)-(-R)-P+O
\]
is principal on the generic fibre. Take for its function
the reciprocal of the group-law chord quotient for
\(P-O-(P-R)+(-R)\), with tangent and vertical limits.
For \(R=0\) take \(F_R=1\); for \(R=P\) take
\(F_R=(x-x(P))^{-1}\).
Write \(\operatorname{div}(F_R)=A_R+W_R\), with integral
vertical \(W_R\). Multiplication by \(F_R\) gives
\[
 t_R^*M\xrightarrow{\ F_R\ }
 \mathcal O(P-O+V_R),\qquad V_R=t_R^*V-W_R,
\]
and multiplication by \(F_R^{-1}\) is its inverse.
The transported multiplier is
\[
             s_R(p)=F_R(p)(t_R^*s)(p)F_R(p+Q)^{-1}.
\]
Its divisor is \(C_0\), with the same ordered sections
\(P,Q\). A change of \(F_R\) by a rational base function
adds its principal whole-fibre divisor to \(W_R\);
the corresponding factor cancels in the difference line
since \(t_Q\) fixes the base. The inverse base translation
and the pulled reciprocal reverse the whole map.
The local charts, units, fan and collars of any chosen
filling are transported by these maps and their inverses.
This is the equivalence used in the enlarged count.
For a direct solution at a fixed \(h\), replace \(W_j\) by
\(W_{j-h}-W_{-h}\) in the preceding orbit formula; its
fractional part remains \(q^{-1}j/n\), so the solution
is integral at that marked component as well.

For two solutions at the same zero put \(Z=V'-V\).
Their ratio makes \(\mathcal O(Z-t_Q^*Z)\) trivial.
A rational function with vertical divisor has neither
zeros nor poles on the generic complete elliptic fibre,
and hence comes from the base. Thus
\[
                       Z-t_Q^*Z=\sum_v k_vF_v.
\]
Summing coefficients on each component orbit gives zero
on the left and \(k_v\) times a positive sum of
multiplicities on the right. Every \(k_v\) is zero,
so \(t_Q^*Z=Z\).
Conversely an invariant divisor cancels in \(J\) by its
canonical tensor identification and inverse. The solution
set is therefore exactly
\(V_0+\operatorname{Div}_{\mathrm{vert}}(S)^Q\).
An integral basis is the sum of the components in each
orbit; its inverse coordinate reads one coefficient at
an orbit representative. This also proves saturation.
Whole regular fibres and principal base divisors are
included in this description.

\item \emph{The complete marked counts.}
For each marked \(I_n\) position, the enlarged ordered
pairs are precisely \((a,d/(na))\) with \(a\mid d/n\).
The preceding constructions prove existence for every pair,
and \eqref{maps:eq:D.FI1} proves exhaustion.
The natural pairs are exactly the earlier ordered-pair
table. In the following table \(i:n\) means the \(i\)-th
fibre in peripheral order, of type \(I_n\); equal fibre
types at different positions remain distinct.
\[
\begin{array}{c|r|l|r|r}
\text{profile}&d&\text{multiplicative markers }i:n&
       V=0&\text{integral }V\\\hline
43&2&2:1,\ 3:1,\ 4:1&6&6\\
43/\mathrm{II}&2&3:1&2&2\\
45&8&1:8,\ 2:1,\ 3:1,\ 4:1,\ 5:1&8&17\\
45/\mathrm{II}&8&1:8,\ 3:1,\ 4:1&4&9\\
46&4&2:1,\ 3:1,\ 4:1&6&9\\
46/\mathrm{II}&4&3:1&2&3\\
47&14&1:7,\ 2:2,\ 3:1,\ 4:1,\ 5:1&14&16\\
47/\mathrm{II}&14&1:7,\ 2:2,\ 3:1&6&8\\
47/\mathrm{III}&14&1:7,\ 2:1,\ 3:1&8&10\\
49&6&1:2,\ 3:1,\ 4:1&10&10\\
49/\mathrm{II}&6&3:2&2&2\\
49/\mathrm{III}&6&3:1&4&4\\
50&12&2:3,\ 3:1,\ 4:1&10&15\\
50/\mathrm{II}&12&2:3&2&3\\
50/\mathrm{IV}&12&3:1&4&6\\
55&20&1:5,\ 2:4,\ 3:1,\ 4:1,\ 5:1&14&23\\
55/\mathrm{II}&20&1:5,\ 2:4,\ 4:1&6&11\\
56&30&1:5,\ 2:3,\ 3:2,\ 4:1,\ 5:1&26&28\\
56/\mathrm{II}&30&1:2,\ 3:5,\ 4:3&10&12\\
56/\mathrm{III}&30&1:5,\ 2:3,\ 4:1&14&16\\
56/\mathrm{IV}&30&1:5,\ 2:2,\ 4:1&14&16\\\hline
\text{total}&&&172&226
\end{array}
\]
Each row is obtained by counting the displayed positive
divisors at its displayed positions and, for \(V=0\),
applying the proven component equations at every fibre.
This proves both counts and the stated marking conventions.
\end{enumerate}
\end{proof}

\subsection{Fundamental groups of the natural section-line families}

\begin{theorem}[Fundamental groups for all logarithmic orders]\hypertarget{tag-02KY}{}\label{maps:src:D.2.23}
{\normalfont\scriptsize tag 02KY}\par
For every one of the $172$ natural initial rows of
Theorem~\ref{maps:src:D.2.22}, use any of the stated original or Mumford
fillings and perform any admissible free affine logarithmic replacements.
At an odd $I_n^*$, a logarithmic replacement means the free cyclic
quotient of its rank-one semistable Mumford cover. Regular logarithmic
sites are allowed. Let $m_i\ge1$ be the orders, let $v_i$ be the integral
power vectors in $H_1(T^4;\mathbb Z)$, and put $n_i=\gamma^*(v_i)$.
For a chosen system of peripheral lifts, let $e$ be the integer in
the global relation of their surviving $\gamma$ coordinates. Then
\begin{equation}
 \pi_1(X)=\left\langle \gamma,c_1,\ldots,c_s\ \left|\
 [\gamma,c_i]=1,\ c_i^{m_i}=\gamma^{n_i},\
 c_1\cdots c_s=\gamma^e\right.\right\rangle .                 \hypertarget{tag-01YQ}{}\label{maps:eq:D.FG1}
\end{equation}
The canonical lifts obtained by unrolling the translation quotient
before filling have $e=0$. Replacing $c_i$ by $c_i\gamma^{t_i}$
replaces $(n_i,e)$ by $(n_i+m_it_i,e+\sum_i t_i)$ and gives an
isomorphic presentation. The first integral homology is the Smith
cokernel of the $(s+1)$-by-$(s+1)$ matrix with rows
\begin{equation}
 (-n_i,0,\ldots,m_i,\ldots,0)\quad(1\le i\le s),
 \qquad (-e,1,\ldots,1).                                    \hypertarget{tag-01YR}{}\label{maps:eq:D.FG2}
\end{equation}
The formula uses the actual power vectors of the admissible affine
translations for every positive order.

Eliminate order-one sites by $c_i=\gamma^{n_i}$ and replace $e$
by $e-\sum_{m_i=1}n_i$. With at least three remaining sites the
group is nontrivial. With zero, one, or two remaining sites it
is abelian and is trivial precisely when, respectively,
\begin{equation}
 |e|=1,\qquad |m_1e-n_1|=1,\qquad
 |m_2n_1+m_1n_2-em_1m_2|=1.                                \hypertarget{tag-01YS}{}\label{maps:eq:D.FG3}
\end{equation}
In the two-site case put
\[
 R=\begin{pmatrix}-n_1&m_1\\m_2e-n_2&-m_2\end{pmatrix},
 \qquad D=\det R,\qquad d=\gcd(m_1,m_2,n_1,n_2).
\]
For $D\ne0$ the invariant factors are $d,|D|/d$; for $D=0$
the group is $\mathbb Z\oplus\mathbb Z/d$.
Simple connectivity forces pairwise coprime remaining
multiplicities and $\gcd(m_i,n_i)=1$ at every remaining site.
\end{theorem}


\begin{proof}
Use the lattice and ordered profile data of
Definition~\ref{mapsIII:def:elliptic-marked-data}. Put
\[
 J=\begin{pmatrix}0&1\\-1&0\end{pmatrix}.
\]
The natural ordered pairs, zero components and their contracting
line isomorphisms are those of Theorem~\ref{maps:src:D.2.22}.
At the chosen $I_n$ write
\[
 A=I+n uu^tJ,\qquad g=\kappa u,\qquad \det(u,r)=1.
\]
The component criterion in Theorem~\ref{maps:src:D.2.22} gives
$ab=d/n$ and, for $n>1$, $a\kappa=b\kappa=\pm1\pmod n$.
For $n=1$ there is one component. At every other reducible fibre
one of $P,Q$ is narrow. The same theorem provides the complete
172 marked choices and the elliptic-inversion map for simultaneous
sign reversal.

\emph{The four-period marking.}
Put $p_i=ag_i$, $q_i=bg_i$, choose
$(A_i-I)x_i=q_i$ over $\mathbb Q$, and write
\[
 \beta_i=-p_i^tJA_i,\qquad
 T_i=\begin{pmatrix}A_i&0&q_i\\
                    \beta_i&1&\beta_i x_i+w_i\\
                    0&0&1\end{pmatrix}.
\]
Here $w_i=1/n$ at the chosen component zero and zero elsewhere.
The Poincar\'e line marking in Theorem~\ref{maps:src:D.2.12} gives this
formula: the cocycles are the negative elliptic-logarithm
increments, the forward period covers $-Q$, and the line
winding has sign $+w_i$.
At an $I_n$ take $x_i=(b\kappa/n)r$, giving central entry
$w_i-ab\kappa^2/n$. A change of solution is fixed by $A_i$
and is annihilated by $\beta_i$, so it leaves that entry
unchanged. The component conditions give integral matrices
of determinant one, and
\[
 \prod_iT_i=I,\qquad T_i\delta=\delta,\qquad \chi T_i=\chi.
\]
Indeed the weight-zero product for $a=b=1$ is
$I-d^{-1}E_{\delta\gamma}$ in each displayed profile;
bilinearity gives $-ab/d$ for general $a,b$, cancelled
by $w=1/n$. The preserved alternating form is
$a\,e_1^*\wedge e_2^*-b\,\delta^*\wedge\gamma^*$.
Replacing a period matrix $\Pi$ by $\Pi H$ for
$H\in SL_4(\mathbb Z)$ uses the identity on the complex
universal cover; $H^{-1}$ is the inverse on period coordinates.

\emph{The primitive cusp discs.}
At the zero put $N=T-I$, $b_0=(r,0,0)$ and
\[
 \lambda_1=Nb_0,\qquad \lambda_2=N\gamma,\qquad
 H=(\lambda_1,b_0,\lambda_2,\gamma).
\]
Relative to $(u,\delta)$ the two short columns form
\[
 C=\begin{pmatrix}
 n&b\kappa\\-a\kappa&(1-ab\kappa^2)/n
 \end{pmatrix},\qquad
 C^{-1}=\begin{pmatrix}
 (1-ab\kappa^2)/n&-b\kappa\\a\kappa&n
 \end{pmatrix}.
\]
Both are integral and $\det C=1$. Hence $H\in SL_4(\mathbb Z)$,
$H^{-1}TH=I+E_{12}+E_{34}$, and the short lattice is exactly
$\mathbb Zu+\mathbb Z\delta$.

Let the short complex periods be the columns of $B(q)$.
Their determinant is the nonzero elliptic $u$-period times
the nonzero vertical line period, multiplied by $\det C$.
Exponentiating the coordinates of $B(q)^{-1}z$ gives the
short-period quotient $(\mathbb C^*)^2$; choosing logarithms
and multiplying by $B(q)$ gives its inverse modulo that lattice.
The long periods act by
\[
 (X,Y,q)\longmapsto
 (c_1(q)q^{l_1}X,c_2(q)q^{l_2}Y,q),
 \qquad (l_1,l_2)\in\mathbb Z^2 ,
\]
with the nonzero unit factors retained.
The triangular fan has rays $(v,1)$, $v\in\mathbb Z^2$.
One chart and its inverse are
\[
 (x,y,w)=(q/(XY),Y,X),\qquad
 (X,Y,q)=(w,y,xyw).
\]
These are the primitive cusp charts of Theorem~\ref{maps:src:D.2.6}.
The angular circles bound their normal toric discs.
The rays $(0,0,1),(1,0,1),(0,1,1)$ span the torus lattice.
Their discs kill the fundamental group of the dense torus;
moving loops off the boundary proves that the toric cover
is simply connected. Its free period quotient has deck
group $\mathbb Z^2$. The cap fibre map and inverse quotient
map are consequently
\[
 H(x_1,x_2,x_3,x_4)\longmapsto(x_2,x_4),\qquad
 (v,w)\longmapsto[H(0,v,0,w)] .
\]
Its kernel is the short lattice and its second coordinate is
$\chi$. The original and smooth Mumford models have these
same normal discs. Under $q=s^m$ the angular lattice stays
fixed and the period translation is multiplied by $m$.
A free affine cyclic quotient maps each disc to a disc.
Its restriction from one chosen regular fibre is an isomorphism
onto the quotient fibre, with inverse the representative over
the chosen root. Both relations therefore persist for every
allowed filling.

\emph{Integral normal closure.}
The fibre subgroup is already abelian before filling.
Conjugation by the $i$th peripheral lift acts by $T_i$,
so the normal subgroup generated by the vanishing vectors
is their integral span under all $T_i^{\pm1}$.
The displayed inverse matrix gives the exact relations
\[
 \widetilde u=\frac{1-ab\kappa^2}{n}\lambda_1
                  +a\kappa\lambda_2=(u,0,0),\qquad
 \delta=-b\kappa\lambda_1+n\lambda_2 .
\]
In the following table $(t_i)$ means $v=\sum_i t_iA_i u$,
and $h$ is the number of accepted section pairs at that
profile and zero position. Indices refer to the ordered
tuples above.
\begingroup\small\setlength{\arraycolsep}{3pt}
\[
\begin{array}{clr|cccrr}
\mathrm{OS}&\mathrm{profile}&j&u&v&(t_i)&\det(u,v)&h\\\hline
43&\mathrm{base}&2&(3,1)&(-4,-1)&(0,0,-1,0)&1&2\\
43&\mathrm{base}&3&(1,0)&(2,1)&(1,0,0,0)&1&2\\
43&\mathrm{base}&4&(1,0)&(2,1)&(1,0,0,0)&1&2\\
43&\mathrm{II}&3&(1,0)&(2,1)&(1,0,0)&1&2\\
45&\mathrm{base}&2&(1,0)&(-3,1)&(0,0,-1,0,0)&1&2\\
45&\mathrm{base}&3&(2,-1)&(-1,1)&(0,-1,0,0,0)&1&2\\
45&\mathrm{base}&4&(1,-2)&(3,-5)&(0,2,-1,0,0)&1&2\\
45&\mathrm{base}&5&(1,1)&(-2,-1)&(0,-1,0,0,0)&1&2\\
45&\mathrm{II}&3&(1,-2)&(-6,13)&(-2,9,0,0)&1&2\\
45&\mathrm{II}&4&(1,1)&(-4,-3)&(-1,1,0,0)&1&2\\
46&\mathrm{base}&2&(1,0)&(-2,1)&(0,0,-1,0)&1&2\\
46&\mathrm{base}&3&(1,-1)&(0,1)&(0,-1,0,0)&1&2\\
46&\mathrm{base}&4&(1,1)&(-2,-1)&(0,-1,0,0)&1&2\\
46&\mathrm{II}&3&(1,1)&(-4,-1)&(1,0,0)&3&2\\
47&\mathrm{base}&2&(1,-1)&(-12,13)&(1,0,-2,0,0)&1&2\\
47&\mathrm{base}&3&(1,-3)&(4,-11)&(1,-8,0,0,0)&1&4\\
47&\mathrm{base}&4&(1,1)&(3,4)&(1,-1,0,0,0)&1&4\\
47&\mathrm{base}&5&(1,0)&(8,1)&(0,4,-1,0,0)&1&4\\
47&\mathrm{II}&2&(1,-1)&(-12,13)&(1,0,-2,0)&1&2\\
47&\mathrm{II}&3&(1,-3)&(4,-11)&(1,-8,0,0)&1&4\\
47&\mathrm{III}&2&(1,-3)&(-8,25)&(1,0,-4,0)&1&4\\
47&\mathrm{III}&3&(1,1)&(-41,-40)&(-7,3,0,0)&1&4\\
49&\mathrm{base}&1&(1,0)&(1,1)&(0,1,0,0)&1&2\\
49&\mathrm{base}&3&(3,1)&(5,2)&(1,1,0,0)&1&4\\
49&\mathrm{base}&4&(1,0)&(1,1)&(0,1,0,0)&1&4\\
49&\mathrm{II}&3&(1,0)&(-1,1)&(-1,0,0)&1&2\\
\end{array}
\]
\endgroup

\begingroup\small\setlength{\arraycolsep}{3pt}
\[
\begin{array}{clr|cccrr}
\mathrm{OS}&\mathrm{profile}&j&u&v&(t_i)&\det(u,v)&h\\\hline
49&\mathrm{III}&3&(1,0)&(1,1)&(1,0,0)&1&4\\
50&\mathrm{base}&2&(1,0)&(-2,1)&(0,0,-1,0)&1&2\\
50&\mathrm{base}&3&(1,-1)&(0,1)&(1,0,0,0)&1&4\\
50&\mathrm{base}&4&(1,1)&(-2,-1)&(1,0,0,0)&1&4\\
50&\mathrm{II}&2&(2,-1)&(1,0)&(-1,0,0)&1&2\\
50&\mathrm{IV}&3&(1,1)&(-2,-1)&(1,0,0)&1&4\\
55&\mathrm{base}&2&(1,0)&(-3,1)&(0,0,-1,0,0)&1&2\\
55&\mathrm{base}&3&(2,-1)&(1,0)&(-1,1,0,0,0)&1&4\\
55&\mathrm{base}&4&(1,-2)&(1,-1)&(2,-2,-1,0,0)&1&4\\
55&\mathrm{base}&5&(1,1)&(-1,0)&(-1,1,0,0,0)&1&4\\
55&\mathrm{II}&2&(1,-1)&(-2,3)&(0,0,-1,0)&1&2\\
55&\mathrm{II}&4&(1,1)&(9,10)&(3,-1,0,0)&1&4\\
56&\mathrm{base}&1&(1,0)&(8,1)&(0,0,4,-1,0)&1&2\\
56&\mathrm{base}&2&(1,0)&(8,1)&(0,0,4,-1,0)&1&4\\
56&\mathrm{base}&3&(1,-1)&(0,1)&(1,-2,0,0,0)&1&4\\
56&\mathrm{base}&4&(1,-3)&(-1,4)&(1,-2,1,0,0)&1&8\\
56&\mathrm{base}&5&(1,1)&(-2,-1)&(1,-2,0,0,0)&1&8\\
56&\mathrm{II}&1&(2,-1)&(-9,5)&(0,-3,1,0)&1&4\\
56&\mathrm{II}&3&(4,-3)&(3,-2)&(0,-1,0,0)&1&2\\
56&\mathrm{II}&4&(1,-1)&(7,-6)&(1,-1,0,0)&1&4\\
56&\mathrm{III}&1&(1,0)&(24,1)&(0,12,-5,0)&1&2\\
56&\mathrm{III}&2&(1,-2)&(-2,5)&(0,0,-1,0)&1&4\\
56&\mathrm{III}&4&(1,1)&(-46,-45)&(-11,2,0,0)&1&8\\
56&\mathrm{IV}&1&(1,0)&(2,1)&(0,1,-1,0)&1&2\\
56&\mathrm{IV}&2&(1,-1)&(0,1)&(0,0,-1,0)&1&4\\
56&\mathrm{IV}&4&(1,1)&(8,9)&(3,-2,0,0)&1&8\\
\end{array}
\]
\endgroup
Except for the indicated $46/\mathrm{II}$ row, $\det(u,v)=1$.
Put
\[
 \widetilde v=\sum_i t_iT_i\widetilde u
               -\left(\sum_i t_i\beta_i u\right)\delta=(v,0,0).
\]
Both terms are vanishing relations. If $u=(u_1,u_2)$ and
$v=(v_1,v_2)$, the inverse matrix gives
\[
 e_1=v_2\widetilde u-u_2\widetilde v,\qquad
 e_2=-v_1\widetilde u+u_1\widetilde v .
\]
These identities cover $170$ rows. For the exceptional profile,
$u=(1,1)$, $v=(-4,-1)$, and
$-\widetilde u-\widetilde v=3e_1$.
The resulting lattice in $\ker\chi$ is
$\{(x,y,z,0):x-y\equiv0\pmod3\}$.

\emph{The two odd-star relations.}
For $(a,b)=(1,4),(4,1)$ the matrices at $I_3^*$ are
\[
 T_*^{(1)}=\begin{pmatrix}-1&-3&0&4\\0&-1&0&4\\-1&-2&1&3\\0&0&0&1\end{pmatrix},
 \qquad
 T_*^{(4)}=\begin{pmatrix}-1&-3&0&1\\0&-1&0&1\\-4&-8&1&3\\0&0&0&1\end{pmatrix}.
\]
Their square differences are $w_1(0,3,0,-6)$ and
$w_4(0,6,0,-3)$, where
$w_1=(2,0,1,0)^t$, $w_4=(1,0,2,0)^t$.
The integral period bases and inverses
\[
 H_1=(w_1,e_2,-e_1,2e_2+\gamma),\qquad
 H_4=(w_4,-\gamma,\delta,e_2+2\gamma),
\]
\[
 H_1^{-1}=
 \begin{pmatrix}0&0&1&0\\0&1&0&-2\\-1&0&2&0\\0&0&0&1\end{pmatrix},
 \qquad
 H_4^{-1}=
 \begin{pmatrix}1&0&0&0\\0&2&0&-1\\-2&0&1&0\\0&1&0&0\end{pmatrix}
\]
have determinant one and satisfy
$H_j^{-1}(T_*^{(j)})^2H_j=I+3E_{12}$.
The first period is the primitive angle of the rank-one
semistable cover. Its toric rays are $(k,1)$, with adjacent
charts
\[
 X=z/s^k,\quad Y=s^{k+1}/z,\quad XY=s;
 \qquad z=Xs^k,\quad s=XY .
\]
The second period translates $k$ by three. The central
normalization consists of three degree-zero ruled surfaces
over an elliptic torus in a cycle.
For a cover of order $2\ell$, the additional base change
gives $H_j^{-1}(T_*^{(j)})^{2\ell}H_j=I+3\ell E_{12}$.
The same adjacent charts repeat in a cycle of $3\ell$
components, with the primitive first period unchanged.
For flat multipliers $\mu_1,\mu_2$, multiplication by
$\exp(-t_1\log\mu_1-t_2\log\mu_2)$ trivializes the smooth
line bundle; the opposite exponent is its inverse.
If the last zero/infinity gluing has translation $a$, choose
a real lift of $a$ and $f:\mathbb P^1\to[0,1]$ with
$f(0)=0,f(\infty)=1$. The map
$(z,x)\mapsto(z,x-f(z)a)$, with inverse given by addition,
removes that translation. Constant translations handle the
remaining components. This identifies the central surface
with its cycle of $\mathbb P^1$'s times $T^2$.
Its plumbing retracts to that surface. The fibre map and
inverse quotient map are
\[
 H_j(x_1,x_2,x_3,x_4)\longmapsto(x_2,x_3,x_4),\qquad
 (v_1,v_2,v_3)\longmapsto[H_j(0,v_1,v_2,v_3)] .
\]
The kernel is $\mathbb Zw_j$. Its normal disc bounds $w_j$,
and its projection bounds the same circle in every allowed
free logarithmic quotient.

For the original cap use the punctured line bundle over the
elliptic $I_3^*$ plumbing before the contracting quotient.
The rational core components have the affine $D_7$ tree as
dual graph. Their disc-bundle neighbourhoods and plumbing
intersections show that the base is simply connected and
that its component spheres generate $H_2$.
For $a=1$, the solution of $(A_*-I)h=g_*$ is $(1/4,-1/2)$,
so $P$ is nonidentity. The line has degree one on the
component met by $P$. The circle-bundle sequence sends that
sphere to a generator of the $\delta$ circle, and the
punctured line bundle is simply connected. On this core the
Hopf model contracts the fibre circle by
$(r e^{it},\sqrt{1-r^2})$ in $S^3$; the dual bundle gives
the opposite degree convention. Projection and path lifting
then contract the elliptic loops as well. In particular $w_1$
is killed.

For $a=4$, $h=(1,-2)$ is integral and the line has degree
zero on every component. A nonzero frame identifies it with
the trivial line over the plumbing; multiplication by the
frame and division by it are inverse maps.
The boundary winding has $\delta$ coefficient one and is
invariant under $T_*^{(4)}$. The equation
\[
 b(A_*-I)+(-4,-8)=0
\]
has the unique solution $b=(-2,-1)=-h^tJ$.
Thus the cap-cover map is $(x,y,z)\mapsto z-2x-y$, with
inverse $k\mapsto k\delta$; it kills $w_4=e_1+2\delta$.
The original and logarithmic arguments both supply the
required relation. For $a=1$ subtract
$w_1-\delta=2e_1$ from $3e_1$ to obtain $e_1$;
for $a=4$, $w_4-2\delta=e_1$ directly.
Then $e_2=\widetilde u-e_1$.

Every relation used has $\chi=0$, and the displayed
identities give the reverse inclusion. The fibre quotient
and its inverse are therefore
\[
 \Lambda/\mathcal N\xrightarrow{\ \sim\ }\mathbb Z\gamma,
 \qquad [v]\longmapsto\chi(v)\gamma,\qquad
 k\gamma\longmapsto[k\gamma],
 \qquad \mathcal N=\ker\chi .
\]

\emph{The local cyclic extension.}
The smooth cover has fibre group $\Lambda$. In a rank-one
Mumford cover factor $T-I=w\lambda$, with $w$ primitive
and $\lambda$ an integral row; put $d_0=\gcd(\lambda)$.
Then $\lambda(w)=0$. Choose $b$ with $\lambda(b)=d_0$
and extend $w$ to a basis of $\ker\lambda$ by integral
Bezout operations. Together with $b$ these give a basis
$K=(w,b,\xi_1,\xi_2)$; reverse the sign of $\xi_1$ when necessary
to make $\det K=1$. Its inverse is $\operatorname{adj}(K)$,
and $K^{-1}TK=I+d_0E_{12}$.
The adjacent toric charts and gluing maps above, with
$d_0$ in place of three, give the cap group $\Lambda/\mathbb Zw$,
the projection $K(x_1,x_2,x_3,x_4)\mapsto(x_2,x_3,x_4)$,
and inverse quotient map $(v_1,v_2,v_3)\mapsto[K(0,v_1,v_2,v_3)]$.
The rank-two cover has group
$\Lambda/(\mathbb Z\lambda_1+\mathbb Z\lambda_2)$,
with its displayed two-coordinate projection and section.

For every cap cover the boundary fibre generates its
fundamental group. In the toric models this follows from
the quotient maps just given. For an original zero-free cap,
the unrolled punctured line bundle extends over the elliptic
central fibre and gives the integral return homomorphism
restricting to $\chi$.
Move a loop off the central divisor by transversality.
Over the punctured disc the group is generated by the regular
fibre and a meridian. A transverse disc at a smooth point
of a multiplicity-one central component, lifted through a
local nonzero line frame, contracts a meridian differing
from the chosen one by a fibre loop. Hence the fibre
generates the cap group. Every local vanishing relation
has zero return. The toric return maps have the same
property and take $\gamma$ to one.
After imposing $\mathcal N$, each cover cap group is
exactly $\mathbb Z\gamma$, with forward map its return
and inverse the loop $\gamma$.

Let the free cyclic quotient have order $m_i$. Write its
base map as $q=s^{m_i}$ and choose a root over a boundary
basepoint. Restriction from that root fibre to the quotient
fibre is an isomorphism; its inverse selects the unique
representative over the chosen root.
Transport along the lifted positive peripheral arc and
then apply the deck identification. By the marked overlap
defining this logarithmic filling, its conjugation on the
boundary lattice is precisely the original $T_i$.
The affine translation affects the power loop and leaves
this linear conjugation unchanged. Since the boundary
fibre surjects onto the whole cover cap group,
$\chi T_i=\chi$ proves preservation of its entire return
homomorphism.

Choose the corresponding deck path $c_i$. Its $m_i$
concatenated deck translates give the actual integral power
vector $v_i$. The covering sequence is an extension by
$\mathbb Z/m_i$.
The image of $\ker\chi$ from the boundary is exactly the
kernel of return in the cover cap group: boundary
surjectivity and the section $\gamma$ prove both inclusions.
It is deck stable. Quotienting by it preserves the injection
of the remaining $\mathbb Z\gamma$ and gives
\[
 \langle\gamma,c_i\mid[\gamma,c_i]=1,\
                 c_i^{m_i}=\gamma^{n_i}\rangle,
 \qquad n_i=\chi(v_i).
\]
The forward map sends the symbols to the geometric return
and deck paths. Its inverse takes the covering-sheet index
$0\le j<m_i$ and remaining return exponent $r$, giving
$\gamma^r c_i^j$. The sheet index and injection of the
cover group prove uniqueness. This is the whole reduced
local group.
An unchanged cap has $m_i=1,n_i=0$ in its contracting
meridian marking. Smooth regular logarithmic sites use
the torus-times-disc cover, and the odd-star logarithmic
sites use the stipulated rank-one semistable cover.

\emph{The global presentation and its inverse.}
Cut the punctured sphere along arcs to its ordered boundary
circles. Its torus family has fibre group $\Lambda$ and
peripheral conjugations $T_i$.
Before quotienting by the contracting translation, the
punctured line bundle is an infinite cyclic cover with
connected fibre. Choose each peripheral lift closed in
that cover. Its return is zero. Their product is a fibre
loop of $\chi$ value zero and is killed by $\mathcal N$.
After removing the unchanged meridians, the product is
$c_1\cdots c_s=1$. A line-frame transition winds only in
$\delta$, so it has the same zero return. This proves
$e=0$ for the canonical lifts.

For general marked lifts write the product as $\gamma^e$.
Van Kampen with the local inverse maps gives
\eqref{maps:eq:D.FG1}. Explicitly, send its generators
to the return and deck paths. In the opposite direction
send a fibre vector $v$ to $\gamma^{\chi(v)}$, an unchanged
meridian to one, and each logarithmic deck path to $c_i$.
Monodromy, vanishing, local power and global product
relations are exactly those already verified. The two
compositions fix every generator.
Replacing $c_i$ by $c_i'=c_i\gamma^{t_i}$ gives
$n_i'=n_i+m_it_i$ and $e'=e+\sum_i t_i$.
The isomorphism sends old $c_i$ to $c_i'\gamma^{-t_i}$;
its inverse sends $c_i'$ to $c_i\gamma^{t_i}$.
Abelianizing these same generator maps gives the row
cokernel \eqref{maps:eq:D.FG2}, with inverse represented
by the corresponding products of geometric loops.

\emph{The zero-, one- and two-site groups.}
For $m_i=1$ substitute $c_i=\gamma^{n_i}$ and change
$e$ to $e-n_i$; reinsertion is the inverse Tietze move.
With no sites the group is $\mathbb Z/e$, including
$\mathbb Z$ for $e=0$. With one site $c_1=\gamma^e$
gives $\mathbb Z/(m_1e-n_1)$ with the same zero convention.
With two sites substitute $c_2=\gamma^e c_1^{-1}$.
The remaining group is abelian with relation matrix $R$
in the statement and
$D=m_2n_1+m_1n_2-em_1m_2$.

These reductions are integral. For the original
three-column relation matrix in the two-site case,
the row and column maps
\[
 U=\begin{pmatrix}1&0&0\\0&1&-m_2\\0&0&1\end{pmatrix},
 \qquad
 V=\begin{pmatrix}1&0&0\\0&1&0\\e&-1&1\end{pmatrix}
\]
give $\operatorname{diag}(R,1)$.
Their inverses reverse the off-diagonal entries.
The gcd of the entries of $R$ is
$d=\gcd(m_1,m_2,n_1,n_2)$.
Start with a nonzero column, exchanging the two columns when
necessary; the exchange is its own inverse. Such a column exists
because $(m_1,-m_2)^t$ is nonzero. Zero columns are skipped.
For a nonzero column $(x,y)^t$, choose
$\alpha x+\beta y=g=\gcd(x,y)>0$. The row map and inverse
\[
 \begin{pmatrix}\alpha&\beta\\-y/g&x/g\end{pmatrix},
 \qquad
 \begin{pmatrix}x/g&-\beta\\y/g&\alpha\end{pmatrix}
\]
send that column to $(g,0)^t$. Transposing gives the
column operation for a row. Alternating these operations
clears the first row and column; when the pivot does not
divide the remaining entry, adding the second column to
the first gives a strictly smaller positive gcd pivot.
This terminates at $d$, and the determinant gives the
other pivot $|D|/d$.
For $D=0$ the rank is one because $m_1,m_2>1$,
leaving one free summand. If $U_0RV_0$ is the resulting
diagonal matrix, the row-cokernel map is
$[z]\mapsto[zV_0]$, with inverse
$[z]\mapsto[zV_0^{-1}]$. This proves
\eqref{maps:eq:D.FG3} and the invariant factors.

\emph{Three or more sites.}
Set $\gamma=1$ and all but three peripheral generators
equal to one. The resulting quotient has triangle orders
$a,b,c\ge2$. On $\mathbb R^3$ take the symmetric form
with diagonal one and
\[
 B_{12}=-\cos(\pi/c),\qquad
 B_{13}=-\cos(\pi/b),\qquad
 B_{23}=-\cos(\pi/a).
\]
Define $R_i(v)=v-2B(v,e_i)e_i$. These are involutions.
On each two-coordinate plane the form is positive
definite, of determinant $\sin^2(\pi/k)$, and the
reflection product has eigenvalues $\exp(\pm2\pi i/k)$.
It fixes the orthogonal complementary line. That direct
decomposition exists even when the full form is degenerate,
because its restriction to the plane is nondegenerate.
Thus $x=R_2R_3$, $y=R_3R_1$, $z=R_1R_2$
have exact orders $a,b,c$ and $xyz=I$.
They give a nontrivial representation of the quotient,
hence of \eqref{maps:eq:D.FG1}.

In the one-site case every common divisor of $m_1,n_1$
divides $m_1e-n_1$. In the two-site case a common divisor
of $m_1,m_2$, or of either pair $(m_i,n_i)$, divides $D$.
Triviality makes the relevant integer a unit, proving
all coprimality conclusions.
\end{proof}

\subsection{Logarithmic families with smooth type \texorpdfstring{$S^2\times S^4$}{S2 x S4}}

\begin{definition}[The two-component OS56 family]
\hypertarget{tag-Q119}{}\label{mapsIII:def:two-component-os56}
{\normalfont\scriptsize tag Q119}\par
Take the OS56 type-$II$ surface with ordered elliptic fibres
\[
                         I_2,\ II,\ I_5,\ I_3.
\]
Let $G$ be its marked Mordell--Weil generator, take $P=30G$, $Q=G$,
and put the degree-one linearization zero at $I_2$. Use the
two-component toroidal filling there, the primitive one-node
fillings at $I_5,I_3$, and the logarithmic filling at $II$.
In the period convention of Definition~\ref{mapsIII:def:elliptic-marked-data}, take
\[
 m=6k,\qquad v=n\gamma+r\delta,\qquad
 k\in\mathbb Z_{>0},\quad n,r\in\mathbb Z,\quad \gcd(n,6k)=1.
\]
\end{definition}

\begin{theorem}[Logarithmic families with smooth type $S^2\times S^4$]
\hypertarget{tag-02KP}{}\label{maps:src:D.2.15}
{\normalfont\scriptsize tag 02KP}\par
The compact smooth threefold of Definition~\ref{mapsIII:def:two-component-os56} has
\[
\begin{aligned}
 \pi_1(X)&=H_1(X)=\mathbb Z/|n|,\qquad H_5(X)=0,\\
 H_2(X)&=H_4(X)=\mathbb Z\oplus\mathbb Z/|n|,
 &H_3(X)&=\mathbb Z/|n|,\\
 H_0(X)&=H_6(X)=\mathbb Z, &\chi(X)&=4.
\end{aligned}
\]
Every simply connected member, equivalently $n=\pm1$, is oriented
diffeomorphic to $S^2\times S^4$. This configuration supplies that
smooth type for every logarithmic multiplicity $6k$.
\end{theorem}

\begin{proof}
\emph{The surface and its line.}
The realization in Theorem~\ref{mapsD:repair:collision-realization} gives the
marked surface with fibres $I_2,II,I_5,I_3$ at
$-5/9,3,\infty,0$. Its section has height $1/30$ and generates the
full Mordell--Weil group. In the primitive component bases its
residues are $-1\pmod2$, $1\pmod5$, and $1\pmod3$.
Thus $P=30G$ is narrow and $\langle P,Q\rangle=1$.
Apply Proposition~\ref{maps:src:D.2.19} to this pair. Its divisor and base-function
maps identify the linearization line with the pullback of
$\mathcal O_{\mathbb P^1}(1)$ and give a section with a simple zero
on the whole $I_2$ fibre. Their inverse maps retain the marked
Poincar\'e line and its rigidifications.

In the integral basis $(e_1,e_2,\delta,\gamma)$, and cyclic order
$II,I_5,I_3,I_2$, the full monodromies are
\[
\begin{aligned}
 T_0&=\begin{pmatrix}27&37&0&0\\-19&-26&0&0\\0&0&1&0\\0&0&0&1\end{pmatrix},
 &T_1&=\begin{pmatrix}61&80&0&4\\-45&-59&0&-3\\-90&-120&1&-6\\0&0&0&1\end{pmatrix},\\
 T_2&=\begin{pmatrix}4&3&0&1\\-3&-2&0&-1\\-30&-30&1&-10\\0&0&0&1\end{pmatrix},
 &T_3&=\begin{pmatrix}5&8&0&-2\\-2&-3&0&1\\30&60&1&-14\\0&0&0&1\end{pmatrix}.
\end{aligned}
\]
Their product is $I$, $T_0^6=I$, and they fix $\delta$ and
$\chi=\gamma^*$. The invariant integral alternating form is
\[
                 \omega=30e_1^*\wedge e_2^*-\delta^*\wedge\gamma^*.
\]
The local bases below specify the attachment maps in this marking.
Put
\[
\begin{array}{lll}
 w_c=(2,-1,15,0),&r_c=(1,0,0,0),\\
 w_1=(4,-3,-6,0),&a_1=(-1,1,2,-5),\\
 w_2=(1,-1,-10,0),&a_2=(1,0,0,-3).
\end{array}
\]
The matrices with columns $(w_i,\gamma,a_i,\delta)$, for $i=1,2$,
are unimodular and conjugate $T_i$ to $I+E_{12}$.
Their inverses are their integral adjugates divided by their
determinants. The primitive one-node construction of
Theorem~\ref{maps:src:D.2.12}, with these changes of basis and their inverses,
therefore gives the specified $I_5$ and $I_3$ fillings.

\emph{The two-component cusp.}
In the basis $(w_c,r_c,\delta,\gamma)$ the difference $T_3-I$ is
\[
 N_c=\begin{pmatrix}0&2&0&-1\\0&0&0&0\\0&0&0&1\\0&0&0&0\end{pmatrix}.
\]
The inverse basis matrix is
\[
 \begin{pmatrix}0&-1&0&0\\1&2&0&0\\0&15&1&0\\0&0&0&1\end{pmatrix}.
\]
The valuation periods are initially $(2,0),(-1,1)$.
Right multiplication by
$\left(\begin{smallmatrix}1&1\\0&1\end{smallmatrix}\right)$,
whose inverse has upper-right entry $-1$, gives the basis
\begin{equation}
                       (2,0),\qquad(1,1).
 \hypertarget{tag-01ZU}{}\label{maps:eq:D.I1}
\end{equation}
The second valuation coordinate is the simple whole-fibre zero.
All holomorphic unit multipliers remain invertible characters on the
toric overlaps. The period lattice $\Gamma$ has index two, with
quotient $(a,b)\mapsto a+b\pmod2$ and representatives $(0,0),(1,0)$.
The unit-triangle fan consequently has two normalized $dP_6$
components, six double curves, and four triple points in its quotient.
Its unimodular charts and the contracting period action give a smooth
total space by the toroidal construction of Theorem~\ref{maps:src:D.2.12}.

The integral involution
\[
 F=\begin{pmatrix}1&1\\0&-1\end{pmatrix},\qquad
 F^2=I,\qquad
 F\begin{pmatrix}2&1\\0&1\end{pmatrix}F
       =\begin{pmatrix}2&0\\0&1\end{pmatrix}
\]
permutes the six fan rays and every adjacent cone. Its torus map is
$(Z,V)\mapsto(ZV,V^{-1})$; the inverse sends
$(Z',V')$ to $(Z'V',V'^{-1})$ and fixes the base coordinate.
It extends on each affine toric chart through the corresponding
character map. Thus the normalization resolution in
Theorem~\ref{maps:src:D.1.2}, at $d=2$, applies to a cusp neighbourhood $U$:
\begin{equation}
 \pi_1(U)=\mathbb Z^2,\qquad
 (H_0,H_1,H_2,H_3,H_4)(U)
   =(\mathbb Z,\mathbb Z^2,\mathbb Z^5,\mathbb Z^2,\mathbb Z^2).
 \hypertarget{tag-01ZV}{}\label{maps:eq:D.I2}
\end{equation}
For $M=\partial U$, the Smith entries of $N_c$ are $(1,2)$,
and those of $\bigwedge^2T_3-I$ are $(1,1)$.
The Wang sequences give
$H_1(M)=\mathbb Z^3\oplus\mathbb Z/2$ and $H_2(M)=\mathbb Z^6$.
The cusp kills $w_c,\delta$ and its base meridian $x_3$, retaining
the two long periods. Its $H_1$ kernel is
$\mathbb Z\langle x_3\rangle\oplus\mathbb Z/2\langle\delta\rangle$.

Let $D_0,D_1$ denote the two cusp components and also use their
indices for the corresponding positive normal discs. The disc
boundaries are $x_3,x_3+w_c$. In the boundary torsion quotient,
$w_c=\delta$ and $2\delta=0$. Under
$H_2(U,M)=H^4(U)=\mathbb Z^2$ the boundary homomorphism is
\[
 (a,b)\longmapsto(a+b,b\bmod2)
       \quad\hbox{in }\mathbb Z\oplus\mathbb Z/2.
\]
Its kernel is $\mathbb Z\langle(-2,2)\rangle$, with inverse
$(a,b)\mapsto b/2$ on that kernel. Therefore
\begin{equation}
           H_2(U)/\operatorname{im}H_2(M)\cong\mathbb Z.
 \hypertarget{tag-01ZW}{}\label{maps:eq:D.I3}
\end{equation}
Sweeping the first normal disc by the $w_c$ and $\delta$ actions
gives the two relative three-cycles. Their boundary tori project to
the full primitive fixed lattice of $T_3$ in the Wang sequence.
That projection has the inverse which takes a fixed loop to its
parametrized sweep. The normalization groups in
\eqref{maps:eq:D.I2} and \eqref{maps:eq:D.I3} give kernel rank two,
so these tori are the full kernel of $H_2(M)\to H_2(U)$.
The remaining boundary image is the regular-fibre specialization image.

\emph{The integral chains of the regular part.}
Let $V$ be the torus bundle over the three-loop retract of the
four-punctured base and put
\[
                    D=[T_0-I\ T_1-I\ T_2-I].
\]
The $12\times3$ integer matrix $C$ with nonzero rows
\[
 C_1=(-27,-37,-4),\quad C_2=(19,26,3),\quad
 C_8=(0,0,-2),\quad C_{12}=(0,0,1)
\]
has $DC=\operatorname{diag}(1,1,2)$ in the first three target rows. For $i<3$ let $j_i(v)$ place $v$
in the $i$th four-coordinate edge block, and put
\[
                  j_3(v)=-(T_1T_2v,T_2v,v).
\]
The following nine columns form the matrix $K$:
\[
\begin{split}
 L_1,L_2,L_3&=j_1(w_1),j_1(a_1),j_1(\delta),\\
 L_4,L_5,L_6&=j_2(w_2),j_2(a_2),j_2(\delta),\\
 L_7,L_8,t&=j_3(w_c),j_3(\delta),j_0(\gamma).
\end{split}
\]
Direct multiplication gives $DK=0$ and $\det[C\ K]=-1$.
Thus $[C\ K]^{-1}=-\operatorname{adj}[C\ K]$ is an integral
inverse, and its last nine rows give the inverse coordinates on
$\ker D$. In those coordinates
\[
                         j_0(\delta)=-L_3-L_6-L_8.
\]
In particular the regular fibre coinvariants in degree one are
$\mathbb Z\langle\gamma\rangle\oplus\mathbb Z/2\langle\delta\rangle$.
The cusp boundary class $-\delta$ maps to that same nonzero
order-two class.

Here is an integral quotient map for the degree-two fibre part.
Order the wedge coordinates as $12,13,14,23,24,34$ and set
\[
 W=[\bigwedge^2T_0-I\ \bigwedge^2T_1-I\ \bigwedge^2T_2-I],
 \qquad \kappa=e_{12}+30e_{34},\qquad s=-e_{34}.
\]
The $18\times5$ matrix $P$ with nonzero rows
\[
\begin{array}{ll}
 P_2=(270,-27,0,-37,0),&P_3=(30,0,-27,0,-37),\\
 P_4=(-190,19,0,26,0),&P_5=(-21,0,19,0,26),\\
 P_{15}=(-1,0,0,0,0)
\end{array}
\]
satisfies $WP=[\kappa,e_{13},e_{14},e_{23},e_{24}]$.
The matrix $[s,\kappa,e_{13},e_{14},e_{23},e_{24}]$ is unimodular,
with first inverse row $\omega$. Since $\omega W=0$, these identities
prove that $\operatorname{im}W=\ker\omega$ integrally. They give
the quotient $\omega$ and inverse $1\mapsto s$, as well as a relation
preimage for every element of its kernel.

The fixed-loop cylinders representing the columns of $K$ are the
geometric sweeps of Theorem~\ref{maps:src:D.2.12}. The invariant closed
form there has fibre restriction $\omega$, value one on $s$, and
value zero on every such cylinder. Its collar prism homotopies
preserve these values. The fibre and graph sequences therefore split
in the displayed geometric basis:
\[
                    H_2(V)=\mathbb Z\langle s,L_1,\ldots,L_8,t\rangle.
\]

\emph{The reduced attachments and their extra cycle.}
The one-node normal-disc sweeps kill $L_1,\ldots,L_6$, and the two cusp
sweeps kill $L_7,L_8$. In base order $(x_0,x_1,x_2)$ the four
reduced normal-disc boundary columns have base parts
\[
 \begin{pmatrix}0&0&-1&-1\\1&0&-1&-1\\0&1&-1&-1\end{pmatrix},
\]
and the last column has the additional fibre term $w_c$.
The base kernel is generated by $(0,0,-1,1)$; its fibre image has
exact order two. The full kernel is consequently
$\mathbb Z\langle(0,0,-2,2)\rangle$.
The chain $C(4,-2,15)^t$ has boundary $2w_c$.
After joining the disc boundaries to their based meridians by their
annular tracks, subtract this chain from twice the difference of the
two cusp normal discs. The resulting closed cycle $z$ has cusp
normal-disc coordinates $(-2,+2)$ and gives an inverse section of
the free kernel. After the three reduced fillings,
\[
                    H_2=\mathbb Z\langle s,t,z\rangle,
                    \qquad H_1=\mathbb Z\langle\gamma\rangle.
\]

\emph{The finite cap and its holomorphic collar.}
At $II$ the elliptic difference matrix is unimodular and the fixed
periods are $\delta,\gamma$. In local base coordinate $q=t-3$,
the good-reduction coordinates are
\[
                q=\sigma^m,\qquad
                x=\sigma^{2k}X,\qquad y=\sigma^{3k}Y.
\]
For the Cartier coordinate of the Poincar\'e line, the identity
\[
 \frac{z_{\rm old}}{z_{\rm old}-p_{\rm old}}
     =\frac{z}{z-p_{\rm new}},\qquad
 z_{\rm old}=\sigma^{-k}z,\quad
 p_{\rm new}=\sigma^kp_{\rm old}(\sigma^m)
\]
shows that the two section characters cancel. The deck action is
the good-reduction linear action followed by translation by
$(n\gamma+r\delta)/m$. On its central torus the $\gamma^*$
coordinate changes by $n/m$; coprimality makes every nonidentity
deck power free there. Away from the central torus the root rotation
is free. This gives a smooth complex cap.

In the marked flat real collar coordinates, the maps from
Theorem~\ref{maps:src:D.2.12} and their inverses are
\[
\begin{split}
 h_u(y,\sigma)&=(y-\theta v/(2\pi)+uB(\sigma),\sigma),\\
 h_u^{-1}(y',\sigma)&=(y'+\theta v/(2\pi)-uB(\sigma),\sigma),
                \qquad 0\leq u\leq1,\\
 B(\sigma)&=\frac{\log|\sigma|}{2\pi}J_\sigma v-u_p(\sigma).
\end{split}
\]
Here $\theta=\arg\sigma$, $J_\sigma$ is the complex structure,
and $u_p$ is the real vector specified by the logarithm of the
nonzero line unit. The function $B$ is angularly periodic and
equivariant for the linear deck action. Hence both maps descend;
$h_1$ is the holomorphic attachment. A radial cutoff of this
translation extends it over the regular collar, with negative
translation as inverse. The resulting prism homotopy preserves
integral homology and the invariant area class used above.

Choose $m\alpha+n\beta=1$ and let $x_0$ be the finite base loop.
The central group has relation $mx_0=n\gamma+r\delta$ and its
abelianization has basis $\delta,z_0$, with inverse formulas
\[
 z_0=\alpha\gamma+\beta x_0,\qquad
 \gamma=mz_0-r\beta\delta,\qquad
 x_0=nz_0+r\alpha\delta.
\]
The normal circle is $-n\gamma-r\delta+mx_0$.
The invariant normal-line trivialization in
Theorem~\ref{maps:src:D.2.12} makes its relative three-cycles exactly the
two normal-disc sweeps along $\delta,z_0$.
Writing $c_L=-L_3-L_6-L_8$, their cylinder boundaries are
\[
\begin{aligned}
 \tau_\delta&=ns+mc_L,& \tau_\gamma&=-rs+mt,\\
 \tau_{x_0}&=nt+rc_L,&
 \tau_{z_0}&=-r\alpha s+t+r\beta c_L.
\end{aligned}
\]
After the reduced sweeps are killed, the two columns on $(s,t)$ are
\begin{equation}
                   B=\begin{pmatrix}n&-r\alpha\\0&1\end{pmatrix}.
 \hypertarget{tag-01ZX}{}\label{maps:eq:D.I4}
\end{equation}
Multiplication by
$U=\left(\begin{smallmatrix}1&r\alpha\\0&1\end{smallmatrix}\right)$
gives $UB=\operatorname{diag}(n,1)$; $U^{-1}$ negates its upper-right
entry. The quotient map is
$(a,b)\mapsto a+r\alpha b\pmod{|n|}$, with inverse $1\mapsto s$.
A kernel vector has relation preimage
$((a+r\alpha b)/n,b)$. Every finite boundary has zero $z$ coordinate,
since its chains avoid the reduced normal discs. This proves
\[
                    H_2(X)=\mathbb Z\langle z\rangle
                                  \oplus\mathbb Z/|n|\langle s\rangle.
\]

\emph{The fundamental group and the remaining homology.}
Van Kampen starts with the commuting fibre lattice and the three
based loops, with conjugations $T_0,T_1,T_2$.
The reduced caps kill their meridians and their geometric vanishing
circles. In $\ker\chi$ the columns $(w_c,w_2,\delta)$ form
\[
 A=\begin{pmatrix}2&1&0\\-1&-1&0\\15&-10&1\end{pmatrix},\qquad
 A^{-1}=\begin{pmatrix}1&1&0\\-1&-2&0\\-25&-35&1\end{pmatrix}.
\]
Thus the angular coordinate discs kill all of $\ker\chi$, leaving
$\mathbb Z\gamma$. Every monodromy is the identity on this quotient.
The product relation kills $x_0$ once the other three meridians
vanish, and the finite cap adds $n\gamma=0$.
Each cap group is generated by its boundary image under the
primitive-node, normalization, and finite mapping-torus maps already
specified. Consequently the mutually inverse maps are
\[
 \pi_1(X)\longrightarrow\mathbb Z/|n|:
       v\longmapsto\chi(v),\quad x_i\longmapsto0,
       \qquad 1\longmapsto\gamma.
\]

For completeness, the degree-one Mayer--Vietoris source has boundary
free ranks $3,4,4,3$, and its target has regular rank four and cap
ranks $2,3,3,2$. Both total free ranks are fourteen, and each group
has one order-two summand. The map on those summands sends
$-\delta$ to the same nonzero class. The total cokernel is the
finite group just computed. Thus the map is injective on the free
part and on the torsion part; its full kernel is zero. The
degree-two cokernel and the explicit lift $z$ above therefore give
the full asserted $H_2$.

The two $dP_6$ components, six double curves, and four triple points
give cusp Euler number $12-12+4=4$. Each remaining central fibre
has a torus factor or a free torus quotient and Euler number zero.
Hence $\chi(X)=4$. Rational duality gives $b_3=0$; the universal
coefficient theorem and oriented integral duality give
\[
 H_0=H_6=\mathbb Z,\quad H_1=H_3=\mathbb Z/|n|,\quad
 H_2=H_4=\mathbb Z\oplus\mathbb Z/|n|,\quad H_5=0.
\]

\emph{Canonical line and characteristic numbers.}
Take $n=\pm1$. Then $X$ is simply connected,
$H_2(X)=H_4(X)=\mathbb Z$, and $H_3(X)=0$.
Let $a\in H^2(X;\mathbb Z)$ satisfy $a(z)=1$.
The positive normal-disc coordinates of $z$ give
\[
 c_1\mathcal O(D_0)=-2a,\qquad c_1\mathcal O(D_1)=2a,
 \qquad c_1f^*\mathcal O_{\mathbb P^1}(1)=0.
\]
The last equality uses the whole reduced cusp fibre $D_0+D_1$.

The meromorphic three-form
\[
                       \Omega=dt\wedge\frac{dx}{y}\wedge\frac{dV}{V}
\]
descends through the line gauges and section translations: their
extra logarithmic differentials lie in the base and elliptic
directions and wedge to zero with the first two factors.
At infinity, $t=1/s$, $x=t^2X$, $y=t^3Y$ give
$\Omega=-s^{-1}ds\wedge(dX/Y)\wedge(dV/V)$.
In every unimodular reduced toric chart the expression
$dq\wedge d\log Z\wedge d\log V$ is a nonzero constant and unit
times the coordinate volume form; the primitive one-node maps have
the same property. At the finite cap,
\[
                 \Omega=m\sigma^{m-1-k}
                    d\sigma\wedge\frac{dX}{Y}\wedge\frac{dV}{V}.
\]
The root characters of $\sigma,X,Y$ are $1,-2k,-3k$.
Together with the coefficient they give the trivial total character
$\zeta^{m-1-k+1+k}=1$. The free quotient is \'etale, so the order
along its reduced central divisor $G_0$ is $m-1-k$.
The form and the base coordinate give the line-bundle maps
\[
 K_X=\mathcal O(-F_\infty+(m-1-k)G_0),\qquad
 \mathcal O(mG_0)=f^*\mathcal O_{\mathbb P^1}(1),\qquad
 K_X^{\otimes m}=f^*\mathcal O_{\mathbb P^1}(-k-1).
\]
The maps are multiplication by the displayed local forms and base
functions, with their reciprocals as inverses. Since $H^2(X)$ is
free, these identities give $c_1(K_X)=0$, and hence $w_2(X)=0$.

Apply the component intersection and tangent-line calculation of
Theorem~\ref{maps:src:D.2.47} to this index-two vertex lattice. The six
ray classes are $(1,1,0,1,1,0)$. For component coefficient differences
$x,y$ the two six-tuples are $(x,x,0,x,x,0)$ and
$(y,y,0,y,y,0)$, and their intersection is
\[
 -\sum_i d_i(x)d_i(y)
 +\sum_i\bigl(d_i(x)d_{i+1}(y)+d_{i+1}(x)d_i(y)\bigr)
                  =-4xy+4xy=0.
\]
Thus $D_0^3=0$. The same result, with its tangent and normal-line
maps, gives $\langle p_1(TX),D_0\rangle=0$.
Since $\operatorname{PD}(D_0)=-2a$, these equalities imply
$a^3=0$ and $\langle p_1(TX)a,[X]\rangle=0$.
The integral $H^4$--$H^2$ pairing is unimodular, so $p_1(TX)=0$.

\emph{Oriented smooth recognition.}
Let $f_X:X\to K(\mathbb Z,2)=\mathbb{CP}^\infty$ classify the line
with first Chern class $a$. Its map on $H_2$ sends $z$ to one,
with inverse $1\mapsto z$. On $S^2\times S^4$ take the line
pulled back from the positive generator on $S^2$.
Both manifolds are closed, simply connected, smooth, and spin,
with free middle-homology rank zero.
The homomorphism $P$ of \cite[Theorem 5, p.~796]{manZhubr75} is one quarter
of the pushed-forward Poincar\'e dual of $p_1$; it vanishes on
both maps. The homomorphism $h$ is the pushed-forward fundamental
class in $H_6(\mathbb{CP}^\infty)$, detected by the cube of the
universal degree-two class. It also vanishes, by $a^3=0$ and by
the square-zero generator of $S^2$.
Injectivity of $(P,h)$ gives the same spin bordism class.
The classification and middle-surgery construction in
\cite[Theorem~3, p.~795; \S4.1, p.~809]{manZhubr75}
give an oriented diffeomorphism
inducing the specified $H_2$ identification. Here the rank-zero
subclass is the one constructed there by killing the free middle
homology; $H_3=0$ has been established for both manifolds.
Therefore $X\cong S^2\times S^4$ as oriented smooth manifolds.
For every $k>0$, the choice $n=1$ satisfies the required coprimality
and completes the asserted existence for multiplicity $6k$.
\end{proof}

\subsection{Odd-star fillings and regular logarithmic transformations}

\begin{definition}[The two odd-star families]
\hypertarget{tag-Q120}{}\label{mapsIII:def:odd-star-logarithmic-data}
{\normalfont\scriptsize tag Q120}\par
Take one of the following marked configurations, with generator $G$
and the period convention of Definition~\ref{mapsIII:def:elliptic-marked-data}:
\[
\begin{array}{c|l|r|r}
 \mathrm{OS}&\text{ordered fibres}&p&b\\\hline
 46&I_3^*,I_1,I_1^\bullet,I_1&4&2\\
 50&I_1^*,I_3,I_1^\bullet,I_1&12&1.
\end{array}
\]
Set $P=pG$, $Q=G$ and $M=\mathcal O(P-O)$, and place the
degree-one linearization zero at $I_1^\bullet$. Use the original
$I_j^*$ filling, the primitive
weight-zero nodal fillings and the hexagonal cusp. At an additional
regular fibre perform a logarithmic transformation of order
$m\in\mathbb Z_{>0}$ with vector
$v=a e_1+b'e_2+c\delta+n\gamma$, where
$a,b',c,n\in\mathbb Z$ and $\gcd(m,a,b',c,n)=1$. Put
\[
 t=\gcd(m,n,a,b')\gcd(m,n,p,c),\qquad
 \mathcal G=\mathbb Z/t\oplus\mathbb Z/m\oplus\mathbb Z/(m|n|/t)
 \quad(n\ne0).
\]
For $n=0$ write $\mathcal T=\mathbb Z/t\oplus\mathbb Z/m$.
\end{definition}

\Needspace{6\baselineskip}
\begin{theorem}[Homology and smooth types of the odd-star families]
\hypertarget{tag-02KQ}{}\label{maps:src:D.2.16}
{\normalfont\scriptsize tag 02KQ}\par
For the family of Definition~\ref{mapsIII:def:odd-star-logarithmic-data},
the logarithmic action is free. For $n\ne0$ the integral groups are
\begin{equation}
\begin{aligned}
 \pi_1(X)&=H_1(X)=\mathbb Z/|n|,&H_5(X)&=0,\\
 H_2(X)&=\mathbb Z^b\oplus\mathcal G,
 &H_3(X)&=\mathbb Z^{2b}\oplus\mathcal G,\\
 H_4(X)&=\mathbb Z^b\oplus\mathbb Z/|n|.
\end{aligned}
\hypertarget{tag-0258}{}\label{maps:eq:D.T1}
\end{equation}
For $n=0$ the integral groups are
\begin{equation}
\begin{aligned}
 \pi_1(X)&=H_1(X)=H_5(X)=\mathbb Z,\\
 H_2(X)&=\mathbb Z^{b+2}\oplus\mathcal T,
 &H_3(X)&=\mathbb Z^{2b+2}\oplus\mathcal T,\\
 H_4(X)&=\mathbb Z^{b+2}.
\end{aligned}
\hypertarget{tag-025A}{}\label{maps:eq:D.T2}
\end{equation}
In both cases $H_0(X)=H_6(X)=\mathbb Z$ and $\chi(X)=2$.
The simply connected members are exactly those with $n=\pm1$.
Their oriented smooth type is
\begin{equation}
 X\cong\#^b(S^2\times S^4)\ \#\ \#^b(S^3\times S^3)
                      \ \#\ \widetilde\Sigma_0N_m.
\hypertarget{tag-025B}{}\label{maps:eq:D.T3}
\end{equation}
Here $N_m$ and its directed twisted suspension are those of
Definition~\ref{mapsIII:def:twisted-smale-suspension}. In particular
$\widetilde\Sigma_0N_1=S^6$, and $m=1$ gives the displayed connected
sums of sphere products.
\end{theorem}

\begin{proof}
In the period convention of Definition~\ref{mapsIII:def:odd-star-logarithmic-data},
write $\chi=\gamma^*$ and
$J=\left(\begin{smallmatrix}0&1\\-1&0\end{smallmatrix}\right)$.
Denote a positive base meridian by $x_i$. We orient fibre products
fibre first and use the complex orientation of $X$.

\emph{The marked periods and the line.}
Put $K(u)=I+uu^tJ$, $A=K(1,0)$, $B=K(1,-1)$ and $C=K(1,1)$.
Then
\[
 BC=\begin{pmatrix}-1&4\\0&-1\end{pmatrix},\qquad
 A^{j+4}BC=\begin{pmatrix}-1&-j\\0&-1\end{pmatrix},\qquad
 (A^{j+4}BC)A^{4-j}BC=I.
\]
These are the two marked words $(A^7BC)ABC$ and
$(A^5BC)(A^3)BC$ in \cite[Table 3, p.~35]{manFukae02}.
The source uses inverse elementary matrices composed from right to
left. Inverting its product gives the displayed left-to-right
positive product; applying the same operation recovers the source
word. In this marking the generator cocycles are
\[
 \begin{array}{c|rrrr}
 &g_1&g_2&g_3&g_4\\\hline
 46&(1,1)&(0,0)&(0,0)&(-1,-1)\\
 50&(-1,1)&(-2,0)&(0,0)&(-1,-1).
 \end{array}
\]
Their heights are $1/p$, their star residues are $3\bmod4$, and
the $I_3$ residue in row 50 is $1\bmod3$.

For the elliptic matrices $A_i$ in this word solve
$(A_i-I)z_i=g_i$ and put
\[
 h_i=pz_i,\quad b_i=-h_i^tJ,\quad
 \beta_i=-p g_i^tJA_i,\quad (w_1,w_2,w_3,w_4)=(0,0,1,0).
\]
The Poincar\'e-line coordinates and their inverses from
Theorem~\ref{maps:src:D.2.12} give
\[
 T_i=\begin{pmatrix}A_i&0&g_i\\\beta_i&1&\beta_i z_i+w_i\\0&0&1\end{pmatrix},
 \quad S_i=\begin{pmatrix}I&0&z_i\\b_i&1&b_i z_i\\0&0&1\end{pmatrix},
 \quad S_i^{-1}=\begin{pmatrix}I&0&-z_i\\-b_i&1&0\\0&0&1\end{pmatrix}.
\]
The map $S_i$ has source lattice $\Lambda$ and target lattice
$S_i\Lambda$; thus these are inverse torus maps also when $z_i$ is
rational. They satisfy
$T_i=S_i^{-1}(\operatorname{diag}(A_i,1,1)+w_iE_{34})S_i$.
In particular the two star matrices are
\begin{equation}
 T_1^{46}=\begin{pmatrix}
 -1&-3&0&1\\0&-1&0&1\\-4&-8&1&3\\0&0&0&1
 \end{pmatrix},\qquad
 T_1^{50}=\begin{pmatrix}
 -1&-1&0&-1\\0&-1&0&1\\-12&-24&1&3\\0&0&0&1
 \end{pmatrix}.
\hypertarget{tag-025N}{}\label{maps:eq:D.T6}
\end{equation}
The remaining matrices are
\[
 \begin{split}
 T_2^{46}&=I+E_{12},\qquad
 T_2^{50}=\begin{pmatrix}1&3&0&-2\\0&1&0&0\\0&24&1&-16\\0&0&0&1\end{pmatrix},\\
 T_3&=\begin{pmatrix}2&1&0&0\\-1&0&0&0\\0&0&1&1\\0&0&0&1\end{pmatrix},\qquad
 T_4=\begin{pmatrix}0&1&0&-1\\-1&2&0&-1\\-p&p&1&-p\\0&0&0&1\end{pmatrix}.
 \end{split}
\]
Multiplication gives $T_1T_2T_3T_4=I$. Each matrix fixes $\delta$
and $\chi$, and preserves the integral exterior covector
\[
 \omega=(p,0,0,0,0,-1)
\]
in the order $(12,13,14,23,24,34)$.

Multiplication by $p$ kills the stated component residues, so $P$ is
narrow at every fibre. Proposition~\ref{maps:src:D.2.19}, with
$(P,Q)=(pG,G)$ and $\langle P,Q\rangle=p(1/p)=1$, identifies
$\operatorname{Hom}(t_Q^*M,M)$ with the pullback of the degree-one
base line. Its multiplication by the specified divisor/base functions
and inverse multiplication by their reciprocals preserve this marking.
Choose the base section with its sole simple zero at $I_1^\bullet$
and multiply it by a sufficiently small constant to make it contracting.
The narrow height formula gives $P.O=p/2-1$ and $\deg(M|O)=p/2$. The
inverse-iteration section construction of Theorem~\ref{maps:src:D.2.12}, with
$p/2$ integral inverse iterates, fixes the reduced original meridians;
forward iteration is its inverse. The additional smooth place uses
the unshifted section frame.

\emph{The original star filling and its boundary marking.}
Let $N$ be the original minimal elliptic-surface neighbourhood over a
closed disc and let $F$ be its $I_j^*$ divisor, where $j=2b-1$.
Its affine $D_{j+4}$ tree has four multiplicity-one outer spheres and
$2b$ multiplicity-two spheres in a chain. The plumbing neighbourhood
retracts onto this tree of spheres: use the normal-disc collapse in
a component chart and the collapse to the two coordinate branches in
a crossing chart, agreeing on their collars. Thus
\[
 \pi_1(N)=0,\qquad H_2(N)=\mathbb Z^{4+2b},\qquad
 H_1(N)=H_3(N)=0.
\]
The oriented components give the indicated basis.

At the star the pairs $(z_*,b_*)$ are
\[
 \left((1/4,-1/2),(-2,-1)\right),\qquad
 \left((3/4,-1/2),(-6,-9)\right).
\]
The component map in this marked basis is the explicit isomorphism
\[
 \mathbb Z^2/(A_1-I)\mathbb Z^2\longrightarrow\mathbb Z/4,
 \quad[c_1,c_2]\longmapsto[2c_1+c_2],\qquad
 [k]\longmapsto[(0,k)]
\]
in the reverse direction. The row annihilates both relation columns
modulo four, is surjective, and $\det(A_1-I)=4$; hence these maps
are inverse. Both $g_1$ have coordinate three. The section cocycle
is minus the Abel--Jacobi increment, so the invariant point $z_*$,
with $(A_1-I)z_*=g_1$, has the component of the forward section
$-Q$, namely $-3=1\bmod4$.

In the convention $e_1=1,e_2=\tau$, the coordinate conversion to
\cite[\S2.4, pp.~6--7; Proposition~3.5(1),(3), pp.~8--10]{manUU21} is
$(u_1,u_2)\mapsto(r,s)=(u_2,u_1)$, inverted by the same swap.
Our two points are its invariant leaves $(1/2,1/4)$ and
$(1/2,3/4)$. They extend to order-four torsion sections representing
the local forward torsion part $q_{\rm tor}$ of $-Q$.
The exact sequence (2.2) there gives a local Lie logarithm $R$ with
$\operatorname{Exp}(R)=-Q-q_{\rm tor}$. Consequently
\[
 \tau_s=t_{q_{\rm tor}+\operatorname{Exp}(sR)},\qquad
 \tau_s^{-1}=t_{-q_{\rm tor}+\operatorname{Exp}(-sR)}
 \quad(0\leq s\leq1)
\]
are inverse translations and give an isotopy on the entire original
smooth $N$. Their extension over its components is the translation
action of \cite[\S\S7.6--7.7, pp.~33--34]{manSS10}. In the notation
$Q_0=-q_{\rm tor}$ the endpoint mapping torus below is $T(t_{Q_0})$.

The zero component degrees imply $c_1(M|N)=0$, since those components
are a basis of $H_2(N)$. In a nonzero smooth frame the contraction is
$(y,w)\mapsto(\tau_1y,a(y)w)$ with $|a(y)|<1$. The angular
$\delta$ action is free: equality with a nonzero iterate of the
contraction would equate a constant norm with a strictly contracted
or expanded norm. Put $u=-\log|w|$ and
$r(y)=-\log|a(\tau_1^{-1}y)|$. The angular orbit space has the
literal description
\[
 \{(y,u):0\leq u\leq r(y)\}/(y,r(y))\sim(\tau_1^{-1}y,0).
\]
Choose $k>0$ with $k\max r<1$ and a smooth nonnegative bump
$b_y$ of integral one supported in the middle third of $[0,r(y)]$.
The increasing diffeomorphism
\[
 f_y(u)=\int_0^u\bigl(k+(1-kr(y))b_y(v)\bigr)\,dv
\]
is $ku$ at the lower end and $1+k(u-r(y))$ at the upper end.
Thus $[y,u]\mapsto[y,f_y(u)]$, inverted by $f_y^{-1}$, identifies
the orbit space smoothly with $T(\tau_1^{-1})$, including the
boundary derivatives. For a cutoff $\zeta$ equal to zero and one
near the two ends, the further inverse maps are
\[
 [y,t]\longmapsto[t_{\operatorname{Exp}(-\zeta(t)R)}y,t],
 \qquad [y,t]\longmapsto[t_{\operatorname{Exp}(\zeta(t)R)}y,t].
\]
Their endpoint identity is
$\tau_0^{-1}t_{\operatorname{Exp}(-R)}=\tau_1^{-1}$.

The angular circle bundle over this mapping torus restricts to the
trivial line bundle on $N$. The Wang sequence makes its restriction
map in $H^2$ injective, because $H^1(N)=0$, so the angular bundle is
trivial. Explicitly remove the integral \v{C}ech cocycle of its transition
logarithms and use a partition of unity to write their differences
as $f_i-f_j$; multiplication by $\exp(2\pi i f_i)$ gives the
trivialization and the reciprocal multipliers give its inverse.
This proves, with the positive $\gamma$ direction,
\begin{equation}
 U\cong S^1_\delta\times T(\tau_0^{-1}:N\to N)
       =S^1_\delta\times T(t_{Q_0}:N\to N).
\hypertarget{tag-025L}{}\label{maps:eq:D.T4}
\end{equation}
The chosen section frame fixes its origin disc and period marking.
An integral phase adjustment along $\gamma$ is multiplication by
$\exp(2\pi i k t)$, with inverse multiplication by its reciprocal.

Label the outer components $L_0,L_1,R_0,R_1$ and the chain
$D_0,\ldots,D_{2b-1}$. Label the arms so that the forward translation
$\tau_0$ acts as
\[
 L_0\mapsto R_0\mapsto L_1\mapsto R_1\mapsto L_0,
 \qquad D_k\mapsto D_{2b-1-k}.
\]
The arm action follows from its component residue one; preservation
of the incidence tree forces the chain action. All orientations are
preserved. If $P_*$ is this permutation, the positive mapping-torus
monodromy is $P_*^{-1}$. The two Wang complexes are identified by
identity on the target and $-P_*$ on the source, since
$(P_*^{-1}-I)(-P_*)=P_*-I$; the source inverse is $-P_*^{-1}$.

There are $\rho=b+1$ orbits, of sizes $4,2,\ldots,2$.
Coinvariants send a component vector to the sum of its coefficients
in each orbit; a section puts that coefficient on its first component.
Invariants repeat an orbit coefficient on all its components; the
inverse reads one representative. Write $U_0$ for the outer-orbit
class and $U_1,\ldots,U_b$ for the chain-orbit classes, and put
$D=\delta\times\gamma$. The product and Wang maps give
\begin{equation}
 \pi_1(U)=\mathbb Z^2,\qquad
 (H_0,H_1,H_2,H_3,H_4)(U)
 = (\mathbb Z,\mathbb Z^2,\mathbb Z^{\rho+1},
                       \mathbb Z^{2\rho},\mathbb Z^\rho),\qquad H_5(U)=0.
\hypertarget{tag-025M}{}\label{maps:eq:D.T5}
\end{equation}
The preceding orbit maps and sections give the integral coordinates.

The local $H_1$ map and its section are
\[
 (x_1,x_2,x_\delta,x_\gamma)
 \longmapsto(b_*x_{\rm ell}+x_\delta,x_\gamma),\qquad
 (d,g)\longmapsto d\delta+g\gamma.
\]
Indeed the Poincar\'e-line coordinate change is
$x_\delta'=x_\delta+b_*x_{\rm ell}$, inverted by subtraction.
Let $W=U_0+\cdots+U_b$. The elliptic fibre specializes to $4W$,
using multiplicity one on the four outer components and multiplicity
two on each size-two central orbit. The complete fibre $H_2$ map is
\[
 x\longmapsto l(x)W+d(x)D,
 \qquad
 \begin{array}{c|rrrrrr|rrrrrr}
 &\multicolumn{6}{c|}{l}&\multicolumn{6}{c}{d}\\\hline
 46&4&0&-2&0&-1&0&0&0&-2&0&-1&1\\
 50&4&0&-2&0&-3&0&0&0&-6&0&-9&1.
 \end{array}
\]
To obtain these rows, the elliptic coordinate change is
$(x_{\rm ell},x_\gamma)\mapsto(x_{\rm ell}+z_*x_\gamma,x_\gamma)$,
with inverse subtraction. Its invariant area functional is
$4(x_{12}+z_{*,2}x_{14}-z_{*,1}x_{24})$, with scale fixed by
the fibre class $4W$. The other row is the cup product of
the two local $H_1$ coordinates,
$x_{34}+b_{*,1}x_{14}+b_{*,2}x_{24}$.
The rational invariant row space has dimension two, and these values
on $e_{12}$ and $e_{34}$ determine every component coordinate.
Integrally, the Smith entries of $\bigwedge^2T_1-I$ are
$(1,1,1,1,0,0)$, and a section of $(l,d)$ sends $(1,0)$ to
$-e_{24}-e_{34}$ in row 46 and to $e_{12}+e_{24}+9e_{34}$ in
row 50; it sends $(0,1)$ to $e_{34}$. Hence
$\operatorname{im}(\bigwedge^2T_1-I)=\ker(l,d)$ and these maps
identify the full fibre quotient with $\mathbb ZW\oplus\mathbb ZD$.

The fixed lattice is
\begin{equation}
 \ker(T_1-I)=\mathbb Z\delta\oplus\mathbb Z\ell_*,\qquad
 \ell_*=
 \begin{cases}(-1,2,0,4),&\mathrm{OS46},\\(-3,2,0,4),&\mathrm{OS50}.
 \end{cases}
\hypertarget{tag-025P}{}\label{maps:eq:D.T8}
\end{equation}
Solving the four equations gives precisely these integral vectors;
the smallest nonzero $\gamma$ coordinate is four. Their geometric
circles are the line circle and the fourth forward orbit over the
invariant torsion section. Both extend over the base disc. Their
normal sweeps therefore bound. The local linearization extends as a
nonzero unit over the disc and has zero winding. The angular row is one on $\delta$ and zero
on $\ell_*$. Thus these two sweeps give the full fixed-lattice
summand of the boundary Wang sequence.

For completeness put $M_*=\partial U$. Direct multiplication gives
\[
 N_*=T_1^2-I=j\bigl(e_1+(p/2)\delta\bigr)(0,2,0,-1),
 \qquad N_*^2=0,\quad\operatorname{rank}N_*=1.
\]
Then $T_1^{2r}=I+rN_*$ and an odd power has eigenvalue $-1$, so
the monodromy has infinite order. The Smith entries of $T_1-I$ are
$(1,1,0,0)$. The section splits the Wang sequences, giving
$H_1(M_*)=\mathbb Z^3$ and $H_2(M_*)=\mathbb Z^4$.
In the two fibre and two sweep coordinates the entire $H_2$ map is
$(u,v,r,s)\mapsto uW+vD$. Its kernel is the primitive sweep summand.
The local $H_1$ map is onto, with kernel the positive meridian of
the extended origin disc. Consequently
\begin{equation}
 H_2(U)/\operatorname{im}H_2(M_*)\xrightarrow{\ \cong\ }\mathbb Z^b,
 \quad [\textstyle\sum_{i=0}^b x_iU_i+yD]
           \longmapsto(x_1-x_0,\ldots,x_b-x_0),
 \quad (z_i)\longmapsto[\textstyle\sum_{i=1}^b z_iU_i].
\hypertarget{tag-025O}{}\label{maps:eq:D.T7}
\end{equation}
These are inverse after imposing $W=D=0$.

The Poincar\'e--Lefschetz map can also be specified integrally.
Let $A_i$ be $\delta$ times the mapping-torus cycle of the constant
invariant vector on orbit $i$; these form a basis of $H_4(U)$.
In the dual coordinates $H_2(U,M_*)=H^4(U)=\mathbb Z^{b+1}$,
the boundary to the meridian kernel is
\[
 (k_0,\ldots,k_b)\longmapsto(k_0+2k_1+\cdots+2k_b)\,x_1.
\]
This follows because the regular four-torus specializes to
$A_0+2\sum_{i=1}^b A_i$. Its section is the positive origin disc,
whose coordinates are $(1,0,\ldots,0)$. Intersecting component
representatives with invariant orbit sums gives the component-to-relative
matrices
\[
 \begin{pmatrix}-2&2&0\\1&-2&1\\0&1&-1\end{pmatrix},\qquad
 \begin{pmatrix}-2&2\\1&-1\end{pmatrix}.
\]
Both kill $W$. Adjoining $e_0$ to the last $b$ columns gives a
unimodular matrix; its inverse has first row $(1,2,\ldots,2)$.
This identifies their images with the complete meridian kernel.
For the product marking $\operatorname{diag}(-A^j,I_2)$ the relation
Smith entries are $(1,4,0,0)$ and boundary $H_1$ is
$\mathbb Z^3\oplus\mathbb Z/4$; the induced homology map therefore
obstructs a boundary diffeomorphism to this $M_*$. An invertible
intertwiner with an order-two monodromy would send the nonzero $N_*$
to zero upon squaring, which is impossible. These computations specify
the boundary information retained by the original star.

\emph{The other fillings and the marked exterior.}
The primitive nodal coordinate matrices are
\[
 H_2^{46}=I,\quad
 H_2^{50}=\bigl[(1,0,8,0),(0,1,0,1),\delta,(0,2,0,3)\bigr],
 \quad
 H_4=\bigl[(1,1,p,0),e_2,\delta,(0,1,0,1)\bigr].
\]
They are unimodular and conjugate $T_2,T_4$ to $I+E_{12}$;
their inverses are their adjugates divided by their determinants.
The component action at the original $I_3$ is a unit cycle.
Thus the ruled-normalization maps, translation isotopy, primitive
vanishing circle and specialization maps of Theorem~\ref{maps:src:D.2.12}
apply in these coordinates. At the cusp use
\[
 H_3=\operatorname{diag}\left(
             \begin{pmatrix}1&0\\-1&1\end{pmatrix},1,1\right),
 \qquad H_3^{-1}T_3H_3=I+E_{12}+E_{34}.
\]
It fixes $\delta,\gamma$. The angular map sends
$H_3(x,0,y,0)$ to $(x,y)$; column inclusion is its inverse.
Here the line shear is zero. The primitive hexagonal
specialization and the normal framing of Theorem~\ref{maps:src:D.2.14}
therefore have exactly this marking.

Remove the four singular discs and the regular-log disc and set
$T_0=I$. The real period map $[x]\mapsto\exp(\Pi_bx)$, with
the group logarithm modulo $\Pi_b\Lambda$ as inverse, identifies
the exterior $V$ with the marked flat torus bundle. Its graph map
and the five based boundary maps are
\[
 D_\Gamma=[0\ T_1-I\ T_2-I\ T_3-I],\qquad
 j_i(x)=\text{block }i\text{ inclusion }(0\leq i\leq3),
\]
\[
 j_4(x)=-(T_1T_2T_3x,T_2T_3x,T_3x,x).
\]
Reversing the ordered boundary walk gives the successive transports
in this formula, and multiplication verifies
$D_\Gamma j_4=I-T_4^{-1}$.
The fixed-loop column lists are
\[
 \begin{split}
 F_1&=(\delta,\ell_*),\\
 F_2^{46}&=(e_1,\delta,\gamma),\qquad
 F_2^{50}=((1,0,8,0),\delta,(0,2,0,3)),\\
 F_3&=((1,-1,0,0),\delta),\qquad
 F_4=((1,1,p,0),\delta,(0,1,0,1)).
 \end{split}
\]
Put $B_\Gamma=[j_1F_1\ j_2F_2\ j_3F_3\ j_4F_4]$ and
\begin{equation}
 K=[B_\Gamma\ j_0e_1\ j_0e_2\ j_0\gamma],\qquad
 LK=I_{13},\qquad
 (Lj_0x)_{11,12,13}=(x_1,x_2,x_4).
\hypertarget{tag-025Q}{}\label{maps:eq:D.T9}
\end{equation}
Here is an explicit integral formula for the inverse used in this
display. For a thirteen-element row set $R$, let $P_R$ be its
coordinate selector and $K_R=P_RK$. Set
\[
 R_a=\{1,\ldots,16\}\setminus\{12,13,16\},\qquad
 R_b=\{1,\ldots,16\}\setminus\{10,14,16\}.
\]
In row 46, $\det K_{R_b}=1$ and take
$L=\operatorname{adj}(K_{R_b})P_{R_b}$. In row 50 the two
determinants are $-2,-3$ and take
$L=\operatorname{adj}(K_{R_a})P_{R_a}-
\operatorname{adj}(K_{R_b})P_{R_b}$.
The adjugate identity proves $LK=I$ in both cases.
If $D_0$ is the first three rows of $D_\Gamma$, an integral right
inverse $C_0$ has the following nonzero rows; every omitted row is zero:
\[
 \begin{array}{c|r|rrr}
 \mathrm{OS}&\text{row}&\multicolumn{3}{c}{C_0}\,\\\hline
 46&6&3&3&-2\\&8&8&8&-5\\&13&2&1&-1\\\hline
 50&5&4&4&0\\&6&-3&-3&0\\&8&-24&-24&3\\
 &13&-18&-27&3\\&14&0&8&0\\&16&48&48&-8
 \end{array}
\]
Multiplication gives
\[
 D_0C_0=I_3,\qquad D_\Gamma K=0,\qquad
 KL(I_{16}-C_0D_0)=I_{16}-C_0D_0.
\]
Thus $K$ and $L$ are inverse on the complete graph kernel, and
$\operatorname{im}D_\Gamma=\ker\chi$. In particular
\[
 Lj_0\delta=(-1,0,0,-1,0,0,-1,0,-1,0,0,0,0)^t.
\]
This is the based chain identity $\sum_{i=0}^4j_i\delta=0$.

For the exterior-square map on the four free graph edges
$B_2=[0\ \bigwedge^2T_1-I\ \bigwedge^2T_2-I\ \bigwedge^2T_3-I]$ put
\[
 F=(e_{12}+pe_{34},e_{13},e_{14},e_{23},e_{24}),\qquad s=-e_{34}.
\]
The six columns $[s,F]$ are unimodular, $\omega s=1$ and
$\omega F=0$. An integer $24\times5$ matrix $Z$ with $B_2Z=F$
is specified by these nonzero rows:
\[
 \begin{array}{c|r|rrrrr}
 \mathrm{OS}&\text{row}&\multicolumn{5}{c}{Z}\\\hline
 46&9&-5&0&2&0&2\\&10&-9&-1&6&1&3\\&11&2&0&-1&0&-1\\
 &12&12&2&-11&-3&-4\\&21&-4&0&2&0&1\\\hline
 50&8&-94&-2&68&1&63\\&9&3&0&-2&0&-2\\
 &10&49&1&-38&-1&-35\\&12&-94&-2&68&1&63\\
 &17&2&0&-1&0&-1\\&21&0&0&0&0&-1
 \end{array}
\]
It follows that the full fibre coinvariant maps are
$[x]\mapsto\omega x$ and $k\mapsto[ks]$, in opposite directions.
The section and the real-period triangles with vertices $0,x,x+y$
split $H_2(V)$ as $\mathbb Zs\oplus\ker D_\Gamma$; subtraction
of those triangles gives the inverse chain map. The constant vertical
form vanishes on every fixed-loop normal sweep, fixing this splitting.

At the star, $(p/4)l-d=\omega$, so its fibre identifications are
$W=(p/4)s$ and $D=-s$. The basis change
$(U_0,U_1,\ldots,U_b,D)\leftrightarrow(W,U_1,\ldots,U_b,D)$
has inverse $U_0=W-\sum_{i=1}^bU_i$. It leaves the independent
classes $U_1,\ldots,U_b$ in the star. The ten full normal
sweeps remove $B_\Gamma$. The four original meridians bound, and
their base relation also bounds $x_0$. The complete local $H_1$
maps have the displayed sections, so these eliminations are identity
blocks of the full Mayer--Vietoris maps, inverted by opposite row
and column additions. Before attaching the smooth log the result is
\[
 \pi_1=\mathbb Z\langle\gamma\rangle,\qquad
 H_2=\mathbb Z\langle U_1,\ldots,U_b,s,t_{e_1},t_{e_2},t_\gamma\rangle.
\]
Indeed the star kills $e_1-b_{*,1}\delta$ and
$e_2-b_{*,2}\delta$, and the cusp kills $\delta$. These generate
the primitive $\ker\chi$; every other vanishing vector lies in it.
The remaining fundamental-group maps are $\chi$ and
$k\mapsto\gamma^k$, in opposite directions.

\emph{The regular logarithm and its full transfer.}
On the root disc $t=z^m$, the deck map in real-period coordinates is
\[
 (y,z)\longmapsto(y+v/m,e^{2\pi i/m}z).
\]
The order of $v/m$ in the central torus is
$m/\gcd(m,a,b',c,n)=m$, so the action is free there; the base
rotation proves freeness off the central fibre. Its finite
holomorphic quotient is therefore smooth for every stated vector.
The exact collar of Theorem~\ref{maps:src:D.2.13} applies in the unshifted
smooth section frame. For clarity write $z=r e^{i\theta}$ and let
$J_z$ be the complex structure in real-period coordinates. Put
\[
 b_{\rm sec}(z)=\frac{\log q_{\rm sec}(z^m)}{2\pi i},\quad
 u_{\rm sec}=\operatorname{Re}(b_{\rm sec})\delta+
                    \operatorname{Im}(b_{\rm sec})J_z\delta,
 \quad B_{\rm col}=\frac{\log r}{2\pi}J_zv-u_{\rm sec}.
\]
The collar isotopy and its inverse are
\[
 h_\tau(y,z)=y-\frac{\theta}{2\pi}v+\tau B_{\rm col}(z),\qquad
 h_\tau^{-1}(y,z)=y+\frac{\theta}{2\pi}v-\tau B_{\rm col}(z).
\]
The remaining correction is deck invariant and periodic, so a radial
cutoff extends it over the old exterior. The positive normal slope
in that exterior is exactly $mx_0-v$.

The central lattice and its quotient map are
\[
 \Lambda_v=\Lambda+\mathbb Z(v/m)=\mathbb Z^5/\mathbb Zq_0,
 \quad q_0=(-a,-b',-c,-n,m)^t,
 \quad[(z,k)]\longmapsto z+kv/m.
\]
Choose $\alpha\in\mathbb Z$ and an integral covector $\ell$ with
$\alpha m+\ell(v)=1$. The smooth normal-line trivialization is
\[
 [(y,z)]\longmapsto([y],e^{-2\pi i\ell(y)}z),\qquad
 ([y],w)\longmapsto[(y,e^{2\pi i\ell(y)}w)]
\]
in the reverse direction. The deck multiplier is
$e^{2\pi i(1-\ell(v))/m}=1$, proving descent and independence of
representatives. The positive normal disc is complex oriented. The
relative groups in degrees $3,2,1$ are $\Lambda_v,\mathbb Z,0$;
orient a relative three-generator as minus its fibre loop times this
disc.

All quotient coordinates can be chosen integrally with inverse maps.
Successive two-row B\'ezout operations give $Uq_0=e_1$ with
$U\in GL_5(\mathbb Z)$. Let $H$ be its last four rows, $E$ the
last four columns of $U^{-1}$, and $r_0$ its first row. Then
\[
 HE=I_4,\qquad Hq_0=0,\qquad EH+q_0r_0=I_5.
\]
Thus $[z]\mapsto Hz$ and $u\mapsto[Eu]$ are inverse maps between
$\Lambda_v$ and $\mathbb Z^4$.

The attaching collar sends a fibre loop $\lambda\in\Lambda$ to
$-\omega(\lambda,v)s+mj_0(\lambda)$, and sends $x_0$ to $j_0(v)$.
The component coordinates $U_i$ are zero, because these collar
chains and their addition triangles are in the regular exterior.
In source order $(e_1,e_2,\delta,\gamma,x_0)$ and target order
$(s,t_{e_1},t_{e_2},t_\gamma)$ the remaining transfer is
\begin{equation}
 R=\begin{pmatrix}
 -pb'&pa&n&-c&0\\
 m&0&0&0&a\\
 0&m&0&0&b'\\
 0&0&0&m&n
 \end{pmatrix},\qquad Rq_0=0.
\hypertarget{tag-0259}{}\label{maps:eq:D.T10}
\end{equation}
The full fourteen-by-fifteen map, before the ten sweep eliminations,
has columns $(0,e_i)$ for $1\leq i\leq10$, followed by
\[
 \bigl(-\omega(e_i,v),mLj_0e_i\bigr)\quad(1\leq i\leq4),
 \qquad (0,Lj_0v)
\]
in $\mathbb Zs\oplus\mathbb Z^{13}$. Subtract their first ten
graph coordinates using the first ten columns. The source change
$\left(\begin{smallmatrix}I&-C_{\rm low}\\0&I\end{smallmatrix}\right)$,
inverted by replacing $-C_{\rm low}$ with $C_{\rm low}$, reduces
this map to ten identity columns and $R$. Every additional $U_i$
row is zero, and the target changes act as identity on that free
summand. The other internal relative star cycles have zero boundary.

The coefficient map
\[
 \mathbb Z[p,m,n,a,b,c]\longrightarrow\mathbb Z[m,n,a,b',c],
 \qquad b\longmapsto b',\quad p\longmapsto4\text{ or }12
\]
carries the transfer and source relation in
\eqref{maps:eq:D.S7} to $R$ and $q_0$ here, in the same ordered
free source and target bases. The letter change $b\leftrightarrow b'$
is inverted by the reverse change. The prime-local calculation in
the proof of Theorem~\ref{maps:src:D.2.13} uses only integrality of
$p$, $m>0$ and the five-entry freeness condition. Its pivot identities
therefore specialize under this coefficient map. At each prime the
same unit cases are exhaustive, and the inverse row additions,
permutations and reciprocal units remain defined.

Put $d_1=\gcd(m,n,a,b')$ and $d_2=\gcd(m,n,p,c)$.
A common prime divisor would divide all five entries of the freeness
condition, so $\gcd(d_1,d_2)=1$ and $t=d_1d_2$ divides $m,n$.
The resulting determinantal divisors are
\[
 1,\quad t,\quad mt,\quad m^2|n|.
\]
Thus the Smith entries are $(1,t,m,m|n|/t)$ for $n\ne0$ and
$(1,t,m,0)$ for $n=0$. Expanding the five signed maximal minors
also gives the polynomial identity
$m^2n(-a,-b',-c,-n,m)$.
For literal quotient maps, keep the integer row and column operations
in a Smith reduction $U_RREV_R=\operatorname{diag}$.
The target map is $[x]\mapsto[U_Rx]$, inverted by $U_R^{-1}$,
and the source map is $EV_R$ with the inverse quotient coordinates
above, with the primitive five-generator relation $q_0$.

\emph{The full degree-one map and the zero-$n$ cycles.}
The five boundary $H_1$ ranks are $(5,3,4,3,4)$ and the five local
ranks are $(4,2,3,2,3)$. The exterior $H_1$ basis is
$(\gamma,x_0,x_1,x_2,x_3)$, so the complete Mayer--Vietoris map is
nineteen by nineteen. Use the smooth boundary coordinates
$(e_1,e_2,\delta,\gamma,x_0)$, the star coordinates
$(\delta,\gamma,x_1)$, and the surviving node or cusp coordinates
followed by their meridians. The fifth meridian maps to
$-x_0-x_1-x_2-x_3$. On a fibre loop $z$, its exterior coordinate
is $\chi(z)\gamma$; its local coordinate is the displayed local
quotient map, with the minus sign in Mayer--Vietoris.

Use the columns $E$ followed by $q_0$ as the new smooth source
basis and the identity local bases at the other places. Order the
fourteen local section columns first and the five normals last.
The resulting matrix is
\[
 \begin{pmatrix}A&G_5\\-I_{14}&0\end{pmatrix},\qquad
 G_5=\begin{pmatrix}
 -n&0&0&0&0\\m&0&0&0&-1\\0&1&0&0&-1\\
 0&0&1&0&-1\\0&0&0&1&-1
 \end{pmatrix}.
\]
The target change
$\left(\begin{smallmatrix}I_5&A\\0&I_{14}\end{smallmatrix}\right)$,
whose inverse replaces $A$ by $-A$, removes the upper left block.
Thus the full map has the same remaining kernel and cokernel as
$G_5$, with $\det G_5=-n$. Its kernel is zero for $n\ne0$ and
is primitively generated by $(1,m,m,m,m)$ for $n=0$.
Van Kampen has the same normal map: the four old fillings leave
$\mathbb Z\langle\gamma\rangle$ and $mx_0-v$ maps to $-n\gamma$.
Hence $\pi_1(X)=\mathbb Z/n\mathbb Z$, with $\mathbb Z/0=\mathbb Z$.
The quotient map is $\chi$ modulo $n$ and the inverse sends the
generator to the loop $\gamma$.

For $n\ne0$ this gives
$H_2(X)=\mathbb Z^b\oplus\operatorname{coker}R$.
At $n=0$ the extra relative two-generator has a closed cycle lift.
The discs $D_0+m(D_1+D_2+D_3+D_4)$ have boundary $-v$.
Since $\chi(v)=0$, the right inverse $C_0$ specifies monodromy
cylinders with boundary $v$. Add them, with their real-period
triangles, to obtain a closed integral cycle $Z_0$ of relative
coefficient one. This splits the extra free summand.
The other free quotient coordinate at $n=0$ is
\[
 f(s,x_1,x_2,x_\gamma)
  =\frac{ms+pb'x_1-pa x_2+c x_\gamma}{d_2}.
\]
It kills all columns of $R$, and its coefficients are primitive:
prime by prime, freeness gives
$\gcd(m,pb',pa,c)=\gcd(m,p,c)=d_2$.
A B\'ezout preimage $z_f$ of one is a section. Together with the
component basis and $Z_0$, these are the free coordinates in
\eqref{maps:eq:D.T1}--\eqref{maps:eq:D.T2}; subtracting their chosen
lifts leaves exactly the specified Smith torsion coordinates.

There is also an explicit primitive relative three-cycle at $n=0$.
Put $d=\gcd(m,a,b')$, choose $A_0,B_0$ with $cA_0+dB_0=1$, and set
\[
 u=(-a/d,-b'/d,0,0,m/d),\qquad
 \lambda=A_0u+B_0e_\delta.
\]
Before the relation, $\ker R=\mathbb Zu+\mathbb Ze_\delta$ and
$q_0=du-ce_\delta$. The quotient coordinate is
$[\alpha u+\beta e_\delta]\mapsto c\alpha+d\beta$, with inverse
$1\mapsto[\lambda]$; in $\Lambda_v$ this is $\delta/d$.
Its full relative vector has the additional ten-sweep coordinates
$-(m/d)Lj_0\delta$. Its fibre boundary is
$(a/d)e_{13}+(b'/d)e_{23}$, bounded by the exact cylinders
$Z(0,a/d,0,b'/d,0)^t$. Subtract these and $m/d$ times the
$\delta$ cylinder over the punctured base from the Thom chains.
The resulting closed three-cycle has primitive relative coefficient
one, using the extending star $\delta$ sweep.

\emph{All integral degrees and the duality maps.}
The original star has Euler characteristic zero by
\eqref{maps:eq:D.T4}. The regular pieces, logarithmic fibres,
boundaries and primitive nodal fibres also have Euler characteristic
zero. The primitive hexagonal normalization and conductor calculation
of Theorem~\ref{maps:src:D.2.12} gives the remaining contribution two.
Thus $\chi(X)=2$, and a point and the complex fundamental class
generate $H_0(X)$ and $H_6(X)$.

Integral duality and universal coefficients give the other groups
with the following explicit coordinates. The integral functional
dual to each free $H_2$ generator, represented by a cellular cocycle
and capped with $[X]$, is the corresponding $H_4$ generator.
For torsion $H_1$, the character $\gamma\mapsto1/|n|$ in
$\mathbb Q/\mathbb Z$, followed by the integral Bockstein and cap
product, is its $H_4$ torsion generator. Linking and free intersection
coordinates are the inverse maps. At $n=0$, $\chi\frown[X]$
generates $H_5$, with evaluation as inverse; otherwise $H_5=0$.

For each Smith $H_2$ summand of order $d$, the character taking its
chosen generator to $1/d$ and the others to zero gives, by Bockstein
and cap product, the corresponding $H_3$ torsion generator.
Its inverse reads the linking coefficient. The Bockstein identifies
these characters with the full $\operatorname{Ext}(H_2,\mathbb Z)$,
so all torsion terms are accounted for.
Free middle coordinates can be obtained with inverse integer
maps as follows. Triangulate the displayed plumbing, collars and
Thom chains. In a Smith reduction of $\partial_3$, the last columns
of its source matrix give the integral cycle basis and the last rows
of its inverse give the coefficient map. Express $\partial_4$ in
that basis and reduce it with inverse row and column operations.
The remaining zero coordinates give representative free cycles and
their dual functionals. On this free quotient the intersection form
$I$ is alternating and unimodular. Choose a primitive $A_1$ and
then $B_1$ with $I(A_1,B_1)=1$ by B\'ezout. The projection
\[
 x\longmapsto x-I(x,B_1)A_1-I(A_1,x)B_1
\]
onto their orthogonal complement is inverted by adding these two
intersection coefficients. Repetition gives pairs $A_i,B_i$, with inverse maps
\[
 (r_i,s_i)\longmapsto\sum_i(r_iA_i+s_iB_i),\qquad
 [x]\longmapsto\bigl(I(x,B_i),I(A_i,x)\bigr)_i.
\]
The zero-$n$ primitive relative three-cycle may be chosen as $A_1$;
its relative coefficient and the integral cycle-basis inverse extend
to the dual cocycle determining $B_1$.
Finally $2-2b_1+2b_2-b_3=\chi(X)=2$ gives $b_3=2b$ for
$n\ne0$ and $b_3=2b+2$ for $n=0$. This proves every group and
isomorphism in \eqref{maps:eq:D.T1}--\eqref{maps:eq:D.T2}.

\emph{The angular action and its normal framing.}
The preceding fundamental-group map proves that the simply connected
members are exactly $n=\pm1$. Here $t=1$ and
$H_2(X)=\mathbb Z^b\oplus(\mathbb Z/m)^2$.
Multiplication of the line coordinate defines one global $\delta$
circle action: it commutes with the line transitions, the contraction
and the logarithmic translations. It is free on the star by the
norm argument above. At a primitive node its short vanishing vector
and $\delta$ span a primitive rank-two lattice, as the displayed
$H_2,H_4$ bases show; $\delta$ therefore remains primitive on the
elliptic base of the ruled normalization. It acts freely on that
surface and its identified sections. At the central logarithmic
torus the integral functional
\[
 (x_1,x_2,x_\delta,x_\gamma,x_0)
       \longmapsto x_\delta-(c/n)x_\gamma
\]
kills $q_0$ and sends $\delta$ to one, for either sign $n=\pm1$.
Thus the action is free there and on every regular fibre.

The cusp angular map and its inverse are
\[
 H_3(x,0,y,0)\longmapsto(x,y),\qquad
 (x,y)\longmapsto H_3(x,0,y,0).
\]
They carry $\delta$ to $(0,1)$ and the full primitive short lattice
to $\mathbb Z^2$. The corresponding character maps, fixing the base
coordinate, identify the triangle cones with those used in the
fixed-curve calculation of Theorem~\ref{maps:src:D.2.14}.
Their inverse character maps extend on the same affine charts.
The complete period group acts by translations of this fan, with
its nonzero unit factors unchanged. That calculation therefore gives
one fixed curve $C\simeq\mathbb P^1$ and trivial stabilizers elsewhere.

In these characters the two adjacent normal charts have transition
$(x',y',w')=(wx,wy,w^{-1})$, with inverse
$(x,y,w)=(w'x',w'y',(w')^{-1})$. Thus the normal bundle is exactly
\eqref{maps:eq:D.S11}. The positive rescaling, negative-weight
conjugation and displayed $SU(2)$ map $U$ in that proof, inverted by
division, conjugation and $U^*$, respectively, supply the scalar
normal framing here. Reverse the core parametrization together with
the negative-weight conjugation, preserving the complex orientation.
The invariant normal exponential and its inverse identify the
neighbourhood with $C\times D^4$ carrying scalar multiplication.

The canonical line on the star has a holomorphic trivialization.
By \cite[Theorem 6.8, p.~26]{manSS10}, $K_N$ is pulled back
from a line on the disc and is trivial. A nonvanishing generator
$\eta$ is translation invariant: on the regular locus it is a base
differential wedged with the invariant elliptic differential, and
the identity extends across the fibre. In a holomorphic line frame
$\Omega=\eta\wedge dw/w$ is independent of that frame. Indeed
$w'=f(y)w$ changes it by $\eta\wedge d\log f=0$ on the
two-dimensional complex surface. The same identity proves
$\phi^*\Omega=\Omega$. Multiplication by $\Omega$, with inverse
division, trivializes the canonical line on the original star quotient.
Every free $H_2(X)$ generator is represented there, so $c_1(X)$
evaluates to zero on all of them. Since $H_1(X)=0$,
$H^2(X;\mathbb Z)=\operatorname{Hom}(H_2(X),\mathbb Z)$;
hence $c_1(X)=0$ and $X$ is spin.

\emph{The smooth quotient and the zero free Euler class.}
The scalar normal quotient is
\[
 h(z_1,z_2)=
 (2\operatorname{Re}(z_1\bar z_2),2\operatorname{Im}(z_1\bar z_2),
                 |z_1|^2-|z_2|^2),\qquad
 \mathcal H(z)=h(z)/|z|\ (z\ne0),\quad\mathcal H(0)=0.
\]
It identifies exactly the circle orbits and has
$|\mathcal H(z)|=|z|$. The Hopf sections
\[
 s_N(v)=\left(\sqrt{\frac{1+v_3}{2}},
                \frac{v_1-iv_2}{\sqrt{2(1+v_3)}}\right),\qquad
 s_S(v)=\left(\frac{v_1+iv_2}{\sqrt{2(1-v_3)}},
                \sqrt{\frac{1-v_3}{2}}\right)
\]
give inverse local orbit charts $(q,t)\mapsto |q|t s_N(q/|q|)$
and its southern version. They specify the smooth quotient charts.
Orient their meridian sphere $\sigma$ so the transition has winding
$+1$. The free part is oriented with the horizontal five-plane
before the positive circle direction; this extends across the charts.
Denote the closed quotient by ${\mathcal Q}$, put
$X_0=X\setminus(C\times\mathring D^4)$ and
${\mathcal Q}_0={\mathcal Q}\setminus(C_{\mathcal Q}\times\mathring D^3)$, and let
$p_0:X_0\to {\mathcal Q}_0$ have Euler class $e$.

Transversality of discs of dimensions at most three to the
codimension-four $C$ gives isomorphisms on $\pi_1,\pi_2,H_2$
for $j:X_0\to X$. The fibration sequence and Hurewicz give
\[
 \pi_1({\mathcal Q}_0)=0,\qquad
 0\longrightarrow H_2(X)\xrightarrow{p_{0*}j_*^{-1}}H_2({\mathcal Q}_0)
   \xrightarrow{\langle e,-\rangle}\mathbb Z\longrightarrow0,
 \qquad k\longmapsto k\sigma
\]
as an explicit splitting. Excision gives
$H_3({\mathcal Q},{\mathcal Q}_0)=\mathbb Z$ with boundary $\sigma$ and
$H_2({\mathcal Q},{\mathcal Q}_0)=0$. Consequently, for $i:{\mathcal Q}_0\to {\mathcal Q}$,
\[
 \Theta=i_*p_{0*}j_*^{-1}:H_2(X)\xrightarrow{\cong}H_2({\mathcal Q}).
\]
Its inverse lifts a class to $a\in H_2({\mathcal Q}_0)$, subtracts
$\langle e,a\rangle\sigma$, lifts the resulting sphere to $X_0$,
and applies $j_*$. The pulled-back circle bundle has Euler number
zero and thus a section. Two sections differ by a map $S^2\to S^1$
and are homotopic; a different first lift differs by a multiple of
$\sigma$ and gives the same corrected class. Van Kampen gives
$\pi_1({\mathcal Q})=0$.

In the residue sequence the next map after
$\operatorname{res}:H^2({\mathcal Q}_0;\mathbb Z)\to H^0(C_{\mathcal Q};\mathbb Z)$
sends one to the Poincar\'e dual of $[C_{\mathcal Q}]$. Since
$\operatorname{res}(e)=1$, exactness proves $[C_{\mathcal Q}]=0$.
On $X_0$, $TX_0=p_0^*T{\mathcal Q}_0\oplus\mathbb R$. The mod-two
Gysin sequence makes $\ker p_0^*$ in degree two the span of
$e\bmod2$. Its nonzero element evaluates to one on $\sigma$,
whereas $w_2({\mathcal Q}_0)$ evaluates to zero there because
$T{\mathcal Q}_0|_\sigma=TS^2\oplus\mathbb R^3$. It follows that
$w_2({\mathcal Q}_0)=0$. Excision makes restriction from $H^2({\mathcal Q};\mathbb Z/2)$
injective, so ${\mathcal Q}$ is spin.

Choose representatives of $U_1,\ldots,U_b$ on central components
at distinct suspension angles. Their quotient images are disjoint
embedded spheres $V_1,\ldots,V_b$ representing the free basis
$\Theta(U_i)$, and their component lifts trivialize the angular
bundle on them. They are disjoint from $C_{\mathcal Q}$.
Set ${\mathcal Q}'={\mathcal Q}\setminus\bigcup_iV_i$. Transversality of loops and
two-discs, using $2+2<5$, makes ${\mathcal Q}'$ simply connected. Excision
gives $H_3({\mathcal Q},{\mathcal Q}')=\mathbb Z^b$ and $H_2({\mathcal Q},{\mathcal Q}')=0$. The map
$H_3({\mathcal Q})\to\mathbb Z^b$ is intersection with the $V_i$ and is
surjective by the unimodular free $H_3$--$H_2$ pairing. Thus
$H_2({\mathcal Q}')\to H_2({\mathcal Q})$ is an isomorphism. It follows that $C_{\mathcal Q}$
is nullhomologous and, by Hurewicz, nullhomotopic already in ${\mathcal Q}'$.

Choose a parametrized standard $C_*\simeq S^2$ in a five-disc
inside ${\mathcal Q}'$. Theorem B(b) of \cite[pp.~367--368]{manBarden65}
applies to this arbitrary ambient manifold: the required connectivity
here is connectedness and simple connectedness, and $2\cdot5>3\cdot2+3$.
It gives a parametrized isotopy from $C_{\mathcal Q}$ to $C_*$ in ${\mathcal Q}'$.
Extend its velocity to a time-dependent vector field supported in a
compact neighbourhood of the isotopy track, disjoint from the $V_i$.
Its flow is the ambient isotopy and the backward flow is its inverse.
It preserves the component parametrizations. Transport the bundle
and all normal maps by these inverse flows.

Now $C_*$ is inside a disc $D^5$ disjoint from every $V_i$.
Restriction $H^2({\mathcal Q};\mathbb Z)\to
H^2({\mathcal Q}\setminus\mathring D^5;\mathbb Z)$ is an isomorphism.
The $V_i$ remain a free homology basis, and the Euler class evaluates
to zero on each of them. Therefore the Euler class is zero on the
entire punctured quotient. The \v{C}ech logarithm and partition-of-unity
construction above gives a bundle trivialization there, inverted by
reciprocal gauges.

\emph{The precisely framed reconstruction.}
Inside $D^5$, the complement of the standard fixed sphere has
$H^2\cong\mathbb Z$ detected by its normal residue. The residue
one bundle is the Hopf bundle in the normal charts above and in the
model of Theorem~\ref{maps:src:D.2.14}. Their \v{C}ech cocycles give an
isomorphism covering identity on the punctured disc. An oriented
normal rotation $S^2\to SO(3)$ lifts to $SU(2)$ by path lifting
on the simply connected sphere; its inverse is the conjugate
transpose. The normal framing adjustment is unique up to homotopy
because $\pi_2(SO(3))=0$. On the overlap $S^4$ times an interval,
the two bundle maps differ by an $S^1$-valued gauge. Its $H^1$ is
zero, so the gauge has a real logarithm. A cutoff exponential makes
the maps equal on a full collar; the reciprocal exponential is the
inverse modification. Thus the maps and their inverses glue smoothly
and preserve the specified orientation.

The resulting model is
\[
 M_-({\mathcal Q})=(({\mathcal Q}\setminus\mathring D^5)\times S^1)
                 \cup_g(S^4\times D^2),\qquad
 g(x,z)=(\alpha_zx,z),
\]
where $g$ goes from complement to attached boundary and $\alpha_z$
rotates the first real two-plane through $\arg z$. Its inverse uses
$\alpha_z^{-1}$. The action on the cap is
$t(x,w)=(\alpha_tx,tw)$. Its boundary orbit coordinate is
$\alpha_z^{-1}x$, whose composition with $g$ is $x$; hence its
quotient restores the disc with the original boundary parametrization.

The rotation loop lifts to
$\cos(\pi t)+\sin(\pi t)e_1e_2$ in $\operatorname{Spin}(5)$,
with endpoints $1,-1$, and represents the nonzero element of
$\pi_1(SO(5))$. The definition of $\widetilde\Sigma_0{\mathcal Q}$ uses
the opposite gluing direction. The exact comparison is
\[
 A=\operatorname{diag}(1,-1,-1,1,1),\qquad
 A\alpha_zA^{-1}=\alpha_z^{-1}.
\]
Let $A_s$ rotate the $(2,3)$ plane through $\pi s$. On a boundary
collar define $F_{\mathcal Q}(r,x)=(r,A_{\zeta(r)}x)$ with $\zeta=1$ near
the boundary and zero off the collar; its inverse replaces $A_s$
by $A_s^{-1}$. The maps $F_{\mathcal Q}\times\mathrm{id}$ and
$A\times\mathrm{id}$ on the two pieces agree because
$A\alpha_z^{-1}=\alpha_z A$. Their determinants are positive,
and their inverses glue. Combined with the preceding reconstruction
they give the oriented diffeomorphism
\[
                       X\longrightarrow\widetilde\Sigma_0{\mathcal Q}
\]
with exactly the rotation loop in Theorem~\ref{maps:src:D.2.14}.

The spin linking form of ${\mathcal Q}$ is alternating by
\cite[Lemma F, p.~372]{manBarden65}. For $m>1$ choose $x$ of
order $m$ and $y$ with linking value $\lambda(x,y)=1/m$.
Nondegeneracy supplies $y$, and pairing a relation with $x,y$
shows that they form a basis of $(\mathbb Z/m)^2$. The inverse
coordinates are
\[
 a\longmapsto
 (\operatorname{num}_m\lambda(a,y),-\operatorname{num}_m\lambda(a,x)),
 \qquad(r,s)\longmapsto rx+sy,
 \quad\operatorname{num}_m([k/m])=k\bmod m.
\]
At $m=1$ these are the unique maps of trivial groups. Use the
component basis for the free factors. These maps preserve the full
linking form and $w_2=0$, so the oriented classification and models
of \cite[Theorems 2.2--2.3, pp.~376--377]{manBarden65} give an
orientation-preserving diffeomorphism
\[
                  {\mathcal Q}\longrightarrow\#^b(S^2\times S^3)\#N_m
\]
inducing them. Its inverse induces their inverse coordinates. The
parametrized disc isotopy theorem makes it agree with the source and
target disc coordinates. It then lifts to the twisted suspensions
by its product with identity on the punctured pieces and identity
on the common cap. The inverse uses the inverse base diffeomorphism.

The suspension distributes over this connected sum with its precise
framing as follows. Write a punctured connected sum as the boundary
connected sum of its punctured factors with handle $[-1,1]\times D^4$.
Choose a rotationally symmetric star-shaped dumbbell for its boundary
$S^4$. Rotation of the first two transverse coordinates has an
equivariant radial diffeomorphism to the round sphere, with radial
inverse, and is the stated loop $\alpha$. The middle separating
hypersurface in the suspension is
\[
             (D^4\times S^1)\cup_{\alpha|S^3}(S^3\times D^2).
\]
The map $(x,z)\mapsto(\alpha_z^{-1}x,z)$ on the first piece and
identity on the second identifies it with
$S^5=\partial(D^4\times D^2)$, after invariant corner rounding.
Its inverse uses $\alpha_z$. The cap-to-complement transition first
sends $y$ to $\alpha_z y$, so these displayed maps agree exactly
on the join. Cutting gives the punctured twisted suspension of each
factor. Use these same boundary parametrizations on both sides and
the orientation reversal in a transverse coordinate fixed by
$\alpha$. The rotations cancel on the joined boundary, giving the
standard oriented connected-sum gluing. The inverse rejoins the
handle collars. Repeating gives
\[
 \widetilde\Sigma_0{\mathcal Q}\cong
   \#^b\widetilde\Sigma_0(S^2\times S^3)\#\widetilde\Sigma_0N_m.
\]
For the product factors, \cite[Theorem~B(2), p.~3; Proposition~5.6, pp.~22--23]{manGR25}
uses the trivial clutching map, with base $S^3$ and fibre $S^2$, to give
\[
 \widetilde\Sigma_0(S^2\times S^3)
                      \cong(S^2\times S^4)\#(S^3\times S^3).
\]
The swaps $S^3\times S^2\to S^2\times S^3$ and
$S^4\times S^2\to S^2\times S^4$ preserve orientation, since
the products of their dimensions are even. The inverse clutching
maps give the inverse diffeomorphisms. Connected-sum factors can be
interchanged by $(t,u)\mapsto(-t,\rho u)$ on the neck, where
$\rho$ is a reflection of $S^5$; this is orientation preserving
and self-inverse. These maps prove \eqref{maps:eq:D.T3} for both
signs $n=\pm1$ with the chosen complex orientation.

\emph{The associated surgery and the sphere case.}
The complete cusp $H_2$ specialization has the section
$(u,v,w,z)\mapsto(u,0,v,0,w,z)$ followed by
$\bigwedge^2H_3$. Composing it with $\omega$ gives the precise
map from a local cusp class to its coefficient of $s$. Thus the
fixed curve is a multiple of $s$, and the $\delta$ column of
\eqref{maps:eq:D.T10} imposes $ns=0$. For $n=\pm1$ it follows
that $[C]=0$. Hurewicz makes the parametrized sphere nullhomotopic.

Use the scalar normal framing to replace $S^2\times D^4$ by
$D^3\times S^3$, with the Hopf action on $S^3$, and call the
result $Y$. The embedding isotopy theorem used in
Theorem~\ref{maps:src:D.2.14} applies since $6>2\cdot2+1$;
$\pi_2(SO(4))=0$ fixes the framing homotopy class. Transport the
surgery to an unknotted sphere in a ball. In
$S^6=\partial(D^3\times D^4)$ the replacement is the double
of $D^3\times S^3$, namely $S^3\times S^3$. Identity outside the
ball and these product maps, transported by the isotopy and its
inverse, give $Y\cong X\#(S^3\times S^3)$.

On the quotient this is the framed replacement of
$S^2\times D^3$ by $D^3\times S^2$, performed in the disc
disjoint from the $V_i$. The product-boundary maps in
$S^5=\partial(D^3\times D^3)$ give
$B_5=Y/S^1\cong {\mathcal Q}\#(S^2\times S^3)$. The Euler class is zero
on the punctured ${\mathcal Q}$ and one on the new Hopf $S^2$. Accordingly
the exact oriented pair is
\[
 (B_5,e_Y)\cong
 \left(\#^{b+1}(S^2\times S^3)\#N_m,(x,0,\ldots,0)\right),
 \qquad H_2(B_5)=\mathbb Z^{b+1}\oplus(\mathbb Z/m)^2.
\]
It is simply connected and spin. Pull back its spin frame bundle
and adjoin the oriented vertical line using an invariant connection.
The circle action lifts to this pulled-back spin bundle, proving
that the spin structure of $Y$ is projectable. The
diffeomorphism already constructed for $X$, extended by identity
over the extra punctured product and its neck, gives the corresponding
decomposition of $Y$ with its inverse. Thus the unstabilized map for
$X$ has been constructed directly.

For $m=1$, Definition~\ref{mapsIII:def:twisted-smale-suspension}
has $N_1=S^5$. The oriented boundary map, its inverse and radial
projection in Theorem~\ref{maps:src:D.2.14} identify
$\widetilde\Sigma_0N_1$ with $S^6$ in this same directed gluing
convention. Removing this connected summand uses the oriented disc
identification, inverted by inserting the same disc. This gives the
stated $m=1$ specialization.
\end{proof}

\section{Original additive profiles}
\subsection{Two additive profiles}
\begin{corollary}[Original profiles of types II and IV]
\hypertarget{tag-02KR}{}\label{maps:src:D.2.17}
{\normalfont\scriptsize tag 02KR}\par
The formulas and oriented diffeomorphism classification
\eqref{maps:eq:D.T1}--\eqref{maps:eq:D.T3}, for every free smooth-fibre
log, hold for
\[
\begin{array}{c|c|l|r|r}
 \mathrm{OS}&\mathrm{profile}&\text{ordered fibres}&p&b\\\hline
 46&II&I_3^*,II,I_1^\bullet&4&2\\
 50&IV&I_1^*,IV,I_1^\bullet&12&1.
\end{array}
\]
Use $P=pG$, $Q=G$, the indicated simple linearization zero, and the
original fillings at both additive fibres.
\end{corollary}

\begin{proof}
We use $\Lambda=\mathbb Z\langle e_1,e_2,\delta,\gamma\rangle$,
with positive physical line circle $\delta$ and positive suspension
period $\gamma$. Put $J=\left(\begin{smallmatrix}0&1\\-1&0\end{smallmatrix}\right)$
and use exterior order $(12,13,14,23,24,34)$.

\emph{The marking and the line.}
Let $N=3$ for the first row and $N=1$ for the second. The primitive
star and nodal cycles $u,v$ have monodromies $-T_u^N,T_v$.
The finite peripheral has trace $1$ at $II$ and $-1$ at $IV$.
Since
$\operatorname{tr}(-T_u^NT_v)=N(u^tJv)^2-2$, one has
$(u^tJv)^2=1$. Choose the sign of $v$ with $u^tJv=1$ and use
$e_1=u,e_2=v-u$, inverted by $u=e_1,v=e_1+e_2$.
The based elliptic tuple in the stated fibre order is
\[
 A_* =\begin{pmatrix}-1&-N\\0&-1\end{pmatrix},\quad
 A_f=\begin{pmatrix}N-2&1\\-1&0\end{pmatrix},\quad
 A_c=\begin{pmatrix}0&1\\-1&2\end{pmatrix},\qquad A_*A_fA_c=I.
\]
The junction--Picard--Abel maps of \ref{mapsD:repair:collision-realization}, applied
in these cycles, give the generator cocycles
\[
 \begin{array}{c|ccc}
 &g_*&g_f&g_c\\\hline
 II&(1,1)&(0,0)&(-1,-1)\\
 IV&(-1,1)&(-2,0)&(-1,-1).
 \end{array}
\]
Here the section cocycle is minus the Abel--Jacobi increment.
For an explicit integral check, every closed admissible cocycle is
\[
 g_f=(a_0,b_0),\quad
 g_*=(a_0+(N-1)b_0+c_0,c_0),\quad
 g_c=(b_0-c_0)(1,1).
\]
The inverse reads $g_f$ and the second coordinate of $g_*$.
Its two coboundary columns are $(N-3,-1,0)$ and $(1,-1,-2)$.
Adjoining the displayed generator as first column gives determinant
$1$ for $II$ and $-1$ for $IV$, so its adjugate gives the inverse
section-quotient map. A cut-lift change by $z\in\mathbb Z^2$
adds $(A_i-I)z$ to each $g_i$; subtraction is its inverse.

The root lattices are $D_7$ and $D_5\oplus A_2$, of rank seven.
Their maximum local height corrections are $7/4$ and $23/12$,
both below two. The height formula excludes nonzero torsion, and
the unimodular rational N\'eron--Severi lattice gives generator
heights $1/4$ and $1/12$. These force $G.O=0$ and the maximal
local corrections. The star residues are those in
\ref{maps:src:D.2.16}; at $IV$ the component coordinate can be
written
\[
 [x,y]\longmapsto[-x-y]\in\mathbb Z/3,\qquad
 [k]\longmapsto[(-k,0)]
\]
on $\mathbb Z^2/(A_f-I)\mathbb Z^2$ in the reverse direction.
The relation columns map to zero and the determinant is three,
so these maps are inverse; $g_f=(-2,0)$ has coordinate two.
Thus $P$ is narrow everywhere, and $\langle P,Q\rangle=1$.
The zero-multidegree argument of \ref{maps:src:D.2.12} gives the
base line of degree one. Its section has its sole simple zero at
the marked $I_1$. The equality $\deg(M|O)=p/2$ fixes the same
inverse-iteration reduced frame as in \ref{maps:src:D.2.16}.

Solving $(A_i-I)z_i=g_i$, the Poincar\'e-line map and its inverse
use $b_i=-p z_i^tJ$ and
\[
 S_i=\begin{pmatrix}I&0&z_i\\b_i&1&0\\0&0&1\end{pmatrix},\qquad
 S_i^{-1}=\begin{pmatrix}I&0&-z_i\\-b_i&1&0\\0&0&1\end{pmatrix}.
\]
Their target lattice is exactly $S_i\Lambda$.
With weights $(0,0,1)$, the star matrix is the corresponding
matrix in \eqref{maps:eq:D.T6}, and the other two matrices are
\[
 T_f^{II}=\begin{pmatrix}1&1&0&0\\-1&0&0&0\\0&0&1&0\\0&0&0&1\end{pmatrix},\qquad
 T_f^{IV}=\begin{pmatrix}-1&1&0&-2\\-1&0&0&0\\-24&0&1&-16\\0&0&0&1\end{pmatrix},
\]
\[
 T_c=\begin{pmatrix}0&1&0&-1\\-1&2&0&-1\\-p&p&1&1-p\\0&0&0&1\end{pmatrix}.
\]
Direct multiplication gives $T_*T_fT_c=I$. All preserve
$\omega=(p,0,0,0,0,-1)$; on the elliptic inclusion
$i(x)=(x,0,0)$, one has $\omega(iu,iv)=p\,u^tJv$.

\emph{The two original additive neighbourhoods.}
Write $E_f$ for the original minimal elliptic-surface neighbourhood.
Its central fibre is a cuspidal rational curve at $II$ and three
spheres meeting at one point at $IV$. Resolve its central divisor
by point blowups. A plumbing retraction fixes that divisor pointwise
and descends through the proper blowdown quotient, because every
collapsed exceptional tree remains fixed. In a crossing chart it
is the collapse of the squared radii
\[
 (r_1^2,r_2^2)\longmapsto
 (r_1^2-t\min(r_1^2,r_2^2),r_2^2-t\min(r_1^2,r_2^2));
\]
the one-component normal collapses agree after choosing nested
collars. Radial base transport, with reverse flow as inverse,
identifies its small neighbourhood with the given fibre-disc
neighbourhood. Thus $E_f$ retracts onto $S^2$ at $II$ and
$\bigvee^3S^2$ at $IV$. In particular $\pi_1(E_f)=H_1(E_f)=0$.

The local logarithm and invariant torsion sections of
\cite[\S2.4, (2.2), p.~6; Proposition 3.5(1),(3), pp.~8--10]{manUU21} give
$-Q=q_{\rm tor}+\operatorname{Exp}(R)$, with $q_{\rm tor}=0$
at $II$ and order three at $IV$. The inverse translations
\[
 t_{q_{\rm tor}+\operatorname{Exp}(sR)},\qquad
 t_{-q_{\rm tor}+\operatorname{Exp}(-sR)}
 \quad(0\leq s\leq1)
\]
extend over the original components by
\cite[\S7.6, pp.~33--34]{manSS10}. At $IV$ they cycle the three
components. The line has zero degree on each component and hence
is topologically trivial on $E_f$. The variable-roof diffeomorphism,
its monotone inverse, and the cutoff translation maps in
\ref{maps:src:D.2.16} apply to these contractions.
Their positive mapping-torus monodromy is the inverse forward
translation. With $Q_0=-q_{\rm tor}$ the resulting marked map is
\[
 U_f\cong S^1_\delta\times T(t_{Q_0}:E_f\to E_f),
\]
which is $E_{II}\times T^2$ in the first case. The physical circle
bundle is trivial because the Wang restriction in $H^2$ is
injective when $H^1(E_f)=0$; its zero restriction is $c_1(M|E_f)$.
The logarithm gauges and their reciprocals retain the boundary frame.

The component action has one orbit. Coinvariants sum coefficients
and have the section choosing one component; invariants repeat a
coefficient on the whole orbit, with inverse reading one component.
The product and Wang maps therefore give
\[
 \pi_1(U_f)=\mathbb Z^2,\qquad
 (H_0,H_1,H_2,H_3,H_4)(U_f)
 =(\mathbb Z,\mathbb Z^2,\mathbb Z^2,\mathbb Z^2,\mathbb Z),
 \qquad H_5(U_f)=0.
\]
Let $W$ be its component-orbit class and $D=\delta\times\gamma$.
The regular elliptic fibre specializes to $rW$, where $r=1,3$.
The full fibre specialization is $x\mapsto l(x)W+d(x)D$, with
\[
 \begin{array}{c|rrrrrr|rrrrrr}
 &\multicolumn{6}{c|}{l}&\multicolumn{6}{c}{d}\\\hline
 II&1&0&0&0&0&0&0&0&0&0&0&1\\
 IV&3&0&-2&0&-2&0&0&0&-8&0&-8&1.
 \end{array}
\]
For $IV$, $z_f=(2/3,-2/3)$ and $b_f=(-8,-8)$.
The elliptic coordinate change $x\mapsto x+z_fx_\gamma$ and
line change $x_\delta\mapsto x_\delta+b_fx$ have subtraction
as inverse. The first specialization row is its invariant elliptic
area multiplied by three; the second is the cup product of the
local $H_1$ coordinates $(x_\delta+b_fx,x_\gamma)$.
At $II$ these changes are identity. A section sends $(W,D)$ to
$(e_{12},e_{34})$ at $II$ and to
$(e_{12}+e_{14}+8e_{34},e_{34})$ at $IV$.

The remaining kernel basis at $IV$ is
\[
 (2,0,3,0,0,24),\quad(0,0,-1,0,1,0),\quad e_{13},\quad e_{23};
\]
together with the two section columns it is unimodular. At $II$
the kernel is $(e_{13},e_{14},e_{23},e_{24})$.
All nonzero determinantal divisors of $T_f-I$ and
$\bigwedge^2T_f-I$ are one, at ranks two and four.
The first image is exactly the kernel of the displayed local $H_1$
map. Explicit preimages of its two kernel generators are
\[
 \begin{array}{c|cc|cc}
 &\multicolumn{2}{c|}{\text{kernel}}&\multicolumn{2}{c}{\text{preimages}}\\\hline
 II&e_1&e_2&(-1,1,0,0)&(-1,0,0,0)\\
 IV&e_1+8\delta&e_2+8\delta&(1,-1,0,-2)&(1,-2,0,-2).
 \end{array}
\]
For the exterior square, the primitive rank-four image is contained
in the displayed primitive kernel and hence equals it. Thus both
integer quotient maps are complete.

The fixed lattices are
\[
 F_f^{II}=(\delta,\gamma),\qquad
 F_f^{IV}=(\delta,(-2,2,0,3)).
\]
Solving $T_fx=x$ gives exactly these bases, with minimal positive
$\gamma$ coordinate three in the second case. The second circle
there is the full order-three torsion orbit. It and the line circle
extend over the base disc, so their full normal sweeps bound.
The boundary Wang maps give $H_1(\partial U_f)=\mathbb Z^3$ and
$H_2(\partial U_f)=\mathbb Z^4$. Its entire $H_2$ map is
$(u,v,s,t)\mapsto uW+vD$, with primitive sweep kernel.
For the relative assertion, the invariant component sum times
$\delta\times\gamma$ generates $H_4(U_f)$; the regular four-torus
has coefficient one on it. The positive origin disc has intersection
one with that generator. Therefore
$H_2(U_f,\partial U_f)=\mathbb Z$ maps isomorphically onto the
meridian kernel, with that disc as section. This also proves the
surjectivity of the entire boundary specialization.

\emph{The complete exterior and regular log.}
The star specialization and its $b$ free component classes are the
same marked maps as in \ref{maps:src:D.2.16}. At the cusp use
\[
 H=[(1,1,p,0),e_2,\delta,e_2+\gamma],\qquad
 \det H=1,\qquad H^{-1}T_cH=I+E_{12}+E_{34}.
\]
Its inverse is its integer adjugate. The cusp fixed columns are
$F_c=((1,1,p,0),\delta)$; its local $H_1$ section columns are
$He_2,He_4$, with $\chi$ coordinates $(0,1)$.
At the star write
$F_*=(\delta,(-1,2,0,4))$ or $(\delta,(-3,2,0,4))$.

For the four punctures including the regular log, set
\[
 D_\Gamma=[0\ T_*-I\ T_f-I],\qquad
 j_i(x)=\text{block }i\text{ inclusion }(i=0,1,2),\qquad
 j_3(x)=-(T_*T_fx,T_fx,x).
\]
Then $D_\Gamma j_3=I-T_c^{-1}$. Put
\[
 B=[j_1F_*\ j_2F_f\ j_3F_c],\qquad
 K=[B\ j_0e_1\ j_0e_2\ j_0\gamma].
\]
Let $\mathsf D_0$ be the first three rows of $D_\Gamma$. An integer
right inverse $C$ has the following nonzero rows:
\[
 \begin{array}{c|r|rrr}
 &\text{row}&\multicolumn{3}{c}{C}\\\hline
 II&6&-3&0&1\\&8&-8&0&3\\&9&-2&-1&1\\\hline
 IV&5&52&88&-8\\&6&27&34&-5\\&8&-72&-120&11\\
 &9&-126&-189&21\\&12&96&144&-16
 \end{array}
\]
The determinant of $(K,C)$ is $-1$ at $II$ and $1$ at $IV$.
Let $L$ be the first nine rows of its inverse adjugate formula.
Multiplication gives
\[
\begin{gathered}{}
 \mathsf D_0C=I_3,\quad D_\Gamma K=0,\quad LK=I_9,\quad
\\{}
 KL(I-C\mathsf D_0)=I-C\mathsf D_0,\quad
\\{}
 Lj_0\delta=(-1,0,-1,0,0,-1,0,0,0)^t.
\end{gathered}
\]
Thus $K,L$ are inverse on the complete integral graph kernel.
The last three coordinates of $Lj_0x$ are $(x_1,x_2,x_4)$.

Set $B_2=[0\ \bigwedge^2T_*-I\ \bigwedge^2T_f-I]$ and
$F=(e_{12}+pe_{34},e_{13},e_{14},e_{23},e_{24})$.
An integer matrix $Z$ with $B_2Z=F$ has only the following nonzero rows:
\[
 \begin{array}{c|r|rrrrr}
 &\text{row}&\multicolumn{5}{c}{Z}\\\hline
 II&9&3&0&-2&0&0\\&10&3&-1&0&1&0\\&11&-2&0&1&0&0\\
 &12&-4&2&-3&-3&0\\&15&4&0&-2&0&-1\\\hline
 IV&8&-43&1&20&1&30\\&9&3&0&-2&0&-2\\
 &10&43&-1&-20&-1&-30\\&11&2&0&-1&0&-2\\
 &12&-72&2&27&1&42\\&15&-4&0&2&0&3
 \end{array}
\]
The basis $[-e_{34},F]$ is unimodular, so the fibre quotient is
$\mathbb Zs$, $s=-e_{34}$, with maps $[x]\mapsto\omega x$ and
$k\mapsto[ks]$. The real-period addition triangles and their inverse
coefficient maps give the exterior splitting, as in
\ref{maps:src:D.2.16}. The six normal sweeps remove $B$.
At the finite additive place $(p/r)l-d=\omega$, so
$W=(p/r)s,D=-s$; the star contributes exactly its free classes
$U_1,\ldots,U_b$. Before the regular log the surviving coordinates are
\[
 \pi_1=\mathbb Z\langle\gamma\rangle,\qquad
 H_2=\mathbb Z\langle U_1,\ldots,U_b,s,t_{e_1},t_{e_2},t_\gamma\rangle.
\]
Indeed $\mathsf D_0C=I$ makes the old vanishing relations generate the
primitive $\ker\chi$, and every original meridian bounds.
The inverse fundamental-group maps are $\chi$ and $k\mapsto\gamma^k$.

For the regular log, the free action
$(y,z)\mapsto(y+v/m,e^{2\pi i/m}z)$, the collar isotopy and its
inverse in \ref{maps:src:D.2.13} give the exact slope $mx_0-v$.
Its central lattice is
\[
 \Lambda_v=\mathbb Z^5/\mathbb Zq_0,\qquad
 q_0=(-a,-b',-c,-n,m)^t,\qquad[(z,k)]\longmapsto z+kv/m.
\]
Choose an integral completion $Uq_0=e_1$. If $H_v$ consists of
its last four rows, $E$ of the last four columns of $U^{-1}$, and
$r_0$ is its first row, then
$H_vE=I_4$ and $EH_v+q_0r_0=I_5$.
Thus $[z]\mapsto H_vz$ and $u\mapsto[Eu]$ are inverse cap coordinates.
The normal-line trivialization is
$[(y,z)]\mapsto([y],e^{-2\pi i\ell(y)}z)$, inverted by the
reciprocal multiplier, where $\alpha m+\ell(v)=1$.
Its deck multiplier is one and its normal orientation is complex.

In $\mathbb Zs\oplus\mathbb Z^9$ the entire transfer has six
columns $(0,e_i)$, followed by
\[
 (-\omega(e_i,v),mLj_0e_i)\quad(1\leq i\leq4),\qquad(0,Lj_0v).
\]
All $U_i$ coordinates are zero because these chains lie in the
regular exterior. Order the six sweep rows first. Subtract the
upper-right block using those six columns, with opposite additions
as inverse. This reduces the full $10\times11$ map to
$\operatorname{diag}(I_6,R)$, where $R$ is exactly
\eqref{maps:eq:D.T10} in the same ordered coordinates and $Rq_0=0$.
The primitive source change $\operatorname{diag}(I_6,E)$ gives
its ten-coordinate cap source. The unrestricted Smith proof
in \ref{maps:src:D.2.16} therefore applies by identity on these
four target and five source coordinates. Its entries are
$(1,t,m,m|n|/t)$ for $n\ne0$ and $(1,t,m,0)$ for $n=0$;
all free star coordinates are fixed by identity.

The boundary $H_1$ ranks are $(5,3,3,3)$ and the local ranks
$(4,2,2,2)$; the exterior basis is $(\gamma,x_0,x_1,x_2)$.
The complete Mayer--Vietoris map is consequently $14\times14$.
Use the ten local section columns first and the four normal columns
last. With the displayed primitive cap completion it has blocks
\[
 \begin{pmatrix}A&G_4\\-I_{10}&0\end{pmatrix},\qquad
 G_4=\begin{pmatrix}-n&0&0&0\\m&0&0&-1\\0&1&0&-1\\0&0&1&-1\end{pmatrix}.
\]
The target change
$\left(\begin{smallmatrix}I_4&A\\0&I_{10}\end{smallmatrix}\right)$,
whose inverse replaces $A$ by $-A$, removes the first block.
Here $\det G_4=n$ and its zero-$n$ kernel is primitively generated
by $(1,m,m,m)$; at nonzero $n$ it has zero kernel.
Van Kampen adds $-n\gamma$ and gives
$\pi_1(X)=\mathbb Z/n\mathbb Z$, with maps $\chi\bmod n$ and
the loop $\gamma$ in the reverse direction.

At $n=0$, the discs $D_0+m(D_1+D_2+D_3)$ have boundary $-v$.
The columns of $C$ give cylinders bounding $v$; add their
real-period triangles to obtain a closed two-cycle $Z_0$ of
relative coefficient one. The other free quotient coordinate is
\[
 (s,x_1,x_2,x_\gamma)\longmapsto
 \frac{ms+pb'x_1-pa x_2+c x_\gamma}{\gcd(m,p,c)},
\]
with a Bezout preimage of one as section. Its coefficients are
primitive because freeness gives
$\gcd(m,pb',pa,c)=\gcd(m,p,c)$. These maps give both zero-$n$
free summands and the same Smith torsion as in \eqref{maps:eq:D.T2}.
For the primitive relative three-cycle put $d=\gcd(m,a,b')$,
choose $A_0,B_0$ with $cA_0+dB_0=1$, and use
\[
 u=(-a/d,-b'/d,0,0,m/d),\qquad
 \lambda=A_0u+B_0e_\delta.
\]
The relation is $q_0=du-ce_\delta$, and the quotient coordinate
$[\alpha u+\beta e_\delta]\mapsto c\alpha+d\beta$ is inverted
by $1\mapsto[\lambda]$. Its sweep correction is
$-(m/d)Lj_0\delta$; its remaining fibre boundary is
$(a/d)e_{13}+(b'/d)e_{23}$, bounded by
$Z(0,a/d,0,b'/d,0)^t$. The corresponding Thom chains and
$\delta$ cylinder over the punctured base give a closed primitive
three-cycle, using the extending sweeps of both additive caps.

Both original additive fillings have Euler characteristic zero by
their free physical circle. The regular pieces and their boundaries
also contribute zero; the primitive hexagonal cusp contributes two.
Thus $\chi(X)=2$. The free $H_2$ coordinates, the cyclic $H_1$
coordinate, and every $H_2$ torsion coordinate coincide with
\eqref{maps:eq:D.T1}--\eqref{maps:eq:D.T2}. Apply the explicit
integral duality maps from \ref{maps:src:D.2.16}: dual free cocycles
capped with $[X]$ give $H_4$, torsion characters followed by
Bockstein and cap product give its torsion and that of $H_3$, and
linking gives the inverse coordinates. At $n=0$, $\chi\frown[X]$
gives $H_5$ and the preceding primitive three-cycle supplies the
extra middle pair. The free alternating middle pairing and
$2-2b_1+2b_2-b_3=2$ give ranks $2b$ and $2b+2$.
A point and the complex fundamental class give $H_0=H_6=\mathbb Z$.
This proves all integral groups in the assertion.

\emph{The oriented reconstruction and its auxiliary surgery.}
The fundamental group proves simple connectivity exactly at
$n=\pm1$. The global physical-circle action is free at each
original additive filling by the contraction norm argument.
At the regular log the integral functional
$x_\delta-(c/n)x_\gamma$ kills $q_0$ and sends $\delta$ to one
for either sign, proving freeness there. The exact cusp basis $H$
identifies its angular $\delta$ with $(0,1)$. It transports the
triangular fan, the two charts
\[
 (x,y,w)=\bigl(q/(X_1X_2),X_2,X_1\bigr),\qquad
 (x',y',w')=\bigl(q/X_2,X_1X_2,1/X_1\bigr),
\]
and their inverses from \ref{maps:src:D.2.14}. Their transition
$(x',y',w')=(wx,wy,w^{-1})$ gives the unique fixed curve and
normal bundle $\mathcal O(-1)_{+1}\oplus\mathcal O(-1)_{-1}$.
The positive rescaling, negative-weight conjugation, simultaneous
core orientation reversal and explicit $SU(2)$ filling there give
the scalar normal framing, with reciprocal rescaling and conjugate
transpose as inverse. Hence all other stabilizers are trivial.

Let ${\mathcal Q}$ be the smooth oriented quotient supplied by those Hopf
charts and ${\mathcal Q}_0$ the complement of the fixed-curve tube. The exact
Thom and Gysin maps of \ref{maps:src:D.2.16} give
\[
 \Theta:H_2(X)\xrightarrow{\cong}H_2({\mathcal Q}),\qquad
 H_2({\mathcal Q})=\mathbb Z^b\oplus(\mathbb Z/m)^2.
\]
Its inverse lifts a class to $a\in H_2({\mathcal Q}_0)$, subtracts
$\langle e,a\rangle\sigma$, lifts the resulting zero-Euler sphere
to the circle bundle, and includes it in $X$.
The residue of $e$ on $\sigma$ is one, so the fixed curve $C_{\mathcal Q}$
is nullhomologous. The star components still represent the full
free $H_2(X)$ basis. On the star the canonical generator
$\eta\wedge dw/w$ descends: $\eta$ is the translation-invariant
canonical form from \cite[Theorem 6.8, p.~26]{manSS10}, and
$\eta\wedge d\log f=0$ for a line-frame multiplier on the surface.
Multiplication and division by that form give the trivialization.
Thus $c_1(X)=0$ and $X$ is spin. The mod-two Gysin kernel is
$\langle e\bmod2\rangle$, whose nonzero element evaluates one
on $\sigma$, whereas $w_2({\mathcal Q}_0)$ evaluates zero there. This proves
that ${\mathcal Q}$ is spin as well.

The projected component spheres $V_i=\Theta(U_i)$ are disjoint,
represent a free basis, and have zero Euler class by their component
lifts. In ${\mathcal Q}'={\mathcal Q}\setminus\bigcup_iV_i$, transversality gives
simple connectivity and the unimodular intersection pairing makes
$H_3({\mathcal Q})\to\mathbb Z^b$ onto. Excision then gives
$H_2({\mathcal Q}')\cong H_2({\mathcal Q})$. Consequently the fixed curve is
nullhomotopic in ${\mathcal Q}'$. Barden's parametrized isotopy theorem
\cite[Theorem B(b), pp.~367--368]{manBarden65} moves it into a five-disc there.
The ambient flow and its reverse preserve every $V_i$.
The Euler class is zero on the whole punctured ${\mathcal Q}$, because it
vanishes on that free basis. Reciprocal \v{C}ech gauges trivialize the
angular bundle there.

These are precisely the hypotheses used in the explicit framed
reconstruction of \ref{maps:src:D.2.16}. The present cusp map $H$
and its inverse carry its Hopf cap maps; the gluing is
$(x,z)\mapsto(\alpha_zx,z)$, where $\alpha$ rotates one real
coordinate plane. Its nontrivial Spin lift has endpoints $1,-1$.
The orientation-preserving gluing-direction comparison uses
$A=\operatorname{diag}(1,-1,-1,1,1)$ and the collar rotation in
its $(2,3)$ plane, with inverse rotations on both pieces.
Thus it gives the map $X\to\widetilde\Sigma_0{\mathcal Q}$.
Barden's oriented linking-coordinate map
\cite[Lemma F, p.~372; Theorem 2.2, p.~376]{manBarden65} gives
${\mathcal Q}\cong\#^b(S^2\times S^3)\#N_m$.
The inverse-rotation neck maps and the product clutching of
\cite[Theorem B(2), p.~3; Proposition 5.6(1), pp.~22--23]{manGR25}, exactly as constructed
in \ref{maps:src:D.2.16}, give \eqref{maps:eq:D.T3}.
All these maps preserve the stated complex orientation for each
of $n=1,-1$. The sphere boundary map there also gives
$\widetilde\Sigma_0N_1=S^6$.

For the surgery assertion, the complete cusp specialization is
surjective, with its section followed by $\bigwedge^2H$.
The fixed curve therefore maps to a multiple of $s$; the
$\delta$ column of $R$ imposes $ns=0$, so it is nullhomologous
at $n=\pm1$. The scalar framed nullsphere replacement and the
inverse ambient isotopy give $Y=X\#(S^3\times S^3)$.
On the quotient they give $B_5={\mathcal Q}\#(S^2\times S^3)$, performed
in the disc disjoint from the $V_i$. Its Euler class is zero on
punctured ${\mathcal Q}$ and one on the new Hopf $S^2$. Pulling back the
spin frame bundle of $B_5$ and adjoining the oriented vertical
line with an invariant connection gives an invariant spin lift;
the spin structure on $Y$ is projectable. This verifies the
associated surgery as well as the direct unstabilized classification.
\end{proof}

\subsection{The three-component cusp}
\begin{definition}[The complex flag threefold]\hypertarget{tag-Q300}{}\label{maps:def:complex-flag-threefold}
{\normalfont\scriptsize tag Q300}\par
The complex flag threefold is the incidence variety
\[
 \operatorname{Fl}(\mathbb C^3)=
 \{(L,P)\in\mathbb P^2\times(\mathbb P^2)^*:L\subset P\},
\]
where $L$ is a line and $P$ is a plane in $\mathbb C^3$.
We give it its complex orientation.
\end{definition}

\begin{theorem}[Three-component cusp with original additive fillings]
\hypertarget{tag-02KS}{}\label{maps:src:D.2.18}
{\normalfont\scriptsize tag 02KS}\par
For OS50 profile $II$, with ordered fibres $II,I_3,I_1^*$,
take $P=12G$, $Q=G$ and the degree-one linearization whose
zero is the whole $I_3$ fibre. Use the original $II$ and $I_1^*$
fillings. At an additional regular fibre perform a free log of order
$m>0$ with vector $v=ae_1+b'e_2+c\delta+n\gamma$, where
$\gcd(m,a,b',c,n)=1$. Put
\[
 t=\gcd(m,n,a,b')\gcd(m,n,12,c),\qquad
 \mathcal G=\mathbb Z/t\oplus\mathbb Z/m\oplus\mathbb Z/(m|n|/t)
 \quad(n\ne0).
\]
For $n\ne0$ the integral homology groups are
\begin{equation}
\begin{aligned}
 \pi_1(X)&=H_1(X)=\mathbb Z/|n|,&H_5(X)&=0,\\
 H_2(X)&=\mathbb Z^3\oplus\mathcal G,
 &H_3(X)&=\mathbb Z^2\oplus\mathcal G,\\
 H_4(X)&=\mathbb Z^3\oplus\mathbb Z/|n|.
\end{aligned}
\hypertarget{tag-0269}{}\label{maps:eq:D.V1}
\end{equation}
For $n=0$, put $\mathcal T=\mathbb Z/t\oplus\mathbb Z/m$; then
\begin{equation}
\begin{aligned}
 \pi_1(X)&=H_1(X)=H_5(X)=\mathbb Z,\\
 H_2(X)&=\mathbb Z^5\oplus\mathcal T,
 &H_3(X)&=\mathbb Z^4\oplus\mathcal T,\\
 H_4(X)&=\mathbb Z^5.
\end{aligned}
\hypertarget{tag-026A}{}\label{maps:eq:D.V2}
\end{equation}
In both cases $H_0(X)=H_6(X)=\mathbb Z$ and $\chi(X)=6$.
The simply connected members are exactly those with $n=\pm1$.
Their oriented smooth type is
\begin{equation}
 X\cong\operatorname{Fl}(\mathbb C^3)\#(S^2\times S^4)
       \#(S^3\times S^3)\#\widetilde\Sigma_0N_m.
\hypertarget{tag-026B}{}\label{maps:eq:D.V3}
\end{equation}
The flag manifold has its complex orientation. The spin five-manifold
$N_m$ and the framed twisted suspension are those of
\ref{maps:src:D.2.14}; in particular $\widetilde\Sigma_0N_1=S^6$.
\end{theorem}


\begin{proof}
\emph{The marked periods and the central fibre.}

Write the marked lattice as
\(\Lambda=\mathbb Z(e_1,e_2,\delta,\gamma)\), put
\(J=\left(\begin{smallmatrix}0&1\\-1&0\end{smallmatrix}\right)\),
and let the positive \(\gamma\) period cover translation by \(-G\).
The marked realization of \ref{mapsD:repair:collision-realization} has elliptic matrices and section columns

\[
\begin{gathered}{}
 A_0=\begin{pmatrix}3&7\\-1&-2\end{pmatrix},\quad g_0=(-5,2)^t;
\\{}
 \quad A_c=\begin{pmatrix}7&12\\-3&-5\end{pmatrix},\quad g_c=(-8,4)^t;
\\{}
 \quad A_*=\begin{pmatrix}-2&-1\\1&0\end{pmatrix},\quad g_*=(0,1)^t.
\end{gathered}
\]

The order is II, I3, I1*. The respective solutions of
\((A_i-I)z_i=g_i\) are
\(z_0=(1,-1)\), \(z_c=(-4/3,0)\), and
\(z_*=(1/4,-3/4)\). Set
\(b_i=-12z_i^tJ\), \(\beta_i=-12g_i^tJA_i\), and
\((w_0,w_c,w_*)=(0,1,0)\). The Poincar\'e-line changes are

\[
 S_i=\begin{pmatrix}I_2&0&z_i\\b_i&1&0\\0&0&1\end{pmatrix},
 \qquad S_i^{-1}=\begin{pmatrix}I_2&0&-z_i\\-b_i&1&0\\0&0&1\end{pmatrix};
 \qquad b_i z_i=0.
\]

Their source and target lattices are \(\Lambda\) and \(S_i\Lambda\).
This specifies the inverse torus maps also for the two rational changes.
Conjugating the diagonal elliptic action and adding the stated line
weight gives

\[
 T_0=\begin{pmatrix}3&7&0&-5\\-1&-2&0&2\\12&48&1&-36\\0&0&0&1\end{pmatrix},
 \quad
 T_c=\begin{pmatrix}7&12&0&-8\\-3&-5&0&4\\48&96&1&-63\\0&0&0&1\end{pmatrix},
 \quad
 T_*=\begin{pmatrix}-2&-1&0&0\\1&0&0&1\\-24&-12&1&3\\0&0&0&1\end{pmatrix}.
\]

The matrices satisfy \(T_0T_cT_*=I\), fix \(\delta\) and
\(\chi=\gamma^*\), and preserve
\(\omega=12e_1^*\wedge e_2^*-\delta^*\wedge\gamma^*\).

For the cusp put
\(\lambda=(-2,1,-16,0)\), \(r=(-1,0,0,0)\), and
\(H=[\lambda,r,\delta,\gamma]\). Then

\[
 H=\begin{pmatrix}-2&-1&0&0\\1&0&0&0\\-16&0&1&0\\0&0&0&1\end{pmatrix},
 \qquad
 H^{-1}=\begin{pmatrix}0&1&0&0\\-1&-2&0&0\\0&16&1&0\\0&0&0&1\end{pmatrix},
 \qquad H^{-1}T_cH-I=3E_{12}+4E_{14}+E_{34}.
\]

The angular coordinate map is
\((X_{\rm ell},V)\mapsto(X_1,X_2)=(X_{\rm ell},X_{\rm ell}^{16}V)\);
its inverse is \((X_1,X_2)\mapsto(X_1,X_1^{-16}X_2)\).
The two multiplicative periods consequently have valuations
\((3,0)\) and \((4,1)\), and have the maps

\[
 \begin{gathered}
 (X_1,X_2,q)\longmapsto
 (u_{1,\rm per}(q)q^{3a+4b}X_1,u_{2,\rm per}(q)q^bX_2,q),\\
 (a,b)\in\mathbb Z^2,
 \end{gathered}
\]

where each \(u_{j,\rm per}\) is the corresponding product of the original
nonzero holomorphic units. The inverse divides by these same units and
powers. The valuation homomorphism is injective. Its image is

\begin{equation}
 \Gamma=\langle(3,0),(4,1)\rangle
        =\ker\bigl(\rho:\mathbb Z^2\to\mathbb Z/3\bigr),
 \qquad \rho(x,y)=x-y\pmod3.
                                                               \hypertarget{tag-026C}{}\label{maps:eq:D.V4}
\end{equation}

The inverse coordinate map on the kernel is
\((x,y)\mapsto((x-4y)/3,y)\). The alternative period basis
\((3,0),(1,1)\) is obtained by the integral matrix
\(\left(\begin{smallmatrix}1&-1\\0&1\end{smallmatrix}\right)\),
whose inverse changes the upper-right entry to one. The unit-triangle fan is quotiented by this lattice.

That fan has rays \((v,1)\) and cones over
\((v,v+e_1,v+e_2)\) and
\((v+e_1+e_2,v+e_1,v+e_2)\). At a ray the quotient lattice maps
to \(\mathbb Z^2\) by \([(w,h)]\mapsto w-hv\), inverted by
\(w\mapsto[(w,0)]\). Its six cyclic rays are

\[
 (1,0),(0,1),(-1,1),(-1,0),(0,-1),(1,-1).
\]

The ray fan is the blow-up of the three-ray projective-plane fan at
all three torus fixed points: insert the sum of the adjacent ray
generators in each cone. The affine blow-up maps are
\((a,t)\mapsto(a,at)\) and \((t,b)\mapsto(tb,b)\), with inverse
ratios on their overlap. They identify the component with
\(dP_6=\mathbb P^2\#3\overline{\mathbb P}^2\), with its complex
orientation. The six boundary curves, in a suitable cyclic starting
place, are

\[
 \begin{gathered}
 E_1^{\rm exc},\ L_{\rm bl}-E_1^{\rm exc}-E_2^{\rm exc},\ E_2^{\rm exc},\\
 L_{\rm bl}-E_2^{\rm exc}-E_3^{\rm exc},\ E_3^{\rm exc},\ L_{\rm bl}-E_3^{\rm exc}-E_1^{\rm exc}.
 \end{gathered}
\]
Here \(L_{\rm bl}\) is the pullback of a projective line and
\(E_i^{\rm exc}\) are the exceptional classes.

The residues of the six neighbours are \((1,2,1,2,1,2)\). Each
triangle has the three different residues. A component therefore meets
none of its distinct period translates. The period action stabilizes no
finite nonempty cell. Its restriction to a component is an embedding;
the inverse on its image chooses its unique lift, using the period maps to compare
charts. Units act on the dense torus of a component by their nonzero
constant values and hence preserve this argument and all boundary
directions. The compact central fibre consists of three embedded
\(dP_6\) components \(D_0,D_1,D_2\), nine double curves and six
triple points. There are no component self-identifications.

\emph{The conductor and the cusp boundary.}

The toroidal neighbourhood retracts onto this central divisor. Here and
at the original additive neighbourhoods we use a regular neighbourhood
compatible with the analytic strata. On its barycentric regular-
neighbourhood model the collapse has the following map: multiply
the coordinates of vertices outside the central subcomplex by \(1-t\),
retain the coordinates on that subcomplex, and renormalize their sum to
one. Each simplex of the derived neighbourhood has a nonempty initial
face in that subcomplex. These maps agree on faces and at \(t=1\)
land in the central divisor; inclusion is the homotopy inverse. The
collar between this neighbourhood and the chosen small tube over a
closed base disc is identified by a horizontal lift on the proper
submersion over an annulus. Reversing that lift gives the inverse collar
map. The construction is made after the period quotient, where the
central strata are compact; its lifted charts retain the period maps.

The conductor chain groups in even fibre degree are

\[
 \begin{array}{c|ccc}
 &p=0&p=1&p=2\\\hline
 q=4&\mathbb Z^3&0&0\\
 q=2&\mathbb Z^{12}&\mathbb Z^9&0\\
 q=0&\mathbb Z^3&\mathbb Z^9&\mathbb Z^6.
 \end{array}
\]

The maps are alternating inclusions of the strata. Take vertex representatives \((i,0)\),
\(i=0,1,2\), and orient the three edge directions
\(e_1,e_2,e_2-e_1\). An edge from colour \(i\) in direction \(d\)
has target colour \(\rho((i,0)+d)\). Its degree-zero column is target
minus source. Its degree-two column is \(-[C_d]_i+[C_{-d}]_{\rho((i,0)+d)}\), where \(C_d\) is the boundary curve corresponding to the ray \(d\) in the six-curve list. Triangles use
their counterclockwise boundary. Elementary integer row and column operations on these three
matrices give respectively two, five, and seven nonzero Smith
factors, all equal to one. Each operation is inverted by the
opposite row or column operation. In particular every quotient in the conductor complex is torsion free.

The degree-zero complex is the cellular chain complex of
\(\mathbb R^2/\Gamma\); the covering map and its period lifts give
its two long generators. The other two rows have respectively a
rank-five cokernel and a rank-two kernel, and rank three in degree
four. The filtration has no nonzero higher differential: its nonzero
rows have even \(q\), and its columns are only \(0,1,2\). All remaining
extensions split over free abelian groups. Thus for the cusp tube \(U_c\)

\[
 \pi_1(U_c)=\mathbb Z^2,\qquad
 (H_0,H_1,H_2,H_3,H_4)(U_c)
 = (\mathbb Z,\mathbb Z^2,\mathbb Z^6,\mathbb Z^2,\mathbb Z^3).
\]

For the fundamental group, the simply connected component and double-
curve spaces have the same van Kampen groupoid as the dual complex.
The chart lifts supply the inverse from the two long-period loops.

Let \(M_c=\partial U_c\). The Smith factors of \(T_c-I\) are
\((1,3)\), and those of \(\wedge^2T_c-I\) are \((1,1)\).
The Wang sequence, split by the marked origin section, gives

\[
 H_1(M_c)=\mathbb Z\langle r,\gamma,x_c\rangle
                    \oplus\mathbb Z/3\langle\lambda\rangle,
 \quad \delta=-4\lambda,\qquad H_2(M_c)=\mathbb Z^6.
\]

The cusp kills the short loops \(\lambda,\delta\) and the reduced
normal meridian, and retains the long loops. These maps are induced by
the period quotient of the fan. A disc for \(\delta\) is,
in the chart \((x,y,w)=(q/(X_1X_2),X_2,X_1)\),
\(z\mapsto(\epsilon\overline z,\epsilon z,c_{\rm disc})\), with fixed nonzero
\(c_{\rm disc}\). Its boundary has fixed \(q=c_{\rm disc}\epsilon^2\), fixed \(X_1=c_{\rm disc}\),
and one positive turn in \(X_2\). Swapping the two angular coordinates
gives the short \(\lambda\) disc. Holding the two nonzero chart
coordinates fixed gives each positive normal disc.

In the component basis of \(H_2(U_c,M_c)=H^4(U_c)=\mathbb Z^3\),
the three positive normal discs have boundaries \(x_c+i\lambda\).
The boundary map to \(\ker(H_1(M_c)\to H_1(U_c))\) is therefore

\[
 (h_0,h_1,h_2)\longmapsto
       (h_0+h_1+h_2,\ h_1+2h_2\pmod3).
\]

The section-frame construction below identifies the colour-zero
normal disc with the stated based meridian. The kernel of this map has
basis and inverse

\[
 v_0=(-1,2,-1)^t,\quad v_1=(-1,-1,2)^t,
 \qquad h\longmapsto ((h_1-h_0)/3,(h_2-h_0)/3).
\]

Indeed the two congruences in the normal boundary map make the two quotients integral;
substitution in the reverse map \((a,b)\mapsto av_0+bv_1\) returns
each coordinate of \(h\).

The angular circle \(\delta\) fixes precisely the three complete
curves on edges in direction \(e_2\). Write \(F_i\) for the edge
starting at colour \(i\); its other endpoint has colour \(i+2\).
The inverse ray matrices give weights \((-1,0,1)\) and
\((1,0,-1)\) in the two adjacent cone orderings. Each triangle contains
one fixed edge, so the three complete fixed curves are disjoint. Their
two normal charts and inverses are exactly

\[
\begin{gathered}{}
 (x,y,w)=(q/(X_1X_2),X_2,X_1),\quad
\\{}
 (x',y',w')=(q/X_2,X_1X_2,1/X_1),
\\{}
 \quad (x',y',w')=(wx,wy,w^{-1}).
\end{gathered}
\]

The inverse transition replaces \(w\) by \(w'\) and has the same
form. Thus each normal bundle is
\(\mathcal O(-1)_{-1}\oplus\mathcal O(-1)_{+1}\). Its intersections
with the three component divisors are respectively the columns

\[
 {\mathsf V}_{\rm int}=\begin{pmatrix}-1&-1&2\\2&-1&-1\\-1&2&-1\end{pmatrix}.
\]

The two endpoint components have normal degree \(-1\), and the
third meets the curve once at each of its two transverse endpoints.
This also follows by multiplying the six-curve intersection matrix
derived from the boundary-curve classes. In particular \(F_0,F_1\) map to the normal-kernel basis.

The two normal-disc sweeps by the short \(\lambda,\delta\) loops
give elements of \(H_3(U_c,M_c)\). Their boundary tori project to
the full primitive fixed lattice in the boundary Wang sequence;
the inverse sends a fixed loop to its parametrized disc sweep.
They are independent. By the cusp homology groups and duality the relative group has
rank two, and the boundary kernel has rank two. Primitivity follows
also directly from that inverse projection. Thus they give the full
kernel of \(H_2(M_c)\to H_2(U_c)\). The fibre quotient
\(\operatorname{coker}(\wedge^2T_c-I)\) maps isomorphically to its
image. Exactness and the normal boundary map give the integral splitting

\[
 H_2(U_c)=\operatorname{im}H_2(M_c)
                   \oplus\mathbb Z\langle F_0,F_1\rangle.
\]
In particular the quotient and its inverse are
\begin{equation}
 H_2(U_c)/\operatorname{im}H_2(M_c)\xrightarrow{\ \cong\ }\mathbb Z^2,
 \qquad [rF_0+sF_1]\longmapsto(r,s),
 \qquad(r,s)\longmapsto[rF_0+sF_1].
\hypertarget{tag-026D}{}\label{maps:eq:D.V5}
\end{equation}

For a class \(z\), compute its normal coordinates \(h\), subtract
the combination of \(F_0,F_1\) with the displayed normal-kernel coordinates, and then lift the remainder
through the fibre quotient. This is the inverse splitting. The boundary description includes its order-three torsion group.

\emph{The original additive fillings.}

The original II surface neighbourhood \(N_0\) retracts onto its
cuspidal rational curve. The normalization parameter has local map
\(u\mapsto(u^2,u^3)\) and inverse \(y/x\) off the cusp, extended
by zero at the cusp; it is a homeomorphism of that curve with
\(\mathbb P^1\). The regular-neighbourhood collapse therefore gives \(\pi_1(N_0)=0\), \(H_2(N_0)=\mathbb Z\), and zero homology in every other positive degree.

The II component group is trivial. The identity-component logarithm
of the original section gives translations
\(\tau_s=t_{\operatorname{Exp}(sR_{\rm log})}\), with inverse
\(t_{\operatorname{Exp}(-sR_{\rm log})}\), on the original smooth surface.
The logarithm follows from \cite[\S2.4, equation (2.2), p.~6]{manUU21}, and the translation extension from \cite[\S7.6, pp.~33--34]{manSS10}. The matrix
\(S_0\) is integral, so its marked angular coordinate is
\(x_\delta-12x_1-12x_2\). The narrow line has degree zero on the
central curve and is topologically trivial on \(N_0\).

Use the nonzero line frame and write the contraction as
\((y,w)\mapsto(\tau_1y,u_{\rm II}(y)w)\), \(|u_{\rm II}(y)|<1\). Set \(u=-\log|w|\),
\(r_{\rm roof}(y)=-\log|u_{\rm II}(\tau_1^{-1}y)|\). Choose a positive \(k_{\rm roof}\) with
\(k_{\rm roof}\max r_{\rm roof}<1\), and a middle-supported bump \(b_y\) of integral
one. The increasing map

\[
 u\longmapsto f_y(u)=\int_0^u[k_{\rm roof}+(1-k_{\rm roof}r_{\rm roof}(y))b_y(t)]\,dt
\]

has a smooth inverse, agrees with \(k_{\rm roof}u\) and
\(1+k_{\rm roof}(u-r_{\rm roof}(y))\) on its end collars, and takes the variable-roof
quotient to its mapping torus. The maps
\([y,t]\mapsto[t_{\operatorname{Exp}(-\zeta(t)R_{\rm log})}y,t]\) and the
same map with the opposite sign identify that torus with
\(N_0\times S^1_\gamma\). Here \(\zeta\) is zero and one on the
two end collars. Their endpoint identity is precisely the mapping-
torus gluing identity.

The angular circle bundle is trivial: its restriction to \(N_0\) is
trivial, and the degree-two Wang restriction is injective because
\(H^1(N_0)=0\). A concrete trivialization uses transition logarithms
and a partition of unity after removing their integral \v{C}ech cocycle;
the inverse gauges are the reciprocal exponentials. Consequently

\[
 U_{\rm II}\cong N_0\times S^1_\gamma\times S^1_\delta,
 \quad H_2(U_{\rm II})=\mathbb Z\langle V_0,D_{\rm II}\rangle,
 \quad D_{\rm II}=\delta\times\gamma.
\]

The fibre specialization, in order \(12,13,14,23,24,34\), is

\[
 x\longmapsto l_0(x)V_0+d_0(x)D_{\rm II},
 \quad l_0=(1,0,-1,0,-1,0),\quad d_0=(0,0,-12,0,-12,1).
\]

The first row is the elliptic area after \(S_0\); its scale is fixed
by the single reduced rational fibre. The second is the product of
the two surviving circle coordinates. Their section sends the two
coordinate generators to \(e_{12}\) and \(e_{34}\).
Their common kernel equals \(\operatorname{im}(\wedge^2T_0-I)\):
the latter has four unit Smith factors and these two rows give a
primitive quotient of rank two. The two fixed loops
\(\delta\) and \((-1,1,0,1)\) extend across the original tube,
so their normal sweeps give the entire kernel in boundary degree
two. The boundary degree-one map is onto, with the extended original
section meridian as its kernel. These statements also follow directly
from the variable-roof and type-II product maps and their inverse maps.

For the star, use the mapping-torus construction over its
original neighbourhood. Let \(N_*\) be its minimal elliptic
surface over the closed base disc, let \(M_*^{\rm lin}\) be its
narrow line, and let \(\phi_*\) be its contracting
linearization.  The basis change is
\(C=\left(\begin{smallmatrix}1&1\\-1&0\end{smallmatrix}\right)\),
\(C^{-1}=\left(\begin{smallmatrix}0&-1\\1&1\end{smallmatrix}\right)\),
and \(P_* =\operatorname{diag}(C,1,1)\). Direct multiplication gives

\[
 P_*^{-1}T_*P_*
 =\begin{pmatrix}-1&-1&0&-1\\0&-1&0&1\\-12&-24&1&3\\0&0&0&1\end{pmatrix},
 \qquad (-9,-3)C=(-6,-9).
\]

This changes coordinates on the marked boundary:
\(u_{\rm old}=Cu_{\rm new}\), inverted by \(C^{-1}\),
with the line angle, suspension angle and base point fixed. The following maps use \(N_*,M_*^{\rm lin},\phi_*\) with this marking.

The original \(I_1^*\) divisor has four multiplicity-one outer
spheres \(L_0,L_1,R_0,R_1\) and two multiplicity-two central
spheres \(B_0,B_1\). The regular-neighbourhood collapse gives \(\pi_1(N_*)=0\), \(H_2(N_*)=\mathbb Z^6\) on
these oriented component classes, and \(H_1(N_*)=H_3(N_*)=0\).
The component homomorphism in the old coordinates is

\[
 \mathbb Z^2/(A_*-I)\mathbb Z^2\longrightarrow\mathbb Z/4,
 \qquad [x_1,x_2]\longmapsto[x_1-x_2],
 \qquad [k]\longmapsto[(k,0)] .
\]

These maps are inverse: the row sends the two relation columns to
\(-4,0\), is onto, and \(\det(A_*-I)=4\). The section
cocycle \(g_*=(0,1)\) has residue three. Since the cocycle is
minus the Abel--Jacobi increment, the invariant point
\(z_*=(1/4,-3/4)\) represents the forward section \(-G\)
of residue one. Its coordinates after \(C^{-1}\) are
\((3/4,-1/2)\). Swapping these two entries, as required by the
convention \(e_1=1,e_2=\tau\), gives Ulmer--Urzua's invariant
leaf \((r,s)=(1/2,3/4)\) modulo \(\mathbb Z^2\).
By \cite[Proposition 3.5(1),(3), pp.~8--10]{manUU21}, this leaf extends to an order-four torsion section \(q_{\rm tor}\) of the surface. It has the same
component as \(-G\), including at the central fibre.

The section \(-G-q_{\rm tor}\) therefore lies in the identity component.
The local exact sequence of \cite[\S2.4, equation (2.2), p.~6]{manUU21}, after shrinking this base disc, supplies its Lie logarithm \(R_*\), with
\(\operatorname{Exp}(R_*)=-G-q_{\rm tor}\). Define

\[
 \tau_s=t_{q_{\rm tor}+\operatorname{Exp}(sR_*)},\qquad
 \tau_s^{-1}=t_{-q_{\rm tor}+\operatorname{Exp}(-sR_*)},
 \quad 0\leq s\leq1.
\]

By \cite[\S\S7.6--7.7, pp.~33--34]{manSS10}, these translations extend over the entire smooth \(N_*\); every map and its displayed
inverse is over the same base disc. Thus \(\tau_1=t_{-G}\)
is the base map of \(\phi_*\), and \(\tau_0=t_{q_{\rm tor}}\).

The degree of \(M_*^{\rm lin}\) on every one of these six
components is zero. Since they are the full integral basis of
\(H_2(N_*)\), the line has zero first Chern class and admits a
nonzero smooth frame. In that very frame write
\(\phi_*(y,w)=(\tau_1y,a_*(y)w)\), using the nowhere-zero multiplier \(a_*\) and a contracting Hermitian
norm. The argument of \(a_*\) enters the angular
circle bundle. The quotient by the line circle, before trivializing
that bundle, is exactly

\[
 \{(y,u):0\leq u\leq r_*(y)\}/
       (y,r_*(y))\sim(\tau_1^{-1}y,0),\qquad
 r_*(y)=-\log|a_*(\tau_1^{-1}y)| .
\]

Use the variable-roof map with this function \(r_*\) and its smoothly varying
middle-supported bump. Its lower and upper collar formulae give
\([y,u]\mapsto[y,f_y(u)]\), with inverse \(f_y^{-1}\),
to \(T(\tau_1^{-1}:N_*\to N_*)\). The further maps

\[
 [y,t]\longmapsto
 [t_{\operatorname{Exp}(-\zeta(t)R_*)}y,t],\qquad
 [y,t]\longmapsto
 [t_{\operatorname{Exp}(\zeta(t)R_*)}y,t]
\]

are inverse maps to and from \(T(\tau_0^{-1})\).
They agree on the glued collars because
\(\tau_0^{-1}t_{\operatorname{Exp}(-R_*)}=\tau_1^{-1}\).
All their maps fix the base coordinate, so they also identify the
boundary collars over the boundary circle of the base disc.

The angular circle bundle restricts to the trivial line on
\(N_*\). The degree-two Wang restriction to \(N_*\) is
injective, because \(H^1(N_*)=0\). Its Euler class is therefore
zero on the entire mapping torus. On a good cover take the logarithms
of its transition functions, including the phase of
\(a_*\). Remove their integral \v{C}ech cocycle by an integral
coboundary, then use a partition of unity to write the remaining real
cocycle as differences of functions. Multiplication by their
exponentials is a smooth angular trivialization; reciprocal
exponentials give its inverse. This constructs the equivariant
diffeomorphism and inverse

\[
 U_*\longleftrightarrow
 S^1_\delta\times T(t_{-q_{\rm tor}}:N_*\to N_*) .
\]

The section-frame construction below fixes the origin disc. The line
coordinate there is the Poincar\'e change
\(x'_\delta=x_\delta-9x_1-3x_2\), with inverse addition.
An integral adjustment in the positive suspension direction is
\(\exp(2\pi i k_{\rm fr}t)\), inverted by its reciprocal; choose its
integer \(k_{\rm fr}\) to match that frame. This fixes the boundary marking and its collar in the
mapping-torus identification.

Residue one makes \(\tau_0\) act on the outer components as
\(L_0\mapsto R_0\mapsto L_1\mapsto R_1\mapsto L_0\).
Preservation of the incidence tree then swaps \(B_0,B_1\).
All component orientations are preserved. The monodromy of the star mapping-torus model
is the inverse permutation. If the forward permutation is \(P_{\rm perm}\),
the two Wang complexes are identified by identity on their targets
and \(-P_{\rm perm}\) on their sources, since
\((P_{\rm perm}^{-1}-I)(-P_{\rm perm})=P_{\rm perm}-I\),
with source inverse \(-P_{\rm perm}^{-1}\).
The coinvariant map adds the coefficients in each of the two orbits;
its section assigns each coefficient to the first component of its
orbit. The invariant map repeats the two orbit coefficients, with
inverse reading one representative. Call the coinvariant generators
\(U_{*,0},U_{*,1}\), put \(W=U_{*,0}+U_{*,1}\), and write
\(D_*=\delta\times\gamma\). Product and Wang give the
complete maps

\[
 H_2(U_*)=\mathbb Z\langle U_{*,0},U_{*,1},D_*\rangle,
 \quad H_1(U_*)=\pi_1(U_*)=\mathbb Z\langle\delta,\gamma\rangle,
\]
\[
 x\longmapsto (-9x_1-3x_2+x_\delta,x_\gamma),\qquad
 x^{(2)}\longmapsto l_*(x^{(2)})W+d_*(x^{(2)})D_*,
\]
\[
 l_*=(4,0,-3,0,-1,0),\qquad d_*=(0,0,-9,0,-3,1).
\]

The first map has section \((d,g)\mapsto d\delta+g\gamma\).
The second has section \((1,0)\mapsto-e_{24}-3e_{34}\),
\((0,1)\mapsto e_{34}\). The fibre specializes to \(4W\):
its four outer multiplicities are one and its two central
multiplicities are two. The row \(l_*\) is obtained by the
elliptic coordinate map
\((x_{\rm ell},x_\gamma)\mapsto
(x_{\rm ell}+z_*x_\gamma,x_\gamma)\), inverted by subtraction,
as \(4(x_{12}+z_{*,2}x_{14}-z_{*,1}x_{24})\).
The row \(d_*\) is the wedge of the two surviving degree-one
coordinates. Their displayed integral section and the four unit
Smith entries of \(\wedge^2T_*-I\) identify their common
kernel with that entire relation image.
The two extending fixed loops are
\(\delta\) and \((-1,3,0,4)\). The former is the line circle; the latter is the fourth forward orbit over the extended
torsion section, with elliptic coordinate \(-4z_*=(-1,3)\)
and zero angular coordinate in the star mapping-torus model. Both extend over the base
disc, so their products with the positive base meridian bound the
corresponding swept discs. Solving \((T_*-I)x=0\) gives
exactly their primitive rank-two span. They are therefore the full
boundary degree-two kernel after the split Wang sequence. The
linearization is a nonzero unit on this original cap, so its local
line winding adds zero to those discs. The original section meridian
generates the degree-one kernel. The quotient of \(H_2(U_*)\) by the boundary image
has free generator \(U=U_{*,1}\), with quotient coordinate
\((u_0,u_1,d)\mapsto u_1-u_0\) and section \(1\mapsto U_{*,1}\).

\emph{The section frame and the global integral map.}

The generator has height \(1/12\) in the realization of \ref{mapsD:repair:collision-realization}.
Label the original \(I_3\) components by Tate valuation modulo
three. Write \(\alpha(q)q^3\) for the elliptic long period and
\(\mu(q)q^4\) for the forward \(\gamma\) multiplier, where
\(\alpha,\mu\) are nonzero units. The latter covers \(-G\),
so its inverse represents \(G\). Multiplication by the square of
the long period gives the same point with coordinate
\(\alpha^2\mu^{-1}q^2\). Thus the component residues of
\(G\) are respectively \(0,2,3\) in the II, I3 and I1* groups. Thus \(12G\) is narrow. The line \(M=\mathcal O(P-O)\) has degree zero on every
component. Its translation difference is fibrewise trivial, hence
vertical. The kernel of each component intersection matrix is generated
by the whole fibre, so its zero multidegrees make it a pullback from
the base. Its degree is the mixed height \(\langle12G,G\rangle=1\).
This gives the degree-one linearization, and \(\deg(M|_O)=6\).

Choose the section of \(M|_O\) with its only zero, of order six,
at the whole I3 fibre. On that fibre write the contracting map as
\(\phi=q\mathcal U\), with \(\mathcal U\) a nowhere-zero line isomorphism. Its six
inverse iterates take \(q^6v\) on \(O\) to
\(\mathcal U_5^{-1}\cdots\mathcal U_0^{-1}v\) on \(6G\); six forward iterates
are the inverse comparison. Every section meets the smooth locus of
the reduced elliptic fibre. The residue of \(6G\) is zero modulo
three. In the cusp angular coordinates the original section has
valuation \((0,6)\), whose colour is zero. Its boundary winding
adds \(6\delta=-24\lambda=0\) in the boundary description. This proves that the
extended section disc has the based meridian used in the normal boundary map.
At the two weight-zero original caps the frame extends as a nonzero
unit. At the additional regular place it is the unshifted frame.
These are the inverse-iteration section maps of \ref{maps:src:D.2.12}, here with six iterates and residue zero.

Start the based boundary order with the additional regular place,
then II, the cusp, and the star. The regular exterior \(V\) retracts
to the flat marked torus bundle over a three-loop graph. Its graph
map and based boundary inclusions are

\[
 \mathsf D_{\rm gr}=[0\ T_0-I\ T_c-I],\quad j_i(v)=\text{block }i\text{ inclusion},
 \quad j_3(v)=-(T_0T_cv,T_cv,v),\quad \mathsf D_{\rm gr}j_3=I-T_*^{-1}.
\]

The inverse period map is the real logarithm modulo the period
lattice; reversing the boundary walk gives the displayed transports.
Let \(K\) have columns

\[
 j_1\delta,\ j_1(-1,1,0,1),\ j_2\lambda,\ j_2\delta,
 j_3\delta,\ j_3(-1,3,0,4),\ j_0e_1,\ j_0e_2,\ j_0\gamma.
\]

Define the \(12\times3\) matrix \({\mathsf C}_{0,\rm gr}\) by its nonzero rows,
numbered from one:

\[
 ({\mathsf C}_{0,\rm gr})_6=(1,2,0),\quad({\mathsf C}_{0,\rm gr})_9=(-1,19,4),\quad
 ({\mathsf C}_{0,\rm gr})_{12}=(0,16,3).
\]

Multiplication gives
\(\mathsf D_{\rm gr}{\mathsf C}_{0,\rm gr}=(\operatorname{diag}(1,1,3),0)^t\), \(\mathsf D_{\rm gr}K=0\), and
\(\det[{\mathsf C}_{0,\rm gr}\ K]=-1\). Hence its inverse is its negative adjugate.
Its last nine rows \({\mathsf L}\) satisfy \({\mathsf L}K=I\),
and on the graph kernel they give the complete inverse coordinates.
In particular

\[
 {\mathsf L}j_0\delta=(-1,0,0,-1,-1,0,0,0,0)^t.
\]

The fibre degree-one quotient is
\(\mathbb Z\gamma\oplus\mathbb Z/3\delta\). The cusp boundary
torsion generator \(\lambda\) maps to \(-16\delta=2\delta\),
an isomorphism of the order-three groups.

For the fibre degree-two map put
\(B_2=[0\ \wedge^2T_0-I\ \wedge^2T_c-I]\),
\(s=-e_{34}\), and
\({\mathsf F}_{\rm bas}=[e_{12}+12e_{34},e_{13},e_{14},e_{23},e_{24}]\).
The six columns \([s,{\mathsf F}_{\rm bas}]\) are unimodular, with inverse first row
\(\omega=(12,0,0,0,0,-1)\). Define a right-preimage matrix
\(P_{\rm cyl}\) by its five columns
\((0,\alpha_i,\beta_i)\) in the three wedge blocks, where
\[
\begin{array}{c|l|l}
i&\alpha_i&\beta_i\\\hline
1&-12e_{13}-3e_{14}-9e_{23}&e_{14}\\
2&-3e_{13}+e_{23}&0\\
3&-12e_{13}-3e_{23}+e_{24}&-e_{14}\\
4&-7e_{13}+2e_{23}&0\\
5&24e_{13}-4e_{14}-39e_{23}+2e_{24}&-e_{14}.
\end{array}
\]
Multiplication gives \(B_2P_{\rm cyl}={\mathsf F}_{\rm bas}\). Together with
\(\omega B_2=0\), this proves

\[
 \operatorname{im}B_2=\ker\omega,\qquad
 [x]\mapsto\omega x,\quad k\mapsto[ks]
\]

as inverse integral quotient maps.
The marked section and the real-period triangles with vertices
\(0,x,x+y\) split the exterior degree-two group as
\(\mathbb Zs\oplus\ker\mathsf D_{\rm gr}\). Subtracting these triangles is the
inverse chain map. The invariant vertical two-form \(\omega\)
vanishes on all fixed-loop normal cylinders, which fixes the splitting.

The six original-cap sweeps kill the first six columns of the graph-kernel basis.
The II fibre map has \(12l_0-d_0=\omega\), so it identifies
\(V_0=12s,D_{\rm II}=-s\). The star map has \(3l_*-d_*=\omega\),
so it identifies \(W=3s,D_*=-s\), leaving \(U=U_{*,1}\).
The cusp splitting leaves exactly \(F_0,F_1\) beside its full
fibre image. Thus the remaining degree-two coordinates before the log
are

\[
 U,F_0,F_1,s,t_1,t_2,t_\gamma.
\]

The degree-one eliminations in this assertion include
the order-three group: use the cusp torsion column \(\lambda\) to
remove the exterior torsion column \(2\delta\). The inverse uses
the inverse automorphism, also multiplication by two modulo three.
Use the local sections in the local boundary computations to remove the long-period target
coordinates; each is an identity block. The three old section
meridians then remove the three base loops, with their single sphere
boundary product relation. The leftover degree-one group is
\(\mathbb Z\gamma\). These inverse additions and the normal-kernel splitting prove that the displayed pre-log coordinates form an integral basis.

The same conclusion for the fundamental group follows by van Kampen.
The star kills \(e_1+9\delta,e_2+3\delta\); the cusp kills \(\delta\).
Their matrix in \((e_1,e_2,\delta)\) is
\(\left(\begin{smallmatrix}1&0&0\\0&1&0\\9&3&1\end{smallmatrix}\right)\),
with inverse obtained by negating nine and three. Hence all of
\(\ker\chi\) is killed. The old based meridians bound by the section-frame construction.
Each cap group is generated by its boundary image: a loop is moved
off the central codimension-two analytic strata and then radially
to the boundary along a horizontal lift of the punctured base.
Reverse radial transport is its inverse. The boundary images generate
each cap group, so the quotient and inverse are
\(\chi\) and \(1\mapsto\gamma\).

\emph{The regular logarithmic filling.}

Write \(v=ae_1+b'e_2+c\delta+n\gamma\). The stated gcd makes
the order of \(v/m\) on the central torus exactly \(m\), so the
root-disc action \((y,z)\mapsto(y+v/m,e^{2\pi i/m}z)\) is free.
The holomorphic collar of \ref{maps:src:D.2.16} has the following isotopy and inverse in real-period coordinates:

\[
 h_\tau(y,z)=y-\frac{\arg z}{2\pi}v+\tau B(z),\qquad
 h_\tau^{-1}(y,z)=y+\frac{\arg z}{2\pi}v-\tau B(z),
\]

where \(B=(\log|z|/2\pi)J_zv-u_{\rm sec}(z)\), and
\(u_{\rm sec}\) is the real vector of the logarithm of the nonzero section unit. The remaining translation is deck invariant;
a radial cutoff extends it to the old exterior, and its negative
translation is the inverse. The normal slope is \(mx_r-v\).

The central lattice has presentation
\(\Lambda_v=\mathbb Z^5/\mathbb Zq_0\),
\(q_0=(-a,-b',-c,-n,m)^t\). The quotient map is
\([(z,k)]\mapsto z+kv/m\). A Bezout row gives a unimodular
matrix \({\mathsf A}_{\rm Bez}\) with \({\mathsf A}_{\rm Bez}q_0=e_1\). Its last four rows \(H_v\),
the last four columns \(E_v\) of \({\mathsf A}_{\rm Bez}^{-1}\), and its first row
\(r_v\) satisfy
\(H_vE_v=I\), \(H_vq_0=0\),
\(E_vH_v+q_0r_v=I\). They give the inverse quotient coordinates.
The smooth normal-line trivialization is
\([(y,z)]\mapsto([y],e^{-2\pi i\ell(y)}z)\), where
\(\ell(v)+m\alpha_{\rm Bez}=1\); its inverse multiplies by the opposite
exponential. The deck multiplier is one. They identify the normal line and its Thom coordinates.

The boundary of a negatively oriented loop-times-normal-disc is
\(-\omega(\lambda,v)s+mj_0\lambda\), and the fifth generator
has boundary \(j_0v\). This follows from the triangle-prism map
of the logarithmic collar, whose fibre triangle has signed area
\(-\omega(\lambda,v)\). On the last four coordinates of the pre-log basis
the full transfer is

\[
 R=\begin{pmatrix}
 -12b'&12a&n&-c&0\\
 m&0&0&0&a\\0&m&0&0&b'\\0&0&0&m&n
 \end{pmatrix},\qquad Rq_0=0.
\]

Its \(U,F_0,F_1\) rows are zero: every collar triangle and cylinder
is in the regular exterior, and the pre-log basis uses the component and
fixed-curve sections of the local quotients. Before the six old sweep
eliminations, adjoin six identity columns and the rows
\(m{\mathsf L}j_0\lambda\).
Subtract their first six graph entries using those identity columns;
the inverse adds them. The resulting log relation matrix fixes the geometric free generators.

This is the \(p=12\) relation matrix of \ref{maps:src:D.2.16}. It has
\(d_1=\gcd(m,n,a,b')\), \(d_2=\gcd(m,n,12,c)\),
\(t=d_1d_2\). A common prime divisor of \(d_1,d_2\) would violate
the freeness gcd, so they are coprime. The primewise pivot calculation in \ref{maps:src:D.2.16} applies to this matrix: a unit \(a\) or \(b'\) leaves
the ideals \((m,n,12,c)\), \(m(m,n,12,c)\), \((m^2n)\);
a unit \(c\) leaves \((m,n,a,b')\), \(m(m,n,a,b')\),
\((m^2n)\). The cases where \(m\) or \(n\) is a unit supply
the remaining primes. Thus the invariant factors are

\[
 (1,t,m,m|n|/t)\quad(n\ne0),\qquad (1,t,m,0)\quad(n=0).
\]

The signed maximal-minor vector is \(m^2nq_0\).
For any allowed integral parameter tuple, integer Smith operations give
\(U_RRV_R=D_R\), with \(U_R,V_R\) unimodular. Applying
\(U_R\) to a target vector and reducing modulo the diagonal entries
is the quotient map; applying \(U_R^{-1}\) to coordinate
representatives is its inverse into the cokernel. The source coordinates
use \(V_R,V_R^{-1}\) and \(H_v,E_v\) on the log lattice.
Thus the diagonal calculation gives the integral maps as well as the
abstract factors.

In degree one the four boundary free ranks are \((5,3,3,3)\),
the four cap ranks are \((4,2,2,2)\), and the exterior free rank
is four. Both source and target additionally have the same order-
three group. Eliminate that group by the cusp torsion isomorphism
in the section-frame construction. Eliminate the ten local section columns with their identity
blocks. In target coordinates \((\gamma,x_r,x_0,x_c)\), the four
remaining normal columns are

\[
 G_4=\begin{pmatrix}
 -n&0&0&0\\m&0&0&-1\\0&1&0&-1\\0&0&1&-1
 \end{pmatrix}.
\]

They are the regular normal, II normal, cusp normal and star normal.
Their determinant is \(n\); their kernel is zero for \(n\ne0\)
and is generated by \((1,m,m,m)\) for \(n=0\). The preceding
identity and torsion block eliminations have explicit inverses, so
this is also the kernel of the full degree-one map. Van Kampen adds
the relation \(n\gamma=0\), and its inverse maps give
\(\pi_1(X)=H_1(X)=\mathbb Z/n\mathbb Z\), with
\(\mathbb Z/0=\mathbb Z\).

At \(n=0\) the extra relative two-generator has an integral closed
lift. Take its four normal discs in coefficients \((1,m,m,m)\),
with their based annular tracks. Their boundary is \(-v\).
The explicit identity

\[
 v=a(e_1+9\delta)+b'(e_2+3\delta)
                 +(c-9a-3b')\delta
\]

fills it using the two original-star kernel discs and the cusp
\(\delta\) disc already constructed. Add these signed chains.
This closes the relative generator with coefficient one and splits
the new free summand. The other free coordinate comes from the
zero Smith factor. Its primitive functional is
\((ms+12b't_1-12at_2+ct_\gamma)/d_2\); a Bezout preimage is
its inverse section. The remaining coordinates are the Smith torsion
coordinates and the three geometric coordinates \(U,F_0,F_1\).

It follows that

\[
 H_2(X)=\mathbb Z^3\oplus
 (\mathbb Z/t\oplus\mathbb Z/m\oplus\mathbb Z/(m|n|/t))
 \quad(n\ne0),
\]
\[
 H_2(X)=\mathbb Z^5\oplus(\mathbb Z/t\oplus\mathbb Z/m)
 \quad(n=0).
\]

The cusp Euler number is \(3\cdot6-9\cdot2+6=6\). Each
other central fibre has a free circle factor or a torus quotient,
and the regular parts and boundaries also have Euler number zero.
Hence \(\chi(X)=6\). Integral duality and universal coefficients
give, for \(n\ne0\),

\[
 H_5=0,\quad H_4=\mathbb Z^3\oplus\mathbb Z/|n|,
 \quad H_3=\mathbb Z^2\oplus
       (\mathbb Z/t\oplus\mathbb Z/m\oplus\mathbb Z/(m|n|/t)),
\]

and, for \(n=0\),

\[
 H_5=\mathbb Z,\quad H_4=\mathbb Z^5,
 \quad H_3=\mathbb Z^4\oplus(\mathbb Z/t\oplus\mathbb Z/m).
\]

In both cases the point and complex fundamental class generate
\(H_0=H_6=\mathbb Z\). To specify the duality identifications,
dual integral cocycles to the free coordinates in the displayed second-homology groups, capped
with \([X]\), give the free \(H_4\) coordinates. The character of
a cyclic torsion generator taking value \(1/d\) in
\(\mathbb Q/\mathbb Z\), followed by the integral Bockstein and
cap product, gives the corresponding dual torsion generator;
linking evaluation is its inverse. This treats \(H_1\) torsion in
\(H_4\) and \(H_2\) torsion in \(H_3\). The free middle ranks
are fixed by \(2-2b_1+2b_2-b_3=6\). Integer cycle reduction and
symplectic Gram--Schmidt on the unimodular alternating free middle
pairing give inverse free coordinates, as in \ref{maps:src:D.2.16}. For \(n=0\), \(\chi\frown[X]\) generates
\(H_5\), with evaluation as inverse. These are \eqref{maps:eq:D.V1} and \eqref{maps:eq:D.V2}.

\emph{The divisor cubic.}

Each \(D_i\) is an embedded effective Cartier divisor of the compact threefold, since the other fillings and the log are in
disjoint base discs. Set \(\eta_i=c_1\mathcal O_X(D_i)\).
The local divisor identity
\(\operatorname{div}(q)=D_0+D_1+D_2\) restricts to the line map

\[
 N_{D_i/X}\otimes\mathcal O_{D_i}(C_0+\cdots+C_5)
       \longrightarrow\mathcal O_{D_i}.
\]

In \(q=u x_0x_1x_2\), multiply the Cartier frames by \(q\);
the remaining multiplier is the nonzero unit \(u\). Its reciprocal
is the inverse, and the maps glue because they use the one function
\(q\). Thus \(N_{D_i/X}=\mathcal O(-\sum C_j)=K_{D_i}\),
with the period units in the multiplier. From the boundary-curve classes, every \(C_j^2=-1\),
adjacent curves have product one and the other products vanish.
Consequently

\[
 \langle\eta_i^3,[X]\rangle
 =\langle c_1(N_{D_i/X})^2,[D_i]\rangle=-6+12=6.
\]

The equality uses the complex Thom homomorphism
\(j_{i!}\): \(j_{i!}(1)=\eta_i\),
\(j_i^*\eta_i=c_1(N_{D_i/X})\), followed by cap-product
naturality. These are the restriction and pushforward maps.
For component coefficients \(d_i^{\rm div}\), the complete cubic is

\[
 \begin{aligned}
 \Big\langle\Big(\sum d_i^{\rm div}\eta_i\Big)^3,[X]\Big\rangle
 &=3(2d_0^{\rm div}-d_1^{\rm div}-d_2^{\rm div})\\
 &\quad (2d_1^{\rm div}-d_0^{\rm div}-d_2^{\rm div})(2d_2^{\rm div}-d_0^{\rm div}-d_1^{\rm div}).
 \end{aligned}
\]

It is obtained by restricting to each \(D_i\), where the six
coefficients alternate between \(d_{i+1}^{\rm div}-d_i^{\rm div}\) and
\(d_{i+2}^{\rm div}-d_i^{\rm div}\), and using the boundary-curve classes. In particular
\(D_i^2D_j=-3\) for distinct indices and \(D_0D_1D_2=6\).

\emph{The circle quotient.}

For the rest of the proof take \(n=\pm1\). Then \(t=1\) and

\[
 H_2(X)=\mathbb Z\langle U,F_0,F_1\rangle\oplus(\mathbb Z/m)^2,
 \qquad H_3(X)=\mathbb Z^2\oplus(\mathbb Z/m)^2.
\]

The \(\delta\) circle is free off the cusp fixed curves. At an
original cap this follows from the contraction norm: a nonzero
iterate cannot preserve a line norm. At the logarithmic central
torus the integer functional
\(x_\delta-(c/n)x_\gamma\) kills its relation and takes
\(\delta\) to one. At the toric strata the residual weights are
\(1,0,-1\), with zero only on the \(e_2\) edges. Thus the action
is semi-free and has exactly the three fixed spheres described above.

The scalar-circle normal trivialization of \ref{maps:src:D.2.14} is given by the following maps. Multiply normal
coordinates by \(\sqrt{1+|w|^2}\), conjugate the negative weight
coordinate, and use the disc extension in \(SU(2)\)

\[
 {\mathsf U}(r,\theta)=\begin{pmatrix}
 \sin\beta(r)e^{-i\theta}&-\cos\beta(r)\\
 \cos\beta(r)&\sin\beta(r)e^{i\theta}
 \end{pmatrix},
\]

where \(\beta=r\) near zero and \(\beta=\pi/2\) on the boundary
collar. The inverse is \({\mathsf U}^*\). Reverse the core parametrization
when conjugating one normal coordinate, so the total orientation
remains the complex orientation. An invariant normal exponential
and its inverse give \(S^2\times D^4\) with scalar action.

Let \({\mathcal Q}=X/S^1\), using the smooth normal quotient

\[
 z\longmapsto\mathcal H(z)=
 \frac{(2\operatorname{Re}(z_1\bar z_2),
          2\operatorname{Im}(z_1\bar z_2),|z_1|^2-|z_2|^2)}{|z|},
 \quad \mathcal H(0)=0.
\]

It has \(|\mathcal H(z)|=|z|\) and precisely the scalar-circle
fibres. The standard northern and southern Hopf sections give the
local inverse with its orbit phase. Orient the base with its
horizontal five-plane before the positive circle direction. Write
\({\mathcal C}_i\) for the three quotient spheres, \(X_0\) for the
complement of their tubular neighbourhoods upstairs, and \({\mathcal Q}_0\)
for the corresponding complement downstairs. The normal meridians
\(\sigma_i\) are oriented so the Euler class \(e\) of
\(X_0\to {\mathcal Q}_0\) has value one on each.

Transversality of discs of dimension at most three to the codimension-
four fixed set makes \(H_2(X_0)\to H_2(X)\) an isomorphism,
and preserves \(\pi_1\). The bundle sequence and Hurewicz give

\[
 \begin{gathered}
 \pi_1({\mathcal Q}_0)=0,\\
 0\longrightarrow {\mathcal A}_{\rm hom}:=H_2(X)
 \xrightarrow{p_*} {\mathcal B}_{\rm hom}:=H_2({\mathcal Q}_0)
 \xrightarrow{\langle e,-\rangle}\mathbb Z\longrightarrow0.
 \end{gathered}
\]

Here \(\sigma_0\) is an explicit splitting. Its inverse sends
\(b\) to \((p_*^{-1}(b-e(b)\sigma_0),e(b))\); the first term
lifts as a sphere with Euler-zero pulled-back circle bundle.
Adjoining the normal \(D^3\) bundles gives
\(\pi_1({\mathcal Q})=0\), \(H_2({\mathcal Q},{\mathcal Q}_0)=0\), and the three meridian
generators of \(H_3({\mathcal Q},{\mathcal Q}_0)=\mathbb Z^3\). Hence the natural
orbit map

\[
 \Theta:H_2(X)\longrightarrow H_2({\mathcal Q})
\]

is onto, and its kernel is generated by the two classes whose
images in the free-circle exact sequence are \(\sigma_1-\sigma_0\) and
\(\sigma_2-\sigma_0\).

The two fixed-curve classes lie in this kernel. Indeed for
each \(F_i\), choose a toric double curve \(E_i\) in direction
\(e_1\) or \(e_2-e_1\) with the same pair of component endpoints.
Its relative intersection column equals that of \(F_i\) in
the fixed-curve intersection matrix: minus one on the two endpoint components and two on the
third. By the cusp splitting, \([F_i]-[E_i]\) lies in the full boundary image.
After the global eliminations its image is a multiple of \(s\).
The \(\delta\) column of the log relation matrix imposes \(ns=0\), so for either
sign under consideration \([F_i]=[E_i]\) in \(H_2(X)\).
On \(E_i\simeq S^2\) the action is the ordinary weight-one
rotation; its orbit map factors through an interval. A small
push-off into \(X_0\) is homotopic in \(X\) to that sphere;
composing this homotopy with the continuous orbit map proves
\(\Theta[E_i]=0\). Thus \(\Theta[F_0]=\Theta[F_1]=0\).

The kernel is determined integrally as follows. The subgroup
\(L=\mathbb ZF_0\oplus\mathbb ZF_1\) is a direct summand of
the simply connected homology groups. The kernel is generated by two elements and contains \(L\).
Project those two generators to \(L\). Their map from
\(\mathbb Z^2\) to \(L\) is onto and hence unimodular.
Accordingly the kernel is the graph of a homomorphism from \(L\)
to its complementary summand. Since it contains the pure subgroup
\(L\), that homomorphism is zero. This supplies the exact quotient
and its inverse on the surviving coordinates:

\[
 H_2({\mathcal Q})=\mathbb Z\langle\Theta U\rangle\oplus(\mathbb Z/m)^2,
 \quad (u,z)\longleftrightarrow\Theta(uU+z),\qquad [{\mathcal C}_i]=0.
\]

The last statement follows by a constant normal push-off of the
fixed curve in the normal trivialization; its quotient is a normal push-off of \({\mathcal C}_i\).
For \(i=2\), the fixed-curve matrix and cusp splitting give \([F_2]=-[F_0]-[F_1]\) modulo
\(s\), which is zero. This proves the statement for all three
components. Hurewicz makes them nullhomotopic in \({\mathcal Q}\).

The manifold \(X\) is spin. On each \(F_i\), the fixed-curve charts give
\(c_1(TX)[F_i]=2-1-1=0\). On \(U\), a component sphere at a
fixed suspension angle in the original star, the canonical form
\(\eta_N\wedge dw/w\) is nonzero and descends. Indeed the
canonical bundle formula \cite[Theorem 6.8, p.~26]{manSS10} pulls
\(K_{N_*}\) back from a line on the base disc; choose its
nonvanishing generator \(\eta_N\). On the regular locus it
is a base differential wedged with the invariant elliptic differential,
so translation fixes it; the identity extends over the central fibre.
A holomorphic line-frame change \(w'=f(y)w\) adds
\(\eta_N\wedge d\log f=0\), since \(N_*\) is a complex
surface. The same computation with the multiplier of
\(\phi_*\) proves \(\phi_*^*(\eta_N\wedge dw/w)
=\eta_N\wedge dw/w\). Multiplication by this descended form,
with inverse division, trivializes the canonical line on the original star quotient.
Thus \(c_1\) vanishes on all three free generators in the simply connected homology groups, and
\(H_1(X)=0\) makes \(H^2(X)=\operatorname{Hom}(H_2(X),\mathbb Z)\).
Hence \(c_1(X)=0\) and \(w_2(X)=0\).

On the free part \(TX_0=p^*T{\mathcal Q}_0\oplus\mathbb R\). The mod-two
Gysin sequence makes the kernel of \(p^*\) in degree two the span
of \(e\bmod2\). This nonzero class evaluates to one on a
meridian \(\sigma_i\), while \(w_2({\mathcal Q}_0)\) evaluates to zero
there, because its restricted tangent bundle is
\(TS^2\oplus\mathbb R^3\). Therefore \(w_2({\mathcal Q}_0)=0\).
Excision makes restriction from \(H^2({\mathcal Q};\mathbb Z/2)\) injective,
so \({\mathcal Q}\) is spin. Its torsion linking form is alternating by \cite[Lemma F, p.~372]{manBarden65}. Choose a class \(x\) of order \(m\)
and \(y\) with linking value \(1/m\). The inverse coordinates
are \(z\mapsto(m\ell(z,y),-m\ell(z,x))\) modulo \(m\).
At \(m=1\) these are the maps of the zero group. Applying \cite[Theorems 2.2--2.3, pp.~376--377]{manBarden65} to these linking coordinates and the free coordinate gives an oriented diffeomorphism

\[
                       {\mathcal Q}\longrightarrow(S^2\times S^3)\#N_m.
\]

It induces exactly these group maps, and its inverse induces their
inverses.

\emph{Localization of the fixed link.}

The quotient of the original-star representative of \(U\) is an
embedded sphere \({\mathcal V}_{\rm orb}\subset {\mathcal Q}\), carrying the free generator in
the quotient homology formula. The component lift is a section of the circle bundle on it.
It is disjoint from the fixed link. Remove \({\mathcal V}_{\rm orb}\) and write
\({\mathcal Q}'={\mathcal Q}\setminus {\mathcal V}_{\rm orb}\). General position of loops and two-discs
gives \(\pi_1({\mathcal Q}')=0\). Excision identifies
\(H_3({\mathcal Q},{\mathcal Q}')\) with \(\mathbb Z\), and the map from
\(H_3({\mathcal Q})\) is intersection with the primitive free generator
\([{\mathcal V}_{\rm orb}]\). Integral duality makes it onto. Consequently
\(H_2({\mathcal Q}')\to H_2({\mathcal Q})\) is an isomorphism, and all \({\mathcal C}_i\)
are nullhomologous already in \({\mathcal Q}'\).

Choose a small oriented five-disc \(D\subset {\mathcal Q}'\). We carry the
entire three-component link into \(D\) by parametrized isotopies
supported in \({\mathcal Q}'\). First move \({\mathcal C}_0\) to a round nullhomotopic
sphere in \(D\). Theorem B(b) of \cite[pp.~367--368]{manBarden65} applies to this ambient manifold: the relative connectivity here is supplied by connectedness
and simple connectivity, and \(2\cdot5>3\cdot2+3\).

For the next stage, remove the already placed components. The
resulting ambient manifold remains simply connected. The kernel of
its \(H_2\) map to \(H_2({\mathcal Q}')\) is freely generated by the
meridian two-spheres of those components. To check this exactly,
use the pair sequence and excision; the preceding map from
\(H_3({\mathcal Q}')\) is intersection with their zero classes and vanishes.
The next component has zero image in \(H_2({\mathcal Q}')\), so its class
is an integral combination of these meridians. The complement of
the same components in \(D\) has precisely these meridian
generators, and its inclusion induces the identical coordinate map.
Represent the required combination by a map from \(S^2\) in this
disc complement. General position, or part (a) of the same theorem,
makes that map an embedding. Hurewicz and the matching meridian
coordinates make it homotopic to the component being moved.
Part (b) therefore gives a parametrized isotopy in the complement
of all components already placed. This argument applies first to
\({\mathcal C}_1\) and then to \({\mathcal C}_2\).

Each compact isotopy track admits an extension of its velocity to a
compactly supported ambient vector field. Its flow is the ambient
map and the backward flow is its inverse. The support avoids \({\mathcal V}_{\rm orb}\)
and all previously placed spheres. Thus the resulting diffeomorphism
fixes \({\mathcal V}_{\rm orb}\), moves the full link into \(D\), and retains every
linking class. Their normal bundles are trivial by the normal trivialization and orbit map,
and their oriented framing adjustments have one homotopy class
because \(\pi_2(SO(3))=0\).

Transport the free bundle and its normal maps under this
diffeomorphism. On \({\mathcal Q}\setminus\mathring D\) its Euler class
is zero: \(H^2({\mathcal Q}\setminus\mathring D;\mathbb Z)=\mathbb Z\)
is detected by \({\mathcal V}_{\rm orb}\), and the preserved component section has
Euler value zero there. The \v{C}ech logarithm/partition-of-unity
construction used in the type-II product model gives an explicit trivialization of
the entire outside bundle, with reciprocal gauge maps as inverse.
Along \(\partial D=S^4\) all circle bundles are trivial; their
two trivializations differ by an \(S^1\)-valued function with a
real logarithm. A collar cutoff of that logarithm makes them agree
on a full collar. Its negative supplies the inverse adjustment.

Let \(W_D\) denote the total-space piece above \(D\), with the
fixed normal pieces restored. Its boundary is parametrized as
\(S^4\times S^1\). Define the closed oriented circle manifold

\[
 M_L=W_D\cup_{\mathrm{id}}(D^5\times S^1).
\]

All gluings here use the same trivialization and the product boundary
orientation. Van Kampen gives \(\pi_1(W_D)=\pi_1(X)=0\): the
outside piece has fundamental group \(\mathbb Z\), and the
boundary circle maps to its generator. The identical argument in
\(M_L\) gives \(\pi_1(M_L)=0\).

For degree two the boundary has zero \(H_2\), and its \(H_1\)
maps isomorphically to the outside circle. Thus the Mayer--Vietoris
maps give

\[
 H_2(X)=H_2(W_D)\oplus H_2({\mathcal Q}),\qquad
 H_2(M_L)=H_2(W_D)=\mathbb Z\langle F_0,F_1\rangle.
\]

The outside summand is represented by the preserved section \({\mathcal V}_{\rm orb}\)
and the torsion lifts. The basis of \(H_2(M_L)\) follows from the computed kernel of \(\Theta\). In degree three the boundary has no
\(H_2\) or \(H_3\), while the outside product has
\(H_3({\mathcal Q})\oplus H_2({\mathcal Q})=\mathbb Z^2\oplus(\mathbb Z/m)^2\).
Comparison with the simply connected homology groups gives \(H_3(W_D)=H_3(M_L)=0\).
The spin structure on \(W_D\) extends over \(D^5\times S^1\):
its two possible spin structures restrict to the two on
\(S^4\times S^1\). Choose the matching one. Hence \(M_L\)
is a closed simply connected spin six-manifold with torsion-free
homology, free \(H_2\) rank two, and zero middle homology.

\emph{The flag-manifold block.}

Each component class \(\eta_i\) has zero restriction to the
outside piece: its value on \({\mathcal V}_{\rm orb}\) is zero, because the original
star sphere and cusp divisor were disjoint and the localization
fixed \({\mathcal V}_{\rm orb}\). The outside cohomology group \(H^2({\mathcal Q}\setminus\mathring D;\mathbb Z)\)
is generated by the functional dual to \({\mathcal V}_{\rm orb}\).
The relative groups of the two decompositions are both
\(H^2(W_D,\partial W_D)\). In each, the kernel of extension
to the closed manifold is the same image of
\(H^1(\partial W_D)\), because the outside product and the
new cap each restrict isomorphically in degree one to
\(S^4\times S^1\). Their quotient therefore identifies the
classes of \(X\) that vanish outside with \(H^2(M_L)\).
Use this quotient map to extend the same relative cocycles by zero
into the cap of \(M_L\).
This gives the precise correspondence to classes, still called
\(\eta_i\), on \(M_L\). Cup-product evaluation and
\(p_1\)-pairing use only \(W_D\), where the orientations and
tangent bundles agree. Consequently the divisor cubic identities and all the
following characteristic pairings transfer unchanged to \(M_L\).

Let \(a_{\rm fl},b_{\rm fl}\) be the integral basis dual to \(F_0,F_1\) in the block homology splitting.
The fixed-curve intersection matrix gives

\[
 \begin{gathered}
 \eta_0=-a_{\rm fl}-b_{\rm fl},\qquad \eta_1=2a_{\rm fl}-b_{\rm fl},\\
 \eta_2=-a_{\rm fl}+2b_{\rm fl}.
 \end{gathered}
\]

Substitution in the component cubic, or polarization of
its four coefficients, gives

\[
 a_{\rm fl}^3=b_{\rm fl}^3=0,\qquad
 \langle a_{\rm fl}^2b_{\rm fl},[M_L]\rangle=
 \langle a_{\rm fl}b_{\rm fl}^2,[M_L]\rangle=-1.
\]

Solving for the four coefficients gives the binary cubic \(-3a_{\rm fl}b_{\rm fl}(a_{\rm fl}+b_{\rm fl})\). Thus the class \(a_{\rm fl}\) or \(b_{\rm fl}\)
has square different from zero, although its own cube is zero;
the two mixed products are essential to the recognition.

For each embedded divisor, the tangent-normal sequence gives
\(p_1(TX)|_{D_i}=p_1(TD_i)+c_1(N_{D_i/X})^2\).
The boundary-curve classes in the blow-up basis give \(c_1(D_i)^2=6\) and
\(c_2(D_i)=\chi(D_i)=6\). Hence

\[
 \langle p_1(TX)\eta_i,[X]\rangle=(6-2\cdot6)+6=0.
\]

The component classes in the divisor basis span the rational degree-two group
of \(M_L\), and its \(H^4\) is torsion free. The perfect
\(H^4\)--\(H^2\) pairing therefore proves \(p_1(M_L)=0\).

Put $F=\operatorname{Fl}(\mathbb C^3)$ as in
Definition~\ref{maps:def:complex-flag-threefold}.

Projection to the line is the projectivization of the dual rank-two
quotient bundle on \(\mathbb P^2\); its fibre is \(\mathbb P^1\).
Let \(h_{\rm fl},k_{\rm fl}\) be the hyperplane classes from the two factors. The
projective-bundle maps give the integral basis and relation

\[
 \begin{gathered}
 H^*(F;\mathbb Z)=
 \mathbb Z[h_{\rm fl},k_{\rm fl}]/(h_{\rm fl}^3,k_{\rm fl}^2-h_{\rm fl}k_{\rm fl}+h_{\rm fl}^2),\\
 \langle h_{\rm fl}^2k_{\rm fl},[F]\rangle=1.
 \end{gathered}
\]

The degree-two basis restricts as one generator from the base and
one of degree one on the fibre; these evaluations give its inverse
coordinates. The bundle sequence makes \(\pi_1(F)=0\), and the flag cohomology ring
gives \(H^3(F)=0\) and torsion-free homology. It also gives
\(k_{\rm fl}^3=h_{\rm fl}^3=0\) and \(\langle h_{\rm fl}k_{\rm fl}^2,[F]\rangle=1\).
As the smooth incidence hypersurface of bidegree \((1,1)\), its
normal line in the product is \(\mathcal O(1,1)\). The tangent
exact sequence therefore yields

\[
 \begin{gathered}
 c_1(F)=2h_{\rm fl}+2k_{\rm fl},\\
 \begin{aligned}
 p_1(F)&=3h_{\rm fl}^2+3k_{\rm fl}^2-(h_{\rm fl}+k_{\rm fl})^2\\
 &=2(h_{\rm fl}^2-h_{\rm fl}k_{\rm fl}+k_{\rm fl}^2)=0.
 \end{aligned}
 \end{gathered}
\]

In particular \(F\) is spin. The group isomorphism specified by
the cohomology pullback \(h_{\rm fl}\mapsto-a_{\rm fl},k_{\rm fl}\mapsto-b_{\rm fl}\) has inverse
\(a_{\rm fl}\mapsto-h_{\rm fl},b_{\rm fl}\mapsto-k_{\rm fl}\). It matches all four cubic
coefficients in the primitive and flag cubic identities and the zero Pontryagin classes.

Apply the injectivity theorem \cite[Theorem 5, p.~796]{manZhubr75} and the construction and uniqueness for zero free third homology in \cite[\S4.1, p.~809]{manZhubr75}. The classifying maps to
\(K(\mathbb Z^2,2)=(\mathbb P^\infty)^2\) are the two line
classes just specified. Their \(P_{\rm bord}\) invariants are one quarter
of the pushed-forward dual \(p_1\), hence zero. Their \(h_{\rm bord}\)
invariants are detected by the four integral degree-six monomials
\(h_{\rm fl}^3,h_{\rm fl}^2k_{\rm fl},h_{\rm fl}k_{\rm fl}^2,k_{\rm fl}^3\); all four agree. Thus the injective
\((P_{\rm bord},h_{\rm bord})\) map gives the same spin bordism class. Both manifolds
have the zero free middle group used in the construction in \cite[\S4.1, p.~809]{manZhubr75}.
That uniqueness gives an orientation-preserving diffeomorphism

\[
                              M_L\longrightarrow F
\]

inducing the stated integral isomorphism. Its inverse induces the
inverse coordinates.

\emph{Reconstruction and the suspension framing.}

The localization also determines the full smooth reconstruction.
The model \(\widetilde\Sigma_0{\mathcal Q}\) has the single nullhomotopic
fixed sphere and quotient \({\mathcal Q}\) constructed in \ref{maps:src:D.2.14}.
Its angular bundle is trivial off the disc containing that sphere.
Its cap-to-complement gluing is
\((x,z)\mapsto(\alpha_{{\rm rot},z}x,z)\), where \(\alpha_{\rm rot}\) rotates the
first real two-plane once. The opposite-direction mapping-torus
model is compared with this one by
\({\mathsf A}_{\rm inv}=\operatorname{diag}(1,-1,-1,1,1)\),
\({\mathsf A}_{\rm inv}\alpha_{{\rm rot},z}{\mathsf A}_{\rm inv}^{-1}=\alpha_{{\rm rot},z}^{-1}\), and the collar rotation in
the (2,3)-plane. These are the framing maps and their inverses from \ref{maps:src:D.2.14}.

Take the equivariant oriented connected sum of \(M_L\) and
\(\widetilde\Sigma_0{\mathcal Q}\) at one point of a fixed sphere in each.
In the local coordinates \(\mathbb R^2\oplus\mathbb C^2\)
with scalar circle action, reflect one real fixed coordinate on the
joining \(S^5\). This is an orientation-reversing equivariant
boundary diffeomorphism and is its own inverse; radial collars and
invariant rounding give the connected-sum map. The quotient is
\(S^5\#{\mathcal Q}\), identified with \({\mathcal Q}\) by the parametrized disc
identification. The fixed link is the link of \(M_L\), with one sphere
connected-summed with the nullhomotopic model sphere in a ball.
That component is homotopic to the original parametrized component
in the complement of the other two: the added sphere has zero
meridian coordinates there. The same embedding-isotopy theorem
used in the link-localization argument carries it back, with an inverse ambient flow. Hence
the quotient pair is precisely the localized pair \(({\mathcal Q},\cup {\mathcal C}_i)\).

The free Euler class also agrees with that of \(X\). It has
residue one on each normal meridian, agrees with the inside-
disc bundle in \(M_L\), and is zero on the outside generator \({\mathcal V}_{\rm orb}\).
The residue sequence shows that the difference of two such classes
comes from \(H^2({\mathcal Q};\mathbb Z)=\mathbb Z\), where evaluation on
\({\mathcal V}_{\rm orb}\) determines it. Therefore this difference is zero. Torsion
in the quotient homology formula contributes no integral degree-two functional.

An isomorphism of these circle bundles can be written from a common
good cover: take logarithms of the ratios of transition functions,
remove their integral cocycle using equality of the Euler classes,
and use a partition of unity to express the remaining cocycle as
differences of functions. Multiplication by their exponentials
is the bundle isomorphism, with reciprocal maps as inverse.
Near every fixed sphere, an oriented normal rotation in \(SO(3)\)
lifts to \(SU(2)\) by path lifting on the simply connected core.
The lift and its conjugate transpose give inverse equivariant
normal maps covering the orbit map. On their boundary collars the bundle
and normal maps differ by a gauge on \(S^2\times S^2\) times
an interval. Its \(H^1\) is zero, so a real logarithm and a collar
cutoff make the maps coincide. Again the negative logarithm is
the inverse correction. The maps glue on full collars to an
orientation-preserving diffeomorphism, with glued inverse,

\[
                        X\longrightarrow M_L\#\widetilde\Sigma_0{\mathcal Q}.
\]

The reconstruction carries the localized fixed link, its linking
classes and its primitive upstairs classes to their specified images.

Use the oriented quotient diffeomorphism, adjusted on a chosen disc by parametrized
disc isotopy, and its inverse to lift to the twisted suspensions.
The neck construction in \ref{maps:src:D.2.16} distributes the
precisely framed suspension over the connected sum. Its separating
hypersurface is
\((D^4\times S^1)\cup(S^3\times D^2)\); the map
\((x,z)\mapsto(\alpha_{{\rm rot},z}^{-1}x,z)\) on the first piece and
identity on the second identifies it with the standard \(S^5\).
The inverse uses \(\alpha_{{\rm rot},z}\). The rotations cancel in the
opposite boundary parametrizations, so this is the oriented
connected-sum splitting. The product-clutching maps of \ref{maps:src:D.2.14} then give

\[
 \widetilde\Sigma_0{\mathcal Q}\cong
 (S^2\times S^4)\#(S^3\times S^3)\#\widetilde\Sigma_0N_m.
\]

Composing the flag identification, reconstruction and suspension maps proves \eqref{maps:eq:D.V3} for every allowed log with
\(n=1\) or \(n=-1\), with the original complex orientation.
At \(m=1\) the last factor is the standard oriented \(S^6\)
by the boundary-of-\(D^5\times D^2\) map in \ref{maps:src:D.2.14}.
Removing that summand uses its oriented disc map. Distinct positive
\(m\) give distinct smooth types because the torsion order in
\(H_2(X)\) is \(m^2\). The smooth type is independent of
\(a,b',c\) and of the two signs of \(n\) within this scope.
\end{proof}

\subsection{Six primitive original families}
\begin{theorem}[Six original additive families]\hypertarget{tag-02KV}{}\label{maps:src:D.2.20}
{\normalfont\scriptsize tag 02KV}\par

Use the following configurations and section choices:
\[
\begin{array}{c|c|l|r}
 \mathrm{OS}&\mathrm{profile}&\text{fibres}&d\\\hline
 43&\mathrm{default}&III^*,3I_1&2\\
 43&II&III^*,II,I_1&2\\
 46&\mathrm{default}&I_3^*,3I_1&4\\
 46&II&I_3^*,II,I_1&4\\
 49&III&IV^*,III,I_1&6\\
 50&IV&I_1^*,IV,I_1&12.
\end{array}
\]
Take $P=G$, $Q=dG$, retain every original additive filling,
and put the simple linearization zero at the first marked
$I_1$. The other $I_1$ fillings are unchanged. At an
additional regular fibre perform any free log
\[
 m>0,\qquad v=a e_1+b e_2+c\delta+n\gamma,
 \qquad\gcd(m,a,b,c,n)=1.
\]
Set
\[
        g=\gcd(m,a,b,n),\qquad
        T_{m,g}=\mathbb Z/g\oplus(\mathbb Z/m)^2.
\]
For $n\ne0$, the integral groups are
\begin{equation}
\begin{aligned}
 \pi_1(X)&=H_1(X)=\mathbb Z/|n|, &H_5(X)&=0,\\
 H_2(X)&=\mathbb Z^7\oplus T_{m,g},
 &H_3(X)&=\mathbb Z^{14}\oplus T_{m,g},\\
 H_4(X)&=\mathbb Z^7\oplus\mathbb Z/|n|.
\end{aligned}
\hypertarget{tag-020M}{}\label{maps:eq:D.O1}
\end{equation}
For $n=0$, they are
\begin{equation}
\begin{aligned}
 \pi_1(X)&=H_1(X)=H_5(X)=\mathbb Z,\\
 H_2(X)&=\mathbb Z^8\oplus T_{m,g},
 &H_3(X)&=\mathbb Z^{14}\oplus T_{m,g},\\
 H_4(X)&=\mathbb Z^8.
\end{aligned}
\hypertarget{tag-020N}{}\label{maps:eq:D.O2}
\end{equation}
In all cases $H_0=H_6=\mathbb Z$ and $\chi(X)=2$.
The simply connected members are exactly $n=\pm1$, and
all of them have oriented diffeomorphism type
\begin{equation}
 X\cong\#^7(S^2\times S^4)\#^7(S^3\times S^3)
                                    \#\widetilde\Sigma_0N_m.
 \hypertarget{tag-020O}{}\label{maps:eq:D.O3}
\end{equation}
In particular none is a six-sphere. For the Engel row,
the trivial regular log $(m,v)=(1,0)$ recovers the
untwisted groups of \SeriesRef{push:src:A.4.2} $(1,1,8,14,8,1,1)$ integrally,
with no torsion.
\end{theorem}


\begin{proof}
Put $\Lambda=\mathbb Z\langle e_1,e_2,\delta,\gamma\rangle$,
$J=\left(\begin{smallmatrix}0&1\\-1&0\end{smallmatrix}\right)$, and
use exterior order $(12,13,14,23,24,34)$. The positive physical
line circle is $\delta$; the positive contracting period $\gamma$
covers translation by $-Q$.

\emph{The marked construction.}
The elliptic monodromies $A_i$ and generator cocycles $g_i$ in the
stated fibre order are as follows. We abbreviate a matrix by its
two rows $(a,b;c,d)$.
\[
\begin{array}{c|l|l}
43/\mathrm{default}&
(2,-5;1,-2),(-2,9;-1,4),(1,1;0,1),(1,1;0,1)&
(-1,0),(0,0),(0,0),(1,0)\\
43/II&(2,-5;1,-2),(-2,7;-1,3),(1,1;0,1)&
(-1,0),(0,0),(1,0)\\
46/\mathrm{default}&
(-1,-3;0,-1),(1,1;0,1),(2,1;-1,0),(0,1;-1,2)&
(1,1),(0,0),(0,0),(-1,-1)\\
46/II&(-1,-3;0,-1),(1,1;-1,0),(0,1;-1,2)&
(1,1),(0,0),(-1,-1)\\
49/III&(1,-3;1,-2),(-2,5;-1,2),(1,1;0,1)&
(-1,0),(-2,-1),(0,0)\\
50/IV&(-1,-1;0,-1),(-1,1;-1,0),(0,1;-1,2)&
(-1,1),(-2,0),(-1,-1).
\end{array}
\]
The height of $G$ is $1/d$. Its component residue is nonzero at
each reducible additive fibre, and multiplication by $d$ kills it.
The components met by $O$ and $G$ have multiplicity one. Thus $Q$
is narrow, and $M=\mathcal O(G-O)$ has component degrees
$e_p-e_0$. At $II$ and $I_1$ its multidegree is zero.
The degree-one criterion of \ref{maps:src:D.2.19}, applied with
$\langle G,dG\rangle=1$ and narrow $Q$, gives
$M\otimes t_Q^*M^{-1}=f^*\mathcal O(1)$.
Choose the section with its simple zero at the specified $I_1$.

Solve $(A_i-I)z_i=g_i$. At a node write
$A_i-I=u_i u_i^tJ$ with $u_i$ primitive, choose $r_i$ with
$u_i^tJr_i=1$, and put $z_i=\kappa_i r_i$ when
$g_i=\kappa_i u_i$. Put $z_{Qi}=dz_i$,
$b_i=-z_i^tJ$, $\beta_i=-g_i^tJA_i$, and let $w_i$ be one at
the marked node and zero elsewhere. The Poincar\'e-line periods are
\[
 T_i=\begin{pmatrix}
 A_i&0&dg_i\\ \beta_i&1&\beta_i z_{Qi}+w_i\\0&0&1
 \end{pmatrix},\qquad
 S_i=\begin{pmatrix}I&0&z_{Qi}\\b_i&1&0\\0&0&1\end{pmatrix}.
\]
Since $b_i z_{Qi}=0$, the inverse $S_i^{-1}$ changes the two
off-diagonal blocks to $-z_{Qi}$ and $-b_i$. Multiplication gives
\[
 S_iT_iS_i^{-1}=\operatorname{diag}(A_i,1,1)+w_iE_{34},
 \qquad T_1\cdots T_s=I.
\]
Every $T_i$ is integral. A rational $S_i$ gives the indicated
map from $\mathbb R^4/\Lambda$ to $\mathbb R^4/S_i\Lambda$,
with that precise target lattice and the displayed inverse.

\emph{The original additive caps.}
Let $E_i$ be the original elliptic-surface neighbourhood, with
components $C_0,\ldots,C_{r-1}$. Resolve its central divisor by
point blowups. The normal-disc collapse in each component and the
crossing collapse
\[
 (a,b)\longmapsto(a-t\min(a,b),b-t\min(a,b)),
 \qquad a=|x|^2,\ b=|y|^2,
\]
give a retraction of a plumbing neighbourhood fixing its divisor.
It descends through the proper blowdown quotient, since each
exceptional tree is fixed throughout. Radial transport identifies
this neighbourhood with a smaller fibre-disc neighbourhood; reverse
transport is the inverse. This also treats the tangency of $III$,
the triple point of $IV$, and the cusp of $II$. Consequently
$\pi_1(E_i)=0$ and $H_2(E_i)=\mathbb Z^r$, with no other positive
homology. At $II$ one has $r=1$.

The local logarithm of $-Q$ in
\cite[\S2.4, equation (2.2), p.~6]{manUU21} contracts the translation
on this very germ. The translations extend over the components by
\cite[\S7.6, pp.~33--34]{manSS10}. Let $E(M)$ be the unit circle
bundle of $M|E_i$. Lift this isotopy by a connection. Its endpoint
differs from the angular contracting lift by a gauge $E_i\to S^1$,
which has a real logarithm because $H^1(E_i;\mathbb Z)=0$.
A cutoff exponential corrects the endpoint, with reciprocal
exponential and inverse translation giving the inverse maps.

The radial quotient has a positive roof $r(y)$. Choose $k_0>0$
with $k_0\max r<1$ and bumps $b_y$ of integral one supported in
the middle of $[0,r(y)]$. The increasing map
\[
 u\longmapsto\int_0^u
 [k_0+(1-k_0r(y))b_y(v)]\,dv
\]
takes the roof to $[0,1]$, has derivative $k_0$ at both ends,
and has a smooth increasing inverse. Together with the angular isotopy,
including its endpoint adjustment, it identifies the original cap with
$E(M)\times S^1_\gamma$, retaining the boundary marking $S_i$.
At a reducible additive fibre its two Chern classes are
$(e_p-e_0,0)$.

The circle-bundle homotopy and Gysin sequences give
$\pi_1(E(M))=0$, $H_2(E(M))=\ker(e_p-e_0)$, and
$H_3(E(M))=\mathbb Z^r$. Projection supplies the first isomorphism
on homology: a representing sphere of degree zero has a section,
and two such sections are homotopic since $\pi_2(S^1)=0$.
K\"unneth therefore gives
\begin{equation}
 \pi_1(V_i)=\mathbb Z\gamma,\qquad
 (H_0,H_1,H_2,H_3,H_4)(V_i)
 =(\mathbb Z,\mathbb Z,\mathbb Z^{r-1},
                  \mathbb Z^{2r-1},\mathbb Z^r),\qquad H_5(V_i)=0.
 \hypertarget{tag-020P}{}\label{maps:eq:D.O4}
\end{equation}
This formula concerns the reducible additive fibres, where the
Euler row is primitive. At $II$ the cap retracts onto
$S^2\times S^1_\delta\times S^1_\gamma$.

Write $F_i=C_0+C_p+\sum_{j\ne0,p}\mu_jC_j$ for the whole elliptic
fibre. The vectors $F_i,C_j$ for $j\ne0,p$ form an integral basis
of the component kernel: its inverse reads
$(x_0,(x_j-\mu_jx_0)_{j\ne0,p})$ from a vector with $x_p=x_0$.
The elliptic projection on the boundary is
$(x_1,x_2)+z_{Qi}x_4$. Its oriented area is
\begin{equation}
 \Omega_i=(1,0,(z_{Qi})_2,0,-(z_{Qi})_1,0).
 \hypertarget{tag-020Q}{}\label{maps:eq:D.O5}
\end{equation}
Thus the fibre maps are $f_{1i}(x)=x_4$ and
$f_{2i}(a)=\Omega_i(a)F_i$.
Every component coefficient is an invariant
exterior covector. Conjugation by $S_i$ gives a two-dimensional
invariant row space, spanned by elliptic area and the transformed
$\delta\wedge\gamma$ row. Its values on $e_{12}$ and $e_{34}$
determine it. The first torus specializes to $F_i$, and the second
projects to zero in $H_2(E_i)$. These two evaluations determine
every component coefficient, giving the full integral map.

The nonzero Smith entries of $T_i-I$ are $1,1$, and those of
$\bigwedge^2T_i-I$ are four ones. The boundary therefore has
$H_1=\mathbb Z^3$ and $H_2=\mathbb Z^4$. Its fixed lattice is
generated by $\delta$ and $(-z_{Qi},0,1)$. The two principal
circle actions sweep the lifted origin disc, so these normal tori
bound and split the Wang projection. The remaining boundary image
is the primitive $\mathbb Z F_i$. Consequently
\begin{equation}
 H_2(V_i)/\operatorname{im}H_2(\partial V_i)
                       =\mathbb Z^{r-2}.
 \hypertarget{tag-020R}{}\label{maps:eq:D.O6}
\end{equation}
The full boundary kernel consists of these two normal tori and
the image of $\ker\Omega_i$ after its Wang relations. The latter
has rank one; before those relations $\ker\Omega_i$ has rank five.
At $II$ the fibre maps are
$(x_3,x_4)$ and $(x_{12},x_{34})$, and the boundary map is onto.

\emph{Nodes, cusp and the global exterior.}
At every node put
$H_i=S_i^{-1}\operatorname{diag}([u_i,r_i],1,1)$.
The matrices $H_i,H_i^{-1}$ are integral, have determinant one,
fix $\delta$, and satisfy
$H_i^{-1}T_iH_i=I+E_{12}+w_iE_{34}$.
In these coordinates an ordinary node has fibre maps retaining
$2,3,4$ in degree one and $12,23,24,34$ in degree two.
They follow from its ruled normalization, its two identified
sections, and their product tori.
At the marked cusp the short directions are $H_ie_1,H_ie_3$,
and the literal valuation columns of the long periods are
$(1,0),(0,1)$. The torus coordinates, their units and their cones
are transported by $H_i$ and its inverse. Thus the primitive cusp
of \ref{maps:src:D.2.16} has the precise fibre maps
\[
 f_{1c}(x)=(x_2,x_4),\qquad
 f_{2c}(a)=(a_{12},a_{14}-a_{23},a_{24},a_{34}),
\]
preceded by $H_i^{-1}$ and $\bigwedge^2H_i^{-1}$ respectively.
An integral section of the latter sends
$(u,v,w,z)$ to $(u,0,v,0,w,z)$ and then applies $\bigwedge^2H_i$.
Its primitive kernel is $\langle e_{13},e_{14}+e_{23}\rangle$.

Geometrically the toric $dP_6$ normalization has its opposite
boundary curves paired by the unit lattice translations.
The conductor differences generate the primitive class $E_1-E_2$
in the presentation as a twice blown-up $\mathbb P^1\times\mathbb P^1$;
the dual two-torus contributes the fourth generator above with
coefficient one. Since the valuation lattice is $\mathbb Z^2$,
all six neighbouring vertices are period translates. The reduced
cusp divisor $D$ is the whole fibre with multiplicity one.
Multiplication by the base coordinate trivializes $\mathcal O(D)$
locally, with reciprocal chart unit as inverse, and globally
$[D]=f^*[\mathrm{point}]$. Hence $[D]^2=0$ and $D^3=0$.
The conductor identifications and the whole-fibre Cartier map
determine these divisor products.

The global exterior coinvariant is $\mathbb Z$, detected by
$\omega=(1,0,0,0,0,-d)$. Indeed the concatenated matrix
$[\bigwedge^2T_i-I]_i$ has five unit Smith entries and is annihilated
by this primitive row; the inverse chooses $e_{12}$.
The basis
$e_{13},e_{14}-(z_Q)_2e_{12},e_{23},
 e_{24}+(z_Q)_1e_{12},e_{34}$
of $\ker\Omega_i$ has $\omega$-values
$0,-(z_Q)_2,0,(z_Q)_1,-d$. Its image is therefore $h_i\mathbb Z$,
where the values are
\begin{equation}
\begin{array}{c|c|r|r|r}
 \mathrm{OS/profile}&\text{fibre}&r&z_Q&h_i\\\hline
 43/\mathrm{default},43/II&III^*&8&(3,1)&1\\
 46/\mathrm{default},46/II&I_3^*&8&(1,-2)&1\\
 49/III&IV^*&7&(6,2)&2\\
       &III&2&(9,3)&3\\
 50/IV&I_1^*&6&(9,-6)&3\\
      &IV&3&(8,-8)&4 .
\end{array}
\hypertarget{tag-020S}{}\label{maps:eq:D.O7}
\end{equation}
Their gcd is one in each full configuration, so these original
caps kill the global area coinvariant. The free local defects sum
to six in the first four rows and five in the last two.

Remove the additional regular-fibre disc and call the resulting
total space $U$. Decompose its base as a boundary connected sum
of the singular discs. The period maps trivialize the intervening
strips; inverse logarithms modulo the period lattice give the
inverse trivializations. The intersections retract onto the
marked four-torus. For $q=1,2$ form $M_q$ by placing $f_{qi}$
in the $i$-th target block and $-f_{qs}$ in the final block for
each of the first $s-1$ source blocks. Then
\[
 0\longrightarrow\operatorname{coker}M_2\longrightarrow H_2(U)
 \longrightarrow\ker M_1\longrightarrow0.
\]
The explicit maps above give the following integral reductions.
All nonzero Smith entries of both matrices are one.
\[
\begin{array}{c|rr}
 &\operatorname{rank}\operatorname{coker}M_2&
   \operatorname{rank}\ker M_1\\\hline
43/\mathrm{default},43/II,46/\mathrm{default},46/II&6&4\\
49/III,50/IV&5&5 .
\end{array}
\]
These reductions are obtained by integer row and column operations
on the displayed fibre maps. In each row every vector
$(f_{21}(a),0,\ldots,0)$ lies in the image of $M_2$: the component
kernel bases and the fibre relations of \eqref{maps:eq:D.O7}
give its integer preimage. Thus every vertical two-torus is zero
in $H_2(U)$. Van Kampen also gives
$\pi_1(U)=\mathbb Z\gamma$: one reducible additive cap kills
$e_1,e_2,\delta$, and $\gamma^*$ extends over every cap.

The outer rim $\mu\times x$ has connecting coordinates
\[
 R_i(x)=(T_i-I)T_{i+1}\cdots T_sx,\qquad i<s.
\]
Their sum, including the final term, is zero. Applying the kernel
coordinates of $M_1$ gives three primitive independent columns
at $x=e_1,e_2,\gamma$; their Smith entries are $1,1,1$.
To see the same saturated coordinates on the regular-base graph,
insert $T_0=I$, put
$D_1=[T_0-I\ \cdots\ T_{s-1}-I]$, and let $j_i$ be block inclusion
for $i<s$. The last injection is
$j_s(x)=-(T_1\cdots T_{s-1}x,\ldots,x)^t$.
If $B$ consists of the unchanged full fixed lattices under these
injections, then
$[B,j_0e_1,j_0e_2,j_0\gamma]$ is an integral basis of $\ker D_1$.
Integer elimination gives the inverse whose last three coordinates
on $j_0x$ are $(x_1,x_2,x_4)$.

The positive base meridian has zero contracting coordinate:
choose a constant fractional roof coordinate on the base circle.
That coordinate returns to its starting value after one turn, so the meridian is
nullhomotopic in $U$. Sweeping its bounding disc by $\delta$ gives
a three-chain with boundary $\mu\times\delta$. Changing the
disc or the annular trivialization changes other rim representatives
only by vertical tori, which are already zero. Hence, with the
three rims denoted $R_1,R_2,R_\gamma$,
\[
 H_2(U)=\mathbb Z^7\oplus
          \mathbb Z\langle R_1,R_2,R_\gamma\rangle,\qquad
 j_2(\mu\times x)=x_1R_1+x_2R_2+x_4R_\gamma,\qquad
 j_2(\bigwedge^2\Lambda)=0.
\]
The inverse first reads the connecting coordinates, subtracts the
rims and chosen lifts of the remaining kernel vectors,
then reads the component-cokernel coordinates. Local fillings
of the loops in $\ker f_{1i}$ give the reverse construction.

\emph{The regular logarithm and all integral groups.}
Its core has lattice
\[
 L_v=(\Lambda\oplus\mathbb Zs)/\langle ms-v\rangle.
\]
The given gcd makes this a free rank-four lattice and the translation
by $v/m$ has exact order $m$. Thus its diagonal cyclic action with
the base-root rotation is free. The boundary sends $\mu$ to $s$
and fibre loops to their indicated classes. Van Kampen kills
$e_1,e_2,\delta,s$, leaving exactly $n\gamma=0$.
The inverse cyclic coordinate sends one to the contracting loop.

In degree two use
$\bigwedge^2L_v=\bigwedge^2\mathbb Z^5/
 \langle(ms-v)\wedge\mathbb Z^5\rangle$.
Eliminate the vertical products and identify
$s\wedge e_1,s\wedge e_2,s\wedge\gamma$ with the rims;
$s\wedge\delta$ vanishes. The residual matrix is
\begin{equation}
 \begin{pmatrix}
 m&0&0&a\\0&m&0&b\\0&0&m&n
 \end{pmatrix}.
 \hypertarget{tag-020T}{}\label{maps:eq:D.O8}
\end{equation}
Its determinantal divisors are $g,mg,m^2g$. More explicitly, a
B\'ezout completion takes $(a,b,n)$ to $(h_0,0,0)$, where
$h_0=\gcd(a,b,n)$, and its inverse returns the rim coordinates.
The quotients are $\mathbb Z/\gcd(m,h_0)$ and two copies of
$\mathbb Z/m$. For $a=b=n=0$ use the identity and $h_0=0$.
The coefficient $c$ remains in the geometric freeness condition.

The following degree-one kernel consists of $(-tv,mt)$ with
$nt=0$. It is zero for $n\ne0$; for $n=0$ it is generated by
$(-v,m)$. Fill its two zero images and glue the fillings. Reading
the boundary coefficient gives an inverse splitting onto this
additional free summand. This proves the $H_2$ formulas for both
values of $n$. In the decomposition with all caps separate, the
degree-one source/target ranks are $(19,18)$ for the default rows,
$(14,13)$ for the $II$ rows, and $(14,12)$ for the last two.
The kernel ranks are respectively $1,1,2$ for $n\ne0$ and increase
by one for $n=0$, agreeing with the same calculation after joining
the unchanged caps first.

The physical $\delta$ action is free on the original caps and
ordinary nodes and has finite stabilizers on a regular log.
In a primitive cusp the two triangle charts give weights
$(-1,0,1)$ and $(1,-1,0)$. The residual weights along
$e_1,e_2,e_2-e_1$ are $1,0,-1$. Thus the fixed curves are the
$e_2$-edges. They are disjoint before the period quotient; the
unit lattice acts transitively, and no nonzero period stabilizes
one. Their quotient is one embedded $C\simeq\mathbb P^1$.
The nonzero unit period factors preserve this description.
The circle-action Euler identity
\cite[Lemma 2.1, p.~6]{manGR25} gives $\chi(X)=2$.

Integral duality and universal coefficients give $H_4,H_5$.
The free $H_4$ coordinates are the $H_2$-dual functionals capped
with the complex fundamental class. The character
$\gamma\mapsto1/|n|$ followed by the integral Bockstein and cap
product gives the torsion coordinate, inverted by the linking
pairing. The same character construction on the Smith summands
of $H_2$ gives its dual torsion in $H_3$. The Euler identity
$2-2b_1+2b_2-b_3=2$ gives free middle rank fourteen.
Successive integral pairs of intersection one give the middle
free coordinates and their inverse sums. This proves
\eqref{maps:eq:D.O1}--\eqref{maps:eq:D.O2}.

\emph{The canonical line and the fixed normal model.}
By \cite[Theorem 6.8, p.~26]{manSS10}, the surface canonical line
is $f^*\mathcal O(-1)$. Choose its meromorphic form $\eta$ with
pole on the removed whole fibre. The form $\eta\wedge dw/w$ is
independent of a line frame and invariant under the
contracting lift: each change adds a base one-form whose wedge
with $\eta$ is zero. At a node $xy=q$ one has
$d\log x\wedge dq=dx\wedge dy$; at a cusp $q=xyw$ one has
$d\log w\wedge d\log y\wedge dq=-dx\wedge dy\wedge dw$.
It therefore extends without zero or pole on $U$.
On the logarithmic cover $q=t^m$, $dq/q=m\,dt/t$; fibre
translations add only $dt$ terms to the other factors.
The free quotient is locally biholomorphic, so the descended form
has precisely a simple pole on the reduced core $F_m$. Hence
\begin{equation}
                         K_X=\mathcal O_X(-F_m).
 \hypertarget{tag-020U}{}\label{maps:eq:D.O9}
\end{equation}
The positive primitive normal meridian is $m\mu-v$, mapping to
$-n\gamma$ in the complement. In the cohomological Thom sequence
$H^1(U)\to H^2(X,U)\to H^2(X)$, the first map is multiplication
by $-n$, and the second takes one to $\operatorname{PD}(F_m)$.
Its order is exactly $|n|$ for $n\ne0$. At $n=0$ its image is
primitive and infinite: the quotient embeds in the torsion-free
$H^2(U)$, since $H_1(U)=\mathbb Z$. In particular $c_1(X)=0$
integrally when $n=\pm1$, and these members are spin.

For these two signs the functional $x_3-(c/n)x_4$ on the log
lattice is one on $\delta$ and kills $ms-v$, so the action is
free on that cap. It is consequently semifree on $X$.
Near its fixed sphere the two charts are
\[
 (x,y,w)=(q/(X_1X_2),X_2,X_1),\qquad
 (x',y',w')=(q/X_2,X_1X_2,1/X_1).
\]
The inverses are $(X_1,X_2,q)=(w,y,xyw)$ and
$(1/w',y'w',x'y'w')$. The transition is
$(x',y',w')=(wx,wy,w^{-1})$, giving
$\mathcal O(-1)_{-1}\oplus\mathcal O(-1)_{+1}$.
Conjugate the negative-weight coordinate and use the positive
metric factor $\sqrt{1+|w|^2}$. The resulting transition loop is
$\operatorname{diag}(e^{-i\theta},e^{i\theta})$ in $SU(2)$.
It extends over the disc as
\[
 \begin{pmatrix}
 \sin\beta(r)e^{-i\theta}&-\cos\beta(r)\\
 \cos\beta(r)&\sin\beta(r)e^{i\theta}
 \end{pmatrix},
 \quad \beta(r)=r\text{ near }0,\quad\beta(r)=\pi/2\text{ near }1.
\]
The inverse is its conjugate transpose. Reversing the core
parametrization together with the normal conjugation preserves
the total orientation. This constructs the framed
$C\times D^4$ scalar-circle model of \ref{maps:src:D.2.14}.

\emph{Direct oriented reconstruction.}
Put $X_0=X\setminus(C\times\mathring D^4)$. The normal orbit map is
$z\mapsto h(z)/|z|$, with value zero at zero, where
\[
 h(z)=(2\Re(z_1\bar z_2),\,2\Im(z_1\bar z_2),\,
                       |z_1|^2-|z_2|^2).
\]
It has precisely the circle orbits as fibres and norm $|z|$.
Adjoin $C\times D^3$ to the free quotient ${\mathcal Q}_0$ to obtain ${\mathcal Q}$.
The Hopf local sections
\[
 s_N(v)=\left(\sqrt{(1+v_3)/2},
       \frac{v_1-iv_2}{\sqrt{2(1+v_3)}}\right),\quad
 s_S(v)=\left(\frac{v_1+iv_2}{\sqrt{2(1-v_3)}},
       \sqrt{(1-v_3)/2}\right)
\]
give inverse local maps $(q,t)\mapsto |q|t\,s_N(q/|q|)$
and the southern version. Orient the meridian sphere $\sigma$
so that its Hopf Euler value is one.

Transversality in codimension four identifies
$\pi_1,\pi_2,H_2$ of $X_0$ with those of $X$. The bundle
homotopy sequence and Hurewicz give
\[
 0\longrightarrow H_2(X)\longrightarrow H_2({\mathcal Q}_0)
 \xrightarrow{\langle e,-\rangle}\mathbb Z\longrightarrow0,
 \qquad 1\longmapsto\sigma .
\]
Excision kills exactly $\mathbb Z\sigma$ in passing to ${\mathcal Q}$.
Thus $\Theta=i_*p_*j_*^{-1}:H_2(X)\to H_2({\mathcal Q})$ is an isomorphism.
Its inverse lifts a class to $a\in H_2({\mathcal Q}_0)$, subtracts
$\langle e,a\rangle\sigma$, lifts its representing sphere by a
section, and includes it into $X$. Changes of lift by $\sigma$
cancel, and section changes are homotopic. Van Kampen gives
$\pi_1({\mathcal Q})=0$. The residue of $e$ is one, so the integral residue
sequence gives $[C_{\mathcal Q}]=0$.
The mod-two Gysin kernel of $p^*$ in degree two is generated by
$e\bmod2$. This class evaluates to one on $\sigma$, whereas
$w_2({\mathcal Q}_0)$ evaluates to zero because
$T{\mathcal Q}_0|_\sigma=TS^2\oplus\mathbb R^3$. Spin of $X$ therefore gives
spin of ${\mathcal Q}_0$ and, by excision, of ${\mathcal Q}$.

Represent a free basis of $H_2(X_0)$ by spheres and project them
to ${\mathcal Q}_0$. General position in dimension five makes these maps
disjoint embeddings $V_i$. Lift the perturbation homotopies
through the principal bundle. This gives sections and proves
$e(V_i)=0$ on the full free basis, including its bridge classes.
The complement ${\mathcal Q}'={\mathcal Q}\setminus\bigcup V_i$ is simply connected.
In its excision sequence the intersection map
$H_3({\mathcal Q})\to\mathbb Z^7$ is onto because the $V_i$ form a primitive
free basis and the free duality pairing is unimodular. Hence
$H_2({\mathcal Q}')\to H_2({\mathcal Q})$ is an isomorphism, and $C_{\mathcal Q}$ is nullhomotopic
already in ${\mathcal Q}'$.

Barden's parametrized embedding theorem
\cite[Theorem B(b), pp.~367--368]{manBarden65} applies since ${\mathcal Q}'$
is connected and simply connected and $2\cdot5>3\cdot2+3$.
It isotopes this sphere into a five-disc disjoint from all $V_i$.
Extending its velocity gives an ambient flow; the backward flow
is its inverse. Transport the bundle by this diffeomorphism.
Its Euler class outside that disc is zero, since $H^2$ there is
the free dual of $H_2({\mathcal Q})$ and all values on the $V_i$ vanish.
On a good cover, remove the integral \v{C}ech cocycle of its
transition functions, and use a partition of unity on their
logarithms. Exponentiation gives a trivialization, with reciprocal
gauges as inverse.

Inside the disc, the residue-one bundle is the Hopf model just
parametrized. Oriented normal maps $S^2\to SO(3)$ lift to
$SU(2)$ and have unique homotopy class because $\pi_2(SO(3))=0$.
The two bundle maps differ on their collar by a gauge on
$S^4\times I$. A real logarithm and cutoff make them agree,
with reciprocal cutoff for the inverse. Their gluing is
\[
 X\cong M_-({\mathcal Q})=(({\mathcal Q}\setminus\mathring D^5)\times S^1)
             \cup_g(S^4\times D^2),\qquad
 g(x,z)=(\alpha_zx,z),
\]
where $g$ runs from complement to cap and $\alpha_z$ rotates
the first real two-plane. The inverse uses $\alpha_z^{-1}$.
The spin lift $\cos(\pi t)+\sin(\pi t)e_1e_2$ has endpoints
$1,-1$, fixing the nontrivial framing loop.
To compare the opposite gluing convention of
$\widetilde\Sigma_0{\mathcal Q}$, use
$A=\operatorname{diag}(1,-1,-1,1,1)$, so
$A\alpha_zA^{-1}=\alpha_z^{-1}$.
Rotate the $2,3$-plane by $\pi$ near the complement boundary,
cut this rotation off along its collar, and use $A$ on the cap.
The identity $A\alpha_z^{-1}=\alpha_zA$ makes these oriented
maps glue. Opposite collar rotation and $A^{-1}$ give the inverse.

We have $H_2({\mathcal Q})=\mathbb Z^7\oplus(\mathbb Z/m)^2$.
Its spin linking form is alternating by
\cite[Lemma F, p.~372]{manBarden65}. Choose $x,y$ of order $m$
with $\lambda(x,y)=1/m$. The inverse torsion coordinates are
\[
 a\longmapsto
 (\operatorname{num}_m\lambda(a,y),
  -\operatorname{num}_m\lambda(a,x)),\qquad
 (r,s)\longmapsto rx+sy .
\]
Nondegeneracy verifies the inverse identities. The oriented
classification \cite[Theorems 2.2--2.3, pp.~376--377]{manBarden65}
therefore gives ${\mathcal Q}\cong\#^7(S^2\times S^3)\#N_m$.
Parametrized disc isotopy makes this map match the removed disc;
its product with the identity extends over the common cap,
and the inverse base diffeomorphism gives the inverse.

The framed suspension distributes over this connected sum.
Use a rotationally symmetric dumbbell for the punctured connected
sum and radial projection to its round boundary. The separator in
the suspension is
$(D^4\times S^1)\cup_{\alpha|S^3}(S^3\times D^2)$.
The map $(x,z)\mapsto(\alpha_z^{-1}x,z)$ on the first part
and identity on the second identifies it with
$\partial(D^4\times D^2)=S^5$; inverse rotation gives its inverse.
Cutting and rejoining along these same collars gives the oriented
connected-sum map. Finally
\cite[Proposition 5.6(1), pp.~22--23]{manGR25}, with base $S^3$,
fibre $S^2$ and trivial clutching, gives
$\widetilde\Sigma_0(S^2\times S^3)\cong
 (S^2\times S^4)\#(S^3\times S^3)$.
Its product surgery and inverse clutching maps preserve these
orientations. This proves \eqref{maps:eq:D.O3}, including its
framing at even $m$.

The same maps also determine the integral degree-two product.
For ${\mathcal P}=({\mathcal Q}\setminus\mathring D^5)\times S^1$,
excision gives $H^k(M_-,{\mathcal P})=H^{k-2}(S^4)$.
The circle generator maps isomorphically from $H^1({\mathcal P})$ to
$H^2(M_-,{\mathcal P})$, so restriction is an isomorphism in degrees two
and four. Every degree-two class on ${\mathcal P}$ comes from the punctured
quotient, whose $H^4$ is zero. Hence every integral product of
two degree-two classes on $X$ is zero.

The auxiliary equivariant surgery can also be read from these maps.
The fixed sphere is a vertical image and is zero in $H_2(U)$;
it is nullhomotopic in the simply connected $X$.
Surgery in its displayed framed neighbourhood gives
$Y=X\#(S^3\times S^3)$ with orbit
${\mathcal Q}\#(S^2\times S^3)$. Its Euler class is zero on the old free
basis and one on the new meridian sphere. The spin structure of
the orbit, pulled back using a connection, gives the projectable
lift. Thus this orbit has
$H_2=\mathbb Z^8\oplus(\mathbb Z/m)^2$ and Euler pair
$(x,0,\ldots,0)$, and $Y$ has the corresponding additional
$S^3\times S^3$ factor.

\emph{The untwisted Engel member.}
At $(m,v)=(1,0)$ the lattice relation is $s=0$ and the log
gluing is the identity. The integral groups
are therefore $(\mathbb Z,\mathbb Z,\mathbb Z^8,\mathbb Z^{14},
\mathbb Z^8,\mathbb Z,\mathbb Z)$.
The $49/III$ row identifies with the model in
\SeriesRef{push:src:A.4.2} by the following maps.
Send its three bad-fibre positions to $0,1,\infty$ by
\[
 t(z)=
 \frac{\det(z,p_0)\det(p_1,p_\infty)}
      {\det(z,p_\infty)\det(p_1,p_0)}.
\]
The adjugate of its invertible two-by-two matrix gives the inverse.
The short Weierstrass degree bounds and fibre orders
\cite[\S\S2.3--2.6, pp.~4--5; Table 1, p.~14; \S4.10, p.~17]{manSS10}
give $A(t)=a t^3(t-1)$ and $B(t)=b t^4(t-1)^2$.
The $I_1$ at infinity implies $4a^3+27b^2=0$.
For $c^2=-3b/(2a)$, the maps
$(X,Y)\mapsto(c^2X,c^3Y)$ and $(x,y)\mapsto(c^{-2}x,c^{-3}y)$
give
\[
 y^2=x^3-3t^3(t-1)x+2t^4(t-1)^2,\qquad j=1728t .
\]
They extend through the resolution charts and preserve the zero
section. Choose the generator $P$ in this model to be the image
of $G$. Pullback of the divisor functions gives inverse line maps
between $\mathcal O(G-O)$ and $\mathcal O(P-O)$, conjugating
the translations by $6G$ and $6P$. The two degree-one
linearizations have the same simple zero; their scalar ratio is
absorbed in the contraction parameter.

In the cusp coordinates of that construction the two periods are
$(Z,W)\mapsto(c_1qZ,c_2W)$ and
$(Z,W)\mapsto(c_3Z,c'_4uqW)$.
Here $H_i=I$, the valuation lattice is $\mathbb Z^2$, and the
chosen $A_2$ diagonal is the edge $e_1-e_2$.
Take the multiplicative coordinates after the line map to be
these very $Z,W$, retaining their unit factors. The chart maps are
\[
 z=Zq^{-k},\quad v=q^{k+1}/Z,\quad
 x=W(qu)^{-l},\quad y=qu/x,
\]
with inverse monomials $Z=zq^k$, $W=x(qu)^l$ on the common torus.
They extend over the same cones, agree on overlaps and commute
with the periods. Thus they identify the filled cusps and glue
to the preceding original-cap maps. This identifies the named
untwisted member itself, and proves the claimed integral
specialization. Finally each member has nonzero free $H_2$.
A diffeomorphism to $S^6$ would induce inverse maps to
$H_2(S^6)=0$, which excludes that possibility.
\end{proof}

\Needspace{9\baselineskip}
\subsection{A mixed original and logarithmic pair}
\begin{theorem}[One original additive fibre and one finite logarithm]\hypertarget{tag-02KW}{}\label{maps:src:D.2.21}
{\normalfont\scriptsize tag 02KW}\par

Use the degree-one data $P=G,Q=dG$ and the simple zero at the
marked $I_1$ in the following five rows. Keep the indicated original
additive filling and replace the other additive fibre by its finite-good
logarithmic model of order $m=qk$, with invariant power vector
$n\ell+r\delta$, where $k>0$, $r\in\mathbb Z$ and $\gcd(n,m)=1$.
\[
\begin{array}{c|c|c|r|r|r|r}
\mathrm{OS/profile}&d&\mathrm{finite\ log}\ (q)&\mathrm{original}&b&h&c\\\hline
49/III&6&IV^*\ (3)&III&1&3&2\\
49/III&6&III\ (4)&IV^*&6&2&3\\
50/IV&12&IV\ (3)&I_1^*&5&3&4\\
43/II&2&II\ (6)&III^*&7&1&2\\
46/II&4&II\ (6)&I_3^*&7&1&4
\end{array}
\]
Then $n\ne0$, $\pi_1(X)=H_1(X)=\mathbb Z/|n|$, and
\[
 H_2(X)=\mathbb Z^b,\qquad H_3(X)=\mathbb Z^{2b},\qquad
 H_4(X)=\mathbb Z^b\oplus\mathbb Z/|n|,\qquad H_5(X)=0.
\]
The endpoint groups are $\mathbb Z$, and $\chi(X)=2$.
Every member with $n=\pm1$ is oriented diffeomorphic to
\[
                 \#^b(S^2\times S^4)\ \#\ \#^b(S^3\times S^3).
\]
In particular none of these five simply connected families is $S^6$.

\end{theorem}
\begin{proof}
Use the marking $\Lambda=\mathbb Z\langle e_1,e_2,\delta,\gamma\rangle$,
the exterior order $(12,13,14,23,24,34)$, and the period
matrices $T_i$ constructed in the proof of \ref{maps:src:D.2.20}.
In particular the positive $\gamma$ period covers $-dG$.
The original line and translation on the original additive
neighbourhood are carried by the logarithm, connection, roof
and inverse maps in that proof to $E(\mathcal O(G-O))\times S^1$.
Its primitive Euler row remains $e_p-e_0$.
The cusp is carried by the integral $H_i$ and its inverse,
with its period units and fan, to the primitive cusp whose
valuation lattice is $\mathbb Z^2$.

\emph{The finite cap in its integral marking.}
Let $A_f,g_f$ be the finite fibre's elliptic matrix and section
cocycle, and put $z_G=(A_f-I)^{-1}g_f$, $z_Q=dz_G$.
The fixed vector with contracting coordinate one and the
integral change of basis are
\[
\begin{gathered}{}
 \ell=(-z_{Q,1},-z_{Q,2},0,1)^t,\qquad
 B=[e_1,e_2,\delta,\ell],\qquad
\\{}
 B^{-1}(x_1,x_2,x_3,x_4)^t
 =(x_1+z_{Q,1}x_4,x_2+z_{Q,2}x_4,x_3,x_4)^t.
\end{gathered}
\]
Thus $\det B=1$ and
\[
 B^{-1}T_fB=\operatorname{diag}(R,1),\qquad
 R=\begin{pmatrix}A_f&0\\-g_f^tJA_f&1\end{pmatrix},\qquad R^q=I.
\]
The fixed lattice is exactly $\mathbb Z\delta\oplus\mathbb Z\ell$,
since $A_f-I$ is invertible over $\mathbb Q$.
Replacing $\ell$ by $\ell+t\delta$ changes the coordinates by
$(n,r)\mapsto(n,r-nt)$, with inverse $(n,r')\mapsto(n,r'+nt)$.

The invariant primitive row on $\mathbb Z^3$ is
$\rho=h(z_{G,2},-z_{G,1},1)$.
Put $\Omega=(1,0,z_{Q,2},0,-z_{Q,1},0)$ and
$\rho_\gamma=(0,0,\rho_1,0,\rho_2,h)$.
Their literal values are
\[
\begin{array}{c|r|r|r|r}
 \text{finite fibre}&\ell&\rho&\Omega&\rho_\gamma\\\hline
 49/IV^*&(-6,-2,0,1)&(1,-3,3)&(1,0,2,0,-6,0)&(0,0,1,0,-3,3)\\
 49/III&(-9,-3,0,1)&(1,-3,2)&(1,0,3,0,-9,0)&(0,0,1,0,-3,2)\\
 50/IV&(-8,8,0,1)&(-2,-2,3)&(1,0,-8,0,-8,0)&(0,0,-2,0,-2,3)\\
 43/II&(0,0,0,1)&(0,0,1)&(1,0,0,0,0,0)&(0,0,0,0,0,1)\\
 46/II&(0,0,0,1)&(0,0,1)&(1,0,0,0,0,0)&(0,0,0,0,0,1).
\end{array}
\]
Both $R-I$ and $\bigwedge^2R-I$ have two unit nonzero Smith
entries. The latter assertion also follows from
$\bigwedge^2\mathbb Z^3\cong(\mathbb Z^3)^*$, under which the
action is $R^{-t}$. Their coinvariants are primitive copies of
$\mathbb Z$, detected by $\rho$ and elliptic area respectively.

On the central universal cover $\mathbb R^3_y\times\mathbb R_t$,
the original lattice acts by $\mathbb Z^3$ translations and by
the unit $t$ translation $\ell$. The additional deck generator is
\[
 x(y,t)=(Ry+(r/m)\delta,t+n/m),\qquad x^m=\ell^n\delta^r.
\]
An equivariant connection, averaged over the finite deck group,
retracts the good-reduction cap to this central flat fibre and
retains the boundary specialization. The action is free: a
stabilizer would have $jn/m\in\mathbb Z$, hence $m\mid j$.
The same calculation shows that the physical $\delta$ circle is
free on the quotient, since a pure angular displacement cannot
compensate a nonintegral $t$ displacement.

Choose integers $A,B_0$ with $An+B_0m=1$ and put
$z=x^A\ell^{B_0}$. The commuting translations give
\[
 x=\delta^{rB_0}z^n,\qquad \ell=\delta^{-rA}z^m.
\]
The coordinate change
$t'=mt$, $y'=y-(Ar/m)t'\delta$, with inverse
$t=t'/m$, $y=y'+(Ar/m)t'\delta$, identifies the central fibre
with the mapping torus of $R^A$:
$z(y',t')=(R^Ay',t'+1)$.
It preserves orientation. Since $An\equiv1\pmod q$,
$R^A$ and $R$ generate the same cyclic group and have equal
integral invariant and coinvariant lattices.
The Wang sequences give
\[
 H_1(V)=\mathbb Z u\oplus\mathbb Z z,\qquad
 H_2(V)=\mathbb Z A_2\oplus\mathbb Z B_2.
\]
Here $\rho$ reads $u$, $\delta=hu$, $A_2$ is the primitive
area coinvariant, and $B_2$ is the torus $\delta\times z$.
B\'ezout preimages of the primitive rows give inverse coordinate
maps, while the fixed $\delta$ circle splits the second Wang
sequence.

For a fibre loop $a$, its cap coordinates are
\[
 f_1(a)=
 (\rho_1a_1+\rho_2a_2+ha_3-rAh\,a_4,\;ma_4),
 \qquad f_1(\mu)=(rB_0h,n).
\]
For a fibre two-class $a$ the map is
\[
 f_2(a)=\Omega(a)A_2+(m/h)\rho_\gamma(a)B_2.
\]
Indeed elliptic area pulls back to $\Omega$, and the invariant
form $(\rho\,dy/h)\wedge dt'$ is one on $\delta\times z$ and
zero on the area generator. Its pullback is
$(m/h)\rho_\gamma$, since $\rho(z_Q,0)=0$.
The affine gauge adds only a multiple of $dt'$.
The two remaining boundary generators are the meridian sweeps,
with maps
\[
 \delta\times\mu\longmapsto nB_2,\qquad
 \ell\times\mu\longmapsto-rB_2.
\]
Their area values vanish in the $B$ coordinates. These formulas
give the entire boundary map, including every integer $r$.

\emph{The complement and its cyclic class.}
Let $U$ be the complement of the finite-log disc. Its base is a
boundary connected sum of the original disc and the
cusp disc, beginning immediately after the removed fibre.
The two fibre maps in degrees one and two are those just
established for the original caps and primitive cusp in
\ref{maps:src:D.2.20}. Put
$M_j=(f_{j,1},-f_{j,2})^t$.
Van Kampen kills $e_1,e_2,\delta$ in the original cap and gives
$\pi_1(U)=\mathbb Z\gamma$. The meridian has zero contracting
coordinate by the fractional-roof section and is nullhomotopic
there. Sweeping a bounding disc proves that its $\delta$ rim
is zero.

The Mayer--Vietoris sequence gives
\[
 0\longrightarrow\operatorname{coker}M_2\longrightarrow H_2(U)
 \longrightarrow\ker M_1\longrightarrow0,\qquad
 \operatorname{coker}M_2=\mathbb Z^{b-1}\oplus\mathbb Z/h,\quad
 \ker M_1=\mathbb Z^2.
\]
The five nonzero Smith entries of $M_2$ are $1,1,1,1,h$.
Integer row and column elimination using the component basis
$F,C_j$ and the cusp maps gives the following vertical row
$\lambda$, where the image is $\lambda(a)\sigma$ in
$\mathbb Z/h\langle\sigma\rangle$ and has zero free component.
The last column gives the $\ell$ rim in the two kernel
coordinates.
\[
\begin{array}{c|c|r|r|r|r}
 \text{removed fibre}&\text{cap order}&b&h&\lambda&\text{rim}\\\hline
 49/IV^*&III,\text{cusp}&1&3&(1,0,3,0,-9,0)&(2,-1)\\
 49/III&\text{cusp},IV^*&6&2&(1,0,2,-2,-6,2)&(-3,1)\\
 50/IV&\text{cusp},I_1^*&5&3&(1,0,-6,0,-9,3)&(4,5)\\
 43/II&\text{cusp},III^*&7&1&(-1,0,-1,0,3,1)&(2,-1)\\
 46/II&\text{cusp},I_3^*&7&1&(-1,0,2,0,1,1)&(-4,-3).
\end{array}
\]
For $h=1$ keep the unit relation on $\sigma$ in this presentation.
The kernel basis is $(e_1,\delta)$ in the first, second and
fourth rows, and $(e_1+e_2,\delta)$ in the third and fifth.
The respective inverses read $(x_1,x_3)$ and $(x_2,x_3)$.
If $T_a,T_b$ are the two remaining matrices, the connecting
column of the rim is $(T_a-I)T_b\ell$. The displayed pairs
are primitive, with determinant-one complements
$(1,0),(-1,0),(-1,-1),(1,0),(-1,-1)$ respectively.
Their adjugates give inverse kernel coordinates.
Choose the rim $w$, including its component-chain part,
as one free lift; choose $z_0$ as a lift of the complementary
kernel vector. Thus the precise exterior coordinates are
\[
 H_2(U)=\mathbb Z^{b-1}\oplus
       \mathbb Z/h\langle\sigma\rangle\oplus\mathbb Zw\oplus\mathbb Zz_0.
\]
To invert, read the connecting coordinates, subtract those
lifts, and read the remaining component coordinates.
This retains the cyclic class until its integral gluing relation
is imposed.

\emph{The complete gluing for unbounded parameters.}
The finite-cap fundamental group is
$\mathbb Z^3\rtimes_{R^A}\mathbb Zz$.
The original cap kills its $\mathbb Z^3$ subgroup. The meridian
gives $nz=0$, and the original contracting loop is $\gamma=mz$.
Thus $\pi_1(X)=\mathbb Z/|n|$. For $tm+sn=1$ the inverse
generator map is $z=t\gamma$. Also $n\ne0$ because $m=qk\ge3$
and $\gcd(n,m)=1$.
In the boundary basis $u,\ell,\mu$, the degree-one map to
$H_1(U)\oplus H_1(V)$ is
\[
 \begin{pmatrix}
 0&1&0\\1&-rAh&rB_0h\\0&m&n
 \end{pmatrix}.
\]
Its Smith factors are $1,1,|n|$ and its kernel is zero.

The unaffected $\mathbb Z^{b-1}$ and $z_0$ remain in degree two.
The other generators $\sigma,w,A_2,B_2$ have relations
\[
 \begin{gathered}
 h\sigma=0,\qquad
 \lambda(a)\sigma-\Omega(a)A_2-(m/h)\rho_\gamma(a)B_2=0
       \quad(a\in\bigwedge^2\Lambda),\\
 nB_2=0,\qquad w-rB_2=0 .
 \end{gathered}
\]
Here $w$ is given the exterior sweep orientation opposite to
the cap convention above; replacing $w$ by $-w$ gives the
inverse change of convention.
The $e_{12}$ relation eliminates $A_2=\lambda_{12}\sigma$.
For $\epsilon=1,1,-1,1,1$ in the five rows, direct subtraction
of the displayed rows gives
\[
 \lambda-\lambda_{12}\Omega=\epsilon\rho_\gamma+h\kappa
 \quad\text{with }\kappa\in\mathbb Z^6.
\]
The row $\rho_\gamma$ is primitive: use coordinate $14$ in the
first two rows, the sum of $14$ and $34$ in the third, and
$34$ in the last two. The fibre relations are therefore
equivalent, by these explicit integer combinations, to
$\epsilon\sigma-(m/h)B_2=0$.
The remaining presentation is
\[
 \langle\sigma,B_2\mid
       h\sigma,\ \epsilon\sigma-(m/h)B_2,\ nB_2\rangle.
\]
Multiplying the second relation by $h$ and subtracting
$\epsilon$ times the first gives $-mB_2$. The combination
\[
 A(nB_2)-B_0
   \{h[\epsilon\sigma-(m/h)B_2]-\epsilon h\sigma\}=B_2
\]
then kills $B_2$, since $An+B_0m=1$. The second relation
kills $\sigma$, and the earlier equations kill $A_2,w$.
The inverse free-coordinate map performs these same subtractions
and reads the $b-1$ component coordinates and $z_0$.
This proves $H_2(X)=\mathbb Z^b$ for every stated integer input,
and proves that $H_2(U)\to H_2(X)$ is onto.

On the full regular-base graph, the integral elimination has the
following coordinates. In cyclic finite-first order its differential
is $[T_0-I,T_1-I]$, with
$j_0,j_1$ the two inclusions and $j_2(a)=-(T_1a,a)^t$.
The two remaining fixed lattices give four columns $B_{\mathrm{fix}}$;
$[B_{\mathrm{fix}},j_0\ell]$ is a primitive basis of its rank-five
kernel. Its last inverse row kills $j_0\delta$ and takes
$j_0\ell$ to one. Choose the vertical coinvariant representatives
$\beta=e_1,e_1,e_1+\delta,\delta,\delta$ in the five rows.
Then $\rho(\beta)=1$, $\delta=h\beta$ in coinvariants, and
\[
 N_m\beta=(m/h)\delta,\quad N_m\ell=m\ell,\quad
 -\omega(\beta,\ell)=d/h=c,\quad-\omega(\ell,\delta)=-d,
 \qquad \omega=(1,0,0,0,0,-d).
\]
These identities follow by summing $I+T_f+\cdots+T_f^{q-1}$
and multiplying by $k$. In boundary coordinates consisting of
$b_0=\beta\wedge\ell$ and the two meridian sweeps, the finite map
is the row $(m/h,n,-r)$. Its kernel has basis
\[
 -nb_0+(m/h)s_\delta,\qquad
 B_0rh\,b_0+Ar\,s_\delta+s_\ell.
\]
Subtract the second coefficient to remove $s_\ell$; the inverse
then uses $\gcd(n,m/h)=1$ on the first two coordinates.
Cutting the swept normal disc into $m$ sectors identifies these
classes with the transfer columns
$(-\omega(\beta,v),j_0N_m\beta)$ and
$(-B_0\omega(\ell,v),j_0(B_0N_m\ell+Av))$.
After the four unchanged columns and the original fibre kernel
are eliminated, their matrix is
\[
 \begin{pmatrix}h&cn&-drB_0\\0&0&1\end{pmatrix}.
\]
Since $\gcd(h,c)=1$ and every prime of $h$ divides $m$,
$\gcd(h,cn)=1$; the two coordinates are eliminated by B\'ezout.
Joining the two unchanged caps first and cutting them apart
again are inverse Mayer--Vietoris subdivision maps, preserving
the rim and its component-chain part. These inverse subdivision maps
identify the cyclic relation in the two presentations.
The full three-boundary degree-one map has ranks nine and eight:
the regular exterior has rank three and the three caps have
ranks two, one and two. Its kernel is one free bridge class,
already included as $z_0$ in the two-piece calculation.

All caps except the primitive cusp have a free physical circle.
The cusp has the one fixed sphere described in
\ref{maps:src:D.2.20}, so $\chi(X)=2$.
Integral duality and universal coefficients give the stated
$H_4,H_5$, while the Euler identity gives $b_3=2b$.
The torsion character $z\mapsto1/|n|$, followed by Bockstein
and cap product, gives the $H_4$ cyclic coordinate, inverted
by linking. Free $H_4$ coordinates are dual to the displayed
$H_2$ coordinates. The middle group is torsion-free and its
integral alternating intersection pairs give inverse cycle
coordinates. A point and the complex fundamental class generate
the endpoint groups.

\emph{Canonical class and oriented smooth type.}
Let $F_{\mathrm{red}}$ be the reduced finite core. Its complement
retracts onto $U$. The primitive positive normal meridian
$m\mu-n\ell-r\delta$ maps to $-n\gamma$ there.
In the cohomological Thom sequence
$H^1(U)\to H^2(X,U)\to H^2(X)$, the first map is multiplication
by $-n$, and the second sends one to $\operatorname{PD}(F_{\mathrm{red}})$.
The divisor class has exact order $|n|$.
Choose the surface canonical form with its pole on the removed
whole fibre, using \cite[Theorem 6.8, p.~26]{manSS10}.
Its product with $dw/w$ is invariant under the original line
quotient and extends across the original caps by the node and
cusp chart identities in the proof of \ref{maps:src:D.2.20}.

At good reduction put $q=u^m$, $x=u^{2L}X$, $y=u^{3L}Y$,
where $L=2k,k,k,k$ for $IV^*,III,IV,II$ respectively.
Substitution in the coefficient orders gives these minimal
scalings, and the meromorphic top form pulls back, up to a
nonzero unit, to
\[
 m u^{-L-1}\,du\wedge(dX/Y)\wedge(dw'/w').
\]
Derivative terms in the coordinate and line changes vanish
in this top wedge. The root action has weights
$\zeta^{-L-1},\zeta,\zeta^L$ on the first three factors,
so their product is invariant; translations preserve $dX/Y$.
The total quotient is free by its $t$-displacement calculation.
It is therefore locally biholomorphic and
$K_X=\mathcal O(-(L+1)F_{\mathrm{red}})$.
For $n=\pm1$ the integral first Chern class is zero, including
on the bridge generator, and $X$ is spin.

For those signs the action is semifree, with exactly the one
fixed sphere $C$. Its normal model, with inverse charts
and the $SU(2)$ contraction given in \ref{maps:src:D.2.20},
is $\mathcal O(-1)_{-1}\oplus\mathcal O(-1)_{+1}$, hence
the oriented framed scalar model $C\times D^4$.
The cusp divisor is the self-identified whole fibre:
its class is the pullback of a base point and its square and
cube vanish. Also $C$ is a vertical image by the surjective
cusp specialization. Its exterior image is a multiple of
$\sigma$, which the integral gluing relations kill. Hurewicz
therefore makes $C$ nullhomotopic in these simply connected
members.

We apply the direct orbit construction with the maps described
here. Remove $C\times\mathring D^4$, take the free quotient
$Q_0$, and adjoin $C\times D^3$ by the Hopf map
$z\mapsto(2\Re(z_1\bar z_2),2\Im(z_1\bar z_2),
|z_1|^2-|z_2|^2)/|z|$. Its local inverse charts are the Hopf
sections in the preceding proof, multiplied by the orbit radius.
The bundle sequence is split by the meridian sphere of Euler
value one. Excision quotients by that sphere and gives
$H_2(X)\cong H_2(Q)$, with inverse obtained by subtracting
the Euler value from any lift and taking a section of the
resulting Euler-zero sphere. Van Kampen gives $\pi_1(Q)=0$.
The Euler residue one gives $[C_Q]=0$. The mod-two Gysin
kernel is generated by $e\bmod2$, which is one on the
meridian, whereas $w_2(Q_0)$ is zero on it. Spin of $X$
therefore proves spin of $Q$.

Represent the free basis of $H_2(X)$ by spheres outside $C$,
project them to $Q_0$, and perturb to disjoint embeddings.
Lifting the perturbation homotopies gives sections on them,
so their Euler values are all zero. Their complement is
simply connected; its $H_2$ maps isomorphically onto $H_2(Q)$,
since the intersection map from $H_3(Q)$ to the meridians
of this primitive free basis is onto. Thus $C_Q$ is
nullhomotopic in that complement. By
\cite[Theorem B(b), pp.~367--368]{manBarden65}, it is isotopic
there into a five-disc, since $10>9$.
Ambient forward and backward flows give inverse maps.
Outside that disc the Euler class is zero on a free basis,
hence is zero. Logarithms of transition functions and a
partition of unity give a bundle trivialization and its
reciprocal inverse.

The residue-one bundle inside the disc is the Hopf model.
Its oriented normal framing is unique since $\pi_2(SO(3))=0$,
and lifts to $SU(2)$. A logarithm of the boundary gauge
on $S^4\times I$ matches the two maps on a collar.
We obtain directly
\[
 X\cong((Q\setminus\mathring D^5)\times S^1)
      \cup_g(S^4\times D^2),\qquad g(x,z)=(\alpha_zx,z).
\]
The inverse rotates by $\alpha_z^{-1}$.
The spin lift of $\alpha$ has endpoints $1,-1$, and
$\operatorname{diag}(1,-1,-1,1,1)$ conjugates $\alpha$ to
$\alpha^{-1}$. Its cutoff collar rotation identifies this
gluing, with inverse rotation, with $\widetilde\Sigma_0Q$.
This fixes the same oriented framing as
\ref{maps:src:D.2.14}.

The quotient $Q$ is simply connected and spin with $H_2(Q)=\mathbb Z^b$.
The oriented classification
\cite[Theorems 2.2--2.3, pp.~376--377]{manBarden65} gives
$Q\cong\#^b(S^2\times S^3)$.
Match the removed disc by parametrized disc isotopy.
On each connected-sum neck the separator in the suspension is
$(D^4\times S^1)\cup_\alpha(S^3\times D^2)$.
Rotating the first piece by $\alpha_z^{-1}$ and fixing the
second identifies it with $S^5$, with inverse rotation.
Cutting and rejoining distributes the precise suspension.
The product surgery of
\cite[Proposition 5.6(1), pp.~22--23]{manGR25}, with base $S^3$,
fibre $S^2$ and trivial clutching, then gives
\[
 X\cong\#^b(S^2\times S^4)\#\#^b(S^3\times S^3)
\]
with the complex orientation.
Restriction to $(Q\setminus\mathring D^5)\times S^1$ is an
isomorphism in cohomological degrees two and four: the
relative groups are $H^{*-2}(S^4)$ and the boundary circle
maps to the relative degree-two generator with coefficient
one. Every degree-two class comes from the punctured quotient,
whose $H^4$ is zero. Thus all integral degree-two products
vanish on the constructed $X$ as well.

Finally the framed equivariant surgery on $C$ gives
$X\#(S^3\times S^3)$ with orbit
$Q\#(S^2\times S^3)$ and Euler class one on the new sphere
and zero on the old free basis. Pulling back the orbit spin
structure by a connection gives the projectable lift.
These maps verify the stabilized orbit description directly.
Each value of $b$ is positive, so a diffeomorphism to $S^6$
would induce inverse maps between $\mathbb Z^b$ and zero in
degree two, which is impossible.
\end{proof}


\Needspace{12\baselineskip}
\section{Sphere fillings and primitive finite caps}
\subsection{The \texorpdfstring{$I_3^*$}{I3*} homology obstruction}
\begin{theorem}[Homology obstruction for the two $I_3^*$ profiles]\hypertarget{tag-02KZ}{}\label{maps:src:D.2.24}
{\normalfont\scriptsize tag 02KZ}\par
Neither OS profile $46$ nor $46/II$ gives a homotopy six-sphere
in the natural section-line construction and filling sector of
\ref{maps:src:D.2.22}, for any admissible logarithmic orders or
affine twists. This includes every permitted original or semistable
logarithmic filling at $I_3^*$. Each completion of the eight
necessary initial rows in that catalogue has positive second
rational Betti number.
\end{theorem}


\begin{proof}
The necessity and catalogue in \ref{maps:src:D.2.22} reduce a
homotopy six-sphere in these profiles to
$(a,b)=(1,4),(4,1)$ at each of the three $I_1$ positions of
profile $46$ and the single $I_1$ position of profile $46/II$.
Indeed $ab=4$ and the natural component equation at $I_3^*$,
whose component group is $\mathbb Z/4$, requires $a$ or $b$
to be divisible by four. These are the eight marked rows.

In the physical period order $(e_1,e_2,\delta,\gamma)$,
their two local matrices, according to $a=1,4$, are
\[
 T_1=\begin{pmatrix}
 -1&-3&0&4\\0&-1&0&4\\-1&-2&1&3\\0&0&0&1
 \end{pmatrix},\qquad
 T_4=\begin{pmatrix}
 -1&-3&0&1\\0&-1&0&1\\-4&-8&1&3\\0&0&0&1
 \end{pmatrix}.
\]
Put $w_1=2e_1+\delta$ and $w_4=e_1+2\delta$.
The integral bases
\[
 H_1=(w_1,e_2,-e_1,2e_2+\gamma),\qquad
 H_4=(w_4,-\gamma,\delta,e_2+2\gamma)
\]
have inverses
\[
 H_1^{-1}=\begin{pmatrix}
 0&0&1&0\\0&1&0&-2\\-1&0&2&0\\0&0&0&1
 \end{pmatrix},\qquad
 H_4^{-1}=\begin{pmatrix}
 1&0&0&0\\0&2&0&-1\\-2&0&1&0\\0&1&0&0
 \end{pmatrix}.
\]
Literal multiplication, with $T$ either matrix, gives
\begin{equation}
 N=T^2-I,\qquad N^2=0,\qquad\operatorname{rank}N=1,
 \qquad\operatorname{Smith}(N)=(3).
 \hypertarget{tag-0207}{}\label{maps:eq:D.IS1}
\end{equation}
More precisely $H_a^{-1}NH_a=3E_{12}$, and
\begin{equation}
 TN=-N,\qquad T=-1\text{ on }\Lambda/\ker N,\qquad
 T=1\text{ on }(\ker N/\operatorname{sat}\operatorname{im}N)
                       \otimes\mathbb Q.
 \hypertarget{tag-0208}{}\label{maps:eq:D.IS2}
\end{equation}
The coordinate on the first quotient is the second $H_a$
coordinate. The last quotient has the third and fourth coordinates,
with inverse their basis vectors modulo the first. The first basis
vector is the primitive vanishing circle and $Tw_a=-w_a$.
The exterior-square matrices also give
\[
 \dim\ker(T-I)=2,\qquad
 \dim\operatorname{coker}(\bigwedge^2T-I)=2.
\]
Thus the boundary has second rational Betti number four.

Consider first a logarithmic cap of order $2k$. Its semistable
cover has $H_a^{-1}T^{2k}H_a=I+3kE_{12}$. In the primitive
toric coordinates its charts and inverse maps are
\[
 X=z/s^j,\quad Y=s^{j+1}/z,\quad XY=s;
 \qquad z=Xs^j,\quad s=XY.
\]
The period sends $j$ to $j+3k$. The central fibre is therefore
a cycle $C_{3k}$ of ruled components over an elliptic curve.
The local maps of \ref{maps:src:D.2.23} remove flat unit factors
by multiplication by their chosen logarithmic exponentials and
remove the gluing translations by subtraction of their interpolated
real lifts. Reciprocal multiplication and addition are the inverse
maps. They identify the smooth core with $C_{3k}\times T^2$.
The cap retracts to this core.

Write $P_j$ for its ruling-sphere classes, $F=\sum_jP_j$,
$\alpha$ for its oriented component-cycle loop, and $u,v$
for the transverse elliptic circles. The product cycles give
\[
 H_2(\widehat U;\mathbb Q)=
 \mathbb Q\langle P_0,\ldots,P_{3k-1}\rangle
 \oplus\mathbb Q\langle\alpha\times u,\alpha\times v\rangle
 \oplus\mathbb Q\langle u\times v\rangle .
\]
Projection and evaluation on complementary factors give the inverse
coordinates. Pinching the primitive circle in $XY=s$ sends the
nearby elliptic area to $F$, with degree one on every component,
and its long circle to $\alpha$, with degree one on the polygon.
Its products with $u,v$ and the transverse area give the four
independent fibre-specialization images.

The boundary Wang sequence has these four fibre-coinvariant classes
and three fixed-circle sweeps. The sweeps bound the vanishing-disc
family and the two products of the transverse circles with a section
disc. Cutting at the same toric chart and subtracting these bounding
chains gives the inverse quotient map. Hence the full boundary image
and its quotient are
\[
 \operatorname{im}H_2(\partial\widehat U;\mathbb Q)
 =\mathbb Q\langle F,\alpha\times u,\alpha\times v,u\times v\rangle,
 \qquad
 \frac{H_2(\widehat U;\mathbb Q)}
      {\operatorname{im}H_2(\partial\widehat U;\mathbb Q)}
 =\frac{\mathbb Q^{3k}}{\mathbb Q F}.
\]

The cyclic deck generator reverses the primitive vanishing angle.
Write $u_{\rm unit}(s)$ for its nonvanishing translated unit.
On the toric character it has the form
$z\mapsto u_{\rm unit}(s)s^c/z$, together with its base rotation, where
preservation of the integral rays forces $c\in\mathbb Z$.
Its inverse uses the inverse base rotation and the reciprocal
translated unit. The induced component action is
\[
                         j\longmapsto c-j\pmod{3k}.
\]
Affine translations and units alter the offset and the maps within
components while preserving this slope and each complex ruling
orientation. On the displayed quotient this is precisely the
permutation action. Let $\rho$ be its number of orbits.
The fixed components solve $2j=c\pmod{3k}$. Thus $\rho=(3k+1)/2$
for odd $3k$, and $\rho=3k/2+1$ or $3k/2$ for even $3k$.
In every case $\rho\ge2$.

The action on the cap is free. Transfer and $(2k)^{-1}$ times
projection are inverse maps between quotient homology and invariant
cover homology, commuting with boundary inclusion. Averaging over
the finite group is exact over $\mathbb Q$. For component orbits
$O_1,\ldots,O_\rho$, put $E_i=\sum_{j\in O_i}P_j$.
The map
\[
 \mathbb Q^\rho/\mathbb Q(1,\ldots,1)
 \longrightarrow
 \left(\mathbb Q^{3k}/\mathbb Q F\right)^{\rm deck},
 \qquad (x_i)\longmapsto\sum_i x_iE_i
\]
is an isomorphism. Its inverse reads one coefficient on each orbit,
modulo their common value. Therefore
\[
 \dim\operatorname{coker}
 [H_2(\partial U;\mathbb Q)\longrightarrow H_2(U;\mathbb Q)]
                           =\rho-1\ge1.
\]
A toric refinement has a degree-one map $f:(U',L')\to(U,L)$,
equal to the identity on the boundary. Its homological push-forward
has right inverse
$\operatorname{PD}_{U'}f^*\operatorname{PD}_{U}^{-1}$,
using relative cohomology and the oriented relative fundamental
classes. The projection formula gives $f_*f^!=1$.
This map commutes with the deck action and also applies on the
free quotient. Consequently the refined local cokernel surjects
onto the displayed one and retains its positive dimension.

For the original cap with $(a,b)=(4,1)$, the line multidegree is
zero. Label the affine $D_7$ outer components met by $O,G,2G,3G$
as $0,6,1,7$, and its interior chain as $2,3,4,5$.
The edges are
\[
 (0,2),(1,2),(2,3),(3,4),(4,5),(5,6),(5,7),
\]
and translation by $G$ is $(0\ 6\ 1\ 7)(2\ 5)(3\ 4)$.
The multiplicity vector is $(1,1,2,2,2,2,1,1)$.
The intersection matrix has diagonal $-2$ and entry one on these
edges; it is negative semidefinite with exactly this radical.
On the three orbit sums its rank is two. Indeed a vector with
zero quadratic value in this subspace already lies in the full
radical. Inclusion of the component classes, followed by the
relative map and pairing with their component four-cycles, realizes
this restricted form. In the finite cover fixing the components
it is the ordinary plumbing pairing; transfer descends it.
Boundary classes have zero relative image. Thus this original
cap has a local degree-two cokernel of dimension at least two.

For the original cap with $(a,b)=(1,4)$, the Euler evaluation
of the punctured line on the component classes is
\[
            \mathbb Z^8\longrightarrow\mathbb Z,
            \qquad x\longmapsto x_6-x_0.
\]
It has right inverse $1\mapsto e_6$. The plumbing retracts to
the tree of eight spheres. The circle-bundle sequence gives
$H_1(E;\mathbb Q)=0$ and identifies $H_2(E;\mathbb Q)$ by
projection with the seven-dimensional Euler kernel. Its basis
is $e_0+e_6,e_1,e_2,e_3,e_4,e_5,e_7$; the inverse forgets
coordinate six. Translation by $4G$ fixes every component class,
so its action on this kernel is the identity. The Wang sequence
of the contracting quotient gives $H_2(U;\mathbb Q)=\mathbb Q^7$.
Its boundary has second Betti number four by the literal matrix
calculation above. Its local cokernel therefore has dimension
at least three.

Finally let $E_0$ be the exterior and $U_i$ all the attached caps.
The degree-two Mayer--Vietoris sequence injects
\[
 \operatorname{coker}\left[
 \bigoplus_i H_2(\partial U_i;\mathbb Q)
 \longrightarrow H_2(E_0;\mathbb Q)\oplus
                         \bigoplus_iH_2(U_i;\mathbb Q)\right]
 \longrightarrow H_2(X;\mathbb Q).
\]
Projection to the star-cap factor surjects onto its local cokernel.
Each permitted star cap has a positive such cokernel by the preceding
calculations. This proves $b_2(X)>0$ for every choice at every other
site, including the additional type $II$ site, regular logarithms,
and the remaining multiplicative fibres.
\end{proof}


\subsection{All-multiplicative profiles}
\begin{theorem}[Integral homology and spheres for multiplicative profiles]\hypertarget{tag-02L0}{}\label{maps:src:D.2.25}
{\normalfont\scriptsize tag 02L0}\par


Consider any of the $62$ natural initial rows of \SeriesRef{maps:src:D.2.22} on the four
all-multiplicative profiles $45,47,55,56$. Write $P=aG,Q=bG$, with
$a,b>0$. Use the primitive Mumford model at each multiplicative fibre.
An original zero-free filling with zero line multidegree and one
component orbit may be used in its place; it has the same marked
primitive rank-one model. Perform any finite number of regular
logarithmic modifications, with data
\begin{equation}
 m_i>0,\qquad v_i=(x_i,y_i,z_i,n_i)\in\mathbb Z^4,
 \qquad\gcd(m_i,x_i,y_i,z_i,n_i)=1.                           \hypertarget{tag-023Z}{}\label{maps:eq:D.SC1}
\end{equation}
The coordinates are $(e_1,e_2,\delta,\gamma)$. Put
$\omega=a e_1^*\wedge e_2^*+b\gamma^*\wedge\delta^*$. The initial rows have
$\gcd(a,b)=1$. Retain the peripheral integer $e$ of \SeriesRef{maps:src:D.2.23};
one may keep every order-one site and use the canonical value $e=0$.

The integral homology is determined as follows for all the stated
orders, numbers of sites and translation coordinates. Let ${\mathsf S_{\rm Sei}}$ be the
Seifert relation matrix \eqref{maps:eq:D.FG2}, and write
\[
 \operatorname{coker}{\mathsf S_{\rm Sei}}=\mathbb Z^r\oplus A_1,\qquad r\in\{0,1\}.
\]
Form an integral presentation $\mathcal C$ with generators
$s$ and $t_i(u)$, linear in $u\in\Lambda=\mathbb Z^4$, and relations
\begin{equation}
 \begin{split}
 m_i t_i(u)&=\omega(u,v_i)s\quad(u=e_1,e_2,\delta,\gamma),\\
 t_i(v_i)&=0,\\
 \sum_i t_i(\delta)+be\,s&=0.
 \end{split}                                                \hypertarget{tag-0240}{}\label{maps:eq:D.SC2}
\end{equation}
Write $T=\operatorname{Tor}\operatorname{coker}\mathcal C$.
Then $\operatorname{coker}\mathcal C$ has free rank $r$ and
\begin{equation}
\begin{array}{c|c}
 j&H_j(X;\mathbb Z)\\\hline
 0,6&\mathbb Z\\
 1&\mathbb Z^r\oplus A_1\\
 2,3&\mathbb Z^{2r}\oplus T\\
 4&\mathbb Z^{2r}\oplus A_1\\
 5&\mathbb Z^r.
\end{array}                                                 \hypertarget{tag-0241}{}\label{maps:eq:D.SC3}
\end{equation}
The fundamental group is given by \eqref{maps:eq:D.FG1}. Smith reduction of the explicit matrix
\eqref{maps:eq:D.SC2} computes $T$, including all extensions.

For a simply connected member, let
\[
 M=\prod_i m_i,\quad
 X_0=\sum_i(M/m_i)x_i,\quad Y_0=\sum_i(M/m_i)y_i,\quad
 g=\gcd(b,M,X_0,Y_0).
\]
For no displayed sites use $M=1$ and $X_0=Y_0=0$.
Then
\begin{equation}
 H_2(X)=H_3(X)=
 \mathbb Z/g\oplus\mathbb Z/M\oplus\mathbb Z/(bM/g).          \hypertarget{tag-0242}{}\label{maps:eq:D.SC4}
\end{equation}
Consequently a sphere occurs precisely when the group criterion
\eqref{maps:eq:D.FG3} holds and $b=M=1$. In the canonical marking with all sites
retained, this says that every regular logarithmic order is one and
$|\sum_i n_i|=1$.

There are exactly $17$ permitted initial rows with $b=1$:
\begin{equation}
\begin{array}{c|c|c|c}
 \text{profile}&\text{zero fibre}&(a,b)&\text{marked positions}\\\hline
 45&I_1&(8,1)&4\\
 47&I_2&(7,1)&1\\
 47&I_1&(14,1)&3\\
 55&I_4&(5,1)&1\\
 55&I_1&(20,1)&3\\
 56&I_5&(6,1)&1\\
 56&I_3&(10,1)&1\\
 56&I_2&(15,1)&1\\
 56&I_1&(30,1)&2.
\end{array}                                                 \hypertarget{tag-0243}{}\label{maps:eq:D.SC5}
\end{equation}
This is the complete natural sphere table for these four profiles,
including the original and Mumford choices. At a zero-free
multiplicative fibre, every other original filling in these natural
rows gives a free direct summand in \(H_2(X;\mathbb Z)\).

More generally, every simply connected threefold in these $17$ rows,
with arbitrary data \eqref{maps:eq:D.SC1}, is oriented diffeomorphic to
\begin{equation}
                              \widetilde\Sigma_0N_M,         \hypertarget{tag-0244}{}\label{maps:eq:D.SC6}
\end{equation}
with the twisted framing and the Smale five-manifold of \SeriesRef{maps:src:D.2.14}.
Thus its full smooth type depends only on $M$, including when $M$
is even. For $b>1$, \eqref{maps:eq:D.SC2}--\eqref{maps:eq:D.SC3} give the complete integral
homology, with \eqref{maps:eq:D.SC4} in the simply connected cases.
Every such threefold is excluded by the sphere criterion.

\end{theorem}


\begin{proof}
\medskip\noindent\emph{Period and orientation conventions.}

Put

\[
 \Lambda=\mathbb Z\langle e_1,e_2,\delta,\gamma\rangle,
 \quad \chi_{\rm ret}=\gamma^*,\quad
 J=\begin{pmatrix}0&1\\-1&0\end{pmatrix},\quad
 \omega=a e_1^*\wedge e_2^*-b\delta^*\wedge\gamma^*.
\]

Thus \(\omega(\delta,\gamma)=-b\). The positive forward period
\(\gamma\) covers translation by \(-Q\), as in the Poincar\'e-line marking of \SeriesRef{maps:src:D.2.23}. The component coordinates used below are the negatives
of the geometric component coordinates; multiplication by minus one,
which is its own inverse, carries those conventions exactly.

The four marked default profiles have the following local data.
Every entry is an ordered pair \((u,\kappa)\), with
\(A_{\rm ell}=I+\nu uu^tJ\), \(g_{\rm ell}=\kappa u\), and fibre \(I_\nu\).

\[
\begin{array}{c|r|c|c}
\mathrm{OS}&d&(\nu_i)_{i=1}^5&(\kappa_i)_{i=1}^5\\\hline
45&8&(8,1,1,1,1)&(3,0,1,1,0)\\
47&14&(7,2,1,1,1)&(5,1,0,1,1)\\
55&20&(5,4,1,1,1)&(-3,-3,-1,0,-1)\\
56&30&(5,3,2,1,1)&(-4,-2,-1,0,-1)
\end{array}
\]
The ordered vectors are
\[
\begin{aligned}
u^{45}=u^{55}&=((1,0),(1,0),(2,-1),(1,-2),(1,1)),\\
u^{47}&=((1,0),(1,-1),(1,-3),(1,1),(1,0)),\\
u^{56}&=((1,0),(1,0),(1,-1),(1,-3),(1,1)).
\end{aligned}
\]

These are the marked local tuples of \SeriesRef{maps:src:D.2.23}.

\medskip\noindent\emph{Direct component equations and all 62 natural rows.}

For \(P=aG,Q=bG\), the degree of the section-linearization line
\(\mathcal J=\mathcal O(P-O-(P-Q)+(-Q))\) on component \(j\)
of \(I_\nu\) is the literal cyclic vector

\[
 d_{j,\rm comp}=[j=a\kappa]-[j=0]-[j=a\kappa-b\kappa]+[j=-b\kappa].
\]

The column of a component \(\Theta_h\) is
\(-2[j=h]+[j=h-1]+[j=h+1]\); when \(\nu=2\), both neighbour
terms are retained, and when \(\nu=1\) the column is zero.
At the selected zero one requires equality of these vectors and
\(\gcd(\nu,b\kappa)=1\). At every unselected fibre the component degree vector must
vanish. \SeriesRef{maps:src:D.2.23} gives
\(ab=d/\nu\) at a selected reduced transitive component.

These tests have a short finite proof. When \(\nu>1\) and
\(b\kappa\) is a unit, the coefficient at zero in the component degree vector is zero
if \(a\kappa=0\), is minus two if
\(a\kappa=b\kappa\ne0\), and is minus one otherwise. Equality
with a component column forces \(h=0\) and
\(a\kappa=b\kappa=\pm1\pmod\nu\), which also suffices.
At an unselected fibre, the component degree vector vanishes exactly when one of
\(a\kappa,b\kappa\) is zero modulo \(\nu\): with both nonzero
its zero coefficient is \(-1-[a\kappa=b\kappa]\). All higher
fibre residues \(\kappa\) in the four rows are units.

Factoring each positive integer \(d/\nu\) and applying these
literal equations gives the complete list below. Each position in
the indicated set is a distinct marked row.

\begingroup\small
\begin{center}
\begin{tabular}{c l p{8cm}}
OS&Selected positions&Ordered pairs \((a,b)\)\\\hline
45&2,3,4,5: \(I_1\)&\((1,8),(8,1)\)\\
47&2: \(I_2\)&\((1,7),(7,1)\)\\
47&3,4,5: \(I_1\)&\((1,14),(2,7),(7,2),(14,1)\)\\
55&2: \(I_4\)&\((1,5),(5,1)\)\\
55&3,4,5: \(I_1\)&\((1,20),(4,5),(5,4),(20,1)\)\\
56&1: \(I_5\)&\((1,6),(6,1)\)\\
56&2: \(I_3\)&\((1,10),(2,5),(5,2),(10,1)\)\\
56&3: \(I_2\)&\((1,15),(3,5),(5,3),(15,1)\)\\
56&4,5: \(I_1\)&\((1,30),(2,15),(3,10),(5,6),\)
\((6,5),(10,3),(15,2),(30,1)\)
\end{tabular}
\end{center}
\endgroup

The counts are \(8,14,14,26\), totalling 62. Inspection of this
complete list proves \(\gcd(a,b)=1\). 

The component calculation gives the following line map. The elliptic
group-law line through \(P,-Q\), divided by the vertical line
through \(P-Q\), has horizontal divisor
\(J_{\rm div}=P-O-(P-Q)+(-Q)\). Its vertical correction \(V_{\rm div}\)
satisfies that \(V_{\rm div}+\Theta_0\) is a sum of whole fibres. Its degree
on the base is zero: the mixed-height formula of \SeriesRef{maps:src:D.2.23} evaluates
\(J_{\rm div}\cdot O=1/\nu+(\nu-1)/\nu=1\), with the second
term zero at \(I_1\). Divide by the base rational function with
divisor \(V_{\rm div}+\Theta_0\). The resulting meromorphic function
\(\rho_{\rm line}\) has divisor \(J_{\rm div}-\Theta_0\), and multiplication
by \(\rho_{\rm line}\), inverted by multiplication by \(\rho_{\rm line}^{-1}\), is
\(\mathcal J\simeq\mathcal O(\Theta_0)\). Its canonical section
gives the required linearization. A small nonzero scalar makes the
forward line map contracting away from its zero. Simultaneous sign
reversal uses elliptic inversion and the self-inverse period map
\(\operatorname{diag}(-I_2,1,1)\).

\medskip\noindent\emph{Periods, normal meridians and local models.}

For each fibre choose \(r_{\rm per}\) with \(\det(u,r_{\rm per})=1\). Let
\(\epsilon_i\) be one at the selected position and zero elsewhere.
The four-period matrix is

\[
 T_i=\begin{pmatrix}
 I+\nu_i u_i u_i^tJ&0&b\kappa_i u_i\\
 -a\kappa_i u_i^tJ&1&(\epsilon_i-ab\kappa_i^2)/\nu_i\\
 0&0&1
 \end{pmatrix}.
\]

All its entries are integral by the component calculation. Direct multiplication gives
\(T_1\cdots T_5=I\), \(T_i\delta=\delta\), \(\chi_{\rm ret} T_i=\chi_{\rm ret}\),
and \(T_i^t\omega T_i=\omega\) in every row. Each equation follows by multiplying the matrices in the displayed four tuples.

On a cut regular base the period exponential
\([x]\mapsto[\Pi(q)x]\), with the group logarithm modulo
\(\Pi(q)\Lambda\) as inverse, identifies the actual family with
its marked real torus bundle. The nonconstant periods and their
nonzero holomorphic unit factors remain in \(\Pi\). A change
\(\Pi\mapsto\Pi H\) is identity on the complex universal cover;
new real coordinates are \(H^{-1}\) times old coordinates.

For the local fibre under discussion write \(\mathsf T=T_i\) for its marked period matrix.
At an unselected place \(N_{\rm mon}=\mathsf T-I\) has primitive rank one. More
explicitly, for \(\nu\mid a\),

\[
 N_{\rm mon}=\bigl(u,-a\kappa/\nu,0\bigr)
                (\nu u^tJ,0,b\kappa),
\]

and for \(\nu\mid b\),

\[
 N_{\rm mon}=\bigl(\nu u,-a\kappa,0\bigr)
                (u^tJ,0,b\kappa/\nu).
\]

The row and column in each factorization are primitive and their
pairing is zero. B\'ezout operations complete the column in the kernel
of the row and choose a long vector on which that row is one. This
gives an integral Jordan basis \(H_{\rm nod}\) with
\(H_{\rm nod}^{-1}\mathsf TH_{\rm nod}=I+E_{12}\); its inverse is the integer adjugate divided
by its unit determinant. These are period changes on the defining
rank-one Mumford charts, so the actual cap maps accompany the basis.

In the resulting coordinates the adjacent node charts and inverses are

\[
 X_j=Z/q^j,\quad Y_j=q^{j+1}/Z,\quad X_jY_j=q;
 \qquad Z=q^jX_j,\quad q=X_jY_j.
\]

The ruled normalization, its two conductor sections and product tori
give the integral fibre specializations

\[
 f_{1,\rm nod}(x)=(x_2,x_3,x_4),\qquad
 f_{2,\rm nod}(y)=(y_{12},y_{23},y_{24},y_{34}).
\]

Their sections insert zeros in the omitted coordinates and apply
\(H_{\rm nod}\) or \(\bigwedge^2H_{\rm nod}\). The kernel of the second map is exactly
the primitive local Wang relation lattice. Every fixed-loop normal
sweep bounds in the adjacent node charts. These give the complete boundary map and
its inverse on the fibre quotient.

At the selected zero define \(S=(u,\delta)\), \(L=(r_{\rm per},\gamma)\),
and

\[
 C_{\rm cusp}=\begin{pmatrix}\nu&b\kappa\\-a\kappa&c\end{pmatrix},
 \quad c=(1-ab\kappa^2)/\nu,
 \quad C_{\rm cusp}^{-1}=\begin{pmatrix}c&-b\kappa\\a\kappa&\nu\end{pmatrix}.
\]

The construction of \SeriesRef{maps:src:D.2.23} uses the Jordan basis
\(H_{\rm cusp}=(SC_{\rm cusp}e_1,Le_1,SC_{\rm cusp}e_2,Le_2)\), for which
\(H_{\rm cusp}^{-1}\mathsf T H_{\rm cusp}=I+E_{12}+E_{34}\). Its angular fan is the
unit-triangle fan in the two short columns \(SC_{\rm cusp}\). We retain this
specified fan. The physical circle has coordinates \(C_{\rm cusp}^{-1}(0,1)^t=(-b\kappa,\nu)^t\). The fan transport below relates this marking to the physical-circle fan.

Exponentiating the two short periods gives coordinates \((X,Y)\).
The exact cusp chart and inverse are

\[
 (x,y,w)=(q/(XY),Y,X),\qquad (X,Y,q)=(w,y,xyw).
\]

The long periods translate the fan by the full integral lattice and
multiply \(X,Y\) by their actual nonzero units. The primitive normalization and conductor calculation gives

\[
 f_{1,\rm cusp}(x)=(x_2,x_4),\qquad
 f_{2,\rm cusp}(y)=(y_{12},y_{14}-y_{23},y_{24},y_{34}).
\]

The section of \(f_{2,\rm cusp}\) sends \((u,v,w,z)\) to
\((u,0,v,0,w,z)\) and applies \(\bigwedge^2H_{\rm cusp}\).
The kernel is generated by \(e_{13}\) and \(e_{14}+e_{23}\),
which are precisely the local Wang relations; the extra
\(e_{13}\) in one literal Wang column is already in this lattice.
Both fixed short loops bound normal toric discs. Thus the cap's
boundary map is onto in degree two and has exactly the fixed-sweep
kernel. These are the integral local maps of \SeriesRef{maps:src:D.2.13} and \SeriesRef{maps:src:D.2.16}, transported through the displayed period and chart maps.

The normal meridians are fixed as follows. Trivialize the actual line
along the origin over the base minus the selected zero, so every other
local origin frame extends. At the remaining collar its factor is
\(q^{k_{\rm orig}}\) times an extending unit for an integer \({k_{\rm orig}}\). The actual
long period

\[
                    L_\delta=H_{\rm cusp,deck}^{-b\kappa}K_{\rm cusp,deck}^\nu
\]

has valuation \(\delta\), because \(C_{\rm cusp}(-b\kappa,\nu)^t=(0,1)^t\).
Here \(H_{\rm cusp,deck},K_{\rm cusp,deck}\) denote the original two long deck transformations. Applying its \(-{k_{\rm orig}}\)-th iterate to a representative
of the framed origin removes \(q^{k_{\rm orig}}\), leaving only nonzero extending
units. Applying the \({k_{\rm orig}}\)-th iterate is the inverse. This is identity
on each quotient fibre, and the resulting section extends in a toric
chart at the selected central component. Its positive normal disc
therefore fills the actual meridian. Changing that representative
adds a multiple of \(\delta\) to a cusp sweep; its fibre coefficient
is zero because \(\omega(\delta,S)=0\). This verifies the normal
framing needed below, including the line's base degree.

\Needspace{8\baselineskip}
\medskip\noindent\emph{The allowed original caps and a global obstruction for every other one.}

\smallskip\noindent\emph{The exact permitted original-to-rank-one comparison.}

For a zero-free original \(I_\nu\), the line multidegree is zero
exactly when \(\nu\mid a\), and one component orbit means
\(\gcd(\nu,b\kappa)=1\). Both hold for each original allowed
in the theorem. Put \({A_{\rm Tate}}=a\kappa/\nu\).

On the infinite Tate plumbing let its elliptic long period act by
\(Z\mapsto q^\nu a_0(q)Z\). In the actual line frame its multiplier
has valuation \(-a\kappa\). The change

\[
                     W=Z^{A_{\rm Tate}} V_{\rm line},\qquad V_{\rm line}=Z^{-{A_{\rm Tate}}}W
\]

removes that valuation. This is the multiplicative version of the
integral Poincar\'e shear
\(S_P(x,z_\delta,x_\gamma)=(x,z_\delta-{A_{\rm Tate}} r_{\rm per}^tJx,x_\gamma)\),
with inverse adding \({A_{\rm Tate}} r_{\rm per}^tJx\). To see that the line map extends,
put \(\mathfrak q=q^\nu a_0(q)\) and use the convergent Tate product

\[
 \theta_{\rm Tate}(Z;\mathfrak q)=\prod_{j\ge0}(1-\mathfrak q^jZ)
                   \prod_{j\ge1}(1-\mathfrak q^j/Z),
 \qquad \theta_{\rm Tate}(\mathfrak q Z;\mathfrak q)=-Z^{-1}\theta_{\rm Tate}(Z;\mathfrak q).
\]

The ratio \(\theta_{\rm Tate}(Z/p_P;\mathfrak q)/\theta_{\rm Tate}(Z;\mathfrak q)\) has divisor \(P-O\)
and multiplier \(p_P\), where the geometric section coordinate has
valuation \(-a\kappa\). On each node chart the product is a finite
monomial times a convergent nonzero unit away from its horizontal
divisor. Comparing its meromorphic section with the canonical
meromorphic section of \(\mathcal O(P-O)\) therefore gives a
meromorphic line isomorphism. Their residual vertical divisor has
zero intersection with every component, because the line has zero
multidegree after the line shear. The cyclic intersection matrix has kernel
the whole fibre, so this divisor is \({k_{\rm vert}}\) times that fibre. Division
by \(q^{k_{\rm vert}}\) removes it. The resulting line map and its reciprocal
are nonvanishing on every chart, including the nodes. This proves that the literal line map extends over the whole cap.

Let \(H_{\rm Tate},K_{\rm Tate}\) now be the actual elliptic and forward periods in these
coordinates. They shift the component index by \(\nu,b\kappa\),
respectively; their residual line multipliers are the retained units.
Choose integers \({l_{\rm per}},{k_{\rm per}}\) with \({l_{\rm per}}\nu+{k_{\rm per}}b\kappa=1\). Define

\[
\begin{gathered}
 L_1=H_{\rm Tate}^{l_{\rm per}}K_{\rm Tate}^{k_{\rm per}},\quad E_{\rm per}=H_{\rm Tate}^{-b\kappa}K_{\rm Tate}^\nu;\\
  H_{\rm Tate}=L_1^\nu E_{\rm per}^{-{k_{\rm per}}},\quad K_{\rm Tate}=L_1^{b\kappa}E_{\rm per}^{l_{\rm per}}.
\end{gathered}
\]

These two formulas are inverse integral period changes. The first
new period shifts the chain by one; the second has valuation zero
and supplies the ruled base with the line circle. The original Tate
charts are exactly the adjacent node charts, with all unit multipliers unchanged.
Their quotient by the inverse period changes has one ruled component and its two sections
identified. Identity on this common infinite plumbing, followed by
the displayed inverse period change, identifies the actual original
quotient with the marked primitive rank-one model. It preserves
the parameter \(q\), its positive normal meridian, the physical
line circle and the outer collar. Flat residual line multipliers
can alternatively be removed smoothly on a period parallelogram by
\(e^{-t_1\log\mu_1-t_2\log\mu_2}\), with the positive exponential
as inverse; the original multipliers can equally be retained.

\smallskip\noindent\emph{Every additional original cap gives a global free two-cycle.}

At an unselected higher fibre the natural component condition and
\(\gcd(a,b)=1\) leave two possibilities: \(\nu\mid a\), already
handled, or \(\nu\mid b\) with \(a\kappa\) a unit modulo \(\nu\).
In the second case the original line has degrees minus one and plus
one on the two distinct components met by \(O\) and \(P\), and
the forward translation is narrow. The following primitive global two-cycle gives the obstruction, including \(\nu=2\).

Put \(\gamma'=\gamma-(b\kappa/\nu)r_{\rm per}\) and choose \(A_0,B_0\)
with \(a\kappa A_0+\nu B_0=1\). In the boundary fibre coinvariant
lattice let \(\mu=A_0u+B_0\delta\). The second rank-one factorization gives

\[
 u=a\kappa\mu,\qquad \delta=\nu\mu,
 \qquad H_1(\partial U_{\rm orig})=
       \mathbb Z\langle\mu,r_{\rm per},\gamma',\beta\rangle,
\]

where \(\beta\) is the actual extending original base meridian.
The determinant \(\nu B_0+a\kappa A_0=1\) proves that \(\mu\)
is primitive in this free abelian group.

Let \(N_{\rm ell}\) be the original elliptic-surface neighbourhood. It retracts
onto its cycle of \(\nu\) spheres, so \(\pi_1(N_{\rm ell})=\mathbb Z\).
The unit circle bundle \(E_{\rm circ}\) of the section line has Euler number
\(\pm1\) on one component. In its homotopy sequence the map
\(\pi_2(N_{\rm ell})\to\pi_1(S^1)\) is therefore onto, and
\(\pi_1(E_{\rm circ})\to\pi_1(N_{\rm ell})\) is an isomorphism. Explicitly the line
circle bounds the Hopf disc
\((r e^{it},\sqrt{1-r^2})\) in the degree-one component's unit
bundle; conjugation gives the opposite Euler sign. The nodal thimble
for \(u\) lifts in a nonzero frame on its disc. Its boundary is
\(u+k_{\rm thim}\delta\) for the actual frame winding \(k_{\rm thim}\), and subtracting
\(k_{\rm thim}\) of these Hopf discs fills \(u\). Hence \(\mu\) bounds in
the actual original cap.

This cap identification also follows directly from its positive
roof suspension: take minus logarithmic norm on the punctured line,
and on each fundamental interval use the increasing smooth map
with constant endpoint derivatives in the norm construction of \SeriesRef{maps:src:D.2.16}. Its inverse is the inverse increasing function.
The cap is the mapping torus of the angular lift on \(E_{\rm circ}\), with
all line units retained. The narrow translation preserves the
oriented dual graph, so it acts as identity on \(\pi_1(E_{\rm circ})=\mathbb Z\).
Thus the preceding bounding discs remain valid in the quotient.

They also bound on the other side of this boundary. For the selected
cusp, its two toric discs, combined using \(C_{\rm cusp}^{-1}\) and the torus
addition triangles, kill \(u_*,\delta\). Integer elimination gives an integral right inverse of the matrix whose columns are

\[
 u_*,\delta,T_1u_*,T_1\delta,\ldots,T_5u_*,T_5\delta
 \quad\hbox{onto }\ker\chi_{\rm ret}.
\]

All 62 matrices have three unit Smith entries. Each translated disc
is an actual disc with a stem transported around the indicated
peripheral loop; the loops exist in the exterior even for the cap
currently omitted. Addition of fibre loops uses the marked torus
triangles. This right inverse therefore supplies an integral
bounding chain for \(\mu\) in the complement of the original cap,
independently of every other chosen filling and every regular log.

Glue the inside and outside chains with their opposite boundary
orientations. The resulting class \(Z_\mu\in H_2(X;\mathbb Z)\)
has Mayer--Vietoris boundary \(\mu\). Reading the \(\mu\)
coefficient in the boundary basis is a homomorphism \(H_2(X)\to\mathbb Z\)
sending \(Z_\mu\) to one. The map \(1\mapsto Z_\mu\) is its
section. This is a global free direct summand and excludes a sphere
for every additional original cap, including every combination of
several such caps. It proves precisely the original-cap exclusion
needed for the claimed complete sphere table.

\medskip\noindent\emph{The integral exterior and its actual sweep splitting.}

First insert one regular hole before the five singular holes. Its
monodromy is \(T_0=I\). Write \(E_{\rm ext}\) for this exterior. Cut the base along an ordered tree and use
parallel period coordinates on its strips. The graph boundary is

\[
 D=[0\ T_1-I\ T_2-I\ T_3-I\ T_4-I]:\Lambda^5\to\Lambda.
\]

For \(0\le i<5\), let \(j_i\) be the indicated block inclusion;
the last based boundary map is

\[
 j_5(x)=-(T_1T_2T_3T_4x,T_2T_3T_4x,T_3T_4x,T_4x,x).
\]

Telescoping gives \(Dj_5(x)=(I-T_5^{-1})x\). Reversing the cut
strips and their transports gives the inverse geometric identification.

For each \(i\), integer Smith reduction of \(N_i=T_i-I\) keeps
its elementary row and column operations and gives an integral basis
\({\mathsf F}_i\) of \(\ker N_i\). Its rank is three at each unselected
place and two at the selected place. Set

\[
 B=[j_1{\mathsf F}_1\ \cdots\ j_5{\mathsf F}_5],\qquad
 K=[B\ j_0e_1\ j_0e_2\ j_0\gamma].
\]

For each row of the displayed complete list, integer elimination gives matrices
\({\mathsf C}_0,{L_K}\) satisfying

\[
 D_0{\mathsf C}_0=I_3,\qquad DK=0,\qquad {L_K}K=I_{17},\qquad
 K{L_K}(I-{\mathsf C}_0D_0)=I-{\mathsf C}_0D_0,
\]

where \(D_0\) is the first three rows of \(D\). These identities
give \(\operatorname{im}D=\ker\chi_{\rm ret}\) and
\(\ker D=K\mathbb Z^{17}\), with actual inverse \({L_K}\).
The nonzero Smith entries of \(D_0\) and \(K\) are respectively three and seventeen units. To construct the matrices, retain every Euclidean row and column operation: if \(U_DD_0V_D=(I_3,0)\), take \({\mathsf C}_0=V_D E_3U_D\), where \(E_3=(I_3,0)^t:\mathbb Z^3\to\mathbb Z^{20}\); if \(U_KKV_K=(I_{17},0)^t\), take \({L_K}=V_KP_{17}U_K\), where \(P_{17}=(I_{17},0):\mathbb Z^{20}\to\mathbb Z^{17}\). These formulas give the displayed integer identities and their inverse coordinates.

In exterior-square order \((12,13,14,23,24,34)\), put
\(B_2=[0,\bigwedge^2T_1-I,\ldots,\bigwedge^2T_4-I]\).
Choose \(x_{\rm fib},y_{\rm fib}\in\mathbb Z\) with \(ax_{\rm fib}+by_{\rm fib}=1\), and put

\[
 s_0=x_{\rm fib}e_{12}-y_{\rm fib}e_{34},\qquad
 \mathsf F=(be_{12}+ae_{34},e_{13},e_{14},e_{23},e_{24}).
\]

The same Euclidean reduction gives an integral \(\mathsf Z\) with
\(B_2\mathsf Z=\mathsf F\). Since \(\omega B_2=0\), \(\omega s_0=1\), and
\([s_0,\mathsf F]\) is unimodular up to the displayed basis order,

\[
 \bigwedge^2\Lambda/\operatorname{im}B_2
 \ \underset{k\mapsto[ks_0]}{\overset{[z]\mapsto\omega(z)}
 {\rightleftarrows}}\ \mathbb Z
\]

are inverse maps. For any difference in the kernel, \(\mathsf Z\) gives
its monodromy-cylinder bounding chain. Write \(s=[s_0]\).

The fixed-loop sweeps give a geometric splitting with the following global relation. A fixed loop \(u\in\ker(T_i-I)\) sweeps its based
boundary circle to an actual two-torus. Give it the fibre-first
orientation. Addition of loop vectors is the boundary of the
real-period triangle with vertices \(0,u,u+v\), and subtraction
gives the inverse operation. The physical \(\delta\)-circle over
the cut base is an actual three-chain, whose boundary is the sum
of all these \(\delta\) sweeps. Its frame changes are line-circle
changes; their potential extra two-torus is \(\delta\wedge\delta=0\).
The local coordinate calculation fixes the remaining origin frame at the cusp.
Consequently the relation among the graph columns is the same
relation among the actual two-cycles, with zero additional coefficient
of \(s\).

For the one-hole graph, the last three coordinates of \({L_K}j_0\delta\)
are zero. The first fourteen coordinates are exactly the negatives
of the representations of \(\delta\) in the five fixed bases:
this is the identity \(j_0\delta+\sum_{i=1}^5j_i\delta=0\),
which follows directly from the last based boundary map. The sequence

\[
 0\longrightarrow\mathbb Zs\longrightarrow H_2(E_{\rm ext})
       \longrightarrow\ker D\longrightarrow0
\]

therefore splits by the actual fixed sweeps and the three displayed
regular sweeps. Its inverse projects to the graph, applies \({L_K}\),
subtracts those actual sweeps, and reads the remaining fibre part
using \(\omega\).

For an arbitrary number \(s_\mathrm{reg}\) of regular holes, the
graph map acquires zero blocks. The last boundary map acquires a
block \(-x\) for each. Subtracting these new coordinates, using
the preceding inverse, and restoring them by their block inclusions
gives the integral inverse for every finite number of holes. With
no regular hole, delete its block and the three independent regular
sweeps; the fourteen fixed columns then have the single primitive
relation \(\sum j_i\delta=0\) and span the thirteen-dimensional
kernel. Equivalently, the one-hole inverse descends by its actual
\(j_0\Lambda\) subgroup. This covers the no-site case.

After the five primitive critical fillings, the rank-one specialization and the cusp specialization remove
their complete fixed-sweep kernels and no further fibre class. For
the remaining regular holes the group is therefore presented by
\(s,t_i(u)\), linear in \(u\in\Lambda\), with the single relation
\(\sum_i t_i(\delta)=0\) in the canonical marking. Surjectivity
of each local boundary map in degree two makes this a complete
Mayer--Vietoris cokernel: lift a cap class to its boundary and
subtract it from the exterior coordinate; the inverse includes the
exterior with zero cap coordinates. Different lifts differ by exactly
the kernels just removed.

\medskip\noindent\emph{Every regular logarithmic cap and the full integral presentation.}

Take arbitrary \(m_i>0\), \(v_i=(x_i,y_i,z_i,n_i)\in\Lambda\)
with \(\gcd(m_i,x_i,y_i,z_i,n_i)=1\). On the root disc
\(q=z^{m_i}\) the deck generator is

\[
       (y,z)\longmapsto(y+v_i/m_i,e^{2\pi i/m_i}z).
\]

Its translation on the central torus has exact order \(m_i\):
the least positive \(k\) with \(kv_i\in m_i\Lambda\) is
\(m_i/\gcd(m_i,x_i,y_i,z_i,n_i)\). Off the central torus the
root rotation is free. This proves the stated precise freeness
condition and gives a smooth proper finite quotient. The actual
period image is holomorphic, so it is the stipulated logarithmic
filling for every such vector, including zero or negative \(n_i\).

Its central lattice has inverse integral descriptions

\[
 \Lambda_i=\Lambda+\mathbb Z(v_i/m_i)
     \cong\mathbb Z^5/\mathbb Zq_{i,\rm lat},
 \quad q_{i,\rm lat}=(-v_i,m_i)^t,
 \quad [(z,k)]\longmapsto z+kv_i/m_i.
\]

Choose \(\alpha_{i,\rm norm}\in\mathbb Z\) and integral \(\ell_i\) with
\(\alpha_{i,\rm norm} m_i+\ell_i(v_i)=1\). The projection
\(P_i=I-q_{i,\rm lat}(-\ell_i,\alpha_{i,\rm norm})\) has image
\(\ker(-\ell_i,\alpha_{i,\rm norm})\). The maps \([z]\mapsto P_i z\)
and \(u\mapsto[u]\) are inverse. This also proves that the
central lattice is free and includes its primitive integral marking.

The normal line and inverse trivialization are

\[
 [(y,z)]\mapsto([y],e^{-2\pi i\ell_i(y)}z),\qquad
 ([y],w)\mapsto[(y,e^{2\pi i\ell_i(y)}w)].
\]

The deck multiplier is
\(e^{2\pi i(1-\ell_i(v_i))/m_i}=e^{2\pi i\alpha_{i,\rm norm}}=1\),
so both maps descend. Thus the local relative groups are
\(H_3=\Lambda_i,H_2=\mathbb Z,H_1=0\), with their Thom
generators in this exact normal framing.

The comparison with the holomorphic collar is also explicit. Write
\(z=r_{\rm col} e^{i\theta}\), let \(J_z\) be the complex structure in real
period coordinates, and let the extending origin factor be
\(q_{\rm sec}(z^{m_i})\ne0\). Put

\[
 b_{\rm sec}=\frac{\log q_{\rm sec}(z^{m_i})}{2\pi i},\quad
 u_{\rm sec}=\Re b_{\rm sec}\,\delta+\Im b_{\rm sec}\,J_z\delta,
 \quad B_{\rm col}=\frac{\log r_{\rm col}}{2\pi}J_zv_i-u_{\rm sec}.
\]

The maps
\(h_\tau(y,z)=y-\theta v_i/(2\pi)+\tau B_{\rm col}(z)\)
and \(h_\tau^{-1}(y,z)=y+\theta v_i/(2\pi)-\tau B_{\rm col}(z)\)
are inverse collar maps. The residual term is periodic and deck
invariant; multiplying it by a radial cutoff extends the isotopy
over the exterior. These maps carry the positive normal to
\(d_i=m_i c_i-v_i\), including the radial and line-unit factors.

Orient a relative three-generator as minus its fibre loop times the
positive normal disc. Its boundary is the fibre-first torus
\(\lambda\times d_i\). Cutting that torus into the \(m_i\)
sectors gives the exact transfer

\[
 R_i(\lambda)=m_i t_i(\lambda)-\omega(\lambda,v_i)s,
 \qquad R_i(c_i)=t_i(v_i).
\]

The second formula follows either by its root-generator torus or
from \(m_ic_i=v_i\): the first gives
\(R_i(v_i)=m_it_i(v_i)\), and the exterior degree-two group is
free. Multiplication by \(m_i\) is therefore injective and determines
the fibre coefficient uniquely. The primitive source relation
\((-v_i,m_i)\) maps to zero because \(\omega(v_i,v_i)=0\).
Normal source columns belonging to the critical fillings are
eliminated by subtracting their already identified unit sweep
columns; adding them back is the inverse elementary operation.
This accounts for the full relative degree-three image.

For noncanonical lifts keep the integer \(e\) in the global
Seifert relation. The attaching polygon for its \(\gamma\)
coordinate adds \(-\omega(\delta,e\gamma)=be\) to the
\(\delta\) sweep. Hence the complete cokernel is

\[
 \mathcal C=\left\langle s,t_i(u)\ \left|\
 m_it_i(u)=\omega(u,v_i)s,\quad t_i(v_i)=0,\quad
 \sum_i t_i(\delta)+be\,s=0\right.\right\rangle.
\]

Every generator is integral. A change of peripheral lift is given by

\[
 v_i'=v_i+m_i\lambda_i,\quad e'=e+\sum_i\chi_{\rm ret}(\lambda_i),
 \qquad s'\mapsto s,\quad
 t_i'(u)\mapsto t_i(u)+\omega(u,\lambda_i)s.
\]

Its inverse subtracts the same term. It is induced by the literal
peripheral replacement \(c_i'=c_i+\lambda_i\) and the addition
triangle in the fibre; angular coordinates change by the opposite
\(\theta\lambda_i\). The local rows are preserved by bilinearity,
the root row by \(\omega(\lambda_i,v_i)=-\omega(v_i,\lambda_i)\),
and the global row by
\(\omega(\delta,\lambda_i)=-b\chi_{\rm ret}(\lambda_i)\). Thus all lift
conventions are related by specified inverse integral maps.

For \(m_i=1\), eliminate
\(t_i(u)=\omega(u,v_i)s\); its root row is zero and the global
row changes \(e\) to \(e-n_i\). The inverse inserts those four
generators with the same values. In particular a presentation with
no displayed sites and arbitrary \(e\) is represented by one
order-one site \(v=-e\gamma\) in the canonical \(e=0\) convention.
The no-site formulas are therefore also covered by actual marked
attachments, as well as by the direct no-hole calculation above.

\medskip\noindent\emph{Fundamental group, every homology group, and all extensions.}

The cusp discs and the right inverse of the matrix of transported loops kill exactly
\(\ker\chi_{\rm ret}\) in the already abelian fibre subgroup. Every other
vanishing loop has zero \(\chi_{\rm ret}\). The marked normal meridians from
the local coordinate calculation and the logarithmic transfer therefore give precisely the Seifert presentation of \SeriesRef{maps:src:D.2.23}

\[
 \pi_1(X)=\left\langle\gamma,c_1,\ldots,c_{s_{\rm reg}}\ \left|\
 [\gamma,c_i]=1,\quad c_i^{m_i}=\gamma^{n_i},\quad
 c_1\cdots c_{s_{\rm reg}}=\gamma^e\right.\right\rangle.
\]

Its abelianization has the relation rows
\((-n_i,0,\ldots,m_i,\ldots,0)\) and \((-e,1,\ldots,1)\).
Call this matrix \({\mathsf S_{\rm Sei}}\), and write its Smith cokernel as
\(\mathbb Z^r\oplus A_1\). Rational elimination gives

\[
 e_{\rm rat}=e-\sum_i n_i/m_i,
 \qquad r=\begin{cases}1&e_{\rm rat}=0,\\0&e_{\rm rat}\ne0.\end{cases}
\]

The exact nonabelian triviality criterion is that of \SeriesRef{maps:src:D.2.23}:
after eliminating order-one sites and subtracting their \(n_i\)
from \(e\), at least three remaining sites give a nontrivial group;
with zero, one or two sites the group is abelian and the respective
unit tests are

\[
 |e|=1,\qquad |m_1e-n_1|=1,\qquad
 |m_2n_1+m_1n_2-em_1m_2|=1.
\]

The zero-determinant
case retains the free summand in the Euler-number formula.

Over \(\mathbb Q\), \eqref{maps:eq:D.SC2} eliminates
\(t_i(u)=\omega(u,v_i)s/m_i\). The root rows then vanish by
alternation and the final row is \(be_{\rm rat}s=0\). Thus
\(\mathcal C\) has exactly free rank \(r\). Let
\(T=\operatorname{Tor}\mathcal C\).

The remaining degree-two term comes from the relative normal discs. Eliminate the four
independent critical meridians by their unit normal-disc columns.
In the resulting integral relative complex the remaining normal
generators \(D_i,D_*\) have boundary columns

\[
 \partial D_i=m_i c_i-n_i\gamma,\qquad
 \partial D_*=\sum_i c_i-e\gamma.
\]

Here \(D_*\) is the base polygon capped by the five critical
normal discs, in the framing already checked in the local coordinate calculation. The
successive unit eliminations are reversed by reinserting those same
disc columns. Thus the reduced normal-disc complex is an actual integral reduction of the
relative complex. It is the transpose of the row convention for \({\mathsf S_{\rm Sei}}\).

For \(e_{\rm rat}\ne0\) the kernel is zero. For \(e_{\rm rat}=0\), put
\({\ell_{\rm mult}}=\operatorname{lcm}_i m_i\), with \({\ell_{\rm mult}}=1\) for no sites. The
kernel is primitively generated by

\[
              ({\ell_{\rm mult}}/m_1,\ldots,{\ell_{\rm mult}}/m_{s_{\rm reg}},-{\ell_{\rm mult}}).
\]

Primitivity follows prime by prime: a site with maximal prime-power
valuation in \({\ell_{\rm mult}}\) makes its corresponding \({\ell_{\rm mult}}/m_i\) a unit at
that prime. For no sites the sole coefficient is minus one.

There is an explicit closed lift of the primitive relative vector. Take that combination of
normal discs and the capped base polygon. Cancel its repeated base
edges with their oriented strip cylinders and torus addition
triangles. The residual fibre loop is, with the chosen boundary sign,

\[
                   w=\sum_i({\ell_{\rm mult}}/m_i)v_i-{\ell_{\rm mult}}e\gamma.
\]

Its \(\chi_{\rm ret}\) coordinate vanishes by \(e_{\rm rat}=0\). The right inverse
in the transported cusp-disc matrix expresses \(w\) as an integer sum of transported cusp
shorts. Attach those actual discs with opposite boundary orientation,
using the same addition triangles. This gives a closed integral
two-cycle whose relative coordinate is the primitive relative vector displayed above. Sending one to this cycle splits the relative extension.
The inverse reads the relative coefficient, subtracts the chosen
cycle, and uses the quotient inverse of the exterior calculation and the logarithmic transfer. Consequently

\[
       H_2(X;\mathbb Z)=\mathcal C\oplus\mathbb Z^r
                        \cong\mathbb Z^{2r}\oplus T.
\]

The closed lift splits the free extension of the Mayer--Vietoris
cokernel. For integral coordinates, make the columns of a matrix \(\mathsf R\) the relations
in \eqref{maps:eq:D.SC2}. Elementary Smith operations give
\(\mathsf U\mathsf R\mathsf V=\operatorname{diag}(d_{1,\rm Smith},\ldots,d_{k_{\rm Smith},\rm Smith},0,\ldots)\).
The forward quotient map is \([z]\mapsto[\mathsf Uz]\), and the inverse
is \([w]\mapsto[\mathsf U^{-1}w]\). Every elementary addition, interchange
and sign change has its stated inverse. Each elementary operation is recorded together with its inverse, giving the displayed product identity. This works for
unbounded integer entries and includes the torsion group \(T\)
with all its extensions.

The total space is compact, connected and oriented by its complex
charts: it is a compact exterior with finitely many proper compact
local fillings attached along full collars. The finite free quotients
in the root action preserve that orientation. The regular caps, primitive nodal
caps and boundary bundles have Euler characteristic zero. The
primitive cusp has Euler characteristic two by its full normalization
and conductor calculation. Therefore \(\chi(X)=2\).

Integral Poincar\'e duality and universal coefficients give all the groups:

\[
\begin{array}{c|c}
 j&H_j(X;\mathbb Z)\\\hline
 0,6&\mathbb Z\\
 1&\mathbb Z^r\oplus A_1\\
 2,3&\mathbb Z^{2r}\oplus T\\
 4&\mathbb Z^{2r}\oplus A_1\\
 5&\mathbb Z^r.
\end{array}
\]

Indeed \(H_4\cong H^2\) has torsion
\(\operatorname{Ext}(H_1,\mathbb Z)=A_1\) and free rank \(2r\),
while \(H_5\cong H^1\) is free of rank \(r\).
The torsion of \(H_3\cong H^3\) is
\(\operatorname{Ext}(H_2,\mathbb Z)=T\). In a Smith factor
\(\mathbb Z/d_{\rm Smith}\), the character \([k]\mapsto k/d_{\rm Smith}\pmod{\mathbb Z}\),
followed by the connecting homomorphism and cap product, gives
this identification and its inverse in Smith coordinates.
Euler characteristic then gives \(b_3=2r\). The free degree-three
intersection form is alternating and unimodular. A primitive
class and a B\'ezout combination pairing to one give a symplectic
pair; subtracting their intersection coordinates gives the inverse
projection to its complement. For \(r=0\) there are no free
classes, and for \(r=1\) this one pair is the full free lattice.
Thus \eqref{maps:eq:D.SC3} records actual integral decompositions, including the free extensions.

\medskip\noindent\emph{The one-site matrix.}
For one site in the canonical marking, the zero delta sweep gives
\begin{equation}
 R_{\rm one}=\begin{pmatrix}
 -ay&ax&bn&-bz&0\\
 m&0&0&0&x\\0&m&0&0&y\\0&0&0&m&n
 \end{pmatrix}.
 \hypertarget{tag-0245}{}\label{maps:eq:D.SC7}
\end{equation}
Its signed cofactor vector is
\(bm^2n(-x,-y,-z,-n,m)\): the last minor is \(bm^3n\),
and multiplication of \(R_{\rm one}\) by the displayed primitive relation
is zero. Direct expansion gives the same formula also at \(n=0\).
For \(n\ne0\), the rank is four, and the gcd of the five maximal
minors is \(bm^2|n|\), since \(\gcd(m,x,y,z,n)=1\).
Thus its degree-two cokernel has exactly that order.

\Needspace{5\baselineskip}
\medskip\noindent\emph{Unbounded simply connected reduction and the exact torsion formula.}

For a simply connected member, criterion \eqref{maps:eq:D.FG3} implies, with
all order-one sites still retained, that the \(m_i\) are pairwise
coprime, \(\gcd(m_i,n_i)=1\), and

\[
 M=\prod_i m_i,\quad q_{i,\rm crt}=M/m_i,\quad
 \Delta=\sum_i q_{i,\rm crt} n_i-eM=\pm1.
\]

These consequences hold for the empty list with \(M=1\), and
the same formulas include any number of order-one sites. They
follow directly by reducing the unit determinant modulo a prime
common to two multiplicities or to one multiplicity and its numerator.

Multiply the last row of \eqref{maps:eq:D.SC2} by \(M\) and use its
\(\delta\) rows. The result is \(-b\Delta s=0\), so
\(bs=0\). The local \(\delta\) rows give
\(m_it_i(\delta)=0\), while their sum is zero. Pairwise
coprimality and the integral Chinese-remainder idempotents isolate
each of them, proving \(t_i(\delta)=0\).

Choose \(\alpha_{i,\rm crt} n_i+\beta_{i,\rm crt} m_i=1\) and put
\(U_i=t_i(e_1), V_i=t_i(e_2)\). The remaining root and
\(\gamma\) relations have the exact inverse elimination

\[
 t_i(\gamma)=-\alpha_{i,\rm crt}(x_iU_i+y_iV_i),\qquad
 bs=0,\quad m_iU_i=ay_i s,\quad m_iV_i=-ax_i s.
\]

To check its converse, the proposed \(t_i(\gamma)\) satisfies
its \(m_i\) relation because
\(x_iay_i-y_iax_i=0\). Its root relation reduces to
\(\beta_{i,\rm crt} m_i(x_iU_i+y_iV_i)=0\) by the same cancellation.
This verifies both substitutions, including \(n_i=0,m_i=1\).

Consolidate the generators by defining

\[
 X_0=\sum_i q_{i,\rm crt} x_i,\qquad Y_0=\sum_i q_{i,\rm crt} y_i,
 \qquad U=\sum_iU_i,\quad V=\sum_iV_i.
\]

Choose \(\eta_i q_{i,\rm crt}\equiv1\pmod{m_i}\), and let
\(r_i=q_{i,\rm crt}\eta_i\). For \(m_i=1\) any B\'ezout choice is allowed.
Set the integers

\[
 A_i^y=\sum_j\frac{r_i-[i=j]}{m_j}y_j,
 \qquad A_i^x=\sum_j\frac{r_i-[i=j]}{m_j}x_j.
\]

Every numerator is divisible by \(m_j\). The forward map from
the three-generator presentation sends \((s,U,V)\) to
\((s,\sum U_i,\sum V_i)\). Its inverse sends

\[
       U_i\longmapsto r_iU-aA_i^y s,
       \qquad V_i\longmapsto r_iV+aA_i^x s.
\]

In fact \(r_iU-U_i=aA_i^y s\) and
\(r_iV-V_i=-aA_i^x s\) by the local inverse elimination. Conversely
\(\sum r_i=1+k_{\rm crt}M\) and \(\sum A_i^y=k_{\rm crt}Y_0\),
\(\sum A_i^x=k_{\rm crt}X_0\), so the composite on \(U,V\) is identity
modulo \(MU=aY_0s\), \(MV=-aX_0s\). These equations also
verify every individual \(m_i\) relation under the inverse consolidation maps.
They are exact integer maps for every finite number of sites.

Finally choose \(ac_{\rm unit}+bd_{\rm unit}=1\), replace \(s\) by
\(\sigma_{\rm tor}=as\), and use \(s=c_{\rm unit}\sigma_{\rm tor}\) as inverse. The reduced
matrix, with columns the three relations, is

\[
 \begin{pmatrix}
 b&-Y_0&X_0\\0&M&0\\0&0&M
 \end{pmatrix}.
\]

Its first determinantal divisor is
\(g=\gcd(b,M,X_0,Y_0)\). Its two-by-two minors generate
\((bM,MX_0,MY_0,M^2)=Mg\,\mathbb Z\), and its determinant
is \(bM^2\). Consequently the invariant factors are
\(g,M,bM/g\). Applying the elementary integer Smith algorithm
to this explicit three-by-three matrix gives the quotient map and
inverse just described after the degree-two splitting. Thus

\[
 H_2(X;\mathbb Z)=H_3(X;\mathbb Z)
 \cong\mathbb Z/g\oplus\mathbb Z/M\oplus\mathbb Z/(bM/g).
\]

For no displayed sites, \(M=1,X_0=Y_0=0\), and the same matrix
has two unit factors and the factor \(b\). The displayed inverse substitutions account for the line coordinate \(z_i\), every sign of \(n_i\), and every peripheral lift change.

For the primary description, let a prime p divide one multiplicity
\(m_i\). The other multiplicities are units in \(\mathbb Z_{(p)}\).
Eliminating their sweeps and the delta sweep, then using the unit
\(n_i\), gives the transpose relation convention
\begin{equation}
 \begin{pmatrix}
 b&0&0\\-ay_i&m_i&0\\ax_i&0&m_i
 \end{pmatrix}.
 \hypertarget{tag-0246}{}\label{maps:eq:D.SC8}
\end{equation}
Its determinantal divisors are \(g_i,m_i g_i,bm_i^2\), where
\(g_i=\gcd(b,m_i,x_i,y_i)\) in this local ring; the equality
uses \(\gcd(a,b)=1\). At primes not dividing \(M\) the sweeps
eliminate to \(\mathbb Z_{(p)}/(b)\). The global inverse
Chinese-remainder maps above identify these local descriptions
with \eqref{maps:eq:D.SC4} over the integers.

\medskip\noindent\emph{Sphere criterion and the 17 rows.}

The order of the finite group in \eqref{maps:eq:D.SC4} is \(bM^2\). It vanishes
as a group exactly when \(b=M=1\). In that case \eqref{maps:eq:D.SC3} and
simple connectivity give an integral homology six-sphere. Hurewicz
and Whitehead give a homotopy six-sphere, and the dimension-six
homotopy-sphere calculation, with the orientation convention detailed
in the fan-transport argument, gives the standard oriented sphere. The smooth reconstruction also
constructs the resulting oriented map explicitly through its
suspension model. Conversely any sphere has trivial fundamental
group and \(H_2=0\), so \eqref{maps:eq:D.SC4} forces these same conditions.

In the canonical convention \(e=0\) with every site retained,
\(M=1\) means every \(m_i=1\). Eliminating them in \eqref{maps:eq:D.FG3} gives
\(|\sum_i n_i|=1\). The empty sum is zero and supplies no sphere.

The complete initial list above has exactly the seventeen \(b=1\) rows of \eqref{maps:eq:D.SC5}.

The original-cap comparison supplies the exact maps for every permitted original cap
and a global free direct summand excluding every additional original
cap. Hence this table includes all original-versus-Mumford choices
in the stated natural semistable sector. 

\medskip\noindent\emph{Keeping and transporting the specified cusp fan.}

\smallskip\noindent\emph{A finite word of actual periodic edge flips.}

In the primitive short lattice \(S=(u,\delta)\), the specified
source fan is the periodic unit-triangle fan \(\mathcal F(C_{\rm cusp})\)
whose basis columns are \(C_{\rm cusp}=(p_{\rm cusp},q_{\rm cusp})\). Its unoriented edge directions
are \(p_{\rm cusp},q_{\rm cusp},q_{\rm cusp}-p_{\rm cusp}\). For any ordered unimodular basis \(E_{\rm fan}=(p_{\rm fan},q_{\rm fan})\),
\(\mathcal F(E_{\rm fan})\) consists of the translates of triangles
\((0,p_{\rm fan},q_{\rm fan})\) and \((p_{\rm fan}+q_{\rm fan},p_{\rm fan},q_{\rm fan})\).

Each operation \(p_{\rm fan}\mapsto p_{\rm fan}\pm q_{\rm fan}\), \(q_{\rm fan}\mapsto q_{\rm fan}\pm p_{\rm fan}\),
or \((p_{\rm fan},q_{\rm fan})\mapsto(q_{\rm fan},-p_{\rm fan})\) retains exactly two of these three
unoriented edge directions and replaces the third by the other
diagonal of its primitive parallelograms. It is one actual periodic
edge flip. Simultaneous negation gives the same triangulation,
because it interchanges the two triangle types by translation.

Here is a terminating word, specified for every integer \(C_{\rm cusp}\).
While its first-row second entry is nonzero, divide the first
entry by the second with remainder; implement the quotient by
unit additions to the first column, then replace \((p_{\rm fan},q_{\rm fan})\) by
\((q_{\rm fan},-p_{\rm fan})\). The absolute value of the new second entry is strictly
smaller. At termination the first row is \((\pm1,0)\). Negate
both columns when necessary. The matrix is now
\(\left(\begin{smallmatrix}1&0\\c_{\rm red}&1\end{smallmatrix}\right)\);
\(|c_{\rm red}|\) further signed unit additions to its first column give
the identity. Each listed inverse reverses the corresponding
addition or rotation. Thus the word is a finite path of actual
fans from the specified \(\mathcal F(C_{\rm cusp})\) to
\(\mathcal F(I)\), whose second short direction is the physical
\(\delta\).

The strict decrease proves termination for every matrix in the complete list. The retained edge directions specify each actual local replacement.

For every intermediate basis \(E_{\rm fan}\), its exact period basis is

\[
\begin{gathered}
 H_{\rm fan}=(SE_{\rm fan}e_1,LC_{\rm cusp}^{-1}E_{\rm fan}e_1,\\
 \quad SE_{\rm fan}e_2,LC_{\rm cusp}^{-1}E_{\rm fan}e_2),\\
  H_{\rm fan}^{-1}\mathsf T H_{\rm fan}=I+E_{12}+E_{34}.
\end{gathered}
\]

All columns and inverses are integral. Exponentiating the short
coordinates and multiplying by their inverse period matrix gives
the torus maps; the long periods keep their actual holomorphic
units. The cusp chart in these coordinates supplies each intermediate
smooth cap. Its period lattice translates the full grid, so the
quotient is the same finite local construction with a compact
central fibre and unchanged punctured family. These maps use the specified punctured family and its marking.

\smallskip\noindent\emph{Each quotient flop is an embedded local replacement.}

Fix the edge direction being flipped. Every triangle contains one
edge in that orbit. Distinct lifted complete edge curves therefore
do not meet; a common triple point would require two such edges
in one triangle. No nonzero lattice translation stabilizes a finite
edge. Translations preserve the two triangle types, so the two
endpoints of one complete edge curve are never identified. Its
quotient is consequently one embedded \(\mathbb P^1\).

In its two affine curve charts the exact transition and inverse are

\[
 (x',y',w')=(wx,wy,w^{-1}),\qquad
 (x,y,w)=(w'x',w'y',(w')^{-1}).
\]

The normal bundle is \(\mathcal O(-1)\oplus\mathcal O(-1)\).
The contraction and base parameter are

\[
 {\mathsf N}=\begin{pmatrix}x&xw\\y&yw\end{pmatrix},\quad \det {\mathsf N}=0,
                    \qquad q=xyw={\mathsf N}_{11}{\mathsf N}_{22}.
\]

On \(x\ne0\), the opposite resolution has coordinates
\(u=y/x,p=x,q'=xw\), with inverse
\(x=p,y=up,w=q'/p\). Choosing another nonzero entry gives the
other inverse charts; they agree since all recover the same matrix
\({\mathsf N}\). Thus the comparison is identity on the punctured cone
and preserves the actual base parameter. Its torus characters
commute with the retained period unit multipliers, so these maps
descend through the period quotient. They are identity on the
unchanged collar of the cusp filling.

On a small matrix-norm link write \({\mathsf N}=\varepsilon \xi\eta^t\),
\(|\xi|=|\eta|=1\), modulo
\((\xi,\eta)\sim(t\xi,t^{-1}\eta)\). The first resolution records the
row line \([\eta]\), the second the column line \([\xi]\). In
ordered base/fibre representatives the boundary comparison swaps
the two factors and its inverse swaps them back. The matrix \({\mathsf N}\) remains fixed while the base and fibre roles change. The product of two three-sphere factors changes
orientation under this swap, and the circle gauge also reverses
orientation. Their quotient comparison is orientation preserving,
as is the identity of the punctured complex cone. The same
positive matrix norm supplies the matching collars.

\smallskip\noindent\emph{The actual curve class remains zero through the entire word.}

For \(b=1\) and trivial fundamental group, the integral torsion reduction proves
\(s=0\) integrally. Initially the cusp specialization is onto from the nearby
\(T^4\) in degree two. Every local curve class therefore has a
lift in \(H_2(T^4)\), whose global image is its \(\omega\)
coefficient times \(s\), hence zero.

This assertion persists through every actual flip. Removing its
exceptional real two-sphere \(C_{\rm flop}\) from a six-dimensional cap \(U_{\rm flop}\) gives
relative groups
\(H_k(U_{\rm flop},U_{\rm flop}\setminus C_{\rm flop})=H_{k-4}(C_{\rm flop})\) by the oriented normal
Thom map. In degrees two and three these are zero. Consequently
\(H_2(U_{\rm flop}\setminus C_{\rm flop})\to H_2(U_{\rm flop})\) is an isomorphism. The same
holds for the opposite resolution. Identity on the punctured
cone identifies the complements, including the nearby \(T^4\).
It therefore transports the surjective fibre specialization and
the global degree-two map to the next stage. The analogous global
complement argument gives the identical integral \(H_2\) group
and preserves the class called \(s\). Codimension four also
preserves fundamental group by perturbing loops and filling discs
off the sphere. Induction proves that every next curve in the
finite word is nullhomologous in a simply connected six-manifold.
Hurewicz makes its parametrized embedding nullhomotopic.

\smallskip\noindent\emph{The exact relative oriented diffeomorphism for each flop.}

We give the relative smooth comparison with its prescribed outer boundary parameterization.
The real oriented normal bundle in the two curve charts is trivial: in the
quaternion coordinate \(z_1+z_2j\), its equatorial clutch
\(e^{-i\theta}\) contracts through

\[
 h_{\rm quat}(t,\theta)=\cos(\pi t/2)e^{-i\theta}
                         +\sin(\pi t/2)j,
\]

followed by a path from the constant \(j\) to one. This specifies
the real trivialization and its inverse. \cite[Theorem B(b), pp.~367--368]{manBarden65} applies to a nullhomotopic sphere in
dimension six: \(12>9\), and its required mapping-cylinder
connectivity holds. An ambient isotopy, with inverse reverse flow,
carries the sphere and this normal neighbourhood into a standard
six-ball.

For an unknotted \(N_{\rm tub}=S^2\times D^4\) in that ball, let
\(C_0=D^6\setminus\mathring N_{\rm tub}\). Mayer--Vietoris gives
\(H_3(C_0)=\mathbb Z\) and \(H_5(C_0)=\mathbb Z\), with
all other positive groups zero. The inner boundary
\(S^2\times S^3\) maps isomorphically onto these two groups
and onto \(H_2(N_{\rm tub})\). The outer \(S^5\) class is retained in
degree five. The complement is simply connected by relative
general position, since a filling two-disc avoids a two-sphere
in dimension six.

Glue the other normal bundle by the link map induced by the matrix-cone comparison.
On each of the three nonzero groups of the inner boundary that
map acts by a unit. Mayer--Vietoris therefore makes the replacement
\(B'\) acyclic, and van Kampen makes it simply connected.
Cap it with the original complementary parametrized standard
\(D^6\). The resulting closed manifold is a homotopy six-sphere.
The dimension-six calculation \(\Theta_6=0\) and the h-cobordism/diffeomorphism comparison in \cite[pp.~504--505 and 512]{manKM63}
give an orientation-preserving diffeomorphism to \(S^6\).

Its image of the attached parametrized disc can be isotoped to
the original complementary disc with its entire parameterization.
For an embedding \(e_{\rm emb}\), shrink via \(e_{\rm emb}(r_{\rm emb}x)\) for positive small
\(r_{\rm emb}\); in a chart interpolate its rescaled map to its invertible
derivative, using the least-singular-value bound to preserve
injectivity; move its centre and positive frame along a path, and
expand the standard disc. The cutoff velocity-field construction of \cite[Theorem 5.8 and proof, pp.~62--64]{mapsMilnor65} extends this isotopy; the reverse
flow is its inverse. The composite diffeomorphism is then exactly
the prescribed map on the cap disc. Restricting to the complement
and straightening a collar gives an oriented \(B'\to D^6\)
relative to the prescribed outer boundary, with inverse obtained
by restriction of the inverse composite.

Conjugating by the sphere isotopy gives an oriented diffeomorphism
between the manifolds before and after this flop, relative to the
transported outer boundary of the chosen ball. Composing the finitely
many maps and their reverse-order inverses gives a global oriented
diffeomorphism between the specified Jordan-fan threefold and the threefold
carrying \(\mathcal F(I)\). In constructing the latter representative,
the other fillings and all their attaching data are those of the
original construction. This proves the smooth fan transport used in
the classification.

\medskip\noindent\emph{The oriented smooth classification for every simply connected \(b=1\) threefold.}

Work on the physical-fan representative obtained by
the fan-transport argument, and compose with that section's oriented diffeomorphism
at the end. Write \(M=\prod_i m_i\) and
\(G_{\rm hom}=(\mathbb Z/M)^2\). The homology calculation and the integral torsion reduction give
\(\pi_1(X)=0\), \(H_2(X)=H_3(X)=G_{\rm hom}\), and all other reduced
homology below degree six zero.

\smallskip\noindent\emph{The global circle and all its stabilizers.}

Multiplication of the section-line coordinate by \(t\in S^1\)
commutes with every line transition, the homogeneous contracting
map, each affine logarithmic map, and the period units. It therefore
gives a single smooth action with physical period \(\delta\).
This is the original action on the unchanged pieces and its toric
extension on the physical fan.

At each unselected higher fibre in these 17 rows, \(\nu\mid a\).
The first rank-one factorization has short vector
\(\lambda_{\rm sh}=(u,-a\kappa/\nu,0)\). The elliptic projection \(u\)
is primitive, so \(\lambda_{\rm sh},\delta\) span a primitive plane.
For example an elliptic covector \(u^*\) taking \(u\) to one
makes \(z_\delta+(a\kappa/\nu)u^*(x)\) take \(\delta\) to
one and \(\lambda_{\rm sh}\) to zero. Its restriction proves that
\(\delta\) remains primitive in
\(\ker(\mathsf T-I)/\mathbb Z\lambda_{\rm sh}\), the ruled base lattice.
Its translation on that elliptic base is free, so the action is
free on the whole rank-one filling, including the identified
sections. The same conclusion applies at \(I_1\) and to the
permitted original caps through the explicit circle-preserving map
of the original-cap comparison.

At a regular logarithmic central torus, simple connectivity gives
\(\gcd(m_i,n_i)=1\). Use the pair chosen in the integral torsion reduction,
\(\alpha_{i,\rm crt} n_i+\beta_{i,\rm crt} m_i=1\). The integer covector on the
five-generator lattice in the central presentation

\[
 F_{i,\rm lat}(x_1,x_2,x_\delta,x_\gamma,k)
       =x_\delta-z_i\alpha_{i,\rm crt} x_\gamma+z_i\beta_{i,\rm crt} k
\]

kills \((-v_i,m_i)\) and takes \(\delta\) to one. It is an
explicit primitive quotient coordinate; its B\'ezout completion
and inverse give a product splitting with that circle. Thus the
action is free there. Away from the central torus a nontrivial
root rotation cannot fix the root coordinate, and the line circle
is primitive on the regular fibre. This proves freeness for every
site, including order one and either sign or zero of a permitted
numerator.

On the physical cusp fan the rays are \((v,1)\), and the two
triangle types are
\((v,v+e_1,v+e_2)\) and
\((v+e_1+e_2,v+e_1,v+e_2)\). The cocharacter is
\((0,1,0)\). Inverting their ray-column matrices gives weights
\((-1,0,1)\) and \((1,-1,0)\). The residual edge weights for
directions \(e_1,e_2,e_2-e_1\) are \(1,0,-1\).
Consequently there is one fixed complete curve \(C_{\rm fix}\), from the
\(e_2\)-edge orbit, and every other stabilizer is trivial.
The embedded-orbit construction in the fan transport proves that this curve
is one \(\mathbb P^1\), with its two endpoints distinct. The
period translations have no stabilizer of a finite fan cell and
their units commute with the action, so the quotient introduces
no additional isotropy.

\smallskip\noindent\emph{The exact scalar normal framing.}

For this edge the two charts are

\[
 (x,y,w)=(q/(X_1X_2),X_2,X_1),\qquad
 (x',y',w')=(q/X_2,X_1X_2,1/X_1).
\]

Their inverses are \((X_1,X_2,q)=(w,y,xyw)\) and
\((1/w',y'w',x'y'w')\). Thus their transition is the two curve charts.
The normal directions have weights \(-1,+1\); in the order
\((y,x)\) this is
\(\mathcal O(-1)_{+1}\oplus\mathcal O(-1)_{-1}\).
Multiplication by \(\sqrt{1+|w|^2}\), with inverse division,
replaces the equatorial multiplier by its unit phase. Conjugate
the negative-weight coordinate and order the two resulting
positive-weight coordinates so that the transition is
\(\operatorname{diag}(e^{-i\theta},e^{i\theta})\).
It extends over a disc through

\[
 U_{\rm fr}(r,\theta)=\begin{pmatrix}
 \sin\beta_{\rm fr}(r)e^{-i\theta}&-\cos\beta_{\rm fr}(r)\\
 \cos\beta_{\rm fr}(r)&\sin\beta_{\rm fr}(r)e^{i\theta}
 \end{pmatrix}\in SU(2),
\]

where \(\beta_{\rm fr}(r)=r\) near zero and \(\beta_{\rm fr}=\pi/2\) near the
boundary. The functions \(\sin r\,e^{\pm i\theta}\) are smooth
at zero, and the inverse is \(U_{\rm fr}^*\). Adjusting the southern
frame by this extension and its inverse gives a scalar circle
action on a trivial real normal \(D^4\) bundle. Conjugation
reverses normal orientation; reverse the core parameterization
as well, so the total map preserves the complex orientation.
An invariant normal exponential and its inverse identify the
neighbourhood with \(C_{\rm fix}\times D^4\), with scalar action on the
second factor. This fixes the normal framing used throughout
the reconstruction.

\smallskip\noindent\emph{The smooth quotient, including the even-M spin calculation.}

The normal quotient is explicitly

\[
 h_{\rm Hopf}(z_1,z_2)=(2\Re(z_1\bar z_2),2\Im(z_1\bar z_2),
                             |z_1|^2-|z_2|^2),
 \quad \mathcal H(z)=h_{\rm Hopf}(z)/|z|\ (z\ne0),\quad\mathcal H(0)=0.
\]

It identifies exactly the circle orbits and preserves radius.
Use it to specify the smooth quotient chart \(C_{\rm fix}\times D^3\).
The inverse orbit charts away from the fixed core are supplied
by the Hopf sections

\[
 s_N(v)=\left(\sqrt{(1+v_3)/2},
                  \frac{v_1-iv_2}{\sqrt{2(1+v_3)}}\right),
 \quad
 s_S(v)=\left(\frac{v_1+iv_2}{\sqrt{2(1-v_3)}},
                  \sqrt{(1-v_3)/2}\right),
\]

namely \((q,t)\mapsto |q|t s_N(q/|q|)\) and its southern
version. These charts glue to the free orbit charts on a full
collar. Orient the free quotient by putting its horizontal
five-plane before the positive circle direction; this extends
over the Hopf quotient map.

Let \({\mathcal Q}=X/S^1\), let
\(X_{\rm free}=X\setminus(C_{\rm fix}\times\mathring D^4)\), and let
\({\mathcal Q}_0={\mathcal Q}\setminus(C_{\mathcal Q}\times\mathring D^3)\). The principal
circle bundle \(p_{\rm quot}:X_{\rm free}\to {\mathcal Q}_0\) has Euler class \(e_0\).
Choose the oriented normal meridian sphere \(\sigma\) so its
Hopf transition has winding plus one. Then
\(\langle e_0,\sigma\rangle=1\).

Relative transversality off the real codimension-four fixed curve
makes \(j:X_{\rm free}\to X\) an isomorphism on \(\pi_1,\pi_2,H_2\).
The bundle sequence and Hurewicz give

\[
 \pi_1({\mathcal Q}_0)=0,\qquad
 0\longrightarrow G_{\rm hom}\xrightarrow{(p_{\rm quot})_*j_*^{-1}}H_2({\mathcal Q}_0)
       \xrightarrow{\langle e_0,-\rangle}\mathbb Z
       \longrightarrow0,
 \qquad k\longmapsto k\sigma
\]

as a specified splitting. Excision for \(({\mathcal Q},{\mathcal Q}_0)\) gives
\(H_3=\mathbb Z\) with boundary \(\sigma\), and \(H_2=0\).
Hence, for the inclusion \(i:{\mathcal Q}_0\to {\mathcal Q}\),

\[
 \Theta=i_*(p_{\rm quot})_*j_*^{-1}:H_2(X)\xrightarrow{\cong}H_2({\mathcal Q})=G_{\rm hom}.
\]

To invert it, lift a class to \(a_{\rm lift}\in H_2({\mathcal Q}_0)\), subtract
\(\langle e_0,a_{\rm lift}\rangle\sigma\), lift the resulting sphere
through its Euler-zero pulled-back circle bundle, and apply
\(j_*\). Different first lifts differ by multiples of
\(\sigma\), and different sections differ by maps
\(S^2\to S^1\), so this inverse is well-defined. Van Kampen
gives \(\pi_1({\mathcal Q})=0\).

In the normal residue sequence, the map after
\(H^2({\mathcal Q}_0;\mathbb Z)\to H^0(C_{\mathcal Q};\mathbb Z)\) sends one to
the Poincar\'e dual of \([C_{\mathcal Q}]\). Its value on \(e_0\) is one,
so \([C_{\mathcal Q}]=0\). Moreover \(H^2({\mathcal Q};\mathbb Z)=0\), because
\(H_1({\mathcal Q})=0\) and \(H_2({\mathcal Q})=G_{\rm hom}\) is finite. The residue map is
therefore an isomorphism, with inverse \(k\mapsto ke_0\).

The complex total space is spin: \(H^2(X;\mathbb Z)=0\)
gives \(c_1(X)=0\), whose reduction is \(w_2(X)\).
On its free part, \(TX_{\rm free}=p_{\rm quot}^*T{\mathcal Q}_0\oplus\mathbb R\).
The mod-two Gysin sequence makes the kernel of \(p_{\rm quot}^*\) in
degree two the span of \(e_0\bmod2\). This nonzero class
evaluates to one on \(\sigma\), while \(w_2({\mathcal Q}_0)\) evaluates
to zero, since \(T{\mathcal Q}_0|_\sigma=TS^2\oplus\mathbb R^3\).
Therefore \(w_2({\mathcal Q}_0)=0\). Excision makes
\(H^2({\mathcal Q};\mathbb Z/2)\to H^2({\mathcal Q}_0;\mathbb Z/2)\) injective,
so \({\mathcal Q}\) is spin. This argument includes every even \(M\) and its mod-two torsion.

\smallskip\noindent\emph{The exact base identification and the gluing maps.}

The torsion linking form on the simply connected spin five-manifold \({\mathcal Q}\) is alternating by \cite[Lemma F, p.~372]{manBarden65}. Choose \(x\in G_{\rm hom}\) of order \(M\)
and \(y\) with \(\lambda(x,y)=1/M\). Nondegeneracy supplies
such a \(y\). Pairing a relation with \(x,y\) proves that

\[
 (r,s)\mapsto rx+sy,\qquad
 z\mapsto\bigl(\operatorname{num}_M\lambda(z,y),
                   -\operatorname{num}_M\lambda(z,x)\bigr)
\]

are inverse maps, where
\(\operatorname{num}_M([k/M])=k\pmod M\).
They identify the full linking form with the standard alternating
form; for \(M=1\) these are the unique maps on the zero group.
\cite[Theorems 2.2--2.3, pp.~376--377]{manBarden65} gives an oriented diffeomorphism
\(f_0:{\mathcal Q}\to N_M\) inducing the linking coordinates, with inverse inducing its
inverse coordinates. Here \(N_1=S^5\).

Let \(C_*\) be the standard parametrized fixed two-sphere in
the model defined below. Since \({\mathcal Q}\) is simply connected and
\([C_{\mathcal Q}]=0\), Hurewicz makes its parametrization nullhomotopic.
\cite[Theorem B(b), pp.~367--368]{manBarden65} applies in dimension five, where \(10>9\) and its required connectivity is satisfied.
It gives an isotopy of \(f_0(C_{\mathcal Q})\) to \(C_*\). Extend its
velocity by a cutoff to an ambient flow \(\mathcal I_t\), with reverse
flow as inverse. Replace \(f_0\) by
\(f=\mathcal I_1f_0\); its inverse is \(f_0^{-1}\mathcal I_1^{-1}\).

Tubular-neighbourhood adjustment makes \(f\) an oriented
orthogonal normal map \(R_{\rm norm}:C_{\mathcal Q}\to SO(3)\). Path lifting from
one chosen point on the simply connected sphere gives a lift
\(\widehat R_{\rm norm}:C_{\mathcal Q}\to SU(2)\) for the double cover characterized
by \(h_{\rm Hopf}(V_{\rm spin}z)=\rho_{\rm spin}(V_{\rm spin})h_{\rm Hopf}(z)\). Its conjugate transpose is the
inverse. The lift and the core map give an equivariant normal
diffeomorphism \(\Phi_N:C_{\rm fix}\times D^4\to C_*\times D^4\).
The normal-framing homotopy has no additional parameter because
\(\pi_2(SO(3))=0\).

Identify the free orbit bases by \(f\). Both principal bundles
have Euler residue one. The residue isomorphism after the quotient homology isomorphism
proves equality of their integral Chern classes. Their bundle
isomorphism can be written explicitly. Take a common good cover
with sections \(s_{i,\rm sec},t_{i,\rm sec}\), transition functions \(a_{ij},b_{ij}\),
and logarithms \(\ell_{ij}\) of \(b_{ij}/a_{ij}\).
Subtract an integral cochain to arrange
\(\ell_{ij}+\ell_{jk}=\ell_{ik}\); equality of the Chern
classes is exactly the obstruction calculation for this step.
For a partition of unity put \(f_{i,\rm gau}=\sum_k\rho_k\ell_{ik}\).
Then

\[
 \Phi_0(s_{i,\rm sec} z)=t_{i,\rm sec} e^{2\pi if_{i,\rm gau}}z,\qquad
 \Phi_0^{-1}(t_{i,\rm sec} z)=s_{i,\rm sec} e^{-2\pi if_{i,\rm gau}}z
\]

are well-defined inverse bundle maps. On the normal-boundary
collar, \(\Phi_N\) and \(\Phi_0\) differ by a circle-valued
gauge. The base there is \(S^2\times S^2\times[0,\varepsilon)\),
whose \(H^1\) is zero, so the gauge is \(e^{2\pi ik_{\rm gau}}\) for
a smooth real \(k_{\rm gau}\). A cutoff exponential modifies
\(\Phi_0\) to equal \(\Phi_N\) on a full collar; the reciprocal
exponential is its inverse modification. The maps and their
inverses therefore glue, preserving the chosen orientations.

The resulting model is

\[
\begin{gathered}
 M_-(N_M)=((N_M\setminus\mathring D^5)\times S^1)
      \cup_{g_{\rm glue}}(S^4\times D^2),\\
 g_{\rm glue}(x,z)=(\alpha_zx,z),\quad g_{\rm glue}^{-1}(x,z)=(\alpha_z^{-1}x,z),
\end{gathered}
\]

where \(g_{\rm glue}\) goes from complement to attached piece and
\(\alpha_{e^{2\pi it}}\) rotates the first real coordinate
plane of \(\mathbb R^5\) through \(2\pi t\).
The actions are \(t(x,z)=(x,tz)\) on the complement and
\(t(x,w)=(\alpha_tx,tw)\) on the attached piece. They agree
along \(g_{\rm glue}\). Its fixed sphere is
\(\{(0,0,u)\in S^4\}\) at \(w=0\), with scalar normal
weights. The boundary orbit coordinate is \(\alpha_z^{-1}x\),
whose composite with \(g_{\rm glue}\) is \(x\); thus its quotient restores
the disc in \(N_M\) with exactly the original boundary map.

The loop \(\alpha\) has spin lift
\(\cos(\pi t)+\sin(\pi t)e_1e_2\), with endpoints one and
minus one. It is the nonzero framing loop in \(\pi_1(SO(5))\).
The definition of \(\widetilde\Sigma_0N_M\) uses the opposite
gluing direction. The exact comparison uses

\[
 A_{\rm rot}=\operatorname{diag}(1,-1,-1,1,1),\qquad
 A_{\rm rot}\alpha_zA_{\rm rot}^{-1}=\alpha_z^{-1}.
\]

Let \(A_{{\rm rot},s}\) be rotation through \(\pi s\) in the \((2,3)\)
plane. On a complement collar use
\(F_{\rm col}(r,x)=(r,A_{{\rm rot},\zeta(r)}x)\), with \(\zeta=1\) at the
boundary and zero away from it. Its inverse uses
\(A_{{\rm rot},\zeta(r)}^{-1}\). The maps \(F_{\rm col}\times\mathrm{id}\)
on the complement and \(A_{\rm rot}\times\mathrm{id}\) on the cap
agree because \(A_{\rm rot}\alpha_z^{-1}=\alpha_zA_{\rm rot}\); their inverses
agree as well. Their determinants are positive. They give an
oriented map from the complementary gluing model to the precisely framed twisted suspension.

Composing the bundle gauges, the normal map, the gluing-direction comparison, and the finite fan
transport gives the promised oriented diffeomorphism

\[
                         X\cong\widetilde\Sigma_0N_M.
\]

The piecewise inverses match on the same collars. The construction verifies every action, spin and framing hypothesis for all seventeen rows and every simply connected regular-log datum.

For \(M=1\), the map from the complementary gluing model to
\(\partial(D^5\times D^2)\) is
\((x,z)\mapsto(\alpha_zx,z)\) on \(D^5\times S^1\) and
identity on \(S^4\times D^2\), with the opposite rotation
as inverse on the first piece. The complement disc has the
negative product-boundary orientation, while the cap has the
positive one. Thus these maps preserve the stated orientation.
Invariant corner rounding and radial projection give the standard
oriented \(S^6\). Different positive \(M\) give different smooth
types because the orders \(|H_2|=M^2\) differ.

The quotient calculation proves spin for every \(M\). Moreover, \(p_1=0\): on the free part it pulls back from
\(p_1({\mathcal Q})\in H^4({\mathcal Q};\mathbb Z)=H_1({\mathcal Q};\mathbb Z)=0\).
The kernel of \(H^4(X)\to H^4(X_{\rm free})\) is generated by
\(\operatorname{PD}[C_{\rm fix}]\), which is zero by the fan-transport argument.
Thus this restriction is injective and the whole class vanishes.
The initial fan representative inherits these conclusions through
the actual oriented diffeomorphism of the fan-transport argument.

\medskip\noindent\emph{The framed surgery.}
The scalar fixed curve is nullhomotopic and has the specified
real normal framing. The embedding isotopy used above places
it in a ball; \(\pi_2(SO(4))=0\) identifies the transported
framing with the product framing. Replace \(S^2\times D^4\)
by \(D^3\times S^3\), with Hopf action on the second factor.
In \(S^6=\partial(D^3\times D^4)\) the replacement is the
double of \(D^3\times S^3\), namely \(S^3\times S^3\).
The product boundary maps, conjugated by the isotopy and its
inverse, therefore give
\(Y\cong X\#(S^3\times S^3)\) with a free circle action.

On the quotient, the corresponding replacement of
\(S^2\times D^3\) by \(D^3\times S^2\) is made in the
fixed disc in \({\mathcal Q}\). The same product-boundary description in
\(S^5=\partial(D^3\times D^3)\) identifies
\(Y/S^1\cong N_M\#(S^2\times S^3)\).
Its Euler class is zero on the punctured \(N_M\), since
\(H^2(N_M;\mathbb Z)=0\), and evaluates to one on the new
Hopf sphere. Thus its oriented bundle pair is
\((N_M\#(S^2\times S^3),(0,x_{\rm Eul}))\), with
\(H_2= (\mathbb Z/M)^2\oplus\mathbb Z\) and primitive \(x_{\rm Eul}\).
The base is spin. Pull back its spin frame bundle and adjoin
the oriented vertical line using an invariant connection.
The circle action lifts to this pulled-back spin bundle, proving
that the spin structure on Y is projectable.
Extending the already constructed unstabilized map for X by
identity over the extra punctured product and its neck gives
\[
 Y\cong(S^3\times S^3)\#\widetilde\Sigma_0N_M.
\]
The inverse uses the inverse map for X and the same product
neck.
\end{proof}


\subsection{Sixteen finite and multiplicative profiles}
\begin{definition}[Marked filling patterns]\hypertarget{tag-Q302}{}\label{maps:def:marked-filling-pattern}
{\normalfont\scriptsize tag Q302}\par
A marked filling pattern consists of a marked elliptic profile,
an ordered section pair, a zero position and a permitted additive
filling choice at each additive fibre. Logarithmic orders and affine
powers are parameters of the chosen fillings.
\end{definition}

\begin{theorem}[Sphere fillings for sixteen elliptic profiles]\hypertarget{tag-02L1}{}\label{maps:src:D.2.26}
{\normalfont\scriptsize tag 02L1}\par

Consider the natural section-line rows of \ref{maps:src:D.2.22}
whose profiles contain no fibre $I_\nu^*$. Write $P=aG$, $Q=bG$,
and use the marking
$\Lambda=\mathbb Z\langle e_1,e_2,\delta,\gamma\rangle$ and
$\omega=a e_1^*\wedge e_2^*-b\delta^*\wedge\gamma^*$.
There are exactly $59$ marked initial/additive-filling patterns
that yield spheres, on $50$ distinct initial rows. They are listed
in the proof. Their counts are
\begin{equation}
\begin{array}{c|r|c|r|c|r|c|r}
43&3&43/II&1&45&4&45/II&4\\
47&4&47/II&4&47/III&4&49&3\\
49/II&1&49/III&2&55&4&55/II&4\\
56&5&56/II&6&56/III&5&56/IV&5 .
\end{array}
\hypertarget{tag-01ZN}{}\label{maps:eq:D.HY1}
\end{equation}
The zero has the index-one Mumford filling. At every other
multiplicative place of these patterns, either its original
filling or its primitive Mumford filling is permitted.

More generally, use original caps only at places where the line
has zero multidegree and the component translation has one orbit,
and use primitive Mumford caps at the other multiplicative places.
At a finite good-reduction log choose a fixed vector $\ell_i$
with $\gamma^*(\ell_i)=1$, and take
\[
 v_i=n_i\ell_i+r_i\delta,\qquad
 m_i=d_i k_i,\quad k_i>0,\qquad \gcd(m_i,n_i)=1,
\]
where $r_i$ is any integer and
\begin{equation}
\begin{array}{c|ccccc}
\text{finite type}&II&III&IV&III^*&IV^*\\ \hline
d_i\text{ (monodromy order)}&6&4&3&4&3\\
q_i\text{ (vertical index)}&1&2&3&2&3 .
\end{array}
\hypertarget{tag-01ZO}{}\label{maps:eq:D.HY2}
\end{equation}
Allow every affine lift with that power. These actions are free.
The existence of $\ell_i$ and the stated coprimality are necessary
for simple connectivity. At any finite number of additional regular
sites allow
\[
 m_j>0,\qquad v_j=x_je_1+y_je_2+z_j\delta+n_j\gamma,\qquad
 \gcd(m_j,x_j,y_j,z_j,n_j)=1 .
\]
Let $e$ be the peripheral integer in the Seifert presentation
\eqref{maps:eq:D.FG3}.

Define $C$ by a generator $s$, one generator $D_i$ at each finite
log and a homomorphism $t_j:\Lambda\to C$ at each regular log,
with relations
\begin{equation}
\begin{aligned}
(m_i/q_i)D_i+(bn_i/q_i)s&=0 &&\text{at finite logs},\\
m_jt_j(u)-\omega(u,v_j)s&=0,\qquad t_j(v_j)=0
                                      &&\text{at regular logs},\\
\sum_iD_i+\sum_jt_j(\delta)+be\,s&=0 .
\end{aligned}
\hypertarget{tag-01ZP}{}\label{maps:eq:D.HY3}
\end{equation}
All coefficients are integral. Put $T=\operatorname{Tor}C$ and
write the abelianization of the complete Seifert presentation as
$H_1(X)=\mathbb Z^r\oplus A_1$, where $r\in\{0,1\}$.
Then
\[
\begin{aligned}
H_0(X)=H_6(X)&=\mathbb Z,& H_5(X)&=\mathbb Z^r,\\
H_2(X)=H_3(X)&=\mathbb Z^{2r}\oplus T,&
H_4(X)&=\mathbb Z^{2r}\oplus A_1 .
\end{aligned}
\]
These formulas apply to every member with the stated finite
coprimality, including those with nontrivial fundamental group.

For a simply connected member set
\[
Q_f=\prod_{i\text{ finite}}q_i,\quad c=b/Q_f,\quad
R=\prod_{j\text{ regular}}m_j,\quad
X_0=\sum_j(R/m_j)x_j,\quad Y_0=\sum_j(R/m_j)y_j .
\]
Empty products are one. Then $c$ is a positive integer and
\begin{equation}
H_2(X)=H_3(X)=
\mathbb Z/g\oplus\mathbb Z/R\oplus\mathbb Z/(cR/g),
\qquad g=\gcd(c,R,X_0,Y_0).
\hypertarget{tag-01ZQ}{}\label{maps:eq:D.HY4}
\end{equation}
The complete sphere criterion, with the fundamental group computed
by \eqref{maps:eq:D.FG3}, is
\begin{equation}
\pi_1(X)=1,\qquad
b=Q_f,\qquad m_j=1\ \text{at every regular log}.
\hypertarget{tag-01ZR}{}\label{maps:eq:D.HY5}
\end{equation}
For one finite log every positive multiple of its monodromy order
occurs. For two, the multiplicities must be coprime and their
Seifert determinant must be $\pm1$. The only sphere pattern with
two finite logs is the $IV^*,III,I_1$ profile, with both additive
places logarithmic. The integers $r_i$ are unrestricted.
Every one of the twelve additive profiles in \eqref{maps:eq:D.HY1}
admits infinitely many sphere multiplicities. The four
all-multiplicative profiles have order one at all regular sites.
For every simply connected member with $c=1$, the oriented smooth
type is
\[
                         X\cong\widetilde\Sigma_0N_R,
\]
with $N_R$ and the boundary rotation defined in
\ref{maps:src:D.2.14}.
\end{theorem}


\begin{proof}
Put $\chi=\gamma^*$ and
$J=\left(\begin{smallmatrix}0&1\\-1&0\end{smallmatrix}\right)$.
Use exterior order $(12,13,14,23,24,34)$ and the complex orientation.
The positive $\gamma$ period covers translation by $-Q$.

\emph{Natural rows and necessary cap choices.}
Let $1/d$ be the generator height, and select a multiplicative
$I_\nu$ with $\mathsf A-I=\nu uu^tJ$, $g_{\rm per}=\kappa u$ and
$\det(u,r_{\rm per})=1$. The component equation of
\ref{maps:src:D.2.22} is
\[
\begin{gathered}{}
ab=d/\nu,\qquad \gcd(b\kappa,\nu)=1,\qquad
\\{}
[j=a\kappa]-[j=0]-[j=a\kappa-b\kappa]+[j=-b\kappa]
 =-2[j=h]+[j=h-1]+[j=h+1].
\end{gathered}
\]
The indices are modulo $\nu$, counting the two neighbours
separately at $\nu=2$. At every other reducible fibre $P$ or $Q$
is narrow. For a finite monodromy matrix this is the integrality
of $a(\mathsf A-I)^{-1}g_{\rm per}$ or $b(\mathsf A-I)^{-1}g_{\rm per}$.
Listing every divisor $a$ of $d/\nu$, with $b=d/(\nu a)$,
and testing the displayed component columns gives $172$ natural
rows. Each has $\gcd(a,b)=1$ and zero component $C_0$.
Exactly $148$ belong to the sixteen profiles here.

We also give the global obstruction to other original caps.
The original elliptic neighbourhood $N_{\rm ell}$ retracts to its central
components. At an additive place it has $H_2(N_{\rm ell})=\mathbb Z^\rho$
and is simply connected; at $I_\nu$ its reduced cycle has
$H_1=\mathbb Z$, $H_2=\mathbb Z^\nu$. The component degrees
of $\mathcal O(P-O)$ are zero or the primitive vector $e_p-e_0$.
For nonzero degree, the circle-bundle homotopy sequence kills
$\delta$; the Hopf bundle over a degree-one component supplies
its bounding disc. Naturalness then makes $Q$ narrow.
The mapping-torus Wang sequence gives cap ranks $b_1=1$ at
an additive place and $b_1=2$ at $I_\nu$.
For zero degree they are respectively two and three.
The component-circle invariant cycles give $b_4=\rho$ in the
nonzero-degree case and the number of component orbits in
the zero-degree case.

The boundary $H_1$ map is onto, by the fibre and transverse
origin-disc maps. Poincar\'e--Lefschetz duality gives the local
free degree-two defect
\[
b_4(V)-b_1(\partial V)+b_1(V).
\]
It is $\rho-2$ for a nonzero-degree additive cap, $\nu-2$
for a nonzero-degree multiplicative cap, and one less than
the number of component orbits for a zero-degree cap.
If $k_{\rm caps}$ original caps have nonzero degree, the degree-one
Mayer--Vietoris kernel has rank $k_{\rm caps}+b_1(X)$: each such cap
lowers its $b_1$ by one, while the boundary and exterior
ranks stay fixed. The degree-two cokernel maps onto the
direct sum of the local defects by projection to the cap
coordinates. Therefore
\[
b_2(X)\ge k_{\rm caps}+b_1(X)+\sum_i
\dim_{\mathbb Q}\operatorname{coker}
       [H_2(\partial V_i;\mathbb Q)\to H_2(V_i;\mathbb Q)].
\]
This excludes all nonzero-degree original caps in a sphere,
including III and $I_2$, and all zero-degree caps with
more than one orbit.

The one-orbit additive originals are II and the transitive
component translations of III and IV. The nontrivial
translations of III* and IV* have respectively five and
three component orbits. At $I_\nu$ the orbit count is
$\gcd(\nu,b\kappa)$. At the selected zero a root cover of
order $m$ has $m^2$ vertices, $3m^2$ edges and $2m^2$
triangles in its periodic central complex. The normalization
and conductor give Euler characteristic
$6m^2-6m^2+2m^2=2m^2$; a free order-$m$ quotient gives
$2m$. All other caps have Euler characteristic zero.
A sphere thus has the index-one cusp.

\emph{The period bundle and integral graph.}
Let $\mathsf A_i$ denote the marked elliptic return at site $i$.
For each marked pair $(\mathsf A_i,g_i)$ put
$z_i=(\mathsf A_i-I)^{-1}g_i$, using $z_i=(\kappa_i/\nu_i)r_{i,\rm per}$
at a multiplicative place. Set
\[
\beta_i=-a g_i^tJ\mathsf A_i,\qquad
T_i=\begin{pmatrix}
\mathsf A_i&0&bg_i\\
\beta_i&1&\beta_i bz_i+w_i\\
0&0&1
\end{pmatrix},
\]
where $w_i=1/\nu_i$ at the zero and zero elsewhere.
The Poincar\'e-line maps of \ref{maps:src:D.2.12}
give these matrices and their period-coordinate inverses.
They are integral, $\prod_iT_i=I$, and preserve
$\delta,\chi,\omega$. On the cut regular base the map
$[x]\mapsto[\Pi(t)x]$ identifies the marked real torus
with the fibre; the group logarithm modulo $\Pi(t)\Lambda$
is its inverse. The period functions include the line units.

With $N$ singular places let
\[
D=[T_1-I\ \cdots\ T_{N-1}-I].
\]
For $i<N$, $j_i$ is block inclusion, while the last based
boundary has
\[
j_N(x)=-(T_2\cdots T_{N-1}x,\ldots,T_{N-1}x,x),
\qquad Dj_N(x)=(I-T_N^{-1})x.
\]
Choose integral fixed bases $F_i$ beginning with $\delta$.
They are constructed by Smith reduction of $T_i-I$:
the kernel columns and their inverse rows give the fixed
lattice, and a B\'ezout change places the primitive $\delta$
first. Integer elimination on the marked tuples gives, for
$B=[j_1F_1\ \cdots\ j_NF_N]$,
\[
DB=0,\qquad \operatorname{rank}D=3,\qquad
\operatorname{Smith}(B)=I_{4N-7},\qquad
\ker B=\mathbb Z\left(\sum_i e_{\delta,i}\right).
\]
There are $4N-6$ columns. These equalities hold on all
$172$ natural rows, hence on the $148$ used here.
To specify the inverse, write $U_BBV_B=\operatorname{diag}(I,0)$.
For $y\in\ker D$ take the first $4N-7$ coordinates of
$U_By$, append zero, and apply $V_B$. Its class modulo
the displayed diagonal relation is the inverse of $B$.
The unit entries prove equality with the entire integral
graph kernel.

The exterior-square relation matrix has five unit entries.
Its kernel covector is the primitive
$\omega=(a,0,0,0,0,-b)$. Choose $au_{\rm fib}-bv_{\rm fib}=1$ and
$s_0=u_{\rm fib}e_{12}+v_{\rm fib}e_{34}$. The inverse quotient maps are
\[
\bigwedge^2\Lambda/\sum_i\operatorname{im}(\bigwedge^2T_i-I)
\ \underset{k\mapsto[ks_0]}{\overset{[x]\mapsto\omega(x)}{\rightleftarrows}}\
\mathbb Z .
\]
Smith coordinates give an integral preimage of
$x-\omega(x)s_0$ in the relation matrix, hence its
monodromy-cylinder bounding chain. Write $s=[s_0]$.

A fixed integral loop sweeps its boundary circle to a
two-torus. Addition uses the real-period triangle with
vertices $0,u,u+v$; subtraction is the inverse.
The physical $\delta$-circle over the cut base is a
three-chain with boundary the sum of its sweeps.
Changes of its line frame contribute $\delta\wedge\delta=0$.
The reduced origin frames of \ref{maps:src:D.2.12}
fix this marking at the cusp. These cycles split
\[
0\to\mathbb Zs\to H_2(E)\to\ker D\to0.
\]
The inverse projects to the graph, applies the inverse of
$B$, subtracts those sweeps and evaluates the remaining
fibre coordinate by $\omega$.

An additional regular disc inserts a zero block in $D$,
and a block $-x$ in $j_N$. Removing these new coordinates,
using the preceding inverse and restoring them by block
inclusion gives the inverse for any number of discs.
The sole relation is the sum of all $\delta$ sweeps.
For the peripheral convention $\prod c_i=\gamma^e$,
the attaching polygon adds
$-\omega(\delta,e\gamma)=be$, giving that relation with
the term $be\,s$.

\emph{Integral maps for the reduced caps.}
In the following local cap formulas write
$(\mathsf A,g_{\rm per},\mathsf T,\beta_{\rm row})=(\mathsf A_i,g_i,T_i,\beta_i)$ for the marked data at that site.

At a permitted original additive place put $q_{\rm vert}=1,2,3$
for II, III, IV, respectively. Write $z=(\mathsf A-I)^{-1}g_{\rm per}$,
$h_{\rm per}=az\in\mathbb Z^2$, $z_Q=bz$, and $\alpha_{\rm line}=-h_{\rm per}^tJ$.
The line map and its inverse are
\[
S_P(x,z_\delta,x_\gamma)=(x,z_\delta+\alpha_{\rm line}x,x_\gamma),
\qquad
S_P^{-1}(x,z_\delta',x_\gamma)
                         =(x,z_\delta'-\alpha_{\rm line}x,x_\gamma).
\]
The identity $\alpha_{\rm line}(\mathsf A-I)+\beta_{\rm row}=0$ conjugates ${\mathsf T}$ to
the product of the line circle with
$\left(\begin{smallmatrix}\mathsf A&bg_{\rm per}\\0&1\end{smallmatrix}\right)$.
Zero component degrees trivialize the line on the
four-dimensional neighbourhood $N_{\rm ell}$.

In that frame the contraction has a positive roof $\varrho(y)$.
Choose $k_{\rm roof}\max\varrho<1$ and a smooth middle-supported bump
$b_y$ of integral one. The increasing function
\[
f_y(t)=\int_0^t\bigl(k_{\rm roof}+(1-k_{\rm roof}\varrho(y))b_y(v)\bigr)\,dv
\]
has matching endpoint derivatives and inverse $f_y^{-1}$.
It identifies the quotient with the mapping torus of
the angular lift. The remaining circle bundle is
trivial: its restriction to $N_{\rm ell}$ is trivial and the
Wang restriction in $H^2$ is injective because
$H^1(N_{\rm ell})=0$. \v{C}ech logarithms and a partition of unity
give its trivialization; reciprocal gauges give the inverse.

The forward torsion section and its identity-component
logarithm are those of
\cite[\S2.4, equation (2.2), p.~6]{manUU21}.
Translation by scalar multiples of this logarithm,
with its negative as inverse and a cutoff at the
ends of the suspension interval, identifies the
mapping torus with the order-$q_{\rm vert}$ component translation.
It permutes the $q_{\rm vert}$ rational components transitively.
The core inclusion induces a homotopy equivalence
of mapping tori; the homology calculation uses this
inclusion and the component permutation on $H_2(N_{\rm ell})$.
Coinvariants sum component coefficients, with a
section on the first component; invariants repeat
a coefficient, with inverse reading the first.
Consequently
\[
H_1(V)=\mathbb Z\delta\oplus\mathbb Z\gamma,\qquad
H_2(V)=\mathbb Z[C_{\rm comp}]\oplus\mathbb Z[\delta\gamma].
\]
The ordinary elliptic fibre represents $q_{\rm vert}[C_{\rm comp}]$.

In the $S_P$ coordinates the full fibre map is
\[
f_{2,\rm add}(x)=\bigl(q_{\rm vert}(x_{12}+z_{Q,2}x_{14}-z_{Q,1}x_{24}),x_{34}\bigr).
\]
The first coordinate follows by the elliptic translation
$x\mapsto x+z_Qx_\gamma$, inverted by subtraction,
and the value $q_{\rm vert}[C_{\rm comp}]$ of the elliptic area.
The second is the cup product of the two $H_1$ coordinates.
Choose integers $A_{\rm Bez},B_{\rm Bez},C_{\rm Bez}$ with
$q_{\rm vert}A_{\rm Bez}+q_{\rm vert}z_{Q,2}B_{\rm Bez}-q_{\rm vert}z_{Q,1}C_{\rm Bez}=1$. Then
\[
[C_{\rm comp}]\mapsto\bigwedge^2S_P^{-1}(A_{\rm Bez}e_{12}+B_{\rm Bez}e_{14}+C_{\rm Bez}e_{24}),
\qquad [\delta\gamma]\mapsto e_{34}
\]
is a right inverse. Such integers exist because
$Q$ generates the component group.
Also $\operatorname{im}(\mathsf A-I)+\mathbb Zbg_{\rm per}=\mathbb Z^2$.
Thus the three-dimensional monodromy has two unit
entries in its difference from identity; oriented
volume identifies its exterior square with its
inverse transpose. It follows that
$\bigwedge^2{\mathsf T}-I$ has four unit entries and the
displayed fibre quotient is exactly $\mathbb Z^2$.

The fixed loops are $\delta$ and
$\ell_0=(-q_{\rm vert}z_{Q,1},-q_{\rm vert}z_{Q,2},0,q_{\rm vert})$.
Their circles extend over the origin disc and the
$q_{\rm vert}$-fold forward orbit tube, respectively; the
closing N\'eron logarithm and the line frame extend
over the disc. Their normal sweeps bound and give
the entire integral boundary kernel.

For a permitted original $I_\nu$, apply $S_P$ with
$h_{\rm per}=(a\kappa/\nu)r_{\rm per}$. Choose $l\nu+k_{\rm Bez}b\kappa=1$ and
replace its long periods by
$t_{\rm per}=lr_{\rm per}+k_{\rm Bez}\gamma$, $w=-b\kappa r_{\rm per}+\nu\gamma$.
The inverse is the inverse of
$\left(\begin{smallmatrix}l&-b\kappa\\k_{\rm Bez}&\nu\end{smallmatrix}\right)$.
The first period has valuation one and the second
valuation zero. Exponentiating the primitive short
period gives the charts
\[
X=z/q^j,\quad Y=q^{j+1}/z,\quad XY=q,
\qquad z=Xq^j,\quad q=XY.
\]
They are charts of the original cycle; the valuation-one
period identifies successive charts. The quotient has
one ruled component with its zero and infinity sections
identified. The inverse period matrix and these chart
inverses identify it with the primitive rank-one cap,
including its nonzero unit factors.

In a Jordan basis $\mathsf H_{\rm nod}$ for that cap, its fibre maps are
\[
\begin{gathered}
f_{1,\rm nod}(x)=(x_2,x_3,x_4),\\
f_{2,\rm nod}(y)=(y_{12},y_{23},y_{24},y_{34}),
\end{gathered}
\]
preceded by $\mathsf H_{\rm nod}^{-1}$ and $\bigwedge^2\mathsf H_{\rm nod}^{-1}$.
Their sections insert the omitted zero coordinates
and apply $\mathsf H_{\rm nod}$. The ruled normalization and its
identified sections give these integral generators.
The fixed sweeps bound the toric normal chains.

At the cusp put $S=(u,\delta)$, $\mathsf L_{\rm per}=(r_{\rm per},\gamma)$ and
\[
\begin{gathered}
\mathsf C_{\rm cusp}=\begin{pmatrix}
\nu&b\kappa\\-a\kappa&(1-ab\kappa^2)/\nu
\end{pmatrix},\\
\mathsf H_{\rm cusp}=(S\mathsf C_{\rm cusp}e_1,\mathsf L_{\rm per}e_1,S\mathsf C_{\rm cusp}e_2,\mathsf L_{\rm per}e_2).
\end{gathered}
\]
The component equation makes $\mathsf C_{\rm cusp},\mathsf H_{\rm cusp}$ unimodular and
$\mathsf H_{\rm cusp}^{-1}\mathsf T\mathsf H_{\rm cusp}=I+E_{12}+E_{34}$.
In this specified short fan, one chart and its inverse are
$(x,y,w)=(q/(XY),Y,X)$ and $(X,Y,q)=(w,y,xyw)$.
The normalization and conductor calculation of
\ref{maps:src:D.2.12}, transported by these period maps, gives
\[
\begin{gathered}
f_{1,\rm cusp}(x)=(x_2,x_4),\\
f_{2,\rm cusp}(y)=(y_{12},y_{14}-y_{23},y_{24},y_{34}).
\end{gathered}
\]
Its section is $(u,v,w,z)\mapsto(u,0,v,0,w,z)$ followed
by $\bigwedge^2\mathsf H_{\rm cusp}$. Its primitive kernel,
$\langle e_{13},e_{14}+e_{23}\rangle$, is the local Wang
image. The two fixed sweeps bound the normal chains.
All these reduced cap maps are onto integrally.

\emph{Finite affine lifts and their full kernel.}
At a finite place use the vertical index $q_{\rm vert}=q_i$
and finite return order $d_{\rm fin}=d_i$;
put $z=(\mathsf A-I)^{-1}g_{\rm per}$ and let $h_Q$
be the least common denominator of $bz$.
Solving $({\mathsf T}-I)x=0$ gives fixed basis
$\delta,(-h_Qbz_1,-h_Qbz_2,0,h_Q)$; the inverse
reads the third coordinate and divides the fourth by $h_Q$.
Thus $\chi(\operatorname{Fix}{\mathsf T})=h_Q\mathbb Z$ and
$h_Q\mid d_{\rm fin}$. If $h_Q>1$, it divides both the finite
multiplicity and the contracting numerator, excluding
trivial abelianization.

For $h_Q=1$ set $\ell=(-bz_1,-bz_2,0,1)$ and
$\Lambda_0=\ker\chi$. The invariant covector on
$\Lambda_0$ is $\rho=q_{\rm vert}(az_2,-az_1,1)$.
Its coordinates are coprime: $g_{\rm per}$ generates
$\mathbb Z^2/(\mathsf A-I)\mathbb Z^2$, and $\gcd(a,q_{\rm vert})=1$.
The two-by-two minors of the three-lattice relation
matrix have the same content. They are therefore
primitive, and B\'ezout supplies $\beta$ with $\rho\beta=1$.
This gives the inverse quotient maps and
\begin{equation}
(\Lambda_0)_{\mathsf T}=\mathbb Z\beta,\qquad
[\delta]=q_{\rm vert}[\beta],\qquad
\operatorname{Fix}({\mathsf T})=\mathbb Z\delta\oplus\mathbb Z\ell .
\hypertarget{tag-01ZS}{}\label{maps:eq:D.HY6}
\end{equation}
The finite matrix has exact order $d_{\rm fin}$ as in
\eqref{maps:eq:D.HY2}. Its average projects $\beta$
to $\delta/q_{\rm vert}$, while fixing $\ell$. Consequently
\[
N_m\beta=(m/q_{\rm vert})\delta,\qquad N_m\ell=m\ell,\qquad
-\omega(\beta,\ell)=b/q_{\rm vert} .
\]
The last identity follows by averaging the first
argument of the invariant alternating form.
It also proves $q_{\rm vert}\mid b$.

There are $62$ such primitive finite positions on
the $148$ nonstar rows, among $94$ finite positions.
The same local proof applies to the finite places of
OS46/II and OS50/II for $(a,b)=(1,4),(4,1)$, and
OS50/IV for $(a,b)=(1,12),(4,3)$. These are the
six additional primitive positions among the eight
finite places on those profiles. Hence the local
count on all $172$ natural rows is $68$ among $102$.

The conjugacy of holomorphic affine lifts can be
written over the whole good root disc. Let
$\sigma(u)=\zeta u$ and
$\mathcal A_{\rm deck}(u):E_u\to E_{\sigma u}$ be the linear deck map
on the holomorphic Lie bundle. On its sections set
\[
(\mathcal Rf)(u)=\mathcal A_{\rm deck}(\sigma^{-1}u)f(\sigma^{-1}u),\qquad
\mathcal N=\sum_{j=0}^{m-1}\mathcal R^j,\qquad
\mathcal R^m=I.
\]
Two affine lifts with the same integral power have
section expressions $\mathcal Rf+b_0$ and
$\mathcal Rf+b_1$, with $\mathcal N(b_1-b_0)=0$.
Put
\[
s_{\rm gau}=-\frac1m\sum_{j=0}^{m-1}j\mathcal R^j(b_1-b_0).
\]
Since $(I-\mathcal R)\sum j\mathcal R^j=\mathcal N-mI$,
translation by $s_{\rm gau}$ conjugates the first lift to the
second; subtraction is the inverse. Both extend
holomorphically over the disc. Changing a lift by
a period $\lambda$ changes its integral power by
$N_m\lambda$; aligning these representatives gives
the same formula. The two-periodic cyclic resolution
identifies the power-class obstruction with
$\Lambda^{\mathsf T}/N_m\Lambda$, and the invariant translation
$v/m$ represents its inverse. Thus all affine lifts
in the statement are included.

Write $v=n\ell+r_{\rm fin}\delta$ and choose
$\alpha_{\rm fin} m+\eta_{\rm fin} n=1$. The central quotient has
\[
H_1(F)=\langle\beta,\ell,c_{\rm quot}\mid mc_{\rm quot}=n\ell+r_{\rm fin}q_{\rm vert}\beta\rangle.
\]
The basis $\beta,z_{\rm quot}=\alpha_{\rm fin}\ell+\eta_{\rm fin} c_{\rm quot}$ has inverse
$\ell=mz_{\rm quot}-\eta_{\rm fin} r_{\rm fin}q_{\rm vert}\beta$, $c_{\rm quot}=nz_{\rm quot}+\alpha_{\rm fin} r_{\rm fin}q_{\rm vert}\beta$.
Thus $H_1(F)=\mathbb Z^2$. Averaging forms over its
finite torus cover gives Betti numbers $(1,2,2,2,1)$;
integral duality gives $H_2(F)=H_3(F)=\mathbb Z^2$.
The flat normal line has torsion Chern class and
free $H^2(F;\mathbb Z)$, so that class is zero.

The normal boundary coordinates and their inverse are
\[
[y,\tau]\mapsto[y+(\tau/m)v,e^{2\pi i\tau/m}],
\qquad [x,z]\mapsto[x-(\tau/m)v,\tau],
\]
where a root argument $\tau$ is chosen; another choice
changes both by the mapping-torus relation.
The positive normal circle is
$(y,\tau)\mapsto(y-tv,\tau+mt)$ and represents $m\mu-v$.
The transfer of a lattice loop $\lambda$ is the torus
$[s\lambda,e^{2\pi it}]$. Cutting it into $m$ sectors
gives area $-\lambda\wedge v$ and sweep $N_m\lambda$;
rejoining the sectors is the inverse chain map.
For $c_{\rm quot}$ use $[sv/m,e^{2\pi i(s/m+t)}]$, whose area is
zero and whose sweep is $v$.
The transfers of the displayed basis are therefore
\begin{equation}
\begin{split}
&\bigl(-\omega(\beta,v),\,t(N_m\beta)\bigr),\\
&\bigl(-\alpha_{\rm fin}\omega(\ell,v),\,
                    t(\alpha_{\rm fin} N_m\ell+\eta_{\rm fin} v)\bigr).
\end{split}
\hypertarget{tag-01ZT}{}\label{maps:eq:D.HY7}
\end{equation}
Thom exactness makes these the full boundary kernel.
They are
\[
(bn/q_{\rm vert})s+(m/q_{\rm vert})t(\delta),\qquad
-\alpha_{\rm fin} br_{\rm fin}s+t(\ell)+\eta_{\rm fin} r_{\rm fin}t(\delta).
\]
The second eliminates $t(\ell)$ with coefficient one;
substitution is its inverse. This gives the finite
relation of \eqref{maps:eq:D.HY3} for every $r_{\rm fin}$.

\emph{Regular logs and Mayer--Vietoris.}
For a regular vector $v$ the central lattice is
$\Lambda_v=\Lambda+\mathbb Z(v/m)
=\mathbb Z^5/\mathbb Z(-v,m)$.
Choose $\alpha_{\rm reg} m+\lambda_{\rm reg}(v)=1$.
The projection $I-(-v,m)(-\lambda_{\rm reg},\alpha_{\rm reg})$ and
inclusion of its kernel give inverse quotient
coordinates. The normal trivialization is
\[
[(y,z)]\mapsto([y],e^{-2\pi i\lambda_{\rm reg}(y)}z),
\qquad ([y],w)\mapsto[(y,e^{2\pi i\lambda_{\rm reg}(y)}w)] .
\]
The deck factors cancel by the B\'ezout identity.
The logarithmic collar isotopy of
\ref{maps:src:D.2.13} has opposite translation as
inverse and carries its positive normal to $m\mu-v$.
The same sector tori give the full kernel
$mt(u)-\omega(u,v)s,\ t(v)$.

Since every cap map in degree two is onto,
lift each cap class to its boundary and subtract
that lift from the exterior coordinate.
The inverse includes the exterior with zero
cap coordinates; two lifts differ by a local
kernel just computed. The reduced caps remove
their fixed sweeps, the finite caps remove
$t(\ell)$, and the regular sweeps remain.
The single global $\delta$ relation proves
\eqref{maps:eq:D.HY3}.
For a regular lift change
$v_j'=v_j+m_j\lambda_j$, the inverse coordinate
maps are
\[
e'=e+\sum_j\chi(\lambda_j),\qquad
t_j'(u)\mapsto t_j(u)+\omega(u,\lambda_j)s,\qquad
t_j(u)\mapsto t_j'(u)-\omega(u,\lambda_j)s .
\]
They preserve the global relation because
$\omega(\delta,\lambda_j)=-b\chi(\lambda_j)$.
At finite places a peripheral period change uses
$v'=v+N_m\lambda$, as proved above.

The boundary $H_1$ groups are free: every local
$T_i-I$ has unit nonzero Smith entries.
If $r_{i,\rm bd}=\operatorname{rank}\operatorname{coker}(T_i-I)$,
the boundary has rank $r_{i,\rm bd}+1$ and its minimal cap
rank $r_{i,\rm bd}$. The exterior rank is the number of
boundary circles. The next Mayer--Vietoris kernel
is thus free of rank $r=b_1(X)$.
Rational elimination in \eqref{maps:eq:D.HY3}
leaves
\[
b\left(e-\sum_i n_i/m_i\right)s=0 ,
\]
where the sum includes both kinds of logs.
This is the Schur complement of the Seifert
abelian matrix, so $C$ has free rank $r$.
Lifting a basis of the free degree-one kernel
to normal discs and exterior joining chains
splits the degree-two sequence, giving
$H_2=\mathbb Z^{2r}\oplus\operatorname{Tor}C$.

The index-one cusp has Euler characteristic two
and all other pieces have Euler characteristic
zero. Integral duality and universal coefficients
give the displayed homology groups. In particular
$b_3=2b_2-2b_1=2r$.
For a torsion Smith generator of order $d_{\rm Smith}$,
the character with value $1/d_{\rm Smith}$ followed by the
integral Bockstein and cap product gives its
dual torsion class; the inverse reads its
linking coordinate. Capping the dual free
cocycles and evaluating intersections gives
the free inverse coordinates.

\emph{Integral reduction for the simply connected members.}
Eliminate an order-one peripheral by
$c_i=\gamma^{n_i}$, changing $e$ to $e-n_i$.
For zero, one or two remaining places the
triviality tests are
\[
|e|=1,\qquad |m_1e-n_1|=1,\qquad
|m_2n_1+m_1n_2-e m_1m_2|=1 .
\]
With at least three, the sphere-orbifold quotient
in \ref{maps:src:D.2.23} is nontrivial.
Thus simple connectivity implies pairwise
coprime multiplicities and
$\gcd(m_i,n_i)=1$ at every site.
Also the complete abelian determinant is
$\bigl(\prod_i m_i\bigr)(\sum_i n_i/m_i-e)=\pm1$.

Put
\[
\begin{aligned}
{F_f}&=\prod_{\mathrm{finite}}m_i,\quad {Q_f}=\prod q_i,\quad
{P_f}={F_f}/{Q_f},\quad c=b/{Q_f},\\
p_i&=m_i/q_i,\quad
\alpha_i=bn_i/q_i,\quad N_f=\sum({F_f}/m_i)n_i.
\end{aligned}
\]
The $p_i$ are pairwise coprime. With ${D_f}=\sum_iD_i$
the finite relations are integrally equivalent to
${P_f}{D_f}+cN_fs=0$. An inverse is obtained by choosing
$E_i({P_f}/p_i)=1+k_{i,\rm crt}p_i$ and setting
\[
D_i=E_i({P_f}/p_i){D_f}+
\left(k_{i,\rm crt}\alpha_i+
 E_i\sum_{\ell\ne i}\frac{{P_f}}{p_ip_\ell}\alpha_\ell\right)s .
\]
For $p_i=1$ take $E_i=0,k_{i,\rm crt}=-1$.
Multiplying this expression by $p_i$ verifies
its relation; summing it differs from ${D_f}$
by a multiple of ${P_f}{D_f}+cN_fs$. These are the
two inverse substitutions.

For the regular sites let
\[
R=\prod m_j,\quad {v_{\rm reg}}=\sum_j(R/m_j)v_j=({X_0},{Y_0},{Z_0},{N_{\rm reg}}),\quad
{T_{\rm reg}}(u)=\sum_jt_j(u).
\]
Their relations are integrally equivalent to
\[
R{T_{\rm reg}}(u)=\omega(u,{v_{\rm reg}})s,\qquad {T_{\rm reg}}({v_{\rm reg}})=0.
\]
Choose $E_j(R/m_j)=1+k_{j,\rm crt}m_j$. The inverse is
\[
t_j(u)=E_j(R/m_j){T_{\rm reg}}(u)-
\left(k_{j,\rm crt}\omega(u,v_j)+
 E_j\sum_{\ell\ne j}\frac{R}{m_jm_\ell}\omega(u,v_\ell)\right)s .
\]
For the extra root rows put
\[
G_j=(R/m_j){T_{\rm reg}}(v_j)-
\sum_{\ell\ne j}\frac{R}{m_jm_\ell}\omega(v_j,v_\ell)s.
\]
The first relation gives $m_jG_j=0$, the inverse
formula gives $t_j(v_j)=E_jG_j$, and alternation
gives $\sum_jG_j={T_{\rm reg}}({v_{\rm reg}})$.
Chinese-remainder idempotents project this sum
onto each $G_j$, since the $m_j$ are pairwise
coprime. Hence all the root rows are preserved
in both directions.

Set $K={F_f}e-N_f$. The global and $\delta$ rows become
\[
{P_f}{T_{\rm reg}}(\delta)+cKs=0,\qquad R{T_{\rm reg}}(\delta)+c{Q_f}{N_{\rm reg}} s=0,\qquad
\Delta={F_f}{N_{\rm reg}}-RK=\pm1 .
\]
Choose $u_{\rm crt}R+v_{\rm crt}{P_f}=1$. The row matrix
$\left(\begin{smallmatrix}u_{\rm crt}&v_{\rm crt}\\{P_f}&-R\end{smallmatrix}\right)$
has determinant $-1$ and transforms these into
\[
{T_{\rm reg}}(\delta)+c(u_{\rm crt}{Q_f}{N_{\rm reg}}+v_{\rm crt}K)s=0,\qquad c\Delta s=0.
\]
Thus $cs=0$ and ${T_{\rm reg}}(\delta)=0$.
Also $\gcd({N_{\rm reg}},R)=1$. If $p_{\rm crt}{N_{\rm reg}}+q_{\rm crt}R=1$, eliminate
\[
{T_{\rm reg}}(\gamma)=-p_{\rm crt}\bigl({X_0}{T_{\rm reg}}(e_1)+{Y_0}{T_{\rm reg}}(e_2)\bigr).
\]
Its inverse satisfies both omitted rows because
$R({X_0}{T_{\rm reg}}(e_1)+{Y_0}{T_{\rm reg}}(e_2))=a{X_0}{Y_0}s-a{Y_0}{X_0}s=0$.
The resulting matrix is
\[
\begin{pmatrix}c&-a{Y_0}&a{X_0}\\0&R&0\\0&0&R\end{pmatrix}.
\]
Its determinantal divisors are $g,Rg,cR^2$,
where $g=\gcd(c,R,{X_0},{Y_0})$ and $\gcd(a,c)=1$.
The Smith factors are $g,R,cR/g$, proving
\eqref{maps:eq:D.HY4}. Empty finite or regular
sets are included by ${F_f}={Q_f}={P_f}=1,N_f={D_f}=0$ or
$R=1,{v_{\rm reg}}={T_{\rm reg}}=0$, respectively; their auxiliary
generators are eliminated by unit relations.

\emph{The complete list and its parameters.}
The order of \eqref{maps:eq:D.HY4} is $cR^2$,
so triviality is exactly $c=R=1$.
This proves \eqref{maps:eq:D.HY5} at the level
of homology and fundamental group.
Each finite pattern is realized by numerator
$n=\pm1$ at one finite place, or by B\'ezout
numerators at two coprime finite multiplicities.
With no finite place use order-one regular
sites with $|e-\sum n_j|=1$.
The affine quotient is free by the primitive
contracting coordinate.

Taking all component-compatible original and
primitive-$\gamma$ finite choices gives $164$
minimal patterns. The two whose finite orders
have an unavoidable common divisor are
43/II with $(a,b)=(1,2)$ and orders $(4,6)$,
and 49/II with $(a,b)=(1,3)$ and orders $(6,3)$.
There remain $162$ patterns admitting simply
connected parameters; exactly $59$ have $b=Q_f$.
The following table lists them. A plain OS number
denotes its default profile; the fibre positions
are those of \ref{maps:src:D.2.22}. Every zero is
$C_0$ of the indicated place. The sets $L_{\rm log},O_{\rm orig}$
give all additive logarithmic and original
positions, respectively.

\begingroup\small
\[
\begin{array}{c|rr|r|c|c}
\mathrm{profile}&a&b&\mathrm{zero\ place}&L_{\rm log}&O_{\rm orig}\\ \hline
43&1&2&2&\{1\}&\varnothing\\
43&1&2&3&\{1\}&\varnothing\\
43&1&2&4&\{1\}&\varnothing\\
43/II&1&2&3&\{1\}&\{2\}\\
45&8&1&2&\varnothing&\varnothing\\
45&8&1&3&\varnothing&\varnothing\\
45&8&1&4&\varnothing&\varnothing\\
45&8&1&5&\varnothing&\varnothing\\
45/II&8&1&3&\varnothing&\{2\}\\
45/II&8&1&3&\{2\}&\varnothing\\
45/II&8&1&4&\varnothing&\{2\}\\
45/II&8&1&4&\{2\}&\varnothing\\
47&7&1&2&\varnothing&\varnothing\\
47&14&1&3&\varnothing&\varnothing\\
47&14&1&4&\varnothing&\varnothing\\
47&14&1&5&\varnothing&\varnothing\\
47/II&7&1&2&\varnothing&\{4\}\\
47/II&7&1&2&\{4\}&\varnothing\\
47/II&14&1&3&\varnothing&\{4\}\\
47/II&14&1&3&\{4\}&\varnothing\\
\end{array}
\]
\[
\begin{array}{c|rr|r|c|c}
\mathrm{profile}&a&b&\mathrm{zero\ place}&L_{\rm log}&O_{\rm orig}\\ \hline
47/III&7&2&2&\{4\}&\varnothing\\
47/III&14&1&2&\varnothing&\{4\}\\
47/III&7&2&3&\{4\}&\varnothing\\
47/III&14&1&3&\varnothing&\{4\}\\
49&1&3&1&\{2\}&\varnothing\\
49&2&3&3&\{2\}&\varnothing\\
49&2&3&4&\{2\}&\varnothing\\
49/II&1&3&3&\{2\}&\{1\}\\
49/III&1&6&3&\{1,2\}&\varnothing\\
49/III&2&3&3&\{1\}&\{2\}\\
55&5&1&2&\varnothing&\varnothing\\
55&20&1&3&\varnothing&\varnothing\\
55&20&1&4&\varnothing&\varnothing\\
55&20&1&5&\varnothing&\varnothing\\
55/II&5&1&2&\varnothing&\{3\}\\
55/II&5&1&2&\{3\}&\varnothing\\
55/II&20&1&4&\varnothing&\{3\}\\
55/II&20&1&4&\{3\}&\varnothing\\
56&6&1&1&\varnothing&\varnothing\\
56&10&1&2&\varnothing&\varnothing\\
\end{array}
\]
\[
\begin{array}{c|rr|r|c|c}
\mathrm{profile}&a&b&\mathrm{zero\ place}&L_{\rm log}&O_{\rm orig}\\ \hline
56&15&1&3&\varnothing&\varnothing\\
56&30&1&4&\varnothing&\varnothing\\
56&30&1&5&\varnothing&\varnothing\\
56/II&15&1&1&\varnothing&\{2\}\\
56/II&15&1&1&\{2\}&\varnothing\\
56/II&6&1&3&\varnothing&\{2\}\\
56/II&6&1&3&\{2\}&\varnothing\\
56/II&10&1&4&\varnothing&\{2\}\\
56/II&10&1&4&\{2\}&\varnothing\\
56/III&6&1&1&\varnothing&\{3\}\\
56/III&5&2&2&\{3\}&\varnothing\\
56/III&10&1&2&\varnothing&\{3\}\\
56/III&15&2&4&\{3\}&\varnothing\\
56/III&30&1&4&\varnothing&\{3\}\\
56/IV&6&1&1&\varnothing&\{3\}\\
56/IV&5&3&2&\{3\}&\varnothing\\
56/IV&15&1&2&\varnothing&\{3\}\\
56/IV&10&3&4&\{3\}&\varnothing\\
56/IV&30&1&4&\varnothing&\{3\}\\
\end{array}
\]
\endgroup

There are $50$ distinct $(\mathrm{profile},a,b,\mathrm{zero})$
rows in this table. At each other multiplicative
place its component test has narrow $P$ and
transitive $Q$, permitting both fillings by the
period and Tate maps above.
Every additive profile has a one-finite pattern,
so arbitrary positive multiples of that finite
order, with $n=1$, realize infinitely many
sphere multiplicities. The unique two-finite
pattern is 49/III with $(a,b)=(1,6)$; its
multiplicities are coprime multiples of three
and four and satisfy the displayed determinant
condition. The $r_i$ disappeared through the
inverse finite-kernel substitution, so remain
unrestricted.

\emph{Transporting the cusp fan when $c=1$.}
The previous integral reduction gives $s=0$.
Every cusp two-class is a specialization of an
integral fibre class, so its global image is
zero. We compare the specified fan with the
fan having the physical $\delta$ as an edge.
In the short lattice $S=(u,\delta)$ its
specified basis is the matrix $\mathsf C_{\rm cusp}$ above.
For an oriented unimodular basis $E_{\rm fan}=(p_{\rm fan},q_{\rm fan})$,
the triangular fan has unoriented slopes
$p_{\rm fan},q_{\rm fan},q_{\rm fan}-p_{\rm fan}$. Each unit column addition and
the rotation $(p_{\rm fan},q_{\rm fan})\mapsto(q_{\rm fan},-p_{\rm fan})$ keeps two
slopes and replaces the third across their
primitive parallelograms.

Apply Euclidean division to the first row of
$\mathsf C_{\rm cusp}$: implement each quotient by unit column
additions and then rotate the columns.
The absolute remainder decreases until the
first row is $(\pm1,0)$. Simultaneous negation
leaves the same fan. A final integral lower
shear, expanded into unit additions, reaches
$I$. Reversing the word gives its inverse.
For each intermediate $E_{\rm fan}$ the period basis is
\[
\begin{aligned}
H_{E_{\rm fan}}={}&(SE_{\rm fan}e_1,\mathsf L_{\rm per}\mathsf C_{\rm cusp}^{-1}E_{\rm fan}e_1,\\
&\quad SE_{\rm fan}e_2,\mathsf L_{\rm per}\mathsf C_{\rm cusp}^{-1}E_{\rm fan}e_2).
\end{aligned}
\]
It is integral unimodular,
$H_{E_{\rm fan}}^{-1}\mathsf TH_{E_{\rm fan}}=I+E_{12}+E_{34}$, and its inverse
gives the period-coordinate map. The long
periods carry their holomorphic unit factors
through the same maps.

Every triangle contains one edge of a given
flipped slope. Distinct period translates of
that edge curve are disjoint, including their
endpoints; a common endpoint would require
two such edges in one triangle. The lattice
stabilizes no finite edge and preserves the
two triangle types. Thus each quotient
replacement concerns one embedded $\mathbb P^1$.
Its charts have inverse transition
\[
(x',y',w')=(wx,wy,w^{-1}),\qquad
(x,y,w)=(w'x',w'y',(w')^{-1}),
\]
so its normal bundle is $\mathcal O(-1)^2$.
The contraction is the rank-one matrix
\[
{\mathsf N}=\begin{pmatrix}x&xw\\y&yw\end{pmatrix},\qquad
q=xyw={\mathsf N}_{11}{\mathsf N}_{22}.
\]
On $x\ne0$ the other resolution has
$(u,p,q')=(y/x,x,xw)$, with inverse
$(x,y,w)=(p,up,q'/p)$. The other charts use
the other nonzero entries of ${\mathsf N}$.
Their comparison is identity on the
punctured cone, preserves the base parameter
and commutes with the torus characters and units.

For an exceptional curve $C_0$ in a cap $U$,
the oriented normal Thom isomorphism gives
$H_k(U,U\setminus C_0)=H_{k-4}(C_0)$.
The relative groups in degrees two and three
vanish. Therefore the punctured-cap inclusion
is an integral $H_2$ isomorphism on both
resolutions. Identity on the punctured cone
preserves the nearby four-torus and transports
its surjective specialization to the next cap.
The same global complement argument preserves
$s$; general position preserves simple
connectivity. By induction every flip curve
is nullhomologous and, by Hurewicz,
nullhomotopic.

We specify the relative smooth comparison.
The real normal clutch of $\mathcal O(-1)^2$
contracts in the unit quaternions through
$\cos(\pi t/2)e^{-i\theta}+\sin(\pi t/2)j$,
followed by a path from $j$ to one.
It gives $S^2\times D^4$ as a real normal
neighbourhood. The parametrized embedding
isotopy theorem
\cite[Theorem B(b), pp.~367--368]{manBarden65},
with $12>9$, places it in a six-ball.
Isotopy extension transports the normal maps
and is inverted by the backward flow.

The complement of that normal neighbourhood
inside the ball is simply connected and has
$H_3=H_5=\mathbb Z$, with other positive
homology zero. Mayer--Vietoris identifies
these groups with the degree-three and
degree-five inner-boundary groups, and
$H_2(S^2\times S^3)\to H_2(S^2\times D^4)$
is an isomorphism. On the matrix-norm link,
write ${\mathsf N}=\varepsilon {\xi}{\eta}^t$ modulo
$({\xi},{\eta})\sim(t{\xi},t^{-1}{\eta})$. The resolution
comparison exchanges the labelled base and
fibre factors; its inverse exchanges them
back. The matrix ${\mathsf N}$ is fixed. The exchange
reverses the product orientation of the two
three-sphere factors and the circle gauge,
so preserves their quotient orientation.

Gluing the opposite normal neighbourhood by
this boundary map gives an acyclic simply
connected ball replacement $B'$. Cap its
unchanged outer parameterization by a
standard six-disc. The result is a homotopy
six-sphere, hence an oriented standard
sphere by
\cite[pp.~504--505 and 512]{manKM63}
and the smooth h-cobordism theorem
\cite[Theorem 9.1, p.~107]{mapsMilnor65}.
Its diffeomorphism can be made equal to the
given parameterization on the attached disc:
shrink a parametrized disc embedding by
$x\mapsto e(rx)$ with positive $r$, interpolate
in a chart to its invertible derivative,
move its centre and positive frame, and
expand the standard disc. A cutoff extension
of the velocity field gives the ambient
isotopy and its inverse
\cite[pp.~62--64]{mapsMilnor65}.
Restriction and collar straightening give
an oriented $B'\to D^6$ with the prescribed
outer map. Glue this map to identity
outside that chosen ball to obtain an
oriented diffeomorphism of the closed
manifolds; the reversed construction is
its inverse. The ball may meet other
caps, since the sphere's nullhomotopy
is in the closed manifold. At each step
the filling construction changes only
the cusp charts and uses the same other
caps and boundary maps. Composing the
finite sequence gives an oriented global
diffeomorphism to the manifold constructed
with the physical-edge fan and those caps.

\emph{The circle quotient and oriented reconstruction.}
Use that physical-edge representative.
Multiplication of the line coordinate
commutes with the line transitions,
contraction and logarithmic actions.
It is free on an original cap because a
nonzero contraction iterate changes norm.
At every other multiplicative place in
the table, narrowness of $P$ makes the
vanishing vector's elliptic projection
primitive. Together with $\delta$ it
spans a primitive plane, so $\delta$
acts freely on the ruled base and its
identified sections.
At a logarithmic central fibre, an equality
with the $j$th deck action implies
$m\mid jn$ by its contracting coordinate.
Simple connectivity gives $\gcd(m,n)=1$,
so that equality gives only the identity.

At the physical cusp the two charts
\[
(x,y,w)=(q/(X_1X_2),X_2,X_1),\qquad
(x',y',w')=(q/X_2,X_1X_2,1/X_1)
\]
give one fixed sphere and all other
stabilizers trivial. The normal weights
are $-1,+1$, with bundle
$\mathcal O(-1)_{-1}\oplus\mathcal O(-1)_{+1}$.
Conjugate the negative-weight coordinate
and remove its $SU(2)$ equatorial clutch
by the disc matrix of \ref{maps:src:D.2.14}.
Its conjugate transpose is the inverse.
Reversing the core parameterization
compensates the normal orientation change.
The invariant normal exponential then
gives the scalar circle action on
$C_0\times D^4$ with this orientation.

The quotient charts are the Hopf maps
\[
h_{\rm Hopf}(z_1,z_2)=(2\Re(z_1\bar z_2),2\Im(z_1\bar z_2),
                  |z_1|^2-|z_2|^2),\qquad
\mathcal H(z)=h_{\rm Hopf}(z)/|z|,\quad\mathcal H(0)=0.
\]
The northern and southern Hopf sections
give the inverse local orbit coordinates.
Let ${\mathcal Q}$ be the closed oriented quotient,
${\mathcal Q}_0={\mathcal Q}\setminus(C_{\mathcal Q}\times\mathring D^3)$,
and $X_{\rm free}=X\setminus(C_0\times\mathring D^4)$.
Orient its meridian $\sigma$ so the Euler
evaluation is one. The fibration and
Hurewicz sequences give
\[
\pi_1({\mathcal Q}_0)=0,\qquad
0\to H_2(X)\to H_2({\mathcal Q}_0)\xrightarrow{\langle e_\delta,-\rangle}
\mathbb Z\to0,\qquad k\mapsto k\sigma
\]
as a splitting. Excision fills precisely
this meridian, so
\[
\pi_1({\mathcal Q})=0,\qquad H_2({\mathcal Q})\cong H_2(X)=(\mathbb Z/R)^2.
\]
To invert the latter map, lift a class to
${\mathcal Q}_0$, subtract its Euler evaluation
times $\sigma$, lift the resulting
zero-Euler sphere to $X_{\rm free}$ and include
it in $X$. A different first lift
changes by $\sigma$ and cancels; sections
of the lifted bundle differ by maps
$S^2\to S^1$ and give the same class.

Integral universal coefficients give
$H^2(X;\mathbb Z)=0$, hence $c_1(X)=0$.
Its reduction modulo two is $w_2(X)$, so $X$ is spin.
Let $p$ be the free circle-quotient projection.
On $X_{\rm free}$, $TX_{\rm free}=p^*T{\mathcal Q}_0\oplus\mathbb R$.
The mod-two Gysin kernel is generated
by $e_\delta\bmod2$, which evaluates
to one on $\sigma$.
The class $w_2({\mathcal Q}_0)$ evaluates to zero there
because $T{\mathcal Q}_0|_\sigma=TS^2\oplus\mathbb R^3$.
Thus ${\mathcal Q}$ is spin, including when $R$ is even.
The residue sequence sends one to
$\operatorname{PD}[C_{\mathcal Q}]$; the residue
of $e_\delta$ is one, so $[C_{\mathcal Q}]=0$.
Hurewicz and the parametrized embedding
isotopy theorem place $C_{\mathcal Q}$ in a five-disc.

The punctured quotient has integral
$H^2=0$, so its circle bundle is trivial.
\v{C}ech logarithms and reciprocal gauges
give the bundle map and inverse.
Inside the disc the residue-one
bundle is the Hopf model above.
An oriented normal rotation lifts
to $SU(2)$ on the simply connected
sphere; $\pi_2(SO(3))=0$ fixes its
framing homotopy class. A difference
of gauges on the common $S^4$ collar
has a real logarithm. A cutoff
exponential makes the maps agree,
and the reciprocal makes their
inverses agree. The resulting model is
\[
M_-({\mathcal Q})=(({\mathcal Q}\setminus\mathring D^5)\times S^1)
\cup_{g_{\rm glue}}(S^4\times D^2),\qquad
g_{\rm glue}(x,z)=(\alpha_zx,z),
\]
where $g_{\rm glue}$ goes from complement to
attached boundary and $\alpha_z$
rotates the first coordinate plane
of $\mathbb R^5$ through $\arg z$.
The inverse uses $\alpha_z^{-1}$.
The cap action is
$t(x,w)=(\alpha_tx,tw)$, whose boundary
orbit coordinate is $\alpha_z^{-1}x$.
Thus the quotient restores the original
disc parameterization.

The loop $\alpha$ lifts in
$\operatorname{Spin}(5)$ to
$\cos(\pi t)+\sin(\pi t)e_1e_2$,
with endpoints $1,-1$.
For the gluing direction in
\eqref{maps:eq:D.S9} use
$\mathsf A_{\rm rot}=\operatorname{diag}(1,-1,-1,1,1)$,
so $\mathsf A_{\rm rot}\alpha_z\mathsf A_{\rm rot}^{-1}=\alpha_z^{-1}$.
A collar rotation through $\pi$
in the $2,3$ plane, equal to identity
away from the collar, and $\mathsf A_{\rm rot}$ on
the cap give inverse oriented maps
from $M_-({\mathcal Q})$ to $\widetilde\Sigma_0{\mathcal Q}$.
They agree on the join by
$\mathsf A_{\rm rot}\alpha_z^{-1}=\alpha_z\mathsf A_{\rm rot}$.

The spin linking form of ${\mathcal Q}$ is
alternating
\cite[Lemma F, p.~372]{manBarden65}.
Choose $x$ of order $R$ and $y$
with $\lambda(x,y)=1/R$.
The inverse coordinates
\[
(r,s)\mapsto rx+sy,\qquad
z\mapsto(\operatorname{num}_R\lambda(z,y),
              -\operatorname{num}_R\lambda(z,x))
\]
identify its full linking form
with the standard form on
$(\mathbb Z/R)^2$.
The oriented classification
\cite[Theorems 2.2--2.3, pp.~376--377]{manBarden65}
gives ${\mathcal Q}\cong N_R$ inducing these maps.
Adjust its disc parameterization
by the disc isotopy above, take
its product with identity on the
punctured pieces and identity
on the common cap. The inverse
uses the inverse base map.
Composing with the cusp comparison
proves $X\cong\widetilde\Sigma_0N_R$.

For $R=1$, $N_1=S^5$.
The map $(x,z)\mapsto(\alpha_zx,z)$
on $D^5\times S^1$, with identity
on $S^4\times D^2$, gives
$\partial(D^5\times D^2)$.
Its inverse uses the opposite
rotation. Corner rounding and
radial projection give the
standard oriented $S^6$.

The associated spin data can also
be read from these maps.
On $X_{\rm free}$, $p_1$ pulls back from
$p_1({\mathcal Q}_0)$, which is zero by
$H^4({\mathcal Q};\mathbb Z)=H_1({\mathcal Q};\mathbb Z)=0$.
The kernel of restriction from
$H^4(X)$ is generated by
$\operatorname{PD}[C_0]=0$,
so $p_1(X)=0$, including torsion.
All positive-degree integral
products vanish: the only
possible product into degree
six has torsion source and
free target.
Replacing the framed
$S^2\times D^4$ by
$D^3\times S^3$ with the Hopf
action gives
$Y\cong X\#(S^3\times S^3)$
by the product-boundary maps
in $\partial(D^3\times D^4)$.
Its quotient and Euler class are
$(N_R\#(S^2\times S^3),x_{\rm Eul})$,
where $x_{\rm Eul}$ evaluates to one
on the Hopf sphere and zero
on the punctured $N_R$.
Pulling back the spin frame
bundle and adjoining the
oriented vertical line gives
the projectable spin lift.
This supplies the surgery
description with the same
boundary framing.
\end{proof}


\subsection{The \texorpdfstring{$I_1^*$}{I1*} integral obstruction}
\begin{theorem}[Integral obstructions for the three $I_1^*$ profiles]\hypertarget{tag-02L2}{}\label{maps:src:D.2.27}
{\normalfont\scriptsize tag 02L2}\par
None of the natural section-line inputs on
\[
 (I_1^*,I_3,I_1,I_1),\qquad (I_1^*,II,I_3),\qquad
 (I_1^*,IV,I_1)
\]
produces a homotopy six-sphere, for any permitted original,
logarithmic or Mumford filling, multiplicity or affine twist in
the sector of \ref{maps:src:D.2.22}.

The fundamental-group condition and the rational homology obstruction
leave the following two star caps. Their return numerator is odd,
their boundary maps in degree two are onto, and their full integral
normal-kernel projection indices are
\begin{equation}
\begin{array}{c|c|c|c}
 m_*&\text{component action}&H_2(U_*;\mathbb Z)
       &\kappa_*\\\hline
 2&\text{one component}&\mathbb Z^2&1\\
 4&\text{two components interchanged}&
                   \mathbb Z^2\oplus\mathbb Z/2&4.
\end{array}                                                   \hypertarget{tag-024V}{}\label{maps:eq:D.ST1}
\end{equation}
For a simply connected complete threefold whose other caps have minimal
rational ranks, $|H_2|$ has exactly one factor of two in the first
row and exactly two factors of two in the second row.
\end{theorem}


\begin{proof}
\emph{The marked inputs and the primitive star periods.}
Write
\[
 \Lambda=\mathbb Z\langle e_1,e_2,\delta,\gamma\rangle,
 \qquad \chi=\gamma^*,\qquad
 \omega=a e_1^*\wedge e_2^*-b\delta^*\wedge\gamma^*.
\]
We use the positive section representatives $P=aG,Q=bG$ of
\ref{maps:src:D.2.22}. Its map $p\mapsto-p$, with itself as inverse,
pulls back $\mathcal O(P-O)$ to $\mathcal O(-P-O)$, conjugates the
forward translations, and transports their multipliers and reciprocal
maps. Its period map is $\operatorname{diag}(-I_2,1,1)$; the charts,
fans and collars are transported by the associated integral character
maps. This covers the simultaneous negative representatives.

Here $\langle G,G\rangle=1/12$. At the selected $I_\nu$ the natural
component equation gives $ab=12/\nu$. Its residues satisfy
$a\kappa=b\kappa=\pm1\pmod\nu$ for $\nu>1$, and one of $aG,bG$
is narrow at every other reducible fibre. Enumerating the positive
divisors of $12/\nu$ gives
\[
\begin{array}{c|c|l}
 \text{profile}&\text{selected position}&(a,b)\\\hline
 50& I_3\ (2)&(1,4),(4,1)\\
 50& I_1\ (3\text{ or }4)&(1,12),(3,4),(4,3),(12,1)\\
 50/II&I_3\ (2)&(1,4),(4,1)\\
 50/IV&I_1\ (3)&(1,12),(3,4),(4,3),(12,1).
\end{array}
\]
At an unselected $I_\nu$ the test is
$\nu\mid a\kappa$ or $\nu\mid b\kappa$; at an additive fibre it is
the integrality of $a(A-I)^{-1}{g_{\rm ell}}$ or $b(A-I)^{-1}{g_{\rm ell}}$. These are the
component equations of \ref{maps:src:D.2.22} in the marked matrices
of \ref{maps:src:D.2.23}. They give all sixteen inputs. Each pair is
coprime and exactly one entry is divisible by four.

For the default and $IV$ profiles, the star data are
\[
 A=\begin{pmatrix}-1&-1\\0&-1\end{pmatrix},\quad
 {g_{\rm ell}}=\binom{-1}{1},\quad (A-I)^{-1}{g_{\rm ell}}=\binom{3/4}{-1/2},
\]
and the four-period matrix is
\[
 {T_{\rm st}}=\begin{pmatrix}
 -1&-1&0&-b\\0&-1&0&b\\-a&-2a&1&ab/4\\0&0&0&1
 \end{pmatrix}.
\]
For $50/II$, its last fibre has
$A'=\left(\begin{smallmatrix}-2&-1\\1&0\end{smallmatrix}\right)$,
${g_{\rm ell}}'=(0,1)^t$. Put
\[
 {\mathsf C_{\rm st}}=\begin{pmatrix}1&1\\-1&0\end{pmatrix},\qquad
 {\mathsf C_{\rm st}}^{-1}=\begin{pmatrix}0&-1\\1&1\end{pmatrix}.
\]
Then ${\mathsf C_{\rm st}}^{-1}A'{\mathsf C_{\rm st}}=A$ and ${\mathsf C_{\rm st}}^{-1}{g_{\rm ell}}'={g_{\rm ell}}$. New coordinates are
$\operatorname{diag}({\mathsf C_{\rm st}}^{-1},1,1)$ times old coordinates; the inverse
is $\operatorname{diag}({\mathsf C_{\rm st}},1,1)$. These maps fix $\delta,\gamma$,
preserve $\omega$, and carry the primitive angular characters, period
units and integer rank-one rays to the displayed marking.

First take $a$ odd and write $b=4B_{\rm per}$, where $B_{\rm per}$ is odd. The possible
pairs are
\begin{equation}
 (a,b)=(1,4),\ (1,12),\ (3,4).                           \hypertarget{tag-0250}{}\label{maps:eq:D.ST2}
\end{equation}
Set ${s_{\rm per}}=(a+1)/2$, ${k_{\rm per}}=(3a-1)/2$, ${h_{\rm per}}=3b/4$ and
\[
 \begin{gathered}
 \alpha=2e_1+a\delta,\qquad \lambda=e_1+{s_{\rm per}}\delta,\\
 \beta=e_2+{k_{\rm per}}\lambda,\qquad \ell=\gamma+(b/2)e_2-{h_{\rm per}}\lambda.
 \end{gathered}
\]
The matrix $H=(\alpha,\beta,\lambda,\ell)$ is integral with
determinant $2{s_{\rm per}}-a=1$, and $H^{-1}=\operatorname{adj}H$.
Multiplication gives
\begin{equation}
\begin{gathered}
 {T_{\rm st}}\alpha=-\alpha,\quad {T_{\rm st}}\beta=-\beta-\lambda,\\
 {T_{\rm st}}\lambda=\lambda-\alpha,\quad {T_{\rm st}}\ell=\ell,\\
 \delta=-\alpha+2\lambda,\qquad
 H^{-1}{T_{\rm st}}^2H=I+E_{12},\qquad \chi(\ell)=1.
\end{gathered}                                               \hypertarget{tag-0251}{}\label{maps:eq:D.ST3}
\end{equation}
In particular $\operatorname{Fix}{T_{\rm st}}=\mathbb Z\delta\oplus\mathbb Z\ell$.
The alternate fixed period
$\ell_0=\gamma+(b/2)e_2-(3b/4)e_1$ satisfies
$\ell_0=\ell+{h_{\rm per}}{s_{\rm per}}\delta$. Adding this multiple of $\delta$, and
subtracting it for the inverse, is an integral shear of the fixed
periods. Since $[\beta\wedge\delta]=0$ and
$\omega(\beta,\delta)=0$, it gives the same exterior generators below.

The primitive exterior-square relation columns in $H$ coordinates are
$e_{13},e_{14},e_{12}+2e_{23},2e_{24}+e_{34}$.
The quotient map and a section are
\[
 (x_{12},x_{13},x_{14},x_{23},x_{24},x_{34})
 \longmapsto(x_{23}-2x_{12},x_{24}-2x_{34}),
 \qquad (x,y)\longmapsto x e_{23}+y e_{24}.
\]
Thus the vertical group is freely generated by
$x=[\beta\wedge\lambda]$, $y=[\beta\wedge\ell]$, with
\begin{equation}
 [\alpha\wedge\beta]=-2x,\quad
 [\lambda\wedge\ell]=-2y,\quad
 [\alpha\wedge\lambda]=[\alpha\wedge\ell]=0.              \hypertarget{tag-0252}{}\label{maps:eq:D.ST4}
\end{equation}
The boundary Wang sequence therefore has free middle group of rank
four. Evaluation gives
\begin{equation}
 \omega(x)=-a,\qquad \omega(y)=b/4.                        \hypertarget{tag-0253}{}\label{maps:eq:D.ST5}
\end{equation}
For example, in this basis
$\omega(\alpha,\beta)=2a$ and
$\omega(\lambda,\ell)=-b/2$; these also verify the displayed relations.

For the other half, $4\mid a$ and $b$ is odd. Use the integral basis
\[
 \begin{gathered}
 {w_{\rm st}}=e_1+(a/2)\delta,\quad {l_{\rm st}}=((b+1)/2)e_2+\gamma,\\
 {z_{\rm st}}=\delta,\quad {t_{\rm st}}=b e_2+2\gamma.
 \end{gathered}
\]
Its determinant is $b+1-b=1$. After collapsing the primitive ${w_{\rm st}}$,
the induced action is
\[
 {l_{\rm st}}\longmapsto-{l_{\rm st}}+{t_{\rm st}}-(3a/4){z_{\rm st}},\qquad {z_{\rm st}}\longmapsto {z_{\rm st}},\qquad
 {t_{\rm st}}\longmapsto {t_{\rm st}}.
\]
Its fixed lattice is $\mathbb Z{z_{\rm st}}\oplus\mathbb Z{t_{\rm st}}$, on which
$\chi$ takes even values. The power of a root generator commutes
with that generator, so its image in the collapsed cover group is
fixed and has even return. The root order is even. The Seifert
presentation of \ref{maps:src:D.2.23} consequently cannot be trivial.
This excludes every logarithmic star cap in this half.

\emph{Original star caps and the rational rank obstruction.}
The original $I_1^*$ has affine $D_5$ component tree, with edges
\[
 (0,2),(1,2),(2,3),(3,4),(3,5)
\]
and multiplicity vector $\mu=(1,1,2,2,1,1)$. Its neighbourhood
retracts onto this tree of spheres. The punctured-line map of
\ref{maps:src:D.2.16}, including its frame and inverse frame, identifies
the cap with the contracting mapping torus of the circle bundle of
$\mathcal O(P-O)$ over this neighbourhood.

For $a$ odd, the line Euler row is the primitive row $e_{p_{\rm comp}}-e_0$,
where ${p_{\rm comp}}=4$ or $5$. The circle bundle has $H_1=0$ and
$H_2=\ker(\mathbb Z^6\to\mathbb Z)=\mathbb Z^5$.
Projection omitting coordinate ${p_{\rm comp}}$ and reconstruction of that
coordinate from the Euler equation give inverse coordinates.
Translation by $bG$ fixes every component and its Euler kernel.
The mapping torus therefore has $b_2=5$, while its boundary has
$b_2=4$ by \eqref{maps:eq:D.ST3}--\eqref{maps:eq:D.ST4}. Its local
degree-two cokernel has positive rational rank.

For $4\mid a$, the line has a nonzero frame, with division by that
frame as inverse. Translation by the odd generator permutes the
components as $(0\ 4\ 1\ 5)(2\ 3)$ or its inverse. The intersection
matrix ${I_{\rm comp}}$, with diagonal $-2$ and edge entries one, satisfies
\[
 -u^t{I_{\rm comp}}u=\sum_{\{i,j\}\text{ edge}}\mu_i\mu_j
       \left(\frac{u_i}{\mu_i}-\frac{u_j}{\mu_j}\right)^2.
\]
Its kernel is $\mathbb Q\mu$. On the two orbit sums it is
$\left(\begin{smallmatrix}-8&4\\4&-2\end{smallmatrix}\right)$,
of rank one. Relative intersection with the component four-cycles
therefore gives a nonzero local degree-two quotient; boundary classes
have zero relative intersection. Transfer through the component
permutation gives the same quotient on the contracting mapping torus.
Original star caps are excluded in both halves.

For a logarithmic star cap, \eqref{maps:eq:D.ST3} gives
${T_{\rm st}}^{2k}=I+kE_{12}$ in the primitive basis. On the root disc the
adjacent charts are
\[
 X_j=Z/s^j,\qquad Y_j=s^{j+1}/Z,\qquad X_jY_j=s;
 \qquad Z=X_js^j,\quad s=X_jY_j.
\]
The long period identifies $j$ with $j+k$, giving $k$ ruled
components. A flat multiplier is removed smoothly by
$\exp(-u\log\mu_1-v\log\mu_2)$ in the transverse universal cover;
the reciprocal exponential is the inverse. If the final pole
identification translates the transverse torus by $\tau$, subtract
$f_{\rm pole}(Z)\tau$, with $f_{\rm pole}(0)=0$, $f_{\rm pole}(\infty)=1$; addition is its inverse.
These are maps of the period charts and their collars.

The root generator reverses the primitive short period. Its toric
form is $Z\mapsto u_{\rm node}(s,y)s^{c_{\rm st}}/Z$, with a nonzero unit; the inverse
uses its reciprocal at the inverse-transformed point. Its action
on the components is exactly $j\mapsto c_{\rm st}-j\pmod k$. Affine changes
alter $c_{\rm st}$ and the units and preserve this slope and cycle length.
The cover has component classes $C_j$, the two cycle-cross-torus
classes, and the transverse class $E$. Specialization sends
$\alpha\wedge\beta$ to $\sum C_j$ and
$\lambda\wedge\ell$ to $E$; the three fixed normal sweeps bound
in the cover. Its local quotient is consequently
$\mathbb Z^k/\mathbb Z(1,\ldots,1)$.

Rational transfer and division by the covering degree identify the
quotient with the invariant part. If $\rho$ is the number of
component orbits, the local cokernel has rank $\rho-1$.
The two mixed classes have eigenvalue $-1$ modulo $E$, while $E$
is fixed, so $b_2(U_*)=\rho+1$ and the vertical coinvariant rank
is two. A positive local cokernel survives globally: projecting the
Mayer--Vietoris cokernel to that local quotient is onto, and this
cokernel injects into $H_2(X;\mathbb Q)$.
Thus $\rho=1$ is necessary. A reflection has this property exactly
for $k=1$, or for $k=2$ with the components interchanged. These give
$m_*=2$ and $m_*=4$. Simple connectivity further requires an odd
return $n_*$ by \ref{maps:src:D.2.23}.

\emph{The normalization and the conductor lattices.}
Write the collapsed cover group multiplicatively as
$\mathbb Z\langle\beta,\lambda,\ell\rangle$, and let $g$ be the
positive root path. Its relations are
\[
 g\beta g^{-1}=\beta^{-1}\lambda^{-1},\qquad
 g\lambda g^{-1}=\lambda,\qquad g\ell g^{-1}=\ell,\qquad
 g^m=\lambda^R\ell^n.
\]
The last relation has no $\beta$ coefficient because its left side
commutes with $g$. Here $n$ is odd and $R$ records the affine choices.
For $d=2$ or $4$ put
\[
 \mathcal L_{\rm lat}(d,R,n)=\mathbb Z^2+\mathbb Z(R/d,n/d).
\]
The identity of $\mathbb R^2$ induces a degree-$d$ map from
${\mathcal E_{\rm cond}}=\mathbb R^2/\mathbb Z^2$ to $\mathbb R^2/\mathcal L_{\rm lat}(d,R,n)$.
For an integral lattice basis matrix $K$, its $H_1$ matrix is
$K^{-1}$ and the inverse on the image is $K$.

Write $W_m$ for the normalization at root order $m$.
At $m=2$, choose the lift fixing a component of the infinite chain.
The normalization ${W_m}$ has base lattice $\mathcal L_{\rm lat}(2,R,n)$. The node
stabilizer is ${h_{\rm stab}}=\beta g$, and ${h_{\rm stab}}^2=\lambda^{R-1}\ell^n$.
The double curve $D$ therefore has lattice $\mathcal L_{\rm lat}(2,R-1,n)$.
The conductor is ${\mathcal E_{\rm cond}}$, a connected bisection of ${W_m}$, with its
second degree-two map to $D$. These two distinct lattices specify
the integral conductor map.

At $m=4$, choose the reflection $j\mapsto1-j$. The normalization
stabilizer is $g^2$, and its base is $\mathcal L_{\rm lat}(2,R,n)$. There are two
double curves with lattices $\mathcal L_{\rm lat}(4,R,n)$ and $\mathcal L_{\rm lat}(4,R-2,n)$, since
$(\beta g)^2=\lambda^{-1}g^2$. The two conductor components are
the pole sections of ${W_m}$. Their $H_1$ matrices are
$K_{D_j}^{-1}K_{W_m}$, with inverse on the image $K_{W_m}^{-1}K_{D_j}$;
both have determinant two. The adjacent-node charts show these
are the gluing maps of the central surface.

All odd $n$ are covered by explicit changes of generators.
Choose ${u_{\rm st}},{v_{\rm st}}$ with ${u_{\rm st}}n+{v_{\rm st}}m=1$, and set $g'=g^{u_{\rm st}}\ell^{v_{\rm st}}$.
Then $(g')^m=\lambda^{{u_{\rm st}}R}\ell$ and ${u_{\rm st}}$ is odd. If
$\alpha_0{u_{\rm st}}+\beta_0m=1$, the inverse substitution is
\[
 g=(g')^{\alpha_0}\lambda^{\beta_0R}
                  \ell^{\beta_0n-\alpha_0{v_{\rm st}}}.
\]
Multiplication of $g'$ by $\lambda$ changes $R$ by a multiple of
$m$. These maps retain the original meridian for the subsequent
normal calculation. Let $q_{\rm col}$ be the quotient killing
$\alpha$. Use the original marked lifts
$g_{\rm lift},\lambda_{\rm lift},\ell_{\rm lift}$, with respective
images $g,\lambda,\ell$. The specified lift
$g^{\prime}_{\rm lift}=g_{\rm lift}^{u_{\rm st}}\ell_{\rm lift}^{v_{\rm st}}$
has image $g'$ under $q_{\rm col}$.

For $m=2$, cut the base in the $\lambda,g'$ directions and let
$\theta$ be its first coordinate and $\phi$ the short-circle angle.
The branch exchange has angle
$-\phi-\theta+u_0(\theta)$, where $u_0$ is periodic. The coefficient
of $-\theta$ is the commutator
\[
 g^{\prime}_{\rm lift}\lambda_{\rm lift}(g^{\prime}_{\rm lift})^{-1}\lambda_{\rm lift}^{-1}=\alpha^{-1}
\]
from
\eqref{maps:eq:D.ST3}; every odd power has this same action on
$\lambda_{\rm lift}$. A periodic logarithm of the degree-zero unit, extended
with a collar cutoff, removes $u_0$, with reciprocal scaling as
inverse. A unitary lift of the projective clutch is then
\[
 \begin{gathered}
 {\mathsf J_{\rm cl}}(\theta)=\begin{pmatrix}0&-e^{-\pi i\theta}\\
                       e^{\pi i\theta}&0\end{pmatrix},\\
 {\mathsf J_{\rm cl}}^{-1}={\mathsf J_{\rm cl}}^*,\qquad {\mathsf J_{\rm cl}}(1)=-{\mathsf J_{\rm cl}}(0).
 \end{gathered}
\]
The lift obstruction is the nonzero class of $H^2(T^2;\mathbb Z/2)$.
The oriented $S^2$ bundle therefore has odd section square. A section
extends across the base cells because $S^2$ is connected and simply
connected. Replacing it over a disc by a degree-$r$ map, constant on
the boundary collar, adds $r$ times the ruling fibre ${f_{\rm two}}$ and changes
its square by $2r$; the reverse replacement has degree $-r$.
Choose a section $\sigma_{\rm two}$ with $\sigma_{\rm two}^2=1$.

The connected conductor bisection has flat normal line on its cover,
so its square is zero. Its class is $2\sigma_{\rm two}+k_{\rm cond}{f_{\rm two}}$; squaring gives
$4+4k_{\rm cond}=0$, hence $[C]=2\sigma_{\rm two}-{f_{\rm two}}$. The normalization pushout gives
\begin{equation}
 \begin{gathered}
 H_2(U_2)=\langle\sigma_{\rm two},{f_{\rm two}},D\mid2\sigma_{\rm two}-{f_{\rm two}}-2D=0\rangle\\
        =\mathbb Z\langle\sigma_{\rm two},D\rangle.
 \end{gathered}                 \hypertarget{tag-0254}{}\label{maps:eq:D.ST6}
\end{equation}
Its inverse coordinate substitution is ${f_{\rm two}}=2\sigma_{\rm two}-2D$.
The conductor embedding has a tubular neighbourhood, so its mapping
cylinder collapses to the gluing and the normalization sequence is
the integral homology sequence of this pushout. The conductor map
on $H_1$ is injective, since its map to $D$ is a covering. In the
$\lambda,g'$ and $\lambda,{h_{\rm stab}}'$ bases the common-cover relation
eliminates $\lambda=2(g'-{h_{\rm stab}}')$, and its other generator is recovered
from the power relation. Thus $H_1(U_2)=\mathbb Z^2$.
The conductor $H_2$ map is injective, giving
$H_3(U_2)=\mathbb Z^2$ and $H_4(U_2)=\mathbb Z$.

At $m=4$ the stabilizer lift $g_{\rm lift}^2$ fixes both branches and
$\alpha,\lambda_{\rm lift},\ell_{\rm lift}$. Its open ruling has abelian
period lattice with relation
$2g_{\rm lift}^2=t_0\alpha+R\lambda_{\rm lift}+n\ell_{\rm lift}$.
Since $n$ is odd, $\alpha$ is primitive in that lattice. The
quotient by $\alpha$ is free of rank two, and a B\'ezout change using
$\gcd(2,n)=1$ gives an integral splitting. The line has zero Chern
class. Its flat multipliers are removed by their logarithms and
reciprocal gauges, preserving the two sections. Both sections have
the same class $\sigma_{\rm four}$. Their degree-two maps give
\begin{equation}
\begin{split}
 H_2(U_4)&=\langle {f_{\rm four}},\sigma_{\rm four},D_0,D_1\mid
                  \sigma_{\rm four}-2D_0=\sigma_{\rm four}-2D_1=0\rangle\\
 &=\mathbb Z {f_{\rm four}}\oplus\mathbb ZD_0
                    \oplus\mathbb Z/2\langle D_1-D_0\rangle.
\end{split}                                                 \hypertarget{tag-0255}{}\label{maps:eq:D.ST7}
\end{equation}
The inverse inserts $\sigma_{\rm four}=2D_0$ and
$D_1=D_0+(D_1-D_0)$. On $H_1$ the two conductor relations use
$g^2=2g$ and $g^2=2{h_{\rm stab}}+\lambda$, eliminating
$\lambda=2(g-{h_{\rm stab}})$ after the preceding base changes. They give
$H_1(U_4)=\mathbb Z^2$. The two conductor $H_2$ columns are
independent in their separate double-curve coordinates, so
$H_3(U_4)=\mathbb Z^2$ and $H_4(U_4)=\mathbb Z$.

These calculations compute the cap groups. To see the retraction
explicitly in a node chart, put $a_0=|X|^2$, $b_0=|Y|^2$ and
${d_{\rm rad}}=a_0-b_0$, and replace the squared radii by
\[
 a_t=\frac{\sqrt{{d_{\rm rad}}^2+4(1-t)^2a_0b_0}+{d_{\rm rad}}}{2},\qquad
 b_t=\frac{\sqrt{{d_{\rm rad}}^2+4(1-t)^2a_0b_0}-{d_{\rm rad}}}{2}.
\]
Scaling $X,Y$ by the nonnegative radius ratios gives
$XY\mapsto(1-t)XY$ and fixes the central axes. It is invariant
under unitary phases and branch exchange. The unitary frames and
an invariant partition of unity extend its normal contraction over
the component collars; the contractions agree on the central surface.
It descends under the free action and gives a homotopy inverse to
the inclusion of the normalization pushout.

\emph{The full normal indices.}
The pair sequence and Poincar\'e--Lefschetz duality give
$H_2(U_*,\partial U_*)=H^4(U_*)=\mathbb Z$, because the computed
$H_3$ is free and $H_4=\mathbb Z$. The map from $H_2(U_*)$ to this
relative group is zero: its rational image is the zero local defect,
and a subgroup of $\mathbb Z$ with zero rationalization is zero.
Equivalently, a positive multiple of the reduced central divisor is
the principal divisor of the base parameter, so all its relative
intersection numbers vanish. It follows that
\begin{equation}
 H_2(\partial U_*)\longrightarrow H_2(U_*)
       \quad\text{is onto integrally}.                    \hypertarget{tag-0256}{}\label{maps:eq:D.ST8}
\end{equation}

For $m=2$, the vertical $\alpha\wedge\beta$ torus maps to ${f_{\rm two}}$;
the transverse $\lambda\wedge\ell$ torus maps with degree two
to $D$. Equations \eqref{maps:eq:D.ST4} and
\eqref{maps:eq:D.ST6} give
\[
 x\longmapsto D-\sigma_{\rm two},\qquad y\longmapsto-D.
\]
The inverse sends $\sigma_{\rm two}$ to $-x-y$ and $D$ to $-y$.
For $m=4$, these same cover tori map to $2{f_{\rm four}}$ and $4D_0$.
On free quotients the map in bases $(x,y)$ and $({f_{\rm four}},D_0)$ is
$\operatorname{diag}(-1,-2)$. The images may have their additional
components in $\mathbb Z/2\langle D_1-D_0\rangle$.
The map is injective on the free quotient, so its image meets
cap torsion trivially. Therefore
\[
 0\longrightarrow\operatorname{Tor}H_2(U_4)
 \longrightarrow H_2(U_4)/\operatorname{im}\langle x,y\rangle
 \longrightarrow H_2(U_4)_{\rm free}/\operatorname{im}
                    \langle x,y\rangle_{\rm free}
 \longrightarrow0
\]
has end groups of orders two and two. Its middle group has order
four for every torsion component of the two images.

Let $V_*=H_2(T^4)_{T_{\rm st}}$ and
$K_*=\ker[H_2(\partial U_*)\to H_2(U_*)]$.
The Wang sequence has quotient $\operatorname{Fix}{T_{\rm st}}$.
Fixed-circle sweeps in the marked frame give its section; addition
of circles is represented by the period triangle and subtraction
gives the inverse. Since $V_*\to H_2(U_*)$ is injective, projection
identifies $K_*$ with a sublattice $N_*$ of $\operatorname{Fix}{T_{\rm st}}$.
The quotient map sends a sweep to its cap image modulo $V_*$;
its inverse lifts a cap class to the boundary and projects that lift
to the fixed lattice. Surjectivity in \eqref{maps:eq:D.ST8} proves
that these maps are inverse. Hence
\begin{equation}
 [\operatorname{Fix}{T_{\rm st}}:N_*]
 =|H_2(U_*)/\operatorname{im}V_*|
 =\begin{cases}1,&m_*=2,\\4,&m_*=4\text{ with swap}.
   \end{cases}                                             \hypertarget{tag-0257}{}\label{maps:eq:D.ST9}
\end{equation}

\emph{Additional multiplicative root quotients.}
We calculate their full integral maps as well. In a primitive
zero-free rank-one cap choose the integral period basis
$({w_{\rm rk}},{l_{\rm rk}},{z_{\rm rk}},{t_{\rm rk}})$ with ${T_{\rm rk}}{l_{\rm rk}}={l_{\rm rk}}+{w_{\rm rk}}$ and the other three periods fixed.
Such bases and their inverse period maps are given by
\ref{maps:src:D.2.26}. On the root disc $q=s^m$ use
\[
 X_j=Z/s^j,\qquad Y_j=s^{j+1}/Z,\qquad X_jY_j=s.
\]
The ${l_{\rm rk}}$ period identifies $j$ with $j+m$.
Write the power vector modulo ${w_{\rm rk}}$ as
${v_{\rm rk}}={c_{\rm rk}} {l_{\rm rk}}+{p_{\rm rk}}{z_{\rm rk}}+q_0{t_{\rm rk}}$. A logarithm of each nonzero period unit exists
on the small disc. Taking its $m$-th part, the short-coordinate
deck map is
\[
 \begin{gathered}
 (s,Z)\longmapsto
 (\zeta s,\eta s^{c_{\rm rk}} U_{\rm per}(q)^{{c_{\rm rk}}/m}V_{\rm per}(q)Z),\\
 \zeta=e^{2\pi i/m},\quad
 \eta^m\zeta^{{c_{\rm rk}} m(m-1)/2}=1.
 \end{gathered}
\]
Here $U_{\rm per}$ is the ${l_{\rm rk}}$-period unit and $V_{\rm per}$ consists of the fractional
${z_{\rm rk}},{t_{\rm rk}}$ multipliers; the transverse coordinate receives their same
fractional period. Its $m$-th power is exactly the period ${v_{\rm rk}}$.
An integral ${w_{\rm rk}}$ part is included by its constant angular multiplier.
In the charts $j$ and $j+{c_{\rm rk}}$, this map multiplies $X_j,Y_j$ by
\[
 \eta\zeta^{-j-{c_{\rm rk}}}U_{\rm per}^{{c_{\rm rk}}/m}V_{\rm per},\qquad
 \eta^{-1}\zeta^{j+{c_{\rm rk}}+1}U_{\rm per}^{-{c_{\rm rk}}/m}V_{\rm per}^{-1},
\]
respectively. They are nonzero units with product $\zeta$.
The inverse uses their reciprocals and shifts the chart by $-{c_{\rm rk}}$,
so it extends over the central nodes. The component action is
$j\mapsto j+{c_{\rm rk}}\pmod m$.

Minimal rational rank requires one component orbit, by the component
quotient computed above. Thus $\gcd({c_{\rm rk}},m)=1$. Every proper deck power
then moves each component and node; away from the central fibre it
moves the base root. The quotient is free. Restriction from a chosen
root fibre is the original fibre isomorphism, inverted by choosing
that representative. Continuation of the long period around $q$
changes it by ${w_{\rm rk}}$, giving the original boundary matrix ${T_{\rm rk}}$.
The positive normal coordinate is $q=s^m$, and the power vector has
return $\chi({v_{\rm rk}})$. These maps specify the permitted affine filling
and its marked boundary.

Its normalization is one ruled surface with trivial component
stabilizer. Its line has flat multipliers; multiplying by
$\exp(-x\log u_{z_{\rm rk}}-y\log u_{t_{\rm rk}})$ on the transverse universal cover,
and by the reciprocal for the inverse, trivializes it smoothly.
The two pole sections have the same class. Their conductor maps
have degree one and differ by a transverse translation $\tau$.
The map $(x,y)\mapsto(x,y-h_{\rm cut}(x)\tau)$, with $h_{\rm cut}$ zero and one at
the two poles, removes that translation; addition gives the inverse.
The central core is a nodal sphere times the transverse torus by
these maps. The node retraction used above gives
\[
 H_1(U_{\rm rk})=\mathbb Z^3,\quad H_2(U_{\rm rk})=\mathbb Z^4,\quad
 H_3(U_{\rm rk})=\mathbb Z^3,\quad H_4(U_{\rm rk})=\mathbb Z.
\]
In the integral normalization sequence, write $B_{\rm pole}$ for
the common pole-section class and $E_{\rm dbl}$ for the double-curve
class. The two degree-one conductor columns are $(B_{\rm pole},-E_{\rm dbl})$;
its $H_1$ kernel gives the two nodal-loop cross tori.

The cap group is
$\langle {l_{\rm rk}},{z_{\rm rk}},{t_{\rm rk}},{g_{\rm rk}}\mid\text{all commute},\,m{g_{\rm rk}}={c_{\rm rk}}{l_{\rm rk}}+{p_{\rm rk}}{z_{\rm rk}}+q_0{t_{\rm rk}}\rangle$.
Choose ${u_{\rm rk,Bez}},{r_{\rm rk,Bez}}$ with ${u_{\rm rk,Bez}}{c_{\rm rk}}+{r_{\rm rk,Bez}}m=1$ and put $\eta_0={u_{\rm rk,Bez}}{g_{\rm rk}}+{r_{\rm rk,Bez}}{l_{\rm rk}}$.
The inverse coordinates are
\[
 \begin{gathered}
 {l_{\rm rk}}=m\eta_0-{u_{\rm rk,Bez}}{p_{\rm rk}}{z_{\rm rk}}-{u_{\rm rk,Bez}}q_0{t_{\rm rk}},\\
 {g_{\rm rk}}={c_{\rm rk}}\eta_0+{r_{\rm rk,Bez}}{p_{\rm rk}}{z_{\rm rk}}+{r_{\rm rk,Bez}}q_0{t_{\rm rk}}.
 \end{gathered}
\]
Thus $(\eta_0,{z_{\rm rk}},{t_{\rm rk}})$ is a free basis. The collapsed fibre inclusion
has matrix
$\left(\begin{smallmatrix}m&0&0\\-{u_{\rm rk,Bez}}{p_{\rm rk}}&1&0\\-{u_{\rm rk,Bez}}q_0&0&1\end{smallmatrix}\right)$,
with its ordinary inverse on the image.

Use the boundary generators
\[
 \begin{gathered}
 {f_{\rm bd}}={w_{\rm rk}}\wedge {l_{\rm rk}},\quad A_0={l_{\rm rk}}\wedge {z_{\rm rk}},\quad B_0={l_{\rm rk}}\wedge {t_{\rm rk}},\\
 E_0={z_{\rm rk}}\wedge {t_{\rm rk}},\quad {W_{\rm bd}}=\operatorname{sw}({w_{\rm rk}}),\\
 Z_0=\operatorname{sw}({z_{\rm rk}}),\quad T_0=\operatorname{sw}({t_{\rm rk}}).
 \end{gathered}
\]
The other two wedge coordinates are the primitive Wang relations.
The root-disc prisms give
\[
 \begin{gathered}
 K_{w,\rm rk}=m{W_{\rm bd}}-{c_{\rm rk}}{f_{\rm bd}},\\
 K_{z,\rm rk}=mZ_0+{c_{\rm rk}}A_0-q_0E_0,\\
 K_{t,\rm rk}=mT_0+{c_{\rm rk}}B_0+{p_{\rm rk}}E_0.
 \end{gathered}
\]
The vertical ${f_{\rm bd}}$ sweeps all $m$ components, and the other vertical
coordinates induce the exterior square of the collapsed lattice
inclusion. An integral cap basis is specified by the boundary lifts
\[
 \begin{gathered}
 {C_{\rm rk}}={r_{\rm rk,Bez}}{f_{\rm bd}}+{u_{\rm rk,Bez}}{W_{\rm bd}},\quad J_{z_{\rm rk}}={r_{\rm rk,Bez}}A_0-{u_{\rm rk,Bez}}Z_0,\\
 J_{t_{\rm rk}}={r_{\rm rk,Bez}}B_0-{u_{\rm rk,Bez}}T_0,\quad {E_{\rm rk}}=E_0.
 \end{gathered}
\]
Their normalization and collapsed-lattice images are the ruling
sphere, $\eta_0\wedge {z_{\rm rk}}$, $\eta_0\wedge {t_{\rm rk}}$ and ${z_{\rm rk}}\wedge {t_{\rm rk}}$.
The full boundary map is therefore
\[
 {\mathsf F_{\rm rk}}=\begin{pmatrix}
 m&0&0&0&{c_{\rm rk}}&0&0\\
 0&m&0&0&0&-{c_{\rm rk}}&0\\
 0&0&m&0&0&0&-{c_{\rm rk}}\\
 0&{u_{\rm rk,Bez}}q_0&-{u_{\rm rk,Bez}}{p_{\rm rk}}&1&0&{r_{\rm rk,Bez}}q_0&-{r_{\rm rk,Bez}}{p_{\rm rk}}
 \end{pmatrix}.
\]
Let ${\mathsf S_{\rm rk}}$ be the four displayed lift columns, ${\mathsf K_{\rm rk}}$ the three kernel
columns, and
\[
 {\mathsf R_{\rm rk}}=\begin{pmatrix}
 -{u_{\rm rk,Bez}}&0&0&0&{r_{\rm rk,Bez}}&0&0\\0&{u_{\rm rk,Bez}}&0&0&0&{r_{\rm rk,Bez}}&0\\0&0&{u_{\rm rk,Bez}}&0&0&0&{r_{\rm rk,Bez}}
 \end{pmatrix}.
\]
Then
\[
\begin{gathered}
 {\mathsf F_{\rm rk}}{\mathsf K_{\rm rk}}=0,\qquad {\mathsf F_{\rm rk}}{\mathsf S_{\rm rk}}=I,\\
 {\mathsf R_{\rm rk}}{\mathsf K_{\rm rk}}=I,\qquad I-{\mathsf S_{\rm rk}}{\mathsf F_{\rm rk}}={\mathsf K_{\rm rk}}{\mathsf R_{\rm rk}},
\end{gathered}
\]
by ${u_{\rm rk,Bez}}{c_{\rm rk}}+{r_{\rm rk,Bez}}m=1$.
These equations give the section and the inverse coordinates of
the entire integral kernel. Its projection is
$m\operatorname{Fix}{T_{\rm rk}}$, so this cap has normal index $m^3$.
At $m=1$ the same formulas give index one. The original
zero-line transitive multiplicative caps have these order-one maps
under the original-to-primitive period isomorphism of
\ref{maps:src:D.2.26}.

\emph{The exterior lattice and the full normal determinant.}
Let $N$ be the number of original singular places and let $T_i$ be
their marked four-period matrices. The period formula of
\ref{maps:src:D.2.23} constructs them from $A_i,g_i$ by
\[
 \beta_i=-a g_i^tJA_i,\qquad
 T_i=\begin{pmatrix}A_i&0&bg_i\\
                  \beta_i&1&\beta_i b h_i+w_i\\
                  0&0&1\end{pmatrix},\qquad
 J=\begin{pmatrix}0&1\\-1&0\end{pmatrix}.
\]
Here $h_i=(A_i-I)^{-1}g_i$ at an additive place; at $I_\nu$,
$A_i-I=\nu {u_{\rm ell}}{u_{\rm ell}}^tJ$, $g_i=\kappa {u_{\rm ell}}$, and
$h_i=(\kappa/\nu){r_{\rm ell}}$ with $\det({u_{\rm ell}},{r_{\rm ell}})=1$.
The weight $w_i$ is $1/\nu$ at the selected zero and zero elsewhere.
All entries are integral, $\prod_iT_i=I$, and $T_i^t\omega T_i=\omega$.

Set ${\mathsf D_{\rm gr}}=[T_1-I,\ldots,T_{N-1}-I]$. For the first $N-1$ boundary
circles let $j_i$ be the block inclusion, and put
\[
 j_N(v)=-(T_2\cdots T_{N-1}v,
                 T_3\cdots T_{N-1}v,\ldots,v).
\]
These are the oriented cut-graph maps. Let $F_i$ be an integral
basis of $\operatorname{Fix}T_i$ starting with $\delta$ and set
${\mathsf B_{\rm gr}}=[j_iF_i]_i$. Integer reduction of the matrices just given yields
three unit Smith entries for ${\mathsf D_{\rm gr}}$; the columns of ${\mathsf B_{\rm gr}}$ span
$\ker {\mathsf D_{\rm gr}}$ integrally with sole relation $\sum_i j_i\delta=0$.
One direct method is to choose ${\mathsf U_{\rm gr}}{\mathsf D_{\rm gr}}{\mathsf V_{\rm gr}}=\operatorname{diag}(I_3,0)$,
express ${\mathsf B_{\rm gr}}$ in the last coordinates of ${\mathsf V_{\rm gr}}^{-1}$, and reduce its
column lattice to Hermite form. That form is the identity of size
$4N-7$. Its inverse column operations give a right inverse from
these kernel coordinates; two lifts differ by the displayed
primitive relation. Substitution of the component equations gives
this calculation for all $172$ natural rows, in particular for the
sixteen rows above.

The matrix with columns $\bigwedge^2T_i-I$ has primitive image
$\ker\omega$ and rank five. This follows by the same integral
reduction: in the kernel coordinates of the primitive row
$(a,0,0,0,0,-b)$, its Hermite form is $I_5$.
Choose $a{x_{\rm fib,Bez}}-b{y_{\rm fib,Bez}}=1$; then ${x_{\rm fib,Bez}} e_{12}+{y_{\rm fib,Bez}} e_{34}$ is an inverse
representative for the generator of the fibre quotient $\mathbb Z{s_{\rm fib}}$.
The image relations are the monodromy cylinders, so these integral
coordinate operations give inverse maps on the homology of the
marked period bundle.

Sweeping $\delta$ over the cut base gives the sole global relation
\[
                 \sum_i t_i(\delta)+be\,{s_{\rm fib}}=0.
\]
Its fibre coefficient is zero in the contracting line frame.
At the cusp the origin multiplier is removed by its $\delta$ deck
period and reciprocal, as in \ref{maps:src:D.2.25}; its short-lattice
pairing with $\delta$ is zero. A change of peripheral lifts has
the stated coefficient $be$ by \ref{maps:src:D.2.23}.
Each additional regular hole inserts an identity block in the graph
complex. Its block projection and insertion are inverse maps, so
the formula covers every finite number of regular sites.

For a cap of root order $m_i$ and power $v_i$, sweep the positive
root disc by $\delta$. Its $m_i$ sectors close along the seam
$(x,y)\mapsto x\delta+yv_i$. The boundary is
$m_i\,t_i(\delta)-\delta\wedge v_i$. Since
$\omega(\delta,v_i)=-b n_i$, its integral relation in the exterior is
\begin{equation}
                  m_i t_i(\delta)+bn_i {s_{\rm fib}}=0.                \hypertarget{tag-024W}{}\label{maps:eq:D.ST10}
\end{equation}
The inverse cut uses those same real period parameters. This fixes
the sign by the positive normal-disc boundary and the ordered sweep
$(\delta,\text{meridian})$. The node charts and unit-frame inverses
used above retain this framing for the star and multiplicative roots;
the finite and regular Thom maps of \ref{maps:src:D.2.26} give the
same sector-and-seam chain. At an unchanged cap $m_i=1,n_i=0$ in
its contracting meridian marking.

For every minimal cap its vertical Wang quotient maps injectively
to $H_2$ and its boundary maps onto $H_2$. These assertions follow
from the star and multiplicative maps above and the finite, regular,
unchanged additive and index-one cusp maps of
\ref{maps:src:D.2.26}. Projection identifies the full kernel $K_i$
with a lattice $N_i\subset\operatorname{Fix}T_i$ of index $\kappa_i$.
As at the star, lifting a cap class to the boundary and projecting
gives the inverse quotient map
\[
 \operatorname{Fix}T_i/N_i\simeq
 H_2(U_i)/\operatorname{im}H_2(T^4)_{T_i}.
\]
Choose an integral basis matrix $P_i$ of $N_i$. Write $f_i$ for
the rational vertical functional of its kernel graph after applying
$\omega$. Its relation columns are $(-f_iP_i,P_i)$, which are
integral. Equation \eqref{maps:eq:D.ST10} gives
$f_i(\delta)=-bn_i/m_i$.

Adjoin the global relation to these columns. Write $N_{\partial}$ for
the total number of boundary holes. On ${s_{\rm fib}}$ and all the fixed
sweep coordinates the square matrix is
\[
 M=\begin{pmatrix}
 be&-f_1P_1&\cdots&-f_{N_{\partial}}P_{N_{\partial}}\\
 d_1&P_1&&0\\
 \vdots&&\ddots&\\
 d_{N_{\partial}}&0&&P_{N_{\partial}}
 \end{pmatrix},
\]
where $d_i$ is the coordinate vector of $\delta$.
Multiplication on the left by the block upper triangular matrix
with first row $(1,f_1,\ldots,f_{N_{\partial}})$ and identity diagonal blocks
has determinant one; its inverse replaces each $f_i$ by $-f_i$.
The new first row is $(be+\sum_i f_i(\delta),0,\ldots,0)$.
Consequently the complete integer determinant is
\begin{equation}
 \det M=\left(\prod_i\det P_i\right)
                   b\left(e-\sum_i n_i/m_i\right),\qquad
 |\det P_i|=\kappa_i.                                    \hypertarget{tag-024X}{}\label{maps:eq:D.ST11}
\end{equation}
This calculation includes arbitrary normal ranks and all affine
coefficients. The indices are those of the full integral kernels,
including the order-two torsion of $U_4$.

For the following index and degree-one calculations, take the complete
threefold to be simply connected and every cap to have minimal rational
rank. The Seifert presentation of \ref{maps:src:D.2.23} has
abelianization determinant
\[
 e\prod_i m_i-\sum_i n_i\prod_{j\ne i}m_j=\pm1.
\]
If a prime divided both $m_i$ and $n_i$, it would divide every
term on the left. Thus $\gcd(m_i,n_i)=1$ for each root, giving
the B\'ezout condition used in the following integral maps.

The local indices are one for unchanged minimal additive and
multiplicative caps and the selected index-one cusp; $m_i^3$ for
an additional transitive multiplicative root; $m_i/q_i$ for a
finite good-reduction cap; and $m_i^3$ for a regular root.
For the finite index, the projected Thom columns are
$(m_i/q_i,0),(0,m_i),(r_i,n_i)$. If $A_{i,\rm Bez} n_i+B_{i,\rm Bez} m_i=1$, their
combination $A_{i,\rm Bez}(r_i,n_i)+B_{i,\rm Bez}(0,m_i)$ has second coordinate one,
and the inverse column substitutions leave $(m_i/q_i,0)$.
For a regular root the projected lattice is
$m_i\Lambda+\mathbb Z v_i$. The primitive class of $v_i$ modulo
$m_i$ has order $m_i$; the quotient therefore has order $m_i^3$.
A B\'ezout functional $\lambda_{i,\rm Bez}(v_i)+\alpha_{i,\rm Bez}m_i=1$ gives the
inverse coset coordinates and the full lattice, as in the regular
Thom map of \ref{maps:src:D.2.26}.

All the degree-one groups here are free integrally.
For the boundaries this follows by reducing the displayed $T_i-I$:
their nonzero Smith factors are one in each retained OS50 row.
Additional roots keep these boundary matrices; regular boundaries
are tori. The exterior is free on its base graph and the surviving
$\gamma$, since its fibre relation lattice is $\ker\chi$ primitively.
The primitive and cusp cap groups are free quotients by their short
periods. The minimal original additive caps have the product or
mapping-torus inverse coordinates of \ref{maps:src:D.2.26}.
The transitive multiplicative group has the free basis
$(\eta_0,z_{\rm rk},t_{\rm rk})$ above. The star conductor calculation gives its free abelian
rank-two group. For finite caps the primitive relation
$m_i g_{i,\rm fin}=v_i$ and the B\'ezout condition $\gcd(m_i,n_i)=1$ give a
free abelian rank-two group; for regular caps the group is the free lattice
$\Lambda+\mathbb Z(v_i/m_i)$, with the same B\'ezout inverse.

If $c_i^{\rm fix}=\operatorname{rank}\operatorname{Fix}T_i$, a boundary has
rank $c_i^{\rm fix}+1$ and its minimal cap rank $c_i^{\rm fix}$. With ${N_{\partial}}$ holes the
exterior first homology has rank ${N_{\partial}}$. The degree-one
Mayer--Vietoris source and target therefore both have rank
$\sum_i c_i^{\rm fix}+{N_{\partial}}$. The degree-zero map is injective, so the cokernel
of this map is $H_1(X)$. For a simply connected threefold the map is
an isomorphism between free abelian groups. Hence $H_2(X)$ is exactly
the finite cokernel of $M$.

Let $Q_o$ be the product of the indices $q_i$ of the other finite
good-reduction caps. Let ${R_{\rm ord}}$ be the product of the orders of all
regular roots and all additional transitive multiplicative roots.
Set $q_*=m_*/\kappa_*$. The Seifert determinant is a unit, so
\[
 \left|e-\sum_i n_i/m_i\right|=\left(\prod_i m_i\right)^{-1}.
\]
Using all the local indices in \eqref{maps:eq:D.ST11} gives
\begin{equation}
 |H_2(X)|=\frac{b}{q_*Q_o}\,{R_{\rm ord}}^2,\qquad
 q_*=
 \begin{cases}2,&m_*=2,\\1,&m_*=4\text{ with swap}.
 \end{cases}                                              \hypertarget{tag-024Y}{}\label{maps:eq:D.ST12}
\end{equation}
For example, take profile $50$, $(a,b)=(3,4)$, and the selected
$I_1$ at position $3$. At its unselected $I_3$ the matrix is
\[
 {T_{\rm ex}}=\begin{pmatrix}1&3&0&-8\\0&1&0&0\\0&6&1&-16\\0&0&0&1\end{pmatrix}.
\]
The basis $(e_1+2\delta,3e_2+\gamma,\delta,8e_2+3\gamma)$
has determinant one and gives $I+E_{12}$. Its second vector $L$
has return one. The order-three root with power $L$ has deck map
$(s,Z)\mapsto(\zeta s,sU_{\rm per}(q)^{1/3}Z)$ and the corresponding
third of its transverse period. Its cube is the $L$ period and it
cyclically permutes the three components, giving the free quotient
above. Take the order-two star with return $-1$ and $e=0$.
The group relations $c_*c_3=1$, $c_*^2=\gamma^{-1}$,
$c_3^3=\gamma$ imply $c_*=\gamma=1$.
The indices are $1$ and $27$, and
$|H_2|=27\cdot4\,|1/2-1/3|=18$. This example also exhibits the
multiplicative-root factors included in ${R_{\rm ord}}$.

Simple connectivity permits at most two nontrivial root orders,
which are pairwise coprime, by \ref{maps:src:D.2.23}. Since $m_*$
is even, all other such orders are odd. The only possible other
finite $q_i$ is three at $IV$; an order-six $II$ root cannot
coexist with the star. Hence $Q_o$ and ${R_{\rm ord}}$ are odd. In
\eqref{maps:eq:D.ST2}, $b$ has precisely two factors of two.
Equation \eqref{maps:eq:D.ST12} gives exactly one factor of two
in $|H_2|$ for $m_*=2$ and exactly two for $m_*=4$.

\emph{Completion over all fillings and refinements.}
An unselected original $I_3$ or $IV$ with nonzero line degree
has local rational defect one by its primitive Euler and component
intersection maps. A zero-line original cap with more than one
component orbit has the corresponding positive component quotient.
The permitted original $II$ has one zero-degree component. These
are the original-cap maps of \ref{maps:src:D.2.26}; the star
originals were settled above. All positive local quotients survive
under the projection of the global Mayer--Vietoris cokernel.

At the selected cusp the natural index is one. A further root of
degree $d$ multiplies both long valuations by $d$; the primitive
triangulation has rays $(v,1)$ for $v\in\mathbb Z^2$, modulo
the long-period translations $d\mathbb Z^2$. These rays give the
central divisors, hence $d^2$ cover components. There are two
maximal triangles in each unit square; the cones over these triangles
give the zero-dimensional toric strata, hence $2d^2$ such strata.
Every positive-dimensional torus orbit has Euler characteristic
zero, so the degree-$d$ free quotient has Euler contribution $2d$.
These are the ray, cone and period maps of \ref{maps:src:D.2.22}.
A homotopy six-sphere therefore requires $d=1$.
Minimal rational rank also forces this value. A quotient of the
$d^2$ components by an order-$d$ group has at least $d$ component
orbits. If their number is
${r_{\rm orb}}$, the normalization sequence gives $b_4(U_{\rm cusp})={r_{\rm orb}}$: the double curves
have zero third and fourth homology. Collapsing the two primitive
short periods gives $b_1(U_{\rm cusp})=2$, while the boundary has $b_1=3$.
The boundary generates the cap group, including its root path, so
the map on $H_1$ is onto. The rational pair sequence and duality
then give local degree-two cokernel rank ${r_{\rm orb}}-3+2={r_{\rm orb}}-1$.
Thus $d>1$ has positive local defect and is excluded from the
minimal-rank scope of \eqref{maps:eq:D.ST12} as well.

An equivariant toric refinement has a proper map to its minimal
cap, equal to the identity on the marked outer collar and of
oriented relative degree one. Extend it by the identity on the
other pieces, obtaining ${f_{\rm ref}}:X'\to X$ of degree one.
On integral degree-two homology it has right inverse
\[
 {f_{\rm ref}}^!=\operatorname{PD}_{X'}\,{f_{\rm ref}}^*\,\operatorname{PD}_X^{-1}.
\]
The projection formula and $(f_{\rm ref})_*[X']=[X]$ give $(f_{\rm ref})_*{f_{\rm ref}}^!=1$,
including torsion. The induced map on fundamental groups is onto:
lift ${f_{\rm ref}}$ to the connected cover ${p_{\rm cov}}:\widetilde X\to X$
corresponding to $\operatorname{im}(f_{\rm ref})_*$. An infinite cover is a
connected noncompact oriented six-manifold, so its top integral
homology is zero; the lifted fundamental class would then give
$\deg {f_{\rm ref}}=0$. For a finite cover of degree ${k_{\rm cov}}$, the lifted map has
integer degree and $1=\deg {f_{\rm ref}}={k_{\rm cov}}\deg\widetilde f_{\rm ref}$, so ${k_{\rm cov}}=1$.
Thus the minimal threefold's obstruction survives
every such refinement; the relative formula gives the same result
for an individual cap's local quotient. The excluded returns have
nontrivial fundamental group, the excluded component actions have
positive rational homology, and the two remaining star caps have
the nonzero two-primary groups just computed. All permitted
fillings and affine parameters are covered.
\end{proof}

\subsection{The classification over twenty-one profiles}
\begin{corollary}[Sphere fillings over twenty-one elliptic profiles]\hypertarget{tag-02L3}{}\label{maps:src:D.2.28}
{\normalfont\scriptsize tag 02L3}\par

For the natural line $M=\mathcal O(P-O)$, consider all section pairs,
component-supported linearization zeros, permitted original and
Mumford fillings, additive and regular logarithms, multiplicities
and affine twists in the sector of \ref{maps:src:D.2.22}.
Use the marked filling patterns of
Definition~\ref{maps:def:marked-filling-pattern}. The complete profile classification is
\begin{equation}
\begin{array}{c|l|r|l}
\text{profile}&\text{singular fibres}&\text{sphere patterns}
                     &\text{sphere multiplicities}\\\hline
43&III^*,3I_1&3&\text{unbounded}\\
43/II&III^*,II,I_1&1&\text{unbounded}\\
45&I_8,4I_1&4&\text{all orders }1\\
45/II&I_8,II,2I_1&4&\text{unbounded}\\
46&I_3^*,3I_1&0&\text{none}\\
46/II&I_3^*,II,I_1&0&\text{none}\\
47&I_7,I_2,3I_1&4&\text{all orders }1\\
47/II&I_7,I_2,II,I_1&4&\text{unbounded}\\
47/III&I_7,III,2I_1&4&\text{unbounded}\\
49&IV^*,I_2,2I_1&3&\text{unbounded}\\
49/II&IV^*,II,I_2&1&\text{unbounded}\\
49/III&IV^*,III,I_1&2&\text{unbounded}\\
50&I_1^*,I_3,2I_1&0&\text{none}\\
50/II&I_1^*,II,I_3&0&\text{none}\\
50/IV&I_1^*,IV,I_1&0&\text{none}\\
55&I_5,I_4,3I_1&4&\text{all orders }1\\
55/II&I_5,I_4,II,I_1&4&\text{unbounded}\\
56&I_5,I_3,I_2,2I_1&5&\text{all orders }1\\
56/II&I_5,I_3,I_2,II&6&\text{unbounded}\\
56/III&I_5,I_3,III,I_1&5&\text{unbounded}\\
56/IV&I_5,I_2,IV,I_1&5&\text{unbounded}.
\end{array}                                                  \hypertarget{tag-024Z}{}\label{maps:eq:D.ST13}
\end{equation}
There are $59$ patterns on $50$ distinct natural initial rows,
listed in \ref{maps:src:D.2.26}. Their full affine and multiplicity
conditions are \eqref{maps:eq:D.HY3}--\eqref{maps:eq:D.HY5}, with
the fundamental group given by \eqref{maps:eq:D.FG1} and its
triviality tests in \eqref{maps:eq:D.FG3}.
Thus twelve profiles admit unbounded sphere multiplicities,
four require every logarithmic order to be one, and five admit none.
Every threefold satisfying these sphere conditions is orientation-preservingly diffeomorphic to
the standard $S^6$.
\end{corollary}


\begin{proof}
The finite initial reduction in \ref{maps:src:D.2.22} applies to
every allowed zero and section pair. For a selected $I_\nu$ place,
write
\[
A-I=\nu uu^tJ,\qquad g=\kappa u,\qquad
J=\begin{pmatrix}0&1\\-1&0\end{pmatrix},\qquad
\langle G,G\rangle=1/d.
\]
Its complete component test is
\[
ab=d/\nu,\qquad \gcd(b\kappa,\nu)=1,
\]
together with the equality in the free abelian group on $\mathbb Z/\nu$
\[
[a\kappa]-[0]-[(a-b)\kappa]+[-b\kappa]
       =-2[h]+[h-1]+[h+1].
\]
At each other reducible place, $P$ or $Q$ is narrow. Enumerating
the positive divisors $a$ of $d/\nu$, setting $b=d/(\nu a)$ and
testing every component $h$ gives the $172$ rows of that theorem;
each has $h=0$. This exhausts the section pairs and zero positions.
The selected cusp has index one, since a further root of degree
$m$ would contribute Euler characteristic $2m$.

The two $I_3^*$ profiles contribute eight initial rows, all excluded
by \ref{maps:src:D.2.24}. The three $I_1^*$ profiles contribute
sixteen rows, all excluded by \ref{maps:src:D.2.27}. Their integral
obstructions include every allowed star return, affine twist and
refinement. The remaining $148$ initial rows lie in the sixteen
profiles of \ref{maps:src:D.2.26}.

At an additive place the eligible original caps have zero line
degree and one component orbit. They are II and the transitive
III and IV cases. For component coordinates $g_j$ modulo $s_j$,
the order of translation by $Q$ is
\[
\operatorname{lcm}_j\frac{s_j}{\gcd(s_j,bg_j)}.
\]
For a finite logarithmic cap, the integral fixed vector of
$\gamma$-character one exists precisely when
$b(A-I)^{-1}g$ is integral. The monodromy orders and vertical
indices are those of \eqref{maps:eq:D.HY2}. The local quotient
and global bridge calculation in \ref{maps:src:D.2.26} excludes
every other original cap from a sphere. At each remaining
multiplicative place its period and Tate maps identify the
original and primitive Mumford choices, with the inverses given there.

For a simply connected threefold put
\[
Q_f=\prod_{i\text{ finite}}q_i,\qquad c=b/Q_f,\qquad
R=\prod_{j\text{ regular}}m_j.
\]
The integral presentation and inverse substitutions of
\ref{maps:src:D.2.26} give \eqref{maps:eq:D.HY4}, whose order is
$cR^2$. Therefore the complete sphere condition is $c=R=1$
together with the fundamental-group criterion
\eqref{maps:eq:D.FG3}. These are exactly
\eqref{maps:eq:D.HY5}, retaining all affine coordinates in
\eqref{maps:eq:D.HY3}.

For each initial row take every eligible additive choice above,
keep $b=Q_f$, and impose the pairwise coprimality of finite
multiplicities required by \eqref{maps:eq:D.FG3}.
This gives the $59$ patterns listed in \ref{maps:src:D.2.26}.
Projecting a pattern to its profile, ordered section pair and
zero position gives $50$ distinct initial rows. Grouping those
patterns by profile gives \eqref{maps:eq:D.ST13}.

Every additive profile has a one-finite pattern. At that place,
take any positive multiple $m$ of its monodromy order, numerator
$n=1$ and effective peripheral integer $e=0$. The determinant
is $-1$, and the primitive contracting character makes every
affine lift with that power free. Its $\delta$ coefficient is
arbitrary. Thus each of the twelve additive profiles has
unbounded sphere multiplicities. The unique two-finite pattern
is $49/III$ with $(a,b)=(1,6)$ and both additive places logarithmic.
Its orders are coprime multiples of three and four; B\'ezout
numerators give
$m_2n_1+m_1n_2=1$ at $e=0$, while the full determinant condition
in \eqref{maps:eq:D.FG3} describes every admissible numerator.
All their remaining affine coordinates are unrestricted by
\ref{maps:src:D.2.26}.

On the four all-multiplicative profiles, $Q_f=1$, so $b=1$ and
every logarithmic order is one. Eliminating an order-one site
by $c_j=\gamma^{n_j}$ changes $e$ to $e-n_j$; its inverse
reintroduces that generator and relation. In the canonical
marking this gives $|\sum_j n_j|=1$. The five excluded profiles
give the remaining zero rows of the table.

Finally \ref{maps:src:D.2.26} gives the oriented global
diffeomorphism to $\widetilde\Sigma_0N_R$ for $c=1$.
Here $R=1$ and $N_1=S^5$. Its explicit boundary rotation
$(x,z)\mapsto(\alpha_zx,z)$, with inverse using
$\alpha_z^{-1}$, identifies the two-piece model with
$\partial(D^5\times D^2)$. Positive radial corner rounding
identifies that boundary with the standard oriented $S^6$.
\end{proof}



\subsection{Finite caps with a primitive contracting period}

\begin{definition}[Primitive fixed finite-cap data]
\hypertarget{tag-Q121}{}\label{mapsIII:def:primitive-fixed-finite-data}
{\normalfont\scriptsize tag Q121}\par
Take a natural initial row of Theorem~\ref{maps:src:D.2.22}, with
section pair $(P,Q)=(aG,bG)$ and the marked period lattice of
Definition~\ref{mapsIII:def:elliptic-marked-data}. Write
$\omega=a\,e_1^*\wedge e_2^*-b\,\delta^*\wedge\gamma^*$.
At a finite additive position with a marked integral fixed vector
$\ell_i$ satisfying $\gamma^*(\ell_i)=1$, take $(d_i,q_i)$ from
\eqref{maps:eq:D.HY2} and set
\[
 m_i=d_i k_i,\qquad v_i=n_i\ell_i+r_i\delta,\qquad
 t_i=\gcd(m_i,n_i,q_i r_i),\qquad
 k_i\in\mathbb Z_{>0},\quad n_i,r_i\in\mathbb Z.
\]
For a global filling, use a nonstar initial row, the minimal
original and Mumford caps of Theorem~\ref{maps:src:D.2.26}, any
number of regular logarithmic caps with their free affine vectors,
and an integer clutch $e\in\mathbb Z$.
\end{definition}

\Needspace{6\baselineskip}
\begin{theorem}[Finite caps with primitive fixed quotient]
\hypertarget{tag-02L4}{}\label{maps:src:D.2.29}
{\normalfont\scriptsize tag 02L4}\par
There are $68$ finite additive positions with the primitive fixed
vector of Definition~\ref{mapsIII:def:primitive-fixed-finite-data}.
The local assertions below hold at all of them. Their restriction
to the nonstar initial rows consists of $62$ positions; the global
assertions concern those rows and the fillings in that definition.
The cyclic affine action with this power vector is free exactly
for the data
\begin{equation}
 \gcd(m_i,n_i,r_i)=1,\qquad d_i\mid m_i/t_i.
 \hypertarget{tag-025Z}{}\label{maps:eq:D.UF1}
\end{equation}
For these data $t_i\in\{1,q_i\}$.
The value $t_i=q_i>1$ is allowed precisely for $q_i\mid k_i$.
Let $F_i$ be the central quotient surface and $L_i$ its complex
normal line. Their integral invariants are
\begin{equation}
\begin{gathered}
 H_1(F_i)=\mathbb Z^2\oplus\mathbb Z/t_i,\qquad
 H_2(F_i)=\mathbb Z^2\oplus\mathbb Z/t_i,\\
 H_3(F_i)=\mathbb Z^2,\qquad H_0(F_i)=H_4(F_i)=\mathbb Z,
 \qquad \operatorname{ord}c_1(L_i)=t_i .
\end{gathered}
\hypertarget{tag-0261}{}\label{maps:eq:D.UF2}
\end{equation}
An order-one Chern class is zero.

The degree-two gluing cokernel is generated by a fibre class $s$,
the fixed sweeps $D_i,L_i'$ of $\delta,\ell_i$ at each finite log,
and the four sweeps $t_j(\Lambda)$ at each regular log.
Its finite relations are
\begin{equation}
\begin{aligned}
 (bn_i/q_i)s+(m_i/q_i)D_i&=0,\\
 -br_i s+m_iL_i'&=0,\\
 r_iD_i+n_iL_i'&=0.
\end{aligned}
\hypertarget{tag-0262}{}\label{maps:eq:D.UF3}
\end{equation}
Its regular relations are
$m_jt_j(u)-\omega(u,v_j)s=0$ for $u\in\Lambda$ and $t_j(v_j)=0$,
as in \eqref{maps:eq:D.SC2}. Its remaining relation is
\begin{equation}
 \sum_iD_i+\sum_jt_j(\delta)+be\,s=0.
 \hypertarget{tag-0263}{}\label{maps:eq:D.UF4}
\end{equation}
Let $T$ be the torsion subgroup of this cokernel.
Write the abelianization of the Seifert presentation in
Theorem~\ref{maps:src:D.2.23} as $\mathbb Z^{r_{\rm glob}}\oplus A_1$, using numerator
$n_i$ at each finite site and $\gamma^*(v_j)$ at each regular site.
Then $r_{\rm glob}\in\{0,1\}$ and
\begin{equation}
\begin{array}{c|c}
 j&H_j(X;\mathbb Z)\\ \hline
 0,6&\mathbb Z\\
 1&\mathbb Z^{r_{\rm glob}}\oplus A_1\\
 2,3&\mathbb Z^{2r_{\rm glob}}\oplus T\\
 4&\mathbb Z^{2r_{\rm glob}}\oplus A_1\\
 5&\mathbb Z^{r_{\rm glob}} .
\end{array}
\hypertarget{tag-0264}{}\label{maps:eq:D.UF5}
\end{equation}
The fundamental group is the full Seifert group of
Theorem~\ref{maps:src:D.2.23}. These assertions include every free affine
twist satisfying \eqref{maps:eq:D.UF1} and every integer clutch $e$.
\end{theorem}


\begin{proof}
Fix the integral marking of one finite site, suppressing
its index. Write $A$ for its monodromy and $\mathsf P_{\rm fix}=N_d/d$ for the
projection to its fixed plane. The integral local calculation in
Theorem~\ref{maps:src:D.2.26} gives
\begin{equation}
\begin{gathered}
 \operatorname{Fix}A=\mathbb Z\delta\oplus\mathbb Z\ell,\qquad
 \mathsf P_{\rm fix}\Lambda=\mathbb Z(\delta/q)\oplus\mathbb Z\ell,\\
 \Lambda/(A-I)\Lambda=\mathbb Z\beta\oplus\mathbb Z\ell,\qquad
 [\delta]=q\beta,\qquad N_m\beta=(m/q)\delta .
\end{gathered}
\hypertarget{tag-0265}{}\label{maps:eq:D.UF6}
\end{equation}
Here the quotient map is integral: for
$\mathsf U_{\rm Sm}(A-I)\mathsf V_{\rm Sm}=\operatorname{diag}(1,1,0,0)$, take the last two
coordinates of $\mathsf U_{\rm Sm}$, followed by the unimodular change making
the quotient coordinates of $\delta$, divided by $q$, and
those of $\ell$ the coordinate vectors.
The last two columns of $\mathsf U_{\rm Sm}^{-1}$ with the inverse change give
a section. The reverse composite differs from the identity
by $(A-I)$ applied to the first two Smith coordinates.
Thus these formulas specify the quotient and its inverse.

The local position set retains the marked tuple consisting of
the profile, section pair, zero and finite site.
The nonstar positions form a subset with the same marked tuples.
The primitive-fixed-quotient complement consists of the two
section pairs $(1,4),(4,1)$ at the type-II site of each of the
profiles $46/\mathrm{II}$ and $50/\mathrm{II}$, and the pairs
$(1,12),(4,3)$ at the type-IV site of $50/\mathrm{IV}$.
This gives the six positions between the local count $68$
and the global count $62$.

Every affine action of power $v=n\ell+r\delta$ has the invariant
translation representative
\[
 h(x)=Ax+v/m.
\]
On the central torus an affine translation difference $w-v/m$
lies in $\ker \mathsf P_{\rm fix}$. Translation by
$(I-A)^{-1}(w-v/m)$ conjugates this representative to $Ax+w$;
translation by its negative is the inverse.
The comparison extends holomorphically over the good-reduction
root disc. To see this, let $\sigma(u)=\zeta u$ and let $R$
be the linear deck operator on holomorphic sections of the
Lie bundle, so $R^m=I$. Two affine lifts with the same
integral $m$th power have translation difference $b_0$
satisfying $N b_0=0$, where $N=\sum_{j=0}^{m-1}R^j$.
Set
\[
 s_0=-\frac1m\sum_{j=0}^{m-1}jR^j b_0 .
\]
The identity
$(I-R)\sum jR^j=N-mI$ gives $(I-R)s_0=b_0$.
Translation by $s_0$, with inverse translation by $-s_0$,
is the required conjugacy. This is a finite sum of holomorphic
sections, so it extends across the central point.
Changing a universal lift by a period $\lambda$ changes its
power by $N_m\lambda$ and gives the same torus action.
This accounts for the power class in
$\Lambda^A/N_m\Lambda$ before applying the conjugacy.
Multiplication by a nonzero constant in the normal coordinate,
and multiplication by its reciprocal, account for unit choices.
These comparisons preserve the root coordinate and its positive
normal orientation.

For $d\nmid j$, the map $A^j-I$ is invertible on $\ker \mathsf P_{\rm fix}$.
The equation $h^j(x)=x$ modulo $\Lambda$ is soluble exactly
when $jv/m\in \mathsf P_{\rm fix}\Lambda$. Indeed, choose $\lambda$ with
$\mathsf P_{\rm fix}\lambda=jv/m$ and solve
$x=(A^j-I)^{-1}(\lambda-jv/m)$ on $\ker \mathsf P_{\rm fix}$.
By \eqref{maps:eq:D.UF6} the obstruction vanishes exactly when
$m\mid jn$ and $m\mid jqr$. Its positive indices are the
multiples of $m/t$.
Put $g_{\rm aff}=\gcd(m,n,r)$. For $g_{\rm aff}>1$ the power $m/g_{\rm aff}$ has integral
translation, and hence has a fixed point. For $g_{\rm aff}=1$, every
proper power divisible by $d$ has identity linear part and
nonintegral translation. The remaining obstruction indices
are all divisible by $d$ exactly when $d\mid m/t$.
This proves \eqref{maps:eq:D.UF1} for every proper power.
Since $q\in\{1,2,3\}$, the first gcd condition implies
$t\in\{1,q\}$; at $t=q>1$ the second condition is $q\mid k$.
For the resulting finite group the closed fixed set of each
nonidentity element misses the compact central torus.
A sufficiently small common root disc therefore gives a
free action on the whole neighborhood.

The quotient surface has the affine fundamental group
generated by $\Lambda$ and $c$, with
$c\lambda c^{-1}=A\lambda$ and $c^m=v$.
The maps to the affine group and back send these generators
to their marked translations and to $h$.
Its abelianization is
\begin{equation}
 H_1(F)=
 \langle\beta,\ell,c\mid mc-n\ell-qr\beta=0\rangle
 =\mathbb Z^2\oplus\mathbb Z/t .
 \hypertarget{tag-0266}{}\label{maps:eq:D.UF7}
\end{equation}
An integral B\'ezout completion of $(qr,n,m)/t$ supplies
the two free coordinates and the inverse Smith map.
The finite torus cover gives rational Betti numbers
$(1,2,2,2,1)$ by averaging its exterior forms.
Integral Poincar\'e duality and the universal coefficient
theorem then give the groups in \eqref{maps:eq:D.UF2}.
The normal root character sends $c$ to $1/m$ in
$\mathbb R/\mathbb Z$ and sends $\Lambda$ to zero.
On the primitive torsion generator
\[
 z=(m/t)c-(n/t)\ell-(qr/t)\beta
\]
it has value $1/t$.
Its Bockstein is $c_1(L)$, so this class has exact order $t$.
The flat character gives no free Chern component.
On each torsion Smith coordinate the duality identification
sends $k$ to the character $j\mapsto kj/t$; evaluation at
the coordinate generator is its inverse.

Let $E$ be the boundary of the normal disk bundle $U$.
The mapping-torus complexes for $A-I$ and
$\bigwedge^2A-I$ have unit nonzero Smith factors, of ranks
two and four. Consequently $H_1(E)=\mathbb Z^3$.
The Wang sequence for $H_2(E)$ has fibre quotient
$\mathbb Z^2$ and fixed-loop quotient
$\mathbb Z\delta\oplus\mathbb Z\ell$.
The sweeps of the two fixed loops give a section; subtracting
them and then reading the fibre quotient is its inverse.
Thus $H_2(E)=\mathbb Z^4$ with the stated marking.
Since $c_1(L)$ is torsion, cap product
$H_2(F)\to H_0(F)$ is zero. The Thom sequence gives a
surjection $H_2(E)\to H_2(F)$, with full kernel the image
of $H_1(F)$ under the normal Thom map.

The marked boundary identification is
\[
 [y,\tau]\longmapsto
 [y+(\tau/m)v,\exp(2\pi i\tau/m)] .
\]
The mapping-cylinder relation $(y,\tau)\sim(Ay,\tau+1)$
is exactly the deck relation on the right.
The inverse chooses a root argument $\tau$ and subtracts
$(\tau/m)v$; changing the argument gives that same relation.
The positive normal circle, with angular parameter $\theta$, is
$(y,\tau)\mapsto(y-\theta v,\tau+m\theta)$.
Write $c_{\partial}$ and $v_{\partial}$ for the marked boundary lifts
of $c$ and the lattice translation $v$. Its central group representative is
$\mu=c_{\partial}^m v_{\partial}^{-1}$. The projection $E\to F$
sends these lifts to $c,v$ and kills $\mu$, by the relation $c^m=v$.

For a lattice loop $\lambda$, use parameters $\xi,\theta$
for its normal torus, represented by
$[\xi\lambda,\exp(2\pi i\theta)]$ before passage to the quotient.
In the displayed boundary coordinates its vertical
parallelogram is $-\lambda\wedge v$. Cutting the normal
circle into its $m$ mapping-cylinder intervals gives
fixed sweep $N_m\lambda$. The inverse chain construction
glues those intervals along their marked $A$-edges.
For the central generator $c$, use
$[\xi v/m,\exp(2\pi i(\xi/m+\theta))]$.
Its boundary coordinates are $y=-\theta v,\ \tau=\xi+m\theta$;
its vertical area is zero and its positive fixed sweep
is $v$. We have therefore obtained
\begin{equation}
\begin{array}{c|c|c}
 \text{generator}&\text{normal projection}&
                  \text{fibre component}\\ \hline
 \beta&N_m\beta&-\beta\wedge v\\
 \ell&m\ell&-\ell\wedge v\\
 c&v&0 .
\end{array}
\hypertarget{tag-0267}{}\label{maps:eq:D.UF8}
\end{equation}
These are the images of generators of all of $H_1(F)$.
The invariant area functional satisfies
$-\omega(\beta,\ell)=b/q$ and $\omega(\ell,\delta)=b$.
Hence the complete kernel on $(s,D,L')$ has columns
\begin{equation}
 K_\beta=\begin{pmatrix}bn/q\\m/q\\0\end{pmatrix},\qquad
 K_\ell=\begin{pmatrix}-br\\0\\m\end{pmatrix},\qquad
 K_c=\begin{pmatrix}0\\r\\n\end{pmatrix}.
 \hypertarget{tag-0268}{}\label{maps:eq:D.UF9}
\end{equation}
Their primitive relation is
$(m/t)K_c-(n/t)K_\ell-(qr/t)K_\beta=0$.
It is the image of the torsion generator under the Thom map.
The gcd of the projected two-by-two minors is
\begin{equation}
 \gcd(m^2/q,mn/q,mr)=(m/q)t .
 \hypertarget{tag-0260}{}\label{maps:eq:D.UF10}
\end{equation}
This computes the index of the full normal-projection image.
For $\gcd(m,n)=1$, choose $A_0n+B_0m=1$.
The column $A_0K_c+B_0K_\ell=(-bB_0r,A_0r,1)$ eliminates
$L'$. The remaining images of $K_\ell,K_c$ are
$-qA_0rK_\beta,qB_0rK_\beta$.
Restoring that eliminated column is the inverse
transformation. This recovers \eqref{maps:eq:D.HY3}.

We now glue the nonstar exterior and its caps.
The integral exterior calculation in Theorem~\ref{maps:src:D.2.26}
identifies its fibre quotient with $\mathbb Zs$ by $\omega$.
Choose $a u_{\rm fib}-b v_{\rm fib}=1$; the fibre class
$u_{\rm fib}(e_1\wedge e_2)+v_{\rm fib}(\delta\wedge\gamma)$ is a section.
The fixed-sweep quotient is the direct sum of the local
fixed lattices modulo their single diagonal $\delta$
relation. The integral graph Smith map and its inverse
prove that this relation is primitive.
The $\delta$ circle over the cut base supplies its
bounding three-chain. A clutch $e$ contributes $be\,s$.
The primitive graph and fibre calculation also holds on
all natural initial rows, while the present gluing uses
the specified nonstar minimal caps.

The degree-two boundary maps onto the unchanged minimal
original, Mumford and regular caps are the integral maps
in Theorems~\ref{maps:src:D.2.25} and \ref{maps:src:D.2.26}.
Together with the finite surjection just proved, they
allow elimination of every cap two-class from
Mayer--Vietoris. Explicitly, lift that class to its
boundary and subtract its exterior image. The inverse
includes the exterior with zero cap coordinates.
Two choices differ by the full local kernel.
Thus the resulting cokernel $C$ has exactly
\eqref{maps:eq:D.UF3}, the regular relations and
\eqref{maps:eq:D.UF4}.

Over $\mathbb Q$ the finite equations give
$D_i=-(bn_i/m_i)s$ and $L_i'=(br_i/m_i)s$.
The regular equations give
$t_j(u)=\omega(u,v_j)s/m_j$.
The last relation is therefore
\[
 b\left(e-\sum_i\frac{n_i}{m_i}
          -\sum_j\frac{\gamma^*(v_j)}{m_j}\right)s=0.
\]
The free rank of $C$ is consequently $r_{\rm glob}=0$ or $1$,
with exactly the same zero criterion as the rational
Seifert relation matrix.

The degree-one boundary groups are free.
The next Mayer--Vietoris kernel is a free subgroup of
their direct sum. At a site with fixed rank $\rho_i$,
the boundary rank is $\rho_i+1$ and the cap rank is $\rho_i$.
The exterior rank equals the number of boundary circles,
so source and target of that degree-one map have equal
free rank. Its cokernel is $H_1(X)$; hence its kernel
has rank $r_{\rm glob}$. Choose an integral basis of that kernel,
represent it by boundary loops, and choose two-chains
bounding the two images of each loop.
Gluing the two-chains produces cycles $Z_1,\ldots,Z_{r_{\rm glob}}$.
The map
\[
 C\oplus\mathbb Z^{r_{\rm glob}}\longrightarrow H_2(X),\qquad
 (c,k)\longmapsto c+\sum k_i Z_i
\]
is an isomorphism. Its inverse first reads the boundary
coordinates, subtracts those $Z_i$, and uses the preceding
exterior-cokernel identification. Thus
$H_2(X)=\mathbb Z^{2r_{\rm glob}}\oplus T$ integrally.

The index-one cusp has Euler contribution two, and each
other cap and the regular exterior has Euler contribution
zero. Therefore $\chi(X)=2$.
Van Kampen with the marked cusp disks gives the Seifert
generators $\gamma,c_i$, with their root numerators;
the inverse sends the symbols to those marked loops.
This is the full group in Theorem~\ref{maps:src:D.2.23}.
Its abelianization is $\mathbb Z^{r_{\rm glob}}\oplus A_1$.
Integral duality and the universal coefficient theorem
give $H_5=\mathbb Z^{r_{\rm glob}}$, $H_4=\mathbb Z^{2r_{\rm glob}}\oplus A_1$,
and $\operatorname{Tor}H_3=T$.
Finally $\chi(X)=2$ gives $b_3=2r_{\rm glob}$.
The torsion dualities have the Smith-coordinate character
maps described above, proving every group in
\eqref{maps:eq:D.UF5}.
\end{proof}

\begin{remark}[Torsion-core and zero-Euler examples]
\hypertarget{tag-Q122}{}\label{mapsIII:rem:finite-torsion-examples}
{\normalfont\scriptsize tag Q122}\par
With clutch $e=0$, take the $IV^*,II,I_2$ sphere initial row
$(a,b)=(1,3)$ of Theorem~\ref{maps:src:D.2.11}, keep $II$ unchanged and use
$(q,m,n,r)=(3,9,3,1)$ at $IV^*$. The action is free, although
$\gcd(m,n)=3$. The result has
\[
\begin{gathered}
 \pi_1=\mathbb Z/3,\qquad H_1=H_4=\mathbb Z/3,\\
 H_2=H_3=(\mathbb Z/3)^2,\qquad H_5=0.
\end{gathered}
\]
Here the global relation is $D=0$, and the full Thom relations
reduce to $3s=3L'=0$. The quotient map sends $xs+yL'$ to
$(x\bmod3,y\bmod3)$; lifting these residues gives its inverse.
For a $III^*$ sphere initial row with $b=2$ and $e=0$, the free
data $(q,m,n,r)=(2,8,0,1)$
give $\pi_1=\mathbb Z$, middle groups
$H_2=H_3=\mathbb Z^2\oplus\mathbb Z/2$, and
$H_4=\mathbb Z^2$, $H_5=\mathbb Z$.
\end{remark}

\begin{remark}[Finite affine-power and attachment calculations]
\hypertarget{tag-Q123}{}\label{mapsIII:rem:primitive-finite-calculations}
{\normalfont\scriptsize tag Q123}\par
At each of the $68$ marked finite positions, the proper-power
fixed-point equations give exactly \eqref{maps:eq:D.UF1}.
The normal tori of the affine-group generators give the full Thom
columns \eqref{maps:eq:D.UF9}. For the stated nonstar fillings,
adjoin the integral graph relation \eqref{maps:eq:D.UF4} and the
regular-cap relations. The Smith cokernel, the degree-one
Mayer--Vietoris kernel and the Seifert abelianization give
\eqref{maps:eq:D.UF5}, including torsion central caps and zero
rational Euler number. These are the integral maps constructed
in the proof of Theorem~\ref{maps:src:D.2.29}.
\end{remark}

\Needspace{12\baselineskip}
\section{Finite cap lattices and natural catalogues}
\subsection{All finite additive lattices and affine twists}

\begin{definition}[General finite additive-cap data]
\hypertarget{tag-Q124}{}\label{mapsIII:def:general-finite-additive-data}
{\normalfont\scriptsize tag Q124}\par
Take a finite additive position on a natural initial row of
Theorem~\ref{maps:src:D.2.22}, with the section pair, lattice and
area form of Definition~\ref{mapsIII:def:primitive-fixed-finite-data}.
Let $A$ have order $d$, let $q$ have the value in
\eqref{maps:eq:D.HY2}, and let $g$ be the positive generator in
\begin{equation}
 \gamma^*(\operatorname{Fix}A)=g\mathbb Z,\qquad q_v=q/g.
 \hypertarget{tag-01XI}{}\label{maps:eq:D.AF1}
\end{equation}
Write $m=dk$ with $k\in\mathbb Z_{>0}$, and use integer
coefficients $r,n$ for the power vector $v=r\delta+n\ell$
in the fixed basis below. Put $t=\gcd(m,q_vr,gn)$.
\end{definition}

\begin{theorem}[Finite additive lattices and free affine twists]
\hypertarget{tag-02L6}{}\label{maps:src:D.2.30}
{\normalfont\scriptsize tag 02L6}\par
There are $102$ finite additive positions on the $172$ natural
initial rows. At each one, the invariant in
Definition~\ref{mapsIII:def:general-finite-additive-data} satisfies
$g\in\{1,q\}$.
The positive generator $g$ equals $1$ at $68$ positions and
equals $q>1$ at $34$ positions. Integral fixed and coinvariant
bases can be chosen with
\begin{equation}
\begin{gathered}
 \operatorname{Fix}A=\mathbb Z\delta\oplus\mathbb Z\ell,
 \qquad \gamma^*(\ell)=g,\\
 \Lambda/(A-I)\Lambda=\mathbb Z\beta\oplus\mathbb Z\eta,
 \qquad [\delta]=q_v\beta,\quad [\ell]=g\eta,\\
 N_m\beta=(m/q_v)\delta,\qquad
 N_m\eta=(m/g)\ell,\qquad \gamma^*(\eta)=1 .
\end{gathered}
\hypertarget{tag-01XJ}{}\label{maps:eq:D.AF2}
\end{equation}
These bases carry their integral quotient coordinates.
For the power vector in Definition~\ref{mapsIII:def:general-finite-additive-data},
the cyclic affine action is free exactly for
\begin{equation}
 \gcd(m,n,r)=1,\qquad d\mid m/t .
 \hypertarget{tag-01XK}{}\label{maps:eq:D.AF3}
\end{equation}
Its central surface and normal line have the integral invariants
in \eqref{maps:eq:D.UF2}, with this value of $t$.
On the fibre generator $s$ and the fixed sweeps $D,L'$ of
$\delta,\ell$, the full normal Thom kernel has columns
\begin{equation}
 \begin{pmatrix}bgn/q_v\\m/q_v\\0\end{pmatrix},\qquad
 \begin{pmatrix}-br\\0\\m/g\end{pmatrix},\qquad
 \begin{pmatrix}0\\r\\n\end{pmatrix}.
 \hypertarget{tag-01XL}{}\label{maps:eq:D.AF4}
\end{equation}
Its entries are integral and its Seifert numerator is $gn$.

The minimal original/log choices on the $148$ nonstar rows give
$198$ filling patterns. At every finite site use
\eqref{maps:eq:D.AF4}; at any number of regular sites use the
full relations in \eqref{maps:eq:D.SC2}; and impose
\eqref{maps:eq:D.UF4}. The torsion of this cokernel, together
with the Seifert matrix using numerators $g_i n_i$ at finite
sites and $\gamma^*(v_j)$ at regular sites, gives the integral
homology in \eqref{maps:eq:D.UF5}. The fundamental group is
the full presentation in Theorem~\ref{maps:src:D.2.23}.
These formulas hold for every logarithmic multiplicity,
every free affine vector and every integer clutch.

For one finite logarithmic site and no regular logarithmic
sites, put
\begin{equation}
 n'=n-(m/g)e,\qquad h=\gcd(b,m/g,n').
 \hypertarget{tag-01XM}{}\label{maps:eq:D.AF5}
\end{equation}
For $n'\ne0$, the groups are
\begin{equation}
\begin{gathered}
 \pi_1(X)=\mathbb Z/(g|n'|),\qquad
 H_1(X)=H_4(X)=\mathbb Z/(g|n'|),\\
 H_2(X)=H_3(X)=
 \mathbb Z/h\oplus\mathbb Z/\bigl(bg|n'|t/(qh)\bigr),
 \qquad H_5(X)=0 .
\end{gathered}
\hypertarget{tag-01XN}{}\label{maps:eq:D.AF6}
\end{equation}
For $n'=0$, the groups are
\begin{equation}
\begin{gathered}
 \pi_1(X)=\mathbb Z,\qquad H_1(X)=H_5(X)=\mathbb Z,\\
 H_2(X)=H_3(X)=\mathbb Z^2\oplus\mathbb Z/\gcd(b,m/g),
 \qquad H_4(X)=\mathbb Z^2 .
\end{gathered}
\hypertarget{tag-01XO}{}\label{maps:eq:D.AF7}
\end{equation}
In both cases $H_0(X)=H_6(X)=\mathbb Z$, and an order-one
cyclic summand is zero. Every threefold containing a finite
site with $g>1$ has nontrivial fundamental group.
The simply connected classification in
Corollary~\ref{maps:src:D.2.28} therefore has the same scope after
these sites are included.
\end{theorem}


\begin{proof}
We use the physical basis $(e_1,e_2,\delta,\gamma)$.
For the elliptic monodromy $A_0$, generator cocycle $c$
and section pair $(a,b)$, put
\[
 J=\begin{pmatrix}0&1\\-1&0\end{pmatrix},\qquad
 z=(A_0-I)^{-1}c,\qquad B=-a c^tJA_0 .
\]
The full monodromy and invariant area form are
\[
 A=\begin{pmatrix}
 A_0&0&bc\\ B&1&Bbz\\0&0&1
 \end{pmatrix},\qquad
 \omega=a\,e_1^*\wedge e_2^*-b\,\delta^*\wedge\gamma^* .
\]
Let $g$ be the least common multiple of the denominators
of the two coordinates of $bz$. A fixed vector has
elliptic part $-x_4bz$, so integrality is equivalent to
$g\mid x_4$. Thus
\[
 \ell=(-gbz_1,-gbz_2,0,g),\qquad
 (x_1,x_2,x_3,x_4)\longmapsto(x_3,x_4/g)
\]
give the fixed basis $(\delta,\ell)$ and its inverse
on the fixed lattice.

Apply integral Euclidean row and column operations to
$A-I$. Its nonzero invariant factors are $1,1$.
From $U(A-I)V=\operatorname{diag}(1,1,0,0)$ take the last
two rows of $U$ as the quotient map. Divide the quotient
coordinates of $\delta$ by $q_v$ and those of $\ell$ by
$g$. They form a unimodular matrix $C$.
Multiplying the quotient map by $C^{-1}$ gives
coordinates $(\beta,\eta)$; the last two columns of
$U^{-1}$ multiplied by $C$ give an integral section.
The two maps compose to the identity on the quotient.
Their other composite differs from the identity by
$(A-I)$ applied to the first two Smith coordinates.
Multiplication by $N_d$ gives
$N_d\beta=(d/q_v)\delta$ and $N_d\eta=(d/g)\ell$.
Multiplication by $k$ proves \eqref{maps:eq:D.AF2}
for every $m=dk$.

These operations applied to the finite elliptic matrices
and section cocycles in the natural catalogue give
$g=1$ at $68$ positions and $g=q>1$ at $34$ positions.
At the positions with $g=q>1$, the integral quotient has the literal basis
$[\delta],[\gamma]$, and the norm computation reads
\begin{equation}
 g=q,\qquad
 \Lambda/(A-I)\Lambda
   =\mathbb Z[\delta]\oplus\mathbb Z[\gamma],
 \qquad \ell=(q/d)N_d\gamma .
 \hypertarget{tag-01XP}{}\label{maps:eq:D.AF8}
\end{equation}
Hence $\beta=\delta,\eta=\gamma$ are integral lifts in
this sector. The nonzero invariant factors of
$\bigwedge^2 A-I$ are four units at every finite position.
The boundary Wang maps and their fixed-sweep sections
from Theorem~\ref{maps:src:D.2.29} consequently give
$H_1(\partial U)=\mathbb Z^3$ and
$H_2(\partial U)=\mathbb Z^4$ throughout this catalogue.

Average the first argument of $\omega$ over $A$ while
keeping the second argument fixed. The norm identities
and $\omega(\delta,\ell)=-bg$ give
\[
 -\omega(\beta,\ell)=bg/q_v,\qquad
 \omega(\eta,\delta)=b,\qquad
 \omega(\eta,\ell)=0.
\]
The integral catalogue has $q_v\mid b$.
These identities prove integrality of the first column
in \eqref{maps:eq:D.AF4}.

The fixed-plane projection has lattice
$P\Lambda=\mathbb Z(\delta/q_v)\oplus\mathbb Z(\ell/g)$.
For $d\nmid j$, the fixed-point equation of
$x\mapsto Ax+v/m$ is soluble precisely when
$m\mid jq_vr$ and $m\mid jgn$.
Choose a lattice lift of that projected displacement
and invert $A^j-I$ on its rotating plane to obtain
the fixed point. The positive obstruction indices
are the multiples of $m/t$.
At a power divisible by $d$, the translation lies
in $\Lambda$ precisely when $m$ divides both $jr$ and
$jn$, because $(\delta,\ell)$ is a primitive fixed
basis. The argument of Theorem~\ref{maps:src:D.2.29} therefore
gives exactly \eqref{maps:eq:D.AF3}.
It also gives $t\in\{1,q\}$.
The holomorphic translation conjugacy there applies
to this fixed basis without alteration.

Abelianization of the affine group now gives
\[
 H_1(F)=
 \langle\beta,\eta,c\mid mc-q_vr\beta-gn\eta=0\rangle
 =\mathbb Z^2\oplus\mathbb Z/t .
\]
The torsion generator is
$(m/t)c-(q_vr/t)\beta-(gn/t)\eta$.
The root character evaluates to $1/t$ on this
generator, so the normal Chern class has exact
order $t$. The integral duality calculation of
Theorem~\ref{maps:src:D.2.29} gives all the other cap groups.

Use the boundary coordinate map and its inverse in
that theorem. The normal tori of $\beta,\eta,c$
have fixed sweeps $(m/q_v)\delta,(m/g)\ell,v$
and vertical classes $-\beta\wedge v,-\eta\wedge v,0$.
Their evaluations by $\omega$ are exactly the three
columns \eqref{maps:eq:D.AF4}.
They generate the full Thom image, with primitive
relation
\[
 (m/t)K_c-(q_vr/t)K_\beta-(gn/t)K_\eta=0.
\]
Their projected minors are
$m^2/q,\ mn/q_v,\ -mr/g$.
Their gcd is
$(m/q)\gcd(m,gn,q_vr)=(m/q)t$.
The normal Chern class is torsion, so the next cap
product is zero and the degree-two boundary map
is onto, including the torsion of the central surface.
Its contracting numerator is $\gamma^*(v)=gn$.

We can therefore use the integral exterior-cokernel
map, its inverse, and the free bridge construction
of Theorem~\ref{maps:src:D.2.29} with these full kernels.
Over $\mathbb Q$ they give
$D_i=-(bg_in_i/m_i)s$ and
$L_i'=(bg_ir_i/m_i)s$.
The remaining scalar relation is
\[
 b\left(e-\sum_i\frac{g_in_i}{m_i}
          -\sum_j\frac{\gamma^*(v_j)}{m_j}\right)s=0.
\]
It is the rational Seifert relation.
The same duality and Euler calculation gives
\eqref{maps:eq:D.UF5} at both possible ranks.

Every finite position admits a free action at every
$m=dk$: use $v=\ell$ for $g=1$ and $v=\delta$
for $g=q>1$. Each choice has $t=1$ in
\eqref{maps:eq:D.AF3}.
An original additive cap is a minimal alternative
precisely when its line multidegree is zero and its
component action is transitive.
At types II, III and IV these are respectively the
conditions $q\mid a$ and $\gcd(b,q)=1$ in the component
coordinates of Theorem~\ref{maps:src:D.2.26}.
The interior vertices of types $\mathrm{III}^*$ and
$\mathrm{IV}^*$ give more than one orbit.
Taking the product of the permitted original/log
choices on the sixteen nonstar profiles gives
\[
 6,\ 4,\ 8,\ 8,\ 14,\ 12,\ 12,\ 10,\ 4,\ 6,\
 14,\ 12,\ 26,\ 20,\ 21,\ 21 ,
\]
whose sum is $198$.
Each pattern retains the profile, section pair,
zero and chosen finite and original sites.
Selecting patterns with $g=1$ at every finite site
gives the $164$ patterns of Theorem~\ref{maps:src:D.2.26};
the inclusion retains those marked coordinates.
The remaining $34$ patterns each have a finite site
with $g>1$.

For one finite site eliminate $D=-be\,s$, keeping
$s,L'$ as generators. Put $K=m/g$.
Replace the third column by the third minus
$e$ times the second. The inverse column operation
adds $e$ times the second, and the inverse generator
substitution restores $D=-be\,s$.
The resulting presentation is
\begin{equation}
 \begin{pmatrix}
 bgn'/q_v&-br&0\\
 0&K&n'
 \end{pmatrix}.
 \hypertarget{tag-01XQ}{}\label{maps:eq:D.AF9}
\end{equation}
The lift change $c'=c\eta^{-e}$, with inverse
$c=c'\eta^e$, changes the power by
$-eN_m\eta=-eK\ell$. It gives the same affine torus
action with parameter $n'$. In particular
$t=\gcd(m,q_vr,gn')$ and its freeness is preserved.

For $n'\ne0$ the matrix has rank two.
Since $\gcd(K,n',r)=1$, choose
$U_0r+V_0K+W_0n'=1$.
Multiplication by $b$ proves
$\gcd(br,K,n')=\gcd(b,K,n')=h$ in both directions.
The other matrix entry is an integer multiple of
$n'$, so $h$ is the first determinantal divisor.
The three two-by-two minors have gcd
\[
 \frac b{q_v}|n'|\gcd(m,gn',q_vr)
 =\frac{bg|n'|t}{q}.
\]
Thus the two Smith factors are $h$ and
$bg|n'|t/(qh)$.
Retaining the unimodular row matrix $U$ in the
Euclidean algorithm gives the quotient map
$x\mapsto Ux$ modulo these two factors; its inverse
is $y\mapsto U^{-1}y$ modulo the original image.
This specifies the group isomorphism integrally.

For $n'=0$, the exact free-action condition gives
$\gcd(m,r)=1$. The only nonzero relation is
$-brs+KL'=0$, of content $\gcd(b,K)$.
A B\'ezout completion of its primitive column gives
the quotient $\mathbb Z\oplus\mathbb Z/\gcd(b,K)$
and its inverse. The free degree-one bridge adds
one further $\mathbb Z$ to $H_2(X)$.

In the one-site Seifert group, eliminate
$c=\gamma^e$. The remaining relation is
$\gamma^{me-gn}=1$, with exponent $-gn'$.
The generator map is the marked loop $\gamma$,
and its inverse is the same Seifert generator.
It gives $\mathbb Z/(g|n'|)$ for $n'\ne0$ and
$\mathbb Z$ for $n'=0$.
Inserting these groups and the two Smith calculations
in \eqref{maps:eq:D.UF5} proves
\eqref{maps:eq:D.AF6} and \eqref{maps:eq:D.AF7}.
Finally, $g_i=q>1$ divides both $m_i$ and its
Seifert numerator $g_i n_i$.
The necessary simple-connectivity condition for
the full Seifert presentation, proved in
Theorem~\ref{maps:src:D.2.23}, excludes this common divisor
at any site. This proves the asserted fundamental
group obstruction for every number of sites.
\end{proof}

\begin{remark}[Finite lattice and one-site calculations]
\hypertarget{tag-Q125}{}\label{mapsIII:rem:general-finite-calculations}
{\normalfont\scriptsize tag Q125}\par
At the $34$ additional positions with $g=q>1$, use the integral
fixed and coinvariant bases in \eqref{maps:eq:D.AF2}. The
proper-power obstruction gives \eqref{maps:eq:D.AF3}, and the
normal tori of $\beta,\eta,c$ give the three columns in
\eqref{maps:eq:D.AF4}. For one finite site, eliminate $D=-be\,s$
and put $n'=n-(m/g)e$. The determinantal divisors of
\eqref{maps:eq:D.AF9} give
\eqref{maps:eq:D.AF6}--\eqref{maps:eq:D.AF7}.
Theorem~\ref{maps:src:D.2.30} constructs these integral maps for
every allowed positive order, every free integer twist and every
integer clutch.
\end{remark}

\subsection{Original additive caps with zero component Chern class}

\begin{definition}[Zero-Chern original additive data]
\hypertarget{tag-Q126}{}\label{mapsIII:def:zero-chern-original-data}
{\normalfont\scriptsize tag Q126}\par
Use the natural initial rows of Theorem~\ref{maps:src:D.2.22},
with $P=aG$, $Q=bG$, $a,b>0$, $\gcd(a,b)=1$, and the periods of
Definition~\ref{mapsIII:def:elliptic-marked-data}. Write
$\omega=a e_1^*\wedge e_2^*-b\delta^*\wedge\gamma^*$.
Retain any original additive cap at which $P$ is narrow, with
any number of component orbits. At $I_n^*$ use the original cap;
at the finite good-reduction additive sites allow all choices and
free affine twists of Theorem~\ref{maps:src:D.2.30}. Use the
marked zero-free multiplicative and unique rank-two zero models
of Theorem~\ref{maps:src:D.2.23}. Allow any number of regular logarithmic sites,
every permitted positive order and an integer clutch $e$.
A filling pattern is the ordered list of options at the additive
sites of one natural initial row.
For an original additive cap let $\rho_i$ be the number of
orbits of $Q$ on its components, and put
\[
 D=\sum_{i\text{ original additive}}(\rho_i-1).
\]
\end{definition}

\begin{theorem}[Zero-Chern original additive caps and component summands]
\hypertarget{tag-02L7}{}\label{maps:src:D.2.31}
{\normalfont\scriptsize tag 02L7}\par
The data of Definition~\ref{mapsIII:def:zero-chern-original-data}
give $228$ additive-filling patterns on $160$ natural initial rows.
Exactly $30$ patterns have a positive original component defect.
The only nonzero local defects in this catalogue are
\begin{equation}
\begin{array}{c|rrrr}
\text{type}&III^*&IV^*&I_3^*&I_1^*\\\hline
\rho_i&5&3&3&2\\
\rho_i-1&4&2&2&1.
\end{array}                                                   \hypertarget{tag-0212}{}\label{maps:eq:D.OD1}
\end{equation}
Every other permitted original cap has $\rho_i=1$.

Let $\mathbb Z^r\oplus {\mathcal A}_1$ be the abelianization of the Seifert
group in Theorem~\ref{maps:src:D.2.23}, treating an unchanged cap
as an order-one meridian with zero numerator in the canonical
marking. Let $T$ be the torsion subgroup of the cokernel of the matrix in
Theorem~\ref{maps:src:D.2.30},
using its finite and regular columns and setting both fixed
sweeps of every original additive cap to zero. Then the full
integral groups are
\begin{equation}
\begin{array}{c|c}
j&H_j(X;\mathbb Z)\\\hline
0,6&\mathbb Z\\
1&\mathbb Z^r\oplus {\mathcal A}_1\\
2&\mathbb Z^{2r+D}\oplus T\\
3&\mathbb Z^{2r+2D}\oplus T\\
4&\mathbb Z^{2r+D}\oplus {\mathcal A}_1\\
5&\mathbb Z^r.
\end{array}                                                   \hypertarget{tag-0213}{}\label{maps:eq:D.OD2}
\end{equation}
The fundamental group is the presentation in
Theorem~\ref{maps:src:D.2.23}, for every permitted order and affine
twist. These homology identifications split the component
summands integrally.

There is a full smooth identification in the $b=1$ sector.
For every simply connected member, let $M$ be the product of
its regular logarithmic multiplicities. Then
\begin{equation}
 X\cong
 \mathop{\#}^{D}(S^2\times S^4)
 \ \#\ \mathop{\#}^{D}(S^3\times S^3)
 \ \#\ \widetilde\Sigma_0 N_M,                             \hypertarget{tag-0214}{}\label{maps:eq:D.OD3}
\end{equation}
with the oriented twisted suspension and the spin five-manifold
$N_M$ defined in Definition~\ref{mapsIII:def:twisted-smale-suspension}. The allowed free defects in this sector
are $D=0,1,2,4$. The identification preserves the oriented framing and the refined
spin data for every $M$, including even $M$.

\end{theorem}


\begin{proof}
\begin{enumerate}[leftmargin=*]
\item \emph{Markings.}
Use Definition~\ref{mapsIII:def:zero-chern-original-data} and put
\(\chi=\gamma^*\) and
\(J=\left(\begin{smallmatrix}0&1\\-1&0\end{smallmatrix}\right)\).
The component point map \(+Q\), forward arrow \(-Q\), and pullback
on line coefficients use the corresponding maps below.

\item \emph{The natural rows and additive choices.}

The $172$ natural rows and their exhaustive divisor/component
criterion are those of Theorem~\ref{maps:src:D.2.22}. The row
identification fixes the ordered profile, selected zero, positive
section pair $(a,b)$ and every component coordinate. It therefore
fixes every marked monodromy
and cocycle in Definition~\ref{mapsIII:def:elliptic-marked-data}.
At each finite additive place its options are
the free finite cap, and additionally the original cap exactly when
\(P\) is narrow. At each odd star there is only the original option;
it is available here exactly when \(P\) is narrow. The Cartesian
product of these site options is the additive-pattern set. The empty
product contributes one pattern.

The resulting complete profile count is as follows. The first count
is the number of natural rows, the second the number admitting these
choices, the third the number of patterns, and the last the number
of patterns with positive original component defect.

\begin{center}
\small
\begin{tabular}{llllll}
\hline
OS & Profile & Natural & Allowed rows & Patterns & Positive \\
\hline
43 & base & 6 & 6 & 9 & 3 \\
43 & II & 2 & 2 & 6 & 2 \\
45 & base & 8 & 8 & 8 & 0 \\
45 & II & 4 & 4 & 8 & 0 \\
46 & base & 6 & 3 & 3 & 3 \\
46 & II & 2 & 1 & 2 & 2 \\
47 & base & 14 & 14 & 14 & 0 \\
47 & II & 6 & 6 & 12 & 0 \\
47 & III & 8 & 8 & 12 & 0 \\
49 & base & 10 & 10 & 15 & 5 \\
49 & II & 2 & 2 & 6 & 2 \\
49 & III & 4 & 4 & 9 & 3 \\
50 & base & 10 & 5 & 5 & 5 \\
50 & II & 2 & 1 & 2 & 2 \\
50 & IV & 4 & 2 & 3 & 3 \\
55 & base & 14 & 14 & 14 & 0 \\
55 & II & 6 & 6 & 12 & 0 \\
56 & base & 26 & 26 & 26 & 0 \\
56 & II & 10 & 10 & 20 & 0 \\
56 & III & 14 & 14 & 21 & 0 \\
56 & IV & 14 & 14 & 21 & 0 \\
\textbf{Total} &  & \textbf{172} & \textbf{160} & \textbf{228} & \textbf{30} \\
\hline
\end{tabular}
\end{center}

The displayed option rule determines all 228 filling patterns,
with the row and site markings fixed. Parts 6--7 include every
affine parameter and every positive allowed order within each list.

\item \emph{The original neighbourhood, section and line.}

Let \(N\to\Delta\) be the minimal elliptic-surface germ at
an original additive place. At a star its rational components have
a tree incidence graph; in types III and IV their underlying spheres
are joined at one common point, and type II is a cuspidal rational
curve. Its component classes
give
\[
 H_1(N)=0,\qquad H_2(N)=\bigoplus_j\mathbb Z[C_j],
 \qquad H_k(N)=0\quad(k>2).
\]
Use the regular-neighbourhood retractions in the original-cap
construction of Theorem~\ref{maps:src:D.2.16}. On a triangulation
adapted to the central fibre, the retraction reduces the barycentric
coordinates off the central subcomplex and renormalizes their sum;
its inclusion is the inverse up to this explicit homotopy. The
star component graph is a tree, so its edge contractions leave the
stated sphere classes. In III and IV the common-point description
gives the same wedge of component spheres. The cusp normalization is a homeomorphism of the
underlying compact rational curve. These maps also identify the
degree coordinates of every line bundle on \(N\).

Put \(z=(A-I)^{-1}g\). In the marking of Theorem~\ref{maps:src:D.2.23} the cocycle is
minus the Abel--Jacobi increment. Thus the invariant Betti point
\(bz\) extends to the torsion section \(T_0\) with the same
component as the forward section \(-Q\). Put \(Q_0=-T_0\).
Both sections lie on this same germ \(N\). The
component order is \(q=1,2,3,2,3,4,4\) for
\(II,III,IV,III^*,IV^*,I_1^*,I_3^*\), respectively. The actual
section \(Q\) is the sum of this torsion section \(Q_0\) and a section
in the identity component. By the local exponential exact sequence,
the latter is \(\operatorname{Exp}L_{\rm log}\) for a holomorphic section of
the inverse Hodge line. The maps
\[
 t_{-Q_s},\quad Q_s=Q_0+\operatorname{Exp}(sL_{\rm log}),\quad0\le s\le1,
 \qquad\text{inverse }t_{Q_s},
\]
are defined on this neighbourhood, including its components.
Translation extends over the minimal surface, and the displayed
isotopy carries the mapping-torus roof to the torsion roof.

Here \(P\) is narrow. Hence \({\mathcal M}=\mathcal O(P-O)\) has zero degree
on every component, and the neighbourhood homology gives \(c_1({\mathcal M})=0\) integrally. A smooth
nowhere-zero frame is constructed from transition logarithms and a
partition of unity; multiplication by the resulting exponential and
its reciprocal are the bundle maps. On the punctured regular fibre
the induced angular change is
\[
 S_P(x,t_\delta,t_\gamma)
 =(x,t_\delta+\alpha x,t_\gamma),\qquad
 \alpha=-a z^tJ,
\]
and the inverse subtracts \(\alpha x\). The row is integral because
\(P\) is narrow. The period exponential implements the line-frame shear: replace the line frame, multiply its
nonzero unit, and use the same group logarithm on the marked
four-torus.

The quotient of \({\mathcal M}^\times\) by the strict contraction is a mapping
torus of the unit-circle bundle over \(N\). Write its arrow in a
frame as \((y,w)\mapsto(f(y),u_{\rm line}(y)w)\), with \(|u_{\rm line}(y)|<1\), and
put \(u=-\log|w|\), \(r_{\rm roof}(y)=-\log|u_{\rm line}(f^{-1}y)|>0\). Its norm
fundamental domain is \(0\le u\le r_{\rm roof}(y)\), with
\((y,r_{\rm roof}(y))\sim(f^{-1}y,0)\). Choose \(\epsilon>0\) with
\(\epsilon\max r_{\rm roof}<1\), and smooth nonnegative \(b_y\) of integral
one supported in the middle third of \([0,r_{\rm roof}(y)]\). Define
\[
 F_y(u)=\int_0^u
       [\epsilon+(1-\epsilon r_{\rm roof}(y))b_y(v)]\,dv .
\]
It is strictly increasing, equals \(\epsilon u\) at the lower
end and \(1+\epsilon(u-r_{\rm roof}(y))\) at the upper end. Consequently
\([y,u]\mapsto[y,F_y(u)]\) respects the gluing with all
derivatives, and its inverse uses the unique smooth \(F_y^{-1}\).
The same formula applies to the unit-circle arrow with its
retained angular multiplier.

Fix the forward-quotient convention
\(T^+(f)=(N\times\mathbb R)/((x,t)\sim(f(x),t+1))\);
increasing \(t\) is the marked positive contracting direction.
For \(f_s=t_{-Q_s}\), put \(\phi_s=f_0f_s^{-1}\).
Then \(\phi_1f_1=f_0\), so the map
\([x,t]\mapsto[\phi_t x,t]\) identifies \(T^+(f_1)\) with
\(T^+(f_0)\), with inverse \([y,t]\mapsto[\phi_t^{-1}y,t]\).
Extend it by \(\phi_{t+1}=f_0\phi_tf_1^{-1}\), and choose a
collar cutoff constant at each interval end. This verifies both
the quotient identification and its smooth endpoint derivatives.
The same construction applies to the lifted angular arrows.
Use the angular parts of all line multipliers. Because \(H^1(N)=0\), their \(S^1\)-parts
have real logarithms, and the cutoff of that logarithm removes
their residual gauge. Reciprocal exponentials are the inverse
gauges. A possible integral choice of the suspension logarithm is
transported by \((t_\delta,t_\gamma)\mapsto
(t_\delta+k t_\gamma,t_\gamma)\), with inverse \(-k\); it fixes
the fibre maps and the full fixed-sweep kernel used below. Choose
the branch matching the given positive contracting period.

Equivalently, the angular circle bundle over the torsion mapping
torus has zero Euler class: the Wang restriction in degree two is
injective since \(H^1(N)=0\), and its restriction to \(N\) is \({\mathcal M}\).
The preceding logarithms give a trivialization with its actual
boundary marking. We have therefore constructed
\[
 U\cong S^1_\delta\times B,\qquad
 B=T^+(t_{-Q_0}:N\to N),
\]
with the inverse maps just described. Collars are carried along
by the same cutoff isotopies. The deck component arrow is the
inverse of the \(+Q\) permutation. If that permutation is \(P_Q\),
then \(P_Q^{-1}-I=-P_Q^{-1}(P_Q-I)\), so the orbit quotient
identification is the identity on component coefficients with
this explicit inverse change of relation generators. Its
invariant component lattice is also unchanged. The positive
suspension direction itself has not been reversed.

The local exponential and extension of invariant torsion leaves are
\cite[Section 2.4, equation (2.2), pp.~6--7;
Proposition 3.5(1),(3), pp.~8--10]{manUU21}. Translation extends over the minimal surface;
see \cite[Sections 7.6--7.7, pp.~33--34]{manSS10}.

\item \emph{The complete integral original-cap map.}

Let \(U_1,\ldots,U_\rho\) be the component-orbit classes in
\(H_2(B)\). The Wang sequence for the product description gives
\begin{equation}
 \begin{aligned}
 H_1(U)&=\mathbb Z\langle\delta,\gamma_B\rangle,\\
 H_2(U)&=\mathbb Z^\rho\langle U_j\rangle
                    \oplus\mathbb Z\langle D_B\rangle,
 \qquad D_B=\delta\times\gamma_B,\\
 H_3(U)&=\mathbb Z^{2\rho},\qquad H_4(U)=\mathbb Z^\rho.
 \end{aligned}
 \hypertarget{tag-0215}{}\label{maps:eq:D.OD4}
\end{equation}
The maps
are the component inclusions, component orbit tubes, and their
products with \(\delta\); inverse coordinates are the Wang
projection and component coefficients. These coordinate maps prove that all the groups are torsion-free.

The ordinary elliptic fibre specializes to
\[
 F=\sum_j f_jU_j,
 \qquad f_j=\sum_{C_k\text{ in orbit }j}\operatorname{mult}(C_k).
\]
In this catalogue \(Q_0\) generates the component group whenever
\(q>1\), because \(q\mid a\) and \(\gcd(a,b)=1\). Hence all
\(f_j\) are divisible by \(q\), and
\[
                         W=F/q
\]
is primitive. More precisely, the orbit of the outer component
met by \(O\) has length \(q\), multiplicity one, and coefficient
one in \(W\). The component group permutations and primitive
vectors, in the indicated orbit order, are

\begin{center}
\small
\begin{tabular}{lp{0.47\textwidth}cl}
\hline
Type & Component action & \(\rho\) & \(W\) \\
\hline
II & one component & 1 & \((1)\) \\
III & exchange two components & 1 & \((1)\) \\
IV & cycle three components & 1 & \((1)\) \\
\(III^*\) & exchange the two long arms; fix centre and short arm & 5 & \((1,2,3,2,1)\) \\
\(IV^*\) & cycle the three arms; fix centre & 3 & \((1,2,1)\) \\
\(I_3^*\) & cycle the four outer components; reverse the four central ones & 3 & \((1,1,1)\) \\
\(I_1^*\) & cycle the four outer components; exchange the central two & 2 & \((1,1)\) \\
\hline
\end{tabular}
\end{center}

Replacing the component permutation by its inverse preserves each
orbit and its coefficient map, through the relation-generator
identity in part 3. Replacing the first orbit basis vector
by \(W=(1,w_2,\ldots,w_\rho)\) has matrix determinant one.
Its inverse sends \((x_1,x_2,\ldots)\) to
\((x_1,x_2-w_2x_1,\ldots)\). Thus
\(\mathbb Z^\rho/\mathbb ZW=\mathbb Z^{\rho-1}\) with explicit
component representatives and no torsion extension.

These permutations are induced by section translation.
On the multiplicity-one components the component group acts by its
regular translation action. Incidence then determines the images
of all remaining components: the \(III^*\) outer exchange forces the
two long-arm exchanges, the \(IV^*\) outer cycle forces its three
middle-arm cycle, and the odd-star outer four-cycle exchanges
the two ends and reverses the central chain. Applying the inverse
outer permutation gives the inverse on the entire graph. Thus
the displayed orbit computation is transported from the section
action on the original fibre, with its multiplicities preserved.

For the full boundary marking put \(z_Q=bz\). In the coordinates
after \(S_P\), the monodromy is the direct product of the angular
circle and
\[
 T_3=\begin{pmatrix}A&bg\\0&1\end{pmatrix}
 \quad\text{on }(e_1,e_2,\gamma).
\]
Indeed \(\alpha(A-I)+\beta=0\), where
\(\beta=-a g^tJA\); it also cancels the central entry
\(\beta z_Q\). The image of \(T_3-I\) is the full elliptic
plane: the class of \(g\) generates the cyclic group
\(\mathbb Z^2/(A-I)\mathbb Z^2\), and \(b\) is a unit modulo
its order \(q\). Consequently the boundary fibre Wang quotient
is torsion-free of rank two. Its full specialization is, in wedge
order \((12,13,14,23,24,34)\),
\[
 f_2(x)=q(x_{12}+z_{Q,2}x_{14}-z_{Q,1}x_{24})W+x_{34}D_B.
\]
In the original marking, evaluate \(f_2\) on
\(\bigwedge^2S_P x\). The first row is the unique invariant area
row in the three nonangular coordinates, with its scale fixed by
\(e_{12}\mapsto F=qW\); the second is the product-circle map
in the product description. Since the component target in \eqref{maps:eq:D.OD4} is torsion-free, these
geometric values determine the integral map, including all mixed
terms and their signs.

The coefficients \(q,qz_{Q,2},-qz_{Q,1}\) have gcd one.
Choose \(A_0,B_0,C_0\) with
\(qA_0+qz_{Q,2}B_0-qz_{Q,1}C_0=1\). Then
\[
 W\longmapsto\bigwedge^2S_P^{-1}
              (A_0e_{12}+B_0e_{14}+C_0e_{24}),\qquad
 D_B\longmapsto\bigwedge^2S_P^{-1}e_{34}
\]
is an integral section. Substitution gives both inverse composites. The entries of the
marked matrices give four unit Smith factors for the boundary
wedge relation, including at each odd star.

For the marked boundary monodromy $\mathsf T$, the fixed lattice is
\[
 \operatorname{Fix}(\mathsf T)=\mathbb Z\delta\oplus\mathbb Z\ell,
 \qquad\ell=(-qz_{Q,1},-qz_{Q,2},0,q).
\]
Both full normal sweeps are zero in \(H_2(U)\). For \(\delta\)
use its line-circle tube over the original base disc in an extending
line frame. For \(\ell\), use the \(q\)-fold orbit tube of \(Q_0\),
close its identity-component part by the translation isotopy, and extend across the same
base disc. The extension has the marked increment \(-qz_Q\), the
positive return count \(q\), and the extending line-frame unit.
The logarithm/cutoff construction in part 3 carries these actual
three-chains back to the given section and contraction. Their inverse
maps reverse the isotopy and gauges.

The boundary Wang exact sequence has fibre quotient of rank two
and fixed quotient \(\mathbb Z\delta\oplus\mathbb Z\ell\).
The displayed section makes the fibre map onto
\(\mathbb ZW\oplus\mathbb ZD_B\) an isomorphism. Hence the full
boundary kernel is exactly its two fixed sweeps, and the full
cokernel is the free component complement of rank \(\rho-1\).
Thus the primitive boundary image and its component complement satisfy
\begin{equation}
 \begin{gathered}
 \operatorname{im}(H_2(\partial U)\to H_2(U))
     =\mathbb ZW\oplus\mathbb ZD_B,\\
 H_2(U)/\operatorname{im}H_2(\partial U)=\mathbb Z^{\rho-1}.
 \end{gathered}
 \hypertarget{tag-0216}{}\label{maps:eq:D.OD5}
\end{equation}


The component-intersection description gives the same primitive
image. Let \(Q_K\) be the Kodaira intersection matrix, let
\(I\) send each orbit to the sum of its components, and let
\(R_0\) choose one representative of each orbit. Intersecting
with the orbit tubes identifies the component part of
\(H_2(U)\to H_2(U,\partial U)=H^4(U)\) with
\(I^tQ_KR_0\). Its columns are independent of the chosen
representatives because section translation preserves \(Q_K\).
In the orbit orders above the nonzero folded matrices are
\[
 \begin{array}{c|c}
 III^*&\begin{pmatrix}
 -2&1&0&0&0\\1&-2&1&0&0\\0&1&-2&2&0\\
 0&0&1&-2&1\\0&0&0&1&-2
 \end{pmatrix}\\[3pt]
 IV^*&\begin{pmatrix}-2&1&0\\1&-2&3\\0&1&-2\end{pmatrix}\\[3pt]
 I_3^*&\begin{pmatrix}-2&2&0\\1&-2&1\\0&1&-1\end{pmatrix}\\[3pt]
 I_1^*&\begin{pmatrix}-2&2\\1&-1\end{pmatrix}.
 \end{array}
\]
They kill the corresponding vector \(W\). Their leading
\((\rho-1)\)-minors are \(2,3,3,-2\), so their rank is
\(\rho-1\); the II, III and IV matrices are zero of size one.
The additional class \(D_B\) is in the boundary image and maps
to zero in the relative group. This also proves the rank of
the entire intersection map. Its image is free because the
relative target is free, in agreement with
\eqref{maps:eq:D.OD5}.

At an unselected multiplicative site, the component equations
for \(P\) narrow give \(\gcd(b\kappa_i,\nu_i)=1\) on each
natural row. Thus \(Q\) acts transitively. The pair
\((u_i,\delta)\) is primitive because \(u_i\) is primitive in
the elliptic plane; a B\'ezout complement to \(u_i\), followed
by \(\gamma\), extends it to an integral basis. Consequently
its original zero-Chern alternative has the same primitive
one-orbit filling maps in Theorem~\ref{maps:src:D.2.26}.
The natural factor-pair enumeration gives 316 such multiplicative
occurrences and 80 additive occurrences, hence 396 unselected
zero-Chern original occurrences in total.

\Needspace{5\baselineskip}
\item \emph{The integral exterior on all natural rows.}

At every marked place use
\[
 T_i=\begin{pmatrix}A_i&0&bg_i\\
            -ag_i^tJA_i&1&-ag_i^tJA_i(bz_i)+w_i\\0&0&1
       \end{pmatrix}.
\]
At a multiplicative place let \(r_{i,\rm per}\) be the marked integral
complement to \(u_i\) in the input period model, with determinant one
in their ordered elliptic basis. Put \(z_i=(\kappa_i/\nu_i)r_{i,\rm per}\).
The weight is \(w_i=1/\nu_i\) at the selected zero and zero
elsewhere. Substitution of each marked profile gives integral matrices with
ordered product identity, and each matrix preserves \(\omega\).

The map \([x]\mapsto[\Pi(t)x]\), for the actual real period matrix
\(\Pi(t)\), identifies the smooth exterior with its marked torus
bundle. The inverse is the group logarithm modulo \(\Lambda\).
Every nonconstant period and nonzero unit remains in this map.
Write \(N_{\partial}\) for the number of marked boundary circles.
Cutting the base along the marked tree gives
\[
 A_\mathrm{gr}=[T_1-I\ \cdots\ T_{{N_{\partial}}-1}-I].
\]
For \(i<{N_{\partial}}\), let \(j_i\) be the block inclusion; put
\[
 j_{N_{\partial}}(x)=-(T_2\cdots T_{{N_{\partial}}-1}x,\ T_3\cdots T_{{N_{\partial}}-1}x,\ldots,x).
\]
Telescoping gives
\(A_\mathrm{gr}j_{N_{\partial}}(x)=(I-T_{N_{\partial}}^{-1})x\). The inverse cut maps
rejoin the same transported strips.

Integer row and column operations give fixed bases
\(F_i\), with first column \(\delta\), and integral left inverses.
For \({\mathsf B}=[j_iF_i]\), these operations give
\[
 \operatorname{Smith}(A_\mathrm{gr})=(1,1,1),\quad
 \operatorname{Smith}({\mathsf B})=(1^{\,4{N_{\partial}}-7}),\quad
 \ker {\mathsf B}=\mathbb Z\rho_\delta,
\]
where \(\rho_\delta\) is one in each first fixed coordinate and
zero elsewhere. There are \(4{N_{\partial}}-6\) columns. Thus
\(\operatorname{im}{\mathsf B}=\ker A_\mathrm{gr}\) integrally. The unimodular matrices
\(U_{\mathsf B},V_{\mathsf B}\) satisfy \(U_{\mathsf B}{\mathsf B}V_{\mathsf B}=\operatorname{diag}(I,0)\);
apply \(U_{\mathsf B}\), retain its first \(4{N_{\partial}}-7\) coordinates, append a
zero, and apply \(V_{\mathsf B}\) to invert \({\mathsf B}\) modulo \(\rho_\delta\).
The displayed inverse composites follow from these Smith identities.

The matrix made from the columns \(\bigwedge^2T_i-I\) has exactly
five unit Smith factors. The fibre quotient and its inverse are
\[
 [x]\longmapsto\omega(x),\qquad
 k\longmapsto[k s_0],\quad
 s_0=u e_{12}+v e_{34},\quad au-bv=1.
\]
The same Smith operations give integral monodromy-cylinder chains
for \(x-\omega(x)s_0\). In particular the quotient is the full
primitive group \(\mathbb Zs\), where \(s=[s_0]\).

For a fixed loop \(u\), use the constant marked loop \(t\mapsto tu\)
over the positive peripheral meridian to define its actual sweep
\(t_i(u)\). Addition is realized by the real-period triangle with
vertices \(0,u,u+v\). The physical \(\delta\)-circle over the cut
base is a three-chain with boundary \(\sum_i t_i(\delta)\).
These chains, the cut-graph inverse and the primitive fibre quotient
give the integral splitting
\[
 H_2(E)=\mathbb Zs\oplus
 \left(\bigoplus_i\operatorname{Fix}(T_i)\right)
                 /\langle\sum_i t_i(\delta)\rangle .
\]
Its inverse projects to the graph chain, uses the inverse of \({\mathsf B}\),
subtracts the resulting actual sweeps, and reads the remaining fibre
coefficient with the primitive fibre quotient. These forward and inverse maps establish the integral splitting.

Each added regular disc contributes a zero block to \(A_\mathrm{gr}\),
a full \(\Lambda\)-sweep block, and the corresponding negative block
to the cut-graph map. Deleting/restoring these coordinates extends the same inverse
to every number of regular sites. Noncanonical peripheral lifts add
\(be\,s\) to the one global relation, with \(e\) as in the complete
Seifert relation. For each boundary, the marked matrix \(T_i-I\) likewise has
only unit nonzero Smith factors. Hence all boundary
\(H_1\) groups are free, including the stars.

\item \emph{All finite and regular affine data, with the norm convention.}

At a finite site, let \(d_i\) be the monodromy order, \(q_i\) the
component-group order, and \(g_{i,\rm den}\) the denominator index of
\(z_{Q,i}\). The fixed and coinvariant bases from Theorem~\ref{maps:src:D.2.30} satisfy
\[
 \ell_i=(-g_{i,\rm den} z_{Q,i,1},-g_{i,\rm den} z_{Q,i,2},0,g_{i,\rm den}),\quad
 q_{v,i}=q_i/g_{i,\rm den},
\]
\[
 [\delta]=q_{v,i}\beta_i,\quad[\ell_i]=g_{i,\rm den}\eta_i,
 \quad N_{m_i}\beta_i=(m_i/q_{v,i})\delta,
 \quad N_{m_i}\eta_i=(m_i/g_{i,\rm den})\ell_i.
\]
These bases have integral quotient sections and inverse coordinates;
their marked pairings are
\(-\omega(\beta_i,\ell_i)=b g_{i,\rm den}/q_{v,i}\) and
\(-\omega(\eta_i,\delta)=-b\).
Take \(m_i=d_i k_i>0\) and \(v_i=r_{i,\rm tw}\delta+n_i\ell_i\).
Every affine lift is conjugated to \(x\mapsto T_ix+v_i/m_i\)
by translation on the rotating plane; the inverse subtracts the
same translation. The complete freeness test is
\[
 \gcd(m_i,r_{i,\rm tw},n_i)=1,\quad
 d_i\mid m_i/\tau_i,\quad
 \tau_i=\gcd(m_i,q_{v,i}r_{i,\rm tw},g_{i,\rm den} n_i).
\]
The finite normal character has exact order \(\tau_i\). Its torsion
Chern class and the core torsion are retained by the full Thom map.

The boundary parametrization is
\([y,\tau]\mapsto[y+(\tau/m_i)v_i,\exp(2\pi i\tau/m_i)]\).
Choosing a root argument and subtracting its translation is the inverse.
The normal circle is \((y,\tau)\mapsto(y-sv_i,\tau+m_is)\).
The Thom torus of a core lattice loop \(\lambda\) has fibre part
\(-\lambda\wedge v_i\) and sweep \(N_{m_i}\lambda\); cutting
its \(m_i\) sectors is the inverse chain construction. The root
generator contributes zero fibre area and sweep \(v_i\).
Consequently the full kernel columns, on
\((s,D_i=t_i(\delta),L_i=t_i(\ell_i))\), are
\[
 \begin{pmatrix}b g_{i,\rm den} n_i/q_{v,i}\\m_i/q_{v,i}\\0\end{pmatrix},
 \quad\begin{pmatrix}-b r_{i,\rm tw}\\0\\m_i/g_{i,\rm den}\end{pmatrix},
 \quad\begin{pmatrix}0\\r_{i,\rm tw}\\n_i\end{pmatrix}.
\]
The primitive relation among them is
\((-q_{v,i}r_{i,\rm tw},-g_{i,\rm den} n_i,m_i)/\tau_i\).
The next Thom map is cap product by a torsion Chern class into
\(H_0\), hence zero. Thus boundary \(H_2\) is onto the core,
and the full kernel index in the finite Thom columns includes all core torsion.

At a regular site take arbitrary \(m_j>0\) and
\(v_j=x_je_1+y_je_2+z_j\delta+\mu_j\gamma\) with
\(\gcd(m_j,x_j,y_j,z_j,\mu_j)=1\). The core torus has lattice
\(\Lambda+\mathbb Z(v_j/m_j)\); B\'ezout completion of \((-v_j,m_j)\)
gives its integral lattice map and inverse. The same normal-circle
parametrization gives the full kernel
\(m_jt_j(u)-\omega(u,v_j)s\), \(t_j(v_j)\), for all \(u\in\Lambda\).

A peripheral replacement \(c_i\mapsto c_i\lambda_i\)
at a finite site changes the power by
\(v_i\mapsto v_i+N_{m_i}\lambda_i\). For an invariant
\(\lambda_i\), this is \(v_i\mapsto v_i+m_i\lambda_i\). It changes
\(e\mapsto e+\sum\chi(\lambda_i)\) and induces
\[
 t'_i(u)\longmapsto t_i(u)+\omega(u,\lambda_i)s,
 \qquad s\longmapsto s.
\]
Subtracting the shear and using \(-\lambda_i\) gives its inverse.
For finite order, \(N_m\) is self-adjoint for \(\omega\), because
the sums of all powers \({\mathsf T}^k\) and \({\mathsf T}^{-k}\) agree. Thus
\(\omega(N_mx,y)=\omega(x,N_my)\) and
\(\omega(N_m\lambda,\lambda)=0\). These identities show that
the sweep shear preserves every Thom column. The root-loop relation transforms
using \(t(N_m\lambda)=\omega(\lambda,v)s\); its two mixed
terms cancel and its last term is zero. The global relation is
preserved since \(\omega(\delta,\lambda)=-b\chi(\lambda)\).
At a regular site \(N_m=mI\), giving the usual \(m\lambda\)
formula. This accounts for every affine lift with its exact inverse.

\item \emph{The full integral global calculation.}

Use the primitive rank-one and index-one cusp maps of
Theorems~\ref{maps:src:D.2.14} and~\ref{maps:src:D.2.26}.
Their fibre specializations are integral surjections, with the
displayed sections there, and their full boundary kernels are
their fixed sweeps. The same statement for a zero-Chern
original, with a free component complement, is proved in part 4.

In degree-two Mayer--Vietoris, split each original cap into its
primitive boundary image and the component complement constructed
from \(W\) in part 4. Lift all image classes to their boundaries using the displayed section of \(f_2\), and
lift finite/regular core classes by their Thom surjections. Subtract
these lifts from the exterior coordinate. The resulting map and
its inverse identify the full cokernel with
\(C\oplus\mathbb Z^D\), where
\(D=\sum_{\mathrm{original}}(\rho_i-1)\), and \(C\) has generators
\(s\), both finite sweeps, and all regular sweeps. Its relations
are the finite Thom columns, the regular kernels of part 6, and
\[
       \sum_{\mathrm{finite}}D_i+
       \sum_{\mathrm{regular}}t_j(\delta)+be\,s=0.
\]
Both full sweeps at each original additive are zero. Different
boundary lifts differ by exactly the kernels already imposed, so
these are inverse isomorphisms of integral groups.

The complete fundamental group, with the actual marked loops and
inverse disc maps of Theorem~\ref{maps:src:D.2.23}, is
\[
 \langle\gamma,c_i\mid [\gamma,c_i]=1,
 c_i^{m_i}=\gamma^{\mu_i},\ \prod_i c_i=\gamma^e\rangle,
 \quad\mu_i=\chi(v_i).
\]
An unchanged cap is \(m_i=1,\mu_i=0\) in the canonical marking.
Its extra component classes add no loops: the original components
and their attached discs lie in the same neighbourhood used by
that van Kampen proof. Write its abelianization as
\(\mathbb Z^r\oplus {\mathcal A}_1\). Rational elimination gives \(r=1\)
exactly when \(e-\sum_i\mu_i/m_i=0\), and \(r=0\) otherwise.

The finite relations in \(C\otimes\mathbb Q\) give
\(D_i=-b\mu_i s/m_i\); the regular relations give
\(t_j(u)=\omega(u,v_j)s/m_j\). Every other sweep is thereby
determined, and the remaining equation is
\(b(e-\sum_i\mu_i/m_i)s=0\). Therefore \(C\) has free rank
\(r\); put \(T=\operatorname{Tor}C\).

The next Mayer--Vietoris kernel is a subgroup of the direct sum
of the free boundary \(H_1\) groups. It is free. With \(c_i^{\rm fix}\)
the fixed-lattice rank, boundary ranks are \(c_i^{\rm fix}+1\), cap ranks
are \(c_i^{\rm fix}\), and exterior rank is the number of boundary circles.
The degree-zero map is injective. Rank exactness therefore gives
this kernel rank \(r\), including zero rational Euler number and
finite cores with torsion in \(H_1\).
Choose an integral basis of that kernel, represent it by boundary
loops, and bound its two images by integral two-chains on the
two sides. Glue those chains to form cycles \(Z_1,\ldots,Z_r\).
Then \((c,u,k)\mapsto c+u+\sum k_jZ_j\) splits the sequence.
Its inverse first reads the Mayer--Vietoris boundary, subtracts
these cycles, and applies the cap-elimination map above. Thus
\(H_2(X)=C\oplus\mathbb Z^D\oplus\mathbb Z^r\) integrally.

Every original cap is a mapping torus and has Euler characteristic
zero; the regular, finite and primitive rank-one caps have Euler
characteristic zero as well. The index-one cusp has Euler
characteristic two by its normalization/conductor maps. Thus
\(\chi(X)=2\). Integral duality and universal coefficients give
\eqref{maps:eq:D.OD2}.
For a torsion Smith generator of order \(d\), choose its character
with value \(1/d\in\mathbb Q/\mathbb Z\), apply the integral
Bockstein and cap with the oriented fundamental class; the inverse
reads its torsion-linking coordinate. Dual integral cocycles capped
with that class give the free coordinates. This specifies the
group identifications in the remaining degrees. In particular
the free component defect cannot conceal a torsion extension.

\item \emph{The simply connected \(b=1\) matrix and the cusp class.}

Assume now that the resulting member is simply connected and
\(b=1\). Triviality of the Seifert group forces
\(\gcd(m_i,\mu_i)=1\) at every place: otherwise the corresponding
column of its square abelian presentation vanishes modulo a prime.
After order-one elimination, more than two nontrivial orders give
a nontrivial sphere-orbifold quotient, by the polygon construction in
Theorem~\ref{maps:src:D.2.23}. With at most two orders, the
determinant is a unit and makes all multiplicities pairwise
coprime. This statement retains every sign of \(\mu_i,e\).

For \(b=1\), every finite additive place with component order
\(q>1\) has \(g_{\rm den}=q>1\), so \(q\mid\mu_i,m_i\). It therefore
cannot occur as a finite cap in this simply connected sector.
Only type II can be finite; there \(q=g_{\rm den}=q_v=1\).
The finite Thom columns reduce, with B\'ezout inverse using
\(A_{i,\rm Bez} n_i+B_{i,\rm Bez} m_i=1\), to
\(m_iD_i+n_i s=0\). This follows directly by solving for \(L_i\)
from \(n_iL_i=-r_{i,\rm tw}D_i\) and \(m_iL_i=r_{i,\rm tw}s\).

For the finite sites put
\({F_f}=\prod_i m_i\), \(N_f=\sum_i({F_f}/m_i)n_i\), and
\(D_f=\sum_iD_i\). Integral CRT, with inverse
\[
 D_i=E_i({F_f}/m_i)D_f+
 \left(k_i n_i+E_i\sum_{\ell\ne i}
                    \frac{{F_f}}{m_i m_\ell}n_\ell\right)s,
 \quad E_i({F_f}/m_i)=1+k_i m_i,
\]
turns these relations into \({F_f} D_f+N_fs=0\). At \(m_i=1\),
take \(E_i=0,k_i=-1\). Multiplication by \(m_i\), followed by
summation, verifies both inverse composites; all cross coefficients
are integral by pairwise coprimality.

For the regular sites put
\(R=\prod_jm_j\), \({v_{\rm reg}}=\sum_j(R/m_j)v_j=({x_{\rm reg}},{y_{\rm reg}},{z_{\rm reg}},{n_{\rm reg}})\), and
\({t_{\rm reg}}(u)=\sum_jt_j(u)\). Their inverse CRT map is
\[
\begin{gathered}
 t_j(u)=E_j(R/m_j){t_{\rm reg}}(u) -
 \left(k_j\omega(u,v_j)+E_j\sum_{\ell\ne j}
       \frac{R}{m_jm_\ell}\omega(u,v_\ell)\right)s,\\
 E_j(R/m_j)=1+k_jm_j.
\end{gathered}
\]
It identifies the relations with \(R{t_{\rm reg}}(u)=\omega(u,{v_{\rm reg}})s\),
\({t_{\rm reg}}({v_{\rm reg}})=0\). To verify the last assertion, write
\(G_j=(R/m_j){t_{\rm reg}}(v_j)-\sum_{\ell\ne j}
R\omega(v_j,v_\ell)s/(m_jm_\ell)\). The first relations
give \(m_jG_j=0\), the regular CRT substitution gives \(t_j(v_j)=E_jG_j\), and
alternation gives \(\sum_jG_j={t_{\rm reg}}({v_{\rm reg}})\). CRT idempotents recover
each \(G_j\) from the sum. These are the exact inverse maps,
including all order-one sites.

Let \(K={F_f}e-N_f\). Eliminating \(D_f\) by the global relation
gives
\[
 {F_f} {t_{\rm reg}}(\delta)+K s=0,
 \qquad R {t_{\rm reg}}(\delta)+{n_{\rm reg}} s=0,
 \qquad\Delta={F_f}{n_{\rm reg}}-RK=\pm1.
\]
Choose \(uR+v{F_f}=1\). The matrix
\(\left(\begin{smallmatrix}u&v\\{F_f}&-R\end{smallmatrix}\right)\)
has determinant \(-1\). Applying it to the second and first relations in that order,
and then its integral inverse, gives precisely
\[
                            s=0,\qquad {t_{\rm reg}}(\delta)=0.
\]
Modulo \(R\), the two Seifert relations give \(\gcd({n_{\rm reg}},R)=1\). Choose \(p{n_{\rm reg}}+qR=1\).
Then \({t_{\rm reg}}(\gamma)=-p({x_{\rm reg}}{t_{\rm reg}}(e_1)+{y_{\rm reg}}{t_{\rm reg}}(e_2))\), with this same expression
as inverse substitution. The remaining two generators have relations
\(R{t_{\rm reg}}(e_1)=R{t_{\rm reg}}(e_2)=0\). Thus
\[
                H_2(X)=\mathbb Z^D\oplus(\mathbb Z/R)^2.
\]
Empty finite or regular sets are handled by \({F_f}=1,N_f=D_f=0\)
or \(R=1,{v_{\rm reg}}={t_{\rm reg}}=0\); the auxiliary generators are then eliminated
by the displayed formulas. The equalities \(s={t_{\rm reg}}(\delta)=0\) and the resulting
second-homology splitting therefore include empty sets of sites.

In particular, \(s=0\) is an equality of the exterior class in
\(H_2(X)\), proved before any smooth comparison. The direct
component splitting of part 7 leaves this equality unchanged.
The integral cusp specialization has a section, so every actual
cusp curve is the image of an integral fibre class and hence a
multiple of \(s\). It is therefore zero in \(H_2(X)\) even when
\(D>0\). Since \(X\) is simply connected, these parametrized
curves are nullhomotopic.

\Needspace{5\baselineskip}
\item \emph{The global comparison of cusp fans.}

The selected cusp has the Jordan fan of
Theorem~\ref{maps:src:D.2.23}. At this zero write \(r_{\rm per}\)
for that same marked complement. In physical short
and long bases \(S=(u,\delta)\), \(L=(r_{\rm per},\gamma)\), put
\[
 {\mathsf C}=\begin{pmatrix}\nu&b\kappa\\
             -a\kappa&(1-ab\kappa^2)/\nu\end{pmatrix}
 \in SL_2(\mathbb Z),\qquad ({\mathsf T}-I)L=S{\mathsf C} .
\]
Its fan edges are \({\mathsf C}_1,{\mathsf C}_2,{\mathsf C}_2-{\mathsf C}_1\). The physical angular
circle has Jordan coordinates \({\mathsf C}^{-1}\delta=(-b\kappa,\nu)\);
it need not be an initial edge. Apply the fan comparison in Theorem~\ref{maps:src:D.2.26}
with the following period and collar maps.


For example, the OS47 base row with the \(I_2\) zero and
\((a,b)=(7,1)\) has \(\kappa=1\) and
\[
 {\mathsf C}=\begin{pmatrix}2&1\\-7&-3\end{pmatrix},\qquad
 {\mathsf C}^{-1}\delta=(-1,2).
\]
Along the first Jordan edge, the first circle collapses and
the scalar action on the remaining angular coordinate is
\(z\mapsto z^2\); its stabilizer is \(\{1,-1\}\).
We therefore make the following global fan comparison before
using a semifree scalar action.

Euclidean reduction of the first row of \({\mathsf C}\), followed by a lower
shear, gives a finite word of unit column additions and quarter
rotations taking its Farey triangle to the one with edge \(\delta\).
At a basis \(E=(p,q)\), the period change is
\[
 H_E=(SE_1,L{\mathsf C}^{-1}E_1,SE_2,L{\mathsf C}^{-1}E_2),
 \qquad H_E^{-1}{\mathsf T}H_E=I+E_{12}+E_{34}.
\]
The identity on the complex universal cover with these period
coordinates, and \(H_E^{-1}\), are the forward and inverse maps.
They retain all nonzero unit factors. Each unit flip orbit is an
embedded complete rational curve with transition
\((x',y',w')=(wx,wy,w^{-1})\), inverse
\((w'x',w'y',(w')^{-1})\); hence its normal bundle is
\(\mathcal O(-1)\oplus\mathcal O(-1)\).

The identity on the punctured rank-one matrix cone, in charts
\((u,p,q')=(y/x,x,xw)\) and inverse
\((x,y,w)=(p,up,q'/p)\), is the flop comparison. It
preserves the base function \(xyw\), the torus characters and all
unit multipliers. The codimension-four Thom groups in degrees
two and three vanish, so the complements identify the two
integral \(H_2\) groups and the marked fibre specialization.
This isomorphism has the reverse complement map as inverse.
It preserves \(s=0\) and simple connectivity at each step; hence
every successive flip curve is nullhomotopic.

The relative-ball construction in Theorem~\ref{maps:src:D.2.26}
applies at each step. Its
quaternion contraction trivializes the real normal clutch;
Barden's parametrized isotopy puts the nullhomotopic curve in
a six-ball. The opposite resolution under the actual boundary
map gives a simply connected acyclic replacement. Capping by
the original complementary disc is a homotopy six-sphere;
\(\Theta_6=0\) and the h-cobordism theorem give its standard
oriented sphere. The specified parametrized-disc isotopy adjusts
the resulting map to identity on the entire outer collar, using
a flow with backward flow as inverse. Thus it is an oriented
diffeomorphism of the replacement relative that collar. Applying
the finite word and reversing it gives an oriented global
diffeomorphism, with the reverse-word inverse, between the completed
threefold using the given cusp fan and the completed threefold using the
\(\delta\)-edge fan.
The relative-ball choices can preserve every original component
section sphere. Choose the component representatives at distinct
suspension angles in their original caps, so that both the
sections and their local quotient spheres are disjoint.
Before each step remove these representatives
from the six-manifold. Their real codimension is four, so the
relative Thom groups in degrees two and three vanish; general
position also preserves simple connectivity. Hence the current
flip curve is still nullhomologous and nullhomotopic in that
complement. Apply the parametrized ball-placement isotopy there.
The resulting ball and the relative-collar diffeomorphism avoid all
these spheres. Repeating the construction fixes the chosen component-section maps
throughout.

\item \emph{The angular quotient, its integral homology and spin.}

Use this oriented global diffeomorphism for the remainder of the
proof, preserving the component sections. On the resulting
completed model the physical \(\delta\)-circle acts
semifreely with one fixed rational curve. It is free on each
original cap, since a nonzero contracting iterate changes the
line norm while an angular rotation does not. At a zero-free
multiplicative cap with \(b=1\), the component equations give \(P\) narrow and \(Q\) transitive. The primitive
Tate vanishing vector has primitive elliptic projection; together
with \(\delta\) it spans a primitive sublattice, giving freeness
on its ruled charts and glued sections. At finite and regular
logs, equality of a rotation with the \(j\)-th deck power would
force \(m_i\mid j\mu_i\); the already proved coprimality forces
\(m_i\mid j\), and then the primitive line circle makes the
rotation identity.

The final cusp fan has one orbit of \(\delta\)-edges, giving
the fixed \({\mathcal C}\simeq S^2\). Its normal weights are \(-1,+1\),
and every other residual edge weight is a unit. Conjugate the
negative-weight coordinate, reverse the core orientation, and
use the \(SU(2)\) clutch contraction in
Theorem~\ref{maps:src:D.2.14}. This gives the
scalar action on \({\mathcal C}\times D^4\), preserving total orientation;
the conjugate-transpose clutch and reverse normal exponential
are the inverse. The Hopf orbit coordinates are
\[
 h(z_1,z_2)=(2\Re z_1\bar z_2,2\Im z_1\bar z_2,
                          |z_1|^2-|z_2|^2),\quad
 \mathcal H(z)=h(z)/|z|,\quad\mathcal H(0)=0 .
\]
The two Hopf sections give their inverse orbit charts.

Let \(\mathcal Q\) be the resulting closed oriented smooth quotient,
\(\mathcal Q_0=\mathcal Q\setminus({\mathcal C}_{\mathcal Q}\times\mathring D^3)\), and \(X_0\)
the free-circle part. Orient the normal meridian sphere
\(\sigma\) so that its Euler evaluation is one. Removing the
codimension-four fixed curve preserves \(\pi_1,\pi_2,H_2\).
The circle fibration and Hurewicz give
\[
 \pi_1(\mathcal Q_0)=0,\qquad
 0\longrightarrow H_2(X)\longrightarrow H_2(\mathcal Q_0)
 \xrightarrow{\langle e_\delta,-\rangle}\mathbb Z
 \longrightarrow0,\quad k\longmapsto k\sigma .
\]
Filling the normal tube kills exactly \(\sigma\), so
\[
        \pi_1(\mathcal Q)=0,\qquad H_2(\mathcal Q)=\mathbb Z^D\oplus(\mathbb Z/R)^2.
\]
The map from \(H_2(X)\) is circle projection and inclusion.
For its inverse, lift a class to \(H_2(\mathcal Q_0)\), subtract its
Euler value times \(\sigma\), lift the zero-Euler sphere to
\(X_0\), and include it in \(X\). Changing the first lift by
\(\sigma\) cancels. The difference of two sections over a sphere
is a map \(S^2\to S^1\), hence homotopic to a constant. This
proves the inverse on integral torsion as well as free classes.

The free basis of \(H_2(X)\) consists of the component-section
spheres of part 4. On each such original neighbourhood the
canonical line has the invariant nonvanishing form
\(\eta_N\wedge dw/w\). Indeed \(K_N\) is pulled back from the disc by
\cite[Theorem 6.8, p.~26]{manSS10}, translation preserves \(\eta_N\), and under any line map
\(w'=u(x)w\), the additional term
\(\eta_N\wedge d\log u\) is zero in complex dimension two.
The contraction has the same cancellation. Hence \(c_1(TX)\)
evaluates to zero on every free basis sphere. Since \(\pi_1(X)=0\),
universal coefficients identify \(H^2(X;\mathbb Z)\) with
\(\operatorname{Hom}(H_2(X),\mathbb Z)\); therefore \(c_1=0\)
integrally and \(w_2(X)=0\). The cusp comparison preserves this
spin structure.

On \(X_0\), \(TX_0=p^*T\mathcal Q_0\oplus\mathbb R\). The mod-two
Gysin kernel is generated by \(e_\delta\bmod2\), whose value on
\(\sigma\) is one. The value of \(w_2(\mathcal Q_0)\) on that sphere is
zero, since its tangent restriction is \(TS^2\oplus\mathbb R^3\).
Thus \(w_2(\mathcal Q_0)=0\), and excision gives \(w_2(\mathcal Q)=0\).
The residue sequence sends the value
\(\operatorname{res}(e_\delta)=1\) to \(\operatorname{PD}[{\mathcal C}_{\mathcal Q}]\).
Exactness gives \( [{\mathcal C}_{\mathcal Q}]=0\) integrally.

\item \emph{Localizing the fixed sphere while fixing every free section.}

The component complement constructed from \(W\) is represented by central
component spheres with a chosen line section. They map under the split circle sequence
to \(D\) embedded quotient spheres \(V_1,\ldots,V_D\), giving
a primitive free basis of \(H_2(\mathcal Q)\), with Euler value zero. Choose
distinct suspension angles for representatives from one original
cap and use a small general-position perturbation when necessary;
two two-dimensional spheres in a five-manifold can thereby be
made disjoint. The bundle sections extend along these isotopies.
They remain disjoint from \({\mathcal C}_{\mathcal Q}\), which lies in the separate cusp.

Set \(\mathcal Q'=\mathcal Q\setminus\bigcup_jV_j\). General position of loops
and filling two-discs gives \(\pi_1(\mathcal Q')=0\). Excision identifies
\(H_3(\mathcal Q,\mathcal Q')=\mathbb Z^D\), and its preceding map from \(H_3(\mathcal Q)\)
is intersection with the primitive free basis \([V_j]\). Integral
duality makes this map onto. Since \(H_2(\mathcal Q,\mathcal Q')=0\),
\[
                       H_2(\mathcal Q')\xrightarrow{\cong}H_2(\mathcal Q).
\]
This includes \(D=0\), when no spheres are removed. Thus \({\mathcal C}_{\mathcal Q}\)
is nullhomologous and, by Hurewicz, nullhomotopic already in \(\mathcal Q'\).

Barden's Theorem B(b) \cite[pp.~367--368]{manBarden65}, for the connected source \(S^2\) and
the arbitrary ambient \(\mathcal Q'\), applies: relative connectivity below
degree two follows from connectedness and simple connectivity,
and \(2\cdot5>3\cdot2+3\). It gives a parametrized isotopy of
this sphere into a small five-disc in \(\mathcal Q'\). Extend the
velocity to a compactly supported vector field; its flow and
backward flow give the ambient maps. They fix every \(V_j\)
and transport the normal framing and circle bundle.

On \(\mathcal Q\setminus\mathring D^5\), the Euler class is now zero:
its \(H^2\) is detected by the preserved free basis \(V_j\),
and each sphere still has a section. Transition logarithms and
a partition of unity give a trivialization of the whole exterior
bundle, with reciprocal gauges for its inverse. On the common
\(S^4\) collar the difference of two gauges has a real logarithm;
a cutoff exponential and its reciprocal make them agree on the
entire collar. The normal framing has one homotopy class because
\(\pi_2(SO(3))=0\). This constructs the global bundle
trivialization in the presence of free \(H_2\).

\item \emph{Exact framed reconstruction and recognition.}

Apply the precisely framed reconstruction in the proof of
Theorem~\ref{maps:src:D.2.16}. Its quotient, component-basis count
and torsion parameter are here \(\mathcal Q,D,R\), respectively;
the count called \(b\) in that reconstruction is \(D\), while
the section coefficient in the present theorem is \(b=1\).
The preceding parts have supplied its scalar normal product,
primitive free section spheres, zero exterior Euler class,
simply connected spin quotient and complete integral linking data.
The geometric maps are the following ones with those actual inputs.

The Hopf charts and the exterior trivialization give
\[
 M_-(\mathcal Q)=((\mathcal Q\setminus\mathring D^5)\times S^1)
                 \cup_g(S^4\times D^2),\qquad
 g(x,z)=(\alpha_zx,z),
\]
where \(\alpha_z\) is the rotation in the first real two-plane
of Definition~\ref{mapsIII:def:twisted-smale-suspension}. The map
\(g\) goes from complement to attached boundary and its
inverse uses \(\alpha_z^{-1}\). The cap action is
\(t(x,w)=(\alpha_tx,tw)\), with boundary orbit coordinate
\(\alpha_z^{-1}x\); this is the same quotient-disc coordinate.
The rotation lifts to \(\cos(\pi t)+\sin(\pi t)e_1e_2\),
with endpoints \(1,-1\), so the framing loop is nontrivial.
The matrix and collar map from that reconstruction satisfy
\[
 {\mathsf A}=\operatorname{diag}(1,-1,-1,1,1),\qquad
 {\mathsf A}\alpha_z{\mathsf A}^{-1}=\alpha_z^{-1},\qquad
 F_{\mathcal Q}|_{\partial}={\mathsf A}.
\]
The maps \(F_{\mathcal Q}\times\mathrm{id}\) and
\({\mathsf A}\times\mathrm{id}\) therefore glue to the oriented map
\(X\to\widetilde\Sigma_0\mathcal Q\); their inverse rotations
glue by the same boundary identity.

For the linking realization, use the component basis \(V_j\)
and, for \(R>1\), choose \(x\) of order \(R\) and \(y\) with
\(\lambda(x,y)=1/R\). Write
\(\operatorname{num}_R([k/R])=k\bmod R\). The explicit mutually inverse coordinates are
\[
 z\longmapsto(\operatorname{num}_R\lambda(z,y),
                  -\operatorname{num}_R\lambda(z,x)),\qquad
 (r,s)\longmapsto rx+sy.
\]
At \(R=1\) use the unique maps of the zero group. The alternating
linking form and \(w_2=0\) meet exactly the oriented realization
conditions of \cite[Lemma F, p.~372; Theorems 2.2--2.3,
pp.~376--377]{manBarden65}. This gives
\(\mathcal Q\to\#^D(S^2\times S^3)\#N_R\) inducing those
coordinates. The parametrized-disc adjustment and product/cap
maps in the cited reconstruction lift it to the suspensions;
the inverse uses the inverse base diffeomorphism on the same disc.

Use its same rotational dumbbell, radial map and inverse to
separate the connected-sum factors. On the separating hypersurface
the maps are \((x,z)\mapsto(\alpha_z^{-1}x,z)\) and identity
on \(S^3\times D^2\); their inverse uses \(\alpha_z\).
Cutting and rejoining the same collars, with the stated transverse
reflection, give the oriented connected sum. This construction
uses one such neck for each of the \(D\) component factors and
uses no neck when \(D=0\).
For each product factor, the trivial-clutch maps of
\cite[Theorem B(2), p.~3; Proposition 5.6, pp.~22--23]{manGR25},
with base \(S^3\) and fibre \(S^2\), give
\((S^2\times S^4)\#(S^3\times S^3)\); inverse clutch maps
and the product swaps whose dimension products are even preserve the orientation.
These maps prove \eqref{maps:eq:D.OD3}, with \(R=M\).
At \(M=1\), the map \((x,z)\mapsto(\alpha_zx,z)\) on
\(D^5\times S^1\), identity on \(S^4\times D^2\), and radial
corner rounding identify the last factor with the oriented
\(S^6\). The inverse uses the inverse rotation.

The enumeration gives precisely \(D=0,1,2,4\) when \(b=1\).
The displayed maps are oriented diffeomorphisms with their inverses;
they preserve the integral group maps and spin structures. Thus
\eqref{maps:eq:D.OD3} includes the full refined data for even \(M\) as well.

\end{enumerate}
\end{proof}

\begin{remark}[Finite zero-Chern original-cap calculations]
\hypertarget{tag-Q127}{}\label{mapsIII:rem:zero-chern-original-calculations}
{\normalfont\scriptsize tag Q127}\par
The natural factor-pair and site-option enumeration gives $396$
unselected zero-Chern original occurrences and $228$ additive
patterns. The component orbits and their multiplicities determine
the primitive fibre vector $W$; the folded intersection matrix
$I^tQ_KR_0$ in part 4 gives the component rank. At the odd stars,
the comparison sends each labelled orbit coefficient to the same
coefficient in Theorems~\ref{maps:src:D.2.16} and
\ref{maps:src:D.2.17}. The integral reduction in part 8 and the
framed reconstruction give the primitive scalar-circle data for
the simply connected $b=1$ sector. Every permitted narrow-$P$
section pair is included in the option rule above.
\end{remark}

\subsection{All original additive caps and their integral homology}

\begin{definition}[General original additive data]
\hypertarget{tag-Q128}{}\label{mapsIII:def:all-original-additive-data}
{\normalfont\scriptsize tag Q128}\par
Use the $172$ natural initial rows and common marking of
Definition~\ref{mapsIII:def:zero-chern-original-data}. At $I_n^*$
choose the original cap; at every finite good-reduction additive
site choose the original cap or any free logarithmic cap. Use the
marked Mumford models at multiplicative sites and the zero-Chern
original alternatives of Theorem~\ref{maps:src:D.2.31}.
All permitted positive orders, free affine twists, regular
logarithmic sites and integer clutches remain parameters.
For a nonzero-Chern original additive cap, let $s_i$ be its
number of Kodaira components and $q_i^c$ its component-group order.
Let $L$ be the number of such nonzero-Chern original caps,
and put
\begin{equation}
 \beta=
 \sum_{\text{zero-Chern originals}}(\rho_i-1)
 +\sum_{\text{nonzero-Chern originals}}(s_i-2).
                                                               \hypertarget{tag-020W}{}\label{maps:eq:D.OC2}
\end{equation}
\end{definition}

\begin{theorem}[All original additive choices on the natural catalogue]
\hypertarget{tag-02L8}{}\label{maps:src:D.2.32}
{\normalfont\scriptsize tag 02L8}\par
The data of Definition~\ref{mapsIII:def:all-original-additive-data}
give exactly $282$ additive-filling patterns on all $172$ initial
rows. Of them, $54$ involve a nonzero component Chern class.
At each nonzero-Chern original additive cap, the full boundary
kernel adds the fibre relation
\begin{equation}
 h_i s=0,\qquad h_i=b/q_i^c,                               \hypertarget{tag-020V}{}\label{maps:eq:D.OC1}
\end{equation}
in addition to its two full fixed-loop normal relations.
Here $s$ is the common exterior fibre generator.
The local cokernel is free of rank $s_i-2$.

Let $T$ be the torsion subgroup of the cokernel of the matrix in
Theorem~\ref{maps:src:D.2.30}, adding
\eqref{maps:eq:D.OC1} for every nonzero-Chern original and setting both
fixed sweeps to zero at every unchanged additive cap.
Let $\mathbb Z^r\oplus A_{1,\rm tor}$ be the abelianization of the
Seifert group in Theorem~\ref{maps:src:D.2.23}, with its finite
and regular numerators.
When $L=0$ the homology is given by Theorem~\ref{maps:src:D.2.31}. When $L>0$ it is
\begin{equation}
\begin{array}{c|c}
j&H_j(X;\mathbb Z)\\\hline
0,6&\mathbb Z\\
1&\mathbb Z^r\oplus A_{1,\rm tor}\\
2&\mathbb Z^{\beta+L+r}\oplus T\\
3&\mathbb Z^{2\beta+2L}\oplus T\\
4&\mathbb Z^{\beta+L+r}\oplus A_{1,\rm tor}\\
5&\mathbb Z^r.
\end{array}                                                   \hypertarget{tag-020X}{}\label{maps:eq:D.OC3}
\end{equation}
All positive logarithmic orders and all free affine twists
are included, with arbitrary regular logarithms. The full
fundamental group is the presentation in Theorem~\ref{maps:src:D.2.23}.
The degree-one Mayer--Vietoris kernel has rank $L+r$, including
the case of zero rational Seifert Euler number.

\end{theorem}


\begin{proof}
\begin{enumerate}[leftmargin=*]
\item \emph{Markings and the zero-Chern cases.}
Use the marking and area form of
Definition~\ref{mapsIII:def:all-original-additive-data}, and put
\(J=\left(\begin{smallmatrix}0&1\\-1&0\end{smallmatrix}\right)\).
Theorem~\ref{maps:src:D.2.31} gives the zero-Chern original maps,
the primitive exterior fibre class \(s\), and the complete
finite and regular Thom maps. We compute the remaining original
maps and their integral assembly.

\item \emph{The nonzero-Chern original neighbourhood and its roof.}

Let \(N\) be the actual original additive neighbourhood, with
central components \(C_0,\ldots,C_{s_i-1}\). The neighbourhood retractions in
Theorem~\ref{maps:src:D.2.31}, part 3, give
\[
 \pi_1(N)=0,\qquad H_2(N)=\mathbb Z^{s_i}\langle C_j\rangle,
 \qquad H_k(N)=0\quad(k>2).
\]
The component equation on a natural row says that \(P\) or \(Q\)
is narrow. At a nonzero-Chern original \(P\) is not narrow,
so \(Q\) is narrow. Put \(q=q_i^c\), the component-group order.
For every additive component group in these profiles the marked
generator has order \(q\). Thus
\[
 q\mid b,\qquad \gcd(a,q)=1,\qquad
 c_1(M)=C_p^*-C_0^*\ne0,\quad M=\mathcal O(P-O),
\]
where \(C_0,C_p\) have multiplicity one and \(p\ne0\).
The last formula follows from the actual section divisors. It is
a primitive integral cohomology class, with value \(+1\) on \(C_p\)
and \(-1\) on \(C_0\).

Let \(E(M)\to N\) be the unit-circle bundle in any smooth Hermitian
metric. The actual strict contraction quotient is the mapping torus
of its angular arrow with the positive variable logarithmic roof.
Use the smooth roof map \(F_y\) of Theorem~\ref{maps:src:D.2.31}, part 3, whose two end
formulas have the same constant positive derivative; its unique
smooth inverse is \(F_y^{-1}\). This respects the variable-roof
gluing with all derivatives. Since \(Q\) is narrow, write
\(Q=\operatorname{Exp}L_{\rm log}\) on this same germ. Translation by
\(-\operatorname{Exp}(tL_{\rm log})\) is an isotopy from identity to the
actual base translation. Lift its vector field horizontally using
a connection on \(E(M)\). The flow is a bundle map covering this
isotopy, and its backward flow is the inverse. At the endpoint its
difference from the actual angular arrow is a gauge
\(u:N\to S^1\). Since \(N\) is simply connected, \(u=\exp(if)\)
for a smooth real \(f\). The gauges \(\exp(itf)\), with their
reciprocals, supply the remaining endpoint isotopy.

Thus the actual angular arrow is isotopic to identity on the
actual \(E(M)\). Write its isotopy as \(f_t\), from identity to
the actual angular arrow \(f_1\). Use the forward quotient
\(T^+(f_1)=(E(M)\times\mathbb R)/((x,t)\sim(f_1(x),t+1))\)
with its forward convention. Put \(\phi_t=f_t^{-1}\). Then
\(\phi_1f_1=I=\phi_0\), so \([x,t]\mapsto[\phi_t x,t]\)
is the mapping-torus identification with the product, and
\([y,t]\mapsto[f_t y,t]\) is its inverse. Extend by the
endpoint identity on consecutive intervals. Collar cutoffs
make the maps constant at the gluing ends.
This constructs
\[
                         U\cong E(M)\times S^1_\gamma.
\]
The line bundle keeps the Chern class \(C_p^*-C_0^*\). An integral branch change of the
suspension logarithm acts by
\((t_\delta,t_\gamma)\mapsto(t_\delta+k_{\rm br} t_\gamma,t_\gamma)\),
with inverse \(-k_{\rm br}\); the complete fixed-sweep kernel below is
unchanged by this displayed map.

The homotopy sequence of the circle bundle, using the surjective
evaluation of this Chern vector on \(\pi_2(N)=H_2(N)\), gives \(\pi_1(E(M))=0\).
The integral Gysin sequence gives
\[
\begin{aligned}
 H_1(E(M))&=0,\qquad
 H_2(E(M))=\ker(c:\mathbb Z^{s_i}\to\mathbb Z),\\
 H_3(E(M))&=\mathbb Z^{s_i},\qquad H_4(E(M))=0,
\end{aligned}
\]
where \(c(x)=x_p-x_0\). The maps are projection to the component
cycles, fibre integration, and the same section lifts on
zero-degree components. The product description now gives
\begin{equation}
 \pi_1(U)=\mathbb Z\gamma,\quad H_1(U)=\mathbb Z\gamma,\quad
 H_2(U)=\ker c,\quad
 H_3(U)=\mathbb Z^{2s_i-1},\quad H_4(U)=\mathbb Z^{s_i}.
                                                               \hypertarget{tag-020Y}{}\label{maps:eq:D.OC4}
\end{equation}
These maps identify every group integrally and prove that it is
torsion-free.

\item \emph{The primitive fibre image and component complement.}

Let \(F=(m_0,\ldots,m_{s_i-1})\) be the actual whole-fibre
multiplicity vector, so \(m_0=m_p=1\). The ordinary elliptic
area maps to \(F\in\ker c\). This map follows by projection from
the line-circle bundle, and the injectivity in the Gysin sequence
of the circle-bundle homology identifies its unique lift. The remaining component
classes \(C_j\), \(j\ne0,p\), have degree zero and therefore
have section lifts to \(E(M)\).

The vectors \(F,C_j\ (j\ne0,p)\) give an integral basis of
\(\ker c\). The forward map is their linear combination. The
inverse is
\[
 x\longmapsto
       \left(x_0,\ (x_j-m_jx_0)_{j\ne0,p}\right).
\]
The identity \(x_p=x_0\) verifies both inverse composites.
In particular \(F\) is primitive and
\[
                    \ker c/\mathbb ZF=\mathbb Z^{s_i-2}.
\]
The zero in the type-III value \(s_i-2=0\) is fully retained.

For the complete boundary map put
\[
 z=(A-I)^{-1}g,\qquad
 \beta_{\rm row}=-a g^tJA,\qquad z_Q=bz\in\mathbb Z^2.
\]
The actual four-period matrix is
\[
 {\mathsf T}=\begin{pmatrix}A&0&bg\\
                  \beta_{\rm row}&1&\beta_{\rm row} z_Q\\0&0&1\end{pmatrix}.
\]
The marked fixed basis is
\[
 \delta,\qquad\ell=(-z_{Q,1},-z_{Q,2},0,1).
\]
Let \(R\) be the integral matrix with columns \(e_1,e_2,\delta,\ell\).
Its inverse adds \(z_Q\) times the fourth coordinate to the first
two; both maps have determinant one. Direct substitution gives
\[
 R^{-1}{\mathsf T}R=\operatorname{diag}({\mathsf T}_3,1),\qquad
 {\mathsf T}_3=\begin{pmatrix}A&0\\\beta_{\rm row}&1\end{pmatrix}.
\]
This is a period-coordinate map on the actual group torus,
implemented by its period exponential and inverse logarithm.

The primitive coinvariant row on the first three coordinates is
\[
                \rho_{\rm cov}=q(a z_2,-a z_1,1),\qquad \rho_{\rm cov}\delta=q.
\]
It is integral, \(\rho_{\rm cov}({\mathsf T}_3-I)=0\), and its entries have gcd one.
The last two statements follow from
\((-az^tJ)(A-I)+\beta_{\rm row}=0\) and the generator component order
\(q\), together with \(\gcd(a,q)=1\). More concretely integer row and column operations give two unit Smith entries for
\({\mathsf T}_3-I\), and the row \(\rho_{\rm cov}\) is the resulting primitive quotient row.
Thus \(\operatorname{coker}({\mathsf T}_3-I)=\mathbb Z\), with its
generator specified by the row \(\rho_{\rm cov}\), not just up to rational scale.

Choose an integral \(v_0=(x,y,z_\delta)\) with \(\rho_{\rm cov} v_0=1\)
by B\'ezout operations, and put
\[
 \beta_0=R(v_0,0),\qquad k=\beta_0\wedge\ell .
\]
Changing this B\'ezout choice changes \(\beta_0\) by
\(\operatorname{im}({\mathsf T}_3-I)\); its wedge with the fixed \(\ell\)
therefore changes \(k\) by a boundary Wang relation. Thus its
class is independent of the choice. The fibre Wang quotient
has the explicit coordinates
\[
\begin{gathered}
 f_{\mathrm{coinv}}(x)=
 \begin{pmatrix}
 1&0&0&0&0&0\\
 0&0&\rho_{{\rm cov},1}&0&\rho_{{\rm cov},2}&\rho_{{\rm cov},3}
 \end{pmatrix}
 \bigwedge\nolimits^2R^{-1}x
 \\[2pt]
 \text{in wedge order }(12,13,14,23,24,34).
\end{gathered}
\]
Its inverse sends the standard basis to \(e_{12}\) and \(k\).
Substitution shows that the fibre quotient map annihilates
\(\bigwedge^2{\mathsf T}-I\) and has this integral right inverse. The
latter difference has four unit Smith factors, so the fibre quotient map gives
its whole quotient \(\mathbb Z^2\).

To identify the second coordinate geometrically, forget the angular
line and project to \(B=T^+(t_{-Q}:N)\). The same local logarithm
gives \(B\cong N\times S^1\). In the induced three-period basis
\((e_1,e_2,\ell)\), its boundary matrix is
\(\operatorname{diag}(A,1)\). The mixed classes
\(e_1\wedge\ell,e_2\wedge\ell\) form the finite coinvariant
group \(\operatorname{coker}(A-I)=\mathbb Z/q\); their images in
the torsion-free \(H_2(B)=H_2(N)\) are therefore zero. The
projection of \(k\) is such a mixed class. Moreover
\(H_2(U)\to H_2(B)\) is the injection \(\ker c\hookrightarrow
H_2(N)\) given by the circle-bundle projection in part 2. Thus \(k\) has zero image in \(H_2(U)\).
This identifies its image by the forgetful map and the marked
period change.

Consequently the full original-cap fibre map takes the first
coordinate of the fibre quotient map to \(F\) and the second to zero. Projection
to the Gysin target in the circle-bundle homology takes \(e_{12}\) to \(F\), and its
rank-one fibre image has the invariant area row
\begin{equation}
 w=(1,0,z_{Q,2},0,-z_{Q,1},0),\qquad
 w(x)=x_{12}+z_{Q,2}x_{14}-z_{Q,1}x_{24}.
 \hypertarget{tag-020Z}{}\label{maps:eq:D.OC5}
\end{equation}
The remaining coordinates are the free component classes with
the inverse constructed above.
This determines the entire map integrally. Its fibre kernel is
the primitive \(\mathbb Z k\), and its fibre image is exactly
\(\mathbb ZF\), with section \(F\mapsto e_{12}\).

Finally, its literal exterior pairing is
\[
 \begin{aligned}
 \omega(k)
 &=a\det((x,y),-bz)-b z_\delta\\
 &=-b(az_2x-az_1y+z_\delta)
 =-b/q.
 \end{aligned}
\]
The denominator disappears because \(q\mid b\). Thus the
boundary fibre kernel adds the relation \(-(b/q)s=0\).
Changing its generator from \(k\) to \(-k\), with the same
sign reversal as inverse, writes it as
\[
                           (b/q)s=0.
\]
This gives the positive coefficient \(h_i=b/q_i^c\) in
\eqref{maps:eq:D.OC1}.


\item \emph{The full normal and fibre kernel.}

The sweep of \(\delta\) bounds its angular tube over a base
disc in an original line frame. For \(\ell\), use the actual
identity-component translation logarithm and its lifted flow
in part 2 to close the positive contracting orbit over that
disc. Its elliptic increment is \(-bz\), its return count is
one, and the angular gauge is the same extending logarithm.
The cutoff maps and inverse flows carry these chains to the
given original arrow. Thus both full fixed-loop normal sweeps
are zero in \(H_2(U)\), in the canonical boundary marking.

The possible suspension branch change by \(k_{\rm br}\delta\) replaces
\(\ell\) by \(\ell+k_{\rm br}\delta\). It replaces one fixed sweep by
that sweep plus \(k_{\rm br}\) times the other. The inverse subtraction
preserves their entire kernel. It also leaves the first area
row in the fibre quotient map unchanged; the class of
\(\beta_0\wedge\delta\) in the fibre quotient is zero, as the fibre quotient map
checks. Hence the primitive class \(k\) and \eqref{maps:eq:D.OC1} are unaffected.
This proves the needed independence from the actual logarithm
branch by its explicit coordinate morphism.

The boundary Wang sequence has fibre part \(\mathbb Z^2\) and
fixed part \(\mathbb Z^2\). The primitive fibre map in part 3
has rank-one image \(\mathbb ZF\), and the two displayed fixed
sweeps bound. Therefore its complete boundary kernel is
\[
 \mathbb Z k\ \oplus\ \mathbb Zt(\delta)\ \oplus\ \mathbb Zt(\ell).
\]
Its complete cokernel is
\begin{equation}
 H_2(U)/\operatorname{im}H_2(\partial U)
       =\ker c/\mathbb ZF=\mathbb Z^{s_i-2}.
 \hypertarget{tag-0210}{}\label{maps:eq:D.OC6}
\end{equation}
These are integral direct
coordinates, proved from the quotient and its sections. This formula supplies both full fixed-loop relations as well as \eqref{maps:eq:D.OC1}.

The kernel of the area row on the original six-dimensional wedge
lattice has basis
\[
 e_{13},\quad e_{14}-z_{Q,2}e_{12},\quad e_{23},\quad
 e_{24}+z_{Q,1}e_{12},\quad e_{34}.
\]
Dropping the first wedge coordinate is its inverse. Its values
under \(\omega\) are \(0,-az_{Q,2},0,az_{Q,1},-b\), so
\begin{equation}
 h=\gcd(b,az_{Q,1},az_{Q,2})=b/q.
 \hypertarget{tag-0211}{}\label{maps:eq:D.OC7}
\end{equation}
Indeed \(q\) is the least denominator of \(z\), whence
\(\gcd(q,qz_1,qz_2)=1\). Write \(b=qh\), use
\(\gcd(a,b)=1\), and factor \(h\) from the three entries.
In the boundary Wang quotient this five-generator kernel becomes
the single primitive class \(k\): for a wedge \(x\) with
\(w(x)=0\), subtract its second quotient coordinate times \(k\).
The difference lies in \(\operatorname{im}(\bigwedge^2{\mathsf T}-I)\),
and the four unit Smith entries give an integral preimage.
The inverse includes \(k\). Thus both kernel descriptions have
the same integral image, including its sign in the global group.

The nonzero-Chern cases and their local free ranks are

\begin{center}
\small
\begin{tabular}{llll}
\hline
Type & \(s_i\) & \(q_i^c\) & Local free cokernel \\
\hline
III & 2 & 2 & 0 \\
IV & 3 & 3 & 1 \\
\(III^*\) & 8 & 2 & 6 \\
\(IV^*\) & 7 & 3 & 5 \\
\(I_1^*\) & 6 & 4 & 4 \\
\(I_3^*\) & 8 & 4 & 6 \\
\hline
\end{tabular}
\end{center}

Type II has trivial component group and never enters this table.
The natural component equations give 46 marked occurrences, with
counts 13, 9, 4, 8, 8 and 4 in the displayed type order.

\item \emph{The complete 282-pattern enumeration.}

Use the exhaustive natural construction in
Theorem~\ref{maps:src:D.2.31}, part 2. At
every finite good-reduction additive place the option set is
now \(\{\mathrm O,\mathrm F\}\), regardless of the component of \(P\). At every
odd star its option set is \(\{\mathrm O\}\). Here \(\mathrm O\) retains the
actual original germ and \(\mathrm F\) uses the full free affine family
of the finite calculation, with all positive allowed orders and
all free twists. The multiplicative models are fixed as in Theorem~\ref{maps:src:D.2.31}. Taking the Cartesian products gives all patterns, indexed by
their row, sites and options.
The finite arithmetic choices have their complete parameter scope
in Theorem~\ref{maps:src:D.2.31}, part 6; a pattern label does not restrict those parameters.

Enumerating the positive factor pairs of the natural equation and
these site options gives

\begin{center}
\small
\begin{tabular}{lllll}
\hline
OS & Profile & Natural rows & Patterns & Patterns with nonzero Chern \\
\hline
43 & base & 6 & 12 & 3 \\
43 & II & 2 & 8 & 2 \\
45 & base & 8 & 8 & 0 \\
45 & II & 4 & 8 & 0 \\
46 & base & 6 & 6 & 3 \\
46 & II & 2 & 4 & 2 \\
47 & base & 14 & 14 & 0 \\
47 & II & 6 & 12 & 0 \\
47 & III & 8 & 16 & 4 \\
49 & base & 10 & 20 & 5 \\
49 & II & 2 & 8 & 2 \\
49 & III & 4 & 16 & 7 \\
50 & base & 10 & 10 & 5 \\
50 & II & 2 & 4 & 2 \\
50 & IV & 4 & 8 & 5 \\
55 & base & 14 & 14 & 0 \\
55 & II & 6 & 12 & 0 \\
56 & base & 26 & 26 & 0 \\
56 & II & 10 & 20 & 0 \\
56 & III & 14 & 28 & 7 \\
56 & IV & 14 & 28 & 7 \\
\textbf{Total} &  & \textbf{172} & \textbf{282} & \textbf{54} \\
\hline
\end{tabular}
\end{center}

Deleting every pattern containing a nonzero-Chern original gives
exactly the 228-pattern set of
Theorem~\ref{maps:src:D.2.31}, with every row and site marking
retained. The complement consists of the 54 patterns containing
a nonzero-Chern original.

\item \emph{The integral global quotient with every original cap.}

Use \(L\) and \(\beta\) from
Definition~\ref{mapsIII:def:all-original-additive-data}.
For zero-Chern original caps use the primitive component split,
fibre map and two fixed-sweep relations in Theorem~\ref{maps:src:D.2.31}, parts 3--4.
For each nonzero-Chern original use the component-basis inverse, the fibre quotient map, and the full boundary-kernel formula.
For finite and regular caps retain the full Thom maps of Theorem~\ref{maps:src:D.2.31}, part 6, including the normal-character torsion and every
free affine twist.

The integral exterior splitting in Theorem~\ref{maps:src:D.2.31},
part 5, extends by zero blocks to an
arbitrary number of regular sites. Its only fibre coordinate is
the primitive class \(s\), with fibre functional \(\omega\). Eliminate each primitive boundary
image in the cap \(H_2\) by its displayed section, retaining
the explicit component complement. This gives an isomorphism
from the degree-two Mayer--Vietoris cokernel to
\[
                             C\oplus\mathbb Z^\beta.
\]
The forward map subtracts the chosen boundary lifts from the
exterior coordinate. The inverse includes the exterior coordinate
and the chosen component cycles. Altering a lift is exactly one
of the complete kernels already imposed, so these maps are
inverse in the integral quotients.

Here \(C\) is presented by the finite/regular matrix of Theorem~\ref{maps:src:D.2.31}, part 7, with both fixed sweeps set to zero at every original
additive and with the additional column \((b/q_i^c)s\) at
every nonzero-Chern original. In particular all finite core
torsion remains encoded by the three complete columns
\[
 (b g_{i,\rm den} n_i/q_{v,i},m_i/q_{v,i},0),\quad
 (-br_{i,\rm tw},0,m_i/g_{i,\rm den}),\quad(0,r_{i,\rm tw},n_i)
\]
in its local \((s,D_i,L_i)\) coordinates. The regular relations
are \(m_jt_j(u)-\omega(u,v_j)s=0\), \(t_j(v_j)=0\), and
the global relation is
\[
 \sum_{\mathrm{finite}}D_i+
 \sum_{\mathrm{regular}}t_j(\delta)+be\,s=0.
\]
All peripheral comparisons use \(v_i\mapsto v_i+N_{m_i}\lambda_i\)
at finite sites, together with the fixed-sweep shear and inverse
in Theorem~\ref{maps:src:D.2.31}, part 6; the regular case uses \(N_m=mI\).
The cokernel splitting and the global relation therefore hold
for every marked affine lift.

The full fundamental group is the presentation of
Theorem~\ref{maps:src:D.2.23}
\[
 \langle\gamma,c_i\mid [\gamma,c_i]=1,\quad
         c_i^{m_i}=\gamma^{\mu_i},\quad\prod_i c_i=\gamma^e\rangle,
 \qquad\mu_i=\chi(v_i).
\]
Indeed every original cap is covered by the same unrolled
contraction neighbourhood in Theorem~\ref{maps:src:D.2.23}. At a nonzero-Chern cap
the extra fibre relation kills \(\delta\), which is already
killed by the global vanishing discs in that proof; its remaining
fundamental group is the same positive \(\gamma\) by \eqref{maps:eq:D.OC4}.
The base meridian bounds the same section disc and has order one
and numerator zero. The maps take actual peripheral loops and
the contraction loop to the displayed generators; their inverse
uses the marked paths and vanishing discs in that theorem. No additional
relation in \(\gamma\) arises from the simply connected component
cycles. This proves the Seifert presentation with its whole nonabelian scope.

Write its abelianization as \(H_1(X)=\mathbb Z^r\oplus A_{1,\rm tor}\);
the rational Seifert matrix gives \(r=0\) or \(1\), with \(r=1\)
exactly at zero rational Euler number. If \(L=0\), the degree-two cokernel splitting and
the rest of the calculation are those of Theorem~\ref{maps:src:D.2.31}.

For \(L>0\), at least one \(b/q_i^c\) is a positive integer.
Consequently \(s=0\) in \(C\otimes\mathbb Q\). Every finite and
regular sweep is a rational multiple of \(s\) by the full
columns above, so \(C\otimes\mathbb Q=0\). Thus
\[
                             C=T=\operatorname{Tor}C
                         \quad\text{when }L>0.
\]
This remains true at zero rational Seifert Euler number; it
does not remove the bridge term computed next.

\item \emph{The full degree-one bridge, including the zero-Euler case.}

Put \(c_i^{\rm fix}=\operatorname{rank}\operatorname{Fix}({\mathsf T}_i)\) at an
original site and \(c_j^{\rm fix}=4\) at a new regular site. Every boundary
\(H_1\) is free of rank \(c_i^{\rm fix}+1\), by the integral graph and boundary matrices
in Theorem~\ref{maps:src:D.2.31}, part 5. The cap \(H_1\) ranks are \(c_i^{\rm fix}\) at every
minimal, finite, regular and zero-Chern original cap. Formula
\eqref{maps:eq:D.OC4} lowers the rank by exactly one at each of the \(L\)
nonzero-Chern originals. The boundary maps are onto cap \(H_1\),
including finite torsion, by the actual fibre-generation and
normal-disc maps. The exterior \(H_1\) is free of rank the
number \(N_{\partial}\) of boundary circles: its fibre coinvariant is the
primitive \(\mathbb Z\gamma\), and its cut base has rank \(N_{\partial}-1\).

The degree-zero Mayer--Vietoris map is injective at the direct
sum of the connected boundaries. Therefore the next kernel
\(K_1\) has rank
\[
 \begin{aligned}
 \operatorname{rank}K_1
 &=\sum_i(c_i^{\rm fix}+1)
   -\left(N_{\partial}+\sum_i c_i^{\rm fix}-L\right)+r\\
 &=L+r .
 \end{aligned}
\]
It is free because it is a subgroup of the direct sum of the
free boundary \(H_1\) groups. Torsion in the cap targets does
not change this assertion or this rational rank equation.
The case \(r=1\) has its full extra summand in the bridge-rank formula.

Choose an integral basis of this kernel using the complete
maps to the exterior and cap \(H_1\) presentations. For each
basis element, take its representing boundary one-cycle.
Both images are zero in integral \(H_1\), so choose integral
two-chains bounding them, and join them along that boundary
cycle. The resulting closed two-cycles split the exact sequence
with cokernel \(C\oplus\mathbb Z^\beta\). The inverse reads its Mayer--Vietoris
boundary coordinates, subtracts these cycles, and applies
the cap-elimination map of part 6. Hence for \(L>0\)
\[
                 H_2(X)=\mathbb Z^{\beta+L+r}\oplus T.
\]
This is an integral splitting; the free bridge cannot produce
an uncomputed torsion extension.

Every original cap is a mapping torus, so its Euler
characteristic is zero even when its component Chern class is
nonzero. All other zero-free caps have Euler characteristic
zero, and the same index-one cusp has Euler characteristic
two. Thus \(\chi(X)=2\). Duality and universal coefficients,
with the explicit Bockstein/linking coordinates described in
Theorem~\ref{maps:src:D.2.31}, part 7, give
\eqref{maps:eq:D.OC3}.
The middle rank follows from
\(b_3=2b_2-2b_1=2\beta+2L\). The torsion identifications use
the full integral Smith cokernel and the torsion linking
pairings, with the indicated inverse coordinates.

For example, at zero rational Euler number, take OS47/III, \(a=1,b=14\),
either permitted \(I_1\) zero, retain its original III cap,
and add no logarithmic sites. In the canonical marking
\(e=0\), the Seifert group is \(\mathbb Z\), so \(r=1\).
Here \(L=1,\beta=0,q^c=2\), and \eqref{maps:eq:D.OC1} is \(7s=0\).
Thus \(C=\mathbb Z/7\), while the bridge has rank two:
\[
 H_2(X)=H_3(X)=\mathbb Z^2\oplus\mathbb Z/7 .
\]
The free rank two comes from the bridge-rank formula, including the summand at
zero rational Euler number. This follows by substitution into the integral cap and bridge maps.

\end{enumerate}
\end{proof}

\begin{remark}[Engel and mixed-filling specializations]
\hypertarget{tag-Q129}{}\label{mapsIII:rem:original-additive-specializations}
{\normalfont\scriptsize tag Q129}\par
For the untwisted Engel row $P=G,Q=6G$ with $e=0$ and no
logarithmic sites, the $IV^*$ and $III$ original caps have
$s_i=7,2$ and $h_i=2,3$. Hence $\beta=5$, $L=2$, $T=0$ and
$r=1$, and \eqref{maps:eq:D.OC3} gives
$(1,1,8,14,8,1,1)$. This is the marked Engel row of
Theorem~\ref{maps:src:D.2.20} with its trivial regular logarithm
$(m,v)=(1,0)$: identity on the period lattice and on every original
cap, together with the ordinary-disc identification at that
order-one site, gives the input identification. That theorem
identifies the named member of \SeriesRef{push:src:A.4.2}.

For a mixed row of Theorem~\ref{maps:src:D.2.21}, write $d_0$
for its height denominator and $q_0$ for its finite monodromy order.
Its parameters here are
\[
 a=1,\quad b=d_0,\quad d_i=q_0,\quad g_{i,\rm den}=1,\quad
 m_i=q_0k,\quad n_i=n,\quad r_{i,\rm tw}=r,\quad e=0,
\]
with no regular logarithmic site and the same original additive cap.
The inverse reads $k=m_i/q_0,n=n_i,r=r_{i,\rm tw}$. Both constructions
use the fixed vector
$\ell_i=(-d_0z_{G,1},-d_0z_{G,2},0,1)$ and the same section line,
so identity on the marked periods carries their finite Thom,
original-cap and exterior maps to each other. Their five homology
specializations are therefore the formulas of that theorem.
\end{remark}

\begin{remark}[Finite original additive calculations]
\hypertarget{tag-Q130}{}\label{mapsIII:rem:all-original-additive-calculations}
{\normalfont\scriptsize tag Q130}\par
For the $46$ marked nonzero-Chern occurrences, compute the
primitive Chern vectors, their logarithm vectors, the full area
kernels and the primitive fibre images as in the preceding proof.
The products of the site-option sets give the $282$ patterns.
Smith reduction of the full graph, normal and fibre-kernel matrix
gives the degree-two cokernel. The integral degree-one
Mayer--Vietoris kernel supplies its free bridge, giving
\eqref{maps:eq:D.OC3} also at zero rational Euler number.
\end{remark}

\subsection{Original multiplicative caps and the full natural catalogue}

\begin{definition}[Natural data with original multiplicative choices]
\hypertarget{tag-Q131}{}\label{mapsIII:def:original-multiplicative-data}
{\normalfont\scriptsize tag Q131}\par
At every zero-free multiplicative fibre of the $172$ natural
initial rows, choose the original filling or the marked primitive
Mumford filling. At additive fibres use the original or finite
logarithmic choices of Definition~\ref{mapsIII:def:all-original-additive-data},
with the original filling at each $I_n^*$. Allow every permitted
positive finite or regular order, every free affine lift and
every integer clutch.
Let $L$ count all nonzero-Chern original fillings, and set
\[
 \beta=
 \sum_{\substack{\text{zero-Chern}\\\text{additive originals}}}(\rho_i-1)
 +\sum_{\substack{\text{nonzero-Chern}\\\text{additive originals}}}(s_i-2)
 +\sum_{\substack{\text{nonzero-Chern}\\\text{multiplicative originals}}}(n_i-2),
\]
where $\rho_i$ is the number of component orbits and $s_i$ the
number of Kodaira components. The original additive contributions
are those of Theorem~\ref{maps:src:D.2.32}.

\end{definition}

\begin{theorem}[Original multiplicative fillings]
\hypertarget{tag-02L9}{}\label{maps:src:D.2.33}
{\normalfont\scriptsize tag 02L9}\par
Identifying the equal marked boundary kernels at zero-Chern
multiplicative fibres in Definition~\ref{mapsIII:def:original-multiplicative-data}
gives $495$ integral homology presentation patterns.
There are $120$ marked zero-free multiplicative positions with
nonzero component Chern class. At such an $I_n$, the original filling
contributes the relation
\begin{equation}\hypertarget{tag-0222}{}\label{maps:eq:D.OM1}
                         h_i s=0,\qquad h_i=b/n,
\end{equation}
has free local degree-two cokernel $\mathbb Z^{n-2}$, and increases
the rank of the degree-one Mayer--Vietoris kernel by one.
Write the Seifert abelianization of Theorem~\ref{maps:src:D.2.23}
as $\mathbb Z^r\oplus A_1$. Let $T$ be the torsion subgroup of the
finite and regular normal presentation, including every additive
relation \eqref{maps:eq:D.OC1} and multiplicative relation
\eqref{maps:eq:D.OM1}. For $L>0$ the integral homology is
\begin{equation}\hypertarget{tag-0224}{}\label{maps:eq:D.OM2}
\begin{array}{c|c}
j&H_j(X;\mathbb Z)\\ \hline
0,6&\mathbb Z\\
1&\mathbb Z^r\oplus A_1\\
2&\mathbb Z^{\beta+L+r}\oplus T\\
3&\mathbb Z^{2\beta+2L}\oplus T\\
4&\mathbb Z^{\beta+L+r}\oplus A_1\\
5&\mathbb Z^r .
\end{array}
\end{equation}
For $L=0$ the free ranks in degrees $2,3,4$ are respectively
$2r+\beta$, $2r+2\beta$, $2r+\beta$, with the same torsion groups.
The full fundamental group is the Seifert group of
Theorem~\ref{maps:src:D.2.23}, for both zero and nonzero rational
Euler number.
\end{theorem}


\begin{proof}
Write the multiplicative fibre as $I_\nu$. Choose the primitive
vanishing vector $u$ and a vector $r_0$ with $\det(u,r_0)=1$.
The local defining data have
\[
 A_{\mathrm{ell}}-I=\nu uu^{\mathsf t}J,\qquad
 g=\kappa u,\qquad
 J=\begin{pmatrix}0&1\\-1&0\end{pmatrix},\qquad
 \gcd(\kappa,\nu)=1.
\]
The natural component equation gives $P=aG$ or $Q=bG$ narrow
at every zero-free multiplicative fibre, and $\gcd(a,b)=1$.

We first treat the nonzero-Chern original filling. Then $\nu>1$,
$\nu\mid b$ and $\gcd(a\kappa,\nu)=1$. Put
\[
 B=b\kappa/\nu,\qquad \ell=\gamma-Br_0.
\]
In the basis $(u,r_0,\delta,\gamma)$ the period return is
\[
 \begin{pmatrix}
 1&\nu&0&b\kappa\\
 0&1&0&0\\
 0&-a\kappa&1&-ab\kappa^2/\nu\\
 0&0&0&1
 \end{pmatrix}.
\]
Replacing $\gamma$ by $\ell$ conjugates it to
$\operatorname{diag}(I+(\nu,0,-a\kappa)^{\mathsf t}(0,1,0),1)$.
The inverse adds $B$ times the last coordinate to the $r_0$
coordinate. Thus, with $u_1=u$, $u_2=r_0$ and $p=-a\kappa$,
\begin{equation}\hypertarget{tag-0228}{}\label{maps:eq:D.OM6}
 Au_2=u_2+\nu u_1+p\delta,\qquad
 Au_1=u_1,\quad A\delta=\delta,\quad A\ell=\ell,\qquad
 \nu>0,\quad\gcd(\nu,p)=1.
\end{equation}
These period changes act by exponentiation on the local torus;
the group logarithm modulo the same lattice is their inverse,
and every period unit is transported with its exponent.

The forward return covers $-Q$. Its valuation is $b\kappa$,
while the elliptic long period has valuation $\nu$. Hence in
the positive Tate indexing the component of $P$ is
$\bar p=-a\kappa\bmod\nu$. In the negative Abel--Jacobi
component coordinate of Definition~\ref{mapsIII:def:elliptic-marked-data}
and Theorem~\ref{maps:src:D.2.22}, the section has coordinate
$a\kappa$. The map $j\mapsto-j$, with itself as inverse,
therefore carries that marking to the positive Tate indexing.
It fixes $C_0$, preserves the component cycle and intersection
matrix, and carries both section components and translation
arrows. The line class on the original neighbourhood $N$ is
\[
                 c_1(M)=C_{\bar p}^*-C_0^*,\qquad \bar p\ne0.
\]

\smallskip
\emph{The original line and its mapping torus.}
The central fibre $K$ is a cycle of $\nu$ rational curves.
Arcs between their nodes form a graph circle. Collapsing a
maximal tree in this graph gives a CW equivalence
\[
 K\simeq S^1\vee\bigvee_{j=0}^{\nu-1}S^2.
\]
Its inverse maps the graph circle around the node cycle and maps
each sphere with degree one to its labelled component, joining
basepoints along the tree. The regular-neighbourhood retraction
$N\to K$ transports these homology coordinates.

Remove the integral long-period part $Br_0$ of the narrow
section $Q$. Its remaining local N\'eron logarithm is holomorphic
on the base disc. Translation by its negative scalar multiples
gives an isotopy from the identity to the forward translation.
On the central cycle it scales each rational component by the
same multiplicative one-parameter group.

Write $c_j=[j=\bar p]-[j=0]$. Assign an integer weight at one
node and subtract $c_j$ across $C_j$. The sum of the $c_j$ is
zero, so the weights agree after traversing the cycle. The
homogeneous lift to the degree-$c_j$ line on $C_j$ has these
endpoint weights. The nonzero node gluing constants commute with
the matching characters. Equivalently, a component scaling by
$t$ changes the product of the line transitions by
$t^{\sum_jc_j}=1$. This constructs the lift on the given line.

In a node chart $xy=q$, the base vector field is a multiple of
$x\partial_x-y\partial_y$ and the line field has the common node
weight. Prescribe the angular normalization of this holomorphic
lift in a Hermitian metric on the central circle bundle. Extend
its vertical difference from a connection lift by a partition
of unity. The resulting flow covers the given base isotopy and
is inverted by the backward flow.

The ratio between its endpoint and the given holomorphic arrow
is constant on each compact rational component and agrees at
the nodes. Its angular part has zero winding on the graph circle.
The retraction $N\to K$ therefore gives a real logarithm $f$ for
that gauge. Multiplication by $\exp(itf)$ completes the isotopy
to the given angular arrow; the reciprocal gauge is the inverse.
The gauge extends over the origin-section disc, so it has zero
normal-meridian winding.

Let this angular isotopy be $f_t$, constant near its ends.
For the forward convention $(y,t)\sim(f_1y,t+1)$, the maps
\[
 [y,t]\longmapsto[f_t^{-1}y,t],\qquad
 [z,t]\longmapsto[f_tz,t]
\]
give inverse product coordinates. The identity
$f_1^{-1}f_1=I$ verifies the gluing. For a positive roof $R(y)$ use the variable-roof map $F_y$
in the proof of Theorem~\ref{maps:src:D.2.31}, part 3, with
$r(y)=R(y)$ and the same forward angular arrow. Its positive
derivative and equal endpoint derivatives give the smooth
identification; $F_y^{-1}$ is its inverse. This map uses the
given circle-bundle point as $y$ and changes only its roof coordinate.
Connection transport along the retraction to $K$ now gives
\begin{equation}\hypertarget{tag-0225}{}\label{maps:eq:D.OM3}
 U\simeq Y_c\times S^1_\ell,\qquad
 S^1_\delta\longrightarrow Y_c\longrightarrow K,\qquad
 c=C_{\bar p}^*-C_0^*.
\end{equation}
The physical contracting loop is $\gamma=\ell+Br_0$.

The evaluation of $c$ on $C_{\bar p}$ is one. The homotopy
sequence kills the angular generator and gives
$\pi_1(Y_c)=\mathbb Zr_0$. Gysin gives
\begin{equation}\hypertarget{tag-0227}{}\label{maps:eq:D.OM5}
 0\longrightarrow\mathbb Z[r_0\delta]
 \longrightarrow H_2(Y_c)
 \longrightarrow
 \ker(c:\mathbb Z^\nu\longrightarrow\mathbb Z)
 \longrightarrow0.
\end{equation}
The whole-fibre vector $F=(1,\ldots,1)$ and the $C_j$ for
$j\ne0,\bar p$ form a basis of the kernel. Its inverse sends
$y$ to $(y_0,(y_j-y_0)_{j\ne0,\bar p})$. Lift $F$ by the nearby
elliptic area and lift the other components by line sections.
Together with the rim torus these give an integral splitting.
Gysin and the product in \eqref{maps:eq:D.OM3} therefore yield
\begin{equation}\hypertarget{tag-0226}{}\label{maps:eq:D.OM4}
\begin{aligned}
 H_1(Y_c)&=\mathbb Z,& H_2(Y_c)&=\mathbb Z^\nu,&
 H_3(Y_c)&=\mathbb Z^\nu,\\
 H_1(U)&=\mathbb Z^2,& H_2(U)&=\mathbb Z^{\nu+1},&
 H_3(U)&=\mathbb Z^{2\nu},&
 H_4(U)&=\mathbb Z^\nu .
\end{aligned}
\end{equation}

\smallskip
\emph{The integral fibre and normal maps.}
Choose $x,z_{\rm Bez}\in\mathbb Z$ with $a\kappa x+\nu z_{\rm Bez}=1$ and set
$\beta_0=xu+z_{\rm Bez}\delta$. The basis
\[
 (\nu u-a\kappa\delta,\ r_0,\ \beta_0,\ \ell)
\]
has determinant one and return $I+E_{12}$. Its inverse on the
short plane has rows $(z_{\rm Bez},-x)$ and $(a\kappa,\nu)$. Thus
$A-I$ has one unit Smith factor, and $\bigwedge^2A-I$ has
two unit Smith factors.

In exterior order $(12,13,14,23,24,34)$ for
$(u,r_0,\delta,\ell)$, the fibre Wang quotient is
\[
 [y]\longmapsto
 (y_{12},y_{23},a\kappa y_{14}+\nu y_{34},y_{24}).
\]
Its inverse section inserts
\[
 ur_0,\quad r_0\delta,\quad k=\beta_0\ell,\quad r_0\ell.
\]
The two primitive relations are
$u\delta$ and $\nu u\ell-a\kappa\delta\ell$.
Their preimages under $\bigwedge^2A-I$ are
$-x\,ur_0+z_{\rm Bez}\,r_0\delta$ and $r_0\ell$, respectively.

Write the marked elliptic long-period multiplier as
$q^\nu U_{\rm hol}(q)$, where $U_{\rm hol}$ is its nonzero
holomorphic unit on the base disc. The nearby long path
$z(t)=z_0\exp(t\log(q^\nu U_{\rm hol}(q)))$ crosses the successive node
annuli once around the graph. Its endpoints differ by the
given elliptic long period. Thus $r_0$ maps with degree one
to the graph circle, while $u$ bounds a node vanishing disc.
The original fibre image consequently has basis
\begin{equation}\hypertarget{tag-0229}{}\label{maps:eq:D.OM7}
 ur_0,\qquad r_0\delta,\qquad r_0\ell .
\end{equation}
Their images are the whole-fibre lift, the primitive rim and
the graph-times-long-circle torus. The component inverse above
and these three tori give an integral image section. The raw
regular-fibre kernel is
\begin{equation}\hypertarget{tag-022A}{}\label{maps:eq:D.OM8}
 K_{\mathrm{fib}}
 =\mathbb Z(u\delta)\oplus\mathbb Z(u\ell)
                    \oplus\mathbb Z(\delta\ell).
\end{equation}
The first torus bounds the circle bundle over the vanishing
disc. The last two vanish because $u$ and $\delta$ are
nullhomotopic in $Y_c$. The basis change from the displayed
kernel to the two Wang relations and $k$ is
\[
 \begin{pmatrix}1&0&0\\0&\nu&x\\0&-a\kappa&z_{\rm Bez}\end{pmatrix},
 \qquad
 \begin{pmatrix}1&0&0\\0&z_{\rm Bez}&-x\\0&a\kappa&\nu\end{pmatrix}
\]
as inverse. Hence on the fibre Wang quotient the kernel is
precisely $\mathbb Zk$.

The normal sweeps in directions $\delta$ and $\ell$ bound
the origin-section and narrow-orbit tubes. Let $t_u$ be the
marked geometric normal torus in direction $u$, with the
origin-disc meridian marking. Its projection is the
vanishing-loop normal torus in the elliptic neighbourhood,
where the vanishing tube bounds it. Its long product
coordinate is zero.
It therefore equals $j[r_0\delta]$ for a unique integer $j$.
To read $j$, use the component inverse of \eqref{maps:eq:D.OM5}
and then the coefficient of the graph-times-angular cell in
the circle-bundle complex. The cycle
\[
               \widehat t_u=t_u-j[r_0\delta]
\]
bounds in the cap; the inverse change adds $j[r_0\delta]$.
In coordinates
$(ur_0,r_0\delta,k,r_0\ell,t_u,t_\delta,t_\ell)$
the map to the primitive fibre image is
\[
 O(j)=
 \begin{pmatrix}
 1&0&0&0&0&0&0\\
 0&1&0&0&j&0&0\\
 0&0&0&1&0&0&0
 \end{pmatrix}.
\]
The shear $Q_j=I-jE_{2,5}$ has inverse $I+jE_{2,5}$ and
$O(j)Q_j=O(0)$. The full boundary kernel is therefore generated
by $k,\widehat t_u,t_\delta,t_\ell$.

For the exterior fibre functional,
\begin{equation}\hypertarget{tag-022B}{}\label{maps:eq:D.OM9}
 \omega=a e_1^*\wedge e_2^*-b\delta^*\wedge\gamma^*,
 \qquad \omega(r_0,\delta)=0.
\end{equation}
The exterior map commutes with $Q_j$. It sends the complete
normal kernel to the three fixed sweeps with zero fibre
coefficient. Its value on $k$ is
\[
 \omega(k)=-aBx-bz_{\rm Bez}
          =-\frac b\nu(a\kappa x+\nu z_{\rm Bez})=-b/\nu.
\]
Replacing $k$ by $-k$ gives \eqref{maps:eq:D.OM1}.
Equivalently, the values on \eqref{maps:eq:D.OM8} are
$0,-ab\kappa/\nu,-b$, whose gcd is $b/\nu$.
The normal image is contained in \eqref{maps:eq:D.OM7},
so the full boundary cokernel is
\begin{equation}\hypertarget{tag-0223}{}\label{maps:eq:D.OM10}
 H_2(U)/\operatorname{im}H_2(\partial U)=\mathbb Z^{\nu-2}.
\end{equation}
The first-homology map kills $\beta_0$ and reads $r_0,\ell$;
the origin disc kills the base meridian. Thus its rank is two.

\smallskip
\emph{Zero-Chern original fillings.}
Here $\nu\mid a$ and $\gcd(b\kappa,\nu)=1$. Choose
$l\nu+k b\kappa=1$. The period basis
\[
 v=u-(a\kappa/\nu)\delta,\quad
 t=lr_0+k\gamma,\quad \delta,\quad
 w=-b\kappa r_0+\nu\gamma
\]
is unimodular. The inverse long matrix is
$\left(\begin{smallmatrix}\nu&b\kappa\\-k&l\end{smallmatrix}\right)$;
the short inverse adds $a\kappa/\nu$ times the short coordinate
to the angular coordinate. Its return is $I+E_{12}$ and its
two long valuations are one and zero.

Trivialize the zero-degree line on each normalized component
and keep its node gluing constants. Their product is its flat
factor around the graph circle. On the Tate cover the line
and translation lift are represented by multipliers independent
of the multiplicative coordinate. Their ratio with the given
holomorphic lift is constant on each compact fibre, and the
central constants agree at the nodes. Hence this ratio is a
nonzero base unit, carried by the period change.

The maps
\[
 X=z/q^j,\qquad Y=q^{j+1}/z,\qquad XY=q,\qquad
 z=Xq^j,\quad q=XY
\]
are the node charts and inverses. A unit in one chart coordinate
and its reciprocal in the other transport the period units.
The valuation-one period identifies consecutive components;
the valuation-zero period acts on the remaining torus coordinates.
The quotient has one ruled normalization, with its zero and
infinity sections glued by the given unit translation. Unrolling
this gluing and applying the inverse period matrix recover the
original charts.

In the basis $(v,t,\delta,w)$ the fibre quotient and section are
\[
 Q_{\rm fib}(y)=(y_{12},y_{23},y_{24},y_{34}),\qquad
 \sigma(z)=z_1e_{12}+z_2e_{23}+z_3e_{24}+z_4e_{34}.
\]
These four classes are the ruled fibre, the two gluing-circle
products and the base torus. Normalization, the gluing tubes
and projection to the base torus give their inverse coordinates.
For both the original and marked primitive filling the boundary
map has this target and kills the three fixed sweeps. If the maps
are $f_O,f_M$ with the displayed sections, the target isomorphism
is $f_M\sigma_O$, with inverse $f_O\sigma_M$. It commutes with
the common marked boundary map.

\smallskip
\emph{Every finite and regular parameter.}
The exterior calculation of Theorems~\ref{maps:src:D.2.30}
and \ref{maps:src:D.2.32} has fibre quotient $\mathbb Zs$
with coordinate $\omega$. Its section sends $s$ to
$ue_{12}+ve_{34}$ for $au-bv=1$. Its graph coordinates are
the integral fixed sweeps with the one sum of their
$\delta$ directions. The present boundary maps enter these
same coordinates by \eqref{maps:eq:D.OM9}.

At each finite and regular logarithmic site, the marked lattice,
power vector, root coordinate and normal orientation are those of
Theorem~\ref{maps:src:D.2.31}, part 6. The identity on these
coordinates carries the full cap diagrams to the ones used here.
In particular a finite site has $m_i=d_i k_i>0$,
$v_i=r_i\delta+n_i\ell_i$ and
\[
 \gcd(m_i,r_i,n_i)=1,\qquad d_i\mid m_i/\tau_i,\qquad
 \tau_i=\gcd(m_i,q_{v,i}r_i,g_i n_i).
\]
Its complete Thom columns, in the identical ordered coordinates
$(s,D_i=t_i(\delta),L_i=t_i(\ell_i))$, are
\[
 \begin{pmatrix}b g_i n_i/q_{v,i}\\m_i/q_{v,i}\\0\end{pmatrix},
 \qquad\begin{pmatrix}-br_i\\0\\m_i/g_i\end{pmatrix},
 \qquad\begin{pmatrix}0\\r_i\\n_i\end{pmatrix}.
\]
The boundary root-argument map and its inverse in that proof give
the normal projection $N_{m_i}\lambda$, vertical class
$-\lambda\wedge v_i$, and root sweep $v_i$. Thus the full
normal character of order $\tau_i$ and the core torsion are retained.
The holomorphic finite-sum conjugacy in
Theorem~\ref{maps:src:D.2.29} applies to the same root-disc operator
and power class, with its inverse translation; it covers every
holomorphic affine lift, including the central point.
At a regular site, $m_j>0$, $v_j=(x_j,y_j,z_j,\mu_j)$ and
$\gcd(m_j,x_j,y_j,z_j,\mu_j)=1$. Its complete relations are
\[
 m_jt_j(u)-\omega(u,v_j)s=0,\qquad t_j(v_j)=0.
\]
They use the same normal-sector maps and integral central-lattice
quotient, including $m_j=1$.
At every original filling set its full fixed sweeps to zero
and impose its fibre relation when its component Chern class
is nonzero. The one remaining global relation is
\[
 \sum_{\text{finite}}D_i+
 \sum_{\text{regular}}t_j(\delta)+be\,s=0.
\]
Call this abelian presentation $C$.

The peripheral-lift maps in Theorem~\ref{maps:src:D.2.31}, part 6,
fix the present original-cap diagrams and act on the unchanged
finite/regular diagrams by
\[
 v_i\mapsto v_i+N_{m_i}\lambda_i,\qquad
 t'_i(u)\mapsto t_i(u)+\omega(u,\lambda_i)s,
 \qquad e\mapsto e+\sum_i\chi(\lambda_i).
\]
Their inverse uses $-\lambda_i$ and subtracts the sweep shear.
The norm-pairing identities proved there preserve every root
column and the global relation; at a regular site $N_m=mI$.
This gives every permitted peripheral marking on the present fillings.

\smallskip
\emph{Global groups.}
Use the primitive image sections of each original filling
and the full Thom surjections at the logarithmic fillings.
Subtracting those boundary lifts from the exterior coordinate
identifies the degree-two Mayer--Vietoris cokernel with
\[
                         C\oplus\mathbb Z^\beta.
\]
Its inverse includes the exterior and the chosen component
classes. Different image lifts differ by the imposed full kernels.

The new original cap kills $u,\delta$ and retains $r_0,\ell$.
The global vanishing discs already kill $r_0$, and
$\ell=\gamma-Br_0$ therefore maps to $\gamma$. The origin
disc kills the meridian. These maps and their inverse marked
loops preserve the full presentation
\[
 \pi_1(X)=
 \left\langle\gamma,c_1,\ldots,c_s\ \middle|\
 [\gamma,c_i]=1,\ c_i^{m_i}=\gamma^{\mu_i},\
 c_1\cdots c_s=\gamma^e\right\rangle ,
 \qquad \mu_i=\chi(v_i).
\]
In particular $\mu_i=g_i n_i$ at a finite site. Its abelian
relation rows are
\[
 (-\mu_i,0,\ldots,m_i,\ldots,0),\qquad (-e,1,\ldots,1).
\]
Writing the cokernel as $\mathbb Z^r\oplus A_1$ gives
$r=1$ exactly when $e-\sum_i\mu_i/m_i=0$.

For $L=0$, rational elimination in $C$ expresses every sweep
as a multiple of $s$. The remaining relation is
$b(e-\sum_i\mu_i/m_i)s=0$, so $C$ has free rank $r$.
For $L>0$, one positive fibre relation kills $s$ rationally,
and then every sweep vanishes rationally. Thus $C$ is finite.
In both cases let $T=\operatorname{Tor}C$.

At boundary $i$, let the fixed rank be $r_i^{\rm fix}$; its first homology
is free of rank $r_i^{\rm fix}+1$. The cap has rank $r_i^{\rm fix}$, except that
a nonzero-Chern original has rank $r_i^{\rm fix}-1$. The exterior has
rank $N_{\rm bd}$ for $N_{\rm bd}$ boundary circles, and the degree-zero
Mayer--Vietoris map is injective. The integral degree-one
kernel is consequently free of rank
\[
 \sum_i(r_i^{\rm fix}+1)-\left(N_{\rm bd}+\sum_i r_i^{\rm fix}-L\right)+r=L+r.
\]
An integral basis of this kernel is represented by boundary
one-cycles. Bound their images by integral two-chains in the
exterior and caps and join the chains across their boundaries.
The resulting closed two-cycles give a section. Its inverse
reads the connecting coordinate and subtracts those cycles.
Therefore
\[
               H_2(X)=C\oplus\mathbb Z^\beta\oplus\mathbb Z^{L+r}.
\]

Every original cap is a mapping torus and has Euler number zero.
Finite, regular and primitive rank-one fillings also have Euler
number zero. The index-one cusp has Euler number two, so
$\chi(X)=2$. Poincar\'e duality and integral universal coefficients
give the stated groups, with
$b_3=2b_2-2b_1$. For a Smith generator of order $h$, its character
of value $1/h$, the integral Bockstein and oriented cap product
give the dual torsion coordinate; the linking coordinate is the
inverse. This proves \eqref{maps:eq:D.OM2} and its $L=0$ case.

\smallskip
\emph{Enumeration.}
Use the $172$ natural rows of Theorem~\ref{maps:src:D.2.22},
identified by the ordered profile, selected zero, section pair and
all marked component coordinates as in Theorem~\ref{maps:src:D.2.31}.
The identification fixes those coordinates.

On a row let $f_{\rm en}$ count finite additive sites, let $c_{\rm en}$ count
nonzero-Chern zero-free multiplicative sites, and let $z_{\rm en}$
count zero-Chern zero-free multiplicative sites. There are
$2^{f_{\rm en}+c_{\rm en}}$ patterns and $2^{f_{\rm en}+c_{\rm en}+z_{\rm en}}$ individual choices.
The quotient map forgets the $z_{\rm en}$ zero-Chern choices, and its
section chooses the primitive filling at those sites. The
target homology isomorphisms constructed above realize each
agreement. At a nonzero-Chern site the original kills $\beta_0\ell$
and the primitive filling maps it to a primitive nonzero class.
At an additive original the two fixed sweeps bound; for a
finite logarithmic filling their rational normal transfer is
$-\omega(-,v_i)/m_i$. The fixed-plane alternating form is
nondegenerate and the free vector is nonzero. Thus some
fixed sweep has a nonzero image, distinguishing these local
kernel choices.

Here are the exhaustive counts. The four numeric columns
are the initial rows, nonzero-Chern multiplicative positions,
kernel patterns and individual choices, respectively:
\[
\begin{array}{cc|rrrr}
\mathrm{OS}&\mathrm{profile}&\mathrm R&\mathrm K&\mathrm P&\mathrm A\\ \hline
43&\mathrm{base}&6&0&12&48\\
43&II&2&0&8&8\\
45&\mathrm{base}&8&4&12&128\\
45&II&4&2&12&32\\
46&\mathrm{base}&6&0&6&24\\
46&II&2&0&4&4\\
47&\mathrm{base}&14&13&30&224\\
47&II&6&5&24&48\\
47&III&8&4&24&64\\
49&\mathrm{base}&10&4&28&80\\
49&II&2&0&8&8\\
49&III&4&0&16&16\\
50&\mathrm{base}&10&4&14&40\\
50&II&2&0&4&4\\
50&IV&4&0&8&8\\
55&\mathrm{base}&14&13&30&224\\
55&II&6&5&24&48\\
56&\mathrm{base}&26&34&77&416\\
56&II&10&10&46&80\\
56&III&14&11&54&112\\
56&IV&14&11&54&112\\ \hline
\mathrm{Total}&{}&172&120&495&1728
\end{array}
\]
There are $316$ zero-Chern multiplicative positions. Projection
to the additive choices gives the $282$ patterns of
Theorem~\ref{maps:src:D.2.32}; adjoining the independent
nonzero-Chern multiplicative choices gives the displayed
$495$. The quotient rule lists every individual representative
and supplies an inverse choice for each kernel pattern.
\end{proof}

\begin{remark}[Finite original multiplicative calculations]
\hypertarget{tag-Q132}{}\label{mapsIII:rem:original-multiplicative-calculations}
{\normalfont\scriptsize tag Q132}\par
At the $120$ marked nonzero-Chern multiplicative positions,
the integral normal forms determine the primitive fibre and
monodromy kernels, their values under $\omega$, and the primitive
Chern and component-multiplicity vectors. The occurrences are
$I_8$ in $6$ cases, $I_7$ in $14$, $I_2$ in $31$, $I_3$ in $23$,
$I_5$ in $38$ and $I_4$ in $8$. Smith reduction of the full graph,
normal and fibre matrix gives the degree-two cokernel. Adjoining
the free degree-one Mayer--Vietoris bridge gives the homology
formulas of Theorem~\ref{maps:src:D.2.33} for the full stated range
of orders and free affine parameters.
\end{remark}

\section{Weighted profiles and saturated cap maps}
\subsection{All positive weights of the Engel construction}

\begin{definition}[Positive-weight Engel data]
\hypertarget{tag-Q133}{}\label{mapsIII:def:positive-engel-data}
{\normalfont\scriptsize tag Q133}\par
On the $IV^*,III,I_1$ rational elliptic surface of
Theorem~\ref{maps:src:D.2.5}, let $G$ generate the
Mordell--Weil group. Use $M_S=\mathcal O(G-O)$, translation
$6dG$ and linearization zero $dI_1$, where $d$ is a positive
integer. At $I_1$ use the period subgroup
$\mathbb Ze_1+d\mathbb Ze_2$ of the infinite marked hexagonal
toric model. Number the $IV^*$ and $III$ sites by $0$ and $1$.
At either site choose its original cap or a finite logarithmic
action with parameters
\[
 m_0=3k,\qquad m_1=4\ell,\qquad k,\ell>0,\qquad
 v_i=r_i\delta+n_i u_i,\quad r_i,n_i\in\mathbb Z,
\]
where, in the marked period basis $(\delta,a,b,\gamma)$ of that theorem,
\[
 u_0=\gamma+2d(a+b),\qquad u_1=\gamma+3d(a+b).
\]
At an original site put $(\overline m_i,\overline n_i)=(1,0)$;
at a logarithmic site put $(\overline m_i,\overline n_i)=(m_i,n_i)$.
Define
\begin{equation}
 P=\overline m_1\overline n_0+
   \overline m_0\overline n_1,\qquad
 R=\begin{pmatrix}
  -\overline n_0&\overline m_0\\
  -\overline n_1&-\overline m_1
 \end{pmatrix},\qquad
 \mathbb Z^r\oplus A_{1,\rm tor}=\operatorname{coker}R.
 \hypertarget{tag-026U}{}\label{maps:eq:D.WE2}
\end{equation}
\end{definition}

\begin{theorem}[Every positive Engel weight]
\hypertarget{tag-02LA}{}\label{maps:src:D.2.34}
{\normalfont\scriptsize tag 02LA}\par
The data of Definition~\ref{mapsIII:def:positive-engel-data}
give a smooth $d$-period hexagonal cap for every $d>0$.
The exact freeness conditions at the chosen finite sites are
\begin{equation}
\begin{aligned}
 t_0&=\gcd(m_0,n_0,3r_0),&
 &\gcd(m_0,n_0,r_0)=1,\quad 3\mid m_0/t_0,\\
 t_1&=\gcd(m_1,n_1,2r_1),&
 &\gcd(m_1,n_1,r_1)=1,\quad 4\mid m_1/t_1 .
\end{aligned}                                             \hypertarget{tag-026O}{}\label{maps:eq:D.WE1}
\end{equation}
All noninvariant parts of affine translations are included
through the marked translation conjugacy. For every free choice,
the full fundamental group is $\mathbb Z^r\oplus A_{1,\rm tor}$ of
\eqref{maps:eq:D.WE2}, and $r=1$ precisely when $P=0$.

The exact degree-two matrix has target $(s,t,t_0',t_1')$.
For two logarithmic caps it is $B=(B_0\ B_1)$, where
\begin{equation}
 B_0=
 \begin{pmatrix}
 -2dn_0&-6dr_0&0\\
 -k&0&r_0\\
 0&m_0&n_0\\
 0&0&0
 \end{pmatrix},
 \qquad
 B_1=
 \begin{pmatrix}
 -3dn_1&-6dr_1&0\\
 2\ell&0&-r_1\\
 0&0&0\\
 0&m_1&n_1
 \end{pmatrix}.                                           \hypertarget{tag-026V}{}\label{maps:eq:D.WE3}
\end{equation}
For an original $IV^*$ or $III$ cap replace its three
columns, respectively, by
\begin{equation}
 O_0=
 \begin{pmatrix}2d&0&0\\0&1&0\\0&0&1\\0&0&0\end{pmatrix},
 \qquad
 O_1=
 \begin{pmatrix}3d&0&0\\0&-1&0\\0&0&0\\0&0&1\end{pmatrix}.
                                                               \hypertarget{tag-026W}{}\label{maps:eq:D.WE4}
\end{equation}
Let $T$ be the torsion of $\operatorname{coker}B$,
let $L$ count original additive caps, and put
$\beta=5$ if $IV^*$ is original, and $\beta=0$ otherwise.
Set
\begin{equation}
 (b_2,b_3)=
 \begin{cases}
 (d-1+2r,\;2r),&L=0,\\
 (d-1+\beta+L+r,\;2\beta+2L),&L>0.
 \end{cases}                                               \hypertarget{tag-026X}{}\label{maps:eq:D.WE5}
\end{equation}
Then all the integral homology is
\begin{equation}
\begin{array}{c|c}
j&H_j(X;\mathbb Z)\\\hline
0,6&\mathbb Z\\
1&\mathbb Z^r\oplus A_{1,\rm tor}\\
2&\mathbb Z^{b_2}\oplus T\\
3&\mathbb Z^{b_3}\oplus T\\
4&\mathbb Z^{b_2}\oplus A_{1,\rm tor}\\
5&\mathbb Z^r .
\end{array}                                                \hypertarget{tag-026Y}{}\label{maps:eq:D.WE6}
\end{equation}
In particular $\chi(X)=2d$ for every choice. At $d=1$ the
construction is the natural row with section multipliers $(1,6)$;
at $d>1$ its Euler characteristic is greater than two.

For two logarithmic caps with $\gcd(m_i,n_i)=1$, put
$g=\gcd(d,k,2\ell)$. Then
\begin{equation}
 T=
 \begin{cases}
 \mathbb Z/g\oplus\mathbb Z/(d|P|/g),&P\ne0,\\
 \mathbb Z/g,&P=0.
 \end{cases}                                               \hypertarget{tag-026Z}{}\label{maps:eq:D.WE7}
\end{equation}
The fundamental group is $\mathbb Z/|P|$ for $P\ne0$
and $\mathbb Z$ for $P=0$. With at least one original
cap the torsion groups, for every free remaining twist,
are
\begin{equation}
\begin{array}{c|c}
\text{original sites}&T\\\hline
IV^*\text{ only}&
 \mathbb Z/(d\gcd(2,n_1))\oplus\mathbb Z/\gcd(m_1,n_1)\\
III\text{ only}&
 \mathbb Z/(d\gcd(3,n_0))\oplus\mathbb Z/\gcd(m_0,n_0)\\
\text{both}&\mathbb Z/d .
\end{array}                                                \hypertarget{tag-0270}{}\label{maps:eq:D.WE8}
\end{equation}
The displayed summands are the independent cyclic coordinates
constructed in the proof.

\end{theorem}


\begin{proof}
We use the marked line, periods and toric charts of
Theorem~\ref{maps:src:D.2.5}. Write the lattice as
$\Lambda=\mathbb Z\langle\delta,a,b,\gamma\rangle$.
The weight-one arrow on $M_S=\mathcal O(G-O)$ covers $-6G$;
its $d$-th iterate covers $-6dG$ and has zero divisor $dI_1$.
The ratio of this iterate and the prescribed weight-$d$ arrow
is a nowhere-zero holomorphic function on the compact connected
surface, hence a constant. Multiplying the weight-one arrow
by a $d$-th root of this constant makes its iterate equal to
the prescribed arrow. Denote the resulting arrow by $\psi$.

Over $S\setminus I_1$, put $M_0=M_S|_{S\setminus I_1}$.
The identity on $M_0^\times$ induces
\[
 M_0^\times/\langle\psi^d\rangle
       \longrightarrow M_0^\times/\langle\psi\rangle,
                    \qquad [z]\longmapsto[z].
\]
A stabilizer for a power of $\psi$ would give a stabilizer
for a power of the given contracting action $\psi^d$.
Properness follows by decomposing the powers of $\psi$ into
the $d$ cosets of $\langle\psi^d\rangle$ and applying
properness to the finitely many translated compact sets.
Thus this is a degree-$d$ cover. Its inverse on an evenly
covered set is the unique lift to the chosen sheet. Its period
inclusion is $P_d=\operatorname{diag}(1,1,1,d)$, whose inverse
on its image divides the last coordinate by $d$.

On the infinite toric cover at $I_1$, take the period subgroup
$\mathbb Ze_1+d\mathbb Ze_2$. The identity of toric charts,
with their original nonzero units, extends the same cover.
A nonzero period translates component vertices and cannot
preserve the finite nonempty set of vertices incident at a
point. The period action is free. The locally finite charts
give evenly covered neighbourhoods, with inverse sheet maps
given by the inverse chart and its chosen period representative.
This constructs the required smooth hexagonal cusp for every $d$.

\smallskip
\emph{The cusp and the primitive exterior maps.}
The cusp normalization consists of $d$ copies of the toric
blowup of $\mathbb P^2$ at its three coordinate points.
The branch maps and their inverses are those of
Theorem~\ref{maps:src:D.1.2}, now on the specified period subgroup.
Over $\mathcal A_d=\mathbb Z[t_{\rm cyc}]/(t_{\rm cyc}^d-1)$ their degree-two restriction matrix is
\[
 \begin{pmatrix}
 -1&-1&-1&-1\\1&1&1&t_{\rm cyc}\\-t_{\rm cyc}&-t_{\rm cyc}&-1&-t_{\rm cyc}
 \end{pmatrix}.
\]
Adding the first row to the second and subtracting $t_{\rm cyc}$ times
the first from the third leave $t_{\rm cyc}-1$ in the fourth and third
columns. Clearing the first-row unit gives $1,t_{\rm cyc}-1,t_{\rm cyc}-1$ and
a zero column. Opposite row additions invert these operations.
On coefficient tuples, partial sums give the inverse of $t_{\rm cyc}-1$
on the zero-sum subgroup, and its kernel is the constant tuples.
Its kernel and cokernel are therefore both $\mathbb Z$.
The degree-zero row is the dual marked torus complex; the
degree-four row is $\mathcal A_d$. The augmented branch-simplex complexes
are exact also at self-identifications. The second differential
has zero odd vertical target, and higher differentials leave
the three columns. Lifting the torus two-cycle in the total
complex splits the remaining free extension. Thus the cusp
neighbourhood $N_d$ has
\begin{equation}\hypertarget{tag-0271}{}\label{maps:eq:D.WE9}
 (H_0,H_1,H_2,H_3,H_4)(N_d;\mathbb Z)
 = (\mathbb Z,\mathbb Z^2,\mathbb Z^{d+3},
                         \mathbb Z^2,\mathbb Z^d).
\end{equation}
The row operations prove this also at weights divisible by
two or three.

The actual period matrices are
\begin{equation}\hypertarget{tag-026P}{}\label{maps:eq:D.WE10}
\begin{aligned}
 A_0&=\begin{pmatrix}1&-1&0&2d\\0&-1&-1&6d\\
                 0&1&0&0\\0&0&0&1\end{pmatrix},\\
 A_1&=\begin{pmatrix}1&0&1&-3d\\0&0&1&0\\
                 0&-1&0&6d\\0&0&0&1\end{pmatrix},&
 A_\infty&=\begin{pmatrix}1&0&0&d\\0&1&1&0\\
                 0&0&1&0\\0&0&0&1\end{pmatrix}.
\end{aligned}
\end{equation}
They satisfy $P_dA_i(d)=A_i(1)P_d$,
$A_0^3=A_1^4=A_0A_1A_\infty=I$. Their fixed lattices have
bases $(\delta,u_0)$, $(\delta,u_1)$ and $(\delta,a)$, where
\[
 u_0=\gamma+2d(a+b),\qquad u_1=\gamma+3d(a+b).
\]
These integral period changes act by the period exponential;
the group logarithm modulo the same lattice gives their inverse.
The invariant fibre functional is
\[
 \xi_d=a^*\wedge b^*+6d\,\gamma^*\wedge\delta^*.
\]
It is the pullback under $P_d$ of its weight-one counterpart.

The images of $A_0-I$ and $A_1-I$ have respective bases
$(\delta+3a,a+b)$ and $(\delta+2a,a+b)$, with preimages
$(-a-b,-b)$ and $(b-a,-a)$. Their common quotient is the
primitive $\mathbb Z\gamma$. In particular
\[
\begin{aligned}
 a&=(A_0-I)(-a-b)-(A_1-I)(b-a),\\
 \delta&=2(A_0-I)(a+b)+3(A_1-I)(b-a),\\
 a+b&=-(A_0-I)b.
\end{aligned}
\]
Tracing these loops along the corresponding base edges gives
integral bounding chains for the three fibre coordinates.

Put $(s_1,\ldots,s_6)=(\delta a,\delta b,\delta\gamma,
ab,a\gamma,b\gamma)$ and $E_i=\bigwedge^2A_i-I$.
Their combined image is the primitive kernel of $\xi_d$.
The following explicit preimages prove that assertion:
\[
\begin{aligned}
 p_1&=E_0s_4=s_1,&p_2&=E_1s_4=s_2,\\
 p_5&=-E_0s_5-E_1s_6-dp_1+4dp_2=s_5,\\
 p_6&=2dp_1-E_0s_6-p_5=s_6,\\
 p_3&=E_1s_6-p_5+p_6-3dp_1-6dp_2=s_3+6ds_4.
\end{aligned}
\]
Conversely $\xi_d$ annihilates every $E_i$. The quotient
coordinate is $\xi_d$, with section $1\mapsto ab$; write
$s=[ab]$. In the kernel basis
$(s_1,s_2,s_3+6ds_4,s_5,s_6)$, columns $0,3,4,5,11$ of
$(E_0\ E_1\ E_\infty)$ have determinant $-1$.

Let $V$ be the torus bundle over the three-holed sphere.
Cutting its section loop $x_\infty=x_1^{-1}x_0^{-1}$ gives
the cusp graph map $j_\infty(\lambda)=-(A_1\lambda,\lambda)$.
The five graph-kernel columns are
\begin{equation}\hypertarget{tag-026Q}{}\label{maps:eq:D.WE11}
 (-\delta,-\delta),\quad(b,-a),\quad
 (\delta,0),\quad(u_0,0),\quad(0,u_1).
\end{equation}
Their minor in rows $0,2,3,4,7$ is $-1$. Here is their
integral inverse. For $(v_0,v_1)$ in the graph kernel, put
$h_{\rm gr}=\gamma^*(v_1)$. Applying the primitive $A_0$ coinvariant
row $(-3,1,-1,0)$ gives $b^*(v_1)=3dh_{\rm gr}$. Hence
$\lambda=v_1-h_{\rm gr}u_1$ lies in $\langle\delta,a\rangle$.
Adding $j_\infty(\lambda)$ leaves
$(v_0-A_1\lambda,h_{\rm gr}u_1)$, whose first component is uniquely
in $\langle\delta,u_0\rangle$. This reads all five coordinates
and proves both inverse composites.

The first two columns are the angular cusp sweeps. The
cellular sequence over the cut base has fibre term
$\mathbb Zs$ and this graph kernel. Actual boundary sweeps
lift the displayed basis. The invariant vertical form
$\xi_d$ is zero on these constant-period sweeps and one on
$s$, so subtracting their coordinates and applying $\xi_d$
gives an integral splitting. The sole sum of the angular
sweeps bounds the physical $\delta$ circle over the cut base.
Writing $L_c$ for the subgroup of the two cusp sweeps, we obtain
\[
 H_2(V)/L_c=\mathbb Z\langle s,t,t'_0,t'_1\rangle,
 \quad t=t_0(\delta),\quad t'_i=t_i(u_i),\quad
 t_1(\delta)=-t.
\]
In particular $L_c$ is primitive in the full $H_2(V)$.

The infinite central toric divisor is simply connected:
its simply connected components and double curves have the
edge-path presentation of the triangulated plane, with the
triple-point triangle relations. The free period quotient
therefore gives $\pi_1(N_d)=\mathbb Z\langle b,\gamma\rangle$.
Its boundary map kills $\delta,a$ and the section meridian,
and has the inverse period coordinates on $b,\gamma$.
In $W=V\cup N_d$, the $A_0$ conjugation relation for $a$
then kills $b$. Thus
\[
 \pi_1(W)=\mathbb Z\langle\gamma,x\rangle,
                         \qquad x=x_0=x_1^{-1}.
\]

The cusp boundary has $H_1=\mathbb Z^3\oplus\mathbb Z/d$,
where the last generator is $\delta$. To retain that torsion,
use the full relative sequence
\[
 H_3(N_d,\partial N_d)\longrightarrow H_2(V)
 \longrightarrow H_2(W)\xrightarrow{{\mathcal I}}\mathbb Z^d
 \xrightarrow{\partial}H_1(V).
\]
By \eqref{maps:eq:D.WE9} and duality its first group is
$\mathbb Z^2$. The two short-circle section-disc thimbles
lift from weight one because their contracting coordinates
are zero. Their images are $L_c$. The entire relative image
has rank at most two; modulo the primitive $L_c$ it is
torsion in the free quotient $H_2(V)/L_c$, hence zero.

The $d$ relative component discs all have boundary
$x_\infty$ in $H_1(V)=\mathbb Z\langle\gamma,x_0,x_1\rangle$.
Their extra short displacements are killed by the explicit
trace chains above. The relative kernel has basis $e_j-e_0$,
$1\le j<d$. Take differences of the component discs and
bound their remaining boundary cycles by those integral
chains in $V$. The closed cycles $\Sigma_j$ obtained this
way give an integral section, with inverse reading ${\mathcal I}$ and
subtracting their coordinates. Consequently
\[
 H_2(W)=\mathbb Z\langle s,t,t'_0,t'_1\rangle
                 \oplus\bigoplus_{j=1}^{d-1}\mathbb Z[\Sigma_j].
\]
The remaining finite-boundary cycles have zero component
intersection ${\mathcal I}$, so their full maps lie in the first summand.

\smallskip
\emph{Every finite affine filling.}
Put $q_0=3,q_1=2$ and $e_0^{\rm ord}=3,e_1^{\rm ord}=4$. The primitive
coinvariant maps and sections are
\[
 \begin{pmatrix}-3&1&-1&0\\0&0&0&1\end{pmatrix},\quad
 \begin{pmatrix}-2&1&-1&0\\0&0&0&1\end{pmatrix},
                         \qquad a,\gamma.
\]
Their inverse relation coordinates are the preimages already
given. Thus $[\delta]=-q_i[a]$ and $[u_i]=[\gamma]$.
The positive vertical coordinate of the full finite-cap
calculation in Theorem~\ref{maps:src:D.2.29} is $-a$;
the inverse is $a\mapsto-\beta_i$. This identifies its
necessary and sufficient free-action test with
\eqref{maps:eq:D.WE1}, including the normal-character order
$t_i$. The parameter $d$ changes only the displayed integral
fixed vectors $u_i$ and adds no condition on that test.

Use the positive-root generator of
Theorem~\ref{maps:src:D.2.29}, with
$(y,\tau)\sim(A_i y,\tau+1)$. In these coordinates the boundary map is
\[
 [y,\tau]\longmapsto
 [y+(\tau/m_i)v_i,\exp(2\pi i\tau/m_i)].
\]
Its inverse chooses the root argument and subtracts the
translation. Cutting a lattice loop into its root sectors
gives fibre coefficient $-\xi_d(\lambda,v_i)$ and sweep
$N_{m_i}\lambda$; the root loop has fibre coefficient zero
and sweep $v_i$. Rejoining the sectors is the inverse chain
construction. Direct multiplication gives
\[
\begin{aligned}
 N_{m_0}a&=-k\delta,&N_{m_1}a&=-2\ell\delta,&
 N_{m_i}\gamma&=m_i u_i,\\
 \xi_d(a,u_0)&=2d,&\xi_d(a,u_1)&=3d,&
 \xi_d(\gamma,\delta)&=6d,
 \qquad\xi_d(\gamma,u_i)=0.
\end{aligned}
\]
The full Thom image uses $a,\gamma,c_i$ in the integral
core presentation
\[
 \langle a,\gamma,c_i\mid
                    q_i r_i a-n_i\gamma+m_i c_i=0\rangle.
\]
It retains its torsion and the normal Chern class. In the
Thom sequence the next map evaluates that torsion Chern
class in free $H_0$, so it is zero; the boundary-to-cap
$H_2$ map is onto. The three columns in the four exterior
coordinates are exactly \eqref{maps:eq:D.WE3}. In particular
the signs of the second block follow from $t_1(\delta)=-t$.

The whole-root-disc translation conjugacy of
Theorem~\ref{maps:src:D.2.29} applies with lattice $\Lambda$,
matrix $A_i(d)$, order $m_i$ and power vector $v_i$ as displayed
here. The identity on these marked data identifies
the affine equations and the positive root coordinate.
In its Lie-bundle notation, the difference $b_0$ of two lifts
with aligned integral powers has zero norm, and the conjugating
translation is
\[
 h_{\rm aff}=-m_i^{-1}\sum_{j=0}^{m_i-1}j\mathcal R^j b_0,
                 \qquad (I-\mathcal R)h_{\rm aff}=b_0.
\]
Its inverse is translation by $-h_{\rm aff}$ over the entire disc.
A period change $\lambda$ changes the power by $N_{m_i}\lambda$,
with inverse subtraction. Thus these maps include every
noninvariant affine part and integral lift.

\smallskip
\emph{The original additive fillings.}
The line degree on either original fibre is the primitive
vector $C_p^*-C_0^*$. The section $6dG$ is narrow for every
$d$. Its identity-component logarithm supplies the base
isotopy; a connection lifts it to the circle bundle. The
endpoint gauge has a real logarithm because the elliptic
neighbourhood is simply connected. The reciprocal gauge,
backward flow and inverse positive-roof coordinate give the
actual product map of Theorem~\ref{maps:src:D.2.32}, now
with the logarithm multiplied by $d$. The line degree stays
$C_p^*-C_0^*$ throughout.

For $s_i$ Kodaira components the cap has $H_1=\mathbb Z\gamma$
and
\[
 H_2=\ker\bigl(c:\mathbb Z^{s_i}\to\mathbb Z\bigr),
                                      \qquad c(y)=y_p-y_0.
\]
The whole-fibre multiplicity vector $F$ and the component
vectors $C_j$ with $j\ne0,p$ form an integral basis of this
kernel. Its inverse reads $y_0$ and $y_j-m_jy_0$; the
entries of $F$ at $0,p$ are both one. The nearby elliptic
area maps to the primitive $F$. Its component complement
has rank five at $IV^*$ and zero at $III$.

The period matrix $(\delta,a,b,u_i)$ is unimodular, with
inverse subtracting $2d(a+b)$ or $3d(a+b)$ times the last
coordinate. It conjugates $A_i(d)$ to its constant
three-dimensional block plus an identity. The two nonzero
Smith factors of that first block minus identity are one,
by the displayed primitive coinvariant maps. Volume duality
gives the same factors on its exterior square. The full
fibre Wang quotient therefore has four unit relations and
basis $[ab],[a u_i]$.

For $\kappa_i=2d,3d$ respectively, its exact quotient rows are
\[
 \begin{pmatrix}0&0&0&1&0&0\\0&0&-q_i&0&1&-1\end{pmatrix}
                  \bigwedge^2(\delta,a,b,u_i)^{-1}.
\]
Their section inserts $ab,a u_i$. The first row is
$(0,0,0,1,-\kappa_i,\kappa_i)$, and its full raw kernel has basis
\[
 \delta a,\quad\delta b,\quad\delta\gamma,\quad
 a\gamma+\kappa_iab,\quad b\gamma-\kappa_iab.
\]
Reading all coordinates except $ab$ is its inverse.
Their $\xi_d$ values are $0,0,-6d,\kappa_i,-\kappa_i$.
Projection forgetting the line sends $a u_i$ to a finite
elliptic coinvariant in a free target, hence to zero.
The circle-bundle Gysin map on $H_2$ is injective, so the
same class has zero image in the actual cap. The two
fixed normal sweeps bound the origin-section tube and
the narrow forward-orbit tube with their extending gauges.
Thus these two sweeps and the primitive mixed fibre class
are the full kernel. They give exactly
\eqref{maps:eq:D.WE4}, with the second sign coming from
$t_1(\delta)=-t$. The primitive area section splits the
free component complement integrally.

\smallskip
\emph{The complete group and the two attachment orders.}
In $W$, a finite filling imposes $m_0x=n_0\gamma$ or
$-m_1x=n_1\gamma$. An original kills its section meridian
and retains $\gamma$, giving effective data $(1,0)$.
The marked loops and their inverse peripheral representatives
therefore give
\[
 \pi_1(X)=\mathbb Z\langle\gamma,x\rangle/
 \langle-\bar n_0\gamma+\bar m_0x,
                   -\bar n_1\gamma-\bar m_1x\rangle.
\]
In a column-image convention this is the cokernel of the
transpose of the matrix $R$ in \eqref{maps:eq:D.WE2}.
For reversible Smith matrices ${\mathsf U}_R R {\mathsf V}_R={\mathsf S}_R$, the map from
$\operatorname{coker}R^{\mathsf t}$ to
$\operatorname{coker}R$ is ${\mathsf U}_R^{-1}{\mathsf V}_R^{\mathsf t}$;
its inverse is ${\mathsf V}_R^{-\mathsf t}{\mathsf U}_R$. Both maps identify through
the same diagonal cokernel. This proves the displayed group
identification by explicit integral maps. Its determinant is
$P$, and its positive multiplicities give $r=1$ exactly at
$P=0$.

Put $C=\operatorname{coker}B$. To give the global integral
extension, attach the additive fillings first and then use
the relative cusp sequence. Write $M=\partial N_d$ and
$E=V\cup U_0\cup U_1$. The image of
$\bigwedge^2A_\infty-I$ has primitive basis
\[
 f_1=e_{\delta a},\qquad
 f_2=-d e_{\delta a}-d e_{\delta b}+e_{a\gamma}.
\]
A left inverse has rows
$e_{\delta a}^*+d e_{a\gamma}^*,e_{a\gamma}^*$.
The quotient and its section are
\[
 Q(y)=(y_{\delta b}+d y_{a\gamma},y_{\delta\gamma},
                                      y_{ab},y_{b\gamma}),
 \quad\sigma(z)=z_1e_{\delta b}+z_2e_{\delta\gamma}
                              +z_3e_{ab}+z_4e_{b\gamma}.
\]
They satisfy $Q\sigma=I$, and subtracting $\sigma Q$ gives
the relation coordinates by the stated left inverse.
The fixed-sweep section adds $\delta,a$, so $H_2(M)$ is
free of rank six. Its two angular sweeps form a primitive
subgroup and bound in $N_d$. Since
$H_3(N_d,M)=H^3(N_d)=\mathbb Z^2$, exactness proves that
they are the entire integral kernel of $H_2(M)\to H_2(N_d)$.

The exact degree-one presentation before filling the cusp is
\[
 H_1(E)=\mathbb Z\langle\gamma,x_0,x_1\rangle/
 \langle-\bar n_0\gamma+\bar m_0x_0,
                        -\bar n_1\gamma+\bar m_1x_1\rangle.
\]
It has free rank one, retaining its possible torsion.
The degree-one Mayer--Vietoris source has rank six, its
target rank $7-L$ and its cokernel rank one. Its free
kernel has rank $L$. Its integral map uses boundary coordinates
$(a_i,\gamma_i,x_i)$. Into $V$ these map to
$(0,\gamma,x_i)$; into a finite cap they map to
$(a,\gamma,c_i)$ modulo the displayed power relation, and
into an original to $(0,\gamma,0)$. Append the power columns
and take the integer kernel by reversible Smith operations.
Projection to the boundary coordinates is injective on the
relation coefficients, giving the full integral kernel.
Bound each basis cycle's images by integral two-chains and
join them across the boundary. These cycles give a section,
with inverse reading the connecting coordinates and
subtracting them. The angular cusp classes come
from $H_2(V)$ and have zero connecting coordinate in that
bridge. Subtract each primitive boundary lift from its cap
coordinate. The inverse includes the four exterior coordinates
and the component complement; different lifts differ by the
complete kernel columns. These maps give
\begin{equation}\hypertarget{tag-026R}{}\label{maps:eq:D.WE12}
 H_2(E)/j_*\ker(H_2(M)\to H_2(N_d))
                 \simeq C\oplus\mathbb Z^{\beta+L}.
\end{equation}

Positive cusp discs use the cut-base relation
$x_\infty=-(x_0+x_1)$. For the inward $E$-boundary
coordinates negate all $d$ relative disc coordinates;
the inverse change is the same $-I_d$. In these coordinates
every disc has boundary $x_0+x_1$: two height-one fan-ray
meridians differ by integral multiples of $\delta,a$,
which vanish in $H_1(E)$ by the explicit trace relations.
The full relative map is
\begin{equation}\hypertarget{tag-026S}{}\label{maps:eq:D.WE13}
 \mathbb Z^d\longrightarrow H_1(E),\qquad
             z\longmapsto\left(\sum_j z_j\right)(x_0+x_1).
\end{equation}
For $P\ne0$ the meridian has infinite order: in the rational
free quotient its coefficient of $\gamma$ is
$P/(\bar m_0\bar m_1)$. Its kernel has basis $e_j-e_0$.
At $P=0$ its exact order is
$h_{\rm mer}=\operatorname{lcm}(\bar m_0,\bar m_1)$, because
\[
 h_{\rm mer}(x_0+x_1)=\frac{h_{\rm mer}}{\bar m_0}
 (-\bar n_0\gamma+\bar m_0x_0)
 +\frac{h_{\rm mer}}{\bar m_1}(-\bar n_1\gamma+\bar m_1x_1).
\]
Conversely, a relation for $k_{\rm mult}(x_0+x_1)$ requires both
positive multiplicities to divide $k_{\rm mult}$. The kernel basis is
then $e_j-e_0$ together with $h_{\rm mer}e_0$; its inverse reads
$z_j$ for $j\ge1$ and $(\sum z_j)/h_{\rm mer}$. The displayed
power-disc combination bounds $h_{\rm mer}$ times the meridian.
Together with differences of component discs it gives
an integral section of the relative sequence, with inverse
reading and subtracting these relative coordinates.
Its free quotient has rank $d-1+r$. Combining this with
\eqref{maps:eq:D.WE12} gives the complete integral group
\[
 H_2(X)=C\oplus\mathbb Z^{d-1+\beta+L+r}.
\]
The boundary $\mathbb Z/d$ and finite meridian torsion
are carried through these exact pair maps. The common
inclusion of $V$ and its four quotient coordinates identify
this $C$ with the matrix computed after attaching the cusp.

For $L=0$, the first columns of \eqref{maps:eq:D.WE3} give
$t=-2dn_0s/k=3dn_1s/(2\ell)$. Clearing the positive
denominators leaves $dPs=0$. The other columns determine
$t'_0,t'_1$, with their root relations agreeing by alternation.
Thus $C$ has free rank $r$. For $L>0$, an original imposes
$2ds=0$ or $3ds=0$; its unit sweeps and the remaining block
then kill every rational coordinate. Thus $C$ is finite.
This proves the degree-two ranks in \eqref{maps:eq:D.WE5}.
The cusp has Euler number $2d$ by \eqref{maps:eq:D.WE9};
the other caps, torus bundle and boundaries have Euler
number zero. Hence $\chi(X)=2d$ and
$b_3=2b_2-2r+2-2d$, giving the remaining rank formula.

Integral duality and universal coefficients give
\eqref{maps:eq:D.WE6}. Explicit torsion coordinates are
the Smith coordinates of $B$ and the group presentation.
A character of value $1/q$ on a generator of order $q$,
followed by the integral Bockstein and oriented cap product,
gives its dual torsion class; the inverse reads its linking
coordinate. Free classes use evaluation and cap product.

\smallskip
\emph{The closed forms.}
In the coprime two-logarithm sector choose
$\alpha_i m_i+\eta_i n_i=1$. Write each block's columns
as ${\mathsf C}_1,{\mathsf C}_2,{\mathsf C}_3$ and put
${\mathsf C}'=\alpha_i{\mathsf C}_2+\eta_i{\mathsf C}_3$. Direct substitution gives
$n_i{\mathsf C}_2-m_i{\mathsf C}_3=q_i r_i {\mathsf C}_1$ and the inverse substitutions
\[
 {\mathsf C}_2=m_i{\mathsf C}'+\eta_iq_i r_i{\mathsf C}_1,\qquad
 {\mathsf C}_3=n_i{\mathsf C}'-\alpha_iq_i r_i{\mathsf C}_1.
\]
The $t'_i$ coordinate of ${\mathsf C}'$ is one. Ordering the two
${\mathsf C}_1$ columns first gives a block matrix
$\left(\begin{smallmatrix}K&R_0\\0&I\end{smallmatrix}\right)$.
The target shear with blocks $(I,-R_0;0,I)$, inverted by
$(I,R_0;0,I)$, identifies its cokernel with
\[
 K=\begin{pmatrix}-2dn_0&-3dn_1\\-k&2\ell\end{pmatrix},
                                      \qquad\det K=-dP.
\]
The quotient subtracts $R_0$ times the last coordinates;
its inverse appends zeros. This proves the removal of
the $r_i$ by integral maps.

The four integers $k,2\ell,2n_0,3n_1$ have gcd one.
For the prime two use that $n_1$ is odd; for three use
$3\nmid n_0$; every other common prime contradicts
$\gcd(k,n_0)=1$. The same divisibility argument at each
prime power gives the entry gcd $g=\gcd(d,k,2\ell)$.
The determinantal divisors are therefore $g,d|P|$.
Reversible Euclidean row and column operations give
\eqref{maps:eq:D.WE7}. The row matrix ${\mathsf U}_K$ gives the quotient
map $[y]\mapsto[{\mathsf U}_K y]$, inverted by ${\mathsf U}_K^{-1}$; the column
matrix supplies the exact relation preimages.

At the rank-one endpoint the map is especially explicit.
Reduced fractions at $P=0$ force
$m_0=m_1=12q$, $k=4q$, $\ell=3q$, and
$n_1=-n_0=-n$, with $\gcd(n,12q)=1$. Then
$K=(dn,2q)^{\mathsf t}(-2,3)$ and $g=\gcd(d,2q)$.
Put $w=(dn/g,2q/g)^{\mathsf t}$ and choose
$\alpha_{\rm Bez} w_1+\beta_{\rm Bez} w_2=1$. The unit-determinant matrix
with rows $(\alpha_{\rm Bez},\beta_{\rm Bez}),(-w_2,w_1)$ sends the image
to $g e_1$. Its coordinates and inverse give
$\mathbb Z/g\oplus\mathbb Z$.
The cyclic group coordinate, including $P=0$, is
\[
 \gamma\longmapsto m_0,\quad x\longmapsto n_0
                 \quad\text{in }\mathbb Z/P\mathbb Z,
 \qquad 1\longmapsto\alpha_0\gamma+\eta_0x.
\]
The power relations and B\'ezout identity verify both
inverse composites.

With $IV^*$ original, its unit columns kill $t,t'_0$
and impose $2ds=0$. Subtracting those columns from the
remaining block leaves independent relations
$(2d,3dn_1)$ on $s$ and $(m_1,n_1)$ on $t'_1$;
the $6dr_1s$ relation is already a multiple of $2ds$.
B\'ezout operations give the first row of
\eqref{maps:eq:D.WE8}. The quotient reads these two
coordinates and its inverse inserts them in the exterior
target with zero eliminated coordinates. With $III$
original the same stated operations leave $(3d,2dn_0)$
on $s$ and $(m_0,n_0)$ on $t'_0$, giving its second row.
With both original, the unit columns kill all sweeps
and $3d-2d=d$ gives $C=\mathbb Z/d$. These computations
retain every free remaining twist, including noncoprime
numerators.

Finally, at $d=1$ the suspension and toric covering maps
are identities, with the same additive period markings.
This is the natural row with section multipliers $1,6$ on
the $IV^*,III,I_1$ profile. For $d>1$, its Euler characteristic $2d$
is greater than the six-sphere's value two. This gives
the stated endpoint and the full positive-weight extension.
\end{proof}

\begin{remark}[Simply connected and original-cap specializations]
\hypertarget{tag-Q134}{}\label{mapsIII:rem:positive-engel-specializations}
{\normalfont\scriptsize tag Q134}\par
For two logarithmic caps, a simply connected member has $P=\pm1$.
The two-by-two presentation has a unit determinant, so each
$\gcd(m_i,n_i)$ and $\gcd(m_0,m_1)$ divides one. Hence $g=1$, and
\begin{equation}
 H_2=\mathbb Z^{d-1}\oplus\mathbb Z/d,\qquad
 H_3=\mathbb Z/d,\qquad H_4=\mathbb Z^{d-1}.
 \hypertarget{tag-026T}{}\label{maps:eq:D.WE14}
\end{equation}
These formulas include even and three-divisible $d$.
At the minimal orders $m_0=3,m_1=4$ and $\gcd(d,6)=1$,
the homology-coordinate comparison with
Theorem~\ref{maps:src:D.1.4} is
\[
 (z_1,\ldots,z_{d-1};[u]_d)
       \longmapsto(z_1,\ldots,z_{d-1};[u]_d)
\]
in the free and Smith coordinates of their respective proofs.
In degree three
use the corresponding cyclic linking coordinate, and in degree
four use the dual free coordinates. Send the unit generator in
degree zero and the complex-orientation generator in degree six
to their counterparts; in degrees one and five use the unique maps
of the zero groups.

With both original caps, \eqref{maps:eq:D.WE6} gives
\[
 \pi_1=\mathbb Z,\qquad
 H_2=\mathbb Z^{d+7}\oplus\mathbb Z/d,\qquad
 H_3=\mathbb Z^{14}\oplus\mathbb Z/d,\qquad
 H_4=\mathbb Z^{d+7},\quad H_5=\mathbb Z.
\]
At $d=1$, the identity suspension/toric maps and unchanged
original caps identify this datum with
Remark~\ref{mapsIII:rem:original-additive-specializations}
in the same marked charts. At every positive $d$ the construction
uses the stated iterated translation.
\end{remark}

\begin{remark}[Integral weighted Engel calculations]
\hypertarget{tag-Q135}{}\label{mapsIII:rem:positive-engel-calculations}
{\normalfont\scriptsize tag Q135}\par
For integer parameters satisfying \eqref{maps:eq:D.WE1}, the
Mayer--Vietoris maps in the preceding proof give the integral group
and homology presentations. The monodromy identities, primitive
exterior-square quotient and graph basis are polynomial identities
in those parameters. The unreduced graph and full area kernels give
\eqref{maps:eq:D.WE3}--\eqref{maps:eq:D.WE4}. These formulas include
nonprimitive free twists, nontrivial cap torsion, every
original/logarithmic choice and zero-Euler cases.
\end{remark}

\subsection{Weighted natural profiles and their integral cap maps}

\begin{definition}[Weighted natural data]
\hypertarget{tag-Q136}{}\label{mapsIII:def:weighted-natural-data}
{\normalfont\scriptsize tag Q136}\par
Choose one of the $172$ natural initial rows of
Theorem~\ref{maps:src:D.2.22}, with $P=aG$, $Q=bG$,
$a,b>0$, $\gcd(a,b)=1$, and weight
\begin{equation}
 d>0,\qquad\gcd(a,d)=1.
 \hypertarget{tag-027O}{}\label{maps:eq:D.WP1}
\end{equation}
Take the $d$th iterate of its contracting line map: the
translation is $dbG$, and the linearization is the norm of the
initial linearization. Use the same unimodular Mumford fan with
the resulting period subgroup and its unchanged chart units.

Let $\mathcal P$ be the marked cap patterns of
Theorem~\ref{maps:src:D.2.33}: original or finite-good fillings
at $\mathrm{II},\mathrm{III},\mathrm{IV},\mathrm{III}^*,\mathrm{IV}^*$,
original fillings at odd $I_n^*$, and the original or Mumford
choices at the specified multiplicative sites. Write
$\mathcal P_{\mathrm{all\ star\ O}}$ for the subset in which
$\mathrm{III}^*$ and $\mathrm{IV}^*$ are also original.
Allow every free finite-good logarithm at an allowed site,
every finite list of free regular logarithms, and every integer
clutch $e$.

Use the periods of Definition~\ref{mapsIII:def:elliptic-marked-data},
in the order $\Lambda=\mathbb Z\langle e_1,e_2,\delta,\gamma\rangle$.
For the initial matrices $A_i$ put
\begin{equation}
 A_i(d)=S_d^{-1}A_iS_d,\qquad S_d=\operatorname{diag}(1,1,1,d),
 \hypertarget{tag-027S}{}\label{maps:eq:D.WP2}
\end{equation}
and set
\begin{equation}
 \omega_d=a e_1^*\wedge e_2^*-db\,\delta^*\wedge\gamma^*.
 \hypertarget{tag-027T}{}\label{maps:eq:D.WP3}
\end{equation}
Each regular site has identity monodromy. Represent clutch $e$
by an auxiliary order-one regular site with power $-e\gamma$.

At a multiplicative site write its type as $I_{n_i}$.
Let $L$ count the nonzero-Chern originals. Define $\beta$ by
summing $\gcd(d,n_i)-1$ at a zero-free rank-one site with
narrow $P$, zero at a non-narrow marked Mumford filling,
$\rho_i(d)-1$ at a zero-Chern additive original, and $s_i-2$
or $n_i-2$ at a nonzero-Chern additive or multiplicative original,
respectively. Here $s_i$ is the number of Kodaira components,
and $\rho_i(d)$ is the number of component orbits under the
$d$th power of the translation by $Q$. The selected rank-two
cusp is excluded from $\beta$.
\end{definition}

\begin{theorem}[Weighted natural profiles]
\hypertarget{tag-02LD}{}\label{maps:src:D.2.37}
{\normalfont\scriptsize tag 02LD}\par
For the data of Definition~\ref{mapsIII:def:weighted-natural-data},
the set $\mathcal P$ has $495$ patterns and
$\mathcal P_{\mathrm{all\ star\ O}}$ has $459$.
Its inclusion preserves every site label; the retraction replaces
only the finite-good $\mathrm{III}^*$ and $\mathrm{IV}^*$ choices
by their original fillings. The following integral algorithm
applies to both sets and every stated parameter.

The matrices $A_i(d)$ are integral and retain their ordered
product identity. The change multiplies precisely the first
three entries of each $\gamma$ column by $d$.
The functional $\omega_d$ is primitive and invariant under them.
The selected cusp has $d$ normalized components.

For $N$ sites, ordered with
$A_0(d)\cdots A_{N-1}(d)=I$, compute a primitive integral
basis of
\begin{equation}
 K=\ker\bigl[(A_0(d)-I)\ \cdots\ (A_{N-2}(d)-I)\bigr].
 \hypertarget{tag-027U}{}\label{maps:eq:D.WP4}
\end{equation}
The degree-two presentation has target $\mathbb Zs\oplus K$.
For $\lambda\in\ker(A_i(d)-I)$, let $J_i\lambda$ be its
inclusion in block $i$ when $i<N-1$. At the last site its
block $j$ is
$-A_{j+1}(d)\cdots A_{N-2}(d)\lambda$.

At a finite-good site of linear order $h$, take a primitive
fixed basis $F_{\rm fix}:\mathbb Z^2\to\Lambda$, put
$N_h=I+\cdots+A^{h-1}$, and let $H_{\rm norm}$ be an integral basis
matrix of $\operatorname{im}(F_{\rm fix}^{-1}N_h)$.
The projected lattice is $F_{\rm fix}(H_{\rm norm}/h)\mathbb Z^2$.
For invariant integral power $v=F_{\rm fix}z$ and order $m=hk>0$, set
\begin{equation}
 \begin{gathered}
 w=hH_{\rm norm}^{-1}z\in\mathbb Z^2,\qquad t=\gcd(m,w_1,w_2),\\
 \text{the action is free}\quad\Longleftrightarrow\quad
 \gcd(m,v_1,v_2,v_3,v_4)=1,\qquad h\mid m/t .
 \end{gathered}
 \hypertarget{tag-027V}{}\label{maps:eq:D.WP5}
\end{equation}
For each free such logarithm use all five columns
\begin{equation}
 \bigl(-\omega_d(\lambda,v),\,J_iN_m\lambda\bigr)
       \quad(\lambda=e_1,e_2,\delta,\gamma),\qquad
 (0,J_iv),\qquad N_m=(m/h)N_h .
 \hypertarget{tag-027W}{}\label{maps:eq:D.WP6}
\end{equation}
At a regular logarithm the same formula uses $A=I$,
$N_m=mI$ and freeness $\gcd(m,v_1,v_2,v_3,v_4)=1$.

At each unchanged original or Mumford filling use
\begin{equation}
             (0,J_iF_i),\qquad
        F_i\text{ an integral basis of }\ker(A_i(d)-I).
 \hypertarget{tag-027X}{}\label{maps:eq:D.WP7}
\end{equation}
For a nonzero-Chern original also include $(dh_i,0)$,
where $h_i=b/q_i^c$ at an additive site and $h_i=b/n_i$
at a multiplicative site, as in
Theorems~\ref{maps:src:D.2.32} and~\ref{maps:src:D.2.33}.
These columns are the full global images of the integral
boundary kernels. Let $B$ be their matrix and put
$T=\operatorname{Tor}\operatorname{coker}B$.

Use $L$ and $\beta$ from
Definition~\ref{mapsIII:def:weighted-natural-data}.

For the logarithms other than the auxiliary clutch site put
$\mu_i=\gamma^*(v_i)$. Then
\[
 \pi_1(X)=
 \left\langle\gamma,c_1,\ldots,c_k\ \middle|\
 [\gamma,c_i]=1,\quad c_i^{m_i}=\gamma^{\mu_i},\quad
                   c_1\cdots c_k=\gamma^e\right\rangle .
\]
Write its abelianization as $\mathbb Z^r\oplus A_{1,\rm tor}$.
The Betti numbers are
\begin{equation}
 b_2=
 \begin{cases}
 d-1+\beta+2r,&L=0,\\
 d-1+\beta+L+r,&L>0,
 \end{cases}
 \qquad
 b_3=
 \begin{cases}
 2\beta+2r,&L=0,\\
 2\beta+2L,&L>0.
 \end{cases}
 \hypertarget{tag-027Y}{}\label{maps:eq:D.WP8}
\end{equation}
The complete integral homology is
\begin{equation}
 \begin{array}{c|c}
 j&H_j(X;\mathbb Z)\\ \hline
 0,6&\mathbb Z\\
 1&\mathbb Z^r\oplus A_{1,\rm tor}\\
 2&\mathbb Z^{b_2}\oplus T\\
 3&\mathbb Z^{b_3}\oplus T\\
 4&\mathbb Z^{b_2}\oplus A_{1,\rm tor}\\
 5&\mathbb Z^r .
 \end{array}
 \hypertarget{tag-027Z}{}\label{maps:eq:D.WP9}
\end{equation}
In particular $\chi(X)=2d$. The two Smith calculations,
for $B$ and for the displayed Seifert abelianization, determine
all finite cyclic factors at every allowed weight, order,
affine power and clutch.
\end{theorem}


\begin{proof}
Let \(d_0\) be the denominator of the generator height. The selected fibre has type \(I_{n_*}\), with \(ab=d_0/n_*\) and \(\gcd(a,b)=1\). Put \(\chi=\gamma^*\).

\emph{The cover and its marked period map.}

Let \(\psi\) denote the original contracting line map, on the nonzero
line bundle over the complement of the selected multiplicative fibre.
The norm linearization is the composition
\[
 \psi^d=\psi\circ\cdots\circ\psi.
\]
In a line frame with multiplier $u_{\rm line}(x)$, its multiplier is the ordered
product of $u_{\rm line}$ over the first $d$ translates. Its divisor is
\(\sum_{j=0}^{d-1}t_{-jQ}^{\,*}\operatorname{div}(u_{\rm line})\), with
all orbit repetitions included. The translation on the
elliptic base is \(-dbG\).

The inclusion of groups \(\langle\psi^d\rangle\subset\langle\psi\rangle\)
defines
\[
 p_d:[z]_{\psi^d}\longmapsto[z]_\psi.
\]
The $d$ disjoint local inverse branches are the lifts to the $d$ chosen
cosets of \(\langle\psi^d\rangle\). Freeness and properness follow from
the original contraction: a compact radial annulus meets only finitely
many of its iterates, and a nonzero iterate cannot fix a point of the
nonzero line. This proves that the covering \(p_d\) is the degree-d cover on
this exterior.

A positive suspension period at weight $d$ consists of $d$ original
suspension periods. Elliptic periods and the angular line circle are unchanged.
Consequently the injection of marked integral fibre lattices is
\[
 S_d=\operatorname{diag}(1,1,1,d).
\]
For every base loop, lift it and then apply the covering \(p_d\). The two endpoint
descriptions give
\[
 A_iS_d=S_d A_i(d),\qquad A_i(d)=S_d^{-1}A_iS_d.
\]
The inverse over a chosen sheet is the period exponential with the weight-d \(\gamma\) coordinate. Since the last row of \(A_i\) is
$(0,0,0,1)$, \eqref{maps:eq:D.WP2} multiplies precisely the first three entries in its
\(\gamma\) column by $d$ and leaves all other entries unchanged. It is
integral for every $d$.

At the selected fibre, the same map extends using the infinite
toric chart system. In the Jordan short marking of Theorem~\ref{maps:src:D.2.23} the full
period lattice is \(\mathbb Z^2\), while the weight-d lattice is
\(\mathbb Ze_1+d\mathbb Ze_2\). Take the identical unimodular fan and
the identical nonzero chart units, then divide by this subgroup. The
identity on every infinite chart descends to the covering \(p_d\). A nonzero period
translates component vertices. At a point of a normal-crossing chart
it therefore cannot fix the finite collection of branches through
that point. Local inverse sheets are obtained by choosing one period
representative and using the inverse chart coordinates. Thus the
period-subgroup construction is an extension of the cover.

For example, in a chart with the chosen \(A_2\) diagonal the map
\[
 (X,Y,q)\longmapsto(x,y,w)=(q/(XY),Y,X)
\]
has inverse \((X,Y,q)=(w,y,xyw)\). These charts use the same original units on both sides. The identity on the infinite charts keeps the fan and every chart unit fixed.

\emph{Raw matrices and complete primitive exterior maps.}

Write the raw elliptic matrix and section cocycle as \(B_i,g_i\),
and put
\[
 J=\begin{pmatrix}0&1\\-1&0\end{pmatrix},\qquad
 \beta_i^1=-g_i^tJB_i,\qquad \beta_i=a\beta_i^1.
\]
The section-line coordinates of Theorem~\ref{maps:src:D.2.23} give
\[
 A_i=
 \begin{pmatrix}
 B_i&0&b g_i\\
 a\beta_i^1&1&c_{i,\rm mat}\\
 0&0&1
 \end{pmatrix}.
\]
At \(I_\nu\), choose \(u,r_{\rm per}\) with determinant one,
\(B_i-I=\nu uu^tJ\), \(g_i=\kappa u\). Then
\[
 c_{i,\rm mat}=-ab\kappa^2/\nu+
       \begin{cases}1/\nu&i=*,\\0&i\ne *.\end{cases}
\]
At every other fibre \(B_i-I\) is invertible over \(\mathbb Q\), and
\[
 c_{i,\rm mat}=\beta_i\,b(B_i-I)^{-1}g_i.
\]
Substitution of the defining elliptic matrices and section cocycles gives
integral entries and the ordered product identity on all 172 rows.

The $172$ natural rows are the complete rows of
Theorem~\ref{maps:src:D.2.22}. Their identification here fixes
the ordered profile, selected zero, positive pair $(a,b)$ and
every component coordinate in Definition~\ref{mapsIII:def:elliptic-marked-data}.
Thus the displayed raw matrices use exactly those marked elliptic
matrices and section
cocycles, including the cases in which component terms coincide.

Write $\kappa_{\mathrm c}$ for the defining component coordinate
of $G$ at $I_\nu$. This coordinate and the positive Tate indexing
are related by the exact map
\(j\mapsto-\kappa\kappa_{\mathrm{c}}^{-1}j\), with inverse
\(j\mapsto-\kappa_{\mathrm{c}}\kappa^{-1}j\), modulo \(\nu\).
For these defining profiles the map is reflection, by the defining residue maps in Theorem~\ref{maps:src:D.2.33}.
The forward arrow covers $-Q$ and has Tate valuation
$b\kappa$; hence $P$ has Tate component \(-a\kappa\).
The displayed reflection transports both section components
and the forward arrow, fixes \(C_0\) and preserves the component
intersection matrix. It preserves the component-degree equation
after transporting its index. Period matrices and fan units
below always retain the Tate marking of the displayed section-line matrix.

\emph{Two universal primitive profile matrices.}

With the last site omitted, define
\[
 U_2=[\,B_i-I\,]_i,\qquad
 U_3=\left[\,\binom{B_i-I}{\beta_i^1}\,\right]_i.
\]
For each of the 21 defining profiles, integer maximal-minor elimination gives
matrices \(R_2,R_3\) satisfying
\[
 U_2R_2=I_2,\qquad U_3R_3=I_3.
\]
Obtain them from a Bézout combination of maximal minors. More
explicitly, for each chosen minor insert its adjugate in the
corresponding column slots. Multiplying by its signed Bézout
coefficient and adding gives the identities \(U_2R_2=I_2,\ U_3R_3=I_3\). On each of the 21 profiles the
gcd of those maximal minors is one in both degrees. Multiplication
therefore gives the displayed integer right inverses.

Let \(G_d\) be the first three rows of
\([A_i(d)-I]_{i<N-1}\). At weight one the same maximal-minor
construction gives \(G_1R_1=I_3\) on all 172 rows. Insert
\[
 R_3\operatorname{diag}(a,a,1)
\]
in the elliptic column slots of \(G_d\), and call the resulting
matrix \(C_a\). The identities \(U_2R_2=I_2,\ U_3R_3=I_3\) give \(G_dC_a=aI_3\).
Let \(L_d\) multiply each nongamma domain column by $d$ and leave each
\(\gamma\) domain column unchanged. The literal column rule \eqref{maps:eq:D.WP2} gives
\[
 G_dL_d=dG_1.
\]
Choose integers \(\alpha_{\rm Bez},\beta_{\rm Bez}\) with \(\alpha_{\rm Bez} a+\beta_{\rm Bez} d=1\).
Then
\[
 R_d=\alpha_{\rm Bez} C_a+\beta_{\rm Bez} L_dR_1,\qquad G_dR_d=I_3.
\]
This is an integer right inverse for every permitted $d$. It proves that the common
coinvariant of \(H_1(T^4)\) is exactly \(\mathbb Z\gamma\), with
quotient \(\chi\) and section \(\gamma\). In particular every integral
short class bounds in the exterior.

\emph{The complete wedge image.}

Order the wedge basis as
\[
 e_{12},e_{1\delta},e_{1\gamma},
 e_{2\delta},e_{2\gamma},\delta\gamma.
\]
Let \(E_d=[\,\bigwedge^2A_i(d)-I\,]_{i<N-1}\).
On \(e_1\delta,e_2\delta\), this map has the block \(B_i-I\).
Inserting \(R_2\) gives exact preimages of \(e_{1\delta}\) and
\(e_{2\delta}\). Call that two-column preimage matrix \(P_{\mathrm{ang}}\).

For an elliptic column \(\lambda\), the components of
\(E_d(\lambda\gamma)\), after dropping the two angular rows, are
\[
 \bigl(db\,\beta_i^1\lambda,\ 
       (B_i-I)\lambda,\ a\beta_i^1\lambda\bigr).
\]
Insert \(R_3\) in these elliptic-\(\gamma\) column slots. Using the identities \(U_2R_2=I_2,\ U_3R_3=I_3\), its
three images, modulo the two angular classes, are respectively
\[
 e_{1\gamma},\quad e_{2\gamma},\quad
 db\,e_{12}+a\,\delta\gamma.
\]
Subtract their angular coordinates times \(P_{\mathrm{ang}}\).
This gives literal integer preimages of the three mixed wedge generators, with inverse
substitution by the displayed mixed wedge coordinates. The coefficients depend linearly on $d$.

Every \(A_i(d)\) preserves
\[
 \omega_d=a e_1^*\wedge e_2^*-db\,\delta^*\wedge\gamma^*.
\]
The mixed wedge calculation proves the exact image
\[
 \operatorname{im}E_d=
 \mathbb Z\langle e_{1\delta},e_{2\delta},e_{1\gamma},
          e_{2\gamma},db\,e_{12}+a\,\delta\gamma\rangle.
\]
For the reverse inclusion, \eqref{maps:eq:D.WP3} kills every monodromy relation,
and the displayed rank-five lattice is precisely its kernel.
Indeed choose \(u_{\rm Bez},v_{\rm Bez}\in\mathbb Z\) with \(u_{\rm Bez}a+v_{\rm Bez}db=1\). Then
\[
 s_0=u_{\rm Bez} e_{12}-v_{\rm Bez}\delta\gamma,\qquad \omega_d(s_0)=1.
\]
On the kernel, the coefficients of the first four displayed
generators are their corresponding wedge coordinates, and the
coefficient of the last is \(v_{\rm Bez} x_{12}+u_{\rm Bez} x_{\delta\gamma}\).
These formulas give both inverse composites. Thus the fibre
coinvariant is exactly \(\mathbb Zs\), with quotient \(\omega_d\) and
section \(s_0\).

\emph{The constant specialization.}

Let \(\widetilde G_0\) be the common degree-one relation matrix at \(d=0\),
with the last site omitted, and let \(E_0\) be its degree-two
exterior relation matrix. These are polynomial specializations;
\(d=0\) denotes only the polynomial specialization.

The primitive profile matrix \(U_3\) and its integer right
inverse above give
\[
 \operatorname{im}\widetilde G_0
   =\mathbb Z e_1+\mathbb Z e_2+a\mathbb Z\delta.
\]
Indeed the elliptic columns equal
\(\operatorname{diag}(1,1,a)U_3\); the \(\gamma\) columns vanish and
the \(\delta\) columns are zero. Multiplication by the stated right
inverse gives every displayed generator. The image isomorphism
\((x_1,x_2,k)\mapsto x_1e_1+x_2e_2+ak\delta\) has inverse
\(x\mapsto(x_1,x_2,x_\delta/a)\) on this image. The cokernel map is
\[
 x\longmapsto(x_\gamma,[x_\delta]_a)
       \in\mathbb Z\gamma\oplus\mathbb Z/a.
\]
Its inverse sends \((n,[k]_a)\) to the class of
\(n\gamma+k\delta\); changing $k$ by $a$ changes this representative
by the displayed relation. The nonzero Smith factors are exactly
\((1,1,a)\).

The five polynomial wedge preimages above specialize to
\[
 \operatorname{im}E_0=
 \mathbb Z\langle e_{1\delta},e_{2\delta},
                 e_{1\gamma},e_{2\gamma},a\,\delta\gamma\rangle.
\]
The coordinate section and inverse are the four indicated unit
coordinates and the last coordinate divided by $a$ on the image.
Every source column has zero elliptic-area coordinate and its
\(\delta\)-\(\gamma\) coordinate is divisible by $a$, so this is also the
reverse inclusion. The nonzero Smith factors are exactly
\((1,1,1,1,a)\).

\emph{The integral graph splitting.}

Retract the punctured base onto the ordered bouquet of $N-1$ loops.
The torus bundle chain complex in degrees one and two gives
\[
 K=\ker[\,A_i(d)-I\,]_{i<N-1}.
\]
The last coordinate of the displayed map is zero, so the formula for \(R_d\) is a
right inverse of its three nonzero rows. The integer projector
\[
 P_d=I-R_dG_d
\]
has image $K$. A Hermite basis of that image gives an integral
basis of $K$. If \(K_{\mathrm{bas}}\) is its basis matrix, the square matrix
\([R_d\ K_{\mathrm{bas}}]\) is unimodular: \(G_d\) reads the first three inverse
coordinates and the remaining inverse rows read kernel coordinates.
This is also an exact algorithm for all integral insertions below.

We use the integral exterior fibre class of Theorem~\ref{maps:src:D.2.31}. Its integral homology splitting defines a
homomorphism \(H_2\to\mathbb Z\). Since \(H_1\) of that exterior is free,
the integral universal coefficient map
\(H^2(V_1;\mathbb Z)\to\operatorname{Hom}(H_2(V_1),\mathbb Z)\)
is an isomorphism, giving the class \(\xi_1\).
It evaluates as \(\omega_1\) on a fibre and zero on
the marked fixed-normal sweeps, with the geometric sections and
origin-disc meridians of Theorem~\ref{maps:src:D.2.31} and Theorem~\ref{maps:src:D.2.33}. Pull it back by the
covering \(p_d\). This gives an integral cohomology class
\(\xi_d=p_d^*\xi_1\), with fibre restriction \(\omega_d\). A marked
weight-d fixed sweep maps to the weight-one sweep with period
\(S_d\lambda\); hence its \(\xi_d\) value is zero. This establishes
integrality and the geometric normalization simultaneously.

The edge projection and the marked fibre class give
\begin{equation}
 H_1(V_d)=\mathbb Z\gamma\oplus\mathbb Z^{N-1},\qquad
 H_2(V_d)=\mathbb Zs\oplus K.
 \hypertarget{tag-027P}{}\label{maps:eq:D.WP10}
\end{equation}

The corresponding map
\[
 H_2(V_d;\mathbb Z)\longrightarrow\mathbb Zs\oplus K,
 \qquad z\longmapsto(\langle\xi_d,z\rangle,\operatorname{edge}(z))
\]
is an isomorphism. To construct its inverse, realize each integral
basis vector of $K$ by the corresponding loop cylinders and the
lattice-addition triangles at their common vertex. Subtract its
integer \(\xi_d\) value times \(s_0\). Together with \(s_0\) these classes
give an inverse. The edge projection reads the $K$ coordinate;
\(\xi_d\) reads the remaining fibre coefficient. The monodromy-cylinder
preimages in the displayed wedge-image basis identify every change of representatives,
so the two composites are the identity.

For a fixed vector \(\lambda\) at site \(i<N-1\), its peripheral
insertion is the block inclusion. At the last site it is
\[
 (J_{N-1}\lambda)_j
       =-A_{j+1}(d)\cdots A_{N-2}(d)\lambda.
\]
Cut the last peripheral path along the bouquet in reverse order;
this gives the last peripheral insertion formula directly. Multiplication by the graph boundary
telescopes to zero because \(\lambda\) is fixed at the last site
and the ordered product is identity. Reading the loop strips
gives its inverse. In the splitting \eqref{maps:eq:D.WP10}, every marked fixed
normal torus has image \((0,J_i\lambda)\).

\emph{The weighted cusp and its integral boundary.}

At the selected fibre \(I_{n_*}\), use
\[
 {\mathsf C}=\begin{pmatrix}
 n_*&b\kappa\\
 -a\kappa&(1-ab\kappa^2)/n_*
 \end{pmatrix},\qquad \det {\mathsf C}=1.
\]
The short lattice is \(S=\mathbb Z\langle u,\delta\rangle\).
The period subgroup at weight $d$ in this fan marking is generated
by the first original period and $d$ times the second. The
normalization of its cusp \(N_d\) therefore has $d$ toric degree-six
del Pezzo components, with the same infinite gluing charts.

The normalization complex of the periodic triangular cusp applies to this subgroup.
Over \(\mathcal R_d=\mathbb Z[t]/(t^d-1)\) its edge-to-curve matrix is
\[
 \begin{pmatrix}
 -1&-1&-1&-1\\
  1& 1& 1&t\\
 -t&-t&-1&-t
 \end{pmatrix}.
\]
Identify $\mathcal R_d$ with the coefficient ring $\mathcal A_d$
in Theorem~\ref{maps:src:D.2.34} by $[p(t)]\mapsto[p(t_{\rm cyc})]$, with inverse $[p(t_{\rm cyc})]\mapsto[p(t)]$.
On the normalization complex, send each marked component,
branch and dual-torus cell with coefficient $t^j$ to the cell
with that same index and coefficient $t_{\rm cyc}^j$. In the displayed ordered
bases the degree-two matrix agrees under this coefficient substitution; the degree-zero row
is the same dual-torus complex and the degree-four row is the
same group ring. This chain isomorphism carries the unit and
$t-1$ reductions,
their partial-sum inverse and the free extension section proved
there to this complex. Consequently
\[
 H_*(N_d;\mathbb Z)=
  (\mathbb Z,\mathbb Z^2,\mathbb Z^{d+3},
       \mathbb Z^2,\mathbb Z^d,0,0).
\]
Write \(S_{\mathrm{bas}}=(u,\delta)\) and put \({\mathsf L}=(r_{\rm per},\gamma)\). The ordered integral basis
\[
 (S_{\mathrm{bas}}{\mathsf C}e_1,{\mathsf L} e_1,S_{\mathrm{bas}}{\mathsf C}e_2,{\mathsf L} e_2)
\]
has determinant one. Its inverse uses \({\mathsf C}^{-1}\) on the short
coordinates and reads the long coordinates. In this basis the
weight-d return is
\[
                         T_d=I+E_{12}+d E_{34}.
\]
This is a coordinate map on the period subgroup; the
same fan and its chart units remain in place.

In wedge order \((12,13,14,23,24,34)\), the complete relation
lattice has the primitive basis
\[
 r_1=e_{13},\qquad r_2=e_{14}+d e_{23}.
\]
The respective preimages under \(\bigwedge^2T_d-I\) are
\(e_{23}\) and \(e_{24}-d e_{23}\).
The quotient, section and relation-coordinate inverse are
\[
 \begin{aligned}
 {\mathsf Q}(y)&=(y_{12},\,y_{23}-d y_{14},\,y_{24},\,y_{34}),\\
 \sigma(x)&=x_1e_{12}+x_2e_{23}+x_3e_{24}+x_4e_{34},\\
 y-\sigma {\mathsf Q}(y)&=y_{13}r_1+y_{14}r_2 .
 \end{aligned}
\]
Hence \({\mathsf Q}\sigma=I\), and the image of the wedge relation matrix
is exactly \(\ker {\mathsf Q}\). Substitution verifies these polynomial identities for every $d$.

The Wang sequence adds the two fixed-loop sweeps \(e_1,e_3\).
They have an integral section given by the corresponding fixed
loops of the linear torus return. Thus
\[
 H_2(\partial N_d;\mathbb Z)=\mathbb Z^4\oplus\mathbb Z^2
                          =\mathbb Z^6.
\]
In degree one, the relation lattice is generated by \(e_1,d e_3\),
so
\[
 H_1(\partial N_d;\mathbb Z)
   =\mathbb Z\langle e_2,e_4,\mu\rangle\oplus
      (\mathbb Z/d)\langle e_3\rangle.
\]
These two statements are compatible through their displayed
Wang quotient maps.

The two angular sweeps form a primitive subgroup of \(H_2\), since
their projection to the fixed lattice has the displayed inverse.
Their thimbles extend in the cusp. Poincaré--Lefschetz duality
and the cusp groups give
\[
 H_3(N_d,\partial N_d)=H^3(N_d)=\mathbb Z^2.
\]
The full boundary kernel has rank at most two and contains that
primitive subgroup. Its quotient inside the free abelian group obtained
by removing that subgroup has rank zero, so is zero. The angular
sweeps are therefore the entire integral kernel. In particular the
rank-four fibre coinvariant injects into \(H_2\) of the cusp.

Put \(W_d=V_d\cup N_d\). Excision and Poincaré--Lefschetz duality
give the segment
\[
 \mathbb Z^2=H_3(N_d,\partial N_d)\longrightarrow
 H_2(V_d)\longrightarrow H_2(W_d)\longrightarrow
 \mathbb Z^d\longrightarrow H_1(V_d).
\]
The $d$ normal component discs all have boundary equal to the
positive base meridian plus an integral short displacement.
The formula for \(R_d\) bounds that short displacement in \(V_d\). The
base meridian is primitive there. Thus the last map in the cusp relative sequence
is the sum of the $d$ coordinates times that meridian.

The two ordinary short thimbles at weight one lift to the d-period
cover: their fundamental loops lie in $S$ and have \(\gamma\) coordinate
zero. Applying ${\mathsf C}^{-1}$ to the fan short thimbles gives the physical
short directions $u$ and \(\delta\); multiplication by ${\mathsf C}$ is the
inverse coordinate map.
Their images in \(H_2(V_d)\) are
\[
 L_c=\mathbb Z(0,J_*u)+\mathbb Z(0,J_*\delta).
\]
This is a primitive rank-two sublattice. The short lattice $S$ is
primitive in \(\Lambda\); for a nonlast site \(J_*\) is a block inclusion,
and for a last site its final block is minus the identity on $S$.
Thus the corresponding projection has an integral left inverse.

The entire relative image in the cusp relative sequence has rank at most two and contains
\(L_c\). Since \(H_2(V_d)/L_c\) is free, its additional image is a
rank-zero subgroup of a free abelian group and is zero. Therefore its
image is exactly \(L_c\), including all possible boundary torsion.

The kernel of the last map is the zero-sum lattice with basis
\(e_j-e_0\), \(1\le j<d\). Take the corresponding differences
of normal discs and bound their short boundary displacements
in \(V_d\) by the formula for \(R_d\). The resulting cycles \(\Sigma_j\) have these
relative coordinates. They split the cusp relative sequence, with inverse reading
the component-disc coordinates and subtracting the \(\Sigma_j\).
We have proved the integral isomorphism
\[
 H_2(W_d)=\bigl(H_2(V_d)/L_c\bigr)
                    \oplus\mathbb Z^{d-1}.
\]
Every other boundary has zero relative component-disc coordinate,
so its degree-two map lands in the first summand. The displayed
relative maps account for the boundary \(\mathbb Z/d\).

The analogous relative degree-one sequence kills exactly the
primitive cusp base meridian. Thus
\(H_1(W_d)=\mathbb Z^{N-1}\). This assertion concerns homology.
The two OS46/II exceptional nonabelian vertical quotients are
handled by their kept original star cap in the fundamental-group calculation below.

\emph{Component powers and unchanged component defects.}

Let $q$ be the order of $G$ in the component group of a zero-Chern
original additive site. Narrow $P$ gives \(q\mid a\). Since
\(\gcd(a,b)=1\), $b$ is a unit modulo $q$; and \(\gcd(a,d)=1\) implies
\(\gcd(d,q)=1\). Let \(\sigma_{\rm perm}\) be the component permutation
of the forward arrow \(-Q\). Its order is $q$, with each component
orbit length dividing $q$. Choose a positive $E_{\rm perm}$ such that
\(E_{\rm perm}d\equiv1\pmod q\). Then
\[
 (\sigma_{\rm perm}^d)^{E_{\rm perm}}=\sigma_{\rm perm}.
\]
Every \(\sigma_{\rm perm}\)-orbit and \(\sigma_{\rm perm}^d\)-orbit are literally the same set.
Their index parametrizations are related by multiplication by $d$
and the inverse multiplication by $E_{\rm perm}$ modulo the orbit length.

This is also an exact comparison of integral component lattices:
the invariant orbit sums coincide, and
\[
 \begin{aligned}
 \sigma_{\rm perm}^d-I
   &=(\sigma_{\rm perm}-I)(I+\cdots+\sigma_{\rm perm}^{d-1}),\\
 \sigma_{\rm perm}-I
   &=(\sigma_{\rm perm}^d-I)(I+\sigma_{\rm perm}^d+\cdots+\sigma_{\rm perm}^{(E_{\rm perm}-1)d}).
 \end{aligned}
\]
Thus their coinvariant relation subgroups coincide. The local mapping torus still uses \(-dQ\); the two component-permutation identities identify its
component invariants and quotients.

At a multiplicative site \(I_\nu\), the raw generator residue
\(\kappa_{\mathrm{c}}\) is a unit modulo \(\nu\). Narrow $P$ gives
\(\nu\mid a\), hence \(\gcd(d,\nu)=1\). The translation in the
defining component indexing is by \(-db\kappa_{\mathrm{c}}\);
the reflection in the component-coordinate comparison carries it to the positive
Tate increment \(db\kappa\).
The number of its orbits is \(\gcd(d,\nu)\), directly by the cyclic
translation map and its Bézout inverse on each orbit. It equals
one throughout the present coprime-a scope. This proves the
printed correction \(\gcd(d,\nu)-1\) and its vanishing here.

At a nonzero-Chern original site the Chern degree vector of the
unchanged line bundle is unchanged. The section $Q$ is narrow, so $dQ$ is narrow
as well. Its component correction is consequently the same
\(s_i-2\) at additive sites and \(\nu_i-2\) at multiplicative
sites, by the component generators and inverse maps
described in the cap constructions.

In particular every \(\beta\) contribution in this theorem equals
its weight-one value, through the explicit lattice maps above.
Raising the component permutation to its $d$th power gives these
orbit invariants and their inverse indexing maps.

\emph{The entire weighted original and Mumford cap maps.}

\emph{Additive lattices and odd-star fillings.}

At every additive site write
\[
 z=(B_{\rm ell}-I)^{-1}g,\qquad \beta_{\rm row}=-ag^tJB_{\rm ell},\qquad z_Q=bz.
\]
Its weighted matrix is
\[
 A(d)=\begin{pmatrix}B_{\rm ell}&0&db g\\
                  \beta_{\rm row}&1&d\beta_{\rm row} z_Q\\0&0&1\end{pmatrix}.
\]
Let \(g_0\) be the least denominator of \(z_Q\). The raw component
maps give \(g_0\in\{1,q\}\subset\{1,2,3,4\}\), and \(g_0\mid a\).
Therefore \(\gcd(d,g_0)=1\). In particular \(d\equiv\epsilon\pmod{g_0}\)
for an \(\epsilon\in\{1,-1\}\). The integer matrix with first
three columns those of the identity and last column
\[
 ((1-\epsilon d)z_Q,0,\epsilon)
\]
will be denoted \(H_d\). Its determinant is \(\epsilon\), its inverse
subtracts \(\epsilon\) times the elliptic last-column entry and
multiplies the final coordinate by \(\epsilon\), and direct
multiplication gives
\[
                 A(d)H_d=H_dA(1).
\]
The elliptic identity is
\((B_{\rm ell}-I)(1-\epsilon d)z_Q=(1-\epsilon d)bg\);
the angular identity is the same equality after applying \(\beta_{\rm row}\).
Thus both the first and second exterior boundary lattices
have exactly the weight-one Smith factors. The weight-one vector and wedge maps of Theorem~\ref{maps:src:D.2.31},
Theorem~\ref{maps:src:D.2.29} and Theorem~\ref{maps:src:D.2.30} have unit nonzero Smith factors.
This is an integer boundary-lattice isomorphism; the local
geometric descriptions below retain the arrow \(-dQ\).

The fixed lattice and its inverse coordinate map are
\[
 F_d=\bigl(\delta,\ell_d\bigr),\qquad
 \ell_d=(-g_0dbz_1,-g_0dbz_2,0,g_0),\qquad
 x\longmapsto(x_\delta,x_\gamma/g_0).
\]
The first coordinate equation for a fixed vector forces its
elliptic part to be \(-x_\gamma dbz\). The denominator remains
\(g_0\) because $d$ is coprime to \(g_0\). This proves all of the displayed fixed basis and inverse,
including saturation.

\emph{Zero-Chern original additive caps.}

Use the zero-degree component line frames of Theorem~\ref{maps:src:D.2.31}. Use their node identifications and linearized translation units. The d-fold
composition multiplies their constants along the
component orbit. The same logarithm and line-gauge maps used
there, now with their d-fold endpoint composition, identify
the cap with the angular circle times the forward mapping
torus of \(t_{-dQ}\) on the Kodaira neighbourhood.
The inverse reverses those gauges and that mapping-torus map.
The positive contraction roof is the sum of the original
roofs on the $d$ translates; the positive-roof map and its inverse of
Theorem~\ref{maps:src:D.2.31} apply to that function.

Let \(K_0\subset N_{\rm loc}\) be the central Kodaira fibre in its plumbing
neighbourhood. Its inclusion is an integral homology equivalence
and commutes with the translation. The induced map of mapping
tori is therefore a map of their Wang complexes
\(\operatorname{Cone}(t_*-I)\), inducing an isomorphism on homology.
After product with the angular circle, its degree-two map sends
the component-orbit cycles and angular-return torus to the same
cycles in the cap. Their inverse coordinates are given below.

Let \(U_j^{\rm orb}\) be its component-orbit basis, \(\rho\) its number
of orbits, and \(F_{\mathrm{ell}}\) the whole-fibre multiplicity
vector. The component-orbit calculation proves that these orbits are the same sets
as at weight one. The component-group order $q$ divides $a$.
The orbit of the outer component met by $O$ has $q$
multiplicity-one components. Consequently
\[
 W=F_{\mathrm{ell}}/q=(1,w_2,\ldots,w_\rho)
\]
in the orbit basis is integral and primitive. The complete
cap degree-two group is
\[
 H_2(U)=\mathbb Z^\rho\langle U_j^{\rm orb}\rangle
                      \oplus\mathbb Z\langle D_B\rangle,
 \qquad D_B=\delta\times\gamma_B.
\]
The forward basis change replaces \(U_1^{\rm orb}\) by $W$. Its inverse
reads \(x_1\), then \(x_j-w_jx_1\). Thus its component
complement is exactly \(\mathbb Z^{\rho-1}\).

The integral line-frame map \(S_P\) of Theorem~\ref{maps:src:D.2.31} is unchanged.
In its coordinates, the fibre map, in wedge order
\((12,13,14,23,24,34)\), is
\[
 f_d(x)=q(x_{12}+dbz_2x_{14}-dbz_1x_{24})W+x_{34}D_B.
\]
In the original marking apply the fibre formula \(f_d\) to
\(\bigwedge^2 S_Px\). Its coefficients have gcd one because
\[
 \gcd(q,qdbz_1,qdbz_2)=1.
\]
Choose integers \(A_{\rm Bez},B_{\rm Bez},C_{\rm Bez}\) with
\(qA_{\rm Bez}+qdbz_2B_{\rm Bez}-qdbz_1C_{\rm Bez}=1\). An inverse section sends $W$ to
\[
 \bigwedge^2S_P^{-1}
       (A_{\rm Bez} e_{12}+B_{\rm Bez}e_{14}+C_{\rm Bez}e_{24})
\]
and \(D_B\) to \(\bigwedge^2S_P^{-1}e_{34}\). Substitution gives
both inverse composites modulo the four primitive wedge
relations supplied by the conjugacy \(A(d)H_d=H_dA(1)\).

The fixed sweeps of \(\delta\) and \(\ell_d\) bound respectively the
extending line-circle tube and the $q$-fold orbit tube closed
by the identity-component translation logarithm. Its
elliptic displacement is \(-qdbz\), its positive return
count is $q$, and its angular unit is the extending
line-frame unit. Reversing the gauges and flow carries these
chains to the given germ. Thus the entire degree-two kernel
is the full fixed-normal lattice, and its global image is
\[
                         (0,J_iF_d).
\]
The cap \(H_1\) has rank two, the boundary \(H_1\) has rank three,
and the remaining component cokernel is the free abelian group
\(\mathbb Z^{\rho-1}\). The component decomposition and fibre map prove these
maps integrally, including the transitive III and IV cases.

For \(\mathrm{III}^*\), \(\mathrm{IV}^*\), \(I_1^*\), \(I_3^*\), the section permutations are
determined by the regular action on the outer components
and the incidence graph. They exchange the two long arms
of \(\mathrm{III}^*\), cycle the three arms of \(\mathrm{IV}^*\), and cycle the four
outer vertices while reversing the central chain of an
odd star. Multiplicity preservation determines all
remaining coefficients. The component-orbit calculation supplies the inverse
exponent on each orbit. Raising these permutations to the $d$th power and using the inverse
exponent gives their component bases.

\emph{Nonzero-Chern original additive caps.}

Now $q$ divides $b$ and \(\gcd(a,q)=1\). The line Chern vector on
the Kodaira neighbourhood $N_{\rm loc}$ is
\(c_{\rm Ch}=C_p^*-C_0^*\), with \(p\ne0\), and both indicated
components have multiplicity one. Since $dQ$ is narrow,
the original logarithm of $Q$, multiplied by $d$, is a logarithm
of that section. Use its lifted flow and the angular
gauge of Theorem~\ref{maps:src:D.2.32}, and compose its endpoint isotopy
$d$ times. If \(f_t\) is that isotopy, \(f_t^d\) joins the
identity to the d-fold angular arrow. The forward map
\[
 [x,t]\longmapsto[(f_t^d)^{-1}x,t]
\]
and inverse \([y,t]\mapsto[f_t^d y,t]\), with collars at
the endpoints, identify the quotient with
\(E(M)\times S^1_{\ell_d}\). The variable roof is treated
by the same positive-roof map and its inverse.

The Gysin projection gives
\[
 H_1(U)=\mathbb Z,\qquad
 H_2(U)=\ker(c_{\rm Ch}:\mathbb Z^{s_i}\to\mathbb Z).
\]
For the whole fibre vector \(F_{\rm fib}=(m_j)\), the basis
\(F_{\rm fib},C_j\ (j\ne0,p)\) has inverse
\[
 x\longmapsto \bigl(x_0,(x_j-m_jx_0)_{j\ne0,p}\bigr).
\]
Hence the fibre image is primitive and the free cokernel
has rank \(s_i-2\), including zero at type III.

Here \(dbz\) is integral and \(g_0=1\). Use the unimodular
period matrix \(\mathcal B_d=(e_1,e_2,\delta,\ell_d)\), whose inverse
adds \(dbz\) times the last coordinate to the elliptic part.
It turns the displayed additive matrix into the product of
\[
 T_3=\begin{pmatrix}B_{\rm ell}&0\\\beta_{\rm row}&1\end{pmatrix}
\]
and the identity on \(\ell_d\). The primitive coinvariant row is
\[
 \rho_{\rm cov}=q(az_2,-az_1,1).
\]
Choose \(v_0=(x,y,z_\delta)\) with \(\rho_{\rm cov} v_0=1\), put
\(\beta_{0,\rm add}=\mathcal B_d(v_0,0)\), and \(k_{d,\rm add}=\beta_{0,\rm add}\wedge\ell_d\).
The two fibre Wang coordinates are ordinary elliptic area
and the \(\rho_{\rm cov}\) coordinate of the three mixed products with
\(\ell_d\). Their integral section is \(e_{12},k_{d,\rm add}\), and their
kernel is exactly the four unit wedge relations from the conjugacy \(A(d)H_d=H_dA(1)\).

The original cap sends \(e_{12}\) to $F_{\rm fib}$ and \(k_{d,\rm add}\) to zero. The latter
claim follows by forgetting the angular circle: the mixed
class then lies in the finite group
\(\operatorname{coker}(B_{\rm ell}-I)\), while the neighbourhood
mapping torus has torsion-free \(H_2\), and its Gysin projection
is injective on \(H_2(U)\). This uses the forgetful map.
The literal global coefficient is
\[
 \omega_d(k_{d,\rm add})
   =-db(az_2x-az_1y+z_\delta)=-db/q.
\]
Both marked fixed sweeps bound the line tube and the
closed logarithmic orbit tube. Reversing the sign of \(k_{d,\rm add}\)
therefore gives the complete global kernel columns
\[
                (0,J_i\delta),\quad(0,J_i\ell_d),
                    \quad(db/q,0).
\]
Their inverse generator change is the same sign reversal.
The boundary \(H_1\) has rank three and the cap \(H_1\) rank one;
this is the additional degree-one bridge recorded later.

\emph{Multiplicative caps with nonzero Chern class.}

At an unselected \(I_\nu\), take the Tate marking
\((u,r_{\rm per},\delta,\gamma)\), \(\det(u,r_{\rm per})=1\), and \(g=\kappa u\).
The matrix is
\[
 A(d)=
 \begin{pmatrix}
 1&\nu&0&db\kappa\\
 0&1&0&0\\
 0&-a\kappa&1&-adb\kappa^2/\nu\\
 0&0&0&1
 \end{pmatrix}.
\]
The nonzero-Chern case has \(\nu\mid b\) and
\(\gcd(a \kappa,\nu)=1\). Put
\[
 \begin{gathered}
 \ell_d=\gamma-(db\kappa/\nu)r_{\rm per},\quad
 p_0=\nu u-a\kappa\delta,\\
 \beta_{0,\rm mult}=xu+z\delta,\quad a\kappa x+\nu z=1 .
 \end{gathered}
\]
The basis \((p_0,r_{\rm per},\beta_{0,\rm mult},\ell_d)\) is unimodular and
takes the displayed multiplicative matrix to \(I+E_{12}\). Its inverse first subtracts
the indicated long shear, then applies the inverse
two-by-two short matrix of determinant one. Hence the
local boundary images are primitive for every $d$.

The original cap is \(E(M)\times S^1_{\ell_d}\).
For this statement the angular endpoint is the d-fold
composition of the lifted holomorphic flow of Theorem~\ref{maps:src:D.2.33},
with its prescribed central unit-circle normalization.
On each normalized \(\mathbb P^1\) the difference from the
line arrow is a nonzero holomorphic scalar; agreement
through the node constants makes it constant
around the cycle. Its phase has zero graph-loop winding.
The logarithm and reciprocal gauge of Theorem~\ref{maps:src:D.2.33} therefore
extend to the d-fold endpoint. The displayed mapping-torus
map in the additive product construction, with this isotopy, is the
product map and has the displayed inverse.

The Tate component of $P$ is \(-a\kappa\bmod\nu\),
so the component Chern vector is \(C_{-a\kappa}^*-C_0^*\)
and is unchanged. Its kernel
basis and inverse are the displayed component-kernel basis with all multiplicities one.
The component complement has rank \(\nu-2\). The fibre Wang
quotient has the four integral coordinates area, rim,
the class \(k_{d,\rm mult}=\beta_{0,\rm mult}\ell_d\), and the graph-times-long
product; the same unimodular basis above supplies its
section. The cap sends area to the whole fibre, rim to
the circle over the graph loop, the long product to
the graph-times-suspension torus, and \(k_{d,\rm mult}\) to zero.
The last claim follows from the circle-bundle
projection and its Gysin inverse in Theorem~\ref{maps:src:D.2.33}.

The coefficient of this primitive kernel is
\[
                  \omega_d(k_{d,\rm mult})=-db/\nu.
\]
The marked geometric normal torus of $u$ can have a rim
image \(j[r_{\rm per}\wedge\delta]\). Its projection is the bounding
vanishing-loop normal torus and its long coordinate
is zero. Subtracting \(j[r_{\rm per}\wedge\delta]\) gives the bounding
normal lift; the inverse adds it back. The rim has
\(\omega_d\) value zero. Therefore the global image of the
entire normal kernel is still \((0,J_i\operatorname{Fix}A(d))\).
This is the precise marked geometric section scope:
an arbitrary Wang section can also add area and long
product classes, and their inverse section changes must include those additional shears.

Thus the complete global kernel consists of all three
fixed-normal columns and \((db/\nu,0)\), its free cap
cokernel has rank \(\nu-2\), its boundary \(H_1\) rank four and
its original cap \(H_1\) rank two. Its primitive Mumford
alternative instead maps all four fibre coordinates
isomorphically to its four cap classes, has only the
three fixed-normal kernel columns, and has \(H_1\) rank
three. These are the distinct original/Mumford choices
at the 120 nonzero-Chern positions.

\emph{Zero-Chern multiplicative caps and all rank-one models.}

For narrow $P$, \(\nu\mid a\) and the physical short vector
\[
 u'=u-(a\kappa/\nu)\delta
\]
is integral and primitive. The rank-one return is
\[
 A(d)-I=u'\otimes(\nu r_{\rm per}^*+db\kappa\gamma^*).
\]
Choose integers $x,y$ with \(\nu x+db\kappa y=1\). Then
\[
 r_{\rm per}'=xr_{\rm per}+y\gamma,\qquad
 \ell'=-db\kappa r_{\rm per}+\nu\gamma
\]
and \((u',r_{\rm per}',\delta,\ell')\) form an integral basis with
return \(I+E_{12}\). The inverse long matrix has entries
\(\nu,db\kappa,-y,x\); the inverse short map restores
\((a\kappa/\nu)\delta\). These are period-coordinate
maps, with the same nonzero Tate and line units.

Use the zero-degree component frames and their node units.
The node charts and their inverses in
Theorem~\ref{maps:src:D.2.33} carry this period basis and its units.
For both the original and the marked Mumford filling, the fibre
map in the basis $(u',r_{\rm per}',\delta,\ell')$ has coordinates
\[
 {\mathsf Q}_d(y)=(y_{12},y_{23},y_{24},y_{34}),\qquad
 \sigma_d(z)=z_1e_{12}+z_2e_{23}+z_3e_{24}+z_4e_{34}.
\]
They are the ruled fibre, two gluing-circle products and base
torus. Cutting the gluing and reading these coordinates gives
the cap-coordinate inverse, and ${\mathsf Q}_d\sigma_d=I$.
The three fixed-normal sweeps bound in both fillings.
For their full boundary maps $f_O,f_M$, let $\sigma_O,\sigma_M$
be the resulting boundary sections after reading each cap's
four coordinates. Then $f_M\sigma_O$ and $f_O\sigma_M$ are
inverse target isomorphisms and commute with the common marked
boundary map. This proves their marked kernel and primitive-image
agreement. Their local free cokernel is
$\gcd(d,\nu)-1=0$ in the present scope.

For non-narrow $P$ the primitive vector is instead
\(p_0=\nu u-a\kappa\delta\), while the long covector
\(r_{\rm per}^*+(db\kappa/\nu)\gamma^*\) is primitive. The
unimodular basis in the multiplicative period construction gives the same node
charts for its marked rank-one Mumford model. It has
one normalized component, four primitive degree-two
coordinates and zero local free cokernel. This proves the rank-one component contribution using its period lattice and fan.
The two factorizations cover all 436 zero-free rank-one positions:
316 have narrow $P$ and 120 have non-narrow $P$. In the narrow case
the long covector has vertical content \(\nu\) and \(\gamma\) entry coprime
to \(\nu\); in the non-narrow case its vertical content is one.
These statements follow directly from the displayed primitive factors.

\emph{Every finite-good affine logarithm and its full kernel.}

This section applies to II, III, IV, \(\mathrm{III}^*\), \(\mathrm{IV}^*\), whose linear
orders $h$ are respectively $6,4,3,4,3$. The odd \(I_n^*\) caps remain
original. The all-star-original subset is described in the enumeration.

Let \(A=A_i(d)\), let ${F_{\rm fix}}$ be any primitive integral basis of
\(\Lambda^A\), and let \(N_h=I+A+\cdots+A^{h-1}\).
Its image lies in \(\Lambda^A\). Let $H_{\rm norm}$ be an integral basis
matrix of the image of \({F_{\rm fix}}^{-1}N_h\), and put
\[
 {\mathsf Q_{\rm norm}}=H_{\rm norm}^{-1}{F_{\rm fix}}^{-1}N_h:\Lambda\longrightarrow\mathbb Z^2.
\]
The definition of $H_{\rm norm}$ makes ${\mathsf Q_{\rm norm}}$ surjective. A Bézout sum of its
maximal-minor adjugates gives an integral section \(R_\pi\)
with \({\mathsf Q_{\rm norm}}R_\pi=I\).
Because \(N_h{F_{\rm fix}}=h{F_{\rm fix}}\), the matrix \(hH_{\rm norm}^{-1}\) is integral.
The fixed-plane projection and its projected lattice are
\[
 {\mathsf P_{\rm fix}}=N_h/h,\qquad
 {\mathsf P_{\rm fix}}\Lambda={F_{\rm fix}}(H_{\rm norm}/h)\mathbb Z^2.
\]
Thus the displayed projection, the ${\mathsf Q_{\rm norm}}$ coordinates, and \(R_\pi\)
give both inverse maps for the projected lattice.

For an invariant integral power \(v={F_{\rm fix}}z\) and \(m=hk>0\), put
\[
             w=hH_{\rm norm}^{-1}z={\mathsf Q_{\rm norm}}v,\qquad t=\gcd(m,w_1,w_2).
\]
Every affine action with that power is conjugate to
\(f(x)=Ax+v/m\) on the central real torus. Indeed the
difference of its translation and $v/m$ lies in \(\ker {\mathsf P_{\rm fix}}\), and
translation by \((I-A)^{-1}\) of that difference gives the
conjugacy; translation by its negative gives its inverse.

Use the finite-operator identity proved in
Theorem~\ref{maps:src:D.2.29} on the good-reduction root disc.
Here its operator is the finite deck operator $R$ on holomorphic
Lie sections, with $R^m=I$, and its zero-norm section is the
difference $b_0$ of two lifts with the same integral power.
Identity on $R,m,b_0$ and the root coordinate identifies the
equations and their inverse translation. Explicitly,
\[
 h_0=-\frac1m\sum_{j=0}^{m-1}jR^j b_0,\qquad
 (I-R)h_0=b_0.
\]
Translation by $h_0$, inverted by translation by $-h_0$,
extends across the centre. A period change $\lambda$ changes
the power by $N_m\lambda$, with inverse subtraction.

\emph{Every-power freeness.}

The moving plane of $A$ has a primitive $h$th-root pair as
its eigenvalues. Consequently \(A^j-I\) is invertible on
\(\ker {\mathsf P_{\rm fix}}\) whenever $h$ does not divide $j$. Its fixed-point
equation is soluble exactly when
\[
                  jv/m\in {\mathsf P_{\rm fix}}\Lambda.
\]
For necessity apply ${\mathsf P_{\rm fix}}$ to the equation. For sufficiency choose
\(\lambda=R_\pi(jw/m)\) in \(\Lambda\), and solve
\[
 x=(A^j-I)^{-1}(\lambda-jv/m)
\]
on \(\ker {\mathsf P_{\rm fix}}\). The chosen \(\lambda\) has the required projection by
the projection and its section; substitution proves the equation. Thus the
obstruction indices are exactly the positive multiples
of $m/t$.

Let \(g_{\rm pow}=\gcd(m,v_1,v_2,v_3,v_4)\). For \(g_{\rm pow}>1\) the power $m/g_{\rm pow}$
has an integral translation, so \(x=0\) is a fixed point of
its linear part. For \(g_{\rm pow}=1\), every proper power divisible
by $h$ has identity linear part and nonintegral translation.
All remaining obstruction indices are divisible by $h$
exactly when \(h\mid m/t\). We have therefore proved, for
every positive order and every invariant integral vector,
\[
 f\text{ is free}\quad\Longleftrightarrow\quad
 \gcd(m,v_1,v_2,v_3,v_4)=1,\quad h\mid m/t.
\]
For a free action, each nonidentity element has closed
fixed set disjoint from the compact central torus.
A sufficiently small common root neighbourhood is
therefore free for all those elements. The conjugacy
above carries this conclusion to the original affine
lift on that neighbourhood.

\emph{Existence at every finite-good position.}

There are 102 finite-good positions in the 172 rows. In the fixed
basis \(F_1=(\delta,\ell_1)\), the norm-coordinate basis \(H_{1,\rm norm}\)
has the following projected coordinates
\(w_\pm=hH_{1,\rm norm}^{-1}(1,\pm1)^t\). Here $q$ is the component-group order.

\begin{center}
\small
\begin{tabular}{lllll}
$q$ & \(g_0\) & Positions & \(w_+\) & \(w_-\) \\
\hline
1 & 1 & 34 & (1,1) & (1,-1) \\
2 & 1 & 17 & (2,1) & (2,-1) \\
2 & 2 & 17 & (1,2) & (1,-2) \\
3 & 1 & 17 & (3,1) & (3,-1) \\
3 & 3 & 17 & (1,3) & (1,-3) \\
\end{tabular}
\end{center}

These follow by summing \(I+A(1)+\cdots+A(1)^{h-1}\) and applying
the fixed-coordinate inverse \((x_\delta,x_\gamma/g_0)\):
\(H_{1,\rm norm}=\operatorname{diag}(h/q,h)\) when \(g_0=1\), and
\(H_{1,\rm norm}=\operatorname{diag}(h,h/q)\) when \(g_0=q>1\).
Both displayed vectors are primitive in every case.

The integral conjugacy \(H_d\) above satisfies
\[
 H_d\delta=\delta,\qquad H_d\ell_1=\epsilon\ell_d,\qquad
 \Pi_dH_d=H_d\Pi_1,\qquad \Pi_d=N_h(A(d))/h .
\]
Its inverse has last column
\((-\epsilon(1-\epsilon d)bz,0,\epsilon)\).
Since \(\epsilon d\equiv1\pmod{g_0}\), both matrices are integral.
Thus it identifies the two projected lattices with inverse
\(H_d^{-1}\), and
\[
       H_d^{-1}(\delta+\ell_d)=\delta+\epsilon\ell_1 .
\]
In the projected basis \(\operatorname{diag}(1,\epsilon)H_{1,\rm norm}\),
the vector \(\delta+\ell_d\) therefore has coordinates
\(w_\epsilon\). Any other integral basis is related by a
unimodular matrix $V_{\rm bas}$, giving the coordinate maps
\(w_\epsilon\mapsto V_{\rm bas}^{-1}w_\epsilon\) and its inverse $V_{\rm bas}$.
Consequently \(\delta+\ell_d\) is primitive in the projected
lattice for every permitted $d$. Its \(\delta\) coordinate in \(\Lambda\)
is one. For every \(m=hk\) both criteria in
\eqref{maps:eq:D.WP5} hold, with \(t=1\). This proves existence
of every finite-good option used in the catalogue.

\emph{The torsion and the normal character.}

Using the conjugacy \(A(d)H_d=H_dA(1)\), the two nonzero Smith entries of \(A-I\) are units.
The kernel of ${\mathsf Q_{\rm norm}}$ is the image of \(A-I\): both are primitive rank-two
subgroups, and \({\mathsf Q_{\rm norm}}(A-I)=0\). Equivalently
\(\operatorname{im}(I-R_\pi {\mathsf Q_{\rm norm}})=\operatorname{im}(A-I)\);
the projector gives its integer image and inverse coordinates. The central
quotient \(F_0\) has the affine group with generators
\(\Lambda\) and $c$, relations
\(c\lambda c^{-1}=A\lambda\), \(c^m=v\). Abelianization
and the map \({\mathsf Q_{\rm norm}}=H_{\rm norm}^{-1}{F_{\rm fix}}^{-1}N_h\) give
\[
 H_1(F_0)=
 (\mathbb Z^2\oplus\mathbb Zc)/
          \langle(-w_1,-w_2,m)\rangle
       =\mathbb Z^2\oplus\mathbb Z/t.
\]
A Bézout completion of \((-w_1,-w_2,m)/t\) gives the
forward Smith coordinates and inverse. The torsion
generator is this primitive vector modulo its $t$-fold
multiple.

The root normal character is zero on \(\Lambda\) and $1/m$
on $c$. Its value on that torsion generator is $1/t$.
Its Bockstein is the integral normal Chern class;
therefore that class has exact order $t$ and no free
component. The finite torus cover, followed by averaging
its exterior forms, gives rational Betti numbers
\((1,2,2,2,1)\). Integral duality and UCT give
\[
 H_*(F_0)=
 (\mathbb Z,\ \mathbb Z^2\oplus\mathbb Z/t,\
  \mathbb Z^2\oplus\mathbb Z/t,\ \mathbb Z^2,\ \mathbb Z).
\]
The torsion identification sends a coordinate $k$ to the
character \(j\mapsto kj/t\); evaluation at its generator
is the inverse. Thus this duality step retains the exact
torsion coordinate.

\emph{Full normal Thom columns, including the root generator.}

The normal disk bundle $U$ retracts onto \(F_0\). Its boundary
is identified with the mapping torus of $A$ by
\[
 [y,\tau]\longmapsto
 [\,y+(\tau/m)v,\ \exp(2\pi i\tau/m)\,].
\]
The relation \((y,\tau)\sim(Ay,\tau+1)\) becomes exactly
the affine deck relation. The inverse chooses a root
argument \(\tau\) and subtracts \((\tau/m)v\); a change of
argument gives that same mapping-torus relation.
The positive normal circle is
\((y,\tau)\mapsto(y-sv,\tau+ms)\).

For a lattice loop \(\lambda\), cut its normal torus into the
$m$ root sectors. Its positive fixed sweep is \(N_m\lambda\);
its fibre parallelogram is \(-\lambda\wedge v\). Rejoining
the sectors along their marked $A$-edges is the inverse
chain map. For the core loop $c$, use the parametrization
\[
 [\,uv/m,\ \exp(2\pi i(u/m+t))\,].
\]
Its coordinates in the root boundary map are \(y=-tv,\ \tau=u+mt\).
Its fixed sweep is $v$ and its vertical area is zero.
These are the normal tori of all generators of \(H_1(F_0)\).
They therefore give exactly the five global columns
\[
 \bigl(-\omega_d(\lambda,v),J_iN_m\lambda\bigr)
      \quad(\lambda=e_1,e_2,\delta,\gamma),\qquad
                      (0,J_iv).
\]
For \(\lambda\) in \((A-I)\Lambda\), both entries vanish:
\(N_m(A-I)=0\) and invariance of \(\omega_d\) with $Av=v$
gives \(\omega_d((A-I)\lambda,v)=0\).
Thus \eqref{maps:eq:D.WP6} factors through the full coinvariant group,
and the inverse section \(R_\pi\) recovers the two essential
lattice generators with the root column included.

For an explicit relation, set \(\lambda_j=R_\pi e_j\).
Then \(R_\pi w-v\in(A-I)\Lambda\), so
\(N_mR_\pi w=mv\) and \(\omega_d(R_\pi w,v)=0\).
Writing \(K_{\lambda}\) and \(K_c\) for the corresponding columns,
these identities give the exact primitive relation
\[
 (m/t)K_c-(w_1/t)K_{\lambda_1}-(w_2/t)K_{\lambda_2}=0.
\]
The two norm projections of \(\lambda_1\) and \(\lambda_2\) are
independent, so their three-column kernel has rank one.
The displayed coefficients have gcd one, hence they
generate that full integral relation lattice. The ${\mathsf Q_{\rm norm}}$
and \(R_\pi\) maps express every original lattice column
in these two coordinates, with inverse modulo \(A-I\).

The boundary \(H_2\) is free, since its four wedge Smith
entries are units by the conjugacy \(A(d)H_d=H_dA(1)\), while its fixed lattice is
free of rank two. Its Thom map consequently kills the
torsion generator in the displayed presentation of \(H_1(F_0)\); its explicit linear
combination of columns is zero. This proves the exact relation in the normal map from the full core group.

Cap product \(H_2(F_0)\to H_0(F_0)\) with the torsion
normal Chern class is zero. The Thom sequence therefore
makes the boundary \(H_2\) map onto all of \(H_2(F_0)\), including
its \(\mathbb Z/t\). Its kernel is precisely the full
image \eqref{maps:eq:D.WP6}. Thus a finite-good cap has no separate
free degree-two defect. The global elimination in
the relative attachment argument uses this surjection with its kernel and
retains its torsion.

\emph{Regular logarithms and clutch.}

At a regular site \(A=I\). The affine action is translation
by $v/m$, and its every-power freeness is exactly
\(\gcd(m,v_1,v_2,v_3,v_4)=1\). The quotient is the
torus with lattice \(\Lambda+\mathbb Z(v/m)\);
its inverse is the same real vector modulo that lattice.
The root boundary map and its inverse remain
valid, and the full kernel is \eqref{maps:eq:D.WP6} with \(N_m=mI\).
A Bézout completion of \((-v,m)\) supplies the integral
quotient coordinates. Primitivity of this vector makes the core homology free.

An extra order-one regular logarithm with \(v_*=-e\gamma\)
encodes any integer peripheral clutch. The boundary shear is
\[
 (y,\theta)\longmapsto(y-\theta v_*,\theta),
\]
with inverse plus \(\theta\) \(v_*\). Its Seifert generator is
\(c_*=\gamma^{-e}\), so deleting it from the product-one
relation gives the required product \(\gamma^e\).
The four unit normal relations eliminate its fixed
sweeps by
\[
                 t_*(u)=\omega_d(u,v_*)s.
\]
In particular \(t_*(\delta)=db\,e\,s\); its remaining
root column follows because \(\omega_d(v_*,v_*)=0\).
The inverse inserts these four sweep coordinates.
This proves the clutch convention for every $e$,
including zero and negative integers.

More generally, changing a finite peripheral lift by
\(\lambda\) changes $v$ to \(v+N_m\lambda\). The normal sweep
map changes by \(t'_i(u)=t_i(u)+\omega_d(u,\lambda)s\);
the inverse subtracts this value. Invariance gives
\(N_m^t\omega_d=\omega_dN_m\) and
\(\omega_d(N_m\lambda,\lambda)=0\). Substitution in
\eqref{maps:eq:D.WP6} verifies every relation under these maps.
At a regular site this reduces to \(v\mapsto v+m\lambda\).
The norm morphism is the change of peripheral power in both cases.

\emph{The complete integral global cokernel and its inverse.}

Form the graph kernel $K$ in \eqref{maps:eq:D.WP10}, including every regular
site and the order-one clutch site. Its primitive basis is
given by the formula for \(R_d\) and the projector following it. Form $B$ with
target \(\mathbb Zs\oplus K\) from the following columns:

\begin{enumerate}
\item
  At the weighted cusp use its two full short-normal
  columns \((0,J_*u),(0,J_*\delta)\).
\item
  At every other unchanged original or Mumford cap use
  all columns \((0,J_iF_i)\).
\item
  At a nonzero-Chern original also use \((dh_i,0)\),
  with \(h_i=b/q_i^c\) or \(b/\nu_i\), proved in
  the additive value \(\omega_d(k_{d,\rm add})=-db/q\) and the multiplicative value \(\omega_d(k_{d,\rm mult})=-db/\nu\).
\item
  At every finite or regular logarithm use all five
  columns \eqref{maps:eq:D.WP6}, with its order and power vector.
\end{enumerate}

These are \eqref{maps:eq:D.WP6}--\eqref{maps:eq:D.WP7}, evaluated in the primitive graph coordinates.
Let \(C_d=\operatorname{coker}B\) and \(T=\operatorname{Tor}C_d\).

Attach the $k$ finite or regular logarithmic caps to $V$ and call the
result $E_{\rm att}$. At least the cusp remains unfilled; write
\(s_{\mathrm{unf}}=N-k\) for the number of unfilled caps.
Choose one unfilled base loop as the
omitted bouquet edge. The integral relative group of each attached
normal disc bundle in degree two is \(H_0\) of its connected core,
namely $\mathbb Z$. Write \(x_i\) for the positive base meridian at site $i$.
A positive normal disc maps to
\[
                         m_ix_i-v_i
\]
in \(H_1\) of its boundary. In \(H_1(V)\), the vertical coordinates vanish
and its \(\gamma\) coordinate is \(-\chi(v_i)\). The resulting relative
boundary columns are
\[
                  m_ix_i-\chi(v_i)\gamma .
\]
Each has its own nonzero base coefficient \(m_i>0\). Their map to
the free abelian group \(H_1(V)=\mathbb Z\gamma\oplus\mathbb Z^{N-1}\)
is injective. The inverse on its image reads the corresponding
base coordinate and divides by \(m_i\). Excision and the Thom
sequence consequently give
\[
 H_2(E_{\rm att})=H_2(V)/(\text{all finite and regular Thom images}),
\]
with the five columns at each site, including the root generator.
The full group \(H_1(E_{\rm att})\) retains the presentation by those power
columns and the one base relation; its free rank is \(s_{\mathrm{unf}}\).

For each remaining original or Mumford cap \(U_j\), duality gives
\[
 H_2(U_j,\partial U_j)=H^4(U_j),
\]
and this is free. The free ranks are $d$ at the cusp, \(\rho_j(d)\)
at a zero-Chern additive original, \(s_j\) at a nonzero-Chern
additive original, \(\nu_j\) at a nonzero-Chern multiplicative
original, and the component count at a rank-one filling.
The circle-bundle Gysin and component-orbit bases above compute these groups. For example a nonzero additive
original has \(H_3\) of its angular circle bundle equal to
\(H_2(N_{\rm loc})=\mathbb Z^{s_j}\); product with the suspension circle
therefore has \(H^4=\mathbb Z^{s_j}\). The multiplicative case
gives \(\mathbb Z^{\nu_j}\) by the same map.

Choose the dual normal discs at the generic points of these
component cycles. In a component of multiplicity \(m_C\), the
base parameter has the form \(q=u^{m_C}\) times a nonvanishing
unit. A disc in a nonzero line frame lifts to the prequotient
line space and has \(\gamma\) displacement zero. Its boundary is
therefore \(m_Cx_j\) plus a vertical class, which vanishes
in \(H_1(E_{\rm att})\). The cusp and primitive Mumford discs have \(m_C=1\).
For a component orbit choose any representative; the quotient
is a local biholomorphism there, so this disc intersects the
orbit cycle once. Its coefficient remains \(m_C\).
Every original has a multiplicity-one outer component, so the
images at site $j$ generate exactly \(\mathbb Zx_j\) in \(H_1(E_{\rm att})\).

The complete degree-three images are the full boundary kernels
proved above. Passing through \(H_2(E_{\rm att})\) gives the exact sequence
\begin{equation}
 0\longrightarrow C_d\longrightarrow H_2(X)
 \longrightarrow
 \ker\!\left\{\bigoplus_jH^4(U_j)\longrightarrow H_1(E_{\rm att})\right\}
 \longrightarrow0,\qquad C_d=\operatorname{coker}B .
\hypertarget{tag-027Q}{}\label{maps:eq:D.WP11}
\end{equation}

An explicit kernel algorithm retains \(H_1(E_{\rm att})\)'s torsion. Let ${\mathsf P_{\rm log}}$ be its
integer power-relation matrix in the omitted-edge coordinates and
${\mathsf M_{\rm disc}}$ the displayed component-disc boundary matrix. Solve \({\mathsf M_{\rm disc}}z={\mathsf P_{\rm log}}z'\) over the integers. Project the
integer kernel of \([{\mathsf M_{\rm disc}},-{\mathsf P_{\rm log}}]\) to $z$. The columns of ${\mathsf P_{\rm log}}$ are independent
by their base coordinates, so this projection is injective on the
relation coefficients and gives exactly the required subgroup of
the component lattice. Hermite or reversible Smith operations
supply a basis and its inverse coordinates.

For each such basis vector choose its relative disc chain. Its
boundary bounds the combination of the power discs specified by
$z'$ and an integral chain in $E_{\rm att}$. Gluing these chains produces a
closed two-cycle with that relative coordinate. Subtracting these
cycles and reading the remaining exterior coordinate gives an
inverse to the splitting of the exact sequence.

Denote the component-disc boundary map by \(\Gamma\). Its rank is
\(s_{\mathrm{unf}}-r\), because its cokernel is \(H_1(X)\) and \(H_1(E_{\rm att})\)
has rank \(s_{\mathrm{unf}}\). Its domain has rank
\(d+(s_{\mathrm{unf}}-1)+\beta+L\). Thus its kernel is free of rank
\begin{equation}
 \operatorname{rank}\ker\Gamma=d-1+\beta+L+r.
\hypertarget{tag-027R}{}\label{maps:eq:D.WP12}
\end{equation}
The integral basis includes every torsion-index effect.
The exact sequence and the integral section constructed above give
\[
              H_2(X)=C_d\oplus\mathbb Z^{d-1+\beta+L+r}.
\]
Its inverse reads the relative component coordinates and then the
remaining exterior class modulo $B$.

\emph{The full group, the unrestricted ranks, and every homology group.}

\emph{The nonabelian vertical quotient.}

The cusp kills its primitive short lattice
\(S=\mathbb Z\langle u,\delta\rangle\) and the origin-disc
meridian. Its normal closure under the base peripheral
conjugations is the same vertical subgroup for all $d$:
on vectors of \(\gamma\) coordinate zero, \eqref{maps:eq:D.WP2} is literally
the original matrix. The vector \(u\) and its conjugates
therefore give the same integer vanishing words as in
Theorem~\ref{maps:src:D.2.23}.

On 170 rows those words give two elliptic vectors $u,v$
of determinant one after subtracting their angular
coordinates. If \(u=(u_1,u_2),v=(v_1,v_2)\), their exact
inverse is
\[
 e_1=v_2\widetilde u-u_2\widetilde v,\qquad
 e_2=-v_1\widetilde u+u_1\widetilde v .
\]
Together with \(\delta\) this kills $\ker\chi$.

The two OS46/II rows have \((a,b)=(1,4),(4,1)\).
Their first short vectors are \(u=(1,1),v=(-4,-1)\),
giving \(3e_1=-\widetilde u-\widetilde v\) and the
index-three elliptic sublattice \(x-y=0\bmod3\).
Their \(I_3^*\) is kept original in both scope readings.

For \(a=1\) its line Chern vector evaluates to one on
the sphere met by $P$. The circle-bundle Hopf disc
therefore kills \(\delta\) and the simply connected affine
\(D_7\) plumbing contracts the elliptic loops by projection
and path lifting. In particular the boundary
loop \(w_1=2e_1+\delta\) bounds. Subtracting
\(w_1-\delta=2e_1\) from \(3e_1\) gives \(e_1\).

For \(a=4\) the line is trivial over that same plumbing.
Its boundary-to-cover map on the first three coordinates
is
\[
                   (x,y,z)\longmapsto z-2x-y,
\]
with section \(k\mapsto k\delta\). It kills
\(w_4=e_1+2\delta\), so \(e_1\) again bounds. Then
\(e_2=\widetilde u-e_1\). These are unchanged maps on the
line bundle for every permitted $d$; only the
suspension arrow has been iterated. The original star cap therefore kills the remaining index-three quotient.

Every vanishing relation used has \(\chi\) value zero.
The unrolled contraction defines a return map
on every cap, with \(\gamma\) mapped to one. This proves
the reverse inclusion of the normal subgroup. Thus
the global vertical quotient is exactly
\[
 \Lambda/\mathcal N\longleftrightarrow\mathbb Z\gamma,
 \qquad [v]\longmapsto\chi(v)\gamma,\quad
       k\gamma\longmapsto[k\gamma].
\]
The two compositions are the identity.

\emph{Every order, numerator and clutch.}

A boundary fibre generates the original or Mumford
cover cap group: move a loop off the central divisor,
then contract its meridian using a transverse disc
at a multiplicity-one point in an extending line
frame. The toric coordinate quotient gives the same
surjection for a Mumford model. After the vertical quotient, the
cover cap group is \(\mathbb Z\gamma\), with return
as forward map and \(\gamma\) as inverse section.

At a free logarithm, select a root over the boundary
basepoint. Its fibre maps isomorphically to the quotient
fibre, with inverse choosing that root representative.
The deck path \(c_i\) has the \(m_i\)-fold power \(v_i\).
Conjugation is \(A_i(d)\), and \(\chi A_i(d)=\chi\). The full
reduced local group is therefore
\[
 \langle\gamma,c_i\mid[\gamma,c_i]=1,\
                      c_i^{m_i}=\gamma^{\mu_i}\rangle,
             \qquad \mu_i=\chi(v_i).
\]
The inverse normal form reads the covering-sheet index
\(0\le j<m_i\), then the remaining return exponent $k_{\rm ret}$,
and writes \(\gamma^{k_{\rm ret}} c_i^j\). The injection of the
cover return group proves uniqueness.

Before taking the contraction quotient choose the
peripheral lifts closed in its infinite cyclic cover.
Their product has return zero, as do all line-frame
windings. Thus the initial canonical clutch is zero.
The order-one map of the regular and clutch construction realizes every integer $e$.

Van Kampen and the exact inverse maps above now give
the full group
\[
 \pi_1(X)=\left\langle\gamma,c_1,\ldots,c_s\ \middle|\
 [\gamma,c_i]=1,\ c_i^{m_i}=\gamma^{\mu_i},\
                  c_1\cdots c_s=\gamma^e\right\rangle.
\]
Its inverse map sends each fibre vector to its \(\chi\)
power of \(\gamma\), each unchanged meridian to one, and
each deck path to \(c_i\). The displayed relations
are exactly their monodromy, vanishing, power and
global product relations, so both compositions fix
every generator.

A finite lift change by \(\lambda\) has
\(\mu_i'=\mu_i+\chi(N_{m_i}\lambda)
        =\mu_i+m_i\chi(\lambda)\).
The global clutch changes by the sum of \(\chi(\lambda)\).
These are the generator changes
\(c'_i=c_i\gamma^{\chi(\lambda)}\), with inverse
\(c_i=c'_i\gamma^{-\chi(\lambda)}\).
They also match the normal-sweep maps in the regular and clutch construction.

The abelianization matrix, on \(\gamma\),\(c_1\),\ldots,\(c_s\), has
rows
\[
 (-\mu_i,0,\ldots,m_i,\ldots,0),\qquad(-e,1,\ldots,1).
\]
Its Smith maps give \(H_1=\mathbb Z^r\oplus A_{1,\rm tor}\) and
inverse coordinates. Over the rationals the first $s$
rows eliminate \(c_i\) as \((\mu_i/m_i)\gamma\), so
\[
 r=1\quad\text{exactly when}\quad
       e-\sum_i\mu_i/m_i=0;
 \qquad r=0\text{ otherwise}.
\]
This statement also covers an empty site set and
order-one sites through the same matrix and its
elimination. The same eliminations apply to every finite number of sites and every positive order.

\Needspace{5\baselineskip}
\emph{The free rank of the full boundary cokernel.}

Over \(\mathbb Q\), the literal map \(S_d\) identifies
the graph complex and fixed-peripheral insertions with
their weight-one counterparts. More explicitly,
\[
 [A_i(d)-I]\operatorname{diag}(S_d^{-1},\ldots,S_d^{-1})
        =S_d^{-1}[A_i-I],
\]
and \(J_i(d)S_d^{-1}\lambda\) is the blockwise inverse
\(S_d\) image of \(J_i(1)\lambda\). The full fixed-normal
map at weight one is onto the graph kernel and has
exactly the one \(\delta\) relation, by the integral map and inverse of Theorem~\ref{maps:src:D.2.31}. These displayed rational
maps prove the corresponding rank and relation
for every $d$. Each added identity site adds its full
four-dimensional fixed space, with the same one
\(\delta\) relation and the deletion/restoration inverse.

At a logarithmic site the full norm is $m$ times the
fixed-plane projection. For a fixed $u$, \eqref{maps:eq:D.WP6} therefore
gives over the rationals
\[
                    t_i(u)=\omega_d(u,v_i)s/m_i.
\]
The root column is redundant there because
\(\omega_d(v_i,v_i)=0\). At every unchanged cap the
fixed sweeps are zero. After deleting the dummy clutch
site, the one \(\delta\) relation consequently becomes
\[
        db\left(e-\sum_i\mu_i/m_i\right)s=0.
\]
Every other graph coordinate has been eliminated by
the explicit fixed-normal maps. If \(L=0\), the Euler criterion and scalar graph relation
give \(\operatorname{rank}C_d=r\).
If \(L>0\), one of the positive coefficients \(dh_i\) kills
$s$ over the rationals, and \(\operatorname{rank}C_d=0\).
This rational calculation is used only for the free
rank; all torsion remains the integral Smith group
of $B$ already constructed.

Substituting these ranks in \eqref{maps:eq:D.WP11}--\eqref{maps:eq:D.WP12} proves
\[
 b_2=
 \begin{cases}
 d-1+\beta+2r,&L=0,\\
 d-1+\beta+L+r,&L>0.
 \end{cases}
\]

\emph{Euler characteristic and integral duality.}

The smooth torus bundle exterior has Euler characteristic
zero. Every noncusp original is a circle mapping-torus
construction and has Euler characteristic zero; the
finite and regular cores are free torus quotients and
also have Euler characteristic zero. Their boundary
mapping tori have Euler characteristic zero. The cusp homology gives \(\chi(N_d)=2d\). Additivity on these pieces therefore proves \(\chi(X)=2d\).

The closed complex threefold is oriented. Integral
Poincaré duality and UCT give the free symmetry
\(b_5=b_1=r,\ b_4=b_2\), torsion in \(H_4\) equal to \(A_{1,\rm tor}\),
and torsion in \(H_3\) equal to $T$. More explicitly, a Smith
torsion generator of order $n$ determines the character
\(j\mapsto j/n\) in \(\mathbb Q/\mathbb Z\); its
Bockstein followed by cap product with the fundamental
class gives the dual torsion coordinate. Evaluating
its linking character on the same generator is the
inverse. Free coordinates use dual integral cocycles
and the same cap product. Thus the identifications
give the integral dual groups and their inverse coordinates.

The Euler identity
\(2-2r+2b_2-b_3=2d\) now gives
\[
 b_3=
 \begin{cases}
 2\beta+2r,&L=0,\\
 2\beta+2L,&L>0.
 \end{cases}
\]
The complete groups are
\[
\begin{array}{c|c}
 j&H_j(X;\mathbb Z)\\ \hline
 0,6&\mathbb Z\\
 1&\mathbb Z^r\oplus A_{1,\rm tor}\\
 2&\mathbb Z^{b_2}\oplus T\\
 3&\mathbb Z^{b_3}\oplus T\\
 4&\mathbb Z^{b_2}\oplus A_{1,\rm tor}\\
 5&\mathbb Z^r .
\end{array}
\]
This proves the displayed homology groups for every allowed parameter.

At \(d=1\) the covering \(p_d\) is the
identity, and every map above restricts to its weight-one counterpart. This includes the ordinary
identity on the subgroup lattice and the empty
zero-sum component lattice in the displayed cusp-relative decomposition.

\emph{Enumeration and the two starred-cap scopes.}

In $\mathcal P$ of Definition~\ref{mapsIII:def:weighted-natural-data},
label all five finite-good types II, III, IV,
\(\mathrm{III}^*\), \(\mathrm{IV}^*\), each with original O or free finite F, and retain
every odd \(I_n^*\) original. At each nonzero-Chern zero-free
multiplicative site take original O or primitive Mumford M.
At each zero-Chern multiplicative site identify O and M through
the marked boundary maps of the zero-Chern multiplicative construction; use M as its
representative. Use C at the selected weighted cusp.

For $\mathcal P_{\mathrm{all\ star\ O}}$ of that definition,
the inclusion $j$ keeps every site
label. The projection \(\pi\) from \(\mathcal P\) replaces each
\(\mathrm{III}^*\)/\(\mathrm{IV}^*\) F by O and keeps every other label. Thus
\(\pi j=I\), with all the changed local labels explicit.
This is a map of the finite pattern sets. The local kernel
change made by \(\pi\) is the change from \eqref{maps:eq:D.WP6} to
\eqref{maps:eq:D.WP7} or the three additive kernel columns, already described by those formulas.

The quotient from all individual multiplicative choices
to these patterns forgets O versus M exactly at the
zero-Chern multiplicative sites. Its canonical section
chooses M there. The full inverse fibre of a pattern is
the independent set of its \(2^{z}\) O/M choices, where $z$
is the number of those sites. The zero-Chern multiplicative construction gives the
marked integral kernel morphisms between these choices.

At a nonzero-Chern multiplicative site, the primitive
class \(k_{d,\rm mult}\) is killed by the original cap and maps to a
primitive nonzero class in the Mumford cap. Its
\(\omega_d\) value is \(-db/\nu\ne0\), so their marked
kernels differ. At an additive site the original
contains all fixed-normal sweeps in its kernel.
For the finite alternative the rational kernel graph
is \(t(u)=\omega_d(u,v)s/m\). The fixed-plane form
is nondegenerate, since
\(\omega_d(\delta,\ell_d)=-db\,g_0\ne0\), and a free
affine power $v$ cannot be zero. Hence some fixed sweep
has nonzero transfer. The finite kernel has no rational
pure fibre element, since its normal projection is an
isomorphism onto the fixed plane. This proves that it
differs from the original kernel. Thus the stated
identifications are exactly the local marked-kernel
identifications used in the enumeration.

For each defining profile the exhaustive factor enumeration of
Theorem~\ref{maps:src:D.2.22} has $a$ among the positive divisors
of $d_0/n_*$, at most $30$ in the marked data. For each resulting
row take the Cartesian product of the site option sets just
stated. Each marked choice is indexed by that row and its ordered
site options, so this product lists every choice exactly once.

The finite factor and site-option enumeration gives the following totals.
Here $\mathrm R$ is the natural-row count, $\mathrm K$ the nonzero-Chern
multiplicative position count, $\mathrm P$ the broad pattern
count, $\mathrm S$ the all-star-original subset, and $\mathrm A$ the
individual original/Mumford choices.

\begin{center}
\small
\begin{tabular}{lllllll}
OS & Profile & R & K & P & S & A \\
\hline
43 & base & 6 & 0 & 12 & 6 & 48 \\
43 & II & 2 & 0 & 8 & 4 & 8 \\
45 & base & 8 & 4 & 12 & 12 & 128 \\
45 & II & 4 & 2 & 12 & 12 & 32 \\
46 & base & 6 & 0 & 6 & 6 & 24 \\
46 & II & 2 & 0 & 4 & 4 & 4 \\
47 & base & 14 & 13 & 30 & 30 & 224 \\
47 & II & 6 & 5 & 24 & 24 & 48 \\
47 & III & 8 & 4 & 24 & 24 & 64 \\
49 & base & 10 & 4 & 28 & 14 & 80 \\
49 & II & 2 & 0 & 8 & 4 & 8 \\
49 & III & 4 & 0 & 16 & 8 & 16 \\
50 & base & 10 & 4 & 14 & 14 & 40 \\
50 & II & 2 & 0 & 4 & 4 & 4 \\
50 & IV & 4 & 0 & 8 & 8 & 8 \\
55 & base & 14 & 13 & 30 & 30 & 224 \\
55 & II & 6 & 5 & 24 & 24 & 48 \\
56 & base & 26 & 34 & 77 & 77 & 416 \\
56 & II & 10 & 10 & 46 & 46 & 80 \\
56 & III & 14 & 11 & 54 & 54 & 112 \\
56 & IV & 14 & 11 & 54 & 54 & 112 \\
\textbf{Total} &  & \textbf{172} & \textbf{120} & \textbf{495} & \textbf{459} & \textbf{1728} \\
\end{tabular}
\end{center}

There are also exactly 316 zero-Chern multiplicative
positions. The 102 finite-good positions comprise
78 of types II/III/IV and 24 of types \(\mathrm{III}^*\)/\(\mathrm{IV}^*\).
The 36 patterns in the broad set outside the all-star-original subset
are distributed as follows. The map \(\pi\) changes exactly the
\(\mathrm{III}^*/\mathrm{IV}^*\) finite label F to O.
The inverse image of a pattern consists of the independently
permitted O/F values at those positions. The inclusion is the all-O section.

\begin{center}
\small
\begin{tabular}{lll}
OS & Profile & Omitted patterns \\
\hline
43 & base & 6 \\
43 & II & 4 \\
49 & base & 14 \\
49 & II & 4 \\
49 & III & 8 \\
\end{tabular}
\end{center}

The local map changed by this projection is exactly the finite
Thom map versus the original fixed-normal map and its possible
fibre column. Thus the formulas hold on all 495 patterns and
restrict through the inclusion to the 459-pattern subset.
\end{proof}

\begin{remark}[Finite weighted-profile calculations]
\hypertarget{tag-Q137}{}\label{mapsIII:rem:weighted-profile-calculations}
{\normalfont\scriptsize tag Q137}\par
For the $172$ natural rows and $495$ patterns, use the matrices
$A_i(d)$ and the complete local kernels in the proof. Its integer
right inverses determine the exterior quotient. Smith reduction
of the full integral boundary presentation gives the cokernel,
and the integral relative splitting gives the homology for every
stated parameter. Definition~\ref{mapsIII:def:weighted-natural-data}
specifies the initial translation iterates and every permitted
cap choice.

Theorem~\ref{maps:src:D.2.41} gives the integral algorithm for all
positive weights, including weights sharing factors with $a$.
The inclusion of the present input set fixes the initial marked
row, $d$, every ordered cap choice, all logarithmic $(m,v)$ data
and the clutch.
\end{remark}

\Needspace{9\baselineskip}
\subsection{Shared odd factors and saturated cap kernels}

\begin{definition}[Weighted data with odd common factor]
\hypertarget{tag-Q138}{}\label{mapsIII:def:odd-weighted-data}
{\normalfont\scriptsize tag Q138}\par
Choose one of the $172$ natural marked rows of
Theorem~\ref{maps:src:D.2.22}, with $P=aG$, $Q=bG$,
$a,b>0$ and $\gcd(a,b)=1$. Take the $d$th positive iterate
of its contracting line map, with
\begin{equation}
                         g=\gcd(a,d)\text{ odd}.
 \hypertarget{tag-022W}{}\label{maps:eq:D.OW1}
\end{equation}
Use the norm linearization, original chart units and period
subgroup of that iterate. Let $\mathcal P$ and
$\mathcal P_{\mathrm{all\ star\ O}}$ be the two sets of ordered
cap choices defined in Definition~\ref{mapsIII:def:weighted-natural-data}:
the latter keeps every starred cap original, while the former
also permits finite-good $\mathrm{III}^*,\mathrm{IV}^*$ fillings.
Odd $I_\nu^*$ caps remain original in both sets. Retain every
permitted original or marked Mumford choice, every free finite
logarithm, every finite list of free regular logarithms, and
every integer clutch $e$.

In the periods of Definition~\ref{mapsIII:def:elliptic-marked-data}, put
\[
 \Lambda=\mathbb Z\langle e_1,e_2,\delta,\gamma\rangle,
 \qquad S_d=\operatorname{diag}(1,1,1,d),
 \qquad A_i(d)=S_d^{-1}A_iS_d.
\]
Include the added regular sites in the ordered list of $N$ sites
and represent clutch $e$ by an auxiliary order-one regular site
with power $-e\gamma$. Put an unfilled site, such as the selected cusp, last in this
ordering and omit it from the bouquet. Write
\[
 K=\ker\bigl([A_i(d)-I]_{i<N-1}:\Lambda^{N-1}\longrightarrow\Lambda\bigr).
\]
For $F_i=\ker(A_i(d)-I)$, let $J_i:F_i\longrightarrow K$ be
the based peripheral insertion with the fixed stems; its block
and inverse-word coordinates are given in the proof.

Let $L$ count nonzero-Chern originals. Let $\beta$ be the sum,
away from the selected cusp, of $\gcd(d,\nu_i)-1$ at a narrow
zero-free rank-one site, zero at a non-narrow Mumford site,
$\rho_i(d)-1$ at a zero-Chern additive original, and $s_i-2$
or $\nu_i-2$ at a nonzero-Chern additive or multiplicative
original. Here $s_i$ is the number of Kodaira components,
$\nu_i$ is the multiplicative index, and $\rho_i(d)$ is the
number of orbits of the actual $d$th component permutation.
\end{definition}

\begin{theorem}[Shared odd factors in weighted profiles]
\hypertarget{tag-02LH}{}\label{maps:src:D.2.40}
{\normalfont\scriptsize tag 02LH}\par
For the data of Definition~\ref{mapsIII:def:odd-weighted-data},
the following integral algorithm applies to both cap sets.
Across the $172$ rows they contain $495$ and $459$ patterns,
respectively. Inclusion
fixes every labelled choice; retraction replaces each
finite-good $\mathrm{III}^*$ or $\mathrm{IV}^*$ filling by its
original cap. Their composite on the smaller set is the identity.
In particular, the algorithm applies to every positive iterate
of every row with odd $a$.

The full relation matrix has target \(\Lambda^2\Lambda\oplus K\).
Include every return column
\begin{equation}
 ((\Lambda^2 A_i(d)-I)\eta,0)
       \qquad(\eta\text{ a standard bivector}).
 \hypertarget{tag-022X}{}\label{maps:eq:D.OW2}
\end{equation}

At a finite-good site let \(h\) be the order of its linear return,
take \(m>0\) divisible by \(h\), and take
\(v=v_i\in\ker(A_i(d)-I)\). Write \(v_{i,1},\ldots,v_{i,4}\)
for its coordinates in \((e_1,e_2,\delta,\gamma)\); their inverse is
\(v_i=v_{i,1}e_1+v_{i,2}e_2+v_{i,3}\delta+v_{i,4}\gamma\).
At a regular site use \(h=1\) and
\(A_i(d)=I\). Put \(N_m=\sum_{j=0}^{m-1}A_i(d)^j\).
At a finite-good site, let \(F:\mathbb Z^2\to\ker(A_i(d)-I)\)
be the isomorphism onto the primitive fixed lattice given by its chosen basis, and put
\[
 H\mathbb Z^2=\operatorname{im}(F^{-1}N_h),\qquad
 Q_{\rm fin}=H^{-1}F^{-1}N_h,\qquad w_{\rm fin}=Q_{\rm fin}v .
\]
The finite action is free precisely when
\[
 \gcd(m,v_{i,1},v_{i,2},v_{i,3},v_{i,4})=1
 \quad\text{and}\quad
 h\mid m/\gcd(m,w_{{\rm fin},1},w_{{\rm fin},2}).
\]
At a regular site the first condition is the complete freeness
criterion. Include all five columns
\begin{equation}
 (-\lambda\wedge v,J_iN_m\lambda)
       \quad(\lambda=e_1,e_2,\delta,\gamma),\qquad
 (0,J_iv).
 \hypertarget{tag-022Y}{}\label{maps:eq:D.OW3}
\end{equation}
An integer clutch \(e\) is represented by the order-one regular
operation with \(v=-e\gamma\).

At every unchanged cap include its full fixed-normal columns
\((0,J_iF_i)\), where \(F_i=\ker(A_i(d)-I)\), and the fibre kernel
\begin{equation}
 \mathcal K_i(d)=
 \{\eta\in\Lambda^2\Lambda:
       (\Lambda^2S_d)\eta\in\mathcal K_i(1)\otimes\mathbb Q\}.
 \hypertarget{tag-022Z}{}\label{maps:eq:D.OW4}
\end{equation}
For a nonzero-Chern original, \(\mathcal K_i(1)\) is its principal-bundle fibre kernel from
Theorems~\ref{maps:src:D.2.32} and~\ref{maps:src:D.2.33}. At each other
unchanged cap its rational span is
\(\operatorname{im}(\Lambda^2A_i(1)-I)\). Here \(A_i(1)\)
is the original marked return, the weight-one member of the same
family \(A_i(d)\), with \(S_1=I\).
The kernel in \eqref{maps:eq:D.OW4} is the saturated integer
kernel of this map. All nonunit factors in
\eqref{maps:eq:D.OW2} remain in the full matrix.

Let \(T\) be the torsion subgroup of its cokernel.
For each logarithm other than the clutch site put
\(\mu_i=\gamma^*(v_i)\). Then
\[
 \pi_1(X)=
 \left\langle\gamma,c_1,\ldots,c_k\ \middle|\
  [\gamma,c_i]=1,\quad c_i^{m_i}=\gamma^{\mu_i},\quad
                         c_1\cdots c_k=\gamma^e\right\rangle .
\]
Write its abelianization as \(\mathbb Z^r\oplus A_1\);
thus \(r=1\) exactly when \(e-\sum_i\mu_i/m_i=0\).
Use $L$ and $\beta$ from
Definition~\ref{mapsIII:def:odd-weighted-data}.
The selected cusp has $d$ components.
The free ranks are
\[
 b_2=\begin{cases}
 d-1+\beta+2r,&L=0,\\
 d-1+\beta+L+r,&L>0,
 \end{cases}
 \qquad
 b_3=\begin{cases}
 2\beta+2r,&L=0,\\
 2\beta+2L,&L>0.
 \end{cases}
\]
The complete integral homology is
\[
 \begin{array}{c|c}
 j&H_j(X;\mathbb Z)\\ \hline
 0,6&\mathbb Z\\
 1&\mathbb Z^r\oplus A_1\\
 2&\mathbb Z^{b_2}\oplus T\\
 3&\mathbb Z^{b_3}\oplus T\\
 4&\mathbb Z^{b_2}\oplus A_1\\
 5&\mathbb Z^r .
 \end{array}
\]
In particular \(\chi(X)=2d\). These are
\eqref{maps:eq:D.WP8}--\eqref{maps:eq:D.WP9}, evaluated at
\(A_{1,\rm tor}=A_1\), with the full integral matrix
\eqref{maps:eq:D.OW2}--\eqref{maps:eq:D.OW4}.
\end{theorem}


\begin{proof}
Let \((a,b,*)\) be the given natural row, and write
\[
 g=\gcd(a,d)=\gcd(a,db),\qquad
 A_i(d)=S_d^{-1}A_i(1)S_d,\qquad
 S_d=\operatorname{diag}(1,1,1,d).
\]
The translation iterate is positive. The equality of
the two gcds uses \(\gcd(a,b)=1\).
Put \(C_d=(\Lambda^2\Lambda)_{\{A_i(d)\}}\) and
\(K_d=K=\ker[A_i(d)-I]_{i<N-1}\).
Write $B_i,g_i$ for the ordered elliptic matrix and section
cocycle in Definition~\ref{mapsIII:def:elliptic-marked-data}. Put
\[
 J_{\mathrm{ell}}=\begin{pmatrix}0&1\\-1&0\end{pmatrix},\qquad
 \beta_i^1=-g_i^tJ_{\mathrm{ell}}B_i,\qquad
 A_i(1)=\begin{pmatrix}
 B_i&0&b g_i\\
 a\beta_i^1&1&c_{i,\rm ret}\\
 0&0&1
 \end{pmatrix}.
\]
At \(I_\nu\), choose the primitive Tate basis \(u_{\rm T},r_{\rm T}\) with
\(B_i-I=\nu u_{\rm T}u_{\rm T}^tJ_{\mathrm{ell}}\), \(g_i=\kappa u_{\rm T}\), and
\[
 c_{i,\rm ret}=-ab\kappa^2/\nu+[i=*]/\nu.
\]
At an additive site,
\(c_{i,\rm ret}=ab\beta_i^1(B_i-I)^{-1}g_i\).
The defining component-degree equations make these matrices
integral and give their ordered product equal to the identity.
The first three entries of the \(\gamma\) column are multiplied by
$d$ under the iterate. All other entries are unchanged.

With the last site omitted, put
\[
 \mathsf U_2=[B_i-I]_i,\qquad
 \mathsf U_3=\left[\binom{B_i-I}{\beta_i^1}\right]_i .
\]
The maximal-minor calculations in Theorem~\ref{maps:src:D.2.37} give
integer sections \(R_2,R_3\) with
\(\mathsf U_2R_2=I_2\), \(\mathsf U_3R_3=I_3\), on each defining profile.
For completeness, take a Bézout combination of the maximal
minors, insert each adjugate in its column positions, and add
with those coefficients. Multiplication gives the stated
sections. Their identities contain no weight parameter.
The weight-one degree-one relation matrix also has its section
\(G_1R_1=I_3\).


The ordered product of the returns is the identity. For lattice
automorphisms \(M_{\rm word},N_{\rm word}\) and a lattice vector \(x\), the identities
\[
 \begin{aligned}
 (M_{\rm word}N_{\rm word}-I)x
   &=(M_{\rm word}-I)N_{\rm word}x+(N_{\rm word}-I)x,\\
 (M_{\rm word}^{-1}-I)x
   &=-(M_{\rm word}-I)M_{\rm word}^{-1}x
 \end{aligned}
\]
express the relation of the last return, the inverse word in the
other returns, as an integer combination of their relations.
Apply the same identities to the exterior-square representation.
Thus the first \(N-1\) sites generate the full relation image in
both degrees. Inclusion of their columns is the forward image
map, and expansion of the inverse word gives its image section.
The two induced quotient maps are inverse because both fix the
classes of every lattice basis vector or standard bivector.

\emph{Exact degree-one exterior image at every positive weight.}

Write \(G_{d}\) for
the first three rows of the common degree-one relation matrix, omitting
the last site. The elliptic columns are
\(\operatorname{diag}(1,1,a)\mathsf U_3\); their image is exactly
\[
       \mathbb Ze_1+\mathbb Ze_2+a\mathbb Z\delta .
\]
Their first two unit preimages are \(R_3e_1\) and \(R_3e_2\) in the elliptic
column slots. Write \(z_{a}\) for the elliptic-column insertion of \(R_3e_3\);
then \(G_dz_a=a\delta\).

At weight one let \(R_1\) be the integer right inverse of \(G_1\). Let \(L_{d}\)
multiply every nongamma domain coordinate by $d$ and fix each \(\gamma\)
coordinate. The literal column rule gives
\[
 G_dL_d=dG_1 .
\]
For integers \(u_g,v_g\) with \(u_g a+v_g d=g\), the vector
\[
                 z_g=u_g z_a+v_g L_dR_1\delta
\]
maps to \(g\delta\). Conversely every angular entry of a source
column is a multiple of $a$ or $d$. Therefore
\[
 \operatorname{im}G_d
      =\mathbb Ze_1+\mathbb Ze_2+g\mathbb Z\delta .
\]
The full degree-one quotient and inverse are
\[
 \Lambda/\operatorname{im}[A_i(d)-I]\longleftrightarrow
       \mathbb Z\gamma\oplus(\mathbb Z/g)\delta,\qquad
 [x]\longmapsto(x_\gamma,[x_\delta]_g),
\]
with inverse \((n,[k])\mapsto[n\gamma+k\delta]\).

Put \(M_d=\operatorname{diag}(1,1,1/g)G_d\). It is an integer matrix,
and the two unit preimages and \(z_{g}\) form an integer right inverse
\(R_{g}\). The projector \(I-R_gM_d\) has image \(K_{d}\). A Hermite basis
\(K_{\mathrm{bas}}\) of that image gives the unimodular square matrix
\([R_g,K_{\mathrm{bas}}]\); its inverse reads \(M_{d}\) in the first three
coordinates and the kernel coordinates in the remaining rows.
This computes every \(J_{i}\) in a primitive integral basis even when
\(g>1\).

\emph{The full fibre quotient and its explicit coordinates.}

Order the wedge basis as
\[
       e_{12},e_{1\delta},e_{1\gamma},
              e_{2\delta},e_{2\gamma},\delta\gamma .
\]
The preimages can be written before dividing by any gcd.
On the two angular columns the relation map is \(B_i-I\), so
inserting \(R_2\) gives preimages of \(e_{1\delta},e_{2\delta}\).
For an elliptic vector \(\lambda\), the remaining coordinates
of the relation on \(\lambda\wedge\gamma\) are
\[
 (db\,\beta_i^1\lambda,\ (B_i-I)\lambda,\
                              a\beta_i^1\lambda).
\]
Insert \(R_3\) in these elliptic-\(\gamma\) column positions. Its
three images, modulo the two angular coordinates, are
\(e_{1\gamma},e_{2\gamma},db\,e_{12}+a\,\delta\gamma\).
Subtract their angular coordinates using the first two
preimages. These integer operations are linear in $d$ and give
\[
 \operatorname{im}[\Lambda^2A_i(d)-I]_{i<N-1}
 =
 \mathbb Z\langle e_{1\delta},e_{2\delta},e_{1\gamma},e_{2\gamma},
                  db\,e_{12}+a\,\delta\gamma\rangle .
\]
For the reverse inclusion at shared factors, inspect a source
bivector. A source with no \(\gamma\) coordinate has zero area and
\(\delta\)-\(\gamma\) coordinates in its relation image. A source
\(\lambda\)-\(\gamma\) has those coordinates
\(db\,\beta_i^1\lambda\) and \(a\,\beta_i^1\lambda\);
\(\delta\)-\(\gamma\) itself has both zero. Thus every source column has
its two remaining coordinates an integer multiple of \((db,a)\).
This proves the reverse inclusion in the displayed full wedge image integrally.

Put \(\alpha_{\mathrm f}=a/g,\ \beta_{\mathrm f}=db/g\), and choose $u_f,v_f$ with
\(u_f\alpha_{\mathrm f}+v_f\beta_{\mathrm f}=1\). The quotient map is
\[
 [x]\longmapsto
 \bigl(\alpha_{\mathrm f} x_{12}-\beta_{\mathrm f} x_{\delta\gamma},\
       [v_f x_{12}+u_f x_{\delta\gamma}]_g\bigr)
       \in\mathbb Z\oplus\mathbb Z/g .
\]
Its inverse sends \((n,[k])\) to
\[
 n(u_f e_{12}-v_f\delta\gamma)
             +k(\beta_{\mathrm f} e_{12}+\alpha_{\mathrm f}\delta\gamma).
\]
Changing $k$ by $g$ adds precisely the last relation in the displayed full wedge image.
The two composites are identity, since the underlying two-by-two
matrix has determinant \(u_f\alpha_{\mathrm f}+v_f\beta_{\mathrm f}=1\).
In particular the full relation target retains its exact \(\mathbb Z/g\)
summand.

The free functional is
\[
 \omega'_d=(a/g)e_1^*\wedge e_2^*
                  -(db/g)\delta^*\wedge\gamma^*.
\]
Put \(\omega_d=g\omega'_d\). The coefficient \(\beta_{\mathrm f}\)
belongs to these quotient coordinates; the cap contribution
$\beta$ is the component count of
Definition~\ref{mapsIII:def:odd-weighted-data}.

\emph{The marked graph and its integral involution section.}

Use the real-period map and its inverse from
Theorems~\ref{maps:src:D.2.26} and~\ref{maps:src:D.2.31}.
The real-period map sends \(x\bmod\Lambda\) to
\([\Pi(t)x]\) in the fibre, retaining the nonconstant periods
and line units. Its inverse is the group logarithm modulo the same
period lattice. The rigidified line over the origin fixes its identity.
After iteration the period matrix is \(\Pi(t)S_d\).

Retract the punctured page along its fixed cut tree. Parallel
coordinates on a cut strip and their inverse logarithms identify
its torus bundle with the mapping cylinders of its literal linear
returns. The universal cover is \(\mathbb R^4\) over the universal
covering tree, hence contractible. Its deck group is the semidirect
group below. These maps identify the exterior with this
\(K(G_{\rm ext},1)\) model up to homotopy, preserving the fibre and all based
peripheral maps.

Write $\mathfrak m_i$ for the based peripheral words in the free
base group, and $\mathfrak m$ when a site is fixed. Every
origin-disc section is exactly $(0,\mathfrak m_i)$. The final
section is the inverse product of the earlier ones; their
semidirect product still has zero translational coordinate.
Its suffix matrices are linear transport along the chosen
stems. Thus the normal torus of a fixed \(\lambda\) is
the commuting pair \((\lambda,1),(0,\mathfrak m_i)\), with its marked
suffix transport. This proves the marking used below.

Let \(\mathfrak F_{N-1}\) be the free group on the first
\(N-1\) based peripheral words. Define the semidirect group
\[
       G_{\rm ext}=\Lambda\rtimes \mathfrak F_{N-1},\qquad
       (x,\mathfrak w)(y,\mathfrak z)
          =(x+A_{\mathfrak w}y,\mathfrak w\mathfrak z),
\]
where \(x,y\in\Lambda\) and \(\mathfrak w,\mathfrak z\in \mathfrak F_{N-1}\).
The linear return \(A_{\mathfrak w}\) is the ordered word action.

Use the period exponential and its inverse above. In the linear coordinates on the cut page, fibre inversion is
\[
 \iota(x,\mathfrak w)=(-x,\mathfrak w),\qquad
 (x,\mathfrak w)(y,\mathfrak z)=(x+A_{\mathfrak w}y,\mathfrak w\mathfrak z).
\]
The first map is a group automorphism, since
\[
 -(x+A_{\mathfrak w}y)=-x+A_{\mathfrak w}(-y).
\]
It preserves all zero-origin peripheral sections, has itself as inverse,
and transports through the period exponential to the exterior
involution. On the fibre two-cycle it sends
\((-x)\wedge(-y)\) to \(x\wedge y\); on an edge cylinder it reverses
the fibre loop and fixes the base direction. The cellular filtration over the base vertex and edges has
degree-two fibre quotient \(C_{d}\) and edge-cycle kernel \(K_{d}\).
The edge boundary is \([A_i(d)-I]_{i<N-1}\); its two-step filtration gives
the exact sequence below, with no higher base cells. Thus the
involution acts on its two ends respectively \(+1\) and \(-1\):
\begin{equation}
 0\longrightarrow C_d\longrightarrow H_2(V_d;\mathbb Z)
       \xrightarrow{\epsilon}K_d\longrightarrow0.
 \hypertarget{tag-0230}{}\label{maps:eq:D.OW5}
\end{equation}
The maps here are the fibre inclusion and the edge projection, with
the loop-cylinder and addition-triangle lifts above.

Write \(\mathfrak a\) for evaluation by the pullback of the integral
weight-one marked class. At weight one that class is invariant under
inversion: the fibre area is fixed, and each generating fixed sweep
has value zero. The weight-one fibre and those sweeps generate its
whole second homology. The covering intertwines inversion,
so \(\mathfrak a\iota=\mathfrak a\). On the fibre subgroup,
\[
 \mathfrak a|_{C_d}=g\omega'_d,\qquad
 \omega'_d=(a/g)e_1^*\wedge e_2^*
                   -(db/g)\delta^*\wedge\gamma^* .
\]
For every \(h_{\rm ext}\in H_2(V_d;\mathbb Z)\), \(h_{\rm ext}+\iota h_{\rm ext}\) lies in \(C_d\), and hence
\[
        2\mathfrak a(h_{\rm ext})=\mathfrak a(h_{\rm ext}+\iota h_{\rm ext})\in g\mathbb Z.
\]
Since \(g\) is odd, \(\mathfrak b=\mathfrak a/g\) is an integral
homomorphism. Choose the fibre class \(c_0\) with free coordinate one,
using the Bézout section of the explicit quotient
\(C_d=\mathbb Z\oplus\mathbb Z/g\).

For \(k_{\rm gr}\in K_d\), first choose a cylinder lift \(h_{\rm ext}\), then set
\(x_{\rm ext}=h_{\rm ext}-\mathfrak b(h_{\rm ext})c_0\). It has edge \(k_{\rm gr}\) and
\(\mathfrak b(x_{\rm ext})=0\). The element \(t_{\rm fib}=x_{\rm ext}+\iota x_{\rm ext}\) belongs to the
torsion subgroup of \(C_d\). The inverse of multiplication by two
on that subgroup is multiplication by the integer
\((g+1)/2\) modulo \(g\). Therefore
\begin{equation}
 s_\iota(k_{\rm gr})=x_{\rm ext}-\tfrac12(x_{\rm ext}+\iota x_{\rm ext})
 \hypertarget{tag-0231}{}\label{maps:eq:D.OW6}
\end{equation}
is an integral anti-invariant lift with zero \(\mathfrak b\).
Changing \(x_{\rm ext}\) by \(c_{\rm fib}\in\operatorname{Tor}C_d\) changes \(t_{\rm fib}\) by
\(2c_{\rm fib}\), so the displayed result is unchanged. Conversely the
difference between two such lifts lies in \(C_d\), is fixed and
anti-invariant, and has zero free coordinate. Twice that difference
is zero, so it is zero. This proves uniqueness; uniqueness applied
to sums proves additivity.

The exact splitting maps are
\[
 h_{\rm ext}\longmapsto(h_{\rm ext}-s_\iota(\epsilon h_{\rm ext}),\epsilon h_{\rm ext}),\qquad
 (c_{\rm fib},k_{\rm gr})\longmapsto c_{\rm fib}+s_\iota(k_{\rm gr}).
\]
Both composites fix their displayed coordinates. They give the integral marking, including every odd torsion class.

Give \(C_{d}\) the trivial $G_{\rm ext}$-action; every \(A_{\mathfrak w}\) acts trivially on this
common coinvariant. For the same semidirect inputs define
\[
 \mathfrak c_{\wedge}((x,\mathfrak w),(y,\mathfrak z))
       =[x\wedge A_{\mathfrak w}y]\in C_d .
\]
This is the integral group 2-cocycle. On the third input
\((z_{\Lambda},\mathfrak v)\), with \(z_{\Lambda}\in\Lambda\)
and \(\mathfrak v\in \mathfrak F_{N-1}\), its coboundary is
\[
 [y\wedge A_{\mathfrak z}z_{\Lambda}
   -A_{\mathfrak w}y\wedge A_{\mathfrak w}A_{\mathfrak z}z_{\Lambda}]=0
\]
by the common coinvariant relation. Thus evaluation gives a
homomorphism
\[
                    q_d:H_2(G_{\rm ext};\mathbb Z)\to C_d .
\]

A fibre torus generated by $x,y$ is represented by the antisymmetric
bar cycle \([x|y]-[y|x]\), and the displayed integral cocycle evaluates to \(2[x\wedge y]\).
Consequently \(q_{d}\) is multiplication by two on the injected fibre
coinvariant. A fixed normal torus represented by
\((\lambda,1)\) and \((0,\mathfrak m_i)\), with \(A_i(d)\lambda=\lambda\),
has \(q_{d}\) value zero.

The cover homomorphism induced by \(S_{d}\) satisfies exact naturality:
its fibre-coordinate injection is \(\Lambda^2S_d\), and applying
that injection to the displayed integral cocycle gives the weight-one cocycle. Under
\(C_d=\mathbb Z\oplus\mathbb Z/g\), the induced map to the
weight-one fibre quotient \(\mathbb Z\) is $g$ times the free
coordinate.

For the weight-one marked graph of Theorem~\ref{maps:src:D.2.31}, the doubled cocycle is
\(2\xi_1\). Both homomorphisms have the same value on the fibre and
vanish on the complete fixed-normal image. That image maps onto
the graph kernel by the exterior construction of Theorems~\ref{maps:src:D.2.31} and~\ref{maps:src:D.2.37}, so
their difference is zero. Pulling back gives
\[
                  g\,\operatorname{free}(q_d)=2\xi_d,
              \qquad \xi_d=p_d^*\xi_1 .
\]
For odd $g$, the covering naturality identity proves that the free coordinate of \(q_{d}\) is even.
Multiplication by two is an automorphism of \(\mathbb Z/g\);
it is injective on all of \(C_{d}\). Therefore \(q_{d}\) has a unique half
\[
                    p_d^{(2)}:H_2(G_{\rm ext};\mathbb Z)\to C_d .
\]
Its explicit coordinates divide the free coordinate by two
and multiply the torsion coordinate by the inverse of two
modulo $g$. It restricts to the identity on the fibre coinvariant
and vanishes on every displayed fixed normal torus.

The cocycle is invariant under inversion, because both wedge
arguments change sign.

If \(y_{\rm ext}=s_\iota(k_{\rm gr})\), then
\(q_d(y_{\rm ext})=q_d(\iota y_{\rm ext})=-q_d(y_{\rm ext})\). Since \(C_d\) has no element
of order two, \(q_d(y_{\rm ext})=0\), and hence \(p_d^{(2)}(y_{\rm ext})=0\).
The cocycle marking therefore sends \(s_\iota(k_{\rm gr})\) to \((0,k_{\rm gr})\)
and sends a fibre class \(c_{\rm fib}\) to \((c_{\rm fib},0)\). Thus it has the same two inverse maps as the involution marking.

The edge projection gives the other coordinate in \(K_{d}\). Its
inverse chooses integral loop cylinders and vertex addition
triangles for a basis of \(K_{d}\), then subtracts a fibre cycle
representing \(p_d^{(2)}\) of that lift. Together with the fibre
inclusion this gives inverse group maps
\[
                     H_2(G_{\rm ext})=C_d\oplus K_d .
\]
The inverse uses the period cylinders on the cut strips.
A lattice-addition triangle has vertices \(0,x,x+y\) in the
real-period cover; its inverse subtracts $y$. The graph boundary
equation in $K$ closes these triangles at the vertex. Their
remaining ambiguity is a fibre cycle, removed exactly by
the subtraction of \(p_d^{(2)}\).

For a fixed vector at a nonlast site, $J_i$ is its block inclusion.
At the omitted last boundary it is
\[
 (J_{N-1}\lambda)_j
       =-A_{j+1}(d)\cdots A_{N-2}(d)\lambda .
\]
Cutting the inverse boundary word gives this formula. Its graph
boundary telescopes to zero; reading the strips in reverse is
the inverse. The inverse rows of \([R_g,K_{\mathrm{bas}}]\) compute
all these insertions in a primitive integral basis.

Every additional regular site inserts a zero block in the graph
boundary, one copy of \(\Lambda\) in $K$, and a minus-identity
block in the last insertion. Remove those coordinates, use the
base inverse, then restore them. This proves the extension for
any finite number of regular sites.

For \(\lambda_{\rm fix}\in\ker(A_i(d)-I)\), the literal fixed normal is the commuting pair
\((\lambda_{\rm fix},1),(0,\mathfrak m_i)\). Inversion reverses the first circle and
fixes the second, so the oriented torus is anti-invariant. Its marked evaluation is zero. Thus its complete coordinates are
\((0,J_i \lambda_{\rm fix})\), including a zero torsion coordinate.

When \(A_i(d)^m=I\), the \(m\)-fold mapping-torus cover has
the product sweep of any lattice vector \(\lambda\). Inversion
commutes with the covering and reverses that sweep. Its pushforward
has graph coordinate
\(N_m\lambda=\sum_{j=0}^{m-1}A_i(d)^j\lambda\) and zero fibre
coordinate by the splitting above. The inverse graph reading
cuts it into the \(m\) transported strips.

\Needspace{5\baselineskip}
\emph{The raw normal images in the same graph coordinates.}

At the logarithmic site $i$, write \(A=A_i(d)\). For its free
logarithm with \(A^m=I\) and invariant $v$, the normal
torus of a covering-lattice generator \(\lambda\) is represented by
the commuting pair
\[
              (\lambda,1),\qquad(-v,\mathfrak m^m).
\]
The two cocycle values differ by
\[
 \lambda\wedge(-v)-(-v)\wedge A^m\lambda=-2\lambda\wedge v.
\]
Its edge projection is the norm \(N_m\lambda\), obtained by
unrolling the $m$ base sectors. Thus its coordinates in the integral graph splitting are
\[
                   (-\lambda\wedge v,\ J_iN_m\lambda).
\]
For the core root generator, the two commuting elements are
\[
                        (0,\mathfrak m),\qquad(-v,\mathfrak m^m).
\]
Both cocycle values are zero; its edge projection is $v$.
Its coordinates are therefore \((0,J_iv)\).
These calculations use the normal parametrizations of Theorem~\ref{maps:src:D.2.29} and give all four lattice
columns and the root column at the integral \(C_{d}\) level.

In the abelian subgroup generated by $(v,\mathfrak m_i)$, the root
normal has ordered generators $(\mathfrak m_i,\mathfrak m_i^m-v)$.
The change from $(v,\mathfrak m_i)$ and its inverse are
\[
 \begin{pmatrix}0&-1\\1&m\end{pmatrix},\qquad
 \begin{pmatrix}m&1\\-1&0\end{pmatrix}.
\]
Their determinant is one, so the normal orientation gives
the same torus class as the fixed \(v\)-sweep.

\emph{Every finite affine action and its complete Thom source.}

Use the root-disc construction of Theorems~\ref{maps:src:D.2.29} and~\ref{maps:src:D.2.34},
with the full weighted matrix $A$. The projected-lattice calculation retains all nonunit factors of \(A-I\). Its linear order is \(h=6,4,3,4,3\) at
\(II,III,IV,III^*,IV^*\), respectively. For a primitive fixed
basis $F$ and its inverse on that lattice, put
\[
 N_h=I+\cdots+A^{h-1},\quad
 H\mathbb Z^2=\operatorname{im}(F^{-1}N_h),\quad
 Q_{\rm fin}=H^{-1}F^{-1}N_h .
\]
Column Hermite reduction with its retained preimages gives $H$.
The integer surjection $Q_{\rm fin}$ has a Smith right inverse \(R_\pi\).
Therefore the projected lattice, its coordinates and section are
\[
 P_{\rm fix}\Lambda=F(H/h)\mathbb Z^2,\quad P_{\rm fix}=N_h/h,\qquad
 F(H/h)z_{\rm fix}\longleftrightarrow z_{\rm fix},\qquad Q_{\rm fin}R_\pi=I.
\]
These maps are recomputed after weighting, including any
reduced denominators at shared component orders.

For \(m=hk_{\rm fin}>0,\ v=Fz_{\rm fix}\), define
\[
 w_{\rm fin}=hH^{-1}z_{\rm fix}=Q_{\rm fin}v,\quad t_{\rm fin}=\gcd(m,w_{{\rm fin},1},w_{{\rm fin},2}).
\]
The central action \(x\mapsto Ax+v/m\) is free exactly when
\[
 \gcd(m,v_{i,1},v_{i,2},v_{i,3},v_{i,4})=1,\qquad h\mid m/t_{\rm fin}.
\]
First, if \(d_0=\gcd(m,v_{i,1},v_{i,2},v_{i,3},v_{i,4})>1\), the power
\(j=m/d_0\) translates by a lattice vector and fixes the
origin, regardless of its linear part. This proves necessity
of the first condition. Under that condition, no proper power
with \(h\mid j\) fixes a point, because its linear part is
identity and \(jv/m\notin\Lambda\).
For \(h\nmid j\), \(A^j-I\) is invertible on the moving plane
\(W=\ker P_{\rm fix}\). A fixed point of the $j$th power exists precisely
when \(jv/m\in P_{\rm fix}\Lambda\), equivalently \(jw_{\rm fin}/m\in\mathbb Z^2\).
Such $j$ are exactly the multiples of \(m/t_{\rm fin}\). Excluding all
those proper powers with \(h\nmid j\) is exactly \(h\mid m/t_{\rm fin}\).
At a violating power with \(h\nmid j\), choose
\(\lambda=R_\pi(jw_{\rm fin}/m)\), then solve
\((A^j-I)x=\lambda-jv/m\) on $W$. This constructs the fixed point and the inverse projection test. An identity-linear
power with \(jv/m\in\Lambda\) fixes the origin directly.

Every affine lift with this power is included. Its projected
displacement is \(v/m\). For residual displacement
\(\alpha-v/m\in W\), translation by
\[
                    y_{\rm aff}=(I-A)^{-1}(\alpha-v/m)
\]
gives \(\tau_{y_{\rm aff}}^{-1}(A+\alpha)\tau_{y_{\rm aff}}=A+v/m\).
The inverse translation is \(-y_{\rm aff}\). For the holomorphic
root-disc family let $R$ include its root rotation and linear
return, so \(R^m=I\). The residual holomorphic affine term
\(b_0\) has norm zero, and
\[
 y_0=-\frac1m\sum_{j=0}^{m-1}jR^j b_0,\qquad
                    (I-R)y_0=b_0.
\]
This is the finite-operator identity in the proof of
Theorem~\ref{maps:src:D.2.29}, with the same operator $R$, order
$m$, norm-zero holomorphic term $b_0$ and root coordinate.
The identification fixes every term of the finite sum and has
itself as inverse. Translation by $y_0$ and by $-y_0$ gives the
holomorphic conjugacy and its inverse. The line gauges and root
collars retain their nonzero units.

The exact core degree-one presentation is
\[
 H_1(F_{\mathrm{core}})=
 (\Lambda\oplus\mathbb Zc_{\rm core})/
 \langle((A-I)\lambda,0),\,(-v,m)\rangle .
\]
A lifted path records its lattice displacement and final root
sheet; the inverse concatenates that lattice path with the
recorded number of root paths. Their changes are precisely the conjugation and power
relations displayed in the core presentation. All nonunit factors of \(A-I\)
and all core torsion are retained.

The normal bundle has finite holonomy, so its first Chern
class is torsion and its cap product to \(H_0=\mathbb Z\)
vanishes. The whole boundary kernel is the Thom image of
the core presentation. That source is generated by the four covering-lattice
loops and the root loop, with all its torsion generators.
Use the positive-root convention of
Theorem~\ref{maps:src:D.2.29}, with boundary identification
$(y,\tau)\sim(Ay,\tau+1)$ for this same $A,m,v$.
The boundary map and its inverse are
\[
 [y,\tau]\longmapsto[y+(\tau/m)v,e^{2\pi i\tau/m}];
 \quad\text{choose that root argument and subtract }(\tau/m)v.
\]
Its positive normal circle is \((y-sv,\tau+ms)\).
The covering-loop normal torus is the commuting pair computed above. The root normal has parametrization
\([s_{\rm root}v/m,e^{2\pi i(s_{\rm root}/m+t_{\rm nor})}]\),
where \(s_{\rm root},t_{\rm nor}\) are its two oriented circle parameters.
It has zero fibre parallelogram and edge displacement $v$.
This proves all five columns.

On a source relation \((A-I)\lambda\), the norm vanishes
and the fibre wedge is a return relation because \(Av=v\).
On \(mc_{\rm core}-v\), the Thom image is
\[
       m(0,J_i v)-(-v\wedge v,J_iN_mv)=0.
\]
Thus these columns induce the full map from the core presentation, including
its torsion. At a regular site the same proof has
\(A=I,h=1,N_m=mI\) and the sole freeness condition
\(\gcd(m,v_{i,1},v_{i,2},v_{i,3},v_{i,4})=1\).

Free data exist at every position.
For a primitive invariant $v$, the gcd of $w_{\rm fin}=Q_{\rm fin}v$ divides $h$:
$hz_{\rm fix}=Hw_{\rm fin}$ and primitive $z_{\rm fix}$ prove this. In particular
\(v=\delta,m=h^2\) satisfies the two freeness conditions. This applies to the projected lattice at every weight.

\Needspace{5\baselineskip}
\emph{Peripheral shears, inverse markings and clutch.}

Changing a peripheral generator to \((\lambda,\mathfrak m_i)\)
changes its fixed normal by
\[
                  t_i(\lambda_{\rm fix})\longmapsto t_i(\lambda_{\rm fix})+[\lambda_{\rm fix}\wedge\lambda].
\]
The cocycle evaluates to \(2[\lambda_{\rm fix}\wedge\lambda]\), and its unique
half proves this in \(C_{d}\), including torsion. The inverse
subtracts the same wedge. The power changes by
\(v\mapsto v+N_m\lambda\), by multiplying all $m$ semidirect
factors; the inverse subtracts this norm. At a regular
site the norm is \(m\lambda\).

The normal columns agree under this exact change because
\[
          [N_m \lambda_{\rm fix}\wedge\lambda]=[\lambda_{\rm fix}\wedge N_m\lambda]
\]
in the coinvariant. For the root column also use
\([N_m\lambda\wedge\lambda]=0\): pair indices $j$ and $m-j$
to show that twice this class is zero, then use the
absence of two-torsion in \(C_{d}\). The transformed root Thom image
is the initial root image plus the Thom image of \(\lambda\).
The sweep and norm identities give exactly that generator change
and its subtraction inverse. Thus these affine changes
retain torsion-valued section corrections.

The order-one operation with \(v_*=-e\gamma\)
has boundary shear
\((y,\theta)\mapsto(y-\theta v_*,\theta)\), inverse addition.
Its relation \(c_*=\gamma^{-e}\) changes the canonical
product relation to \(\prod_i c_i=\gamma^e\) on deleting
the dummy root generator. Keeping its five columns
implements every integer clutch, including zero.

\emph{The period cover and its transfer.}

Let \(\psi\) be the original contracting arrow on the line complement.
Use its given infinite toric completion at a Mumford place. The arrow
permutes those charts by the marked integral periods and their
nonzero units. Replacing \(\langle\psi\rangle\) by
\(\langle\psi^d\rangle\) gives the map
\[
 p_d:U_d\longrightarrow U_1,\qquad [x]_{\psi^d}\longmapsto[x]_\psi .
\]
On a small lifted chart the inverse sheets are \(x,\psi x,\ldots,
\psi^{d-1}x\). The positive contraction roof makes the action free
and proper before completion; on the completion the same chart
translations give the free period-subgroup action. These maps and
inverse sheets use the arrow. On the line its iterate multiplier
is the product of the $d$ translated weight-one multipliers.

The nearby fibre maps by its real periods. Its lattice injection is
\[
 S_d=\operatorname{diag}(1,1,1,d).
\]
The exponential of the period matrix is the forward map;
the group logarithm modulo that period lattice is its inverse.
Thus the square of nearby-fibre inclusions and the covering \(p_d\) commutes as
maps of spaces, including the marked angular and long directions.
This construction applies at every positive iterate.

For a singular simplex in a target chart, sum its $d$ inverse-sheet
lifts. Compatibility on common faces defines the integral transfer
\(\operatorname{tr}_U\). The same construction on the nearby torus \(\mathcal F_1\)
gives
\[
 \operatorname{tr}_F=d(\Lambda^2S_d)^{-1}
 \quad\text{on }H_2(\mathcal F_1).
\]
This matrix is integral: its entries are $d$ on the three nongamma
bivectors and one on the three \(\gamma\) bivectors. If \(f_d\) and
\(f_1\) denote the fibre specializations, the exact identities are
\[
 p_{d*}f_d=f_1\Lambda^2S_d,\qquad
 \operatorname{tr}_Uf_1=f_d\operatorname{tr}_F .
\]
Moreover
\[
 \operatorname{tr}_U p_{d*}f_d=d f_d.
\]
Indeed the left side sums the deck translates of the fibre cycle;
on that torus every deck map is translation by a fractional long
period, hence is homotopic to the identity by the explicit straight
translation isotopy. This establishes the transfer identity on the fibre image.

The cap calculations below give free integral targets for \(f_d\). Consequently
\[
 \ker f_d
 =\{\eta\in\Lambda^2\Lambda:
       (\Lambda^2S_d)\eta\in\ker f_1\otimes\mathbb Q\}.
\]
For the nontrivial inclusion, a vector in the right side maps to
zero under \(f_1\), because \(H_2(U_1)\) is free. The covering-square and transfer identities give \(d f_d(\eta)=0\), and freeness of \(H_2(U_d)\)
then gives \(f_d(\eta)=0\). The reverse inclusion follows directly
from the first equation in the two covering-square identities.

Here is an integer implementation with inverse coordinates. Choose
an integer quotient \(Q_1\) with kernel \(\ker f_1\), using its
primitive basis below. Write \(r_{\rm SNF}\) for the number of nonzero
Smith entries and compute the reversible decomposition
\[
 U_{\rm SNF}(Q_1\Lambda^2S_d)V_{\rm SNF}=\operatorname{diag}(d_1,\ldots,d_{r_{\rm SNF}},0).
\]
The last columns of $V_{\rm SNF}$ are the full kernel, and the corresponding
last rows of \(V_{\rm SNF}^{-1}\) are its inverse coordinates. Every zero
Smith coordinate is retained; the nonzero values \(d_i\) are
retained in the map, with their values. The equivalent
saturated-image construction gives a quotient \(Q_d\), section
\(\sigma_d\), kernel basis \(\mathsf K_d\), and inverse \(\mathsf L_d\) satisfying
\[
 Q_d\sigma_d=I,\quad \mathsf L_d\mathsf K_d=I,\quad
 \mathsf K_d\mathsf L_d+\sigma_dQ_d=I.
\]
The cover and transfer on these fibre-image coordinates are
\[
 Q_1\Lambda^2S_d\,\sigma_d,\qquad
 Q_dd(\Lambda^2S_d)^{-1}\sigma_1.
\]
Their two composites are multiplication by $d$. On their respective
images, division by $d$ in these displayed composites gives the
inverse rational coordinates, with integrality checked by the reversible kernel Smith map.

\emph{The cusp: its free basis and each normal disc.}

At the selected \(I_{n_*}\), let
\[
 C_{\rm cusp}=\begin{pmatrix}n_*&b\kappa\\
       -a\kappa&(1-ab\kappa^2)/n_*\end{pmatrix},\qquad \det C_{\rm cusp}=1,
 \quad S_{\mathrm{sh}}=(u_{\rm T},\delta),\quad S_{\mathrm{lg}}=(r_{\rm T},\gamma).
\]
Write \(\varepsilon_1^{(2)},\varepsilon_2^{(2)}\) for the standard
two-coordinate columns used by the short and long period maps.
The infinite fan is the original fan in this marking. Its
period subgroup has columns \(C_{\rm cusp}\varepsilon_1^{(2)},d C_{\rm cusp}\varepsilon_2^{(2)}\). It has $d$ normalized
degree-six del Pezzo components.

The normalization complex of Theorem~\ref{maps:src:D.2.34} is defined over
\(\mathcal R_d=\mathbb Z[t_{\rm ring}]/(t_{\rm ring}^d-1)\), with the edge-to-curve map
\[
 \begin{pmatrix}-1&-1&-1&-1\\
 1&1&1&t_{\rm ring}\\-t_{\rm ring}&-t_{\rm ring}&-1&-t_{\rm ring}\end{pmatrix}.
\]
Identify this coefficient ring with $\mathcal A_d$ in
Theorem~\ref{maps:src:D.2.34} by $[p(t_{\rm ring})]\mapsto[p(t_{\rm cyc})]$, with inverse $[p(t_{\rm cyc})]\mapsto[p(t_{\rm ring})]$.
Send every ordered component, branch and dual-torus cell with
coefficient $t_{\rm ring}^j$ to the identically indexed cell with that
coefficient $t_{\rm cyc}^j$. The displayed matrix, the degree-zero dual-torus
complex and the degree-four group-ring row agree under this
chain map. The unit and $t_{\rm ring}-1$ reductions, their partial-sum
inverse and the free
extension section proved there therefore give
\[
 H_*(U_d)=
 (\mathbb Z,\mathbb Z^2,\mathbb Z^{d+3},
  \mathbb Z^2,\mathbb Z^d,0,0).
\]
These operations are over integers and hold for every $d$.

In this marking the ordered basis is
\[
 \begin{aligned}
 \mathsf e_1^{\rm C}&=S_{\mathrm{sh}}C_{\rm cusp}\varepsilon_1^{(2)},&
 \mathsf e_2^{\rm C}&=S_{\mathrm{lg}}\varepsilon_1^{(2)},\\
 \mathsf e_3^{\rm C}&=S_{\mathrm{sh}}C_{\rm cusp}\varepsilon_2^{(2)},&
 \mathsf e_4^{\rm C}&=S_{\mathrm{lg}}\varepsilon_2^{(2)}.
 \end{aligned}
\]
In coordinates relative to \((u_{\rm T},r_{\rm T},\delta,\gamma)\), its map and inverse read
\[
 \begin{gathered}
 \binom{x_u}{x_\delta}=C_{\rm cusp}\binom{x_1^{\rm C}}{x_3^{\rm C}},
 \qquad x_r=x_2^{\rm C},\quad x_\gamma=x_4^{\rm C},\\
 \binom{x_1^{\rm C}}{x_3^{\rm C}}=C_{\rm cusp}^{-1}\binom{x_u}{x_\delta}.
 \end{gathered}
\]
The return in this basis is \(I+E_{12}+d E_{34}\).
Put \(\mathsf e_{jk}^{\rm C}=\mathsf e_j^{\rm C}\wedge\mathsf e_k^{\rm C}\).
The following \(y_{jk}\) are coordinates in its wedge basis, in order
\(12,13,14,23,24,34\):
the relation image has the primitive basis
\[
 r_1=\mathsf e_{13}^{\rm C},\quad r_2=\mathsf e_{14}^{\rm C}+d \mathsf e_{23}^{\rm C}.
\]
Preimages are \(\mathsf e_{23}^{\rm C}\) and \(\mathsf e_{24}^{\rm C}-d \mathsf e_{23}^{\rm C}\). The quotient,
section, and relation-coordinate inverse are
\[
 Q_{\rm cusp}(y)=(y_{12},y_{23}-d y_{14},y_{24},y_{34}),\qquad
 \sigma_{\rm cusp}(x)=x_1\mathsf e_{12}^{\rm C}+x_2\mathsf e_{23}^{\rm C}+x_3\mathsf e_{24}^{\rm C}+x_4\mathsf e_{34}^{\rm C},
\]
\[
 y-\sigma_{\rm cusp} Q_{\rm cusp}(y)=y_{13}r_1+y_{14}r_2.
\]
The boundary has \(H_2=\mathbb Z^4\oplus\mathbb Z^2\): its last
two coordinates are the marked short normal sweeps. In degree
one it has
\[
 H_1(\partial U_d)=
 \mathbb Z\langle \mathsf e_2^{\rm C},\mathsf e_4^{\rm C},\mu\rangle
          \oplus(\mathbb Z/d)\langle \mathsf e_3^{\rm C}\rangle .
\]
The two short normal sweeps bound their extending angular
thimbles. They are primitive in boundary \(H_2\), since the Wang
projection reads their two fixed coordinates. Duality gives
\(H_3(U_d,\partial U_d)=H^3(U_d)=\mathbb Z^2\). Its image contains
these two independent primitive classes. After removing them the
boundary quotient is free; a further subgroup of rank zero is
zero. They are therefore the entire integral boundary kernel.
The fibre quotient with coordinates $Q_{\rm cusp}$ and section $\sigma_{\rm cusp}$
injects into $H_2(U_d)$.

For completeness the component-disc map and a free complement
can be read explicitly. Represent the component vertices by
\[
 v_{j,\rm C}=j C_{\rm cusp}\varepsilon_2^{(2)},\qquad 0\leq j<d.
\]
The positive toric normal disc at the vertex \(v_{j,\rm C}\) has boundary
\[
                    \mu+j S_{\mathrm{sh}}C_{\rm cusp}\varepsilon_2^{(2)}.
\]
To verify the sign and both coordinates, the normal ray in the
height-one toric lattice is \((v_{j,\rm C},1)\). A character
\((m_{\rm ch},k_{\rm ht})\) restricts to the normal disc, with coordinate
\(t_{\rm disc}\), as a nonzero unit times
\(t_{\rm disc}^{\langle m_{\rm ch},v_{j,\rm C}\rangle+k_{\rm ht}}\). On the positive boundary its
winding is that exponent. Pairing with the dual height and short
basis gives exactly the ray-character normal-disc map; conversely these pairings recover the
ray and the positive disc coordinate. A unit extends nonvanishingly
over the disc and contributes zero winding. Replacing \(v_{j,\rm C}\) by
a period translate changes the ray-character normal-disc map by an image of \(A_*(d)-I\), so
the boundary class is independent of this representative.

In the cusp boundary degree-one marking, the $d$ columns are \(\mu+j\mathsf e_3^{\rm C}\).
Their kernel consists of the vectors $x$ with
\[
 \sum_jx_j=0,\qquad \sum_jj x_j=0\pmod d .
\]
Let \(\varepsilon_j^{\rm disc}\) be the coefficient vector of the
\(j\)th positive component disc. For \(d>1\), an integral kernel basis is
\[
 \begin{gathered}
 d(\varepsilon_1^{\rm disc}-\varepsilon_0^{\rm disc}),\\
 \varepsilon_j^{\rm disc}-j\varepsilon_1^{\rm disc}
        +(j-1)\varepsilon_0^{\rm disc}\quad(2\leq j<d).
 \end{gathered}
\]
The inverse reads \(x_j\) for \(j\geq2\) and
\((\sum jx_j)/d\) for the first coefficient. The disc differences
in these combinations close using the boundary relation
\(d\mathsf e_3^{\rm C}=(A_*(d)-I)\mathsf e_4^{\rm C}\); the corresponding monodromy cylinder is
an explicit bounding chain. Together with the four fibre
classes specified by $Q_{\rm cusp}$ and $\sigma_{\rm cusp}$, they give a free integral basis of
\(H_2(U_d)\). Reading relative disc coordinates, subtracting
these closed chains, and applying $Q_{\rm cusp}$ is the inverse.

In the exterior, \(e_1,e_2\) vanish and the physical angular
class has order \(g=\gcd(a,d)\). The ray-character normal-disc map becomes
\[
                   \mu+j c_{22}\delta,\qquad
 c_{22}=(1-ab\kappa^2)/n_* .
\]
Naturality at the selected zero gives \(\gcd(a,n_*)=1\):
for \(n_*>1\), \(a\kappa\equiv\pm1\pmod {n_*}\), and
the \(n_*=1\) assertion is immediate. Thus
\(n_*c_{22}\equiv1\pmod g\), so \(\gcd(c_{22},g)=1\).
Choose \(s_{\rm Bez},t_{\rm Bez}\) with \(s_{\rm Bez}c_{22}+t_{\rm Bez}g=1\). When \(g>1\),
the difference of disc 1 and disc 0, multiplied by $s_{\rm Bez}$,
together with $t_{\rm Bez}$ times the known \(g\delta\) boundary, gives
exactly \(\delta\). The inverse coefficients are these Bézout
coefficients. When \(g=1\), the angular class is already zero.
This is the full angular boundary map used in the finite-first
sequence. Its different columns are retained.

\emph{Zero-Chern originals and all rank-one Mumford caps.}

At a zero-Chern original, use the frames on each
normalized component, their node constants, and the translated line units. The line and arrow are flat on the
central mapping-torus core. On a normalized \(\mathbb P^1\),
the ratio of two nonvanishing holomorphic line maps is constant.
Agreement through the node identifications makes the ratio
constant around a multiplicative cycle as well. Its angular
restriction has zero graph-loop winding. For a $d$-fold arrow the
ratio is the product of those $d$ constants, so the same statement
holds for every positive iterate.

Here are the gauges. On the universal cover of a graph
circle with flat holonomy $h_{\rm hol}$, choose $\ell_{\rm hol}$ with
$e^{\ell_{\rm hol}}=h_{\rm hol}$ and replace the line coordinate $w_{\rm line}$ by
$e^{-s_{\rm gau}\ell_{\rm hol}}w_{\rm line}$.
The inverse multiplies by $e^{s_{\rm gau}\ell_{\rm hol}}$.
For an additional return holonomy $h'_{\rm hol}$, choose
$e^{\ell'_{\rm hol}}=h'_{\rm hol}$ and use
$e^{-s_{\rm gau}\ell_{\rm hol}-t_{\rm gau}\ell'_{\rm hol}}w_{\rm line}$.
Here $s_{\rm gau},t_{\rm gau}$ are the two real deck coordinates.
The two scalar holonomies commute, so this
descends under both deck generators. On a tree, start with one
frame and propagate it along the node constants; the inverse
restores those constants. Extend the prescribed angular frames
to the plumbing by a connection agreeing with them on the
central core. The endpoint ratio has the zero winding just
proved and hence a global logarithm.

For the line mapping torus, write $T_{\rm base}$ for its base
return and $u_{\rm hol}(x)$ for its nonvanishing endpoint line multiplier
at a plumbing point $x$. Its endpoint is
\((w_{\rm line},x,1)\sim(u_{\rm hol}(x)w_{\rm line},T_{\rm base}x,0)\).
The map
\[
 [w_{\rm line},x,s]\longmapsto
   [e^{\zeta(s)\log u_{\rm hol}(x)}w_{\rm line},x,s]
\]
removes the endpoint line multiplier. Here \(\zeta\) is smooth,
zero near 0 and one near 1. The inverse uses the negative
exponent. Constancy at the two ends makes these maps smooth
in the mapping-torus collar coordinates. Let \(r_d(x)\) be the sum of the $d$ positive contraction roofs.
On the compact plumbing collar choose \(c_{\rm roof}>0\) with
\(c_{\rm roof}\max r_d<1\), and a smooth nonnegative function
\(\phi\) supported in \((0,1)\), with integral one.
Put \(\Phi(s)=\int_0^s\phi(t)\,dt\). Here $s$ is the unit-interval
coordinate, and $t$ in the roof map ranges from $0$ to $r_d(x)$:
\[
 f_x(t)=c_{\rm roof}t+(1-c_{\rm roof}r_d(x))\Phi(t/r_d(x)),\qquad
              0\leq t\leq r_d(x).
\]
Then \(f_x(0)=0\), \(f_x(r_d(x))=1\), and
\[
 f_x'(t)=c_{\rm roof}+(1-c_{\rm roof}r_d(x))\phi(t/r_d(x))/r_d(x)>0.
\]
The roof map is \((x,t)\mapsto(x,f_x(t))\), with inverse
\((x,s)\mapsto(x,f_x^{-1}(s))\).
Near the two ends its coordinates are \(c_{\rm roof}t\) and
\(1-c_{\rm roof}(r_d(x)-t)\), so the mapping-torus transition has the
same constant slope $c_{\rm roof}$ on both sides. The maps therefore
descend smoothly and carry the variable roof. Therefore
these are maps on the original germ, with the chosen constant
units and all $d$ iterates retained.

The central fibre is an integral homology retract of its
plumbing. Its inclusion commutes with the translation, so the
induced map of mapping-torus cone complexes is an integral
homology isomorphism. An additive central fibre has
\(H_1=0,\ H_2=\mathbb Z^{s_i}\); the translation permutes its
component basis. If \(\rho(d)\) is the number of orbits of the
permutation raised to $d$, the resulting cap groups are
\[
 H_1=\mathbb Z^2,\quad H_2=\mathbb Z^{\rho(d)+1},\quad
 H_3=\mathbb Z^{2\rho(d)},\quad H_4=\mathbb Z^{\rho(d)}.
\]
The basis is the component coinvariant basis, the angular-return
torus, and the corresponding invariant orbit sums in the next
degree. On each orbit, the coinvariant map sums coefficients;
an inverse picks its specified representative. The invariant
map repeats a coefficient on the orbit, and its inverse reads
one representative. These give all free bases in the additive orbit homology.
The orbit count is taken after raising the permutation to $d$.

The nearby elliptic area specializes to the whole central
fibre \(\sum_C\mu_C[C]\), followed by the orbit quotient.
Its coefficient on an orbit $O$ is
\(\sum_{C\in O}\mu_C>0\). The fixed angular-return torus
\(\delta\wedge\ell_{i,d}\) has angular-return coefficient \(g_{i,d}>0\),
where \(g_{i,d}\) is its primitive return length. The flat angular
gauge gives its product map; projecting to that coordinate and
dividing by \(g_{i,d}\) on its image is its inverse. The component
coefficients are read in the orbit basis. These two images are
independent. The rational fibre coinvariant has dimension two,
so its map to the cap is injective. This applies to every
zero-Chern additive original, including the original odd stars.

For a multiplicative fibre, its central cycle has
\(H_1=\mathbb Z\), preserved by translation, and
\(H_2=\mathbb Z^\nu\). The same cone maps and flat angular
gauge give
\[
 H_1=\mathbb Z^3,\quad H_2=\mathbb Z^{\rho(d)+3},\quad
 H_3=\mathbb Z^{2\rho(d)+1},\quad H_4=\mathbb Z^{\rho(d)}.
\]
The additional coordinates are the graph circle times the two
residual periods and their torus product. Cutting that circle
at the specified node and reading these product coordinates
is the inverse.

At a narrow unselected \(I_\nu\), the natural row has
\(\nu\mid a\) and \(\gcd(b\kappa,\nu)=1\). Put
\[
 r_0=\gcd(d,\nu),\quad u'_N=u_{\rm T}-(a\kappa/\nu)\delta,\quad
 A_i(d)-I=u'_N\otimes(\nu r_{\rm T}^*+db\kappa\gamma^*).
\]
Choose $x_N,y_N$ with
\((\nu/r_0)x_N+(db\kappa/r_0)y_N=1\). The integral period
basis and its long inverse are
\[
 (l_N,h_N,p_N,w_N)=
 (u'_N,x_N r_{\rm T}+y_N\gamma,\delta,-(db\kappa/r_0)r_{\rm T}+(\nu/r_0)\gamma),
\]
\[
 (r_{\rm T},\gamma)=(h_N,w_N)
 \begin{pmatrix}\nu/r_0&db\kappa/r_0\\-y_N&x_N\end{pmatrix}
 .
\]
The short inverse restores \((a\kappa/\nu)\delta\).
The determinant is one and the return is \(I+r_0E_{12}\).
These formulas change the periods of the chart model;
the existing Tate and line units are carried by the same
period exponential and its inverse.

In the one-dimensional fan the period is \(r_0\) primitive
edges. Its normalized ruled components are pullbacks of the flat degree-zero ruled component of Theorem~\ref{maps:src:D.2.31}; they have the same
flat transition characters restricted to the new period
subgroup. The flat gauges above trivialize each ruled bundle
topologically, preserving its zero and infinity sections.
Choose successive translations on the residual torus to
remove the first \(r_0-1\) node translations. The remaining
translation has a lift \(\tau_{\rm node}\in\mathbb R^2\); the map
\((y,s)\mapsto(y+s\tau_{\rm node},s)\), with inverse subtraction,
removes it on the gluing circle. Constant nonzero multipliers
on the ruled coordinate are removed by their logarithm
isotopies, fixing zero and infinity. Thus the normalization
chain diagram, with these maps and inverse homotopies, is
the diagram of \(I_{r_0}\times T^2\).

Its degree-two specialization is explicitly
\[
 l_N\wedge h_N\longmapsto\sum_{j=0}^{r_0-1}[C_j],\quad
 p_N\wedge w_N\longmapsto[T^2],\quad
 h_N\wedge p_N,\ h_N\wedge w_N\longmapsto
       [\text{graph}]\times p_N,\ [\text{graph}]\times w_N.
\]
The long path crosses each of the \(r_0\) edges once, hence is
the primitive graph circle. The first image has every component
coefficient equal to one. These four classes are a primitive
direct coordinate subgroup. A section reads the first component
coefficient for their sum, and the three indicated product
coordinates; the remaining component coordinates are
\(x_j-x_0\). The kernel is
\(\mathbb Z\langle l_N\wedge p_N,l_N\wedge w_N\rangle\).
The return relation has image
\(r_0\mathbb Z\langle l_N\wedge p_N,l_N\wedge w_N\rangle\);
the additional boundary torsion is therefore killed in the cap,
exactly by this saturated kernel.

The ordinary zero-degree original and the marked rank-one
Mumford cap have this same normalization diagram via their
node-chart maps. These maps use the flat gauges just described,
with their inverse multiplication and inverse translations.
They give the multiplicative orbit homology with \(\rho(d)=r_0\) and local free defect
\(r_0-1\).

For non-narrow $P$, \(\nu\mid b\) and
\(\gcd(a\kappa,\nu)=1\). The period basis
\[
 p_0=\nu u_{\rm T}-a\kappa\delta,\quad
 \beta_0=x_P u_{\rm T}+z_P\delta,\quad a\kappa x_P+\nu z_P=1,\quad
 \ell_{i,d}=\gamma-(db\kappa/\nu)r_{\rm T}
\]
is unimodular for every $d$ and gives return \(I+E_{12}\).
Its inverse subtracts the long shear and applies the inverse
short matrix. The corresponding Mumford diagram has one
component and the four primitive classes displayed above; its defect is
zero. This proves freeness of its \(H_2,H_3,H_4\) for every $d$.

Changing any chosen logarithm in these gauges by \(2\pi i k_{\rm branch}\)
changes a period path by $k_{\rm branch}$ times its angular loop. On the
integer period basis this is the corresponding elementary
shear; its inverse subtracts $k_{\rm branch}$. Applying its exterior square
to the four rank-one specialization maps and its inverse to the period map preserves the
marked specialization. These paired shears preserve the marked kernel.

\emph{Nonzero-Chern originals.}

At such an additive site, the component-group order $q_{\rm comp}$ divides $b$
and \(\gcd(a,q_{\rm comp})=1\). At a multiplicative site the corresponding
facts are \(\nu\mid b,\ \gcd(a\kappa,\nu)=1\).
Thus the translation \(dQ\) remains in the identity component
for every $d$. For the respective cap, use the lifted flow in
Theorem~\ref{maps:src:D.2.32} or Theorem~\ref{maps:src:D.2.33}.
Its endpoint line ratio, logarithm and reciprocal gauge apply
to the $d$-fold endpoint composition. The component Chern vector stays
\[
                     c_{\rm Ch}=C_{p_{\rm Ch}}^*-C_0^* ,
\]
with $p_{\rm Ch}$ nonzero and both indicated components of multiplicity one.
The product map is \(E(M)\times S^1_{\ell_{i,d}}\), with
the inverse line flow of Theorems~\ref{maps:src:D.2.32} and~\ref{maps:src:D.2.33} and the inverse roof map displayed above.
For the multiplicative cycle the central angular normalization
is the phase of that lifted holomorphic flow. Its endpoint
ratio is constant around all nodes, so it has zero graph
winding at every $d$.

The component map $c_{\rm Ch}$ is primitive. If $F_{\rm mult}$ is the full multiplicity
vector, a basis of its kernel is $F_{\rm mult}$ and \(C_j\), \(j\ne0,p_{\rm Ch}\).
The inverse on \(\ker c_{\rm Ch}\) is \(x_0\), followed by
\(x_j-\mu_jx_0\), where \(\mu_j\) is the multiplicity of \(C_j\).
The Gysin/Serre chain map is exactly
\(c_{\rm Ch}:\mathbb Z^{s_i}\to\mathbb Z\), with inverse section \(C_{p_{\rm Ch}}\).
The remaining base graph coordinate, when present, is unchanged.
It follows directly that all the homology of \(E(M)\) used here
is free. For an additive tree its \(H_1,H_2,H_3\) are
\(0,\mathbb Z^{s_i-1},\mathbb Z^{s_i}\); for \(I_\nu\) they are
\(\mathbb Z,\mathbb Z^\nu,\mathbb Z^\nu\).
Taking the suspension product gives
\[
\begin{array}{c|cccc}
 &H_1&H_2&H_3&H_4\\ \hline
\text{additive}&\mathbb Z&\mathbb Z^{s_i-1}&
       \mathbb Z^{2s_i-1}&\mathbb Z^{s_i}\\
I_\nu&\mathbb Z^2&\mathbb Z^{\nu+1}&
       \mathbb Z^{2\nu}&\mathbb Z^\nu .
\end{array}
\]
All bases and inverses are the displayed component-kernel map,
its Gysin section, and the ordinary product coordinate maps.

At the additive site use its indexed elliptic return \(B_i\) and put \(z_{\rm ell}=(B_i-I)^{-1}g_i\).
Since $b z_{\rm ell}$ is integral, the basis
\[
 \mathsf R_d=(e_1,e_2,\delta,\ell_{i,d}),\quad
 \ell_{i,d}=(-db z_{{\rm ell},1},-db z_{{\rm ell},2},0,1)
\]
is unimodular, with inverse subtracting its last-column shear.
The primitive covector on its first three coordinates is
\[
             \rho_{\rm Gys}=q_{\rm comp}(a z_{{\rm ell},2},-a z_{{\rm ell},1},1).
\]
Choose \(v_{\rm Gys}\in\mathbb Z^3\) with \(\rho_{\rm Gys}(v_{\rm Gys})=1\). The Gysin
projection and its inverse show that, in addition to the
rational return image, the fibre kernel has generator
\[
                  k_{i,d}=\mathsf R_d(v_{\rm Gys},0)\wedge\ell_{i,d} .
\]
For a multiplicative original the analogous generator is
\[
                  k_{i,d}=\beta_0\wedge\ell_{i,d}
\]
from the non-narrow period basis. These kernel maps use the displayed Gysin projection and its section. They satisfy
\[
 (\Lambda^2S_d)k_{i,d}=dk_{i,1},\qquad
 \omega_d(k_{i,d})=-db/q_{\rm comp}\ \text{or}\ -db/\nu .
\]
The other component classes give local defects \(s_i-2\) and
\(\nu-2\). Taking the saturated span of the return image and the extra
generator gives the displayed quotient and kernel inverse identities.

\emph{The elliptic image and the exact angular rim.}

For each defining profile form the literal elliptic relation matrix
\[
 E_{\rm ell}=[B_0-I\mid B_1-I\mid\cdots\mid B_{N-1}-I].
\]
For each pair \(J\) of columns let \(d_J=\det E_{{\rm ell},J}\). Integer
Bézout coefficients \(c_J\), obtained by the Euclidean algorithm,
satisfy \(\sum_Jc_Jd_J=1\) on every one of the twenty-one
profiles. Let \(\iota_J:\mathbb Z^2\to\mathbb Z^{2N}\)
insert the two coordinates in those columns. Then
\[
 R_{\rm ell}=\sum_J c_J\iota_J\operatorname{adj}(E_{{\rm ell},J}),\qquad E_{\rm ell}R_{\rm ell}=I_2.
\]
The forward quotient is the zero group; the displayed \(R_{\rm ell}\)
is the inverse section of its full image. Selecting the same
elliptic matrices for each of the 172 natural rows proves
\begin{equation}
                 \mathbb Z^2/\operatorname{im}E_{\rm ell}=0.
 \hypertarget{tag-0232}{}\label{maps:eq:D.OW7}
\end{equation}
The displayed minor sections give every integral preimage.

On the two OS46/II rows the inverse is particularly explicit.
Their type-II return satisfies
\[
 B_{\mathrm{II}}-I=
 \begin{pmatrix}0&1\\-1&-1\end{pmatrix},\qquad
 (B_{\mathrm{II}}-I)^{-1}=
 \begin{pmatrix}-1&-1\\1&0\end{pmatrix}.
\]
Their two selected vanishing words, expressed in the same
elliptic lattice, are
\[
 u_{\rm van}=(1,1)=(B_{\mathrm{II}}-I)(-2,1),\qquad
 v_{\rm van}=(-4,-1)=(B_{\mathrm{II}}-I)(5,-4).
\]
The subgroup generated by \(u_{\rm van},v_{\rm van}\) has index three, with quotient
\((x,y)\mapsto[x-y]_3\). Its inclusion into the image of \(E_{\rm ell}\)
is given by these two displayed preimages. The original-star discs in the fundamental-group calculation below
supply the remaining vanishing relation. Thus the full elliptic relation
image and the subgroup of these two vanishing words are related
by an explicit index-three inclusion. \eqref{maps:eq:D.OW7} holds on all 172 natural rows.

Every return fixes \(\delta\), and a vector \(z\) in the elliptic
plane has no \(\gamma\) coordinate. Consequently
\[
 (\Lambda^2A_i(d)-I)(\delta\wedge z)
                    =\delta\wedge(B_i-I)z .
\]
For \(e_j\), insert the corresponding blocks of \(R_{\rm ell}e_j\);
the sum of these wedge relations is exactly
\(\delta\wedge e_j\). This is an integral preimage for each
mixed angular generator at every weight. The marked
multiplicative normal rim is therefore zero in the full
coinvariant even when it contains \(\mathbb Z/g\).
The geometric normal here has bounding projection and zero
long-product coordinate. An arbitrary Wang section can also
add area and long-product classes and carries those shears.

\emph{All unchanged fixed-normal images, including section corrections.}

For the cusp, angular orbit thimbles bound both fixed short
normals. In the rank-one models, the short angular thimble and
the two residual-torus families of normal discs bound the
three fixed normals. The flat gauges in the zero-Chern and rank-one constructions define these
families on the original charts; the inverse gauges restore
their exact unit factors.

At an additive original, put \(z_{\rm ell}=(B_i-I)^{-1}g_i\)
for this site's indexed return and cocycle. Its fixed lattice is
\[
 F_{i,d}=\mathbb Z\langle\delta,\ell_{i,d}\rangle,\qquad
 \ell_{i,d}=(-g_{i,d}db z_{{\rm ell},1},-g_{i,d}db z_{{\rm ell},2},0,g_{i,d}),
 \quad g_{i,d}=\operatorname{den}(db z_{\rm ell}).
\]
The inverse reads \((x_\delta,x_\gamma/g_{i,d})\). The component
orbit has this same return length: the generator section has
component order $q_{\rm comp}$ and $z_{\rm ell}$ has denominator $q_{\rm comp}$ in the defining profiles, so \(g_{i,d}=q_{\rm comp}/\gcd(q_{\rm comp},db)\). Its orbit tube closes by the
identity-component translation logarithm after \(g_{i,d}\) turns.
The extending line-circle tube bounds the \(\delta\) normal. The
orbit tube with its endpoint unit bounds the $\ell_{i,d}$ normal.
For a nonzero-Chern original, \(q_{\rm comp}\mid b\), so \(g_{i,d}=1\);
the product maps in the principal-bundle construction give the same bounding families.

At a nonzero-Chern multiplicative original, the geometric $u_{\rm T}$-normal of Theorem~\ref{maps:src:D.2.33} can have the rim term
\(j[r_{\rm T}\wedge\delta]\). Subtracting that rim gives the bounding
normal; the inverse adds it. This precise normal has bounding
projection and zero long-product coordinate. The remaining
fixed normals use the product families. An arbitrary Wang
section may also add area and long-product classes, and those
extra shears are retained when it is used.

In the common exterior quotient, every rim
\(r_{\rm T}\wedge\delta\) is exactly zero, because the four mixed unit
relations above include both elliptic-angular
coordinates. This is an integral assertion even when the
quotient contains \(\mathbb Z/g\). Hence the full global
normal image is
\[
                         (0,J_iF_{i,d}),\qquad F_{i,d}=\ker(A_i(d)-I)=F_i.
\]
The same statement holds after any marked period shear:
apply its exterior square to the fibre correction and its
linear inverse to the normal period. The formulas are inverse
changes of the whole boundary map.

Finally, the Wang sequence of a boundary mapping torus is
\[
 0\longrightarrow\operatorname{coker}(\Lambda^2A_i(d)-I)
 \longrightarrow H_2(\partial U_d)
 \longrightarrow\ker(A_i(d)-I)\longrightarrow0.
\]
Its fixed-normal section is integral. Subtracting the full
normal part of a class in the cap kernel leaves a fibre class.
The cap \(H_2\) groups in the cusp homology table, the additive orbit homology, the multiplicative orbit homology, the nonzero-Chern cap homology and the
rank-one model are free. Thus the remaining fibre kernel is
exactly \eqref{maps:eq:D.OW4}. This proves that \eqref{maps:eq:D.OW4} and the marked global fixed-normal image give the entire
global boundary-kernel image, including every torsion class
in the boundary Wang sequence. The full Wang quotient retains its torsion in this calculation.

\emph{Relative normal discs at the other unchanged caps.}

The groups \(H_2(U_i,\partial U_i)=H^4(U_i)\) are free in all
the displayed bases. A dual normal disc is chosen at a generic
point of the indicated component, or at one representative
of the indicated component orbit. The quotient is a local
biholomorphism there, so this disc intersects that orbit
cycle once. On a component of multiplicity \(\mu_C\), a local
base equation is \(q_{\rm base}=t_{\rm disc}^{\mu_C}\epsilon_{\rm unit}(t_{\rm disc})\),
where \(t_{\rm disc}\) is the normal-disc coordinate and
\(\epsilon_{\rm unit}\) is nonvanishing. The normal disc is taken in a nonzero line
frame. It has \(\gamma\) displacement zero.

The exact remaining vertical coordinate is defined by the
period inverse. Restrict the normal
disc to its positive boundary, transport it along the fixed
stem to the cut exterior, and apply the group logarithm in
the period matrix. Its endpoint difference is
\((u_{C,1},u_{C,2},\rho_C,0)\in\Lambda\). The inverse period
exponential takes this lattice path and the
\(\mu_C\)-fold meridian back to that boundary path. Equivalently
\(\rho_C\) is the winding of the marked line coordinate
after that stem transport, computed by
\((2\pi i)^{-1}\int d\log w_{\rm line}\), including the marked
period-transition jumps, in its nonvanishing cut frame.
Thus the exterior boundary column is exactly
\[
                      \mu_C x_i+\rho_C\delta .
\]
The integers \(u_{C,j}\) vanish under the proved exterior
elliptic quotient; \(\rho_C\) is retained modulo $g$. A change
of the local nonzero frame has a nonvanishing extension
over the disc and does not change its winding. Changing a
stem or its peripheral lift is instead its marked
integral period shear, applied to the general component-disc boundary with its inverse.

The general component-disc boundary supplies the component-disc matrix for
the relative splitting. The free-rank calculation uses its rank, and the integral
splitting uses all these integer columns. At the cusp they are the
explicit exterior cusp normal-disc columns. All angular entries enter the general component-disc boundary.

\emph{The complete relation matrix.}

The full target is \(\Lambda^2\Lambda\oplus K_d\).
Its matrix $B_{\rm rel}$ has all the following integer columns:
\[
       ((\Lambda^2A_i(d)-I)\eta,0),
\]
\[
 (-\lambda\wedge v_i,J_iN_{m_i}\lambda),\quad
       (0,J_iv_i)\quad\text{at a logarithmic site},
\]
\[
 (\mathcal{K}_i(d),0),\quad(0,J_iF_i)
                     \quad\text{at an unchanged cap}.
\]
\eqref{maps:eq:D.OW2} uses all six standard bivectors, \eqref{maps:eq:D.OW3} all four
standard lattice vectors, and the full unchanged-cap columns use primitive integer
bases with the inverse coordinates displayed above. Write \(\mathcal C=\operatorname{coker}B_{\rm rel}\) and
\(T=\operatorname{Tor}\mathcal C\).
The nonunit factors of \eqref{maps:eq:D.OW2} remain in this target.

\emph{Finite-first attachment retaining the angular torsion.}

Attach the $k$ logarithmic caps first and call the result $E_{\rm log}$.
At least the cusp is unfilled. Omit an unfilled base
meridian when choosing the bouquet. The exterior group is
\[
 H_1(V_d)=\mathbb Z\gamma\oplus(\mathbb Z/g)\delta
                         \oplus\mathbb Z^{N-1}.
\]
Write $F_{\mathrm{core},i}$ for the connected central core of
the $i$th logarithmic cap. Its relative degree-two group is
$\mathbb Z=H_0(F_{\mathrm{core},i})$. Its positive normal disc maps to
\[
                m_i x_i-v_{i,\gamma}\gamma-v_{i,\delta}\delta .
\]
The distinct positive coefficients \(m_i\) on their own
base coordinates prove injectivity, including the
torsion in the exterior degree-one presentation. The inverse on the image divides the
corresponding base coordinates by \(m_i\). The relative
sequence gives
\[
 H_2(E_{\rm log})=(C_d\oplus K_d)/
                    \langle\text{all logarithmic columns}\rangle .
\]
In a free presentation of the exterior degree-one group, retain \(g\delta=0\)
alongside the logarithmic normal-disc columns, and call their matrix ${\mathsf P_{\rm log}}$. Its columns
are independent: read the log base coordinates, then
the \(\delta\) coordinate. This supplies its image inverse.

For every remaining cap, duality and the cap bases above give the free abelian group
\[
                    H_2(U_j,\partial U_j)=H^4(U_j).
\]
Its rank is $d$ at the cusp, \(\rho_j(d)\) at a zero-Chern
additive original, \(s_j\) or \(\nu_j\) at a nonzero-Chern
original, and the component count at a rank-one cap.
Use the dual positive normal discs. Their boundary
columns are
\(\mu_C x_j+\rho_C\delta\), with the angular
winding defined by the inverse period map above. At the cusp they are explicitly
\[
                     x_*+j c_{22}\delta,\quad0\leq j<d,
\]
where \(c_{22}=(1-ab\kappa^2)/n_*\) is a unit modulo $g$.
The ray/character map, inverse, and the Bézout
reconstruction of the angular generator are proved
there.

The full degree-three images of the remaining caps
are exactly the full unchanged-cap columns, by their complete Wang kernels.
Thus the integral relative sequence is
\[
 0\longrightarrow\mathcal C\longrightarrow H_2(X)
 \longrightarrow\ker\{\bigoplus_jH^4(U_j)
                   \xrightarrow{\Gamma}H_1(E_{\rm log})\}\longrightarrow0 .
\]
This is \eqref{maps:eq:D.WP11} with its complete angular coordinates.

Let ${\mathsf M_{\rm disc}}$ be the component-disc boundary matrix. Solve
\({\mathsf M_{\rm disc}}z_{\rm disc}={\mathsf P_{\rm log}}z_{\rm log}\) over integers by reversible Smith operations
on \([{\mathsf M_{\rm disc}},-{\mathsf P_{\rm log}}]\). Projection of that kernel to $z_{\rm disc}$ is
injective on relation coefficients because ${\mathsf P_{\rm log}}$ is
injective. Its inverse reads the unique \(z_{\rm log}\).
A Hermite basis gives a basis of \(\ker\Gamma\)
and its integer inverse coordinates.

For each basis vector take its relative disc chain,
subtract the power discs specified by \(z_{\rm log}\), and
bound the remaining homologically zero chain in the
exterior. The \(g\delta\) relation is bounded by
the explicit Bézout monodromy chain of the degree-one image construction;
its coefficient is retained in ${\mathsf P_{\rm log}}$. Gluing gives a
closed two-cycle. Reading the relative coordinates
and subtracting these cycles is the inverse splitting.
Consequently
\[
                   H_2(X)=\mathcal C\oplus\ker\Gamma .
\]
The second summand is free. Its integer kernel retains
all torsion-index effects, while $T$ is exactly the
torsion of \(H_2(X)\).

\emph{The full fundamental group for all weights and orders.}

The cusp kills its short lattice
\(\mathbb Z\langle u_{\rm T},\delta\rangle\) and its origin-disc
meridian by the extending toric discs. On vectors of
\(\gamma\) coordinate zero, \(A_i(d)\) is literally
\(A_i(1)\). The conjugate vanishing words are therefore
the identical words in the group proof of Theorem~\ref{maps:src:D.2.37}.
On 170 rows two elliptic words have determinant one.
For \(u_{\rm van}=(u_{{\rm van},1},u_{{\rm van},2}),v_{\rm van}=(v_{{\rm van},1},v_{{\rm van},2})\), after subtracting
their killed angular coordinates, their inverse is
\[
 e_1=v_{{\rm van},2}\widetilde u_{\rm van}-u_{{\rm van},2}\widetilde v_{\rm van},\qquad
 e_2=-v_{{\rm van},1}\widetilde u_{\rm van}+u_{{\rm van},1}\widetilde v_{\rm van}.
\]
On the two OS46/II rows \((a,b)=(1,4),(4,1)\),
use \(u_{\rm van}=(1,1),v_{\rm van}=(-4,-1)\), giving \(3e_1=-u_{\rm van}-v_{\rm van}\).
Their \(I_3^*\) is original in both pattern scopes.
At \(a=1\) its unchanged Hopf/plumbing disc bounds
\(2e_1+\delta\), which gives \(e_1\) by subtraction.
At \(a=4\) its unchanged line-to-cover map
\((x,y,z)\mapsto z-2x-y\), with section \(k\delta\),
kills \(e_1+2\delta\), again giving \(e_1\).
Then \(e_2=u_{\rm van}-e_1\). The line plumbing maps are
unchanged for every $d$; only the return is iterated.

The return \(\chi_{\rm ret}=\gamma^*\) vanishes on every used
relation and extends over all original and Mumford
caps. It proves that the vertical quotient is exactly
\(\mathbb Z\gamma\), with the \(\gamma\) period as inverse.
At a logarithmic site, reading the root sheet and
return of a lifted path gives
\[
 \langle\gamma,c_i\mid[\gamma,c_i]=1,\
                      c_i^{m_i}=\gamma^{\mu_i}\rangle,
                  \qquad\mu_i=v_{i,\gamma}.
\]
The inverse records \(0\leq j<m_i\) and the remaining
return exponent \(n_{\rm ret}\in\mathbb Z\), then writes \(\gamma^{n_{\rm ret}}c_i^j\).
Boundary fibres surject onto unchanged cap groups
after the multiplicity-one meridians are contracted,
with return as inverse.

The chosen origin lifts have canonical product
return zero. The order-one clutch gives
\[
 \pi_1(X)=\langle\gamma,c_1,\ldots,c_k\mid
 [\gamma,c_i]=1,\ c_i^{m_i}=\gamma^{\mu_i},\
                         c_1\cdots c_k=\gamma^e\rangle .
\]
The forward map sends each fibre vector to its \(\chi_{\rm ret}\)
power, unchanged meridians to one, and root paths to
the displayed generators. The inverse uses the
return and root paths. Their monodromy,
vanishing, power and product relations are exactly
the displayed path relations; both composites fix
all generators.

The abelianization matrix has rows
\[
 (-\mu_i,0,\ldots,m_i,\ldots,0),\qquad(-e,1,\ldots,1).
\]
Smith maps and their inverses give
\(H_1=\mathbb Z^r\oplus A_1\). Rational elimination
gives \(r=1\) exactly when
\[
                         e-\sum_i\mu_i/m_i=0,
\]
and \(r=0\) otherwise. The same equations cover all
positive orders, allowed numerators, empty site
sets and order-one operations.

\Needspace{5\baselineskip}
\emph{Free ranks and every integral homology group.}

Over \(\mathbb Q\), the literal map \(S_d\) identifies
the graph and all fixed-peripheral maps with their
weight-one maps:
\[
 [A_i(d)-I]\operatorname{diag}(S_d^{-1},\ldots)
                         =S_d^{-1}[A_i(1)-I].
\]
The weight-one inverse of Theorem~\ref{maps:src:D.2.31} has exactly one
relation among full fixed normals, the \(\delta\) sweep
over the punctured page. The identity-block inverse
of the marked graph construction preserves this for any number of
regular sites.

The free coordinate of \(C_{d}\) is
\(\omega'_d=\omega_d/g\). Let \(s_{\rm fib}=c_0\otimes1\) be its
rational generator of free coordinate one. For the fixed-lattice input
\(\lambda_{\rm fix}\), each rational logarithmic block gives
\(t_i(\lambda_{\rm fix})=\omega'_d(\lambda_{\rm fix},v_i)s_{\rm fib}/m_i\). The root column
is redundant there because \(\omega'_d(v_i,v_i)=0\).
After deleting the dummy clutch, the one \(\delta\)
relation becomes
\[
 \frac{db}{g}\left(e-\sum_i\mu_i/m_i\right)s_{\rm fib}=0.
\]
At a nonzero-Chern original its extra
generator pairs as \(-db/(gq_{\rm comp})\) or
\(-db/(g\nu)\), a nonzero rational number. Thus
\[
 \operatorname{rank}\mathcal C=r\ (L=0),\qquad
 \operatorname{rank}\mathcal C=0\ (L>0).
\]
Only the free rank is read here; all torsion remains
the integer cokernel of the full relation matrix.

Read $L$ and $\beta$ in Definition~\ref{mapsIII:def:odd-weighted-data}
from the actual $d$th component permutations.
A narrow rank-one place contributes
\(\gcd(d,\nu)-1\), and a non-narrow Mumford place
contributes zero. A zero-Chern additive original
contributes \(\rho_i(d)-1\). A nonzero-Chern original
contributes \(s_i-2\) or \(\nu_i-2\) and one to $L$.
The cusp contributes $d$ components separately.

With \(s_{\mathrm{unf}}\) unfilled caps, the domain of
\(\Gamma\) in the integral relative sequence has rank
\[
                  d+(s_{\mathrm{unf}}-1)+\beta+L.
\]
The independent base coefficients in the logarithmic normal-disc columns give
\(\operatorname{rank}H_1(E_{\rm log})=s_{\mathrm{unf}}\). The
cokernel of \(\Gamma\) is \(H_1(X)\), already computed;
therefore \(\operatorname{rank}\Gamma=s_{\mathrm{unf}}-r\).
Its free kernel has rank
\[
                         d-1+\beta+L+r.
\]
The relative splitting and the two rank formulas give
\[
 b_2=\begin{cases}d-1+\beta+2r,&L=0,\\
                  d-1+\beta+L+r,&L>0.\end{cases}
\]

The exterior and all boundaries have Euler number
zero. Original caps have their circle or
mapping-torus maps, finite cores are torus quotients,
and rank-one diagrams have zero Euler number.
The cusp groups in the cusp homology table have Euler
number \(2d\). Additivity on the gluing
therefore gives \(\chi(X)=2d\).
Duality and UCT give
\[
 b_3=\begin{cases}2\beta+2r,&L=0,\\
                  2\beta+2L,&L>0,\end{cases}
\]
and the complete integral groups
\[
\begin{array}{c|c}
j&H_j(X;\mathbb Z)\\ \hline
0,6&\mathbb Z\\
1&\mathbb Z^r\oplus A_1\\
2&\mathbb Z^{b_2}\oplus T\\
3&\mathbb Z^{b_3}\oplus T\\
4&\mathbb Z^{b_2}\oplus A_1\\
5&\mathbb Z^r .
\end{array}
\]
For the torsion identifications, UCT gives
\(\operatorname{Tor}H^3=\operatorname{Ext}(T,\mathbb Z)\)
and \(\operatorname{Tor}H^2=\operatorname{Ext}(A_1,\mathbb Z)\).
The cyclic Smith coordinates identify these Ext
groups with the same displayed cyclic groups;
cap product with the oriented fundamental class
maps them to \(H_3,H_4\). The inverse uses the same
cyclic coordinates. Thus the integral homology table includes explicit
group isomorphisms.

\emph{The two pattern sets.}

The all-star-original set has 459 patterns. The broad set has 495,
allowing finite-good \(III^*,IV^*\) replacements while keeping odd
\(I_\nu^*\) caps original. Their inclusion fixes each labelled
choice, and the retraction replaces only a finite-good \(III^*\)
or \(IV^*\) choice by its original cap. Reading the allowed
independent choices at those positions describes every inverse
image of the retraction. Its composite with inclusion is identity.

At each of the 172 natural rows, form the product of its allowed
choices away from the selected cusp. The zero-Chern original and
Mumford choice have the same complete boundary kernel by their
node-chart and flat-gauge maps and their common saturated fibre
kernel. At a nonzero-Chern original the extra Gysin generator
has nonzero rational value \(-db/(gq_{\rm comp})\) or \(-db/(g\nu)\);
the quotient functional therefore detects its class outside the
Mumford kernel. The invertible rational map \(\Lambda^2S_d\)
preserves both assertions. Finite choices keep their order,
power vector and all five Thom columns. These rules give the
following finite product enumeration, in which the two columns
retain their literal pattern maps.

\begin{center}
\small
\begin{tabular}{lrr}
Profile & All star caps original & Broad finite-good scope \\
\hline
OS43/II & 4 & 8 \\
OS43/base & 6 & 12 \\
OS45/II & 12 & 12 \\
OS45/base & 12 & 12 \\
OS46/II & 4 & 4 \\
OS46/base & 6 & 6 \\
OS47/II & 24 & 24 \\
OS47/III & 24 & 24 \\
OS47/base & 30 & 30 \\
OS49/II & 4 & 8 \\
OS49/III & 8 & 16 \\
OS49/base & 14 & 28 \\
OS50/II & 4 & 4 \\
OS50/IV & 8 & 8 \\
OS50/base & 14 & 14 \\
OS55/II & 24 & 24 \\
OS55/base & 30 & 30 \\
OS56/II & 46 & 46 \\
OS56/III & 54 & 54 \\
OS56/IV & 54 & 54 \\
OS56/base & 77 & 77 \\
Total & 459 & 495 \\
\end{tabular}
\end{center}

The formulas for the graph, finite actions, kernel transfer and
relative attachment hold for every positive $d$ with odd $g$ and
every finite number of regular sites. The integer operations
retain every allowed order, numerator, affine power and clutch.

\emph{The role of oddness.}

Take the natural OS43/base row
\((a,b,\text{selected site})=(2,1,2)\) and \(d=2\).
The all-weight coefficient maps, obtained before taking an odd half, give
\[
 C_d=\mathbb Z^6/
 \langle e_{1\delta},e_{2\delta},e_{1\gamma},e_{2\gamma},
                       2e_{12}+2\delta\gamma\rangle.
\]
The quotient map reads
\[
        (x_{12}-x_{\delta\gamma},[x_{\delta\gamma}]_2).
\]
Its inverse is the same two-coordinate Bézout map as above.
Thus \(\tau_2=e_{12}+\delta\gamma\) has exact order two.
The graph sequence injects this nonzero class into \(H_2(V_d)\).

Choose the literal \(\delta\)-normal at a nonlast boundary.
Its edge vector \(J_i\delta\) is a primitive nonzero vector in
the integer graph kernel: its source block has \(\delta\) coordinate
one, and the graph-boundary equation is zero. Let $t_{\rm tor}$ be that
torus. Both $t_{\rm tor}$ and \(t_{\rm tor}+\tau_2\) have this edge, have zero marked
evaluation, and are anti-invariant, since
\(\iota t_{\rm tor}=-t_{\rm tor}\) and \(\iota\tau_2=\tau_2=-\tau_2\).
They are distinct by the fibre injection. Adding \(\tau_2\)
interchanges the two lifts and is its own inverse.
Thus the anti-invariance and zero-evaluation conditions can have
two different integral lifts when both $a$ and $d$ are even.
\end{proof}

\begin{remark}[Integral comparisons and finite calculations]
\hypertarget{tag-Q139}{}\label{mapsIII:rem:odd-weighted-calculations}
{\normalfont\scriptsize tag Q139}\par
Compute each saturated preimage in \eqref{maps:eq:D.OW4} before
forming the full Smith quotient. All nonunit return factors and
the complete graph coordinates remain in that presentation.
For $g=1$, its comparison with Theorem~\ref{maps:src:D.2.37}
is induced by
\[
 (\eta,k)\longmapsto(\omega_d(\eta),k),\qquad
 (n,k)\longmapsto
 \bigl(n(u e_{12}-v\delta\gamma),k\bigr),\quad ua+vdb=1.
\]
The second map is the inverse modulo the return columns, by
the explicit fibre quotient proved above. Each logarithmic
column maps to $(-\omega_d(\lambda,v_i),J_iN_{m_i}\lambda)$,
and every fixed-normal column keeps its graph coordinate.
At a nonzero-Chern original, let $q_i$ be the additive component
order, or $\nu_i$ the multiplicative index. Its fibre-kernel
generator has value $-db/q_i$ or $-db/\nu_i$; multiplying
that source generator by $-1$ gives
the positive generator used there, with the same sign change
as inverse. These maps commute with every boundary column.

These saturated kernels and full boundary columns give the
integral presentation for every allowed positive weight and all
finite, regular and clutch parameters in \eqref{maps:eq:D.OW1}.

Theorem~\ref{maps:src:D.2.41} supplies the integral algorithm
for every positive weight. The inclusion of the present input
set fixes the natural row, weight, ordered cap choices, every
logarithmic pair $(m_i,v_i)$ and the clutch.
\end{remark}

\section{Integral bar coordinates and initial pairs}
\subsection{Integral bar coordinates at every positive weight}

\begin{definition}[All-weight data and integral cochains]
\hypertarget{tag-Q140}{}\label{mapsIII:def:all-weight-integral-data}
{\normalfont\scriptsize tag Q140}\par
Choose a natural initial row of Theorem~\ref{maps:src:D.2.22},
with $P=aG$, $Q=bG$, $a,b>0$ and $\gcd(a,b)=1$, and choose
any positive integer $d$. Use its $d$th translation iterate,
norm linearization, original chart units and period subgroup.
Take either cap set $\mathcal P$ or
$\mathcal P_{\mathrm{all\ star\ O}}$ of
Definition~\ref{mapsIII:def:weighted-natural-data}, every free
finite-good affine logarithm, every finite list of free regular
logarithms and every integer clutch $e$. Keep the same formulas
for the component counts $L$ and $\beta$ as in
Definition~\ref{mapsIII:def:odd-weighted-data}, evaluated at this $d$.

Use the ordered periods
$\Lambda=\mathbb Z\langle e_1,e_2,\delta,\gamma\rangle$
of Definition~\ref{mapsIII:def:elliptic-marked-data}, with
$e_3=\delta$ and $e_4=\gamma$. Index the boundary matrices,
including added regular and auxiliary clutch sites, by
$A_0,\ldots,A_r$. These denote the weighted matrices
$A_i(d)=S_d^{-1}A_i(1)S_d$, where
$S_d=\operatorname{diag}(1,1,1,d)$ and $A_0\cdots A_r=I$.
Place an unfilled site last, and take the first $r$ base loops
as independent generators. Put
\[
 K=\ker\bigl([A_0-I\ \cdots\ A_{r-1}-I]:
              \Lambda^r\longrightarrow\Lambda\bigr).
\]
The based insertions $J_i:\operatorname{Fix}A_i\longrightarrow K$
use these same stems and the final inverse word. At a logarithmic
site write $m>0$ and $v\in\Lambda$ for its order and invariant
power, so $A_i^m=I$ and $A_iv=v$, and put
$N_m=\sum_{\nu=0}^{m-1}A_i^\nu$. At a regular site $A_i=I$.

For $A\in\{A_0,\ldots,A_r\}$ and $u,v\in\Lambda$, set
\begin{equation}
 p(u,v)=\sum_{\alpha<\beta}u_\alpha v_\beta\,
                  e_\alpha\wedge e_\beta,\qquad
 D_A(u,v)=p(Au,Av)-(\Lambda^2A)p(u,v).
                                                        \hypertarget{tag-01XR}{}\label{maps:eq:D.BC1}
\end{equation}
Define the integer-valued quadratic map
\begin{equation}
 f_A(u)=-\sum_\alpha D_A(e_\alpha,e_\alpha)
                       \binom{u_\alpha}{2}
        -\sum_{\alpha<\beta}D_A(e_\alpha,e_\beta)
                       u_\alpha u_\beta .
                                                        \hypertarget{tag-01XY}{}\label{maps:eq:D.BC2}
\end{equation}
Negative coordinates use the same integral polynomial
$\binom n2=n(n-1)/2$.
Define the peripheral correction
\begin{equation}
 \begin{split}
 c_i(u)&=f_{A_i}(u) \quad (i<r),\\
 c_r(u)&=-\sum_{j<r}
       f_{A_j}(A_{j+1}\cdots A_{r-1}A_r u).
 \end{split}                                            \hypertarget{tag-01XZ}{}\label{maps:eq:D.BC3}
\end{equation}
The transported reversed word in the second line uses the
displayed order, including the final factor $A_r$.
For commuting based loops $x,y$ in the graph exterior, write
$\mathsf T(x,y)$ for the image of the positively oriented fundamental
class of $T^2$ under the map induced by
$(1,0)\mapsto x$, $(0,1)\mapsto y$ on fundamental groups.
The ordered pair fixes its orientation.
\end{definition}

\Needspace{6\baselineskip}
\begin{theorem}[Integral homology at every positive weight]
\hypertarget{tag-02LI}{}\label{maps:src:D.2.41}
{\normalfont\scriptsize tag 02LI}\par
For every datum in Definition~\ref{mapsIII:def:all-weight-integral-data},
the following integral presentation computes the full homology.
Across the $172$ natural rows, the broad cap set has $495$
patterns and the all-star-original subset has $459$. Their
inclusion $j$ and the map $\pi$ replacing finite-good
$\mathrm{III}^*,\mathrm{IV}^*$ caps by their originals satisfy
$\pi j=1$; $j\pi$ performs precisely those replacements.

The map $D_A$ is symmetric, and $f_A$ is its integral quadratic
primitive in \eqref{maps:eq:D.BC8}. Keep every fibre relation \eqref{maps:eq:D.OW2} and every full
saturated cap kernel \eqref{maps:eq:D.OW4}. Use the five
finite or regular logarithmic normal columns
\begin{equation}
 \begin{gathered}
 \left(\sum_{\nu=0}^{m-1}c_i(A_i^\nu\lambda)
                  -\lambda\wedge v,\quad J_iN_m\lambda\right)
       \quad(\lambda=e_1,e_2,\delta,\gamma),\\
                       (c_i(v),J_iv).
 \end{gathered}                                          \hypertarget{tag-01Y0}{}\label{maps:eq:D.BC4}
\end{equation}
At every other cap the fixed-normal columns are
\begin{equation}
                         (c_i(v),J_iv),
                         \qquad v\text{ in a basis of }\operatorname{Fix} A_i.
                                                        \hypertarget{tag-01Y1}{}\label{maps:eq:D.BC5}
\end{equation}
The finite freeness test is still exactly \eqref{maps:eq:D.WP5}.
For $A_i$ of order $h\mid m$, the sum in \eqref{maps:eq:D.BC4}
is $(m/h)\sum_{\nu=0}^{h-1}c_i(A_i^\nu\lambda)$;
thus this is a finite algorithm for unbounded orders.
Let $T$ be the torsion cokernel of this full matrix.
Enumerate the logarithmic sites other than the auxiliary clutch
in base order as $(m_j,v_j)$, $1\leq j\leq k$, and put
$\mu_j=\gamma^*(v_j)$. The full fundamental group is
\[
 \pi_1(X)=\left\langle\gamma,\zeta_1,\ldots,\zeta_k\ \middle|\
 [\gamma,\zeta_j]=1,\quad
 \zeta_j^{m_j}=\gamma^{\mu_j},\quad
 \zeta_1\cdots\zeta_k=\gamma^e\right\rangle .
\]
Write its abelianization as $\mathbb Z^{r_1}\oplus\mathcal A_1$.
The free ranks and complete integral homology are
\eqref{maps:eq:D.WP8} and~\eqref{maps:eq:D.WP9}, using $r=r_1$,
$A_{1,\rm tor}=\mathcal A_1$ and the full torsion group $T$ just defined. These conclusions
hold for every $d>0$, including all shared even factors.

\end{theorem}


\begin{proof}
\emph{The integral comparison.}
Use the bar chains of $\Lambda$ with trivial coefficients and with
degenerate simplices factored out. Their low-dimensional boundaries are
\begin{equation}
 \partial[u]=0,\qquad
 \partial[u|v]=[v]-[u+v]+[u].
                                                        \hypertarget{tag-01Y2}{}\label{maps:eq:D.BC6}
\end{equation}
The standard resolution and its contracting homotopy give this
functorial model of $B\Lambda$; the contraction preserves the
degenerate subcomplex
\cite[Chapter I, Section 5, pp.~18--19]{manBrown82}.
Put
\begin{equation}
             p_1([u])=u,\qquad p_2([u|v])=p(u,v).
                                                        \hypertarget{tag-01Y3}{}\label{maps:eq:D.BC7}
\end{equation}
The first map kills two-boundaries. Bilinearity of $p$ makes
$p_2\partial[u|v|w]=0$. The cycle $[u|v]-[v|u]$ maps to
$u\wedge v$. In particular the six cycles
$[e_\alpha|e_\beta]-[e_\beta|e_\alpha]$, $\alpha<\beta$,
are the coordinate two-tori, and give the inverse on
$H_2(T^4;\mathbb Z)=\Lambda^2\Lambda$.

Antisymmetrizing \eqref{maps:eq:D.BC1} gives zero, proving the
symmetry of $D_A$. The binomial identity, valid for every pair
of integers, gives
\begin{equation}
 f_A(u)+f_A(v)-f_A(u+v)
             =p(Au,Av)-(\Lambda^2A)p(u,v).
                                                        \hypertarget{tag-01Y4}{}\label{maps:eq:D.BC8}
\end{equation}
Extend $f_A$ linearly on bar one-chains. Then
$f_A\partial=p_2A-(\Lambda^2A)p_2$ on two-chains.

The marked period maps identify the exterior with the torus bundle
over the based graph obtained by cutting along the chosen stems.
At weight $d$ their lattice map is
$S_d=\operatorname{diag}(1,1,1,d)$ and their monodromies are
$A_i(d)=S_d^{-1}A_i(1)S_d$. The period-subgroup cover sends
$[x]_{\psi^d}$ to $[x]_\psi$; on an open fundamental strip its
inverse sheets are $x,\psi x,\ldots,\psi^{d-1}x$.
The line map is the product of the original line maps on these
sheets, and its positive roof $r_d(x)$ is the sum of their
positive roofs. Use the smooth roof coordinate $f_x$ in the
zero-Chern cap part of the proof of Theorem~\ref{maps:src:D.2.40}:
the input map fixes $x$, $r_d$, the collar and its line units,
and sends $(x,t)$ to $(x,f_x(t))$. Its inverse is
$(x,s)\mapsto(x,f_x^{-1}(s))$. The identical endpoint slopes
proved there give the mapping-torus coordinates. Thus the graph
and peripheral maps carry the same line units and origin sections.

With fibre-first orientation the integral mapping-cone complex is
\begin{equation}
 {\cal C}_n=B_n(\Lambda)\oplus
                   \bigoplus_{i<r}B_{n-1}(\Lambda),
 \qquad
 \partial(c;g_i)=
 \left(\partial c+\sum_{i<r}(A_i-I)g_i;
                                     -\partial g_i\right).
                                                        \hypertarget{tag-01Y5}{}\label{maps:eq:D.BC9}
\end{equation}
For the convention $(x,1)\sim(A_i^{-1}x,0)$, a cylinder labelled
by its terminal chain $g$ starts at $A_i g$, giving the displayed
sign. The universal cover is real four-space over a tree and is
contractible. The graph of bar models has the same marked
$\Lambda\rtimes F_r$ group; the period maps induce the indicated
maps on each vertex and cylinder, hence on each peripheral torus.

On two-cycles define
\begin{equation}
 P_2(c;g_i)=
 \left(p_2c+\sum_{i<r}f_{A_i}(g_i);
                                  (p_1g_i)_{i<r}\right).
                                                        \hypertarget{tag-01XS}{}\label{maps:eq:D.BC10}
\end{equation}
The target coordinates are
\begin{equation}
 C\oplus K,\qquad
 C=\Lambda^2\Lambda/
          \sum_{i<r}(\Lambda^2A_i-I)\Lambda^2\Lambda,
 \quad
 K=\ker[A_0-I\ \cdots\ A_{r-1}-I].
                                                        \hypertarget{tag-01XT}{}\label{maps:eq:D.BC11}
\end{equation}
The entire relation lattice in $C$ is retained. On a three-boundary
the fibre coordinate of \eqref{maps:eq:D.BC10} is
\begin{equation}
 \sum_{i<r}\bigl(p_2(A_i-I)h_i-f_{A_i}\partial h_i\bigr)
       =\sum_{i<r}(\Lambda^2A_i-I)p_2h_i .
                                                        \hypertarget{tag-01XU}{}\label{maps:eq:D.BC12}
\end{equation}
Here $p_2\partial c_3=0$. Its differential to degree one also
agrees, by \eqref{maps:eq:D.BC8}, so $P_2$ induces a homomorphism.

We specify its inverse. For $k=(u_0,\ldots,u_{r-1})\in K$,
list $v_1^{\rm bar}=A_0u_0,v_2^{\rm bar}=-u_0,\ldots,v_{2r}^{\rm bar}=-u_{r-1}$ and put
$s_j=\sum_{h\le j}v_h^{\rm bar}$, with $s_0=s_{2r}=0$. Set
\[
 T_{\rm bar}(k)=\sum_{j=1}^{2r}[s_{j-1}|v_j^{\rm bar}]
                           -\sum_{i<r}[u_i|-u_i],
 \qquad Z(k)=(-T_{\rm bar}(k);[u_i]_{i<r}).
\]
The boundary of $T_{\rm bar}(k)$ is $\sum_i([A_iu_i]-[u_i])$, so $Z(k)$
is a cycle of graph coordinate $k$. Its fibre coordinate is
\[
 z(k)=\sum_i f_{A_i}(u_i)-\sum_jp(s_{j-1},v_j^{\rm bar})
                                      +\sum_i p(u_i,-u_i)\quad\text{in }C.
\]
Choose an integer basis $k_j$ of $K$, with its inverse coordinate
matrix. For a bivector $x=\sum x_{\alpha\beta}e_\alpha\wedge
e_\beta$ let
$F(x)=\sum x_{\alpha\beta}([e_\alpha|e_\beta]
-[e_\beta|e_\alpha])$ in the fibre. Define
$s(k_j)=[Z(k_j)]-[F(z(k_j))]$ and extend linearly.
Changing the representative of $z(k_j)$ changes $F$ by the
fibre relation bounding in the corresponding mapping cylinder.
Thus
\[
 (c,k)\longmapsto [F(c)]+s(k)
\]
is well-defined. Its composite with $P_2$ is the identity.
Conversely, subtract $s$ of the graph coordinate from a two-cycle.
The mapping-cone exact sequence puts the remainder in the fibre;
$p_2$ and $F$ are inverse there. This proves the other composite
and the integral isomorphism in \eqref{maps:eq:D.BC11}.

\emph{Peripheral and logarithmic tori.}
At an independent boundary a fixed vector $A_iv=v$ gives the
cycle with $g_i=[v]$ and every other component zero. Its coordinates
are $(f_{A_i}(v),J_iv)$. For $A_i^m\lambda=\lambda$, the torus
in the $m$-fold mapping-torus cover has edge chain
\begin{equation}
                        \sum_{\nu=0}^{m-1}[A_i^\nu\lambda].
                                                        \hypertarget{tag-01XV}{}\label{maps:eq:D.BC13}
\end{equation}
Its endpoints telescope. The graph coordinate is $J_iN_m\lambda$
and its fibre coordinate is the sum in \eqref{maps:eq:D.BC4}.

The last loop is $t_r=t_{r-1}^{-1}\cdots t_0^{-1}$.
The inverse cylinders carry its local one-chain $[u]$ to
\begin{equation}
         -[A_{j+1}\cdots A_{r-1}A_ru],
                            \qquad 0\le j<r .
                                                        \hypertarget{tag-01XW}{}\label{maps:eq:D.BC14}
\end{equation}
Applying $\sum(A_j-I)$ gives $[A_ru]-[u]$ by telescoping.
There is no residual fibre chain. For a fixed vector this is
$J_rv$; applying $P_2$ gives $c_r(v)$. Applying the same map to
each summand of \eqref{maps:eq:D.BC13} gives its full norm sum.
This proves both lines of \eqref{maps:eq:D.BC3} with their
ordered suffixes.

The positive normal meridian of a logarithmic cap is
$\mu=t_i^m v^{-1}$, where $A_i^m=I$ and $A_iv=v$.
The commuting fibre and normal loops give
\begin{equation}
              \mathsf T(\lambda,\mu)=\mathsf T(\lambda,t_i^m)-\mathsf T(\lambda,v).
                                                        \hypertarget{tag-01XX}{}\label{maps:eq:D.BC15}
\end{equation}
The last term is the fibre bivector $\lambda\wedge v$.
The root pair gives $\mathsf T(t_i,\mu)=\mathsf T(v,t_i)$, and hence the fifth
column $(c_i(v),J_iv)$. These signs are fixed by the boundary map
\[
 (y,\tau)\longmapsto
       \bigl(y+(\tau/m)v,\exp(2\pi i\tau/m)\bigr).
\]
Its inverse chooses a root sheet and subtracts $(\tau/m)v$;
changing the sheet is the deck identification. The positive normal
circle is $(y-sv,\tau+ms)$.

The core group in degree one is generated by its four fibre loops
and its root, with relations $(A_i-I)\lambda$ and $m\,\mathrm{root}-v$.
The five displayed columns respect these relations as follows.
Every $\Lambda^2A_i$ acts trivially on $C$, including the final
inverse word. For an independent site, \eqref{maps:eq:D.BC8}
therefore gives
\[
 c_i(u)+c_i(w)-c_i(u+w)=p(A_iu,A_iw)-p(u,w)\quad\text{in }C.
\]
For the last site put $B_j=A_{j+1}\cdots A_{r-1}A_r$.
The polar defect of the second line of \eqref{maps:eq:D.BC3}
is, in the same group,
\[
 -\sum_{j<r}\bigl(p(A_jB_ju,A_jB_jw)-p(B_ju,B_jw)\bigr)
             =p(A_ru,A_rw)-p(u,w),
\]
by telescoping the ordered suffixes. Thus the same polar formula
holds for every site. Put
$s_{i,m}(u)=\sum_{\nu<m}c_i(A_i^\nu u)$ in $C$.
Its polar defect telescopes to
$p(A_i^mu,A_i^mw)-p(u,w)=0$, so $s_{i,m}$ is additive.
Cyclic reindexing gives $s_{i,m}(A_iu)=s_{i,m}(u)$, and hence
$s_{i,m}((A_i-I)\lambda)=0$. The remaining wedge term is a
return relation because $A_iv=v$:
\[
 ((A_i-I)\lambda)\wedge v
              =(\Lambda^2A_i-I)(\lambda\wedge v).
\]
The root relation reads $s_{i,m}(v)=mc_i(v)$.
On $\operatorname{Fix}A_i$ the same polar formula is zero,
so the fixed-normal map is additive there as well.
The normal line pulls back to the trivial normal line on the
finite cover, so its Chern class is torsion. Its integer evaluation
on each integral two-cycle is zero. The Thom sequence therefore
identifies the full boundary kernel with this whole core $H_1$
image. This proves all five columns, including every torsion
relation in the normal-line case, as in
Theorems~\ref{maps:src:D.2.29} and~\ref{maps:src:D.2.30}.
Writing $\nu=qh+j$, $0\le j<h$, proves the stated period
compression of the sum and $N_m=(m/h)N_h$ over the integers.

The affine freeness criterion is \eqref{maps:eq:D.WP5} at the
weighted matrix.
Here \(v=v_i\) is the power at the chosen logarithmic site,
and \(v_{i,\alpha}=e_\alpha^*(v_i)\) are its coordinates in the
ordered marked lattice basis. Let ${F_{\rm fix}}$ be a primitive
fixed basis and $H$ a basis of the image lattice of ${F_{\rm fix}}^{-1}N_h$.
Then $P_{\rm avg}\Lambda={F_{\rm fix}}(H/h)\mathbb Z^2$, where $P_{\rm avg}=N_h/h$; an integer
right inverse of $H^{-1}{F_{\rm fix}}^{-1}N_h$ gives every lattice preimage.
Write $v={F_{\rm fix}}z_{\rm fix}$, $w=hH^{-1}z_{\rm fix}$ and $t=\gcd(m,w_1,w_2)$.
The fixed-plane condition for the $j$th power is $m/t\mid j$.
On a nonidentity moving plane, $(A_i^j-I)^{-1}$ gives the fixed
point once that condition holds. On an identity moving plane,
the condition is $jv/m\in\Lambda$. Consequently freeness is
exactly
\[
 \gcd(m,v_{i,1},v_{i,2},v_{i,3},v_{i,4})=1,\qquad h\mid m/t.
\]
The first failed condition gives the fixed power
$j=m/\gcd(m,v_{i,1},v_{i,2},v_{i,3},v_{i,4})$; the second gives $j=m/t$.
For an affine lift $A_ix+\eta$ with $P_{\rm avg}\eta=v/m$, put
$y=(I-A_i)^{-1}(\eta-v/m)$ on the moving plane.
The coordinate change $x'=x-y$, with inverse $x=x'+y$,
carries it to $A_ix'+v/m$.
On the root family the same change is analytic: for its root
operator $R$ and norm-zero moving section $b_0$, use
$y_0=-m^{-1}\sum_{j=0}^{m-1}jR^jb_0$.
The finite-operator identity in the proof of
Theorem~\ref{maps:src:D.2.29} gives $(I-R)y_0=b_0$.
Its input identification fixes $R,m,b_0$ and the root coordinate.
The coordinate map subtracts
$y_0$ and its inverse adds $y_0$; the finite sum carries the
original translation and line units.

The integral marking is also explicit under a change
$t_i'=(\lambda,t_i)$. On $\Lambda\times C$ use
\[
 (x,z)(y,w)=(x+y,z+w+p(x,y)),\qquad
 (x,z)^{-1}=(-x,-z+p(x,x)).
\]
The lift of $A_i$ is $(x,z)\mapsto(A_ix,z-c_i(x))$;
\eqref{maps:eq:D.BC8} proves multiplicativity in $C$.
Conjugation gives
$c_i'(u)=c_i(u)+(A_iu)\wedge\lambda$ and
$v'=v+N_m\lambda$. Commuting $t_i'$ with its $m$th power gives
\[
 \sum_{\nu<m}c_i(A_i^\nu\lambda)
       =c_i(N_m\lambda)+(N_m\lambda)\wedge\lambda.
\]
Substitution shows that the four norm columns are fixed and the
new root column is the old root plus the old $\lambda$ norm column.
The column change and its inverse are
$\left(\begin{smallmatrix}I_4&\lambda\\0&1\end{smallmatrix}\right)$
and the same matrix with $-\lambda$.
On core groups the map sends the new root to
$\mathrm{root}+\lambda$ and fixes the fibre, with inverse
subtraction. It respects the power relations because
\[
 mI-N_m=(I-A_i)\sum_{j=0}^{m-1}(m-1-j)A_i^j.
\]
In the central lift the inverse generator change retains
$(-\lambda,p(\lambda,\lambda))$. All these maps are integral
also on the even part of $C$.

At a regular site $A_i=I$ and $h=1$; the same five columns
apply. An added identity edge adds a direct $\Lambda$ block to
$K$, with inverse given by reading that block and projecting to
the others. The signed clutch $e$ is an order-one regular site
of power $-e\gamma$. Its coordinate shear and inverse are
$(y,\theta)\mapsto(y+\theta e\gamma,\theta)$ and subtraction.
Deleting its root $\zeta_*=\gamma^{-e}$ gives the asserted product
relation. Replacing a chosen independent base loop by the inverse
of the ordered product gives the change to any last boundary;
the original product relation gives its inverse.

\emph{Unchanged caps.}
Use the geometric local maps in the proof of
Theorem~\ref{maps:src:D.2.40}: the period cover and transfer,
the cusp normalization, the zero-Chern and rank-one component
maps, and the nonzero-Chern line flow and Gysin sections.
Their input identification fixes the natural row, positive $d$,
fan, node units, powered component permutation and line norm.
These displayed local
maps are defined for every positive $d$, and their component
and Gysin bases give free cap $H_2$ at every such weight.
The bounding geometric normal families in that proof are
evaluated here by $P_2$ below. Put $W_d=\Lambda^2S_d$ and let
$f_d:\Lambda^2\Lambda\to H_2(U_d)$ be the fibre map.
The cover and simplex transfer satisfy
\[
 p_*f_d=f_1W_d,\qquad
 \operatorname{tr}_U f_1=f_d\,dW_d^{-1}.
\]
The transfer matrix is integral: it multiplies nongamma wedges
by $d$ and gamma wedges by one. On the fibre image the composite
$\operatorname{tr}_Up_*$ is multiplication by $d$, since a deck
translation of a fibre is homotopic to the identity. Freeness of
$H_2(U_d)$ gives the two containments in
\[
 \ker f_d=(\Lambda^2\Lambda)\cap
          W_d^{-1}(\ker f_1\otimes\mathbb Q).
\]
Choose primitive quotient rows $Q_1$ and a section $\sigma_1$
for the weight-one fibre kernel. A unimodular Smith decomposition
of $Q_1W_d$ gives a kernel basis $K_d$, inverse kernel rows $L_d$,
quotient rows $Q_d$ and section $\sigma_d$, satisfying
\[
 Q_d\sigma_d=I,\quad L_dK_d=I,\quad
 K_dL_d+\sigma_dQ_d=I.
\]
The cover and transfer in these coordinates are
$Q_1W_d\sigma_d$ and $Q_d dW_d^{-1}\sigma_1$.
Both composites are multiplication by $d$. These identities
give every saturated kernel and its inverse coordinates in
\eqref{maps:eq:D.OW4}.

The boundary Wang sequence has fibre term
$\operatorname{coker}(\Lambda^2A_i-I)$ and graph term
$\operatorname{Fix}A_i$. The fixed normal tori constructed in the
local cap maps bound their node thimbles or original cap discs.
Subtracting those literal tori from a boundary class leaves a
fibre class. The latter bounds precisely on the kernel just
computed. Their images under $P_2$ are the full columns
\eqref{maps:eq:D.BC5}. Write \((u_{\rm T},r_{\rm T})\) for the same ordered primitive
Tate frame in the imported multiplicative cap marking. Identity
on this marked frame fixes \(\delta,\gamma\) and carries the original
cap calculation to these coordinates. The possible rim term, with
its original order, is $\delta\wedge r_{\rm T}$. It vanishes by the integral unit elliptic
preimages in \eqref{maps:eq:D.OW7}; those preimages are polynomial
in $d$. Hence every cap kernel, with its torsion, is represented
by the stated fibre and fixed-normal columns.
At a narrow zero-Chern multiplicative site, use the integral
period basis $(l_N,h_N,p_N,w_N)$ in the rank-one part of the proof of
Theorem~\ref{maps:src:D.2.40}, with $r_0=\gcd(d,\nu)$.
Identity on this basis carries the two marked boundary maps
to their common normalization diagram. Both fibre kernels are
$\mathbb Z\langle l_N\wedge p_N,l_N\wedge w_N\rangle$, and all three
fixed normals bound. Write
\[
 \begin{gathered}
 (\mathsf e_1^{\rm N},\mathsf e_2^{\rm N},
  \mathsf e_3^{\rm N},\mathsf e_4^{\rm N})=(l_N,h_N,p_N,w_N),\\
 \mathsf e_{jk}^{\rm N}=\mathsf e_j^{\rm N}\wedge\mathsf e_k^{\rm N}.
 \end{gathered}
\]
The period map \(\Phi_N:\mathbb Z^4\longrightarrow\Lambda\) sends
\((z_1,z_2,z_3,z_4)\) to \(\sum_j z_j\mathsf e_j^{\rm N}\).
Its inverse, in physical coordinates \((x_u,x_r,x_\delta,x_\gamma)\)
relative to \((u_{\rm T},r_{\rm T},\delta,\gamma)\), is
\[
 \begin{aligned}
 z_1&=x_u,&z_3&=x_\delta+(a\kappa/\nu)x_u,\\
 z_2&=(\nu/r_0)x_r+(db\kappa/r_0)x_\gamma,&
 z_4&=-y_Nx_r+x_Nx_\gamma.
 \end{aligned}
\]
Here \(x_N,y_N\) are the same imported integers with
\((\nu/r_0)x_N+(db\kappa/r_0)y_N=1\). The two composites
are identity by that equation and the displayed short shear.
Use \(\Lambda^2\Phi_N\) to reconstruct wedge vectors and
\(\Lambda^2\Phi_N^{-1}\) to extract their coefficients \(y_{jk}\)
in order \(12,13,14,23,24,34\). The four complementary coordinates
and their section, for \(z\in\mathbb Z^4\), are
\[
 \begin{aligned}
 Q(y)&=(y_{12},y_{23},y_{24},y_{34}),\\
 \sigma(z)&=z_1\mathsf e_{12}^{\rm N}+z_2\mathsf e_{23}^{\rm N}
           +z_3\mathsf e_{24}^{\rm N}+z_4\mathsf e_{34}^{\rm N}.
 \end{aligned}
\]
The inverse cap coordinates read the first coefficient of the
component sum and the three graph-circle/torus products in that
diagram. Thus $Q\sigma=I$ gives the same complete marked
boundary kernel for these two choices, with every node and line
unit carried by the stated period maps. The finite choices
at $\mathrm{III}^*$ and $\mathrm{IV}^*$ account for the extra
$36$ patterns. Sitewise inclusion and replacement give the
maps $j,\pi$ stated above and their two displayed composites.

\emph{The complete gluing and homology.}
Let $\mathcal B$ be the full relation matrix with target
$\Lambda^2\Lambda\oplus K$: include all fibre relations,
all primitive unchanged-cap fibre kernels and all the normal
columns just proved. Put $\mathcal C=\operatorname{coker}\mathcal B$.
In a Smith decomposition $U\mathcal B V=\Delta$, the quotient
generator for a diagonal entry $d_j$ is $U^{-1}e_j$, and its
bounding chain is $Ve_j$, since
$\mathcal B Ve_j=d_jU^{-1}e_j$. This supplies the cyclic
coordinates and their inverse representatives.

Fill the finite and regular logarithms first and leave the
selected cusp unfilled; call the result $E$. With
$g=\gcd(a,d)$, the exterior degree-one group is
\[
 \mathbb Z\gamma\oplus(\mathbb Z/g)\delta
                   \oplus\bigoplus_{i\ne *}\mathbb Z x_i.
\]
The logarithmic root discs have boundaries
$m_ix_i-v_{i,\gamma}\gamma-v_{i,\delta}\delta$.
Their distinct nonzero base coefficients prove injectivity
of the disc boundary map, with inverse obtained by reading
and dividing those coefficients by $m_i$.
Use a free lift of this group and let $P$ contain these columns
and $g\delta$. The same argument proves that $P$ is injective.

For each remaining component use its positive normal disc;
its boundary is $\mu_Cx_i+\rho_C\delta$, where $\mu_C$ is its
multiplicity and $\rho_C$ is the integer winding read from its
marked period and line transitions. Exponentiating the period
logarithm recovers the path. Use the selected-cusp matrix
$C_{\rm cusp}$ of the imported cap construction in its same
short-period marking: its coordinate map is multiplication by
$C_{\rm cusp}$ and its inverse is $C_{\rm cusp}^{-1}$.
Thus $(C_{\rm cusp})_{22}=(1-ab\kappa^2)/n_*$, and the $d$
discs have boundaries $x_*+j(C_{\rm cusp})_{22}\delta$, $0\le j<d$.
Let $M$ contain all these boundaries. The kernel $\ker\Gamma$
of their map to $H_1(E)$ is computed by the integer equation $Mz=Pz'$.
Injectivity of $P$ makes $z'$ unique. A Hermite basis of the
projected $z$-lattice and its inverse coordinates give exactly
this kernel, including its index in the component-disc lattice.

The relative sequence \eqref{maps:eq:D.WP11} is therefore
\[
 0\longrightarrow\mathcal C\longrightarrow H_2(X)
        \longrightarrow\ker\Gamma\longrightarrow0.
\]
For a kernel basis vector take its component discs, subtract
the logarithmic power discs given by $z'$, and close the remaining
boundary by the exterior monodromy chains. The coefficient of
$g\delta$ is bounded by the B\'ezout combination of the
$a\delta$ and $d\delta$ monodromy chains. Reading the component
coordinates and subtracting these closed cycles is the inverse
splitting. Thus $H_2(X)=\mathcal C\oplus\ker\Gamma$,
and the free abelian group $\ker\Gamma$ adds no torsion. Its actual
integral construction retains every common even factor in
$T=\operatorname{Tor}\mathcal C$.

The nongamma vanishing words are unchanged by $S_d$.
The determinant-one elliptic pairs kill both elliptic generators.
In the two $\mathrm{OS}46/\mathrm{II}$ cases, use the imported
words $u_{\rm van}=(1,1)$ and $v_{\rm van}=(-4,-1)$ in the
same marked elliptic basis $(e_1,e_2)$. The original $I_3^*$
discs give $2e_1+\delta$ for $a=1$ and $e_1+2\delta$ for $a=4$,
together with $3e_1=-u_{\rm van}-v_{\rm van}$.
The selected cusp kills $\delta$. These relations kill the
elliptic lattice in those cases as well. Reading the remaining
period exponent and representing it by the return path are
inverse maps onto $\mathbb Z\gamma$.
The actual logarithmic root paths give the Seifert presentation
in the statement: their powers are
$\zeta_i^{m_i}=\gamma^{\mu_i}$ and their product is $\gamma^e$.
Conversely the return and root paths satisfy exactly these
vanishing, power and product relations; the maps fix every
generator on composition.

For free ranks, the rational inverse of $S_d$ carries the graph
and fixed-normal diagrams to weight one. Equivalently, let
$\rho=P_2$ on the fibre coordinate and let $\tau$ be evaluation
of the alternating cocycle $x\wedge A_wy$. On the fibre
$\tau=2\rho$, so $2\rho-\tau=\mathfrak h\circ\operatorname{graph}$.
On the basis above,
$\mathfrak h(k_j)=2z(k_j)+\sum_t s_{t-1}\wedge v_t^{\rm bar}$.
The maps $(c,k)\mapsto(c-\mathfrak h(k)/2,k)$ and addition of the same
term are inverse over $\mathbb Q$. They send fixed-normal
coordinates to zero and norm coordinates to
$-\lambda\wedge v$, giving the free-rank calculation of
\eqref{maps:eq:D.WP9}. At odd $g$ these maps are also integral:
the free coordinate of $\tau$ is even under the weight-one
cover, and multiplication by two is invertible on $\mathbb Z/g$.
At every $g$, the finite factors remain those of $\mathcal B$.

In detail, write $r_1=\operatorname{rank}H_1(X)$, let $L$
count the nonzero-Chern originals and let $\beta$ be the sum
of the component contributions in
Definition~\ref{mapsIII:def:all-weight-integral-data}.
The component-disc group has rank $d+(u-1)+\beta+L$ for
$u$ unchanged caps, and $\operatorname{rank}\Gamma=u-r_1$.
The rational cokernel $\mathcal C$ has rank $r_1$ at $L=0$
and zero at $L>0$. Hence
\[
 b_2=\begin{cases}d-1+\beta+2r_1,&L=0,\\
                  d-1+\beta+L+r_1,&L>0.
       \end{cases}
\]
The local component complexes give $\chi(X)=2d$, and duality
gives $b_3=2\beta+2r_1$ at $L=0$ and $2\beta+2L$ at $L>0$.
For each cyclic Smith factor use its two-term resolution
$\mathbb Z\xrightarrow{d_j}\mathbb Z$. The dual unit cochain
represents its Ext generator modulo $d_j$, with inverse the
same residue; cap product with the oriented fundamental class
gives the corresponding $H_3$ and $H_4$ coordinates.
This proves the complete integral table \eqref{maps:eq:D.WP9}.

\end{proof}

\Needspace{5\baselineskip}
\begin{remark}[Integral computation and the odd-factor comparison]
\hypertarget{tag-Q141}{}\label{mapsIII:rem:all-weight-integral-calculation}
{\normalfont\scriptsize tag Q141}\par
Form the full quadratic coefficient vectors of
\eqref{maps:eq:D.BC2}, the ordered suffixes of
\eqref{maps:eq:D.BC3}, every fibre relation and all five normal
columns. Integer polynomial identities, saturated integer
kernels, the period sum of length $h$, and Smith or Hermite
decompositions of the full matrices give the complete Seifert
abelianization and integral homology for every positive $d$,
every permitted affine power, every finite regular-logarithm
list and every integer clutch.

The odd-factor comparison uses the map
$(c,k)\mapsto(c-\mathfrak h(k)/2,k)$ and its addition inverse constructed
in the proof. It identifies the fibre, fixed-normal and full
norm columns with their involution coordinates in
Theorem~\ref{maps:src:D.2.40}. At every weight the integer
matrix of the present theorem retains all finite factors.
\end{remark}

\subsection{Whole-fibre initial pairs}

\begin{definition}[Whole-fibre data]\hypertarget{tag-Q200}{}\label{maps:def:whole-fibre-data}
{\normalfont\scriptsize tag Q200}\par
Choose one of the $21$ marked rational elliptic profiles of
Definition~\ref{mapsIII:def:elliptic-marked-data}. Write
$\langle G,G\rangle=1/d_0$, so $d_0$ is its height denominator $d$. Let
\begin{equation}
 P=aG,\quad Q=bG,\quad a,b>0,
 \qquad \overline P_i=0\ \hbox{or}\ \overline Q_i=0
                         \quad\hbox{at every fibre}.       \hypertarget{tag-01ZC}{}\label{maps:eq:D.GI1}
\end{equation}
Choose any nonnegative integral whole-fibre linearization
weights $w_i$, of total $ab/d_0$, supported on multiplicative
or regular fibres. Distinct regular support points are allowed.
At a positive weight use the marked periodic Mumford filling.
At a weight-zero multiplicative fibre choose either its original
cap or its marked Mumford cap. At a finite good-reduction
additive fibre choose either its original cap or any free affine
logarithmic cap. Use the original cap at $I_n^*$.
Arbitrary regular logarithms at additional weight-zero regular sites
and an arbitrary ordinary clutching integer are allowed. These choices constitute \emph{whole-fibre data}.
\end{definition}

\begin{theorem}[Homology of whole-fibre constructions]\hypertarget{tag-02LN}{}\label{maps:src:D.2.46}
{\normalfont\scriptsize tag 02LN}\par
Every threefold determined by Definition~\ref{maps:def:whole-fibre-data}
has the following finite integral presentation of its fundamental group
and integral homology. The construction includes positive weights at
several fibres and noncyclic toric period lattices. It also includes
marked pairs beyond the scalar iterates of the $172$ natural initial
constructions of \ref{maps:src:D.2.22}.

Use the lattice
$\Lambda=\mathbb Z\langle e_1,e_2,\delta,\gamma\rangle$.
Write $B_i$ for the elliptic matrices $A_i$ of
Definition~\ref{mapsIII:def:elliptic-marked-data}, and retain its
cocycles $g_i$ in the identical ordered basis $(e_1,e_2)$.
Put $p_i=ag_i$, $q_i=bg_i$,
choose any rational $z_i$ with $(B_i-I)z_i=q_i$, and set
\begin{equation}
 r_i=((p_i)_2,-(p_i)_1)B_i,\qquad
 A_i=\begin{pmatrix}
 B_i&0&q_i\\ r_i&1&r_i z_i+w_i\\0&0&1
 \end{pmatrix}.                                           \hypertarget{tag-01ZF}{}\label{maps:eq:D.GI2}
\end{equation}
The block sizes are $2,1,1$. At a regular zero use
$A_i=I+w_i\delta\gamma^*$. At an ordinary or regular
logarithm the boundary matrix is $I$. These are integral
matrices, their ordered product is $I$, and they are the
marked monodromies.

Apply the integral bar comparison \eqref{maps:eq:D.BC1}--\eqref{maps:eq:D.BC14} to these
matrices. Thus the regular-piece degree-two group is
$C\oplus K$, with its full unsaturated fibre presentation
and its transported last-boundary marking. At each finite or
regular logarithm use all five columns \eqref{maps:eq:D.BC4}, with its
order $m$ and fixed power $v$, and the exact freeness
test of \ref{maps:src:D.2.29}--\ref{maps:src:D.2.30}. At every other cap use the full fixed-normal
columns \eqref{maps:eq:D.BC5} and the following fibre kernel $F_i$ in
$\Lambda^2\Lambda$. Here $R_i=\Lambda^2A_i-I$.
For a sublattice $L$ of a lattice $W$, write
$\operatorname{sat}L=(L\otimes\mathbb Q)\cap W$; the ambient lattice
is $\Lambda^2\Lambda$ for $F_i$ and $\Lambda$ for the degree-one kernels.
\begin{equation}
\begin{array}{c|c|c}
\hbox{cap}&F_i& h_i=\operatorname{rank}H^4(U_i)\\ \hline
w_i>0\hbox{ at }I_n, n>0&
 \operatorname{im}R_i+\Lambda^2\operatorname{Fix}A_i&nw_i\\
\hbox{rank-one Mumford}&\operatorname{sat}(\operatorname{im}R_i)&\rho_i\\
\hbox{original, }\overline P_i=0&
 \operatorname{sat}(\operatorname{im}R_i)&\rho_i\\
\hbox{original additive, }\overline P_i\ne0&\ker\omega_i&s_i\\
\hbox{original }I_n,\ \overline P_i\ne0&
 \langle u_1\delta,u_1\ell_i,\delta\ell_i\rangle&n .
\end{array}                                                \hypertarget{tag-01ZG}{}\label{maps:eq:D.GI3}
\end{equation}
Products in the last row are exterior products. Coordinates on
$\Lambda^2\Lambda$ are ordered as $12,13,14,23,24,34$.
For a rank-one Mumford cap, $\rho_i$ is the unique nonzero
Smith factor of $A_i-I$. For a zero-Chern original cap it is
the number of orbits of translation by $Q$ on the Kodaira
components. The marked permutations were given in
\ref{maps:src:D.2.22}. In the last two rows, $Q$ is narrow, so $z_i$ may
and must be integral, and
\begin{equation}
 \ell_i=(-z_i,0,1),\qquad
 \omega_i=(1,0,(z_i)_2,0,-(z_i)_1,0).
                                                               \hypertarget{tag-01ZH}{}\label{maps:eq:D.GI4}
\end{equation}
Here $u_1$ is the primitive elliptic vanishing circle and
$s_i$ is the number of Kodaira components. At a positive
regular zero the second row applies, with $\rho_i=w_i$.

Let $\mathcal B$ be the resulting full integral degree-two
matrix, in $\Lambda^2\Lambda\oplus K$, and write
$\operatorname{coker}\mathcal B=\mathbb Z^g\oplus T$.
Let $E$ be the regular piece with all finite and regular
logarithmic caps already filled. Compute $H_1(E)$ from its
marked graph-bundle presentation and the relations
$\tau_i^{m_i}=v_i$. Compute $H_1(X)=\mathbb Z^r\oplus T_1$
by additionally imposing, at each other cap, $\tau_i=1$
and the following fibre relations; here $T_1=\operatorname{Tor}H_1(X)$:
\begin{equation}
\begin{array}{c|c}
\hbox{cap}&\hbox{killed subgroup of }\Lambda\\\hline
\hbox{positive rank-two Mumford}&\operatorname{Fix}A_i\\
\hbox{rank-one Mumford or zero-Chern original}&
                  \operatorname{sat}(\operatorname{im}(A_i-I))\\
\hbox{nonzero-Chern additive original}&
                  \mathbb Z\langle e_1,e_2,\delta\rangle\\
\hbox{nonzero-Chern multiplicative original}&
                  \mathbb Z\langle u_1,\delta\rangle .
\end{array}                                                \hypertarget{tag-01ZI}{}\label{maps:eq:D.GI5}
\end{equation}
Set
\begin{equation}
 q=\sum_{\text{nonlog caps}}h_i-b_1(E)+r,\qquad
 b_2=g+q,\qquad
 \chi=2\sum_{I_n,\ n>0}nw_i,\qquad
 b_3=2+2b_2-2r-\chi .                                    \hypertarget{tag-01ZJ}{}\label{maps:eq:D.GI6}
\end{equation}
Then $q\ge0$ and
\begin{equation}
\begin{array}{c|c}
j&H_j(X;\mathbb Z)\\\hline
0,6&\mathbb Z\\
1&\mathbb Z^r\oplus T_1\\
2&\mathbb Z^{b_2}\oplus T\\
3&\mathbb Z^{b_3}\oplus T\\
4&\mathbb Z^{b_2}\oplus T_1\\
5&\mathbb Z^r .
\end{array}                                                \hypertarget{tag-01ZK}{}\label{maps:eq:D.GI7}
\end{equation}
For the nonabelian fundamental group, take four commuting
fibre generators, all peripheral generators $\tau_i$, the
relations $\tau_i x\tau_i^{-1}=A_i x$ and
$\prod_i\tau_i=1$, all finite-power relations, and
\eqref{maps:eq:D.GI5}. This is a finite presentation of the fundamental group.

Except in the case described next, the entire vertical
subgroup $\langle e_1,e_2,\delta\rangle$ is killed and this
presentation reduces to the Seifert group of \ref{maps:src:D.2.23}, with
the gamma numerators. The exceptional possibility is
\begin{equation}
\begin{gathered}
 \mathrm{OS49/II}=(II,IV^*,I_2),\quad a\text{ odd},\quad
 b\text{ even},\quad\text{both additive caps logarithmic},\\
 w_{I_2}=0,\quad I_2\text{ filled by Mumford},\quad
 \text{all positive weights at regular fibres}.
\end{gathered}                                             \hypertarget{tag-01ZL}{}\label{maps:eq:D.GI8}
\end{equation}
The compatibility condition \eqref{maps:eq:D.GI1} remains required; it
can impose an additional factor $3$ on $a$ or $b$.
Before the logarithmic power relations the remaining
vertical group in \eqref{maps:eq:D.GI8} is $(\mathbb Z/2)^2$. The two
additive monodromies act by inverse order-three
automorphisms. With no additional regular logarithms,
the group is the split semidirect product of this Klein
four group by the two-fibre Seifert group. With arbitrary
regular logarithms use the full presentation above.
Every group in \eqref{maps:eq:D.GI8} surjects onto $\mathbb Z/3$, so none
is simply connected.

\end{theorem}


\begin{proof}
\smallskip\noindent\emph{The line and its return.}
Write \(\pi:N_{\rm ell}\to\mathbb P^1\) for the chosen elliptic surface and
\(M=\mathcal O(P-O)\). The matrices \(B_i\), cocycle \(g_i\),
and component markings are those of \ref{maps:src:D.2.19}.

Write
\[
 J_{\rm div}=P-O-(P-Q)+(-Q).
\]
The group-law function (line through \(P,-Q\), divided by the vertical
line through \(P-Q\), with tangent and vertical limits) has divisor
\(J_{\rm div}+V_{\rm div}\), where \(V_{\rm div}\) is integral vertical. At a reducible
fibre its component degree vector is
\((I-S_{q_{\rm comp}})(\varepsilon_{p_{\rm comp}}^{\rm cmp}-\varepsilon_0^{\rm cmp})\), where \(p_{\rm comp},q_{\rm comp}\) are the component labels of \(P,Q\), \(S_{q_{\rm comp}}\) is the same component translation, and \(\varepsilon_j^{\rm cmp}\) is the component-degree column with coefficient one at component \(j\) and zero at the others. Coefficient reading and insertion in these columns are inverse. If \(p_{\rm comp}=0\), this is zero. If \(q_{\rm comp}=0\), the
translation is identity on the component graph, so it is again zero.
The intersection kernel is the integral multiplicity vector, with
a component of multiplicity one. Consequently \(V_{\rm div}\) is a sum of
whole fibres, with integer coefficients.

The height formula \cite[Section 11.8, arXiv:0907.0298v3, pp.~50--52]{manSS10} gives
\[
 J_{\rm div}\cdot O=\langle P,Q\rangle+
      \sum_i\operatorname{contr}_i(P,Q)=ab/d_0.
\]
Each local correction is zero because one of its two component
classes is zero. Intersecting the principal divisor with \(O\) gives
\(-\deg V_{\rm div}=ab/d_0\). Thus the divisor \(D_{\rm div}=\sum_iw_i\mathcal F_i\) formed from the assumed integral weights,
where \(\mathcal F_i\) is the whole elliptic fibre at the chosen site, has the same base degree as \(J_{\rm div}\).
The divisor \(V_{\rm div}+D_{\rm div}\) on the base has degree zero. A product of
integer powers of linear factors supplies a rational function \(b_{\rm base}\)
with that divisor. If \(f_{\rm grp}\) is the group-law function, then
\[
 \rho_{\rm line}=f_{\rm grp}/(b_{\rm base}\circ\pi),\qquad \operatorname{div}(\rho_{\rm line})=J_{\rm div}-D_{\rm div}.
\]
Multiplication by \(\rho_{\rm line}\) maps \(\mathcal O(J_{\rm div})\) to
\(\mathcal O(D_{\rm div})\); multiplication by \(\rho_{\rm line}^{-1}\) is its inverse.
The canonical section therefore gives the required linearization,
and a small nonzero scalar gives the contracting arrow. This constructs the line for every stated weight distribution.

Put \(J=\left(\begin{smallmatrix}0&1\\-1&0\end{smallmatrix}\right)\).
The Poincar\'e period map in \ref{maps:src:D.2.22} gives
\(r_i=-p_i^tJB_i\) and the return matrix \eqref{maps:eq:D.GI2}.
The forward arrow covers \(-Q\); its short-period exponential and
inverse logarithm are those of the original line. Its zero winding
is \(+w_i\). A regular zero has \(I+w_iE_{\delta\gamma}\).

Compatibility implies integrality as follows. If \(Q\) is narrow, take integral \(z_i\).
If \(P\) is narrow, write \(p_i=(B_i-I)h_{i,\rm ell}\) with \(h_{i,\rm ell}\) integral.
Symplecticity gives
\[
 r_i=h_{i,\rm ell}^tJ(B_i-I),\qquad r_iz_i=h_{i,\rm ell}^tJq_i\in\mathbb Z.
\]
A second solution differs by a \(B_i\)-fixed vector, which \(r_i\)
annihilates by the same identity over the rationals. This proves
integrality and independence of the chosen rational solution.
Every determinant is one, every \(A_i\) fixes \(\delta\), and every
\(A_i\) preserves \(\gamma^*\).

The section cocycle has zero ordered total. The product of
\eqref{maps:eq:D.GI2} consequently has only a possible \(\delta\gamma\) entry.
Its weight-zero part is bilinear in \(a,b\). The 21
products at formal \(a,b\), computed from the defining \(B_i,g_i\),
are
\[
 \prod_i A_i=I+\left(\sum_iw_i-ab/d_0\right)E_{\delta\gamma}.
\]
It follows geometrically from the base-line map above: its degree
is precisely the remaining line winding. Direct multiplication of
the given matrices gives the same bilinear identity in each of
the 21 profiles. The prescribed total weight makes the ordered product equal to the identity.
An arbitrary ordinary clutch is represented by a separate
order-one identity edge with power vector \(-e\gamma\); it does
not change the product identity above.

\smallskip\noindent\emph{Positive-weight period lattices.}

For each local cap calculation, evaluate the marked data at its site by
\[
 \begin{aligned}
 (A,B)&=(A_i,B_i),\\
 (p_{\rm ell},q_{\rm ell},z_{\rm ell},g_{\rm ell})&=(p_i,q_i,z_i,g_i),\\
 (w,U)&=(w_i,U_i).
 \end{aligned}
\]
Write \(q_{\rm base}\) for the given local base coordinate.
At \(I_\nu\), choose primitive \(u_{\rm per},r_{\rm per}\) with \(\det(u_{\rm per},r_{\rm per})=1\), so
\[
 B=I+\nu u_{\rm per}u_{\rm per}^tJ,\qquad g_{\rm ell}=\kappa u_{\rm per}.
\]
The coordinate inverse is the matrix \((u_{\rm per},r_{\rm per})^{-1}\). Here \(u_{\rm per}\) is the marked primitive elliptic vanishing circle used in the statement. In the
short/long decomposition
\(\mathcal S_{\rm per}=\mathbb Z(u_{\rm per},\delta),\ \mathcal L_{\rm per}=\mathbb Z(r_{\rm per},\gamma)\), the difference \(A-I\) is zero on \(\mathcal S_{\rm per}\), and its long columns are
\[
 \mathsf C=\begin{pmatrix}
       \nu&b\kappa\\
       -a\kappa&w-ab\kappa^2/\nu
    \end{pmatrix},\qquad \det \mathsf C=\nu w.
\]
This determinant is computed on the short lattice.
The primitive elliptic period and the line period form an
invertible short complex period matrix \(\Pi_{\rm sh}(q_{\rm base})\).
Exponentiating \(\Pi_{\rm sh}(q_{\rm base})^{-1}z\) identifies its quotient with
\((\mathbb C^*)^2\); logarithms followed by \(\Pi_{\rm sh}(q_{\rm base})\) give
the inverse modulo \(\mathcal S_{\rm per}\). Let \(\varepsilon_j^{\rm long}\) denote the coordinate columns of the ordered long basis \((r_{\rm per},\gamma)\); insertion in that basis and coefficient reading are inverse. The two long maps are
\[
 X\longmapsto q_{\rm base}^{\mathsf C\varepsilon_j^{\rm long}}u_{j,\rm hol}(q_{\rm base})X,\qquad j=1,2,
\]
with the original nonvanishing vector units. Their inverses divide
by precisely those multipliers. These are the same period maps
used in the construction of \ref{maps:src:D.2.22}. 

For \(w>0\), the stipulated unit-triangle fan has period subgroup
\(L_{\rm per}=\mathsf C\mathbb Z^2\subset \mathcal S_{\rm per}\), index
\({D_{\rm cusp}}=\nu w\). Apply \ref{maps:src:D.2.43} to this same fan, including its integral boundary maps. Put
\[
 s=\gcd(\mathsf C_{11},\mathsf C_{12},\mathsf C_{21},\mathsf C_{22}),\qquad t={D_{\rm cusp}}/s.
\]
Unimodular Smith maps on short and long periods are accompanied by
their inverse matrices. They relabel the fan and its period subgroup. In an interleaved
Smith marking the return is \(I+sE_{12}+tE_{34}\).
Let \(\varepsilon_{ij}^{\rm snf}\) be the wedge columns of that interleaved Smith basis.
The short and long Smith basis maps transport these columns from the original periods;
their inverse matrices, after undoing the interleaving, recover the original period columns.
Set, in wedge order \(12,13,14,23,24,34\),
\[
 \begin{gathered}
 a_0=\varepsilon_{13}^{\rm snf},\\
 b_0=\varepsilon_{14}^{\rm snf}+(t/s)\varepsilon_{23}^{\rm snf}+t\varepsilon_{13}^{\rm snf}.
 \end{gathered}
\]
The exact maps in that marking are
\[
 \operatorname{im}(\Lambda^2A-I)=s\mathbb Za_0+s\mathbb Zb_0,
 \qquad
 \ker\{H_2(T^4)\to H_2(U)\}=\mathbb Za_0+s\mathbb Zb_0.
\]
The second equality adds precisely \(\Lambda^2\operatorname{Fix}A\).
Thus, in the original marking, its kernel is
\[
 F_{\rm ker}=\operatorname{im}(\Lambda^2A-I)+\Lambda^2\operatorname{Fix}A.
\]
Its quotient is \(\mathbb Z^4\oplus\mathbb Z/s\).

The cap groups are
\[
 H_1(U)=\mathbb Z^2,\quad
 H_2(U)=\mathbb Z^{{D_{\rm cusp}}+3}\oplus\mathbb Z/s,\quad
 H_3(U)=\mathbb Z^2,\quad H_4(U)=\mathbb Z^{D_{\rm cusp}}.
\]
The boundary group in degree two is
\(\mathbb Z^6\oplus(\mathbb Z/s)^2\). The two fixed-normal
generators bound; among its fibre torsion coordinates, \(a_0\)
is killed and \(b_0\) survives with order \(s\). These are the
whole boundary kernel maps from the normalization and duality calculation in \ref{maps:src:D.2.43}, including its integral inverse maps.

The degree-one map is
\[
 (x_{\rm short},x_{\rm long})\longmapsto \mathsf Cx_{\rm long}\in L_{\rm per}.
\]
A vector of \(L_{\rm per}\) has a unique integral long coordinate \(\mathsf C^{-1}v\);
the corresponding marked long loop is the inverse section.
Its kernel is the primitive short lattice
\(\operatorname{Fix}A=\mathcal S_{\rm per}\). The ordinary meridian bounds the
normal disc at ray \(0\).

For every
\([v]\in G_{\rm cusp}=\mathcal S_{\rm per}/L_{\rm per}\), the ray \((v,1)\) gives
\[
 \partial D_v=\mu+[v]\in\mathbb Z\mu\oplus G_{\rm cusp}.
\]
Pairing a character with the height-one ray gives its exact positive
winding; the same character values recover the ray and disc.
A nonzero unit extending across the disc has zero winding.
The full map on \({D_{\rm cusp}}\) disc coordinates is
\[
 (x_v)\longmapsto\left(\sum_vx_v,\ \sum_vx_v[v]\right).
\]
It has both cyclic coordinates when \(G_{\rm cusp}\) is noncyclic.
The kernel and its inverse are obtained from the Smith period
basis and the two congruence equations, or equivalently from
the explicit disc basis below. Closing the disc differences uses
the long-period cylinders and commuting torus triangles.

The local comparison with its units of \ref{maps:src:D.2.43} is
\[
 Z_\tau(q_{\rm base};x,y)=
 \exp\!\left(2\pi ix+\log q_{\rm base}\,\mathsf Cy+
       \tau\sum_j y_j\log u_{j,\rm hol}(q_{\rm base})\right).
\]
Its inverse is the group logarithm modulo the same periods.
Changing the branch of \(\log q_{\rm base}\) changes the angular coordinate
by \(\mathsf Cy\) on both sides. The compact collar interpolation and
inverse flow preserve the given period action and all disc
markings. These maps verify the period hypotheses of the cusp calculation.

\smallskip\noindent\emph{Rank-one Mumford caps.}

In the rank-one and zero-Chern original cases, let \(\rho=\rho_i\)
be the cap invariant assigned in the statement. The following
calculations give its stated Smith factor and orbit count.
For a rank-one return write
\[
 A-I=l_{\rm rk}\lambda,\qquad l_{\rm rk}\in\Lambda\text{ primitive},
 \quad\lambda(l_{\rm rk})=0,\quad \rho=\gcd(\lambda).
\]
An integral B\'ezout construction extends \(l_{\rm rk}\) to a basis
\((l_{\rm rk},h_{\rm rk},p_{\rm rk},t_{\rm rk})\), with \(\lambda(h_{\rm rk})=\rho\) and \(p_{\rm rk},t_{\rm rk}\in\ker\lambda\).
Its inverse is its integer adjugate (with the chosen positive
orientation). The return is \(I+\rho E_{12}\).
The one-dimensional fan has \(\rho\) consecutive primitive
edges. The node charts are
\[
 \begin{gathered}
 x_j=z_{\rm Tate}/q_{\rm base}^j,\quad y_j=q_{\rm base}^{j+1}/z_{\rm Tate},\quad x_jy_j=q_{\rm base},\\
 z_{\rm Tate}=q_{\rm base}^jx_j,\quad q_{\rm base}=x_jy_j.
 \end{gathered}
\]
Their line-unit changes multiply \(x_j\) and \(y_j\) by reciprocal
units. Flat residual-torus characters are removed by
\(\exp(-s_1L_1-s_2L_2)\), with inverse the positive exponential.
Successive translations remove all but the final node translation.
Writing that last translation as a real lift \(v_{\rm node}\), the map
\((\zeta,y)\mapsto(\zeta,y-f_{\rm cut}(\zeta)v_{\rm node})\), with
\(f_{\rm cut}(0)=0,f_{\rm cut}(\infty)=1\), and its inverse addition remove it.
These are the ruled-component and node maps. Their diagram
is \(I_\rho\times T^2\).

The fundamental-group map and its inverse quotient section are
\[
 \begin{gathered}
 l_{\rm rk} x_1+h_{\rm rk} x_2+p_{\rm rk} x_3+t_{\rm rk} x_4\longmapsto(x_2,x_3,x_4),\\
 (a,b,c)\longmapsto[ah_{\rm rk}+bp_{\rm rk}+ct_{\rm rk}].
 \end{gathered}
\]
Let \(C_j^{\rm rk}\) be the ruled component indexed by the consecutive fan edge \(j\). The fibre specialization in degree two sends
\[
 \begin{gathered}
 l_{\rm rk}\wedge h_{\rm rk}\mapsto\sum_{j=0}^{\rho-1}[C_j^{\rm rk}],\\
 p_{\rm rk}\wedge t_{\rm rk}\mapsto[T^2],\\
 h_{\rm rk}\wedge p_{\rm rk},\ h_{\rm rk}\wedge t_{\rm rk}\mapsto[\text{graph}]\times p_{\rm rk},\
                              [\text{graph}]\times t_{\rm rk}.
 \end{gathered}
\]
The long path crosses each edge once and hence is the primitive
graph circle. Reading the first component coefficient and the
three product coordinates is a section of these four primitive
classes. The remaining component coordinates are differences.
The kernel is
\[
 \mathbb Z(l_{\rm rk}\wedge p_{\rm rk})+\mathbb Z(l_{\rm rk}\wedge t_{\rm rk})
       =\operatorname{sat}\operatorname{im}(\Lambda^2A-I).
\]
The return image itself is \(\rho\) times this lattice.
All boundary torsion in those two coordinates is therefore
killed exactly by the rank-one fibre kernel. The three fixed normals bound by the
short thimble and the two residual-torus families of origin discs.

The cap groups needed below are
\[
 H_1=\mathbb Z^3,\quad H_2=\mathbb Z^{\rho+3},\quad
 H_3=\mathbb Z^{2\rho+1},\quad H_4=\mathbb Z^\rho.
\]
For a positive regular weight these node charts use the
line coordinate itself; \(\rho=w\) and \(l_{\rm rk}=\delta\).
The normalized components and their seams are products with the
residual elliptic torus, after the displayed inverse
translation gauges.

\smallskip\noindent\emph{Original caps with zero component degree.}

Suppose \(P\) is narrow. The line degree vanishes on every
normalized component. Trivialize the line there and use its
node constants. On a tree the constants are removed successively.
On a graph circle with flat holonomy \(h_{\rm hol}=e^{L_{\rm hol}}\), the coordinate
map is \(w\mapsto e^{-sL_{\rm hol}}w\), with inverse \(e^{sL_{\rm hol}}\).
The same construction on a return circle uses an additional
commuting factor \(e^{-tL_{\rm ret,hol}}\). Ratios of the line maps on
a normalized \(\mathbb P^1\) are constant; the node relations
preserve that statement around a cycle. These gauges use the unit transitions and have zero graph winding.

The contraction roof identifies the angular orbit space with the
mapping torus of the translation. For a roof \(r_{\rm roof}(x)>0\),
use its normalized orbit parameter; the inverse multiplies by
that same roof. An endpoint scalar \(u_{\rm end}(x)\) of zero winding is
removed by \([w,x,s]\mapsto[e^{s\log u_{\rm end}(x)}w,x,s]\);
the inverse has the opposite exponent. The component retract
and its mapping-torus cone have the permutation induced by \(Q\).

For an additive fibre with \(s_i\) components, let the orbits
of this permutation be \(\mathcal O_1,\ldots,\mathcal O_\rho\).
Coinvariants sum coefficients in each orbit; an inverse section
chooses one representative. Invariants repeat a coefficient on
each orbit; their inverse reads a representative. The cap groups
are
\[
 H_1=\mathbb Z^2,\quad H_2=\mathbb Z^{\rho+1},\quad
 H_3=\mathbb Z^{2\rho},\quad H_4=\mathbb Z^\rho.
\]
Here and below all orbit permutations are the powers
corresponding to \(bG\).

We specify the integral fibre map, including a possible
nontrivial component orbit. Choose integral \(h_{\rm ell}\) with
\((B-I)h_{\rm ell}=p_{\rm ell}\), and let \(z_{\rm ell}=(B-I)^{-1}q_{\rm ell}\).
Put \(k=\operatorname{den}(z_{\rm ell})\), the exact component order of \(Q\).
The line-coordinate change is
\[
 \phi_{\rm line}(x)=x_\delta-h_{\rm ell}^tJx_{\rm ell},\qquad\chi_{\rm long}(x)=x_\gamma,
\]
with inverse adding \(h_{\rm ell}^tJx_{\rm ell}\). It conjugates the
return to its elliptic block with translation column \(q_{\rm ell}\),
and identity on the line circle. Hence the cap \(H_1\) map
is \((\phi_{\rm line},\chi_{\rm long})\), with section \((d,g)\mapsto d\delta+g\gamma\).
Its kernel is \(\operatorname{sat}\operatorname{im}(A-I)\).

Let \(\mathcal U_j\) be the class of the chosen representative of \(\mathcal O_j\) under the coinvariant map above, and let \(\mu_C\) be the Kodaira multiplicities. On each orbit
the product \( |\mathcal O_j|\mu_C\) is divisible by \(k\).
There is a simple-component orbit of length \(k\); in the
explicit III, IV, III*, IV*, and odd-star permutations this
follows directly by cycling the simple labels. Put
\[
 W=\sum_j\frac{|\mathcal O_j|\mu_C}{k}\,\mathcal U_j .
\]
This is primitive, and the full elliptic fibre maps to \(kW\).
Define the two integer rows on \(\Lambda^2\Lambda\)
\[
 \begin{gathered}
 L_{\rm fib}(x)=k(x_{12}+z_{\rm ell,2}x_{14}-z_{\rm ell,1}x_{24}),\\
 D_{\rm fib}(x)=x_{34}+h_{\rm ell,2}x_{14}-h_{\rm ell,1}x_{24}.
 \end{gathered}
\]
The map is \(x\mapsto L_{\rm fib}(x)W+D_{\rm fib}(x)(\delta'\times\gamma)\).
The rational invariant rows, the complete fibre image \(kW\),
and the two \(H_1\) coordinates determine this formula, including
the mixed elliptic-return terms.

Since \(k\) is the least common denominator,
\(\gcd(k,kz_{\rm ell,2},-kz_{\rm ell,1})=1\). Choose B\'ezout coefficients giving one,
use their combination of \(e_{12},e_{14},e_{24}\), and subtract
its \(D_{\rm fib}\)-value times \(e_{34}\). This is a preimage of \((1,0)\);
\(e_{34}\) is a preimage of \((0,1)\). Reading \((L_{\rm fib},D_{\rm fib})\)
and these two representatives gives inverse quotient coordinates.
The rational kernel is the return image and the target is free.
Therefore its full integer kernel is
\[
                  \operatorname{sat}\operatorname{im}(\Lambda^2A-I).
\]
This verifies the arbitrary nontransitive additive case directly.

For a multiplicative zero-Chern original, \(\nu\mid a\kappa\).
Set \(u_{\rm per}'=u_{\rm per}-(a\kappa/\nu)\delta\),
\(\rho=\gcd(\nu,b\kappa)\), and choose
\(\alpha_{\rm mult}\nu+\beta_{\rm mult} b\kappa=\rho\). The periods are
\[
 \begin{gathered}
 (u_{\rm per}',h_{\rm mult},\delta,t_{\rm mult}),\\
 h_{\rm mult}=\alpha_{\rm mult} r_{\rm per}+\beta_{\rm mult}\gamma,\\
 t_{\rm mult}=-(b\kappa/\rho)r_{\rm per}+(\nu/\rho)\gamma.
 \end{gathered}
\]
The long inverse is
\[
 \begin{gathered}
 r_{\rm per}=(\nu/\rho)h_{\rm mult}-\beta_{\rm mult} t_{\rm mult},\\
 \gamma=(b\kappa/\rho)h_{\rm mult}+\alpha_{\rm mult} t_{\rm mult},
 \end{gathered}
\]
and the short inverse restores \((a\kappa/\nu)\delta\).
The return is \(I+\rho E_{12}\). The unwrapped Tate chain
is divided by precisely the subgroup
\(\nu\mathbb Z+b\kappa\mathbb Z=\rho\mathbb Z\);
\(h_{\rm mult}\) advances by \(\rho\) edges and \(t_{\rm mult}\) has zero edge shift.
The preceding rank-one formulas are evaluated in the ordered columns
\((l_{\rm rk},h_{\rm rk},p_{\rm rk},t_{\rm rk})=(u_{\rm per}',h_{\rm mult},\delta,t_{\rm mult})\),
with the short and long coordinate inverse displayed above.
The flat gauges above give the \(I_\rho\times T^2\)
diagram and the rank-one quotient, specialization and cap-group maps above. In particular its orbit count
is exactly \(\rho\), and its fibre kernel is the saturated return kernel.

\smallskip\noindent\emph{Original caps with nonzero component degree.}

Compatibility now makes \(Q\) narrow. Choose integral \(z_{\rm ell}\) with
\((B-I)z_{\rm ell}=q_{\rm ell}\), and write
\[
 \ell=(-z_{\rm ell,1},-z_{\rm ell,2},0,1),\qquad
 R=(e_1,e_2,\delta,\ell).
\]
This basis is unimodular. Its inverse restores \(z_{\rm ell}\) to the
elliptic coordinates of a vector with a gamma component.

The identity-component section \(Q\) has a logarithm in the
local elliptic group. At a multiplicative central fibre this
is the chosen logarithm in \(\mathbb C^*\); at an additive
fibre it is the coordinate in \(\mathbb C\). The implicit
function theorem extends it to the small base disc. Its
translation flow extends to the minimal elliptic neighbourhood
and has \(Q\) as endpoint. Lift it to the line using a
connection agreeing with the normalized component frames.
The endpoint ratio to the line arrow is a nonzero
function. On an additive tree it has a logarithm because
\(H_1\) of the tree neighbourhood is zero. On a multiplicative
cycle its component ratios are constant and the node
constants show zero graph winding. Multiply by the reciprocal
exponential along the flow to correct this endpoint exactly.
Write \(E_{\rm circ}(M)\) for the unit-circle bundle of the restricted line in this local original cap.
The roof and inverse roof then give the product
\[
                         E_{\rm circ}(M)\times S^1_\ell .
\]
The inverse restores the flow, the same endpoint units and
the same roof. These mutually inverse maps identify the original germ with the displayed product.

If the component met by \(P\) is \(C_{p_{\rm comp}}\ne C_0\), the circle
bundle has Chern vector \(c=C_{p_{\rm comp}}^*-C_0^*\). Both components
have multiplicity one, even when \(p_{\rm comp}\) is not a generator
of the component group. Thus \(c\) is primitive, with
section \(C_{p_{\rm comp}}\). If \(F_{\rm cyc}=\sum\mu_j C_j\) is the component-coordinate vector of the full fibre \(\mathcal F_i\),
a basis of \(\ker c\) is
\[
 F_{\rm cyc},\quad C_j\ (j\ne0,p_{\rm comp}).
\]
The inverse reads \(x_0\), then \(x_j-\mu_jx_0\).
These are the complete Gysin maps and their inverse coordinates.

For an additive tree, \(E_{\rm circ}(M)\) has
\[
 H_1=0,\quad H_2=\mathbb Z^{s_i-1},\quad H_3=\mathbb Z^{s_i}.
\]
The elliptic area maps to the primitive full-fibre vector \(F_{\rm cyc}\);
both angular-elliptic tori vanish. Taking the product \(E_{\rm circ}(M)\times S^1_\ell\)
kills every wedge with \(\ell\). Therefore the fibre map is
the primitive area row
\[
 \begin{gathered}
 \omega=(1,0,z_{\rm ell,2},0,-z_{\rm ell,1},0),\\
 f_{\rm fib}(x)=\omega(x)F_{\rm cyc},\qquad \ker f_{\rm fib}=\ker\omega .
 \end{gathered}
\]
Its section is \(1\mapsto e_{12}\). The cap has
\[
 H_1=\mathbb Z,\quad H_2=\mathbb Z^{s_i-1},\quad
 H_3=\mathbb Z^{2s_i-1},\quad H_4=\mathbb Z^{s_i}.
\]
Its degree-one map kills precisely \(e_1,e_2,\delta\)
and sends \(\ell\) to its return circle.

For \(I_\nu\), the underlying central cycle additionally has
one graph circle \(r_{\rm gr}\), the image of \(r_{\rm per}\) under elliptic specialization. The map kills \(u_{\rm per}\), sends \(r_{\rm per}\) to \(r_{\rm gr}\), and has quotient inverse \(r_{\rm gr}\mapsto[r_{\rm per}]\). The Gysin maps give
\[
 \begin{gathered}
 H_1(E_{\rm circ}(M))=\mathbb Zr_{\rm gr},\\
 H_2(E_{\rm circ}(M))=\ker c\oplus\mathbb Z(r_{\rm gr}\delta)=\mathbb Z^\nu,\\
 H_3(E_{\rm circ}(M))=\mathbb Z^\nu.
 \end{gathered}
\]
The fibre images are
\[
 \begin{gathered}
 u_{\rm per}\wedge r_{\rm per}\mapsto F_{\rm cyc},\\
 r_{\rm per}\wedge\delta\mapsto r_{\rm gr}\times\delta,\\
 r_{\rm per}\wedge\ell\mapsto r_{\rm gr}\times\ell.
 \end{gathered}
\]
They are primitive independent coordinates, with inverse given
by those three coefficient readings. Exactly
\[
             u_{\rm per}\wedge\delta,\quad u_{\rm per}\wedge\ell,\quad\delta\wedge\ell
\]
are killed. The cap groups are
\[
 H_1=\mathbb Z^2,\quad H_2=\mathbb Z^{\nu+1},\quad
 H_3=\mathbb Z^{2\nu},\quad H_4=\mathbb Z^\nu,
\]
and its \(H_1\) kernel is \(\mathbb Z(u_{\rm per},\delta)\).

Every map in the rank-one and original-cap calculations above preserves the component action, the line
units and the integral long period. The orbit and B\'ezout maps prove these assertions for every compatible
\(a,b\), including nontransitive component permutations.

\smallskip\noindent\emph{The exact global rim relation and the fixed normals.}

In the exterior, the return on \(x\wedge\delta\) is
\((B_i-I)x\wedge\delta\), independently of \(a,b,w_i\).
For every one of the 21 marked profiles,
\[
                         \sum_i\operatorname{im}(B_i-I)=\mathbb Z^2.
\]
An integer inverse for both basis vectors is obtained as follows.
For the monodromy columns \(v^{\rm mon}_j\), choose integers \(c_{jk}\) such that
\(\sum c_{jk}\det(v^{\rm mon}_j,v^{\rm mon}_k)=1\). Then
\[
 \sum c_{jk}(v^{\rm mon}_j,v^{\rm mon}_k)\operatorname{adj}(v^{\rm mon}_j,v^{\rm mon}_k)=I_2.
\]
The gcd of the two-by-two minors of the monodromy column matrix is
one in each of the 21 profiles. The extended Euclidean algorithm
gives the integers \(c_{jk}\); the displayed identity is the inverse
on both basis vectors. All these columns are the monodromy images.
Thus both elliptic-line rims vanish integrally in the
full unsaturated exterior group. 
For the cusp the two fixed-normal tori bound the angular
thimbles of its origin disc. For rank one they bound the
short thimble and the two residual-torus families in the node charts.
For a zero-Chern additive original the fixed lattice is
\[
 \mathbb Z\delta\oplus\mathbb Z(-kz_{\rm ell,1},-kz_{\rm ell,2},0,k).
\]
The line circle extends over the origin disc. The second
circle is the \(k\)-fold return orbit, closed by the
identity-component logarithm of \(kQ\). The flat line gauges
of the zero-Chern calculation remove its endpoint unit; an integral phase
adjustment subtracts its delta component. Its normal tube
therefore bounds with this exact integral marking.
The product map \(E_{\rm circ}(M)\times S^1_\ell\) gives the corresponding tubes when
\(Q\) is narrow.

At a nonzero-Chern multiplicative original the geometric
bounding normal of \(u_{\rm per}\) can differ from the origin-section
Wang normal by an integer elliptic-line rim. Indeed its
projection to the elliptic neighbourhood bounds; it has
zero \(\ell\)-product coordinate; the Gysin description
then places its entire difference in
\(\mathbb Z(r_{\rm gr}\wedge\delta)\). Its coefficient is read by
the line winding along that normal tube, and the
inverse correction adds that rim. The integral rim identity kills
this precise difference in the exterior. The other fixed
normals use the product families. Consequently the global
unchanged-cap columns are exactly the fixed-normal
columns of the integral bar comparison, together with
the fibre kernels of the local cap calculations above.

To see completeness, subtract the fixed-normal coordinate
of any boundary class by these normal maps. The
remaining class is a fibre class by the integral Wang
sequence. At the cusp its kernel is exactly the first kernel in \eqref{maps:eq:D.GI3}, with
its surviving order-\(s\) coordinate. At every other cap
the explicit free target and the maps in the rank-one and original-cap calculations above
give exactly the stated saturated or primitive kernel.
This proves the whole global boundary-kernel image,
including its boundary torsion.

\smallskip\noindent\emph{Bar coordinates and logarithmic caps.}

We apply the integral comparison in the proof of
\ref{maps:src:D.2.41}. Its mapping-cone differential
\eqref{maps:eq:D.BC9}, comparison \eqref{maps:eq:D.BC10},
and ordered bar-polygon inverse use the lattice \(\Lambda\),
integral matrices \(A_i\), and their ordered product relation.
Those are the lattice and matrices constructed above. The period
exponential and its inverse logarithm identify each vertex torus,
each edge cylinder, and its origin section with this graph model.
On the fibre lattice and peripheral paths the identification
is the identity. Thus the comparison carries the
markings of \eqref{maps:eq:D.GI2} to the bar markings.

Put \(B_{\rm fib}=\Lambda^2\Lambda\). For \(N_{\partial}\) boundary sites set
\[
 D_\Gamma=(A_0-I,\ldots,A_{N_{\partial}-2}-I),\qquad K=\ker D_\Gamma,
 \qquad C=B_{\rm fib}/\sum_i\operatorname{im}(\Lambda^2A_i-I).
\]
An integer Smith map \(U_{\rm gr}D_\Gamma V_{\rm gr}=\Delta_{\rm gr}\) gives a primitive
basis \(K_0\) from the zero columns of \(V_{\rm gr}\), with inverse the
corresponding rows \(K_0^\vee\) of \(V_{\rm gr}^{-1}\). Consequently
\(K_0^\vee K_0=I\), and \(K_0K_0^\vee\) is the identity on \(K\).
The comparison and its bar-polygon inverse identify the exterior
degree-two group with \(C\oplus K\).

Use the integer polynomial \(f_A\) of \eqref{maps:eq:D.BC2},
whose coboundary is \eqref{maps:eq:D.BC8}. The peripheral data are
\[
 \begin{aligned}
 \mathfrak c_i(x)&=f_{A_i}(x) &&(i<N_{\partial}-1),\\
 \mathfrak c_{N_{\partial}-1}(x)&=-\sum_{j<N_{\partial}-1}f_{A_j}(A_{j+1}\cdots A_{N_{\partial}-1}x)
                                                   &&\text{in }C,\\
 J_i(x)&=(0,\ldots,x,\ldots,0) &&(i<N_{\partial}-1),\\
 J_{N_{\partial}-1}(x)&=-(A_1\cdots A_{N_{\partial}-1}x,\ldots,A_{N_{\partial}-1}x).
 \end{aligned}
\]
The identity \(D_\Gamma J_i(x)=(A_i-I)x\) follows by telescoping.
These are \eqref{maps:eq:D.BC3} and \eqref{maps:eq:D.BC14}
with their graph-size input \(r\) evaluated at \(N_{\partial}-1\). The indexed peripheral and last-boundary maps are the displayed evaluations of that formula family.
A fixed normal therefore has coordinates
\((\mathfrak c_i(v),J_i(v))\), as in \eqref{maps:eq:D.BC5}.

At a finite-good or regular logarithm, use its root order \(m=m_i>0\)
and fixed power \(v=v_i\in\Lambda\). An affine lift
\(x\mapsto A_ix+\eta\) is carried to \(x\mapsto A_ix+v/m\)
by subtracting \((I-A_i)^{-1}(\eta-v/m)\) on the moving subspace;
addition is the inverse. The averaging map in
\ref{maps:src:D.2.29} gives the same change on the root disc
with its line units, and preserves the marked root and power.

Write \(\operatorname{coord}_{\Lambda}\) for coefficient reading in the stated ordered basis and \(\operatorname{ins}_{\Lambda}\) for its inverse insertion into those same four columns. Thus
\[
 (v_{i,\alpha})_\alpha=\operatorname{coord}_{\Lambda}(v_i),\qquad
 v_i=\operatorname{ins}_{\Lambda}((v_{i,\alpha})_\alpha).
\]
For the arithmetic test let \(h_{\rm ord}\mid m\) be the order of \(A_i\),
let \(\mathsf F\) be a primitive basis of its fixed lattice, with coordinate inverse \(\mathsf F^\vee\) on that lattice, and put
\(\mathsf N_{\rm av}=\sum_{j<h_{\rm ord}}A_i^j\). If \(\mathsf H\) is a basis of the integer
image \(\mathsf F^\vee \mathsf N_{\rm av}\), the projected lattice is
\(\mathsf F(\mathsf H/h_{\rm ord})\mathbb Z^{\operatorname{rank}\mathsf F}\).
Write \(v=\mathsf Fz_{\rm fix}\) and \(w_{\rm fix}=h_{\rm ord}\mathsf H^{-1}z_{\rm fix}\). The freeness criterion is
\[
 \gcd(m,v_{i,1},v_{i,2},v_{i,3},v_{i,4})=1,\qquad h_{\rm ord}\mid m/\gcd(m,w_{\rm fix}).
\]
Indeed the order of the projected translation is
\(m/\gcd(m,w_{\rm fix})\). A nonidentity moving-plane power has a fixed
point precisely when this projected translation vanishes; its
fixed point is obtained with \((A_i^j-I)^{-1}\).
An identity moving-plane power is a pure translation and has a
fixed point precisely when \(jv/m\in\Lambda\).
These are the two tests proved in \ref{maps:src:D.2.41}, applied
to the present conjugacy. For a regular logarithm \(h_{\rm ord}=1\) and
\(\mathsf F=I_4\); all four coordinates enter the gcd.

The degree-three Thom image is given by all five columns
\eqref{maps:eq:D.BC4}. The ordered cochains and maps above
identify its four fibre columns and its root column with the
present boundary. The norm sums have at most \(h_{\rm ord}\) distinct
terms and are multiplied by \(m/h_{\rm ord}\). The positive normal
convention is fixed by
\[
 (y,\tau)\longmapsto
 \bigl(y+(\tau/m)v,\exp(2\pi i\tau/m)\bigr),
 \qquad s\longmapsto(y-sv,\tau_0+ms).
\]
The inverse chooses a root sheet and subtracts \((\tau/m)v\);
changing the sheet is the deck identification. The radial
coordinate is \(q_{\rm base}=\zeta_{\rm root}^m\). Thus the Thom map uses the whole
core \(H_1\), including its torsion, with the same normal sign
as \eqref{maps:eq:D.BC4}.

The core group has the presentation
\[
 \langle\Lambda,\tau\mid[\Lambda,\Lambda]=1,\
       \tau x\tau^{-1}=A_ix,\ \tau^m=v\rangle.
\]
The generators map to the fibre loops and deck path; the inverse
reads the unique sheet index \(0\le j<m\) and the remaining
cover lattice vector. Both composites fix these generators.
An ordinary clutch has \(m=1, A=I, v=-e\gamma\).

\smallskip\noindent\emph{The fundamental group.}

Cut the punctured sphere to its ordered polygon and use the
origin section to mark peripheral lifts. The exterior
has four commuting fibre generators, peripheral conjugations
by \(A_i\), and ordered product one. Each finite cap adds
precisely the finite-core presentation. Each other cap has an abelian fundamental
group with the surjective quotient \eqref{maps:eq:D.GI5}, and its chosen
origin meridian bounds. Van Kampen therefore gives exactly
the presentation in the theorem.

The forward homomorphism sends the generators to the specified
fibre and peripheral paths. For the inverse, use the local
sheet-and-return coordinates of the finite-core presentation, or the inverse quotient
sections of the local cap calculations above, and then the cut-base graph
coordinates. They agree on every boundary generator. Both
composites fix every generator, proving the full nonabelian
presentation.

A positive weight kills \(\delta\): a rank-two cap kills
its short lattice \(\mathbb Z(u_{\rm per},\delta)\), and a positive
regular cap has primitive vanishing circle \(\delta\).
An original additive cap then kills both remaining elliptic
generators. In the zero-Chern case this follows from the map \((\phi_{\rm line},\chi_{\rm long})\),
whose two kernel vectors differ from \(e_1,e_2\) only by
integer multiples of \(\delta\); in the nonzero-Chern case
it is explicit in \eqref{maps:eq:D.GI5}. In particular every profile with
an \(I_n^*\) cap has trivial remaining vertical subgroup.

For the other profiles consider a zero-weight \(I_\nu\).
Its primitive rank-one generator, in the earlier ordered-basis convention, is
\[
 l_{\rm rk}=\frac{\nu}{g_{\rm van}}u_{\rm per}-\frac{a\kappa}{g_{\rm van}}\delta,\qquad
                          g_{\rm van}=\gcd(\nu,a\kappa).
\]
For every multiplicative fibre in the marked profiles,
\(\gcd(\nu,\kappa)=1\). Thus the projected primitive relation is
\begin{equation}
       \frac{\nu}{\gcd(\nu,a)}u_{\rm per}.
       \hypertarget{tag-01ZM}{}\label{maps:eq:D.GI9}
\end{equation}
After \(\delta\) is killed, it contains \(\nu u_{\rm per}\). An original or positive-weight cap kills \(u_{\rm per}\)
itself, so it also contains \(\nu u_{\rm per}\). All elliptic conjugates
under the marked \(B_i^{\pm1}\) of these relations are killed.

For each of the 16 profiles without a required original odd-star cap,
form the two-row matrix whose columns are \(\nu u_{\rm per}\) and their first
\(B_i^{\pm1}\) translates. The gcd of its two-by-two minors is one
in every profile except OS49/II. Hence the same B\'ezout-adjugate formula
above supplies an integer right inverse onto \(\mathbb Z^2\).
Here the profiles are OS43/default, OS43/II, OS45/default, OS45/II,
OS47/default, OS47/II, OS47/III, OS49/default, OS49/II, OS49/III,
OS55/default, OS55/II, and OS56/default, OS56/II, OS56/III, OS56/IV.
The five remaining profiles contain a required original odd-star cap.
The column containment used here is independent of \(a,b\), the weights,
and the other filling choices; it proves the stated unbounded scope.

For OS49/II the matrices are
\[
 B_0=\begin{pmatrix}1&1\\-1&0\end{pmatrix},\quad
 B_1=\begin{pmatrix}0&-1\\1&-1\end{pmatrix},\quad
 B_2=\begin{pmatrix}1&2\\0&1\end{pmatrix},\qquad u_{\rm per}=e_1.
\]
The coarse relation is \(2e_1\), and its \(B_0\)-translate
is \(2(e_1-e_2)\). Their difference gives \(2e_2\);
all further translates are even. The resulting lattice is
exactly \(2\mathbb Z^2\). If \(a\) is even, the primitive vanishing-vector formula kills
\(e_1\); its \(B_0\)-translate also kills \(e_2\).
An original or positive \(I_2\) cap has the same effect.
An original additive cap kills both by the preceding paragraph.
Thus the only surviving possibility is precisely
\[
 a\text{ odd},\quad b\text{ even},\quad
 \text{both additive caps logarithmic},\quad
 w_{I_2}=0,\quad I_2\text{ Mumford},
\]
with every positive weight at regular fibres. Compatibility
at IV* still requires \(3\mid a\) or \(3\mid b\).

Before logarithmic power relations the vertical group is
\(V_{\rm vert}=(\mathbb Z/2)^2\). In the displayed basis the two actions are
\[
 S=\begin{pmatrix}1&1\\1&0\end{pmatrix},\qquad
 S^{-1}=\begin{pmatrix}0&1\\1&1\end{pmatrix}.
\]
They have order three, and both \(S-I,S^{-1}-I\) are
invertible over \(\mathbb F_2\). Since \(b\) is even,
the elliptic part of \(A_i\gamma-\gamma=bg_i\) is zero
in \(V_{\rm vert}\). Thus \(\gamma\) is fixed. A fixed logarithmic
power \(v_i\) satisfies
\((B_i-I)v_{i,\rm ell}+bg_i v_{i,\gamma}=0\);
the displayed Klein-four actions therefore make its \(V_{\rm vert}\)-coordinate zero.

With no additional regular logarithms, let
\[
 \mathcal S=\langle\gamma,c_0,c_1\mid
 [\gamma,c_i]=1,\ c_i^{m_i}=\gamma^{n_i},\
                   c_0c_1=\gamma^e\rangle.
\]
Here \(6\mid m_0\) and \(3\mid m_1\).
The action \(\gamma\mapsto I,c_0\mapsto S,c_1\mapsto S^{-1}\)
respects every relator. The full presentation maps to
\(V_{\rm vert}\rtimes\mathcal S\) by these generators; its inverse maps
the two basis elements of \(V_{\rm vert}\), \(\gamma,c_0,c_1\) back
to their fibre and peripheral symbols. The powers
have zero \(V_{\rm vert}\)-coordinate as just proved, and the origin marking has no additional vertical product term.
Both composites fix all generators, proving the split semidirect
product and the injection of \(V_{\rm vert}\).

Additional regular powers may have nonzero \(V_{\rm vert}\)-coordinates,
so the full presentation applies to them. In every
case the assignment
\[
 c_0\longmapsto1,\quad c_1\longmapsto-1
                  \quad\text{in }\mathbb Z/3,
 \qquad \gamma,V_{\rm vert},\text{every regular peripheral}\longmapsto0
\]
respects conjugation, all powers, the ordinary clutch and
the product relation. It is onto. No exceptional group is
simply connected.

Outside this case, the exact quotient and inverse are
\(\Lambda/\langle e_1,e_2,\delta\rangle\to\mathbb Z\gamma\),
\([x]\mapsto x_\gamma\gamma\), and
\(n\gamma\mapsto[n\gamma]\). Substituting them in the full
presentation gives the Seifert group with the \(n_i=(v_i)_\gamma\) and the clutch. These generator
maps prove the claimed reduction.

\smallskip\noindent\emph{The integral algorithm.}

\begin{enumerate}
\item Form \eqref{maps:eq:D.GI2}, including every regular-weight and identity site.
   Use the ordered product, the last-boundary maps,
   and the integer kernel maps \(K_0,K_0^\vee\).

\item Work in the free target \(B_{\rm fib}\oplus\mathbb Z^{\operatorname{rank}K}\).
   Include every column of \((\Lambda^2A_i-I,0)\).
   Include every nonlogarithmic fibre kernel of the local cap calculations above.
   At a cusp this is the first kernel in \eqref{maps:eq:D.GI3}, using the integer return image.
   Include all columns of the fixed-normal formula, expressed using
   \(K_0^\vee J_i\).

\item At each logarithm, including the ordinary clutch, include
   all five columns of \eqref{maps:eq:D.BC4}. The resulting integer matrix is
   \(\mathcal B\). A Smith decomposition
   \(U_{\rm rel}\mathcal BV_{\rm rel}=\Delta_{\rm rel}\) gives
   \[
                    \operatorname{coker}\mathcal B=\mathbb Z^g\oplus T.
   \]
   Let \(\varepsilon_j^{\rm rel,out}\) and \(\varepsilon_j^{\rm rel,ch}\) be the standard coordinate columns of the free presentation target and the relation-column source, respectively.
   The generator of a nonunit diagonal factor \(d_{j,\rm rel}\) is
   \(U_{\rm rel}^{-1}\varepsilon_j^{\rm rel,out}\), with bounding relation
   \(\mathcal BV_{\rm rel}\varepsilon_j^{\rm rel,ch}=d_{j,\rm rel}U_{\rm rel}^{-1}\varepsilon_j^{\rm rel,out}\).
   The zero diagonal coordinates and the corresponding inverse
   rows give the free quotient and its representatives.

\item Compute \(H_1(E)\) on the free abelian group with four fibre
   coordinates and all peripheral coordinates. Its columns are
   every \((A_i-I)e_\alpha\), the sum of the peripheral generators,
   and every logarithmic \(m_i\tau_i-v_i\).
   For \(H_1(X)\), additionally kill each nonlogarithmic meridian
   and the fibre subgroup in \eqref{maps:eq:D.GI5}.
   The short/long quotient, the rank-one quotient and \((\phi_{\rm line},\chi_{\rm long})\)
   prove the first two rows; the additive and multiplicative Gysin maps
   prove the last two. Their inverse sections were specified above.

\end{enumerate}

All operations are finite for every integer input. Smith and
B\'ezout operations give their inverse maps. The finite-order
sums use at most six powers. A noncyclic cusp uses a finite
quotient \(\mathcal S_{\rm per}/\mathsf C\mathbb Z^2\) of exactly \(\nu w\) elements.

\smallskip\noindent\emph{Relative splitting, including every component boundary.}

At least one weight is positive because \(ab/d_0>0\).
At such a site the prescribed filling is the nonlogarithmic
Mumford cap. The additional regular logarithms are at
zero-weight regular sites. Thus at least one cap remains
outside the first logarithmic filling step.
Leave all nonlogarithmic caps unfilled and first attach every
finite and regular logarithm, including the dummy ordinary
clutch. Call the result \(E\).

The exterior abelianization is its full fibre coinvariant group
plus the free base group with the single sum relation. Each
logarithmic relative \(H_2\) generator is its normal root disc,
with boundary \(m_i\tau_i-v_i\). These boundaries are jointly
injective. Indeed their peripheral coefficient vector is zero
modulo the sum relation only when it is a constant vector; at
a nonlogarithmic site that coefficient is zero, so the constant
is zero and every logarithmic coefficient vanishes. Its inverse
on the image reads the independent peripheral coefficient and
divides by \(m_i\). Therefore \(H_2(E)\) is exactly the exterior
\(H_2\) modulo the full five-column formula.

Each nonlogarithmic cap retracts to the displayed component complex
of real dimension at most four. Thus its fifth cohomology vanishes.
Excision and Poincar\'e--Lefschetz duality give
\begin{equation}
 H_2(X,E)=\bigoplus_iH^4(U_i)=\mathbb Z^{\sum h_i},
 \qquad H_1(X,E)=0.
 \hypertarget{tag-01ZD}{}\label{maps:eq:D.GI10}
\end{equation}
The local cap calculations give the ranks \(h_i\) in
\eqref{maps:eq:D.GI3}.
For the cusp, \(H_2(U_i)\) can contain \(\mathbb Z/s\), but
\(H_3(U_i)=\mathbb Z^2\); thus its \(H^4\), used here, is still
free. The torsion in its \(H^3\) is included in
the displayed Smith boundary maps.

Use positive normal discs dual to the component or component-orbit
basis of \(H_4(U_i)\). At a cusp their boundaries are exactly \(\mu+[v]\), with the
complete disc map displayed above.
In Smith coordinates \(G_{\rm cusp}=\mathbb Z/s\oplus\mathbb Z/t\), let \(\varepsilon_{ij}^{\rm disc}\) denote the normal-disc columns indexed by this finite group, and write
\(w_{ij}=\varepsilon_{ij}^{\rm disc}-\varepsilon_{00}^{\rm disc}\). For \(s,t>1\), an explicit basis
of the local disc kernel is
\[
 s w_{10},\quad t w_{01},\quad
       w_{ij}-i w_{10}-j w_{01}
       \quad((i,j)\ne(0,0),(1,0),(0,1)).
\]
The inverse reads the latter component coordinates and the
two moments divided by \(s,t\). If \(s=1\), use \(t w_{01}\)
and \(w_{0j}-j w_{01}\); if \(s=t=1\), this kernel is zero.
The inverse short Smith map transports these ray
representatives back into the original fan.

At another cap, a generic component point has local equation
\(q_{\rm base}=t_{\rm disc}^{\mu_C}\epsilon(t_{\rm disc})\), with \(\epsilon\ne0\).
Take the normal disc in its extending nonzero line frame.
Continue its boundary through the prescribed stems and the
period logarithms. If its resulting fibre coordinate is
\(\lambda_C\), its full boundary is
\[
                       \mu_C\tau_i+\lambda_C.
\]
The integer transition jumps of the Cartier frames
are included. The inverse is their period exponential followed
by the inverse chart transitions. The line unit on the normal
disc has zero winding; changes of stems or peripheral lifts
are the explicit integral shears of the bar comparison.
These maps are computable from the finite normalized-component
diagram, its monomial Cartier exponents and its unit
transitions. Thus they compute the full component-disc boundary map.

Let \(\Gamma:\bigoplus_i\mathbb Z^{h_i}\to H_1(E)\) be the
matrix of these complete boundaries. The relative exact sequence
and the complete boundary kernels proved above give
\begin{equation}
 0\longrightarrow\operatorname{coker}\mathcal B
   \longrightarrow H_2(X)
   \longrightarrow\ker\Gamma\longrightarrow0.\hypertarget{tag-01ZE}{}\label{maps:eq:D.GI11}
\end{equation}
This includes the cusp's surviving torsion coordinate.

Here is an inverse splitting over the integers. Replace a free
presentation of \(H_1(E)\) by an injective relation matrix \(\mathsf P\):
take an integer basis of its relation image, together with a preimage
in the original relation columns. Let \(\mathsf M\) be the lift of
the full component boundary formula to that free presentation. Solve
\[
                              \mathsf Mz=\mathsf Pz'
\]
by reversible Smith operations. Since \(\mathsf P\) is injective,
projection onto \(z\) is one-to-one on the kernel of
\([\mathsf M,-\mathsf P]\); its inverse recovers the unique \(z'\).
An integral basis of this projected lattice is therefore a
basis of \(\ker\Gamma\), even when this kernel has a
nontrivial index in its rational saturation.

For a basis vector take its component normal discs, subtract
the logarithmic power discs prescribed by \(z'\), and close
the remaining boundary by the exterior graph cylinders and
the relation preimages used to construct \(\mathsf P\). Abelian
commutator corrections are bounded by the corresponding
torus polygons. This is a closed integral two-cycle.
Reading relative disc coordinates and subtracting these cycles
is the inverse splitting of \eqref{maps:eq:D.GI11}. In particular
\[
 H_2(X)\cong\operatorname{coker}\mathcal B\oplus\ker\Gamma.
\]
The second summand is free. 
Each boundary fibre and meridian surjects onto the local cap
group by the explicit maps above. Thus
\(\operatorname{coker}\Gamma=H_1(X)\). If
\(H_1(X)=\mathbb Z^r\oplus T_1\), rank in \(\mathsf Mz=\mathsf Pz'\) gives
\[
 \operatorname{rank}\ker\Gamma
           =\sum_i h_i-b_1(E)+r=:q\ge0.
\]
The Smith quotient \(\mathbb Z^g\oplus T\), together with these
splitting and rank identities, proves
\(H_2(X)=\mathbb Z^{g+q}\oplus T\).

\smallskip\noindent\emph{Euler number, ranks and every integral homology group.}

All constructions carry the complex orientation, and the
ordinary collar shear has determinant one. The result is a
connected closed oriented six-manifold.

The positive rank-two cap has exactly \({D_{\rm cusp}}=\nu w\)
component rays and \(2{D_{\rm cusp}}\) maximal triangles. Its
zero-dimensional toric strata are the latter; every other
toric stratum has a positive-dimensional torus factor and
Euler number zero. Thus its local Euler number is \(2\nu w\).
This counts the same fan used in the positive-cusp calculation. A positive
regular rank-one cap, a rank-one ruled wheel, an original
zero-free quotient and a free finite torus quotient all
have Euler number zero by their displayed torus or circle
strata. The exterior and all boundary mapping tori have
Euler number zero as well. Additivity gives
\[
                         \chi(X)=2\sum_{I_\nu,\ \nu>0}\nu w_i.
\]

Put \(b_2=g+q\). Poincar\'e duality and the integral universal
coefficient sequence identify \(\operatorname{Tor}H_3(X)\) with
\(\operatorname{Tor}H_2(X)=T\), and \(\operatorname{Tor}H_4(X)\)
with \(\operatorname{Tor}H_1(X)=T_1\). Their free ranks are paired
in complementary dimensions; \(H_5(X)\) is free of rank \(r\).
The Euler identity gives \(b_3=2+2b_2-2r-\chi(X)\).
For each Smith factor \(\mathbb Z/d\), the resolution
\(\mathbb Z\xrightarrow{d}\mathbb Z\) identifies its Ext
generator with \(1\bmod d\); evaluation on the oriented
dual chain gives the Poincar\'e map and its inverse.
These maps give the torsion identifications in \eqref{maps:eq:D.GI7}.


\smallskip\noindent\emph{Positive support at the \(I_8\) fibre.}
In OS45/default take \(P=2G,Q=8G\), weight two at the \(I_8\)
fibre, all other weights zero, Mumford fillings throughout, and \(e=1\).
The positive period matrix and its Smith factors are
\[
 \mathsf C=\begin{pmatrix}8&24\\-6&-16\end{pmatrix},\qquad
 \det \mathsf C=16,\qquad(s,t)=(2,8).
\]
The preceding integral algorithm gives
\(\operatorname{coker}\mathcal B=\mathbb Z/8\),
\(\sum h_i=20\), \(b_1(E)=5\), and \(H_1(X)=0\).
The vertical quotient is trivial by the OS45 vanishing relations;
the ordinary power \(\tau=-\gamma\), together with the product relation,
then kills \(\gamma\). Thus \(\pi_1(X)=1\), and
\[
 H_2(X)=\mathbb Z^{15}\oplus\mathbb Z/8,\quad
 H_3(X)=\mathbb Z/8,\quad H_4(X)=\mathbb Z^{15},\quad\chi(X)=32.
\]
The scalar iteration map on the original construction is the period
subgroup inclusion \(\gamma_{D_{\rm it}}\mapsto {D_{\rm it}}\gamma\), covering the identity
of the elliptic base. It changes the positive weight to \({D_{\rm it}}\) and
preserves its support. Every natural OS45 positive fibre is \(I_1\),
by \ref{maps:src:D.2.22}; hence none of these scalar iterates has this
positive \(I_8\) support. Thus this pair extends the
marked construction family. \end{proof}

\subsection{Canonical star logarithms}

\begin{definition}[Canonical star data]\hypertarget{tag-Q201}{}\label{maps:def:canonical-star-data}
{\normalfont\scriptsize tag Q201}\par
Let a marked star boundary with nonvanishing line linearization have
the integral basis
$(\alpha,\beta,\lambda,\ell)$ and monodromy
\begin{equation}\hypertarget{tag-0247}{}\label{maps:eq:D.SL1}
\begin{gathered}
 A\alpha=-\alpha,\quad A\beta=-\beta-n\lambda,\quad
 A\lambda=\lambda-\alpha,\quad A\ell=\ell,\\
 \delta=-\alpha+2\lambda,\qquad\gamma^*(\ell)=1,
 \qquad n\in\{1,3\}.
\end{gathered}
\end{equation}
Choose integers $k>0$, $j,r,N$, set $m=2k$ and $q=nk$, and
prescribe the component action $i\mapsto j-i$ on $\mathbb Z/q$.
The integer lift $j$ is part of the peripheral marking.
The prescribed collapsed power is
\begin{equation}\hypertarget{tag-024I}{}\label{maps:eq:D.SL2}
v_{\rm col}=R\lambda+N\ell,\qquad R=2r-j,
                \qquad \gcd(N,2k)=1 .
\end{equation}
These choices constitute \emph{canonical star data}.
\end{definition}

\begin{theorem}[Canonical star logarithms]\hypertarget{tag-02LQ}{}\label{maps:src:D.2.49}
{\normalfont\scriptsize tag 02LQ}\par
Every datum of Definition~\ref{maps:def:canonical-star-data} is realized
by a free affine root action on a rank-one semistable Mumford cap with
$q=nk$ components. The action has the prescribed component reflection
and collapsed power, with its holomorphic unit factors and integer twists.

There are twelve such marked star positions in the natural catalogue:
the eight primitive-gamma $I_1^*$ positions of
\ref{maps:src:D.2.27} and the four primitive-gamma $I_3^*$ positions.
Their unimodular bases are given in the proof.
Put
\[
 f=\#\{i:2i=j\pmod q\},\qquad \rho=(q+f)/2.
\]
Thus $f=1$ for odd $k$, $f=2$ for even $k$ and even $j$,
and $f=0$ for even $k$ and odd $j$.
For the quotient cap $U$, the complete integral groups are
\begin{equation}\hypertarget{tag-024O}{}\label{maps:eq:D.SL3}
\begin{array}{c|c}
a&H_a(U;\mathbb Z)\\\hline
0&\mathbb Z\\
1&\mathbb Z^2\\
2&\mathbb Z^{\rho+1}\oplus
                    (\mathbb Z/2\text{ if }f=0)\\
3&\mathbb Z^{2\rho}\\
4&\mathbb Z^\rho,
\end{array}
\end{equation}
and higher groups vanish.

The boundary has $H_2(\partial U)=\mathbb Z^4$. Its
degree-two fibre coinvariant has the primitive basis
\begin{equation}\hypertarget{tag-024P}{}\label{maps:eq:D.SL4}
x=[\beta\lambda],\qquad
 w=[\beta\ell]+\frac{n-1}{2}[\lambda\ell],\qquad
 [\alpha\beta]=-2x,\quad [\beta\ell]=nw,\quad
 [\lambda\ell]=-2w.
\end{equation}
Let $D,L$ be the fixed normal tori defined by the marked origin
section in the $\delta,\ell$ directions. In coordinates $(x,w,D,L)$,
the boundary kernel is
\begin{equation}\hypertarget{tag-024Q}{}\label{maps:eq:D.SL5}
\begin{array}{c|c}
f&\ker[H_2(\partial U)\longrightarrow H_2(U)]\\\hline
1&\{(0,b,c,d):kb=2Nc-Rd\}\\
2&\{(0,b,c,d):(k/2)b=Nc-(R/2)d\}\\
0&\{(0,b,c,d):kb=2Nc-Rd,\ b\equiv0\pmod2\}.
\end{array}
\end{equation}
These equations specify integral sublattices with projected indices
$k,k/2,2k$ in the full fixed lattice. The boundary cokernel is free
of rank $\rho-1$. The last row includes the cap summand
$\mathbb Z/2$ and the nonzero torsion coordinates of both fibre images.

These columns enter the integral
bar presentation of \ref{maps:src:D.2.41} or \ref{maps:src:D.2.46}. A vector $(0,b,c,d)$
contributes the regular-fibre representative of
$bw$, together with $c$ times the marked $\delta$ sweep
and $d$ times the marked $\ell$ sweep. The cap contributes
$\rho$ to the free relative rank in \eqref{maps:eq:D.GI10} and zero to
the Euler number. The same presentation includes the other cap
choices and their full finite Thom columns.

At a natural position, take this as the sole logarithmic cap
and use the minimal natural marking of \ref{maps:src:D.2.26} at
every other site, with no additional regular logarithms.
For the $\mathrm{OS50/IV}$ position with $(a,b_0)=(1,12)$, the
minimal $IV$ filling is a second finite logarithm; the general
bar insertion above computes that construction. The six default
examples below use minimal multiplicative fillings at their other sites.

For the sole-logarithm construction write $b_0=4B$, with
$B\in\mathbb Z_{>0}$, and let $e$
be the ordinary clutching integer. Put
\begin{equation}\hypertarget{tag-024R}{}\label{maps:eq:D.SL6}
s_{\rm ord}=N-2ke .
\end{equation}
It is an odd nonzero integer, prime to $2k$. Then
\begin{equation}\hypertarget{tag-024S}{}\label{maps:eq:D.SL7}
\begin{gathered}
 \pi_1(X)=\mathbb Z/|s_{\rm ord}|,\\
 T=\operatorname{Tor}H_2(X)=
 \begin{cases}
 \mathbb Z/(2B|s_{\rm ord}|),&f=1,\\
 \mathbb Z/(B|s_{\rm ord}|),&f=2,\\
 \mathbb Z/2\oplus\mathbb Z/(2B|s_{\rm ord}|),&f=0.
 \end{cases}
\end{gathered}
\end{equation}
The other groups are
\begin{equation}\hypertarget{tag-024T}{}\label{maps:eq:D.SL8}
\begin{gathered}
 H_2=\mathbb Z^{\rho-1}\oplus T,\qquad
 H_3=\mathbb Z^{2\rho-2}\oplus T,\\
 H_1=\mathbb Z/|s_{\rm ord}|,\qquad
 H_4=\mathbb Z^{\rho-1}\oplus\mathbb Z/|s_{\rm ord}|,\qquad H_5=0,\\
 H_0=H_6=\mathbb Z .
\end{gathered}
\end{equation}
Cyclic summands of order one are omitted. For example,
the six default natural positions with $b_0=4$ admit the
simply connected choice $s_{\rm ord}=\pm1$, $k=2h$ with $h\in\mathbb Z_{>0}$, and $j$ even.
It has
\begin{equation}\hypertarget{tag-024U}{}\label{maps:eq:D.SL9}
H_2=H_4=\mathbb Z^{nh},\qquad
                         H_3=\mathbb Z^{2nh}.
\end{equation}
Here $n=1$ on OS50/default and $n=3$ on OS46/default.


\end{theorem}


\begin{proof}
\emph{The twelve marked boundaries.}
Use the physical period lattice
$\Lambda=\mathbb Z\langle e_1,e_2,\delta,\gamma\rangle$ and
$\Omega=a e_1^*\wedge e_2^*-b_0\delta^*\wedge\gamma^*$.
The natural component equations of \ref{maps:src:D.2.26} require
$ab_0=d/\nu$ at the selected $I_\nu$, with $a,b_0$ coprime,
and require one of $aG,b_0G$ to be narrow at every other reducible
fibre. Testing the component residues gives
\[
\begin{array}{c|c|c|c}
\text{profile}&\text{selected position}&(a,b_0)&n\\\hline
\mathrm{OS46}&2,3,4&(1,4)&3\\
\mathrm{OS46/II}&3&(1,4)&3\\
\mathrm{OS50}&2&(1,4)&1\\
\mathrm{OS50}&3,4&(1,12),(3,4)&1\\
\mathrm{OS50/II}&2&(1,4)&1\\
\mathrm{OS50/IV}&3&(1,12),(3,4)&1 .
\end{array}
\]
The selected fibre labels are their positions, starting at one, in the ordered table of Definition~\ref{mapsIII:def:elliptic-marked-data}.
These are four $I_3^*$ and eight $I_1^*$ marked positions.
For $n=1$ put
\[
\alpha=2e_1+a\delta,\quad
\lambda=e_1+\frac{a+1}{2}\delta,\quad
\beta=e_2+\frac{3a-1}{2}\lambda,\quad
\ell=\gamma+\frac{b_0}{2}e_2-\frac{3b_0}{4}\lambda.
\]
For $n=3$ put
$\alpha=2e_1+\delta$, $\lambda=e_1+\delta$,
$\beta=e_2-\lambda$, $\ell=\gamma+2e_2-\lambda$.
At $\mathrm{OS50/II}$ apply
$\operatorname{diag}(C,1,1)$, where
$C=\left(\begin{smallmatrix}1&1\\-1&0\end{smallmatrix}\right)$
and $C^{-1}=\left(\begin{smallmatrix}0&-1\\1&1\end{smallmatrix}\right)$.
These column matrices $H$ have determinant one, so their adjugates
are integral inverses. Multiplication by the marked monodromy
gives \eqref{maps:eq:D.SL1}, and multiplication by $\bigwedge^2H$
gives $\Omega(x)=-a$, $\Omega(w)=b_0/4$.
The real period maps $[v]\mapsto[Hv]$ and $[u]\mapsto[H^{-1}u]$
also transport their character lattices, toric rays, units and collars.
Thus the displayed bases identify marked models.

\emph{Flat sections and the holomorphic root.}
The matrix in \eqref{maps:eq:D.SL1} satisfies
$A^2=I+nE_{12}$ and $(I+A)E_{12}=0$.
The relevant flat translation is
\begin{equation}\hypertarget{tag-0248}{}\label{maps:eq:D.SL10}
t=\frac{j}{q}\beta+\frac r k\lambda+\frac{N}{2k}\ell,
 \qquad q=nk .
\end{equation}
It is single-valued because $(A^{2k}-I)t=j\alpha$ is integral,
and its exact forward norm is
\begin{equation}\hypertarget{tag-0249}{}\label{maps:eq:D.SL11}
\sum_{a=0}^{2k-1}A^at=k(I+A)t
                  =-r\alpha+(2r-j)\lambda+N\ell.
\end{equation}
The inverse-power norm $k(I+A^{-1})t$ differs by $j\alpha$.
Conversely, single-valuedness and integrality of the norm give
$qt_\beta\in\mathbb Z$, $kt_\lambda\in\mathbb Z$ and
$2kt_\ell\in\mathbb Z$. An additional angular translation
$c\alpha$ is removed by conjugating with $(c/2)\alpha$;
the inverse conjugacy subtracts that vector.

Here is a period construction with the analytic units.
Set $m=2k$, $\zeta=e^{2\pi i/m}$ and
$l(s)=\log(s)/(2\pi i)$ on the universal cover of the punctured
root disc. In $\mathbb C^2$ take
\[
\begin{split}
 \alpha&=(1,0),\qquad \lambda=(0,1),\\
 \ell&=(-\tau/2+\eta(s),\tau),\\
 \beta&=(-ql(s)-\epsilon(s)/2+B_{\rm per}(s^m),-n/2+\epsilon(s)),
\end{split}
\]
where $\Im\tau>0$, $\tau=\tau(s^m)$, and
$\eta(\zeta s)=-\eta(s)$, $\epsilon(\zeta s)=-\epsilon(s)$.
These functions are obtained from the periods:
$\epsilon=\beta_w+n/2$, $\eta=\ell_z+\tau/2$, and
$B_{\rm per}=\beta_z+ql+\epsilon/2$. Substitution is the inverse map.
Indeed, in the frame given by $\alpha,\lambda$ the finite linear
part is $\mathsf L(z,w)=(-z-w,w)$. The period equations give
$\beta_w(\zeta s)+\beta_w(s)=-n$,
$\beta_z(\zeta s)-\beta_z(s)=\beta_w(s)$,
and $\ell_z(\zeta s)+\ell_z(s)=-\tau$.
An equivariant frame can be obtained by averaging a central
intertwining frame against this finite linear action; its value
at the centre is unchanged and its inverse matrix exists on a
smaller disc. All nonzero line units have holomorphic logarithms
there, with the integer changes of logarithm retained as periods.

For $R=2r-j$ set $t_0=(R+N\tau)/m$. The map
\[
\begin{gathered}
g(s,z,w)=\\
\left(\zeta s,-z-w-jl(s)+F_{\rm hol}(s^m)-\frac{N\eta(s)}m,\ w+t_0\right)
\end{gathered}
\]
conjugates the four periods exactly by $A$.
Its inverse first recovers $s=\zeta^{-1}s'$, then $w=w'-t_0$,
and finally
$z=-z'-w-jl(s)+F_{\rm hol}(s^m)-N\eta(s)/m$.
Two applications translate $z$ by
$-t_0-j/m+2N\eta(s)/m$ and $w$ by $2t_0$.
Thus $g^m$ is exactly the period
$-r\alpha+R\lambda+N\ell$, proving the required collapsed power.

With $Z=e^{2\pi iz}$ this is
$Z'=U_{\rm node}(s,w)s^{-j}/Z$, where
$U_{\rm node}=e^{2\pi i(F_{\rm hol}-N\eta/m-w)}$ is nonzero.
On the edge joining the components $i,i+1$ use
$X=Zs^{i+1}$, $Y=s^{-i}/Z$, so $XY=s$.
On the target edge put $h_{\rm edge}=j-i$. The exact chart maps are
\[
 \begin{gathered}
 X'=\zeta^{h_{\rm edge}}U_{\rm node}Y,\\
 Y'=\zeta^{1-h_{\rm edge}}U_{\rm node}^{-1}X,\\
 w'=w+t_0 .
 \end{gathered}
\]
Their product is $\zeta s$, and the displayed inverse recovers
both branches. Thus the action extends across every node and
pole, sends the component label to $j-i$, and includes
every unit and the integer lift $j$.

The original translation \eqref{maps:eq:D.SL10} and this formula
are explicitly conjugate. Composing the zero-translation deck map
with that flat section at its endpoint changes the displayed
translation by
\[
\left(-\frac jm+\frac{j\epsilon}{2q}
       +\frac{jB_{\rm per}}{q}-\frac{N\tau}{2m}-F_{\rm hol},\
                         -\frac{j\epsilon}{q}\right).
\]
Its norm is zero: its invariant first-coordinate part is killed
by $I+\mathsf L$, and its anti-invariant part lies in the positive
eigenspace of $\mathsf L$. In general, for such a difference $b(s)$,
put $\mathcal T_{\rm av}h(s)=\mathsf L^{-1}h(\zeta s)$ and $c_{\rm av}=\mathsf L^{-1}b$. Then
$h_0=m^{-1}\sum_{a=0}^{m-1}a\mathcal T_{\rm av}^{a}c_{\rm av}$ satisfies
$h_0(\zeta s)-\mathsf Lh_0(s)=b(s)$.
Translation by $h_0$ is the conjugacy and subtraction its inverse.
The integer winding is encoded in $j,r,N$ before applying this
holomorphic gauge.

On the central infinite chain let its graph coordinate
$\theta_{\rm gr}$ have vertices $i/q$. The long period changes
$(\theta_{\rm gr},w)$ by $(1,-n/2)$.
The map $u=w+(n/2)\theta_{\rm gr}$, with inverse
$w=u-(n/2)\theta_{\rm gr}$, gives central coordinates
$u=x+\tau y$ in which
\[
g(\theta_{\rm gr},x,y)=
(-\theta_{\rm gr}+j/q,\ x-n\theta_{\rm gr}+r/k,\ y+N/m).
\]
Every proper power shifts $y$ by a nonintegral multiple of $1/m$,
because $\gcd(N,m)=1$. This proves freeness on all strata.

The period units give the same marked smooth diagrams by explicit
maps. Interpolate their logarithms while fixing integer winding,
and interpolate $\tau$ in the upper half-plane, keeping the
anti-invariance of $\eta,\epsilon$. In a node equation $XY=sV$,
divide $Y$ by $V$; multiplication is the inverse. In line frames
use the exponential of the chosen logarithm and its reciprocal.
At a pole translation subtract its real lift multiplied
by a collar function equal to zero and one at the two ends;
addition is the inverse. The derivatives of these chart maps
give lifts of the interval direction tangent to all central strata.
An invariant partition of unity and finite averaging give one
such lift on the compact proper tube. Its flow and reverse flow
are inverse maps, preserving the prescribed outer marking.

The cap retracts onto this central gluing. For an invariant
Hermitian metric put $f_0=|s|^2$. On $XY=s$ one has
$\|\nabla f_0\|\ge c_{\rm grad} f_0^{3/4}$, with one positive $c_{\rm grad}$
after metric comparison on the finite atlas. The normalized
negative-gradient flow has tail length at most
$4c_{\rm grad}^{-1}f_0^{1/4}$, so its limiting collapse is continuous and
fixes the central fibre. A compact proper tube prevents escape
through a side boundary. Reverse flow gives the inverse on
positive-radius annular collars. The invariant construction
descends under the free action.

\emph{Normalization, conductor and integral groups.}
The quotient of the $q$-cycle by reflection is an interval.
Its $\rho$ vertex orbits and $\epsilon_e$ edge orbits satisfy
$\rho+\epsilon_e=q+1$. The number of fixed edges is $2-f$.
There are exactly $q$ half-edge orbits, since reflection
interchanges the signs of the $2q$ half-edges.
This distinguishes the branches also at $q=1,2$.

At a paired vertex the residual torus lattice is
$\mathcal L_{0,\rm pair}=\mathbb Z^2+\mathbb Z(R/k,N/k)$.
Choose $u_{\rm pair}N+v_{\rm pair}k=1$. A basis is
$\lambda=(1,0)$, $\eta_0^{\rm lat}=(u_{\rm pair}R/k,1/k)$, with inverse
$\ell=k\eta_0^{\rm lat}-u_{\rm pair}R\lambda$ and
$(R/k,N/k)=N\eta_0^{\rm lat}+v_{\rm pair}R\lambda$.
Its ruled normalization is the trivial oriented sphere bundle:
the open ruling period group is abelian, its short circle is
primitive by $\gcd(N,k)=1$, and the logarithms of the remaining
flat multipliers give the line trivialization and its inverse.
Its two pole sections have common class $\sigma_v$.

At a fixed vertex $i$ the lattice is
$\mathcal L_{i,\rm fix}=\mathbb Z^2+\mathbb Z((2r-2i)/m,N/m)$ and $\mathcal L_{0,\rm pair}$
has index two in it. The pole-interchanging clutch has angle
$-\phi-\theta$ after its periodic logarithmic unit is removed.
A unitary lift is
\[
\begin{pmatrix}0&-e^{-\pi i\theta}\\e^{\pi i\theta}&0\end{pmatrix},
\]
whose two endpoint values differ by $-I$. Thus the oriented
sphere bundle has odd section square. Choose a section
$\sigma_v^2=1$: a section extends over the two base cells,
and changing it over a disc by degree $d$ adds $dF_v$ and
changes its square by $2d$, with inverse degree $-d$.
Here $F_v$ is the complex-oriented ruling sphere.
The connected conductor bisection has flat normal line on its
cover, so it has square zero and class $2\sigma_v-F_v$.

At a fixed edge its conductor-to-double-curve map has degree two;
at a paired edge it has degree one. All degree-one maps on
torus periods are the literal lattice inclusions
$K_D^{-1}K_C$, with inverse on the image $K_C^{-1}K_D$.
The fixed-edge stabilizers at the centres $j/2$ and $(j+q)/2$
are $g$ and $\beta g$ respectively. Their lattices are
$\mathcal L_{\rm lat}(m,R,N)$ and $\mathcal L_{\rm lat}(m,R-q,N)$, where
$\mathcal L_{\rm lat}(d,A,N)=\mathbb Z^2+\mathbb Z(A/d,N/d)$:
$(\beta g)^2=\lambda^{-n}g^2$.
The same B\'ezout basis and its inverse give both inclusions.

Use the following integral change from the complex-oriented classes
to homology coordinates:
\[
\widehat F_v=-F_v,\qquad
\widehat\sigma_v=
\begin{cases}S-\sigma_v,&v\text{ fixed},\\
\sigma_v,&v\text{ paired},\end{cases}
\]
where $S$ is the common conductor section class.
Fix $S$ and all $D_e$ under this change. It is an involution.
A fixed-vertex relation changes to its negative plus
$2(S-d_eD_e)$, an integral sum of the conductor relations
along the interval; paired-vertex relations are fixed.
Thus this is an exact integral change of presentation with
the same inverse. Hatted vectors are homology coordinates;
the complex orientation and section-square assertion above
refer to the unhatted geometric classes.

In these coordinates the complete normalization presentation is
\begin{equation}\hypertarget{tag-024A}{}\label{maps:eq:D.SL12}
\begin{cases}
 2\widehat\sigma_v-\widehat F_v-d_eD_e=0,&v\text{ fixed},\\
 \widehat\sigma_v-d_eD_e=0,&v\text{ paired},
\end{cases}
\qquad
d_e=\begin{cases}2,&e\text{ fixed},\\1,&e\text{ paired}.\end{cases}
\end{equation}
The same equations hold in the geometric unhatted coordinates.
The conductor mapping cylinder collapses to the gluing.
Its $H_1$ map is injective by peeling the incidence interval:
each end map is an injective finite-covering lattice map.
Its $H_2$ map is likewise injective, using the end coefficient,
including the unit ruling coefficient at a fixed vertex.
Consequently \eqref{maps:eq:D.SL12} is the entire $H_2$
presentation.

When $f=1$, all paired sections are $2D_0$ and
$\widehat F_{\rm fixed}=2\widehat\sigma_{\rm fixed}-2D_0$.
When $f=2$, each fixed end satisfies
$\widehat F_i=2\widehat\sigma_i-S$.
These substitutions, reversed for the inverse, give free rank
$\rho+1$. When $f=0$ all sections equal $2D_0$ and
\begin{equation}\hypertarget{tag-024B}{}\label{maps:eq:D.SL13}
S=2D_0=2D_1,\qquad \tau_2=D_1-D_0,\qquad 2\tau_2=0 .
\end{equation}
For section coefficients $c_i$ and endpoint coefficients
$d_0,d_1$, the exact residual quotient is
$(d_0+d_1+2\sum c_i,[d_1]_2)$.
Its inverse includes $D_0$ and $D_1-D_0$.
Thus $\tau_2$ has exact order two.
The injective $H_2$ conductor map leaves
$H_3(U)=\mathbb Z^{2\rho}$ and $H_4(U)=\mathbb Z^\rho$.
The $H_3$ generators are the sphere fibres over the two base
circles on each normalization, with inverse integration over
the sphere fibre; the $H_4$ generators are its fundamental
classes. All higher groups vanish.

The free-cover path lifting gives
\begin{equation}\hypertarget{tag-024C}{}\label{maps:eq:D.SL14}
\begin{split}
\Gamma=\langle\beta,\lambda,\ell,g\mid&
[\beta,\lambda]=[\beta,\ell]=[\lambda,\ell]=1,\\
&g\beta g^{-1}=\beta^{-1}\lambda^{-n},\
[g,\lambda]=[g,\ell]=1,\
g^{2k}=\lambda^R\ell^N\rangle .
\end{split}
\end{equation}
The inverse writes a lifted path in its root sheet followed by
its three cover periods. In the abelianization put
$\eta_0=\beta+(n-1)\lambda/2$; then
$\beta=n\eta_0$, $\lambda=-2\eta_0$.
For $uN+2kv=1$ put $h_{\rm ab}=ug+v\ell$. The inverse is
$g=Nh_{\rm ab}-2vR\eta_0$, $\ell=2kh_{\rm ab}+2uR\eta_0$.
Hence $H_1(U)=\mathbb Z^2$, in the ordered basis
$(\eta_0,h_{\rm ab})$.
In boundary coordinates $(\eta_0,\ell,\mu)$ the map is
\[
(a,b,c)\longmapsto(a+2uRb-2vRc,\,2kb+Nc),
\qquad (x,y)\longmapsto(x,vy,uy).
\]
The second is a section. The identity minus its composite with
the first map is $(ub-vc)(-2R,N,-2k)$, the exact primitive kernel.
This completes \eqref{maps:eq:D.SL3}.

\emph{The fibre and the complete torsion marking.}
The return relations in $\bigwedge^2\Lambda$ are generated
primitively by
$\alpha\lambda,\alpha\ell,\alpha\beta+2\beta\lambda,
2\beta\ell+n\lambda\ell$.
For the first one use the $n\alpha\lambda$ and
$2\alpha\lambda$ columns and $\gcd(n,2)=1$.
In wedge order $(12,13,14,23,24,34)$ the quotient and section are
\[
z\longmapsto(z_{23}-2z_{12},\,nz_{24}-2z_{34}),\qquad
(x,w)\longmapsto xe_{23}+w(e_{24}+(n-1)e_{34}/2).
\]
Their composites differ exactly by those four relation columns.
The fixed lattice is $\mathbb Z\delta+\mathbb Z\ell$.
Its marked origin-normal tori give the Wang section, so the
complete boundary group is the stated $\mathbb Z^4$.

In the displayed holomorphic chart, $\alpha$ is the positive
angular coordinate and $\beta$ moves radially outward.
Thus the ordered $\alpha\beta$ torus specializes to minus the
sum of the complex-oriented ruling spheres. In the hatted
coordinates this gives exactly
\begin{equation}\hypertarget{tag-024D}{}\label{maps:eq:D.SL15}
2x=-\sum_v|v|\widehat F_v,\qquad 2w=-kS .
\end{equation}
For the second equation the transverse $\lambda\ell$ torus
has degree $k$ to a paired normalization section.
At a fixed vertex its degree is $2k$, and its class is
$k(2\sigma_v-F_v)$: its coefficient on the base is $2k$,
and disjointness from the conductor bisection fixes the ruling
coefficient to $-k$. Hence it is $kS$ in either case.

Put
$\widehat T=\sum_{v\ {\rm paired}}\widehat F_v+
\sum_{v\ {\rm fixed}}\widehat\sigma_v$.
The torsion-free cases of \eqref{maps:eq:D.SL15} give
\begin{equation}\hypertarget{tag-024E}{}\label{maps:eq:D.SL16}
\begin{array}{c|c|c}
f&x&w\\\hline
1&-\widehat T+D_0&-kD_0\\
2&-\widehat T+S&-(k/2)S .
\end{array}
\end{equation}
The conductor substitutions above are inverse coordinate maps.
They show that $x$ and the indicated conductor class form a
primitive free pair, with $\rho-1$ unused free coordinates.

For the remaining parity take $uN+2kv=1$ and
$h_\Gamma=g^u\ell^v$. The exact inverse words are
$g=h_\Gamma^N\lambda^{vR}$ and
$\ell=h_\Gamma^{2k}\lambda^{-uR}$.
Under abelianization, $h_\Gamma$ maps to
$u[g]+v[\ell]=h_{\rm ab}$ for the same $u,v$.
Thus $\Gamma$ is the mapping-torus group of
$\beta\mapsto-\beta-n\lambda$, $\lambda\mapsto\lambda$.
Its integral Wang injection and literal torus section give
\begin{equation}\hypertarget{tag-024F}{}\label{maps:eq:D.SL17}
H_2(\Gamma)=\mathbb Z F\oplus\mathbb Z/2\,T_{\rm W},\qquad
 F=[\lambda,h_\Gamma],\quad T_{\rm W}=[\beta,\lambda].
\end{equation}
An exact central extension realizing these markings adjoins
central $F,T_{\rm W}$ with $[\beta,\lambda]=T_{\rm W}$, $2T_{\rm W}=0$, and lets
$h_\Gamma$ act by
$\beta\mapsto\beta^{-1}\lambda^{-n}$,
$\lambda\mapsto\lambda F^{-1}$.
This is an automorphism since it sends $T_{\rm W}$ to $-T_{\rm W}=T_{\rm W}$.
The semidirect product realizes both central generators.
It evaluates the two-sector bar prism as
$[\beta,h_\Gamma^2]=-nF+nT_{\rm W}$. Therefore
\begin{equation}\hypertarget{tag-024G}{}\label{maps:eq:D.SL18}
[\beta,\ell]\longmapsto-knF-(kn+uR)T_{\rm W},\qquad
 [\lambda,\ell]\longmapsto2kF .
\end{equation}
The equivalent torsion coefficient $kn-uR$ differs by $2knT_{\rm W}$.
This also gives the ordered quadratic correction
$-n\binom{x_\beta}{2}T_{\rm W}$ by collecting powers of
$\beta^{-1}\lambda^{-n}$.

Choose the fixed edges at the centres $j/2$ and
$(j+q)/2$. Their stabilizer words are $g$ and $\beta g$.
Their primitive transverse periods are
$h_\Gamma$ and $(\beta g)^u\ell^v
=\beta h_\Gamma-n(u-1)\lambda/2$.
The inverse solves that last equality for $\beta h_\Gamma$.
Since $u$ is odd, their oriented torus classes satisfy
\begin{equation}\hypertarget{tag-024H}{}\label{maps:eq:D.SL19}
D_0\longmapsto F,\qquad D_1\longmapsto F-T_{\rm W} .
\end{equation}
The universal cover of the central core is the infinite
sphere chain times $\mathbb R^2$. Its $H_2$ is generated by
the lifted ruling spheres. The Hopf sequence therefore
identifies $H_2(\Gamma)$ with the quotient of $H_2(U)$ by
their images. The exact normalization quotient here is
$\mathbb ZD_0\oplus\mathbb Z/2\,\tau_2$; by
\eqref{maps:eq:D.SL19} its inverse identifies $D_0$ with
$F$ and $\tau_2$ with $T_{\rm W}$.
This detects the full order-two ambiguity in the cap.

When $f=0$, the integers $k$ are even and $j,R,u$ are odd.
Thus \eqref{maps:eq:D.SL18} sends $w$ to $-kF+T_{\rm W}$ and $x$
to $T_{\rm W}$. Combining this with their free coordinates yields
\begin{equation}\hypertarget{tag-024J}{}\label{maps:eq:D.SL20}
x=-\sum_v\widehat F_v+\tau_2,\qquad
 w=-kD_0+\tau_2 .
\end{equation}
Killing the two fibre images sets $\tau_2=kD_0$ and
$2kD_0=0$, so the finite part of the fibre cokernel is
cyclic $\mathbb Z/(2k)$, with these exact inverse coordinates.
For $f=1,2$ its corresponding finite part is $\mathbb Z/k$
and $\mathbb Z/(k/2)$.

\emph{Normal discs, all integer lifts and full kernels.}
The root-disc boundary is
$m\mu-(-r\alpha+R\lambda+N\ell)$.
Sweeping it by $\delta$ and $\ell$, with the ordered seam
parallelogram, gives the exact cap relations
\begin{equation}\hypertarget{tag-024K}{}\label{maps:eq:D.SL21}
2kD+4Nw=0,\qquad 2kL-2Rw=0 .
\end{equation}
The omitted angular wedges vanish in the explicitly computed
fibre coinvariant. Cutting the same $m$ sectors and seam
parallelogram is the inverse chain construction.
These equations determine all free coordinates.
For $f=0$ the classifying tori give
$D=2[\lambda,g]\mapsto2NF$ and
$L=[\ell,g]\mapsto-RF$, with zero $T_{\rm W}$ coordinate.
The determinant is $-R(uN+2kv)=-R$.
Consequently, in the chosen-centre conductor marking,
\begin{equation}\hypertarget{tag-024L}{}\label{maps:eq:D.SL22}
\begin{array}{c|c|c}
f&D&L\\\hline
1&2ND_0&-RD_0\\
2&NS&-(R/2)S\\
0&2ND_0&-RD_0 .
\end{array}
\end{equation}

In the reduced-component marking,
write $j=j_0+h_{\rm lift}q$, $0\le j_0<q$.
Then $g=\beta^{h_{\rm lift}}g_0$, $R=R_0-h_{\rm lift}q$,
and the exact normal-torus change is
\begin{equation}\hypertarget{tag-024N}{}\label{maps:eq:D.SL24}
D=D_{\rm red},\qquad L=L_{\rm red}-h_{\rm lift}nw .
\end{equation}
Its inverse adds $h_{\rm lift}nw$, since
$[\delta\beta]=0$ and $[\beta\ell]=nw$.
At $f=0$ the chosen-centre class for $j$ is
$D_0=D_{0,\rm red}+h_{\rm lift}\tau_2$.
Thus \eqref{maps:eq:D.SL22}, in the reduced-component conductor marking,
has $L=-RD_{0,\rm red}+h_{\rm lift}\tau_2$.
The free relation already forces the coefficient of $L$
to be even, so its contribution to a kernel relation is zero.
These formulas give the correction for every positive or
negative integer lift.

The independent free coordinate $x$ forces its coefficient
in a boundary kernel vector to vanish. The remaining three
columns in \eqref{maps:eq:D.SL16},
\eqref{maps:eq:D.SL20} and \eqref{maps:eq:D.SL22} give exactly
\eqref{maps:eq:D.SL5}, including the even $w$ coefficient at $f=0$.
Its projection to $(c,d)$ is injective, with inverse
$b=(2Nc-Rd)/k$ or the displayed halved formula.
The projected lattices are the kernels of the surjective rows
\[
(2N,-R)\bmod k,\quad
(N,-R/2)\bmod(k/2),\quad
(2N,-R)\bmod 2k
\]
respectively. Their surjectivity follows from
$\gcd(k,2N)=1$, $\gcd(k/2,N)=1$, or
$\gcd(2N,R,2k)=1$ in the corresponding cases.
These are the three projected indices in \eqref{maps:eq:D.SL5}.

These kernels have explicit inverse basis coordinates.
For a row $(A,B)$ modulo $M$ put $d_{\rm row}=\gcd(A,B)$,
$A=d_{\rm row}A_0$, $B=d_{\rm row}B_0$, and choose $p_{\rm row}A_0+t_{\rm row}B_0=1$.
Here $\gcd(d_{\rm row},M)=1$. The unimodular coordinates
$\xi=A_0c+B_0d'$, $\eta=-t_{\rm row}c+p_{\rm row}d'$ have inverse
$c=p_{\rm row}\xi-B_0\eta$, $d'=t_{\rm row}\xi+A_0\eta$.
Membership is exactly $\xi\in M\mathbb Z$.
Thus $(Mp_{\rm row},Mt_{\rm row})$ and $(-B_0,A_0)$ are a basis of the
projected kernel, and the displayed formula recovers $b$.

The full boundary image has primitive complement
$\mathbb Z^{\rho-1}$. Here are explicit lifts for its
remaining conductor coordinates.
For $f=1$ choose $p_{{\rm lift},f}k+t_{{\rm lift},f}(2N)=1$ and use $-p_{{\rm lift},f}w+t_{{\rm lift},f}D$ for $D_0$.
For $f=2$ choose $p_{{\rm lift},f}(k/2)+t_{{\rm lift},f}N=1$ and use the same expression
for $S$. For $f=0$ choose $p_{{\rm lift},f}k+t_{{\rm lift},f}N=1$ and put
$V_{\rm lift}=-2p_{{\rm lift},f}w+t_{{\rm lift},f}D$. Then
$D_0=-L+(1-R)V_{\rm lift}/2$ and $\tau_2=w+kD_0$.
Both expressions are integral and have the asserted free
and torsion coordinates. Together with the primitive $x$
coordinate and the normalization substitutions, these are
inverse quotient maps proving the entire free cokernel.
Equivalently it injects into
$H_2(U,\partial U)=H^4(U)=\mathbb Z^\rho$.

In the integral bar splitting of \ref{maps:src:D.2.41},
a marked origin normal at the site has coordinates
$(c_i(u),J_i u)$. Thus a local kernel column $(0,b,c,d)$
is inserted as
\[
\bigl(b[w]+cc_i(\delta)+dc_i(\ell),\
                  cJ_i(\delta)+dJ_i(\ell)\bigr).
\]
The representative of $[w]$ is the explicit fibre section
above, transported by $\bigwedge^2H$.
Restricting to the boundary, reading its two Wang coordinates
and applying that fibre quotient is the inverse test for
the complete kernel. Hence the other cap and finite
Thom columns belong to the same chain group.
The component fundamental classes give precisely $\rho$
free relative-disc coordinates and the cap has Euler number zero.

\emph{The global one-star calculation.}
For the minimal natural exterior the fibre quotient
is $\mathbb Zs_0$ with evaluation $\Omega$.
Its fixed-normal columns span the primitive graph kernel
with sole relation
$\sum_i t_i(\delta)+b_0e\,s_0=0$.
These maps and their inverses are the cut-tree
period maps of \ref{maps:src:D.2.26} and \ref{maps:src:D.2.27}.
The other minimal caps kill their fixed normals and have
zero fibre-kernel evaluation. Thus
$D=-4Be\,s_0$, $w=Bs_0$, and a star kernel column gives
\begin{equation}\hypertarget{tag-024M}{}\label{maps:eq:D.SL23}
\bigl(B(b-4ec),d\bigr)\quad\hbox{on }(s_0,L).
\end{equation}
Set $u_{\rm rel}=b-4ec$ and $s_{\rm ord}=N-2ke$.
For $f=1$ the condition is $ku_{\rm rel}+Rd=2s_{\rm ord}c$.
Choose $v_{{\rm inv},f} k=1\bmod 2|s_{\rm ord}|$. The relation lattice has basis
$(2Bs_{\rm ord},0),(-Bv_{{\rm inv},f}R,1)$.
Its quotient map is
$(x,y)\mapsto[x+Bv_{{\rm inv},f}Ry]_{2B|s_{\rm ord}|}$, with inverse $[z]\mapsto zs_0$.
For $f=2$, choose $v_{{\rm inv},f}(k/2)=1\bmod |s_{\rm ord}|$ and replace
$2s_{\rm ord},R$ by $s_{\rm ord},R/2$ in those maps. This gives
$\mathbb Z/(B|s_{\rm ord}|)$, also when $|s_{\rm ord}|=1$.
For $f=0$ both $u_{\rm rel},d$ are even. Write $u_{\rm rel}=2u_{0,\rm rel},d=2d_0$,
choose $v_{{\rm inv},f}k=1\bmod |s_{\rm ord}|$, and obtain the columns
$(2Bs_{\rm ord},0),(-2Bv_{{\rm inv},f}R,2)$.
With $Y_0=L-Bv_{{\rm inv},f}R s_0$ the inverse quotient maps are
\[
\begin{gathered}
(x,y)\longmapsto
([x+Bv_{{\rm inv},f}Ry]_{2B|s_{\rm ord}|},[y]_2),\\
([z],[t])\longmapsto zs_0+tY_0 .
\end{gathered}
\]
These prove the full torsion groups at all parameters.

For the fundamental group the selected cusp and its transported
vanishing words kill the two elliptic generators and $\delta$.
At the exceptional $\mathrm{OS46/II}$ position the first vanishing words
have index three; the star disc bounding
$\alpha=2e_1+\delta$ supplies the coprime twofold relation
and kills $e_1$, then $e_2$. The remaining fibre loop is
the gamma return, whose exponent gives the inverse
quotient map. The star group is generated by its boundary,
so van Kampen gives
$\langle\gamma,c\mid[\gamma,c]=1,\
c^{2k}=\gamma^N,\ c=\gamma^e\rangle$.
Sending $\gamma$ to $1\bmod |s_{\rm ord}|$ and $c$ to $e$,
with inverse the gamma loop, proves the group claim.
$s_{\rm ord}$ is odd, nonzero and prime to $2k$.

Every boundary $H_1$ in these twelve natural diagrams is free:
the return matrices have only unit nonzero degree-one
Smith factors. With $t_{\partial}$ boundary components the exterior
has $H_1$ rank $t_{\partial}$, and each boundary rank exceeds its
minimal cap rank by one. The full degree-one Mayer--Vietoris
source and target therefore have equal rank. Its cokernel
is the finite group just computed, so its integer kernel
is zero. Thus $H_2$ is exactly the free component complement
$\mathbb Z^{\rho-1}$ plus the displayed torsion group.

Finally the selected primitive rank-two cusp contributes
Euler number two, and the exterior, boundaries and all
other caps contribute zero. Oriented integral duality and
UCT give \eqref{maps:eq:D.SL8}.
On each finite Smith factor the two-term resolution
$\mathbb Z\xrightarrow{d}\mathbb Z$ identifies its Ext
generator with the dual unit cochain; cap product gives
the corresponding $H_3$ or $H_4$ generator, with inverse
its residue modulo $d$. This supplies the torsion isomorphisms.
The Euler identity then gives $b_3=2\rho-2$.
For the six default $b_0=4$ positions choose
$k=2h$ with $h\in\mathbb Z_{>0}$, $j$ even and $N=2ke\pm1$.
Then $f=2$, $\rho=nh+1$, $s_{\rm ord}=\pm1$ and $B=1$,
which gives \eqref{maps:eq:D.SL9} and simple connectivity.
\end{proof}
\subsection{Vertical corrections}

\begin{definition}[Vertical correction data]\hypertarget{tag-Q202}{}\label{maps:def:vertical-correction-data}
{\normalfont\scriptsize tag Q202}\par
Let $f:S\to\mathbb P^1$ be one of the twenty-one marked rational
elliptic surfaces of \ref{maps:src:D.2.19}, with zero section $O$,
Mordell--Weil group $\mathbb ZG$ and height
$\langle G,G\rangle=1/d_0$.
For integers $a,b$, put $P=aG$ and $Q=bG$.
For an integral divisor $D$ supported on fibre components, define
\[
 M_D=\mathcal O_S(P-O+D),\qquad
 J_D=M_D\otimes T_Q^*M_D^{-1},
\]
where $T_Q:S\to S$ is section translation. A \emph{vertical correction}
is such a divisor $D$ for which $J_D$ is a pullback from the base.
\end{definition}

\begin{theorem}[Vertical corrections and fillings]\hypertarget{tag-02LS}{}\label{maps:src:D.2.50}
{\normalfont\scriptsize tag 02LS}\par
For the data of Definition~\ref{maps:def:vertical-correction-data},
a vertical correction exists precisely when $d_0\mid ab$. For every such correction the exact identity is
\begin{equation}
                  J_D\simeq f^*\mathcal O_{\mathbb P^1}(ab/d_0).
                                                        \hypertarget{tag-026E}{}\label{maps:eq:D.VC1}
\end{equation}
The local equation \eqref{maps:eq:D.VC3} determines every correction;
its ambiguity is the invariant component-orbit lattice
\eqref{maps:eq:D.VC4}.
For $a,b>0$ and any nonnegative whole-fibre weights of total $ab/d_0$,
the corrected punctured family has exactly the marked monodromies
\eqref{maps:eq:D.GI2}.

The following OS45 fillings are simply connected.
On OS45/default, choose $(a,b)=(2,4)$ or $(4,2)$, a simple
linearization zero at any of its four $I_1$ fibres, multiplicative
Mumford fillings at every singular fibre, no extra affine boundary
translations, and ordinary gamma clutch $e=\pm1$. Use the
vertical corrections \eqref{maps:eq:D.VC10} with the cap origins
and affine boundary maps defined by the tensor transport in the proof. All sixteen choices have
\begin{equation}
 \begin{gathered}
 \pi_1(X)=1,\qquad H_0(X)=H_6(X)=\mathbb Z,\\
 H_2(X)=H_3(X)=\mathbb Z/b,\qquad
 H_1(X)=H_4(X)=H_5(X)=0 .
 \end{gathered}                                             \hypertarget{tag-026L}{}\label{maps:eq:D.VC7}
\end{equation}

\end{theorem}

\begin{proof}
\begin{enumerate}
\item \emph{Actual component maps and the integral equation.}
Let $C_0,\ldots,C_{v_{\rm comp}-1}$ be the components of a reducible fibre,
with $C_0$ meeting $O$, let $\mu$ be its primitive positive
multiplicity column, and put $C=(C_iC_j)$. In particular $\mu_0=1$.
Write $p_{\rm comp},q_{\rm comp}$ for the component labels of $P,Q$.
If the actual translation $T_Q$ sends component $i$ to $\sigma(i)$,
write $\mathsf S$ for pullback on columns: $(\mathsf S z)_i=z_{\sigma(i)}$.
Then $C\mathsf S=\mathsf S C$, $\mathsf S\mu=\mu$, and $\mathsf S^{-1}=\mathsf S^t$.
Write $z_i$ for the coefficient of $C_i$ in $D$. The degree column
of the section line and its correction are
\begin{equation}
 m_P=e_{p_{\rm comp}}-e_0,\qquad m=m_P+Cz,\qquad \mu^t m=0.
                                                        \hypertarget{tag-026G}{}\label{maps:eq:D.VC2}
\end{equation}
Thus $J_D$ has zero degree on every component precisely when
\begin{equation}
       (\mathsf S-I)Cz=-(\mathsf S-I)(e_{p_{\rm comp}}-e_0).
                                                        \hypertarget{tag-026H}{}\label{maps:eq:D.VC3}
\end{equation}
The pushforward convention has matrix $\mathsf S^{-1}$; its equation is
obtained by the unimodular left multiplication $-\mathsf S^{-1}$, whose
inverse is $-\mathsf S$. This proves the exact equivalence of the two
permutation conventions. All lattices in this proof are integral.

If $z,z'$ solve this equation, then
$C(\mathsf S-I)(z-z')=0$. The integral kernel of $C$ is $\mathbb Z\mu$,
so $(\mathsf S-I)(z-z')=k\mu$. Pairing with $\mu$ gives
$0=k\mu^t\mu$ and hence $k=0$. Conversely an invariant column
can be added to a solution. The precise ambiguity is therefore
\begin{equation}
           z+\ker(\mathsf S-I)\subset\mathbb Z^{v_{\rm comp}}.
                                                        \hypertarget{tag-026I}{}\label{maps:eq:D.VC4}
\end{equation}
The kernel has the basis consisting of the sums of the component basis
vectors in each orbit. Reading the common coefficient on a chosen
orbit representative is the inverse coordinate map. Whole fibres
are included, as are the other invariant orbit sums.

\item \emph{A rational function and the exact base line.}
On the generic elliptic curve the divisor
$P-O-(P-Q)+(-Q)$ is principal. For distinct nondegenerate points,
take the line through $P,-Q$, divided by the vertical line through
$P-Q$. If $P=Q$, use $x-x(P)$; if $P=-Q$, use the tangent line
at $P$, divided by the vertical line at $2P$. If $P=0$ or $Q=0$,
the divisor is zero and the function is $1$. These are the actual
group-law functions, with reciprocal inverse maps.

Multiplying by the canonical rational sections of $\mathcal O(D)$
and $\mathcal O(-T_Q^*D)$ gives a rational section of $J_D$ whose
remaining divisor is vertical. If \eqref{maps:eq:D.VC3} holds,
the coefficient column of that divisor at each fibre belongs to
$\ker C=\mathbb Z\mu$. Since $\mu_0=1$, its whole-fibre
coefficient is an integer. Therefore $J_D=f^*\mathcal O(t_{\rm base})$ for
an integer $t_{\rm base}$. The isomorphism is multiplication by the displayed
function, the canonical vertical sections, and a degree-zero base
rational function identifying the whole-fibre divisor with $t_{\rm base}[\infty]$.
Multiplication by their reciprocals is its inverse. Restriction to
$O$ makes pullback from the base injective. Conversely a base
pullback has zero component degrees, proving necessity of the local
equations at the level of actual line bundles.

To compute $t_{\rm base}$, let $f_{q_{\rm comp}}$ be the rational column determined by
$Cf_{q_{\rm comp}}=e_{q_{\rm comp}}-e_0$ and $f_{q_{\rm comp}}(0)=0$. The norm of $f_{q_{\rm comp}}$ under $\mathsf S$ has
$C$-image zero, because the orbit sum of $e_{q_{\rm comp}}-e_0$ telescopes.
It is thus a rational multiple of $\mu$. Since $m$ is invariant
and $m^t\mu=0$, we obtain $m^tf_{q_{\rm comp}}=0$, and hence
\[
             z_{q_{\rm comp}}-z_0=-f_{q_{\rm comp}}(p_{\rm comp})=\operatorname{contr}_v(P,Q).
\]
The mixed height formula on the rational elliptic surface is
\[
 1+P\mathbin{\cdot}O+Q\mathbin{\cdot}O-P\mathbin{\cdot}Q
   =\langle P,Q\rangle+\sum_v\operatorname{contr}_v(P,Q).
\]
Here the local terms are the negative inverse-Cartan pairings;
see \cite[author version arXiv:0907.0298v3, Sections 7, 8 and 11.8, pp.32--34 and 51--52]{manSS10}.
Intersecting the rational section of $J_D$ with $O$ adds
$\sum_v(z_0-z_{q_{\rm comp}})$ to the left side. The terms cancel literally,
so $t_{\rm base}=\langle P,Q\rangle=ab/d_0$. This proves the asserted
identity \eqref{maps:eq:D.VC1}, its necessity, and independence
from every invariant-orbit ambiguity. The argument includes signed
$a,b$ and the zero cases; positivity is not used here.

\item \emph{Integral solutions for every component pair.}
At an $I_n$ fibre use the positive-valuation Tate labels
$\mathbb Z/n$. The actual translation sends $i$ to $i+q_{\rm comp}$.
The cycle matrix has diagonal $-2$ and one $+1$ for each incident
edge. At $n=2$ both parallel edges are retained, giving off-diagonal
entry $2$. Its image is exactly
\[
 \left\{h\in\mathbb Z^n:\ \sum_i h_i=0,
                  \ \sum_i i h_i=0\pmod n\right\}.
\]
Indeed, for such $h$, put
\[
 \begin{gathered}
 r_{{\rm lap},0}=-\frac1n\sum_{i=0}^{n-1}(n-1-i)h_i,
 \quad r_{{\rm lap},i}=r_{{\rm lap},0}+\sum_{j<i}h_j,\\
 z_0=0,\quad z_i=\sum_{k=1}^{i}r_{{\rm lap},k}.
 \end{gathered}
\]
The moment condition makes $r_{{\rm lap},0}$ integral; augmentation zero
makes the differences cyclic. Then
$(Cz)_i=r_{{\rm lap},i+1}-r_{{\rm lap},i}=h_i$, including the two-edge case.
Subtracting the common coefficient $z_0$ is the inverse choice
modulo the whole-fibre constant. At $n=1$ the degree-zero lattice
is zero.

Put ${g_{\rm comp}}=\gcd(n,q_{\rm comp})$ and $\ell=n/{g_{\rm comp}}$. An invariant multidegree is
constant on the ${g_{\rm comp}}$ residue classes modulo ${g_{\rm comp}}$. Its augmentation
condition cancels the common part of the orbit moments, giving
\[
 \sum_i i m_i=\ell\sum_{j=0}^{{g_{\rm comp}}-1}j m_j\pmod n.
\]
For $m=e_{p_{\rm comp}}-e_0+Cz$ this moment is $p_{\rm comp}$, so $\ell\mid p_{\rm comp}$,
equivalently $n\mid p_{\rm comp}q_{\rm comp}$. If $p_{\rm comp}=0$ or $q_{\rm comp}=0$, take $z=0$.
Otherwise ${g_{\rm comp}}>1$, and the invariant column
\[
 m_i=(p_{\rm comp}/\ell)
       \bigl({\bf1}_{i=1\bmod {g_{\rm comp}}}-{\bf1}_{i=0\bmod {g_{\rm comp}}}\bigr)
\]
has the required moment. Apply the displayed inverse to
$h=m-e_{p_{\rm comp}}+e_0$. This constructs an integral solution for every
admissible signed residue pair.

For the other types, the component orders and generator
self-pairings from the deleted-zero-component matrix are
\[
\begin{array}{c|rrrrrrr}
 \text{type}&II&III&IV&III^*&IV^*&I_1^*&I_3^*\\\hline
 n&1&2&3&2&3&4&4\\
 \operatorname{contr}(1,1)&0&1/2&2/3&3/2&4/3&5/4&7/4.
\end{array}
\]
The preceding identity $z_{q_{\rm comp}}-z_0=\operatorname{contr}(p_{\rm comp},q_{\rm comp})$
shows that this pairing must be integral. Each numerator is a
unit modulo its component order; therefore the exact condition is
\begin{equation}
                  \operatorname{ord}(\Phi_v)\mid p_{\rm comp}q_{\rm comp}.
                                                        \hypertarget{tag-026K}{}\label{maps:eq:D.VC6}
\end{equation}
For orders $2$ and $3$ this makes one section narrow, and $z=0$ works.
For either odd star the only further pair is $(2,2)$. Label the
near leaves $0,1$, their adjacent inner vertex $2$, the inner
chain $2,\ldots,r-2$, and the far leaves $r-1,r$, where $r=5$ or $7$.
The multiplicities are $1$ on the four leaves and $2$ on the chain.
In $D_r=\{x\in\mathbb Z^r:\sum x_i\text{ even}\}$ the near
weight is $e_1$, while a far weight is $(e_1+\cdots+e_r)/2$.
Its square is $r/4$ and twice it equals the near class modulo
$D_r$. This identifies the simple-component labels. Actual
translation by class $2$ swaps the two near leaves, swaps the two
far leaves and fixes the chain. Taking $z=e_1$ gives
$m=-e_0-e_1+e_2$, invariant under that literal permutation.
Its inverse is the same permutation. This proves sufficiency
without replacing the marked components by an abstract graph.

\item \emph{The twenty-one profiles and all residues.}
The marked component coordinates of $G$ are units in every
nontrivial cyclic component group. Their nontrivial orders are
pairwise coprime and have the following products:
\[
\begin{array}{c|c|c}
 \mathrm{OS}&\text{nontrivial orders}&d_0\\\hline
 43&(2)&2\\45&(8)&8\\46&(4)&4\\47&(7,2)&14\\
 49&(2,3)&6\\50&(4,3)&12\\55&(5,4)&20\\56&(5,3,2)&30.
\end{array}
\]
Every collision profile in \ref{maps:src:D.2.19} retains these marked component
orders. The exact defining matrices and cocycles determine the
orders too: in a Smith decomposition ${\mathsf U}_{\rm ell}(B-I){\mathsf V}_{\rm ell}=\Delta_{\rm ell}$, the class of
$g$ has order
\[
 \operatorname{lcm}_{j\leq\operatorname{rank}(B-I)}
       \frac{|(\Delta_{\rm ell})_{jj}|}{\gcd(|(\Delta_{\rm ell})_{jj}|,({\mathsf U}_{\rm ell}g)_j)}.
\]
The zero Smith rows of ${\mathsf U}_{\rm ell}g$ vanish. For each marked pair $(B,g)$ this gives the stated component order $h_{\rm comp}$.
An integral preimage of $h_{\rm comp}g$ is ${\mathsf V}_{\rm ell}x_{\rm ell}$, where
$(x_{\rm ell})_j=h_{\rm comp}({\mathsf U}_{\rm ell}g)_j/(\Delta_{\rm ell})_{jj}$ on the nonzero Smith rows and $(x_{\rm ell})_j=0$
on the zero rows. Thus \eqref{maps:eq:D.VC6}
at all fibres is equivalent to $d_0\mid ab$. The constructed
local divisors sum to an integral global $D$.

Reduction to the component residues includes every signed integer.
The positive height-one pairs having a fibre where both section
classes are nonzero are
\[
\begin{array}{c|c|c}
 \mathrm{OS}&\text{profiles}&(a,b)\\\hline
45&\mathrm{default},II&(2,4),(4,2)\\
46&\mathrm{default},II&(2,2)\\
50&\mathrm{default},II,IV&(2,6),(6,2)\\
55&\mathrm{default},II&(2,10),(10,2).
\end{array}
\]
The table gives sixteen choices of profile and initial pair.
The OS45/default filling parameter set in the theorem is
$\{(2,4),(4,2)\}\times\{1,2,3,4\}\times\{-1,1\}$,
where the last two coordinates specify the nodal zero position
and ordinary clutch.

For example, on $I_4$ with $(p_{\rm comp},q_{\rm comp})=(2,2)$, $D=-C_0$ changes
$(-1,0,1,0)$ to $(1,-1,1,-1)$. The inverse ambiguity coordinates
are the coefficients of $C_0+C_2$ and $C_1+C_3$.

\item \emph{The punctured tensor and period maps.}
Off the singular fibres the canonical section $\sigma_D$ is
nonzero. The actual tensor map and inverse are
\[
 I_D:(p,v)\longmapsto(p,v\sigma_D(p)),\qquad
 I_D^{-1}:(p,v')\longmapsto(p,v'/\sigma_D(p)).
\]
On the arrow line it multiplies by
$\sigma_D(p)/\sigma_D(p+Q)$. An arrow is a map from the line
at $p+Q$ to that at $p$ and thus covers $-Q$. In actual local
frames the conjugated arrow multiplier is
$j_D(p)\sigma_D(p+Q)/\sigma_D(p)$.
The group-law function above identifies it with the uncorrected
Poincare arrow times its actual base section. A ratio constant
on each compact elliptic fibre is a base unit; retain it as
$q^{w_i}u_i(q)$, with $u_i$ nowhere zero.

For a local elliptic matrix $B$, choose a solution $z_G$ of
$(B-I)z_G=g$ with least positive denominator, and put
$J=\left(\begin{smallmatrix}0&1\\-1&0\end{smallmatrix}\right)$.
The actual logarithmic comparison and inverse on the transported
period lattice are
\[
 L=\begin{pmatrix}I_2&0&-bz_G\\a z_G^tJ&1&0\\0&0&1\end{pmatrix},
 \qquad
 L^{-1}=\begin{pmatrix}I_2&0&bz_G\\-a z_G^tJ&1&0\\0&0&1\end{pmatrix}.
\]
The cross term is zero because $z_G^tJz_G=0$.
These matrices describe the actual exponential and logarithmic
line-frame maps, with reciprocal multipliers. When denominators
occur, the domain and range are the transported lattices, equivalently
the lattices on the corresponding denominator cover. The displayed
matrices are inverse integral maps between those lattices.

The line row satisfies
$r_G=-g^tJB=z_G^tJ(B-I)$. In the literal marking
$(e_1,e_2,\delta,\gamma)$ the resulting monodromy and inverse are
\[
 T_i=\begin{pmatrix}B_i&0&bg_i\\r_i&1&c_i\\0&0&1\end{pmatrix},
 \quad r_i=-a g_i^tJB_i,\quad c_i=r_i(bz_{G,i})+w_i,
\]
\[
 T_i^{-1}=\begin{pmatrix}
 B_i^{-1}&0&-B_i^{-1}bg_i\\
 -r_iB_i^{-1}&1&r_iB_i^{-1}bg_i-c_i\\0&0&1
 \end{pmatrix}.
\]
The denominator of $z_G$ is the actual component order $h_i$.
It divides $d_0$, so $d_0\mid ab$ makes $c_i$ integral.
At an $I_\nu$ site, write $g=\kappa u_{\rm per}$,
$B-I=\nu u_{\rm per}(u_{\rm per}^tJ)$, and choose $u_{\rm per}^tJv_{\rm per}=1$; then
$z_G=(\kappa/\nu)v_{\rm per}$. This is the transverse period, with
denominator $\nu$ when $\kappa$ is a unit. Changing $z_G$ by
a fixed elliptic vector leaves $r_i z_G$ unchanged.

For each defining profile the ordered product is
\[
 \prod_i T_i=I+\left(\sum_i w_i-ab/d_0\right)E_{\delta\gamma}.
\]
This is also the exact base-degree cocycle of the constructed
line map. The coefficient matrix $E_{\delta\gamma}$ is fixed
by all blocks and has square zero, so distributing the prescribed
weights among singular or regular sites changes the product by
their sum. The total in the theorem makes this product the identity.

Choose a positive integer $h$ such that $hQ$
is narrow everywhere. It fixes every component divisor. Evaluation
of adjacent tensor factors gives the literal identity
\begin{equation}
 \bigotimes_{j=0}^{h-1}T_Q^{j*}J_D
 =M_D\otimes T_{hQ}^*M_D^{-1}
 =M_0\otimes T_{hQ}^*M_0^{-1}.
                                                        \hypertarget{tag-026J}{}\label{maps:eq:D.VC5}
\end{equation}
All canonical sections of $D$ telescope, with reciprocal inverse
maps. On the logarithmic cover the deck generator is exactly
$\gamma_h=h\gamma$, and its lattice inclusion is
\begin{equation}
 S_h=\operatorname{diag}(1,1,1,h),\qquad
                 T_iS_h=S_hT_{i,h}.
                                                        \hypertarget{tag-026M}{}\label{maps:eq:D.VC8}
\end{equation}
The inverse on the image divides the last coordinate by $h$.
Here $bg_i$, $bz_G$ and $w_i$ are multiplied by $h$, while $r_i$
is fixed. This proves the actual marked period comparison.
The multiplier units and any constant scalar are retained in the
contracting roof. Its normalized positive density has a strictly
increasing integral, whose inverse gives the chosen circle
parameter of the punctured quotient.

\item \emph{The OS45 data and its two actual corrections.}
The defining matrices and component labels are
\begin{equation}
\begin{gathered}
 (I_8,I_1,I_1,I_1,I_1),\quad d_0=8,
 \quad \overline G_{I_8}=3\pmod8,\\
 g_i=(3,0),(0,0),(2,-1),(1,-2),(0,0),\\
 (B_0,\ldots,B_4)=
 \begin{pmatrix}1&8\\0&1\end{pmatrix},
 \begin{pmatrix}1&1\\0&1\end{pmatrix},
 \begin{pmatrix}3&4\\-1&-1\end{pmatrix},
 \begin{pmatrix}3&1\\-4&-1\end{pmatrix},
 \begin{pmatrix}0&1\\-1&2\end{pmatrix}.
\end{gathered}
                                                        \hypertarget{tag-026N}{}\label{maps:eq:D.VC9}
\end{equation}
Use the following component coefficients:
\begin{equation}
\begin{array}{c|c|c}
 (a,b)&(D_{C_0},\ldots,D_{C_7})&m\\\hline
 (2,4)&(-1,-1,-1,0,0,0,0,0)&(0,0,1,-1,0,0,1,-1)\\
 (4,2)&(-1,0,0,1,1,1,0,0)&(1,-1,1,-1,1,-1,1,-1).
\end{array}
                                                        \hypertarget{tag-026F}{}\label{maps:eq:D.VC10}
\end{equation}
The multidegree column $m$ is invariant under the respective
shifts $4$ and $6$. At the four irreducible fibres take $D=0$.
Both height pairings are $1$.

The positive-valuation Tate labels are obtained from the stated
ones by $j\mapsto-j\pmod8$, an involution with matrix $F$ and
$FCF=C$. It conjugates the shift by $q_{\rm comp}$ to the shift by $-q_{\rm comp}$.
Thus the pairs in the Tate labels are $(p_{\rm comp},q_{\rm comp})=(2,4),(4,2)$.
The explicit Laplacian inverse above gives the Tate coefficient choices
\[
 z^{\rm Tate}_{2,4}=(0,0,1,1,1,0,0,0),\qquad
 z^{\rm Tate}_{4,2}=(0,0,1,1,2,1,1,0).
\]
Subtracting these from the reflected choices in \eqref{maps:eq:D.VC10} gives
\[
 E_{2,4}=(-1,0,-1,-1,-1,0,-1,-1),\qquad
 E_{4,2}=(-1,0,-1,0,-1,0,-1,0),
\]
invariant under shifts $4$ and $2$, respectively. The exact punctured
comparison is $v\mapsto v\sigma_E(p)$, with inverse division.
It conjugates the arrow by
$\sigma_E(p)/\sigma_E(p+Q)$, retaining all units.
Along $O$ its local expression is $q^{-1}u(q)$, with $u$ a
nowhere-zero holomorphic unit. On $q=\epsilon e^{2\pi it}$,
the angular displacement is $-t+\arg u(q)/(2\pi)$; the inverse
subtracts it. The radial multipliers are reciprocal. In particular
the peripheral origin shifts by $-\delta$, with inverse $+\delta$.

For each correction we use the punctured period map just proved
to define the cap origin and the specified affine boundary map.
Comparing corrections transports that origin and gluing by the
displayed tensor map and its inverse, including this shear.
We now compute the fillings of \eqref{maps:eq:D.VC10} with these
cap origins and affine boundary maps.

\item \emph{Rank-one cap maps.}
The multiplicative data are
\[
\begin{array}{c|c|c|c|c}
 i&\nu_i&\kappa_i&u_{{\rm per},i}&v_{{\rm per},i}\\\hline
0&8&3&(1,0)&(0,1)\\
1&1&0&(1,0)&(0,1)\\
2&1&1&(2,-1)&(1,0)\\
3&1&1&(1,-2)&(1,-1)\\
4&1&0&(1,1)&(-1,0).
\end{array}
\]
Embed $u_{\rm per},v_{\rm per}$ in the first two coordinates and put
$\mathcal S=(u_{\rm per},\delta)$, $\mathcal L=(v_{\rm per},\gamma)$.
At a weight-zero site the off-diagonal block between these
two lattices is
\[
 N=\begin{pmatrix}\nu&b\kappa\\
             -a\kappa&-ab\kappa^2/\nu\end{pmatrix}.
\]
Set $d=\gcd(\nu,a\kappa)$, $n_{\rm rk}=\nu/d$,
$A=a\kappa/d$, $b_{\rm rk}=b\kappa/n_{\rm rk}$. For every site in the two
specified pairs these are integers, $\gcd(n_{\rm rk},A)=1$ and
$\gcd(d,b_{\rm rk})=1$. Choose $n_{\rm rk}\theta+A\sigma_{\rm Bez}=1$ and $dx_{\rm Bez}+b_{\rm rk}y_{\rm Bez}=1$,
and put
\[
 \begin{gathered}
 C_S=\begin{pmatrix}n_{\rm rk}&\sigma_{\rm Bez}\\-A&\theta\end{pmatrix},
 \quad C_S^{-1}=\begin{pmatrix}\theta&-\sigma_{\rm Bez}\\A&n_{\rm rk}\end{pmatrix},\\
 C_L=\begin{pmatrix}x_{\rm Bez}&-b_{\rm rk}\\y_{\rm Bez}&d\end{pmatrix},
 \quad C_L^{-1}=\begin{pmatrix}d&b_{\rm rk}\\-y_{\rm Bez}&x_{\rm Bez}\end{pmatrix}.
 \end{gathered}
\]
Then $N=(n_{\rm rk},-A)^t(d,b_{\rm rk})$ and $NC_L=C_S\operatorname{diag}(1,0)$.
The ordered integral period matrix
\[
 H_{\rm rk}=(\mathcal S C_Se_1,\mathcal L C_Le_1,
                   \mathcal S C_Se_2,\mathcal L C_Le_2)
\]
has determinant $1$ and $H_{\rm rk}^{-1}TH_{\rm rk}=I+E_{12}$.
At $I_8$ one may use, respectively,
\[
\begin{array}{c|c|c}
(a,b)&C_S&C_L\\\hline
(2,4)&\begin{pmatrix}4&-1\\-3&1\end{pmatrix}
     &\begin{pmatrix}-1&-3\\1&2\end{pmatrix}\\
(4,2)&\begin{pmatrix}2&1\\-3&-1\end{pmatrix}
     &\begin{pmatrix}1&-3\\-1&4\end{pmatrix}.
\end{array}
\]
For both initial pairs, $\gcd(a,b)=2$; the displayed B\'ezout
relations give the required primitive period bases.

These period changes are realized on the actual torus charts.
Characters transform by $C_S^{-t}$, with inverse $C_S^t$;
the long-period multipliers are the indicated integer products
with exponents $C_L$, whose inverse products have exponents
$C_L^{-1}$. Retain the actual unit multipliers in both directions.
On a one-dimensional fan the node charts are
\[
 X=z/q^j,\qquad Y=q^{j+1}/z,\qquad XY=q;
 \qquad z=Xq^j
\]
on their overlap, with the displayed reciprocal inverse.
The primitive period shift identifies consecutive components.
The resulting normalization is a degree-zero ruled surface over
the torus of the remaining long periods. If its flat multipliers
are $\lambda_1,\lambda_2$, multiplication by
$\exp(-t_1\log\lambda_1-t_2\log\lambda_2)$ on the universal
cover is the smooth trivializing gauge; the positive exponent
is its inverse. A pole-gluing translation $x\mapsto x+a_0$
is removed by $(z,x)\mapsto(z,x-f_{\rm cut}(z)a_0)$, where $f_{\rm cut}$ is $0$
near one pole and $1$ near the other; the inverse adds that term.
This describes the actual nodal ruled normalization and its
unit maps, yielding a core $(\text{nodal }\mathbb P^1)\times T^2$.

In the $H_{\rm rk}$-coordinates, write
$z=(z_{12},z_{13},z_{14},z_{23},z_{24},z_{34})$.
Its specialization maps are
\[
 f_1^{\rm rk}(x)=(x_2,x_3,x_4),\qquad
 f_2^{\rm rk}(z)=(z_{12},z_{23},z_{24},z_{34}).
\]
Their sections insert zero in the omitted coordinates, and are
transported physically by $H_{\rm rk}$ and $\Lambda^2H_{\rm rk}$.
The integral fibre kernel is
$\langle e_{13},e_{14}\rangle$, with exact Wang preimages
$e_{23},e_{24}$. Thus the core has
$\pi_1=\mathbb Z^3$, $H_2=\mathbb Z^4$, $H_3=\mathbb Z^3$,
$H_4=\mathbb Z$, and Euler characteristic $0$. These identifications
are induced by the displayed normalization and node maps.

\item \emph{The selected triangular cap and all boundary kernels.}
At the selected $I_1$ site, $w=1$, $\nu=1$, and
\[
 C_+=\begin{pmatrix}1&b\kappa\\-a\kappa&1-ab\kappa^2\end{pmatrix},
 \qquad C_+^{-1}=\begin{pmatrix}1-ab\kappa^2&-b\kappa\\a\kappa&1\end{pmatrix}.
\]
Set $H_{\rm cusp}=(\mathcal S C_+e_1,\mathcal Le_1,
\mathcal S C_+e_2,\mathcal Le_2)$. This has determinant $1$,
and $H_{\rm cusp}^{-1}TH_{\rm cusp}=I+E_{12}+E_{34}$.
The prescribed fan lies in the actual short lattice $\mathcal S C_+$;
its characters transform by $C_+^{-t}$, with inverse $C_+^t$.
The primitive period cell has one vertex, three edges and two
triangles. Its normalization is the toric hexagonal surface,
with the three actual conductor identifications and all period
unit multipliers retained. In each triangle the product-coordinate
chart has $XYZ=q$; the lattice character maps above and their
inverse monomials identify adjacent charts and the period translates.
They construct the specified cap, including its boundary marking.

The normalization/conductor chain maps give
\[
 f_1^{\rm cusp}(x)=(x_2,x_4),\qquad
 f_2^{\rm cusp}(z)=(z_{12},z_{14}-z_{23},z_{24},z_{34}).
\]
Sections are $(s,t)\mapsto(0,s,0,t)$ and
$(r,s,t,u)\mapsto(r,0,s,0,t,u)$.
The kernel of $f_2^{\rm cusp}$ is the primitive lattice
$\langle e_{13},e_{14}+e_{23}\rangle$, with Wang preimages
$e_{23},e_{24}-e_{23}$.
The transverse double tori and the primitive conductor differences
give the four displayed $H_2$ coordinates; the two graph cycles
of the triangulation give $H_3=\mathbb Z^2$.
The universal dual triangulation is a plane, and its toric
components have trivial fundamental group. Thus the only
fundamental-group generators are the two long period deck
translations: $\pi_1=\mathbb Z^2$.
The whole fibre generates $H_4=\mathbb Z$. Open positive-dimensional
torus orbits have Euler characteristic zero; the two triangle
vertices contribute $2$, so this cap has Euler characteristic $2$.

For both cap types the actual fixed circles extend over the
toric charts. Sweeping an extending origin disc by such a circle
gives a three-chain whose boundary is the fixed-normal torus,
with sign fixed by the positive base meridian. The kernel of
boundary $H_2$ specialization is consequently generated by the
displayed primitive fibre kernel and these fixed-normal tori.
There are no additional quotient classes: $f_2^{\rm rk}$ and its
triangular-cap counterpart both have their displayed integral sections, and the normalization/conductor complex
is exact in the remaining degrees. Boundary $H_1$ specializes
onto core $H_1$ with kernel the positive normal meridian; the
primitive image $(T-I)\Lambda$ is exactly the short vanishing
lattice and the origin disc fixes its affine lift.

\item \emph{The full integral exterior and its bar inverse.}
Let $\Lambda=\mathbb Z^4$ and order $\Lambda^2\Lambda$ as
$(12,13,14,23,24,34)$. For the five singular matrices put
\[
 \begin{gathered}
 R=(\Lambda^2T_0-I\mid\cdots\mid\Lambda^2T_4-I),\\
 {\mathsf D}_{\rm gr}=(T_0-I\mid\cdots\mid T_4-I),\qquad
 K=\ker {\mathsf D}_{\rm gr}.
 \end{gathered}
\]
The exact Smith factors of ${\mathsf D}_{\rm gr}$ are $1,1,1$, so $K$ is primitive
of rank $17$. If ${\mathsf U}_{\rm gr}{\mathsf D}_{\rm gr}{\mathsf V}_{\rm gr}$ is its Smith form, let
$K_0:\mathbb Z^{17}\to K\subset\mathbb Z^{20}$ be the matrix of
columns $4$ through $20$ of ${\mathsf V}_{\rm gr}$, and let $K_0^{\vee}$ be the
matrix of the corresponding rows of ${\mathsf V}_{\rm gr}^{-1}$. Thus
\[
 K_0^{\vee}K_0=I_{17},\qquad
 K_0K_0^{\vee}|_K=\operatorname{id}_K.
\]
We use $K_0^{\vee}$ to read the chosen graph coordinates below.
The nonzero Smith factors of $R$ are $1,1,1,1,2$. Thus the fibre coinvariant is
$\mathbb Z\oplus\mathbb Z/2$; its torsion is retained.
If $U_RRV_R$ is the Smith form, the quotient map reads the last
coordinate of $U_Rz$ and the fifth modulo $2$, with inverse induced
by $U_R^{-1}$ on those two coordinate vectors.

Here is the actual bar section and inverse, with the conventions
of \ref{maps:src:D.2.41}. Put
$p(x,y)=\sum_{i<j}x_i y_j e_{ij}$ and use the central group
\[
 (x,z)(y,w)=(x+y,z+w+p(x,y)),\qquad
 (x,z)^{-1}=(-x,-z+p(x,x)).
\]
Lift a signed basis word for $x$ by the ordered images of its
generators under $T$, and denote the central coordinate of that
lifted word by $z_T(x)$. With this literal word convention,
\[
 f_T(x)=\Lambda^2T\sum_{i<j}x_i x_j e_{ij}-z_T(x).
\]
Equivalently, with
$A_{ij}=p(Te_i,Te_j)-\Lambda^2T\,p(e_i,e_j)$,
\[
 f_T(x)=-\sum_i\binom{x_i}{2}A_{ii}
        -\sum_{i<j}x_i x_j A_{ij}.
\]
The symmetry $A_{ij}=A_{ji}$ follows from the exterior-square
identity. Signed integer powers use the group inverse just given,
so this formula applies without a positivity restriction.
The lifted map is
$(x,z)\mapsto(Tx,\Lambda^2Tz-f_T(x))$; its inverse sends
$(y,w)$ to
$(T^{-1}y,\Lambda^2T^{-1}(w+f_T(T^{-1}y)))$.

For $k=(u_0,\ldots,u_4)\in K$, let $\mathcal B_i(u_i)$ be the
normalized prism over the $i$-th graph edge. Close its fibre
boundary by the ordered polygon with sides
$T_0u_0,-u_0,\ldots,T_4u_4,-u_4$, retaining the self terms
$p(u_i,-u_i)$. Its lifted central coordinate is $z_{\rm pr}(k)$.
The actual integral splitting is
\[
 (c,k)\longmapsto \operatorname{fibre}(c)+Z(k)
                         -\operatorname{fibre}(z_{\rm pr}(k)),
 \qquad H_2(V)\simeq\operatorname{coker}R\oplus K.
\]
The inverse projects a cycle to its primitive graph coordinate,
subtracts this exact prism/polygon section, and reads the remaining
fibre coinvariant. The normalized bar boundary formula and the
central group law prove both composites are the identity.
In these coordinates a fixed-normal torus at site $i$ is
\[
        (f_{T_i}(u),K_0^{\vee}J_i u),\qquad (T_i-I)u=0,
\]
where $J_i$ inserts $u$ in the $i$-th graph block. The displayed
$f_{T_i}(u)$ already includes the polygon correction. This specifies the integral
fixed-normal columns and preserves the order-two fibre class.

\item \emph{The ordinary boundary word and its Thom columns.}
The positive ordinary word is
$\theta_o=(\theta_0\cdots\theta_4)^{-1}$.
Set $P_i=T_{i+1}\cdots T_4$ and
\[
 J_o u=(-P_0u,\ldots,-P_4u),\qquad
 c_o(u)=-\sum_{i=0}^4 f_{T_i}(P_i u)\pmod{\operatorname{im}R}.
\]
Applying the actual inverse lifted maps in their word order gives
the same correction modulo $\operatorname{im}R$: the signed
central group law expresses the difference as a sum of columns of $R$. As the product monodromy is identity, this
map is additive in the fibre coinvariant; the signed word formula
proves the assertion for every integral $u$.

Let $v_{\rm cl}=-e\gamma$. The ordinary $m_{\rm ord}=1$ cap has boundary map
\[
 \begin{gathered}
 (y,\tau)\longmapsto(x=y-\tau v_{\rm cl},e^{2\pi i\tau}),\\
 (x,\tau)\longmapsto(y=x+\tau v_{\rm cl},\tau)
 \end{gathered}
\]
as its inverse. Its meridian is $\mu_o-v_{\rm cl}=\mu_o+e\gamma$.
The four fixed-normal Thom columns and the core column in the
above exterior splitting are
\[
 \begin{gathered}
 \left(c_o(u)-u\wedge v_{\rm cl},K_0^{\vee}J_o u\right)
       \quad(u=e_1,e_2,\delta,\gamma),\\
 \left(c_o(v_{\rm cl}),K_0^{\vee}J_o v_{\rm cl}\right).
 \end{gathered}
\]
They are induced by the displayed cap map, including the sign
of its meridian.

\item \emph{The exact finite presentation and its inverse coordinates.}
Use six free fibre coordinates and $17$ graph coordinates. The
matrix ${\mathsf B}_{\rm rel}$ has $23$ rows and $59$ columns, ordered as follows:
the $30$ columns of $R$, the ten cap fibre-kernel columns, the $14$
singular fixed-normal columns, and the five ordinary columns.
The preceding formulas give every entry from the matrices in \eqref{maps:eq:D.VC9}, the two pairs, the selected site and $e$.
The resulting exact Smith factors are
\[
\begin{array}{c|c|c|c|c}
 (a,b)&\text{selected }I_1&e&\operatorname{rank}{\mathsf B}_{\rm rel}
                           &\text{nonzero Smith factors}\\\hline
 (2,4)&1,2,3,4&-1,+1&23&1^{22},4\\
 (4,2)&1,2,3,4&-1,+1&23&1^{22},2.
\end{array}
\]
Every row of this table denotes its eight separate choices.
Write $\Delta_{\rm rel}={\mathsf U}_{\rm rel}{\mathsf B}_{\rm rel}{\mathsf V}_{\rm rel}$ for the integral Smith reduction, with ${\mathsf U}_{\rm rel},{\mathsf V}_{\rm rel}$
unimodular. The resulting quotient map and inverse use
${\mathsf U}_{\rm rel}^{-1}$ and ${\mathsf V}_{\rm rel}^{-1}$ as follows.
\[
 \begin{gathered}
 [z]\longmapsto ({\mathsf U}_{\rm rel}z)_{23}\pmod b,\\
 1\pmod b\longmapsto[{\mathsf U}_{\rm rel}^{-1}e_{23}].
 \end{gathered}
\]
The relation $b{\mathsf U}_{\rm rel}^{-1}e_{23}={\mathsf B}_{\rm rel}{\mathsf V}_{\rm rel} e_{23}$ and the unit Smith
columns prove both composites. Thus this calculation is an
integral presentation proof, including the marking and inverse,
of $\operatorname{coker}{\mathsf B}_{\rm rel}=\mathbb Z/b$.

Taking the ordinary word as the first graph generator and retaining
all triangular-cap fibre columns gives a $23\times70$ presentation
${\mathsf B}_{o,\rm rel}$. It has six zero ordinary fibre columns and five redundant
triangular-cap fibre columns. The free-group change from
$(\theta_o,\theta_0,\ldots,\theta_3)$ to
$(\theta_0,\ldots,\theta_4)$ has inverse word $\theta_4=(\theta_o\theta_0\cdots\theta_3)^{-1}$. On graph coordinates it is
\[
 (u_o,u_0,u_1,u_2,u_3)
     \longmapsto J_o u_o+(u_0,u_1,u_2,u_3,0).
\]
Write ${\mathsf J_{\Gamma}}$ for this ambient graph map. Its integral inverse is
\[
 {\mathsf J_{\Gamma}}^{-1}(v_0,\ldots,v_4)
       =(-v_4,v_0-P_0v_4,\ldots,v_3-P_3v_4).
\]
The two composites cancel term by term. Let $D_o,D_s$ denote the
boundary maps in the ordinary-first and singular-first graph bases.
The ordered product identity gives $D_s{\mathsf J_{\Gamma}}=D_o$.
Let $K_o,K_s$ be the matrices of the chosen
primitive kernel bases, with their coordinate inverses on those
kernels. The graph coordinate maps are therefore exactly
\[
 {\mathsf Q_{\Gamma}}=K_s^{-1}{\mathsf J_{\Gamma}}K_o,\qquad
 {\mathsf Q_{\Gamma}}^{-1}=K_o^{-1}{\mathsf J_{\Gamma}}^{-1}K_s.
\]
They are integral and inverse on the graph-kernel coordinates.
On the $23$ generators its additional central term is
$C_oE_oK_o$, where $C_o$ has columns $c_o(e_j)$
and $E_o$ selects the ordinary block in the ordinary-first basis. Hence the marked
row map is
${\mathsf M_{\Gamma}}=\left(\begin{smallmatrix}I_6&C_oE_oK_o\\0&{\mathsf Q_{\Gamma}}\end{smallmatrix}\right)$,
with inverse
$\left(\begin{smallmatrix}I_6&-C_oE_oK_o{\mathsf Q_{\Gamma}}^{-1}\\0&{\mathsf Q_{\Gamma}}^{-1}\end{smallmatrix}\right)$.
Both are integral. Applying the bar lift to these words and
including the zero and redundant fibre columns gives integral
relation matrices ${\mathsf C}_{\rm rel},{\mathsf D}_{\rm rel}$ with
${\mathsf B}_{\rm rel} {\mathsf C}_{\rm rel}={\mathsf M_{\Gamma}} {\mathsf B}_{o,\rm rel}$ and ${\mathsf B}_{o,\rm rel}{\mathsf D}_{\rm rel}={\mathsf M_{\Gamma}}^{-1}{\mathsf B}_{\rm rel}$. These identities make
the quotient maps induced by ${\mathsf M_{\Gamma}}$ and ${\mathsf M_{\Gamma}}^{-1}$ well-defined
and inverse. The $59$-column presentation above therefore computes
the same marked quotient in the singular-first graph basis.

\item \emph{The relative map and full fundamental group.}
Each cap $H_2$ specialization is onto, as its explicit section
shows. The only possible additional degree-two contribution
is the kernel of the relative component-disc map to $H_1(V)$.
After the vertical fibre relations, take the positive basis
$(\gamma,\mu_0,\ldots,\mu_4)$ of $H_1(V)$.
The five singular discs map to $\mu_i$, and the ordinary disc
maps to $e\gamma-\sum_i\mu_i$. Thus the full relative map is
\[
 (t_o,t_0,\ldots,t_4)
       \longmapsto(e t_o,-t_o+t_0,\ldots,-t_o+t_4).
\]
For $e=\pm1$ its integral inverse is
$t_o=y_\gamma/e$, $t_i=y_i+y_\gamma/e$.
Changing all five peripheral row and disc signs gives the
matrix with first column $(e,1,1,1,1,1)$;
that sign change is its own inverse. The relative kernel is zero.
Consequently $H_2(X)=\operatorname{coker}{\mathsf B}_{\rm rel}$.

The exterior fundamental group is the actual abelian-fibre
semidirect-product presentation with the five graph words and
the matrices $T_i$. The six primitive toric vanishing vectors
are the first columns of the four rank-one period matrices and
the two short columns of the selected triangular matrix.
Their nongamma coordinates have an integer right inverse
onto $\mathbb Z\langle e_1,e_2,\delta\rangle$ in every one
of the eight geometric choices. Thus their unsaturated normal
closure kills that entire abelian normal fibre subgroup.
All singular peripheral words bound the specified origin discs.
The last ordinary relation is then $\gamma^e=1$, with $e=\pm1$.
This proves $\pi_1(X)=1$ in the actual group presentation.

\item \emph{The remaining homology.}
All regular and boundary torus pieces have Euler characteristic $0$.
The four rank-one caps have Euler characteristic $0$ and the unique
triangular cap contributes $2$, so $\chi(X)=2$.
The manifold is closed and oriented. Since $H_1=0$ and $H_2$
is finite, universal coefficients and integral duality give
$H_5=0$, $H_4=H^2=0$ and
$\operatorname{Tor}H_3\simeq\operatorname{Ext}(H_2,\mathbb Z)
\simeq\mathbb Z/b$. More concretely, send the displayed
$H_2$ generator to $1/b$ in $\mathbb Q/\mathbb Z$, take its
Bockstein in $H^3$, and cap with the oriented fundamental class.
The inverse is the perfect torsion linking evaluation against
that generator. The Euler equation now makes the free rank of
$H_3$ zero. Together with $H_0=H_6=\mathbb Z$, these are
exactly all the groups in \eqref{maps:eq:D.VC7}.

\end{enumerate}
\end{proof}
\section{Invariant-Chern caps and height-integral pairs}
\subsection{Invariant-Chern additive caps}

\begin{definition}[Invariant-Chern cap data]\hypertarget{tag-Q203}{}\label{maps:def:invariant-chern-data}
{\normalfont\scriptsize tag Q203}\par
Let $K$ be the reduced central fibre of an original additive
Kodaira neighbourhood with its line bundle $M$, section $Q$
and original six-dimensional cap $U$. The fibre type is
$II,III,IV,IV^*,III^*,I_1^*$ or $I_3^*$.
Write its rational components as $C_i$, its primitive
multiplicity vector as $\mu=(\mu_i)$, and let $\sigma$
be the actual component permutation induced by the section $Q$.
The component degrees $c_i$ of $M$ satisfy
\begin{equation}
                  \sigma c=c,\qquad \sum_i\mu_i c_i=0.
                                                        \hypertarget{tag-01ZY}{}\label{maps:eq:D.IC1}
\end{equation}
The corrected line bundles of \ref{maps:src:D.2.50} give these
data at the listed additive fibres. Let the component orbits be
$O_1,\ldots,O_\rho$, with lengths $l_j$, common degrees $c_j$
and common multiplicities $\mu_j$ on $O_j$.
For $c\ne0$ put
\begin{equation}
                  g=\gcd_j(c_j),\qquad
                  h=\gcd_j(l_jc_j),
                                                        \hypertarget{tag-01ZZ}{}\label{maps:eq:D.IC2}
\end{equation}
where both gcds are positive. These are the \emph{invariant-Chern
cap data}.
\end{definition}

\begin{theorem}[Invariant-Chern additive caps]\hypertarget{tag-02LT}{}\label{maps:src:D.2.51}
{\normalfont\scriptsize tag 02LT}\par
The cap $U$ of Definition~\ref{maps:def:invariant-chern-data} has
the following integral cellular cone, with its scalar-circle and
regular-area classes:
\begin{equation}
\begin{gathered}
 \mathsf C_0=\mathbb Z,\quad \mathsf C_1=\mathbb Z^2,\quad
 \mathsf C_2=\mathbb Z^v\oplus\mathbb Zt,\quad
 \mathsf C_3=\mathbb Z^v\oplus\mathbb Z^v,\quad \mathsf C_4=\mathbb Z^v,\\
 d_2=\begin{pmatrix}c^t&0\\0&0\end{pmatrix},\qquad
 d_3=\begin{pmatrix}0&\sigma-I\\0&c^t\end{pmatrix},\qquad
 d_4=\begin{pmatrix}I-\sigma\\0\end{pmatrix}.
\end{gathered}                                             \hypertarget{tag-0200}{}\label{maps:eq:D.IC3}
\end{equation}
The degree-one generators are the scalar circle $\delta$
and a positive long circle $\gamma$. For $c\ne0$,
\begin{equation}
\begin{array}{c|ccccc}
i&0&1&2&3&4\\\hline
H_i(U;\mathbb Z)&\mathbb Z&
\mathbb Z\oplus\mathbb Z/g&
\mathbb Z^{\rho-1}\oplus\mathbb Z/h&
\mathbb Z^{2\rho-1}&\mathbb Z^\rho .
\end{array}                                               \hypertarget{tag-0201}{}\label{maps:eq:D.IC4}
\end{equation}
Moreover $\pi_1(U)=\mathbb Z\oplus\mathbb Z/g$.
When $c=0$, $\pi_1(U)=\mathbb Z^2$ and the homology groups are
\begin{equation}
 (H_0,H_1,H_2,H_3,H_4)=
 (\mathbb Z,\mathbb Z^2,\mathbb Z^{\rho+1},
                         \mathbb Z^{2\rho},\mathbb Z^\rho).
                                                        \hypertarget{tag-0202}{}\label{maps:eq:D.IC5}
\end{equation}
The actual regular-fibre torus $\delta\wedge\gamma$
maps to the class of $t$ in \eqref{maps:eq:D.IC3}, which generates
the cyclic torsion $\mathbb Z/h$ when
$c\ne0$. The elliptic tori crossed with $\delta$ map
to zero. The complex elliptic area maps to the literal
cycle $(\mu,0)$ in \eqref{maps:eq:D.IC3}.

Let $B$ be the elliptic monodromy, $q$ its marked
$Q$-cocycle, and $z=(B-I)^{-1}q$. Put
\begin{equation}
        q_Q=\gcd_j(l_j\mu_j)
             =\operatorname{ord}(\overline Q).
                                                        \hypertarget{tag-0203}{}\label{maps:eq:D.IC6}
\end{equation}
The actual logarithm $z$ has exact denominator $q_Q$.
For $c\ne0$, the full regular-fibre image in $H_2(U)$
is $\mathbb Z\oplus\mathbb Z/h$ with free cokernel
of rank $\rho-2$. For $c=0$ it is a primitive
$\mathbb Z^2$ with free cokernel of rank $\rho-1$.
In both cases this is already the full boundary image.
Consequently the full boundary kernel projects onto
all of $\operatorname{Fix}(A)$ in the Wang sequence.


In the sector $c\ne0$, $h=1$, the complete marked local kernels
are as follows. In the regular basis
$(e_1,e_2,\delta,\gamma)$ the actual fibre kernel is
\begin{equation}
 \ker\bigl(q_Q(1,0,z_2,0,-z_1,0):\mathbb Z^6\to\mathbb Z\bigr),
                                                        \hypertarget{tag-0204}{}\label{maps:eq:D.IC7}
\end{equation}
where the six coordinates are ordered as
$(12,13,14,23,24,34)$ in that basis. Both literal fixed-normal tori bound.
The cap kills every vertical fundamental-group
coordinate and retains the primitive $\gamma$.
This includes the following star degrees from Theorem~\ref{maps:src:D.2.50}:
\begin{equation}
 \begin{array}{c|c}
 I_1^*&(1,1,0,0,-1,0)\\
 I_3^*&(1,1,0,0,-1,0,0,0)
 \end{array}
                                                        \hypertarget{tag-0205}{}\label{maps:eq:D.IC8}
\end{equation}
for component classes $P=Q=2$ in $\mathbb Z/4$.
Their actual permutation has two paired outer orbits
and fixed inner components, and $g=h=1$.

For $h>1$, the torsion coordinate of a specified cellular
representative is given by \eqref{maps:eq:IC_local_15}. This
representative includes the line-frame contributions of the mixed
elliptic--long and fixed-normal tori. The image and cokernel formulas
apply to every invariant degree vector; the marked kernels
\eqref{maps:eq:D.IC7} apply in the sector $h=1$.

\end{theorem}

\begin{proof}
Set $A_0=\sigma-I$. Every orbit is ordered by the return in the
positive roof direction; the scalar circle has its complex orientation.
\begin{enumerate}
\item \emph{The actual cap and its central mapping-torus model.}

Let $N$ be a sufficiently small closed original Kodaira neighbourhood,
$K$ its reduced central fibre, and $M$ the given line bundle.
The central configurations in the statement are simply connected.
For II, normalization is a continuous bijection from a sphere to the
unibranch cuspidal component, and compactness makes it a homeomorphism.
For III and IV, respectively two and three spheres meet at one point.
For the remaining listed types, the rational components have a tree
incidence graph. Contract chosen connecting arcs on the components.
The resulting map is a homotopy equivalence
\begin{equation}
 K\longrightarrow \bigvee_{i=1}^{v}S_i^2.
 \label{maps:eq:IC_local_1}
\end{equation}
Its inverse maps each sphere to its normalized component with the
chosen connecting arcs; the arc contractions give both homotopies.
Each sphere is oriented by its complex orientation.

Take the neighbourhood retraction for the original filling in
Theorem~\ref{maps:src:D.2.31}. We describe its lift and the return maps. In a Hermitian norm, write the original nonzero line arrow
as its unit-circle map $\phi$ and its positive roof length $\varrho(p)$.
For a positive density $\nu_p(s)$ on that interval, chosen with the
matching endpoint jets, the roof coordinate is
\begin{equation}
 \tau_p(s)=
 \frac{\int_0^s\nu_p(u)\,du}{\int_0^{\varrho(p)}\nu_p(u)\,du}.
 \label{maps:eq:IC_local_2}
\end{equation}
Its inverse is the unique increasing inverse function. The original
line vector is recovered from its unit vector and its original norm
at this inverse roof coordinate. Thus every unit in the actual arrow
is retained in the circle map or in the roof; the maps are defined on the given analytic germ.

Let $X$ be the unit circle bundle over $N$, and $E$ its restriction
over $K$. Lift the central retraction by connection parallel transport.
Reverse parallel transport is the inverse on each path. We have an
inclusion $i:E\to X$, a retraction $r_{\rm ret}:X\to E$, and a homotopy
$R_s:\operatorname{id}_X\simeq i r_{\rm ret}$, fixed on $E$.
The actual arrow preserves $E$; denote its restriction by $\phi_0$.
The homotopy
\begin{equation}
 H_s=r_{\rm ret}\phi R_s,\qquad H_0=r_{\rm ret}\phi,\quad H_1=\phi_0 r_{\rm ret}
 \label{maps:eq:IC_local_3}
\end{equation}
gives a map of mapping tori with a specified homotopy inverse to inclusion.
Use the convention $(x,1)\sim(\phi x,0)$. The map sends $[x,t]$ to
\begin{equation}
 \begin{cases}
 [r_{\rm ret}(x),2t],&0\leq t\leq\tfrac12,\\
 [H_{2-2t}(x),0],&\tfrac12\leq t\leq1.
 \end{cases}
 \label{maps:eq:IC_local_4}
\end{equation}
At the midpoint both values are $[\phi_0 r_{\rm ret}(x),0]$; at the endpoint
the value is $[r_{\rm ret}\phi(x),0]$, as required by the source identification.

Here is also a concrete homotopy for the inverse composites. On $E$,
\eqref{maps:eq:IC_local_4} is traversal of the original roof followed by a constant path;
compressing that constant subinterval gives the identity. On $X$,
use $R_s(x)$ on the traversed roof and join $\phi R_s(x)$ to
$R_s\phi(x)$ by the path
\begin{equation}
 C_s(x,u)=
 \begin{cases}
 R_{2su}\phi R_s(x),&0\leq u\leq\tfrac12,\\
 R_s\phi R_{2s(1-u)}(x),&\tfrac12\leq u\leq1.
 \end{cases}
 \label{maps:eq:IC_local_5}
\end{equation}
The two pieces agree. At $s=0$ this is a constant endpoint path.
At $s=1$ its first piece is constant, since $\phi i r_{\rm ret}(x)\in E$;
its second piece is $iH_{2-2u}(x)$. Reparametrizing the constant
pieces gives exactly inclusion followed by \eqref{maps:eq:IC_local_4}. This proves both
homotopy-inverse identities with the displayed homotopies.

It remains to determine the circle-bundle map over the wedge \eqref{maps:eq:IC_local_1}.
Choose reference Möbius maps carrying each marked node to its actual
permuted node on the target component. After composing with the inverse
reference map, the residual automorphism preserving
two marked nodes has the form $z\mapsto a_{\rm aff}z$; preserving just the
marked point at infinity it has the form $z\mapsto a_{\rm aff}z+b_{\rm aff}$.
The paths
\begin{equation}
 z\mapsto e^{s\log a_{\rm aff}}z,\qquad
 z\mapsto e^{s\log a_{\rm aff}}z+s b_{\rm aff}
 \label{maps:eq:IC_local_6}
\end{equation}
and their inverse affine maps remove these residual maps. A residual
map fixing three marked points is the identity. The reference maps
retain the actual permutation of the nodes, including the three nodes
on the fixed central component of IV*. After the arc contractions in
\eqref{maps:eq:IC_local_1}, their component two-cell degrees are precisely $\sigma$, and
their component-times-scalar three-cell degrees are also $\sigma$.
Thus their cellular map has exactly the degrees used in \eqref{maps:eq:IC_local_7}--\eqref{maps:eq:IC_local_8}.
Line frames are joined from the root of the component tree outward;
their inverse changes are the reciprocal nonzero frame functions.

Two resulting unit circle-bundle maps covering the same component map
differ by a function $u_{\rm gau}:K\to S^1$. Since \eqref{maps:eq:IC_local_1} has $H^1(K;\mathbb Z)=0$,
the covering $\mathbb R\to S^1$ gives a global lift $u_{\rm gau}=e^{2\pi i f_{\rm gau}}$.
The explicit gauge homotopy is multiplication by $e^{2\pi i s f_{\rm gau}}$,
with inverse multiplication by $e^{-2\pi i s f_{\rm gau}}$.
Consequently the actual map is homotopic, through principal-circle
maps, to the component permutation on identical degree-$c_i$ clutch
models. Invariance of $c$ is exactly what makes these clutch models
agree around every orbit. All gauge homotopies have specified inverses.

The displayed maps give a homotopy equivalence from the cap to
the mapping torus of the circle bundle over \eqref{maps:eq:IC_local_1},
with the stated marked cell maps and both inverse homotopies.
The component clutch degree is the given integer $c_i$ at each component.

\item \emph{The integral cells, signs, and marked cycles.}

The circle bundle over a sphere of degree $c_i$ has one cell in
each dimension 0,1,2,3, with degree-two boundary $c_i\delta$.
This can be read directly by gluing the two disc trivializations
with transition $e^{2\pi i c_i\theta}$: the component disc boundary
traverses the scalar circle $c_i$ times. Its three-cell is the
component times that scalar circle and has zero cellular boundary.
Glue these models along their common scalar fibre over the wedge point.
The circle-bundle chain groups are therefore
\begin{equation}
 E_0=\mathbb Z,\quad E_1=\mathbb Z\delta,\quad
 E_2=\mathbb Z^v,\quad E_3=\mathbb Z^v,\qquad
 \partial_2=c^t .
 \label{maps:eq:IC_local_7}
\end{equation}
The return acts as the identity on degrees 0 and 1, and as $\sigma$
on both component groups.

Attach the positive roof prisms, with boundary target minus source.
Their exact cone is
\begin{equation}
\begin{aligned}
 \mathsf C_0&=\mathbb Z,& \mathsf C_1&=\mathbb Z\delta\oplus\mathbb Z\gamma,\\
 \mathsf C_2&=\mathbb Z^v\oplus\mathbb Zt,&
 \mathsf C_3&=\mathbb Z^v\oplus\mathbb Z^v,& \mathsf C_4&=\mathbb Z^v,\\
 d_2(x,t)&=(c^t x,0),&
 d_3(u,y)&=(A_0y,c^ty),&
 d_4(z)&=(-A_0z,0).
\end{aligned}
 \label{maps:eq:IC_local_8}
\end{equation}
The equalities $c^t\sigma=c^t$ imply $d_2d_3=d_3d_4=0$.
The inverse cellular map reads the oriented component cells and roof
prisms. A changed line frame contributes its literal prism coefficient;
it is included in $t$.

The cell $t$ is the actual scalar circle times the positive long
circle. Hence the regular torus $\delta\wedge\gamma$ maps to $(0,1)$.
The elliptic area specializes to the complex cycle
\begin{equation}
 (\mu,0).
 \label{maps:eq:IC_local_9}
\end{equation}
To see the zero last coordinate, place that fibre at roof coordinate
zero before taking its specialization. The specialization stays in $E$;
its component projection is the divisor of the base parameter,
$\sum_i\mu_i C_i$. By \eqref{maps:eq:IC_local_7}, $H_2(E)$ injects into
$H_2(K)=\mathbb Z^v$ with image $\ker c^t$, so this lift is unique.
The condition $c^t\mu=0$ makes \eqref{maps:eq:IC_local_9} a cycle.

An elliptic loop crossed with the scalar circle also lies in $E$.
Its projection to $H_2(K)$ is zero, because that projection factors
through the elliptic loop. The same integral injection from \eqref{maps:eq:IC_local_7}
proves that both elliptic-scalar tori vanish. This argument applies
for every $g$, including $g>1$.

\item \emph{Orbit maps and the augmentation kernel.}

Write $\varepsilon_i$ for the standard component-coordinate column;
its degree-two copy is the oriented cell over $C_i$. Write each orbit as
\begin{equation}
 O_j=(i_{j,0},\ldots,i_{j,l_j-1}),\qquad
 \sigma \varepsilon_{i_{j,k}}=\varepsilon_{i_{j,k+1}},
 \label{maps:eq:IC_local_10}
\end{equation}
where indices are cyclic. Define the following integral maps:

\begin{itemize}
\item $R:\mathbb Z^v\to\mathbb Z^\rho$ takes orbit sums.
\item $s:\mathbb Z^\rho\to\mathbb Z^v$ puts the $j$-th coordinate at
  $i_{j,0}$ and puts zero at the other orbit positions.
\item $\iota:\mathbb Z^\rho\to\mathbb Z^v$ puts the $j$-th coordinate
  at every position of that orbit.
\end{itemize}

Then $Rs=I$, $\ker A_0=\operatorname{im}\iota$, and $RA_0=0$.
For $Rz=0$, put
\begin{equation}
 (\mathcal K z)_{i_{j,k}}=-\sum_{r=0}^k z_{i_{j,r}}
 \quad(0\leq k<l_j-1),\qquad
 (\mathcal K z)_{i_{j,l_j-1}}=0.
 \label{maps:eq:IC_local_11}
\end{equation}
Since $(A_0y)_{i_{j,k}}=y_{i_{j,k-1}}-y_{i_{j,k}}$,
direct subtraction gives
\begin{equation}
 A_0\mathcal K z=z,\qquad
 y-\mathcal K A_0y=\iota\bigl(y_{i_{j,l_j-1}}\bigr)_{j=1}^{\rho}.
 \label{maps:eq:IC_local_12}
\end{equation}
In particular $\operatorname{im}A_0=\ker R$ integrally, with the displayed right inverse.

Let $d=(c_1,\ldots,c_\rho)$ denote the orbit degree vector, so
$c=\iota d$, $c^tx=d^tRx$, and
\begin{equation}
 c^t\iota a=\sum_j l_jc_j a_j.
 \label{maps:eq:IC_local_13}
\end{equation}
For $c\ne0$, choose a Bézout vector $a_{\rm orb}$ with
$\sum_j l_jc_j a_{{\rm orb},j}=h$, and set $\eta=\iota a_{\rm orb}$. Then
\begin{equation}
 A_0\eta=0,\qquad c^t\eta=h,\qquad d_3(0,\eta)=(0,h).
 \label{maps:eq:IC_local_14}
\end{equation}
Both $g$ and $h$ are positive and $g\mid h$.

\item \emph{$H_2$ and its quotient maps.}

For $c\ne0$, define
\begin{equation}
 \Phi[(x,t)] =
 \left(Rx,\,
 \left[t-c^t\mathcal K(x-sRx)\right]_h\right)
 \in \ker d^t\oplus\mathbb Z/h .
 \label{maps:eq:IC_local_15}
\end{equation}
Its domain consists of the cycles $c^tx=0$.
For a boundary $(A_0y,c^ty)$, the first coordinate vanishes and,
by \eqref{maps:eq:IC_local_12}, the second coordinate before reduction is
\begin{equation}
 c^t(y-\mathcal K A_0y)
   =\sum_j l_jc_j\,y_{i_{j,l_j-1}}\in h\mathbb Z.
 \label{maps:eq:IC_local_16}
\end{equation}
Thus \eqref{maps:eq:IC_local_15} is defined on homology.

Its inverse is
\begin{equation}
 \Psi(u,[b]_h)=[(su,b)].
 \label{maps:eq:IC_local_17}
\end{equation}
This is a cycle when $d^tu=0$; changing $b$ by $h$ changes it
by the boundary \eqref{maps:eq:IC_local_14}. Clearly $\Phi\Psi=I$.
For the other composite choose any integer representative $b_{\rm hom}$ of
the second coordinate in \eqref{maps:eq:IC_local_15}, put $z_{\rm aug}=x-sRx$, and write
\begin{equation}
 k_{\rm hom}=\frac{t-b_{\rm hom}-c^t\mathcal K z_{\rm aug}}{h}\in\mathbb Z.
 \label{maps:eq:IC_local_18}
\end{equation}
The error has the exact integral boundary preimage
\begin{equation}
 (x,t)-(sRx,b_{\rm hom})=d_3(0,\mathcal K z_{\rm aug}+k_{\rm hom}\eta).
 \label{maps:eq:IC_local_19}
\end{equation}
This proves $\Psi\Phi=I$. It also supplies a preimage for every
zero class: take $Rx=0$ and $b_{\rm hom}=0$ in \eqref{maps:eq:IC_local_18}--\eqref{maps:eq:IC_local_19}.

Consequently
\begin{equation}
 H_2(U)=\ker d^t\oplus\mathbb Z/h
       \cong\mathbb Z^{\rho-1}\oplus\mathbb Z/h .
 \label{maps:eq:IC_local_20}
\end{equation}
The element $t=(0,1)$ maps to $(0,[1]_h)$, so it has exact
order $h$ and generates all torsion. Formula \eqref{maps:eq:IC_local_16} proves the
lower bound on its order; \eqref{maps:eq:IC_local_14} proves the upper bound.
The maps identify the orbit part and the scalar-long subgroup.

For $c=0$, put $g=h=0$ in this branch. The same construction is
read with an integer second coordinate and without division by $h$:
\begin{equation}
 [(x,t)]\longmapsto (Rx,t),\qquad
 (u,b)\longmapsto[(su,b)].
 \label{maps:eq:IC_local_21}
\end{equation}
The error is $d_3(0,\mathcal K(x-sRx))$. Thus
$H_2(U)=\mathbb Z^\rho\oplus\mathbb Zt$.

Under the splitting \eqref{maps:eq:IC_local_15}, the area \eqref{maps:eq:IC_local_9} has coordinates
\begin{equation}
 \left(R\mu,\,
   \left[-\sum_j c_j\mu_j\,\frac{l_j(l_j-1)}2\right]_h\right).
 \label{maps:eq:IC_local_22}
\end{equation}
Indeed \eqref{maps:eq:IC_local_11} has values $(l_j-1-k)\mu_j$ on that orbit.
This is an exact evaluation of the displayed splitting. The
geometric class itself remains the literal cycle $(\mu,0)$.


\item \emph{The two extension maps.}
On $O_j$, write $\varepsilon_{j,i}=\varepsilon_{i_{j,i}}$.
For $c\ne0$, the orbit relations are
\begin{equation}\hypertarget{tag-0206}{}\label{maps:eq:D.IC9}
 \varepsilon_{j,i+1}-\varepsilon_{j,i}+c_jt=0 .
\end{equation}
Elimination gives $\varepsilon_{j,i}=\varepsilon_{j,0}-i c_jt$ and $l_jc_jt=0$ on each orbit.
The partial-sum formula gives exactly the same scalar coordinate,
since
\begin{equation}
 c^t\mathcal K(x-sRx)=\sum_jc_j\sum_{i=0}^{l_j-1}i x_{j,i}.
\end{equation}
The projection onto component orbits has exact sequence
\begin{equation}
0\longrightarrow\mathbb Z/h
\xrightarrow{[a]\mapsto[(0,a)]}H_2(U)
\xrightarrow{[(x,t)]\mapsto Rx}\ker d^t
\longrightarrow0 ,
\end{equation}
with section $u\mapsto[(su,0)]$.

For the mapping torus of the circle bundle $E$, put
$D=h/g$ and choose $\beta$ with $c^t\beta=g$.
Its Wang sequence has left term
$\mathcal C_{\rm Wang}=\ker c^t/A_0(\ker c^t)$.
The following maps are inverse:
\begin{equation}
 [x]\longmapsto
 \left(Rx,\left[-c^t\mathcal K(x-sRx)/g\right]_D\right),\qquad
 (u,[a]_D)\longmapsto[su-aA_0\beta].
\end{equation}
The division is integral. Replacing $x$ by $x+A_0y$ with $c^ty=0$
changes the residue by a multiple of $D$, since $y-\mathcal K A_0y$ is
orbit-constant. Replacing $a$ by $a+D$ changes its image by
$A_0(\eta-D\beta)$, whose preimage has zero Chern evaluation.
For $z_{\rm aug}=x-sRx$ and a representative $a_{\rm Wang}$ of the displayed residue,
$k_{\rm Wang}=(c^t\mathcal K z_{\rm aug}+a_{\rm Wang}g)/h$ is integral and
\begin{equation}
\begin{gathered}
 x-(sRx-a_{\rm Wang}A_0\beta)
   =A_0(\mathcal K z_{\rm aug}+a_{\rm Wang}\beta-k_{\rm Wang}\eta),\\
 c^t(\mathcal K z_{\rm aug}+a_{\rm Wang}\beta-k_{\rm Wang}\eta)=0 .
\end{gathered}
\end{equation}
These are the two inverse identities.
In these coordinates the whole sequence is
\begin{equation}
0\longrightarrow\ker d^t\oplus\mathbb Z/D
\xrightarrow{(u,[a])\mapsto(u,[ga]_h)}
\ker d^t\oplus\mathbb Z/h
\xrightarrow{(u,[b])\mapsto[b]_g}\mathbb Z/g
\longrightarrow0 .
\end{equation}
Its torsion generator has the sign
$[(-A_0\beta,0)]=[(0,g)]$, by $d_3(0,\beta)=(A_0\beta,g)$.
A section exists exactly when $\gcd(g,D)=1$:
the scalar image of $1$ must satisfy $b_{\rm sec}\equiv1\pmod g$ and
$D\mid b_{\rm sec}$. For coprime $g,D$, choose $a_{\rm split}D+b_{\rm split}g=1$. The section is
$[n]_g\mapsto(0,[a_{\rm split}Dn]_h)$; the full splitting and inverse are
\begin{equation}
\begin{aligned}
 ((u,[v]_D),[n]_g)&\longmapsto(u,[gv+a_{\rm split}Dn]_h),\\
 (u,[t]_h)&\longmapsto
 \left((u,[(t-a_{\rm split}Dn)/g]_D),[n]_g\right),\quad n\equiv t\pmod g .
\end{aligned}
\end{equation}
All changes of representatives give the indicated moduli.
For $c=0$, the Wang sequence is
$0\to\mathbb Z^\rho\to\mathbb Z^\rho\oplus\mathbb Zt
\to\mathbb Z\to0$, with section given by the infinite class $t$.
For example, at type $III$ with $Q=0$ and $M=\mathcal O(C_1)$,
$c=(2,-2)$ and $g=h=2$, so the Wang sequence splits.
For $I_1^*$ in the affine $D_5$ order with near leaves $0,1$
and far leaves $4,5$, the invariant divisor $(0,0,2,4,3,3)$
under $(0\,1)(4\,5)$ has degrees $(2,2,0,0,-2,-2)$.
Its canonical line arrow, multiplied by a constant contraction,
has $g=2,h=4$ and gives a nonsplit Wang sequence.

\item \emph{The other groups and the fundamental group.}

From \eqref{maps:eq:IC_local_8}, the degree-one quotient and its inverse are
\begin{equation}
 [a\delta+b\gamma]\longmapsto([a]_g,b),\qquad
 ([a]_g,b)\longmapsto[a\delta+b\gamma].
 \label{maps:eq:IC_local_23}
\end{equation}
A Bézout vector $b_{i,{\rm deg}}$ satisfying $\sum_i b_{i,{\rm deg}} c_i=g$ gives the
explicit boundary $g\delta=d_2(\sum_i b_{i,{\rm deg}} \varepsilon_i,0)$.
For $c=0$ retain the integer scalar coordinate instead.

The two-cell attaching loops in the circle bundle are
$\delta^{c_i}$. The roof prism over the scalar circle attaches
the commutator $[\delta,\gamma]$. All other added cells have
dimension at least three. Therefore van Kampen gives the actual
presentation
\begin{equation}
 \pi_1(U)=
 \langle\delta,\gamma\mid
    \delta^{c_i}=1\ (1\leq i\leq v),\ [\delta,\gamma]=1\rangle.
 \label{maps:eq:IC_local_24}
\end{equation}
Its isomorphism with $\mathbb Z/g\oplus\mathbb Z$ sends a word to
its scalar exponent modulo $g$ and its positive roof exponent.
The inverse sends the two coordinates to those powers of the
literal loops. Bézout reduces exactly the subgroup generated by
the $c_i$ to $g\mathbb Z$, proving both composites.
For $c=0$ this is $\mathbb Z^2$.

A degree-three cycle $(u,y)$ has
$A_0y=0$ and $c^ty=0$. Thus $y=\iota a$ with
$\sum_j l_jc_j a_j=0$. The exact quotient map is
\begin{equation}
 [(u,y)]\longmapsto
 \left(Ru,\,(y_{i_{j,l_j-1}})_j\right)
 \in\mathbb Z^\rho\oplus\ker(l_jc_j)_j .
 \label{maps:eq:IC_local_25}
\end{equation}
Its inverse is $(b,a)\mapsto[(sb,\iota a)]$.
The error in $u$ equals $d_4(-\mathcal K(u-sRu))$; the $y$ error is zero.
Hence $H_3(U)=\mathbb Z^{2\rho-1}$ for $c\ne0$, and
$\mathbb Z^{2\rho}$ for $c=0$.

Finally
\begin{equation}
 H_4(U)=\ker A_0=\iota\mathbb Z^\rho.
 \label{maps:eq:IC_local_26}
\end{equation}
The inverse reads one coefficient from each orbit; its basis is
the literal orbit sums of the oriented three-cell prisms.
$H_0(U)=\mathbb Z$, and all groups above degree four vanish
because of the homotopy equivalence and the four-dimensional cone.
This gives \eqref{maps:eq:D.IC4} and \eqref{maps:eq:D.IC5}.

\item \emph{Component order, weighted orbit gcd, and the marked logarithm.}

The marked section cocycle reads the component class by
\begin{equation}
 \overline Q=[q]\in\mathbb Z^2/(B-I)\mathbb Z^2 .
 \label{maps:eq:IC_local_27}
\end{equation}
The inverse change of a logarithm adds an integral elliptic period.
Since $B-I$ is invertible over $\mathbb Q$, for
$z=(B-I)^{-1}q$ we have, for every positive integer $k$,
\begin{equation}
 kz\in\mathbb Z^2
 \ \Longleftrightarrow\
 kq\in(B-I)\mathbb Z^2 .
 \label{maps:eq:IC_local_28}
\end{equation}
Thus the exact denominator of the actual logarithm is the exact
order $q_Q$ of \eqref{maps:eq:IC_local_27}. This uses the actual marked component map,
with inverse given by an integral change of logarithm.

Here are the weighted diagrams and the component action giving
the second formula for $q_Q$.

\begin{itemize}
\item Type $II$ has one component of multiplicity one and the identity action.
\item Type $III$ has two multiplicity-one components; the nontrivial
component class interchanges them.
\item Type $IV$ has three multiplicity-one components; a generator cycles them.
\item Type $IV^*$ has a central multiplicity-three vertex and three
arms with multiplicities two and one. A generator cycles the three arms.
\item Type $III^*$ has chain multiplicities $(1,2,3,4,3,2,1)$ and
one extra multiplicity-two vertex at the centre. The nontrivial
component class reflects the chain and fixes the extra vertex.
\item For $I_n^*$, $n=1,3$, the four outer multiplicities are one and
the $n+1$ inner multiplicities are two. A component generator cycles
the four outer vertices and reflects the inner chain; its square
pairs the outer vertices and fixes the inner chain.
\end{itemize}

The multiplicity-one components carry the component-group torsor.
Translation by $\overline Q$ has an orbit of length $q_Q$ there.
On III* the fixed vertices of the nontrivial action have even
multiplicity. On IV* its fixed central vertex has multiplicity 3.
On an odd star, order four pairs the inner multiplicity-2 vertices,
so their weighted orbit length is 4; order two fixes them, giving
weighted length 2. In each listed action every $l_j\mu_j$ is
divisible by $q_Q$, and at least one multiplicity-one orbit has
weighted length exactly $q_Q$. Hence
\begin{equation}
 \gcd_j(l_j\mu_j)=q_Q .
 \label{maps:eq:IC_local_29}
\end{equation}
These checks also give the actual positive inverse coefficient:
choose one such multiplicity-one orbit and read its coordinate
after division by $q_Q$.

The diagram descriptions are unchanged by following the inverse
cyclic direction. At the cellular-coordinate level the exact
conversion is $i_{j,k}\mapsto i_{j,-k}$, with inverse itself;
it conjugates the permutation to its inverse and preserves $c,\mu$
and the scalar-long coordinate. We still use the actual forward
direction \eqref{maps:eq:IC_local_10} for all partial sums and all marked period maps.

Put
\begin{equation}
 w=R\mu/q_Q\in\mathbb Z^\rho.
 \label{maps:eq:IC_local_30}
\end{equation}
It is primitive by \eqref{maps:eq:IC_local_29}, and $d^tw=0$. For $c\ne0$, necessarily
$\rho\geq2$: a single orbit would give the positive multiple
$(\sum_i\mu_i)c_1=0$, hence $c=0$.

\item \emph{Integral bases and their inverses.}

For a nonzero row $r=(r_1,\ldots,r_k)$, first interchange the first
column and the least-indexed column with nonzero entry. The inverse is
the same transposition. Multiply the first column by the sign of its
entry; that diagonal sign map is also its own inverse. The first entry
is now positive. For $k=1$ these operations already give the required
matrix and its inverse. For $k\geq2$, successively combine the positive
first entry $a_{\rm col}$ and each next entry $b_{\rm col}$. Choose integers
$s_{\rm Bez},t_{\rm Bez}$ with $s_{\rm Bez}a_{\rm col}+t_{\rm Bez}b_{\rm col}=d_{\rm Bez}=\gcd(a_{\rm col},b_{\rm col})>0$, and multiply those two columns
by
\begin{equation}
\begin{gathered}
 E_{\rm Bez}=\begin{pmatrix}s_{\rm Bez}&-b_{\rm col}/d_{\rm Bez}\\t_{\rm Bez}&a_{\rm col}/d_{\rm Bez}\end{pmatrix},\\
 E_{\rm Bez}^{-1}=\begin{pmatrix}a_{\rm col}/d_{\rm Bez}&b_{\rm col}/d_{\rm Bez}\\-t_{\rm Bez}&s_{\rm Bez}\end{pmatrix}.
\end{gathered}
 \label{maps:eq:IC_local_31}
\end{equation}
Both matrices are integral and have determinant 1. The product
$V_{\rm alg}$ and the product of these inverse matrices satisfy
\begin{equation}
 rV_{\rm alg}=(\gcd(r),0,\ldots,0),\qquad V_{\rm alg}V_{\rm alg}^{-1}=V_{\rm alg}^{-1}V_{\rm alg}=I .
 \label{maps:eq:IC_local_32}
\end{equation}
The initial transposition and sign map are included in $V_{\rm alg}$ and $V_{\rm alg}^{-1}$.
The last columns of $V_{\rm alg}$, with the last coordinates of $V_{\rm alg}^{-1}$
as inverse, give an integral basis and inverse for $\ker r$.
For the zero row take $V_{\rm alg}=I$ and retain all columns.

Apply this first to $d^t$. Denote the resulting kernel basis by
$J:\mathbb Z^{k_{\rm ker}}\to\ker d^t$, with inverse $J^{-1}$ on that kernel;
$k_{\rm ker}=\rho-1$ for $c\ne0$, and $k_{\rm ker}=\rho$ for $c=0$.
The vector $b_{\rm prim}=J^{-1}w$ is primitive. Indeed a common divisor of
its coordinates would divide every coordinate of $w=Jb_{\rm prim}$.
Apply \eqref{maps:eq:IC_local_31}--\eqref{maps:eq:IC_local_32} to $b_{\rm prim}^t$, obtaining $b_{\rm prim}^tV_{\rm prim}=(1,0,\ldots,0)$.
Then
\begin{equation}
 L=J\,V_{\rm prim}^{-t},\qquad L^{-1}=V_{\rm prim}^tJ^{-1}
 \label{maps:eq:IC_local_33}
\end{equation}
are inverse integral bases of $\ker d^t$, and the first column
of $L$ is exactly $w$.
Dropping the first coordinate of $L^{-1}$, with inverse given by
the remaining columns of $L$, is the exact quotient
\begin{equation}
 \ker d^t/\mathbb Zw\cong\mathbb Z^{k_{\rm ker}-1}.
 \label{maps:eq:IC_local_34}
\end{equation}
These maps give the quotient and its inverse over the integers.

\item \emph{The long cover and the component image.}

Let $q_0=q_Q$. The $q_0$-fold long cover has the $q_0$-th power
of the actual return. Its underlying section lies in the identity
component, with integral logarithm $q_0z$ in the actual period marking.
Use the section logarithm and lift its translation to the line bundle.
Explicitly, the narrow section has a logarithm in the smooth
identity group over a sufficiently small base disc. At the additive
central fibre this identity group is additive; the local exponential
has invertible derivative. Shrinking the disc around the compact
path of logarithms gives the translations
\begin{equation}
 p\longmapsto p+t\mathcal L,\qquad
 p\longmapsto p-t\mathcal L
 \quad(0\leq t\leq1),
 \label{maps:eq:IC_local_35}
\end{equation}
as mutually inverse maps. The translations extend on each marked
rational normalization by the affine maps \eqref{maps:eq:IC_local_6}.
The line lift is parallel transport along this specified path;
reverse transport is its inverse. Its endpoint is compared with
the actual line arrow by their nonzero unit ratio. The latter is
kept, with its reciprocal inverse; it contributes only scalar
coordinates to a marked torus.

On the regular elliptic coordinates the resulting covering map is
\begin{equation}
 (x_1,x_2,s)\longmapsto(x_1+q_0z_1s,\ x_2+q_0z_2s,\ s),
 \label{maps:eq:IC_local_36}
\end{equation}
with inverse subtraction. On one original long period this reads
$x\mapsto x+zs$. The actual logarithm contains its analytic unit
part as well; that part is carried through \eqref{maps:eq:IC_local_35} and has no extra
integral period. Thus \eqref{maps:eq:IC_local_36} is the induced marked integral-period
map of an actual germ map, with the displayed inverse.

Forget the scalar coordinate and take the component part of a
two-cycle. For the long cover the area specializes to $\mu$,
and the two mixed elliptic-long tori have component parts
\begin{equation}
 q_0z_2\,\mu,\qquad -q_0z_1\,\mu .
 \label{maps:eq:IC_local_37}
\end{equation}
The signs follow by expanding
\begin{equation}
 (dx_1+q_0z_1ds)\wedge(dx_2+q_0z_2ds).
\end{equation}
Under covering projection a component basis vector maps to its
orbit coordinate, namely $R$. A mixed torus with covering long
generator maps to $q_0$ times the original mixed torus.
Consequently, if $f$ is the component projection of the original
fibre map, naturality gives
\begin{equation}
 \begin{aligned}
 f(e_1\wedge e_2)&=q_0w,\\
 q_0 f(e_1\wedge\gamma)&=q_0z_2\,R\mu=q_0^2z_2w,\\
 q_0 f(e_2\wedge\gamma)&=-q_0z_1\,R\mu=-q_0^2z_1w.
 \end{aligned}
 \label{maps:eq:IC_local_38}
\end{equation}
The target here is the free subgroup $\ker d^t\subset\mathbb Z^\rho$.
Cancellation of $q_0$ in this free abelian group gives exactly
\begin{equation}
 f(y)=q_0(y_{12}+z_2y_{14}-z_1y_{24})\,w.
 \label{maps:eq:IC_local_39}
\end{equation}
The coordinates are ordered as
$(12,13,14,23,24,34)$ for $(e_1,e_2,\delta,\gamma)$.
The scalar-long coordinate is read separately by \eqref{maps:eq:IC_local_15}.

The row of three nonzero possible coefficients is
\begin{equation}
 (q_0,\ q_0z_2,\ -q_0z_1).
 \label{maps:eq:IC_local_40}
\end{equation}
It is integral and primitive. A common divisor would contradict
the least denominator assertion \eqref{maps:eq:IC_local_28}: if a prime divided all
three entries, division of $q_0$ by that prime would still clear
both denominators. Formula \eqref{maps:eq:IC_local_31} supplies a Bézout section for \eqref{maps:eq:IC_local_40}.
Thus the component image of the entire regular fibre is exactly
$\mathbb Zw$.

\item \emph{Fibre and boundary images.}

First let $c\ne0$. The scalar-long torus maps to the entire
torsion factor by \eqref{maps:eq:IC_local_14}--\eqref{maps:eq:IC_local_20}, and \eqref{maps:eq:IC_local_39}--\eqref{maps:eq:IC_local_40} give precisely the
primitive free direction $w$. Choose the three integer coefficients
in a Bézout section of \eqref{maps:eq:IC_local_40}, and form that literal combination $Y_{\rm fib}$
of the actual area and mixed fibre tori. Cellularize its image using
the component and roof maps above. Let $\theta\in\mathbb Z/h$ be its
second coordinate under the explicit map \eqref{maps:eq:IC_local_15}.
Then
\begin{equation}
 Y_{\rm fib}-\theta(\delta\wedge\gamma)\longmapsto
        \Psi(w,0).
 \label{maps:eq:IC_local_41}
\end{equation}
The expression is independent of the representative of $\theta$
at the level of its image, by the boundary \eqref{maps:eq:IC_local_14}.
This gives a section onto both the free and the torsion directions,
read in the cellular coordinates of the specified torus.

Hence the full image is exactly
\begin{equation}
 \Psi(\mathbb Zw\oplus\mathbb Z/h)
       \cong\mathbb Z\oplus\mathbb Z/h .
 \label{maps:eq:IC_local_42}
\end{equation}
Its cokernel is \eqref{maps:eq:IC_local_34}, of rank $\rho-2$.
The cokernel quotient reads the last coordinates of $L^{-1}Rx$.
Its inverse inserts the corresponding last columns of $L$ by
$s$, and sets $t=0$. Formula \eqref{maps:eq:IC_local_19} and \eqref{maps:eq:IC_local_41} show directly that
both composites are identities in that quotient.

For $c=0$, use the integer second coordinate of \eqref{maps:eq:IC_local_21}.
In \eqref{maps:eq:IC_local_41} subtract its actual integer $\theta$.
The image is
\begin{equation}
 s(\mathbb Zw)\oplus\mathbb Zt\cong\mathbb Z^2
 \label{maps:eq:IC_local_43}
\end{equation}
under the homology identification \eqref{maps:eq:IC_local_21}, and its cokernel is
$\mathbb Z^{\rho-1}$ with exactly the quotient and inverse \eqref{maps:eq:IC_local_33}--\eqref{maps:eq:IC_local_34}.
This proves the asserted primitive embedding.

To compute the boundary image for every invariant degree vector,
forget the scalar circle, obtaining the actual base mapping torus
$Y$ of $N$. The maps \eqref{maps:eq:IC_local_3}--\eqref{maps:eq:IC_local_5} also apply to this base.
Its central cone is the mapping torus of $K$, and \eqref{maps:eq:IC_local_12} gives
\begin{equation}
 H_2(Y;\mathbb Z)=\mathbb Z^v/A_0\mathbb Z^v
                 \xrightarrow[\cong]{R}\mathbb Z^\rho .
 \label{maps:eq:IC_local_43a}
\end{equation}
The inverse inserts orbit representatives. The projection of the
circle bundle sends $[(x,t)]$ to $Rx$, exactly as its cell map
shows; this also follows by forgetting the scalar component in \eqref{maps:eq:IC_local_8}.

The regular three-torus in $Y$ has monodromy
$\bar A=\left(\begin{smallmatrix}B&q\\0&1\end{smallmatrix}\right)$.
Its primitive fixed vector is $(-q_0z,q_0)$, by \eqref{maps:eq:IC_local_28}.
Take its literal base normal torus. On the $q_0$-fold long cover,
\eqref{maps:eq:IC_local_35}--\eqref{maps:eq:IC_local_36} take it to the boundary of the covering long circle
times the actual section disc, with every base translation retained.
Transfer and covering projection multiply its image in \eqref{maps:eq:IC_local_43a} by
$q_0$; that image is therefore zero since \eqref{maps:eq:IC_local_43a} is free.

This also determines the scalar entry of the full monodromy at
the stated generality. Its integral block has the form
\begin{equation}
 A=\begin{pmatrix}
       B&0&q\\
       r_{\rm mon}&1&\alpha\\
       0&0&1
    \end{pmatrix}
 \label{maps:eq:IC_local_43b}
\end{equation}
because the actual principal-circle map preserves the scalar
generator. Lift the base vector $(-q_0z,q_0)$ with scalar coordinate
zero. Its endpoint differs from its start by
$q_0(\alpha-r_{\rm mon}z)\delta$.
By the same oriented clutch construction as \eqref{maps:eq:IC_local_7}, this integer is
the Chern pairing of the circle bundle on the base normal torus:
it is the winding of the lifted edge around the final scalar fibre.
That torus has zero image in $H_2(Y)$, so the pairing is zero.
Thus $\alpha=r_{\rm mon}z$. This argument retains the original line units
and applies to every invariant $c$, independently of a particular
Mordell--Weil parametrization of the line.

The equation $(B-I)z=q$ now shows directly that a fixed vector
satisfies $x=-z\,x_\gamma$, and the scalar equation vanishes.
By \eqref{maps:eq:IC_local_28}, the integral fixed lattice has basis
\begin{equation}
 \delta,\qquad
 \ell=(-q_0z_1,-q_0z_2,0,q_0).
 \label{maps:eq:IC_local_44}
\end{equation}
The integer vector is primitive because of \eqref{maps:eq:IC_local_40}.
The normal torus for $\delta$ bounds the scalar circle bundle
over the actual section disc in an extending nonzero line frame.
The second literal fixed normal projects to the base normal torus
just treated. Its component projection is zero by \eqref{maps:eq:IC_local_43a} and
naturality. Every line-frame adjustment is a scalar prism and is
still recorded in its second coordinate under \eqref{maps:eq:IC_local_15} or \eqref{maps:eq:IC_local_21}.
Its actual cellular image $(x,t)$ therefore satisfies $Rx=0$.
Both fixed normals have image in $\Psi(0,\mathbb Z/h)$
for $c\ne0$, or in the integer scalar-long factor for $c=0$.
These factors are already in the fibre image by \eqref{maps:eq:IC_local_42}--\eqref{maps:eq:IC_local_43}.
The Wang quotient of the boundary is generated by these two
literal normal tori, so the boundary image equals the fibre image.

For an integral fixed class, take the corresponding integral sum
of the two normal tori. Its cap image has scalar coordinate $b_{\rm nor}$
read by \eqref{maps:eq:IC_local_15} or \eqref{maps:eq:IC_local_21}. Subtract $b_{\rm nor}(\delta\wedge\gamma)$ in the
boundary. Its cap image is zero, and its Wang projection is the
prescribed fixed class because a fibre torus has zero Wang projection.
This is an integral lift of every element of $\operatorname{Fix}(A)$
to the boundary kernel.


\item \emph{The zero-degree scalar-long coordinate.}
The scalar-long subgroup for $c=0$ can contain a nonzero normal image.
Take a type $II$ neighbourhood $N$ with section and base coordinate $u$,
put $F=\operatorname{div}(u)$, $Q=0$, and $M=\mathcal O(F)$.
The global nonvanishing frame is $\tau=1/u$, and its canonical section
is $1=u\tau$. A constant contraction $\lambda$, $0<|\lambda|<1$,
gives
\begin{equation}
 U=N\times(\mathbb C^*/\lambda^{\mathbb Z})
   \simeq N\times S^1_\delta\times S^1_\gamma .
\end{equation}
The inverse of the frame coordinate sends $W$ to $W\tau$.
From the trivial line with its constant frame, the canonical tensor map
on the punctured neighbourhood and its inverse are
\begin{equation}
 I_{\rm tens}(p,w)=(p,u(p)w),\qquad I_{\rm tens}^{-1}(p,W)=(p,W/u(p)).
\end{equation}
Writing $\mu_{\rm mer}$ for the source meridian and $\mu_{0,{\rm mer}}$
for the meridian in the extending frame, these maps give
\begin{equation}
\begin{gathered}
 (I_{\rm tens})_*\mu_{\rm mer}=\mu_{0,{\rm mer}}+\delta,\quad
 (I_{\rm tens})_*\delta=\delta,\quad (I_{\rm tens})_*\gamma=\gamma,\\
 (I_{\rm tens})_*^{-1}\mu_{0,{\rm mer}}=\mu_{\rm mer}-\delta .
\end{gathered}
\end{equation}
Here $q_Q=1,z=0,\ell=\gamma$. The torus
$\gamma\wedge\mu_{0,{\rm mer}}$ bounds the section disc times the long circle, so
the marked second normal has image
\begin{equation}
 [\gamma\wedge\mu_{\rm mer}]=[\gamma\wedge(\mu_{0,{\rm mer}}+\delta)]=-t .
\end{equation}
The inverse meridian change contributes $+t$.
Projection to the two circle factors and inclusion at a fixed point
of $N$ are inverse on this subgroup, proving that $t$ is primitive
and has infinite order. This class is already in the fibre image.

\item \emph{Complete marked maps when $h=1$.}

Here $g=1$, the homology class $t$ is zero, and
$H_2(U)=\ker d^t$ is free. Formula~\eqref{maps:eq:IC_local_39}
therefore gives the complete fibre map
\begin{equation}
 F_2:\mathbb Z^6\longrightarrow\ker d^t,\qquad
 F_2(y)=q_0(y_{12}+z_2y_{14}-z_1y_{24})\,w.
 \label{maps:eq:IC_local_45}
\end{equation}
The vector $w$ is primitive, so its integral span embeds in
$\ker d^t$ with inverse given by the first coordinate of $L^{-1}$.
Thus $\ker F_2$ is exactly \eqref{maps:eq:D.IC7}.
Apply \eqref{maps:eq:IC_local_31} to its six-entry primitive row. Its last five columns
give a basis of the kernel, and the last five coordinates of the
inverse matrix give the inverse marking; its first column is an
integral section of the image map.

For each literal normal, let $(x,t)$ be its cellularized image,
including every actual line-frame prism. The preceding projection calculation proves $Rx=0$.
Choose $\eta$ as in \eqref{maps:eq:IC_local_14}, now with $c^t\eta=1$.
The exact bounding chain is
\begin{equation}
 d_3\bigl(0,\mathcal K x+(t-c^t\mathcal K x)\eta\bigr)=(x,t).
 \label{maps:eq:IC_local_46}
\end{equation}
Thus both literal fixed-normal tori bound. Together with the
kernel of the fibre map \eqref{maps:eq:IC_local_45}, these give the complete boundary kernel:
their Wang projections give the two fixed basis vectors, and
subtracting those two normal multiples reduces every kernel
element to the fibre Wang quotient, where \eqref{maps:eq:IC_local_45} applies.

Equation \eqref{maps:eq:IC_local_24} is now the infinite cyclic group generated by the
positive roof loop. Every vertical regular-fibre loop projects to
a nullhomotopic loop in the simply connected $K$; its lift lies
in the scalar subgroup. That subgroup is killed by the degree
Bézout relation. Hence the complete degree-one map is the literal
$\gamma$ coordinate, with section the positive roof loop.
Its kernel is $\mathbb Ze_1\oplus\mathbb Ze_2\oplus\mathbb Z\delta$,
with inverse marking given by these three coordinates.

\item \emph{The two star degree vectors and their component markings.}

Use the affine $D_{r_{\rm Dyn}}$ diagram with $r_{\rm Dyn}=n+4$, $n=1,3$.
The near leaves 0 and 1 meet vertex 2; the inner chain runs from
2 to $r_{\rm Dyn}-2$; the far leaves $r_{\rm Dyn}-1,r_{\rm Dyn}$ meet $r_{\rm Dyn}-2$.
The intersection matrix $C$ has diagonal $-2$ and entry $1$
on each of these edges. The outer multiplicities are 1 and all
inner multiplicities are 2.

For component classes $P=Q=2$, the marked section $P$ meets
the near leaf 1. The actual $Q$ permutation is
$(0\,1)(r_{\rm Dyn}-1\,r_{\rm Dyn})$, with all inner vertices fixed.
In this chain order write $\varepsilon_i^{\rm ch}=\varepsilon_i$.
The component divisor of Theorem~\ref{maps:src:D.2.50} is $d_{\mathrm{can}}=-\varepsilon_0^{\rm ch}$. Its literal degree vector is
\begin{equation}
 c=(\varepsilon_1^{\rm ch}-\varepsilon_0^{\rm ch})+C(-\varepsilon_0^{\rm ch})=\varepsilon_0^{\rm ch}+\varepsilon_1^{\rm ch}-\varepsilon_2^{\rm ch}.
 \label{maps:eq:IC_local_47}
\end{equation}
This reads the degrees of the specified line
$M_0\otimes\mathcal O(-C_0)$. Its transition functions are those
of $M_0$ multiplied by those of $\mathcal O(-C_0)$; their
units are included in the component clutch maps above.

Pass from the chain order to the outer-components-first order
by the index map
\begin{equation}
 p=(0,1,r_{\rm Dyn}-1,r_{\rm Dyn},2,3,\ldots,r_{\rm Dyn}-2).
 \label{maps:eq:IC_local_48}
\end{equation}
For $n=1$ its tuple is $(0,1,4,5,2,3)$, with the same tuple
as inverse. For $n=3$ it is $(0,1,6,7,2,3,4,5)$, with inverse
$(0,1,4,5,6,7,2,3)$. Let $\varepsilon_k^{\rm out}$ be the outer-first coordinate column.
Then $\mathsf P\varepsilon_k^{\rm out}=\varepsilon_{p(k)}^{\rm ch}$ maps
outer-first coordinates to chain coordinates. With $c_{\rm ch}$ as in
\eqref{maps:eq:IC_local_47}, the degree and return maps are
\begin{equation}
 c_{\rm out}=\mathsf P^{-1}c_{\rm ch},\qquad
 \sigma_{\rm out}=\mathsf P^{-1}\sigma_{\rm ch}\mathsf P,\qquad
 \mathsf P^{-1}=\mathsf P^t.
 \label{maps:eq:IC_local_49}
\end{equation}
These give the two degree vectors in \eqref{maps:eq:D.IC8}.
The outer-first permutation is $(0\,1)(2\,3)$, with all other vertices
fixed. It has an orbit degree 1 on the first pair and degree
$-1$ at the fixed vertex 4. Therefore $g=h=1$.
In these outer-first coordinates, explicitly
$\eta_{\rm st}=-\varepsilon_4^{\rm out}$ satisfies
$A_0\eta_{\rm st}=0$ and $c^t\eta_{\rm st}=1$.
Use $\eta=\eta_{\rm st}$ for these two star caps.
The multiplicity pairing is $1+1-2=0$, as required.
All the marked maps for $h=1$ apply to these two actual
corrected line bundles.
\end{enumerate}
\end{proof}
\subsection{Height-integral initial pairs}

\begin{definition}[Height-integral filling data]\hypertarget{tag-Q204}{}\label{maps:def:height-integral-data}
{\normalfont\scriptsize tag Q204}\par
Choose one of the twenty-one marked Mordell--Weil rank-one
profiles of \ref{maps:src:D.2.19}, write
$\langle G,G\rangle=1/d_0$, and take
$P=aG,Q=bG$ with $a,b>0$ and $d_0\mid ab$.
Choose nonnegative integral whole-fibre weights of total
$ab/d_0$, supported at multiplicative or additional regular
fibres, and use the corrected line of \ref{maps:src:D.2.50}.
At each site choose a filling from the following list, retaining
its indicated markings:
\begin{enumerate}[label=(\roman*)]
\item marked rank-one or triangular Mumford caps, with their
prescribed periodic fan, at multiplicative or regular sites;
\item free finite-good additive logarithmic caps, with their
affine powers and all five integral Thom columns;
\item original additive caps whose corrected multidegree has
$h=1$ in \eqref{maps:eq:D.IC2}, and unchanged zero-degree
original additive caps with their extending frames;
\item original multiplicative caps with the uncorrected line,
degree vector and relative frame of \ref{maps:src:D.2.46};
\item any finite number of regular logarithms at additional
weight-zero sites, and any ordinary integral clutch.
\end{enumerate}
Positive weights use the prescribed Mumford filling. Odd-star
fibres use their original caps. A supplied integral component
divisor must satisfy \eqref{maps:eq:D.VC3}; its component
permutation, degree vector and extending frames enter the
construction. An original additive cap has either $h=1$, or
zero degree with its unchanged extending frame and relative marking.
These choices are the \emph{height-integral filling data}.
\end{definition}

\begin{theorem}[Height-integral fillings]\hypertarget{tag-02LU}{}\label{maps:src:D.2.52}
{\normalfont\scriptsize tag 02LU}\par
For every datum of Definition~\ref{maps:def:height-integral-data},
a finite exact integral algorithm computes a fundamental-group
presentation and every integral homology group. The canonical corrections of \ref{maps:src:D.2.50} supply these
fillings for every positive pair of integral height. The only
additive positions where both section classes are nonzero have component pair
$(2,2)$ at an odd star, with degree vector
\eqref{maps:eq:D.IC8}; multiplicative positions where both classes are nonzero
may use the marked Mumford filling. The initial-pair range is all positive $(a,b)$ with $d_0\mid ab$. Across the twenty-one profiles, exactly $45$
height-integral residue pairs modulo $d_0$ have a fibre at which
both section classes are nonzero.

The algorithm applies to every finite list and all section multiples,
weights, logarithmic orders and affine numerators in the defined input.
The line and relative frame are part of each original-cap datum;
a whole-fibre tensor change transports its meridian by the explicit
frame map in the proof.
\end{theorem}

\begin{proof}
\begin{enumerate}
\item \emph{Corrected lines and marked periods.}
Use Theorem~\ref{maps:src:D.2.50} with the same marked surface,
$P=aG$, $Q=bG$, divisor $D$, and whole-fibre weights as in
Definition~\ref{maps:def:height-integral-data}. Its local coefficient
column $z$ is the present column $d$, its component label
$p_{\rm comp}$ is the present $p$, and its pullback matrix
$\mathsf S$ is the present $S$. Write $\varepsilon_j$ for the same ordered component basis.
Thus the identity coordinate maps
carry \eqref{maps:eq:D.VC3} to
\[
 C(I-S)d=-(I-S)(\varepsilon_p-\varepsilon_0).
\]
The remaining left multiplication is $-I$, with inverse $-I$.
The entire solution ambiguity is therefore the same integral
orbit-constant lattice, with its representative-coordinate inverse.

The line bundles are literally
$M_0=\mathcal O(P-O)$ and $M_D=M_0\otimes\mathcal O(D)$.
On the complement of the singular fibres, the comparison is
\[
 I_{D,p}(v)=v\otimes\sigma_D(p),\qquad
 I_{D,p}^{-1}(v\otimes\sigma_D(p))=v,
\]
and the difference-line map multiplies by
$\sigma_D(p)/\sigma_D(p+Q)$, with reciprocal inverse.
The arrow $A_{D,p}:M_{D,p+Q}\to M_{D,p}$ covers $-Q$; its
conjugate is
\[
 I_{D,p}^{-1}A_{D,p}I_{D,p+Q}
   =j_D(p)\frac{\sigma_D(p+Q)}{\sigma_D(p)}.
\]
These are exactly the tensor, group-law and base-line maps in
Theorem~\ref{maps:src:D.2.50}. A coefficient $q^{w_i}u_i(q)$
uses the same nonzero base unit $u_i$ in both directions.
Let $B,g_{\rm ell}$ be the marked elliptic monodromy and cocycle of $G$,
and put $J=\left(\begin{smallmatrix}0&1\\-1&0\end{smallmatrix}\right)$.
The given boundary marking $(e_1,e_2,\delta,\gamma)$ then has
precisely the monodromy $T_i$ of \eqref{maps:eq:D.GI2}, now on the entire height-integral range
proved there. Its period shear and inverse act on the same
transported lattices: $z_G$ solves $(B-I)z_G=g_{\rm ell}$, the elliptic
translation is $bz_G$, and the line row is
$-ag_{\rm ell}^{\mathsf t}JB=az_G^{\mathsf t}J(B-I)$.
The component order divides $d_0$, so $abz_G$ is integral.
The proved inverse monodromy is integral, and the same ordered
base-degree cocycle gives
\[
 \prod_iT_i=I+\left(\sum_iw_i-ab/d_0\right)\delta\gamma^*=I.
\]
The integer bar construction therefore applies to this tuple
with its actual periods and units.

For an invariant orbit divisor $E_{\rm orb}$, multiplication by $\sigma_{E_{\rm orb}}$
and division by it are inverse maps on the punctured family.
Its arrow ratio cancels through the canonical trivialization
$\sigma_{E_{\rm orb}}/T_Q^*\sigma_{E_{\rm orb}}$ of the difference line. The period
blocks and specified base weights therefore agree under that map
on the complement of the zeros and poles of $\sigma_{E_{\rm orb}}$.

\item \emph{Original additive caps and their marked maps.}
At an original additive site in the $h=1$ sector, the weight is
zero and the same line $M_D$ has invariant degree vector $c$. Theorem~\ref{maps:src:D.2.51} therefore applies
to this actual original cap with its line arrow and extending frames.
Here is the exact transport between return conventions.

For the unit-circle return $\phi$, write $E_{\rm circ}$ for its
circle-bundle fibre and put
$\mathcal T(\phi)=(E_{\rm circ}\times[0,1])/((x,1)\sim(\phi x,0))$.
The two roof presentations are related by
\[
 j:\mathcal T(\phi)\longrightarrow\mathcal T(\phi^{-1}),
 \quad [x,t]\longmapsto[\phi x,1-t],\qquad
 j^{-1}[y,s]=[\phi^{-1}y,1-s].
\]
The endpoint relations prove that both maps are well-defined and
that both composites are the identity. At roof zero the target
identification makes $j$ the identity on the component and scalar
cells. It reverses the roof marking: $\widehat\gamma=-\gamma$,
$\widehat\delta=\delta$, and $\widehat t=-t$.
For the cone of a component return $\sigma$, the induced chain maps
to the cone of $\sigma^{-1}$, in degrees zero through four, are
\[
 1,\quad \operatorname{diag}(1,-1),\quad
 \operatorname{diag}(I_v,-1),\quad
 \operatorname{diag}(I_v,-\sigma),\quad -\sigma.
\]
Their inverses replace $\sigma$ by $\sigma^{-1}$.
For example
$(\sigma^{-1}-I)(-\sigma)=\sigma-I$ and
$c^{\mathsf t}(-\sigma)=-c^{\mathsf t}$; these are the two
degree-three boundary identities, and the degree-four identity is
$(I-\sigma^{-1})(-\sigma)=I-\sigma$.
Write $\mathbf q$ for the marked section cocycle in this return
convention. Return inversion carries it and its logarithm to
$\widehat{\mathbf q}=-\mathbf q$ and $\widehat z=-z$.
The corresponding fibre basis change is
$P_\gamma=\operatorname{diag}(I_2,1,-1)$, its own inverse;
the monodromy changes by $\widehat T=P_\gamma^{-1}TP_\gamma$.
Its elliptic--long column is therefore $-\mathbf q$.
Consequently each mixed term has two cancelling sign changes:
\[
 \widehat y_{12}+\widehat z_2\widehat y_{14}
       -\widehat z_1\widehat y_{24}
 =y_{12}+z_2y_{14}-z_1y_{24}.
\]
The scalar normal is fixed and the other fixed-normal generator
changes sign. Both retain the bounding chains of
Theorem~\ref{maps:src:D.2.51}, with the displayed inverse chain maps.
Thus that theorem's positive-return calculation and a forward-arrow
calculation give the same marked maps in the chosen period basis.
All line units occur in $\phi$ and its actual inverse.

In this basis put $z=bz_G$ and let $q_0$ be its least positive
denominator. Use the orbit-sum map $R$, representative section $s$ and
partial-sum inverse $\mathcal K$ of
\eqref{maps:eq:IC_local_10}--\eqref{maps:eq:IC_local_12}, and put
$w=R\mu/q_0$. These are exactly
\eqref{maps:eq:IC_local_27}--\eqref{maps:eq:IC_local_30};
the weighted diagrams there prove $q_0=\gcd(\ell_j\mu_j)$.
The primitive class $W$ has orbit coordinates $w$; the section
$w\mapsto[(sw,0)]$ and its inverse $[(x,t)]\mapsto Rx$
are the $h=1$ specialization of
\eqref{maps:eq:IC_local_15}--\eqref{maps:eq:IC_local_19}.
The complete fibre map is therefore
\[
 \bigwedge\nolimits^2\Lambda\longrightarrow H_2(U),\qquad
 y\longmapsto q_0(y_{12}+z_2y_{14}-z_1y_{24})W,
 \quad\Lambda=\mathbb Z\langle e_1,e_2,\delta,\gamma\rangle.
\]
Its integral primitive row has the B\'ezout section and inverse
kernel coordinates of \eqref{maps:eq:IC_local_31}--\eqref{maps:eq:IC_local_34}.
The actual fixed lattice has basis
$\delta$ and $(-q_0z_1,-q_0z_2,0,q_0)$.
For a normal image $(x,t)$, its component projection is $Rx=0$.
With $y_{\rm pr}=\mathcal Kx$ and invariant $\eta$ satisfying
$c^{\mathsf t}\eta=1$, the complete bounding chain is
\[
 d_3\bigl(0,y_{\rm pr}+(t-c^{\mathsf t}y_{\rm pr})\eta\bigr)=(x,t),
\]
as in \eqref{maps:eq:IC_local_46}; it includes every frame prism.
The fibre kernel and these two normals are the complete boundary
kernel. The degree-one map reads the primitive $\gamma$ coordinate,
with its roof-loop section. The groups are
\[
 H_0(U)=H_1(U)=\mathbb Z,\quad H_2(U)=\mathbb Z^{\rho-1},
 \quad H_3(U)=\mathbb Z^{2\rho-1},\quad H_4(U)=\mathbb Z^\rho,
 \qquad\chi(U)=0.
\]
The cone's orbit inclusions and inverse coefficient readings provide
the component-disc and higher-degree maps used below.

\item \emph{Odd-star corrections and the residue count.}
For the affine $D_5,D_7$ diagrams, use the chain and outer-first
component maps \eqref{maps:eq:IC_local_47}--\eqref{maps:eq:IC_local_49}.
Use $\varepsilon_i^{\rm ch}$ and $\varepsilon_i^{\rm out}$ for the ordered
component columns in these two markings; the index-permutation matrices
and their transposes give the inverse changes. They give the actual divisor $d_{\rm can}=-\varepsilon_0^{\rm ch}$ at
$(\overline P,\overline Q)=(2,2)$, its degree
$\varepsilon_0^{\rm ch}+\varepsilon_1^{\rm ch}-\varepsilon_2^{\rm ch}$, both inverse index permutations and the degree vectors
\eqref{maps:eq:D.IC8}. The outer-first return is $(01)(23)$,
and $\eta_{\rm st}=-\varepsilon_4^{\rm out}$ proves $h=1$.
The alternative divisor $\varepsilon_1^{\rm ch}$ differs by $\varepsilon_0^{\rm ch}+\varepsilon_1^{\rm ch}$.
Multiplication by the canonical section of
$\mathcal O(-C_0-C_1)$ and its reciprocal compare these punctured
lines, conjugating the arrows by the same section ratio.
Each original cap keeps its own supplied line and extending frame.

The nontrivial component orders $n_i$ in the marked profiles of
Theorem~\ref{maps:src:D.2.50} are pairwise coprime
and their product is $d_0$. The CRT coordinate map is
$a\mapsto(a\bmod n_i)_i$. If $M_{i,\rm CRT}=d_0/n_i$ and
$M_{i,\rm CRT}t_{i,\rm CRT}\equiv1\pmod{n_i}$, its inverse is
$\sum_i M_{i,\rm CRT}t_{i,\rm CRT} a_i\bmod d_0$. Both composites are the
identity on every residue coordinate. Applying it to
$a$ and $b$ makes $d_0\mid ab$ equivalent to
$n_i\mid a_i b_i$. There are
\[
 H(h)=\sum_{a\bmod h}\gcd(h,a)
 \quad\hbox{such local pairs, and}\quad N(h)=2h-1
\]
pairs with $a_i=0$ or $b_i=0$. The number of additional
pairs per profile is $\prod_iH(n_i)-\prod_iN(n_i)$.
For a prime $h$ these counts agree. At $h=4$ the only
additional pair is $(2,2)$; at $h=8$ the pairs are
$(2,4),(4,2),(4,4),(4,6),(6,4)$.
The CRT inverses give the complete nonzero contributions:
\[
\begin{array}{c|l|r}
 \text{profiles}&\text{additional pairs in each profile}&\text{total}\\ \hline
 \mathrm{OS45/default},\mathrm{OS45/II}
   &(2,4),(4,2),(4,4),(4,6),(6,4)&10\\
 \mathrm{OS46/default},\mathrm{OS46/II}&(2,2)&2\\
 \mathrm{OS50/default},\mathrm{OS50/II},\mathrm{OS50/IV}
   &(2,6),(6,2),(6,6),(6,10),(10,6)&15\\
 \mathrm{OS55/default},\mathrm{OS55/II}
   &(2,10),(6,10),(10,2),(10,6),(10,10),&18\\
   &\quad(10,14),(10,18),(14,10),(18,10)&
\end{array}
\]
The other profiles contribute zero, so the sum is $45$.
At an additive prime-order position height integrality
forces one section to be narrow. The only remaining
additive case is the displayed odd-star pair $(2,2)$.
It has the required literal correction with $h=1$.
This proves the stated initial-pair range. At a multiplicative
site where both section classes are nonzero, use its prescribed
Mumford cap.

\item \emph{Global relations and the fundamental group.}
Apply the integral ordered bar model of
\ref{maps:src:D.2.41} to the tuple $T_i$ just constructed.
It keeps the six fibre bivectors, their return columns,
and the integral graph kernel with its Smith coordinate
inverse. The binomial quadratic cochains and the suffixes
of the ordered last-boundary word give every normal column.
These maps apply to all integral determinant-one matrices
with the displayed product, so they apply to this tuple.

At the corrected $h=1$ original additive caps insert the full kernel,
both literal normals and their bounding chains proved above.
At every other permitted cap use the full marked maps of
\ref{maps:src:D.2.46} and \ref{maps:src:D.2.43}:
the prescribed-fan period and logarithm inverse, all component
normal discs, the nonsaturated fibre kernel and boundary
torsion, or the finite sheet/norm maps and all five Thom
columns. These local maps are functions of the displayed
periods, fan, extending frames and affine power; each input
is retained here. At an uncorrected multiplicative site,
invariance of $\varepsilon_p-\varepsilon_0$ forces $p=0$ or $q=0$: for $p\ne0$
the coefficient $-1$ occurs uniquely at zero, which must
be fixed by the component translation. Thus this cap has
exactly the unchanged geometry and frame covered by
\ref{maps:src:D.2.46}.

The fundamental-group presentation retains the four
commuting fibre generators, each peripheral generator,
conjugation by $T_i$, the ordered peripheral product,
every affine power relation, and every nonlogarithmic cap
quotient. At an $h=1$ additive cap it kills
$e_1,e_2,\delta$ and the origin meridian and retains
the positive $\gamma$ loop. The forward maps are these
paths; the inverses read the roof, finite sheet or short/long
period coordinates and then the cut-base graph coordinates.
They agree on every boundary generator and their two
composites fix every generator. Van Kampen therefore gives
the asserted finite presentation, including its nonabelian
relations. Abelianization computes $H_1(X)$.

The vertical normal closure is a terminating
integer procedure with explicit subgroup maps. Let
$V_{\rm vert}=\mathbb Z\langle e_1,e_2,\delta\rangle$, let
$M_{i,\rm vert}=T_i|_{V_{\rm vert}}\in\mathrm{SL}(3,\mathbb Z)$, and let $L_0$
be the integer span of the killed vertical vectors. Set
\[
 L_{r+1}=L_r+\sum_iM_{i,\rm vert}L_r.
\]
The union is a subgroup of $\mathbb Z^3$, hence finitely
generated. Each of its finitely many generators occurs
at some stage; their maximal stage contains the union.
This proves termination. Hermite column operations retain
the exact subgroup and the integer word for each basis
vector, together with the inverse coordinate reading.
If the characteristic polynomial of $M_{i,\rm vert}$ is
$x^3-t_ix^2+s_ix-1$, Cayley--Hamilton gives
\[
 M_{i,\rm vert}^{-1}=M_{i,\rm vert}^2-t_iM_{i,\rm vert}+s_iI.
\]
Positive stability thus gives inverse stability, so the
stable lattice contains all positive and negative conjugates.

Before the cap and power relators, the exterior group is
the semidirect product of the abelian fibre with the
cut-base free group. Send $M_{w,\rm vert}v$ to
$t_wx_vt_w^{-1}$ for each initially killed vector $v$.
Products of conjugates add their fibre exponent vectors.
Reading that exponent is the inverse map, so the stable
lattice is precisely this normal subgroup. Taking its
quotient first or adding its generators as relators gives
inverse homomorphisms on every fibre and peripheral
generator. All remaining power and product relations are
then retained. This gives the unsaturated normal-closure
quotient within the full fundamental-group presentation.

\item \emph{Integral splitting and homology.}
Let $\mathcal B$ be the complete fibre-and-graph relation
matrix. Write $e_{j,\rm in}$ and $e_{j,\rm out}$ for the standard
columns in its domain and codomain, respectively. Keep the Smith transformations
$U_{\rm SNF}\mathcal B V_{\rm SNF}=\Delta$. The quotient generators are
$U_{\rm SNF}^{-1}e_{j,\rm out}$, and the relations satisfy
\[
 \mathcal B V_{\rm SNF} e_{j,\rm in}=\Delta_{jj}U_{\rm SNF}^{-1}e_{j,\rm out}.
\]
The inverse map reads $U_{\rm SNF}$ times a represented cycle,
modulo the corresponding diagonal factor. This retains
every torsion factor of a nonsaturated local kernel.

At least one weight is positive, since $ab/d_0>0$.
Its prescribed Mumford cap is nonlogarithmic. Attach the
finite and regular logarithmic caps first and call the
result $E_{\rm ext}$. Their root-disc boundaries are
$m_i\tau_i-v_i$. A relation among them in peripheral
coordinates is a multiple of the single sum relation.
Its coefficient at the remaining nonlogarithmic site is
zero, so that multiple is zero; each coefficient of a
root boundary is then zero. Reading a peripheral
coordinate and dividing by its $m_i$ gives the inverse
on the image. Thus this first stage is exactly the
quotient by the five Thom columns at each logarithmic cap.

For each remaining cap $U_i$, take its component or orbit
normal discs in the extending Cartier frames. Their
boundaries follow the stems, period logarithms and integer
frame jumps already used in the bar model. They define
\[
 \Gamma:\mathbb Z^{\sum_i h_i}\longrightarrow H_1(E_{\rm ext}),
 \qquad h_i=b_4(U_i).
\]
At a corrected $h=1$ cap, $h_i=\rho_i$, and its free $H_3$
shows directly that $H^4(U_i)$ is the free component-orbit
target. The other allowed caps have the component-disc
targets of \ref{maps:src:D.2.46} and
\ref{maps:src:D.2.43}. Exponentiating the period coordinate
and reversing the chart transitions invert each boundary
construction, with its meridian coefficient.

The pair sequence is therefore
\[
 0\longrightarrow\operatorname{coker}\mathcal B
  \longrightarrow H_2(X)
  \longrightarrow\ker\Gamma\longrightarrow0.
\]
Here is an integral section and its inverse. Choose an
injective basis matrix $P_{\rm rel}$ for the relation image in a
free presentation of $H_1(E_{\rm ext})$ and retain its preimages
in the original relation columns. Lift $\Gamma$ to a
matrix $M_{\rm disc}$ and solve $M_{\rm disc}z_{\rm disc}=P_{\rm rel}z_{\rm rel}$ by integer Smith operations.
Projection to $z_{\rm disc}$ has inverse recovering the unique $z_{\rm rel}$
because $P_{\rm rel}$ is injective. For a resulting kernel basis
vector, join the component discs to the retained relation
chains, exterior cylinders and ordered torus polygons.
This is a closed integral two-cycle. Reading the disc
coordinates and subtracting these cycles is the inverse
splitting. In particular the full integral kernel,
including its index in its rational span, is retained.

The local degree-one quotient maps are onto, so
$\operatorname{coker}\Gamma=H_1(X)$. The relative free
rank is consequently
\begin{equation}
 \sum_i b_4(U_i)-b_1(E_{\rm ext})+b_1(X).
                         \hypertarget{tag-01YN}{}\label{maps:eq:D.CI1}
\end{equation}
Write $\operatorname{coker}\mathcal B=\mathbb Z^g\oplus T_{\rm tor}$
and $H_1(X)=\mathbb Z^r\oplus T_{H_1}$, with $T_{\rm tor},T_{H_1}$ finite.
If $q_{\rm rel}$ is \eqref{maps:eq:D.CI1}, then
\[
 b_2=g+q_{\rm rel},\qquad
 \chi(X)=2\sum_{\text{positive multiplicative }i}\nu_iw_i,
 \qquad b_3=2+2b_2-2r-\chi(X).
\]
The corrected $h=1$ additive caps have Euler number zero by the cone
calculation. The other cap indices give the displayed
Euler sum, and the exterior and boundaries are mapping
tori with zero Euler number. The complex orientation and
the determinant-one ordinary clutch orient the closed
six-manifold. Integral duality and universal coefficients
give
\[
\begin{array}{c|c}
 j&H_j(X;\mathbb Z)\\ \hline
 0,6&\mathbb Z\\
 1&\mathbb Z^r\oplus T_{H_1}\\
 2&\mathbb Z^{b_2}\oplus T_{\rm tor}\\
 3&\mathbb Z^{b_3}\oplus T_{\rm tor}\\
 4&\mathbb Z^{b_2}\oplus T_{H_1}\\
 5&\mathbb Z^r.
\end{array}
\]
For a torsion factor the resolution
$\mathbb Z\xrightarrow{d_{\rm tor}}\mathbb Z$, its dual unit
cochain and oriented cap product give the torsion map;
reading the residue modulo $d_{\rm tor}$ is its inverse. This also
computes the Euler-zero cases.

\item \emph{Relative frames.}
The algorithm forms $C(I-S)d$ and compares it integrally
with $-(I-S)(\varepsilon_p-\varepsilon_0)$. It reads the literal degree $c$,
its orbits and $h=\gcd(\ell_jc_j)$; the preceding maps
apply to $h=1$. At an unchanged zero-degree original
cap it uses its given extending frame and the original
maps of \ref{maps:src:D.2.46}.

To specify the whole-fibre change, write $F=\operatorname{div}(q)$
and $D=kF$. The punctured tensor map is
$I_D(v)=v\sigma_D$, with reciprocal inverse. The extending
isomorphism is $E_D(v)=vq^{-k}\sigma_D$, since
$q^{-k}\sigma_D$ is its nonzero extending frame. Translation
preserves $q$, and the line arrows conjugate by these exact
multipliers. Hence
\[
 I_D^{-1}E_D(v)=q^{-k}v,\qquad E_D^{-1}I_D(v)=q^kv.
\]
On the boundary the first sends
$\tau\mapsto\tau-k\delta$ and the inverse sends
$\tau\mapsto\tau+k\delta$; both fix $\delta$.
At $h=1$ the resulting scalar-long addition bounds by
the displayed $\eta$ chain. The zero-degree cap retains
this meridian shear as part of its relative marking.
\end{enumerate}
\end{proof}
\section{Minimal primary smooth types}

\subsubsection*{Order two}

\begin{definition}[Minimal order-two data]\hypertarget{tag-Q205}{}\label{maps:def:minimal-two-data}
{\normalfont\scriptsize tag Q205}\par
Use the marked elliptic data of
Definition~\ref{mapsIII:def:elliptic-marked-data} and choose one
of the patterns below. Use the prescribed primitive Mumford fan
at each multiplicative site, the listed original or logarithmic
additive cap, and the ordinary clutching $e\gamma$ of
\eqref{maps:eq:D.HY3}. The logarithmic sites are exactly those
listed in the pattern. At a finite site of linear monodromy order
$d_i$, take $m=d_i k$ with $k>0$, and any free affine lift.
The ordinary clutch parameter is $e\in\mathbb Z$. The selected zero's cocycle is $\kappa u$,
with primitive $u$ whose first nonzero coordinate is positive;
put $\varepsilon=1+\kappa\pmod2$.

The selected positions in the following table are $i-1$, where $i$
is the peripheral index in Definition~\ref{mapsIII:def:elliptic-marked-data}. The two additive choices at an $II$ site
are separate patterns:
\begin{equation}
\begin{array}{c|r|r|l|l|r|c}
\text{OS/profile}&a&b&\text{zero}&\text{additive choice}
      &\kappa&\varepsilon\\\hline
\mathrm{OS47}&7&2&2&\text{none}&0&1\\
\mathrm{OS47}&7&2&3,4&\text{none}&1&0\\
\mathrm{OS47}/\mathrm{II}&7&2&2&II\text{ original or log}&0&1\\
\mathrm{OS49}&1&6&2,3&IV^*\text{ log}&0&1\\
\mathrm{OS56}&5&2&1&\text{none}&-2&1\\
\mathrm{OS56}&15&2&3&\text{none}&0&1\\
\mathrm{OS56}&15&2&4&\text{none}&-1&0\\
\mathrm{OS56}/\mathrm{II}&5&2&3&II\text{ original or log}&1&0\\
\mathrm{OS56}/\mathrm{IV}&1&6&0&IV\text{ log}&-4&1\\
\mathrm{OS56}/\mathrm{IV}&5&6&3&IV\text{ log}&-1&0\\
\mathrm{OS56}/\mathrm{IV}&15&2&3&IV\text{ original}&-1&0.
\end{array}                                               \hypertarget{tag-01YF}{}\label{maps:eq:D.C22}
\end{equation}

Let $M_1=M_{\mathrm{conic}}$ be the oriented conic surgery of
\ref{maps:src:D.2.8}. Define $M_0$ by joining two positively oriented
parallel spheres $S^2\times\{p_+\}$ and $S^2\times\{p_-\}$ in
$S^2\times S^4$ by an oriented tube along a fixed arc, and performing
framed two-sphere surgery on the resulting sphere. Its oriented normal
framing is unique up to homotopy. These choices constitute the \emph{minimal order-two data}.
\end{definition}

\begin{theorem}[Order-two smooth types]\hypertarget{tag-02LB}{}\label{maps:src:D.2.35}
{\normalfont\scriptsize tag 02LB}\par
Every simply connected threefold $X$ determined by
Definition~\ref{maps:def:minimal-two-data}, with any permitted finite
multiplicity and affine twist, satisfies
\begin{equation}
 H_2(X;\mathbb Z)=H_3(X;\mathbb Z)=\mathbb Z/2,
 \qquad \langle\alpha^3,[X]_2\rangle=\varepsilon,
 \qquad X\cong M_\varepsilon
                                                        \hypertarget{tag-01Y7}{}\label{maps:eq:D.C21}
\end{equation}
by an orientation-preserving diffeomorphism, where $\alpha$ is the
unique nonzero element of $H^2(X;\mathbb F_2)$.


The fifteen patterns occupy thirteen marked initial rows. Eight
patterns on seven rows give $M_1$, and seven patterns on six rows give
$M_0$. The two oriented models have distinct mod-two cubic forms.
\end{theorem}


\begin{proof}
\begin{enumerate}
\item \emph{The marked inputs and the common piece.}
Use the actual lattice
$\Lambda=\mathbb Z\langle e_1,e_2,\delta,\gamma\rangle$,
where positive $\gamma$ covers translation by $-Q$. Write $P=aG$,
$Q=bG$. The invariant alternating form is
\begin{equation}
 \Omega=a\,e_1^*\wedge e_2^*
                -b\,\delta^*\wedge\gamma^*.
                                                        \hypertarget{tag-01YG}{}\label{maps:eq:D.C23}
\end{equation}
Substitution in the forward period formula of \ref{maps:src:D.2.23}
gives $T_i^t\Omega T_i=\Omega$ and $T_i\delta=\delta$.
Every row in the table has $a$ odd, $b=2$ or $6$, and $b/2$ odd.
At a multiplicative place choose primitive $u,r$ with
\[
 A-I=\nu uu^tJ,\qquad g=\kappa u,\qquad
 J=\begin{pmatrix}0&1\\-1&0\end{pmatrix},\qquad \det(u,r)=1.
\]
Fix the sign of $u$ by taking its first nonzero coordinate positive.
These are the ordered elliptic matrix and section cocycle of
Definition~\ref{mapsIII:def:elliptic-marked-data}. At the selected zero put $S=(u,\delta)$,
$L=(r,\gamma)$ and
\[
 C_0=\begin{pmatrix}\nu&b\kappa\\
        -a\kappa&(1-ab\kappa^2)/\nu\end{pmatrix}.
\]
Its determinant is one. The actual short basis is $SC_0$, and its long
basis is $L$. The period-basis map is the identity on the complex
universal cover, with coordinates changed by
$H=(SC_0e_1,Le_1,SC_0e_2,Le_2)$ and inverse $H^{-1}$.
Exponentiation in the short basis $SC_0$ gives the prescribed integer
square lattice divided along its antidiagonals.

Here is the finite input determination. Write the generator height as
$1/d$; the natural-row equations of \ref{maps:src:D.2.22} are
\[
 ab=d/\nu,\quad a,b>0,\quad\gcd(a,b)=1,
 \quad\gcd(b\kappa,\nu)=1,
\]
together with the literal component equation
\[
 [j=a\kappa]-[j=0]-[j=(a-b)\kappa]+[j=-b\kappa]
       =-2[j=0]+[j=1]+[j=-1]
\]
and narrowness of $P$ or $Q$ at every other reducible fibre.
The repeated terms at $\nu=1,2$ are retained. For a finite additive
option, $(A-I)^{-1}bg$ must be integral. Its monodromy order and
vertical index are respectively
\[
\begin{array}{c|ccccc}
 &II&III&IV&III^*&IV^*\\\hline
 d_i&6&4&3&4&3\\ q_i&1&2&3&2&3.
\end{array}
\]
An original minimal additive option is available at $II,III,IV$
exactly when $P$ is narrow and $Q$ is transitive on its components,
by the cap maps of \ref{maps:src:D.2.26}. Simple connectivity requires
the finite orders to be pairwise coprime, by the Seifert presentation
of \ref{maps:src:D.2.23}. Apply these finite conditions to the
natural rows without an $I_\nu^*$ fibre in \ref{maps:src:D.2.22}
and retain $b=2\prod q_i$. The resulting choices are precisely
\eqref{maps:eq:D.C22}. Each has at most one finite option. The
original additive options have one component orbit and zero line
multidegree. Thus all their original-cap maps have the stated scope.

For one finite option write $m=d_i k>0$ and
$v=n\ell+r_0\delta$, with $\chi(\ell)=1$ for
$\chi=\gamma^*$. The group presentation is
\[
 \langle\gamma,c\mid[\gamma,c]=1,
                 c^m=\gamma^n, c=\gamma^e\rangle.
\]
The substitution $c\mapsto\gamma^e$, with inverse inclusion of
$\gamma$, identifies it with $\mathbb Z/(me-n)$. Hence the full
simply connected parameter range is $me-n=\pm1$, which implies
$\gcd(m,n)=1$. Every positive multiple of $d_i$ occurs, for example
with $e=1,n=m-1$; $r_0$ and every affine lift with this power are
allowed. With no finite option the group is $\mathbb Z/e$, and
simple connectivity means $e=\pm1$.

Remove the additive caps and one ordinary disc carrying the clutching.
Let $W$ be the resulting piece with all multiplicative caps filled,
and let $V$ be its regular torus-bundle part. At an unselected
multiplicative place put $d_0=\gcd(\nu,a\kappa)$. Its actual
vanishing column and primitive row are
\[
 l=(\nu/d_0)u-(a\kappa/d_0)\delta,
 \qquad f=(d_0u^tJ,0,b\kappa d_0/\nu),\qquad T-I=lf.
\]
At the selected cusp the image of $T-I$ is $SC_0\mathbb Z^2$,
which equals $S\mathbb Z^2$. Thus the multiplicative angular
columns generate
\begin{equation}
 \Lambda_0=\ker\gamma^*
        =\mathbb Ze_1\oplus\mathbb Ze_2\oplus\mathbb Z\delta.
                                                        \hypertarget{tag-01YH}{}\label{maps:eq:D.C24}
\end{equation}
For an explicit integral verification, order the columns as
$u_{\rm zero},\delta$, followed by the other multiplicative $l$'s
in profile order. The following submatrices $B_{\rm ang}$ in their first three
coordinates have unit determinant. Column indices are zero-based:
\[
\begin{array}{c|c|c|c|r}
\text{profile}&(a,b)&\text{zero}&\text{columns of }B_{\rm ang}&\det B_{\rm ang}\\\hline
\mathrm{OS47}&(7,2)&2&(0,2,3)&-1\\
\mathrm{OS47}&(7,2)&3&(0,1,2)&1\\
\mathrm{OS47}&(7,2)&4&(0,1,5)&-1\\
\mathrm{OS47}/\mathrm{II}&(7,2)&2&(0,2,3)&-1\\
\mathrm{OS49}&(1,6)&2&(0,1,3)&1\\
\mathrm{OS49}&(1,6)&3&(0,1,3)&-1\\
\mathrm{OS56}&(5,2)&1&(0,1,5)&-1\\
\mathrm{OS56}&(15,2)&3&(1,2,5)&1\\
\mathrm{OS56}&(15,2)&4&(0,1,2)&1\\
\mathrm{OS56}/\mathrm{II}&(5,2)&3&(0,1,3)&-1\\
\mathrm{OS56}/\mathrm{IV}&(1,6)&0&(0,1,3)&-1\\
\mathrm{OS56}/\mathrm{IV}&(5,6)&3&(0,1,2)&1\\
\mathrm{OS56}/\mathrm{IV}&(15,2)&3&(0,1,2)&1.
\end{array}
\]
Inserting $(\det B_{\rm ang})\operatorname{adj}B_{\rm ang}$ in those columns gives
the right inverse. Since $f$ is primitive, a B\'ezout column $h$
with $fh=1$ gives $(T-I)h=l$. At the cusp the two long columns
map to $SC_0$. These are inverse preimages of all angular columns.
Every monodromy relation has zero $\gamma$ coordinate, so the sum
of the images of all $T_i-I$ is also exactly $\Lambda_0$.

Filling a multiplicative cap kills its base meridian and angular
lattice. The section and the remaining base loops therefore give
\begin{equation}
 \pi_1(W)=\mathbb Z\gamma\times F_{r_B-1},
                                                        \hypertarget{tag-01YI}{}\label{maps:eq:D.C25}
\end{equation}
where $r_B$ counts the unfilled discs, including the ordinary one.
The maps send $\gamma$ to the first generator and each remaining
additive meridian to its free generator; inverse inclusion of these
loops recovers the group. The ordinary meridian is the inverse of
their product. Each remaining matrix acts identically on the
quotient $\mathbb Z\gamma$. Here $r_B=1$ or $2$.

Let $N$ count all holes before filling, including the ordinary disc.
Then $H_1(V)=\mathbb Z^N$. Each filled boundary has first rank one
greater than its cap, with both groups free. In the degree-one
Mayer--Vietoris map the target rank minus source rank is therefore
$N-(N-r_B)=r_B=b_1(W)$. Its kernel is a rank-zero subgroup of a free abelian
group, hence zero. The primitive cap degree-two boundary maps are
onto by their sections in \ref{maps:src:D.2.26}. Subtracting those
sections and taking the exterior coordinate gives
\begin{equation}
 H_2(V;\mathbb Z)\longrightarrow H_2(W;\mathbb Z)
                         \text{ onto}.                 \hypertarget{tag-01YJ}{}\label{maps:eq:D.C26}
\end{equation}
Inclusion modulo the local kernels is the inverse quotient map.
The free $H_1$ groups and the universal coefficient theorem give
injective restriction in degree-two cohomology, with both integral
and mod-two coefficients. The generators of $H_2(W)$ are the
primitive fibre class $s$ and the unfilled fixed sweeps, subject to
the primitive global $\delta$-sum. Its fibre coefficient is zero
before the ordinary clutching is filled, so these groups are free.

\item \emph{One line with the actual factors, periods and inverse gauges.}
On a multiplicative chart use the physical coordinates $(z,v)$,
elliptic period $h_1$, section factor $h_G$, and contraction factor
$R$. The following maps on an auxiliary line coordinate $\zeta_o$
will be used:
\begin{equation}
\begin{aligned}
 h(z,v,\zeta_o)&=(h_1z,h_G^av,(-1)^az^{-a}\zeta_o),\\
 F(z,v,\zeta_o)&=(h_G^{-b}z,Rv,
                     h_G^{-ab/2}R^{b/2}v^b\zeta_o),\\
 J(z,v,\zeta_o)&=(z^{-1},v^{-1},z^{-a}\zeta_o).
\end{aligned}                                             \hypertarget{tag-01YK}{}\label{maps:eq:D.C27}
\end{equation}
All powers of the nonzero unit functions are integral. Multiplication
gives $hF=Fh$, $J^2=1$, $JFJ=F^{-1}$, and $JhJ=h^{-1}$.
For the last identity the difference in logarithmic line coordinates
is the integer $a$, whose exponential is one.

Here are the global line and its explicit comparison with these
factors. On the real period cover set
\begin{equation}
 \mathcal Q(\lambda)=\lambda_1\lambda_2+\lambda_1+\lambda_2\pmod2.
                                                        \hypertarget{tag-01YL}{}\label{maps:eq:D.C28}
\end{equation}
Its polar form is $\Omega\bmod2$. It is one on each nonzero
elliptic vector modulo two, and the elliptic $\gamma$ column $bg$
is even; hence every actual monodromy preserves $\mathcal Q$. Put
\[
 J_\lambda(x,\zeta)=
 (x+\lambda,(-1)^{\mathcal Q(\lambda)}
                         e^{\pi i\Omega(\lambda,x)}\zeta).
\]
The quadratic identity proves its exact cocycle law. The connection
$d-\pi i\Omega(x,dx)$ descends and has $c_1=\Omega$.
The matrix lifts $(x,\zeta)\mapsto(T_ix,\zeta)$, with inverses
$(x,\zeta)\mapsto(T_i^{-1}x,\zeta)$, preserve the factors and
have product the identity. They construct one line on $V$, with
value one on $s$ and zero on every fixed-loop sweep: its curvature
is vertical and the sweep closes by the scalar-one matrix lift.

For $\tau=\delta/2$, the map
\[
 T_\tau(x,\zeta)=(x+\tau,e^{\pi i\Omega(\tau,x)}\zeta)
\]
commutes with every $J_\lambda$, since
$\Omega(\tau,\Lambda)=-(b/2)\mathbb Z$ is integral. Its square
is $J_\delta$, and for integral $\lambda$ one has
$T_\lambda=(-1)^{\mathcal Q(\lambda)}J_\lambda$. This fixes its
integral-period phase.

For the actual nonzero short period $\omega_s$, use the additive
coordinate $z_E=z_{\rm add}/\omega_s$, with inverse
$z_{\rm add}=\omega_s z_E$. Put $\tau_E=\omega_l/\omega_s$
and choose $\zeta_G=G_{\rm lift}/\omega_s$ for a local lift of $G$.
The multiplicative coordinate is $z=e^{2\pi i z_E}$, with inverse
logarithm modulo the same short period.
The ratio $\vartheta(z_E-a\zeta_G)/\vartheta(z_E)$ has divisor
$[a\zeta_G]-[0]$, short multiplier one and long multiplier
$e^{2\pi ia\zeta_G}$, by reindexing the theta series. It and its
reciprocal identify the divisor line with its Poincar\'e factors.
A forward lift of $-bG$ then has constant fibre multiplier: its
ratio with the constant-factor lift is a holomorphic function on
the compact elliptic curve. Thus the actual long-period matrix is
\[
 \Pi_{\rm phys}=\begin{pmatrix}\tau_E&-b\zeta_G\\a\zeta_G&\rho_0\end{pmatrix},
 \qquad D=\operatorname{diag}(a,-b),
 \qquad D\Pi_{\rm phys}=(D\Pi_{\rm phys})^t.
\]
This includes the complete nonzero unit functions.

In real short/long coordinates $x=S U+L v_{\rm per}$, put
$Z=U+\Pi_{\rm phys}v_{\rm per}$. Multiplication of the line coordinate by
\[
 G_D(U,v_{\rm per})=\exp\!\left(\pi i[-U^tDv_{\rm per}-v_{\rm per}^tD\Pi_{\rm phys}v_{\rm per}+Z_1]\right)
\]
is an invertible gauge, with reciprocal inverse. Its short factors
are one and its long factors are
\[
 (-1)^a h_1^{(1-a)/2}z^{-a},\qquad
 h_G^{-b/2}R^{b/2}v^b.
\]
The inversion of the real period cover becomes
$(z,v,\zeta_p)\mapsto(z^{-1},v^{-1},z^{-1}\zeta_p)$.
The further gauge
\[
 \zeta_o=z^{(a-1)/2}\zeta_p,
 \qquad \zeta_p=z^{-(a-1)/2}\zeta_o
\]
conjugates these three maps exactly to \eqref{maps:eq:D.C27}.
It fixes the rigidifying origin and commutes with the half-$\delta$
lift. In the physical presentation that lift is
$(z,v,\zeta_o)\mapsto(z,-v,\zeta_o)$. The explicit inverse
gauges identify its overlaps with the global period line and prove
the stated pure-fibre and zero-sweep Chern values on $V$.

\item \emph{The prescribed cap extensions.}
In the chosen elliptic marking the actual valuations are
\begin{equation}
 (\operatorname{ord}h_1,\operatorname{ord}h_G,
        \operatorname{ord}R)=(n,j,\rho),
 \qquad P_{\rm val}=
        \begin{pmatrix}n&-bj\\aj&\rho\end{pmatrix}.
                                                        \hypertarget{tag-01YM}{}\label{maps:eq:D.C29}
\end{equation}
The monodromy of the first long period gives $n=\nu,j=-\kappa$;
that of the second gives the displayed last column. At the zero,
$P_{\rm val}=C_0$ and $n\rho+abj^2=1$. At an unselected
multiplicative site its determinant is zero.

In the physical lattice define
\begin{equation}
 \varphi(x,y)=
 \frac{-a\rho x^2-2abjxy+bny^2+ax}{2}.                  \hypertarget{tag-01Y8}{}\label{maps:eq:D.C210}
\end{equation}
At the zero $n,\rho$ are odd and $b$ even, so its values at integer
vertices are integral. Direct expansion gives
\begin{equation}
\begin{aligned}
 \varphi(x+n,y+aj)-\varphi(x,y)&=-ax,\\
 \varphi(x-bj,y+\rho)-\varphi(x,y)
                          &=by+(b\rho-abj)/2.
\end{aligned}                                             \hypertarget{tag-01Y9}{}\label{maps:eq:D.C211}
\end{equation}
These are the two valuations of the factors in \eqref{maps:eq:D.C27}.

To extend on the stipulated fan, take the short basis $SC_0$ and put
\[
 K=C_0^tD=
 \begin{pmatrix}a\nu&ab\kappa\\
            ab\kappa&b(ab\kappa^2-1)/\nu\end{pmatrix},
 \qquad \Pi=C_0^{-1}\Pi_{\rm phys}.
\]
Then $K=K^t$ and $K\Pi=D\Pi_{\rm phys}$ is symmetric. The actual
period matrix has the form
$\Pi=(\log q)/(2\pi i)I+B_{\rm hol}(q)$, retaining its holomorphic unit
matrix. Let $\epsilon_j=\mathcal Q(SC_0e_j)$ and $\eta_j=\mathcal Q(Le_j)$, with
representatives zero or one. In this table $\epsilon=(1,0)^t$.
For $x=SC_0U+Lv_{\rm per}$, $Z=U+\Pi v_{\rm per}$, the gauge
\[
 G_K(U,v_{\rm per})=\exp\!\left(\pi i[-U^tKv_{\rm per}-v_{\rm per}^tK\Pi v_{\rm per}+
                                           \epsilon^tZ]\right)
\]
has reciprocal inverse. Its short factors are one and its long
factor for $h\in\mathbb Z^2$ is
\[
 (-1)^{\eta^th}
 \exp\!\left(-2\pi i h^tKZ-\pi i h^tK\Pi h
                                  +\pi i\epsilon^t\Pi h\right).
\]
Substitution verifies the parameter-monodromy identity
$G_K(U,v_{\rm per},\Pi+I)=G_K(U+v_{\rm per},v_{\rm per},\Pi)$.
The long factor is its exact nonzero unit times
$q^{(-h^tKh+\epsilon^th)/2}X^{-Kh}$.

The integral Cartier function on the Jordan lattice is
\[
 \psi(v)=\tfrac12(v^tKv-\epsilon^tv).
\]
Its integrality follows from $K_{jj}\equiv\epsilon_j\pmod2$,
and its translation differences reproduce the long monomials.
Interpolate on each prescribed unimodular triangle. If $m_\sigma$
is the integral affine interpolation and $e$ the meromorphic frame,
the divisor, regular frame and coordinate maps are
\[
 D_{\mathcal L}=\sum_v\psi(v)D_v,
 \qquad e_\sigma=\chi^{-m_\sigma}e,
 \qquad \zeta_\sigma=\chi^{m_\sigma}\zeta_C,
 \qquad \zeta_C=\chi^{-m_\sigma}\zeta_\sigma.
\]
Translation changes the divisor by the reciprocal long multiplier's
divisor, so its lift is precisely the displayed long factor. All
unit functions remain in that lift and have zero divisor.

The exact comparison with \eqref{maps:eq:D.C210} is
\[
 x=C_0v,\qquad v=C_0^{-1}x,\qquad
 L_0=\tfrac12(a\operatorname{row}_1C_0-\epsilon^t),\qquad
 \varphi(C_0v)=-\psi(v)+L_0v.
\]
$L_0$ is integral. On dense line coordinates the comparison and
inverse are $\zeta_o=X^{L_0}\zeta_C$ and
$\zeta_C=X^{-L_0}\zeta_o$. Thus
\[
 \zeta_\sigma=
 \chi^{m_{\psi,\sigma}}\zeta_C
 =\chi^{-m_{\varphi\circ C_0,\sigma}}\zeta_o.
\]
This proves that the physical support function is interpolated on
the images under $C_0$ of the prescribed triangles. It gives the same
extension and every inverse chart map, with all units retained.

At rank one write the valuation columns as
$H_{\rm val}=\alpha_0\nu_0$, $F_{\rm val}=\beta_0\nu_0$, where
$\nu_0=(\nu_1,\nu_2)$ is primitive and
$\gcd(\alpha_0,\beta_0)=1$. Then
\[
 n=\alpha_0\nu_1,\quad aj=\alpha_0\nu_2,\quad
 -bj=\beta_0\nu_1,\quad\rho=\beta_0\nu_2.
\]
The relation $-a\beta_0\nu_1=b\alpha_0\nu_2$ and coprimality
give the integer $k=a\nu_1/\alpha_0$. The function
$\varphi_1(t)=-kt(t-\alpha_0)/2$ has period differences
$-a\nu_1t$ and $b\nu_2t+(b\rho-abj)/2$.
At every odd-isotropy site $\nu_1$ is odd, so $\beta_0$ is even
and $\alpha_0$ odd. At the unique even-isotropy site
$\nu_1=2,\alpha_0=1$. Thus all these functions are integral.

The common prequantum line has the same extension directly. Choose
primitive short basis $S'=(l,p)$ and long basis $L'=(h,w)$ with
$f(h)=1,f(w)=0$. Put $K_1=S'^t\Omega L'$. Its first row is
$(n_{\rm per},0)$, where $n_{\rm per}=\Omega(l,h)$, and $K_1^t\Pi_1$ is symmetric.
The Gaussian gauge uses
\[
 -U^tK_1v_{\rm per}-v_{\rm per}^tK_1^t\Pi_1v_{\rm per}+\epsilon^t(U+\Pi_1v_{\rm per})
\]
in its exponent. Its long factor is
\[
 (-1)^{\eta^th}
 \exp\!\left(-2\pi i h^tK_1^tZ-\pi i h^tK_1^t\Pi_1h
                                  +\pi i\epsilon^t\Pi_1h\right).
\]
This keeps the full off-diagonal term and unit matrix. The integer
interpolation is $(n_{\rm per}j^2-\mathcal Q(l)j)/2$, since $n_{\rm per}\equiv \mathcal Q(l)\pmod2$.
To compare the two rank-one supports, write $S'=SB_0$, with
first column $\nu_0$ and $\det B_0=1$, and set
$L_1=(a\operatorname{row}_1B_0-\epsilon^t)/2$. This is integral.
The dense coordinate maps are $\zeta_o=X^{L_1}\zeta_1$ and
$\zeta_1=X^{-L_1}\zeta_o$. Here $n_{\rm per}=k$ and
\[
 \varphi_1(t)=-\frac{n_{\rm per}t^2-\epsilon_1t}{2}+(L_1)_1t,
 \qquad n_{\rm per}\alpha_0=a\nu_1.
\]
Thus the regular Cartier coordinates agree after this gauge;
its transverse monomial is invertible on the node charts and is
retained together with the long-period unit factors.
The node charts are $x_+=X/q^j$, $x_-=q^{j+1}/X$, with inverse
$X=x_+q^j$, $q=x_+x_-$. The same Cartier coordinate maps extend
the line through each chart. These extensions all restrict by the
displayed gauge and its inverse to the one line on $V$. They
therefore glue to one equivariant line $\mathcal L$ on $W$.

\item \emph{The actual fixed curves and their characters.}
The physical scalar circle extends through these period and cap
maps; let $\iota$ be its half-period $\delta/2$. At a rank-one
place the surviving elliptic lattice is $\ker f/\mathbb Zl$.
Its half-period vanishes exactly when $\delta-l\in2\Lambda$.
The primitive Jordan projection and insertion give the quotient
map and its inverse. The actual vectors in the thirteen rows give
exactly one such place, an $I_2$ with its local $\kappa=\pm1$.
The node equation $x_+x_-=q$ has both normal coordinates negated,
and its elliptic double curve $E$ is fixed. At every other rank-one
place the surviving half-period is a nonzero elliptic translation,
so gives no fixed point.

At the cusp the physical circle has short coordinate
$d_\delta=C_0^{-1}\delta=(-b\kappa,\nu)$. Its reduction is $(0,1)$.
The adjacent charts at the vertical edge are
\[
 (x,y,w)=(q/(X_1X_2),X_2,X_1),\qquad
 (x',y',w')=(q/X_2,X_1X_2,X_1^{-1}).
\]
Their overlap is $(x',y',w')=(wx,wy,w^{-1})$, with inverse
$(x,y,w)=(w'x',w'y',(w')^{-1})$. The curve $C$ is given by
$x=y=0$ and $x'=y'=0$. Its two normal transitions are those of
$\mathcal O(-1)\oplus\mathcal O(-1)$.
The full physical circle weights at its two endpoints are
\[
 (b\kappa-\nu,\nu,-b\kappa),\qquad
 (-\nu,\nu-b\kappa,b\kappa),
\]
respectively. Their reductions are $(1,1,0)$, so the involution
fixes $C$ and acts negatively on both normal directions.
There is one edge orbit and its endpoints lie in the two different
triangle orbits. A nonzero lattice translation does not stabilize a
finite edge. The horizontal and diagonal directions have different
nonzero reductions from $(0,1)$, so give no further fixed strata.

The half-circle lift in the Jordan Gaussian coordinate has dense
multiplier $\exp(\pi i\epsilon^td_\delta/2)$. On the triangle
$0,e_1,e_2$ the linear Cartier slope is
\[
 m=\bigl((K_{11}-\epsilon_1)/2,
                    (K_{22}-\epsilon_2)/2\bigr).
\]
The forward coordinate therefore has weight
\[
\begin{split}
 c_C&=\epsilon^td_\delta/2+m d_\delta
       =(K_{11}(d_\delta)_1+K_{22}(d_\delta)_2)/2\\
    &=\frac b2\bigl(a\kappa(b\kappa-\nu)-1\bigr)
       =1+\kappa\pmod2.
\end{split}
\]
The adjacent-cone relation gives the degree
$K_{11}-K_{12}=a(\nu-b\kappa)$, with the positive Cartier sign.
Equivalently the regular line coordinates have transition
$\zeta_{\sigma'}=w^{-(K_{11}-K_{12})}\zeta_\sigma$.
Thus
\begin{equation}
 \operatorname{char}_{\iota}(\mathcal L|_C)=1+\kappa,
 \qquad \deg(\mathcal L|_C)=a(\nu-b\kappa)\equiv1\pmod2.
                                                        \hypertarget{tag-01YA}{}\label{maps:eq:D.C212}
\end{equation}
For the other unit-square diagonal the same wall calculation gives
$K_{11}+K_{12}$, differing by $2K_{12}$; the odd-degree assertion
is unchanged. This comparison is an equality of the adjacent-cone
calculations. The stipulated fan remains the one used above.

At the fixed $I_2$ use its local odd $\kappa$. The exact bases are
\[
 l=2u-a\kappa\delta,\quad
 p=u+\frac{1-a\kappa}{2}\delta,\quad
 h=r,\quad w=\gamma-\frac{b\kappa}{2}r.
\]
The short matrix and its inverse are
\[
 \begin{pmatrix}2&1\\-a\kappa&(1-a\kappa)/2\end{pmatrix},
 \qquad
 \begin{pmatrix}(1-a\kappa)/2&-1\\a\kappa&2\end{pmatrix}.
\]
The long matrix is
$\left(\begin{smallmatrix}1&-b\kappa/2\\0&1\end{smallmatrix}\right)$,
with inverse obtained by reversing the upper-right sign. They give
\[
 Th=h+l,\quad Tw=w,\quad
 K_1=\begin{pmatrix}2a&0\\a&-b/2\end{pmatrix},\quad
 \Omega(p,w)=-b/2.
\]
Hence $\deg\mathcal L|E=-b/2$, an odd integer. Since
$\delta=-l+2p$, $\mathcal Q(p)=1$ and the first nodal Cartier slope is
$a$, its half-character is
\begin{equation}
 \operatorname{char}_{\iota}(\mathcal L|_E)
                  =1-a=0\pmod2.                         \hypertarget{tag-01YB}{}\label{maps:eq:D.C213}
\end{equation}
In the physical factors, the alternative node coordinates are
\[
 x=z^{(1-aj)/2}v,\quad y=z^{-aj}v^2;\qquad
 z=x^2/y,\quad v=x^{aj}y^{(1-aj)/2}.
\]
They give short basis $(l,p-l)$, related to $(l,p)$ by
$(x,y)\mapsto(xy,y)$ with inverse division. The long change
$F\mapsto Fh^{bj/2}$ has inverse multiplication by $h^{-bj/2}$.
Its multiplier has zero $x$ exponent and $y$ exponent $b/2$;
the $h$ multiplier is $(-1)^a x^{-2a}y^a$ up to its retained unit.
The support function $-at(t-1)$ has even slopes. Thus this
presentation gives the same zero character with explicit coordinate
and line inverses.

There are no fixed points in the additive or ordinary caps.
For an original additive cap its zero multidegree product description
and contracting factor give a free scalar circle, as in
\ref{maps:src:D.2.26}. At a finite cap write the affine generator
as $x\mapsto Tx+t_0$, where $N_mt_0=v$.
The invariant character gives $\chi(t_0)=n/m$. A half-circle
fixed point would require $jn/m\in\mathbb Z$, $0\le j<m$.
The proved condition $\gcd(m,n)=1$ forces $j=0$, and the
remaining nonintegral half-period cannot fix a point. The base root
gives the same restriction off the central torus. The ordinary
torus cap has primitive $\delta$ and a free half-translation.
Consequently $X^\iota=C\amalg E$ for every parameter in the
statement, including even finite multiplicities.

\item \emph{The integral attachment and the global ordinary class.}
The complete maps of \ref{maps:src:D.2.26} give, at a finite site,
\[
 (bn/q)s+(m/q)D_\delta=0,\qquad D_\delta+be\,s=0,\qquad D_\delta=t(\delta).
\]
The second relation eliminates $D_\delta$ with inverse substitution
$D_\delta=-be\,s$. The first becomes $(b/q)(n-me)s=\mp2s$.
Without a finite site the presentation is $\langle s\mid be\,s=0\rangle$
with $b=2,e=\pm1$. Therefore $H_2(X)=\mathbb Z/2$, generated
by the actual fibre class $s$. The degree-one bridge vanishes by
simple connectivity and the integral cap maps.
The index-one cusp has Euler contribution two; the remaining caps
have contribution zero. Integral duality now gives
\[
 H_0=H_6=\mathbb Z,\quad H_2=H_3=\mathbb Z/2,\quad
 H_1=H_4=H_5=0.
\]
Thus the nonzero class $\alpha$ is specified by $\alpha(s)=1$.

Choose $\alpha_{\rm fin}m+\eta_0n=1$. The full two finite kernel columns
are
\[
 (bn/q)s+(m/q)t(\delta),\qquad
 -\alpha_{\rm fin}br_0s+t(\ell)+\eta_0r_0t(\delta).
\]
Their fibre coefficients are
\begin{equation}
 bn/q=2n,\qquad -\alpha_{\rm fin}br_0,                            \hypertarget{tag-01YC}{}\label{maps:eq:D.C214}
\end{equation}
both even. These columns come from the sector map with area
$-\lambda\wedge v$ and sweep $N_m\lambda$; cutting and rejoining
the same sectors are inverse maps. Changing the peripheral lift
changes $v$ by the actual norm $N_m\lambda$, so all affine lifts
are included. Original additive fixed sweeps bound with zero fibre
coefficient in their product frames.

The ordinary cap is retained by translating the lifted meridian by
$-e\gamma$, with inverse translation by $e\gamma$. Its Thom
kernel is
\[
 t(u)-\Omega(u,-e\gamma)s,\qquad t(-e\gamma),
\]
from $(y,t)\mapsto(y+te\gamma,t)$ and its inverse. Its only
nonzero fibre coefficient is $be$. These full relations imply
$D_\delta=-be\,s$ and
$t(\ell)=br_0(\alpha_{\rm fin}+\eta_0e)s$, both even.
Thus the functional one on $s$ and zero on every sweep descends
to the specified $\alpha$. The line already has exactly these
values on $V$, and \eqref{maps:eq:D.C26} gives
\begin{equation}
                   \alpha|_W=\rho_2c_1(\mathcal L).       \hypertarget{tag-01YD}{}\label{maps:eq:D.C215}
\end{equation}
The free $H_1(W)$ gives the exact universal-coefficient identification.

\item \emph{The global equivariant comparison and the cubic.}
Since $H^2(X;\mathbb Z)=0$, the complex tangent bundle has $c_1=0$,
so $X$ is spin. Also $H^4(X;\mathbb Z)=\mathbb Z/2$ and
$p_1=c_1^2-2c_2=0$ integrally.
Every scalar-circle element is isotopic to the identity. The ordinary
class $\alpha$ survives the circle Borel spectral sequence: its
only outgoing targets are $H^2(BS^1;H^1(X;\mathbb F_2))=0$ and
$H^3(BS^1;\mathbb F_2)=0$. Restrict its extension to
$G=\langle\iota\rangle$ and call it $\widetilde\alpha$.
Equivalently, for $G$ the first obstruction vanishes by $H^1(X)=0$,
and a fixed-point section makes pullback from $BG$ injective,
so the obstruction into its third cohomology also vanishes.

A toric vertex is fixed by the entire circle. Its based circle
isotopy induces the identity on \eqref{maps:eq:D.C25}, and its
section splits the Borel extension. Hence
\[
 \pi_1(W_G)=\bigl(\mathbb Z\times F_{r_B-1}\bigr)\times G.
\]
The first factor is torsion-free, so every fixed-point section sends
the nontrivial element to the same order-two element $(0,1,\iota)$.
Every integral degree-one class from the first factor consequently
restricts to zero at a fixed point. All degree-one mod-two classes
of $W$ lift integrally. The degree-two Borel filtration gives
\[
 \widetilde\alpha|_{W_G}-\rho_2c_1^G(\mathcal L)
        =h\tau^2+\sum_j h_j\lambda_j\tau,
 \qquad \tau\in H^1(BG;\mathbb F_2),
\]
where the $\lambda_j$ are reductions of those integral classes.
The section preserves $\tau^2$ and the mixed terms are products
of actual degree-one group classes, so these representatives survive.
One may also obtain this kernel from the universal-cover sequence:
$H^1$ of the universal cover is zero, so group degree-two
cohomology injects; its ordinary component is zero by
\eqref{maps:eq:D.C215}. This gives the same two types of terms.

At either fixed curve the mixed terms have zero point weight, and
$h\tau^2$ shifts both weights by the same constant. After subtracting
that one global class, the restrictions are
\begin{equation}
 \widetilde\alpha|_C=\alpha|_C+(1+\kappa)\tau^2,
 \qquad \widetilde\alpha|_E=\alpha|_E+\tau\eta,
 \qquad \eta\in H^1(E;\mathbb F_2).                      \hypertarget{tag-01YE}{}\label{maps:eq:D.C216}
\end{equation}
Both ordinary curve degrees are odd, so the sum of their
degree-weight contributions is unchanged by the common shift as well.

The rank-two complex normal bundle has the sign representation on
both fibres. Its equivariant Euler class is
\[
 \tau^4+\tau^2\rho_2c_1(\nu_F)+\rho_2c_2(\nu_F)=\tau^4.
\]
Indeed $c_1(\nu_F)=-c_1(TF)$ is even, and the second class
vanishes on a curve. Remove invariant normal discs. The complement
has free involution, so its Borel space is homotopy equivalent to
its finite-dimensional quotient and its cohomology vanishes after
inverting $\tau$. Thom excision gives the localized restriction
isomorphism, with inverse the sum of the two Gysin maps divided by
their Euler classes. Consequently
\[
\begin{split}
 \langle\alpha^3,[X]_2\rangle
 &=\int_C\frac{(\alpha|_C+(1+\kappa)\tau^2)^3}{\tau^4}
   +\int_E\frac{(\alpha|_E+\tau\eta)^3}{\tau^4}\\
 &=(1+\kappa)\int_C\alpha|_C=1+\kappa.
\end{split}
\]
Every degree-one class on the oriented elliptic curve lifts integrally
and has square zero, so its displayed contribution is zero.
Equivariant integration in total degree six takes values in
$H^0(BG;\mathbb F_2)$ and restricts to the ordinary evaluation.

For example, OS47 with $(a,b)=(7,2)$ and zero position $3$ has
\[
 K=\begin{pmatrix}7&14\\14&26\end{pmatrix},\qquad
 d_\delta=(-2,1),\qquad m=(3,13).
\]
Its dense half-phase exponent is $-1$ and its Cartier exponent is
$7$, giving weight zero. Its degree on $C$ is $-7$, its endpoint
circle weights are $(1,1,-2)$ and $(-1,-1,2)$, and its elliptic
degree and weight are $-1$ and zero. Taking $e=1$ gives a simply
connected member with cubic zero. Every row in
\eqref{maps:eq:D.C22} follows from the same integral formula.

\item \emph{The two framed models and the oriented recognition.}
The sphere defining $M_0$ represents $2A$ in $Z_0=S^2\times S^4$.
Its oriented normal rank-four bundle has $w_2=0$, since both
$Z_0$ and the sphere are spin. Its clutching loop in $SO(4)$ is
therefore null-homotopic, and $\pi_2(SO(4))=0$ gives the unique
homotopy class of normal framing. Choose its oriented tubular
embedding and replace $S^2\times D^4$ by $D^3\times S^3$ using
the identity on the framed boundary with the surgery-trace
orientations.

The complement is simply connected: a loop and a disc null-homotopy
can be moved off the codimension-four sphere and then out of its
tubular interior by radial retraction. Thom excision sends the
positive $H_4(Z_0)$ generator to twice the normal-disc class and
the top fundamental class to the top Thom class with degree one.
Thus the complement has $H_2=\mathbb Z$, $H_3=\mathbb Z/2$,
and all its other positive reduced groups below six vanish.
The boundary $H_2$ map is multiplication by two, and its $H_3$
map to the two pieces is $j\mapsto([j]_2,j)$. The latter
cokernel is identified with $\mathbb Z/2$ by
$([u]_2,v)\mapsto[u-v]_2$, with inverse
$[u]_2\mapsto([u]_2,0)$. Mayer--Vietoris and van Kampen give
the stated homology and simple connectivity of $M_0$.

The index-three trace $W_0$ has relative boundary
$\mathbb Z\xrightarrow{2}\mathbb ZA$. Hence
$H_2(W_0)=\mathbb Z/2$, $H_3(W_0)=0$, and the incoming
degree-two map is reduction modulo two. Viewed from the outgoing
end its relative handle has index four, so the outgoing $H_2$
map is an isomorphism. The common degree-two mod-two class
restricts to the reduction of the $S^2$ class on the incoming
boundary. Its square there is zero; evaluation on
$\partial[W_0]_2$ gives cubic zero on $M_0$.
The unique spin structure of the attaching sphere bounds the handle
core, so the spin structure extends. The incoming degree-four
homology map is an isomorphism and $H_3(W_0)=0$, giving an
integral degree-four cohomology isomorphism. Since $TZ_0$ is
stably trivial, tangent restriction gives $p_1(W_0)=0$, hence
$p_1(M_0)=0$.

The conic construction of \ref{maps:src:D.2.8} has the same trace
maps starting from $\mathbb CP^3$. Its conic represents twice the
line class and has trivial oriented real normal bundle, with the
same unique framing class. The incoming hyperplane cube is one.
The trace consequently gives cubic one, simple connectivity, spin,
and $H_2=H_3=\mathbb Z/2$ for $M_1$. Its degree-four cohomology
restricts isomorphically to the incoming group; $p_1$ there is
$4h^2$, whose outgoing restriction is zero in $\mathbb Z/2$.

The exact recognition input is \cite[Theorem 5.2, p.~42]{manAndrus77},
with the torsion-manifold definition on p.~17 and augmented-ring
definition on pp.~19--21. It applies to closed oriented simply
connected spin six-manifolds with $H_2=H_3=\mathbb Z/n$, for
$4\nmid n$; the additional Pontryagin cubing operation is needed
only when $3\mid n$. At $n=2$ the augmented even ring has basis
$(1,\alpha,d,z)$ with
\[
 \alpha d=z,\qquad \alpha^2=\varepsilon d,\qquad
 \overline p_1=0,
\]
where $z$ is the integral-orientation reduction. Sending these four
classes to the corresponding classes of $M_\varepsilon$, and back
for the inverse, preserves the ring, $\overline p_1$ and orientation.
The theorem gives the orientation-preserving diffeomorphism in
\eqref{maps:eq:D.C21}. The two models have different cubes of their
unique nonzero degree-two classes and therefore are not diffeomorphic.

The full spin-bordism description gives the same classification
with an explicit conversion. By \cite[Theorem 2.10, p.~803,
and \S3.12, p.~807]{manZhubr75}, at order two
\[
 \Omega^{\mathrm{Spin}}_6(K(\mathbb Z/2,2))
   =\mathbb Z\langle\epsilon,\zeta\rangle/
                        \langle\epsilon,2\zeta\rangle.
\]
The quotient map sends $(A,B)$ to $B\bmod2$, and the inverse sends
one to the reduction of
$\zeta=[\mathbb CP^3\longrightarrow\mathbb CP^\infty]$.
The universal mod-two degree-two class pulls back to the hyperplane
class, whose cube is one. Thus cubic evaluation is the inverse
detector. The marked degree-positive recognition is the map constructed in
\cite[Theorem 3, p.~795, and its proof in \S4.1, p.~809]{manZhubr75}.
That proof preserves $H_2\cong G$ and kills $H_3/\operatorname{Tors}H_3$;
its vanishing-free-third-homology scope applies here, and
$\operatorname{Aut}(\mathbb Z/2)=1$.
The augmented ring above maps to $\varepsilon$ times that reduced
bordism generator; the inverse is its cubic and the displayed ring.
The two surgery traces represent its two values, proving the stated
conversion of the complete invariants.
\end{enumerate}
\end{proof}



\subsubsection*{Orders five and seven}

\begin{definition}[Minimal odd-primary data]\hypertarget{tag-Q206}{}\label{maps:def:minimal-odd-data}
{\normalfont\scriptsize tag Q206}\par
Use the ordered profiles of Definition~\ref{mapsIII:def:elliptic-marked-data}
and the natural initial pairs of \ref{maps:src:D.2.22}. Choose the
minimal additive fillings of \ref{maps:src:D.2.26} with residual
$c=5$ or $7$ and trivial orbifold obstruction. Use the primitive
Mumford cap at every multiplicative position and its prescribed fan
at the selected zero. The logarithmic sites are exactly the chosen
additive sites. At every finite site allow each positive multiple
of its monodromy order and every free affine twist; retain the
ordinary clutch $e\gamma$ of \eqref{maps:eq:D.HY3}.
Put $p=c$, let $I_\nu$ be the selected zero, and set $Q_f=b/p$.
Selected positions are $i-1$ in the input definition's peripheral
order. In the tables, $O$ and $F$ denote the original and finite
additive caps, and a dash denotes an absent additive site.
These choices are the \emph{minimal odd-primary data}.
\end{definition}

\begin{theorem}[Five- and seven-primary smooth types]\hypertarget{tag-02LC}{}\label{maps:src:D.2.36}
{\normalfont\scriptsize tag 02LC}\par
The data of Definition~\ref{maps:def:minimal-odd-data} give $36$
patterns on $30$ marked initial rows. Every simply connected threefold
$X$ has, in its complex orientation, the following complete marked
data, where $s$ is the primitive fibre generator of
\ref{maps:src:D.2.26}:
\begin{equation}
\begin{gathered}
 H_2(X;\mathbb Z)=H_3(X;\mathbb Z)=\mathbb Z/p,
 \qquad b_3(X)=0,\qquad w_2(X)=0,\qquad \mu(X)=0,\\
 \operatorname{PD}p_1(X)=2(Q_f-a\nu)s,\qquad
 \mathfrak p(X)=P_Xs,\qquad
 P_X=\frac{Q_f-a\nu}{2}\in\mathbb F_p,\qquad
 4\mathfrak p(X)=\operatorname{PD}p_1(X).
\end{gathered}
 \hypertarget{tag-022C}{}\label{maps:eq:D.OP1}
\end{equation}
There is exactly one unselected multiplicative position whose double
curve is fixed by the physical scalar subgroup $C_p$; it has type
$I_p$. In its positive-valuation Tate marking write
\[
 (\operatorname{ord}h_1,\operatorname{ord}h_G,
       \operatorname{ord}R)=(p,r,\rho),\qquad
 w=(ar)^{-1}\in\mathbb F_p.
\]
Then
\[
 P_X=\frac{Q_f}{2}\left(1-\frac{\nu^2}{w^2}\right),
 \qquad \mathfrak p(X)=0\quad\Longleftrightarrow\quad w^2=\nu^2.
\]
The cusp normal weights for this scalar generator are $(-\nu,\nu)$.
Precomposing the action with $\zeta\mapsto\zeta^{\nu^{-1}}$, whose
inverse is $\zeta\mapsto\zeta^\nu$, changes the elliptic normal
weights to $(w/\nu,-w/\nu)$ and the cusp weights to $(-1,1)$.

Among all $36$ patterns, twelve have $(p,P_X)=(5,0)$, twelve have
$p=5$ and $P_X\ne0$, and twelve have $p=7$ and $P_X\ne0$.
There are exactly two oriented diffeomorphism types of order five
and one of order seven. For a nonzero $P_X$, the marking
$1\mapsto P_Xs$, with inverse evaluation $P_X^{-1}\alpha$ where
$\alpha(s)=1$, gives refined coordinate $1$ and full fundamental-class
coordinate $0$. The zero case uses $1\mapsto s$ and inverse $\alpha$.
Thus the three types are specified by
$(H_2,b_3,\mathfrak p,\mu)=(\mathbb Z/5,0,0,0)$,
$(\mathbb Z/5,0,1,0)$, and $(\mathbb Z/7,0,1,0)$.
\end{theorem}


\begin{proof}
\begin{enumerate}[label={\arabic*.},leftmargin=*]
\item \emph{The finite list and its parameter range.}
Use the ordered elliptic matrices and section cocycles of
Definition~\ref{mapsIII:def:elliptic-marked-data}. At a multiplicative position choose
\[
 A_i-I=\nu_i u_i u_i^tJ_0,\qquad g_i=\kappa_i u_i,
 \qquad J_0=\begin{pmatrix}0&1\\-1&0\end{pmatrix},
 \qquad\det(u_i,v_i)=1,
\]
with $u_i$ primitive and its first nonzero entry positive.
The content of $A_i-I$ is $\nu_i$, a primitive column determines
$u_i$, and a Bezout identity determines $v_i$. These rules include
$I_1$ and $\kappa_i=0$.

For height $1/d$, the natural-row criterion at a selected $I_\nu$
is
\[
 ab=d/\nu,\quad a,b>0,\quad\gcd(a,b)=1,
 \quad\gcd(b\kappa,\nu)=1,
\]
together with the component equation
\[
 [j=a\kappa]-[j=0]-[j=(a-b)\kappa]+[j=-b\kappa]
       =-2[j=0]+[j=1]+[j=-1]
\]
on $\mathbb Z/\nu$, and narrowness of at least one of $P,Q$ at
every other reducible position. Repeated terms in this equation
are added, including when $\nu=1$ or $2$.
Enumerate the divisors $b$ of $d/\nu$, put $a=d/(\nu b)$, and
apply these component equations. At an additive site a minimal
original option requires $P$ narrow and $Q$ transitive on the
components. A finite option requires
$b(A_i-I)^{-1}g_i\in\mathbb Z^2$. Its vertical index $q_i$ and
linear order $d_i$ are
\[
\begin{array}{c|ccccc}
 &\mathrm{II}&\mathrm{III}&\mathrm{IV}&\mathrm{III}^*&\mathrm{IV}^*\\\hline
q_i&1&2&3&2&3\\
d_i&6&4&3&4&3 .
\end{array}
\]
These are the complete minimal options of
Theorem~\ref{maps:src:D.2.26}; its component and bridge defects exclude
the other original choices. Remove an option product with nontrivial
orbifold obstruction and compute $c=b/\prod q_i$.

The height denominators in the defining profiles are
$2,8,4,14,6,12,20,30$. Since $p\mid b\mid d$, only the
families $47,55,56$, of heights $1/14,1/20,1/30$, enter here.
Their ordered profiles are
\[
\begin{array}{c|l}
\mathrm{OS47}&I_7,I_2,I_1,I_1,I_1\\
\mathrm{OS47}/\mathrm{II}&I_7,I_2,I_1,\mathrm{II}\\
\mathrm{OS47}/\mathrm{III}&I_7,I_1,I_1,\mathrm{III}\\
\mathrm{OS55}&I_5,I_4,I_1,I_1,I_1\\
\mathrm{OS55}/\mathrm{II}&I_5,I_4,\mathrm{II},I_1\\
\mathrm{OS56}&I_5,I_3,I_2,I_1,I_1\\
\mathrm{OS56}/\mathrm{II}&I_2,\mathrm{II},I_5,I_3\\
\mathrm{OS56}/\mathrm{III}&I_5,I_3,\mathrm{III},I_1\\
\mathrm{OS56}/\mathrm{IV}&I_5,I_2,\mathrm{IV},I_1 .
\end{array}
\]
The divisor and component calculation gives the complete table below.
The last column is computed by the Borel comparison and Pontryagin
evaluation below. The cap labels $\mathrm O$ and $\mathrm F$
mean original and finite choices, respectively; those choices
at the same site are separate patterns.
\[
\begin{array}{r|c|c|r|c|r|r|r|r}
 &\text{profile}&(a,b)&\text{zero}&\text{cap}&p&\nu&w&P_X\\\hline
1&\mathrm{OS47}&(1,7)&1&-&7&2&4&3\\
2&\mathrm{OS47}&(2,7)&2&-&7&1&2&3\\
3&\mathrm{OS47}&(2,7)&3&-&7&1&2&3\\
4&\mathrm{OS47}&(2,7)&4&-&7&1&2&3\\
5&\mathrm{OS47}/\mathrm{II}&(1,7)&1&\mathrm O&7&2&4&3\\
6&\mathrm{OS47}/\mathrm{II}&(1,7)&1&\mathrm F&7&2&4&3\\
7&\mathrm{OS47}/\mathrm{II}&(2,7)&2&\mathrm O&7&1&2&3\\
8&\mathrm{OS47}/\mathrm{II}&(2,7)&2&\mathrm F&7&1&2&3\\
9&\mathrm{OS47}/\mathrm{III}&(2,7)&1&\mathrm O&7&1&2&3\\
10&\mathrm{OS47}/\mathrm{III}&(1,14)&1&\mathrm F&7&1&4&4\\
11&\mathrm{OS47}/\mathrm{III}&(2,7)&2&\mathrm O&7&1&2&3\\
12&\mathrm{OS47}/\mathrm{III}&(1,14)&2&\mathrm F&7&1&4&4
\end{array}
\]
\[
\begin{array}{r|c|c|r|c|r|r|r|r}
 &\text{profile}&(a,b)&\text{zero}&\text{cap}&p&\nu&w&P_X\\\hline
13&\mathrm{OS55}&(1,5)&1&-&5&4&2&1\\
14&\mathrm{OS55}&(4,5)&2&-&5&1&3&1\\
15&\mathrm{OS55}&(4,5)&3&-&5&1&3&1\\
16&\mathrm{OS55}&(4,5)&4&-&5&1&3&1\\
17&\mathrm{OS55}/\mathrm{II}&(1,5)&1&\mathrm O&5&4&2&1\\
18&\mathrm{OS55}/\mathrm{II}&(1,5)&1&\mathrm F&5&4&2&1\\
19&\mathrm{OS55}/\mathrm{II}&(4,5)&3&\mathrm O&5&1&3&1\\
20&\mathrm{OS55}/\mathrm{II}&(4,5)&3&\mathrm F&5&1&3&1\\
21&\mathrm{OS56}&(2,5)&1&-&5&3&2&0\\
22&\mathrm{OS56}&(3,5)&2&-&5&2&3&0\\
23&\mathrm{OS56}&(6,5)&3&-&5&1&4&0\\
24&\mathrm{OS56}&(6,5)&4&-&5&1&4&0
\end{array}
\]
\[
\begin{array}{r|c|c|r|c|r|r|r|r}
 &\text{profile}&(a,b)&\text{zero}&\text{cap}&p&\nu&w&P_X\\\hline
25&\mathrm{OS56}/\mathrm{II}&(3,5)&0&\mathrm O&5&2&3&0\\
26&\mathrm{OS56}/\mathrm{II}&(3,5)&0&\mathrm F&5&2&3&0\\
27&\mathrm{OS56}/\mathrm{II}&(2,5)&3&\mathrm O&5&3&2&0\\
28&\mathrm{OS56}/\mathrm{II}&(2,5)&3&\mathrm F&5&3&2&0\\
29&\mathrm{OS56}/\mathrm{III}&(2,5)&1&\mathrm O&5&3&2&0\\
30&\mathrm{OS56}/\mathrm{III}&(1,10)&1&\mathrm F&5&3&4&2\\
31&\mathrm{OS56}/\mathrm{III}&(6,5)&3&\mathrm O&5&1&4&0\\
32&\mathrm{OS56}/\mathrm{III}&(3,10)&3&\mathrm F&5&1&3&2\\
33&\mathrm{OS56}/\mathrm{IV}&(3,5)&1&\mathrm O&5&2&3&0\\
34&\mathrm{OS56}/\mathrm{IV}&(1,15)&1&\mathrm F&5&2&4&3\\
35&\mathrm{OS56}/\mathrm{IV}&(6,5)&3&\mathrm O&5&1&4&0\\
36&\mathrm{OS56}/\mathrm{IV}&(2,15)&3&\mathrm F&5&1&2&3
\end{array}
\]
There are $36$ patterns and $30$ initial rows; each has at most one
additive site. The enumeration is exhaustive because the positive
factor pairs are bounded and the displayed site options are complete.

\item \emph{The marked fibre and all finite parameters.}
The positive forward period $\gamma$ covers translation by $-Q$.
Let $\widehat s$ denote the primitive exterior fibre class, with image
$s$ after filling. On the physical lattice put
\begin{equation}
 \Lambda=\mathbb Z(e_1,e_2,\delta,\gamma),\qquad
 \xi=a e_1^*\wedge e_2^*-b\delta^*\wedge\gamma^*,
 \qquad \xi(\widehat s)=1.
 \hypertarget{tag-022F}{}\label{maps:eq:D.OP2}
\end{equation}
At a finite site take $m=d_i k>0$ and
$v=n\ell+r_0\delta$, with $\chi(\ell)=1$ and
$\chi=\gamma^*$. An affine lift $R_{\rm aff}(x)=T_i x+t_0$ has
$N_m t_0=v$, where $N_m=I+T_i+\cdots+T_i^{m-1}$.
Its sector construction identifies the boundary with the power $v$.
The ordinary clutching integer $e$ gives the presentation
\[
 \langle\gamma,c\mid[\gamma,c]=1, c^m=\gamma^n,
                                   c=\gamma^e\rangle
       \simeq\mathbb Z/(me-n).
\]
The forward substitution is $c\mapsto\gamma^e$ and its inverse
retains $\gamma$. Thus simple connectivity is exactly $me-n=\pm1$,
which implies $\gcd(m,n)=1$. Every $k>0$ occurs, for example with
$e=1,n=m-1$, and any $r_0$. Without a finite site the group is
$\mathbb Z/e$, so $e=\pm1$.

The integral specialization maps of
Theorem~\ref{maps:src:D.2.26} leave a fibre generator $s$ and
$D_\delta=t(\delta)$, with relations
\[
 (bn/q_i)s+(m/q_i)D_\delta=0,\qquad D_\delta+be\,s=0.
\]
The substitution $D_\delta=-be\,s$, with its evident inverse inclusion,
gives $(b/q_i)(n-me)s=0$. Here $b/q_i=p$.
Without a finite site the relation is $be\,s=0$ with $b=p$.
The cap maps in degree two are surjective: lift a cap coordinate
by its boundary section, subtract it, and read the remaining
exterior coordinate. The inverse includes the exterior with zero
cap coordinates. Changes of section lie in the full local kernel.
The next Mayer--Vietoris bridge has rank $b_1(X)$ and is free;
it vanishes in the simply connected range. This proves
$H_2(X)=\mathbb Z/p\langle s\rangle$ integrally.

The selected cusp contributes two to the Euler characteristic and
the other caps contribute zero, by Theorem~\ref{maps:src:D.2.3}.
Poincare duality and the universal coefficient sequence give
\[
 H_0=H_6=\mathbb Z,\quad H_2=H_3=\mathbb Z/p,\quad
 H_1=H_4=H_5=0.
\]
In particular $b_3=0$. The class
$\alpha\in H^2(X;\mathbb F_p)$ is uniquely characterized by
$\alpha(s)=1$. Since $p$ is odd, $H^2(X;\mathbb Z/2)=0$,
so $X$ is spin with its unique spin structure. Also
$H^2(X;\mathbb Z)=0$, hence $c_1(TX)=0$.

\item \emph{The complement and its integral restriction maps.}
Remove the additive caps and one existing ordinary base disc, and
call the remaining manifold $Y$. It contains all the multiplicative
caps. Move the ordinary clutching to that boundary by the torus map
\[
 (y,t)\longmapsto(y+te\gamma,t),\qquad
 (y,t)\longmapsto(y-te\gamma,t)
\]
as inverse. Its normal vector is $-e\gamma$, and its full Thom
kernel is
\[
 t(u)-\xi(u,-e\gamma)s,\qquad t(-e\gamma),\qquad u\in\Lambda.
\]
Every fibre coefficient is divisible by $b$, hence by $p$.
At a finite site choose $\alpha_0m+\eta_0n=1$.
Before eliminating a sweep, the full normal kernel is generated by
\[
 (bn/q_i)s+(m/q_i)t(\delta),\qquad
 -\alpha_0br_0s+t(\ell)+\eta_0r_0t(\delta).
\]
The sector swept through $\lambda$ has area $-\lambda\wedge v$
and swept boundary vector $N_m\lambda$; rejoining its sectors is
the inverse map. These formulas give
\[
 t(\delta)=-be\,s,\qquad
 t(\ell)=br_0(\alpha_0+\eta_0e)s.
\]
They both vanish modulo $p$. At an original minimal additive
cap both fixed sweeps bound, by its boundary-to-cap map.

Replacing a finite peripheral generator by its translate through
$\lambda$ changes its power to $v+N_m\lambda$ and its fixed sweep by
\[
 t'(u)\longmapsto t(u)+\xi(u,\lambda)s;
\]
subtraction is the inverse. The changes of the gamma numerators are
$n'=n+m\chi(\lambda)$ and $e'=e+\chi(\lambda)$, which preserve
$me-n$. In the present profiles the fixed vector $\ell$ has
elliptic coordinates divisible by $p$: the denominators of
$(A_i-I)^{-1}g_i$ are prime to $p$, and multiplication by
$b=pQ_f$ makes those coordinates integral. Thus the added fibre
coefficient is divisible by $p$ for $u=\delta,\ell$. This proves
the restriction statements for every peripheral lift.

The exterior degree-two splitting is
\[
 \mathbb Zs\oplus
 \left(\bigoplus_i\operatorname{Fix}(T_i)\right)
            \big/\left\langle\sum_i t_i(\delta)\right\rangle.
\]
Its inverse projects to the graph kernel, applies the integral
inverse of the boundary injection, subtracts those section sweeps,
and evaluates the remaining fibre class by $\xi$. The single
relation is the boundary of the physical delta circle over the
punctured base, with zero fibre coefficient in this framing.
Filling the multiplicative caps removes their fixed sweeps through
the same surjective maps and sections. The next bridge is zero
because their base meridians are independent while the ordinary
disc is unfilled. Hence $H_2(Y)$ is free and generated by $s$ and
the remaining section sweeps, with their delta-sum relation.
The restriction of $\alpha$ is one on $s$ and zero on all these
sweeps. These generators will also identify the class of a line
on $Y$.

\item \emph{Period coordinates, fan coordinates, and their inverses.}
The forward matrices in the order of \eqref{maps:eq:D.OP2} are
\[
 T_i=\begin{pmatrix}
 A_i&0&b g_i\\
 \beta_i&1&\beta_i b z_i+\omega_i\\
 0&0&1
 \end{pmatrix},\qquad
 \beta_i=-a g_i^tJ_0A_i,
\]
where $z_i=(A_i-I)^{-1}g_i$ at a finite place and
$z_i=(\kappa_i/\nu_i)v_i$ at a multiplicative place.
Only the selected zero has $\omega_i\ne0$, with value $1/\nu$.
The matrices have integral entries, product identity, and preserve
$\xi$ and $\delta$.

In a Tate chart use a positive short period and a positive long
period with determinant one. For its short period $\omega_s$ and
long period $\omega_l$, put $\tau_E=\omega_l/\omega_s$ and
$\zeta_G=G_{\rm lift}/\omega_s$. Let $\rho_L$ be the logarithm
of the forward line multiplier in this period coordinate. The physical logarithmic columns are
\[
 \Omega_{\rm ph}=
 \begin{pmatrix}\tau_E&-b\zeta_G\\a\zeta_G&\rho_L\end{pmatrix}.
\]
For the additive coordinate $z_E=z_{\rm add}/\omega_s$, division
by the nonzero short period has inverse multiplication by $\omega_s$;
the physical coordinate is $z=e^{2\pi i z_E}$.
The ratio $\theta(z_E-a\zeta_G)/\theta(z_E)$ representing
$\mathcal O(aG-O)$ gains $e^{2\pi i a\zeta_G}$ under the
long elliptic period. Its reciprocal is the inverse line map.
Continuation of $\zeta_G$ adds $-\kappa_i u_i$, so the two
off-diagonal columns agree with $T_i$, including the $-Q$ arrow.
Consequently
\[
 \operatorname{ord}h_1=\nu_i,\qquad
 \operatorname{ord}h_G=-\kappa_i,\qquad
 \operatorname{ord}R=\omega_i-ab\kappa_i^2/\nu_i.
\]
The nonzero holomorphic unit factors remain in these periods,
with their logarithms on the base disc.

At the cusp put $S=(u,\delta)$, $L=(v,\gamma)$ and
\[
 C_0=\begin{pmatrix}
 \nu&b\kappa\\-a\kappa&(1-ab\kappa^2)/\nu
 \end{pmatrix},\qquad
 C_0^{-1}=\begin{pmatrix}
 (1-ab\kappa^2)/\nu&-b\kappa\\a\kappa&\nu
 \end{pmatrix}.
\]
Both are integral and $\det C_0=1$. The ordered periods
$(SC_0e_1,Le_1,SC_0e_2,Le_2)$ give $I+E_{12}+E_{34}$.
The map of period quotients is induced by the identity on the
complex universal cover and this integral basis change; the inverse
uses the inverse basis matrix. Applying $C_0^{-1}$ to the period coordinates gives
\[
 Z=U+\Omega W,\qquad
 \Omega=C_0^{-1}\Omega_{\rm ph}
       =\frac{\log q}{2\pi i}I+B(q).
\]
All the original units are retained in $B(q)$.
The prescribed antidiagonal fan has the chart and inverse
\[
 (x,y,z)=\bigl(q/(X_1X_2),X_2,X_1\bigr),\qquad
 (X_1,X_2,q)=(z,y,xyz).
\]
At a nodal rank-one place the corresponding charts are
\[
 X=X_1/q^j,\qquad Y=q^{j+1}/X_1,\qquad XY=q,\qquad
 (X_1,q)=(Xq^j,XY).
\]
Translations by the long periods, with their units, identify the
charts. These maps retain the original outer collars.

\item \emph{The physical scalar action and its full fixed set.}
Translation by $\mathbb R\delta/\mathbb Z\delta$ commutes with
the line multipliers, every forward arrow, every affine finite
arrow, and the ordinary clutching maps. The displayed toric charts
extend it to $X$. Specify the generator of its subgroup $C_p$ to
be $\tau=\delta/p$, and put $\zeta_p=e^{2\pi i/p}$.

At an unselected multiplicative place set
\[
 d_0=\gcd(\nu_i,a\kappa_i),\quad n_i=\nu_i/d_0,\quad
 A'_i=a\kappa_i/d_0,\quad l_i=n_i u_i-A'_i\delta.
\]
Then
\[
 T_i-I=l_i\bigl(d_0u_i^tJ_0, 0, b\kappa_i d_0/\nu_i\bigr).
\]
Its row and column are primitive for these minimal caps.
Choose a short complement $k_i$ with
$\delta=-\sigma_i l_i+n_i k_i$, in the nodal basis constructed below.
On the elliptic double curve, translation by $\tau$ is the
class of $(n_i/p)k_i$ in the surviving primitive elliptic lattice.
It is zero exactly when $p\mid n_i$. A nonzero translation of
that elliptic curve has no fixed point. Applying the formula to
the complete table gives exactly one such site, of type $I_p$;
there $d_0=1$ and $n_i=p$. Its double curve $E$ is embedded:
the two sections in the normalization are distinct and form the
two branches of the nodal chart, and no nonzero period translation
stabilizes one adjacent-chart pair.

At the cusp the scalar vector in short coordinates is
\[
 d_\delta=C_0^{-1}\binom01=\binom{-b\kappa}{\nu}.
\]
Its reduction modulo $p$ is nonzero and vertical. Only the vertical
edge direction of the prescribed triangular fan is parallel to it
modulo $p$. The vertical edges have one translation orbit, giving
an embedded $C\simeq\mathbb P^1$. Every triangle has such an edge;
its fixed vertex is an endpoint of $C$, so there is no further
isolated fixed point. A translation stabilizing an edge would fix
its two endpoints and is zero. The two triangle orbits give the
two different endpoints of $C$. In the cusp chart the action is
\[
 (x,y,z)\longmapsto
       (\zeta_p^{-\nu}x,\zeta_p^\nu y,z).
\]
The two normal characters are $-\nu,\nu$, and the adjacent chart
transition gives normal degrees $-1,-1$.

An original additive cap has a strictly contracting forward arrow
on its nonzero line bundle. A unit-modulus scalar can agree with
an iterate of that arrow only for the zero iterate; a nontrivial
scalar then fixes no nonzero line vector. In a finite cap, a
putative fixed point satisfies $x+\delta/p=R_{\rm aff}^{\,j}x$ modulo the
covering torus lattice. Applying the integral covector $\chi$
gives $jn/m\in\mathbb Z$. Since $\gcd(m,n)=1$, this forces
$m\mid j$, and a nonzero torus translation cannot fix a point.
This includes finite orders divisible by $p$. The ordinary cap
is free for the same torus-translation reason. Thus
\begin{equation}
                       X^{C_p}=C\sqcup E.
 \hypertarget{tag-022L}{}\label{maps:eq:D.OP7}
\end{equation}

\item \emph{One equivariant line, with explicit comparison maps.}
On the real period cover define
\[
 J_\lambda(x,\zeta_R)=
   (x+\lambda,e^{2\pi i\xi(\lambda,x)}\zeta_R),\qquad\lambda\in\Lambda.
\]
The quotient of $J_\lambda J_\mu$ by $J_{\lambda+\mu}$ is
$e^{2\pi i\xi(\lambda,\mu)}=1$, and the inverse is
$J_{-\lambda}$. The connection
$d-2\pi i\xi(x,dx)$ has first Chern class $2\xi$ in the
positive complex-line convention. Each monodromy lifts by
$(x,\zeta_R)\mapsto(T_i x,\zeta_R)$, with inverse $T_i^{-1}$, because
$T_i^*\xi=\xi$. Choose the origin framing using the zero elliptic
section and a nonzero frame of its line over the punctured base;
the latter exists since that base has zero second cohomology.
The ordinary clutching is on the removed disc. These maps define
one line on the regular part of $Y$.

The scalar generator lifts by
\[
 \mathcal T_\tau(x,\zeta_R)=
   (x+\tau,e^{2\pi i\xi(\tau,x)}\zeta_R),\qquad\tau=\delta/p.
\]
Its commutator with $J_\lambda$ is
$e^{4\pi i\xi(\tau,\lambda)}=1$, since
$\xi(\tau,\Lambda)=-(b/p)\mathbb Z$.
Moreover $\mathcal T_\tau^p=J_\delta$, and the monodromy lifts
commute with it. The inverse uses $-\tau$.
On a constant fixed-loop sweep $x=tu$, the connection form is
zero, and the base gluing adds no multiplier. The line has degree
zero on every section sweep and degree two on $s$.

Here is also a holomorphic description retaining the physical
period functions. Put $D=\operatorname{diag}(a,-b)$ and
$Z=U+\Omega_{\rm ph}W$. The gauge
\[
 G_0=\exp\{-2\pi i(U^tDW+W^tD\Omega_{\rm ph}W)\}
\]
has reciprocal as inverse. Since $D\Omega_{\rm ph}$ is symmetric,
its short factors are one and its long factors are
$h_1^{-a}z^{-2a}$ and $R^b v^{2b}$. Inversion has factor one.
The additional map $(z,v,\zeta_0)\mapsto(z,v,z^a\zeta_0)$,
with inverse multiplication by $z^{-a}$, gives
\begin{equation}
\begin{aligned}
 h(z,v,\zeta)&=(h_1z,h_G^a v,z^{-2a}\zeta),\\
 F(z,v,\zeta)&=(h_G^{-b}z,Rv,h_G^{-ab}R^b v^{2b}\zeta),\\
 J(z,v,\zeta)&=(z^{-1},v^{-1},z^{-2a}\zeta).
\end{aligned}
 \hypertarget{tag-022H}{}\label{maps:eq:D.OP4}
\end{equation}
The inverses of the first two maps have auxiliary factors
$h_1^{-2a}z^{2a}$ and $h_G^{ab}R^b v^{-2b}$, respectively,
and $J^{-1}=J$. Substitution gives $hF=Fh$, $JhJ=h^{-1}$,
and $JFJ=F^{-1}$. The gauge $z^a$ is scalar invariant.
On an overlap, compose these gauges with the integral period-basis
map and the inverse gauge in the other chart. Their triple
composites cancel exactly. Thus \eqref{maps:eq:D.OP4} describes the
same global equivariant line, including all swept-cycle values.
A rigidified holomorphic line automorphism on a compact connected
torus is the identity, since its multiplier is a constant fixed at
the origin; the explicit gauges supply the overlap isomorphisms.

\item \emph{Nodal extension, including nonprimitive first periods.}
Retain $d_0,n_i,A'_i$ from the primitive vanishing vector $l_i$, and set
$B_i=b\kappa_i/n_i$. Choose integers $\theta,\sigma,x_0,y_0$
with
\[
 n_i\theta+A'_i\sigma=1,\qquad d_0x_0+B_i y_0=1.
\]
The primitive-row condition makes $B_i$ integral and the second
choice possible. The short and long changes and their inverses are
\[
\begin{gathered}
 C_S=\begin{pmatrix}n_i&\sigma\\-A'_i&\theta\end{pmatrix},\quad
 C_L=\begin{pmatrix}x_0&-B_i\\y_0&d_0\end{pmatrix},\\
 C_S^{-1}=\begin{pmatrix}\theta&-\sigma\\A'_i&n_i\end{pmatrix},\quad
 C_L^{-1}=\begin{pmatrix}d_0&B_i\\-y_0&x_0\end{pmatrix}.
\end{gathered}
\]
With $S=(u_i,\delta)$ and $L=(v_i,\gamma)$ write
$(l,k)=SC_S$ and $(h,w_0)=LC_L$. The ordered basis
$(l,h,k,w_0)$ has determinant one and changes $T_i$ to
$I+E_{12}$. The identity on the universal cover, with this period
change and the displayed inverse, identifies the original germ
and its collar. Its alternating pairing matrix is
\[
 K_N=(SC_S)^t\xi(LC_L)
       =\begin{pmatrix}N&0\\M&R_0\end{pmatrix},\qquad
 N=\frac{a\nu_i}{d_0^2},\qquad R_0=-\frac{bd_0^2}{\nu_i}.
\]
These entries are integral, and $\delta=-\sigma l+n_i k$.
The matrix $\Omega=C_S^{-1}\Omega_{\rm ph}C_L$ changes by
$E_{11}$ around the base; $K_N^t\Omega$ is symmetric.
For $Z=U+\Omega W$, the gauge
\[
 G_N=\exp\{-2\pi i(U^tK_NW+W^tK_N^t\Omega W)\}
\]
and its reciprocal give short factors one and long $h$-factor
\[
 \exp\{-4\pi i h^tK_N^tZ-2\pi i h^tK_N^t\Omega h\}.
\]
The shift of $U$ cancels its original short factor; a shift of
$W$ gives this expression by the stated symmetry. The integral
support $Nj^2$ has difference $2Nhj+Nh^2$, which supplies
exactly its normal-coordinate factor on the adjacent nodal
charts. The remaining unit is the corresponding unit of $\Omega$.
Around the base, replacing $\Omega$ by $\Omega+E_{11}$ agrees
with replacing $U$ by $U+E_{11}W$ in $G_N$; both changes add
$NW_1^2$ to its exponent before multiplication by $-2\pi i$.
This proves the extension and its compatibility with the collar.

For this rank-one paragraph put $r=\operatorname{ord}h_G=-\kappa_i$
and $\rho=\operatorname{ord}R=-ab\kappa_i^2/\nu_i$.
The physical support can also be retained explicitly. Write
$H_{\rm val}=\alpha_0 u$, $F_{\rm val}=\beta_0u$ with $u=(u_1,u_2)$ primitive,
$\alpha_0>0$ and $\gcd(\alpha_0,\beta_0)=1$. Here
\[
 \alpha_0=d_0,\quad u_1=n_i,\quad u_2=-A'_i,\quad\beta_0=B_i,\quad
 a\beta_0u_1+b\alpha_0u_2=0.
\]
Hence $\alpha_0\mid au_1$, and
\begin{equation}
 \psi_1(j)=-\frac{au_1}{\alpha_0}j(j-\alpha_0)
          =-Nj^2+N\alpha_0j.
 \hypertarget{tag-022J}{}\label{maps:eq:D.OP5a}
\end{equation}
Its differences under $\alpha_0,\beta_0$ are
$-2au_1j$ and $2bu_2j+b\rho-ab r$.
The second difference follows by factoring its residual as
$(a\beta_0u_1+b\alpha_0u_2)(\alpha_0-\beta_0-2j)/\alpha_0$.
The elliptic combination $\alpha_0F_{\rm val}-\beta_0H_{\rm val}$ has zero
normal monomial exponent. Its remaining elliptic multiplier is
retained. The linear term in \eqref{maps:eq:D.OP5a} is precisely
the gauge $z^a$; the sign of the quadratic term is the opposite
frame convention. Thus the same extension includes nonprimitive
$H_{\rm val}$, as well as $\beta_0=0$.

At the fixed $I_p$ one has
\[
 d_0=1,\quad N=ap,\quad R_0=-b/p=-Q_f,\quad
 \delta=-\sigma l+p k,\quad a\kappa_i\sigma\equiv1\pmod p.
\]
Its normal weights are $(-\sigma,\sigma)$. The genuine Tate
valuation is $r=-\kappa_i$, so $w=-\sigma=(ar)^{-1}$.
The line character is $-\sigma N(2j+1)=0$ modulo $p$.
The integral shift by $k$ in $\delta/p=-\sigma l/p+k$
adds no character, since the short factors are one.
The surviving period pair $(k,w_0)$ has the complex orientation
of $E$: the whole ordered basis has determinant one and the
normal pair has the positive orientation from $XY=q$.
Specialization sends this torus to $E$ with coefficient one.
Therefore
\[
                 \deg(\mathcal L|_E)=-2Q_f.
\]

\item \emph{Cusp extension on the prescribed fan.}
With the displayed cusp period matrix $C_0$ put
\[
 K=C_0^t\operatorname{diag}(a,-b)
   =\begin{pmatrix}
       a\nu&ab\kappa\\
       ab\kappa&b(ab\kappa^2-1)/\nu
     \end{pmatrix}.
\]
Both $K$ and $K\Omega$ are symmetric, and $Kd_\delta=(0,-b)^t$.
The gauge
\[
 G_C=\exp\{-2\pi i(U^tKW+W^tK\Omega W)\}
\]
and its reciprocal give short factors one and long factors
\[
 \exp\{-4\pi i h^tKZ-2\pi i h^tK\Omega h\}
   =q^{-h^tKh}X^{-2Kh}\exp\{-2\pi i h^tKB(q)h\}.
\]
The final factor is retained. The identity
$G_C(U,W,\Omega+I)=G_C(U+W,W,\Omega)$ gives compatibility
around the original base circle.
The integral support $\psi(t)=t^tKt$ has difference
$2h^tKt+h^tKh$. On a unit triangle let $m_\sigma$ be its
affine slope. A regular frame and vector coordinate are
$e_\sigma=\chi^{-m_\sigma}e$ and
$\zeta_\sigma=\chi^{m_\sigma}\zeta_C$, with inverse $\zeta_C=\chi^{-m_\sigma}\zeta_\sigma$; the opposite exponents give
the inverse maps. They extend the line with its complete period
factors. A Cartier line requires no convexity of this support.

To relate this to the physical expression, write $r=-\kappa$ and
$\rho=(1-ab r^2)/\nu$. Then
\begin{equation}
\begin{gathered}
 \psi_{\rm ph}(x,y)=-a\rho x^2-2ab rxy+b\nu y^2+ax,\\
 C_0=\begin{pmatrix}\nu&-br\\ar&\rho\end{pmatrix},\qquad
 \psi_{\rm ph}(C_0t)=-t^tKt+a(\nu t_1-brt_2).
\end{gathered}
 \hypertarget{tag-022I}{}\label{maps:eq:D.OP5}
\end{equation}
Its physical differences under $H_{\rm val}=(\nu,ar)$ and $F_{\rm val}=(-br,\rho)$
are $-2ax$ and $2by+b\rho-ab r$, the factors in
\eqref{maps:eq:D.OP4}. The last equality is their pullback to
the specified short lattice; its affine term is the same
scalar-invariant monomial gauge $z^a$.

For the edge $[0,e_2]$ the two triangle slopes of $\psi$ are
\[
 m_+=(K_{11},K_{22}),\qquad m_-=(2K_{12}-K_{11},K_{22}).
\]
The dense-torus scalar lift in this gauge has factor one.
Its fixed-line characters are $m_+\cdot d_\delta$ and $m_-\cdot d_\delta$.
The first equals $b\{a\kappa(b\kappa-\nu)-1\}$; their
difference is $2(K_{11}-K_{12})(d_\delta)_1$. Both are divisible by $p$.
The coordinates on the two patches of $C$ are $X_1$ and
$X_1^{-1}$. Their frame transition has exponent
$2(K_{11}-K_{12})$. With the complex orientation this yields
\begin{equation}
\begin{split}
 &\psi_{\rm ph}(C_0e_1)+\psi_{\rm ph}(C_0(-e_1+e_2))
               -\psi_{\rm ph}(0)-\psi_{\rm ph}(C_0e_2)
                  =-2a(\nu+br),\\
 &\deg(\mathcal L|_C)=2(K_{11}-K_{12})
                   =2a(\nu+br)=2a(\nu-b\kappa).
\end{split}
 \hypertarget{tag-022K}{}\label{maps:eq:D.OP6}
\end{equation}
The matrix $C_0$ and its inverse identify these four physical vertices
with the wall of the prescribed fan. The sign between the two
lines follows from the frame and vector-coordinate maps above.
The construction gives one $C_p$-equivariant line $\mathcal L$
on all of $Y$, with zero characters at both fixed curves.
Its class has value two on $s$ and zero on every section sweep.
Nonzero unit functions on a base disc have logarithms, are scalar
invariant, and have zero divisor degree on a central curve;
the displayed factors retain them throughout.

\item \emph{The global Borel comparison.}
The cusp kills precisely $u,\delta$ in degree one, since its
two short columns are $SC_0$ and $C_0$ is unimodular. Each rank-one
cap kills its primitive loop $l_i$ and its base meridian.
Collect these angular columns in the coordinates $e_1,e_2,\delta$
as a matrix $K_1$. For each of the thirty marked rows in the table,
the greatest common divisor of its three-column minors is one.
This is obtained directly from the primitive $u_i$, $\nu_i$ and
$\kappa_i$ in the ordered input and the displayed formula for $l_i$.
It gives a constructive integral right inverse: if $K_{1,J}$
is a three-column submatrix and
$\sum_J c_J\det K_{1,J}=1$, take
\[
 R_1=\sum_J c_J E_J\operatorname{adj}(K_{1,J}),\qquad
                         K_1R_1=I_3,
\]
where $E_J$ includes the chosen three columns. Thus all three
angular generators vanish integrally. The remaining monodromy
relations make $\gamma$ commute with every unfilled base meridian.
The ordinary meridian is the inverse product of the additive ones.
Consequently
\begin{equation}
                    \pi_1(Y)=\mathbb Z\gamma\times F_{r_Y},
                    \qquad r_Y=0\ \text{or }1.
 \hypertarget{tag-022G}{}\label{maps:eq:D.OP3}
\end{equation}
The inverse map kills $e_1,e_2,\delta$, retains $\gamma$ and the
unfilled additive meridians, and replaces the ordinary meridian
by their inverse product. It satisfies every matrix and cap
relation; the right inverse $R_1$ gives expressions for the killed
generators. In particular $H_1(Y)$ is free. The degree-two
universal coefficient map and the exterior generators and their section sweeps now give
\[
                  \tfrac12\rho_p c_1(\mathcal L)=\alpha|_Y.
\]

The connected scalar circle acts trivially on cohomology. In its
Borel spectral sequence, the only possible outgoing differentials
from $\alpha$ have targets
$H^2(BS^1;H^1(X;\mathbb F_p))$ and $H^3(BS^1;\mathbb F_p)$;
both are zero. Hence $\alpha$ lifts to circle-equivariant
cohomology, and then to $A\in H^2_{C_p}(X;\mathbb F_p)$.
Write
\[
 H^*(BC_p;\mathbb F_p)=\mathbb F_p[t]\otimes\Lambda(\eta),
 \qquad |t|=2,\quad|\eta|=1,\quad\eta^2=0,
\]
where $t$ is the first Chern class of the character taking the
specified generator to $\zeta_p$.

Choose a toric endpoint of $C$ as basepoint. The whole scalar
circle fixes this point, so its isotopy is based and gives
$\pi_1(Y_{C_p})=\pi_1(Y)\times C_p$.
Any fixed point supplies an order-$p$ section. The first factor
is torsion free by \eqref{maps:eq:D.OP3}, so every such section has
the same generator $(1,\zeta_p)$. Extend an integral basis
$\lambda_j$ of $H^1(Y)$ to be zero on that group factor.
The kernel of degree-two Borel restriction to $Y$ is
\[
                 \mathbb F_p t+
                 \sum_j\mathbb F_p\rho_p(\lambda_j)\eta.
\]
Indeed the positive-base-degree terms on that diagonal are
$(2,0)$ and $(1,1)$. A fixed-point section preserves the base
class $t$, and the displayed degree-one products represent every
$(1,1)$ term; hence they span the kernel with no further possible
term. The difference
$A|_Y-\tfrac12\rho_p c_1^{C_p}(\mathcal L)$ lies in it.
At every fixed point the products restrict to zero, so that
difference has one common coefficient of $t$ on both curves.
Subtract that base multiple from $A$, with addition as inverse.
Since both line characters are zero, the resulting restrictions are
\begin{equation}
 A|_C=\alpha_C[C]^*,\qquad
 A|_E=\alpha_E[E]^*+\eta\zeta,\qquad
 \alpha_C=a\nu,\quad\alpha_E=-Q_f,\quad
 \zeta\in H^1(E;\mathbb F_p).
 \hypertarget{tag-022M}{}\label{maps:eq:D.OP8}
\end{equation}
The degrees follow from \eqref{maps:eq:D.OP6} and the nodal degree $-2Q_f$, using
$p\mid b$. This comparison fixes the characters globally and
retains the possible mixed term.

\item \emph{Localization with the physical weights.}
Average a metric over $C_p$ and use its exponential map to identify
tubular neighborhoods of the fixed curves with their normal bundles.
The complement is free, because $p$ is prime and
\eqref{maps:eq:D.OP7} is the whole fixed set. Its Borel space is
homotopy equivalent to its finite-dimensional quotient, so a high
power of $t$ annihilates its cohomology. The same holds on the
punctured normal boundaries. Thom excision therefore makes
restriction to $C\sqcup E$ an isomorphism after inverting $t$.

For a fixed curve $F_{\rm fix}$ with normal characters $\lambda,-\lambda$
and ordinary normal roots $x_1,x_2$, its Euler class is
\[
 e_{F_{\rm fix}}=(\lambda t+x_1)(-\lambda t+x_2)
          =-\lambda^2t^2+\lambda(x_2-x_1)t.
\]
The two character projectors over the complex bundle give its
normal line subbundles. The leading coefficient is a unit, and
the inverse of $e_{F_{\rm fix}}$ is the finite geometric-series inverse,
since products of positive-degree classes vanish on a curve.
The inverse restriction map is
\[
 (z_C,z_E)\longmapsto
       i_{C*}(z_C/e_C)+i_{E*}(z_E/e_E).
\]
The equivariant Thom identity gives both inverse composites:
$i_{F_{\rm fix}}^*i_{F_{\rm fix}*}$ is multiplication by $e_{F_{\rm fix}}$, and the supports are
disjoint. All orientations are the complex orientations.

Apply this map to $A$ and integrate, obtaining zero in negative
degree. With $\lambda=\nu$ at $C$ and $\lambda=w$ at $E$,
\begin{equation}
                     \frac{\alpha_C}{\nu^2}
                         +\frac{\alpha_E}{w^2}=0.
 \hypertarget{tag-022N}{}\label{maps:eq:D.OP9}
\end{equation}
A mixed term $\eta\zeta$ has ordinary degree one; multiplying
it by an ordinary degree-two Euler correction has degree three
on a curve. It cannot contribute to this integral.

The tangent line has zero character and its square is zero on
the fixed curve. The normal roots give the exact identity
\[
 p_1^{C_p}(TX)|_{F_{\rm fix}}=-2e_{F_{\rm fix}},
\]
because the difference after adding $2e_{F_{\rm fix}}$ is
$(x_1+x_2)^2=0$. This holds with unequal ordinary normal degrees.
Equivariant $c_1$ is also zero: it is a base class because
$H^1(X;\mathbb Z)=H^2(X;\mathbb Z)=0$, and its restriction at
$C$ is zero since the two normal characters sum to zero.
Localization now gives
\begin{equation}
\begin{split}
 \langle\alpha\,\rho_p p_1(X),[X]_p\rangle
   &=-2(\alpha_C+\alpha_E)=2(Q_f-a\nu),\\
 \operatorname{PD}p_1(X)&=2(Q_f-a\nu)s.
\end{split}
 \hypertarget{tag-022D}{}\label{maps:eq:D.OP10}
\end{equation}
For the second equality, $H^4(X;\mathbb Z)=\mathbb Z/p$ and
its reduction to $H^4(X;\mathbb F_p)$ is an isomorphism:
multiplication by $p$ on the former is zero and the next
Bockstein target $H^5(X;\mathbb Z)$ is zero. Pairing with
$\alpha(s)=1$ is therefore the coordinate map on the integral
Poincare dual.

Each restriction in \eqref{maps:eq:D.OP8} has cube zero: a
degree-two class squares to zero on a curve, $\eta^2=0$, and
the remaining product has ordinary degree three. The same
localization isomorphism and integration yield
\begin{equation}
                         \langle\alpha^3,[X]_p\rangle=0.
 \hypertarget{tag-022E}{}\label{maps:eq:D.OP11}
\end{equation}
The additive caps lie in the free complement; their role in
extending the ordinary class was accounted for by the complete finite and ordinary cap kernels
and the global comparison above.

\item \emph{The full integral fundamental-class invariant.}
Let $A_X=H_2(X;\mathbb Z)$ and let
$f_X:X\to K(A_X,2)$ induce the identity on $H_2$.
Existence and uniqueness follow from the degree-two universal
coefficient isomorphism, since $H_1(X)=0$.
The marking and its inverse are
\[
 \psi_s:\mathbb Z/p\longrightarrow A_X,\quad[j]\longmapsto js,
 \qquad\phi_s(z)=\alpha(z).
\]
The marked map $f_s=K(\phi_s,2)\circ f_X$ pulls the universal
class $u$ back to $\alpha$. Lemma~\ref{maps:src:D.1.10} gives
\[
 H_6(K(\mathbb Z/p,2);\mathbb Z)
       =\mathbb Z/p\langle z_3\rangle,\qquad
 z_3=i_*[\mathbb {CP}^3],\qquad\langle u^3,z_3\rangle=1,
\]
where $i:K(\mathbb Z,2)\to K(\mathbb Z/p,2)$ induces reduction
on $H_2$, and $\mathbb {CP}^3$ has its complex orientation.
Thus the evaluation map is an isomorphism, with inverse
$[j]\mapsto jz_3$. Equation~\eqref{maps:eq:D.OP11} implies
$(f_s)_*[X]=0$ in this entire integral group. Applying
$K(\psi_s,2)_*$ proves $\mu(X)=0$ with its intrinsic marking.

The refined-class map in Theorem~\ref{maps:src:D.1.11}, for
a group $A$ without two-torsion, evaluates on the marked
spin-bordism representative $[N_{\rm bord},f_{\rm bord}]$ as
\[
 P_A([N_{\rm bord},f_{\rm bord}])
   =\tfrac14 (f_{\rm bord})_*(\operatorname{PD}_{N_{\rm bord}}p_1(TN_{\rm bord})).
\]
This is the map of the preparatory lemma on
\cite[pp.~795--796]{manZhubr75}; Section~3.13 on p.~807
identifies the intrinsic refined class with its value on the
marked spin-bordism class. Hence
\[
 P_{\mathbb Z/p}(\gamma(X,\phi_s))
      =\phi_s(\mathfrak p(X))
      =\tfrac14\phi_s(\operatorname{PD}p_1(X)).
\]
Multiplication by four is an automorphism with inverse $4^{-1}$.
Equation~\eqref{maps:eq:D.OP10} gives
$\mathfrak p(X)=(Q_f-a\nu)s/2$. For any other marking $\phi$,
naturality gives $P_A(\gamma(X,\phi))=\phi(\mathfrak p(X))$.
This fixes the factor and the integral torsion class in
\eqref{maps:eq:D.OP1}.

\item \emph{Tate coordinates and all generator changes.}
At the unique $I_p$ the original-germ valuations are
\[
 (\operatorname{ord}h_1,\operatorname{ord}h_G,
       \operatorname{ord}R)=(p,r,-ab r^2/p),\qquad r=-\kappa_i.
\]
Write $d=pD_0$. In families $47,55,56$ the pairs
$(D_0,r\bmod p)$ are $(2,2),(4,3),(6,4)$, respectively,
with the same residues in their displayed collision profiles.
Thus $D_0r^2=1$ in $\mathbb F_p$.
The same identity follows from the height formula: multiply
$1/d$ by $p$; only the contribution of the unique $I_p$ has
a denominator divisible by $p$, and its residue is $r^2$.
Since $ab\nu=d$ and $b=pQ_f$, we obtain
\[
 D_0=a\nu Q_f,\qquad
 w^2=(ar)^{-2}=D_0/a^2=\nu Q_f/a,\qquad
 P_X=\frac{Q_f}{2}\left(1-\frac{\nu^2}{w^2}\right).
\]
Both $Q_f$ and $\nu$ are units modulo $p$, so the zero test is
exactly $w^2=\nu^2$.

The line degrees and the global comparison give
$[E]=-Q_fs$ and $[C]=a\nu s$. The following markings of
$H_2(X)$ have the displayed inverse evaluation classes and
coordinates of the same refined class:
\[
\begin{array}{c|c|c}
\text{image of }1&\text{inverse evaluation}&\text{refined coordinate}\\\hline
s&\alpha&\tfrac{Q_f}{2}(1-\nu^2/w^2)\\
{[E]}&-Q_f^{-1}\alpha&\tfrac12(\nu^2/w^2-1)\\
w^{-2}[E]&-w^2Q_f^{-1}\alpha&\tfrac12(\nu^2-w^2)\\
(\nu^2/w^2)[E]&-w^2(Q_f\nu^2)^{-1}\alpha&\tfrac12(1-w^2/\nu^2)\\
{[C]}& (a\nu)^{-1}\alpha&\tfrac12(w^2/\nu^2-1).
\end{array}
\]
Each coefficient multiplying $s$ in the first column is a unit,
and evaluation of its inverse on
the image of $1$ is one; the reverse composite fixes every $js$.
The subgroup automorphism
$\zeta_p\mapsto\zeta_p^{\nu^{-1}}$, with inverse exponent
$\nu$, pulls both $t$ and $\eta$ back to $\nu^{-1}$ times
themselves. It changes the normal weights to $(-1,1)$ at $C$
and $(w/\nu,-w/\nu)$ at $E$.
For the cohomology class $u_C=(a\nu)^{-1}\alpha$, which has
$u_C[C]=1$, equations \eqref{maps:eq:D.OP9}--\eqref{maps:eq:D.OP10}
therefore read
\[
 1+\frac{u_C[E]}{(w/\nu)^2}=0,\qquad
 \langle u_Cp_1,[X]\rangle=-2\bigl(1-w^2/\nu^2\bigr).
\]
The acting-group map changes both normal characters together;
the Tate value $w$ in the statement refers to the physical generator.

For any further marking $1\mapsto\lambda s$, the inverse is
$\lambda^{-1}\alpha$, the refined coordinate is
$\lambda^{-1}P_X$, and the full degree-six coordinate is
$\lambda^{-3}\langle\alpha^3,[X]_p\rangle=0$.
Indeed multiplication by $\lambda$ on $H_2(K(\mathbb Z/p,2))$
pulls $u$ back to $\lambda u$ and sends $z_3$ to
$\lambda^3z_3$. This gives the induced map and inverse on
the entire fundamental-class group.

In particular the eight patterns numbered
$21,22,25,26,27,28,29,33$ have $w^2=\nu^2=4$ and $P_X=0$.
Pattern $30$ has $w^2=1$, $\nu^2=4$ and $P_X=2$;
pattern $34$ has the same two squares and $P_X=3$.
For example, pattern $21$ with $e=1$ has cusp data
\[
 C_0=\begin{pmatrix}3&-10\\4&-13\end{pmatrix},\quad
 K=\begin{pmatrix}6&-20\\-20&65\end{pmatrix},\quad
 d_\delta=\binom{10}{3}.
\]
Its line characters are $255,-265$, both zero modulo five,
and its half-line degree is $26$; the elliptic half-line
degree is $-1$. Thus the Pontryagin pairing is $-50=0$.
Pattern $30$ is realized by $m=4,e=1,n=3,r_0=0$; its two
half-line degrees are $13,-2$, giving pairing $-22=3$ and
refined coefficient $2$. The preceding maps prove the values
for every allowed parameter in these patterns.

\item \emph{Realization and oriented classification.}
The table gives twelve zero and twelve nonzero order-five
patterns, and twelve nonzero order-seven patterns.
Patterns $21,13,1$, with $e=1$ and no finite cap, realize
the three data sets in the statement. For the first two take
$1\mapsto s$; for pattern $1$ take $1\mapsto3s$, with
inverse $3^{-1}\alpha$. This makes its refined coordinate one.
For any nonzero $P_X$, the maps
\[
 \psi_X([j])=jP_Xs_X,\qquad
 \phi_X(z)=P_X^{-1}\alpha_X(z)
\]
are mutually inverse and give refined coordinate one and
fundamental-class coordinate zero. When $P_X=0$, use
the already defined inverse marking maps $\psi_s,\phi_s$.

Theorem~5 of \cite[p.~796]{manZhubr75} makes
\[
 (P,h):\Omega_6^{\mathrm{spin}}(K(\mathbb Z/p,2))
       \longrightarrow\mathbb Z/p\oplus
                H_6(K(\mathbb Z/p,2);\mathbb Z),\qquad
 h([N_{\rm bord},f_{\rm bord}])=(f_{\rm bord})_*[N_{\rm bord}],
\]
injective. The marked Theorem~3 on p.~795, with its
Section~4.1 construction and uniqueness proof on p.~809,
then gives an orientation-preserving diffeomorphism $F_{\rm dif}:X\to X'$
whenever the marked bordism classes agree and the free
third-homology ranks vanish. Its marking equation is
$\phi_{X'}(F_{\rm dif})_* =\phi_X$. The integral homology calculation gives these hypotheses;
both free ranks $H_2,H_3$ vanish in the present application.
The equality of $P$ and $h$ follows from
\eqref{maps:eq:D.OP10}--\eqref{maps:eq:D.OP11} and the displayed markings, so the theorem applies to any two members
in the same one of the three data sets.

For two nonzero refined classes the induced map and inverse are
\[
 (F_{\rm dif})_*(s_X)=\frac{P_{X'}}{P_X}s_{X'},\qquad
 (F_{\rm dif}^{-1})_*(s_{X'})=\frac{P_X}{P_{X'}}s_X.
\]
For two zero classes they send $s_X$ to $s_{X'}$ and back.
Naturality of $P$ verifies the refined class through these maps;
the cube of each scalar gives the induced map on the complete
degree-six coordinate. A group isomorphism preserves element
orders, so orders five and seven give different types.
Within order five every automorphism multiplies by a unit and
therefore preserves zero versus nonzero refined class. Conversely
the displayed maps identify every two nonzero cases. This proves
the stated three oriented types with every marked assignment.
\end{enumerate}
\end{proof}



\subsubsection*{Orders ten and fourteen}

\begin{definition}[Minimal mixed-primary data]\hypertarget{tag-Q207}{}\label{maps:def:minimal-mixed-data}
{\normalfont\scriptsize tag Q207}\par
Use the ordered profiles of Definition~\ref{mapsIII:def:elliptic-marked-data}
and a pattern in \eqref{maps:eq:D.MP2}. Use the prescribed primitive
Mumford fan at every multiplicative position, its listed additive
choice and the ordinary clutch $e\gamma$. The logarithmic sites are
exactly the listed additive sites. At a finite site of linear
monodromy order $d_i$, allow every multiplicity $m=d_i k$, $k>0$,
and every free affine twist. The ordinary parameter is $e\in\mathbb Z$. A selected position is
$i-1$ in the input definition's peripheral order.

Use the framed models $M_0,M_1$ of
Definition~\ref{maps:def:minimal-two-data}. Define the marked models
$Z_{5,0},Z_{5,1},Z_{7,1}$ by the respective constructions
$\mathrm{OS56}$ with $(a,b)=(2,5)$ and selected position $1$,
$\mathrm{OS55}$ with $(a,b)=(1,5)$ and selected position $1$,
and $\mathrm{OS47}$ with $(a,b)=(1,7)$ and selected position $1$.
In each case take $e=1$ and its listed multiplicative caps.
Their marking sends the odd generator $e_p$ to $s,s,3s$,
respectively; the inverse evaluations are $\alpha,\alpha,3^{-1}\alpha$,
where $\alpha(s)=1$. Theorem~\ref{maps:src:D.2.36} proves
$H_2=\mathbb Z/p$, $b_3=0$, $\mu=0$ and
$\mathfrak p=\epsilon e_p$ for these actual models.
The preceding filling choices constitute the \emph{minimal mixed-primary data}.
\end{definition}

\begin{theorem}[Ten- and fourteen-torsion smooth types]\hypertarget{tag-02LE}{}\label{maps:src:D.2.38}
{\normalfont\scriptsize tag 02LE}\par
Every simply connected threefold $X$ of
Definition~\ref{maps:def:minimal-mixed-data} satisfies
\begin{equation}
\begin{aligned}
 H_2(X)&=\mathbb Z/(2p)\langle s\rangle,\quad H_3(X)=\mathbb Z/(2p),\quad b_3(X)=0,\\
 X&\cong M_t\# Z_{p,\epsilon}
 \quad\hbox{orientation-preservingly}.
\end{aligned}
 \hypertarget{tag-0209}{}\label{maps:eq:D.MP1}
\end{equation}
Here $s$ is the primitive fibre class. If the selected zero is $I_\nu$
and its marked cocycle is $\kappa u$, put
\[
\begin{aligned}
 t&=1+\kappa\pmod2,\qquad P_p=\frac{b/p-a\nu}{2}\pmod p,\\
 P_{2p}&\equiv t\pmod2,\quad P_{2p}\equiv P_p\pmod p.
\end{aligned}
\]
Then $X$ is spin, $\mathfrak p(X)=P_{2p}s$ and
$\mu(X)=pt\,\gamma_3(s)$. The generator $\gamma_3(s)$ is the
reduced positive complex $\mathbb{CP}^3$ class. Set
$\epsilon=0$ for $P_p=0$ and $\epsilon=1$ otherwise.
Writing $\mathrm F$ for a finite cap and $\mathrm O$ for an original cap, the
complete table, where $\mu_{2p}$ is the coefficient of $\gamma_3(s)$, is
\begin{equation}
\begin{array}{c|c|r|c|r|r|r|r|r|r}
\text{OS/profile}&(a,b)&\text{zero}&\text{cap}&2p&t&P_p&P_{2p}&\mu_{2p}&\epsilon\\\hline
\mathrm{OS56}&(1,10)&1&\text{none}&10&1&2&7&5&1\\
\mathrm{OS56}&(3,10)&3&\text{none}&10&1&2&7&5&1\\
\mathrm{OS56}&(3,10)&4&\text{none}&10&0&2&2&0&1\\
\mathrm{OS56}/\mathrm{II}&(1,10)&3&II:\mathrm F&10&0&2&2&0&1\\
\mathrm{OS56}/\mathrm{II}&(1,10)&3&II:\mathrm O&10&0&2&2&0&1\\
\mathrm{OS56}/\mathrm{IV}&(1,30)&3&IV:\mathrm F&10&0&0&0&0&0\\
\mathrm{OS56}/\mathrm{IV}&(3,10)&3&IV:\mathrm O&10&0&2&2&0&1\\
\mathrm{OS47}&(1,14)&2&\text{none}&14&1&4&11&7&1\\
\mathrm{OS47}&(1,14)&3&\text{none}&14&0&4&4&0&1\\
\mathrm{OS47}&(1,14)&4&\text{none}&14&0&4&4&0&1\\
\mathrm{OS47}/\mathrm{II}&(1,14)&2&II:\mathrm F&14&1&4&11&7&1\\
\mathrm{OS47}/\mathrm{II}&(1,14)&2&II:\mathrm O&14&1&4&11&7&1\\
\end{array}
\hypertarget{tag-020B}{}\label{maps:eq:D.MP2}
\end{equation}
There are seven order-ten patterns on six initial rows, giving exactly
the three types $M_1\#Z_{5,1}$, $M_0\#Z_{5,1}$ and $M_0\#Z_{5,0}$
with multiplicities $2,4,1$. There are five order-fourteen patterns
on four initial rows, giving exactly $M_1\#Z_{7,1}$ and $M_0\#Z_{7,1}$
with multiplicities $3,2$.
\end{theorem}


\begin{proof}
\begin{enumerate}
\item \emph{Models and markings.}
The two order-two models of Definition~\ref{maps:def:minimal-two-data}
have the following full marked invariants, by
Theorem~\ref{maps:src:D.2.35} and its reduced spin-bordism quotient:
\begin{equation}
 H_2=H_3=\mathbb Z/2,\quad b_3=0,\quad
 \mathfrak p(M_t)=t e_{2,\rm coef},\quad
 \mu(M_t)=t\gamma_3(e_{2,\rm coef}),\qquad t=0,1.                 \hypertarget{tag-Q208}{}\label{maps:eq:MP_local_1}
\end{equation}

Here $e_{2,\rm coef}$ is the unique nonzero second-homology class. Let

\begin{equation}
 Z_{p,\epsilon},\quad
 (H_2,b_3,\mathfrak p,\mu)
       =(\mathbb Z/p,0,\epsilon e_{p,\rm coef},0)                 \hypertarget{tag-Q209}{}\label{maps:eq:MP_local_2}
\end{equation}

be the marked odd-primary models of Definition~\ref{maps:def:minimal-mixed-data}.
Here $e_{p,\rm coef}$ denotes the marked $e_p$ of that definition.
Write $s_{\rm mod}$ for the fibre generator $s$ of that model,
and $\alpha_{\rm mod}$ for its coordinate evaluation.
In particular, its order-seven marking sends $e_{7,\rm coef}$ to $3s_{\rm mod}$,
with inverse evaluation $3^{-1}\alpha_{\rm mod}$.
Both values $\epsilon=0,1$ exist for $p=5$; the value $1$ exists
for $p=7$. The marking for a nonzero original odd coefficient $P_{\rm odd}$
is $e_{p,\rm coef}\mapsto P_{\rm odd} s_{\rm mod}$, with inverse $P_{\rm odd}^{-1}\alpha_{\rm mod}$, where
$\alpha_{\rm mod}(s_{\rm mod})=1$. This is the exact meaning of the coordinate one in \eqref{maps:eq:MP_local_2}.


\item \emph{Exhaustion of minimal fillings.}

Use the 21 defining ordered profiles, their elliptic monodromies $A_i$,
generator cocycles $g_i$, component coordinates and generator heights.
These are the ordered data of Definition~\ref{mapsIII:def:elliptic-marked-data}.

At a multiplicative position choose primitive $u_i,v_i\in\mathbb Z^2$
by

\begin{equation}
 A_i-I=\nu_i u_i u_i^tJ,\quad g_i=\kappa_i u_i,\quad
 J=\begin{pmatrix}0&1\\-1&0\end{pmatrix},\quad
 \det(u_i,v_i)=1.                                      \hypertarget{tag-Q20A}{}\label{maps:eq:MP_local_5}
\end{equation}

The content of $A_i-I$ is $\nu_i$. A nonzero column, divided by
its content and with first nonzero entry positive, gives $u_i$.
A Bezout identity gives $v_i=(-y_{\rm ell},x_{\rm ell})$ from
$x_{\rm ell}u_{i,1}+y_{\rm ell}u_{i,2}=1$.
This determines all signs, including $I_1$ and $\kappa_i=0$.

For height $1/d$, the necessity and sufficiency criterion of \ref{maps:src:D.2.22}
at the selected $I_\nu$ is

\begin{equation}
 ab=d/\nu,\quad a,b>0,\quad \gcd(a,b)=1,
 \quad \gcd(b\kappa,\nu)=1,                            \hypertarget{tag-Q20B}{}\label{maps:eq:MP_local_6}
\end{equation}

together with narrowness of $aG$ or $bG$ at every other reducible
fibre and the literal component identity

\begin{equation}
 [j=a\kappa]-[j=0]-[j=(a-b)\kappa]+[j=-b\kappa]
       =-2[j=0]+[j=1]+[j=-1]\quad(j\in\mathbb Z/\nu). \hypertarget{tag-Q20C}{}\label{maps:eq:MP_local_7}
\end{equation}

For every positive divisor $a$ of $d/\nu$, set
$b=d/(a\nu)$ and test \eqref{maps:eq:MP_local_6}--\eqref{maps:eq:MP_local_7} and the component congruences.
Repeated residues at $\nu=1,2$ are added with their integer signs.
The per-profile counts, in the order of that definition, are

\begin{equation}
 6,2,8,4,6,2,14,6,8,10,2,4,10,2,4,14,6,26,10,14,14,
                                                               \hypertarget{tag-Q20D}{}\label{maps:eq:MP_local_8}
\end{equation}

whose sum is 172. The divisor bounds and the complete component congruences prove exhaustion.

Discard the profiles with $I_n^*$. At an additive site the finite
option is present exactly when $b(A-I)^{-1}g$ is integral. Its
monodromy order and vertical index are

\begin{equation}
\begin{array}{c|ccccc}
 &II&III&IV&III^*&IV^*\\\hline
 d_i&6&4&3&4&3\\q_i&1&2&3&2&3.
\end{array}                                             \hypertarget{tag-Q20E}{}\label{maps:eq:MP_local_9}
\end{equation}

The original minimal option is present for $II,III,IV$ when the section $P_{\rm sec}=aG$
is narrow and $Q$ is transitive on the components. The exact cap maps of \ref{maps:src:D.2.26} prove this criterion and exclude the additional local defects.
In the marked component coordinates it is narrowness of $aG$ and
$\gcd(b,q_i)=1$. Every retained original has zero line multidegree
and one component orbit; a nonzero-Chern original is not substituted.

Choose one option at each additive site. The finite orders must be
pairwise coprime by the Seifert presentation in \ref{maps:src:D.2.26}. Compute the
residual $c=b/\prod_{i\in\mathcal I_{\rm fin}}q_i$, where $\mathcal I_{\rm fin}$
is the set of finite additive sites, and retain $c=10,14$. These finite
tests give exactly the twelve rows in the theorem. Each has at most one
finite site. The finite local monodromy orders $d_i$ are six at $II$ and three at
the finite $IV$; the multiplicity $m=d_i k$ remains arbitrary. The displayed pairs satisfy

\begin{equation}
 a\text{ odd},\quad b/2\text{ odd},\quad
 \nu=1\text{ or }3,\quad \gcd(\nu,2p)=1.              \hypertarget{tag-Q20F}{}\label{maps:eq:MP_local_10}
\end{equation}

The ordered profiles determining these positions are

\begin{equation}
\begin{array}{c|l}
\mathrm{OS56}&I_5,I_3,I_2,I_1,I_1\\
\mathrm{OS56}/\mathrm{II}&I_2,II,I_5,I_3\\
\mathrm{OS56}/\mathrm{IV}&I_5,I_2,IV,I_1\\
\mathrm{OS47}&I_7,I_2,I_1,I_1,I_1\\
\mathrm{OS47}/\mathrm{II}&I_7,I_2,I_1,II.
\end{array}                                             \hypertarget{tag-Q20G}{}\label{maps:eq:MP_local_11}
\end{equation}

Use their ordered defining matrices and cocycles in the following period maps.

\item \emph{The actual period and integral exterior maps.}

Write

\begin{equation}
 \Lambda=\mathbb Z\langle e_1,e_2,\delta,\gamma\rangle,
 \quad\chi=\gamma^*,\quad
 \Omega=a e_1^*\wedge e_2^*-b\delta^*\wedge\gamma^*.
                                                               \hypertarget{tag-020C}{}\label{maps:eq:D.MP3}
\end{equation}

Positive $\gamma$ is the actual forward period covering $-Q$.
For a finite elliptic matrix put $z_{\rm per}=(A-I)^{-1}g$.
At a rank-one site use $(u,v,n,\kappa)=(u_i,v_i,\nu_i,\kappa_i)$.
At $I_n$ put
$z_{\rm per}=(\kappa/n)v$, with \eqref{maps:eq:MP_local_5}. Then the actual period matrix is

\begin{equation}
\begin{gathered}
 T=\begin{pmatrix}A&0&bg\\
             \beta&1&\beta(bz_{\rm per})+w_{\rm per}\\0&0&1\end{pmatrix},\\
 \beta=-a g^tJA,\\
 w_{\rm per}=1/\nu\text{ at the selected zero, otherwise }0.
\end{gathered}
                                                               \hypertarget{tag-Q20H}{}\label{maps:eq:MP_local_13}
\end{equation}

Put $c_{\rm per}=\beta(bz_{\rm per})+w_{\rm per}$. The inverse on the same marked lattice is
\[
 T^{-1}=
 \begin{pmatrix}
 A^{-1}&0&-bA^{-1}g\\
 -\beta A^{-1}&1&b\beta A^{-1}g-c_{\rm per}\\
 0&0&1
 \end{pmatrix}.
\]
Multiplication on either side gives the identity.
Every entry is integral; the ordered product is identity, $T\delta=\delta$,
$\chi T=\chi$, and $T^t\Omega T=\Omega$. These identities follow by substitution in the five defining profiles.
The identity of the complex universal cover, followed by the actual
period exponential, maps real period coordinates to the given torus.
The inverse is its group logarithm modulo this same lattice. The
rigidified origin fixes the group map. Thus \eqref{maps:eq:MP_local_13} acts on the actual
bundle, including its original period functions, rather than on an
unidentified abstract torus bundle.

For the wedge order $12,13,14,23,24,34$, the primitive fibre row is

\[
        \omega=(a,0,0,0,0,-b).
\]

Choose $x_{\rm fib},y_{\rm fib}\in\mathbb Z$ with $ax_{\rm fib}+by_{\rm fib}=1$ and set
$s_0=x_{\rm fib} e_{12}-y_{\rm fib} e_{34}$. The matrix

\begin{equation}
 P_{\rm fib}=(s_0,\ b e_{12}+a e_{34},\ e_{13},e_{14},e_{23},e_{24})
                                                               \hypertarget{tag-Q20I}{}\label{maps:eq:MP_local_14}
\end{equation}

is unimodular. The integral exterior sequence of \ref{maps:src:D.2.26} supplies an integer matrix $R_{\rm cyl}$ for the literal monodromy-cylinder matrix

\[
 B_2=[\wedge^2T_1-I\ \cdots\ \wedge^2T_N-I]
\]

and verifies

\begin{equation}
 B_2R_{\rm cyl}=I_6-s_0\omega.                                  \hypertarget{tag-Q20J}{}\label{maps:eq:MP_local_15}
\end{equation}

The columns of $R_{\rm cyl}$ give the integer cylinder chains in that exterior sequence.
Equation \eqref{maps:eq:MP_local_15}, together with $\omega B_2=0$, proves the exact inverse
maps

\begin{equation}
 \wedge^2\Lambda/\operatorname{im}B_2
   \underset{n\mapsto[n s_0]}{\overset{[v]\mapsto\omega(v)}
   {\rightleftarrows}}\mathbb Z .                     \hypertarget{tag-Q20K}{}\label{maps:eq:MP_local_16}
\end{equation}

Indeed the difference $v-\omega(v)s_0$ is the boundary of the
explicit monodromy cylinders with coefficient column $R_{\rm cyl}v$. Denote
the quotient generator by $s$. This is the primitive marking used
throughout; the individual elliptic torus $e_{12}$ represents $a s$.

For a rank-one site, set

\begin{equation}
 d_0=\gcd(n,a\kappa),\quad n_0=n/d_0,\quad A_0=a\kappa/d_0,
 \quad l=n_0u-A_0\delta,
 \quad f=(d_0u^tJ,0,b\kappa d_0/n).
                                                               \hypertarget{tag-Q20L}{}\label{maps:eq:MP_local_17}
\end{equation}

The exact identity is $T-I=lf$. The row $f$ is primitive. A Bezout column $h_f$ has $f h_f=1$, so $(T-I)h_f=l$.

At the selected cusp put $S=(u,\delta)$, $L=(v,\gamma)$, and

\begin{equation}
 C=\begin{pmatrix}\nu&b\kappa\\-a\kappa&\rho\end{pmatrix},
 \quad\rho=(1-ab\kappa^2)/\nu,
 \quad C^{-1}=\begin{pmatrix}\rho&-b\kappa\\a\kappa&\nu\end{pmatrix}.
                                                               \hypertarget{tag-Q20M}{}\label{maps:eq:MP_local_18}
\end{equation}

Its determinant is one. Write $\mathbf e_j$ for the standard columns
in the coordinate domains of $S$ and $L$. In the actual basis
$H=(SC\mathbf e_1,L\mathbf e_1,SC\mathbf e_2,L\mathbf e_2)$, the return sends each long vector to
itself plus the corresponding short vector. Thus
$(T-I)L C^{-1}=S$ provides inverse preimages of $u,\delta$.

Take the columns $u,\delta$ from this cusp and then the other $l$'s
in profile order. Their first-three-coordinate matrix $B_1$ has an
integer right inverse $R_1$ from the same integral exterior maps, with

\begin{equation}
                  B_1R_1=I_3.                         \hypertarget{tag-Q20N}{}\label{maps:eq:MP_local_19}
\end{equation}

Every column has its actual $(T-I)$-preimage just described. Hence
these angular columns generate exactly $\ker\chi$ and the full
monodromy coinvariant lattice is the primitive $\mathbb Z\gamma$.

Remove all additive caps and the ordinary clutch disc, retaining every
multiplicative cap; call this common piece $W$, with regular part $V$.
Filling the multiplicative discs kills their meridians and angular
columns. The remaining section loops therefore give explicit inverse
group maps

\begin{equation}
 \pi_1(W)=\mathbb Z\gamma\times F_{r_B},
       \qquad r_B=0\text{ or }1.                        \hypertarget{tag-Q20O}{}\label{maps:eq:MP_local_20}
\end{equation}

The remaining ordinary meridian is the inverse of the product of the
additive meridians. The maps send those literal loops to the displayed
generators and include them back. Every remaining return induces the
identity on $\mathbb Z\gamma$. In particular \eqref{maps:eq:MP_local_20} is torsion-free.

The primitive boundary maps of \ref{maps:src:D.2.26} are onto in degree two. The
degree-one Mayer--Vietoris kernel for filling only the multiplicative
caps has rank zero: each such boundary has first Betti number one more
than its cap, and the regular exterior's first rank is its number of
holes. The groups are free, so that kernel is zero. Subtract the cap
boundary sections to obtain the surjection $H_2(V)\to H_2(W)$,
with inverse given by the quotient of exterior coordinates. Its
generators are $s$ and the unfilled fixed sweeps, with the one
primitive delta-sum relation, whose fibre coefficient is zero before
the ordinary clutch is filled. Thus $H_1(W)$ and $H_2(W)$ are
free and restriction in degree-two cohomology from $W$ to $V$
is injective over the integers and over each finite field used below.
The exterior retracts over a graph. Its homology sequence has the
free fibre coinvariant group \eqref{maps:eq:MP_local_16} and a subgroup of the free
degree-one cylinder group, so $H_2(V)$ is free as well.
The same graph presentation gives free $H_1(V)$; consequently
$H^2(V;\mathbb Z)$ is free.

\item \emph{Every finite parameter and the integral homology.}

At a finite site let $m=d_i k>0$, $k>0$, and use the actual fixed
lattice $\mathbb Z\delta\oplus\mathbb Z\ell$, with
$\chi(\ell)=1$. It is given explicitly by

\begin{equation}
       \ell=(-b(A-I)^{-1}g,0,1)^t.                    \hypertarget{tag-Q20P}{}\label{maps:eq:MP_local_21}
\end{equation}

The fixed columns are primitive and the finite coinvariant section
$\beta_{0,\rm fin}$ satisfies

\begin{equation}
 [\delta]=q_{\rm fin}\beta_{0,\rm fin},\quad
 N_m\beta_{0,\rm fin}=(m/q_{\rm fin})\delta,\quad N_m\ell=m\ell,
 \quad-\Omega(\beta_{0,\rm fin},\ell)=b/q_{\rm fin}.
                                                               \hypertarget{tag-Q20Q}{}\label{maps:eq:MP_local_22}
\end{equation}

These are the finite-section maps of \ref{maps:src:D.2.26}, applied to the literal matrices \eqref{maps:eq:MP_local_13}, with their integer inverse coordinates; here $q_{\rm fin}=1$ at $II$ and $q_{\rm fin}=3$ at $IV$.

Write the power $v_{\rm fin}=n_{\rm fin}\ell+r\delta$. The affine action
$x\mapsto Tx+v_{\rm fin}/m$ has that power. Its invariant gamma translation is
$n_{\rm fin}/m$, so every proper power is free when $\gcd(m,n_{\rm fin})=1$.
For any other affine vector $t_0$ with $N_m t_0=v_{\rm fin}$, set
$\pi_{\rm av}=N_{d_i}/d_i$. The difference $t_0-v_{\rm fin}/m$ lies in
$\ker\pi_{\rm av}$. Translation by

\begin{equation}
          (I-T)^{-1}(t_0-v_{\rm fin}/m)\quad\hbox{on }\ker\pi_{\rm av}
                                                               \hypertarget{tag-Q20R}{}\label{maps:eq:MP_local_23}
\end{equation}

conjugates one action to the other; opposite translation is its inverse.
The map retains the same actual root-disc boundary parametrization.
Write $t_{\rm sw}(u)$ for the marked degree-two sweep of a fixed lattice
vector $u$, with the same exterior sweep coordinates as above.
Changing a peripheral lift by an integral $\lambda$ instead changes
the power by $N_m\lambda$ and changes a fixed sweep by
$t_{\rm sw}'(u)=t_{\rm sw}(u)+\Omega(u,\lambda)s$. These are the full norm changes;
an unrestricted $m\lambda$ formula is not used at a finite site.

The Seifert presentation, with the ordinary clutch included, is

\begin{equation}
 \langle\gamma,c_0\mid[\gamma,c_0]=1,
          c_0^m=\gamma^{n_{\rm fin}},\ c_0=\gamma^e\rangle
       =\mathbb Z/(me-n_{\rm fin}).                             \hypertarget{tag-Q20S}{}\label{maps:eq:MP_local_24}
\end{equation}

Eliminate $c_0$ by $c_0\mapsto\gamma^e$; inclusion of $\gamma$
is the inverse presentation map. Thus the entire simply connected
range is $me-n_{\rm fin}=\pm1$. This includes arbitrary integer $e,r$, every
positive multiple $m$ of the stated order, and all affine lifts \eqref{maps:eq:MP_local_23}.
It implies $\gcd(m,n_{\rm fin})=1$. With no finite site, the presentation is
$\mathbb Z/e$, so the exact range is $e=\pm1$.

For completeness retain the full degree-two Thom kernel before
elimination. In coordinates $(s,D_\delta=t_{\rm sw}(\delta),L_\ell=t_{\rm sw}(\ell))$, it is

\begin{equation}
\begin{gathered}
 \mathsf C_1=(bn_{\rm fin}/q_{\rm fin},m/q_{\rm fin},0),\\
 \mathsf C_2=(-br,0,m),\quad \mathsf C_3=(0,r,n_{\rm fin}).
\end{gathered}                    \hypertarget{tag-Q20T}{}\label{maps:eq:MP_local_25}
\end{equation}

The actual normal meridian is $m\mu-v_{\rm fin}$: on the root boundary the
maps are $[y,\tau]\mapsto[y+(\tau/m)v_{\rm fin},e^{2\pi i\tau/m}]$ and
the inverse subtracts $(\tau/m)v_{\rm fin}$. Cutting a central-loop Thom torus
into its $m$ sectors gives area $-\lambda\wedge v_{\rm fin}$ and sweep
$N_m\lambda$; rejoining the sectors is its inverse chain map.
The root generator has zero area and sweep $v_{\rm fin}$. These give exactly
the three columns in \eqref{maps:eq:MP_local_25}, not only a multiple of their span.

Choose $\alpha_{\rm fin} m+\eta_{\rm fin} n_{\rm fin}=1$. Then

\begin{equation}
\begin{gathered}
 K_0=\alpha_{\rm fin} \mathsf C_2+\eta_{\rm fin} \mathsf C_3
     =(-\alpha_{\rm fin} br,\eta_{\rm fin} r,1),\\
 \mathsf C_2-mK_0=-\eta_{\rm fin} q_{\rm fin} r \mathsf C_1,\qquad
 \mathsf C_3-n_{\rm fin}K_0=\alpha_{\rm fin} q_{\rm fin} r \mathsf C_1.
\end{gathered}                            \hypertarget{tag-Q20U}{}\label{maps:eq:MP_local_26}
\end{equation}

These identities provide inverse integer substitutions. The ordinary
cap has normal vector $-e\gamma$ and full kernel

\begin{equation}
          t_{\rm sw}(u)-\Omega(u,-e\gamma)s,\qquad t_{\rm sw}(-e\gamma).
                                                               \hypertarget{tag-Q20V}{}\label{maps:eq:MP_local_27}
\end{equation}

Its collar translation is $(y,t)\mapsto(y+te\gamma,t)$, with opposite
translation inverse. In particular it gives $D_\delta=-be\,s$. Eliminating
$L_\ell$ by \eqref{maps:eq:MP_local_26} gives $L_\ell=br(\alpha_{\rm fin}+\eta_{\rm fin} e)s$, and the remaining
relation is

\begin{equation}
          (b/q_{\rm fin})(n_{\rm fin}-me)s=\pm2p\,s=0.                  \hypertarget{tag-Q20W}{}\label{maps:eq:MP_local_28}
\end{equation}

Without a finite site it is $be\,s=\pm2p\,s=0$. Original additives
in this table have zero line multidegree and one component orbit; their
actual product/quotient maps remove both fixed sweeps with zero fibre
coefficient. The complete primitive boundary maps of \ref{maps:src:D.2.26} give no further relation. The degree-one bridge has rank $b_1(X)=0$ and is free, so
it vanishes. Therefore $H_2(X)=\mathbb Z/(2p)\langle s\rangle$.

The index-one cusp has Euler contribution two, and the other retained
caps have contribution zero. Integral duality and universal
coefficients now give

\begin{equation}
 H_0=H_6=\mathbb Z,\quad H_2=H_3=\mathbb Z/(2p),
 \quad H_1=H_4=H_5=0,\quad b_3=0.                   \hypertarget{tag-Q20X}{}\label{maps:eq:MP_local_29}
\end{equation}

In particular $H^2(X;\mathbb Z)=0$, so the complex tangent bundle
has $c_1=0$ and $X$ is spin. Its spin structure is unique since
$H^1(X;\mathbb Z/2)=0$. For $\varpi=2,p$, let $\alpha_{\varpi}(s)=1$.
The full columns \eqref{maps:eq:MP_local_25}--\eqref{maps:eq:MP_local_27} have fibre coefficients divisible by
$2p$: $bn_{\rm fin}/q_i=2pn_{\rm fin}$, $br$, and $be$. Thus the functional one
on $s$ and zero on each exterior sweep descends to both $\alpha_{\varpi}$,
through every actual finite twist, with no omitted sweep coefficient.

\item \emph{One undoubled line on the common piece.}

By \eqref{maps:eq:MP_local_10}, the form $\Omega\bmod2$ has the quadratic refinement

\begin{equation}
       \mathcal Q(\lambda)=\lambda_1\lambda_2+
                    \lambda_1+\lambda_2\pmod2.       \hypertarget{tag-Q20Y}{}\label{maps:eq:MP_local_30}
\end{equation}

It is one on every nonzero vector in the elliptic plane modulo two.
Every $A_i$ acts invertibly on that plane, and the elliptic gamma
column $bg_i$ is even. Hence every actual $T_i$ preserves
$\mathcal Q$. Its polar form is $\Omega\bmod2$. The factors

\begin{equation}
 J_\lambda(x,\zeta_R)=
 (x+\lambda,(-1)^{\mathcal Q(\lambda)}e^{\pi i\Omega(\lambda,x)}\zeta_R)
                                                               \hypertarget{tag-Q20Z}{}\label{maps:eq:MP_local_31}
\end{equation}

satisfy their exact cocycle law; $J_{-\lambda}$ is the inverse.
The connection $d-\pi i\Omega(x,dx)$ descends and has Chern class
$\Omega$. The lifts $(x,\zeta_R)\mapsto(T_i x,\zeta_R)$ preserve it and its
origin frame, have the stated inverse matrices and ordered product
identity. Thus they define one line on the actual exterior $V$.
Its degree on $s$ is one, and its degree on every marked fixed-loop
sweep is zero: on $x=t u$, $\Omega(u,u)=0$, and the scalar-one
matrix lift closes the line with no base winding.

The finite scalar lift is

\begin{equation}
\begin{gathered}
 T_{\tau_{\rm sc}}(x,\zeta_R)=(x+\tau_{\rm sc},e^{\pi i\Omega(\tau_{\rm sc},x)}\zeta_R),\\
 \tau_{\rm sc}=\delta/N_{\rm sc},\qquad
 T_{\tau_{\rm sc}}^{-1}=T_{-\tau_{\rm sc}}.
\end{gathered}                   \hypertarget{tag-Q210}{}\label{maps:eq:MP_local_32}
\end{equation}

Its commutator with $J_\lambda$ is
$e^{-2\pi i b\lambda_\gamma/N_{\rm sc}}$. It descends exactly for $N_{\rm sc}\mid b$,
and its $N_{\rm sc}$-th power is $J_\delta$, since $\mathcal Q(\delta)=0$.
Every actual monodromy fixes $\delta$ and preserves $\Omega$, so
these lifts commute with its return. They give the required $C_2$,
$C_p$ and $C_{2p}$ lifts; only the first two will be used for
localization. Proper subgroups have their own fixed elliptic curves.

Here are the maps identifying this line with the actual holomorphic
period and cap charts. For the actual nonzero short period $\omega_s$, use
$z_E=z_{\rm add}/\omega_s$, with inverse
$z_{\rm add}=\omega_s z_E$. Put $\tau_E=\omega_l/\omega_s$
and $\zeta_G=G_{\rm lift}/\omega_s$; the physical coordinate
is $z=e^{2\pi i z_E}$, with inverse logarithm modulo that period.
The theta ratio $\vartheta(z_E-a\zeta_G)/\vartheta(z_E)$,
with reciprocal inverse,
identifies the actual divisor line $\mathcal O(aG-O)$ with its
Poincare factors. Its short multiplier is one and its long multiplier
is $e^{2\pi i a\zeta_G}$, by shifting the theta sum. A forward lift
of $-bG$ has a constant multiplier in this trivialization: the ratio
with the constant-factor lift is a holomorphic function on the compact
elliptic curve. This retains the actual base-dependent nonzero units.

In the physical bases $S=(u,\delta)$, $L=(v,\gamma)$, the period
matrix and pairing are

\begin{equation}
 \Pi_{\rm ph}=\begin{pmatrix}\tau_E&-b\zeta_G\\
                                  a\zeta_G&\rho_L\end{pmatrix},
 \qquad D=\operatorname{diag}(a,-b),\qquad
                  D\Pi_{\rm ph}=(D\Pi_{\rm ph})^t.    \hypertarget{tag-Q211}{}\label{maps:eq:MP_local_33}
\end{equation}

Write $x=SU+Lv_{\rm per}$, $Z=U+\Pi_{\rm ph}v_{\rm per}$. The invertible gauge

\begin{equation}
 G_D(U,v_{\rm per})=\exp\{\pi i[-U^tDv_{\rm per}-v_{\rm per}^tD\Pi_{\rm ph}v_{\rm per}+Z_1]\}
                                                               \hypertarget{tag-Q212}{}\label{maps:eq:MP_local_34}
\end{equation}

has reciprocal inverse. The short factors become one and the long
factors become

\[
 (-1)^a h_1^{(1-a)/2}z^{-a},\qquad
              h_G^{-b/2}R^{b/2}v^b.
\]

The further gauge $z^{(a-1)/2}$, with inverse $z^{-(a-1)/2}$, gives
the literal maps

\begin{equation}
\begin{aligned}
 h_{\rm hol}(z,v,\zeta)&=(h_1z,h_G^av,(-1)^a z^{-a}\zeta),\\
 F(z,v,\zeta)&=(h_G^{-b}z,Rv,h_G^{-ab/2}R^{b/2}v^b\zeta),\\
 J_{\rm inv}(z,v,\zeta)&=(z^{-1},v^{-1},z^{-a}\zeta).
\end{aligned}                                             \hypertarget{tag-Q213}{}\label{maps:eq:MP_local_35}
\end{equation}

All exponents of the actual unit functions are integral. Composition
gives $h_{\rm hol}F=Fh_{\rm hol}$, $J_{\rm inv}^2=1$, $J_{\rm inv}FJ_{\rm inv}=F^{-1}$, $J_{\rm inv}h_{\rm hol}J_{\rm inv}=h_{\rm hol}^{-1}$.
For the last comparison the extra logarithmic phase is the integer
$a$, so its exponential is one. The inverse of $h_{\rm hol}$ has line factor
$(-1)^a h_1^{-a}z^a$, and the inverse of $F$ has line factor
$h_G^{ab/2}R^{b/2}v^{-b}$, evaluated on the target coordinates.
The scalar map in \eqref{maps:eq:MP_local_35} is $(z,v,\zeta)\mapsto(z,e^{2\pi i/N_{\rm sc}}v,\zeta)$.
These gauges, reciprocal inverses and theta factors identify the
actual line and all its period gluings.

\item \emph{Extension on the prescribed cusp fan.}

In \eqref{maps:eq:MP_local_18} the actual physical valuation matrix is $C$; the positive
Tate valuation of $G$ is $-\kappa$. The prescribed fan's primitive short basis is $SC$; the map
from its coordinates to physical short coordinates is $v\mapsto Cv$,
with inverse $x\mapsto C^{-1}x$.
Use $x=SCU+Lv_{\rm per}$, $Z=U+\Pi v_{\rm per}$, where

\begin{equation}
 K=C^tD=
 \begin{pmatrix}a\nu&ab\kappa\\
       ab\kappa&b(ab\kappa^2-1)/\nu\end{pmatrix},
 \quad\Pi=C^{-1}\Pi_{\rm ph},\quad K\Pi=D\Pi_{\rm ph}.
                                                               \hypertarget{tag-Q214}{}\label{maps:eq:MP_local_36}
\end{equation}

Both displayed pairings are symmetric. The actual period matrix is
$\Pi=(\log q)/(2\pi i)I+B_{\rm hol}(q)$, with its full holomorphic unit matrix.
The short refinement vector is $\varepsilon=(1,0)^t$, by \eqref{maps:eq:MP_local_10}.
The Gaussian

\begin{equation}
 G_K=\exp\{\pi i[-U^tKv_{\rm per}-v_{\rm per}^tK\Pi v_{\rm per}+\varepsilon^t Z]\}
                                                               \hypertarget{tag-Q215}{}\label{maps:eq:MP_local_37}
\end{equation}

has reciprocal inverse. For a long period $h\in\mathbb Z^2$, its
factor is

\begin{equation}
\begin{gathered}
 (-1)^{\eta_{\rm quad}^th}
 e^{-2\pi i h^tKZ-\pi i h^tK\Pi h+\pi i\varepsilon^t\Pi h},\\
 \eta_{{\rm quad},j}=\mathcal Q(L\mathbf e_j).
\end{gathered}                \hypertarget{tag-Q216}{}\label{maps:eq:MP_local_38}
\end{equation}

Substitution gives $G_K(U,v_{\rm per},\Pi+I)=G_K(U+v_{\rm per},v_{\rm per},\Pi)$, including the
units. For an integral exponent column $\lambda$, write
$X_{\rm C}^{\lambda}=\exp(2\pi i\lambda^tZ)$ for the toric monomial,
with $X_{{\rm C},j}=X_{\rm C}^{\mathbf e_j}$.
The factor is its exact unit times
$q^{(-h^tKh+\varepsilon^th)/2}X_{\rm C}^{-Kh}$.

The integer vertex support on the prescribed antidiagonal triangulation
is

\begin{equation}
          \psi(v)=\tfrac12(v^tKv-\varepsilon^tv).     \hypertarget{tag-Q217}{}\label{maps:eq:MP_local_39}
\end{equation}

Its integrality follows from $K_{jj}\equiv\varepsilon_j\pmod2$.
On every unimodular triangle its affine slope is integral. If $m_\sigma$
is that affine interpolation, the regular line coordinate is
$\zeta_\sigma=\chi_{\rm tor}^{m_\sigma}\zeta_C$, with inverse multiplication
by $\chi_{\rm tor}^{-m_\sigma}$. Period translation changes the divisor by
the reciprocal long factor's divisor, so \eqref{maps:eq:MP_local_38}, with its units, is
the actual line gluing. The physical-coordinate support gives the same
line: putting $j_{\rm cusp}=-\kappa$,

\begin{equation}
\begin{gathered}
 \varphi(x,y)=\tfrac12(-a\rho x^2-2abj_{\rm cusp}xy+b\nu y^2+ax),\\
 \varphi(Cv)=-\psi(v)+L_0v,\\
 L_0=\tfrac12(a\operatorname{row}_1C-\varepsilon^t).
\end{gathered}
                                                               \hypertarget{tag-020D}{}\label{maps:eq:D.MP4}
\end{equation}

$L_0$ is integral. The dense maps are $\zeta=X_{\rm C}^{L_0}\zeta_C$
and its reciprocal. Their regular coordinates coincide after \eqref{maps:eq:MP_local_39},
so they extend as mutually inverse line maps on each chart.
Every nonzero holomorphic unit in \eqref{maps:eq:MP_local_35} remains in the period gluing.

The physical scalar vector is $S\mathbf e_2=\delta$ and has actual short coordinates

\begin{equation}
       d_\delta=C^{-1}\mathbf e_2=(-b\kappa,\nu)^t,
             \qquad Kd_\delta=(0,-b)^t.              \hypertarget{tag-Q218}{}\label{maps:eq:MP_local_41}
\end{equation}

For the vertical edge, the two adjacent toric charts are

\begin{equation}
\begin{gathered}
 (x,y,w)=(q/(X_{{\rm C},1}X_{{\rm C},2}),X_{{\rm C},2},X_{{\rm C},1}),\\
 (x',y',w')=(q/X_{{\rm C},2},X_{{\rm C},1}X_{{\rm C},2},X_{{\rm C},1}^{-1}).
\end{gathered}                 \hypertarget{tag-Q219}{}\label{maps:eq:MP_local_42}
\end{equation}

Their overlap is $(x',y',w')=(wx,wy,w^{-1})$, with inverse
$(w'x',w'y',(w')^{-1})$. The rational curve $C_f$ is the vanishing
of the two normal coordinates. Its normal transition has degree
$-1$ in each factor. There is one such vertical-edge orbit; its
endpoints belong to different triangle orbits. A nonzero lattice
translation cannot stabilize a finite edge.

The endpoint circle weights are

\begin{equation}
 (b\kappa-\nu,\nu,-b\kappa),\qquad
          (-\nu,\nu-b\kappa,b\kappa).               \hypertarget{tag-Q21A}{}\label{maps:eq:MP_local_43}
\end{equation}

Modulo either $2$ or $p$, $d_\delta=(0,\nu)$ and $\nu$ is a
unit. Thus the finite subgroup fixes exactly $C_f$ in this cap.
The other edge directions have distinct nonzero reductions, and the
toric fixed vertices are the endpoints already contained in $C_f$.
The normal weights are both one modulo two, and $(-\nu,\nu)$
modulo $p$. These are the actual scalar weights; no change of one
local generator without its effect at the other component is made.

On the triangle $0,\mathbf e_1,\mathbf e_2$, the slope in \eqref{maps:eq:MP_local_39} is

\[
          m_+=((K_{11}-1)/2,K_{22}/2).
\]

On its adjacent triangle it is
$m_-=((2K_{12}-K_{11}-1)/2,K_{22}/2)$. The Gaussian scalar phase
contributes $\varepsilon^td_\delta/2$. Write $d_\delta=(d_1,d_2)^t$. The total line character is

\begin{equation}
 c_C=\tfrac12(K_{11}d_1+K_{22}d_2)
       =\frac{b}{2}\{a\kappa(b\kappa-\nu)-1\},
 \quad c_C\equiv1+\kappa\pmod2,
 \quad c_C\equiv0\pmod p.                            \hypertarget{tag-Q21B}{}\label{maps:eq:MP_local_44}
\end{equation}

The two slopes give the same character modulo $2p$; their difference
is paired with $-b\kappa$. The positive Cartier wall gives

\begin{equation}
           \deg(\mathcal L|C_f)=K_{11}-K_{12}
                              =a(\nu-b\kappa).       \hypertarget{tag-Q21C}{}\label{maps:eq:MP_local_45}
\end{equation}

It is odd and reduces to $a\nu$ modulo $p$. These calculations use the prescribed fan under $C$ and $C^{-1}$.

\item \emph{Every rank-one period and the square-line comparison.}

For \eqref{maps:eq:MP_local_17}, choose integers $\theta_N,\sigma,x_{\rm long},y_{\rm long}$ satisfying

\[
 n_0\theta_N+A_0\sigma=1,\quad B_0=b\kappa/n_0\in\mathbb Z,
       \quad d_0x_{\rm long}+B_0y_{\rm long}=1.
\]

The two actual primitive changes are

\begin{equation}
\begin{gathered}
 C_S=\begin{pmatrix}n_0&\sigma\\-A_0&\theta_N\end{pmatrix},
 \quad C_L=\begin{pmatrix}x_{\rm long}&-B_0\\y_{\rm long}&d_0\end{pmatrix},\\
 C_S^{-1}=\begin{pmatrix}\theta_N&-\sigma\\A_0&n_0\end{pmatrix},
 \quad C_L^{-1}=\begin{pmatrix}d_0&B_0\\-y_{\rm long}&x_{\rm long}\end{pmatrix}.
\end{gathered}                                                               \hypertarget{tag-Q21D}{}\label{maps:eq:MP_local_46}
\end{equation}

Both have determinant one. Put $S_N=(u,\delta)C_S=(l,p_{\rm gen})$,
$L_N=(v,\gamma)C_L=(h_{\rm gen},w_0)$. In the actual period basis
$H_N=(l,h_{\rm gen},p_{\rm gen},w_0)$, $Th_{\rm gen}=h_{\rm gen}+l$, and the other three
vectors are fixed. The period exponential and the inverse logarithm
in this basis identify the same original germ, section, contraction
and collars. They retain all holomorphic units in these inverse germ maps.

The full, possibly nonsymmetric pairing is

\begin{equation}
\begin{aligned}
 K_N&=S_N^t\Omega L_N,\quad
 (K_N)_{11}=a n/d_0^2=:N_0,\quad (K_N)_{12}=0,\\
 (K_N)_{22}&=-b d_0^2/n,
 \quad C_S^{-1}\mathbf e_2=(-\sigma,n_0)^t.
\end{aligned}
\hypertarget{tag-Q21E}{}\label{maps:eq:MP_local_47}
\end{equation}

The Gaussian formula for this basis uses

\[
 -U^tK_Nv_{\rm per}-v_{\rm per}^tK_N^t\Pi_Nv_{\rm per}+\epsilon_N^t(U+\Pi_Nv_{\rm per}),
                  \quad \epsilon_{N,j}=\mathcal Q(S_N\mathbf e_j).
\]

The symmetry is $K_N^t\Pi_N=(K_N^t\Pi_N)^t$. Thus the long factor
has $K_N^t$, and the entire off-diagonal pairing and unit matrix
are retained. The node support is
$(N_0j^2-\mathcal Q(l)j)/2$, an integer since
$N_0\equiv\mathcal Q(l)\pmod2$. Its affine interpolants and
reciprocal monomials extend the same line through the exact Tate maps, with character $X_{\rm N}$,

\begin{equation}
       x_+=X_{\rm N}/q^j,\quad x_-=q^{j+1}/X_{\rm N},
       \qquad X_{\rm N}=x_+q^j,\quad q=x_+x_- .               \hypertarget{tag-Q21F}{}\label{maps:eq:MP_local_48}
\end{equation}

For comparison with \eqref{maps:eq:MP_local_35}, write the original physical valuation
columns as $\alpha_0 v_0,\beta_0v_0$, with $v_0$ primitive and
$\gcd(\alpha_0,\beta_0)=1$. The line-support integer is
$k_{\rm supp}=a(v_0)_1/\alpha_0$, because
$-a\beta_0(v_0)_1=b\alpha_0(v_0)_2$ and coprimality imply that
divisibility. Its physical support is
\begin{equation}
 \phi_1(j)=-\frac{k_{\rm supp}}{2}j(j-\alpha_0),\qquad k_{\rm supp}=\frac{a(v_0)_1}{\alpha_0}.
\hypertarget{tag-020E}{}\label{maps:eq:D.MP5}
\end{equation}

In these rows the physical first-period content is $\alpha_0=d_0$,
equal to one or three. It is odd, so $j(j-\alpha_0)$ is even for
every integer vertex $j$; this proves the asserted integrality of
the physical support, including the nonprimitive column.
Take $B_S=C_S$, whose first column is
$v_0=(n_0,-A_0)^t$ and whose inverse is displayed in
\eqref{maps:eq:MP_local_46}. The difference from the Gaussian
support is the integral linear form
$(a\operatorname{row}_1B_S-\epsilon_N^t)/2$. Its monomial and reciprocal
give the dense comparison and inverse, and the regular Cartier
coordinates then agree. The transverse monomial stays invertible on
\eqref{maps:eq:MP_local_48}. This includes the nonprimitive first period at the unselected
$I_3$ for $a=3$; no primitivity of that original column is assumed.

All these line extensions restrict to \eqref{maps:eq:MP_local_31} on $V$, through the
specified gauges. They therefore glue to a single equivariant line
$\mathcal L$ on $W$. By the exterior quotient and the sweep values, for both
$\varpi=2,p$,

\begin{equation}
                \alpha_{\varpi}|_W=\rho_{\varpi} c_1(\mathcal L).   \hypertarget{tag-Q21G}{}\label{maps:eq:MP_local_49}
\end{equation}

There is also an exact comparison to the doubled-line construction of \ref{maps:src:D.2.36} used for odd-primary localization. On the real cover,
tensor-square of \eqref{maps:eq:MP_local_31} has factors $e^{2\pi i\Omega(\lambda,x)}$,
because its quadratic signs square to one. The fibre tensor map gives
$\mathcal L^{\otimes2}\simeq\mathcal D$, where $\mathcal D$ is
this doubled line; no square root is chosen. The finite actions also
square exactly to its prequantum action.

In the cusp Gaussian coordinates, the doubled support is $v^tKv$.
The tensor-square support $2\psi(v)$ from
\eqref{maps:eq:MP_local_39} differs from it by $-\varepsilon^tv$; the dense
square-line coordinate differs by the opposite monomial $X_{\rm C}^{\varepsilon}$.
Their regular Cartier coordinates are therefore identical. The same
calculation at a node uses $\mathcal Q(l)j$ in place of
$\varepsilon^tv$. The reciprocal monomial is the inverse map.
This proves the tensor comparison on the completed caps and their
collars, with the actual finite linearizations and unit functions.

\item \emph{The full $C_2$ and $C_p$ fixed sets.}

At a rank-one site, the surviving elliptic lattice is
$\ker f/\mathbb Z l$. In the primitive coordinates \eqref{maps:eq:MP_local_46}, the
scalar period $\delta/\varpi$ has surviving translation
$(n_0/\varpi)p_{\rm gen}$. It fixes that elliptic curve exactly when $\varpi\mid n_0$.
Otherwise it is a nonzero finite translation of an elliptic curve and
has no fixed point; subtracting its actual surviving lattice coordinate
is the inverse quotient map. The profile lists and this divisibility criterion give exactly one
$C_2$-fixed rank-one place, $I_2$, and exactly one $C_p$-fixed
rank-one place, $I_p$, in each row. Both have $d_0=1$. Call their
elliptic double curves $E_2,E_p$.

At $I_2$, its local $\kappa$ is odd. The basis below is the evaluation of \eqref{maps:eq:MP_local_46} at
$\sigma=1$, $\theta_N=(1-a\kappa)/2$, $x_{\rm long}=1$ and
$y_{\rm long}=0$.
A useful exact basis is

\begin{equation}
 l=2u-a\kappa\delta,\quad
 p_0=u+(1-a\kappa)\delta/2,\quad
 h_0=v,\quad w_0=\gamma-(b\kappa/2)v.                \hypertarget{tag-Q21H}{}\label{maps:eq:MP_local_50}
\end{equation}

The short and long matrices, and their inverses, are

\begin{equation}
 \begin{pmatrix}2&1\\-a\kappa&(1-a\kappa)/2\end{pmatrix},
 \begin{pmatrix}(1-a\kappa)/2&-1\\a\kappa&2\end{pmatrix};
 \qquad
 \begin{pmatrix}1&-b\kappa/2\\0&1\end{pmatrix},
 \begin{pmatrix}1&b\kappa/2\\0&1\end{pmatrix}.        \hypertarget{tag-Q21I}{}\label{maps:eq:MP_local_51}
\end{equation}

The generic Bezout basis in \eqref{maps:eq:MP_local_46} is related to \eqref{maps:eq:MP_local_50} by
$p_{\rm gen}=p_0+k_{\rm bas} l$, $k_{\rm bas}=(\sigma-1)/2\in\mathbb Z$, and
$h_{\rm gen}=h_0+y_{\rm long} w_0$. The inverse subtracts these multiples.
The equalities follow from $2\theta_N+a\kappa\sigma=1$ and
$x_{\rm long}+(b\kappa/2)y_{\rm long}=1$. Thus the generic inverse matrices
and \eqref{maps:eq:MP_local_50} describe the same chart and line with their integral phases.

The pairing in \eqref{maps:eq:MP_local_50} is

\[
          \begin{pmatrix}2a&0\\a&-b/2\end{pmatrix}.
\]

Therefore $\deg\mathcal L|E_2=-b/2$, which is odd. Since
$\delta=-l+2p_0$, $\mathcal Q(p_0)=1$, and the first nodal slope
is $a$, the actual half-period character is $1-a=0\pmod2$.
Both normal coordinates in $x_+x_-=q$ have sign character. Hence

\begin{equation}
 \deg\mathcal L|E_2=-b/2,\qquad c_{E_2}=0\pmod2.     \hypertarget{tag-Q21J}{}\label{maps:eq:MP_local_52}
\end{equation}

At $I_p$, put $r_{\rm Tate}=-\kappa_i$, the genuine positive-valuation Tate
component, and define the physical weight by
\begin{equation}
 w_{\rm Tate}=(a r_{\rm Tate})^{-1}\pmod p.
\hypertarget{tag-020F}{}\label{maps:eq:D.MP6}
\end{equation}
Formula \eqref{maps:eq:MP_local_46} gives
$\sigma a\kappa_i=1\pmod p$, so $w_{\rm Tate}=-\sigma$. The normal pair
is $(w_{\rm Tate},-w_{\rm Tate})$ and

\begin{equation}
            \deg\mathcal L|E_p=-b/p.                 \hypertarget{tag-Q21K}{}\label{maps:eq:MP_local_53}
\end{equation}

The doubled-line first slope is $N_0=ap$; its finite character is
$-\sigma N_0=0\pmod p$. The exact tensor-square comparison above
therefore gives $2c_{E_p}=0\pmod p$, so $c_{E_p}=0$, with inverse
division by two in $\mathbb F_p$. At the cusp its undoubled character
is also zero modulo $p$, directly from \eqref{maps:eq:MP_local_44}.

There are no fixed points in any additive cap or the ordinary cap.
For an original cap, the actual zero-multidegree product and contracting
line quotient give a free scalar circle: a nonzero iterate has
nonunit contracting norm, and the zero iterate cannot equal a
nonintegral scalar period. At a finite cap, a scalar fixed point for
any nontrivial subgroup would require $j n_{\rm fin}/m\in\mathbb Z$ with
$0\le j<m$. Equation \eqref{maps:eq:MP_local_24} forces $j=0$; the remaining primitive
scalar translation is nontrivial. The root-disc argument gives the
same conclusion away from the core. These statements are unchanged
under the exact affine conjugacies \eqref{maps:eq:MP_local_23}. The ordinary torus cap has
primitive $\delta$, so its finite scalar translations are free.

Consequently the complete fixed sets, for every allowed parameter, are

\begin{equation}
             X^{C_2}=C_f\amalg E_2,\qquad
             X^{C_p}=C_f\amalg E_p.                 \hypertarget{tag-Q21L}{}\label{maps:eq:MP_local_54}
\end{equation}

Each prime subgroup has free complement. Both share the prescribed
rational curve, while their elliptic curves lie in different actual
caps. No free-complement assertion is made for $C_{2p}$.

\item \emph{A single global Borel character comparison.}

For $\varpi=2,p$, $\alpha_{\varpi}$ has a full-circle Borel extension. The
only possible outgoing differentials in total degree two have targets
$H^2(BS^1;H^1(X;\mathbb F_{\varpi}))$ and $H^3(BS^1;\mathbb F_{\varpi})$,
both zero. Restrict this extension to $C_{\varpi}$, calling it $U_{\varpi}$.
The line $\mathcal L$ has the actual finite lift \eqref{maps:eq:MP_local_32}, and \eqref{maps:eq:MP_local_49}
identifies their ordinary restrictions on $W$.

Choose an endpoint of $C_f$. It is fixed by the whole scalar circle,
even when the full circle rotates the interior of $C_f$. The based
scalar isotopy acts trivially on \eqref{maps:eq:MP_local_20}, and its section gives

\begin{equation}
       \pi_1(W_{C_{\varpi}})=(\mathbb Z\times F_{r_B})\times C_{\varpi}.  \hypertarget{tag-Q21M}{}\label{maps:eq:MP_local_55}
\end{equation}

Any section from a finite fixed point has the same induced inclusion
of $C_{\varpi}$: its component in the torsion-free first factor is zero,
as its image has order dividing $\varpi$. Thus all integral degree-one
classes of $W$ restrict with zero point weight on every fixed curve.
Every finite degree-one class of $W$ lifts integrally, since $H_1(W)$
is free.

The difference $U_{\varpi}-\rho_{\varpi}c_1^{C_{\varpi}}(\mathcal L)$ has zero ordinary
restriction. In the universal-cover sequence it lies in group
degree-two cohomology, because the cover's $H^1$ is zero. The ordinary
group component is zero by the injective low-degree map for $W$.
Its remaining terms are one base class and products of the actual
integral degree-one classes with the group degree-one generator.
Equation \eqref{maps:eq:MP_local_55} shows that the base coefficient is a single common
constant on all fixed components. Subtract that one global base class;
addition is the inverse. The mixed term on the elliptic component
is retained, rather than discarded by a pointwise normalization.

Write $\tau_{\rm Bor}\in H^1(BC_2;\mathbb F_2)$, and write
$\eta_{\rm Bor}\in H^1(BC_p;\mathbb F_p)$, $t_p\in H^2(BC_p;\mathbb F_p)$
with $\eta_{\rm Bor}^2=0$. After this one global character adjustment,

\begin{equation}
\begin{aligned}
 U_2|C_f&=\alpha_2|C_f+(1+\kappa)\tau_{\rm Bor}^2,&
 U_2|E_2&=\alpha_2|E_2+\tau_{\rm Bor}\xi_2,\\
 U_p|C_f&=\alpha_p|C_f,&
 U_p|E_p&=\alpha_p|E_p+\eta_{\rm Bor}\xi_p,
\end{aligned}                                             \hypertarget{tag-Q21N}{}\label{maps:eq:MP_local_56}
\end{equation}

where $\xi_{\varpi}\in H^1(E_{\varpi};\mathbb F_{\varpi})$ lifts integrally and has
square zero. The sphere has no degree-one mixed term. This proves the
comparison globally on the actual common piece, with its group section
and inverse character change. The equality is not inferred by choosing
unrelated lifts at its fixed curves.

\item \emph{The full coprime homotopy and bordism maps.}
Put $K_{\varpi}=K(\mathbb Z/\varpi,2)$ for $\varpi=2,p$. The cyclic-group
resolution makes the positive homology of $K(\mathbb Z/\varpi,1)$
zero over $\mathbb Z[1/\varpi]$. The path-loop Serre sequence, whose
total space is contractible, then makes the positive homology
of $K_{\varpi}$ zero over that ring. Exactness of localization implies
that each positive integral homology element is killed by a
power of $\varpi$. For elements of the two factors, a Bezout identity
between their annihilating powers makes their tensor zero.
Likewise $\operatorname{Tor}(\mathbb Z/2^{n_{\rm Tor}},B)=B[2^{n_{\rm Tor}}]=0$ for every positive integer $n_{\rm Tor}$ and
$p$-primary $B$; finite direct sums and filtered colimits give
the assertion for all the positive homology groups.

The reduced integral Kunneth sequence therefore makes
$K_2\wedge K_p$ acyclic. The based CW inclusion
\begin{equation}
 i:K_2\vee K_p\longrightarrow K_2\times K_p
\hypertarget{tag-020I}{}\label{maps:eq:D.MP9}
\end{equation}
is an integral homology equivalence of simply connected spaces.
Every relative homotopy group vanishes: a first nonzero one
would be the corresponding zero relative homology group by
relative Hurewicz. The inverse can be constructed on cells.
Start with the identity of the wedge; injectivity on homotopy
groups extends its retraction across each attaching cell, and
surjectivity adjusts that extension so that its composite with
$i$ is homotopic to the original cell map relative to its
boundary. The homotopy-extension property assembles the choices
over successive skeleta. This gives a based map $r_{\rm CW}$ with
$r_{\rm CW}i=\mathrm{id}$ and $i r_{\rm CW}\simeq\mathrm{id}$ relative to the wedge.
Thus \eqref{maps:eq:D.MP9} is a full CW homotopy equivalence,
with the stated inverse at that homotopy scope. The coefficient maps
below may be implemented by the Eilenberg--Mac Lane functor on
degree-two chain complexes in simplicial abelian groups; in that
model their two composites are literally the identity.

The coefficient maps
\[
 \theta(n)=(n\bmod2,n\bmod p),\qquad
 \theta^{-1}(x,y)=px+(1-p)y\pmod{2p}
\]
are exact inverses. Apply the Eilenberg--Mac Lane construction
to these degree-two coefficient maps and then use $i,r_{\rm CW}$.
The finite wedge axiom for the spin-bordism Thom spectrum gives
\begin{equation}
 \Omega^{\mathrm{Spin}}_6(K(\mathbb Z/(2p),2))
  \cong
 \Omega^{\mathrm{Spin}}_6(K_2)\oplus
 \Omega^{\mathrm{Spin}}_6(K_p).
\hypertarget{tag-020A}{}\label{maps:eq:D.MP10}
\end{equation}
Here $\Omega^{\mathrm{Spin}}_6(\mathrm{pt})=0$, by
\ref{maps:src:D.1.11} for coefficient group zero: its injective pair
$(P,h)$ has zero target.
The forward map in \eqref{maps:eq:D.MP10} is the pair of
coefficient projections. Its inverse is the sum of the two
factor inclusions followed by $K(\theta^{-1},2)_*$.
An automorphism on the odd factor, multiplication by $u_{\rm aut}$, lifts
to the unit $p+(1-p)u_{\rm aut}$ modulo $2p$; its inverse is
$p+(1-p)u_{\rm aut}^{-1}$.

\item \emph{The full fundamental-class invariant.}

For the involution the normal representation on each fixed curve $F_{\rm fix}$ is
the sign representation of complex rank two. Its equivariant Euler
class is

\begin{equation}
 \tau_{\rm Bor}^4+\tau_{\rm Bor}^2\rho_2 c_1(\nu_{F_{\rm fix}})+\rho_2 c_2(\nu_{F_{\rm fix}})
                         =\tau_{\rm Bor}^4.                    \hypertarget{tag-Q21O}{}\label{maps:eq:MP_local_57}
\end{equation}

Here $c_1(\nu_{F_{\rm fix}})=-c_1(TF_{\rm fix})$ is even and $c_2(\nu_{F_{\rm fix}})=0$ on a
curve. Remove invariant normal discs. The complement is free, so its
Borel space has the homotopy type of a finite-dimensional quotient;
after inverting $\tau_{\rm Bor}$ it has zero cohomology. Thom excision and
self-intersection give the exact restriction inverse
$\sum_{F_{\rm fix}} i_{F_{\rm fix}*}(z/e_{F_{\rm fix}})$. Applying it to \eqref{maps:eq:MP_local_56},

\begin{equation}
 \langle\alpha_2^3,[X]_2\rangle
     =(1+\kappa)\deg(\mathcal L|C_f)=1+\kappa\pmod2. \hypertarget{tag-020G}{}\label{maps:eq:D.MP7}
\end{equation}

The elliptic term is zero because $\xi_2^2=0$ and every product of
curve degree greater than two vanishes. The pure-base term on $C_f$
integrates to zero; its mixed cube coefficient is $3=1\pmod2$.
The common shift removed in \eqref{maps:eq:MP_local_56} changes the sum by the sum of the
two odd ordinary degrees, hence by zero as well.

For the odd subgroup, the two normal characters on each component are
$\lambda_{\rm nor},-\lambda_{\rm nor}$, where $\lambda_{\rm nor}=\nu$ on $C_f$ and
$\lambda_{\rm nor}=w_{\rm Tate}$ on $E_p$, up to their displayed order. The distinct
characters split the normal bundle into its eigensubbundles. Its Euler
class has leading term $-\lambda_{\rm nor}^2t_p^2$, which is a unit after
inverting $t_p$. The same free-complement argument gives the inverse
restriction. Both restrictions $U_p$ in \eqref{maps:eq:MP_local_56} have cube zero: their
ordinary terms square to zero and the mixed terms have square zero.
Therefore

\begin{equation}
                    \langle\alpha_p^3,[X]_p\rangle=0. \hypertarget{tag-Q21P}{}\label{maps:eq:MP_local_59}
\end{equation}

For the complete integral class at order two, Cartan's elementary
complex has basic factor $\Gamma(x_{\rm Car})\otimes\Lambda(y_{\rm Car})$,
with generators in degrees $2$ and $3$ and differential
$d_{\rm Car}y_{\rm Car}=2x_{\rm Car}$. Its degree-six boundary is
$d_{\rm Car}(y_{\rm Car}\gamma_2(x_{\rm Car}))=6\gamma_3(x_{\rm Car})$,
so its two-primary degree-six group has order two.
The next elementary factor is
$\Lambda(z_{\rm Car})\otimes\Gamma(v_{\rm Car})$, in generator degrees
$5$ and $6$, with $d_{\rm Car}v_{\rm Car}=2z_{\rm Car}$; it has no
degree-six homology. Products with a positive basic class begin
in degree seven, and their Tor terms begin in degree eight.
All further factors begin later. This is the complete integral
model in \cite[Theorem~1, pp.~11-09--11-10; Theorem~4 and its
corollary, p.~11-13; Appendix~1, pp.~11-22--11-23]{manCartan1955EML}.
Thus $H_6(K(\mathbb Z/2,2);\mathbb Z)=\mathbb Z/2$, generated
by $\gamma_3(1)$. Reduction of the positive complex
$\mathbb{CP}^3$ represents this divided-power generator:
the positive degree-two generator of $\mathbb{CP}^{\infty}$ has
third divided power equal to its positive degree-six generator,
and integer reduction induces the divided-power map.
The mod-two cube evaluates to one on this image. Evaluation
therefore has inverse $a\mapsto a\gamma_3(1)$.
For $p=5,7$, \ref{maps:src:D.1.10} gives the corresponding
cyclic group and cube detector.
The preceding coprime homotopy maps then give, with
$s=1\in\mathbb Z/(2p)$,

\begin{equation}
 H_6(K(\mathbb Z/(2p),2);\mathbb Z)
        =\mathbb Z/(2p)\langle\gamma_3(s)\rangle .    \hypertarget{tag-Q21Q}{}\label{maps:eq:MP_local_60}
\end{equation}

The positive generator is the image of the complex-oriented
$\mathbb{CP}^3$ under integer reduction. Its two primary projections
are the generators just specified. Thus the pair of ordinary cube
evaluations detects the complete integral group; the inverse maps
are explicit below. The projections to $\mathbb Z/2$ and
$\mathbb Z/p$ send $s\mapsto1$, and send $\gamma_3(s)\mapsto \gamma_3(1)$. Use the models' given markings for these coefficient coordinates.
Their inverse on homology uses the inclusions

\begin{equation}
\begin{aligned}
 e_{2,\rm coef}&\mapsto p s,\qquad e_{p,\rm coef}\mapsto(1-p)s,\\
 \gamma_3(e_{2,\rm coef})&\mapsto p^3\gamma_3(s)=p\gamma_3(s),\\
 \gamma_3(e_{p,\rm coef})&\mapsto(1-p)^3\gamma_3(s)
                               =(1-p)\gamma_3(s).
\end{aligned}
\hypertarget{tag-Q21R}{}\label{maps:eq:MP_local_61}
\end{equation}

These are integer Chinese-remainder maps with inverse reduction in
the two factors; the coefficients are idempotent modulo $2p$.
Equations \eqref{maps:eq:D.MP7}--\eqref{maps:eq:MP_local_59} consequently give the full invariant

\begin{equation}
               \mu(X)=p(1+\kappa\bmod2)\gamma_3(s). \hypertarget{tag-Q21S}{}\label{maps:eq:MP_local_62}
\end{equation}

No coordinate has been lost by reducing only modulo $p$.

\item \emph{The refined Pontryagin class and exact Tate conversion.}

On a fixed curve $F_{\rm fix}$ for $C_p$, let the ordinary normal roots be
$x_1,x_2$. They and the tangent root are curve classes. The equivariant
normal roots are $x_1+\lambda_{\rm nor} t_p,x_2-\lambda_{\rm nor} t_p$, and

\begin{equation}
 p_1(TX)|F_{\rm fix}=-2e_{F_{\rm fix}}.                                    \hypertarget{tag-Q21T}{}\label{maps:eq:MP_local_63}
\end{equation}

Indeed $x_1^2=x_2^2=x_1x_2=c_1(TF_{\rm fix})^2=0$ on a curve; expanding
both sides gives $2\lambda_{\rm nor}^2t_p^2+2\lambda_{\rm nor} t_p(x_1-x_2)$.
The identity allows unequal ordinary normal degrees. Hence the exact
localized pairing is

\begin{equation}
 \langle\alpha_p\,\rho_p p_1(TX),[X]_p\rangle
    =-2\left(a(\nu-b\kappa)-{b\over p}\right)
    =2\left({b\over p}-a\nu\right)\pmod p.          \hypertarget{tag-Q21U}{}\label{maps:eq:MP_local_64}
\end{equation}

The degree-one mixed term from $E_p$ integrates to zero. Naturality
of the refined invariant under reduction to $\mathbb Z/p$, followed
by the quarter map of \ref{maps:src:D.1.11} and the Bezout inverse of four modulo $p$,
gives

\begin{equation}
            P_p={b/p-a\nu\over2}\in\mathbb F_p.     \hypertarget{tag-Q21V}{}\label{maps:eq:MP_local_65}
\end{equation}

The action conversion retains every generator. In the genuine Tate
marking at the unique $I_p$, $r_{\rm Tate}=-\kappa_i$ and $w_{\rm Tate}=(a r_{\rm Tate})^{-1}$.
The zero-degree equivariant integration of $U_p$ gives

\begin{equation}
\begin{gathered}
 {a\nu\over\nu^2}-{b/p\over w_{\rm Tate}^2}=0,
 \quad w_{\rm Tate}^2={\nu(b/p)\over a},\\
 P_p={b/p\over2}\left(1-{\nu^2\over w_{\rm Tate}^2}\right).
\end{gathered}    \hypertarget{tag-Q21W}{}\label{maps:eq:MP_local_66}
\end{equation}

All denominators are units in these rows. The subgroup automorphism
$z\mapsto z^{\nu^{-1}}$, with inverse $z\mapsto z^\nu$, changes
both normal pairs: cusp to $(-1,1)$, elliptic to
$(w_{\rm Tate}/\nu,-w_{\rm Tate}/\nu)$. The Borel character and degree-one generator
are both pulled back by the same exponent. Thus \eqref{maps:eq:MP_local_65}--\eqref{maps:eq:MP_local_66} use one
action on the actual common space. The zero criterion is
$w_{\rm Tate}^2=\nu^2$, not a rescaling that changes only one component.


For comparison with the marking that evaluates to one on $C_f$, put
$z_{\rm coh}=(a\nu)^{-1}\alpha_p$. Its homology generator is
$g_{z_{\rm coh}}=(a\nu)s$, with inverse coordinate evaluation by
$z_{\rm coh}$. The simultaneous action conversion gives
\begin{equation}
 z_{\rm coh}[E_p]=-\frac{w_{\rm Tate}^2}{\nu^2},\qquad
 \langle z_{\rm coh}^3,[X]_p\rangle=0,\qquad
 \langle z_{\rm coh}\,\rho_p p_1,[X]_p\rangle
       =-2\left(1-\frac{w_{\rm Tate}^2}{\nu^2}\right).
\hypertarget{tag-020H}{}\label{maps:eq:D.MP8}
\end{equation}
The corresponding refined coefficient is $(w_{\rm Tate}^2/\nu^2-1)/2$.
Multiplication by $a\nu$, with inverse $(a\nu)^{-1}$, recovers
the primitive coefficient $P_p$. The physical weights are $w_{\rm Tate}=4$
at the order-five rows with $a=1$, $w_{\rm Tate}=3$ there when $a=3$, and
$w_{\rm Tate}=4$ at every order-seven row. The selected $\nu$ is three
on $56(1,10)$ and $56/II(1,10)$, and one on the other rows.
Thus the common normalized weights $w_{\rm Tate}/\nu$ are three in all
nonzero order-five cases and four in its zero case. This accounts
for every assignment in \eqref{maps:eq:D.MP2}.

For the two-primary coordinate, use the precise primary quotient on
\cite[Sections~3.7 and~3.12, pp.~803, 807]{manZhubr75}:

\begin{equation}
 \Omega^{\mathrm{Spin}}_6(K(\mathbb Z/2,2))
   =\mathbb Z\langle\epsilon_{\rm bord},\zeta_{\rm bord}\rangle/
                        \langle\epsilon_{\rm bord},2\zeta_{\rm bord}\rangle,
       \quad P(\zeta_{\rm bord})=1,\quad h(\zeta_{\rm bord})=\gamma_3(1).   \hypertarget{tag-Q21X}{}\label{maps:eq:MP_local_67}
\end{equation}

The quotient sends $(A,B)\mapsto B\bmod2$; its inverse chooses the
reduction of $\zeta_{\rm bord}=[\mathbb CP^3\to\mathbb CP^\infty]$.
Its cubic is one. Therefore the marked classifying map
$X\to K(\mathbb Z/2,2)$ has refined coefficient exactly \eqref{maps:eq:D.MP7}.
The map $P$ is natural, so this is the actual mod-two projection
of $\mathfrak p(X)$. No division by four in $\mathbb Z/(2p)$
is used.

Set $t=1+\kappa\pmod2$. The unique coefficient $P_{2p}$ is
specified by

\begin{equation}
        P_{2p}\equiv t\pmod2,\qquad
        P_{2p}\equiv P_p\pmod p.                    \hypertarget{tag-Q21Y}{}\label{maps:eq:MP_local_68}
\end{equation}

The inverse CRT map is $ (t,P_p)\mapsto p t+(1-p)P_p\pmod{2p}$.
Its composites are identity on both coordinates. This proves the
full marked data in the theorem:

\begin{equation}
 (H_2,b_3,\mathfrak p,\mu)
  =\left(\mathbb Z/(2p),0,P_{2p}s,p t\gamma_3(s)\right).
                                                               \hypertarget{tag-Q21Z}{}\label{maps:eq:MP_local_69}
\end{equation}

\item \emph{Exact connected-sum marking and oriented classification.}

The surgery traces of Theorem~\ref{maps:src:D.2.35} use exactly
the embeddings and orientations in Definition~\ref{maps:def:minimal-two-data}.
Their incoming generator maps to $e_{2,\rm coef}$ and their outgoing second-homology
map is an isomorphism. The respective incoming cubes are zero and one.
The reduced CP$^3$ spin-bordism quotient above therefore gives the
full marked invariants \eqref{maps:eq:MP_local_1}, including the refined coordinate.

The oriented connected sum $M_t\# Z_{p,\epsilon}$ is simply connected,
spin and has $b_3=0$, with

\begin{equation}
\begin{gathered}
 H_2=\mathbb Z/2\langle e_{2,\rm coef}\rangle\oplus
                  \mathbb Z/p\langle e_{p,\rm coef}\rangle,\\
 \mathfrak p=t e_{2,\rm coef}+\epsilon e_{p,\rm coef},\qquad
 \mu=t\gamma_3(e_{2,\rm coef}).
\end{gathered}                            \hypertarget{tag-Q220}{}\label{maps:eq:MP_local_70}
\end{equation}

These are the actual addition maps: puncture both manifolds in discs
where their classifying maps are constant, attach the oriented neck,
and use the bordism with one one-handle from the disjoint union to
the connected sum. The classifying map extends by the constant map
on that handle. Its boundary fundamental classes add and $P$ is a
bordism homomorphism. The homology inclusions of the summands and
collapse to them are inverse coordinates. The finite cross terms
between the two torsion factors vanish by the integer Kunneth maps,
whose tensor and Tor groups have coprime orders. This is also the
same exact CRT equivalence of the two Eilenberg--Mac Lane markings
used in \eqref{maps:eq:MP_local_61}.

Here is the marking that compares \eqref{maps:eq:MP_local_70} to \eqref{maps:eq:MP_local_69}. If $P_p\ne0$, put

\begin{equation}
\begin{aligned}
 k_{\rm CRT}&=P_{2p}-tp\pmod{2p},\\
 \Phi(e_{2,\rm coef})&=p s,\quad\Phi(e_{p,\rm coef})=k_{\rm CRT}s,\\
 \Phi^{-1}(n s)&=(n\bmod2)e_{2,\rm coef}+
                       (P_p^{-1}n\bmod p)e_{p,\rm coef}.
\end{aligned}
\hypertarget{tag-Q221}{}\label{maps:eq:MP_local_71}
\end{equation}

Then $k_{\rm CRT}$ is even and $k_{\rm CRT}\bmod p=P_p\ne0$, so it has order $p$.
Both compositions in \eqref{maps:eq:MP_local_71} are identity. Moreover

\begin{equation}
 \Phi(t e_{2,\rm coef}+e_{p,\rm coef})=P_{2p}s,\quad
 K(\Phi,2)_*(t\gamma_3(e_{2,\rm coef}))=tp^3\gamma_3(s)
                                        =tp\gamma_3(s).        \hypertarget{tag-Q222}{}\label{maps:eq:MP_local_72}
\end{equation}

If $P_p=0$, take $\Phi(e_{p,\rm coef})=(1-p)s$ and replace $P_p^{-1}$
in the inverse by one. Then $P_{2p}=tp\pmod{2p}$, and the same
identities hold with $\epsilon=0$. These give every generator
rescaling and its inverse. The concrete maps in the table are
$(e_{2,\rm coef},e_{p,\rm coef})\mapsto(5s,2s)$ for the nonzero order-five cases,
$(5s,6s)$ for its zero case, and $(7s,4s)$ for order seven.
Their inverse odd coefficients are respectively three, one and two.

Use the marked recognition theorem. Theorem~5 of \cite[p.~796]{manZhubr75} gives the injection

\[
 (P,h):\Omega^{\mathrm{Spin}}_6(K(G,2))
                    \longrightarrow G\oplus H_6(K(G,2);\mathbb Z).
\]

The construction in \cite[Theorem~3, p.~795; Section~4.1, p.~809]{manZhubr75} preserves the
specified $H_2\cong G$, kills the free part of $H_3$, and its
marked uniqueness proof yields a degree-positive diffeomorphism from
equality of the resulting bordism classes. Here both free ranks vanish, and both manifolds
have the unique spin structure required. By \eqref{maps:eq:MP_local_69}--\eqref{maps:eq:MP_local_72}, their marked
$(P,h)$ agree. Injection and marked uniqueness give the oriented
diffeomorphism in \eqref{maps:eq:D.MP1}, inducing precisely $\Phi$ on second homology.
Its inverse induces \eqref{maps:eq:MP_local_71}'s inverse. Thus the connected-sum conclusion
is an actual oriented identification with its prescribed homology map.

Every type listed in \eqref{maps:eq:D.MP2} is realized.
For a row without a finite cap choose $e=1$. At a finite cap choose
$m=d_i$, $e=1$, $n_{\rm fin}=m-1$, $r=0$, and the affine vector $v_{\rm fin}/m$.
It is free, has $me-n_{\rm fin}=1$, and is an actual allowed construction.
Equations \eqref{maps:eq:MP_local_23}--\eqref{maps:eq:MP_local_69} prove the same result for every other allowed
finite multiplicity, numerator, residue and affine vector.

Finally, an automorphism of $\mathbb Z/(2p)$ is multiplication by a
unit $u_{\rm cls}$. It sends $P_{2p}$ to $u_{\rm cls}P_{2p}$ and the full $\mu$
coefficient to $u_{\rm cls}^3\mu_{2p}$. Modulo two, $u_{\rm cls}=1$, so $t$ is
preserved. Modulo $p$, a unit preserves zero versus nonzero. These
two properties distinguish the five listed types. Conversely, any
two rows with the same pair $(t,\epsilon)$ are related by the unit
$u_{\rm cls}\equiv1\pmod2$, $u_{\rm cls}\equiv P_p'/P_p\pmod p$ when $P_p\ne0$,
or by one on the odd factor when both are zero. Its inverse is the
corresponding inverse unit, and it preserves the full $(P,\mu)$.
The same marked recognition gives their oriented diffeomorphism.
This proves exactly three and two types, with the counts in \eqref{maps:eq:D.MP2}.


\end{enumerate}
\end{proof}

\section{Three-primary and residual smooth types}
\subsection{Three-primary smooth types}

\subsubsection{Order three}

\begin{definition}[Minimal order-three data]\hypertarget{tag-Q223}{}\label{maps:def:minimal-three-data}
{\normalfont\scriptsize tag Q223}\par
Use the ordered elliptic data of
Definition~\ref{mapsIII:def:elliptic-marked-data} and choose a pattern
below. Use the prescribed primitive Mumford fan at every multiplicative
site, the listed additive cap and ordinary clutch $e\gamma$,
$e\in\mathbb Z$. The logarithmic sites are precisely the listed
additive sites. At a finite site of linear monodromy order $d_i$,
allow every multiplicity $m=d_i k$, $k>0$, and every free affine lift.
Selected positions are $i-1$ in the input definition's peripheral order.

\begin{equation}
\begin{array}{c|r|r|l|l}
\text{OS/profile}&a&b&\text{zero positions}&\text{additive choice}\\\hline
\mathrm{OS56}&5&3&2&\text{none}\\
\mathrm{OS56}&10&3&3,4&\text{none}\\
\mathrm{OS56}/\mathrm{II}&5&3&0&II\text{ original or log}\\
\mathrm{OS56}/\mathrm{III}&1&6&0&III\text{ log}\\
\mathrm{OS56}/\mathrm{III}&5&6&3&III\text{ log}\\
\mathrm{OS56}/\mathrm{III}&10&3&3&III\text{ original}.
\end{array}                                                   \hypertarget{tag-025C}{}\label{maps:eq:D.T31}
\end{equation}

\begin{equation}
 M_{3,0}=\bigl((S^2\times S^4)\setminus
          \operatorname{int}(S^2\times D^4)\bigr)
               \mathop{\cup}_{S^2\times S^3}(D^3\times S^3),     \hypertarget{tag-025D}{}\label{maps:eq:D.T32}
\end{equation}
where the embedded sphere represents three times the generator of
$H_2(S^2\times S^4;\mathbb Z)$ and has its oriented normal framing.
Take the graph of the degree-three map
$z\mapsto z^3$ on $\mathbb {CP}^1=S^2$ and an equatorial embedding
$S^2\hookrightarrow S^4$.
These choices are the \emph{minimal order-three data}.
\end{definition}

\Needspace{6\baselineskip}
\begin{theorem}[Order-three smooth type]\hypertarget{tag-02LF}{}\label{maps:src:D.2.39}
{\normalfont\scriptsize tag 02LF}\par
Every simply connected threefold $X$ of
Definition~\ref{maps:def:minimal-three-data} is orientation-preservingly
diffeomorphic to $M_{3,0}$ with the marking defined in the proof.
The complete Zhubr data are
\begin{equation}
 H_2=H_3=\mathbb Z/3,\qquad b_3=0,\qquad
 \mathfrak p=0,\qquad
 \mu=0\in H_6(K(\mathbb Z/3,2);\mathbb Z)=\mathbb Z/9.          \hypertarget{tag-025E}{}\label{maps:eq:D.T33}
\end{equation}

\end{theorem}


\begin{proof}
\begin{enumerate}
\item \emph{The integral marking and the cap maps.}
Write
\[
 \Lambda=\mathbb Z\langle e_1,e_2,\delta,\gamma\rangle,
 \qquad J_{\rm ell}=\begin{pmatrix}0&1\\-1&0\end{pmatrix},
 \qquad\omega=a e_1^*\wedge e_2^*-b\delta^*\wedge\gamma^*,
 \qquad\chi=\gamma^*.
\]
The positive scalar circle is translation by $t\delta$, $t\in\mathbb R/\mathbb Z$.
The arrow $\gamma$ is the forward contracting arrow covering $-bG$.
All curves have their complex orientations. The local data at the selected
place in the seven marked rows underlying \eqref{maps:eq:D.T31} are
\[
\begin{array}{c|r|r|r|r|r}
\text{profile}&a&b&\text{zero}&\nu&\kappa\\\hline
\mathrm{OS56}/\mathrm{II}&5&3&0&2&-1\\
\mathrm{OS56}/\mathrm{III}&1&6&0&5&-4\\
\mathrm{OS56}/\mathrm{III}&5&6&3&1&-1\\
\mathrm{OS56}/\mathrm{III}&10&3&3&1&-1\\
\mathrm{OS56}&5&3&2&2&-1\\
\mathrm{OS56}&10&3&3&1&0\\
\mathrm{OS56}&10&3&4&1&-1
\end{array}
\]
The first row has the two additive choices in \eqref{maps:eq:D.T31}. In particular
$b\in\{3,6\}$, $a\nu=30/b$, and $3\nmid a\nu$.
In the defining elliptic marking, a multiplicative monodromy and section
cocycle are
\[
 A_i-I=\nu_i u_i u_i^tJ_{\rm ell},\qquad g_i=\kappa_i u_i,
 \qquad\det(u_i,v_i)=1,
\]
where the first nonzero entry of $u_i$ is positive. At a finite elliptic
matrix set $z_i=(A_i-I)^{-1}g_i$; at a multiplicative place set
$z_i=(\kappa_i/\nu_i)v_i$. The four-period monodromies are
\[
\begin{gathered}
 \beta_i=-a g_i^tJ_{\rm ell}A_i,\qquad c_i^{\rm mon}=\beta_i b z_i+w_i,
 \\
 T_i=\begin{pmatrix}A_i&0&bg_i\\\beta_i&1&c_i^{\rm mon}\\0&0&1\end{pmatrix},
 \qquad
 w_i=\begin{cases}1/\nu_i&i\text{ selected},\\0&\text{otherwise}.
 \end{cases}
\end{gathered}
\]
Their inverse maps are
\[
 T_i^{-1}=\begin{pmatrix}
 A_i^{-1}&0&-bA_i^{-1}g_i\\
 -\beta_iA_i^{-1}&1&b\beta_iA_i^{-1}g_i-c_i^{\rm mon}\\0&0&1
 \end{pmatrix}.
\]
Substitution gives integral matrices, $T_i^*\omega=\omega$,
$T_i\delta=\delta$, $\chi T_i=\chi$, and the ordered product $\prod_iT_i=I$.
At the cap under consideration write $A=A_i$, $g=g_i$,
$\beta=\beta_i$ and $T_{\rm mon}=T_i$; at a multiplicative
place write $u=u_i$ and $v=v_i$.
On a cut base the period map is
$x\bmod\Lambda\mapsto\Pi(q)x\bmod\Pi(q)\Lambda$, with inverse the group
logarithm modulo that lattice. Here $\Pi$ consists of the elliptic and line
period functions, including their holomorphic units. For an integral basis
$H_{\rm change}$, replacing $\Pi$ by $\Pi H_{\rm change}$ replaces coordinates by $H_{\rm change}^{-1}x$.
The inverse basis change reverses this map, preserving the base coordinate,
zero section, and line frames.

Let $\mathcal E$ be the regular torus bundle after all cap discs have been
removed. We use the integral period-graph and cap maps of
\ref{maps:src:D.2.26}, in the following coordinates. In wedge order
$12,13,14,23,24,34$, put
\[
 B_2=[\wedge^2T_1-I\ \cdots\ \wedge^2T_N-I],\qquad
 \omega=(a,0,0,0,0,-b).
\]
Choose $a x_{\rm Bez}+b y_{\rm Bez}=1$ and
$s_0=x_{\rm Bez}e_{12}-y_{\rm Bez}e_{34}$. Integer row and column
reduction gives five unit nonzero Smith factors and an integer matrix
$R_2$ with
\[
 B_2R_2=I-s_0\omega,\qquad\omega B_2=0,\qquad\omega(s_0)=1.
\]
Thus the fibre quotient and its inverse are
\[
 [v_{\rm wedge}]\longmapsto\omega(v_{\rm wedge}),\qquad n\longmapsto[ns_0].
\]
The chain bounding $v_{\rm wedge}-\omega(v_{\rm wedge})s_0$ is the sum of the monodromy cylinders
with coefficients $R_2v_{\rm wedge}$. Write $s=[s_0]$.

Include the removed ordinary disc as the final site, with monodromy $I$.
For $M_{\rm sites}$ sites set $D=[T_1-I\ \cdots\ T_{M_{\rm sites}-1}-I]$. The first boundary
maps are block inclusions, and the last is
\[
 \begin{gathered}
 j_{M_{\rm sites}}(x)=-(T_2\cdots T_{M_{\rm sites}-1}x,\ldots,T_{M_{\rm sites}-1}x,x),\\
  Dj_{M_{\rm sites}}(x)=(I-T_{M_{\rm sites}}^{-1})x.
 \end{gathered}
\]
Choose each fixed basis to begin with $\delta$. These bases are
$(\delta,u)$ at the cusp, $(\delta,u,w_0)$ at a rank-one place,
$(\delta,\ell)$ at a finite additive cap, $(\delta,\ell_0)$ at an
original additive cap, and $(\delta,e_1,e_2,\gamma)$ at the ordinary disc.
The vectors $w_0,\ell,\ell_0$ are given below. If
$B=[j_1F_1\ \cdots\ j_{M_{\rm sites}}F_{M_{\rm sites}}]$, then $\operatorname{im}B=\ker D$
integrally, and
\[
 \ker B=\mathbb Z\sum_i e_{\delta,i}.
\]
An explicit inverse on $\ker D$ is obtained by reducing
$U_DDV_D=\operatorname{diag}(I_3,0)$, denoting the final columns of $V_D$
by $V_0$, and solving $BS_{\rm graph}=V_0$ by integer operations. It is
$R_{\rm graph}=S_{\rm graph}(V_D^{-1})_{\text{rows}>3}$, with
$BR_{\rm graph}y=y$ for $Dy=0$. The reduction of $B$ has only unit
nonzero factors and one relation. The fixed loops sweep the boundary
circles; their addition and subtraction are the period triangles
with vertices $0,u,u+v$. Projection to the graph, this inverse, and
the fibre quotient give
\[
 H_2(\mathcal E;\mathbb Z)=\mathbb Zs\oplus
 \left(\bigoplus_i\operatorname{Fix}(T_i)\right)
             /\mathbb Z\sum_i\delta_i.
\]
The last relation is the scalar circle over the cut base. Write
$c_i^{\rm base}$ for the marked base meridians. The ordinary
clutch $\prod c_i^{\rm base}=\gamma^e$ adds
$-\omega(\delta,e\gamma)=be$ to it.

At an original additive cap put $z_{\rm ell}=(A-I)^{-1}g$, $h_{\rm line}=az_{\rm ell}\in\mathbb Z^2$,
$z_Q=bz_{\rm ell}$, and $\lambda_{\rm line}=-h_{\rm line}^tJ_{\rm ell}$. Its line-frame maps are
\[
 S_P(x,d,g_0)=(x,d+\lambda_{\rm line} x,g_0),\qquad
 S_P^{-1}(x,d',g_0)=(x,d'-\lambda_{\rm line} x,g_0).
\]
The identity $\lambda_{\rm line}(A-I)+\beta=0$ removes the line shear. Zero
multidegree trivializes the smooth complex line bundle on that elliptic neighbourhood.
The contracting quotient is a mapping torus with a positive roof $r_{\rm roof}(y)$.
Choose $0<k_{\rm roof}<(\max r_{\rm roof})^{-1}$ and a nonnegative smooth function $b_y$ supported away from the
ends of $[0,r_{\rm roof}(y)]$ with integral one. The increasing function
\[
 f_y(t)=\int_0^t\{k_{\rm roof}+(1-k_{\rm roof}r_{\rm roof}(y))b_y(v)\}\,dv
\]
and its inverse identify that roof with a unit interval, with matching
end derivatives. Translation by the identity-component N\'eron logarithm
times the interval parameter, with a cutoff at its ends, removes that
part of the return; the negative translation is the inverse. The
remaining finite translation has one orbit on the rational components.
Reciprocal line gauges preserve their node constants.
For type $II$ put $q_{\rm orig}=1$, and for type $III$ put $q_{\rm orig}=2$. The cap has
$H_2=\mathbb Z[C_{\rm orig}]\oplus\mathbb Z[\delta\gamma]$, with elliptic fibre
$q_{\rm orig}[C_{\rm orig}]$. In the $S_P$ coordinates its fibre map is
\[
 f_{2,\rm orig}(y)=\bigl(q_{\rm orig}(y_{12}+z_{Q,2}y_{14}-z_{Q,1}y_{24}),y_{34}\bigr).
\]
The first coordinate is obtained by $x\mapsto x+z_Qx_\gamma$;
subtraction is the inverse. Choose
$q_{\rm orig}a_0+q_{\rm orig}z_{Q,2}b_0-q_{\rm orig}z_{Q,1}c_0=1$. A section sends $[C_{\rm orig}]$ to
$\wedge^2S_P^{-1}(a_0e_{12}+b_0e_{14}+c_0e_{24})$ and
$[\delta\gamma]$ to $e_{34}$. Component transitivity gives these integers.
The fixed lattice has basis
\[
 \delta,\qquad\ell_0=(-q_{\rm orig}z_{Q,1},-q_{\rm orig}z_{Q,2},0,q_{\rm orig}).
\]
The two sweeps bound the origin disc and the $q_{\rm orig}$-fold forward orbit tube,
closed using the same logarithm and line frame. At original $III$ the
second vector is $(-3,9,0,2)$.

At a primitive rank-one cap, in the ordered basis $(l,h_{\rm per},p_0,w_0)$ below,
the fibre maps and their sections are
\[
 f_{1,N}(x)=(x_2,x_3,x_4),\qquad
 f_{2,N}(y)=(y_{12},y_{23},y_{24},y_{34}),
\]
with sections inserting zeros in the omitted positions and then applying
the period basis. The normalization is one ruled component, with its
zero and infinity sections identified along the elliptic core. Its
normalization/conductor sequence gives these generators; fixed sweeps
bound the normal toric chains. At the cusp, use
$H_C=(S_{\rm sh}Ce_1,L_{\rm long}e_1,S_{\rm sh}Ce_2,L_{\rm long}e_2)$, with $S_{\rm sh},L_{\rm long},C$ defined below. The maps are
\[
 f_{1,C}(x)=(x_2,x_4),\qquad
 f_{2,C}(y)=(y_{12},y_{14}-y_{23},y_{24},y_{34}).
\]
On a target vector $(a_{C,1},a_{C,2},a_{C,3},a_{C,4})$, the section
of $f_{2,C}$ inserts $(a_{C,1},0,a_{C,2},0,a_{C,3},a_{C,4})$
and applies $\wedge^2H_C$.
Its primitive kernel is generated by $e_{13},e_{14}+e_{23}$, the local
Wang image. The cusp charts below give the corresponding normal chains.

\item \emph{Finite caps and arbitrary affine parameters.}
For finite $II$ and $III$, respectively, write $(d_{\rm fin},q_{\rm fin})=(6,1)$ and $(4,2)$.
At this finite site set $z_{\rm fin}=z_i=(A-I)^{-1}g$.
The vector $\ell=(-b z_{{\rm fin},1},-b z_{{\rm fin},2},0,1)$ is integral. The power and
multiplicity of a simply connected finite filling have the form
\[
 \begin{gathered}
 v_{\rm fin}=n_{\rm fin}\ell+r_{\rm fin}\delta,\qquad m_{\rm fin}=d_{\rm fin}k_{\rm fin},\\
 k_{\rm fin}\in\mathbb Z_{>0},\quad r_{\rm fin}\in\mathbb Z,\\
 m_{\rm fin}e-n_{\rm fin}=\varepsilon_{\rm det}\in\{1,-1\}.
 \end{gathered}
\]
In particular $\gcd(m_{\rm fin},n_{\rm fin})=1$. The return is $x\mapsto T_{\rm mon}x+v_{\rm fin}/m_{\rm fin}$.
For a lift $T_{\rm mon}x+t_{\rm aff}$ with $N_{m_{\rm fin}}t_{\rm aff}=v_{\rm fin}$, where $N_{m_{\rm fin}}=\sum_{j=0}^{m_{\rm fin}-1}T_{\rm mon}^j$,
put $P_{\rm fin}=N_{m_{\rm fin}}/m_{\rm fin}$. On $\ker P_{\rm fin}$, the operator $I-T_{\rm mon}$ is invertible;
solving $(I-T_{\rm mon})h_{\rm aff}=t_{\rm aff}-v_{\rm fin}/m_{\rm fin}$ gives the conjugacy $x\mapsto x+h_{\rm aff}$, with inverse
$x\mapsto x-h_{\rm aff}$. It commutes with scalar translation.
For base-dependent holomorphic lifts on the good root disc, let
$\sigma_{\rm root}(u_{\rm root})=\zeta_{\rm root}u_{\rm root}$ and let $A_{\rm hol}(u_{\rm root}):E_{u_{\rm root}}\to E_{\sigma_{\rm root} u_{\rm root}}$ be the linear
deck map on its Lie bundle. Put
\[
 \begin{gathered}
 (\mathcal Rf_{\rm hol})(u_{\rm root})=A_{\rm hol}(\sigma_{\rm root}^{-1}u_{\rm root})f_{\rm hol}(\sigma_{\rm root}^{-1}u_{\rm root}),\\
 \mathcal N=\sum_{j=0}^{m_{\rm fin}-1}\mathcal R^j,\qquad\mathcal R^{m_{\rm fin}}=I.
 \end{gathered}
\]
Two lifts with the same integral power have terms $b_{\rm lift,0},b_{\rm lift,1}$ with
$\mathcal N(b_{\rm lift,1}-b_{\rm lift,0})=0$. The holomorphic section
\[
 h_{\rm hol}=-\frac{1}{m_{\rm fin}}\sum_{j=0}^{m_{\rm fin}-1}j\mathcal R^j(b_{\rm lift,1}-b_{\rm lift,0})
\]
satisfies $(I-\mathcal R)h_{\rm hol}=b_{\rm lift,1}-b_{\rm lift,0}$. Translation by $h_{\rm hol}$, inverted by
translation by $-h_{\rm hol}$, gives the conjugacy on the disc and collar.
Changing a period representative changes the power by $N_{m_{\rm fin}}\lambda$;
aligning those representatives gives the same comparison.

Choose a vertical coinvariant $\beta_{\rm fin}$ and a period $y_{\rm coin}$ with
\[
 \begin{gathered}
 q_{\rm fin}\beta_{\rm fin}+(T_{\rm mon}-I)y_{\rm coin}=\delta,\\
 N_{d_{\rm fin}}\beta_{\rm fin}=(d_{\rm fin}/q_{\rm fin})\delta,\\ N_{d_{\rm fin}}\ell=d_{\rm fin}\ell,
 \qquad-\omega(\beta_{\rm fin},\ell)=b/q_{\rm fin}.
 \end{gathered}
\]
For $II$, take $\beta_{\rm fin}=\delta$, $\ell=\gamma$. For $III$ with
$a=1$ or $5$ and $b=6$, take
\[
 \ell=(-3,9,0,1),\quad
 \beta_{\rm fin}=(0,-1,(1-a)/2,0),\quad y_{\rm coin}=(-5,11,0,0).
\]
The inverse coinvariant coordinates are $(-3a,-a,2,0)$ and $(0,0,0,1)$.
They give a primitive section of the quotient by $(T_{\rm mon}-I)\Lambda$.
An integer right inverse of $[T_{\rm mon}-I\mid\beta_{\rm fin}\mid\ell]$ expresses any
change of these representatives as an explicit $(T_{\rm mon}-I)$-boundary;
negating the boundary reverses the change.

Let $c_{\rm bd}$ be the forward mapping-torus root loop and
$\mu_{\rm mer}=[c_{\rm bd}]$ its boundary homology class.
Writing $D_\delta=\operatorname{sw}(\delta)$ and $L_\ell=\operatorname{sw}(\ell)$ for fixed sweeps, the meridian
$m_{\rm fin}\mu_{\rm mer}-v_{\rm fin}$ gives the Thom columns in coordinates $(s,D_\delta,L_\ell)$:
\[
 \begin{gathered}
 C^{\rm Th}_1=(bn_{\rm fin}/q_{\rm fin},m_{\rm fin}/q_{\rm fin},0),\\ C^{\rm Th}_2=(-br_{\rm fin},0,m_{\rm fin}),\qquad C^{\rm Th}_3=(0,r_{\rm fin},n_{\rm fin}).
 \end{gathered}
\]
The boundary map and its inverse are
\[
 \begin{gathered}
 [y_{\rm bd},\tau_{\rm bd}]\longmapsto[y_{\rm bd}+(\tau_{\rm bd}/m_{\rm fin})v_{\rm fin},e^{2\pi i\tau_{\rm bd}/m_{\rm fin}}],\\
 [x_{\rm bd},z_{\rm bd}]\longmapsto[x_{\rm bd}-(\tau_{\rm bd}/m_{\rm fin})v_{\rm fin},\tau_{\rm bd}].
 \end{gathered}
\]
A change of root argument is exactly the mapping-torus relation on both
sides. The normal circle is $(y_{\rm bd},\tau_{\rm bd})\mapsto(y_{\rm bd}-t_{\rm nor}v_{\rm fin},\tau_{\rm bd}+m_{\rm fin}t_{\rm nor})$.
The boundary-to-core map sends $\beta_{\rm fin}$ and $\ell$
to their coinvariant classes and $\mu_{\rm mer}$ to $c_{\rm core}$.
Thus $m_{\rm fin}\mu_{\rm mer}-v_{\rm fin}$ maps to zero.
The central surface satisfies
\[
 H_1(F_{\rm fin})=\langle\beta_{\rm fin},\ell,c_{\rm core}\mid m_{\rm fin}c_{\rm core}=n_{\rm fin}\ell+r_{\rm fin}q_{\rm fin}\beta_{\rm fin}\rangle.
\]
For $A_0m_{\rm fin}+B_0n_{\rm fin}=1$, a basis and inverse are
\[
 \begin{gathered}
 \beta_{\rm fin},\ z_{\rm cen}=A_0\ell+B_0c_{\rm core};\\
 \ell=m_{\rm fin}z_{\rm cen}-B_0r_{\rm fin}q_{\rm fin}\beta_{\rm fin},\\ c_{\rm core}=n_{\rm fin}z_{\rm cen}+A_0r_{\rm fin}q_{\rm fin}\beta_{\rm fin}.
 \end{gathered}
\]
Thus $H_1(F_{\rm fin})$ is free of rank two. Its finite torus cover and transfer
give rational Betti numbers $(1,2,2,2,1)$; integral duality then makes
$H_2(F_{\rm fin})$ and $H_3(F_{\rm fin})$ free. The flat normal line has torsion Chern class
in the free abelian group $H^2(F_{\rm fin};\mathbb Z)$, so this class is zero.
Transfer of a loop $\lambda$ is the torus
$[t_{\rm loop}\lambda,\exp(2\pi i t_{\rm arg})]$. Cutting into $m_{\rm fin}$ sectors gives fibre area
$-\lambda\wedge v_{\rm fin}$ and sweep $N_{m_{\rm fin}}\lambda$; rejoining gives the inverse
chain map. The torus for $c_{\rm core}$ has area zero and sweep $v_{\rm fin}$.
Thom exactness on the displayed basis gives the complete integral kernel
$\langle C^{\rm Th}_1,A_0C^{\rm Th}_2+B_0C^{\rm Th}_3\rangle$.

There is at most one finite cap. Gluing sends the marked root loop
$c_{\rm bd}$ to $c_{\rm tot}$. Its Seifert presentation is
\[
 \langle\gamma,c_{\rm tot}\mid[\gamma,c_{\rm tot}]=1,
                  \ c_{\rm tot}^{m_{\rm fin}}=\gamma^{n_{\rm fin}},\ c_{\rm tot}=\gamma^e\rangle,
\]
which is trivial exactly when $m_{\rm fin}e-n_{\rm fin}=\varepsilon_{\rm det}=\pm1$.
In the other rows the group is trivial exactly when $e=\pm1$.
Take $A_0=\varepsilon_{\rm det} e$, $B_0=-\varepsilon_{\rm det}$. Then
\[
 \begin{aligned}
 K_0&=A_0C^{\rm Th}_2+B_0C^{\rm Th}_3=(-A_0br_{\rm fin},B_0r_{\rm fin},1),\\
 C^{\rm Th}_2-m_{\rm fin}K_0&=-B_0q_{\rm fin}r_{\rm fin}C^{\rm Th}_1,\\
 C^{\rm Th}_3-n_{\rm fin}K_0&=A_0q_{\rm fin}r_{\rm fin}C^{\rm Th}_1.
 \end{aligned}
\]
These identities give both directions of the generator change.
The ordinary polygon relation $D_\delta=-be\,s$ gives
\[
 L_\ell=br_{\rm fin}(A_0+B_0e)s=0,\qquad
 (b/q_{\rm fin})(n_{\rm fin}-m_{\rm fin}e)s=-3\varepsilon_{\rm det} s=0.
\]
Without a finite cap the remaining relation is $be\,s=0$, with
$b=3$ and $e=\pm1$. A peripheral change $c_{\rm bd}\mapsto c_{\rm bd}\lambda$ changes
the power by $N_{m_{\rm fin}}\lambda$ and changes a fixed sweep by
$\operatorname{sw}(u)\mapsto \operatorname{sw}(u)+\omega(u,\lambda)s$; the opposite shear and inverse
period change reverse the map. This proves the same quotient for all
the affine representatives and all integers $k_{\rm fin},e,r_{\rm fin}$ above.
Mayer--Vietoris with the fibre sections and graph inverse therefore gives
\[
 H_2(X;\mathbb Z)=\mathbb Z/3\langle s\rangle.
\]
The two zero-dimensional toric strata at the cusp contribute Euler
number two. All other torus strata and caps have Euler number zero.
Simple connectivity, duality, and the universal coefficient sequence give
\[
 H_0=H_6=\mathbb Z,\qquad H_1=H_4=H_5=0,\qquad
 H_2=H_3=\mathbb Z/3,\qquad b_3=0.
\]
Consequently $H^2(X;\mathbb Z)=H^2(X;\mathbb F_2)=0$; the complex first
Chern class vanishes and the spin structure is unique. Let
$\alpha\in H^2(X;\mathbb F_3)$ satisfy $\alpha(s)=1$.

\item \emph{The primitive period charts and circle isotropy.}
At an unselected multiplicative place set
\[
 d_0=\gcd(\nu,a\kappa),\quad n_0=\nu/d_0,
 \quad A'=a\kappa/d_0,\quad B'=b\kappa/n_0.
\]
Here $B'$ is integral and $\gcd(d_0,B')=1$. Choose
$n_0\theta_{\rm Bez}+A'\sigma_{\rm Bez}=1$ and $d_0x_{\rm Bez,N}+B'y_{\rm Bez,N}=1$. With
$S_{\rm sh}=(u,\delta)$ and $L_{\rm long}=(v,\gamma)$ put
\[
 C_S=\begin{pmatrix}n_0&\sigma_{\rm Bez}\\-A'&\theta_{\rm Bez}\end{pmatrix},\qquad
 C_L=\begin{pmatrix}x_{\rm Bez,N}&-B'\\y_{\rm Bez,N}&d_0\end{pmatrix},
\]
\[
 C_S^{-1}=\begin{pmatrix}\theta_{\rm Bez}&-\sigma_{\rm Bez}\\A'&n_0\end{pmatrix},\qquad
 C_L^{-1}=\begin{pmatrix}d_0&B'\\-y_{\rm Bez,N}&x_{\rm Bez,N}\end{pmatrix}.
\]
Define $(l,p_0)=S_{\rm sh}C_S$, $(h_{\rm per},w_0)=L_{\rm long}C_L$. The ordered basis
$(l,h_{\rm per},p_0,w_0)$ has determinant one and return $h_{\rm per}\mapsto h_{\rm per}+l$.
Its pairing and scalar vector are
\[
 \begin{gathered}
 K_N=(S_{\rm sh}C_S)^t\omega(L_{\rm long}C_L)=\begin{pmatrix}n_{\rm per}&0\\M_N&r_{\rm per}\end{pmatrix},\\
 n_{\rm per}=a\nu/d_0^2,\quad r_{\rm per}=-bd_0^2/\nu,\\
 \delta=-\sigma_{\rm Bez} l+n_0p_0.
 \end{gathered}
\]
The identity on the period cover, followed by these basis changes and
exponentiation of $l$, gives adjacent Tate charts
\[
 \begin{gathered}
 X_{\rm Tate}=z_{\rm Tate}/q^j,\quad Y_{\rm Tate}=q^{j+1}/z_{\rm Tate},\quad X_{\rm Tate}Y_{\rm Tate}=q;\\
 z_{\rm Tate}=X_{\rm Tate}q^j,\quad q=X_{\rm Tate}Y_{\rm Tate}.
 \end{gathered}
\]
The primitive long period identifies successive charts, with its period
units. The inverse uses the inverse basis matrices and these inverse
chart formulas. The elliptic conductor has primitive periods $p_0,w_0$.
The scalar circle translates it by $n_0t\,p_0$.
In every row there is exactly one site with $3\mid n_0$: the $I_3$
site, where $d_0=1$, $n_0=3$, $n_{\rm per}=3a$, and $r_{\rm per}=-b/3$.
Call its elliptic conductor $E$. Its stabilizer is $C_3$ and its two
normal characters are $-\sigma_{\rm Bez},\sigma_{\rm Bez}$, both units modulo three.
Away from the node a nonzero normal coordinate is rotated by that
character. Every other rank-one site has $n_0=1$, so the translation
by $n_0/3$ on its conductor is nontrivial.

At the selected cusp put
\[
 \rho_{\rm cusp}=(1-ab\kappa^2)/\nu,\qquad
 C=\begin{pmatrix}\nu&b\kappa\\-a\kappa&\rho_{\rm cusp}\end{pmatrix},\qquad
 C^{-1}=\begin{pmatrix}\rho_{\rm cusp}&-b\kappa\\a\kappa&\nu\end{pmatrix}.
\]
This integral matrix has determinant one, and $(T_{\rm mon}-I)L_{\rm long}=S_{\rm sh}C$.
Thus the short basis of the prescribed Jordan fan is $S_{\rm sh}C$, and the
scalar circle has short coordinates
\[
 d_\delta=C^{-1}(0,1)^t=(-b\kappa,\nu)^t.
\]
Write $k_{\rm edge}=b\kappa$. Then $3\mid k_{\rm edge}$ and $3\nmid\nu$. The vertical edge
of the prescribed unit triangulation is the unique edge direction
parallel to $d_\delta$ modulo three. Its curve $C_0$ has charts
\[
 (x,y,w)=(q/(X_1X_2),X_2,X_1),\qquad
 (X_1,X_2,q)=(w,y,xyw),
\]
\[
 (x',y',w')=(q/X_2,X_1X_2,X_1^{-1}),\qquad
 (X_1,X_2,q)=(1/w',y'w',x'y'w').
\]
Their transition and inverse are
\[
 (x',y',w')=(wx,wy,w^{-1}),\qquad
 (x,y,w)=(w'x',w'y',(w')^{-1}).
\]
Hence $C_0\cong\mathbb P^1$ with normal bundle
$\mathcal O(-1)\oplus\mathcal O(-1)$. The circle weights in these charts are
\[
 (k_{\rm edge}-\nu,\nu,-k_{\rm edge}),\qquad(-\nu,\nu-k_{\rm edge},k_{\rm edge}).
\]
It fixes $C_0$ pointwise for $k_{\rm edge}=0$ and rotates it with tangent weight
$-k_{\rm edge}$ for $k_{\rm edge}\ne0$. In either case $C_3$ fixes $C_0$, and its two normal
characters are nonzero. Every triangle contains a vertical edge; its
two vertex orbits are the two endpoints of this curve. A translation
fixing a bounded edge fixes both endpoints and is zero, so the edge
quotient is embedded.

On an original additive cap, a norm strictly decreasing under the
positive contracting arrow excludes equality of a nonzero vector with
a nonzero iterate times a unit scalar. The zero iterate fixes that
vector only for the identity scalar. Thus its scalar action is free.
At a finite cap, an equality between scalar translation and the $j$th
deck iterate implies $jn_{\rm fin}/m_{\rm fin}\in\mathbb Z$ after applying $\chi$.
Since $\gcd(m_{\rm fin},n_{\rm fin})=1$, this forces $m_{\rm fin}\mid j$; primitivity of $\delta$
then makes the scalar trivial. Off the core the base rotation gives
the same divisibility. The ordinary cap is free for the same primitive
torus-translation reason. We have proved
\[
                         X^{C_3}=C_0\sqcup E.
\]
Every nontrivial subgroup of $C_9$ contains $C_3$. Consequently $C_9$
is free on the complement and on the two normal sphere bundles.

\item \emph{One equivariant line on the partial completion.}
Let $Y$ denote the completion at all multiplicative places, with the
ordinary disc and the additive discs still removed. On its real period
cover define
\[
 J_\lambda(x,\zeta_R)=(x+\lambda,e^{2\pi i\omega(\lambda,x)}\zeta_R).
\]
The cocycle quotient is $e^{2\pi i\omega(\lambda,\mu_{\rm lat})}=1$ and the
inverse is $J_{-\lambda}$. The connection
$d-2\pi i\omega(x,dx)$ has first Chern class $2\omega$ in the positive
line convention. Monodromy lifts by $(x,\zeta_R)\mapsto(T_ix,\zeta_R)$, with
inverse $(x,\zeta_R)\mapsto(T_i^{-1}x,\zeta_R)$.
The scalar generator of $C_3$ lifts as
\[
 \mathcal T(x,\zeta_R)=
 (x+\delta/3,e^{2\pi i\omega(\delta/3,x)}\zeta_R).
\]
Its commutator with each period is one because
$\omega(\delta/3,\Lambda)=-(b/3)\mathbb Z$, and its third power is
$J_\delta$. These formulas define a $C_3$-equivariant line $\mathcal D$.

For the nonzero elliptic short period $\omega_s$, set
$z_E=z_{\rm add}/\omega_s$, with inverse $z_{\rm add}=\omega_s z_E$.
Put $\tau=\omega_l/\omega_s$, $\zeta_G=G_{\rm lift}/\omega_s$,
and $z_{\rm ph}=e^{2\pi i z_E}$; its inverse logarithm is taken modulo the
same short period. Let $\rho_L$ be the forward line-period logarithm.
For the holomorphic comparison put $D_0=\operatorname{diag}(a,-b)$
and use physical logarithmic coordinates
$Z_{\rm ph}=(z_E,z_{\rm line})^t=u_{\rm ph}+\Omega_{\rm ph}W_{\rm ph}$,
with $v_{\rm ph}=e^{2\pi i z_{\rm line}}$. Both inverse logarithms
are taken modulo their integral short periods. The entries of $\Omega_{\rm ph}$ are
$\tau,-b\zeta_G;a\zeta_G,\rho_L$, so $D_0\Omega_{\rm ph}$ is symmetric.
The gauge
\[
 G_0=\exp\{-2\pi i(u_{\rm ph}^tD_0W_{\rm ph}+W_{\rm ph}^tD_0\Omega_{\rm ph}W_{\rm ph})\}
\]
and its reciprocal give short factors one and long factors
$h_1^{-a}z_{\rm ph}^{-2a}$ and $R^bv_{\rm ph}^{2b}$. Multiplication of the vector
coordinate by $z_{\rm ph}^a$, with inverse $z_{\rm ph}^{-a}$, gives
\[
 \begin{aligned}
 h_{\rm deck}(z_{\rm ph},v_{\rm ph},\zeta_{\rm aux})&=(h_1z_{\rm ph},h_G^av_{\rm ph},z_{\rm ph}^{-2a}\zeta_{\rm aux}),\\
 F_{\rm deck}(z_{\rm ph},v_{\rm ph},\zeta_{\rm aux})&=(h_G^{-b}z_{\rm ph},Rv_{\rm ph},h_G^{-ab}R^bv_{\rm ph}^{2b}\zeta_{\rm aux}),\\
 J_{\rm inv}(z_{\rm ph},v_{\rm ph},\zeta_{\rm aux})&=(z_{\rm ph}^{-1},v_{\rm ph}^{-1},z_{\rm ph}^{-2a}\zeta_{\rm aux}).
 \end{aligned}
\]
For $h_{\rm deck},F_{\rm deck}$ the inverse auxiliary factors, in target coordinates, are
$h_1^{-2a}z_{\rm ph}^{2a}$ and $h_G^{ab}R^bv_{\rm ph}^{-2b}$; $J_{\rm inv}$ is its own inverse.
Substitution gives $h_{\rm deck}F_{\rm deck}=F_{\rm deck}h_{\rm deck}$, $J_{\rm inv}h_{\rm deck}J_{\rm inv}=h_{\rm deck}^{-1}$, and $J_{\rm inv}F_{\rm deck}J_{\rm inv}=F_{\rm deck}^{-1}$.
The functions $h_1,h_G,R$ include their units. Overlaps are $G_0$,
the integral period change, and the reciprocal gauge of the next
chart. Their triple compositions cancel. The origin section and a
nonzero line frame on the punctured base fix the origin frames.
On a fixed-loop sweep the connection integral and monodromy multiplier
are zero; hence $\mathcal D$ has degree two on $s$ and zero on every
section sweep.

The rank-one coordinates and their inverse are
\[
 \begin{gathered}
 (u_N,W_N,Z_N)=(C_S^{-1}u_{\rm ph},C_L^{-1}W_{\rm ph},
                 C_S^{-1}Z_{\rm ph}),\\
 (u_{\rm ph},W_{\rm ph},Z_{\rm ph})=(C_Su_N,C_LW_N,C_SZ_N).
 \end{gathered}
\]
Set $\Omega_N=C_S^{-1}\Omega_{\rm ph}C_L$, so
$Z_N=u_N+\Omega_NW_N$. Since
$K_N=C_S^tD_0C_L$, the identities
\[
 u_N^tK_NW_N=u_{\rm ph}^tD_0W_{\rm ph},\qquad
 K_N^t\Omega_N=C_L^tD_0\Omega_{\rm ph}C_L
\]
identify the following gauge with $G_0$ in these coordinates.
The return of $\Omega_N$ adds $E_{11}$ and $K_N^t\Omega_N$ is symmetric. The gauge
\[
 G_N=\exp\{-2\pi i(u_N^tK_NW_N+W_N^tK_N^t\Omega_N W_N)\}
\]
and its reciprocal have long factors, for $h_N\in\mathbb Z^2$,
$\exp(-4\pi i h_N^tK_N^tZ_N-2\pi i h_N^tK_N^t\Omega_N h_N)$.
For the first long-shift coordinate $h_{N,1}$, the integral support
$n_{\rm per}j^2$ has difference
$2n_{\rm per}h_{N,1}j+n_{\rm per}h_{N,1}^2$, giving
the normal monomial factors in the Tate charts. The other factors
are the units of $\Omega_N$. Replacing $\Omega_N$ by $\Omega_N+E_{11}$
and replacing $u_N$ by $u_N+E_{11}W_N$ both add $n_{\rm per}W_{N,1}^2$ to the exponent
before multiplication by $-2\pi i$, proving compatibility around the base.
In physical coordinates let $r_{\rm Tate}=\operatorname{ord}h_G=-\kappa$,
$\rho_{\rm Tate}=\operatorname{ord}R=-ab\kappa^2/\nu$, and write the two
valuation periods as $\alpha_0u_0,\beta_{\rm val}u_0$, where
$\alpha_0=d_0$, $u_0=(n_0,-A')$, and $\beta_{\rm val}=B'$.
Then $a\beta_{\rm val}u_{0,1}+b\alpha_0u_{0,2}=0$ and the physical support is
\[
 \psi_1(j)=-\frac{a u_{0,1}}{\alpha_0}j(j-\alpha_0)
          =-n_{\rm per}j^2+n_{\rm per}\alpha_0j.
\]
Its period differences are $-2au_{0,1}j$ and
$2bu_{0,2}j+b\rho_{\rm Tate}-ab r_{\rm Tate}$. The residual of the second equality factors as
$(a\beta_{\rm val}u_{0,1}+b\alpha_0u_{0,2})(\alpha_0-\beta_{\rm val}-2j)/\alpha_0$.
The linear term is the gauge $z_{\rm ph}^a$, and the quadratic sign is the
change from frame to vector coordinates. At the $I_3$ core,
$\delta/3=-\sigma_{\rm Bez} l/3+p_0$; its integral short shift adds no character.
The local character is $-\sigma_{\rm Bez} n_{\rm per}(2j+1)$, divisible by three.
The complex-oriented torus $(p_0,w_0)$ specializes with coefficient one
to $E$, giving
\[
                  \deg\mathcal D|_E=-2b/3.
\]

At the cusp use
\[
 \begin{gathered}
 (u_C,W_C,Z_C)=(C^{-1}u_{\rm ph},W_{\rm ph},C^{-1}Z_{\rm ph}),\\
 (u_{\rm ph},W_{\rm ph},Z_{\rm ph})=(Cu_C,W_C,CZ_C),\\
 \Omega_C=C^{-1}\Omega_{\rm ph}.
 \end{gathered}
\]
Thus $Z_C=u_C+\Omega_CW_C$. Put
\[
 K_C=C^tD_0=
 \begin{pmatrix}a\nu&ab\kappa\\ab\kappa&b(ab\kappa^2-1)/\nu\end{pmatrix},
 \qquad K_Cd_\delta=(0,-b)^t.
\]
The equality $C^tD_0=D_0C$ gives
\[
 K_C\Omega_C=D_0\Omega_{\rm ph},\qquad
 u_C^tK_CW_C=u_{\rm ph}^tD_0W_{\rm ph}.
\]
Hence $K_C$ and $K_C\Omega_C$ are symmetric. On the chosen logarithm branch write
\[
 \Omega_C={\log q\over2\pi i}I+B_C(q),\qquad
 (X_1,X_2)=\bigl(e^{2\pi i(Z_C)_1},e^{2\pi i(Z_C)_2}\bigr).
\]
Here $B_C(q)$ retains all the single-valued period units; multiplication
by $C$ recovers the complete physical period matrix. The gauge
\[
 G_C=\exp\{-2\pi i(u_C^tK_CW_C+W_C^tK_C\Omega_C W_C)\}
\]
and its reciprocal give long factors
$q^{-h_C^tK_Ch_C}X_C^{-2K_Ch_C}\exp(-2\pi i h_C^tK_CB_C(q)h_C)$,
where $X_C^{\lambda}=X_1^{\lambda_1}X_2^{\lambda_2}$ and $h_C\in\mathbb Z^2$.
The integral support $\psi(v_{\rm lat})=v_{\rm lat}^tK_Cv_{\rm lat}$ on the prescribed unit triangles
has difference $2h_C^tK_Cv_{\rm lat}+h_C^tK_Ch_C$. On a triangle of slope $m_{\rm slope}$, the regular
frame and holomorphic vector coordinate are
$e_{m_{\rm slope}}=X_C^{-m_{\rm slope}}e_{\rm frame}$ and
$\zeta_{m_{\rm slope}}=X_C^{m_{\rm slope}}\zeta_{\rm aux}$;
the opposite exponents give their inverses. Around the base,
$G_C(u_C,W_C,\Omega_C+I)=G_C(u_C+W_C,W_C,\Omega_C)$.
For the physical support use the selected cusp value
$\rho_{\rm cusp}=(1-ab\kappa^2)/\nu$ and $r_{\rm Tate}=-\kappa$.
Its quadratic matrix is $-D_0C^{-1}$; conjugation by $C$ gives
$C^t(-D_0C^{-1})C=-K_C$. Hence the exact pullback is
\[
 \begin{aligned}
 \psi_{\rm ph}(x,y)&=-a\rho_{\rm cusp} x^2-2ab r_{\rm Tate}xy+b\nu y^2+ax,\\
 \psi_{\rm ph}(Cv_{\rm lat})&=-v_{\rm lat}^tK_Cv_{\rm lat}+a(\nu v_{{\rm lat},1}-b r_{\rm Tate}v_{{\rm lat},2}).
 \end{aligned}
\]
Its two period differences are $-2ax$ and $2by+b\rho_{\rm cusp}-ab r_{\rm Tate}$.
The affine term is the same $z_{\rm ph}^a$ gauge. The slopes adjacent to
$[0,e_2]$ are $(K_{C,11},K_{C,22})$ and $(2K_{C,12}-K_{C,11},K_{C,22})$.
The first character is $b\{a\kappa(b\kappa-\nu)-1\}$ and their
character difference is $2(K_{C,11}-K_{C,12})(d_\delta)_1$, both divisible
by three. The transition in $X_1$ and $X_1^{-1}$ gives
\[
 \deg\mathcal D|_{C_0}=2(K_{C,11}-K_{C,12})=2a(\nu-b\kappa).
\]
Nonzero base units have local logarithms, are scalar invariant, and
have degree zero on these curves; their factors remain in the gauges.
This constructs one line on $Y$ with zero $C_3$ characters on both
components and the stated degrees.

\item \emph{The global Borel class and its finite characters.}
The cusp kills the two short columns $S_{\rm sh}C$, equivalently $u,\delta$.
Every other multiplicative cap kills its primitive $l$.
Their columns in the first three period coordinates form a matrix
$K_1$ with an integer right inverse $K_1R_1=I_3$ in each row.
Thus $e_1,e_2,\delta$ vanish in the fundamental group, by the cap
discs and the corresponding integer combinations. The remaining
matrix relations make $\gamma$ commute with each unfilled base meridian.
There are zero or one unfilled additive meridians, and
\[
                \pi_1(Y)=\mathbb Z_\gamma\times F_{r_Y},
                  \qquad r_Y=0\text{ or }1.
\]
The inverse group map kills the first three periods, fixes $\gamma$
and the additive meridian, and replaces the ordinary meridian by the
inverse product. This proves torsion freeness.
The graph and cap maps above send every section sweep to zero in
$H_2(X;\mathbb F_3)$: for a finite cap $L_\ell=0$ and $D_\delta=-be\,s=0$,
and for an original cap both sweeps bound. Since $H_1(Y)$ is free,
the cycle values determine the ordinary degree-two class, giving
\[
                  \tfrac12\rho_3c_1(\mathcal D)=\alpha|_Y.
\]
Here $1/2=2$ in $\mathbb F_3$.
The connected circle acts trivially on ordinary cohomology. The only
possible outgoing differentials on $\alpha$ in the Borel sequence have
targets $H^2(BS^1;H^1(X;\mathbb F_3))$ and $H^3(BS^1;\mathbb F_3)$,
both zero. Hence a class $U_{\rm raw}\in H^2_{S^1}(X;\mathbb F_3)$ restricts to
$\alpha$.

Use a toric endpoint of $C_0$ as basepoint. The circle fixes it, so
its isotopy to any group element is based and
$\pi_1(Y_{C_3})=\pi_1(Y)\times C_3$. Every finite isotropy section
has generator $(1,\zeta_3)$, since its first component lies in a
torsion-free group. Let $T_{\rm circ}=c_1(\text{positive circle character})$ and $t_3=\rho_3(T_{\rm circ})$.
In degree two the kernel of restriction to $Y$
is spanned by the base class $t_3$ and
$\rho_3(\lambda_j)\eta$, where the integral $\lambda_j$ form a basis
of $H^1(Y)$ and vanish on the finite factor, and
$\eta\in H^1(BC_3;\mathbb F_3)$. These realize the two positive-base
terms $(2,0)$ and $(1,1)$ in total degree two. Equivalently, the
universal-cover five-term sequence gives those same group-cohomology
terms for $\pi_1(Y)\times C_3$; its ordinary group-cohomology summand
injects into $H^2(Y)$.
Thus
$U_{\rm raw}|_{Y_{C_3}}-\tfrac12\rho_3c_1^{C_3}(\mathcal D)$
has one common base coefficient $c_B$ on both components. The mixed
products vanish on each finite isotropy section. Set
$U=U_{\rm raw}-c_B\rho_3(T_{\rm circ})$, with inverse
$U_{\rm raw}=U+c_B\rho_3(T_{\rm circ})$. Both finite point
characters of $U$ vanish.

Put ${a_{\rm edge}}=a(\nu-k_{\rm edge})$ and let $z_C=c_1^{S^1}(\mathcal O(1))$ on $C_0$,
with value zero at the first endpoint and $k_{\rm edge}T_{\rm circ}$ at the second.
The chart weights give $z_C(z_C-k_{\rm edge}T_{\rm circ})=0$ and $\int_{C_0}z_C=1$. The restriction is
\begin{equation}
             U|_{(C_0)_{S^1}}=\rho_3({a_{\rm edge}}z_C),\qquad z_C(z_C-k_{\rm edge}T_{\rm circ})=0.
             \hypertarget{tag-025F}{}\label{maps:eq:D.T34}
\end{equation}
The generator satisfying $u_{\rm coh}([C_0])=1$ is
$u_{\rm coh}={a_{\rm edge}}^{-1}\alpha$ in $\mathbb F_3$; its inverse change is $\alpha={a_{\rm edge}}u_{\rm coh}$.

\item \emph{The coefficient-one detector and spin bordism.}
Cartan's divided-power computation gives
\begin{equation}
 H_5(K(\mathbb Z/3,2);\mathbb Z)=0,\qquad
 H_6(K(\mathbb Z/3,2);\mathbb Z)=
       \mathbb Z/9\langle\gamma_3(e_{\rm EM})\rangle,
\hypertarget{tag-Q224}{}\label{maps:eq:three-primary-homology}
\end{equation}
where positive reduction $\mathbb {CP}^{\infty}\to K(\mathbb Z/3,2)$
sends $[\mathbb {CP}^3]$ to $\gamma_3(e_{\rm EM})$
\cite[Expos\'e~11, Theorems~1 and~4, pp.~11-09--11-14]{manCartan1955EML}.
Universal coefficients identify universal degree-six cohomology with
$\operatorname{Hom}(\mathbb Z/9,\mathbb Z/9)$. The detector is therefore
fixed by positive $\mathbb {CP}^3$ evaluation one:
\begin{equation}
 Q_3:H^2(-;\mathbb F_3)\longrightarrow H^6(-;\mathbb Z/9),\qquad
 Q_3(\rho_3x)=\rho_9(x^3).
\hypertarget{tag-Q225}{}\label{maps:eq:third-power-detector}
\end{equation}
Its coefficient map is $\gamma_3(e_{\rm EM})\mapsto1$, with inverse
$j\mapsto j\gamma_3(e_{\rm EM})$.
Thomas's odd-prime operation, for the coefficient groups
$\mathbb Z/(pm_{\rm op})$ and $\mathbb Z/(p^2m_{\rm op})$, has the cochain representative
\begin{equation}
 \begin{aligned}
 P_p(v_{\rm cochain})&=(v_{\rm cochain}^p,0)+(\Sigma^*v_{\rm cochain}^{p-1}\delta_{\rm cochain} v_{\rm cochain},1),\\
 P_3(v_{\rm cochain})&=(v_{\rm cochain}^3,0)+(\Sigma^*v_{\rm cochain}^2\delta_{\rm cochain} v_{\rm cochain},1)\qquad(p=3,\ m_{\rm op}=1).
 \end{aligned}
\hypertarget{tag-Q226}{}\label{maps:eq:third-power-cochain}
\end{equation}
Here $\delta_{\rm cochain}$ is the integral cochain differential,
$p$ is an odd prime and $m_{\rm op}$ is a positive integer
\cite[(2.1),(2.6),(2.7), p.~267; (3.3)--(3.6), p.~268]{ffThomas56}.
For an integral cocycle the second term vanishes. The degree-zero
part of the diagonal carrier may be the iterated Alexander--Whitney
diagonal with augmentation one, and the carried homotopy identifies
any other such choice with it. Its product is consequently the
ordinary $p$-th power with coefficient one. For $p=3,m_{\rm op}=1$ this
identifies Thomas's operation with $Q_3$, including its coefficient map.
The odd-prime relative construction gives vanishing on relative
boundaries and suspensions; the external-product identity reduces the
suspension to the zero product on the dimension-one interval factor.
Its order-three specialization is
\begin{equation}
                      Q_3(\sigma v_{\rm susp})=0,\qquad
                      v_{\rm susp}\in H^1(K_{\rm susp};\mathbb F_3)
\hypertarget{tag-Q227}{}\label{maps:eq:third-power-suspension}
\end{equation}
\cite[Theorem~I(ii), pp.~897--898; (3.1), p.~900; \S5, pp.~904--905]{ffThomas1959Suspension}.
Here suspension is the left cross product by the positive generator
of $H^1(I,\partial I;\mathbb Z)$; slant product by the interval is
its inverse. The odd-prime proof uses the relative external-product
identity and the vanishing on the dimension-one factor.
For $e_{\rm EM}\mapsto v_{\rm unit}e_{\rm EM}$ with $v_{\rm unit}\in(\mathbb Z/3)^\times$, the coefficient
map is multiplication by $v_{\rm unit}^3$ modulo nine. It is independent of the
integer lift, and its inverse is multiplication by $w_{\rm unit}^3$ when
$v_{\rm unit}w_{\rm unit}=1\pmod3$.

We also identify the refined invariant in this marking.
Zhubr's free-lift epimorphism and injective pair $(P,h_{\rm bord})$ imply that
$\Omega^{\mathrm{Spin}}_6(K(\mathbb Z/3,2))$ is generated by the images
of $\epsilon_{\rm bord},\zeta_{\rm bord}$, with $(P,h_{\rm bord})(\epsilon_{\rm bord})=(6,0)$ and
$\zeta_{\rm bord}=[\mathbb {CP}^3,x]$
\cite[\S\S3.11--3.14, pp.~806--808; Theorem~5, p.~796]{manZhubr75}.
Reduction kills $\epsilon_{\rm bord}$ because both invariants vanish.
The class $\zeta_{\rm bord}$ has $h_{\rm bord}=\gamma_3(e_{\rm EM})$ and $P=1\pmod3$:
$c(T\mathbb {CP}^3)=(1+x)^4$ gives $p_1=4x^2$.
Thus $9\zeta_{\rm bord}$ has both invariants zero and vanishes by injectivity,
while its $h_{\rm bord}$-image has exact order nine. We obtain the exact maps
\begin{equation}
 \begin{gathered}
 \Omega^{\mathrm{Spin}}_6(K(\mathbb Z/3,2))
   =\mathbb Z/9\langle\zeta_{\rm bord}\rangle,\\
 h_{\rm bord}(j\zeta_{\rm bord})=j\gamma_3(e_{\rm EM}),\qquad P(j\zeta_{\rm bord})=j\pmod3.
 \end{gathered}
\hypertarget{tag-Q228}{}\label{maps:eq:three-primary-bordism}
\end{equation}
In particular the kernel of the fundamental-class map is zero.

\item \emph{The $C_9$ localization and the two component contributions.}
Let $U_9$ be the restriction of $U$ to $X_{C_9}$ and put
$R_9=\mathbb Z/9$. The periodic cyclic resolution gives
$H^*(BC_9;R_9)=R_9[T_9]\otimes\Lambda(\eta_9)$.
Its coefficient constants inject after $T_9$ is inverted.
Choose disjoint invariant tubular discs around $C_0$ and $E$,
and denote their union by $N_{\rm tub}$. The complement and their sphere bundles are free $C_9$ spaces, whose
Borel spaces have the cohomology of finite-dimensional quotients.
After inverting $T_9$ these groups vanish. Relative excision and
the Thom sequences consequently give
\begin{equation}
 H^*_{C_9}(X,N_{\rm tub};R_9)[T_9^{-1}]=0,\qquad
 H^*_{C_9}(\partial N_{\rm tub};R_9)[T_9^{-1}]=0.
 \hypertarget{tag-025I}{}\label{maps:eq:D.T37}
\end{equation}
The normal Euler classes become invertible, and restriction has
the explicit inverse
\begin{equation}
 H^*_{C_9}(X;R_9)[T_9^{-1}]
    \xrightarrow{\ \cong\ }
 H^*_{C_9}(C_0\sqcup E;R_9)[T_9^{-1}],\qquad
 (v_C,v_E)\longmapsto
       i_{C*}(v_C/e_C)+i_{E*}(v_E/e_E).
 \hypertarget{tag-025J}{}\label{maps:eq:D.T38}
\end{equation}
These maps use the invariant tubular components with their induced
actions. Equivariant integration of $Q_3(U_9)$ lies in $H^0(BC_9;R_9)$
and equals the ordinary evaluation $\langle Q_3(\alpha),[X]\rangle$.

The normal roots on $C_0$ are $(k_{\rm edge}-\nu)T_{\rm circ}-z_C$ and $\nu T_{\rm circ}-z_C$.
Using \eqref{maps:eq:D.T34}, their product is
\[
 e_C=\nu(k_{\rm edge}-\nu)T_{\rm circ}^2,\qquad
 Q_3(U)|_{(C_0)_{S^1}}=\rho_9({a_{\rm edge}}^3z_C^3)
                         =\rho_9({a_{\rm edge}}^3k_{\rm edge}^2T_{\rm circ}^2z_C).
\]
The contribution after restriction to $C_9$ and integration is
\[
                     \frac{{a_{\rm edge}}^3k_{\rm edge}^2}{\nu(k_{\rm edge}-\nu)}=0\pmod9.
\]
The denominator factors are units modulo nine and $3\mid k_{\rm edge}$.

Write $t_{\rm circ}=\exp(2\pi i t)$ for the multiplicative circle
parameter; its inverse is the argument divided by $2\pi$, modulo
$\mathbb Z$. The action on the elliptic component, in its primitive periods, is
$t_{\rm circ}\cdot(z_{\rm circ},w_{\rm circ})=(t_{\rm circ}^3z_{\rm circ},w_{\rm circ})$. The map
$[e_{\rm ES},z_{\rm circ}]\mapsto[e_{\rm ES} r_{\rm root}]\in ES^1/C_3$, where $r_{\rm root}^3=z_{\rm circ}$, has inverse
$[e_{\rm ES}]\mapsto[e_{\rm ES},1]$. Changing the root multiplies by $C_3$ and leaves
the map unchanged; the associated relation
$[e_{\rm ES}t_{\rm circ},z_{\rm circ}]=[e_{\rm ES},t_{\rm circ}^3z_{\rm circ}]$ verifies it on the Borel quotient. With the second
circle this gives
\begin{equation}
 E_{S^1}\simeq BC_3\times S^1,\qquad
 U|_{E_{S^1}}=d_E\,\eta\,\rho_3(\xi),\qquad d_E=-b/3.
 \hypertarget{tag-025G}{}\label{maps:eq:D.T35}
\end{equation}
Here the positive first-period loop maps to the chosen positive $C_3$ generator,
and $\xi$ is the integral generator on the second circle.
The degree-two base coefficient vanishes by the common-character
calculation. The mixed coefficient is $d_E$, by the ordinary degree
of $\tfrac12c_1(\mathcal D)$.
Flip the factors and collapse the wedge to obtain
$c_{\rm susp}:BC_3\times S^1\to S^1\wedge BC_3$.
The left suspension pulls back as
$c_{\rm susp}^*\sigma(\eta)=\rho_3(\xi)\eta=-\eta\rho_3(\xi)$.
Thus \eqref{maps:eq:D.T35} is $c_{\rm susp}^*\sigma(-d_E\eta)$ and
\begin{equation}
 Q_3(U)|_{E_{S^1}}=c_{\rm susp}^*Q_3(\sigma(-d_E\eta))=0
                  \quad\text{in }H^6(E_{S^1};\mathbb Z/9).
 \hypertarget{tag-025H}{}\label{maps:eq:D.T36}
\end{equation}
Its $C_9$ restriction also vanishes. Formula \eqref{maps:eq:D.T38} and
injectivity of constants give
\begin{equation}
                  \langle Q_3(\alpha),[X]\rangle=0
                                   \quad\text{in }\mathbb Z/9.
 \hypertarget{tag-025K}{}\label{maps:eq:D.T39}
\end{equation}
The coefficient-one identification proves $\mu(X)=0$. The spin-bordism
coefficient map then gives $\mathfrak p(X)=0$ and the zero marked
spin-bordism class.

\item \emph{The oriented surgery model.}
Let $f_{\rm surg}:S^2\to S^2$ have degree three and let
$\iota:S^2\hookrightarrow S^4$ be an equatorial embedding. The map
$x\mapsto(f_{\rm surg}(x),\iota(x))$ is an embedding representing three times
the positive first-factor generator. Its oriented rank-four normal
bundle has $w_2=0$, since $S^2\times S^4$ is spin and
$w_2(TS^2)=0$. Oriented rank-four bundles on $S^2$ are classified by
this class, so the normal bundle is trivial. Its framing is unique
up to homotopy because $\pi_2SO(4)=0$.
The framed index-three surgery trace is a spin cobordism. The attaching
map on $H_2$ is multiplication by three, giving
$H_2(W_{\rm surg})=\mathbb Z/3$, $H_3(W_{\rm surg})=0$, and $H_4(W_{\rm surg})=\mathbb Z$.
From the other end its dual handle has index four; intersection of
the first-factor sphere with the second-factor sphere gives
$H_4(W_{\rm surg})\to H_4(W_{\rm surg},M_{3,0})$ equal to multiplication by three.
The two exact sequences therefore give
$H_2(M_{3,0})=H_3(M_{3,0})=\mathbb Z/3$ and $b_3=0$.
General position in real codimension four and van Kampen preserve
simple connectivity. The spin structure extends across the framed handle.
The mod-three class of the first factor extends over the trace since
the attaching sphere evaluates to zero; the relative extension is
unique because $\pi_3K(\mathbb Z/3,2)=0$.
The initial marked pair bounds $S^2\times D^5$ over $K(\mathbb Z/3,2)$.
Consequently $M_{3,0}$ has zero marked spin-bordism class.

Zhubr's marked recognition theorem for zero free third homology,
with the construction and uniqueness in \S4.1, applies to $X$
and $M_{3,0}$
\cite[Theorem~3, p.~795; Theorem~5, p.~796; \S4.1, pp.~809--810]{manZhubr75}.
Both classifying maps induce the given isomorphism on $H_2$, both
have $b_3=0$, and their spin-bordism classes coincide. The resulting
oriented diffeomorphism carries the positive surgery generator to $s$.
This proves \eqref{maps:eq:D.T32} and \eqref{maps:eq:D.T33}.
\end{enumerate}
\end{proof}

\subsubsection{Orders six, fifteen and thirty}

\begin{definition}[Mixed three-primary data]\hypertarget{tag-Q229}{}\label{maps:def:mixed-three-data}
{\normalfont\scriptsize tag Q229}\par
Use the ordered data of Definition~\ref{mapsIII:def:elliptic-marked-data}
and the natural initial pairs of \ref{maps:src:D.2.22}. Choose minimal
additive fillings of \ref{maps:src:D.2.26} with residual
$c\in\{6,15,30\}$ and trivial orbifold obstruction. Use the prescribed
primitive Mumford fan at every multiplicative site. The logarithmic
sites are precisely the chosen additive sites. At a finite site of
linear monodromy order $d_i$, allow $m=d_i k$ for every $k>0$ and
every free affine lift; allow every ordinary integer clutch $e\gamma$.
All orientations are the complex orientations. Selected positions are
$i-1$ in the input definition's peripheral order.

Use $M_0,M_1$ from Definition~\ref{maps:def:minimal-two-data},
$M_{3,0}$ from Definition~\ref{maps:def:minimal-three-data}, and
$Z_{5,\epsilon}$ from Definition~\ref{maps:def:minimal-mixed-data},
with their stated primitive markings. Put
\begin{equation}
 Y_{6,t}=M_t\#M_{3,0},\qquad
 Y_{15,\epsilon}=M_{3,0}\#Z_{5,\epsilon},\qquad
 Y_{30,t}=M_t\#M_{3,0}\#Z_{5,0}.
\hypertarget{tag-Q22D}{}\label{maps:eq:MT_local_4}
\end{equation}

These choices are the \emph{mixed three-primary data}.
\end{definition}

\begin{theorem}[Mixed three-primary smooth types]\hypertarget{tag-02LJ}{}\label{maps:src:D.2.42}
{\normalfont\scriptsize tag 02LJ}\par
The data of Definition~\ref{maps:def:mixed-three-data} give fourteen
patterns. Every simply connected threefold has the oriented type below
for every permitted finite multiplicity, affine lift and ordinary clutch.
The complete oriented table is
\begin{equation}
\begin{array}{r|l|r|r|r|l|r|l}
 &\mathrm{OS/profile}&a&b&\mathrm{zero}&\mathrm{additive}&c&\mathrm{type}\\\hline
1&\mathrm{OS56}/\mathrm{II}&1&6&2&II\ \mathrm{log}&6&Y_{6,0}\\
2&\mathrm{OS56}/\mathrm{II}&1&6&2&II\ \mathrm{original}&6&Y_{6,0}\\
3&\mathrm{OS56}&1&6&0&\mathrm{none}&6&Y_{6,1}\\
4&\mathrm{OS56}&5&6&3&\mathrm{none}&6&Y_{6,1}\\
5&\mathrm{OS56}&5&6&4&\mathrm{none}&6&Y_{6,0}\\
6&\mathrm{OS56}/\mathrm{II}&1&15&0&II\ \mathrm{log}&15&Y_{15,1}\\
7&\mathrm{OS56}/\mathrm{II}&1&15&0&II\ \mathrm{original}&15&Y_{15,1}\\
8&\mathrm{OS56}/\mathrm{III}&1&30&3&III\ \mathrm{log}&15&Y_{15,0}\\
9&\mathrm{OS56}/\mathrm{III}&2&15&3&III\ \mathrm{original}&15&Y_{15,1}\\
10&\mathrm{OS56}&1&15&2&\mathrm{none}&15&Y_{15,1}\\
11&\mathrm{OS56}&2&15&3&\mathrm{none}&15&Y_{15,1}\\
12&\mathrm{OS56}&2&15&4&\mathrm{none}&15&Y_{15,1}\\
13&\mathrm{OS56}&1&30&3&\mathrm{none}&30&Y_{30,1}\\
14&\mathrm{OS56}&1&30&4&\mathrm{none}&30&Y_{30,0}
\end{array}
\hypertarget{tag-020J}{}\label{maps:eq:D.MT1}
\end{equation}
Selected positions use the index $i-1$ of
Definition~\ref{mapsIII:def:elliptic-marked-data}.
There are two oriented types in each sector. At \(c=6\) their
multiplicities are three for \(t=0\) and two for \(t=1\);
at \(c=15\) they are one for \(\epsilon=0\) and six for
\(\epsilon=1\); at \(c=30\) each occurs once.

Every threefold is spin and has
\begin{equation}
\begin{aligned}
 H_2(X)=H_3(X)&=\mathbb Z/c,\quad
 H_1(X)=H_4(X)=H_5(X)=0,\\
 H_0(X)=H_6(X)&=\mathbb Z,\quad b_3(X)=0.
\end{aligned}
\hypertarget{tag-Q22A}{}\label{maps:eq:MT_local_1}
\end{equation}
In the physical fibre marking \(s\) defined in the proof, let
\(\gamma_3(s)\) be the positive \(\mathbb{CP}^3\) reduction generator
of \(H_6(K(\mathbb Z/c,2);\mathbb Z)=\mathbb Z/(3c)\).
At the selected cusp write its marked data as \((\nu,\kappa)\).
The complete primary coefficients are
\begin{equation}
 t=1+\kappa\pmod2\quad(2\mid c),\qquad
 P_p=\frac{b/p-a\nu}{2}\pmod p\quad(p=3,5;\ p\mid c).
\hypertarget{tag-Q22B}{}\label{maps:eq:MT_local_2}
\end{equation}
Here \(P_3=0\), and \(P_5\) is either \(0\) or \(3\);
write \(\epsilon=[P_5\ne0]\). The full invariants are
\begin{equation}
\begin{array}{c|c|c}
c&P(X)&\mu(X)\\\hline
6&3t\,s&9t\,\gamma_3(s)\pmod{18}\\
15&3\epsilon\,s&0\pmod{45}\\
30&15t\,s&45t\,\gamma_3(s)\pmod{90}.
\end{array}
\hypertarget{tag-Q22C}{}\label{maps:eq:MT_local_3}
\end{equation}
In particular the full three-primary coordinate is zero in
\(\mathbb Z/9\), for every allowed finite parameter.
\end{theorem}


\begin{proof}
\begin{enumerate}
\item \emph{Exhaustion of minimal fillings.}

Use the ordered matrices, section cocycles, component coordinates
and heights of Definition~\ref{mapsIII:def:elliptic-marked-data}. For a multiplicative site,
choose the primitive vector \(u\) with first nonzero entry positive and
the B\'ezout complement \(v\) so that
\begin{equation}
 A-I=\nu u u^tJ,\quad g=\kappa u,\quad
 J=\begin{pmatrix}0&1\\-1&0\end{pmatrix},\quad\det(u,v)=1.
\hypertarget{tag-Q22E}{}\label{maps:eq:MT_local_5}
\end{equation}
The integer \(\nu\) is the content of \(A-I\). If
\(x_u u_1+y_u u_2=1\), take \(v=(-y_u,x_u)^t\).
For height \(1/d\), a natural selected zero requires \(ab=d/\nu\)
with \(a,b>0\), \(\gcd(a,b)=1\), \(\gcd(b\kappa,\nu)=1\), and the literal
component equality
\begin{equation}
 e_{a\kappa}-e_0-e_{(a-b)\kappa}+e_{-b\kappa}
                =-2e_0+e_1+e_{-1}\quad\text{in }\mathbb Z[\mathbb Z/\nu].
\hypertarget{tag-Q22F}{}\label{maps:eq:MT_local_6}
\end{equation}
Repeated residues are added. The formula also applies to \(\nu=1,2\).
At each other reducible site one of \(aG,bG\) is narrow. In the marked
component group this means every listed coordinate vanishes modulo its
listed modulus. These are the necessary and sufficient natural-row
conditions of \ref{maps:src:D.2.22}. Its map
multiplies the original line vector by the chord and vertical-divisor
quotient divided by the rational base function; its inverse multiplies
by the reciprocal. The small contracting constant and the nonzero
period-unit functions stay in that construction.

The positive-valuation Tate labelling sends the input component label
\(j\) to \(-j\). This bijection is its own inverse. Applying it to every
term of \eqref{maps:eq:MT_local_6} gives the physical component equation; its right side is
fixed by the reflection. The physical component coordinates of $G$ are obtained by this reflection.

An exact enumeration proceeds through every divisor pair and evaluates
\eqref{maps:eq:MT_local_6} as an integer vector, then checks every unselected component. It
recovers the following 21 natural-row counts, summing to 172:
\[
 \begin{array}{c|rrrrrrr}
 &\mathrm{OS43}&\mathrm{OS43}/\mathrm{II}&\mathrm{OS45}&\mathrm{OS45}/\mathrm{II}&\mathrm{OS46}&\mathrm{OS46}/\mathrm{II}&\mathrm{OS47}\\
 \#&6&2&8&4&6&2&14\\ \hline
 &\mathrm{OS47}/\mathrm{II}&\mathrm{OS47}/\mathrm{III}&\mathrm{OS49}&\mathrm{OS49}/\mathrm{II}&\mathrm{OS49}/\mathrm{III}&\mathrm{OS50}&\mathrm{OS50}/\mathrm{II}\\
 \#&6&8&10&2&4&10&2\\ \hline
 &\mathrm{OS50}/\mathrm{IV}&\mathrm{OS55}&\mathrm{OS55}/\mathrm{II}&\mathrm{OS56}&\mathrm{OS56}/\mathrm{II}&\mathrm{OS56}/\mathrm{III}&\mathrm{OS56}/\mathrm{IV}\\
 \#&4&14&6&26&10&14&14
 \end{array}
\]
At a finite additive site of type \(II,III,IV,III^*,IV^*\), respectively,
the finite order and vertical coinvariant index are
\begin{equation}
 (d_i,q_i)=(6,1),(4,2),(3,3),(4,2),(3,3).
\hypertarget{tag-Q22G}{}\label{maps:eq:MT_local_7}
\end{equation}
The exact finite lattice condition is
\(b(A_i-I)^{-1}g_i\in\mathbb Z^2\). A minimal original additive cap
occurs at \(II,III,IV\) when \(aG\) is narrow and \(bG\) has one component
orbit; the latter is \(\gcd(b,q_i)=1\) in these cyclic cases.
The cap and rank maps of \ref{maps:src:D.2.26} prove these criteria. Finite orders
at different sites are pairwise coprime for a simply connected threefold.
Writing $\mathcal I_{\rm fin}$ for the set of finite additive sites gives
\begin{equation}
                    c=b/\prod_{i\in\mathcal I_{\rm fin}}q_i.
\hypertarget{tag-Q22H}{}\label{maps:eq:MT_local_8}
\end{equation}
Expanding these choices over the sixteen nonstar profiles and
restricting to $c\in\{6,15,30\}$ gives exactly the fourteen patterns
in \eqref{maps:eq:D.MT1}.

Only three defining profiles occur. Their ordered fibre lists and
primitive multiplicative markings are as follows:
\begin{equation}
 \begin{array}{c|l|l}
 \text{profile}&\text{fibres}&(\nu,\kappa)\text{ at multiplicative sites}\\\hline
 \mathrm{OS56}&I_5,I_3,I_2,I_1,I_1&(5,-4),(3,-2),(2,-1),(1,0),(1,-1)\\
 \mathrm{OS56}/\mathrm{II}&I_2,II,I_5,I_3&(2,-1),\ --,\ (5,1),(3,1)\\
 \mathrm{OS56}/\mathrm{III}&I_5,I_3,III,I_1&(5,-4),(3,-1),\ --,\ (1,-1).
 \end{array}
\hypertarget{tag-Q22I}{}\label{maps:eq:MT_local_9}
\end{equation}
Each of these profiles has generator height $1/30$ and at most
one additive site. Its matrices and cocycles are the ordered entries
of Definition~\ref{mapsIII:def:elliptic-marked-data}.

\item \emph{Period maps and integral exterior coordinates.}

Keep
\begin{equation}
 \Lambda=\mathbb Z\langle e_1,e_2,\delta,\gamma\rangle,\qquad
 \omega=a e_1^*\wedge e_2^*-b\delta^*\wedge\gamma^*,\quad
 \chi=\gamma^*.
\hypertarget{tag-Q22J}{}\label{maps:eq:MT_local_10}
\end{equation}
Positive \(\gamma\) is the forward contracting arrow covering \(-bG\).
The scalar circle is \(x\mapsto x+\theta\delta\), \(\theta\in\mathbb R/\mathbb Z\).
At a multiplicative site put \(z_i=(\kappa_i/\nu_i)v_i\); at an additive
site put \(z_i=(A_i-I)^{-1}g_i\). Then
\begin{equation}
 \beta_i=-a g_i^tJA_i,\quad
 c_i=\beta_i b z_i+w_i,\quad
 T_i=\begin{pmatrix}A_i&0&bg_i\\\beta_i&1&c_i\\0&0&1\end{pmatrix},
 \qquad
 w_i=\begin{cases}1/\nu_i&i\text{ selected},\\0&\text{otherwise}.\end{cases}
\hypertarget{tag-Q22K}{}\label{maps:eq:MT_local_11}
\end{equation}
The inverse has blocks
\begin{equation}
 T_i^{-1}=\begin{pmatrix}
 A_i^{-1}&0&-bA_i^{-1}g_i\\
 -\beta_iA_i^{-1}&1&b\beta_iA_i^{-1}g_i-c_i\\0&0&1
 \end{pmatrix}.
\hypertarget{tag-Q22L}{}\label{maps:eq:MT_local_12}
\end{equation}
Substitution gives integral matrices and both inverse products,
$T_i^*\omega=\omega$, $T_i\delta=\delta$, $\chi T_i=\chi$,
and the ordered product $\prod T_i=I$.

These matrices act on the same germs. On a cut base, the period
map is
\begin{equation}
 x\bmod\Lambda\longmapsto\Pi(q)x\bmod\Pi(q)\Lambda.
\hypertarget{tag-Q22M}{}\label{maps:eq:MT_local_13}
\end{equation}
Its inverse is the elliptic and line logarithm modulo that same lattice.
Here \(\Pi(q)\) contains the original nonzero holomorphic units. An
integral basis \(H_{\rm bas}\) replaces \(\Pi\) by \(\Pi H_{\rm bas}\) and coordinates by
\(H_{\rm bas}^{-1}x\). The inverse uses \(H_{\rm bas}\) and the inverse logarithm. These maps
preserve the base coordinate, zero section, contraction and collars.

Use wedge order \(12,13,14,23,24,34\). Choose \(a x_{\wedge}+b y_{\wedge}=1\) and put
\(s_0=x_{\wedge} e_{12}-y_{\wedge} e_{34}\). Set
\[
 B_2=[\wedge^2T_1-I\ \cdots\ \wedge^2T_N-I].
\]
Tracked integer row and column reduction gives a section $R_2$ satisfying
\begin{equation}
 B_2R_2=I-s_0\omega,\qquad \omega B_2=0,\qquad\omega(s_0)=1.
\hypertarget{tag-Q22N}{}\label{maps:eq:MT_local_14}
\end{equation}
Its columns give the monodromy-cylinder preimages. The inverse quotient maps
are
\begin{equation}
 \wedge^2\Lambda/\operatorname{im}B_2
 \ \underset{n\mapsto[ns_0]}{\overset{[v]\mapsto\omega(v)}{\rightleftarrows}}
 \mathbb Z\langle s\rangle.
\hypertarget{tag-Q22O}{}\label{maps:eq:MT_local_15}
\end{equation}
The difference \(v-\omega(v)s_0\) bounds the monodromy cylinders
with coefficients \(R_2v\). The five unit Smith entries make this the integral quotient.

Include the removed ordinary disc as the last site, with \(T=I\).
For \(D=[T_1-I\ \cdots\ T_N-I]\), the first \(N\) boundary maps are
block inclusions. The last is
\begin{equation}
 j_{N+1}(u)=-(T_2\cdots T_Nu,\ldots,T_Nu,u).
\hypertarget{tag-Q22P}{}\label{maps:eq:MT_local_16}
\end{equation}
It is obtained by tracing the polygon boundary, with its boundary
orientation. Its \(D\)-image is \(I-T_1\cdots T_N=0\).
At every cap take its fixed integral basis beginning with \(\delta\).
For $B=[j_iF_i]$, the tracked integer reductions give $V_0,Q_0,R_G$ with
\begin{equation}
 DV_0=0,\quad Q_0V_0=I,\quad BR_G=V_0Q_0,\quad
 \operatorname{im}V_0=\ker D,\quad
 \ker B=\mathbb Z\sum_i\delta_i.
\hypertarget{tag-Q22Q}{}\label{maps:eq:MT_local_17}
\end{equation}
Tracked elementary integer operations give a saturated kernel \(V_0\)
of $D$; its three nonzero Smith entries are one. The same reduction
of $B$ has unit nonzero Smith entries and nullity one. The displayed kernel vector has coefficient one in every
scalar position and is primitive.

Fixed loops sweep base boundary circles to two-tori. Addition is realized
by the period triangle with vertices \(0,u,u+v\); the opposite triangle
is the inverse chain. Projecting to the graph, using \(R_G\), subtracting
those sweeps and then using \eqref{maps:eq:MT_local_14} gives inverse coordinates
\begin{equation}
 H_2(\mathcal E)=\mathbb Zs\oplus
 \left(\bigoplus_i\operatorname{Fix}(T_i)\right)/
                    \mathbb Z\sum_i\delta_i.
\hypertarget{tag-Q22R}{}\label{maps:eq:MT_local_18}
\end{equation}
A line-frame change on the scalar sweep contributes
\(\delta\wedge\delta=0\). Clutching by \(e\gamma\) changes the polygon
relation by \(-\omega(\delta,e\gamma)=be\).

\item \emph{Cap maps, arbitrary finite parameters, and homology.}

In the cap formulas, \(A,g,\beta,T\) denote
\(A_i,g_i,\beta_i,T_i\) at the site under discussion, in its fixed physical marking.
At a rank-one multiplicative cap use the basis \(H_N=(l,h_{\rm per},p_0,w_0)\)
constructed in the primitive rank-one construction. Here \(x\) and \(y\)
denote coordinate columns in \(H_N\) and \(\wedge^2H_N\), respectively.
The normalized ruled component with its two sections identified gives
\begin{equation}
 f^{\rm N}_1(x)=(x_2,x_3,x_4),\qquad
 f^{\rm N}_2(y)=(y_{12},y_{23},y_{24},y_{34}).
\hypertarget{tag-Q22S}{}\label{maps:eq:MT_local_19}
\end{equation}
Inserting zero in the omitted coordinates and applying \(H_N\) or
\(\wedge^2H_N\) gives sections. The fibre kernel is generated by
\(e_{13},e_{14}\), with monodromy-cylinder preimages \(e_{23},e_{24}\).
Fixed sweeps bound the normal toric chains.

At the cusp use \(H_C=(SCe_1,Le_1,SCe_2,Le_2)\) of the prescribed-cusp construction.
Here the input columns are in \(H_C\) and \(\wedge^2H_C\) coordinates. Then
\begin{equation}
 f^{\rm C}_1(x)=(x_2,x_4),\qquad
 f^{\rm C}_2(y)=(y_{12},y_{14}-y_{23},y_{24},y_{34}).
\hypertarget{tag-Q22T}{}\label{maps:eq:MT_local_20}
\end{equation}
A section of the second map inserts \((u,0,v,0,w,z)\) and applies
\(\wedge^2H_C\). The primitive kernel is \(e_{13},e_{14}+e_{23}\);
monodromy-cylinder preimages are \(e_{23}\) and \(e_{24}-e_{23}\).
The normalization/conductor sequence gives these cycles and
the normal chains bounding both fixed sweeps. Substitution in the
displayed sections gives the identity and their Wang-image vanishing.

For an original additive cap, \(h_{\rm line}=a(A-I)^{-1}g\) is integral. Set
\(\lambda_{\rm line}=-h_{\rm line}^tJ\). The line-frame map and inverse are
\begin{equation}
 S_P(x,d,g_0)=(x,d+\lambda_{\rm line} x,g_0),\qquad
 S_P^{-1}(x,d',g_0)=(x,d'-\lambda_{\rm line} x,g_0).
\hypertarget{tag-Q22U}{}\label{maps:eq:MT_local_21}
\end{equation}
The identity \(\lambda_{\rm line}(A-I)+\beta=0\) removes the line shear on
the neighbourhood. Zero line multidegree gives a topological
line trivialization there. The positive-roof mapping torus for the
contracting arrow has a unit-roof parametrization: choose a smooth
nonnegative bump \(b_y\) of integral one on a roof of length \(r_{\rm roof}(y)\),
choose \(0<k_{\rm roof}<(\max r_{\rm roof})^{-1}\), and use the increasing map
\begin{equation}
 f_y(t)=\int_0^t\{k_{\rm roof}+(1-k_{\rm roof}r_{\rm roof}(y))b_y(v)\}\,dv
\hypertarget{tag-Q22V}{}\label{maps:eq:MT_local_22}
\end{equation}
and its inverse. The ends have matching derivatives. Translation by
the identity-component N\'eron logarithm times the roof coordinate,
with inverse negative translation and a cutoff at the ends, removes
that part of the return. The finite component translation has one
orbit. The reciprocal line gauges have the original node constants.
These inverse maps identify the original cap and its indicated mapping torus.

Put \(z_Q=b(A-I)^{-1}g\) and \(q_{\rm vert}=1\) for \(II\), \(q_{\rm vert}=2\) for \(III\).
The original core has \(H_2\) basis \([C_{\rm core}],[\delta\gamma]\), with
elliptic fibre \(q_{\rm vert}[C_{\rm core}]\). In coordinates after \eqref{maps:eq:MT_local_21}, its map is
\begin{equation}
 f^{\rm orig}_2(y)=\bigl(q_{\rm vert}(y_{12}+z_{Q,2}y_{14}-z_{Q,1}y_{24}),y_{34}\bigr).
\hypertarget{tag-Q22W}{}\label{maps:eq:MT_local_23}
\end{equation}
The first coordinate is the pullback under \(x\mapsto x+z_Qx_\gamma\),
whose inverse subtracts that term. An integral B\'ezout solution of
\(q_{\rm vert}A_{\rm Bez}+q_{\rm vert}z_{Q,2}B_{\rm Bez}-q_{\rm vert}z_{Q,1}C_{\rm Bez}=1\), followed by \(\wedge^2S_P^{-1}\),
gives a section on \([C_{\rm core}]\); \(e_{34}\) gives the other section.
These two sections recover the stated generators. The fixed basis is
\begin{equation}
                  \delta,\quad\ell_0=(-q_{\rm vert}z_{Q,1},-q_{\rm vert}z_{Q,2},0,q_{\rm vert}).
\hypertarget{tag-Q22X}{}\label{maps:eq:MT_local_24}
\end{equation}
For original III in row 9 it is
\(\ell_0=(-15,45,0,2)\). Its two sweeps bound the origin disc and
the two-fold forward orbit tube closed by the same logarithm and line
frame. The gamma coordinate of this primitive return is two.

At a finite cap, write \(z_{\rm site}=z_i=(A_i-I)^{-1}g_i\).
Then \(\ell=(-b(z_{\rm site})_1,-b(z_{\rm site})_2,0,1)\) is integral and
\begin{equation}
 m=d_i k>0,\qquad v_{\rm fin}=n\ell+r\delta,\qquad
 \gcd(m,n)=1,\quad r\in\mathbb Z.
\hypertarget{tag-Q22Y}{}\label{maps:eq:MT_local_25}
\end{equation}
For finite II, \(\ell=\gamma,\beta_{\rm fin}=\delta,q_{\rm vert}=1\). For finite III in
row 8,
\begin{equation}
 \ell=(-15,45,0,1),\quad\beta_{\rm fin}=(0,-1,0,0),\quad q_{\rm vert}=2.
\hypertarget{tag-Q22Z}{}\label{maps:eq:MT_local_26}
\end{equation}
Integer reduction of $[T-I\mid\beta_{\rm fin}\mid\ell]$ gives a right
inverse and an integral $y_{\rm norm}$ with
\begin{equation}
 q_{\rm vert}\beta_{\rm fin}+(T-I)y_{\rm norm}=\delta,\quad
 N_{d_i}\beta_{\rm fin}=(d_i/q_{\rm vert})\delta,\quad
 N_{d_i}\ell=d_i\ell,\quad-\omega(\beta_{\rm fin},\ell)=b/q_{\rm vert}.
\hypertarget{tag-Q230}{}\label{maps:eq:MT_local_27}
\end{equation}
Its last two inverse rows give the primitive coinvariant coordinates.
This distinguishes the gamma-one finite lift \eqref{maps:eq:MT_local_26} from \eqref{maps:eq:MT_local_24}.

The affine representative is \(R_{\rm aff}(x)=Tx+v_{\rm fin}/m\). Another lift
\(Tx+t_{\rm aff}\) with \(N_m t_{\rm aff}=v_{\rm fin}\) is conjugate to this one by \(x\mapsto x+h_{\rm aff}\),
where \((I-T)h_{\rm aff}=t_{\rm aff}-v_{\rm fin}/m\) on the complement of
\(\pi_{\rm fix}=N_m/m\); \(I-T\) is invertible on that complement. Translation
by \(-h_{\rm aff}\) is the inverse and commutes with scalar translation.
For the holomorphic units on the whole good root disc, write \(A_{\rm root}(u)\)
for the linear part of the retained good-root action, and let
\(\mathcal Rf(u)=A_{\rm root}(\sigma_{\rm root}^{-1}u)f(\sigma_{\rm root}^{-1}u)\), where
\(\sigma_{\rm root}(u)=\zeta_{\rm root}u\) and \(\mathcal R^m=I\).
Put \(\mathcal N=\sum_{j=0}^{m-1}\mathcal R^j\).
For the two affine terms with difference \(d_{\rm hol}\) and
\(\mathcal N d_{\rm hol}=0\), use
\begin{equation}
 h_{\rm hol}=-\frac1m\sum_{j=0}^{m-1}j\mathcal R^j d_{\rm hol},\quad
 (I-\mathcal R)h_{\rm hol}=d_{\rm hol}.
\hypertarget{tag-Q231}{}\label{maps:eq:MT_local_28}
\end{equation}
The inverse is \(-h_{\rm hol}\), so the comparison is over the disc and its collar,
with the units present. A peripheral period change uses the norm
\(v_{\rm fin}\mapsto v_{\rm fin}+N_m\lambda\), with fixed-sweep shear
\(t_{\rm swp}(u)\mapsto t_{\rm swp}(u)+\omega(u,\lambda)s\). The inverse uses \(-\lambda\).
On the fixed lattice, $N_m\lambda=m\lambda$.

The boundary map and its inverse in the chosen affine representative are
\begin{equation}
 [y,\tau]\longmapsto[y+(\tau/m)v_{\rm fin},e^{2\pi i\tau/m}],\qquad
 [x,z]\longmapsto[x-(\tau/m)v_{\rm fin},\tau].
\hypertarget{tag-Q232}{}\label{maps:eq:MT_local_29}
\end{equation}
A different root argument changes both sides by the mapping-torus
relation. Write \(\mu_{\partial}\) for the positive boundary-base
generator in these mapping-torus coordinates. The positive normal meridian is
\(m\mu_{\partial}-v_{\rm fin}\).
Sectors of its meridional disc and the period parallelograms give,
in \((s,D_\delta=t_{\rm swp}(\delta),L_\ell=t_{\rm swp}(\ell))\), the full integral Thom columns
\begin{equation}
 C_{1,\rm Th}=(bn/q_{\rm vert},m/q_{\rm vert},0),\quad C_{2,\rm Th}=(-br,0,m),\quad C_{3,\rm Th}=(0,r,n).
\hypertarget{tag-Q233}{}\label{maps:eq:MT_local_30}
\end{equation}
The opposite sectors give the inverse boundary chains.

To check their integral scope, the finite central surface \(F_{\rm fin}\) has
\begin{equation}
 H_1(F_{\rm fin})=\langle\beta_{\rm fin},\ell,c_0\mid
                  mc_0=n\ell+r q_{\rm vert}\beta_{\rm fin}\rangle.
\hypertarget{tag-Q234}{}\label{maps:eq:MT_local_31}
\end{equation}
For \(A_{\rm fin}m+B_{\rm fin}n=1\), the basis
\(\beta_{\rm fin},z_{\rm fin}=A_{\rm fin}\ell+B_{\rm fin}c_0\) has inverse
\(\ell=m z_{\rm fin}-B_{\rm fin}r q_{\rm vert}\beta_{\rm fin}\), \(c_0=n z_{\rm fin}+A_{\rm fin}r q_{\rm vert}\beta_{\rm fin}\).
Thus \(H_1(F_{\rm fin})\) is free. The finite torus cover and transfer give rational
Betti numbers \(1,2,2,2,1\); integral duality then makes \(H_2,H_3\)
free as well. The flat normal line has zero Chern class in free
\(H^2(F_{\rm fin};\mathbb Z)\). Transfer of \(\lambda\) cuts into sectors with
fibre area \(-\lambda\wedge v_{\rm fin}\) and sweep \(N_m\lambda\); the torus for
\(c_0\) has fibre area zero and sweep \(v_{\rm fin}\). Rejoining those sectors
is the inverse chain map. Thom exactness on the basis above gives
exactly the integral subgroup generated by \(C_{1,\rm Th},A_{\rm fin}C_{2,\rm Th}+B_{\rm fin}C_{3,\rm Th}\).

There is at most one finite cap. Its full fundamental-group presentation,
after the angular relations of the common global Borel comparison, is
\begin{equation}
 \langle\gamma,c_0\mid[\gamma,c_0]=1,\ c_0^m=\gamma^n,\ c_0=\gamma^e\rangle.
\hypertarget{tag-Q235}{}\label{maps:eq:MT_local_32}
\end{equation}
Hence simple connectivity is precisely \(me-n=\varepsilon_{\rm det}=\pm1\).
In the rows without a finite cap it is \(e=\pm1\).
Choose \(A_{\rm fin}=\varepsilon_{\rm det} e,B_{\rm fin}=-\varepsilon_{\rm det}\). With
\(K_0=A_{\rm fin}C_{2,\rm Th}+B_{\rm fin}C_{3,\rm Th}\), the exact change of column generators is
\begin{equation}
 \begin{split}
 K_0&=(-A_{\rm fin}br,B_{\rm fin}r,1),\\
 C_{2,\rm Th}-mK_0&=-B_{\rm fin}q_{\rm vert}rC_{1,\rm Th},\qquad
 C_{3,\rm Th}-nK_0=A_{\rm fin}q_{\rm vert}rC_{1,\rm Th}.
 \end{split}
\hypertarget{tag-Q236}{}\label{maps:eq:MT_local_33}
\end{equation}
The ordinary polygon relation is \(D_\delta=-be\,s\). It gives \(L_\ell=0\) and
\begin{equation}
                   (b/q_{\rm vert})(n-me)s=-\varepsilon_{\rm det} c\,s=0.
\hypertarget{tag-Q237}{}\label{maps:eq:MT_local_34}
\end{equation}
Substituting $A_{\rm fin}=\varepsilon_{\rm det} e$, $B_{\rm fin}=-\varepsilon_{\rm det}$ and
$me-n=\varepsilon_{\rm det}$ proves \eqref{maps:eq:MT_local_33}--\eqref{maps:eq:MT_local_34} for both signs and all
integers $m,e,r$ in the stated domain. For no finite cap the remaining relation is
\(be\,s=0\), and \(c=b\). All cap maps are surjective and \eqref{maps:eq:MT_local_17} is
integral, so Mayer--Vietoris gives \(H_2(X)=\mathbb Z/c\langle s\rangle\).
The two cusp vertices give Euler number two, and the other torus strata
and caps have Euler number zero. Duality and universal coefficients give
\eqref{maps:eq:MT_local_1}. Since \(H^2(X;\mathbb Z)=0\), the complex first Chern class is
zero, so \(w_2=0\); the spin structure is unique because \(H^1(X;\mathbb F_2)=0\).

\item \emph{Primitive rank-one periods and all scalar isotropy.}

At an unselected multiplicative site put
\begin{equation}
 d_0=\gcd(\nu,a\kappa),\quad n_0=\nu/d_0,\quad
 A_0=a\kappa/d_0,\quad B_0=b\kappa/n_0.
\hypertarget{tag-Q238}{}\label{maps:eq:MT_local_35}
\end{equation}
Primitivity makes \(B_0\) integral and \(\gcd(d_0,B_0)=1\).
Choose \(n_0\theta_N+A_0\sigma_{\rm Bez}=1\), \(d_0x_L+B_0y_L=1\), and set
\begin{equation}
 C_S=\begin{pmatrix}n_0&\sigma_{\rm Bez}\\-A_0&\theta_N\end{pmatrix},\qquad
 C_L=\begin{pmatrix}x_L&-B_0\\y_L&d_0\end{pmatrix}.
\hypertarget{tag-Q239}{}\label{maps:eq:MT_local_36}
\end{equation}
Their literal inverses are
\begin{equation}
 C_S^{-1}=\begin{pmatrix}\theta_N&-\sigma_{\rm Bez}\\A_0&n_0\end{pmatrix},\qquad
 C_L^{-1}=\begin{pmatrix}d_0&B_0\\-y_L&x_L\end{pmatrix}.
\hypertarget{tag-Q23A}{}\label{maps:eq:MT_local_37}
\end{equation}
With \(S=(u,\delta),L=(v,\gamma)\), write
\((l,p_0)=SC_S\), \((h_{\rm per},w_0)=LC_L\). The ordered period basis
\(H_N=(l,h_{\rm per},p_0,w_0)\) has determinant one and return \(h_{\rm per}\mapsto h_{\rm per}+l\).
The other columns are fixed. Its pairing is
\begin{equation}
\begin{gathered}
 K_N=(SC_S)^t\omega(LC_L)
       =\begin{pmatrix}n_{\rm per}&0\\M&r_{\rm per}\end{pmatrix},\\
 n_{\rm per}=\frac{a\nu}{d_0^2},\quad r_{\rm per}=-\frac{bd_0^2}{\nu},\qquad
 \delta=-\sigma_{\rm Bez} l+n_0p_0.
\end{gathered}
\hypertarget{tag-Q23B}{}\label{maps:eq:MT_local_38}
\end{equation}
Their determinants are one; the displayed products with their inverse matrices are the identity.

The chart identification is \eqref{maps:eq:MT_local_13} in this basis, followed by
exponentiation of \(l\). Consecutive charts and their inverses are
\begin{equation}
 X_N=z_N/q^j,\quad Y_N=q^{j+1}/z_N,\quad X_NY_N=q;\qquad z_N=X_Nq^j,\ q=X_NY_N.
\hypertarget{tag-Q23C}{}\label{maps:eq:MT_local_39}
\end{equation}
The primitive long period identifies successive charts with the original
unit multipliers transformed by \eqref{maps:eq:MT_local_36}. The inverse uses \eqref{maps:eq:MT_local_37} and \eqref{maps:eq:MT_local_39}.
The core is one ruled normalization with zero and infinity sections
identified along the elliptic curve of periods \(p_0,w_0\).

On that elliptic curve, the scalar circle translates the first factor
by \(n_0\theta p_0\). Thus a subgroup \(C_p\) fixes it exactly when
\(p\mid n_0\). Otherwise it acts there by a nonzero translation.
At a fixed elliptic curve, away from the double locus, its nonzero
normal coordinate has a unit character, so there is no extra fixed
point on the ruled normalization. Substitution of the ordered
component data gives, for each $p\in\{2,3,5\}$ dividing $c$, exactly one such site,
of type \(I_p\), with
\begin{equation}
 d_0=1,\quad n_0=p,\quad n_{\rm per}=ap,\quad r_{\rm per}=-b/p.
\hypertarget{tag-Q23D}{}\label{maps:eq:MT_local_40}
\end{equation}
Call its elliptic curve \(E_p\). In the subgroup generator its
normal characters are \((-\sigma_{\rm Bez},\sigma_{\rm Bez})\). For odd \(p\), the
positive-valuation Tate residue is \(r_{\rm Tate}=-\kappa_i\) and
\begin{equation}
             w_{\rm norm}=(a r_{\rm Tate})^{-1}=-\sigma_{\rm Bez}\pmod p.
\hypertarget{tag-Q23E}{}\label{maps:eq:MT_local_41}
\end{equation}
This uses \(\sigma_{\rm Bez} a\kappa_i=1\pmod p\).

For example, the unselected
\(I_5\) at \(a=5,b=6\) has \(d_0=5,n_0=1\). Its scalar translation
has no nontrivial finite stabilizer. The displayed basis and its inverse identify that primitive
conductor and its nontrivial scalar translation.

\item \emph{The prescribed cusp and the remaining isotropy.}

At the selected site set
\begin{equation}
 \rho=\frac{1-ab\kappa^2}{\nu},\qquad
 C=\begin{pmatrix}\nu&b\kappa\\-a\kappa&\rho\end{pmatrix},\qquad
 C^{-1}=\begin{pmatrix}\rho&-b\kappa\\a\kappa&\nu\end{pmatrix}.
\hypertarget{tag-Q23F}{}\label{maps:eq:MT_local_42}
\end{equation}
Equation \eqref{maps:eq:MT_local_6} makes \(C\) integral with determinant one.
The return satisfies \((T-I)L=SC\), \((T-I)S=0\).
The prescribed Jordan fan therefore has short basis \(SC\).
The physical scalar direction in that fan is
\begin{equation}
                  d_\delta=C^{-1}(0,1)^t=(-b\kappa,\nu)^t.
\hypertarget{tag-Q23G}{}\label{maps:eq:MT_local_43}
\end{equation}
Let \(k_{\rm edge}=b\kappa\). For each prime \(p\mid c\), \(p\mid k_{\rm edge}\) and
\(\nu\) is a unit modulo \(p\). Thus \eqref{maps:eq:MT_local_43} is vertical and nonzero
modulo \(p\). Among the three edge directions of the prescribed unit
triangulation, only the vertical one is parallel to it modulo \(p\).
There is one such edge orbit. Its compact curve \(C_f\) is described
by the adjacent charts
\begin{equation}
 (x,y,w)=(q/(X_1X_2),X_2,X_1),\qquad
 (x',y',w')=(q/X_2,X_1X_2,X_1^{-1}).
\hypertarget{tag-Q23H}{}\label{maps:eq:MT_local_44}
\end{equation}
Their inverses are
\((X_1,X_2,q)=(w,y,xyw)\) and
\((1/w',y'w',x'y'w')\). Their overlap and inverse are
\begin{equation}
 (x',y',w')=(wx,wy,w^{-1}),\qquad
 (x,y,w)=(w'x',w'y',(w')^{-1}).
\hypertarget{tag-Q23I}{}\label{maps:eq:MT_local_45}
\end{equation}
Consequently \(C_f\simeq\mathbb P^1\) is embedded and its normal
bundle is \(\mathcal O(-1)\oplus\mathcal O(-1)\).
A translation stabilizing a bounded edge fixes its endpoints and is
zero. Every point of the quotient curve therefore has a unique
representative on this chosen edge curve. Choosing that representative
is the inverse of the quotient map.

The full-circle endpoint weights are
\begin{equation}
           (k_{\rm edge}-\nu,\nu,-k_{\rm edge}),\qquad(-\nu,\nu-k_{\rm edge},k_{\rm edge}).
\hypertarget{tag-Q23J}{}\label{maps:eq:MT_local_46}
\end{equation}
For \(k_{\rm edge}=0\) the circle fixes \(C_f\); otherwise it rotates \(C_f\)
with tangent weight \(-k_{\rm edge}\). Its generic stabilizer then has order
\(|k_{\rm edge}|\). In every case \(C_p\) fixes precisely this curve within the
cusp, with opposite nonzero normal characters. The two triangle
orbits give the two endpoints already on \(C_f\).

The scalar circle is free on every original additive cap. In an adapted
norm the forward arrow strictly contracts. A unit-modulus scalar
cannot equal a nonzero iterate on a nonzero line vector; the inverse
iterate strictly expands. The zero iterate fixes only the identity
scalar. The roof and logarithm maps of the cap-map construction intertwine
this scalar action.

At a finite cap a scalar fixed point would satisfy, on its good cover,
\(x+\theta\delta=R_{\rm aff}^{\,j}x\) modulo the lattice. Applying the integral
covector \(\chi\) gives \(jn/m\in\mathbb Z\); coprimality forces
\(m\mid j\). The remaining scalar translation is primitive and is
fixed only for integral \(\theta\). Away from the central fibre the
base rotation already gives \(m\mid j\). This proves freeness for
every allowed \(m,n,r\), including \(3\mid m\), and every affine
lift through \eqref{maps:eq:MT_local_28}. The ordinary cap is free by the same primitive
scalar translation.

Therefore, on each relevant row,
\begin{equation}
                     X^{C_p}=C_f\amalg E_p\qquad(p\mid c).
\hypertarget{tag-Q23K}{}\label{maps:eq:MT_local_47}
\end{equation}
Each prime subgroup has free complement. Every nontrivial subgroup
of \(C_9\) contains \(C_3\), so \(C_9\) is free off
\(C_f\amalg E_3\). Its normal sphere bundles there are free, since
their \(C_3\) normal characters are nonzero. These statements apply to every $k_{\rm edge}=b\kappa$ in the table.

\item \emph{One line, its fan support, and its inverse gauges.}

Let \(Y\) be the partial completion at all multiplicative sites,
leaving the ordinary and additive discs open. It contains all fixed
curves in \eqref{maps:eq:MT_local_47}.

For every row there is a doubled line \(\mathcal D\) on the regular
period cover with factors
\begin{equation}
 J_\lambda(x,\zeta_R)=(x+\lambda,e^{2\pi i\omega(\lambda,x)}\zeta_R).
\hypertarget{tag-Q23L}{}\label{maps:eq:MT_local_48}
\end{equation}
The inverse is \(J_{-\lambda}\); the cocycle quotient is
\(e^{2\pi i\omega(\lambda,\mu)}=1\).
The connection \(d-2\pi i\omega(x,dx)\) gives \(c_1=2\omega\).
The monodromy lift is \((x,\zeta_R)\mapsto(T_i x,\zeta_R)\), with
inverse \(T_i^{-1}\). For \(p=3,5\) dividing \(c\), the lift of
\(\delta/p\) is
\begin{equation}
 (x,\zeta_R)\longmapsto(x+\delta/p,e^{2\pi i\omega(\delta/p,x)}\zeta_R).
\hypertarget{tag-Q23M}{}\label{maps:eq:MT_local_49}
\end{equation}
Its commutator with \eqref{maps:eq:MT_local_48} is
\(e^{-4\pi i b\lambda_\gamma/p}=1\), and its \(p\)-th power is
\(J_\delta\). It commutes with every monodromy lift.

For the nonzero short period $\omega_s$, use
$z_E=z_{\rm add}/\omega_s$, with inverse $z_{\rm add}=\omega_s z_E$.
Put $\tau_E=\omega_l/\omega_s$, $\zeta_G=G_{\rm lift}/\omega_s$,
and $z=e^{2\pi i z_E}$; the inverse logarithm is modulo the short
period. Let $\rho_L$ be the forward line-period logarithm.
The theta ratio
$\vartheta(z_E-a\zeta_G)/\vartheta(z_E)$, with reciprocal inverse,
identifies \(\mathcal O(aG-O)\) with its Poincar\'e factors. Shifting
the theta sum gives short factor one and long factor
\(e^{2\pi i a\zeta_G}\). The forward lift of \(-bG\) has constant
multiplier in this trivialization, because its ratio with the
constant-factor lift is a holomorphic function on the compact elliptic
curve. This constant includes its base-dependent holomorphic unit.

In physical bases \(S=(u,\delta),L=(v,\gamma)\), write
\begin{equation}
 D_{\rm pol}=\operatorname{diag}(a,-b),\quad
 \Pi_{\rm ph}=\begin{pmatrix}\tau_E&-b\zeta_G\\a\zeta_G&\rho_L\end{pmatrix},
 \qquad D_{\rm pol}\Pi_{\rm ph}=(D_{\rm pol}\Pi_{\rm ph})^t.
\hypertarget{tag-Q23N}{}\label{maps:eq:MT_local_50}
\end{equation}
Write \(x=SU_{\rm ph}+LW_{\rm ph}\) and
\(Z_{\rm ph}=U_{\rm ph}+\Pi_{\rm ph}W_{\rm ph}\).
These are the real coordinates in the ordered physical short and long bases
and the corresponding normalized complex coordinates. Thus
\(z=\exp(2\pi i(Z_{\rm ph})_1)\) and
\(v_{\rm lin}=\exp(2\pi i(Z_{\rm ph})_2)\), with inverse logarithms
modulo the same marked periods. Put
\(h_1=\exp(2\pi i\tau_E)\), \(h_G=\exp(2\pi i\zeta_G)\) and
\(R=\exp(2\pi i\rho_L)\); these retain the original holomorphic units.
The invertible gauge
\begin{equation}
 G_D=\exp\{-2\pi i(U_{\rm ph}^tD_{\rm pol}W_{\rm ph}+W_{\rm ph}^tD_{\rm pol}\Pi_{\rm ph}W_{\rm ph})\}
\hypertarget{tag-Q23O}{}\label{maps:eq:MT_local_51}
\end{equation}
and its reciprocal give short factors one and long factors
\(h_1^{-a}z^{-2a}\), \(R^b v_{\rm lin}^{2b}\). Multiplication by \(z^a\),
inverted by \(z^{-a}\), gives the exact period and inversion maps
\begin{equation}
 \begin{split}
 h_{\mathcal D}(z,v_{\rm lin},\zeta)&=(h_1z,h_G^a v_{\rm lin},z^{-2a}\zeta),\\
 F_{\mathcal D}(z,v_{\rm lin},\zeta)&=(h_G^{-b}z,R v_{\rm lin},h_G^{-ab}R^b v_{\rm lin}^{2b}\zeta),\\
 J_{\mathcal D}(z,v_{\rm lin},\zeta)&=(z^{-1},v_{\rm lin}^{-1},z^{-2a}\zeta).
 \end{split}
\hypertarget{tag-Q23P}{}\label{maps:eq:MT_local_52}
\end{equation}
The line factors of the inverses of \(h_{\mathcal D},F_{\mathcal D}\), in target coordinates,
are \(h_1^{-2a}z^{2a}\) and \(h_G^{ab}R^b v_{\rm lin}^{-2b}\);
\(J_{\mathcal D}^{-1}=J_{\mathcal D}\). Substitution gives \(h_{\mathcal D}F_{\mathcal D}=F_{\mathcal D}h_{\mathcal D}\), \(J_{\mathcal D}h_{\mathcal D}J_{\mathcal D}=h_{\mathcal D}^{-1}\),
\(J_{\mathcal D}F_{\mathcal D}J_{\mathcal D}=F_{\mathcal D}^{-1}\). The functions $h_1,h_G,R$ include their holomorphic units. On overlaps compose one Gaussian, the 
integral period change, and the reciprocal Gaussian of the other
chart. Triple composites cancel, so this is one line on the regular
part of \(Y\). Its degree on \(s\) is two and on every constant fixed
loop sweep is zero.

At a rank-one cap use \eqref{maps:eq:MT_local_36} and put
\[
 \begin{aligned}
 U_N&=C_S^{-1}U_{\rm ph},& W_N&=C_L^{-1}W_{\rm ph},\\
 Z_N&=C_S^{-1}Z_{\rm ph},& \Pi_N&=C_S^{-1}\Pi_{\rm ph}C_L.
 \end{aligned}
\]
The inverse is
\((U_{\rm ph},W_{\rm ph},Z_{\rm ph})=(C_SU_N,C_LW_N,C_SZ_N)\),
and \(Z_N=U_N+\Pi_NW_N\). The rank-one chart coordinate is
\(z_N=\exp(2\pi i(Z_N)_1)\). The identity
\(K_N=C_S^tD_{\rm pol}C_L\) gives
\[
 U_N^tK_NW_N=U_{\rm ph}^tD_{\rm pol}W_{\rm ph},\qquad
 K_N^t\Pi_N=C_L^tD_{\rm pol}\Pi_{\rm ph}C_L.
\]
Thus the same Gaussian in these coordinates is
\begin{equation}
 \exp\{-2\pi i(U_N^tK_NW_N+W_N^tK_N^t\Pi_NW_N)\}.
\hypertarget{tag-Q23Q}{}\label{maps:eq:MT_local_53}
\end{equation}
The symmetry required is that of \(K_N^t\Pi_N\); the full off-diagonal
entry \(M\) remains. For a long-period coefficient column
\(\mathbf n_N\in\mathbb Z^2\) in \((h_{\rm per},w_0)\), its long factors have exponent
\(-4\pi i\mathbf n_N^tK_N^tZ_N-2\pi i\mathbf n_N^tK_N^t\Pi_N\mathbf n_N\).
Write \(r_N=(\mathbf n_N)_1\). The first long generator advances the fan index
by one, so the integer support \(n_{\rm per}j^2\) has difference
\(2n_{\rm per}r_Nj+n_{\rm per}r_N^2\) under \(j\mapsto j+r_N\), and
its affine monomials with reciprocal inverses extend the line through
\eqref{maps:eq:MT_local_39}. The remaining factors are the unit functions.
Replacing \(\Pi_N\) by \(\Pi_N+E_{11}\) and replacing \(U_N\) by
\(U_N+E_{11}W_N\) add the same \(n_{\rm per}(W_N)_1^2\), proving around-base compatibility.

For the nonprimitive physical first period put
\(\alpha_0=d_0\), \(v_0=(n_0,-A_0)\), \(\beta_{\rm val}=B_0\).
The two physical valuation columns are \(\alpha_0v_0,\beta_{\rm val}v_0\)
and
\begin{equation}
 a\beta_{\rm val}(v_0)_1+b\alpha_0(v_0)_2=0,\qquad
 \psi_1(j)=-\frac{a(v_0)_1}{\alpha_0}j(j-\alpha_0)
                         =-n_{\rm per}j^2+n_{\rm per}\alpha_0j.
\hypertarget{tag-Q23R}{}\label{maps:eq:MT_local_54}
\end{equation}
Its two differences are exactly the physical exponents of \eqref{maps:eq:MT_local_52}.
The linear term is the \(z^a\) gauge and the quadratic sign changes
between frame and vector coordinates. This support comparison uses the physical first column
$\alpha_0v_0$, including the case $d_0=5$.
At \(E_p\) the doubled line has degree \(-2b/p\), and its character
is \(-\sigma_{\rm Bez} n_{\rm per}(2j+1)=0\pmod p\) for odd \(p\).

At the cusp the short and long bases are \(SC,L\). Put
\[
 \begin{aligned}
 U_C&=C^{-1}U_{\rm ph},& W_C&=W_{\rm ph},\\
 Z_C&=C^{-1}Z_{\rm ph},& \Pi_C&=C^{-1}\Pi_{\rm ph}.
 \end{aligned}
\]
The inverse is
\((U_{\rm ph},W_{\rm ph},Z_{\rm ph})=(CU_C,W_C,CZ_C)\),
and \(Z_C=U_C+\Pi_CW_C\). Since \(C^tD_{\rm pol}=D_{\rm pol}C\),
the symmetric pairing is
\begin{equation}
 K=C^tD_{\rm pol}=
 \begin{pmatrix}a\nu&ab\kappa\\ab\kappa&b(ab\kappa^2-1)/\nu\end{pmatrix},
 \qquad Kd_\delta=(0,-b)^t.
\hypertarget{tag-Q23S}{}\label{maps:eq:MT_local_55}
\end{equation}
In particular
\[
 U_C^tKW_C=U_{\rm ph}^tD_{\rm pol}W_{\rm ph},\qquad
 K\Pi_C=D_{\rm pol}\Pi_{\rm ph}.
\]
The Gaussian \(\exp\{-2\pi i(U_C^tKW_C+W_C^tK\Pi_CW_C)\}\)
is therefore the same physical gauge. Define
\(B_C(q)=\Pi_C-(\log q)/(2\pi i)I\) on the cut disc; it retains the
full holomorphic period-unit matrix. Let
\(\mathbf X_C=(X_1,X_2)=(e^{2\pi i(Z_C)_1},e^{2\pi i(Z_C)_2})\),
using the same dense coordinates as the adjacent cusp charts. For
\(\mathbf m\in\mathbb Z^2\) write
\(\mathbf X_C^{\mathbf m}=\exp(2\pi i\mathbf m^tZ_C)\).
For a long-period coefficient \(\mathbf n_C\in\mathbb Z^2\), the long factors are
\[
 q^{-\mathbf n_C^tK\mathbf n_C}\mathbf X_C^{-2K\mathbf n_C}
 \exp(-2\pi i\mathbf n_C^tK B_C(q)\mathbf n_C).
\]
Thus the original holomorphic factors remain present. The support
\(\psi(v)=v^tKv\) on the prescribed unit triangles has difference
\(2\mathbf n_C^tKv+\mathbf n_C^tK\mathbf n_C\) under
\(v\mapsto v+\mathbf n_C\).
Let \(e_{\mathcal D}\) be the reference frame with fibre coordinate \(\zeta\).
For a triangle slope \(\mathbf m_C\), its regular frame is
\(e_{\mathbf m_C}=\mathbf X_C^{-\mathbf m_C}e_{\mathcal D}\) and vector coordinate is
\(\zeta_{\mathbf m_C}=\mathbf X_C^{\mathbf m_C}\zeta\);
the opposite monomials give the inverse.

The physical support is
\begin{equation}
 \psi_{\rm ph}(x,y)=-a\rho x^2+2ab\kappa xy+b\nu y^2+ax,\qquad
 \psi_{\rm ph}(Cv)=-v^tKv+a\,\operatorname{row}_1(C)v.
 \hypertarget{tag-020K}{}\label{maps:eq:D.MT2}
\end{equation}
Direct expansion proves the identity. The inverse period map
is \(C^{-1}\), the reciprocal frame changes the quadratic sign, and
the displayed integral affine term is the \(z^a\) monomial gauge,
with its reciprocal as inverse. These maps identify the four physical
wall vertices with the wall of the prescribed fan. On the edge \([0,e_2]\), the slopes are
\begin{equation}
 m_+=(K_{11},K_{22}),\qquad
 m_-=(2K_{12}-K_{11},K_{22}).
\hypertarget{tag-Q23T}{}\label{maps:eq:MT_local_57}
\end{equation}
Their character and degree are
\begin{equation}
 c_{\mathcal D,C}=b\{a\kappa(b\kappa-\nu)-1\},\qquad
 \deg\mathcal D|C_f=2(K_{11}-K_{12})=2a(\nu-b\kappa).
\hypertarget{tag-Q23U}{}\label{maps:eq:MT_local_58}
\end{equation}
The character is zero modulo every odd \(p\mid c\).
All base units have local logarithms, are scalar invariant, and have
zero degree on the central curve; their factors remain in the
maps above.

For the even sectors \(c=6,30\), the table gives \(a\) odd and \(b\)
even. There is an undoubled line \(\mathcal L\) with factors
\begin{equation}
 \mathcal Q(\lambda)=\lambda_1\lambda_2+\lambda_1+\lambda_2\pmod2,
 \quad
 J^\mathcal L_\lambda(x,\zeta_R)=
 (x+\lambda,(-1)^{\mathcal Q(\lambda)}e^{\pi i\omega(\lambda,x)}\zeta_R).
\hypertarget{tag-Q23V}{}\label{maps:eq:MT_local_59}
\end{equation}
Its polar form is \(\omega\bmod2\). The form \(\mathcal Q\) is one
on every nonzero elliptic vector modulo two; each \(A_i\) is
invertible there, and the elliptic gamma column \(bg_i\) is even.
Thus every \(T_i\) preserves \(\mathcal Q\). This proves the
cocycle law and the same scalar-one monodromy lift. Its connection
has Chern class \(\omega\), degree one on \(s\), zero on sweeps.
The scalar lift by \(\delta/N_{\rm sc}\), \(N_{\rm sc}\mid b\), has factor
\(e^{\pi i\omega(\delta/N_{\rm sc},x)}\), commutator
\(e^{-2\pi i b\lambda_\gamma/N_{\rm sc}}=1\), and \(N_{\rm sc}\)-th power \(J^\mathcal L_\delta\).

The physical Gaussian is
\begin{equation}
 \exp\{\pi i[-U_{\rm ph}^tD_{\rm pol}W_{\rm ph}-W_{\rm ph}^tD_{\rm pol}\Pi_{\rm ph}W_{\rm ph}+(Z_{\rm ph})_1]\},
\hypertarget{tag-Q23W}{}\label{maps:eq:MT_local_60}
\end{equation}
followed by \(z^{(a-1)/2}\). Its reciprocal is the inverse. The 
maps corresponding to \eqref{maps:eq:MT_local_52} have factors
\begin{equation}
 h_{\mathcal L}:\ (-1)^a z^{-a},\qquad
 F_{\mathcal L}:\ h_G^{-ab/2}R^{b/2}v_{\rm lin}^b,\qquad
 J_{\mathcal L}:\ z^{-a}.
\hypertarget{tag-Q23X}{}\label{maps:eq:MT_local_61}
\end{equation}
Their inverse factors are
\((-1)^a h_1^{-a}z^a\), \(h_G^{ab/2}R^{b/2}v_{\rm lin}^{-b}\), and the
same \(J_{\mathcal L}\). All unit exponents are integers.

In the cusp coordinates the short refinement is
\(\epsilon_{\rm ref}=(1,0)^t\), and the support is
\begin{equation}
 \psi_{\mathcal L}(v)=\frac{v^tKv-v_1}{2},\qquad
 \frac12\psi_{\rm ph}(Cv)=-\psi_{\mathcal L}(v)+L_0v,\quad
 L_0=\frac{a\,\operatorname{row}_1C-(1,0)}2.
\hypertarget{tag-Q23Y}{}\label{maps:eq:MT_local_62}
\end{equation}
Both the support and \(L_0\) are integral: \(K_{11}\) is odd,
\(K_{22}\) is even, and the same parity gives \(L_0\).
After the reciprocal frame identification, the dense monomial
\(\mathbf X_C^{L_0^t}\), with inverse \(\mathbf X_C^{-L_0^t}\), identifies the regular
Cartier coordinates. The Gaussian scalar phase adds
\(\epsilon_{\rm ref}^td_\delta/2\). Hence
\begin{equation}
 \deg\mathcal L|C_f=a(\nu-b\kappa),\quad
 c_{\mathcal L,C}=\frac b2\{a\kappa(b\kappa-\nu)-1\}
                    =1+\kappa\pmod2.
\hypertarget{tag-Q23Z}{}\label{maps:eq:MT_local_63}
\end{equation}
The degree is odd. Here \(b/2\) and \(\nu\) are odd, as seen in every
even-sector row.

At a rank-one site the short refinement is
\(\epsilon_{N,j}=\mathcal Q(SC_Se_j)\); use the Gaussian with
\(\epsilon_N^tZ_N\) and support
\((n_{\rm per}j^2-\mathcal Q(l)j)/2\).
The physical first content \(d_0\) is odd in these sectors
(one or five), so \(-n_{\rm per}j(j-d_0)/2\) is integral.
The difference of the two supports is the integral linear form
\((a\,\operatorname{row}_1C_S-\epsilon_N^t)/2\).
Its monomial and reciprocal give the same-line comparison.
Squaring \eqref{maps:eq:MT_local_59} gives exactly \eqref{maps:eq:MT_local_48}. The supports of the tensor squares differ by
the linear refinement term; the opposite monomial identifies their
regular coordinates. This gives
\(\mathcal L^{\otimes2}\simeq\mathcal D\), including finite actions.

At \(I_2\) use the explicit primitive basis
\begin{equation}
 l=2u-a\kappa_i\delta,\quad
 p_2=u+(1-a\kappa_i)\delta/2,\quad
 h_2=v,\quad w_0=\gamma-(b\kappa_i/2)v.
\hypertarget{tag-Q240}{}\label{maps:eq:MT_local_64}
\end{equation}
Its short matrix is
\(\left(\begin{smallmatrix}2&1\\-a\kappa_i&(1-a\kappa_i)/2\end{smallmatrix}\right)\),
inverse
\(\left(\begin{smallmatrix}(1-a\kappa_i)/2&-1\\a\kappa_i&2\end{smallmatrix}\right)\);
the long matrix and inverse are
\(\left(\begin{smallmatrix}1&-b\kappa_i/2\\0&1\end{smallmatrix}\right)\)
and
\(\left(\begin{smallmatrix}1&b\kappa_i/2\\0&1\end{smallmatrix}\right)\).
Compared with \eqref{maps:eq:MT_local_36}, \(p_0=p_2+(\sigma_{\rm Bez}-1)l/2\) and
\(h_{\rm per}=h_2+y_Lw_0\); subtracting these multiples is the inverse.
The pairing in \eqref{maps:eq:MT_local_64} is
\(\left(\begin{smallmatrix}2a&0\\a&-b/2\end{smallmatrix}\right)\).
Thus
\begin{equation}
                  \deg\mathcal L|E_2=-b/2,\qquad c_{\mathcal L,E_2}=0.
\hypertarget{tag-Q241}{}\label{maps:eq:MT_local_65}
\end{equation}
Indeed \(\delta=-l+2p_2\), \(\mathcal Q(p_2)=1\), and the nodal
slope is \(a\); its half-period character is \(1-a=0\pmod2\).
This completes the equivariant-line construction on \(Y\).

\item \emph{The common global Borel comparison.}

The cusp kills both columns \(SC\), equivalently \(u,\delta\) through
\eqref{maps:eq:MT_local_42}. Each rank-one cap kills its primitive \(l\).
Let \(K_1\) collect these angular columns in the first three coordinates.
Integer reduction of these angular columns gives the section
\begin{equation}
                              K_1R_1=I_3.
\hypertarget{tag-Q242}{}\label{maps:eq:MT_local_66}
\end{equation}
Thus \(e_1,e_2,\delta\) vanish in the fundamental group:
the fibre starts abelian, and these are cap-disc relations.
The remaining period relations make \(\gamma\) commute with every
unfilled base meridian. With zero or one additive meridian,
\begin{equation}
                         \pi_1(Y)=\mathbb Z\gamma\times F_{r_Y},
                         \qquad r_Y=0\text{ or }1.
\hypertarget{tag-Q243}{}\label{maps:eq:MT_local_67}
\end{equation}
The inverse presentation kills the three angular generators, fixes
\(\gamma\) and the additive meridian, and replaces the ordinary
meridian by their inverse product. Equation \eqref{maps:eq:MT_local_66} supplies exact
words for the killed generators. In particular \(H_1(Y)\) is free.

Write \(\alpha_p(s)=1\). Every section sweep has zero image in
\(H_2(X;\mathbb F_p)\) for \(p\mid c\): the toric/original sweeps
bound; at a finite site \eqref{maps:eq:MT_local_33}--\eqref{maps:eq:MT_local_34} gives \(L_\ell=0,D_\delta=-be\,s=0\) modulo
\(p\). This includes the gamma-two original III lift. By the cycle
coordinates \eqref{maps:eq:MT_local_18} and the torsion-free \(H_1(Y)\),
\begin{equation}
 \alpha_p|Y=\tfrac12\rho_p c_1(\mathcal D)\quad(p=3,5),\qquad
 \alpha_2|Y=\rho_2c_1(\mathcal L)\quad(2\mid c).
\hypertarget{tag-Q244}{}\label{maps:eq:MT_local_68}
\end{equation}
The inverses of the factors two are the usual B\'ezout inverses modulo
the odd primes.

The connected scalar circle acts trivially on ordinary cohomology.
The possible outgoing differentials from \(\alpha_p\) in its Borel
spectral sequence have targets
\(H^2(BS^1;H^1(X;\mathbb F_p))=0\) and \(H^3(BS^1;\mathbb F_p)=0\).
Thus each \(\alpha_p\) extends to \(U_p\in H^2_{S^1}(X;\mathbb F_p)\).
Choose a toric endpoint as basepoint; the circle fixes it. The based
isotopy and that section give
\begin{equation}
                 \pi_1(Y_{C_p})=(\mathbb Z\times F_{r_Y})\times C_p.
\hypertarget{tag-Q245}{}\label{maps:eq:MT_local_69}
\end{equation}
A section from any finite fixed point sends the generator of \(C_p\)
to the same finite factor: its component in the torsion-free first
factor must vanish.

The difference between \(U_p|Y_{C_p}\) and the line class in \eqref{maps:eq:MT_local_68}
has zero ordinary restriction. In the universal-cover low-degree
sequence it lies in group cohomology; its ordinary group component
vanishes by the injective low-degree map for \(Y\). The remaining
terms are a base degree-two class and products of a degree-one group
class with the integral degree-one classes of \(Y\).
All these integral classes vanish on each finite isotropy section
by \eqref{maps:eq:MT_local_69}. Consequently the base coefficient is one common number on
all fixed curves. Put
\(T_{\rm Bor}=c_1(\text{positive weight-one character})\in H^2(BS^1;\mathbb Z)\).
Subtract that multiple of the whole-circle base class \(T_{\rm Bor}\);
adding it back is the inverse. The degree-one mixed terms restrict to the products displayed below.

For odd \(p\), after this one global adjustment,
\begin{equation}
 U_p|C_f=\alpha_p|C_f,\qquad
 U_p|E_p=\alpha_p|E_p+\eta_p\xi_p,
\hypertarget{tag-Q246}{}\label{maps:eq:MT_local_70}
\end{equation}
where \(\eta_p\in H^1(BC_p;\mathbb F_p)\) and
\(\xi_p\in H^1(E_p;\mathbb F_p)\) lifts integrally.
For \(p=2\), writing \(\tau\in H^1(BC_2;\mathbb F_2)\),
\begin{equation}
 U_2|C_f=\alpha_2|C_f+(1+\kappa)\tau^2,\qquad
 U_2|E_2=\alpha_2|E_2+\tau\xi_2.
\hypertarget{tag-Q247}{}\label{maps:eq:MT_local_71}
\end{equation}
There is no mixed term on the sphere. Equation \eqref{maps:eq:MT_local_70} is subsequently
refined on the full-circle Borel space of \(E_3\), so its order-nine
contribution is also controlled.

\item \emph{Full coprime homotopy, detector, and bordism maps.}

Put \(K_p=K(\mathbb Z/p,2)\). The cyclic resolution makes
\(K(\mathbb Z/p,1)\) acyclic in positive homology over
\(\mathbb Z[1/p]\). The path-loop Serre sequence, with contractible
total space and simply connected base \(K_p\), then makes \(K_p\)
acyclic in positive homology over that ring. Exactness of localization
implies that each positive integral homology element is killed by a
power of \(p\). For different primes \(p,\varpi\), B\'ezout between the two
annihilating powers kills the tensor product. For a positive integer
\(n_{\rm Tor}\) and a \(\varpi\)-primary group \(B_{\varpi}\),
\(\operatorname{Tor}(\mathbb Z/p^{n_{\rm Tor}},B_{\varpi})
 =B_{\varpi}[p^{n_{\rm Tor}}]=0\).
Direct sums and filtered colimits prove the corresponding
Tor vanishing in every degree.

The reduced Kunneth sequence therefore makes \(K_p\wedge K_{\varpi}\)
acyclic. The CW inclusion
\begin{equation}
                         i:K_p\vee K_{\varpi}\longrightarrow K_p\times K_{\varpi}
\hypertarget{tag-Q24P}{}\label{maps:eq:MT_local_94}
\end{equation}
is an integral homology equivalence of simply connected spaces.
A first nonzero relative homotopy group would, by relative Hurewicz,
be the corresponding zero relative homology group. All relative
homotopy groups thus vanish. One obtains a based inverse \(r_{\rm CW}\) on cells:
start with the identity on the wedge, extend its retraction over each
attaching cell using the zero relative groups, and adjust the extension
so that \(i r_{\rm CW}\) is homotopic to the original cell map relative boundary.
The homotopy-extension property assembles the maps on skeleta.
Thus \(r_{\rm CW}i=\mathrm{id}\) and \(i r_{\rm CW}\simeq\mathrm{id}\) relative to the
wedge. This is the inverse at its full CW homotopy scope.
For three primes apply the same construction successively; the same
tensor and Tor argument applies to the product of the first two.

On coefficient groups, let
\begin{equation}
 \Theta_{\rm CRT}:\mathbb Z/c\to\bigoplus_{p\mid c}\mathbb Z/p,\quad
 x\mapsto(x\bmod p)_p,\qquad
 j_p(e_{p,\rm coef})=\frac cp\left(\frac cp\right)^{-1}_{p}s.
\hypertarget{tag-Q24Q}{}\label{maps:eq:MT_local_95}
\end{equation}
The inverse of \(\Theta_{\rm CRT}\) is \(\sum j_p\). The two composites are
identity by the residue calculation. Apply the Eilenberg--Mac Lane
functor to these degree-two coefficient complexes; in a simplicial
abelian-group model its composites are literally identity.
Composing with \eqref{maps:eq:MT_local_94} and its inverse gives the corresponding full
wedge/product homotopy maps.

The finite wedge axiom for the spin-bordism spectrum gives
\begin{equation}
 \Omega_6^{\rm Spin}(K(\mathbb Z/c,2))
             \simeq\bigoplus_{p\mid c}\Omega_6^{\rm Spin}(K_p).
\hypertarget{tag-Q24R}{}\label{maps:eq:MT_local_96}
\end{equation}
Write \(h_{\rm bord}([M,f])=f_*[M]\) for the fundamental-class map
in the cited spin-bordism pair. Here $\Omega_6^{\rm Spin}(\mathrm{pt})=0$, by
\ref{maps:src:D.1.11} for coefficient group zero, whose injective
pair $(P,h_{\rm bord})$ has zero target. The forward map
is the tuple of coefficient projections. Its inverse is the sum
of the factor inclusions followed by \(K(\Theta_{\rm CRT}^{-1},2)_*\).
These maps identify the complete spin-bordism group with the
displayed direct sum.

Let \(D_2=2,D_3'=9,D_5=5\). The same maps on degree-six homology
and the positive generators give
\begin{equation}
 H_6(K(\mathbb Z/c,2);\mathbb Z)=
 \mathbb Z/(3c)\langle\gamma_3(s)\rangle.
\hypertarget{tag-Q24S}{}\label{maps:eq:MT_local_97}
\end{equation}
Naturality from positive \(\mathbb{CP}^3\) reduction gives
\(\gamma_3(ms)=m^3\gamma_3(s)\). For the subgroup injection \eqref{maps:eq:MT_local_95},
the cube of its coefficient is an idempotent modulo \(3c\).
In the three cases the inverse detector maps are exactly
\begin{equation}
 \begin{array}{c|c|c}
 c&\text{forward coordinates}&\text{inverse coefficient modulo }3c\\\hline
 6&(\mu_2,\mu_3)\in\mathbb Z/2\oplus\mathbb Z/9&
                    9\mu_2+10\mu_3\pmod{18}\\
 15&(\mu_3,\mu_5)\in\mathbb Z/9\oplus\mathbb Z/5&
                    10\mu_3+36\mu_5\pmod{45}\\
 30&(\mu_2,\mu_3,\mu_5)&
                    45\mu_2+10\mu_3+36\mu_5\pmod{90}.
 \end{array}
\hypertarget{tag-Q24T}{}\label{maps:eq:MT_local_98}
\end{equation}
For example, at \(c=6\) the coefficient injections are
\(j_2=3\), \(j_3=4\); their cubes are \(9,10\) modulo eighteen.
At \(c=15\), \(j_3=10,j_5=6\), giving cubes \(10,36\) modulo
forty-five. At \(c=30\), \(j_2=15,j_3=10,j_5=6\), giving
\(45,10,36\) modulo ninety. Reductions to the listed detector
orders are the inverse maps: each displayed coefficient has residue
one in its own factor and zero in the other factors. The computed primary evaluations will determine \eqref{maps:eq:MT_local_3} through these maps.

For \(P\), the inverse coefficient maps are
\begin{equation}
 \begin{array}{c|c}
 c&\Theta_{\rm CRT}^{-1}\\\hline
 6&3P_2+4P_3\pmod6\\
 15&10P_3+6P_5\pmod{15}\\
 30&15P_2+10P_3+6P_5\pmod{30}.
 \end{array}
\hypertarget{tag-Q24U}{}\label{maps:eq:MT_local_99}
\end{equation}
These are the inverse maps for the refined primary coefficients.
An automorphism \(s\mapsto us\) acts by \(u\) on \(P\) and \(u^3\)
on \eqref{maps:eq:MT_local_97}. The latter is independent of the integer lift of \(u\):
\((u+c)^3-u^3\) is divisible by \(3c\). If \(uv=1\pmod c\),
then \((uv)^3=1\pmod{3c}\), since \(3\mid c\).
Thus the inverse automorphism is multiplication by \(v\) on \(P\)
and $v^3$ on $\mu$.


\item \emph{The full order-nine detector and localization.}

Use the coefficient map $e_{\rm EM}\mapsto e_{3,\rm coef}$ on $\mathbb Z/3$,
with inverse $e_{3,\rm coef}\mapsto e_{\rm EM}$. The induced Eilenberg--Mac Lane map
identifies the positive reduction generator in
\eqref{maps:eq:three-primary-homology} with $\gamma_3(e_{3,\rm coef})$.
Under this same map the detector is
\eqref{maps:eq:third-power-detector}, with coefficient map
$\gamma_3(e_{3,\rm coef})\mapsto1$ and inverse $j\mapsto j\gamma_3(e_{3,\rm coef})$.
The coefficient-one cochain construction is
\eqref{maps:eq:third-power-cochain}, and the left-interval suspension
and its slant inverse are those of
\eqref{maps:eq:third-power-suspension}. These are statements about
the same universal target and natural operation. Applying them to the
classifying map for $U_3$ therefore gives a class in $\mathbb Z/9$. The fixed curves and their restrictions for the
present data are computed below.

Restrict the whole-circle \(U_3\) of the common global Borel comparison to \(C_9\), obtaining
\(U_9\).
For the chosen positive generator, let
\(T_9=c_1(\text{positive character})\) modulo nine.
Write \(g_9\) for that positive generator of \(C_9\).
The cyclic resolution has differentials \(g_9-1,N_9\), where
\(N_9=\sum_{j=0}^{8}g_9^j\); with trivial
\(\mathbb Z/9\) coefficients they are zero. Its period-two shift gives
\begin{equation}
 H^*(BC_9;\mathbb Z/9)
       =(\mathbb Z/9)[T_9]\otimes\Lambda(\eta_9).
\hypertarget{tag-Q248}{}\label{maps:eq:MT_local_76}
\end{equation}
In particular the constants inject after \(T_9\) is inverted.

Use invariant tubular neighborhoods of \(C_f,E_3\). Their sphere
boundaries and the compact complement are free for \(C_9\), by
the prescribed-cusp construction. Their Borel spaces have the homotopy type of the
finite-dimensional quotients, so their cohomology is killed by
inverting \(T_9\). Thom excision and self-intersection give the 
inverse restriction map
\begin{equation}
            (z_C,z_E)\longmapsto
             i_{C*}(z_C/e_C)+i_{E*}(z_E/e_E).
\hypertarget{tag-Q249}{}\label{maps:eq:MT_local_77}
\end{equation}
Here \(i_C:C_f\hookrightarrow X\) and \(i_E:E_3\hookrightarrow X\)
are the fixed-curve inclusions, and \(e_C,e_E\) are their equivariant
normal Euler classes. Pullback verifies the inverse through
\[
 i_C^*i_{C*}(z_C)=e_Cz_C,\qquad
 i_E^*i_{E*}(z_E)=e_Ez_E.
\]
The remaining relative class vanishes after localization. This works for the
invariant curves with their \(C_9\) actions, even when they
are rotated or translated.

The equivariant integral of \(Q_3(U_9)\) is a degree-zero constant,
namely \(\langle Q_3(\alpha_3),[X]\rangle\in\mathbb Z/9\).
It therefore suffices to compute the two terms of \eqref{maps:eq:MT_local_77}.

On \(C_f\), set \({a_{\rm edge}}=a(\nu-k_{\rm edge})\), \(k_{\rm edge}=b\kappa\), and let
\(z_C=c_1^{S^1}(\mathcal O(1))\) with north character zero.
The two affine coordinates in \eqref{maps:eq:MT_local_45} identify its action with homogeneous
weights \(0,-k_{\rm edge}\), so
\begin{equation}
                 z_C(z_C-k_{\rm edge}T_{\rm Bor})=0,\qquad\int_{C_f}z_C=1.
\hypertarget{tag-Q24A}{}\label{maps:eq:MT_local_78}
\end{equation}
The south character is \(k_{\rm edge}\). After the one common adjustment in
the common global Borel comparison the north character of \(U_3\) is zero, and its degree is
\({a_{\rm edge}}\), so \(U_3|C_f=\rho_3({a_{\rm edge}}z_C)\) on the whole-circle Borel space.
Its two normal roots, by \eqref{maps:eq:MT_local_46}, are
\((k_{\rm edge}-\nu)T_{\rm Bor}-z_C\) and \(\nu T_{\rm Bor}-z_C\). Their product is
\begin{equation}
                             e_C=\nu(k_{\rm edge}-\nu)T_{\rm Bor}^2.
\hypertarget{tag-Q24B}{}\label{maps:eq:MT_local_79}
\end{equation}
The coefficient \(\nu(k_{\rm edge}-\nu)\) is a unit modulo nine. The integral
lift formula \eqref{maps:eq:third-power-detector} therefore gives
\begin{equation}
 \int_{C_f}\frac{Q_3(U_9|C_f)}{e_C}
       =\frac{{a_{\rm edge}}^3k_{\rm edge}^2}{\nu(k_{\rm edge}-\nu)}=0\quad\text{in }\mathbb Z/9,
\hypertarget{tag-Q24C}{}\label{maps:eq:MT_local_80}
\end{equation}
because $3\mid k_{\rm edge}$. The same formula holds at $k_{\rm edge}=0$.

At \(E_3\), the primitive periods in the primitive rank-one construction give
\begin{equation}
 E_3=(\mathbb R/\mathbb Z)^2,\qquad
 \theta\cdot(z,w)=(e^{2\pi i3\theta}z,w),
\hypertarget{tag-Q24D}{}\label{maps:eq:MT_local_81}
\end{equation}
with positive first and second directions \(p_0,w_0\).
There is an explicit Borel identification
\begin{equation}
 ES^1\times_{S^1}S^1_{(3)}\longrightarrow ES^1/C_3,\quad
 [e,z]\longmapsto[e r_{\rm root}],\quad r_{\rm root}^3=z.
\hypertarget{tag-Q24E}{}\label{maps:eq:MT_local_82}
\end{equation}
Changing the root changes \(e r_{\rm root}\) by \(C_3\). The associated-bundle
relation \([et,z]=[e,t^3z]\) proves it is well-defined.
The inverse is \([e]\mapsto[e,1]\); both composites agree. Thus
\((E_3)_{S^1}\simeq BC_3\times S^1\). A positive \(p_0\) turn
lifts to a root path ending at the chosen positive generator of \(C_3\),
while the second-period map is the identity on $w_0$. If \(\eta\) is its degree-one generator
and \(\xi\) the positive integral class of the second circle, then
\begin{equation}
 H^2((E_3)_{S^1};\mathbb F_3)
          =\mathbb F_3t_3\oplus\mathbb F_3\eta\rho_3\xi.
\hypertarget{tag-Q24F}{}\label{maps:eq:MT_local_83}
\end{equation}
Here $t_3=\rho_3(T_{\rm Bor})$ is the restricted positive circle character.
The common zero isotropy character removes the first coefficient.
Ordinary pullback of \(\eta\rho_3\xi\) is the positive area class by
the exact generator map just described. The line degree gives
\begin{equation}
                  U_3|E_3=d_E\eta\rho_3\xi,\qquad d_E=-b/3\pmod3.
\hypertarget{tag-Q24G}{}\label{maps:eq:MT_local_84}
\end{equation}
Let \(q_{\rm susp}:BC_3\times S^1\to S^1\wedge BC_3\) flip factors and then
collapse the wedge. Write \(\Sigma_{\rm susp}\) for the left-interval suspension of
\eqref{maps:eq:third-power-suspension}, with its stated slant inverse. Then
\(q_{\rm susp}^*\Sigma_{\rm susp}\eta=\rho_3\xi\,\eta=-\eta\rho_3\xi\).
Thus \eqref{maps:eq:MT_local_84} is \(q_{\rm susp}^*\Sigma_{\rm susp}(-d_E\eta)\).
Naturality and \eqref{maps:eq:third-power-suspension} show \(Q_3(U_3|E_3)=0\), and hence its pullback
\(Q_3(U_9|E_3)=0\). This kills the entire order-nine class, including
the subgroup invisible to ordinary cubing modulo three.

Combining \eqref{maps:eq:MT_local_76}--\eqref{maps:eq:MT_local_84} proves
\begin{equation}
                     \langle Q_3(\alpha_3),[X]\rangle=0
                             \quad\text{in }\mathbb Z/9.
\hypertarget{tag-Q24H}{}\label{maps:eq:MT_local_85}
\end{equation}
The global comparison \eqref{maps:eq:MT_local_66}--\eqref{maps:eq:MT_local_71}
and the finite cap maps \eqref{maps:eq:MT_local_25}--\eqref{maps:eq:MT_local_34}
hold on all fourteen patterns for every allowed multiplicity and affine lift.

\item \emph{The other primary cubics.}

For \(C_2\), the normal representation on each fixed curve \(F_{\rm fix}\) is complex
rank two with sign characters. Its equivariant Euler class is
\begin{equation}
 e_{F_{\rm fix}}=\tau^4+\tau^2\rho_2c_1(\nu_{F_{\rm fix}})+\rho_2c_2(\nu_{F_{\rm fix}})=\tau^4.
\hypertarget{tag-Q24I}{}\label{maps:eq:MT_local_86}
\end{equation}
Indeed \(c_1(\nu_{F_{\rm fix}})=-c_1(TF_{\rm fix})\) is even and a curve has no degree-four
cohomology. Its normal sphere bundles and complement are free; the
same Thom restriction/inverse argument, after inverting \(\tau\),
applies to \eqref{maps:eq:MT_local_71}.
On \(E_2\), the integral lift of \(\xi_2\) has square zero, and curve
classes of degree greater than two vanish. Its cubic contribution is
therefore zero. On \(C_f\), put \(\alpha_{2,C}=\alpha_2|C_f\). The only
integrating term of \((\alpha_{2,C}+t\tau^2)^3/\tau^4\) is \(t\alpha_{2,C}\).
The degree in \eqref{maps:eq:MT_local_63} is odd. Thus
\begin{equation}
                         \langle\alpha_2^3,[X]_2\rangle
                           =t=1+\kappa\pmod2.
\hypertarget{tag-Q24J}{}\label{maps:eq:MT_local_87}
\end{equation}
The removed common character changes this total by its value times
the sum of the two odd ordinary degrees, which is zero modulo two.

For \(C_5\), the normal characters on each fixed curve are opposite
units. Their Euler class is a unit after the base generator is inverted.
Both restrictions in \eqref{maps:eq:MT_local_70} have cube zero: the ordinary curve term
squares to zero, and the mixed product has square zero because
\(\eta_5^2=0\) and \(\xi_5\) has an integral lift. Hence
\begin{equation}
                           \langle\alpha_5^3,[X]_5\rangle=0.
\hypertarget{tag-Q24K}{}\label{maps:eq:MT_local_88}
\end{equation}
The order-two projection in \eqref{maps:eq:MP_local_61} and the
order-five calculation of \ref{maps:src:D.1.10} give
\(H_6(K(\mathbb Z/2,2);\mathbb Z)=\mathbb Z/2\), detected by the
mod-two cube, and \(H_6(K(\mathbb Z/5,2);\mathbb Z)=\mathbb Z/5\),
detected by the mod-five cube. Both generators are reduced positive
\(\mathbb{CP}^3\). Thus \eqref{maps:eq:MT_local_85},\eqref{maps:eq:MT_local_87},\eqref{maps:eq:MT_local_88} are full primary coordinates.


\item \emph{The complete three-primary bordism class.}

Compose the classifying map $X\to K(\mathbb Z/c,2)$
with reduction $s\mapsto e_{3,\rm coef}$. The target group and maps are exactly
\eqref{maps:eq:three-primary-bordism}. Write \(\zeta_{3,\rm bord}\)
for the image of its positive \(\mathbb{CP}^3\) reduction generator
\(\zeta_{\rm bord}\) under the coefficient map \(e_{\rm EM}\mapsto e_{3,\rm coef}\).
The inverse fundamental-class map is
$j\gamma_3(e_{3,\rm coef})\mapsto j\zeta_{3,\rm bord}$. This applies to the resulting
bordism class of this map because the calculation is a target-group
isomorphism. The full zero coordinate just computed forces $j=0$,
hence also $P=0$. The resulting three-primary bordism class is the
class of the $M_{3,0}$ of
Definition~\ref{maps:def:minimal-three-data}.

\item \emph{Refined Pontryagin coordinates and exact weight conversion.}

For \(p=3,5\), let the ordinary roots of the two normal eigenlines at
a fixed curve \(F_{\rm fix}\) be \(x_1,x_2\), and their characters be \(\lambda_{\rm norm},-\lambda_{\rm norm}\).
Put $t_p=\rho_p(T_{\rm Bor})$ for the restricted positive circle character.
The tangent root and all products of ordinary degree exceeding two
vanish on the curve. Expanding the two normal roots gives
\begin{equation}
                         p_1(TX)|F_{\rm fix}=-2e_{F_{\rm fix}}.
\hypertarget{tag-Q24L}{}\label{maps:eq:MT_local_89}
\end{equation}
Both sides equal
\(2\lambda_{\rm norm}^2t_p^2+2\lambda_{\rm norm} t_p(x_1-x_2)\).
This identity holds for both ordinary normal-degree pairs
$(x_1,x_2)$. Localization gives
\begin{equation}
 \langle\alpha_p\rho_pp_1,[X]_p\rangle
          =-2\{a(\nu-b\kappa)-b/p\}
          =2(b/p-a\nu)\pmod p.
\hypertarget{tag-Q24M}{}\label{maps:eq:MT_local_90}
\end{equation}
The mixed degree-one term integrates to zero. The primary definition
of \(P\) on groups without two-torsion is the quarter of the pushed
forward Poincar\'e dual of \(p_1\). Its identification with the intrinsic
class is proved in \cite[Section~3.13, p.~807]{manZhubr75}. Multiplication by four has its
B\'ezout inverse modulo \(p\). By naturality under the coefficient
projection \(\mathbb Z/c\to\mathbb Z/p\), \eqref{maps:eq:MT_local_90} gives \eqref{maps:eq:MT_local_2}.
The quarter-map is defined for every odd-order coefficient group,
including $\mathbb Z/3$.

Here is the complete conversion to the Tate coordinates.
At \(E_p\), \(r_{\rm Tate}=-\kappa_i\) and \(w_{\rm norm}=(a r_{\rm Tate})^{-1}=-\sigma_{\rm Bez}\).
The ordinary degrees in \eqref{maps:eq:MT_local_70} and the degree-zero equivariant
integration give
\begin{equation}
       \frac{a\nu}{\nu^2}-\frac{b/p}{w_{\rm norm}^2}=0,\qquad
       w_{\rm norm}^2=\frac{\nu(b/p)}a,\qquad
       P_p=\frac{b/p}{2}\left(1-\frac{\nu^2}{w_{\rm norm}^2}\right).
\hypertarget{tag-Q24N}{}\label{maps:eq:MT_local_91}
\end{equation}
All denominators are units on these rows. The physical-weight computation
gives the middle equality from the marked $\kappa_i$ and recovers
\(P_p\) by the last expression and the marking below, independently
of the first expression in \eqref{maps:eq:MT_local_2}.

The subgroup automorphism \(z\mapsto z^{\nu^{-1}}\), inverted by
\(z\mapsto z^\nu\), changes both normal pairs simultaneously:
the cusp pair becomes \((-1,1)\), and the elliptic pair becomes
\((w_{\rm norm}/\nu,-w_{\rm norm}/\nu)\). Its pullback multiplies the Borel character and
degree-one generator by the same exponent. If
\(\alpha_C=(a\nu)^{-1}\alpha_p\), its dual homology generator is
\(g_C=(a\nu)s\), with inverse coordinate evaluation by \(\alpha_C\).
Then
\begin{equation}
 \alpha_C[E_p]=-\frac{w_{\rm norm}^2}{\nu^2},\qquad
 P_C=\frac{w_{\rm norm}^2/\nu^2-1}{2},\qquad P_p=(a\nu)P_C.
 \hypertarget{tag-020L}{}\label{maps:eq:D.MT3}
\end{equation}
Multiplication by \((a\nu)^{-1}\) is the inverse rescaling.
The normalizing action and the homology rescaling are distinct
displayed maps with their displayed inverses.

The physical weights in the fourteen patterns are:

\begin{center}\begin{tabular}{l|l|l}
rows & \(p=3\): \(w_{\rm norm},\nu\bmod3,w_{\rm norm}/\nu\) & \(p=5\): \(w_{\rm norm},\nu\bmod5,w_{\rm norm}/\nu\)\\
\hline
1,2,3 & \(2,2,1\) & not present\\
4,5 & \(1,1,1\) & not present\\
6,7 & \(2,2,1\) & \(4,2,2\)\\
8 & \(1,1,1\) & \(4,1,4\)\\
9 & \(2,1,2\) & \(2,1,2\)\\
10 & \(2,2,1\) & \(4,2,2\)\\
11,12 & \(1,1,1\) & \(2,1,2\)\\
13,14 & \(2,1,2\) & \(4,1,4\)\\
\end{tabular}\end{center}

Thus \(P_3=0\) on all fourteen rows. At \(p=5\), the six nonzero
rows have primitive coefficient \(3\) and curve coefficient \(4\);
row 8 and both \(c=30\) rows have zero coefficient. In particular
the criterion is \(w_{\rm norm}^2=\nu^2\) in one common action marking.

For the two-primary projection use Zhubr's exact quotient
\cite[Sections~3.7 and~3.12, pp.~803,807]{manZhubr75}:
\begin{equation}
\begin{gathered}
 \Omega_6^{\rm Spin}(K(\mathbb Z/2,2))
       =\mathbb Z\langle\epsilon_{2,\rm bord},\zeta_{2,\rm bord}\rangle/
                        \langle\epsilon_{2,\rm bord},2\zeta_{2,\rm bord}\rangle,\\
 P(\zeta_{2,\rm bord})=1,\quad h_{\rm bord}(\zeta_{2,\rm bord})=\gamma_3(e_{2,\rm coef}).
\end{gathered}
\hypertarget{tag-Q24O}{}\label{maps:eq:MT_local_93}
\end{equation}
Its forward map is \((a_{\rm bord},b_{\rm bord})\mapsto b_{\rm bord}\bmod2\), with inverse the
reduction of \(b_{\rm bord}\zeta_{2,\rm bord}\). The generator is positive spin
\(\mathbb{CP}^3\) reduction and has cube one. Hence \eqref{maps:eq:MT_local_87} gives
the mod-two projection of \(P(X)\), namely \(t\).
The exact CRT maps give $P_6=3t$, $P_{15}=3\epsilon$ and
$P_{30}=15t$. Substitution of \eqref{maps:eq:MT_local_85},\eqref{maps:eq:MT_local_87},\eqref{maps:eq:MT_local_88} into \eqref{maps:eq:MT_local_98} gives
the full fundamental-class coefficients in \eqref{maps:eq:MT_local_3}.

\Needspace{5\baselineskip}
\item \emph{Marked connected sums, recognition, and the refined count.}

The models in Definitions~\ref{maps:def:minimal-two-data},
\ref{maps:def:minimal-three-data} and
\ref{maps:def:minimal-mixed-data} have their original orientations,
normal framings and primitive generators. The degree-two tube,
degree-three graph and positive conic are the identical embeddings
used in their defining proofs. Their generator identifications are
the identity. Those proofs give
$(P,h_{\rm bord})(M_t)=(te_{2,\rm coef},t\gamma_3(e_{2,\rm coef}))$, $(P,h_{\rm bord})(M_{3,0})=(0,0)$,
and $(P,h_{\rm bord})(Z_{5,\epsilon})=(\epsilon e_{5,\rm coef},0)$, with simple
connectivity, spin and zero free third homology. We transport
these complete marked classes through the connected sums.

Puncture the model manifolds in oriented discs on which their
classifying maps are constant. Attach the oriented neck, and extend
the maps constantly over the one-handle bordism from the disjoint
union. Its boundary classes add and \(P\) is a bordism homomorphism.
The homology inclusions and the summand collapse maps give inverse
coordinates. The coprime maps \eqref{maps:eq:MT_local_94}--\eqref{maps:eq:MT_local_96} prove that this addition
determines every spin-bordism coordinate.

The precise marking \(\Phi\) from each model's \(H_2\) to the 
\(\mathbb Z/c\langle s\rangle\), together with its inverse, is
\begin{equation}
 \begin{array}{c|c|c}
 c&\Phi&\Phi^{-1}(xs)\\\hline
 6&(e_{2,\rm coef},e_{3,\rm coef})\mapsto(3s,2s)&(x\bmod2,\ 2x\bmod3)\\
 15&(e_{3,\rm coef},e_{5,\rm coef})\mapsto(5s,3s)&(2x\bmod3,\ 2x\bmod5)\\
 30&(e_{2,\rm coef},e_{3,\rm coef},e_{5,\rm coef})\mapsto(15s,10s,6s)&
                         (x\bmod2,\ x\bmod3,\ x\bmod5).
 \end{array}
\hypertarget{tag-Q24V}{}\label{maps:eq:MT_local_100}
\end{equation}
Substitution in each row gives $\Phi^{-1}\Phi=1$ and
$\Phi\Phi^{-1}=1$ on every coefficient.
These markings send the model Pontryagin classes to \eqref{maps:eq:MT_local_99}.
Their degree-six maps send the model's sole nonzero fundamental
class \(t\gamma_3(e_{2,\rm coef})\) to \(27t=9t\pmod{18}\) at \(c=6\),
and \(3375t=45t\pmod{90}\) at \(c=30\).
At \(c=15\) all fundamental-class coordinates are zero.
Thus both full marked invariants agree with \eqref{maps:eq:MT_local_1}--\eqref{maps:eq:MT_local_3}.

The recognition scope is that of Theorem~5 and the marked
construction and uniqueness proof in Section~4.1
\cite[pp.~795--796,809--810]{manZhubr75}. Zhubr Theorem 5, printed p.796,
gives the injection
\begin{equation}
 (P,h_{\rm bord}):\Omega_6^{\rm Spin}(K(G_{\rm coeff},2))
                  \longrightarrow G_{\rm coeff}\oplus H_6(K(G_{\rm coeff},2);\mathbb Z).
\hypertarget{tag-Q24W}{}\label{maps:eq:MT_local_101}
\end{equation}
Theorem 3, printed p.795, with the construction and marked uniqueness
proof of §4.1, printed p.809, constructs the representative by killing
the free part of \(H_3\) while preserving its marking $H_2\simeq G_{\rm coeff}$.
Equality of marked bordism classes then yields a degree-positive
diffeomorphism for the simply connected spin representatives with
zero free \(H_3\). The proof changes the marked bordism into an
\(h\)-cobordism.

Both the constructed threefolds and their models satisfy that scope,
with unique spin structures. Equations \eqref{maps:eq:MT_local_98}--\eqref{maps:eq:MT_local_100} and injection
\eqref{maps:eq:MT_local_101} give equality of their marked bordism classes. The marked
uniqueness gives an oriented diffeomorphism with induced second
homology map exactly \(\Phi\), and its inverse induces the inverse
in \eqref{maps:eq:MT_local_100}. This proves \eqref{maps:eq:MT_local_4}.

All listed types are realized: choose \(e=1\) without a finite cap.
At a finite cap take \(m=d_i,e=1,n=m-1,r=0\) and the affine vector
\(v_{\rm fin}/m\). It is free, satisfies \(me-n=1\), and is an allowed 
construction. The cap, line and localization constructions prove the same invariant formulas for
every other allowed finite parameter and affine lift.

\Needspace{10\baselineskip}
Within an even sector, every marking automorphism is a unit, hence
is one modulo two and preserves \(t\). Its inverse is the inverse
unit. Thus the two values of the mod-two cube distinguish the two
types. At \(c=15\), a unit preserves zero versus nonzero \(P\);
the two orders are one and five. Conversely, the explicit models
and maps \eqref{maps:eq:MT_local_100} identify every row with the same listed value.
Different values of \(c\) have nonisomorphic \(H_2\), as their
orders differ, so all six sector types are distinct. This proves
the counts in the theorem.
\end{enumerate}
\end{proof}


\Needspace{12\baselineskip}
\subsection{Residual orders four, eight and twenty}

\subsubsection{Semistable orders four and eight}

\begin{definition}[Semistable residual data]\hypertarget{tag-Q24X}{}\label{maps:def:semistable-residual-data}
{\normalfont\scriptsize tag Q24X}\par
Use the ordered elliptic data of
Definition~\ref{mapsIII:def:elliptic-marked-data}. Choose OS55 with
$(a,b)=(5,4)$ and selected position $i_0\in\{2,3,4\}$, or OS45
with $(a,b)=(1,8)$ and $i_0\in\{1,2,3,4\}$. Both profiles are
the base semistable profiles. The selected position is $i_0=i-1$
in the input definition's peripheral order.
Use the prescribed primitive Mumford cap at every multiplicative
site and ordinary clutch $e\gamma$, $e=\pm1$. These are all the
fillings in this definition. Put $c=b$.

Define two oriented spin surgeries, using the unique homotopy class
of normal framing of each specified embedded sphere.
For $Z_{c,0}$, tube $c$ disjoint parallel copies of $S^2$ in
$S^2\times S^4$ and perform surgery on that sphere. For $Z_{c,1}$,
put the conic of $M_{\rm conic}$ in the plane $z_3=0$ with equation
$z_0z_2-z_1^2=0$. The line spanned by $(0,0,0,1)$ and $(1,0,1,0)$
is disjoint from it. Tube its image in $M_{\rm conic}$ to $c/2$
negatively oriented parallel spheres in the $S^2\times S^4$
summand and perform surgery. Its class is
$(t,-c/2)\in\mathbb Z/2\oplus\mathbb Z$, where $t$ is the
conic-surgery generator of \ref{maps:src:D.2.8}.
These choices are the \emph{semistable residual data}.
\end{definition}

\Needspace{6\baselineskip}
\begin{theorem}[Semistable residual smooth types]\hypertarget{tag-02LL}{}\label{maps:src:D.2.44}
{\normalfont\scriptsize tag 02LL}\par
Every threefold of Definition~\ref{maps:def:semistable-residual-data}
is simply connected and satisfies
\begin{equation}
 \begin{gathered}
 H_2(X;\mathbb Z)=H_3(X;\mathbb Z)=\mathbb Z/c,
 \qquad \lambda(X)=0,\qquad \mu(X)=0,\\
 P(X)=
 \begin{cases}
 0,&\mathrm{OS55},\ i_0=2,4,\\
 2,&\mathrm{OS55},\ i_0=3,\\
 0,&\mathrm{OS45},\ i_0=1,4,\\
 4,&\mathrm{OS45},\ i_0=2,3.
 \end{cases}
 \end{gathered}
 \hypertarget{tag-0282}{}\label{maps:eq:D44.1}
\end{equation}
Here $P$ is Zhubr's refined quarter-Pontryagin class. Every cyclic
generator gives the displayed values. The threefolds with $P=0$
are oriented diffeomorphic to $Z_{c,0}$, and those with $P=c/2$
are oriented diffeomorphic to $Z_{c,1}$, with the markings in the proof.
\end{theorem}
\begin{proof}\leavevmode\par
\begin{enumerate}[leftmargin=*]
\item\emph{The marked quotient and the integral auxiliary line.}
Use the physical lattice and alternating form
\[
 \Lambda=\mathbb Z\langle e_1,e_2,\delta,\gamma\rangle,
 \qquad
 \Omega=a e_1^*\wedge e_2^*-c\delta^*\wedge\gamma^*.
\]
Positive $\gamma$ covers translation by $-cG$. The marked graph and
cap calculation of \ref{maps:src:D.2.26}, in the all-weight marking
of \ref{maps:src:D.2.41}, gives the following explicit maps. On the
exterior $W$ with the ordinary disc unfilled,
\[
 \begin{gathered}
 \pi_1(W)=\mathbb Z\gamma,\\
 H_2(W)=\mathbb Zs\oplus
 \mathbb Z\langle t_{\rm sw}(e_1),t_{\rm sw}(e_2),t_{\rm sw}(\gamma)\rangle,\quad t_{\rm sw}(\delta)=0.
 \end{gathered}
\]
The fibre quotient is $[v]\mapsto\Omega(v)$, with inverse
$1\mapsto s$. Representatives are
$s=e_1\wedge e_2+\delta\wedge\gamma$ for OS55 and
$s=e_1\wedge e_2$ for OS45. The other three inverse coordinates are
the origin-section sweeps. The sole fixed-sweep relation is the
primitive global $\delta$-sum, which gives $t_{\rm sw}(\delta)=0$ after the
other caps are filled.

The ordinary boundary shear is
$(y,t)\mapsto(y+te\gamma,t)$, with inverse
$(y,t)\mapsto(y-te\gamma,t)$. Its complete Thom kernel is
\[
 t_{\rm sw}(v)-\Omega(v,-e\gamma)s\quad(v\in\Lambda),\qquad t_{\rm sw}(-e\gamma).
\]
It kills $\gamma$ in $\pi_1$ and eliminates the three sweeps in
$H_2$, with inverse substitution $t_{\rm sw}(v)=\Omega(v,-e\gamma)s$.
The remaining relation is $ce\,s=0$. Thus $X$ is simply connected
and $H_2(X)=\mathbb Z/c$, with inverse marking $1\mapsto[s]$ for
both signs. The prescribed cusp contributes Euler number two and
the other caps contribute zero. Integral duality then gives
\[
 H_0=H_6=\mathbb Z,\qquad H_2=H_3=\mathbb Z/c,\qquad
 H_1=H_4=H_5=0.
\]
In particular $H^2(X;\mathbb Z)=0$, so $c_1(TX)=0$ and the spin
structure is unique. Let $\alpha(s)=1$ define
$\alpha\in H^2(X;\mathbb Z/c)$.

We construct one line on $W$ together with its entire finite
$G_c=C_c$ lift. For
$Q(\lambda)=\lambda_1\lambda_2+\lambda_1+\lambda_2\pmod2$, put
\[
 J_\lambda(x,\zeta)=
 (x+\lambda,(-1)^{Q(\lambda)}e^{\pi i\Omega(\lambda,x)}\zeta).
\]
The polar form of $Q$ is $\Omega\bmod2$, so these are exact lattice
factors. The connection $d-\pi i\Omega(x,dx)$ descends and has
Chern class $\Omega$. The period monodromies preserve $Q$:
the elliptic quadratic form is the unique odd one over $\mathbb F_2$,
and every added elliptic translation column is divisible by the even
integer $c$. Their scalar-one lifts, and their matrix inverses, have
product one; hence every section-sweep value is zero. The map
\[
 (x,\zeta)\longmapsto
 (x+\delta/c,e^{\pi i\Omega(\delta/c,x)}\zeta)
\]
commutes with all periods because $\Omega(\delta/c,\Lambda)$ is
integral. Its $c$th power is $J_\delta$, so it gives an order-$c$ lift.

The local extension across the prescribed caps uses the following marking. At an $I_\nu$ place write
$A-I=\nu uu^tJ$, $g=\kappa u$, and choose $r$ with $\det(u,r)=1$.
Set $S=(u,\delta)$ and $B=(r,\gamma)$. The valuation matrix is
\begin{equation}
 C_i=\begin{pmatrix}
 \nu&c\kappa\\-a\kappa&(\varepsilon_i-ac\kappa^2)/\nu
 \end{pmatrix},\qquad
 \varepsilon_i=\begin{cases}1&i=i_0,\\0&i\ne i_0.\end{cases}
 \hypertarget{tag-0283}{}\label{maps:eq:D44.10}
\end{equation}
At the selected place $\nu=1$ and $\det C_{i_0}=1$. Write $C=C_{i_0}$.
The short basis of the prescribed fan is $SC$, its long basis is
$B$, and the interleaved four-dimensional basis has determinant one.
Applying the integer inverse recovers the physical coordinates. In these coordinates the scalar vector is
$d=C^{-1}(0,1)^t=(-c\kappa,1)^t$, and
\[
 K=C^t\operatorname{diag}(a,-c)
 =\begin{pmatrix}a&ac\kappa\\ac\kappa&c(ac\kappa^2-1)\end{pmatrix}.
\]

For real period coordinates $x=SCU+BB'$ and complex coordinates
$Z=U+\Pi B'$, use the full period matrix
$\Pi=(\log q)/(2\pi i)I+\Pi_0(q)$, including all units. The
Gaussian coordinate change is multiplication by
\[
 \exp\!\left(\pi i[-U^tKB'-(B')^tK\Pi B'+(1,0)Z]\right),
\]
and its reciprocal is the inverse. Its long factors are
\[
 (-1)^{\eta_{\rm char}^th}
 \exp\!\left(-2\pi ih^tKZ-\pi ih^tK\Pi h
                     +\pi i(1,0)\Pi h\right).
\]
Here $\eta_{\rm char}=(Q(r),Q(\gamma))^t=(1,0)^t$: $r$ is primitive
and the elliptic quadratic form is one on every nonzero vector
of $\mathbb F_2^2$. Substitution verifies the base-loop identity
obtained by replacing $\Pi$ by $\Pi+I$ and $U$ by $U+B'$. The integral support function
is $\psi(j)=(j^tKj-j_1)/2$. On a prescribed triangle, let
$m_\sigma$ be its affine interpolation. The Cartier coordinate and
inverse are $\zeta_\sigma=\chi^{m_\sigma}\zeta_C$ and
$\zeta_C=\chi^{-m_\sigma}\zeta_\sigma$; their translation
differences are exactly the long factors just displayed. In physical
coordinates the extra monomial is $X^{L_0}$, with reciprocal inverse,
where $L_0=(a\operatorname{row}_1C-(1,0))/2$. Its coefficients
are integers. Thus every period unit and the physical scalar lift
are transported into the fan charts.

At rank one let $T_{\rm loc}$ denote the full period return and choose
short basis $(l,p)$ and long basis $(h_{\rm per},w)$ with
$(T_{\rm loc}-I)h_{\rm per}=l$, $(T_{\rm loc}-I)w=0$, and set
$K_1=(l,p)^t\Omega(h_{\rm per},w)$.
Put $\epsilon=(Q(l),Q(p))^t$. Since the short plane is isotropic,
these two values specify its restricted semicharacter. The Gaussian
exponent is
$-U^tK_1B'-(B')^tK_1^t\Pi_1B'+\epsilon^t(U+\Pi_1B')$;
its reciprocal is again the inverse. The first row of $K_1$ is
$(N,0)$, and the integral support is $(Nj^2-Q(l)j)/2$.
The node maps
\[
 x_+=X/q^j,\quad x_-=q^{j+1}/X,\qquad
 X=x_+q^j,\quad q=x_+x_-
\]
give inverse identifications, preserving the transverse coordinate
and the period units. The Cartier maps extend the same line and
agree with its lattice factors on the annuli. We have therefore
constructed a smooth $G_c$-equivariant line $\mathcal L$ on $W$ with
\(\rho_c c_1(\mathcal L)=\alpha|_W\), including zero swept values.

\item\emph{The fixed data and the full cyclic invariants.}
The two cusp charts adjacent to the relevant edge are
\[
 (x,y,w)=(q/(X_1X_2),X_2,X_1),\qquad
 (x',y',w')=(q/X_2,X_1X_2,X_1^{-1}).
\]
Their overlap is $(x',y',w')=(wx,wy,w^{-1})$, with inverse
$(x,y,w)=(w'x',w'y',(w')^{-1})$. Thus the edge curve $C_0$ has
normal bundle $\mathcal O(-1)\oplus\mathcal O(-1)$.
The full scalar weights on the two charts are
\[
 (c\kappa-1,1,-c\kappa),\qquad(-1,1-c\kappa,c\kappa).
\]
Modulo $c$, the normal characters are $(-1,1)$. Every other edge
isotropy has order $1$ or $|1-c\kappa|$, relatively prime to $c$.
The Gaussian half-phase contributes $d_1/2$ to the line character,
and its first Cartier slope is
$((K_{11}-1)/2,K_{22}/2)$. Hence
\[
 r_C=\frac c2[a\kappa(c\kappa-1)-1]\pmod c,\qquad
 D_C=\deg\mathcal L|_{C_0}=K_{11}-K_{12}=a(1-c\kappa).
\]
The latter degree follows directly from the transition
$\zeta_{\sigma'}=w^{-(K_{11}-K_{12})}\zeta_\sigma$.

The only nodal isotropy divisible by a factor of $c$ occurs at
$I_c$, at position1 in OS55 and position0 in OS45. Its
$\kappa_E$ is respectively $-3$ and $3$. Choose
$a\kappa_Et_{\rm Bez}+cs_{0,\rm Bez}=1$, using $(t_{\rm Bez},s_{0,\rm Bez})=(1,4)$ or $(3,-1)$.
The bases are
\[
 \begin{gathered}
 l=cu-a\kappa_E\delta,\quad p=t_{\rm Bez}u+s_{0,\rm Bez}\delta,\\
 h_{\rm per}=r,\quad w=\gamma-\kappa_Er.
 \end{gathered}
\]
The short matrix has determinant one and adjugate inverse, with
$\delta=-t_{\rm Bez}l+cp$. Its pairing matrix is
$K_1=\left(\begin{smallmatrix}ac&0\\at_{\rm Bez}&-1\end{smallmatrix}\right)$.
Thus the elliptic double curve $E$ has line degree $-1$, normal
characters $(-t_{\rm Bez},t_{\rm Bez})$, and line character $c(1-at_{\rm Bez})/2=0\pmod c$.
The surviving period multiplies one normal coordinate by a constant
nonzero unit and the other by its reciprocal. The two normal lines
are flat and have degree zero. All other nodal isotropies are one,
and the ordinary cap is free. For every nonidentity element of $G_c$
the entire fixed set is $C_0\amalg E$.

The complete data are
\[
\begin{array}{c|c|r|r|r|r}
 \mathrm{OS}&i_0&\kappa&r_C&D_C&t_{\rm Bez}\\\hline
 \mathrm{OS55}&2,4&-1&0&25&1\\
 \mathrm{OS55}&3&0&2&5&1\\
 \mathrm{OS45}&1,4&0&4&1&3\\
 \mathrm{OS45}&2,3&1&0&-7&3
\end{array}
\]
In particular $D_C\equiv1$, $t_{\rm Bez}^2\equiv1$ and
$2r_C=r_C^2=0$ in $\mathbb Z/c$.

The class $\alpha$ extends to the circle Borel space: its degree-two
Serre obstructions have targets determined by $H^1(X;\mathbb Z/c)=0$
and $H^3(BS^1;\mathbb Z/c)=0$. Restrict such an extension to $G_c$.
On $W_{G_c}$ the based fixed-vertex section gives
$\pi_1(W_{G_c})=\mathbb Z\gamma\times G_c$. Every fixed-point section
has zero projection to the torsion-free factor. The degree-two
universal-cover sequence therefore writes the difference from
$\rho_c c_1^{G_c}(\mathcal L)$ as
$h_{\rm Bor}\tau+k_{\rm mix}\eta_W\theta$, where $\eta_W$ is integral in origin,
$\theta$ is the degree-one $G_c$ class, and $\tau$ is the
character-one degree-two class. The same $h_{\rm Bor}$ occurs on both fixed
curves. Adjust this common character to zero. Since $D_C\equiv1$,
the normalized restrictions are
\begin{equation}
 \begin{gathered}
 \alpha|_{C_0}=x_C+r_C\tau,\qquad
 \alpha|_E=-x_E+\eta_E\theta,\\
 e_{G_c}(N_{C_0})=-\tau^2,\quad e_{G_c}(N_E)=-t_{\rm Bez}^2\tau^2.
 \end{gathered}
 \hypertarget{tag-0280}{}\label{maps:eq:D44.0a}
\end{equation}
Here $x_C,x_E$ integrate to one, and $\eta_E$ is an integral degree-one
reduction on $E$. Thus $\eta_E^2=0$. Although
$\theta^2=(c/2)\tau$, the mixed term $(\eta_E\theta)^2$ is zero.
The free complement has finite-dimensional quotient and disappears
after inverting $\tau$; both normal Euler leading coefficients are
units. Localization therefore gives the full coefficient cube zero:
the rational contribution is $-3r_C^2=0$, and the elliptic cube
has no nonzero curve-degree-two term. Equivalently, use
the common $h_{\rm Bor}$; its contribution is
$-3[D_C(r_C+h_{\rm Bor})^2-h_{\rm Bor}^2/t_{\rm Bez}^2]=0$.

The equivariant determinant of $TX$ is trivial. Its ordinary first
Chern class is zero, and its residual base character is zero at a
fixed vertex, where the sum of the three tangent weights is zero.
The resulting equivariant $SU(3)$ reduction gives
$\lambda_{G_c}=-c_{2,G_c}$ by the integral universal spin identity.
On the two fixed curves $c_{2,G_c}$ is the normal Euler class, so
\begin{equation}
 \langle\alpha\lambda,[X]\rangle
       =1-D_C=1-a=0\quad\text{in }\mathbb Z/c.
 \hypertarget{tag-0281}{}\label{maps:eq:D44.0b}
\end{equation}
Reduction identifies $H^4(X;\mathbb Z)$ with
$H^4(X;\mathbb Z/c)$, and the marked $\alpha$ detects its cyclic
coordinate; hence $\lambda=0$ integrally. For $c=4,8$ the positive
$\mathbb {CP}^3$ generator identifies
$H_6(K(\mathbb Z/c,2);\mathbb Z)=\mathbb Z/c$, with cube evaluation
one. Its inverse is multiplication of that generator by the cube
value. Thus the computed cube proves the full $\mu=0$
\cite[Sections~3.11--3.14, pp.~806--808]{manZhubr75}.
The exact spin-bordism coefficient morphism gives
$2P=\operatorname{PD}(\lambda)$; its kernel here is
$\{0,c/2\}$. We next determine which of these two values occurs.

\item\emph{A characteristic trace and an integral index.}
Fill the ordinary cap of $W$ with zero shear, obtaining $Y$.
The origin-section maps give $\pi_1(Y)=\mathbb Z\gamma$ and
\[
 H_2(Y)=\mathbb Zs\oplus\mathbb ZB_0.
\]
Indeed the boundary $T^4\times S^1_\beta$ maps in degree one by
$(v,k\beta)\mapsto(\gamma^*(v)\gamma,-v)$, whose kernel is
$\mathbb Z\beta$. In degree two the fibre summand eliminates
the cap's identical fibre coordinate; the boundary sweeps kill
the three section sweeps, leaving $\mathbb Zs$. The two origin
discs join to a sphere $B_0$ mapping to $\beta$ and split the
resulting exact sequence. Let $\beta_{\rm coh}$ denote the positive
base cohomology class, so $\beta_{\rm coh}(B_0)=1$, $\beta_{\rm coh}(s)=0$.
The line $\mathcal L$ extends across this zero-shear cap; write
$\ell_{\rm line}=c_1(\mathcal L)$.

The tangent determinant satisfies $c_1(TY)=\beta_{\rm coh}$. To fix its sign
and all cap contributions, in the degree-one Weierstrass bundle
use $t'=1/t$, $x'=t^{-2}x$, $y'=t^{-3}y$. Then
\[
 dt'\wedge dx'/y'=-t^{-1}dt\wedge dx/y.
\]
Appending $d\log v$ gives the three-form transition: line-frame
units add only surface differentials, and elliptic translations add
only base differentials, which wedge to zero. At a node $xy=q$,
$d\log x\wedge dz\wedge dq=dx\wedge dz\wedge dy$; at the cusp
$q=xyw$, $X_1=w$, $X_2=y$, the corresponding form is
$dw\wedge dy\wedge dx$. These are nonzero in every prescribed
unimodular chart. Long periods multiply $X_j$ by $u_j(q)q^{h_j}$;
their added logarithmic differentials are multiples of $dq$ and do
not change the volume form. The inverse uses reciprocal units and
powers. Thus $K_Y=\pi^*\mathcal O(-1)$ with no additional cap
divisor. The scalar action preserves this determinant comparison.

Write the ordinary cap as
$A_{\rm ord}=T^3_{e_1,e_2,\delta}\times S^1_\gamma\times D^2_\beta$.
Attach the round handle $T^3\times D^2_\gamma\times D^2_\beta$
to $Y\times I_{\rm tr}$, where $I_{\rm tr}$ is the trace interval, by
\[
 (q,w,z)\mapsto(q,w,w^ez),\qquad
 (q,w,z)\mapsto(q,w,w^{-e}z)\quad\text{for its inverse}.
\]
Its attaching boundary matrix in $(\gamma,\beta)$ is
$F_e=\left(\begin{smallmatrix}1&0\\e&1\end{smallmatrix}\right)$.
The prescribed ordinary shear is
$G_e=\left(\begin{smallmatrix}1&e\\0&1\end{smallmatrix}\right)$.
The other solid-torus face is identified with the prescribed cap by
\[
 (z',w')\mapsto(w'z'^{[e]},(w')^{-e}),\qquad
 (z_\gamma,w_\beta)\mapsto
       (w_\beta z_\gamma^{[e]},w_\beta^{-e}),
\]
where $z^{[1]}=z$ and $z^{[-1]}=\overline z$. These maps are
smooth inverses including the disc centres. Their boundary matrix
$R_e=\left(\begin{smallmatrix}1&e\\-e&0\end{smallmatrix}\right)$
satisfies $F_eR_e=G_e$. The attaching face has negative old product
orientation; the other face has negative product orientation on
$T^3\times D^2_\gamma\times S^1_\beta$. The displayed solid-torus
map reverses the circle for $e=1$ and conjugates the disc for $e=-1$,
so it has sign $-1$ in both cases and preserves the outgoing
orientation. We have constructed an oriented cobordism $Z_e$ from
$Y$ to the stipulated $X_e$.

Set $\widehat\omega=\beta_{\rm coh}+2\ell_{\rm line}$, an integral characteristic class on
$Y$. On $A_{\rm ord}$ it is
$2a e_1^*\wedge e_2^*-2c\delta^*\wedge\gamma^*$.
Modulo $2c$ this is the restriction of $2a e_1^*\wedge e_2^*$
from the round handle. The Mayer--Vietoris map in degree two is
injective: $H^1(Y)$ supplies $\gamma^*$ on $A_{\rm ord}$, and the handle
supplies the other three generators. Thus the glued class is
uniquely characteristic, because its mod-two restrictions equal
$w_2$ on both pieces. Its outgoing restriction evaluates to two
on $s$, and the marked universal-coefficient isomorphism identifies
it with $i_2\alpha$, where $i_2(j)=2j$ in $\mathbb Z/(2c)$.
The inverse on its even image is $2j\mapsto j\pmod c$.

To obtain a simply connected integral endpoint, perform framed
circle surgery on the primitive $\gamma$ in the ordinary product
cap of $Y$. Use its product rank-five normal frame. Its trace $T_{\rm surg}$
has $H_2(T_{\rm surg},Y)=\mathbb Z\to H_1(Y)=\mathbb Z$ equal to the
isomorphism $1\mapsto\gamma$, so $H_2(Y)\to H_2(T_{\rm surg})$ is an
isomorphism. The reverse handle has index five, giving the same
isomorphism from the outgoing simply connected $Y_1$. The integral
class $\widehat\omega$ extends uniquely, since $H^1(Y)$ maps onto
the attaching circle's $H^1$. Its mod-two extension is $w_2(T_{\rm surg})$.

In the composite from $X_e$ to $Y_1$, first glue the matching
mod-$2c$ classes. A possible difference of their mod-two reduction
from $w_2$ is the connecting class of $\epsilon_{\rm conn}\gamma^*$ in
$H^1(Y;\mathbb Z/2)$. Lift $\epsilon_{\rm conn}$ to $0$ or $1$ modulo $2c$
and add its connecting class. Reduction commutes with this map,
and connecting classes vanish on both pieces and on both outer
boundaries. This makes the complete class characteristic while
preserving $i_2\alpha$ and the integral endpoint class $\omega_1$.
Both traces and their composite are simply connected by van Kampen.

Zhubr's characteristic invariant obeys
\[
 \Gamma_c(X,i_2\alpha)=\alpha(P(X))-
                         \langle\alpha^3,[X]\rangle.
\]
Lift the marked spin bordism class through the
free-coefficient surjection. Write $(Y_{\rm int},v_{\rm int})$ for a representative
of this integral lift, with $Y_{\rm int}$ a closed spin six-manifold.
The integral normalization is
\[
 \begin{gathered}
 \frac{\langle p_1(2v_{\rm int})-(2v_{\rm int})^3,[Y_{\rm int}]\rangle}{8}\\
       =\frac{\langle p_1v_{\rm int},[Y_{\rm int}]\rangle}{4}-\langle v_{\rm int}^3,[Y_{\rm int}]\rangle.
 \end{gathered}
\]
The first term is the free refined class; its coefficient pushforward
is $\alpha(P)$. The second is the coefficient cube. This proves
the map, whose inverse adds that cube, including even coefficients
\cite[Theorem~5, p.~796]{manZhubr75}; \cite[pp.~471--472]{manZhubr2010Gamma}.
The characteristic trace and integral normalization therefore give
\[
 P(X_e)=
 \frac{\langle p_1(Y)(\beta_{\rm coh}+2\ell_{\rm line})-(\beta_{\rm coh}+2\ell_{\rm line})^3,[Y]\rangle}{8}
 \pmod c.
\]
The number is integral: the same numerator is obtained on $Y_1$
with its integral characteristic lift. The characteristic-bordism axiom applies to these simply connected
endpoints and this characteristic oriented cobordism
\cite[pp.~471,479]{manZhubr2010Gamma}.

Use the canonical almost-complex spin-c symbol on $Y$ twisted by
the smooth equivariant line $\mathcal L$. Average a Hermitian
metric and unitary connection over $G_c$. Its even and odd spinors
are $\Lambda^{0,\mathrm{even}}T^*Y\otimes\mathcal L$ and
$\Lambda^{0,\mathrm{odd}}T^*Y\otimes\mathcal L$; Clifford
multiplication by $\xi$ has inverse $-\mathsf c(\xi)/|\xi|^2$ off zero.
Thus its index $I_{\rm ind}$ is an integral virtual $G_c$-character. The general
equivariant symbol theorem and the $U(3)$ symbol calculation give
\[
 \begin{gathered}
 I_{\rm ind}(1)=\int_Y e^{\ell_{\rm line}}e^{\beta_{\rm coh}/2}\widehat A(TY),\\
 P(X_e)=-6I_{\rm ind}(1)\pmod c
 \end{gathered}
\]
\cite[Theorem~3.9 and formulas~(4.1)--(4.4), pp.~562--567]{AtiyahSinger1968IIIE}.
This equality follows by extracting degree six in the integral index formula.

The fixed curves and their line data on $Y$ are unchanged, because
the zero-shear ordinary cap is free. Let $g_c$ be the generator of $G_c$ given by the lift of $\delta/c$.
At every nonidentity $g_c^j$,
put $z=\exp(2\pi ij/c)$. The fixed-symbol formula is
\[
 I_{\rm ind}(g_c^j)=\frac{D_C z^{r_C}}{2-z-z^{-1}}
                      -\frac{1}{2-z^{t_{\rm Bez}}-z^{-t_{\rm Bez}}}.
\]
On $C_0$ the Todd factor $1+x_C$ cancels the factor $1+x_C$ in
$(1-ze^{x_C})(1-z^{-1}e^{x_C})=(2-z-z^{-1})(1+x_C)$; its remaining degree
is $D_C$. On $E$ the normal lines are flat and the line degree is
$-1$. This proves the formula with its signs at every nonidentity
element, including all proper orders.

Let $R_c(z)=1+z+\cdots+z^{c-1}$. Restriction of integral virtual
characters to all nonidentity elements has kernel $\mathbb ZR_c$.
Every quotient class has a unique coefficient vector with constant
coefficient zero, obtained by subtracting that coefficient from all
entries. Inclusion of this representative is its inverse. Augmentation
modulo $c$ is well defined because $R_c(1)=c$. An exact Green inverse is
\[
 q_c(z)=-\frac1{2c}\sum_{j=0}^{c-1}j(c-j)z^j,\qquad
 (2-z-z^{-1})q_c(z)=1-R_c(z)/c\pmod{z^c-1}.
\]
The coefficient at zero is $(c-1)/c$ and every other coefficient is
$-1/c$, proving the identity. Apply it to
$D_Cz^{r_C}q_c(z)-q_c(z^{t_{\rm Bez}})$ and subtract its constant multiple of
$R_c$. This gives
\[
\begin{array}{c|c|r|r}
(c,r_C,D_C,t_{\rm Bez})&\text{integral coefficient vector}&I_{\rm ind}(1)\pmod c&-6I_{\rm ind}(1)\pmod c\\\hline
(4,0,25,1)&(0,-9,-12,-9)&-30&0\\
(4,2,5,1)&(0,1,3,1)&5&2\\
(8,4,1,3)&(0,1,1,1,2,1,1,1)&8&0\\
(8,0,-7,3)&(0,4,6,7,8,7,6,4)&42&4
\end{array}
\]
The regular-character ambiguity changes augmentation by a multiple
of $c$ and has no effect on the last column. This proves all the
refined values in \eqref{maps:eq:D44.1} for both ordinary signs.

\item \emph{A framing parity that depends only on the three-cycle.}
Brown's framed construction compares a regular three-torus with
a torsion three-sphere as follows.
Let a closed oriented six-manifold $Y_{\rm fr}$ have a stable tangent framing.
For an embedded oriented three-manifold $N_{\rm fr}\subset Y_{\rm fr}$ with a chosen
normal framing, denote by $[N_{\rm fr}]_{\rm fr}\in\Omega^{\rm fr}_3
\cong\mathbb Z/24$ its induced stable framed-bordism class. Then
\begin{equation}
 [N_{\rm fr}]_{\rm fr}\pmod2
   =q_{Y_{\rm fr}}\bigl(\operatorname{PD}_{\mathbb F_2}[N_{\rm fr}]\bigr)             \hypertarget{tag-0288}{}\label{maps:eq:D44.2}
\end{equation}
for a function of the indicated mod-two cohomology class alone.
This statement includes torsion classes in integral homology.
Here is a derivation directly from Brown's construction
\cite[pp.~369--372]{ffBrown72}. Set $K_3=K(\mathbb F_2,3)$ and let
$j:S^3\to K_3$ represent its fundamental class. Brown's $F(u_{\rm Br},0)$
is the stabilization of the map $u_{\rm Br}:Y_{\rm fr}\to K_3$.
The framed Pontryagin--Thom collapse $Y_{\rm fr}\to S^3$, followed by $j$,
therefore represents $F(\operatorname{PD}[N_{\rm fr}],0)$.
Apply the framed duality map $W_*A_\alpha$ of his Theorem 1.7 and
Corollary 1.11. Its value is the image of $[N_{\rm fr}]_{\rm fr}$ under
\begin{equation}
 j_*:\pi_3^{\rm S}\longrightarrow\pi_6^{\rm S}K_3
          \cong\mathbb Z/2.                                  \hypertarget{tag-0289}{}\label{maps:eq:D44.3}
\end{equation}
The target is Brown's Lemma 1.2. The remark immediately following
his Corollary 1.13 identifies $j_*\nu_{\rm H}$ with its nonzero generator:
the coefficient is the mod-two Hopf invariant, and the quaternionic
Hopf element $\nu_{\rm H}$ has Hopf invariant one. Dividing Brown's
$\mathbb Z/4$-valued function, whose values here are $0,2$, by two
gives \eqref{maps:eq:D44.2} for the specified normal framing.


\begin{enumerate}[leftmargin=*]
\item The homomorphism $\pi_3SO(3)\to\pi_3SO(6)$ has image
$2\mathbb Z$ and fixes the sphere-framing convention.
Indeed, under the double cover $SU(2)\to SO(3)$ the adjoint
bundle of a rank-two complex bundle $E_{\rm bun}$ has $p_1=-4c_2(E_{\rm bun})$,
whereas its underlying real rank-four bundle has $p_1=-2c_2(E_{\rm bun})$
and represents the stable generator. This proves the factor two;
the quotient map to $\pi_3V_{6,3}$ is reduction modulo two.
Consequently changing the normal framing of an embedded $S^3$
does not change the parity of its induced stable framed class.
For the standard sphere bounding a ball, the induced stable framed
class is zero and its tangent three-frame represents the nonzero
class in $\pi_3V_{6,3}=\mathbb Z/2$ relative to a constant ambient
three-frame. Equivalently, the tangent Lie framing on $S^3$ has
class $\nu_{\rm H}$, and a change of ambient frame by degree $d_{\rm fr}$ changes
that class by $d_{\rm fr}\nu_{\rm H}$, while the relative three-frame class changes
by $d_{\rm fr}\bmod2$. Thus Zhubr's required nontrivial relative
three-frame is equivalent to \emph{even} induced stable framed
class on the sphere. These conventions agree with the quaternionic
calculation in \cite[Section~5, pp.~8--9]{ffZhubr13}.


\item  Since $\lambda(X)=0$, $TX$ is stably trivial: the only obstruction
for a spin bundle on a six-dimensional complex is its class in
$H^4(-;\pi_3\mathrm{Spin})$, namely $\lambda$.
Moreover $TX$ is trivial on $X^\circ=X\setminus\{p_{\rm punct}\}$.
For example its complex bundle has $c_1=c_2=0$ and its last
obstruction lies in $H^6(X^\circ;\mathbb Z)=0$.
Any framing there extends stably across $p_{\rm punct}$, since
$\pi_5SO=0$. All such full framings have the same relative
three-frame parity on an embedded torsion $S^3$: their difference
lifts to $\mathrm{Spin}(6)$ and its degree-three class is a class
of $H^3(X^\circ;\mathbb Z)$, which evaluates to zero on a torsion
three-cycle.
An embedded sphere in the integral class of the regular torus
$T=\langle e_1,e_2,\gamma\rangle$ used below is constructed as follows.
Apply Kaplan's construction to
the spin structure induced on the normally framed $T$:
it gives a four-dimensional filling with one zero-handle and
only two-handles, preserving that almost framing
\cite[Definition~1.7 and Remark~1.8, pp.~239--240; Theorem~3.1, pp.~243--246]{mapsKaplan79}.
Delete the zero-handle and reverse the cobordism. Realize the
resulting two-handle attachments as ambient surgeries in
$X^\circ$. Their attaching circles bound embedded disks
with interiors disjoint from $T$ and from one another, since
$\pi_1(X^\circ)=0$, $2+3<6$ and $2+2<6$.
Prescribe a boundary collar along the first normal vector.
The two tangent framing directions of each attachment and
the other two normal directions give a four-frame around
the boundary of its disk. Its extension obstruction in
$\pi_1SO(4)=\mathbb Z/2$ is exactly the spin obstruction
of that two-handle, and vanishes. Thickening these disks
therefore performs the prescribed surgeries inside $X^\circ$.
The resulting embedded $S^3$ has the same integral homology
class as $T$. The embedded sphere has Brown's parity and the prescribed spin
structure by these framed ambient surgeries.
Zhubr's Theorem 1 now implies
\begin{equation}
 P(X)=0\quad\Longleftrightarrow\quad
 \text{a full ambient framing is admissible on the torsion }S^3.
                                                               \hypertarget{tag-028A}{}\label{maps:eq:D44.4}
\end{equation}
For the forward implication take the admissible three-frame in that
theorem. If $P=0$, its spin rank-three complement has vanishing
primary obstruction $c_2$; the only possible further obstruction
on the punctured six-manifold is in $H^5(X^\circ;\mathbb Z/2)=0$.
The complement is therefore trivial, and the three-frame extends
to a full framing. Conversely an admissible full framing has
trivial complement and hence refined class zero. Independence
of the full framing proves \eqref{maps:eq:D44.4} for the particular framing
used below.

 
\end{enumerate}


\item \emph{The reference frame and the critical-cycle defect.}
Write $f:X\to\mathbb P^1$ for the torus fibration and
$\mathcal V=\ker(df)$ for its holomorphic differential kernel.
This is a locally free rank-two sheaf, including at the critical
curves. In the nodal chart $f=xy$ it has basis
$x\partial_x-y\partial_y,\partial_z$; in the triangular cusp
chart $f=xyz$ it has basis
$x\partial_x-y\partial_y,y\partial_y-z\partial_z$.
For the first assertion, a relation among $(y,x,0)$ forces the
first two coefficients to be divisible by $x,y$, with opposite
quotients. For the second, a relation among $(yz,xz,xy)$ forces
the coefficients to be divisible by $x,y,z$, and their quotients
sum to zero. These arguments prove generation of the kernels,
including at the critical strata.
The scalar $\delta$-action defines a nowhere-zero section of this
\emph{abstract} vector bundle, although its image as a vector
field vanishes on the fixed curve. In a toric chart its two
coefficients are those of the nonzero primitive scalar lattice
vector, so they cannot both vanish. This description is unchanged
by the translation clutch.
As $H^2(X;\mathbb Z)=0$, the quotient line $\mathcal V/\mathcal O\delta$
is topologically trivial; hence $\mathcal V$ has a global smooth
frame whose first vector is $\delta$. The line
$\mathcal T=f^*T\mathbb P^1$ is also topologically trivial.


\begin{enumerate}[leftmargin=*]
\item  Let $Z_{\rm crit}$ denote the scheme defined by the coefficients of $df$.
There is an exact sequence
\begin{equation}
 0\longrightarrow\mathcal V\longrightarrow TX
  \xrightarrow{df}\mathcal T
  \longrightarrow\mathcal T\otimes\mathcal O_{Z_{\rm crit}}
  \longrightarrow0.                                         \hypertarget{tag-028B}{}\label{maps:eq:D44.5}
\end{equation}
Its codimension-two cycle is the sum $C_{\rm crit}$ of
the complex double curves, each with coefficient one. This
includes the three axes at a triple point: the ideal
$(xy,xz,yz)$ has multiplicity one on each axis. Since
$\mathcal V$ and $\mathcal T$ are topologically trivial,
\eqref{maps:eq:D44.5} gives $c_2(TX)=[C_{\rm crit}]$. Thus this cycle is
integrally nullhomologous.


\item  Off $Z_{\rm crit}$, splitting \eqref{maps:eq:D44.5} identifies $TX$ with the framed
bundle $\mathcal V\oplus\mathcal T$. Its comparison with a
global frame on $X^\circ$ has an integral degree-three class
$\mathfrak a$ whose local residue on every smooth double curve is one,
up to the same overall sign. In the model $xy$ this is the
explicit Bott matrix
\begin{equation}
 \frac1{\sqrt{|x|^2+|y|^2}}
 \begin{pmatrix}x&\overline y\\-y&\overline x\end{pmatrix}
       \in SU(2).                                            \hypertarget{tag-028C}{}\label{maps:eq:D44.6}
\end{equation}
It has degree one on the normal three-sphere; adding the
tangent line to the critical curve sends this to the stable
generator, since $\lambda=-c_2$ for the real bundle underlying
this complex rank-two bundle. Under the adjoint map $SO(3)\to SO$
the class is instead $-2c_2$, giving the factor two computed above.
Removing the finitely many triple points does not affect this
degree-three residue calculation. The integral relative
cohomology sequence, or equivalently the Thom intersection
construction, therefore gives
\begin{equation}
 \langle\mathfrak a,[T]\rangle
      =\pm\operatorname{lk}(T,C_{\rm crit})                   \hypertarget{tag-028D}{}\label{maps:eq:D44.7}
\end{equation}
for any regular three-torus $T$ disjoint from the critical
cycle. The ambiguity of adding a class from $H^3(X^\circ;\mathbb Z)$
evaluates to zero on our torsion three-cycle. Thus the integer
in \eqref{maps:eq:D44.7} is well defined in $\mathbb Z$.


\item  Take specifically $T=\langle e_1,e_2,\gamma\rangle$ in a regular
fibre inside the ordinary clutch disc. The reference frame
$\mathcal V\oplus\mathcal T$ has zero degree-three difference
from the flat translation frame there. Indeed the vertical
frames have the same first vector $\delta$; their change lies
in the stabilizer $\mathbb C\rtimes\mathbb C^*$, which retracts
to a circle. The base line contributes another $\mathbb C^*$,
and changing the horizontal splitting contributes a contractible
space. None contributes a degree-three class. The induced flat
framing of $T^3$ is its product Lie framing, with stable class
$\eta^3=12\nu_{\rm H}$, an even element of $\mathbb Z/24$.
A possible determinant-circle transition cannot change this
parity. A map $T^3\to S^1$ is homotopic to a character and,
after an integral change of torus basis, depends on just one
circle. Its real rotation can be placed stably in that circle
and one normal direction. It changes just that circle's framed
class, so the resulting product-framed torus has class either
$0$ or $\eta^3$, both even. After splitting off the determinant,
the remaining $SU(3)$ transition is a degree-$k$ map on a
three-dimensional complex and changes the framed class by
$\pm k\nu_{\rm H}$. Thus its exact parity is the integer in \eqref{maps:eq:D44.7}.
Equations \eqref{maps:eq:D44.2} and \eqref{maps:eq:D44.7} consequently give
\begin{equation}
 q_X(\operatorname{PD}[T])
       =\operatorname{lk}(T,C_{\rm crit})\pmod2.              \hypertarget{tag-028E}{}\label{maps:eq:D44.8}
\end{equation}
The torus represents an odd generator of $H_3(X)=\mathbb Z/c$
by the explicit linking calculation below. Its reduction is
therefore the same as that of an embedded torsion generator
$S^3$. Combine \eqref{maps:eq:D44.4} and \eqref{maps:eq:D44.8} to obtain
\begin{equation}
 P(X)=\frac c2
       \bigl(\operatorname{lk}(T,C_{\rm crit})\bmod2\bigr).
                                                               \hypertarget{tag-028F}{}\label{maps:eq:D44.9}
\end{equation}

 
\end{enumerate}

\item\emph{The integral critical-cycle calculation in the prescribed fan.}
Compute the integer in \eqref{maps:eq:D44.9} in the prescribed
fan marking of \eqref{maps:eq:D44.10}. At a rank-one place set
\[
 d_0=\gcd(\nu,a\kappa),\qquad
 \mathsf a=(\nu/d_0,-a\kappa/d_0)^t,\qquad
 \mathsf r=(d_0,c\kappa d_0/\nu).
\]
Both vectors are primitive integral in these rows. Choose
$\sigma_{\rm per},\beta_{\rm per}$ with $\det(\mathsf a,\sigma_{\rm per})=1$ and
$\mathsf r\beta_{\rm per}=1$, and put
$\ell=(-\mathsf r_2,\mathsf r_1)^t$. The Euclidean algorithm gives
these choices and their inverse coordinate maps. The physical basis
\[
 H=(S\mathsf a,B\beta_{\rm per},S\sigma_{\rm per},B\ell)
\]
has determinant one and $H^{-1}T_{\rm loc}H=I+E_{12}$. The product chart
$xy=q$ identifies its positive transverse elliptic double curve
with the specialization of $(S\sigma_{\rm per})\wedge(B\ell)$, with inverse
marking given by the last two coordinates of $H^{-1}$. Its degree is
\[
 \Omega\bigl((S\sigma_{\rm per})\wedge(B\ell)\bigr)
       =-c d_0^2/\nu.
\]
For example, the two determinant-one bases give
$\Omega(S\mathsf a,B\beta_{\rm per})=a\nu/d_0^2$ and
$\Omega(S\mathsf a,B\ell)=0$; the Pfaffian is $-ac$, which gives
the displayed remaining pairing. Changing $\sigma_{\rm per}$ by an integral
multiple of $\mathsf a$ adds a vanishing-circle torus bounded by
its normal disc.

At the cusp write $A_j=SCe_j$, $B_j=Be_j$. Let $T_{\rm cusp}$
be the positive period return. The displayed full period matrix changes
by $\Pi\mapsto\Pi+I$, so $T_{\rm cusp}$ fixes $A_j$ and sends
$B_j$ to $B_j+A_j$. Thus $(T_{\rm cusp}-I)B_j=A_j$.
In the positive wall convention, the degree of the line with
support $\psi$ is the sum at the two opposite vertices minus
the sum at the two common edge vertices. The vertical, horizontal
and diagonal walls have four-vertex data
\[
 \{00,01;10,(-1,1)\},\quad
 \{00,10;01,(1,-1)\},\quad
 \{10,01;00,11\}.
\]
Evaluation of $r(r-1)/2,rs,s(s-1)/2$ gives the matrix
$M=\left(\begin{smallmatrix}1&0&0\\-1&-1&1\\0&1&0\end{smallmatrix}\right)$,
whose inverse is
$\left(\begin{smallmatrix}1&0&0\\0&0&1\\1&1&1\end{smallmatrix}\right)$.
Equivalently, for the dual Cartier lines, whose supports are the
negatives of those three polynomials, and in horizontal, vertical,
diagonal order, the degree matrix is
\begin{equation}
 \begin{pmatrix}0&-1&0\\1&1&-1\\-1&0&0\end{pmatrix}.
 \hypertarget{tag-0284}{}\label{maps:eq:D44.11}
\end{equation}
The conversion is $-M$ followed by the interchange of its first
two columns; applying the same interchange and negation is the
inverse. Thus all support and curve orientation conventions are
fixed explicitly.

The complete integral fibre specialization at this primitive cusp
is the free quotient with basis
\[
 [A_1\wedge B_1],\quad
 [A_2\wedge B_1]=[A_1\wedge B_2],\quad
 [A_2\wedge B_2],\quad[B_1\wedge B_2].
\]
The two primitive relations are $A_1\wedge A_2$ and
$A_1\wedge B_2-A_2\wedge B_1$. They are the short normal-disc
relation and the long-period monodromy-cylinder relation.
Subtracting these two columns gives the inverse quotient coordinates;
the period-square cell decomposition of the primitive cap gives
the same four coordinates in its $H_2$. The three mixed coordinate
functionals are the integral Cartier lines above. The fourth is
projection to the dual long-period two-torus and evaluates to one
on $B_1\wedge B_2$.
Each double curve projects into a one-dimensional edge of that
dual complex, so its fourth coordinate is zero. Denote the first
three quotient coordinates by $x,y,z$ in the displayed order.
The Cartier evaluations and the zero fourth coordinate determine
the curves in the entire integral cap group. In vertical, horizontal
and diagonal order, the resulting classes are
\begin{equation}
 [C_v]=x-y,\qquad [C_h]=z-y,\qquad [C_d]=y.
\hypertarget{tag-Q24Z}{}\label{maps:eq:primitive-cusp-wall-classes}
\end{equation}
Indeed their coordinate columns are exactly $M$. Its displayed
integral inverse sends a coordinate column $(r,s,t)^t$ to
$(r,t,r+s+t)^t$ in the wall basis. Composing either way gives the identity. The representatives
$A_1\wedge B_1-A_2\wedge B_1$,
$A_2\wedge B_2-A_2\wedge B_1$ and $A_2\wedge B_1$ have these
coordinates; subtracting the two primitive relations gives the
inverse specialization coordinates. Their sum is $(1,-1,1)$.
Hence the three complex double curves together have integral preimage
\begin{equation}
 A_1\wedge B_1+B_1\wedge A_2+A_2\wedge B_2.
 \hypertarget{tag-0285}{}\label{maps:eq:D44.12}
\end{equation}
Replacing the middle term by $B_2\wedge A_1$ adds exactly the
displayed mixed relation. This also proves independence of that
alternative preimage. In particular the fourth coordinate has
been determined as well as the three Cartier degrees.

Since $K=C^t\operatorname{diag}(a,-c)$, the three degrees of the
auxiliary line on these curves are $K_{11}-K_{12}$,
$K_{22}-K_{12}$ and $K_{12}$. Their sum is
$a-ac\kappa-c+ac^2\kappa^2$. Transport all five local preimages to
the marked regular fibre along the origin-section cuts.
Let $v$ be their sum and $R$ the full exterior monodromy relation
matrix. The marked fibre quotient supplies an integral basis
$\mathsf F_{\rm fib}=(s_0,K_0)$ and an integer relation inverse $\mathsf Q_{\rm rel}$ with
$R\mathsf Q_{\rm rel}=K_0$, $\Omega(s_0)=1$. Let $\pi_{\rm rel}$ be projection onto
the five $K_0$ coordinates, with inclusion $z\mapsto(0,z)$
as its section. Therefore
\[
 \begin{gathered}
 q_{\rm rel}=\mathsf Q_{\rm rel}\pi_{\rm rel}\mathsf F_{\rm fib}^{-1}(v-\Omega(v)s_0),\\
 Rq_{\rm rel}=v-\Omega(v)s_0.
 \end{gathered}
\]
These relation columns are the transported monodromy cylinders,
all outside the ordinary disc. The ordinary $\delta$ Thom chain
has boundary $-ce\,s_0$, as its complete normal-map formula shows.
Consequently the full bounding equation and the linking number are
\[
 v=Rq_{\rm rel}+L(-ce\,s_0),\qquad
 L=\operatorname{lk}(T,C_{\rm crit})=-\frac{\Omega(v)}{ce}.
\]
The normal chain meets $T=e_1e_2\gamma$ once with its oriented
$\delta$ coefficient. A zero relation first has zero non-$\delta$
sweep coordinates; its remaining coefficient $h_{\rm rel}$ satisfies
$ceh_{\rm rel}=0$ in $\mathbb Z$, and hence is zero. Thus the integer is
independent of the relation solution. The same argument shows
that a cap-kernel change has no effect.

The exact values are
\begin{equation}
 \begin{array}{c|c|c|c|c}
 \mathrm{OS}&c&i_0&\Omega(v)&
       \operatorname{lk}(T,C_{\rm crit})\ (e=1)\\\hline
 \mathrm{OS45}&8&1,2,3,4&-32,24,24,-32&4,-3,-3,4\\
 \mathrm{OS55}&4&2,3,4&72,-28,72&-18,7,-18
 \end{array}
 \hypertarget{tag-0286}{}\label{maps:eq:D44.13}
\end{equation}
For OS55 at zero position $2$ the five preimages in the ordered physical
bivector coordinates are
\begin{equation}
 \begin{gathered}
 (0,0,0,-12,0,5),\quad(3,0,1,-12,0,4),\\
 (5,-24,-8,0,4,-19),\quad
 (0,0,0,0,0,1),\quad(0,4,0,0,0,1).
 \end{gathered}
 \hypertarget{tag-0287}{}\label{maps:eq:D44.14}
\end{equation}
Their $\Omega$ values are $-20,-1,101,-4,-4$, summing to $72$,
so their bounding ordinary loop has $\delta$ coefficient $-18$.
For $e=-1$ the same normal formula reverses all the linking
signs in \eqref{maps:eq:D44.13}. For $c$ times the marked generator
$s$, its ordinary $\delta$ coefficient is exactly $-e$; hence
its torsion linking with $T$ is the unit $-e/c$ in $\mathbb Q/\mathbb Z$.
This proves that $T$ is an odd generator of $H_3(X)$ and verifies
the generator assertion used in the framing argument. Applying
\eqref{maps:eq:D44.9} to these integers gives the same four
refined values as the characteristic-index calculation.

\item\emph{The concrete surgery models and oriented recognition.}
Represent $c[S^2]$ in $S^2\times S^4$ by tubing $c$ disjoint
parallel spheres along disjoint arcs. Its oriented rank-four normal
bundle has zero $w_2$, and so is trivial; its framing is unique
because $\pi_2SO(4)=0$. The spin surgery trace has marked quotient
$H_2=\mathbb Z/c$, with inverse $1\mapsto[S^2]$.
The incoming map to $K(\mathbb Z/c,2)$ extends across the trace,
since the attaching sphere maps to zero. The incoming product
bounds $S^2\times D^5$ with the same coefficient map. Thus
$(P,\mu)(Z_{c,0})=(0,0)$.

For the second model, use coordinates in which the conic lies in
$z_3=0$ and has equation $z_0z_2-z_1^2=0$. The projective line
spanned by $(0,0,0,1)$ and $(1,0,1,0)$ misses the conic: its
only intersection with the containing plane is $(1,0,1,0)$,
where that equation is nonzero. It therefore survives as an
embedded sphere representing $t$ in $M_{\rm conic}$.
Tube it to $c/2$ negatively oriented parallel spheres in the
$S^2\times S^4$ summand. This constructs the stated embedded
representative of $(t,-c/2)$, again with its unique spin framing.
The exact quotient and inverse are
\[
 \frac{\mathbb Z/2\langle t\rangle\oplus\mathbb Z\langle a_{\rm sph}\rangle}
      {\langle t-(c/2)a_{\rm sph}\rangle}
 \ \underset{1\mapsto[a_{\rm sph}]}{\overset{t\mapsto c/2,\ a_{\rm sph}\mapsto1}
 {\rightleftarrows}}\ \mathbb Z/c.
\]
The attaching sphere maps to zero, so the coefficient map extends
over this surgery trace as well. Naturality transports the conic
invariants under $1\mapsto c/2$: its refined class becomes $c/2$
and its complete cubic becomes $(c/2)^3=0\pmod c$.
The product summand contributes zero. Hence
$(P,\mu)(Z_{c,1})=(c/2,0)$
\cite[Theorem~5, p.~796]{manZhubr75}.

These surgeries preserve simple connectivity. Their Euler number
is two and their second homology is the displayed finite cyclic
group; integral duality gives third homology $\mathbb Z/c$ and
zero free third Betti number. Both have $\lambda=2P=0$.
The marked classification in the zero-free-third-homology scope
of Zhubr's Theorem 3 and its Section 4.1 construction therefore gives
the asserted oriented diffeomorphisms
\cite[Theorem~3, p.~795; Section~4.1, pp.~809--810]{manZhubr75}.
Every automorphism of $\mathbb Z/c$ is an odd unit, which fixes
both $0$ and $c/2$, so the conclusion is independent of the chosen
generator and includes both ordinary signs.
\end{enumerate}
\end{proof}


\Needspace{9\baselineskip}
\subsubsection{Additive caps and order twenty}

\begin{definition}[Residual additive data]\hypertarget{tag-Q24Y}{}\label{maps:def:additive-residual-data}
{\normalfont\scriptsize tag Q24Y}\par
Use Definition~\ref{mapsIII:def:elliptic-marked-data}, the primitive
Mumford cap at every multiplicative site, the prescribed fan at
the selected position, and the additive fillings in this table:
\begin{equation}
\begin{array}{c|c|c|c}
 \mathrm{OS/profile}&(a,b)&\text{zero indices}&\text{type-}II\text{ choice}\\\hline
 \mathrm{OS45}/\mathrm{II}&(1,8)&2,3&\text{original or good log}\\
 \mathrm{OS55}/\mathrm{II}&(5,4)&3&\text{original or good log}\\
 \mathrm{OS55}/\mathrm{base}&(1,20)&2,3,4&\text{none}\\
 \mathrm{OS55}/\mathrm{II}&(1,20)&3&\text{original or good log}.
\end{array}
\hypertarget{tag-028G}{}\label{maps:eq:D45.1}
\end{equation}
Selected positions are $i_0=i-1$ in
Definition~\ref{mapsIII:def:elliptic-marked-data}. In an original-only pattern the ordinary
clutch is $e\gamma$ with $e=\pm1$. At a good $II$ logarithm use
every order $m=6k$, $k>0$, every integer ordinary clutch $e$, and
every power vector
\begin{equation}
 v=r\delta+n\ell,\qquad r\in\mathbb Z,\qquad
 n-me=\sigma\in\{1,-1\},
\hypertarget{tag-028R}{}\label{maps:eq:D45.2}
\end{equation}
where $\ell$ is the primitive fixed lift with
$\gamma^*(\ell)=1$. Every affine twist is allowed.


These listed sites are all the logarithmic sites in this definition.
Use the identical oriented models $Z_{c,0},Z_{c,1}$, $c=4,8$,
of Definition~\ref{maps:def:semistable-residual-data}, with their
specified embeddings, normal framings and primitive generators.
Use $Z_{5,1}$ from Definition~\ref{maps:def:minimal-mixed-data},
marked by its refined generator, with fundamental-class invariant zero.
These choices are the \emph{residual additive data}.
\end{definition}

\Needspace{6\baselineskip}
\begin{theorem}[Residual additive smooth types]\hypertarget{tag-02LM}{}\label{maps:src:D.2.45}
{\normalfont\scriptsize tag 02LM}\par
The data of Definition~\ref{maps:def:additive-residual-data} are
exactly the simply connected data in its listed one-log sector, and
all are free. Every resulting $X$ is spin and satisfies
$H_2(X)=H_3(X)=\mathbb Z/b$ and $b_3(X)=0$.
Let $s$ be the primitive fibre generator specified in the proof.
The full invariants, with $P$ the coefficient of the refined class
in $\mathbb Z/b\langle s\rangle$ and $\lambda$ the coefficient of
$\operatorname{PD}(p_1/2)$, are
\begin{equation}
\begin{array}{c|c|c|c|c|c|c|c}
 \mathrm{OS/profile}&(a,b)&\text{zero}&II&P&\lambda&\mu&
       \text{oriented type}\\\hline
 \mathrm{OS45}/\mathrm{II}&(1,8)&2&O\text{ or }F&4&0&0&Z_{8,1}\\
 \mathrm{OS45}/\mathrm{II}&(1,8)&3&O\text{ or }F&0&0&0&Z_{8,0}\\
 \mathrm{OS55}/\mathrm{II}&(5,4)&3&O\text{ or }F&0&0&0&Z_{4,0}\\
 \mathrm{OS55}/\mathrm{base}&(1,20)&2&-&4&8&0&Z_{4,0}\#Z_{5,1}\\
 \mathrm{OS55}/\mathrm{base}&(1,20)&3&-&14&8&0&Z_{4,1}\#Z_{5,1}\\
 \mathrm{OS55}/\mathrm{base}&(1,20)&4&-&4&8&0&Z_{4,0}\#Z_{5,1}\\
 \mathrm{OS55}/\mathrm{II}&(1,20)&3&O\text{ or }F&4&8&0&Z_{4,0}\#Z_{5,1}.
\end{array}
\hypertarget{tag-028V}{}\label{maps:eq:D45.3}
\end{equation}
Here $O$ means the original cap and $F$ means any cap in
\eqref{maps:eq:D45.2}. Expanding those alternatives gives eleven
patterns. Together with \ref{maps:src:D.2.44}, this identifies all
eighteen minimal patterns of residual order $4,8,20$, throughout
their stated finite-parameter scopes.
\end{theorem}
\begin{proof}
\begin{enumerate}
\item \emph{The periods, fixed lift and filling parameters.}
Use the physical lattice
$\Lambda=\mathbb Z\langle e_1,e_2,\delta,\gamma\rangle$.
Positive long return covers translation by $-bG$.
The primitive alternating functional is
\begin{equation}
 \Omega=a e_1^*\wedge e_2^*-b\delta^*\wedge\gamma^*.
\hypertarget{tag-028N}{}\label{maps:eq:D45.16}
\end{equation}
At a multiplicative site the defining matrix and section column satisfy
$B-I=\nu uu^tJ$, $G=\kappa u$, where
$J=\left(\begin{smallmatrix}0&1\\-1&0\end{smallmatrix}\right)$.
The complete rank-one data needed here are
\[
\begin{array}{c|c|c|c|r}
\text{profile}&i&\text{fibre}&u&\kappa\\\hline
\mathrm{OS45}/\mathrm{II}&0&I_8&(1,0)&3\\
\mathrm{OS45}/\mathrm{II}&2&I_1&(1,-2)&1\\
\mathrm{OS45}/\mathrm{II}&3&I_1&(1,1)&0\\
\mathrm{OS55}/\mathrm{II}&0&I_5&(1,0)&-3\\
\mathrm{OS55}/\mathrm{II}&1&I_4&(1,-1)&-1\\
\mathrm{OS55}/\mathrm{II}&3&I_1&(1,1)&-1\\
\mathrm{OS55}/\mathrm{base}&0&I_5&(1,0)&-3\\
\mathrm{OS55}/\mathrm{base}&1&I_4&(1,0)&-3\\
\mathrm{OS55}/\mathrm{base}&2&I_1&(2,-1)&-1\\
\mathrm{OS55}/\mathrm{base}&3&I_1&(1,-2)&0\\
\mathrm{OS55}/\mathrm{base}&4&I_1&(1,1)&-1.
\end{array}
\]
These equalities follow by dividing a nonzero column of $B-I$
by its coordinate gcd, choosing its positive first nonzero entry,
and substituting the displayed $u,\nu,\kappa$.
If $pu_1+qu_2=1$, set $r_u=(-q,p)$; the basis $(u,r_u)$
has determinant one and inverse its integral adjugate.

Put $\beta=-G^tJB$. The full return has elliptic block $B$,
last elliptic column $bG$, angular row $a\beta$, and last
angular entry
\[
 c_i=\frac{[i=*]-ab\kappa_i^2}{\nu_i}
 \quad\text{at a multiplicative site},\qquad
 c_i=ab\beta_i(B_i-I)^{-1}G_i\quad\text{at }II.
\]
Substitution in the ordered defining tuple gives an integral return,
fixing $\delta$, preserving $\Omega$, with ordered product identity.
At the $II$ site the pairs $(B,G)$ are
\[
 \left(\begin{pmatrix}2&3\\-1&-1\end{pmatrix},
             \binom{1}{-1}\right)
 \quad\text{and}\quad
 \left(\begin{pmatrix}5&3\\-7&-4\end{pmatrix},\binom{0}{0}\right)
\]
for $45/II$ and $55/II$, respectively.
Both have $B^6=I$ and $\det(B-I)=1$.
Writing $q_0=bG$, the fixed lift and its geometric curve are
\begin{equation}
 z=(B-I)^{-1}q_0\in\mathbb Z^2,\qquad
 \ell=(-z_1,-z_2,0,1),\qquad [E_{II}]=\delta\wedge\ell.
\hypertarget{tag-028O}{}\label{maps:eq:D45.17}
\end{equation}
Thus $\ell=(-8,0,0,1)$ for $45/II$ and $\ell=\gamma$
for $55/II$. The angular entry of $T\ell-\ell$ is zero
by the displayed formula for $c_i$.
With $\xi_{\rm ang}=a\beta(B-I)^{-1}$, the integral matrix
\[
 H_{II}=(e_1+(\xi_{\rm ang})_1\delta,e_2+(\xi_{\rm ang})_2\delta,\delta,\ell)
\]
has determinant one and conjugates the return to $B\oplus I_2$.
Its adjugate is the inverse. Since $B-I$ is unimodular, the fixed
and coinvariant coordinates are precisely $\delta,\ell$.
For every positive multiple $m$ of six the norm is $m$ on
these coordinates and zero on the moving elliptic plane.
The translation conjugacy on that plane is translation by
$(I-B)^{-1}$ times its affine constant, with negative translation
as inverse; all surviving line multipliers are carried with it.

The primitive cap and graph maps of \ref{maps:src:D.2.26}
and \ref{maps:src:D.2.41} give the fibre quotient
$\mathbb Zs$, measured by $\Omega$.
For its integral section, choose
$\alpha a-\beta_0b=1$ and set
$s=\alpha e_1\wedge e_2+\beta_0\delta\wedge\gamma$.
Together with $b e_1\wedge e_2+a\delta\wedge\gamma$
and the four mixed bivectors, this is a unimodular basis $\mathsf F_{\rm fib}$;
the first row of $\mathsf F_{\rm fib}^{-1}$ is $\Omega$.
For $a=1$ take $s=e_1\wedge e_2$; for $(a,b)=(5,4)$
take $s=e_1\wedge e_2+\delta\wedge\gamma$.
The five remaining coordinates are the complete integral return
image. Indeed, for
$R=(\bigwedge^2T_i-I)_i$, integral column Euclid on the last five
coordinates of $\mathsf F_{\rm fib}^{-1}R$ gives
$\mathsf U_{\rm rel}\pi \mathsf F_{\rm fib}^{-1}R\mathsf V_{\rm rel}=(I_5,0)$, where $\pi$ projects onto those five
coordinates. Every pivot is one for the seven displayed rows.
Let $\iota$ include $\mathbb Z^5$ as the first five coordinates
of the relation-column module. The section is $\mathsf V_{\rm rel}\iota\mathsf U_{\rm rel}$;
the displayed identity gives $\pi \mathsf F_{\rm fib}^{-1}R\mathsf V_{\rm rel}\iota\mathsf U_{\rm rel}=I_5$.
Their inverses are $\operatorname{adj}(\mathsf U_{\rm rel})/\det\mathsf U_{\rm rel}$ and
$\operatorname{adj}(\mathsf V_{\rm rel})/\det\mathsf V_{\rm rel}$, including the determinant signs.
Thus each claimed image element is a cylinder boundary.

The remaining attachment maps are as follows.
Let $D_{\rm sw},L_{\rm sw}$ be the fixed sweeps of $\delta,\ell$.
For a finite $II$ cap their normal and root columns, in
coordinates $(s,D_{\rm sw},L_{\rm sw})$, are
\[
 C_1=(bn,m,0),\quad C_2=(-br,0,m),\quad C_3=(0,r,n).
\]
Let $G_e=(be,1,0)$ be the negative of the ordinary
$\delta$ normal column. The ordinary gluing is
$(y,t)\mapsto(y+te\gamma,t)$, with inverse subtraction.
The root-loop group has relations $mc=n\gamma$, $c=e\gamma$,
after the angular discs kill $e_1,e_2,\delta$.
Its quotient map is $\gamma\mapsto1$, $c\mapsto e$ in
$\mathbb Z/(n-me)$, with inverse $1\mapsto\gamma$.
Consequently simple connectivity is exactly $n-me=\pm1$.
The full freeness conditions of \ref{maps:src:D.2.30} then
hold because $\gcd(m,n)=1$, independently of $r$.

For $n=me+\sigma$ the changes of relation generators are
\[
 E_1=bs=\sigma(C_1-mG_e),\quad
 E_2=D_{\rm sw}+be\,s=G_e,\quad
 E_3=L_{\rm sw}=\sigma(C_3-rG_e-eC_2),
\]
with inverse
\[
 C_1=\sigma E_1+mE_2,\quad C_2=-rE_1+mE_3,\quad
 C_3=-erE_1+rE_2+nE_3,\quad G_e=E_2.
\]
These formulas identify the relation sublattices.
On the four-column source the primitive syzygy is
$v_0=(r,n,-m,0)^t$. Let $\mathsf A_{\rm rel}$ express the three $E$ columns
in the four $C_1,C_2,C_3,G_e$ columns, and let $\mathsf B_{\rm rel}$ express
those four columns in the $E$ coordinates as displayed.
Then $\mathsf B_{\rm rel}\mathsf A_{\rm rel}=I_3$ and
$\mathsf A_{\rm rel}\mathsf B_{\rm rel}=I_4-\sigma v_0(0,1,e,0)$.
Thus they are inverse maps between $\mathbb Z^3$ and
$\mathbb Z^4/\mathbb Zv_0$; $v_0$ is primitive since
$\gcd(m,n)=1$. The quotient in the ambient cycle module is
therefore $\mathbb Z/b$, sending $s$
to one and $D_{\rm sw},L_{\rm sw}$ to zero. Its inverse sends one to $s$.
In an original-only pattern the group quotient is $\mathbb Z/e$;
the ordinary normal has fibre value $-be$, and the other normal
columns eliminate the remaining sweeps. Thus precisely $e=\pm1$
is allowed and the same marked second homology results.
There is one cusp of Euler number two and all other caps have
Euler number zero. The complete cap maps, duality and Euler
characteristic give
\[
 H_0=H_6=\mathbb Z,\quad H_2=H_3=\mathbb Z/b,\quad
 H_1=H_4=H_5=0.
\]
In particular $H^2(X;\mathbb Z)=0$; the complex determinant
line has $c_1=0$, and $X$ has its unique spin structure.

\item \emph{The corrected differential and its reference frame.}
At an original $II$ cap the transverse map has coordinates
$q=y^2-x^3$. To obtain them, solve the local Weierstrass equation
for the base variable using its nonzero linear base coefficient,
split off the nondegenerate square in $y$, and take the analytic
cubic root of the remaining nonzero unit multiplying $x^3$.
These changes and their analytic inverses preserve the scalar
product coordinate. The full locally free kernel is
\begin{equation}
 \ker(dq)=\mathcal O(2y,3x^2,0)\oplus\mathcal O(0,0,1).
\hypertarget{tag-028W}{}\label{maps:eq:D45.4}
\end{equation}
Indeed, $-3x^2a_{\rm II}+2yb_{\rm II}=0$ implies
$y\mid a_{\rm II}$, then $a_{\rm II}=2yh_{\rm II}$ and
$b_{\rm II}=3x^2h_{\rm II}$; the third coordinate is arbitrary.
The ideal $(x^2,y)$ has basis $1,x$ in its quotient, so the
reduced critical elliptic curve has multiplicity two.
On the complement of that curve the comparison matrix is
\begin{equation}
 \frac1{\sqrt{4|y|^2+9|x|^4}}
 \begin{pmatrix}2y&-3\overline{x}^{\,2}\\
                 3x^2&2\overline y\end{pmatrix}.
\hypertarget{tag-028X}{}\label{maps:eq:D45.5}
\end{equation}
Its adjoint is its inverse and its determinant is one.
The first column is the normalized map $(2y,3x^2)$.
A sufficiently small regular value has two simple positive
complex preimages, so its degree is two.
The second abstract kernel generator is the nonzero scalar section,
including along the reduced critical curve.

At a free good logarithm the covering space has $q=s^m$ and
the quotient covering is \'etale. Let $D_{\rm fin}$ be the reduced
central divisor; its pullback to this cover is $\{s=0\}$.
The differential has target
\begin{equation}
 \mathcal L=f^*T\mathbb P^1\otimes\mathcal O(-(m-1)D_{\rm fin}).
\hypertarget{tag-028Y}{}\label{maps:eq:D45.6}
\end{equation}
Dividing the common factor $s^{m-1}$ leaves $m\,ds$.
With a nonzero unit $u$, the expression is $mu\,ds+s\,du$,
which remains surjective at $s=0$.
The target twist changes by the same deck character as the common
factor; hence the divided map descends. It has no corrected
critical cycle and its vertical kernel preserves the scalar section.
On $D_{\rm fin}$ this kernel is $TD_{\rm fin}$, since its pullback
is the kernel of $m\,ds$ along $s=0$.

For the node $xy$ the kernel is generated by $(x,-y,0)$ and
$(0,0,1)$; for $xyz$ it is generated by $(x,-y,0)$ and
$(x,0,-z)$. In the latter case the equation
$yz a_{\rm xyz}+xz b_{\rm xyz}+xy c_{\rm xyz}=0$ forces
$x\mid a_{\rm xyz}$, $y\mid b_{\rm xyz}$ and
$z\mid c_{\rm xyz}$; the remaining coefficients sum to zero.
These divisibility arguments give the full kernels.
The ideal $(xy,xz,yz)$ gives multiplicity one on each of its
three axes. The scalar section stays nonzero in the abstract
kernel. Its quotient line and $\mathcal L$ are smooth trivial
lines because $H^2(X;\mathbb Z)=0$.
Choose their trivializations and Hermitian splittings.
They give a reference tangent frame $\rho$ off the complete
critical cycle $\mathcal C$, with node multiplicity one,
triple-axis multiplicity one, and original-$II$ multiplicity two.
The quaternion comparison in \eqref{maps:eq:D45.5}, and its
node counterpart, fixes the relative orientation and yields
\begin{equation}
 \lambda(TX,\rho)=-\operatorname{PD}[\mathcal C].
\hypertarget{tag-028Z}{}\label{maps:eq:D45.7}
\end{equation}
Here $\lambda=p_1/2$ is the integral spin class, including its
relative version. Its image in absolute cohomology is computed from the complete
critical class in the next step.

\item \emph{The specified fan and the critical chain.}
At the selected $I_1$ write $S=(u,\delta)$ and
$L=(r_u,\gamma)$, and put
\[
 C_* =\begin{pmatrix}1&b\kappa\\-a\kappa&1-ab\kappa^2\end{pmatrix},
 \qquad H=(SC_*e_1,Le_1,SC_*e_2,Le_2).
\]
Both matrices are integral of determinant one, with adjugate
inverses. The stipulated unit triangulation is in the short basis
$SC_*$. In that basis the Cartier matrix is
\[
 K_* = C_*^t\operatorname{diag}(a,-b)
      =\begin{pmatrix}a&ab\kappa\\ab\kappa&ab^2\kappa^2-b\end{pmatrix}.
\]
On every unimodular triangle interpolate
$(j^tK_*j-j_1)/2$ by its affine integral function.
The jumps at the vertical, horizontal and diagonal edges are
$K_{11}-K_{12}$, $K_{22}-K_{12}$ and $K_{12}$.
The inverse sends the third degree to $K_{12}$ and adds it to
the first and second to recover $K_{11},K_{22}$.
Apply the primitive-cap specialization of
\eqref{maps:eq:primitive-cusp-wall-classes} with
$(A_1,B_1,A_2,B_2)=(H_1,H_2,H_3,H_4)$. This is the identity
on the ordered unit fan, its normal-disc and monodromy-cylinder
relations, and its dual long-period torus.
Thus the vertical, horizontal and diagonal curves have representatives
\[
 H_1\wedge H_2+H_2\wedge H_3,\qquad
 H_3\wedge H_4+H_2\wedge H_3,\qquad
 -H_2\wedge H_3.
\]
Their sum is $H_1\wedge H_2+H_2\wedge H_3+H_3\wedge H_4$,
with $\Omega$-value $a-ab\kappa-b+ab^2\kappa^2$.
The local annuli lie in the two smooth charts along each edge.
The charts $q/(X_1X_2),X_2,X_1$ and
$q/X_2,X_1X_2,X_1^{-1}$ overlap by
$(x,y,w)\mapsto(wx,wy,w^{-1})$, whose inverse is
$(x',y',w')\mapsto(w'x',w'y',(w')^{-1})$.
Their oriented collapsing annuli cap the adjacent triangles.
The complete quotient and its Cartier dual identify their other
boundaries with the three classes above, including the zero
long-period two-torus coordinate.

The alternative Jordan basis $(S,LC_*^{-1})$
is related to $(SC_*,L)$ by the coordinate map
$\operatorname{diag}(C_*,C_*)$, with inverse
$\operatorname{diag}(C_*^{-1},C_*^{-1})$.
On character tori the map is
$\chi^m\mapsto\chi^{C_*^tm}$, with inverse exponent
$C_*^{-t}$. It carries each cone to its image under $C_*$.
Thus the fan in $S$ coordinates is $C_*\Sigma$, when $\Sigma$
is the unit fan in $SC_*$ coordinates. The Cartier gradients in $S$ coordinates are the inverse-transpose
images of those on $\Sigma$, with inverse the transpose map.
The line-coordinate comparison below includes the corresponding
integral monomial and all period units. These are the chart and
line morphisms used for the displayed critical curves.

At any other multiplicative site let $d_0=\gcd(\nu,a\kappa)$.
The primitive collapsing vector is
$l=(\nu/d_0)u-(a\kappa/d_0)\delta$.
Choose an integral complement $p$ in the short plane and an
integral long vector $h_{\rm per}$ on which the primitive return covector
is one; let $\ell_0$ generate its kernel.
B\'ezout gives $(l,h_{\rm per},p,\ell_0)$ of determinant one and adjugate
inverse, carrying the return to $I+E_{12}$.
The surviving elliptic double curve is the degree-one image of
$p\wedge\ell_0$, whose value is
\[
 \Omega(p,\ell_0)=-\frac{b d_0^2}{\nu}.
\]
The local inverse reads its residual product coordinate.
In node charts $x_+=X/q^j$, $x_-=q^{j+1}/X$, the inverse is
$q=x_+x_-$, $X=x_+q^j$; the remaining elliptic period and its
nonzero unit remain unchanged. These chart maps give the degree-one
specialization, with the positive complex torus orientation.
At an original $II$ site \eqref{maps:eq:D45.17} is the product
$\delta,\ell$ torus through the elliptic cusp. Its component
group is one, so the specialization has degree one. Its critical
contribution is $2\delta\wedge\ell$, of value $-2b$.

Transport these local cycles along the fixed base paths and let
$K=\Omega(\mathcal C)$. The complete integers are
\begin{equation}
\begin{array}{c|c|c|c|r|c}
\text{profile}&(a,b)&\text{zero}&II&K&\operatorname{ord}[\mathcal C]\\\hline
\mathrm{OS45}/\mathrm{II}&(1,8)&2&O&24&1\\
\mathrm{OS45}/\mathrm{II}&(1,8)&2&F&40&1\\
\mathrm{OS45}/\mathrm{II}&(1,8)&3&O&-32&1\\
\mathrm{OS45}/\mathrm{II}&(1,8)&3&F&-16&1\\
\mathrm{OS55}/\mathrm{II}&(5,4)&3&O&72&1\\
\mathrm{OS55}/\mathrm{II}&(5,4)&3&F&80&1\\
\mathrm{OS55}/\mathrm{base}&(1,20)&2&-&352&5\\
\mathrm{OS55}/\mathrm{base}&(1,20)&3&-&-68&5\\
\mathrm{OS55}/\mathrm{base}&(1,20)&4&-&352&5\\
\mathrm{OS55}/\mathrm{II}&(1,20)&3&O&352&5\\
\mathrm{OS55}/\mathrm{II}&(1,20)&3&F&392&5.
\end{array}
\hypertarget{tag-028P}{}\label{maps:eq:D45.18}
\end{equation}
For example, with the $II$ contribution set to zero, the seven
ordered degree lists are
\[
\begin{gathered}
(-1,0,49,-8),\quad(-1,0,-8,-7),\quad(-20,-1,0,101),\\
(-4,-5,401,-20,-20),\quad(-4,-5,-20,-19,-20),\\
(-4,-5,-20,-20,401),\quad(-4,-5,0,401).
\end{gathered}
\]
Their sums are $40,-16,80,352,-68,352,392$.
The original cap subtracts $2b$. Every $K$ is even.
Apply the integral section of $R$ above to the difference between
the sum of the bivectors and $Ks$. It gives an integral
monodromy-cylinder chain with that boundary.
Together with the local annuli, this is a chain comparing
$\mathcal C$ with $Ks$, supported outside the finite and ordinary
detector cores. This support preserves the rational link before halving.
The absolute image of \eqref{maps:eq:D45.7} gives
$\lambda=-K\pmod b$, hence zero at $b=4,8$ and eight at $b=20$.

\item \emph{The finite detector and its tangent framing.}
At a finite $II$ core the central two-torus has presentation
\begin{equation}
 \langle\delta,\ell,g\mid mg=r\delta+n\ell\rangle.
\hypertarget{tag-028J}{}\label{maps:eq:D45.12}
\end{equation}
The elliptic return on $g$ is $B$, with identity return on
$\delta,\ell$. Use the primitive character
\begin{equation}
 \chi(\delta)=1,\qquad\chi(\ell)=-\sigma r,\qquad
 \chi(g)=-\sigma re.
\hypertarget{tag-028K}{}\label{maps:eq:D45.13}
\end{equation}
It respects \eqref{maps:eq:D45.12}, since
$m(-\sigma re)=r+n(-\sigma r)$.
Extend this character by zero on the moving elliptic lattice.
Its zero level $N\subset D_{\rm fin}$ is connected and transverse
to $\delta$.
Put $k_0=\ell+\sigma r\delta$ and
$h_0=g+\sigma re\delta$; then $mh_0=nk_0$.
The exact primitive kernel generator and inverse coordinates are
\[
 t=-\sigma e k_0+\sigma h_0=\sigma(g-e\ell),\qquad
 k_0=mt,\quad h_0=nt,
\]
because $(-\sigma e)m+\sigma n=1$.
Equivalently $\ell=-\sigma r\delta+mt$ and
$g=-\sigma re\delta+nt$; these maps are inverse integral
maps of the presented lattices.
The elliptic monodromy on $N$ is $B^\sigma$.
Both $B-I$ and $B^{-1}-I$ are unimodular, so their explicit
adjugate inverses and the Wang sequence give
\begin{equation}
 N\text{ is an affine }T^2\text{ mapping torus with return }B^\sigma,
 \qquad H_1(N;\mathbb Z)=\mathbb Z.
\hypertarget{tag-028L}{}\label{maps:eq:D45.14}
\end{equation}
For a return $z\mapsto B^\sigma z+c$, translation by
$(I-B^\sigma)^{-1}c$ conjugates it to its linear part;
negative translation is the inverse. This preserves every affine twist.

The reference frame induces a tangent parallelization on $N$.
Along $N$, write $TD_{\rm fin}|N=\mathbb R\delta\oplus\mathcal S_{\rm fin}$,
where $\mathcal S_{\rm fin}$ is spanned by
$J\delta$ and the two real directions of the second complex
vertical frame. Let $\ell_\delta$ be the projection onto the
$\delta$ coefficient in this decomposition.
The mutually inverse projections are
\[
 \mathcal S_{\rm fin}\longrightarrow TN,\quad w\longmapsto w-\chi(w)\delta,
 \qquad
 TN\longrightarrow\mathcal S_{\rm fin},\quad v\longmapsto v-\ell_\delta(v)\delta.
\]
Their composites are identity because $\chi(\delta)=1$,
$\chi(v)=0$ and $\ell_\delta(w)=0$.
The triangular shears are homotopic through invertible shears
to identity. The normal frame is $\delta$ followed by the two
real directions of the corrected target line.

For any tangent parallelization $\tau$ of a closed
oriented three-manifold $M$, the resulting stable framed parity is
\begin{equation}
 [M,\tau]_{\mathrm{fr}}\pmod2
        =1+\dim H_1(M;\mathbb F_2).
\hypertarget{tag-028M}{}\label{maps:eq:D45.15}
\end{equation}
Here is the parity argument with its spin scope.
Kaplan's prescribed-spin filling $W_{\rm spin}$
\cite[Definition~1.7 and Remark~1.8, pp.~239--240; Theorem~3.1, pp.~243--246]{mapsKaplan79}
has one zero-handle and $j$ two-handles and induces the spin
structure determined by $\tau$.
Its tangent bundle is trivial on its two-dimensional spine.
For the boundary outward-normal-plus-tangent frame the two
$\operatorname{Spin}(4)$ degrees have first-column degree $a_{\rm fr}-b_{\rm fr}$
and stable degree $a_{\rm fr}+b_{\rm fr}$. Their parities coincide with the relative
Euler number, hence with $\chi(W_{\rm spin})=1+j$.
The spin intersection matrix is alternating modulo two and has
even rank. Its cokernel is $H_1(M;\mathbb F_2)$, so
$j\equiv\dim H_1(M;\mathbb F_2)\pmod2$.
This proves \eqref{maps:eq:D45.15} for the specified induced spin
structure. A change of tangent frame through $\operatorname{Spin}(3)$
changes the stable degree by an even integer, consistently with
this formula. Thus \eqref{maps:eq:D45.14} gives even parity for $N$.
The ordinary detector $T^3=\langle e_1,e_2,\gamma\rangle$
is the zero level of the primitive $\delta$ character and has the
same projection construction; its parity is $1+3=0$.

\item \emph{The all-parameter chains and the full generator pairing.}
Every return fixes the single loop $\delta$. Its bar correction
is zero on that loop, including in the transported last peripheral
word: each two-letter correction contains a moving displacement,
which is zero for $\delta$.
The boundary of the punctured-base fundamental polygon gives
$\sum_iJ_i\delta=0$ for the oriented peripheral sweep maps.
This is an equality of cylinders, with paired sides cancelling.
Take the finite normal Thom column for $\delta$, add $m$ times
the ordinary $\delta$ normal, and add $m$ times the remaining
filled-cap $\delta$ sweeps. Their graph boundaries cancel.
The fibre part is
\begin{equation}
 -\delta\wedge(r\delta+n\ell)+me\delta\wedge\gamma,
 \qquad \Omega=b(n-me)=b\sigma.
\hypertarget{tag-028Q}{}\label{maps:eq:D45.19}
\end{equation}
Its difference from $b\sigma s$ is an integral return boundary,
obtained by the section of $R$. The two constant differences
$-\delta\wedge\ell-bs$ and $\delta\wedge\gamma+bs$
have zero $\Omega$ value; the same section supplies their
integral preimages. Multiplying the first by $n$ and the second
by $me$ gives the exact difference in \eqref{maps:eq:D45.19}
for every integer $m,n,e,r$.

The normal disc bundle on the central loop intersects the
level $N$ with multiplicity $\chi(\delta)=1$.
All ordinary and other-cap sweeps miss this detector.
The local critical comparison chains above also miss it.
With $B_4$ an oriented four-chain bounding $bN$, define the
rational critical link by
\begin{equation}
 \mathcal L(N,\mathcal C)=\frac{B_4\cdot\mathcal C}{b}.
\hypertarget{tag-0290}{}\label{maps:eq:D45.8}
\end{equation}
The value is independent of the bounding chain: their difference
is a closed integral four-cycle, and its intersection with the
torsion class $[\mathcal C]$ is zero.
The same oriented Thom chain and its complementary intersection
pairing therefore give
\begin{equation}
 \operatorname{lk}(N,s)=\frac{\sigma}{b}\pmod{\mathbb Z},
 \qquad \mathcal L(N,\mathcal C)=\frac{\sigma K}{b}.
\hypertarget{tag-028S}{}\label{maps:eq:D45.20}
\end{equation}
The unit first pairing proves that $N$ has exact order $b$,
including $b=20$. If a bounding column uses the affine root,
its intersection is its character $\chi(g)=-\sigma re$.
This term is included in the same chain comparison.
The multiplicity $m$ is already in the normal norm column;
the displayed primitive level has no further covering factor.

In an original-only pattern the ordinary normal has fibre value
$-be$ and meets its $T^3$ level once. With $e=\pm1$ this gives
\begin{equation}
 \operatorname{lk}(T^3,s)=-\frac e b\pmod{\mathbb Z},
 \qquad\mathcal L(T^3,\mathcal C)=-\frac{eK}{b}.
\hypertarget{tag-028T}{}\label{maps:eq:D45.21}
\end{equation}
It is a full generator. At $b=20$ the critical class has order
five, so one may instead bound $5\mathcal C$ and divide its
intersection by five; this is the same rational value.

\item \emph{Admissibility on the two-primary basis and full transgression.}
Let $N$ denote either preceding detector and let
$b=2^{\nu_2}q_{\rm odd}$ with $q_{\rm odd}$ odd. Thus
$q_{\rm odd}=1$ at $b=4,8$ and $q_{\rm odd}=5$
at $b=20$. Remove a disjoint six-ball.
The complement of the critical curves is simply connected by
relative general position: a two-dimensional nullhomotopy can
avoid their two-dimensional strata in a six-manifold.
Take $q_{\rm odd}$ positive normally framed parallel copies of $N$ and
join them by framed tubes. Their framed bordism class is the
sum of the $q_{\rm odd}$ incoming classes. Puncture Kaplan's prescribed-spin
filling and reverse it. Its two-handles turn that connected
three-manifold into a sphere. The attaching circles bound disjoint
embedded discs in the critical-curve complement by simple
connectivity and relative general position. The first normal
direction fixes the collar; the remaining normal four-frame extends
exactly when its spin obstruction in $\pi_1(SO(4))$ vanishes.
The prescribed spin handles have that extension. Their traces
are disjoint from the critical curves.
The outgoing embedded sphere represents $q_{\rm odd}[N]$, of exact order
$2^{\nu_2}$, and therefore is a basis of $\operatorname{Tors}_2H_3(X)$.

Zhubr's proved criterion
\cite[Theorem~1, p.~6; proof, pp.~8--9]{ffZhubr13}
chooses a three-frame $\mathfrak v$ on $X$ minus the ball, nontrivial
relative to this two-primary sphere. Its spin three-plane complement
$\xi$ has obstruction $o_4(\xi)=\operatorname{PD}P$.
The sphere has the required two-primary basis class, and the
complementary obstruction is the full integral refined class.

The surgery handle fixes the parity convention.
For $q_*=a_*+x_*i+y_*j+z_*k\in S^3$ let
$t_1=q_*i,t_2=q_*j,t_3=q_*k$ and let
$e^{\rm nor}_1,e^{\rm nor}_2,e^{\rm nor}_3$ be its normal frame.
The extending triple is
\[
 (-a_*e^{\rm nor}_3-x_*t_1-y_*t_2-z_*t_3,e^{\rm nor}_1,e^{\rm nor}_2).
\]
The coefficient map has degree one and surjects onto
$\pi_3(V_{6,3})=\mathbb Z/2$, giving the nontrivial relative
class. Stabilization of the complementary $\operatorname{Spin}(3)$
frame changes its degree by two: the adjoint and quaternion-real
representations have $p_1=-4c_2$ and $p_1=-2c_2$, respectively.
Thus this admissible sphere has even induced stable framed parity.

Choose a trivialization $\tau$ of $\xi|N$.
Such a trivialization exists because $B\operatorname{Spin}(3)$
has no homotopy in degrees at most three.
Choose its parallel copies and extend them across the tubes and
the two-handle trace. The relative handles have index one or two;
the extension obstructions are $\pi_0\operatorname{Spin}(3)$ and
$\pi_1\operatorname{Spin}(3)$, both zero.
Consequently the local full frames $\mathfrak v+\tau$ on the incoming
copies and the outgoing sphere are stably framed bordant.
Their parities satisfy $q_{\rm odd}\epsilon(\mathfrak v+\tau|N)=0$ modulo two.
Multiplication by the odd integer $q_{\rm odd}$ is the identity
on $\mathbb Z/2$; hence the full-order $N$
also has even parity. The global three-frame and these local
complementary trivializations therefore determine the full-order
comparison, including the order-twenty manifolds.

Compare $\mathfrak v+\tau|N$ with the reference frame $\rho|N$.
They induce the same spin structure, so their difference lifts to
$\operatorname{Spin}(6)$. Its degree-three class is an integer
$k_{\rm fr}$ on $N$. The stable $J$ homomorphism sends a generator to the
generator of $\Omega^{\mathrm{fr}}_3=\mathbb Z/24$.
Both frames have even parity, so $k_{\rm fr}$ is even.
In the relative pair $(X\setminus B^6,N)$ the obstruction identity is
\begin{equation}
 \lambda(TX,\rho|N)=2o_4(\xi,\tau)+\partial\mathfrak a,
 \qquad\langle\mathfrak a,[N]\rangle=k_{\rm fr}.
\hypertarget{tag-028H}{}\label{maps:eq:D45.10}
\end{equation}
The factor two is the same $\operatorname{Spin}(3)$ stabilization
map. Put $q_{\rm obs}=\langle o_4(\xi,\tau),B_4\rangle$. Evaluation gives
the integer identity
\begin{equation}
 -B_4\cdot\mathcal C=2q_{\rm obs}+bk_{\rm fr},\qquad k_{\rm fr}\equiv0\pmod2.
\hypertarget{tag-028I}{}\label{maps:eq:D45.11}
\end{equation}
Here $q_{\rm obs}/b$ is the torsion linking with the full $P$.
Dividing this integer identity and then reducing in $\mathbb Q/\mathbb Z$
gives
\begin{equation}
 \operatorname{lk}(N,P)=-\tfrac12\mathcal L(N,\mathcal C)
                      \pmod{\mathbb Z}.
\hypertarget{tag-0291}{}\label{maps:eq:D45.9}
\end{equation}
This formula includes the order-five coordinate at $b=20$.
Insert \eqref{maps:eq:D45.20} or \eqref{maps:eq:D45.21} and
cancel their displayed unit generator pairing. Since $K$ is even,
\begin{equation}
                   P=-K/2\quad\text{in }\mathbb Z/b.
\hypertarget{tag-028U}{}\label{maps:eq:D45.22}
\end{equation}
The division is of the integer in
\eqref{maps:eq:D45.11}, before passing to the cyclic group.
The original-$II$ change $K\mapsto K-2b$
changes $-K/2$ by $b$, so every finite and original alternative
in \eqref{maps:eq:D45.3} has the stated same value.

\item \emph{One exterior line with its full finite lift.}
Let $W$ be the union of the regular piece and all multiplicative caps,
omitting the ordinary cap and, when present, the $II$ cap.
At the selected cusp both short vectors $SC_*e_1,SC_*e_2$
bound. At every other multiplicative cap $l$ bounds.
The matrix of their first three coordinates has a right inverse
over $\mathbb Z$: column Euclid gives $(I_3,0)$, with unit
row and column transformations. Its section expresses each of
$e_1,e_2,\delta$ as a combination of angular disc boundaries.
The exponent $\gamma^*$ vanishes on all those relations and gives
the inverse coordinate on the surviving loop. The multiplicative
base meridians bound. If the $II$ cap is omitted its base meridian
$x$ remains and commutes with $\gamma$ by the residual return.
Thus the exponent maps and the loops are inverse presentations
\[
 \pi_1(W)=\mathbb Z\langle\gamma\rangle
 \quad\text{or}\quad
 \pi_1(W)=\mathbb Z\langle\gamma\rangle\oplus\mathbb Z\langle x\rangle.
\]
The cap relative sequence gives surjectivity from regular-piece
$H_2$ onto $H_2(W)$: each weight-one cap has one component,
whose disc maps with coefficient one to its base meridian.
The remaining coordinates are the fibre coordinate and
boundary sweeps. The integral class with fibre form $\Omega$
and zero swept values reduces to the unique
$\mathfrak u\in H^2(X;\mathbb Z/b)$ with $\mathfrak u(s)=1$, restricted to $W$.
Indeed the attachment relations above send every sweep to zero
modulo $b$, and $H_1(W)$ is free, so UCT has no Ext ambiguity.

We realize that integral class and its scalar lift explicitly.
On the real period cover set
\[
 Q(\lambda)=\lambda_1\lambda_2+\lambda_1+\lambda_2\pmod2,
 \quad
 J_\lambda(x,z)=
 (x+\lambda,(-1)^{Q(\lambda)}e^{\pi i\Omega(\lambda,x)}z).
\]
Since $a$ is odd and $b$ is even, the polar form of $Q$ is
$\Omega$ modulo two. Its elliptic part is the unique odd
quadratic refinement and each $bG$ shift is even, so every
displayed return preserves $Q$ and $\Omega$.
Direct composition proves the period cocycle law.
The connection $d-\pi i\Omega(x,dx)$ has first Chern class
$\Omega$. The scalar-one return lifts and their inverse matrix
lifts have ordered product identity and zero swept values.
For $v_\delta=\delta/b$ define
\[
 T_{v_\delta}(x,z)=(x+v_\delta,e^{\pi i\Omega(v_\delta,x)}z).
\]
The commutator with $J_\lambda$ is
$e^{2\pi i\Omega(v_\delta,\lambda)}=1$.
Also $T_{v_\delta}^b=J_\delta$, and the return lifts commute with $T_{v_\delta}$
because they fix $\delta$. The inverse is $T_{-v_\delta}$.
This is the full order-$b$ lift of the scalar subgroup.

At the cusp write $x=SC_*U+LB_{\rm per}$, $Z=U+\Pi B_{\rm per}$, preserving
the period matrix
$\Pi=(\log q)/(2\pi i)I+\Pi_{\rm hol}(q)$ and all its holomorphic units.
The line-coordinate change is
\[
 G_{K_*}=\exp\{\pi i[-U^tK_*B_{\rm per}-B_{\rm per}^tK_*\Pi B_{\rm per}+(1,0)Z]\},
\]
with reciprocal inverse. The long-period multiplier is
\[
 (-1)^{\eta^th}
 \exp\{-2\pi i h^tK_*Z-\pi i h^tK_*\Pi h
                         +\pi i(1,0)\Pi h\}.
\]
Here $\eta=(1,0)$: the two long generators are the primitive
elliptic vector $r_u$ and $\gamma$, with $Q$-values one and zero.
The function $\psi(j)=(j^tK_*j-j_1)/2$ is integral because
$\operatorname{diag}K_*=(1,0)$ modulo two.
If $m_\tau$ is its affine interpolation on a prescribed triangle,
the regular line coordinate is $z_\tau=\chi^{m_\tau}z_C$,
with reciprocal inverse on the overlap. Translation of $\psi$
gives exactly the divisor of the preceding long multiplier,
including its unit.
In physical short coordinates the support $\varphi$ satisfies
\[
 \varphi(C_*j)=-\psi(j)+L_0j,\qquad
 L_0=\frac{a\operatorname{row}_1(C_*)-(1,0)}2\in\mathbb Z^2.
\]
Multiplication by $X^{L_0}$ and by its reciprocal identifies the
physical and prescribed-fan regular frames. This is the line map
corresponding to the displayed torus and cone map for the prescribed fan.

At rank one use the basis $(l,h_{\rm per},p,\ell_0)$.
The line pairing matrix $(l,p)^t\Omega(h_{\rm per},\ell_0)$ has first row
$(n_{\rm pair},0)$. Return-invariance of $Q$ gives $n_{\rm pair}\equiv Q(l)\pmod2$.
The one-dimensional support is therefore
$(n_{\rm pair}j^2-Q(l)j)/2$. Its affine slopes and their reciprocal
monomials, together with the same period-coordinate Gaussian,
extend the line in the node charts already displayed.
All remaining elliptic factors are nonzero and keep their original
unit multipliers. The constructions agree on the dense collars,
so they define one line $\mathcal A$ on $W$, with
$\mathfrak u|W=\rho_b c_1(\mathcal A)$ and the finite lift $T_{v_\delta}$.

\item \emph{Every fixed component and its character.}
Let $C$ now denote the rational curve at the vertical cusp edge;
the critical cycle continues to be denoted $\mathcal C$.
The scalar vector in the stipulated short basis is
$C_*^{-1}\delta=(-b\kappa,1)$.
In the two edge charts its full circle weights are
$(b\kappa-1,1,-b\kappa)$ and $(-1,1-b\kappa,b\kappa)$.
The chart transition proves
$\nu_C=\mathcal O(-1)\oplus\mathcal O(-1)$.
For every $q\mid b$, $C_q$ fixes $C$ with normal characters
$-1,1$. The other edge isotropies divide $1$ or $1-b\kappa$,
which are units at the relevant primes.
The Gaussian weight and the adjacent Cartier slope give the
integer character and degree
\[
 r_C=\frac b2\{a\kappa(b\kappa-1)-1\},\qquad
 D_C=a(1-b\kappa).
\]
The other chart changes the character by $b\kappa D_C$,
giving the same character modulo every $q\mid b$.

At a rank-one cap the surviving scalar coordinate in $(l,p)$
is $(\nu/d_0)p$. Thus $\delta/q$ fixes its elliptic double
curve exactly when $q\mid\nu/d_0$.
For the coefficient groups needed here there is exactly one such
curve $E_q$, as follows:
\[
\begin{array}{c|c|c|c|r|r|r}
\text{profile}&(a,b)&q&\text{site}&\kappa&t&D_E\\\hline
\mathrm{OS45}/\mathrm{II}&(1,8)&8&I_8&3&3&-1\\
\mathrm{OS55}/\mathrm{II}&(5,4)&4&I_4&-1&3&-1\\
\mathrm{OS55}/\mathrm{base}&(1,20)&4&I_4&-3&1&-5\\
\mathrm{OS55}/\mathrm{II}&(1,20)&4&I_4&-1&3&-5\\
\mathrm{OS55}/\mathrm{base}\text{ or }\mathrm{II}&(1,20)&5&I_5&-3&3&-4.
\end{array}
\]
Here $d_0=1$, $\nu=q$ and $a\kappa t+qs_0=1$.
The exact primitive short basis is
$l=qu-a\kappa\delta$, $p=tu+s_0\delta$, with adjugate inverse;
take $h_{\rm per}=r_u$ and $\ell_0=\gamma-(b\kappa/q)r_u$.
Then
\[
 \delta=-tl+qp,\qquad
 (l,p)^t\Omega(h_{\rm per},\ell_0)=
 \begin{pmatrix}aq&0\\at&-b/q\end{pmatrix}.
\]
The normal characters are $-t,t$, and the two normal lines have
reciprocal constant unit multipliers on the elliptic period.
They are flat and of integral degree zero.
The line degree is $D_E=-b/q$.
Its quadratic-refinement values are $Q(l)=q$ and $Q(p)=t$
modulo two. The dense and Cartier factors give
\[
 r_E=\frac q2\{(t\bmod2)-at\}\equiv0\pmod q.
\]
The braces are even, including when $q=5$, so this is an
integral character before reduction.

The cusp characters for the seven rows, at $b=4,8$
and the four-part of $b=20$, are respectively
$0,4,0,0,2,0,0$. Thus at $b=20$ the character is two only
at base-profile position $3$; the other four-part characters are zero.
Every five-part cusp character is zero.
At $q=5$ the pair $(-3,3)$ is the same unordered pair as
$(2,-2)$; exchanging the two complex normal summands is the
orientation-preserving bundle map, with itself as inverse.

The original $II$ scalar action is free.
At a finite $II$ cap a compensating deck power $j$ has surviving
$\ell$ displacement $jn/m$. Its integrality and $\gcd(n,m)=1$
force $m\mid j$; then the scalar compensator is identity.
This proves freeness for every root and affine twist.
The ordinary cap is scalar-free.
For $q=4,8$ the elliptic isotropy is exactly $q$, and its normal
characters are units. The other nodal isotropies are prime to $q$.
At $q=5$ the same assertions follow from the table.
Consequently each nonidentity element of $C_q$ has fixed set
$C\sqcup E_q$, and the complement is $C_q$-free.

\item \emph{The common Borel character and the complete fundamental class.}
The full coefficient splitting at order twenty is as follows.
Put $K_q=K(\mathbb Z/q,2)$ for $q=4,5$.
The cyclic resolution makes the positive homology of
$K(\mathbb Z/q,1)$ vanish over $\mathbb Z[1/q]$.
The path-loop Serre sequence then gives the same statement for
$K_q$. Exactness of localization says that every positive
integral homology element is killed by a power of $q$.
B\'ezout therefore kills all positive tensor products between the
two factors. Also
$\operatorname{Tor}(\mathbb Z/4^j,B_{\rm five})=B_{\rm five}[4^j]=0$ for
five-primary $B_{\rm five}$; finite sums and filtered colimits give all
the required Tor vanishings. Reduced K\"unneth makes
$K_4\wedge K_5$ acyclic. Thus the based inclusion
\[
 i:K_4\vee K_5\longrightarrow K_4\times K_5
\]
is an integral homology equivalence of simply connected CW spaces.
A first nonzero relative homotopy group would, by relative
Hurewicz, give nonzero relative homology; hence all vanish.
Extend the identity of the wedge across successive cells of the
product, using this vanishing and the homotopy-extension property.
It gives a based map
$r_{\rm CW}:K_4\times K_5\longrightarrow K_4\vee K_5$ with
$r_{\rm CW}i=\mathrm{id}$ and
$ir_{\rm CW}\simeq\mathrm{id}$ relative to the wedge.
This is the inverse at the full CW homotopy scope.

The group maps $x\mapsto(x\bmod4,x\bmod5)$ and
$(x_4,x_5)\mapsto5x_4+16x_5$ induce inverse Eilenberg--Mac Lane
maps; in the simplicial abelian-group model the composites are
identity. Combining them with $i,r_{\rm CW}$ and the finite
wedge axiom for the spin Thom spectrum gives the reduced
spin-bordism splitting. Its forward map is the two coefficient
projections; its inverse is the sum of factor inclusions followed
by the inverse coefficient map.
Here $\Omega^{\mathrm{Spin}}_6(\mathrm{pt})=0$ by
\ref{maps:src:D.1.11} at coefficient group zero: its injective
pair $(P,h)$ has zero target.

The class $\mathfrak u_q$ extends to the full-circle Borel space.
Its first possible differential has target containing
$H^1(X;\mathbb Z/q)=0$; its next target is
$H^3(BS^1;\mathbb Z/q)=0$.
Restrict an extension to $C_q$ and call it $U_{\rm Bor}$.
A cusp vertex is fixed by the whole circle. Its section and the
based scalar isotopy give
$\pi_1(W_{C_q})=\pi_1(W)\times C_q$, with the two projections
and their loops as inverse coordinates.
Every fixed-point section has zero free abelian component,
since a homomorphism from a finite group to a free abelian group
is zero. In particular the same cyclic subgroup is used at both curves.

The ordinary restrictions of $U_{\rm Bor}$ and
$\rho_q c_1^{C_q}(\mathcal A)$ agree.
In the universal-cover degree-two sequence their difference comes
from group cohomology because the cover has zero $H^1$.
The ordinary group component vanishes by its injective map to
$H^2(W;\mathbb Z/q)$. K\"unneth leaves one base class $h_{\rm Bor}T_{\rm Bor}$ and
products of integral degree-one classes of $W$ with the degree-one
cyclic class $\theta$; there is no extra Tor term from the free
abelian factor. Every fixed-point section pulls back this same $h_{\rm Bor}$.
Here $T_{\rm Bor}$ is the pullback from $BC_q$ of the modulo-$q$
first Chern class of the weight-one character. With positive curve
generators $x_C,x_E$ on $C,E_q$, respectively, the restrictions are
\[
\begin{aligned}
 U_{\rm Bor}|C&=D_Cx_C+(r_C+h_{\rm Bor})T_{\rm Bor},\\
 U_{\rm Bor}|E_q&=D_Ex_E+h_{\rm Bor}T_{\rm Bor}+\eta\theta,
\end{aligned}
\]
where $\eta$ lifts integrally.
Thus $\eta^2=0$ on the elliptic curve and
$(\eta\theta)^2=0$, also for even $q$ when
$\theta^2=(q/2)T_{\rm Bor}$. The other terms containing $\eta$ have
curve degree different from two and integrate to zero.

The normal splittings give
\[
\begin{aligned}
 e_{C_q}(\nu_C)&=(-x_C-T_{\rm Bor})(-x_C+T_{\rm Bor})=-T_{\rm Bor}^2,\\
 e_{C_q}(\nu_E)&=-t^2T_{\rm Bor}^2.
\end{aligned}
\]
Their coefficients are units in $\mathbb Z/q$.
Remove invariant normal discs; the free complement has a
finite-dimensional quotient, so its Borel cohomology vanishes
after inverting $T_{\rm Bor}$. Thom excision and self-intersection give
the inverse to restriction as the sum of $i_{F*}(z/e_F)$.
The positive even powers of $T_{\rm Bor}$ have coefficient group
$\mathbb Z/q$, so its constant coefficient injects into the
localized ring. Apply this exact inverse and pushforward to $U_{\rm Bor}^3$:
\[
 \langle \mathfrak u_q^3,[X]_q\rangle
 =-3\{D_C(r_C+h_{\rm Bor})^2+D_Eh_{\rm Bor}^2/t^2\}\quad\text{in }\mathbb Z/q.
\]
At $q=4,8$, $2r_C=r_C^2=0$ and the table gives
$at^2-b/q=0$ modulo $q$.
At $q=5$, $r_C=0$ and $at^2-b/5=0$.
Since $D_C=a$ modulo $q$, the constant, linear and quadratic
coefficients in $h_{\rm Bor}$ all vanish. Hence the full cube is zero
at $b=4,8$, and at both coefficient factors of $b=20$.

The coefficient projections at $b=20$ have inverse
$(x_4,x_5)\mapsto5x_4+16x_5\pmod{20}$;
both compositions are identity.
For $3\nmid b$, the integral group
$H_6(K(\mathbb Z/b,2);\mathbb Z)$ is $\mathbb Z/b$, generated
by the image $\gamma_3(s)$ of positive complex $\mathbb{CP}^3$.
The cube evaluates as one on this generator, and therefore is
the full detector, including $b=4,8$;
see \cite[Sections~3.11--3.14, pp.~806--808]{manZhubr75}.
The Ext part of cohomology pairs as zero with integral cycles.
It follows that the complete marked fundamental-class invariant is
$\mu(X)=0$ in every row, for every finite datum in Definition~\ref{maps:def:additive-residual-data}.

\item \emph{The model traces, coprime maps and oriented recognition.}
For $c=4,8$ use the embeddings, framings and trace maps of
\ref{maps:src:D.2.44}. The sphere defining $Z_{c,0}$ is its
$c$-fold tube; the sphere defining $Z_{c,1}$ is its conic-disjoint
projective line tubed to $c/2$ negatively oriented product spheres.
Their rank-four spin framings are identical, with unique homotopy
class because $\pi_2SO(4)=0$. The coefficient maps are
$x\mapsto s$, and $t\mapsto(c/2)s$, $x\mapsto s$, with inverse
$s\mapsto[x]$ on each quotient. These are exactly the earlier
trace maps, so the identity identifies the marked
manifolds. That proof gives simple connectivity, the unique spin
structure, $H_2=H_3=\mathbb Z/c$, zero free third Betti number,
and the full values $(P,\mu)=(0,0)$ or $((c/2)s,0)$.

For $b=20$ the connected sum $Z_{4,\epsilon}\#Z_{5,1}$
has marked group $\mathbb Z/4\langle e_4\rangle\oplus
\mathbb Z/5\langle e_5\rangle$, refined class
$2\epsilon e_4+e_5$, and $\mu=0$.
Its homology marking and inverse are
\[
 \Phi(e_4)=5s,\quad\Phi(e_5)=4s,\qquad
 \Phi^{-1}(xs)=(x\bmod4)e_4+(4x\bmod5)e_5.
\]
Both compositions are identity; $\Phi(P)=4s$ for
$\epsilon=0$ and $14s$ for $\epsilon=1$.
This model marking differs from the coefficient projection above
by the explicit multiplication by four, with inverse four, on
the order-five coordinate.

For the connected sum make the classifying maps constant on
the removed balls and collapse the neck to the basepoint.
The oriented spin one-handle bordism from the disjoint union
extends these maps. Its second-homology coordinates are the
summand inclusions, with inverse the two collapse maps, and its
bordism class is their sum. Composing with $\Phi$ gives the
same marked $P,\mu$ as $X$.

Finally, \cite[Theorem~5, p.~796]{manZhubr75} gives the injection
of marked spin bordism by $(P,h)$.
The construction and marked uniqueness in
\cite[Theorem~3, p.~795; Section~4.1, pp.~809--810]{manZhubr75}
preserve the specified $H_2$ and remove the free part of $H_3$.
They apply here with both free ranks zero and the unique spin
structure. Equality of the full invariants established above
therefore gives the oriented diffeomorphisms in
\eqref{maps:eq:D45.3}, inducing the displayed markings.
Their inverses induce the corresponding inverse group maps.
The two alternatives at each residual order are separated by
the order of $P$: one and two at orders four and eight, and
five and ten at order twenty. Multiplication by a unit, with
inverse the inverse unit, preserves these orders.
Each original pattern is realized with $e=1$; each finite pattern
is realized with $m=6$, $e=0$, $n=1$, $r=0$.
The unbounded relation maps, framing construction and full-coefficient
calculation prove the same conclusions for every datum in
\eqref{maps:eq:D45.2}.
\end{enumerate}
\end{proof}


\section{Weighted Engel and OS45 manifolds}
\subsection{Weighted Engel smooth types}

\begin{definition}[Weighted Engel data]\hypertarget{tag-Q250}{}\label{maps:def:weighted-engel-data}
{\normalfont\scriptsize tag Q250}\par
For every positive integer $d$, use the weight-$d$ Engel construction
of \ref{maps:src:D.2.34} and its prescribed cusp fan. Replace both
additive sites by good logarithmic caps, allowing every positive
integer $k,\ell$, every integer $r_i,n_i$, and every affine twist
in that construction, subject to
\begin{equation}
 m_0=3k,\qquad m_1=4\ell,\qquad
 v_i=r_i\delta+n_i u_i,\qquad
 P=m_1n_0+m_0n_1=\pm1 .
 \hypertarget{tag-0272}{}\label{maps:eq:D.WI1}
\end{equation}
Here $u_i$ is the marked fixed direction at site $i$ of that construction.
Define $M_d^{\rm w}$ to be the member with $k=\ell=1$, $n_0=1$, $n_1=-1$
and $r_0=r_1=0$, with the complex orientation. Its determinant is one.
These are the \emph{weighted Engel data}.
\end{definition}

\begin{theorem}[Weighted Engel smooth classification]\hypertarget{tag-02LP}{}\label{maps:src:D.2.48}
{\normalfont\scriptsize tag 02LP}\par
Every threefold $X$ of Definition~\ref{maps:def:weighted-engel-data}
is simply connected. Its complex first and second
Chern classes, spin characteristic class, and refined class satisfy
\begin{equation}
 c_1(X)=c_2(X)=\lambda(X)=0,\qquad P_{\mathrm{ref}}(X)=0.
 \hypertarget{tag-027D}{}\label{maps:eq:D.WI2}
\end{equation}
Here $P_{\mathrm{ref}}\in H_2(X;\mathbb Z)$ is Zhubr's refined
class, and $P$ in \eqref{maps:eq:D.WI1} is the Seifert determinant.
Write $\mathcal D_X:H^4(X;\mathbb Z)\to H_2(X;\mathbb Z)$ for
$c\mapsto c\frown[X]$, with inverse Poincar\'e duality, and put
$\mathfrak p_X=\mathcal D_X^{-1}(P_{\mathrm{ref}})$.
Thus $\mathfrak p_X\frown[X]=P_{\mathrm{ref}}$.

The classes $x_i=\operatorname{PD}(D_i)$ of the $d$ cusp components,
indexed cyclically, generate $H^2(X;\mathbb Z)$ with sole relation
$\sum_i x_i=0$.  The complex ruling $f=H-E_1$ has image
$s=[a\wedge b]$, a generator of $\operatorname{Tor}H_2(X)=\mathbb Z/d$.
Put $\kappa=\operatorname{PD}(f)$.  Counting both directed incidences when
$d=2$, the complete degree-two product is
\begin{equation}
 x_i x_j=
 \bigl(\mathbf1_{j=i+1}+\mathbf1_{j=i-1}-2\mathbf1_{j=i}\bigr)\kappa.
 \hypertarget{tag-027H}{}\label{maps:eq:D.WI3}
\end{equation}
For $d=1$ the sole component class is zero.  For $d=2$ the two
off-diagonal incidences both occur, and every product in \eqref{maps:eq:D.WI3}
is zero since $2\kappa=0$.  The integral cubic is zero for every $d$.
The complete fundamental-class invariant includes the following
two- and three-primary coordinates. Put
$r=d-1$, $D_2=d\gcd(d,2)$, and $D_3=d\gcd(d,3)$.
For $d>1$ there is a class $u\in H^2(X;\mathbb Z/d)$
which, together with $x_0,\ldots,x_{d-2}$, fixes the splitting
$H_2(X)=\mathbb Z^r\oplus\mathbb Z/d$ described below. In
\begin{equation}
 H_6(K(\mathbb Z^r\oplus\mathbb Z/d,2);\mathbb Z)
 \simeq
 \mathbb Z^{\binom{r+2}{3}}\oplus
 (\mathbb Z/d)^{\binom{r+1}{2}}\oplus
 (\mathbb Z/D_2)^r\oplus\mathbb Z/D_3
 \hypertarget{tag-0277}{}\label{maps:eq:D.WI14}
\end{equation}
the complete class $\mu$ is determined by
\begin{equation}
 \begin{aligned}
 \langle x_i x_j x_k,[X]\rangle&=0,\\
 \langle u x_i x_j,[X]\rangle
    &=\mathbf1_{j=i+1}+\mathbf1_{j=i-1}-2\mathbf1_{j=i}
                                                     &&\pmod d,\\
 \langle \mathcal Q_2(u)x_i,[X]\rangle&=-12             &&\pmod {D_2},\\
 \langle \mathcal Q_3(u),[X]\rangle&=0                 &&\pmod {D_3}.
 \end{aligned}
 \hypertarget{tag-0278}{}\label{maps:eq:D.WI15}
\end{equation}
The indices in the second line are cyclic modulo $d$; those
used as independent coordinates run from zero to $d-2$.
Here $\mathcal Q_2$ is the ordinary square if $d$ is odd and
the Pontryagin square with values modulo $2d$ if $d$ is even.
Likewise $\mathcal Q_3$ is the ordinary cube if $3\nmid d$
and Thomas's Pontryagin third power with values modulo $3d$
if $3\mid d$.

Consequently, for every $d>0$ all the simply connected data
\eqref{maps:eq:D.WI1} give one oriented diffeomorphism type $M_d^{\rm w}$, namely
the unique type with
\begin{equation}
 H_2=\mathbb Z^{d-1}\oplus\mathbb Z/d,\qquad b_3=0,\qquad
 P_{\rm ref}=0,\qquad \mu\ \text{given by \eqref{maps:eq:D.WI15}}.
 \hypertarget{tag-0279}{}\label{maps:eq:D.WI16}
\end{equation}
One has $M_1^{\rm w}=S^6$ and, for $d$ prime to six,
$M_d^{\rm w}\cong M_d$ with the complex orientations and the
marking proved below; here $M_d$ is the model of
\cref{maps:src:D.1.11}. The types for different
$d$ are distinct, since their Euler numbers are $2d$.

\end{theorem}

\begin{proof}
\begin{enumerate}[leftmargin=*]
\item \emph{The marked groups and primitive component classes.}
Use the physical lattice $(\delta,a,b,\gamma)$ and the alternating form
$\xi=a^*\wedge b^*+6d\gamma^*\wedge\delta^*$.
Put $s=[a\wedge b]$. The full normal maps of
\ref{maps:src:D.2.34} give, in coordinates $(s,t,t'_0,t'_1)$,
the six relation columns
\[
\begin{aligned}
C^{\rm nor}_{01}&=(-2dn_0,-k,0,0),&
C^{\rm nor}_{02}&=(-6dr_0,0,m_0,0),&
C^{\rm nor}_{03}&=(0,r_0,n_0,0),\\
C^{\rm nor}_{11}&=(-3dn_1,2\ell,0,0),&
C^{\rm nor}_{12}&=(-6dr_1,0,0,m_1),&
C^{\rm nor}_{13}&=(0,-r_1,0,n_1).
\end{aligned}
\]
Every common divisor of $m_i,n_i$ or of $m_0,m_1$ divides $P$.
Thus these pairs are coprime and $\gcd(k,2\ell)=1$.
The surviving Seifert relation matrix and its integral inverse are
\[
 R=\begin{pmatrix}-n_0&m_0\\-n_1&-m_1\end{pmatrix},
 \qquad R^{-1}=P\begin{pmatrix}-m_1&-m_0\\n_1&-n_0\end{pmatrix}.
\]
They eliminate every loop, proving simple connectivity.

Let $B_{\rm nor}$ be the six-column normal matrix and $e_{ij}$ its source basis.
The following formulas give the full inverse quotient marking.
Set $R_d=-P(2\ell e_{01}+k e_{11})$, so $B_{\rm nor}R_d=ds$.
Choose $a_{\rm nor},b_{\rm nor}$ with $-ka_{\rm nor}+2\ell b_{\rm nor}=1$, and put
$q_{\rm nor}=2a_{\rm nor}n_0+3b_{\rm nor}n_1$ and $R_t=a_{\rm nor}e_{01}+b_{\rm nor}e_{11}+q_{\rm nor}R_d$.
Then $B_{\rm nor}R_t=t$. Choose $\alpha_i m_i+\eta_i n_i=1$ and put
$C_i'=\alpha_i e_{i2}+\eta_i e_{i3}$. The remaining inverse columns are
\[
 R_{t_0}=C_0'+6r_0\alpha_0R_d-r_0\eta_0R_t,\qquad
 R_{t_1}=C_1'+6r_1\alpha_1R_d+r_1\eta_1R_t,
\]
with images $t'_0,t'_1$. Since every first coordinate of $B_{\rm nor}$ is
divisible by $d$, the quotient and its inverse are
\[
 \operatorname{coker}B_{\rm nor}\longrightarrow\mathbb Z/d,\quad
 [(z_s,z_t,z_0,z_1)]\longmapsto[z_s],
 \qquad [j]\longmapsto[js].
\]
For a kernel vector the exact preimage is
$(z_s/d)R_d+z_tR_t+z_0R_{t_0}+z_1R_{t_1}$.

Let $E$ be the exterior after filling the two finite caps and before
adjoining the cusp. Its first homology is
\[
 \langle\gamma,x'_0,x'_1\mid m_0x'_0=n_0\gamma,\
                                  m_1x'_1=n_1\gamma\rangle.
\]
The quotient map sends $(z_\gamma,z_0,z_1)$ to
$m_0m_1z_\gamma+n_0m_1z_0+m_0n_1z_1$.
Its inverse sends $1$ to $P(x'_0+x'_1)$: the two relation columns
together with this vector have determinant $P^2=1$.
Thus the cusp meridian $x'_0+x'_1$ has primitive value $P$.
The normal-disc differences in the cusp give a section of
the relative augmentation map $\mathbb Z^d\to\mathbb Z$.
Together with the quotient above they give
\[
 H_2(X)=\mathbb Z^{d-1}\oplus(\mathbb Z/d)\langle s\rangle,
 \quad H_3(X)=\mathbb Z/d,\quad H_1=H_5=0,\quad b_3=0.
\]
The maps include all cusp boundary torsion and hold for every
$k,\ell,n_i,r_i$ in \eqref{maps:eq:D.WI1}.

The dual relative Thom map is therefore
\begin{equation}
 H^1(E)\longrightarrow H^2(X,E)=\mathbb Z^d,\qquad
                   1\longmapsto\pm(1,\ldots,1).
 \hypertarget{tag-027I}{}\label{maps:eq:D.WI4}
\end{equation}
Its cokernel injects into $H^2(X)$ with rank-zero quotient inside
the torsion-free group $H^2(E)$. Hence it is all of $H^2(X)$.
The classes are the Cartier divisors $x_i=\operatorname{PD}(D_i)$;
at a self-identification the local equation is the product of the
branch coordinates. The inverse reads evaluation on the normal-disc
differences. Their sole relation is $\sum_i x_i=0$, so
$f_X^*[\mathrm{point}]=0$. Each reduced finite multiple-fibre divisor
also has zero class, since its positive multiple is this fibre class
and $H^2(X;\mathbb Z)$ is free.

\item \emph{The ruling and the divisor products.}
In short coordinates $(a,\delta)$ and long coordinates $(b,\gamma)$,
the valuation sends $b$ to $e_1$ and $\gamma$ to $de_2$.
The prescribed rays
$e_1,e_2,e_2-e_1,-e_1,-e_2,e_1-e_2$ have boundary curves
\[
 H-E_A-E_C,\quad E_A,\quad H-E_A-E_B,\quad
 E_B,\quad H-E_B-E_C,\quad E_C.
\]
The notation in the statement is
$E_1=E_B,E_2=E_C,E_3=E_A$.
In the toric chain construction of \ref{maps:src:D.1.9}, substitute
$(\widehat u,\widehat w,\widehat\delta,\widehat\gamma)
=(b,a,\delta,-\gamma_1)$, with inverse
$(a,b,\delta,\gamma_1)
=(\widehat w,\widehat u,\widehat\delta,-\widehat\gamma)$.
Apply its normal-disc and period-cylinder maps to the
holomorphic period functions of this construction, including all units.
The ruling has relative vector $n_{2e}-2n_e+n_0$, and the diagonal
has $n_{e_1+e_2}-n_{e_1}-n_{e_2}+n_0$.
Injectivity into the toric relative group, followed by cancellation
of translated normal discs, gives the chains
\[
 E_A+E_C=-T(b,a)=T(a,b),\quad
 E_B+E_C=T(-\gamma_1,\delta)=T(\delta,\gamma_1),\quad
 E_C=-T(b,\delta).
\]
In particular the positive ruling $f=H-E_B$ has in the period
quotient the same image as $E_A+E_C$, and maps to $s$, with sign fixed by
$T(b,a)=-T(a,b)$. Its $\xi$ value is one; $E_C$ has value zero.
The inverse marking reads the same short and long coordinates.
This is a unit-period chain comparison and applies in every
$\gamma$ component of the degree-$d$ period cover.

The two branch curves at each neighbouring component intersection
sum to $f$. The two directed incidences remain distinct when $d=2$.
Multiplying $\sum_i x_i=0$ by $x_j$ fixes the diagonal coefficient,
giving \eqref{maps:eq:D.WI3}. At $d=1$ the only class is zero;
at $d=2$ both off-diagonal coefficients are two and $2\kappa=0$.
Every integral cubic is zero, since its degree-four product is
torsion and $H^6(X;\mathbb Z)=\mathbb Z$.

\item \emph{The tangent classes and normalized components.}
The normalization immersion $j_i:dP_6\longrightarrow X$ has
normal line $K_{dP_6}$.  Indeed the principal local central
divisor is the product of the normal coordinates to all its
branches; on a fixed normalized branch its normal transition
is the inverse of the six toric boundary transitions.  This
includes both conductor branches of each self-identification.
Consequently
\begin{equation}
 j_i^*c_1(TX)=c_1(TdP_6)+c_1(K_{dP_6})=0,\qquad
 \langle j_i^*c_2(TX),[dP_6]\rangle=6-6=0 .
 \hypertarget{tag-027J}{}\label{maps:eq:D.WI5}
\end{equation}
Restriction of $H^2(X)$ to the normalized components is injective.
For a component combination with coefficients $u_i$, its
restriction to component $i$ is
$\sum_r(u_{i+r}-u_i)C_r$.  Pairing with the anticanonical ample
class, which has degree one on each of the six boundary curves,
gives the component-graph Laplacian.  The graph is connected, so
zero restrictions force the $u_i$ to be constant.  That is
precisely the already removed principal relation.  Equation
\eqref{maps:eq:D.WI5} now gives $c_1(TX)=0$ integrally.

\item \emph{The corrected differential and reference frame.}
Remove the divisorial multiple-fibre factors from $df_X$.  Its
target is the line
\[
 L=f_X^*T\mathbb P^1\otimes
                      \mathcal O\bigl(-\sum (m_i-1)D_i^{\rm fin}\bigr).
\]
It is topologically trivial by \eqref{maps:eq:D.WI4}.  The abstract rank-two
kernel $V$ is locally free, including at $xy$ and $xyz$ charts,
and has the nowhere-zero abstract section given by $\delta$.
At a node the kernel columns are $(x,-y,0)$ and $(0,0,1)$;
at a triple point they are $(x,-y,0)$ and $(x,0,-z)$.
Indeed $yzA+xzB+xyC=0$ implies $A=xA'$, $B=yB'$, $C=zC'$
by reduction modulo each coordinate in the convergent power-series
domain, and then $A'+B'+C'=0$. These are exactly the two displayed
columns, with inverse coefficient extraction. The node calculation
is the same divisibility argument. The constant scalar-circle weight vector is a nonzero constant
combination in these abstract bases.  At the good finite caps
the corrected differential is a submersion and $\delta$ remains
a nonzero tangent vector.  Since
$c_1(V)=c_1(TX)-c_1(L)=0$, its quotient by this section is a
trivial complex line.  Thus $V$ and $L$ admit full smooth frames.
A Hermitian complement to
$\delta$ and a unit section give a global frame $(\delta,v_2)$ of $V$.
Outside the critical cycle, a unit frame of $L$ lifts through
$df^*/\|df\|^2$; kernel projection and $df$ give its inverse.
They give the reference frame $\rho$ on the complement of the
critical curve cycle $C$.

\item \emph{The complete critical cycle.}
The local Bott clutch at $xy$ has coefficient one:
its two normal columns are
$\bigl(\begin{smallmatrix}x&\overline y\\-y&\overline x\end{smallmatrix}\bigr)$.
The $xyz$ Hilbert--Burch minors give multiplicity one on each
of its three axes and no further codimension-two contribution.
The good finite logarithmic caps have no residual critical
cycle after correcting the target.  Hence, with the complex
orientation and the spin normalization of \ref{maps:src:D.2.44},
\begin{equation}
                 \lambda(TX,\rho)=-\operatorname{PD}(C).
 \hypertarget{tag-027K}{}\label{maps:eq:D.WI6}
\end{equation}

The cusp period covering computes $[C]$.
In the positive period basis
$N_{\rm mon} b=a$, $N_{\rm mon}\gamma_1=\delta$, the three chain equalities above identify the weight-one cusp's
three-curve sum with
$a\wedge b+b\wedge\delta+\delta\wedge\gamma_1$.
The weight-$d$ cusp is its degree-$d$ \'etale period cover, with
$\gamma_d=d\gamma_1$.  Transfer therefore gives the preimage of its full $3d$-curve sum as
\begin{equation}
 d\,a\wedge b+d\,b\wedge\delta+\delta\wedge\gamma_d.
 \hypertarget{tag-027L}{}\label{maps:eq:D.WI7}
\end{equation}
The scalar-circle parameter is the angular vector $\delta$
fixed by $N_{\rm mon}\gamma_d=d\delta$.
Since $\xi_d$ has values $1,0,-6d$ on these three bivectors,
\eqref{maps:eq:D.WI7} has image $-5d\,s=0$.  Thus $C$ is integrally
nullhomologous and $c_2(TX)=\lambda(TX)=0$.

\item \emph{One integral auxiliary line and all period maps.}
Let $W$ consist of the regular core and the cusp, leaving both finite
discs unfilled. Its group and the exponent marking are
\[
 \pi_1(W)=\langle\gamma,x_{\rm bas}\mid[\gamma,x_{\rm bas}]=1\rangle\simeq\mathbb Z^2,
 \qquad (u,v)\longmapsto\gamma^ux_{\rm bas}^v.
\]
The inverse reads the two exponents in this presentation.
For the core monodromies in the basis $(\delta,a,b,\gamma)$,
\[
A_0=\begin{pmatrix}1&-1&0&2d\\0&-1&-1&6d\\0&1&0&0\\0&0&0&1\end{pmatrix},
\quad
A_1=\begin{pmatrix}1&0&1&-3d\\0&0&1&0\\0&-1&0&6d\\0&0&0&1\end{pmatrix},
\quad A_\infty=I+E_{23}+dE_{14},
\]
the quadratic refinement $q_\xi(\lambda)=\lambda_a\lambda_b+
\lambda_a+\lambda_b\bmod2$ is invariant. Its polarization is
$\xi\bmod2$; direct substitution in both displayed matrices proves
invariance, with every additional term even as an integer-valued
polynomial. Consequently
\[
 J_\lambda(x_{\rm per},z)=
 (x_{\rm per}+\lambda,(-1)^{q_\xi(\lambda)}e^{\pi i\xi(\lambda,x_{\rm per})}z)
\]
satisfies the exact lattice cocycle identity.
The connection $d-\pi i\xi(x_{\rm per},d x_{\rm per})$ has Chern class $\xi$.
The scalar-one lifts $(x_{\rm per},z)\mapsto(A_i x_{\rm per},z)$ preserve it and have
inverse $A_i^{-1}$. They fix the origin frame, so all marked
fixed-loop sweep values are zero. This constructs the integral
core class at every weight.

For the cusp write the real coordinates as $x_{\rm per}=(a,\delta)u_{\rm per}+
(b,\gamma)v_{\rm per}$ and the complex coordinates as $Z=u_{\rm per}+\Omega v_{\rm per}$, with
\[
 K_{\rm pol}=\operatorname{diag}(1,-6d),\qquad
 \Omega=\frac{\log q_{\rm base}}{2\pi i}\operatorname{diag}(1,d)
                  +B_{\rm hol}(q_{\rm base}),\qquad K_{\rm pol}\Omega=(K_{\rm pol}\Omega)^t.
\]
The holomorphic units are included in $B_{\rm hol}$.
Multiply the line coordinate by
\[
 G(u_{\rm per},v_{\rm per})=
 \exp\bigl(-\pi i(u_{\rm per}^tK_{\rm pol}v_{\rm per}+v_{\rm per}^tK_{\rm pol}\Omega v_{\rm per})+\pi iZ_1\bigr);
\]
its reciprocal is the inverse. Short factors become trivial, and
the long $b,\gamma$ factors are
\[
 -z_1^{-1},\qquad
 \exp\bigl(\pi i(6d\Omega_{22}+\Omega_{12})\bigr)z_2^{6d}.
\]
The second has base order $3d^2$ and a nonzero holomorphic unit.
The relation $\Omega_{12}=-6d\Omega_{21}$ verifies commutation.
Replacing $\Omega$ by $\Omega+\operatorname{diag}(1,d)$ and
$u_{\rm per}$ by $u_{\rm per}+\operatorname{diag}(1,d)v_{\rm per}$ gives the same Gaussian
comparison, proving the base return and its inverse.

Use $\phi(r_{\rm tor},s_{\rm tor})=-r_{\rm tor}(r_{\rm tor}-1)/2+3s_{\rm tor}^2$. Its period differences are
$\phi(r_{\rm tor}+1,s_{\rm tor})-\phi(r_{\rm tor},s_{\rm tor})=-r_{\rm tor}$ and
$\phi(r_{\rm tor},s_{\rm tor}+d)-\phi(r_{\rm tor},s_{\rm tor})=6d s_{\rm tor}+3d^2$, exactly the two factor valuations.
Interpolate on the given unimodular triangles. The Cartier monomials
and their reciprocals give regular local frames and inverse maps;
the holomorphic period unit is kept in the long identification.
The equivalent support used below is
\begin{equation}
 \begin{aligned}
 \psi(r_{\rm tor},s_{\rm tor})&=-\frac{r_{\rm tor}(r_{\rm tor}-1)}2+3s_{\rm tor}(s_{\rm tor}-1)\\
 &=\phi(r_{\rm tor},s_{\rm tor})-3s_{\rm tor}.
 \end{aligned}
 \hypertarget{tag-027A}{}\label{maps:eq:D.WI17}
\end{equation}
The exact line map is $\zeta_\psi=z_2^{-3}\zeta_\phi$, with inverse
$\zeta_\phi=z_2^3\zeta_\psi$. If a period multiplies $z_2$ by
$q_{\rm base}^{\lambda_2}u_\lambda(q_{\rm base})$, the line factor
changes by $q_{\rm base}^{-3\lambda_2}u_\lambda^{-3}$.
For Cartier slopes $m_\psi=m_\phi-3e_2$, this map is the identity
between the regular frames $\chi^{-m_\phi}\zeta_\phi$ and
$\chi^{-m_\psi}\zeta_\psi$. It therefore extends through every chart
and collar, including every period unit.

The weight-$d$ core and cusp are the period cover
$\gamma_d=d\gamma_1$. Pullback preserves the first three core graph
generators and multiplies the two fixed long sweeps by $d$.
The weight-one core class has fibre value $\xi_1$ and zero on all
five graph generators, and its cusp boundary class is the line just
constructed: the two angular Wang sweeps bound in the cusp and
$H_1$ of that boundary is free. Pullback therefore identifies the
whole integral boundary line at weight $d$, including its torsion
cohomology class. This also identifies the direct weight-$d$ cocycle
construction above with that pullback. The cover is used on these
core and cusp pieces.
Denote the first Chern class of this line on $W$ by $\alpha_d$.

\item \emph{Extension across the finite caps.}
At a finite cap let $\pi:Y_i\longrightarrow Q_i$ be its
normal circle bundle. The obstruction for
$\alpha_d\bmod d$ to extend across the disk bundle is its
Gysin image
\begin{equation}
                  \pi_!(\alpha_d\bmod d)\in
                         H^1(Q_i;\mathbb Z/d).
 \hypertarget{tag-027B}{}\label{maps:eq:D.WI18}
\end{equation}
Its values on the complete marked generators
$a,\gamma,g$ of $H_1(Q_i)$ are precisely the first-row
coefficients of \eqref{maps:eq:D.WE3}:
\begin{equation}
 \begin{array}{c|ccc}
       &a&\gamma&g\\ \hline
 i=0&-2dn_0&-6dr_0&0\\
 i=1&-3dn_1&-6dr_1&0.
 \end{array}
 \hypertarget{tag-027C}{}\label{maps:eq:D.WI19}
\end{equation}
The zero $g$ column is part of the Thom
definition of fibre integration: a loop in $Q_i$ sweeps its
normal circle, and pairing with $\alpha_d$ gives the displayed
normal-column coefficient. Every coefficient is zero
modulo $d$. Since $H^1(Q_i;\mathbb Z/d)$ is
$\operatorname{Hom}(H_1(Q_i),\mathbb Z/d)$, \eqref{maps:eq:D.WI18} vanishes,
even when the integral normal line has torsion. The circle
Gysin sequence therefore supplies an extension over each
finite cap. Mayer--Vietoris now glues them with the pulled-back
cusp line to a class $u\in H^2(X;\mathbb Z/d)$.
For an exact global extension map, the two relative degree-two
columns into $H_1(W)$ are the Seifert matrix $R$ above, whose inverse
is integral. Hence $H_2(W)\to H_2(X)$ is onto and its kernel is the
complete normal Thom image. The functional given by the line modulo
$d$ annihilates it by \eqref{maps:eq:D.WI19}; it therefore descends
uniquely. The universal-coefficient isomorphism for $H_1(X)=0$
identifies this functional with the same glued class $u$.

\item \emph{Component restrictions and the homology splitting.}
On every normalized cusp component the integral representative
of its restriction is
\begin{equation}
              \eta=-5H+6E_1-E_3,\qquad
       \eta(E_1,E_2,E_3)=(-6,0,1),\qquad
       \eta(f)=1,\qquad \eta^2=-12.
 \hypertarget{tag-027E}{}\label{maps:eq:D.WI20}
\end{equation}
The three degrees follow directly by taking the differences
of \eqref{maps:eq:D.WI17} across the internal, diagonal, and scalar-fixed
walls, as in \ref{maps:src:D.1.9}. In particular $u(s)=1$.
Choose free lifts in $H_2(X)$ annihilated by $u$, subtracting
appropriate multiples of $s$. This specifies the splitting
in \eqref{maps:eq:D.WI14}--\eqref{maps:eq:D.WI16}. The integral classes $x_i$ form its
dual free basis, and $u$ is the standard cyclic coordinate.

The six wall differences for $\phi$ are $(-6,1,0,-6,1,0)$.
In the $(H,E_A,E_B,E_C)$ marking they give degrees $(-5,1,-6,0)$,
so the displayed $\eta$ is exactly $-5H-E_A+6E_B$.
The inverse of the adjusted homology splitting reads the evaluations
of $x_0,\ldots,x_{d-2}$ and then $u$. Changing an integer representative
of a $u$-value changes the adjustment by $ds=0$, so both maps are
well defined and mutually inverse.

\item \emph{The complete integral target and inverse detectors.}
Put $B_d=K(\mathbb Z/d,2)$. Its marked low-degree homology is
\begin{equation}
 H_2(B_d;\mathbb Z)=\mathbb Z/d,\quad
 H_4(B_d;\mathbb Z)=\mathbb Z/D_2,\quad
 H_6(B_d;\mathbb Z)=\mathbb Z/D_3 .
 \hypertarget{tag-027G}{}\label{maps:eq:D.WI22}
\end{equation}
Here the positive reductions of $\mathbb{CP}^j$ give the generators
$\gamma_j(g)$, with $g$ the Hurewicz image of $1\bmod d$.
In Cartan's integral model for $d=p^f$ the basic pair is
$\Gamma(x_2)\otimes\Lambda(y_3)$ with $dy=p^fx$.
The differential on $y\gamma_{j-1}(x)$ is
$jp^f\gamma_j(x)$. Its $p$-primary homology in degree $2j$ has
order $p^{f+v_p(j)}$. The first further pair has degrees
$2p+1,2p+2$ and differential $dt=pz$.
For $p=2$ it contributes $\mathbb Z/2$ in degree five and
zero in degree six; for $p\ge3$ it starts above degree six.
Its products and torsion products with the basic positive classes
start in higher degrees. These are the complete low-degree groups
of the integral model
\cite[Expos\'e~11, Theorems~1 and~4, pp.~11-09--11-13; Appendix~1, pp.~11-22--11-23]{manCartan1955EML}.
Distinct prime components have relatively prime positive homology
orders, so their tensor and torsion products vanish. Write
$d=\prod_{p\mid d}p^{f_p}$ and, for $j=1,2,3$, put
$D_{j,p}=p^{f_p+v_p(j)}$ and $D_1=d$.
The coefficient reduction sends $g$ to $g_p$ and hence
$\gamma_j(g)$ to $\gamma_j(g_p)$, the positive reduced
$\mathbb{CP}^j$ generator in that factor. Choose
$c_{j,p}(D_j/D_{j,p})\equiv1\pmod{D_{j,p}}$.
The complete coordinate maps and inverse are
\[
 z\longmapsto(z\bmod D_{j,p})_{p\mid d},\qquad
 (z_p)_p\longmapsto
 \sum_{p\mid d}(D_j/D_{j,p})c_{j,p}z_p\pmod{D_j}.
\]
Each summand is one on its own factor and zero on the others,
so both composites are identity. At $j=1$ these are the underlying
coefficient maps; their induced Eilenberg--Mac Lane maps are inverse.
This proves \eqref{maps:eq:D.WI22} for all $d$; at $d=1$ each
group is zero. The argument retains the positive divided-power
generator at each stage.

The quadratic generator has relations $d^2q_{\rm quad}(1)=2dq_{\rm quad}(1)=0$.
Conversely $q_{\rm quad}(j)=j^2\bmod D_2$ is well defined and quadratic,
giving the inverse detector and proving these are its complete
relations. For the degree-six generator $z$, the maps with
integral classes $dh$ on $\mathbb{CP}^3$ and $dx+y$ on
$\mathbb{CP}^1\times\mathbb{CP}^2$ push forward to $d^3z$ and
$3dz$. Their reductions factor through a point or
$\mathbb{CP}^2$, respectively, so both images vanish.
Thus its order divides $\gcd(d^3,3d)=D_3$.
The ordinary cube or Thomas operation evaluates to one on the
positive reduced hyperplane class, giving the opposite bound
\cite[equations~(2.8), (3.4), pp.~267--268]{ffThomas56}.
Divided-power functoriality identifies
$\gamma_j(g)$ with the positive $\mathbb{CP}^j$ pushforward:
addition on $\mathbb{CP}^\infty$ gives
$[\mathbb{CP}^i]*[\mathbb{CP}^j]=\binom{i+j}{i}[\mathbb{CP}^{i+j}]$,
and its free homology permits cancellation of $j!$ before applying
integer reduction. The detector and inverse are therefore
$n\gamma_j(g)\mapsto n$ and $n\mapsto n\gamma_j(g)$.

The free factor in
$K(\mathbb Z^{d-1}\oplus\mathbb Z/d,2)
=(\mathbb{CP}^\infty)^{d-1}\times B_d$
has free even homology and no odd homology. K\"unneth gives
\eqref{maps:eq:D.WI14}, with inverse coordinate map
\[
 (C_\alpha,J_\beta,K_i,L)\longmapsto
 \sum C_\alpha\gamma_\alpha(e)+
 \sum J_\beta\gamma_\beta(e)g+
 \sum K_i e_i\gamma_2(g)+L\gamma_3(g).
\]
The forward evaluations are respectively the free cubic monomials,
$x^\beta u$, $x_i\mathcal Q_2(u)$ and $\mathcal Q_3(u)$.
They pair to one with their own displayed generator and to zero
with the others. The possible degree-six cohomology
$\operatorname{Ext}$ ambiguity from $H_5(B_d)$ pairs to zero
with integral six-cycles. Thus the evaluations in \eqref{maps:eq:D.WI15} determine the
entire fundamental-class invariant, including its two- and
three-primary coordinates. We compute these evaluations next.

\item \emph{The full mixed coordinates.}
The first mixed value in \eqref{maps:eq:D.WI15} follows from \eqref{maps:eq:D.WI3} and
$\langle u\kappa,[X]\rangle=1$. The Pontryagin square of the
reduction of an integral degree-two class is its ordinary
square with the full target modulus $2d$; this is the
$p=2$ case of the cochain construction
\cite[equations~(2.1), (2.8) and~(3.4), pp.~267--268]{ffThomas56}.
By naturality and the projection formula for the normalized
component fundamental class,
\begin{equation}
 \langle \mathcal Q_2(u)x_i,[X]\rangle
       =\langle\mathcal Q_2(j_i^*u),[dP_6]\rangle
       =\eta^2=-12\pmod {D_2}.
 \hypertarget{tag-027F}{}\label{maps:eq:D.WI21}
\end{equation}
This computes the coefficient modulo $2d$ when $d$ is even;
reduction of this target modulo $d$ has its order-two kernel. The relation
$\sum_i x_i=0$ is consistent with \eqref{maps:eq:D.WI21}, since
$12d$ is divisible by $D_2$.

\item \emph{The finite lift and the common fixed-point characters.}
For every $N\mid6d$ the magnetic translation by $\delta/N$ is
\[
 (x_{\rm per},z)\longmapsto
 (x_{\rm per}+\delta/N,e^{\pi i\xi(\delta/N,x_{\rm per})}z).
\]
Its commutator with $J_\lambda$ is
$\exp(-2\pi i\,6d\lambda_\gamma/N)=1$, and its $N$th power is
$J_\delta$. It commutes with all return lifts, since they fix
$\delta$ and preserve $\xi$. In the $\phi$ Gaussian coordinate
it acts by $(z_1,z_2,\zeta)\mapsto(z_1,e^{2\pi i/N}z_2,\zeta)$.
In the $\psi$ coordinate the dense line character is $-3$.
Thus the global finite class is represented on each component by
\begin{equation}
 j_i^*u=\rho_d(-5H+6E_1-E_3),\qquad
 \psi(r_{\rm tor},s_{\rm tor})=\phi(r_{\rm tor},s_{\rm tor})-3s_{\rm tor}.
 \hypertarget{tag-025R}{}\label{maps:eq:D.TP1}
\end{equation}
The degree on every vertical fixed curve is one.
The $\phi$ slope on the edge $(0,j)$ to $(0,j+1)$ is $6j+3$,
so the regular line character is $-6j-3$.
The period $j\mapsto j+d$ changes it by $-6d$, zero modulo $N$.
The reciprocal gauges give the inverse identifications.
This proves finite lifts for $N=d,2d,3d$ whenever those groups
are used.

\item \emph{The fixed curves and the free complement.}
The scalar circle action is semifree. On a finite cap of
order $m_i$ its gamma translation has primitive numerator
$n_i$, with $\gcd(m_i,n_i)=1$. A nontrivial root power cannot
be compensated by the scalar direction, because that would
require $j n_i/m_i$ integral. Thus these caps are scalar-free.
The cusp has exactly the $d$ fixed curves $C_j$; their normal
bundles are $\mathcal O(-1)\oplus\mathcal O(-1)$ with weights
$(1,-1)$, and the other cusp orbits are free. These are the
marked geometric inputs of the weighted family.

On the given cusp fan every nonvertical edge has determinant one
with $\delta$; the vertical edges give exactly the $d$ curves.
The adjacent charts have overlap $(x',y',w')=(wx,wy,w^{-1})$,
with inverse $(x,y,w)=(w'x',w'y',(w')^{-1})$.
It gives the two $\mathcal O(-1)$ normal transitions and weights
$1,-1$. Thus every nonidentity subgroup element has precisely
these fixed curves and a free complement.
At $d=1$, the groups $H_2,H_3$ and all coefficient coordinates
under consideration are zero directly. Subsequent finite
localization is used for nontrivial groups.

\item \emph{The common Borel character.}
Since $H_1(X;\mathbb Z/d)=0$ and $H^3(BS^1;\mathbb Z/d)=0$,
the degree-two Serre sequence gives a full-circle Borel lift
$\widehat u$ of $u$. Its restriction is
\begin{equation}
 \widehat u|_{C_j}=x_C+\lambda_j t,\qquad
 \lambda_j=h-6j-3\quad\text{in }\mathbb Z/d,
 \quad \langle x_C,[C_j]\rangle=1.
 \hypertarget{tag-025S}{}\label{maps:eq:D.TP2}
\end{equation}
Here the same $h$ occurs on every component. To prove this, restrict
to $C_N$, $d\mid N\mid6d$, and compare with the constructed finite
line. The fixed-point section and the scalar isotopy give
$\pi_1(W_{C_N})=\mathbb Z^2\times C_N$.
Every fixed-point section has zero map into $\mathbb Z^2$, since
multiplication by $N$ is injective there. The universal-cover
degree-two sequence places the difference of the two classes in
group cohomology. Restriction along these identical finite subgroups
therefore gives one common character. Each fixed curve has no
degree-one mixed term. The map
$H^2(BS^1;\mathbb Z/d)\to H^2(BC_N;\mathbb Z/d)$ sends the
character-one generator to that generator and is an isomorphism;
the inverse does the reverse assignment. This proves the full-circle
formula as well.

In the $\psi$ support convention,
\begin{equation}
 \psi(0,j+1)-\psi(0,j)=6j.
 \hypertarget{tag-025T}{}\label{maps:eq:D.TP3}
\end{equation}
Its usual dense-character-zero linearization differs from the
transported one by the common character $3$.
Thus, on putting $\lambda_0=h-3$ and taking $\epsilon=-1$,
\begin{equation}
 \lambda_j=\lambda_0+6\epsilon j\quad\text{in }\mathbb Z/d.
 \hypertarget{tag-025U}{}\label{maps:eq:D.TP4}
\end{equation}
The circle automorphism $z\mapsto z^{-1}$, with itself as inverse,
induces $t\mapsto-t$ and negates every weight and character.
It sends $(\lambda_0,\epsilon)$ to $(-\lambda_0,1)$ in this
formula. The monomial frame maps above compare the two support
conventions, and both descriptions respect $j\mapsto j+d$.

\item \emph{The full third-power operation.}
For the weights divisible by three, Thomas's operation is defined
with precisely the required full coefficients:
\begin{equation}
 \mathcal P_3:H^2(-;\mathbb Z/d)
                    \longrightarrow H^6(-;\mathbb Z/(3d)).
                                                               \hypertarget{tag-025V}{}\label{maps:eq:D.TP5}
\end{equation}
This is \cite[equation~(2.1), p.~267]{ffThomas56} with
$p=3$ and $m=d/3$. On the reduction of an integral
class $z$, it equals the reduction of $z^3$ modulo
$3d$. Indeed the correction term in Thomas's cochain formula (3.4) is zero for an integral cocycle,
and its first term is the ordinary threefold diagonal
product \cite[equations~(3.4)--(3.6), p.~268]{ffThomas56}.
Each restriction \eqref{maps:eq:D.TP2} has such an integral lift.

\item \emph{Localization at the full third-power modulus.}
Restrict to $C_N$ with $N=3d$ and apply equivariant localization
with coefficients $\mathbb Z/(3d)$ to $\mathcal P_3(\widehat u)$.
The finite line lift exists by $N\mid6d$. The complement is free
and has finite-dimensional quotient, so its cohomology is killed
by a sufficiently high power of $t$. The inverse of restriction
after localization is $\sum_j i_{j*}(z/e_j)$, by the Thom
self-intersection formula. The normal Euler class on each fixed curve is
\[
                    (-x_C+t)(-x_C-t)=-t^2,
\]
because $x_C^2=0$. Integrating the term of ordinary
curve degree two gives
\begin{equation}
 \left\langle\mathcal P_3(u),[X]\right\rangle
                      =-3\sum_{j=0}^{d-1}\lambda_j^2
                                  \quad\text{in }\mathbb Z/(3d).
                                                               \hypertarget{tag-025W}{}\label{maps:eq:D.TP6}
\end{equation}
The right side is independent of the integer
representatives of the characters: $3\mid d$ makes
the cube of an integral lift well defined modulo
$3d$. Since $N=3d$, multiplication by $t$ sends
$H^0(BC_N;\mathbb Z/(3d))$ isomorphically to $H^2$, and the
positive even powers have the same coefficient group. Thus
localization preserves the constant coefficient in this formula.

\item \emph{The pure cyclic coordinate.}
For arbitrary integers $d>0,\lambda_0$ and either
sign in \eqref{maps:eq:D.TP4}, the exact identity is
\begin{equation}
 \sum_{j=0}^{d-1}(\lambda_0+6\epsilon j)^2
   =d\left(\lambda_0^2+6\epsilon\lambda_0(d-1)
                         +6(d-1)(2d-1)\right).              \hypertarget{tag-025X}{}\label{maps:eq:D.TP7}
\end{equation}
Substitution into \eqref{maps:eq:D.TP6} proves
\begin{equation}
                 \left\langle\mathcal P_3(u),[X]\right\rangle=0
                                  \quad\text{in }\mathbb Z/(3d).
                                                               \hypertarget{tag-025Y}{}\label{maps:eq:D.TP8}
\end{equation}
Thus the complete pure cyclic fundamental-class coordinate
is zero, including the order-three subgroup invisible to
the ordinary mod-$d$ cube. For $3\nmid d$ take $N=d$ and use the ordinary cube to obtain
zero modulo $d$. At $d=1$ this value lies in the zero coefficient group. The mixed coordinates have the full modulus $D_2$ computed above.

\item \emph{A finite-cap detector with its reference parity.}
If $d$ is odd, $2P_{\rm ref}=\mathcal D_X(\lambda)=0$ already gives
$P_{\rm ref}=0$. For even $d$, at the finite cap with
order-four elliptic monodromy, put $\beta=[a]$ and $\eta_{\rm bas}=[u_1]$ in its two
central base directions.  Its marked base lattice is
\begin{equation}
 \langle\beta,\eta_{\rm bas},g\mid m_1g=-2r_1\beta+n_1\eta_{\rm bas}\rangle .
 \hypertarget{tag-027M}{}\label{maps:eq:D.WI8}
\end{equation}
Choose integers ${b_{\chi}},{c_{\chi}}$ with $m_1{c_{\chi}}=-2r_1+n_1{b_{\chi}}$ and define the
primitive character by $\chi(\beta)=1$, $\chi(\eta_{\rm bas})={b_{\chi}}$,
$\chi(g)={c_{\chi}}$.  Such integers exist because $\gcd(m_1,n_1)=1$.
Its connected fibre $N_{\rm det}$ in the central four-manifold is an
elliptic-torus mapping torus.  In its base kernel
$h_{\rm ker}=g-{c_{\chi}}\beta$ and $q_{\rm ker}=\eta_{\rm bas}-{b_{\chi}}\beta$ obey $m_1h_{\rm ker}=n_1q_{\rm ker}$.
Choose ${a_{\rm ker}},b_{\rm ker}$ with ${a_{\rm ker}}m_1+b_{\rm ker}n_1=1$ and put $t_{\rm ker}={a_{\rm ker}}q_{\rm ker}+b_{\rm ker}h_{\rm ker}$.
Then $q_{\rm ker}=m_1t_{\rm ker}$, $h_{\rm ker}=n_1t_{\rm ker}$, so the inverse base coordinates are
$\eta_{\rm bas}={b_{\chi}}\beta+m_1t_{\rm ker}$, $g={c_{\chi}}\beta+n_1t_{\rm ker}$.
The root-loop exponent is $b_{\rm ker}=n_1^{-1}\bmod m_1$, odd modulo four.  The elliptic
monodromy of $N_{\rm det}$ is therefore the order-four matrix or its
inverse.  An affine translation in this return is conjugated away by
translation through $(I-A_{\rm ell})^{-1}$ of that vector, with the opposite
translation as inverse. For $A_{\rm ell}=\left(\begin{smallmatrix}0&1\\-1&0\end{smallmatrix}\right)$,
the matrices
$U_{\rm snf}=\left(\begin{smallmatrix}-1&0\\1&-1\end{smallmatrix}\right)$,
$V_{\rm snf}=\left(\begin{smallmatrix}1&1\\0&1\end{smallmatrix}\right)$ satisfy
$U_{\rm snf}(A_{\rm ell}-I)V_{\rm snf}=\operatorname{diag}(1,2)$.
For $A_{\rm ell}^{-1}$ use
$U_{\rm snf}=\left(\begin{smallmatrix}-1&0\\-1&-1\end{smallmatrix}\right)$ and
$V_{\rm snf}=\left(\begin{smallmatrix}1&-1\\0&1\end{smallmatrix}\right)$.
Every determinant is one and each adjugate is its integer inverse.
The Wang sequence therefore gives
\begin{equation}
             H_1(N_{\rm det};\mathbb Z)=\mathbb Z\oplus\mathbb Z/2.
 \hypertarget{tag-027N}{}\label{maps:eq:D.WI9}
\end{equation}
Moreover $\chi(\delta)=-2$, so the scalar tangent
vector is transverse to $N_{\rm det}$ in the reduced finite fibre.
Projecting $J\delta,\operatorname{Re}v_2,\operatorname{Im}v_2$
to $TN_{\rm det}$ along $\delta$ identifies the stable tangent frame
induced by $\rho$ with a tangent parallelization.
The normal three-frame is $\delta$ followed by the two real
directions of $L$.

Put $S=\operatorname{span}_{\mathbb R}
\{J\delta,\operatorname{Re}v_2,\operatorname{Im}v_2\}$, and let
$\ell_\delta$ read the $\delta$ coefficient in
$TQ_1=\mathbb R\delta\oplus S$. Since $\chi(\delta)=-2$, the
projection and its inverse are
\[
 S\longrightarrow TN_{\rm det},\quad w\longmapsto w+\tfrac12\chi(w)\delta,
 \qquad
 TN_{\rm det}\longrightarrow S,\quad v\longmapsto v-\ell_\delta(v)\delta.
\]
For $w\in S$ one has $\ell_\delta(w)=0$, and for $v\in TN_{\rm det}$
one has $\chi(v)=0$; these identities verify both composites.
The ambient change is a triangular shear, homotopic through
invertible shears to identity, so it identifies the two stated framings.

\item \emph{The tangent-parallelization parity.}
Apply \eqref{maps:eq:D45.15} to the closed oriented three-manifold
$N_{\rm det}$ with the tangent parallelization obtained by the two projections
above. The identification with that statement's input is the identity
on $(N_{\rm det},TN_{\rm det})$ and on its projected frame. Its induced spin
structure is the one used by the prescribed-spin filling
there. The coordinate maps in \eqref{maps:eq:D.WI9} and their
inverses reduce modulo two to identify
$H_1(N_{\rm det};\mathbb F_2)=\mathbb F_2^2$.

The explicit convention maps in that parity identity are as follows.
For $\operatorname{Spin}(4)=SU(2)_L\times SU(2)_R$, the action is
$(q_L,q_R)v=q_Lvq_R^{-1}$, with inverse action
$v\mapsto q_L^{-1}vq_R$. The factors are two-connected, so
cellular obstruction theory identifies the homotopy classes of maps
from $N_{\rm det}$ with $H^3(N_{\rm det};\mathbb Z^2)$. Evaluation on $[N_{\rm det}]$ gives
the two integer degrees $d_L,d_R$. For the inverse, collapse the
complement of a positive three-ball and compose with the degree
$d_L$ and $d_R$ maps to the two three-spheres. The first column is $q_Lq_R^{-1}$, of degree $d_L-d_R$.
Stabilization has parity $d_L+d_R\pmod2$, because the two half-spin
summands have equal stable Pontryagin magnitude and opposite Euler
contributions. The quaternionic generator has framed coefficient one.
These are precisely the first-column and stable maps in the proof of
\eqref{maps:eq:D45.15}; its prescribed-spin handle and even-rank
intersection argument proves their parity is
$1+\dim H_1(N_{\rm det};\mathbb F_2)$ for every tangent parallelization and
its specified spin structure
\cite[Definition~1.7, Remark~1.8 and Theorem~3.1, pp.~239--240,243--246]{mapsKaplan79}.
For the present $N_{\rm det}$ this gives
\begin{equation}
 \epsilon_{\rm det}=[N_{\rm det}]_\rho=1\pmod2.
 \hypertarget{tag-0273}{}\label{maps:eq:D.WI10}
\end{equation}

\item \emph{The integral normal chain and linking number.}
The following normal columns are realized by their Thom chains.
Let $K_{a0},K_{a1}$ be the first columns of the two blocks in
\eqref{maps:eq:D.WE3}.  Directly in its four target coordinates,
\begin{equation}
                    d s=-2\ell P\,K_{a0}-kP\,K_{a1}.
 \hypertarget{tag-0274}{}\label{maps:eq:D.WI11}
\end{equation}
The chain at cap zero misses $N_{\rm det}$, whereas the $a$-Thom
chain in cap one pairs with $N_{\rm det}$ by $\chi(a)=1$.
Thus $\operatorname{lk}(N_{\rm det},s)=-kP/d$.
The determinant equation implies that $k$ is odd, so $N_{\rm det}$
generates the two-primary summand of $H_3(X)$.
The specialization chains for \eqref{maps:eq:D.WI7} are supported in the
cusp and miss $N_{\rm det}$.  Equations \eqref{maps:eq:D.WI7}--\eqref{maps:eq:D.WI11} then give
the integer link
\begin{equation}
                          \ell(N_{\rm det},C)=5kP.
 \hypertarget{tag-0275}{}\label{maps:eq:D.WI12}
\end{equation}
The orientations here are those of \eqref{maps:eq:D.WI11}, with
the cap-one $a$ normal intersecting $N_{\rm det}$ positively. Reversing the
orientation of $N_{\rm det}$ negates both linking values by linearity on chains;
the inverse is the same orientation reversal. In either orientation
this integer is odd.

\item \emph{The admissible sphere and relative transgression.}
An odd multiple of $N_{\rm det}$ kills its odd-primary homology component;
join parallel copies by tubes in $X\setminus C$.
Its induced stable parity and critical-link parity stay odd.
Use Kaplan's filling for its prescribed almost framing, delete
the zero-handle, and reverse the resulting two-handle cobordism.
Deleting the real-codimension-four critical cycle preserves simple
connectivity. Each attaching circle therefore bounds an embedded
disc in $X\setminus C$. General position makes these discs mutually
disjoint and disjoint from the three-manifold in their
interiors, since $2+2<6$ and $2+3<6$.
The first normal gives the boundary collar. The four remaining
directions extend over each disc precisely when their
$\pi_1SO(4)=\mathbb Z/2$ spin obstruction vanishes; it does by
the prescribed handle spin structure. Thickening gives ambient surgeries to an embedded torsion $S^3$, preserving its
integral class and induced framing parity.

For $q_{\mathbb H}=a_{\mathbb H}+x_{\mathbb H}i+y_{\mathbb H}j+z_{\mathbb H}k$, put
\[
 t_{{\mathbb H},1}=q_{\mathbb H}i,\qquad t_{{\mathbb H},2}=q_{\mathbb H}j,\qquad t_{{\mathbb H},3}=q_{\mathbb H}k.
\]
The frame
\[
 (-a_{\mathbb H} e_3-x_{\mathbb H} t_{{\mathbb H},1}-y_{\mathbb H} t_{{\mathbb H},2}-z_{\mathbb H} t_{{\mathbb H},3},e_1,e_2)
\]
is the nontrivial class relative to the tangent sphere frame,
and extends over the four-handle. Its coefficient map has degree
one on $S^3$ and maps onto $\pi_3V_{6,3}=\mathbb Z/2$.
This fixes the convention in Zhubr's proved three-frame criterion
\cite[Theorem~1 and Section~5, pp.~6,8--9]{ffZhubr13}.
The adjoint $SU(2)$ bundle has $p_1=-4c_2$, whereas its real
quaternion bundle has $p_1=-2c_2$. Hence
$\pi_3\operatorname{Spin}(3)\to\pi_3\operatorname{Spin}(6)$
is multiplication by two. The tangent quaternion sphere frame
is odd, while the extending admissible class has even induced
stable coefficient.

Choose the admissible three-frame of that theorem; its spin
rank-three complement has obstruction $\mathfrak p_X$.
Trivialize the complement on the detector. Its choices change
the stable degree by even integers, and extension over the
reversed two-handle trace preserves that parity.
The resulting full frame $\sigma$ has even detector parity.
Thus $\sigma/\rho$ has an integral $\operatorname{Spin}(6)$
degree $d_\rho\equiv\epsilon_{\rm det}\bmod2$.
If $B_4$ bounds $\operatorname{ord}(N_{\rm det})$ copies of $N_{\rm det}$, the relative
four-cell obstruction and the index-two inclusion just computed give
\[
 -B_4\cdot C=2q_{\rm ref}+\operatorname{ord}(N_{\rm det})d_\rho.
\]
Here $q_{\rm ref}$ is the complementary obstruction cochain evaluated with
its $\sigma$ boundary trivialization. Division by
$2\operatorname{ord}(N_{\rm det})$ yields the exact formula
\begin{equation}
 \operatorname{lk}(N_{\rm det},P_{\rm ref})
 =-\frac{\ell(N_{\rm det},C)}2-\frac{\epsilon_{\rm det}}2\quad\text{in }\mathbb Q/\mathbb Z.
 \hypertarget{tag-0276}{}\label{maps:eq:D.WI13}
\end{equation}
The complementary cochain represents $\mathfrak p_X$.
Its cap product with $[X]$ is $P_{\rm ref}$, and the inverse map
is $\mathcal D_X^{-1}$. Thus $q_{\rm ref}/\operatorname{ord}(N_{\rm det})$ is the
linking with the displayed homology class.
Changing $d_\rho$ by an even integer changes the right side by an
integer. Multiplication by the chosen odd number is invertible
on the order-two values, so the formula also holds for the
initial detector. The free part vanishes from
$2P_{\rm ref}=\mathcal D_X(\lambda)=0$.
Equations \eqref{maps:eq:D.WI10} and \eqref{maps:eq:D.WI12}
make the two odd halves cancel. The linking is primitive on the
two-primary summand, and multiplication by two is invertible on
the odd summand. This proves $P_{\rm ref}=0$ with this marked detector.

\item \emph{Complete marked oriented recognition.}
Let $A_d=\mathbb Z^{d-1}\oplus\mathbb Z/d$ with its displayed basis.
The isomorphism $\psi_X:A_d\to H_2(X)$ sends the free basis to the
lifts adjusted by their $u$ values and the cyclic generator to $s$.
Its inverse evaluates $x_0,\ldots,x_{d-2}$ and then $u$.
The universal-coefficient isomorphism
$H^2(X;A_d)=\operatorname{Hom}(H_2(X),A_d)$ represents this inverse
by a classifying map $f_{\rm cls}:X\to K(A_d,2)$.
The full marked pair is
\[
 P(f_{\rm cls})=0,\qquad
 \mu_d=\sum_{0\le i\le j\le d-2}c^{\rm inc}_{ij}\gamma_{ij}(e)g
                  -12\sum_{i=0}^{d-2}e_i\gamma_2(g),
\]
where $c^{\rm inc}_{ij}$ is the incidence coefficient of
\eqref{maps:eq:D.WI3}. All free cubic and pure cyclic coordinates
are zero. The inverse detectors above fix every coefficient,
so this pair depends on $d$ alone for all allowed data.

Zhubr's Theorem 5 is injective on the marked spin-bordism group
for the finitely generated group $A_d$. The marked recognition
of Theorem 3, with its Section 4.1 construction and uniqueness
at zero free $H_3$, applies here
\cite[Theorems~3,5 and Section~4.1, pp.~795--796,809--810]{manZhubr75}.
For any two members it gives an oriented diffeomorphism $F$ with
$F_*=\psi_{X'}\psi_X^{-1}$ and inverse
$(F^{-1})_*=\psi_X\psi_{X'}^{-1}$.
Thus every adjusted free generator and the positive cyclic generator
have the displayed images in both directions.

Existence at every $d$ follows from $k=\ell=1$, $m_0=3,m_1=4$,
$n_0=1,n_1=-1,r_0=r_1=0$, whose determinant is $1$.
The coprimalities give the finite-cap freeness conditions
of \ref{maps:src:D.2.34}, so they realize the member $M_d^{\rm w}$
of Definition~\ref{maps:def:weighted-engel-data}.
For $d=1$ all its positive homology below degree six is zero.
The sphere's stable tangent map
$(x;v,a)\mapsto(x;v+ax)$ has inverse
$(x;y)\mapsto(x;y-\langle y,x\rangle x,\langle y,x\rangle)$.
It gives the same zero spin-bordism data; marked recognition gives
the orientation-preserving identification $M_1^{\rm w}=S^6$.
When $d$ is prime to six, $D_2=D_3=d$ and the marked invariants
coincide with \ref{maps:src:D.1.11}, under the same basis map and
its inverse. Finally $\chi(X)=2+2(d-1)=2d$; diffeomorphism-induced
inverse homology maps preserve this alternating sum, so distinct
positive weights have distinct types.
\end{enumerate}
\end{proof}

\subsection{The weighted OS45 manifold}

\begin{definition}[Weighted OS45 data]\hypertarget{tag-Q251}{}\label{maps:def:weighted-os45-data}
{\normalfont\scriptsize tag Q251}\par
Use the OS45 base profile and its ordered markings in
Definition~\ref{mapsIII:def:elliptic-marked-data}, with generator $G$,
initial sections and whole-fibre weights
\begin{equation}
 \mathrm{OS45/base},\quad P=2G,\quad Q=8G,\quad
 (w_0,w_1,w_2,w_3,w_4)=(2,0,0,0,0),\quad e=1,              \hypertarget{tag-023A}{}\label{maps:eq:D.PC3}
\end{equation}
The weight two is supported at the $I_8$ site. Fill all multiplicative
sites by the prescribed Mumford caps and use the ordinary clutch $e=1$.
Write $X$ for this marked member with its complex orientation.
These are the \emph{weighted OS45 data}.
\end{definition}

\Needspace{6\baselineskip}
\begin{theorem}[Invariants of the weighted OS45 manifold]\hypertarget{tag-02LV}{}\label{maps:src:D.2.53}
{\normalfont\scriptsize tag 02LV}\par
The manifold $X$ of Definition~\ref{maps:def:weighted-os45-data}
is simply connected and has
\begin{equation}
 H_2(X;\mathbb Z)=\mathbb Z^{15}\oplus\mathbb Z/8,\qquad
 H_3(X;\mathbb Z)=\mathbb Z/8,\qquad
 c_1(X)=c_2(X)=\lambda(X)=0.                              \hypertarget{tag-022O}{}\label{maps:eq:D.OS1}
\end{equation}
Its sixteen normalized cusp components are indexed by
\((r,s)\in\mathbb Z/8\times\mathbb Z/2\), with \(i=2r+s\).
Their divisor classes generate an index-two sublattice in
\(H^2(X;\mathbb Z)\). The full primitive lattice is generated
by \(x_0,\ldots,x_{13},y\), with
\begin{equation}
             2y=\sum_{r=0}^7x_{r,1},\qquad
             \sum_{r,s}x_{r,s}=0.                        \hypertarget{tag-022P}{}\label{maps:eq:D.OS2}
\end{equation}
The integral cubic is \eqref{maps:eq:D.PC11}--\eqref{maps:eq:D.PC12}.
There is a primitive cyclic coordinate \(u\in H^2(X;\mathbb Z/8)\)
with the full boundary marking \eqref{maps:eq:D.OF1}--\eqref{maps:eq:D.OF14}. In the splitting
fixed by \(u\) and this primitive basis,
\begin{equation}
 H_6(K(H_2(X),2);\mathbb Z)
   =\mathbb Z^{680}\oplus(\mathbb Z/8)^{120}
                       \oplus(\mathbb Z/16)^{15}
                       \oplus\mathbb Z/8,                \hypertarget{tag-022Q}{}\label{maps:eq:D.OS3}
\end{equation}
and the complete fundamental-class invariant is specified by
the integral cubic and \eqref{maps:eq:D.OF17}--\eqref{maps:eq:D.OF18}.
In particular
\begin{equation}
 \begin{gathered}
 y^2=0,\qquad
 \langle\mathcal P_2(u)x_i,[X]\rangle=7\pmod{16}
                   \ (0\le i<14),\\
 \langle\mathcal P_2(u)y,[X]\rangle=0,\qquad
 \langle u^3,[X]\rangle=0.
 \end{gathered}                                           \hypertarget{tag-022R}{}\label{maps:eq:D.OS4}
\end{equation}
The refined characteristic class is
\[
                         P_{\mathrm{ref}}(X)=0.
\]
Here $P_{\mathrm{ref}}\in H_2(X;\mathbb Z)$ is the refined class.
The cap-product map $\mathcal D_X:H^4(X;\mathbb Z)\to H_2(X;\mathbb Z)$,
$c\mapsto c\frown[X]$, has inverse Poincar\'e duality.
The primitive cubic and the finite coordinates above determine the
complete fundamental-class invariant of this marked member.  Its source-line circle bundle
embeds as a primitive torsion sphere.  For the differential reference
frame its framed parity is one, and its intersection with an integral
nullchain of the complete critical cycle is odd.

\end{theorem}
\begin{proof}

\begin{enumerate}[leftmargin=*]
\item \emph{The primitive lattice, cubic and critical cycle.}
Let $N$ be the sixteen-component cusp of the marked manifold in the
statement and put $E=X\setminus\operatorname{int}N$.  The integral
cap maps of \ref{maps:src:D.2.46}, evaluated on its five marked
returns and the ordinary clutch, give $\pi_1(X)=1$,
$b_2(X)=15$ and the following complement marking, with
$D=16$, $A=\mathbb Z/2$ and $a_{r,s}=s\bmod2$:
\begin{equation}
 H_1(E;\mathbb Z)=\mathbb Z\oplus A,\qquad
                       \mu_i=(1,a_i),                      \hypertarget{tag-0233}{}\label{maps:eq:D.PC1}
\end{equation}
where $A$ is finite and $\mu_i$ is the positively oriented meridian of
component $i$. The component classes $x_i$ then generate a sublattice
\[
 L_{\rm cmp}=\mathbb Z^D/\mathbb Z(1,\ldots,1)
                       \subset H^2(X;\mathbb Z)
\]
whose full primitive extension, inside $L_{\rm cmp}\otimes\mathbb Q$, is
\begin{equation}
 H^2(X;\mathbb Z)
   =L_{\rm cmp}+\left\{\sum_i\widetilde{\chi(a_i)}x_i:
                         \chi\in\operatorname{Hom}(A,\mathbb Q/\mathbb Z)
        \right\}.                                        \hypertarget{tag-0239}{}\label{maps:eq:D.PC2}
\end{equation}
A tilde means any rational lift of a value in $\mathbb Q/\mathbb Z$;
changing it adds an element of $L_{\rm cmp}$. In particular the component
lattice has index $|A|$.

\begin{enumerate}[label=(\alph*),leftmargin=*]
\item \emph{The relative primitive extension.}
Poincar\'e--Lefschetz duality identifies
$H^2(X,E;\mathbb Z)=H_4(N;\mathbb Z)=\mathbb Z^D$,
with its positive component Thom basis, and
$H^3(X,E;\mathbb Z)=H_3(N;\mathbb Z)$ is free by the
normalization complex for this triangular cusp in
\ref{maps:src:D.2.43}. The connecting
map from $H^1(E;\mathbb Z)$ evaluates on the meridians \eqref{maps:eq:D.PC1},
so its image is precisely $\mathbb Z(1,\ldots,1)$.
The relative exact sequence therefore injects $L_{\rm cmp}$ into $H^2(X)$.
Since $X$ is simply connected, the latter group is free, of rank fifteen. Its quotient by $L_{\rm cmp}$ is finite and embeds in
$H^2(E)$, whose torsion is
$\operatorname{Ext}(H_1(E),\mathbb Z)=\operatorname{Hom}(A,\mathbb Q/\mathbb Z)$.
All this torsion lifts to $H^2(X)$ because the next relative group is
free. Thus the quotient has exactly that order.

\item \emph{The character embedding and its inverse.}
For a character $\chi$ of order $m$, extend it by zero on the free
summand of $H_1(E)$. Regard it as a class in
$H^1(E;\mathbb Z/m)$ by multiplying its values by $m$.
Its relative coboundary has component coordinates
$m\widetilde{\chi(a_i)}\pmod m$. Choose the displayed integers
as an integral relative lift. Its image in $H^2(X;\mathbb Z/m)$
is zero by exactness. Hence its integral image is divisible by $m$
in $H^2(X;\mathbb Z)$, giving the class in \eqref{maps:eq:D.PC2}.
With the connecting map given by the positive normal-disc boundary,
the coefficient exact sequence identifies its exterior restriction
with the corresponding torsion Bockstein.  Reversing that connecting
map negates both the character and its rational component coefficients,
and the two negations give inverse descriptions of the same extension.
These characters exhaust the torsion of $H^2(E)$ and therefore give
the full primitive extension. This proves
\eqref{maps:eq:D.PC2}, including every finite denominator.

\item \emph{The marked returns and period lattice.}
For the data of Definition~\ref{maps:def:weighted-os45-data},
\ref{maps:src:D.2.46} gives
\begin{equation}
 H_2(X)=\mathbb Z^{15}\oplus\mathbb Z/8,\qquad
 H_3(X)=\mathbb Z/8,\qquad \pi_1(X)=1.
                                                               \hypertarget{tag-023B}{}\label{maps:eq:D.PC4}
\end{equation}
The elliptic monodromy at $I_8$ is
$\left(\begin{smallmatrix}1&8\\0&1\end{smallmatrix}\right)$
and its section cocycle is $(3,0)$. Formula \eqref{maps:eq:D.GI2} therefore gives
\begin{equation}
 A_0=\begin{pmatrix}
 1&8&0&24\\0&1&0&0\\0&-6&1&-16\\0&0&0&1
 \end{pmatrix},\qquad
 L_{\rm per}=\left\langle(8,-6),(24,-16)\right\rangle
                 =8\mathbb Z\oplus2\mathbb Z.             \hypertarget{tag-023C}{}\label{maps:eq:D.PC5}
\end{equation}
The short basis is $(e_1,\delta)$ and the long basis is $(e_2,\gamma)$.
The equality of sublattices follows from the exact factorization
\[
 \begin{pmatrix}8&24\\-6&-16\end{pmatrix}
 =\operatorname{diag}(8,2)T_{\rm per},\qquad
 T_{\rm per}=\begin{pmatrix}1&3\\-3&-8\end{pmatrix},\quad
 T_{\rm per}^{-1}=\begin{pmatrix}-8&-3\\3&1\end{pmatrix}.
\]
Both products of the last two matrices are identity. The two
resulting long periods are $-8e_2+3\gamma$ and $-3e_2+\gamma$.
The determinant is sixteen and the Smith factors are two and eight. Label the normalized components by
\begin{equation}
             (r,s)\in\mathbb Z/8\times\mathbb Z/2,
             \qquad i=2r+s.                              \hypertarget{tag-023D}{}\label{maps:eq:D.PC6}
\end{equation}

\item \emph{The complement and component meridians.}
The cusp complement satisfies
\begin{equation}
 H_1(E)=\mathbb Z\oplus\mathbb Z/2,
                      \qquad \mu_{r,s}=(1,s\bmod2).      \hypertarget{tag-023E}{}\label{maps:eq:D.PC7}
\end{equation}
To verify the marking, use all peripheral generators and the
ordinary clutch in the exact degree-one presentation. Its Smith
form has seven unit factors, one factor two, and one free generator.
The torsion character in the nine coordinates of that marked presentation is
\begin{equation}
                  (0,-2,1,2,0,0,2,-2,2)\pmod2,
                                                               \hypertarget{tag-023F}{}\label{maps:eq:D.PC8}
\end{equation}
which is precisely $\delta^*$ modulo two. The free character
evaluates to one on the cusp meridian. A height-one ray at
$(r,s)$ has meridian $\theta+r e_1+s\delta$; it
therefore gives exactly \eqref{maps:eq:D.PC7}. The inverse sends the free generator to the cusp meridian and the
order-two generator to $\delta$.  In the marked peripheral presentation
the first rank-one cap kills $e_1$, the next return gives
$e_2=-2\delta$, and the cusp relations give $2\delta=0$.
The remaining relations impose no further relation on $\delta$.
The four other meridians vanish, and the ordinary clutch and the
peripheral product identify $\gamma$ with the cusp meridian.
These substitutions verify both inverse composites.

\item \emph{The primitive basis.}
Applying \eqref{maps:eq:D.PC2}, there is a unique integral class $y$ with
\begin{equation}
                    2y=\sum_{r=0}^{7}x_{r,1}.             \hypertarget{tag-023G}{}\label{maps:eq:D.PC9}
\end{equation}
The component lattice has exact index two. An integral primitive basis is
\begin{equation}
 x_0,x_1,\ldots,x_{13},y,\qquad
 x_{14}=-2y-\sum_{\substack{i<14\\i\ \mathrm{even}}}x_i,
 \qquad
 x_{15}=2y-\sum_{\substack{i<14\\i\ \mathrm{odd}}}x_i.
                                                               \hypertarget{tag-0234}{}\label{maps:eq:D.PC10}
\end{equation}
The integral matrix of the sixteen component classes in this basis
has Smith factors $1^{14},2$ and kernel generated by the all-one
vector. Equivalently, the relative kernel consists of integer vectors $n$
with $\sum n_i=0$ and $\sum_{i\ {\rm odd}}n_i$ even.
Its basis consists of $e_i-e_{14}$ for even $i<14$,
$e_i-e_{15}$ for odd $i<14$, and $2e_{15}-2e_{14}$.
The inverse reads $n_0,\ldots,n_{13}$ and half the odd-coordinate
sum. Dual evaluation gives exactly \eqref{maps:eq:D.PC10}.

\item \emph{The complete integral cubic.}
For component classes,
\ref{maps:src:D.2.47} specializes to the following symmetric integral trilinear
form. With distinct vertices as appropriate,
\begin{equation}
 \begin{aligned}
 T(x_i,x_i,x_i)&=6,\\
 T(x_i,x_i,x_j)&=
 \begin{cases}
 -2,&i=(r,s),\ j=(r,1-s),\\
 -1,&i=(r,s),\ j=(r\pm1,t),\quad t=0,1,\\
 0,&\text{otherwise},
 \end{cases}\\
 T(x_i,x_j,x_k)&=
 \begin{cases}
 1,&\{i,j,k\}\text{ is three of the four vertices in}\\[-2pt]
   &\hspace{12pt}\{(r,0),(r,1),(r+1,0),(r+1,1)\}
                   \text{ for some }r,\\
 0,&\text{otherwise}.
 \end{cases}
 \end{aligned}                                             \hypertarget{tag-0235}{}\label{maps:eq:D.PC11}
\end{equation}
All horizontal indices are cyclic modulo eight. Repeated geometric
edges and triangle cells are counted with their geometric multiplicities.
The remaining primitive entries are
\begin{equation}
 \begin{aligned}
 T(y,x_{r,s},x_{r',s'})
 &=(-1)^s\left(
       \mathbf1_{s'=s,\ r'=r+1}
      +\mathbf1_{s'=s,\ r'=r-1}
      -2\mathbf1_{(r',s')=(r,s)}\right),\\
 T(y,y,x_i)&=0,\qquad T(y,y,y)=0.
 \end{aligned}                                             \hypertarget{tag-0236}{}\label{maps:eq:D.PC12}
\end{equation}
These formulas determine every entry on the primitive basis
\eqref{maps:eq:D.PC10}, and every resulting entry is integral.

For a geometric verification of the half-class entries, on a
normalized $dP_6$ take the cyclic boundary classes
\begin{equation}
 E_1,\ H-E_1-E_2,\ E_2,\ H-E_2-E_3,\ E_3,\ H-E_1-E_3
                                                               \hypertarget{tag-0237}{}\label{maps:eq:D.PC13}
\end{equation}
on the ordered rays
$(1,0),(0,1),(-1,1),(-1,0),(0,-1),(1,-1)$
of \ref{maps:src:D.2.47}. The four nonhorizontal boundary classes
sum to $2f$, where $f=H-E_1$ and $f^2=0$. Restricting \eqref{maps:eq:D.PC9}
therefore gives
\begin{equation}
                         y|_{P_{r,s}}=(-1)^s f.            \hypertarget{tag-0238}{}\label{maps:eq:D.PC14}
\end{equation}
This division is valid in the torsion-free group $H^2(dP_6;\mathbb Z)$.
The fibre $f$ has degree one on each horizontal boundary curve
and zero on the other four. The Cartier restriction formula in
\ref{maps:src:D.2.47} now gives the first line of \eqref{maps:eq:D.PC12}. Squaring \eqref{maps:eq:D.PC14}
gives $y^2x_i=0$; multiplying \eqref{maps:eq:D.PC9} by $y^2$ gives $y^3=0$.
This proves the displayed form on the full integral primitive lattice.

Finally $c_1(X)=0$ and both $c_2$ and $p_1$ pair to zero with
every class in the primitive free basis. For $c_1$, the normalized
component tangent calculation is zero and the component restriction
is rationally injective by the connected toric graph Laplacian;
$H^2(X)$ is torsion-free, so this proves integral vanishing.
For $c_2,p_1$, \ref{maps:src:D.2.47} gives zero on every component. The component
classes have finite index, and the evaluations are integers, so they
also vanish on $y$. The torsion class of $c_2$, the refined class and the full
finite-coefficient invariant are computed in the following steps.

\medskip\noindent\emph{Integral vanishing of the critical class.}
The primitive-lattice argument above gives the relation
\(\sum x_{r,s}=0\). Therefore the target line of \(df\) is
topologically trivial. All fibres in this example are reduced
semistable fibres. The exact complex of \ref{maps:src:D.2.44} consequently
applies without any additive correction: its abstract kernel
has the nowhere-zero scalar section, and its quotient line is
topologically trivial because \(c_1(X)=0\). The local unit Bott
clutch gives
\begin{equation}
                     c_2(X)=\operatorname{PD}(C_{\rm crit}).
                                                               \hypertarget{tag-022S}{}\label{maps:eq:D.OS5}
\end{equation}
Here every complex double curve occurs once.

The integral specialization of the complete cycle is as follows.
In the ordered bivector basis
\((e_1e_2,e_1\delta,e_1\gamma,e_2\delta,e_2\gamma,\delta\gamma)\),
the five fibre preimages, in their site order, are
\begin{equation}
 \begin{array}{c|rrrrrr}
 I_8\text{ weighted cusp}&-16&0&6&8&0&2\\
 I_1&0&0&0&0&0&1\\
 I_1&0&8&0&0&0&1\\
 I_1&0&0&0&8&0&1\\
 I_1&0&0&0&0&0&1
 \end{array}.                                             \hypertarget{tag-022T}{}\label{maps:eq:D.OS6}
\end{equation}
For the first row, the local period cover has matrix
\(M_{\rm per}=\left(\begin{smallmatrix}8&24\\-6&-16\end{smallmatrix}\right)\)
and degree sixteen. Transfer the minimal cusp's three-curve preimage
\(a_1b_1+b_1a_2+a_2b_2\), expressing \(b_1,b_2\) by \(M_{\rm per}^{-1}\)
in the long basis, and multiply by sixteen. This gives
the first integral row of \eqref{maps:eq:D.OS6}; each lifted simplex
has its integer transfer coefficient on this local period cover.
Each other row is the determinant-one transverse torus of the
corresponding primitive node, as in \ref{maps:src:D.2.44}.
Their sum is
\begin{equation}
                        z=(-16,8,6,16,0,6).              \hypertarget{tag-022U}{}\label{maps:eq:D.OS7}
\end{equation}
For the second rank-one site take the transverse elliptic period
$r_{\rm per}=e_1$, for which $\det(2e_1-e_2,r_{\rm per})=1$; its fixed long period is
$\gamma-8e_1$.  For the third take $r_{\rm per}'=e_2$.  These choices give
the last four rows of \eqref{maps:eq:D.OS6} by the primitive nodal
specialization maps in \ref{maps:src:D.2.46}.

Let $T_i(v)$ be the marked monodromy prism with boundary
$(\bigwedge^2A_i-I)v$.  In the ordered bivectors of
\eqref{maps:eq:D.OS6}, put
\[
\begin{aligned}
q_0&=(48,0,-16,0,2,0),&q_1&=(0,0,0,8,2,0),\\
q_2&=(-2,-52,-8,0,0,0),&q_3&=(0,0,2,0,0,0),&q_4&=0.
\end{aligned}
\]
Direct exterior multiplication of the five marked returns gives
\begin{equation}
 \sum_{i=0}^4(\bigwedge^2A_i-I)q_i+10\delta\wedge\gamma=z.
 \hypertarget{tag-022V}{}\label{maps:eq:D.OS8}
\end{equation}
The pure $\delta\wedge\gamma$ term has a global bounding
chain.  At each nonordinary cap the origin $\delta$-normal chain
has boundary $T(\delta,\mu_i)$; at the ordinary cap it has boundary
$T(\delta,\mu_o)+\delta\wedge\gamma$.  The product of the
$\delta$ circle with the punctured base has boundary
$-\sum_iT(\delta,\mu_i)$, including the ordinary boundary.
Their sum $D_\delta$ therefore bounds $\delta\wedge\gamma$.
The inverse coordinates read the physical scalar angle, the marked
origin section and each peripheral collar.

For the five critical-specialization chains $S_i$, the local charts give $\partial S_i=C_i-v_i$, with $v_i$ the rows of
\eqref{maps:eq:D.OS6}.  Thus
\[
 B_C=\sum_i S_i+\sum_iT_i(q_i)+10D_\delta,\qquad
                          \partial B_C=C.
\]
The cusp chain is the degree-sixteen transfer of the primitive
three-curve chain: every lifted simplex is taken once.  On each lifted
simplex the covering chart and its inverse identify its coefficient;
globally $p_*\operatorname{tr}=16\,\mathrm{id}$.  The two long periods are
$-8e_2+3\gamma$ and $-3e_2+\gamma$, whose valuations are
$(8,0)$ and $(0,2)$.  Their reciprocal unit-period maps retain the
source line units.  This proves the integral critical-class
vanishing in \eqref{maps:eq:D.OS1} with the full cycle and its chain.

\end{enumerate}

\item \emph{The cyclic coordinate and component evaluations.}
Use the datum \eqref{maps:eq:D.PC3}, its period marking \eqref{maps:eq:D.PC5}, and the
primitive free basis \eqref{maps:eq:D.PC10}. Put \(k=\mathbb Z/8\).
There is a primitive torsion coordinate
\(u\in H^2(X;k)\) whose complete cusp support is
\begin{equation}
 \Psi(r,s)=4\left\lfloor r/8\right\rfloor
       +3(r\bmod8)\left\lfloor s/2\right\rfloor
                                      \quad\hbox{in }k.       \hypertarget{tag-0217}{}\label{maps:eq:D.OF1}
\hypertarget{tag-021Q}{}
\end{equation}
\SsixDisplayAliases{maps:eq:D.OF1}{\label{maps:eq:D.OL1}}
Here and below \(r\bmod8\) means the representative between zero and
seven. The assertion includes the entire boundary cohomology class,
including its coefficient-extension part. The full boundary comparison is proved below.

\medskip\noindent\emph{The specified core cochain.}
Use the degree-two core model of \ref{maps:src:D.2.46}. Its fibre cells are
ordered
\[
 e_1e_2,\ e_1\delta,\ e_1\gamma,\ e_2\delta,\
                                      e_2\gamma,\ \delta\gamma.
\]
For the six boundary monodromies, including the ordinary identity cap,
let
\[
 D_{\rm gr}=(A_0-I,\ldots,A_4-I),\qquad \mathcal K=\ker_{\mathbb Z}D_{\rm gr}.
\]
Let $K:\mathbb Z^{17}\to\mathbb Z^{20}$ be the integral basis embedding of $\mathcal K$, with integral left inverse
$K^+$ obtained from the marked graph differential in
\ref{maps:src:D.2.46}; thus $D_{\rm gr}K=0$ and $K^+K=1$.
The ordinary degree-two cycles are the six fibre cells followed by
these seventeen graph cycles. Its relation matrix \(M\) is
the marked relation matrix of that construction.  Integer row and
column reduction gives $U_{\rm snf}MV_{\rm snf}=\Delta$, where $\Delta$ has diagonal $(1^{22},8)$ and
its remaining columns are zero.  The quotient map is $z\mapsto U_{\rm snf}z$
modulo this diagonal and its inverse sends a standard generator to
$U_{\rm snf}^{-1}e_j$; the relation $\Delta_{jj}U_{\rm snf}^{-1}e_j$ is bounded by the
prism and cap chain whose coefficients are $V_{\rm snf}e_j$.

We specify a cochain on the twenty ambient graph coordinates,
ordered in five blocks \((e_1,e_2,\delta,\gamma)\). Put
\[
\begin{aligned}
\theta_{\rm gr}={}&(0,4,0,12;\ 0,0,0,0;\ 0,7,0,4;\ 5,0,0,4;\ 1,6,0,0),\\
r_{\rm gr}={}&(0,6,0,17;\ 0,0,0,0;\ 0,0,0,0;\ 0,0,0,0;\ 0,0,0,0),\\
\theta_{\rm gr}'={}&\theta_{\rm gr}+8r_{\rm gr},\qquad
\varphi=((1,0,0,0,0,0),\theta_{\rm gr}'K).
\end{aligned}
\]
Substitution in the marked bar-prism and cap boundary maps gives
\begin{equation}
\varphi M=0\pmod8,\qquad
\varphi|_{\Lambda^2\Lambda}=(1,0,0,0,0,0).              \hypertarget{tag-021I}{}\label{maps:eq:D.OF2}
\end{equation}
The core quotient has order eight, and the fibre class
$g=-[e_1\wedge e_2]$ has $\varphi(g)=-1\pmod8$.
Its order is therefore eight.  Hence this cochain is
primitive on the order-eight core image.

The ambient graph row is \(\theta_{\rm gr}'\) itself: its restriction to the
integral cycle lattice is exactly \(\theta_{\rm gr}'K\). The fixed bar-prism
comparisons of \ref{maps:src:D.2.46} transport this cochain to
the torus bundle. Its finite-coefficient class is unchanged
by the displayed addition \(8r_{\rm gr}\), and its boundary restriction is
specified on all ambient generators, including the coefficient-extension
part arising from the core's degree-one torsion.

Let \(J_{\rm raw}\) be the rank-two cusp peripheral map into
the twenty graph coordinates, with its bar-prism comparison.
The explicit ambient row gives
\begin{equation}
\theta_{\rm gr}'J_{\rm raw}=(0,52,0,148)
       \quad\hbox{on }(e_1,e_2,\delta,\gamma).             \hypertarget{tag-021J}{}\label{maps:eq:D.OF3}
\end{equation}
The quadratic correction on each of these four basis loops is zero
at this cusp. The two long-strip constants are therefore
\begin{equation}
                      (c_1^{\rm str},c_2^{\rm str})=(4,4)\quad\hbox{in }k.   \hypertarget{tag-021K}{}\label{maps:eq:D.OF4}
\end{equation}
\medskip\noindent\emph{A full boundary comparison for finite-coefficient supports.}
The cusp boundary has torsion in \(H_1\); its full comparison
uses the following cochain maps. Let \(\widetilde N\) be the common universal
cover of the triangular Mumford cusp and let
\(\widetilde M\) be its boundary. Its normalized components are
indexed by \(v\in\mathbb Z^2\). The boundary cover retracts to the
three-torus generated by the two short angular loops and the positive
\(q\)-meridian \(\theta\). \ref{maps:src:D.2.43} gives
\begin{equation}
 H_1(\widetilde N;\mathbb Z)=0,\quad
 H_2(\widetilde N;\mathbb Z)=I^2,\quad
 \operatorname{im}\{H_2(\widetilde M;\mathbb Z)\to
                         H_2(\widetilde N;\mathbb Z)\}=0. \hypertarget{tag-021L}{}\label{maps:eq:D.OF5}
\end{equation}
Indeed the three boundary generators are the
angular area and the two angular-normal tori, all with the bounding sweeps described there.

Let $e_v$ be the positive relative Thom class of the component $v$.
Our support coordinate is the explicit involution
$\Psi\longmapsto-\sum_v\Psi(v)e_v$, with inverse the negative
of the Thom coefficient at each component.  Thus these coordinates
identify
\[
 H^2(\widetilde N,\widetilde M;k)
                  =\operatorname{Map}(\mathbb Z^2,k).
\]
One can obtain this directly from locally finite
component chains: a degree-four locally finite cycle is an arbitrary
family of fundamental cycles of the normalized components, and
degree-four intersection strata do not occur. Equivalently,
Poincar\'e--Lefschetz duality uses Borel--Moore degree-four homology.
The locally finite component-chain map and its coefficient-reading
inverse give this product of coefficient groups.

The positive meridian of component \(v\) is
\(\theta+v_1e_1+v_2\delta\). Therefore the relative connecting map
from \(H^1(\widetilde M;k)\) sends a covector \((c,m_1,m_2)\) to
the support function $-(c+m_1v_1+m_2v_2)$: its positive Thom
coefficients are $c+m_1v_1+m_2v_2$.  Negation is its own inverse,
and the image is exactly $\operatorname{Aff}$. By \eqref{maps:eq:D.OF5}, the relative
sequence is exactly
\begin{equation}
 0\longrightarrow\operatorname{Aff}(\mathbb Z^2,k)
 \longrightarrow\operatorname{Map}(\mathbb Z^2,k)
 \longrightarrow H^2(\widetilde N;k)\longrightarrow0.     \hypertarget{tag-021M}{}\label{maps:eq:D.OF6}
\end{equation}
Here \(\operatorname{Aff}\) consists of these three-parameter
affine functions. This also identifies the quotient with
\(\operatorname{Hom}(I^2,k)\), consistently with the normalization computation of \ref{maps:src:D.2.43}.

Let the period lattice \(L_{\rm per}\) act by translation on the
argument \(v\). If a function \(\Psi\) satisfies
\begin{equation}
                  \Psi(v+l)-\Psi(v)=a_l(v)
                       \in\operatorname{Aff},           \hypertarget{tag-021N}{}\label{maps:eq:D.OF7}
\end{equation}
its absolute class on \(\widetilde N\) is \(L_{\rm per}\)-invariant.
It descends to \(N_{L_{\rm per}}\): in total degree two the covering
cochain double complex has \(H^1(\widetilde N;k)=0\), and
the only possible obstruction lies in \(H^3(L_{\rm per};k)=0\).
The choices of descent form a torsor for \(H^2(L_{\rm per};k)=k\).

The boundary group and the two marked long loops give the split extension
\begin{equation}
       \pi_1(M_{L_{\rm per}})=\mathbb Z^3\rtimes L_{\rm per},                    \hypertarget{tag-021O}{}\label{maps:eq:D.OF8}
\end{equation}
where the normal \(\mathbb Z^3\) is generated by the two shorts
and \(\theta\), and the two long loops give the specified splitting.
Its dual module \(H^1(\widetilde M;k)\) is exactly
\(\operatorname{Aff}\): a covector $\alpha$ corresponds to the
function $v\mapsto-\alpha(\theta+v)$.  Evaluation and negation give
its inverse, and translating $v$ gives the deck action.
For boundary classes with zero restriction to \(\widetilde M\),
the degree-two filtration is the short exact sequence
\begin{equation}
 0\longrightarrow H^2(L_{\rm per};k)\longrightarrow
 \ker\{H^2(M_{L_{\rm per}};k)\to H^2(\widetilde M;k)\}
 \longrightarrow H^1(L_{\rm per};\operatorname{Aff})\longrightarrow0.
                                                               \hypertarget{tag-021P}{}\label{maps:eq:D.OF9}
\end{equation}
The splitting of \eqref{maps:eq:D.OF8} kills the degree-two transgression;
there is no degree-three group cohomology of \(L_{\rm per}\).
These two observations also prove \eqref{maps:eq:D.OF9} directly from the
two-dimensional free resolution of the free abelian group \(L_{\rm per}\) in the ordered basis \((8,-6),(24,-16)\).
Restriction to the long two-torus splits its leftmost term.

The image in the rightmost group of the boundary value of a
descended support is represented by the full affine cocycle
\(l\mapsto a_l\) in \eqref{maps:eq:D.OF7}. This follows either by chasing
\eqref{maps:eq:D.OF6} in that resolution, or by choosing the relative
cochain with positive Thom coefficients $-\Psi$.  Its deck difference
has Thom coefficients $-a_l$, so the boundary covector has value
$-a_l(v)$ on $\theta+v$.  Under the negative-evaluation identification
just specified, its function is exactly $a_l(v)$.
Thus the gradients, constants, and one long-torus value determine
the entire boundary class. This statement retains every
\(\operatorname{Ext}(H_1(M_{L_{\rm per}}),k)\) ambiguity.

In our positive fan convention a core cochain with short-long
matrix \(E_{\rm sl}\), zero angular area and zero short-normal strip values
has affine cocycle
\begin{equation}
 a_{l_j}(v)=-E_{\rm sl}(v,b_j)-c_j^{\rm str},\qquad j=1,2,                  \hypertarget{tag-0218}{}\label{maps:eq:D.OF10}
\end{equation}
where \(b_1=e_2,b_2=\gamma\) and \(l_j=(A-I)b_j\).
The line comparison fixes both signs as follows.  Put $h=-\Psi$
and $D_h=\sum_v h(v)D_v$.  Write $r=8a+r_0$, $0\leq r_0<8$.
For $l_1=(8,-6)$ and $l_2=(24,-16)$ the exact integral identities are
\[
\begin{array}{c|cc}
 j&m_j(v)&k_j(v)\\\hline
 1&r-4&r_0-a\\
 2&-4&3r_0-1
\end{array},\qquad h(v+l_j)-h(v)=m_j(v)+8k_j(v).
\]
Let $F_{{\rm per},j}$ be the marked period map on the common cover, retaining
its unit factors.  Pullback gives the equality of divisor
line bundles
\[
 F_{{\rm per},j}^*\mathcal O(D_h)=\mathcal O(D_h)\otimes
 \mathcal O(\operatorname{div}\chi^{m_j})\otimes
 \mathcal O(D_{k_j})^{\otimes8}.
\]
The explicit line map is
\[
 \chi^{-m_j}:\mathcal O(D_h)\otimes\mathcal O(D_{k_j})^{\otimes8}
                  \longrightarrow F_{{\rm per},j}^*\mathcal O(D_h),
\]
and its inverse is multiplication by $\chi^{m_j}$.  Local frames
retain reciprocal eighth powers from the displayed tensor factor.
Its first Chern class contributes zero modulo eight, giving the
finite-coefficient descent with its eighth tensor factor.  The two period comparisons
commute with their full tensor factors: for example
$\Delta_{l_2}m_1=24$ and $\Delta_{l_2}k_1=-3$, whose sum with
the factor eight is zero.

For the oriented short-angle/long rectangle $(\vartheta,t)$,
an end-to-start transition $\exp(\mathrm i \nu\vartheta)$, $\nu\in\mathbb Z$, has
connection $A_{\rm conn}=\mathrm i \nu t\,d\vartheta$ and curvature
$F_{\rm conn}=-\mathrm i \nu\,d\vartheta\wedge dt$.  Its Chern form
$\mathrm iF_{\rm conn}/(2\pi)$ evaluates to $\nu$.  The normal-angle/long
rectangle reads the constant in $m_j$ by the same formula.
Consequently the positive Thom potential has difference
$E_{\rm sl}(v,b_j)+c_j^{\rm str}$, and the support $\Psi=-h$ has difference
$-E_{\rm sl}(v,b_j)-c_j^{\rm str}$.  This is the full boundary cocycle in
\eqref{maps:eq:D.OF10}.  The resulting fibre and strip evaluations include $u(e_1\wedge e_2)=1$.
The already included basis-loop prism correction gives
$c_1^{\rm str}=c_2^{\rm str}=4$, for which $-c_j^{\rm str}=c_j^{\rm str}$ in $k$.

For \eqref{maps:eq:D.OF2}--\eqref{maps:eq:D.OF4}, \(E_{\rm sl}=\left(\begin{smallmatrix}1&0\\0&0
\end{smallmatrix}\right)\). The required affine functions are
therefore
\begin{equation}
             a_{(8,-6)}=-r+4,\qquad a_{(24,-16)}=4.
                                                               \hypertarget{tag-0219}{}\label{maps:eq:D.OF11}
\hypertarget{tag-021U}{}
\end{equation}
\SsixDisplayAliases{maps:eq:D.OF11}{\label{maps:eq:D.OL2}}
The support \eqref{maps:eq:D.OF1} has exactly these differences.
For all integers, write \(r=8a+r_0,s=2b+s_0\) with \(0\le r_0<8\) and
\(0\le s_0<2\), and substitute. The two differences are
\(4-9r_0\) and \(12-24r_0\), respectively, which reduce to
\eqref{maps:eq:D.OF11}.
Choose the remaining descent parameter so its value on the
long torus \(\langle e_2,\gamma\rangle\) is zero,
as it is for \eqref{maps:eq:D.OF2}.
Equations \eqref{maps:eq:D.OF9}--\eqref{maps:eq:D.OF11} now prove equality on the full
cusp boundary, including both its finite-coefficient and
integral-torsion contributions.

The short-long marking also follows from the annulus map
of \ref{maps:src:D.2.43}.
If \(x_{\rm gr}^p y_{\rm gr}^q-1=(x_{\rm gr}-1)S_{x_{\rm gr}}+(y_{\rm gr}-1)S_{y_{\rm gr}}\), the angular vector \((a,b)\)
closes to the \(I^2\)-class with coefficients
\begin{equation}
             (aS_{x_{\rm gr}},\ bS_{x_{\rm gr}}+aS_{y_{\rm gr}},\ bS_{y_{\rm gr}})                   \hypertarget{tag-021A}{}\label{maps:eq:D.OF12}
\end{equation}
on \(((x_{\rm gr}-1)^2,(x_{\rm gr}-1)(y_{\rm gr}-1),(y_{\rm gr}-1)^2)\).
Its Picard columns are
\begin{equation}
 x_{\rm gr}^{-1}y_{\rm gr}^{-2}E_2+y_{\rm gr}^{-1}E_3,\quad
       -x_{\rm gr}^{-1}y_{\rm gr}^{-1}E_2,\quad
       x_{\rm gr}^{-1}E_1+x_{\rm gr}^{-2}y_{\rm gr}^{-1}E_2.                        \hypertarget{tag-021B}{}\label{maps:eq:D.OF13}
\end{equation}
In the present positive component labels \(x_{\rm gr}\) translates by
\((-1,0)\), and \(y_{\rm gr}\) by \((0,1)\); the path and angular exponents
are changed simultaneously. Substituting \eqref{maps:eq:D.OF1} into all the
shifted columns gives the integer values
\begin{equation}
                         \begin{pmatrix}9&24\\0&0\end{pmatrix}
                  \equiv E_{\rm sl}\pmod8.                       \hypertarget{tag-021C}{}\label{maps:eq:D.OF14}
\end{equation}
The path dependence in \(S_{x_{\rm gr}},S_{y_{\rm gr}}\) is a
Hilbert--Burch relation, which identifies the same integral annulus
class for every path choice.

\medskip\noindent\emph{Extension through the other caps.}
Each of the four other multiplicative fibres has primitive
rank-one monodromy. Its boundary \(H_1\) is free; so is the
ordinary cap boundary \(H_1\).
The regular-fibre specialization is onto the degree-two group
of each such cap: in a primitive nodal marking this is the
specialization from \(T^2\times T^2\) to
\((S^2\vee S^1)\times T^2\); for the ordinary cap it is the
identity on the central \(T^4\).
Their degree-two boundary kernels
are precisely the columns used in \(M\), including the
ordinary cap's affine-meridian column.
Equation \eqref{maps:eq:D.OF2} therefore makes the specified core cochain
extend across each such cap. The free boundary \(H_1\)
ensures that agreement of the resulting integral-cycle
functional is agreement of the full finite-coefficient
cohomology class. These choices and the cusp comparison
glue to the asserted \(u\in H^2(X;k)\).
It has $u(g)=-1$ for $g=-[e_1\wedge e_2]$.
Put $\tau=-g=[e_1\wedge e_2]$; negation is its own inverse,
and $u(\tau)=1$. Let $h_i$ be free lifts dual to the fifteen
classes in \eqref{maps:eq:D.PC10}. Send the $i$th free basis vector
to $h_i-u(h_i)\tau$ and $1\bmod8$ to $\tau$.
Changing an integer representative of $u(h_i)$ changes this image
by $8\tau=0$. The inverse reads the fifteen integral evaluations
and then $u$. Both composites are identity, giving the splitting
\begin{equation}
                         H_2(X)=\mathbb Z^{15}\oplus\mathbb Z/8
                                                               \hypertarget{tag-021D}{}\label{maps:eq:D.OF15}
\end{equation}
used for the fundamental-class coordinates.

\medskip\noindent\emph{The complete coordinate values.}
Let \(R_0,\ldots,R_5\) be the six positive fan rays in cyclic
order and define the wall values
\[
 d_j(v)=\Psi(v)+\Psi(v+R_j)
                    -\Psi(v+R_{j-1})-\Psi(v+R_{j+1}).
\]
An integral representative for \(u|_{P_v}\) is
\begin{equation}
 \eta_v=(d_0+d_1+d_2)H-d_0E_1-d_2E_2-d_4E_3.             \hypertarget{tag-021E}{}\label{maps:eq:D.OF16}
\hypertarget{tag-021T}{}
\end{equation}
\SsixDisplayAliases{maps:eq:D.OF16}{\label{maps:eq:D.OL12}}
Opposite wall restrictions agree modulo eight. The six wall differences
specify every component restriction and its integral lift. Changing
a lift by $8z$ changes its square by
$16\eta z+64z^2$, zero modulo sixteen.

For component labels \(i,j\), all mixed ordinary coordinates are
specified by
\begin{equation}
 \begin{aligned}
 \langle ux_i^2,[X]\rangle
                    &=-\sum_{a=0}^{5}d_a(v_i)\pmod8,\\
 \langle ux_ix_j,[X]\rangle
                    &=\sum_{\{a:\ v_i+R_a=v_j\ {\rm mod}\ L_{\rm per}\}}
                                       d_a(v_i)\pmod8\quad(i\ne j),\\
 \langle ux_i y,[X]\rangle
                    &=(-1)^{s_i}\eta_i\cdot(H-E_1)\pmod8.
 \end{aligned}                                             \hypertarget{tag-021F}{}\label{maps:eq:D.OF17}
\end{equation}
The last value is four for component labels \(0,1,14,15\)
and zero for the other twelve. These formulas use Cartier restrictions and retain the repeated vertical edges.
On the primitive free basis \(x_0,\ldots,x_{13},y\), the remaining
power coordinates are
\begin{equation}
 \begin{split}
 \langle\mathcal P_2(u)x_i,[X]\rangle&=7\pmod{16}
                                             \quad(0\le i<14),\\
 \langle\mathcal P_2(u)y,[X]\rangle&=0\pmod{16},\\
 \langle u^3,[X]\rangle&=0\pmod8,\qquad
 \langle uy^2,[X]\rangle=0\pmod8.
 \end{split}                                               \hypertarget{tag-021G}{}\label{maps:eq:D.OF18}
\hypertarget{tag-021V}{}
\end{equation}
\SsixDisplayAliases{maps:eq:D.OF18}{\label{maps:eq:D.OL3}}
The first line is \(\eta_i^2\) on the normalized component.
The two omitted component labels have square fifteen modulo
sixteen, consistently with their primitive-basis relations.
The remaining three evaluations, including the full
order-sixteen value which cannot be obtained by division,
are proved in the following localization step.

\item \emph{The finite-coefficient evaluations.}
Use the global class $u$ and the entire boundary comparison just proved,
with support \eqref{maps:eq:D.OF1} and affine differences
\eqref{maps:eq:D.OF11}. Both differences have zero slope in the
scalar $s$ direction. Let $y$ be the primitive class in
\eqref{maps:eq:D.PC9}. It remains to prove the last three evaluations
in \eqref{maps:eq:D.OF18}, including the full order-sixteen value.

The scalar circle action is semifree. Its fixed locus consists of
the sixteen vertical toric curves $C_{r,s}$ joining vertices
$(r,s)$ and $(r,s+1)$, with $(r,s)$ indexed as in \eqref{maps:eq:D.PC6}.
Each has normal bundle $\mathcal O(-1)\oplus\mathcal O(-1)$
with weights $1,-1$. In every other triangular toric chart the
nonzero weights are also primitive. At each of the four zero-free
$I_1$ caps the primitive angular collapsing vector has first
coordinate one in its elliptic/scalar plane; the scalar circle
therefore projects primitively to the transverse torus and is
free. The ordinary torus cap is free as well. Thus the complement of these sixteen fixed curves is free.

Since $X$ is simply connected, both $u$ and $y$ have circle-equivariant
lifts. Indeed the only outgoing differentials from ordinary degree
two in the Borel spectral sequence have targets involving
$H^1(X)$ or $H^3(BS^1)$, both zero. Let $x_C$ be the positive
generator of $H^2(C_{r,s};\mathbb Z)$ and let $t$ be the generator
of $H^2(BS^1;\mathbb Z)$. Write
\begin{equation}
 \widehat u|_{C_{r,s}}=d_{r,s}x_C+\lambda_{r,s}t,
 \qquad
 \widehat y|_{C_{r,s}}=\mu_{r,s}t.                         \hypertarget{tag-021W}{}\label{maps:eq:D.OL4}
\end{equation}
There is no $x_C$ term in the second expression: \eqref{maps:eq:D.PC14} and
$f\cdot C_{r,s}=0$ give $y(C_{r,s})=0$.

The ordinary degrees and fibre characters have the following formulas. The degree of the support on a vertical wall is
\begin{equation}
 d_{r,s}=\Psi(r,s)+\Psi(r,s+1)
              -\Psi(r+1,s)-\Psi(r-1,s+1)\pmod8 .         \hypertarget{tag-021X}{}\label{maps:eq:D.OL5}
\end{equation}
Choosing representatives from zero to seven gives
\begin{equation}
\begin{array}{c|cc}
 r&d_{r,0}&d_{r,1}\\\hline
 0,7&4&7\\
 1,\ldots,6&0&3.
\end{array}                                                \hypertarget{tag-021Y}{}\label{maps:eq:D.OL6}
\end{equation}
The circle used here is generated by $\delta$, the difference
of the upper and lower vertical ray vectors.  In the source
coordinate $W$ it acts by $W\mapsto g^{-1}W$, as verified in
\eqref{maps:eq:D.OR3}; this fixes its relation to physical line rotation.
The positive divisor coefficients on the common cover are $h=-\Psi$,
as proved by the full boundary comparison.  On a cone let $m_\sigma$
take these values $h$ on its three rays.  The rational Cartier frame
is $e_\sigma=\chi^{-m_\sigma}$.  The canonical fibre map is
\[
 g\,e_\sigma(p)=\chi^{m_\sigma}(g)e_\sigma(gp);
\]
it follows by applying $(g\cdot s)(p)=s(g^{-1}p)$ to sections.
Its inverse is the same formula with $g^{-1}$.  Therefore the
scalar weight is $m_\sigma(\delta)$, which here equals
$-\bigl(\Psi(r,s+1)-\Psi(r,s)\bigr)$.  The eighth tensor factors
in the period comparison do not change this weight modulo eight.
The characters are
\begin{equation}
                 \lambda_{r,0}=0,\qquad
                 \lambda_{r,1}=-3r .                      \hypertarget{tag-021Z}{}\label{maps:eq:D.OL7}
\end{equation}
This is periodic under the period lattice because of
\eqref{maps:eq:D.OL2}. The common character is determined by comparison along the
toric orbit curves. Their degrees determine the differences
of the characters at their fixed endpoints. Horizontal and diagonal
orbit curves connect the sixteen fixed curves in a connected graph:
they join $(r,s)$ to $(r+1,s-1)$ and to $(r-1,s)$ respectively.
The support realizes all these differences. Consequently any global
equivariant lift has the characters in \eqref{maps:eq:D.OL7}, with one common additive character $\ell_u\in\mathbb Z/8$.
Subtracting $\ell_u t$ from $\widehat u$, with inverse addition,
chooses the displayed characters without changing $u$.
Each connecting orbit curve has scalar tangent weight $1$ or $-1$.
Its integral equivariant degree-two group has the two generators
given by the curve class and the character class, so reduction
modulo eight is onto.  For a coordinate $z\mapsto gz$ at the first endpoint, the degree-$d$
frames satisfy $e_0=z^{-d}e_\infty$, with inverse
$e_\infty=z^d e_0$.  Equivariance gives
$\lambda_0-\lambda_\infty=d$.  Reversing the coordinate gives the
same relation multiplied by its tangent weight.  The inverse
reads the degree and either endpoint character.  For example the
horizontal curve through $(1,0)$ has degree six and endpoint
characters $0,-6$; their difference is six.  These exact maps prove the same
assertion for the finite coefficient class $u$.  The horizontal
and diagonal connections above form a connected graph: one changes
$r$ with $s$ fixed and the other changes both $r$ and the parity of
$s$.  Hence their differences leave exactly one common character.

The integral component divisors have their canonical equivariant
linearizations. Along $C_{r,s}$ the two adjacent component normal
weights are $-1,+1$ in the lower/upper vertex order. Thus the weight
of the odd-component sum in \eqref{maps:eq:D.PC9} is $+1$ for $s=0$ and $-1$
for $s=1$. The difference
$2\widehat y-\sum_r\widehat x_{r,1}$ is a global multiple $ct$.
The existence of integral fixed-point weights forces $c$ odd.
Replace $\widehat y$ by $\widehat y-(c+1)t/2$; this is integral,
and its inverse adds $(c+1)t/2$. It sends $c$ to $-1$ and leaves
the ordinary class $y$ fixed. The resulting weights are
\begin{equation}
                       \mu_{r,0}=0,\qquad
                       \mu_{r,1}=-1.                     \hypertarget{tag-0220}{}\label{maps:eq:D.OL8}
\end{equation}
Every other lift adds one common integer $\ell_y$.

Each restriction of $\widehat u$ in \eqref{maps:eq:D.OL4} lifts integrally on
$C_{r,s}\times BS^1$. Thomas's cochain formula therefore gives
$\mathcal P_2(\widehat u)=(d x_C+\lambda t)^2$ with the full
modulus sixteen \cite[equations~(2.1), (3.4), pp.~267--268]{ffThomas56}.
The ordinary cube uses modulus eight. The equivariant normal Euler
class is
\[
                           (-x_C+t)(-x_C-t)=-t^2.
\]
Localization is valid over these coefficient rings because the
complement of the fixed locus is free and all normal weights are
units. The inverse of restriction after inverting $t$ is
$\sum i_{r,s*}(z_{r,s}/(-t^2))$, by the normal Thom formula;
the free complement is killed by a power of $t$.
Taking the coefficient of $x_C$ gives
\begin{equation}
 \begin{aligned}
 \langle u^3,[X]\rangle&=-3\sum_{r,s}d_{r,s}\lambda_{r,s}^2
                                                        &&\pmod8,\\
 \langle\mathcal P_2(u)y,[X]\rangle
       &=-2\sum_{r,s}d_{r,s}\lambda_{r,s}\mu_{r,s}
                                                        &&\pmod{16},\\
 \langle uy^2,[X]\rangle&=-\sum_{r,s}d_{r,s}\mu_{r,s}^2
                                                        &&\pmod8.
 \end{aligned}                                             \hypertarget{tag-0221}{}\label{maps:eq:D.OL9}
\end{equation}
The degree-zero base ring injects into its localization, so these
are identities before inverting $t$ as well.

The representatives \eqref{maps:eq:D.OL6}--\eqref{maps:eq:D.OL8} satisfy the exact integer sums
\begin{equation}
 \sum d_{r,s}=40,\qquad \sum d_{r,s}\lambda_{r,s}=-336,
 \qquad \sum d_{r,s}\lambda_{r,s}^2=5544,
 \qquad \sum_r d_{r,1}=32.                               \hypertarget{tag-021R}{}\label{maps:eq:D.OL10}
\end{equation}
The three values in \eqref{maps:eq:D.OL9} are consequently $-16632$, $-672$,
and $-32$, proving \eqref{maps:eq:D.OL3} with their full moduli.
Adding characters sends $\lambda\mapsto \ell_u+\lambda$ and
$\mu\mapsto \ell_y+\mu$.  The following integer identities also give
the algebraic substitution $\lambda\mapsto \ell_u+\epsilon\lambda$,
$\epsilon=\pm1$:
\begin{equation}
 \begin{aligned}
 -120\ell_u^2+2016\epsilon \ell_u-16632,\\
 -80\ell_u\ell_y+64\ell_u+672\epsilon \ell_y-672\epsilon,\\
 -40\ell_y^2+64\ell_y-32.
 \end{aligned}                                             \hypertarget{tag-021S}{}\label{maps:eq:D.OL11}
\end{equation}
Every coefficient is divisible by eight, sixteen, and eight,
respectively.  In particular all global character additions preserve
the three evaluations in their full coefficient groups.  Replacing
the lift of $y$ by its character-one twist gives the equally explicit
weights $(1,0)$ in place of $(0,-1)$; the inverse subtracts that
single character.

The component restrictions in \eqref{maps:eq:D.OF16} and their
Cartier pairings give \eqref{maps:eq:D.OF17}; the preceding localization
completes \eqref{maps:eq:D.OF18} in its full coefficient groups.

In particular \(y^2=0\) integrally. Its free part is zero by
\eqref{maps:eq:D.PC12}, and its torsion part is detected injectively by the
primitive \(u\); the last equality of \eqref{maps:eq:D.OF18} kills that part.

Use the cyclic homology calculation \eqref{maps:eq:D.WI22} at $d=8$:
its degrees two, four and six have orders $8,16,8$, with their
positive reduced projective-space generators. The splitting
\eqref{maps:eq:D.OF15} sends the positive cyclic generator to $\tau$.
Its fifteen integral coordinate projections and cyclic projection induce
\[
 K(\mathbb Z^{15}\oplus\mathbb Z/8,2)
 \longrightarrow(\mathbb{CP}^{\infty})^{15}\times K(\mathbb Z/8,2).
\]
The inverse is induced by the sum of the fifteen free inclusions and
the cyclic inclusion. Both composites induce identity on $\pi_2$;
the Eilenberg--Mac Lane classification makes them inverse up to based
homotopy. The free factor has free even homology and no odd homology.
The integral K\"unneth map has no torsion-product term and gives
\begin{equation}
 H_6(K(\mathbb Z^{15}\oplus\mathbb Z/8,2);\mathbb Z)
       =\mathbb Z^{680}\oplus(\mathbb Z/8)^{120}
                         \oplus(\mathbb Z/16)^{15}
                         \oplus\mathbb Z/8.              \hypertarget{tag-021H}{}\label{maps:eq:D.OF19}
\end{equation}
The free coordinates are exactly the primitive integral cubic
\eqref{maps:eq:D.PC11}--\eqref{maps:eq:D.PC12}. Equations \eqref{maps:eq:D.OF17}--\eqref{maps:eq:D.OF18} give all
remaining coordinates: the first mixed group is detected by
\(ux_ix_j\), the second by the full Pontryagin square, and the
last by \(u^3\). More explicitly, write $\varepsilon_0,\ldots,\varepsilon_{14}$ for
the marked free generators, dual to $x_0,\ldots,x_{13},y$, and
write $g_8$ for $1\bmod8$. The classifying map for the inverse
homology marking takes values in $K(\mathbb Z^{15}\oplus\mathbb Z/8,2)$.
The inverse coordinate map sends $(T_\alpha,J_\beta,K_i,c_{\rm cyc})$ to
\[
 \sum_{|\alpha|=3}T_\alpha\gamma_\alpha(\varepsilon)
 +\sum_{|\beta|=2}J_\beta\gamma_\beta(\varepsilon)g_8
 +\sum_{i=0}^{14}K_i\varepsilon_i\gamma_2(g_8)+c_{\rm cyc}\gamma_3(g_8).
\]
The forward map is the four displayed detector families, whose
pairings with this cross-product basis are identity. Here $T$ is
the primitive cubic, $J$ is \eqref{maps:eq:D.OF17},
$(K_0,\ldots,K_{14})=(7,\ldots,7,0)$ in $(\mathbb Z/16)^{15}$,
and $c_{\rm cyc}=0$. This gives the complete $\mu$ and both coordinate maps.
\item \emph{The source sphere and refined class.}
\hypertarget{tag-03LK}{}\label{maps:proof:D.OS.refined}
{\normalfont\scriptsize tag 03LK}\par
Use the source sections, fan, line linearization and peripheral
markings of Definition~\ref{maps:def:weighted-os45-data}.  Write $C$ for the integral sum of the
complex critical curves, with the chain $B_C$ constructed above.
The boundary class used to construct the sphere has the following
integral prism description. Put \(\alpha_{\partial}=e_1\wedge\delta\)
and use \(A_0\) from \eqref{maps:eq:D.PC5}. At the first finite
\(I_1\) site the marked matrix is \(A_1=I+E_{12}\).
Direct exterior multiplication gives
\[
 (\wedge^2A_0-I)(e_1\wedge e_2+e_2\wedge\delta)=2\alpha_{\partial},
 \qquad
 (\wedge^2A_1-I)(e_2\wedge\delta)=\alpha_{\partial}.
\]
In the ordered six-cell basis above, every entry of
\(\wedge^2A_0-I\) is even. The row
\((0,1,0,0,0,0)\bmod2\) kills its image and evaluates to one on
\(\alpha_{\partial}\); hence \(\alpha_{\partial}\) has exact order two in the cusp
fibre coinvariant group and in its injected term of the integral
mapping-torus sequence. The prism $T_1(e_2\wedge\delta)$ in the finite-filled exterior
has boundary $\alpha_{\partial}$ by the second displayed equality.
At the weighted cusp the fixed lattice is
$\mathbb Ze_1\oplus\mathbb Z\delta$. In the order $(\delta,e_1)$,
its exterior generator is $-\alpha_{\partial}$. The angular-area sweep in
\eqref{maps:eq:D.OF5}, projected from the common cover, bounds this
fixed torus in the cusp; reversing that chain bounds $\alpha_{\partial}$.
The maps read the two marked angular coordinates and the radial
sweep coordinate, with the inverse parametrizations of
\ref{maps:src:D.2.43}. Thus the exterior and cusp images of this
same boundary class are both zero with their explicit chains.
\begin{enumerate}[label=(\alph*),leftmargin=*]
\item \emph{The source line and the selected cusp chart.}
Let $q_E$ be the based multiplicative elliptic period and use the multiplicative lifts of $P$ and $Q$.  The calculation underlying
\eqref{maps:eq:D.FI2} gives
\begin{equation}
 \sigma(Z)=\frac{\theta(Z/P;q_E)}{\theta(Z;q_E)},\qquad
 \sigma(q_E Z)=P\sigma(Z),\qquad e_{\rm flat}=s_{P-O}/\sigma .
 \hypertarget{tag-03LL}{}\label{maps:eq:D.OR1}
\end{equation}
Here $s_{P-O}$ is the canonical meromorphic section of
$\mathcal O(P-O)$.  If $W$ is its nonzero-line coordinate in $e_{\rm flat}$,
the quotient coordinate is $V=W^{-1}$, with inverse $W=V^{-1}$.  This is the tautological map from
the nonzero line to its dual; it inverts every line multiplier.

At the $I_8$ site the lifts have valuations $6$ and $24$.
The extending-frame heights are $F_i=\max(6-i,0)$ for $0\leq i<8$,
with $F_{i+8}=F_i-6$.  Thus the source-line graph has the edge
$(0,6)$--$(1,5)$.  It is an edge of the prescribed triangular fan.
The adjacent ordered cone
\[
 (1,0,5),\quad(1,1,5),\quad(1,0,6)
\]
has determinant one.  Its coordinates, and the extending
physical line coordinate, are
\begin{equation}
 q=xyz,\quad Z=y,\quad V=x^5y^5z^6,\quad
 W_{\rm ext}=\frac{q^6}{ZV}=x,\qquad
 (X,Y,W_{\rm ext})=(xz,y,x).
 \hypertarget{tag-03LM}{}\label{maps:eq:D.OR2}
\end{equation}
The inverse is $(x,y,z)=(W_{\rm ext},Y,X/W_{\rm ext})$ and its forward
Jacobian is $-x$.  The $W_{\rm ext}$ coordinate therefore records the line winding.
The extending frame is exactly $(Z/q^6)e_{\rm flat}$: its orders on both
branches vanish, $P$ and $O$ miss the node, and the frame and its inverse
extend over the remaining codimension-two point.

The marked weights in this cone are
\begin{equation}
 e_1=(-1,1,0),\qquad \delta=(-1,0,1),\qquad
 a_0=e_1-\delta=(0,1,-1).
 \hypertarget{tag-03LN}{}\label{maps:eq:D.OR3}
\end{equation}
Thus $\delta$ is inverse physical line rotation and $a_0$ is the source
vanishing circle at constant $W_{\rm ext}$.  At the first $I_1$ site its
zero section cocycle gives the extending flat frame and the vanishing
circle $a_1=e_1$.  Both oriented boundary areas are
$a_i\wedge\delta=\alpha_{\partial}=e_1\wedge\delta$; changing the common boundary
orientation changes both signs.  The two meridians $(1,-1),(1,0)$ have
determinant one.

\item \emph{The based collar and its embedded odd sphere.}
Choose a narrow neighborhood of the prescribed based arc from this cusp
to the first $I_1$ site, containing no other singular site.  The two based
elliptic matrices are $T^8,T$ and their $G$ cocycles are $(3,0),(0,0)$.
They fix the same short circle.  Thus the based exponentials $q_E,P,Q$
are single-valued on this punctured neighborhood, with valuations
$(8,6,24)$ and $(1,0,0)$ at the two ends.  If the forward suspension
lift has coefficient $c_F(Z)$ in the meromorphic section, its coefficient
in the flat frame is
\begin{equation}
 \kappa(t)=c_F(Z)\frac{\sigma(QZ)}{\sigma(Z)}.
 \hypertarget{tag-03LO}{}\label{maps:eq:D.OR4}
\end{equation}
The cancellation of its divisor and automorphy makes this a nonzero
holomorphic function on the compact elliptic fibre, hence independent of
$Z$.  Its valuations are $16,0$, with the inverse-lift convention of
\eqref{maps:eq:D.FI2}: $w=-16$ and $8w-(-6)24=16$ at the cusp.
The quotient maps are therefore
\begin{equation}
 (Z,V)\longmapsto(q_E Z,P^{-1}V),\qquad
 (Z,V)\longmapsto(QZ,\kappa^{-1}V).
 \hypertarget{tag-03LP}{}\label{maps:eq:D.OR5}
\end{equation}
These maps include all phases and nonzero holomorphic units.  The origin
is the marked section $O$, namely $Z=1$, and the frame is
\eqref{maps:eq:D.OR1}; these are also the endpoint choices.

Parametrize the embedded based arc by a smooth regular $p:[0,1]\to B$,
linear in affine coordinates near its ends.  On a sufficiently thin strip
the explicit map $t(s+iu)=p(s)+iu p'(s)$ is an orientation-preserving
tubular diffeomorphism; along the real interval its Jacobian is
$|p'(s)|^2$.  Write the nonzero endpoint units as
$q_0(t(f))=fc_0(f)$ and $q_1(t(f))=(1-f)c_1(f)$.
On $\mathcal Q=\{ad-bc=1\}$, put $f=ad$ and define
\begin{equation}
 \Phi(a,b,c,d)=[\,t(f),\ Z=ab,\ V=f^6/a\,],
 \hypertarget{tag-03LQ}{}\label{maps:eq:D.OR6}
\end{equation}
where the brackets mean exactly \eqref{maps:eq:D.OR5}.  At the cusp this
extends in the selected fan chart by
\[
 (x,y,z)=(c_0^6/b,\ ab,\ c_0^{-5}d);
\]
at the other endpoint it extends in the source-node/line chart by
\begin{equation}
 (X,Y,W)=(-c_1c/a,\ ab,\ a/f^6).
 \hypertarget{tag-03LR}{}\label{maps:eq:D.OR7}
\end{equation}
The inverse at the cusp reads
$b=c_0^6/x$, $a=xy/c_0^6$, $d=c_0^5z$ and
$c=(ad-1)/b$.  At the other end it reads
$a=Wf^6$, $b=Y/a$, $c=-Xa/c_1$ and $d=f/a$.
The inverse base coordinate determines $f$ in each formula.
Along the two central core circles the Jacobians are respectively
$c_0/b$ and $c_1$, hence nonzero; the endpoint neighborhoods are
chosen small enough to retain this property.
On the middle collar the fibre character matrix on $(a,b)$ is
$\left(\begin{smallmatrix}1&1\\-1&0\end{smallmatrix}\right)$, of determinant
one, and the base derivative preserves orientation.

The compact real form
\begin{equation}
 N_{{\rm sph},0}=\left\{
 \begin{pmatrix}\alpha&\beta\\-\overline\beta&\overline\alpha\end{pmatrix}:
 |\alpha|^2+|\beta|^2=1\right\}
 \hypertarget{tag-03LS}{}\label{maps:eq:D.OR8}
\end{equation}
has generic short two-torus fibre, with the two meridians in
\eqref{maps:eq:D.OR3}.  The physical line transition is
$W_1=(Z/f^6)W_0$; its $Z$ factor has winding one and its $f^6$ factor
is in \eqref{maps:eq:D.OR6}.  These are the two Hopf charts of
the displayed $S^3$.
The map is injective on $N_{{\rm sph},0}$: on a regular fibre the real short plane
intersects the full period lattice in exactly
$\mathbb Ze_1+\mathbb Z\delta$, by real independence of the four marked
periods; distinct arc parameters have distinct base images.  At either
end no nonzero period translation fixes the selected finite edge or core
circle.  The endpoint Jacobians prove immersion.  Compactness then gives
an embedded sphere and a neighborhood on which $\Phi$ is an embedding.

Left multiplication by
$L_\varepsilon=\left(\begin{smallmatrix}1&\varepsilon\\0&1\end{smallmatrix}\right)$,
for sufficiently small positive $\varepsilon$, gives
\begin{equation}
 a=\alpha-\varepsilon\overline\beta,\quad
 b=\beta+\varepsilon\overline\alpha,\quad
 c=-\overline\beta,\quad d=\overline\alpha.
 \hypertarget{tag-03LT}{}\label{maps:eq:D.OR9}
\end{equation}
It misses $a=d=0$ and $b=c=0$, and the neighborhood meets no other
critical curve.  It commutes with the scalar right action, so its image
$N_{\rm sph}$ is a free scalar-invariant embedded sphere disjoint from $C$.
The unperturbed sphere has Mayer--Vietoris boundary $\alpha_{\partial}$.
Its oddness is detected by the following integral boundary character.
The matrix $A_0$ of \eqref{maps:eq:D.PC5} is the identity modulo
$2$. In the ordered bivector basis
$(e_1\wedge e_2,e_1\wedge\delta,e_1\wedge\gamma,
 e_2\wedge\delta,e_2\wedge\gamma,\delta\wedge\gamma)$, the row
$(0,1,0,0,0,0)$ modulo $2$ annihilates every column of
$\bigwedge^2A_0-I$ and takes $\alpha_{\partial}$ to $1$.
The cusp boundary $\partial N\to S^1$, over the positive
$q$-meridian and with marked fibre $T^4$, has return $A_0$.
Its Wang sequence is
\[
0\longrightarrow\operatorname{coker}(\bigwedge^2A_0-I)
\longrightarrow H_2(\partial N;\mathbb Z)
\longrightarrow\ker(A_0-I\mid H_1(T^4;\mathbb Z))
\longrightarrow0.
\]
Its right group is the fixed lattice
$\mathbb Z(e_1,\delta)$.  On these two columns the ordered bar
quadratic correction of $A_0$ is zero.  Their marked prism sweeps
therefore split the sequence: the forward map reads the fixed
sweep, and the inverse subtracts that sweep to recover the fibre
coinvariant.  Extend the displayed mod-two character by zero on
these two fixed-sweep generators.  It is then a homomorphism on
all of $H_2(\partial N;\mathbb Z)$. In particular $\alpha_{\partial}$ is not twice
any boundary class. Applying the character to the Mayer--Vietoris
image proves $[N_{\rm sph}]\notin2H_3(X;\mathbb Z)$. Since this last group is
$\mathbb Z/8$, the sphere is an odd generator. The perturbation
preserves its class.

\item \emph{The exact reference frame and its local degree.}
Let $\rho$ be the reference from the abstract kernel and target bundles
in \eqref{maps:eq:D44.5}.  Fix the nowhere-zero abstract section
$\delta$.  Comparing a global second kernel vector with the local one
changes the quotient complex-line frame by a map to $\mathbb C^*$;
the choices of its lift form an affine space.  The target-frame change
also maps to $\mathbb C^*$.  On $N_{\rm sph}\simeq S^3$ these changes have zero
degree-three homotopy class.  Splittings of the differential sequence
form an affine space as well.  Thus the global reference and the
local reference below have zero integral degree-three discrepancy.

This comparison remains valid along the smooth middle base map.
In vertical/horizontal splittings the derivative of \eqref{maps:eq:D.OR6}
has complex-linear vertical block and real base block
$a+b\,\text{conjugation}$ with $|a|>|b|$.  Replacing the latter by
$a+(1-r)b\,\text{conjugation}$ and contracting the horizontal-to-vertical
block gives invertible oriented bundle maps throughout, fixed at the
holomorphic ends and retaining $\delta$ and the vertical kernel.

At $g=\left(\begin{smallmatrix}a&b\\c&d\end{smallmatrix}\right)$,
left multiplication identifies the Lie algebra with the tangent space;
its inverse multiplies by
$g^{-1}=\left(\begin{smallmatrix}d&-b\\-c&a\end{smallmatrix}\right)$.
The scalar vector is $g\operatorname{diag}(-1,1)$, and
\[
 df\bigl(g(\zeta_1E_{12}+\zeta_2E_{21})\bigr)
                          =ac\zeta_1+bd\zeta_2.
\]
This is the left complex Lie frame, with its two off-diagonal coordinates.
Set $a_{\rm B}=ac$, $b_{\rm B}=bd$.  The Bott matrix with first column
$(b_{\rm B},-a_{\rm B})/\sqrt{|a_{\rm B}|^2+|b_{\rm B}|^2}$ gives the local reference, after the affine
splitting changes just described.  On \eqref{maps:eq:D.OR9},
\begin{equation}
 a_{\rm B}=-\alpha\overline\beta+\varepsilon\overline\beta^2,
 \qquad b_{\rm B}=\beta\overline\alpha+\varepsilon\overline\alpha^2.
 \hypertarget{tag-03LU}{}\label{maps:eq:D.OR10}
\end{equation}
For $0\leq s\leq1$, set
\[
 F^{\rm hom}_s=\bigl((1-s)\beta\overline\alpha+\varepsilon\overline\alpha^2,
          (1-s)\alpha\overline\beta-\varepsilon\overline\beta^2\bigr).
\]
At $s=0$ this is $(b_{\rm B},-a_{\rm B})$. It has no common zero: off the axes
the equations imply $((1-s)^2+\varepsilon^2)\overline\alpha=0$,
and on either axis the remaining square term is nonzero.
The polar map $z\mapsto(z/\|z\|,\|z\|)$ has inverse
$(w,r)\mapsto rw$. Its first component applied to $F^{\rm hom}_s$ gives
the homotopy to the radial map of
$(\overline\alpha^2,-\overline\beta^2)$.
The proper holomorphic quadratic map $(z_1,z_2)\mapsto(z_1^2,-z_2^2)$
has four regular preimages, each of positive real local degree.
Simultaneous conjugation of both complex coordinates has real
determinant one. These maps therefore give degree four for the Bott matrix.  Consequently $\rho$ induces the odd framed
parity of a Lie-parallelized $S^3$.

The coefficient convention is fixed by the quaternionic left-multiplication
frame: its stable difference from the bounding frame is one unit.
An oriented rank-three normal-frame change stabilizes with factor two.
Indeed the adjoint complex roots are $2x,0,-2x$, so
$c_2(\operatorname{ad}_{\mathbb C})=4c_2(\mathrm{fund})$ and
$\lambda(\operatorname{ad}_{\mathbb R})=-2c_2(\mathrm{fund})$;
the quaternionic rank-four representation instead has
$\lambda=-c_2(\mathrm{fund})$.
These are the precise maps in
\cite[arXiv version~2, Remark~2.17, p.~12; \S6.1, pp.~24--25]{Redden2011StringStructures}.
Thus the resulting parity is independent of the oriented normal frame.

\item \emph{The integral equivariant characteristic class.}
Use $t=c_1(\mathcal O(1))\in H^2(BS^1;\mathbb Z)$.  On each of the
sixteen fixed curves let $x_C$ be its positive generator.  The tangent and normal data give roots
$2x_C,-x_C+t,-x_C-t$, with $x_C^2=0$.
The ordinary first Chern class is zero, and the integral Gysin sequence
of $X\times ES^1\to X_{S^1}$ makes the kernel in degree two the
multiples of $t$.  The sum of the fixed weights is zero; hence
$c_{1,S^1}(TX)=0$ and $\lambda_{S^1}=-c_{2,S^1}$.
Its restriction to every fixed curve is $t^2$.

As the ordinary $\lambda$ vanishes, that same sequence gives
\begin{equation}
 \lambda_{S^1}(TX)-t^2=t\eta,\qquad
 0\longrightarrow\mathbb Zt\longrightarrow H^2_{S^1}(X;\mathbb Z)
 \longrightarrow\mathbb Z^{15}\longrightarrow0.
 \hypertarget{tag-03LV}{}\label{maps:eq:D.OR11}
\end{equation}
For this Gysin sequence and its Thom-pair construction see
\cite[\S4.D, pp.~438,444]{Hatcher2002}; its total space here is
homotopy equivalent to $X$ and its Euler class is the displayed $t$.
Restriction to a fixed curve gives $t\eta=0$ and hence $\eta=0$ there.
Over the rationals the Borel spectral sequence has only even entries,
so all its differentials vanish.  Multiplication by $t$ is injective
on its associated graded and hence there is no $t$-torsion.
Localization, in the precise form
\cite[Theorem~3.5 and the circle specialization, p.~8]{AtiyahBott1984},
now gives $\eta=0$ rationally.  The integral group in
\eqref{maps:eq:D.OR11} is free, so $\eta=0$ integrally:
\begin{equation}
                         \lambda_{S^1}(TX)=t^2.
 \hypertarget{tag-03LW}{}\label{maps:eq:D.OR12}
\end{equation}
This is an integral characteristic-class identity.

\item \emph{The Borel pair and the framing transport.}
The scalar action on the sphere coordinates is
$(\alpha,\beta)\mapsto(s^{-1}\alpha,s\beta)$.
The explicit map $v(\alpha,\beta)=(\overline\alpha,\beta)$ identifies
it with the scalar Hopf sphere and is its own inverse. Conjugating
its first complex coordinate has real determinant $-1$, which is
its orientation sign on the unit sphere. The orientations below are
fixed by the positive scalar circle and quotient. Form
\[
 W_{\rm B}=(N_{\rm sph}\times D^2)/((n,z)\sim(sn,sz)),\qquad
 \partial W_{\rm B}\longrightarrow N_{\rm sph},\quad[n,z]\longmapsto z^{-1}n.
\]
With the diagonal convention for the Borel quotient the map of
pairs into $(X_{S^1},X)$ is
\begin{equation}
 [n,z]\longmapsto
 [\,n,(z,\sqrt{1-|z|^2}\,v(n))\,].
 \hypertarget{tag-03LX}{}\label{maps:eq:D.OR13}
\end{equation}
At $|z|=1$ it is $[z^{-1}n,(1,0,0)]$, exactly the ordinary inclusion
after the displayed boundary identification.  The pulled-back tangent
bundle is the pullback of $(TX|_{N_{\rm sph}})/S^1$ from $N_{\rm sph}/S^1$; on the boundary
its vector identification is the derivative $d(z^{-1})$.

The collapsed map
\[
 W_{\rm B}/N_{\rm sph}\longrightarrow\mathbb{CP}^2,\qquad
 [n,z]\longmapsto[z:\sqrt{1-|z|^2}\,v(n)]
\]
is a homeomorphism.  Off $[1:0:0]$, a unit representative $(Z_0,Z')$
has inverse $[v^{-1}(Z'/\|Z'\|),Z_0]$; changing the unit representative
by a phase is exactly the defining equivalence in $W_{\rm B}$.
The collapsed boundary supplies the remaining point.
Thus the square of the base generator evaluates to a unit on $(W_{\rm B},N_{\rm sph})$.

Orient $N_{\rm sph}$ by its positive scalar circle followed by the positive
orientation of the quotient sphere $\Sigma$.  Use the character
generator $t$ for which $\langle t,[\Sigma]\rangle=1$ under the
displayed Hopf identification.  The inverse chart of the collapsed
pair above has the corresponding relative unit.  All tangent
and normal quotient maps here come from the explicit equivariant
embedding \eqref{maps:eq:D.OR6}--\eqref{maps:eq:D.OR9}.

\item \emph{The complete critical-chain intersection.}
Write $\mathcal B=X_{S^1}$ and $\mathcal B^\circ=(X\setminus C)_{S^1}$.
The scalar section is nonzero in the abstract differential kernel,
including its $xy$ and $xyz$ syzygy frames.  Its quotient line and
the corrected target line have ordinary first Chern class zero.
Their fixed-point weights are zero, so their equivariant first
Chern classes vanish as well: the kernel of restriction to
ordinary degree-two cohomology is the character line, and its
coefficient is read at a fixed point.  The equivariant critical-ideal
sequence therefore gives
\[
 [C]_{S^1}=c_{2,S^1}(TX)=-t^2.
\]
This identity is integral.  At $xy$ the two normal columns are
$(x,-y)$ and $(\overline y,\overline x)/(|x|^2+|y|^2)$;
their determinant is one and their inverse is
\[
 \begin{pmatrix}
 \overline x/(|x|^2+|y|^2)&-\overline y/(|x|^2+|y|^2)\\
 y&x
 \end{pmatrix}.
\]
Positive column rescaling gives the unit Bott map.  It has second
Chern residue $+1$ on the positive complex normal sphere of each
critical component.  At $xyz$ the syzygies
$(x,-y,0)$ and $(x,0,-z)$ give the same unit residue on each
axis.  The finitely many triple strata have real codimension six,
so their relative cohomology below degree six vanishes.  They
contribute no degree-four term.  This proves the equivariant
Thom equality, including its torsion part; the final equality
uses \eqref{maps:eq:D.OR12}.

Let $a_C\in H^3(X\setminus C;\mathbb Z)$ be the residue class
obtained by intersecting with the complete chain $B_C$.
Its residue on every positive normal three-sphere is one.
For the principal circle bundle
$X\times ES^1\longrightarrow \mathcal B$, fibre integration sends its
restriction over $\mathcal B^\circ$ to a class extending uniquely to
$b\in H^2(\mathcal B;\mathbb Z)$.  Indeed deleting the codimension-four
critical strata induces an isomorphism in degree two; the
triple strata do not change that range.

The Thom-pair construction of the integral Gysin sequence
\cite[\S4.D, pp.~438,442--444]{Hatcher2002} gives the exact
boundary equation
\[
                           t b=[C]_{S^1}.
\]
More explicitly its cochain cone has differential
$D_{\rm Gys}(\alpha,\beta)=(d\alpha+t\smile\beta,-d\beta)$, and fibre
integration reads $\beta$.  The positive Hopf connection
$\eta_{\rm H}=(2\pi)^{-1}\operatorname{Im}(\overline z_0dz_0+
\overline z_1dz_1)$ has fibre integral one and fixes the sign
of this cone.  The first component extends across the critical
strata with their positive Thom boundary, and the second extends
to $b$.  Their boundary is exactly the displayed cup product.
The relative pullback from the Borel pair is an isomorphism in
degree four because the lower relative Thom groups vanish;
its inverse extends the same oriented normal-disc cochain.

Multiplication by $t$ on $H^2(\mathcal B;\mathbb Z)$ is injective.
Its kernel in the Gysin sequence is the image of
$H^3(X;\mathbb Z)=\mathbb Z/8$, while
\eqref{maps:eq:D.OR11} makes $H^2(\mathcal B;\mathbb Z)$ free.
Consequently $b=-t$.  Naturality for the Hopf sphere gives
\begin{equation}
 \ell(N_{\rm sph},C)=\langle a_C,[N_{\rm sph}]\rangle
 =\langle b,[\Sigma]\rangle=-1.
 \hypertarget{tag-Q252}{}\label{maps:eq:weighted-os45-total-link}
\end{equation}
This is the intersection with the entire prescribed chain $B_C$.
Changing that chain by a closed integral three-cycle changes the
intersection by zero, since $[N_{\rm sph}]$ is torsion in $H_3(X)$.
Reversing the sphere orientation gives $+1$ and preserves the
odd intersection parity used below.

\item \emph{The refined class.}
The reference parity is one by the Bott degree-four calculation
\eqref{maps:eq:D.OR10}, and the total link is $-1$ by
\eqref{maps:eq:weighted-os45-total-link}. Use the integral relative
cochain argument of \eqref{maps:eq:D.WI13} on this $X$ and its
embedded sphere $N_{\rm sph}$. The manifold is simply connected and spin with
$\lambda=0$; the unit Thom residues and the index-two normal-frame
map have been proved above. The admissible three-frame and its
rank-three spin complement in that argument give the obstruction
$\mathcal D_X^{-1}(P_{\rm ref})$ on this manifold. Its full frame
has even sphere parity, so its degree $d_\rho$ relative to $\rho$
satisfies $d_\rho\equiv\epsilon_\rho\pmod2$.
For an integral four-chain $D_4$ with $\partial D_4=8N_{\rm sph}$, the same
oriented cell obstruction is
\[
 -D_4\cdot C=2q_{\rm ref}+8d_\rho,
 \qquad q_{\rm ref}/8=\operatorname{lk}(N_{\rm sph},P_{\rm ref})\quad\text{in }\mathbb Q/\mathbb Z.
\]
Here the second equality is the cap-product identification
$\mathcal D_X(c)=c\frown[X]$, with its declared inverse, and
$D_4\cdot C=8\ell(N_{\rm sph},C)$ because $\partial B_C=C$.
Division by sixteen gives
\[
 \operatorname{lk}(N_{\rm sph},P_{\rm ref})
 =-\tfrac12\ell(N_{\rm sph},C)-\tfrac12\epsilon_\rho=0
                  \quad\hbox{in }\mathbb Q/\mathbb Z.
\]
The sphere is odd in $H_3(X)=\mathbb Z/8$ and the torsion linking
pairing is nondegenerate.  Also $2P_{\rm ref}=\mathcal D_X(\lambda)=0$ kills the
free part.  We obtain
\begin{equation}
                              P_{\rm ref}(X)=0.
 \hypertarget{tag-03LY}{}\label{maps:eq:D.OR14}
\end{equation}
\end{enumerate}
\end{enumerate}
\end{proof}

\section{The OS49 family}\hypertarget{tag-034L}{}\label{sec:III:os49-night}



\subsection{An explicit sphere filling}
\begin{theorem}[An explicit OS49 sphere]\hypertarget{tag-02S2}{}\label{night:P3F.26}\label{stmt:P3F.26}
{\normalfont\scriptsize tag 02S2}\par

The explicit $IV^*+II+I_2$ source with $M=\mathcal O(G-O)$ and $Q=3G$ has component-supported linearization divisor $C_0$. Its original type-$II$ filling, specified full-lattice toric filling at $I_2$ and order-three affine cap give a simply connected integral homology six-sphere, hence the standard smooth $S^6$.
\end{theorem}
\begin{proof}
\leavevmode
\paragraph{The source and its line.}
\hypertarget{tag-0351}{}\label{the-source-section-and-component-divisor}
The source and its marked section are
\begin{equation}
 y^2=x^3-3t^3(t-1)x+2t^4(t-1)(t-\tfrac12),
 \qquad G=(t^2,t^2).
 \hypertarget{tag-Q303}{}\label{os49:eq:source-cubic}
\end{equation}
Use the two Weierstrass charts related by
\(q=t^{-1}, X=q^2x,Y=q^3y\), with inverse
\(x=q^{-2}X,y=q^{-3}Y\).  For the displayed cubic,
\[
 \Delta=-432t^8(t-1)^2,\qquad J=j/1728=4t(1-t).
\]
The orders at zero and one are \((3,4,8)\) and \((1,1,2)\);
at infinity they are \((0,0,2)\).  The minimal degree-\((4,6)\)
model is therefore the rational elliptic surface with the stated
fibres; see \cite[\S4.5 and Table~1, p.~14; \S8.4, p.~39]{manSS10}.
The group law on this cubic gives
\[
 2G=\left(\frac{t^2}{4},\frac{t^2(9t-8)}8\right),\qquad
 3G=\left(\frac{9t^2-72t+64}{9},
          -\frac{351t^2-864t+512}{27}\right).
\]
Substitution verifies both points.  Their coordinates and those of
\(G\) have degrees at most \((2,3)\), so all three sections are
disjoint from \(O\), including in the second Weierstrass chart.

At infinity put \(X=1+r\).  The equation becomes
\[
 Y^2=3r^2+r^3+3qr+q^2.
\]
The blow-up chart \(r=qA,Y=qB\) has equation
\(B^2=3A^2+qA^3+3A+1\).
Here \(G\) and \(Q=3G\) meet the exceptional component \(C_1\)
at \((0,1)\) and \((-8,-13)\), away from its two intersections
with \(C_0\).  The section \(2G\) meets \(C_0\) at
\((X,Y)=(1/4,9/8)\).  At zero, the map
\(t=s^3,x=s^4X,y=s^6Y\) gives
\[
 Y^2=X^3-3s(s^3-1)X+2(s^3-1)(s^3-1/2).
\]
The central curve is \(Y^2=X^3+1\), and \(G=(s^2,1)\)
specializes to \((0,1)\).  The line \(Y=1\) meets it three times
at that point, so it has exact order three.  Thus \(G\) is on a
nonidentity component of \(IV^*\), and \(Q\) is narrow there.
In particular
\(\langle G,G\rangle=2-4/3-1/2=1/6\).
Shioda--Tate gives Mordell--Weil rank \(10-2-6-1=1\).
For any section other than \(O\), the height is
\(2+2P\cdot O-\operatorname{contr}_{IV^*}(P)
 -\operatorname{contr}_{I_2}(P)>0\), since the two corrections
are at most \(4/3\) and \(1/2\).  Hence there is no torsion.
All heights belong to \(\frac16\mathbb Z\); the height of \(G\)
then proves that it is primitive.  These are the rank and height
formulas of \cite[\S6.11, p.~29; \S11.8 and Table~4, pp.~51--52]{manSS10}, applied to
the components just located.

Set \(T=t_Q\), \(M=\mathcal O(G-O)\), and
\(L=\operatorname{Hom}(T^*M,M)\).
Translation fixes the generic degree-zero line class, so a rational
trivialization of \(L\) has a vertical divisor.  At \(IV^*\) its
component degrees vanish because \(Q\) is narrow.  At \(II\) the
fibre is irreducible.  At \(I_2\) the degrees of \(M\) and
\(T^*M\) are \((-1,1)\) and \((1,-1)\), respectively.
Thus \(L\) has the component degrees \((-2,2)\) of \(C_0\).
The kernel of each fibre intersection matrix is generated by its
full fibre, so
\(L=\mathcal O(C_0)\otimes\pi^*\mathcal O(e)\).
There are no multiple fibres, and restriction to \(O\) gives
\[
 \deg L|_O=G\cdot O-O^2-(-2G)\cdot O+(-3G)\cdot O=1=C_0\cdot O.
\]
Hence \(e=0\).  The section with divisor \(C_0\), multiplied by
\(\lambda\ne0\), is precisely the arrow
\(M_{p+Q}\to M_p\) covering \(T^{-1}\).
Its two successive divisors are \(C_0,C_1\); their sum is the
full reduced fibre.  This establishes the literal component
linearization used below.

\paragraph{The integral marking.}
\hypertarget{tag-02QG}{}\label{marked-monodromy-and-its-completeness}
Choose the following period marking of this source.
Let \(\rho_{\rm ell}=(1+i\sqrt3)/2\).  In the modular uniformization, the
arc from \(\rho_{\rm ell}\) to \(i\) maps to \([0,1]\) under \(J\);
use its point over \(1/2\) as base point.  The quotient coordinates
at its endpoints are the cube of
\((\tau-\rho_{\rm ell})/(\tau-\overline{\rho_{\rm ell}})\) and the square of
\((\tau-i)/(\tau+i)\).  Positive loops \(a_{\rm orb},b_{\rm orb}\) around \(0,1\)
therefore lift with endpoint transformations
\[
 U=\begin{pmatrix}0&1\\-1&1\end{pmatrix},\qquad
 V=\begin{pmatrix}0&1\\-1&0\end{pmatrix}.
\]
Indeed \(U'(\rho_{\rm ell})=e^{2\pi i/3}\) and \(V'(i)=-1\).
This uses only the fundamental domain and normalized modular
coordinate in \cite[Theorem~2.12, p.~32; Proposition~4.3 and
Remark~4.4, p.~49]{mapsMilne2017}.
For \(G_0=\left(\begin{smallmatrix}a&b\\c&d\end{smallmatrix}\right)\),
the torus comparison is
\[
 f_{G_0}:\mathbb C/(\mathbb Z+\mathbb Z\tau)
 \longrightarrow\mathbb C/(\mathbb Z+\mathbb ZG_0\tau),
 \qquad z\longmapsto z/(c\tau+d),
\]
with inverse multiplication by \(c\tau+d\).
Its matrix in the oriented cycle basis \((1,\tau)\) is
\(M_{G_0}=\left(\begin{smallmatrix}a&-b\\-c&d\end{smallmatrix}\right)\).
The endpoint identification is \(f_{G_0}^{-1}\), so cohomology
transport is \(M_{G_0}^t\).  Direct substitution gives
\(M_U^t=U,M_V^t=V\), fixing the dual convention.

This projective local system pulls back to the displayed cubic by
an actual pointed isomorphism.  Write its coefficients as \(a_{\rm W}(t),b_{\rm W}(t)\).
On a simply connected open set away from \(J=0,1,\infty\), lift
\(J(t)\) to \(\tau(t)\), and put
\[
 A_4=-4a_{\rm W}(t)/g_2(\tau),\quad B_6=-4b_{\rm W}(t)/g_3(\tau),\quad
 c_{\rm unif}^2=B_6/A_4.
\]
The identity of \(J\)'s gives \(A_4^3=B_6^2\), so a local square
root \(c_{\rm unif}\) satisfies \(c_{\rm unif}^4=A_4,c_{\rm unif}^6=B_6\).
The isomorphism and inverse are
\[
\begin{gathered}
 x=c_{\rm unif}^2\wp(z),\qquad y=c_{\rm unif}^3\wp'(z)/2,\\
 z=c_{\rm unif}\int_O^{(x,y)}\frac{dx}{2y}\pmod{\mathbb Z+\mathbb Z\tau};
\end{gathered}
\]
see \cite[Propositions~3.12--3.13, p.~47]{mapsMilne2017}.
Changing the sign of \(c_{\rm unif}\) changes only the central lattice sign.

The two sheets of the actual cover are
\(t_\pm=(1\pm\sqrt{1-J})/2\), with positive real square root
at \(J=1/2\).  Start on \(t_-\).  The lift of \(a_{\rm orb}\) is the
positive peripheral at \(t=0\).  Lifting \(b_{\rm orb}^{-1}\), then \(a_{\rm orb}\),
then \(b_{\rm orb}\) gives the positive peripheral at \(t=1\); in
rightmost-first path notation it is \(b_{\rm orb}a_{\rm orb}b_{\rm orb}^{-1}\).
The connecting stem goes above \(t=1/2\), as is seen by continuing
the displayed square root through a clockwise loop about \(J=1\).
The third Schreier generator \(b_{\rm orb}^2\) bounds the filled ordinary
point, since \(J-1=-4(t-1/2)^2\).  The infinity peripheral is
the inverse product of the two finite peripherals.
Their actual traces \(-1,1\) fix the central signs as
\(-U,VUV^{-1}\).  The determinant-one change
\[
 \mathsf H=\begin{pmatrix}0&-1\\1&1\end{pmatrix},\qquad
 \mathsf H^{-1}=\begin{pmatrix}1&1\\-1&0\end{pmatrix}
\]
gives the elliptic cohomology matrices
\[
 E_A=\begin{pmatrix}-2&-3\\1&1\end{pmatrix},\quad
 E_B=\begin{pmatrix}1&1\\-1&0\end{pmatrix},\quad
 E_C=\begin{pmatrix}1&2\\0&1\end{pmatrix},\quad E_AE_BE_C=I.
\]

In the dual homology basis let \(A_i=E_i^{-t}\).  The negative
Abel--Jacobi cocycle \(c_i\) lies in the saturated image of
\(A_i-I\) and obeys \(c_A+A_Ac_B+A_AA_Bc_C=0\).
Because \((A_B-I)^{-1}=\left(\begin{smallmatrix}0&-1\\1&-1\end{smallmatrix}\right)\),
the unique integral logarithm shift makes \(c_B=0\).
The other equations then give
\(c_A=(0,a_{\rm MW})^t,c_C=(0,-a_{\rm MW})^t\).
The Poincar\'e comparison in \SeriesRef{maps:src:D.2.12}
identifies the full integer height with this cocycle pairing:
for \(J_0=\left(\begin{smallmatrix}0&1\\-1&0\end{smallmatrix}\right)\),
\(\beta_i=-c_i^tJ_0A_i\), and
\((A_i-I)z_{G,i}=c_i\), the sum is \(a_{\rm MW}^2/6\).
Explicitly \(z_{G,A}=(a_{\rm MW}/3,0)^t,z_{G,C}=(a_{\rm MW}/2,0)^t\), the cross term
vanishes, and \(-\beta_Az_{G,A}-\beta_Cz_{G,C}=-a_{\rm MW}^2/3+a_{\rm MW}^2/2\).
The scope of that comparison is its universal rigidified line maps:
\(\eta=e^{-2\pi izy_2}\xi\), with inverse
\(\xi=e^{2\pi izy_2}\eta\), and circle lattice rule
\(\theta\mapsto\theta+r^tJ_0y\).
Multiplying the present section by six makes it narrow at both
reducible fibres; contraction of the zero-axis rigidifications fixes
the integer lift there.  Division by \(36\) gives the displayed
pairing for this source.  Thus \(a_{\rm MW}^2=1\).
Negating both elliptic basis cycles, with inverse \(-I_2\), gives
\(a_{\rm MW}=1\) without changing the actual section \(G\).

For the weights \(M=\mathcal O(G-O)\), \(Q=3G\), the full
homology lift in order \((e_1,e_2,\delta,\gamma)\) is
\[
 L_i=\begin{pmatrix}A_i&0&3c_i\\
              \beta_i&1&\beta_i z_i+w_i\\0&0&1\end{pmatrix},
 \quad
 z_A=(1,0)^t,\ z_B=0,\ z_C=(3/2,0)^t,
 \quad(w_A,w_B,w_C)=(0,0,1/2).
\]
The last weight is fixed by the actual two-iterate divisor
\(C_0+C_1\); it combines with \(\beta_Cz_C=-3/2\) to give
the integral entry \(-1\).  Cyclically reorder homology to
\((\gamma,e_1,e_2,\delta)\) and take inverse transpose.
In the dual order \((\psi,u,w,\delta)\), this gives
\begin{equation}
 T_A=\begin{pmatrix}1&-3&-3&-1\\0&-2&-3&-1\\0&1&1&0\\0&0&0&1\end{pmatrix},\quad
 T_B=\begin{pmatrix}1&0&0&0\\0&1&1&0\\0&-1&0&0\\0&0&0&1\end{pmatrix},
 \quad
 T_C=\begin{pmatrix}1&0&3&1\\0&1&2&1\\0&0&1&0\\0&0&0&1\end{pmatrix}.
 \hypertarget{tag-Q304}{}\label{os49:eq:cohomology-monodromy}
\end{equation}
Their product and \(T_A^3\) are the identity.
On homology \(A=T_A^{-t}\) fixes
\(\ell_0=(1,-1,0,0)^t\) and \(\delta=(0,0,0,1)^t\).

\paragraph{The reduced fillings and the exterior section.}
\hypertarget{tag-0007}{}\label{actual-local-fillings-and-their-gluing}
At infinity the actual period maps, with their holomorphic units, are
\begin{equation}
 h(Z,V)=(q^2aZ,qbV),\qquad
 g(Z,V)=(q^{-3}cZ,\lambda q^{-1}d_0V).
 \hypertarget{tag-Q305}{}\label{os49:eq:cusp-deck}
\end{equation}
The valuation of \(G\) is one, and the two-iterate divisor just
computed fixes the second valuation of \(g\).
The exact deck change and its inverse are
\[
 H=h^{-1}g^{-1},\quad K=h^3g^2,
 \qquad h=H^2K,\quad g=H^{-3}K^{-1}.
\]
Their exponent matrices are inverse integral matrices, and \(H,K\)
have valuation columns \((1,0)^t,(0,1)^t\).
These identities retain all units and \(\lambda\).
The specified negative-diagonal height-one triangulation is therefore
the full-lattice Mumford filling.  Each cone is unimodular, its
chart is \(\mathbb C^3\), and \(q\) is the product of the three
coordinates.  Adjacent charts glue by their common integral monomials.
A compact annular fundamental region in the logarithmic moduli,
completed by these charts, proves properness and the Hausdorff
property.  The normalization of the central fibre is the hexagonal
toric surface, with its opposite boundaries identified by the actual
unit periods.  Modulo the full period lattice there are one vertex,
three edges and two triangles.  Its open torus and double-curve
strata have Euler number zero, while the two triple points each
contribute one, giving central Euler number two.
The determinant-one cohomology basis
\[
 (a_1,a_2,b_1,b_2)=(3\psi+2u,\psi+u,w,\delta),
 \qquad \psi=a_1-2a_2,\quad u=-a_1+3a_2
\]
puts \(T_C\) in the form \(I+E_{13}+E_{24}\).
Thus the full deck-group map and this lattice map identify the cusp
with the primitive rank-two model of
\cite[Proposition~3.1, pp.~6--7]{manViaclovskyTwo26}; its specialization is
the full invariant lattice in every degree.  Equivalently one may
use the Mumford clause of \cite[Lemma~7.4, pp.~18--19]{E26} with this
index-one marking.  The local fundamental-group kernel consists of
the primitive circles \(w,\delta\) and the section meridian.

At \(II\), use the original smooth total \(M^\times\) and its
contracting lift.  Let \(N_E\) be a small closed elliptic
neighbourhood.  It retracts to its cuspidal rational curve
\(C\simeq S^2\), by normalization, whose inverse is \(y/x\)
away from the cusp and zero there.  The degree of \(M|_C\) is zero,
so \(M|_{N_E}\) is topologically trivial.  A local relative
logarithm of \(Q\) gives translations \(T_s\), with inverse
\(T_{-s}\), and \(T_1=T\).  In a line trivialization write the
lift as \(F_{\rm line}(p,v)=(T^{-1}p,a(p)v)\), where \(0<|a(p)|<1\).
Since \(N_E\) is simply connected, \(a=e^{A_0}\) with
\(\operatorname{Re}A_0<0\).  The fundamental annulus is
\[
 (p,\theta,s)\longmapsto
   (p,e^{sA_0(Tp)}e^{2\pi i\theta}),\qquad 0\le s\le1,
\]
with inverse
\(s=\log|v|/\operatorname{Re}A_0(Tp)\) and
\(\theta=(\arg v-s\operatorname{Im}A_0(Tp))/(2\pi)\).
Its boundary identifies \((p,\theta,1)\) with \((Tp,\theta,0)\).
The inverse maps
\[
 [p,\theta,s]\longmapsto(T_sp,\theta,e^{2\pi is}),\qquad
 (x,\theta,e^{2\pi is})\longmapsto[T_{-s}x,\theta,s]
\]
give its product description, fibrewise over the elliptic base.
The central fibre is \(S^2\times T^2\).
Near the cusp the base function is analytically \(Y^2-X^3\):
split its nondegenerate quadratic direction and absorb the remaining
cubic unit by its analytic cube root.  Its two Milnor thimbles
collapse primitive elliptic cycles meeting once; these are a
one-skeleton of the torus.  The quotient of its single oriented
two-cell is a sphere of degree one.  Taking the product with the
two retained circles proves full integral specialization onto the
\(T_B\)-invariants in all degrees.  The local group kernel is
generated by \(u,w\) and the section meridian.

There is an actual section over the complement of the finite cap.
The line \(M|_O\) has degree one; choose its section with its only
zero at infinity.  Near \(O\) there write \(\phi=qU_{\rm line}\), with
\(U_{\rm line}\) a unit.  Applying \(\phi^{-1}\) to \(qv\) gives
\(U_{\rm line}^{-1}v\) along \(Q\), at its already computed smooth point of
\(C_1\).  This nonzero representative is a source-line chart in
the toroidal filling.  One forward iterate is the inverse
comparison, so the section extends through the cusp.  It is
nonzero at \(II\).  These extending representatives fix both reduced meridians
and the exterior attachment.

\paragraph{The cap and its complete specialization maps.}
Use the clockwise root deck \(s\mapsto e^{-2\pi i/3}s\).
The actual good-reduction coordinates give \(X\mapsto\zeta_3X\),
\(Y\mapsto Y\), inducing \(A=T_A^{-t}\) on homology.
Its elliptic matrix is conjugate to
\(\left(\begin{smallmatrix}0&-1\\1&-1\end{smallmatrix}\right)\)
by \(\left(\begin{smallmatrix}-1&0\\-1&-1\end{smallmatrix}\right)\),
whose inverse is \(\left(\begin{smallmatrix}-1&0\\1&-1\end{smallmatrix}\right)\).
A natural line character \(c\delta/3\) is removed by translating
the origin by \(-c\gamma/3\), since
\((A-I)(\gamma/3)=w+\delta/3\); the remainder is an integral
period.  The inverse origin translation restores the original line
frame.  This gives a fixed central origin.  Choose an invariant
linear complement to the base tangent by averaging a splitting.
Averaging local coordinates by
\(\frac13\sum_{j=0}^2(D\rho)^{-j}\rho^j\), retaining the base
coordinate, has identity derivative and an analytic inverse.  Its
base axis supplies an equivariant holomorphic origin on the root
disc.  Hence the specified affine generator is
\[
 g_{\rm cap}(z,s)=(Az+\ell_0/3,e^{-2\pi i/3}s).
\]
It has order three.  Its two nonidentity powers change the invariant
\(\psi\)-circle by \(1/3,2/3\), so the central action is free;
the base rotation proves freeness off the central fibre.

The punctured overlap and inverse are
\[
 h_s(z)=z+\sigma(s),\qquad h_s^{-1}(z)=z-\sigma(s),\qquad
 \sigma(s)=\frac{\log s}{2\pi i}\Pi(s)\ell_0.
\]
The clockwise identity
\(\sigma(e^{-2\pi i/3}s)=\rho\sigma(s)-\Pi(e^{-2\pi i/3}s)\ell_0/3\)
gives \(h_{e^{-2\pi i/3}s}g_{\rm cap}h_s^{-1}=\rho\).
Branch changes add integral periods.  The positive exterior
meridian is represented by \(-\sigma(s)\) on the twisted cover;
its endpoint is \(g_{\rm cap}^{-1}(-\sigma(s))\).  Consequently the normal
circle has kernel vector \((\ell_0,3)\) in fibre and positive
exterior-meridian coordinates.
The finite quotient is smooth and proper.  Together with the
reduced charts and the compact regular-base part it gives the
specified compact smooth complex threefold.
For the covering map
\(p:\Delta_r^*\to\Delta_\lambda^*,\ p(r)=r^3\), pull the
scalar family back along \(p\).  The caps just constructed give
a holomorphic family over \(\Delta_r^*\).
The scalar \(\lambda\) occurs
only in the line part of the suspension period, so
\(\Pi\gamma=\Gamma(s)+(\log\lambda/(2\pi i))\Pi\delta\),
with \(\Gamma\) independent of that scalar.  Its one-third
translation therefore contains \((\log r/(2\pi i))\Pi\delta\),
whose branch changes are integral periods.  The cap and origin
translations thus exponentiate to single-valued holomorphic maps
on the punctured \(r\)-disc.  On each compact parameter subannulus,
uniform contraction and bounds on the retained units permit one
choice of sufficiently small base discs and compact fundamental
regions.  Exhaustion by such subannuli, using the same maps on the
overlaps, proves properness of the resulting parameter family.

For the integral specialization, use homology coordinates
\((u,w,\delta,\ell_0)\).  Then
\[
 A=\operatorname{diag}(F,1),\quad
 F=\begin{pmatrix}1&-1&0\\3&-2&0\\1&-1&1\end{pmatrix},\quad
 P=\begin{pmatrix}0&1&-1\\-1&2&-1\\0&1&0\end{pmatrix},\quad
 P^{-1}=\begin{pmatrix}1&-1&1\\0&0&1\\-1&0&1\end{pmatrix}.
\]
Both changes have determinant one and \(P^{-1}FP\) cyclically
permutes three circle factors.  Write \((a,b,c,t)\) for the
coordinates before \(P\).  The quotient of the central torus is
\[
 (T^3\times[0,1])/((x,0)\sim(Fx,1));
\]
the covering from \(T^4\) is its threefold unwrapping, with base
coordinate \(3t\).  This displayed endpoint convention fixes the
mapping-torus direction.
The one-vertex, one-edge cell structure on each circle gives a
product cell structure on \(T^3\), with zero differential.  In
degrees one and two the cyclic permutation permutes three oriented
cells.  Its invariant generator is their sum; its coinvariants
are freely generated by one cell.  In degrees zero and three it
is the identity.  The mapping-torus cochain complex is the cone of
this cellular map minus the identity, so its cohomology has free
ranks \((1,2,2,2,1)\).  Pullback of the unwrapping is identity on
invariant fibre cochains and the norm \(1+F^*+F^{2*}\) on the
edge cochains.

Put \(\alpha_2=da\wedge db\), \(\lambda_1=3dc-db\),
\(v_3=da\wedge db\wedge dc\).  In cyclic coordinates these are
the invariant sums and volume just described.  Moreover
\[
 \psi=dt,\quad u=da-dt,\quad w=db,\quad\delta=dc,
\]
\[
 \omega=uw+3\psi\delta=\alpha_2+dt\lambda_1,
 \quad\phi_{\rm coh}=\psi uw=dt\alpha_2,
 \quad\nu=\psi uw\delta=dtv_3.
\]
The full specialization images are therefore
\[
 \begin{array}{c|c}
 1&\langle\lambda_1,3dt\rangle\\
 2&\langle\alpha_2,dt\lambda_1\rangle\\
 3&\langle v_3,dt\alpha_2\rangle\\
 4&\langle3dtv_3\rangle.
 \end{array}
\]
Their full cokernels have orders \((3,1,1,3)\), generated by
\(\psi,\omega,\phi_{\rm coh},\nu\), respectively.

\paragraph{The integral global complex and differentials.}
\hypertarget{tag-01RP}{}\label{integral-topology}
Write \(f:X\to\mathbb P^1\) for the base projection of the
glued threefold.
For \(\mathcal L_j=\bigwedge^j\mathbb Z^4\), put
\(\mathsf E_{j,i}=\bigwedge^jT_i\) and use the exact integral images in
\[
 \mathcal L_j\xrightarrow{D_j^0}\bigoplus_{i=A,B,C}(\mathsf E_{j,i}-I)\mathcal L_j
 \xrightarrow{D_j^1}\mathcal L_j,
\]
\[
 D_j^0v=((\mathsf E_{j,i}-I)v)_i,\qquad
 D_j^1(y_A,y_B,y_C)=y_A+\mathsf E_{j,A}y_B+\mathsf E_{j,A}\mathsf E_{j,B}y_C.
\]
The explicit integer contractions recorded below prove
\(H^0=\mathbb Z,H^1=0,H^2=\mathbb Z\) for \(j=1,2,3\).
Their inclusions \(i_j\) and quotient maps \(q_j\), in the
lexicographic exterior bases, are
\[
 \begin{array}{c|c|c}
 j&i_j(1)&q_j\\\hline
 1&\psi&(0,0,0,1)\\
 2&\omega&(0,0,1,3,0,0)\\
 3&\phi_{\rm coh}&(0,0,0,1).
 \end{array}
\]
Degrees zero and four are constant sheaves.  The reduced
specializations are full invariant lattices, and the finite
specialization is injective with the complete images above.  Hence \(R^jf_*\mathbb Z\) is the kernel of the
invariant-extension sheaf's map to a single cyclic skyscraper.
Its global invariant maps onto that skyscraper.  The long exact
sequence gives zero \(H^1\), the same free \(H^2\), and
\(H^0\) generators \(1,3\psi,\omega,\phi_{\rm coh},3\nu\).

The intrinsic differential is determined by the following auxiliary filling.  Replace only the affine cap by the
coarse quotient of the same root torus family by its linear action.
Inversion about the exterior section extends to this auxiliary
filling.  At the cusp it is the monomial map
\((Z,V)\mapsto(f_1(q)^2/Z,f_2(q)^2/V)\), where \((f_1,f_2)\)
represents that section.  Its integral valuation sends a fan ray
\((i,j,1)\) to \((2v_1-i,2v_2-j,1)\); the centrally symmetric
periodic triangulation is preserved, and the unit factors invert
the full deck maps.  At \(II\), the zero-rigidified identity
\([-1]^*M\simeq M^{-1}\) gives the fibrewise reciprocal line map
over elliptic inversion.  This identity extends over the whole
irreducible-fibre neighbourhood: its generic difference is vertical,
and its base-disc factor is fixed by the zero-axis rigidification.
It conjugates the contracting group translation to its inverse, so
it descends to the actual quotient.  At the linear cap it commutes
with the root action.  These maps agree on their common group-torus
locus and hence on the overlaps.

The rational cohomology of a finite coarse quotient is the invariant
cohomology upstairs: an equivariant subdivision and averaging give
the cochain comparison.  Inversion consequently acts by \((-1)^j\)
on every rational direct image of the auxiliary filling.  Naturality
of \(d_2\) then gives \(d_2=-d_2\), so this intrinsic differential
is zero over \(\mathbb Q\).
The actual collar differs by the translation path
\(x\mapsto x+\theta\ell_0/(2\pi\cdot3)\).
Its prism acts by \(\iota_{\ell_0}/3\) on exterior cochains.
The radial logarithmic term is removed by scaling that translation
to zero; its inverse is subtraction, and its deck equivariance
follows from invariance of \(\ell_0\) and equivariance of the
complex structure.  Logarithms of nonvanishing unit factors have
the same explicit isotopy.  Thus this is the full translation
contribution.  With a cutoff \(\chi_{\rm cut}\) increasing from the inner chart to
the outer chart, \(d\chi_{\rm cut}\wedge d\theta/(2\pi)\) integrates to
\(+1\) in the complex base orientation.  The resulting formula is
\[
 d_2(a)=\frac{H_B}{3}[\iota_{\ell_0}a].
\]
Its rational determination fixes the integral map because all its
actual sources and targets above are free.  The projected
contractions of \(\psi,\omega,\phi_{\rm coh},\nu\) are
\(1,3[\delta],3\eta_{\rm coh},[uw\delta]\), with
\([\psi\delta]=\eta_{\rm coh}\), \([uw]=3\eta_{\rm coh}\).
After the source factors \((3,1,1,3)\), all four differentials
are multiplication by \(+1\).  The two-dimensional base permits only these differentials, and
each surviving total-degree filtration has one term. Hence \(H^0=H^6=\mathbb Z\) and all
other integral cohomology vanishes.

Van Kampen uses the same actual boundary maps.  The type-II filling
kills \(u,w\) and its section meridian; the cusp kills
\(w,\delta\) and its section meridian.  The base relation kills
the remaining exterior meridian.  The resulting fibre quotient is
\(\mathbb Z\langle\gamma\rangle\), with quotient
\(\gamma\mapsto1,u,w,\delta\mapsto0\) and inverse generator
inclusion.  The finite normal kernel \((\ell_0,3)\) maps to
\(\gamma\).  Its boundary map is onto, since the disk-bundle
boundary contains the affine generator and every fibre generator.
Thus \(\pi_1(X)=1\).  Hurewicz and Whitehead give a homotopy
equivalence to \(S^6\), and \(\Theta_6=0\), in
\cite[table on p.~504 and first remark on p.~505]{manKM63},
gives the standard smooth sphere.

The physical unit circle is semifree on this
construction.  It is free on the contracting line-complement
quotient and on the cap, since a cap deck identification of a line
circle would have nonintegral \(\psi\)-displacement.
In the cusp, its cocharacter is \((0,1,0)\).  A vertical wall has
rays \(v_0=(0,0,1),v_1=(0,1,1)\), with adjacent rays
\(v_- =(-1,1,1),v_+=(1,0,1)\).  The cocharacter is
\(v_1-v_0\), so the normal weights are \((-1,+1)\).
The relation \(v_-+v_+=v_0+v_1\) gives normal bundle
\(\mathcal O(-1)\oplus\mathcal O(-1)\) on the wall's
\(\mathbb P^1\).  There is one such wall modulo the full period
lattice.  The remaining cones have primitive nonzero weights and
no further fixed component.  Nonzero deck translations move their
cones, so they create no additional finite stabilizer.
\paragraph{Integer contractions for the parabolic complex.}
The following integral maps give the parabolic calculation. Columns use
the lexicographic exterior bases; the middle basis is the concatenation
of the columns of $B_{j,A},B_{j,B},B_{j,C}$ below.  The middle groups use the exact image lattices.
Write $D_j^0$ in its three corresponding row blocks $C_{j,i}$.  The
column sets in the following table have determinant $\pm1$:
\[\begin{array}{c|ccc}
j&I_A&I_B&I_C\\\hline
1&(2,4)&(2,3)&(3,4)\\
2&(1,3,5,6)&(1,2,5,6)&(3,6)\\
3&(2,3)&(2,3)&(3,4)\\
\end{array}\]
Indeed $B_{j,i}C_{j,i}=\mathsf E_{j,i}-I$.  With $\iota_{I_i}$ the indicated
column inclusion, the integer right inverse
$R_{j,i}=\iota_{I_i}(C_{j,i}|_{I_i})^{-1}$ satisfies
$(\mathsf E_{j,i}-I)R_{j,i}=B_{j,i}$, proving equality with each exact image lattice.
For $j=1,3$ put $p_j=(1,0,0,0)$ and $s_j=(0,0,0,1)^t$;
put $p_2=(0,0,0,1,0,0)$ and $s_2=(0,0,1,0,0,0)^t$.
The inclusions $i_j$ and quotients $q_j$ are the preceding displayed
ones.  Direct multiplication of the following integer matrices gives
\[
 h_j^1D_j^0=I-i_jp_j,\quad
 D_j^0h_j^1+h_j^2D_j^1=I,\quad
 D_j^1h_j^2=I-s_jq_j,\quad p_ji_j=q_js_j=1.
\]
These identities prove both integral exactness in the middle and the
stated primitive end groups.
\smallskip\noindent Degree $1$:
\[
 B_{1,A}=\begin{pmatrix}1&0\\1&0\\0&1\\0&0\end{pmatrix},\quad B_{1,B}=\begin{pmatrix}0&0\\1&0\\0&1\\0&0\end{pmatrix},\quad B_{1,C}=\begin{pmatrix}1&0\\0&1\\0&0\\0&0\end{pmatrix}.
\]
\[
D^0_{1}=\begin{pmatrix}0&-3&-3&-1\\0&1&0&0\\0&0&1&0\\0&-1&-1&0\\0&0&3&1\\0&0&2&1\end{pmatrix}.
\]
\[
D^1_{1}=\begin{pmatrix}1&0&-3&-3&1&0\\1&0&-2&-3&0&1\\0&1&1&1&0&0\\0&0&0&0&0&0\end{pmatrix}.
\]
\[
h^1_{1}=\begin{pmatrix}0&0&0&0&0&0\\0&0&0&-1&-1&1\\0&0&0&0&1&-1\\0&0&0&0&-2&3\end{pmatrix}.
\]
\[
h^2_{1}=\begin{pmatrix}-2&3&0&0\\1&-1&1&0\\-1&1&0&0\\0&0&0&0\\0&0&0&0\\0&0&0&0\end{pmatrix}.
\]
\smallskip\noindent Degree $2$:
\[
 B_{2,A}=\begin{pmatrix}1&0&0&0\\0&1&0&0\\0&0&3&0\\0&0&-1&0\\0&0&3&0\\0&0&0&1\end{pmatrix},\quad B_{2,B}=\begin{pmatrix}1&0&0&0\\0&1&0&0\\0&0&0&0\\0&0&0&0\\0&0&1&0\\0&0&0&1\end{pmatrix},\quad B_{2,C}=\begin{pmatrix}1&0\\0&-1\\0&3\\0&-1\\0&2\\0&0\end{pmatrix}.
\]
\[
D^0_{2}=\begin{pmatrix}-3&-3&-1&3&1&0\\1&0&0&0&1&1\\0&0&0&0&-1&-1\\0&0&0&0&1&0\\0&1&0&0&0&0\\-1&-1&0&0&0&0\\0&0&0&0&0&1\\0&0&0&0&-1&-1\\0&2&1&-3&-1&1\\0&0&0&0&0&1\end{pmatrix}.
\]
\[
D^1_{2}=\begin{pmatrix}1&0&0&0&-2&-3&1&0&1&-2\\0&1&0&0&1&1&1&1&0&-1\\0&0&3&0&0&0&-3&-3&0&3\\0&0&-1&0&0&0&1&1&0&-1\\0&0&3&0&0&0&-2&-3&0&2\\0&0&0&1&0&0&1&1&0&0\end{pmatrix}.
\]
\[
h^1_{2}=\begin{pmatrix}0&0&0&0&-1&-1&0&0&0&0\\0&0&0&0&1&0&0&0&0&0\\0&0&0&0&-2&0&0&-1&1&-2\\0&0&0&0&0&0&0&0&0&0\\0&0&0&0&0&0&0&-1&0&-1\\0&0&0&0&0&0&0&0&0&1\end{pmatrix}.
\]
\[
h^2_{2}=\begin{pmatrix}1&0&0&-3&-1&0\\0&1&0&-3&-1&0\\0&0&0&2&1&0\\0&0&0&-3&-1&1\\0&0&0&0&0&0\\0&0&0&0&0&0\\0&0&0&3&1&0\\0&0&0&0&0&0\\0&0&0&0&0&0\\0&0&0&0&0&0\end{pmatrix}.
\]
\smallskip\noindent Degree $3$:
\[
 B_{3,A}=\begin{pmatrix}-1&0\\3&0\\0&1\\0&0\end{pmatrix},\quad B_{3,B}=\begin{pmatrix}0&0\\1&0\\0&1\\0&0\end{pmatrix},\quad B_{3,C}=\begin{pmatrix}1&0\\0&1\\0&0\\0&0\end{pmatrix}.
\]
\[
D^0_{3}=\begin{pmatrix}0&-1&-1&1\\0&1&0&0\\0&0&1&0\\0&-1&-1&0\\0&0&-1&1\\0&0&2&-3\end{pmatrix}.
\]
\[
D^1_{3}=\begin{pmatrix}-1&0&1&1&1&0\\3&0&-2&-3&0&1\\0&1&1&1&0&0\\0&0&0&0&0&0\end{pmatrix}.
\]
\[
h^1_{3}=\begin{pmatrix}0&0&0&0&0&0\\0&0&0&-1&3&1\\0&0&0&0&-3&-1\\0&0&0&0&-2&-1\end{pmatrix}.
\]
\[
h^2_{3}=\begin{pmatrix}2&1&0&0\\-3&-1&1&0\\3&1&0&0\\0&0&0&0\\0&0&0&0\\0&0&0&0\end{pmatrix}.
\]
\end{proof}

\subsection{The angular clutch}
\begin{theorem}[Angular clutch for the OS49 family]\hypertarget{tag-02S3}{}\label{night:P3F.27}\label{stmt:P3F.27}
{\normalfont\scriptsize tag 02S3}\par
The positive OS49 \(r\)-loop, with covering map \(\lambda=r^3\),
has cusp angular clutch
\[
 \theta K+3\theta(\theta-1)\delta,\qquad K=3u+2\gamma,
 \qquad 0\leq\theta\leq1.
\]
The primitive page marking and its inverse are
\(H=-u-\gamma\), \(K=3u+2\gamma\),
\(u=2H+K\), \(\gamma=-3H-K\).
Its physical line clutch has winding \(+6\) on the cusp cover
\(q=Q_0^6\).
\end{theorem}
\begin{proof}
Use the homology marking of \SeriesRef{night:P3F.26}
\((e_\gamma,e_u,e_w,e_\delta)\), dual to
\((\psi,u,w,\delta)\) in cohomology; thus
\(\ell_0=e_\gamma-e_u=(1,-1,0,0)^t\).
In the cusp formulas below, \(u=u_h\) is the literal logarithmic
period of \(h\), and \(\gamma=\gamma_g=\Pi(e_\gamma)\) is the
forward suspension period. We give the period calculation for \(m=3k\) and
\(\ell=n\ell_0+d\delta\), and then take \(m=3,n=1,d=0\).
The family is constructed over \(r\), with covering map
\(\lambda=r^m\).  Put
\[
 q=e^{2\pi is}=t^{-1},\qquad r=e^{2\pi iv},
 \qquad \delta=(0,1)^t.
\]
Use the unit logarithms in the actual marking of
\SeriesRef{night:P3F.26}.  Its full periods are
\[
 u=(2s,s)^t+U(q),\qquad
 \gamma=(-3s,-s+mv)^t+C(q).
\]
Their exact relation to the homology
marking is
\[
 \Pi(e_w)=w_Z=(1,0)^t,\quad \Pi(e_\delta)=\delta=(0,1)^t,\quad
 \Pi(e_u)=-u_h+\kappa,\quad
 \kappa=pw_Z+q_\delta\delta,\qquad p,q_\delta\in\mathbb Z.
\]
The gamma column of \(T_C^{-t}\) fixes \(w_Z\); comparison of the
two u increments leaves the displayed integral invariant vector.
Its coefficients are defined at the based regular fibre and
continued as constant integers. In geometric period order
\((\gamma_g,u_h,w_Z,\delta)\), the map and inverse are
\begin{equation}
 F_{p,q_\delta}=
 \begin{pmatrix}1&0&0&0\\0&-1&0&0\\0&p&1&0\\0&q_\delta&0&1\end{pmatrix},
 \qquad F_{p,q_\delta}^{-1}=F_{p,q_\delta}.
 \hypertarget{tag-Q30O}{}\label{os49:eq:based-cusp-period-map}
\end{equation}
For the full unit columns displayed below,
\[
 \Pi(s)=
 \begin{pmatrix}
 -3s+C_1&-2s-U_1+p&1&0\\
 -s+mv+C_2&-s-U_2+q_\delta&0&1
 \end{pmatrix},
 \quad
 \Pi(s+1)=\Pi(s)T_C^{-t}.
\]
The inverse arrow uses \(s-1\) and \(T_C^t\).
In particular the actual finite vector has geometric coordinates
\((n,n,-np,d-nq_\delta)^t\), and its period is
\(n(\gamma_g+u_h-\kappa)+d\delta\). The same \(F_{p,q_\delta}\) conjugates
its linear root action and transports its inverse.
The logarithms \(U,C\) are holomorphic on the cusp disc.
The full integral change of marking and its inverse are
\begin{equation}
 H=-u-\gamma,\qquad K=3u+2\gamma,
 \qquad u=2H+K,\qquad\gamma=-3H-K.
 \hypertarget{tag-Q306}{}\label{os49:eq:primitive-periods}
\end{equation}
The two matrices are
\(\left(\begin{smallmatrix}-1&-1\\3&2\end{smallmatrix}\right)\)
and
\(\left(\begin{smallmatrix}2&1\\-3&-1\end{smallmatrix}\right)\);
they have determinant one and are inverse.  Multiplicatively these
are the literal identities \(H=h^{-1}g^{-1}\), \(K=h^3g^2\),
\(h=H^2K\), \(g=H^{-3}K^{-1}\), including every unit.
The valuation columns of \(H,K\) are \((1,0)^t,(0,1)^t\), so
\(K\) is primitive.  The point coordinate of the \(K\)-circle evaluates
to one on \(u\), to minus one on \(\gamma\), and to zero on
\(H,w\).  With \(\varepsilon=3U+2C\), the full columns are
\begin{equation}
 H=(s,-mv)^t-U-C,\qquad
 K=(0,s+2mv)^t+\varepsilon(q).
 \hypertarget{tag-Q307}{}\label{os49:eq:full-cusp-periods}
\end{equation}

The finite collar is translation by
\(\log\sigma\,\Pi(\sigma)\ell/(2\pi i)\), with
\(t=\sigma^m\).  Under the positive parameter change \(v\mapsto
v+1\), the period of \(\gamma\) increases by \(m\delta\);
those of \(u,\delta\) do not change.  The collar change is
therefore \(n\log t\,\delta/(2\pi i)\).  Thus the ordinary
positive closing and its inverse are
\[
 D_n(p_E,V)=(p_E,t^nV),\qquad D_n^{-1}(p_E,V)=(p_E,t^{-n}V).
\]
At the cusp this subtracts \(ns\delta\).  On rays its map is
\((i,j,1)\mapsto(i,j-n,1)\), with inverse
\((i,j,1)\mapsto(i,j+n,1)\); both maps preserve each of the two
periodic triangulations.  Thus the full monomial map extends through
the toric filling.  At the type-II fibre \(t\) is a nonzero
holomorphic unit.  The finite logarithmic collar supplies the map
at the finite cap.  These maps agree on their common ordinary
locus.  The delta twist contributes zero to this parameter difference.

For an additive vector \(\Omega\), let \(\tau_\Omega\) multiply
the two coordinates by \(e^{2\pi i\Omega_1}\) and
\(e^{2\pi i\Omega_2}\).  Its inverse is \(\tau_{-\Omega}\).
On the logarithmic cover of the punctured cusp disc take
\[
 \Omega_{\mathrm{ord}}=nsv\delta.
\]
This is the actual map \((Z,V)\mapsto(Z,q^{nv}V)\), with
reciprocal inverse, on that logarithmic domain.  It satisfies
\(\Omega_{\mathrm{ord}}(s,v+1)-ns\delta
=\Omega_{\mathrm{ord}}(s,v)\).
On the whole cusp disc and a logarithmic parameter chart take
\[
 \Omega_{\mathrm{ext}}
   =-nv\varepsilon(q)-nmv(v-1)\delta.
\]
Its map and inverse are
\[
 \begin{split}
 (Z,V)&\longmapsto
 \bigl(e^{-2\pi inv\varepsilon_1}Z,
 e^{-2\pi i(nv\varepsilon_2+nmv(v-1))}V\bigr),\\
 (Z',V')&\longmapsto
 \bigl(e^{2\pi inv\varepsilon_1}Z',
 e^{2\pi i(nv\varepsilon_2+nmv(v-1))}V'\bigr).
 \end{split}
\]
All multipliers are holomorphic units at \(q=0\).  On every
fan chart a character \(Z^{a_1}V^{a_2}q^{a_3}\) is multiplied
by the corresponding unit \(A^{a_1}B^{a_2}\); its inverse is
multiplied by \(A^{-a_1}B^{-a_2}\).  These algebra isomorphisms
preserve products, agree on faces, and commute with the full
\(h,g\) deck maps.  They prove extension through every divisor
and node of the chosen fan.  Moreover
\[
 \Omega_{\mathrm{ext}}(q,v+1)-ns\delta
   -\Omega_{\mathrm{ext}}(q,v)=-nK(s,v).
\]
The extending chart therefore closes by the full period \(K^{-n}\),
including its units; the inverse closing uses \(K^n\).

For the logarithmic point \(z\), set
\(\xi_{\mathrm{ord}}=z+\Omega_{\mathrm{ord}}\) and
\(\xi_{\mathrm{ext}}=z+\Omega_{\mathrm{ext}}\).
Then \(\xi_{\mathrm{ord}}=\xi_{\mathrm{ext}}+F\); the transition
and its inverse are addition and subtraction of
\[
 F(s,v)=nv\{s\delta+\varepsilon(q)\}
       +nmv(v-1)\delta.
\]
It obeys \(F(s,v+1)-F(s,v)=nK(s,v)\).  Hence the based transition
along \(r_\theta=e^{2\pi i\theta}r\) is
\[
 C_\theta=F(s,v+\theta)-F(s,v)
   =n\theta K(s,v)+nm\theta(\theta-1)\delta.
\]
Its endpoints are \(0\) and \(nK(s,v)\), and
\(C_{\theta+1}-C_\theta=nK(s,v+\theta)\).  This identity uses the moving period \(K(s,v+\theta)\).  For \(m=3,n=1\) its
primitive angular endpoint is exactly \(K\).

Use the cusp cover \(q=Q_0^{2m}\).  A finite root-deck generator
changes \(s\) by one; a positive full \(Q_0\) loop changes it by
\(2m\).  For either operation, with integer increment \(j\),
\[
 F(s+j,v)-F(s,v)=njv\delta,\qquad
 C_\theta(s+j,v)-C_\theta(s,v)=nj\theta\delta.
\]
For a fixed physical point \(z\), its coordinate values satisfy
\[
 \xi_{\mathrm{ord}}^{\,s+j}(z)
   =\xi_{\mathrm{ord}}^{\,s}(z)+njv\delta.
\]
The directed chart identification from the \((s+j)\)-chart to the
\(s\)-chart therefore subtracts \(njv\delta\), with line
multiplier \(r^{-nj}\).  Its inverse, from the \(s\)-chart to
the \((s+j)\)-chart, adds \(njv\delta\) and multiplies by
\(r^{nj}\).  The extending coordinate is unchanged.  The
positively oriented root-deck clutch uses this latter, forward
identification; dividing its multiplier at \(v+\theta\) by its
multiplier at \(v\) gives the based line gauge
\[
 G_j(\theta)=e^{2\pi inj\theta},\qquad
 G_j^{-1}(\theta)=e^{-2\pi inj\theta}.
\]
With the positive physical delta circle its degree is
\((2\pi i)^{-1}\int_0^1G_j^{-1}G_j'\,d\theta=nj\).
For the full root loop this is \(2mn\), and for the specified
sixth-root cover it is \(+6\).

The ordinary chart also gives the parameter extension at fixed
\(q\ne0\).  Its suspension line multiplier is
\(r^m q^{-1}d_0(q)\), with the unchanged source translation.
On the infinite rank-one chain put
\[
 X_j=V_{\mathrm{ord}}/r^j,\qquad
 Y_j=r^{j+1}/V_{\mathrm{ord}},\qquad X_jY_j=r.
\]
The full suspension sends chart \(j\) to chart \(j+m\), with
nonzero fixed-\(q\) source and line multipliers.  The quotient has
\(m\) components.  The zero origin is
\(Z_{\mathrm{ord}}=1, X_0=1,Y_0=r\), and therefore extends over
\(r=0\) in this ordinary chart.

For the cyclic cusp descent one can equivalently use the ordinary
coordinate \(\widetilde\xi_{\mathrm{ord}}
=\xi_{\mathrm{ord}}+R\), where
\[
 R=\frac{ns\varepsilon}{2m}+\frac{ns^2}{4m}\delta.
\]
Its inverse subtracts \(R\); both are defined on the logarithmic
punctured cusp chart.  Subtraction of the two origins
\[
 \sigma_{\rm ext}=nv\varepsilon+nmv(v-1)\delta,\qquad
 o=-\frac{ns\varepsilon}{2m}-\frac{ns^2}{4m}\delta-nsv\delta
\]
gives these extending and ordinary coordinates.  Their transition is
\(\widetilde F=\sigma_{\rm ext}-o=F+R\), and
\[
 \widetilde F(s+2m,v)-\widetilde F(s,v)=nK+nm\delta.
\]
This is an integral period.  More precisely, the forward root arrow
adds \(2mnv\delta\) to \(F\) and
\(n\varepsilon+n(s+m)\delta\) to \(R\); hence it adds
\(nK+nm\delta\) to \(F+R\).  The reverse arrow negates all
three increments.  These formulas give the conjugacy in each
direction.  Since \(R\) is independent of \(v\), it preserves
the based angular clutch.  The origin \(\sigma_{\rm ext}\) extends at \(q=0\);
\(o\) provides the ordinary parameter extension just described at
fixed \(q\ne0\).  The degree \(2mn\) belongs to the full line
clutch.  Its quadratic angular term has zero endpoints and is
independent of \(q\).

For the based lift \(u_*=u_h-\kappa=-\Pi(e_u)\), put
\(K_*=K-3\kappa\). Its extending origin is
\(\Omega_{\rm ext,*}=\Omega_{\rm ext}+3nv\kappa\).
The coordinate map and inverse are
\[
 (Z_{\rm ext},V_{\rm ext})\longmapsto
 (r^{3np}Z_{\rm ext},r^{3nq_\delta}V_{\rm ext}),\qquad
 (Z_*,V_*)\longmapsto(r^{-3np}Z_*,r^{-3nq_\delta}V_*).
\]
For \(r\ne0\) these units extend through every cusp chart.
They carry the complete based angular path to
\(n\theta K_*+nm\theta(\theta-1)\delta\), with the same inverse
endpoint arrows. Their q-root difference is zero, so the full
line winding remains \(2mn\).
\end{proof}

\subsection{Iterates of the section line}
\begin{theorem}[Section-line iterates and direct images]\hypertarget{tag-02S4}{}\label{night:P3F.28}\label{stmt:P3F.28}
{\normalfont\scriptsize tag 02S4}\par

The OS49 linearization gives the exact sequence
$0\to T^*M\to M\to\mathcal O_{C_0}(-1)\to0$.
Its iterates alternate the two components, and
$R\pi_*(T^{j*}M)=0$ for every $j\in\mathbb Z$.
On the compact degree-three regular-binding model constructed in the proof,
the section line $M' =\mathcal O(G'-O')$ and its translate
$T'^*M'=M'(-\sum_i C_{0,i})$ both have square $-2$
and zero derived direct image.
\end{theorem}
\begin{proof}
Use the actual surface, section and line of
\SeriesRef{night:P3F.26}.  Translation \(T=t_{3G}\) and its
inverse preserve the elliptic projection \(\pi\).
The linearization is the sheaf map \(\phi:T^*M\to M\), whose
source at \(p\) is \(M_{p+3G}\); its suspension therefore covers
\(T^{-1}\).  The proved line identity
\(\operatorname{Hom}(T^*M,M)\simeq\mathcal O(C_0)\)
and the actual simple divisor of \(\phi\) give an isomorphism
\(T^*M\simeq M(-C_0)\) onto its image.
In local line frames \(\phi\) is multiplication by \(zu\), with
\(z=0\) defining \(C_0\) and \(u\) a unit.  Dividing by
\(u\), with inverse multiplication by \(u\), identifies the
image with the ideal \(zM\).  Its quotient is \(M|_{C_0}\).
Since \(C_0\simeq\mathbb P^1\), \(G\cdot C_0=0\), and
\(O\cdot C_0=1\), this is precisely \(\mathcal O_{C_0}(-1)\).
Thus
\begin{equation}
 0\longrightarrow T^*M\xrightarrow{\phi}M
 \longrightarrow\mathcal O_{C_0}(-1)\longrightarrow0.
 \hypertarget{tag-Q308}{}\label{os49:eq:one-step-line}
\end{equation}

The section \(3G\) meets the nonidentity component of the \(I_2\)
fibre, so \(T\) and \(T^{-1}\) exchange \(C_0,C_1\).
For \(j\ge1\), the actual composition
\[
 \phi_j=\phi\circ T^*\phi\circ\cdots\circ T^{(j-1)*}\phi
 :T^{j*}M\longrightarrow M
\]
has divisor \(D_j=\sum_{r=0}^{j-1}T^{r*}C_0\).  Hence
\[
 D_{2r}=r(C_0+C_1),\qquad
 D_{2r+1}=(r+1)C_0+rC_1,
\]
and \(\phi_j\) is the explicit isomorphism onto \(M(-D_j)\).
In particular
\begin{equation}
 T^*M\simeq M(-C_0),\qquad T^{2*}M\simeq M(-F_\infty).
 \hypertarget{tag-Q309}{}\label{os49:eq:two-step-line}
\end{equation}
For \(I_j=\phi_j(T^{j*}M)\), the successive quotient is
\(I_j/I_{j+1}=T^{j*}\mathcal O_{C_0}(-1)\).
The restriction \(T^j:C_{j\bmod2}\to C_0\) is an isomorphism,
so the successive quotient lines have degree \(-1\) on the
alternating components.  The same pulled-back one-step sequence
is defined for every integer \(j\), including negative \(j\).
The two affine charts of \(\mathbb P^1\) give
\(H^0(\mathcal O(-1))=H^1(\mathcal O(-1))=0\).
Thus every one-step quotient and every finite iterate quotient
has zero derived direct image and zero derived global sections.

To prove the source-line vanishing at every fibre, use the other
actual divisor sequence
\[
 0\longrightarrow\mathcal O_S(-O)\longrightarrow M
 \longrightarrow\mathcal O_G(G-O)\longrightarrow0.
\]
Here \(G^2=-1\) and \(G\cdot O=0\), so its last line is
\(\mathcal O_{\mathbb P^1}(-1)\).
Connected fibres and the section identify restriction
\(\pi_*\mathcal O_S\to\pi_*\mathcal O_O\) with the identity
of \(\mathcal O_{\mathbb P^1}\).  The sequence along \(O\)
therefore gives
\[
 \pi_*\mathcal O_S(-O)=0,\qquad
 R^1\pi_*\mathcal O_S(-O)=R^1\pi_*\mathcal O_S.
\]
The canonical-bundle formula gives
\(\omega_{S/\mathbb P^1}=\pi^*\mathcal O(1)\).
Relative duality gives \(R^1\pi_*\mathcal O_S=\mathcal O(-1)\).
Its hypotheses include the singular fibres here: \(\pi\) is
proper and flat, and its one-dimensional fibres are Cartier
divisors on the smooth surface, hence Gorenstein.  This is the
representing property after Definition~109.19.2 and
Lemma~109.19.3, tag~0E6N, in \cite{StacksProject}.
Applying \(R\pi_*\) yields
\[
 0\longrightarrow\pi_*M\longrightarrow\mathcal O(-1)
 \xrightarrow{\partial}\mathcal O(-1)
 \longrightarrow R^1\pi_*M\longrightarrow0.
\]
On a smooth fibre, the disjoint sections give a nontrivial
degree-zero line \(M_t\).  A nonzero section of that line would
have no zero and would trivialize it.  Thus its \(H^0\), and by
Serre duality its \(H^1\), vanish.  Consequently \(\partial\)
is nonzero.  An endomorphism of \(\mathcal O(-1)\) on
\(\mathbb P^1\) is a constant scalar, so it is everywhere an
isomorphism, with inverse that scalar's reciprocal.
The connecting map has zero kernel and cokernel at every fibre,
so \(R\pi_*M=0\).

For any \(j\in\mathbb Z\), the inverse surface maps
\(T^j,T^{-j}\) induce inverse direct-image cochain maps.
Explicitly, over \(\pi^{-1}U\), a section \(s\) of \(M\)
goes to \(x\mapsto s(T^jx)\) of \(T^{j*}M\); the inverse
evaluates at \(T^{-j}x\).  These formulas agree on intersections
and commute with the \v{C}ech differentials.  Hence
\begin{equation}
 R\pi_*(T^{j*}M)\simeq R\pi_*M=0
 \qquad (j\in\mathbb Z).
 \hypertarget{tag-Q30A}{}\label{os49:eq:translated-direct-images}
\end{equation}
\paragraph{The compact regular-binding pullback.}
For a regular binding value \(b\notin\{0,1,\infty\}\),
the base cover and inverse are
\[
 z^3=\frac{t}{t-b},\qquad t=\frac{bz^3}{z^3-1}.
\]
It is totally ramified at \(t=0,b\), and unramified elsewhere.
Normalize the Weierstrass pullback and take its relatively minimal
smooth model \(S'\).  Near zero choose \(t=s^3\); the actual
minimalization is \(x=s^4X,y=s^6Y\).  Thus its fundamental line is
\[
 \mathcal L'=h^*\mathcal O(1)\otimes\mathcal O(-2p_0),
 \qquad \deg\mathcal L'=1.
\]
There is no further minimalization, since the other ramified fibre
is smooth.  The lifted sections \(G',O'\) are disjoint; at
\(p_0\), \(G'=(s^2,1)\) tends to \((0,1)\), not to \(O'\).
The section line on this model is \(M'=\mathcal O(G'-O')\).
Its relative canonical line is \(\pi'^*\mathcal L'\), so
adjunction gives \(G'^2=O'^2=-1\) and \(M'^2=-2\).

Compare this line with the ordinary degree-three
pullback of \(M\) on the following common resolution.  The clockwise action on the good-reduction
chart is \((s,X,Y)\mapsto(\zeta^{-1}s,\zeta X,Y)\).
Its three central fixed points are \(O'\) and \((0,\pm1)\).
At each one choose local coordinates \((a,b_0)\) of characters
\((-1,1)\).  At \((0,1)\), take \(a=s,b_0=X-s^2\), so
\(b_0=0\) is \(G'\); at \(O'\), a relative formal-group
coordinate cuts out \(O'\).
The local quotient invariants \(a^3,b_0^3,ab_0\) satisfy
\(a^3b_0^3=(ab_0)^3\).  Its resolution lattice is
\(N_{\rm down}=\mathbb Z^2+\mathbb Z(2/3,1/3)\), with rays
\[
 (1,0),\quad(2/3,1/3),\quad(1/3,2/3),\quad(0,1).
\]
The base function \(a^3\) has multiplicities \(3,2,1\) on
these fibre rays; the three resolved quotient points give the
three arms of the original \(IV^*\) fibre.
The common smooth resolution \(Z\) has upstairs rays
\[
 (1,0),\quad(2,1),\quad(1,1),\quad(1,2),\quad(0,1).
\]
Let \(\beta:Z\to S'\) be their blow-down and
\(\rho:Z\to S\) the quotient-resolution map.  In the downstairs
basis \(((1,0),(2/3,1/3))\), its lattice map is
\(\left(\begin{smallmatrix}1&-2\\0&3\end{smallmatrix}\right)\).
The rays \((2,1),(1,2)\) map to three times the two resolving
rays, while \((1,1)\) maps to their common cone and is contracted.
All adjacent cones upstairs and downstairs are unimodular in
their stated lattices, giving these morphisms in their monomial
charts.

Over a fixed point \(P\), the three exceptional curves have
intersection matrix
\[
 Q=\begin{pmatrix}-1&1&0\\1&-3&1\\0&1&-1\end{pmatrix}.
\]
The total transform of the section \(b_0=0\) differs from its
strict transform by
\(V_P=E_{21,P}+E_{11,P}+2E_{12,P}\), with \(V_P^2=-2\).
These coefficients are exactly the valuations of \(b_0\) on the
three upstairs rays.  Downstairs, the section has Cartier equation
\(b_0/a^2\) on the cone adjacent to \((0,1)\), and coefficient
zero on both resolving rays.  Its pullback to \(Z\) is therefore
just its strict transform.  Consequently
\[
 \beta^*M'\simeq\rho^*M\otimes
       \mathcal O_Z(V_{G'}-V_{O'}).
\]
The stated monomial frames give this isomorphism and their
reciprocals its inverse; it is the identity on the common punctured
locus.  The two exceptional sets are disjoint, and \(\beta^*M'\)
has zero intersection with them.  Thus the projection formula gives
\[
 3M^2=(\rho^*M)^2=M'^2-4=-6.
\]
This verifies the square \(-2\) on the minimalized compact model
with its exact correction to degree-three pullback.

Infinity has three preimages \(z_i^3=1\).  Their local inverse
coordinate maps are \(z=z_i(1-bq)^{-1/3}\), with the branch
equal to one at \(q=0\); hence the cover and the line comparison
are local isomorphisms through both components and both nodes of
each \(I_2\) fibre.  The three identity components satisfy
\[
 M'\cdot C_{0,i}=-1,\qquad C_{0,i}^2=-2,\qquad
 C_{0,i}\cdot C_{0,j}=0\quad(i\ne j).
\]
For \(D=\sum_i C_{0,i}\), the actual Cartier modification is
\[
 0\longrightarrow M'(-D)\longrightarrow M'
 \longrightarrow\bigoplus_i\mathcal O_{C_{0,i}}(-1)
 \longrightarrow0.
\]
It gives an isomorphism of derived section complexes and equality
of their Dolbeault indices.
It is also the translated line for \(T'=t_{3G'}\).
The component calculation determines
\(\operatorname{Hom}(T'^*M',M')=\mathcal O(D)\) up to a base
line.  The section \(2G'\) tends to \((0,-1)\), whereas
the local coordinate of \(3G'\) at \(O'\) is
\[
 -X_{3G'}/Y_{3G'}=
 s^2\frac{3(9s^6-72s^3+64)}{351s^6-864s^3+512}
 =\frac38s^2+O(s^5).
\]
Thus \((2G')\cdot O'=0\), \((3G')\cdot O'=2\), and that
Hom line has degree \(1+2=3=D\cdot O'\) on \(O'\).
The base twist is zero, proving \(T'^*M'\simeq M'(-D)\).
Finally
\[
 (M'-D)^2=-2-2(-3)-6=-2.
\]
The disjoint sections and degree-one fundamental line permit the
same degree-minus-one connecting-map argument on \(S'\), giving
\(R\pi'_*M'=0\).  The modification sequence gives the same
vanishing for \(M'(-D)\).

\end{proof}

\subsection{Coprime finite caps}
\begin{remark}[The component-supported Type C marking]
\hypertarget{tag-Q30N}{}\label{os49:rem:type-c-marking}
{\normalfont\scriptsize tag Q30N}\par
The source of \eqref{os49:eq:source-cubic} is the source of
\cref{maps:src:D.2.11}: the identity on both Weierstrass charts
carries \(P\) to \(G\), \(Q=3P\) to \(3G\), \(O\) to \(O\), and
\(\mathcal O(P-O)\) to \(M\). The two simple-divisor linearizations
have divisor \(C_0\), so their ratio is a nonzero constant \(c\)
on the compact surface. The scalar map
\(\lambda_{\rm KL}=c^{-1}\lambda\), with inverse
\(\lambda=c\lambda_{\rm KL}\), makes the arrows identical.
The map on the line complement is then the identity.

In the homology order \((\gamma,u,w,\delta)\), the exact change
from the marking of \cref{night:P3F.26} to the index-one Type C
marking of \cref{mapsIII:def:yang-finite-data} is
\[
 G_0=\begin{pmatrix}
 -1&0&0&0\\0&-1&0&0\\0&-1&1&0\\0&0&0&1
 \end{pmatrix},\qquad G_0^{-1}=G_0.
\]
Writing \(A=T_A^{-t}\), \(B=T_B^{-t}\), its products are
\(A_1G_0=G_0A\) and \(A_2G_0=G_0B\).
The period matrix is \(\Pi_Y=\Pi G_0^{-1}\); the map on the
universal \(\mathbb C^2\) is the identity.
The actual periods are
\[
 \gamma_Y=-\gamma_g,\qquad u_Y=u_*-w_Z,\qquad
 w_Y=w_Z,\qquad\delta_Y=\delta_V,
\]
where \(u_*=-\Pi(e_u)\) is the based lift in
\eqref{os49:eq:based-cusp-period-map}. Its two branch shifts are
carried by that equation and its inverse.

The rigidified Poincar\'e frame of \cref{maps:src:D.2.12}
has the theta quotient
\(\vartheta(z-\mu)/\vartheta(z)\), whose divisor is \(G-O\)
and whose \(\tau_0\)-multiplier is \(e^{2\pi i\mu}\).
Consequently \(u_*=(\tau_0,\mu)\) and
\(\gamma_g=(-3\mu,\beta_P)\), because the forward arrow covers
\(-3G\). A remaining integral source-period correction would be
fixed by both finite elliptic matrices; the second matrix minus
the identity has determinant one, so this correction is zero.
The line-frame maps are
\(\eta=e^{-2\pi izy_2}\xi\) and
\(\xi=e^{2\pi izy_2}\eta\), with the rigidified lattice
multipliers proved there. Thus
\[
 \Pi_Y=
 \begin{pmatrix}3\mu&\tau&1&0\\\beta&\mu&0&1\end{pmatrix},
 \qquad \tau=\tau_0-1,\quad\beta=-\beta_P .
\]
The cycle change \((w_Z,u_*)\mapsto(w_Z,u_*-w_Z)\) has inverse
\((w_Z,u_Y)\mapsto(w_Z,u_Y+w_Z)\) and determinant one.
The source point \(G\) has value \(\tau_0/3=(\tau+1)/3\)
at \(IV^*\): its negative Abel--Jacobi fixed representative is
\((1/3,0)\) in \((e_u,e_w)=(-u_*,w_Z)\).
Its value at \(II\) is zero. At the cusp both
\(\tau_0-2s\) and \(\mu-s\) are holomorphic, so
\(\mu-\tau/2\) is holomorphic there.

The last two columns of the period equivariance give
\[
 R_1=\begin{pmatrix}-1/\tau&0\\-\mu/\tau&1\end{pmatrix},
 \qquad
 R_2=\begin{pmatrix}1/(\tau+1)&0\\-\mu/(\tau+1)&1\end{pmatrix}.
\]
The first two columns then give
\[
 \begin{aligned}
 \tau\circ g_1&=(\tau-1)/\tau,&
 \mu\circ g_1&=1-\mu/\tau,&
 \beta\circ g_1&=\beta+1-3\mu^2/\tau,\\
 \tau\circ g_2&=-1/(\tau+1),&
 \mu\circ g_2&=\mu/(\tau+1),&
 \beta\circ g_2&=\beta-3\mu^2/(\tau+1).
 \end{aligned}
\]
The finite values and cusp condition therefore identify this
\(\mu\) with \(\mu_C\) by
\cite[Proposition~3.4, p.~13]{manYang26}; its second affine coordinate
differs from the particular solution by the constant of
\cite[Proposition~3.6, p.~14]{manYang26}.
At the based fibre, write
\(\beta_P=\Gamma_2+\log\lambda/(2\pi i)\).
The parameter map and inverse are
\[
 c_Y=c_*-\frac{\log\lambda}{2\pi i},\qquad
 \lambda=e^{2\pi i(c_*-c_Y)},\qquad
 c_*=-\Gamma_2-\beta_{C,\mathrm{part}}(z_*).
\]
They use the chosen logarithmic chart; changes of logarithm are
the corresponding integral delta-period changes.
Small \(|\lambda|\) gives the admissible Type C half-plane.

The intermediate \((II,IV^*,I_2)\) marking
in \cref{maps:src:D.2.11} is obtained by the inverse Hurwitz map
\[
 (U,V)\longmapsto(UVU^{-1},U),\qquad
 (M_{II},M_{IV^*})\longmapsto
 (M_{IV^*},M_{IV^*}^{-1}M_{II}M_{IV^*}).
\]
These specify its based stems from those of \cref{night:P3F.26}.
The lattice map to its period order \((e_1,e_2,\delta,\gamma)\)
and its inverse are
\[
 F_0=\begin{pmatrix}3&2&-1&0\\0&1&0&0\\0&-1&0&1\\1&0&0&0\end{pmatrix},
 \quad
 F_0^{-1}=\begin{pmatrix}0&0&0&1\\0&1&0&0\\-1&2&0&3\\0&1&1&0\end{pmatrix}.
\]
For its displayed \({\mathsf S}\), one has \({\mathsf S}^tF_0=G_0\); these identities
carry the full period columns of this source.

For every \(m=3k\), the finite vector map is
\[
 G_0\{n(e_\gamma-e_u)+d e_\delta\}
 =-n(\gamma_Y-u_Y-w_Y)+d\delta_Y .
\]
Thus \((k,n_Y,r_Y)=(k,-n,d)\), with inverse
\((k,n,d)=(k,-n_Y,r_Y)\). The complete coprime domains agree.
In root coordinates the complex point map is the identity, the
real marking map is \(G_0\), and the root coordinate is fixed.
The identities \(A_1G_0=G_0A\) and
\(\Pi_Y(gs)G_0\ell=\Pi(gs)\ell\) carry the affine generator,
its inverse and both logarithmic collar translations exactly.
The positive exterior meridian is fixed; the positive normal
kernel \((\ell,m)\) maps to \((G_0\ell,m)\).

At the cusp the full multiplicative generators become
\(g_Y=g^{-1}\), \(h_Y=h\), with inverse
\(g=g_Y^{-1}\), \(h=h_Y\). Hence their subgroup, all its units
and the scalar are unchanged. Yang's triangles
\(\operatorname{conv}(v,v+e_1,v+e_2)\) and
\(\operatorname{conv}(v+e_1+e_2,v+e_1,v+e_2)\) are the same
negative-diagonal fan used here. The identity on every character
semigroup extends across every toric divisor and node. At \(II\) the source-line quotient is the identical one.
These maps agree on the ordinary overlaps and identify the
completed threefolds with all their specified attaching maps.
\end{remark}

\begin{theorem}[Integral topology of coprime OS49 caps]\hypertarget{tag-02S5}{}\label{night:P3F.29}\label{stmt:P3F.29}
{\normalfont\scriptsize tag 02S5}\par

For every \(k\in\mathbb Z_{>0}\), \(n,d\in\mathbb Z\), put \(m=3k\) and
\(\ell=n(1,-1,0,0)+d\delta\), with \(\gcd(n,m)=1\).
The finite specialization indices in degrees one through four
are \((m,k,k,m)\), and all four integral Leray differentials
are multiplication by \(+n\) in the marking of
\cref{night:P3F.26}. These caps give a proper holomorphic family
on the punctured parameter disc with \(\lambda=r^m\).
The fundamental group is \(\mathbb Z/|n|\), and
\begin{equation}
 H_0=H_6=\mathbb Z,\quad H_5=0,\quad
 H_1=H_2=H_3=H_4=\mathbb Z/|n|.
 \hypertarget{tag-Q30D}{}\label{os49:eq:coprime-homology}
\end{equation}
Every \(n=\pm1\) member is the standard smooth \(S^6\), for all
these \(k\) and \(d\).
\end{theorem}
\begin{proof}
Keep the actual source, component divisor, reduced fillings, exterior
section and cohomology marking of \SeriesRef{night:P3F.26}.
Write \(\ell_0=(1,-1,0,0)^t\), \(m=3k>0\), and
\begin{equation}
 \ell=n\ell_0+d\delta,\qquad \gcd(n,m)=1.
 \hypertarget{tag-Q30B}{}\label{os49:eq:coprime-cap-vector}
\end{equation}
The homology matrix \(A=T_A^{-t}\) fixes both \(\ell_0\)
and \(\delta\), and \(A^3=I\).

\paragraph{The actual cap and its boundary map.}
Pull the cubic good reduction back by the degree-\(k\) map of the
root disc.  Its explicit coordinates are
\(t=s^m,x=s^{4k}X,y=s^{6k}Y\).
The clockwise root action is the pullback of the order-three action
already marked in \SeriesRef{night:P3F.26}.
With the equivariant origin constructed there, its affine generator is
\[
 g(z,s)=(Az+\ell/m,e^{-2\pi i/m}s).
\]
Its \(m\)-th power is an integral translation on the torus, hence
the identity.  For \(0<j<m\), the \(j\)-th power changes the
invariant \(\psi\)-circle by \(jn/m\), which is nonintegral.
The action is therefore free on the central torus, for every
integer \(d\); the base rotation proves freeness elsewhere.
Its finite quotient is a smooth proper holomorphic cap.

Its overlap with the specified punctured family, and its inverse, are
\[
 h_s(z)=z+\frac{\log s}{2\pi i}\Pi(s)\ell,
 \qquad h_s^{-1}(z)=z-\frac{\log s}{2\pi i}\Pi(s)\ell.
\]
The clockwise logarithm change is \(-1/m\); invariance of
\(\ell\) gives \(h_{e^{-2\pi i/m}s}gh_s^{-1}=\rho\), where
\(\rho\) is the untwisted root action.  Branch changes add an
integral period.  Thus these maps give the marked
attachment, for every displayed integer choice.
The exterior section is represented by the negative logarithmic
translation, so its positive meridian maps to \(g^{-1}\).
The positive normal-circle kernel is consequently
\((\ell,m)\) in fibre and positive exterior-meridian coordinates.
The reduced pieces and the compact regular-base part are unchanged,
and give a compact smooth complex threefold after this gluing.
Pull the scalar family back by the covering map
\(p_m:\Delta_r^*\to\Delta_\lambda^*,\ p_m(r)=r^m\).
The fractional suspension
period \(n\Pi\gamma/m\) contains
\(n\log r\,\Pi\delta/(2\pi i)\); its branch changes are
integral.  The \(d\delta/m\) term is constant.  The same uniform
subannulus construction in \SeriesRef{night:P3F.26} therefore
gives the actual proper holomorphic family over \(\Delta_r^*\),
with all these cap maps single-valued there.

\paragraph{The quotient cells and the full specialization lattices.}
Use the determinant-one matrix \(P\), its inverse, and the block
\(F\) constructed in \SeriesRef{night:P3F.26}, so that in
homology coordinates \((u,w,\delta,\ell_0)\) one has
\(A=\operatorname{diag}(F,1)\), and \(P^{-1}FP\) cyclically
permutes the three circle factors.
Choose integers \(\alpha,\beta\) with \(m\alpha+n\beta=1\).
In coordinates \((a,b,c,t)\) the integral shear
\[
 (a,b,c,t)\longmapsto(a,b,c-\beta d\,t,t)
\]
has inverse with \(+\beta d\,t\).  It transforms the affine
generator to \((Fx,t+n/m)\), since its remaining \(\delta\)
translation is the integral vector \(\alpha d\delta\).
The element \(\ell_0^\alpha g^\beta\) advances \(t\) by
\(1/m\).  Thus the central quotient, with its endpoint direction
specified, is
\[
 Q_m=(T^3\times[0,1])/((x,0)\sim(F^\beta x,1)).
\]
The map from the original \(T^4\) is its \(m\)-fold unwrapping,
followed by the inverse shear.  The inverse coordinate maps are the
displayed shear and \(P^{-1}\).  The base coordinate is \(mt\),
so the oriented covering degree is \(+m\), including when \(n<0\).

The product cell structure with one vertex and one edge on each
circle has zero fibre differential.  In degrees one and two, the
action is a three-cycle on the oriented cell bases; in degrees zero
and three it is the identity.  Its mapping-torus cone therefore has
torsion-free cohomology of ranks \((1,2,2,2,1)\).
The cochain pullback is identity on invariant fibre cells and the
norm on cells meeting the base edge.  As \(\beta\) is prime to
three, this norm is exactly
\[
 1+C^\beta+\cdots+C^{(m-1)\beta}=k(1+C+C^2),
\]
where \(C\) is the three-cycle.

With \(\alpha_2=da\wedge db\), \(\lambda_1=3dc-db\),
\(v_3=da\wedge db\wedge dc\), undoing the integral shear gives
the complete specialization images inside the local invariant
lattices:
\[
 \begin{array}{c|l}
 j&\operatorname{im}(H^j(Q_m;\mathbb Z)\to H^j(T^4;\mathbb Z))\\\hline
 1&\langle\lambda_1-3\beta d\,dt,\ mdt\rangle\\
 2&\langle\alpha_2,\ kdt\lambda_1\rangle\\
 3&\langle v_3-\beta d\,dt\alpha_2,\ kdt\alpha_2\rangle\\
 4&\langle mdtv_3\rangle.
 \end{array}
\]
The full invariant bases before these inclusions are respectively
\((\lambda_1,dt)\), \((\alpha_2,dt\lambda_1)\),
\((v_3,dt\alpha_2)\), and \(dtv_3\).
Hence all four cokernels are cyclic.
The global invariant classes are
\[
 \psi=dt,\quad\omega=\alpha_2+dt\lambda_1,
 \quad\phi=dt\alpha_2,\quad\nu=dtv_3.
\]
Each generates its corresponding cokernel, proving both the full
local images and the claimed global specialization indices.
\begin{equation}
 (m,k,k,m).
 \hypertarget{tag-Q30C}{}\label{os49:eq:coprime-specialization-indices}
\end{equation}
In the boundary basis the explicit change is
\[
 h=\lambda_1,\quad u_2=dt\lambda_1,\quad v_2=\alpha_2,
 \quad\chi=v_3,\quad\phi=dt\alpha_2;
\]
thus the degree-three congruence is
\(a+d\beta b\equiv0\pmod k\), with inverse identification
\((a,b)\mapsto a\phi+b\chi\).

\paragraph{Every integral differential and the resulting topology.}
The exact image-lattice parabolic complexes, their integer
contractions, and their \(H^2\) quotient maps are those in
\SeriesRef{night:P3F.26}.  Only the just computed finite stalk
inclusion changes.  Since the global invariant reaches every full
local cokernel, the skyscraper exact sequence has surjective global
map, gives \(H^1(\mathbb P^1,R^jf_*\mathbb Z)=0\), and keeps
the same free \(H^2\) groups.  Its \(H^0\) generators are now
\[
 1,\quad m\psi,\quad k\omega,\quad k\phi,\quad m\nu.
\]

The linear coarse-cap comparison, its explicit inversion maps and
the exterior section in \SeriesRef{night:P3F.26} apply to the same
root pullback.  They give zero intrinsic rational differential.
The actual logarithmic collar has the translation prism
\(\iota_\ell/m\).  Its radial and unit terms have the same
deck-equivariant translation isotopies, with inverse subtraction.
The positive annular base orientation therefore gives
\[
 d_2(a)=\frac{H_B}{m}[\iota_\ell a].
\]
In the four exact coinvariant targets, contraction of
\(\psi,\omega,\phi,\nu\) by \(\ell_0\) is
\(1,3[\delta],3\eta,[uw\delta]\), respectively; contraction
by \(\delta\) maps to zero in each case.
Multiplying by the actual source indices gives
\[
 \frac1{3k}(3k,k,k,3k)\cdot n(1,3,3,1)=(n,n,n,n).
\]
The actual sources and targets are free, so this rational
determination is their unique integral map.  These are all four
Leray differentials, including their common sign.

Remark~\ref{os49:rem:type-c-marking} identifies this entire
family with \cref{maps:src:D.2.11} by
\((k,n_Y,r_Y)=(k,-n,d)\). Its hypothesis is
\(\gcd(n_Y,3k)=1\), exactly the present hypothesis. Its original
type-\(II\) and index-one toroidal \(I_2\) pieces are the pieces
carried by the displayed root, collar and fan maps. The group
generator maps are \(\gamma\mapsto-\gamma_Y\) and
\(\gamma_Y\mapsto-\gamma\).
That theorem therefore gives the stated fundamental group,
\eqref{os49:eq:coprime-homology} and the smooth sphere conclusion
for \(n=\pm1\). Equivalently, integral duality gives
\[
 H^0=H^6=\mathbb Z,\quad H^1=0,\quad
 H^2=H^3=H^4=H^5=\mathbb Z/|n|.
\]
Together with the complete specialization lattices and the four
signed differentials above, this proves every stated assertion.
\end{proof}

\subsection{Return of the positive-diagonal family}
\begin{theorem}[Return of the positive-diagonal OS49 family]\hypertarget{tag-02S8}{}\label{night:P3F.31}\label{stmt:P3F.31}
{\normalfont\scriptsize tag 02S8}\par
For the positive-diagonal OS49 filling over a sufficiently small
punctured \(r\)-disc, with
\(\lambda=r^3\), the fibres are standard smooth six-spheres.
The positive \(r\)-loop has return integer zero and oriented
return class zero in \(\Theta_7\), and its mapping torus is
diffeomorphic to \(S^6\times S^1\).
The page constructed below carries the source section line
through both \(I_2\) components and both nodes.
The negative-diagonal filling has the native C2 comparison of
\SeriesRef{night:P3K.65} and positive \(r\)-loop class \(-3\).
\end{theorem}
\begin{proof}
We use the source and section of \SeriesRef{night:P3F.26}.
On \(q=t^{-1}\) its periods and physical line are
\[
 \begin{aligned}
 h(Z,V)&=(q^2a_{\rm un}Z,qb_{\rm un}V),\\
 g(Z,V)&=(q^{-3}c_{\rm un}Z,\lambda q^{-1}d_0V),\qquad M=\mathcal O(G-O),
 \end{aligned}
\]
with holomorphic units \(a_{\rm un},b_{\rm un},c_{\rm un},d_0\).
The forward arrow covers \(-3G\), has divisor \(C_0\) in the
target variable and zero \(C_1\) in the source variable.

\hypertarget{tag-034S}{}\label{the-exact-line-subfan}
\paragraph{The positive fan and source line.}
The height-one triangles are
\begin{equation}
 (i,j),(i+1,j),(i+1,j+1),\qquad
 (i,j),(i,j+1),(i+1,j+1),\quad (i,j)\in\mathbb Z^2.
 \hypertarget{tag-Q30E}{}\label{os49:eq:positive-diagonal-fan}
\end{equation}
Each cone is unimodular.  The period valuation lattice is all
\(\mathbb Z^2\), so the compact fundamental-region construction
of \SeriesRef{night:P3F.26} gives a smooth proper filling.

The common refinement of its two diagonals adds \((1,1,2)\).
On the conifold \(AD=BC\), its map is
\[
 \operatorname{Tot}\mathcal O_{\mathbb P^1\times\mathbb P^1}(-1,-1)
 \longrightarrow\{AD=BC\},\qquad
 \begin{pmatrix}A&B\\C&D\end{pmatrix}
  =\xi\begin{pmatrix}s\\t\end{pmatrix}\begin{pmatrix}p&q_1\end{pmatrix}.
\]
One contraction keeps \([s:t]\) and the vector \((\xi p,\xi q_1)\);
the other keeps \([p:q_1]\) and \((\xi s,\xi t)\).
Off the zero sections their inverses recover the discarded
projective coordinate from the nonzero vector.  On the exceptional
divisor they contract its two rulings.  Both centres have normal
bundle \(\mathcal O(-1)^{\oplus2}\).
The full period maps preserve the corresponding monomial ideals,
including their unit multipliers, so these maps descend; there is
one orbit of centres.
For either blowup,
\[
 H^j(\widetilde X;\mathbb Z)
   =H^j(X;\mathbb Z)\oplus H^{j-2}(\mathbb P^1;\mathbb Z).
\]
The exceptional-divisor degree \(-1\) on its ruling supplies an
integral left inverse of the Gysin summand.  Van Kampen preserves
the fundamental group, since the local pieces and their sphere-bundle
boundary are simply connected.  Applying both maps to the same
refinement transfers the sphere homology and simple connectivity
from \SeriesRef{night:P3F.26}.  The fibres are standard smooth
six-spheres, using \(\Theta_6=0\)
\cite[the table and Theorem~5.1, p.~512]{manKM63}.

The source line has height \(k_j=\lceil j/2\rceil\).
Indeed \(k_0=0\), \(k_{j+2}=k_j+1\), and
\(2k_j-k_{j-1}-k_{j+1}\) is \(-1,+1\) for even, odd \(j\).
Its rays \((j,k_j,1)\) and consecutive-pair cones are contained
in the positive fan.  The edge from \((2l,l)\) to
\((2l+1,l+1)\) has direction \((1,1)\).
The negative-diagonal fan has edge directions
\((1,0),(0,1),(1,-1)\), so that particular source edge is
not one of its cones.  The common refinement above gives
the toric morphisms relating the two fillings.
The source coordinates, with \(xy=q\), are
\[
\begin{array}{c|ccc}
 j&x&y&\eta\\ \hline
 2l&Z/q^{2l}&q^{2l+1}/Z&q^lV/Z\\
 2l+1&Z/q^{2l+1}&q^{2l+2}/Z&q^{-l-1}V .
\end{array}
\]
The consecutive line transitions are \(\eta_{j+1}=\eta_j/y_j\)
for even \(j\) and \(\eta_{j+1}=y_j\eta_j\) for odd \(j\),
with reciprocal inverses.  They have component degrees \((-1,+1)\).
The generic line isomorphism, rigidified at the zero section,
extends through these charts: any remaining vertical difference
of degree zero on both components is a whole-fibre multiple,
whose multiplier \(q^{-e}\) has \(e=0\) by that rigidification.

The full \(g\)-maps into chart \(j-3\) and their inverses are
\[
\begin{array}{c|c|c}
 &g(x,y,\eta)&g^{-1}(x',y',\eta')\\ \hline
 j\text{ even}
 &(c_{\rm un}x,c_{\rm un}^{-1}y,\lambda d_0x\eta)
 &(x'/c_{\rm un},c_{\rm un}y',c_{\rm un}\eta'/(\lambda d_0x'))\\
 j\text{ odd}
 &(c_{\rm un}x,c_{\rm un}^{-1}y,\lambda(d_0/c_{\rm un})y\eta)
 &(x'/c_{\rm un},c_{\rm un}y',\eta'/(\lambda d_0y')) .
\end{array}
\]
Each inverse is used on its nonzero-multiplier overlap.
The \(h\)-map shifts \(j\) by two and has line multipliers \(b_{\rm un}/a_{\rm un},b_{\rm un}\).
Thus the nonzero arrow on the even source component maps to the
odd component, and its simple source zero is \(C_1\).
Choose a Hermitian norm flat along every regular elliptic fibre
on an outer cusp annulus.  The global operator norm of the line
arrow is bounded on the compact source.  For small \(|\lambda|\)
it is bounded by \(\kappa<1\).  Choose three close norm levels
with \(\kappa R_3<R_1\).  Iterates staying in the line complement
cannot identify them.  Iterates reaching a zero cannot re-enter,
because their height defect is
\[
 k_j-N-k_{j-3N}
 =\begin{cases}\lfloor N/2\rfloor&j\text{ even},\\
                 \lceil N/2\rceil&j\text{ odd}.
   \end{cases}
\]
At a node both endpoint defects apply.  These injective immersions
of compact shells are embeddings; their free physical-circle
quotients contain each complete source \(I_2\) neighbourhood.

\paragraph{The page and its compact model.}\leavevmode
\hypertarget{tag-0003}{}\label{a-page-carrying-the-actual-source-line}
On a regular cut band write the arrow in a unitary line coordinate as
\[
 F_t(z,\eta)=(z-q_Q(t),e^{L_{\rm line}(t)}\eta),\qquad \Re L_{\rm line}<0,
\]
where \(q_Q\) is the logarithm of \(3G\).
Translation pulls a flat metric to a flat metric; its logarithmic
norm ratio is harmonic and hence constant on the compact fibre.
Choose \(\vartheta_{\rm sus}(a_{\rm sus}+1)=\vartheta_{\rm sus}(a_{\rm sus})+1\), equal to zero on a short interval near
zero and to one near one.  The maps
\[
 \begin{aligned}
 \Phi_{a_{\rm sus}}(z,\eta)
 &=
   \bigl(z-\vartheta_{\rm sus}(a_{\rm sus})q_Q,\\
 &\quad e^{a_{\rm sus}\Re L_{\rm line}+i\vartheta_{\rm sus}(a_{\rm sus})\Im L_{\rm line}}\eta\bigr)
 \end{aligned}
\]
have inverse given by addition and reciprocal multiplication and
satisfy \(\Phi_{a_{\rm sus}+1}=F_t\Phi_{a_{\rm sus}}\).
They give the marked suspension product, including the physical
line.  The inverse reads \(a_{\rm sus}=\log|\eta|/\Re L_{\rm line}\pmod1\) and applies
\(\Phi_{a_{\rm sus}}^{-1}\).  Where \(\vartheta_{\rm sus}=0\) it is the source radial shell chart.
An increasing circle reparametrization joins the three ordered
page levels to the shell levels.  Its lift is
\[
 \begin{gathered}
 \eta\longmapsto c_{\rm sh}(t)^{\vartheta_{{\rm sh},t}(a_{\rm sus})-a_{\rm sus}}\eta,\\
 a_{\rm sus}=\log\|\eta\|/\log c_{\rm sh}(t),\\
 \vartheta_{{\rm sh},t}(a_{\rm sus}+1)=\vartheta_{{\rm sh},t}(a_{\rm sus})+1,
 \end{gathered}
\]
with inverse using \(\vartheta_{{\rm sh},t}^{-1}\); the positive multiplier is periodic
and commutes with the arrow.  The cyclic cover
\(\zeta^3=t/(t-t_b)\) supplies the three sheets.
The angular cut deformation has zero-winding ratio, so its real
logarithm divided by three gives the cover map and subtraction its
inverse.

At the finite end the affine deck advances the invariant cap
coordinate by \(1/3\) and rotates \(\sigma\) by \(e^{-2\pi i/3}\).
Its three page levels form one orbit; selecting level zero gives
the good-reduction chart and its inverse, also on unit line vectors.
Its logarithmic attachment is
\[
 z\longmapsto z+\frac{\log\sigma}{2\pi i}\Pi(\sigma)\ell_0,
 \qquad t=\sigma^3,\qquad \ell_0=e_\gamma-e_u,
\]
with inverse subtraction.  The Weierstrass map
\(x=\sigma^4X,y=\sigma^6Y\) and its punctured inverse extend by
the good-reduction model of \SeriesRef{night:P3F.28}.
The type-\(II\) charts use the original line-complement quotient.

For smooth comparison at the regular binding, write
\(q-\epsilon=\rho e^{2\pi ie\beta+i\alpha_0}\), \(e=3\).
Choose a smooth positive-derivative function \(s(\rho)\), with
\(s(0)=s_0>0\), equal to \(\rho^{1/e}\) for \(\rho\geq\rho_1\).
One explicit choice is
\[
 \begin{aligned}
 s(\rho)&=s_0+\int_0^\rho
 \bigl((1-\chi(t))c_{\rm sm}+\chi(t)e^{-1}t^{1/e-1}\bigr)\,dt,\\
 c_{\rm sm}&=\frac{\rho_1^{1/e}-s_0-\int_0^{\rho_1}
              \chi(t)e^{-1}t^{1/e-1}\,dt}
          {\int_0^{\rho_1}(1-\chi(t))\,dt},
 \end{aligned}
\]
where \(\chi=0\) near zero and \(\chi=1\) near \(\rho_1\);
take \(s_0\) small enough that \(c_{\rm sm}>0\).
The map to the compact-cover annulus is
\(\zeta_{\rm loc}=s(\rho)e^{2\pi i\beta}\);
its inverse reads \(\beta\) and \(s^{-1}(|\zeta_{\rm loc}|)\).
Transport the elliptic point and line vector between the two
regular base values using their marked period coordinates.
The inverse exchanges those base values.  These maps identify
\[
 V_3=S'_3\setminus(E_{t_b}\times\operatorname{int}D^2),\qquad
 M_{\rm page}=\mathcal O(G'_3-O'_3)|_{V_3}.
\]
The compact model has three \(II\) and three \(I_2\) fibres,
disjoint sections of square \(-1\), \(M_{\rm page,comp}^2=-2\),
\(c_2=12\), \(c_1^2=0\), \(p_1=-24\), and signature \(-8\).
For its section height, the class
\(G'_3-O'_3-F+\frac12(C_1+C_2+C_3)\) has square \(-1/2\),
so the height is \(1/2\) and the component correction is \(3/2\).

\paragraph{The cusp product and based marking.}\leavevmode
\hypertarget{tag-034V}{}\label{the-full-graph-to-page-comparison-by-regular-bands}
The full period \(K=h^3g^2\) is
\[
 (Z,V)\longmapsto(A_{\rm un}(q)Z,q\lambda^2B_{\rm un}(q)V),\qquad
 A_{\rm un}=a_{\rm un}^3c_{\rm un}^2,\quad B_{\rm un}=b_{\rm un}^3d_0^2 .
\]
It covers \(-6G\), and its inverse covers \(+6G\).
On a height cover its unit straightening is
\[
 (Z,V,\tau)\longmapsto
 (e^{-\tau\log A_{\rm un}}Z,e^{-\tau\log(\lambda^2B_{\rm un})}V,\tau),
\]
with inverse using the two positive exponents.  It conjugates
the period to \((Z',V',\tau)\mapsto(Z',qV',\tau+1)\).
The positive-fan fixed-edge charts are
\[
 (A_0,B_0,Z)=(q/V,V/Z,Z),\qquad
 (A_1,B_1,Z_1)=(qZ/V,V,1/Z),
\]
with \((A_1,B_1)=(ZA_0,ZB_0)\) and physical weights \((-1,+1,0)\).
Their quotient map is
\[
 (A_i,B_i)\longmapsto
 (Q_i,\mu_i)=(2A_iB_i,|B_i|^2-|A_i|^2),
\]
with inverse orbits determined by
\[
 |B_i|^2=\frac{\sqrt{\mu_i^2+|Q_i|^2}+\mu_i}{2},\quad
 |A_i|^2=\frac{\sqrt{\mu_i^2+|Q_i|^2}-\mu_i}{2},\quad
 A_iB_i=Q_i/2.
\]
Successive real-height charts have positive transition derivative:
in an ordinary chart they are \(R_{\rm ht}/X_{\rm ht}-Q_{\rm ht}/R_{\rm ht}\) and
\(R_{\rm ht}/(Q_{\rm ht}X_{\rm ht})-Q_{\rm ht}^2/R_{\rm ht}\), where \(R_{\rm ht}=|V|^2,X_{\rm ht}=|Z|^2,Q_{\rm ht}=|q|^2\).
At the central overlap it is \(\mu'=-1/(X_{\rm ht}\mu)\), also increasing.
An invariant increasing vector field on their height cover,
divided by its positive period integral, supplies the circle
product; its reverse flow gives the inverse.  Fix it on the outer
collar and retain the unit straightening above.

On \(|Z|=1\) the fixed normal transition is
\((Q_1,\mu_1)=(Z^2Q_0,\mu_0)\).  Its lift to \(SU(2)\) is filled by
\[
 \begin{pmatrix}
 \sin(\pi\rho/2)e^{i\phi}&-\cos(\pi\rho/2)\\
 \cos(\pi\rho/2)&\sin(\pi\rho/2)e^{-i\phi}
 \end{pmatrix},\qquad 0\leq\rho\leq1 .
\]
Conjugation and its inverse trivialize the real fixed normal bundle.
There is one fixed-edge orbit \({\mathcal F}=S^2\); all other physical-circle
orbits, including the finite cap, are free.
For \(E=X\setminus {\mathcal F}\), general position gives \(\pi_1(E)=0\)
and Alexander duality gives the homology of \(S^3\).
The circle quotient \({\mathcal Q}_0\) is simply connected; its Gysin sequence,
with the degree-one normal Hopf meridian, gives the homology of \(S^2\).
In \(Y={\mathcal Q}_0\cup(S^2\times D^3)\), the boundary degree-two map has
matrix \(\left(\begin{smallmatrix}\alpha_{\rm bd}&1\\1&0\end{smallmatrix}\right)\),
of determinant \(-1\).  Mayer--Vietoris and van Kampen give a
simply connected homology five-sphere, hence \(Y\cong S^5\)
\cite[the introductory table and comparison, pp.~504--505]{manKM63}.

Use the complete based period map
\eqref{os49:eq:based-cusp-period-map}, and write its retained vector
as \(\kappa_0=pw_Z+q_\delta\delta_V\).
Thus \(u_*=u_h-\kappa_0=-\Pi(e_u)\).
Put \(H_*=-u_*-\gamma\), \(K_*=3u_*+2\gamma\).
The change of period columns and its inverse are
\[
 L=\begin{pmatrix}1&0&0&0\\0&1&0&0\\0&-p&1&0\\0&-q_\delta&0&1\end{pmatrix},
 \quad
 L^{-1}=\begin{pmatrix}1&0&0&0\\0&1&0&0\\0&p&1&0\\0&q_\delta&0&1\end{pmatrix}.
\]
Here the new coordinates are \(L^{-1}\) times the old coordinates.
In particular
\(L^{-1}(1,1,-p,-q_\delta)^t=(1,1,0,0)^t\);
the new linear part is \(L^{-1}A_{\rm old}L\).
The ordinary meridian is \(\beta_0=-\ell_0=H_*\), with
\[
 u_*=2H_*+K_*,\quad\gamma=-3H_*-K_*,
 \qquad h_*=2\xi+\beta,\quad\tau_*=\xi,\quad
 \xi=\tau_*,\quad\beta=h_*-2\tau_* .
\]
The quotient point map is \(w_*=w-ph+3p\tau,h_*=h,\tau_*=\tau\),
with inverse \(w=w_*+ph_*-3p\tau_*\).
The physical coordinate has
\(\delta_*=\delta-q_\delta h+3q_\delta\tau\).
On the compression cap the maps are
\[
 (z_w,V)\longmapsto
 (e^{-2\pi iph_*+6\pi ip\tau_*}z_w,
  e^{-2\pi iq_\delta h_*+6\pi iq_\delta\tau_*}V),
\]
with reciprocal inverses extending across \(z_w=0\).
For \(r\ne0\), the parameter-chart comparison is
\((Z,V)\mapsto(r^{3p}Z,r^{3q_\delta}V)\), with reciprocal inverse.
It carries the entire angular path on the punctured parameter
disc to the based \(K_*\) path.
The raw pair below is \((e_1,e_2)=(-w_Z,-u_*)\).
Its determinant is \(+1\) in the source complex pair
\((w_Z,u_*)\), whose second period has positive imaginary part.

We make the same page construction for \(m=3k>0\), with \(m\) close
shells and
\[
 \zeta^m=t/(t-t_b),\quad t=t_b\zeta^m/(\zeta^m-1),\quad
 t_b=1/\epsilon>2.
\]
The finite root is \(t=\sigma^{3k},x=\sigma^{4k}X,y=\sigma^{6k}Y\).
Its affine generator advances the invariant suspension coordinate
by \(1/m\), and selection of level zero gives its good-reduction
page chart.  Apply the preceding source-shell, line and smooth
binding maps with \(e=m\).  This constructs
\(V_m=S'_m\setminus(E_{t_b}\times\operatorname{int}D^2)\).
The sphere topology for these orders follows from
\SeriesRef{night:P3F.29} through the same common-refinement maps.
The preceding circle-quotient calculation then gives \(Y_m\cong S^5\).
The compact model is the degree-\(k\) pullback of \(S'_3\) under
\(\zeta_3=\zeta^k\), branched at its two smooth end fibres.
Its disjoint sections remain disjoint; it has \(m\) \(II\) and \(m\) \(I_2\)
fibres and fundamental-line degree \(k\).  Thus
\[
 c_2(S'_m)=12k,\quad c_1(S'_m)^2=0,\quad p_1(S'_m)=-24k,
 \quad \sigma(S'_m)=-8k
\]
by the four-dimensional signature formula
\cite[p.~173]{MilnorDifferentialTopology1964}.

The compression in the based cusp product is
\[
 (z,\phi)\longmapsto
 (x=z\cos\phi,y=\bar z\sin\phi,\tau_*=g(z,\phi)).
\]
Off \(z=0\) its inverse reads \(\phi\) from \(|y|/|x|\), then \(z\).
On its interior core it reads \(\phi\) from the strictly monotone
height.  Both that core and the normalization ends are included.
Its base value is \(|z|^2\cos\phi\sin\phi\), and its boundary is
\({\mathcal F}-\bar E_\epsilon\).  The incoming graph is
\(\tau_*=g_\epsilon(p)\).
The suspension class satisfies
\(\psi(H_*)=-1\) and \(\psi(K_*)=2\).
The displayed change of marking fixes \(\psi\), with its inverse
fixing the same covector. On the incoming graph its phase is
\(-h_*+2g_\epsilon(p)\). Its restriction to the \(H_*\)-circle
has degree \(-1\), and therefore takes every phase.
The page graph, outer trace and normal quotient are
\[
\begin{gathered}
 f(p,\tau_*)=r_0e^{2\pi i(\tau_*-g_\epsilon(p))},\qquad
 ((y,r),\theta)\longmapsto(R_\theta y,re^{2\pi i\theta}),\\
 [(v_T+a n_+,z_0)]\longmapsto(a,z_0-df(v_T)),\qquad
 (a,z_0)\longmapsto[(a n_+,z_0)] .
\end{gathered}
\]
The graph cutoff is flat at zero and vanishes off its binding collar.
The outer inverse reads \(r,\theta\) and applies \(R_{-\theta}\).
It attaches at \(\theta=\tau_*-g_\epsilon(p)\) and has outgoing
map \((p,\theta)\mapsto(R_\theta p,e^{2\pi i\theta})\).
These give the trace \(X_m\) and its fixed meridian-one class \(x\).
The outer line is the physical dual through \(D_{\rm dual}(v)(v)=1\),
\(D_{\rm dual}(av)=a^{-1}D_{\rm dual}(v)\).
In source log-norm \(c\) the normal is \(-\partial_c\);
the ambient orientation in
\((x_R,x_I,y_R,y_I,c,\Re z_0,\Im z_0)\) is the negative coordinate
orientation.  The compression orientation is
\(-(z_R,z_I,\phi)\), and the outer trace frame is
\((-W_{z_R},W_{z_I},W_\phi,X_\theta)\).

\paragraph{The compatible core arc.}
The even source node has fixed-edge substitution
\((A_0,B_0,Z)=(y/\eta,\eta,x)\) and core value \(\mu=|\eta|^2>0\).
At the odd node the full \(H^{-1}\)-period gives
\[
 (A'_0,B'_0,Z')=
 \bigl((\lambda b_{\rm un}d_0\eta)^{-1},
       (\lambda b_{\rm un}d_0/(a_{\rm un}c_{\rm un}))y\eta,a_{\rm un}c_{\rm un}x\bigr),\qquad
 \mu=-|\lambda b_{\rm un}d_0\eta|^{-2}<0 .
\]
The quotient-node change \((x,y)\mapsto(a_{\rm un}c_{\rm un}x,(a_{\rm un}c_{\rm un})^{-1}y)\) has
reciprocal inverse and positive real determinant one.
The two nodes lie in the complementary free core arcs, with
log-height derivatives of signs \(+1\) and \(-1\).
Every complete source shell meets each arc at its corresponding node.

On any regular \(q\)-band, the field
\[
 T=\frac{z_R\partial_{z_R}+z_I\partial_{z_I}}
          {2|z|^2\cos\phi\sin\phi}
\]
has \(dq(T)=1\) and preserves the \(w\) phase.
Patch it to the exterior product field and average over \(w\) rotation.
Its flow and reverse flow give inverse maps
\(\bar E\times[q_0,q_1]\leftrightarrow W_{[q_0,q_1]}\)
preserving the same primitive \(w\).
A height-graph isotopy, fixed near the core, makes the incoming
height \(g_*\) constant.  Change the graph angle simultaneously
to \(\tau_*-g_{\epsilon,s}(p)\).
The quotient-point and trace inverses prove embedding throughout.
Its real-logarithm incoming gauge extends over the cap by a cutoff.

The collar equation is
\[
 q-\epsilon=\rho e^{2\pi im\beta+i\alpha_0},
 \qquad \beta=h_*-2\tau_* .
\]
Its positive coorientation agrees with the source normal.
Indeed, with \(L_q=\log|q|\), put
\[
\begin{aligned}
 Z_H&=L_q-\log|a_{\rm un}c_{\rm un}|,&Z_K&=\log|a_{\rm un}^3c_{\rm un}^2|,\\
 V_H&=-\log|\lambda b_{\rm un}d_0|,&
 V_K&=L_q+2\log|\lambda|+\log|b_{\rm un}^3d_0^2|,\\
 D_{\rm rad}&=Z_HV_K-Z_KV_H .
\end{aligned}
\]
Inverting this radial period matrix gives
\[
 \frac{d\beta}{dc}
  =-\frac{Z_K+2Z_H}{D_{\rm rad}}
  =-\frac{\log|q^2a_{\rm un}|}{D_{\rm rad}}>0,
\]
since \(D_{\rm rad}=L_q^2+O(|L_q|)\) and
\(\log|q^2a_{\rm un}|=2L_q+O(1)<0\) on a sufficiently small cusp disc.
For
\(G_{\rm bind}=\Im(e^{-2\pi im\beta-i\alpha_0}(q-\epsilon))\),
the derivative on \(-\partial_c\) is positive for \(\rho>0\).
Near \(\rho=0\), interpolate in that positive half-space to
\(i e^{2\pi im\beta+i\alpha_0}\partial_q\), on which \(dG_{\rm bind}=1\).
This gives a smooth positive binding normal and its inverse homotopy.

Make \(V_m\cap W\) transverse by a relative interior height
perturbation on a regular band, and orient the intersection by
coorientation first.  At a selected node its local disc has
orientation \(\varepsilon_{\rm node}\) times complex \(z\) orientation, where
\(\varepsilon_{\rm node}=\operatorname{sign}(c_\phi)\); its small positive-\(q\)
boundary is \(\varepsilon_{\rm node} w\).
There are exactly \(m\) such discs: the inverse cusp charts are
\(\zeta=\zeta_j(1-t_bq)^{-1/m}\), \(\zeta_j=e^{2\pi ij/m}\),
and each shell contributes one selected node.
The page misses the normalization ends; after removing these
node neighbourhoods, compactness excludes further lower-level
intersections.
At the upper collar the positive-real equation is
\(2\pi m(h_*-2g_*)+\alpha_0=\pi\pmod{2\pi}\).
Its \(m\) circles have coorientation derivative
\(-2\pi m(\epsilon-q)dh_*\).
Since \(\Omega_W=-dw\wedge dh_*\wedge dq\), their upper
boundary orientation is \(+w\).
The intersection surface in the compact band is proper, with
exactly these two boundary lists; extra closed components
contribute zero boundary.
Projection through the product inverse gives
\((m-\varepsilon_{\rm node} m)w=0\) in the free abelian group \(H_1(\bar E;\mathbb Z)\).
Thus \(\varepsilon_{\rm node}=+1\), and the compatible core is the even-node arc.

\paragraph{The global graph class.}
The type-\(II\) values satisfy \(\zeta^m=-1/(t_b-1)\), at radius
\(R_{\rm II}=(t_b-1)^{-1/m}<1\), between consecutive cusp arguments.
The projected compression cuts are the rays
\(\zeta=r\zeta_j\), \(r\geq1\).
Choose a ring tree in \(R_{\rm II}<|\zeta|<1\), branching inward to
type-\(II\) values and outward to cusps, and enter it from the
binding between two rays.
Resolving its edges into parallel lanes gives disjoint stems;
the inverse isotopy collapses the lanes to the tree.
They avoid the compression cuts away from their cusp endpoints.

In the based raw convention put
\[
 A=\begin{pmatrix}0&-1\\1&-1\end{pmatrix},\quad
 B=\begin{pmatrix}1&1\\-1&0\end{pmatrix},\quad
 C=\begin{pmatrix}1&2\\0&1\end{pmatrix},\quad BAC=I,
\]
with negative Abel--Jacobi cocycle
\(g_A=(0,-1)^t,g_B=0,g_C=e_1\).
For the peripheral words define
\(\mathfrak b_j=\mathfrak a^j\mathfrak b\mathfrak a^{-j}\),
\(\mathfrak c_j=\mathfrak a^j\mathfrak c\mathfrak a^{-j}\).
The stems realize the exact free-word identity
\[
 (\mathfrak b\mathfrak a\mathfrak c)^m
   =(\mathfrak b_0\mathfrak c_1)(\mathfrak b_1\mathfrak c_2)
     \cdots(\mathfrak b_{m-1}\mathfrak c_m)\mathfrak a^m .
\]
The last factor fills smoothly, since \(A^m=I\) and
\((I+A+A^2)g_A=0\).
For \(J=\left(\begin{smallmatrix}0&1\\-1&0\end{smallmatrix}\right)\)
write \(T_v=I+vv^tJ\).
Then \(B=T_{e_2}T_{e_1}\), \(C=T_{e_1}T_{e_1}\).
The cover transport and inverse are
\[
 P_j(z)=A^jz+t_j,\quad P_j^{-1}(z)=A^{-j}(z-t_j),\quad
 t_j=\sum_{l=0}^{j-1}A^lg_A .
\]
A local complex deformation splits a type-\(II\) value into its two
positive nodes and separates the two \(I_2\) node values.
The inverse deformations restore the same source fibres and line.
The four vanishing vectors in block \(j\) are
\(A^je_2,A^je_1,A^{j+1}e_1,A^{j+1}e_1\).

Let \(\ell_i\) be the physical line's relative Chern coefficient
on its oriented thimble.
The null lift \(v_i-\ell_i\delta\) gives lower transport row
\(-\ell_i v_i^tJ\).
The Poincar\'e row \(-c_i^tJT_{v_i}\) therefore gives
\(c_i=\ell_i v_i\).
At type \(II\) the independent vectors and zero affine cocycle give
both coefficients zero.
At \(I_2\) the frames \(V/Z,V/q\) have geometric \(w\)-disc windings
\(-1,0\), hence raw coefficients \(1,0\).
Their difference is the physical degree \(+1\) on \(C_1\),
whose raw chain is \(\Delta_{\rm even}-\Delta_{\rm odd}\);
its two punctures have the opposite boundary orientations of
the intervening complex annulus.
Affine conjugation gives
\[
 A^j(\ell v)+(I-T_{A^jv})t_j
   =\bigl(\ell-\det(A^jv,t_j)\bigr)A^jv .
\]
The values \(t_0=0,t_1=(0,-1),t_2=(1,0),t_3=0\) give
the following twelve entries, repeated \(k\) times:
\[
\begin{split}
 (v_i)={}&(e_2,e_1,e_2,e_2,-e_1-e_2,e_2,-e_1-e_2,-e_1-e_2,\\
         &\hspace{29mm}e_1,-e_1-e_2,e_1,e_1),\\
 (\ell_i)={}&(0,0,1,0,-1,0,0,-1,0,-1,1,0).
\end{split}
\]

Fix the reference-fibre comparison of graph meridian and physical
dual by a positive-normal radial sweep away from all compression
cuts; both have fibre degree zero.
The stems confine each remaining sweep intersection to its
possible node centre.
For a transverse geometric thimble with tangent
\((P_{\rm tan}e^{i\alpha}z,Q_{\rm tan}e^{i\alpha}\bar z)\), the oriented determinant is
\[
 \begin{gathered}
 \begin{aligned}
 2\pi r c_\phi\bigl(&
 \cos^2\alpha(P_{\rm tan}s_{\rm tr}-Q_{\rm tan}c_{\rm tr})^2\\
 &+\sin^2\alpha(P_{\rm tan}s_{\rm tr}+Q_{\rm tan}c_{\rm tr})^2\bigr),
 \end{aligned}\\
 c_{\rm tr}=\cos\phi,\quad s_{\rm tr}=\sin\phi .
 \end{gathered}
\]
A small nonzero stem-angle perturbation makes the bracket
strictly positive.
The even-node contribution is \(+1\) geometrically and \(-1\)
in raw orientation; the odd node is on the other free arc.
Type-\(II\) sweeps miss the compression.
The punctured-sweep boundary is the outer dual thimble minus
the graph thimble minus its signed meridian spheres.
Thus \(a=j_+^*x\) has relative cochain
\[
 a=-\ell+\sum e_{\rm even}
   =(0,0,0,0,1,0,1,1,0,1,0,0)
\]
per block and zero fibre value.
Reference-frame changes add \(V^t\lambda\), zero on closed chains.
Its entries on each \(I_2\) pair agree, and both component values are zero.

For the boundary rim, start with \(y_0=e_2\) and put
\[
 x_i=-\det(v_i,y_i),\qquad
 y_{i+1}=y_i+x_iv_i=T_{v_i}^{-1}y_i .
\]
The strips and these thimbles give \(R_u=(0,-1,-1,-1)\)
repeated \(3k\) times.
The critical-value loop is \(-\beta_0\), whereas
\(H_*\wedge K_*=-u_*\wedge\beta_0=e_2\wedge\beta_0\).
Hence \(T_{H_*K_*}=-R_u\).
The strips miss the zero section, fixing their added fibre
coefficient to zero.  Therefore
\[
 a|_{\partial V_m}=m[dh_*\wedge d\tau_*].
\]
On the capped normal sphere bundle, define \(b\) by
\(e_v=2x-\pi^*b\).
The graph framing gives \(b|_{V_m}=2a\).
The cap \(H_w=D^2_w\times T^2_{H_*K_*}\) has injective
degree-two restriction to its boundary, and consequently
\[
 b|_{H_w}=2m[dh_*\wedge d\tau_*].
\]
\paragraph{The cap sphere and its normal Euler number.}
Choose a point \(\phi_*\) on the compatible free core.
An isotopy of the outer real radius, fixed at both boundaries,
makes its core radius have a strict minimum \(r_*>r_0\) there,
with \(r_\phi=0\).  The compression projection is embedded,
so convex interpolation of the radius remains embedded.
Together with the preceding height isotopy, its velocity extends
by tubular projection and a cutoff to an ambient vector field.
Its flow and reverse flow transport the trace, its normal
bundle and its outward meridian.

The compression disc near its centre is
\[
 D_W:z\longmapsto
 (z\cos\phi_*,\bar z\sin\phi_*,c_W(z),r_W(z)).
\]
Continue it to the incoming boundary using the regular product.
The source map on a regular fibre sends \(w\) with degree one and
\(u_*\) to \(2H_*\) after quotient by \(K_*\).
Consequently its punctured quotient disc has two source lifts.
The reciprocal node charts above extend them at their respective
nodes.  Inserting a full \(H_*\) circuit in a regular band changes
the lift; pairs of circuits change its integer height lift.
The maps
\[
 (w,q)\longmapsto \Phi((w,h(q)+N\chi(q)),q)
\]
and the product inverse give these disc changes, with fixed
incoming parametrized boundary and strictly ordered \(q\) coordinate.
Choose the even lift and its integer height so that its page
disc \(D_P\) has the same quotient disc \(P(z)\) and incoming
height as \(D_W\).  Its graph radius \(f_P\) is nonnegative real
on that collar and zero away from it.

The two discs bound the ball
\[
 {\mathcal B}(z,t)=\bigl(P(z),(1-t)c_P(z)+tc_W(z),
                    (1-t)f_P(z)+tr_W(z)\bigr),\quad 0\leq t\leq1,
\]
with the incoming boundary intervals collapsed.
Away from the centre the inverse reads \(q\) and \(w\), and then the
strictly increasing real radius.  At the centre it reads
\(tr_*\), which determines \(t\).
The two node-disc tangents and the real-radius tangent are
independent, so this is an embedded ball after product-collar
rounding.  The height interval lies in the chosen free core arc.

We check all graph sheets in this ball's collar.
If \(\eta=\arg(q-\epsilon)/(2\pi)-\alpha_0/(2\pi)\), their
heights, modulo one, are
\[
 \begin{gathered}
 \beta_j=(\eta+j)/m,\\
 \tau_{j,\varepsilon_{\rm sh}}=(h_*-\beta_j)/2+\varepsilon_{\rm sh}/2,\\
 0\leq j<m,\quad \varepsilon_{\rm sh}=0,1 .
 \end{gathered}
\]
Equality of two heights implies
\(-(j-j')+m(\varepsilon_{\rm sh}-\varepsilon_{\rm sh}')=0\pmod {2m}\), hence \(j=j'\), \(\varepsilon_{\rm sh}=\varepsilon_{\rm sh}'\).
Exactly one sheet has parameter angle zero at the selected
incoming circle.  It is \(D_P\), at the ball's smaller radius.
The outer trace on this ray has angle zero and the unique
compression quotient point, at its larger radius.
Off the collar every page graph radius is zero.
At the core the ball's interior radius is less than the minimum
core radius \(r_*\).  Thus the ball's interior misses the trace.
Its inward normal is nonzero modulo the trace tangent bundle:
on the graph it has positive real-radius component, on the
compression it has negative real-radius component, and at the
centre \(r_\phi=0\) retains this property.
An inward push of
\(S=D_W-D_P\), with the complex orientation on \(D_W\),
bounds an inner ball in the complement.  Therefore \(x(S_+)=0\).

Project the positive imaginary parameter vector to the normal
two-plane modulo that inward normal.  On \(D_P\) it is nonzero:
on the nonzero graph collar the ball displacement is purely in
real radius, and elsewhere \(df_P=0\).
On \(D_W\) away from the centre the projection of the compression
to its quotient node or regular base has rank three.
A relation expressing this imaginary vector in the trace and ball
directions would therefore leave a nonzero pure-height vector
after its angular term is removed; the ball direction's nonzero
real-radius component cannot cancel it.
The sole zero is the centre.

Put \(z=s+it\), \(c_{{\rm tr},*}=\cos\phi_*\), \(s_{{\rm tr},*}=\sin\phi_*\), and
\[
 N_1=(s_{{\rm tr},*},0,-c_{{\rm tr},*},0),\qquad N_2=(0,s_{{\rm tr},*},0,c_{{\rm tr},*}).
\]
The quotient component of \(W_\phi\) has first jet
\(-sN_1-tN_2\), and its height component is \(c_\phi\).
Writing the angular tangent as
\(K_c\partial_c+2\pi r_*\partial_{\Im z_0}\), the section's
first jet in the normal quotient is
\[
 -\frac{K_c}{2\pi r_*c_\phi}(sN_1+tN_2).
\]
Its coordinate Jacobian is positive.
With the ambient and trace orientations fixed above, the
normal-first frame
\[
 (-\partial_{\Re z_0},N_1,N_2,-W_s,W_t,W_\phi,X_\theta)
\]
has determinant
\(-2\pi r_*c_\phi(c_{{\rm tr},*}^2+s_{{\rm tr},*}^2)^2\).
Terms from the other height and radius derivatives cancel in
this determinant; \(r_\phi=0\) removes their first-order
contribution to the quotient section.
Its Euler index is consequently
\(-\operatorname{sign}(c_\phi)=-1\).
Since \(e_v=2x-\pi^*b\) and \(x(S_+)=0\), we obtain \(b(S)=1\).

This evaluates the ordered cap disc without a path ambiguity.
The selected cusp has the unramified chart
\(\zeta=(1-t_bq)^{-1/m}\), with inverse
\(q=(1-\zeta^{-m})/t_b\).
All inserted \(H_*\) circuits stay in its source \(I_2\)
neighbourhood \(U_0\) and its regular product collar.
The two component normal disc bundles have Euler numbers \(-2\):
a regular fibre is \(C_0+C_1\), and its intersection with either
component is zero, while their two complex nodes contribute two.
The plumbing overlaps are two disjoint four-balls.
Mayer--Vietoris therefore gives the integral isomorphism
\[
 \mathbb Z[C_0]\oplus\mathbb Z[C_1]\ \xrightarrow{\ \cong\ }\
 H_2(U_0;\mathbb Z).
\]
The collar collapses onto its inner endpoint and adds no generator.
Join the ordinary even thimble \(\Delta_{\rm ev}\) to the same
parametrized incoming \(w\) circle by its marked product annulus.
Then \(D_P-\Delta_{\rm ev}\) is a closed chain in \(U_0\);
its \(b\) pairing is zero, since both component pairings are zero.

The compression disc and the product cap meridian \(D_c\),
joined to that same boundary circle, lie in the angle-zero
slice of the compression and its outgoing cap.
This slice is \(D^2_w\times S^1_{H_*}\), with its explicit
product inverse; its second homology is zero by contraction
of the disc factor.  Thus \(D_W-D_c\) bounds there.
The based shear multiplies \(z_w\) by
\(e^{-2\pi iph_*+6\pi ipg_*}\) on this slice, and its reciprocal
inverse preserves the slice.
It follows that
\[
 b(D_c-\Delta_{\rm ev})=b(S)=1 .
\]
In raw coordinates \(e_1=-w_Z\), both disc orientations reverse.
The raw cap-minus-even-thimble value is \(-1\).
The last even graph coefficient is zero, so the relative cap
coefficient of \(b\) is exactly \(-1\).

\paragraph{The integral square and return.}\leavevmode
\hypertarget{tag-01GY}{}\label{characteristic-calculation-and-the-scope-of-the-result}
Put \(N=12k\), and let
\(V:\mathbb Z^{N+1}\to\mathbb Z^2\) have the displayed nodal
columns and final column \(e_1\).
The cap is \(D^2_w\times T^2_{H_*K_*}\): it adds one
two-handle of product framing zero, two three-handles and
one four-handle.
Along an attaching \(w\) circle the based shear multiplier is
constant, so it adds no framing winding.
The first columns \(e_2,e_1\) split \(V\) integrally, and
\[
 e_j-v_{j,2}e_0-v_{j,1}e_1,\qquad 2\leq j\leq N,
\]
are an integral basis of \(\ker V\).
Before the three-handles, \(H_2\) is
\(\mathbb ZF\oplus\ker V\), where \(F\) is the reference fibre.
The product boundary tori are \(wH_*\) and
\(wK_*=wu_*-2wH_*\).
Their boundary span is generated integrally by \(F\) and the rim
\[
 R=(-1,1,0,0)\quad\hbox{repeated \(3k\) times},
 \qquad R_{\rm cap}=0 .
\]
The change from these two tori to \(R,F\) is unimodular.
In the displayed kernel basis \(R\) has a coefficient \(-1\),
so the quotient is free.
Both \(b\) pairings vanish, and
\[
 H_2(\bar X_m;\mathbb Z)
   =(\mathbb ZF\oplus\ker V)/\langle F,R\rangle .
\]

Define the integral ordered matrix \(U\) by
\[
 U_{ii}=-1\ (i<N),\quad U_{NN}=0,\quad
 U_{ij}=\det(v_i,v_j)\ (i<j),\quad U_{ij}=0\ (i>j).
\]
The complex nodal framing is \(-1\)
\cite[Theorem~1.7 and the surface case, pp.~93--96]{Kas1980}.
For the off-diagonal sign, carry the \(i\)th boundary cycle along
the annulus oriented \((v_i,\partial_\theta)\).
Its crossing with an earlier inward core
\((v_j,-\partial_r)\), in orientation
\((e_1,e_2,\partial_r,\partial_\theta)\), has determinant
\(-\det(v_i,v_j)=\det(v_j,v_i)\).
This is the ordered relative pairing \(U^t\).
Since \(U-U^t=V^tJV\), the pairing on closed chains is the
symmetric integral form \(Q(x,y)=x^tUy\).
Its radical before the three-handle quotient is generated by
the fibre and \(R\), so it descends to the intersection form of
the closed trace.

The \(b\) cochain is \(2a\) on the old handles and \(-1\) on the cap.
For block \(j=0,\ldots,k-1\), take
\[
 \begin{split}
 x_{{\rm PD},j}={}&(-6j,6j-1,-1,-1,-6j-2,6j+1,-1,-1,\\
       &\hspace{24mm}-6j-4,6j+3,-1,-1),\\
 x_{{\rm PD},{\rm cap}}=&6k .
 \end{split}
\]
A block has period sum \(-6e_1\) and \(b\) pairing \(-2\).
Thus \(Vx_{\rm PD}=0\).
If \(L_i=\sum_{h<i}x_{{\rm PD},h}v_h\) and \(\lambda_0=(1,0)\), substitution
in the twelve displayed entries gives
\[
 -x_{{\rm PD},i}-\det(v_i,L_i)=b_i+\lambda_0(v_i)\quad(i<N).
\]
At the cap \(L_N=-6ke_1\); both sides of its corresponding
equation are zero, because \(-1+\lambda_0(e_1)=0\).
Equivalently \(U^tx_{\rm PD}=b+V^t\lambda_0\).
For every \(y\in\ker V\) this gives \(Q(x_{\rm PD},y)=b(y)\).
Both sides vanish on the killed fibre and rim.
Hence the image of \(x_{\rm PD}\) is the integral Poincar\'e dual of \(b\), and
\[
 \langle b^2,[\bar X_m]\rangle=b(x_{\rm PD})=-2k-6k=-8k .
\]

For the tangent Pontryagin number compare the two marked caps.
The ordinary base meridian is \(\beta_0=H_*\), and
\(u_*=2H_*+K_*\).
The point-coordinate maps
\((\xi,\beta)\mapsto(h_*,\tau_*)=(2\xi+\beta,\xi)\) and
\((h_*,\tau_*)\mapsto(\tau_*,h_*-2\tau_*)\)
identify their surviving torus coordinates integrally.
In these marked coordinates the caps are
\[
 H_w=D^2_w\times S^1_{u_*}\times S^1_{\beta_0},\qquad
 H_{\beta_0}=S^1_w\times S^1_{u_*}\times D^2_{\beta_0}.
\]
Their common-boundary gluing is the two solid-torus gluing
of \(S^3\), times \(S^1_{u_*}\); the displayed maps specify
the boundary meridians and longitudes, including their inverse.
Each cap has zero signature.
Novikov additivity along this entire closed three-torus
\cite[Proposition~(7.1), pp.~587--589]{AtiyahSinger1968IIIE}
therefore gives
\(\sigma(\bar X_m)=\sigma(S'_m)=-8k\).
The four-dimensional signature formula yields
\[
 \langle p_1(T\bar X_m),[\bar X_m]\rangle=-24k .
\]

Let \(\pi:B_\nu=S(\bar\nu)\to\bar X_m\) be the capped oriented
normal sphere bundle.
The graph quotient frame above, its outgoing product cap, and
the same meridian-one class give
\[
 e_v=2x-\pi^*b,\quad \pi_*x=1,\quad
 e_v^2=\pi^*p_1(\bar\nu),\quad
 p_1(TB_\nu)=\pi^*\bigl(p_1(T\bar X_m)+p_1(\bar\nu)\bigr).
\]
Fibre integration of the first squared identity and then its
product with \(x\) gives
\[
 \pi_*x^2=b,\qquad
 \pi_*x^3=\frac{p_1(\bar\nu)+3b^2}{4}.
\]
The meridian index is therefore
\[
 A_U=\frac{\langle4x^3-p_1(TB_\nu)x,[B_\nu]\rangle}{24}
     =\frac{3(-8k)-(-24k)}{24}=0,\qquad A_M=-A_U=0 .
\]
These equations use the nonzero cap class computed above.

For \(k=1\), the full parameter clutch is
\(\theta K_*+3\theta(\theta-1)\delta\),
by \SeriesRef{night:P3F.27} and the based maps already given.
Its integer fan translations preserve the positive triangles;
the bounded torus origin multiplies the chart coordinates by
nonzero units and extends with reciprocal inverse.
The page graph, compression, fixed normal frame and outgoing
cap are thus the construction of
\SeriesRef{iv:correction:P3K.6} and \SeriesRef{night:P3K.7}.
The latter's deleted-tube bordism has boundary
\(-P_f\sqcup-B_\nu\).
Its physical closing is \((z,1)\sim(f(z),0)\), while the
inverse-closing convention is \((z,0)\sim(f(z),1)\).
Keeping these two reversals, the Hopf comparison
\SeriesRef{night:P3K.68} gives \([f]=A_M=0\) in \(\Theta_7\).
Cerf's identification of oriented sphere diffeomorphism
components with \(\Theta_7\)
\cite[Introduction, Theorem~0 and Corollaries~1--2, pp.~7--8]{Cerf1970}
supplies a smooth isotopy \(f_t\) from the identity to \(f\).
Reparametrize it to be constant near its endpoints.
Then
\[
 [z,t]\longmapsto(f_t(z),e^{2\pi it})
\]
is a diffeomorphism \(P_f\to S^6\times S^1\);
its inverse is \((y,e^{2\pi it})\mapsto[f_t^{-1}(y),t]\).
The endpoint identification and their constant-end collars
verify that both maps are well defined and smooth.

Finally, the negative comparison in
\SeriesRef{night:P3K.65} has exactly the source cubic and section,
the line \(\mathcal O(G-O)\), divisor \(C_0\), the same type-\(II\)
quotient and finite cap \((m,n,d)=(3,1,0)\).
Its two fan contractions are the conifold ruling maps displayed
at the beginning of this proof.
Under its oriented Cox ruling map, the native free pair for
the positive \(r\) loop is \((-3,3)\), giving its stated class \(-3\).
Its positive-fan calculation uses the present page calculation;
only its negative-fan comparison is invoked here.
The fan construction of
\cite[Proposition~2.8, pp.~6--7]{E26} permits any periodic
unimodular triangulation, so both displayed diagonals meet
that hypothesis.
For the whole-fibre hypothesis in
\cite[Remark~8.5, p.~24]{E26}, the arrow line here is
\(\operatorname{Hom}(T^*M,M)=\mathcal O(C_0)\).
Its degree on \(C_0\) is \(-2\), whereas every base-pulled line
has degree zero on \(C_0\).
Restriction of an isomorphism would preserve that degree;
hence this arrow line has no base-pulled whole-fibre
presentation.  The construction above supplies its component
divisor directly.
\end{proof}

\subsection{Return for all sphere caps}
\begin{theorem}[All finite caps of the positive-diagonal OS49 family]\hypertarget{tag-02S9}{}\label{night:P3F.32}\label{stmt:P3F.32}
{\normalfont\scriptsize tag 02S9}\par
Let \(m=3k\), \(k>0\), \(n\in\{1,-1\}\), and \(d\in\mathbb Z\).
The positive-diagonal OS49 filling with cap vector
\(\ell=n(e_\gamma-e_u)+d\delta\) is a family of standard smooth
six-spheres over a sufficiently small punctured parameter
disc \(\lambda=r^m\).
Its positive r-loop has return integer zero, oriented return
class zero in \(\Theta_7\), and mapping torus diffeomorphic to
\(S^6\times S^1\).
In the based cusp marking its angular clutch is
\[
 n\theta K_*+nm\theta(\theta-1)\delta,\qquad 0\leq\theta\leq1 .
\]
On the cusp cover \(q=Q_0^{2m}\), the physical line gauge of
the positive full root loop has winding \(2mn\).
With the native C2 comparison of \SeriesRef{night:P3K.65}, the
negative-diagonal filling has positive-loop class \(-3k\) for
\(\lambda=r^{3k}\), and class \(-k\) for
\(\lambda=w^k\), \(w=r^3\), for both numerator signs and every
fixed integral delta character in this cap family.
\end{theorem}
\begin{proof}
Use the positive fan and source line of
\SeriesRef{night:P3F.31}.
The common-refinement contractions there, applied to
\SeriesRef{night:P3F.29}, give the integral sphere homology
and simple connectivity for each \(m,n,d\).
Their exceptional ruling has primitive degree \(-1\), so the
same integral blowup splitting applies for every cap.
Thus these fibres are standard smooth six-spheres.
The unit-preserving cusp quotient and normal trivialization
of \SeriesRef{night:P3F.31} give their orbit five-spheres.

\hypertarget{tag-034P}{}\label{the-actual-closing-map-and-its-origins}
\paragraph{The closing map and its origins.}
Put \(v=\log r/(2\pi i)\), \(s=\log q/(2\pi i)\).
Retain the based period lift \(u_*=u_h-pw_Z-q_\delta\delta\)
of \SeriesRef{night:P3F.31}, and write
\[
 u_*=(2s,s)^t+U_*(q),\quad
 \gamma=(-3s,-s+mv)^t+C(q),\quad D_*=3U_*+2C .
\]
Then
\[
 H_*=-u_*-\gamma,\qquad
 K_*=3u_*+2\gamma=D_*+(s+2mv)\delta .
\]
The finite cap covector has value n on \(\ell\).
The logarithmic finite attachment
\[
 z\longmapsto z+\frac{\log\sigma}{2\pi i}\Pi(\sigma)\ell,
 \qquad t=\sigma^m,
\]
with inverse subtraction, changes by
\(n\log t\,\delta/(2\pi i)\) when \(v\) increases by one.
The fixed displacement \(d\delta/m\) has zero parameter difference.
Thus multiplication of the ordinary physical line coordinate by
\(t^n\), with inverse \(t^{-n}\), closes the family.
It is a unit at the type-II value and is \(q^{-n}\) at the cusp.
Its lattice arrow is integral vertical translation by n,
which preserves the positive fan.

On \(q=Q_0^{2m}\), the two origins are
\[
 \begin{split}
 \sigma_{\rm ext}(s,v)&=nvD_*+nmv(v-1)\delta,\\
 o_{\rm ord}(s,v)&=-\frac{nsD_*}{2m}
              -\frac{ns^2}{4m}\delta-nsv\delta .
 \end{split}
\]
Their translations and opposite translations are the chart
maps and inverses of \SeriesRef{night:P3F.27}.
The first is bounded in the cusp charts.
The second has parameter difference \(-ns\delta\), the
logarithm of the same cusp closing \(q^{-n}\).
Their difference is
\[
 F_*(s,v)=\frac{nD_*(s+2mv)}{2m}
       +\left(\frac{n(s+2mv)^2}{4m}-nmv\right)\delta .
\]
Direct subtraction gives
\[
 \begin{split}
 F_*(s,v+\theta)-F_*(s,v)
     &=n\theta K_*+nm\theta(\theta-1)\delta,\\
 F_*(s+2m,v)-F_*(s,v)&=nK_*+nm\delta .
 \end{split}
\]
The second arrow is an integral period.
The forward based line gauge of the full root loop is
\(e^{2\pi i(2mn)\theta}\), with reciprocal inverse, as
computed from the same ordinary and extending charts in
\SeriesRef{night:P3F.27}.
Its oriented winding is \(2mn\).
The quadratic angular gauge has zero endpoint values,
and the positive parameter loop makes n primitive \(K_*\) turns.

These formulas retain both invariant branch integers.
In the geometric coordinates the affine vector is
\((n,n,-np,d-nq_\delta)^t\).
The new coordinates are \(L^{-1}\) times the old, where L is
the matrix of \SeriesRef{night:P3F.31}, so
\[
 L^{-1}(n,n,-np,d-nq_\delta)^t=(n,n,0,d)^t,
 \qquad A_{\rm new}=L^{-1}A_{\rm old}L .
\]
On the punctured parameter disc the accompanying chart map is
\((Z,V)\mapsto(r^{3np}Z,r^{3nq_\delta}V)\), with reciprocal
inverse.  It is nonzero also at \(q=0\) for each \(r\ne0\).
Its q-root difference is zero, and it carries the full angular
path and its endpoint arrows to the displayed based formulas.

\paragraph{The finite cap and page line.}\leavevmode
\hypertarget{tag-034Z}{}\label{the-page-including-the-finite-cap-line}
First take \(n=1\).
In the finite origin of \cref{night:P3F.26}, the cap generator is
\[
 g_\ell(z,\sigma)
   =(A_{\rm cap}z+\ell/m,e^{-2\pi i/m}\sigma),
 \qquad A_{\rm cap}=T_A^{-t}.
\]
It fixes \(\ell\), its linear part has order three, and its
jth power has suspension displacement \(j/m\).
For \(0<j<m\), this is nonintegral.
Thus the affine cap is free, including its central fibre.
A physical delta rotation cannot compensate that displacement,
so the physical circle is free on the cap.

For \(d=0\), let \(a_{\rm cap}\) be the coefficient along the invariant
\(\ell_0=e_\gamma-e_u\) direction.
The quotient solid-torus coordinates and their inverse are
\[
 S=\sigma e^{2\pi i a_{\rm cap}},\quad B=e^{2\pi i m a_{\rm cap}},\quad
 t=S^m/B;\qquad
 e^{2\pi i a_{\rm cap}}=B^{1/m},\quad \sigma=S e^{-2\pi i a_{\rm cap}}.
\]
The m choices of root in the inverse are exactly the deck orbit.
Translation of \(a_{\rm cap}\) has weights \((1,m)\) on \((S,B)\).
With the ordinary exterior solid torus and its unchanged
logarithmic collar, these give the weighted Hopf
three-sphere charts.
The regular-binding page is the cyclic cover
\[
 \zeta^m=\frac{t}{t-t_b},\qquad
 t=\frac{t_b\zeta^m}{\zeta^m-1}.
\]
Its compact base is the \(\zeta\)-sphere.
The page equation \(m a_{\rm cap}=0\pmod1\) selects m levels permuted
transitively by the cap generator; the unique level-zero
representative gives the good-reduction page chart and inverse.
The same selection applies to unit physical line vectors.

On the finite root the explicit maps are
\[
 t=\sigma^{3k},\qquad x=\sigma^{4k}X,\qquad y=\sigma^{6k}Y,
 \qquad X=\sigma^{-4k}x,\quad Y=\sigma^{-6k}y
\]
on the punctured overlap, with the smooth good-reduction
extension of \SeriesRef{night:P3F.28}.
The section line extends there as \(\mathcal O(G'_m-O'_m)\).
In the logarithmic attachment the radial term is
\(-\log|\sigma|/(2\pi)\) times the complex structure applied
to \(\Pi\ell_0\).  A root-equivariant collar cutoff extends
this smooth translation to zero on a smaller disc; its
opposite is the inverse.
Its angular term is the same \(a_{\rm cap}\) translation as in the level
selection.  This retains the entire attachment and its units.

At the cusp choose m norm shells satisfying
\(\kappa R_m<R_1\).
The full line maps and height defects in
\SeriesRef{night:P3F.31} prove their embeddings and prevent
re-entry after a source zero.
The positive radial suspension maps there join them, in their
cyclic order, to the m outer page levels.
The type-II pieces use the original source-line quotient.
The smooth binding comparison there uses \(e=m\).
Thus the page is
\(V_m=S'_m\setminus(E_{t_b}\times\operatorname{int}D^2)\),
with all source components, nodes and line frames retained.

The compact model is pulled back from \(S'_3\) by
\(\zeta_3=\zeta^k\), branched only at two smooth fibres.
Its pulled-back sections are disjoint and have square \(-k\);
it has m type-II and m \(I_2\) fibres.  Consequently
\[
 \chi(\mathcal O_{S'_m})=k,\quad
 (\mathcal O(G'_m-O'_m))^2=-2k,\quad
 \sigma(S'_m)=-8k,\quad p_1(S'_m)=-24k .
\]

For the twist d, the actual line overlap is
\[
 V_{\rm out}=\sigma^dV_{\rm in},\qquad
 V_{\rm in}=\sigma^{-d}V_{\rm out}.
\]
Equivalently its frames satisfy \(e_{\rm in}=\sigma^d e_{\rm out}\).
Hence the page line is
\[
 L_d=\bigl(\mathcal O(G'_m-O'_m)\otimes
                    \mathcal O(-dF_0)\bigr)|_{V_m},
\]
where \(F_0\) is the smooth finite-end fibre.
Both \(F_0^2\) and \((G'_m-O'_m)\cdot F_0\) vanish.
Thus the stated compact extension, including every d, has
square \(-2k\).

\paragraph{Transport of the graph and the return.}\leavevmode
\hypertarget{tag-034Q}{}\label{the-actual-graph-line-and-return}
For \(d=0,n=1\), the m-sheet construction just given is the
one used in \SeriesRef{night:P3F.31}.
That proof calculates its graph meridian from the ordered
thimbles and its cap entry from the local normal sphere.
In its notation,
\[
 a=(0,0,0,0,1,0,1,1,0,1,0,0)
 \quad\hbox{in each twelve-handle block},\qquad
 b_{\rm old}=2a,\quad b_{\rm cap}=-1.
\]
The boundary is \(b=2m[dh_*\wedge d\tau_*]\).
Its integral dual cycle in that proof gives
\[
 \langle b^2,[\bar X_m]\rangle=-8k,\qquad
 \langle p_1(T\bar X_m),[\bar X_m]\rangle=-24k .
\]

We transport the full normal problem to every d.
After quotient by the physical circle, the cap displacement
\(d\delta/m\) disappears from the displayed affine action.
The identity on the remaining cap coordinates, including
\(\sigma=0\), maps its orbit cap to the zero-twist cap.
On an attaching annulus \(a_0<|\sigma|<b_0\), \(a_0>0\),
write \(c=\log|V|\) and choose a base cutoff \(\chi=1\)
on the entire attachment, zero at its outer edge.
The exterior map and inverse are
\[
 (p,c)\longmapsto(p,c-d\chi(p)\log|\sigma(p)|),
 \qquad
 (p,c)\longmapsto(p,c+d\chi(p)\log|\sigma(p)|).
\]
Since \(c_{\rm out}=c_{\rm in}+d\log|\sigma|\), these maps
agree with the identity cap map on their common overlap.
They glue to a smooth orbit diffeomorphism \(F_d:Y_d\to Y_0\).
The derivative is triangular with diagonal entries one
and sends \(-\partial_c\) to itself.
The function is a single-valued base function
\(\log|\sigma|=\log|t|/m\) on its support.
Thus \(F_d\) commutes with the suspension and parameter
closing; it is the identity near the cusp, binding,
compression and fixed sphere.

Choose the d-page norm by transporting the zero-page norm.
The resulting graph maps satisfy
\[
 (F_d\times{\rm id}_{D^2})(X_d)=X_0
\]
with the same positive graph normal, incoming inverse angle
and outgoing cap.
The map extends over the prescribed \(D^3\times S^1\) cap
by the identity.
For any other admissible norm the positive ratio is \(e^{f_{\rm norm}}\);
the interpolation \(e^{s f_{\rm norm}}\) and its reciprocal give the
source-preserving radial isotopy.
On the compact source, a common sufficiently small parameter
radius keeps every contraction and shell-separation inequality
strict during this isotopy.
The graph quotient maps and their inverses then give the
same marked normal transport.

The induced complement map preserves each oriented meridian
sphere.  Uniqueness of the meridian-one class gives its
pullback as \(x_d\), and the oriented normal-bundle map
carries the vertical Euler class as well.
Pulling back \(e_v=2x-\pi^*b\) therefore carries \(b_0\) to
\(b_d\) on the closed trace.
Here \(\pi^*\) is injective, since
\(\pi_*(x\,\pi^*u)=u\).
Thus the closed b square and tangent Pontryagin number
are the displayed values for every d.

The same conclusion has an explicit page-line boundary map.
Near \(F_0\), write
\(\sigma=u(\zeta)\zeta\), with \(u\) the chosen nonzero
root of \(t_b/(\zeta^m-1)\).
The isomorphism \(L_d\to L_0\) is
\[
 V_{0,\rm out}=\zeta^{-d}V_{d,\rm out},\qquad
 V_{0,\rm in}=u(\zeta)^dV_{d,\rm in}.
\]
Both factors are nonzero on their respective charts and have
reciprocal inverses; on the overlap
\(\zeta^{-d}\sigma^d=u^d\).
At the binding put \(s_0=1/\zeta\).
The exact relation \(q-\epsilon=-\epsilon s_0^m\) identifies
\(\arg s_0/(2\pi)\) with
\(\beta=h_*-2\tau_*\), up to the fixed binding phase.
The d-to-zero boundary character is therefore
\[
 e^{2\pi id(h_*-2\tau_*)}.
\]
This same formula, independent of \(z_w\), extends across
\(D^2_w\times T^2_{H_*K_*}\), with reciprocal inverse.
The fixed constant unit and positive modulus extend by
their product collar and a positive cutoff.
Thus its Mayer--Vietoris connecting class is zero because
the character is the restriction of a cap character.
This retains the nonzero value \(b(H_*K_*)=2m\).

It remains to transport to \(n=-1\).
Use the exterior origin of \SeriesRef{night:P3F.26}.
On its regular real four-torus the map is \(z\mapsto-z\).
On a regular or type-II source fibre put \(f_{\rm inv}=x-t^2\).
The source cubic gives
\[
 \operatorname{div}(f_{\rm inv})=G+(-G)-2O .
\]
Indeed substitution of \(x=t^2\) gives \(y^2=t^4\);
the two zeros are \(G,-G\), and the pole at O has order two.
On the type-II fibre these sections are smooth and \(f_{\rm inv}\)
does not vanish identically on the fibre, so no fibre
component is added.
In rational frames \(s_D\) of \(M=\mathcal O(G-O)\),
the line-pairing and reciprocal map are
\[
 s_D(p)\otimes s_D(-p)\longmapsto f_{\rm inv}(p),\qquad
 a s_D(p)\longmapsto(af_{\rm inv}(p))^{-1}s_D(-p).
\]
The divisor equation proves extension through
\(G,-G,O\) in their line frames.
Since \(f_{\rm inv}(-p)=f_{\rm inv}(p)\), the map is its own inverse.
The zero-section rigidification fixes its scalar.
It is the semiabelian inversion in the chosen exterior origin
and conjugates every complete suspension translation to its
inverse, including the type-II quotient.

At the \(I_2\) cusp we use the full toric extension.
Represent that exterior origin by
\((f_1(q),f_2(q))\), with \(f_i=q^{v_i}u_i(q)\),
\(v_i\in\mathbb Z\) and \(u_i\) nonvanishing.
The map and its inverse are
\[
 J_f(q,Z,V)=(q,f_1(q)^2/Z,f_2(q)^2/V).
\]
The lattice arrow
\((i,j,1)\mapsto(2v_1-i,2v_2-j,1)\) interchanges
the two positive triangle types.
For a cone ray matrix C and this lattice matrix T,
the target matrix is TC and
\((TC)^{-1}T=C^{-1}\).
Thus each target monomial coordinate equals the corresponding
source coordinate times an integral power of the nonzero
units \(u_1^2,u_2^2\).
The reciprocal units give its inverse on the entire chart.
This proves extension over all divisors, nodes and triple points.
For the full multiplier \(D_a(Z,V)=(a_1Z,a_2V)\),
direct substitution gives \(J_fD_aJ_f=D_a^{-1}\),
retaining every unit and \(\lambda\).
These maps agree with the preceding inversion on the common
group-torus locus and hence on their chart overlaps.

In the finite origin, \(-I\) commutes with \(A_{\rm cap}\), and
\[
 (-I)g_\ell=g_{-\ell}(-I).
\]
The root coordinate is unchanged.
The logarithmic attachment is linear in \(\ell\), so
\((-I)h_\ell=h_{-\ell}(-I)\), also for their inverse arrows.
This identity carries the complete based vector and both
invariant root characters.
Moreover \(J_f(t^nV)=t^{-n}J_f(V)\).
The maps therefore give a family diffeomorphism
\[
 X_{m,1,d}\longrightarrow X_{m,-1,-d}
\]
over the identity of the punctured parameter disc, intertwining
the two positive-loop endpoint maps.
On the regular real four-torus its derivative is \(-I_4\),
of determinant \(+1\), and it fixes the complex base.
The smooth six-dimensional fibre map is orientation preserving.
It reverses the physical circle and therefore the orbit
orientation in the convention
\(\Omega_X=d\delta\wedge\Omega_Y\).
It carries the meridian line class to its negative and
reverses the fundamental class of the capped normal sphere
bundle \(B_\nu=S(\bar\nu)\).
Both signs reverse in the integral of
\(4x^3-p_1(TB_\nu)x\), so its return index is preserved.
The return diffeomorphisms are conjugate by the displayed
orientation-preserving fibre map.
Every negative-numerator twist occurs as \(-d\), proving
the whole stated scope.

For \(n=1\) the sphere-bundle identity and the physical
duality from \SeriesRef{night:P3F.31} give
\[
 A_U=\frac{3(-8k)-(-24k)}{24}=0,\qquad A_M=-A_U=0 .
\]
The endpoint-zero angular gauge extends along the real path
specified above; it changes neither closing arrow.
The cap, normal frame, spin extension and both closing
orientation reversals are precisely those in
\SeriesRef{night:P3K.7}.
The Hopf comparison \SeriesRef{night:P3K.68} gives oriented
return class zero, and the negative-numerator family map
gives the same conclusion there.
An isotopy \(f_t\) from the identity to the return \(f\),
constant near its endpoints, exists by
\cite[Introduction, Theorem~0 and Corollaries~1--2, pp.~7--8]{Cerf1970}.
For the physical convention \((z,1)\sim(f(z),0)\), the
diffeomorphism and inverse are
\[
 [z,t]\longmapsto(f_t(z),e^{2\pi it}),\qquad
 (y,e^{2\pi it})\longmapsto[f_t^{-1}(y),t].
\]
Thus every stated mapping torus is \(S^6\times S^1\).

For the negative-diagonal comparison,
\SeriesRef{night:P3K.65} has the same source, line, component
divisor, type-II quotient and cap vector
\(n(e_\gamma-e_u)+d\delta\).
Its oriented native C2 comparison has free pair
\((-3k,3k)\) for \(\lambda=r^{3k}\), hence class \(-3k\).
For \(\lambda=w^k\), \(w=r^3\), its pair is \((-k,k)\)
and class \(-k\).
Those statements apply to both numerator signs and every
fixed integral delta character in their stated cap scope.
The comparison uses the conifold ruling maps of
\SeriesRef{night:P3F.31}; the positive-fan computation is
the graph and cap calculation above.
\end{proof}

\subsection{The complete free affine sector}
\begin{theorem}[The free OS49 affine caps]\hypertarget{tag-02SA}{}\label{night:P3F.33}\label{stmt:P3F.33}
{\normalfont\scriptsize tag 02SA}\par
Let \(m=3k\), \(k>0\), and let \(n,d\in\mathbb Z\).
The explicit OS49 affine cap with
\(\ell=n(e_\gamma-e_u)+d\delta\) is free exactly when
\begin{equation}
 \gcd(k,n,d)=1,\qquad
 3\nmid n\ \text{or}\ 3\mid k.
 \hypertarget{tag-Q30F}{}\label{os49:eq:free-affine-cap-criterion}
\end{equation}
Every such completed filling has fundamental group
\(\mathbb Z/|n|\), with \(\mathbb Z/0=\mathbb Z\).
Its sphere sector is exactly \(n=\pm1\); these members are
standard smooth six-spheres.
\end{theorem}
\begin{proof}
Keep the marked root charts and integral order-three matrix
\(A=T_A^{-t}\) of \cref{night:P3F.26}. On the central torus
the generator is
\[
 \alpha(z)=Az+\ell/m,\qquad
 \ell=n(1,-1,0,0)^t+d\delta.
\]
The root-disc coordinate is multiplied by \(e^{-2\pi i/m}\).

\paragraph{All stabilizers on the root chart.}
\hypertarget{tag-0008}{}\label{all-stabilizers-on-the-actual-root-chart}
A nonidentity root rotation has no fixed point away from the
central fibre. For \(j=3r\), one has \(A^j=I\), and
\(\alpha^j\) is translation by \(r\ell/k\). It has a fixed
point precisely when
\[
 k\mid rn,\qquad k\mid rd.
\]
No \(1\le r<k\) satisfies both conditions exactly when
\(\gcd(k,n,d)=1\).

For \(3\nmid j\), the images of \(A^j-I\) and \(A-I\)
agree: \(A^2-I=-A^2(A-I)\), with inverse factor \(-A\).
The primitive quotient \(\Lambda/(A-I)\Lambda\) has coordinates
\[
 \psi,\qquad h=-w+3\delta.
\]
Its coordinate map and inverse are
\[
 [v]\longmapsto(\psi(v),h(v)),\qquad
 (a,b)\longmapsto[a e_\gamma-b e_w].
\]
Indeed the displayed matrix gives
\((A-I)\Lambda=\mathbb Z e_u+\mathbb Z(3e_w+e_\delta)\),
which is exactly the kernel of these two coordinates.
Their values on \(\ell\) are \(n,3d\). The fixed-point
equation is therefore soluble exactly when
\begin{equation}
 m\mid jn,\qquad m\mid3jd.
 \hypertarget{tag-Q30G}{}\label{os49:eq:affine-fixed-point-test}
\end{equation}
To prove sufficiency, choose the integral vector with these two
coordinates using the displayed inverse, and subtract
\(j\ell/m\). The resulting vector \(y\) lies in the real image
of \(B-I\), where \(B=A^j\). On that image
\((B-I)^{-1}y=-(B+2I)y/3\), since \(I+B+B^2=0\) there.
This gives an actual point of the underlying complex torus
solving the affine fixed-point equation.

Put \(g=\gcd(m,n,3d)\). The integers in
\eqref{os49:eq:affine-fixed-point-test} are precisely the multiples
of \(m/g\). There is no such \(j<m\) with \(3\nmid j\)
exactly when \(3\mid m/g\). Under \(\gcd(k,n,d)=1\), every
prime other than three dividing g would divide k,n,d and is
excluded. If \(3\nmid n\), then \(g=1\). If \(3\mid n\)
and \(3\mid k\), the gcd condition gives \(3\nmid d\), so
\(g=3\) and \(3\mid m/g=k\). If \(3\mid n\) and
\(3\nmid k\), one has \(g=3\) and \(3\nmid m/g=k\).
These cases prove exactly
\eqref{os49:eq:free-affine-cap-criterion} for every group element.
The free triple \((m,n,d)=(9,3,1)\) exhibits the case
\(3\mid n\); its delta displacement supplies the required
nonintegral translation for the remaining stabilizers.

\paragraph{The fundamental-group presentation.}
\hypertarget{tag-034U}{}\label{the-full-fundamental-group-presentation}
The cap universal cover has deck generators \(\Lambda\) and
the lifted root rotation \(\alpha\), with relations
\[
 [\Lambda,\Lambda]=1,\qquad
 \alpha v\alpha^{-1}=Av,\qquad \alpha^m=\ell.
\]
Attach it by the punctured torus maps of \cref{night:P3F.26}.
The original type-II and cusp pieces impose its proved primitive
relation lattice: \(u,w,\delta\) and their section meridians
vanish, leaving \(\gamma\). The base relation kills the remaining
meridian \(\alpha\). The power relation becomes \(n\gamma=0\),
and the conjugation relations vanish because \(\psi A=\psi\).
The quotient and inverse maps are
\[
 v\longmapsto\psi(v)\pmod n,\quad\alpha\longmapsto0,
 \qquad 1\pmod n\longmapsto[e_\gamma].
\]
They give \(\pi_1(X)=\mathbb Z/|n|\), also for \(n=0\).
For \(|n|\ne1\) this group excludes a sphere. For \(n=\pm1\),
\cref{night:P3F.29} proves integral sphere homology and simple
connectivity for every k,d; the common blowup maps of
\cref{night:P3F.31} carry this topology to the positive fan.
Both displayed diagonal fillings are therefore standard smooth
six-spheres exactly in the stated sector.
\end{proof}

\Needspace{9\baselineskip}
\subsection{Integral homology of every free cap}
\begin{theorem}[Homology of the free OS49 affine caps]\hypertarget{tag-02SB}{}\label{night:P3F.34}\label{stmt:P3F.34}
{\normalfont\scriptsize tag 02SB}\par
Use the complete free sector of \cref{night:P3F.33}, and put
\(m=3k\), \(g=\gcd(m,n,3d)\in\{1,3\}\).
For \(n\ne0\), the completed threefold has
\begin{equation}
 \begin{aligned}
 H_1(X)&=H_4(X)=\mathbb Z/|n|,\\
 H_2(X)&=H_3(X)=\mathbb Z/|n|\oplus\mathbb Z/g,\\
 H_5(X)&=0,\qquad H_0(X)=H_6(X)=\mathbb Z.
 \end{aligned}
 \hypertarget{tag-Q30H}{}\label{os49:eq:free-cap-nonzero-homology}
\end{equation}
For \(n=0\), the free criterion is \(3\mid k\) and
\(\gcd(k,d)=1\), and the groups are
\begin{equation}
 H_1=H_5=\mathbb Z,\quad
 H_2=H_3=\mathbb Z^2\oplus\mathbb Z/3,\quad
 H_4=\mathbb Z^2,\quad H_0=H_6=\mathbb Z.
 \hypertarget{tag-Q30I}{}\label{os49:eq:free-cap-zero-homology}
\end{equation}
The free triple \((m,n,d)=(9,3,1)\) has
\(H_2=H_3=(\mathbb Z/3)^2\). On the entire coprime sector
\(\gcd(n,3k)=1\), one has \(g=1\), giving the homology of
\cref{maps:src:D.2.11} through the exact family map of
\cref{os49:rem:type-c-marking}.
\end{theorem}
\Needspace{8\baselineskip}
\begin{proof}
\leavevmode
\paragraph{The fixed exterior and the oriented boundary circle.}
\hypertarget{tag-034T}{}\label{the-fixed-exterior-and-the-actual-boundary-circle}
Remove a closed finite-cap neighbourhood and fix the same scalar
\(\lambda\) and ordinary, type-II and cusp charts. The exterior
\(V\) and its marked boundary \(E\) are then fixed for every
\(m,n,d\). Begin with the positive exterior meridian \(\mu\)
of \cref{night:P3F.26,night:P3F.29}. For the affine roof
\[
 T_{A,a}=(T^4\times[0,1])/((x,1)\sim(Ax+a,0)),
\]
the linear roof map and inverse are
\[
 U([x,t])=[x+tA^{-1}a,t],\qquad
 U^{-1}([y,t])=[y-tA^{-1}a,t].
\]
At the endpoints \(a=A(A^{-1}a)\) verifies the gluing.
The based lift \(x(t)=-tA^{-1}a\) maps to the zero-section
base loop. The vector a is the actual root-gluing lift; the
map fixes every marked fibre period and adds no integral period.

Use the clockwise base generator \(\alpha=-\mu\). Its roof
has linear part \(A^{-1}\), through the maps
\[
 R([y,t])=[Ay,1-t],\qquad
 R^{-1}([z,s])=[A^{-1}z,1-s].
\]
The endpoint relation \([Ay,1]=[y,0]\) in the target verifies
both arrows, and their composites are the identity. The full
based map and inverse are therefore
\[
 RU([x,t])=[Ax+ta,1-t],\qquad
 U^{-1}R^{-1}([z,s])=[A^{-1}(z-(1-s)a),1-s].
\]
On the marked fibre over the base point they are the identity.
They preserve its complex orientation. The coordinate s has
clockwise base orientation; the physical positive base orientation
is \(-ds\). If \(\eta_+\) is dual to \(\mu\), and \(\eta\)
is dual to \(\alpha\), the induced maps and their inverses are
\[
 \psi\longmapsto\psi,\quad h\longmapsto h,\quad
 \eta_+\longmapsto-\eta,\qquad h=-w+3\delta.
\]
The invariant forms and integral image lattices for A and its
inverse agree by \(A^{-1}-I=-A^{-1}(A-I)\), with inverse
factor \(-A\). These based maps specify one fixed pair \((V,E)\);
only the finite-cap logarithmic attachment varies.

The root generator is
\[
 g_{\rm root}(z,s)=(Az+\ell/m,e^{-2\pi i/m}s),\qquad
 \ell=n(e_\gamma-e_u)+d\delta.
\]
Its logarithmic collar and inverse are addition and subtraction
of \(\log s\,\Pi(s)\ell/(2\pi i)\).
The positive exterior meridian maps to \(g_{\rm root}^{-1}\), and
\(\alpha\) maps to \(g_{\rm root}\). Thus the complex-positive normal circle is
\(\nu_+=m\mu+\ell\). Hold the cap point fixed and traverse
\(s(t)=s_0e^{-2\pi it}\). The logarithm drops by t, the ordinary
fibre coordinate changes by \(-t\Pi\ell\), and the ordinary
base winds clockwise m times. This clockwise normal circle is
\begin{equation}
 \beta_{m,n,d}=m\alpha-\ell=-\nu_+.
 \hypertarget{tag-Q30J}{}\label{os49:eq:clockwise-normal-circle}
\end{equation}
Its values on \((\psi,h,\eta)\) are \((-n,-3d,m)\).
The circle-parameter map \(e^{2\pi it}\mapsto e^{-2\pi it}\)
and its identical inverse have degree \(-1\).
The root quotient gives a free geometric normal circle action,
including the cases where this abelianized vector is divisible.

Keep the complex orientations of the completed X and the central
cap B. If \(\Theta_+\) is the complex-normal Thom generator, use
\(\Theta_-=-\Theta_+\) for the clockwise parameter. The maps
between Thom identifications multiply every shifted group by
\(-1\), with the same inverse. Accordingly
\((\pi_\beta)_*=-(\pi_{\nu_+})_*\) in fibre integration.
These choices fix the signs of both the transfers and the pair
sequence used below.

The cap retracts onto the compact flat four-manifold
\(B=B_{m,n,d}=T^4/\langle g_{\rm root}\rangle\).
Its rational cohomology is the invariant exterior algebra, with
Betti numbers \((1,2,2,2,1)\), so its Euler characteristic is zero.
Abelianizing \(g_{\rm root}^{\,m}=\ell\) and the conjugation relations gives
\[
 H_1(B)=\mathbb Z^3/\langle(-n,-3d,m)\rangle
       =\mathbb Z^2\oplus\mathbb Z/g,
\]
and its first-cohomology pullback is the exact lattice
\begin{equation}
 L_{m,n,d}=H^1(B)
 =\ker\bigl[(-n,-3d,m):\mathbb Z^3\longrightarrow\mathbb Z\bigr].
 \hypertarget{tag-Q30K}{}\label{os49:eq:cap-dual-lattice}
\end{equation}

\paragraph{An integral basis of the boundary two-cohomology.}
\hypertarget{tag-000A}{}\label{an-integral-basis-of-the-boundary-two-cohomology}
In the clockwise roof put
\[
 u_2=\psi h,\qquad v_2=(\psi+u)w,\qquad
 p=\eta\psi,\qquad q=\eta h/3.
\]
Under the preceding roof reversal, \(u_2,v_2\) are fixed,
and \(p_+=\eta_+\psi,q_+=\eta_+h/3\) map to \(-p,-q\);
the inverse has the same signs. We establish that
\((u_2,v_2,p,q)\) is an integral basis, retaining the division
by three in q.

Use the reference cap \((m,n,d)=(3,1,0)\), with projection
\(\pi_0\). The integral class \(-\psi\) evaluates to one on
\(\beta=3\alpha-\ell_0\). The function
\(f=e^{-2\pi i\psi}\) is well defined on the linear roof and
has weight one under this oriented circle action. The maps
\[
 E\longrightarrow B\times S^1,\qquad x\longmapsto(\pi_0(x),f(x)),
 \qquad (b,z)\longmapsto z\cdot\sigma(b),
\]
are inverse, where \(\sigma(b)=f(x)^{-1}\cdot x\) for any
point x over b. Weight-one equivariance makes this section
independent of the chosen x. This proves the integral product
splitting with precisely the marked circle.

The first-cohomology pullback from B has basis
\(h,3\psi+\eta\). The complete finite specialization in
\cref{night:P3F.29} identifies its integral two-cohomology
with the full fibre invariant lattice \(\langle u_2,v_2\rangle\).
Its boundary pullback is
\[
 u_2+q,\qquad v_2.
\]
Indeed these invariant forms are horizontal for
\(3\alpha-\ell_0\). Averaging forms over the finite torus
cover computes all rational cohomology. The kernel of restriction
is the rational span of p,q, and contraction determines the
displayed correction uniquely. The product classes supplied by
\(-\psi\) are \((-\psi)h=-u_2\) and
\((-\psi)(3\psi+\eta)=p\). Hence
\((v_2,u_2+q,u_2,p)\) is an integral basis. The inverse change
recovers \(q=(u_2+q)-u_2\) and fixes the other three classes,
giving the asserted basis without discarding the factor three.

\paragraph{The exterior restriction lattice.}
\hypertarget{tag-01O2}{}\label{determine-the-exterior-restriction-without-guessing-a-splitting}
Fill V by the reference cap. The resulting manifold is the
sphere of \cref{night:P3F.26}. Integral Mayer--Vietoris gives
\[
 H^2(V)=\mathbb Z^2,\qquad H_1(V)=\mathbb Z,\qquad
 H_2(V)=\mathbb Z^2.
\]
Restriction embeds \(H^2(V)\) as a direct summand of
\(H^2(E)\), complementary to the cap pullback just computed.
The generator of \(H_1(V)\) is \(\gamma\) by the same marked
van Kampen quotient.

Project two exterior restriction basis vectors modulo \(v_2\) to
\(U_1,U_2\in\mathbb Z^3\), in the order \((u_2,p,q)\), and
write \(W=U_1\times U_2\). For the period-one clockwise normal
action, direct contraction of the wedge factors gives the integral
fibre integrations
\begin{equation}
 \begin{aligned}
 (\pi_\beta)_*u_2&=3d\psi-nh,&
 (\pi_\beta)_*v_2&=0,\\
 (\pi_\beta)_*p&=m\psi+n\eta,&
 (\pi_\beta)_*q&=kh+d\eta.
 \end{aligned}
 \hypertarget{tag-Q30L}{}\label{os49:eq:normal-circle-integration}
\end{equation}
The factors \(\psi+u\) and \(w\) both vanish on \(\beta\), giving the
second formula. The division by three in \(q\) gives exactly \(k,d\)
in the fourth. These are the clockwise transfers associated to
\(\Theta_-\); the positive-normal transfers are their negatives.

For every \(k>0,d\in\mathbb Z\), the \(n=1\) filling is the
sphere proved in \cref{night:P3F.29}. Its pair sequence makes
\(H_1(B_{3k,1,d})\to H_2(V)\) an integral isomorphism.
The dual is the restriction of
\eqref{os49:eq:normal-circle-integration} to \(H^2(V)\).
The determinant of its two columns, in the primitive kernel
lattice, is \(\pm(W_1k-W_2d+W_3)\). Its absolute value is one
for every positive k and every integer d. Varying d and then k
therefore gives
\[
 W_1=W_2=0,\qquad W_3=\varepsilon\in\{1,-1\}.
\]
Write the projected columns as \((a_1,b_1,0),(a_2,b_2,0)\).
They have no q component: the first two components of their
cross product vanish and their \((u_2,p)\) minor is the unit \(\varepsilon\).
The exact exterior basis map and inverse use
\[
 M=\begin{pmatrix}a_1&a_2\\b_1&b_2\end{pmatrix},\qquad
 M^{-1}=\varepsilon\begin{pmatrix}b_2&-a_2\\-b_1&a_1\end{pmatrix}.
\]
Multiply the row of exterior basis classes by \(M^{-1}\); the
inverse multiplies by M. Their restrictions are \(u_2,p\) modulo
multiples of \(v_2\), whose transfer is zero.

\paragraph{The Smith invariants and the full homology.}
\hypertarget{tag-0353}{}\label{the-two-smith-invariants-and-every-homology-group}
In this explicit exterior basis the dual transfer has columns
\begin{equation}
 v_1^{\rm tr}=(3d,-n,0),\qquad v_2^{\rm tr}=(m,0,n)
 \quad\text{in }L_{m,n,d}.
 \hypertarget{tag-Q30M}{}\label{os49:eq:free-cap-transfer-columns}
\end{equation}
For \(n\ne0\), their cross product is
\[
 v_1^{\rm tr}\times v_2^{\rm tr}=n(-n,-3d,m).
\]
The primitive normal of \(L_{m,n,d}\) is
\((-n,-3d,m)/g\), so their index is \(g|n|\).
The gcd of their entries is g. This kernel lattice is primitive
in \(\mathbb Z^3\), so a vector's integral divisibility inside
it agrees with its ambient divisibility. Thus its first Smith
invariant is g and its second is \(|n|\). The freeness proof
gives \(g\in\{1,3\}\) and \(g\mid n\).

Use the pair sequence for \(X=V\cup N(B)\), with the clockwise
Thom identification \(H_j(N(B),E)=H_{j-2}(B)\) specified above.
Its relevant part is
\[
 H_1(B)\longrightarrow H_2(V)\longrightarrow H_2(X)
 \longrightarrow H_0(B)\xrightarrow{\ -n\ }H_1(V).
\]
The last arrow is \(-n\), since \(\beta\) maps to
\(-n\gamma\) in the exterior. Its comparison to the
complex-positive Thom sequence multiplies \(H_0(B)\) by
\(-1\) and fixes \(H_1(V)\), with identical inverse.
For \(n\ne0\) it is injective. Torsion in \(H_1(B)\) maps
to zero in free \(H_2(V)\); the transpose of the free transfer
has the computed Smith invariants. Consequently
\[
 H_2(X)=\mathbb Z/g\oplus\mathbb Z/|n|.
\]
The fundamental group calculation of \cref{night:P3F.33}
gives \(H_1(X)=\mathbb Z/|n|\). Euler additivity gives
\(\chi(X)=2\), since V and E are fixed, every flat cap has
Euler characteristic zero and the reference filling is a sphere.
The rational groups in degrees one and two vanish. Poincar\'e
duality and this Euler characteristic give zero in every middle
rational degree. Integral duality and the universal coefficient
theorem then give
\[
 H_3(X)\cong\operatorname{Ext}(H_2(X),\mathbb Z),\qquad
 H_4(X)\cong\operatorname{Ext}(H_1(X),\mathbb Z),\qquad H_5(X)=0.
\]
This proves \eqref{os49:eq:free-cap-nonzero-homology}.

For \(n=0\), freeness gives \(3\mid k\) and \(\gcd(k,d)=1\).
The target kernel has basis \(\psi,kh+d\eta\). Choose integers
\(a_{\rm Bez},b_{\rm Bez}\) with \(a_{\rm Bez}d+b_{\rm Bez}k=1\). The transfer columns are
\(3d\psi,3k\psi\), and
\(a_{\rm Bez}(3d\psi)+b_{\rm Bez}(3k\psi)=3\psi\). Their image is therefore
exactly \(3\mathbb Z\psi\), with cokernel
\(C=\mathbb Z\oplus\mathbb Z/3\).
The pair sequence gives \(0\to C\xrightarrow{i}H_2(X)
\xrightarrow{p_{\rm spl}}\mathbb Z\to0\).
Choose z with \(p_{\rm spl}(z)=1\). Its splitting and inverse are
\[
 (c,r)\longmapsto i(c)+rz,\qquad
 x\longmapsto\bigl(i^{-1}(x-p_{\rm spl}(x)z),p_{\rm spl}(x)\bigr).
\]
Thus \(H_2(X)=\mathbb Z^2\oplus\mathbb Z/3\).
The group calculation gives \(H_1(X)=\mathbb Z\);
duality gives \(H_5(X)=\mathbb Z,H_4(X)=\mathbb Z^2\).
Euler characteristic two makes the third Betti number two, and
integral duality identifies its torsion with that of \(H_2(X)\). These
are exactly \eqref{os49:eq:free-cap-zero-homology}.
\end{proof}
\section{Conifold smoothings and transverse pullbacks}
\subsection{Moving the zero of the linearization}

\begin{theorem}[The moving-zero family]\hypertarget{tag-02M0}{}\label{maps:src:D.3.1}
{\normalfont\scriptsize tag 02M0}\par
Let $S\to\mathbb P^1$ be Engel's rational elliptic surface, with
fibres $IV^*,III,I_1$ at $0,1,\infty$, respectively. Put
$L=\mathcal O_S(P-O)$, $T=t_{6P}$ and
$J=L\otimes T^*L^{-1}\simeq\pi^*\mathcal O_{\mathbb P^1}(1)$,
with the direction of $T$ fixed in \SeriesRef{push:src:A.1.2}.
Write $X^{\log}(u)$ for the logarithmically transformed sphere,
and $X(u)$ for the untwisted fibre of
\SeriesRef{push:src:A.1.2} and \SeriesRef{push:src:A.2.2}.
Choose sections $\phi_0,\phi_1$ of $\mathcal O(1)$ with
divisors $\infty,0$, respectively, and with
$\phi_0=q\phi_1$ near infinity, where $q=1/t$.
Set $\phi_s=\pi^*(\phi_0-s\phi_1)$.

For every sufficiently small fixed $u\ne0$ there is a proper flat
holomorphic family $\mathcal Y_u\to\Delta_s$, with a holomorphic
map to $\mathbb P^1$. Its joint map to $\mathbb P^1\times\Delta_s$
is also proper and flat. Its finite fillings have the same
multiplicities $3,4$ and the same admissible sphere twist data as
the chosen $X^{\log}(u)$. The central threefold $Y_{u,0}$ has
exactly one ordinary double point and is the contraction of the
diagonal $(-1,-1)$-curve of $X^{\log}(u)$. Every fibre $Y_{u,s}$
with $s\ne0$ is smooth. Near the contracted point the family is
\[
 xy=u(zv-s),\qquad q=zv.
\]
The parameter $s$ is its smoothing coordinate. On $u\ne0$, the
map from $(u,s)$ to the section coefficients $(a,b)=(u,-us)$
is a holomorphic coordinate change, and the original $u$-family
is the locus $b=0$.
\end{theorem}

\begin{proof}
We use the line bundle and its linearization in
\SeriesRef{push:src:A.1.2}. Their primary construction is given in
\cite[Proposition~2.4, Definition~2.5 and Propositions~2.7--2.8,
pp.~5--7]{E26}; see also the zero/pole construction in
\cite[Remark~8.5, p.~24]{E26}. The cusp compactification is described in
\cite[Proposition~3.1 and Corollary~3.3, pp.~8--9]{E26}.
In $\mathbb P(\mathcal O\oplus L)$ the normal line at the zero
section over $p$ is $L_p$, and the normal line at the infinity
section over $Tp$ is $L_{Tp}^{-1}$. Their product is $J_p$.

Choose Hermitian norms and restrict the two parameter discs so that
\[
 \|u\phi_s(p)\|<\delta=\frac1{32}
 \quad\hbox{for every }p\in S.
\]
This uniform choice is possible because $S$ is compact. Use the
annular core and neck
\[
 C=\{(p,\xi,s):\xi\in L_p,\ 1/4<\|\xi\|<4\},
\]
\[
 N=\{(p,x,y,s):x\in L_p,\ y\in L_{Tp}^{-1},\
        x\otimes y=u\phi_s(p),\ \|x\|,\|y\|<1/2\}.
\]
On $1/4<\|x\|<1/2$, glue by
$(p,x,y,s)\mapsto(p,x,s)$; its inverse has
$y=u\phi_s(p)/x$. On $1/4<\|y\|<1/2$, glue by
$(p,x,y,s)\mapsto(Tp,y^{-1},s)$; its inverse has
$x=u\phi_s(p)/y$. The two neck collars are disjoint since
$\delta<(1/4)^2$.

For line frames $e_j=Ae_i$ and $E_j=BE_i$ of $L_p$ and
$L_{Tp}$, respectively, the coordinates transform as
\[
 x_j=A^{-1}x_i,\qquad y_j=By_i,\qquad
 (\phi_s)_j=(B/A)(\phi_s)_i.
\]
These transformations preserve the equation and both collar
maps. On a triple overlap the multipliers multiply in the same
order in all three formulas, proving the cocycle identity.
Since $\pi(Tp)=\pi(p)$, the map to the elliptic base also
descends through both collar identifications.

Each equivalence class has at most one neck and one core
representative: the neck collars are separated, and their core
images have norms less than $1/2$ and greater than $2$.
The graph of each gluing is closed relative to the product of
the two charts. Indeed, a convergent pair of representatives
whose two limits remain in their charts still satisfies the
strict collar bounds; at either missing collar endpoint one
of those limits lies outside its chart. The quotient map is
open, so the closed equivalence relation gives a Hausdorff
complex space.

It has compact representatives
\[
 1/3\leq\|\xi\|\leq3,
 \qquad \|x\|,\|y\|\leq5/12.
\]
For example, a core representative with $\|\xi\|<1/3$
maps to the neck with $\|y\|<\delta/(1/4)=1/8$;
the other tail is identical using the second collar.
A neck representative outside the displayed closed neck
maps to the closed core using the coordinate whose norm
exceeds $5/12$. These closed subsets are compact over
$S$ times any compact parameter set. Hence the resulting
family, and its joint map to $\mathbb P^1\times\Delta_s$,
are proper.

Away from the zero divisor of $\phi_s$, the two collars give
the generator
\[
 L_{Tp}^{\times}\longrightarrow L_p^{\times},\qquad
 \xi\longmapsto u\phi_s(p)\xi,
\]
covering $T^{-1}$. Its norm contracts by a factor less than
$\delta$. Every orbit has adjacent iterates whose norms
straddle $1/2$. When the larger has norm at most $2$, it is
a core representative; otherwise the smaller iterate and
the reciprocal of the larger give a neck representative.
The disjoint-collar bound proves uniqueness up to precisely
the two displayed identifications. This supplies the
biholomorphism with the punctured line-bundle quotient.
At $u=0$ the neck is $xy=0$; its normalization is the
projective bundle, with the zero section over $p$ identified
with the infinity section over $Tp$.

Near the two finite points, write the original base coordinate
as $t$. The exact section ratio is
\[
 \phi_s=(1-st)\phi_0.
\]
Thus the new punctured family is the pullback of the original
two-parameter family by $(t,s)\mapsto(t,u(1-st))$.
The second coordinate is nonzero on these discs. On the
root covers $\sigma_1^3=t$ and $\sigma_2^4=t-1$ it is
invariant under the root rotation. Pull back the smooth
good-reduction families and their fixed affine twist data.
The invariant circle coordinate changes by a primitive
multiple of $1/3$ or $1/4$, respectively; every nonidentity
power therefore acts freely on the central torus. Away from
the centre the root rotation gives freeness. The quotients
are smooth proper families with the required multiplicities.

The punctured gluing is explicit. With the original marked
period map $\Pi$ and invariant vector $v_j$, put
\[
 \sigma_j^{\mathrm{tr}}=
       \frac{\log\sigma_j}{2\pi i}\Pi v_j,
 \qquad h_j(\zeta)=\zeta+\sigma_j^{\mathrm{tr}}.
\]
Its inverse subtracts the same translation. A change of
logarithm adds a period. For the clockwise root generator,
the compatible logarithm decreases by $2\pi i/m_j$;
period equivariance and invariance of $v_j$ give
$h_jg_j^{\log}h_j^{-1}=g_j^0$.
These are holomorphic overlap maps on the pulled-back
families, retaining the full period functions. Gluing the
proper pieces over the base discs gives the asserted global
space. At $s=0$ the parameter pullback is exactly $u$, so
the finite pieces and their overlaps are the original ones.

Flatness over the $s$-disc follows directly from the neck
equation. In its regular ambient local ring $A$, write
$F=xy-u f_s$. Modulo $s-s_0$, the element $F_0$ is nonzero
in the domain $A/(s-s_0)$, since its $xy$ coefficient is one.
An equality $(s-s_0)h=Fk$ implies $k=(s-s_0)k'$ by
reduction modulo $s-s_0$, and then $h=Fk'$ by cancellation.
Thus $A/(F)$ is torsion free over the discrete valuation
ring $\mathbb C\{s-s_0\}$. Each finitely generated submodule
is free, and the whole module is their directed union;
tensor products commute with that union, proving flatness.
The core is a product in $s$. The finite pieces are locally
free quotients of the smooth pulled-back root families,
so their maps to the $s$-disc are also flat.

The joint map to $(t,s)$ is flat as well. Away from a zero
of $\phi_s$, its local form is the product of
$S\to\mathbb P^1$ with a nonzero line coordinate and the
$s$-disc. The surface map is flat by the same torsion-free
criterion. At a smooth moving zero, local surface coordinates
$(q,r)$ give $q=s+xy/u$, so $(q-s,s)=(xy/u,s)$ is a product
of the flat nodal-curve map with smooth factors. At the
elliptic node, the base coordinates
\[
 (Q,R)=(q,u(q-s))=(zv,xy)
\]
have inverse $q=Q$, $s=Q-R/u$. This is the product of the
two flat nodal-curve maps. Finally, on a finite root disc,
$\mathbb C\{\sigma,s\}$ is free over
$\mathbb C\{t,s\}$ with basis $1,\sigma,\ldots,\sigma^{m-1}$;
compose this flat map with the smooth root family and use
the local isomorphisms of the free finite quotient.

The zero of $\phi_s$ near infinity is $q=s$.
At the unique node of $S_\infty$ choose $q=zv$ and absorb
the nonvanishing factor $\phi_1$ into a line coordinate.
The equation is then $xy=u(zv-s)$. For fixed $u\ne0$ set
${\widehat x}=x/u$. Its relative gradient is
\[
 (y,{\widehat x},-v,-z)
\]
for the equation ${\widehat x}y-zv+s=0$. The singular ideal of a
fibre is $({\widehat x},y,z,v,s)$, so its only singular point occurs
at $s=0$ and is the origin. The Hessian of ${\widehat x}y-zv$ is
nondegenerate. At a smooth point of the moving zero the
coordinate $q$ has nonzero differential, so the neck is
smooth. The finite pieces have already been shown smooth.
This proves every singularity assertion, including uniqueness.

At $s=0$ the small resolution of
$\det\left(\begin{smallmatrix}{\widehat x}&v\\z&y\end{smallmatrix}\right)=0$
has the two charts
\[
 ({\widehat x},v,z,y)=(a,a\tau,b,b\tau),\qquad
 ({\widehat x},v,z,y)=(a'\rho,a',b'\rho,b'),
\]
with overlap
\[
 \rho=\tau^{-1},\qquad a'=a\tau,\qquad b'=b\tau.
\]
Both map to the given node and have
$q=ab\tau=a'b'\rho$. The exceptional locus is
$a=b=0$, with coordinates $\tau,\rho=1/\tau$.
The displayed normal-coordinate changes identify its normal
bundle with $\mathcal O(-1)\oplus\mathcal O(-1)$.
The incidence description inside the node times
$\mathbb P^1$ makes the resolution proper and an
isomorphism off the origin.

In character coordinates $({\widehat x},z,q)$ the chart cones have
height-one vertices
\[
 (1,0),(0,1),(0,0),\qquad
 (1,0),(0,1),(1,1).
\]
Their determinants are $1$ and $-1$, and their common
edge is the prescribed $e_1-e_2$ diagonal. The punctured
period generators are the actual line-bundle generators
\[
 (\zeta,\xi)\longmapsto
   (c_1(q)q\zeta,c_2(q)\xi),\qquad
 (\zeta,\xi)\longmapsto
   (c_3(q)\zeta,u c_4(q)q\xi).
\]
Here all four holomorphic units are those of the same
surface, line bundle and linearization as in
\SeriesRef{push:src:A.1.2}. The two annular plumbings
give the unsubdivided square. The displayed subdivision
commutes with the full unit actions: multiplication by a
unit acts by the corresponding invertible character
factor in each chart. Its chart monomials and collar
maps are therefore precisely those of the chosen model
in \SeriesRef{push:src:A.2.2}. Passing to periods contracts
its single orbit of diagonal curves. The finite pieces
already agree at $s=0$, giving the claimed global map
$X^{\log}(u)\to Y_{u,0}$.

Finally, the coefficients of $u(\phi_0-s\phi_1)$ in
$(\phi_0,\phi_1)$ are $(a,b)=(u,-us)$. Their Jacobian
determinant is $-u$, with inverse $u=a$, $s=-b/a$ on
$u\ne0$. In the node equation $Xy-zv+s=0$, the derivative
with respect to $s$ is $1$ in the local deformation
quotient
$\mathbb C\{X,y,z,v\}/(X,y,z,v)$.
This proves the asserted independence of the smoothing
and nonzero scale parameters. Before the finite
replacements, the $u=0$ equation is $xy=0$ for every
$s$, giving the pinched projective bundle described above
\cite[Remark~8.1, p.~22]{E26}.
\end{proof}


\subsection{Topology of the smoothing}
\begin{theorem}[Topology of the smooth conifold fibre]
\hypertarget{tag-02M1}{}\label{maps:src:D.3.2}
{\normalfont\scriptsize tag 02M1}\par
For the sphere twist data of \cref{maps:src:D.3.1}, fix its
sufficiently small nonzero parameter $u$. For $0<|s|\ll1$, the
complex-oriented manifold $Y_{u,s}$ is orientation-preservingly
diffeomorphic to $S^3\times S^3$. Its integral homology is
$(\mathbb Z,0,0,\mathbb Z^2,0,0,\mathbb Z)$, its middle intersection
form has matrix
$\left(\begin{smallmatrix}0&1\\-1&0\end{smallmatrix}\right)$
in an integral basis, and $w_2=p_1=0$.
\end{theorem}

\begin{proof}
Let $E\simeq\mathbb P^1$ be the contracted curve in
$X^{\log}(u)\simeq S^6$. Its normal bundle is
$\mathcal O(-1)^{\oplus2}$. In the real identification
$\mathbb C^2\to\mathbb H$, $(z_1,z_2)\mapsto z_1+z_2j$,
its clutching loop acts by left multiplication by $e^{-i\theta}$.
This loop contracts in the unit quaternions $S^3$, giving an
oriented real trivialization. A closed tubular neighbourhood is
therefore $A\simeq S^2\times D^4$, with boundary
$L\simeq S^2\times S^3$. Put $C=S^6\setminus\operatorname{int}A$.

For the fixed nonzero \(u,s\), choose \(c\in\mathbb C\) with
\(|c|=1\) and \(-c^2s=\epsilon=|s|\). The complex-linear map
and its inverse are
\[
 (a_1,a_2,a_3,a_4)
 =\frac c2\left(\frac{x}{u}+y,\frac{x/u-y}{i},
                         z-v,\frac{z+v}{i}\right),
\]
\[
 x=\frac u c(a_1+ia_2),\quad y=\frac1c(a_1-ia_2),\quad
 z=\frac1c(a_3+ia_4),\quad v=\frac1c(-a_3+ia_4).
\]
They preserve the complex orientations, and
\(\sum a_j^2=c^2(xy/u-zv)=\epsilon\).
Writing \(a=p+iq\), the quadric equations are
\(|p|^2-|q|^2=\epsilon\) and \(p\cdot q=0\).
The mutually inverse maps
\[
 a\longmapsto(p/|p|,q),\qquad
 (v,q)\longmapsto\sqrt{\epsilon+|q|^2}\,v+iq
\]
identify it with $TS^3$. A ball truncation becomes a disc bundle;
the quaternion frame $(vi,vj,vk)$ identifies that bundle with
$B=S^3\times D^3$. On a sufficiently small fixed boundary sphere,
the local family is a submersion. Horizontal lifts tangent to that
boundary extend over the compact exterior, and their flows identify
the complements as $s$ varies. Thus
\[
 Y_{u,s}=C\cup_\varphi B
\]
for the boundary diffeomorphism supplied by the family. These
identifications preserve the complex orientations.

Relative general position in $S^6$ makes null-homotopies of loops
and two-spheres disjoint from $E$: their filling domains have
dimensions two and three, whose sums with $\dim E=2$ are less
than six. Radial retraction of the punctured tube gives
$\pi_1(C)=\pi_2(C)=0$. Excision and the oriented Thom map give
\[
 H_j(S^6,C;\mathbb Z)\simeq H_j(A,L;\mathbb Z)
       \simeq H_{j-4}(S^2;\mathbb Z).
\]
The fundamental class maps to $[S^2]$ in degree six. The pair
sequence consequently gives $H_3(C)=\mathbb Z$ and no other
positive reduced homology. Its degree-four boundary map sends a
normal disc to its boundary sphere, proving that
$H_3(L)\to H_3(C)$ is an isomorphism with that meridian as generator.

The projection $\partial B=S^3\times S^2\to S^3$ also induces
an isomorphism in degree three. In these marked generators the
Mayer--Vietoris homomorphism has column $(1,e)^{\mathsf t}$,
where $e=\pm1$. Its integral cokernel coordinate is
$(a,b)\mapsto b-ea$, with inverse represented by $(0,1)$.
The remaining degree-two boundary map has zero target. Hence
\[
 0\longrightarrow\mathbb Z\longrightarrow H_3(Y_{u,s};\mathbb Z)
       \longrightarrow\mathbb Z\longrightarrow0,
 \qquad H_2(Y_{u,s};\mathbb Z)=0.
\]
To split the extension, take the fibre $S^2$ in $\partial B$.
It bounds its normal three-disc in $B$ and, after applying
$\varphi$, a mapped three-ball in $C$ because $\pi_2(C)=0$.
Their difference is a cycle whose connecting image is the
generator of $H_2(L)$. Together with the image of the zero-section
of $B$, it gives an integral basis of $H_3(Y_{u,s})\simeq\mathbb Z^2$.
Van Kampen gives $\pi_1(Y_{u,s})=0$, and Hurewicz then gives
$\pi_2(Y_{u,s})=0$. Duality and universal coefficients give the
remaining asserted homology groups.

The middle form is skew-symmetric and unimodular. Its matrix in
the obtained basis is
$\left(\begin{smallmatrix}0&k\\-k&0\end{smallmatrix}\right)$
with $k=\pm1$; changing the second generator by $k$ gives the
displayed matrix. Sending this basis to the factor-sphere basis
of $S^3\times S^3$ is an isometry. Both normal invariants
$H_3\to\pi_2SO(3)$ vanish, since the covering
$S^3\to SO(3)$ identifies $\pi_2SO(3)$ with zero.
Wall's classification in
\cite[Recollection~2.2 and Definition~2.5, pp.~7--8;
Theorem~2.6, p.~8]{BHS19} therefore gives an oriented
diffeomorphism up to connected sum with a homotopy six-sphere.
The dimension-six table and its smooth interpretation in
\cite[pp.~504--505]{manKM63} give $\Theta_6=0$, which removes
that connected-sum factor.

Finally $H^2(Y_{u,s};\mathbb Z/2)=0$ and
$H^4(Y_{u,s};\mathbb Z)=0$, so $w_2=p_1=0$.
\end{proof}

\Needspace{12\baselineskip}
\subsection{Complex invariants of the smoothing}
\begin{theorem}[Complex invariants of the smoothing]
\hypertarget{tag-02M2}{}\label{maps:src:D.3.3}
{\normalfont\scriptsize tag 02M2}\par
For $u\ne0$ sufficiently small and $0<|s|\ll1$, let
$g:Y_{u,s}\to B=\mathbb P^1$ be the map of
\cref{maps:src:D.3.1}, and
$S'_2=(g^{-1}(1))_{\mathrm{red}}$. Then
\[
 K_{Y_{u,s}}\simeq g^*\mathcal O_B(-1)
 \otimes\mathcal O_{Y_{u,s}}(2S'_2),\qquad a(Y_{u,s})=1,
\]
\[
 (h^{0,0},h^{0,1},h^{0,2},h^{0,3})=(1,1,0,0).
\]
In particular $K^{\otimes12}\simeq g^*\mathcal O_B(-6)$
and the Kodaira dimension is $-\infty$.
\end{theorem}

\begin{proof}
Fix the stated nonzero $u$ and write $Y=Y_{u,s}$.
We may shrink its $s$-disc throughout.

\paragraph{The canonical form.}
Write $(x_E,y_E,t)$ for the coordinates of the elliptic
surface in \cite[\S2, p.~3]{E26}. On its regular part take
\[
 \omega_S=dt\wedge\frac{dx_E}{y_E},\qquad
 \Omega=\frac{d\xi}{\xi}\wedge\omega_S
\]
on the punctured line bundle. Replacing a line frame by
$\xi'=h\xi$ adds $dh/h$, pulled back from the surface;
its wedge with $\omega_S$ is zero. Elliptic translation
preserves the relative differential $dx_E/y_E$, and the
horizontal term in its pullback disappears after wedging
with $dt$. The lifted translation and the logarithmic
overlaps therefore preserve $\Omega$. Its pullback
characterizes a meromorphic form on the quotient.

At $t=0$, the good-reduction coordinates
$t=r^3$, $x_E=r^4X_{\rm W}$, $y_E=r^6Y_{\rm W}$ give
$\omega_S=3\,dr\wedge dX_{\rm W}/Y_{\rm W}$; the central elliptic curve
is $Y_{\rm W}^2=X_{\rm W}^3+2$. At $t=1$, the coordinates
$t-1=r^4$, $x_E=r^2X_{\rm W}$, $y_E=r^3Y_{\rm W}$ give
$\omega_S=4r^2\,dr\wedge dX_{\rm W}/Y_{\rm W}$, with smooth central
curve $Y_{\rm W}^2=X_{\rm W}^3-3X_{\rm W}$. The total-space finite quotients are
free and hence \'etale. Thus the orders of $\Omega$ on
the reduced finite fibres are respectively zero and two.

At infinity put $q=1/t$, $X_{\rm W}=q^2x_E$, $Y_{\rm W}=q^3y_E$.
The equation and form become
\[
 Y_{\rm W}^2=X_{\rm W}^3-3(1-q)X_{\rm W}+2(1-q)^2,\qquad
 \omega_S=-\frac{dq\wedge dX_{\rm W}}{qY_{\rm W}}.
\]
At $(q,X_{\rm W},Y_{\rm W})=(0,1,0)$ the derivative of the defining
equation with respect to $q$ is nonzero. In nodal
coordinates $q=zv$, the numerator $dq\wedge dX_{\rm W}/Y_{\rm W}$
is a nonvanishing holomorphic unit times $dz\wedge dv$.
On each unchanged toric chart the height-one monomial
map has determinant $\pm1$ and $q=z_0z_1z_2$.
It takes the logarithmic three-form to
$\pm dz_0\wedge dz_1\wedge dz_2/q$, up to a
nonvanishing holomorphic unit.

On the moving chart $xy=u(zv-s)$ the form
\[
 \eta=\frac{dx\wedge dz\wedge dv}{x}
\]
extends as a nonvanishing holomorphic form. Its other
coordinate expressions are
\[
 -\frac{dy\wedge dz\wedge dv}{y},\qquad
 \frac{dx\wedge dz\wedge dy}{uz},\qquad
 \frac{dx\wedge dy\wedge dv}{uv}.
\]
They follow by eliminating $x$, $v$, and $z$,
respectively. The four charts cover the hypersurface
when $us\ne0$. On $x\ne0$, the plumbing coordinate
is the line coordinate and $\Omega$ is a unit times
$\eta/q$. This equality extends across the charts.
At $q=0$ one has $xy=-us\ne0$, so the pole divisor
there is precisely $g^*(\infty)$; at $q=s$ both $q$
and the extended $\eta$ are nonzero. Consequently
\[
 \operatorname{div}(\Omega)=-g^*(\infty)+2S'_2.
\]
Sending the meromorphic frame of this divisor to
$\Omega$ gives the asserted canonical-bundle isomorphism.
Since $4S'_2=g^*(1)$, its twelfth power is
$g^*\mathcal O_B(-6)$.

A holomorphic function on $g^{-1}(V)$ is constant on
each ordinary torus fibre. Local sections of the
submersion make its descended values holomorphic away
from the exceptional base values. Properness bounds
the function over any smaller closed disc in $V$,
so those values extend across its finitely many
exceptional points. Thus $g_*\mathcal O_Y=\mathcal O_B$,
with inverse given by pullback. It follows that
$H^0(Y,g^*\mathcal O_B(-6m))=0$ for every $m>0$.
The twelfth power of a nonzero section of $K_Y^m$
would contradict this equality. All plurigenera vanish.

\paragraph{Meromorphic functions.}
By \cref{maps:src:D.3.2}, $H_1(Y;\mathbb Z)$ and
$H_2(Y;\mathbb Z)$ vanish; universal coefficients give
$H^2(Y;\mathbb Z)=0$. An effective divisor $D$
dominating $B$ restricts to a nonzero effective curve
on an ordinary general torus fibre $F$. For a
translation-invariant K\"ahler form $\omega$ on $F$,
\[
 0<\int_{D|_F}\omega
 =\bigl\langle c_1(\mathcal O_Y(D)|_F)\smile[\omega],
 [F]\bigr\rangle=0,
\]
a contradiction. Hence every effective divisor is vertical.

The polar divisor of a meromorphic function on $Y$
is therefore vertical. Its restriction to a general
torus fibre has no polar divisor, extends across its
indeterminacy by normality, and is constant. The
closed graph has proper analytic image in $B\times
\mathbb P^1$ by \cite[\S7, Satz~1, p.~60]{Grauert1960Coherent}.
This image is an irreducible curve of generic degree
one over $B$. Its normalization maps finitely and
birationally to the normal curve $B$, hence
isomorphically. The other projection gives a
meromorphic function on $B$ whose pullback is the
original function. Reading its value on an ordinary
fibre is inverse to pullback. Thus
$\mathbb C(Y)=g^*\mathbb C(B)$ and $a(Y)=1$.

\paragraph{Structure-sheaf cohomology.}
For the sphere used in the contraction,
\cite[Proposition~9.13, pp.~74--76]{AC26} gives
\[
 f_*\mathcal O=\mathcal O_B,\qquad
 R^1f_*\mathcal O=\mathcal O_B\oplus\mathcal O_B(-1),
 \qquad R^2f_*\mathcal O=\mathcal O_B(-1).
\]
The precise identification with $X^{\log}(u)$ is
the map
$[\xi,\eta]\mapsto[\xi,e^{2\pi i\eta}]$ of
\SeriesRef{fam:src:B.1.1}, together with the transported
reference lift of \SeriesRef{fam:src:B.1.4}
and the filled-space comparison of
\SeriesRef{fam:src:B.2.2}. The periods are
$(6\mu,\beta_c)$, $(\tau,\mu)$, $(1,0)$, $(0,1)$;
the contracting line-bundle generator is the inverse
of the first period generator. Local logarithms
provide the inverse quotient map. These comparisons
include the specified finite fillings and cusp, so
their pullback maps identify the structure sheaves
and their cohomology. Leray gives dimensions
$(1,1,0,0)$ on $X^{\log}(u)$.

Let $\rho:X^{\log}(u)\to Y_{u,0}$ be the small
resolution. Its exceptional curve $E$ has normal
bundle $\mathcal O_E(-1)^{\oplus2}$. For its ideal $I$,
\[
 I^m/I^{m+1}=\operatorname{Sym}^m(N_E^*)
 =\mathcal O_E(m)^{\oplus(m+1)}.
\]
The pullback of the node's maximal ideal is $I$:
in a resolution chart the four rank-one matrix
entries are $a,b,at,bt$, generating $(a,b)$.
The displayed graded pieces have no positive
cohomology. The exact sequences of successive
thickenings give the same vanishing on every
infinitesimal neighbourhood of $E$.
The analytic formal-functions and coherence
theorems \cite[\S6, Haupts\"atze~I and IIa,
p.~59]{Grauert1960Coherent}, with
\cite{Grauert1963Correction}, identify the completed
stalks of $R^j\rho_*\mathcal O$ with these zero
inverse limits for $j>0$. Nakayama's lemma makes
the stalks zero. Away from the node $\rho$ is
an isomorphism. Normality of the node also gives
$\rho_*\mathcal O=\mathcal O_{Y_{u,0}}$.
Thus the Leray pullback isomorphisms identify
$H^q(Y_{u,0},\mathcal O)$ with
$H^q(X^{\log}(u),\mathcal O)$.

The family in \cref{maps:src:D.3.1} is proper and
flat over the reduced connected $s$-disc.
Upper semicontinuity and constancy of the holomorphic
Euler characteristic
\cite[\S7, Satz~3, pp.~62--63, and Satz~6,
p.~64]{Grauert1960Coherent} give
$h^{0,2}=h^{0,3}=0$ and $\chi(\mathcal O_Y)=0$
near the centre. Connectedness and compactness
give $h^{0,0}=1$, hence $h^{0,1}=1$.
\end{proof}




\subsection{Comparison with Calabi--Eckmann deformations}
\begin{proposition}[Comparison with Calabi--Eckmann deformations]
\hypertarget{tag-02M3}{}\label{maps:src:D.3.4}
{\normalfont\scriptsize tag 02M3}\par
The smooth threefolds $Y_{u,s}$ of \cref{maps:src:D.3.1}, for the
parameters of \cref{maps:src:D.3.3}, are not biholomorphic to any
member of the algebraic-dimension-one subfamily on $S^3\times S^3$
displayed in \cite[Example~2.3, p.~83]{manCDP98}.
\end{proposition}

\begin{proof}
That subfamily consists of the quotients $Z$ of
$(\mathbb C^2\setminus\{0\})^2$ by
\[
 t\cdot(x,y,z,w)=(e^t x,e^{at}y,e^{ibt}z,e^{ibt}w),
 \qquad a,b>0,\quad a\notin\mathbb Q.
\]
The invariant map $h:Z\to\mathbb P^1$, given by $[z:w]$, has
connected fibres. Over $[1:v]$, setting $z=1$ leaves the periods
$t=2\pi k/b$. Thus its fibre is
\[
 (\mathbb C^2\setminus\{0\})/\langle A\rangle,
 \qquad A(x,y)=(e^{-2\pi/b}x,e^{-2\pi a/b}y).
\]
The other affine chart gives the same description at infinity.
The covering space retracts onto $S^3$ and is simply connected;
its deck group is $\langle A\rangle\simeq\mathbb Z$.
Consequently every fibre has $H_1\simeq\mathbb Z$.

For $g:Y_{u,s}\to\mathbb P^1$, a general fibre is the quotient
of the degree-zero extension
$1\to\mathbb C^*\to L_t^\times\to S_t\to0$
by the contracting lift of $t_{6P}^{-1}$. Its universal cover is
$\mathbb C^2$, and its quotient period group is a rank-four lattice
\cite[Definition~2.3 and Proposition~2.6, p.~5]{E26}.
Replacing $\phi_0$ by $\phi_0-s\phi_1$ multiplies this lift by
a nonzero central scalar away from the moving zero, so it remains
a translation of the same extension. Hence the general fibre
$\mathcal F_t$ is a complex two-torus and $H_1(\mathcal F_t;\mathbb Z)\simeq\mathbb Z^4$.

Both total spaces have algebraic dimension one, by
\cite[Example~2.3, p.~83]{manCDP98} and \cref{maps:src:D.3.3}.
Their connected general fibres identify their meromorphic function
fields with $\mathbb C(\mathbb P^1)$ through $h^*$ and $g^*$.
Indeed, a meromorphic function is algebraic over this base field,
so on a general connected fibre it takes values in a finite set and
is constant. Its value descends meromorphically to the base; near an
omitted point its single-valued algebraic branch has at most a pole.

A biholomorphism $F:Y_{u,s}\to Z$ would therefore induce
\[
 (g^*)^{-1}F^*h^*:\mathbb C(\mathbb P^1)
       \xrightarrow{\sim}\mathbb C(\mathbb P^1).
\]
This automorphism is induced by a M\"obius transformation $\rho$:
its rational function and that of its inverse have compositional
degree one. It follows that $hF=\rho g$. For a regular value $t$,
restriction gives inverse biholomorphisms
\[
 F_t:g^{-1}(t)\longrightarrow h^{-1}(\rho(t)),
 \qquad (F^{-1})_{\rho(t)}:h^{-1}(\rho(t))\longrightarrow g^{-1}(t).
\]
They would induce inverse homomorphisms between $\mathbb Z^4$
and $\mathbb Z$. After tensoring with $\mathbb Q$, a map from
dimension four to dimension one has nonzero kernel and cannot
have a left inverse. This proves the asserted obstruction.
\end{proof}

\subsection{Transverse pullbacks of the sphere fibration}
\begin{theorem}[Transverse pullbacks and products]\hypertarget{tag-02MK}{}\label{maps:src:E.2.2}
{\normalfont\scriptsize tag 02MK}\par


Let $f:X\to B$ be the original sphere fibration of
\cref{maps:src:B.2.1}, with exceptional values $p_0,p_1,p_2$
as in \cref{mapsIII:def:marked-family}. Let $Y$ be a compact
complex manifold and $h:Y\to B$ a holomorphic map which is a
submersion at every point over these three values.
Then $P=X\times_B Y$ is a compact complex manifold, of dimension
$\dim_{\mathbb C}Y+2$. For $Y=B\times N$ with its projection,
$P\simeq X\times N$ holomorphically, is diffeomorphic to
$S^6\times N$, and
\begin{equation}
 \pi_1(P)=\pi_1(N),\quad
 H_k(P;\mathbb Z)=H_k(N;\mathbb Z)\oplus H_{k-6}(N;\mathbb Z),
 \quad \chi(P)=2\chi(N).
 \hypertarget{tag-029L}{}\label{maps:eq:E.15}
\end{equation}

\end{theorem}


\begin{proof}
The pullback is the inverse image of the diagonal under
\((f,h):X\times Y\to B\times B\). At a point of this inverse
image, either \(df\) is onto, or the common base value is one
of the three exceptional values and \(dh\) is onto. Hence
\(df-dh:T_xX\oplus T_yY\to T_{f(x)}B\) is onto.
The holomorphic implicit-function theorem gives a complex
submanifold of dimension \(3+\dim_{\mathbb C}Y-1\).
It is closed in the compact product, hence compact.

For the projection \(B\times N\to B\), the mutually inverse
holomorphic maps are
\[
 (x,(t,n))\longmapsto(x,n),\qquad
 (x,n)\longmapsto(x,(f(x),n)).
\]
Choose the diffeomorphism \(\sigma:X\to S^6\) from the sphere
construction. The smooth product map is
\((x,n)\mapsto(\sigma(x),n)\), with inverse
\((v,n)\mapsto(\sigma^{-1}(v),n)\).

The integral K\"unneth map has no Tor terms because
\(H_*(X;\mathbb Z)=H_*(S^6;\mathbb Z)\) is free. Its map
for the displayed direct sum and its inverse are
\[
 (a,b)\longmapsto[x_0]\times a+[X]\times b,\qquad
 c\longmapsto\bigl((\operatorname{pr}_N)_*c,
  (\operatorname{pr}_N)_*(\operatorname{pr}_X^*\omega\frown c)\bigr),
\]
where \(x_0\in X\) and \(\omega\in H^6(X;\mathbb Z)\)
satisfies \(\langle\omega,[X]\rangle=1\). The two composites
are the identity on the K\"unneth summands. Projection to N
and the section \(n\mapsto(x_0,n)\) induce inverse fundamental
group maps, since \(\pi_1(X)=0\). The Euler characteristic of
the product is \(\chi(X)\chi(N)=2\chi(N)\).
\end{proof}
\section{Local covers of the cusp}\hypertarget{tag-034I}{}\label{sec:III:carriers}

\subsection{Completed Kummer covers}
\hypertarget{tag-0385}{}\label{wb:cusp-cover-section}

\begin{definition}[Completed cusp root data]
\hypertarget{tag-Q019}{}\label{mapsIII:def:completed-cusp-root-data}
{\normalfont\scriptsize tag Q019}\par
Use the completed target cusp and period correction of
Definition~\ref{mapsIII:def:kummer-data}. Write $t=q$ for its base coordinate. In these coordinates the data are
\[
 B_0=\begin{pmatrix}0&1\\-1&0\end{pmatrix},\qquad
 Z(s)=sB_0+C(t),\qquad t=e^{2\pi i s},
\]
where $C$ is holomorphic on a disc $\Delta_t$ about zero.  Put
\[
 c_{\bar\lambda}(t)
 :=\exp\!\bigl(2\pi iC(t)\bar\lambda\bigr)
 \in(\mathbb C^*)^2,
 \qquad \bar\lambda\in\bar\Lambda\simeq\mathbb Z^2.
\]
Let $\Sigma$ be the fan in
$N'=\mathbb Z^2\oplus\mathbb Z$ consisting of the cones over the cells of
the translation-invariant $A_2$ triangulation at height one, and write
$Y=Y_\Sigma$.  On the dense torus, with coordinates $(x_1,x_2,t)$, the
target action is
\begin{equation}
 \Psi_{\bar\lambda}(x,t)
 =\bigl(c_{\bar\lambda}(t)t^{B_0\bar\lambda}x,t\bigr).
 \hypertarget{tag-0388}{}\label{wb:cusp-cover-target-action}
\end{equation}
Its quotient is the completed cusp family
$f_0\colon N_0=Y/\bar\Lambda\to\Delta_t$.

Fix an integer $N\geq1$. For a sufficiently small $r>0$, take
\[
 \Delta_v=\{|v|<r\},\qquad \Delta_t=\{|t|<r^N\},\qquad
 \Delta_v=\rho_N^{-1}(\Delta_t),
\]
and restrict the target cusp to $\Delta_t$. The germ $C$ is represented
on a disc containing $\{|z|<r\}$. Let $v$ be the root coordinate and set
\begin{equation}
 \rho_N\colon\Delta_v\longrightarrow\Delta_t,\qquad
 \rho_N(v)=v^N,\qquad t=v^N.
 \hypertarget{tag-037Q}{}\label{wb:cusp-cover-base-change}
\end{equation}
On a copy $Y^{(N)}$ of the same toric variety, with dense-torus coordinates
$(\xi_1,\xi_2,v)$, define
\begin{equation}
 C_N(v):=\frac1N C(v^N),\qquad
 c^{(N)}_{\bar\lambda}(v)
 :=\exp\!\bigl(2\pi iC_N(v)\bar\lambda\bigr)
 \hypertarget{tag-0387}{}\label{wb:cusp-cover-source-factor}
\end{equation}
and
\begin{equation}
 \Psi^{(N)}_{\bar\lambda}(\xi,v)
 :=\bigl(c^{(N)}_{\bar\lambda}(v)
          v^{B_0\bar\lambda}\xi,v\bigr).
 \hypertarget{tag-0386}{}\label{wb:cusp-cover-source-action}
\end{equation}

Also write
\[
 \Psi^{\rm un}_{\bar\lambda}(\xi,v)
 =\bigl(c_{\bar\lambda}(v)v^{B_0\bar\lambda}\xi,v\bigr)
\]
for the action with the unrooted correction at $v$.
\end{definition}

\begin{theorem}[Central fibres and deck transformations of the Kummer map]
\hypertarget{tag-037X}{}\label{wb:cusp-cover-extension-theorem}
{\normalfont\scriptsize tag 037X}\par
For the data of Definition~\ref{mapsIII:def:completed-cusp-root-data},
the local Kummer maps of Theorem~\ref{maps:src:C.1.1} have the
following central-fibre and whole-family descriptions.

\begin{enumerate}
\item The source action extends freely and properly discontinuously to
$Y^{(N)}$.  Its quotient
\[
 f_N\colon N_0^{\langle N\rangle}:=Y^{(N)}/\bar\Lambda
 \longrightarrow\Delta_v
\]
is a proper holomorphic map from a smooth threefold.  Its central fibre
$W_N$ is reduced, irreducible and nonnormal, with normalization the
toric del Pezzo surface of degree six.

\item Multiplication by $N$ on the cocharacter lattice defines a finite
flat toric map of degree $N^3$,
\begin{equation}
 P_N\colon Y^{(N)}\longrightarrow Y,\qquad
 P_N(\xi_1,\xi_2,v)=(\xi_1^N,\xi_2^N,v^N).
 \hypertarget{tag-0389}{}\label{wb:cusp-cover-toric-power-map}
\end{equation}
It is equivariant for the source and target lattice actions and descends
to a finite flat map of degree $N^3$,
\begin{equation}
 \kappa_N\colon N_0^{\langle N\rangle}\longrightarrow N_0,
 \qquad f_0\circ\kappa_N=\rho_N\circ f_N.
 \hypertarget{tag-037T}{}\label{wb:cusp-cover-completed-map}
\end{equation}

\item The map $\kappa_N$ is \textup{\'etale} away from $W_N$, and
\begin{equation}
 \kappa_N^*W=NW_N,\qquad
 \operatorname{Ram}(\kappa_N)=(N-1)W_N.
 \hypertarget{tag-037U}{}\label{wb:cusp-cover-direct-ramification}
\end{equation}
Its branch support is all of $W$ for $N>1$ and is empty for $N=1$.

\item The base-changed target
\begin{equation}
 N_0[N]:=N_0\times_{\Delta_t}\Delta_v,\qquad t=v^N,
 \hypertarget{tag-037R}{}\label{wb:cusp-cover-base-changed-target}
\end{equation}
is normal.  There is a unique factorization
\begin{equation}
 N_0^{\langle N\rangle}
 \xrightarrow{\ Q_N\ }N_0[N]\longrightarrow N_0,
 \qquad Q_N(x)=(\kappa_N(x),f_N(x)).
 \hypertarget{tag-0384}{}\label{wb:cusp-cover-relative-factorization}
\end{equation}
The finite map $Q_N$ has generic degree $N^2$.  For $N>1$ it is
quasi-\textup{\'etale}: its non-\textup{\'etale} and nonflat loci are
exactly the inverse images of the double and triple strata of the
base-changed central fibre, and it is \textup{\'etale} over the smooth
locus of $N_0[N]$.  The degree $N^2>1$ then excludes a normalization
morphism.  At $N=1$, both $\kappa_1$ and $Q_1$ are identities under
the coordinate identification of the source and target families and the
base-change isomorphism
$N_0\xrightarrow{\sim}N_0[1]$, $y\mapsto(y,f_0(y))$.

\item The map on reduced central fibres has generic degree $N^2$ and
lifts to the commutative diagram
\begin{equation}
 \begin{CD}
 E_N\simeq\mathrm{dP}_6 @>{[N]}>> E\simeq\mathrm{dP}_6\\
 @V{\nu_N}VV @VV{\nu}V\\
 W_N @>>> W.
 \end{CD}
 \hypertarget{tag-0383}{}\label{wb:cusp-cover-normalization-diagram}
\end{equation}
Here $[N]$ is toric multiplication by $N$.  It has degree $N^2$ and
ramification index $N$ along each of the six toric boundary curves.
The square respects the paired conductor curves and both triple-point
classes.

\item Over $v\ne0$, the map $Q_N$ is the connected cover with fibre
lattice inclusion
\[
 \mathbb Z\langle\widehat\gamma,\widehat u,
 N\widehat w,N\widehat\delta\rangle
 \subset
 \mathbb Z\langle\widehat\gamma,\widehat u,
 \widehat w,\widehat\delta\rangle.
\]
Its deck group on the punctured family is $(\mathbb Z/N\mathbb Z)^2$.
The fibre-linear matrix of the pulled-back clockwise analytic deck
generator is $M_0^N$, and its matrix in the cover basis is $M_0$:
\begin{equation}
 S_N^{-1}M_0^NS_N=M_0,\qquad
 S_N=\operatorname{diag}(1,1,N,N).
 \hypertarget{tag-0381}{}\label{wb:cusp-cover-monodromy-conjugacy}
\end{equation}
\end{enumerate}
\end{theorem}

\begin{proof}
The data agree with Theorem~\ref{maps:src:C.1.1} through the maps
\[
 (x,t)\longmapsto(x,q=t),\qquad
 (\xi,v)\longmapsto(\xi,v),\qquad
 \bar\lambda\longmapsto\lambda
\]
in the fixed ordered bases. These identifications fix every cone
and every character coordinate. The holomorphic
correction is the same $C$, and both source actions use
$C(v^N)/N$. Thus these maps extend through every toric chart, commute
with the period actions and the power maps, and induce inverse maps
on their quotients and fibre products. With radius $r=\epsilon_N$,
we obtain the full inverse-image discs of that construction.

Under these identifications, the earlier theorem proves the free and
properly discontinuous source action, proper smooth quotient, both
finite degrees, finite flatness of $\kappa_N$, normality of $N_0[N]$,
and the precise relative non-\'etale and nonflat loci. Its complete
chart and invariant-ring proof applies to these identical maps. In
particular its maximal-chart formula is
\begin{equation}
 \mathbb C\{z_0,z_1,z_2\}\longrightarrow
 \mathbb C\{w_0,w_1,w_2\},\qquad z_i\longmapsto w_i^N.
 \hypertarget{tag-0380}{}\label{wb:cusp-cover-maximal-chart-map}
\end{equation}
The equality of the two coordinate descriptions of equivariance is
\begin{align}
 P_N\Psi^{(N)}_{\bar\lambda}(\xi,v)
 &=\bigl((c^{(N)}_{\bar\lambda}(v))^N
             v^{NB_0\bar\lambda}\xi^N,v^N\bigr)\\
 &=\bigl(c_{\bar\lambda}(v^N)
             (v^N)^{B_0\bar\lambda}\xi^N,v^N\bigr)
 =\Psi_{\bar\lambda}P_N(\xi,v).
 \hypertarget{tag-037W}{}\label{wb:cusp-cover-equivariance-computation}
\end{align}
It uses the same identity $NC_N(v)=C(v^N)$ on the whole root disc.

We describe the normalization and conductor maps explicitly.
The local equation $v=w_0w_1w_2$ makes $W_N$ reduced. Before the
lattice quotient its components are the toric divisors $D_p$,
$p\in\mathbb Z^2$, and the action sends $D_p$ to
$D_{p+B_0\bar\lambda}$. The equality $B_0\bar\Lambda=\mathbb Z^2$
gives one component orbit with trivial component stabilizer, as in
the earlier construction. The branch-separating map is the injection
\[
 \frac{\mathbb C\{w_0,w_1,w_2\}}{(w_0\cdots w_{k-1})}
 \longrightarrow
 \prod_{i=0}^{k-1}\frac{\mathbb C\{w_0,w_1,w_2\}}{(w_i)},
 \qquad 1\le k\le3.
\]
Its kernel is zero because $\bigcap_{i<k}(w_i)=(w_0\cdots w_{k-1})$.
It is finite, has the same total ring of fractions, and has normal
target, so it is the integral closure. These maps glue; the quotient
of the separated components is $D_0$. The identity of the star fans
identifies this $D_0$ with the target normalization $E\simeq\mathrm{dP}_6$.
The two branches at an open double
stratum prove that $W_N$ is nonnormal. The ramification formula of
Theorem~\ref{maps:src:C.1.1} gives index $N$ on each central component.
Each maps onto its target component. Hence the branch support is all
of $W$ for $N>1$ and empty for $N=1$.

The relative map has the same local-ring homomorphism as
\eqref{maps:eq:C.1.2}, now written
\begin{equation}
 \begin{gathered}
 A_N:=\frac{\mathbb C\{z_0,z_1,z_2,v\}}{(z_0z_1z_2-v^N)},
 \qquad B:=\mathbb C\{w_0,w_1,w_2\},\\
 A_N\longrightarrow B,\qquad z_i\longmapsto w_i^N,
 \qquad v\longmapsto w_0w_1w_2.
 \end{gathered}
 \hypertarget{tag-037S}{}\label{wb:cusp-cover-basechange-local-ring}
\end{equation}
Write
\[
 \mathcal D_N=\{(\varepsilon_0,\varepsilon_1,\varepsilon_2)\in\mu_N^3:
                   \varepsilon_0\varepsilon_1\varepsilon_2=1\}.
\]
The identity on the three diagonal multipliers identifies this with
the group in that earlier proof, and the invariant-ring map gives
\begin{equation}
 A_N=B^{\mathcal D_N}.
 \hypertarget{tag-037Z}{}\label{wb:cusp-cover-invariant-ring}
\end{equation}
The inverse on an invariant monomial takes its common exponent residue
$r$ modulo $N$ and writes it uniquely as
$(w_0w_1w_2)^r\prod_i(w_i^N)^{b_i}$, $0\le r<N$.
The finite residue grouping is the convergent-series inverse proved
there. On the dense torus the group is the phase action
\[
 (\eta_1,\eta_2)(\xi_1,\xi_2,v)
       =(\eta_1\xi_1,\eta_2\xi_2,v),\qquad (\eta_1,\eta_2)\in\mu_N^2.
\]
Its first-chart identification is
$(\eta_1,\eta_2)\mapsto((\eta_1\eta_2)^{-1},\eta_1,\eta_2)$,
with inverse $(\varepsilon_0,\varepsilon_1,\varepsilon_2)
\mapsto(\varepsilon_1,\varepsilon_2)$.
It extends torically and commutes with every period action. The chart
invariant-ring isomorphisms therefore glue to identify its whole
quotient with $N_0[N]$.

The closed-fibre computations in the earlier proof give, at a triple point,
\begin{align}
 \dim_{\mathbb C}
 \frac{\mathbb C\{w_0,w_1,w_2\}}
      {(w_0^N,w_1^N,w_2^N,w_0w_1w_2)}
 &=N^3-(N-1)^3\\
 &=3N^2-3N+1.
 \hypertarget{tag-038A}{}\label{wb:cusp-cover-triple-fibre-length}
\end{align}
At a double stratum, the $N$ roots of the nonzero transverse coordinate
separate $N$ local branches. Each has generic degree $N$ and the
two-variable map $z_i=w_i^N$, $v=w_0w_1$, with closed-fibre length
\begin{equation}
 \dim_{\mathbb C}
 \frac{\mathbb C\{w_0,w_1\}}{(w_0^N,w_1^N,w_0w_1)}=2N-1.
 \hypertarget{tag-037V}{}\label{wb:cusp-cover-double-fibre-length}
\end{equation}
Thus the $N$ branches have total generic degree $N^2$.
For $N>1$ the excesses over the respective local generic ranks are
$(N-1)(2N-1)$ and $N-1$, exactly as in the earlier nonflatness proof.
At $N=1$ both excesses vanish and both maps are identities.

For the regular fibres use $v=e^{2\pi i s_v}$, so $t=e^{2\pi iNs_v}$.
The log coordinate $s_v$ is the coordinate $r$ of the earlier period
calculation, and
\[
 Z_N(s_v)=s_vB_0+\frac1N C(v^N),\qquad
 Z(Ns_v)=NZ_N(s_v),\qquad
 N\Pi_{\rm source}=\Pi_{\rm target}S_N.
\]
The same isomorphism $[\zeta]\mapsto[N\zeta]$, with inverse
$[\widetilde\zeta]\mapsto[\widetilde\zeta/N]$, puts the fibre map in the form
\begin{equation}
 \frac{\mathbb C^2}{N\mathbb Z^2+Z(Ns_v)\mathbb Z^2}
 \longrightarrow
 \frac{\mathbb C^2}{\mathbb Z^2+Z(Ns_v)\mathbb Z^2},\qquad
 [\widetilde\zeta]\longmapsto[\widetilde\zeta].
 \hypertarget{tag-037P}{}\label{wb:cusp-cover-additive-fibre-map}
\end{equation}
It has the same image lattice of index $\det S_N=N^2$.
In source coordinates its deck maps are
$\zeta\mapsto\zeta+(a,b)/N$, with inverse translation by $-(a,b)/N$;
their exponentials are the phase maps above.

The clockwise analytic deck generators have fibre-linear matrices
$M_0$ on the source and $M_0^N$ on the pulled-back target.
Write $M_0=I+R$, where $R(a,b,c,d)=(0,0,b,-a)$ and $R^2=0$.
Then $M_0^N=I+NR$ and $M_0^NS_N=S_NM_0$.
Their point-return maps have the inverse matrices by
Definition~\ref{mapsIII:def:three-filling-datum}; inversion gives
$M_0^{-N}S_N=S_NM_0^{-1}$. All these affine maps fix zero,
represented by $(1,1)$ after exponentiation, and
$\zeta\mapsto N\zeta$ preserves it. Both intertwining identities
therefore hold for the affine maps. On $\Lambda/S_N\Lambda$
the two target actions are trivial: their changes are respectively
$NR$ and $-NR$, whose images are $N$ times short periods.
Thus the deck group acts on the whole punctured relative family.

Finally, $P_N$ preserves every ray and sends $D_p$ to $D_p$.
On $D_0$ its induced cocharacter map on $N'/\mathbb Z(0,0,1)$
is multiplication by $N$. This is $[N]:E_N\to E$, of degree $N^2$.
Its boundary charts are $(a,b)\mapsto(a^N,b^N)$, giving index $N$
on each of the six boundary curves. The normalization maps are the
inclusions of $D_0$ followed by the period quotient maps.
Equation~\eqref{wb:cusp-cover-equivariance-computation} makes the
normalization square commute pointwise. Every source conductor gluing
by $\bar\lambda$ maps to the target gluing by the same $\bar\lambda$,
including its holomorphic unit factors. Each of the two translation
orbits of maximal triangles is preserved; these are the two
triple-point classes. The generic isomorphisms of the normalizations
give degree $N^2$ on the reduced central fibres.

At $N=1$ the actions coincide, $P_1$ is the coordinate identity,
and $\kappa_1$ is its quotient. The projection $N_0[1]\to N_0$ has
inverse $y\mapsto(y,f_0(y))$, identifying $Q_1$ with the identity.
All ramification and branch supports are empty and $S_1=I$.
\end{proof}


\begin{remark}[Comparison of the two source actions]
\hypertarget{tag-0382}{}\label{wb:cusp-cover-naive-action-failure}
{\normalfont\scriptsize tag 0382}\par
The two action formulas of
Definition~\ref{mapsIII:def:completed-cusp-root-data} have the following
comparison for every holomorphic matrix $C$ on the compatible discs.
Direct substitution gives
\[
 \begin{aligned}
 P_N\Psi^{\mathrm{un}}_{\bar\lambda}(\xi,v)
   &=\bigl(c_{\bar\lambda}(v)^N
               v^{NB_0\bar\lambda}\xi^N,v^N\bigr),\\
 \Psi_{\bar\lambda}P_N(\xi,v)
   &=\bigl(c_{\bar\lambda}(v^N)
               v^{NB_0\bar\lambda}\xi^N,v^N\bigr).
 \end{aligned}
\]
More precisely, the first map is the second followed by multiplication
in the two fibre coordinates by
\[
 d_{\bar\lambda}(v)
   =\exp\!\bigl(2\pi i\{NC(v)-C(v^N)\}\bar\lambda\bigr).
\]
Consequently the unmodified action is equivariant exactly when
$c_{\bar\lambda}(v)^N=c_{\bar\lambda}(v^N)$ for every
$\bar\lambda\in\bar\Lambda$.  Equality on the dense torus extends
through all toric charts.

The source factor \eqref{wb:cusp-cover-source-factor} satisfies
\[
 (c^{(N)}_{\bar\lambda}(v))^N=c_{\bar\lambda}(v^N),\qquad
 c^{(N)}_{\bar\lambda+\bar\mu}
   =c^{(N)}_{\bar\lambda}c^{(N)}_{\bar\mu}.
\]
Both identities follow from its holomorphic logarithm
$2\pi iC(v^N)\bar\lambda/N$.  At $N=1$ the two actions coincide
for every holomorphic $C$.

For $N>1$, the equality criterion can also be written
\[
 C(v)=C_0\text{ is constant},\qquad
 (N-1)C_0\in\operatorname{Mat}_2(\mathbb Z).
\]
To prove this, equality for the two basis vectors of $\bar\Lambda$
makes every entry of $NC(v)-C(v^N)$ integer-valued.  It is a
holomorphic matrix on a connected disc, so it is constant, with value
$(N-1)C(0)$.  Write $C(v)=\sum_{k\geq0}C_kv^k$.
The positive coefficients satisfy $NC_k=0$ when $N\nmid k$ and
$NC_k=C_{k/N}$ when $N\mid k$.  Writing $k=N^a k_0$ with
$N\nmid k_0$ proves $C_k=0$ for every $k>0$.  Conversely, a constant
$C_0$ with the displayed integrality has $d_{\bar\lambda}=1$.
In this equality case, putting $L=(N-1)C_0$ gives the exact relation
between the two source actions:
\[
 \Psi^{\mathrm{un}}_{\bar\lambda}
   =m_{\exp(2\pi iL\bar\lambda/N)}
       \circ\Psi^{(N)}_{\bar\lambda}.
\]
Here $m_\eta(\xi,v)=(\eta\xi,v)$ has inverse
$m_{\eta^{-1}}$, and the multiplier lies in the deck group
$\mu_N^2$ of the relative toric power map.
\end{remark}






\begin{remark}[Monodromy of the marked fibre subgroup]
\hypertarget{tag-037Y}{}\label{wb:cusp-cover-global-scope}
{\normalfont\scriptsize tag 037Y}\par
Use the maps and compatible discs of
Definition~\ref{mapsIII:def:completed-cusp-root-data} and
Theorem~\ref{wb:cusp-cover-extension-theorem}.

For $N>1$, the connected regular-fibre cover with image
\[
 \Lambda_N=\mathbb Z\langle\widehat\gamma,\widehat u,
                       N\widehat w,N\widehat\delta\rangle
\]
does not extend as a fibrewise covering over the original regular
base $B^\circ$, and consequently does not extend over that same base
through both finite fillings. The counterclockwise based loop
$a_1^{-1}$ has point-return matrix $A_1$ by
Definition~\ref{mapsIII:def:three-filling-datum}, and
$A_1$ satisfies
\[
 A_1\widehat u=-\widehat w+\widehat\delta,\qquad
 S_N^{-1}A_1S_N\widehat u=(0,0,-1/N,1/N)^{\mathsf T}.
\]
Thus $A_1\widehat u\notin\Lambda_N$.  A connected fibrewise torus
cover over $B^\circ$ induces an inclusion of integral $H_1$ local
systems; its image lattice must be invariant under every base loop.
This necessary condition fails for $\Lambda_N$.

On the changed local cusp base the intertwining is instead
$M_0^NS_N=S_NM_0$, as proved in the theorem.  The obstruction just
given concerns the original base and its prescribed fibre subgroup.
\end{remark}





\section{Global torus quotients and their extension obstructions}
\subsection{Invariant characters and compactification}
\hypertarget{tag-Q006}{}\label{mapsIII:sec:regular-quotients}
\begin{definition}[Regular quotient data]
\hypertarget{tag-Q310}{}\label{mapsIII:def:regular-quotient-data}
{\normalfont\scriptsize tag Q310}\par
Let $f:X\to B$ be the three-fibre family of \cref{maps:src:B.5.2},
with $B^\circ=B\setminus\{p_0,p_1,p_2\}$ and
$J=f^{-1}(B^\circ)$. Use its period marking and the unrescaled
radial Euclidean fibre metrics of
\cref{wb:cusp-spec-radial-notation,wb:cusp-spec-native-quotient}.
Write
\[
 \Lambda=\mathbb Z\langle\widehat\gamma,\widehat u,
                   \widehat w,\widehat\delta\rangle,\qquad
 K_\Lambda=\mathbb Z\langle\widehat w,\widehat\delta\rangle,
\]
and define
\begin{equation}
 \varepsilon_\gamma
 (a\widehat\gamma+b\widehat u+c\widehat w+d\widehat\delta)=a,\qquad
 L_\Lambda=\ker\varepsilon_\gamma
 =\mathbb Z\langle\widehat u,\widehat w,\widehat\delta\rangle.
\end{equation}
\end{definition}

\begin{lemma}[Length metric after a closed-set collapse]
\hypertarget{tag-Q311}{}\label{mapsIII:lem:collapsed-subset-distance}
{\normalfont\scriptsize tag Q311}\par
For every compact length space $(E,d_E)$ and nonempty closed subset
$A\subset E$, collapsing $A$ to one point gives the quotient length metric
\[
 d_{E/A}([x],[y])
 =\min\{d_E(x,y),d_E(x,A)+d_E(y,A)\},\qquad
 d_{E/A}([x],*)=d_E(x,A).
\]
This metric induces the quotient topology. The quotient map identifies
only the points of $A$, so off the collapsed point its inverse takes
the unique representative in $E\setminus A$.
\end{lemma}
\begin{proof}
Indeed, a distance chain without a free jump in $A$ has length at
least $d_E(x,y)$.  A chain containing such jumps has initial and final
pieces of total length at least $d_E(x,A)+d_E(y,A)$; deleting the
intermediate pieces gives that upper bound through projected paths.
Taking infima proves the formula, including the distance to the
collapsed point.  The projected paths also prove that it is a length
metric, and compactness identifies its topology with the quotient topology.
\end{proof}

\begin{theorem}[Global circle quotients and compactification obstructions]
\hypertarget{tag-039U}{}\label{wb:cusp-spec-globalization-obstruction}
{\normalfont\scriptsize tag 039U}\par
For the data of \cref{mapsIII:def:regular-quotient-data}, the following
statements hold.

\smallskip
\noindent
\emph{(i) Regular group quotient.}
Over \(B^\circ\), the period family \(J=X^\circ=f^{-1}(B^\circ)\) is a
holomorphic family of complex Lie groups.  The formula
\begin{equation}
 \chi_\gamma\bigl([z,\Pi(z)x]\bigr)
 =\bigl([z],\varepsilon_\gamma(x)\bmod\mathbb Z\bigr)
 \hypertarget{tag-039R}{}\label{wb:cusp-spec-global-circle-map}
\end{equation}
defines a global smooth real-analytic fibrewise group epimorphism
\begin{equation}
 \chi_\gamma:J\longrightarrow B^\circ\times(\mathbb R/\mathbb Z).
 \hypertarget{tag-039S}{}\label{wb:cusp-spec-global-circle-quotient}
\end{equation}
Write
\begin{equation}
 \theta_\gamma:=\operatorname{pr}_{\mathbb R/\mathbb Z}\circ\chi_\gamma:
 J\longrightarrow\mathbb R/\mathbb Z
 \hypertarget{tag-039Q}{}\label{wb:cusp-spec-global-circle-coordinate}
\end{equation}
for its circle coordinate.  Thus \(\chi_\gamma=(f,\theta_\gamma)\).
Its kernel is the connected real three-torus family
\begin{equation}
 \mathcal H:=
 \frac{\mathfrak h_z^\circ\times
       ((L_\Lambda\otimes\mathbb R)/L_\Lambda)}{\Delta},
 \hypertarget{tag-039T}{}\label{wb:cusp-spec-global-kernel-family}
\end{equation}
where \(\Delta\) is the triangle deck group of
\(\mathfrak h_z^\circ\to B^\circ\).  It acts diagonally by its deck action
and by the restricted integral monodromy on \(L_\Lambda\).
The deck representation is the quotient homomorphism
\(\pi_1(B^\circ)\twoheadrightarrow\Delta\).  Equivalently, one can replace
\(\mathfrak h_z^\circ\) by the universal cover of \(B^\circ\) and
take the quotient by \(\pi_1(B^\circ)\).
The invariant-lattice calculation uses the image of the monodromy
representation, while finite base covers correspond to finite-index
subgroups of \(\pi_1(B^\circ)\).  Thus every fibre of \(\mathcal H\) is
\((L_\Lambda\otimes\mathbb R)/L_\Lambda\); the family need not be a
product.
The quotient is real and is not a holomorphic quotient of complex tori.

\smallskip
\noindent
\emph{(ii) Exact radial quotient metric.}
On the positive radial cusp, put
\(\chi_{\gamma,T}:=\theta_\gamma|_{F_T}\), where
\(F_T=\mathbb R^4/P_T\mathbb Z^4\). Then
\(\chi_\gamma|_{F_T}=(e^{-2\pi T},\chi_{\gamma,T})\), and
\begin{equation}
 \chi_{\gamma,T}([y])
 =\left\langle P_T^{-\mathsf T}e_\gamma,y\right\rangle
   \bmod\mathbb Z,\qquad e_\gamma=(1,0,0,0)^{\mathsf T}.
 \hypertarget{tag-03A2}{}\label{wb:cusp-spec-radial-circle-character}
\end{equation}
For the Riemannian quotient of the Euclidean metric, the target circle has
circumference
\begin{equation}
 \ell_T
 =\frac1{\lVert P_T^{-\mathsf T}e_\gamma\rVert}
 =\frac{D_T}{\sqrt{L_T^2+m_T^2}}
 =T+H-B_{\rm rad}+O(T^{-1}).
 \hypertarget{tag-039P}{}\label{wb:cusp-spec-global-circle-circumference}
\end{equation}
Here \(B_{\rm rad}=\operatorname{Im}b(0)\) denotes the bounded radial coefficient. In particular, \(\ell_T\to\infty\).

\smallskip
\noindent
\emph{(iii) Four copies and minimality.}
The source
\begin{equation}
 J^{[4]}=J\times_{B^\circ}J\times_{B^\circ}J\times_{B^\circ}J
 \hypertarget{tag-039N}{}\label{wb:cusp-spec-fourfold-fibre-product}
\end{equation}
is the fourfold fibre product over \(B^\circ\).  The product
map
\begin{equation}
 \chi_\gamma^{[4]}:
 J^{[4]}\longrightarrow B^\circ\times(\mathbb R/\mathbb Z)^4
 \hypertarget{tag-039M}{}\label{wb:cusp-spec-fourfold-circle-map}
\end{equation}
is the fibrewise product, explicitly
\[
 (x_1,x_2,x_3,x_4)\longmapsto
 (f(x_1),\theta_\gamma(x_1),\ldots,\theta_\gamma(x_4)).
\]
It
has kernel
\(\mathcal H^{[4]}=\mathcal H\times_{B^\circ}\cdots
\times_{B^\circ}\mathcal H\), whose fibres are
\(\mathbb T^{12}\).  Along the synchronized positive
radial cusp, its quotient fibre, with product quotient metric, is
\begin{equation}
 G_T=(\mathbb R/\ell_T\mathbb Z)^4,\qquad
 \operatorname{sys}(G_T)=\ell_T.
 \hypertarget{tag-039J}{}\label{wb:cusp-spec-four-circle-fibre}
\end{equation}
Consequently
\begin{equation}
 (G_T,0)\longrightarrow(\mathbb R^4,0)
 \hypertarget{tag-039K}{}\label{wb:cusp-spec-four-circle-limit}
\end{equation}
in pointed local-isometric, pointed smooth, and pointed
Gromov--Hausdorff topology.  If \(\ell_T>4R\), the radius-\(R\) ball
about zero is exactly Euclidean.

Equivalently, the external Cartesian product gives a quotient
\[
 J^4\longrightarrow(B^\circ)^4\times(\mathbb R/\mathbb Z)^4,
\]
given by \((x_i)_i\mapsto((f(x_i))_i,(\theta_\gamma(x_i))_i)\),
and \(J^{[4]}\) is its pullback along the diagonal
\(B^\circ\to(B^\circ)^4\).  At four independent cusp parameters its target
lattice is
\(\operatorname{diag}(\ell_{T_1},\ldots,\ell_{T_4})\mathbb Z^4\), so the
same pointed limit holds whenever
\(\min_i\ell_{T_i}\to\infty\).

For \(r\in\mathbb Z_{\geq1}\), among connected torus quotients obtained from \(r\) copies by a primitive
full-monodromy-invariant kernel, with the quotient norm induced by the
unrescaled Euclidean fibre metric, and whose quotient systole tends to
infinity on the radial cusp, the quotient rank is at most \(r\).  Thus four copies
are minimal for a four-dimensional all-expanding quotient, and
\cref{wb:cusp-spec-fourfold-circle-map} attains the bound.

The cusp-local two-copy map
\begin{equation}
 \pi_T\times\pi_T:F_T\times F_T\longrightarrow Q_T\times Q_T
 \hypertarget{tag-0398}{}\label{wb:cusp-spec-doubled-local-quotient}
\end{equation}
from \cref{wb:cusp-spec-native-quotient} also has pointed limit
\(\mathbb R^4\).  It is not a quotient over the full \(B^\circ\), however,
and no finite base cover makes it one.

\smallskip
\noindent
\emph{(iv) Finite fillings.}
At \(p_j\), let \(J'_j\to\Delta_j\) be the period family on a
linearizing root disc with coordinate \(s\), and let
\(t_j=s^{m_j}\) be the coordinate on \(D_j\).
Write \(N_j=J'_j/\langle\Phi_j\rangle\) for the finite filling defined
by the action below, and put
\[
 (m_1,v_1,\kappa_1)
 =\bigl(3,\widehat\gamma+2\widehat u-4\widehat w,1\bigr),
\]
\begin{equation}
 (m_2,v_2,\kappa_2)
 =\bigl(4,-\widehat\gamma-3\widehat u+3\widehat w,-1\bigr),\qquad
 \kappa_j=\varepsilon_\gamma(v_j).
\end{equation}
In the flat period marking, the finite filling generator is
\begin{equation}
 \Phi_j(s,x)=
 \left(e^{-2\pi i/m_j}s,A_jx+\frac{v_j}{m_j}\right).
 \hypertarget{tag-039A}{}\label{wb:cusp-spec-finite-affine-action}
\end{equation}
The character transforms by:
\begin{equation}
 \varepsilon_\gamma\left(A_jx+\frac{v_j}{m_j}\right)
 =\varepsilon_\gamma(x)+\frac{\kappa_j}{m_j}\pmod{\mathbb Z}.
 \hypertarget{tag-039B}{}\label{wb:cusp-spec-finite-character-shift}
\end{equation}
It descends equivariantly to the Seifert circle family
\begin{equation}
 C_j=
 \frac{\Delta_j\times(\mathbb R/\mathbb Z)}
 {(s,\theta)\sim
  (e^{-2\pi i/m_j}s,\theta+\kappa_j/m_j)}.
 \hypertarget{tag-0399}{}\label{wb:cusp-spec-finite-Seifert-target}
\end{equation}
The total space \(C_j\) is a smooth solid torus, the map
\(C_j\to D_j\) has a multiple central circle of order \(m_j\), and the
induced map \(N_j\to C_j\) is a smooth proper \(\mathbb T^3\)-bundle.
The target overlap induced by the logarithmic gluing map \(h_j\) is
the translation \(h_j^{S^1}\) defined in
\cref{wb:cusp-spec-finite-target-conjugacy,wb:cusp-spec-finite-overlap-identity}.

For four synchronized copies, the smooth simultaneous-root extension is
\begin{equation}
 \widetilde N_j^{[4]}
 =\frac{(J'_j)^{\times_{\Delta_j}4}}
 {\langle(\widetilde g_j^{\log})^{\times4}\rangle},
 \hypertarget{tag-039D}{}\label{wb:cusp-spec-finite-fourfold-normalization}
\end{equation}
with target
\begin{equation}
 C_j^{[4]}=
 \frac{\Delta_j\times(\mathbb R/\mathbb Z)^4}
 {(s,\theta_1,\ldots,\theta_4)\sim
  (e^{-2\pi i/m_j}s,
   \theta_1+\kappa_j/m_j,\ldots,\theta_4+\kappa_j/m_j)}.
 \hypertarget{tag-039E}{}\label{wb:cusp-spec-finite-fourfold-target}
\end{equation}
The map \(\widetilde N_j^{[4]}\to C_j^{[4]}\) is a smooth proper
\(\mathbb T^{12}\)-bundle.  There is a canonical finite map
\begin{equation}
 \nu_j:\widetilde N_j^{[4]}
 \longrightarrow N_j\times_{D_j}N_j\times_{D_j}N_j\times_{D_j}N_j
 \hypertarget{tag-039F}{}\label{wb:cusp-spec-finite-normalization-map}
\end{equation}
which is an isomorphism over \(D_j^*\), is the normalization of the
reduced ordinary filled fibre product, and has degree \(m_j^3\) on the
reduced central fibre.  These two degrees are \(27\) and \(64\).  Gluing
these two local normalizations to \(J^{[4]}\) by
\((h_j)^{\times4}\), and gluing their targets to the product
\(B^\circ\times(\mathbb R/\mathbb Z)^4\) by
\((h_j^{S^1})^{\times4}\), gives the corresponding normalization and
Seifert target over \(B\setminus\{p_0\}\).

\smallskip
\noindent
\emph{(v) Cusp nonextension.}
The circle coordinate \(\theta_\gamma\) has no continuous extension
through \(W\). Consequently \(\chi_\gamma\) has no continuous
extension to \(N_0\to D_0\times(\mathbb R/\mathbb Z)\), and
\(\chi_\gamma^{[4]}\) has no continuous extension to the compact
fourfold fibre product, or to \(X^4\), with target
\(B\times(\mathbb R/\mathbb Z)^4\) and the stated regular-fibre
restriction. No map on \(X^{\wedge4}\cong S^{24}\) induces that
restriction after precomposition with \(X^4\to X^{\wedge4}\).
Thus the four-copy quotient is defined over \(B^\circ\), and its
finite-filling normalization over \(B\setminus\{p_0\}\).
The integral class \([\theta_\gamma]\in H^1(N_0\setminus W;\mathbb Z)\)
extends over \(N_0\), with a representative whose restriction is
homotopic to \(\theta_\gamma\). The specified quotient map and its
expanding circle metric have no such extension.

\smallskip
\noindent
\emph{(vi) Products, smashes, and sphere constructions.}
The homotopy-sphere recognition in \cref{maps:src:B.5.2}, together with
$\Theta_6=0$ \cite[remark on p.~505; Theorem~5.1 and table on p.~512]{manKM63},
gives a based diffeomorphism $h:X\to S^6$. Its Cartesian powers give
\begin{equation}
 X^2\cong_{\mathrm{diff}}S^6\times S^6\not\simeq S^{12},\qquad
 H_6(X^2;\mathbb Z)\cong\mathbb Z^2,
 \hypertarget{tag-03A1}{}\label{wb:cusp-spec-product-not-sphere}
\end{equation}
and
\begin{equation}
 \begin{aligned}
 X^4&\cong_{\mathrm{diff}}(S^6)^4\not\simeq S^{24},\\
 H_6(X^4;\mathbb Z)&\cong\mathbb Z^4,&
 H_{12}(X^4;\mathbb Z)&\cong\mathbb Z^6,&
 H_{18}(X^4;\mathbb Z)&\cong\mathbb Z^4.
 \end{aligned}
 \hypertarget{tag-039L}{}\label{wb:cusp-spec-four-product-not-sphere}
\end{equation}
The diagonal pullback \(J^{[4]}\subset J^4\) has real dimension \(18\)
and codimension six.  Choosing a base point gives based homeomorphisms
\begin{equation}
X\wedge X\cong S^{12},\qquad X^{\wedge4}\cong S^{24}.
\hypertarget{tag-03A4}{}\label{wb:cusp-spec-smash-spheres}
\end{equation}
The corresponding source maps are the Cartesian-to-smash collapses.
For $A_{\rm tor}=(\mathbb R/\mathbb Z)^2$ or $\mathbb R/\mathbb Z$, the
projection $[r,z,a_{\rm cone}]\mapsto a_{\rm cone}$ from the positive part of the cone over
$S^1\times A_{\rm tor}$ recovers the fixed quotient coordinate and has no
continuous extension to the vertex. Neither \(S^{12}\) nor \(S^{24}\), with any differentiable structure,
admits an almost-complex structure
\cite[\S15, Proposition~15.1 and the following remark,
p.~434]{BorelSerre1953}; hence neither admits a complex structure.
For the fixed-target quotient maps,
\(\pi_T\times\pi_T\) is a map on \(F_T^2\), and it fails to factor
through \(F_T\wedge F_T\); \(\chi_{\gamma,T}^{[4]}\) is a map on
\(F_T^4\), and it fails to factor through \(F_T^{\wedge4}\).
No map on \(X^2\) or \(X^4\) has the stated fibre restriction, and
precomposition with the Cartesian-to-smash collapse excludes such
extensions from \(X\wedge X\) or \(X^{\wedge4}\) as well.

For \(C_T=\mathbb R/\ell_T\mathbb Z\) based at zero, let
\(\vartheta_T:\mathbb R/\mathbb Z\to C_T\),
\([\theta]\mapsto[\ell_T\theta]\), be the arclength isometry,
with inverse $[u]\mapsto[u/\ell_T]$. The domain metric is
$\ell_T^2d\theta^2$.
Postcomposition with the target fat-wedge collapse gives
\begin{equation}
 F_T^4\xrightarrow{\vartheta_T^{\times4}\circ\chi_{\gamma,T}^{[4]}}C_T^4
 \longrightarrow C_T^{\wedge4}\cong S^4.
 \hypertarget{tag-03A5}{}\label{wb:cusp-spec-target-smash-map}
\end{equation}
This map does factor through \(F_T^{\wedge4}\).  With the quotient length
metric and base point
\([\ell_T/2,\ell_T/2,\ell_T/2,\ell_T/2]\), its pointed limit is
\(\mathbb R^4\), with every radius-\(R\) ball exactly Euclidean once
\(\ell_T>4R\).  The metric is singular at the collapsed fat wedge, and
the map still does not extend to \(X^4\) or \(X^{\wedge4}\).
\end{theorem}

\begin{proof}

 
\medskip\noindent\emph{Invariant lattices and the regular group law.}\par
The monodromy matrices used in \cref{maps:src:B.5.2} give
\[
 (M_0-I)(a,b,c,d)^{\mathsf T}=(0,0,b,-a)^{\mathsf T},
\]
so
\begin{equation}
 K_\Lambda=\ker(M_0-I)=\operatorname{im}(M_0-I).
 \hypertarget{tag-0394}{}\label{wb:cusp-spec-cusp-kernel}
\end{equation}
They also give
\begin{equation}
 \begin{aligned}
 A_1\widehat u&=-\widehat w+\widehat\delta,&
 A_1\widehat w&=\widehat u-\widehat w,&
 A_1\widehat\delta&=\widehat\delta,\\
 A_2\widehat u&=\widehat w,&
 A_2\widehat w&=-\widehat u+\widehat\delta,&
 A_2\widehat\delta&=\widehat\delta.
 \end{aligned}
 \hypertarget{tag-039X}{}\label{wb:cusp-spec-invariant-flag-calculation}
\end{equation}
Thus \(L_\Lambda\) is invariant.  It is the smallest invariant lattice
containing \(K_\Lambda\), because
\(\widehat u=A_1\widehat w+\widehat w\).  Both lattices are primitive.

For a row covector \((a,b,c,d)\), invariance under \(A_1\) gives
\(b=2c\) and \(d=3c\), whereas invariance under \(A_2\) gives
\(b=c\) and \(d=2c\).  Hence \(b=c=d=0\), and
\begin{equation}
 \operatorname{Hom}(\Lambda,\mathbb Z)^\Delta
 =\mathbb Z\varepsilon_\gamma.
 \hypertarget{tag-039W}{}\label{wb:cusp-spec-invariant-character-lattice}
\end{equation}

The family over the cover is
\((\mathfrak h_z^\circ\times\mathbb C^2)/\Lambda\), and the action used
to form \(J\) is
\[
 [z,\zeta]\longmapsto[gz,R_g(z)\zeta].
\]
It is complex linear in every fibre and fixes the zero section.
Fibrewise addition, inversion, and zero are therefore equivariant and
descend to \(J\).  In the flat real-analytic marking
\[
 \Xi(z,x+\Lambda)=[z,\Pi(z)x],
\]
the period identity \(R_g(z)\Pi(z)=\Pi(gz)A_g\) makes the action
\[
 (z,x+\Lambda)\longmapsto(gz,A_gx+\Lambda).
\]
Equation \eqref{wb:cusp-spec-invariant-character-lattice} proves that
\cref{wb:cusp-spec-global-circle-map} is well defined.  Primitivity of
\(\varepsilon_\gamma\) proves surjectivity and connectedness of the
kernel. The induced action on \(\Lambda/L_\Lambda\simeq\mathbb Z\) is
trivial, so the target is the product circle bundle. The free covering
\(\mathfrak h_z^\circ\to B^\circ\) gives the exact sequence
\[
 \pi_1(\mathfrak h_z^\circ)\longrightarrow\pi_1(B^\circ)
 \longrightarrow\Delta\longrightarrow1.
\]
The middle group is free on the two puncture loops, whose images have
orders three and four. The cube of the first free generator is a
nontrivial element of the kernel. Let
$p:\widetilde B^\circ\to\mathfrak h_z^\circ$ be the lift of the
universal base cover. In the flat marking the comparison of the
two family quotients is
\[
 [\widetilde z,x]\longmapsto[p(\widetilde z),x].
\]
Its inverse chooses any lift $\widetilde z$ of $z$ and sends
$[z,x]$ to $[\widetilde z,x]$. Two such lifts differ by the
kernel of $\pi_1(B^\circ)\to\Delta$, which acts trivially on the
marked fibre. The inverse is therefore independent of that choice;
local covering charts prove continuity and real analyticity in both
directions. These maps identify the same group family.
The kernel family has nonidentity monodromy, since
\(A_1\widehat w=\widehat u-\widehat w\); a product group family would
conjugate this action on \(H_1\) to the identity. Finally a complex
structure on the real tangent line of the target circle would be a real
scalar with square \(-1\). Such a scalar does not exist, so this
fibrewise quotient cannot be holomorphic.

\medskip\noindent\emph{The quotient metric and the four-copy limit.}\par
The first column of
\cref{wb:cusp-spec-full-inverse-transpose} is
\begin{equation}
 P_T^{-\mathsf T}e_\gamma
 =\left(0,\frac{m_T}{D_T},0,-\frac{L_T}{D_T}\right)^{\mathsf T}.
 \hypertarget{tag-0393}{}\label{wb:cusp-spec-circle-dual-vector}
\end{equation}
Write \(a_T=P_T^{-\mathsf T}e_\gamma\). Arclength projection to the
orthogonal line is \(y\mapsto\langle a_T,y\rangle/\|a_T\|\).
Its period subgroup is \(\|a_T\|^{-1}\mathbb Z\), since
\(\langle a_T,P_Tn\rangle=n_1\). Thus the quotient metric on
\(\mathbb R/\mathbb Z\) is \(\ell_T^2d\theta^2\), and
\([\theta]\mapsto[\ell_T\theta]\) is its isometry to
\(\mathbb R/\ell_T\mathbb Z\). Equivalently,
\[
 P_T(L_\Lambda\otimes\mathbb R)
 =\operatorname{span}_{\mathbb R}\{e_1,e_3,v_{u,T}\}
\]
has unit normal
\[
 \nu_T=\frac{(0,m_T,0,-L_T)}{\sqrt{L_T^2+m_T^2}},
\]
and
\[
 \operatorname{dist}
 \bigl(v_{\gamma,T},P_T(L_\Lambda\otimes\mathbb R)\bigr)
 =\left|\langle v_{\gamma,T},\nu_T\rangle\right|
 =\frac{L_Tq_T+6m_T^2}{\sqrt{L_T^2+m_T^2}}
 =\ell_T.
\]
The expansions in
\cref{wb:cusp-spec-radial-notation,wb:cusp-spec-holomorphic-errors} give
\[
 D_T=T^2+(2H-B_{\rm rad})T+O(1),\qquad
 \sqrt{L_T^2+m_T^2}=T+H+O(T^{-1}),
\]
which proves \cref{wb:cusp-spec-global-circle-circumference}.

After rescaling the circle coordinate by \(\ell_T\), the product target is
\(\mathbb R^4/\ell_T\mathbb Z^4\).  If
\(x,y\in B_R(0)\) and \(0\ne\lambda\in\ell_T\mathbb Z^4\), then
\[
 \lVert x-y+\lambda\rVert
 \ge\lVert\lambda\rVert-\lVert x-y\rVert
 \ge\ell_T-2R>2R\ge\lVert x-y\rVert
\]
when \(\ell_T>4R\).  Hence the covering projection is an exact isometry
on \(B_R(0)\). Every point of the quotient radius-\(R\) ball has a
representative in this Euclidean ball, by choosing a shortest lattice
representative. Hence the entire balls are isometric, proving all three
convergences in \cref{wb:cusp-spec-four-circle-limit} with the indicated
radial metrics. The external product is defined over \((B^\circ)^4\),
and the synchronized product is its pullback along the base diagonal.
At independent parameters, the shortest vector of
\(\operatorname{diag}(\ell_{T_1},\ldots,\ell_{T_4})\mathbb Z^4\)
has length \(\min_i\ell_{T_i}\); the same ball argument proves the
independent-parameter limit.

\medskip\noindent\emph{Minimality and the finite-base-cover obstruction.}\par
Let \(N\subset\Lambda^{\oplus r}\) be a primitive invariant kernel whose
connected quotient has radial systole tending to infinity.  Every element
of \(K_\Lambda^{\oplus r}\) has uniformly bounded physical length, since
\[
 P_T\widehat w=e_1,\qquad P_T\widehat\delta=e_3.
\]
If such an element survived in the quotient lattice, its quotient norm
would be no greater than its source norm and would give a nonzero uniformly
bounded period. Here a marked vector has zero quotient norm precisely
when it lies in \(N\): orthogonal projection has kernel
\(P_T(N\otimes\mathbb R)\), and primitivity gives
\((N\otimes\mathbb R)\cap\Lambda^{\oplus r}=N\).
Therefore \(K_\Lambda^{\oplus r}\subset N\).
Invariance and
\cref{wb:cusp-spec-invariant-flag-calculation} then force
\(L_\Lambda^{\oplus r}\subset N\), and
\[
 \operatorname{rank}(\Lambda^{\oplus r}/N)\le4r-3r=r.
\]
The quotient is related to the maximal one by the homomorphisms
\[
 \Lambda^{\oplus r}\longrightarrow(\Lambda/L_\Lambda)^{\oplus r}
 \longrightarrow\Lambda^{\oplus r}/N
\]
and their induced torus maps. This proves the lower bound \(r\ge4\)
for rank four, while \(N=L_\Lambda^{\oplus4}\) proves sharpness.

For the finite-cover assertion, a connected fibrewise torus quotient
induces the exact sequence of integral first-homology local systems
\[
 0\longrightarrow\mathscr K\longrightarrow\mathscr L
 \longrightarrow\mathscr Q\longrightarrow0,
\]
where the terms are the first homology of the kernel, source and target
fibres. Naturality under transport around a base loop makes the kernel
lattice invariant, including when the target monodromy is nontrivial.
For the native quotient its fibre is \(K_\Lambda\), and for two copies
it is \(K_\Lambda\oplus K_\Lambda\). Put
\begin{equation}
 H_*=A_2A_1
 =\begin{pmatrix}
 1&0&0&0\\
 6&1&1&0\\
 0&0&1&0\\
 -5&0&-1&1
 \end{pmatrix}=I+N_*.
 \hypertarget{tag-039Z}{}\label{wb:cusp-spec-no-finite-cover-matrix}
\end{equation}
Direct multiplication gives \(N_*^2=0\), and hence
\begin{equation}
 H_*^n\widehat w
 =\widehat w+n(\widehat u-\widehat\delta),\qquad
 H_*^n\widehat\delta=\widehat\delta.
 \hypertarget{tag-03A0}{}\label{wb:cusp-spec-no-finite-cover-orbit}
\end{equation}
Thus \(H_*^nK_\Lambda\ne K_\Lambda\) for every \(n\ne0\).
Every finite-index subgroup of the base fundamental group contains a
positive power of the loop with monodromy \(H_*\).  Such a subgroup cannot
preserve \(K_\Lambda\), proving that no finite base cover globalizes the
native two-torus quotient. On two copies the diagonal action sends
\((\widehat w,0)\) to
\((\widehat w+n(\widehat u-\widehat\delta),0)\), outside
\(K_\Lambda\oplus K_\Lambda\) for every \(n\ne0\).
The same local-system obstruction proves the two-copy assertion.

\medskip\noindent\emph{Affine finite fillings and logarithmic overlaps.}\par
Equation \eqref{wb:cusp-spec-invariant-character-lattice} gives
\(\varepsilon_\gamma A_j=\varepsilon_\gamma\), so
\cref{wb:cusp-spec-finite-character-shift} follows.  Define
\[
 \overline\Phi_j(s,\theta)
 =\left(e^{-2\pi i/m_j}s,\theta+\frac{\kappa_j}{m_j}\right).
\]
Then \((s,x)\mapsto(s,\varepsilon_\gamma(x))\) is
\((\Phi_j,\overline\Phi_j)\)-equivariant; denote the descended map by
\(\overline\chi_j:N_j\to C_j\).  A nonidentity power which fixes
the base must lie over \(s=0\), where it translates the circle by
\(k\kappa_j/m_j\).  Since \(\gcd(\kappa_j,m_j)=1\), this is nonintegral
for \(0<k<m_j\). The target action is free. Since \(A_j^{m_j}=I\)
and \(A_jv_j=v_j\), the \(m_j\)-th power of the source action translates
by \(v_j\in\Lambda\), hence is the identity on the torus.
Equivariance with the free target action proves that its lower powers
have no fixed point.

For an explicit solid-torus chart, put
\(\lambda_j=e^{2\pi i/m_j}\), \(u=e^{2\pi i\theta}\), and choose
\(r_j\in\mathbb Z\) with \(r_j\kappa_j\equiv1\pmod {m_j}\).  Then
\[
 w=su^{r_j},\qquad U=u^{m_j}
\]
are invariant and identify \(C_j\) real-analytically with a solid torus.
The inverse chooses an \(m_j\)-th root \(u\) of \(U\) and sets
\(s=wu^{-r_j}\); changing the root is exactly a power of the displayed
cyclic action. In these coordinates
\[
 t_j=s^{m_j}=w^{m_j}U^{-r_j},
\]
which proves the central multiplicity.  Equivariance and the primitive
kernel give the \(\mathbb T^3\)-bundle \(N_j\to C_j\): local lifts
of the free finite target covering reduce it to the product of the
primitive character quotient with a disc. Its compact fibre proves
properness.

For the logarithmic gluing, let
\(\Pi_{\mathbb R}(z):\Lambda\otimes\mathbb R\to\mathbb C^2\) be the
underlying real isomorphism.  On a logarithm branch \(\ell\), define
\(b_j^\ell(z)\) uniquely by
\begin{equation}
 \Pi_{\mathbb R}(z)b_j^\ell(z)
 =\frac{\ell(z)}{2\pi i}\Pi(z)v_j,\qquad
 \beta_j(z)=\varepsilon_\gamma(b_j^\ell(z))\bmod\mathbb Z.
 \hypertarget{tag-039I}{}\label{wb:cusp-spec-finite-target-conjugacy}
\end{equation}
Changing \(\ell\) by \(2\pi in\) adds \(nv_j\) to \(b_j^\ell\), hence
adds \(n\kappa_j\in\mathbb Z\) to its
\(\varepsilon_\gamma\)-value.  Thus \(\beta_j\) is branch independent.
In flat coordinates the holomorphic overlap \(h_j\), from the logarithmic
filling to the regular family, is translation by \(b_j^\ell\). The
clockwise root generator has compatible branch
\(\ell'(g_jz)=\ell(z)-2\pi i/m_j\). The period equivariance and
\(A_jv_j=v_j\) give
\[
 b_j^{\ell'}(g_jz)=A_jb_j^\ell(z)-v_j/m_j.
\]
Applying \(\varepsilon_\gamma\) gives
\begin{equation}
 \beta_j(g_jz)=\beta_j(z)-\frac{\kappa_j}{m_j}\pmod{\mathbb Z}.
 \hypertarget{tag-039H}{}\label{wb:cusp-spec-finite-target-cocycle}
\end{equation}
Thus the lifted translation \((z,\theta)\mapsto(z,\theta+\beta_j(z))\)
conjugates \(\overline\Phi_j\) to the linear target action
\((z,\theta)\mapsto(g_jz,\theta)\). Its descended map is
\[
 h_j^{S^1}:C_j|_{D_j^*}\longrightarrow D_j^*\times(\mathbb R/\mathbb Z),
 \qquad [z,\theta]\longmapsto([z],\theta+\beta_j(z)).
\]
The inverse chooses a root lift and subtracts \(\beta_j\); the cocycle
identity makes that inverse independent of the lift. The overlap identity is
\begin{equation}
 \chi_\gamma\circ h_j=h_j^{S^1}\circ\overline\chi_j.
 \hypertarget{tag-039G}{}\label{wb:cusp-spec-finite-overlap-identity}
\end{equation}
Let \(\mathcal J_z\) be the complex structure in flat real coordinates. Writing
\(\ell=\log|s|+i\arg s\) displays the full translation:
\[
 \varepsilon_\gamma(b_j^\ell)
 =\frac{\kappa_j\arg s}{2\pi}
  -\frac{\log|s|}{2\pi}\varepsilon_\gamma(\mathcal J_zv_j).
\]
The radial and angular terms both enter the target overlap.

The diagonal action on
\((J'_j)^{\times_{\Delta_j}4}\) is free, and the same coprimality argument
proves that \(C_j^{[4]}\) is smooth.  Every fibre of the resulting map is
a \(\mathbb T^{12}\)-torsor; over \(D_j^*\) its structure torus is the
fourfold fibre product of \(\mathcal H\). Local lifts of the target's
free cyclic covering give bundle charts, and compact fibres give
properness. The canonical map is
\[
 [s,x_1,x_2,x_3,x_4]_{\rm diag}
 \longmapsto([s,x_1],[s,x_2],[s,x_3],[s,x_4]).
\]
Over \(t_j\ne0\), choose one root \(s\) of \(t_j\), adjust each representative
to that root by the cyclic action, and then divide by the residual diagonal
action. This is its holomorphic inverse.
In local holomorphic coordinates along the central fibre, the ordinary filled
fibre product has equations
\begin{equation}
 s_i^{m_j}=s_1^{m_j},\qquad i=2,3,4.
 \hypertarget{tag-039C}{}\label{wb:cusp-spec-finite-fibre-product-branches}
\end{equation}
With \(m=m_j\) and \(z\) denoting the eight fibre coordinates, its local
ring is
\[
 \mathcal A_{\rm loc}=\frac{\mathbb C\{s_1,s_2,s_3,s_4,z\}}
 {(s_2^m-s_1^m,s_3^m-s_1^m,s_4^m-s_1^m)}.
\]
It is free over \(\mathcal R_{\rm loc}=\mathbb C\{s_1,z\}\), with basis
\(s_2^as_3^bs_4^c\), \(0\leq a,b,c<m\). After inverting \(s_1\), it
splits into the \(m^3\) branches \(s_i=\epsilon_i s_1\),
\(\epsilon_i\in\mu_m\). The evaluation map \(\mathcal A_{\rm loc}\to\prod_{\mu_m^3}\mathcal R_{\rm loc}\)
is injective, since its kernel vanishes after inverting \(s_1\) and
\(\mathcal A_{\rm loc}\) is \(s_1\)-torsion-free. Hence \(\mathcal A_{\rm loc}\) is reduced. The product is
finite over \(\mathcal A_{\rm loc}\), is normal and has the same total ring of fractions;
it is therefore the integral closure. The simultaneous-root model has
exactly these separated smooth branches, so its canonical map is the
normalization. Globally this map is proper because both spaces are
proper over \(D_j\), and its finite fibres also prove finiteness.
On the reduced central fibre it is
\[
 (A_0^{(j)})^4/G_{\rm diag}\longrightarrow(A_0^{(j)}/G)^4.
\]
Its relative root choices form \(G^4/G_{\rm diag}\), where
\(G=\mathbb Z/m_j\); the degree is \(m_j^3\), namely \(27\) and \(64\).

On \(D_j^*\), the source identification is \((h_j)^{\times4}\), and
\cref{wb:cusp-spec-finite-overlap-identity} identifies the target overlap
with \((h_j^{S^1})^{\times4}\).  These componentwise overlap maps satisfy the cocycle identities.
They glue the two local normalizations and Seifert targets to
\(J^{[4]}\) and to
\(B^\circ\times(\mathbb R/\mathbb Z)^4\), respectively.  Equivalently,
the source gluing follows from uniqueness of normalization.

The fixed product-target trivialization does not extend over \(X\), even if
base winding is allowed. In the attachment presentation of
\cref{maps:src:B.5.2}, let a homomorphism to \(\mathbb Z\) restrict to
\(n_{\rm att}\varepsilon_\gamma\) on \(\Lambda\) and take values \(a_{\rm att},b_{\rm att}\) on
\(\widehat\rho_1,\widehat\rho_2\).  The three attachment relations give
\[
 3a_{\rm att}=n_{\rm att},\qquad4b_{\rm att}=-n_{\rm att},\qquad a_{\rm att}+b_{\rm att}=0.
\]
The first two give \(a_{\rm att}=n_{\rm att}/3\), \(b_{\rm att}=-n_{\rm att}/4\), while the third gives
\(n_{\rm att}/12=0\).  Hence \(n_{\rm att}=0\).  Thus a fixed-circle map on \(X\) cannot restrict to a nonzero multiple
of the regular character, even after changing its winding on the base.

\medskip\noindent\emph{The explicit toric cusp obstruction.}\par
On the positive radial ray, write
\[
 Z_T=iTB_0+C(t_c),\qquad
 B_0=\begin{pmatrix}0&1\\-1&0\end{pmatrix},\qquad
 t_c=e^{-2\pi T}.
\]
For \(a_{\rm tor}=(a_{1,\rm tor},a_{2,\rm tor})^{\mathsf T}\in\mathbb R^2\), take the point of the
regular fibre represented by \(Z_Ta_{\rm tor}\), with zero short-period coordinates,
and exponentiate it to the dense torus of the cusp model:
\[
 x_T(a_{\rm tor})=\exp(2\pi iZ_Ta_{\rm tor})\in(\mathbb C^*)^2.
\]
Its logarithmic position is
\begin{equation}
 y_T(a_{\rm tor})=
 \frac{(\log|x_{T,1}(a_{\rm tor})|,\log|x_{T,2}(a_{\rm tor})|)}{\log|t_c|}
 =B_0a_{\rm tor}+\frac{\operatorname{Im}C(t_c)a_{\rm tor}}{T}.
 \hypertarget{tag-0397}{}\label{wb:cusp-spec-cusp-nonextension-position}
\end{equation}
The circle character has the constant value
\begin{equation}
 \chi_{\gamma,T}(x_T(a_{\rm tor}))=a_{1,\rm tor}\bmod\mathbb Z.
 \hypertarget{tag-0395}{}\label{wb:cusp-spec-cusp-nonextension-character}
\end{equation}

Let \(e_{1,\rm fan},e_{2,\rm fan}\) be the ordered standard basis
of this two-dimensional fan lattice \(\mathbb Z^2\). Use the logarithmic
base coordinate \(s_{\rm cusp}\), with \(t=e^{2\pi i s_{\rm cusp}}\);
on the positive radial ray choose \(s_{\rm cusp}=iT\), so \(t=t_c\).
The prescribed fan consists of the integer translates of
\(\sigma=\operatorname{conv}(0,e_{1,\rm fan},e_{2,\rm fan})\) and
\(\sigma'=\operatorname{conv}(e_{1,\rm fan},e_{2,\rm fan},e_{1,\rm fan}+e_{2,\rm fan})\).
In the native long marking, the period matrix is \(s_{\rm cusp}I+C(t)B_0^{-1}\):
the coordinate change is $a_{\rm native}=B_0a_{\rm tor}$, with inverse
$a_{\rm tor}=B_0^{-1}a_{\rm native}$ and with the physical complex coordinate fixed. Equivalently
\[
 [s_{\rm cusp}I+C(t)B_0^{-1}\mid I]\operatorname{diag}(B_0,I)
 =[s_{\rm cusp}B_0+C(t)\mid I].
\]
The long-period action translates the fan by \(B_0\mathbb Z^2=\mathbb Z^2\).
The chart of \(\sigma\) has coordinates
\begin{equation}
 z_0=\frac{t_c}{x_1x_2},\qquad z_1=x_1,\qquad z_2=x_2.
 \hypertarget{tag-0396}{}\label{wb:cusp-spec-cusp-nonextension-chart}
\end{equation}
If \(B_0a_{\rm tor}=(y_1,y_2)\) is in the interior of \(\sigma\), then
\(y_1>0\), \(y_2>0\), and \(1-y_1-y_2>0\).  For all sufficiently large
\(T\), the three absolute values in
\cref{wb:cusp-spec-cusp-nonextension-chart} are
\[
 |t_c|^{1-y_{T,1}-y_{T,2}},\qquad
 |t_c|^{y_{T,1}},\qquad |t_c|^{y_{T,2}},
\]
and all tend to zero.  Every such path converges to the same toric fixed
point and hence to the same triple point of \(W\).

Evaluate the same path family at
\[
 a_{\rm tor}=\left(-\frac14,\frac14\right),\qquad
 a_{\rm tor}'=\left(-\frac12,\frac14\right).
\]
Then
\[
 B_0a_{\rm tor}=\left(\frac14,\frac14\right),\qquad
 B_0a_{\rm tor}'=\left(\frac14,\frac12\right)
\]
are both in the interior of \(\sigma\).  The paths converge to the same
point of \(W\), whereas
\[
 \chi_{\gamma,T}(x_T(a_{\rm tor}))=\frac34,\qquad
 \chi_{\gamma,T}(x_T(a_{\rm tor}'))=\frac12.
\]
No continuous extension of the specified map \(\chi_\gamma\) exists.
Choosing identical marked radial paths in the other three factors and
using these two paths in one factor proves the same assertion for
\(\chi_\gamma^{[4]}\).

The radial family of translations by \(\widehat u/2\) in \(\mathcal H\)
fails to extend as well.  Take two paths
whose limiting logarithmic positions are
\[
 y_{\rm tr}=\left(\frac18,\frac18\right),\qquad
 y_{\rm tr}'=\left(\frac38,\frac38\right).
\]
Both approach the fixed point of
\(\operatorname{conv}(0,e_{1,\rm fan},e_{2,\rm fan})\).  Translation by
\(\widehat u/2\in\mathcal H_T\) adds the exact displacement
\[
 \frac12B_0e_{2,\rm fan}+
 \frac1{2T}\operatorname{Im}C(t_c)e_{2,\rm fan}
 \longrightarrow\left(\frac12,0\right)
\]
to logarithmic position.  The translated limits are
\[
 \left(\frac58,\frac18\right)
 \quad\hbox{and}\quad
 \left(\frac78,\frac38\right).
\]
The first lies in the interior of
\(\operatorname{conv}(0,e_{1,\rm fan},e_{2,\rm fan})\), while the second lies in the interior
of \(\operatorname{conv}(e_{1,\rm fan},e_{2,\rm fan},e_{1,\rm fan}+e_{2,\rm fan})\).  The two triangles have barycentres \((1/3,1/3)\) and \((2/3,2/3)\),
whose difference is not integral; no deck translation identifies them.
They give the two distinct triple points of \(W\).  Thus one regular-fibre translation
sends two paths with the same limit to paths with different limits; it has
no continuous extension to \(N_0\).

The cusp calculation in \cref{maps:src:B.5.2} gives the inclusion homomorphism
\[
 \Lambda\rtimes_{M_0}\mathbb Z\longrightarrow\pi_1(N_0)=\Lambda/K_\Lambda,
 \qquad(\lambda,n)\longmapsto[\lambda].
\]
Since \(\varepsilon_\gamma\) kills \(K_\Lambda\), the class
\([\theta_\gamma]\) extends to an integral \(H^1\)-class on \(N_0\).
Represent it by \(h_{\rm cusp}:N_0\to S^1\). The difference
\(h_{\rm cusp}|_{N_0\setminus W}-\theta_\gamma\) lifts to a real function \(\alpha_{\rm htpy}\),
and \(\theta_\gamma+u\alpha_{\rm htpy}\), \(0\leq u\leq1\), is a homotopy from
\(\theta_\gamma\) to \(h_{\rm cusp}|_{N_0\setminus W}\). The explicit paths prove that the
specified map has no continuous extension. The metric
\(\ell_T^2d\theta^2\), with \(\ell_T\to\infty\), likewise has no finite
continuous extension on the fixed circle at \(t=0\).

The phase torus \((K_\Lambda\otimes\mathbb R)/K_\Lambda\) acts on the
chart by
\[
 (\eta_1,\eta_2)\cdot(z_0,z_1,z_2)
 =((\eta_1\eta_2)^{-1}z_0,\eta_1z_1,\eta_2z_2).
\]
These toric actions agree on chart overlaps and commute with the
long-period multipliers, so they descend to \(N_0\). At a triple point
the whole phase torus is a stabilizer. This gives the extending action
and its exact relation to the short-period subgroup; the preceding
translated paths exclude extension of the \(\widehat u/2\)-translation.

\medskip\noindent\emph{Products and smash quotients.}\par
The table of \cite[p.~512]{manKM63} gives \(\Theta_6/bP_7=0\).
Theorem 5.1 there gives a contractible filling for an element of \(bP_7\);
removing an interior ball gives an h-cobordism with the standard sphere.
The diffeomorphism conclusion recorded on p.~505 therefore gives
\(\Theta_6=0\). Together with \cref{maps:src:B.5.2}, this supplies the
diffeomorphism $h_0:X\to S^6$.
Write $b_d$ for the north pole of $S^d$, indexed by its sphere dimension $d$.
For a fixed fibre inclusion choose
its zero $x_0$ as basepoint, and let ${p_{\rm sph}}=h_0(x_0)$. Use $b_6$, the
north pole of $S^6\subset\mathbb R^7$. The following rotation sends
${p_{\rm sph}}$ to $b_6$. If ${p_{\rm sph}}\ne\pm b_6$, put $c_{\rm rot}=\langle {p_{\rm sph}},b_6\rangle$,
$s_{\rm rot}=\sqrt{1-c_{\rm rot}^2}$ and $w_{\rm rot}=(b_6-c_{\rm rot}{p_{\rm sph}})/s_{\rm rot}$. On the orthonormal plane
$({p_{\rm sph}},w_{\rm rot})$ use the matrix
\[
 O_{\rm rot}|_{\langle {p_{\rm sph}},w_{\rm rot}\rangle}
   =\begin{pmatrix}c_{\rm rot}&-s_{\rm rot}\\s_{\rm rot}&c_{\rm rot}\end{pmatrix},
 \qquad O_{\rm rot}|_{\langle {p_{\rm sph}},w_{\rm rot}\rangle^\perp}=I.
\]
For ${p_{\rm sph}}=b_6$ take $O_{\rm rot}=I$; for ${p_{\rm sph}}=-b_6$ take $O_{\rm rot}=-I$ on
$\langle {p_{\rm sph}},w_{\rm rot}\rangle$ for any unit $w_{\rm rot}\perp {p_{\rm sph}}$, and $I$ on its
orthogonal complement. In every case $O_{\rm rot}{p_{\rm sph}}=b_6$ and $O_{\rm rot}^{-1}=O_{\rm rot}^{\mathsf T}$.
Thus $h=O_{\rm rot}h_0$ is based, with inverse $h_0^{-1}O_{\rm rot}^{\mathsf T}$.
In dimension $d$, use the same pole $b_d$.
Stereographic projection identifies $(S^d,b_d)$ with
$(\mathbb R^d)^+$: writing $x=a_{\rm st}b_d+v$ with $v\perp b_d$, its maps are
\[
 x\longmapsto\frac{v}{1-a_{\rm st}},\qquad
 y\longmapsto\frac{2y+(\|y\|^2-1)b_d}{1+\|y\|^2},
 \qquad b_d\longleftrightarrow\infty.
\]
An orthonormal basis identifies $b_d^\perp$ with $\mathbb R^d$;
the two displayed substitutions are inverse.
The Cartesian powers of $h$, with coordinatewise inverse $h^{-1}$,
and the integral K\"unneth theorem give
\[
 H_k(X^n;\mathbb Z)\cong
 \begin{cases}
  \mathbb Z^{\binom nj},&k=6j,\quad0\le j\le n,\\
  0,&\text{otherwise},
 \end{cases}
 \qquad \chi(X^n)=2^n,
\]
which proves
\cref{wb:cusp-spec-product-not-sphere,wb:cusp-spec-four-product-not-sphere}.
The synchronized regular \(n\)-fold fibre product has dimension
\[
 6n-2(n-1)=4n+2,
\]
which gives \(18\) for four copies. The map is the diagonal pullback
into \(J^n\), of codimension \(2(n-1)\). The Cartesian-to-smash quotient,
followed by \(h^{\wedge n}\), identifies its target with
\[
 ((\mathbb R^6)^+)^{\wedge n}\cong(\mathbb R^{6n})^+=S^{6n},
\]
by the identity on $(\mathbb R^6)^n$ and collapse of its complement.
The inverse splits a finite vector into its $n$ ordered six-dimensional
blocks and sends $\infty$ to the smash basepoint; the inverse of
$h^{\wedge n}$ applies $h^{-1}$ to every nonbasepoint factor.
The block identification and its inverse extend continuously across
infinity because they are proper linear maps. Both $h^{\wedge n}$
and its inverse descend from the coordinatewise based diffeomorphisms.
The nonzero displayed middle homology groups exclude homotopy
equivalences from the Cartesian products to \(S^{12}\) or \(S^{24}\).

A product map factors through a smash quotient if and only if it is
constant on the corresponding fat wedge.  On
\(F_T\times\{0\}\), the local two-copy quotient is
\[
 x\longmapsto(\pi_T(x),0),
\]
which is nonconstant.  Hence \(\pi_T\times\pi_T\) does not factor through
\(F_T\wedge F_T\).  The restriction of
\(\chi_{\gamma,T}^{[4]}\) to \(F_T\times\{0\}^3\) is similarly
\[
 x\longmapsto(\chi_{\gamma,T}(x),0,0,0),
\]
so it does not factor through \(F_T^{\wedge4}\).

For an extension to the Cartesian powers,
\(H^1(X^n;\mathbb Z)=0\), whereas both local quotient maps pull nonzero
target \(H^1\)-classes back to nonzero classes on the cusp fibre:
the \(\widehat\gamma\)-loop maps to the first lattice generator in
\(Q_T\), and to a degree-one loop in the circle quotient. Hence no
map on \(X^2\) or \(X^4\) has the required restriction, and precomposition
with the smash collapses in \cref{wb:cusp-spec-smash-spheres} cannot produce
the local quotient maps.

Identify the metric quotient circle with
\(C_T=\mathbb R/\ell_T\mathbb Z\) by
\(\theta\mapsto\ell_T\theta\).  For
\cref{wb:cusp-spec-target-smash-map}, target-smash collapse is constant
whenever one source coordinate is the identity, so the modified map does
factor through \(F_T^{\wedge4}\).  The complement of the target fat wedge \(F\) is \((0,\ell_T)^4\).
Its collapse is the one-point compactification of that cube, identified
with \(S^4\) by the coordinate maps
$\theta_i\mapsto\tan(\pi(\theta_i/\ell_T-1/2))$,
with inverse $y_i\mapsto\ell_T(1/2+\arctan(y_i)/\pi)$, followed
by the stereographic inverse in dimension four. Infinity corresponds
to the collapsed fat wedge. For the quotient length metric,
\cref{mapsIII:lem:collapsed-subset-distance}, with $E=C_T^4$ and
$A=F$, gives
\[
 d_Q([x],[y])=\min\{d_{C_T^4}(x,y),
                   d_{C_T^4}(x,F)+d_{C_T^4}(y,F)\}.
\]
The comparison with that lemma's quotient sends $[x]$ to $[x]$,
with $F$ representing the basepoint.
The chosen centre has distance \(\ell_T/2\) from \(F\). Within radius
\(R\), every route through the collapsed point has length at least
\(\ell_T-2R\), while direct distances are at most \(2R\). When
\(\ell_T>4R\), the entire quotient ball is the Euclidean ball: reaching
it through the collapsed point already costs more than \(R\), and the
flat-torus covering introduces no shorter pairwise distance.

The collapsed point is a metric singularity. For \(0<r<\ell_T/2\),
its radius-\(r\) ball has four-dimensional Hausdorff measure
\[
 \ell_T^4-(\ell_T-2r)^4=8\ell_T^3r+O(r^2),
\]
because the metric is locally Euclidean away from that point. This
linear growth contradicts the order-\(r^4\) volume growth at a smooth
four-dimensional Riemannian point.

Let \(u_i\in H^1(F_T^4;\mathbb Z)\) be the pullback of the generator from
the \(i\)-th quotient circle.  The orientation class of
\(S^4=C_T^{\wedge4}\) pulls back under
\cref{wb:cusp-spec-target-smash-map} to
\[
 u_1\smile u_2\smile u_3\smile u_4\ne0,
\]
after choosing the target orientation from the open cube. The pullback
evaluates to one on the product of the four \(\widehat\gamma\)-loops.
But \(H^4(X^4;\mathbb Z)=0\).  Thus even this modified \(S^4\)-target map
does not extend to \(X^4\), and a fortiori does not arise from a map on
\(X^{\wedge4}\cong S^{24}\).

Borel--Serre's Proposition 15.1 and its following remark on p.~434
exclude almost-complex structures on \(S^{2n}\), \(n\geq4\), for every
differentiable structure \cite{BorelSerre1953}. Taking \(n=6,12\) proves
the assertions for \(S^{12}\) and \(S^{24}\). A complex structure would
induce an almost-complex structure.

The join comparison is the homeomorphism
\[
 (S^6)^{*n}\longrightarrow S^{7n-1},\qquad
 [(t_i,x_i)_{i=1}^n]\longmapsto(\sqrt{t_i}\,x_i)_{i=1}^n.
\]
Together with \(h\), it gives \(X*X\cong S^{13}\) and
\(X^{*4}\cong S^{27}\): the inverse assigns to each vector block its
squared norm \(t_i\) and, for \(t_i>0\), its unit vector \(x_i\).
Reduced suspension is smash product with a
sphere, so \(\Sigma^k h\) gives \(\Sigma^k X\cong S^{6+k}\), in
particular the dimensions \(12\) and \(24\) at \(k=6,18\).
For $A_{\rm tor}=(\mathbb R/\mathbb Z)^2$ or $\mathbb R/\mathbb Z$, form
the cone by collapsing $\{0\}\times S^1\times A_{\rm tor}$ in
$[0,\epsilon]\times S^1\times A_{\rm tor}$. Its positive part has projection
$p_{\rm cone}([r,z,a_{\rm cone}])=a_{\rm cone}$. For any $0<r_0\leq\epsilon$ and $z_0\in S^1$, the map
$a_{\rm cone}\mapsto[r_0,z_0,a_{\rm cone}]$ is a section: its composite with $p_{\rm cone}$ is
the identity. Thus $p_{\rm cone}$ recovers exactly the fixed quotient coordinate
from a radial quotient map. For distinct $a_{\rm cone},a_{\rm cone}'\in A_{\rm tor}$, the paths
$[r,z_0,a_{\rm cone}]$ and $[r,z_0,a_{\rm cone}']$ approach the same vertex as $r\to0$,
while their projections are constantly $a_{\rm cone}$ and $a_{\rm cone}'$. Continuity at
that vertex is therefore impossible.

\end{proof}

\subsection{External smash products of sphere maps}

\begin{definition}[Ordered cube spheres]\hypertarget{tag-Q253}{}\label{wb:def:ordered-cube-spheres}
{\normalfont\scriptsize tag Q253}\par
For every integer $d\geq1$, put $C_d=I^d/\partial I^d$, with the
collapsed boundary as basepoint. The based homeomorphism
$h_d:C_d\longrightarrow(S^d,N)$, where $N=(0,\ldots,0,1)$, is
\[
 v_j=\tan\bigl(\pi(u_j-\tfrac12)\bigr),\qquad
 h_d([u])=\left(\frac{2v}{1+\lVert v\rVert^2},
                  \frac{\lVert v\rVert^2-1}{1+\lVert v\rVert^2}\right)
 \quad(u\in(0,1)^d),\qquad h_d(*)=N.
\]
Its inverse off $N$ sends $z=(z',z_{d+1})$ to the class of
\[
 u_j=\tfrac12+\frac1\pi\arctan\frac{z'_j}{1-z_{d+1}};
\]
it sends $N$ to $*$. Approaching the cube boundary makes
$\lVert v\rVert$ tend to infinity, so $h_d$ is continuous there.
The two formulas are inverse on the open cells; compactness of $C_d$
and the Hausdorff property of $S^d$ give continuity of the inverse.

For positive integers $d_1,\ldots,d_k$, define the ordered based
homeomorphism
\[
 \mu_{d_1,\ldots,d_k}:
 C_{d_1}\wedge\cdots\wedge C_{d_k}\longrightarrow C_{d_1+\cdots+d_k},
 \qquad [[u^{(1)}],\ldots,[u^{(k)}]]\longmapsto
 [u^{(1)},\ldots,u^{(k)}].
\]
Its inverse splits an interior point into the displayed blocks and
sends the basepoint to the smash basepoint. The product cube boundary
is exactly the inverse image of the fat wedge, so these maps descend
from the same compact quotient and are inverse homeomorphisms.
The first-coordinate pinch $p_d:C_d\to C_d\vee C_d$ sends
$u_1\leq\tfrac12$ to the first copy with coordinate $2u_1$ and
$u_1\geq\tfrac12$ to the second copy with coordinate $2u_1-1$.
The middle face and boundary map to the basepoint.
Thus $(f\vee g)p_d$ is the representative of addition in
$\pi_d$ in these coordinates \cite[pp.~340--341]{Hatcher2002}.
The cube representative of a based map $f:S^m\to S^n$ is
$f_C=h_n^{-1}fh_m$; the inverse transport is $f=h_nf_Ch_m^{-1}$.
The same formulas transport based homotopies at each time.
Below, $x,y,g$ also denote these representatives, and the displayed
$\mu$ maps specify the source and target coordinates of every smash.
\end{definition}

\Needspace{6\baselineskip}
\begin{theorem}[The full external smash pairing is zero]
\hypertarget{tag-038V}{}\label{wb:cusp-s12-full-carrier-null}
{\normalfont\scriptsize tag 038V}\par
For every pair of based maps $x,y:S^6\to S^3$,
\[
 [x\wedge y]=0\qquad\text{in }\pi_{12}(S^6).
\]
\end{theorem}

\begin{proof}
Use the based sphere models, concatenation maps and pinch maps of
Definition~\ref{wb:def:ordered-cube-spheres}.

For $x,y:C_6\to C_3$ put
\[
 R(x)=\mu_{3,3}(x\wedge\mathrm{id}_{C_3})\mu_{6,3}^{-1}
       :C_9\longrightarrow C_6,
 \qquad
 L(y)=\mu_{6,3}(\mathrm{id}_{C_6}\wedge y)\mu_{6,6}^{-1}
       :C_{12}\longrightarrow C_9.
\]
These are the right threefold and left sixfold suspensions in the
displayed block order.  The exact smash factorization is
\begin{equation}
 x\wedge y=(x\wedge\mathrm{id}_{C_3})
              (\mathrm{id}_{C_6}\wedge y):
 C_6\wedge C_6\longrightarrow C_6\wedge C_3
                \longrightarrow C_3\wedge C_3.
 \hypertarget{tag-0390}{}\label{wb:cusp-s12-smash-factorization}
\end{equation}
After the ordered identifications it gives
\[
 R(x)L(y)=\mu_{3,3}(x\wedge y)\mu_{6,6}^{-1}.
\]
To express this in one suspension convention, write
$E_R^r[f]=[\mu_{n,r}(f\wedge\mathrm{id}_{C_r})\mu_{m,r}^{-1}]$
for $f:C_m\to C_n$.
The block twist $\tau_{a,b}:C_{a+b}\to C_{b+a}$,
$\tau_{a,b}[u,v]=[v,u]$, has inverse $\tau_{b,a}$ and degree
$(-1)^{ab}$, since its coordinate permutation has $ab$ inversions.
In particular,
\[
 L(y)=\tau_{3,6}\,
       \bigl(\mu_{3,6}(y\wedge\mathrm{id}_{C_6})\mu_{6,6}^{-1}\bigr)
       \,\tau_{6,6}.
\]
Both twists here have degree one and hence are based homotopic to the
identity by \cite[Corollary~4.25, p.~361]{Hatcher2002}.  Therefore
\begin{equation}
 [\mu_{3,3}(x\wedge y)\mu_{6,6}^{-1}]
       =(E_R^3[x])\circ(E_R^6[y]).
 \hypertarget{tag-0392}{}\label{wb:cusp-s12-toda-smash-formula}
\end{equation}
Pinching the first coordinate of the six-coordinate input block of
$x$ before concatenation gives
\[
 R((x\vee x')p_6)=(R(x)\vee R(x'))p_9.
\]
Hence $R_*:\pi_6(C_3)\to\pi_9(C_6)$, $[x]\mapsto[R(x)]$,
is a homomorphism.  Based homotopies descend through these smash
quotients because the fat wedge stays at the basepoint at every time;
the product of each compact quotient with the homotopy interval is
again a quotient map.

For fixed $y$, the map $L(y)$ leaves its first six coordinates
unchanged.  Consequently it satisfies the based-map identity
\[
 p_9L(y)=(L(y)\vee L(y))p_{12}.
\]
Indeed, both sides double the first coordinate in the first half and
replace it by $2u_1-1$ in the second half.  Both send the middle face,
the boundary, and every point whose $y$-coordinate is the basepoint
to the basepoint.  Thus
\[
 D_y:\pi_9(C_6)\longrightarrow\pi_{12}(C_6),
 \qquad [z]\longmapsto[zL(y)]
\]
is a homomorphism: applying the displayed pinch identity to
$(z\vee z')p_9$ proves additivity.

The integral computations in \cite[p.~339]{Hatcher2002} give the first,
second and fourth groups below:
\begin{equation}
 \begin{gathered}
 \pi_6(S^3)\cong\mathbb Z/12,\qquad
 \pi_9(S^6)\cong\mathbb Z/24,\\
 \pi_{12}(S^9)\cong\mathbb Z/24,\qquad
 \pi_{12}(S^6)\cong\mathbb Z/2.
 \end{gathered}
 \hypertarget{tag-0391}{}\label{wb:cusp-s12-three-stem-groups}
\end{equation}
The remaining group is identified with $\pi_9(S^6)$ by the exact
suspension isomorphism
\[
 E_R^3:\pi_9(S^6)\xrightarrow{\ \cong\ }\pi_{12}(S^9).
\]
Indeed, the three Freudenthal inequalities are $9<11$, $10<13$ and
$11<15$ \cite[Corollary~4.24, pp.~360--361]{Hatcher2002}.
Choose a generator $\nu$ of $\pi_9(C_6)$, so that
$[k]_{24}\mapsto k\nu$ is its cyclic coordinate isomorphism.
Since $12[x]=0$ and $R_*$ is a homomorphism, writing
$R_*[x]=a_{\rm cyc}\nu$ gives $24\mid12a_{\rm cyc}$, hence $2\mid a_{\rm cyc}$.
Equivalently, the kernel of multiplication by $12$ on $\mathbb Z/24$
is the image of the isomorphism
$[j]_{12}\mapsto[2j]_{24}$, whose inverse sends $[2j]_{24}$ to
$[j]_{12}$.  Therefore $R_*[x]=2\zeta$ for some
$\zeta\in\pi_9(C_6)$, and
\[
 [\mu_{3,3}(x\wedge y)\mu_{6,6}^{-1}]
   =D_y(R_*[x])=2D_y(\zeta)=0.
\]
All groups and homotopies here are based, so this is a homotopy to
the constant map that fixes the basepoint throughout.
\end{proof}

\Needspace{6\baselineskip}
\begin{corollary}[Higher smash powers]
\hypertarget{tag-038T}{}\label{wb:cusp-s12-all-higher-smashes-null}
{\normalfont\scriptsize tag 038T}\par
For every continuous based map $g:S^6\to S^3$ and every integer
$r\geq2$, the ordered external smash
\[
 g^{\wedge r}:S^{6r}\longrightarrow S^{3r}
\]
is based null-homotopic.
\end{corollary}

\begin{proof}
Use $C_d$ and the ordered block homeomorphisms
$\mu_{d_1,\ldots,d_k}$ of Definition~\ref{wb:def:ordered-cube-spheres}.
By \cref{wb:cusp-s12-full-carrier-null}, there is a based homotopy
$H_2:C_{12}\times I\to C_6$ with initial map
$\mu_{3,3}(g\wedge g)\mu_{6,6}^{-1}$ and constant terminal map.
This proves the case $r=2$.

For $r>2$, split the source by
$a_r=\mu_{12,6,\ldots,6}^{-1}$ and put
$b_r=\mu_{6,3,\ldots,3}$, with $r-2$ remaining blocks in each.
Define
\[
 H_r\bigl(a_r^{-1}[w,u_3,\ldots,u_r],s\bigr)
 =b_r[H_2(w,s),g(u_3),\ldots,g(u_r)].
\]
The fat wedge maps to the target basepoint at every time, so the
formula descends continuously to $C_{6r}\times I$.
Indeed, the product-to-smash quotient is a closed surjection between
compact Hausdorff spaces, and its product with $I$ is also a quotient.
Concatenating the first two blocks and then the remaining blocks is
the same map as concatenating all blocks at once. Thus
$H_r(-,0)=\mu_{3,\ldots,3}g^{\wedge r}\mu_{6,\ldots,6}^{-1}$.
At time one the first smash factor is the basepoint, so $H_r(-,1)=*$;
the source basepoint remains fixed for every time.
\end{proof}

\subsection{Recognition of a two-sphere and three-sphere wedge}

\begin{lemma}[The two-generator wedge]
\hypertarget{tag-02QC}{}\label{mapsIII:carrier-three-complex-wedge}
{\normalfont\scriptsize tag 02QC}\par
Let $C$ be a simply connected CW complex with
$\widetilde H_2(C;\mathbb Z)=\widetilde H_3(C;\mathbb Z)=\mathbb Z$
and zero reduced integral homology in every other degree.
For any chosen generators $u\in H_2(C;\mathbb Z)$ and
$v\in H_3(C;\mathbb Z)$, there are based maps
$a:S^2\to C$ and $b:S^3\to C$ carrying the oriented sphere
classes to $u$ and $v$. The map $F=a\vee b:S^2\vee S^3\to C$
is a based homotopy equivalence.
\end{lemma}

\begin{proof}
The Hurewicz isomorphism $\pi_2(C)\to H_2(C;\mathbb Z)$
supplies a based cellular map $a$ representing $u$.
Naturality shows that $a$ induces an isomorphism on $\pi_2$
as well as on $H_2$. Form its mapping cylinder
\[
 M=(S^2\times[0,1])\amalg C/((z,1)\sim a(z)),
 \qquad A=S^2\times\{0\}.
\]
The retraction $r:M\to C$ has $r([z,t])=a(z)$ and is the
identity on $C$. Its deformation retraction sends
$[z,t]$ to $[z,(1-s)t+s]$ and fixes $C$. The basepoint path
$[z_0,t]$ identifies the bottom basepoint with $a(z_0)$;
all homotopy groups below use this path. If $i:A\to M$ is
the inclusion, then $ri=a$ and
$i_\#=r_\#^{-1}a_\#$, with inverse basepoint transport along
the reversed path. Both $A$ and $M$ are simply connected.
The pair exact sequence therefore gives
$\pi_1(M,A)=\pi_2(M,A)=0$.

Relative Hurewicz and the homology exact sequence give the
isomorphisms
\[
 h:\pi_3(M,A)\xrightarrow{\cong}H_3(M,A;\mathbb Z),
 \qquad q_H:H_3(M;\mathbb Z)\xrightarrow{\cong}H_3(M,A;\mathbb Z).
\]
Here $H_3(A)=0$ and $i_*:H_2(A)\to H_2(M)$ is an isomorphism,
which prove both sides of the second assertion. The next
homotopy boundary has image
$\ker(i_\#: \pi_2(A)\to\pi_2(M))=0$; thus
$q_\pi:\pi_3(M)\to\pi_3(M,A)$ is surjective. Choose
$[g]\in\pi_3(M)$ with
\[
 q_\pi[g]=h^{-1}q_H(r_*^{-1}v),\qquad b=rg.
\]
Give $S^3=D^3/\partial D^3$ its disk orientation.
Naturality gives
$q_H(g_*[S^3])=h(q_\pi[g])=q_H(r_*^{-1}v)$.
The injectivity of $q_H$ proves $b_*[S^3]=v$.
These are the canonical relative maps of
\cite[Theorem~4.37 and Lemma~4.38, pp.~370--372]{Hatcher2002}.

The wedge map $F=a\vee b$ induces the chosen generator maps
in degrees two and three, the identity on the connected
degree-zero groups, and the unique isomorphisms between zero
groups in every other degree. It is an integral homology
equivalence between simply connected CW complexes. To specify
its homotopy inverse, let $W=S^2\vee S^3$ and form the reduced
based mapping cylinder
\[
 N=(W\times[0,1])\amalg C/\bigl((w,1)\sim F(w),\
                    (w_0,t)\sim c_0\bigr).
\]
Write $j_W,j_C$ for its bottom and top inclusions and let
$p([w,t])=F(w)$, $p(c)=c$. The pair $(N,W)$
is acyclic. Relative Hurewicz, starting with its simply
connected pair and proceeding through every degree, gives
$\pi_n(N,W)=0$ for all $n\ge1$. The CW Whitehead theorem
\cite[Theorem~4.5 and Corollary~4.33, pp.~346--347 and~367--368]{Hatcher2002}
therefore gives a retraction $k:N\to W$ and a deformation
$T_s$ from $1_N$ to $j_Wk$, fixed on $W$. Set $G=kj_C$.
The cylinder paths $U_t(w)=[w,t]$ give
$kU_t:1_W\simeq GF$, and $pT_sj_C$ gives
$1_C\simeq FG$. The defining collapsed arc and the condition that $T_s$ fixes
$W$ make both homotopies fix the basepoint. This constructs the inverse
for the actual chosen generator map, without imposing a bound
on the dimension or number of cells of $C$.
\end{proof}

\Needspace{12\baselineskip}
\section{Cusp quotients and conical targets}
\hypertarget{tag-037I}{}\label{wb:cusp-cone-section}
\subsection{The marked quotient over the punctured cusp}

\begin{definition}[Cusp disc and quotient marking]
\hypertarget{tag-Q30P}{}\label{mapsIII:def:cone-cusp-data}
{\normalfont\scriptsize tag Q30P}\par
Use the original sphere fibration $f:X\to B$ of
\cref{maps:src:B.2.1}. Fix its cusp coordinate and a closed coordinate disc
\begin{equation}
 D_0=\{t\in\mathbb C:|t|\leq\epsilon\}\subset B,\qquad
 p_0=\{t=0\},\qquad N_0=f^{-1}(D_0),\qquad W=f^{-1}(0),
 \hypertarget{tag-036L}{}\label{wb:cusp-cone-disc}
\end{equation}
small enough that it contains no singular value other than \(p_0\).  Put
\begin{equation}
 K_\Lambda=\mathbb Z\langle\widehat w,\widehat\delta\rangle,\qquad
 \overline\Lambda=\Lambda/K_\Lambda,\qquad
 \overline T^2=
 \frac{\overline\Lambda\otimes\mathbb R}{\overline\Lambda}.
 \hypertarget{tag-037F}{}\label{wb:cusp-cone-quotient-lattice}
\end{equation}
The ordered basis of \(\overline\Lambda\) is the image of
\((\widehat\gamma,\widehat u)\). In these sections,
\[
 B_R^Z(z)=\{w\in Z:d_Z(z,w)\leq R\}
\]
denotes the closed metric ball. The center and ambient space may be
suppressed when they are fixed by the displayed map.
\end{definition}

\begin{lemma}[The native quotient over the whole punctured cusp]
\hypertarget{tag-0370}{}\label{wb:cusp-cone-native-descent}
{\normalfont\scriptsize tag 0370}\par
There is a real-analytic fibrewise group epimorphism
\begin{equation}
 \overline\pi:N_0\setminus W
 \longrightarrow D_0^*\times\overline T^2
 \hypertarget{tag-0371}{}\label{wb:cusp-cone-native-family-map}
\end{equation}
whose kernel is the connected two-torus family defined by \(K_\Lambda\).
The target local system is trivial. On the positive radial fibre \(F_T\),
at \(t=e^{-2\pi T}\), put
\[
 \overline\pi_T:=\operatorname{pr}_{\overline T^2}
     \circ\overline\pi|_{F_T}.
\]
The isomorphism
\begin{equation}
 \iota_T:\overline T^2\longrightarrow Q_T,\qquad
 [(a,b)]\longmapsto[Y_T(a,b)^{\mathsf T}]
 \hypertarget{tag-037G}{}\label{wb:cusp-cone-radial-isometry}
\end{equation}
satisfies
\begin{equation}
 \pi_T=\iota_T\circ\overline\pi_T.
 \hypertarget{tag-0372}{}\label{wb:cusp-cone-native-restriction}
\end{equation}
Give \(\overline T^2\) the pullback metric \(g_T\) under \(\iota_T\) of
the Euclidean quotient metric on \(Q_T=\mathbb R^2/Y_T\mathbb Z^2\).
Then \(\iota_T\) is an isometry and, for every \(T\geq T_0\),
\begin{equation}
 \operatorname{sys}(\overline T^2,g_T)
 =\operatorname{sys}(Q_T)\geq T-C_{\rm sys},
 \hypertarget{tag-0373}{}\label{wb:cusp-cone-native-systole}
\end{equation}
where \(C_{\rm sys}\) is the constant in \cref{wb:cusp-spec-native-quotient}.
\end{lemma}

\begin{proof}
On a logarithmic cusp lift, write
\[
 t=e^{2\pi i s},\qquad
 \Pi(s)=[sB_0+C(t)\mid I_2],\qquad
 B_0=\begin{pmatrix}0&1\\-1&0\end{pmatrix}.
\]
The flat marking is \(\Xi(s,x+\Lambda)=[s,\Pi(s)x]\).
The period laws and monodromy give
\[
 s(g_0z)=s(z)-1,\qquad
 \Pi(s-1)M_0=\Pi(s),\qquad
 (s,x+\Lambda)\longmapsto(s-1,M_0x+\Lambda).
\]
Thus the complex fibre coordinate is fixed by this cusp generator;
its inverse is \((s,x)\mapsto(s+1,M_0^{-1}x)\), and both preserve the
zero section. Exponentiation
\[
 (s,\zeta_1,\zeta_2)\longmapsto
 (e^{2\pi i\zeta_1},e^{2\pi i\zeta_2},e^{2\pi i s})
\]
identifies \(\mathbb C^2/\mathbb Z^2\) with \((\mathbb C^*)^2\), and
takes a long period \(\lambda\) to the multiplier
\(e^{2\pi iC(t)\lambda}t^{B_0\lambda}\). These are the defining
long-period identifications of \(N_0\setminus W\). The cusp generator
fixes \(\zeta\) and changes \(s\) by an integer, so exponentiation also
identifies the base quotients. It induces a biholomorphism of the period
quotient with \(N_0\setminus W\): a local inverse chooses logarithms of
the two multiplicative coordinates, whose branch changes are exactly
the short integer periods. Thus it transports the group law and marking,
with zero section represented by \((1,1,t)\).

Let \(p:\Lambda\to\overline\Lambda\) be the quotient projection, in the
ordered basis \((\widehat\gamma,\widehat u)\). Direct multiplication gives
\begin{equation}
 (M_0-I)(a,b,c,d)^{\mathsf T}=(0,0,b,-a)^{\mathsf T},\qquad
 \operatorname{im}(M_0-I)=K_\Lambda.
 \hypertarget{tag-036X}{}\label{wb:cusp-cone-monodromy-image}
\end{equation}
In particular \(pM_0=pM_0^{-1}=p\), and \(M_0\) fixes \(K_\Lambda\)
pointwise. Define
\[
 \overline\pi([s,\Pi(s)x])=(t,[p x]).
\]
Adding an element of \(\Lambda\) to \(x\) adds an element of
\(\overline\Lambda\) to \(p x\); changing the logarithmic lift replaces
\(x\) by a power of \(M_0\), which fixes \(p x\). These identities prove
descent over the entire punctured cusp. Let \(P_s\) be the realification
of \(\Pi(s)\) in the order
\((\operatorname{Re}\zeta_1,\operatorname{Im}\zeta_1,
\operatorname{Re}\zeta_2,\operatorname{Im}\zeta_2)\).
Locally the formula is \((t,y)\mapsto(t,[pP_s^{-1}y])\), so it is real
analytic. The formulas hold on a neighbourhood of the chosen closed
cusp disc and hence also at its outer boundary.

The map preserves fibrewise addition and is surjective because
\(p(a,b,0,0)=(a,b)\). Its fibrewise kernel is
\((K_\Lambda\otimes\mathbb R)/K_\Lambda\): subtracting a lattice vector
whose image is \(p x\) puts every kernel element in this subspace, and
primitivity identifies its lattice with \(K_\Lambda\). It is connected.
The kernel family has the explicit inclusion
\[
 (t,[c,d])\longmapsto[(e^{2\pi ic},e^{2\pi id},t)]
\]
in the long-period quotient. The identity action of \(M_0\) on
\(\overline\Lambda\) proves that the target local system is trivial.

On the radial fibre, \(y=P_Tx\) and
\(\rho P_T=Y_Tp\). Hence
\[
 \pi_T([y])=[\rho y]=[Y_Tp x]
   =\iota_T(\overline\pi_T([y])).
\]
The inverse of \(\iota_T\) is \([v]\mapsto[Y_T^{-1}v]\), and
\(\overline\pi|_{F_T}=(t,\overline\pi_T)\); these formulas also identify
the full family-valued restriction. In the ordered quotient basis,
\[
 g_T=Y_T^{\mathsf T}Y_T
 =\begin{pmatrix}
 36m_T^2+q_T^2&m_T(6L_T-q_T)\\
 m_T(6L_T-q_T)&L_T^2+m_T^2
 \end{pmatrix},\qquad \det g_T=D_T^2>0.
\]
Thus \(\iota_T\) is a global Riemannian isometry. The systolic bound
follows from \cref{wb:cusp-spec-native-quotient}; when \(T-C_{\rm sys}\leq0\),
it also follows directly from the positive systole. This covers every
\(T\geq T_0\).
\end{proof}

\subsection{The two-torus suspension carrier}
\begin{definition}[The two-torus suspension carrier]
\hypertarget{tag-Q30Q}{}\label{mapsIII:def:two-torus-suspension}
{\normalfont\scriptsize tag Q30Q}\par
Set
\begin{equation}
 \mathcal A=
 \frac{[0,\epsilon]\times S^1\times\overline T^2}
 {\{0\}\times S^1\times\overline T^2\sim v_0,\quad
  \{\epsilon\}\times S^1\times\overline T^2\sim v_1}.
 \hypertarget{tag-0354}{}\label{wb:cusp-cone-A-definition}
\end{equation}
Write $\operatorname{Susp}$ for the unreduced suspension. The radial
coordinate maps $[r,z,a]\mapsto[r/\epsilon,z,a]$ and
$[s,z,a]\mapsto[\epsilon s,z,a]$ give
\begin{equation}
 \mathcal A\cong
 \operatorname{Susp}(S^1\times\overline T^2)
 \cong\operatorname{Susp}(T^3).
 \hypertarget{tag-0358}{}\label{wb:cusp-cone-A-suspension}
\end{equation}
The first circle is the angular coordinate of the punctured base.
\end{definition}

\Needspace{6\baselineskip}
\begin{theorem}[Localized two-torus quotient with sphere source]
\hypertarget{tag-035X}{}\label{wb:cusp-cone-S12-route}
{\normalfont\scriptsize tag 035X}\par
Choose $x_0\in W$ and base $\mathcal A$ at $v_0$.
There is a continuous surjective based quotient map
\begin{equation}
 q_2:(X,x_0)\longrightarrow(\mathcal A,v_0).
 \hypertarget{tag-037D}{}\label{wb:cusp-cone-q2}
\end{equation}
Let $\overline\pi_f=\operatorname{pr}_{\overline T^2}\circ
\overline\pi$ be the fibre component of the native family quotient.
On the punctured cusp interior write $t(f(x))=re^{i\phi}$,
where $0<r<\epsilon$.  Then
\begin{equation}
 \begin{gathered}
 q_2(x)=[r,e^{i\phi},\overline\pi_f(x)],\qquad
 q_2(W)=v_0,\\
 q_2\bigl(X\setminus\operatorname{int}N_0\bigr)=v_1.
 \end{gathered}
 \hypertarget{tag-037E}{}\label{wb:cusp-cone-q2-formula}
\end{equation}
In particular, $f^{-1}(|t|=\epsilon)$ maps to $v_1$.
The smash square is a global based map
\begin{equation}
 q_2^{\wedge2}:X\wedge X\longrightarrow\mathcal A\wedge\mathcal A.
 \hypertarget{tag-035V}{}\label{wb:cusp-cone-S12-map}
\end{equation}
A based diffeomorphism $X\to S^6$ identifies its source as
\begin{equation}
 X\wedge X\cong S^6\wedge S^6\cong S^{12}.
 \hypertarget{tag-035Y}{}\label{wb:cusp-cone-S12-source}
\end{equation}
For $0<e^{-2\pi T}<\epsilon$, the canonical product-to-smash quotients
embed $F_T^2$ and the corresponding interior target product.
Write
\[
 \overline\pi_T=\operatorname{pr}_{\overline T^2}\circ
             \overline\pi|_{F_T},\qquad
 \pi_T=\iota_T\circ\overline\pi_T.
\]
Under those embeddings and the isomorphism
$\iota_T:\overline T^2\to Q_T$, the restricted map is exactly
\begin{equation}
 \pi_T\times\pi_T:F_T\times F_T\longrightarrow Q_T\times Q_T.
 \hypertarget{tag-035U}{}\label{wb:cusp-cone-S12-fibre-map}
\end{equation}
Use the constant \(C_{\rm sys}\) fixed in
\cref{wb:cusp-cone-native-descent}.
With the intrinsic Euclidean $\ell^2$ product flat quotient metric,
\begin{equation}
 \operatorname{sys}(Q_T\times Q_T)\geq T-C_{\rm sys},\qquad
 (Q_T\times Q_T,(0,0))\longrightarrow(\mathbb R^4,0)
 \hypertarget{tag-035Q}{}\label{wb:cusp-cone-S12-R4-limit}
\end{equation}
in pointed local-isometric, pointed smooth and pointed
Gromov--Hausdorff topology.  The radius-$R$ ball about $(0,0)$ is
exactly Euclidean once $T-C_{\rm sys}>4R$.

The integral target homology groups are
\begin{equation}
 \widetilde H_k(\mathcal A;\mathbb Z)\cong
 \begin{cases}
  \mathbb Z^3,&k=2,3,\\
  \mathbb Z,&k=4,\\
  0,&\text{otherwise},
 \end{cases}
 \hypertarget{tag-0355}{}\label{wb:cusp-cone-A-homology}
\end{equation}
and
\begin{equation}
 \widetilde H_k(\mathcal A\wedge\mathcal A;\mathbb Z)\cong
 \begin{cases}
  \mathbb Z^9,&k=4,\\
  \mathbb Z^{18},&k=5,\\
  \mathbb Z^{15},&k=6,\\
  \mathbb Z^6,&k=7,\\
  \mathbb Z,&k=8,\\
  0,&\text{otherwise}.
 \end{cases}
 \hypertarget{tag-0356}{}\label{wb:cusp-cone-A-smash-homology}
\end{equation}
Neither $\mathcal A$ nor $\mathcal A\wedge\mathcal A$ admits a
homeomorphism to $S^4$, $S^{12}$, or a compact four-manifold,
with or without boundary.
\end{theorem}

\begin{proof}
In the defining cylinder for $\mathcal A$, the equivalence relation
is the diagonal together with the squares of the two closed end tori.
It is closed, so the quotient is compact Hausdorff.
The maps $[r,z,a]\mapsto[r/\epsilon,z,a]$ and
$[u,z,a]\mapsto[\epsilon u,z,a]$ identify it with the unreduced
suspension of $S^1\times\overline T^2=T^3$.

By Lemma~\ref{wb:cusp-cone-native-descent},
$\overline\pi:N_0\setminus W\to D_0^*\times\overline T^2$
is single-valued and fibrewise onto.  Its trivial target local system
comes from the identity action of $M_0$ on $\Lambda/K_\Lambda$.
Composing with the displayed target projection therefore gives the
continuous interior formula for $q_2$ on the whole punctured cusp.
An open neighbourhood $U$ of $v_0$ pulls back to an open subset of the
cylinder containing $\{0\}\times T^3$.  Compactness of $T^3$ gives
$\delta>0$ such that $[0,\delta)\times T^3$ lies in that preimage.
Thus $q_2^{-1}(U)$ contains the neighbourhood
$f^{-1}(|t|<\delta)$ of $W$, proving continuity there.
At $v_1$ the same compactness argument gives an outer collar
$(\epsilon-\delta,\epsilon]\times T^3$.  Source points in this
collar and points outside $\operatorname{int}N_0$ map into the same
neighbourhood of $v_1$, which proves continuity at the outer boundary.
Together with interior continuity and the constant exterior map,
these arguments prove continuity everywhere.

Every interior point $[r,z,a]$ is attained from the regular fibre
over $t=rz$, and $W$ and the outer boundary attain the two poles.
Consequently $q_2$ is a surjection.  A continuous surjection from
compact $X$ to Hausdorff $\mathcal A$ is closed and hence a quotient.
Since $q_2(x_0)=v_0$, its product sends the source fat wedge into
the target fat wedge and descends to \eqref{wb:cusp-cone-S12-map}.
Take the diffeomorphism $h:X\to S^6$ from
\cref{wb:cusp-spec-globalization-obstruction} and a rotation $O_{\rm rot}$
with $O_{\rm rot}h(x_0)=\infty$. Apply the ordered source homeomorphism of
\eqref{wb:cusp-spec-smash-spheres} to $h_0=O_{\rm rot}h$.
On nonbasepoint factors it concatenates the stereographic coordinates
of $h_0$; its inverse splits the two blocks and applies inverse
stereographic projection followed by $h^{-1}O_{\rm rot}^{-1}$.
The basepoints correspond, giving \eqref{wb:cusp-cone-S12-source}.

The compact product $F_T^2$ avoids the source fat wedge; the smash
quotient is injective there and has Hausdorff target, so its
restriction is an embedding $i_T$.  The native target isomorphism
and inverse are
\[
 \iota_T[a]=[Y_Ta],\qquad
 \iota_T^{-1}[y]=[Y_T^{-1}y].
\]
They are well defined because $Y_T$ maps $\mathbb Z^2$ bijectively
onto the period lattice of $Q_T$.
For the target embedding use
\[
 j_T:Q_T\longrightarrow\mathcal A,\qquad
 j_T[q]=[e^{-2\pi T},1,\iota_T^{-1}(q)].
\]
Its image lies in an interior cross-section and avoids $v_0$.
Its product followed by the target smash quotient is consequently
an embedding $j_T^{[2]}$.  On each radial fibre the defining identity
$\pi_T=\iota_T\overline\pi_T$ gives
\[
 q_2^{\wedge2}\circ i_T
   =j_T^{[2]}\circ(\pi_T\times\pi_T),
\]
since both sides have coordinates
$[e^{-2\pi T},1,\overline\pi_T(x)]$ in each target factor.
This proves \eqref{wb:cusp-cone-S12-fibre-map} with these domains and targets.

The intrinsic product period lattice is
$Y_T\mathbb Z^2\oplus Y_T\mathbb Z^2\subset\mathbb R^4$.
Every nonzero vector has a nonzero component of length at least
$\operatorname{sys}(Q_T)$, and a shortest vector in one component
attains that bound.  Thus its systole equals
$\operatorname{sys}(Q_T)\geq T-C_{\rm sys}$ by
\eqref{wb:cusp-cone-native-systole}.
For $T-C_{\rm sys}>4R$, any $x,y\in B_R(0)\subset\mathbb R^4$ and nonzero
product period $\lambda$ satisfy
\[
 \lVert x-y+\lambda\rVert\geq\lVert\lambda\rVert-2R
       \geq T-C_{\rm sys}-2R>2R\geq\lVert x-y\rVert.
\]
Every point in the quotient radius-$R$ ball has a closest lattice
representative of norm at most $R$.  The displayed inequality
therefore identifies the entire quotient ball isometrically with
the Euclidean ball.  On open balls the covering map also pulls the
metric tensor back to the Euclidean tensor.  These identifications
for every fixed $R$ prove all three limits in
\eqref{wb:cusp-cone-S12-R4-limit}.

The product CW structure on $T^3$ has one vertex, three one-cells,
three two-cells and one three-cell, with zero cellular differential.
The unreduced suspension has vertices $v_0,v_1$, one edge $e$
with $\partial e=v_1-v_0$, three two-cells, three three-cells and
one four-cell.  Its remaining differentials vanish, proving
\eqref{wb:cusp-cone-A-homology}.
In the based reduced chain complex, $e\mapsto v_1$ is a contractible
summand.  Projection to the seven remaining cells, their inclusion,
and the homotopy $v_1\mapsto e$ split off this summand.
The cellular complex of the smash square is the tensor square of
that reduced complex.  The tensor factors containing the
contractible summand contract, so its homology is free with rank
polynomial
\[
 (3z^2+3z^3+z^4)^2
   =9z^4+18z^5+15z^6+6z^7+z^8.
\]
This proves \eqref{wb:cusp-cone-A-smash-homology} integrally.

In particular, $H_8(\mathcal A\wedge\mathcal A)=\mathbb Z$,
whereas $H_8(S^4)=H_8(S^{12})=0$.  A homeomorphism induces an
isomorphism on these groups, which excludes both spheres.
For the manifold assertion, choose two points in
$(0,\epsilon)\times T^3$ and small product coordinate neighbourhoods.
Their product avoids the fat wedge, and its image in the smash is
an open eight-dimensional chart.  At its centre $p$, excision and
the retraction $\mathbb R^8\setminus\{0\}\to S^7$ give
\[
 H_8(\mathcal A\wedge\mathcal A,
       (\mathcal A\wedge\mathcal A)\setminus\{p\};\mathbb Z)
       \cong\mathbb Z.
\]
At an interior point of a four-manifold this local group is
$H_8(\mathbb R^4,\mathbb R^4\setminus\{0\})=0$.
At a boundary point the half-space and its complement of the
boundary origin are both contractible, so this group is again zero.
The local homology maps induced by a homeomorphism are isomorphisms;
the displayed groups therefore exclude every compact four-manifold,
including one with boundary.

For the one-factor target, $H_2(\mathcal A)=\mathbb Z^3$ excludes
both spheres through the induced global homology maps.
A small pole neighbourhood is an open cone on $T^3$; it is
contractible and its punctured neighbourhood retracts onto $T^3$.
Excision and the exact sequence of this pair give
\[
 H_k(\mathcal A,\mathcal A\setminus\{v_0\};\mathbb Z)
 \cong\widetilde H_{k-1}(T^3;\mathbb Z).
\]
In degree two this is $\mathbb Z^3$, whereas the degree-two local
group at an interior or boundary point of a four-manifold is zero.
The same local homology obstruction therefore proves the asserted
exclusion for $\mathcal A$.
\end{proof}

\subsection{Circle suspensions and finite fillings}
\begin{definition}[The circle suspension carrier]
\hypertarget{tag-Q30R}{}\label{mapsIII:def:circle-suspension}
{\normalfont\scriptsize tag Q30R}\par
Define
\begin{equation}
 \mathcal B=
 \frac{[0,\epsilon]\times S^1\times S^1}
 {\{0\}\times S^1\times S^1\sim b_0,\quad
  \{\epsilon\}\times S^1\times S^1\sim b_1}
 \cong\operatorname{Susp}(T^2),
 \hypertarget{tag-0359}{}\label{wb:cusp-cone-B-definition}
\end{equation}
where the two circle coordinates in an interior cross-section are the base
angle and the value of \(\theta_\gamma\). The suspension coordinates are
related by $[r,z,\theta]\mapsto[r/\epsilon,z,\theta]$, with inverse
$[s,z,\theta]\mapsto[\epsilon s,z,\theta]$.
\end{definition}

\begin{proposition}[Localized circle quotient with sphere source]
\hypertarget{tag-0365}{}\label{wb:cusp-cone-S24-local-route}
{\normalfont\scriptsize tag 0365}\par
Choose $x_0\in W$.  There is a continuous surjective based quotient
\begin{equation}
 q_{1,\mathrm{cusp}}:(X,x_0)\longrightarrow(\mathcal B,b_0),
 \hypertarget{tag-037B}{}\label{wb:cusp-cone-q1-local}
\end{equation}
given on the punctured cusp interior by
\begin{equation}
 \begin{gathered}
 q_{1,\mathrm{cusp}}(x)=[r,e^{i\phi},\theta_\gamma(x)]
       \quad(0<|t|=r<\epsilon),\\
 q_{1,\mathrm{cusp}}(W)=b_0,\qquad
 q_{1,\mathrm{cusp}}(X\setminus\operatorname{int}N_0)=b_1.
 \end{gathered}
 \hypertarget{tag-037C}{}\label{wb:cusp-cone-q1-local-formula}
\end{equation}
Its fourth smash power has source $X^{\wedge4}\cong S^{24}$.
For $0<e^{-2\pi T}<\epsilon$, put
\[
 C_T=\mathbb R/\ell_T\mathbb Z,\qquad
 \vartheta_T:\mathbb R/\mathbb Z\longrightarrow C_T,
 \quad \vartheta_T[\theta]=[\ell_T\theta].
\]
The canonical product-to-smash quotients embed the source radial
product and the corresponding interior target product.  In the
target arclength coordinates, the restricted map is
\begin{equation}
 (\vartheta_T\circ\chi_{\gamma,T})^{\times4}:F_T^4\longrightarrow C_T^4.
 \hypertarget{tag-0364}{}\label{wb:cusp-cone-S24-local-fibre-map}
\end{equation}
The intrinsic target slices converge to $\mathbb R^4$ in pointed
local-isometric, smooth and Gromov--Hausdorff topology, by
\eqref{wb:cusp-spec-four-circle-limit}. The integral homology groups of
$\mathcal B$ and its fourth smash are
\begin{equation}
 \widetilde H_k(\mathcal B;\mathbb Z)\cong
 \begin{cases}
  \mathbb Z^2,&k=2,\\
  \mathbb Z,&k=3,\\
  0,&\text{otherwise},
 \end{cases}
 \hypertarget{tag-035A}{}\label{wb:cusp-cone-B-homology}
\end{equation}
and
\begin{equation}
 \widetilde H_k(\mathcal B^{\wedge4};\mathbb Z)\cong
 \begin{cases}
  \mathbb Z^{16},&k=8,\\
  \mathbb Z^{32},&k=9,\\
  \mathbb Z^{24},&k=10,\\
  \mathbb Z^8,&k=11,\\
  \mathbb Z,&k=12,\\
  0,&\text{otherwise}.
 \end{cases}
 \hypertarget{tag-035B}{}\label{wb:cusp-cone-B-smash-homology}
\end{equation}
\end{proposition}

\begin{proof}
The invariant character $\varepsilon_\gamma M_0=\varepsilon_\gamma$
makes the interior formula single-valued around the cusp.
An open neighbourhood of $b_0$ pulls back to an open subset of
$[0,\epsilon]\times T^2$ containing $\{0\}\times T^2$.
Compactness of $T^2$ gives $\delta>0$ such that
$[0,\delta)\times T^2$ lies in that preimage.  The inverse image
under $q_{1,\mathrm{cusp}}$ therefore contains a neighbourhood
$f^{-1}(|t|<\delta)$ of $W$.  At $b_1$ the same argument gives
an outer collar $(\epsilon-\delta,\epsilon]\times T^2$.
The formula on that collar and the constant value on
$X\setminus\operatorname{int}N_0$ prove continuity across the outer
boundary.  Continuity elsewhere follows from the character map.

The character is onto in every regular fibre, so every interior
target point is attained.  The fibre $W$ and the outer boundary
give the two poles.  Thus the map is a continuous surjection from
compact $X$ onto Hausdorff $\mathcal B$, and hence is a quotient.
It is based at $x_0\in W$, so its fourth product descends to the
fourth smash.  Use the based diffeomorphism $h_0$ of
\cref{wb:cusp-cone-S12-route}. The four-block source homeomorphism
of \eqref{wb:cusp-spec-smash-spheres}, with its block-splitting
inverse followed by $h_0^{-1}$ in sphere coordinates, gives
$X^{\wedge4}\cong S^{24}$.

Write $r=e^{-2\pi T}$.  The source $F_T^4$ is compact and avoids the
source fat wedge.  Its restriction under the compact Hausdorff smash
quotient is injective, hence an embedding.  The corresponding target
embedding in one factor is
\[
 j_T:C_T\longrightarrow\mathcal B,
 \qquad j_T[u]=[r,1,[u/\ell_T]].
\]
Its image is an interior circle avoiding $b_0$, so its fourth product
followed by the target smash quotient is also an embedding.
The identity $j_T(\vartheta_T(\chi_{\gamma,T}(x)))
=[r,1,\chi_{\gamma,T}(x)]$ proves
\eqref{wb:cusp-cone-S24-local-fibre-map} with its stated target.

The isometry $\vartheta_T$ identifies the character circle with
$C_T$; its inverse divides arclength by $\ell_T$. On four factors the
coordinate map $([u_i])_{i=1}^4\mapsto[(u_1,u_2,u_3,u_4)]$ and its
coordinatewise inverse identify the product isometrically with
$\mathbb R^4/\ell_T\mathbb Z^4$. This is exactly the metric and lattice
of \eqref{wb:cusp-spec-four-circle-fibre}. Its entire-ball isometries
and the three pointed limits are those of
\eqref{wb:cusp-spec-four-circle-limit}.

Finally, the coordinate map $[r,z,\theta]\mapsto[r/\epsilon,z,\theta]$
identifies $\mathcal B$ with the unreduced suspension of $T^2$.
Its CW structure has vertices $b_0,b_1$, an edge $e$ with
$\partial e=b_1-b_0$, two cells in degree two and one in degree
three; all other cellular differentials vanish.  This proves
\eqref{wb:cusp-cone-B-homology}.  In the based reduced cellular
complex the pair $e\mapsto b_1$ is contractible.  Projection onto
the remaining cells, inclusion, and the homotopy $b_1\mapsto e$
split off that pair.  Tensoring four copies therefore gives free
homology with rank polynomial
\[
 (2z^2+z^3)^4=16z^8+32z^9+24z^{10}+8z^{11}+z^{12}.
\]
This proves \eqref{wb:cusp-cone-B-smash-homology} integrally.
\end{proof}

\begin{definition}[The finite-filled circle carrier]
\hypertarget{tag-Q30S}{}\label{mapsIII:def:finite-filled-circle-carrier}
{\normalfont\scriptsize tag Q30S}\par
Choose pairwise disjoint closed coordinate discs \(D_0,D_1,D_2\) about the
three singular values and put
\begin{equation}
 P=B\setminus\bigcup_{j=0}^2\operatorname{int}D_j.
 \hypertarget{tag-0377}{}\label{wb:cusp-cone-pair-of-pants}
\end{equation}
For \(j=1,2\), let \(C_j\) be the solid-torus Seifert target in
\cref{wb:cusp-spec-finite-Seifert-target}, and let \(h_j^{S^1}\) be the
target overlap map in \cref{wb:cusp-spec-finite-overlap-identity}.  Glue
\begin{equation}
 \mathcal M=(P\times S^1)
 \underset{h_1^{S^1}}{\cup}C_1
 \underset{h_2^{S^1}}{\cup}C_2.
 \hypertarget{tag-035P}{}\label{wb:cusp-cone-M-definition}
\end{equation}
The gluing is along boundary tori over \(\partial D_1,\partial D_2\).
The remaining boundary is
\begin{equation}
 \partial\mathcal M=\partial D_0\times S^1\cong T^2.
 \hypertarget{tag-035N}{}\label{wb:cusp-cone-M-boundary}
\end{equation}
Define
\begin{equation}
 \mathcal C=\mathcal M\cup_{\partial\mathcal M}
 \operatorname{Cone}(\partial\mathcal M),
 \hypertarget{tag-035C}{}\label{wb:cusp-cone-C-definition}
\end{equation}
and write \(c_0\) for the cone vertex. Its radius runs from zero to
$\epsilon$; the outer torus is attached to $\partial\mathcal M$ by the identity.
\end{definition}

\begin{theorem}[Circle-cone quotient with sphere source]
\hypertarget{tag-036A}{}\label{wb:cusp-cone-S24-route}
{\normalfont\scriptsize tag 036A}\par
There is a continuous surjective based quotient
\begin{equation}
 q_1:(X,x_0)\longrightarrow(\mathcal C,c_0),\qquad x_0\in W,
 \hypertarget{tag-0379}{}\label{wb:cusp-cone-q1}
\end{equation}
with the following restrictions:
\begin{enumerate}
\item Over $P$ it is
$(f,\theta_\gamma):f^{-1}(P)\to P\times S^1$.
\item Over $D_j$, $j=1,2$, it is the finite-filling map
$N_j\to C_j$ of \eqref{wb:cusp-spec-finite-Seifert-target}.
\item On the punctured cusp, writing $t=re^{i\phi}$, it is
\begin{equation}
 x\longmapsto[r,(e^{i\phi},\theta_\gamma(x))]
       \in\operatorname{Cone}(\partial D_0\times S^1),
 \hypertarget{tag-037A}{}\label{wb:cusp-cone-q1-cusp-formula}
\end{equation}
and every point of $W$ maps to $c_0$.  Here
$(e^{i\phi},\theta)$ denotes the boundary point
$(\epsilon e^{i\phi},\theta)$, the cone radius is $r\in[0,\epsilon]$,
and its outer boundary is glued to $\partial\mathcal M$.
\end{enumerate}
The fourth smash power is a global based map
\begin{equation}
 q_1^{\wedge4}:X^{\wedge4}\longrightarrow\mathcal C^{\wedge4}.
 \hypertarget{tag-0366}{}\label{wb:cusp-cone-S24-map}
\end{equation}
A based diffeomorphism $X\to S^6$ identifies its source as
\begin{equation}
 X^{\wedge4}\cong(S^6)^{\wedge4}\cong S^{24}.
 \hypertarget{tag-036D}{}\label{wb:cusp-cone-S24-source}
\end{equation}
For $0<e^{-2\pi T}<\epsilon$, put
$\vartheta_T[\theta]=[\ell_T\theta]$ in
$C_T=\mathbb R/\ell_T\mathbb Z$, and give $\mathbb R/\mathbb Z$
the metric $\ell_T^2\,d\theta^2$ pulled back by $\vartheta_T$.
The canonical product-to-smash quotients embed the four radial fibres
and their interior target product.  Under these embeddings and the
arclength isometry $\vartheta_T$, the restriction is
\begin{equation}
 (\vartheta_T\circ\chi_{\gamma,T})^{\times4}:F_T^4\longrightarrow C_T^4.
 \hypertarget{tag-0363}{}\label{wb:cusp-cone-S24-fibre-map}
\end{equation}
For the intrinsic Euclidean $\ell^2$ product quotient metric,
\begin{equation}
 \operatorname{sys}(C_T^4)=\ell_T
 =T+H-B+O(T^{-1}),\qquad
 (C_T^4,0)\longrightarrow(\mathbb R^4,0)
 \hypertarget{tag-0360}{}\label{wb:cusp-cone-S24-R4-limit}
\end{equation}
in pointed local-isometric, pointed smooth and pointed
Gromov--Hausdorff topology.
The integral target homology is
\begin{equation}
 \widetilde H_k(\mathcal C;\mathbb Z)\cong
 \begin{cases}
  \mathbb Z,&k=2,3,\\
  0,&\text{otherwise},
 \end{cases}
 \hypertarget{tag-035D}{}\label{wb:cusp-cone-C-homology}
\end{equation}
and
\begin{equation}
 \widetilde H_k(\mathcal C^{\wedge4};\mathbb Z)\cong
 \begin{cases}
  \mathbb Z,&k=8,12,\\
  \mathbb Z^4,&k=9,11,\\
  \mathbb Z^6,&k=10,\\
  0,&\text{otherwise}.
 \end{cases}
 \hypertarget{tag-035E}{}\label{wb:cusp-cone-C-smash-homology}
\end{equation}
\end{theorem}

\begin{proof}
The compact core $P\times S^1$ and the two solid tori $C_j$ are
glued along closed boundary tori by the homeomorphisms $h_j^{S^1}$.
Their union $\mathcal M$ is a compact three-manifold with remaining
boundary $T_0=\partial D_0\times S^1$, and attaching its cone gives
a compact Hausdorff $\mathcal C$.
The identity
\[
 \chi_\gamma h_j=h_j^{S^1}\overline\chi_j
\]
makes the finite-filling pieces of $q_1$ agree with the core map.
The cusp formula agrees with that core map on the outer cone boundary.
A neighbourhood of $c_0$ contains the image of
$[0,\delta)\times T_0$ for some $\delta>0$, by compactness of $T_0$.
Its inverse image contains $f^{-1}(|t|<\delta)$, which proves
continuity at $W$.  Closed pasting proves continuity everywhere.
Each regular character and finite-filling map is onto, and $W$ gives
the cone vertex.  Compactness of $X$ then makes $q_1$ a quotient map.

Its based fourth product descends to the smash.  The based source homeomorphism in
Proposition~\ref{wb:cusp-cone-S24-local-route}
gives \eqref{wb:cusp-cone-S24-source}.  The compact product $F_T^4$
avoids the source fat wedge.  In the target use the embedding
\[
 j_T:C_T\longrightarrow\mathcal C,
 \qquad j_T[u]=[e^{-2\pi T},(1,[u/\ell_T])]
\]
in the cone coordinates of the statement.  Its image avoids $c_0$,
so the two product-to-smash quotients restrict to embeddings.
The identity $j_T(\vartheta_T(\chi_{\gamma,T}(x)))
=[e^{-2\pi T},(1,\chi_{\gamma,T}(x))]$ proves
\eqref{wb:cusp-cone-S24-fibre-map}.  In the marked fibre lattice,
the kernel of the angle character is identified with
$(L\otimes_{\mathbb Z}\mathbb R)/L$ by $[v]\mapsto[P_Tv]$, where
$L=\mathbb Z\langle\widehat u,\widehat w,\widehat\delta\rangle$.
The isomorphism $\vartheta_T$ preserves this kernel, so the displayed product
map has kernel $\bigl((L\otimes_{\mathbb Z}\mathbb R)/L\bigr)^4
\cong(\mathbb T^3)^4$.
The target product metric and the entire-radius-ball argument are
the same $\mathbb R^4/\ell_T\mathbb Z^4$ calculation proved in that
proposition.  Together with
\eqref{wb:cusp-spec-global-circle-circumference}, this proves
\eqref{wb:cusp-cone-S24-R4-limit}.

We compute the integral boundary maps of the finite fillings.
For $(m,\kappa)=(3,1)$ or $(4,-1)$, the quotient
\eqref{wb:cusp-spec-finite-Seifert-target} has solid-torus coordinates
\[
 z=s e^{2\pi i\kappa\theta},\qquad U=e^{2\pi im\theta},
 \qquad t=z^mU^{-\kappa}.
\]
They are invariant under its generator because $\kappa^2=1$.
An inverse chooses $\theta$ with $e^{2\pi im\theta}=U$ and sets
$s=z e^{-2\pi i\kappa\theta}$; a different choice changes the
pair by a power of the quotient generator.  Let $l_j$ be the positive
$U$-circle in the central solid torus.

On $|s|=R_j$, the logarithmic overlap has the real lift
\[
 \beta_j^\ell(s)=\frac{\kappa_j\arg s}{2\pi}+e_j(t),\qquad
 e_j(t)=-\frac{\log R_j}{2\pi}\varepsilon_\gamma(J_sv_j).
\]
Under the clockwise root generator, the compatible logarithm changes
the first term by $-\kappa_j/m_j$; the same change in
$\beta_j^\ell$ follows from
\eqref{wb:cusp-spec-finite-target-cocycle} and its real lift in
\eqref{wb:cusp-spec-finite-target-conjugacy}.
Their difference $e_j$ is therefore a single-valued real function of
$t$ on the boundary.  Write $a=\arg z/(2\pi)$,
$b=\arg U/(2\pi)$, $\tau=\arg t/(2\pi)$, and let $\Theta$
denote the regular circle coordinate.  The full boundary map is
\[
 (a,b)\longmapsto
 (\tau,\Theta)=(m_ja-\kappa_jb,\kappa_ja+e_j(\tau)).
\]
Its linear matrix and inverse are
\[
 G_j=\begin{pmatrix}m_j&-\kappa_j\\\kappa_j&0\end{pmatrix},
 \qquad
 G_j^{-1}=\begin{pmatrix}0&\kappa_j\\-\kappa_j&m_j\end{pmatrix}.
\]
The remaining shear $(\tau,\Theta)\mapsto
(\tau,\Theta+e_j(\tau))$ is isotopic to the identity through
the same formula with $u e_j$, $0\leq u\leq1$.
In disjoint core collars, the coordinate change
$(p,\Theta)\mapsto(p,\Theta-\chi(u)e_j(\tau))$, with
$\chi=1$ at the boundary and $\chi=0$ at the inner collar edge,
changes the full gluing to $G_j$ and is the identity at $T_0$.
Its inverse adds the same function.  Together with the identity
on the solid tori and cusp cone, it is a homeomorphism of the
corresponding filled targets.  The integer linear gluings have
finite cellular subdivisions, which give a finite CW structure on
$\mathcal C$ with $c_0$ a vertex.  These formulas also preserve the
product orientation on the common annuli, so $\mathcal M$ is orientable.

Let $\rho_j$ be the specified clockwise base loops and let $\eta$
be the positive regular circle generator.  Since $\rho_j$ has
$\tau$-degree $-1$, the inverse matrix gives
\[
 \rho_j\longmapsto\kappa_jl_j,\qquad
 \eta\longmapsto m_jl_j
\]
on inclusion into the filling.  The actual negative meridian
$(a,b)=(-u,0)$ maps to
$(\tau,\Theta)=(-m_ju,-\kappa_ju+e_j(-m_ju))$; its homology
class is $m_j\rho_j-\kappa_j\eta$.  Hence the meridian relations are
\begin{equation}
 3\rho_1-\eta=0,\qquad4\rho_2+\eta=0.
 \hypertarget{tag-036F}{}\label{wb:cusp-cone-Seifert-meridians}
\end{equation}
Reversing the circle coordinate gives generators $\eta'$ and
relations $3\rho_1+\eta'=0$, $4\rho_2-\eta'=0$.
The maps $(s,\theta)\mapsto(s,-\theta)$ on the solid tori and
$(p,\Theta)\mapsto(p,-\Theta)$ on the core conjugate the actions
$\kappa_j\mapsto-\kappa_j$ and the overlaps
$\beta_j\mapsto-\beta_j$.  They extend over the cusp cone and
give a homeomorphism to the reflected target, inducing
$\eta\mapsto-\eta'$ on these presentations.  Thus their integral
homology groups agree.

The core has $H_1=\mathbb Z\langle\rho_1,\rho_2,\eta\rangle$
and $H_2=\mathbb Z\langle\rho_1\times\eta,\rho_2\times\eta\rangle$.
In the Mayer--Vietoris map from the two gluing tori to the core and
solid tori, the $H_2$ map is the identity onto this $\mathbb Z^2$.
On $H_1$, using source basis $(\rho_1,\eta_1,\rho_2,\eta_2)$
and target basis $(\rho_1,\rho_2,\eta,l_1,l_2)$, the matrix is
\[
 D_1=\begin{pmatrix}
 1&0&0&0\\0&0&1&0\\0&1&0&1\\-1&-3&0&0\\0&0&1&-4
 \end{pmatrix}.
\]
Its columns are independent.  On $H_0$ the map is
$(u,v)\mapsto(u+v,-u,-v)$, also injective.
The exact sequence therefore gives $H_3(\mathcal M)=H_2(\mathcal M)=0$
and identifies $H_1(\mathcal M)$ with the cokernel of $D_1$.
Its relations are $\rho_1=l_1$, $\rho_2=-l_2$ and
$\eta=3l_1=4l_2$.  The columns $(3,-4)$ and $(1,-1)$ form a
unimodular basis of $\mathbb Z^2$, so $g=l_1-l_2$ is a free generator:
\begin{equation}
 H_1(\mathcal M;\mathbb Z)
 =\frac{\mathbb Z\langle\eta,\rho_1,\rho_2\rangle}
        {\langle3\rho_1-\eta,\;4\rho_2+\eta\rangle}
 \cong\mathbb Z\langle g\rangle.
 \hypertarget{tag-035M}{}\label{wb:cusp-cone-M-H1}
\end{equation}
Since $\rho_0+\rho_1+\rho_2=0$ in the pair of pants, the remaining
boundary map is
\begin{equation}
 \rho_0=-g,\qquad\rho_1=4g,\qquad
 \rho_2=-3g,\qquad\eta=12g.
 \hypertarget{tag-035O}{}\label{wb:cusp-cone-M-boundary-map}
\end{equation}
Thus $H_1(T_0)\to H_1(\mathcal M)$ is the surjective row $(-1,12)$,
with primitive kernel generated by $12\rho_0+\eta$.

The cone is contractible.  The Mayer--Vietoris sequence for its
attachment identifies $H_3(\mathcal C)$ with $H_2(T_0)=\mathbb Z$
and $H_2(\mathcal C)$ with that primitive kernel.  The same sequence
gives $H_1(\mathcal C)=0$; in degree zero the intersection map is
$1\mapsto(1,-1)$.  The finite CW target has dimension three, so its
higher homology vanishes.  This proves
\eqref{wb:cusp-cone-C-homology} integrally.
Since these groups are free and $c_0$ is a CW basepoint, the reduced
cellular tensor product for the smash has no Tor terms.  Its rank
polynomial is
\[
 (z^2+z^3)^4=z^8+4z^9+6z^{10}+4z^{11}+z^{12},
\]
which proves \eqref{wb:cusp-cone-C-smash-homology}.
\end{proof}

\Needspace{9\baselineskip}
\subsection{Radial embeddings and intrinsic metrics}
\begin{definition}[Radial target-slice maps]
\hypertarget{tag-037N}{}\label{wb:cusp-cone-target-slice-embeddings}
{\normalfont\scriptsize tag 037N}\par
Take \(T\) in the native admissible range, with \(Y_T\) invertible and
\(\ell_T>0\), and impose
\[
 0<r_T=e^{-2\pi T}<\epsilon,
 \qquad T> -\frac{\log\epsilon}{2\pi}.
\]
The second inequality is exactly the upper radial condition.  The native
radial coordinate has interval \([0,\epsilon]\).  The basepoints of
\(\mathcal A,\mathcal B,\mathcal C\) are their radial-zero points.
The two endpoint tori of each suspension are collapsed separately.
The cone collapses its zero torus and attaches its outer torus to
\(\partial\mathcal M=T_0\) by the identity.
The marked angular homeomorphisms and their inverses are
\[
\begin{aligned}
 \iota_T:(\mathbb R/\mathbb Z)^2&\longrightarrow
 Q_T=\mathbb R^2/Y_T\mathbb Z^2,&
 [a]&\longmapsto[Y_Ta],&
 [y]&\longmapsto[Y_T^{-1}y],\\
 \vartheta_T:\mathbb R/\mathbb Z&\longrightarrow
 C_T=\mathbb R/\ell_T\mathbb Z,&
 [a]&\longmapsto[\ell_Ta],&
 [u]&\longmapsto[u/\ell_T].
\end{aligned}
\]
For \(y_1,y_2\in Q_T\), put
\[
 j_T(y_1,y_2)=
 \bigl[[r_T,1,\iota_T^{-1}(y_1)],
       [r_T,1,\iota_T^{-1}(y_2)]\bigr]_{\wedge}
 \in\mathcal A^{\wedge2}.
\]
For \(c_1,\ldots,c_4\in C_T\), put
\[
 k_T^{\mathcal B}(c_1,\ldots,c_4)=
 \bigl[[r_T,1,\vartheta_T^{-1}(c_1)],\ldots,
       [r_T,1,\vartheta_T^{-1}(c_4)]\bigr]_{\wedge}
 \in\mathcal B^{\wedge4}
\]
and
\[
 k_T^{\mathcal C}(c_1,\ldots,c_4)=
 \bigwedge_{j=1}^4[r_T,(1,\vartheta_T^{-1}(c_j))]
 \in\mathcal C^{\wedge4}.
\]
In the last formula the boundary coordinates are related to the boundary
torus \(T_0=\partial D_0\times S^1\) by
\[
 \rho_0(z,[a])=(\epsilon z,e^{2\pi ia}),\qquad
 \rho_0^{-1}(p,w)=(p/\epsilon,\arg(w)/(2\pi)\bmod\mathbb Z).
\]
Here the argument is a class modulo \(2\pi\), so the inverse involves no
choice of a real argument branch.
\end{definition}

\Needspace{6\baselineskip}
\begin{lemma}[Radial embeddings and endpoint maps]
\hypertarget{tag-Q30T}{}\label{mapsIII:lem:radial-slice-embeddings}
{\normalfont\scriptsize tag Q30T}\par
The three maps of Definition~\ref{wb:cusp-cone-target-slice-embeddings}
are topological embeddings for $0<r_T<\epsilon$.
At radius zero all three are constant. At radius $\epsilon$ the
two suspension maps are constant, whereas the cone map remains an embedding.
\end{lemma}
\begin{proof}
Changing \(y\) by \(Y_Tn\) changes \(Y_T^{-1}y\) by \(n\), and changing
\(u\) by \(\ell_Tm\) changes \(u/\ell_T\) by \(m\).  Thus both pairs of
angular maps descend continuously from inverse real linear maps and have
the displayed inverse compositions.

The quotient maps defining the two suspensions restrict to homeomorphisms
on their open cylinders
\[
 (0,\epsilon)\times S^1\times(\mathbb R/\mathbb Z)^2,
 \qquad (0,\epsilon)\times S^1\times(\mathbb R/\mathbb Z).
\]
Indeed, each open cylinder and every open subset of it are open and
saturated in the defining product, and every equivalence class there is
a singleton.  The restrictions are therefore continuous open bijections;
their inverses take the unique interior representatives.  In the
presentation of \(\mathcal C=\mathcal M\cup_{T_0}\operatorname{Cone}_\epsilon(T_0)\),
the subset \((0,\epsilon)\times T_0\), with empty \(\mathcal M\)-part, has the same
open and saturated property.  Its quotient restriction is a homeomorphism
as well.  Composing its inverse with \(\rho_0^{-1}\) gives the displayed
cone coordinates.

Consequently the single-factor maps
\[
 \alpha_T(y)=[r_T,1,\iota_T^{-1}(y)],\quad
 \beta_T(c)=[r_T,1,\vartheta_T^{-1}(c)],\quad
 \gamma_T(c)=[r_T,(1,\vartheta_T^{-1}(c))]
\]
are embeddings into those cylinders.  Their inverses on their images are,
respectively,
\[
 [r_T,1,a]\longmapsto\iota_T(a),\qquad
 [r_T,1,a]\longmapsto\vartheta_T(a),\qquad
 [r_T,(1,a)]\longmapsto\vartheta_T(a).
\]
All of these maps are continuous chart projections followed by the
angular homeomorphisms.  Their images avoid the radial-zero basepoints.

For each of these carriers \((X,x_0)\), the singleton \(\{x_0\}\) is
closed: its full preimage in the defining quotient is exactly the closed
radial-zero torus, with empty \(\mathcal M\)-part in the cone attachment.  Hence the
fat wedge
\[
 F_n=\bigcup_{i=1}^n\operatorname{pr}_i^{-1}(\{x_0\})\subset X^n
\]
is closed.  The smash quotient restricts to the homeomorphism
\[
 q_n:(X\setminus\{x_0\})^n
 \xrightarrow{\ \cong\ }X^{\wedge n}\setminus\{*\}.
\]
To verify the topology explicitly, every open subset of this complement
is open and saturated in \(X^n\); its classes are singletons.  The inverse
of this restriction sends a nonbasepoint smash class to its unique
ordered tuple of nonbasepoint factors.  Therefore
\[
 j_T=q_2\circ(\alpha_T\times\alpha_T),\qquad
 k_T^{\mathcal B}=q_4\circ\beta_T^4,\qquad
 k_T^{\mathcal C}=q_4\circ\gamma_T^4
\]
are embeddings.  On their images the inverses first take that ordered
tuple and then apply the corresponding single-factor inverse above.
This proves both inverse compositions and inverse continuity.

The endpoint maps are as follows. At
\(r=0\) all three formulas take the single smash basepoint value.  At
\(r=\epsilon\), the two suspension formulas are constant at the ordered
tuple of their radial-\(\epsilon\) poles.  The distinct classes
\([0]\) and \([Y_T(1/2,0)]\) in \(Q_T\), and \([0]\) and
\([\ell_T/2]\) in \(C_T\), give explicit failures of injectivity there.
For the cone, \([\epsilon,(1,a)]\) is instead the boundary point
\((\epsilon,e^{2\pi ia})\) in \(\mathcal M\).  The inclusion \(\mathcal M\to\mathcal C\)
is an embedding: an open \(O\subset\mathcal M\) extends to the open saturated
set with \(\mathcal M\)-part \(O\) and cone part
\((\epsilon/2,\epsilon]\times(O\cap T_0)\).  The outer cone circle and its
fourfold smash map therefore remain embeddings.  Finally, passage to
unit cone radius is the homeomorphism \([r,t]\mapsto[r/\epsilon,t]\),
with inverse \([u,t]\mapsto[\epsilon u,t]\), and is the identity on \(\mathcal M\)
for the attachment.  The unit-coordinate slice therefore has radius $r_T/\epsilon$.
\end{proof}



\Needspace{6\baselineskip}
\begin{proposition}[Intrinsic and ambient slice metrics]
\hypertarget{tag-036W}{}\label{wb:cusp-cone-metric-scope}
{\normalfont\scriptsize tag 036W}\par

Let \(D_{12}\), \(D_{24}^{\mathcal B}\), and
\(D_{24}^{\mathcal C}\) be arbitrary fixed metrics compatible with the
topologies of \(\mathcal A^{\wedge2}\), \(\mathcal B^{\wedge4}\), and
\(\mathcal C^{\wedge4}\), respectively.  In their corresponding ambient
metrics, the radial target
slices satisfy
\begin{equation}
 \operatorname{diam}_{\mathrm{amb}}
 \bigl(j_T(Q_T^2)\bigr)\longrightarrow0,\qquad
 \operatorname{diam}_{\mathrm{amb}}
 \bigl(k_T^{\mathcal B}(C_T^4)\bigr)\longrightarrow0,\qquad
 \operatorname{diam}_{\mathrm{amb}}
 \bigl(k_T^{\mathcal C}(C_T^4)\bigr)\longrightarrow0.
 \hypertarget{tag-036I}{}\label{wb:cusp-cone-ambient-collapse}
\end{equation}
Use the constant \(C_{\rm sys}\) fixed in
\cref{wb:cusp-cone-native-descent}.
In the intrinsic quotient metrics,
\begin{equation}
 \operatorname{diam}(Q_T^2)\geq\tfrac12(T-C_{\rm sys})\longrightarrow\infty,\qquad
 \operatorname{diam}(C_T^4)\geq\tfrac12\ell_T\longrightarrow\infty.
 \hypertarget{tag-036S}{}\label{wb:cusp-cone-intrinsic-expansion}
\end{equation}
Thus \cref{wb:cusp-cone-S12-R4-limit} and \cref{wb:cusp-cone-S24-R4-limit} assert
pointed local-isometric limits for the specified flat quotient metrics on the varying slices. For any of the slice embeddings \(e_T\), the ambient pullback metric is
\[
 d_{\mathrm{amb},T}(x,y)=D(e_T(x),e_T(y)).
\]
The identity from the intrinsic slice to this pullback metric is a
homeomorphism. For all sufficiently large \(T\), it is not an isometry,
as the respective diameters tend to infinity and zero.
\end{proposition}



\begin{proof}
Write \(r_T=e^{-2\pi T}\) and \(K=(\mathbb R/\mathbb Z)^4\).
For \(u=(u_1,u_2)\in(\overline T^2)^2=K\), define
\[
 H_{12}(r,u)=\bigl[[r,1,u_1],[r,1,u_2]\bigr]_{\wedge}.
\]
For \(\theta=(\theta_1,\ldots,\theta_4)\in K\), define
\[
 H_{24}^{\mathcal B}(r,\theta)
   =\bigl[[r,1,\theta_1],\ldots,[r,1,\theta_4]\bigr]_{\wedge},
 \qquad
 H_{24}^{\mathcal C}(r,\theta)
   =\bigl[[r,(1,\theta_1)],\ldots,[r,(1,\theta_4)]\bigr]_{\wedge}.
\]
These maps are continuous on \([0,r_0]\times K\), for
\(0<r_0<\epsilon\), and their zero cross-sections have the respective
smash basepoints as their constant values. Their images at \(r_T\)
are exactly the three slices in
\cref{wb:cusp-cone-ambient-collapse}. Indeed, the inverse markings
\[
 \iota_T^{-1}[y]=[Y_T^{-1}y],\qquad
 \vartheta_T^{-1}[b]=[b/\ell_T]
\]
parametrize their entire angular factors; here
\(\vartheta_T[\theta]=[\ell_T\theta]\).

For any one of the three maps \(H\), write \(D\) for the specified
fixed compatible ambient metric and \(*\) for its smash basepoint.
Continuity on the compact parameter space gives
\[
 a(r):=\max_{z\in K}D(H(r,z),*)\longrightarrow0
 \quad\text{as }r\longrightarrow0.
\]
To see the uniformity directly, a sequence \(r_j\to0\) and maximizing
points \(z_j\) has a subsequence with \(z_j\to z\); continuity then
gives \(D(H(r_j,z_j),*)\to D(H(0,z),*)=0\).
The triangle inequality yields
\(\operatorname{diam}_D H(r,K)\le2a(r)\), proving all three limits.

Let \(\Lambda\subset\mathbb R^d\) be a lattice with shortest nonzero
length \(s\), and choose \(\lambda\in\Lambda\) of length \(s\).
For every \(\mu\in\Lambda\),
\[
 |\lambda/2-\mu|\ge s/2.
\]
For \(\mu=0\) this is equality; for \(\mu\ne0\) it follows from
\(|\mu|-|\lambda|/2\ge s/2\).
Thus \([0]\) and \([\lambda/2]\) have quotient distance exactly
\(s/2\). Taking this pair in one factor of \(Q_T^2\), and fixing
the other factor, gives
\(\operatorname{diam}(Q_T^2)\ge\operatorname{sys}(Q_T)/2
\ge(T-C_{\rm sys})/2\). Taking an antipodal pair in one factor of \(C_T^4\)
gives \(\operatorname{diam}(C_T^4)\ge\ell_T/2\).
The bounds and divergence in
\cref{wb:cusp-cone-native-systole,wb:cusp-cone-S24-R4-limit}
therefore give \cref{wb:cusp-cone-intrinsic-expansion}.

For any of the slice embeddings \(e_T\), the pullback is the metric
\[
 d_{\mathrm{amb},T}(x,y)=D(e_T(x),e_T(y)).
\]
The identity on the slice, viewed as a map from its intrinsic metric
to this pullback metric, is a homeomorphism: both metrics give the
specified compact torus topology. Its target diameter tends to zero,
while its source diameter tends to infinity. For every sufficiently
large \(T\) these diameters are unequal, so the identity fails to be
an isometry at every such \(T\). The pointed limits use precisely the
intrinsic metrics specified in the two cited statements.
\end{proof}

\section{Monodromy constraints and sphere maps}
\subsection{Finite-cover and extension obstructions}
\begin{theorem}[Global quotient and topology obstructions]
\hypertarget{tag-0374}{}\label{wb:cusp-cone-obstructions}
{\normalfont\scriptsize tag 0374}\par
For the established sphere \(X\) and its native radial fibre metrics,
take the admissible radial parameters \(T\) with
\(0<r_T=e^{-2\pi T}<\epsilon\), equivalently
\(T>-\log\epsilon/(2\pi)\) within that native admissible range.
The following assertions hold.
\begin{enumerate}
\item Let \(r\geq1\), and let \(N\subset\Lambda^{\oplus r}\) be a primitive
kernel defining a connected fibrewise torus quotient over the full regular
base \(B^\circ\), or over a connected finite cover of \(B^\circ\).
Give the quotient the norm induced by the unrescaled Euclidean metric
on the original marked fibres.  Divergence of its quotient systole on
the radial cusp implies
\begin{equation}
 \operatorname{rank}(\Lambda^{\oplus r}/N)\leq r.
 \hypertarget{tag-037H}{}\label{wb:cusp-cone-rank-bound}
\end{equation}
For two copies this excludes a full-base four-torus quotient with
diverging systole.  The \(S^{12}\) construction of
\cref{wb:cusp-cone-S12-route} uses the quotient on the cusp and collapses
the exterior to the outer suspension pole.

\item The subgroup \(K_\Lambda\) is not a sub-local system over any
connected finite cover of \(B^\circ\), including covers obtained by
branching over the original singular values.  For an integer \(M>1\), the local
change \(t=s^M\) in the degree-\(M^2\) fibre cover cannot extend as a
connected cover of \(B\cong\mathbb P^1\) branched only at \(p_0\).
For \(M=1\), that base change and the local fibre cover are identities.

\item No map from \(X^2\) or \(X\wedge X\) to a fixed \(T^4\) retains on
a regular fibre the four characters of \(Q_T^2\).  No map from \(X^4\)
or \(X^{\wedge4}\) to a fixed \(T^4\) retains the four characters of
the four-circle quotient.  In either route, no map to a fixed \(S^4\)
restricts to the degree-one collapse of the indicated four-torus,
even at one regular fibre.  On the conical targets the exact restriction
maps in degrees \(k=1,4\) are
\[
\begin{aligned}
 j_T^*:H^k(\mathcal A^{\wedge2};\mathbb Z)&\longrightarrow
             H^k(Q_T^2;\mathbb Z),\\
 (k_T^{\mathcal B})^*:H^k(\mathcal B^{\wedge4};\mathbb Z)&\longrightarrow
             H^k(C_T^4;\mathbb Z),\\
 (k_T^{\mathcal C})^*:H^k(\mathcal C^{\wedge4};\mathbb Z)&\longrightarrow
             H^k(C_T^4;\mathbb Z);
\end{aligned}
\qquad\text{all three maps are zero.}
\]
Thus the four primitive character classes and their orientation product
do not extend as global classes on these conical targets.

\item The Cartesian products satisfy
\begin{equation}
 X^2\cong S^6\times S^6\not\simeq S^{12},\qquad
 X^4\cong(S^6)^4\not\simeq S^{24}.
 \hypertarget{tag-0378}{}\label{wb:cusp-cone-products-not-spheres}
\end{equation}
The two- and four-fold synchronized fibre products have real dimensions
\(10\) and \(18\) over the regular base.  Their completed cusp fibres
are singular, and their total completions contain the singular loci
described in the proof.  The join homeomorphisms are
\begin{equation}
 X*X\cong S^{13},\qquad X^{*4}\cong S^{27}.
 \hypertarget{tag-036T}{}\label{wb:cusp-cone-join-dimensions}
\end{equation}

\item Every connected unbranched finite cover of \(X\), \(X^r\) for
\(r\geq1\), \(X\wedge X\cong S^{12}\), or \(X^{\wedge4}\cong S^{24}\)
has degree one.  For \(M>1\), the local degree-\(M^2\) fibre cover cannot
be the restriction of such a cover; for \(M=1\) the identity extends
as the global identity.  The finite flat normal covers in
\cref{maps:src:C.1.5} have the explicit regular-base factorization and
integral-closure completion given below.

\item No smooth manifold homeomorphic to \(S^{12}\) or \(S^{24}\),
with any differentiable structure, admits an almost-complex structure.
In particular these smash sources cannot carry a compact complex
structure.
\end{enumerate}
\end{theorem}

\begin{proof}
Use the native marking
\[
 (g,u,w,\delta)
 =(\widehat\gamma,\widehat u,\widehat w,\widehat\delta),\qquad
 K_\Lambda=\mathbb Z\langle w,\delta\rangle,\qquad
 L_\Lambda=\mathbb Z\langle u,w,\delta\rangle.
\]
On the positive radial ray the physical period matrix is
\[
 P_T=\begin{pmatrix}
 6\alpha_T&\eta_T&1&0\\
 6m_T&L_T&0&0\\
 c_T&\alpha_T&0&1\\
 -q_T&m_T&0&0
 \end{pmatrix},\qquad
 Y_T=\begin{pmatrix}6m_T&L_T\\-q_T&m_T\end{pmatrix},
 \quad D_T=L_Tq_T+6m_T^2>0,
\]
where \(L_T=T+H_T\) and \(q_T=L_T-B_T\). The radial coefficients
\(\alpha_T,\eta_T,m_T,c_T,H_T,B_T\) are bounded.
The marking is \(x_T([a])=[P_Ta]\).  In particular
\(P_Tw=(1,0,0,0)\) and \(P_T\delta=(0,0,1,0)\), both of norm one.

\emph{The global rank bound.}
The quotient is the same physical quotient used in
\cref{wb:cusp-spec-globalization-obstruction}: its real-coordinate
map is $[a]\mapsto[P_T^{\oplus r}a]$, with inverse induced by
$(P_T^{\oplus r})^{-1}$. Rank zero satisfies the bound immediately.
The bounded-period argument in that theorem uses only the displayed
length-one vectors and primitivity, before imposing any monodromy
condition. It therefore applies on every lifted radial ray and gives
\begin{equation}
 K_\Lambda^{\oplus r}\subset N.
 \hypertarget{tag-036U}{}\label{wb:cusp-cone-kernel-contained}
\end{equation}

The left-deck matrices give
\[
 U=A_2M_0A_2^{-1}
   =\begin{pmatrix}
      1&0&0&0\\-6&1&-1&0\\0&0&1&0\\5&0&1&1
     \end{pmatrix}=I+L_0,\qquad L_0^2=0.
\]
Consequently, for every integer \(k_{\rm loop}\),
\[
 U^{k_{\rm loop}}=I+k_{\rm loop}L_0,\qquad
 U^{k_{\rm loop}} w=w-k_{\rm loop}u+k_{\rm loop}\delta,\qquad U^{k_{\rm loop}}\delta=\delta.
\]
The opposite loop has the exact word and matrix
\[
 H_*=A_2A_1=U^{-1}=I-L_0,\qquad
 U=A_2(A_1A_2)^{-1}A_2^{-1}=A_1^{-1}A_2^{-1}.
\]
Thus \(N_*=H_*-I=-L_0\) satisfies \(N_*^2=0\), and
\begin{equation}
 H_*^n w=w+n(u-\delta)\qquad(n\in\mathbb Z).
 \hypertarget{tag-036Q}{}\label{wb:cusp-cone-finite-index-orbit}
\end{equation}
The matrix convention is fixed by
\((x,1)\sim(A^{-1}x,0)\): forward interval transport is the inverse
of the corresponding left-deck operator.  Invariance under either
operator is equivalent to invariance under its inverse.

For a finite-index based covering subgroup \(\mathcal H\), two nonnegative
powers of the loop with matrix \(U\) have the same \(\mathcal H\)-coset.
Their difference is a positive power in \(\mathcal H\).  Thus \(N\) is invariant
under \(U^{k_{\rm loop}}\) for some \(k_{\rm loop}>0\), including for a nonnormal base cover.
Applying \(U^{k_{\rm loop}}\) to \(w_i\in N\) gives \(k_{\rm loop}u_i\in N\); primitivity
then gives \(u_i\in N\).  Hence
\[
 \bigl(\mathbb Z\langle u,w,\delta\rangle\bigr)^{\oplus r}
 \subset N,\qquad
 \operatorname{rank}(\Lambda^{\oplus r}/N)\leq r.
\]
On the original base the same inclusion also follows directly from
\(A_1w=u-w\).
A cover branched over the original singular values restricts to a
connected finite cover of \(B^\circ\), so this argument also applies there.

For $N=L_\Lambda^{\oplus r}$ the quotient is the product of
the character circles of \cref{wb:cusp-spec-globalization-obstruction}.
The map $[P_Ta]\mapsto[\ell_Ta_g]$, with the angle/arclength inverse
$[u]\mapsto[u/\ell_T]$, is precisely that theorem's isometry.
Its product lattice is $\ell_T\mathbb Z^r$ and its systole tends to infinity.

\emph{The local quotient and the base-cover obstruction.}
The local compactification is exactly $q_2$ of
\cref{wb:cusp-cone-S12-route}: its fibre component is
$\overline\pi_f$ and $\pi_T=\iota_T\overline\pi_T$.
That theorem's endpoint values and collar proof give the whole
continuous map, with restriction
$q_2^{\wedge2}|_{F_T^2}=j_T(\pi_T\times\pi_T)$.

For every \(k_{\rm loop}>0\), \(U^{k_{\rm loop}}w\notin K_\Lambda\).  The same finite-index
power argument proves that \(K_\Lambda\) never becomes a sub-local
system over a finite full-base cover.  Its planes
\(\operatorname{span}(\delta,w-k_{\rm loop}u)\) are pairwise distinct:
their projections to the \((u,w)\)-plane have coefficients \((-k_{\rm loop},1)\).

For each positive integer \(M\), the finite-index cover lattice is
related to this kernel by the exact diagram
\[
\begin{array}{ccccccccc}
0&\longrightarrow&MK_\Lambda&\longrightarrow&
 \Lambda_M&\longrightarrow&\Lambda/K_\Lambda&\longrightarrow&0\\
&&\big\downarrow&&\big\downarrow&&\big\Vert\\
0&\longrightarrow&K_\Lambda&\longrightarrow&
 \Lambda&\longrightarrow&\Lambda/K_\Lambda&\longrightarrow&0,
\end{array}
\quad
\Lambda_M=\mathbb Z\langle g,u,Mw,M\delta\rangle.
\]
The two downward maps are inclusions.  On tori they give the
degree-\(M^2\) map
\(\mathbb R^4/\Lambda_M\to\mathbb R^4/\Lambda\), with identity on
the long quotient.  In the basis \(S_M=\operatorname{diag}(1,1,M,M)\),
\[
 M_0^MS_M=S_MM_0,\qquad
 S_M^{-1}A_2S_Mu=(0,0,1/M,0)^{\mathsf T}.
\]
This proves the local changed-base compatibility and the failure on
the original full base for \(M>1\).  A congruence base subgroup does
preserve \(\Lambda_M\), since
\[
 S_M^{-1}(I+M\mathsf C)S_M
 =\begin{pmatrix}I+M\mathsf C_{11}&M^2\mathsf C_{12}\\
                  \mathsf C_{21}&I+M\mathsf C_{22}\end{pmatrix},
\]
and the same formula applies to the inverse.  These explicit maps
retain the separate finite-index and kernel statements.

A connected compact base cover branched only at \(p_0\) restricts,
after removing the finite inverse image of \(p_0\), to a connected
unbranched cover of \(\mathbb P^1\setminus\{p_0\}\cong\mathbb C\).
That punctured source is connected, and the target is simply connected,
so its degree is one.  Thus the local change \(t=s^M\) has the stated
obstruction precisely for \(M>1\).  At \(M=1\), \(S_1=I\) and the
base-change isomorphism \(y\mapsto(y,f_0(y))\), inverse to projection,
identifies the local factors with identities.  The global identity
of \(\mathbb P^1\) extends the base map and has empty branch locus.

\emph{Character and orientation restrictions.}
The marked identifications are
\(\iota_T([a])=[Y_Ta]\) and
\(\vartheta_T([a])=[\ell_Ta]\), with inverses induced by \(Y_T^{-1}\)
and division by \(\ell_T\), respectively.
The fixed-fibre quotient maps and their sections are
\[
\begin{aligned}
 \pi_T([P_Ta])&=[Y_T(a_g,a_u)],&
 \sigma_Q([Y_T(a,b)])&=[P_T(a,b,0,0)],\\
 \pi_{\gamma,T}([P_Ta])&=[\ell_Ta_g],&
 \sigma_C([\ell_Ta])&=[P_T(a,0,0,0)].
\end{aligned}
\]
They are well defined under integral representative changes and satisfy
\(\pi_T\sigma_Q=\operatorname{id}\) and
\(\pi_{\gamma,T}\sigma_C=\operatorname{id}\).
Under \(\iota_T^{-1}\) on each \(Q_T\) factor and \(\vartheta_T^{-1}\)
on each circle factor, the four primitive character pullbacks are
\[
 (g_1^*,u_1^*,g_2^*,u_2^*)\quad\text{and}\quad
 (g_1^*,g_2^*,g_3^*,g_4^*),
\]
respectively.  Their exterior products evaluate to one on
\(\sigma_Q^2\) and \(\sigma_C^4\), so they are nonzero.
In particular, let \(\omega_T=(\iota_T^{-1})^*(x_1\smile x_2)\),
where \(x_1,x_2\) are the ordered primitive characters of \(T^2\).
This is the positive integral orientation class of \(Q_T\), and
\begin{equation}
 (\pi_T\times\pi_T)^*(\omega_T\times\omega_T)
 =g_1^*\smile u_1^*\smile g_2^*\smile u_2^*
 \ne0\quad\text{in }H^4(F_T^2;\mathbb Z).
 \hypertarget{tag-036P}{}\label{wb:cusp-cone-fibre-H4-nonzero}
\end{equation}

For each listed product or smash source \(Z\),
\[
 H^1(Z;\mathbb Z)=H^4(Z;\mathbb Z)=0.
\]
Write \(i_T:F_T\to X\) for the fibre inclusion and
\(q_{X,r}:X^r\to X^{\wedge r}\) for the source smash quotient.
The two-torus route uses \(i_T^2\) and \(q_{X,2}\circ i_T^2\);
the four-circle route uses \(i_T^4\) and \(q_{X,4}\circ i_T^4\).
Composing these maps with \(\sigma_Q^2\) or \(\sigma_C^4\) gives
the corresponding section-torus restrictions.
Hence the pullbacks of the four characters
of any fixed \(T^4\) map restrict to zero.  A degree-one collapse of the
target four-torus pulls the orientation generator of \(S^4\) back to
the exterior product just computed, with the chosen positive
orientations.  Its pullback to the section torus has evaluation one.
Reversing the sphere orientation multiplies this class and its
evaluation by \(-1\), so both choices are nonzero, whereas restriction
of any global degree-four class from \(Z\) is zero.
This proves both nonextension assertions,
also when the restriction is taken directly along the section torus.

For the three conical targets, replace every \(r_T\) in \(j_T\),
\(k_T^{\mathcal B}\), or \(k_T^{\mathcal C}\) by \((1-s)r_T\),
\(0\leq s\leq1\), leaving all angular coordinates fixed.
These are continuous homotopies in the literal radial quotient
presentations.  At \(s=1\) each factor is its radial-zero basepoint,
so each smash map is constant.  Their positive-degree cohomology
restrictions factor through the zero groups of a point.  This proves
the displayed zero maps in degrees one and four.

\emph{Products, synchronized products and joins.}
Use the based diffeomorphism $h:X\to S^6$ and its ordered
product and smash maps from \cref{wb:cusp-spec-globalization-obstruction}.
Their inverses apply $h^{-1}$ to each coordinate and split the
Euclidean coordinate blocks, respectively. The K\"unneth groups and
the resulting homotopy obstructions there prove
\eqref{wb:cusp-cone-products-not-spheres} for these identical sources.

Over the regular base, the synchronized product has coordinates
\((t,z_{11},z_{12},\ldots,z_{r1},z_{r2})\), hence real dimension \(4r+2\).
At two cusp double points the local equation is
\[
 x_1y_1=x_2y_2,
\]
with two further complex coordinates free.  Its complex dimension is
five and its tangent dimension at the crossing is six.  With two
additional smooth factors the equations are
\[
 x_1y_1=x_2y_2,\qquad t_3=x_1y_1,\qquad t_4=x_1y_1.
\]
Eliminating \(t_3,t_4\) gives the same quadric node times six free
complex coordinates, with dimension nine and tangent dimension ten.
These equations exhibit singular loci in both completed products.
The cusp fibres \(W^r\) themselves contain products of \(xy=0\)
and are singular, while the regular-base products are smooth.

The join maps in \cref{wb:cusp-spec-globalization-obstruction}
send $[(\lambda_i,x_i)]$ to $(\sqrt{\lambda_i}\,h(x_i))_i$.
Their inverses use squared block norms and $h^{-1}$ on each nonzero
unit block; the join relation removes every zero-weight coordinate.
For two and four factors these are the homeomorphisms in
\eqref{wb:cusp-cone-join-dimensions}.

\emph{The global covers and their finite normal counterparts.}
The products and smash sources above are simply connected.  For a
connected unbranched cover, choose one point over a basepoint and lift
paths from it.  Simple connectivity makes the endpoint independent of
the path.  Covering charts make this assignment continuous, and
connectedness makes it the two-sided inverse of the cover.  Thus its
degree is one, including every finite cover in the statement.

For the construction of \cref{maps:src:C.1.5}, let its integer cover
parameter be \(N_{\rm cov}\geq1\), and write
\[
 \Gamma=\Lambda\rtimes F(\rho_1,\rho_2),\quad
 G_{N_{\rm cov}}=\langle(A_1\bmod N_{\rm cov},-1\bmod N_{\rm cov}),(A_2\bmod N_{\rm cov},0)\rangle,
\]
and let \(b_{N_{\rm cov}}:F\to G_{N_{\rm cov}}\) be the displayed base homomorphism.
The subgroup defining its regular cover has normal form
\(K_{N_{\rm cov}}=\Lambda_{N_{\rm cov}}\rtimes\ker b_{N_{\rm cov}}\).
Indeed its finite quotient requires trivial \(G_{N_{\rm cov}}\) coordinate and
a translation in the long two-generator submodule modulo \(N_{\rm cov}\).
Let \(\widetilde B^\circ\to B^\circ\) be the universal base cover.
In the common vector cover \(\widetilde B^\circ\times\mathbb C^2\)
the exact factorization is
\[
 \frac{\widetilde B^\circ\times\mathbb C^2}{\Lambda_{N_{\rm cov}}\rtimes\ker b_{N_{\rm cov}}}
 \longrightarrow
 \frac{\widetilde B^\circ\times\mathbb C^2}{\Lambda\rtimes\ker b_{N_{\rm cov}}}
       =X^\circ\times_{B^\circ}B_{N_{\rm cov}}^\circ
 \longrightarrow X^\circ.
\]
Both maps are induced by the identity on the vector cover; their
degrees are \(N_{\rm cov}^2\) and \(|G_{N_{\rm cov}}|\).
For a local component subgroup \(S\leq\mathbb Z^k\), choose
\(m\mathbb Z^k\subset S\).  The completion uses \(x_i=z_i^m\)
and quotient by \(S/m\mathbb Z^k\).  Its invariant algebra is the
integral closure in the given cover's meromorphic-function sheaf.
The unique local closures glue to the finite flat normal map \(Y\to X\).
On its selected cusp component the map is
\((s_v,\zeta_v)\mapsto(N_{\rm cov}s_v,N_{\rm cov}\zeta_v)\), with
\(N_{\rm cov}Z_{N_{\rm cov}}(s_v)=Z(N_{\rm cov}s_v)\), image lattice \(\Lambda_{N_{\rm cov}}\), and base
generator \(\rho_0^{N_{\rm cov}}\).
At \(N_{\rm cov}=1\) these maps and the integral closure are identities.
For \(N_{\rm cov}>1\) the connected completed map has degree
\(N_{\rm cov}^2|G_{N_{\rm cov}}|>1\); the preceding path-lifting inverse excludes its being
globally unbranched.  Its failure occurs over the deleted divisor,
since its restriction over \(X^\circ\) is a covering.

\emph{Almost-complex structures.}
The arbitrary-smoothing assertion is the Borel--Serre application in
\cref{wb:cusp-spec-globalization-obstruction}, from
\cite[\S15, Proposition~15.1 and the following remark,
p.~434]{BorelSerre1953}.
To apply this precise scope, let \(h_{\rm sm}:M_{\rm sm}\to S^{2n}\) be a homeomorphism
from a smooth manifold \(M_{\rm sm}\).  Transport its smooth charts
\((V,\phi)\) to \((h_{\rm sm}(V),\phi\circ h_{\rm sm}^{-1})\).  In this transported
smooth structure \(h_{\rm sm}\) is a diffeomorphism, and an almost-complex
endomorphism \(J\) on \(TM_{\rm sm}\) transports to
\(dh_{\rm sm}\circ J\circ dh_{\rm sm}^{-1}\) on the sphere.
For \(2n=12,24\), the resulting \(n=6,12\) are both in the cited
range, which excludes \(J\).
Finally, multiplication by \(i\) on each holomorphic coordinate
tangent space is preserved by holomorphic transition derivatives;
it gives a global endomorphism with square \(-I\).
Every compact complex structure claimed here would therefore give
the excluded almost-complex structure.
\end{proof}

\subsection{Minimum smash dimensions and compactification maps}
\hypertarget{tag-0376}{}\label{wb:cusp-cone-optimal-strengthening}

\begin{theorem}[Optimality within the native connected torus-quotient class]
\hypertarget{tag-0375}{}\label{wb:cusp-cone-optimal-exponents}
{\normalfont\scriptsize tag 0375}\par
Let \(r\geq1\) be an integer.  Consider connected fibrewise torus quotients
of \(r\) copies of the original marked fibres.  Their kernels are primitive
integral sub-local systems \(N\subset\Lambda^{\oplus r}\), constant in the
marking along a lifted radial ray.  Give every quotient the metric induced
by the original, unscaled Euclidean product metric.  Use the native
admissible parameters \(T\) with
\[
 0<r_T=e^{-2\pi T}<\epsilon,
 \qquad T>-\frac{\log\epsilon}{2\pi},
\]
and take the limit \(T\to+\infty\) within that range.
\begin{enumerate}
\item Over the punctured cusp, divergence of the quotient systole forces
\(\operatorname{rank}(\Lambda^{\oplus r}/N)\leq2r\).
A rank-four quotient therefore requires \(r\geq2\).  The least value
\(r=2\) is attained by the primitive kernel
\[
 K_\Lambda^{\oplus2},\qquad
 K_\Lambda=\mathbb Z\langle\widehat w,\widehat\delta\rangle,
\]
whose physical quotient period lattice is
\[
 L_T^{(12)}=\operatorname{diag}(Y_T,Y_T)\mathbb Z^4.
\]
Its systole equals \(\operatorname{sys}(Y_T\mathbb Z^2)\) and tends
to infinity.  The associated compact sphere source is
\(X^{\wedge2}\cong S^{12}\).

\item Over the full regular base \(B^\circ\), or over any connected finite
cover of \(B^\circ\) of positive degree, divergence of the quotient
systole forces
\(\operatorname{rank}(\Lambda^{\oplus r}/N)\leq r\).
A rank-four quotient therefore requires \(r\geq4\).  The least value
\(r=4\) is attained already over \(B^\circ\) by four copies of the
full-monodromy-invariant character \(\chi_\gamma\), and hence also after
pullback to each such cover.  On the radial fibre the metric quotient is
given by \(\vartheta_T\circ\chi_{\gamma,T}\), where
\[
 \vartheta_T:\mathbb R/\mathbb Z\longrightarrow
 C_T=\mathbb R/\ell_T\mathbb Z,\qquad [a]\longmapsto[\ell_Ta].
\]
Its physical product period lattice and systole are
\[
 L_T^{(24)}=\ell_T\mathbb Z^4,\qquad
 \operatorname{sys}(L_T^{(24)})=\ell_T
 =T+H-B_{\rm rad}+O(T^{-1}).
\]
Here \(B_{\rm rad}=\operatorname{Im}b(0)\) is the limiting radial coefficient.
The associated compact source of the original sphere construction is
\(X^{\wedge4}\cong S^{24}\).
\end{enumerate}
Thus the least sphere dimensions obtained from these original
\(r\)-fold smash sources and rank-four quotients are \(12\) in the cusp
class and \(24\) in the full-base class, with the stated metrics.
\end{theorem}

\begin{proof}
Write
\[
 (g,u,w,\delta)
 =(\widehat\gamma,\widehat u,\widehat w,\widehat\delta),\qquad
 L_\Lambda=\mathbb Z\langle u,w,\delta\rangle.
\]
In physical coordinate order
\((\operatorname{Re}z_1,\operatorname{Im}z_1,
  \operatorname{Re}z_2,\operatorname{Im}z_2)\), the native marking is
\[
 x_T([a])=[P_Ta],\qquad
 P_T=\begin{pmatrix}
 6\alpha_T&\eta_T&1&0\\
 6m_T&L_T&0&0\\
 c_T&\alpha_T&0&1\\
 -q_T&m_T&0&0
 \end{pmatrix},\qquad
 Y_T=\begin{pmatrix}6m_T&L_T\\-q_T&m_T\end{pmatrix},
\]
where \(L_T=T+H_T\), \(q_T=L_T-B_T\),
\(D_T=\det Y_T=L_Tq_T+6m_T^2>0\), and the residual coefficients are bounded
on a sufficiently small radial tail.

\emph{The quotient and its character lattice.}
For a primitive kernel \(N\), put \(N_{\mathbb R}=N\otimes\mathbb R\).
The physical quotient real space and its period lattice are
\[
 V_{N,T}=(\mathbb R^4)^r/P_T^{\oplus r}N_{\mathbb R},
 \qquad
 \bigl\{[P_T^{\oplus r}a]:a\in\Lambda^{\oplus r}\bigr\}.
\]
The map
\[
 (\Lambda^{\oplus r}\otimes\mathbb R)/N_{\mathbb R}
 \longrightarrow V_{N,T},\qquad
 [a]\longmapsto[P_T^{\oplus r}a]
\]
has inverse induced by \((P_T^{\oplus r})^{-1}\).
Primitivity gives
\(N_{\mathbb R}\cap\Lambda^{\oplus r}=N\), so the displayed period
lattice is exactly \(\Lambda^{\oplus r}/N\).  The torus quotient map
is induced by this real quotient; its connected kernel is
\((P_T^{\oplus r}N_{\mathbb R})/(P_T^{\oplus r}N)\).
Its metric is the infimum of the original product norm over each
real coset.

Pullback along the lattice quotient identifies its character lattice
with the integral annihilator
\[
 N^\perp=\{a\in(\Lambda^{\oplus r})^*:a(N)=0\}.
\]
In physical coordinates the character is
\([y]\mapsto a((P_T^{\oplus r})^{-1}y)\pmod{\mathbb Z}\).
A basis of \(N\) extends to an integral basis of
\(\Lambda^{\oplus r}\): lift a basis of the free quotient and use it
together with a basis of \(N\).  The corresponding dual basis proves
\[
 \operatorname{rank}N^\perp
 =\operatorname{rank}(\Lambda^{\oplus r}/N),
 \qquad \operatorname{Ann}(N^\perp)=N
\]
over \(\mathbb Z\).
Rank zero satisfies both asserted upper bounds immediately.  For
positive rank, the systole is the least length of a nonzero quotient
period.

\emph{Rank bounds and attaining quotients.}
The physical marking and primitive kernel are the ones in
\cref{wb:cusp-cone-obstructions}. The proof of
\eqref{wb:cusp-cone-kernel-contained} uses only the bounded short
periods on a radial ray, so it also applies over the punctured cusp.
It gives $K_\Lambda^{\oplus r}\subset N$ and hence rank at most $2r$.
For the full base and every connected finite cover, the same theorem
gives $L_\Lambda^{\oplus r}\subset N$ and rank at most $r$.
Its loop operators satisfy $G_*=A_2A_1=H_*=U^{-1}$, so its all-power
formula is exactly the one for this marking. The dual form is
\[
 (a_i,b_i,0,0)G_*^{j_{\rm loop}}=(a_i+6j_{\rm loop}b_i,b_i,j_{\rm loop}b_i,0);
\]
the integral double-annihilator identification above gives the same
kernel inclusion. The covering argument includes degree one and
nonnormal covers. A branched compact base cover restricts to that
class over $B^\circ$, with metric pulled back from
$T=-\log|t|/(2\pi)$.

For $N=K_\Lambda^{\oplus2}$, the quotient is the product of the
orthogonal quotients of \cref{wb:cusp-cone-native-descent}.
Its map and section are
\[
 [P_Ta]\longmapsto[Y_T(a_g,a_u)],\qquad
 [Y_T(a,b)]\longmapsto[P_T(a,b,0,0)].
\]
The quotient marking is $[a]\mapsto[Y_Ta]$, with inverse induced
by $Y_T^{-1}$. The product lattice and its exact systole are those
proved in \cref{wb:cusp-cone-S12-route}.

For $N=L_\Lambda^{\oplus4}$, use the circle quotient of
\cref{wb:cusp-spec-globalization-obstruction}. The lattice isomorphism
$\Lambda/L_\Lambda\to\mathbb Z$ is induced by
$\gamma=(1,0,0,0)$ and has inverse $n\mapsto[ng]$.
For every affine transition $a\mapsto Aa+b$ its coordinate is
$\gamma(Aa+b)=\gamma(a)+\gamma(b)$; this also gives its pullback
transition over any finite base cover. The metric map is
$[P_Ta]\mapsto[\ell_Ta_g]$ and the inverse arclength marking is
$[u]\mapsto[u/\ell_T]$.
The unit covector $\ell_TP_T^{-\mathsf T}\gamma^{\mathsf T}$ and
its least-norm lift $s\mapsto s\ell_TP_T^{-\mathsf T}\gamma^{\mathsf T}$
give the same quotient isometry as that theorem. Equivalently,
for $G=Y_T^{\mathsf T}Y_T$,
\[
 G_{11}-\frac{G_{12}^2}{G_{22}}
 =\frac{D_T^2}{L_T^2+m_T^2}=\ell_T^2.
\]
The four covectors have disjoint supports, giving exactly
$\ell_T\mathbb Z^4$ and systole $\ell_T$. With
$R_T=\sqrt{L_T^2+m_T^2}$, the earlier circumference expansion has the
explicit remainder
\[
 \ell_T-(L_T-B_T)=\frac{m_T^2}{R_T}
       \left(6-\frac{L_T-B_T}{R_T+L_T}\right).
\]
For bounded $m_T,H_T,B_T$ and $L_T=T+H_T>0$ this is $O(T^{-1})$;
at $m_T=0$ it is zero. This is the stated unscaled quotient metric.

\emph{Attained minima and compact sphere maps.}
The inequalities for rank four give \(r\geq2\) and \(r\geq4\),
and the constructed quotients attain the two endpoints.
For \(r>2\) in the cusp class one may retain the first two
quotients using
\(K_\Lambda^{\oplus2}\oplus\Lambda^{\oplus(r-2)}\).
For \(r>4\) in the full-base class use
\(L_\Lambda^{\oplus4}\oplus\Lambda^{\oplus(r-4)}\).
At the minimal values the kernels are forced: \(N/K_\Lambda^{\oplus2}\),
respectively \(N/L_\Lambda^{\oplus4}\), is a rank-zero subgroup of
the corresponding free quotient, and is therefore zero.

Set $c_{\mathcal A}=q_2$ and
$c_{\mathcal B}=q_{1,\mathrm{cusp}}$ from
\cref{wb:cusp-cone-S12-route,wb:cusp-cone-S24-local-route}.
Thus $v_\epsilon=v_1$ and $b_\epsilon=b_1$ denote their outer poles.
These are identities of the whole maps, including $W$, the outer
boundary and the exterior; their continuity and surjectivity were
proved there. Since $\varepsilon_\gamma=\operatorname{pr}_1p$,
the following map satisfies $c_{\mathcal B}=\kappa c_{\mathcal A}$.
The product formula
\[
 \kappa:\mathcal A\to\mathcal B,\qquad
 [s,z,(a_g,a_u)]\longmapsto[s,z,a_g]
\]
respects both collapsed endpoints and is continuous by the quotient
property. Its section is $[s,z,a_g]\mapsto[s,z,(a_g,0)]$;
composing with $\kappa$ is the identity.

For the full-base construction let
\[
 \mathcal C=\mathcal M\cup_{T_0}\operatorname{Cone}_\epsilon(T_0),
 \qquad T_0=\partial D_0\times S^1,
\]
and let \(c_{\mathcal C}=q_1:X\to\mathcal C\) be the based quotient
of Theorem~\ref{wb:cusp-cone-S24-route}.
It includes the core \(P\times S^1\), both finite Seifert fillings
\(C_1,C_2\), and the cusp cone.  Its defining finite overlap identity
\(\chi_\gamma h_j=h_j^{S^1}\overline\chi_j\) gives the gluing
to \(\mathcal M\); on the cusp it is
\[
 x\longmapsto[|t|,(\epsilon\,t/|t|,\theta_\gamma(x))],
 \qquad W\longmapsto c_0.
\]
There is an explicit based quotient
\[
 p_{\mathcal C\mathcal B}:\mathcal C\to\mathcal B,\qquad
 \mathcal M\longmapsto b_\epsilon,\qquad
 [s,(\epsilon z,\theta)]\longmapsto[s,z,\theta].
\]
At \(s=\epsilon\) the cone formula is the same constant as the core
formula, and at \(s=0\) it maps the cone vertex to \(b_0\).
The pushout property therefore proves continuity.
On \(0<s<\epsilon\) its inverse coordinates are
\([s,z,\theta]\mapsto[s,(\epsilon z,\theta)]\);
the two poles correspond to \(c_0\) and the collapsed \(\mathcal M\).
These formulas induce a continuous bijection
\(\mathcal C/\mathcal M\to\mathcal B\), hence a homeomorphism
by compactness and the Hausdorff property.
The definitions on the cusp, on \(W\), and on the exterior give the
equality of global based maps
\[
 p_{\mathcal C\mathcal B}c_{\mathcal C}=c_{\mathcal B}.
\]
In particular the two finite fillings are retained in
\(c_{\mathcal C}\) and are sent to the outer pole by
\(p_{\mathcal C\mathcal B}\).

For a radial slice put
\[
 \alpha_{T,\rm sl}(y)=[r_T,1,\iota_T^{-1}(y)]\in\mathcal A,\qquad
 \beta_T(z)=[r_T,1,\vartheta_T^{-1}(z)]\in\mathcal B,\qquad
 \zeta_T(z)=[r_T,(\epsilon,\vartheta_T^{-1}(z))]\in\mathcal C.
\]
Then \(p_{\mathcal C\mathcal B}\zeta_T=\beta_T\), and, writing
\(i_T:F_T\to X\) for the fibre inclusion,
\[
 c_{\mathcal A}i_T=\alpha_{T,\rm sl}\pi_T,\qquad
 c_{\mathcal B}i_T=\beta_T\pi_{\gamma,T},\qquad
 c_{\mathcal C}i_T=\zeta_T\pi_{\gamma,T}.
\]
With the product-to-smash quotients \(q_{X,r}:X^r\to X^{\wedge r}\)
and the interior carrier embeddings of
Definition~\ref{wb:cusp-cone-target-slice-embeddings}, these identities give
\[
 \begin{aligned}
 c_{\mathcal A}^{\wedge2}q_{X,2}i_T^{\times2}
   &=j_T(\pi_T\times\pi_T),\\
 c_{\mathcal B}^{\wedge4}q_{X,4}i_T^{\times4}
   &=k_T^{\mathcal B}\pi_{\gamma,T}^{\times4},\\
 c_{\mathcal C}^{\wedge4}q_{X,4}i_T^{\times4}
   &=k_T^{\mathcal C}\pi_{\gamma,T}^{\times4},\qquad
 p_{\mathcal C\mathcal B}^{\wedge4}k_T^{\mathcal C}=k_T^{\mathcal B}.
 \end{aligned}
\]
The based maps descend to smash powers because their product maps
send the source fat wedge to the target fat wedge.

Let $h_r:X^{\wedge r}\to S^{6r}$ be the ordered smash
homeomorphism of \cref{wb:cusp-spec-globalization-obstruction},
using the same based $h:X\to S^6$. Its finite-coordinate map
concatenates the $r$ Euclidean blocks; its inverse splits those
blocks and applies $h^{-1}$, and the two basepoints correspond.
The compact source maps are
\[
 c_{\mathcal A}^{\wedge2}h_2^{-1}:S^{12}\to\mathcal A^{\wedge2},
 \qquad
 c_{\mathcal C}^{\wedge4}h_4^{-1}:S^{24}\to\mathcal C^{\wedge4},
\]
and the latter composed with \(p_{\mathcal C\mathcal B}^{\wedge4}\)
is \(c_{\mathcal B}^{\wedge4}h_4^{-1}\).
Their radial restrictions are precisely the two attaining rank-four
quotients above.  Thus the least factor numbers give the claimed
dimensions \(12\) and \(24\).  Different values of \(6r\) have distinct
top reduced integral homology degrees, so these dimension comparisons
also respect the stated sphere identifications.
\end{proof}

\subsection{Collapsed sphere targets and their intrinsic metrics}
\begin{theorem}[Sphere quotient maps with moving four-spheres]
\hypertarget{tag-036J}{}\label{wb:cusp-cone-angle-forgotten-sphere-carriers}
{\normalfont\scriptsize tag 036J}\par
Choose $x_0\in W$.  There are continuous surjective based quotient maps
\[
 \widehat q_3:(X,x_0)\longrightarrow(S^3,*),\qquad
 \widehat q_\gamma:(X,x_0)\longrightarrow(S^2,*).
\]
Their smash powers, under the based sphere identifications, are
continuous surjective quotient maps
\begin{equation}
 \widehat q_3^{\wedge2}:S^{12}\twoheadrightarrow S^6,\qquad
 \widehat q_\gamma^{\wedge4}:S^{24}\twoheadrightarrow S^8.
 \hypertarget{tag-037L}{}\label{wb:cusp-cone-sphere-to-sphere-maps}
\end{equation}

Put $\overline T^2=\mathbb R^2/\mathbb Z^2$ and
\[
 K^{(1)}=((\mathbb R/\mathbb Z)\times\{0\})
          \cup(\{0\}\times(\mathbb R/\mathbb Z)),\qquad
 Z_2=\overline T^2/K^{(1)}.
\]
The quotient $c:\overline T^2\to Z_2$ has degree one when $Z_2\cong S^2$
is oriented by the image of the oriented open two-cell.
For the native radial parameters with $0<r_T=e^{-2\pi T}<\epsilon$, set
\begin{equation}
 \begin{gathered}
 K_T^{(1)}=\iota_T(K^{(1)}),\qquad
 S_T^2=Z_2\cong Q_T/K_T^{(1)},\\
 c_T=c\circ\iota_T^{-1}:Q_T\longrightarrow S_T^2.
 \end{gathered}
 \hypertarget{tag-037O}{}\label{wb:cusp-cone-transported-skeleton-collapse}
\end{equation}
The displayed identification is induced by $\iota_T[a]=[Y_Ta]$.
In this theorem $c_T$ denotes this quotient map.
Write
\[
 \pi_{\gamma,T}=\vartheta_T\circ\chi_{\gamma,T},\qquad
 \vartheta_T[a]=[\ell_Ta]\in C_T=\mathbb R/\ell_T\mathbb Z.
\]

At height $r_T$ in the reduced suspensions, the based sections give
topological embeddings
\[
 \mathbb S_{12,T}^4=S_T^2\wedge S_T^2\hookrightarrow S^6,\qquad
 \mathbb S_{24,T}^4=C_T^{\wedge4}\hookrightarrow S^8.
\]
Give $S_T^2$ the quotient length metric of the Euclidean torus $Q_T$,
give $(S_T^2)^2$ the Euclidean $\ell^2$ product length metric, and give
its smash the quotient length metric under $\rho_{12,T}$.
Give $C_T^4$ its Euclidean product length metric and its smash the
quotient length metric under $\rho_{24,T}$.
Here $\rho_{12,T}$ and $\rho_{24,T}$ are the product-to-smash quotient maps.
The radial restrictions are
\begin{align}
 F_T^2&\xrightarrow{\ \pi_T\times\pi_T\ }Q_T^2
 \xrightarrow{\ c_T\times c_T\ }(S_T^2)^2
 \xrightarrow{\ \rho_{12,T}\ }\mathbb S_{12,T}^4,
 \hypertarget{tag-035R}{}\label{wb:cusp-cone-S12-S4-radial-map}\\
 F_T^4&\xrightarrow{\ \pi_{\gamma,T}^{\times4}\ }C_T^4
 \xrightarrow{\ \rho_{24,T}\ }\mathbb S_{24,T}^4.
 \hypertarget{tag-0361}{}\label{wb:cusp-cone-S24-S4-radial-map}
\end{align}
The basepoint fibres of the maps from $Q_T^2$ and $C_T^4$ are exactly
\begin{equation}
 W_{12,T}=(K_T^{(1)}\times Q_T)\cup(Q_T\times K_T^{(1)}),
 \hypertarget{tag-035T}{}\label{wb:cusp-cone-S12-exact-collapsed-locus}
\end{equation}
\begin{equation}
 W_{24,T}=\{(u_1,u_2,u_3,u_4):u_j=0\pmod{\ell_T}
                       \text{ for some }j\}.
 \hypertarget{tag-0362}{}\label{wb:cusp-cone-S24-exact-collapsed-locus}
\end{equation}
Each map is injective off its basepoint fibre.  The full radial maps
have respective basepoint fibres
\[
 (\pi_T^{-1}(K_T^{(1)})\times F_T)
 \cup(F_T\times\pi_T^{-1}(K_T^{(1)}))
\]
and
\[
 \{(x_1,x_2,x_3,x_4):\pi_{\gamma,T}(x_j)=0\text{ in }C_T
                                      \text{ for some }j\}.
\]
Off these fibres, equality of final images is equivalent to equality
under the corresponding first arrow.  Every fibre of a first arrow
is contained in a fibre of the full radial composite.

Put $y_T=\iota_T[(1/2,1/2)]$, $z_T=c_T(y_T)$, and
\[
 \zeta_{12,T}=[z_T,z_T],\qquad
 \zeta_{24,T}=[\ell_T/2,\ell_T/2,\ell_T/2,\ell_T/2].
\]
On a fixed sufficiently small radial tail choose
$C\geq\sup_T\|E_T\|_{\mathrm{op}}$, where
$Y_T=T\bigl(\begin{smallmatrix}0&1\\-1&0\end{smallmatrix}\bigr)+E_T$.
For the two row covectors $l_{j,T}$ of $Y_T^{-1}$,
\begin{equation}
 d_T:=\operatorname{dist}_{Q_T}(y_T,K_T^{(1)})
     =\frac{1}{2\max_j\|l_{j,T}\|}
     \geq\frac{T-C}{2}.
 \hypertarget{tag-036M}{}\label{wb:cusp-cone-distance-to-skeleton}
\end{equation}
For each $R>0$, the parameters on that tail satisfying $T-C>4R$ give
an isometry onto the entire indicated ball,
\begin{equation}
 \rho_{12,T}\circ(c_T\times c_T):
 B_R^{Q_T^2}((y_T,y_T))
 \xrightarrow{\ \cong_{\mathrm{isom}}\ }
 B_R^{\mathbb S_{12,T}^4}(\zeta_{12,T}).
 \hypertarget{tag-035W}{}\label{wb:cusp-cone-S12-onto-exact-ball}
\end{equation}
The parameters satisfying $\ell_T>4R$ similarly give
\begin{equation}
 \rho_{24,T}:B_R^{C_T^4}((\ell_T/2,\ell_T/2,\ell_T/2,\ell_T/2))
 \xrightarrow{\ \cong_{\mathrm{isom}}\ }
 B_R^{\mathbb S_{24,T}^4}(\zeta_{24,T}).
 \hypertarget{tag-0369}{}\label{wb:cusp-cone-S24-onto-exact-ball}
\end{equation}
All four balls are Euclidean four-balls.  Consequently
\begin{equation}
 (\mathbb S_{12,T}^4,\zeta_{12,T})\longrightarrow(\mathbb R^4,0),
 \qquad
 (\mathbb S_{24,T}^4,\zeta_{24,T})\longrightarrow(\mathbb R^4,0)
 \hypertarget{tag-036Y}{}\label{wb:cusp-cone-moving-S4-R4-limits}
\end{equation}
in pointed local-isometric and pointed Gromov--Hausdorff topology,
and smoothly on the eventually Euclidean basepoint neighbourhoods.
\end{theorem}

\begin{proof}
We use the reduced suspension
\[
 \Sigma_{\mathrm{red},\epsilon}Z
 =([0,\epsilon]\times Z)/
   ((\{0,\epsilon\}\times Z)\cup([0,\epsilon]\times\{*\})).
\]
The characteristic square gives a based homeomorphism
$[0,1]^2/\partial[0,1]^2\to Z_2$; its inverse off the basepoint
chooses the two representatives in $(0,1)$.
On the interior of $[0,1]^d$, the coordinates
\[
 u_j=\tan\bigl(\pi(a_j-1/2)\bigr),\qquad
 a_j=\frac12+\frac{\arctan u_j}{\pi}
\]
identify its interior with $\mathbb R^d$.  Inverse stereographic
projection sends the collapsed boundary to the point at infinity
and gives $[0,1]^d/\partial[0,1]^d\cong S^d$.
The map $c$ sends the oriented torus two-cell to this quotient
two-cell with coefficient one, proving its stated degree.
The maps induced by $\iota_T$ and $\iota_T^{-1}$ give the displayed
inverse homeomorphisms $Z_2\leftrightarrow Q_T/K_T^{(1)}$.

The reduced suspensions of $Z_2$ and $\mathbb R/\mathbb Z$
are respectively $[0,1]^3/\partial[0,1]^3$ and
$[0,1]^2/\partial[0,1]^2$, using $r/\epsilon$ as first coordinate.
Coordinate concatenation on product cubes identifies their smash
products with the spheres of the sums of the dimensions; the inverse
splits the nonboundary coordinates into the corresponding blocks.

For the maps $q_2$ and $q_{1,\mathrm{cusp}}$ of
\cref{wb:cusp-cone-S12-route,wb:cusp-cone-S24-local-route}, define
\[
 \widehat p_2:\mathcal A\to\Sigma_{\mathrm{red},\epsilon}Z_2,
 \quad[r,z,a]\mapsto[r,c(a)],
\]
\[
 \widehat p_1:\mathcal B\to
       \Sigma_{\mathrm{red},\epsilon}(\mathbb R/\mathbb Z),
 \quad[r,z,\theta]\mapsto[r,\theta].
\]
Both end sections in each source map to the reduced-suspension
basepoint.  The formulas therefore descend to continuous surjections.
Set $\widehat q_3=\widehat p_2q_2$ and
$\widehat q_\gamma=\widehat p_1q_{1,\mathrm{cusp}}$.
They are based, and compactness of $X$ makes them quotient maps.
Their smash powers are surjective by choosing preimages in each factor
and are quotient maps by the same compactness argument.
The based sphere identifications in
\cref{wb:cusp-cone-S12-route,wb:cusp-cone-S24-local-route}
and the target cube maps above prove
\eqref{wb:cusp-cone-sphere-to-sphere-maps}.

For $0<r<\epsilon$, the height section $j_r(z)=[r,z]$ is a
based continuous injection into the reduced suspension, hence an
embedding.  Its smash powers are injective: equality away from the
smash basepoint means equality in each factor, while the basepoint
has precisely the fat wedge as preimage.  They are therefore
embeddings of the two displayed four-spheres.
For the circle section use $\vartheta_T^{-1}[u]=[u/\ell_T]$.
The pullback of $du^2$ is $\ell_T^2\,d\theta^2$.
The radial formulas are
\[
 \widehat q_3|_{F_T}=j_{r_T}c_T\pi_T,\qquad
 \widehat q_\gamma|_{F_T}
       =j_{r_T}\vartheta_T^{-1}\pi_{\gamma,T}.
\]
Their products followed by the smash quotients give the two
factorizations in the statement.  The multiplicative character
coordinate is $e^{2\pi i\theta}$; its basepoint $1$ corresponds to
$0\in\mathbb R/\mathbb Z$ under this map and its inverse.

Since $c_T^{-1}(*)=K_T^{(1)}$ and $c_T$ is injective elsewhere,
the basepoint fibre of $\rho_{12,T}(c_T\times c_T)$ is $W_{12,T}$,
and this composite is injective on its complement.  The definition
of the fourfold smash gives the same assertions for $W_{24,T}$.
Taking inverse images under the first arrows proves all the
full-source fibre and equality assertions.  It also gives
homeomorphisms
\[
 Q_T^2/W_{12,T}\cong\mathbb S_{12,T}^4,\qquad
 C_T^4/W_{24,T}\cong\mathbb S_{24,T}^4.
\]
Their inverses off the basepoint recover the unique uncollapsed
coordinates and send the basepoint to the corresponding collapsed set.

For the metric comparison use
\cref{mapsIII:lem:collapsed-subset-distance} on
$Q_T/K_T^{(1)}$, on the smash of its twofold product, and on
$Q_T^2/W_{12,T}$ and $C_T^4/W_{24,T}$. In each case the map
to that lemma's quotient is $[x]\mapsto[x]$; off the collapsed
point it takes the unique uncollapsed tuple, and
the collapsed subset represents the basepoint. The lemma supplies
the quotient distance, length and quotient-topology assertions.

For $x,y\in Q_T^2$, write
\[
 d_i=d_{Q_T}(x_i,y_i),\quad
 \alpha_i=d_{Q_T}(x_i,K_T^{(1)}),\quad
 \beta_i=d_{Q_T}(y_i,K_T^{(1)}),\quad
 A=\min_i \alpha_i,\quad B=\min_i \beta_i.
\]
The two successive quotient metrics give
\[
 \min\left\{
 \sqrt{\sum_{i=1}^2\min\{d_i,\alpha_i+\beta_i\}^2},A+B\right\}
 =\min\left\{\sqrt{d_1^2+d_2^2},A+B\right\}.
\]
Every branch in the first square root that uses $\alpha_i+\beta_i$ is at
least $A+B$; the remaining branch uses both $d_i$.
The right side is exactly the single-collapse metric of $Q_T^2/W_{12,T}$,
because the distance to $W_{12,T}$ is the smaller of the two
coordinate distances to $K_T^{(1)}$.  Thus the displayed homeomorphism
and its inverse preserve distances.  The four-circle quotient already
has the single-collapse metric of $C_T^4/W_{24,T}$.

The inverse image of $K_T^{(1)}$ in $\mathbb R^2$ is the union of
the hyperplanes $l_{j,T}(y)\in\mathbb Z$.  Orthogonal projection onto
$l_{j,T}(y)=k$ is
\[
 y\longmapsto y-
       \frac{l_{j,T}(y)-k}{\|l_{j,T}\|^2}\,l_{j,T}^{\mathsf T}.
\]
Consequently
\[
 \operatorname{dist}_{Q_T}([y],K_T^{(1)})
 =\min_j\frac{\operatorname{dist}(l_{j,T}(y),\mathbb Z)}
                    {\|l_{j,T}\|}.
\]
At $y=Y_T(1/2,1/2)^{\mathsf T}$ this is
$1/(2\max_j\|l_{j,T}\|)$.  The bound
$\|Y_Tv\|\geq(T-C)\|v\|$ and
$\|l_{j,T}\|\leq\|Y_T^{-1}\|_{\mathrm{op}}$
prove \eqref{wb:cusp-cone-distance-to-skeleton}.
The two product midpoints have respective distances
$d_T$ and $\ell_T/2$ from $W_{12,T}$ and $W_{24,T}$.
The product period lattices are
$\operatorname{diag}(Y_T,Y_T)\mathbb Z^4$ and $\ell_T\mathbb Z^4$;
their systoles are at least $T-C$ and equal to $\ell_T$, respectively.

In either construction write $b$ for the product midpoint, $W_{\rm col}$ for
the collapsed subset, and $s$ for the product systole.  At the
stated thresholds, $s>4R$ and $d_0=d(b,W_{\rm col})>2R$.
Choose its displayed lift $\widetilde b\in\mathbb R^4$.
For $u,v\in B_R(0)$ and a nonzero period $\lambda$,
\[
 \|u-v+\lambda\|\geq s-2R>2R\geq\|u-v\|.
\]
The covering map $u\mapsto[\widetilde b+u]$ therefore identifies
the Euclidean radius-$R$ ball with the entire product-torus ball.
Its inverse chooses the unique lift of norm at most $R$ relative to
$\widetilde b$; a closest lift exists by discreteness of the lattice.
For $x,y$ in that ball,
\[
 d(x,W_{\rm col})+d(y,W_{\rm col})\geq2(d_0-R)>2R\geq d(x,y).
\]
The quotient-distance formula proves that the collapse preserves
every distance on the ball.  Conversely, for a target point $[y]$
within distance $R$ of $[b]$, the term $d_0+d(y,W_{\rm col})$ exceeds $R$.
Thus the same formula forces $d(b,y)\leq R$; the basepoint itself
is at distance $d_0>R$.  This proves surjectivity onto the entire
target ball and both asserted isometries.
With one $Q_T$ factor, the same estimates give an onto isometry
\[
 c_T:B_R^{Q_T}(y_T)\longrightarrow B_R^{S_T^2}(z_T)
\]
of Euclidean two-balls whenever $T-C>4R$.
On the corresponding open coordinate neighbourhoods the pullback
metric tensor is the Euclidean tensor.  Since $T-C$ and $\ell_T$
tend to infinity, these exact ball identifications prove all three
convergences in \eqref{wb:cusp-cone-moving-S4-R4-limits}.
\end{proof}

\subsection{Homotopy classes and suspension factorizations}
\begin{remark}[Homotopy and graded invariants of the sphere maps]
\hypertarget{tag-037K}{}\label{wb:cusp-cone-sphere-map-homotopy-status}
{\normalfont\scriptsize tag 037K}\par
Use the based sphere identifications of
\cref{wb:cusp-cone-angle-forgotten-sphere-carriers}, and write
\[
 \begin{aligned}
 G_{12}&=\mu_{3,2}\widehat q_3^{\wedge2}\mu_{6,2}^{-1}:
                  S^{12}\longrightarrow S^6,\\
 G_{24}&=\mu_{2,4}\widehat q_\gamma^{\wedge4}\mu_{6,4}^{-1}:
                  S^{24}\longrightarrow S^8.
 \end{aligned}
\]
Here \(\mu_{n,r}:(S^n)^{\wedge r}\to S^{nr}\) is the
ordered characteristic-cube concatenation of that theorem, with
inverse given by splitting the nonboundary coordinates into blocks
and sending the boundary class to the smash basepoint.

The ordered quotient basis \((\widehat\gamma,\widehat u)\)
gives the mutually inverse torus maps
\[
 \beta:[a\widehat\gamma+b\widehat u]\longmapsto[a,b],\qquad
 \beta^{-1}:[a,b]\longmapsto[a\widehat\gamma+b\widehat u].
\]
They preserve the coordinate circle subgroups and hence induce
inverse skeleton-quotient maps and reduced-suspension maps.
In these coordinates the lattice projection is
\(p(a,b,c,d)=(a,b)\), and
\(\varepsilon_\gamma=\operatorname{pr}_1p\).
Let $\alpha$ denote the long-period map of
Definition~\ref{mapsIII:def:cusp-suspension-maps}, viewed in its
quotient lattice: a lift with coefficients $(a,b,c,d)$ has
$\alpha=[a\widehat\gamma+b\widehat u]$. The displayed $\beta$
and $\beta^{-1}$ give its coefficient coordinates. Consequently
\(\overline\pi_f=\beta\alpha\) and
\(\theta_\gamma=\operatorname{pr}_1\beta\alpha\).
On the punctured cusp the defining formulas are therefore
\[
 \widehat q_3(x)=[|t|,c(\beta\alpha(x))],\qquad
 \widehat q_\gamma(x)=[|t|,\operatorname{pr}_1\beta\alpha(x)].
\]
Both maps equal the suspension basepoint on \(W\) and outside
the cusp interior. These formulas identify them on all of \(X\)
with the two maps of Definition~\ref{mapsIII:def:cusp-suspension-maps},
through the displayed target homeomorphisms.
Proposition~\ref{mapsIII:angle-forgotten-one-copy-null} supplies their
based null-homotopies, fixing the chosen \(x_0\in W\), and the based
null-homotopy of every positive smash power. Thus
\[
 [\widehat q_3]=0\in\pi_6(S^3),\qquad
 [\widehat q_\gamma]=0\in\pi_6(S^2),\qquad
 [G_{12}]=0\in\pi_{12}(S^6),\qquad
 [G_{24}]=0\in\pi_{24}(S^8).
\]

The positive-degree integral homology maps with a nonzero source
or target group are precisely
\[
 \begin{aligned}
 (G_{12})_* &: H_6(S^{12})=0\longrightarrow H_6(S^6)=\mathbb Z,\\
 (G_{12})_* &: H_{12}(S^{12})=\mathbb Z\longrightarrow H_{12}(S^6)=0,\\
 (G_{24})_* &: H_8(S^{24})=0\longrightarrow H_8(S^8)=\mathbb Z,\\
 (G_{24})_* &: H_{24}(S^{24})=\mathbb Z\longrightarrow H_{24}(S^8)=0.
 \end{aligned}
\]
Every displayed map is zero; all other positive-degree homology maps
have both groups zero. The integral cohomology pullbacks are zero
in every positive degree as well: the potentially nonzero source
groups \(H^6(S^6)\) and \(H^8(S^8)\) map into
\(H^6(S^{12})=0\) and \(H^8(S^{24})=0\), respectively.
In degree zero, all four induced maps are the identity on \(\mathbb Z\).
The ordinary Hurewicz maps are
\[
 \begin{aligned}
 h&:\pi_{12}(S^6)\longrightarrow H_{12}(S^6;\mathbb Z)=0,\\
 h&:\pi_{24}(S^8)\longrightarrow H_{24}(S^8;\mathbb Z)=0,
 \end{aligned}
\]
and send these two classes to \((G_{12})_*[S^{12}]=0\) and
\((G_{24})_*[S^{24}]=0\).

For maps \(g:M^n\to N^n\) between connected closed oriented
\(n\)-manifolds, ordinary degree is defined by
\(g_*[M]=\deg(g)[N]\) in \(H_n(N;\mathbb Z)=\mathbb Z[N]\).
Here the source fundamental classes map into
\(H_{12}(S^6)=0\) and \(H_{24}(S^8)=0\), while the target
fundamental classes generate \(H_6(S^6)\) and \(H_8(S^8)\).
These are the exact graded pushforwards; the ordinary degree
definition assigns no integer degree to either of the two maps.
\end{remark}

\Needspace{6\baselineskip}
\begin{remark}[Squares of the rank-two sphere carrier]
\hypertarget{tag-035Z}{}\label{wb:cusp-cone-S12-sphere-map-null}
{\normalfont\scriptsize tag 035Z}\par
Taking $x=y=g$ in \cref{wb:cusp-s12-full-carrier-null} shows that
$g\wedge g:S^{12}\to S^6$ is based null-homotopic for every based
map $g:S^6\to S^3$.  With $g=\widehat q_3$ in the based sphere
identifications of \cref{wb:cusp-cone-sphere-to-sphere-maps}, this is
the continuous surjective quotient
$\widehat q_3^{\wedge2}:S^{12}\twoheadrightarrow S^6$.
\end{remark}

\begin{remark}[Suspension factorization of the circle carrier]
\hypertarget{tag-036E}{}\label{wb:cusp-cone-S24-sphere-map-null}
{\normalfont\scriptsize tag 036E}\par
Let $f=\widehat q_\gamma:S^6\to S^2$ in the based coordinates of
\cref{wb:cusp-cone-angle-forgotten-sphere-carriers}. The identification
in \cref{wb:cusp-cone-sphere-map-homotopy-status} applies
\cref{mapsIII:angle-forgotten-one-copy-null} to this entire map, giving
\begin{equation}
 [f^{\wedge r}]=0\quad\text{in }\pi_{6r}(S^{2r})\qquad(r\geq2),
 \hypertarget{tag-0367}{}\label{wb:cusp-cone-S24-null-maps}
\end{equation}
and in particular
\begin{equation}
 [\widehat q_\gamma^{\wedge2}]=0\in\pi_{12}(S^4),\qquad
 [\widehat q_\gamma^{\wedge4}]=0\in\pi_{24}(S^8).
 \hypertarget{tag-0368}{}\label{wb:cusp-cone-S24-null-special-cases}
\end{equation}
Here the sphere-coordinate map is
$F_r=\mu_{2,r}f^{\wedge r}\mu_{6,r}^{-1}$, with
$\mu_{n,1}=\operatorname{id}$ and $F_1=f$.
The coordinate concatenations $\mu$ and their block-splitting inverses
are those of \cref{wb:cusp-cone-angle-forgotten-sphere-carriers}.

The two-factor map has the factorization
\begin{equation}
 \begin{tikzcd}[column sep=large]
 S^6\wedge S^6\arrow[r,"f\wedge\operatorname{id}_{S^6}"]&
 S^2\wedge S^6\arrow[r,"\operatorname{id}_{S^2}\wedge f"]&
 S^2\wedge S^2.
 \end{tikzcd}
 \hypertarget{tag-036C}{}\label{wb:cusp-cone-S24-smash-square-factorization}
\end{equation}
Write $\beta_{p,q}$ for two-block concatenation and
$\tau_{p,q}[x,y]=[y,x]$, whose inverse is $\tau_{q,p}$. On every
representative, including the fat wedge,
\begin{equation}
 \begin{aligned}
 f\wedge\operatorname{id}_{S^6}
   &=\tau_{6,2}(\operatorname{id}_{S^6}\wedge f)\tau_{6,6},\\
 \beta_{6,2}(\operatorname{id}_{S^6}\wedge f)\beta_{6,6}^{-1}
   &=\Sigma^6f:S^{12}\longrightarrow S^8.
 \end{aligned}
 \hypertarget{tag-036B}{}\label{wb:cusp-cone-S24-sixfold-suspension}
\end{equation}
The first composite sends $[x,y]$ successively to $[y,x]$,
$[y,f(x)]$ and $[f(x),y]$. The intermediate homotopy group is
\phantomsection\hypertarget{tag-033R}{}\label{qa:III:layout:11}
\begin{equation}
 \pi_{12}(S^8)\cong\pi_4^{\mathrm S}=0,
 \hypertarget{tag-036R}{}\label{wb:cusp-cone-fourth-stable-stem}
\end{equation}
by \cite[Corollary~4.24, p.~360]{Hatcher2002} and
\cite[Corollary~1.3, p.~112, and Table~1, p.~113]{IsaksenWangXu2023}.
The suspension range holds since $12+j<2(8+j)-1$ for every $j\geq0$.

Let $H_4$ be the based homotopy supplied by
\cref{mapsIII:angle-forgotten-one-copy-null}, transported through the
specified source and target homeomorphisms. Its initial map is the
continuous surjective quotient $F_4$ of
\cref{wb:cusp-cone-angle-forgotten-sphere-carriers}. With $\kappa_T$
the native radial product inclusion followed by source coordinates,
\[
 H_4(\kappa_T(x),0)
   =i_T\rho_{24,T}\bigl(\pi_{\gamma,T}(x_1),\ldots,
                         \pi_{\gamma,T}(x_4)\bigr),
 \qquad \pi_{\gamma,T}=\vartheta_T\chi_{\gamma,T}.
\]
Here $i_T$ is the fourth smash of $[u]\mapsto[r_T,[u/\ell_T]]$
in the chosen target coordinates. The quotient metric, midpoint
$[\ell_T/2,\ldots,\ell_T/2]$, and entire-ball isometries for
$\ell_T>4R$ are exactly those of
\eqref{wb:cusp-cone-S24-onto-exact-ball}. Thus the radial
restrictions of this initial representative have the stated pointed
local-isometric and Gromov--Hausdorff limits.
\end{remark}

\subsection{Open four-cells and moving spheres}
\begin{corollary}[Open four-cells in the sphere targets]
\hypertarget{tag-036V}{}\label{wb:cusp-cone-literal-R4-open-cells}
{\normalfont\scriptsize tag 036V}\par
Use the quotients, metrics and height sections of
\cref{wb:cusp-cone-angle-forgotten-sphere-carriers}.
For every admissible $T$, with $0<r_T<\epsilon$, put
\[
 U_{12,T}=\mathbb S_{12,T}^4\setminus\{*\},\qquad
 U_{24,T}=\mathbb S_{24,T}^4\setminus\{*\}.
\]
The restrictions of the specified quotient maps are homeomorphisms
\begin{equation}
 \rho_{12,T}\circ(c_T\times c_T):
 (Q_T\setminus K_T^{(1)})^2\xrightarrow{\ \cong\ }U_{12,T},
 \hypertarget{tag-036G}{}\label{wb:cusp-cone-U12-homeomorphism}
\end{equation}
\begin{equation}
 \rho_{24,T}:(C_T\setminus\{0\})^4
          \xrightarrow{\ \cong\ }U_{24,T}.
 \hypertarget{tag-036H}{}\label{wb:cusp-cone-U24-homeomorphism}
\end{equation}
Each $U_{\nu,T}$ is homeomorphic to $\mathbb R^4$ and is
topologically embedded in its ambient sphere, $S^6$ for $\nu=12$
and $S^8$ for $\nu=24$.  It is open in its moving four-sphere and
nonclosed both there and in its ambient sphere.  Its ambient closure
is precisely $\mathbb S_{\nu,T}^4$, and it is not open in that
ambient sphere.

For each fixed admissible $T$, the restriction to $U_{\nu,T}$ of the
assigned compact quotient metric is not globally isometric to
Euclidean $\mathbb R^4$.  The same assertion holds for the restriction
of any compatible metric on its compact ambient sphere.
For the assigned quotient metrics, every fixed radius-$R$ Euclidean
four-ball occurs isometrically as the entire midpoint ball in
$U_{12,T}$ whenever $T-C>4R$, and in $U_{24,T}$ whenever
$\ell_T>4R$, with the tail constant $C$ of
\cref{wb:cusp-cone-angle-forgotten-sphere-carriers}.
The isometries are the restrictions of the two displayed quotient maps,
after the marked Euclidean covering charts at the specified midpoints.
Both thresholds hold eventually as $T\to+\infty$.
\end{corollary}

\begin{proof}
The basepoint fibres of the two quotient maps are
\[
 W_{12,T}=(K_T^{(1)}\times Q_T)\cup(Q_T\times K_T^{(1)}),\qquad
 W_{24,T}=\{(u_1,u_2,u_3,u_4):u_j=0\text{ for some }j\}.
\]
They identify no distinct points off those fibres.  If a closed subset
$A$ of $E$ is collapsed to one point, the restriction
$E\setminus A\to(E/A)\setminus\{*\}$ is a homeomorphism:
it is a continuous bijection, and every open subset of $E\setminus A$
is an open saturated subset of $E$.  Apply this observation to the
quotient identifications of
\cref{wb:cusp-cone-angle-forgotten-sphere-carriers}.
This proves both displayed homeomorphisms with their prescribed maps.

Here are coordinates and inverses.  The maps
\[
 \begin{aligned}
 u_T:(0,1)^2&\longrightarrow Q_T\setminus K_T^{(1)},\quad
        a\longmapsto[Y_Ta],\\
 v_T:(0,1)&\longrightarrow C_T\setminus\{0\},\quad
        a\longmapsto[\ell_Ta].
 \end{aligned}
\]
are homeomorphisms.  The inverse of $u_T$ takes the two representatives
of $Y_T^{-1}y$ in $(0,1)$; changing the lift $y$ changes these
coordinates by integers.  The inverse of $v_T$ chooses the representative
in $(0,\ell_T)$ and divides by $\ell_T$.
Set
\[
 \eta(t)=\frac12\left(1+\frac{t}{\sqrt{1+t^2}}\right),\qquad
 \eta^{-1}(a)=\frac{2a-1}{2\sqrt{a(1-a)}}.
\]
The derivative of $\eta$ is positive, its limits are $0$ and $1$,
and $4\eta(t)(1-\eta(t))=(1+t^2)^{-1}$ proves the displayed inverse.
The homeomorphisms from $\mathbb R^4$ to the two open cells are
\[
 t\longmapsto\rho_{12,T}
   \bigl(c_Tu_T(\eta(t_1),\eta(t_2)),
         c_Tu_T(\eta(t_3),\eta(t_4))\bigr)
\]
and
\[
 t\longmapsto\rho_{24,T}
   \bigl([\ell_T\eta(t_1)],\ldots,[\ell_T\eta(t_4)]\bigr).
\]
Their inverses recover the uncollapsed coordinates by the two quotient
homeomorphisms, use $u_T^{-1}$ or $v_T^{-1}$, and then apply
$\eta^{-1}$ to each coordinate.

The compact height-section embeddings of
\cref{wb:cusp-cone-angle-forgotten-sphere-carriers} restrict to the
claimed embeddings.  Their images in the interiors of the ambient
cube quotients $[0,1]^6/\partial[0,1]^6$ and
$[0,1]^8/\partial[0,1]^8$ have coordinates
\[
 (r_T/\epsilon,a_1,a_2,r_T/\epsilon,a_3,a_4)
\]
and
\[
 (r_T/\epsilon,a_1,r_T/\epsilon,a_2,
                 r_T/\epsilon,a_3,r_T/\epsilon,a_4),
 \qquad a_j\in(0,1).
\]
The inverse on either image projects to the four $a_j$ coordinates.
Varying the first radial coordinate by nonzero quantities tending to
zero, while keeping it in $(0,1)$, gives ambient points outside the
moving sphere converging to each point of the open cell.  Thus neither
open cell is open in its ambient sphere.

Each open cell is the complement of a closed point in its moving
four-sphere.  The coordinates $a_1=1/N$, $a_2=a_3=a_4=1/2$
give points in it tending to that omitted point: the first torus
coordinate reaches $K_T^{(1)}$, or the first circle coordinate reaches
$0$, under the defining continuous quotient.  Hence the open cells
are nonclosed.  Since each compact moving four-sphere is closed in
its ambient sphere, its whole image is exactly the ambient closure
of the corresponding open cell.

For the metric assertion, choose representatives of the four marked
coordinate differences in $[-1/2,1/2]$.  The Euclidean product
diameters satisfy
\[
 \operatorname{diam}(Q_T^2)\leq\|Y_T\|_{\mathrm{op}},\qquad
 \operatorname{diam}(C_T^4)\leq\ell_T.
\]
The quotient maps decrease distance.  Their restricted open-cell
metrics therefore have these finite diameter bounds.  Any metric
compatible with the topology of the compact ambient sphere also has
finite diameter.  A surjective isometry from $\mathbb R^4$ would
preserve the unbounded distances between $0$ and $(n,0,0,0)$, which
contradicts each finite bound.

Finally let
\[
 \widetilde y_T=Y_T(1/2,1/2)^{\mathsf T},\qquad
 b_{12,T}=(\widetilde y_T,\widetilde y_T),\qquad
 b_{24,T}=(\ell_T/2,\ell_T/2,\ell_T/2,\ell_T/2).
\]
The Euclidean covering charts are $u\mapsto[b_{12,T}+u]$ in
$Q_T^2$ and $u\mapsto[b_{24,T}+u]$ in $C_T^4$.
Their compositions with the two quotient maps are the onto isometries
of \eqref{wb:cusp-cone-S12-onto-exact-ball} and
\eqref{wb:cusp-cone-S24-onto-exact-ball}.
The inverses first recover the unique uncollapsed product point and
then its unique lift within Euclidean distance $R$ of $b_{\nu,T}$.
The midpoint distances from the collapsed loci are at least
$(T-C)/2$ and equal to $\ell_T/2$, respectively, so these entire
balls lie in the open cells at the stated thresholds.  The same
isometries apply to closed or open radius-$R$ balls.  The period
growth $\ell_T\to\infty$ and the fixed residual bound $C$ prove
the final eventuality assertion.
\end{proof}


\begin{lemma}[Null-homotopies of the moving four-spheres]
\hypertarget{tag-036Z}{}\label{wb:cusp-cone-moving-S4-null-homotopy}
{\normalfont\scriptsize tag 036Z}\par
For every native parameter with \(0<r_T=e^{-2\pi T}<\epsilon\),
the height sections
\[
 a_T:S_T^2\longrightarrow S^3,\qquad
 b_T:C_T\longrightarrow S^2
\]
of \cref{wb:cusp-cone-angle-forgotten-sphere-carriers} are based
null-homotopic. Their smash powers defining the moving four-spheres,
\[
 I_{12,T}=a_T^{\wedge2}:S_T^2\wedge S_T^2\longrightarrow S^6,
 \qquad
 I_{24,T}=b_T^{\wedge4}:C_T^{\wedge4}\longrightarrow S^8,
\]
are based null-homotopic. The induced integral homology maps are
\[
 \begin{aligned}
 (I_{12,T})_* &:H_4(S_T^2\wedge S_T^2;\mathbb Z)
                    \longrightarrow H_4(S^6;\mathbb Z)=0,\\
 (I_{24,T})_* &:H_4(C_T^{\wedge4};\mathbb Z)
                    \longrightarrow H_4(S^8;\mathbb Z)=0.
 \end{aligned}
\]
Thus the fundamental classes of both moving four-spheres have zero
image in their ambient targets.
\end{lemma}

\begin{proof}
Use the fixed reduced-suspension models
\[
 S^3=\Sigma_{\mathrm{red},\epsilon}Z_2,\qquad
 S^2=\Sigma_{\mathrm{red},\epsilon}(\mathbb R/\mathbb Z),
 \qquad Z_2=S_T^2,
\]
with both end sections and the based meridian collapsed to one point.
The exact sections, in the markings of the cited theorem, are
\[
 a_T(z)=[r_T,z],\qquad
 b_T([u])=[r_T,[u/\ell_T]],\qquad C_T=\mathbb R/\ell_T\mathbb Z.
\]
For \(0\leq s\leq1\), define
\[
 A_T(z,s)=[(1-s)r_T,z],\qquad
 B_T([u],s)=[(1-s)r_T,[u/\ell_T]].
\]
These are continuous compositions with the suspension quotients.
Replacing \(u\) by \(u+n\ell_T\), \(n\in\mathbb Z\), adds
an integer to its circle coordinate, so the second formula is
well defined. At \(s=0\) they are the specified sections, at
\(s=1\) they are constant on the collapsed zero-height section,
and at every time they take their source basepoints to the
collapsed based meridian.

The two smash homotopies are
\[
 ([z_1,z_2],s)\longmapsto[A_T(z_1,s),A_T(z_2,s)],
\]
\[
 ([u_1,u_2,u_3,u_4],s)\longmapsto
 [B_T([u_1],s),B_T([u_2],s),B_T([u_3],s),B_T([u_4],s)].
\]
Every source fat wedge maps to the target basepoint throughout.
The product-to-smash quotient, multiplied by the time interval, is
a continuous surjection from a compact space to a Hausdorff space,
and hence a quotient map. The displayed continuous product formulas
therefore descend to jointly continuous based homotopies on the
whole smash products. Their initial maps are \(I_{12,T}\) and
\(I_{24,T}\), and their terminal maps are constant.

The characteristic-square and characteristic-interval identifications
in the cited theorem orient both source smashes as four-spheres.
Homotopy invariance identifies their homology pushforwards with
those of constant maps, which factor through
\(H_4(\{*\};\mathbb Z)=0\). This proves the fundamental-class assertions.
\end{proof}

\section{Compact carriers and Euclidean boundary fibres}
\subsection{Conical carriers and their homotopy types}
\begin{proposition}[Compact cones on the four-torus]
\hypertarget{tag-036O}{}\label{wb:cusp-cone-explicit-carriers}
{\normalfont\scriptsize tag 036O}\par
Fix $T_0$ sufficiently large that $r_0=e^{-2\pi T_0}<\epsilon$,
and define
\[
 \mathcal D_{12}
 =\frac{[0,r_0]\times(\overline T^2)^2}
        {\{0\}\times(\overline T^2)^2\sim d_{12}},
 \qquad
 \mathcal D_{24}
 =\frac{[0,r_0]\times(\mathbb R/\mathbb Z)^4}
        {\{0\}\times(\mathbb R/\mathbb Z)^4\sim d_{24}}.
\]
Both spaces are homeomorphic to $\operatorname{Cone}(T^4)$, with
the outer height-$r_0$ torus retained.  There are topological embeddings
\begin{align}
 \delta_{12}:\mathcal D_{12}&\hookrightarrow\mathcal A\wedge\mathcal A,
 & [r,u_1,u_2]&\longmapsto
       \bigl[[r,1,u_1],[r,1,u_2]\bigr],
 \hypertarget{tag-035G}{}\label{wb:cusp-cone-D12-embedding}\\
 \delta_{24}:\mathcal D_{24}&\hookrightarrow\mathcal C^{\wedge4},
 & [r,\theta_1,\ldots,\theta_4]&\longmapsto
       \bigwedge_{j=1}^4[r,(1,\theta_j)].
 \hypertarget{tag-035I}{}\label{wb:cusp-cone-D24-embedding}
\end{align}
Here $1$ fixes the positive radial base angle, and the fibre circle
coordinate $\theta\in\mathbb R/\mathbb Z$ corresponds to
$e^{2\pi i\theta}\in S^1$.

For $T\geq T_0$ put $r_T=e^{-2\pi T}$ and
$\vartheta_T[\theta]=[\ell_T\theta]\in C_T$.
The maps
\[
 \begin{aligned}
 k_{12,T}(q_1,q_2)
   &=[r_T,\iota_T^{-1}(q_1),\iota_T^{-1}(q_2)],\\
 k_{24,T}([b_1],\ldots,[b_4])
   &=[r_T,[b_1/\ell_T],\ldots,[b_4/\ell_T]]
 \end{aligned}
\]
identify the two flat quotient tori with the corresponding complete
height-$r_T$ cross-sections.  Under the embeddings above, the sphere
quotient maps restrict exactly to
\begin{align}
 F_T^2&\xrightarrow{\ \pi_T\times\pi_T\ }Q_T^2
       \xrightarrow{\ k_{12,T}\ }\{r_T\}\times T^4
       \subset\mathcal D_{12},
 \hypertarget{tag-035H}{}\label{wb:cusp-cone-D12-slice}\\
 F_T^4&\xrightarrow{\ (\vartheta_T\circ\chi_{\gamma,T})^{\times4}\ }C_T^4
       \xrightarrow{\ k_{24,T}\ }\{r_T\}\times T^4
       \subset\mathcal D_{24}.
 \hypertarget{tag-035J}{}\label{wb:cusp-cone-D24-slice}
\end{align}
Each of these embedded four-torus slices is null-homotopic in its
carrier and in the full smash target.

For $\nu\in\{12,24\}$ the local integral homology at the vertex is
\begin{equation}
 H_k(\mathcal D_\nu,\mathcal D_\nu\setminus\{d_\nu\};\mathbb Z)
 \cong
 \begin{cases}
  \mathbb Z^4,&k=2,\\
  \mathbb Z^6,&k=3,\\
  \mathbb Z^4,&k=4,\\
  \mathbb Z,&k=5,\\
  0,&\text{otherwise}.
 \end{cases}
 \hypertarget{tag-036K}{}\label{wb:cusp-cone-carrier-local-homology}
\end{equation}
These groups preclude an integral homology-manifold structure at
$d_\nu$, with or without boundary, in every dimension.  They also
preclude a topological manifold chart there, with or without boundary,
in every dimension.
\end{proposition}

\begin{proof}
Set $T^4=\mathbb R^4/\mathbb Z^4$.  Concatenating the two marked
pairs gives the isomorphism
\[
 c_{12}:(\overline T^2)^2\longrightarrow T^4,
 \qquad ([a_1,a_2],[a_3,a_4])\longmapsto[a_1,a_2,a_3,a_4].
\]
Its inverse regroups the four coordinates into pairs.  For the other
link use $c_{24}([\theta_1],\ldots,[\theta_4])
=[\theta_1,\ldots,\theta_4]$, whose inverse takes the four separate
residue classes.  The circle coordinate isomorphism to $S^1$ has
inverse $z\mapsto[\arg z/(2\pi)]$.
Use $\operatorname{Cone}(T^4)=([0,1]\times T^4)/(\{0\}\times T^4)$.
The cone homeomorphism for each carrier and its inverse are
\[
 [r,a]\longmapsto[r/r_0,c_\nu(a)],\qquad
 [s,b]\longmapsto[r_0s,c_\nu^{-1}(b)].
\]
These formulas respect the zero-end collapse and are continuous
in both directions.  Each carrier is compact Hausdorff, since its
defining equivalence relation is the diagonal together with the
square of the closed zero-end torus.

The cylinder formula defining $\delta_{12}$ is continuous into
$\mathcal A^2$ before its smash quotient.  At height zero both
entries are $v_0$, so the formula descends to the carrier.  For
$0<r\leq r_0<\epsilon$, both entries are interior points of
$\mathcal A$ and avoid its basepoint.  Equality of their smash
images therefore gives equality of the two interior coordinate
pairs, recovering $r,u_1,u_2$.  These positive-height images
also avoid the smash basepoint.  Thus $\delta_{12}$ is injective,
including at $d_{12}$.

The same cylinder construction defines $\delta_{24}$ continuously.
At positive height at most $r_0$, all four factors lie in the
interior of the cusp cone of $\mathcal C$, with their unique
coordinates $[r,(1,\theta_j)]$.  They avoid $c_0$ and the outer
attaching torus.  Equality in the smash consequently recovers $r$
and all four circle coordinates, proving injectivity on the whole
carrier.  The two smash targets are Hausdorff: they are quotients
of compact Hausdorff products by their closed fat wedges.
Each $\delta_\nu$ is therefore a continuous injection from a compact
space into a Hausdorff space, and hence a topological embedding.

For the cross-sections, the inverse of
$\iota_T[a]=[Y_Ta]$ is $[y]\mapsto[Y_T^{-1}y]$, and the inverse
of $\vartheta_T$ is $[b]\mapsto[b/\ell_T]$.
These formulas are well defined on these period lattices,
and give the displayed maps $k_{\nu,T}$ and their inverse maps
onto the full cross-sections.  Write $i_{12,T}$ and $i_{24,T}$
for the source product-to-smash embeddings.  The exact identities
of Theorems~\ref{wb:cusp-cone-S12-route} and
\ref{wb:cusp-cone-S24-route} give
\[
 \begin{aligned}
 q_2^{\wedge2}\circ i_{12,T}
   &=\delta_{12}\circ k_{12,T}\circ(\pi_T\times\pi_T),\\
 q_1^{\wedge4}\circ i_{24,T}
   &=\delta_{24}\circ k_{24,T}\circ
                    (\vartheta_T\circ\chi_{\gamma,T})^{\times4}.
 \end{aligned}
\]
Indeed, using
$\overline\pi_T=\operatorname{pr}_{\overline T^2}\circ
\overline\pi|_{F_T}$ and $\pi_T=\iota_T\overline\pi_T$,
both sides of the first equality have factors
$[r_T,1,\overline\pi_T(x)]$.  Both sides of the second equality
have factors $[r_T,(1,\chi_{\gamma,T}(x))]$, since
$\vartheta_T^{-1}\circ(\vartheta_T\circ\chi_{\gamma,T})
=\chi_{\gamma,T}$.
This proves the asserted restrictions as exact maps; the factor
epimorphisms attain the entire indicated slices.

The formula
\[
 H_\nu([r,a],s)=[(1-s)r,a],\qquad 0\leq s\leq1,
\]
defines a continuous contraction of the carrier fixing its vertex.
Continuity follows from its continuous cylinder formula and the
product of the defining compact quotient with $[0,1]$, which is
again a quotient.  At $s=0$ it is the identity and at $s=1$ it is
constant at $d_\nu$.  The homotopy
$H_\nu(k_{\nu,T}(q),s)$ contracts each stated slice; composing with
$\delta_\nu$ contracts its image in the full smash target.
This also covers $T=T_0$, whose slice is the outer carrier torus.

The punctured carrier has coordinates $(0,r_0]\times T^4$.
The homotopy $[r,a]\mapsto[(1-s)r+sr_0,a]$ retracts it onto the
outer torus while staying at positive height.  The pair exact
sequence and the contraction of the full carrier give
\[
 \partial:H_k(\mathcal D_\nu,
        \mathcal D_\nu\setminus\{d_\nu\};\mathbb Z)
        \xrightarrow{\ \cong\ }H_{k-1}(T^4;\mathbb Z),\qquad k\geq2.
\]
Both spaces are connected and their $H_0$ inclusion is the identity,
so the relative groups in degrees zero and one vanish.
The product CW structure on $T^4$ has one cell for every subset
of its four marked circle directions and has zero differential.
Thus $H_j(T^4;\mathbb Z)=\Lambda^j\mathbb Z^4$ with that subset
basis.  Orienting each cone cell so that its outer boundary is the
corresponding torus cell, $\partial$ sends its relative class to
that basis element.  This proves every group and vanishing in
\eqref{wb:cusp-cone-carrier-local-homology}.

At an interior point of an integral homology $n$-manifold, the local
integral homology is $\mathbb Z$ in degree $n$ and zero otherwise;
at a boundary point it vanishes in all degrees.  The four nonzero
degrees computed above satisfy neither pattern for any $n$.
For a topological manifold chart, excision identifies its local
homology with that of $(\mathbb R^n,\mathbb R^n\setminus\{0\})$,
or of the corresponding half-space pair at a boundary point.
The punctured Euclidean space retracts onto $S^{n-1}$ for $n\geq1$;
for $n=0$ the pair is a point relative to the empty set.
The half-space and its punctured version are contractible, the latter
retracting onto a hemisphere.  These pairs have precisely the
interior and boundary homology patterns just stated.
A chart induces isomorphisms on local homology, so the computed
vertex groups exclude every such chart.
\end{proof}

\begin{theorem}[Homotopy types of the conical targets]
\hypertarget{tag-037M}{}\label{wb:cusp-cone-target-homotopy-types}
{\normalfont\scriptsize tag 037M}\par

The smash maps are continuous surjective quotient maps
\[
 q_2^{\wedge2}:S^{12}\twoheadrightarrow\mathcal A^{\wedge2},
 \qquad
 q_1^{\wedge4}:S^{24}\twoheadrightarrow\mathcal C^{\wedge4}.
\]
Their targets have homotopy types
\begin{align}
 \mathcal A^{\wedge2}
 &\simeq
 \bigvee^9S^4\vee\bigvee^{18}S^5\vee
 \bigvee^{15}S^6\vee\bigvee^6S^7\vee S^8,
 \hypertarget{tag-0357}{}\label{wb:cusp-cone-A-smash-wedge}\\
 \mathcal C^{\wedge4}
 &\simeq
 S^8\vee\bigvee^4S^9\vee\bigvee^6S^{10}
 \vee\bigvee^4S^{11}\vee S^{12}.
 \hypertarget{tag-035F}{}\label{wb:cusp-cone-C-smash-wedge}
\end{align}
Both maps induce zero on every positive-degree integral homology and
cohomology group and have zero ordinary Hurewicz image.  Neither map has a
continuous homotopy section.  These conclusions do not determine whether
the unstable homotopy class of either map is zero.
\end{theorem}



\begin{proof}
Choose the basepoint of $X$ in the cusp fibre. Each target is a
compact Hausdorff finite CW space, with its indicated basepoint a
vertex. The maps $q_2$ and $q_1$ are based continuous surjections.
For a point of a target smash product, choose a preimage of each
coordinate in $X$. This proves surjectivity of both smash maps.
Their sources are compact, and the target fat wedges are closed,
so the target smash products are Hausdorff. Both maps are therefore
closed quotient maps. The source homeomorphisms and their inverses are the ordered
two- and four-block maps of
\cref{wb:cusp-cone-S12-route,wb:cusp-cone-S24-local-route},
based at the same point of $W$.

The target $\mathcal A$ is the unreduced suspension of $T^3$.
Let $I$ be its interval with all circle coordinates at their
basepoints, joining the two poles. The quotient
\[
 c_A:\mathcal A\longrightarrow\Sigma T^3=\mathcal A/I
\]
is based at the zero pole, and $\Sigma$ on the right denotes reduced
suspension. Contract $I$ to that pole and extend this contraction
to $H_t:\mathcal A\to\mathcal A$, starting at the identity,
by the homotopy extension property. The map $H_1$ descends to
$h_A:\Sigma T^3\to\mathcal A$. The homotopy $H_t$ proves
$h_Ac_A\simeq\mathrm{id}$, and $c_AH_t$ descends to the homotopy
$c_Ah_A\simeq\mathrm{id}$. Both homotopies fix the basepoint.

For each nonempty subset $S\subseteq\{1,2,3\}$, let
$p_S:T^3\to\bigwedge_{i\in S}S^1=S^{|S|}$ be coordinate
projection followed by the smash quotient, with the product
orientation. List the seven subsets by cardinality and then
lexicographically. Pinch the suspension coordinate into seven
equal subintervals and apply $\Sigma p_S$ to the corresponding
summand. This defines a based map
\[
 F:\Sigma T^3\longrightarrow
 W_A=\bigvee^3S^2\vee\bigvee^3S^3\vee S^4.
\]
On the $j$th subinterval its formula is
$[t,z]\mapsto[7t-j,p_{S_j}(z)]$; its subdivision levels and
base line map to the wedge point. On the suspended homology
classes of the coordinate subtori, this map has identity matrices
in degrees two, three and four. Indeed the oriented quotient
$T^S\to\bigwedge_{i\in S}S^1$ has degree one; the component
indexed by a different subset of the same size has a constant
circle coordinate and degree zero. Both spaces are simply
connected, so the homological Whitehead theorem
\cite[Corollary~4.33, pp.~367--368]{Hatcher2002} makes $F$ a based
homotopy equivalence. Thus $Fc_A$ is the required equivalence from
$\mathcal A$ to $W_A$, with inverse $h_AK$ for a based homotopy
inverse $K$ of $F$.

For $\mathcal C$, put $a=\rho_1$, $b=\rho_2$ and $h=\eta$.
The pair-of-pants product has fundamental group generated by
$a,b,h$, with $h$ central. In the two solid-torus fillings the
boundary maps send $(\rho_1,h)$ to $(\ell_1,3\ell_1)$ and
$(\rho_2,h)$ to $(-\ell_2,4\ell_2)$. Their kernels give
\[
 \pi_1(\mathcal M)=
 \langle a,b,h\mid[a,h]=[b,h]=1,\ a^3=h,\ b^4=h^{-1}\rangle.
\]
The boundary torus coned in $\mathcal C$ has generators
$\rho_0=(ab)^{-1}$ and $h$. Van Kampen therefore imposes
$h=1$ and $b=a^{-1}$. The remaining relations are $a^3=a^4=1$,
and $a=a^4(a^3)^{-1}=1$. Consequently $\pi_1(\mathcal C)=1$.

The integral boundary calculation in
\cref{wb:cusp-cone-S24-route} identifies $H_1(\mathcal M)$
with $\mathbb Z\tau$, the boundary map with
$(\rho_0,h)\mapsto(-\tau,12\tau)$, and all higher homology of
$\mathcal M$ with zero. The cone exact sequence identifies
\[
 H_2(\mathcal C)=\mathbb Z u,\qquad
 H_3(\mathcal C)=\mathbb Z v,
\]
where $u$ corresponds to the primitive boundary kernel
$12\rho_0+h$ and $v$ to the oriented fundamental class of the
boundary torus. All other reduced homology groups vanish.

Apply \cref{mapsIII:carrier-three-complex-wedge} with
$C=\mathcal C$, based at its cusp vertex, and with the ordered
generators $u,v$ just computed. The comparison of the complex in
that lemma with this target is the identity.
Denote the lemma's sphere maps for $u$ and $v$ by
$\alpha:S^2\to\mathcal C$ and $\beta:S^3\to\mathcal C$,
respectively. Their wedge equivalence is
\[
 G=\alpha\vee\beta:S^2\vee S^3\longrightarrow\mathcal C,
 \qquad G_*[S^2]=u,\quad G_*[S^3]=v.
\]
In the lemma's second mapping cylinder its inverse is $L=kj_C$.
The specified homotopies $kU_t$ and $pT_sj_C$ give
$LG\simeq\operatorname{id}_{S^2\vee S^3}$ and
$GL\simeq\operatorname{id}_{\mathcal C}$ through based maps. Thus
\[
 L_*u=[S^2],\qquad L_*v=[S^3].
\]
These are the same primitive boundary-kernel and oriented-torus
generators under the cone exact sequence above.

The smashes $(Fc_A)^{\wedge2}$ and $G^{\wedge4}$, with inverse
smashes of $h_AK$ and $L$, realize the target equivalences in the
statement. Their homotopies descend because they fix the
basepoints. Finite smash products distribute over finite wedges:
the summand inclusions give the map to the smash, and each
non-basepoint pair belongs to unique summands, giving its inverse.
These continuous inverse maps identify compact Hausdorff
quotients. The homeomorphism $S^p\wedge S^q\cong S^{p+q}$ is
induced by one-point compactification of
$\mathbb R^p\times\mathbb R^q$.
For $W_A^{\wedge2}$ the ordered pairs $(S,T)$ of nonempty subsets
have dimension $|S|+|T|+2$, giving multiplicities
$9,18,15,6,1$ in degrees $4,5,6,7,8$.
For $(S^2\vee S^3)^{\wedge4}$, the words of length four in
$2,3$ give multiplicities $1,4,6,4,1$ in degrees $8,9,10,11,12$.

Write $Y_A=\mathcal A^{\wedge2}$ and
$Y_C=\mathcal C^{\wedge4}$. Their positive homology and
cohomology degrees are respectively $4$--$8$ and $8$--$12$;
the sphere sources have positive homology and cohomology only
in degrees $12$ and $24$. Every positive-degree induced map
therefore vanishes. The ordinary Hurewicz images lie in
$H_{12}(Y_A)=0$ and $H_{24}(Y_C)=0$. A right homotopy inverse
would factor the identity of $H_4(Y_A)=\mathbb Z^9$ through
$H_4(S^{12})=0$, or that of $H_8(Y_C)=\mathbb Z$ through
$H_8(S^{24})=0$. These contradictions exclude homotopy sections.

The scope of these conclusions can be seen through maps on the
same source and target pairs. In the wedge model for $Y_A$ take
two distinct $S^4$ summands and an $S^6$ summand. Their adjoint
loop classes have degrees $3,3,5$ and homology generators
$e_3,f_3,g_5$. The iterated Samelson product has loop Hurewicz image
\[
 [e_3,[f_3,g_5]]=e_3f_3g_5+e_3g_5f_3-f_3g_5e_3-g_5f_3e_3\ne0.
\]
For $Y_C$ use its $S^8$ summand and an $S^{10}$ summand, with
loop generators $e_7,g_9$ of degrees $7,9$. Here
\[
 [e_7,[e_7,g_9]]=e_7e_7g_9-g_9e_7e_7\ne0.
\]
These computations take place in the free integral tensor
algebra on the desuspended sphere generators, identified with
loop homology by the James construction
\cite[Proposition~3C.8, p.~288; Theorem~4J.1, p.~471, and
 p.~473]{Hatcher2002}. The sphere-inclusion generators are
primitive. The loop commutator acts on their cross-product
class by $xy-(-1)^{|x||y|}yx$, giving the displayed expressions;
their tensor words are linearly independent. Adjointness gives
nonzero classes in $\pi_{12}(Y_A)$ and $\pi_{24}(Y_C)$ through
the based equivalences above.

Each target also has a based null-homotopic surjection from its
corresponding sphere. To construct it, consider each of its
finitely many characteristic maps $\phi:D^d\to Y$, where
$d\leq n$ and $(n,Y)$ is $(12,Y_A)$ or $(24,Y_C)$.
Choose a path $\gamma$ from the basepoint of $Y$ to $\phi(0)$.
For $x\in D^n$, write $u_{\rm cell}$ for its first $d$ coordinates and
$r=|x|$. The map
\[
 p(x)=
 \begin{cases}
  2u_{\rm cell},&r\leq\tfrac12,\\
  2(1-r)u_{\rm cell}/r,&\tfrac12\leq r\leq1
 \end{cases}
\]
maps onto $D^d$, is zero on $\partial D^n$, and contracts to
zero through $(1-s)p$. On $|x|\leq1/2$ use $\phi(p(2x))$;
on the outer annulus use $\gamma(2-2r)$. The two formulas agree
at $r=1/2$, and the boundary has the target basepoint value.
They define a based map $S^n\to Y$ covering the closure of the
cell. Contract $p$ on the inner ball, keeping the annulus fixed.
The resulting map is $\gamma(t(r))$, with $t=1$ on that ball
and $t=2-2r$ on the annulus; $\gamma((1-s)t(r))$ contracts it
through based maps. The same construction with constant inner
map includes zero-cells.

Pinch $S^n$ onto a finite wedge with one summand for every cell
and use these maps on the summands. This gives a based
null-homotopic surjection $h_{\rm null}:S^n\to Y$. For either nonzero
class $f$ just constructed, the twofold pinch followed by
$f\vee h_{\rm null}$ is surjective and represents $[f]+[h_{\rm null}]=[f]$.
All these surjections are quotient maps by compactness of the
source and the Hausdorff property of the target. Thus both zero
and nonzero homotopy classes for each of the two source and
target pairs have the stated ordinary invariants and section
obstruction. Those conclusions alone leave the values of the
two specified geometric homotopy classes undetermined.
\end{proof}

\Needspace{9\baselineskip}
\subsection{Decompactification spaces}
\begin{theorem}[Decompactification spaces with Euclidean boundary]
\hypertarget{tag-036N}{}\label{wb:cusp-cone-exact-decompactification-space}
{\normalfont\scriptsize tag 036N}\par
Fix $T_0>0$ so that the radial periods are defined and $D_T>0$
for all $T\geq T_0$, with $r_0=e^{-2\pi T_0}<\epsilon$,
and put $s_0=T_0^{-1}$.
For $\nu\in\{12,24\}$ set
\[
 A_T^{(12)}=\operatorname{diag}(Y_T,Y_T),\qquad
 A_T^{(24)}=\ell_T I_4.
\]
On $[0,s_0]\times\mathbb R^4$ define the fibrewise equivalence
relation
\[
 \begin{aligned}
 (0,x)\sim_\nu(0,y)&\quad\Longleftrightarrow\quad x=y,\\
 (s,x)\sim_\nu(s,y)&\quad\Longleftrightarrow\quad
       x-y\in A_{1/s}^{(\nu)}\mathbb Z^4\qquad(0<s\leq s_0).
 \end{aligned}
\]
Let
\[
 \mathfrak E_\nu=([0,s_0]\times\mathbb R^4)/\!\sim_\nu,
 \qquad q_\nu(s,x)=[s,x],\qquad p_\nu([s,x])=s.
\]
Then $\mathfrak E_\nu$ is a Hausdorff, second-countable topological
five-manifold with boundary, with fibres
\begin{align}
 p_{12}^{-1}(s)&\cong Q_{1/s}^2 &(s>0),&&
 p_{12}^{-1}(0)\cong\mathbb R^4,
 \hypertarget{tag-035K}{}\label{wb:cusp-cone-E12-fibres}\\
 p_{24}^{-1}(s)&\cong C_{1/s}^4 &(s>0),&&
 p_{24}^{-1}(0)\cong\mathbb R^4.
 \hypertarget{tag-035L}{}\label{wb:cusp-cone-E24-fibres}
\end{align}
Its boundary has precisely two components: this zero fibre and
the compact four-torus over $s=s_0$.

Give each positive fibre its Euclidean quotient metric and the zero
fibre its Euclidean metric.  For every $R>0$ there is $s_R>0$ such
that, for $0\leq s<s_R$, the map $q_{\nu,s}(x)=[s,x]$ restricts to
an exact isometry
\[
 q_{\nu,s}:B_R^{\mathbb R^4}(0)
   \xrightarrow{\ \cong_{\mathrm{isom}}\ }
   B_R^{p_\nu^{-1}(s)}([s,0])
\]
onto the entire quotient ball.  The positive fibres converge to the
zero fibre in pointed local-isometric, pointed smooth and pointed
Gromov--Hausdorff topology.

There are continuous surjective comparison maps in the diagrams
\[
 \begin{tikzcd}[column sep=large]
 \mathbb R^4=p_{12}^{-1}(0)
 \arrow[r,hook,"q_{12,0}"]&
 \mathfrak E_{12}\arrow[r,"\beta_{12}"]&
 \mathcal D_{12}\arrow[r,hook,"\delta_{12}"]&
 \mathcal A^{\wedge2},
 \end{tikzcd}
\]
and
\[
 \begin{tikzcd}[column sep=large]
 \mathbb R^4=p_{24}^{-1}(0)
 \arrow[r,hook,"q_{24,0}"]&
 \mathfrak E_{24}\arrow[r,"\beta_{24}"]&
 \mathcal D_{24}\arrow[r,hook,"\delta_{24}"]&
 \mathcal C^{\wedge4}.
 \end{tikzcd}
\]
At positive $s$, $\beta_\nu$ is the marked homeomorphism from the
fibre to the full torus slice at radius $r=e^{-2\pi/s}$.
Its positive-stratum restriction is a homeomorphism
\[
 \beta_\nu:\mathfrak E_\nu\setminus p_\nu^{-1}(0)
       \xrightarrow{\ \cong\ }\mathcal D_\nu\setminus\{d_\nu\},
\]
and
\[
 \beta_\nu^{-1}(d_\nu)=p_\nu^{-1}(0)\cong\mathbb R^4.
\]
Neither comparison map is proper or a quotient map.
The induced continuous bijection
\[
 \overline\beta_\nu:
 \mathfrak E_\nu/p_\nu^{-1}(0)\longrightarrow\mathcal D_\nu
\]
is not a homeomorphism: the cone topology is strictly coarser than
the quotient topology induced by $\beta_\nu$.
Thus the limiting $\mathbb R^4$ is a boundary fibre of
$\mathfrak E_\nu$, and the displayed comparison sends that whole
fibre to the cone vertex.
\end{theorem}

\begin{proof}
The exact radial formulas give continuous invertible matrices
$Y_T$ and positive continuous $\ell_T$ on $T\geq T_0$.
Write
\[
 Y_T=TJ+E_T,\qquad
 J=\begin{pmatrix}0&1\\-1&0\end{pmatrix},\qquad
 E_T=\begin{pmatrix}6m_T&H_T\\B_T-H_T&m_T\end{pmatrix}.
\]
The holomorphic residual functions are bounded on the radial ray,
so the finite constant
\[
 C=\sup_{T\geq T_0}
       \sqrt{37m_T^2+H_T^2+(B_T-H_T)^2}
\]
bounds $\lVert E_T\rVert$.  Orthogonality of $J$ gives
$\lVert Y_Tv\rVert\geq(T-C)\lVert v\rVert$.
The same constant also bounds the circle period below.
With the subscripts $T$ suppressed, direct completion of the square gives
\[
 \left\lVert Y\binom{1}{z}\right\rVert^2
 =(L^2+m^2)\left(z+\frac{m(6L-q)}{L^2+m^2}\right)^2
       +\frac{D^2}{L^2+m^2}.
\]
By \eqref{wb:cusp-spec-global-circle-circumference},
$\ell_T=D_T/\sqrt{L_T^2+m_T^2}$; since $D>0$, the last minimum
is $\ell_T^2$.
For $T>C$ this implies $\ell_T\geq T-C$ by the norm bound and
$\lVert(1,z)\rVert\geq1$; for $T\leq C$ the inequality follows
from positivity of $\ell_T$.
Consequently the shortest nonzero period $\lambda_\nu(T)$ of either
$A_T^{(\nu)}\mathbb Z^4$ satisfies
\[
 \lambda_\nu(T)\geq T-C\longrightarrow\infty.
\]

Put $M=[0,s_0]\times\mathbb R^4$, $M_+=(0,s_0]\times\mathbb R^4$,
and abbreviate $A(s)=A_{1/s}^{(\nu)}$ for $s>0$.
The saturation of an open subset $U\subset M$ is
\[
 U\ \cup\!\bigcup_{n\in\mathbb Z^4}
   \{(s,x+A(s)n):(s,x)\in U\cap M_+\}.
\]
Every positive-stratum translation is a homeomorphism with inverse
translation by $-A(s)n$, and $M_+$ is open in $M$.
Thus the saturation is open and $q_\nu$ is an open quotient map.

The equivalence relation is closed.  Indeed, in a convergent sequence
of related pairs whose parameters tend to $s>0$, the integer vectors
$A(s_j)^{-1}(x_j-y_j)$ converge and hence are eventually constant.
The limiting pair is therefore related.  When $s_j\to0$, the
differences $x_j-y_j$ are bounded, whereas every nonzero period has
length at least $\lambda_\nu(1/s_j)\to\infty$ at positive indices.
The differences are eventually zero, including the zero-level
indices by definition, so $x=y$.  This proves closedness because
$M\times M$ is metrizable.
For unrelated representatives, choose a product neighbourhood disjoint
from the relation.  The images of its two factors under the open
map $q_\nu$ are disjoint open neighbourhoods.  Hence the quotient
is Hausdorff.  Images of a countable base of $M$ under $q_\nu$ form
a countable base, proving second countability.

The projection $p_\nu$ is continuous by the quotient property.
For an open $V\subset\mathbb R^4$ and a relative open interval
$I\subset[0,s_0]$ containing $s$,
\[
 q_{\nu,s}(V)=p_\nu^{-1}(s)\cap q_\nu(I\times V).
\]
Thus $q_{\nu,s}$ is open onto its fibre.
At zero it is bijective, giving the stated homeomorphism with
$\mathbb R^4$.  At positive $s$ it induces a homeomorphism from
$\mathbb R^4/A(s)\mathbb Z^4$ onto the fibre.  Splitting the four
coordinates into two pairs for $\nu=12$, or taking their individual
residues modulo $\ell_{1/s}$ for $\nu=24$, gives
\eqref{wb:cusp-cone-E12-fibres} and \eqref{wb:cusp-cone-E24-fibres}.

For a zero-fibre point represented by $x_0$, choose $a_{\rm ch}>0$ and
$0<\eta<\min\{s_0,(2a_{\rm ch}+C+1)^{-1}\}$.
The map $q_\nu$ is injective on $[0,\eta)\times B_{a_{\rm ch}}^\circ(x_0)$,
because a difference of two vectors in that ball has length less
than $2a_{\rm ch}$, while every nonzero period has greater length.
Its restriction is an open continuous bijection, giving a half-space
chart of dimension five.  Near any $s_*>0$, continuity of $A(s)^{-1}$
gives a uniform positive lower bound on $\lVert A(s)v\rVert/\lVert v\rVert$.
An interval around $s_*$ times a sufficiently small ball therefore
gives the same injective open quotient chart.  For $0<s_*<s_0$
this is an interior chart; at $s_*=s_0$ the coordinate $s_0-s$
gives a half-space chart.  These charts cover the entire space.
Their local homology identifies the boundary precisely as the two
end fibres.  Both end fibres are connected and are separated inside
the boundary by their $p_\nu$ values, proving the boundary assertion.

For the metric conclusion take
$s_R=\min\{s_0,(4R+C+1)^{-1}\}$.
For $0<s<s_R$ and $x,y\in B_R(0)$, every nonzero period $\lambda$
satisfies
\[
 \lVert x-y+\lambda\rVert
 \geq\lVert\lambda\rVert-2R
 \geq1/s-C-2R>2R\geq\lVert x-y\rVert.
\]
The quotient therefore preserves all distances on this ball.
Every lattice coset has a closest representative, since the full
rank lattice has finitely many points in bounded sets and its
vectors tend to infinity with their integer coefficients.
A point in the quotient radius-$R$ ball consequently has a
representative of norm at most $R$, proving surjectivity onto the
whole ball.  The same argument applies to open balls with strict
inequalities.  At $s=0$ the given Euclidean identification is the
identity.  On open balls the quotient maps are smooth local
isometries and pull the metric tensor back exactly to the Euclidean
tensor.  For every fixed radius these identifications give the
local-isometric and smooth limits; exact isometry of the closed
balls also gives the pointed Gromov--Hausdorff limit.

Let $c_{12}:(\overline T^2)^2\to\mathbb R^4/\mathbb Z^4$ concatenate
the two coordinate pairs, and let $c_{24}$ take the four individual
circle coordinates, as in Proposition~\ref{wb:cusp-cone-explicit-carriers}.
The comparison formulas are
\[
 \beta_\nu([s,x])=
 \begin{cases}
 [e^{-2\pi/s},c_\nu^{-1}([A(s)^{-1}x])],&s>0,\\
 d_\nu,&s=0.
 \end{cases}
\]
Adding $A(s)n$ to $x$ adds $n$ to the marked coordinate, so the
formula is well defined.  Its positive-stratum cylinder map is
continuous.  Every vertex neighbourhood in the cone contains a
full collar of sufficiently small radius by compactness of its
link.  The limit $e^{-2\pi/s}\to0$ therefore proves continuity
at the zero fibre, uniformly in $x$, and the formula descends
continuously through $q_\nu$.

The map $[s,x]\mapsto(s,[A(s)^{-1}x])$ identifies the positive
stratum with $(0,s_0]\times(\mathbb R^4/\mathbb Z^4)$.
Its inverse is $(s,[a])\mapsto[s,A(s)a]$; both maps descend from
continuous cylinder maps through the open torus quotient.
The radius map $s\mapsto e^{-2\pi/s}$ has inverse
$r\mapsto-2\pi/\log r$ on $(0,r_0]$, where $r_0<1$.
Together with $c_\nu^{-1}$ these formulas give the full continuous
inverse of $\beta_\nu$ on the positive stratum.
In the first family they apply $\iota_T^{-1}$ to each of the two
$Q_T$ factors; in the second they apply $\vartheta_T^{-1}[b]
=[b/\ell_T]$ to each $C_T$ factor.  Thus these are precisely the
marked slice maps of that proposition, and composition with the
displayed $\delta_\nu$ gives the full diagrams.
The zero embedding is $q_{\nu,0}$; it is closed because it is the
fibre of the continuous $p_\nu$ over zero.  The positive-stratum
homeomorphism proves surjectivity and shows that the vertex inverse
image is exactly that zero fibre.  Its homeomorphism with the
noncompact $\mathbb R^4$ makes the inverse image of the compact
singleton $\{d_\nu\}$ noncompact, proving that $\beta_\nu$ is not proper.

For any $a\in\mathbb R^4\setminus\mathbb Z^4$, the attained distance
$d=\operatorname{dist}(a,\mathbb Z^4)$ is positive.  Applying the
$Y_T$ bound to both blocks gives
\[
 \inf_{k\in\mathbb Z^4}\|A_T^{(12)}(a+k)\|\geq(T-C)d,
 \qquad
 \inf_{k\in\mathbb Z^4}\|A_T^{(24)}(a+k)\|=\ell_Td.
\]
For $T\leq C$ the first inequality follows from nonnegativity.
Both bounds tend to infinity, and subtracting any fixed $x_0$
decreases them by at most $\|x_0\|$.  Thus every such nonintegral
marked coset escapes all bounded zero-fibre charts as $s=1/T\to0$.

For a concrete quotient-topology witness, take $v=\tfrac12 e_1$ and,
for integers $n\geq1$, put
\[
 T_n=T_0+n,\qquad s_n=T_n^{-1},\qquad
 z_n=q_\nu(s_n,A_{T_n}^{(\nu)}v),\qquad
 S_\nu=\{z_n:n\geq1\}.
\]
Every representative of $z_n$ is $A_{T_n}^{(\nu)}(v+k)$ for
$k\in\mathbb Z^4$.  Twice this vector is
$A_{T_n}^{(\nu)}(e_1+2k)$, a nonzero period because its first
integer coefficient is odd.  Thus every representative has norm
at least $(T_n-C)/2\to\infty$.
The full inverse image $q_\nu^{-1}(S_\nu)$ is closed: at a positive
base level a sufficiently small interval meets only finitely many
of these closed lattice cosets; at zero every bounded spatial
neighbourhood misses the tail, and a smaller base interval excludes
the remaining finitely many levels.  The quotient definition then
makes $S_\nu$ closed in $\mathfrak E_\nu$.

Its image is
\[
 B_\nu=\{[e^{-2\pi T_n},c_\nu^{-1}([v])]:n\geq1\}.
\]
This fixed-link sequence converges to $d_\nu$, which it omits.
Consequently $U_\nu=\mathcal D_\nu\setminus B_\nu$ contains the
vertex but is not open.  On the other hand,
$\beta_\nu^{-1}(U_\nu)=\mathfrak E_\nu\setminus S_\nu$ is open,
by the exact positive-stratum bijection and the vertex fibre.
This proves that $\beta_\nu$ is not a quotient map.

Finally let $c:\mathfrak E_\nu\to
\mathfrak E_\nu/p_\nu^{-1}(0)$ be the collapse quotient.
The quotient universal property gives the continuous
$\overline\beta_\nu$ with $\overline\beta_\nu c=\beta_\nu$;
the computed fibres make it bijective.
The set $\overline\beta_\nu^{-1}(U_\nu)$ is open because its
inverse image under $c$ is the open set
$\mathfrak E_\nu\setminus S_\nu$, while its image $U_\nu$ is not
open.  Hence the inverse of this bijection is not continuous.
Continuity of $\overline\beta_\nu$ gives the inclusion of cone-open
sets among quotient-open sets, and this $U_\nu$ makes the inclusion
strict.  This proves the stated comparison of the two topologies.
\end{proof}
\subsection{A framed nullbordism of the cusp fibre}

\begin{definition}[Marked cusp data]\hypertarget{tag-Q258}{}\label{wb:def:marked-cusp-nullbordism-data}
{\normalfont\scriptsize tag Q258}\par
Use the original smooth sphere member and its intrinsic closed cusp
neighbourhood $N_0$ from Definition~\ref{mapsIII:def:cusp-suspension-maps}.
Write $f:X\to B_{\mathrm{base}}$ for its map to the base and $p_0$
for the cusp. Let $t$ be its specified local coordinate, with
$t(p_0)=0$, put $W=f^{-1}(p_0)$, and choose $x_0\in W$.
The ordered quotient basis is induced by
$(\widehat\gamma,\widehat u)$ and the kernel lattice is
$\mathbb Z\langle\widehat w,\widehat\delta\rangle$.
The quotient torus is identified with
$\overline T^2=\mathbb R^2/\mathbb Z^2$ by the inverse marked maps
\[
 \beta:[a\widehat\gamma+b\widehat u]\longmapsto[a,b],\qquad
 \beta^{-1}:[a,b]\longmapsto[a\widehat\gamma+b\widehat u].
\]
Write $h_{\mathrm{per}}$ for the same period function denoted $h$
in that definition. The actual periods are
\begin{equation}
 Z(s)=sB_0+C(t),\qquad e^{2\pi i s}=t,\qquad
 B_0=\begin{pmatrix}0&1\\-1&0\end{pmatrix},\qquad
 C(t)=\begin{pmatrix}6\mu(t)&h_{\mathrm{per}}(t)\\
 b(t)-h_{\mathrm{per}}(t)&\mu(t)\end{pmatrix}.
 \hypertarget{tag-038N}{}\label{wb:cusp-null:eq:periods}
\end{equation}
Every entry of $C$ is holomorphic at $t=0$. On the dense torus
$(x_1,x_2,t)$, with $x=\exp(2\pi i\zeta)$, the residual period
action is
\begin{equation}
 \Psi_\lambda(x,t)=
 \bigl(t^{B_0\lambda}\exp(2\pi iC(t)\lambda)x,t\bigr),
 \qquad \lambda\in\mathbb Z^2.
 \hypertarget{tag-038E}{}\label{wb:cusp-null:eq:action}
\end{equation}
Its inverse is $\Psi_{-\lambda}$, as direct multiplication of the
two componentwise factors shows. It is the free properly
discontinuous action on the unimodular filling of
\SeriesRef{fam:src:A.3.1}; its quotient projection is a local
biholomorphism at every toric fixed point. The explicit injective
chart needed below is proved in \cref{wb:h24:ret:chart-lemma}.

Put
\[
 K^{(1)}=((\mathbb R/\mathbb Z)\times\{0\})
          \cup(\{0\}\times(\mathbb R/\mathbb Z)),\qquad
 Z_2=\overline T^2/K^{(1)},\qquad
 c:\overline T^2\to Z_2.
\]
The open two-cell gives $Z_2$ its orientation, so $c$ has degree
one. Give it the ordered smooth open-cell coordinates $(c_1,c_2)$.
The reduced-suspension target and the native map are
\[
 B=\Sigma_{\mathrm{red},\epsilon}Z_2
 =\frac{[0,\epsilon]\times Z_2}
 {(\{0,\epsilon\}\times Z_2)\cup([0,\epsilon]\times\{*\})},
\]
\begin{equation}
 g(x)=[|t|,c(\overline\pi_f(x))]\quad(0<|t|<\epsilon),\qquad
 g(W)=g(X\setminus\operatorname{int}N_0)=*.
 \hypertarget{tag-038J}{}\label{wb:cusp-null:eq:g}
\end{equation}
The cube target map and its inverse are
\[
 \phi[r,c(u_1,u_2)]=[r/\epsilon,u_1,u_2],\qquad
 \phi^{-1}[s,u_1,u_2]=[\epsilon s,c([u_1,u_2])]
\]
on the open cells; both send the basepoint to the basepoint.
The same collapsed faces give both inverse identities and continuous
quotient descent. Thus $\phi:B\to C_3$ is a based homeomorphism.
The round target is $h_3\phi:B\to S^3$, with inverse
$\phi^{-1}h_3^{-1}$ from Definition~\ref{wb:def:ordered-cube-spheres}.
For the fixed based diffeomorphism $d_S:S^6\to X$, put
$d=d_Sh_6:C_6\to X$, with inverse $h_6^{-1}d_S^{-1}$.
In homotopy groups, $[g]$ denotes the class of $g_C=\phi gd$;
the inverse coordinate transport is $g=\phi^{-1}g_Cd^{-1}$.
\end{definition}

\begin{theorem}[A framed nullbordism of the cusp fibre]
\hypertarget{tag-038S}{}\label{wb:cusp-null:thm}
{\normalfont\scriptsize tag 038S}\par
The based map $g:(X,x_0)\to(B,*)$ of
\eqref{wb:cusp-null:eq:g}, with the sphere identification
$h_3\phi:B\to S^3$ of
Definition~\ref{wb:def:marked-cusp-nullbordism-data}, is based
null-homotopic.
There is an interior regular value whose inverse image, with the
ordered normal framing induced by $(r,c_1,c_2)$, is framed nullbordant
in one cusp coordinate chart times $[0,1]$, disjoint from
$\{x_0\}\times[0,1]$.
\end{theorem}

\begin{proof}
The inverse marking maps $\beta,\beta^{-1}$ in
Definition~\ref{wb:def:marked-cusp-nullbordism-data} carry the two
marked coordinate circles onto the two coordinate circles.
They induce inverse maps
$\bar\beta$ and $\bar\beta^{-1}$ on the skeleton collapses, and
inverse suspension maps $[r,z]\mapsto[r,\bar\beta(z)]$ and
$[r,z]\mapsto[r,\bar\beta^{-1}(z)]$.
Under the first suspension map the angle-forgetting map of
\cref{mapsIII:angle-forgotten-one-copy-null} is exactly $g$ on the
punctured cusp and at both collapsed radial ends.
Write $S_\beta$ for the first suspension homeomorphism and
$H_0$ for that proved based null-homotopy. Then
$H_g(x,s)=S_\beta H_0(x,s)$ has initial map $g$, constant terminal
map and $H_g(x_0,s)=*$; its inverse transport is $S_\beta^{-1}H_g=H_0$.
This proves the first assertion.

\emph{1. The complete inverse image.}
Set
\[
 u^0=\binom{-1/4}{1/4},\qquad
 \overline u^0=(3/4,1/4),\qquad
 p^0=B_0u^0=\binom{1/4}{1/4}.
\]
Use the injective neighbourhood constructed in the proof of
\cref{wb:h24:ret:chart-lemma}. Its toric coordinates are
\begin{equation}
 z_0=t/(x_1x_2),\qquad z_1=x_1,\qquad z_2=x_2,
 \qquad t=z_0z_1z_2.
 \hypertarget{tag-038H}{}\label{wb:cusp-null:eq:chart}
\end{equation}
Write $\ell=\log|t|$, $R(t)=-2\pi\operatorname{Im}C(t)$ and
$A(t)=\ell B_0+R(t)$. The marking $\zeta=Z(s)u+k$, with
$u,k\in\mathbb R^2$ and $\operatorname{Im}s=-\ell/(2\pi)$, gives
\begin{equation}
 \binom{\log|x_1|}{\log|x_2|}=A(t)u.
 \hypertarget{tag-038L}{}\label{wb:cusp-null:eq:moduli}
\end{equation}
Since $B_0^TB_0=I_2$, the reverse triangle inequality gives
\begin{equation}
 \|A(t)v\|\geq\bigl(|\ell|-\|R(t)\|_{\mathrm{op}}\bigr)\|v\|.
 \hypertarget{tag-038K}{}\label{wb:cusp-null:eq:invertibility}
\end{equation}
The chart uses the full holomorphic units
\begin{align}
 h_j^0(t)&=\exp\bigl(-2\pi i(C(t)u^0)_j\bigr),\qquad j=1,2,\\
 \mathcal F(z_0,z_1,z_2)&=
 \left(\frac{z_0}{h_1^0(t)h_2^0(t)},
 h_1^0(t)z_1,h_2^0(t)z_2\right)=(w_0,w_1,w_2).
 \hypertarget{tag-038Q}{}\label{wb:cusp-null:eq:straightening}
\end{align}
Its inverse is
\[
 (w_0,w_1,w_2)\longmapsto
 \left(w_0h_1^0(t)h_2^0(t),\frac{w_1}{h_1^0(t)},
 \frac{w_2}{h_2^0(t)}\right),\qquad t=w_0w_1w_2.
\]
Both formulas hold on the central coordinate planes as well.
The cited chart lemma gives inverse coordinate maps
$\iota_U:W_\delta\to U\subset X$ and $\iota_U^{-1}$, where
$W_\delta=\{|w_j|<\delta\}$, $0<\delta<1$; all nonidentity
period translates of $\mathcal F^{-1}(W_\delta)$ are disjoint.

On the closed period disc $|t|\leq a_0$ of that lemma, put
$M_R=\max\|R(t)\|_{\mathrm{op}}$ and choose
\[
 0<r_0<\min\{\delta^4,e^{-(2M_R+1)}\},\qquad
 \ell_0=\log r_0.
\]
The relation $\delta^3<a_0<\epsilon$ ensures that this choice is
inside the original cusp disc. We will also impose a basepoint
bound below. Set $y_0=[r_0,c(\overline u^0)]$.
Because $y_0$ is in the open target cell, its inverse image has
$|t|=r_0$ and $[u]=[u^0]$.
Applying the period translate $\Psi_{-\lambda}$ to a lift
$u=u^0+\lambda$ gives $u=u^0$. A residual translate changes
the logarithmic moduli by $A(t)\lambda$; invertibility from
\eqref{wb:cusp-null:eq:invertibility} makes this lift unique.
The real kernel coordinates $k\bmod\mathbb Z^2$ supply all
independent phases of $x_1,x_2$.

Since $\log|h^0|=-R(t)u^0$, the complete inverse image is
\begin{equation}
 T_{r_0}=\{|w_0|=\rho_0,\ |w_1|=\rho_1,\ |w_2|=\rho_2\},
 \qquad (\rho_0,\rho_1,\rho_2)
 =(r_0^{1/2},r_0^{1/4},r_0^{1/4}).
 \hypertarget{tag-038R}{}\label{wb:cusp-null:eq:torus}
\end{equation}
Conversely these moduli give $A(t)(u-u^0)=0$ and hence the same
quotient value. Its angular coordinates and their inverse are
\begin{equation}
 (\varphi,a,b)\longmapsto(\varphi-a-b,a,b),\qquad
 (\theta_0,a,b)\longmapsto(\theta_0+a+b,a,b)
 \quad\pmod{2\pi}.
 \hypertarget{tag-033Q}{}\label{qa:III:layout:1}
\end{equation}
Thus the full inverse image is the embedded three-torus
$(\rho_0e^{i(\varphi-a-b)},\rho_1e^{ia},\rho_2e^{ib})$.
Both it and the compact spatial set
\begin{equation}
 \{|w_0|\leq\rho_0,\ |w_1|=\rho_1,\ |w_2|=\rho_2\}
 \subset W_\delta
 \hypertarget{tag-038P}{}\label{wb:cusp-null:eq:spatial-bound}
\end{equation}
lie in the injective chart: $\rho_1=\rho_2<\delta$ and
$\rho_0<\delta^2<\delta$.

\emph{2. The ordered differential and its frame homotopy.}
Put $L_j=\log|w_j|$ and use the real lift $u$ near $u^0$.
The coordinate formulas give
\begin{equation}
 \binom{L_1}{L_2}=\ell B_0u^0+A(t)(u-u^0),\qquad
 L_0=\ell-L_1-L_2.
 \hypertarget{tag-038M}{}\label{wb:cusp-null:eq:nearby-moduli}
\end{equation}
On the fibre $u=u^0$, the derivative term $(dA)(u-u^0)$ vanishes.
Consequently
\begin{equation}
 \begin{pmatrix}dL_0\\dL_1\\dL_2\end{pmatrix}
 =J(t)\begin{pmatrix}d\ell\\du_1\\du_2\end{pmatrix},\qquad
 J(t)=\begin{pmatrix}
 1/2&-A_{11}-A_{21}&-A_{12}-A_{22}\\
 1/4&A_{11}&A_{12}\\
 1/4&A_{21}&A_{22}
 \end{pmatrix}.
 \hypertarget{tag-038B}{}\label{wb:cusp-null:eq:J}
\end{equation}
Adding the last two rows to the first gives $(1,0,0)$, so
$\det J(t)=\det A(t)\ne0$.
Let $D=Dc_{\overline u^0}$ in the ordered open target cell,
and set $H=\operatorname{diag}(r_0,D)$.
For the open-cell coordinates $(u_1+1,u_2)$, $D=I_2$; the matrix
$H$ records the derivative of $(\ell,u)$ to $(r,c_1,c_2)$.
Let $n_j$ be the radial vector $w_j$ in complex coordinate $j$,
and put $N=(n_0,n_1,n_2)$. Since $dL_i(n_j)=\delta_{ij}$,
$dg(N)=HJ(t)^{-1}$, and the induced ordered normal framing is
\begin{equation}
 \mathfrak n(t)=NJ(t)H^{-1}.
 \hypertarget{tag-038F}{}\label{wb:cusp-null:eq:actual-frame}
\end{equation}
These vectors are normal in the Euclidean chart metric. For another
metric, orthogonal projection identifies their quotient-normal classes
with the corresponding metric normal vectors; the inverse sends a
normal vector to its class modulo the tangent bundle.

For its Pontryagin--Thom interpretation, apply
Lemma~\ref{mapsIII:based-regular-fibre-smoothing} to the actual map
$h_3\phi g:X\to S^3$, source basepoint $x_0$ and value
$h_3\phi(y_0)$. The full inverse image is $T_{r_0}$, disjoint from
$x_0$ because $|t|=r_0>0$ there. The displayed differential proves
smoothness and rank three near it. The ordered target chart is
$(r,c_1,c_2)\phi^{-1}h_3^{-1}$, with inverse
$h_3\phi(r,c_1,c_2)^{-1}$, so its composite with $h_3\phi g$
is exactly $(r,c_1,c_2)g$. The lemma preserves this entire inverse
image and its actual frame. Postcomposition of its homotopy by
$\phi^{-1}h_3^{-1}$ returns to $B$ and fixes its basepoint.

For $0\leq\sigma\leq1$, let
$A_\sigma(t)=\ell_0B_0+\sigma R(t)$ on $|t|=r_0$ and form
$J_\sigma$ from the same matrix formula.
For $v\ne0$, the bound
$\|A_\sigma v\|\geq(|\ell_0|-M_R)\|v\|>0$
makes every $J_\sigma$ invertible. At $\sigma=0$,
\begin{equation}
 J_0=\begin{pmatrix}
 1/2&\ell_0&-\ell_0\\
 1/4&0&\ell_0\\
 1/4&-\ell_0&0
 \end{pmatrix},\qquad \det J_0=\ell_0^2>0.
 \hypertarget{tag-038C}{}\label{wb:cusp-null:eq:J0}
\end{equation}
Run $NJ_\sigma H^{-1}$ from $\sigma=1$ to $\sigma=0$, using
a smooth time parameter constant near its endpoints.
This gives a framed product cylinder from $\mathfrak n$ to
$NP$, where
\begin{equation}
 P=J_0H^{-1}.
 \hypertarget{tag-038D}{}\label{wb:cusp-null:eq:P}
\end{equation}

\emph{3. The collared framed cap.}
Use a new time coordinate $v\in[0,1]$ and define
\begin{equation}
 q(v)=\begin{cases}
 0,&0\leq v\leq1/4,\\
 \rho_0^2\exp\bigl(4-1/(v-1/4)\bigr),&1/4<v\leq1.
 \end{cases}
 \hypertarget{tag-038O}{}\label{wb:cusp-null:eq:q}
\end{equation}
All right derivatives at $1/4$ vanish. The function increases for
$v>1/4$, with $q(1/2)=\rho_0^2$ and
$q'(1/2)=16\rho_0^2$.
In $W_\delta\times[0,1]$, set
\begin{equation}
 \mathcal W=\{|w_0|^2+q(v)=\rho_0^2,
 |w_1|^2=\rho_1^2,
 |w_2|^2=\rho_2^2\}.
 \hypertarget{tag-038G}{}\label{wb:cusp-null:eq:bordism}
\end{equation}
The three defining differentials are independent. The first has a
nonzero $w_0$ derivative off $w_0=0$, and at $w_0=0$ its time
derivative is $q'(1/2)>0$; the last two have nonzero derivatives in
their separate complex factors. This is a compact smooth
four-manifold with only boundary $T_{r_0}\times\{0\}$ and exact
collar $T_{r_0}\times[0,1/4]$; there are no points with $v>1/2$.
For an explicit parametrization of its disk factor, take
$z\in\overline{\mathbb D}$ and put
\[
 v=\frac{1-|z|^2}{2},\qquad w_0=f_D(|z|^2)z,
 \qquad
 f_D(s)=\sqrt{\frac{\rho_0^2-q((1-s)/2)}{s}},\quad
 f_D(0)=\sqrt{q'(1/2)/2}.
\]
The numerator has a simple zero with positive derivative at $s=0$,
so $f_D$ extends smoothly and positively. At $z=0$ the
$w_0$ derivative is $f_D(0)I_2$; elsewhere the radial time derivative
and angular $w_0$ derivative are nonzero.
The inverse is
\[
 z=\sqrt{1-2v}\,\frac{w_0}{|w_0|}\quad(w_0\ne0),\qquad
 z=0\quad(w_0=0).
\]
It is smooth at the centre by the inverse-function theorem applied
to the derivative $f_D(0)I_2$; away from the centre the displayed
formula is smooth. Thus this is an embedding onto the disk factor, giving
$\mathcal W\cong\overline{\mathbb D}\times T^2$.
Its spatial coordinates satisfy \eqref{wb:cusp-null:eq:spatial-bound}.

In the chart product metric, let
\[
 V_0=(w_0,0,0,q'(v)/2),\qquad
 V_1=(0,w_1,0,0),\qquad V_2=(0,0,w_2,0).
\]
They are half the three defining gradients, mutually orthogonal and
nonzero. In particular the first defining differential takes the
positive value $2|w_0|^2+q'(v)^2/2$ on $V_0$.
On the entire bottom collar they equal $(N,0)$, so
\begin{equation}
 (V_0,V_1,V_2)P
 \hypertarget{tag-038I}{}\label{wb:cusp-null:eq:extended-frame}
\end{equation}
extends exactly the constant-matrix framing $NP$.
Attach the preceding framed cylinder to this cap and rescale their
time intervals. The product collars and their frames agree exactly,
so this gives a smooth framed nullbordism of the inverse image.

Finally $t(x_0)=0$. Every point of the cylinder has $|t|=r_0$;
the cap meets $t=0$ only at $w_0=0$, $v=1/2$, where
$|w_1|=|w_2|=r_0^{1/4}$.
For $x_0\notin U$ disjointness follows from containment in $U$.
For $x_0\in U$, write $w^*=\iota_U^{-1}(x_0)$.
A zero $w_1^*$ or $w_2^*$ excludes equality with the cap.
When both are nonzero, also require
\[
 r_0<\frac{|w_1^*|^4}{16},\qquad
 r_0<\frac{|w_2^*|^4}{16}.
\]
These positive bounds are compatible with the earlier choices and
make the cap moduli smaller than half the corresponding basepoint
moduli. The whole framed bordism lies in $U\times[0,1]$, avoids
$\{x_0\}\times[0,1]$, and has empty terminal boundary.
\end{proof}

\Needspace{6\baselineskip}
\begin{corollary}[Coefficients and characters of the cusp map]
\hypertarget{tag-03A6}{}\label{wb:cusp_s12_carrier_homotopy:L567}
{\normalfont\scriptsize tag 03A6}\par
Under the fixed based coordinates of
Definition~\ref{wb:def:marked-cusp-nullbordism-data},
\[
 [g]=0\in\pi_6(S^3)\cong\mathbb Z/12,\qquad
 [E^6g_C]=0\in\pi_{12}(S^9)\cong\mathbb Z/24.
\]
For a chosen generator $\nu_9$ of the target group, define its
even-coordinate map and inverse by
\[
 \delta:\mathbb Z/12\longrightarrow2\langle\nu_9\rangle,
 \quad[a]_{12}\longmapsto2a\nu_9,\qquad
 \delta^{-1}(2a\nu_9)=[a]_{12}.
\]
The native coefficient is
$a=\delta^{-1}([E^6g_C])=0\pmod{12}$.
The homomorphism
\[
 \Theta_g:\mathbb Z/12\longrightarrow\pi_6(S^3),
 \qquad \Theta_g([n])=n[g],
\]
is zero.  For every additive character
$\chi:\pi_6(S^3)\to\mathbb R/\mathbb Z$,
$\chi\circ\Theta_g=0$.
The associated multiplicative character
$e^{2\pi i\chi}\circ\Theta_g$ is identically $1$.
\end{corollary}

\begin{proof}
Use $d:C_6\to(X,x_0)$ and $\phi:B\to C_3$ from
Definition~\ref{wb:def:marked-cusp-nullbordism-data}, and let
$H_g$ be the based null-homotopy of \cref{wb:cusp-null:thm}.
The formula $H(u,s)=\phi H_g(d(u),s)$ has initial map $g_C$,
constant terminal map and fixed basepoint, so $[g]=0$.
The inverse coordinate changes are $d^{-1}$ and $\phi^{-1}$.
The right and left sixfold suspension homotopies have formulas
$([u,v],s)\mapsto[H(u,s),v]$ and
$([v,u],s)\mapsto[v,H(u,s)]$, respectively.
They descend because each fat wedge maps to the basepoint at every
time. Transport through the fixed ordered smash identifications
therefore gives $[E^6g_C]=0$ in either placement convention.
For the displayed $\delta$, adding $12$ to $a$ changes the image
by $24\nu_9=0$, and
$2a\nu_9=2b\nu_9$ is equivalent to $24\mid2(a-b)$, hence to
$12\mid(a-b)$. Thus $\delta$ is an isomorphism onto the even
subgroup, with the stated well-defined inverse. The zero class is
in this image, so its actual coefficient exists and is
$\delta^{-1}(0)=0$.
Since $[g]=0$, the formula defining $\Theta_g$ is well-defined and
has value zero on every class.  Each additive character sends this
value to $0\in\mathbb R/\mathbb Z$, whose exponential is $1$.
\end{proof}

\subsection{Coefficients and framed fibres of the sphere carriers}
\hypertarget{tag-038U}{}\label{wb:cusp-s12-carrier-homotopy}

\begin{definition}[The native sphere carriers]\hypertarget{tag-Q256}{}\label{wb:def:native-sphere-carriers}
{\normalfont\scriptsize tag Q256}\par
Use the original sphere member, cusp neighbourhood, ordered period
marking, skeleton collapse $c:\overline T^2\to Z_2$ and map $g$
of Definition~\ref{wb:def:marked-cusp-nullbordism-data}. Put
\[
 K=S^1\times\overline T^2,\qquad
 \mathcal A=([0,\epsilon]\times K)/
 (\{0\}\times K=v_0,\ \{\epsilon\}\times K=v_1),\qquad
 B=\Sigma_{\mathrm{red},\epsilon}Z_2.
\]
The two poles of $\mathcal A$ are distinct; its basepoint is $v_0$.
The map $q_2$ of \cref{wb:cusp-cone-S12-route} is
\[
 q_2(x)=[r,e^{i\varphi},\overline\pi_f(x)]
 \quad(t(f(x))=re^{i\varphi},\ 0<r<\epsilon),\qquad
 q_2(W)=v_0,\quad
 q_2(X\setminus\operatorname{int}N_0)=v_1.
\]
The continuous based postcomposition
\[
 p=\widehat p_2:\mathcal A\to B,\qquad
 [r,z,a]\longmapsto[r,c(a)]
\]
sends both poles to the reduced-suspension basepoint, and $g=pq_2$
by its defining formula on the interior and both collapsed ends.
This is the continuous surjective based quotient $\widehat q_3$
of \cref{wb:cusp-cone-angle-forgotten-sphere-carriers}.

Use the same $d,\phi$ and $g_C=\phi gd$ of
Definition~\ref{wb:def:marked-cusp-nullbordism-data}, with inverse
transport $g=\phi^{-1}g_Cd^{-1}$. In homotopy groups $[g]$ denotes
$[g_C]$.
All smash coordinates use Definition~\ref{wb:def:ordered-cube-spheres}.
In particular, define
\[
 Q=(q_2\wedge q_2)(d\wedge d)\mu_{6,6}^{-1}:C_{12}\to\mathcal A\wedge\mathcal A,
 \qquad
 P=\mu_{3,3}(\phi\wedge\phi)(p\wedge p):\mathcal A\wedge\mathcal A\to C_6.
\]
\end{definition}

\Needspace{6\baselineskip}
\begin{corollary}[The two-copy sphere carrier]\hypertarget{tag-Q254}{}\label{wb:cusp-s12-native-smash-null}
{\normalfont\scriptsize tag Q254}\par
The specified continuous surjective based quotient
$\widehat q_3^{\wedge2}:S^{12}\twoheadrightarrow S^6$ of
\cref{wb:cusp-cone-sphere-to-sphere-maps} is based null-homotopic.
\end{corollary}

\begin{proof}
Use $d,\phi,g_C$ of Definition~\ref{wb:def:native-sphere-carriers}.
The displayed two-copy quotient is exactly
\[
 \mu_{3,3}(g_C\wedge g_C)\mu_{6,6}^{-1}:C_{12}\longrightarrow C_6.
\]
The inverse coordinate maps are $d^{-1}$, $\phi^{-1}$ and the block
splittings in Definition~\ref{wb:def:ordered-cube-spheres}.
Apply \cref{wb:cusp-s12-full-carrier-null} to $x=y=g_C$.
The original quotient's continuity, surjectivity and quotient property
are those proved in \cref{wb:cusp-cone-angle-forgotten-sphere-carriers};
the displayed based homeomorphisms transport its based null-homotopy.
\end{proof}

\Needspace{6\baselineskip}
\begin{proposition}[The coefficient and framed fibre]\hypertarget{tag-038X}{}\label{wb:cusp-s12-one-copy-invariant}
{\normalfont\scriptsize tag 038X}\par
The sixfold suspension identifies $\pi_6(S^3)$ with the even subgroup
of $\pi_{12}(S^9)\cong\mathbb Z/24$. For every chosen generator
$\nu_9$, the map
\[
 \operatorname{coef}_{\nu_9}:\pi_6(S^3)\longrightarrow\mathbb Z/12,
 \qquad E^6x=2\operatorname{coef}_{\nu_9}(x)\nu_9,
\]
is an isomorphism. Its inverse is $[a]_{12}\mapsto a\gamma$, where
$E^6\gamma=2\nu_9$. The native $g$ of
Definition~\ref{wb:def:native-sphere-carriers} has coefficient zero.

For every $0<r_0<\epsilon$ and every
$u_0\in\overline T^2\setminus K^{(1)}$, the value
$y_0=[r_0,c(u_0)]$ is regular. Its entire inverse image is an
embedded three-torus with the ordered normal framing induced by
$(r,c_1,c_2)$. That framed fibre represents $[g]$. Its product
with an interior point of the ordered six-cell, with the six
coordinate normal vectors, represents $E^6[g]$. Thus the framing
determines the displayed coefficient.

For the particular embedded torus $N$ of \cref{wb:cusp-null:thm},
there are based smooth surjective quotient maps
$F_0,F_1:X\to S^3_{\mathbb H}$, where $S^3_{\mathbb H}$ is the
unit quaternion sphere based at $1$, and one based homeomorphism
$b:C_3\to S^3_{\mathbb H}$, such that
\[
 F_0^{-1}(-1)=F_1^{-1}(-1)=N,\qquad
 \operatorname{coef}_{\nu_9}[b^{-1}F_0d]=0,\quad
 \operatorname{coef}_{\nu_9}[b^{-1}F_1d]=1.
\]
The two maps have the same unframed embedded regular fibre and the
same maps on ordinary integral homology: all positive-degree maps
are zero, and both degree-zero maps are the identity.
\end{proposition}

\begin{proof}
\emph{1. Suspension coordinates.}
Use the exact cube representative $g_C=\phi gd$ of
Definition~\ref{wb:def:native-sphere-carriers}, with inverse
geometric transport $g=\phi^{-1}g_Cd^{-1}$.
The suspension calculation in
\cite[Theorem~5.41(b), p.~573, and pp.~577--578]{HatcherSpectralSequences}
gives an injection
$\pi_6(S^3)=\mathbb Z/12\to\pi_7(S^4)=\mathbb Z\oplus\mathbb Z/12$
onto the torsion summand.  The next suspension is the quotient by
the relation $2e=t$, where $e$ and $t$ generate the two summands.
Its explicit coordinate map is
\[
 L:\mathbb Z\oplus\mathbb Z/12\longrightarrow\mathbb Z/24,
 \qquad(n,[m]_{12})\longmapsto[n+2m]_{24}.
\]
Its kernel is $\langle(2,[-1]_{12})\rangle$: the congruence
$n+2m\equiv0\pmod{24}$ gives $n=2k$ and $[m]_{12}=[-k]_{12}$.
The inverse on this quotient sends $[j]_{24}$ to the class of
$(j,0)$; changing $j$ by $24$ adds $12(2,-1)$, and
$(n,m)-(n+2m,0)=m(-2,1)$ proves both inverse identities.
The restriction to the torsion summand is therefore
$[m]_{12}\mapsto[2m]_{24}$, with inverse
$[2m]_{24}\mapsto[m]_{12}$ on its even image.
The next four suspensions are isomorphisms by
\cite[Corollary~4.24, pp.~360--361]{Hatcher2002}, in the strict ranges
$8<9$, $9<11$, $10<13$ and $11<15$.  Thus $E^6$ is injective
and its image is exactly the even subgroup of $\pi_{12}(S^9)$.

For the chosen generator $\nu_9$, let $\gamma$ be the unique class
with $E^6\gamma=2\nu_9$.  It has order $12$.  There is a unique
$a\in\mathbb Z/12$ such that
\begin{equation}
 [E^6g_C]=2a\nu_9\in\pi_{12}(S^9)\cong\mathbb Z/24.
 \hypertarget{tag-038Z}{}\label{wb:cusp-s12-sixfold-suspension-class}
\end{equation}
The coefficient map in the statement has inverse
$[a]_{12}\mapsto a\gamma$: apply the injective $E^6$ and the
inverse on its even image to verify both composites. Consequently
$[g]=0$ precisely when $[E^6g]=0$ and precisely when $a=0$.
For the native map these values are zero by \cref{wb:cusp-null:thm}.

These suspensions use the ordered cube homeomorphisms $\mu_{p,q}$,
whose inverse splits the coordinates into blocks of lengths $p,q$.
The left and right representatives are
\[
 E_L=\mu_{6,3}(\mathrm{id}_{C_6}\wedge g_C)\mu_{6,6}^{-1},
 \qquad
 E_R=\mu_{3,6}(g_C\wedge\mathrm{id}_{C_6})\mu_{6,6}^{-1}.
\]
For the block interchanges, evaluation on representatives gives
$E_L=\tau_{3,6}E_R\tau_{6,6}$.  The degrees of the two interchanges
are $(-1)^{18}=1$ and $(-1)^{36}=1$, respectively.  The degree
classification \cite[Corollary~4.25, p.~361]{Hatcher2002} makes both based
homotopic to the identity.  Thus these two actual models give the
same coefficient convention.

\emph{2. The entire framed inverse image.}
For $0<r_0<\epsilon$ and
$u_0\in\overline T^2\setminus K^{(1)}$, put $y_0=[r_0,c(u_0)]$.
The literal map on the open cell gives the entire inverse image
\begin{equation}
 g^{-1}(y_0)
 =\{x\in f^{-1}(|t|=r_0):\overline\pi_f(x)=u_0\}.
 \hypertarget{tag-038Y}{}\label{wb:cusp-s12-regular-inverse-image}
\end{equation}
The real marking has local coordinates $(r,\theta,v_1,v_2,u_1,u_2)$,
and there $g=(r,u_1,u_2)$ in the target open-cell coordinates.
Hence $y_0$ is a regular value.  To identify its full fibre, choose
a real lift $\widetilde u^0$ of $u_0$, set
$p^0=B_0\widetilde u^0$ and $T_0=-\log r_0/(2\pi)$, and use
$s=\theta+iT_0$.  In $\zeta=v+Z(s)\widetilde u^0$ the monodromy
$Z(s+1)=Z(s)+B_0$ gives exactly
$(1,v)\sim(0,v+p^0)$.  This is the affine translation of the kernel
torus with lattice
$K_\Lambda=\mathbb Z\langle\widehat w,\widehat\delta\rangle$;
the monodromy fixes its two lattice generators.  The coordinates
$\alpha=v+\theta p^0\pmod{\mathbb Z^2}$ agree at the gluing and give
the product parametrization
\[
 t=r_0e^{2\pi i\theta},\qquad
 x_j=e^{2\pi i\alpha_j}r_0^{p^0_j}
       e^{2\pi i(C(t)\widetilde u^0)_j},\quad j=1,2,
\]
followed by the original period quotient.  Its inverse reads
$t/r_0$ and $\alpha=v+\theta p^0$.  Replacing $\theta$ by
$\theta+n$ replaces $v$ by $v-np^0$, and integer logarithmic
changes alter $\alpha$ only by $\mathbb Z^2$.  The long-period
translation selects the lift $u=\widetilde u^0$ uniquely modulo the
kernel lattice.  These formulas are smooth inverse maps, proving
that the whole fibre is the embedded $S^1\times T^2\cong T^3$.

The inverse differential of the ordered coordinates $(r,c_1,c_2)$
gives its normal three-framing. A complement to the tangent bundle
makes this inverse literal. Two choices of lifts differ by tangent
vectors; their affine interpolation still maps to the same three
target coordinate vectors and gives the framing homotopy.

Apply Lemma~\ref{mapsIII:based-regular-fibre-smoothing} with source
$X$, basepoint $x_0$, map $h_3\phi g:X\to S^3$, and value
$h_3\phi(y_0)$. Its entire inverse image is the torus just proved,
and the displayed local coordinates prove smoothness and rank three
near it. The target chart is
$(r,c_1,c_2)\circ\phi^{-1}h_3^{-1}$, whose inverse is
$h_3\phi\circ(r,c_1,c_2)^{-1}$ on the chosen open cell.
Thus its composite with $h_3\phi g$ is exactly the original ordered
coordinate map. The lemma first makes the map constant near $x_0$
and then supplies the relative smoothing with the same full fibre
and differential. Composing its homotopy with $\phi^{-1}h_3^{-1}$
gives the corresponding based homotopy in $B$.
Its framed Pontryagin--Thom class is therefore exactly $[g]$.
In the right suspension model, an interior value $(y_0,z_0)$ has
fibre $g^{-1}(y_0)\times\{z_0\}$ and frame consisting of the original
triple followed by six coordinate normals.  In the left model the
order is six normals followed by that triple.  The displayed block
maps relate these framed products.  Their class is $E^6[g_C]$, and
$\operatorname{coef}_{\nu_9}$ recovers $a$ exactly.
The product framing calculation takes place on $X$ times the open
six-cell; precomposition by $d\wedge\mathrm{id}_{C_6}$ is the
displayed source map defining its sphere cohomotopy class.

\emph{3. Two framings on the same embedded fibre.}
Take the particular embedded torus $N$ and bounding derivative
framing $\tau_0$ of \cref{wb:cusp-null:thm}. Keep the chosen
$\nu_9$ fixed, and use that theorem's regular value $y_0$.

We specify the target coordinate change used in the construction.
Let $S^3_{\mathbb H}$ be the unit quaternions, based at $1$, and
write imaginary quaternions as vectors in $\mathbb R^3$.  Define
\[
 \mathcal Q(w)=\frac{|w|^2-1+2w}{1+|w|^2},\qquad
 b_i(t)=\lambda_i\frac{t-y_i}{t(1-t)},
\]
where $(y_1,y_2,y_3)$ are the cube coordinates of $\phi(y_0)$,
$\lambda_1=\epsilon y_1(1-y_1)/2$, and
$\lambda_i=y_i(1-y_i)/2$ for $i=2,3$.
The map $b:C_3\to S^3_{\mathbb H}$ sends an open-cube point
$(t_1,t_2,t_3)$ to $\mathcal Q(b_1(t_1),b_2(t_2),b_3(t_3))$ and
the basepoint to $1$.  Each $b_i$ is an increasing diffeomorphism
onto $\mathbb R$, since
\[
 b_i'(t)=\lambda_i
 \frac{(t-y_i)^2+y_i(1-y_i)}{t^2(1-t)^2}>0.
\]
The inverse of $b$ away from $1$ first takes
$w=\mathcal Q^{-1}(q_{\mathbb H})=\operatorname{Im}(q_{\mathbb H})/(1-\operatorname{Re}(q_{\mathbb H}))$,
then in each coordinate takes the unique root of $b_i(t)=w_i$ in $(0,1)$.
At the cube boundary at least one $|b_i|$ tends to infinity, so
the map tends to $1$.  This proves the continuous quotient descent
and both inverse identities.  Moreover $b(\phi(y_0))=-1$ and the
derivative of $b\phi$ sends $(\partial_r,\partial_{c_1},\partial_{c_2})$
exactly to $(i,j,k)$ at $-1$.

Choose a framed tubular embedding
$\Phi:N\times D^3_\rho\to X$ with
$d_v\Phi(x,0)=\tau_0(x)$ and image disjoint from $x_0$.
For $r=|v|<\rho$, put
\[
 \sigma(r)=\exp\!\left(-\frac{r^2}{\rho^2-r^2}\right),\qquad
 q(v)=\frac{r^2-4\sigma(r)^2+4\sigma(r)v}{r^2+4\sigma(r)^2},
 \qquad q|_{\partial D^3_\rho}=1.
\]
This is $\mathcal Q(v/(2\sigma(r)))$.  The radial function
$r/(2\sigma(r))$ increases from $0$ to infinity, so $q$ induces a
homeomorphism $D^3_\rho/\partial D^3_\rho\to S^3_{\mathbb H}$.
Its inverse first applies $\mathcal Q^{-1}$ and then the inverse
radial function.  The formula is smooth and flat at the boundary,
has derivative $I$ at $0$, and takes only $0$ to $-1$.
It commutes with rotations of the imaginary coordinates.  Orient
$S^3_{\mathbb H}$ by $q$, so $(i,j,k)$ at $-1$ is positive; orient
$D^4$ to induce this boundary orientation.

Define $F_0(\Phi(x,v))=q(v)$ and $F_0=1$ outside the tube.  Its
full fibre at $-1$ is $N$ with frame $\tau_0$, so
$[F_0d]=b_*[g]=0$.  The conjugation map
$\operatorname{Ad}:S^3_{\mathbb H}\to\mathrm{SO}(3)$ is the double
cover with kernel $\{1,-1\}$: the rotation of angle $\vartheta$
about a unit imaginary axis $\omega_{\rm ax}$ is represented by
$\cos(\vartheta/2)+\omega_{\rm ax}\sin(\vartheta/2)$.
Based maps from $S^3$ and their based homotopies lift uniquely after
choosing the lift $1$, so this cover induces an isomorphism on
$\pi_3$.  Thus $\xi=[\operatorname{Ad}]$ is a generator.

In the boundary definition of $J$
\cite[\S4.1, p.~108]{CuspInvariantHatcherVB}, for a based map
$f_{\rm rot}:S^3\to\mathrm{SO}(3)$ the map
$Jf_{\rm rot}:\partial(D^4\times D^3_\rho)\to S^3_{\mathbb H}$ is
$q(f_{\rm rot}(x)v)$ on $S^3\times D^3_\rho$ and $1$ on
$D^4\times\partial D^3_\rho$.  The identity-family map $J1$
extends over the product by $q(v)$.  Pointwise multiplication and
source-pinch addition agree in $\pi_6(S^3_{\mathbb H})$: they have
the same unit and satisfy the interchange rule on cube loops.
Hence $[J\operatorname{Ad}]-[J1]$ is represented on the side by
\[
 xq(v)x^{-1}q(v)^{-1}.
\]
It equals $1$ on the cap and on $\{1\}\times D^3_\rho$, so it
descends to $S^3\wedge(D^3_\rho/\partial D^3_\rho)$.
The collapse from the boundary sphere has degree $+1$: its open
top cell has the same ordered tangent-then-normal coordinates.
The degree classification therefore identifies this commutator
with exactly $J\xi$ in these sphere conventions.  The classical
calculation states that $J\xi$ generates $\pi_6(S^3)=\mathbb Z/12$
\cite[p.~59, paragraph following Proposition~3.1]{CuspInvariantBaues1998}.

Let $u_{\rm coef}=\operatorname{coef}_{\nu_9}(b_*^{-1}J\xi)$, a unit modulo
$12$, and choose $k=u_{\rm coef}^{-1}\pmod{12}$ in $\{1,5,7,11\}$.
Orient a three-ball $D\subset N$ so its tangent orientation followed
by $\tau_0$ agrees with the orientation transported from $C_6$ by
$d$.  Choose $k$ disjoint positively oriented smaller balls in
$\operatorname{int}D$, put a rescaled copy of $q$ on each, and put
$\kappa=1$ on their complement.  Flatness makes
$\kappa:D\to S^3_{\mathbb H}$ smooth and constant on a boundary
collar; its quotient map has degree $k$ by addition of the $k$
positive relative fundamental classes.  Extend $\kappa=1$ over
$N\setminus D$, and set
\[
 \Theta_{\rm fr}=\operatorname{Ad}_{\kappa},\qquad
 \tau_1=\tau_0\Theta_{\rm fr}^{-1},\qquad
 F_1(\Phi(x,v))=q(\Theta_{\rm fr}(x)v),\quad F_1=1\text{ outside the tube}.
\]
The full inverse image $F_1^{-1}(-1)$ is the same $N$.  Its normal
derivative in $\tau_0$ coordinates is $\Theta_{\rm fr}$, so its inverse
derivative frame is exactly $\tau_1$.  Both $F_i$ are based smooth
surjections, since each normal ball maps onto $S^3_{\mathbb H}$;
compactness makes them quotient maps.  Their pointwise quotient is
$1$ outside $\Phi(D\times D^3_\rho)$ and inside that product is
\[
 F_1F_0^{-1}=\kappa(x)q(v)\kappa(x)^{-1}q(v)^{-1}.
\]
The collapse of the complement of this product ball, composed with
$d$, has degree $+1$ by excision and the chosen ordered orientation.
Its quotient is exactly
$(D/\partial D)\wedge(D^3_\rho/\partial D^3_\rho)$, with the inverse
product coordinates on its open cell.  The map
$\bar\kappa\wedge\bar q$ has degree $k$.  The degree theorem
identifies precomposition by it with multiplication by $k$ on
$\pi_6(S^3)$, through the $k$-fold source pinch and fold.
It follows that $[F_1d]-[F_0d]=kJ\xi$.
For $f_i=b^{-1}F_id$ this gives
\[
 [f_0]=0,\qquad [f_1]=\gamma,\qquad
 \operatorname{coef}_{\nu_9}[f_0]=0,\qquad
 \operatorname{coef}_{\nu_9}[f_1]=1,\qquad E^6[f_1]=2\nu_9\ne0.
\]
Thus the coefficient comparison uses the original fixed $\nu_9$;
the inverse unit degree specifies the required framing change.

The two maps have the identical unframed embedded torus $N$ and are
both surjective quotient maps.  Its fundamental class has zero
image under the same inclusion $N\to X$, because $H_3(X;\mathbb Z)=0$.
All positive ordinary homology maps induced by $f_0$ and $f_1$ are
also zero: in degree $3$ the source group is zero, in degree $6$
the target group is zero, and in every other positive degree both
groups are zero. Both $H_0$ maps are the identity.
Since $F_i=b f_i d^{-1}$, functoriality gives
$(F_i)_*=b_*(f_i)_*(d_*)^{-1}$ in every degree, proving the
same statements for the two smooth maps on $X$.
The two coefficients are $0$ and $1$ with the same unframed
fibre and homology maps. For the native $g$, the derivative framing
is $\tau_0$ and its coefficient is $0$.
\end{proof}

\begin{proposition}[The postcomposition kernel and metric slices]\hypertarget{tag-038W}{}\label{wb:cusp-s12-nullity-scope}
{\normalfont\scriptsize tag 038W}\par
For $P,Q$ of Definition~\ref{wb:def:native-sphere-carriers}, the
postcomposition homomorphism
\[
 P_*:\pi_{12}(\mathcal A\wedge\mathcal A,*)\longrightarrow\pi_{12}(C_6,*)
\]
satisfies $P_*[Q]=0$ and has a nonzero kernel. More precisely, there
are based maps $J:C_6\to\mathcal A\wedge\mathcal A$ and
$R_{\rm ret}:\mathcal A\wedge\mathcal A\to C_6$ with
$R_{\rm ret}J\simeq\mathrm{id}_{C_6}$ through based maps and $PJ=*$.
For a generator $[h]\in\pi_{12}(C_6)\cong\mathbb Z/2$, the class
$\beta=[Jh]$ is nonzero, with $(R_{\rm ret})_*\beta=[h]$ and $P_*\beta=0$.

The initial quotient maps $\widehat q_3^{\wedge2}$ and
$\widehat q_\gamma^{\wedge4}$ of
\cref{wb:cusp-cone-angle-forgotten-sphere-carriers} are based
null-homotopic and have the specified moving embedded four-spheres.
Their intrinsic radius-$R$ basepoint balls are exactly Euclidean
at the respective thresholds $T-C>4R$ and $\ell_T>4R$ on the
stated radial tail. The same initial maps therefore have the
pointed $\mathbb R^4$ limits proved there.
\end{proposition}

\begin{proof}
\emph{1. The actual postcomposition and its kernel.}
Use $P,Q,d,\phi,p$ and the ordered $\mu$ of
Definition~\ref{wb:def:native-sphere-carriers}.
For $u_1,u_2\in C_6$, the composite sends
\[
 \mu_{6,6}[u_1,u_2]\longmapsto
 \mu_{3,3}[\phi p q_2d(u_1),\phi p q_2d(u_2)]
 =\mu_{3,3}[g_C(u_1),g_C(u_2)].
\]
Each side sends the fat wedge to the basepoint, proving the literal
identity $PQ=\mu_{3,3}(g_C\wedge g_C)\mu_{6,6}^{-1}$.
Postcomposition therefore gives the homomorphism
\[
 P_*:\pi_{12}(\mathcal A\wedge\mathcal A,*)\longrightarrow
       \pi_{12}(C_6,*),\qquad [f]\longmapsto[Pf],
 \qquad P_*[Q]=0
\]
by \cref{wb:cusp-s12-native-smash-null}.  The basepoint in the source
group is the image of
$(\{v_0\}\times\mathcal A)\cup(\mathcal A\times\{v_0\})$.
Postcomposition respects addition because it commutes with the
source-sphere pinch.

This homomorphism has a nonzero kernel in the indicated dimension.
Here are the maps that exhibit it.  Choose the coordinate basepoint
$a_*\in K^{(1)}\subset\overline T^2$.  On the unreduced suspension
$S(S^1)\cong C_2$, based at its zero pole, define
\[
 j_2[r,z]=[r,z,a_*],\qquad r_2[r,z,a]=[r,z].
\]
These formulas respect both pole quotients, and
$r_2j_2=\mathrm{id}_{C_2}$ and $pj_2=*$ exactly.  The sphere
homeomorphism, after $s=r/\epsilon$, is
$[s,e^{i\theta}]\mapsto
(\sin(\pi s)\cos\theta,\sin(\pi s)\sin\theta,\cos(\pi s))$;
off the poles its inverse reads $s=\pi^{-1}\arccos x_3$ and the
phase of $x_1+ix_2$. Denote this sphere map by $\chi_2$.
The specified identification with $C_2$ is
$\psi_2=h_2^{-1}\chi_2$, with inverse $\chi_2^{-1}h_2$.
Thus the displayed $j_2,r_2$ formulas mean the suspension formulas
composed with $\psi_2^{-1}$ on the source and $\psi_2$ on the
target, respectively; their two strict composite identities follow
by cancellation of these inverse maps.

To construct the four-sphere homotopy retract, give $K=(S^1)^3$
its product cell structure. The based circle homeomorphism
$\vartheta:C_1\to S^1$ is $[a]\mapsto e^{2\pi ia}$; its inverse
chooses $0<a<1$ off $1$ and sends $1$ to the cube basepoint.
For each nonempty $I\subset\{0,1,2\}$, define
$k_I:K\to C_{|I|}$ by the ordered coordinate projection, the
product-to-smash quotient, the smash of the maps $\vartheta^{-1}$,
and $\mu_{1,\ldots,1}$. Its nonbasepoint cell coordinates are
exactly the selected ordered circle coordinates. Order the subsets as
\[
 \{0\},\{1\},\{2\},\{0,1\},\{0,2\},\{1,2\},\{0,1,2\}
\]
and put $V=\bigvee_{I\ne\varnothing}\Sigma_{\mathrm{red}}C_{|I|}$.
The summand $\Sigma_{\mathrm{red}}C_{|I|}$ is $C_{|I|+1}$ by
putting its suspension coordinate first; the inverse splits that
coordinate from the remaining block.
On the $j$th closed radial subinterval define
\[
 F:\mathcal A\longrightarrow V,\qquad
 F[r,k]=\iota_{I_j}[7r/\epsilon-(j-1),k_{I_j}(k)],
 \quad (j-1)\epsilon/7\leq r\leq j\epsilon/7.
\]
At every adjoining endpoint both expressions are the wedge
basepoint, so finite pasting and quotient descent make $F$ a
continuous based map.  The product cells $e_I$ of $K$ have zero
cellular boundary.  Suspending them gives
$H_2(\mathcal A)=\mathbb Z^3$, $H_3(\mathcal A)=\mathbb Z^3$ and
$H_4(\mathcal A)=\mathbb Z$; the edge joining the two poles has
boundary $v_1-v_0$, so $H_1(\mathcal A)=0$ and $H_0(\mathcal A)=\mathbb Z$.
All other homology groups vanish.  For $|I|=|I'|$, the map $k_I$ on
$e_{I'}$ has degree $1$ for $I=I'$ and $0$ otherwise, since in the latter
case a projected coordinate is the basepoint.  Each increasing
radial reparametrization has degree $1$.  Hence $F_*$ is the identity
matrix in these bases and the corresponding sphere bases of $V$.
The two overlapping cones on the connected space $K$ show that
$\mathcal A$ is simply connected; $V$ is a wedge of spheres of
dimensions at least two.  Thus $F$ is a homotopy equivalence by
\cite[Corollary~4.33, pp.~367--368]{Hatcher2002}.

Choose a based homotopy inverse $G:V\to\mathcal A$.  An initially
unbased inverse can be moved at the basepoint along a path by the
CW homotopy extension property.  The basepoint tracks of its two
inverse homotopies are loops in simply connected spaces, so filling
those loops and applying the same relative extension property makes
both homotopies based.  Their inverse relations determine the based
class of $G$, since $G\simeq G(FG')\simeq G'$ for any other such
inverse.  For the inclusion $i_4:C_4\to V$ of its last summand and
its wedge projection $p_4$, set $j_4=Gi_4$ and $r_4=p_4F$.  Then
$r_4j_4\simeq\mathrm{id}_{C_4}$ through based maps.  Consequently
\[
 J=(j_4\wedge j_2)\mu_{4,2}^{-1}:C_6\longrightarrow
        \mathcal A\wedge\mathcal A,
 \qquad R_{\rm ret}=\mu_{4,2}(r_4\wedge r_2)
\]
satisfy $R_{\rm ret}J\simeq\mathrm{id}_{C_6}$ through based maps and $PJ=*$
exactly.  Take $h:C_{12}\to C_6$ representing the nonzero element
of $\pi_{12}(S^6)=\mathbb Z/2$
\cite[p.~339]{Hatcher2002}.  Then
\[
 \beta=[Jh]\ne0,\qquad (R_{\rm ret})_*\beta=[h],\qquad P_*\beta=0.
\]
Thus $(R_{\rm ret})_*\beta=[h]\ne0$ detects a nonzero element of
$\ker P_*$, and the native $[Q]$ lies in the same kernel by the
proved equation $P_*[Q]=0$.

\emph{2. The initial maps and their intrinsic metric slices.}
Use the quotient maps and metrics of
\cref{wb:cusp-cone-angle-forgotten-sphere-carriers}.
In particular, $PQ$ is the continuous surjective quotient in that
theorem.  At height $r_T$ its restriction factors through
$Q_T^2\to S_T^2\wedge S_T^2$ and the two height sections
$j_{r_T}:Z_2\to B$, $z\mapsto[r_T,z]$, followed by the same $\phi$
and $\mu_{3,3}$.  The inverse of each section on its image reads the
$z$ coordinate.  For $T-C>4R$, the intrinsic entire-ball isometry is
$u\mapsto[\widetilde b_T+u]$, first in $Q_T^2$ and then in its
collapsed-fat-wedge quotient; its inverse is the unique period lift
of norm at most $R$ relative to $\widetilde b_T$.  These are the
moving embedded four-spheres and pointed $\mathbb R^4$ limits of the
same initial map of the based null-homotopy.

For the four-circle carrier, the corresponding restriction is
$F_T^4\to C_T^4\to C_T^{\wedge4}$, followed in each factor by
$[u]\mapsto[u/\ell_T]$, its height section and the ordered
$\mu_{2,2,2,2}$.  The inverse circle coordinate multiplies by
$\ell_T$.  For $\ell_T>4R$, the same formula
$u\mapsto[\widetilde b_T+u]$ and its unique period-lift inverse give
the second entire-ball isometry.  Smashing the based homotopy in
\cref{mapsIII:angle-forgotten-one-copy-null} four times gives
$[x_1,x_2,x_3,x_4],s\mapsto
[H_\gamma(x_1,s),H_\gamma(x_2,s),H_\gamma(x_3,s),H_\gamma(x_4,s)]$.
It descends through the fat wedge at every time, starts at this same
four-circle carrier and ends at the constant map.  Thus both moving
four-sphere families and their intrinsic limits occur at the
specified initial quotient maps of based null-homotopies.
\end{proof}

\Needspace{6\baselineskip}
\begin{remark}[The two cusp postcompositions]\hypertarget{tag-Q257}{}\label{wb:remark:two-cusp-postcompositions}
{\normalfont\scriptsize tag Q257}\par
Use the same $q_2:X\to\mathcal A$ and ordered coordinates of
Definition~\ref{wb:def:native-sphere-carriers}. For
$K=S^1\times\overline T^2$, put
\[
 K^{(2)}=(\{1\}\times\overline T^2)
 \cup(S^1\times\{u_1=0\})\cup(S^1\times\{u_2=0\}),
 \qquad c_3:K\to K/K^{(2)}.
\]
The two actual based maps from $\mathcal A$ are
\[
 p_3[r,z,a]=[r,c_3(z,a)]\in\Sigma_{\mathrm{red},\epsilon}(K/K^{(2)})\cong S^4,
 \qquad
 \widehat p_2[r,z,a]=[r,c(a)]\in B\cong S^3.
\]
Both send $v_0,v_1$ to their target basepoints. The quotient relations
give the continuous maps and the exact composites
$H=p_3q_2$ and $g=\widehat p_2q_2$, on the punctured cusp and on
both collapsed ends.

At the common small-radius value with
$u^0=(-1/4,1/4)$ and $\varphi_0=\pi$, the coordinates of
\cref{wb:h24:ret:chart-lemma,wb:cusp-null:thm} give
\[
 N=\{(\rho_0e^{i(\varphi-a-b)},\rho_1e^{ia},\rho_2e^{ib})\},
 \qquad
 M=\{(\rho_0e^{i(\pi-a-b)},\rho_1e^{ia},\rho_2e^{ib})\}\subset N,
\]
where all phase variables are modulo $2\pi$ and
$(\rho_0,\rho_1,\rho_2)=(r_0^{1/2},r_0^{1/4},r_0^{1/4})$.
The inclusion in parameter coordinates is $(a,b)\mapsto(\pi,a,b)$;
projection to $(a,b)$ is its inverse on the image. In the open
target cells, the map from the retained coordinates to the
angle-forgotten coordinates is
$(r,\varphi,u_1,u_2)\mapsto(r,u_1,u_2)$.
The full differential calculations in the cited proofs give the
actual ordered normal four-frame of $M$ and three-frame of $N$.
In the displayed complex chart and normal-first order, the former
induces the Lie stable tangent frame $(-\partial_a,\partial_b)$
with Arf invariant one, with the
orientation comparison of \cref{wb:h24:ret:orientation-lemma}.
The latter is the derivative frame $\mathfrak n$ of
\eqref{wb:cusp-null:eq:actual-frame}. Its proved frame homotopy
joins it to $NP$; the cap frame
\eqref{wb:cusp-null:eq:extended-frame} extends $NP$.
The cylinder followed by that cap is its framed nullbordism.

Consequently the two postcomposition homomorphisms, with the fixed
source $d$, have the proved values
\[
 (p_3)_*[q_2d]=[Hd]=\eta_4\eta_5\ne0\in\pi_6(S^4),\qquad
 (\widehat p_2)_*[q_2d]=[gd]=0\in\pi_6(S^3).
\]
The first equality is \cref{maps:src:B.3.1}, and the second is
\cref{wb:cusp-null:thm} with the target coordinates above.
In particular $[q_2d]\ne0$, since its image under $(p_3)_*$ is
nonzero. These two values are computed from the specified normal
frames and the same native quotient. The point-set quotient and
intrinsic metric applications use these same initial maps, as
proved in \cref{wb:cusp-s12-nullity-scope}.

The Hopf smash has an explicit representative in the same models.
Let $\eta_2$ be the complex Hopf map of
Definition~\ref{mapsIII:def:framing-conventions} and put
$\eta_{2,C}=h_2^{-1}\eta_2h_3:C_3\to C_2$.
On the retained target use
\[
 \phi_4[r,c_3(e^{i\varphi},[u_1,u_2])]
      =[r/\epsilon,\varphi/(2\pi),u_1,u_2],\qquad
 \phi_4^{-1}[s,a,u_1,u_2]
      =[\epsilon s,c_3(e^{2\pi ia},[u_1,u_2])]
\]
on the open cells and the basepoint elsewhere. Both maps respect
exactly the collapsed faces and are inverse. Composing $\phi_4$
with $h_4$ gives the same ordered cotangent-coordinate sphere
identification as \eqref{wb:h24:ret:eq:four-cell}, since
$\tan(\pi(a-1/2))=-\cot(\pi a)$. Hence the actual smash is
\[
 \mu_{4,2}\bigl((\phi_4Hd)\wedge\eta_{2,C}\bigr)
             \mu_{6,3}^{-1}:C_9\longrightarrow C_6.
\]
Its stable class is $\eta^2\eta=\eta^3=12\nu$ by the product
framed-collapse convention and \eqref{maps:eq:B.1.2}.
Suspension from $\pi_9(S^6)$ is an isomorphism at every subsequent
stage: $9+k<2(6+k)-1$ for every $k\geq0$
\cite[Corollary~4.24, pp.~360--361]{Hatcher2002}. Let $\nu_6$ be the inverse
image of the fixed stable generator $\nu$ under this isomorphism.
In cyclic coordinates its inverse sends $a\nu$ to $a\nu_6$ modulo
$24$. The displayed Hopf smash therefore represents
$12\nu_6\ne0$ in $\pi_9(S^6)$, as asserted for the retained map.
\end{remark}
\section{Metrics on the conical targets}
\hypertarget{tag-03D9}{}\label{wb:sphere-metric-limit-note}
\subsection{Fixed targets and radial collapse}

\begin{definition}[Radial metric data]
\hypertarget{tag-Q30W}{}\label{mapsIII:def:fixed-target-metric-data}
{\normalfont\scriptsize tag Q30W}\par
Use the based quotient maps and their target carriers
\[
 q_2^{\wedge2}:X^{\wedge2}\longrightarrow
 Z_{12}:=\mathcal A^{\wedge2},\qquad
 q_{1,\mathrm{cusp}}^{\wedge4}:X^{\wedge4}\longrightarrow
 Z_{24}^{\mathcal B}:=\mathcal B^{\wedge4},\qquad
 q_1^{\wedge4}:X^{\wedge4}\longrightarrow
 Z_{24}^{\mathcal C}:=\mathcal C^{\wedge4}.
\]
The based source identifications are the maps $\mu_{6,2}$ and
$\mu_{6,4}$ of \cref{wb:cusp-cone-sphere-map-homotopy-status}.
Write
\[
 r_T=e^{-2\pi T},\qquad
 J=\begin{pmatrix}0&1\\-1&0\end{pmatrix},\qquad
 Y_T=TJ+E(r_T),\qquad
 T(r)=-\frac{\log r}{2\pi}.
\]
Fix $T_0>0$ in the native cusp range so that $r_0=e^{-2\pi T_0}<\epsilon$,
$Y_T$ is invertible and $\ell_T>0$ for every $T\geq T_0$.
The residual matrix $E$ is the matrix of the radial cusp expansion;
write $C_E=\sup_{0\leq t\leq r_0}\|E(t)\|$.
On $\overline T^2=(\overline\Lambda\otimes\mathbb R)/\overline\Lambda$ set
\[
 h_T=\iota_T^*g_{Q_T}=dx^{\mathsf T}Y_T^{\mathsf T}Y_Tdx,
 \qquad\iota_T[a]=[Y_Ta],\quad\vartheta_T[\theta]=[\ell_T\theta].
\]
The inverse markings are $[y]\mapsto[Y_T^{-1}y]$ and
$[c]\mapsto[c/\ell_T]$. The period-one base and character coordinates
$u,\theta\in\mathbb R/\mathbb Z$ represent the circle values
$e^{2\pi iu},e^{2\pi i\theta}$; the inverse sends a circle value to
its argument divided by $2\pi$ modulo $\mathbb Z$.
Use the slice embeddings $j_T,k_T^{\mathcal B},k_T^{\mathcal C}$ of
\cref{wb:cusp-cone-target-slice-embeddings}. Thus $u=0$ represents
the positive radial angle $1\in S^1$.
\end{definition}

\begin{corollary}[Collapse in fixed compatible metrics]
\hypertarget{tag-03DG}{}\label{wb:sphere-metric-uniform-collapse}
{\normalfont\scriptsize tag 03DG}\par


Let \(D_{12},D_{24}^{\mathcal B},D_{24}^{\mathcal C}\) be arbitrary fixed
metrics compatible with the respective compact target topologies, and let
\(*\) denote the corresponding smash base point.  Then
\begin{align}
 d_H^{D_{12}}\bigl(j_T(Q_T^2),\{*\}\bigr)&\longrightarrow0,\\
 d_H^{D_{24}^{\mathcal B}}
  \bigl(k_T^{\mathcal B}(C_T^4),\{*\}\bigr)&\longrightarrow0,\\
 d_H^{D_{24}^{\mathcal C}}
  \bigl(k_T^{\mathcal C}(C_T^4),\{*\}\bigr)&\longrightarrow0.
 \hypertarget{tag-03D2}{}\label{wb:sphere-metric-Hausdorff-collapse}
\end{align}
Consequently, with the metrics restricted from any one fixed ambient
target, all three target slices converge in Gromov--Hausdorff topology to
a point, not to \(\mathbb R^4\).
\end{corollary}



\begin{proof}
Use the continuous maps $H_{12},H_{24}^{\mathcal B},H_{24}^{\mathcal C}$
of \cref{wb:cusp-cone-metric-scope} on
$[0,r_0]\times P$, $P=(\mathbb R/\mathbb Z)^4$.
For the first map the identification of $P$ with
$(\overline T^2)^2$ concatenates the two ordered coordinate pairs,
with inverse splitting them; for the other maps it is the identity
on the four circle coordinates. The inverse markings in
\cref{mapsIII:def:fixed-target-metric-data} give exactly
$H_{12}(r_T,P)=j_T(Q_T^2)$ and
$H_{24}^{\mathcal D}(r_T,P)=k_T^{\mathcal D}(C_T^4)$.
That proposition proves, for each of these maps and every compatible $D$,
\[
 A(r):=\max_{a\in P}D(H(r,a),*)\longrightarrow0.
\]
For the nonempty compact set $S_r=H(r,P)$, the definition of
Hausdorff distance gives
\[
 d_H^D(S_r,\{*\})
 =\max\left\{\sup_{x\in S_r}D(x,*),
             \inf_{x\in S_r}D(x,*)\right\}=A(r).
\]
This proves all three limits in
\eqref{wb:sphere-metric-Hausdorff-collapse} since $r_T\to0$.
The inclusion of $(S_{r_T},D|_{S_{r_T}})$ into the ambient metric
space is isometric.  Together with the point inclusion, it gives
\[
 d_{\mathrm{GH}}\bigl((S_{r_T},D|_{S_{r_T}}),\{*\}\bigr)
 \leq A(r_T)\longrightarrow0.
\]
The triangle inequality also gives
$\delta_T=\operatorname{diam}_D(S_{r_T})\leq2A(r_T)$.

For the Euclidean limit assertion, fix $R>0$.  The two points
$0$ and $(R/2,0,0,0)$ belong to the Euclidean radius-$R$ ball
and have distance $R/2$.  Every map from this ball to $S_{r_T}$
has distortion at least $R/2-\delta_T$, whose limit is $R/2$.
In any common metric space placing both Euclidean
points within $\eta$ of points of $S_{r_T}$, the triangle
inequality requires $R/2\leq2\eta+\delta_T$.
Thus the errors required for pointed Gromov--Hausdorff convergence
to $\mathbb R^4$ cannot tend to zero on that fixed ball.
\end{proof}

\subsection{Compatible length metrics and slice forms}
\begin{lemma}[Conical metrics and their induced slice forms]
\hypertarget{tag-Q30X}{}\label{mapsIII:lem:fixed-conical-metrics}
{\normalfont\scriptsize tag Q30X}\par
In the data of \cref{mapsIII:def:fixed-target-metric-data},
\begin{equation}
 T^{-2}h_T\longrightarrow h_\infty=dx^{\mathsf T}dx
 \hypertarget{tag-03DA}{}\label{wb:sphere-metric-normalized-Q-limit}
\end{equation}
in every $C^k$ norm in the fixed flat coordinates.
There are fixed compatible length metrics
\[
 D_{12}\text{ on }Z_{12},\qquad
 D_{24}^{\mathcal B}\text{ on }Z_{24}^{\mathcal B},\qquad
 D_{24}^{\mathcal C}\text{ on }Z_{24}^{\mathcal C}
\]
with the following properties.

On the inner cone of \(\mathcal A\), the metric is induced by
\begin{equation}
 g_{\mathcal A}
 =dr^2+r^2du^2+\frac{r^2}{T(r)^2}h_{T(r)}.
 \hypertarget{tag-03CW}{}\label{wb:sphere-metric-A-fixed-metric}
\end{equation}
On the cusp cone of either \(\mathcal B\) or \(\mathcal C\), it is induced
by
\begin{equation}
 g_{\mathcal D}=dr^2+r^2(du^2+d\theta^2),
 \qquad\mathcal D\in\{\mathcal B,\mathcal C\}.
 \hypertarget{tag-03CY}{}\label{wb:sphere-metric-BC-fixed-metric}
\end{equation}
The metrics on the smash powers are the metric quotients of the product
length metrics after the closed fat wedge is collapsed.  Write
\(g_{Z_{12}}\) and \(g_{Z_{24}^{\mathcal D}}\) for the resulting smooth
Riemannian tensors on the open regular strata.

Let
\[
 p_T^{12}:=j_T(0,0),\qquad
 p_T^{24,\mathcal D}:=k_T^{\mathcal D}(0,0,0,0).
\]
Put
\[
 \Sigma_T^{12}:=j_T(Q_T^2),\qquad
 \Sigma_T^{24,\mathcal D}:=k_T^{\mathcal D}(C_T^4).
\]
Let \(d_{T,12}^{\mathrm{slice}}\) and
\(d_{T,24,\mathcal D}^{\mathrm{slice}}\) be the intrinsic Riemannian path
metrics of the first fundamental forms on these embedded slices.  Define
also the restricted ambient distance functions
\[
 d_{T,12}^{\mathrm{amb}}
 :=D_{12}|_{\Sigma_T^{12}\times\Sigma_T^{12}},\qquad
 d_{T,24,\mathcal D}^{\mathrm{amb}}
 :=D_{24}^{\mathcal D}|_{\Sigma_T^{24,\mathcal D}
                         \times\Sigma_T^{24,\mathcal D}}.
\]
For all \(T\geq T_0\),
\[
 d_{T,12}^{\mathrm{amb}}\leq d_{T,12}^{\mathrm{slice}},\qquad
 d_{T,24,\mathcal D}^{\mathrm{amb}}
 \leq d_{T,24,\mathcal D}^{\mathrm{slice}}.
\]
For every sufficiently large \(T\), each of these inequalities is strict
on some pairs of points.

Then the induced first fundamental forms on the target slices satisfy the
exact identities
\begin{equation}
 \left(\frac{T}{r_T}\right)^2
  j_T^*g_{Z_{12}}
 =g_{Q_T}\oplus g_{Q_T},
 \hypertarget{tag-03CX}{}\label{wb:sphere-metric-A-slice-exact}
\end{equation}
and
\begin{equation}
 \left(\frac{\ell_T}{r_T}\right)^2
  (k_T^{\mathcal D})^*g_{Z_{24}^{\mathcal D}}
 =g_{C_T}^{\oplus4}.
 \hypertarget{tag-03CZ}{}\label{wb:sphere-metric-BC-slice-exact}
\end{equation}
\end{lemma}

\begin{proof}
Write the period matrix on the radial locus as
\[
 Y_T=TJ+E(r_T),\qquad
 J=\begin{pmatrix}0&1\\-1&0\end{pmatrix},\qquad
 E(t)=\begin{pmatrix}6m(t)&H(t)\\B(t)-H(t)&m(t)\end{pmatrix}.
\]
The entries of \(E\) are the imaginary parts on the real \(t\)-axis of
the holomorphic residual functions.  They and all their derivatives of
each fixed order are bounded on the closed cusp interval.  Put
\(C_E=\sup_{0\leq t\leq r_0}\|E(t)\|\).  Since \(J\) is orthogonal,
\[
 (T-C_E)|v|\leq |Y_Tv|\leq(T+C_E)|v|.
\]
In particular $T^{-2}h_T$ has coefficient matrix
$(J+E(r_T)/T)^{\mathsf T}(J+E(r_T)/T)$, which tends to $I$.
These coefficients are constant in the fixed torus coordinates, so
the convergence holds in every $C^k$ norm.
Invertibility on the remaining compact part of the parameter interval
gives constants \(0<c_-\leq1\leq c_+\) such that
\[
 c_-|v|\leq |M(r)v|\leq c_+|v|,
 \qquad M(r):=J+\frac{E(r)}{T(r)},\quad 0<r\leq r_0.
\]
The exact quotient-circle formula also gives
\[
 \ell_T=\min_{z\in\mathbb R}|Y_T(1,z)|\geq T-C_E.
\]
Thus the shortest nonzero periods of \(Q_T^2\) and \(C_T^4\) tend to
infinity.

\medskip\noindent
\emph{The factor metrics and their global extensions.}
On \(\mathcal A\) use
\[
 dr^2+r^2du^2+r^2dx^{\mathsf T}M(r)^{\mathsf T}M(r)dx
\]
for \(r\leq r_0\).  Extend its positive angular coefficient matrix
smoothly over the remainder of the cylinder up to \(r=\epsilon\),
keeping the radial summand \(dr^2\).  For example, continue it over a
short interval beyond \(r_0\) and take a fixed smooth convex combination
with a constant positive angular matrix, using a cutoff equal to one
on a neighborhood of the prescribed inner interval.  The choices are
made once and are independent of the slice parameter.

Complete this cylinder at \(r=0\), retaining its outer boundary torus.
The bounds for \(M(r)\) show that a level torus has diameter \(O(r)\).
A radial segment has length \(r\), and \(r\) is a \(1\)-Lipschitz
function.  Hence all sequences whose radii tend to zero converge to
one new point, uniformly in their angular coordinates; points whose
radii stay bounded below cannot converge to it.  Away from that point
the smooth tensor gives the usual cylinder topology.  These statements
identify the completion with the compact cone with its outer torus
retained.  Radial paths to its apex and approximating cylinder paths
show directly that its completed distance is a length metric.  Collapse
the outer torus to the other suspension pole, using the metric quotient
described below.  This gives the prescribed topology of \(\mathcal A\)
and its fixed factor distance \(d_{\mathcal A}\).

For \(\mathcal B\) use
\(dr^2+r^2(du^2+d\theta^2)\) on the whole open cylinder, add the
inner apex by completion, and collapse the outer torus to its separate
pole.  For \(\mathcal C\), give its attached cone the same tensor.
The boundary-coordinate map to \(\partial\mathcal M=\partial D_0\times S^1\) is
\[
 \rho_0(z,[\theta])=(\epsilon z,e^{2\pi i\theta}),\qquad
 \rho_0^{-1}(p,w)
 =\left(p/\epsilon,\arg(w)/(2\pi)\bmod\mathbb Z\right).
\]
On a collar of its compact smooth core \(\mathcal M\), with coordinate
\(\tau\geq0\) directed into the core, prescribe
\[
 d\tau^2+(\epsilon+\tau)^2(du^2+d\theta^2).
\]
Under \(r=\epsilon+\tau\), this matches the cone tensor and every
derivative at the boundary.  A fixed cutoff and convex combination with
a smooth positive tensor on the rest of \(\mathcal M\) extend it over the core.
Completing at the inner apex gives the compact cone-attachment topology,
by the same radial and diameter estimates.  The tensor is smooth across
the attaching torus.

The radial function descends to each suspension and remains
\(1\)-Lipschitz after the outer collapse: a route through that pole
has length at least
\((\epsilon-r(x))+(\epsilon-r(y))\geq|r(x)-r(y)|\).
For \(\mathcal C\), extend \(r\) by the constant \(\epsilon\) on
the core.  The extension is \(1\)-Lipschitz, as follows by integrating
its derivative along each piece of a path.  Consequently, for every
cone point of each factor \(X\),
\[
 d_X(x_0,[r,a])=r.
\]
The radial path proves the upper bound and the radial function proves
the lower bound.  This equality holds on the entire prescribed interval,
including \(r=r_0\).  The same function bounds the distance from an
inner point to the outer pole or to the core below by \(\epsilon-r\).

\medskip\noindent
\emph{The smash metrics.}
For a compact length space \((P,d)\) and a closed subset \(F\) that
is collapsed to a point, the quotient distance is
\[
 D([x],[y])=\min\{d(x,y),d(x,F)+d(y,F)\},\qquad
 D([x],*)=d(x,F).
\]
This is the closed-subset quotient metric proved in
\cref{mapsIII:lem:collapsed-subset-distance}, applied here
to the quotient map $P\to P/F$, $x\mapsto[x]$. Its restriction
from $P\setminus F$ to $(P/F)\setminus\{*\}$ has inverse taking
the unique representative. That proof gives the metric,
length and compatible quotient-topology properties, and also
applies to each outer-torus collapse above.

Give \(P=X^n\) the product distance
\[
 d_P(x,y)^2=\sum_{i=1}^n d_X(x_i,y_i)^2.
\]
It is the length metric of the product tensor on the regular stratum:
the lower bound follows by integrating the vector of component speeds,
and the upper bound follows by traversing constant-speed minimizing
factor paths simultaneously.  Such minimizing paths exist in these
compact length spaces, by taking a uniformly convergent subsequence of
constant-speed paths whose lengths decrease to the distance and using
lower semicontinuity of length.

For the fat wedge \(F=\bigcup_i\{x_i=x_0\}\), the product formula
gives
\[
 d_P(x,F)=\min_i d_X(x_i,x_0).
\]
Replacing only a minimizing factor by \(x_0\) gives the upper bound,
and every point of \(F\) gives the matching lower bound.  Applying
the quotient formula defines \(D_{12},D_{24}^{\mathcal B},
D_{24}^{\mathcal C}\) with their prescribed compact topologies.
If \(d_P(x,F)=\delta_{\rm dist}>0\), two points in a product ball of radius
\(\eta<\delta_{\rm dist}/2\) about \(x\) have direct distance at most
\(2\eta\), whereas a route through \(F\) costs at least
\(2(\delta_{\rm dist}-\eta)>2\eta\).  Thus the quotient is locally isometric
to the product off the fat wedge.  Restricting also away from the factor
singularities gives exactly the regular tensors in the statement.

At each synchronized slice point, every factor has inward distance
$r_T$, so the distance to the fat wedge is exactly $r_T$.

\medskip\noindent
\emph{The slice forms.}
The maps of Definition~\ref{wb:cusp-cone-target-slice-embeddings} are
\[
\begin{aligned}
 j_T(y_1,y_2)
 &=\bigl[[r_T,1,[Y_T^{-1}y_1]],
          [r_T,1,[Y_T^{-1}y_2]]\bigr]_{\wedge},\\
 k_T^{\mathcal B}(c_1,\ldots,c_4)
 &=\bigwedge_{i=1}^4[r_T,1,[c_i/\ell_T]],\\
 k_T^{\mathcal C}(c_1,\ldots,c_4)
 &=\bigwedge_{i=1}^4[r_T,(1,[c_i/\ell_T])].
\end{aligned}
\]
Their inverses take the unique nonbasepoint product representative
and apply $\iota_T$ or $\vartheta_T$ to its marked angular classes.
The representative independence, endpoint scope and regular
embedding assertions are proved for these exact maps and inverses
in \cref{mapsIII:lem:radial-slice-embeddings}.
On \(j_T\), \(dr=du=0\) and
\(dx=Y_T^{-1}dy\).  Its single-factor first fundamental form is
\[
 \frac{r_T^2}{T^2}
 dy^{\mathsf T}(Y_T^{-1})^{\mathsf T}Y_T^{\mathsf T}Y_TY_T^{-1}dy
 =\frac{r_T^2}{T^2}\,dy^{\mathsf T}dy.
\]
On \(k_T^{\mathcal D}\), \(d\theta=dc/\ell_T\) and
\(dr=du=0\), so its single-factor form is
\(r_T^2dc^2/\ell_T^2\).  Summing factors proves
\eqref{wb:sphere-metric-A-slice-exact} and
\eqref{wb:sphere-metric-BC-slice-exact}, including the exact equalities
of intrinsic path metrics with those of \(Q_T^2\) and \(C_T^4\).

\medskip\noindent
\emph{Intrinsic and ambient paths on the slices.}
Inclusion of a slice preserves the length of every path it contains.
Taking infima over slice paths and over all ambient paths proves the
two distance inequalities in the statement.  Their strictness can be
seen directly.  In a \(\mathcal B\)- or \(\mathcal C\)-factor,
keep \(u=0\) and take character-angle endpoints \(0,\delta_{\rm ang}\), with
\(0<\delta_{\rm ang}<1/2\), at radius \(r_T\).  The map
\[
 (r,\theta)\longmapsto(r\cos\theta,r\sin\theta)
\]
pulls the Euclidean tensor back to \(dr^2+r^2d\theta^2\) on this
sector.  The straight chord lifts to a cone path with radius between
\(r_T\cos(\delta_{\rm ang}/2)\) and \(r_T\), and has length
\[
 2r_T\sin(\delta_{\rm ang}/2)<r_T\delta_{\rm ang}.
\]
The right side is the intrinsic slice distance.  Holding the other
three factors fixed gives the strict inequality in each smash target.

For \(\mathcal A\) put \(G(r)=r^2M(r)^{\mathsf T}M(r)\).
Differentiation gives
\[
 rM'(r)=\frac{rE'(r)}{T(r)}+\frac{E(r)}{2\pi T(r)^2},\qquad
 r^{-1}G'(r)
 =2M^{\mathsf T}M+(rM')^{\mathsf T}M+M^{\mathsf T}(rM')
 \longrightarrow2I.
\]
Thus \(G'(r)\) is positive definite for every sufficiently small
positive \(r\).  Take such \(r_T\), with \(T-C_E>4\), and set
\(v=Y_T^{-1}(1,0)\).  The physical segment \(Y_Tsv=s(1,0)\)
is minimizing on the flat torus: every nonzero translated difference
has norm at least \(T-C_E-1>3\).  Choose a smooth nonnegative function
\(f\), supported in \((0,1)\) and positive somewhere.  The path
\[
 s\longmapsto[r_T-hf(s),1,[sv]],\qquad 0\leq s\leq1,
\]
stays in the cone for small \(h>0\) and has energy
\[
 \mathcal E(h)
 =\int_0^1\bigl(h^2f'(s)^2+v^{\mathsf T}G(r_T-hf(s))v\bigr)\,ds.
\]
Its derivative at zero is
\[
 \mathcal E'(0)
 =-\left(\int_0^1f(s)\,ds\right)v^{\mathsf T}G'(r_T)v<0.
\]
The unperturbed energy is the square of the intrinsic slice distance.
The perturbed path has length at most \(\sqrt{\mathcal E(h)}\),
which is strictly smaller for small \(h>0\).  Keeping the other
factor fixed proves strictness on \(\Sigma_T^{12}\).

\end{proof}

\Needspace{9\baselineskip}
\subsection{Arbitrary radial scaling}
\begin{theorem}[Complete intrinsic scalar-rescaling trichotomy]
\hypertarget{tag-03DC}{}\label{wb:sphere-metric-scaling-trichotomy}
{\normalfont\scriptsize tag 03DC}\par
For the fixed compatible length metrics of
\cref{mapsIII:lem:fixed-conical-metrics}, use the native range
\(T\geq T_0>0\), put \(r_T=e^{-2\pi T}\), and let
\(\lambda_T>0\) be arbitrary.  Set
\[
 \beta_T=\lambda_Tr_T.
\]
Choose one of the following slice and ambient metrics, with its marked
point:
\[
\begin{split}
 (\Sigma_T,d_T^{\mathrm{slice}},d_T^{\mathrm{amb}},p_T;D)
 &=(\Sigma_T^{12},d_{T,12}^{\mathrm{slice}},
 d_{T,12}^{\mathrm{amb}},p_T^{12};D_{12}),\\
 (\Sigma_T,d_T^{\mathrm{slice}},d_T^{\mathrm{amb}},p_T;D)
 &=(\Sigma_T^{24,\mathcal D},d_{T,24,\mathcal D}^{\mathrm{slice}},
 d_{T,24,\mathcal D}^{\mathrm{amb}},p_T^{24,\mathcal D};D_{24}^{\mathcal D}),
 \qquad \mathcal D\in\{\mathcal B,\mathcal C\}.
\end{split}
\]
For each choice let \(Z\) denote its ambient target.
Along every sequence \(T\to\infty\), a subsequence has
\(\beta_T\to\beta_{\rm lim}\in[0,\infty]\).  Its intrinsic pointed slice limit is
\begin{equation}
 (\Sigma_T,\lambda_Td_T^{\mathrm{slice}},p_T)
 \longrightarrow
 \begin{cases}
 \{\mathrm{pt}\},&\beta_T\longrightarrow0,\\
 ((\mathbb R/L\mathbb Z)^4,0),
       &\beta_T\longrightarrow L\in(0,\infty),\\
 (\mathbb R^4,0),&\beta_T\longrightarrow\infty.
 \end{cases}
 \hypertarget{tag-03DD}{}\label{wb:sphere-metric-scaling-trichotomy-result}
\end{equation}
All three convergences are pointed Gromov--Hausdorff.  In the middle
case the flat tensors converge smoothly in the fixed marked torus.
The physical limiting lattice for the \(Q_T^2\) route is
\(L\operatorname{diag}(J,J)\mathbb Z^4\), with its explicit isometry
\([z]\mapsto[\operatorname{diag}(J,J)^{\mathsf T}z]\) to
\(\mathbb R^4/L\mathbb Z^4\).  In the last case the intrinsic balls
are eventually isometric to the whole Euclidean balls, with exactly
Euclidean pullback tensors.  Intrinsic pointed convergence to
\(\mathbb R^4\) holds precisely when \(\beta_T\to\infty\).
For every admissible \(T\) and every \(\lambda_T>0\),
\begin{equation}
 \lambda_TD(p_T,*)=\beta_T.
 \hypertarget{tag-03DI}{}\label{wb:sphere-metric-vertex-distance}
\end{equation}

The inclusion of a slice preserves its path lengths and gives
\(d_T^{\mathrm{amb}}\leq d_T^{\mathrm{slice}}\).
In the regime \(\beta_T\to\infty\), for every \(R>0\),
\begin{equation}
 \sup_{x,y\in\overline B_{\lambda_Td_T^{\mathrm{slice}}}(p_T,R)}
 \left|\lambda_Td_T^{\mathrm{amb}}(x,y)
            -\lambda_Td_T^{\mathrm{slice}}(x,y)\right|
 \longrightarrow0,
 \hypertarget{tag-03D7}{}\label{wb:sphere-metric-intrinsic-ambient-agreement}
\end{equation}
and there are \(\varepsilon_{T,R}\to0\) with
\begin{equation}
 \overline B_{\lambda_Td_T^{\mathrm{slice}}}(p_T,R)
 \subset
 \overline B_{\lambda_Td_T^{\mathrm{amb}}}(p_T,R)
 \subset
 \overline B_{\lambda_Td_T^{\mathrm{slice}}}
                (p_T,(1+\varepsilon_{T,R})R).
 \hypertarget{tag-03D8}{}\label{wb:sphere-metric-intrinsic-ambient-ball-comparison}
\end{equation}
Here every ball is a subset of the same slice.  The ambient-restricted
slices therefore have the pointed \(\mathbb R^4\) limit as well.
For \(\beta_T\to0\), both slice distances have diameter tending to
zero and their absolute discrepancy tends uniformly to zero.
For \(\beta_T\to L\in(0,\infty)\), the explicit paths in the proof
give a positive limiting discrepancy and a fixed ball radius for
which the second inclusion in
\eqref{wb:sphere-metric-intrinsic-ambient-ball-comparison} fails with
any error tending to zero.

Every intrinsic unrolling sends the marked fat-wedge point to distance
\(\beta_T\to\infty\).  Moreover, for each of these fixed smash metrics,
\((Z,qD,*)\) has no proper pointed Gromov--Hausdorff convergent
subsequence as \(q\to\infty\).  In particular the Euclidean limits
here are not proper pointed Gromov--Hausdorff tangents of the fixed
metric at \(*\); the proof gives exact marked-distance maps and a
packing obstruction at that point.

In the regime \(\beta_T\to\infty\), the full ambient limits are
\begin{align}
 (Z_{12},\lambda_TD_{12},p_T^{12})
 &\longrightarrow(\mathbb R^8,0),\\
 (Z_{24}^{\mathcal D},\lambda_TD_{24}^{\mathcal D},p_T^{24,\mathcal D})
 &\longrightarrow(\mathbb R^{12},0),
 \qquad \mathcal D\in\{\mathcal B,\mathcal C\}.
 \hypertarget{tag-03D4}{}\label{wb:sphere-metric-arbitrary-unrolling-ambient-bubbles}
\end{align}
These are pointed Gromov--Hausdorff limits of the whole metric targets.
There are explicit regular coordinate neighborhoods whose domains
exhaust the Euclidean limits, whose images eventually contain every
fixed rescaled target ball, and on which the pulled-back tensors
converge in every \(C^k\) norm.  On these maps the slice is exactly
the vertical coordinate subspace.
\end{theorem}

\begin{proof}
\medskip\noindent
\emph{The marked flat tori.}
Write the period matrix and residual as
\[
 Y_T=TJ+E(r_T),\qquad
 J=\begin{pmatrix}0&1\\-1&0\end{pmatrix},\qquad
 E(t)=\begin{pmatrix}6m(t)&H(t)\\B(t)-H(t)&m(t)\end{pmatrix},
 \qquad M_T=Y_T/T.
\]
The residual entries are the imaginary parts of the defining functions
holomorphic on a neighborhood of the fixed smaller cusp disc.
Their derivatives of every finite order are bounded there.  For
\(C_E=\sup_{0\leq t\leq r_0}\|E(t)\|\),
\[
 \sigma_{\min}(M_T)\geq1-C_E/T,\qquad
 \|M_T\|\leq1+C_E/T.
\]
The first bound is useful as a positive lower bound for \(T>C_E\);
the markings are invertible throughout the native
range.  They are
\[
 \iota_T[a]=[Y_Ta],\qquad
 \iota_T^{-1}[y]=[Y_T^{-1}y],\qquad
 \vartheta_T[a]=[\ell_Ta],\qquad
 \vartheta_T^{-1}[c]=[c/\ell_T].
\]
Here \(\vartheta_T\) is the arclength marking, and \(r_T\) is the
radial parameter.

On \(N=\mathbb R^4/\mathbb Z^4\) define the slice markings
\[
\begin{aligned}
 \Psi_{\mathcal A,T}([a_1],[a_2])
 &=j_T(\iota_T[a_1],\iota_T[a_2]),\\
 \Psi_{\mathcal D,T}([\theta_1],\ldots,[\theta_4])
 &=k_T^{\mathcal D}(\vartheta_T[\theta_1],\ldots,
                    \vartheta_T[\theta_4]).
\end{aligned}
\]
Their target representatives are the smash tuples of
\([r_T,1,[a_i]]\), \([r_T,1,[\theta_i]]\), or
\([r_T,(1,[\theta_i])]\), respectively.  The inverse takes the unique
nonbasepoint smash tuple and reads those native angle classes.
Both markings send zero to \(p_T\).  The exact slice pullbacks from
\cref{mapsIII:lem:fixed-conical-metrics} give
\[
\begin{aligned}
 \Psi_{\mathcal A,T}^*(\lambda_T^2g_T^{\mathrm{slice}})
 &=\beta_T^2\operatorname{diag}(M_T^{\mathsf T}M_T,
                               M_T^{\mathsf T}M_T),\\
 \Psi_{\mathcal D,T}^*(\lambda_T^2g_T^{\mathrm{slice}})
 &=\beta_T^2 I_4.
\end{aligned}
\]
In the second line the factors \(\ell_T\) cancel because the
single-factor metric is \((r_T/\ell_T)^2g_{C_T}\) and the marking
pulls \(g_{C_T}\) back to \(\ell_T^2d\theta^2\).

Put
\[
 A_{\mathcal A,T}=\beta_T\operatorname{diag}(M_T,M_T),
 \qquad A_{\mathcal D,T}=\beta_T I_4,
 \qquad O=\operatorname{diag}(J,J).
\]
The map \([a]\mapsto[A_{\nu,T}a]\) identifies the corresponding
marked tensor on \(N\) isometrically with the Euclidean quotient
\(\mathbb R^4/A_{\nu,T}\mathbb Z^4\).  Its inverse is
\([z]\mapsto[A_{\nu,T}^{-1}z]\); changing a representative by an
integer vector changes the other representative by its displayed
period.  The physical maps from \(\mathbb R^4\) to the slices
are consequently
\[
\begin{aligned}
 z=(z_1,z_2)&\longmapsto
 j_T([(T/\beta_T)z_1],[(T/\beta_T)z_2]),\quad z_i\in\mathbb R^2,\\
 z&\longmapsto
 k_T^{\mathcal D}([(\ell_T/\beta_T)z_1],\ldots,
                 [(\ell_T/\beta_T)z_4]).
\end{aligned}
\]
The bracketed arguments lie in \(Q_T\) and \(C_T\), respectively.
These maps have exactly the period lattices \(A_{\nu,T}\mathbb Z^4\).

Every sequence of \(\beta_T/(1+\beta_T)\) has a subsequence converging
in \([0,1]\); applying \(u\mapsto u/(1-u)\) on \([0,1)\), and
assigning infinity to the endpoint \(1\), gives the required
subsequence for \(\beta_T\).
For \(\beta_T\to0\), choose differences of native representatives
in \([-1/2,1/2]^4\).  Their Euclidean norm is at most one, so
\[
 \operatorname{diam}(\lambda_Td_{T,12}^{\mathrm{slice}})
 \leq\beta_T(1+C_E/T),\qquad
 \operatorname{diam}(\lambda_Td_{T,24,\mathcal D}^{\mathrm{slice}})
 =\beta_T.
\]
The maps to a single marked point have distortions bounded by these
diameters.  For \(\beta_T\to L\in(0,\infty)\), the native tensors
converge to \(L^2I_4\).  All positive-order derivatives in the native
variables vanish, and length comparison by factors tending to one
proves both smooth and Gromov--Hausdorff convergence.  In physical
coordinates the \(\mathcal A\) limiting lattice is \(LO\mathbb Z^4\).
The exact isometry to the square torus is \([z]\mapsto[O^{\mathsf T}z]\),
with inverse \([w]\mapsto[Ow]\).  For \(\mathcal B,\mathcal C\)
it is the identity.  In native coordinates the limiting square-torus
map is \([a]\mapsto[La]\), with inverse \([w]\mapsto[w/L]\).

For \(\beta_T\to\infty\), the shortest nonzero physical periods obey
\[
 \operatorname{sys}(A_{\mathcal A,T}\mathbb Z^4)
 \geq\beta_T(1-C_E/T),\qquad
 \operatorname{sys}(A_{\mathcal D,T}\mathbb Z^4)=\beta_T.
\]
Fix \(R>0\).  Eventually both systoles exceed \(4R\).
A nonzero translated difference of two vectors in the Euclidean
\(R\)-ball then has norm greater than \(2R\); their direct difference
has norm at most \(2R\).  Hence the Euclidean quotient map is an
isometry on that whole ball.  Every quotient point at distance at
most \(R\) has a shortest representative in the same ball, since
the lattice is discrete.  The physical slice maps above therefore
give onto isometries of whole balls, with exactly Euclidean tensors.
This proves the last line of
\eqref{wb:sphere-metric-scaling-trichotomy-result}.
Conversely, a sequence for which \(\beta_T\) fails to tend to
infinity has a bounded subsequence.  The displayed diameter estimates
bound its diameters by a constant \(D_0\).  Any correspondence to a
Euclidean \(R\)-ball that covers its two opposite boundary points
has distortion at least \(2R-D_0\).  Taking \(2R>D_0\) rules out
pointed convergence to \(\mathbb R^4\) on that subsequence and
proves the stated equivalence.

\medskip\noindent
\emph{The marked point and the ambient charts.}
The global radial functions constructed in
\cref{mapsIII:lem:fixed-conical-metrics} are \(1\)-Lipschitz and
vanish at the inward factor vertices.  Radial segments attain their
lower bounds, giving
\[
 d_X(x_0,[r,a])=r\qquad(0\leq r\leq r_0).
\]
This includes \(r=r_0\).  In the product the distance to the fat wedge
is the minimum of the factor distances; the metric quotient satisfies
\[
 D([x],[y])=\min\{d_P(x,y),d_P(x,F)+d_P(y,F)\}.
\]
Every synchronized slice point has \(d_P(x,F)=r_T\), which proves
\eqref{wb:sphere-metric-vertex-distance} for the full native range.
In particular, radial travel to the fixed outer cusp level has the
exact rescaled length
\begin{equation}
 \lambda_T(r_0-r_T)
 =\beta_T\left(\frac{r_0}{r_T}-1\right)\longrightarrow\infty
 \quad\text{as }T\to\infty\text{ and }\beta_T\to\infty.
 \hypertarget{tag-03DB}{}\label{wb:sphere-metric-outward-shortcut-exclusion}
\end{equation}

For the ambient unrolling claims take \(\beta=\beta_T\to\infty\).
The native angles \(u,\theta\) have period one, with the base angle
\(e^{2\pi iu}\).  In one \(\mathcal A\)-factor define
\[
 \Phi_{\mathcal A,T}(a,b,\xi)
 =\left[r_T(1+a/\beta),e^{2\pi ib/\beta},
               [(T/\beta)Y_T^{-1}\xi]\right],\qquad \xi\in\mathbb R^2,
\]
on
\[
 |a|<\beta/2,\quad |b|<\beta/4,\quad
 (T/\beta)Y_T^{-1}\xi\in(-1/4,1/4)^2.
\]
For all large \(T\), its radial interval lies in \((0,r_0)\).
The inverse takes the unique native lifts and sends them to
\[
 a=\beta(r/r_T-1),\qquad b=\beta u_{\mathrm{lift}},\qquad
 \xi=(\beta/T)Y_Tx_{\mathrm{lift}}.
\]
In the other two factors use
\[
\begin{aligned}
 \Phi_{\mathcal B,T}(a,b,\xi)
 &=[r_T(1+a/\beta),e^{2\pi ib/\beta},[\xi/\beta]],\\
 \Phi_{\mathcal C,T}(a,b,\xi)
 &=[r_T(1+a/\beta),(e^{2\pi ib/\beta},[\xi/\beta])],
\end{aligned}
\]
on \(|a|<\beta/2\), \(|b|,|\xi|<\beta/4\).  Their inverse is
\(a=\beta(r/r_T-1)\), \(b=\beta u_{\mathrm{lift}}\),
\(\xi=\beta\theta_{\mathrm{lift}}\).
The \(\mathcal C\) boundary coordinate uses the map
\(\rho_0\) to \(\partial\mathcal M\), with inverse \(\rho_0^{-1}\).
Each pair of displayed maps composes to the identity.  The products
of two \(\mathcal A\) charts or four \(\mathcal D\) charts followed
by the smash quotient are open regular coordinate neighborhoods:
off the fat wedge that quotient has singleton classes and is open.
Their inverses first take the unique product tuple and then the
factor inverses.  Their domains contain every fixed coordinate ball
for all sufficiently large \(T\), by the preceding period bounds.
On \(a=b=0\), they are exactly the physical slice maps already given.

Set
\[
 \alpha=1+a/\beta,\quad
 \tau_T(a)=T-\frac{\log(1+a/\beta)}{2\pi},\quad
 F_T(a)=T/\tau_T(a),\quad U_T(a)=F_T(a)Y_{\tau_T(a)}Y_T^{-1}.
\]
Substitution in the fixed tensors gives the exact single-factor
pullbacks
\[
\begin{aligned}
 \lambda_T^2\Phi_{\mathcal A,T}^*g_{\mathcal A}
 &=da^2+\alpha^2db^2+\alpha^2d\xi^{\mathsf T}U_T^{\mathsf T}U_Td\xi,\\
 \lambda_T^2\Phi_{\mathcal D,T}^*g_{\mathcal D}
 &=da^2+\alpha^2(db^2+d\xi^2).
\end{aligned}
\]
At \(a=0\), \(U_T=I\), so these pullbacks are exactly Euclidean
on the slice tangent spaces.

\medskip\noindent
\emph{All derivative orders at arbitrary unrolling rates.}
For \(j\geq1\),
\[
 \tau_T^{(j)}(a)=\frac{(-1)^j(j-1)!}{2\pi(\beta+a)^j},\qquad
 e^{-2\pi\tau_T(a)}=r_T(1+a/\beta),
\]
and the residual argument satisfies
\[
 \frac{d^j}{da^j}E(r_T(1+a/\beta))
 =\left(\frac{r_T}{\beta}\right)^jE^{(j)}(r_T(1+a/\beta)).
\]
On each fixed interval \(|a|\leq L_0\),
\(\tau_T-T=O_{L_0}(\beta^{-1})\) and
\(F_T-1=O_{L_0}((\beta T)^{-1})\).
Differentiating \(F_T\tau_T=T\) yields
\[
 F_T^{(j)}=-\frac1{\tau_T}
 \sum_{i=0}^{j-1}\binom ji F_T^{(i)}\tau_T^{(j-i)}.
\]
Induction gives \(F_T^{(j)}=O_{L_0,j}(T^{-1}\beta^{-j})\).
The term with \(i=0\) has this order directly; every term with
\(i\geq1\) has the additional factor \(T^{-1}\).
The exact cancellation
\[
 U_T-I=\bigl[F_T E(r_T(1+a/\beta))-E(r_T)\bigr]Y_T^{-1}
\]
and \(\|Y_T^{-1}\|\leq(T-C_E)^{-1}\) now give, for each fixed \(k\),
\[
 \|U_T-I\|_{C^k([-L_0,L_0])}
 =O_{L_0,k}\left(\frac1{\beta T^2}
                       +\frac{r_T}{\beta T}\right).
\]
Indeed the order-zero residual difference is \(O(r_T/\beta)\).
At order \(j\geq1\) the product rule gives the bound
\(O_{L_0,j}(\beta^{-j}T^{-2}
                  +r_T^j\beta^{-j}T^{-1})\);
the mixed terms have an additional inverse power of \(T\).
Since eventually \(\beta\geq1\) and \(0<r_T\leq1\), these bounds
imply the displayed common \(C^k\) bound.
For the raw matrix ratio the identity
\[
 Y_{\tau_T(a)}Y_T^{-1}-I
 =\bigl[(\tau_T(a)-T)J+E(r_T(1+a/\beta))-E(r_T)\bigr]Y_T^{-1}
\]
also gives \(O_{L_0,k}((\beta T)^{-1})\) in every fixed \(C^k\) norm.
Thus the exact pullbacks converge to their Euclidean tensors in every
order.  Their coefficients depend only on \(a\), so all other
spatial derivatives vanish.  These estimates require only
\(T\to\infty\) and \(\beta_T\to\infty\) separately; the function
\(T\mapsto\lambda_T\) may be arbitrary.

\medskip\noindent
\emph{Whole ambient balls.}
Let \(c_->0\) be a uniform lower bound for
\(\sigma_{\min}(Y_{T(r)}/T(r))\) on the inner cusp, furnished by
\cref{mapsIII:lem:fixed-conical-metrics}, and put
\(c_*=\min(1,c_-)\), \(\rho_{\rm exit}=\min(1/2,c_*/8)\).
A path from the center leaving the native window
\[
 r_T/2<r<3r_T/2,\qquad |u_{\mathrm{lift}}|<1/4,
\]
with the corresponding vertical lifts in \((-1/4,1/4)\), first
reaches a radial boundary or an angular boundary.  The first event
costs at least \(r_T/2\).  Before a radial exit, reaching an angular
boundary costs at least \(c_*r_T/8\), by the variation \(1/4\) of
one lifted angle and the radius lower bound \(r_T/2\).
This includes a path that subsequently goes to an inward vertex, an
outer pole or the core.  Product projections are \(1\)-Lipschitz,
and a quotient route through the fat wedge costs \(r_T\) from the
center.  The rescaled distance to the complement of the product
chart is therefore at least \(\rho_{\rm exit}\beta\).
The inward and fat-wedge distances are \(\beta\); the outward
radial distance is given by
\eqref{wb:sphere-metric-outward-shortcut-exclusion}.
The coordinate period lattices are
\[
 \beta\mathbb Z\oplus(\beta/T)Y_T\mathbb Z^2
 \quad\text{and}\quad \beta\mathbb Z^2.
\]
Their least nonzero Euclidean lengths are at least
\(\beta(1-C_E/T)\) and equal to \(\beta\), respectively.
The first-exit estimate controls paths whose radii vary while they
move in these angular directions.

Fix \(R>0\).  Every rescaled \(R\)-ball is in the regular chart
once \(\rho_{\rm exit}\beta>R\).  The compact length targets are geodesic.
A minimizing curve between two of its points has length at most
\(2R\) and lies in the rescaled \(3R\)-ball.  For
\(\rho_{\rm exit}\beta>3R\), that entire curve stays in the chart.  A path
through the fat wedge has length at least \(2(\beta-R)>2R\) for
all large \(T\), so it cannot be minimizing.  This proves exclusion
of both quotient and global excursions before any local tensor
comparison is used.

On the full native window, \(1/2<\alpha<3/2\) and
\(\tau_T=T+O(1)\).  The singular-value estimates imply, for all
large \(T\),
\[
 \frac14\leq\sigma_{\min}(U_T)\leq\|U_T\|\leq3.
\]
For each rescaled product tensor this gives the common bounds
\[
 \frac18|v|\leq\sqrt{g_T(v,v)}\leq5|v|.
\]
A minimizing origin path to a point of the \(3R\)-ball has
coordinate length at most \(24R\).  Thus the \(R\)-ball has
coordinates in \(B_{8R}(0)\), and all minimizing paths between its
points lie in \(B_{24R}(0)\).
Choose \(0<e_{T,R}<1/2\), tending to zero, so that on the fixed
coordinate ball of radius \(24R+1\) the rescaled norms lie between
\((1-e_{T,R})|v|\) and \((1+e_{T,R})|v|\).
Minimizing paths give the lower bound and straight coordinate
segments give the upper bound in
\[
 (1-e_{T,R})|z-w|
 \leq\lambda_TD(\Phi_Tz,\Phi_Tw)
 \leq(1+e_{T,R})|z-w|
\]
for points in the target \(R\)-ball.  Comparing with the origin also
gives
\[
 B_{R/(1+e_{T,R})}(0)
 \subset\Phi_T^{-1}(B_R(p_T;\lambda_TD))
 \subset B_{R/(1-e_{T,R})}(0),
\]
with the corresponding inclusions for closed balls.
Radial projection of the middle set to the Euclidean \(R\)-ball
moves coordinates by \(O_R(e_{T,R})\); its distortion and its image
density error have the same bound.  These are Gromov--Hausdorff
approximations of whole balls, proving
\eqref{wb:sphere-metric-arbitrary-unrolling-ambient-bubbles}.
The chart domains exhaust \(\mathbb R^8\) or \(\mathbb R^{12}\),
their images eventually contain each fixed rescaled target ball,
and the tensor estimates prove the claimed smooth exhaustion.

\medskip\noindent
\emph{The two distances on a slice.}
Taking infima over the included slice paths gives
\(d_T^{\mathrm{amb}}\leq d_T^{\mathrm{slice}}\), for every native
\(T\).  In the unrolling regime, the physical slice map identifies
its intrinsic \(R\)-ball exactly with the Euclidean \(R\)-ball in
\(a=b=0\), for all large \(T\).  For two points of this ball the
preceding ambient estimate consequently gives
\[
 0\leq\lambda_Td_T^{\mathrm{slice}}(x,y)
           -\lambda_Td_T^{\mathrm{amb}}(x,y)
 \leq e_{T,R}|\xi-\eta|\leq2Re_{T,R}.
\]
This proves \eqref{wb:sphere-metric-intrinsic-ambient-agreement}.
The first inclusion in
\eqref{wb:sphere-metric-intrinsic-ambient-ball-comparison} follows
from the path inequality.  A point of its middle ball lies in the
same full ambient chart; its exact slice equations are \(a=b=0\).
The origin comparison gives \(|\xi|\leq R/(1-e_{T,R})\).
Apply the intrinsic whole-ball isometry on the Euclidean \(2R\)-ball,
which also holds eventually.  Its distance from the intrinsic origin
is therefore \(|\xi|\).  Taking
\[
 \varepsilon_{T,R}=\frac{e_{T,R}}{1-e_{T,R}}
\]
proves the second inclusion.  The restricted ambient distance
estimates and these inclusions supply whole-ball approximations to
\(\mathbb R^4\).

For \(\beta_T\to0\), the same path inequality bounds both diameters
and their absolute discrepancy by the intrinsic diameter, which tends
to zero.  Every fixed positive-radius ball eventually contains the
whole slice.
For \(\beta_T\to L\in(0,\infty)\), put \(\delta=1/4\).
In a \(\mathcal B\)- or \(\mathcal C\)-factor take the two
character coordinates \(0,\delta\) at radius \(r_T\), with all
other factor coordinates at their marked origins.  Their rescaled
intrinsic distance is \(\beta_T\delta\).
The inverse of the map
\((r,\theta)\mapsto(r\cos\theta,r\sin\theta)\) on the sector
\(0\leq\theta\leq\delta\) takes the straight chord to an
ambient path of rescaled length \(2\beta_T\sin(\delta/2)\).
Its radius lies between \(r_T\cos(\delta/2)\) and \(r_T\), so the
path stays in the native cone.

For \(\mathcal A\), use the first native vertical displacement
\(\delta e_1\) in one factor.  Its rescaled intrinsic distance is
\(\beta_T\delta|M_Te_1|\) for all large \(T\): a nonzero integer
translate has norm at least
\(\beta_T(1-C_E/T)(1-\delta)\), eventually greater than the direct
norm \(\beta_T(1+C_E/T)\delta\).
Set, for \(0\leq s\leq1\),
\[
 w(s)=(1-s+s\cos\delta,s\sin\delta),\qquad
 \rho(s)=|w(s)|,\qquad\theta(s)=\arg w(s).
\]
Then
\[
\begin{gathered}
 \rho^2=1-2s(1-s)(1-\cos\delta)\geq\cos^2(\delta/2)>0,\\
 \theta'=\frac{\sin\delta}{\rho^2},\qquad
 (\rho')^2+\rho^2(\theta')^2=4\sin^2(\delta/2).
\end{gathered}
\]
The target path is
\([r_T\rho(s),1,[(\theta(s),0)]]\) in that factor, with the other
factor fixed.  Writing \(M(r)=Y_{T(r)}/T(r)\), its rescaled
squared speed is
\[
 \beta_T^2\bigl((\rho')^2+
          \rho^2|M(r_T\rho)e_1|^2(\theta')^2\bigr).
\]
Here \(T(r_T\rho)=T-\log\rho/(2\pi)=T+O(1)\) uniformly, so
\(M(r_T\rho)\to J\).  The path length tends to
\(2L\sin(\delta/2)\), while the intrinsic endpoint distance tends
to \(L\delta\).  In all three routes these points thus satisfy
\[
 \liminf_{T\to\infty}
 \lambda_T\bigl(d_T^{\mathrm{slice}}(p_T,x_T)
                  -d_T^{\mathrm{amb}}(p_T,x_T)\bigr)
 \geq L\bigl(\delta-2\sin(\delta/2)\bigr)>0.
\]
For any fixed \(R>L\delta\), both endpoints eventually lie in the
intrinsic \(R\)-ball.  For a fixed \(R\) strictly between
\(2L\sin(\delta/2)\) and \(L\delta\), the point \(x_T\) lies
in the ambient \(R\)-ball and remains outside every intrinsic
\((1+\varepsilon_T)R\)-ball with \(\varepsilon_T\to0\).
These exact paths prove the asserted finite-positive scope.

\medskip\noindent
\emph{The marked fat-wedge point and proper tangents.}
The intrinsic equivalence already proved and
\eqref{wb:sphere-metric-vertex-distance} show that every intrinsic
unrolling has \(\lambda_TD(p_T,*)=\beta_T\to\infty\).
Each displayed chart sends its origin to \(p_T\), and each bounded
chart image eventually avoids the fat wedge.  This describes the
marked-point behavior under the displayed coordinate maps.

For the obstruction at \(*\), put \(r_c=r_0/2\), \(L_c=r_0/4\).
In a factor \(X=\mathcal A,\mathcal B,\mathcal C\), respectively,
define the radial paths
\[
\begin{aligned}
 \gamma_{\mathcal A}(s)&=[r_c+s,1,[0,0]],\\
 \gamma_{\mathcal B}(s)&=[r_c+s,1,[0]],\\
 \gamma_{\mathcal C}(s)&=[r_c+s,(1,[0])],
 \qquad 0\leq s\leq L_c.
\end{aligned}
\]
The global radial lower bound and these radial paths prove
\[
 d_X(\gamma(s),\gamma(t))=|s-t|,\qquad
 d_X(x_0,\gamma(s))=r_c+s.
\]
For \(q>1/r_c\), let \(z_q\) be the corresponding point of radius
\(1/q\) and all zero angular coordinates, and form
\[
 x_j=[z_q,\gamma(2j/q),\gamma(0),\ldots,\gamma(0)]_{\wedge},
 \qquad 0\leq j\leq\lfloor L_cq/2\rfloor.
\]
For the two-factor target the tuple ends after its second entry;
for the four-factor targets its last two entries are \(\gamma(0)\).
Each product point has distance \(1/q\) to the fat wedge.
The product distance between its \(i\)-th and \(j\)-th members is
\(2|i-j|/q\).  The exact quotient formula therefore gives
\[
 qD(*,x_j)=1,\qquad qD(x_i,x_j)=2\quad(i\ne j).
\]
The number \(\lfloor L_cq/2\rfloor+1\) tends to infinity.
A compact metric ball has a finite bound on the cardinality of its
sets with pairwise distances at least one, by a finite cover with
balls of radius less than \(1/2\).  Proper pointed Gromov--Hausdorff
convergence would send these unit-ball points, with arbitrarily small
distortion, into the compact radius-two ball of the limit, retaining
pairwise separation greater than one.  Their unbounded cardinalities
contradict that bound.  Hence no sequence of scales \(q\to\infty\)
has a proper pointed limit at \(*\).  For the unrolling scales,
\(\lambda_T=\beta_T/r_T\to\infty\), so this obstruction applies
in addition to their exact marked-point escape.  This proves the
stated tangent scope for these fixed metrics and this marked point.
\end{proof}

\subsection{Native and unit-volume scaling}
\begin{corollary}[Native Euclidean slice and ambient limits]
\hypertarget{tag-03D6}{}\label{wb:sphere-metric-fixed-realization}
{\normalfont\scriptsize tag 03D6}\par
For the metrics, slices and marked points of
\cref{mapsIII:lem:fixed-conical-metrics},
\begin{align}
 \left(\Sigma_T^{12},\frac{T}{r_T}d_{T,12}^{\mathrm{slice}},
       p_T^{12}\right)
 &\longrightarrow(\mathbb R^4,0),\\
 \left(\Sigma_T^{24,\mathcal D},
       \frac{\ell_T}{r_T}d_{T,24,\mathcal D}^{\mathrm{slice}},
       p_T^{24,\mathcal D}\right)
 &\longrightarrow(\mathbb R^4,0)
 \hypertarget{tag-03DE}{}\label{wb:sphere-metric-slice-fixed-blowups}
\end{align}
in pointed local-isometric and pointed Gromov--Hausdorff topology, and
smoothly on the exhausting Euclidean slice balls.

The full ambient pointed blow-ups are
\begin{align}
 \left(Z_{12},\frac{T}{r_T}D_{12},p_T^{12}\right)
 &\longrightarrow(\mathbb R^8,0),\\
 \left(Z_{24}^{\mathcal D},
       \frac{\ell_T}{r_T}D_{24}^{\mathcal D},
       p_T^{24,\mathcal D}\right)
 &\longrightarrow(\mathbb R^{12},0).
 \hypertarget{tag-03D3}{}\label{wb:sphere-metric-ambient-fixed-blowups}
\end{align}
as pointed metric spaces in pointed Gromov--Hausdorff topology.  More
precisely, for every fixed radius the corresponding rescaled target ball is
eventually contained in an open regular-stratum coordinate neighborhood,
and the pulled-back Riemannian tensors converge there in every \(C^k\) norm.
The smooth convergence uses an exhaustion by these regular-stratum
neighborhoods.  Under these limits, the slice \(\mathbb R^4\) is
respectively the coordinate subspace
obtained by fixing the two radial and two base-angle coordinates, or by
fixing the four radial and four base-angle coordinates.
\end{corollary}

\begin{proof}
In \cref{wb:sphere-metric-scaling-trichotomy}, set
\[
 \lambda_T=T/r_T,\quad\beta_T=T\quad\text{on }Z_{12},
 \qquad
 \lambda_T=\ell_T/r_T,\quad\beta_T=\ell_T
       \quad\text{on }Z_{24}^{\mathcal D}.
\]
These substitutions have inverse $\lambda_T=\beta_T/r_T$.
Both parameters tend to infinity, since $\ell_T\geq T-C_E$.
The exact forms \eqref{wb:sphere-metric-A-slice-exact} and
\eqref{wb:sphere-metric-BC-slice-exact} identify the resulting
intrinsic metrics with $g_{Q_T}^{\oplus2}$ and $g_{C_T}^{\oplus4}$.
The whole-ball isometries and ambient convergence in that theorem
therefore give both displayed limits in their stated categories.

More explicitly, its factor coordinate maps become
\[
 (a,b,\xi)\longmapsto
 [r_T(1+a/T),e^{2\pi ib/T},[Y_T^{-1}\xi]]
\]
on $\mathcal A$, with inverse
$(r,u,x)\mapsto(T(r/r_T-1),Tu,Y_Tx)$ on the specified lift window.
On $\mathcal B$ and $\mathcal C$ they become
\[
 (a,b,\xi)\longmapsto
 [r_T(1+a/\ell_T),e^{2\pi ib/\ell_T},[\xi/\ell_T]],
\]
using the cone's paired representative in $\mathcal C$.
Their inverse is $(r,u,\theta)\mapsto
(\ell_T(r/r_T-1),\ell_Tu,\ell_T\theta)$ on that window.
Both inverse compositions are identities. The products of two or
four factors are exactly the general theorem's maps after the
displayed substitutions, so its all-order estimates and whole-ball
exhaustion apply to these parameter values.
The slice in each factor has $a=b=0$. Its four physical coordinates
therefore give the asserted coordinate four-plane, including the
ambient-restricted slice limit proved there.
\end{proof}

\begin{corollary}[Unit-volume flat quotients]
\hypertarget{tag-03DH}{}\label{wb:sphere-metric-unit-volume}
{\normalfont\scriptsize tag 03DH}\par
Write $D_T=\det Y_T>0$ on the native admissible radial range.
The four-volume of $Q_T^2$ is $D_T^2\sim T^4$, and its
unit-volume length factor is $D_T^{-1/2}\sim T^{-1}$.
Under the marked torus identifications, these unit-volume metrics
converge smoothly and in Gromov--Hausdorff topology to the square
unit four-torus.  They have no pointed $\mathbb R^4$ limit.
The four-volume of $C_T^4$ is $\ell_T^4$.  For every admissible
$T$, its unit-volume metric is isometric to the square unit
four-torus by
\[
 ([y_1],\ldots,[y_4])\longmapsto
 ([y_1/\ell_T],\ldots,[y_4/\ell_T])\in\mathbb R^4/\mathbb Z^4.
\]
\end{corollary}

\begin{proof}
The Euclidean fundamental parallelepiped of a period matrix $A$
has volume $|\det A|$.  Multiplication of a length metric by
$c>0$ multiplies four-volume by $c^4$.
For $A_T=\operatorname{diag}(Y_T,Y_T)$ this gives volume $D_T^2$
and the unique positive unit-volume scale $c_T=D_T^{-1/2}$.
Put $B_T=c_TA_T$.  The map
\[
 (Q_T^2,c_Td_{Q_T^2})\longrightarrow\mathbb R^4/B_T\mathbb Z^4,
 \qquad[y]\longmapsto[c_Ty]
\]
is an isometry with inverse $[z]\mapsto[c_T^{-1}z]$:
it carries every period $A_Tn$ to $B_Tn$, and multiplies every
term in the quotient-distance infimum by $c_T$.
The marked map $[a]\mapsto[B_Ta]$ from $\mathbb R^4/\mathbb Z^4$
has inverse $[z]\mapsto[B_T^{-1}z]$ and pulls back the metric to
\[
 G_T=D_T^{-1}\operatorname{diag}(Y_T^{\mathsf T}Y_T,
                                 Y_T^{\mathsf T}Y_T).
\]

Since $Y_T=TJ+E(r_T)$ with bounded $E(r_T)$ and
$J=\left(\begin{smallmatrix}0&1\\-1&0\end{smallmatrix}\right)$,
we have $D_T/T^2\to1$ and
$B_T\to R=\operatorname{diag}(J,J)$.
This proves both stated volume and scale asymptotics.
The matrix $R$ is orthogonal and $R\mathbb Z^4=\mathbb Z^4$;
the map $[z]\mapsto[R^{\mathsf T}z]$, inverse $[a]\mapsto[Ra]$,
identifies its quotient with the square unit four-torus.
Also $G_T\to I_4$ in every $C^k$ norm, because the tensors have
constant coefficients in the fixed marked coordinates.

The relation to \cref{wb:sphere-metric-scaling-trichotomy} is
the exact choice
\[
 \lambda_T=\frac{T}{r_T\sqrt{D_T}},\quad
 \beta_T=\frac{T}{\sqrt{D_T}}\longrightarrow1
 \quad\text{on }Z_{12},\qquad
 \lambda_T=\frac1{r_T},\quad\beta_T=1
 \quad\text{on }Z_{24}^{\mathcal D}.
\]
Indeed the slice factors $r_T/T$ and $r_T/\ell_T$ turn these
scales into $D_T^{-1/2}$ and $\ell_T^{-1}$ on the physical flat
tori. The inverse relation is $\lambda_T=\beta_T/r_T$.
Thus that theorem's finite-positive case proves the smooth and
Gromov--Hausdorff limit. For the quantitative diameter bounds below,
put $\eta_T=\|B_T-R\|\to0$. The inequalities
$(1-\eta_T)|v|\leq|B_Tv|\leq(1+\eta_T)|v|$, followed by
infima over $v=a-b+n$, give diameter at most $1+\eta_T$.

For $C_T^4$, the period matrix is $\ell_TI_4$, so its volume is
$\ell_T^4$ and the scale is $\ell_T^{-1}$.
The displayed division map carries its periods to $\mathbb Z^4$;
its inverse multiplies all four coordinates by $\ell_T$.
The quotient-distance formula proves exact isometry at every $T$.

Finally the first family has diameter at most $1+\eta_T<2$ for
large $T$, and the second has diameter one.  The Euclidean
points $0$ and $3e_1$ in the radius-four ball have distance three.
Every map of that ball into either family has distortion at
least three minus its diameter.  These lower bounds stay positive,
so the approximations required for a pointed $\mathbb R^4$ limit
cannot have errors tending to zero.
\end{proof}



\subsection{Vertex links and pointed limits}
\begin{proposition}[Vertex links and compact metric limits]
\hypertarget{tag-03DF}{}\label{wb:sphere-metric-tangent-scope}
{\normalfont\scriptsize tag 03DF}\par
Put \(L_3=T^3\), \(L_2=T^2\), and \(W=\mathcal M^4\).
The carriers have based homeomorphisms
\[
 Z_{12}\cong\operatorname{Susp}(L_3*L_3),\qquad
 Z_{24}^{\mathcal B}\cong\operatorname{Susp}(L_2^{*4}),\qquad
 Z_{24}^{\mathcal C}\cong W/\partial W.
\]
Their links at the chosen base point are respectively
\(L_3*L_3\), \(L_2^{*4}\), and \(\partial W\).
The topological dimensions of these links are \(7,11,11\).
A fixed piecewise Riemannian metric on each link gives the same
Hausdorff dimensions, and an exact conic length metric near the chosen
point gives vertex tangent cones of dimensions \(8,12,12\).

The complete integral local homology of the third carrier is
\begin{equation}
 H_i\bigl(Z_{24}^{\mathcal C},
       Z_{24}^{\mathcal C}\setminus\{*\};\mathbb Z\bigr)
 \cong
 \begin{cases}
  \mathbb Z^{\binom4{i-1}},&2\leq i\leq5,\\
  0,&i=6,7,\\
  \mathbb Z^{\binom4{i-8}},&8\leq i\leq12,\\
  0,&\text{otherwise}.
 \end{cases}
 \hypertarget{tag-03D1}{}\label{wb:sphere-metric-C-smash-local-homology}
\end{equation}
For the first two carriers the highest nonzero local homology groups
are \(\mathbb Z\) in degrees \(8\) and \(12\), respectively.
These groups obstruct a pointed homeomorphism from any of the three
vertex models to \((\mathbb R^4,0)\).

On each carrier, the conic tensors \(dr^2+r^2h\) and
\(dr^2+4r^2h\) extend to compatible length metrics.
The identity on the carrier is a homeomorphism between their metric
topologies.  For a link of dimension \(d\), the unit balls of their
vertex tangent cones have Hausdorff volumes in ratio \(2^d\).
Consequently these marked tangent cones have no pointed isometry.

Every blow-down of a fixed compact metric target has one-point limit.
For the fixed metrics and moving points of
\cref{wb:sphere-metric-fixed-realization}, the chosen vertex is at
scaled distance \(T\) in the first family and \(\ell_T\) in the
other two.  Both distances tend to infinity.  The limits there are
moving-base-point blow-ups, with their intrinsic \(\mathbb R^4\)
slices and ambient \(\mathbb R^8\) or \(\mathbb R^{12}\) limits
as stated in that corollary.
For these fixed quotient metrics,
\cref{wb:sphere-metric-scaling-trichotomy} excludes every proper
pointed Gromov--Hausdorff tangent at the chosen vertex.
\end{proposition}

\begin{proof}
For compact Hausdorff spaces \(V_i\) and closed nonempty subspaces
\(A_i\), the product of the quotient maps induces
\[
 \frac{\prod_iV_i}
 {\bigcup_i V_1\times\cdots\times A_i\times\cdots\times V_m}
 \longrightarrow\bigwedge_i(V_i/A_i).
\]
It maps the collapsed union to the smash base point and maps every
remaining tuple to the tuple of its quotient classes.  The inverse
away from the base point takes each class to its unique representative
in \(V_i\setminus A_i\).  The map is therefore a continuous bijection
from a compact space to a Hausdorff space, hence a homeomorphism.

For \(\mathcal A\) and \(\mathcal B\), delete the chosen inner
suspension pole.  The coordinate change
\[
 \rho=\frac{\epsilon-r}{r},\qquad
 r=\frac{\epsilon}{1+\rho}
\]
gives a homeomorphism to the infinite cone on \(L_3\) or \(L_2\).
The outer suspension pole becomes its cone vertex \(\rho=0\).
The complement of the smash base point is the product of these
infinite cones.  The homeomorphism from that product to the infinite
cone on the join is
\[
 R=\Bigl(\sum_i\rho_i^2\Bigr)^{1/2},\qquad
 \lambda_i=\frac{\rho_i^2}{R^2},\qquad
 \rho_i=R\sqrt{\lambda_i}.
\]
At \(R=0\) all link coordinates are discarded; at
\(\lambda_i=0\) the \(i\)-th link coordinate is discarded.
Thus these formulas respect precisely the cone and join equivalence
relations.  They are inverse continuous maps, including the vertices.
Their one-point compactifications, with suspension coordinate
\(\tau=1/(1+R)\), give the asserted based suspension models;
the chosen base point is \(\tau=0\).

Let \(L=\partial\mathcal M\), and choose a collar
\(c:L\times[0,1]\hookrightarrow\mathcal M\) with
\(c(y,0)=y\).  On the attached cone write \(t=r/\epsilon\),
whose inverse is \(r=\epsilon t\).  Define a map from
\(\mathcal C\) to \(\mathcal M/L\) by
\[
 [y,t]\longmapsto[c(y,t/2)],\qquad
 c(y,s)\longmapsto[c(y,(s+1)/2)],
\]
and by the identity on the complement of the open collar.
The two formulas agree on the glued boundary \(t=1,s=0\),
and the latter formula agrees with the identity at \(s=1\).
The cone vertex maps to \([L]\).  For a collar coordinate
\(u\in(0,1]\), its inverse is the cone point \([y,2u]\)
for \(u\leq1/2\), and the core collar point \(c(y,2u-1)\)
for \(u\geq1/2\).  These inverses agree at \(u=1/2\).
Compactness gives the based homeomorphism
\begin{equation}
 (\mathcal C,c_0)\xrightarrow{\ \cong\ }
 (\mathcal M/\partial\mathcal M,[\partial\mathcal M]).
 \hypertarget{tag-03D0}{}\label{wb:sphere-metric-C-boundary-quotient-homeomorphism}
\end{equation}
Applying the product quotient map gives
\[
 (\mathcal C^{\wedge4},*)\cong
 \left(\mathcal M^4\Big/
 \bigcup_{i=1}^4\mathcal M^{i-1}\times\partial\mathcal M
                   \times\mathcal M^{4-i},[\partial W]\right)
 =(W/\partial W,[\partial W]).
\]
On the complement of the collapsed set the inverse is the ordered
product of the preceding inverse maps.

We specify a boundary collar for this product.  Set
\({d_{\rm col}}(c(y,u))=u\) on the collar and \({d_{\rm col}}=1\) on its complement.
For \(0\leq t<1\), let
\[
 H_t(c(y,u))=c(y,t+(1-t)u)
\]
and let \(H_t\) be the identity outside the open collar.
It maps \(\mathcal M\) homeomorphically onto \(\{{d_{\rm col}}\geq t\}\).
The map
\[
 \partial W\times[0,1)\longrightarrow
 \{(x_1,\ldots,x_4):\min_i {d_{\rm col}}(x_i)<1\},\qquad
 (y,t)\longmapsto H_t^{\times4}(y)
\]
is a collar: for a given tuple set \(t=\min_i {d_{\rm col}}(x_i)\)
and apply \(H_t^{-1}\) in each factor, using
\(u=(u'-t)/(1-t)\) in the collar.  At least one resulting
factor lies in \(\partial\mathcal M\).

At a point with \(k\) boundary factors, a corner chart of
\(\partial W\) is
\[
 \{y\in[0,\infty)^k:\min_i y_i=0\}\times\mathbb R^{12-k}.
\]
The first factor is homeomorphic to the hyperplane
\(\sum_i v_i=0\) by the mutually inverse maps
\[
 v\longmapsto v-\min_i(v_i)\mathbf1,\qquad
 y\longmapsto y-\frac1k\sum_i y_i\mathbf1.
\]
Thus \(\partial W\) is a closed eleven-dimensional manifold.
The product orientation induces its boundary orientation.
The collar descends after collapse to a cone neighborhood of
\([\partial W]\).  Excision and the boundary map of the cone pair give
\begin{equation}
 H_i(W/\partial W,(W/\partial W)\setminus\{*\};\mathbb Z)
 \xrightarrow{\ \cong\ }\widetilde H_{i-1}(\partial W;\mathbb Z).
 \hypertarget{tag-03D5}{}\label{wb:sphere-metric-boundary-link-local-shift}
\end{equation}
Explicitly, the punctured cone retracts by radial projection to a
positive-radius copy of the link.  The displayed map is the relative
boundary followed by this projection.  Its inverse sends a reduced
link cycle to its cone chain.  For degree zero one uses the reduced
cone sequence; the connected links give zero there.

By \cref{wb:cusp-cone-S24-route}, the integral homology of
\(\mathcal M\) is \(\mathbb Z\) in degrees zero and one and
zero in every other degree.  The integral cross product therefore
gives
\[
 H_j(W;\mathbb Z)=\mathbb Z^{\binom4j}\quad(0\leq j\leq4),
 \qquad H_j(W;\mathbb Z)=0\quad(j>4),
\]
by the K\"unneth theorem
\cite[Theorem~3B.6, p.~275]{Hatcher2002}.
Let \(\gamma\) generate \(H_1(\mathcal M;\mathbb Z)\),
and let \(\alpha\in H^1(\mathcal M;\mathbb Z)\) satisfy
\(\langle\alpha,\gamma\rangle=1\).
For each subset \(S\subset\{1,2,3,4\}\), the ordered cross product
\(\gamma_S\), with a point in every other factor, and the ordered
cup product \(\alpha_S=\prod_{i\in S}\operatorname{pr}_i^*\alpha\) are dual bases.  This identifies the
cohomology cross product of
\cite[Theorem~3.15, p.~216]{Hatcher2002} explicitly.
Cap product with the relative fundamental class gives
\[
 \beta_S=\alpha_S\frown[W,\partial W]
 \in H_{12-|S|}(W,\partial W;\mathbb Z).
\]
These classes form a basis by Poincar\'e--Lefschetz duality
\cite[Theorem~3.43, p.~254]{Hatcher2002}, applied with the
boundary pieces \(A=\varnothing\), \(B=\partial W\).
The product corner charts and collar above give the required
compact oriented manifold with boundary.  Consequently
\[
 H_j(W,\partial W;\mathbb Z)=
 \begin{cases}
  \mathbb Z^{\binom4{j-8}},&8\leq j\leq12,\\
  0,&\text{otherwise}.
 \end{cases}
\]
In the long exact sequence of this pair, inclusion gives
\(H_j(\partial W)\xrightarrow{\cong}H_j(W)\) for
\(0\leq j\leq4\); the boundary groups in degrees five and six
vanish; and the relative boundary gives
\[
 H_{j+1}(W,\partial W)\xrightarrow{\ \cong\ }H_j(\partial W)
 \quad(7\leq j\leq11).
\]
These maps give the low-degree generators from \(\gamma_S\)
and the high-degree generators \(\partial\beta_S\).
Together with \eqref{wb:sphere-metric-boundary-link-local-shift}
they give \eqref{wb:sphere-metric-C-smash-local-homology}.
In the high range the indexing bijection is
\(S\mapsto\{1,2,3,4\}\setminus S\): the sizes \(12-i\)
and \(i-8\) sum to four, and taking complements twice is the identity.

For the other two links, the augmented cellular complex of a join
is the tensor product of the augmented cellular complexes of its
factors, with the degree shifted by one for each join.
In detail, an ordered join of cells of degrees \(p_i\geq-1\)
has degree \(\sum_i(p_i+1)-1\), and its boundary differentiates
the \(i\)-th factor with sign
\((-1)^{\sum_{j<i}(p_j+1)}\).
The torus has one cell for each subset of its circle factors,
and its cellular differentials vanish in positive degrees.
In its augmented complex the degree-zero cell maps by the identity
to the augmentation generator in degree minus one.  This two-term
summand contracts by the inverse identity; the positive cells give
the reduced complex.  Tensor terms containing that contractible
summand contract with the displayed tensor sign.  Hence the
remaining join complex has free generators from positive-degree
cells in every torus factor and zero differential.
After the additional local-homology shift, the rank polynomials in the formal variable \(z\) are
\[
 z^2(3z+3z^2+z^3)^2
  =9z^4+18z^5+15z^6+6z^7+z^8,
\]
\[
 z^4(2z+z^2)^4
  =16z^8+32z^9+24z^{10}+8z^{11}+z^{12}.
\]
In particular, the top groups are the claimed copies of \(\mathbb Z\).
The identity of pairs from a cone with its positive-radius link to
the cone with its vertex deleted induces the same groups: the
inclusion of that link in the punctured cone has radial retraction
and a radial homotopy to the identity.
For \(\mathbb R^4\), radial retraction of the punctured space to
\(S^3\) gives local homology \(\mathbb Z\) in degree four and
zero in degrees eight and twelve.  A pointed homeomorphism would
induce an isomorphism in those degrees, where the displayed
\(\mathbb Z\) would have to map isomorphically to zero.

The join cell dimensions give \(3+3+1=7\) and \(4\cdot2+3=11\).
Their top-dimensional open strata have these dimensions, and the
finite triangulations give the matching upper bounds.
The boundary charts above give dimension eleven for the third link.
On each top-dimensional simplex a piecewise Riemannian metric is
locally comparable with a Euclidean metric.  Lower-dimensional
strata have the dimensions of their simplices.  Thus the Hausdorff
dimensions are \(7,11,11\), and those of the cones are \(8,12,12\).

Equip a cone collar with \(dr^2+r^2h\) and extend it over the
compact carrier by a piecewise Riemannian length metric.  For a
suspension one may use \(dt^2+\min(t,1-t)^2h\) on the whole
suspension.  For \(W/\partial W\), the boundary collar allows
extension over the compact core, interpolating positive tensors
on an overlap away from the vertex.  Replacing \(h\) by \(4h\)
gives a second metric through the same construction.  Both induce
the given topology: on the punctured collar this follows from the
positive tensors, and distance to the vertex is the collar radius.

In an exact cone, dilation \((r,l)\mapsto(\lambda r,l)\)
has inverse \((r,l)\mapsto(r/\lambda,l)\) and multiplies all
lengths by \(\lambda\).  To compare the compact carrier with
the complete cone, fix a conic collar of radius \(\delta\).
For points of radius at most \(R/\lambda\), every path leaving
that collar has length at least \(2(\delta-R/\lambda)\),
whereas the path through the vertex has length at most
\(2R/\lambda\).  For \(\lambda>2R/\delta\), such an exit
cannot shorten any distance in this ball.  Dilation therefore
gives exact isometries of every fixed-radius ball for all
sufficiently large \(\lambda\).  The vertex tangent is the
complete proper cone.

For the chosen link \(L_{\rm link}\in\{L_3*L_3,L_2^{*4},\partial W\}\), let \(d=\dim L_{\rm link}\) and let \(V_h>0\) be the \(d\)-dimensional
volume of the link.  On each top-dimensional stratum the cone
volume density is \(r^d\,dr\,dV_h\).  The lower strata have
zero \((d+1)\)-dimensional measure.  Thus the unit cone ball
has volume \(V_h/(d+1)\).  For the tensor \(4h\), the angular
volume density is multiplied by \(2^d\), while the radial
distance and the unit-ball radius remain unchanged.
Pointed isometries preserve balls and their Hausdorff measures,
so the two cones admit no pointed isometry.  The identity of
their underlying cones supplies the stated homeomorphism.

Finally, for any fixed compact metric space \((Z,D)\) and
\(R_n\to\infty\),
\[
 \operatorname{diam}(Z,R_n^{-1}D)
   =R_n^{-1}\operatorname{diam}(Z,D)\longrightarrow0.
\]
The map to a point and the section choosing the marked point
have distortion at most this diameter, which proves the pointed
Gromov--Hausdorff limit.  Every pair of sequences also has
rescaled distance bounded by this diameter, so its metric
ultralimit is the same point.
For the metrics in \cref{wb:sphere-metric-fixed-realization},
the quotient distance from the moving point to the vertex is
exactly \(r_T\).  Multiplication by \(T/r_T\) or
\(\ell_T/r_T\) gives \(T\) or \(\ell_T\), respectively.
The marked-period estimate \(\ell_T\geq T-C_E\) makes both
distances divergent.  The exact charts in that corollary place
the slices in coordinate four-planes in the ambient Euclidean
limits.  The chosen vertex leaves every bounded ball about the
moving point, while the compact blow-down maps above have
one-point image.
For a blow-up based at the fixed vertex, the unit-ball configurations
in \cref{wb:sphere-metric-scaling-trichotomy} have mutual distance
two and unbounded cardinality.  The finite covering numbers of a
compact ball in a proper pointed limit exclude such convergence.
This obstruction applies to all three fixed quotient metrics.
\end{proof}

\subsection{Source pullbacks and moving sphere metrics}
\begin{remark}[Quotient differentials and moving quotient metrics]
\hypertarget{tag-Q30Y}{}\label{mapsIII:rem:source-and-moving-metrics}
{\normalfont\scriptsize tag Q30Y}\par
On the regular cusp chart, write a fibre vector as $P_Ta$ in the
ordered basis $(\widehat\gamma,\widehat u,\widehat w,\widehat\delta)$.
The two-torus quotient sends it to $Y_T(a_\gamma,a_u)$, and the
circle quotient sends it to $\ell_Ta_\gamma$. Their differential
kernels are exactly
\[
 P_T\operatorname{span}_{\mathbb R}\{\widehat w,\widehat\delta\},
 \qquad
 P_T\operatorname{span}_{\mathbb R}
       \{\widehat u,\widehat w,\widehat\delta\},
\]
because $Y_T$ is invertible and $\ell_T>0$. The base coordinates
are retained, so these are also the kernels on the full regular
charts. The product maps have the direct sums of these kernels,
of dimensions four and twelve. The smash quotient is a local
diffeomorphism off its fat wedge. Thus for every such nonzero
kernel vector $v$ the pulled-back target form has
$q^*g(v,v)=g(dq(v),dq(v))=0$. A smooth Riemannian metric on the
source sphere has positive value on that same $v$. For an orthonormal
frame $(e_i)$, the tangent isometry to $\mathbb R^{12}$ or
$\mathbb R^{24}$ is $v\mapsto(\langle v,e_i\rangle)_i$, with
inverse $(c_i)\mapsto\sum_i c_ie_i$.

For the moving sphere carriers, the maps
\[
 \widehat p_2[r,z,a]=[r,c(a)],\qquad
 \widehat p_1[r,z,\theta]=[r,\theta]
\]
of \cref{wb:cusp-cone-angle-forgotten-sphere-carriers} relate the
conical carriers to the angle-forgotten sphere targets. Here $c$
collapses the marked one-skeleton in $\overline T^2$.
On the radial slices their products are exactly the quotient maps
from $Q_T^2$ and $C_T^4$ defining the metrics on
$\mathbb S_{12,T}^4$ and $\mathbb S_{24,T}^4$ in that theorem.
On those radial slices, the inverse off the collapsed sets takes
the unique marked tuple.
The origins $p_T$ map to the collapsed points. The marked midpoint
classes are
\[
 \zeta_{12,T}=[c(1/2,1/2),c(1/2,1/2)],\qquad
 \zeta_{24,T}=[\ell_T/2,\ell_T/2,\ell_T/2,\ell_T/2].
\]
Their quotient distances to the collapsed points are at least
$(T-C_E)/2$ and exactly $\ell_T/2$, respectively, by the
closed-set distance formula of
\cref{mapsIII:lem:collapsed-subset-distance} and the period estimates in
\cref{wb:cusp-cone-angle-forgotten-sphere-carriers}.
These distances tend to infinity. For any fixed compatible metric
on the compact ambient $S^6$ or $S^8$, the corresponding distances
are bounded by that sphere's diameter. The identity on each moving
sphere is therefore a homeomorphism between these two metric
topologies, and the displayed pair of points proves that it is
not an isometry for every sufficiently large $T$.
\end{remark}

\subsection{Euclidean cells, intrinsic limits and boundary fibres}
\begin{corollary}[Three Euclidean realizations and their exact scope]
\hypertarget{tag-037J}{}\label{wb:cusp-cone-sharp-global-R4-scope}
{\normalfont\scriptsize tag 037J}\par
The two native routes realize $\mathbb R^4$ in the following ways:
\begin{enumerate}
\item for each admissible $T$, as the nonclosed topologically
embedded open cells $U_{12,T}\subset S^6$ and $U_{24,T}\subset S^8$;
\item as the pointed intrinsic Euclidean limit of the four-torus
slices in the global quotient carriers of $S^{12}$ and $S^{24}$;
\item as the Euclidean boundary fibre at $s=0$ of
$\mathfrak E_{12}$ and $\mathfrak E_{24}$.
\end{enumerate}
The cells in the first item are open in their moving four-spheres.
No whole compact sphere or whole compact conical target in these
routes is homeomorphic to $\mathbb R^4$, and no continuous map from
any of those compact spaces to $\mathbb R^4$ is surjective.
In particular $\mathbb R^4$ is not a continuous quotient target
of any of them.  For every fixed compatible ambient target metric,
the unrescaled restricted compact slices converge to a singleton.
Their slice limit is therefore not Euclidean $\mathbb R^4$.
At each fixed $T$, the restricted open-cell metrics admit no
distance-preserving map of all Euclidean $\mathbb R^4$ into the cell.
\end{corollary}

\begin{proof}
For the first realization use the homeomorphisms in
\cref{wb:cusp-cone-literal-R4-open-cells}. Explicitly, set
\[
 u_T(a)=[Y_Ta],\quad v_T(a)=[\ell_Ta],\quad
 \eta(t)=\tfrac12(1+t/\sqrt{1+t^2}).
\]
The maps are $h_{12,T}(u_T\eta^2,u_T\eta^2)$ and
$h_{24,T}(v_T\eta,v_T\eta,v_T\eta,v_T\eta)$, with the arguments
split into the indicated ordered pairs or single coordinates.
Their inverse takes the unique uncollapsed tuple, applies the
representative maps $Y_T^{-1}$ or division by $\ell_T$ in $(0,1)$,
and applies $\eta^{-1}(a)=(2a-1)/(2\sqrt{a(1-a)})$.
The cited theorem proves both inverse identities, the exact
synchronized cube coordinates, and the open-in-the-moving-sphere
and nonclosed-in-the-ambient-sphere assertions for these maps.

For the second realization put
\[
 A_T^{(12)}=\operatorname{diag}(Y_T,Y_T),\qquad
 A_T^{(24)}=\ell_TI_4,\qquad
 P_{\nu,T}=\mathbb R^4/A_T^{(\nu)}\mathbb Z^4.
\]
The maps splitting two coordinate pairs or four coordinates
identify $P_{12,T}$ with $Q_T^2$ and $P_{24,T}$ with $C_T^4$;
their inverse concatenates the coordinates. They preserve the
quotient-distance formula
$d_T([x],[y])=\min_{n\in\mathbb Z^4}|x-y+A_T^{(\nu)}n|$.
The exact maps in \cref{wb:cusp-cone-explicit-carriers} are
\[
 k_{\nu,T}[x]=[r_T,c_\nu^{-1}([(A_T^{(\nu)})^{-1}x])],
 \qquad j_{\nu,T}=\delta_\nu k_{\nu,T}.
\]
Their inverse at positive height sends a lift $b$ of the marked
angle to $[A_T^{(\nu)}b]$; changing $b$ by an integer changes its
image by a period. The commutative sphere-fibre restrictions in
that theorem identify these maps with the target slices.
The complete whole-ball isometry proof in
\cref{wb:cusp-cone-S12-route,wb:cusp-cone-S24-route}, transported
by these maps and inverses, supplies the intrinsic Euclidean
limits. Their diameters diverge by
\cref{wb:cusp-cone-intrinsic-expansion}.
For the fixed metrics of \cref{mapsIII:lem:fixed-conical-metrics},
the exact first forms give
\[
 j_{12,T}^*d_T^{\mathrm{slice}}=(r_T/T)d_T,\qquad
 j_{24,T}^*d_T^{\mathrm{slice}}=(r_T/\ell_T)d_T.
\]
The corresponding unrescaled path-metric diameters are at most
$r_T(1+C_E/T)$ and $r_T$. Inclusion bounds the restricted ambient
distance by the intrinsic path distance.

For the third realization use $e_\nu(x)=[0,x]$ and its
unique-coordinate inverse from
\cref{wb:cusp-cone-exact-decompactification-space}.
That theorem proves that the zero fibre is a closed Euclidean
boundary fibre. Its exact comparison map is
\[
 \beta_\nu[s,x]=
 \begin{cases}
 [e^{-2\pi/s},c_\nu^{-1}([(A_{1/s}^{(\nu)})^{-1}x])],&s>0,\\
 d_\nu,&s=0.
 \end{cases}
\]
The positive-height inverse sends $(r,a)$ to
$[s,A_{1/s}^{(\nu)}b]$, where $s=-2\pi/\log r$ and $b$ is any
lift of $c_\nu(a)$. Both compositions, representative independence
and continuity at zero are proved there. Thus
$\delta_\nu\beta_\nu e_\nu$ is the constant vertex map, while
the positive-fibre restrictions are precisely the slices above.
The closed noncompact zero fibre also proves that
$\mathfrak E_\nu$ is noncompact.

The whole sphere and conical spaces in these routes are compact
through their defining quotient maps: the cube-sphere maps, the
two-end cylinder quotients defining $\mathcal A,\mathcal B$, the
attachment of a compact cone to the compact core of $\mathcal C$,
the finite products and their smash quotients, and
$[0,r_0]\times T^4\to\mathcal D_\nu$ all have compact domains.
For any such space $K$ and continuous $F:K\to\mathbb R^4$, the
continuous function $|F|$ has a finite maximum $M$.  Its image
omits $(M+1,0,0,0)$, so $F$ cannot be surjective.  This excludes
both a quotient onto $\mathbb R^4$ and the continuous surjective
direction of any proposed homeomorphism with it.

The local cone comparison is the map
$[r,a]\mapsto[r/r_0,c_\nu(a)]$, inverse
$[s,b]\mapsto[r_0s,c_\nu^{-1}(b)]$, of
\cref{wb:cusp-cone-explicit-carriers}. Its complete pair-homology
calculation in \eqref{wb:cusp-cone-carrier-local-homology} applies
to these same vertices, with the coned-cell boundary map as its
specified isomorphism. It gives the stated local chart obstruction.

Finally let $D_0$ be the fixed compatible metric of
\cref{mapsIII:lem:fixed-conical-metrics} on a conical smash target.
The synchronized torus slices satisfy $D_0(z,*)=r_T$ exactly:
the product distance to the fat wedge is the minimum of the
factor distances to their inward vertices, each of which is $r_T$.
For any fixed compatible metric $D$ on that same compact target,
the identity is a homeomorphism.  Define its modulus
\[
 \omega_D(t)=\sup\{D(x,y):D_0(x,y)\leq t\}.
\]
It is finite and tends to zero with $t$.  Indeed a contradicting
sequence of pairs with $D_0(x_n,y_n)\to0$ has a convergent
subsequence whose two components have the same limit; compatibility
then forces $D(x_n,y_n)\to0$.  Consequently
\[
 \sup_{z\in\Sigma_T}D(z,*)\leq\omega_D(r_T),\qquad
 \operatorname{diam}_D\Sigma_T\leq2\omega_D(r_T)\longrightarrow0.
\]
The corresponding sphere-target assertion follows from the cube
quotient metric $d_0([x],*)=\operatorname{dist}(x,\partial I^n)$:
the displayed synchronized coordinates give distance at most
$r_T/\epsilon$.  The identity modulus for any fixed compatible
sphere metric thus bounds the entire moving four-sphere's diameter
by $2\omega_D(r_T/\epsilon)\to0$.

These compact slices converge in ambient Hausdorff distance to
the vertex.  The correspondence sending every slice point to one
point has distortion equal to the slice diameter and respects any
chosen slice basepoint, so the unrescaled pointed restricted limit
is a singleton.  A Euclidean radius-two ball contains $0,e_1$ at
distance one; no correspondences of vanishing distortion from
spaces of vanishing diameter can approximate that pair.  Hence
this restricted limit cannot be $\mathbb R^4$.
For each fixed $T$, the open cell's restricted distance has the
finite diameter bound of its compact ambient sphere.  A
distance-preserving map from all $\mathbb R^4$ would send $0,ne_1$
to points at distance $n$, exceeding that bound for large $n$.
This proves the last assertion for every compatible ambient metric.
\end{proof}
\part*{Open problems}
\phantomsection
\addcontentsline{toc}{part}{Open problems}
\section{Further classification problems}
\subsection{Section lines and sphere multiplicities}
\begin{remark}[Natural sphere data and further classifications]\hypertarget{tag-Q301}{}\label{maps:rem:natural-sphere-sector-scope}
{\normalfont\scriptsize tag Q301}\par
The sphere table \eqref{maps:eq:D.ST13} concerns the natural line
$M_0=\mathcal O(P-O)$. In the vertically corrected family
$M_V=\mathcal O(P-O+V)$ of \ref{maps:src:D.2.22}, the specialization
$V=0$ is identity on this line bundle, its linearization and the
marked quotient with the same filling data. The finite reduction
for integral
vertical $V$ has $226$ marked rows, with the solution lattices
specified in that theorem.

The nontrivial fundamental-group obstruction at fixed-character
index $g>1$ in \ref{maps:src:D.2.30} preserves the sphere
classification of Corollary~\ref{maps:src:D.2.28} in the full
finite-lattice family.

The unrestricted finite-log homology is given in
\ref{maps:src:D.2.29}--\ref{maps:src:D.2.33}. The specified minimal
simply connected smooth sector is classified by
\ref{maps:src:D.2.35}--\ref{maps:src:D.2.36},
\ref{maps:src:D.2.38}--\ref{maps:src:D.2.39},
\ref{maps:src:D.2.42} and
\ref{maps:src:D.2.44}--\ref{maps:src:D.2.45}.
The remaining weighted and additional-regular-log smooth
classifications are included in Problem~\ref{maps:src:D.2.4}.
\end{remark}

\begin{problem}[Degree-zero lines and sphere multiplicities]
\hypertarget{tag-Q30Z}{}\label{mapsIII:problem:degree-zero-lines-and-multiplicities}
{\normalfont\scriptsize tag Q30Z}\par
In the elliptic-surface logarithmic construction, allow the arbitrary
fibrewise degree-zero line bundles of
\cite[Remark~8.5, p.~24]{E26}. Complete the integral cap-map and
filling classification for the remaining choices, and determine
which initial configurations admit infinitely many sphere-producing
multiplicities, as asked in \cite[Question~4, p.~2]{manViaclovskyLogs26}.

The natural and vertically corrected section-line families are the
specializations $M=\mathcal O(P-O)$ and
$M=\mathcal O(P-O+V)$ of \cref{maps:src:D.2.22}.
With the marked section pair, divisor, linearization and filling data
retained, the comparison on either specialization is the identity on
the line bundle and its quotient construction.
The finite initial-data reduction in that theorem and the integral
algorithms and smooth sectors cited in
\cref{maps:rem:natural-sphere-sector-scope,maps:src:D.2.4}
give the established reductions for those specializations and their
stated filling families.
\end{problem}

\subsection{Cusp compactifications and quotient classes}
\begin{problem}[Smooth compactifications of finite cusp covers]
\hypertarget{tag-Q30U}{}\label{mapsIII:problem:finite-cusp-cover-compactification}
{\normalfont\scriptsize tag Q30U}\par
For an integer $M>1$, consider the local degree-$M^2$ fibre cover
with lattice $\Lambda_M$ and base change $t=s^M$ in
\cref{wb:cusp-cone-obstructions}. Determine whether it admits a smooth
compactification after a connected global base cover with additional
branch values away from $p_0$, and determine the topology of the
resulting total space, including whether it can be a sphere.
The construction must specify the global base cover, the normalization
and resolution of every pulled-back singular fibre, and the lifted
metric. Its restriction to the chosen cusp must be the stated local
cover with that metric.

The finite flat normal completion in \cref{maps:src:C.1.5} gives the
established global covering construction. The one-branch-value
obstruction and the identity endpoint $M=1$ are proved in
\cref{wb:cusp-cone-obstructions}; the question concerns the smooth
completion with the additional branch values just specified.
\end{problem}

\begin{problem}[Unstable classes of the conical quotient maps]
\hypertarget{tag-Q30V}{}\label{mapsIII:problem:raw-conical-quotient-classes}
{\normalfont\scriptsize tag Q30V}\par
Use the fixed cusp disc, marked quotients, original finite fillings
and basepoint $x_0\in W$ of
\cref{wb:cusp-cone-S12-route,wb:cusp-cone-S24-route}.
For the based source homeomorphisms $h_2:X^{\wedge2}\to S^{12}$
and $h_4:X^{\wedge4}\to S^{24}$ of
\cref{wb:cusp-cone-optimal-exponents}, determine
\[
 [q_2^{\wedge2}h_2^{-1}]
       \in\pi_{12}(\mathcal A^{\wedge2}),\qquad
 [q_1^{\wedge4}h_4^{-1}]
       \in\pi_{24}(\mathcal C^{\wedge4}).
\]
The target homotopy equivalences are
\eqref{wb:cusp-cone-A-smash-wedge} and
\eqref{wb:cusp-cone-C-smash-wedge}.

The target maps $p_A=\widehat p_2$ and
$p_C=\widehat p_1p_{\mathcal C\mathcal B}$ of
\cref{wb:cusp-cone-optimal-exponents,wb:cusp-cone-angle-forgotten-sphere-carriers}
satisfy $p_Aq_2=\widehat q_3$ and
$p_Cq_1=\widehat q_\gamma$ on all of $X$.
Their smash powers send the displayed classes to the zero classes
computed in \cref{wb:cusp-cone-sphere-map-homotopy-status}.
The zero and nonzero surjective representatives in
\cref{wb:cusp-cone-target-homotopy-types} prove that the listed
ordinary invariants and homotopy-section obstruction do not
determine the two classes asked for here.
\end{problem}

\subsection{Remaining cap maps and smooth types}
\begin{problem}[Remaining filling and smooth classification]\hypertarget{tag-02LG}{}\label{maps:src:D.2.4}
{\normalfont\scriptsize tag 02LG}\par
Determine the marked integral cap maps and oriented smooth types in
the following remaining sectors of the section-line construction:
\begin{enumerate}[label=(\roman*)]
\item logarithmic star caps with nonprimitive long circle, and the
other logarithmic star data beyond the canonical primitive-long
family of \ref{maps:src:D.2.49};
\item original multiplicative caps for component-corrected line
bundles beyond the uncorrected relative model of
\ref{maps:src:D.2.46};
\item the individual mixed elliptic--$\gamma$ and second
fixed-normal torsion columns of invariant-Chern original additive
caps with $h>1$ in \ref{maps:src:D.2.51};
\item the complete smooth invariants of the remaining weighted
fillings and fillings with additional regular logarithms.
\end{enumerate}

The natural catalogue has $172$ initial section/zero rows on $21$
profiles. The original and logarithmic additive choices and the
original and Mumford multiplicative choices give $495$ integral
presentation patterns in the scope of
\ref{maps:src:D.2.29}--\ref{maps:src:D.2.33}. This count includes
the permitted finite-good choices at $III^*$ and $IV^*$; the
$I_n^*$ caps are original. The original-star subset in
\ref{maps:src:D.2.41} consists of the $459$ patterns that keep
every starred cap original. The natural sphere sector consists
of $59$ marked patterns on $50$ initial rows, with the logarithmic
orders and affine parameters of
\ref{maps:src:D.2.24}--\ref{maps:src:D.2.28}.

For the original-star sector,
\ref{maps:src:D.2.37}, \ref{maps:src:D.2.40} and
\ref{maps:src:D.2.41} give the integral algorithm for every
positive translation iterate, including shared even factors,
the permitted free finite logarithms, regular logarithms and
the ordinary clutch. The full nonsaturated triangular-cusp
boundary kernel is computed in \ref{maps:src:D.2.43}.
For positive section pairs for which at least one section is
narrow at each singular fibre, \ref{maps:src:D.2.46} gives the integral algorithm
for every effective whole-fibre weight distribution on
multiplicative or regular fibres, with original star caps.
Its fundamental-group presentation includes the OS49/II
Klein-four vertical subgroup.

The minimal simply connected sector with primitive Mumford
multiplicative caps and no additional regular logarithms has
$162$ patterns. Its smooth classification includes all permitted
finite logarithmic orders, affine twists and simply connected
ordinary clutches. In the $18$ patterns of residual orders
$4,8,20$, the refined class depends on the marked zero position
as in \eqref{maps:eq:D44.1} and \eqref{maps:eq:D45.3}.
In the primitive fibre coordinates of those tables, the values
and models occurring for $\varepsilon=0,1$ are
\[
\begin{array}{c|c|c}
 c&P&\text{oriented smooth type}\\\hline
 4&2\varepsilon&Z_{4,\varepsilon}\\
 8&4\varepsilon&Z_{8,\varepsilon}\\
 20&4+10\varepsilon&Z_{4,\varepsilon}\# Z_{5,1}.
\end{array}
\]
Their complete fundamental-class invariants are zero.
The relative spin calculation uses the rational linking of
the critical cycle and the primitive finite-$II$ core class
of \ref{maps:src:D.2.45}.

The simply connected weighted Engel family is classified for
every positive weight in \ref{maps:src:D.2.48}, including its
complete finite-coefficient fundamental-class invariant.
The canonical primitive-long odd-star logarithmic caps and
their marked integral boundary kernels are computed in
\ref{maps:src:D.2.49}. Component corrections to the initial
line exist exactly when the height product is integral;
\ref{maps:src:D.2.50} constructs sixteen simply connected
fillings from those corrections.

For every invariant-Chern original additive cap,
\ref{maps:src:D.2.51} computes the integral cellular cone,
homology and full boundary image, including the orbit extension.
Its complete marked columns apply to nonzero corrected
multidegree with $h=1$. The unchanged zero-degree caps use
their specified extending frames. With these original
additive choices and the other filling choices in its
statement, \ref{maps:src:D.2.52} gives the integral global
algorithm for every positive initial pair of integral height.

The weighted OS45/default member with $P=2G$, $Q=8G$,
weight $2$ at $I_8$, all other weights zero, all Mumford
caps and ordinary clutch $e=1$ has
\[
 H_2=\mathbb Z^{15}\oplus\mathbb Z/8,\qquad
 H_3=\mathbb Z/8,\qquad \chi=32.
\]
Its primitive integral cubic, full fundamental-class invariant
and refined class are computed in \ref{maps:src:D.2.53}.
The component lattice has index two, and
\[
 c_1=c_2=\lambda=0,\qquad P_{\mathrm{ref}}=0.
\]
The source-line circle bundle gives a primitive embedded torsion
sphere. For the differential reference it has parity
$\varepsilon_\rho=1$, and its linking with the complete critical
nullchain is $-1$. These two values give zero in the refined
linking formula.
\end{problem}


\clearpage
\section{Programs}

The following programs compute the integral, polynomial and finite-coefficient
calculations associated with the indicated results. The stated finite ranges
and the calculation remarks specify the inputs of the enumerations.

\begingroup
\small
\setlength{\tabcolsep}{4pt}
\renewcommand{\arraystretch}{1.12}
\begin{longtable}{@{}p{0.49\linewidth}p{0.47\linewidth}@{}}
\textbf{Program path} & \textbf{Calculation and mathematical scope} \\
\hline
\endfirsthead
\textbf{Program path} & \textbf{Calculation and mathematical scope} \\
\hline
\endhead
\hline
\endfoot
\nolinkurl{code/S6_SERIES_PUBLISH_III_02JA_new_exact.py}
& Exact integral and mod-two invariant spaces, exterior powers, primitive kernels and marking matrices for the three named sphere families.
See \cref{maps:src:B.4.1,maps:src:B.1.1}. \\[3pt]
\nolinkurl{code/S6_SERIES_PUBLISH_III_02JC_fresh_arithmetic.py}
\par\nolinkurl{code/S6_SERIES_PUBLISH_III_02JC_comparison_arithmetic.py}
& Monodromy and integral coordinate maps, coinvariants, core kernels, cusp mixed cycles, normal circles and the tube matrix of determinant $-1$.
See \cref{maps:src:B.5.2}. \\[3pt]
\nolinkurl{code/kummer_normal_cover.py}
\par\nolinkurl{code/kummer_orbits.py}
& Normal-cover and finite Kummer-orbit calculations.
See \cref{maps:src:C.1.4,maps:src:C.1.5}. \\[3pt]
\nolinkurl{code/S6_SERIES_PUBLISH_III_02K5_fresh_orbit.py}
& Integral Kummer identities for $1\leq N\leq64$.
See \cref{maps:src:C.1.4}. \\[3pt]
\nolinkurl{code/S6_SERIES_PUBLISH_III_02K5_exterior_orbit.py}
\par\nolinkurl{code/S6_SERIES_PUBLISH_III_02K6_schreier_kernel.py}
& Exterior and unit-projective orbits for $2\leq N\leq8$; Schreier kernels and characters for $N=2,3,4$.
See \cref{maps:src:C.1.4,maps:src:C.1.5}. \\[3pt]
\nolinkurl{code/S6_SERIES_PUBLISH_III_02KC_fresh_linear.py}
\par\nolinkurl{code/S6_SERIES_PUBLISH_III_02KC_postrelease_matrix_comparison.py}
& Exact rational attachment and character matrices, their full inverses, and the symbolic $B/D$ coordinate transport.
See \cref{maps:src:D.1.3}. \\[3pt]
\nolinkurl{code/S6_SERIES_PUBLISH_III_02KD_fresh_arithmetic.py}
\par\nolinkurl{code/S6_SERIES_PUBLISH_III_02KD_comparison_arithmetic.py}
& Integral intersections, image and quotient inverses, norm averages and signed finite-order coefficient maps of orders three and four.
See \cref{maps:src:D.1.4}. \\[3pt]
\nolinkurl{code/S6_SERIES_PUBLISH_III_02KF_fresh_ring.py}
& Curve intersections, normal-disc chain and ruling inverse; both signs for $d\in\{1,5,7,11,13,17,19,25,31,49,101\}$, with cyclic ring matrices for $d>1$.
See \cref{maps:src:D.1.6}. \\[3pt]
\nolinkurl{code/S6_SERIES_PUBLISH_III_02KG_fresh_family.py}
& Symbolic finite orders, invariant columns, integral period change, $-d$ scaling, cusp cocycle, coherent even shifts, fan maps and parameter radii.
See \cref{maps:src:D.1.7}. \\[3pt]
\nolinkurl{code/S6_SERIES_PUBLISH_III_02KH_fresh_scalar.py}
& Symbolic covering-torus and scalar-line coordinate identities, including the marked quadratic expression.
See \cref{maps:src:D.1.8}. \\[3pt]
\nolinkurl{code/S6_SERIES_PUBLISH_III_02KI_fresh_cubic.py}
& Both ruling vectors, diagonal relative vector, three fibre evaluations, residue values, square $-12$ and the divisor relation.
See \cref{maps:src:D.1.9}. \\[3pt]
\nolinkurl{code/S6_SERIES_PUBLISH_III_02K8_fresh_homology.py}
& Symbolic $d,2d,3d$ relations; $d\in\{1,-1,5,-5,7,11,25,35,49,55,125,175\}$; tensor detectors in ranks $0$--$8$ and cubic coordinate inverses in ranks $0$--$4$.
See \cref{maps:src:D.1.10}. \\[3pt]
\nolinkurl{code/S6_SERIES_PUBLISH_III_02K9_fresh_classification.py}
& Full relation matrices, all ordered cubic triples and mixed/diagonal coefficient transports for $d=1,5,7,11,25,35,49$, including mod-two and factor-four maps.
See \cref{maps:src:D.1.11}. \\[3pt]
\nolinkurl{code/S6_SERIES_PUBLISH_III_02KA_fresh_action.py}
& Symbolic conjugacies, powers, affine sums, commutators, quotient inverses and fan characters; $30240$ freeness cases and $17550$ representatives for positive $d\leq161$ coprime to six.
See \cref{maps:src:D.1.12}. \\[3pt]
\nolinkurl{code/S6_SERIES_PUBLISH_III_02LO_fresh_intersections.py}
& The $61$ Hermite forms with entries at most five and index at most sixteen, cyclic cases $d=1,\ldots,30$, and all $64$ entries of the four-vertex intersection tensor.
See \cref{maps:src:D.2.47}. \\[3pt]
\nolinkurl{code/S6_SERIES_PUBLISH_III_02LO_polarization.py}
& All $127$ column-Hermite forms of index at most twelve and $23276$ unordered basis triples; seven-value cubic polarization compared with the surface restriction formula.
See \cref{maps:src:D.2.47}. \\[3pt]
\nolinkurl{code/III/manifolds_general_cusp.py}
& Laurent row/column operations and $75$ Hermite lattices $(a,b,c)$ with $1\leq a,c\leq5$ and $0\leq b<a$; normalization and fibre quotients, exterior returns and component meridians.
See \cref{maps:src:D.2.43}. \\[3pt]
\nolinkurl{code/S6_SERIES_PUBLISH_c7_lk_lattices_v1.py}
\par\nolinkurl{code/S6_SERIES_PUBLISH_c7_lk_saved_check_v1.py}
& The $127$ column-Hermite lattices of index at most twelve and twelve specified larger lattices; all $139$ lattice quotients, translation permutations, torsion generators, period/wedge inverses and disc-kernel maps.
See \cref{maps:src:D.2.43}. \\[3pt]
\nolinkurl{code/S6_SERIES_PUBLISH_c7_lk_post_laurent_v1.py}
\par\nolinkurl{code/S6_SERIES_PUBLISH_c7_lk_post_tor_v1.py}
& Literal Laurent identities, full chain homotopies and $139$ oriented Smith comparisons; derivative/division inverses on $278$ oriented bases, including $139$ negative-period bases and $2594$ free-section vectors.
See \cref{maps:src:D.2.43}. \\[3pt]
\nolinkurl{code/S6_SERIES_PUBLISH_c7_lk_historical_validate_v2.py}
& Ten symbolic coefficient fields and all $750$ finite coefficient fields in the specified $75$-case box, with exact kernel and component coordinate maps and inverses.
See \cref{maps:src:D.2.43}. \\[3pt]
\nolinkurl{code/III/receiving_final_certificates.py}
& Finite coset actions, peripheral words, matrices, cocycles and rational heights for the marked covering construction.
See \cref{mapsD:repair:collision-realization}. \\[3pt]
\nolinkurl{code/S6_SERIES_PUBLISH_III_02OG_comparison_arithmetic.py}
& Integral arithmetic comparison for the marked collision-cover construction.
See \cref{mapsD:repair:collision-realization}. \\[3pt]
\nolinkurl{code/S6_SERIES_PUBLISH_c7_ku_archive_audit.py}
& Monodromy reduction orders, occurrence permutations with inverses and switch-coordinate maps for $21$ profiles and $50$ inputs; four pairs of integral clutch coordinates.
See \cref{maps:src:D.2.2}. \\[3pt]
\path{code/S6_SERIES_PUBLISH_c9_ku_matrix_replay.py}
& The $50$ marked filling inputs and eight integral-clutch inputs:
full integral cochain cones, an independent block assembly,
Smith calculations in every degree, and full peripheral presentations
with integral centrality witnesses. The calculation uses the cited
local models and their marked boundary maps.
See \cref{maps:src:D.2.2}. \\[3pt]
\nolinkurl{code/III/manifolds_unbounded_logs.py}
& Exact logarithmic boundary matrices and determinant cases.
See \cref{maps:src:D.2.5}. \\[3pt]
\nolinkurl{code/III/manifolds_iii_star.py}
& The integral $\mathrm{III}^{*}$ boundary matrix.
See \cref{maps:src:D.2.6}. \\[3pt]
\nolinkurl{code/III/manifolds_wang.py}
& Wang boundary matrices and the associated additive-support calculation.
See \cref{maps:src:D.2.7,maps:src:D.2.8}. \\[3pt]
\nolinkurl{code/III/manifolds_liu_rollenske.py}
& Marked boundary matrices and affine-shift shears.
See \cref{maps:src:D.2.9}. \\[3pt]
\nolinkurl{code/III/manifolds_yang.py}
& Integral attachment matrices for all stated finite multiplicities.
See \cref{maps:src:D.2.10}. \\[3pt]
\nolinkurl{code/III/manifolds_type_c_ii.py}
& Period conjugacy and $168$ Smith-form examples for the coprime attachment.
See \cref{maps:src:D.2.11}. \\[3pt]
\nolinkurl{code/S6_SERIES_PUBLISH_c8_foundations_kl_maps.py}
& Symbolic $S/S^{-1}$, homological and cohomological conjugacies, the full augmented reduction, quotient map and cyclic-kernel preimage.
See \cref{maps:src:D.2.11}. \\[3pt]
\nolinkurl{code/III/manifolds_rankone_families.py}
& Seven logarithmic profiles and $234$ integral-attachment examples; the exterior-coordinate wedge used in the general cusp calculation.
See \cref{maps:src:D.2.12,maps:src:D.2.43}. \\[3pt]
\nolinkurl{code/S6_SERIES_PUBLISH_III_02KM_analytic_cover.py}
& Exact rational whole-segment Rouch\'e bounds and root continuation for the seven displayed Weierstrass models.
See \cref{maps:src:D.2.12}. \\[3pt]
\nolinkurl{code/S6_SERIES_PUBLISH_III_02KM_dessin_transport.py}
\par\nolinkurl{code/S6_SERIES_PUBLISH_III_02KM_section_index.py}
& Ribbon and peripheral marking words with their inverses; integral height and section-index maps for the eleven KM/KN sections.
See \cref{maps:src:D.2.12,maps:src:D.2.13}. \\[3pt]
\nolinkurl{code/S6_SERIES_PUBLISH_III_02KM_poincare_gauge.py}
\par\nolinkurl{code/S6_SERIES_PUBLISH_III_02KM_component_labels.py}
& Poincar\'e line and affine gauge maps with inverses, and the signed component-label involution.
See \cref{maps:src:D.2.12}. \\[3pt]
\nolinkurl{code/S6_SERIES_PUBLISH_III_02KM_fresh_arithmetic.py}
& Integral model matrices and exact lattice functions for the seven marked families.
See \cref{maps:src:D.2.12}. \\[3pt]
\nolinkurl{code/manifolds_semistable.py}
& The four semistable integral exterior and logarithmic boundary calculations.
See \cref{maps:src:D.2.13}. \\[3pt]
\nolinkurl{code/S6_SERIES_PUBLISH_III_02KN_preexposure_audit.py}
\par\nolinkurl{code/S6_SERIES_PUBLISH_III_02KN_comparison_arithmetic.py}
& Direct minors and full integral comparison maps, including reversed graph-block order, all-zero endpoints and the peripheral sign.
See \cref{maps:src:D.2.13}. \\[3pt]
\nolinkurl{code/S6_SERIES_PUBLISH_III_02KO_fresh_geometry.py}
\par\nolinkurl{code/S6_SERIES_PUBLISH_III_02KO_comparison.py}
& Primitive angular, fan, normal, Hopf and directed-gluing maps, and their oriented comparison.
See \cref{maps:src:D.2.14}. \\[3pt]
\nolinkurl{code/S6_SERIES_PUBLISH_III_02KP_dessin_fresh.py}
\par\nolinkurl{code/S6_SERIES_PUBLISH_III_02KP_modular_cosets_fresh.py}
\par\nolinkurl{code/S6_SERIES_PUBLISH_III_02KP_geometry_fresh.py}
& Exact rational surface covering, modular cosets and ordered peripheral marking.
See \cref{maps:src:D.2.15}. \\[3pt]
\nolinkurl{code/S6_SERIES_PUBLISH_III_02KP_boundary_order_fresh.py}
\par\nolinkurl{code/S6_SERIES_PUBLISH_III_02KP_lattice_fresh.py}
\par\nolinkurl{code/S6_SERIES_PUBLISH_III_02KP_comparison.py}
& Integral boundary order, complete lattice matrices, inverses, kernels and their marked comparison.
See \cref{maps:src:D.2.15}. \\[3pt]
\nolinkurl{code/S6_SERIES_PUBLISH_III_02KQ_fresh_arithmetic.py}
\par\nolinkurl{code/S6_SERIES_PUBLISH_III_02KQ_post_exposure.py}
& Odd-star lattices, full attachment matrices, cap maps and their integral comparison.
See \cref{maps:src:D.2.16}. \\[3pt]
\nolinkurl{code/S6_SERIES_PUBLISH_c7_kr_review_arithmetic.py}
& Original type-II and type-IV cap maps, primitive kernels and transfer constants.
See \cref{maps:src:D.2.17}. \\[3pt]
\nolinkurl{code/S6_SERIES_PUBLISH_c7_ks_integral.py}
& Literal periods, conductor Smith maps and inverses, star data, primitive kernels, both cubics, symbolic log minors and $405$ allowed Smith tuples.
See \cref{maps:src:D.2.18}. \\[3pt]
\nolinkurl{code/S6_SERIES_PUBLISH_c7_ks_comparison.py}
& Twelve triangular fan symmetries, wall columns, surgery trace and volume Jacobian.
See \cref{maps:src:D.2.18}. \\[3pt]
\nolinkurl{code/S6_SERIES_PUBLISH_III_02KT_fresh_calculation.py}
& Complete whole-fibre enumeration with totals $104$, $88$ and $348$ from the ordered profile data.
See \cref{maps:src:D.2.19}. \\[3pt]
\nolinkurl{code/S6_SERIES_PUBLISH_III_02KT_block_inverse.py}
\par\nolinkurl{code/S6_SERIES_PUBLISH_III_02KT_period_product.py}
& Integral block inverse and period products for all $21$ profiles, with symbolic free coefficients and weights.
See \cref{maps:src:D.2.19}. \\[3pt]
\nolinkurl{code/S6_SERIES_PUBLISH_c7_original_caps.py}
& Six original cap and Mayer--Vietoris maps, Smith kernels, zero-vertical endpoints and primitive rims.
See \cref{maps:src:D.2.20,maps:src:D.2.21}. \\[3pt]
\nolinkurl{code/S6_SERIES_PUBLISH_c7_original_log_arithmetic.py}
\par\nolinkurl{code/S6_SERIES_PUBLISH_c7_original_post_exposure.py}
& Literal presentations for $96$ regular-log and $468$ finite-log cases; graph-basis inverses for six original and five mixed cap-constant cases.
See \cref{maps:src:D.2.20,maps:src:D.2.21}. \\[3pt]
\nolinkurl{code/S6_SERIES_PUBLISH_c7_kx_vertical.py}
\par\nolinkurl{code/S6_SERIES_PUBLISH_c7_kx_auxiliary.py}
& Integral vertical solutions, the $172/226$ marked counts and exact cap-rank input maps.
See \cref{maps:src:D.2.22}. \\[3pt]
\nolinkurl{code/S6_SERIES_PUBLISH_III_02KY_fresh_enumeration.py}
\par\nolinkurl{code/S6_SERIES_PUBLISH_III_02KY_fresh_group_arithmetic.py}
& Natural-row and capped-profile enumeration from the ordered component data; symbolic group matrices for all orders.
See \cref{maps:src:D.2.23,mapsIII:def:elliptic-marked-data}. \\[3pt]
\nolinkurl{code/S6_SERIES_PUBLISH_c7_kz_star.py}
\par\nolinkurl{code/S6_SERIES_PUBLISH_c7_kz_auxiliary.py}
& Eight marked stars in both integral bases, boundary dimensions, $D_7$ intersection/Euler data, $234$ reflection offsets in twelve covering orders and quotient-action identities.
See \cref{maps:src:D.2.24}. \\[3pt]
\nolinkurl{code/S6_SERIES_PUBLISH_c7_semistable_initial_and_lattice.py}
\par\nolinkurl{code/S6_SERIES_PUBLISH_c7_semistable_auxiliary.py}
& The $62$ semistable rows, graph/exterior inverses, fan words and CRT maps; cofactor, primary-presentation and physical-weight identities.
See \cref{maps:src:D.2.25}. \\[3pt]
\nolinkurl{code/S6_SERIES_PUBLISH_c7_l1_catalogue.py}
\par\nolinkurl{code/S6_SERIES_PUBLISH_c7_l1_inverse_checks.py}
& The $172/148$ row catalogues and exact $164\to162\to59$ count maps on $50$ underlying rows.
See \cref{maps:src:D.2.26}. \\[3pt]
\nolinkurl{code/S6_SERIES_PUBLISH_c7_l1_post_exposure.py}
& All $172$ natural graph/fibre maps and $102$ finite positions.
See \cref{maps:src:D.2.26,maps:src:D.2.29,maps:src:D.2.30}. \\[3pt]
\nolinkurl{code/S6_SERIES_PUBLISH_c7_l2_star_integral.py}
\par\nolinkurl{code/S6_SERIES_PUBLISH_c7_l2_aux_integral.py}
& Sixteen marked star rows, both universal star bases, conductor lattices and torsion-sensitive indices; component-degree enumeration, Hermite comparisons, universal $\mathsf F_{\rm rk}$, $\mathsf S_{\rm rk}$, $\mathsf K_{\rm rk}$ and $\mathsf R_{\rm rk}$ identities, and the order-$18$ example.
See \cref{maps:src:D.2.27}. \\[3pt]
\nolinkurl{code/S6_SERIES_PUBLISH_c7_l3_profiles.py}
& All $21$ profiles and their $172$ natural rows; the $59$ selected patterns on $50$ rows.
See \cref{maps:src:D.2.28}. \\[3pt]
\nolinkurl{code/III/manifolds_unrestricted_finite_logs.py}
& All $68$ primitive-fixed positions, $1\leq k\leq6$ and every $(n,r)\bmod m$: $150059$ affine-power tests; $305$ full-matrix comparisons, including $99$ torsion-cap and $55$ zero-Euler cases.
See \cref{mapsIII:rem:primitive-finite-calculations}. \\[3pt]
\nolinkurl{code/III/manifolds_all_finite_logs.py}
& The $34$ additional lattice bases, $1\leq k\leq6$ and every residue pair: $38675$ affine-power tests; $175$ full normal/graph matrices, including $35$ torsion-cap cases; $725$ one-site Smith forms at five integer clutch values.
See \cref{mapsIII:rem:general-finite-calculations,maps:src:D.2.4}. \\[3pt]
\nolinkurl{code/S6_SERIES_PUBLISH_c7_finite_lattices.py}
\par\nolinkurl{code/S6_SERIES_PUBLISH_c7_finite_patterns.py}
\par\nolinkurl{code/S6_SERIES_PUBLISH_c7_one_finite_minors.py}
& All $102$ finite positions, $198$ filling patterns and $163200$ one-site minors, including zero-Euler and torsion-core endpoints.
See \cref{maps:src:D.2.29,maps:src:D.2.30}. \\[3pt]
\nolinkurl{code/S6_SERIES_PUBLISH_c7_finite_auxiliary.py}
& All $102$ boundary comparisons and $34$ literal norm/lift maps.
See \cref{maps:src:D.2.29,maps:src:D.2.30}. \\[3pt]
\nolinkurl{code/III/manifolds_original_zero_degree.py}
& Zero-Chern original-cap data: $396$ local occurrences, $228$ patterns and $144$ full-matrix comparisons.
See \cref{mapsIII:rem:zero-chern-original-calculations,maps:src:D.2.4}. \\[3pt]
\nolinkurl{code/III/manifolds_original_chern.py}
& All $46$ additive Chern vectors and $282$ patterns; $74$ full matrices, nine torsion and $54$ zero-Euler comparisons.
See \cref{mapsIII:rem:all-original-additive-calculations,maps:src:D.2.4}. \\[3pt]
\nolinkurl{code/III/manifolds_original_multiplicative.py}
& The $120$ local multiplicative normal forms, with the type subtotals in the cited remark; $213$ Smith comparisons, including $18$ torsion-cap and $172$ zero-Euler cases.
See \cref{mapsIII:rem:original-multiplicative-calculations,maps:src:D.2.4}. \\[3pt]
\nolinkurl{code/S6_SERIES_PUBLISH_c7_all_original.py}
\par\nolinkurl{code/S6_SERIES_PUBLISH_c7_all_original_auxiliary.py}
& All $172$ natural rows, $228/282$ patterns, inverse Smith maps, $396$ zero-Chern occurrences and $46$ additive Chern positions.
See \cref{maps:src:D.2.31,maps:src:D.2.32}. \\[3pt]
\nolinkurl{code/S6_SERIES_PUBLISH_c7_l9_catalogue.py}
\par\nolinkurl{code/S6_SERIES_PUBLISH_c7_l9_post_exposure.py}
& All $120$ nonzero- and $316$ zero-Chern positions, $495$ kernels, $1728$ representatives and $990$ complete finite/regular matrices; arbitrary integral rim shears.
See \cref{maps:src:D.2.33}. \\[3pt]
\nolinkurl{code/III/manifolds_weighted_engel.py}
& The arbitrary-integer free-action algorithm, polynomial identities and full group/homology calculations in the stated Engel domain.
See \cref{mapsIII:rem:positive-engel-calculations}. \\[3pt]
\nolinkurl{code/S6_SERIES_PUBLISH_c7_weighted_engel.py}
\par\nolinkurl{code/S6_SERIES_PUBLISH_c7_weighted_engel_auxiliary.py}
& All-weight symbolic identities; $2376$ free cases ($1536$ noncoprime, $840$ coprime), $216$ zero-Euler, $132/108/6$ original-cap cases and $46$ finite-meridian examples.
See \cref{maps:src:D.2.34}. \\[3pt]
\nolinkurl{code/III/manifolds_weighted_profiles.py}
& The $172$ constant inputs and $1626$ graph presentations, including $1131$ weights greater than one; the quotient-group helper used in the general cusp calculation.
See \cref{mapsIII:rem:weighted-profile-calculations,maps:src:D.2.43}. \\[3pt]
\nolinkurl{code/S6_SERIES_PUBLISH_c7_ld_exterior.py}
\par\nolinkurl{code/S6_SERIES_PUBLISH_c7_ld_finite.py}
& Twenty-one profile inverses, all $172$ rows and all-weight polynomial maps; $739$ conjugacies, $598$ projected lattices and $44850$ power comparisons.
See \cref{maps:src:D.2.37}. \\[3pt]
\nolinkurl{code/S6_SERIES_PUBLISH_c7_ld_catalogue.py}
\par\nolinkurl{code/S6_SERIES_PUBLISH_c7_ld_primitive_examples.py}
\par\nolinkurl{code/S6_SERIES_PUBLISH_c7_ld_post_exposure.py}
& The $495/459/36$ pattern maps and $4242$ full matrices; $102$ signed fixed examples, all-weight $H_d$ transport, $436$ rank-one factors, $172$ constant matrix pairs and cusp maps.
See \cref{maps:src:D.2.37}. \\[3pt]
\nolinkurl{code/III/manifolds_odd_weighted_profiles.py}
& Odd-weighted saturated kernels and coprime full-presentation comparisons in the domain of the cited remark.
See \cref{mapsIII:rem:odd-weighted-calculations}. \\[3pt]
\nolinkurl{code/S6_SERIES_PUBLISH_c7_lh_exterior_local_v3.py}
\par\nolinkurl{code/S6_SERIES_PUBLISH_c7_lh_component_denominators_v1.py}
& Allowed odd-common-factor weights among $1,\allowbreak2,\allowbreak3,\allowbreak5,\allowbreak7,\allowbreak9,\allowbreak15,\allowbreak21,\allowbreak25,\allowbreak35,\allowbreak45,\allowbreak105,\allowbreak315,\allowbreak1001$: $2334$ exterior and $586$ shared-factor cases, $12249$ saturated kernels, $4265$ basis maps, $2334$ cusp-disc maps, $1386$ lattices, $33264$ finite actions/proper powers and $126/1152$ denominator/orbit cases.
See \cref{maps:src:D.2.40}. \\[3pt]
\nolinkurl{code/S6_SERIES_PUBLISH_c7_lh_catalogue_v1.py}
\par\nolinkurl{code/S6_SERIES_PUBLISH_c7_lh_saved_audit_v2.py}
& The $495/459$ patterns and $6326$ complete global matrices: $4950$ on the two Euler values and $1376$ with two additional full-vector regular logarithms; $24$ reversible Smith equivalences and $112$ fibre-torsion comparisons.
See \cref{maps:src:D.2.40}. \\[3pt]
\nolinkurl{code/S6_SERIES_PUBLISH_c7_lh_post_exposure_v1.py}
& All $172$ elliptic image sections, all-weight angular preimages, both OS46/II index-three maps, the order-two example and $13359$ bounded odd-half residue equations.
See \cref{maps:src:D.2.40}. \\[3pt]
\nolinkurl{code/III/manifolds_all_weighted_profiles.py}
& The $6606$ coefficient and $1548$ word cases; $1953$ presentations with $426$ shared-even and $1527$ odd comparisons.
See \cref{mapsIII:rem:all-weight-integral-calculation}. \\[3pt]
\nolinkurl{code/S6_SERIES_PUBLISH_c7_li_ordered_words_v02.py}
& The $4002$ global and $1074$ shared-even cases, $17192$ ordered lifts, $15468$ transfer maps, $5726$ five-column maps and $1380$ power/shear cases.
See \cref{maps:src:D.2.41}. \\[3pt]
\nolinkurl{code/S6_SERIES_PUBLISH_c7_li_saved_output_check_v001.py}
& The specified $35$ integral matrices, $141$ quadratic blocks and $65$ affine-power identities.
See \cref{maps:src:D.2.41}. \\[3pt]
\nolinkurl{code/III/manifolds_general_initial.py}
\par\nolinkurl{code/III/manifolds_general_initial_audit.py}
& Whole-fibre integral presentations, examples and integer-closure data.
See \cref{maps:src:D.2.46}. \\[3pt]
\nolinkurl{code/III/manifolds_star_log_caps.py}
& Canonical-star local and global finite comparisons, together with the symbolic local expansion.
See \cref{maps:src:D.2.49}. \\[3pt]
\nolinkurl{code/S6_SERIES_PUBLISH_c7_ls_os45_bar_v001.py}
& All sixteen OS45 fillings: marked cap/bar maps, Smith quotients and inverse matrices.
See \cref{maps:src:D.2.50}. \\[3pt]
\nolinkurl{code/S6_SERIES_PUBLISH_c7_ls_post_local_v001.py}
\par\nolinkurl{code/S6_SERIES_PUBLISH_c7_ls_post_obstructions_v001.py}
\par\nolinkurl{code/S6_SERIES_PUBLISH_c7_ls_post_relations_v001.py}
& All $227$ local component pairs, $135$ literal Smith obstruction inputs and sixteen relation maps from the $23\times70$ to the $23\times59$ presentations.
See \cref{maps:src:D.2.50}. \\[3pt]
\nolinkurl{code/III/root_vertical_corrections.py}
\par\nolinkurl{code/III/fourfolds_corrected_Mumford_examples.py}
& Integral vertical-correction data and all sixteen corrected all-Mumford fillings.
See \cref{maps:src:D.2.50}. \\[3pt]
\nolinkurl{code/V/replay_chains/code/fourfolds_vertical_correction_audit.py}
& Auxiliary integral component and residue calculations.
See \cref{maps:src:D.2.50}. \\[3pt]
\nolinkurl{code/III_invariant_chern_cone_fresh.py}
& Integral cone, scalar-long order and multiplicity-cycle calculations on $390$ formal chains, $172$ torsion models and nineteen marked additive-component models.
See \cref{maps:src:D.2.51}. \\[3pt]
\nolinkurl{code/III/manifolds_corrected_initial.py}
& Height-integral corrected inputs and the exact original-multiplicative cap and frame criterion.
See \cref{maps:src:D.2.52}. \\[3pt]
\nolinkurl{code/III/manifolds_corrected_initial_audit.py}
& The $42$ overlap tables, sixteen OS45 inputs, $212$ new-residue inputs as a multiset and $45$ higher-weight/log inputs.
See \cref{maps:src:D.2.52}. \\[3pt]
\nolinkurl{code/III/manifolds_tate_markings.py}
& Integral Tate markings reconstructed from the actual monodromies for the listed minimal and three-primary rows.
See \cref{maps:src:D.2.35,maps:src:D.2.36,maps:src:D.2.39}. \\[3pt]
\nolinkurl{code/III/fourfolds_D2_c2_AH.py}
& Fifteen patterns and thirteen marked rows: $928$ semicharacter substitutions, both Hermite lattices, conjugated monodromy, toric characters and degree/normal-column parities.
See \cref{maps:src:D.2.35}. \\[3pt]
\nolinkurl{code/S6_SERIES_PUBLISH_c7_lb_conic.py}
\par\nolinkurl{code/S6_SERIES_PUBLISH_c7_lb_auxiliary_maps.py}
& The complete fifteen-pattern conic table and order-two invariant; exact auxiliary lattice, framing and surgery maps.
See \cref{maps:src:D.2.35}. \\[3pt]
\nolinkurl{code/III/manifolds_odd_prime_AH.py}
& Doubled-line identities, support functions, primitive angular lattices, scalar isotropy, fixed elliptic curves and the listed signed characteristic inputs.
See \cref{maps:src:D.2.36}. \\[3pt]
\nolinkurl{code/S6_SERIES_PUBLISH_c7_lc_odd_cubic.py}
\par\nolinkurl{code/S6_SERIES_PUBLISH_c7_lc_post_exposure.py}
& Complete marked odd-primary tables, physical weights and the simultaneous subgroup/coefficient coordinate transports.
See \cref{maps:src:D.2.36}. \\[3pt]
\nolinkurl{code/III/manifolds_mixed_prime_AH.py}
& Twelve patterns and ten rows: primitive angular lattices, $736$ semicharacters, period differences, scalar characters, fixed curves and odd normal weights.
See \cref{maps:src:D.2.38}. \\[3pt]
\nolinkurl{code/S6_SERIES_PUBLISH_c7_le_mixed_v001.py}
\par\nolinkurl{code/S6_SERIES_PUBLISH_c7_le_post_exposure_v001.py}
& Mixed-primary integral lattices, marked invariants and exact auxiliary maps, including $864$ semicharacter checks.
See \cref{maps:src:D.2.38}. \\[3pt]
\nolinkurl{code/III/fourfolds_D2_c3_AH.py}
& The listed three-primary row data and degree auxiliaries, with the explicit fan-coordinate transport in the cited proof.
See \cref{maps:src:D.2.39}. \\[3pt]
\nolinkurl{code/III/manifolds_three_mixed_AH.py}
& Mixed three-primary inputs and invariant data.
See \cref{maps:src:D.2.42}. \\[3pt]
\nolinkurl{code/S6_SERIES_PUBLISH_c7_lf_three_primary_v003.py}
\par\nolinkurl{code/S6_SERIES_PUBLISH_c7_lf_auxiliary_maps_v001.py}
& Eight patterns and their complete integral maps; $27$ site and seven fan comparisons.
See \cref{maps:src:D.2.39}. \\[3pt]
\nolinkurl{code/S6_SERIES_PUBLISH_c7_lj_mixed_three_v001.py}
\par\nolinkurl{code/S6_SERIES_PUBLISH_c7_lj_action_coefficients_v001.py}
\par\nolinkurl{code/S6_SERIES_PUBLISH_c7_lj_post_exposure_v001.py}
& Fourteen patterns, full CRT and $528$ action/refinement identities; twelve keys, $464$ refinements and $52$ adapted matrices.
See \cref{maps:src:D.2.42}. \\[3pt]
\nolinkurl{code/III/fourfolds_D2_2power_AH.py}
\par\nolinkurl{code/III/manifolds_even_critical_cycles.py}
\par\nolinkurl{code/III/fourfolds_D2_frame_certificate.py}
& Seven selected two-primary invariant rows, fourteen signed critical-cycle rows and fourteen frame rows, with the integral Cartier and frame maps of the cited proof.
See \cref{maps:src:D.2.44}. \\[3pt]
\nolinkurl{code/III/maps_D2_II_frame_checks.py}
& Integral type-II kernel and degree-two frame calculations.
See \cref{maps:src:D.2.44}. \\[3pt]
\nolinkurl{code/III/fourfolds_D2_remaining_AH.py}
\par\nolinkurl{code/III/manifolds_extended_critical_cycles.py}
& The $1056$ semicharacters, $62$ valuations, $158$ pairings and eleven critical preimages used in the residual comparison.
See \cref{maps:src:D.2.45}. \\[3pt]
\nolinkurl{code/S6_SERIES_PUBLISH_c7_ll_spin_localization.py}
\par\nolinkurl{code/S6_SERIES_PUBLISH_c7_ll_gamma_character.py}
\par\nolinkurl{code/S6_SERIES_PUBLISH_c7_ll_critical_cycles_v01.py}
& Seven-row cyclic localization, the integral characteristic-index character and all seven critical chains with both signs.
See \cref{maps:src:D.2.44}. \\[3pt]
\nolinkurl{code/S6_SERIES_PUBLISH_c7_lm_fresh_geometry_v003.py}
\par\nolinkurl{code/S6_SERIES_PUBLISH_c7_lm_receive_aux_v001.py}
& All eleven residual critical rows and integral sections; $1056$ semicharacter, $62$ valuation, $158$ pairing and eleven integral-preimage comparisons.
See \cref{maps:src:D.2.45}. \\[3pt]
\nolinkurl{code/S6_SERIES_PUBLISH_c8_residual_wall_arithmetic.py}
& Exact wall matrix and inverse, the eleven-row total-invariant substitutions, and both-sign primitive-cap projection identities.
See \cref{maps:src:D.2.45}. \\[3pt]
\nolinkurl{code/III/fourfolds_D48_thomas.py}
& The $4920$ weighted cusp cases in the stated range.
See \cref{maps:src:D.2.48}. \\[3pt]
\nolinkurl{code/S6_SERIES_PUBLISH_c7_lp_fresh_geometry.py}
& Exact primitive lattices, divisor restrictions and weighted cusp geometry.
See \cref{maps:src:D.2.48}. \\[3pt]
\nolinkurl{code/S6_SERIES_PUBLISH_c7_lp_finite_character.py}
& The $24$-convolution finite Fourier character obstruction.
See \cref{maps:src:D.2.48}. \\[3pt]
\nolinkurl{code/S6_SERIES_PUBLISH_c7_lp_circle_bordism.py}
\par\nolinkurl{code/S6_SERIES_PUBLISH_c7_lp_aux_comparison_v002.py}
& Symbolic circle endpoints, $160$ finite cases and expanded-difference comparisons.
See \cref{maps:src:D.2.48}. \\[3pt]
\nolinkurl{code/III_OS45_core_independent.py}
\par\nolinkurl{code/III_OS45_primitive_cubic_independent.py}
\par\nolinkurl{code/III_OS45_annulus_independent.py}
& Integral core relations, primitive cubic coefficients and annular boundary maps for the OS45 family.
See \cref{maps:src:D.2.53}. \\[3pt]
\nolinkurl{code/III_OS45_boundary_prisms_07.py}
\par\nolinkurl{code/III_weighted_OS45_checks.py}
& Boundary-prism and cap null witnesses; weighted coefficient identities.
See \cref{maps:src:D.2.53}. \\[3pt]
\nolinkurl{code/S6_SERIES_PUBLISH_c7_lv_fresh_lattices_v001.py}
\par\nolinkurl{code/S6_SERIES_PUBLISH_c7_lv_toric_coefficients_v001.py}
\par\nolinkurl{code/S6_SERIES_PUBLISH_c7_lv_graph_link_v001.py}
& Exact marked lattices, toric coefficient identities and integral graph/link calculations for the weighted OS45 family.
See \cref{maps:src:D.2.53}. \\[3pt]
\nolinkurl{code/S6_SERIES_PUBLISH_c8_lv_compare_arithmetic_v001.py}
\par\nolinkurl{code/S6_SERIES_PUBLISH_c8_lv_character_transport_v001.py}
& Relation, Picard and critical-chain coordinate maps; all-integer floor identities, $32$ orbit transitions and the physical character/coefficient polynomials.
See \cref{maps:src:D.2.53}. \\[3pt]
\nolinkurl{code/III/manifolds_sphere_initial.py}
\par\nolinkurl{code/III/manifolds_nonstar_catalogue.py}
& Sphere-pattern and nonstar catalogues compared by the exact count maps in the cited synthesis.
See \cref{maps:src:D.2.4}. \\[3pt]
\nolinkurl{code/S6_SERIES_PUBLISH_c7_lg_catalogue_v001.py}
& Enumeration from the proved per-pattern formulas: $172$ raw rows, $495$ patterns, $162$ minimal patterns and $59$ sphere patterns on $50$ rows.
See \cref{maps:src:D.2.4}. \\[3pt]
\nolinkurl{code/S6_SERIES_PUBLISH_c8_lg_coefficients_v001.py}
& All $162$ minimal marked assignments and their $24$ complete-invariant orbits; full linear/cubic CRT inverses for factors $(2,3),(2,5),(2,7),(3,5),(4,5),(2,3,5)$.
See \cref{maps:src:D.2.4}. \\[3pt]
\nolinkurl{code/S6_SERIES_PUBLISH_c8_lg_residual_comparison_v001.py}
\par\nolinkurl{code/S6_SERIES_PUBLISH_c8_lg_history_counts_v001.py}
& All eighteen residual assignments (twelve corrected, six unchanged), cyclic units and CRT inverses; the $495$-pattern bijection, $459$-pattern retraction and $36$-pattern complement, and $164\to162\to59$ on $50$ rows.
See \cref{maps:src:D.2.4}. \\[3pt]
\nolinkurl{code/S6_SERIES_PUBLISH_c8_later_os49_marking.py}
\par\nolinkurl{code/S6_SERIES_PUBLISH_c8_later_os49_bridge_v002.py}
& Integral OS49/Type-C marking matrices and inverses; period affine laws, $G_0$, parameter inverse, contractions and cross product.
See \cref{os49:rem:type-c-marking,night:P3F.29}. \\[3pt]
\nolinkurl{code/S6_SERIES_PUBLISH_c7_graph_integral_pairing.py}
& Ordered handle form, integral dual cochain and the symbolic all-$k$ pairing $-8k$; dense integral matrices for $k=1,2,3$.
See \cref{night:P3F.31,night:P3F.32}. \\[3pt]
\nolinkurl{code/S6_SERIES_PUBLISH_c7_graph_cap_orientation.py}
& Sparse exterior-algebra expansion of the cap normal frame, arbitrary height/radial derivatives and the positive first-jet determinant.
See \cref{night:P3F.31,night:P3F.32}. \\[3pt]
\nolinkurl{code/S6_SERIES_PUBLISH_III_037X_fresh_arithmetic.py}
& Symbolic Kummer characters, units and fan identities; stabilizers and monomial bases for $N=1,\ldots,8$.
See \cref{wb:cusp-cover-extension-theorem}. \\[3pt]
\nolinkurl{code/S6_SERIES_PUBLISH_III_0382_fresh_arithmetic.py}
& Exponent and character identities: polynomials through degree ten for $N=1,\ldots,6$, and coefficient indices $1$--$128$ for $N=2,\ldots,8$.
See \cref{wb:cusp-cover-naive-action-failure}. \\[3pt]
\nolinkurl{code/S6_SERIES_PUBLISH_III_037Y_fresh_arithmetic.py}
& Symbolic $A_1/A_2$ witnesses, changed-base intertwining, index and block maps; integrality examples for $N=1,\ldots,8$.
See \cref{wb:cusp-cover-global-scope}. \\[3pt]
\end{longtable}
\endgroup

\clearpage
\begin{thebibliography}{99}
\providecommand{\natexlab}[1]{#1}
\providecommand{\doi}[1]{\url{https://doi.org/#1}}

\bibitem[Andrus(1977)]{manAndrus77}
D.~G. Andrus.
\newblock \emph{Classification of Certain 6-Manifolds}.
\newblock M.Sc. thesis, McMaster University, January 1977.
\newblock \url{https://hdl.handle.net/11375/9107}.

\bibitem[Atiyah and Bott(1984)]{AtiyahBott1984}
M.~F. Atiyah and R.~Bott.
\newblock The moment map and equivariant cohomology.
\newblock \emph{Topology} \textbf{23} (1984), no.~1, 1--28.
\newblock \doi{10.1016/0040-9383(84)90021-1}.

\bibitem[Atiyah et~al.(1975)Atiyah, Patodi, and Singer]{mapsAPS2}
M.~F. Atiyah, V.~K. Patodi, and I.~M. Singer.
\newblock Spectral asymmetry and Riemannian geometry.~II.
\newblock \emph{Mathematical Proceedings of the Cambridge Philosophical Society}
\textbf{78} (1975), no.~3, 405--432.
\newblock \doi{10.1017/S0305004100051872}.

\bibitem[Atiyah and Singer(1968)]{AtiyahSinger1968IIIE}
M.~F. Atiyah and I.~M. Singer.
\newblock The index of elliptic operators.~III.
\newblock \emph{Annals of Mathematics} \textbf{87} (1968), no.~3, 546--604.
\newblock \doi{10.2307/1970717}.

\bibitem[Barden(1965)]{manBarden65}
D.~Barden.
\newblock Simply connected five-manifolds.
\newblock \emph{Annals of Mathematics} \textbf{82} (1965), no.~3, 365--385.
\newblock \doi{10.2307/1970702}.

\bibitem[Baues(1998)]{CuspInvariantBaues1998}
Hans-Joachim Baues.
\newblock The degree of maps between certain 6-manifolds.
\newblock \emph{Compositio Mathematica} \textbf{110} (1998), 51--64.
\newblock \doi{10.1023/A:1000203804884}.

\bibitem[Besser and Livn\'e(2012)]{manBesserLivne12}
Amnon Besser and Ron Livn\'e.
\newblock Picard--Fuchs equations of families of QM abelian surfaces.
\newblock arXiv:1202.2808v2, 14 February 2012.
\newblock \url{https://arxiv.org/abs/1202.2808v2}.

\bibitem[Borel and Serre(1953)]{BorelSerre1953}
Armand Borel and Jean-Pierre Serre.
\newblock Groupes de Lie et puissances r\'eduites de Steenrod.
\newblock \emph{American Journal of Mathematics} \textbf{75} (1953), no.~3, 409--448.
\newblock \doi{10.2307/2372495}.

\bibitem[Bott(1959)]{mapsBott59}
Raoul Bott.
\newblock The stable homotopy of the classical groups.
\newblock \emph{Annals of Mathematics} \textbf{70} (1959), no.~2, 313--337.
\newblock \doi{10.2307/1970106}.

\bibitem[Brown(1972)]{ffBrown72}
Edgar~H. Brown, Jr.
\newblock Generalizations of the Kervaire invariant.
\newblock \emph{Annals of Mathematics} \textbf{95} (1972), no.~2, 368--383.
\newblock \doi{10.2307/1970804}.

\bibitem[Brown(1982)]{manBrown82}
Kenneth~S. Brown.
\newblock \emph{Cohomology of Groups}.
\newblock Graduate Texts in Mathematics, vol.~87, Springer-Verlag, New York, 1982.
\newblock \doi{10.1007/978-1-4684-9327-6}.

\bibitem[Burklund et~al.(2023)Burklund, Hahn, and Senger]{BHS19}
Robert Burklund, Jeremy Hahn, and Andrew Senger.
\newblock On the boundaries of highly connected, almost closed manifolds.
\newblock \emph{Acta Mathematica} \textbf{231} (2023), no.~2, 205--344.
\newblock \doi{10.4310/ACTA.2023.v231.n2.a1}.
\newblock Page references use arXiv:1910.14116v3, 10 May 2022.
\newblock \url{https://arxiv.org/abs/1910.14116v3}.

\bibitem[Campana et~al.(1998)Campana, Demailly, and Peternell]{manCDP98}
Fr\'ed\'eric Campana, Jean-Pierre Demailly, and Thomas Peternell.
\newblock The algebraic dimension of compact complex threefolds with vanishing
second Betti number.
\newblock \emph{Compositio Mathematica} \textbf{112} (1998), no.~1, 77--91.
\newblock \doi{10.1023/A:1000313214795}.

\bibitem[Cartan(1954--1955)]{manCartan1955EML}
Henri Cartan.
\newblock D\'etermination des alg\`ebres \(H_*(\pi,n;\mathbf Z)\).
\newblock \emph{S\'eminaire Henri Cartan} \textbf{7} (1954--1955),
no.~1, expos\'e~11, 11-01--11-24, 31 January 1955.
\newblock \url{https://www.numdam.org/item/SHC_1954-1955__7_1_A11_0/}.

\bibitem[Cerf(1970)]{Cerf1970}
Jean Cerf.
\newblock La stratification naturelle des espaces de fonctions
diff\'erentiables r\'eelles et le th\'eor\`eme de la pseudo-isotopie.
\newblock \emph{Publications Math\'ematiques de l'IH\'ES} \textbf{39} (1970), 5--173.
\newblock \doi{10.1007/BF02684687}.

\bibitem[Claude and Alp{\"o}ge(2026)]{AC26}
Claude and Levent Alp{\"o}ge.
\newblock The $(3,4,\infty)$ modular family of $2$-tori, completed at its three
special points, is a complex structure on $S^6$.
\newblock Preprint, 2026, 108~pp.
\newblock \url{https://alpo.ge/s6.pdf}.

\bibitem[Debray and Gunningham(2018)]{DebrayGunningham2018Root}
Arun Debray and Sam Gunningham.
\newblock The Arf--Brown TQFT of $\mathrm{Pin}^-$ surfaces.
\newblock In \emph{Topology and Quantum Theory in Interaction},
Contemporary Mathematics, vol.~718, American Mathematical Society,
Providence, 2018, 49--87.
\newblock \doi{10.1090/conm/718/14478}.
\newblock Page references use arXiv:1803.11183v1, 29 March 2018,
whose title page is dated 30 March 2018.
\newblock \url{https://arxiv.org/abs/1803.11183v1}.

\bibitem[Demailly(2012)]{mapsDemailly2012CADG}
Jean-Pierre Demailly.
\newblock \emph{Complex Analytic and Differential Geometry}.
\newblock Version of 21 June 2012.
\newblock \url{https://www-fourier.univ-grenoble-alpes.fr/~demailly/manuscripts/agbook.pdf}.

\bibitem[Edixhoven and Lido(2023)]{manEdixhovenLido23}
B.~Edixhoven and G.~Lido.
\newblock Geometric quadratic Chabauty.
\newblock \emph{Journal of the Institute of Mathematics of Jussieu}
\textbf{22} (2023), no.~1, 279--333.
\newblock \doi{10.1017/S1474748021000244}.

\bibitem[Engel(2026{\natexlab{a}})]{E26}
Philip Engel.
\newblock Complex structures on $S^6$.
\newblock arXiv:2609.38442v1, 29 September 2026.
\newblock \url{https://arxiv.org/abs/2609.38442v1}.

\bibitem[Engel(2026{\natexlab{b}})]{manExplorer26}
Philip Engel.
\newblock Threefold explorer, 2026.
\newblock Documentation, marked local-model database and source,
consulted 26 November 2026.
\newblock \url{https://philip-engel.github.io/threefold-explorer/}.
\newblock Documentation:
\url{https://philip-engel.github.io/threefold-explorer/docs/}.

\bibitem[Fukae(2002)]{manFukae02}
Mitsuaki Fukae.
\newblock Monodromies of rational elliptic surfaces and extremal elliptic K3 surfaces.
\newblock arXiv:math/0205062v1, 7 May 2002.
\newblock \url{https://arxiv.org/abs/math/0205062v1}.

\bibitem[Galaz-Garc\'{i}a and Reiser(2025)]{manGR25}
Fernando Galaz-Garc\'{i}a and Philipp Reiser.
\newblock Free torus actions and twisted suspensions.
\newblock \emph{Forum of Mathematics, Sigma} \textbf{13} (2025), e3.
\newblock \doi{10.1017/fms.2024.141}.

\bibitem[Grauert(1960)]{Grauert1960Coherent}
Hans Grauert.
\newblock Ein Theorem der analytischen Garbentheorie und die
Modulr\"aume komplexer Strukturen.
\newblock \emph{Publications Math\'ematiques de l'IH\'ES} \textbf{5} (1960), 5--64.
\newblock \doi{10.1007/BF02684746}.

\bibitem[Grauert(1963)]{Grauert1963Correction}
Hans Grauert.
\newblock Berichtigung: Ein Theorem der analytischen Garbentheorie und die
Modulr\"aume komplexer Strukturen.
\newblock \emph{Publications Math\'ematiques de l'IH\'ES} \textbf{16} (1963), 35--36.
\newblock \doi{10.1007/BF02684291}.

\bibitem[Grothendieck(2003)]{SGA1Recomposed2004}
A.~Grothendieck (dir.), with two additional expos\'es by M.~Raynaud.
\newblock \emph{Rev\^etements \'etales et groupe fondamental (SGA~1)}.
\newblock Recomposed and annotated edition, Documents Math\'ematiques, vol.~3,
Soci\'et\'e Math\'ematique de France, 2003.
\newblock \url{https://smf.emath.fr/node/43215}.
\newblock Page references use arXiv:math/0206203v2, 4 January 2004.
\newblock \url{https://arxiv.org/abs/math/0206203v2}.

\bibitem[Hatcher(2002)]{Hatcher2002}
Allen Hatcher.
\newblock \emph{Algebraic Topology}.
\newblock Cambridge University Press, Cambridge, 2002.
\newblock \url{https://pi.math.cornell.edu/~hatcher/AT/AT.pdf}.

\bibitem[Hatcher(2017)]{CuspInvariantHatcherVB}
Allen Hatcher.
\newblock \emph{Vector Bundles and K-Theory}.
\newblock Version~2.2, November 2017.
\newblock \url{https://pi.math.cornell.edu/~hatcher/VBKT/VB.pdf}.

\bibitem[Hatcher(2022)]{HatcherSpectralSequences}
Allen Hatcher.
\newblock Spectral sequences.
\newblock Preliminary additional Chapter~5 of \emph{Algebraic Topology},
PDF version of 29 April 2022, 117 pages.
\newblock \url{https://pi.math.cornell.edu/~hatcher/AT/ATch5.pdf}.
\newblock Page references use this version, consulted 26 November 2026.

\bibitem[Herfurtner(1990)]{manHerfurtner90}
Stephan Herfurtner.
\newblock Elliptic surfaces with four singular fibres.
\newblock Max-Planck-Institut f\"ur Mathematik preprint MPI/90-28, 1990.
\newblock \url{https://archive.mpim-bonn.mpg.de/id/eprint/860/1/preprint_1990_28.pdf}.

\bibitem[Isaksen et~al.(2023)Isaksen, Wang, and Xu]{IsaksenWangXu2023}
Daniel~C. Isaksen, Guozhen Wang, and Zhouli Xu.
\newblock Stable homotopy groups of spheres: from dimension 0 to 90.
\newblock \emph{Publications Math\'ematiques de l'IH\'ES} \textbf{137} (2023), 107--243.
\newblock \doi{10.1007/s10240-023-00139-1}.

\bibitem[Iwase(2022)]{IwaseWhitney2022}
N.~Iwase.
\newblock Whitney approximation for smooth CW complex.
\newblock \emph{Kyushu Journal of Mathematics} \textbf{76} (2022), no.~1, 177--186.
\newblock \doi{10.2206/kyushujm.76.177}.
\newblock Page references use arXiv:2001.02893v4, 5 July 2021.
\newblock \url{https://arxiv.org/abs/2001.02893v4}.

\bibitem[Johnson(1980)]{Johnson1980Root}
Dennis Johnson.
\newblock Spin structures and quadratic forms on surfaces.
\newblock \emph{Journal of the London Mathematical Society} (2)
\textbf{22} (1980), no.~2, 365--373.
\newblock \doi{10.1112/jlms/s2-22.2.365}.

\bibitem[Kaplan(1979)]{mapsKaplan79}
Steve~J. Kaplan.
\newblock Constructing framed 4-manifolds with given almost framed boundaries.
\newblock \emph{Transactions of the American Mathematical Society}
\textbf{254} (1979), 237--263.
\newblock \doi{10.1090/S0002-9947-1979-0539917-X}.

\bibitem[Kas(1980)]{Kas1980}
Arnold Kas.
\newblock On the handlebody decomposition associated to a Lefschetz fibration.
\newblock \emph{Pacific Journal of Mathematics} \textbf{89} (1980), no.~1, 89--104.
\newblock \doi{10.2140/pjm.1980.89.89}.

\bibitem[Kervaire and Milnor(1963)]{manKM63}
Michel~A. Kervaire and John~W. Milnor.
\newblock Groups of homotopy spheres.~I.
\newblock \emph{Annals of Mathematics} \textbf{77} (1963), no.~3, 504--537.
\newblock \doi{10.2307/1970128}.

\bibitem[Li(1995)]{mapsLi1995Cycles}
Yingchen Li.
\newblock Analytic cycles and generically finite holomorphic maps.
\newblock \emph{Bulletin of the Australian Mathematical Society}
\textbf{52} (1995), 457--460.
\newblock \doi{10.1017/S0004972700014933}.

\bibitem[Liu and Rollenske(2026)]{manLR26}
Wenfei Liu and S{\"o}nke Rollenske.
\newblock Yet another family of complex structures on $S^6$.
\newblock arXiv:2609.33826v2, 9 October 2026.
\newblock \url{https://arxiv.org/abs/2609.33826v2}.

\bibitem[Milne(2017)]{mapsMilne2017}
James~S. Milne.
\newblock \emph{Modular Functions and Modular Forms (Elliptic Modular Curves)}.
\newblock Version~1.31, 22 March 2017.
\newblock \url{https://www.jmilne.org/math/CourseNotes/MF.pdf}.

\bibitem[Milnor(1964)]{MilnorDifferentialTopology1964}
John~W. Milnor.
\newblock Differential topology.
\newblock In Thomas~L. Saaty (ed.), \emph{Lectures on Modern Mathematics},
vol.~2, John Wiley and Sons, New York, 1964, 165--183.

\bibitem[Milnor(1965{\natexlab{a}})]{mapsMilnor65}
John~W. Milnor.
\newblock \emph{Lectures on the $h$-Cobordism Theorem}.
\newblock Notes by L.~Siebenmann and J.~Sondow.
\newblock Princeton University Press, Princeton, New Jersey, 1965.

\bibitem[Milnor(1965{\natexlab{b}})]{CuspNullMilnor1965}
John~W. Milnor.
\newblock \emph{Topology from the Differentiable Viewpoint}.
\newblock University Press of Virginia, Charlottesville, 1965.

\bibitem[Milnor and Stasheff(1974)]{MilnorStasheff1974}
John~W. Milnor and James~D. Stasheff.
\newblock \emph{Characteristic Classes}.
\newblock Annals of Mathematics Studies, vol.~76,
Princeton University Press, Princeton, 1974.
\newblock \doi{10.1515/9781400881826}.

\bibitem[Miranda(1990)]{manMiranda90}
Rick Miranda.
\newblock Persson's list of singular fibers for a rational elliptic surface.
\newblock \emph{Mathematische Zeitschrift} \textbf{205} (1990), 191--211.
\newblock \doi{10.1007/BF02571235}.
\newblock \url{https://www.math.colostate.edu/~miranda/preprints/Perssonslist.pdf}.

\bibitem[Oguiso and Shioda(1991)]{manOguisoShioda91}
Keiji Oguiso and Tetsuji Shioda.
\newblock The Mordell--Weil lattice of a rational elliptic surface.
\newblock \emph{Commentarii Mathematici Universitatis Sancti Pauli}
\textbf{40} (1991), 83--99.
\newblock \url{https://rikkyo.repo.nii.ac.jp/record/10006/files/AA00610867_40-01_08.pdf}.

\bibitem[Pontryagin(1950)]{Pontryagin1950Root}
L.~S. Pontryagin.
\newblock Homotopy classification of the mappings of an $(n+2)$-dimensional
sphere on an $n$-dimensional one.
\newblock \emph{Doklady Akademii Nauk SSSR} \textbf{70} (1950), 957--959 (Russian).

\bibitem[Pontryagin(1955, 1959)]{Pontryagin1955Root}
L.~S. Pontryagin.
\newblock Smooth manifolds and their applications in homotopy theory.
\newblock \emph{Trudy Matematicheskogo Instituta imeni V.~A.~Steklova}
\textbf{45}, Academy of Sciences of the USSR, Moscow, 1955, 3--139 (Russian).
\newblock \url{https://www.mathnet.ru/eng/tm1181}.
\newblock English translation by P.~J. Hilton, \emph{American Mathematical
Society Translations}, Series~2, vol.~11, 1959, 1--114.
Page references use this English translation.
\newblock \url{https://bookstore.ams.org/trans2-11}.

\bibitem[Redden(2011)]{Redden2011StringStructures}
Corbett Redden.
\newblock String structures and canonical 3-forms.
\newblock \emph{Pacific Journal of Mathematics} \textbf{249} (2011), no.~2, 447--484.
\newblock \doi{10.2140/pjm.2011.249.447}.
\newblock Page references use arXiv:0912.2086v2, 8 February 2011.
\newblock \url{https://arxiv.org/abs/0912.2086v2}.

\bibitem[Reiser(2024)]{manReiser24}
Philipp Reiser.
\newblock Positive Ricci curvature on twisted suspensions.
\newblock \emph{International Mathematics Research Notices}
\textbf{2024} (2024), no.~22, 14115--14137.
\newblock \doi{10.1093/imrn/rnae231}.

\bibitem[Sch\"utt and Shioda(2010)]{manSS10}
Matthias Sch\"utt and Tetsuji Shioda.
\newblock Elliptic surfaces.
\newblock In \emph{Algebraic Geometry in East Asia---Seoul 2008},
Advanced Studies in Pure Mathematics, vol.~60,
Mathematical Society of Japan, Tokyo, 2010, 51--160.
\newblock \doi{10.2969/aspm/06010051}.
\newblock Page references use arXiv:0907.0298v3, 9 July 2010.
\newblock \url{https://arxiv.org/abs/0907.0298v3}.

\bibitem[The Stacks Project Authors(2026)]{StacksProject}
The Stacks Project Authors.
\newblock \emph{The Stacks Project}, 2026.
\newblock \url{https://stacks.math.columbia.edu}.

\bibitem[Telen(2022)]{TelenToric2022}
S.~Telen.
\newblock Introduction to toric geometry.
\newblock arXiv:2203.01690v1, 3 March 2022.
\newblock \url{https://arxiv.org/abs/2203.01690v1}.

\bibitem[Teschl(2012)]{TeschlODE2012}
Gerald Teschl.
\newblock \emph{Ordinary Differential Equations and Dynamical Systems}.
\newblock Graduate Studies in Mathematics, vol.~140,
American Mathematical Society, Providence, 2012.
\newblock Page references use the author's preliminary version.
\newblock \url{https://www.mat.univie.ac.at/~gerald/ftp/book-ode/ode.pdf}.

\bibitem[Thomas(1956)]{ffThomas56}
Emery Thomas.
\newblock A generalization of the Pontrjagin square cohomology operation.
\newblock \emph{Proceedings of the National Academy of Sciences of the
United States of America} \textbf{42} (1956), no.~5, 266--269.
\newblock \doi{10.1073/pnas.42.5.266}.

\bibitem[Thomas(1959)]{ffThomas1959Suspension}
Emery Thomas.
\newblock The suspension of the generalized Pontrjagin cohomology operations.
\newblock \emph{Pacific Journal of Mathematics} \textbf{9} (1959), no.~3, 897--911.
\newblock \url{https://msp.org/pjm/1959/9-3/pjm-v9-n3-p26-p.pdf}.

\bibitem[Toda(1952)]{mapsToda52}
Hirosi Toda.
\newblock Some relations in homotopy groups of spheres.
\newblock \emph{Journal of the Institute of Polytechnics, Osaka City
University, Series~A} \textbf{2} (1952), no.~2, 71--80.
\newblock \url{https://www.i-repository.net/contents/osakacu/sugaku/DB-M-002-02-01.pdf}.

\bibitem[Ulmer and Urz\'ua(2021)]{manUU21}
Douglas Ulmer and Giancarlo Urz\'ua.
\newblock Bounding tangencies of sections on elliptic surfaces.
\newblock \emph{International Mathematics Research Notices}
\textbf{2021} (2021), no.~6, 4768--4802.
\newblock \doi{10.1093/imrn/rnaa222}.
\newblock Page references use arXiv:2002.01906v2, 26 May 2020.
\newblock \url{https://arxiv.org/abs/2002.01906v2}.

\bibitem[Viaclovsky(2026{\natexlab{a}})]{manViaclovskyLogs26}
Jeff Viaclovsky.
\newblock Logarithmic transformations on the six-sphere.
\newblock arXiv:2610.39102v1, 29 October 2026.
\newblock \url{https://arxiv.org/abs/2610.39102v1}.

\bibitem[Viaclovsky(2026{\natexlab{b}})]{manViaclovskyTwo26}
Jeff Viaclovsky.
\newblock A two-parameter family of complex structures on $S^6$.
\newblock arXiv:2609.33785v1, 27 September 2026.
\newblock \url{https://arxiv.org/abs/2609.33785v1}.

\bibitem[Waldhausen(1968)]{mapsWaldhausen68}
Friedhelm Waldhausen.
\newblock On irreducible 3-manifolds which are sufficiently large.
\newblock \emph{Annals of Mathematics} \textbf{87} (1968), no.~1, 56--88.
\newblock \doi{10.2307/1970594}.

\bibitem[Wang(2026)]{manWang26}
Zichang Wang.
\newblock A complex structure on a rational homology six-sphere with two-torsion.
\newblock arXiv:2609.11541v2, 2026.
\newblock \url{https://arxiv.org/abs/2609.11541v2}.

\bibitem[Whitehead(1950)]{mapsWhitehead50}
J.~H.~C. Whitehead.
\newblock The secondary boundary operator.
\newblock \emph{Proceedings of the National Academy of Sciences of the
United States of America} \textbf{36} (1950), no.~1, 55--60.
\newblock \doi{10.1073/pnas.36.1.55}.

\bibitem[Yang(2026)]{manYang26}
Xiufan Yang.
\newblock Compact complex threefolds fibred by two-dimensional complex tori.
\newblock arXiv:2609.29666v1, 13 September 2026.
\newblock \url{https://arxiv.org/abs/2609.29666v1}.

\bibitem[Zhubr(1975)]{manZhubr75}
A.~V. Zhubr.
\newblock Classification of simply-connected six-dimensional spinor manifolds.
\newblock \emph{Mathematics of the USSR-Izvestiya} \textbf{9} (1975), no.~4, 793--812.
\newblock \doi{10.1070/IM1975v009n04ABEH001498}.
\newblock English edition.

\bibitem[Zhubr(2010)]{manZhubr2010Gamma}
A.~V. Zhubr.
\newblock An exotic invariant for 6-manifolds: the direct construction.
\newblock \emph{St. Petersburg Mathematical Journal} \textbf{21} (2010), no.~3, 469--482.
\newblock \doi{10.1090/S1061-0022-10-01104-0}.

\bibitem[Zhubr(2013)]{ffZhubr13}
A.~V. Zhubr.
\newblock On the definition of ``exotic'' invariant for simply connected closed
6-dimensional spin manifolds.
\newblock \emph{Izvestia Komi Science Centre of the Ural Division of the
Russian Academy of Sciences} (2013), no.~4(16), 4--10.
\newblock \url{https://izvestia.komisc.ru/Archive/i16.pdf}.

\end{thebibliography}

\end{document}
